% Compile from the repository root with pdfLaTeX.
\documentclass[11pt,letterpaper]{article}
\usepackage[T1]{fontenc}
\usepackage{lmodern,amsmath,amssymb,mathtools,geometry,fancyhdr,xcolor}
\geometry{left=72bp,right=72bp,top=72bp,bottom=66bp,headheight=15pt,headsep=15pt,footskip=30pt}
\pdfminorversion=5
\pdfinclusioncopyfonts=1
\pdfdecimaldigits=4
\nofiles
\definecolor{ink}{RGB}{31,65,77}
\setlength{\parindent}{0pt}
\setlength{\parskip}{4.5pt plus .6pt minus .4pt}
\setlength{\abovedisplayskip}{8pt plus 1pt minus 2pt}
\setlength{\belowdisplayskip}{8pt plus 1pt minus 2pt}
\setlength{\emergencystretch}{1em}
\pagestyle{fancy}
\fancyhf{}
\fancyhead[L]{\small\scshape The Music of the Primes}
\fancyhead[R]{\small\scshape\rightmarktext}
\fancyfoot[L]{\footnotesize\color{ink}Revised mathematical text}
\fancyfoot[R]{\small\bfseries\printfolio}
\renewcommand{\headrulewidth}{.3pt}
\newcommand{\rightmarktext}{}
\newcommand{\printfolio}{}
\newcommand{\leaf}[4]{\clearpage\gdef\printfolio{#1}\gdef\rightmarktext{#2}%
  {\small\sffamily\bfseries\color{ink}#3\par}\vspace{3pt}%
  {\LARGE\bfseries #4\par}\vspace{5pt}}
\newcommand{\point}[1]{\par\smallskip\textbf{#1.}\enspace\ignorespaces}
\newcommand{\R}{\mathbb R}
\newcommand{\C}{\mathbb C}
\newcommand{\Q}{\mathbb Q}
\newcommand{\Z}{\mathbb Z}
\newcommand{\A}{\mathbb A}
\newcommand{\dd}{\,\mathrm d}
\newcommand{\citetitle}[2]{\leavevmode\pdfstartlink attr{/Border [0 0 0]} user{%
 /Subtype /Link /A << /S /URI /URI (\pdfescapestring{\detokenize{#1}}) >>}{\color{ink}#2}\pdfendlink}
\DeclareMathOperator{\Tr}{Tr}
\DeclareMathOperator{\supp}{supp}
\DeclareMathOperator{\Span}{span}
\newcommand{\abs}[1]{\lvert#1\rvert}

% Map the original destinations and the lettered leaves to physical PDF pages.
\makeatletter
\newcommand{\afterfolio}[2]{\expandafter\def\csname leaves@#1\endcsname{#2}}
\afterfolio{426}{a}
\afterfolio{446}{a,b,c,d}
\afterfolio{463}{a,b}
\afterfolio{476}{a}
\afterfolio{480}{a}
\afterfolio{484}{a,b}
\afterfolio{488}{a,b,c,d,e,f}
\afterfolio{489}{a,b}
\afterfolio{495}{a,b}
\afterfolio{508}{a}
\afterfolio{514}{a,b}
\afterfolio{518}{a}
\afterfolio{520}{a,b,c,d,e,f,g,h,i,j,k,l,m,n}
\afterfolio{522}{a,b,c,d,e,f,g,h,i,j,k,l,m,n,o,p,q,r,s,t,u,v,w,x,y,z,aa,ab,ac,ad,ae,af,ag,ah,ai,aj}
\afterfolio{558}{a}
\newcount\originalfolio
\newcount\physicalfolio
\newcommand{\baselinePDF}{output/pdf/the-music-of-the-primes-reproduction.pdf}
\pdfximage page 1 mediabox {\baselinePDF}
\edef\baselinecount{\the\pdflastximagepages}
\def\pdfpagelabels{0 << /S /D >> }
\originalfolio=0 \physicalfolio=0
\loop\ifnum\originalfolio<\baselinecount\relax
  \advance\originalfolio by 1 \advance\physicalfolio by 1
  \expandafter\xdef\csname physical@\the\originalfolio\endcsname{\the\physicalfolio}%
  \ifcsname leaves@\the\originalfolio\endcsname
    \edef\letterlist{\csname leaves@\the\originalfolio\endcsname}%
    \@for\letterfolio:=\letterlist\do{%
      \xdef\pdfpagelabels{\pdfpagelabels\the\physicalfolio\space << /P (\the\originalfolio\letterfolio) >> }%
      \advance\physicalfolio by 1
      \expandafter\xdef\csname physical@\the\originalfolio\letterfolio\endcsname{\the\physicalfolio}}%
    \xdef\pdfpagelabels{\pdfpagelabels\the\physicalfolio\space << /S /D /St \the\numexpr\originalfolio+1\relax\space >> }%
  \fi
\repeat
\newcommand{\MusicPDFPageRef}[1]{\pdfpageref\csname physical@#1\endcsname\space\space 0 R}
% Register the object references used by the manuscript's PDF navigation.
\newcount\MusicPDFObjectIndex
\newcommand*{\MusicPDFReserveObjects}[1]{%
  \MusicPDFObjectIndex=0
  \loop\ifnum\MusicPDFObjectIndex<#1\relax
    \advance\MusicPDFObjectIndex by 1
    \pdfobj reserveobjnum
    \expandafter\xdef
      \csname MusicPDFObject\the\MusicPDFObjectIndex\endcsname{\the\pdflastobj}%
  \repeat
}
\newcommand*{\MusicPDFRef}[1]{\csname MusicPDFObject#1\endcsname\space 0 R}
\providecommand*{\MusicPDFPageRef}[1]{\pdfpageref#1 \space 0 R}
\newcommand*{\MusicPDFWriteObject}[2]{%
  \immediate\pdfobj useobjnum \csname MusicPDFObject#1\endcsname {#2}%
}
\newcommand*{\MusicPDFSetAnnotations}[2]{%
  \expandafter\gdef\csname MusicPDFAnnotations#1\endcsname{/Annots #2}%
}
\endinput

\newcommand{\pagelink}[2]{\leavevmode\pdfstartlink attr{/Border [0 0 0]} user{%
  /Subtype /Link /A << /S /GoTo /D [\MusicPDFPageRef{#1} /Fit] >>}%
  {\color{ink}#2}\pdfendlink}
\newcommand{\pg}[1]{\pagelink{#1}{p.\,#1}}
\newcommand{\sourcepages}[3]{\clearpage
  \begingroup
  \pdfhorigin=0pt \pdfvorigin=0pt
  \originalfolio=#1\relax
  \loop\ifnum\originalfolio<\numexpr#2+1\relax
    \pdfximage page \originalfolio\space mediabox {\baselinePDF}%
    \edef\pageimage{\the\pdflastximage}%
    \edef\pageattributes{\ifcsname MusicPDFAnnotations\the\originalfolio\endcsname
      \csname MusicPDFAnnotations\the\originalfolio\endcsname\fi}%
    \expandafter\pdfpageattr\expandafter{\pageattributes}%
    \shipout\hbox{\rlap{\pdfrefximage\pageimage}%
      \vbox to 792bp{\kern 15bp\hbox to 612bp{\hfil
        {\footnotesize\sffamily\color{ink}#3}\hfil}\vss}}%
    \advance\originalfolio by 1 \stepcounter{page}%
  \repeat
  \endgroup\pdfpageattr{}}
\makeatother
\begin{document}
\let\editionpdfinfo\pdfinfo
\def\pdfinfo#1{}
% Navigation data extracted from the-music-of-the-primes.pdf.
% Read by the-music-of-the-primes-reproduction.tex.
\newcommand*{\MusicPDFPageCount}{940}
\MusicPDFReserveObjects{17886}
\MusicPDFWriteObject{1}{<< /Type /Outlines /First \MusicPDFRef{8817} /Last \MusicPDFRef{8818} /Count 21 >>}
\MusicPDFWriteObject{2}{<< /Dests \MusicPDFRef{8819} >>}
\MusicPDFWriteObject{3}{<< /S /GoTo /D [\MusicPDFPageRef{1} /Fit] >>}
\MusicPDFWriteObject{4}{<< /Type /Annot /Subtype /Link /Border [0 0 0] /H /I /C [1 0 0] /Rect [71.004 571.612 148.306 582.403] /A << /S /GoTo /D <636861707465722A2E32> >> >>}
\MusicPDFWriteObject{5}{<< /Type /Annot /Subtype /Link /Border [0 0 0] /H /I /C [1 0 0] /Rect [71.004 549.38 183.997 560.171] /A << /S /GoTo /D <636861707465722E31> >> >>}
\MusicPDFWriteObject{6}{<< /Type /Annot /Subtype /Link /Border [0 0 0] /H /I /C [1 0 0] /Rect [82.959 536.23 219.497 547.02] /A << /S /GoTo /D <73656374696F6E2E312E31> >> >>}
\MusicPDFWriteObject{7}{<< /Type /Annot /Subtype /Link /Border [0 0 0] /H /I /C [1 0 0] /Rect [82.959 525.016 286.535 533.87] /A << /S /GoTo /D <73656374696F6E2E312E32> >> >>}
\MusicPDFWriteObject{8}{<< /Type /Annot /Subtype /Link /Border [0 0 0] /H /I /C [1 0 0] /Rect [82.959 509.928 502.671 520.719] /A << /S /GoTo /D <73656374696F6E2E312E33> >> >>}
\MusicPDFWriteObject{9}{<< /Type /Annot /Subtype /Link /Border [0 0 0] /H /I /C [1 0 0] /Rect [82.959 496.777 382.042 507.568] /A << /S /GoTo /D <73656374696F6E2E312E34> >> >>}
\MusicPDFWriteObject{10}{<< /Type /Annot /Subtype /Link /Border [0 0 0] /H /I /C [1 0 0] /Rect [82.959 483.627 292.788 494.417] /A << /S /GoTo /D <73656374696F6E2E312E35> >> >>}
\MusicPDFWriteObject{11}{<< /Type /Annot /Subtype /Link /Border [0 0 0] /H /I /C [1 0 0] /Rect [82.959 470.476 283.464 481.267] /A << /S /GoTo /D <73656374696F6E2E312E36> >> >>}
\MusicPDFWriteObject{12}{<< /Type /Annot /Subtype /Link /Border [0 0 0] /H /I /C [1 0 0] /Rect [82.959 457.325 279.092 468.116] /A << /S /GoTo /D <73656374696F6E2E312E37> >> >>}
\MusicPDFWriteObject{13}{<< /Type /Annot /Subtype /Link /Border [0 0 0] /H /I /C [1 0 0] /Rect [82.959 444.174 333.735 454.965] /A << /S /GoTo /D <73656374696F6E2E312E38> >> >>}
\MusicPDFWriteObject{14}{<< /Type /Annot /Subtype /Link /Border [0 0 0] /H /I /C [1 0 0] /Rect [71.004 421.942 207.792 432.79] /A << /S /GoTo /D <636861707465722E32> >> >>}
\MusicPDFWriteObject{15}{<< /Type /Annot /Subtype /Link /Border [0 0 0] /H /I /C [1 0 0] /Rect [82.959 410.728 226.58 419.582] /A << /S /GoTo /D <73656374696F6E2E322E31> >> >>}
\MusicPDFWriteObject{16}{<< /Type /Annot /Subtype /Link /Border [0 0 0] /H /I /C [1 0 0] /Rect [82.959 397.577 218.682 406.488] /A << /S /GoTo /D <73656374696F6E2E322E32> >> >>}
\MusicPDFWriteObject{17}{<< /Type /Annot /Subtype /Link /Border [0 0 0] /H /I /C [1 0 0] /Rect [82.959 384.427 220.04 393.338] /A << /S /GoTo /D <73656374696F6E2E322E33> >> >>}
\MusicPDFWriteObject{18}{<< /Type /Annot /Subtype /Link /Border [0 0 0] /H /I /C [1 0 0] /Rect [82.959 368.425 379.15 381.201] /A << /S /GoTo /D <73656374696F6E2E322E34> >> >>}
\MusicPDFWriteObject{19}{<< /Type /Annot /Subtype /Link /Border [0 0 0] /H /I /C [1 0 0] /Rect [71.004 347.107 221.237 357.898] /A << /S /GoTo /D <636861707465722E33> >> >>}
\MusicPDFWriteObject{20}{<< /Type /Annot /Subtype /Link /Border [0 0 0] /H /I /C [1 0 0] /Rect [82.959 333.956 321.174 344.747] /A << /S /GoTo /D <73656374696F6E2E332E31> >> >>}
\MusicPDFWriteObject{21}{<< /Type /Annot /Subtype /Link /Border [0 0 0] /H /I /C [1 0 0] /Rect [82.959 322.742 234.341 331.596] /A << /S /GoTo /D <73656374696F6E2E332E32> >> >>}
\MusicPDFWriteObject{22}{<< /Type /Annot /Subtype /Link /Border [0 0 0] /H /I /C [1 0 0] /Rect [82.959 309.591 239.156 318.445] /A << /S /GoTo /D <73656374696F6E2E332E33> >> >>}
\MusicPDFWriteObject{23}{<< /Type /Annot /Subtype /Link /Border [0 0 0] /H /I /C [1 0 0] /Rect [82.959 294.504 276.187 305.295] /A << /S /GoTo /D <73656374696F6E2E332E34> >> >>}
\MusicPDFWriteObject{24}{<< /Type /Annot /Subtype /Link /Border [0 0 0] /H /I /C [1 0 0] /Rect [71.004 272.272 331.132 283.063] /A << /S /GoTo /D <636861707465722E34> >> >>}
\MusicPDFWriteObject{25}{<< /Type /Annot /Subtype /Link /Border [0 0 0] /H /I /C [1 0 0] /Rect [82.959 261.058 176.115 269.912] /A << /S /GoTo /D <73656374696F6E2E342E31> >> >>}
\MusicPDFWriteObject{26}{<< /Type /Annot /Subtype /Link /Border [0 0 0] /H /I /C [1 0 0] /Rect [82.959 245.885 453.767 256.761] /A << /S /GoTo /D <73656374696F6E2E342E32> >> >>}
\MusicPDFWriteObject{27}{<< /Type /Annot /Subtype /Link /Border [0 0 0] /H /I /C [1 0 0] /Rect [82.959 232.82 263.267 243.61] /A << /S /GoTo /D <73656374696F6E2E342E33> >> >>}
\MusicPDFWriteObject{28}{<< /Type /Annot /Subtype /Link /Border [0 0 0] /H /I /C [1 0 0] /Rect [82.959 219.583 417.821 230.46] /A << /S /GoTo /D <73656374696F6E2E342E34> >> >>}
\MusicPDFWriteObject{29}{<< /Type /Annot /Subtype /Link /Border [0 0 0] /H /I /C [1 0 0] /Rect [82.959 206.518 346.628 217.309] /A << /S /GoTo /D <73656374696F6E2E342E35> >> >>}
\MusicPDFWriteObject{30}{<< /Type /Annot /Subtype /Link /Border [0 0 0] /H /I /C [1 0 0] /Rect [82.959 193.367 349.976 204.158] /A << /S /GoTo /D <73656374696F6E2E342E36> >> >>}
\MusicPDFWriteObject{31}{<< /Type /Annot /Subtype /Link /Border [0 0 0] /H /I /C [1 0 0] /Rect [71.004 171.135 349.049 181.983] /A << /S /GoTo /D <636861707465722E35> >> >>}
\MusicPDFWriteObject{32}{<< /Type /Annot /Subtype /Link /Border [0 0 0] /H /I /C [1 0 0] /Rect [82.959 157.984 429.398 168.832] /A << /S /GoTo /D <73656374696F6E2E352E31> >> >>}
\MusicPDFWriteObject{33}{<< /Type /Annot /Subtype /Link /Border [0 0 0] /H /I /C [1 0 0] /Rect [82.959 144.834 393.94 155.681] /A << /S /GoTo /D <73656374696F6E2E352E32> >> >>}
\MusicPDFWriteObject{34}{<< /Type /Annot /Subtype /Link /Border [0 0 0] /H /I /C [1 0 0] /Rect [82.959 131.683 253.777 142.474] /A << /S /GoTo /D <73656374696F6E2E352E33> >> >>}
\MusicPDFWriteObject{35}{<< /Type /Annot /Subtype /Link /Border [0 0 0] /H /I /C [1 0 0] /Rect [82.959 118.532 306.471 131.343] /A << /S /GoTo /D <73656374696F6E2E352E34> >> >>}
\MusicPDFWriteObject{36}{<< /Type /Annot /Subtype /Link /Border [0 0 0] /H /I /C [1 0 0] /Rect [82.959 105.381 369.592 116.172] /A << /S /GoTo /D <73656374696F6E2E352E35> >> >>}
\MusicPDFWriteObject{37}{<< /Type /Annot /Subtype /Link /Border [0 0 0] /H /I /C [1 0 0] /Rect [71.004 83.149 231.931 93.94] /A << /S /GoTo /D <636861707465722E36> >> >>}
\MusicPDFWriteObject{38}{<< /Type /Annot /Subtype /Link /Border [0 0 0] /H /I /C [1 0 0] /Rect [82.959 685.187 294.475 695.977] /A << /S /GoTo /D <73656374696F6E2E362E31> >> >>}
\MusicPDFWriteObject{39}{<< /Type /Annot /Subtype /Link /Border [0 0 0] /H /I /C [1 0 0] /Rect [82.959 672.036 310.467 682.827] /A << /S /GoTo /D <73656374696F6E2E362E32> >> >>}
\MusicPDFWriteObject{40}{<< /Type /Annot /Subtype /Link /Border [0 0 0] /H /I /C [1 0 0] /Rect [82.959 660.822 274.016 669.676] /A << /S /GoTo /D <73656374696F6E2E362E33> >> >>}
\MusicPDFWriteObject{41}{<< /Type /Annot /Subtype /Link /Border [0 0 0] /H /I /C [1 0 0] /Rect [71.004 636.233 300.504 647.024] /A << /S /GoTo /D <636861707465722E37> >> >>}
\MusicPDFWriteObject{42}{<< /Type /Annot /Subtype /Link /Border [0 0 0] /H /I /C [1 0 0] /Rect [82.959 623.083 194.182 633.873] /A << /S /GoTo /D <73656374696F6E2E372E31> >> >>}
\MusicPDFWriteObject{43}{<< /Type /Annot /Subtype /Link /Border [0 0 0] /H /I /C [1 0 0] /Rect [82.959 611.869 263.144 620.723] /A << /S /GoTo /D <73656374696F6E2E372E32> >> >>}
\MusicPDFWriteObject{44}{<< /Type /Annot /Subtype /Link /Border [0 0 0] /H /I /C [1 0 0] /Rect [82.959 596.781 369.564 607.572] /A << /S /GoTo /D <73656374696F6E2E372E33> >> >>}
\MusicPDFWriteObject{45}{<< /Type /Annot /Subtype /Link /Border [0 0 0] /H /I /C [1 0 0] /Rect [82.959 583.63 155.531 594.421] /A << /S /GoTo /D <73656374696F6E2E372E34> >> >>}
\MusicPDFWriteObject{46}{<< /Type /Annot /Subtype /Link /Border [0 0 0] /H /I /C [1 0 0] /Rect [82.959 570.48 199.134 581.27] /A << /S /GoTo /D <73656374696F6E2E372E35> >> >>}
\MusicPDFWriteObject{47}{<< /Type /Annot /Subtype /Link /Border [0 0 0] /H /I /C [1 0 0] /Rect [71.004 547.828 273.916 558.619] /A << /S /GoTo /D <636861707465722E38> >> >>}
\MusicPDFWriteObject{48}{<< /Type /Annot /Subtype /Link /Border [0 0 0] /H /I /C [1 0 0] /Rect [82.959 534.677 214.545 545.468] /A << /S /GoTo /D <73656374696F6E2E382E31> >> >>}
\MusicPDFWriteObject{49}{<< /Type /Annot /Subtype /Link /Border [0 0 0] /H /I /C [1 0 0] /Rect [82.959 523.463 323.498 532.317] /A << /S /GoTo /D <73656374696F6E2E382E32> >> >>}
\MusicPDFWriteObject{50}{<< /Type /Annot /Subtype /Link /Border [0 0 0] /H /I /C [1 0 0] /Rect [82.959 508.376 218.418 519.166] /A << /S /GoTo /D <73656374696F6E2E382E33> >> >>}
\MusicPDFWriteObject{51}{<< /Type /Annot /Subtype /Link /Border [0 0 0] /H /I /C [1 0 0] /Rect [82.959 497.162 173.321 506.016] /A << /S /GoTo /D <73656374696F6E2E382E34> >> >>}
\MusicPDFWriteObject{52}{<< /Type /Annot /Subtype /Link /Border [0 0 0] /H /I /C [1 0 0] /Rect [82.959 482.074 177.941 492.865] /A << /S /GoTo /D <73656374696F6E2E382E35> >> >>}
\MusicPDFWriteObject{53}{<< /Type /Annot /Subtype /Link /Border [0 0 0] /H /I /C [1 0 0] /Rect [71.004 459.423 240.494 470.213] /A << /S /GoTo /D <636861707465722E39> >> >>}
\MusicPDFWriteObject{54}{<< /Type /Annot /Subtype /Link /Border [0 0 0] /H /I /C [1 0 0] /Rect [82.959 445.718 214.994 458.311] /A << /S /GoTo /D <73656374696F6E2E392E31> >> >>}
\MusicPDFWriteObject{55}{<< /Type /Annot /Subtype /Link /Border [0 0 0] /H /I /C [1 0 0] /Rect [82.959 432.314 367.993 444.522] /A << /S /GoTo /D <73656374696F6E2E392E32> >> >>}
\MusicPDFWriteObject{56}{<< /Type /Annot /Subtype /Link /Border [0 0 0] /H /I /C [1 0 0] /Rect [82.959 419.416 340.105 431.372] /A << /S /GoTo /D <73656374696F6E2E392E33> >> >>}
\MusicPDFWriteObject{57}{<< /Type /Annot /Subtype /Link /Border [0 0 0] /H /I /C [1 0 0] /Rect [71.004 397.319 259.584 408.109] /A << /S /GoTo /D <636861707465722E3130> >> >>}
\MusicPDFWriteObject{58}{<< /Type /Annot /Subtype /Link /Border [0 0 0] /H /I /C [1 0 0] /Rect [82.959 381.821 447.467 396.676] /A << /S /GoTo /D <73656374696F6E2E31302E31> >> >>}
\MusicPDFWriteObject{59}{<< /Type /Annot /Subtype /Link /Border [0 0 0] /H /I /C [1 0 0] /Rect [82.959 371.017 242.461 381.808] /A << /S /GoTo /D <73656374696F6E2E31302E32> >> >>}
\MusicPDFWriteObject{60}{<< /Type /Annot /Subtype /Link /Border [0 0 0] /H /I /C [1 0 0] /Rect [82.959 357.866 199.549 368.657] /A << /S /GoTo /D <73656374696F6E2E31302E33> >> >>}
\MusicPDFWriteObject{61}{<< /Type /Annot /Subtype /Link /Border [0 0 0] /H /I /C [1 0 0] /Rect [82.959 344.716 240.524 355.506] /A << /S /GoTo /D <73656374696F6E2E31302E34> >> >>}
\MusicPDFWriteObject{62}{<< /Type /Annot /Subtype /Link /Border [0 0 0] /H /I /C [1 0 0] /Rect [71.004 322.064 220.269 332.855] /A << /S /GoTo /D <636861707465722E3131> >> >>}
\MusicPDFWriteObject{63}{<< /Type /Annot /Subtype /Link /Border [0 0 0] /H /I /C [1 0 0] /Rect [82.959 310.85 201.209 319.704] /A << /S /GoTo /D <73656374696F6E2E31312E31> >> >>}
\MusicPDFWriteObject{64}{<< /Type /Annot /Subtype /Link /Border [0 0 0] /H /I /C [1 0 0] /Rect [82.959 295.762 458.266 306.553] /A << /S /GoTo /D <73656374696F6E2E31312E32> >> >>}
\MusicPDFWriteObject{65}{<< /Type /Annot /Subtype /Link /Border [0 0 0] /H /I /C [1 0 0] /Rect [82.959 284.548 268.607 293.402] /A << /S /GoTo /D <73656374696F6E2E31312E33> >> >>}
\MusicPDFWriteObject{66}{<< /Type /Annot /Subtype /Link /Border [0 0 0] /H /I /C [1 0 0] /Rect [71.004 259.96 268.784 270.751] /A << /S /GoTo /D <636861707465722E3132> >> >>}
\MusicPDFWriteObject{67}{<< /Type /Annot /Subtype /Link /Border [0 0 0] /H /I /C [1 0 0] /Rect [82.959 246.809 259.587 257.6] /A << /S /GoTo /D <73656374696F6E2E31322E31> >> >>}
\MusicPDFWriteObject{68}{<< /Type /Annot /Subtype /Link /Border [0 0 0] /H /I /C [1 0 0] /Rect [82.959 233.659 304.989 244.449] /A << /S /GoTo /D <73656374696F6E2E31322E32> >> >>}
\MusicPDFWriteObject{69}{<< /Type /Annot /Subtype /Link /Border [0 0 0] /H /I /C [1 0 0] /Rect [82.959 220.508 483.209 231.299] /A << /S /GoTo /D <73656374696F6E2E31322E33> >> >>}
\MusicPDFWriteObject{70}{<< /Type /Annot /Subtype /Link /Border [0 0 0] /H /I /C [1 0 0] /Rect [82.959 209.294 220.868 218.148] /A << /S /GoTo /D <73656374696F6E2E31322E34> >> >>}
\MusicPDFWriteObject{71}{<< /Type /Annot /Subtype /Link /Border [0 0 0] /H /I /C [1 0 0] /Rect [82.959 196.143 320.012 204.997] /A << /S /GoTo /D <73656374696F6E2E31322E35> >> >>}
\MusicPDFWriteObject{72}{<< /Type /Annot /Subtype /Link /Border [0 0 0] /H /I /C [1 0 0] /Rect [71.004 171.555 250.926 182.345] /A << /S /GoTo /D <636861707465722E3133> >> >>}
\MusicPDFWriteObject{73}{<< /Type /Annot /Subtype /Link /Border [0 0 0] /H /I /C [1 0 0] /Rect [82.959 158.404 283.602 169.195] /A << /S /GoTo /D <73656374696F6E2E31332E31> >> >>}
\MusicPDFWriteObject{74}{<< /Type /Annot /Subtype /Link /Border [0 0 0] /H /I /C [1 0 0] /Rect [82.959 145.253 281.029 156.044] /A << /S /GoTo /D <73656374696F6E2E31332E32> >> >>}
\MusicPDFWriteObject{75}{<< /Type /Annot /Subtype /Link /Border [0 0 0] /H /I /C [1 0 0] /Rect [82.959 134.039 306.179 142.893] /A << /S /GoTo /D <73656374696F6E2E31332E33> >> >>}
\MusicPDFWriteObject{76}{<< /Type /Annot /Subtype /Link /Border [0 0 0] /H /I /C [1 0 0] /Rect [71.004 109.451 193.432 120.241] /A << /S /GoTo /D <636861707465722E3134> >> >>}
\MusicPDFWriteObject{77}{<< /Type /Annot /Subtype /Link /Border [0 0 0] /H /I /C [1 0 0] /Rect [82.959 96.3 258.84 107.091] /A << /S /GoTo /D <73656374696F6E2E31342E31> >> >>}
\MusicPDFWriteObject{78}{<< /Type /Annot /Subtype /Link /Border [0 0 0] /H /I /C [1 0 0] /Rect [82.959 83.149 368.264 93.94] /A << /S /GoTo /D <73656374696F6E2E31342E32> >> >>}
\MusicPDFWriteObject{79}{<< /Type /Annot /Subtype /Link /Border [0 0 0] /H /I /C [1 0 0] /Rect [82.959 685.187 207.075 695.977] /A << /S /GoTo /D <73656374696F6E2E31342E33> >> >>}
\MusicPDFWriteObject{80}{<< /Type /Annot /Subtype /Link /Border [0 0 0] /H /I /C [1 0 0] /Rect [71.004 664.066 335.641 674.856] /A << /S /GoTo /D <617070656E6469782E41> >> >>}
\MusicPDFWriteObject{81}{<< /Type /Annot /Subtype /Link /Border [0 0 0] /H /I /C [1 0 0] /Rect [82.959 650.001 431.249 662.316] /A << /S /GoTo /D <73656374696F6E2E412E31> >> >>}
\MusicPDFWriteObject{82}{<< /Type /Annot /Subtype /Link /Border [0 0 0] /H /I /C [1 0 0] /Rect [82.959 639.701 280.863 648.555] /A << /S /GoTo /D <73656374696F6E2E412E32> >> >>}
\MusicPDFWriteObject{83}{<< /Type /Annot /Subtype /Link /Border [0 0 0] /H /I /C [1 0 0] /Rect [82.959 626.55 176.973 635.404] /A << /S /GoTo /D <73656374696F6E2E412E33> >> >>}
\MusicPDFWriteObject{84}{<< /Type /Annot /Subtype /Link /Border [0 0 0] /H /I /C [1 0 0] /Rect [82.959 611.463 389.141 622.253] /A << /S /GoTo /D <73656374696F6E2E412E34> >> >>}
\MusicPDFWriteObject{85}{<< /Type /Annot /Subtype /Link /Border [0 0 0] /H /I /C [1 0 0] /Rect [82.959 598.312 429.409 609.103] /A << /S /GoTo /D <73656374696F6E2E412E35> >> >>}
\MusicPDFWriteObject{86}{<< /Type /Annot /Subtype /Link /Border [0 0 0] /H /I /C [1 0 0] /Rect [71.004 577.191 331.049 587.982] /A << /S /GoTo /D <617070656E6469782E42> >> >>}
\MusicPDFWriteObject{87}{<< /Type /Annot /Subtype /Link /Border [0 0 0] /H /I /C [1 0 0] /Rect [82.959 563.127 239.084 575.442] /A << /S /GoTo /D <73656374696F6E2E422E31> >> >>}
\MusicPDFWriteObject{88}{<< /Type /Annot /Subtype /Link /Border [0 0 0] /H /I /C [1 0 0] /Rect [82.959 550.89 274.859 561.68] /A << /S /GoTo /D <73656374696F6E2E422E32> >> >>}
\MusicPDFWriteObject{89}{<< /Type /Annot /Subtype /Link /Border [0 0 0] /H /I /C [1 0 0] /Rect [82.959 537.739 242.24 548.529] /A << /S /GoTo /D <73656374696F6E2E422E33> >> >>}
\MusicPDFWriteObject{90}{<< /Type /Annot /Subtype /Link /Border [0 0 0] /H /I /C [1 0 0] /Rect [82.959 524.588 286.813 535.379] /A << /S /GoTo /D <73656374696F6E2E422E34> >> >>}
\MusicPDFWriteObject{91}{<< /Type /Annot /Subtype /Link /Border [0 0 0] /H /I /C [1 0 0] /Rect [82.959 511.437 252.477 522.228] /A << /S /GoTo /D <73656374696F6E2E422E35> >> >>}
\MusicPDFWriteObject{92}{<< /Type /Annot /Subtype /Link /Border [0 0 0] /H /I /C [1 0 0] /Rect [82.959 498.287 321.287 509.077] /A << /S /GoTo /D <73656374696F6E2E422E36> >> >>}
\MusicPDFWriteObject{93}{<< /Type /Annot /Subtype /Link /Border [0 0 0] /H /I /C [1 0 0] /Rect [82.959 485.136 340.819 495.926] /A << /S /GoTo /D <73656374696F6E2E422E37> >> >>}
\MusicPDFWriteObject{94}{<< /Type /Annot /Subtype /Link /Border [0 0 0] /H /I /C [1 0 0] /Rect [71.004 464.015 383.257 474.806] /A << /S /GoTo /D <617070656E6469782E43> >> >>}
\MusicPDFWriteObject{95}{<< /Type /Annot /Subtype /Link /Border [0 0 0] /H /I /C [1 0 0] /Rect [82.959 450.864 288.029 461.655] /A << /S /GoTo /D <73656374696F6E2E432E31> >> >>}
\MusicPDFWriteObject{96}{<< /Type /Annot /Subtype /Link /Border [0 0 0] /H /I /C [1 0 0] /Rect [82.959 437.713 449.965 448.504] /A << /S /GoTo /D <73656374696F6E2E432E32> >> >>}
\MusicPDFWriteObject{97}{<< /Type /Annot /Subtype /Link /Border [0 0 0] /H /I /C [1 0 0] /Rect [82.959 424.563 190.225 435.353] /A << /S /GoTo /D <73656374696F6E2E432E33> >> >>}
\MusicPDFWriteObject{98}{<< /Type /Annot /Subtype /Link /Border [0 0 0] /H /I /C [1 0 0] /Rect [82.959 411.412 255.16 422.203] /A << /S /GoTo /D <73656374696F6E2E432E34> >> >>}
\MusicPDFWriteObject{99}{<< /Type /Annot /Subtype /Link /Border [0 0 0] /H /I /C [1 0 0] /Rect [82.959 398.261 413.03 409.052] /A << /S /GoTo /D <73656374696F6E2E432E35> >> >>}
\MusicPDFWriteObject{100}{<< /Type /Annot /Subtype /Link /Border [0 0 0] /H /I /C [1 0 0] /Rect [71.004 377.14 255.157 387.931] /A << /S /GoTo /D <617070656E6469782E44> >> >>}
\MusicPDFWriteObject{101}{<< /Type /Annot /Subtype /Link /Border [0 0 0] /H /I /C [1 0 0] /Rect [82.959 363.989 412.2 374.78] /A << /S /GoTo /D <73656374696F6E2E442E31> >> >>}
\MusicPDFWriteObject{102}{<< /Type /Annot /Subtype /Link /Border [0 0 0] /H /I /C [1 0 0] /Rect [82.959 350.839 324.743 361.629] /A << /S /GoTo /D <73656374696F6E2E442E32> >> >>}
\MusicPDFWriteObject{103}{<< /Type /Annot /Subtype /Link /Border [0 0 0] /H /I /C [1 0 0] /Rect [82.959 337.688 341.442 348.479] /A << /S /GoTo /D <73656374696F6E2E442E33> >> >>}
\MusicPDFWriteObject{104}{<< /Type /Annot /Subtype /Link /Border [0 0 0] /H /I /C [1 0 0] /Rect [71.004 316.567 217.392 327.358] /A << /S /GoTo /D <617070656E6469782E45> >> >>}
\MusicPDFWriteObject{105}{<< /Type /Annot /Subtype /Link /Border [0 0 0] /H /I /C [1 0 0] /Rect [82.959 303.416 235.461 314.207] /A << /S /GoTo /D <73656374696F6E2E452E31> >> >>}
\MusicPDFWriteObject{106}{<< /Type /Annot /Subtype /Link /Border [0 0 0] /H /I /C [1 0 0] /Rect [82.959 290.266 313.843 301.056] /A << /S /GoTo /D <73656374696F6E2E452E32> >> >>}
\MusicPDFWriteObject{107}{<< /Type /Annot /Subtype /Link /Border [0 0 0] /H /I /C [1 0 0] /Rect [82.959 277.115 338.743 287.906] /A << /S /GoTo /D <73656374696F6E2E452E33> >> >>}
\MusicPDFWriteObject{108}{<< /Type /Annot /Subtype /Link /Border [0 0 0] /H /I /C [1 0 0] /Rect [82.959 263.964 207.739 274.755] /A << /S /GoTo /D <73656374696F6E2E452E34> >> >>}
\MusicPDFWriteObject{109}{<< /Type /Annot /Subtype /Link /Border [0 0 0] /H /I /C [1 0 0] /Rect [82.959 250.813 224.672 261.604] /A << /S /GoTo /D <73656374696F6E2E452E35> >> >>}
\MusicPDFWriteObject{110}{<< /Type /Annot /Subtype /Link /Border [0 0 0] /H /I /C [1 0 0] /Rect [82.959 236.749 207.697 248.453] /A << /S /GoTo /D <73656374696F6E2E452E36> >> >>}
\MusicPDFWriteObject{111}{<< /Type /Annot /Subtype /Link /Border [0 0 0] /H /I /C [1 0 0] /Rect [71.004 216.542 128.773 227.332] /A << /S /GoTo /D <617070656E6469782A2E33> >> >>}
\MusicPDFWriteObject{112}{<< /Type /Annot /Subtype /Link /Border [0 0 0] /H /I /C [0 1 0] /Rect [172.726 608.219 190.199 617.087] /A << /S /GoTo /D <636974652E4346> >> >>}
\MusicPDFWriteObject{113}{<< /Type /Annot /Subtype /Link /Border [0 0 0] /H /I /C [0 1 0] /Rect [150.895 555.616 168.367 564.484] /A << /S /GoTo /D <636974652E4346> >> >>}
\MusicPDFWriteObject{114}{<< /Type /Annot /Subtype /Link /Border [0 0 0] /H /I /C [1 0 0] /Rect [120.307 533.356 127.279 544.147] /A << /S /GoTo /D <706167652E36> >> >>}
\MusicPDFWriteObject{115}{<< /Type /Annot /Subtype /Link /Border [0 0 0] /H /I /C [1 0 0] /Rect [235.002 520.205 241.975 531.053] /A << /S /GoTo /D <706167652E36> >> >>}
\MusicPDFWriteObject{116}{<< /Type /Annot /Subtype /Link /Border [0 0 0] /H /I /C [1 0 0] /Rect [245.977 506.141 252.949 517.902] /A << /S /GoTo /D <706167652E36> >> >>}
\MusicPDFWriteObject{117}{<< /Type /Annot /Subtype /Link /Border [0 0 0] /H /I /C [1 0 0] /Rect [160.331 467.602 167.304 479.464] /A << /S /GoTo /D <706167652E36> >> >>}
\MusicPDFWriteObject{118}{<< /Type /Annot /Subtype /Link /Border [0 0 0] /H /I /C [1 0 0] /Rect [508.12 557.31 515.093 572.165] /A << /S /GoTo /D <706167652E36> >> >>}
\MusicPDFWriteObject{119}{<< /Type /Annot /Subtype /Link /Border [0 0 0] /H /I /C [1 0 0] /Rect [455.255 546.507 462.228 557.297] /A << /S /GoTo /D <706167652E36> >> >>}
\MusicPDFWriteObject{120}{<< /Type /Annot /Subtype /Link /Border [0 0 0] /H /I /C [1 0 0] /Rect [464.508 533.798 471.48 544.147] /A << /S /GoTo /D <706167652E36> >> >>}
\MusicPDFWriteObject{121}{<< /Type /Annot /Subtype /Link /Border [0 0 0] /H /I /C [1 0 0] /Rect [88.936 444.433 94.899 454.8] /A << /S /GoTo /D <48666F6F746E6F74652E31> >> >>}
\MusicPDFWriteObject{122}{<< /Type /Annot /Subtype /Link /Border [0 0 0] /H /I /C [1 0 0] /Rect [409.831 343.822 415.794 356.679] /A << /S /GoTo /D <48666F6F746E6F74652E32> >> >>}
\MusicPDFWriteObject{123}{<< /Type /Annot /Subtype /Link /Border [0 0 0] /H /I /C [1 0 0] /Rect [360.652 672.036 366.615 684.34] /A << /S /GoTo /D <48666F6F746E6F74652E33> >> >>}
\MusicPDFWriteObject{124}{<< /Type /Annot /Subtype /Link /Border [0 0 0] /H /I /C [1 0 0] /Rect [221.126 593.131 227.089 605.435] /A << /S /GoTo /D <48666F6F746E6F74652E34> >> >>}
\MusicPDFWriteObject{125}{<< /Type /Annot /Subtype /Link /Border [0 0 0] /H /I /C [1 0 0] /Rect [371.904 508.898 402.117 520.851] /A << /S /GoTo /D <6571756174696F6E2E312E332E312E31> >> >>}
\MusicPDFWriteObject{126}{<< /Type /Annot /Subtype /Link /Border [0 0 0] /H /I /C [1 0 0] /Rect [387.455 272.738 393.417 285.042] /A << /S /GoTo /D <48666F6F746E6F74652E35> >> >>}
\MusicPDFWriteObject{127}{<< /Type /Annot /Subtype /Link /Border [0 0 0] /H /I /C [1 0 0] /Rect [294.907 474.388 325.12 486.341] /A << /S /GoTo /D <6571756174696F6E2E312E382E332E35> >> >>}
\MusicPDFWriteObject{128}{<< /Type /Annot /Subtype /Link /Border [0 0 0] /H /I /C [1 0 0] /Rect [278.171 269.995 313.365 281.948] /A << /S /GoTo /D <6571756174696F6E2E31322E322E322E34> >> >>}
\MusicPDFWriteObject{129}{<< /Type /Annot /Subtype /Link /Border [0 0 0] /H /I /C [1 0 0] /Rect [161.65 122.048 196.844 134.001] /A << /S /GoTo /D <6571756174696F6E2E312E332E312E3132> >> >>}
\MusicPDFWriteObject{130}{<< /Type /Annot /Subtype /Link /Border [0 0 0] /H /I /C [1 0 0] /Rect [73.712 95.747 108.905 107.699] /A << /S /GoTo /D <6571756174696F6E2E312E332E312E3132> >> >>}
\MusicPDFWriteObject{131}{<< /Type /Annot /Subtype /Link /Border [0 0 0] /H /I /C [1 0 0] /Rect [170.978 338.448 206.171 350.4] /A << /S /GoTo /D <6571756174696F6E2E31312E332E312E31> >> >>}
\MusicPDFWriteObject{132}{<< /Type /Annot /Subtype /Link /Border [0 0 0] /H /I /C [1 0 0] /Rect [192.745 466.623 227.938 478.576] /A << /S /GoTo /D <6571756174696F6E2E312E332E312E3132> >> >>}
\MusicPDFWriteObject{133}{<< /Type /Annot /Subtype /Link /Border [0 0 0] /H /I /C [1 0 0] /Rect [506.695 401.62 541.888 413.573] /A << /S /GoTo /D <6571756174696F6E2E312E332E312E3132> >> >>}
\MusicPDFWriteObject{134}{<< /Type /Annot /Subtype /Link /Border [0 0 0] /H /I /C [1 0 0] /Rect [233.55 275.148 239.513 287.452] /A << /S /GoTo /D <48666F6F746E6F74652E36> >> >>}
\MusicPDFWriteObject{135}{<< /Type /Annot /Subtype /Link /Border [0 0 0] /H /I /C [1 0 0] /Rect [325.049 197.701 331.012 210.005] /A << /S /GoTo /D <48666F6F746E6F74652E37> >> >>}
\MusicPDFWriteObject{136}{<< /Type /Annot /Subtype /Link /Border [0 0 0] /H /I /C [0 1 0] /Rect [201.82 274.527 219.292 283.395] /A << /S /GoTo /D <636974652E4346> >> >>}
\MusicPDFWriteObject{137}{<< /Type /Annot /Subtype /Link /Border [0 0 0] /H /I /C [1 0 0] /Rect [203.991 592.578 234.204 604.531] /A << /S /GoTo /D <6571756174696F6E2E312E372E312E31> >> >>}
\MusicPDFWriteObject{138}{<< /Type /Annot /Subtype /Link /Border [0 0 0] /H /I /C [1 0 0] /Rect [369.677 309.311 399.89 321.263] /A << /S /GoTo /D <6571756174696F6E2E312E372E312E32> >> >>}
\MusicPDFWriteObject{139}{<< /Type /Annot /Subtype /Link /Border [0 0 0] /H /I /C [1 0 0] /Rect [386.387 296.16 416.6 308.113] /A << /S /GoTo /D <6571756174696F6E2E312E372E312E32> >> >>}
\MusicPDFWriteObject{140}{<< /Type /Annot /Subtype /Link /Border [0 0 0] /H /I /C [1 0 0] /Rect [191.438 256.708 221.651 268.66] /A << /S /GoTo /D <6571756174696F6E2E312E372E312E32> >> >>}
\MusicPDFWriteObject{141}{<< /Type /Annot /Subtype /Link /Border [0 0 0] /H /I /C [0 1 0] /Rect [174.843 206.595 194.072 215.462] /A << /S /GoTo /D <636974652E4D53> >> >>}
\MusicPDFWriteObject{142}{<< /Type /Annot /Subtype /Link /Border [0 0 0] /H /I /C [0 1 0] /Rect [225.185 193.444 242.547 202.312] /A << /S /GoTo /D <636974652E4753> >> >>}
\MusicPDFWriteObject{143}{<< /Type /Annot /Subtype /Link /Border [0 0 0] /H /I /C [1 0 0] /Rect [211.451 445.96 242.229 457.913] /A << /S /GoTo /D <74686D2E312E372E312E31> >> >>}
\MusicPDFWriteObject{144}{<< /Type /Annot /Subtype /Link /Border [0 0 0] /H /I /C [1 0 0] /Rect [171.281 407.061 201.494 417.852] /A << /S /GoTo /D <74686D2E312E372E312E31> >> >>}
\MusicPDFWriteObject{145}{<< /Type /Annot /Subtype /Link /Border [0 0 0] /H /I /C [1 0 0] /Rect [137.079 279.857 167.292 291.81] /A << /S /GoTo /D <6571756174696F6E2E312E372E322E31> >> >>}
\MusicPDFWriteObject{146}{<< /Type /Annot /Subtype /Link /Border [0 0 0] /H /I /C [1 0 0] /Rect [296.291 253.556 326.505 265.508] /A << /S /GoTo /D <6571756174696F6E2E312E372E322E31> >> >>}
\MusicPDFWriteObject{147}{<< /Type /Annot /Subtype /Link /Border [0 0 0] /H /I /C [1 0 0] /Rect [419.015 626.653 449.229 638.605] /A << /S /GoTo /D <6571756174696F6E2E312E382E322E32> >> >>}
\MusicPDFWriteObject{148}{<< /Type /Annot /Subtype /Link /Border [0 0 0] /H /I /C [1 0 0] /Rect [401.296 489.915 431.509 501.868] /A << /S /GoTo /D <6571756174696F6E2E312E382E322E32> >> >>}
\MusicPDFWriteObject{149}{<< /Type /Annot /Subtype /Link /Border [0 0 0] /H /I /C [1 0 0] /Rect [308.256 477.318 338.469 488.109] /A << /S /GoTo /D <74686D2E312E372E312E31> >> >>}
\MusicPDFWriteObject{150}{<< /Type /Annot /Subtype /Link /Border [0 0 0] /H /I /C [1 0 0] /Rect [195.522 299.104 225.171 309.895] /A << /S /GoTo /D <6571756174696F6E2E312E382E322E32> >> >>}
\MusicPDFWriteObject{151}{<< /Type /Annot /Subtype /Link /Border [0 0 0] /H /I /C [1 0 0] /Rect [434.326 444.063 456.382 456.016] /A << /S /GoTo /D <6578706563746174696F6E2E312E372E31> >> >>}
\MusicPDFWriteObject{152}{<< /Type /Annot /Subtype /Link /Border [0 0 0] /H /I /C [1 0 0] /Rect [323.671 418.315 346.547 429.105] /A << /S /GoTo /D <6578706563746174696F6E2E312E372E31> >> >>}
\MusicPDFWriteObject{153}{<< /Type /Annot /Subtype /Link /Border [0 0 0] /H /I /C [1 0 0] /Rect [303.984 391.46 326.86 403.413] /A << /S /GoTo /D <6578706563746174696F6E2E312E372E31> >> >>}
\MusicPDFWriteObject{154}{<< /Type /Annot /Subtype /Link /Border [0 0 0] /H /I /C [1 0 0] /Rect [494.144 365.158 516.201 377.131] /A << /S /GoTo /D <6578706563746174696F6E2E312E372E31> >> >>}
\MusicPDFWriteObject{155}{<< /Type /Annot /Subtype /Link /Border [0 0 0] /H /I /C [1 0 0] /Rect [262.301 312.195 284.767 324.51] /A << /S /GoTo /D <6578706563746174696F6E2E312E372E31> >> >>}
\MusicPDFWriteObject{156}{<< /Type /Annot /Subtype /Link /Border [0 0 0] /H /I /C [0 1 0] /Rect [343.104 568.766 366.947 577.634] /A << /S /GoTo /D <636974652E4D6531> >> >>}
\MusicPDFWriteObject{157}{<< /Type /Annot /Subtype /Link /Border [0 0 0] /H /I /C [1 0 0] /Rect [395.862 579.427 426.075 591.38] /A << /S /GoTo /D <6571756174696F6E2E312E382E312E31> >> >>}
\MusicPDFWriteObject{158}{<< /Type /Annot /Subtype /Link /Border [0 0 0] /H /I /C [1 0 0] /Rect [236.201 562.743 266.414 574.695] /A << /S /GoTo /D <6571756174696F6E2E312E382E312E32> >> >>}
\MusicPDFWriteObject{159}{<< /Type /Annot /Subtype /Link /Border [0 0 0] /H /I /C [1 0 0] /Rect [128.386 132.072 158.599 144.025] /A << /S /GoTo /D <6571756174696F6E2E312E382E322E32> >> >>}
\MusicPDFWriteObject{160}{<< /Type /Annot /Subtype /Link /Border [0 0 0] /H /I /C [0 1 0] /Rect [370.37 124.538 394.213 133.406] /A << /S /GoTo /D <636974652E4D6531> >> >>}
\MusicPDFWriteObject{161}{<< /Type /Annot /Subtype /Link /Border [0 0 0] /H /I /C [0 1 0] /Rect [339.925 647.671 357.398 656.539] /A << /S /GoTo /D <636974652E4346> >> >>}
\MusicPDFWriteObject{162}{<< /Type /Annot /Subtype /Link /Border [0 0 0] /H /I /C [0 1 0] /Rect [371.852 279.812 387.844 288.68] /A << /S /GoTo /D <636974652E436F> >> >>}
\MusicPDFWriteObject{163}{<< /Type /Annot /Subtype /Link /Border [0 0 0] /H /I /C [0 1 0] /Rect [239.366 266.661 255.358 275.529] /A << /S /GoTo /D <636974652E436F> >> >>}
\MusicPDFWriteObject{164}{<< /Type /Annot /Subtype /Link /Border [0 0 0] /H /I /C [1 0 0] /Rect [156.074 212.951 186.287 224.904] /A << /S /GoTo /D <6571756174696F6E2E312E382E332E32> >> >>}
\MusicPDFWriteObject{165}{<< /Type /Annot /Subtype /Link /Border [0 0 0] /H /I /C [1 0 0] /Rect [511.675 212.951 541.888 224.904] /A << /S /GoTo /D <6571756174696F6E2E312E382E332E32> >> >>}
\MusicPDFWriteObject{166}{<< /Type /Annot /Subtype /Link /Border [0 0 0] /H /I /C [1 0 0] /Rect [73.712 131.515 103.925 143.468] /A << /S /GoTo /D <6571756174696F6E2E372E322E302E31> >> >>}
\MusicPDFWriteObject{167}{<< /Type /Annot /Subtype /Link /Border [0 0 0] /H /I /C [1 0 0] /Rect [132.366 131.515 162.579 143.468] /A << /S /GoTo /D <6571756174696F6E2E312E382E332E32> >> >>}
\MusicPDFWriteObject{168}{<< /Type /Annot /Subtype /Link /Border [0 0 0] /H /I /C [1 0 0] /Rect [505.935 671.482 536.148 683.435] /A << /S /GoTo /D <6571756174696F6E2E312E382E322E33> >> >>}
\MusicPDFWriteObject{169}{<< /Type /Annot /Subtype /Link /Border [0 0 0] /H /I /C [1 0 0] /Rect [232.42 511.662 262.633 523.615] /A << /S /GoTo /D <6571756174696F6E2E372E322E302E31> >> >>}
\MusicPDFWriteObject{170}{<< /Type /Annot /Subtype /Link /Border [0 0 0] /H /I /C [1 0 0] /Rect [148.201 122.6 178.414 134.553] /A << /S /GoTo /D <6571756174696F6E2E312E382E332E34> >> >>}
\MusicPDFWriteObject{171}{<< /Type /Annot /Subtype /Link /Border [0 0 0] /H /I /C [1 0 0] /Rect [508.926 531.39 539.14 543.343] /A << /S /GoTo /D <6571756174696F6E2E312E382E332E34> >> >>}
\MusicPDFWriteObject{172}{<< /Type /Annot /Subtype /Link /Border [0 0 0] /H /I /C [1 0 0] /Rect [232.366 286.815 262.579 298.768] /A << /S /GoTo /D <6571756174696F6E2E312E382E332E34> >> >>}
\MusicPDFWriteObject{173}{<< /Type /Annot /Subtype /Link /Border [0 0 0] /H /I /C [1 0 0] /Rect [442.762 406.37 472.975 418.323] /A << /S /GoTo /D <6571756174696F6E2E312E382E332E36> >> >>}
\MusicPDFWriteObject{174}{<< /Type /Annot /Subtype /Link /Border [0 0 0] /H /I /C [1 0 0] /Rect [170.978 498.401 206.171 510.353] /A << /S /GoTo /D <6571756174696F6E2E31312E332E312E31> >> >>}
\MusicPDFWriteObject{175}{<< /Type /Annot /Subtype /Link /Border [0 0 0] /H /I /C [1 0 0] /Rect [70.506 543.493 100.719 555.446] /A << /S /GoTo /D <74686D2E312E382E312E31> >> >>}
\MusicPDFWriteObject{176}{<< /Type /Annot /Subtype /Link /Border [0 0 0] /H /I /C [1 0 0] /Rect [354.945 411.932 390.138 423.884] /A << /S /GoTo /D <6571756174696F6E2E31312E332E312E33> >> >>}
\MusicPDFWriteObject{177}{<< /Type /Annot /Subtype /Link /Border [0 0 0] /H /I /C [1 0 0] /Rect [139.539 332.286 174.732 344.239] /A << /S /GoTo /D <6571756174696F6E2E31312E332E312E33> >> >>}
\MusicPDFWriteObject{178}{<< /Type /Annot /Subtype /Link /Border [0 0 0] /H /I /C [1 0 0] /Rect [506.695 509.373 541.888 521.326] /A << /S /GoTo /D <6571756174696F6E2E312E332E312E3132> >> >>}
\MusicPDFWriteObject{179}{<< /Type /Annot /Subtype /Link /Border [0 0 0] /H /I /C [1 0 0] /Rect [430.86 117.805 436.822 130.663] /A << /S /GoTo /D <48666F6F746E6F74652E38> >> >>}
\MusicPDFWriteObject{180}{<< /Type /Annot /Subtype /Link /Border [0 0 0] /H /I /C [0 1 0] /Rect [104.153 560.398 121.626 569.266] /A << /S /GoTo /D <636974652E4346> >> >>}
\MusicPDFWriteObject{181}{<< /Type /Annot /Subtype /Link /Border [0 0 0] /H /I /C [0 1 0] /Rect [287.934 474.484 305.406 483.352] /A << /S /GoTo /D <636974652E4346> >> >>}
\MusicPDFWriteObject{182}{<< /Type /Annot /Subtype /Link /Border [0 0 0] /H /I /C [0 1 0] /Rect [525.687 448.182 543.159 457.05] /A << /S /GoTo /D <636974652E4346> >> >>}
\MusicPDFWriteObject{183}{<< /Type /Annot /Subtype /Link /Border [0 0 0] /H /I /C [0 1 0] /Rect [143.189 421.881 160.661 430.749] /A << /S /GoTo /D <636974652E4346> >> >>}
\MusicPDFWriteObject{184}{<< /Type /Annot /Subtype /Link /Border [0 0 0] /H /I /C [0 1 0] /Rect [137.302 235.678 154.774 244.546] /A << /S /GoTo /D <636974652E4346> >> >>}
\MusicPDFWriteObject{185}{<< /Type /Annot /Subtype /Link /Border [0 0 0] /H /I /C [1 0 0] /Rect [82.959 687.123 226.58 695.977] /A << /S /GoTo /D <73656374696F6E2E322E31> >> >>}
\MusicPDFWriteObject{186}{<< /Type /Annot /Subtype /Link /Border [0 0 0] /H /I /C [1 0 0] /Rect [82.959 673.973 234.327 682.827] /A << /S /GoTo /D <73756273656374696F6E2E322E312E31> >> >>}
\MusicPDFWriteObject{187}{<< /Type /Annot /Subtype /Link /Border [0 0 0] /H /I /C [1 0 0] /Rect [82.959 660.822 279.839 669.676] /A << /S /GoTo /D <73756273756273656374696F6E2E322E312E312E31> >> >>}
\MusicPDFWriteObject{188}{<< /Type /Annot /Subtype /Link /Border [0 0 0] /H /I /C [1 0 0] /Rect [82.959 645.734 301.088 656.525] /A << /S /GoTo /D <73756273756273656374696F6E2E322E312E312E32> >> >>}
\MusicPDFWriteObject{189}{<< /Type /Annot /Subtype /Link /Border [0 0 0] /H /I /C [1 0 0] /Rect [82.959 632.584 442.274 643.374] /A << /S /GoTo /D <73756273756273656374696F6E2E322E312E312E33> >> >>}
\MusicPDFWriteObject{190}{<< /Type /Annot /Subtype /Link /Border [0 0 0] /H /I /C [1 0 0] /Rect [82.959 621.369 261.302 630.223] /A << /S /GoTo /D <73756273656374696F6E2E322E312E32> >> >>}
\MusicPDFWriteObject{191}{<< /Type /Annot /Subtype /Link /Border [0 0 0] /H /I /C [1 0 0] /Rect [82.959 608.219 324.135 617.073] /A << /S /GoTo /D <73756273756273656374696F6E2E322E312E322E31> >> >>}
\MusicPDFWriteObject{192}{<< /Type /Annot /Subtype /Link /Border [0 0 0] /H /I /C [1 0 0] /Rect [82.959 592.218 443.632 604.532] /A << /S /GoTo /D <73756273756273656374696F6E2E322E312E322E32> >> >>}
\MusicPDFWriteObject{193}{<< /Type /Annot /Subtype /Link /Border [0 0 0] /H /I /C [1 0 0] /Rect [82.959 579.427 431.83 591.842] /A << /S /GoTo /D <73756273756273656374696F6E2E322E312E322E33> >> >>}
\MusicPDFWriteObject{194}{<< /Type /Annot /Subtype /Link /Border [0 0 0] /H /I /C [1 0 0] /Rect [82.959 568.766 264.761 577.62] /A << /S /GoTo /D <73756273656374696F6E2E322E312E33> >> >>}
\MusicPDFWriteObject{195}{<< /Type /Annot /Subtype /Link /Border [0 0 0] /H /I /C [1 0 0] /Rect [82.959 553.679 367.213 564.47] /A << /S /GoTo /D <73756273756273656374696F6E2E322E312E332E31> >> >>}
\MusicPDFWriteObject{196}{<< /Type /Annot /Subtype /Link /Border [0 0 0] /H /I /C [1 0 0] /Rect [82.959 539.974 247.406 551.929] /A << /S /GoTo /D <73756273656374696F6E2E322E312E34> >> >>}
\MusicPDFWriteObject{197}{<< /Type /Annot /Subtype /Link /Border [0 0 0] /H /I /C [1 0 0] /Rect [82.959 529.314 218.682 538.225] /A << /S /GoTo /D <73656374696F6E2E322E32> >> >>}
\MusicPDFWriteObject{198}{<< /Type /Annot /Subtype /Link /Border [0 0 0] /H /I /C [1 0 0] /Rect [82.959 514.227 238.56 525.017] /A << /S /GoTo /D <73756273656374696F6E2E322E322E31> >> >>}
\MusicPDFWriteObject{199}{<< /Type /Annot /Subtype /Link /Border [0 0 0] /H /I /C [1 0 0] /Rect [82.959 500.522 381.728 512.938] /A << /S /GoTo /D <73756273656374696F6E2E322E322E32> >> >>}
\MusicPDFWriteObject{200}{<< /Type /Annot /Subtype /Link /Border [0 0 0] /H /I /C [1 0 0] /Rect [82.959 487.925 261.358 498.716] /A << /S /GoTo /D <73756273756273656374696F6E2E322E322E322E31> >> >>}
\MusicPDFWriteObject{201}{<< /Type /Annot /Subtype /Link /Border [0 0 0] /H /I /C [1 0 0] /Rect [82.959 474.22 321.461 486.176] /A << /S /GoTo /D <73756273656374696F6E2E322E322E33> >> >>}
\MusicPDFWriteObject{202}{<< /Type /Annot /Subtype /Link /Border [0 0 0] /H /I /C [1 0 0] /Rect [82.959 461.624 284.847 472.414] /A << /S /GoTo /D <73756273756273656374696F6E2E322E322E332E31> >> >>}
\MusicPDFWriteObject{203}{<< /Type /Annot /Subtype /Link /Border [0 0 0] /H /I /C [1 0 0] /Rect [82.959 448.473 412.671 459.321] /A << /S /GoTo /D <73756273656374696F6E2E322E322E34> >> >>}
\MusicPDFWriteObject{204}{<< /Type /Annot /Subtype /Link /Border [0 0 0] /H /I /C [1 0 0] /Rect [82.959 435.322 370.284 446.113] /A << /S /GoTo /D <73756273756273656374696F6E2E322E322E342E31> >> >>}
\MusicPDFWriteObject{205}{<< /Type /Annot /Subtype /Link /Border [0 0 0] /H /I /C [1 0 0] /Rect [82.959 421.258 411.815 434.21] /A << /S /GoTo /D <73756273656374696F6E2E322E322E35> >> >>}
\MusicPDFWriteObject{206}{<< /Type /Annot /Subtype /Link /Border [0 0 0] /H /I /C [1 0 0] /Rect [82.959 409.021 347.57 419.868] /A << /S /GoTo /D <73756273756273656374696F6E2E322E322E352E31> >> >>}
\MusicPDFWriteObject{207}{<< /Type /Annot /Subtype /Link /Border [0 0 0] /H /I /C [1 0 0] /Rect [82.959 395.063 256.033 407.909] /A << /S /GoTo /D <73756273756273656374696F6E2E322E322E352E32> >> >>}
\MusicPDFWriteObject{208}{<< /Type /Annot /Subtype /Link /Border [0 0 0] /H /I /C [1 0 0] /Rect [82.959 384.656 220.04 393.567] /A << /S /GoTo /D <73656374696F6E2E322E33> >> >>}
\MusicPDFWriteObject{209}{<< /Type /Annot /Subtype /Link /Border [0 0 0] /H /I /C [1 0 0] /Rect [82.959 369.568 239.739 380.416] /A << /S /GoTo /D <73756273656374696F6E2E322E332E31> >> >>}
\MusicPDFWriteObject{210}{<< /Type /Annot /Subtype /Link /Border [0 0 0] /H /I /C [1 0 0] /Rect [82.959 355.504 212.336 367.208] /A << /S /GoTo /D <73756273756273656374696F6E2E322E332E312E31> >> >>}
\MusicPDFWriteObject{211}{<< /Type /Annot /Subtype /Link /Border [0 0 0] /H /I /C [1 0 0] /Rect [82.959 343.267 274.876 354.115] /A << /S /GoTo /D <73756273756273656374696F6E2E322E332E312E32> >> >>}
\MusicPDFWriteObject{212}{<< /Type /Annot /Subtype /Link /Border [0 0 0] /H /I /C [1 0 0] /Rect [82.959 327.953 236.246 342.155] /A << /S /GoTo /D <73756273756273656374696F6E2E322E332E312E33> >> >>}
\MusicPDFWriteObject{213}{<< /Type /Annot /Subtype /Link /Border [0 0 0] /H /I /C [1 0 0] /Rect [82.959 316.965 300.037 327.756] /A << /S /GoTo /D <73756273656374696F6E2E322E332E32> >> >>}
\MusicPDFWriteObject{214}{<< /Type /Annot /Subtype /Link /Border [0 0 0] /H /I /C [1 0 0] /Rect [82.959 302.901 379.15 315.676] /A << /S /GoTo /D <73656374696F6E2E322E34> >> >>}
\MusicPDFWriteObject{215}{<< /Type /Annot /Subtype /Link /Border [0 0 0] /H /I /C [1 0 0] /Rect [82.959 290.664 324.688 301.455] /A << /S /GoTo /D <73756273656374696F6E2E322E342E31> >> >>}
\MusicPDFWriteObject{216}{<< /Type /Annot /Subtype /Link /Border [0 0 0] /H /I /C [1 0 0] /Rect [82.959 277.513 190.723 288.304] /A << /S /GoTo /D <73756273656374696F6E2E322E342E32> >> >>}
\MusicPDFWriteObject{217}{<< /Type /Annot /Subtype /Link /Border [0 0 0] /H /I /C [1 0 0] /Rect [82.959 263.449 255.2 276.224] /A << /S /GoTo /D <73756273656374696F6E2E322E342E33> >> >>}
\MusicPDFWriteObject{218}{<< /Type /Annot /Subtype /Link /Border [0 0 0] /H /I /C [0 1 0] /Rect [203.723 628.543 221.922 637.411] /A << /S /GoTo /D <636974652E4A6172> >> >>}
\MusicPDFWriteObject{219}{<< /Type /Annot /Subtype /Link /Border [0 0 0] /H /I /C [1 0 0] /Rect [364.398 573.45 396.489 585.402] /A << /S /GoTo /D <74686D2E432E322E312E35> >> >>}
\MusicPDFWriteObject{220}{<< /Type /Annot /Subtype /Link /Border [0 0 0] /H /I /C [0 1 0] /Rect [514.281 270.773 543.159 279.641] /A << /S /GoTo /D <636974652E5363682D61> >> >>}
\MusicPDFWriteObject{221}{<< /Type /Annot /Subtype /Link /Border [0 0 0] /H /I /C [0 1 0] /Rect [455.541 244.472 480.601 253.34] /A << /S /GoTo /D <636974652E53636861> >> >>}
\MusicPDFWriteObject{222}{<< /Type /Annot /Subtype /Link /Border [0 0 0] /H /I /C [0 1 0] /Rect [93.197 231.321 111.396 240.189] /A << /S /GoTo /D <636974652E4A6172> >> >>}
\MusicPDFWriteObject{223}{<< /Type /Annot /Subtype /Link /Border [0 0 0] /H /I /C [0 1 0] /Rect [196.675 231.321 225.553 240.189] /A << /S /GoTo /D <636974652E5363682D61> >> >>}
\MusicPDFWriteObject{224}{<< /Type /Annot /Subtype /Link /Border [0 0 0] /H /I /C [0 1 0] /Rect [211.749 191.869 240.627 200.737] /A << /S /GoTo /D <636974652E5363682D61> >> >>}
\MusicPDFWriteObject{225}{<< /Type /Annot /Subtype /Link /Border [0 0 0] /H /I /C [0 1 0] /Rect [326.649 178.718 355.527 187.586] /A << /S /GoTo /D <636974652E5363682D61> >> >>}
\MusicPDFWriteObject{226}{<< /Type /Annot /Subtype /Link /Border [0 0 0] /H /I /C [0 1 0] /Rect [441.96 98.237 459.862 107.105] /A << /S /GoTo /D <636974652E536869> >> >>}
\MusicPDFWriteObject{227}{<< /Type /Annot /Subtype /Link /Border [0 0 0] /H /I /C [0 1 0] /Rect [514.459 523.351 543.22 532.219] /A << /S /GoTo /D <636974652E47726F2D66> >> >>}
\MusicPDFWriteObject{228}{<< /Type /Annot /Subtype /Link /Border [0 0 0] /H /I /C [0 1 0] /Rect [289.382 502.881 309.04 511.748] /A << /S /GoTo /D <636974652E4B47> >> >>}
\MusicPDFWriteObject{229}{<< /Type /Annot /Subtype /Link /Border [0 0 0] /H /I /C [0 1 0] /Rect [524.96 454.323 543.159 463.19] /A << /S /GoTo /D <636974652E4A6172> >> >>}
\MusicPDFWriteObject{230}{<< /Type /Annot /Subtype /Link /Border [0 0 0] /H /I /C [0 1 0] /Rect [115.27 401.72 133.469 410.587] /A << /S /GoTo /D <636974652E4A6172> >> >>}
\MusicPDFWriteObject{231}{<< /Type /Annot /Subtype /Link /Border [0 0 0] /H /I /C [0 1 0] /Rect [324.096 401.72 343.201 410.587] /A << /S /GoTo /D <636974652E4E42> >> >>}
\MusicPDFWriteObject{232}{<< /Type /Annot /Subtype /Link /Border [0 0 0] /H /I /C [1 0 0] /Rect [457.407 399.229 490.442 411.182] /A << /S /GoTo /D <74686D2E432E322E312E35> >> >>}
\MusicPDFWriteObject{233}{<< /Type /Annot /Subtype /Link /Border [0 0 0] /H /I /C [0 1 0] /Rect [204.911 388.569 223.11 397.437] /A << /S /GoTo /D <636974652E4A6172> >> >>}
\MusicPDFWriteObject{234}{<< /Type /Annot /Subtype /Link /Border [0 0 0] /H /I /C [0 1 0] /Rect [434.547 388.569 459.608 397.437] /A << /S /GoTo /D <636974652E53636861> >> >>}
\MusicPDFWriteObject{235}{<< /Type /Annot /Subtype /Link /Border [0 0 0] /H /I /C [0 1 0] /Rect [157.007 291.453 175.206 300.321] /A << /S /GoTo /D <636974652E4A6172> >> >>}
\MusicPDFWriteObject{236}{<< /Type /Annot /Subtype /Link /Border [0 0 0] /H /I /C [0 1 0] /Rect [256.666 291.453 274.865 300.321] /A << /S /GoTo /D <636974652E4A6172> >> >>}
\MusicPDFWriteObject{237}{<< /Type /Annot /Subtype /Link /Border [0 0 0] /H /I /C [0 1 0] /Rect [92.16 278.302 110.358 287.17] /A << /S /GoTo /D <636974652E4A6172> >> >>}
\MusicPDFWriteObject{238}{<< /Type /Annot /Subtype /Link /Border [0 0 0] /H /I /C [1 0 0] /Rect [203.197 242.895 236.758 251.749] /A << /S /GoTo /D <74686D2E322E312E312E36> >> >>}
\MusicPDFWriteObject{239}{<< /Type /Annot /Subtype /Link /Border [0 0 0] /H /I /C [0 1 0] /Rect [353.531 242.895 372.504 251.763] /A << /S /GoTo /D <636974652E547265> >> >>}
\MusicPDFWriteObject{240}{<< /Type /Annot /Subtype /Link /Border [0 0 0] /H /I /C [1 0 0] /Rect [221.868 227.254 252.081 239.207] /A << /S /GoTo /D <6571756174696F6E2E322E312E312E34> >> >>}
\MusicPDFWriteObject{241}{<< /Type /Annot /Subtype /Link /Border [0 0 0] /H /I /C [1 0 0] /Rect [184.847 200.953 215.06 212.905] /A << /S /GoTo /D <6571756174696F6E2E322E312E312E34> >> >>}
\MusicPDFWriteObject{242}{<< /Type /Annot /Subtype /Link /Border [0 0 0] /H /I /C [0 1 0] /Rect [415.62 203.443 456.229 212.311] /A << /S /GoTo /D <636974652E4B6F7468652D49> >> >>}
\MusicPDFWriteObject{243}{<< /Type /Annot /Subtype /Link /Border [0 0 0] /H /I /C [1 0 0] /Rect [444.799 187.802 475.012 199.755] /A << /S /GoTo /D <6571756174696F6E2E322E312E312E33> >> >>}
\MusicPDFWriteObject{244}{<< /Type /Annot /Subtype /Link /Border [0 0 0] /H /I /C [0 1 0] /Rect [129.299 137.689 146.453 146.557] /A << /S /GoTo /D <636974652E466F6C> >> >>}
\MusicPDFWriteObject{245}{<< /Type /Annot /Subtype /Link /Border [0 0 0] /H /I /C [0 1 0] /Rect [498.121 608.219 517.094 617.087] /A << /S /GoTo /D <636974652E547265> >> >>}
\MusicPDFWriteObject{246}{<< /Type /Annot /Subtype /Link /Border [0 0 0] /H /I /C [0 1 0] /Rect [73.77 581.917 91.969 590.785] /A << /S /GoTo /D <636974652E4A6172> >> >>}
\MusicPDFWriteObject{247}{<< /Type /Annot /Subtype /Link /Border [0 0 0] /H /I /C [0 1 0] /Rect [107.254 568.766 126.358 577.634] /A << /S /GoTo /D <636974652E4E42> >> >>}
\MusicPDFWriteObject{248}{<< /Type /Annot /Subtype /Link /Border [0 0 0] /H /I /C [0 1 0] /Rect [164.569 555.616 183.542 564.484] /A << /S /GoTo /D <636974652E547265> >> >>}
\MusicPDFWriteObject{249}{<< /Type /Annot /Subtype /Link /Border [0 0 0] /H /I /C [0 1 0] /Rect [73.77 542.465 92.744 551.333] /A << /S /GoTo /D <636974652E547265> >> >>}
\MusicPDFWriteObject{250}{<< /Type /Annot /Subtype /Link /Border [0 0 0] /H /I /C [0 1 0] /Rect [326.421 542.465 355.299 551.333] /A << /S /GoTo /D <636974652E5363682D61> >> >>}
\MusicPDFWriteObject{251}{<< /Type /Annot /Subtype /Link /Border [0 0 0] /H /I /C [0 1 0] /Rect [173.818 529.314 192.016 538.182] /A << /S /GoTo /D <636974652E4A6172> >> >>}
\MusicPDFWriteObject{252}{<< /Type /Annot /Subtype /Link /Border [0 0 0] /H /I /C [0 1 0] /Rect [524.993 529.314 543.192 538.182] /A << /S /GoTo /D <636974652E4A6172> >> >>}
\MusicPDFWriteObject{253}{<< /Type /Annot /Subtype /Link /Border [0 0 0] /H /I /C [1 0 0] /Rect [489.39 451.612 519.66 462.403] /A << /S /GoTo /D <74686D2E322E312E312E34> >> >>}
\MusicPDFWriteObject{254}{<< /Type /Annot /Subtype /Link /Border [0 0 0] /H /I /C [0 1 0] /Rect [142.912 400.946 164.133 409.814] /A << /S /GoTo /D <636974652E427275> >> >>}
\MusicPDFWriteObject{255}{<< /Type /Annot /Subtype /Link /Border [0 0 0] /H /I /C [0 1 0] /Rect [255.515 400.946 276.736 409.814] /A << /S /GoTo /D <636974652E427275> >> >>}
\MusicPDFWriteObject{256}{<< /Type /Annot /Subtype /Link /Border [0 0 0] /H /I /C [1 0 0] /Rect [137.274 372.154 169.498 384.107] /A << /S /GoTo /D <74686D2E442E312E332E37> >> >>}
\MusicPDFWriteObject{257}{<< /Type /Annot /Subtype /Link /Border [0 0 0] /H /I /C [1 0 0] /Rect [123.571 332.701 153.22 344.656] /A << /S /GoTo /D <74686D2E322E312E312E36> >> >>}
\MusicPDFWriteObject{258}{<< /Type /Annot /Subtype /Link /Border [0 0 0] /H /I /C [0 1 0] /Rect [335.226 256.287 356.448 265.155] /A << /S /GoTo /D <636974652E427275> >> >>}
\MusicPDFWriteObject{259}{<< /Type /Annot /Subtype /Link /Border [0 0 0] /H /I /C [0 1 0] /Rect [478.695 229.937 496.998 238.805] /A << /S /GoTo /D <636974652E4445> >> >>}
\MusicPDFWriteObject{260}{<< /Type /Annot /Subtype /Link /Border [0 0 0] /H /I /C [0 1 0] /Rect [93.138 216.786 111.44 225.654] /A << /S /GoTo /D <636974652E4445> >> >>}
\MusicPDFWriteObject{261}{<< /Type /Annot /Subtype /Link /Border [0 0 0] /H /I /C [0 1 0] /Rect [186.39 216.786 204.692 225.654] /A << /S /GoTo /D <636974652E4445> >> >>}
\MusicPDFWriteObject{262}{<< /Type /Annot /Subtype /Link /Border [0 0 0] /H /I /C [0 1 0] /Rect [439.525 203.635 460.746 212.503] /A << /S /GoTo /D <636974652E427275> >> >>}
\MusicPDFWriteObject{263}{<< /Type /Annot /Subtype /Link /Border [0 0 0] /H /I /C [1 0 0] /Rect [417.352 109.089 449.978 121.044] /A << /S /GoTo /D <74686D2E442E312E352E37> >> >>}
\MusicPDFWriteObject{264}{<< /Type /Annot /Subtype /Link /Border [0 0 0] /H /I /C [1 0 0] /Rect [262.506 95.938 294.761 107.894] /A << /S /GoTo /D <74686D2E442E312E352E36> >> >>}
\MusicPDFWriteObject{265}{<< /Type /Annot /Subtype /Link /Border [0 0 0] /H /I /C [1 0 0] /Rect [354.776 81.101 387.617 94.761] /A << /S /GoTo /D <74686D2E442E312E352E37> >> >>}
\MusicPDFWriteObject{266}{<< /Type /Annot /Subtype /Link /Border [0 0 0] /H /I /C [0 1 0] /Rect [301.183 199.816 322.404 208.684] /A << /S /GoTo /D <636974652E427275> >> >>}
\MusicPDFWriteObject{267}{<< /Type /Annot /Subtype /Link /Border [0 0 0] /H /I /C [1 0 0] /Rect [124.623 207.058 157.049 219.011] /A << /S /GoTo /D <74686D2E432E322E352E31> >> >>}
\MusicPDFWriteObject{268}{<< /Type /Annot /Subtype /Link /Border [0 0 0] /H /I /C [1 0 0] /Rect [280.423 685.187 311.201 695.977] /A << /S /GoTo /D <74686D2E322E312E322E35> >> >>}
\MusicPDFWriteObject{269}{<< /Type /Annot /Subtype /Link /Border [0 0 0] /H /I /C [1 0 0] /Rect [287.335 413.183 318.112 427.734] /A << /S /GoTo /D <74686D2E322E312E322E38> >> >>}
\MusicPDFWriteObject{270}{<< /Type /Annot /Subtype /Link /Border [0 0 0] /H /I /C [1 0 0] /Rect [186.874 387.873 219.301 398.664] /A << /S /GoTo /D <74686D2E432E322E352E31> >> >>}
\MusicPDFWriteObject{271}{<< /Type /Annot /Subtype /Link /Border [0 0 0] /H /I /C [1 0 0] /Rect [499.918 374.168 532.831 386.123] /A << /S /GoTo /D <74686D2E432E322E352E31> >> >>}
\MusicPDFWriteObject{272}{<< /Type /Annot /Subtype /Link /Border [0 0 0] /H /I /C [1 0 0] /Rect [448.044 230.754 480.714 242.71] /A << /S /GoTo /D <74686D2E432E322E352E31> >> >>}
\MusicPDFWriteObject{273}{<< /Type /Annot /Subtype /Link /Border [0 0 0] /H /I /C [1 0 0] /Rect [215.891 347.027 245.54 358.982] /A << /S /GoTo /D <74686D2E322E312E322E34> >> >>}
\MusicPDFWriteObject{274}{<< /Type /Annot /Subtype /Link /Border [0 0 0] /H /I /C [1 0 0] /Rect [435.301 320.428 467.727 332.681] /A << /S /GoTo /D <74686D2E432E342E332E34> >> >>}
\MusicPDFWriteObject{275}{<< /Type /Annot /Subtype /Link /Border [0 0 0] /H /I /C [1 0 0] /Rect [155.616 307.575 187.434 320.439] /A << /S /GoTo /D <74686D2E432E322E352E31> >> >>}
\MusicPDFWriteObject{276}{<< /Type /Annot /Subtype /Link /Border [0 0 0] /H /I /C [1 0 0] /Rect [155.748 254.275 185.961 266.228] /A << /S /GoTo /D <6571756174696F6E2E322E312E322E33> >> >>}
\MusicPDFWriteObject{277}{<< /Type /Annot /Subtype /Link /Border [0 0 0] /H /I /C [1 0 0] /Rect [235.79 240.776 266.003 252.729] /A << /S /GoTo /D <6571756174696F6E2E322E312E322E33> >> >>}
\MusicPDFWriteObject{278}{<< /Type /Annot /Subtype /Link /Border [0 0 0] /H /I /C [1 0 0] /Rect [303.05 175.205 338.908 185.995] /A << /S /GoTo /D <74686D2E322E312E322E3130> >> >>}
\MusicPDFWriteObject{279}{<< /Type /Annot /Subtype /Link /Border [0 0 0] /H /I /C [1 0 0] /Rect [198.909 672.036 231.751 682.827] /A << /S /GoTo /D <74686D2E442E312E392E32> >> >>}
\MusicPDFWriteObject{280}{<< /Type /Annot /Subtype /Link /Border [0 0 0] /H /I /C [1 0 0] /Rect [205.601 578.15 235.249 594.037] /A << /S /GoTo /D <74686D2E322E312E322E34> >> >>}
\MusicPDFWriteObject{281}{<< /Type /Annot /Subtype /Link /Border [0 0 0] /H /I /C [1 0 0] /Rect [307.114 578.15 336.763 594.037] /A << /S /GoTo /D <74686D2E322E312E322E38> >> >>}
\MusicPDFWriteObject{282}{<< /Type /Annot /Subtype /Link /Border [0 0 0] /H /I /C [1 0 0] /Rect [190.707 551.081 225.701 567.259] /A << /S /GoTo /D <74686D2E322E312E322E3130> >> >>}
\MusicPDFWriteObject{283}{<< /Type /Annot /Subtype /Link /Border [0 0 0] /H /I /C [1 0 0] /Rect [397.382 540.277 427.567 551.068] /A << /S /GoTo /D <74686D2E342E332E332E31> >> >>}
\MusicPDFWriteObject{284}{<< /Type /Annot /Subtype /Link /Border [0 0 0] /H /I /C [1 0 0] /Rect [400.677 473.169 415.295 484.017] /A << /S /GoTo /D <73656374696F6E2E322E31> >> >>}
\MusicPDFWriteObject{285}{<< /Type /Annot /Subtype /Link /Border [0 0 0] /H /I /C [0 1 0] /Rect [469.183 461.955 490.404 470.823] /A << /S /GoTo /D <636974652E427275> >> >>}
\MusicPDFWriteObject{286}{<< /Type /Annot /Subtype /Link /Border [0 0 0] /H /I /C [1 0 0] /Rect [465.82 324.434 498.923 340.36] /A << /S /GoTo /D <74686D2E412E312E312E33> >> >>}
\MusicPDFWriteObject{287}{<< /Type /Annot /Subtype /Link /Border [0 0 0] /H /I /C [1 0 0] /Rect [111.403 196.347 141.052 212.273] /A << /S /GoTo /D <74686D2E322E312E322E35> >> >>}
\MusicPDFWriteObject{288}{<< /Type /Annot /Subtype /Link /Border [0 0 0] /H /I /C [1 0 0] /Rect [70.755 166.487 100.968 182.413] /A << /S /GoTo /D <74686D2E322E312E322E38> >> >>}
\MusicPDFWriteObject{289}{<< /Type /Annot /Subtype /Link /Border [0 0 0] /H /I /C [1 0 0] /Rect [254.405 658.332 284.618 670.284] /A << /S /GoTo /D <6571756174696F6E2E322E312E332E31> >> >>}
\MusicPDFWriteObject{290}{<< /Type /Annot /Subtype /Link /Border [0 0 0] /H /I /C [1 0 0] /Rect [264.928 615.287 294.746 631.213] /A << /S /GoTo /D <74686D2E322E312E332E32> >> >>}
\MusicPDFWriteObject{291}{<< /Type /Annot /Subtype /Link /Border [0 0 0] /H /I /C [0 1 0] /Rect [73.77 580.138 98.831 589.006] /A << /S /GoTo /D <636974652E53636861> >> >>}
\MusicPDFWriteObject{292}{<< /Type /Annot /Subtype /Link /Border [0 0 0] /H /I /C [1 0 0] /Rect [303.05 564.497 340.456 576.45] /A << /S /GoTo /D <6571756174696F6E2E432E322E332E3232> >> >>}
\MusicPDFWriteObject{293}{<< /Type /Annot /Subtype /Link /Border [0 0 0] /H /I /C [1 0 0] /Rect [374.197 564.497 399.33 576.45] /A << /S /GoTo /D <73756273656374696F6E2E432E322E33> >> >>}
\MusicPDFWriteObject{294}{<< /Type /Annot /Subtype /Link /Border [0 0 0] /H /I /C [1 0 0] /Rect [330.095 551.9 367.005 562.69] /A << /S /GoTo /D <74686D2E432E322E332E3133> >> >>}
\MusicPDFWriteObject{295}{<< /Type /Annot /Subtype /Link /Border [0 0 0] /H /I /C [1 0 0] /Rect [493.664 511.894 525.95 523.847] /A << /S /GoTo /D <74686D2E442E312E382E31> >> >>}
\MusicPDFWriteObject{296}{<< /Type /Annot /Subtype /Link /Border [0 0 0] /H /I /C [1 0 0] /Rect [350.166 499.297 380.379 510.087] /A << /S /GoTo /D <74686D2E322E312E332E33> >> >>}
\MusicPDFWriteObject{297}{<< /Type /Annot /Subtype /Link /Border [0 0 0] /H /I /C [1 0 0] /Rect [511.675 375.103 541.888 387.056] /A << /S /GoTo /D <6571756174696F6E2E322E312E332E32> >> >>}
\MusicPDFWriteObject{298}{<< /Type /Annot /Subtype /Link /Border [0 0 0] /H /I /C [0 1 0] /Rect [518.149 253.021 543.209 261.889] /A << /S /GoTo /D <636974652E53636861> >> >>}
\MusicPDFWriteObject{299}{<< /Type /Annot /Subtype /Link /Border [0 0 0] /H /I /C [1 0 0] /Rect [333.164 671.482 366.622 683.435] /A << /S /GoTo /D <74686D2E442E312E382E31> >> >>}
\MusicPDFWriteObject{300}{<< /Type /Annot /Subtype /Link /Border [0 0 0] /H /I /C [1 0 0] /Rect [199.989 645.181 225.123 657.134] /A << /S /GoTo /D <73756273656374696F6E2E432E322E33> >> >>}
\MusicPDFWriteObject{301}{<< /Type /Annot /Subtype /Link /Border [0 0 0] /H /I /C [0 1 0] /Rect [318.05 647.671 348.401 656.539] /A << /S /GoTo /D <636974652E47726F2D63> >> >>}
\MusicPDFWriteObject{302}{<< /Type /Annot /Subtype /Link /Border [0 0 0] /H /I /C [0 1 0] /Rect [122.289 634.52 147.35 643.388] /A << /S /GoTo /D <636974652E53636861> >> >>}
\MusicPDFWriteObject{303}{<< /Type /Annot /Subtype /Link /Border [0 0 0] /H /I /C [0 1 0] /Rect [214.996 634.52 233.195 643.388] /A << /S /GoTo /D <636974652E4A6172> >> >>}
\MusicPDFWriteObject{304}{<< /Type /Annot /Subtype /Link /Border [0 0 0] /H /I /C [1 0 0] /Rect [483.131 618.879 520.537 630.832] /A << /S /GoTo /D <6571756174696F6E2E432E322E332E3233> >> >>}
\MusicPDFWriteObject{305}{<< /Type /Annot /Subtype /Link /Border [0 0 0] /H /I /C [1 0 0] /Rect [190.803 606.282 228.209 617.073] /A << /S /GoTo /D <74686D2E432E322E332E3133> >> >>}
\MusicPDFWriteObject{306}{<< /Type /Annot /Subtype /Link /Border [0 0 0] /H /I /C [0 1 0] /Rect [182.033 492.483 213.658 501.351] /A << /S /GoTo /D <636974652E47726F2D62> >> >>}
\MusicPDFWriteObject{307}{<< /Type /Annot /Subtype /Link /Border [0 0 0] /H /I /C [1 0 0] /Rect [259.975 479.332 290.188 488.186] /A << /S /GoTo /D <74686D2E342E352E312E33> >> >>}
\MusicPDFWriteObject{308}{<< /Type /Annot /Subtype /Link /Border [0 0 0] /H /I /C [1 0 0] /Rect [262.303 464.245 292.968 475.035] /A << /S /GoTo /D <74686D2E322E312E332E34> >> >>}
\MusicPDFWriteObject{309}{<< /Type /Annot /Subtype /Link /Border [0 0 0] /H /I /C [1 0 0] /Rect [203.197 429.003 236.758 437.857] /A << /S /GoTo /D <74686D2E322E312E332E32> >> >>}
\MusicPDFWriteObject{310}{<< /Type /Annot /Subtype /Link /Border [0 0 0] /H /I /C [0 1 0] /Rect [388.637 415.852 409.858 424.72] /A << /S /GoTo /D <636974652E427275> >> >>}
\MusicPDFWriteObject{311}{<< /Type /Annot /Subtype /Link /Border [0 0 0] /H /I /C [1 0 0] /Rect [442.64 251.816 475.219 266.367] /A << /S /GoTo /D <74686D2E432E322E352E31> >> >>}
\MusicPDFWriteObject{312}{<< /Type /Annot /Subtype /Link /Border [0 0 0] /H /I /C [1 0 0] /Rect [253.746 238.262 283.959 252.813] /A << /S /GoTo /D <74686D2E322E312E322E38> >> >>}
\MusicPDFWriteObject{313}{<< /Type /Annot /Subtype /Link /Border [0 0 0] /H /I /C [1 0 0] /Rect [154.728 223.458 184.376 239.258] /A << /S /GoTo /D <74686D2E322E312E322E38> >> >>}
\MusicPDFWriteObject{314}{<< /Type /Annot /Subtype /Link /Border [0 0 0] /H /I /C [0 1 0] /Rect [423.649 227.635 442.623 236.503] /A << /S /GoTo /D <636974652E547265> >> >>}
\MusicPDFWriteObject{315}{<< /Type /Annot /Subtype /Link /Border [0 0 0] /H /I /C [0 1 0] /Rect [180.947 214.484 199.146 223.352] /A << /S /GoTo /D <636974652E4A6172> >> >>}
\MusicPDFWriteObject{316}{<< /Type /Annot /Subtype /Link /Border [0 0 0] /H /I /C [0 1 0] /Rect [500.957 214.484 519.93 223.352] /A << /S /GoTo /D <636974652E547265> >> >>}
\MusicPDFWriteObject{317}{<< /Type /Annot /Subtype /Link /Border [0 0 0] /H /I /C [0 1 0] /Rect [272.947 111.387 291.146 120.255] /A << /S /GoTo /D <636974652E4A6172> >> >>}
\MusicPDFWriteObject{318}{<< /Type /Annot /Subtype /Link /Border [0 0 0] /H /I /C [1 0 0] /Rect [442.77 95.386 472.419 107.701] /A << /S /GoTo /D <74686D2E322E312E312E36> >> >>}
\MusicPDFWriteObject{319}{<< /Type /Annot /Subtype /Link /Border [0 0 0] /H /I /C [0 1 0] /Rect [354.566 85.086 397.589 93.954] /A << /S /GoTo /D <636974652E426F752D4754> >> >>}
\MusicPDFWriteObject{320}{<< /Type /Annot /Subtype /Link /Border [0 0 0] /H /I /C [0 1 0] /Rect [165.626 660.822 183.824 669.69] /A << /S /GoTo /D <636974652E4A6172> >> >>}
\MusicPDFWriteObject{321}{<< /Type /Annot /Subtype /Link /Border [0 0 0] /H /I /C [0 1 0] /Rect [411.791 647.671 429.99 656.539] /A << /S /GoTo /D <636974652E4A6172> >> >>}
\MusicPDFWriteObject{322}{<< /Type /Annot /Subtype /Link /Border [0 0 0] /H /I /C [0 1 0] /Rect [445.396 634.52 463.594 643.388] /A << /S /GoTo /D <636974652E4A6172> >> >>}
\MusicPDFWriteObject{323}{<< /Type /Annot /Subtype /Link /Border [0 0 0] /H /I /C [0 1 0] /Rect [149.025 620.109 167.224 628.977] /A << /S /GoTo /D <636974652E4A6172> >> >>}
\MusicPDFWriteObject{324}{<< /Type /Annot /Subtype /Link /Border [0 0 0] /H /I /C [1 0 0] /Rect [411.373 590.476 444.408 605.027] /A << /S /GoTo /D <74686D2E432E322E352E31> >> >>}
\MusicPDFWriteObject{325}{<< /Type /Annot /Subtype /Link /Border [0 0 0] /H /I /C [0 1 0] /Rect [263.706 579.849 281.905 588.717] /A << /S /GoTo /D <636974652E4A6172> >> >>}
\MusicPDFWriteObject{326}{<< /Type /Annot /Subtype /Link /Border [0 0 0] /H /I /C [1 0 0] /Rect [71.004 527.611 108.587 545.178] /A << /S /GoTo /D <74686D2E432E322E332E3133> >> >>}
\MusicPDFWriteObject{327}{<< /Type /Annot /Subtype /Link /Border [0 0 0] /H /I /C [1 0 0] /Rect [317.037 477.811 347.25 489.763] /A << /S /GoTo /D <6571756174696F6E2E322E312E332E32> >> >>}
\MusicPDFWriteObject{328}{<< /Type /Annot /Subtype /Link /Border [0 0 0] /H /I /C [1 0 0] /Rect [125.649 451.149 156.427 463.464] /A << /S /GoTo /D <74686D2E322E312E312E36> >> >>}
\MusicPDFWriteObject{329}{<< /Type /Annot /Subtype /Link /Border [0 0 0] /H /I /C [0 1 0] /Rect [153.578 414.547 178.638 423.415] /A << /S /GoTo /D <636974652E53636861> >> >>}
\MusicPDFWriteObject{330}{<< /Type /Annot /Subtype /Link /Border [0 0 0] /H /I /C [1 0 0] /Rect [512.709 386.309 545.592 397.099] /A << /S /GoTo /D <74686D2E432E312E362E31> >> >>}
\MusicPDFWriteObject{331}{<< /Type /Annot /Subtype /Link /Border [0 0 0] /H /I /C [0 1 0] /Rect [350.37 361.944 375.43 370.812] /A << /S /GoTo /D <636974652E53636861> >> >>}
\MusicPDFWriteObject{332}{<< /Type /Annot /Subtype /Link /Border [0 0 0] /H /I /C [0 1 0] /Rect [518.209 361.944 543.27 370.812] /A << /S /GoTo /D <636974652E53636861> >> >>}
\MusicPDFWriteObject{333}{<< /Type /Annot /Subtype /Link /Border [0 0 0] /H /I /C [0 1 0] /Rect [289.908 348.793 314.968 357.661] /A << /S /GoTo /D <636974652E53636861> >> >>}
\MusicPDFWriteObject{334}{<< /Type /Annot /Subtype /Link /Border [0 0 0] /H /I /C [0 1 0] /Rect [384.017 277.601 402.99 286.468] /A << /S /GoTo /D <636974652E547265> >> >>}
\MusicPDFWriteObject{335}{<< /Type /Annot /Subtype /Link /Border [0 0 0] /H /I /C [0 1 0] /Rect [448.14 277.601 479.765 286.468] /A << /S /GoTo /D <636974652E47726F2D62> >> >>}
\MusicPDFWriteObject{336}{<< /Type /Annot /Subtype /Link /Border [0 0 0] /H /I /C [0 1 0] /Rect [518.099 277.601 543.159 286.468] /A << /S /GoTo /D <636974652E53636861> >> >>}
\MusicPDFWriteObject{337}{<< /Type /Annot /Subtype /Link /Border [0 0 0] /H /I /C [0 1 0] /Rect [138.236 264.45 163.296 273.318] /A << /S /GoTo /D <636974652E53636861> >> >>}
\MusicPDFWriteObject{338}{<< /Type /Annot /Subtype /Link /Border [0 0 0] /H /I /C [0 1 0] /Rect [287.523 264.45 319.147 273.318] /A << /S /GoTo /D <636974652E47726F2D62> >> >>}
\MusicPDFWriteObject{339}{<< /Type /Annot /Subtype /Link /Border [0 0 0] /H /I /C [0 1 0] /Rect [124.34 227.618 142.539 236.486] /A << /S /GoTo /D <636974652E4A6172> >> >>}
\MusicPDFWriteObject{340}{<< /Type /Annot /Subtype /Link /Border [0 0 0] /H /I /C [0 1 0] /Rect [433.634 214.064 465.258 222.932] /A << /S /GoTo /D <636974652E47726F2D62> >> >>}
\MusicPDFWriteObject{341}{<< /Type /Annot /Subtype /Link /Border [0 0 0] /H /I /C [1 0 0] /Rect [173.227 197.581 203.44 212.132] /A << /S /GoTo /D <74686D2E322E312E332E33> >> >>}
\MusicPDFWriteObject{342}{<< /Type /Annot /Subtype /Link /Border [0 0 0] /H /I /C [0 1 0] /Rect [378.683 142.797 397.657 151.665] /A << /S /GoTo /D <636974652E547265> >> >>}
\MusicPDFWriteObject{343}{<< /Type /Annot /Subtype /Link /Border [0 0 0] /H /I /C [1 0 0] /Rect [212.517 583.551 242.73 595.504] /A << /S /GoTo /D <6571756174696F6E2E322E312E342E32> >> >>}
\MusicPDFWriteObject{344}{<< /Type /Annot /Subtype /Link /Border [0 0 0] /H /I /C [1 0 0] /Rect [186.571 432.695 216.785 444.648] /A << /S /GoTo /D <6571756174696F6E2E322E312E342E32> >> >>}
\MusicPDFWriteObject{345}{<< /Type /Annot /Subtype /Link /Border [0 0 0] /H /I /C [1 0 0] /Rect [157.392 212.239 187.238 226.719] /A << /S /GoTo /D <74686D2E322E312E322E34> >> >>}
\MusicPDFWriteObject{346}{<< /Type /Annot /Subtype /Link /Border [0 0 0] /H /I /C [1 0 0] /Rect [242.879 212.239 274.91 226.719] /A << /S /GoTo /D <74686D2E432E322E352E32> >> >>}
\MusicPDFWriteObject{347}{<< /Type /Annot /Subtype /Link /Border [0 0 0] /H /I /C [0 1 0] /Rect [502.566 203.773 520.765 212.641] /A << /S /GoTo /D <636974652E4A6172> >> >>}
\MusicPDFWriteObject{348}{<< /Type /Annot /Subtype /Link /Border [0 0 0] /H /I /C [1 0 0] /Rect [269.815 188.132 300.028 200.085] /A << /S /GoTo /D <6571756174696F6E2E322E312E342E31> >> >>}
\MusicPDFWriteObject{349}{<< /Type /Annot /Subtype /Link /Border [0 0 0] /H /I /C [0 1 0] /Rect [364.59 164.321 382.789 173.189] /A << /S /GoTo /D <636974652E4A6172> >> >>}
\MusicPDFWriteObject{350}{<< /Type /Annot /Subtype /Link /Border [0 0 0] /H /I /C [1 0 0] /Rect [71.004 135.169 108.41 147.484] /A << /S /GoTo /D <74686D2E432E322E332E3132> >> >>}
\MusicPDFWriteObject{351}{<< /Type /Annot /Subtype /Link /Border [0 0 0] /H /I /C [0 1 0] /Rect [108.118 124.868 133.178 133.736] /A << /S /GoTo /D <636974652E53636861> >> >>}
\MusicPDFWriteObject{352}{<< /Type /Annot /Subtype /Link /Border [0 0 0] /H /I /C [1 0 0] /Rect [411.968 119.854 450.082 134.333] /A << /S /GoTo /D <74686D2E432E322E332E3133> >> >>}
\MusicPDFWriteObject{353}{<< /Type /Annot /Subtype /Link /Border [0 0 0] /H /I /C [0 1 0] /Rect [152.382 111.387 170.58 120.255] /A << /S /GoTo /D <636974652E4A6172> >> >>}
\MusicPDFWriteObject{354}{<< /Type /Annot /Subtype /Link /Border [0 0 0] /H /I /C [0 1 0] /Rect [232.746 98.237 250.945 107.105] /A << /S /GoTo /D <636974652E4A6172> >> >>}
\MusicPDFWriteObject{355}{<< /Type /Annot /Subtype /Link /Border [0 0 0] /H /I /C [1 0 0] /Rect [174.374 82.596 204.587 94.549] /A << /S /GoTo /D <6571756174696F6E2E322E312E342E31> >> >>}
\MusicPDFWriteObject{356}{<< /Type /Annot /Subtype /Link /Border [0 0 0] /H /I /C [1 0 0] /Rect [189.983 611.084 220.196 623.037] /A << /S /GoTo /D <6571756174696F6E2E322E312E342E36> >> >>}
\MusicPDFWriteObject{357}{<< /Type /Annot /Subtype /Link /Border [0 0 0] /H /I /C [1 0 0] /Rect [444.963 611.084 475.176 623.037] /A << /S /GoTo /D <6571756174696F6E2E322E312E342E31> >> >>}
\MusicPDFWriteObject{358}{<< /Type /Annot /Subtype /Link /Border [0 0 0] /H /I /C [1 0 0] /Rect [430.598 570.939 460.811 583.254] /A << /S /GoTo /D <74686D2E322E312E312E36> >> >>}
\MusicPDFWriteObject{359}{<< /Type /Annot /Subtype /Link /Border [0 0 0] /H /I /C [0 1 0] /Rect [107.44 560.639 126.545 569.507] /A << /S /GoTo /D <636974652E4E42> >> >>}
\MusicPDFWriteObject{360}{<< /Type /Annot /Subtype /Link /Border [0 0 0] /H /I /C [1 0 0] /Rect [429.084 558.149 459.297 570.102] /A << /S /GoTo /D <6571756174696F6E2E322E312E342E31> >> >>}
\MusicPDFWriteObject{361}{<< /Type /Annot /Subtype /Link /Border [0 0 0] /H /I /C [1 0 0] /Rect [300.299 518.365 330.512 530.317] /A << /S /GoTo /D <6571756174696F6E2E322E312E342E31> >> >>}
\MusicPDFWriteObject{362}{<< /Type /Annot /Subtype /Link /Border [0 0 0] /H /I /C [1 0 0] /Rect [406.59 518.365 443.289 530.317] /A << /S /GoTo /D <74686D2E432E322E332E3133> >> >>}
\MusicPDFWriteObject{363}{<< /Type /Annot /Subtype /Link /Border [0 0 0] /H /I /C [0 1 0] /Rect [457.288 481.053 482.348 489.92] /A << /S /GoTo /D <636974652E53636861> >> >>}
\MusicPDFWriteObject{364}{<< /Type /Annot /Subtype /Link /Border [0 0 0] /H /I /C [0 1 0] /Rect [111.973 467.902 133.498 476.77] /A << /S /GoTo /D <636974652E4D56> >> >>}
\MusicPDFWriteObject{365}{<< /Type /Annot /Subtype /Link /Border [0 0 0] /H /I /C [0 1 0] /Rect [364.581 467.902 389.641 476.77] /A << /S /GoTo /D <636974652E53636861> >> >>}
\MusicPDFWriteObject{366}{<< /Type /Annot /Subtype /Link /Border [0 0 0] /H /I /C [1 0 0] /Rect [370.842 439.11 402.66 451.063] /A << /S /GoTo /D <74686D2E432E312E362E31> >> >>}
\MusicPDFWriteObject{367}{<< /Type /Annot /Subtype /Link /Border [0 0 0] /H /I /C [0 1 0] /Rect [87.604 428.45 119.228 437.317] /A << /S /GoTo /D <636974652E47726F2D62> >> >>}
\MusicPDFWriteObject{368}{<< /Type /Annot /Subtype /Link /Border [0 0 0] /H /I /C [1 0 0] /Rect [326.861 412.809 357.074 424.762] /A << /S /GoTo /D <6571756174696F6E2E322E312E342E31> >> >>}
\MusicPDFWriteObject{369}{<< /Type /Annot /Subtype /Link /Border [0 0 0] /H /I /C [0 1 0] /Rect [288.498 362.696 313.558 371.564] /A << /S /GoTo /D <636974652E53636861> >> >>}
\MusicPDFWriteObject{370}{<< /Type /Annot /Subtype /Link /Border [0 0 0] /H /I /C [1 0 0] /Rect [224.371 320.754 254.584 332.706] /A << /S /GoTo /D <6571756174696F6E2E322E312E342E34> >> >>}
\MusicPDFWriteObject{371}{<< /Type /Annot /Subtype /Link /Border [0 0 0] /H /I /C [1 0 0] /Rect [234.885 182.113 265.662 194.065] /A << /S /GoTo /D <74686D2E322E312E342E31> >> >>}
\MusicPDFWriteObject{372}{<< /Type /Annot /Subtype /Link /Border [0 0 0] /H /I /C [1 0 0] /Rect [357.712 182.113 388.489 194.065] /A << /S /GoTo /D <74686D2E322E312E332E32> >> >>}
\MusicPDFWriteObject{373}{<< /Type /Annot /Subtype /Link /Border [0 0 0] /H /I /C [1 0 0] /Rect [286.451 122.048 316.664 134.001] /A << /S /GoTo /D <6571756174696F6E2E322E312E342E38> >> >>}
\MusicPDFWriteObject{374}{<< /Type /Annot /Subtype /Link /Border [0 0 0] /H /I /C [1 0 0] /Rect [246.361 109.451 279.062 120.241] /A << /S /GoTo /D <74686D2E432E322E352E31> >> >>}
\MusicPDFWriteObject{375}{<< /Type /Annot /Subtype /Link /Border [0 0 0] /H /I /C [1 0 0] /Rect [240.134 96.3 272.975 107.091] /A << /S /GoTo /D <74686D2E442E312E382E31> >> >>}
\MusicPDFWriteObject{376}{<< /Type /Annot /Subtype /Link /Border [0 0 0] /H /I /C [1 0 0] /Rect [238.201 82.596 268.414 94.549] /A << /S /GoTo /D <6571756174696F6E2E322E312E342E37> >> >>}
\MusicPDFWriteObject{377}{<< /Type /Annot /Subtype /Link /Border [0 0 0] /H /I /C [1 0 0] /Rect [448.041 613.219 478.254 625.171] /A << /S /GoTo /D <6571756174696F6E2E322E312E342E35> >> >>}
\MusicPDFWriteObject{378}{<< /Type /Annot /Subtype /Link /Border [0 0 0] /H /I /C [1 0 0] /Rect [291.265 600.068 321.478 612.021] /A << /S /GoTo /D <6571756174696F6E2E322E312E342E34> >> >>}
\MusicPDFWriteObject{379}{<< /Type /Annot /Subtype /Link /Border [0 0 0] /H /I /C [1 0 0] /Rect [259.262 464.484 289.475 476.436] /A << /S /GoTo /D <6571756174696F6E2E322E312E342E33> >> >>}
\MusicPDFWriteObject{380}{<< /Type /Annot /Subtype /Link /Border [0 0 0] /H /I /C [1 0 0] /Rect [181.379 414.905 211.592 426.858] /A << /S /GoTo /D <6571756174696F6E2E322E312E342E32> >> >>}
\MusicPDFWriteObject{381}{<< /Type /Annot /Subtype /Link /Border [0 0 0] /H /I /C [0 1 0] /Rect [524.054 318.282 543.159 327.15] /A << /S /GoTo /D <636974652E4E42> >> >>}
\MusicPDFWriteObject{382}{<< /Type /Annot /Subtype /Link /Border [0 0 0] /H /I /C [1 0 0] /Rect [373.209 223.869 402.858 235.574] /A << /S /GoTo /D <74686D2E322E312E342E32> >> >>}
\MusicPDFWriteObject{383}{<< /Type /Annot /Subtype /Link /Border [0 0 0] /H /I /C [1 0 0] /Rect [317.811 658.332 350.238 670.284] /A << /S /GoTo /D <74686D2E432E322E312E35> >> >>}
\MusicPDFWriteObject{384}{<< /Type /Annot /Subtype /Link /Border [0 0 0] /H /I /C [0 1 0] /Rect [325.404 568.766 342.876 577.634] /A << /S /GoTo /D <636974652E4346> >> >>}
\MusicPDFWriteObject{385}{<< /Type /Annot /Subtype /Link /Border [0 0 0] /H /I /C [1 0 0] /Rect [343.319 493 373.532 504.952] /A << /S /GoTo /D <6571756174696F6E2E322E312E342E35> >> >>}
\MusicPDFWriteObject{386}{<< /Type /Annot /Subtype /Link /Border [0 0 0] /H /I /C [1 0 0] /Rect [433.639 419.722 464.417 431.677] /A << /S /GoTo /D <74686D2E322E312E342E31> >> >>}
\MusicPDFWriteObject{387}{<< /Type /Annot /Subtype /Link /Border [0 0 0] /H /I /C [0 1 0] /Rect [518.099 422.213 543.159 431.081] /A << /S /GoTo /D <636974652E53636861> >> >>}
\MusicPDFWriteObject{388}{<< /Type /Annot /Subtype /Link /Border [0 0 0] /H /I /C [0 1 0] /Rect [126.288 235.845 151.348 244.713] /A << /S /GoTo /D <636974652E53636861> >> >>}
\MusicPDFWriteObject{389}{<< /Type /Annot /Subtype /Link /Border [0 0 0] /H /I /C [1 0 0] /Rect [309.4 220.204 341.827 232.157] /A << /S /GoTo /D <6571756174696F6E2E432E322E312E34> >> >>}
\MusicPDFWriteObject{390}{<< /Type /Annot /Subtype /Link /Border [0 0 0] /H /I /C [1 0 0] /Rect [273.447 672.036 280.389 682.827] /A << /S /GoTo /D <636861707465722E32> >> >>}
\MusicPDFWriteObject{391}{<< /Type /Annot /Subtype /Link /Border [0 0 0] /H /I /C [1 0 0] /Rect [338.784 476.337 368.997 488.289] /A << /S /GoTo /D <6571756174696F6E2E322E322E312E32> >> >>}
\MusicPDFWriteObject{392}{<< /Type /Annot /Subtype /Link /Border [0 0 0] /H /I /C [0 1 0] /Rect [314.558 426.224 343.596 435.091] /A << /S /GoTo /D <636974652E5363686D> >> >>}
\MusicPDFWriteObject{393}{<< /Type /Annot /Subtype /Link /Border [0 0 0] /H /I /C [0 1 0] /Rect [514.122 398.051 543.159 406.919] /A << /S /GoTo /D <636974652E5363686D> >> >>}
\MusicPDFWriteObject{394}{<< /Type /Annot /Subtype /Link /Border [0 0 0] /H /I /C [0 1 0] /Rect [277.733 382.963 342.226 393.768] /A << /S /GoTo /D <636974652E446176656E706F72742D31> >> >>}
\MusicPDFWriteObject{395}{<< /Type /Annot /Subtype /Link /Border [0 0 0] /H /I /C [0 1 0] /Rect [156.772 98.237 185.81 107.105] /A << /S /GoTo /D <636974652E5363686D> >> >>}
\MusicPDFWriteObject{396}{<< /Type /Annot /Subtype /Link /Border [0 0 0] /H /I /C [0 1 0] /Rect [305.242 98.237 327.535 107.105] /A << /S /GoTo /D <636974652E47616F> >> >>}
\MusicPDFWriteObject{397}{<< /Type /Annot /Subtype /Link /Border [0 0 0] /H /I /C [1 0 0] /Rect [161.75 598.367 191.963 610.32] /A << /S /GoTo /D <6571756174696F6E2E322E322E312E32> >> >>}
\MusicPDFWriteObject{398}{<< /Type /Annot /Subtype /Link /Border [0 0 0] /H /I /C [1 0 0] /Rect [248.381 260.1 278.2 273.324] /A << /S /GoTo /D <74686D2E322E322E312E31> >> >>}
\MusicPDFWriteObject{399}{<< /Type /Annot /Subtype /Link /Border [0 0 0] /H /I /C [1 0 0] /Rect [178.671 646.149 208.884 658.102] /A << /S /GoTo /D <6571756174696F6E2E322E322E312E33> >> >>}
\MusicPDFWriteObject{400}{<< /Type /Annot /Subtype /Link /Border [0 0 0] /H /I /C [0 1 0] /Rect [85.391 479.686 103.361 488.553] /A << /S /GoTo /D <636974652E5256> >> >>}
\MusicPDFWriteObject{401}{<< /Type /Annot /Subtype /Link /Border [0 0 0] /H /I /C [1 0 0] /Rect [241.063 379.48 278.746 391.432] /A << /S /GoTo /D <6571756174696F6E2E412E322E332E3134> >> >>}
\MusicPDFWriteObject{402}{<< /Type /Annot /Subtype /Link /Border [0 0 0] /H /I /C [1 0 0] /Rect [142.557 326.516 180.954 338.831] /A << /S /GoTo /D <74686D2E412E322E332E3131> >> >>}
\MusicPDFWriteObject{403}{<< /Type /Annot /Subtype /Link /Border [0 0 0] /H /I /C [1 0 0] /Rect [472.374 326.516 510.771 338.831] /A << /S /GoTo /D <74686D2E412E322E332E3130> >> >>}
\MusicPDFWriteObject{404}{<< /Type /Annot /Subtype /Link /Border [0 0 0] /H /I /C [1 0 0] /Rect [419.969 564.134 449.617 576.089] /A << /S /GoTo /D <74686D2E322E322E322E31> >> >>}
\MusicPDFWriteObject{405}{<< /Type /Annot /Subtype /Link /Border [0 0 0] /H /I /C [1 0 0] /Rect [360.249 537.833 390.462 549.786] /A << /S /GoTo /D <6571756174696F6E2E322E322E322E33> >> >>}
\MusicPDFWriteObject{406}{<< /Type /Annot /Subtype /Link /Border [0 0 0] /H /I /C [1 0 0] /Rect [501.638 510.364 539.321 522.317] /A << /S /GoTo /D <6571756174696F6E2E412E322E332E3134> >> >>}
\MusicPDFWriteObject{407}{<< /Type /Annot /Subtype /Link /Border [0 0 0] /H /I /C [1 0 0] /Rect [260.435 497.213 293.138 509.166] /A << /S /GoTo /D <6571756174696F6E2E412E322E312E32> >> >>}
\MusicPDFWriteObject{408}{<< /Type /Annot /Subtype /Link /Border [0 0 0] /H /I /C [1 0 0] /Rect [323.418 484.063 356.121 496.015] /A << /S /GoTo /D <6571756174696F6E2E412E322E322E39> >> >>}
\MusicPDFWriteObject{409}{<< /Type /Annot /Subtype /Link /Border [0 0 0] /H /I /C [1 0 0] /Rect [209.976 313.38 240.189 325.333] /A << /S /GoTo /D <6571756174696F6E2E322E322E322E34> >> >>}
\MusicPDFWriteObject{410}{<< /Type /Annot /Subtype /Link /Border [0 0 0] /H /I /C [1 0 0] /Rect [386.513 313.379 416.726 325.335] /A << /S /GoTo /D <74686D2E322E312E332E32> >> >>}
\MusicPDFWriteObject{411}{<< /Type /Annot /Subtype /Link /Border [0 0 0] /H /I /C [1 0 0] /Rect [457.918 287.079 488.131 299.031] /A << /S /GoTo /D <6571756174696F6E2E322E322E322E34> >> >>}
\MusicPDFWriteObject{412}{<< /Type /Annot /Subtype /Link /Border [0 0 0] /H /I /C [1 0 0] /Rect [403.276 260.777 440.959 272.73] /A << /S /GoTo /D <6571756174696F6E2E412E322E332E3134> >> >>}
\MusicPDFWriteObject{413}{<< /Type /Annot /Subtype /Link /Border [0 0 0] /H /I /C [1 0 0] /Rect [214.737 221.325 247.44 233.277] /A << /S /GoTo /D <6571756174696F6E2E412E322E312E32> >> >>}
\MusicPDFWriteObject{414}{<< /Type /Annot /Subtype /Link /Border [0 0 0] /H /I /C [1 0 0] /Rect [70.625 208.173 108.309 220.129] /A << /S /GoTo /D <74686D2E412E322E332E3130> >> >>}
\MusicPDFWriteObject{415}{<< /Type /Annot /Subtype /Link /Border [0 0 0] /H /I /C [1 0 0] /Rect [330.776 195.023 360.989 206.976] /A << /S /GoTo /D <6571756174696F6E2E322E322E322E31> >> >>}
\MusicPDFWriteObject{416}{<< /Type /Annot /Subtype /Link /Border [0 0 0] /H /I /C [1 0 0] /Rect [318.765 635.055 351.468 647.008] /A << /S /GoTo /D <6571756174696F6E2E412E322E312E32> >> >>}
\MusicPDFWriteObject{417}{<< /Type /Annot /Subtype /Link /Border [0 0 0] /H /I /C [1 0 0] /Rect [141.739 414.813 179.422 426.765] /A << /S /GoTo /D <6571756174696F6E2E412E322E332E3136> >> >>}
\MusicPDFWriteObject{418}{<< /Type /Annot /Subtype /Link /Border [0 0 0] /H /I /C [1 0 0] /Rect [328.861 414.813 366.544 426.765] /A << /S /GoTo /D <6571756174696F6E2E412E322E332E3134> >> >>}
\MusicPDFWriteObject{419}{<< /Type /Annot /Subtype /Link /Border [0 0 0] /H /I /C [1 0 0] /Rect [315.975 562.643 351.168 574.596] /A << /S /GoTo /D <6571756174696F6E2E322E322E322E3131> >> >>}
\MusicPDFWriteObject{420}{<< /Type /Annot /Subtype /Link /Border [0 0 0] /H /I /C [1 0 0] /Rect [370.697 549.492 405.89 561.445] /A << /S /GoTo /D <6571756174696F6E2E322E322E322E3131> >> >>}
\MusicPDFWriteObject{421}{<< /Type /Annot /Subtype /Link /Border [0 0 0] /H /I /C [1 0 0] /Rect [281.841 508.71 317.035 520.663] /A << /S /GoTo /D <6571756174696F6E2E322E322E322E3131> >> >>}
\MusicPDFWriteObject{422}{<< /Type /Annot /Subtype /Link /Border [0 0 0] /H /I /C [1 0 0] /Rect [418.924 495.559 454.117 507.512] /A << /S /GoTo /D <6571756174696F6E2E322E322E322E3131> >> >>}
\MusicPDFWriteObject{423}{<< /Type /Annot /Subtype /Link /Border [0 0 0] /H /I /C [1 0 0] /Rect [352.467 469.258 387.66 481.21] /A << /S /GoTo /D <6571756174696F6E2E322E322E322E3131> >> >>}
\MusicPDFWriteObject{424}{<< /Type /Annot /Subtype /Link /Border [0 0 0] /H /I /C [1 0 0] /Rect [355.186 410.673 388.503 423.266] /A << /S /GoTo /D <74686D2E412E322E332E39> >> >>}
\MusicPDFWriteObject{425}{<< /Type /Annot /Subtype /Link /Border [0 0 0] /H /I /C [1 0 0] /Rect [329.541 397.285 367.224 409.238] /A << /S /GoTo /D <6571756174696F6E2E412E322E332E3133> >> >>}
\MusicPDFWriteObject{426}{<< /Type /Annot /Subtype /Link /Border [0 0 0] /H /I /C [1 0 0] /Rect [216.809 357.833 249.512 369.785] /A << /S /GoTo /D <6571756174696F6E2E412E322E312E31> >> >>}
\MusicPDFWriteObject{427}{<< /Type /Annot /Subtype /Link /Border [0 0 0] /H /I /C [1 0 0] /Rect [237.009 331.531 269.712 343.484] /A << /S /GoTo /D <6571756174696F6E2E412E322E322E39> >> >>}
\MusicPDFWriteObject{428}{<< /Type /Annot /Subtype /Link /Border [0 0 0] /H /I /C [1 0 0] /Rect [352.029 150.096 387.222 162.049] /A << /S /GoTo /D <6571756174696F6E2E322E322E322E3131> >> >>}
\MusicPDFWriteObject{429}{<< /Type /Annot /Subtype /Link /Border [0 0 0] /H /I /C [1 0 0] /Rect [332 612.003 367.193 623.956] /A << /S /GoTo /D <6571756174696F6E2E322E322E322E3133> >> >>}
\MusicPDFWriteObject{430}{<< /Type /Annot /Subtype /Link /Border [0 0 0] /H /I /C [1 0 0] /Rect [124.623 524.112 154.836 536.425] /A << /S /GoTo /D <74686D2E322E322E312E33> >> >>}
\MusicPDFWriteObject{431}{<< /Type /Annot /Subtype /Link /Border [0 0 0] /H /I /C [1 0 0] /Rect [207.071 524.472 242.264 536.425] /A << /S /GoTo /D <6571756174696F6E2E322E322E322E3134> >> >>}
\MusicPDFWriteObject{432}{<< /Type /Annot /Subtype /Link /Border [0 0 0] /H /I /C [1 0 0] /Rect [171.47 511.322 206.663 523.274] /A << /S /GoTo /D <6571756174696F6E2E322E322E322E3131> >> >>}
\MusicPDFWriteObject{433}{<< /Type /Annot /Subtype /Link /Border [0 0 0] /H /I /C [1 0 0] /Rect [321.882 511.322 352.095 523.274] /A << /S /GoTo /D <6571756174696F6E2E322E322E322E31> >> >>}
\MusicPDFWriteObject{434}{<< /Type /Annot /Subtype /Link /Border [0 0 0] /H /I /C [1 0 0] /Rect [347.118 212.899 377.077 224.852] /A << /S /GoTo /D <74686D2E322E322E322E31> >> >>}
\MusicPDFWriteObject{435}{<< /Type /Annot /Subtype /Link /Border [0 0 0] /H /I /C [0 1 0] /Rect [525.767 215.389 543.239 224.257] /A << /S /GoTo /D <636974652E4346> >> >>}
\MusicPDFWriteObject{436}{<< /Type /Annot /Subtype /Link /Border [0 0 0] /H /I /C [1 0 0] /Rect [140.649 173.446 171.201 185.402] /A << /S /GoTo /D <74686D2E322E322E322E31> >> >>}
\MusicPDFWriteObject{437}{<< /Type /Annot /Subtype /Link /Border [0 0 0] /H /I /C [1 0 0] /Rect [197.653 638.633 227.867 650.586] /A << /S /GoTo /D <6571756174696F6E2E322E322E332E31> >> >>}
\MusicPDFWriteObject{438}{<< /Type /Annot /Subtype /Link /Border [0 0 0] /H /I /C [0 1 0] /Rect [215.012 432.045 232.484 440.913] /A << /S /GoTo /D <636974652E4346> >> >>}
\MusicPDFWriteObject{439}{<< /Type /Annot /Subtype /Link /Border [0 0 0] /H /I /C [1 0 0] /Rect [391.484 319.895 424.943 330.686] /A << /S /GoTo /D <74686D2E442E312E322E31> >> >>}
\MusicPDFWriteObject{440}{<< /Type /Annot /Subtype /Link /Border [0 0 0] /H /I /C [1 0 0] /Rect [226.443 293.594 259.146 304.384] /A << /S /GoTo /D <74686D2E412E332E342E32> >> >>}
\MusicPDFWriteObject{441}{<< /Type /Annot /Subtype /Link /Border [0 0 0] /H /I /C [1 0 0] /Rect [332.298 293.594 365.001 304.384] /A << /S /GoTo /D <74686D2E412E332E322E35> >> >>}
\MusicPDFWriteObject{442}{<< /Type /Annot /Subtype /Link /Border [0 0 0] /H /I /C [0 1 0] /Rect [306.917 282.38 324.389 291.248] /A << /S /GoTo /D <636974652E4346> >> >>}
\MusicPDFWriteObject{443}{<< /Type /Annot /Subtype /Link /Border [0 0 0] /H /I /C [0 1 0] /Rect [409.268 282.38 425.26 291.248] /A << /S /GoTo /D <636974652E436F> >> >>}
\MusicPDFWriteObject{444}{<< /Type /Annot /Subtype /Link /Border [0 0 0] /H /I /C [1 0 0] /Rect [312.44 266.739 342.173 278.692] /A << /S /GoTo /D <74686D2E322E322E332E32> >> >>}
\MusicPDFWriteObject{445}{<< /Type /Annot /Subtype /Link /Border [0 0 0] /H /I /C [0 1 0] /Rect [525.789 269.229 543.261 278.097] /A << /S /GoTo /D <636974652E4346> >> >>}
\MusicPDFWriteObject{446}{<< /Type /Annot /Subtype /Link /Border [0 0 0] /H /I /C [0 1 0] /Rect [269.414 236.38 286.887 245.248] /A << /S /GoTo /D <636974652E4346> >> >>}
\MusicPDFWriteObject{447}{<< /Type /Annot /Subtype /Link /Border [0 0 0] /H /I /C [1 0 0] /Rect [431.578 233.89 462.243 245.843] /A << /S /GoTo /D <74686D2E322E322E322E31> >> >>}
\MusicPDFWriteObject{448}{<< /Type /Annot /Subtype /Link /Border [0 0 0] /H /I /C [0 1 0] /Rect [226.853 223.229 244.326 232.097] /A << /S /GoTo /D <636974652E4346> >> >>}
\MusicPDFWriteObject{449}{<< /Type /Annot /Subtype /Link /Border [0 0 0] /H /I /C [1 0 0] /Rect [511.675 184.152 541.888 196.104] /A << /S /GoTo /D <6571756174696F6E2E322E322E332E33> >> >>}
\MusicPDFWriteObject{450}{<< /Type /Annot /Subtype /Link /Border [0 0 0] /H /I /C [1 0 0] /Rect [113.031 134.573 143.244 146.526] /A << /S /GoTo /D <6571756174696F6E2E322E322E332E32> >> >>}
\MusicPDFWriteObject{451}{<< /Type /Annot /Subtype /Link /Border [0 0 0] /H /I /C [1 0 0] /Rect [134.621 509.276 164.835 521.229] /A << /S /GoTo /D <6571756174696F6E2E322E322E332E34> >> >>}
\MusicPDFWriteObject{452}{<< /Type /Annot /Subtype /Link /Border [0 0 0] /H /I /C [1 0 0] /Rect [176.874 509.276 207.088 521.229] /A << /S /GoTo /D <6571756174696F6E2E322E322E332E35> >> >>}
\MusicPDFWriteObject{453}{<< /Type /Annot /Subtype /Link /Border [0 0 0] /H /I /C [1 0 0] /Rect [235.785 509.276 265.998 521.229] /A << /S /GoTo /D <6571756174696F6E2E322E322E332E36> >> >>}
\MusicPDFWriteObject{454}{<< /Type /Annot /Subtype /Link /Border [0 0 0] /H /I /C [1 0 0] /Rect [178.626 430.777 208.36 444.586] /A << /S /GoTo /D <74686D2E322E322E322E31> >> >>}
\MusicPDFWriteObject{455}{<< /Type /Annot /Subtype /Link /Border [0 0 0] /H /I /C [1 0 0] /Rect [259.516 432.632 289.729 444.584] /A << /S /GoTo /D <6571756174696F6E2E322E322E322E31> >> >>}
\MusicPDFWriteObject{456}{<< /Type /Annot /Subtype /Link /Border [0 0 0] /H /I /C [1 0 0] /Rect [146.65 406.883 176.863 417.674] /A << /S /GoTo /D <74686D2E322E322E332E32> >> >>}
\MusicPDFWriteObject{457}{<< /Type /Annot /Subtype /Link /Border [0 0 0] /H /I /C [1 0 0] /Rect [346.489 369.305 376.702 381.258] /A << /S /GoTo /D <6571756174696F6E2E322E322E332E35> >> >>}
\MusicPDFWriteObject{458}{<< /Type /Annot /Subtype /Link /Border [0 0 0] /H /I /C [1 0 0] /Rect [423.056 369.305 453.269 381.258] /A << /S /GoTo /D <6571756174696F6E2E322E322E332E34> >> >>}
\MusicPDFWriteObject{459}{<< /Type /Annot /Subtype /Link /Border [0 0 0] /H /I /C [1 0 0] /Rect [222.587 316.618 252.8 328.571] /A << /S /GoTo /D <6571756174696F6E2E322E322E332E34> >> >>}
\MusicPDFWriteObject{460}{<< /Type /Annot /Subtype /Link /Border [0 0 0] /H /I /C [1 0 0] /Rect [154.759 290.317 184.972 302.27] /A << /S /GoTo /D <6571756174696F6E2E322E322E332E34> >> >>}
\MusicPDFWriteObject{461}{<< /Type /Annot /Subtype /Link /Border [0 0 0] /H /I /C [1 0 0] /Rect [510.024 277.166 539.955 289.119] /A << /S /GoTo /D <74686D2E322E322E322E31> >> >>}
\MusicPDFWriteObject{462}{<< /Type /Annot /Subtype /Link /Border [0 0 0] /H /I /C [1 0 0] /Rect [236.342 250.864 274.026 262.817] /A << /S /GoTo /D <6571756174696F6E2E412E322E322E3131> >> >>}
\MusicPDFWriteObject{463}{<< /Type /Annot /Subtype /Link /Border [0 0 0] /H /I /C [1 0 0] /Rect [149.363 224.563 179.576 236.516] /A << /S /GoTo /D <74686D2E322E322E322E31> >> >>}
\MusicPDFWriteObject{464}{<< /Type /Annot /Subtype /Link /Border [0 0 0] /H /I /C [1 0 0] /Rect [501.602 684.633 539.285 696.586] /A << /S /GoTo /D <6571756174696F6E2E412E322E322E3131> >> >>}
\MusicPDFWriteObject{465}{<< /Type /Annot /Subtype /Link /Border [0 0 0] /H /I /C [1 0 0] /Rect [418.357 671.482 456.04 683.435] /A << /S /GoTo /D <6571756174696F6E2E412E322E332E3134> >> >>}
\MusicPDFWriteObject{466}{<< /Type /Annot /Subtype /Link /Border [0 0 0] /H /I /C [1 0 0] /Rect [158.167 480.271 193.36 492.223] /A << /S /GoTo /D <6571756174696F6E2E322E322E332E3130> >> >>}
\MusicPDFWriteObject{467}{<< /Type /Annot /Subtype /Link /Border [0 0 0] /H /I /C [1 0 0] /Rect [514.604 364.622 544.845 381.018] /A << /S /GoTo /D <74686D2E322E322E312E33> >> >>}
\MusicPDFWriteObject{468}{<< /Type /Annot /Subtype /Link /Border [0 0 0] /H /I /C [1 0 0] /Rect [180.715 262.873 211.07 277.424] /A << /S /GoTo /D <74686D2E322E312E322E34> >> >>}
\MusicPDFWriteObject{469}{<< /Type /Annot /Subtype /Link /Border [0 0 0] /H /I /C [1 0 0] /Rect [285.752 262.873 316.106 277.424] /A << /S /GoTo /D <74686D2E322E312E322E38> >> >>}
\MusicPDFWriteObject{470}{<< /Type /Annot /Subtype /Link /Border [0 0 0] /H /I /C [1 0 0] /Rect [219.736 249.319 250.344 263.869] /A << /S /GoTo /D <74686D2E322E312E332E32> >> >>}
\MusicPDFWriteObject{471}{<< /Type /Annot /Subtype /Link /Border [0 0 0] /H /I /C [1 0 0] /Rect [278.002 235.764 310.429 250.315] /A << /S /GoTo /D <74686D2E432E322E352E31> >> >>}
\MusicPDFWriteObject{472}{<< /Type /Annot /Subtype /Link /Border [0 0 0] /H /I /C [1 0 0] /Rect [318.319 222.691 347.967 235.006] /A << /S /GoTo /D <74686D2E322E312E332E32> >> >>}
\MusicPDFWriteObject{473}{<< /Type /Annot /Subtype /Link /Border [0 0 0] /H /I /C [1 0 0] /Rect [410.501 209.136 447.199 223.687] /A << /S /GoTo /D <74686D2E432E322E332E3133> >> >>}
\MusicPDFWriteObject{474}{<< /Type /Annot /Subtype /Link /Border [0 0 0] /H /I /C [0 1 0] /Rect [270.976 185.762 296.036 194.63] /A << /S /GoTo /D <636974652E53636861> >> >>}
\MusicPDFWriteObject{475}{<< /Type /Annot /Subtype /Link /Border [0 0 0] /H /I /C [0 1 0] /Rect [291.403 135.858 309.602 144.725] /A << /S /GoTo /D <636974652E4A6172> >> >>}
\MusicPDFWriteObject{476}{<< /Type /Annot /Subtype /Link /Border [0 0 0] /H /I /C [0 1 0] /Rect [499.876 673.973 518.074 682.84] /A << /S /GoTo /D <636974652E4A6172> >> >>}
\MusicPDFWriteObject{477}{<< /Type /Annot /Subtype /Link /Border [0 0 0] /H /I /C [0 1 0] /Rect [236.902 660.822 256.007 669.69] /A << /S /GoTo /D <636974652E4E42> >> >>}
\MusicPDFWriteObject{478}{<< /Type /Annot /Subtype /Link /Border [0 0 0] /H /I /C [1 0 0] /Rect [266.074 500.523 296.287 512.475] /A << /S /GoTo /D <6571756174696F6E2E322E332E312E33> >> >>}
\MusicPDFWriteObject{479}{<< /Type /Annot /Subtype /Link /Border [0 0 0] /H /I /C [1 0 0] /Rect [126.394 447.92 156.607 459.872] /A << /S /GoTo /D <6571756174696F6E2E322E322E332E38> >> >>}
\MusicPDFWriteObject{480}{<< /Type /Annot /Subtype /Link /Border [0 0 0] /H /I /C [1 0 0] /Rect [200.568 434.769 230.781 446.722] /A << /S /GoTo /D <6571756174696F6E2E322E332E312E33> >> >>}
\MusicPDFWriteObject{481}{<< /Type /Annot /Subtype /Link /Border [0 0 0] /H /I /C [0 1 0] /Rect [377.613 410.957 395.812 419.825] /A << /S /GoTo /D <636974652E4A6172> >> >>}
\MusicPDFWriteObject{482}{<< /Type /Annot /Subtype /Link /Border [0 0 0] /H /I /C [0 1 0] /Rect [109.094 397.807 128.198 406.674] /A << /S /GoTo /D <636974652E4E42> >> >>}
\MusicPDFWriteObject{483}{<< /Type /Annot /Subtype /Link /Border [0 0 0] /H /I /C [1 0 0] /Rect [264.153 395.317 294.367 407.269] /A << /S /GoTo /D <6571756174696F6E2E322E322E332E34> >> >>}
\MusicPDFWriteObject{484}{<< /Type /Annot /Subtype /Link /Border [0 0 0] /H /I /C [1 0 0] /Rect [147.744 355.864 177.957 367.817] /A << /S /GoTo /D <6571756174696F6E2E322E322E332E35> >> >>}
\MusicPDFWriteObject{485}{<< /Type /Annot /Subtype /Link /Border [0 0 0] /H /I /C [1 0 0] /Rect [207.226 355.864 237.439 367.817] /A << /S /GoTo /D <6571756174696F6E2E322E322E332E36> >> >>}
\MusicPDFWriteObject{486}{<< /Type /Annot /Subtype /Link /Border [0 0 0] /H /I /C [0 1 0] /Rect [435.906 150.64 460.966 159.508] /A << /S /GoTo /D <636974652E53636861> >> >>}
\MusicPDFWriteObject{487}{<< /Type /Annot /Subtype /Link /Border [0 0 0] /H /I /C [1 0 0] /Rect [219.553 684.633 254.746 696.586] /A << /S /GoTo /D <6571756174696F6E2E322E322E332E3134> >> >>}
\MusicPDFWriteObject{488}{<< /Type /Annot /Subtype /Link /Border [0 0 0] /H /I /C [1 0 0] /Rect [405.595 571.144 443.278 583.096] /A << /S /GoTo /D <6571756174696F6E2E412E322E332E3134> >> >>}
\MusicPDFWriteObject{489}{<< /Type /Annot /Subtype /Link /Border [0 0 0] /H /I /C [1 0 0] /Rect [501.696 570.89 531.345 583.736] /A << /S /GoTo /D <74686D2E322E322E322E31> >> >>}
\MusicPDFWriteObject{490}{<< /Type /Annot /Subtype /Link /Border [0 0 0] /H /I /C [1 0 0] /Rect [223.691 557.993 258.884 569.946] /A << /S /GoTo /D <6571756174696F6E2E322E322E332E3135> >> >>}
\MusicPDFWriteObject{491}{<< /Type /Annot /Subtype /Link /Border [0 0 0] /H /I /C [1 0 0] /Rect [124.623 378.529 154.836 390.485] /A << /S /GoTo /D <74686D2E322E322E312E33> >> >>}
\MusicPDFWriteObject{492}{<< /Type /Annot /Subtype /Link /Border [0 0 0] /H /I /C [1 0 0] /Rect [260.334 312.316 290.604 324.631] /A << /S /GoTo /D <74686D2E322E312E332E32> >> >>}
\MusicPDFWriteObject{493}{<< /Type /Annot /Subtype /Link /Border [0 0 0] /H /I /C [1 0 0] /Rect [506.695 233.446 541.888 245.398] /A << /S /GoTo /D <6571756174696F6E2E322E322E332E3134> >> >>}
\MusicPDFWriteObject{494}{<< /Type /Annot /Subtype /Link /Border [0 0 0] /H /I /C [1 0 0] /Rect [184.089 547.959 214.302 559.912] /A << /S /GoTo /D <6571756174696F6E2E322E322E342E31> >> >>}
\MusicPDFWriteObject{495}{<< /Type /Annot /Subtype /Link /Border [0 0 0] /H /I /C [1 0 0] /Rect [329.488 364.247 359.701 376.2] /A << /S /GoTo /D <6571756174696F6E2E322E312E332E32> >> >>}
\MusicPDFWriteObject{496}{<< /Type /Annot /Subtype /Link /Border [0 0 0] /H /I /C [1 0 0] /Rect [125.958 351.096 156.735 363.051] /A << /S /GoTo /D <74686D2E322E312E332E32> >> >>}
\MusicPDFWriteObject{497}{<< /Type /Annot /Subtype /Link /Border [0 0 0] /H /I /C [1 0 0] /Rect [226.077 271.77 256.291 283.722] /A << /S /GoTo /D <6571756174696F6E2E322E322E342E32> >> >>}
\MusicPDFWriteObject{498}{<< /Type /Annot /Subtype /Link /Border [0 0 0] /H /I /C [1 0 0] /Rect [141.512 205.655 172.036 217.417] /A << /S /GoTo /D <74686D2E322E322E342E31> >> >>}
\MusicPDFWriteObject{499}{<< /Type /Annot /Subtype /Link /Border [0 0 0] /H /I /C [1 0 0] /Rect [453.473 601.006 483.686 612.958] /A << /S /GoTo /D <6571756174696F6E2E322E322E342E34> >> >>}
\MusicPDFWriteObject{500}{<< /Type /Annot /Subtype /Link /Border [0 0 0] /H /I /C [1 0 0] /Rect [198.35 542.29 228.563 554.243] /A << /S /GoTo /D <6571756174696F6E2E322E322E342E34> >> >>}
\MusicPDFWriteObject{501}{<< /Type /Annot /Subtype /Link /Border [0 0 0] /H /I /C [1 0 0] /Rect [99.714 502.838 129.927 514.791] /A << /S /GoTo /D <6571756174696F6E2E322E312E342E35> >> >>}
\MusicPDFWriteObject{502}{<< /Type /Annot /Subtype /Link /Border [0 0 0] /H /I /C [1 0 0] /Rect [124.566 489.327 154.751 501.642] /A << /S /GoTo /D <74686D2E322E312E342E31> >> >>}
\MusicPDFWriteObject{503}{<< /Type /Annot /Subtype /Link /Border [0 0 0] /H /I /C [0 1 0] /Rect [305.693 479.027 330.753 487.894] /A << /S /GoTo /D <636974652E53636861> >> >>}
\MusicPDFWriteObject{504}{<< /Type /Annot /Subtype /Link /Border [0 0 0] /H /I /C [1 0 0] /Rect [98.422 413.808 128.635 425.76] /A << /S /GoTo /D <6571756174696F6E2E322E312E332E32> >> >>}
\MusicPDFWriteObject{505}{<< /Type /Annot /Subtype /Link /Border [0 0 0] /H /I /C [1 0 0] /Rect [198.799 413.808 229.012 425.76] /A << /S /GoTo /D <74686D2E322E312E332E32> >> >>}
\MusicPDFWriteObject{506}{<< /Type /Annot /Subtype /Link /Border [0 0 0] /H /I /C [1 0 0] /Rect [162.77 336.218 192.983 348.17] /A << /S /GoTo /D <6571756174696F6E2E322E322E342E36> >> >>}
\MusicPDFWriteObject{507}{<< /Type /Annot /Subtype /Link /Border [0 0 0] /H /I /C [1 0 0] /Rect [472.399 322.706 502.612 335.021] /A << /S /GoTo /D <74686D2E322E322E342E31> >> >>}
\MusicPDFWriteObject{508}{<< /Type /Annot /Subtype /Link /Border [0 0 0] /H /I /C [1 0 0] /Rect [124.623 187.801 154.836 199.755] /A << /S /GoTo /D <74686D2E322E312E332E32> >> >>}
\MusicPDFWriteObject{509}{<< /Type /Annot /Subtype /Link /Border [0 0 0] /H /I /C [1 0 0] /Rect [307.01 122.048 344.693 134.001] /A << /S /GoTo /D <74686D2E412E322E332E3131> >> >>}
\MusicPDFWriteObject{510}{<< /Type /Annot /Subtype /Link /Border [0 0 0] /H /I /C [1 0 0] /Rect [126.149 108.537 156.927 120.852] /A << /S /GoTo /D <74686D2E322E322E322E31> >> >>}
\MusicPDFWriteObject{511}{<< /Type /Annot /Subtype /Link /Border [0 0 0] /H /I /C [1 0 0] /Rect [475.637 82.596 513.32 94.549] /A << /S /GoTo /D <6571756174696F6E2E412E322E332E3134> >> >>}
\MusicPDFWriteObject{512}{<< /Type /Annot /Subtype /Link /Border [0 0 0] /H /I /C [1 0 0] /Rect [248.699 621.965 279.477 634.811] /A << /S /GoTo /D <74686D2E322E322E322E33> >> >>}
\MusicPDFWriteObject{513}{<< /Type /Annot /Subtype /Link /Border [0 0 0] /H /I /C [1 0 0] /Rect [213.61 542.982 243.823 554.935] /A << /S /GoTo /D <6571756174696F6E2E322E322E342E38> >> >>}
\MusicPDFWriteObject{514}{<< /Type /Annot /Subtype /Link /Border [0 0 0] /H /I /C [1 0 0] /Rect [445.986 455.817 476.764 468.13] /A << /S /GoTo /D <74686D2E322E322E322E31> >> >>}
\MusicPDFWriteObject{515}{<< /Type /Annot /Subtype /Link /Border [0 0 0] /H /I /C [1 0 0] /Rect [118.983 416.725 149.196 428.678] /A << /S /GoTo /D <6571756174696F6E2E322E322E342E38> >> >>}
\MusicPDFWriteObject{516}{<< /Type /Annot /Subtype /Link /Border [0 0 0] /H /I /C [1 0 0] /Rect [205.305 416.725 235.518 428.678] /A << /S /GoTo /D <6571756174696F6E2E322E322E342E39> >> >>}
\MusicPDFWriteObject{517}{<< /Type /Annot /Subtype /Link /Border [0 0 0] /H /I /C [1 0 0] /Rect [332.742 416.725 367.935 428.678] /A << /S /GoTo /D <6571756174696F6E2E322E322E342E3130> >> >>}
\MusicPDFWriteObject{518}{<< /Type /Annot /Subtype /Link /Border [0 0 0] /H /I /C [1 0 0] /Rect [247.642 338.149 282.835 350.102] /A << /S /GoTo /D <6571756174696F6E2E322E322E342E3131> >> >>}
\MusicPDFWriteObject{519}{<< /Type /Annot /Subtype /Link /Border [0 0 0] /H /I /C [1 0 0] /Rect [190.471 571.965 220.684 583.727] /A << /S /GoTo /D <74686D2E322E322E342E32> >> >>}
\MusicPDFWriteObject{520}{<< /Type /Annot /Subtype /Link /Border [0 0 0] /H /I /C [1 0 0] /Rect [101.05 459.782 136.244 471.735] /A << /S /GoTo /D <6571756174696F6E2E322E322E342E3134> >> >>}
\MusicPDFWriteObject{521}{<< /Type /Annot /Subtype /Link /Border [0 0 0] /H /I /C [1 0 0] /Rect [418.058 459.782 453.251 471.735] /A << /S /GoTo /D <6571756174696F6E2E322E322E342E3132> >> >>}
\MusicPDFWriteObject{522}{<< /Type /Annot /Subtype /Link /Border [0 0 0] /H /I /C [1 0 0] /Rect [146.453 347.315 181.647 359.268] /A << /S /GoTo /D <6571756174696F6E2E322E322E342E3133> >> >>}
\MusicPDFWriteObject{523}{<< /Type /Annot /Subtype /Link /Border [0 0 0] /H /I /C [1 0 0] /Rect [159.843 321.014 195.036 332.966] /A << /S /GoTo /D <6571756174696F6E2E322E322E342E3131> >> >>}
\MusicPDFWriteObject{524}{<< /Type /Annot /Subtype /Link /Border [0 0 0] /H /I /C [1 0 0] /Rect [136.96 266.226 166.609 281.332] /A << /S /GoTo /D <74686D2E322E322E312E33> >> >>}
\MusicPDFWriteObject{525}{<< /Type /Annot /Subtype /Link /Border [0 0 0] /H /I /C [1 0 0] /Rect [169.966 684.633 205.159 696.586] /A << /S /GoTo /D <6571756174696F6E2E322E322E342E3137> >> >>}
\MusicPDFWriteObject{526}{<< /Type /Annot /Subtype /Link /Border [0 0 0] /H /I /C [1 0 0] /Rect [217.476 684.633 252.669 696.586] /A << /S /GoTo /D <6571756174696F6E2E322E322E342E3230> >> >>}
\MusicPDFWriteObject{527}{<< /Type /Annot /Subtype /Link /Border [0 0 0] /H /I /C [1 0 0] /Rect [282.104 684.633 317.297 696.586] /A << /S /GoTo /D <6571756174696F6E2E322E322E342E3139> >> >>}
\MusicPDFWriteObject{528}{<< /Type /Annot /Subtype /Link /Border [0 0 0] /H /I /C [1 0 0] /Rect [110.381 556.447 145.575 568.4] /A << /S /GoTo /D <6571756174696F6E2E322E322E342E3232> >> >>}
\MusicPDFWriteObject{529}{<< /Type /Annot /Subtype /Link /Border [0 0 0] /H /I /C [1 0 0] /Rect [170.299 516.995 205.492 528.948] /A << /S /GoTo /D <6571756174696F6E2E322E322E342E3232> >> >>}
\MusicPDFWriteObject{530}{<< /Type /Annot /Subtype /Link /Border [0 0 0] /H /I /C [1 0 0] /Rect [96.645 503.844 126.858 515.797] /A << /S /GoTo /D <6571756174696F6E2E322E322E342E37> >> >>}
\MusicPDFWriteObject{531}{<< /Type /Annot /Subtype /Link /Border [0 0 0] /H /I /C [1 0 0] /Rect [296.961 503.484 327.456 515.797] /A << /S /GoTo /D <74686D2E322E322E342E31> >> >>}
\MusicPDFWriteObject{532}{<< /Type /Annot /Subtype /Link /Border [0 0 0] /H /I /C [1 0 0] /Rect [352.6 490.333 382.813 502.648] /A << /S /GoTo /D <74686D2E322E322E322E31> >> >>}
\MusicPDFWriteObject{533}{<< /Type /Annot /Subtype /Link /Border [0 0 0] /H /I /C [1 0 0] /Rect [70.755 314.884 100.968 326.837] /A << /S /GoTo /D <74686D2E332E322E322E31> >> >>}
\MusicPDFWriteObject{534}{<< /Type /Annot /Subtype /Link /Border [0 0 0] /H /I /C [1 0 0] /Rect [163.025 314.884 193.238 326.837] /A << /S /GoTo /D <74686D2E332E322E322E34> >> >>}
\MusicPDFWriteObject{535}{<< /Type /Annot /Subtype /Link /Border [0 0 0] /H /I /C [1 0 0] /Rect [354.778 544.084 389.971 556.036] /A << /S /GoTo /D <6571756174696F6E2E322E322E342E3238> >> >>}
\MusicPDFWriteObject{536}{<< /Type /Annot /Subtype /Link /Border [0 0 0] /H /I /C [0 1 0] /Rect [241.821 464.216 266.881 473.084] /A << /S /GoTo /D <636974652E53636861> >> >>}
\MusicPDFWriteObject{537}{<< /Type /Annot /Subtype /Link /Border [0 0 0] /H /I /C [0 1 0] /Rect [409.645 464.216 434.705 473.084] /A << /S /GoTo /D <636974652E53636861> >> >>}
\MusicPDFWriteObject{538}{<< /Type /Annot /Subtype /Link /Border [0 0 0] /H /I /C [1 0 0] /Rect [461.073 349.189 496.266 361.142] /A << /S /GoTo /D <6571756174696F6E2E322E322E342E3239> >> >>}
\MusicPDFWriteObject{539}{<< /Type /Annot /Subtype /Link /Border [0 0 0] /H /I /C [1 0 0] /Rect [395.429 245.199 430.622 257.152] /A << /S /GoTo /D <6571756174696F6E2E322E322E342E3239> >> >>}
\MusicPDFWriteObject{540}{<< /Type /Annot /Subtype /Link /Border [0 0 0] /H /I /C [0 1 0] /Rect [107.335 160.767 125.306 169.635] /A << /S /GoTo /D <636974652E5256> >> >>}
\MusicPDFWriteObject{541}{<< /Type /Annot /Subtype /Link /Border [0 0 0] /H /I /C [0 1 0] /Rect [333.463 98.237 351.434 107.105] /A << /S /GoTo /D <636974652E5256> >> >>}
\MusicPDFWriteObject{542}{<< /Type /Annot /Subtype /Link /Border [0 0 0] /H /I /C [1 0 0] /Rect [290.086 82.595 322.636 94.55] /A << /S /GoTo /D <74686D2E412E332E312E31> >> >>}
\MusicPDFWriteObject{543}{<< /Type /Annot /Subtype /Link /Border [0 0 0] /H /I /C [0 1 0] /Rect [93.426 687.123 111.396 695.991] /A << /S /GoTo /D <636974652E5256> >> >>}
\MusicPDFWriteObject{544}{<< /Type /Annot /Subtype /Link /Border [0 0 0] /H /I /C [1 0 0] /Rect [473.7 533.466 508.893 545.419] /A << /S /GoTo /D <6571756174696F6E2E322E322E342E3333> >> >>}
\MusicPDFWriteObject{545}{<< /Type /Annot /Subtype /Link /Border [0 0 0] /H /I /C [1 0 0] /Rect [122.188 395.818 152.401 409.268] /A << /S /GoTo /D <74686D2E322E322E342E31> >> >>}
\MusicPDFWriteObject{546}{<< /Type /Annot /Subtype /Link /Border [0 0 0] /H /I /C [1 0 0] /Rect [167.926 384.163 203.12 396.115] /A << /S /GoTo /D <6571756174696F6E2E322E322E342E3335> >> >>}
\MusicPDFWriteObject{547}{<< /Type /Annot /Subtype /Link /Border [0 0 0] /H /I /C [1 0 0] /Rect [231.561 384.163 266.754 396.115] /A << /S /GoTo /D <6571756174696F6E2E322E322E342E3332> >> >>}
\MusicPDFWriteObject{548}{<< /Type /Annot /Subtype /Link /Border [0 0 0] /H /I /C [0 1 0] /Rect [314.01 386.653 331.98 395.52] /A << /S /GoTo /D <636974652E5256> >> >>}
\MusicPDFWriteObject{549}{<< /Type /Annot /Subtype /Link /Border [0 0 0] /H /I /C [1 0 0] /Rect [215.555 314.958 250.748 326.911] /A << /S /GoTo /D <6571756174696F6E2E322E322E342E3336> >> >>}
\MusicPDFWriteObject{550}{<< /Type /Annot /Subtype /Link /Border [0 0 0] /H /I /C [1 0 0] /Rect [463.126 314.958 498.319 326.911] /A << /S /GoTo /D <6571756174696F6E2E322E322E342E3336> >> >>}
\MusicPDFWriteObject{551}{<< /Type /Annot /Subtype /Link /Border [0 0 0] /H /I /C [1 0 0] /Rect [128.273 207.406 163.466 219.359] /A << /S /GoTo /D <6571756174696F6E2E322E322E342E3238> >> >>}
\MusicPDFWriteObject{552}{<< /Type /Annot /Subtype /Link /Border [0 0 0] /H /I /C [1 0 0] /Rect [365.951 194.255 401.144 206.208] /A << /S /GoTo /D <6571756174696F6E2E322E322E342E3239> >> >>}
\MusicPDFWriteObject{553}{<< /Type /Annot /Subtype /Link /Border [0 0 0] /H /I /C [1 0 0] /Rect [170.811 181.105 206.004 193.057] /A << /S /GoTo /D <6571756174696F6E2E322E322E342E3238> >> >>}
\MusicPDFWriteObject{554}{<< /Type /Annot /Subtype /Link /Border [0 0 0] /H /I /C [1 0 0] /Rect [230.85 436.229 263.553 448.182] /A << /S /GoTo /D <74686D2E412E332E322E35> >> >>}
\MusicPDFWriteObject{555}{<< /Type /Annot /Subtype /Link /Border [0 0 0] /H /I /C [1 0 0] /Rect [229.294 223.392 264.487 235.345] /A << /S /GoTo /D <6571756174696F6E2E322E322E342E3338> >> >>}
\MusicPDFWriteObject{556}{<< /Type /Annot /Subtype /Link /Border [0 0 0] /H /I /C [1 0 0] /Rect [335.044 210.241 370.237 222.194] /A << /S /GoTo /D <6571756174696F6E2E322E322E342E3337> >> >>}
\MusicPDFWriteObject{557}{<< /Type /Annot /Subtype /Link /Border [0 0 0] /H /I /C [1 0 0] /Rect [259.502 671.482 289.941 683.437] /A << /S /GoTo /D <74686D2E322E322E342E32> >> >>}
\MusicPDFWriteObject{558}{<< /Type /Annot /Subtype /Link /Border [0 0 0] /H /I /C [1 0 0] /Rect [365.063 671.482 400.256 683.435] /A << /S /GoTo /D <6571756174696F6E2E322E322E342E3336> >> >>}
\MusicPDFWriteObject{559}{<< /Type /Annot /Subtype /Link /Border [0 0 0] /H /I /C [1 0 0] /Rect [351.06 658.332 386.253 670.284] /A << /S /GoTo /D <6571756174696F6E2E322E322E342E3336> >> >>}
\MusicPDFWriteObject{560}{<< /Type /Annot /Subtype /Link /Border [0 0 0] /H /I /C [1 0 0] /Rect [86.497 632.03 121.691 643.983] /A << /S /GoTo /D <6571756174696F6E2E322E322E342E3336> >> >>}
\MusicPDFWriteObject{561}{<< /Type /Annot /Subtype /Link /Border [0 0 0] /H /I /C [1 0 0] /Rect [282.405 618.719 317.598 630.672] /A << /S /GoTo /D <6571756174696F6E2E322E322E342E3239> >> >>}
\MusicPDFWriteObject{562}{<< /Type /Annot /Subtype /Link /Border [0 0 0] /H /I /C [1 0 0] /Rect [435.585 557.677 468.519 570.523] /A << /S /GoTo /D <74686D2E442E312E322E31> >> >>}
\MusicPDFWriteObject{563}{<< /Type /Annot /Subtype /Link /Border [0 0 0] /H /I /C [1 0 0] /Rect [487.987 544.78 520.69 557.372] /A << /S /GoTo /D <74686D2E412E332E322E31> >> >>}
\MusicPDFWriteObject{564}{<< /Type /Annot /Subtype /Link /Border [0 0 0] /H /I /C [1 0 0] /Rect [196.558 339.834 226.771 351.787] /A << /S /GoTo /D <74686D2E322E322E352E32> >> >>}
\MusicPDFWriteObject{565}{<< /Type /Annot /Subtype /Link /Border [0 0 0] /H /I /C [1 0 0] /Rect [168.92 95.747 199.133 107.699] /A << /S /GoTo /D <6571756174696F6E2E322E322E322E33> >> >>}
\MusicPDFWriteObject{566}{<< /Type /Annot /Subtype /Link /Border [0 0 0] /H /I /C [1 0 0] /Rect [261.864 95.746 291.513 107.719] /A << /S /GoTo /D <74686D2E322E322E322E31> >> >>}
\MusicPDFWriteObject{567}{<< /Type /Annot /Subtype /Link /Border [0 0 0] /H /I /C [1 0 0] /Rect [315.851 610.951 353.534 622.904] /A << /S /GoTo /D <6571756174696F6E2E412E322E332E3134> >> >>}
\MusicPDFWriteObject{568}{<< /Type /Annot /Subtype /Link /Border [0 0 0] /H /I /C [1 0 0] /Rect [265.05 597.8 302.734 609.753] /A << /S /GoTo /D <6571756174696F6E2E412E322E332E3136> >> >>}
\MusicPDFWriteObject{569}{<< /Type /Annot /Subtype /Link /Border [0 0 0] /H /I /C [1 0 0] /Rect [185.812 534.58 216.025 546.533] /A << /S /GoTo /D <6571756174696F6E2E322E322E352E35> >> >>}
\MusicPDFWriteObject{570}{<< /Type /Annot /Subtype /Link /Border [0 0 0] /H /I /C [1 0 0] /Rect [375.396 482.881 408.793 495.727] /A << /S /GoTo /D <74686D2E442E312E322E31> >> >>}
\MusicPDFWriteObject{571}{<< /Type /Annot /Subtype /Link /Border [0 0 0] /H /I /C [1 0 0] /Rect [509.841 469.984 543.159 482.577] /A << /S /GoTo /D <74686D2E412E332E322E31> >> >>}
\MusicPDFWriteObject{572}{<< /Type /Annot /Subtype /Link /Border [0 0 0] /H /I /C [1 0 0] /Rect [160.817 443.683 191.03 455.636] /A << /S /GoTo /D <6571756174696F6E2E322E322E352E37> >> >>}
\MusicPDFWriteObject{573}{<< /Type /Annot /Subtype /Link /Border [0 0 0] /H /I /C [1 0 0] /Rect [138.353 392.776 171.133 405.191] /A << /S /GoTo /D <74686D2E442E312E322E34> >> >>}
\MusicPDFWriteObject{574}{<< /Type /Annot /Subtype /Link /Border [0 0 0] /H /I /C [0 1 0] /Rect [215.135 395.266 250.189 404.134] /A << /S /GoTo /D <636974652E426F752D4949> >> >>}
\MusicPDFWriteObject{575}{<< /Type /Annot /Subtype /Link /Border [0 0 0] /H /I /C [1 0 0] /Rect [118.955 351.169 151.044 365.616] /A << /S /GoTo /D <74686D2E412E322E332E38> >> >>}
\MusicPDFWriteObject{576}{<< /Type /Annot /Subtype /Link /Border [0 0 0] /H /I /C [1 0 0] /Rect [511.675 353.023 541.888 364.976] /A << /S /GoTo /D <6571756174696F6E2E322E322E352E38> >> >>}
\MusicPDFWriteObject{577}{<< /Type /Annot /Subtype /Link /Border [0 0 0] /H /I /C [1 0 0] /Rect [328.962 339.873 359.175 351.825] /A << /S /GoTo /D <6571756174696F6E2E322E322E352E33> >> >>}
\MusicPDFWriteObject{578}{<< /Type /Annot /Subtype /Link /Border [0 0 0] /H /I /C [1 0 0] /Rect [446.925 339.873 480.26 351.825] /A << /S /GoTo /D <74686D2E442E312E322E33> >> >>}
\MusicPDFWriteObject{579}{<< /Type /Annot /Subtype /Link /Border [0 0 0] /H /I /C [0 1 0] /Rect [524.868 342.363 543.17 351.231] /A << /S /GoTo /D <636974652E4445> >> >>}
\MusicPDFWriteObject{580}{<< /Type /Annot /Subtype /Link /Border [0 0 0] /H /I /C [1 0 0] /Rect [380.494 326.722 410.707 338.675] /A << /S /GoTo /D <6571756174696F6E2E322E322E352E33> >> >>}
\MusicPDFWriteObject{581}{<< /Type /Annot /Subtype /Link /Border [0 0 0] /H /I /C [1 0 0] /Rect [324.042 287.27 354.256 299.222] /A << /S /GoTo /D <6571756174696F6E2E322E322E352E33> >> >>}
\MusicPDFWriteObject{582}{<< /Type /Annot /Subtype /Link /Border [0 0 0] /H /I /C [1 0 0] /Rect [493.448 260.968 523.661 272.921] /A << /S /GoTo /D <6571756174696F6E2E322E322E352E32> >> >>}
\MusicPDFWriteObject{583}{<< /Type /Annot /Subtype /Link /Border [0 0 0] /H /I /C [1 0 0] /Rect [205.223 550.12 238.682 560.91] /A << /S /GoTo /D <74686D2E442E312E322E31> >> >>}
\MusicPDFWriteObject{584}{<< /Type /Annot /Subtype /Link /Border [0 0 0] /H /I /C [1 0 0] /Rect [208.597 454.459 242.055 466.412] /A << /S /GoTo /D <74686D2E442E312E322E31> >> >>}
\MusicPDFWriteObject{585}{<< /Type /Annot /Subtype /Link /Border [0 0 0] /H /I /C [0 1 0] /Rect [287.916 456.949 306.218 465.817] /A << /S /GoTo /D <636974652E4445> >> >>}
\MusicPDFWriteObject{586}{<< /Type /Annot /Subtype /Link /Border [0 0 0] /H /I /C [1 0 0] /Rect [441.008 454.459 471.221 466.412] /A << /S /GoTo /D <6571756174696F6E2E322E322E352E34> >> >>}
\MusicPDFWriteObject{587}{<< /Type /Annot /Subtype /Link /Border [0 0 0] /H /I /C [1 0 0] /Rect [293.72 416.368 323.51 429.214] /A << /S /GoTo /D <74686D2E322E322E352E32> >> >>}
\MusicPDFWriteObject{588}{<< /Type /Annot /Subtype /Link /Border [0 0 0] /H /I /C [1 0 0] /Rect [117.346 405.961 147.559 414.815] /A << /S /GoTo /D <74686D2E322E322E352E31> >> >>}
\MusicPDFWriteObject{589}{<< /Type /Annot /Subtype /Link /Border [0 0 0] /H /I /C [1 0 0] /Rect [205.143 342.401 235.666 353.191] /A << /S /GoTo /D <74686D2E342E332E332E31> >> >>}
\MusicPDFWriteObject{590}{<< /Type /Annot /Subtype /Link /Border [0 0 0] /H /I /C [1 0 0] /Rect [402.489 280.776 432.702 291.567] /A << /S /GoTo /D <74686D2E322E322E332E34> >> >>}
\MusicPDFWriteObject{591}{<< /Type /Annot /Subtype /Link /Border [0 0 0] /H /I /C [1 0 0] /Rect [137.499 569.301 167.713 581.254] /A << /S /GoTo /D <6571756174696F6E2E322E322E352E39> >> >>}
\MusicPDFWriteObject{592}{<< /Type /Annot /Subtype /Link /Border [0 0 0] /H /I /C [1 0 0] /Rect [180.623 569.301 215.817 581.254] /A << /S /GoTo /D <6571756174696F6E2E322E322E352E3130> >> >>}
\MusicPDFWriteObject{593}{<< /Type /Annot /Subtype /Link /Border [0 0 0] /H /I /C [1 0 0] /Rect [228.728 569.301 263.921 581.254] /A << /S /GoTo /D <6571756174696F6E2E322E322E352E3131> >> >>}
\MusicPDFWriteObject{594}{<< /Type /Annot /Subtype /Link /Border [0 0 0] /H /I /C [1 0 0] /Rect [276.832 569.301 312.025 581.254] /A << /S /GoTo /D <6571756174696F6E2E322E322E352E3132> >> >>}
\MusicPDFWriteObject{595}{<< /Type /Annot /Subtype /Link /Border [0 0 0] /H /I /C [1 0 0] /Rect [323.219 509.351 353.997 522.197] /A << /S /GoTo /D <74686D2E322E322E352E32> >> >>}
\MusicPDFWriteObject{596}{<< /Type /Annot /Subtype /Link /Border [0 0 0] /H /I /C [1 0 0] /Rect [419.416 509.351 450.194 522.197] /A << /S /GoTo /D <74686D2E322E322E332E34> >> >>}
\MusicPDFWriteObject{597}{<< /Type /Annot /Subtype /Link /Border [0 0 0] /H /I /C [1 0 0] /Rect [73.712 483.303 103.925 495.256] /A << /S /GoTo /D <6571756174696F6E2E322E322E332E34> >> >>}
\MusicPDFWriteObject{598}{<< /Type /Annot /Subtype /Link /Border [0 0 0] /H /I /C [1 0 0] /Rect [132.366 483.303 162.579 495.256] /A << /S /GoTo /D <6571756174696F6E2E322E322E332E35> >> >>}
\MusicPDFWriteObject{599}{<< /Type /Annot /Subtype /Link /Border [0 0 0] /H /I /C [1 0 0] /Rect [180.197 469.792 210.128 482.745] /A << /S /GoTo /D <74686D2E322E322E342E33> >> >>}
\MusicPDFWriteObject{600}{<< /Type /Annot /Subtype /Link /Border [0 0 0] /H /I /C [1 0 0] /Rect [254.054 470.153 289.247 482.105] /A << /S /GoTo /D <6571756174696F6E2E322E322E352E3131> >> >>}
\MusicPDFWriteObject{601}{<< /Type /Annot /Subtype /Link /Border [0 0 0] /H /I /C [1 0 0] /Rect [73.712 372.317 108.905 384.27] /A << /S /GoTo /D <6571756174696F6E2E322E322E352E3133> >> >>}
\MusicPDFWriteObject{602}{<< /Type /Annot /Subtype /Link /Border [0 0 0] /H /I /C [1 0 0] /Rect [121.861 269.895 152.639 282.21] /A << /S /GoTo /D <74686D2E322E322E342E33> >> >>}
\MusicPDFWriteObject{603}{<< /Type /Annot /Subtype /Link /Border [0 0 0] /H /I /C [1 0 0] /Rect [296.901 257.105 332.094 269.057] /A << /S /GoTo /D <6571756174696F6E2E322E322E352E3134> >> >>}
\MusicPDFWriteObject{604}{<< /Type /Annot /Subtype /Link /Border [0 0 0] /H /I /C [1 0 0] /Rect [200.521 191.351 235.715 203.304] /A << /S /GoTo /D <6571756174696F6E2E322E322E352E3132> >> >>}
\MusicPDFWriteObject{605}{<< /Type /Annot /Subtype /Link /Border [0 0 0] /H /I /C [1 0 0] /Rect [314.703 178.2 349.896 190.153] /A << /S /GoTo /D <6571756174696F6E2E322E322E352E3131> >> >>}
\MusicPDFWriteObject{606}{<< /Type /Annot /Subtype /Link /Border [0 0 0] /H /I /C [1 0 0] /Rect [275.658 164.689 306.435 177.002] /A << /S /GoTo /D <74686D2E322E322E332E34> >> >>}
\MusicPDFWriteObject{607}{<< /Type /Annot /Subtype /Link /Border [0 0 0] /H /I /C [1 0 0] /Rect [295.162 96.3 320.577 107.148] /A << /S /GoTo /D <73756273656374696F6E2E412E312E35> >> >>}
\MusicPDFWriteObject{608}{<< /Type /Annot /Subtype /Link /Border [0 0 0] /H /I /C [1 0 0] /Rect [271.26 687.123 304.454 695.977] /A << /S /GoTo /D <74686D2E412E332E342E31> >> >>}
\MusicPDFWriteObject{609}{<< /Type /Annot /Subtype /Link /Border [0 0 0] /H /I /C [1 0 0] /Rect [512.929 553.125 543.707 565.541] /A << /S /GoTo /D <74686D2E322E322E332E31> >> >>}
\MusicPDFWriteObject{610}{<< /Type /Annot /Subtype /Link /Border [0 0 0] /H /I /C [1 0 0] /Rect [331.802 556.894 362.016 568.846] /A << /S /GoTo /D <6571756174696F6E2E322E332E312E31> >> >>}
\MusicPDFWriteObject{611}{<< /Type /Annot /Subtype /Link /Border [0 0 0] /H /I /C [1 0 0] /Rect [253.586 454.712 283.799 466.665] /A << /S /GoTo /D <6571756174696F6E2E322E332E312E33> >> >>}
\MusicPDFWriteObject{612}{<< /Type /Annot /Subtype /Link /Border [0 0 0] /H /I /C [0 1 0] /Rect [237.828 390.711 253.82 399.579] /A << /S /GoTo /D <636974652E436F> >> >>}
\MusicPDFWriteObject{613}{<< /Type /Annot /Subtype /Link /Border [0 0 0] /H /I /C [0 1 0] /Rect [399.521 390.711 415.513 399.579] /A << /S /GoTo /D <636974652E436F> >> >>}
\MusicPDFWriteObject{614}{<< /Type /Annot /Subtype /Link /Border [0 0 0] /H /I /C [0 1 0] /Rect [73.721 377.561 89.713 386.429] /A << /S /GoTo /D <636974652E436F> >> >>}
\MusicPDFWriteObject{615}{<< /Type /Annot /Subtype /Link /Border [0 0 0] /H /I /C [1 0 0] /Rect [226.466 361.92 256.679 373.873] /A << /S /GoTo /D <74686D2E322E332E312E32> >> >>}
\MusicPDFWriteObject{616}{<< /Type /Annot /Subtype /Link /Border [0 0 0] /H /I /C [0 1 0] /Rect [203.534 351.259 219.526 360.127] /A << /S /GoTo /D <636974652E436F> >> >>}
\MusicPDFWriteObject{617}{<< /Type /Annot /Subtype /Link /Border [0 0 0] /H /I /C [1 0 0] /Rect [179.354 265.192 209.567 277.145] /A << /S /GoTo /D <6571756174696F6E2E322E312E342E31> >> >>}
\MusicPDFWriteObject{618}{<< /Type /Annot /Subtype /Link /Border [0 0 0] /H /I /C [1 0 0] /Rect [293.814 264.832 324.591 277.147] /A << /S /GoTo /D <74686D2E322E312E342E31> >> >>}
\MusicPDFWriteObject{619}{<< /Type /Annot /Subtype /Link /Border [0 0 0] /H /I /C [1 0 0] /Rect [358.515 125.308 388.728 137.261] /A << /S /GoTo /D <6571756174696F6E2E322E332E312E32> >> >>}
\MusicPDFWriteObject{620}{<< /Type /Annot /Subtype /Link /Border [0 0 0] /H /I /C [1 0 0] /Rect [124.871 107.402 157.727 124.51] /A << /S /GoTo /D <74686D2E412E322E342E34> >> >>}
\MusicPDFWriteObject{621}{<< /Type /Annot /Subtype /Link /Border [0 0 0] /H /I /C [1 0 0] /Rect [256.845 96.045 289.549 107.998] /A << /S /GoTo /D <6571756174696F6E2E412E322E342E34> >> >>}
\MusicPDFWriteObject{622}{<< /Type /Annot /Subtype /Link /Border [0 0 0] /H /I /C [1 0 0] /Rect [182.152 572.474 212.366 584.427] /A << /S /GoTo /D <6571756174696F6E2E322E332E312E32> >> >>}
\MusicPDFWriteObject{623}{<< /Type /Annot /Subtype /Link /Border [0 0 0] /H /I /C [1 0 0] /Rect [320.073 448.336 350.286 460.289] /A << /S /GoTo /D <6571756174696F6E2E322E332E312E35> >> >>}
\MusicPDFWriteObject{624}{<< /Type /Annot /Subtype /Link /Border [0 0 0] /H /I /C [1 0 0] /Rect [334.935 422.034 364.584 433.99] /A << /S /GoTo /D <74686D2E322E322E312E33> >> >>}
\MusicPDFWriteObject{625}{<< /Type /Annot /Subtype /Link /Border [0 0 0] /H /I /C [1 0 0] /Rect [87.328 408.884 117.541 420.837] /A << /S /GoTo /D <6571756174696F6E2E322E332E312E34> >> >>}
\MusicPDFWriteObject{626}{<< /Type /Annot /Subtype /Link /Border [0 0 0] /H /I /C [1 0 0] /Rect [186.91 395.733 217.123 407.686] /A << /S /GoTo /D <6571756174696F6E2E322E332E312E32> >> >>}
\MusicPDFWriteObject{627}{<< /Type /Annot /Subtype /Link /Border [0 0 0] /H /I /C [1 0 0] /Rect [134.694 382.583 167.397 394.535] /A << /S /GoTo /D <6571756174696F6E2E412E322E342E35> >> >>}
\MusicPDFWriteObject{628}{<< /Type /Annot /Subtype /Link /Border [0 0 0] /H /I /C [1 0 0] /Rect [441.684 154.88 471.897 166.833] /A << /S /GoTo /D <6571756174696F6E2E322E332E312E37> >> >>}
\MusicPDFWriteObject{629}{<< /Type /Annot /Subtype /Link /Border [0 0 0] /H /I /C [1 0 0] /Rect [217.819 464.916 248.032 476.869] /A << /S /GoTo /D <6571756174696F6E2E322E332E312E34> >> >>}
\MusicPDFWriteObject{630}{<< /Type /Annot /Subtype /Link /Border [0 0 0] /H /I /C [1 0 0] /Rect [70.755 438.615 100.968 450.567] /A << /S /GoTo /D <74686D2E322E322E312E33> >> >>}
\MusicPDFWriteObject{631}{<< /Type /Annot /Subtype /Link /Border [0 0 0] /H /I /C [1 0 0] /Rect [227.445 267.655 260.148 279.607] /A << /S /GoTo /D <74686D2E412E322E322E35> >> >>}
\MusicPDFWriteObject{632}{<< /Type /Annot /Subtype /Link /Border [0 0 0] /H /I /C [1 0 0] /Rect [124.623 672.036 154.836 682.827] /A << /S /GoTo /D <74686D2E322E322E312E33> >> >>}
\MusicPDFWriteObject{633}{<< /Type /Annot /Subtype /Link /Border [0 0 0] /H /I /C [1 0 0] /Rect [294.712 623.831 324.925 635.784] /A << /S /GoTo /D <6571756174696F6E2E322E332E312E33> >> >>}
\MusicPDFWriteObject{634}{<< /Type /Annot /Subtype /Link /Border [0 0 0] /H /I /C [1 0 0] /Rect [365.337 610.32 396.114 622.635] /A << /S /GoTo /D <74686D2E322E312E342E31> >> >>}
\MusicPDFWriteObject{635}{<< /Type /Annot /Subtype /Link /Border [0 0 0] /H /I /C [1 0 0] /Rect [494.099 595.005 530.797 609.484] /A << /S /GoTo /D <74686D2E432E322E332E3133> >> >>}
\MusicPDFWriteObject{636}{<< /Type /Annot /Subtype /Link /Border [0 0 0] /H /I /C [1 0 0] /Rect [170.451 522.374 201.229 534.689] /A << /S /GoTo /D <74686D2E322E312E322E34> >> >>}
\MusicPDFWriteObject{637}{<< /Type /Annot /Subtype /Link /Border [0 0 0] /H /I /C [1 0 0] /Rect [364.329 509.583 396.755 521.538] /A << /S /GoTo /D <74686D2E432E322E352E31> >> >>}
\MusicPDFWriteObject{638}{<< /Type /Annot /Subtype /Link /Border [0 0 0] /H /I /C [0 1 0] /Rect [376.744 432.839 394.942 441.707] /A << /S /GoTo /D <636974652E4A6172> >> >>}
\MusicPDFWriteObject{639}{<< /Type /Annot /Subtype /Link /Border [0 0 0] /H /I /C [0 1 0] /Rect [110.754 419.688 129.859 428.556] /A << /S /GoTo /D <636974652E4E42> >> >>}
\MusicPDFWriteObject{640}{<< /Type /Annot /Subtype /Link /Border [0 0 0] /H /I /C [1 0 0] /Rect [168.625 377.416 198.838 389.369] /A << /S /GoTo /D <6571756174696F6E2E322E332E312E34> >> >>}
\MusicPDFWriteObject{641}{<< /Type /Annot /Subtype /Link /Border [0 0 0] /H /I /C [1 0 0] /Rect [462.066 324.481 492.279 336.433] /A << /S /GoTo /D <6571756174696F6E2E322E332E312E35> >> >>}
\MusicPDFWriteObject{642}{<< /Type /Annot /Subtype /Link /Border [0 0 0] /H /I /C [1 0 0] /Rect [272.596 271.213 302.809 283.166] /A << /S /GoTo /D <6571756174696F6E2E322E332E312E33> >> >>}
\MusicPDFWriteObject{643}{<< /Type /Annot /Subtype /Link /Border [0 0 0] /H /I /C [1 0 0] /Rect [126.394 218.61 156.607 230.563] /A << /S /GoTo /D <6571756174696F6E2E322E332E312E39> >> >>}
\MusicPDFWriteObject{644}{<< /Type /Annot /Subtype /Link /Border [0 0 0] /H /I /C [1 0 0] /Rect [204.688 205.46 234.901 217.412] /A << /S /GoTo /D <6571756174696F6E2E322E332E312E33> >> >>}
\MusicPDFWriteObject{645}{<< /Type /Annot /Subtype /Link /Border [0 0 0] /H /I /C [0 1 0] /Rect [379.904 180.986 398.102 189.854] /A << /S /GoTo /D <636974652E4A6172> >> >>}
\MusicPDFWriteObject{646}{<< /Type /Annot /Subtype /Link /Border [0 0 0] /H /I /C [0 1 0] /Rect [109.094 167.835 128.198 176.703] /A << /S /GoTo /D <636974652E4E42> >> >>}
\MusicPDFWriteObject{647}{<< /Type /Annot /Subtype /Link /Border [0 0 0] /H /I /C [1 0 0] /Rect [264.153 165.345 294.367 177.298] /A << /S /GoTo /D <6571756174696F6E2E322E332E312E33> >> >>}
\MusicPDFWriteObject{648}{<< /Type /Annot /Subtype /Link /Border [0 0 0] /H /I /C [1 0 0] /Rect [192.118 125.563 222.331 137.515] /A << /S /GoTo /D <6571756174696F6E2E322E332E312E34> >> >>}
\MusicPDFWriteObject{649}{<< /Type /Annot /Subtype /Link /Border [0 0 0] /H /I /C [1 0 0] /Rect [133.836 98.737 166.539 110.689] /A << /S /GoTo /D <6571756174696F6E2E412E322E342E34> >> >>}
\MusicPDFWriteObject{650}{<< /Type /Annot /Subtype /Link /Border [0 0 0] /H /I /C [1 0 0] /Rect [418.4 684.273 448.923 696.588] /A << /S /GoTo /D <74686D2E322E312E342E31> >> >>}
\MusicPDFWriteObject{651}{<< /Type /Annot /Subtype /Link /Border [0 0 0] /H /I /C [1 0 0] /Rect [383.491 655.015 414.268 669.954] /A << /S /GoTo /D <74686D2E322E312E322E34> >> >>}
\MusicPDFWriteObject{652}{<< /Type /Annot /Subtype /Link /Border [0 0 0] /H /I /C [1 0 0] /Rect [490.895 655.015 523.931 669.954] /A << /S /GoTo /D <74686D2E432E322E352E31> >> >>}
\MusicPDFWriteObject{653}{<< /Type /Annot /Subtype /Link /Border [0 0 0] /H /I /C [1 0 0] /Rect [210.112 527.932 240.325 539.884] /A << /S /GoTo /D <6571756174696F6E2E322E332E312E34> >> >>}
\MusicPDFWriteObject{654}{<< /Type /Annot /Subtype /Link /Border [0 0 0] /H /I /C [1 0 0] /Rect [163.988 501.63 194.201 513.583] /A << /S /GoTo /D <6571756174696F6E2E322E332E312E34> >> >>}
\MusicPDFWriteObject{655}{<< /Type /Annot /Subtype /Link /Border [0 0 0] /H /I /C [1 0 0] /Rect [352.269 400.54 382.115 412.493] /A << /S /GoTo /D <74686D2E322E322E332E34> >> >>}
\MusicPDFWriteObject{656}{<< /Type /Annot /Subtype /Link /Border [0 0 0] /H /I /C [1 0 0] /Rect [324.403 373.879 355.181 386.193] /A << /S /GoTo /D <74686D2E322E312E312E36> >> >>}
\MusicPDFWriteObject{657}{<< /Type /Annot /Subtype /Link /Border [0 0 0] /H /I /C [1 0 0] /Rect [515.197 346.083 544.845 359.892] /A << /S /GoTo /D <74686D2E322E312E312E36> >> >>}
\MusicPDFWriteObject{658}{<< /Type /Annot /Subtype /Link /Border [0 0 0] /H /I /C [0 1 0] /Rect [245.726 324.126 264.831 332.994] /A << /S /GoTo /D <636974652E4E42> >> >>}
\MusicPDFWriteObject{659}{<< /Type /Annot /Subtype /Link /Border [0 0 0] /H /I /C [1 0 0] /Rect [418.96 567.639 448.609 581.841] /A << /S /GoTo /D <74686D2E322E332E312E34> >> >>}
\MusicPDFWriteObject{660}{<< /Type /Annot /Subtype /Link /Border [0 0 0] /H /I /C [1 0 0] /Rect [347.873 454.252 383.066 466.205] /A << /S /GoTo /D <6571756174696F6E2E322E332E312E3132> >> >>}
\MusicPDFWriteObject{661}{<< /Type /Annot /Subtype /Link /Border [0 0 0] /H /I /C [1 0 0] /Rect [480.666 365.221 515.86 377.174] /A << /S /GoTo /D <6571756174696F6E2E322E332E312E3134> >> >>}
\MusicPDFWriteObject{662}{<< /Type /Annot /Subtype /Link /Border [0 0 0] /H /I /C [1 0 0] /Rect [318.241 293.367 351.558 307.569] /A << /S /GoTo /D <74686D2E412E312E352E32> >> >>}
\MusicPDFWriteObject{663}{<< /Type /Annot /Subtype /Link /Border [0 0 0] /H /I /C [1 0 0] /Rect [70.755 268.151 100.968 280.104] /A << /S /GoTo /D <74686D2E322E332E312E32> >> >>}
\MusicPDFWriteObject{664}{<< /Type /Annot /Subtype /Link /Border [0 0 0] /H /I /C [1 0 0] /Rect [429.571 268.151 459.784 280.104] /A << /S /GoTo /D <74686D2E322E322E352E32> >> >>}
\MusicPDFWriteObject{665}{<< /Type /Annot /Subtype /Link /Border [0 0 0] /H /I /C [1 0 0] /Rect [337.679 254.64 367.327 266.953] /A << /S /GoTo /D <74686D2E322E332E312E32> >> >>}
\MusicPDFWriteObject{666}{<< /Type /Annot /Subtype /Link /Border [0 0 0] /H /I /C [1 0 0] /Rect [124.623 227.034 154.836 241.236] /A << /S /GoTo /D <74686D2E322E332E312E32> >> >>}
\MusicPDFWriteObject{667}{<< /Type /Annot /Subtype /Link /Border [0 0 0] /H /I /C [1 0 0] /Rect [165.384 196.918 196.162 207.766] /A << /S /GoTo /D <74686D2E322E322E342E32> >> >>}
\MusicPDFWriteObject{668}{<< /Type /Annot /Subtype /Link /Border [0 0 0] /H /I /C [1 0 0] /Rect [255.362 196.918 286.14 207.766] /A << /S /GoTo /D <74686D2E322E322E342E33> >> >>}
\MusicPDFWriteObject{669}{<< /Type /Annot /Subtype /Link /Border [0 0 0] /H /I /C [1 0 0] /Rect [319.662 651.16 349.311 663.475] /A << /S /GoTo /D <74686D2E322E322E342E32> >> >>}
\MusicPDFWriteObject{670}{<< /Type /Annot /Subtype /Link /Border [0 0 0] /H /I /C [1 0 0] /Rect [130.176 638.37 160.389 650.322] /A << /S /GoTo /D <6571756174696F6E2E322E312E342E35> >> >>}
\MusicPDFWriteObject{671}{<< /Type /Annot /Subtype /Link /Border [0 0 0] /H /I /C [1 0 0] /Rect [236.066 371.639 265.715 383.592] /A << /S /GoTo /D <74686D2E322E332E322E32> >> >>}
\MusicPDFWriteObject{672}{<< /Type /Annot /Subtype /Link /Border [0 0 0] /H /I /C [1 0 0] /Rect [499.626 371.639 529.274 383.592] /A << /S /GoTo /D <74686D2E322E322E342E33> >> >>}
\MusicPDFWriteObject{673}{<< /Type /Annot /Subtype /Link /Border [0 0 0] /H /I /C [1 0 0] /Rect [203.197 341.016 236.758 349.87] /A << /S /GoTo /D <74686D2E322E332E322E32> >> >>}
\MusicPDFWriteObject{674}{<< /Type /Annot /Subtype /Link /Border [0 0 0] /H /I /C [1 0 0] /Rect [378.521 338.526 408.734 350.479] /A << /S /GoTo /D <6571756174696F6E2E322E332E322E33> >> >>}
\MusicPDFWriteObject{675}{<< /Type /Annot /Subtype /Link /Border [0 0 0] /H /I /C [1 0 0] /Rect [132.355 325.376 162.568 337.328] /A << /S /GoTo /D <6571756174696F6E2E322E332E322E32> >> >>}
\MusicPDFWriteObject{676}{<< /Type /Annot /Subtype /Link /Border [0 0 0] /H /I /C [1 0 0] /Rect [397.42 312.17 427.633 324.122] /A << /S /GoTo /D <6571756174696F6E2E322E332E322E34> >> >>}
\MusicPDFWriteObject{677}{<< /Type /Annot /Subtype /Link /Border [0 0 0] /H /I /C [1 0 0] /Rect [120.687 284.258 150.336 298.46] /A << /S /GoTo /D <74686D2E322E322E342E33> >> >>}
\MusicPDFWriteObject{678}{<< /Type /Annot /Subtype /Link /Border [0 0 0] /H /I /C [1 0 0] /Rect [211.67 284.258 241.319 298.46] /A << /S /GoTo /D <74686D2E322E322E352E34> >> >>}
\MusicPDFWriteObject{679}{<< /Type /Annot /Subtype /Link /Border [0 0 0] /H /I /C [1 0 0] /Rect [269.726 204.696 299.375 216.649] /A << /S /GoTo /D <74686D2E322E332E322E32> >> >>}
\MusicPDFWriteObject{680}{<< /Type /Annot /Subtype /Link /Border [0 0 0] /H /I /C [1 0 0] /Rect [130.379 82.596 152.845 94.549] /A << /S /GoTo /D <73756273656374696F6E2E322E342E32> >> >>}
\MusicPDFWriteObject{681}{<< /Type /Annot /Subtype /Link /Border [0 0 0] /H /I /C [1 0 0] /Rect [216.23 82.596 246.443 94.549] /A << /S /GoTo /D <74686D2E322E342E322E31> >> >>}
\MusicPDFWriteObject{682}{<< /Type /Annot /Subtype /Link /Border [0 0 0] /H /I /C [0 1 0] /Rect [302.696 687.123 318.688 695.991] /A << /S /GoTo /D <636974652E436F> >> >>}
\MusicPDFWriteObject{683}{<< /Type /Annot /Subtype /Link /Border [0 0 0] /H /I /C [1 0 0] /Rect [181.147 672.036 211.36 682.827] /A << /S /GoTo /D <74686D2E322E342E312E31> >> >>}
\MusicPDFWriteObject{684}{<< /Type /Annot /Subtype /Link /Border [0 0 0] /H /I /C [1 0 0] /Rect [148.268 172.561 178.481 184.514] /A << /S /GoTo /D <6571756174696F6E2E322E342E312E34> >> >>}
\MusicPDFWriteObject{685}{<< /Type /Annot /Subtype /Link /Border [0 0 0] /H /I /C [1 0 0] /Rect [206.923 172.561 237.136 184.514] /A << /S /GoTo /D <6571756174696F6E2E322E342E312E35> >> >>}
\MusicPDFWriteObject{686}{<< /Type /Annot /Subtype /Link /Border [0 0 0] /H /I /C [1 0 0] /Rect [391.852 376.759 422.066 388.712] /A << /S /GoTo /D <6571756174696F6E2E322E342E312E39> >> >>}
\MusicPDFWriteObject{687}{<< /Type /Annot /Subtype /Link /Border [0 0 0] /H /I /C [1 0 0] /Rect [454.151 376.759 489.344 388.712] /A << /S /GoTo /D <6571756174696F6E2E322E342E312E3132> >> >>}
\MusicPDFWriteObject{688}{<< /Type /Annot /Subtype /Link /Border [0 0 0] /H /I /C [1 0 0] /Rect [96.458 363.608 126.671 375.561] /A << /S /GoTo /D <6571756174696F6E2E322E342E312E33> >> >>}
\MusicPDFWriteObject{689}{<< /Type /Annot /Subtype /Link /Border [0 0 0] /H /I /C [1 0 0] /Rect [225.585 338.489 255.798 349.28] /A << /S /GoTo /D <74686D2E322E342E312E31> >> >>}
\MusicPDFWriteObject{690}{<< /Type /Annot /Subtype /Link /Border [0 0 0] /H /I /C [1 0 0] /Rect [203.197 314.753 236.758 323.607] /A << /S /GoTo /D <74686D2E322E342E312E31> >> >>}
\MusicPDFWriteObject{691}{<< /Type /Annot /Subtype /Link /Border [0 0 0] /H /I /C [1 0 0] /Rect [252.827 299.112 283.04 311.065] /A << /S /GoTo /D <6571756174696F6E2E322E342E312E31> >> >>}
\MusicPDFWriteObject{692}{<< /Type /Annot /Subtype /Link /Border [0 0 0] /H /I /C [1 0 0] /Rect [414.87 299.112 444.518 311.065] /A << /S /GoTo /D <74686D2E322E342E312E32> >> >>}
\MusicPDFWriteObject{693}{<< /Type /Annot /Subtype /Link /Border [0 0 0] /H /I /C [1 0 0] /Rect [419.7 285.962 449.913 297.914] /A << /S /GoTo /D <6571756174696F6E2E322E342E312E31> >> >>}
\MusicPDFWriteObject{694}{<< /Type /Annot /Subtype /Link /Border [0 0 0] /H /I /C [1 0 0] /Rect [406.805 684.633 441.998 696.586] /A << /S /GoTo /D <6571756174696F6E2E322E342E312E3133> >> >>}
\MusicPDFWriteObject{695}{<< /Type /Annot /Subtype /Link /Border [0 0 0] /H /I /C [1 0 0] /Rect [486.29 509.832 516.503 521.785] /A << /S /GoTo /D <6571756174696F6E2E322E342E312E31> >> >>}
\MusicPDFWriteObject{696}{<< /Type /Annot /Subtype /Link /Border [0 0 0] /H /I /C [1 0 0] /Rect [400.885 433.514 433.588 445.467] /A << /S /GoTo /D <6571756174696F6E2E412E322E342E35> >> >>}
\MusicPDFWriteObject{697}{<< /Type /Annot /Subtype /Link /Border [0 0 0] /H /I /C [1 0 0] /Rect [363.754 420.364 398.947 432.316] /A << /S /GoTo /D <6571756174696F6E2E322E342E312E3135> >> >>}
\MusicPDFWriteObject{698}{<< /Type /Annot /Subtype /Link /Border [0 0 0] /H /I /C [1 0 0] /Rect [191.384 219.636 224.087 231.588] /A << /S /GoTo /D <6571756174696F6E2E412E322E342E35> >> >>}
\MusicPDFWriteObject{699}{<< /Type /Annot /Subtype /Link /Border [0 0 0] /H /I /C [1 0 0] /Rect [322.05 113.149 351.953 125.104] /A << /S /GoTo /D <74686D2E322E342E312E32> >> >>}
\MusicPDFWriteObject{700}{<< /Type /Annot /Subtype /Link /Border [0 0 0] /H /I /C [1 0 0] /Rect [154.369 99.999 184.299 111.952] /A << /S /GoTo /D <74686D2E322E342E312E31> >> >>}
\MusicPDFWriteObject{701}{<< /Type /Annot /Subtype /Link /Border [0 0 0] /H /I /C [0 1 0] /Rect [309.441 637.897 325.433 646.765] /A << /S /GoTo /D <636974652E436F> >> >>}
\MusicPDFWriteObject{702}{<< /Type /Annot /Subtype /Link /Border [0 0 0] /H /I /C [1 0 0] /Rect [330.185 622.256 362.888 634.209] /A << /S /GoTo /D <6571756174696F6E2E412E322E342E37> >> >>}
\MusicPDFWriteObject{703}{<< /Type /Annot /Subtype /Link /Border [0 0 0] /H /I /C [1 0 0] /Rect [374.217 432.479 404.43 444.794] /A << /S /GoTo /D <74686D2E322E332E312E32> >> >>}
\MusicPDFWriteObject{704}{<< /Type /Annot /Subtype /Link /Border [0 0 0] /H /I /C [1 0 0] /Rect [436.672 419.689 466.886 431.642] /A << /S /GoTo /D <6571756174696F6E2E322E342E312E31> >> >>}
\MusicPDFWriteObject{705}{<< /Type /Annot /Subtype /Link /Border [0 0 0] /H /I /C [1 0 0] /Rect [70.755 406.538 100.403 418.491] /A << /S /GoTo /D <74686D2E322E342E312E31> >> >>}
\MusicPDFWriteObject{706}{<< /Type /Annot /Subtype /Link /Border [0 0 0] /H /I /C [1 0 0] /Rect [171.962 406.538 201.61 418.491] /A << /S /GoTo /D <74686D2E322E332E312E32> >> >>}
\MusicPDFWriteObject{707}{<< /Type /Annot /Subtype /Link /Border [0 0 0] /H /I /C [1 0 0] /Rect [320.197 406.538 351.879 418.491] /A << /S /GoTo /D <74686D2E422E322E312E37> >> >>}
\MusicPDFWriteObject{708}{<< /Type /Annot /Subtype /Link /Border [0 0 0] /H /I /C [1 0 0] /Rect [171.005 393.387 203.899 405.34] /A << /S /GoTo /D <74686D2E422E332E322E31> >> >>}
\MusicPDFWriteObject{709}{<< /Type /Annot /Subtype /Link /Border [0 0 0] /H /I /C [0 1 0] /Rect [260.676 395.877 278.521 404.745] /A << /S /GoTo /D <636974652E484C> >> >>}
\MusicPDFWriteObject{710}{<< /Type /Annot /Subtype /Link /Border [0 0 0] /H /I /C [0 1 0] /Rect [379.547 395.877 397.392 404.745] /A << /S /GoTo /D <636974652E484C> >> >>}
\MusicPDFWriteObject{711}{<< /Type /Annot /Subtype /Link /Border [0 0 0] /H /I /C [1 0 0] /Rect [98.745 380.237 128.958 392.189] /A << /S /GoTo /D <6571756174696F6E2E322E342E312E31> >> >>}
\MusicPDFWriteObject{712}{<< /Type /Annot /Subtype /Link /Border [0 0 0] /H /I /C [0 1 0] /Rect [525.687 369.576 543.159 378.444] /A << /S /GoTo /D <636974652E4346> >> >>}
\MusicPDFWriteObject{713}{<< /Type /Annot /Subtype /Link /Border [0 0 0] /H /I /C [1 0 0] /Rect [287.251 314.483 317.464 326.436] /A << /S /GoTo /D <6571756174696F6E2E322E342E312E31> >> >>}
\MusicPDFWriteObject{714}{<< /Type /Annot /Subtype /Link /Border [0 0 0] /H /I /C [1 0 0] /Rect [110.087 288.735 117.104 299.525] /A << /S /GoTo /D <636861707465722E35> >> >>}
\MusicPDFWriteObject{715}{<< /Type /Annot /Subtype /Link /Border [0 0 0] /H /I /C [1 0 0] /Rect [345.495 236.132 375.708 246.922] /A << /S /GoTo /D <74686D2E322E342E312E31> >> >>}
\MusicPDFWriteObject{716}{<< /Type /Annot /Subtype /Link /Border [0 0 0] /H /I /C [1 0 0] /Rect [132.612 159.085 163.39 169.876] /A << /S /GoTo /D <74686D2E322E342E322E31> >> >>}
\MusicPDFWriteObject{717}{<< /Type /Annot /Subtype /Link /Border [0 0 0] /H /I /C [1 0 0] /Rect [203.687 145.381 233.901 157.334] /A << /S /GoTo /D <6571756174696F6E2E352E312E312E32> >> >>}
\MusicPDFWriteObject{718}{<< /Type /Annot /Subtype /Link /Border [0 0 0] /H /I /C [0 1 0] /Rect [190.202 108.419 206.194 117.287] /A << /S /GoTo /D <636974652E436F> >> >>}
\MusicPDFWriteObject{719}{<< /Type /Annot /Subtype /Link /Border [0 0 0] /H /I /C [1 0 0] /Rect [370.975 588.367 401.075 600.682] /A << /S /GoTo /D <74686D2E322E312E342E31> >> >>}
\MusicPDFWriteObject{720}{<< /Type /Annot /Subtype /Link /Border [0 0 0] /H /I /C [1 0 0] /Rect [492.063 573.052 530.177 587.531] /A << /S /GoTo /D <74686D2E432E322E332E3133> >> >>}
\MusicPDFWriteObject{721}{<< /Type /Annot /Subtype /Link /Border [0 0 0] /H /I /C [1 0 0] /Rect [184.436 561.734 222.551 574.049] /A << /S /GoTo /D <74686D2E432E322E332E3132> >> >>}
\MusicPDFWriteObject{722}{<< /Type /Annot /Subtype /Link /Border [0 0 0] /H /I /C [0 1 0] /Rect [246.47 564.584 264.669 573.452] /A << /S /GoTo /D <636974652E4A6172> >> >>}
\MusicPDFWriteObject{723}{<< /Type /Annot /Subtype /Link /Border [0 0 0] /H /I /C [1 0 0] /Rect [200.379 482.988 230.507 494.943] /A << /S /GoTo /D <74686D2E322E312E322E34> >> >>}
\MusicPDFWriteObject{724}{<< /Type /Annot /Subtype /Link /Border [0 0 0] /H /I /C [1 0 0] /Rect [467.419 469.478 499.845 481.793] /A << /S /GoTo /D <74686D2E432E322E352E31> >> >>}
\MusicPDFWriteObject{725}{<< /Type /Annot /Subtype /Link /Border [0 0 0] /H /I /C [0 1 0] /Rect [376.744 393.093 394.942 401.961] /A << /S /GoTo /D <636974652E4A6172> >> >>}
\MusicPDFWriteObject{726}{<< /Type /Annot /Subtype /Link /Border [0 0 0] /H /I /C [0 1 0] /Rect [110.754 379.942 129.859 388.81] /A << /S /GoTo /D <636974652E4E42> >> >>}
\MusicPDFWriteObject{727}{<< /Type /Annot /Subtype /Link /Border [0 0 0] /H /I /C [1 0 0] /Rect [376.013 351.151 406.226 363.104] /A << /S /GoTo /D <74686D2E322E332E312E32> >> >>}
\MusicPDFWriteObject{728}{<< /Type /Annot /Subtype /Link /Border [0 0 0] /H /I /C [1 0 0] /Rect [305.469 297.98 335.202 309.933] /A << /S /GoTo /D <74686D2E322E312E342E31> >> >>}
\MusicPDFWriteObject{729}{<< /Type /Annot /Subtype /Link /Border [0 0 0] /H /I /C [1 0 0] /Rect [388.209 632.285 418.422 644.238] /A << /S /GoTo /D <6571756174696F6E2E322E342E322E32> >> >>}
\MusicPDFWriteObject{730}{<< /Type /Annot /Subtype /Link /Border [0 0 0] /H /I /C [1 0 0] /Rect [330.295 356.803 362.998 368.756] /A << /S /GoTo /D <6571756174696F6E2E412E322E342E35> >> >>}
\MusicPDFWriteObject{731}{<< /Type /Annot /Subtype /Link /Border [0 0 0] /H /I /C [1 0 0] /Rect [425.885 356.443 458.588 368.758] /A << /S /GoTo /D <74686D2E412E322E342E33> >> >>}
\MusicPDFWriteObject{732}{<< /Type /Annot /Subtype /Link /Border [0 0 0] /H /I /C [1 0 0] /Rect [357.827 270.146 390.53 281.022] /A << /S /GoTo /D <74686D2E412E342E312E31> >> >>}
\MusicPDFWriteObject{733}{<< /Type /Annot /Subtype /Link /Border [0 0 0] /H /I /C [1 0 0] /Rect [332.7 229.866 363.478 242.18] /A << /S /GoTo /D <74686D2E322E332E312E32> >> >>}
\MusicPDFWriteObject{734}{<< /Type /Annot /Subtype /Link /Border [0 0 0] /H /I /C [1 0 0] /Rect [424.247 217.075 456.95 229.028] /A << /S /GoTo /D <6571756174696F6E2E412E322E342E34> >> >>}
\MusicPDFWriteObject{735}{<< /Type /Annot /Subtype /Link /Border [0 0 0] /H /I /C [1 0 0] /Rect [345.957 153.822 376.17 165.774] /A << /S /GoTo /D <6571756174696F6E2E322E342E322E31> >> >>}
\MusicPDFWriteObject{736}{<< /Type /Annot /Subtype /Link /Border [0 0 0] /H /I /C [1 0 0] /Rect [514.815 95.386 545.592 108.162] /A << /S /GoTo /D <74686D2E322E342E322E31> >> >>}
\MusicPDFWriteObject{737}{<< /Type /Annot /Subtype /Link /Border [0 0 0] /H /I /C [0 1 0] /Rect [377.628 687.123 395.1 695.991] /A << /S /GoTo /D <636974652E4346> >> >>}
\MusicPDFWriteObject{738}{<< /Type /Annot /Subtype /Link /Border [0 0 0] /H /I /C [1 0 0] /Rect [449.198 543.695 479.412 555.648] /A << /S /GoTo /D <6571756174696F6E2E322E342E332E31> >> >>}
\MusicPDFWriteObject{739}{<< /Type /Annot /Subtype /Link /Border [0 0 0] /H /I /C [1 0 0] /Rect [444.4 504.243 474.613 516.196] /A << /S /GoTo /D <6571756174696F6E2E322E342E332E31> >> >>}
\MusicPDFWriteObject{740}{<< /Type /Annot /Subtype /Link /Border [0 0 0] /H /I /C [0 1 0] /Rect [73.787 445.476 91.903 454.343] /A << /S /GoTo /D <636974652E4D65> >> >>}
\MusicPDFWriteObject{741}{<< /Type /Annot /Subtype /Link /Border [0 0 0] /H /I /C [0 1 0] /Rect [411.278 445.476 429.394 454.343] /A << /S /GoTo /D <636974652E4D65> >> >>}
\MusicPDFWriteObject{742}{<< /Type /Annot /Subtype /Link /Border [0 0 0] /H /I /C [1 0 0] /Rect [248.484 416.684 278.725 428.637] /A << /S /GoTo /D <74686D2E352E342E312E31> >> >>}
\MusicPDFWriteObject{743}{<< /Type /Annot /Subtype /Link /Border [0 0 0] /H /I /C [1 0 0] /Rect [71.004 377.232 101.217 389.184] /A << /S /GoTo /D <74686D2E352E312E312E31> >> >>}
\MusicPDFWriteObject{744}{<< /Type /Annot /Subtype /Link /Border [0 0 0] /H /I /C [1 0 0] /Rect [130.517 377.232 152.983 389.184] /A << /S /GoTo /D <73756273656374696F6E2E352E312E31> >> >>}
\MusicPDFWriteObject{745}{<< /Type /Annot /Subtype /Link /Border [0 0 0] /H /I /C [1 0 0] /Rect [203.197 355.869 236.758 364.723] /A << /S /GoTo /D <74686D2E322E342E332E31> >> >>}
\MusicPDFWriteObject{746}{<< /Type /Annot /Subtype /Link /Border [0 0 0] /H /I /C [1 0 0] /Rect [157.13 303.44 186.863 316.393] /A << /S /GoTo /D <6571756174696F6E2E322E332E322E32> >> >>}
\MusicPDFWriteObject{747}{<< /Type /Annot /Subtype /Link /Border [0 0 0] /H /I /C [1 0 0] /Rect [162.576 134.839 193.354 147.153] /A << /S /GoTo /D <6571756174696F6E2E322E332E322E32> >> >>}
\MusicPDFWriteObject{748}{<< /Type /Annot /Subtype /Link /Border [0 0 0] /H /I /C [1 0 0] /Rect [168.918 108.897 199.131 120.85] /A << /S /GoTo /D <6571756174696F6E2E322E342E332E32> >> >>}
\MusicPDFWriteObject{749}{<< /Type /Annot /Subtype /Link /Border [0 0 0] /H /I /C [1 0 0] /Rect [342.93 108.897 372.945 121.761] /A << /S /GoTo /D <74686D2E322E332E312E34> >> >>}
\MusicPDFWriteObject{750}{<< /Type /Annot /Subtype /Link /Border [0 0 0] /H /I /C [1 0 0] /Rect [300.703 95.746 331.481 108.339] /A << /S /GoTo /D <74686D2E322E342E322E31> >> >>}
\MusicPDFWriteObject{751}{<< /Type /Annot /Subtype /Link /Border [0 0 0] /H /I /C [1 0 0] /Rect [369.062 684.633 399.275 696.586] /A << /S /GoTo /D <6571756174696F6E2E322E342E332E32> >> >>}
\MusicPDFWriteObject{752}{<< /Type /Annot /Subtype /Link /Border [0 0 0] /H /I /C [1 0 0] /Rect [348.494 658.332 378.707 670.284] /A << /S /GoTo /D <6571756174696F6E2E322E342E332E32> >> >>}
\MusicPDFWriteObject{753}{<< /Type /Annot /Subtype /Link /Border [0 0 0] /H /I /C [1 0 0] /Rect [173.82 600.534 204.033 612.487] /A << /S /GoTo /D <6571756174696F6E2E322E342E332E33> >> >>}
\MusicPDFWriteObject{754}{<< /Type /Annot /Subtype /Link /Border [0 0 0] /H /I /C [1 0 0] /Rect [260.104 587.384 290.317 599.336] /A << /S /GoTo /D <6571756174696F6E2E322E342E332E31> >> >>}
\MusicPDFWriteObject{755}{<< /Type /Annot /Subtype /Link /Border [0 0 0] /H /I /C [1 0 0] /Rect [158.188 516.607 188.401 527.54] /A << /S /GoTo /D <6571756174696F6E2E322E332E322E32> >> >>}
\MusicPDFWriteObject{756}{<< /Type /Annot /Subtype /Link /Border [0 0 0] /H /I /C [1 0 0] /Rect [156.242 400.854 185.891 413.169] /A << /S /GoTo /D <6571756174696F6E2E322E332E322E32> >> >>}
\MusicPDFWriteObject{757}{<< /Type /Annot /Subtype /Link /Border [0 0 0] /H /I /C [1 0 0] /Rect [168.918 374.913 199.131 386.866] /A << /S /GoTo /D <6571756174696F6E2E322E342E332E34> >> >>}
\MusicPDFWriteObject{758}{<< /Type /Annot /Subtype /Link /Border [0 0 0] /H /I /C [1 0 0] /Rect [342.93 374.912 372.945 387.777] /A << /S /GoTo /D <74686D2E322E332E312E34> >> >>}
\MusicPDFWriteObject{759}{<< /Type /Annot /Subtype /Link /Border [0 0 0] /H /I /C [1 0 0] /Rect [365.233 348.611 395.447 360.566] /A << /S /GoTo /D <74686D2E322E342E322E31> >> >>}
\MusicPDFWriteObject{760}{<< /Type /Annot /Subtype /Link /Border [0 0 0] /H /I /C [1 0 0] /Rect [166.201 335.461 196.414 347.414] /A << /S /GoTo /D <6571756174696F6E2E322E342E332E34> >> >>}
\MusicPDFWriteObject{761}{<< /Type /Annot /Subtype /Link /Border [0 0 0] /H /I /C [1 0 0] /Rect [441.258 293.477 472.036 305.43] /A << /S /GoTo /D <74686D2E322E342E322E31> >> >>}
\MusicPDFWriteObject{762}{<< /Type /Annot /Subtype /Link /Border [0 0 0] /H /I /C [0 1 0] /Rect [354.311 251.571 373.415 260.439] /A << /S /GoTo /D <636974652E4E42> >> >>}
\MusicPDFWriteObject{763}{<< /Type /Annot /Subtype /Link /Border [0 0 0] /H /I /C [1 0 0] /Rect [179.348 235.93 209.082 247.883] /A << /S /GoTo /D <74686D2E322E312E342E31> >> >>}
\MusicPDFWriteObject{764}{<< /Type /Annot /Subtype /Link /Border [0 0 0] /H /I /C [0 1 0] /Rect [330.517 238.42 355.577 247.288] /A << /S /GoTo /D <636974652E53636861> >> >>}
\MusicPDFWriteObject{765}{<< /Type /Annot /Subtype /Link /Border [0 0 0] /H /I /C [0 1 0] /Rect [93.314 225.269 112.418 234.137] /A << /S /GoTo /D <636974652E4E42> >> >>}
\MusicPDFWriteObject{766}{<< /Type /Annot /Subtype /Link /Border [0 0 0] /H /I /C [0 1 0] /Rect [429.288 212.119 448.393 220.986] /A << /S /GoTo /D <636974652E4E42> >> >>}
\MusicPDFWriteObject{767}{<< /Type /Annot /Subtype /Link /Border [0 0 0] /H /I /C [0 1 0] /Rect [228.237 195.003 247.341 203.871] /A << /S /GoTo /D <636974652E4E42> >> >>}
\MusicPDFWriteObject{768}{<< /Type /Annot /Subtype /Link /Border [0 0 0] /H /I /C [1 0 0] /Rect [448.031 144.213 478.809 156.528] /A << /S /GoTo /D <74686D2E322E312E342E31> >> >>}
\MusicPDFWriteObject{769}{<< /Type /Annot /Subtype /Link /Border [0 0 0] /H /I /C [0 1 0] /Rect [360.972 133.912 386.033 142.78] /A << /S /GoTo /D <636974652E53636861> >> >>}
\MusicPDFWriteObject{770}{<< /Type /Annot /Subtype /Link /Border [0 0 0] /H /I /C [1 0 0] /Rect [171.52 104.76 202.297 117.074] /A << /S /GoTo /D <74686D2E322E342E322E31> >> >>}
\MusicPDFWriteObject{771}{<< /Type /Annot /Subtype /Link /Border [0 0 0] /H /I /C [1 0 0] /Rect [166.201 572.823 196.414 584.775] /A << /S /GoTo /D <6571756174696F6E2E322E342E332E32> >> >>}
\MusicPDFWriteObject{772}{<< /Type /Annot /Subtype /Link /Border [0 0 0] /H /I /C [1 0 0] /Rect [249.015 473.265 279.228 485.218] /A << /S /GoTo /D <6571756174696F6E2E322E342E332E36> >> >>}
\MusicPDFWriteObject{773}{<< /Type /Annot /Subtype /Link /Border [0 0 0] /H /I /C [0 1 0] /Rect [474.734 573.549 492.206 582.417] /A << /S /GoTo /D <636974652E4346> >> >>}
\MusicPDFWriteObject{774}{<< /Type /Annot /Subtype /Link /Border [0 0 0] /H /I /C [0 1 0] /Rect [472.636 534.097 490.108 542.965] /A << /S /GoTo /D <636974652E4346> >> >>}
\MusicPDFWriteObject{775}{<< /Type /Annot /Subtype /Link /Border [0 0 0] /H /I /C [0 1 0] /Rect [448.194 507.795 495.049 516.663] /A << /S /GoTo /D <636974652E5269656D616E6E> >> >>}
\MusicPDFWriteObject{776}{<< /Type /Annot /Subtype /Link /Border [0 0 0] /H /I /C [0 1 0] /Rect [297.681 442.042 315.154 450.909] /A << /S /GoTo /D <636974652E4346> >> >>}
\MusicPDFWriteObject{777}{<< /Type /Annot /Subtype /Link /Border [0 0 0] /H /I /C [0 1 0] /Rect [458.26 634.52 482.719 643.388] /A << /S /GoTo /D <636974652E435053> >> >>}
\MusicPDFWriteObject{778}{<< /Type /Annot /Subtype /Link /Border [0 0 0] /H /I /C [1 0 0] /Rect [82.959 527.377 321.174 538.168] /A << /S /GoTo /D <73656374696F6E2E332E31> >> >>}
\MusicPDFWriteObject{779}{<< /Type /Annot /Subtype /Link /Border [0 0 0] /H /I /C [1 0 0] /Rect [82.959 516.163 169.765 525.017] /A << /S /GoTo /D <73756273656374696F6E2E332E312E31> >> >>}
\MusicPDFWriteObject{780}{<< /Type /Annot /Subtype /Link /Border [0 0 0] /H /I /C [1 0 0] /Rect [82.959 501.076 354.237 511.867] /A << /S /GoTo /D <73756273656374696F6E2E332E312E32> >> >>}
\MusicPDFWriteObject{781}{<< /Type /Annot /Subtype /Link /Border [0 0 0] /H /I /C [1 0 0] /Rect [82.959 487.925 225.529 498.716] /A << /S /GoTo /D <73756273756273656374696F6E2E332E312E322E31> >> >>}
\MusicPDFWriteObject{782}{<< /Type /Annot /Subtype /Link /Border [0 0 0] /H /I /C [1 0 0] /Rect [82.959 474.774 291.363 485.565] /A << /S /GoTo /D <73756273756273656374696F6E2E332E312E322E32> >> >>}
\MusicPDFWriteObject{783}{<< /Type /Annot /Subtype /Link /Border [0 0 0] /H /I /C [1 0 0] /Rect [82.959 461.07 460.463 473.025] /A << /S /GoTo /D <73756273756273656374696F6E2E332E312E322E33> >> >>}
\MusicPDFWriteObject{784}{<< /Type /Annot /Subtype /Link /Border [0 0 0] /H /I /C [1 0 0] /Rect [82.959 450.41 234.341 459.264] /A << /S /GoTo /D <73656374696F6E2E332E32> >> >>}
\MusicPDFWriteObject{785}{<< /Type /Annot /Subtype /Link /Border [0 0 0] /H /I /C [1 0 0] /Rect [82.959 437.259 270.434 446.113] /A << /S /GoTo /D <73756273656374696F6E2E332E322E31> >> >>}
\MusicPDFWriteObject{786}{<< /Type /Annot /Subtype /Link /Border [0 0 0] /H /I /C [1 0 0] /Rect [82.959 424.108 242.088 432.962] /A << /S /GoTo /D <73756273656374696F6E2E332E322E32> >> >>}
\MusicPDFWriteObject{787}{<< /Type /Annot /Subtype /Link /Border [0 0 0] /H /I /C [1 0 0] /Rect [82.959 410.957 239.156 419.811] /A << /S /GoTo /D <73656374696F6E2E332E33> >> >>}
\MusicPDFWriteObject{788}{<< /Type /Annot /Subtype /Link /Border [0 0 0] /H /I /C [1 0 0] /Rect [82.959 397.807 228.963 406.661] /A << /S /GoTo /D <73756273656374696F6E2E332E332E31> >> >>}
\MusicPDFWriteObject{789}{<< /Type /Annot /Subtype /Link /Border [0 0 0] /H /I /C [1 0 0] /Rect [82.959 382.719 450.694 393.51] /A << /S /GoTo /D <73756273756273656374696F6E2E332E332E312E31> >> >>}
\MusicPDFWriteObject{790}{<< /Type /Annot /Subtype /Link /Border [0 0 0] /H /I /C [1 0 0] /Rect [82.959 369.568 325.638 380.359] /A << /S /GoTo /D <73756273756273656374696F6E2E332E332E312E32> >> >>}
\MusicPDFWriteObject{791}{<< /Type /Annot /Subtype /Link /Border [0 0 0] /H /I /C [1 0 0] /Rect [82.959 358.354 289.22 367.208] /A << /S /GoTo /D <73756273656374696F6E2E332E332E32> >> >>}
\MusicPDFWriteObject{792}{<< /Type /Annot /Subtype /Link /Border [0 0 0] /H /I /C [1 0 0] /Rect [82.959 343.267 430.476 354.058] /A << /S /GoTo /D <73756273756273656374696F6E2E332E332E322E31> >> >>}
\MusicPDFWriteObject{793}{<< /Type /Annot /Subtype /Link /Border [0 0 0] /H /I /C [1 0 0] /Rect [82.959 330.116 449.843 340.907] /A << /S /GoTo /D <73756273756273656374696F6E2E332E332E322E32> >> >>}
\MusicPDFWriteObject{794}{<< /Type /Annot /Subtype /Link /Border [0 0 0] /H /I /C [1 0 0] /Rect [82.959 316.965 276.187 327.756] /A << /S /GoTo /D <73656374696F6E2E332E34> >> >>}
\MusicPDFWriteObject{795}{<< /Type /Annot /Subtype /Link /Border [0 0 0] /H /I /C [1 0 0] /Rect [82.959 303.815 273.78 314.605] /A << /S /GoTo /D <73756273656374696F6E2E332E342E31> >> >>}
\MusicPDFWriteObject{796}{<< /Type /Annot /Subtype /Link /Border [0 0 0] /H /I /C [1 0 0] /Rect [82.959 290.664 350.474 301.455] /A << /S /GoTo /D <73756273656374696F6E2E332E342E32> >> >>}
\MusicPDFWriteObject{797}{<< /Type /Annot /Subtype /Link /Border [0 0 0] /H /I /C [1 0 0] /Rect [82.959 277.513 273.975 288.304] /A << /S /GoTo /D <73756273756273656374696F6E2E332E342E322E31> >> >>}
\MusicPDFWriteObject{798}{<< /Type /Annot /Subtype /Link /Border [0 0 0] /H /I /C [1 0 0] /Rect [82.959 264.362 475.725 275.21] /A << /S /GoTo /D <73756273756273656374696F6E2E332E342E322E32> >> >>}
\MusicPDFWriteObject{799}{<< /Type /Annot /Subtype /Link /Border [0 0 0] /H /I /C [1 0 0] /Rect [82.959 251.212 360.379 262.002] /A << /S /GoTo /D <73756273656374696F6E2E332E342E33> >> >>}
\MusicPDFWriteObject{800}{<< /Type /Annot /Subtype /Link /Border [0 0 0] /H /I /C [1 0 0] /Rect [82.959 238.061 404.19 248.852] /A << /S /GoTo /D <73756273756273656374696F6E2E332E342E332E31> >> >>}
\MusicPDFWriteObject{801}{<< /Type /Annot /Subtype /Link /Border [0 0 0] /H /I /C [1 0 0] /Rect [82.959 224.91 266.034 235.701] /A << /S /GoTo /D <73756273756273656374696F6E2E332E342E332E32> >> >>}
\MusicPDFWriteObject{802}{<< /Type /Annot /Subtype /Link /Border [0 0 0] /H /I /C [1 0 0] /Rect [298.443 127.136 331.901 139.089] /A << /S /GoTo /D <74686D2E442E322E312E31> >> >>}
\MusicPDFWriteObject{803}{<< /Type /Annot /Subtype /Link /Border [0 0 0] /H /I /C [0 1 0] /Rect [441.831 358.685 459.304 367.553] /A << /S /GoTo /D <636974652E4346> >> >>}
\MusicPDFWriteObject{804}{<< /Type /Annot /Subtype /Link /Border [0 0 0] /H /I /C [0 1 0] /Rect [93.647 345.534 111.617 354.402] /A << /S /GoTo /D <636974652E5256> >> >>}
\MusicPDFWriteObject{805}{<< /Type /Annot /Subtype /Link /Border [0 0 0] /H /I /C [1 0 0] /Rect [327.059 294.575 357.272 306.528] /A << /S /GoTo /D <6571756174696F6E2E332E312E312E32> >> >>}
\MusicPDFWriteObject{806}{<< /Type /Annot /Subtype /Link /Border [0 0 0] /H /I /C [1 0 0] /Rect [70.625 645.181 103.328 657.134] /A << /S /GoTo /D <74686D2E412E312E352E32> >> >>}
\MusicPDFWriteObject{807}{<< /Type /Annot /Subtype /Link /Border [0 0 0] /H /I /C [1 0 0] /Rect [112.727 595.003 145.43 607.318] /A << /S /GoTo /D <74686D2E412E312E352E32> >> >>}
\MusicPDFWriteObject{808}{<< /Type /Annot /Subtype /Link /Border [0 0 0] /H /I /C [1 0 0] /Rect [374.35 595.003 396.816 607.318] /A << /S /GoTo /D <73756273656374696F6E2E332E312E32> >> >>}
\MusicPDFWriteObject{809}{<< /Type /Annot /Subtype /Link /Border [0 0 0] /H /I /C [0 1 0] /Rect [425.761 505.798 443.233 514.666] /A << /S /GoTo /D <636974652E4346> >> >>}
\MusicPDFWriteObject{810}{<< /Type /Annot /Subtype /Link /Border [0 0 0] /H /I /C [0 1 0] /Rect [191.955 376.29 209.428 385.158] /A << /S /GoTo /D <636974652E4346> >> >>}
\MusicPDFWriteObject{811}{<< /Type /Annot /Subtype /Link /Border [0 0 0] /H /I /C [0 1 0] /Rect [292.31 376.29 310.28 385.158] /A << /S /GoTo /D <636974652E5256> >> >>}
\MusicPDFWriteObject{812}{<< /Type /Annot /Subtype /Link /Border [0 0 0] /H /I /C [0 1 0] /Rect [328.086 363.14 345.559 372.007] /A << /S /GoTo /D <636974652E4346> >> >>}
\MusicPDFWriteObject{813}{<< /Type /Annot /Subtype /Link /Border [0 0 0] /H /I /C [0 1 0] /Rect [402.957 349.989 420.429 358.857] /A << /S /GoTo /D <636974652E4346> >> >>}
\MusicPDFWriteObject{814}{<< /Type /Annot /Subtype /Link /Border [0 0 0] /H /I /C [1 0 0] /Rect [314.962 150.237 337.326 162.192] /A << /S /GoTo /D <73756273656374696F6E2E332E312E31> >> >>}
\MusicPDFWriteObject{815}{<< /Type /Annot /Subtype /Link /Border [0 0 0] /H /I /C [0 1 0] /Rect [525.797 595.995 543.27 604.863] /A << /S /GoTo /D <636974652E4346> >> >>}
\MusicPDFWriteObject{816}{<< /Type /Annot /Subtype /Link /Border [0 0 0] /H /I /C [1 0 0] /Rect [202.984 82.596 233.197 94.549] /A << /S /GoTo /D <6571756174696F6E2E332E312E322E39> >> >>}
\MusicPDFWriteObject{817}{<< /Type /Annot /Subtype /Link /Border [0 0 0] /H /I /C [0 1 0] /Rect [481.564 618.992 528.419 627.859] /A << /S /GoTo /D <636974652E5269656D616E6E> >> >>}
\MusicPDFWriteObject{818}{<< /Type /Annot /Subtype /Link /Border [0 0 0] /H /I /C [0 1 0] /Rect [195.374 579.539 212.846 588.407] /A << /S /GoTo /D <636974652E4346> >> >>}
\MusicPDFWriteObject{819}{<< /Type /Annot /Subtype /Link /Border [0 0 0] /H /I /C [0 1 0] /Rect [346.763 460.412 364.734 469.28] /A << /S /GoTo /D <636974652E5256> >> >>}
\MusicPDFWriteObject{820}{<< /Type /Annot /Subtype /Link /Border [0 0 0] /H /I /C [1 0 0] /Rect [219.106 411.529 249.32 423.481] /A << /S /GoTo /D <6571756174696F6E2E332E322E312E32> >> >>}
\MusicPDFWriteObject{821}{<< /Type /Annot /Subtype /Link /Border [0 0 0] /H /I /C [1 0 0] /Rect [353.741 358.846 383.954 370.799] /A << /S /GoTo /D <6571756174696F6E2E332E322E312E32> >> >>}
\MusicPDFWriteObject{822}{<< /Type /Annot /Subtype /Link /Border [0 0 0] /H /I /C [0 1 0] /Rect [285.585 209.538 303.555 218.406] /A << /S /GoTo /D <636974652E5256> >> >>}
\MusicPDFWriteObject{823}{<< /Type /Annot /Subtype /Link /Border [0 0 0] /H /I /C [0 1 0] /Rect [235.761 158.333 253.732 167.201] /A << /S /GoTo /D <636974652E5256> >> >>}
\MusicPDFWriteObject{824}{<< /Type /Annot /Subtype /Link /Border [0 0 0] /H /I /C [0 1 0] /Rect [393.214 158.333 411.184 167.201] /A << /S /GoTo /D <636974652E5256> >> >>}
\MusicPDFWriteObject{825}{<< /Type /Annot /Subtype /Link /Border [0 0 0] /H /I /C [0 1 0] /Rect [73.793 145.182 91.763 154.05] /A << /S /GoTo /D <636974652E5256> >> >>}
\MusicPDFWriteObject{826}{<< /Type /Annot /Subtype /Link /Border [0 0 0] /H /I /C [0 1 0] /Rect [150.091 132.032 168.062 140.899] /A << /S /GoTo /D <636974652E5256> >> >>}
\MusicPDFWriteObject{827}{<< /Type /Annot /Subtype /Link /Border [0 0 0] /H /I /C [1 0 0] /Rect [71.004 116.391 103.292 128.344] /A << /S /GoTo /D <74686D2E422E362E312E31> >> >>}
\MusicPDFWriteObject{828}{<< /Type /Annot /Subtype /Link /Border [0 0 0] /H /I /C [1 0 0] /Rect [159.453 632.555 190.23 644.87] /A << /S /GoTo /D <74686D2E332E322E312E33> >> >>}
\MusicPDFWriteObject{829}{<< /Type /Annot /Subtype /Link /Border [0 0 0] /H /I /C [1 0 0] /Rect [73.712 513.909 103.925 525.861] /A << /S /GoTo /D <6571756174696F6E2E332E322E312E34> >> >>}
\MusicPDFWriteObject{830}{<< /Type /Annot /Subtype /Link /Border [0 0 0] /H /I /C [1 0 0] /Rect [151.733 513.909 181.947 525.861] /A << /S /GoTo /D <6571756174696F6E2E332E322E312E35> >> >>}
\MusicPDFWriteObject{831}{<< /Type /Annot /Subtype /Link /Border [0 0 0] /H /I /C [0 1 0] /Rect [527.845 478.601 543.159 487.468] /A << /S /GoTo /D <636974652E494B> >> >>}
\MusicPDFWriteObject{832}{<< /Type /Annot /Subtype /Link /Border [0 0 0] /H /I /C [1 0 0] /Rect [259.529 462.96 294.723 474.912] /A << /S /GoTo /D <6571756174696F6E2E332E322E312E3130> >> >>}
\MusicPDFWriteObject{833}{<< /Type /Annot /Subtype /Link /Border [0 0 0] /H /I /C [1 0 0] /Rect [438.15 462.96 454.945 474.912] /A << /S /GoTo /D <73656374696F6E2E422E31> >> >>}
\MusicPDFWriteObject{834}{<< /Type /Annot /Subtype /Link /Border [0 0 0] /H /I /C [0 1 0] /Rect [300.531 439.148 324.81 448.016] /A << /S /GoTo /D <636974652E54617465> >> >>}
\MusicPDFWriteObject{835}{<< /Type /Annot /Subtype /Link /Border [0 0 0] /H /I /C [0 1 0] /Rect [458.618 439.148 477.826 448.016] /A << /S /GoTo /D <636974652E44656C> >> >>}
\MusicPDFWriteObject{836}{<< /Type /Annot /Subtype /Link /Border [0 0 0] /H /I /C [0 1 0] /Rect [124.227 424.061 176.762 434.865] /A << /S /GoTo /D <636974652E4C616E676C616E6473> >> >>}
\MusicPDFWriteObject{837}{<< /Type /Annot /Subtype /Link /Border [0 0 0] /H /I /C [1 0 0] /Rect [158.267 410.91 185.714 421.758] /A << /S /GoTo /D <73756273656374696F6E2E31342E312E31> >> >>}
\MusicPDFWriteObject{838}{<< /Type /Annot /Subtype /Link /Border [0 0 0] /H /I /C [0 1 0] /Rect [293.754 355.411 311.226 364.279] /A << /S /GoTo /D <636974652E4346> >> >>}
\MusicPDFWriteObject{839}{<< /Type /Annot /Subtype /Link /Border [0 0 0] /H /I /C [1 0 0] /Rect [125.081 109.057 155.301 121.012] /A << /S /GoTo /D <74686D2E332E322E322E32> >> >>}
\MusicPDFWriteObject{840}{<< /Type /Annot /Subtype /Link /Border [0 0 0] /H /I /C [1 0 0] /Rect [350.731 95.747 380.944 107.699] /A << /S /GoTo /D <6571756174696F6E2E332E322E322E33> >> >>}
\MusicPDFWriteObject{841}{<< /Type /Annot /Subtype /Link /Border [0 0 0] /H /I /C [1 0 0] /Rect [109.128 671.482 139.341 683.435] /A << /S /GoTo /D <6571756174696F6E2E332E322E322E33> >> >>}
\MusicPDFWriteObject{842}{<< /Type /Annot /Subtype /Link /Border [0 0 0] /H /I /C [1 0 0] /Rect [333.599 557.808 363.812 569.761] /A << /S /GoTo /D <6571756174696F6E2E332E322E322E34> >> >>}
\MusicPDFWriteObject{843}{<< /Type /Annot /Subtype /Link /Border [0 0 0] /H /I /C [1 0 0] /Rect [214.319 531.507 244.533 543.459] /A << /S /GoTo /D <6571756174696F6E2E332E322E322E34> >> >>}
\MusicPDFWriteObject{844}{<< /Type /Annot /Subtype /Link /Border [0 0 0] /H /I /C [1 0 0] /Rect [511.675 457.008 541.888 468.96] /A << /S /GoTo /D <6571756174696F6E2E332E322E322E31> >> >>}
\MusicPDFWriteObject{845}{<< /Type /Annot /Subtype /Link /Border [0 0 0] /H /I /C [1 0 0] /Rect [343.61 383.271 373.823 395.223] /A << /S /GoTo /D <6571756174696F6E2E332E322E322E31> >> >>}
\MusicPDFWriteObject{846}{<< /Type /Annot /Subtype /Link /Border [0 0 0] /H /I /C [0 1 0] /Rect [231.415 146.184 256.475 155.052] /A << /S /GoTo /D <636974652E53636861> >> >>}
\MusicPDFWriteObject{847}{<< /Type /Annot /Subtype /Link /Border [0 0 0] /H /I /C [1 0 0] /Rect [476.746 143.693 506.536 155.648] /A << /S /GoTo /D <74686D2E322E312E332E32> >> >>}
\MusicPDFWriteObject{848}{<< /Type /Annot /Subtype /Link /Border [0 0 0] /H /I /C [0 1 0] /Rect [179.93 133.033 204.99 141.901] /A << /S /GoTo /D <636974652E53636861> >> >>}
\MusicPDFWriteObject{849}{<< /Type /Annot /Subtype /Link /Border [0 0 0] /H /I /C [1 0 0] /Rect [112.727 102.496 142.94 116.214] /A << /S /GoTo /D <74686D2E322E322E332E34> >> >>}
\MusicPDFWriteObject{850}{<< /Type /Annot /Subtype /Link /Border [0 0 0] /H /I /C [1 0 0] /Rect [187.794 623.986 221.355 632.84] /A << /S /GoTo /D <74686D2E332E322E322E31> >> >>}
\MusicPDFWriteObject{851}{<< /Type /Annot /Subtype /Link /Border [0 0 0] /H /I /C [1 0 0] /Rect [301.324 621.496 331.791 633.449] /A << /S /GoTo /D <74686D2E322E322E322E31> >> >>}
\MusicPDFWriteObject{852}{<< /Type /Annot /Subtype /Link /Border [0 0 0] /H /I /C [0 1 0] /Rect [144.055 597.684 162.357 606.552] /A << /S /GoTo /D <636974652E4445> >> >>}
\MusicPDFWriteObject{853}{<< /Type /Annot /Subtype /Link /Border [0 0 0] /H /I /C [0 1 0] /Rect [386.074 584.534 403.547 593.401] /A << /S /GoTo /D <636974652E4346> >> >>}
\MusicPDFWriteObject{854}{<< /Type /Annot /Subtype /Link /Border [0 0 0] /H /I /C [0 1 0] /Rect [92.581 558.232 110.053 567.1] /A << /S /GoTo /D <636974652E4346> >> >>}
\MusicPDFWriteObject{855}{<< /Type /Annot /Subtype /Link /Border [0 0 0] /H /I /C [0 1 0] /Rect [335.619 545.081 353.092 553.949] /A << /S /GoTo /D <636974652E4346> >> >>}
\MusicPDFWriteObject{856}{<< /Type /Annot /Subtype /Link /Border [0 0 0] /H /I /C [1 0 0] /Rect [501.017 529.027 530.751 543.588] /A << /S /GoTo /D <74686D2E322E312E332E32> >> >>}
\MusicPDFWriteObject{857}{<< /Type /Annot /Subtype /Link /Border [0 0 0] /H /I /C [0 1 0] /Rect [89.402 458.476 107.373 467.344] /A << /S /GoTo /D <636974652E5256> >> >>}
\MusicPDFWriteObject{858}{<< /Type /Annot /Subtype /Link /Border [0 0 0] /H /I /C [1 0 0] /Rect [383.334 405.057 412.982 417.372] /A << /S /GoTo /D <74686D2E322E342E332E31> >> >>}
\MusicPDFWriteObject{859}{<< /Type /Annot /Subtype /Link /Border [0 0 0] /H /I /C [1 0 0] /Rect [262.893 362.755 293.106 374.708] /A << /S /GoTo /D <6571756174696F6E2E332E322E322E36> >> >>}
\MusicPDFWriteObject{860}{<< /Type /Annot /Subtype /Link /Border [0 0 0] /H /I /C [1 0 0] /Rect [223.662 336.453 253.875 348.406] /A << /S /GoTo /D <6571756174696F6E2E332E322E322E31> >> >>}
\MusicPDFWriteObject{861}{<< /Type /Annot /Subtype /Link /Border [0 0 0] /H /I /C [0 1 0] /Rect [525.687 325.793 543.159 334.661] /A << /S /GoTo /D <636974652E4346> >> >>}
\MusicPDFWriteObject{862}{<< /Type /Annot /Subtype /Link /Border [0 0 0] /H /I /C [0 1 0] /Rect [195.889 312.642 213.859 321.51] /A << /S /GoTo /D <636974652E5256> >> >>}
\MusicPDFWriteObject{863}{<< /Type /Annot /Subtype /Link /Border [0 0 0] /H /I /C [1 0 0] /Rect [491.285 270.539 521.498 282.492] /A << /S /GoTo /D <6571756174696F6E2E332E322E322E32> >> >>}
\MusicPDFWriteObject{864}{<< /Type /Annot /Subtype /Link /Border [0 0 0] /H /I /C [1 0 0] /Rect [124.429 243.639 154.642 257.316] /A << /S /GoTo /D <74686D2E322E322E332E34> >> >>}
\MusicPDFWriteObject{865}{<< /Type /Annot /Subtype /Link /Border [0 0 0] /H /I /C [1 0 0] /Rect [144.933 229.31 174.582 244.636] /A << /S /GoTo /D <74686D2E322E322E332E34> >> >>}
\MusicPDFWriteObject{866}{<< /Type /Annot /Subtype /Link /Border [0 0 0] /H /I /C [1 0 0] /Rect [213.065 117.44 243.278 129.393] /A << /S /GoTo /D <74686D2E322E322E332E34> >> >>}
\MusicPDFWriteObject{867}{<< /Type /Annot /Subtype /Link /Border [0 0 0] /H /I /C [1 0 0] /Rect [211.834 684.379 245.045 697.225] /A << /S /GoTo /D <74686D2E442E312E322E31> >> >>}
\MusicPDFWriteObject{868}{<< /Type /Annot /Subtype /Link /Border [0 0 0] /H /I /C [1 0 0] /Rect [324.113 671.482 356.755 684.075] /A << /S /GoTo /D <74686D2E412E332E322E31> >> >>}
\MusicPDFWriteObject{869}{<< /Type /Annot /Subtype /Link /Border [0 0 0] /H /I /C [1 0 0] /Rect [149.158 645.181 179.371 657.134] /A << /S /GoTo /D <6571756174696F6E2E332E322E322E37> >> >>}
\MusicPDFWriteObject{870}{<< /Type /Annot /Subtype /Link /Border [0 0 0] /H /I /C [1 0 0] /Rect [373.662 585.124 403.875 597.077] /A << /S /GoTo /D <6571756174696F6E2E332E322E322E39> >> >>}
\MusicPDFWriteObject{871}{<< /Type /Annot /Subtype /Link /Border [0 0 0] /H /I /C [1 0 0] /Rect [248.578 558.823 278.791 570.776] /A << /S /GoTo /D <6571756174696F6E2E332E322E322E39> >> >>}
\MusicPDFWriteObject{872}{<< /Type /Annot /Subtype /Link /Border [0 0 0] /H /I /C [1 0 0] /Rect [137.691 494.314 169.915 506.267] /A << /S /GoTo /D <74686D2E442E312E322E33> >> >>}
\MusicPDFWriteObject{873}{<< /Type /Annot /Subtype /Link /Border [0 0 0] /H /I /C [0 1 0] /Rect [210.93 496.804 229.232 505.672] /A << /S /GoTo /D <636974652E4445> >> >>}
\MusicPDFWriteObject{874}{<< /Type /Annot /Subtype /Link /Border [0 0 0] /H /I /C [1 0 0] /Rect [422.142 494.314 452.356 506.267] /A << /S /GoTo /D <6571756174696F6E2E332E322E322E33> >> >>}
\MusicPDFWriteObject{875}{<< /Type /Annot /Subtype /Link /Border [0 0 0] /H /I /C [1 0 0] /Rect [471.299 468.012 501.512 479.965] /A << /S /GoTo /D <6571756174696F6E2E332E322E322E39> >> >>}
\MusicPDFWriteObject{876}{<< /Type /Annot /Subtype /Link /Border [0 0 0] /H /I /C [1 0 0] /Rect [73.712 427.293 103.925 439.245] /A << /S /GoTo /D <6571756174696F6E2E332E322E322E33> >> >>}
\MusicPDFWriteObject{877}{<< /Type /Annot /Subtype /Link /Border [0 0 0] /H /I /C [1 0 0] /Rect [243.542 400.99 276.507 412.946] /A << /S /GoTo /D <74686D2E442E312E322E31> >> >>}
\MusicPDFWriteObject{878}{<< /Type /Annot /Subtype /Link /Border [0 0 0] /H /I /C [0 1 0] /Rect [320.707 403.481 339.01 412.349] /A << /S /GoTo /D <636974652E4445> >> >>}
\MusicPDFWriteObject{879}{<< /Type /Annot /Subtype /Link /Border [0 0 0] /H /I /C [1 0 0] /Rect [511.881 400.99 544.845 412.946] /A << /S /GoTo /D <74686D2E442E312E322E33> >> >>}
\MusicPDFWriteObject{880}{<< /Type /Annot /Subtype /Link /Border [0 0 0] /H /I /C [0 1 0] /Rect [112.557 390.33 130.859 399.198] /A << /S /GoTo /D <636974652E4445> >> >>}
\MusicPDFWriteObject{881}{<< /Type /Annot /Subtype /Link /Border [0 0 0] /H /I /C [1 0 0] /Rect [127.418 341.526 157.631 353.479] /A << /S /GoTo /D <6571756174696F6E2E332E322E322E35> >> >>}
\MusicPDFWriteObject{882}{<< /Type /Annot /Subtype /Link /Border [0 0 0] /H /I /C [1 0 0] /Rect [184.758 328.376 214.971 340.328] /A << /S /GoTo /D <6571756174696F6E2E332E322E322E31> >> >>}
\MusicPDFWriteObject{883}{<< /Type /Annot /Subtype /Link /Border [0 0 0] /H /I /C [1 0 0] /Rect [71.004 315.224 103.845 327.179] /A << /S /GoTo /D <74686D2E442E312E322E31> >> >>}
\MusicPDFWriteObject{884}{<< /Type /Annot /Subtype /Link /Border [0 0 0] /H /I /C [1 0 0] /Rect [182.756 227.26 212.969 239.213] /A << /S /GoTo /D <74686D2E332E322E322E31> >> >>}
\MusicPDFWriteObject{885}{<< /Type /Annot /Subtype /Link /Border [0 0 0] /H /I /C [1 0 0] /Rect [453.391 81.101 486.094 97.259] /A << /S /GoTo /D <74686D2E412E332E322E32> >> >>}
\MusicPDFWriteObject{886}{<< /Type /Annot /Subtype /Link /Border [0 0 0] /H /I /C [0 1 0] /Rect [172.153 596.446 187.841 605.314] /A << /S /GoTo /D <636974652E5431> >> >>}
\MusicPDFWriteObject{887}{<< /Type /Annot /Subtype /Link /Border [0 0 0] /H /I /C [0 1 0] /Rect [294.947 596.446 305.903 605.314] /A << /S /GoTo /D <636974652E48> >> >>}
\MusicPDFWriteObject{888}{<< /Type /Annot /Subtype /Link /Border [0 0 0] /H /I /C [1 0 0] /Rect [321.291 552.527 353.38 565.856] /A << /S /GoTo /D <74686D2E412E332E312E36> >> >>}
\MusicPDFWriteObject{889}{<< /Type /Annot /Subtype /Link /Border [0 0 0] /H /I /C [1 0 0] /Rect [363.966 497.975 394.433 509.928] /A << /S /GoTo /D <74686D2E332E322E322E31> >> >>}
\MusicPDFWriteObject{890}{<< /Type /Annot /Subtype /Link /Border [0 0 0] /H /I /C [1 0 0] /Rect [395.519 415.294 425.732 427.246] /A << /S /GoTo /D <6571756174696F6E2E332E322E322E32> >> >>}
\MusicPDFWriteObject{891}{<< /Type /Annot /Subtype /Link /Border [0 0 0] /H /I /C [1 0 0] /Rect [483.351 415.293 513 427.248] /A << /S /GoTo /D <74686D2E332E322E322E31> >> >>}
\MusicPDFWriteObject{892}{<< /Type /Annot /Subtype /Link /Border [0 0 0] /H /I /C [1 0 0] /Rect [491.794 254.109 524.988 264.9] /A << /S /GoTo /D <74686D2E412E332E312E38> >> >>}
\MusicPDFWriteObject{893}{<< /Type /Annot /Subtype /Link /Border [0 0 0] /H /I /C [1 0 0] /Rect [114.009 240.045 147.326 252.36] /A << /S /GoTo /D <74686D2E412E332E312E39> >> >>}
\MusicPDFWriteObject{894}{<< /Type /Annot /Subtype /Link /Border [0 0 0] /H /I /C [1 0 0] /Rect [141.823 214.657 174.987 225.448] /A << /S /GoTo /D <74686D2E412E332E342E31> >> >>}
\MusicPDFWriteObject{895}{<< /Type /Annot /Subtype /Link /Border [0 0 0] /H /I /C [1 0 0] /Rect [217.077 200.953 249.627 212.905] /A << /S /GoTo /D <74686D2E412E332E312E32> >> >>}
\MusicPDFWriteObject{896}{<< /Type /Annot /Subtype /Link /Border [0 0 0] /H /I /C [1 0 0] /Rect [403.421 187.442 436.124 199.756] /A << /S /GoTo /D <74686D2E412E332E342E32> >> >>}
\MusicPDFWriteObject{897}{<< /Type /Annot /Subtype /Link /Border [0 0 0] /H /I /C [1 0 0] /Rect [506.509 187.442 539.212 199.756] /A << /S /GoTo /D <74686D2E412E332E322E35> >> >>}
\MusicPDFWriteObject{898}{<< /Type /Annot /Subtype /Link /Border [0 0 0] /H /I /C [1 0 0] /Rect [71.004 135.199 103.845 147.152] /A << /S /GoTo /D <74686D2E442E312E322E31> >> >>}
\MusicPDFWriteObject{899}{<< /Type /Annot /Subtype /Link /Border [0 0 0] /H /I /C [1 0 0] /Rect [174.119 108.897 211.16 120.85] /A << /S /GoTo /D <74686D2E412E332E312E3131> >> >>}
\MusicPDFWriteObject{900}{<< /Type /Annot /Subtype /Link /Border [0 0 0] /H /I /C [0 1 0] /Rect [210.97 98.237 228.94 107.105] /A << /S /GoTo /D <636974652E5256> >> >>}
\MusicPDFWriteObject{901}{<< /Type /Annot /Subtype /Link /Border [0 0 0] /H /I /C [0 1 0] /Rect [422.39 98.237 440.36 107.105] /A << /S /GoTo /D <636974652E5256> >> >>}
\MusicPDFWriteObject{902}{<< /Type /Annot /Subtype /Link /Border [0 0 0] /H /I /C [0 1 0] /Rect [335.382 486.126 352.854 494.994] /A << /S /GoTo /D <636974652E4346> >> >>}
\MusicPDFWriteObject{903}{<< /Type /Annot /Subtype /Link /Border [0 0 0] /H /I /C [1 0 0] /Rect [511.675 247.082 541.888 259.035] /A << /S /GoTo /D <6571756174696F6E2E332E332E312E31> >> >>}
\MusicPDFWriteObject{904}{<< /Type /Annot /Subtype /Link /Border [0 0 0] /H /I /C [1 0 0] /Rect [348.778 211.199 379.556 223.151] /A << /S /GoTo /D <74686D2E332E332E312E31> >> >>}
\MusicPDFWriteObject{905}{<< /Type /Annot /Subtype /Link /Border [0 0 0] /H /I /C [1 0 0] /Rect [448.336 211.199 479.113 223.151] /A << /S /GoTo /D <74686D2E332E322E322E31> >> >>}
\MusicPDFWriteObject{906}{<< /Type /Annot /Subtype /Link /Border [0 0 0] /H /I /C [1 0 0] /Rect [187.794 151.504 221.355 160.358] /A << /S /GoTo /D <74686D2E332E332E312E31> >> >>}
\MusicPDFWriteObject{907}{<< /Type /Annot /Subtype /Link /Border [0 0 0] /H /I /C [1 0 0] /Rect [297.443 148.653 327.966 160.968] /A << /S /GoTo /D <74686D2E322E312E342E31> >> >>}
\MusicPDFWriteObject{908}{<< /Type /Annot /Subtype /Link /Border [0 0 0] /H /I /C [0 1 0] /Rect [131.052 687.123 149.023 695.991] /A << /S /GoTo /D <636974652E5256> >> >>}
\MusicPDFWriteObject{909}{<< /Type /Annot /Subtype /Link /Border [0 0 0] /H /I /C [0 1 0] /Rect [270.511 625.455 287.984 634.323] /A << /S /GoTo /D <636974652E4346> >> >>}
\MusicPDFWriteObject{910}{<< /Type /Annot /Subtype /Link /Border [0 0 0] /H /I /C [1 0 0] /Rect [290.347 596.664 320.56 608.617] /A << /S /GoTo /D <6571756174696F6E2E332E312E322E34> >> >>}
\MusicPDFWriteObject{911}{<< /Type /Annot /Subtype /Link /Border [0 0 0] /H /I /C [1 0 0] /Rect [364.08 596.304 394.293 609.079] /A << /S /GoTo /D <74686D2E332E312E322E31> >> >>}
\MusicPDFWriteObject{912}{<< /Type /Annot /Subtype /Link /Border [0 0 0] /H /I /C [1 0 0] /Rect [453.307 596.304 483.52 609.079] /A << /S /GoTo /D <74686D2E332E312E322E32> >> >>}
\MusicPDFWriteObject{913}{<< /Type /Annot /Subtype /Link /Border [0 0 0] /H /I /C [1 0 0] /Rect [394.375 524.64 427.601 537.055] /A << /S /GoTo /D <74686D2E412E332E322E32> >> >>}
\MusicPDFWriteObject{914}{<< /Type /Annot /Subtype /Link /Border [0 0 0] /H /I /C [1 0 0] /Rect [511.675 432.584 541.888 444.537] /A << /S /GoTo /D <6571756174696F6E2E332E332E312E37> >> >>}
\MusicPDFWriteObject{915}{<< /Type /Annot /Subtype /Link /Border [0 0 0] /H /I /C [1 0 0] /Rect [73.712 339.644 103.925 351.596] /A << /S /GoTo /D <6571756174696F6E2E332E332E312E31> >> >>}
\MusicPDFWriteObject{916}{<< /Type /Annot /Subtype /Link /Border [0 0 0] /H /I /C [1 0 0] /Rect [278.918 313.895 309.47 324.686] /A << /S /GoTo /D <74686D2E332E332E312E31> >> >>}
\MusicPDFWriteObject{917}{<< /Type /Annot /Subtype /Link /Border [0 0 0] /H /I /C [1 0 0] /Rect [246.085 229.733 276.834 242.048] /A << /S /GoTo /D <74686D2E332E322E322E33> >> >>}
\MusicPDFWriteObject{918}{<< /Type /Annot /Subtype /Link /Border [0 0 0] /H /I /C [1 0 0] /Rect [487.11 229.733 517.86 242.048] /A << /S /GoTo /D <74686D2E322E312E342E31> >> >>}
\MusicPDFWriteObject{919}{<< /Type /Annot /Subtype /Link /Border [0 0 0] /H /I /C [1 0 0] /Rect [203.354 216.942 233.567 228.895] /A << /S /GoTo /D <6571756174696F6E2E332E332E312E31> >> >>}
\MusicPDFWriteObject{920}{<< /Type /Annot /Subtype /Link /Border [0 0 0] /H /I /C [1 0 0] /Rect [431.693 177.13 462.471 189.445] /A << /S /GoTo /D <74686D2E332E332E312E31> >> >>}
\MusicPDFWriteObject{921}{<< /Type /Annot /Subtype /Link /Border [0 0 0] /H /I /C [1 0 0] /Rect [233.654 163.979 266.218 176.294] /A << /S /GoTo /D <74686D2E442E312E322E31> >> >>}
\MusicPDFWriteObject{922}{<< /Type /Annot /Subtype /Link /Border [0 0 0] /H /I /C [1 0 0] /Rect [130.543 117.292 161.321 129.245] /A << /S /GoTo /D <74686D2E322E342E322E31> >> >>}
\MusicPDFWriteObject{923}{<< /Type /Annot /Subtype /Link /Border [0 0 0] /H /I /C [1 0 0] /Rect [238.162 117.292 268.939 129.245] /A << /S /GoTo /D <74686D2E322E332E312E32> >> >>}
\MusicPDFWriteObject{924}{<< /Type /Annot /Subtype /Link /Border [0 0 0] /H /I /C [1 0 0] /Rect [303.192 611.509 333.969 624.49] /A << /S /GoTo /D <74686D2E332E332E312E31> >> >>}
\MusicPDFWriteObject{925}{<< /Type /Annot /Subtype /Link /Border [0 0 0] /H /I /C [1 0 0] /Rect [152.119 599.385 182.332 611.338] /A << /S /GoTo /D <6571756174696F6E2E332E332E312E32> >> >>}
\MusicPDFWriteObject{926}{<< /Type /Annot /Subtype /Link /Border [0 0 0] /H /I /C [1 0 0] /Rect [282.249 599.385 312.462 611.338] /A << /S /GoTo /D <6571756174696F6E2E332E332E312E33> >> >>}
\MusicPDFWriteObject{927}{<< /Type /Annot /Subtype /Link /Border [0 0 0] /H /I /C [0 1 0] /Rect [272.945 588.564 290.418 597.432] /A << /S /GoTo /D <636974652E4346> >> >>}
\MusicPDFWriteObject{928}{<< /Type /Annot /Subtype /Link /Border [0 0 0] /H /I /C [1 0 0] /Rect [437.809 533.311 468.023 545.264] /A << /S /GoTo /D <6571756174696F6E2E332E332E312E35> >> >>}
\MusicPDFWriteObject{929}{<< /Type /Annot /Subtype /Link /Border [0 0 0] /H /I /C [1 0 0] /Rect [212.169 519.8 242.946 532.575] /A << /S /GoTo /D <74686D2E322E332E312E32> >> >>}
\MusicPDFWriteObject{930}{<< /Type /Annot /Subtype /Link /Border [0 0 0] /H /I /C [1 0 0] /Rect [401.606 519.8 432.384 532.575] /A << /S /GoTo /D <74686D2E322E342E332E31> >> >>}
\MusicPDFWriteObject{931}{<< /Type /Annot /Subtype /Link /Border [0 0 0] /H /I /C [1 0 0] /Rect [112.727 507.009 142.94 518.962] /A << /S /GoTo /D <74686D2E322E342E322E31> >> >>}
\MusicPDFWriteObject{932}{<< /Type /Annot /Subtype /Link /Border [0 0 0] /H /I /C [1 0 0] /Rect [124.563 480.708 154.776 492.661] /A << /S /GoTo /D <6571756174696F6E2E332E332E312E35> >> >>}
\MusicPDFWriteObject{933}{<< /Type /Annot /Subtype /Link /Border [0 0 0] /H /I /C [1 0 0] /Rect [380.113 480.347 410.891 492.662] /A << /S /GoTo /D <74686D2E322E332E312E34> >> >>}
\MusicPDFWriteObject{934}{<< /Type /Annot /Subtype /Link /Border [0 0 0] /H /I /C [1 0 0] /Rect [269.049 467.197 299.77 479.972] /A << /S /GoTo /D <74686D2E322E342E332E31> >> >>}
\MusicPDFWriteObject{935}{<< /Type /Annot /Subtype /Link /Border [0 0 0] /H /I /C [1 0 0] /Rect [149.164 454.406 179.377 466.359] /A << /S /GoTo /D <74686D2E322E342E322E31> >> >>}
\MusicPDFWriteObject{936}{<< /Type /Annot /Subtype /Link /Border [0 0 0] /H /I /C [1 0 0] /Rect [157.598 441.256 187.812 453.208] /A << /S /GoTo /D <6571756174696F6E2E332E332E312E36> >> >>}
\MusicPDFWriteObject{937}{<< /Type /Annot /Subtype /Link /Border [0 0 0] /H /I /C [1 0 0] /Rect [151.266 428.105 181.48 440.057] /A << /S /GoTo /D <6571756174696F6E2E332E332E312E34> >> >>}
\MusicPDFWriteObject{938}{<< /Type /Annot /Subtype /Link /Border [0 0 0] /H /I /C [1 0 0] /Rect [296.104 428.105 328.946 440.057] /A << /S /GoTo /D <74686D2E442E312E322E31> >> >>}
\MusicPDFWriteObject{939}{<< /Type /Annot /Subtype /Link /Border [0 0 0] /H /I /C [1 0 0] /Rect [88.936 411.1 94.899 421.468] /A << /S /GoTo /D <48666F6F746E6F74652E39> >> >>}
\MusicPDFWriteObject{940}{<< /Type /Annot /Subtype /Link /Border [0 0 0] /H /I /C [1 0 0] /Rect [515.296 270.006 544.945 282.321] /A << /S /GoTo /D <74686D2E322E312E322E34> >> >>}
\MusicPDFWriteObject{941}{<< /Type /Annot /Subtype /Link /Border [0 0 0] /H /I /C [1 0 0] /Rect [143.99 257.215 176.416 269.168] /A << /S /GoTo /D <74686D2E432E322E352E31> >> >>}
\MusicPDFWriteObject{942}{<< /Type /Annot /Subtype /Link /Border [0 0 0] /H /I /C [1 0 0] /Rect [180.924 671.122 211.646 683.437] /A << /S /GoTo /D <74686D2E332E332E312E33> >> >>}
\MusicPDFWriteObject{943}{<< /Type /Annot /Subtype /Link /Border [0 0 0] /H /I /C [0 1 0] /Rect [301.474 633.856 320.448 642.724] /A << /S /GoTo /D <636974652E547265> >> >>}
\MusicPDFWriteObject{944}{<< /Type /Annot /Subtype /Link /Border [0 0 0] /H /I /C [1 0 0] /Rect [478.389 616.72 510.737 630.17] /A << /S /GoTo /D <74686D2E442E332E362E31> >> >>}
\MusicPDFWriteObject{945}{<< /Type /Annot /Subtype /Link /Border [0 0 0] /H /I /C [1 0 0] /Rect [394.223 532.019 424.436 543.972] /A << /S /GoTo /D <6571756174696F6E2E332E332E312E33> >> >>}
\MusicPDFWriteObject{946}{<< /Type /Annot /Subtype /Link /Border [0 0 0] /H /I /C [1 0 0] /Rect [115.548 518.868 145.761 530.821] /A << /S /GoTo /D <6571756174696F6E2E332E332E312E32> >> >>}
\MusicPDFWriteObject{947}{<< /Type /Annot /Subtype /Link /Border [0 0 0] /H /I /C [1 0 0] /Rect [216.589 518.868 246.802 530.821] /A << /S /GoTo /D <74686D2E332E332E312E31> >> >>}
\MusicPDFWriteObject{948}{<< /Type /Annot /Subtype /Link /Border [0 0 0] /H /I /C [1 0 0] /Rect [417.2 473.419 447.413 485.372] /A << /S /GoTo /D <6571756174696F6E2E332E332E312E33> >> >>}
\MusicPDFWriteObject{949}{<< /Type /Annot /Subtype /Link /Border [0 0 0] /H /I /C [1 0 0] /Rect [250.381 261.448 280.594 273.401] /A << /S /GoTo /D <6571756174696F6E2E332E332E312E33> >> >>}
\MusicPDFWriteObject{950}{<< /Type /Annot /Subtype /Link /Border [0 0 0] /H /I /C [0 1 0] /Rect [205.322 250.787 222.794 259.655] /A << /S /GoTo /D <636974652E4346> >> >>}
\MusicPDFWriteObject{951}{<< /Type /Annot /Subtype /Link /Border [0 0 0] /H /I /C [0 1 0] /Rect [107.335 237.637 124.808 246.504] /A << /S /GoTo /D <636974652E4346> >> >>}
\MusicPDFWriteObject{952}{<< /Type /Annot /Subtype /Link /Border [0 0 0] /H /I /C [1 0 0] /Rect [169.056 132.479 199.27 144.432] /A << /S /GoTo /D <6571756174696F6E2E332E332E312E38> >> >>}
\MusicPDFWriteObject{953}{<< /Type /Annot /Subtype /Link /Border [0 0 0] /H /I /C [1 0 0] /Rect [488.547 118.587 518.76 130.54] /A << /S /GoTo /D <6571756174696F6E2E332E332E312E38> >> >>}
\MusicPDFWriteObject{954}{<< /Type /Annot /Subtype /Link /Border [0 0 0] /H /I /C [1 0 0] /Rect [243.006 684.633 278.199 696.586] /A << /S /GoTo /D <6571756174696F6E2E332E332E312E3134> >> >>}
\MusicPDFWriteObject{955}{<< /Type /Annot /Subtype /Link /Border [0 0 0] /H /I /C [1 0 0] /Rect [309.121 684.633 339.334 696.586] /A << /S /GoTo /D <6571756174696F6E2E332E332E312E39> >> >>}
\MusicPDFWriteObject{956}{<< /Type /Annot /Subtype /Link /Border [0 0 0] /H /I /C [1 0 0] /Rect [124.734 592.627 159.927 604.58] /A << /S /GoTo /D <6571756174696F6E2E332E332E312E3135> >> >>}
\MusicPDFWriteObject{957}{<< /Type /Annot /Subtype /Link /Border [0 0 0] /H /I /C [0 1 0] /Rect [341.5 525.733 358.972 534.601] /A << /S /GoTo /D <636974652E4346> >> >>}
\MusicPDFWriteObject{958}{<< /Type /Annot /Subtype /Link /Border [0 0 0] /H /I /C [0 1 0] /Rect [138.465 624.394 155.937 633.262] /A << /S /GoTo /D <636974652E4346> >> >>}
\MusicPDFWriteObject{959}{<< /Type /Annot /Subtype /Link /Border [0 0 0] /H /I /C [0 1 0] /Rect [257.742 551.655 275.215 560.523] /A << /S /GoTo /D <636974652E4346> >> >>}
\MusicPDFWriteObject{960}{<< /Type /Annot /Subtype /Link /Border [0 0 0] /H /I /C [0 1 0] /Rect [420.211 551.655 437.683 560.523] /A << /S /GoTo /D <636974652E4346> >> >>}
\MusicPDFWriteObject{961}{<< /Type /Annot /Subtype /Link /Border [0 0 0] /H /I /C [1 0 0] /Rect [277.013 536.014 312.206 547.967] /A << /S /GoTo /D <6571756174696F6E2E332E332E312E3139> >> >>}
\MusicPDFWriteObject{962}{<< /Type /Annot /Subtype /Link /Border [0 0 0] /H /I /C [1 0 0] /Rect [73.712 278.613 108.905 290.566] /A << /S /GoTo /D <6571756174696F6E2E332E332E312E3230> >> >>}
\MusicPDFWriteObject{963}{<< /Type /Annot /Subtype /Link /Border [0 0 0] /H /I /C [1 0 0] /Rect [458.226 265.463 493.419 277.415] /A << /S /GoTo /D <6571756174696F6E2E332E332E312E3230> >> >>}
\MusicPDFWriteObject{964}{<< /Type /Annot /Subtype /Link /Border [0 0 0] /H /I /C [1 0 0] /Rect [283.853 252.312 319.046 264.265] /A << /S /GoTo /D <6571756174696F6E2E332E332E312E3231> >> >>}
\MusicPDFWriteObject{965}{<< /Type /Annot /Subtype /Link /Border [0 0 0] /H /I /C [1 0 0] /Rect [193.365 189.158 224.029 201.113] /A << /S /GoTo /D <74686D2E332E332E312E35> >> >>}
\MusicPDFWriteObject{966}{<< /Type /Annot /Subtype /Link /Border [0 0 0] /H /I /C [0 1 0] /Rect [73.826 526.424 92.128 535.292] /A << /S /GoTo /D <636974652E4445> >> >>}
\MusicPDFWriteObject{967}{<< /Type /Annot /Subtype /Link /Border [0 0 0] /H /I /C [0 1 0] /Rect [314.559 526.424 332.862 535.292] /A << /S /GoTo /D <636974652E4445> >> >>}
\MusicPDFWriteObject{968}{<< /Type /Annot /Subtype /Link /Border [0 0 0] /H /I /C [1 0 0] /Rect [224.059 446.689 254.272 457.48] /A << /S /GoTo /D <74686D2E332E332E312E31> >> >>}
\MusicPDFWriteObject{969}{<< /Type /Annot /Subtype /Link /Border [0 0 0] /H /I /C [1 0 0] /Rect [267.408 520.328 297.621 532.281] /A << /S /GoTo /D <6571756174696F6E2E332E332E312E33> >> >>}
\MusicPDFWriteObject{970}{<< /Type /Annot /Subtype /Link /Border [0 0 0] /H /I /C [1 0 0] /Rect [70.755 429.032 101.137 442.26] /A << /S /GoTo /D <74686D2E332E332E312E37> >> >>}
\MusicPDFWriteObject{971}{<< /Type /Annot /Subtype /Link /Border [0 0 0] /H /I /C [1 0 0] /Rect [419.666 429.032 452.553 442.26] /A << /S /GoTo /D <74686D2E412E332E312E36> >> >>}
\MusicPDFWriteObject{972}{<< /Type /Annot /Subtype /Link /Border [0 0 0] /H /I /C [1 0 0] /Rect [322.558 415.881 355.875 428.556] /A << /S /GoTo /D <74686D2E412E332E312E36> >> >>}
\MusicPDFWriteObject{973}{<< /Type /Annot /Subtype /Link /Border [0 0 0] /H /I /C [1 0 0] /Rect [169.739 404.004 202.964 415.957] /A << /S /GoTo /D <74686D2E412E332E312E35> >> >>}
\MusicPDFWriteObject{974}{<< /Type /Annot /Subtype /Link /Border [0 0 0] /H /I /C [1 0 0] /Rect [438.204 404.004 471.429 415.957] /A << /S /GoTo /D <74686D2E412E332E322E35> >> >>}
\MusicPDFWriteObject{975}{<< /Type /Annot /Subtype /Link /Border [0 0 0] /H /I /C [1 0 0] /Rect [432.69 391.406 464.779 402.197] /A << /S /GoTo /D <74686D2E412E332E312E32> >> >>}
\MusicPDFWriteObject{976}{<< /Type /Annot /Subtype /Link /Border [0 0 0] /H /I /C [0 1 0] /Rect [296.768 353.891 314.241 362.759] /A << /S /GoTo /D <636974652E4346> >> >>}
\MusicPDFWriteObject{977}{<< /Type /Annot /Subtype /Link /Border [0 0 0] /H /I /C [1 0 0] /Rect [176.323 322.635 209.884 331.489] /A << /S /GoTo /D <74686D2E332E332E312E37> >> >>}
\MusicPDFWriteObject{978}{<< /Type /Annot /Subtype /Link /Border [0 0 0] /H /I /C [1 0 0] /Rect [150.209 264.201 180.987 276.157] /A << /S /GoTo /D <74686D2E322E342E332E32> >> >>}
\MusicPDFWriteObject{979}{<< /Type /Annot /Subtype /Link /Border [0 0 0] /H /I /C [1 0 0] /Rect [306.45 264.201 337.227 276.157] /A << /S /GoTo /D <74686D2E322E332E312E34> >> >>}
\MusicPDFWriteObject{980}{<< /Type /Annot /Subtype /Link /Border [0 0 0] /H /I /C [1 0 0] /Rect [495.867 251.051 527.549 263.004] /A << /S /GoTo /D <74686D2E422E322E312E31> >> >>}
\MusicPDFWriteObject{981}{<< /Type /Annot /Subtype /Link /Border [0 0 0] /H /I /C [1 0 0] /Rect [273.583 162.02 305.871 173.975] /A << /S /GoTo /D <74686D2E422E322E312E35> >> >>}
\MusicPDFWriteObject{982}{<< /Type /Annot /Subtype /Link /Border [0 0 0] /H /I /C [1 0 0] /Rect [191.345 135.199 221.559 147.152] /A << /S /GoTo /D <6571756174696F6E2E332E332E312E32> >> >>}
\MusicPDFWriteObject{983}{<< /Type /Annot /Subtype /Link /Border [0 0 0] /H /I /C [1 0 0] /Rect [292.248 135.198 322.405 149.506] /A << /S /GoTo /D <74686D2E332E332E312E31> >> >>}
\MusicPDFWriteObject{984}{<< /Type /Annot /Subtype /Link /Border [0 0 0] /H /I /C [1 0 0] /Rect [201.367 120.996 231.213 134.003] /A << /S /GoTo /D <74686D2E332E332E312E31> >> >>}
\MusicPDFWriteObject{985}{<< /Type /Annot /Subtype /Link /Border [0 0 0] /H /I /C [1 0 0] /Rect [297.402 122.048 327.615 134.001] /A << /S /GoTo /D <6571756174696F6E2E332E332E312E32> >> >>}
\MusicPDFWriteObject{986}{<< /Type /Annot /Subtype /Link /Border [0 0 0] /H /I /C [1 0 0] /Rect [317.485 108.897 347.698 120.85] /A << /S /GoTo /D <6571756174696F6E2E332E332E312E33> >> >>}
\MusicPDFWriteObject{987}{<< /Type /Annot /Subtype /Link /Border [0 0 0] /H /I /C [0 1 0] /Rect [107.377 98.237 124.85 107.105] /A << /S /GoTo /D <636974652E4346> >> >>}
\MusicPDFWriteObject{988}{<< /Type /Annot /Subtype /Link /Border [0 0 0] /H /I /C [0 1 0] /Rect [364.876 98.237 382.348 107.105] /A << /S /GoTo /D <636974652E4346> >> >>}
\MusicPDFWriteObject{989}{<< /Type /Annot /Subtype /Link /Border [0 0 0] /H /I /C [1 0 0] /Rect [447.905 684.633 478.118 696.586] /A << /S /GoTo /D <6571756174696F6E2E332E332E312E32> >> >>}
\MusicPDFWriteObject{990}{<< /Type /Annot /Subtype /Link /Border [0 0 0] /H /I /C [1 0 0] /Rect [108.553 610.362 143.746 622.315] /A << /S /GoTo /D <6571756174696F6E2E332E332E312E3333> >> >>}
\MusicPDFWriteObject{991}{<< /Type /Annot /Subtype /Link /Border [0 0 0] /H /I /C [1 0 0] /Rect [538.832 480.763 544.795 493.981] /A << /S /GoTo /D <48666F6F746E6F74652E3130> >> >>}
\MusicPDFWriteObject{992}{<< /Type /Annot /Subtype /Link /Border [0 0 0] /H /I /C [1 0 0] /Rect [479.83 467.973 515.024 479.926] /A << /S /GoTo /D <6571756174696F6E2E332E332E312E3237> >> >>}
\MusicPDFWriteObject{993}{<< /Type /Annot /Subtype /Link /Border [0 0 0] /H /I /C [0 1 0] /Rect [262.35 457.312 277.665 466.18] /A << /S /GoTo /D <636974652E494B> >> >>}
\MusicPDFWriteObject{994}{<< /Type /Annot /Subtype /Link /Border [0 0 0] /H /I /C [1 0 0] /Rect [338.382 441.672 368.595 453.624] /A << /S /GoTo /D <6571756174696F6E2E332E332E312E32> >> >>}
\MusicPDFWriteObject{995}{<< /Type /Annot /Subtype /Link /Border [0 0 0] /H /I /C [1 0 0] /Rect [400.483 441.672 435.677 453.624] /A << /S /GoTo /D <6571756174696F6E2E332E332E312E3333> >> >>}
\MusicPDFWriteObject{996}{<< /Type /Annot /Subtype /Link /Border [0 0 0] /H /I /C [0 1 0] /Rect [525.227 444.162 543.198 453.029] /A << /S /GoTo /D <636974652E5256> >> >>}
\MusicPDFWriteObject{997}{<< /Type /Annot /Subtype /Link /Border [0 0 0] /H /I /C [1 0 0] /Rect [207.808 324.217 240.511 336.632] /A << /S /GoTo /D <74686D2E412E332E322E32> >> >>}
\MusicPDFWriteObject{998}{<< /Type /Annot /Subtype /Link /Border [0 0 0] /H /I /C [0 1 0] /Rect [186.242 313.557 203.714 322.424] /A << /S /GoTo /D <636974652E4346> >> >>}
\MusicPDFWriteObject{999}{<< /Type /Annot /Subtype /Link /Border [0 0 0] /H /I /C [0 1 0] /Rect [337.279 313.557 354.751 322.424] /A << /S /GoTo /D <636974652E4346> >> >>}
\MusicPDFWriteObject{1000}{<< /Type /Annot /Subtype /Link /Border [0 0 0] /H /I /C [1 0 0] /Rect [408.971 297.915 439.185 309.87] /A << /S /GoTo /D <74686D2E332E332E312E31> >> >>}
\MusicPDFWriteObject{1001}{<< /Type /Annot /Subtype /Link /Border [0 0 0] /H /I /C [1 0 0] /Rect [153.756 271.614 183.97 283.567] /A << /S /GoTo /D <6571756174696F6E2E332E332E322E31> >> >>}
\MusicPDFWriteObject{1002}{<< /Type /Annot /Subtype /Link /Border [0 0 0] /H /I /C [1 0 0] /Rect [156.803 166.408 181.482 178.361] /A << /S /GoTo /D <73756273656374696F6E2E432E352E32> >> >>}
\MusicPDFWriteObject{1003}{<< /Type /Annot /Subtype /Link /Border [0 0 0] /H /I /C [1 0 0] /Rect [206.588 166.408 231.267 178.361] /A << /S /GoTo /D <73756273656374696F6E2E432E352E33> >> >>}
\MusicPDFWriteObject{1004}{<< /Type /Annot /Subtype /Link /Border [0 0 0] /H /I /C [0 1 0] /Rect [364.725 129.446 401.343 138.314] /A << /S /GoTo /D <636974652E5361726E616B> >> >>}
\MusicPDFWriteObject{1005}{<< /Type /Annot /Subtype /Link /Border [0 0 0] /H /I /C [1 0 0] /Rect [507.046 672.035 529.922 682.884] /A << /S /GoTo /D <73756273656374696F6E2E332E312E31> >> >>}
\MusicPDFWriteObject{1006}{<< /Type /Annot /Subtype /Link /Border [0 0 0] /H /I /C [1 0 0] /Rect [180.224 592.577 202.281 604.532] /A << /S /GoTo /D <73756273656374696F6E2E332E312E31> >> >>}
\MusicPDFWriteObject{1007}{<< /Type /Annot /Subtype /Link /Border [0 0 0] /H /I /C [0 1 0] /Rect [87.503 582.527 104.975 591.395] /A << /S /GoTo /D <636974652E4346> >> >>}
\MusicPDFWriteObject{1008}{<< /Type /Annot /Subtype /Link /Border [0 0 0] /H /I /C [1 0 0] /Rect [260.626 566.886 295.819 578.839] /A << /S /GoTo /D <6571756174696F6E2E332E332E322E3131> >> >>}
\MusicPDFWriteObject{1009}{<< /Type /Annot /Subtype /Link /Border [0 0 0] /H /I /C [1 0 0] /Rect [118.215 527.434 153.408 539.387] /A << /S /GoTo /D <6571756174696F6E2E332E332E322E3131> >> >>}
\MusicPDFWriteObject{1010}{<< /Type /Annot /Subtype /Link /Border [0 0 0] /H /I /C [1 0 0] /Rect [510.29 474.83 543.173 486.785] /A << /S /GoTo /D <74686D2E432E352E332E33> >> >>}
\MusicPDFWriteObject{1011}{<< /Type /Annot /Subtype /Link /Border [0 0 0] /H /I /C [0 1 0] /Rect [354.828 437.869 372.3 446.737] /A << /S /GoTo /D <636974652E4346> >> >>}
\MusicPDFWriteObject{1012}{<< /Type /Annot /Subtype /Link /Border [0 0 0] /H /I /C [0 1 0] /Rect [269.591 398.416 287.063 407.284] /A << /S /GoTo /D <636974652E4346> >> >>}
\MusicPDFWriteObject{1013}{<< /Type /Annot /Subtype /Link /Border [0 0 0] /H /I /C [0 1 0] /Rect [139.62 293.203 156.774 302.071] /A << /S /GoTo /D <636974652E466F6C> >> >>}
\MusicPDFWriteObject{1014}{<< /Type /Annot /Subtype /Link /Border [0 0 0] /H /I /C [0 1 0] /Rect [245.53 293.203 262.684 302.071] /A << /S /GoTo /D <636974652E466F6C> >> >>}
\MusicPDFWriteObject{1015}{<< /Type /Annot /Subtype /Link /Border [0 0 0] /H /I /C [1 0 0] /Rect [164.237 178.096 194.45 190.048] /A << /S /GoTo /D <6571756174696F6E2E332E332E322E36> >> >>}
\MusicPDFWriteObject{1016}{<< /Type /Annot /Subtype /Link /Border [0 0 0] /H /I /C [1 0 0] /Rect [149.44 507.348 179.653 518.138] /A << /S /GoTo /D <74686D2E332E332E322E31> >> >>}
\MusicPDFWriteObject{1017}{<< /Type /Annot /Subtype /Link /Border [0 0 0] /H /I /C [0 1 0] /Rect [88.195 374.969 105.667 383.836] /A << /S /GoTo /D <636974652E4346> >> >>}
\MusicPDFWriteObject{1018}{<< /Type /Annot /Subtype /Link /Border [0 0 0] /H /I /C [0 1 0] /Rect [310.922 361.818 328.395 370.686] /A << /S /GoTo /D <636974652E4346> >> >>}
\MusicPDFWriteObject{1019}{<< /Type /Annot /Subtype /Link /Border [0 0 0] /H /I /C [1 0 0] /Rect [152.152 159.742 187.345 171.695] /A << /S /GoTo /D <6571756174696F6E2E332E332E322E3235> >> >>}
\MusicPDFWriteObject{1020}{<< /Type /Annot /Subtype /Link /Border [0 0 0] /H /I /C [1 0 0] /Rect [332.404 159.742 367.597 171.695] /A << /S /GoTo /D <6571756174696F6E2E332E332E322E3233> >> >>}
\MusicPDFWriteObject{1021}{<< /Type /Annot /Subtype /Link /Border [0 0 0] /H /I /C [1 0 0] /Rect [345.13 95.493 375.005 107.701] /A << /S /GoTo /D <74686D2E332E322E322E31> >> >>}
\MusicPDFWriteObject{1022}{<< /Type /Annot /Subtype /Link /Border [0 0 0] /H /I /C [0 1 0] /Rect [379.406 687.123 413.285 695.991] /A << /S /GoTo /D <636974652E5765696C2D61> >> >>}
\MusicPDFWriteObject{1023}{<< /Type /Annot /Subtype /Link /Border [0 0 0] /H /I /C [1 0 0] /Rect [240.497 672.036 270.146 682.827] /A << /S /GoTo /D <74686D2E332E322E322E35> >> >>}
\MusicPDFWriteObject{1024}{<< /Type /Annot /Subtype /Link /Border [0 0 0] /H /I /C [0 1 0] /Rect [117.291 660.822 151.17 669.69] /A << /S /GoTo /D <636974652E5765696C2D61> >> >>}
\MusicPDFWriteObject{1025}{<< /Type /Annot /Subtype /Link /Border [0 0 0] /H /I /C [1 0 0] /Rect [210.975 468.61 227.609 480.563] /A << /S /GoTo /D <73656374696F6E2E432E35> >> >>}
\MusicPDFWriteObject{1026}{<< /Type /Annot /Subtype /Link /Border [0 0 0] /H /I /C [1 0 0] /Rect [344.532 468.61 361.166 480.563] /A << /S /GoTo /D <73656374696F6E2E432E35> >> >>}
\MusicPDFWriteObject{1027}{<< /Type /Annot /Subtype /Link /Border [0 0 0] /H /I /C [1 0 0] /Rect [487.347 468.61 519.165 480.563] /A << /S /GoTo /D <74686D2E432E352E322E36> >> >>}
\MusicPDFWriteObject{1028}{<< /Type /Annot /Subtype /Link /Border [0 0 0] /H /I /C [1 0 0] /Rect [266.6 403.409 281.32 414.257] /A << /S /GoTo /D <73656374696F6E2E352E35> >> >>}
\MusicPDFWriteObject{1029}{<< /Type /Annot /Subtype /Link /Border [0 0 0] /H /I /C [0 1 0] /Rect [319.333 392.195 353.212 401.063] /A << /S /GoTo /D <636974652E5765696C2D61> >> >>}
\MusicPDFWriteObject{1030}{<< /Type /Annot /Subtype /Link /Border [0 0 0] /H /I /C [0 1 0] /Rect [525.731 304.471 543.203 313.338] /A << /S /GoTo /D <636974652E4346> >> >>}
\MusicPDFWriteObject{1031}{<< /Type /Annot /Subtype /Link /Border [0 0 0] /H /I /C [1 0 0] /Rect [193.354 194.796 215.82 205.644] /A << /S /GoTo /D <73756273656374696F6E2E332E312E31> >> >>}
\MusicPDFWriteObject{1032}{<< /Type /Annot /Subtype /Link /Border [0 0 0] /H /I /C [1 0 0] /Rect [394.961 145.963 425.174 157.916] /A << /S /GoTo /D <6571756174696F6E2E332E322E322E31> >> >>}
\MusicPDFWriteObject{1033}{<< /Type /Annot /Subtype /Link /Border [0 0 0] /H /I /C [1 0 0] /Rect [484.4 145.963 514.98 157.916] /A << /S /GoTo /D <74686D2E332E322E322E31> >> >>}
\MusicPDFWriteObject{1034}{<< /Type /Annot /Subtype /Link /Border [0 0 0] /H /I /C [1 0 0] /Rect [282.575 132.452 312.788 144.767] /A << /S /GoTo /D <74686D2E332E322E312E34> >> >>}
\MusicPDFWriteObject{1035}{<< /Type /Annot /Subtype /Link /Border [0 0 0] /H /I /C [1 0 0] /Rect [510.373 119.662 540.021 131.614] /A << /S /GoTo /D <74686D2E332E342E332E33> >> >>}
\MusicPDFWriteObject{1036}{<< /Type /Annot /Subtype /Link /Border [0 0 0] /H /I /C [1 0 0] /Rect [241.403 107.064 271.56 117.855] /A << /S /GoTo /D <74686D2E332E342E312E31> >> >>}
\MusicPDFWriteObject{1037}{<< /Type /Annot /Subtype /Link /Border [0 0 0] /H /I /C [1 0 0] /Rect [163.274 257.647 193.487 268.495] /A << /S /GoTo /D <74686D2E352E312E312E31> >> >>}
\MusicPDFWriteObject{1038}{<< /Type /Annot /Subtype /Link /Border [0 0 0] /H /I /C [1 0 0] /Rect [305.568 257.647 328.034 268.495] /A << /S /GoTo /D <73756273656374696F6E2E352E312E31> >> >>}
\MusicPDFWriteObject{1039}{<< /Type /Annot /Subtype /Link /Border [0 0 0] /H /I /C [1 0 0] /Rect [499.372 214.845 529.585 226.798] /A << /S /GoTo /D <6571756174696F6E2E332E342E312E32> >> >>}
\MusicPDFWriteObject{1040}{<< /Type /Annot /Subtype /Link /Border [0 0 0] /H /I /C [1 0 0] /Rect [219.607 530.98 252.925 541.771] /A << /S /GoTo /D <74686D2E412E332E322E32> >> >>}
\MusicPDFWriteObject{1041}{<< /Type /Annot /Subtype /Link /Border [0 0 0] /H /I /C [1 0 0] /Rect [304.297 317.371 334.51 329.324] /A << /S /GoTo /D <6571756174696F6E2E332E342E312E31> >> >>}
\MusicPDFWriteObject{1042}{<< /Type /Annot /Subtype /Link /Border [0 0 0] /H /I /C [0 1 0] /Rect [290.056 280.309 307.528 289.177] /A << /S /GoTo /D <636974652E4346> >> >>}
\MusicPDFWriteObject{1043}{<< /Type /Annot /Subtype /Link /Border [0 0 0] /H /I /C [1 0 0] /Rect [163.093 251.517 193.307 263.472] /A << /S /GoTo /D <74686D2E332E312E322E31> >> >>}
\MusicPDFWriteObject{1044}{<< /Type /Annot /Subtype /Link /Border [0 0 0] /H /I /C [1 0 0] /Rect [408.944 188.788 439.157 200.741] /A << /S /GoTo /D <6571756174696F6E2E332E342E312E35> >> >>}
\MusicPDFWriteObject{1045}{<< /Type /Annot /Subtype /Link /Border [0 0 0] /H /I /C [1 0 0] /Rect [467.485 188.788 497.698 200.741] /A << /S /GoTo /D <6571756174696F6E2E332E342E312E36> >> >>}
\MusicPDFWriteObject{1046}{<< /Type /Annot /Subtype /Link /Border [0 0 0] /H /I /C [1 0 0] /Rect [141.8 173.59 147.763 188.495] /A << /S /GoTo /D <48666F6F746E6F74652E3131> >> >>}
\MusicPDFWriteObject{1047}{<< /Type /Annot /Subtype /Link /Border [0 0 0] /H /I /C [1 0 0] /Rect [511.675 134.089 541.888 146.042] /A << /S /GoTo /D <6571756174696F6E2E332E342E312E37> >> >>}
\MusicPDFWriteObject{1048}{<< /Type /Annot /Subtype /Link /Border [0 0 0] /H /I /C [1 0 0] /Rect [213.391 120.938 243.604 132.891] /A << /S /GoTo /D <6571756174696F6E2E332E342E312E35> >> >>}
\MusicPDFWriteObject{1049}{<< /Type /Annot /Subtype /Link /Border [0 0 0] /H /I /C [1 0 0] /Rect [270.84 120.938 301.053 132.891] /A << /S /GoTo /D <6571756174696F6E2E332E342E312E36> >> >>}
\MusicPDFWriteObject{1050}{<< /Type /Annot /Subtype /Link /Border [0 0 0] /H /I /C [1 0 0] /Rect [113.516 672.036 144.294 682.827] /A << /S /GoTo /D <74686D2E332E342E312E31> >> >>}
\MusicPDFWriteObject{1051}{<< /Type /Annot /Subtype /Link /Border [0 0 0] /H /I /C [1 0 0] /Rect [276.704 658.332 306.917 670.284] /A << /S /GoTo /D <6571756174696F6E2E332E322E322E31> >> >>}
\MusicPDFWriteObject{1052}{<< /Type /Annot /Subtype /Link /Border [0 0 0] /H /I /C [1 0 0] /Rect [376.153 658.332 405.802 670.284] /A << /S /GoTo /D <74686D2E332E322E322E31> >> >>}
\MusicPDFWriteObject{1053}{<< /Type /Annot /Subtype /Link /Border [0 0 0] /H /I /C [1 0 0] /Rect [499.926 465.456 530.139 477.409] /A << /S /GoTo /D <6571756174696F6E2E332E342E322E31> >> >>}
\MusicPDFWriteObject{1054}{<< /Type /Annot /Subtype /Link /Border [0 0 0] /H /I /C [1 0 0] /Rect [298.393 452.305 328.042 464.26] /A << /S /GoTo /D <74686D2E322E312E332E32> >> >>}
\MusicPDFWriteObject{1055}{<< /Type /Annot /Subtype /Link /Border [0 0 0] /H /I /C [0 1 0] /Rect [107.335 415.011 125.306 423.879] /A << /S /GoTo /D <636974652E5256> >> >>}
\MusicPDFWriteObject{1056}{<< /Type /Annot /Subtype /Link /Border [0 0 0] /H /I /C [1 0 0] /Rect [335.467 412.521 365.68 424.474] /A << /S /GoTo /D <6571756174696F6E2E332E342E322E31> >> >>}
\MusicPDFWriteObject{1057}{<< /Type /Annot /Subtype /Link /Border [0 0 0] /H /I /C [1 0 0] /Rect [185.174 361.955 215.387 373.908] /A << /S /GoTo /D <6571756174696F6E2E332E342E322E32> >> >>}
\MusicPDFWriteObject{1058}{<< /Type /Annot /Subtype /Link /Border [0 0 0] /H /I /C [1 0 0] /Rect [98.422 335.653 128.635 347.606] /A << /S /GoTo /D <6571756174696F6E2E332E342E322E31> >> >>}
\MusicPDFWriteObject{1059}{<< /Type /Annot /Subtype /Link /Border [0 0 0] /H /I /C [1 0 0] /Rect [295.334 322.503 325.547 334.455] /A << /S /GoTo /D <6571756174696F6E2E332E342E322E32> >> >>}
\MusicPDFWriteObject{1060}{<< /Type /Annot /Subtype /Link /Border [0 0 0] /H /I /C [1 0 0] /Rect [332.249 253.423 362.462 265.376] /A << /S /GoTo /D <6571756174696F6E2E332E342E312E36> >> >>}
\MusicPDFWriteObject{1061}{<< /Type /Annot /Subtype /Link /Border [0 0 0] /H /I /C [1 0 0] /Rect [436.561 253.063 467.339 265.376] /A << /S /GoTo /D <74686D2E332E342E312E31> >> >>}
\MusicPDFWriteObject{1062}{<< /Type /Annot /Subtype /Link /Border [0 0 0] /H /I /C [1 0 0] /Rect [238.497 148.403 270.179 160.356] /A << /S /GoTo /D <74686D2E422E372E322E35> >> >>}
\MusicPDFWriteObject{1063}{<< /Type /Annot /Subtype /Link /Border [0 0 0] /H /I /C [0 1 0] /Rect [459.249 150.893 474.936 159.761] /A << /S /GoTo /D <636974652E5432> >> >>}
\MusicPDFWriteObject{1064}{<< /Type /Annot /Subtype /Link /Border [0 0 0] /H /I /C [1 0 0] /Rect [306.606 83.149 336.255 93.94] /A << /S /GoTo /D <74686D2E332E342E312E31> >> >>}
\MusicPDFWriteObject{1065}{<< /Type /Annot /Subtype /Link /Border [0 0 0] /H /I /C [1 0 0] /Rect [205.106 684.633 234.755 697.048] /A << /S /GoTo /D <74686D2E322E322E352E34> >> >>}
\MusicPDFWriteObject{1066}{<< /Type /Annot /Subtype /Link /Border [0 0 0] /H /I /C [1 0 0] /Rect [334.568 658.885 364.244 669.676] /A << /S /GoTo /D <74686D2E322E322E352E34> >> >>}
\MusicPDFWriteObject{1067}{<< /Type /Annot /Subtype /Link /Border [0 0 0] /H /I /C [1 0 0] /Rect [124.888 632.03 155.214 643.985] /A << /S /GoTo /D <74686D2E322E322E352E34> >> >>}
\MusicPDFWriteObject{1068}{<< /Type /Annot /Subtype /Link /Border [0 0 0] /H /I /C [1 0 0] /Rect [306.438 605.729 336.086 617.681] /A << /S /GoTo /D <74686D2E322E322E352E34> >> >>}
\MusicPDFWriteObject{1069}{<< /Type /Annot /Subtype /Link /Border [0 0 0] /H /I /C [1 0 0] /Rect [112.727 568.766 142.94 577.62] /A << /S /GoTo /D <74686D2E322E322E352E34> >> >>}
\MusicPDFWriteObject{1070}{<< /Type /Annot /Subtype /Link /Border [0 0 0] /H /I /C [1 0 0] /Rect [222.615 275.98 252.828 287.933] /A << /S /GoTo /D <6571756174696F6E2E332E342E322E33> >> >>}
\MusicPDFWriteObject{1071}{<< /Type /Annot /Subtype /Link /Border [0 0 0] /H /I /C [1 0 0] /Rect [282.342 275.98 312.555 287.933] /A << /S /GoTo /D <6571756174696F6E2E332E342E322E34> >> >>}
\MusicPDFWriteObject{1072}{<< /Type /Annot /Subtype /Link /Border [0 0 0] /H /I /C [1 0 0] /Rect [431.755 242.405 461.573 254.718] /A << /S /GoTo /D <74686D2E322E322E352E34> >> >>}
\MusicPDFWriteObject{1073}{<< /Type /Annot /Subtype /Link /Border [0 0 0] /H /I /C [1 0 0] /Rect [337.296 216.464 367.509 228.417] /A << /S /GoTo /D <6571756174696F6E2E322E312E342E35> >> >>}
\MusicPDFWriteObject{1074}{<< /Type /Annot /Subtype /Link /Border [0 0 0] /H /I /C [1 0 0] /Rect [319.761 187.802 349.974 199.755] /A << /S /GoTo /D <6571756174696F6E2E332E342E312E36> >> >>}
\MusicPDFWriteObject{1075}{<< /Type /Annot /Subtype /Link /Border [0 0 0] /H /I /C [1 0 0] /Rect [419.948 187.802 449.738 199.755] /A << /S /GoTo /D <74686D2E332E342E312E31> >> >>}
\MusicPDFWriteObject{1076}{<< /Type /Annot /Subtype /Link /Border [0 0 0] /H /I /C [1 0 0] /Rect [158.716 109.45 175.51 120.298] /A << /S /GoTo /D <73656374696F6E2E422E31> >> >>}
\MusicPDFWriteObject{1077}{<< /Type /Annot /Subtype /Link /Border [0 0 0] /H /I /C [1 0 0] /Rect [312.407 109.45 344.695 120.298] /A << /S /GoTo /D <74686D2E422E312E312E33> >> >>}
\MusicPDFWriteObject{1078}{<< /Type /Annot /Subtype /Link /Border [0 0 0] /H /I /C [1 0 0] /Rect [393.641 95.747 423.91 107.699] /A << /S /GoTo /D <74686D2E322E322E332E34> >> >>}
\MusicPDFWriteObject{1079}{<< /Type /Annot /Subtype /Link /Border [0 0 0] /H /I /C [1 0 0] /Rect [449.114 95.747 479.327 107.699] /A << /S /GoTo /D <6571756174696F6E2E322E322E332E34> >> >>}
\MusicPDFWriteObject{1080}{<< /Type /Annot /Subtype /Link /Border [0 0 0] /H /I /C [1 0 0] /Rect [507.794 95.747 538.007 107.699] /A << /S /GoTo /D <6571756174696F6E2E322E322E332E35> >> >>}
\MusicPDFWriteObject{1081}{<< /Type /Annot /Subtype /Link /Border [0 0 0] /H /I /C [1 0 0] /Rect [348.62 671.122 378.297 683.437] /A << /S /GoTo /D <74686D2E322E322E352E34> >> >>}
\MusicPDFWriteObject{1082}{<< /Type /Annot /Subtype /Link /Border [0 0 0] /H /I /C [1 0 0] /Rect [391.608 632.584 422.386 643.374] /A << /S /GoTo /D <74686D2E322E312E332E32> >> >>}
\MusicPDFWriteObject{1083}{<< /Type /Annot /Subtype /Link /Border [0 0 0] /H /I /C [1 0 0] /Rect [213.139 579.427 242.788 592.019] /A << /S /GoTo /D <74686D2E322E322E352E34> >> >>}
\MusicPDFWriteObject{1084}{<< /Type /Annot /Subtype /Link /Border [0 0 0] /H /I /C [1 0 0] /Rect [187.794 463.515 221.355 472.369] /A << /S /GoTo /D <74686D2E332E342E322E32> >> >>}
\MusicPDFWriteObject{1085}{<< /Type /Annot /Subtype /Link /Border [0 0 0] /H /I /C [1 0 0] /Rect [283.336 461.024 312.984 472.979] /A << /S /GoTo /D <74686D2E332E322E322E31> >> >>}
\MusicPDFWriteObject{1086}{<< /Type /Annot /Subtype /Link /Border [0 0 0] /H /I /C [1 0 0] /Rect [244.595 433.279 274.808 445.231] /A << /S /GoTo /D <6571756174696F6E2E332E312E312E31> >> >>}
\MusicPDFWriteObject{1087}{<< /Type /Annot /Subtype /Link /Border [0 0 0] /H /I /C [1 0 0] /Rect [114.005 406.723 144.783 418.932] /A << /S /GoTo /D <74686D2E332E322E322E31> >> >>}
\MusicPDFWriteObject{1088}{<< /Type /Annot /Subtype /Link /Border [0 0 0] /H /I /C [1 0 0] /Rect [429.534 346.217 461.823 358.17] /A << /S /GoTo /D <6571756174696F6E2E422E322E312E32> >> >>}
\MusicPDFWriteObject{1089}{<< /Type /Annot /Subtype /Link /Border [0 0 0] /H /I /C [1 0 0] /Rect [144.988 332.347 174.636 344.662] /A << /S /GoTo /D <74686D2E332E322E322E31> >> >>}
\MusicPDFWriteObject{1090}{<< /Type /Annot /Subtype /Link /Border [0 0 0] /H /I /C [1 0 0] /Rect [146.106 306.405 176.009 318.361] /A << /S /GoTo /D <74686D2E332E322E322E31> >> >>}
\MusicPDFWriteObject{1091}{<< /Type /Annot /Subtype /Link /Border [0 0 0] /H /I /C [1 0 0] /Rect [183.221 187.226 213.434 199.179] /A << /S /GoTo /D <6571756174696F6E2E332E342E312E36> >> >>}
\MusicPDFWriteObject{1092}{<< /Type /Annot /Subtype /Link /Border [0 0 0] /H /I /C [1 0 0] /Rect [283.829 187.226 313.844 199.179] /A << /S /GoTo /D <74686D2E332E342E312E31> >> >>}
\MusicPDFWriteObject{1093}{<< /Type /Annot /Subtype /Link /Border [0 0 0] /H /I /C [1 0 0] /Rect [256.813 97.079 287.026 109.031] /A << /S /GoTo /D <6571756174696F6E2E332E342E312E33> >> >>}
\MusicPDFWriteObject{1094}{<< /Type /Annot /Subtype /Link /Border [0 0 0] /H /I /C [1 0 0] /Rect [349.475 684.633 379.688 696.586] /A << /S /GoTo /D <6571756174696F6E2E332E342E322E37> >> >>}
\MusicPDFWriteObject{1095}{<< /Type /Annot /Subtype /Link /Border [0 0 0] /H /I /C [1 0 0] /Rect [296.907 671.482 329.528 683.435] /A << /S /GoTo /D <74686D2E422E332E322E31> >> >>}
\MusicPDFWriteObject{1096}{<< /Type /Annot /Subtype /Link /Border [0 0 0] /H /I /C [1 0 0] /Rect [366.014 645.181 398.302 657.134] /A << /S /GoTo /D <74686D2E422E332E322E31> >> >>}
\MusicPDFWriteObject{1097}{<< /Type /Annot /Subtype /Link /Border [0 0 0] /H /I /C [1 0 0] /Rect [112.727 569.047 145.015 581.256] /A << /S /GoTo /D <74686D2E422E332E322E31> >> >>}
\MusicPDFWriteObject{1098}{<< /Type /Annot /Subtype /Link /Border [0 0 0] /H /I /C [1 0 0] /Rect [71.004 506.572 103.898 518.525] /A << /S /GoTo /D <74686D2E422E332E312E31> >> >>}
\MusicPDFWriteObject{1099}{<< /Type /Annot /Subtype /Link /Border [0 0 0] /H /I /C [1 0 0] /Rect [125.679 427.285 158.573 441.482] /A << /S /GoTo /D <74686D2E422E312E312E33> >> >>}
\MusicPDFWriteObject{1100}{<< /Type /Annot /Subtype /Link /Border [0 0 0] /H /I /C [1 0 0] /Rect [158.937 414.135 189.15 426.088] /A << /S /GoTo /D <6571756174696F6E2E332E342E312E33> >> >>}
\MusicPDFWriteObject{1101}{<< /Type /Annot /Subtype /Link /Border [0 0 0] /H /I /C [0 1 0] /Rect [418.541 416.625 436.013 425.493] /A << /S /GoTo /D <636974652E4346> >> >>}
\MusicPDFWriteObject{1102}{<< /Type /Annot /Subtype /Link /Border [0 0 0] /H /I /C [0 1 0] /Rect [73.715 403.474 91.187 412.342] /A << /S /GoTo /D <636974652E4346> >> >>}
\MusicPDFWriteObject{1103}{<< /Type /Annot /Subtype /Link /Border [0 0 0] /H /I /C [1 0 0] /Rect [365.979 362 397.661 372.876] /A << /S /GoTo /D <74686D2E422E322E322E33> >> >>}
\MusicPDFWriteObject{1104}{<< /Type /Annot /Subtype /Link /Border [0 0 0] /H /I /C [1 0 0] /Rect [112.001 348.381 141.678 360.334] /A << /S /GoTo /D <74686D2E322E332E312E34> >> >>}
\MusicPDFWriteObject{1105}{<< /Type /Annot /Subtype /Link /Border [0 0 0] /H /I /C [1 0 0] /Rect [214.769 348.381 244.446 360.334] /A << /S /GoTo /D <74686D2E322E332E322E32> >> >>}
\MusicPDFWriteObject{1106}{<< /Type /Annot /Subtype /Link /Border [0 0 0] /H /I /C [1 0 0] /Rect [412.333 348.381 442.546 360.334] /A << /S /GoTo /D <6571756174696F6E2E332E332E312E33> >> >>}
\MusicPDFWriteObject{1107}{<< /Type /Annot /Subtype /Link /Border [0 0 0] /H /I /C [1 0 0] /Rect [512.285 348.381 541.962 360.334] /A << /S /GoTo /D <74686D2E332E332E312E31> >> >>}
\MusicPDFWriteObject{1108}{<< /Type /Annot /Subtype /Link /Border [0 0 0] /H /I /C [1 0 0] /Rect [378.589 333.62 408.435 347.185] /A << /S /GoTo /D <74686D2E332E332E312E35> >> >>}
\MusicPDFWriteObject{1109}{<< /Type /Annot /Subtype /Link /Border [0 0 0] /H /I /C [1 0 0] /Rect [439.485 155.923 471.709 167.876] /A << /S /GoTo /D <74686D2E442E332E312E31> >> >>}
\MusicPDFWriteObject{1110}{<< /Type /Annot /Subtype /Link /Border [0 0 0] /H /I /C [1 0 0] /Rect [255.054 141.162 287.039 154.727] /A << /S /GoTo /D <74686D2E422E322E322E33> >> >>}
\MusicPDFWriteObject{1111}{<< /Type /Annot /Subtype /Link /Border [0 0 0] /H /I /C [1 0 0] /Rect [198.47 116.471 228.683 128.424] /A << /S /GoTo /D <6571756174696F6E2E332E342E322E37> >> >>}
\MusicPDFWriteObject{1112}{<< /Type /Annot /Subtype /Link /Border [0 0 0] /H /I /C [1 0 0] /Rect [172.153 684.633 202.366 696.586] /A << /S /GoTo /D <6571756174696F6E2E332E342E322E37> >> >>}
\MusicPDFWriteObject{1113}{<< /Type /Annot /Subtype /Link /Border [0 0 0] /H /I /C [1 0 0] /Rect [186.52 615.544 216.734 627.497] /A << /S /GoTo /D <6571756174696F6E2E332E342E322E38> >> >>}
\MusicPDFWriteObject{1114}{<< /Type /Annot /Subtype /Link /Border [0 0 0] /H /I /C [1 0 0] /Rect [98.422 602.393 128.635 614.346] /A << /S /GoTo /D <6571756174696F6E2E332E342E322E38> >> >>}
\MusicPDFWriteObject{1115}{<< /Type /Annot /Subtype /Link /Border [0 0 0] /H /I /C [1 0 0] /Rect [433.184 511.21 468.377 523.163] /A << /S /GoTo /D <6571756174696F6E2E332E342E322E3130> >> >>}
\MusicPDFWriteObject{1116}{<< /Type /Annot /Subtype /Link /Border [0 0 0] /H /I /C [1 0 0] /Rect [373.863 484.909 406.757 496.862] /A << /S /GoTo /D <74686D2E422E332E322E31> >> >>}
\MusicPDFWriteObject{1117}{<< /Type /Annot /Subtype /Link /Border [0 0 0] /H /I /C [1 0 0] /Rect [179.564 397.605 214.757 409.558] /A << /S /GoTo /D <6571756174696F6E2E332E342E322E3131> >> >>}
\MusicPDFWriteObject{1118}{<< /Type /Annot /Subtype /Link /Border [0 0 0] /H /I /C [1 0 0] /Rect [442.463 397.245 474.146 409.56] /A << /S /GoTo /D <74686D2E422E322E312E37> >> >>}
\MusicPDFWriteObject{1119}{<< /Type /Annot /Subtype /Link /Border [0 0 0] /H /I /C [1 0 0] /Rect [71.004 382.373 102.837 398.241] /A << /S /GoTo /D <74686D2E422E322E312E38> >> >>}
\MusicPDFWriteObject{1120}{<< /Type /Annot /Subtype /Link /Border [0 0 0] /H /I /C [1 0 0] /Rect [323.532 328.731 358.725 340.684] /A << /S /GoTo /D <6571756174696F6E2E332E342E322E3132> >> >>}
\MusicPDFWriteObject{1121}{<< /Type /Annot /Subtype /Link /Border [0 0 0] /H /I /C [1 0 0] /Rect [126.127 140.262 156.763 152.215] /A << /S /GoTo /D <74686D2E322E322E342E33> >> >>}
\MusicPDFWriteObject{1122}{<< /Type /Annot /Subtype /Link /Border [0 0 0] /H /I /C [1 0 0] /Rect [488.278 140.262 521.582 152.215] /A << /S /GoTo /D <74686D2E442E332E342E32> >> >>}
\MusicPDFWriteObject{1123}{<< /Type /Annot /Subtype /Link /Border [0 0 0] /H /I /C [1 0 0] /Rect [467.883 127.665 498.661 138.455] /A << /S /GoTo /D <74686D2E332E342E322E32> >> >>}
\MusicPDFWriteObject{1124}{<< /Type /Annot /Subtype /Link /Border [0 0 0] /H /I /C [1 0 0] /Rect [112.727 116.45 142.94 125.304] /A << /S /GoTo /D <74686D2E322E322E332E34> >> >>}
\MusicPDFWriteObject{1125}{<< /Type /Annot /Subtype /Link /Border [0 0 0] /H /I /C [1 0 0] /Rect [267.198 611.618 297.411 623.57] /A << /S /GoTo /D <6571756174696F6E2E332E342E322E36> >> >>}
\MusicPDFWriteObject{1126}{<< /Type /Annot /Subtype /Link /Border [0 0 0] /H /I /C [1 0 0] /Rect [134.346 598.141 164.559 610.094] /A << /S /GoTo /D <6571756174696F6E2E332E342E322E36> >> >>}
\MusicPDFWriteObject{1127}{<< /Type /Annot /Subtype /Link /Border [0 0 0] /H /I /C [1 0 0] /Rect [418.787 572.393 449.565 583.184] /A << /S /GoTo /D <74686D2E332E342E312E31> >> >>}
\MusicPDFWriteObject{1128}{<< /Type /Annot /Subtype /Link /Border [0 0 0] /H /I /C [1 0 0] /Rect [181.474 559.242 211.969 570.033] /A << /S /GoTo /D <74686D2E322E322E332E34> >> >>}
\MusicPDFWriteObject{1129}{<< /Type /Annot /Subtype /Link /Border [0 0 0] /H /I /C [1 0 0] /Rect [70.755 506.086 100.968 518.038] /A << /S /GoTo /D <74686D2E322E322E352E34> >> >>}
\MusicPDFWriteObject{1130}{<< /Type /Annot /Subtype /Link /Border [0 0 0] /H /I /C [1 0 0] /Rect [272.532 506.086 302.745 518.038] /A << /S /GoTo /D <6571756174696F6E2E332E322E322E32> >> >>}
\MusicPDFWriteObject{1131}{<< /Type /Annot /Subtype /Link /Border [0 0 0] /H /I /C [1 0 0] /Rect [385.667 469.498 420.86 481.451] /A << /S /GoTo /D <6571756174696F6E2E332E342E322E3130> >> >>}
\MusicPDFWriteObject{1132}{<< /Type /Annot /Subtype /Link /Border [0 0 0] /H /I /C [1 0 0] /Rect [449.194 469.498 484.387 481.451] /A << /S /GoTo /D <6571756174696F6E2E332E342E322E3132> >> >>}
\MusicPDFWriteObject{1133}{<< /Type /Annot /Subtype /Link /Border [0 0 0] /H /I /C [1 0 0] /Rect [73.712 456.348 108.905 468.3] /A << /S /GoTo /D <6571756174696F6E2E332E342E322E3133> >> >>}
\MusicPDFWriteObject{1134}{<< /Type /Annot /Subtype /Link /Border [0 0 0] /H /I /C [1 0 0] /Rect [137.346 456.348 172.54 468.3] /A << /S /GoTo /D <6571756174696F6E2E332E342E322E3134> >> >>}
\MusicPDFWriteObject{1135}{<< /Type /Annot /Subtype /Link /Border [0 0 0] /H /I /C [0 1 0] /Rect [438.425 394.411 456.727 403.278] /A << /S /GoTo /D <636974652E4445> >> >>}
\MusicPDFWriteObject{1136}{<< /Type /Annot /Subtype /Link /Border [0 0 0] /H /I /C [1 0 0] /Rect [215.06 352.468 250.253 364.421] /A << /S /GoTo /D <6571756174696F6E2E332E342E322E3135> >> >>}
\MusicPDFWriteObject{1137}{<< /Type /Annot /Subtype /Link /Border [0 0 0] /H /I /C [1 0 0] /Rect [279.895 352.468 315.088 364.421] /A << /S /GoTo /D <6571756174696F6E2E332E342E322E3136> >> >>}
\MusicPDFWriteObject{1138}{<< /Type /Annot /Subtype /Link /Border [0 0 0] /H /I /C [1 0 0] /Rect [390.123 352.468 425.316 364.421] /A << /S /GoTo /D <6571756174696F6E2E332E342E322E3130> >> >>}
\MusicPDFWriteObject{1139}{<< /Type /Annot /Subtype /Link /Border [0 0 0] /H /I /C [1 0 0] /Rect [437.763 352.468 472.956 364.421] /A << /S /GoTo /D <6571756174696F6E2E332E342E322E3132> >> >>}
\MusicPDFWriteObject{1140}{<< /Type /Annot /Subtype /Link /Border [0 0 0] /H /I /C [1 0 0] /Rect [485.402 352.468 520.595 364.421] /A << /S /GoTo /D <6571756174696F6E2E332E342E322E3133> >> >>}
\MusicPDFWriteObject{1141}{<< /Type /Annot /Subtype /Link /Border [0 0 0] /H /I /C [1 0 0] /Rect [73.712 339.318 108.905 351.27] /A << /S /GoTo /D <6571756174696F6E2E332E342E322E3134> >> >>}
\MusicPDFWriteObject{1142}{<< /Type /Annot /Subtype /Link /Border [0 0 0] /H /I /C [1 0 0] /Rect [338.246 185.196 368.882 197.149] /A << /S /GoTo /D <74686D2E322E322E352E34> >> >>}
\MusicPDFWriteObject{1143}{<< /Type /Annot /Subtype /Link /Border [0 0 0] /H /I /C [1 0 0] /Rect [242.873 158.534 273.086 170.238] /A << /S /GoTo /D <74686D2E322E322E352E34> >> >>}
\MusicPDFWriteObject{1144}{<< /Type /Annot /Subtype /Link /Border [0 0 0] /H /I /C [1 0 0] /Rect [312.625 133.146 342.273 143.937] /A << /S /GoTo /D <74686D2E322E322E332E34> >> >>}
\MusicPDFWriteObject{1145}{<< /Type /Annot /Subtype /Link /Border [0 0 0] /H /I /C [1 0 0] /Rect [403.514 133.146 433.163 143.937] /A << /S /GoTo /D <74686D2E322E322E352E34> >> >>}
\MusicPDFWriteObject{1146}{<< /Type /Annot /Subtype /Link /Border [0 0 0] /H /I /C [1 0 0] /Rect [294.373 119.442 324.981 131.395] /A << /S /GoTo /D <74686D2E322E342E322E31> >> >>}
\MusicPDFWriteObject{1147}{<< /Type /Annot /Subtype /Link /Border [0 0 0] /H /I /C [1 0 0] /Rect [400.192 119.442 430.8 131.395] /A << /S /GoTo /D <74686D2E322E332E312E32> >> >>}
\MusicPDFWriteObject{1148}{<< /Type /Annot /Subtype /Link /Border [0 0 0] /H /I /C [1 0 0] /Rect [128.701 106.291 159.478 118.244] /A << /S /GoTo /D <74686D2E322E322E332E34> >> >>}
\MusicPDFWriteObject{1149}{<< /Type /Annot /Subtype /Link /Border [0 0 0] /H /I /C [1 0 0] /Rect [227.085 106.291 257.862 118.244] /A << /S /GoTo /D <74686D2E322E322E352E34> >> >>}
\MusicPDFWriteObject{1150}{<< /Type /Annot /Subtype /Link /Border [0 0 0] /H /I /C [1 0 0] /Rect [373.881 671.122 406.722 683.437] /A << /S /GoTo /D <74686D2E442E332E332E33> >> >>}
\MusicPDFWriteObject{1151}{<< /Type /Annot /Subtype /Link /Border [0 0 0] /H /I /C [1 0 0] /Rect [71.004 645.181 103.352 657.134] /A << /S /GoTo /D <74686D2E422E312E312E33> >> >>}
\MusicPDFWriteObject{1152}{<< /Type /Annot /Subtype /Link /Border [0 0 0] /H /I /C [1 0 0] /Rect [391.738 645.181 422.008 657.134] /A << /S /GoTo /D <74686D2E332E342E312E31> >> >>}
\MusicPDFWriteObject{1153}{<< /Type /Annot /Subtype /Link /Border [0 0 0] /H /I /C [1 0 0] /Rect [170.661 632.03 200.874 643.983] /A << /S /GoTo /D <74686D2E332E342E332E33> >> >>}
\MusicPDFWriteObject{1154}{<< /Type /Annot /Subtype /Link /Border [0 0 0] /H /I /C [1 0 0] /Rect [114.002 606.282 144.78 617.073] /A << /S /GoTo /D <74686D2E322E322E332E34> >> >>}
\MusicPDFWriteObject{1155}{<< /Type /Annot /Subtype /Link /Border [0 0 0] /H /I /C [1 0 0] /Rect [200.791 514.227 231.004 525.017] /A << /S /GoTo /D <74686D2E322E342E322E31> >> >>}
\MusicPDFWriteObject{1156}{<< /Type /Annot /Subtype /Link /Border [0 0 0] /H /I /C [1 0 0] /Rect [305.318 514.227 335.531 525.017] /A << /S /GoTo /D <74686D2E322E332E312E32> >> >>}
\MusicPDFWriteObject{1157}{<< /Type /Annot /Subtype /Link /Border [0 0 0] /H /I /C [1 0 0] /Rect [412.066 487.925 441.827 498.716] /A << /S /GoTo /D <74686D2E322E342E322E31> >> >>}
\MusicPDFWriteObject{1158}{<< /Type /Annot /Subtype /Link /Border [0 0 0] /H /I /C [1 0 0] /Rect [515.084 487.925 544.845 498.716] /A << /S /GoTo /D <74686D2E322E332E312E32> >> >>}
\MusicPDFWriteObject{1159}{<< /Type /Annot /Subtype /Link /Border [0 0 0] /H /I /C [1 0 0] /Rect [493.217 447.559 522.865 459.264] /A << /S /GoTo /D <74686D2E322E342E322E31> >> >>}
\MusicPDFWriteObject{1160}{<< /Type /Annot /Subtype /Link /Border [0 0 0] /H /I /C [1 0 0] /Rect [180.183 434.408 210.396 446.113] /A << /S /GoTo /D <74686D2E322E332E312E32> >> >>}
\MusicPDFWriteObject{1161}{<< /Type /Annot /Subtype /Link /Border [0 0 0] /H /I /C [1 0 0] /Rect [124.487 325.531 154.701 337.483] /A << /S /GoTo /D <6571756174696F6E2E332E342E312E36> >> >>}
\MusicPDFWriteObject{1162}{<< /Type /Annot /Subtype /Link /Border [0 0 0] /H /I /C [0 1 0] /Rect [73.77 288.568 107.649 297.436] /A << /S /GoTo /D <636974652E5765696C2D61> >> >>}
\MusicPDFWriteObject{1163}{<< /Type /Annot /Subtype /Link /Border [0 0 0] /H /I /C [1 0 0] /Rect [135.279 82.596 165.492 94.549] /A << /S /GoTo /D <74686D2E332E342E332E31> >> >>}
\MusicPDFWriteObject{1164}{<< /Type /Annot /Subtype /Link /Border [0 0 0] /H /I /C [1 0 0] /Rect [332.989 82.596 363.202 94.549] /A << /S /GoTo /D <74686D2E352E352E312E31> >> >>}
\MusicPDFWriteObject{1165}{<< /Type /Annot /Subtype /Link /Border [0 0 0] /H /I /C [1 0 0] /Rect [308.561 462.922 338.382 474.877] /A << /S /GoTo /D <74686D2E332E342E332E34> >> >>}
\MusicPDFWriteObject{1166}{<< /Type /Annot /Subtype /Link /Border [0 0 0] /H /I /C [1 0 0] /Rect [267.596 380.455 303.453 392.411] /A << /S /GoTo /D <6571756174696F6E2E332E342E332E3134> >> >>}
\MusicPDFWriteObject{1167}{<< /Type /Annot /Subtype /Link /Border [0 0 0] /H /I /C [1 0 0] /Rect [227.767 328.406 264.465 339.197] /A << /S /GoTo /D <74686D2E432E322E332E3132> >> >>}
\MusicPDFWriteObject{1168}{<< /Type /Annot /Subtype /Link /Border [0 0 0] /H /I /C [0 1 0] /Rect [236.518 297.466 270.397 306.334] /A << /S /GoTo /D <636974652E5765696C2D61> >> >>}
\MusicPDFWriteObject{1169}{<< /Type /Annot /Subtype /Link /Border [0 0 0] /H /I /C [1 0 0] /Rect [274.826 281.825 307.861 293.78] /A << /S /GoTo /D <74686D2E432E352E322E36> >> >>}
\MusicPDFWriteObject{1170}{<< /Type /Annot /Subtype /Link /Border [0 0 0] /H /I /C [1 0 0] /Rect [371.971 281.825 405.006 293.78] /A << /S /GoTo /D <74686D2E432E352E322E34> >> >>}
\MusicPDFWriteObject{1171}{<< /Type /Annot /Subtype /Link /Border [0 0 0] /H /I /C [1 0 0] /Rect [511.81 252.32 544.845 266.982] /A << /S /GoTo /D <74686D2E432E352E322E36> >> >>}
\MusicPDFWriteObject{1172}{<< /Type /Annot /Subtype /Link /Border [0 0 0] /H /I /C [1 0 0] /Rect [132.453 239.656 164.879 253.316] /A << /S /GoTo /D <74686D2E432E352E322E34> >> >>}
\MusicPDFWriteObject{1173}{<< /Type /Annot /Subtype /Link /Border [0 0 0] /H /I /C [0 1 0] /Rect [227.567 230.683 260.015 239.551] /A << /S /GoTo /D <636974652E4B75646C61> >> >>}
\MusicPDFWriteObject{1174}{<< /Type /Annot /Subtype /Link /Border [0 0 0] /H /I /C [0 1 0] /Rect [404.931 230.683 437.378 239.551] /A << /S /GoTo /D <636974652E4B75646C61> >> >>}
\MusicPDFWriteObject{1175}{<< /Type /Annot /Subtype /Link /Border [0 0 0] /H /I /C [0 1 0] /Rect [322.173 217.532 354.621 226.4] /A << /S /GoTo /D <636974652E4B75646C61> >> >>}
\MusicPDFWriteObject{1176}{<< /Type /Annot /Subtype /Link /Border [0 0 0] /H /I /C [0 1 0] /Rect [190.721 178.08 223.168 186.948] /A << /S /GoTo /D <636974652E4B75646C61> >> >>}
\MusicPDFWriteObject{1177}{<< /Type /Annot /Subtype /Link /Border [0 0 0] /H /I /C [0 1 0] /Rect [387.377 164.929 419.825 173.797] /A << /S /GoTo /D <636974652E4B75646C61> >> >>}
\MusicPDFWriteObject{1178}{<< /Type /Annot /Subtype /Link /Border [0 0 0] /H /I /C [1 0 0] /Rect [241.95 121.3 274.985 134.959] /A << /S /GoTo /D <74686D2E432E352E322E36> >> >>}
\MusicPDFWriteObject{1179}{<< /Type /Annot /Subtype /Link /Border [0 0 0] /H /I /C [1 0 0] /Rect [340.366 121.3 373.402 134.959] /A << /S /GoTo /D <74686D2E432E352E322E34> >> >>}
\MusicPDFWriteObject{1180}{<< /Type /Annot /Subtype /Link /Border [0 0 0] /H /I /C [1 0 0] /Rect [284.927 109.836 315.14 121.789] /A << /S /GoTo /D <74686D2E322E312E332E32> >> >>}
\MusicPDFWriteObject{1181}{<< /Type /Annot /Subtype /Link /Border [0 0 0] /H /I /C [1 0 0] /Rect [82.959 560.398 176.115 569.252] /A << /S /GoTo /D <73656374696F6E2E342E31> >> >>}
\MusicPDFWriteObject{1182}{<< /Type /Annot /Subtype /Link /Border [0 0 0] /H /I /C [1 0 0] /Rect [82.959 545.311 247.994 556.102] /A << /S /GoTo /D <73756273656374696F6E2E342E312E31> >> >>}
\MusicPDFWriteObject{1183}{<< /Type /Annot /Subtype /Link /Border [0 0 0] /H /I /C [1 0 0] /Rect [82.959 532.16 296.938 542.951] /A << /S /GoTo /D <73756273656374696F6E2E342E312E32> >> >>}
\MusicPDFWriteObject{1184}{<< /Type /Annot /Subtype /Link /Border [0 0 0] /H /I /C [1 0 0] /Rect [82.959 519.009 317.882 529.8] /A << /S /GoTo /D <73756273756273656374696F6E2E342E312E322E31> >> >>}
\MusicPDFWriteObject{1185}{<< /Type /Annot /Subtype /Link /Border [0 0 0] /H /I /C [1 0 0] /Rect [82.959 505.773 453.767 516.649] /A << /S /GoTo /D <73656374696F6E2E342E32> >> >>}
\MusicPDFWriteObject{1186}{<< /Type /Annot /Subtype /Link /Border [0 0 0] /H /I /C [1 0 0] /Rect [82.959 492.154 293.878 505.037] /A << /S /GoTo /D <73756273656374696F6E2E342E322E31> >> >>}
\MusicPDFWriteObject{1187}{<< /Type /Annot /Subtype /Link /Border [0 0 0] /H /I /C [1 0 0] /Rect [82.959 479.557 524.224 490.348] /A << /S /GoTo /D <73756273756273656374696F6E2E342E322E312E31> >> >>}
\MusicPDFWriteObject{1188}{<< /Type /Annot /Subtype /Link /Border [0 0 0] /H /I /C [1 0 0] /Rect [82.959 465.852 509.173 477.808] /A << /S /GoTo /D <73756273756273656374696F6E2E342E322E312E32> >> >>}
\MusicPDFWriteObject{1189}{<< /Type /Annot /Subtype /Link /Border [0 0 0] /H /I /C [1 0 0] /Rect [82.959 452.702 367.584 465.585] /A << /S /GoTo /D <73756273656374696F6E2E342E322E32> >> >>}
\MusicPDFWriteObject{1190}{<< /Type /Annot /Subtype /Link /Border [0 0 0] /H /I /C [1 0 0] /Rect [82.959 440.105 263.267 450.896] /A << /S /GoTo /D <73656374696F6E2E342E33> >> >>}
\MusicPDFWriteObject{1191}{<< /Type /Annot /Subtype /Link /Border [0 0 0] /H /I /C [1 0 0] /Rect [82.959 426.954 256.82 437.745] /A << /S /GoTo /D <73756273656374696F6E2E342E332E31> >> >>}
\MusicPDFWriteObject{1192}{<< /Type /Annot /Subtype /Link /Border [0 0 0] /H /I /C [1 0 0] /Rect [82.959 413.249 318.871 425.842] /A << /S /GoTo /D <73756273756273656374696F6E2E342E332E312E31> >> >>}
\MusicPDFWriteObject{1193}{<< /Type /Annot /Subtype /Link /Border [0 0 0] /H /I /C [1 0 0] /Rect [82.959 400.653 309.655 412.691] /A << /S /GoTo /D <73756273656374696F6E2E342E332E32> >> >>}
\MusicPDFWriteObject{1194}{<< /Type /Annot /Subtype /Link /Border [0 0 0] /H /I /C [1 0 0] /Rect [82.959 387.502 479.047 398.293] /A << /S /GoTo /D <73756273656374696F6E2E342E332E33> >> >>}
\MusicPDFWriteObject{1195}{<< /Type /Annot /Subtype /Link /Border [0 0 0] /H /I /C [1 0 0] /Rect [82.959 374.351 337.194 385.142] /A << /S /GoTo /D <73756273756273656374696F6E2E342E332E332E31> >> >>}
\MusicPDFWriteObject{1196}{<< /Type /Annot /Subtype /Link /Border [0 0 0] /H /I /C [1 0 0] /Rect [82.959 361.2 544.596 371.991] /A << /S /GoTo /D <73756273756273656374696F6E2E342E332E332E32> >> >>}
\MusicPDFWriteObject{1197}{<< /Type /Annot /Subtype /Link /Border [0 0 0] /H /I /C [1 0 0] /Rect [71.004 348.05 193.566 358.37] /A << /S /GoTo /D <73756273756273656374696F6E2E342E332E332E32> >> >>}
\MusicPDFWriteObject{1198}{<< /Type /Annot /Subtype /Link /Border [0 0 0] /H /I /C [1 0 0] /Rect [82.959 334.813 417.821 345.69] /A << /S /GoTo /D <73656374696F6E2E342E34> >> >>}
\MusicPDFWriteObject{1199}{<< /Type /Annot /Subtype /Link /Border [0 0 0] /H /I /C [1 0 0] /Rect [82.959 321.194 327.305 333.149] /A << /S /GoTo /D <73756273656374696F6E2E342E342E31> >> >>}
\MusicPDFWriteObject{1200}{<< /Type /Annot /Subtype /Link /Border [0 0 0] /H /I /C [1 0 0] /Rect [82.959 308.043 276.569 320.927] /A << /S /GoTo /D <73756273656374696F6E2E342E342E32> >> >>}
\MusicPDFWriteObject{1201}{<< /Type /Annot /Subtype /Link /Border [0 0 0] /H /I /C [1 0 0] /Rect [82.959 294.893 284.316 307.776] /A << /S /GoTo /D <73756273756273656374696F6E2E342E342E322E31> >> >>}
\MusicPDFWriteObject{1202}{<< /Type /Annot /Subtype /Link /Border [0 0 0] /H /I /C [1 0 0] /Rect [82.959 282.296 346.628 293.087] /A << /S /GoTo /D <73656374696F6E2E342E35> >> >>}
\MusicPDFWriteObject{1203}{<< /Type /Annot /Subtype /Link /Border [0 0 0] /H /I /C [1 0 0] /Rect [82.959 269.145 315.641 279.936] /A << /S /GoTo /D <73756273656374696F6E2E342E352E31> >> >>}
\MusicPDFWriteObject{1204}{<< /Type /Annot /Subtype /Link /Border [0 0 0] /H /I /C [1 0 0] /Rect [82.959 255.909 351.143 266.785] /A << /S /GoTo /D <73756273756273656374696F6E2E342E352E312E31> >> >>}
\MusicPDFWriteObject{1205}{<< /Type /Annot /Subtype /Link /Border [0 0 0] /H /I /C [1 0 0] /Rect [82.959 241.93 302.705 254.245] /A << /S /GoTo /D <73756273756273656374696F6E2E342E352E312E32> >> >>}
\MusicPDFWriteObject{1206}{<< /Type /Annot /Subtype /Link /Border [0 0 0] /H /I /C [1 0 0] /Rect [82.959 229.693 411.785 240.483] /A << /S /GoTo /D <73756273656374696F6E2E342E352E32> >> >>}
\MusicPDFWriteObject{1207}{<< /Type /Annot /Subtype /Link /Border [0 0 0] /H /I /C [1 0 0] /Rect [82.959 216.542 466.94 227.333] /A << /S /GoTo /D <73756273656374696F6E2E342E352E33> >> >>}
\MusicPDFWriteObject{1208}{<< /Type /Annot /Subtype /Link /Border [0 0 0] /H /I /C [1 0 0] /Rect [82.959 203.391 349.976 214.182] /A << /S /GoTo /D <73656374696F6E2E342E36> >> >>}
\MusicPDFWriteObject{1209}{<< /Type /Annot /Subtype /Link /Border [0 0 0] /H /I /C [1 0 0] /Rect [82.959 190.241 275.911 201.031] /A << /S /GoTo /D <73756273656374696F6E2E342E362E31> >> >>}
\MusicPDFWriteObject{1210}{<< /Type /Annot /Subtype /Link /Border [0 0 0] /H /I /C [1 0 0] /Rect [82.959 177.09 446.147 187.88] /A << /S /GoTo /D <73756273756273656374696F6E2E342E362E312E31> >> >>}
\MusicPDFWriteObject{1211}{<< /Type /Annot /Subtype /Link /Border [0 0 0] /H /I /C [1 0 0] /Rect [82.959 163.939 334.399 174.73] /A << /S /GoTo /D <73756273656374696F6E2E342E362E32> >> >>}
\MusicPDFWriteObject{1212}{<< /Type /Annot /Subtype /Link /Border [0 0 0] /H /I /C [1 0 0] /Rect [82.959 150.234 196.605 162.189] /A << /S /GoTo /D <73756273656374696F6E2E342E362E33> >> >>}
\MusicPDFWriteObject{1213}{<< /Type /Annot /Subtype /Link /Border [0 0 0] /H /I /C [1 0 0] /Rect [194.234 553.679 211.331 564.527] /A << /S /GoTo /D <73656374696F6E2E432E33> >> >>}
\MusicPDFWriteObject{1214}{<< /Type /Annot /Subtype /Link /Border [0 0 0] /H /I /C [1 0 0] /Rect [306.227 553.679 331.156 564.527] /A << /S /GoTo /D <73756273656374696F6E2E432E332E36> >> >>}
\MusicPDFWriteObject{1215}{<< /Type /Annot /Subtype /Link /Border [0 0 0] /H /I /C [0 1 0] /Rect [148.187 542.465 171.885 551.333] /A << /S /GoTo /D <636974652E4B7562> >> >>}
\MusicPDFWriteObject{1216}{<< /Type /Annot /Subtype /Link /Border [0 0 0] /H /I /C [1 0 0] /Rect [161.903 501.076 194.208 511.867] /A << /S /GoTo /D <74686D2E432E332E342E31> >> >>}
\MusicPDFWriteObject{1217}{<< /Type /Annot /Subtype /Link /Border [0 0 0] /H /I /C [0 1 0] /Rect [345.186 358.354 363.385 367.222] /A << /S /GoTo /D <636974652E4A6172> >> >>}
\MusicPDFWriteObject{1218}{<< /Type /Annot /Subtype /Link /Border [0 0 0] /H /I /C [0 1 0] /Rect [175.842 345.204 198.792 354.071] /A << /S /GoTo /D <636974652E424A4C> >> >>}
\MusicPDFWriteObject{1219}{<< /Type /Annot /Subtype /Link /Border [0 0 0] /H /I /C [0 1 0] /Rect [341.163 266.299 360.136 275.167] /A << /S /GoTo /D <636974652E547265> >> >>}
\MusicPDFWriteObject{1220}{<< /Type /Annot /Subtype /Link /Border [0 0 0] /H /I /C [0 1 0] /Rect [96.099 253.148 115.072 262.016] /A << /S /GoTo /D <636974652E547265> >> >>}
\MusicPDFWriteObject{1221}{<< /Type /Annot /Subtype /Link /Border [0 0 0] /H /I /C [0 1 0] /Rect [516.46 687.123 543.159 695.991] /A << /S /GoTo /D <636974652E466F6C2D32> >> >>}
\MusicPDFWriteObject{1222}{<< /Type /Annot /Subtype /Link /Border [0 0 0] /H /I /C [0 1 0] /Rect [516.515 647.671 543.214 656.539] /A << /S /GoTo /D <636974652E466F6C2D32> >> >>}
\MusicPDFWriteObject{1223}{<< /Type /Annot /Subtype /Link /Border [0 0 0] /H /I /C [0 1 0] /Rect [116.433 621.369 162.024 630.237] /A << /S /GoTo /D <636974652E4361726C65736F6E> >> >>}
\MusicPDFWriteObject{1224}{<< /Type /Annot /Subtype /Link /Border [0 0 0] /H /I /C [1 0 0] /Rect [146.941 543 177.154 554.952] /A << /S /GoTo /D <6571756174696F6E2E342E312E312E31> >> >>}
\MusicPDFWriteObject{1225}{<< /Type /Annot /Subtype /Link /Border [0 0 0] /H /I /C [0 1 0] /Rect [291.506 519.188 309.705 528.056] /A << /S /GoTo /D <636974652E4A6172> >> >>}
\MusicPDFWriteObject{1226}{<< /Type /Annot /Subtype /Link /Border [0 0 0] /H /I /C [0 1 0] /Rect [226.794 506.037 244.993 514.905] /A << /S /GoTo /D <636974652E4A6172> >> >>}
\MusicPDFWriteObject{1227}{<< /Type /Annot /Subtype /Link /Border [0 0 0] /H /I /C [1 0 0] /Rect [446.432 312.551 478.657 324.504] /A << /S /GoTo /D <74686D2E442E322E312E31> >> >>}
\MusicPDFWriteObject{1228}{<< /Type /Annot /Subtype /Link /Border [0 0 0] /H /I /C [0 1 0] /Rect [510.778 137.689 543.225 146.557] /A << /S /GoTo /D <636974652E4B61747A6E> >> >>}
\MusicPDFWriteObject{1229}{<< /Type /Annot /Subtype /Link /Border [0 0 0] /H /I /C [1 0 0] /Rect [297.026 122.048 327.239 134.001] /A << /S /GoTo /D <6571756174696F6E2E342E312E312E31> >> >>}
\MusicPDFWriteObject{1230}{<< /Type /Annot /Subtype /Link /Border [0 0 0] /H /I /C [0 1 0] /Rect [357.498 111.387 385.318 120.255] /A << /S /GoTo /D <636974652E48756E74> >> >>}
\MusicPDFWriteObject{1231}{<< /Type /Annot /Subtype /Link /Border [0 0 0] /H /I /C [0 1 0] /Rect [196.645 687.123 229.092 695.991] /A << /S /GoTo /D <636974652E4B61747A6E> >> >>}
\MusicPDFWriteObject{1232}{<< /Type /Annot /Subtype /Link /Border [0 0 0] /H /I /C [0 1 0] /Rect [408.323 673.973 427.884 682.84] /A << /S /GoTo /D <636974652E4B6F6C> >> >>}
\MusicPDFWriteObject{1233}{<< /Type /Annot /Subtype /Link /Border [0 0 0] /H /I /C [1 0 0] /Rect [455.092 605.729 485.305 617.681] /A << /S /GoTo /D <6571756174696F6E2E342E312E312E31> >> >>}
\MusicPDFWriteObject{1234}{<< /Type /Annot /Subtype /Link /Border [0 0 0] /H /I /C [0 1 0] /Rect [181.196 542.465 210.538 551.333] /A << /S /GoTo /D <636974652E4C69742D4949> >> >>}
\MusicPDFWriteObject{1235}{<< /Type /Annot /Subtype /Link /Border [0 0 0] /H /I /C [0 1 0] /Rect [344.688 516.163 369.368 525.031] /A << /S /GoTo /D <636974652E4C69742D49> >> >>}
\MusicPDFWriteObject{1236}{<< /Type /Annot /Subtype /Link /Border [0 0 0] /H /I /C [0 1 0] /Rect [158.819 479.313 177.018 488.18] /A << /S /GoTo /D <636974652E4A6172> >> >>}
\MusicPDFWriteObject{1237}{<< /Type /Annot /Subtype /Link /Border [0 0 0] /H /I /C [0 1 0] /Rect [233.85 387.257 252.049 396.125] /A << /S /GoTo /D <636974652E4A6172> >> >>}
\MusicPDFWriteObject{1238}{<< /Type /Annot /Subtype /Link /Border [0 0 0] /H /I /C [0 1 0] /Rect [187.742 360.956 205.941 369.824] /A << /S /GoTo /D <636974652E4A6172> >> >>}
\MusicPDFWriteObject{1239}{<< /Type /Annot /Subtype /Link /Border [0 0 0] /H /I /C [0 1 0] /Rect [410.357 360.956 428.556 369.824] /A << /S /GoTo /D <636974652E4A6172> >> >>}
\MusicPDFWriteObject{1240}{<< /Type /Annot /Subtype /Link /Border [0 0 0] /H /I /C [0 1 0] /Rect [277.806 347.805 296.004 356.673] /A << /S /GoTo /D <636974652E4A6172> >> >>}
\MusicPDFWriteObject{1241}{<< /Type /Annot /Subtype /Link /Border [0 0 0] /H /I /C [0 1 0] /Rect [472.079 347.805 491.8 356.673] /A << /S /GoTo /D <636974652E4B616C> >> >>}
\MusicPDFWriteObject{1242}{<< /Type /Annot /Subtype /Link /Border [0 0 0] /H /I /C [0 1 0] /Rect [164.049 334.654 183.769 343.522] /A << /S /GoTo /D <636974652E4B616C> >> >>}
\MusicPDFWriteObject{1243}{<< /Type /Annot /Subtype /Link /Border [0 0 0] /H /I /C [1 0 0] /Rect [230.55 95.747 260.763 107.699] /A << /S /GoTo /D <6571756174696F6E2E342E312E322E31> >> >>}
\MusicPDFWriteObject{1244}{<< /Type /Annot /Subtype /Link /Border [0 0 0] /H /I /C [0 1 0] /Rect [501.588 687.123 519.787 695.991] /A << /S /GoTo /D <636974652E4A6172> >> >>}
\MusicPDFWriteObject{1245}{<< /Type /Annot /Subtype /Link /Border [0 0 0] /H /I /C [0 1 0] /Rect [365.523 673.973 385.243 682.84] /A << /S /GoTo /D <636974652E4B616C> >> >>}
\MusicPDFWriteObject{1246}{<< /Type /Annot /Subtype /Link /Border [0 0 0] /H /I /C [0 1 0] /Rect [87.637 597.849 107.357 606.717] /A << /S /GoTo /D <636974652E4B616C> >> >>}
\MusicPDFWriteObject{1247}{<< /Type /Annot /Subtype /Link /Border [0 0 0] /H /I /C [1 0 0] /Rect [93.688 440.331 126.114 452.284] /A << /S /GoTo /D <6571756174696F6E2E432E342E312E31> >> >>}
\MusicPDFWriteObject{1248}{<< /Type /Annot /Subtype /Link /Border [0 0 0] /H /I /C [1 0 0] /Rect [186.234 440.331 218.661 452.284] /A << /S /GoTo /D <74686D2E432E342E312E31> >> >>}
\MusicPDFWriteObject{1249}{<< /Type /Annot /Subtype /Link /Border [0 0 0] /H /I /C [1 0 0] /Rect [123.475 425.493 153.124 439.153] /A << /S /GoTo /D <74686D2E342E312E322E34> >> >>}
\MusicPDFWriteObject{1250}{<< /Type /Annot /Subtype /Link /Border [0 0 0] /H /I /C [1 0 0] /Rect [355.582 241.888 385.795 253.841] /A << /S /GoTo /D <6571756174696F6E2E342E312E322E32> >> >>}
\MusicPDFWriteObject{1251}{<< /Type /Annot /Subtype /Link /Border [0 0 0] /H /I /C [1 0 0] /Rect [206.678 155.667 236.891 167.619] /A << /S /GoTo /D <6571756174696F6E2E342E312E322E33> >> >>}
\MusicPDFWriteObject{1252}{<< /Type /Annot /Subtype /Link /Border [0 0 0] /H /I /C [1 0 0] /Rect [318.32 95.747 348.533 107.699] /A << /S /GoTo /D <6571756174696F6E2E342E312E322E34> >> >>}
\MusicPDFWriteObject{1253}{<< /Type /Annot /Subtype /Link /Border [0 0 0] /H /I /C [1 0 0] /Rect [73.712 82.596 103.925 94.549] /A << /S /GoTo /D <6571756174696F6E2E342E312E322E34> >> >>}
\MusicPDFWriteObject{1254}{<< /Type /Annot /Subtype /Link /Border [0 0 0] /H /I /C [1 0 0] /Rect [108.815 635.055 139.028 647.008] /A << /S /GoTo /D <6571756174696F6E2E342E312E322E34> >> >>}
\MusicPDFWriteObject{1255}{<< /Type /Annot /Subtype /Link /Border [0 0 0] /H /I /C [0 1 0] /Rect [161.339 611.243 179.537 620.111] /A << /S /GoTo /D <636974652E4A6172> >> >>}
\MusicPDFWriteObject{1256}{<< /Type /Annot /Subtype /Link /Border [0 0 0] /H /I /C [0 1 0] /Rect [386.279 598.093 404.477 606.961] /A << /S /GoTo /D <636974652E4A6172> >> >>}
\MusicPDFWriteObject{1257}{<< /Type /Annot /Subtype /Link /Border [0 0 0] /H /I /C [1 0 0] /Rect [108.776 569.301 138.989 581.254] /A << /S /GoTo /D <6571756174696F6E2E342E312E322E34> >> >>}
\MusicPDFWriteObject{1258}{<< /Type /Annot /Subtype /Link /Border [0 0 0] /H /I /C [0 1 0] /Rect [425.594 545.49 444.567 554.357] /A << /S /GoTo /D <636974652E547265> >> >>}
\MusicPDFWriteObject{1259}{<< /Type /Annot /Subtype /Link /Border [0 0 0] /H /I /C [0 1 0] /Rect [87.994 532.339 106.193 541.207] /A << /S /GoTo /D <636974652E4A6172> >> >>}
\MusicPDFWriteObject{1260}{<< /Type /Annot /Subtype /Link /Border [0 0 0] /H /I /C [0 1 0] /Rect [407.299 532.339 425.498 541.207] /A << /S /GoTo /D <636974652E4A6172> >> >>}
\MusicPDFWriteObject{1261}{<< /Type /Annot /Subtype /Link /Border [0 0 0] /H /I /C [0 1 0] /Rect [192.586 425.136 210.784 434.004] /A << /S /GoTo /D <636974652E4A6172> >> >>}
\MusicPDFWriteObject{1262}{<< /Type /Annot /Subtype /Link /Border [0 0 0] /H /I /C [0 1 0] /Rect [409.411 425.136 440.744 434.004] /A << /S /GoTo /D <636974652E4D634172> >> >>}
\MusicPDFWriteObject{1263}{<< /Type /Annot /Subtype /Link /Border [0 0 0] /H /I /C [0 1 0] /Rect [73.77 410.049 101.466 420.853] /A << /S /GoTo /D <636974652E44794D> >> >>}
\MusicPDFWriteObject{1264}{<< /Type /Annot /Subtype /Link /Border [0 0 0] /H /I /C [0 1 0] /Rect [458.617 382.438 489.951 391.306] /A << /S /GoTo /D <636974652E4D634172> >> >>}
\MusicPDFWriteObject{1265}{<< /Type /Annot /Subtype /Link /Border [0 0 0] /H /I /C [0 1 0] /Rect [502.233 297.398 520.432 306.266] /A << /S /GoTo /D <636974652E4A6172> >> >>}
\MusicPDFWriteObject{1266}{<< /Type /Annot /Subtype /Link /Border [0 0 0] /H /I /C [0 1 0] /Rect [134.858 271.096 153.057 279.964] /A << /S /GoTo /D <636974652E4A6172> >> >>}
\MusicPDFWriteObject{1267}{<< /Type /Annot /Subtype /Link /Border [0 0 0] /H /I /C [0 1 0] /Rect [321.804 244.795 340.003 253.663] /A << /S /GoTo /D <636974652E4A6172> >> >>}
\MusicPDFWriteObject{1268}{<< /Type /Annot /Subtype /Link /Border [0 0 0] /H /I /C [0 1 0] /Rect [356.939 218.493 375.137 227.361] /A << /S /GoTo /D <636974652E4A6172> >> >>}
\MusicPDFWriteObject{1269}{<< /Type /Annot /Subtype /Link /Border [0 0 0] /H /I /C [0 1 0] /Rect [73.77 205.343 91.969 214.21] /A << /S /GoTo /D <636974652E4A6172> >> >>}
\MusicPDFWriteObject{1270}{<< /Type /Annot /Subtype /Link /Border [0 0 0] /H /I /C [0 1 0] /Rect [107.901 165.89 126.099 174.758] /A << /S /GoTo /D <636974652E4A6172> >> >>}
\MusicPDFWriteObject{1271}{<< /Type /Annot /Subtype /Link /Border [0 0 0] /H /I /C [0 1 0] /Rect [239.811 152.74 258.01 161.607] /A << /S /GoTo /D <636974652E4A6172> >> >>}
\MusicPDFWriteObject{1272}{<< /Type /Annot /Subtype /Link /Border [0 0 0] /H /I /C [1 0 0] /Rect [131.965 122.601 168.664 134.64] /A << /S /GoTo /D <74686D2E432E322E332E3132> >> >>}
\MusicPDFWriteObject{1273}{<< /Type /Annot /Subtype /Link /Border [0 0 0] /H /I /C [1 0 0] /Rect [296.731 95.747 334.137 107.699] /A << /S /GoTo /D <6571756174696F6E2E432E322E332E3232> >> >>}
\MusicPDFWriteObject{1274}{<< /Type /Annot /Subtype /Link /Border [0 0 0] /H /I /C [1 0 0] /Rect [423.635 95.747 460.652 107.699] /A << /S /GoTo /D <74686D2E432E322E332E3133> >> >>}
\MusicPDFWriteObject{1275}{<< /Type /Annot /Subtype /Link /Border [0 0 0] /H /I /C [1 0 0] /Rect [477.644 95.747 502.074 107.699] /A << /S /GoTo /D <73756273656374696F6E2E432E322E33> >> >>}
\MusicPDFWriteObject{1276}{<< /Type /Annot /Subtype /Link /Border [0 0 0] /H /I /C [1 0 0] /Rect [156.389 684.633 193.088 697.225] /A << /S /GoTo /D <74686D2E432E322E332E3133> >> >>}
\MusicPDFWriteObject{1277}{<< /Type /Annot /Subtype /Link /Border [0 0 0] /H /I /C [0 1 0] /Rect [501.282 673.973 519.481 682.84] /A << /S /GoTo /D <636974652E4A6172> >> >>}
\MusicPDFWriteObject{1278}{<< /Type /Annot /Subtype /Link /Border [0 0 0] /H /I /C [0 1 0] /Rect [447.008 608.219 478.341 617.087] /A << /S /GoTo /D <636974652E4D634172> >> >>}
\MusicPDFWriteObject{1279}{<< /Type /Annot /Subtype /Link /Border [0 0 0] /H /I /C [0 1 0] /Rect [274.801 595.068 292.999 603.936] /A << /S /GoTo /D <636974652E4A6172> >> >>}
\MusicPDFWriteObject{1280}{<< /Type /Annot /Subtype /Link /Border [0 0 0] /H /I /C [0 1 0] /Rect [185.789 499.138 203.988 508.005] /A << /S /GoTo /D <636974652E4A6172> >> >>}
\MusicPDFWriteObject{1281}{<< /Type /Annot /Subtype /Link /Border [0 0 0] /H /I /C [0 1 0] /Rect [129.457 485.987 147.656 494.855] /A << /S /GoTo /D <636974652E4A6172> >> >>}
\MusicPDFWriteObject{1282}{<< /Type /Annot /Subtype /Link /Border [0 0 0] /H /I /C [0 1 0] /Rect [137.29 459.685 155.489 468.553] /A << /S /GoTo /D <636974652E4A6172> >> >>}
\MusicPDFWriteObject{1283}{<< /Type /Annot /Subtype /Link /Border [0 0 0] /H /I /C [0 1 0] /Rect [354.469 446.534 372.667 455.402] /A << /S /GoTo /D <636974652E4A6172> >> >>}
\MusicPDFWriteObject{1284}{<< /Type /Annot /Subtype /Link /Border [0 0 0] /H /I /C [0 1 0] /Rect [73.759 433.384 91.958 442.252] /A << /S /GoTo /D <636974652E4A6172> >> >>}
\MusicPDFWriteObject{1285}{<< /Type /Annot /Subtype /Link /Border [0 0 0] /H /I /C [0 1 0] /Rect [353.008 433.384 371.207 442.252] /A << /S /GoTo /D <636974652E4A6172> >> >>}
\MusicPDFWriteObject{1286}{<< /Type /Annot /Subtype /Link /Border [0 0 0] /H /I /C [0 1 0] /Rect [525.027 433.384 543.225 442.252] /A << /S /GoTo /D <636974652E4A6172> >> >>}
\MusicPDFWriteObject{1287}{<< /Type /Annot /Subtype /Link /Border [0 0 0] /H /I /C [0 1 0] /Rect [500.258 407.082 518.456 415.95] /A << /S /GoTo /D <636974652E4A6172> >> >>}
\MusicPDFWriteObject{1288}{<< /Type /Annot /Subtype /Link /Border [0 0 0] /H /I /C [0 1 0] /Rect [395.563 393.931 420.623 402.799] /A << /S /GoTo /D <636974652E53636861> >> >>}
\MusicPDFWriteObject{1289}{<< /Type /Annot /Subtype /Link /Border [0 0 0] /H /I /C [1 0 0] /Rect [208.183 358.266 238.086 369.057] /A << /S /GoTo /D <74686D2E342E312E322E36> >> >>}
\MusicPDFWriteObject{1290}{<< /Type /Annot /Subtype /Link /Border [0 0 0] /H /I /C [0 1 0] /Rect [343.711 347.052 361.91 355.92] /A << /S /GoTo /D <636974652E4A6172> >> >>}
\MusicPDFWriteObject{1291}{<< /Type /Annot /Subtype /Link /Border [0 0 0] /H /I /C [0 1 0] /Rect [221.237 254.997 240.21 263.865] /A << /S /GoTo /D <636974652E547265> >> >>}
\MusicPDFWriteObject{1292}{<< /Type /Annot /Subtype /Link /Border [0 0 0] /H /I /C [1 0 0] /Rect [116.433 182.658 148.657 194.611] /A << /S /GoTo /D <74686D2E442E322E312E31> >> >>}
\MusicPDFWriteObject{1293}{<< /Type /Annot /Subtype /Link /Border [0 0 0] /H /I /C [1 0 0] /Rect [205.219 182.658 237.444 194.611] /A << /S /GoTo /D <74686D2E442E332E322E35> >> >>}
\MusicPDFWriteObject{1294}{<< /Type /Annot /Subtype /Link /Border [0 0 0] /H /I /C [1 0 0] /Rect [298.964 182.658 331.189 194.611] /A << /S /GoTo /D <74686D2E442E332E332E37> >> >>}
\MusicPDFWriteObject{1295}{<< /Type /Annot /Subtype /Link /Border [0 0 0] /H /I /C [1 0 0] /Rect [270.618 156.91 293.493 167.758] /A << /S /GoTo /D <73756273656374696F6E2E342E362E33> >> >>}
\MusicPDFWriteObject{1296}{<< /Type /Annot /Subtype /Link /Border [0 0 0] /H /I /C [1 0 0] /Rect [218.852 82.596 251.693 94.549] /A << /S /GoTo /D <6571756174696F6E2E442E312E372E32> >> >>}
\MusicPDFWriteObject{1297}{<< /Type /Annot /Subtype /Link /Border [0 0 0] /H /I /C [1 0 0] /Rect [337.859 685.187 371.317 696.034] /A << /S /GoTo /D <74686D2E442E312E372E31> >> >>}
\MusicPDFWriteObject{1298}{<< /Type /Annot /Subtype /Link /Border [0 0 0] /H /I /C [1 0 0] /Rect [127.892 240.697 161.329 254.891] /A << /S /GoTo /D <74686D2E442E312E322E31> >> >>}
\MusicPDFWriteObject{1299}{<< /Type /Annot /Subtype /Link /Border [0 0 0] /H /I /C [0 1 0] /Rect [348.758 98.237 366.957 107.105] /A << /S /GoTo /D <636974652E4A6172> >> >>}
\MusicPDFWriteObject{1300}{<< /Type /Annot /Subtype /Link /Border [0 0 0] /H /I /C [1 0 0] /Rect [222.596 392.205 252.809 404.157] /A << /S /GoTo /D <6571756174696F6E2E342E322E312E36> >> >>}
\MusicPDFWriteObject{1301}{<< /Type /Annot /Subtype /Link /Border [0 0 0] /H /I /C [1 0 0] /Rect [480.391 364.924 510.689 377.808] /A << /S /GoTo /D <74686D2E342E312E322E35> >> >>}
\MusicPDFWriteObject{1302}{<< /Type /Annot /Subtype /Link /Border [0 0 0] /H /I /C [1 0 0] /Rect [87.542 338.624 117.755 350.576] /A << /S /GoTo /D <6571756174696F6E2E342E322E312E36> >> >>}
\MusicPDFWriteObject{1303}{<< /Type /Annot /Subtype /Link /Border [0 0 0] /H /I /C [1 0 0] /Rect [475.859 325.472 506.496 338.356] /A << /S /GoTo /D <74686D2E342E312E322E36> >> >>}
\MusicPDFWriteObject{1304}{<< /Type /Annot /Subtype /Link /Border [0 0 0] /H /I /C [1 0 0] /Rect [254.206 299.171 284.673 312.054] /A << /S /GoTo /D <74686D2E342E312E322E34> >> >>}
\MusicPDFWriteObject{1305}{<< /Type /Annot /Subtype /Link /Border [0 0 0] /H /I /C [1 0 0] /Rect [260.928 109.09 291.141 121.043] /A << /S /GoTo /D <6571756174696F6E2E342E322E312E38> >> >>}
\MusicPDFWriteObject{1306}{<< /Type /Annot /Subtype /Link /Border [0 0 0] /H /I /C [1 0 0] /Rect [217.651 684.633 247.864 696.586] /A << /S /GoTo /D <6571756174696F6E2E342E322E312E38> >> >>}
\MusicPDFWriteObject{1307}{<< /Type /Annot /Subtype /Link /Border [0 0 0] /H /I /C [0 1 0] /Rect [290.94 636.373 310.045 645.241] /A << /S /GoTo /D <636974652E4E42> >> >>}
\MusicPDFWriteObject{1308}{<< /Type /Annot /Subtype /Link /Border [0 0 0] /H /I /C [0 1 0] /Rect [106.885 623.222 125.989 632.09] /A << /S /GoTo /D <636974652E4E42> >> >>}
\MusicPDFWriteObject{1309}{<< /Type /Annot /Subtype /Link /Border [0 0 0] /H /I /C [0 1 0] /Rect [263.323 623.222 282.428 632.09] /A << /S /GoTo /D <636974652E4E42> >> >>}
\MusicPDFWriteObject{1310}{<< /Type /Annot /Subtype /Link /Border [0 0 0] /H /I /C [0 1 0] /Rect [438.261 610.071 457.365 618.939] /A << /S /GoTo /D <636974652E4E42> >> >>}
\MusicPDFWriteObject{1311}{<< /Type /Annot /Subtype /Link /Border [0 0 0] /H /I /C [1 0 0] /Rect [289.041 594.431 319.254 606.383] /A << /S /GoTo /D <6571756174696F6E2E342E322E312E38> >> >>}
\MusicPDFWriteObject{1312}{<< /Type /Annot /Subtype /Link /Border [0 0 0] /H /I /C [0 1 0] /Rect [129.413 568.682 164.088 579.487] /A << /S /GoTo /D <636974652E4170696F6C61> >> >>}
\MusicPDFWriteObject{1313}{<< /Type /Annot /Subtype /Link /Border [0 0 0] /H /I /C [0 1 0] /Rect [106.971 557.468 138.042 566.336] /A << /S /GoTo /D <636974652E4D562D32> >> >>}
\MusicPDFWriteObject{1314}{<< /Type /Annot /Subtype /Link /Border [0 0 0] /H /I /C [0 1 0] /Rect [106.525 544.318 124.723 553.185] /A << /S /GoTo /D <636974652E4A6172> >> >>}
\MusicPDFWriteObject{1315}{<< /Type /Annot /Subtype /Link /Border [0 0 0] /H /I /C [0 1 0] /Rect [287.821 544.318 306.794 553.185] /A << /S /GoTo /D <636974652E547265> >> >>}
\MusicPDFWriteObject{1316}{<< /Type /Annot /Subtype /Link /Border [0 0 0] /H /I /C [0 1 0] /Rect [304.879 531.167 333.508 540.035] /A << /S /GoTo /D <636974652E56616C2D62> >> >>}
\MusicPDFWriteObject{1317}{<< /Type /Annot /Subtype /Link /Border [0 0 0] /H /I /C [0 1 0] /Rect [518.165 504.865 543.225 513.733] /A << /S /GoTo /D <636974652E53636861> >> >>}
\MusicPDFWriteObject{1318}{<< /Type /Annot /Subtype /Link /Border [0 0 0] /H /I /C [0 1 0] /Rect [264.582 491.715 283.555 500.582] /A << /S /GoTo /D <636974652E547265> >> >>}
\MusicPDFWriteObject{1319}{<< /Type /Annot /Subtype /Link /Border [0 0 0] /H /I /C [0 1 0] /Rect [518.176 491.715 543.236 500.582] /A << /S /GoTo /D <636974652E53636861> >> >>}
\MusicPDFWriteObject{1320}{<< /Type /Annot /Subtype /Link /Border [0 0 0] /H /I /C [0 1 0] /Rect [73.715 425.345 98.775 434.213] /A << /S /GoTo /D <636974652E53636861> >> >>}
\MusicPDFWriteObject{1321}{<< /Type /Annot /Subtype /Link /Border [0 0 0] /H /I /C [0 1 0] /Rect [262.721 425.345 287.781 434.213] /A << /S /GoTo /D <636974652E53636861> >> >>}
\MusicPDFWriteObject{1322}{<< /Type /Annot /Subtype /Link /Border [0 0 0] /H /I /C [0 1 0] /Rect [108.376 399.044 127.481 407.912] /A << /S /GoTo /D <636974652E4E42> >> >>}
\MusicPDFWriteObject{1323}{<< /Type /Annot /Subtype /Link /Border [0 0 0] /H /I /C [1 0 0] /Rect [355.975 357.102 386.188 369.054] /A << /S /GoTo /D <6571756174696F6E2E342E322E312E38> >> >>}
\MusicPDFWriteObject{1324}{<< /Type /Annot /Subtype /Link /Border [0 0 0] /H /I /C [0 1 0] /Rect [129.739 320.139 158.368 329.007] /A << /S /GoTo /D <636974652E56616C2D62> >> >>}
\MusicPDFWriteObject{1325}{<< /Type /Annot /Subtype /Link /Border [0 0 0] /H /I /C [0 1 0] /Rect [366.794 293.838 391.854 302.706] /A << /S /GoTo /D <636974652E53636861> >> >>}
\MusicPDFWriteObject{1326}{<< /Type /Annot /Subtype /Link /Border [0 0 0] /H /I /C [0 1 0] /Rect [518.151 293.838 543.211 302.706] /A << /S /GoTo /D <636974652E53636861> >> >>}
\MusicPDFWriteObject{1327}{<< /Type /Annot /Subtype /Link /Border [0 0 0] /H /I /C [0 1 0] /Rect [186.951 280.687 212.011 289.555] /A << /S /GoTo /D <636974652E53636861> >> >>}
\MusicPDFWriteObject{1328}{<< /Type /Annot /Subtype /Link /Border [0 0 0] /H /I /C [0 1 0] /Rect [73.715 267.536 92.82 276.404] /A << /S /GoTo /D <636974652E4E42> >> >>}
\MusicPDFWriteObject{1329}{<< /Type /Annot /Subtype /Link /Border [0 0 0] /H /I /C [1 0 0] /Rect [208.696 265.046 238.909 276.999] /A << /S /GoTo /D <6571756174696F6E2E342E322E312E38> >> >>}
\MusicPDFWriteObject{1330}{<< /Type /Annot /Subtype /Link /Border [0 0 0] /H /I /C [0 1 0] /Rect [73.77 190.966 104.281 199.834] /A << /S /GoTo /D <636974652E47726F2D65> >> >>}
\MusicPDFWriteObject{1331}{<< /Type /Annot /Subtype /Link /Border [0 0 0] /H /I /C [1 0 0] /Rect [387.716 175.324 417.929 187.297] /A << /S /GoTo /D <74686D2E342E322E312E33> >> >>}
\MusicPDFWriteObject{1332}{<< /Type /Annot /Subtype /Link /Border [0 0 0] /H /I /C [0 1 0] /Rect [466.038 164.664 485.142 173.532] /A << /S /GoTo /D <636974652E4E42> >> >>}
\MusicPDFWriteObject{1333}{<< /Type /Annot /Subtype /Link /Border [0 0 0] /H /I /C [1 0 0] /Rect [352.647 148.542 382.86 160.495] /A << /S /GoTo /D <6571756174696F6E2E342E322E312E38> >> >>}
\MusicPDFWriteObject{1334}{<< /Type /Annot /Subtype /Link /Border [0 0 0] /H /I /C [0 1 0] /Rect [107.685 98.429 132.745 107.297] /A << /S /GoTo /D <636974652E53636861> >> >>}
\MusicPDFWriteObject{1335}{<< /Type /Annot /Subtype /Link /Border [0 0 0] /H /I /C [0 1 0] /Rect [450.069 85.278 469.174 94.146] /A << /S /GoTo /D <636974652E4E42> >> >>}
\MusicPDFWriteObject{1336}{<< /Type /Annot /Subtype /Link /Border [0 0 0] /H /I /C [0 1 0] /Rect [73.823 687.123 92.928 695.991] /A << /S /GoTo /D <636974652E4E42> >> >>}
\MusicPDFWriteObject{1337}{<< /Type /Annot /Subtype /Link /Border [0 0 0] /H /I /C [0 1 0] /Rect [285.09 687.123 304.194 695.991] /A << /S /GoTo /D <636974652E4E42> >> >>}
\MusicPDFWriteObject{1338}{<< /Type /Annot /Subtype /Link /Border [0 0 0] /H /I /C [0 1 0] /Rect [408.027 687.123 433.087 695.991] /A << /S /GoTo /D <636974652E53636861> >> >>}
\MusicPDFWriteObject{1339}{<< /Type /Annot /Subtype /Link /Border [0 0 0] /H /I /C [0 1 0] /Rect [190.443 673.973 215.503 682.84] /A << /S /GoTo /D <636974652E53636861> >> >>}
\MusicPDFWriteObject{1340}{<< /Type /Annot /Subtype /Link /Border [0 0 0] /H /I /C [0 1 0] /Rect [107.194 647.671 135.823 656.539] /A << /S /GoTo /D <636974652E56616C2D62> >> >>}
\MusicPDFWriteObject{1341}{<< /Type /Annot /Subtype /Link /Border [0 0 0] /H /I /C [0 1 0] /Rect [519.932 621.369 543.187 630.237] /A << /S /GoTo /D <636974652E466C4D> >> >>}
\MusicPDFWriteObject{1342}{<< /Type /Annot /Subtype /Link /Border [0 0 0] /H /I /C [1 0 0] /Rect [232.017 579.981 263.834 590.771] /A << /S /GoTo /D <74686D2E432E312E352E39> >> >>}
\MusicPDFWriteObject{1343}{<< /Type /Annot /Subtype /Link /Border [0 0 0] /H /I /C [1 0 0] /Rect [192.84 492.683 226.299 503.531] /A << /S /GoTo /D <74686D2E442E332E312E33> >> >>}
\MusicPDFWriteObject{1344}{<< /Type /Annot /Subtype /Link /Border [0 0 0] /H /I /C [1 0 0] /Rect [301.338 492.683 326.895 503.531] /A << /S /GoTo /D <73756273656374696F6E2E442E332E31> >> >>}
\MusicPDFWriteObject{1345}{<< /Type /Annot /Subtype /Link /Border [0 0 0] /H /I /C [1 0 0] /Rect [318.515 466.382 351.11 477.172] /A << /S /GoTo /D <74686D2E442E332E312E32> >> >>}
\MusicPDFWriteObject{1346}{<< /Type /Annot /Subtype /Link /Border [0 0 0] /H /I /C [1 0 0] /Rect [327.756 453.231 360.597 464.021] /A << /S /GoTo /D <74686D2E442E332E322E35> >> >>}
\MusicPDFWriteObject{1347}{<< /Type /Annot /Subtype /Link /Border [0 0 0] /H /I /C [1 0 0] /Rect [317.076 394.317 341.618 405.166] /A << /S /GoTo /D <73756273656374696F6E2E422E322E31> >> >>}
\MusicPDFWriteObject{1348}{<< /Type /Annot /Subtype /Link /Border [0 0 0] /H /I /C [1 0 0] /Rect [218.885 316.08 249.099 328.033] /A << /S /GoTo /D <6571756174696F6E2E342E322E312E31> >> >>}
\MusicPDFWriteObject{1349}{<< /Type /Annot /Subtype /Link /Border [0 0 0] /H /I /C [0 1 0] /Rect [93.483 203.443 112.587 212.311] /A << /S /GoTo /D <636974652E4E42> >> >>}
\MusicPDFWriteObject{1350}{<< /Type /Annot /Subtype /Link /Border [0 0 0] /H /I /C [1 0 0] /Rect [105.891 161.5 138.116 173.453] /A << /S /GoTo /D <74686D2E442E332E322E33> >> >>}
\MusicPDFWriteObject{1351}{<< /Type /Annot /Subtype /Link /Border [0 0 0] /H /I /C [1 0 0] /Rect [224.478 122.048 254.691 134.001] /A << /S /GoTo /D <74686D2E342E322E312E33> >> >>}
\MusicPDFWriteObject{1352}{<< /Type /Annot /Subtype /Link /Border [0 0 0] /H /I /C [0 1 0] /Rect [290.326 124.538 313.581 133.406] /A << /S /GoTo /D <636974652E466C4D> >> >>}
\MusicPDFWriteObject{1353}{<< /Type /Annot /Subtype /Link /Border [0 0 0] /H /I /C [0 1 0] /Rect [204.194 85.086 232.823 93.954] /A << /S /GoTo /D <636974652E56616C2D62> >> >>}
\MusicPDFWriteObject{1354}{<< /Type /Annot /Subtype /Link /Border [0 0 0] /H /I /C [1 0 0] /Rect [426.177 685.187 456.954 695.977] /A << /S /GoTo /D <74686D2E342E322E312E35> >> >>}
\MusicPDFWriteObject{1355}{<< /Type /Annot /Subtype /Link /Border [0 0 0] /H /I /C [1 0 0] /Rect [176.323 628.035 209.884 636.889] /A << /S /GoTo /D <74686D2E342E322E312E35> >> >>}
\MusicPDFWriteObject{1356}{<< /Type /Annot /Subtype /Link /Border [0 0 0] /H /I /C [1 0 0] /Rect [234.352 395.508 264.565 407.461] /A << /S /GoTo /D <6571756174696F6E2E342E322E312E39> >> >>}
\MusicPDFWriteObject{1357}{<< /Type /Annot /Subtype /Link /Border [0 0 0] /H /I /C [0 1 0] /Rect [256.237 211.691 275.341 220.558] /A << /S /GoTo /D <636974652E4E42> >> >>}
\MusicPDFWriteObject{1358}{<< /Type /Annot /Subtype /Link /Border [0 0 0] /H /I /C [1 0 0] /Rect [196.483 157.151 229.105 167.942] /A << /S /GoTo /D <74686D2E422E332E322E31> >> >>}
\MusicPDFWriteObject{1359}{<< /Type /Annot /Subtype /Link /Border [0 0 0] /H /I /C [1 0 0] /Rect [222.825 143.446 255.666 155.401] /A << /S /GoTo /D <74686D2E442E332E312E32> >> >>}
\MusicPDFWriteObject{1360}{<< /Type /Annot /Subtype /Link /Border [0 0 0] /H /I /C [1 0 0] /Rect [106.758 103.994 137.14 115.949] /A << /S /GoTo /D <74686D2E342E322E312E33> >> >>}
\MusicPDFWriteObject{1361}{<< /Type /Annot /Subtype /Link /Border [0 0 0] /H /I /C [1 0 0] /Rect [456.344 103.994 489.371 115.949] /A << /S /GoTo /D <74686D2E442E332E322E35> >> >>}
\MusicPDFWriteObject{1362}{<< /Type /Annot /Subtype /Link /Border [0 0 0] /H /I /C [1 0 0] /Rect [71.004 684.633 103.228 696.606] /A << /S /GoTo /D <74686D2E442E332E322E33> >> >>}
\MusicPDFWriteObject{1363}{<< /Type /Annot /Subtype /Link /Border [0 0 0] /H /I /C [0 1 0] /Rect [436.711 617.783 455.816 626.651] /A << /S /GoTo /D <636974652E4E42> >> >>}
\MusicPDFWriteObject{1364}{<< /Type /Annot /Subtype /Link /Border [0 0 0] /H /I /C [0 1 0] /Rect [423.043 578.331 441.241 587.198] /A << /S /GoTo /D <636974652E4A6172> >> >>}
\MusicPDFWriteObject{1365}{<< /Type /Annot /Subtype /Link /Border [0 0 0] /H /I /C [0 1 0] /Rect [112.69 565.18 137.75 574.048] /A << /S /GoTo /D <636974652E53636861> >> >>}
\MusicPDFWriteObject{1366}{<< /Type /Annot /Subtype /Link /Border [0 0 0] /H /I /C [0 1 0] /Rect [287.586 473.125 314.493 481.992] /A << /S /GoTo /D <636974652E47646D> >> >>}
\MusicPDFWriteObject{1367}{<< /Type /Annot /Subtype /Link /Border [0 0 0] /H /I /C [0 1 0] /Rect [399.219 473.125 426.126 481.992] /A << /S /GoTo /D <636974652E47646D> >> >>}
\MusicPDFWriteObject{1368}{<< /Type /Annot /Subtype /Link /Border [0 0 0] /H /I /C [0 1 0] /Rect [151.073 459.974 177.98 468.842] /A << /S /GoTo /D <636974652E47646D> >> >>}
\MusicPDFWriteObject{1369}{<< /Type /Annot /Subtype /Link /Border [0 0 0] /H /I /C [0 1 0] /Rect [73.782 304.552 108.546 313.42] /A << /S /GoTo /D <636974652E436E772D62> >> >>}
\MusicPDFWriteObject{1370}{<< /Type /Annot /Subtype /Link /Border [0 0 0] /H /I /C [0 1 0] /Rect [224.17 183.664 243.274 192.531] /A << /S /GoTo /D <636974652E4E42> >> >>}
\MusicPDFWriteObject{1371}{<< /Type /Annot /Subtype /Link /Border [0 0 0] /H /I /C [1 0 0] /Rect [178.217 682.946 208.43 697.516] /A << /S /GoTo /D <74686D2E342E322E312E32> >> >>}
\MusicPDFWriteObject{1372}{<< /Type /Annot /Subtype /Link /Border [0 0 0] /H /I /C [0 1 0] /Rect [107.382 636.624 126.487 645.491] /A << /S /GoTo /D <636974652E4E42> >> >>}
\MusicPDFWriteObject{1373}{<< /Type /Annot /Subtype /Link /Border [0 0 0] /H /I /C [0 1 0] /Rect [435.343 610.322 454.447 619.19] /A << /S /GoTo /D <636974652E4E42> >> >>}
\MusicPDFWriteObject{1374}{<< /Type /Annot /Subtype /Link /Border [0 0 0] /H /I /C [0 1 0] /Rect [499.367 578.362 517.565 587.23] /A << /S /GoTo /D <636974652E4A6172> >> >>}
\MusicPDFWriteObject{1375}{<< /Type /Annot /Subtype /Link /Border [0 0 0] /H /I /C [1 0 0] /Rect [98.801 511.047 129.015 523] /A << /S /GoTo /D <6571756174696F6E2E342E322E312E34> >> >>}
\MusicPDFWriteObject{1376}{<< /Type /Annot /Subtype /Link /Border [0 0 0] /H /I /C [0 1 0] /Rect [494.118 474.168 543.159 483.036] /A << /S /GoTo /D <636974652E426F752D545653> >> >>}
\MusicPDFWriteObject{1377}{<< /Type /Annot /Subtype /Link /Border [0 0 0] /H /I /C [0 1 0] /Rect [446.356 461.017 465.33 469.885] /A << /S /GoTo /D <636974652E547265> >> >>}
\MusicPDFWriteObject{1378}{<< /Type /Annot /Subtype /Link /Border [0 0 0] /H /I /C [0 1 0] /Rect [190.401 447.866 208.599 456.734] /A << /S /GoTo /D <636974652E4A6172> >> >>}
\MusicPDFWriteObject{1379}{<< /Type /Annot /Subtype /Link /Border [0 0 0] /H /I /C [1 0 0] /Rect [418.625 430.538 448.274 445.108] /A << /S /GoTo /D <74686D2E342E312E322E36> >> >>}
\MusicPDFWriteObject{1380}{<< /Type /Annot /Subtype /Link /Border [0 0 0] /H /I /C [0 1 0] /Rect [156.357 421.565 175.33 430.433] /A << /S /GoTo /D <636974652E547265> >> >>}
\MusicPDFWriteObject{1381}{<< /Type /Annot /Subtype /Link /Border [0 0 0] /H /I /C [1 0 0] /Rect [139.872 405.564 169.916 418.807] /A << /S /GoTo /D <74686D2E342E322E312E35> >> >>}
\MusicPDFWriteObject{1382}{<< /Type /Annot /Subtype /Link /Border [0 0 0] /H /I /C [0 1 0] /Rect [88.532 395.263 111.787 404.131] /A << /S /GoTo /D <636974652E466C4D> >> >>}
\MusicPDFWriteObject{1383}{<< /Type /Annot /Subtype /Link /Border [0 0 0] /H /I /C [0 1 0] /Rect [186.428 382.113 205.401 390.981] /A << /S /GoTo /D <636974652E547265> >> >>}
\MusicPDFWriteObject{1384}{<< /Type /Annot /Subtype /Link /Border [0 0 0] /H /I /C [0 1 0] /Rect [349.387 382.113 368.361 390.981] /A << /S /GoTo /D <636974652E547265> >> >>}
\MusicPDFWriteObject{1385}{<< /Type /Annot /Subtype /Link /Border [0 0 0] /H /I /C [1 0 0] /Rect [424.89 364.784 455.103 379.355] /A << /S /GoTo /D <74686D2E342E322E312E33> >> >>}
\MusicPDFWriteObject{1386}{<< /Type /Annot /Subtype /Link /Border [0 0 0] /H /I /C [1 0 0] /Rect [442.888 352.898 473.101 364.851] /A << /S /GoTo /D <6571756174696F6E2E342E322E312E34> >> >>}
\MusicPDFWriteObject{1387}{<< /Type /Annot /Subtype /Link /Border [0 0 0] /H /I /C [1 0 0] /Rect [352.179 324.844 385.422 338.145] /A << /S /GoTo /D <74686D2E442E312E322E31> >> >>}
\MusicPDFWriteObject{1388}{<< /Type /Annot /Subtype /Link /Border [0 0 0] /H /I /C [1 0 0] /Rect [400.286 250.867 433.18 262.822] /A << /S /GoTo /D <74686D2E422E332E322E31> >> >>}
\MusicPDFWriteObject{1389}{<< /Type /Annot /Subtype /Link /Border [0 0 0] /H /I /C [1 0 0] /Rect [197.384 138.87 232.577 150.823] /A << /S /GoTo /D <6571756174696F6E2E342E322E312E3139> >> >>}
\MusicPDFWriteObject{1390}{<< /Type /Annot /Subtype /Link /Border [0 0 0] /H /I /C [1 0 0] /Rect [356.884 124.032 390.342 137.692] /A << /S /GoTo /D <74686D2E442E332E312E32> >> >>}
\MusicPDFWriteObject{1391}{<< /Type /Annot /Subtype /Link /Border [0 0 0] /H /I /C [0 1 0] /Rect [518.099 128.21 543.159 137.077] /A << /S /GoTo /D <636974652E53636861> >> >>}
\MusicPDFWriteObject{1392}{<< /Type /Annot /Subtype /Link /Border [0 0 0] /H /I /C [0 1 0] /Rect [73.77 101.485 98.831 110.353] /A << /S /GoTo /D <636974652E53636861> >> >>}
\MusicPDFWriteObject{1393}{<< /Type /Annot /Subtype /Link /Border [0 0 0] /H /I /C [1 0 0] /Rect [106.556 672.036 136.769 682.827] /A << /S /GoTo /D <74686D2E342E322E312E35> >> >>}
\MusicPDFWriteObject{1394}{<< /Type /Annot /Subtype /Link /Border [0 0 0] /H /I /C [1 0 0] /Rect [441.385 635.608 471.598 646.399] /A << /S /GoTo /D <74686D2E342E352E332E31> >> >>}
\MusicPDFWriteObject{1395}{<< /Type /Annot /Subtype /Link /Border [0 0 0] /H /I /C [0 1 0] /Rect [107.197 584.942 126.917 593.81] /A << /S /GoTo /D <636974652E4B616C> >> >>}
\MusicPDFWriteObject{1396}{<< /Type /Annot /Subtype /Link /Border [0 0 0] /H /I /C [1 0 0] /Rect [437.345 582.451 467.304 594.406] /A << /S /GoTo /D <74686D2E342E312E322E31> >> >>}
\MusicPDFWriteObject{1397}{<< /Type /Annot /Subtype /Link /Border [0 0 0] /H /I /C [0 1 0] /Rect [523.519 584.942 543.239 593.81] /A << /S /GoTo /D <636974652E4B616C> >> >>}
\MusicPDFWriteObject{1398}{<< /Type /Annot /Subtype /Link /Border [0 0 0] /H /I /C [1 0 0] /Rect [474.515 569.301 505.236 581.254] /A << /S /GoTo /D <74686D2E342E312E322E31> >> >>}
\MusicPDFWriteObject{1399}{<< /Type /Annot /Subtype /Link /Border [0 0 0] /H /I /C [0 1 0] /Rect [199.087 558.64 218.807 567.508] /A << /S /GoTo /D <636974652E4B616C> >> >>}
\MusicPDFWriteObject{1400}{<< /Type /Annot /Subtype /Link /Border [0 0 0] /H /I /C [0 1 0] /Rect [206.407 545.49 226.128 554.357] /A << /S /GoTo /D <636974652E4B616C> >> >>}
\MusicPDFWriteObject{1401}{<< /Type /Annot /Subtype /Link /Border [0 0 0] /H /I /C [0 1 0] /Rect [107.24 388.96 136.118 397.828] /A << /S /GoTo /D <636974652E5363682D61> >> >>}
\MusicPDFWriteObject{1402}{<< /Type /Annot /Subtype /Link /Border [0 0 0] /H /I /C [0 1 0] /Rect [220.253 388.96 245.313 397.828] /A << /S /GoTo /D <636974652E53636861> >> >>}
\MusicPDFWriteObject{1403}{<< /Type /Annot /Subtype /Link /Border [0 0 0] /H /I /C [0 1 0] /Rect [516.667 231.151 543.159 240.019] /A << /S /GoTo /D <636974652E52752D32> >> >>}
\MusicPDFWriteObject{1404}{<< /Type /Annot /Subtype /Link /Border [0 0 0] /H /I /C [0 1 0] /Rect [188.249 653.181 214.74 662.049] /A << /S /GoTo /D <636974652E52752D32> >> >>}
\MusicPDFWriteObject{1405}{<< /Type /Annot /Subtype /Link /Border [0 0 0] /H /I /C [1 0 0] /Rect [363.023 650.691 392.671 662.644] /A << /S /GoTo /D <74686D2E342E322E322E31> >> >>}
\MusicPDFWriteObject{1406}{<< /Type /Annot /Subtype /Link /Border [0 0 0] /H /I /C [1 0 0] /Rect [185.601 248.803 218.443 260.756] /A << /S /GoTo /D <74686D2E442E332E332E31> >> >>}
\MusicPDFWriteObject{1407}{<< /Type /Annot /Subtype /Link /Border [0 0 0] /H /I /C [1 0 0] /Rect [71.004 198.69 103.845 207.544] /A << /S /GoTo /D <74686D2E442E332E312E32> >> >>}
\MusicPDFWriteObject{1408}{<< /Type /Annot /Subtype /Link /Border [0 0 0] /H /I /C [0 1 0] /Rect [73.77 159.238 92.875 168.105] /A << /S /GoTo /D <636974652E4E42> >> >>}
\MusicPDFWriteObject{1409}{<< /Type /Annot /Subtype /Link /Border [0 0 0] /H /I /C [1 0 0] /Rect [203.197 137.689 236.758 146.543] /A << /S /GoTo /D <74686D2E342E322E322E32> >> >>}
\MusicPDFWriteObject{1410}{<< /Type /Annot /Subtype /Link /Border [0 0 0] /H /I /C [0 1 0] /Rect [450.278 124.538 469.383 133.406] /A << /S /GoTo /D <636974652E4E42> >> >>}
\MusicPDFWriteObject{1411}{<< /Type /Annot /Subtype /Link /Border [0 0 0] /H /I /C [1 0 0] /Rect [73.712 82.596 103.925 94.549] /A << /S /GoTo /D <6571756174696F6E2E342E322E322E32> >> >>}
\MusicPDFWriteObject{1412}{<< /Type /Annot /Subtype /Link /Border [0 0 0] /H /I /C [1 0 0] /Rect [148.851 671.482 179.064 683.435] /A << /S /GoTo /D <6571756174696F6E2E342E322E322E32> >> >>}
\MusicPDFWriteObject{1413}{<< /Type /Annot /Subtype /Link /Border [0 0 0] /H /I /C [0 1 0] /Rect [298.063 481.915 333.118 490.782] /A << /S /GoTo /D <636974652E426F752D4949> >> >>}
\MusicPDFWriteObject{1414}{<< /Type /Annot /Subtype /Link /Border [0 0 0] /H /I /C [0 1 0] /Rect [92.027 350.407 118.519 359.275] /A << /S /GoTo /D <636974652E52752D32> >> >>}
\MusicPDFWriteObject{1415}{<< /Type /Annot /Subtype /Link /Border [0 0 0] /H /I /C [1 0 0] /Rect [390.419 196.441 423.313 207.232] /A << /S /GoTo /D <74686D2E422E332E312E31> >> >>}
\MusicPDFWriteObject{1416}{<< /Type /Annot /Subtype /Link /Border [0 0 0] /H /I /C [1 0 0] /Rect [246.129 182.376 277.811 194.691] /A << /S /GoTo /D <74686D2E422E332E322E31> >> >>}
\MusicPDFWriteObject{1417}{<< /Type /Annot /Subtype /Link /Border [0 0 0] /H /I /C [1 0 0] /Rect [513.621 182.376 543.27 194.691] /A << /S /GoTo /D <74686D2E342E322E312E35> >> >>}
\MusicPDFWriteObject{1418}{<< /Type /Annot /Subtype /Link /Border [0 0 0] /H /I /C [1 0 0] /Rect [407.962 96.748 438.175 108.701] /A << /S /GoTo /D <6571756174696F6E2E342E322E322E35> >> >>}
\MusicPDFWriteObject{1419}{<< /Type /Annot /Subtype /Link /Border [0 0 0] /H /I /C [1 0 0] /Rect [356.802 671.482 387.016 683.435] /A << /S /GoTo /D <6571756174696F6E2E342E322E322E33> >> >>}
\MusicPDFWriteObject{1420}{<< /Type /Annot /Subtype /Link /Border [0 0 0] /H /I /C [0 1 0] /Rect [516.667 611.243 543.159 620.111] /A << /S /GoTo /D <636974652E52752D32> >> >>}
\MusicPDFWriteObject{1421}{<< /Type /Annot /Subtype /Link /Border [0 0 0] /H /I /C [1 0 0] /Rect [202.898 508.239 233.111 519.03] /A << /S /GoTo /D <74686D2E342E322E312E35> >> >>}
\MusicPDFWriteObject{1422}{<< /Type /Annot /Subtype /Link /Border [0 0 0] /H /I /C [1 0 0] /Rect [307.425 508.239 337.638 519.03] /A << /S /GoTo /D <74686D2E342E322E322E32> >> >>}
\MusicPDFWriteObject{1423}{<< /Type /Annot /Subtype /Link /Border [0 0 0] /H /I /C [1 0 0] /Rect [188.813 326.244 219.138 338.197] /A << /S /GoTo /D <74686D2E342E322E312E35> >> >>}
\MusicPDFWriteObject{1424}{<< /Type /Annot /Subtype /Link /Border [0 0 0] /H /I /C [1 0 0] /Rect [323.107 326.244 353.433 338.197] /A << /S /GoTo /D <74686D2E342E322E322E32> >> >>}
\MusicPDFWriteObject{1425}{<< /Type /Annot /Subtype /Link /Border [0 0 0] /H /I /C [1 0 0] /Rect [197.597 263.515 227.81 275.468] /A << /S /GoTo /D <6571756174696F6E2E342E322E322E37> >> >>}
\MusicPDFWriteObject{1426}{<< /Type /Annot /Subtype /Link /Border [0 0 0] /H /I /C [1 0 0] /Rect [452.882 263.515 482.756 275.468] /A << /S /GoTo /D <74686D2E342E322E322E32> >> >>}
\MusicPDFWriteObject{1427}{<< /Type /Annot /Subtype /Link /Border [0 0 0] /H /I /C [1 0 0] /Rect [320.917 215.375 350.791 226.165] /A << /S /GoTo /D <74686D2E342E322E312E35> >> >>}
\MusicPDFWriteObject{1428}{<< /Type /Annot /Subtype /Link /Border [0 0 0] /H /I /C [1 0 0] /Rect [424.309 215.375 454.184 226.165] /A << /S /GoTo /D <74686D2E342E322E322E32> >> >>}
\MusicPDFWriteObject{1429}{<< /Type /Annot /Subtype /Link /Border [0 0 0] /H /I /C [0 1 0] /Rect [125.602 85.086 158.776 93.954] /A << /S /GoTo /D <636974652E53696D6F6E> >> >>}
\MusicPDFWriteObject{1430}{<< /Type /Annot /Subtype /Link /Border [0 0 0] /H /I /C [1 0 0] /Rect [71.004 523.081 104.462 533.957] /A << /S /GoTo /D <74686D2E442E322E312E31> >> >>}
\MusicPDFWriteObject{1431}{<< /Type /Annot /Subtype /Link /Border [0 0 0] /H /I /C [1 0 0] /Rect [231.388 496.864 256.945 507.712] /A << /S /GoTo /D <73756273656374696F6E2E442E332E33> >> >>}
\MusicPDFWriteObject{1432}{<< /Type /Annot /Subtype /Link /Border [0 0 0] /H /I /C [1 0 0] /Rect [477.487 496.864 510.945 507.712] /A << /S /GoTo /D <74686D2E442E332E332E37> >> >>}
\MusicPDFWriteObject{1433}{<< /Type /Annot /Subtype /Link /Border [0 0 0] /H /I /C [1 0 0] /Rect [425.235 361.355 455.448 372.145] /A << /S /GoTo /D <74686D2E342E332E302E31> >> >>}
\MusicPDFWriteObject{1434}{<< /Type /Annot /Subtype /Link /Border [0 0 0] /H /I /C [0 1 0] /Rect [258.655 137.689 275.809 146.557] /A << /S /GoTo /D <636974652E466F6C> >> >>}
\MusicPDFWriteObject{1435}{<< /Type /Annot /Subtype /Link /Border [0 0 0] /H /I /C [1 0 0] /Rect [172.91 456.374 205.613 468.326] /A << /S /GoTo /D <6571756174696F6E2E412E312E332E32> >> >>}
\MusicPDFWriteObject{1436}{<< /Type /Annot /Subtype /Link /Border [0 0 0] /H /I /C [1 0 0] /Rect [513.621 391.088 543.27 401.964] /A << /S /GoTo /D <74686D2E342E332E302E31> >> >>}
\MusicPDFWriteObject{1437}{<< /Type /Annot /Subtype /Link /Border [0 0 0] /H /I /C [1 0 0] /Rect [185.278 351.167 218.119 363.122] /A << /S /GoTo /D <74686D2E442E322E312E32> >> >>}
\MusicPDFWriteObject{1438}{<< /Type /Annot /Subtype /Link /Border [0 0 0] /H /I /C [0 1 0] /Rect [296.304 278.31 314.606 287.178] /A << /S /GoTo /D <636974652E4445> >> >>}
\MusicPDFWriteObject{1439}{<< /Type /Annot /Subtype /Link /Border [0 0 0] /H /I /C [0 1 0] /Rect [526.005 278.31 543.159 287.178] /A << /S /GoTo /D <636974652E466F6C> >> >>}
\MusicPDFWriteObject{1440}{<< /Type /Annot /Subtype /Link /Border [0 0 0] /H /I /C [1 0 0] /Rect [338.282 241.936 368.326 253.892] /A << /S /GoTo /D <74686D2E342E332E312E33> >> >>}
\MusicPDFWriteObject{1441}{<< /Type /Annot /Subtype /Link /Border [0 0 0] /H /I /C [0 1 0] /Rect [251.925 218.126 269.079 226.993] /A << /S /GoTo /D <636974652E466F6C> >> >>}
\MusicPDFWriteObject{1442}{<< /Type /Annot /Subtype /Link /Border [0 0 0] /H /I /C [0 1 0] /Rect [332.728 218.126 351.03 226.993] /A << /S /GoTo /D <636974652E4445> >> >>}
\MusicPDFWriteObject{1443}{<< /Type /Annot /Subtype /Link /Border [0 0 0] /H /I /C [1 0 0] /Rect [134.27 625.034 164.484 636.986] /A << /S /GoTo /D <6571756174696F6E2E342E332E312E34> >> >>}
\MusicPDFWriteObject{1444}{<< /Type /Annot /Subtype /Link /Border [0 0 0] /H /I /C [1 0 0] /Rect [124.082 460.943 154.295 472.896] /A << /S /GoTo /D <6571756174696F6E2E342E332E312E34> >> >>}
\MusicPDFWriteObject{1445}{<< /Type /Annot /Subtype /Link /Border [0 0 0] /H /I /C [0 1 0] /Rect [509.19 374.403 543.159 383.27] /A << /S /GoTo /D <636974652E436E772D61> >> >>}
\MusicPDFWriteObject{1446}{<< /Type /Annot /Subtype /Link /Border [0 0 0] /H /I /C [0 1 0] /Rect [518.118 256.046 543.178 264.914] /A << /S /GoTo /D <636974652E53636861> >> >>}
\MusicPDFWriteObject{1447}{<< /Type /Annot /Subtype /Link /Border [0 0 0] /H /I /C [0 1 0] /Rect [215.863 242.895 249.832 251.763] /A << /S /GoTo /D <636974652E436E772D61> >> >>}
\MusicPDFWriteObject{1448}{<< /Type /Annot /Subtype /Link /Border [0 0 0] /H /I /C [0 1 0] /Rect [420.865 216.594 439.195 225.461] /A << /S /GoTo /D <636974652E4555> >> >>}
\MusicPDFWriteObject{1449}{<< /Type /Annot /Subtype /Link /Border [0 0 0] /H /I /C [0 1 0] /Rect [395.407 190.292 429.376 199.16] /A << /S /GoTo /D <636974652E436E772D61> >> >>}
\MusicPDFWriteObject{1450}{<< /Type /Annot /Subtype /Link /Border [0 0 0] /H /I /C [0 1 0] /Rect [395.012 150.84 428.981 159.708] /A << /S /GoTo /D <636974652E436E772D61> >> >>}
\MusicPDFWriteObject{1451}{<< /Type /Annot /Subtype /Link /Border [0 0 0] /H /I /C [0 1 0] /Rect [211.442 137.689 245.411 146.557] /A << /S /GoTo /D <636974652E436E772D61> >> >>}
\MusicPDFWriteObject{1452}{<< /Type /Annot /Subtype /Link /Border [0 0 0] /H /I /C [0 1 0] /Rect [509.301 111.387 543.27 120.255] /A << /S /GoTo /D <636974652E436E772D61> >> >>}
\MusicPDFWriteObject{1453}{<< /Type /Annot /Subtype /Link /Border [0 0 0] /H /I /C [0 1 0] /Rect [185.179 98.237 219.148 107.105] /A << /S /GoTo /D <636974652E436E772D61> >> >>}
\MusicPDFWriteObject{1454}{<< /Type /Annot /Subtype /Link /Border [0 0 0] /H /I /C [0 1 0] /Rect [124.512 85.086 158.481 93.954] /A << /S /GoTo /D <636974652E436E772D61> >> >>}
\MusicPDFWriteObject{1455}{<< /Type /Annot /Subtype /Link /Border [0 0 0] /H /I /C [0 1 0] /Rect [196.124 637.545 230.093 646.413] /A << /S /GoTo /D <636974652E436E772D61> >> >>}
\MusicPDFWriteObject{1456}{<< /Type /Annot /Subtype /Link /Border [0 0 0] /H /I /C [0 1 0] /Rect [341.587 637.545 375.556 646.413] /A << /S /GoTo /D <636974652E436E772D61> >> >>}
\MusicPDFWriteObject{1457}{<< /Type /Annot /Subtype /Link /Border [0 0 0] /H /I /C [1 0 0] /Rect [499.4 635.055 530.177 647.008] /A << /S /GoTo /D <74686D2E342E332E312E35> >> >>}
\MusicPDFWriteObject{1458}{<< /Type /Annot /Subtype /Link /Border [0 0 0] /H /I /C [0 1 0] /Rect [140.22 598.093 159.006 606.961] /A << /S /GoTo /D <636974652E4850> >> >>}
\MusicPDFWriteObject{1459}{<< /Type /Annot /Subtype /Link /Border [0 0 0] /H /I /C [0 1 0] /Rect [441.655 598.093 475.623 606.961] /A << /S /GoTo /D <636974652E436E772D61> >> >>}
\MusicPDFWriteObject{1460}{<< /Type /Annot /Subtype /Link /Border [0 0 0] /H /I /C [0 1 0] /Rect [104.288 584.942 123.075 593.81] /A << /S /GoTo /D <636974652E4850> >> >>}
\MusicPDFWriteObject{1461}{<< /Type /Annot /Subtype /Link /Border [0 0 0] /H /I /C [0 1 0] /Rect [92.783 558.64 111.113 567.508] /A << /S /GoTo /D <636974652E4555> >> >>}
\MusicPDFWriteObject{1462}{<< /Type /Annot /Subtype /Link /Border [0 0 0] /H /I /C [0 1 0] /Rect [421.12 558.64 439.45 567.508] /A << /S /GoTo /D <636974652E4555> >> >>}
\MusicPDFWriteObject{1463}{<< /Type /Annot /Subtype /Link /Border [0 0 0] /H /I /C [0 1 0] /Rect [204.358 532.339 229.418 541.207] /A << /S /GoTo /D <636974652E53636861> >> >>}
\MusicPDFWriteObject{1464}{<< /Type /Annot /Subtype /Link /Border [0 0 0] /H /I /C [0 1 0] /Rect [526.237 532.339 543.253 541.207] /A << /S /GoTo /D <636974652E4153> >> >>}
\MusicPDFWriteObject{1465}{<< /Type /Annot /Subtype /Link /Border [0 0 0] /H /I /C [1 0 0] /Rect [223.982 504.1 248.207 514.948] /A << /S /GoTo /D <73756273656374696F6E2E432E332E33> >> >>}
\MusicPDFWriteObject{1466}{<< /Type /Annot /Subtype /Link /Border [0 0 0] /H /I /C [1 0 0] /Rect [272.804 504.1 297.03 514.948] /A << /S /GoTo /D <73756273656374696F6E2E432E332E34> >> >>}
\MusicPDFWriteObject{1467}{<< /Type /Annot /Subtype /Link /Border [0 0 0] /H /I /C [1 0 0] /Rect [176.323 471.797 209.884 480.651] /A << /S /GoTo /D <74686D2E342E332E312E35> >> >>}
\MusicPDFWriteObject{1468}{<< /Type /Annot /Subtype /Link /Border [0 0 0] /H /I /C [1 0 0] /Rect [70.625 456.709 103.175 467.5] /A << /S /GoTo /D <74686D2E412E352E322E39> >> >>}
\MusicPDFWriteObject{1469}{<< /Type /Annot /Subtype /Link /Border [0 0 0] /H /I /C [1 0 0] /Rect [442.355 456.709 472.427 467.5] /A << /S /GoTo /D <74686D2E342E332E302E31> >> >>}
\MusicPDFWriteObject{1470}{<< /Type /Annot /Subtype /Link /Border [0 0 0] /H /I /C [1 0 0] /Rect [70.625 443.005 108.88 454.96] /A << /S /GoTo /D <74686D2E412E322E332E3130> >> >>}
\MusicPDFWriteObject{1471}{<< /Type /Annot /Subtype /Link /Border [0 0 0] /H /I /C [1 0 0] /Rect [349.412 443.005 380.077 454.96] /A << /S /GoTo /D <74686D2E342E332E312E33> >> >>}
\MusicPDFWriteObject{1472}{<< /Type /Annot /Subtype /Link /Border [0 0 0] /H /I /C [0 1 0] /Rect [73.77 432.344 96.126 441.212] /A << /S /GoTo /D <636974652E436F68> >> >>}
\MusicPDFWriteObject{1473}{<< /Type /Annot /Subtype /Link /Border [0 0 0] /H /I /C [0 1 0] /Rect [469.773 328.94 491.368 339.744] /A << /S /GoTo /D <636974652E426F67> >> >>}
\MusicPDFWriteObject{1474}{<< /Type /Annot /Subtype /Link /Border [0 0 0] /H /I /C [1 0 0] /Rect [145.005 302.085 175.218 314.038] /A << /S /GoTo /D <6571756174696F6E2E342E332E312E35> >> >>}
\MusicPDFWriteObject{1475}{<< /Type /Annot /Subtype /Link /Border [0 0 0] /H /I /C [0 1 0] /Rect [195.1 265.123 213.402 273.991] /A << /S /GoTo /D <636974652E4445> >> >>}
\MusicPDFWriteObject{1476}{<< /Type /Annot /Subtype /Link /Border [0 0 0] /H /I /C [1 0 0] /Rect [458.729 210.029 488.378 221.984] /A << /S /GoTo /D <74686D2E342E332E312E33> >> >>}
\MusicPDFWriteObject{1477}{<< /Type /Annot /Subtype /Link /Border [0 0 0] /H /I /C [1 0 0] /Rect [70.655 196.879 100.868 208.832] /A << /S /GoTo /D <74686D2E342E332E312E34> >> >>}
\MusicPDFWriteObject{1478}{<< /Type /Annot /Subtype /Link /Border [0 0 0] /H /I /C [1 0 0] /Rect [207.666 144.276 237.879 156.229] /A << /S /GoTo /D <6571756174696F6E2E342E332E312E35> >> >>}
\MusicPDFWriteObject{1479}{<< /Type /Annot /Subtype /Link /Border [0 0 0] /H /I /C [1 0 0] /Rect [226.886 685.186 251.981 696.034] /A << /S /GoTo /D <73756273656374696F6E2E442E322E31> >> >>}
\MusicPDFWriteObject{1480}{<< /Type /Annot /Subtype /Link /Border [0 0 0] /H /I /C [1 0 0] /Rect [70.655 657.971 101.433 669.676] /A << /S /GoTo /D <74686D2E342E332E302E31> >> >>}
\MusicPDFWriteObject{1481}{<< /Type /Annot /Subtype /Link /Border [0 0 0] /H /I /C [1 0 0] /Rect [124.429 618.879 157.132 630.832] /A << /S /GoTo /D <74686D2E412E312E342E33> >> >>}
\MusicPDFWriteObject{1482}{<< /Type /Annot /Subtype /Link /Border [0 0 0] /H /I /C [1 0 0] /Rect [357.279 548.328 387.126 559.12] /A << /S /GoTo /D <74686D2E342E332E302E31> >> >>}
\MusicPDFWriteObject{1483}{<< /Type /Annot /Subtype /Link /Border [0 0 0] /H /I /C [1 0 0] /Rect [145.576 521.473 179.035 534.066] /A << /S /GoTo /D <74686D2E442E312E322E31> >> >>}
\MusicPDFWriteObject{1484}{<< /Type /Annot /Subtype /Link /Border [0 0 0] /H /I /C [1 0 0] /Rect [175.369 482.022 208.211 493.974] /A << /S /GoTo /D <74686D2E442E322E312E31> >> >>}
\MusicPDFWriteObject{1485}{<< /Type /Annot /Subtype /Link /Border [0 0 0] /H /I /C [1 0 0] /Rect [262.549 482.022 295.39 493.974] /A << /S /GoTo /D <74686D2E442E322E312E33> >> >>}
\MusicPDFWriteObject{1486}{<< /Type /Annot /Subtype /Link /Border [0 0 0] /H /I /C [1 0 0] /Rect [353.602 482.022 386.443 493.974] /A << /S /GoTo /D <74686D2E442E322E312E34> >> >>}
\MusicPDFWriteObject{1487}{<< /Type /Annot /Subtype /Link /Border [0 0 0] /H /I /C [0 1 0] /Rect [215.012 687.123 248.981 695.991] /A << /S /GoTo /D <636974652E436E772D61> >> >>}
\MusicPDFWriteObject{1488}{<< /Type /Annot /Subtype /Link /Border [0 0 0] /H /I /C [1 0 0] /Rect [239.793 640.253 272.465 652.566] /A << /S /GoTo /D <74686D2E412E312E342E33> >> >>}
\MusicPDFWriteObject{1489}{<< /Type /Annot /Subtype /Link /Border [0 0 0] /H /I /C [1 0 0] /Rect [497.83 640.253 530.503 652.566] /A << /S /GoTo /D <74686D2E412E312E352E32> >> >>}
\MusicPDFWriteObject{1490}{<< /Type /Annot /Subtype /Link /Border [0 0 0] /H /I /C [0 1 0] /Rect [233.073 599.9 267.042 608.768] /A << /S /GoTo /D <636974652E436E772D61> >> >>}
\MusicPDFWriteObject{1491}{<< /Type /Annot /Subtype /Link /Border [0 0 0] /H /I /C [0 1 0] /Rect [282.926 559.557 300.592 568.425] /A << /S /GoTo /D <636974652E4154> >> >>}
\MusicPDFWriteObject{1492}{<< /Type /Annot /Subtype /Link /Border [0 0 0] /H /I /C [0 1 0] /Rect [174.58 532.365 192.883 541.233] /A << /S /GoTo /D <636974652E4445> >> >>}
\MusicPDFWriteObject{1493}{<< /Type /Annot /Subtype /Link /Border [0 0 0] /H /I /C [0 1 0] /Rect [525.604 506.064 543.27 514.932] /A << /S /GoTo /D <636974652E4154> >> >>}
\MusicPDFWriteObject{1494}{<< /Type /Annot /Subtype /Link /Border [0 0 0] /H /I /C [1 0 0] /Rect [450.527 378.256 480.74 390.208] /A << /S /GoTo /D <6571756174696F6E2E342E332E322E37> >> >>}
\MusicPDFWriteObject{1495}{<< /Type /Annot /Subtype /Link /Border [0 0 0] /H /I /C [0 1 0] /Rect [360.923 349.877 378.257 358.745] /A << /S /GoTo /D <636974652E4B6E> >> >>}
\MusicPDFWriteObject{1496}{<< /Type /Annot /Subtype /Link /Border [0 0 0] /H /I /C [1 0 0] /Rect [161.988 227.417 192.201 239.37] /A << /S /GoTo /D <6571756174696F6E2E342E332E322E38> >> >>}
\MusicPDFWriteObject{1497}{<< /Type /Annot /Subtype /Link /Border [0 0 0] /H /I /C [1 0 0] /Rect [273.746 214.266 303.959 226.219] /A << /S /GoTo /D <6571756174696F6E2E342E332E322E38> >> >>}
\MusicPDFWriteObject{1498}{<< /Type /Annot /Subtype /Link /Border [0 0 0] /H /I /C [1 0 0] /Rect [98.422 82.596 128.635 94.549] /A << /S /GoTo /D <6571756174696F6E2E342E332E322E38> >> >>}
\MusicPDFWriteObject{1499}{<< /Type /Annot /Subtype /Link /Border [0 0 0] /H /I /C [1 0 0] /Rect [161.897 448.221 192.11 460.174] /A << /S /GoTo /D <6571756174696F6E2E342E332E322E37> >> >>}
\MusicPDFWriteObject{1500}{<< /Type /Annot /Subtype /Link /Border [0 0 0] /H /I /C [1 0 0] /Rect [323.444 383.848 340.078 395.801] /A << /S /GoTo /D <73656374696F6E2E432E34> >> >>}
\MusicPDFWriteObject{1501}{<< /Type /Annot /Subtype /Link /Border [0 0 0] /H /I /C [1 0 0] /Rect [438.635 383.848 468.283 395.801] /A << /S /GoTo /D <74686D2E342E332E332E31> >> >>}
\MusicPDFWriteObject{1502}{<< /Type /Annot /Subtype /Link /Border [0 0 0] /H /I /C [1 0 0] /Rect [355.676 278.642 388.102 290.595] /A << /S /GoTo /D <74686D2E432E342E332E34> >> >>}
\MusicPDFWriteObject{1503}{<< /Type /Annot /Subtype /Link /Border [0 0 0] /H /I /C [1 0 0] /Rect [371.199 252.894 395.583 263.742] /A << /S /GoTo /D <73756273656374696F6E2E432E342E33> >> >>}
\MusicPDFWriteObject{1504}{<< /Type /Annot /Subtype /Link /Border [0 0 0] /H /I /C [0 1 0] /Rect [314.112 241.68 363.153 250.548] /A << /S /GoTo /D <636974652E426F752D545653> >> >>}
\MusicPDFWriteObject{1505}{<< /Type /Annot /Subtype /Link /Border [0 0 0] /H /I /C [1 0 0] /Rect [493.739 226.039 526.774 237.992] /A << /S /GoTo /D <74686D2E432E342E312E32> >> >>}
\MusicPDFWriteObject{1506}{<< /Type /Annot /Subtype /Link /Border [0 0 0] /H /I /C [1 0 0] /Rect [520.461 213.441 544.845 224.29] /A << /S /GoTo /D <73756273656374696F6E2E432E342E33> >> >>}
\MusicPDFWriteObject{1507}{<< /Type /Annot /Subtype /Link /Border [0 0 0] /H /I /C [1 0 0] /Rect [131.896 199.738 164.323 211.691] /A << /S /GoTo /D <74686D2E432E342E332E35> >> >>}
\MusicPDFWriteObject{1508}{<< /Type /Annot /Subtype /Link /Border [0 0 0] /H /I /C [1 0 0] /Rect [192.142 187.14 208.851 197.988] /A << /S /GoTo /D <73656374696F6E2E432E32> >> >>}
\MusicPDFWriteObject{1509}{<< /Type /Annot /Subtype /Link /Border [0 0 0] /H /I /C [1 0 0] /Rect [217.344 187.14 241.684 197.988] /A << /S /GoTo /D <73756273656374696F6E2E432E352E31> >> >>}
\MusicPDFWriteObject{1510}{<< /Type /Annot /Subtype /Link /Border [0 0 0] /H /I /C [1 0 0] /Rect [266.586 187.14 290.925 197.988] /A << /S /GoTo /D <73756273656374696F6E2E432E332E33> >> >>}
\MusicPDFWriteObject{1511}{<< /Type /Annot /Subtype /Link /Border [0 0 0] /H /I /C [1 0 0] /Rect [407.051 359.607 423.984 370.455] /A << /S /GoTo /D <73656374696F6E2E432E34> >> >>}
\MusicPDFWriteObject{1512}{<< /Type /Annot /Subtype /Link /Border [0 0 0] /H /I /C [1 0 0] /Rect [527.183 214.948 543.817 225.796] /A << /S /GoTo /D <73656374696F6E2E432E34> >> >>}
\MusicPDFWriteObject{1513}{<< /Type /Annot /Subtype /Link /Border [0 0 0] /H /I /C [1 0 0] /Rect [73.712 181.373 103.925 193.325] /A << /S /GoTo /D <6571756174696F6E2E342E332E332E32> >> >>}
\MusicPDFWriteObject{1514}{<< /Type /Annot /Subtype /Link /Border [0 0 0] /H /I /C [1 0 0] /Rect [176.323 163.99 209.884 172.844] /A << /S /GoTo /D <74686D2E342E332E332E31> >> >>}
\MusicPDFWriteObject{1515}{<< /Type /Annot /Subtype /Link /Border [0 0 0] /H /I /C [1 0 0] /Rect [294.412 161.5 326.23 173.453] /A << /S /GoTo /D <74686D2E432E342E332E35> >> >>}
\MusicPDFWriteObject{1516}{<< /Type /Annot /Subtype /Link /Border [0 0 0] /H /I /C [1 0 0] /Rect [140.277 96.3 177.826 107.091] /A << /S /GoTo /D <74686D2E432E342E332E3132> >> >>}
\MusicPDFWriteObject{1517}{<< /Type /Annot /Subtype /Link /Border [0 0 0] /H /I /C [1 0 0] /Rect [303.462 684.633 340.443 696.586] /A << /S /GoTo /D <74686D2E432E342E332E3132> >> >>}
\MusicPDFWriteObject{1518}{<< /Type /Annot /Subtype /Link /Border [0 0 0] /H /I /C [0 1 0] /Rect [73.826 673.973 92.931 682.84] /A << /S /GoTo /D <636974652E4E42> >> >>}
\MusicPDFWriteObject{1519}{<< /Type /Annot /Subtype /Link /Border [0 0 0] /H /I /C [0 1 0] /Rect [377.62 660.822 396.725 669.69] /A << /S /GoTo /D <636974652E4E42> >> >>}
\MusicPDFWriteObject{1520}{<< /Type /Annot /Subtype /Link /Border [0 0 0] /H /I /C [0 1 0] /Rect [311.548 647.671 330.653 656.539] /A << /S /GoTo /D <636974652E4E42> >> >>}
\MusicPDFWriteObject{1521}{<< /Type /Annot /Subtype /Link /Border [0 0 0] /H /I /C [0 1 0] /Rect [453.151 647.671 472.256 656.539] /A << /S /GoTo /D <636974652E4E42> >> >>}
\MusicPDFWriteObject{1522}{<< /Type /Annot /Subtype /Link /Border [0 0 0] /H /I /C [1 0 0] /Rect [277.142 632.584 309.234 643.374] /A << /S /GoTo /D <74686D2E432E342E332E34> >> >>}
\MusicPDFWriteObject{1523}{<< /Type /Annot /Subtype /Link /Border [0 0 0] /H /I /C [1 0 0] /Rect [419.212 543 449.425 554.952] /A << /S /GoTo /D <6571756174696F6E2E342E332E332E33> >> >>}
\MusicPDFWriteObject{1524}{<< /Type /Annot /Subtype /Link /Border [0 0 0] /H /I /C [1 0 0] /Rect [368.245 498.029 398.458 509.982] /A << /S /GoTo /D <6571756174696F6E2E342E332E332E34> >> >>}
\MusicPDFWriteObject{1525}{<< /Type /Annot /Subtype /Link /Border [0 0 0] /H /I /C [0 1 0] /Rect [89.714 357.842 132.737 366.71] /A << /S /GoTo /D <636974652E426F752D4754> >> >>}
\MusicPDFWriteObject{1526}{<< /Type /Annot /Subtype /Link /Border [0 0 0] /H /I /C [0 1 0] /Rect [298.817 292.088 317.922 300.956] /A << /S /GoTo /D <636974652E4E42> >> >>}
\MusicPDFWriteObject{1527}{<< /Type /Annot /Subtype /Link /Border [0 0 0] /H /I /C [1 0 0] /Rect [407.373 247.614 437.586 259.567] /A << /S /GoTo /D <6571756174696F6E2E342E332E332E35> >> >>}
\MusicPDFWriteObject{1528}{<< /Type /Annot /Subtype /Link /Border [0 0 0] /H /I /C [1 0 0] /Rect [257.104 590.89 287.317 602.842] /A << /S /GoTo /D <6571756174696F6E2E342E332E332E36> >> >>}
\MusicPDFWriteObject{1529}{<< /Type /Annot /Subtype /Link /Border [0 0 0] /H /I /C [0 1 0] /Rect [107.816 441.62 126.921 450.488] /A << /S /GoTo /D <636974652E4E42> >> >>}
\MusicPDFWriteObject{1530}{<< /Type /Annot /Subtype /Link /Border [0 0 0] /H /I /C [0 1 0] /Rect [111.238 428.469 129.437 437.337] /A << /S /GoTo /D <636974652E4A6172> >> >>}
\MusicPDFWriteObject{1531}{<< /Type /Annot /Subtype /Link /Border [0 0 0] /H /I /C [1 0 0] /Rect [328.092 381.654 358.305 392.445] /A << /S /GoTo /D <74686D2E342E332E332E31> >> >>}
\MusicPDFWriteObject{1532}{<< /Type /Annot /Subtype /Link /Border [0 0 0] /H /I /C [1 0 0] /Rect [412.976 645.18 443.19 657.136] /A << /S /GoTo /D <74686D2E342E332E332E33> >> >>}
\MusicPDFWriteObject{1533}{<< /Type /Annot /Subtype /Link /Border [0 0 0] /H /I /C [1 0 0] /Rect [196.747 567.435 230.308 576.289] /A << /S /GoTo /D <74686D2E342E332E332E33> >> >>}
\MusicPDFWriteObject{1534}{<< /Type /Annot /Subtype /Link /Border [0 0 0] /H /I /C [1 0 0] /Rect [321.1 498.831 350.748 511.144] /A << /S /GoTo /D <74686D2E342E332E332E31> >> >>}
\MusicPDFWriteObject{1535}{<< /Type /Annot /Subtype /Link /Border [0 0 0] /H /I /C [1 0 0] /Rect [266.41 390.979 296.623 402.932] /A << /S /GoTo /D <74686D2E342E332E332E31> >> >>}
\MusicPDFWriteObject{1536}{<< /Type /Annot /Subtype /Link /Border [0 0 0] /H /I /C [1 0 0] /Rect [171.74 314.739 201.953 327.052] /A << /S /GoTo /D <74686D2E342E332E332E31> >> >>}
\MusicPDFWriteObject{1537}{<< /Type /Annot /Subtype /Link /Border [0 0 0] /H /I /C [1 0 0] /Rect [294.754 302.502 324.967 313.292] /A << /S /GoTo /D <74686D2E342E332E332E31> >> >>}
\MusicPDFWriteObject{1538}{<< /Type /Annot /Subtype /Link /Border [0 0 0] /H /I /C [1 0 0] /Rect [170.418 238.662 200.632 250.615] /A << /S /GoTo /D <6571756174696F6E2E342E332E322E31> >> >>}
\MusicPDFWriteObject{1539}{<< /Type /Annot /Subtype /Link /Border [0 0 0] /H /I /C [1 0 0] /Rect [260.752 238.301 290.965 250.615] /A << /S /GoTo /D <74686D2E342E332E322E32> >> >>}
\MusicPDFWriteObject{1540}{<< /Type /Annot /Subtype /Link /Border [0 0 0] /H /I /C [1 0 0] /Rect [356.674 238.301 386.887 250.615] /A << /S /GoTo /D <74686D2E342E332E332E31> >> >>}
\MusicPDFWriteObject{1541}{<< /Type /Annot /Subtype /Link /Border [0 0 0] /H /I /C [1 0 0] /Rect [451.628 238.301 481.841 250.615] /A << /S /GoTo /D <74686D2E342E332E332E33> >> >>}
\MusicPDFWriteObject{1542}{<< /Type /Annot /Subtype /Link /Border [0 0 0] /H /I /C [1 0 0] /Rect [383.775 122.048 413.988 134.001] /A << /S /GoTo /D <6571756174696F6E2E342E332E322E31> >> >>}
\MusicPDFWriteObject{1543}{<< /Type /Annot /Subtype /Link /Border [0 0 0] /H /I /C [1 0 0] /Rect [254.286 108.897 289.479 120.85] /A << /S /GoTo /D <6571756174696F6E2E342E332E332E3137> >> >>}
\MusicPDFWriteObject{1544}{<< /Type /Annot /Subtype /Link /Border [0 0 0] /H /I /C [0 1 0] /Rect [397.58 673.973 416.684 682.84] /A << /S /GoTo /D <636974652E4E42> >> >>}
\MusicPDFWriteObject{1545}{<< /Type /Annot /Subtype /Link /Border [0 0 0] /H /I /C [1 0 0] /Rect [73.712 658.332 103.925 670.284] /A << /S /GoTo /D <6571756174696F6E2E342E332E322E35> >> >>}
\MusicPDFWriteObject{1546}{<< /Type /Annot /Subtype /Link /Border [0 0 0] /H /I /C [1 0 0] /Rect [405.462 645.181 435.731 657.134] /A << /S /GoTo /D <74686D2E342E332E312E35> >> >>}
\MusicPDFWriteObject{1547}{<< /Type /Annot /Subtype /Link /Border [0 0 0] /H /I /C [1 0 0] /Rect [124.489 618.242 154.702 630.223] /A << /S /GoTo /D <74686D2E342E332E332E31> >> >>}
\MusicPDFWriteObject{1548}{<< /Type /Annot /Subtype /Link /Border [0 0 0] /H /I /C [1 0 0] /Rect [401.985 618.242 432.198 630.223] /A << /S /GoTo /D <74686D2E342E332E332E33> >> >>}
\MusicPDFWriteObject{1549}{<< /Type /Annot /Subtype /Link /Border [0 0 0] /H /I /C [0 1 0] /Rect [526.005 608.219 543.159 617.087] /A << /S /GoTo /D <636974652E466F6C> >> >>}
\MusicPDFWriteObject{1550}{<< /Type /Annot /Subtype /Link /Border [0 0 0] /H /I /C [1 0 0] /Rect [413.921 591.941 444.473 605.17] /A << /S /GoTo /D <74686D2E342E332E332E33> >> >>}
\MusicPDFWriteObject{1551}{<< /Type /Annot /Subtype /Link /Border [0 0 0] /H /I /C [0 1 0] /Rect [132.781 581.917 166.75 590.785] /A << /S /GoTo /D <636974652E436E772D61> >> >>}
\MusicPDFWriteObject{1552}{<< /Type /Annot /Subtype /Link /Border [0 0 0] /H /I /C [0 1 0] /Rect [133.951 555.616 151.285 564.484] /A << /S /GoTo /D <636974652E4B6E> >> >>}
\MusicPDFWriteObject{1553}{<< /Type /Annot /Subtype /Link /Border [0 0 0] /H /I /C [1 0 0] /Rect [133.837 495.098 164.615 505.889] /A << /S /GoTo /D <74686D2E342E332E332E34> >> >>}
\MusicPDFWriteObject{1554}{<< /Type /Annot /Subtype /Link /Border [0 0 0] /H /I /C [1 0 0] /Rect [343.321 495.098 374.098 505.889] /A << /S /GoTo /D <74686D2E342E332E322E32> >> >>}
\MusicPDFWriteObject{1555}{<< /Type /Annot /Subtype /Link /Border [0 0 0] /H /I /C [0 1 0] /Rect [294.082 183.264 312.384 192.132] /A << /S /GoTo /D <636974652E4445> >> >>}
\MusicPDFWriteObject{1556}{<< /Type /Annot /Subtype /Link /Border [0 0 0] /H /I /C [1 0 0] /Rect [207.478 298.07 239.691 310.026] /A << /S /GoTo /D <74686D2E432E342E332E35> >> >>}
\MusicPDFWriteObject{1557}{<< /Type /Annot /Subtype /Link /Border [0 0 0] /H /I /C [0 1 0] /Rect [512.791 287.41 543.184 296.278] /A << /S /GoTo /D <636974652E426F752D49> >> >>}
\MusicPDFWriteObject{1558}{<< /Type /Annot /Subtype /Link /Border [0 0 0] /H /I /C [0 1 0] /Rect [307.17 274.26 337.563 283.127] /A << /S /GoTo /D <636974652E426F752D49> >> >>}
\MusicPDFWriteObject{1559}{<< /Type /Annot /Subtype /Link /Border [0 0 0] /H /I /C [1 0 0] /Rect [140.168 193.418 177.61 204.209] /A << /S /GoTo /D <74686D2E432E342E332E3132> >> >>}
\MusicPDFWriteObject{1560}{<< /Type /Annot /Subtype /Link /Border [0 0 0] /H /I /C [1 0 0] /Rect [463.095 166.563 492.744 178.516] /A << /S /GoTo /D <74686D2E342E332E332E31> >> >>}
\MusicPDFWriteObject{1561}{<< /Type /Annot /Subtype /Link /Border [0 0 0] /H /I /C [1 0 0] /Rect [71.004 153.413 108.41 165.365] /A << /S /GoTo /D <74686D2E432E342E332E3132> >> >>}
\MusicPDFWriteObject{1562}{<< /Type /Annot /Subtype /Link /Border [0 0 0] /H /I /C [0 1 0] /Rect [267.245 155.903 286.35 164.771] /A << /S /GoTo /D <636974652E4E42> >> >>}
\MusicPDFWriteObject{1563}{<< /Type /Annot /Subtype /Link /Border [0 0 0] /H /I /C [1 0 0] /Rect [419.212 684.633 454.405 696.586] /A << /S /GoTo /D <6571756174696F6E2E342E332E332E3132> >> >>}
\MusicPDFWriteObject{1564}{<< /Type /Annot /Subtype /Link /Border [0 0 0] /H /I /C [1 0 0] /Rect [368.245 631.829 403.438 643.782] /A << /S /GoTo /D <6571756174696F6E2E342E332E332E3133> >> >>}
\MusicPDFWriteObject{1565}{<< /Type /Annot /Subtype /Link /Border [0 0 0] /H /I /C [1 0 0] /Rect [210.153 573.322 239.83 585.915] /A << /S /GoTo /D <74686D2E342E332E332E31> >> >>}
\MusicPDFWriteObject{1566}{<< /Type /Annot /Subtype /Link /Border [0 0 0] /H /I /C [1 0 0] /Rect [403.398 447.168 438.591 459.12] /A << /S /GoTo /D <6571756174696F6E2E342E332E332E3134> >> >>}
\MusicPDFWriteObject{1567}{<< /Type /Annot /Subtype /Link /Border [0 0 0] /H /I /C [0 1 0] /Rect [107.335 660.822 126.44 669.69] /A << /S /GoTo /D <636974652E4E42> >> >>}
\MusicPDFWriteObject{1568}{<< /Type /Annot /Subtype /Link /Border [0 0 0] /H /I /C [1 0 0] /Rect [320.677 638.241 350.89 649.031] /A << /S /GoTo /D <74686D2E342E332E332E33> >> >>}
\MusicPDFWriteObject{1569}{<< /Type /Annot /Subtype /Link /Border [0 0 0] /H /I /C [1 0 0] /Rect [429.182 379.136 459.056 389.927] /A << /S /GoTo /D <74686D2E342E332E332E33> >> >>}
\MusicPDFWriteObject{1570}{<< /Type /Annot /Subtype /Link /Border [0 0 0] /H /I /C [1 0 0] /Rect [117.9 365.985 148.113 376.776] /A << /S /GoTo /D <74686D2E342E332E332E31> >> >>}
\MusicPDFWriteObject{1571}{<< /Type /Annot /Subtype /Link /Border [0 0 0] /H /I /C [1 0 0] /Rect [258.92 365.985 289.133 376.776] /A << /S /GoTo /D <74686D2E342E332E332E37> >> >>}
\MusicPDFWriteObject{1572}{<< /Type /Annot /Subtype /Link /Border [0 0 0] /H /I /C [1 0 0] /Rect [237.314 345.341 269.539 356.132] /A << /S /GoTo /D <74686D2E442E312E392E32> >> >>}
\MusicPDFWriteObject{1573}{<< /Type /Annot /Subtype /Link /Border [0 0 0] /H /I /C [1 0 0] /Rect [172.31 318.485 186.991 330.441] /A << /S /GoTo /D <73656374696F6E2E322E31> >> >>}
\MusicPDFWriteObject{1574}{<< /Type /Annot /Subtype /Link /Border [0 0 0] /H /I /C [1 0 0] /Rect [70.655 279.587 100.868 290.378] /A << /S /GoTo /D <74686D2E342E332E332E31> >> >>}
\MusicPDFWriteObject{1575}{<< /Type /Annot /Subtype /Link /Border [0 0 0] /H /I /C [1 0 0] /Rect [140.515 279.587 170.728 290.378] /A << /S /GoTo /D <74686D2E342E332E332E37> >> >>}
\MusicPDFWriteObject{1576}{<< /Type /Annot /Subtype /Link /Border [0 0 0] /H /I /C [1 0 0] /Rect [210.772 235.907 240.421 246.698] /A << /S /GoTo /D <74686D2E342E322E312E35> >> >>}
\MusicPDFWriteObject{1577}{<< /Type /Annot /Subtype /Link /Border [0 0 0] /H /I /C [1 0 0] /Rect [312.639 235.907 342.288 246.698] /A << /S /GoTo /D <6571756174696F6E2E342E322E322E37> >> >>}
\MusicPDFWriteObject{1578}{<< /Type /Annot /Subtype /Link /Border [0 0 0] /H /I /C [1 0 0] /Rect [266.09 685.187 298.931 695.977] /A << /S /GoTo /D <74686D2E442E332E322E35> >> >>}
\MusicPDFWriteObject{1579}{<< /Type /Annot /Subtype /Link /Border [0 0 0] /H /I /C [1 0 0] /Rect [345.107 685.187 377.949 695.977] /A << /S /GoTo /D <74686D2E442E332E332E37> >> >>}
\MusicPDFWriteObject{1580}{<< /Type /Annot /Subtype /Link /Border [0 0 0] /H /I /C [1 0 0] /Rect [440.725 685.187 478.547 695.977] /A << /S /GoTo /D <74686D2E442E332E332E3133> >> >>}
\MusicPDFWriteObject{1581}{<< /Type /Annot /Subtype /Link /Border [0 0 0] /H /I /C [1 0 0] /Rect [193.921 658.885 227.379 669.676] /A << /S /GoTo /D <74686D2E442E322E312E31> >> >>}
\MusicPDFWriteObject{1582}{<< /Type /Annot /Subtype /Link /Border [0 0 0] /H /I /C [1 0 0] /Rect [203.366 582.451 236.825 594.406] /A << /S /GoTo /D <74686D2E442E332E332E31> >> >>}
\MusicPDFWriteObject{1583}{<< /Type /Annot /Subtype /Link /Border [0 0 0] /H /I /C [1 0 0] /Rect [71.004 477.246 104.462 489.199] /A << /S /GoTo /D <74686D2E442E322E312E32> >> >>}
\MusicPDFWriteObject{1584}{<< /Type /Annot /Subtype /Link /Border [0 0 0] /H /I /C [1 0 0] /Rect [249.838 437.793 282.248 449.748] /A << /S /GoTo /D <74686D2E442E332E332E37> >> >>}
\MusicPDFWriteObject{1585}{<< /Type /Annot /Subtype /Link /Border [0 0 0] /H /I /C [1 0 0] /Rect [366.961 388.215 404.368 400.168] /A << /S /GoTo /D <6571756174696F6E2E432E322E332E3232> >> >>}
\MusicPDFWriteObject{1586}{<< /Type /Annot /Subtype /Link /Border [0 0 0] /H /I /C [1 0 0] /Rect [87.215 361.914 112.207 373.866] /A << /S /GoTo /D <73756273656374696F6E2E422E322E31> >> >>}
\MusicPDFWriteObject{1587}{<< /Type /Annot /Subtype /Link /Border [0 0 0] /H /I /C [1 0 0] /Rect [120.782 361.914 146.339 373.866] /A << /S /GoTo /D <73756273656374696F6E2E442E332E32> >> >>}
\MusicPDFWriteObject{1588}{<< /Type /Annot /Subtype /Link /Border [0 0 0] /H /I /C [1 0 0] /Rect [389.227 361.914 422.686 373.866] /A << /S /GoTo /D <74686D2E442E332E322E33> >> >>}
\MusicPDFWriteObject{1589}{<< /Type /Annot /Subtype /Link /Border [0 0 0] /H /I /C [1 0 0] /Rect [116.301 322.461 153.707 334.414] /A << /S /GoTo /D <6571756174696F6E2E432E322E332E3232> >> >>}
\MusicPDFWriteObject{1590}{<< /Type /Annot /Subtype /Link /Border [0 0 0] /H /I /C [1 0 0] /Rect [245.329 322.461 282.028 334.414] /A << /S /GoTo /D <74686D2E432E322E332E3133> >> >>}
\MusicPDFWriteObject{1591}{<< /Type /Annot /Subtype /Link /Border [0 0 0] /H /I /C [0 1 0] /Rect [85.944 311.801 117.568 320.669] /A << /S /GoTo /D <636974652E47726F2D62> >> >>}
\MusicPDFWriteObject{1592}{<< /Type /Annot /Subtype /Link /Border [0 0 0] /H /I /C [1 0 0] /Rect [480.567 211.991 513.192 223.943] /A << /S /GoTo /D <74686D2E442E322E312E31> >> >>}
\MusicPDFWriteObject{1593}{<< /Type /Annot /Subtype /Link /Border [0 0 0] /H /I /C [1 0 0] /Rect [127.594 199.393 159.88 211.432] /A << /S /GoTo /D <74686D2E442E322E312E33> >> >>}
\MusicPDFWriteObject{1594}{<< /Type /Annot /Subtype /Link /Border [0 0 0] /H /I /C [1 0 0] /Rect [146.65 685.187 184.472 695.977] /A << /S /GoTo /D <74686D2E442E332E332E3133> >> >>}
\MusicPDFWriteObject{1595}{<< /Type /Annot /Subtype /Link /Border [0 0 0] /H /I /C [1 0 0] /Rect [483.576 635.054 515.801 647.009] /A << /S /GoTo /D <74686D2E442E322E312E32> >> >>}
\MusicPDFWriteObject{1596}{<< /Type /Annot /Subtype /Link /Border [0 0 0] /H /I /C [1 0 0] /Rect [313.588 447.055 346.429 459.01] /A << /S /GoTo /D <74686D2E442E332E322E35> >> >>}
\MusicPDFWriteObject{1597}{<< /Type /Annot /Subtype /Link /Border [0 0 0] /H /I /C [1 0 0] /Rect [409.206 447.055 447.027 459.01] /A << /S /GoTo /D <74686D2E442E332E332E3133> >> >>}
\MusicPDFWriteObject{1598}{<< /Type /Annot /Subtype /Link /Border [0 0 0] /H /I /C [1 0 0] /Rect [123.769 225.269 160.981 238.222] /A << /S /GoTo /D <74686D2E442E332E332E3132> >> >>}
\MusicPDFWriteObject{1599}{<< /Type /Annot /Subtype /Link /Border [0 0 0] /H /I /C [1 0 0] /Rect [190.582 658.332 220.795 670.284] /A << /S /GoTo /D <6571756174696F6E2E342E342E312E33> >> >>}
\MusicPDFWriteObject{1600}{<< /Type /Annot /Subtype /Link /Border [0 0 0] /H /I /C [1 0 0] /Rect [475.194 575.665 505.407 587.618] /A << /S /GoTo /D <6571756174696F6E2E342E342E312E33> >> >>}
\MusicPDFWriteObject{1601}{<< /Type /Annot /Subtype /Link /Border [0 0 0] /H /I /C [1 0 0] /Rect [165.958 334.56 196.171 346.513] /A << /S /GoTo /D <6571756174696F6E2E342E342E312E33> >> >>}
\MusicPDFWriteObject{1602}{<< /Type /Annot /Subtype /Link /Border [0 0 0] /H /I /C [1 0 0] /Rect [196.633 295.107 229.475 307.062] /A << /S /GoTo /D <74686D2E442E332E332E37> >> >>}
\MusicPDFWriteObject{1603}{<< /Type /Annot /Subtype /Link /Border [0 0 0] /H /I /C [1 0 0] /Rect [396.528 198.704 426.402 211.587] /A << /S /GoTo /D <74686D2E342E342E312E31> >> >>}
\MusicPDFWriteObject{1604}{<< /Type /Annot /Subtype /Link /Border [0 0 0] /H /I /C [1 0 0] /Rect [259.626 172.403 289.275 185.286] /A << /S /GoTo /D <74686D2E342E322E312E35> >> >>}
\MusicPDFWriteObject{1605}{<< /Type /Annot /Subtype /Link /Border [0 0 0] /H /I /C [1 0 0] /Rect [361.767 172.403 391.416 185.286] /A << /S /GoTo /D <6571756174696F6E2E342E322E322E37> >> >>}
\MusicPDFWriteObject{1606}{<< /Type /Annot /Subtype /Link /Border [0 0 0] /H /I /C [1 0 0] /Rect [289.377 159.252 319.59 171.205] /A << /S /GoTo /D <6571756174696F6E2E342E322E322E37> >> >>}
\MusicPDFWriteObject{1607}{<< /Type /Annot /Subtype /Link /Border [0 0 0] /H /I /C [1 0 0] /Rect [387.082 527.866 419.924 542.06] /A << /S /GoTo /D <74686D2E442E312E322E31> >> >>}
\MusicPDFWriteObject{1608}{<< /Type /Annot /Subtype /Link /Border [0 0 0] /H /I /C [1 0 0] /Rect [164.716 237.424 194.929 249.377] /A << /S /GoTo /D <6571756174696F6E2E342E342E322E33> >> >>}
\MusicPDFWriteObject{1609}{<< /Type /Annot /Subtype /Link /Border [0 0 0] /H /I /C [1 0 0] /Rect [277.964 158.183 310.805 168.974] /A << /S /GoTo /D <74686D2E442E322E312E32> >> >>}
\MusicPDFWriteObject{1610}{<< /Type /Annot /Subtype /Link /Border [0 0 0] /H /I /C [1 0 0] /Rect [71.004 671.482 103.845 683.435] /A << /S /GoTo /D <74686D2E442E312E342E36> >> >>}
\MusicPDFWriteObject{1611}{<< /Type /Annot /Subtype /Link /Border [0 0 0] /H /I /C [1 0 0] /Rect [390.117 521.078 420.33 533.031] /A << /S /GoTo /D <6571756174696F6E2E342E342E322E34> >> >>}
\MusicPDFWriteObject{1612}{<< /Type /Annot /Subtype /Link /Border [0 0 0] /H /I /C [1 0 0] /Rect [335.486 408.102 365.699 420.054] /A << /S /GoTo /D <6571756174696F6E2E342E342E322E34> >> >>}
\MusicPDFWriteObject{1613}{<< /Type /Annot /Subtype /Link /Border [0 0 0] /H /I /C [1 0 0] /Rect [402.062 394.95 435.521 407.892] /A << /S /GoTo /D <74686D2E442E312E342E36> >> >>}
\MusicPDFWriteObject{1614}{<< /Type /Annot /Subtype /Link /Border [0 0 0] /H /I /C [1 0 0] /Rect [276.488 296.656 307.265 307.447] /A << /S /GoTo /D <74686D2E342E322E312E32> >> >>}
\MusicPDFWriteObject{1615}{<< /Type /Annot /Subtype /Link /Border [0 0 0] /H /I /C [1 0 0] /Rect [280.13 249.523 311.413 260.314] /A << /S /GoTo /D <74686D2E342E342E322E33> >> >>}
\MusicPDFWriteObject{1616}{<< /Type /Annot /Subtype /Link /Border [0 0 0] /H /I /C [1 0 0] /Rect [158.965 96.3 192.913 108.629] /A << /S /GoTo /D <74686D2E442E312E322E31> >> >>}
\MusicPDFWriteObject{1617}{<< /Type /Annot /Subtype /Link /Border [0 0 0] /H /I /C [1 0 0] /Rect [275.729 553.693 308.971 565.732] /A << /S /GoTo /D <74686D2E442E312E322E31> >> >>}
\MusicPDFWriteObject{1618}{<< /Type /Annot /Subtype /Link /Border [0 0 0] /H /I /C [1 0 0] /Rect [496.15 553.693 529.393 565.732] /A << /S /GoTo /D <74686D2E442E322E312E31> >> >>}
\MusicPDFWriteObject{1619}{<< /Type /Annot /Subtype /Link /Border [0 0 0] /H /I /C [0 1 0] /Rect [97.058 479.072 125.936 487.94] /A << /S /GoTo /D <636974652E5363682D61> >> >>}
\MusicPDFWriteObject{1620}{<< /Type /Annot /Subtype /Link /Border [0 0 0] /H /I /C [1 0 0] /Rect [70.655 437.13 100.868 449.082] /A << /S /GoTo /D <74686D2E342E342E322E36> >> >>}
\MusicPDFWriteObject{1621}{<< /Type /Annot /Subtype /Link /Border [0 0 0] /H /I /C [0 1 0] /Rect [287.29 439.62 306.263 448.487] /A << /S /GoTo /D <636974652E547265> >> >>}
\MusicPDFWriteObject{1622}{<< /Type /Annot /Subtype /Link /Border [0 0 0] /H /I /C [0 1 0] /Rect [467.924 358.09 487.028 366.957] /A << /S /GoTo /D <636974652E4E42> >> >>}
\MusicPDFWriteObject{1623}{<< /Type /Annot /Subtype /Link /Border [0 0 0] /H /I /C [1 0 0] /Rect [254.679 301.865 285.456 314.18] /A << /S /GoTo /D <74686D2E342E342E322E36> >> >>}
\MusicPDFWriteObject{1624}{<< /Type /Annot /Subtype /Link /Border [0 0 0] /H /I /C [1 0 0] /Rect [124.591 262.773 154.888 274.726] /A << /S /GoTo /D <74686D2E342E342E322E36> >> >>}
\MusicPDFWriteObject{1625}{<< /Type /Annot /Subtype /Link /Border [0 0 0] /H /I /C [0 1 0] /Rect [351.287 643.151 369.485 652.019] /A << /S /GoTo /D <636974652E4A6172> >> >>}
\MusicPDFWriteObject{1626}{<< /Type /Annot /Subtype /Link /Border [0 0 0] /H /I /C [1 0 0] /Rect [260.935 506.966 293.776 522.954] /A << /S /GoTo /D <74686D2E442E312E342E34> >> >>}
\MusicPDFWriteObject{1627}{<< /Type /Annot /Subtype /Link /Border [0 0 0] /H /I /C [0 1 0] /Rect [441.297 381.043 459.6 389.911] /A << /S /GoTo /D <636974652E4445> >> >>}
\MusicPDFWriteObject{1628}{<< /Type /Annot /Subtype /Link /Border [0 0 0] /H /I /C [0 1 0] /Rect [73.765 367.892 92.067 376.76] /A << /S /GoTo /D <636974652E4445> >> >>}
\MusicPDFWriteObject{1629}{<< /Type /Annot /Subtype /Link /Border [0 0 0] /H /I /C [1 0 0] /Rect [514.169 259.442 543.817 277.98] /A << /S /GoTo /D <74686D2E342E342E322E35> >> >>}
\MusicPDFWriteObject{1630}{<< /Type /Annot /Subtype /Link /Border [0 0 0] /H /I /C [1 0 0] /Rect [243.305 218.69 273.518 230.643] /A << /S /GoTo /D <6571756174696F6E2E342E342E322E38> >> >>}
\MusicPDFWriteObject{1631}{<< /Type /Annot /Subtype /Link /Border [0 0 0] /H /I /C [0 1 0] /Rect [499.367 208.029 517.566 216.897] /A << /S /GoTo /D <636974652E4A6172> >> >>}
\MusicPDFWriteObject{1632}{<< /Type /Annot /Subtype /Link /Border [0 0 0] /H /I /C [0 1 0] /Rect [363.316 194.879 381.515 203.747] /A << /S /GoTo /D <636974652E4A6172> >> >>}
\MusicPDFWriteObject{1633}{<< /Type /Annot /Subtype /Link /Border [0 0 0] /H /I /C [0 1 0] /Rect [181.95 181.728 200.149 190.596] /A << /S /GoTo /D <636974652E4A6172> >> >>}
\MusicPDFWriteObject{1634}{<< /Type /Annot /Subtype /Link /Border [0 0 0] /H /I /C [1 0 0] /Rect [238.436 632.03 273.63 643.983] /A << /S /GoTo /D <6571756174696F6E2E342E342E322E3131> >> >>}
\MusicPDFWriteObject{1635}{<< /Type /Annot /Subtype /Link /Border [0 0 0] /H /I /C [1 0 0] /Rect [484.805 618.879 519.998 630.832] /A << /S /GoTo /D <6571756174696F6E2E342E342E322E3131> >> >>}
\MusicPDFWriteObject{1636}{<< /Type /Annot /Subtype /Link /Border [0 0 0] /H /I /C [1 0 0] /Rect [163.912 577.753 199.106 589.706] /A << /S /GoTo /D <6571756174696F6E2E342E342E322E3131> >> >>}
\MusicPDFWriteObject{1637}{<< /Type /Annot /Subtype /Link /Border [0 0 0] /H /I /C [0 1 0] /Rect [295.116 459.933 313.419 468.801] /A << /S /GoTo /D <636974652E4445> >> >>}
\MusicPDFWriteObject{1638}{<< /Type /Annot /Subtype /Link /Border [0 0 0] /H /I /C [1 0 0] /Rect [73.712 370.903 108.905 382.856] /A << /S /GoTo /D <6571756174696F6E2E342E342E322E3132> >> >>}
\MusicPDFWriteObject{1639}{<< /Type /Annot /Subtype /Link /Border [0 0 0] /H /I /C [1 0 0] /Rect [139.978 309.059 175.171 321.011] /A << /S /GoTo /D <6571756174696F6E2E342E342E322E3131> >> >>}
\MusicPDFWriteObject{1640}{<< /Type /Annot /Subtype /Link /Border [0 0 0] /H /I /C [0 1 0] /Rect [88.878 285.247 107.181 294.115] /A << /S /GoTo /D <636974652E4445> >> >>}
\MusicPDFWriteObject{1641}{<< /Type /Annot /Subtype /Link /Border [0 0 0] /H /I /C [1 0 0] /Rect [205.314 256.456 240.507 268.408] /A << /S /GoTo /D <6571756174696F6E2E342E342E322E3131> >> >>}
\MusicPDFWriteObject{1642}{<< /Type /Annot /Subtype /Link /Border [0 0 0] /H /I /C [1 0 0] /Rect [388.753 185.857 418.966 197.81] /A << /S /GoTo /D <6571756174696F6E2E342E342E322E38> >> >>}
\MusicPDFWriteObject{1643}{<< /Type /Annot /Subtype /Link /Border [0 0 0] /H /I /C [0 1 0] /Rect [291.941 112.528 310.139 121.396] /A << /S /GoTo /D <636974652E4A6172> >> >>}
\MusicPDFWriteObject{1644}{<< /Type /Annot /Subtype /Link /Border [0 0 0] /H /I /C [0 1 0] /Rect [161.026 99.377 179.225 108.245] /A << /S /GoTo /D <636974652E4A6172> >> >>}
\MusicPDFWriteObject{1645}{<< /Type /Annot /Subtype /Link /Border [0 0 0] /H /I /C [0 1 0] /Rect [498.96 671.984 517.159 680.852] /A << /S /GoTo /D <636974652E4A6172> >> >>}
\MusicPDFWriteObject{1646}{<< /Type /Annot /Subtype /Link /Border [0 0 0] /H /I /C [1 0 0] /Rect [235.628 643.193 265.841 655.146] /A << /S /GoTo /D <6571756174696F6E2E342E342E322E38> >> >>}
\MusicPDFWriteObject{1647}{<< /Type /Annot /Subtype /Link /Border [0 0 0] /H /I /C [0 1 0] /Rect [524.96 645.683 543.159 654.551] /A << /S /GoTo /D <636974652E4A6172> >> >>}
\MusicPDFWriteObject{1648}{<< /Type /Annot /Subtype /Link /Border [0 0 0] /H /I /C [1 0 0] /Rect [385.474 630.042 416.252 642.925] /A << /S /GoTo /D <74686D2E342E312E322E36> >> >>}
\MusicPDFWriteObject{1649}{<< /Type /Annot /Subtype /Link /Border [0 0 0] /H /I /C [1 0 0] /Rect [478.035 616.891 511.493 629.774] /A << /S /GoTo /D <74686D2E442E312E352E36> >> >>}
\MusicPDFWriteObject{1650}{<< /Type /Annot /Subtype /Link /Border [0 0 0] /H /I /C [1 0 0] /Rect [139.2 590.589 168.849 603.473] /A << /S /GoTo /D <74686D2E342E312E322E34> >> >>}
\MusicPDFWriteObject{1651}{<< /Type /Annot /Subtype /Link /Border [0 0 0] /H /I /C [1 0 0] /Rect [71.004 574.187 103.907 589.394] /A << /S /GoTo /D <74686D2E442E312E352E37> >> >>}
\MusicPDFWriteObject{1652}{<< /Type /Annot /Subtype /Link /Border [0 0 0] /H /I /C [1 0 0] /Rect [498.109 558.458 530.642 575.183] /A << /S /GoTo /D <74686D2E442E312E352E37> >> >>}
\MusicPDFWriteObject{1653}{<< /Type /Annot /Subtype /Link /Border [0 0 0] /H /I /C [0 1 0] /Rect [305.389 533.756 323.588 542.624] /A << /S /GoTo /D <636974652E4A6172> >> >>}
\MusicPDFWriteObject{1654}{<< /Type /Annot /Subtype /Link /Border [0 0 0] /H /I /C [1 0 0] /Rect [112.727 517.057 142.94 529.01] /A << /S /GoTo /D <74686D2E342E312E322E36> >> >>}
\MusicPDFWriteObject{1655}{<< /Type /Annot /Subtype /Link /Border [0 0 0] /H /I /C [1 0 0] /Rect [392.103 459.321 422.316 471.636] /A << /S /GoTo /D <74686D2E342E342E322E31> >> >>}
\MusicPDFWriteObject{1656}{<< /Type /Annot /Subtype /Link /Border [0 0 0] /H /I /C [1 0 0] /Rect [118.884 223.39 149.662 234.181] /A << /S /GoTo /D <74686D2E342E342E322E31> >> >>}
\MusicPDFWriteObject{1657}{<< /Type /Annot /Subtype /Link /Border [0 0 0] /H /I /C [1 0 0] /Rect [281.634 176.575 313.153 187.366] /A << /S /GoTo /D <74686D2E342E342E322E33> >> >>}
\MusicPDFWriteObject{1658}{<< /Type /Annot /Subtype /Link /Border [0 0 0] /H /I /C [1 0 0] /Rect [226.229 632.03 256.443 643.983] /A << /S /GoTo /D <6571756174696F6E2E342E342E322E31> >> >>}
\MusicPDFWriteObject{1659}{<< /Type /Annot /Subtype /Link /Border [0 0 0] /H /I /C [1 0 0] /Rect [414.926 469.725 447.244 482.318] /A << /S /GoTo /D <74686D2E442E332E332E37> >> >>}
\MusicPDFWriteObject{1660}{<< /Type /Annot /Subtype /Link /Border [0 0 0] /H /I /C [0 1 0] /Rect [200.449 459.065 229.21 467.933] /A << /S /GoTo /D <636974652E47726F2D66> >> >>}
\MusicPDFWriteObject{1661}{<< /Type /Annot /Subtype /Link /Border [0 0 0] /H /I /C [0 1 0] /Rect [363.665 459.065 380.985 467.933] /A << /S /GoTo /D <636974652E534E> >> >>}
\MusicPDFWriteObject{1662}{<< /Type /Annot /Subtype /Link /Border [0 0 0] /H /I /C [1 0 0] /Rect [273.749 368.098 311.156 378.888] /A << /S /GoTo /D <74686D2E432E322E332E3132> >> >>}
\MusicPDFWriteObject{1663}{<< /Type /Annot /Subtype /Link /Border [0 0 0] /H /I /C [1 0 0] /Rect [164.136 356.884 194.349 365.738] /A << /S /GoTo /D <74686D2E342E352E312E33> >> >>}
\MusicPDFWriteObject{1664}{<< /Type /Annot /Subtype /Link /Border [0 0 0] /H /I /C [1 0 0] /Rect [107.691 307.421 137.339 319.374] /A << /S /GoTo /D <74686D2E342E352E312E32> >> >>}
\MusicPDFWriteObject{1665}{<< /Type /Annot /Subtype /Link /Border [0 0 0] /H /I /C [0 1 0] /Rect [316.039 309.911 347.663 318.779] /A << /S /GoTo /D <636974652E47726F2D62> >> >>}
\MusicPDFWriteObject{1666}{<< /Type /Annot /Subtype /Link /Border [0 0 0] /H /I /C [1 0 0] /Rect [387.853 240.404 421.034 252.36] /A << /S /GoTo /D <74686D2E442E332E332E37> >> >>}
\MusicPDFWriteObject{1667}{<< /Type /Annot /Subtype /Link /Border [0 0 0] /H /I /C [1 0 0] /Rect [512.381 148.349 543.159 160.304] /A << /S /GoTo /D <74686D2E342E352E312E32> >> >>}
\MusicPDFWriteObject{1668}{<< /Type /Annot /Subtype /Link /Border [0 0 0] /H /I /C [1 0 0] /Rect [425.594 95.747 463.071 107.699] /A << /S /GoTo /D <74686D2E432E322E332E3131> >> >>}
\MusicPDFWriteObject{1669}{<< /Type /Annot /Subtype /Link /Border [0 0 0] /H /I /C [1 0 0] /Rect [206.563 684.633 238.787 696.586] /A << /S /GoTo /D <74686D2E442E332E332E37> >> >>}
\MusicPDFWriteObject{1670}{<< /Type /Annot /Subtype /Link /Border [0 0 0] /H /I /C [1 0 0] /Rect [278.849 516.338 308.498 528.653] /A << /S /GoTo /D <74686D2E342E322E312E35> >> >>}
\MusicPDFWriteObject{1671}{<< /Type /Annot /Subtype /Link /Border [0 0 0] /H /I /C [1 0 0] /Rect [220.877 451.412 251.09 462.288] /A << /S /GoTo /D <74686D2E342E322E312E35> >> >>}
\MusicPDFWriteObject{1672}{<< /Type /Annot /Subtype /Link /Border [0 0 0] /H /I /C [1 0 0] /Rect [338.835 424.643 369.613 436.596] /A << /S /GoTo /D <74686D2E322E322E342E33> >> >>}
\MusicPDFWriteObject{1673}{<< /Type /Annot /Subtype /Link /Border [0 0 0] /H /I /C [1 0 0] /Rect [106.556 321.927 136.769 330.781] /A << /S /GoTo /D <74686D2E342E322E312E35> >> >>}
\MusicPDFWriteObject{1674}{<< /Type /Annot /Subtype /Link /Border [0 0 0] /H /I /C [0 1 0] /Rect [85.385 216.721 120.44 225.589] /A << /S /GoTo /D <636974652E426F752D4949> >> >>}
\MusicPDFWriteObject{1675}{<< /Type /Annot /Subtype /Link /Border [0 0 0] /H /I /C [0 1 0] /Rect [375.07 216.721 410.125 225.589] /A << /S /GoTo /D <636974652E426F752D4949> >> >>}
\MusicPDFWriteObject{1676}{<< /Type /Annot /Subtype /Link /Border [0 0 0] /H /I /C [1 0 0] /Rect [260.786 174.778 266.749 187.636] /A << /S /GoTo /D <48666F6F746E6F74652E3132> >> >>}
\MusicPDFWriteObject{1677}{<< /Type /Annot /Subtype /Link /Border [0 0 0] /H /I /C [1 0 0] /Rect [199.445 135.326 234.638 147.279] /A << /S /GoTo /D <6571756174696F6E2E342E352E312E3132> >> >>}
\MusicPDFWriteObject{1678}{<< /Type /Annot /Subtype /Link /Border [0 0 0] /H /I /C [1 0 0] /Rect [309.078 135.325 342.536 147.613] /A << /S /GoTo /D <74686D2E442E312E352E33> >> >>}
\MusicPDFWriteObject{1679}{<< /Type /Annot /Subtype /Link /Border [0 0 0] /H /I /C [1 0 0] /Rect [317.927 632.03 350.769 643.985] /A << /S /GoTo /D <74686D2E442E332E332E37> >> >>}
\MusicPDFWriteObject{1680}{<< /Type /Annot /Subtype /Link /Border [0 0 0] /H /I /C [1 0 0] /Rect [241.09 328.799 247.052 342.017] /A << /S /GoTo /D <48666F6F746E6F74652E3133> >> >>}
\MusicPDFWriteObject{1681}{<< /Type /Annot /Subtype /Link /Border [0 0 0] /H /I /C [1 0 0] /Rect [194.413 316.008 227.872 327.963] /A << /S /GoTo /D <74686D2E442E332E332E37> >> >>}
\MusicPDFWriteObject{1682}{<< /Type /Annot /Subtype /Link /Border [0 0 0] /H /I /C [1 0 0] /Rect [70.655 210.802 100.868 222.757] /A << /S /GoTo /D <74686D2E342E322E312E35> >> >>}
\MusicPDFWriteObject{1683}{<< /Type /Annot /Subtype /Link /Border [0 0 0] /H /I /C [1 0 0] /Rect [229.513 474.165 260.291 488.341] /A << /S /GoTo /D <74686D2E342E332E332E31> >> >>}
\MusicPDFWriteObject{1684}{<< /Type /Annot /Subtype /Link /Border [0 0 0] /H /I /C [1 0 0] /Rect [373.543 462.597 406.384 474.552] /A << /S /GoTo /D <74686D2E442E322E322E31> >> >>}
\MusicPDFWriteObject{1685}{<< /Type /Annot /Subtype /Link /Border [0 0 0] /H /I /C [1 0 0] /Rect [234.586 449.446 264.234 462.33] /A << /S /GoTo /D <74686D2E342E332E332E31> >> >>}
\MusicPDFWriteObject{1686}{<< /Type /Annot /Subtype /Link /Border [0 0 0] /H /I /C [1 0 0] /Rect [317.34 449.446 349.564 462.33] /A << /S /GoTo /D <74686D2E442E322E322E31> >> >>}
\MusicPDFWriteObject{1687}{<< /Type /Annot /Subtype /Link /Border [0 0 0] /H /I /C [1 0 0] /Rect [419.389 409.995 449.602 421.947] /A << /S /GoTo /D <6571756174696F6E2E342E352E312E37> >> >>}
\MusicPDFWriteObject{1688}{<< /Type /Annot /Subtype /Link /Border [0 0 0] /H /I /C [1 0 0] /Rect [450.734 357.392 480.947 369.344] /A << /S /GoTo /D <6571756174696F6E2E342E352E312E38> >> >>}
\MusicPDFWriteObject{1689}{<< /Type /Annot /Subtype /Link /Border [0 0 0] /H /I /C [1 0 0] /Rect [73.712 344.241 103.925 356.194] /A << /S /GoTo /D <6571756174696F6E2E342E352E312E31> >> >>}
\MusicPDFWriteObject{1690}{<< /Type /Annot /Subtype /Link /Border [0 0 0] /H /I /C [1 0 0] /Rect [210.343 331.09 241.12 343.682] /A << /S /GoTo /D <74686D2E342E332E332E31> >> >>}
\MusicPDFWriteObject{1691}{<< /Type /Annot /Subtype /Link /Border [0 0 0] /H /I /C [1 0 0] /Rect [352.947 317.939 386.313 329.892] /A << /S /GoTo /D <74686D2E442E332E312E32> >> >>}
\MusicPDFWriteObject{1692}{<< /Type /Annot /Subtype /Link /Border [0 0 0] /H /I /C [1 0 0] /Rect [449.977 317.939 483.343 329.892] /A << /S /GoTo /D <74686D2E442E332E332E37> >> >>}
\MusicPDFWriteObject{1693}{<< /Type /Annot /Subtype /Link /Border [0 0 0] /H /I /C [1 0 0] /Rect [492.2 304.788 525.165 317.381] /A << /S /GoTo /D <74686D2E442E312E392E32> >> >>}
\MusicPDFWriteObject{1694}{<< /Type /Annot /Subtype /Link /Border [0 0 0] /H /I /C [1 0 0] /Rect [124.623 291.637 162.444 303.592] /A << /S /GoTo /D <74686D2E442E332E332E3132> >> >>}
\MusicPDFWriteObject{1695}{<< /Type /Annot /Subtype /Link /Border [0 0 0] /H /I /C [1 0 0] /Rect [150.012 239.034 183.471 251.627] /A << /S /GoTo /D <74686D2E442E332E332E37> >> >>}
\MusicPDFWriteObject{1696}{<< /Type /Annot /Subtype /Link /Border [0 0 0] /H /I /C [1 0 0] /Rect [269.05 199.583 299.263 211.535] /A << /S /GoTo /D <6571756174696F6E2E342E352E312E37> >> >>}
\MusicPDFWriteObject{1697}{<< /Type /Annot /Subtype /Link /Border [0 0 0] /H /I /C [1 0 0] /Rect [512.381 197.999 543.159 211.535] /A << /S /GoTo /D <74686D2E342E332E332E31> >> >>}
\MusicPDFWriteObject{1698}{<< /Type /Annot /Subtype /Link /Border [0 0 0] /H /I /C [0 1 0] /Rect [514.398 162.62 543.159 171.488] /A << /S /GoTo /D <636974652E47726F2D66> >> >>}
\MusicPDFWriteObject{1699}{<< /Type /Annot /Subtype /Link /Border [0 0 0] /H /I /C [0 1 0] /Rect [204.914 149.47 222.234 158.337] /A << /S /GoTo /D <636974652E534E> >> >>}
\MusicPDFWriteObject{1700}{<< /Type /Annot /Subtype /Link /Border [0 0 0] /H /I /C [1 0 0] /Rect [393.057 683.05 423.835 697.225] /A << /S /GoTo /D <74686D2E342E342E312E31> >> >>}
\MusicPDFWriteObject{1701}{<< /Type /Annot /Subtype /Link /Border [0 0 0] /H /I /C [1 0 0] /Rect [511.707 671.093 541.892 684.046] /A << /S /GoTo /D <74686D2E342E332E332E31> >> >>}
\MusicPDFWriteObject{1702}{<< /Type /Annot /Subtype /Link /Border [0 0 0] /H /I /C [1 0 0] /Rect [71.004 645.152 103.898 657.107] /A << /S /GoTo /D <74686D2E422E332E322E31> >> >>}
\MusicPDFWriteObject{1703}{<< /Type /Annot /Subtype /Link /Border [0 0 0] /H /I /C [1 0 0] /Rect [106.556 632.001 139.398 643.954] /A << /S /GoTo /D <74686D2E442E332E312E32> >> >>}
\MusicPDFWriteObject{1704}{<< /Type /Annot /Subtype /Link /Border [0 0 0] /H /I /C [1 0 0] /Rect [351.104 576.368 380.866 588.32] /A << /S /GoTo /D <74686D2E372E312E342E32> >> >>}
\MusicPDFWriteObject{1705}{<< /Type /Annot /Subtype /Link /Border [0 0 0] /H /I /C [0 1 0] /Rect [184.882 539.405 203.081 548.273] /A << /S /GoTo /D <636974652E4A6172> >> >>}
\MusicPDFWriteObject{1706}{<< /Type /Annot /Subtype /Link /Border [0 0 0] /H /I /C [1 0 0] /Rect [295.673 510.614 325.886 522.567] /A << /S /GoTo /D <6571756174696F6E2E342E342E322E36> >> >>}
\MusicPDFWriteObject{1707}{<< /Type /Annot /Subtype /Link /Border [0 0 0] /H /I /C [1 0 0] /Rect [377.434 510.253 407.252 522.568] /A << /S /GoTo /D <74686D2E342E342E322E36> >> >>}
\MusicPDFWriteObject{1708}{<< /Type /Annot /Subtype /Link /Border [0 0 0] /H /I /C [1 0 0] /Rect [444.772 495.776 474.929 510.346] /A << /S /GoTo /D <74686D2E342E342E322E36> >> >>}
\MusicPDFWriteObject{1709}{<< /Type /Annot /Subtype /Link /Border [0 0 0] /H /I /C [1 0 0] /Rect [275.767 442.329 310.96 454.281] /A << /S /GoTo /D <6571756174696F6E2E342E352E312E3130> >> >>}
\MusicPDFWriteObject{1710}{<< /Type /Annot /Subtype /Link /Border [0 0 0] /H /I /C [1 0 0] /Rect [308.605 388.942 347.143 401.257] /A << /S /GoTo /D <74686D2E442E332E332E3132> >> >>}
\MusicPDFWriteObject{1711}{<< /Type /Annot /Subtype /Link /Border [0 0 0] /H /I /C [0 1 0] /Rect [518.099 391.792 543.159 400.66] /A << /S /GoTo /D <636974652E53636861> >> >>}
\MusicPDFWriteObject{1712}{<< /Type /Annot /Subtype /Link /Border [0 0 0] /H /I /C [0 1 0] /Rect [391.293 352.34 410.398 361.208] /A << /S /GoTo /D <636974652E4E42> >> >>}
\MusicPDFWriteObject{1713}{<< /Type /Annot /Subtype /Link /Border [0 0 0] /H /I /C [1 0 0] /Rect [119.59 230.1 149.267 243.401] /A << /S /GoTo /D <74686D2E342E342E322E36> >> >>}
\MusicPDFWriteObject{1714}{<< /Type /Annot /Subtype /Link /Border [0 0 0] /H /I /C [1 0 0] /Rect [409.31 165.674 439.748 178.267] /A << /S /GoTo /D <74686D2E342E332E332E31> >> >>}
\MusicPDFWriteObject{1715}{<< /Type /Annot /Subtype /Link /Border [0 0 0] /H /I /C [0 1 0] /Rect [371.818 141.863 390.016 150.731] /A << /S /GoTo /D <636974652E4A6172> >> >>}
\MusicPDFWriteObject{1716}{<< /Type /Annot /Subtype /Link /Border [0 0 0] /H /I /C [1 0 0] /Rect [241.661 615.876 274.502 627.829] /A << /S /GoTo /D <74686D2E442E332E332E37> >> >>}
\MusicPDFWriteObject{1717}{<< /Type /Annot /Subtype /Link /Border [0 0 0] /H /I /C [0 1 0] /Rect [107.236 605.215 126.341 614.083] /A << /S /GoTo /D <636974652E4E42> >> >>}
\MusicPDFWriteObject{1718}{<< /Type /Annot /Subtype /Link /Border [0 0 0] /H /I /C [1 0 0] /Rect [255.541 575.999 288.999 588.313] /A << /S /GoTo /D <74686D2E442E332E332E37> >> >>}
\MusicPDFWriteObject{1719}{<< /Type /Annot /Subtype /Link /Border [0 0 0] /H /I /C [1 0 0] /Rect [372.378 523.756 402.591 535.709] /A << /S /GoTo /D <6571756174696F6E2E342E352E312E31> >> >>}
\MusicPDFWriteObject{1720}{<< /Type /Annot /Subtype /Link /Border [0 0 0] /H /I /C [1 0 0] /Rect [166.774 497.454 197.552 510.047] /A << /S /GoTo /D <74686D2E342E332E332E31> >> >>}
\MusicPDFWriteObject{1721}{<< /Type /Annot /Subtype /Link /Border [0 0 0] /H /I /C [1 0 0] /Rect [70.655 344.896 100.868 353.28] /A << /S /GoTo /D <74686D2E342E342E322E33> >> >>}
\MusicPDFWriteObject{1722}{<< /Type /Annot /Subtype /Link /Border [0 0 0] /H /I /C [1 0 0] /Rect [298.514 263.501 335.92 276.702] /A << /S /GoTo /D <74686D2E432E322E332E3131> >> >>}
\MusicPDFWriteObject{1723}{<< /Type /Annot /Subtype /Link /Border [0 0 0] /H /I /C [0 1 0] /Rect [227.682 252.841 246.787 261.709] /A << /S /GoTo /D <636974652E4E42> >> >>}
\MusicPDFWriteObject{1724}{<< /Type /Annot /Subtype /Link /Border [0 0 0] /H /I /C [1 0 0] /Rect [246.762 214.213 278.281 225.004] /A << /S /GoTo /D <74686D2E342E342E322E33> >> >>}
\MusicPDFWriteObject{1725}{<< /Type /Annot /Subtype /Link /Border [0 0 0] /H /I /C [1 0 0] /Rect [426.868 200.725 458.977 215.21] /A << /S /GoTo /D <74686D2E342E342E322E33> >> >>}
\MusicPDFWriteObject{1726}{<< /Type /Annot /Subtype /Link /Border [0 0 0] /H /I /C [1 0 0] /Rect [290.009 187.021 320.222 198.974] /A << /S /GoTo /D <6571756174696F6E2E342E342E322E36> >> >>}
\MusicPDFWriteObject{1727}{<< /Type /Annot /Subtype /Link /Border [0 0 0] /H /I /C [1 0 0] /Rect [199.639 136.263 230.417 149.506] /A << /S /GoTo /D <74686D2E342E352E312E31> >> >>}
\MusicPDFWriteObject{1728}{<< /Type /Annot /Subtype /Link /Border [0 0 0] /H /I /C [1 0 0] /Rect [292.535 82.596 322.748 94.549] /A << /S /GoTo /D <74686D2E342E332E332E31> >> >>}
\MusicPDFWriteObject{1729}{<< /Type /Annot /Subtype /Link /Border [0 0 0] /H /I /C [0 1 0] /Rect [105.927 687.123 151.198 695.991] /A << /S /GoTo /D <636974652E4B6F7468652D4949> >> >>}
\MusicPDFWriteObject{1730}{<< /Type /Annot /Subtype /Link /Border [0 0 0] /H /I /C [0 1 0] /Rect [361.608 660.822 380.581 669.69] /A << /S /GoTo /D <636974652E547265> >> >>}
\MusicPDFWriteObject{1731}{<< /Type /Annot /Subtype /Link /Border [0 0 0] /H /I /C [1 0 0] /Rect [325.825 619.433 356.602 630.223] /A << /S /GoTo /D <74686D2E342E322E312E37> >> >>}
\MusicPDFWriteObject{1732}{<< /Type /Annot /Subtype /Link /Border [0 0 0] /H /I /C [1 0 0] /Rect [71.004 608.219 108.825 617.073] /A << /S /GoTo /D <74686D2E442E332E332E3133> >> >>}
\MusicPDFWriteObject{1733}{<< /Type /Annot /Subtype /Link /Border [0 0 0] /H /I /C [0 1 0] /Rect [109.065 595.068 134.125 603.936] /A << /S /GoTo /D <636974652E53636861> >> >>}
\MusicPDFWriteObject{1734}{<< /Type /Annot /Subtype /Link /Border [0 0 0] /H /I /C [0 1 0] /Rect [322.093 581.917 347.153 590.785] /A << /S /GoTo /D <636974652E53636861> >> >>}
\MusicPDFWriteObject{1735}{<< /Type /Annot /Subtype /Link /Border [0 0 0] /H /I /C [1 0 0] /Rect [285.068 565.916 315.846 579.159] /A << /S /GoTo /D <74686D2E342E352E312E31> >> >>}
\MusicPDFWriteObject{1736}{<< /Type /Annot /Subtype /Link /Border [0 0 0] /H /I /C [0 1 0] /Rect [521.938 568.766 543.159 577.634] /A << /S /GoTo /D <636974652E427275> >> >>}
\MusicPDFWriteObject{1737}{<< /Type /Annot /Subtype /Link /Border [0 0 0] /H /I /C [1 0 0] /Rect [513.797 490.353 543.784 501.229] /A << /S /GoTo /D <74686D2E342E342E322E33> >> >>}
\MusicPDFWriteObject{1738}{<< /Type /Annot /Subtype /Link /Border [0 0 0] /H /I /C [1 0 0] /Rect [106.518 435.857 136.703 448.74] /A << /S /GoTo /D <74686D2E342E332E332E31> >> >>}
\MusicPDFWriteObject{1739}{<< /Type /Annot /Subtype /Link /Border [0 0 0] /H /I /C [1 0 0] /Rect [281.634 342.721 313.153 353.511] /A << /S /GoTo /D <74686D2E342E342E322E33> >> >>}
\MusicPDFWriteObject{1740}{<< /Type /Annot /Subtype /Link /Border [0 0 0] /H /I /C [1 0 0] /Rect [508.815 243.158 539.028 255.111] /A << /S /GoTo /D <6571756174696F6E2E342E352E312E31> >> >>}
\MusicPDFWriteObject{1741}{<< /Type /Annot /Subtype /Link /Border [0 0 0] /H /I /C [1 0 0] /Rect [197.44 175.409 227.653 187.722] /A << /S /GoTo /D <74686D2E342E352E312E31> >> >>}
\MusicPDFWriteObject{1742}{<< /Type /Annot /Subtype /Link /Border [0 0 0] /H /I /C [1 0 0] /Rect [373.727 112.68 403.941 125.923] /A << /S /GoTo /D <74686D2E342E342E322E36> >> >>}
\MusicPDFWriteObject{1743}{<< /Type /Annot /Subtype /Link /Border [0 0 0] /H /I /C [1 0 0] /Rect [451.371 99.889 481.584 111.842] /A << /S /GoTo /D <6571756174696F6E2E342E342E322E36> >> >>}
\MusicPDFWriteObject{1744}{<< /Type /Annot /Subtype /Link /Border [0 0 0] /H /I /C [1 0 0] /Rect [228.513 670.318 258.726 679.172] /A << /S /GoTo /D <74686D2E342E342E322E33> >> >>}
\MusicPDFWriteObject{1745}{<< /Type /Annot /Subtype /Link /Border [0 0 0] /H /I /C [1 0 0] /Rect [447.713 626.018 478.49 636.894] /A << /S /GoTo /D <74686D2E342E342E322E33> >> >>}
\MusicPDFWriteObject{1746}{<< /Type /Annot /Subtype /Link /Border [0 0 0] /H /I /C [1 0 0] /Rect [70.655 614.889 100.868 623.273] /A << /S /GoTo /D <74686D2E342E322E322E32> >> >>}
\MusicPDFWriteObject{1747}{<< /Type /Annot /Subtype /Link /Border [0 0 0] /H /I /C [1 0 0] /Rect [458.883 599.716 489.096 610.592] /A << /S /GoTo /D <74686D2E342E342E322E33> >> >>}
\MusicPDFWriteObject{1748}{<< /Type /Annot /Subtype /Link /Border [0 0 0] /H /I /C [0 1 0] /Rect [106.702 549.135 135.581 558.003] /A << /S /GoTo /D <636974652E5363682D61> >> >>}
\MusicPDFWriteObject{1749}{<< /Type /Annot /Subtype /Link /Border [0 0 0] /H /I /C [0 1 0] /Rect [217.732 549.135 242.792 558.003] /A << /S /GoTo /D <636974652E53636861> >> >>}
\MusicPDFWriteObject{1750}{<< /Type /Annot /Subtype /Link /Border [0 0 0] /H /I /C [0 1 0] /Rect [522.048 549.135 543.27 558.003] /A << /S /GoTo /D <636974652E427275> >> >>}
\MusicPDFWriteObject{1751}{<< /Type /Annot /Subtype /Link /Border [0 0 0] /H /I /C [1 0 0] /Rect [71.004 441.439 104.462 454.38] /A << /S /GoTo /D <74686D2E442E332E332E37> >> >>}
\MusicPDFWriteObject{1752}{<< /Type /Annot /Subtype /Link /Border [0 0 0] /H /I /C [0 1 0] /Rect [174.33 417.628 193.304 426.495] /A << /S /GoTo /D <636974652E547265> >> >>}
\MusicPDFWriteObject{1753}{<< /Type /Annot /Subtype /Link /Border [0 0 0] /H /I /C [0 1 0] /Rect [361.661 417.628 379.859 426.495] /A << /S /GoTo /D <636974652E4A6172> >> >>}
\MusicPDFWriteObject{1754}{<< /Type /Annot /Subtype /Link /Border [0 0 0] /H /I /C [1 0 0] /Rect [139.554 388.836 169.202 400.789] /A << /S /GoTo /D <74686D2E342E322E322E32> >> >>}
\MusicPDFWriteObject{1755}{<< /Type /Annot /Subtype /Link /Border [0 0 0] /H /I /C [1 0 0] /Rect [290.543 312.661 322.062 323.452] /A << /S /GoTo /D <74686D2E342E342E322E33> >> >>}
\MusicPDFWriteObject{1756}{<< /Type /Annot /Subtype /Link /Border [0 0 0] /H /I /C [1 0 0] /Rect [497.414 579.427 530.41 592.019] /A << /S /GoTo /D <74686D2E442E332E332E37> >> >>}
\MusicPDFWriteObject{1757}{<< /Type /Annot /Subtype /Link /Border [0 0 0] /H /I /C [1 0 0] /Rect [71.004 539.975 103.845 551.928] /A << /S /GoTo /D <74686D2E442E332E322E33> >> >>}
\MusicPDFWriteObject{1758}{<< /Type /Annot /Subtype /Link /Border [0 0 0] /H /I /C [1 0 0] /Rect [207.486 539.975 244.893 551.928] /A << /S /GoTo /D <74686D2E432E322E332E3133> >> >>}
\MusicPDFWriteObject{1759}{<< /Type /Annot /Subtype /Link /Border [0 0 0] /H /I /C [0 1 0] /Rect [150.158 516.163 168.357 525.031] /A << /S /GoTo /D <636974652E4A6172> >> >>}
\MusicPDFWriteObject{1760}{<< /Type /Annot /Subtype /Link /Border [0 0 0] /H /I /C [0 1 0] /Rect [186.736 361.379 205.841 370.247] /A << /S /GoTo /D <636974652E4E42> >> >>}
\MusicPDFWriteObject{1761}{<< /Type /Annot /Subtype /Link /Border [0 0 0] /H /I /C [1 0 0] /Rect [348.213 200.028 380.592 213.672] /A << /S /GoTo /D <74686D2E442E312E322E31> >> >>}
\MusicPDFWriteObject{1762}{<< /Type /Annot /Subtype /Link /Border [0 0 0] /H /I /C [1 0 0] /Rect [414.742 135.326 444.955 147.279] /A << /S /GoTo /D <6571756174696F6E2E342E352E322E31> >> >>}
\MusicPDFWriteObject{1763}{<< /Type /Annot /Subtype /Link /Border [0 0 0] /H /I /C [1 0 0] /Rect [256.68 122.175 286.328 134.128] /A << /S /GoTo /D <74686D2E342E342E312E31> >> >>}
\MusicPDFWriteObject{1764}{<< /Type /Annot /Subtype /Link /Border [0 0 0] /H /I /C [1 0 0] /Rect [158.504 95.873 188.858 107.828] /A << /S /GoTo /D <74686D2E342E352E312E31> >> >>}
\MusicPDFWriteObject{1765}{<< /Type /Annot /Subtype /Link /Border [0 0 0] /H /I /C [1 0 0] /Rect [248.762 671.482 279.539 683.435] /A << /S /GoTo /D <74686D2E342E322E322E32> >> >>}
\MusicPDFWriteObject{1766}{<< /Type /Annot /Subtype /Link /Border [0 0 0] /H /I /C [0 1 0] /Rect [516.667 673.973 543.159 682.84] /A << /S /GoTo /D <636974652E52752D32> >> >>}
\MusicPDFWriteObject{1767}{<< /Type /Annot /Subtype /Link /Border [0 0 0] /H /I /C [1 0 0] /Rect [345.688 622.513 376.466 634.468] /A << /S /GoTo /D <74686D2E342E352E322E31> >> >>}
\MusicPDFWriteObject{1768}{<< /Type /Annot /Subtype /Link /Border [0 0 0] /H /I /C [1 0 0] /Rect [168.875 556.76 206.99 568.713] /A << /S /GoTo /D <74686D2E432E322E332E3132> >> >>}
\MusicPDFWriteObject{1769}{<< /Type /Annot /Subtype /Link /Border [0 0 0] /H /I /C [0 1 0] /Rect [299.296 546.099 317.494 554.967] /A << /S /GoTo /D <636974652E4A6172> >> >>}
\MusicPDFWriteObject{1770}{<< /Type /Annot /Subtype /Link /Border [0 0 0] /H /I /C [1 0 0] /Rect [438.203 530.458 475.609 542.413] /A << /S /GoTo /D <74686D2E432E322E332E3133> >> >>}
\MusicPDFWriteObject{1771}{<< /Type /Annot /Subtype /Link /Border [0 0 0] /H /I /C [1 0 0] /Rect [266.859 442.541 298.968 456.185] /A << /S /GoTo /D <74686D2E342E342E322E33> >> >>}
\MusicPDFWriteObject{1772}{<< /Type /Annot /Subtype /Link /Border [0 0 0] /H /I /C [1 0 0] /Rect [487.536 429.054 519.056 443.538] /A << /S /GoTo /D <74686D2E342E342E322E33> >> >>}
\MusicPDFWriteObject{1773}{<< /Type /Annot /Subtype /Link /Border [0 0 0] /H /I /C [1 0 0] /Rect [192.486 290.265 227.679 302.217] /A << /S /GoTo /D <6571756174696F6E2E342E352E312E3130> >> >>}
\MusicPDFWriteObject{1774}{<< /Type /Annot /Subtype /Link /Border [0 0 0] /H /I /C [1 0 0] /Rect [311.665 290.265 343.774 302.217] /A << /S /GoTo /D <74686D2E342E352E312E31> >> >>}
\MusicPDFWriteObject{1775}{<< /Type /Annot /Subtype /Link /Border [0 0 0] /H /I /C [0 1 0] /Rect [350.029 137.689 368.331 146.557] /A << /S /GoTo /D <636974652E4445> >> >>}
\MusicPDFWriteObject{1776}{<< /Type /Annot /Subtype /Link /Border [0 0 0] /H /I /C [1 0 0] /Rect [299.12 108.897 329.333 120.85] /A << /S /GoTo /D <6571756174696F6E2E342E352E322E31> >> >>}
\MusicPDFWriteObject{1777}{<< /Type /Annot /Subtype /Link /Border [0 0 0] /H /I /C [1 0 0] /Rect [414.993 108.537 445.77 120.85] /A << /S /GoTo /D <74686D2E342E352E322E31> >> >>}
\MusicPDFWriteObject{1778}{<< /Type /Annot /Subtype /Link /Border [0 0 0] /H /I /C [1 0 0] /Rect [214.838 684.633 250.031 696.586] /A << /S /GoTo /D <6571756174696F6E2E342E352E312E3130> >> >>}
\MusicPDFWriteObject{1779}{<< /Type /Annot /Subtype /Link /Border [0 0 0] /H /I /C [1 0 0] /Rect [271.824 671.482 307.018 683.435] /A << /S /GoTo /D <6571756174696F6E2E342E352E312E3131> >> >>}
\MusicPDFWriteObject{1780}{<< /Type /Annot /Subtype /Link /Border [0 0 0] /H /I /C [1 0 0] /Rect [367.716 671.122 398.013 683.437] /A << /S /GoTo /D <74686D2E342E352E312E36> >> >>}
\MusicPDFWriteObject{1781}{<< /Type /Annot /Subtype /Link /Border [0 0 0] /H /I /C [1 0 0] /Rect [509.571 671.122 539.869 683.437] /A << /S /GoTo /D <74686D2E342E342E322E36> >> >>}
\MusicPDFWriteObject{1782}{<< /Type /Annot /Subtype /Link /Border [0 0 0] /H /I /C [1 0 0] /Rect [127.805 618.879 158.018 630.832] /A << /S /GoTo /D <6571756174696F6E2E342E352E322E31> >> >>}
\MusicPDFWriteObject{1783}{<< /Type /Annot /Subtype /Link /Border [0 0 0] /H /I /C [1 0 0] /Rect [285.422 581.364 318.11 592.212] /A << /S /GoTo /D <74686D2E442E312E372E31> >> >>}
\MusicPDFWriteObject{1784}{<< /Type /Annot /Subtype /Link /Border [0 0 0] /H /I /C [1 0 0] /Rect [347.307 581.364 372.286 592.212] /A << /S /GoTo /D <73756273656374696F6E2E442E312E37> >> >>}
\MusicPDFWriteObject{1785}{<< /Type /Annot /Subtype /Link /Border [0 0 0] /H /I /C [1 0 0] /Rect [228.347 568.214 258.56 579.005] /A << /S /GoTo /D <74686D2E342E362E322E31> >> >>}
\MusicPDFWriteObject{1786}{<< /Type /Annot /Subtype /Link /Border [0 0 0] /H /I /C [1 0 0] /Rect [304.737 568.214 334.95 579.005] /A << /S /GoTo /D <74686D2E342E362E312E35> >> >>}
\MusicPDFWriteObject{1787}{<< /Type /Annot /Subtype /Link /Border [0 0 0] /H /I /C [1 0 0] /Rect [391.197 568.214 421.41 579.005] /A << /S /GoTo /D <74686D2E342E342E322E36> >> >>}
\MusicPDFWriteObject{1788}{<< /Type /Annot /Subtype /Link /Border [0 0 0] /H /I /C [1 0 0] /Rect [137.566 519.639 169.085 530.43] /A << /S /GoTo /D <74686D2E342E342E322E33> >> >>}
\MusicPDFWriteObject{1789}{<< /Type /Annot /Subtype /Link /Border [0 0 0] /H /I /C [1 0 0] /Rect [448.611 492.784 478.824 504.737] /A << /S /GoTo /D <6571756174696F6E2E342E352E312E31> >> >>}
\MusicPDFWriteObject{1790}{<< /Type /Annot /Subtype /Link /Border [0 0 0] /H /I /C [1 0 0] /Rect [70.342 480.187 101.861 490.978] /A << /S /GoTo /D <74686D2E342E352E312E31> >> >>}
\MusicPDFWriteObject{1791}{<< /Type /Annot /Subtype /Link /Border [0 0 0] /H /I /C [1 0 0] /Rect [300.271 339.652 333.113 351.604] /A << /S /GoTo /D <6571756174696F6E2E442E312E372E32> >> >>}
\MusicPDFWriteObject{1792}{<< /Type /Annot /Subtype /Link /Border [0 0 0] /H /I /C [1 0 0] /Rect [391.628 339.652 425.577 352.534] /A << /S /GoTo /D <74686D2E442E312E372E31> >> >>}
\MusicPDFWriteObject{1793}{<< /Type /Annot /Subtype /Link /Border [0 0 0] /H /I /C [1 0 0] /Rect [278.32 233.352 311.162 245.305] /A << /S /GoTo /D <6571756174696F6E2E442E312E372E32> >> >>}
\MusicPDFWriteObject{1794}{<< /Type /Annot /Subtype /Link /Border [0 0 0] /H /I /C [1 0 0] /Rect [368.619 232.507 401.928 246.235] /A << /S /GoTo /D <74686D2E442E312E372E31> >> >>}
\MusicPDFWriteObject{1795}{<< /Type /Annot /Subtype /Link /Border [0 0 0] /H /I /C [1 0 0] /Rect [473.916 171.626 506.757 183.581] /A << /S /GoTo /D <74686D2E442E332E332E37> >> >>}
\MusicPDFWriteObject{1796}{<< /Type /Annot /Subtype /Link /Border [0 0 0] /H /I /C [1 0 0] /Rect [158.188 158.115 188.401 171.359] /A << /S /GoTo /D <74686D2E342E352E312E31> >> >>}
\MusicPDFWriteObject{1797}{<< /Type /Annot /Subtype /Link /Border [0 0 0] /H /I /C [1 0 0] /Rect [191.122 108.537 221.505 120.852] /A << /S /GoTo /D <74686D2E342E352E312E31> >> >>}
\MusicPDFWriteObject{1798}{<< /Type /Annot /Subtype /Link /Border [0 0 0] /H /I /C [1 0 0] /Rect [502.997 526.697 532.646 537.502] /A << /S /GoTo /D <74686D2E342E352E322E33> >> >>}
\MusicPDFWriteObject{1799}{<< /Type /Annot /Subtype /Link /Border [0 0 0] /H /I /C [1 0 0] /Rect [316.961 263.724 346.609 274.515] /A << /S /GoTo /D <74686D2E342E352E312E36> >> >>}
\MusicPDFWriteObject{1800}{<< /Type /Annot /Subtype /Link /Border [0 0 0] /H /I /C [1 0 0] /Rect [272.483 250.236 302.47 264.72] /A << /S /GoTo /D <74686D2E342E342E322E33> >> >>}
\MusicPDFWriteObject{1801}{<< /Type /Annot /Subtype /Link /Border [0 0 0] /H /I /C [0 1 0] /Rect [262.032 125.766 280.334 134.634] /A << /S /GoTo /D <636974652E4445> >> >>}
\MusicPDFWriteObject{1802}{<< /Type /Annot /Subtype /Link /Border [0 0 0] /H /I /C [0 1 0] /Rect [524.904 673.38 543.206 682.248] /A << /S /GoTo /D <636974652E4445> >> >>}
\MusicPDFWriteObject{1803}{<< /Type /Annot /Subtype /Link /Border [0 0 0] /H /I /C [0 1 0] /Rect [333.661 610.602 351.963 619.47] /A << /S /GoTo /D <636974652E4445> >> >>}
\MusicPDFWriteObject{1804}{<< /Type /Annot /Subtype /Link /Border [0 0 0] /H /I /C [1 0 0] /Rect [514.228 491.239 545.156 502.115] /A << /S /GoTo /D <74686D2E342E362E312E31> >> >>}
\MusicPDFWriteObject{1805}{<< /Type /Annot /Subtype /Link /Border [0 0 0] /H /I /C [1 0 0] /Rect [312.505 179.714 342.718 191.667] /A << /S /GoTo /D <74686D2E342E352E312E31> >> >>}
\MusicPDFWriteObject{1806}{<< /Type /Annot /Subtype /Link /Border [0 0 0] /H /I /C [1 0 0] /Rect [362.739 166.563 401.277 178.518] /A << /S /GoTo /D <74686D2E442E332E332E3132> >> >>}
\MusicPDFWriteObject{1807}{<< /Type /Annot /Subtype /Link /Border [0 0 0] /H /I /C [1 0 0] /Rect [370.722 153.412 403.564 165.367] /A << /S /GoTo /D <74686D2E442E332E322E33> >> >>}
\MusicPDFWriteObject{1808}{<< /Type /Annot /Subtype /Link /Border [0 0 0] /H /I /C [1 0 0] /Rect [405.496 140.262 435.71 152.215] /A << /S /GoTo /D <6571756174696F6E2E342E362E312E35> >> >>}
\MusicPDFWriteObject{1809}{<< /Type /Annot /Subtype /Link /Border [0 0 0] /H /I /C [0 1 0] /Rect [122.216 129.601 147.276 138.469] /A << /S /GoTo /D <636974652E53636861> >> >>}
\MusicPDFWriteObject{1810}{<< /Type /Annot /Subtype /Link /Border [0 0 0] /H /I /C [1 0 0] /Rect [224.511 113.96 254.724 125.913] /A << /S /GoTo /D <6571756174696F6E2E342E362E312E36> >> >>}
\MusicPDFWriteObject{1811}{<< /Type /Annot /Subtype /Link /Border [0 0 0] /H /I /C [0 1 0] /Rect [157.072 687.123 176.176 695.991] /A << /S /GoTo /D <636974652E4E42> >> >>}
\MusicPDFWriteObject{1812}{<< /Type /Annot /Subtype /Link /Border [0 0 0] /H /I /C [1 0 0] /Rect [513.156 629.498 543.228 641.453] /A << /S /GoTo /D <74686D2E342E342E312E31> >> >>}
\MusicPDFWriteObject{1813}{<< /Type /Annot /Subtype /Link /Border [0 0 0] /H /I /C [1 0 0] /Rect [323.614 541.495 356.508 553.448] /A << /S /GoTo /D <74686D2E422E332E322E31> >> >>}
\MusicPDFWriteObject{1814}{<< /Type /Annot /Subtype /Link /Border [0 0 0] /H /I /C [1 0 0] /Rect [438.235 514.43 468.448 526.382] /A << /S /GoTo /D <6571756174696F6E2E342E342E312E33> >> >>}
\MusicPDFWriteObject{1815}{<< /Type /Annot /Subtype /Link /Border [0 0 0] /H /I /C [1 0 0] /Rect [70.655 501.278 100.388 513.234] /A << /S /GoTo /D <74686D2E342E342E312E31> >> >>}
\MusicPDFWriteObject{1816}{<< /Type /Annot /Subtype /Link /Border [0 0 0] /H /I /C [1 0 0] /Rect [345.239 501.278 382.452 513.234] /A << /S /GoTo /D <74686D2E442E332E332E3133> >> >>}
\MusicPDFWriteObject{1817}{<< /Type /Annot /Subtype /Link /Border [0 0 0] /H /I /C [1 0 0] /Rect [284.656 488.128 315.434 500.083] /A << /S /GoTo /D <74686D2E342E332E332E31> >> >>}
\MusicPDFWriteObject{1818}{<< /Type /Annot /Subtype /Link /Border [0 0 0] /H /I /C [1 0 0] /Rect [249.951 435.031 280.164 446.983] /A << /S /GoTo /D <6571756174696F6E2E342E362E312E38> >> >>}
\MusicPDFWriteObject{1819}{<< /Type /Annot /Subtype /Link /Border [0 0 0] /H /I /C [1 0 0] /Rect [193.903 398.603 224.116 410.556] /A << /S /GoTo /D <6571756174696F6E2E342E362E312E37> >> >>}
\MusicPDFWriteObject{1820}{<< /Type /Annot /Subtype /Link /Border [0 0 0] /H /I /C [1 0 0] /Rect [397.367 162.412 432.56 174.365] /A << /S /GoTo /D <6571756174696F6E2E342E362E312E3131> >> >>}
\MusicPDFWriteObject{1821}{<< /Type /Annot /Subtype /Link /Border [0 0 0] /H /I /C [1 0 0] /Rect [199.89 148.864 235.083 160.817] /A << /S /GoTo /D <6571756174696F6E2E342E362E312E3131> >> >>}
\MusicPDFWriteObject{1822}{<< /Type /Annot /Subtype /Link /Border [0 0 0] /H /I /C [1 0 0] /Rect [178.166 685.187 207.843 695.977] /A << /S /GoTo /D <74686D2E342E362E312E33> >> >>}
\MusicPDFWriteObject{1823}{<< /Type /Annot /Subtype /Link /Border [0 0 0] /H /I /C [1 0 0] /Rect [269.62 685.187 299.297 695.977] /A << /S /GoTo /D <74686D2E342E352E322E33> >> >>}
\MusicPDFWriteObject{1824}{<< /Type /Annot /Subtype /Link /Border [0 0 0] /H /I /C [1 0 0] /Rect [250.535 565.615 281.463 577.57] /A << /S /GoTo /D <74686D2E342E362E312E31> >> >>}
\MusicPDFWriteObject{1825}{<< /Type /Annot /Subtype /Link /Border [0 0 0] /H /I /C [1 0 0] /Rect [70.755 673.973 100.968 682.356] /A << /S /GoTo /D <74686D2E372E312E342E32> >> >>}
\MusicPDFWriteObject{1826}{<< /Type /Annot /Subtype /Link /Border [0 0 0] /H /I /C [1 0 0] /Rect [222.958 358.236 255.458 370.189] /A << /S /GoTo /D <74686D2E422E332E312E31> >> >>}
\MusicPDFWriteObject{1827}{<< /Type /Annot /Subtype /Link /Border [0 0 0] /H /I /C [1 0 0] /Rect [473.126 358.236 505.627 370.189] /A << /S /GoTo /D <74686D2E422E332E322E31> >> >>}
\MusicPDFWriteObject{1828}{<< /Type /Annot /Subtype /Link /Border [0 0 0] /H /I /C [1 0 0] /Rect [70.655 264.477 100.868 276.429] /A << /S /GoTo /D <74686D2E342E342E312E31> >> >>}
\MusicPDFWriteObject{1829}{<< /Type /Annot /Subtype /Link /Border [0 0 0] /H /I /C [1 0 0] /Rect [73.712 424.698 103.925 436.65] /A << /S /GoTo /D <6571756174696F6E2E342E362E312E37> >> >>}
\MusicPDFWriteObject{1830}{<< /Type /Annot /Subtype /Link /Border [0 0 0] /H /I /C [1 0 0] /Rect [199.487 424.698 229.7 436.65] /A << /S /GoTo /D <74686D2E342E362E312E33> >> >>}
\MusicPDFWriteObject{1831}{<< /Type /Annot /Subtype /Link /Border [0 0 0] /H /I /C [1 0 0] /Rect [325.722 411.547 360.915 423.5] /A << /S /GoTo /D <6571756174696F6E2E342E362E312E3232> >> >>}
\MusicPDFWriteObject{1832}{<< /Type /Annot /Subtype /Link /Border [0 0 0] /H /I /C [1 0 0] /Rect [98.171 398.396 128.384 410.349] /A << /S /GoTo /D <6571756174696F6E2E342E362E312E37> >> >>}
\MusicPDFWriteObject{1833}{<< /Type /Annot /Subtype /Link /Border [0 0 0] /H /I /C [1 0 0] /Rect [179.974 397.344 209.848 410.351] /A << /S /GoTo /D <74686D2E342E362E312E33> >> >>}
\MusicPDFWriteObject{1834}{<< /Type /Annot /Subtype /Link /Border [0 0 0] /H /I /C [1 0 0] /Rect [432.835 385.245 468.028 397.198] /A << /S /GoTo /D <6571756174696F6E2E342E362E312E3232> >> >>}
\MusicPDFWriteObject{1835}{<< /Type /Annot /Subtype /Link /Border [0 0 0] /H /I /C [0 1 0] /Rect [152.49 348.283 171.594 357.151] /A << /S /GoTo /D <636974652E4E42> >> >>}
\MusicPDFWriteObject{1836}{<< /Type /Annot /Subtype /Link /Border [0 0 0] /H /I /C [1 0 0] /Rect [386.853 319.492 417.066 331.444] /A << /S /GoTo /D <74686D2E342E362E312E33> >> >>}
\MusicPDFWriteObject{1837}{<< /Type /Annot /Subtype /Link /Border [0 0 0] /H /I /C [1 0 0] /Rect [434.521 278.615 465.186 291.497] /A << /S /GoTo /D <74686D2E342E352E312E31> >> >>}
\MusicPDFWriteObject{1838}{<< /Type /Annot /Subtype /Link /Border [0 0 0] /H /I /C [1 0 0] /Rect [175.781 251.953 206.558 265.196] /A << /S /GoTo /D <74686D2E342E332E332E31> >> >>}
\MusicPDFWriteObject{1839}{<< /Type /Annot /Subtype /Link /Border [0 0 0] /H /I /C [1 0 0] /Rect [273.948 163.283 309.142 175.235] /A << /S /GoTo /D <6571756174696F6E2E342E362E312E3136> >> >>}
\MusicPDFWriteObject{1840}{<< /Type /Annot /Subtype /Link /Border [0 0 0] /H /I /C [1 0 0] /Rect [239.346 671.244 277.884 685.164] /A << /S /GoTo /D <74686D2E442E332E332E3133> >> >>}
\MusicPDFWriteObject{1841}{<< /Type /Annot /Subtype /Link /Border [0 0 0] /H /I /C [1 0 0] /Rect [119.726 656.488 149.939 669.495] /A << /S /GoTo /D <74686D2E342E342E312E31> >> >>}
\MusicPDFWriteObject{1842}{<< /Type /Annot /Subtype /Link /Border [0 0 0] /H /I /C [1 0 0] /Rect [168.678 626.234 198.891 638.186] /A << /S /GoTo /D <6571756174696F6E2E342E342E322E36> >> >>}
\MusicPDFWriteObject{1843}{<< /Type /Annot /Subtype /Link /Border [0 0 0] /H /I /C [1 0 0] /Rect [112.685 565.256 145.496 578.121] /A << /S /GoTo /D <74686D2E442E312E322E31> >> >>}
\MusicPDFWriteObject{1844}{<< /Type /Annot /Subtype /Link /Border [0 0 0] /H /I /C [1 0 0] /Rect [185.308 405.373 214.957 418.38] /A << /S /GoTo /D <74686D2E342E342E322E36> >> >>}
\MusicPDFWriteObject{1845}{<< /Type /Annot /Subtype /Link /Border [0 0 0] /H /I /C [1 0 0] /Rect [145.75 366.973 175.964 378.926] /A << /S /GoTo /D <6571756174696F6E2E342E342E322E35> >> >>}
\MusicPDFWriteObject{1846}{<< /Type /Annot /Subtype /Link /Border [0 0 0] /H /I /C [1 0 0] /Rect [140.121 354.376 170.334 365.166] /A << /S /GoTo /D <74686D2E342E342E322E36> >> >>}
\MusicPDFWriteObject{1847}{<< /Type /Annot /Subtype /Link /Border [0 0 0] /H /I /C [1 0 0] /Rect [199.241 291.093 229.454 303.046] /A << /S /GoTo /D <6571756174696F6E2E342E342E322E36> >> >>}
\MusicPDFWriteObject{1848}{<< /Type /Annot /Subtype /Link /Border [0 0 0] /H /I /C [1 0 0] /Rect [281.433 290.733 311.562 303.046] /A << /S /GoTo /D <74686D2E342E342E322E36> >> >>}
\MusicPDFWriteObject{1849}{<< /Type /Annot /Subtype /Link /Border [0 0 0] /H /I /C [1 0 0] /Rect [155.859 209.859 186.072 220.65] /A << /S /GoTo /D <74686D2E342E362E312E34> >> >>}
\MusicPDFWriteObject{1850}{<< /Type /Annot /Subtype /Link /Border [0 0 0] /H /I /C [1 0 0] /Rect [294.868 680.978 330.061 692.931] /A << /S /GoTo /D <6571756174696F6E2E342E362E312E3236> >> >>}
\MusicPDFWriteObject{1851}{<< /Type /Annot /Subtype /Link /Border [0 0 0] /H /I /C [1 0 0] /Rect [289.66 628.375 324.853 640.328] /A << /S /GoTo /D <6571756174696F6E2E342E362E312E3235> >> >>}
\MusicPDFWriteObject{1852}{<< /Type /Annot /Subtype /Link /Border [0 0 0] /H /I /C [1 0 0] /Rect [159.555 614.864 189.203 626.569] /A << /S /GoTo /D <74686D2E342E342E322E36> >> >>}
\MusicPDFWriteObject{1853}{<< /Type /Annot /Subtype /Link /Border [0 0 0] /H /I /C [1 0 0] /Rect [176.206 601.022 206.984 614.028] /A << /S /GoTo /D <74686D2E342E342E322E36> >> >>}
\MusicPDFWriteObject{1854}{<< /Type /Annot /Subtype /Link /Border [0 0 0] /H /I /C [1 0 0] /Rect [168.766 490.391 203.959 502.343] /A << /S /GoTo /D <6571756174696F6E2E342E362E312E3237> >> >>}
\MusicPDFWriteObject{1855}{<< /Type /Annot /Subtype /Link /Border [0 0 0] /H /I /C [1 0 0] /Rect [273.659 477.24 308.852 489.193] /A << /S /GoTo /D <6571756174696F6E2E342E362E312E3237> >> >>}
\MusicPDFWriteObject{1856}{<< /Type /Annot /Subtype /Link /Border [0 0 0] /H /I /C [1 0 0] /Rect [73.712 424.637 108.905 436.59] /A << /S /GoTo /D <6571756174696F6E2E342E362E312E3236> >> >>}
\MusicPDFWriteObject{1857}{<< /Type /Annot /Subtype /Link /Border [0 0 0] /H /I /C [1 0 0] /Rect [389.699 385.185 424.892 397.137] /A << /S /GoTo /D <6571756174696F6E2E342E362E312E3234> >> >>}
\MusicPDFWriteObject{1858}{<< /Type /Annot /Subtype /Link /Border [0 0 0] /H /I /C [1 0 0] /Rect [120.947 265.293 156.141 277.245] /A << /S /GoTo /D <6571756174696F6E2E342E362E312E3237> >> >>}
\MusicPDFWriteObject{1859}{<< /Type /Annot /Subtype /Link /Border [0 0 0] /H /I /C [0 1 0] /Rect [107.544 116.255 134.035 125.123] /A << /S /GoTo /D <636974652E52752D33> >> >>}
\MusicPDFWriteObject{1860}{<< /Type /Annot /Subtype /Link /Border [0 0 0] /H /I /C [1 0 0] /Rect [476.585 632.03 511.778 643.983] /A << /S /GoTo /D <6571756174696F6E2E342E362E312E3135> >> >>}
\MusicPDFWriteObject{1861}{<< /Type /Annot /Subtype /Link /Border [0 0 0] /H /I /C [1 0 0] /Rect [511.675 491.513 541.888 503.466] /A << /S /GoTo /D <6571756174696F6E2E342E362E312E37> >> >>}
\MusicPDFWriteObject{1862}{<< /Type /Annot /Subtype /Link /Border [0 0 0] /H /I /C [1 0 0] /Rect [161.365 478.362 191.578 490.315] /A << /S /GoTo /D <74686D2E342E362E312E33> >> >>}
\MusicPDFWriteObject{1863}{<< /Type /Annot /Subtype /Link /Border [0 0 0] /H /I /C [1 0 0] /Rect [434.496 373.71 464.709 384.5] /A << /S /GoTo /D <74686D2E342E362E312E34> >> >>}
\MusicPDFWriteObject{1864}{<< /Type /Annot /Subtype /Link /Border [0 0 0] /H /I /C [1 0 0] /Rect [383.852 360.006 419.045 371.958] /A << /S /GoTo /D <6571756174696F6E2E342E362E312E3137> >> >>}
\MusicPDFWriteObject{1865}{<< /Type /Annot /Subtype /Link /Border [0 0 0] /H /I /C [1 0 0] /Rect [131.152 346.855 166.345 358.808] /A << /S /GoTo /D <6571756174696F6E2E342E362E312E3136> >> >>}
\MusicPDFWriteObject{1866}{<< /Type /Annot /Subtype /Link /Border [0 0 0] /H /I /C [0 1 0] /Rect [491.672 296.742 518.164 305.61] /A << /S /GoTo /D <636974652E52752D33> >> >>}
\MusicPDFWriteObject{1867}{<< /Type /Annot /Subtype /Link /Border [0 0 0] /H /I /C [1 0 0] /Rect [190.309 251.163 220.522 261.954] /A << /S /GoTo /D <74686D2E342E362E312E35> >> >>}
\MusicPDFWriteObject{1868}{<< /Type /Annot /Subtype /Link /Border [0 0 0] /H /I /C [1 0 0] /Rect [364.149 213.201 394.221 224.078] /A << /S /GoTo /D <74686D2E342E362E312E35> >> >>}
\MusicPDFWriteObject{1869}{<< /Type /Annot /Subtype /Link /Border [0 0 0] /H /I /C [1 0 0] /Rect [387.485 183.448 418.262 199.25] /A << /S /GoTo /D <74686D2E342E342E312E31> >> >>}
\MusicPDFWriteObject{1870}{<< /Type /Annot /Subtype /Link /Border [0 0 0] /H /I /C [1 0 0] /Rect [254.843 122.048 285.056 134.001] /A << /S /GoTo /D <6571756174696F6E2E342E362E312E37> >> >>}
\MusicPDFWriteObject{1871}{<< /Type /Annot /Subtype /Link /Border [0 0 0] /H /I /C [1 0 0] /Rect [337.325 122.048 367.679 134.003] /A << /S /GoTo /D <74686D2E342E362E312E33> >> >>}
\MusicPDFWriteObject{1872}{<< /Type /Annot /Subtype /Link /Border [0 0 0] /H /I /C [1 0 0] /Rect [336.169 95.747 371.362 107.699] /A << /S /GoTo /D <6571756174696F6E2E342E362E312E3239> >> >>}
\MusicPDFWriteObject{1873}{<< /Type /Annot /Subtype /Link /Border [0 0 0] /H /I /C [1 0 0] /Rect [359.567 685.186 390.345 695.977] /A << /S /GoTo /D <74686D2E342E362E312E33> >> >>}
\MusicPDFWriteObject{1874}{<< /Type /Annot /Subtype /Link /Border [0 0 0] /H /I /C [1 0 0] /Rect [192.592 634.99 222.805 646.943] /A << /S /GoTo /D <6571756174696F6E2E342E322E312E31> >> >>}
\MusicPDFWriteObject{1875}{<< /Type /Annot /Subtype /Link /Border [0 0 0] /H /I /C [1 0 0] /Rect [348.499 634.99 383.692 646.943] /A << /S /GoTo /D <6571756174696F6E2E342E362E312E3330> >> >>}
\MusicPDFWriteObject{1876}{<< /Type /Annot /Subtype /Link /Border [0 0 0] /H /I /C [1 0 0] /Rect [234.438 582.387 269.632 594.34] /A << /S /GoTo /D <6571756174696F6E2E342E362E312E3330> >> >>}
\MusicPDFWriteObject{1877}{<< /Type /Annot /Subtype /Link /Border [0 0 0] /H /I /C [1 0 0] /Rect [320.311 492.682 350.524 504.635] /A << /S /GoTo /D <6571756174696F6E2E342E322E312E34> >> >>}
\MusicPDFWriteObject{1878}{<< /Type /Annot /Subtype /Link /Border [0 0 0] /H /I /C [1 0 0] /Rect [402.468 492.682 432.569 504.637] /A << /S /GoTo /D <74686D2E342E322E312E32> >> >>}
\MusicPDFWriteObject{1879}{<< /Type /Annot /Subtype /Link /Border [0 0 0] /H /I /C [1 0 0] /Rect [323.1 385.716 354.619 396.507] /A << /S /GoTo /D <74686D2E342E362E312E31> >> >>}
\MusicPDFWriteObject{1880}{<< /Type /Annot /Subtype /Link /Border [0 0 0] /H /I /C [1 0 0] /Rect [221.475 95.747 251.689 107.699] /A << /S /GoTo /D <6571756174696F6E2E342E362E312E37> >> >>}
\MusicPDFWriteObject{1881}{<< /Type /Annot /Subtype /Link /Border [0 0 0] /H /I /C [1 0 0] /Rect [349.413 95.747 380.191 107.699] /A << /S /GoTo /D <74686D2E342E362E312E33> >> >>}
\MusicPDFWriteObject{1882}{<< /Type /Annot /Subtype /Link /Border [0 0 0] /H /I /C [1 0 0] /Rect [70.655 671.482 100.868 683.435] /A << /S /GoTo /D <74686D2E342E352E322E31> >> >>}
\MusicPDFWriteObject{1883}{<< /Type /Annot /Subtype /Link /Border [0 0 0] /H /I /C [1 0 0] /Rect [187.909 671.482 218.122 683.435] /A << /S /GoTo /D <74686D2E342E352E312E31> >> >>}
\MusicPDFWriteObject{1884}{<< /Type /Annot /Subtype /Link /Border [0 0 0] /H /I /C [1 0 0] /Rect [427.863 585.477 458.076 597.429] /A << /S /GoTo /D <6571756174696F6E2E342E362E322E35> >> >>}
\MusicPDFWriteObject{1885}{<< /Type /Annot /Subtype /Link /Border [0 0 0] /H /I /C [1 0 0] /Rect [511.675 585.477 541.888 597.429] /A << /S /GoTo /D <6571756174696F6E2E342E362E322E34> >> >>}
\MusicPDFWriteObject{1886}{<< /Type /Annot /Subtype /Link /Border [0 0 0] /H /I /C [1 0 0] /Rect [335.05 530.446 365.263 542.398] /A << /S /GoTo /D <6571756174696F6E2E342E362E322E33> >> >>}
\MusicPDFWriteObject{1887}{<< /Type /Annot /Subtype /Link /Border [0 0 0] /H /I /C [1 0 0] /Rect [430.389 485.93 460.603 497.883] /A << /S /GoTo /D <6571756174696F6E2E342E362E322E35> >> >>}
\MusicPDFWriteObject{1888}{<< /Type /Annot /Subtype /Link /Border [0 0 0] /H /I /C [1 0 0] /Rect [180.226 472.78 210.439 484.732] /A << /S /GoTo /D <6571756174696F6E2E342E362E322E36> >> >>}
\MusicPDFWriteObject{1889}{<< /Type /Annot /Subtype /Link /Border [0 0 0] /H /I /C [0 1 0] /Rect [242.485 462.119 261.59 470.987] /A << /S /GoTo /D <636974652E4E42> >> >>}
\MusicPDFWriteObject{1890}{<< /Type /Annot /Subtype /Link /Border [0 0 0] /H /I /C [1 0 0] /Rect [206.509 446.478 236.722 458.431] /A << /S /GoTo /D <6571756174696F6E2E342E362E322E33> >> >>}
\MusicPDFWriteObject{1891}{<< /Type /Annot /Subtype /Link /Border [0 0 0] /H /I /C [1 0 0] /Rect [446.276 234.286 477.795 246.241] /A << /S /GoTo /D <74686D2E342E362E312E31> >> >>}
\MusicPDFWriteObject{1892}{<< /Type /Annot /Subtype /Link /Border [0 0 0] /H /I /C [1 0 0] /Rect [125.342 194.474 156.713 206.788] /A << /S /GoTo /D <74686D2E342E362E312E31> >> >>}
\MusicPDFWriteObject{1893}{<< /Type /Annot /Subtype /Link /Border [0 0 0] /H /I /C [1 0 0] /Rect [421.976 194.834 452.189 206.787] /A << /S /GoTo /D <6571756174696F6E2E342E352E332E33> >> >>}
\MusicPDFWriteObject{1894}{<< /Type /Annot /Subtype /Link /Border [0 0 0] /H /I /C [1 0 0] /Rect [186.648 95.746 217.172 107.701] /A << /S /GoTo /D <74686D2E342E362E312E31> >> >>}
\MusicPDFWriteObject{1895}{<< /Type /Annot /Subtype /Link /Border [0 0 0] /H /I /C [1 0 0] /Rect [501.574 81.267 531.787 95.188] /A << /S /GoTo /D <74686D2E342E362E312E36> >> >>}
\MusicPDFWriteObject{1896}{<< /Type /Annot /Subtype /Link /Border [0 0 0] /H /I /C [1 0 0] /Rect [148.615 539.486 178.828 551.439] /A << /S /GoTo /D <74686D2E342E352E332E31> >> >>}
\MusicPDFWriteObject{1897}{<< /Type /Annot /Subtype /Link /Border [0 0 0] /H /I /C [1 0 0] /Rect [328.12 539.486 358.333 551.439] /A << /S /GoTo /D <6571756174696F6E2E342E362E332E32> >> >>}
\MusicPDFWriteObject{1898}{<< /Type /Annot /Subtype /Link /Border [0 0 0] /H /I /C [1 0 0] /Rect [394.419 525.975 424.067 538.288] /A << /S /GoTo /D <74686D2E342E352E312E31> >> >>}
\MusicPDFWriteObject{1899}{<< /Type /Annot /Subtype /Link /Border [0 0 0] /H /I /C [1 0 0] /Rect [514.169 525.975 543.817 538.288] /A << /S /GoTo /D <74686D2E342E352E332E31> >> >>}
\MusicPDFWriteObject{1900}{<< /Type /Annot /Subtype /Link /Border [0 0 0] /H /I /C [1 0 0] /Rect [477.42 499.673 507.068 511.988] /A << /S /GoTo /D <74686D2E342E362E312E33> >> >>}
\MusicPDFWriteObject{1901}{<< /Type /Annot /Subtype /Link /Border [0 0 0] /H /I /C [1 0 0] /Rect [331.152 461.135 361.365 471.926] /A << /S /GoTo /D <74686D2E342E352E332E31> >> >>}
\MusicPDFWriteObject{1902}{<< /Type /Annot /Subtype /Link /Border [0 0 0] /H /I /C [1 0 0] /Rect [341.74 419.903 371.502 431.856] /A << /S /GoTo /D <74686D2E342E332E332E31> >> >>}
\MusicPDFWriteObject{1903}{<< /Type /Annot /Subtype /Link /Border [0 0 0] /H /I /C [1 0 0] /Rect [272.579 392.931 302.792 404.884] /A << /S /GoTo /D <74686D2E352E322E332E32> >> >>}
\MusicPDFWriteObject{1904}{<< /Type /Annot /Subtype /Link /Border [0 0 0] /H /I /C [1 0 0] /Rect [330.292 366.629 360.505 379.222] /A << /S /GoTo /D <74686D2E342E362E312E33> >> >>}
\MusicPDFWriteObject{1905}{<< /Type /Annot /Subtype /Link /Border [0 0 0] /H /I /C [1 0 0] /Rect [293.028 353.479 323.241 365.432] /A << /S /GoTo /D <6571756174696F6E2E342E362E332E32> >> >>}
\MusicPDFWriteObject{1906}{<< /Type /Annot /Subtype /Link /Border [0 0 0] /H /I /C [1 0 0] /Rect [175.017 340.328 208.414 352.283] /A << /S /GoTo /D <74686D2E442E332E332E37> >> >>}
\MusicPDFWriteObject{1907}{<< /Type /Annot /Subtype /Link /Border [0 0 0] /H /I /C [1 0 0] /Rect [71.004 327.177 108.646 339.132] /A << /S /GoTo /D <74686D2E442E332E332E3131> >> >>}
\MusicPDFWriteObject{1908}{<< /Type /Annot /Subtype /Link /Border [0 0 0] /H /I /C [1 0 0] /Rect [292.895 314.027 323.109 325.979] /A << /S /GoTo /D <6571756174696F6E2E342E362E332E32> >> >>}
\MusicPDFWriteObject{1909}{<< /Type /Annot /Subtype /Link /Border [0 0 0] /H /I /C [0 1 0] /Rect [93.41 303.366 122.171 312.234] /A << /S /GoTo /D <636974652E47726F2D66> >> >>}
\MusicPDFWriteObject{1910}{<< /Type /Annot /Subtype /Link /Border [0 0 0] /H /I /C [0 1 0] /Rect [264.231 303.366 281.551 312.234] /A << /S /GoTo /D <636974652E534E> >> >>}
\MusicPDFWriteObject{1911}{<< /Type /Annot /Subtype /Link /Border [0 0 0] /H /I /C [0 1 0] /Rect [219.682 263.914 244.742 272.782] /A << /S /GoTo /D <636974652E53636861> >> >>}
\MusicPDFWriteObject{1912}{<< /Type /Annot /Subtype /Link /Border [0 0 0] /H /I /C [0 1 0] /Rect [87.002 250.763 112.062 259.631] /A << /S /GoTo /D <636974652E53636861> >> >>}
\MusicPDFWriteObject{1913}{<< /Type /Annot /Subtype /Link /Border [0 0 0] /H /I /C [1 0 0] /Rect [484.322 235.122 514.535 247.075] /A << /S /GoTo /D <6571756174696F6E2E342E362E332E32> >> >>}
\MusicPDFWriteObject{1914}{<< /Type /Annot /Subtype /Link /Border [0 0 0] /H /I /C [0 1 0] /Rect [362.063 224.461 381.168 233.329] /A << /S /GoTo /D <636974652E4E42> >> >>}
\MusicPDFWriteObject{1915}{<< /Type /Annot /Subtype /Link /Border [0 0 0] /H /I /C [1 0 0] /Rect [82.959 558.462 429.398 569.309] /A << /S /GoTo /D <73656374696F6E2E352E31> >> >>}
\MusicPDFWriteObject{1916}{<< /Type /Annot /Subtype /Link /Border [0 0 0] /H /I /C [1 0 0] /Rect [82.959 545.311 493.186 556.159] /A << /S /GoTo /D <73756273656374696F6E2E352E312E31> >> >>}
\MusicPDFWriteObject{1917}{<< /Type /Annot /Subtype /Link /Border [0 0 0] /H /I /C [1 0 0] /Rect [82.959 532.16 398.908 543.008] /A << /S /GoTo /D <73756273656374696F6E2E352E312E32> >> >>}
\MusicPDFWriteObject{1918}{<< /Type /Annot /Subtype /Link /Border [0 0 0] /H /I /C [1 0 0] /Rect [82.959 519.009 313.704 529.8] /A << /S /GoTo /D <73756273656374696F6E2E352E312E33> >> >>}
\MusicPDFWriteObject{1919}{<< /Type /Annot /Subtype /Link /Border [0 0 0] /H /I /C [1 0 0] /Rect [82.959 505.859 393.94 516.706] /A << /S /GoTo /D <73656374696F6E2E352E32> >> >>}
\MusicPDFWriteObject{1920}{<< /Type /Annot /Subtype /Link /Border [0 0 0] /H /I /C [1 0 0] /Rect [82.959 490.361 376.833 505.216] /A << /S /GoTo /D <73756273656374696F6E2E352E322E31> >> >>}
\MusicPDFWriteObject{1921}{<< /Type /Annot /Subtype /Link /Border [0 0 0] /H /I /C [1 0 0] /Rect [82.959 476.502 372.24 491.357] /A << /S /GoTo /D <73756273656374696F6E2E352E322E32> >> >>}
\MusicPDFWriteObject{1922}{<< /Type /Annot /Subtype /Link /Border [0 0 0] /H /I /C [1 0 0] /Rect [82.959 464.784 413.543 477.099] /A << /S /GoTo /D <73756273656374696F6E2E352E322E33> >> >>}
\MusicPDFWriteObject{1923}{<< /Type /Annot /Subtype /Link /Border [0 0 0] /H /I /C [1 0 0] /Rect [82.959 454.484 288.361 463.338] /A << /S /GoTo /D <73756273656374696F6E2E352E322E34> >> >>}
\MusicPDFWriteObject{1924}{<< /Type /Annot /Subtype /Link /Border [0 0 0] /H /I /C [1 0 0] /Rect [82.959 439.397 253.777 450.187] /A << /S /GoTo /D <73656374696F6E2E352E33> >> >>}
\MusicPDFWriteObject{1925}{<< /Type /Annot /Subtype /Link /Border [0 0 0] /H /I /C [1 0 0] /Rect [82.959 426.246 313.787 437.037] /A << /S /GoTo /D <73756273656374696F6E2E352E332E31> >> >>}
\MusicPDFWriteObject{1926}{<< /Type /Annot /Subtype /Link /Border [0 0 0] /H /I /C [1 0 0] /Rect [82.959 413.095 394.244 423.886] /A << /S /GoTo /D <73756273656374696F6E2E352E332E32> >> >>}
\MusicPDFWriteObject{1927}{<< /Type /Annot /Subtype /Link /Border [0 0 0] /H /I /C [1 0 0] /Rect [82.959 399.944 336.723 410.735] /A << /S /GoTo /D <73756273656374696F6E2E352E332E33> >> >>}
\MusicPDFWriteObject{1928}{<< /Type /Annot /Subtype /Link /Border [0 0 0] /H /I /C [1 0 0] /Rect [82.959 386.794 439.55 397.584] /A << /S /GoTo /D <73756273656374696F6E2E352E332E34> >> >>}
\MusicPDFWriteObject{1929}{<< /Type /Annot /Subtype /Link /Border [0 0 0] /H /I /C [1 0 0] /Rect [82.959 373.643 306.471 386.453] /A << /S /GoTo /D <73656374696F6E2E352E34> >> >>}
\MusicPDFWriteObject{1930}{<< /Type /Annot /Subtype /Link /Border [0 0 0] /H /I /C [1 0 0] /Rect [82.959 362.429 234.991 371.283] /A << /S /GoTo /D <73756273656374696F6E2E352E342E31> >> >>}
\MusicPDFWriteObject{1931}{<< /Type /Annot /Subtype /Link /Border [0 0 0] /H /I /C [1 0 0] /Rect [82.959 347.341 218.185 360.152] /A << /S /GoTo /D <73756273656374696F6E2E352E342E32> >> >>}
\MusicPDFWriteObject{1932}{<< /Type /Annot /Subtype /Link /Border [0 0 0] /H /I /C [1 0 0] /Rect [82.959 334.191 369.592 344.981] /A << /S /GoTo /D <73656374696F6E2E352E35> >> >>}
\MusicPDFWriteObject{1933}{<< /Type /Annot /Subtype /Link /Border [0 0 0] /H /I /C [1 0 0] /Rect [82.959 318.693 509.22 333.548] /A << /S /GoTo /D <73756273656374696F6E2E352E352E31> >> >>}
\MusicPDFWriteObject{1934}{<< /Type /Annot /Subtype /Link /Border [0 0 0] /H /I /C [1 0 0] /Rect [82.959 307.335 376.302 319.29] /A << /S /GoTo /D <73756273656374696F6E2E352E352E32> >> >>}
\MusicPDFWriteObject{1935}{<< /Type /Annot /Subtype /Link /Border [0 0 0] /H /I /C [1 0 0] /Rect [82.959 296.675 229.485 305.529] /A << /S /GoTo /D <73756273656374696F6E2E352E352E33> >> >>}
\MusicPDFWriteObject{1936}{<< /Type /Annot /Subtype /Link /Border [0 0 0] /H /I /C [1 0 0] /Rect [123.482 669.689 153.131 684.544] /A << /S /GoTo /D <74686D2E322E312E342E31> >> >>}
\MusicPDFWriteObject{1937}{<< /Type /Annot /Subtype /Link /Border [0 0 0] /H /I /C [1 0 0] /Rect [324.323 618.171 355.1 630.126] /A << /S /GoTo /D <74686D2E332E322E322E31> >> >>}
\MusicPDFWriteObject{1938}{<< /Type /Annot /Subtype /Link /Border [0 0 0] /H /I /C [1 0 0] /Rect [371.316 605.02 403.24 616.973] /A << /S /GoTo /D <74686D2E422E332E322E31> >> >>}
\MusicPDFWriteObject{1939}{<< /Type /Annot /Subtype /Link /Border [0 0 0] /H /I /C [0 1 0] /Rect [202.131 554.907 225.973 563.775] /A << /S /GoTo /D <636974652E4D6531> >> >>}
\MusicPDFWriteObject{1940}{<< /Type /Annot /Subtype /Link /Border [0 0 0] /H /I /C [0 1 0] /Rect [423.436 554.907 447.278 563.775] /A << /S /GoTo /D <636974652E4D6531> >> >>}
\MusicPDFWriteObject{1941}{<< /Type /Annot /Subtype /Link /Border [0 0 0] /H /I /C [0 1 0] /Rect [203.487 528.606 227.33 537.474] /A << /S /GoTo /D <636974652E4D6531> >> >>}
\MusicPDFWriteObject{1942}{<< /Type /Annot /Subtype /Link /Border [0 0 0] /H /I /C [0 1 0] /Rect [329.377 528.606 353.22 537.474] /A << /S /GoTo /D <636974652E4D6531> >> >>}
\MusicPDFWriteObject{1943}{<< /Type /Annot /Subtype /Link /Border [0 0 0] /H /I /C [1 0 0] /Rect [321.924 513.518 344.042 524.366] /A << /S /GoTo /D <73756273656374696F6E2E352E332E31> >> >>}
\MusicPDFWriteObject{1944}{<< /Type /Annot /Subtype /Link /Border [0 0 0] /H /I /C [1 0 0] /Rect [169.441 500.367 200.219 511.215] /A << /S /GoTo /D <74686D2E352E342E312E31> >> >>}
\MusicPDFWriteObject{1945}{<< /Type /Annot /Subtype /Link /Border [0 0 0] /H /I /C [1 0 0] /Rect [381.021 500.367 411.799 511.215] /A << /S /GoTo /D <74686D2E352E342E322E32> >> >>}
\MusicPDFWriteObject{1946}{<< /Type /Annot /Subtype /Link /Border [0 0 0] /H /I /C [1 0 0] /Rect [441.585 500.367 464.461 511.215] /A << /S /GoTo /D <73756273656374696F6E2E352E342E32> >> >>}
\MusicPDFWriteObject{1947}{<< /Type /Annot /Subtype /Link /Border [0 0 0] /H /I /C [1 0 0] /Rect [112.727 486.664 142.94 498.616] /A << /S /GoTo /D <74686D2E352E342E312E31> >> >>}
\MusicPDFWriteObject{1948}{<< /Type /Annot /Subtype /Link /Border [0 0 0] /H /I /C [0 1 0] /Rect [296.576 489.154 320.419 498.021] /A << /S /GoTo /D <636974652E4D6531> >> >>}
\MusicPDFWriteObject{1949}{<< /Type /Annot /Subtype /Link /Border [0 0 0] /H /I /C [0 1 0] /Rect [356.159 476.003 380.002 484.871] /A << /S /GoTo /D <636974652E4D6531> >> >>}
\MusicPDFWriteObject{1950}{<< /Type /Annot /Subtype /Link /Border [0 0 0] /H /I /C [0 1 0] /Rect [444.391 423.4 462.507 432.268] /A << /S /GoTo /D <636974652E4D65> >> >>}
\MusicPDFWriteObject{1951}{<< /Type /Annot /Subtype /Link /Border [0 0 0] /H /I /C [0 1 0] /Rect [321.499 383.948 345.342 392.815] /A << /S /GoTo /D <636974652E4D6531> >> >>}
\MusicPDFWriteObject{1952}{<< /Type /Annot /Subtype /Link /Border [0 0 0] /H /I /C [0 1 0] /Rect [391.811 383.948 415.654 392.815] /A << /S /GoTo /D <636974652E4D6532> >> >>}
\MusicPDFWriteObject{1953}{<< /Type /Annot /Subtype /Link /Border [0 0 0] /H /I /C [0 1 0] /Rect [441.692 383.948 465.535 392.815] /A << /S /GoTo /D <636974652E4D6533> >> >>}
\MusicPDFWriteObject{1954}{<< /Type /Annot /Subtype /Link /Border [0 0 0] /H /I /C [0 1 0] /Rect [188.811 344.495 207.916 353.363] /A << /S /GoTo /D <636974652E4E42> >> >>}
\MusicPDFWriteObject{1955}{<< /Type /Annot /Subtype /Link /Border [0 0 0] /H /I /C [1 0 0] /Rect [259.975 318.194 290.188 327.048] /A << /S /GoTo /D <74686D2E352E322E332E34> >> >>}
\MusicPDFWriteObject{1956}{<< /Type /Annot /Subtype /Link /Border [0 0 0] /H /I /C [1 0 0] /Rect [263.373 248.328 293.586 260.281] /A << /S /GoTo /D <6571756174696F6E2E332E322E322E31> >> >>}
\MusicPDFWriteObject{1957}{<< /Type /Annot /Subtype /Link /Border [0 0 0] /H /I /C [1 0 0] /Rect [371.358 248.328 401.571 260.281] /A << /S /GoTo /D <6571756174696F6E2E332E332E312E31> >> >>}
\MusicPDFWriteObject{1958}{<< /Type /Annot /Subtype /Link /Border [0 0 0] /H /I /C [1 0 0] /Rect [332.274 201.279 362.713 213.593] /A << /S /GoTo /D <74686D2E332E322E322E31> >> >>}
\MusicPDFWriteObject{1959}{<< /Type /Annot /Subtype /Link /Border [0 0 0] /H /I /C [1 0 0] /Rect [463.909 201.279 494.348 213.593] /A << /S /GoTo /D <74686D2E332E332E312E31> >> >>}
\MusicPDFWriteObject{1960}{<< /Type /Annot /Subtype /Link /Border [0 0 0] /H /I /C [1 0 0] /Rect [136.172 162.187 166.385 174.139] /A << /S /GoTo /D <6571756174696F6E2E332E332E312E33> >> >>}
\MusicPDFWriteObject{1961}{<< /Type /Annot /Subtype /Link /Border [0 0 0] /H /I /C [1 0 0] /Rect [254.173 570.655 284.387 582.607] /A << /S /GoTo /D <6571756174696F6E2E332E322E322E31> >> >>}
\MusicPDFWriteObject{1962}{<< /Type /Annot /Subtype /Link /Border [0 0 0] /H /I /C [1 0 0] /Rect [350.29 494.775 380.503 506.727] /A << /S /GoTo /D <6571756174696F6E2E352E312E312E31> >> >>}
\MusicPDFWriteObject{1963}{<< /Type /Annot /Subtype /Link /Border [0 0 0] /H /I /C [1 0 0] /Rect [295.108 468.473 325.321 480.426] /A << /S /GoTo /D <6571756174696F6E2E332E332E312E31> >> >>}
\MusicPDFWriteObject{1964}{<< /Type /Annot /Subtype /Link /Border [0 0 0] /H /I /C [1 0 0] /Rect [297.107 379.443 327.32 391.395] /A << /S /GoTo /D <6571756174696F6E2E332E332E312E33> >> >>}
\MusicPDFWriteObject{1965}{<< /Type /Annot /Subtype /Link /Border [0 0 0] /H /I /C [1 0 0] /Rect [511.675 379.443 541.888 391.395] /A << /S /GoTo /D <6571756174696F6E2E352E312E312E32> >> >>}
\MusicPDFWriteObject{1966}{<< /Type /Annot /Subtype /Link /Border [0 0 0] /H /I /C [1 0 0] /Rect [237.454 267.135 267.667 279.088] /A << /S /GoTo /D <6571756174696F6E2E352E312E312E31> >> >>}
\MusicPDFWriteObject{1967}{<< /Type /Annot /Subtype /Link /Border [0 0 0] /H /I /C [1 0 0] /Rect [497.197 267.135 527.41 279.088] /A << /S /GoTo /D <6571756174696F6E2E352E312E312E32> >> >>}
\MusicPDFWriteObject{1968}{<< /Type /Annot /Subtype /Link /Border [0 0 0] /H /I /C [1 0 0] /Rect [446.291 253.984 476.504 265.937] /A << /S /GoTo /D <6571756174696F6E2E352E312E312E34> >> >>}
\MusicPDFWriteObject{1969}{<< /Type /Annot /Subtype /Link /Border [0 0 0] /H /I /C [0 1 0] /Rect [165.239 161.909 190.299 170.777] /A << /S /GoTo /D <636974652E53636861> >> >>}
\MusicPDFWriteObject{1970}{<< /Type /Annot /Subtype /Link /Border [0 0 0] /H /I /C [1 0 0] /Rect [349.143 146.268 379.356 158.221] /A << /S /GoTo /D <74686D2E352E312E312E33> >> >>}
\MusicPDFWriteObject{1971}{<< /Type /Annot /Subtype /Link /Border [0 0 0] /H /I /C [1 0 0] /Rect [212.868 519.735 243.646 530.526] /A << /S /GoTo /D <74686D2E352E312E312E33> >> >>}
\MusicPDFWriteObject{1972}{<< /Type /Annot /Subtype /Link /Border [0 0 0] /H /I /C [1 0 0] /Rect [513.621 448.004 543.27 458.851] /A << /S /GoTo /D <74686D2E352E312E312E33> >> >>}
\MusicPDFWriteObject{1973}{<< /Type /Annot /Subtype /Link /Border [0 0 0] /H /I /C [1 0 0] /Rect [508.961 309.745 539.174 321.697] /A << /S /GoTo /D <6571756174696F6E2E352E312E312E31> >> >>}
\MusicPDFWriteObject{1974}{<< /Type /Annot /Subtype /Link /Border [0 0 0] /H /I /C [1 0 0] /Rect [450.573 150.293 480.532 162.248] /A << /S /GoTo /D <74686D2E332E322E322E31> >> >>}
\MusicPDFWriteObject{1975}{<< /Type /Annot /Subtype /Link /Border [0 0 0] /H /I /C [1 0 0] /Rect [445.277 669.976 474.926 682.291] /A << /S /GoTo /D <74686D2E322E322E352E34> >> >>}
\MusicPDFWriteObject{1976}{<< /Type /Annot /Subtype /Link /Border [0 0 0] /H /I /C [1 0 0] /Rect [404.977 641.793 434.88 657.822] /A << /S /GoTo /D <74686D2E342E332E322E32> >> >>}
\MusicPDFWriteObject{1977}{<< /Type /Annot /Subtype /Link /Border [0 0 0] /H /I /C [1 0 0] /Rect [513.621 602.778 543.27 614.733] /A << /S /GoTo /D <74686D2E322E322E332E34> >> >>}
\MusicPDFWriteObject{1978}{<< /Type /Annot /Subtype /Link /Border [0 0 0] /H /I /C [1 0 0] /Rect [400.274 484.793 431.052 497.108] /A << /S /GoTo /D <74686D2E322E322E352E34> >> >>}
\MusicPDFWriteObject{1979}{<< /Type /Annot /Subtype /Link /Border [0 0 0] /H /I /C [1 0 0] /Rect [293.333 439.363 324.111 451.067] /A << /S /GoTo /D <74686D2E352E312E322E31> >> >>}
\MusicPDFWriteObject{1980}{<< /Type /Annot /Subtype /Link /Border [0 0 0] /H /I /C [1 0 0] /Rect [211.26 426.572 241.474 438.525] /A << /S /GoTo /D <6571756174696F6E2E352E312E312E32> >> >>}
\MusicPDFWriteObject{1981}{<< /Type /Annot /Subtype /Link /Border [0 0 0] /H /I /C [1 0 0] /Rect [339.34 381.404 369.553 393.356] /A << /S /GoTo /D <6571756174696F6E2E352E312E312E32> >> >>}
\MusicPDFWriteObject{1982}{<< /Type /Annot /Subtype /Link /Border [0 0 0] /H /I /C [1 0 0] /Rect [327.949 221.909 360.791 232.7] /A << /S /GoTo /D <74686D2E442E312E322E31> >> >>}
\MusicPDFWriteObject{1983}{<< /Type /Annot /Subtype /Link /Border [0 0 0] /H /I /C [1 0 0] /Rect [404.575 95.747 434.958 107.699] /A << /S /GoTo /D <74686D2E322E332E312E34> >> >>}
\MusicPDFWriteObject{1984}{<< /Type /Annot /Subtype /Link /Border [0 0 0] /H /I /C [1 0 0] /Rect [512.815 95.747 543.198 107.699] /A << /S /GoTo /D <74686D2E322E342E332E32> >> >>}
\MusicPDFWriteObject{1985}{<< /Type /Annot /Subtype /Link /Border [0 0 0] /H /I /C [1 0 0] /Rect [333.141 82.236 362.789 94.55] /A << /S /GoTo /D <74686D2E322E332E322E32> >> >>}
\MusicPDFWriteObject{1986}{<< /Type /Annot /Subtype /Link /Border [0 0 0] /H /I /C [1 0 0] /Rect [70.755 684.273 100.545 696.588] /A << /S /GoTo /D <74686D2E322E332E322E31> >> >>}
\MusicPDFWriteObject{1987}{<< /Type /Annot /Subtype /Link /Border [0 0 0] /H /I /C [1 0 0] /Rect [513.404 667.884 543.25 685.269] /A << /S /GoTo /D <74686D2E342E332E322E32> >> >>}
\MusicPDFWriteObject{1988}{<< /Type /Annot /Subtype /Link /Border [0 0 0] /H /I /C [1 0 0] /Rect [327.949 544.889 360.791 555.679] /A << /S /GoTo /D <74686D2E442E312E322E31> >> >>}
\MusicPDFWriteObject{1989}{<< /Type /Annot /Subtype /Link /Border [0 0 0] /H /I /C [1 0 0] /Rect [70.755 371.047 101.25 383.002] /A << /S /GoTo /D <74686D2E332E322E322E31> >> >>}
\MusicPDFWriteObject{1990}{<< /Type /Annot /Subtype /Link /Border [0 0 0] /H /I /C [1 0 0] /Rect [409.968 371.047 440.181 383] /A << /S /GoTo /D <6571756174696F6E2E352E312E312E31> >> >>}
\MusicPDFWriteObject{1991}{<< /Type /Annot /Subtype /Link /Border [0 0 0] /H /I /C [0 1 0] /Rect [428.921 218.753 447.119 227.621] /A << /S /GoTo /D <636974652E4A6172> >> >>}
\MusicPDFWriteObject{1992}{<< /Type /Annot /Subtype /Link /Border [0 0 0] /H /I /C [1 0 0] /Rect [274.722 189.601 305.499 201.916] /A << /S /GoTo /D <74686D2E332E332E312E31> >> >>}
\MusicPDFWriteObject{1993}{<< /Type /Annot /Subtype /Link /Border [0 0 0] /H /I /C [1 0 0] /Rect [155.678 176.811 185.891 188.763] /A << /S /GoTo /D <6571756174696F6E2E352E312E312E32> >> >>}
\MusicPDFWriteObject{1994}{<< /Type /Annot /Subtype /Link /Border [0 0 0] /H /I /C [1 0 0] /Rect [144.52 124.208 179.713 136.16] /A << /S /GoTo /D <6571756174696F6E2E332E332E312E3332> >> >>}
\MusicPDFWriteObject{1995}{<< /Type /Annot /Subtype /Link /Border [0 0 0] /H /I /C [1 0 0] /Rect [86.995 111.057 117.209 123.01] /A << /S /GoTo /D <6571756174696F6E2E332E332E312E32> >> >>}
\MusicPDFWriteObject{1996}{<< /Type /Annot /Subtype /Link /Border [0 0 0] /H /I /C [1 0 0] /Rect [188.036 110.696 218.249 123.011] /A << /S /GoTo /D <74686D2E332E332E312E31> >> >>}
\MusicPDFWriteObject{1997}{<< /Type /Annot /Subtype /Link /Border [0 0 0] /H /I /C [1 0 0] /Rect [362.811 684.633 398.004 696.586] /A << /S /GoTo /D <6571756174696F6E2E332E332E312E3238> >> >>}
\MusicPDFWriteObject{1998}{<< /Type /Annot /Subtype /Link /Border [0 0 0] /H /I /C [1 0 0] /Rect [452.86 684.633 483.638 696.588] /A << /S /GoTo /D <74686D2E332E332E312E37> >> >>}
\MusicPDFWriteObject{1999}{<< /Type /Annot /Subtype /Link /Border [0 0 0] /H /I /C [0 1 0] /Rect [269.611 595.068 287.81 603.936] /A << /S /GoTo /D <636974652E4A6172> >> >>}
\MusicPDFWriteObject{2000}{<< /Type /Annot /Subtype /Link /Border [0 0 0] /H /I /C [0 1 0] /Rect [363.601 555.616 382.125 564.484] /A << /S /GoTo /D <636974652E4745> >> >>}
\MusicPDFWriteObject{2001}{<< /Type /Annot /Subtype /Link /Border [0 0 0] /H /I /C [0 1 0] /Rect [508.395 529.314 543.159 538.182] /A << /S /GoTo /D <636974652E436E772D62> >> >>}
\MusicPDFWriteObject{2002}{<< /Type /Annot /Subtype /Link /Border [0 0 0] /H /I /C [0 1 0] /Rect [206.644 516.163 241.408 525.031] /A << /S /GoTo /D <636974652E436E772D62> >> >>}
\MusicPDFWriteObject{2003}{<< /Type /Annot /Subtype /Link /Border [0 0 0] /H /I /C [0 1 0] /Rect [472.462 503.013 490.986 511.881] /A << /S /GoTo /D <636974652E4745> >> >>}
\MusicPDFWriteObject{2004}{<< /Type /Annot /Subtype /Link /Border [0 0 0] /H /I /C [0 1 0] /Rect [395.894 463.56 436.503 472.428] /A << /S /GoTo /D <636974652E4B6F7468652D49> >> >>}
\MusicPDFWriteObject{2005}{<< /Type /Annot /Subtype /Link /Border [0 0 0] /H /I /C [0 1 0] /Rect [105.927 450.41 125.032 459.277] /A << /S /GoTo /D <636974652E4E42> >> >>}
\MusicPDFWriteObject{2006}{<< /Type /Annot /Subtype /Link /Border [0 0 0] /H /I /C [0 1 0] /Rect [231.597 410.957 250.121 419.825] /A << /S /GoTo /D <636974652E4745> >> >>}
\MusicPDFWriteObject{2007}{<< /Type /Annot /Subtype /Link /Border [0 0 0] /H /I /C [0 1 0] /Rect [435.882 384.656 454.406 393.524] /A << /S /GoTo /D <636974652E4745> >> >>}
\MusicPDFWriteObject{2008}{<< /Type /Annot /Subtype /Link /Border [0 0 0] /H /I /C [1 0 0] /Rect [428.949 286.959 459.162 298.911] /A << /S /GoTo /D <6571756174696F6E2E352E312E312E31> >> >>}
\MusicPDFWriteObject{2009}{<< /Type /Annot /Subtype /Link /Border [0 0 0] /H /I /C [1 0 0] /Rect [447.466 184.777 477.679 196.73] /A << /S /GoTo /D <6571756174696F6E2E352E312E312E32> >> >>}
\MusicPDFWriteObject{2010}{<< /Type /Annot /Subtype /Link /Border [0 0 0] /H /I /C [1 0 0] /Rect [337.901 145.325 368.114 157.278] /A << /S /GoTo /D <6571756174696F6E2E332E332E312E33> >> >>}
\MusicPDFWriteObject{2011}{<< /Type /Annot /Subtype /Link /Border [0 0 0] /H /I /C [0 1 0] /Rect [88.585 560.986 106.058 569.854] /A << /S /GoTo /D <636974652E4346> >> >>}
\MusicPDFWriteObject{2012}{<< /Type /Annot /Subtype /Link /Border [0 0 0] /H /I /C [0 1 0] /Rect [277.734 560.986 295.206 569.854] /A << /S /GoTo /D <636974652E4346> >> >>}
\MusicPDFWriteObject{2013}{<< /Type /Annot /Subtype /Link /Border [0 0 0] /H /I /C [1 0 0] /Rect [376.953 485.638 407.166 497.591] /A << /S /GoTo /D <6571756174696F6E2E352E312E332E33> >> >>}
\MusicPDFWriteObject{2014}{<< /Type /Annot /Subtype /Link /Border [0 0 0] /H /I /C [1 0 0] /Rect [112.727 459.336 142.94 471.289] /A << /S /GoTo /D <74686D2E322E322E352E34> >> >>}
\MusicPDFWriteObject{2015}{<< /Type /Annot /Subtype /Link /Border [0 0 0] /H /I /C [1 0 0] /Rect [149.094 393.222 179.871 405.535] /A << /S /GoTo /D <74686D2E322E322E332E34> >> >>}
\MusicPDFWriteObject{2016}{<< /Type /Annot /Subtype /Link /Border [0 0 0] /H /I /C [1 0 0] /Rect [189.588 343.644 219.491 355.957] /A << /S /GoTo /D <74686D2E322E322E332E34> >> >>}
\MusicPDFWriteObject{2017}{<< /Type /Annot /Subtype /Link /Border [0 0 0] /H /I /C [1 0 0] /Rect [277.82 317.702 310.631 329.657] /A << /S /GoTo /D <74686D2E442E312E322E31> >> >>}
\MusicPDFWriteObject{2018}{<< /Type /Annot /Subtype /Link /Border [0 0 0] /H /I /C [1 0 0] /Rect [252.217 211.012 282.43 222.965] /A << /S /GoTo /D <6571756174696F6E2E352E312E332E36> >> >>}
\MusicPDFWriteObject{2019}{<< /Type /Annot /Subtype /Link /Border [0 0 0] /H /I /C [1 0 0] /Rect [398.531 195.089 428.745 207.042] /A << /S /GoTo /D <6571756174696F6E2E352E312E332E34> >> >>}
\MusicPDFWriteObject{2020}{<< /Type /Annot /Subtype /Link /Border [0 0 0] /H /I /C [1 0 0] /Rect [235.536 182.492 268.995 193.283] /A << /S /GoTo /D <74686D2E442E312E322E31> >> >>}
\MusicPDFWriteObject{2021}{<< /Type /Annot /Subtype /Link /Border [0 0 0] /H /I /C [1 0 0] /Rect [238.783 168.788 268.996 180.741] /A << /S /GoTo /D <6571756174696F6E2E352E312E332E33> >> >>}
\MusicPDFWriteObject{2022}{<< /Type /Annot /Subtype /Link /Border [0 0 0] /H /I /C [1 0 0] /Rect [174.59 129.336 204.803 141.288] /A << /S /GoTo /D <6571756174696F6E2E352E312E332E33> >> >>}
\MusicPDFWriteObject{2023}{<< /Type /Annot /Subtype /Link /Border [0 0 0] /H /I /C [1 0 0] /Rect [70.755 89.883 100.742 101.836] /A << /S /GoTo /D <74686D2E322E312E332E32> >> >>}
\MusicPDFWriteObject{2024}{<< /Type /Annot /Subtype /Link /Border [0 0 0] /H /I /C [0 1 0] /Rect [209.572 92.373 234.632 101.241] /A << /S /GoTo /D <636974652E53636861> >> >>}
\MusicPDFWriteObject{2025}{<< /Type /Annot /Subtype /Link /Border [0 0 0] /H /I /C [0 1 0] /Rect [365.373 92.373 384.478 101.241] /A << /S /GoTo /D <636974652E4E42> >> >>}
\MusicPDFWriteObject{2026}{<< /Type /Annot /Subtype /Link /Border [0 0 0] /H /I /C [0 1 0] /Rect [220.867 660.822 239.971 669.69] /A << /S /GoTo /D <636974652E4E42> >> >>}
\MusicPDFWriteObject{2027}{<< /Type /Annot /Subtype /Link /Border [0 0 0] /H /I /C [0 1 0] /Rect [382.939 660.822 402.044 669.69] /A << /S /GoTo /D <636974652E4E42> >> >>}
\MusicPDFWriteObject{2028}{<< /Type /Annot /Subtype /Link /Border [0 0 0] /H /I /C [0 1 0] /Rect [172.187 647.671 190.385 656.539] /A << /S /GoTo /D <636974652E4A6172> >> >>}
\MusicPDFWriteObject{2029}{<< /Type /Annot /Subtype /Link /Border [0 0 0] /H /I /C [0 1 0] /Rect [501.671 647.671 519.869 656.539] /A << /S /GoTo /D <636974652E4A6172> >> >>}
\MusicPDFWriteObject{2030}{<< /Type /Annot /Subtype /Link /Border [0 0 0] /H /I /C [1 0 0] /Rect [157.798 618.879 188.011 630.832] /A << /S /GoTo /D <6571756174696F6E2E352E312E332E31> >> >>}
\MusicPDFWriteObject{2031}{<< /Type /Annot /Subtype /Link /Border [0 0 0] /H /I /C [1 0 0] /Rect [337.23 545.169 370.133 557.484] /A << /S /GoTo /D <74686D2E442E312E322E31> >> >>}
\MusicPDFWriteObject{2032}{<< /Type /Annot /Subtype /Link /Border [0 0 0] /H /I /C [1 0 0] /Rect [383.879 519.228 414.092 531.18] /A << /S /GoTo /D <6571756174696F6E2E352E312E332E37> >> >>}
\MusicPDFWriteObject{2033}{<< /Type /Annot /Subtype /Link /Border [0 0 0] /H /I /C [1 0 0] /Rect [323.418 492.926 353.067 504.879] /A << /S /GoTo /D <74686D2E322E332E312E32> >> >>}
\MusicPDFWriteObject{2034}{<< /Type /Annot /Subtype /Link /Border [0 0 0] /H /I /C [1 0 0] /Rect [491.163 492.926 520.811 504.879] /A << /S /GoTo /D <74686D2E322E322E332E34> >> >>}
\MusicPDFWriteObject{2035}{<< /Type /Annot /Subtype /Link /Border [0 0 0] /H /I /C [1 0 0] /Rect [70.755 427.172 100.968 439.125] /A << /S /GoTo /D <74686D2E322E312E342E31> >> >>}
\MusicPDFWriteObject{2036}{<< /Type /Annot /Subtype /Link /Border [0 0 0] /H /I /C [0 1 0] /Rect [197.83 429.662 222.89 438.53] /A << /S /GoTo /D <636974652E53636861> >> >>}
\MusicPDFWriteObject{2037}{<< /Type /Annot /Subtype /Link /Border [0 0 0] /H /I /C [0 1 0] /Rect [354.56 429.662 373.664 438.53] /A << /S /GoTo /D <636974652E4E42> >> >>}
\MusicPDFWriteObject{2038}{<< /Type /Annot /Subtype /Link /Border [0 0 0] /H /I /C [0 1 0] /Rect [123.213 390.21 142.318 399.078] /A << /S /GoTo /D <636974652E4E42> >> >>}
\MusicPDFWriteObject{2039}{<< /Type /Annot /Subtype /Link /Border [0 0 0] /H /I /C [0 1 0] /Rect [170.61 377.059 189.715 385.927] /A << /S /GoTo /D <636974652E4E42> >> >>}
\MusicPDFWriteObject{2040}{<< /Type /Annot /Subtype /Link /Border [0 0 0] /H /I /C [1 0 0] /Rect [244.378 361.418 274.591 373.371] /A << /S /GoTo /D <6571756174696F6E2E332E332E312E32> >> >>}
\MusicPDFWriteObject{2041}{<< /Type /Annot /Subtype /Link /Border [0 0 0] /H /I /C [1 0 0] /Rect [337.313 361.418 368.09 373.371] /A << /S /GoTo /D <74686D2E332E332E312E31> >> >>}
\MusicPDFWriteObject{2042}{<< /Type /Annot /Subtype /Link /Border [0 0 0] /H /I /C [1 0 0] /Rect [166.799 335.117 197.012 347.07] /A << /S /GoTo /D <6571756174696F6E2E352E312E332E32> >> >>}
\MusicPDFWriteObject{2043}{<< /Type /Annot /Subtype /Link /Border [0 0 0] /H /I /C [1 0 0] /Rect [196.396 296.502 226.778 308.817] /A << /S /GoTo /D <74686D2E352E312E312E33> >> >>}
\MusicPDFWriteObject{2044}{<< /Type /Annot /Subtype /Link /Border [0 0 0] /H /I /C [1 0 0] /Rect [187.794 219.842 221.355 228.696] /A << /S /GoTo /D <74686D2E352E312E312E33> >> >>}
\MusicPDFWriteObject{2045}{<< /Type /Annot /Subtype /Link /Border [0 0 0] /H /I /C [1 0 0] /Rect [332.727 217.352 363.505 229.305] /A << /S /GoTo /D <74686D2E352E312E332E31> >> >>}
\MusicPDFWriteObject{2046}{<< /Type /Annot /Subtype /Link /Border [0 0 0] /H /I /C [1 0 0] /Rect [311.814 204.201 341.462 216.156] /A << /S /GoTo /D <74686D2E352E312E322E31> >> >>}
\MusicPDFWriteObject{2047}{<< /Type /Annot /Subtype /Link /Border [0 0 0] /H /I /C [1 0 0] /Rect [210.024 190.69 239.814 203.005] /A << /S /GoTo /D <74686D2E352E312E332E31> >> >>}
\MusicPDFWriteObject{2048}{<< /Type /Annot /Subtype /Link /Border [0 0 0] /H /I /C [1 0 0] /Rect [255.738 177.54 285.951 189.854] /A << /S /GoTo /D <74686D2E352E312E322E32> >> >>}
\MusicPDFWriteObject{2049}{<< /Type /Annot /Subtype /Link /Border [0 0 0] /H /I /C [1 0 0] /Rect [208.033 149.434 237.681 161.749] /A << /S /GoTo /D <74686D2E352E312E322E31> >> >>}
\MusicPDFWriteObject{2050}{<< /Type /Annot /Subtype /Link /Border [0 0 0] /H /I /C [1 0 0] /Rect [215.297 108.537 246.074 120.852] /A << /S /GoTo /D <74686D2E352E312E322E32> >> >>}
\MusicPDFWriteObject{2051}{<< /Type /Annot /Subtype /Link /Border [0 0 0] /H /I /C [1 0 0] /Rect [88.936 687.123 94.899 697.491] /A << /S /GoTo /D <48666F6F746E6F74652E3134> >> >>}
\MusicPDFWriteObject{2052}{<< /Type /Annot /Subtype /Link /Border [0 0 0] /H /I /C [1 0 0] /Rect [308.029 644.539 337.678 655.33] /A << /S /GoTo /D <74686D2E352E312E332E31> >> >>}
\MusicPDFWriteObject{2053}{<< /Type /Annot /Subtype /Link /Border [0 0 0] /H /I /C [1 0 0] /Rect [376.183 644.539 405.832 655.33] /A << /S /GoTo /D <74686D2E352E312E322E31> >> >>}
\MusicPDFWriteObject{2054}{<< /Type /Annot /Subtype /Link /Border [0 0 0] /H /I /C [1 0 0] /Rect [460.378 644.539 490.026 655.33] /A << /S /GoTo /D <74686D2E352E312E322E32> >> >>}
\MusicPDFWriteObject{2055}{<< /Type /Annot /Subtype /Link /Border [0 0 0] /H /I /C [1 0 0] /Rect [513.001 590.991 543.214 605.846] /A << /S /GoTo /D <74686D2E352E312E312E31> >> >>}
\MusicPDFWriteObject{2056}{<< /Type /Annot /Subtype /Link /Border [0 0 0] /H /I /C [1 0 0] /Rect [354.928 516.544 385.141 528.306] /A << /S /GoTo /D <74686D2E322E322E332E31> >> >>}
\MusicPDFWriteObject{2057}{<< /Type /Annot /Subtype /Link /Border [0 0 0] /H /I /C [1 0 0] /Rect [452.178 516.544 482.391 528.306] /A << /S /GoTo /D <74686D2E322E332E312E31> >> >>}
\MusicPDFWriteObject{2058}{<< /Type /Annot /Subtype /Link /Border [0 0 0] /H /I /C [1 0 0] /Rect [194.561 503.393 225.338 515.708] /A << /S /GoTo /D <74686D2E352E312E312E33> >> >>}
\MusicPDFWriteObject{2059}{<< /Type /Annot /Subtype /Link /Border [0 0 0] /H /I /C [1 0 0] /Rect [475.83 367.089 506.043 379.042] /A << /S /GoTo /D <6571756174696F6E2E322E312E342E33> >> >>}
\MusicPDFWriteObject{2060}{<< /Type /Annot /Subtype /Link /Border [0 0 0] /H /I /C [1 0 0] /Rect [508.824 188.922 539.037 200.875] /A << /S /GoTo /D <6571756174696F6E2E352E322E312E32> >> >>}
\MusicPDFWriteObject{2061}{<< /Type /Annot /Subtype /Link /Border [0 0 0] /H /I /C [1 0 0] /Rect [88.936 687.123 94.899 697.491] /A << /S /GoTo /D <48666F6F746E6F74652E3135> >> >>}
\MusicPDFWriteObject{2062}{<< /Type /Annot /Subtype /Link /Border [0 0 0] /H /I /C [1 0 0] /Rect [210.606 660.913 241.383 675.768] /A << /S /GoTo /D <74686D2E352E312E312E33> >> >>}
\MusicPDFWriteObject{2063}{<< /Type /Annot /Subtype /Link /Border [0 0 0] /H /I /C [1 0 0] /Rect [208.459 610.104 238.221 622.056] /A << /S /GoTo /D <74686D2E322E312E332E32> >> >>}
\MusicPDFWriteObject{2064}{<< /Type /Annot /Subtype /Link /Border [0 0 0] /H /I /C [0 1 0] /Rect [113.896 599.443 138.956 608.311] /A << /S /GoTo /D <636974652E53636861> >> >>}
\MusicPDFWriteObject{2065}{<< /Type /Annot /Subtype /Link /Border [0 0 0] /H /I /C [0 1 0] /Rect [162.698 586.292 187.758 595.16] /A << /S /GoTo /D <636974652E53636861> >> >>}
\MusicPDFWriteObject{2066}{<< /Type /Annot /Subtype /Link /Border [0 0 0] /H /I /C [1 0 0] /Rect [124.623 570.651 154.836 582.604] /A << /S /GoTo /D <74686D2E322E312E342E31> >> >>}
\MusicPDFWriteObject{2067}{<< /Type /Annot /Subtype /Link /Border [0 0 0] /H /I /C [0 1 0] /Rect [106.406 559.991 131.466 568.858] /A << /S /GoTo /D <636974652E53636861> >> >>}
\MusicPDFWriteObject{2068}{<< /Type /Annot /Subtype /Link /Border [0 0 0] /H /I /C [0 1 0] /Rect [351.476 559.991 376.536 568.858] /A << /S /GoTo /D <636974652E53636861> >> >>}
\MusicPDFWriteObject{2069}{<< /Type /Annot /Subtype /Link /Border [0 0 0] /H /I /C [0 1 0] /Rect [396.066 546.84 421.126 555.708] /A << /S /GoTo /D <636974652E53636861> >> >>}
\MusicPDFWriteObject{2070}{<< /Type /Annot /Subtype /Link /Border [0 0 0] /H /I /C [1 0 0] /Rect [124.623 518.048 154.836 530.001] /A << /S /GoTo /D <74686D2E322E312E332E32> >> >>}
\MusicPDFWriteObject{2071}{<< /Type /Annot /Subtype /Link /Border [0 0 0] /H /I /C [0 1 0] /Rect [264.425 520.538 289.486 529.406] /A << /S /GoTo /D <636974652E53636861> >> >>}
\MusicPDFWriteObject{2072}{<< /Type /Annot /Subtype /Link /Border [0 0 0] /H /I /C [1 0 0] /Rect [167.534 491.747 197.267 503.699] /A << /S /GoTo /D <74686D2E322E312E342E31> >> >>}
\MusicPDFWriteObject{2073}{<< /Type /Annot /Subtype /Link /Border [0 0 0] /H /I /C [0 1 0] /Rect [292.783 494.237 317.843 503.105] /A << /S /GoTo /D <636974652E53636861> >> >>}
\MusicPDFWriteObject{2074}{<< /Type /Annot /Subtype /Link /Border [0 0 0] /H /I /C [0 1 0] /Rect [431.377 494.237 456.437 503.105] /A << /S /GoTo /D <636974652E53636861> >> >>}
\MusicPDFWriteObject{2075}{<< /Type /Annot /Subtype /Link /Border [0 0 0] /H /I /C [0 1 0] /Rect [108.163 467.935 133.223 476.803] /A << /S /GoTo /D <636974652E53636861> >> >>}
\MusicPDFWriteObject{2076}{<< /Type /Annot /Subtype /Link /Border [0 0 0] /H /I /C [0 1 0] /Rect [337.202 467.935 358.728 476.803] /A << /S /GoTo /D <636974652E4D56> >> >>}
\MusicPDFWriteObject{2077}{<< /Type /Annot /Subtype /Link /Border [0 0 0] /H /I /C [0 1 0] /Rect [147.107 418.357 166.212 427.225] /A << /S /GoTo /D <636974652E4E42> >> >>}
\MusicPDFWriteObject{2078}{<< /Type /Annot /Subtype /Link /Border [0 0 0] /H /I /C [0 1 0] /Rect [297.209 418.357 316.314 427.225] /A << /S /GoTo /D <636974652E4E42> >> >>}
\MusicPDFWriteObject{2079}{<< /Type /Annot /Subtype /Link /Border [0 0 0] /H /I /C [0 1 0] /Rect [88.367 405.206 107.471 414.074] /A << /S /GoTo /D <636974652E4E42> >> >>}
\MusicPDFWriteObject{2080}{<< /Type /Annot /Subtype /Link /Border [0 0 0] /H /I /C [0 1 0] /Rect [333.449 405.206 352.553 414.074] /A << /S /GoTo /D <636974652E4E42> >> >>}
\MusicPDFWriteObject{2081}{<< /Type /Annot /Subtype /Link /Border [0 0 0] /H /I /C [0 1 0] /Rect [239.686 392.055 258.791 400.923] /A << /S /GoTo /D <636974652E4E42> >> >>}
\MusicPDFWriteObject{2082}{<< /Type /Annot /Subtype /Link /Border [0 0 0] /H /I /C [1 0 0] /Rect [254.748 361.576 284.397 375.679] /A << /S /GoTo /D <74686D2E322E312E332E32> >> >>}
\MusicPDFWriteObject{2083}{<< /Type /Annot /Subtype /Link /Border [0 0 0] /H /I /C [0 1 0] /Rect [426.165 365.754 451.225 374.622] /A << /S /GoTo /D <636974652E53636861> >> >>}
\MusicPDFWriteObject{2084}{<< /Type /Annot /Subtype /Link /Border [0 0 0] /H /I /C [1 0 0] /Rect [244.611 336.602 274.711 348.917] /A << /S /GoTo /D <74686D2E322E312E342E31> >> >>}
\MusicPDFWriteObject{2085}{<< /Type /Annot /Subtype /Link /Border [0 0 0] /H /I /C [1 0 0] /Rect [73.712 310.661 108.905 322.614] /A << /S /GoTo /D <6571756174696F6E2E322E312E342E3132> >> >>}
\MusicPDFWriteObject{2086}{<< /Type /Annot /Subtype /Link /Border [0 0 0] /H /I /C [1 0 0] /Rect [290.42 310.66 321.198 322.615] /A << /S /GoTo /D <74686D2E322E312E342E33> >> >>}
\MusicPDFWriteObject{2087}{<< /Type /Annot /Subtype /Link /Border [0 0 0] /H /I /C [1 0 0] /Rect [118.371 270.848 148.02 283.163] /A << /S /GoTo /D <74686D2E352E312E322E32> >> >>}
\MusicPDFWriteObject{2088}{<< /Type /Annot /Subtype /Link /Border [0 0 0] /H /I /C [1 0 0] /Rect [231.713 270.848 261.361 283.163] /A << /S /GoTo /D <74686D2E332E342E312E31> >> >>}
\MusicPDFWriteObject{2089}{<< /Type /Annot /Subtype /Link /Border [0 0 0] /H /I /C [1 0 0] /Rect [254.282 258.057 285.031 270.012] /A << /S /GoTo /D <74686D2E332E312E322E32> >> >>}
\MusicPDFWriteObject{2090}{<< /Type /Annot /Subtype /Link /Border [0 0 0] /H /I /C [1 0 0] /Rect [239.145 645.175 268.794 658.464] /A << /S /GoTo /D <74686D2E322E312E342E34> >> >>}
\MusicPDFWriteObject{2091}{<< /Type /Annot /Subtype /Link /Border [0 0 0] /H /I /C [1 0 0] /Rect [394.779 597.555 400.741 612.41] /A << /S /GoTo /D <48666F6F746E6F74652E3136> >> >>}
\MusicPDFWriteObject{2092}{<< /Type /Annot /Subtype /Link /Border [0 0 0] /H /I /C [1 0 0] /Rect [194.958 586.751 217.425 597.599] /A << /S /GoTo /D <73756273656374696F6E2E332E312E31> >> >>}
\MusicPDFWriteObject{2093}{<< /Type /Annot /Subtype /Link /Border [0 0 0] /H /I /C [1 0 0] /Rect [323.376 534.547 345.76 546.862] /A << /S /GoTo /D <73756273656374696F6E2E332E312E31> >> >>}
\MusicPDFWriteObject{2094}{<< /Type /Annot /Subtype /Link /Border [0 0 0] /H /I /C [1 0 0] /Rect [304.079 481.481 334.292 493.434] /A << /S /GoTo /D <6571756174696F6E2E352E322E322E31> >> >>}
\MusicPDFWriteObject{2095}{<< /Type /Annot /Subtype /Link /Border [0 0 0] /H /I /C [1 0 0] /Rect [235.101 417.041 265.314 428.994] /A << /S /GoTo /D <74686D2E352E312E312E31> >> >>}
\MusicPDFWriteObject{2096}{<< /Type /Annot /Subtype /Link /Border [0 0 0] /H /I /C [1 0 0] /Rect [339.783 417.04 369.714 428.996] /A << /S /GoTo /D <6571756174696F6E2E352E312E312E34> >> >>}
\MusicPDFWriteObject{2097}{<< /Type /Annot /Subtype /Link /Border [0 0 0] /H /I /C [1 0 0] /Rect [488.004 377.589 518.217 389.541] /A << /S /GoTo /D <6571756174696F6E2E332E332E312E33> >> >>}
\MusicPDFWriteObject{2098}{<< /Type /Annot /Subtype /Link /Border [0 0 0] /H /I /C [1 0 0] /Rect [217.098 350.927 247.424 363.242] /A << /S /GoTo /D <74686D2E352E312E332E31> >> >>}
\MusicPDFWriteObject{2099}{<< /Type /Annot /Subtype /Link /Border [0 0 0] /H /I /C [1 0 0] /Rect [237.108 684.273 267.886 695.977] /A << /S /GoTo /D <74686D2E342E332E332E31> >> >>}
\MusicPDFWriteObject{2100}{<< /Type /Annot /Subtype /Link /Border [0 0 0] /H /I /C [1 0 0] /Rect [254.355 672.036 284.568 682.827] /A << /S /GoTo /D <74686D2E352E322E332E34> >> >>}
\MusicPDFWriteObject{2101}{<< /Type /Annot /Subtype /Link /Border [0 0 0] /H /I /C [1 0 0] /Rect [445.363 651.994 451.326 665.212] /A << /S /GoTo /D <48666F6F746E6F74652E3137> >> >>}
\MusicPDFWriteObject{2102}{<< /Type /Annot /Subtype /Link /Border [0 0 0] /H /I /C [1 0 0] /Rect [475.83 590.177 506.043 602.13] /A << /S /GoTo /D <6571756174696F6E2E352E322E312E31> >> >>}
\MusicPDFWriteObject{2103}{<< /Type /Annot /Subtype /Link /Border [0 0 0] /H /I /C [1 0 0] /Rect [363.632 369.533 393.845 381.485] /A << /S /GoTo /D <6571756174696F6E2E352E322E332E32> >> >>}
\MusicPDFWriteObject{2104}{<< /Type /Annot /Subtype /Link /Border [0 0 0] /H /I /C [1 0 0] /Rect [335.663 260.627 365.876 272.58] /A << /S /GoTo /D <6571756174696F6E2E352E322E312E31> >> >>}
\MusicPDFWriteObject{2105}{<< /Type /Annot /Subtype /Link /Border [0 0 0] /H /I /C [1 0 0] /Rect [155.527 163.741 185.74 176.056] /A << /S /GoTo /D <74686D2E352E312E322E32> >> >>}
\MusicPDFWriteObject{2106}{<< /Type /Annot /Subtype /Link /Border [0 0 0] /H /I /C [1 0 0] /Rect [73.712 104.317 103.925 116.27] /A << /S /GoTo /D <6571756174696F6E2E352E322E332E34> >> >>}
\MusicPDFWriteObject{2107}{<< /Type /Annot /Subtype /Link /Border [0 0 0] /H /I /C [1 0 0] /Rect [258.266 658.332 288.479 670.284] /A << /S /GoTo /D <6571756174696F6E2E352E322E332E34> >> >>}
\MusicPDFWriteObject{2108}{<< /Type /Annot /Subtype /Link /Border [0 0 0] /H /I /C [1 0 0] /Rect [283.761 618.879 313.974 630.832] /A << /S /GoTo /D <6571756174696F6E2E352E322E312E31> >> >>}
\MusicPDFWriteObject{2109}{<< /Type /Annot /Subtype /Link /Border [0 0 0] /H /I /C [1 0 0] /Rect [123.759 605.729 153.973 617.681] /A << /S /GoTo /D <6571756174696F6E2E352E322E332E31> >> >>}
\MusicPDFWriteObject{2110}{<< /Type /Annot /Subtype /Link /Border [0 0 0] /H /I /C [1 0 0] /Rect [107.917 592.578 138.13 604.531] /A << /S /GoTo /D <6571756174696F6E2E352E322E332E34> >> >>}
\MusicPDFWriteObject{2111}{<< /Type /Annot /Subtype /Link /Border [0 0 0] /H /I /C [1 0 0] /Rect [341.067 579.427 371.28 591.38] /A << /S /GoTo /D <6571756174696F6E2E352E322E332E34> >> >>}
\MusicPDFWriteObject{2112}{<< /Type /Annot /Subtype /Link /Border [0 0 0] /H /I /C [1 0 0] /Rect [419.064 565.916 449.277 578.231] /A << /S /GoTo /D <74686D2E352E312E322E32> >> >>}
\MusicPDFWriteObject{2113}{<< /Type /Annot /Subtype /Link /Border [0 0 0] /H /I /C [1 0 0] /Rect [163.217 553.126 193.43 565.078] /A << /S /GoTo /D <6571756174696F6E2E352E322E312E31> >> >>}
\MusicPDFWriteObject{2114}{<< /Type /Annot /Subtype /Link /Border [0 0 0] /H /I /C [1 0 0] /Rect [298.722 526.824 328.935 538.777] /A << /S /GoTo /D <6571756174696F6E2E352E322E332E31> >> >>}
\MusicPDFWriteObject{2115}{<< /Type /Annot /Subtype /Link /Border [0 0 0] /H /I /C [1 0 0] /Rect [269.68 511.88 275.643 526.956] /A << /S /GoTo /D <48666F6F746E6F74652E3138> >> >>}
\MusicPDFWriteObject{2116}{<< /Type /Annot /Subtype /Link /Border [0 0 0] /H /I /C [1 0 0] /Rect [73.712 355.864 103.925 367.817] /A << /S /GoTo /D <6571756174696F6E2E322E312E342E32> >> >>}
\MusicPDFWriteObject{2117}{<< /Type /Annot /Subtype /Link /Border [0 0 0] /H /I /C [1 0 0] /Rect [330.776 355.864 360.961 367.817] /A << /S /GoTo /D <74686D2E322E312E342E31> >> >>}
\MusicPDFWriteObject{2118}{<< /Type /Annot /Subtype /Link /Border [0 0 0] /H /I /C [0 1 0] /Rect [148.737 305.036 177.497 313.904] /A << /S /GoTo /D <636974652E47726F2D66> >> >>}
\MusicPDFWriteObject{2119}{<< /Type /Annot /Subtype /Link /Border [0 0 0] /H /I /C [0 1 0] /Rect [331.047 305.036 348.367 313.904] /A << /S /GoTo /D <636974652E534E> >> >>}
\MusicPDFWriteObject{2120}{<< /Type /Annot /Subtype /Link /Border [0 0 0] /H /I /C [1 0 0] /Rect [245.188 289.396 275.402 301.348] /A << /S /GoTo /D <6571756174696F6E2E322E312E342E34> >> >>}
\MusicPDFWriteObject{2121}{<< /Type /Annot /Subtype /Link /Border [0 0 0] /H /I /C [1 0 0] /Rect [345.893 289.035 376.275 301.348] /A << /S /GoTo /D <74686D2E322E312E342E31> >> >>}
\MusicPDFWriteObject{2122}{<< /Type /Annot /Subtype /Link /Border [0 0 0] /H /I /C [1 0 0] /Rect [158.217 236.432 188.995 248.747] /A << /S /GoTo /D <74686D2E352E312E322E32> >> >>}
\MusicPDFWriteObject{2123}{<< /Type /Annot /Subtype /Link /Border [0 0 0] /H /I /C [1 0 0] /Rect [396.978 236.432 427.756 248.747] /A << /S /GoTo /D <74686D2E352E312E322E31> >> >>}
\MusicPDFWriteObject{2124}{<< /Type /Annot /Subtype /Link /Border [0 0 0] /H /I /C [1 0 0] /Rect [208.647 223.204 239.198 236.88] /A << /S /GoTo /D <74686D2E352E322E312E31> >> >>}
\MusicPDFWriteObject{2125}{<< /Type /Annot /Subtype /Link /Border [0 0 0] /H /I /C [1 0 0] /Rect [427.64 223.204 458.192 236.88] /A << /S /GoTo /D <74686D2E352E312E322E32> >> >>}
\MusicPDFWriteObject{2126}{<< /Type /Annot /Subtype /Link /Border [0 0 0] /H /I /C [1 0 0] /Rect [515.041 223.204 545.592 236.88] /A << /S /GoTo /D <74686D2E352E312E322E31> >> >>}
\MusicPDFWriteObject{2127}{<< /Type /Annot /Subtype /Link /Border [0 0 0] /H /I /C [1 0 0] /Rect [322.123 184.743 351.771 195.534] /A << /S /GoTo /D <74686D2E352E322E312E31> >> >>}
\MusicPDFWriteObject{2128}{<< /Type /Annot /Subtype /Link /Border [0 0 0] /H /I /C [1 0 0] /Rect [113.678 144.377 144.455 156.692] /A << /S /GoTo /D <74686D2E352E312E312E33> >> >>}
\MusicPDFWriteObject{2129}{<< /Type /Annot /Subtype /Link /Border [0 0 0] /H /I /C [0 1 0] /Rect [518.099 687.123 543.159 695.991] /A << /S /GoTo /D <636974652E53636861> >> >>}
\MusicPDFWriteObject{2130}{<< /Type /Annot /Subtype /Link /Border [0 0 0] /H /I /C [0 1 0] /Rect [234.724 673.973 259.784 682.84] /A << /S /GoTo /D <636974652E53636861> >> >>}
\MusicPDFWriteObject{2131}{<< /Type /Annot /Subtype /Link /Border [0 0 0] /H /I /C [0 1 0] /Rect [95.45 660.822 114.423 669.69] /A << /S /GoTo /D <636974652E547265> >> >>}
\MusicPDFWriteObject{2132}{<< /Type /Annot /Subtype /Link /Border [0 0 0] /H /I /C [0 1 0] /Rect [416.509 660.822 441.569 669.69] /A << /S /GoTo /D <636974652E53636861> >> >>}
\MusicPDFWriteObject{2133}{<< /Type /Annot /Subtype /Link /Border [0 0 0] /H /I /C [0 1 0] /Rect [214.429 634.52 239.49 643.388] /A << /S /GoTo /D <636974652E53636861> >> >>}
\MusicPDFWriteObject{2134}{<< /Type /Annot /Subtype /Link /Border [0 0 0] /H /I /C [0 1 0] /Rect [241.13 621.369 266.19 630.237] /A << /S /GoTo /D <636974652E53636861> >> >>}
\MusicPDFWriteObject{2135}{<< /Type /Annot /Subtype /Link /Border [0 0 0] /H /I /C [0 1 0] /Rect [183.755 608.219 208.816 617.087] /A << /S /GoTo /D <636974652E53636861> >> >>}
\MusicPDFWriteObject{2136}{<< /Type /Annot /Subtype /Link /Border [0 0 0] /H /I /C [0 1 0] /Rect [342.842 608.219 361.041 617.087] /A << /S /GoTo /D <636974652E4A6172> >> >>}
\MusicPDFWriteObject{2137}{<< /Type /Annot /Subtype /Link /Border [0 0 0] /H /I /C [0 1 0] /Rect [436.131 581.917 461.191 590.785] /A << /S /GoTo /D <636974652E53636861> >> >>}
\MusicPDFWriteObject{2138}{<< /Type /Annot /Subtype /Link /Border [0 0 0] /H /I /C [0 1 0] /Rect [291.936 555.616 311.041 564.484] /A << /S /GoTo /D <636974652E4E42> >> >>}
\MusicPDFWriteObject{2139}{<< /Type /Annot /Subtype /Link /Border [0 0 0] /H /I /C [0 1 0] /Rect [267.716 542.465 286.821 551.333] /A << /S /GoTo /D <636974652E4E42> >> >>}
\MusicPDFWriteObject{2140}{<< /Type /Annot /Subtype /Link /Border [0 0 0] /H /I /C [1 0 0] /Rect [170.597 526.464 200.331 538.779] /A << /S /GoTo /D <74686D2E352E312E312E33> >> >>}
\MusicPDFWriteObject{2141}{<< /Type /Annot /Subtype /Link /Border [0 0 0] /H /I /C [0 1 0] /Rect [435.644 529.314 454.749 538.182] /A << /S /GoTo /D <636974652E4E42> >> >>}
\MusicPDFWriteObject{2142}{<< /Type /Annot /Subtype /Link /Border [0 0 0] /H /I /C [0 1 0] /Rect [140.563 516.163 159.667 525.031] /A << /S /GoTo /D <636974652E4E42> >> >>}
\MusicPDFWriteObject{2143}{<< /Type /Annot /Subtype /Link /Border [0 0 0] /H /I /C [1 0 0] /Rect [196.896 487.372 227.11 499.325] /A << /S /GoTo /D <6571756174696F6E2E352E322E332E31> >> >>}
\MusicPDFWriteObject{2144}{<< /Type /Annot /Subtype /Link /Border [0 0 0] /H /I /C [1 0 0] /Rect [398.673 461.07 428.886 473.023] /A << /S /GoTo /D <6571756174696F6E2E322E312E342E35> >> >>}
\MusicPDFWriteObject{2145}{<< /Type /Annot /Subtype /Link /Border [0 0 0] /H /I /C [1 0 0] /Rect [495.892 460.71 525.541 473.023] /A << /S /GoTo /D <74686D2E322E312E342E31> >> >>}
\MusicPDFWriteObject{2146}{<< /Type /Annot /Subtype /Link /Border [0 0 0] /H /I /C [1 0 0] /Rect [282.858 447.559 312.507 461.202] /A << /S /GoTo /D <74686D2E322E312E342E33> >> >>}
\MusicPDFWriteObject{2147}{<< /Type /Annot /Subtype /Link /Border [0 0 0] /H /I /C [1 0 0] /Rect [306.2 434.408 335.849 446.113] /A << /S /GoTo /D <74686D2E322E312E342E34> >> >>}
\MusicPDFWriteObject{2148}{<< /Type /Annot /Subtype /Link /Border [0 0 0] /H /I /C [1 0 0] /Rect [73.712 421.618 103.925 433.571] /A << /S /GoTo /D <6571756174696F6E2E352E322E332E31> >> >>}
\MusicPDFWriteObject{2149}{<< /Type /Annot /Subtype /Link /Border [0 0 0] /H /I /C [0 1 0] /Rect [180.955 424.108 206.015 432.976] /A << /S /GoTo /D <636974652E53636861> >> >>}
\MusicPDFWriteObject{2150}{<< /Type /Annot /Subtype /Link /Border [0 0 0] /H /I /C [1 0 0] /Rect [88.936 387.493 94.899 397.86] /A << /S /GoTo /D <48666F6F746E6F74652E3139> >> >>}
\MusicPDFWriteObject{2151}{<< /Type /Annot /Subtype /Link /Border [0 0 0] /H /I /C [1 0 0] /Rect [153.844 259.859 183.493 270.65] /A << /S /GoTo /D <74686D2E352E322E332E32> >> >>}
\MusicPDFWriteObject{2152}{<< /Type /Annot /Subtype /Link /Border [0 0 0] /H /I /C [1 0 0] /Rect [172.543 159.35 202.756 171.303] /A << /S /GoTo /D <6571756174696F6E2E352E322E332E37> >> >>}
\MusicPDFWriteObject{2153}{<< /Type /Annot /Subtype /Link /Border [0 0 0] /H /I /C [1 0 0] /Rect [231.197 159.35 261.41 171.303] /A << /S /GoTo /D <6571756174696F6E2E352E322E332E38> >> >>}
\MusicPDFWriteObject{2154}{<< /Type /Annot /Subtype /Link /Border [0 0 0] /H /I /C [1 0 0] /Rect [261.678 146.199 291.892 158.152] /A << /S /GoTo /D <6571756174696F6E2E352E322E332E37> >> >>}
\MusicPDFWriteObject{2155}{<< /Type /Annot /Subtype /Link /Border [0 0 0] /H /I /C [1 0 0] /Rect [320.437 146.199 350.65 158.152] /A << /S /GoTo /D <6571756174696F6E2E352E322E332E38> >> >>}
\MusicPDFWriteObject{2156}{<< /Type /Annot /Subtype /Link /Border [0 0 0] /H /I /C [1 0 0] /Rect [451.772 382.538 481.986 394.49] /A << /S /GoTo /D <6571756174696F6E2E352E322E322E32> >> >>}
\MusicPDFWriteObject{2157}{<< /Type /Annot /Subtype /Link /Border [0 0 0] /H /I /C [1 0 0] /Rect [254.662 191.894 284.734 202.685] /A << /S /GoTo /D <74686D2E352E322E332E32> >> >>}
\MusicPDFWriteObject{2158}{<< /Type /Annot /Subtype /Link /Border [0 0 0] /H /I /C [0 1 0] /Rect [151.792 180.68 176.852 189.548] /A << /S /GoTo /D <636974652E53636861> >> >>}
\MusicPDFWriteObject{2159}{<< /Type /Annot /Subtype /Link /Border [0 0 0] /H /I /C [1 0 0] /Rect [128.189 157.452 134.152 171.919] /A << /S /GoTo /D <48666F6F746E6F74652E3230> >> >>}
\MusicPDFWriteObject{2160}{<< /Type /Annot /Subtype /Link /Border [0 0 0] /H /I /C [1 0 0] /Rect [515.69 117.816 545.592 132.671] /A << /S /GoTo /D <74686D2E342E332E332E31> >> >>}
\MusicPDFWriteObject{2161}{<< /Type /Annot /Subtype /Link /Border [0 0 0] /H /I /C [1 0 0] /Rect [191.384 103.957 222.161 118.812] /A << /S /GoTo /D <74686D2E352E322E312E31> >> >>}
\MusicPDFWriteObject{2162}{<< /Type /Annot /Subtype /Link /Border [0 0 0] /H /I /C [0 1 0] /Rect [278.927 687.123 303.987 695.991] /A << /S /GoTo /D <636974652E53636861> >> >>}
\MusicPDFWriteObject{2163}{<< /Type /Annot /Subtype /Link /Border [0 0 0] /H /I /C [1 0 0] /Rect [411.993 671.482 442.206 683.435] /A << /S /GoTo /D <6571756174696F6E2E352E312E312E31> >> >>}
\MusicPDFWriteObject{2164}{<< /Type /Annot /Subtype /Link /Border [0 0 0] /H /I /C [1 0 0] /Rect [73.712 658.332 103.925 670.284] /A << /S /GoTo /D <6571756174696F6E2E352E312E312E32> >> >>}
\MusicPDFWriteObject{2165}{<< /Type /Annot /Subtype /Link /Border [0 0 0] /H /I /C [1 0 0] /Rect [217.56 658.332 248.338 670.284] /A << /S /GoTo /D <74686D2E352E312E322E32> >> >>}
\MusicPDFWriteObject{2166}{<< /Type /Annot /Subtype /Link /Border [0 0 0] /H /I /C [1 0 0] /Rect [479.188 645.18 509.825 657.136] /A << /S /GoTo /D <74686D2E352E312E332E31> >> >>}
\MusicPDFWriteObject{2167}{<< /Type /Annot /Subtype /Link /Border [0 0 0] /H /I /C [1 0 0] /Rect [208.648 619.433 238.861 630.223] /A << /S /GoTo /D <74686D2E342E332E332E31> >> >>}
\MusicPDFWriteObject{2168}{<< /Type /Annot /Subtype /Link /Border [0 0 0] /H /I /C [1 0 0] /Rect [124.489 605.368 154.702 617.681] /A << /S /GoTo /D <74686D2E342E332E332E31> >> >>}
\MusicPDFWriteObject{2169}{<< /Type /Annot /Subtype /Link /Border [0 0 0] /H /I /C [1 0 0] /Rect [263.271 592.578 298.464 604.531] /A << /S /GoTo /D <6571756174696F6E2E352E322E332E3133> >> >>}
\MusicPDFWriteObject{2170}{<< /Type /Annot /Subtype /Link /Border [0 0 0] /H /I /C [1 0 0] /Rect [220.093 579.067 250.87 592.019] /A << /S /GoTo /D <74686D2E342E332E322E32> >> >>}
\MusicPDFWriteObject{2171}{<< /Type /Annot /Subtype /Link /Border [0 0 0] /H /I /C [1 0 0] /Rect [299.187 552.765 328.977 564.47] /A << /S /GoTo /D <74686D2E322E322E332E34> >> >>}
\MusicPDFWriteObject{2172}{<< /Type /Annot /Subtype /Link /Border [0 0 0] /H /I /C [1 0 0] /Rect [173.926 538.181 203.575 553.036] /A << /S /GoTo /D <74686D2E322E332E312E32> >> >>}
\MusicPDFWriteObject{2173}{<< /Type /Annot /Subtype /Link /Border [0 0 0] /H /I /C [1 0 0] /Rect [247.632 361.389 282.825 373.342] /A << /S /GoTo /D <6571756174696F6E2E352E322E332E3135> >> >>}
\MusicPDFWriteObject{2174}{<< /Type /Annot /Subtype /Link /Border [0 0 0] /H /I /C [0 1 0] /Rect [395.509 673.973 419.352 682.84] /A << /S /GoTo /D <636974652E4D6533> >> >>}
\MusicPDFWriteObject{2175}{<< /Type /Annot /Subtype /Link /Border [0 0 0] /H /I /C [0 1 0] /Rect [374.981 660.822 398.824 669.69] /A << /S /GoTo /D <636974652E4D6533> >> >>}
\MusicPDFWriteObject{2176}{<< /Type /Annot /Subtype /Link /Border [0 0 0] /H /I /C [0 1 0] /Rect [379.504 647.609 403.347 656.477] /A << /S /GoTo /D <636974652E4D6533> >> >>}
\MusicPDFWriteObject{2177}{<< /Type /Annot /Subtype /Link /Border [0 0 0] /H /I /C [0 1 0] /Rect [156.701 595.006 177.922 603.874] /A << /S /GoTo /D <636974652E427275> >> >>}
\MusicPDFWriteObject{2178}{<< /Type /Annot /Subtype /Link /Border [0 0 0] /H /I /C [0 1 0] /Rect [348.984 595.006 372.826 603.874] /A << /S /GoTo /D <636974652E4D6533> >> >>}
\MusicPDFWriteObject{2179}{<< /Type /Annot /Subtype /Link /Border [0 0 0] /H /I /C [0 1 0] /Rect [219.703 489.8 243.546 498.668] /A << /S /GoTo /D <636974652E4D6531> >> >>}
\MusicPDFWriteObject{2180}{<< /Type /Annot /Subtype /Link /Border [0 0 0] /H /I /C [0 1 0] /Rect [175.217 462.79 195.138 471.658] /A << /S /GoTo /D <636974652E484E> >> >>}
\MusicPDFWriteObject{2181}{<< /Type /Annot /Subtype /Link /Border [0 0 0] /H /I /C [0 1 0] /Rect [244.87 462.79 268.713 471.658] /A << /S /GoTo /D <636974652E4D6532> >> >>}
\MusicPDFWriteObject{2182}{<< /Type /Annot /Subtype /Link /Border [0 0 0] /H /I /C [0 1 0] /Rect [409.319 449.639 433.162 458.507] /A << /S /GoTo /D <636974652E4D6533> >> >>}
\MusicPDFWriteObject{2183}{<< /Type /Annot /Subtype /Link /Border [0 0 0] /H /I /C [1 0 0] /Rect [284.816 342.257 314.465 358.286] /A << /S /GoTo /D <74686D2E352E322E332E32> >> >>}
\MusicPDFWriteObject{2184}{<< /Type /Annot /Subtype /Link /Border [0 0 0] /H /I /C [1 0 0] /Rect [198.215 381.11 227.864 393.425] /A << /S /GoTo /D <74686D2E352E312E322E32> >> >>}
\MusicPDFWriteObject{2185}{<< /Type /Annot /Subtype /Link /Border [0 0 0] /H /I /C [1 0 0] /Rect [88.849 321.273 119.062 333.226] /A << /S /GoTo /D <6571756174696F6E2E332E332E312E32> >> >>}
\MusicPDFWriteObject{2186}{<< /Type /Annot /Subtype /Link /Border [0 0 0] /H /I /C [1 0 0] /Rect [189.89 321.273 220.103 333.226] /A << /S /GoTo /D <74686D2E332E332E312E31> >> >>}
\MusicPDFWriteObject{2187}{<< /Type /Annot /Subtype /Link /Border [0 0 0] /H /I /C [1 0 0] /Rect [277.773 266.584 307.986 278.537] /A << /S /GoTo /D <6571756174696F6E2E352E322E342E32> >> >>}
\MusicPDFWriteObject{2188}{<< /Type /Annot /Subtype /Link /Border [0 0 0] /H /I /C [1 0 0] /Rect [160.467 226.246 190.116 238.201] /A << /S /GoTo /D <74686D2E352E312E322E31> >> >>}
\MusicPDFWriteObject{2189}{<< /Type /Annot /Subtype /Link /Border [0 0 0] /H /I /C [1 0 0] /Rect [88.849 174.971 119.062 186.924] /A << /S /GoTo /D <6571756174696F6E2E332E322E322E32> >> >>}
\MusicPDFWriteObject{2190}{<< /Type /Annot /Subtype /Link /Border [0 0 0] /H /I /C [1 0 0] /Rect [189.89 174.971 220.103 186.924] /A << /S /GoTo /D <74686D2E332E322E322E31> >> >>}
\MusicPDFWriteObject{2191}{<< /Type /Annot /Subtype /Link /Border [0 0 0] /H /I /C [1 0 0] /Rect [147.67 122.048 177.883 134.001] /A << /S /GoTo /D <6571756174696F6E2E352E322E342E34> >> >>}
\MusicPDFWriteObject{2192}{<< /Type /Annot /Subtype /Link /Border [0 0 0] /H /I /C [1 0 0] /Rect [247.825 82.596 278.038 94.549] /A << /S /GoTo /D <6571756174696F6E2E332E332E312E32> >> >>}
\MusicPDFWriteObject{2193}{<< /Type /Annot /Subtype /Link /Border [0 0 0] /H /I /C [1 0 0] /Rect [348.866 82.596 379.079 94.549] /A << /S /GoTo /D <74686D2E332E332E312E31> >> >>}
\MusicPDFWriteObject{2194}{<< /Type /Annot /Subtype /Link /Border [0 0 0] /H /I /C [1 0 0] /Rect [312.722 483.686 345.757 495.641] /A << /S /GoTo /D <74686D2E432E322E352E31> >> >>}
\MusicPDFWriteObject{2195}{<< /Type /Annot /Subtype /Link /Border [0 0 0] /H /I /C [1 0 0] /Rect [179.798 468.849 210.124 482.508] /A << /S /GoTo /D <74686D2E322E312E322E34> >> >>}
\MusicPDFWriteObject{2196}{<< /Type /Annot /Subtype /Link /Border [0 0 0] /H /I /C [1 0 0] /Rect [239.498 252.779 269.711 264.732] /A << /S /GoTo /D <6571756174696F6E2E352E332E312E31> >> >>}
\MusicPDFWriteObject{2197}{<< /Type /Annot /Subtype /Link /Border [0 0 0] /H /I /C [1 0 0] /Rect [199.845 239.628 230.51 252.044] /A << /S /GoTo /D <74686D2E332E342E332E33> >> >>}
\MusicPDFWriteObject{2198}{<< /Type /Annot /Subtype /Link /Border [0 0 0] /H /I /C [1 0 0] /Rect [384.872 224.79 415.085 238.45] /A << /S /GoTo /D <74686D2E332E342E332E33> >> >>}
\MusicPDFWriteObject{2199}{<< /Type /Annot /Subtype /Link /Border [0 0 0] /H /I /C [1 0 0] /Rect [195.724 134.113 226.501 154.342] /A << /S /GoTo /D <74686D2E332E342E332E33> >> >>}
\MusicPDFWriteObject{2200}{<< /Type /Annot /Subtype /Link /Border [0 0 0] /H /I /C [1 0 0] /Rect [125.679 682.946 156.457 696.606] /A << /S /GoTo /D <74686D2E352E332E312E31> >> >>}
\MusicPDFWriteObject{2201}{<< /Type /Annot /Subtype /Link /Border [0 0 0] /H /I /C [1 0 0] /Rect [124.905 631.607 155.485 643.56] /A << /S /GoTo /D <74686D2E352E332E312E31> >> >>}
\MusicPDFWriteObject{2202}{<< /Type /Annot /Subtype /Link /Border [0 0 0] /H /I /C [1 0 0] /Rect [124.623 605.305 154.836 617.258] /A << /S /GoTo /D <74686D2E322E312E322E34> >> >>}
\MusicPDFWriteObject{2203}{<< /Type /Annot /Subtype /Link /Border [0 0 0] /H /I /C [1 0 0] /Rect [294.977 366.133 329.506 378.089] /A << /S /GoTo /D <74686D2E322E312E322E3130> >> >>}
\MusicPDFWriteObject{2204}{<< /Type /Annot /Subtype /Link /Border [0 0 0] /H /I /C [1 0 0] /Rect [451.172 366.133 480.821 378.089] /A << /S /GoTo /D <74686D2E342E332E332E31> >> >>}
\MusicPDFWriteObject{2205}{<< /Type /Annot /Subtype /Link /Border [0 0 0] /H /I /C [0 1 0] /Rect [129.658 245.95 154.718 254.818] /A << /S /GoTo /D <636974652E53636861> >> >>}
\MusicPDFWriteObject{2206}{<< /Type /Annot /Subtype /Link /Border [0 0 0] /H /I /C [1 0 0] /Rect [372.956 176.272 402.605 190.768] /A << /S /GoTo /D <74686D2E352E332E312E31> >> >>}
\MusicPDFWriteObject{2207}{<< /Type /Annot /Subtype /Link /Border [0 0 0] /H /I /C [1 0 0] /Rect [147.67 637.72 177.883 649.672] /A << /S /GoTo /D <6571756174696F6E2E322E312E342E31> >> >>}
\MusicPDFWriteObject{2208}{<< /Type /Annot /Subtype /Link /Border [0 0 0] /H /I /C [1 0 0] /Rect [260.445 636.413 290.743 649.672] /A << /S /GoTo /D <74686D2E322E312E342E31> >> >>}
\MusicPDFWriteObject{2209}{<< /Type /Annot /Subtype /Link /Border [0 0 0] /H /I /C [1 0 0] /Rect [393.358 540.804 423.571 552.757] /A << /S /GoTo /D <6571756174696F6E2E352E332E312E36> >> >>}
\MusicPDFWriteObject{2210}{<< /Type /Annot /Subtype /Link /Border [0 0 0] /H /I /C [1 0 0] /Rect [451.182 540.804 481.395 552.757] /A << /S /GoTo /D <6571756174696F6E2E352E332E312E38> >> >>}
\MusicPDFWriteObject{2211}{<< /Type /Annot /Subtype /Link /Border [0 0 0] /H /I /C [1 0 0] /Rect [252.139 527.653 282.352 539.606] /A << /S /GoTo /D <6571756174696F6E2E322E312E342E32> >> >>}
\MusicPDFWriteObject{2212}{<< /Type /Annot /Subtype /Link /Border [0 0 0] /H /I /C [1 0 0] /Rect [352.516 527.293 382.729 539.606] /A << /S /GoTo /D <74686D2E322E312E342E31> >> >>}
\MusicPDFWriteObject{2213}{<< /Type /Annot /Subtype /Link /Border [0 0 0] /H /I /C [1 0 0] /Rect [329.597 514.142 359.5 526.457] /A << /S /GoTo /D <74686D2E322E312E342E31> >> >>}
\MusicPDFWriteObject{2214}{<< /Type /Annot /Subtype /Link /Border [0 0 0] /H /I /C [1 0 0] /Rect [482.922 514.142 512.825 526.457] /A << /S /GoTo /D <74686D2E322E312E342E31> >> >>}
\MusicPDFWriteObject{2215}{<< /Type /Annot /Subtype /Link /Border [0 0 0] /H /I /C [0 1 0] /Rect [73.715 503.842 91.914 512.709] /A << /S /GoTo /D <636974652E4A6172> >> >>}
\MusicPDFWriteObject{2216}{<< /Type /Annot /Subtype /Link /Border [0 0 0] /H /I /C [1 0 0] /Rect [244.067 488.201 273.716 500.154] /A << /S /GoTo /D <74686D2E352E332E312E33> >> >>}
\MusicPDFWriteObject{2217}{<< /Type /Annot /Subtype /Link /Border [0 0 0] /H /I /C [0 1 0] /Rect [295.933 490.691 314.131 499.559] /A << /S /GoTo /D <636974652E4A6172> >> >>}
\MusicPDFWriteObject{2218}{<< /Type /Annot /Subtype /Link /Border [0 0 0] /H /I /C [1 0 0] /Rect [71.004 477.54 101.217 485.924] /A << /S /GoTo /D <74686D2E352E332E312E33> >> >>}
\MusicPDFWriteObject{2219}{<< /Type /Annot /Subtype /Link /Border [0 0 0] /H /I /C [1 0 0] /Rect [255.999 448.749 286.212 460.701] /A << /S /GoTo /D <6571756174696F6E2E352E322E312E31> >> >>}
\MusicPDFWriteObject{2220}{<< /Type /Annot /Subtype /Link /Border [0 0 0] /H /I /C [1 0 0] /Rect [446.364 329.27 476.577 341.223] /A << /S /GoTo /D <6571756174696F6E2E352E332E312E39> >> >>}
\MusicPDFWriteObject{2221}{<< /Type /Annot /Subtype /Link /Border [0 0 0] /H /I /C [1 0 0] /Rect [506.695 329.27 541.888 341.223] /A << /S /GoTo /D <6571756174696F6E2E352E332E312E3130> >> >>}
\MusicPDFWriteObject{2222}{<< /Type /Annot /Subtype /Link /Border [0 0 0] /H /I /C [1 0 0] /Rect [124.574 290.371 156.093 301.162] /A << /S /GoTo /D <74686D2E352E322E312E31> >> >>}
\MusicPDFWriteObject{2223}{<< /Type /Annot /Subtype /Link /Border [0 0 0] /H /I /C [1 0 0] /Rect [205.99 174.651 236.203 186.604] /A << /S /GoTo /D <6571756174696F6E2E352E332E312E36> >> >>}
\MusicPDFWriteObject{2224}{<< /Type /Annot /Subtype /Link /Border [0 0 0] /H /I /C [1 0 0] /Rect [354.251 174.651 384.464 186.604] /A << /S /GoTo /D <6571756174696F6E2E352E332E312E38> >> >>}
\MusicPDFWriteObject{2225}{<< /Type /Annot /Subtype /Link /Border [0 0 0] /H /I /C [1 0 0] /Rect [194.605 161.5 224.818 173.453] /A << /S /GoTo /D <6571756174696F6E2E352E332E312E39> >> >>}
\MusicPDFWriteObject{2226}{<< /Type /Annot /Subtype /Link /Border [0 0 0] /H /I /C [1 0 0] /Rect [303.793 161.5 338.987 173.453] /A << /S /GoTo /D <6571756174696F6E2E352E332E312E3130> >> >>}
\MusicPDFWriteObject{2227}{<< /Type /Annot /Subtype /Link /Border [0 0 0] /H /I /C [1 0 0] /Rect [161.552 135.199 191.765 147.152] /A << /S /GoTo /D <6571756174696F6E2E352E322E312E31> >> >>}
\MusicPDFWriteObject{2228}{<< /Type /Annot /Subtype /Link /Border [0 0 0] /H /I /C [1 0 0] /Rect [418.477 135.199 449.254 147.152] /A << /S /GoTo /D <74686D2E352E322E312E31> >> >>}
\MusicPDFWriteObject{2229}{<< /Type /Annot /Subtype /Link /Border [0 0 0] /H /I /C [1 0 0] /Rect [382.956 83.149 413.169 93.94] /A << /S /GoTo /D <74686D2E352E332E312E33> >> >>}
\MusicPDFWriteObject{2230}{<< /Type /Annot /Subtype /Link /Border [0 0 0] /H /I /C [0 1 0] /Rect [263.86 621.369 282.059 630.237] /A << /S /GoTo /D <636974652E4A6172> >> >>}
\MusicPDFWriteObject{2231}{<< /Type /Annot /Subtype /Link /Border [0 0 0] /H /I /C [0 1 0] /Rect [239.227 608.219 279.836 617.087] /A << /S /GoTo /D <636974652E4B6F7468652D49> >> >>}
\MusicPDFWriteObject{2232}{<< /Type /Annot /Subtype /Link /Border [0 0 0] /H /I /C [0 1 0] /Rect [284.867 608.219 330.138 617.087] /A << /S /GoTo /D <636974652E4B6F7468652D4949> >> >>}
\MusicPDFWriteObject{2233}{<< /Type /Annot /Subtype /Link /Border [0 0 0] /H /I /C [0 1 0] /Rect [357.148 608.219 397.757 617.087] /A << /S /GoTo /D <636974652E4B6F7468652D49> >> >>}
\MusicPDFWriteObject{2234}{<< /Type /Annot /Subtype /Link /Border [0 0 0] /H /I /C [0 1 0] /Rect [518.132 608.219 543.192 617.087] /A << /S /GoTo /D <636974652E53636861> >> >>}
\MusicPDFWriteObject{2235}{<< /Type /Annot /Subtype /Link /Border [0 0 0] /H /I /C [0 1 0] /Rect [192.748 515.649 223.694 524.517] /A << /S /GoTo /D <636974652E486F726D> >> >>}
\MusicPDFWriteObject{2236}{<< /Type /Annot /Subtype /Link /Border [0 0 0] /H /I /C [1 0 0] /Rect [511.045 498.321 543.27 511.98] /A << /S /GoTo /D <74686D2E442E312E382E32> >> >>}
\MusicPDFWriteObject{2237}{<< /Type /Annot /Subtype /Link /Border [0 0 0] /H /I /C [0 1 0] /Rect [213.601 130.839 231.799 139.707] /A << /S /GoTo /D <636974652E4A6172> >> >>}
\MusicPDFWriteObject{2238}{<< /Type /Annot /Subtype /Link /Border [0 0 0] /H /I /C [1 0 0] /Rect [142.196 115.198 172.409 128.14] /A << /S /GoTo /D <74686D2E342E322E312E33> >> >>}
\MusicPDFWriteObject{2239}{<< /Type /Annot /Subtype /Link /Border [0 0 0] /H /I /C [0 1 0] /Rect [148.531 98.237 179.477 107.105] /A << /S /GoTo /D <636974652E486F726D> >> >>}
\MusicPDFWriteObject{2240}{<< /Type /Annot /Subtype /Link /Border [0 0 0] /H /I /C [0 1 0] /Rect [161.09 85.086 192.036 93.954] /A << /S /GoTo /D <636974652E486F726D> >> >>}
\MusicPDFWriteObject{2241}{<< /Type /Annot /Subtype /Link /Border [0 0 0] /H /I /C [0 1 0] /Rect [512.213 85.086 543.159 93.954] /A << /S /GoTo /D <636974652E486F726D> >> >>}
\MusicPDFWriteObject{2242}{<< /Type /Annot /Subtype /Link /Border [0 0 0] /H /I /C [0 1 0] /Rect [318.537 660.822 349.483 669.69] /A << /S /GoTo /D <636974652E486F726D> >> >>}
\MusicPDFWriteObject{2243}{<< /Type /Annot /Subtype /Link /Border [0 0 0] /H /I /C [0 1 0] /Rect [445.951 608.219 476.897 617.087] /A << /S /GoTo /D <636974652E486F726D> >> >>}
\MusicPDFWriteObject{2244}{<< /Type /Annot /Subtype /Link /Border [0 0 0] /H /I /C [0 1 0] /Rect [144.488 592.64 175.434 601.507] /A << /S /GoTo /D <636974652E486F726D> >> >>}
\MusicPDFWriteObject{2245}{<< /Type /Annot /Subtype /Link /Border [0 0 0] /H /I /C [1 0 0] /Rect [73.712 523.81 103.925 535.762] /A << /S /GoTo /D <6571756174696F6E2E352E332E322E31> >> >>}
\MusicPDFWriteObject{2246}{<< /Type /Annot /Subtype /Link /Border [0 0 0] /H /I /C [1 0 0] /Rect [369.779 523.81 399.992 535.762] /A << /S /GoTo /D <6571756174696F6E2E352E332E322E31> >> >>}
\MusicPDFWriteObject{2247}{<< /Type /Annot /Subtype /Link /Border [0 0 0] /H /I /C [0 1 0] /Rect [140.172 459.804 159.277 468.672] /A << /S /GoTo /D <636974652E4E42> >> >>}
\MusicPDFWriteObject{2248}{<< /Type /Annot /Subtype /Link /Border [0 0 0] /H /I /C [1 0 0] /Rect [423.03 404.711 453.215 416.666] /A << /S /GoTo /D <74686D2E322E312E312E36> >> >>}
\MusicPDFWriteObject{2249}{<< /Type /Annot /Subtype /Link /Border [0 0 0] /H /I /C [0 1 0] /Rect [518.157 407.201 543.217 416.069] /A << /S /GoTo /D <636974652E53636861> >> >>}
\MusicPDFWriteObject{2250}{<< /Type /Annot /Subtype /Link /Border [0 0 0] /H /I /C [0 1 0] /Rect [246.618 394.051 271.678 402.918] /A << /S /GoTo /D <636974652E53636861> >> >>}
\MusicPDFWriteObject{2251}{<< /Type /Annot /Subtype /Link /Border [0 0 0] /H /I /C [0 1 0] /Rect [312.289 380.9 331.262 389.768] /A << /S /GoTo /D <636974652E547265> >> >>}
\MusicPDFWriteObject{2252}{<< /Type /Annot /Subtype /Link /Border [0 0 0] /H /I /C [0 1 0] /Rect [132.785 367.749 150.983 376.617] /A << /S /GoTo /D <636974652E4A6172> >> >>}
\MusicPDFWriteObject{2253}{<< /Type /Annot /Subtype /Link /Border [0 0 0] /H /I /C [0 1 0] /Rect [279.545 288.844 307.379 297.712] /A << /S /GoTo /D <636974652E56616C2D61> >> >>}
\MusicPDFWriteObject{2254}{<< /Type /Annot /Subtype /Link /Border [0 0 0] /H /I /C [0 1 0] /Rect [166.737 249.392 184.935 258.26] /A << /S /GoTo /D <636974652E4A6172> >> >>}
\MusicPDFWriteObject{2255}{<< /Type /Annot /Subtype /Link /Border [0 0 0] /H /I /C [0 1 0] /Rect [256.596 221.83 281.656 230.698] /A << /S /GoTo /D <636974652E53636861> >> >>}
\MusicPDFWriteObject{2256}{<< /Type /Annot /Subtype /Link /Border [0 0 0] /H /I /C [0 1 0] /Rect [302.916 168.231 321.114 177.099] /A << /S /GoTo /D <636974652E4A6172> >> >>}
\MusicPDFWriteObject{2257}{<< /Type /Annot /Subtype /Link /Border [0 0 0] /H /I /C [0 1 0] /Rect [88.148 140.973 113.208 149.841] /A << /S /GoTo /D <636974652E53636861> >> >>}
\MusicPDFWriteObject{2258}{<< /Type /Annot /Subtype /Link /Border [0 0 0] /H /I /C [0 1 0] /Rect [133.482 113.104 151.681 121.972] /A << /S /GoTo /D <636974652E4A6172> >> >>}
\MusicPDFWriteObject{2259}{<< /Type /Annot /Subtype /Link /Border [0 0 0] /H /I /C [0 1 0] /Rect [465.533 86.803 482.853 95.67] /A << /S /GoTo /D <636974652E534E> >> >>}
\MusicPDFWriteObject{2260}{<< /Type /Annot /Subtype /Link /Border [0 0 0] /H /I /C [0 1 0] /Rect [132.124 647.19 151.097 656.058] /A << /S /GoTo /D <636974652E547265> >> >>}
\MusicPDFWriteObject{2261}{<< /Type /Annot /Subtype /Link /Border [0 0 0] /H /I /C [0 1 0] /Rect [343.779 647.19 362.752 656.058] /A << /S /GoTo /D <636974652E547265> >> >>}
\MusicPDFWriteObject{2262}{<< /Type /Annot /Subtype /Link /Border [0 0 0] /H /I /C [0 1 0] /Rect [524.271 647.19 543.245 656.058] /A << /S /GoTo /D <636974652E547265> >> >>}
\MusicPDFWriteObject{2263}{<< /Type /Annot /Subtype /Link /Border [0 0 0] /H /I /C [0 1 0] /Rect [73.77 620.888 92.875 629.756] /A << /S /GoTo /D <636974652E4E42> >> >>}
\MusicPDFWriteObject{2264}{<< /Type /Annot /Subtype /Link /Border [0 0 0] /H /I /C [0 1 0] /Rect [174.438 620.888 193.542 629.756] /A << /S /GoTo /D <636974652E4E42> >> >>}
\MusicPDFWriteObject{2265}{<< /Type /Annot /Subtype /Link /Border [0 0 0] /H /I /C [1 0 0] /Rect [91.42 341.266 121.633 353.219] /A << /S /GoTo /D <6571756174696F6E2E352E332E322E34> >> >>}
\MusicPDFWriteObject{2266}{<< /Type /Annot /Subtype /Link /Border [0 0 0] /H /I /C [1 0 0] /Rect [260.578 242.884 290.791 254.837] /A << /S /GoTo /D <6571756174696F6E2E352E332E322E36> >> >>}
\MusicPDFWriteObject{2267}{<< /Type /Annot /Subtype /Link /Border [0 0 0] /H /I /C [1 0 0] /Rect [416.444 242.884 446.657 254.837] /A << /S /GoTo /D <6571756174696F6E2E352E332E322E36> >> >>}
\MusicPDFWriteObject{2268}{<< /Type /Annot /Subtype /Link /Border [0 0 0] /H /I /C [1 0 0] /Rect [147.808 82.596 185.63 94.549] /A << /S /GoTo /D <6571756174696F6E2E442E312E31302E36> >> >>}
\MusicPDFWriteObject{2269}{<< /Type /Annot /Subtype /Link /Border [0 0 0] /H /I /C [1 0 0] /Rect [193.803 645.181 224.017 657.134] /A << /S /GoTo /D <6571756174696F6E2E352E332E322E37> >> >>}
\MusicPDFWriteObject{2270}{<< /Type /Annot /Subtype /Link /Border [0 0 0] /H /I /C [0 1 0] /Rect [146.536 621.369 177.483 630.237] /A << /S /GoTo /D <636974652E486F726D> >> >>}
\MusicPDFWriteObject{2271}{<< /Type /Annot /Subtype /Link /Border [0 0 0] /H /I /C [1 0 0] /Rect [200.67 533.454 230.884 545.407] /A << /S /GoTo /D <6571756174696F6E2E352E332E322E38> >> >>}
\MusicPDFWriteObject{2272}{<< /Type /Annot /Subtype /Link /Border [0 0 0] /H /I /C [1 0 0] /Rect [263.325 533.454 293.538 545.407] /A << /S /GoTo /D <6571756174696F6E2E352E332E322E34> >> >>}
\MusicPDFWriteObject{2273}{<< /Type /Annot /Subtype /Link /Border [0 0 0] /H /I /C [1 0 0] /Rect [400.757 533.454 430.97 545.407] /A << /S /GoTo /D <6571756174696F6E2E352E332E322E35> >> >>}
\MusicPDFWriteObject{2274}{<< /Type /Annot /Subtype /Link /Border [0 0 0] /H /I /C [1 0 0] /Rect [460.025 533.454 490.238 545.407] /A << /S /GoTo /D <6571756174696F6E2E352E332E322E36> >> >>}
\MusicPDFWriteObject{2275}{<< /Type /Annot /Subtype /Link /Border [0 0 0] /H /I /C [0 1 0] /Rect [325.133 522.793 344.106 531.661] /A << /S /GoTo /D <636974652E547265> >> >>}
\MusicPDFWriteObject{2276}{<< /Type /Annot /Subtype /Link /Border [0 0 0] /H /I /C [0 1 0] /Rect [178.768 483.341 197.872 492.209] /A << /S /GoTo /D <636974652E4E42> >> >>}
\MusicPDFWriteObject{2277}{<< /Type /Annot /Subtype /Link /Border [0 0 0] /H /I /C [1 0 0] /Rect [343.402 480.851 373.615 492.804] /A << /S /GoTo /D <6571756174696F6E2E352E332E322E36> >> >>}
\MusicPDFWriteObject{2278}{<< /Type /Annot /Subtype /Link /Border [0 0 0] /H /I /C [0 1 0] /Rect [499.258 447.189 518.231 456.057] /A << /S /GoTo /D <636974652E547265> >> >>}
\MusicPDFWriteObject{2279}{<< /Type /Annot /Subtype /Link /Border [0 0 0] /H /I /C [1 0 0] /Rect [247.502 605.702 277.715 618.295] /A << /S /GoTo /D <74686D2E352E332E322E33> >> >>}
\MusicPDFWriteObject{2280}{<< /Type /Annot /Subtype /Link /Border [0 0 0] /H /I /C [0 1 0] /Rect [500.754 555.589 519.728 564.457] /A << /S /GoTo /D <636974652E547265> >> >>}
\MusicPDFWriteObject{2281}{<< /Type /Annot /Subtype /Link /Border [0 0 0] /H /I /C [1 0 0] /Rect [465.253 385.031 500.447 396.984] /A << /S /GoTo /D <6571756174696F6E2E352E332E322E3134> >> >>}
\MusicPDFWriteObject{2282}{<< /Type /Annot /Subtype /Link /Border [0 0 0] /H /I /C [1 0 0] /Rect [257.545 357.595 292.738 369.548] /A << /S /GoTo /D <6571756174696F6E2E352E332E322E3134> >> >>}
\MusicPDFWriteObject{2283}{<< /Type /Annot /Subtype /Link /Border [0 0 0] /H /I /C [0 1 0] /Rect [512.213 360.085 543.159 368.953] /A << /S /GoTo /D <636974652E486F726D> >> >>}
\MusicPDFWriteObject{2284}{<< /Type /Annot /Subtype /Link /Border [0 0 0] /H /I /C [1 0 0] /Rect [148.539 329.606 179.316 343.266] /A << /S /GoTo /D <74686D2E352E332E322E33> >> >>}
\MusicPDFWriteObject{2285}{<< /Type /Annot /Subtype /Link /Border [0 0 0] /H /I /C [1 0 0] /Rect [272.619 684.633 307.812 696.586] /A << /S /GoTo /D <6571756174696F6E2E352E332E322E3135> >> >>}
\MusicPDFWriteObject{2286}{<< /Type /Annot /Subtype /Link /Border [0 0 0] /H /I /C [1 0 0] /Rect [213.239 609.319 248.432 621.272] /A << /S /GoTo /D <6571756174696F6E2E352E332E322E3136> >> >>}
\MusicPDFWriteObject{2287}{<< /Type /Annot /Subtype /Link /Border [0 0 0] /H /I /C [1 0 0] /Rect [273.858 555.939 309.051 567.892] /A << /S /GoTo /D <6571756174696F6E2E352E332E322E3135> >> >>}
\MusicPDFWriteObject{2288}{<< /Type /Annot /Subtype /Link /Border [0 0 0] /H /I /C [1 0 0] /Rect [373.087 554.887 403.385 567.892] /A << /S /GoTo /D <74686D2E352E332E322E33> >> >>}
\MusicPDFWriteObject{2289}{<< /Type /Annot /Subtype /Link /Border [0 0 0] /H /I /C [1 0 0] /Rect [244.182 394.152 279.375 406.105] /A << /S /GoTo /D <6571756174696F6E2E352E332E322E3137> >> >>}
\MusicPDFWriteObject{2290}{<< /Type /Annot /Subtype /Link /Border [0 0 0] /H /I /C [1 0 0] /Rect [226.725 381.002 261.918 392.954] /A << /S /GoTo /D <6571756174696F6E2E352E332E322E3131> >> >>}
\MusicPDFWriteObject{2291}{<< /Type /Annot /Subtype /Link /Border [0 0 0] /H /I /C [1 0 0] /Rect [73.712 315.248 108.905 327.201] /A << /S /GoTo /D <6571756174696F6E2E352E332E322E3131> >> >>}
\MusicPDFWriteObject{2292}{<< /Type /Annot /Subtype /Link /Border [0 0 0] /H /I /C [1 0 0] /Rect [374.759 315.248 409.952 327.201] /A << /S /GoTo /D <6571756174696F6E2E352E332E322E3131> >> >>}
\MusicPDFWriteObject{2293}{<< /Type /Annot /Subtype /Link /Border [0 0 0] /H /I /C [0 1 0] /Rect [173.978 238.833 193.083 247.701] /A << /S /GoTo /D <636974652E4E42> >> >>}
\MusicPDFWriteObject{2294}{<< /Type /Annot /Subtype /Link /Border [0 0 0] /H /I /C [1 0 0] /Rect [422.881 210.041 453.122 221.996] /A << /S /GoTo /D <74686D2E322E312E312E36> >> >>}
\MusicPDFWriteObject{2295}{<< /Type /Annot /Subtype /Link /Border [0 0 0] /H /I /C [0 1 0] /Rect [518.151 212.532 543.211 221.4] /A << /S /GoTo /D <636974652E53636861> >> >>}
\MusicPDFWriteObject{2296}{<< /Type /Annot /Subtype /Link /Border [0 0 0] /H /I /C [0 1 0] /Rect [246.618 199.381 271.678 208.249] /A << /S /GoTo /D <636974652E53636861> >> >>}
\MusicPDFWriteObject{2297}{<< /Type /Annot /Subtype /Link /Border [0 0 0] /H /I /C [0 1 0] /Rect [312.289 186.23 331.262 195.098] /A << /S /GoTo /D <636974652E547265> >> >>}
\MusicPDFWriteObject{2298}{<< /Type /Annot /Subtype /Link /Border [0 0 0] /H /I /C [0 1 0] /Rect [135.173 173.08 153.371 181.947] /A << /S /GoTo /D <636974652E4A6172> >> >>}
\MusicPDFWriteObject{2299}{<< /Type /Annot /Subtype /Link /Border [0 0 0] /H /I /C [0 1 0] /Rect [262.24 133.627 287.3 142.495] /A << /S /GoTo /D <636974652E53636861> >> >>}
\MusicPDFWriteObject{2300}{<< /Type /Annot /Subtype /Link /Border [0 0 0] /H /I /C [1 0 0] /Rect [339.922 117.986 370.699 129.939] /A << /S /GoTo /D <74686D2E352E332E322E33> >> >>}
\MusicPDFWriteObject{2301}{<< /Type /Annot /Subtype /Link /Border [0 0 0] /H /I /C [0 1 0] /Rect [137.128 107.326 164.962 116.194] /A << /S /GoTo /D <636974652E56616C2D61> >> >>}
\MusicPDFWriteObject{2302}{<< /Type /Annot /Subtype /Link /Border [0 0 0] /H /I /C [0 1 0] /Rect [459.443 687.123 488.203 695.991] /A << /S /GoTo /D <636974652E47726F2D66> >> >>}
\MusicPDFWriteObject{2303}{<< /Type /Annot /Subtype /Link /Border [0 0 0] /H /I /C [0 1 0] /Rect [112.557 673.973 129.877 682.84] /A << /S /GoTo /D <636974652E534E> >> >>}
\MusicPDFWriteObject{2304}{<< /Type /Annot /Subtype /Link /Border [0 0 0] /H /I /C [1 0 0] /Rect [263.096 618.951 293.874 629.742] /A << /S /GoTo /D <74686D2E342E322E312E33> >> >>}
\MusicPDFWriteObject{2305}{<< /Type /Annot /Subtype /Link /Border [0 0 0] /H /I /C [1 0 0] /Rect [270.104 507.577 305.298 519.53] /A << /S /GoTo /D <6571756174696F6E2E352E332E322E3133> >> >>}
\MusicPDFWriteObject{2306}{<< /Type /Annot /Subtype /Link /Border [0 0 0] /H /I /C [0 1 0] /Rect [329.388 510.067 348.493 518.935] /A << /S /GoTo /D <636974652E4E42> >> >>}
\MusicPDFWriteObject{2307}{<< /Type /Annot /Subtype /Link /Border [0 0 0] /H /I /C [0 1 0] /Rect [358.586 483.766 377.559 492.633] /A << /S /GoTo /D <636974652E547265> >> >>}
\MusicPDFWriteObject{2308}{<< /Type /Annot /Subtype /Link /Border [0 0 0] /H /I /C [1 0 0] /Rect [71.004 443.073 103.228 455.028] /A << /S /GoTo /D <74686D2E442E332E332E31> >> >>}
\MusicPDFWriteObject{2309}{<< /Type /Annot /Subtype /Link /Border [0 0 0] /H /I /C [1 0 0] /Rect [119.894 443.073 136.934 455.028] /A << /S /GoTo /D <73656374696F6E2E442E33> >> >>}
\MusicPDFWriteObject{2310}{<< /Type /Annot /Subtype /Link /Border [0 0 0] /H /I /C [1 0 0] /Rect [275.895 381.843 309.354 392.634] /A << /S /GoTo /D <74686D2E442E322E312E31> >> >>}
\MusicPDFWriteObject{2311}{<< /Type /Annot /Subtype /Link /Border [0 0 0] /H /I /C [1 0 0] /Rect [233.732 263.486 267.191 274.277] /A << /S /GoTo /D <74686D2E442E332E332E37> >> >>}
\MusicPDFWriteObject{2312}{<< /Type /Annot /Subtype /Link /Border [0 0 0] /H /I /C [0 1 0] /Rect [94.035 519.106 122.913 527.974] /A << /S /GoTo /D <636974652E5363682D61> >> >>}
\MusicPDFWriteObject{2313}{<< /Type /Annot /Subtype /Link /Border [0 0 0] /H /I /C [0 1 0] /Rect [224.717 492.805 241.872 501.673] /A << /S /GoTo /D <636974652E466F6C> >> >>}
\MusicPDFWriteObject{2314}{<< /Type /Annot /Subtype /Link /Border [0 0 0] /H /I /C [0 1 0] /Rect [467.344 479.654 496.222 488.522] /A << /S /GoTo /D <636974652E5363682D61> >> >>}
\MusicPDFWriteObject{2315}{<< /Type /Annot /Subtype /Link /Border [0 0 0] /H /I /C [0 1 0] /Rect [108.106 453.352 136.984 462.22] /A << /S /GoTo /D <636974652E5363682D61> >> >>}
\MusicPDFWriteObject{2316}{<< /Type /Annot /Subtype /Link /Border [0 0 0] /H /I /C [0 1 0] /Rect [269.451 279.752 288.556 288.62] /A << /S /GoTo /D <636974652E4E42> >> >>}
\MusicPDFWriteObject{2317}{<< /Type /Annot /Subtype /Link /Border [0 0 0] /H /I /C [0 1 0] /Rect [384.111 279.752 409.171 288.62] /A << /S /GoTo /D <636974652E53636861> >> >>}
\MusicPDFWriteObject{2318}{<< /Type /Annot /Subtype /Link /Border [0 0 0] /H /I /C [0 1 0] /Rect [315.91 253.451 335.015 262.318] /A << /S /GoTo /D <636974652E4E42> >> >>}
\MusicPDFWriteObject{2319}{<< /Type /Annot /Subtype /Link /Border [0 0 0] /H /I /C [0 1 0] /Rect [514.398 240.3 543.159 249.168] /A << /S /GoTo /D <636974652E47726F2D66> >> >>}
\MusicPDFWriteObject{2320}{<< /Type /Annot /Subtype /Link /Border [0 0 0] /H /I /C [0 1 0] /Rect [201.575 227.149 218.896 236.017] /A << /S /GoTo /D <636974652E534E> >> >>}
\MusicPDFWriteObject{2321}{<< /Type /Annot /Subtype /Link /Border [0 0 0] /H /I /C [0 1 0] /Rect [392.067 200.848 420.945 209.715] /A << /S /GoTo /D <636974652E5363682D61> >> >>}
\MusicPDFWriteObject{2322}{<< /Type /Annot /Subtype /Link /Border [0 0 0] /H /I /C [0 1 0] /Rect [472.755 161.395 491.86 170.263] /A << /S /GoTo /D <636974652E4E42> >> >>}
\MusicPDFWriteObject{2323}{<< /Type /Annot /Subtype /Link /Border [0 0 0] /H /I /C [1 0 0] /Rect [302.8 146.308 308.763 158.612] /A << /S /GoTo /D <48666F6F746E6F74652E3231> >> >>}
\MusicPDFWriteObject{2324}{<< /Type /Annot /Subtype /Link /Border [0 0 0] /H /I /C [0 1 0] /Rect [369.269 585.736 398.148 594.604] /A << /S /GoTo /D <636974652E5363682D61> >> >>}
\MusicPDFWriteObject{2325}{<< /Type /Annot /Subtype /Link /Border [0 0 0] /H /I /C [1 0 0] /Rect [116.962 491.19 149.526 503.145] /A << /S /GoTo /D <74686D2E442E332E332E31> >> >>}
\MusicPDFWriteObject{2326}{<< /Type /Annot /Subtype /Link /Border [0 0 0] /H /I /C [1 0 0] /Rect [419.967 167.271 450.18 179.224] /A << /S /GoTo /D <6571756174696F6E2E352E332E332E31> >> >>}
\MusicPDFWriteObject{2327}{<< /Type /Annot /Subtype /Link /Border [0 0 0] /H /I /C [1 0 0] /Rect [148.339 140.97 178.552 152.922] /A << /S /GoTo /D <6571756174696F6E2E352E332E332E31> >> >>}
\MusicPDFWriteObject{2328}{<< /Type /Annot /Subtype /Link /Border [0 0 0] /H /I /C [1 0 0] /Rect [71.004 547.961 103.228 559.916] /A << /S /GoTo /D <74686D2E442E332E332E37> >> >>}
\MusicPDFWriteObject{2329}{<< /Type /Annot /Subtype /Link /Border [0 0 0] /H /I /C [1 0 0] /Rect [399.043 610.229 432.501 622.182] /A << /S /GoTo /D <74686D2E442E332E322E33> >> >>}
\MusicPDFWriteObject{2330}{<< /Type /Annot /Subtype /Link /Border [0 0 0] /H /I /C [1 0 0] /Rect [213.071 523.179 243.284 535.132] /A << /S /GoTo /D <6571756174696F6E2E352E332E332E38> >> >>}
\MusicPDFWriteObject{2331}{<< /Type /Annot /Subtype /Link /Border [0 0 0] /H /I /C [1 0 0] /Rect [106.556 510.028 139.398 521.981] /A << /S /GoTo /D <74686D2E442E332E322E33> >> >>}
\MusicPDFWriteObject{2332}{<< /Type /Annot /Subtype /Link /Border [0 0 0] /H /I /C [0 1 0] /Rect [156.472 486.217 181.532 495.085] /A << /S /GoTo /D <636974652E53636861> >> >>}
\MusicPDFWriteObject{2333}{<< /Type /Annot /Subtype /Link /Border [0 0 0] /H /I /C [0 1 0] /Rect [270.409 486.217 295.469 495.085] /A << /S /GoTo /D <636974652E53636861> >> >>}
\MusicPDFWriteObject{2334}{<< /Type /Annot /Subtype /Link /Border [0 0 0] /H /I /C [1 0 0] /Rect [203.36 470.576 238.554 482.529] /A << /S /GoTo /D <6571756174696F6E2E352E332E332E3131> >> >>}
\MusicPDFWriteObject{2335}{<< /Type /Annot /Subtype /Link /Border [0 0 0] /H /I /C [0 1 0] /Rect [212.549 420.463 237.609 429.331] /A << /S /GoTo /D <636974652E53636861> >> >>}
\MusicPDFWriteObject{2336}{<< /Type /Annot /Subtype /Link /Border [0 0 0] /H /I /C [0 1 0] /Rect [107.466 289.449 125.665 298.316] /A << /S /GoTo /D <636974652E4A6172> >> >>}
\MusicPDFWriteObject{2337}{<< /Type /Annot /Subtype /Link /Border [0 0 0] /H /I /C [0 1 0] /Rect [316.422 289.449 334.621 298.316] /A << /S /GoTo /D <636974652E4A6172> >> >>}
\MusicPDFWriteObject{2338}{<< /Type /Annot /Subtype /Link /Border [0 0 0] /H /I /C [1 0 0] /Rect [145.394 247.506 177.618 259.459] /A << /S /GoTo /D <74686D2E442E332E332E37> >> >>}
\MusicPDFWriteObject{2339}{<< /Type /Annot /Subtype /Link /Border [0 0 0] /H /I /C [1 0 0] /Rect [296.138 197.928 328.979 209.881] /A << /S /GoTo /D <74686D2E442E322E322E32> >> >>}
\MusicPDFWriteObject{2340}{<< /Type /Annot /Subtype /Link /Border [0 0 0] /H /I /C [1 0 0] /Rect [278.738 148.35 311.734 160.302] /A << /S /GoTo /D <74686D2E442E322E322E32> >> >>}
\MusicPDFWriteObject{2341}{<< /Type /Annot /Subtype /Link /Border [0 0 0] /H /I /C [1 0 0] /Rect [473.025 135.198 511.14 147.791] /A << /S /GoTo /D <74686D2E432E322E332E3131> >> >>}
\MusicPDFWriteObject{2342}{<< /Type /Annot /Subtype /Link /Border [0 0 0] /H /I /C [1 0 0] /Rect [224.336 122.048 257.177 134.001] /A << /S /GoTo /D <74686D2E442E322E322E32> >> >>}
\MusicPDFWriteObject{2343}{<< /Type /Annot /Subtype /Link /Border [0 0 0] /H /I /C [1 0 0] /Rect [481.669 122.048 516.862 134.001] /A << /S /GoTo /D <6571756174696F6E2E352E332E332E3134> >> >>}
\MusicPDFWriteObject{2344}{<< /Type /Annot /Subtype /Link /Border [0 0 0] /H /I /C [0 1 0] /Rect [430.476 111.387 449.45 120.255] /A << /S /GoTo /D <636974652E547265> >> >>}
\MusicPDFWriteObject{2345}{<< /Type /Annot /Subtype /Link /Border [0 0 0] /H /I /C [0 1 0] /Rect [367.633 687.123 416.674 695.991] /A << /S /GoTo /D <636974652E426F752D545653> >> >>}
\MusicPDFWriteObject{2346}{<< /Type /Annot /Subtype /Link /Border [0 0 0] /H /I /C [1 0 0] /Rect [106.556 671.482 139.398 684.075] /A << /S /GoTo /D <74686D2E442E322E322E32> >> >>}
\MusicPDFWriteObject{2347}{<< /Type /Annot /Subtype /Link /Border [0 0 0] /H /I /C [1 0 0] /Rect [341.003 658.332 373.844 670.284] /A << /S /GoTo /D <74686D2E442E332E332E37> >> >>}
\MusicPDFWriteObject{2348}{<< /Type /Annot /Subtype /Link /Border [0 0 0] /H /I /C [1 0 0] /Rect [139.572 547.839 172.013 559.792] /A << /S /GoTo /D <74686D2E442E312E382E31> >> >>}
\MusicPDFWriteObject{2349}{<< /Type /Annot /Subtype /Link /Border [0 0 0] /H /I /C [0 1 0] /Rect [245.291 550.329 263.49 559.197] /A << /S /GoTo /D <636974652E4A6172> >> >>}
\MusicPDFWriteObject{2350}{<< /Type /Annot /Subtype /Link /Border [0 0 0] /H /I /C [1 0 0] /Rect [214.722 533.994 248.18 548.276] /A << /S /GoTo /D <74686D2E442E312E352E37> >> >>}
\MusicPDFWriteObject{2351}{<< /Type /Annot /Subtype /Link /Border [0 0 0] /H /I /C [1 0 0] /Rect [303.82 520.709 334.287 534.99] /A << /S /GoTo /D <74686D2E322E312E312E36> >> >>}
\MusicPDFWriteObject{2352}{<< /Type /Annot /Subtype /Link /Border [0 0 0] /H /I /C [0 1 0] /Rect [378.629 523.893 397.602 532.761] /A << /S /GoTo /D <636974652E547265> >> >>}
\MusicPDFWriteObject{2353}{<< /Type /Annot /Subtype /Link /Border [0 0 0] /H /I /C [1 0 0] /Rect [378.474 458.451 413.667 470.404] /A << /S /GoTo /D <6571756174696F6E2E352E332E332E3133> >> >>}
\MusicPDFWriteObject{2354}{<< /Type /Annot /Subtype /Link /Border [0 0 0] /H /I /C [0 1 0] /Rect [441.666 460.941 459.864 469.809] /A << /S /GoTo /D <636974652E4A6172> >> >>}
\MusicPDFWriteObject{2355}{<< /Type /Annot /Subtype /Link /Border [0 0 0] /H /I /C [1 0 0] /Rect [424.163 418.999 459.356 430.951] /A << /S /GoTo /D <6571756174696F6E2E352E332E332E3136> >> >>}
\MusicPDFWriteObject{2356}{<< /Type /Annot /Subtype /Link /Border [0 0 0] /H /I /C [1 0 0] /Rect [474.728 277.364 507.569 289.319] /A << /S /GoTo /D <74686D2E442E332E332E31> >> >>}
\MusicPDFWriteObject{2357}{<< /Type /Annot /Subtype /Link /Border [0 0 0] /H /I /C [0 1 0] /Rect [500.131 221.289 518.329 230.156] /A << /S /GoTo /D <636974652E4A6172> >> >>}
\MusicPDFWriteObject{2358}{<< /Type /Annot /Subtype /Link /Border [0 0 0] /H /I /C [1 0 0] /Rect [309.131 205.648 346.537 217.6] /A << /S /GoTo /D <6571756174696F6E2E432E322E332E3134> >> >>}
\MusicPDFWriteObject{2359}{<< /Type /Annot /Subtype /Link /Border [0 0 0] /H /I /C [0 1 0] /Rect [390.522 208.138 408.721 217.006] /A << /S /GoTo /D <636974652E4A6172> >> >>}
\MusicPDFWriteObject{2360}{<< /Type /Annot /Subtype /Link /Border [0 0 0] /H /I /C [0 1 0] /Rect [263.339 147.891 282.312 156.758] /A << /S /GoTo /D <636974652E547265> >> >>}
\MusicPDFWriteObject{2361}{<< /Type /Annot /Subtype /Link /Border [0 0 0] /H /I /C [1 0 0] /Rect [292.848 92.798 328.041 104.75] /A << /S /GoTo /D <6571756174696F6E2E352E332E332E3133> >> >>}
\MusicPDFWriteObject{2362}{<< /Type /Annot /Subtype /Link /Border [0 0 0] /H /I /C [0 1 0] /Rect [73.715 608.219 98.775 617.087] /A << /S /GoTo /D <636974652E53636861> >> >>}
\MusicPDFWriteObject{2363}{<< /Type /Annot /Subtype /Link /Border [0 0 0] /H /I /C [0 1 0] /Rect [465.887 567.471 484.991 576.339] /A << /S /GoTo /D <636974652E4E42> >> >>}
\MusicPDFWriteObject{2364}{<< /Type /Annot /Subtype /Link /Border [0 0 0] /H /I /C [0 1 0] /Rect [263.568 514.868 281.766 523.736] /A << /S /GoTo /D <636974652E4A6172> >> >>}
\MusicPDFWriteObject{2365}{<< /Type /Annot /Subtype /Link /Border [0 0 0] /H /I /C [1 0 0] /Rect [504.35 512.378 537.191 524.331] /A << /S /GoTo /D <74686D2E442E332E362E31> >> >>}
\MusicPDFWriteObject{2366}{<< /Type /Annot /Subtype /Link /Border [0 0 0] /H /I /C [0 1 0] /Rect [107.479 501.718 126.452 510.585] /A << /S /GoTo /D <636974652E547265> >> >>}
\MusicPDFWriteObject{2367}{<< /Type /Annot /Subtype /Link /Border [0 0 0] /H /I /C [1 0 0] /Rect [254.465 473.479 287.307 484.27] /A << /S /GoTo /D <74686D2E442E332E362E31> >> >>}
\MusicPDFWriteObject{2368}{<< /Type /Annot /Subtype /Link /Border [0 0 0] /H /I /C [0 1 0] /Rect [314.386 299.781 333.359 308.649] /A << /S /GoTo /D <636974652E547265> >> >>}
\MusicPDFWriteObject{2369}{<< /Type /Annot /Subtype /Link /Border [0 0 0] /H /I /C [0 1 0] /Rect [525.071 299.781 543.27 308.649] /A << /S /GoTo /D <636974652E4A6172> >> >>}
\MusicPDFWriteObject{2370}{<< /Type /Annot /Subtype /Link /Border [0 0 0] /H /I /C [1 0 0] /Rect [265.952 270.99 301.145 282.942] /A << /S /GoTo /D <6571756174696F6E2E352E332E332E3138> >> >>}
\MusicPDFWriteObject{2371}{<< /Type /Annot /Subtype /Link /Border [0 0 0] /H /I /C [0 1 0] /Rect [124.161 234.028 149.221 242.895] /A << /S /GoTo /D <636974652E53636861> >> >>}
\MusicPDFWriteObject{2372}{<< /Type /Annot /Subtype /Link /Border [0 0 0] /H /I /C [0 1 0] /Rect [351.793 234.028 369.991 242.895] /A << /S /GoTo /D <636974652E4A6172> >> >>}
\MusicPDFWriteObject{2373}{<< /Type /Annot /Subtype /Link /Border [0 0 0] /H /I /C [0 1 0] /Rect [354.919 207.726 373.892 216.594] /A << /S /GoTo /D <636974652E547265> >> >>}
\MusicPDFWriteObject{2374}{<< /Type /Annot /Subtype /Link /Border [0 0 0] /H /I /C [1 0 0] /Rect [306.263 192.085 341.456 204.038] /A << /S /GoTo /D <6571756174696F6E2E352E332E332E3138> >> >>}
\MusicPDFWriteObject{2375}{<< /Type /Annot /Subtype /Link /Border [0 0 0] /H /I /C [1 0 0] /Rect [470.499 684.633 505.692 696.586] /A << /S /GoTo /D <6571756174696F6E2E352E332E332E3138> >> >>}
\MusicPDFWriteObject{2376}{<< /Type /Annot /Subtype /Link /Border [0 0 0] /H /I /C [0 1 0] /Rect [431.453 673.973 450.557 682.84] /A << /S /GoTo /D <636974652E4E42> >> >>}
\MusicPDFWriteObject{2377}{<< /Type /Annot /Subtype /Link /Border [0 0 0] /H /I /C [1 0 0] /Rect [439.203 658.332 474.396 670.284] /A << /S /GoTo /D <6571756174696F6E2E352E332E332E3138> >> >>}
\MusicPDFWriteObject{2378}{<< /Type /Annot /Subtype /Link /Border [0 0 0] /H /I /C [0 1 0] /Rect [73.715 647.671 91.914 656.539] /A << /S /GoTo /D <636974652E4A6172> >> >>}
\MusicPDFWriteObject{2379}{<< /Type /Annot /Subtype /Link /Border [0 0 0] /H /I /C [1 0 0] /Rect [506.695 605.729 541.888 617.681] /A << /S /GoTo /D <6571756174696F6E2E352E332E332E3138> >> >>}
\MusicPDFWriteObject{2380}{<< /Type /Annot /Subtype /Link /Border [0 0 0] /H /I /C [1 0 0] /Rect [213.573 505.525 248.766 517.478] /A << /S /GoTo /D <6571756174696F6E2E352E332E332E3138> >> >>}
\MusicPDFWriteObject{2381}{<< /Type /Annot /Subtype /Link /Border [0 0 0] /H /I /C [1 0 0] /Rect [354.918 492.374 390.111 504.327] /A << /S /GoTo /D <6571756174696F6E2E352E332E332E3230> >> >>}
\MusicPDFWriteObject{2382}{<< /Type /Annot /Subtype /Link /Border [0 0 0] /H /I /C [1 0 0] /Rect [176.918 426.62 209.76 438.573] /A << /S /GoTo /D <74686D2E442E332E322E35> >> >>}
\MusicPDFWriteObject{2383}{<< /Type /Annot /Subtype /Link /Border [0 0 0] /H /I /C [0 1 0] /Rect [233.221 429.11 250.694 437.978] /A << /S /GoTo /D <636974652E4346> >> >>}
\MusicPDFWriteObject{2384}{<< /Type /Annot /Subtype /Link /Border [0 0 0] /H /I /C [0 1 0] /Rect [243.856 259.348 261.329 268.216] /A << /S /GoTo /D <636974652E4346> >> >>}
\MusicPDFWriteObject{2385}{<< /Type /Annot /Subtype /Link /Border [0 0 0] /H /I /C [1 0 0] /Rect [259.848 243.707 290.626 255.66] /A << /S /GoTo /D <74686D2E332E332E312E35> >> >>}
\MusicPDFWriteObject{2386}{<< /Type /Annot /Subtype /Link /Border [0 0 0] /H /I /C [1 0 0] /Rect [391.625 243.707 426.818 255.66] /A << /S /GoTo /D <6571756174696F6E2E352E332E332E3231> >> >>}
\MusicPDFWriteObject{2387}{<< /Type /Annot /Subtype /Link /Border [0 0 0] /H /I /C [1 0 0] /Rect [458.63 243.707 493.824 255.66] /A << /S /GoTo /D <6571756174696F6E2E352E332E332E3232> >> >>}
\MusicPDFWriteObject{2388}{<< /Type /Annot /Subtype /Link /Border [0 0 0] /H /I /C [1 0 0] /Rect [293.443 230.556 326.337 242.511] /A << /S /GoTo /D <74686D2E422E332E322E31> >> >>}
\MusicPDFWriteObject{2389}{<< /Type /Annot /Subtype /Link /Border [0 0 0] /H /I /C [1 0 0] /Rect [430.699 161.972 463.571 174.388] /A << /S /GoTo /D <74686D2E442E332E322E35> >> >>}
\MusicPDFWriteObject{2390}{<< /Type /Annot /Subtype /Link /Border [0 0 0] /H /I /C [1 0 0] /Rect [353.974 148.821 387.309 162.022] /A << /S /GoTo /D <74686D2E442E332E332E37> >> >>}
\MusicPDFWriteObject{2391}{<< /Type /Annot /Subtype /Link /Border [0 0 0] /H /I /C [1 0 0] /Rect [379.395 109.369 410.059 121.785] /A << /S /GoTo /D <74686D2E352E332E332E32> >> >>}
\MusicPDFWriteObject{2392}{<< /Type /Annot /Subtype /Link /Border [0 0 0] /H /I /C [1 0 0] /Rect [181.416 658.331 213.098 670.747] /A << /S /GoTo /D <74686D2E422E322E312E34> >> >>}
\MusicPDFWriteObject{2393}{<< /Type /Annot /Subtype /Link /Border [0 0 0] /H /I /C [1 0 0] /Rect [236.763 491.663 271.956 503.616] /A << /S /GoTo /D <6571756174696F6E2E352E332E332E3233> >> >>}
\MusicPDFWriteObject{2394}{<< /Type /Annot /Subtype /Link /Border [0 0 0] /H /I /C [1 0 0] /Rect [467.978 491.663 503.171 503.616] /A << /S /GoTo /D <6571756174696F6E2E352E332E332E3233> >> >>}
\MusicPDFWriteObject{2395}{<< /Type /Annot /Subtype /Link /Border [0 0 0] /H /I /C [1 0 0] /Rect [424.686 478.539 454.984 489.913] /A << /S /GoTo /D <74686D2E352E332E332E32> >> >>}
\MusicPDFWriteObject{2396}{<< /Type /Annot /Subtype /Link /Border [0 0 0] /H /I /C [1 0 0] /Rect [73.712 425.909 108.905 437.862] /A << /S /GoTo /D <6571756174696F6E2E352E332E332E3233> >> >>}
\MusicPDFWriteObject{2397}{<< /Type /Annot /Subtype /Link /Border [0 0 0] /H /I /C [1 0 0] /Rect [124.623 215.099 154.836 227.052] /A << /S /GoTo /D <74686D2E352E332E332E32> >> >>}
\MusicPDFWriteObject{2398}{<< /Type /Annot /Subtype /Link /Border [0 0 0] /H /I /C [1 0 0] /Rect [370.905 158.35 401.118 170.302] /A << /S /GoTo /D <6571756174696F6E2E332E312E322E33> >> >>}
\MusicPDFWriteObject{2399}{<< /Type /Annot /Subtype /Link /Border [0 0 0] /H /I /C [1 0 0] /Rect [432.998 158.35 463.211 170.302] /A << /S /GoTo /D <6571756174696F6E2E332E312E322E35> >> >>}
\MusicPDFWriteObject{2400}{<< /Type /Annot /Subtype /Link /Border [0 0 0] /H /I /C [0 1 0] /Rect [301.6 134.824 319.073 143.692] /A << /S /GoTo /D <636974652E4346> >> >>}
\MusicPDFWriteObject{2401}{<< /Type /Annot /Subtype /Link /Border [0 0 0] /H /I /C [1 0 0] /Rect [455.515 290.131 486.293 302.084] /A << /S /GoTo /D <74686D2E352E342E312E31> >> >>}
\MusicPDFWriteObject{2402}{<< /Type /Annot /Subtype /Link /Border [0 0 0] /H /I /C [1 0 0] /Rect [451.967 569.885 482.744 582.2] /A << /S /GoTo /D <74686D2E352E332E332E33> >> >>}
\MusicPDFWriteObject{2403}{<< /Type /Annot /Subtype /Link /Border [0 0 0] /H /I /C [1 0 0] /Rect [70.755 543.944 101.306 555.897] /A << /S /GoTo /D <74686D2E322E312E322E34> >> >>}
\MusicPDFWriteObject{2404}{<< /Type /Annot /Subtype /Link /Border [0 0 0] /H /I /C [0 1 0] /Rect [158.028 546.434 176.227 555.302] /A << /S /GoTo /D <636974652E4A6172> >> >>}
\MusicPDFWriteObject{2405}{<< /Type /Annot /Subtype /Link /Border [0 0 0] /H /I /C [1 0 0] /Rect [471.413 530.793 501.626 542.746] /A << /S /GoTo /D <6571756174696F6E2E352E332E342E31> >> >>}
\MusicPDFWriteObject{2406}{<< /Type /Annot /Subtype /Link /Border [0 0 0] /H /I /C [1 0 0] /Rect [228.895 515.118 259.249 530.058] /A << /S /GoTo /D <74686D2E322E312E342E31> >> >>}
\MusicPDFWriteObject{2407}{<< /Type /Annot /Subtype /Link /Border [0 0 0] /H /I /C [1 0 0] /Rect [231.667 504.162 264.094 516.114] /A << /S /GoTo /D <74686D2E432E322E352E32> >> >>}
\MusicPDFWriteObject{2408}{<< /Type /Annot /Subtype /Link /Border [0 0 0] /H /I /C [1 0 0] /Rect [124.645 227.995 154.999 239.95] /A << /S /GoTo /D <74686D2E332E332E322E31> >> >>}
\MusicPDFWriteObject{2409}{<< /Type /Annot /Subtype /Link /Border [0 0 0] /H /I /C [1 0 0] /Rect [211.472 227.995 241.826 239.95] /A << /S /GoTo /D <74686D2E332E332E322E32> >> >>}
\MusicPDFWriteObject{2410}{<< /Type /Annot /Subtype /Link /Border [0 0 0] /H /I /C [1 0 0] /Rect [254.991 200.953 285.204 212.905] /A << /S /GoTo /D <6571756174696F6E2E332E332E322E36> >> >>}
\MusicPDFWriteObject{2411}{<< /Type /Annot /Subtype /Link /Border [0 0 0] /H /I /C [1 0 0] /Rect [313.188 200.953 343.401 212.905] /A << /S /GoTo /D <6571756174696F6E2E332E332E322E37> >> >>}
\MusicPDFWriteObject{2412}{<< /Type /Annot /Subtype /Link /Border [0 0 0] /H /I /C [1 0 0] /Rect [394.487 200.768 424.136 215.841] /A << /S /GoTo /D <74686D2E332E332E322E31> >> >>}
\MusicPDFWriteObject{2413}{<< /Type /Annot /Subtype /Link /Border [0 0 0] /H /I /C [1 0 0] /Rect [340.951 187.802 371.164 199.755] /A << /S /GoTo /D <6571756174696F6E2E332E332E322E38> >> >>}
\MusicPDFWriteObject{2414}{<< /Type /Annot /Subtype /Link /Border [0 0 0] /H /I /C [1 0 0] /Rect [399.605 187.802 429.818 199.755] /A << /S /GoTo /D <6571756174696F6E2E332E332E322E39> >> >>}
\MusicPDFWriteObject{2415}{<< /Type /Annot /Subtype /Link /Border [0 0 0] /H /I /C [1 0 0] /Rect [481.915 187.801 512.128 200.394] /A << /S /GoTo /D <74686D2E332E332E322E31> >> >>}
\MusicPDFWriteObject{2416}{<< /Type /Annot /Subtype /Link /Border [0 0 0] /H /I /C [1 0 0] /Rect [139.336 174.651 169.041 186.606] /A << /S /GoTo /D <74686D2E332E332E322E32> >> >>}
\MusicPDFWriteObject{2417}{<< /Type /Annot /Subtype /Link /Border [0 0 0] /H /I /C [1 0 0] /Rect [73.712 95.747 103.925 107.699] /A << /S /GoTo /D <6571756174696F6E2E332E332E312E33> >> >>}
\MusicPDFWriteObject{2418}{<< /Type /Annot /Subtype /Link /Border [0 0 0] /H /I /C [1 0 0] /Rect [165.509 95.386 196.23 107.701] /A << /S /GoTo /D <74686D2E352E312E332E31> >> >>}
\MusicPDFWriteObject{2419}{<< /Type /Annot /Subtype /Link /Border [0 0 0] /H /I /C [0 1 0] /Rect [285.611 687.123 303.81 695.991] /A << /S /GoTo /D <636974652E4A6172> >> >>}
\MusicPDFWriteObject{2420}{<< /Type /Annot /Subtype /Link /Border [0 0 0] /H /I /C [1 0 0] /Rect [399.319 474.22 429.532 486.176] /A << /S /GoTo /D <74686D2E352E332E332E33> >> >>}
\MusicPDFWriteObject{2421}{<< /Type /Annot /Subtype /Link /Border [0 0 0] /H /I /C [0 1 0] /Rect [405.138 463.56 433.898 472.428] /A << /S /GoTo /D <636974652E47726F2D66> >> >>}
\MusicPDFWriteObject{2422}{<< /Type /Annot /Subtype /Link /Border [0 0 0] /H /I /C [0 1 0] /Rect [92.638 450.41 109.958 459.277] /A << /S /GoTo /D <636974652E534E> >> >>}
\MusicPDFWriteObject{2423}{<< /Type /Annot /Subtype /Link /Border [0 0 0] /H /I /C [0 1 0] /Rect [524.96 370.713 543.159 379.58] /A << /S /GoTo /D <636974652E4A6172> >> >>}
\MusicPDFWriteObject{2424}{<< /Type /Annot /Subtype /Link /Border [0 0 0] /H /I /C [1 0 0] /Rect [158.279 120.883 187.928 135.823] /A << /S /GoTo /D <74686D2E352E332E332E33> >> >>}
\MusicPDFWriteObject{2425}{<< /Type /Annot /Subtype /Link /Border [0 0 0] /H /I /C [0 1 0] /Rect [312.124 112.417 330.322 121.285] /A << /S /GoTo /D <636974652E4A6172> >> >>}
\MusicPDFWriteObject{2426}{<< /Type /Annot /Subtype /Link /Border [0 0 0] /H /I /C [1 0 0] /Rect [336.44 96.776 366.653 108.729] /A << /S /GoTo /D <6571756174696F6E2E352E332E342E33> >> >>}
\MusicPDFWriteObject{2427}{<< /Type /Annot /Subtype /Link /Border [0 0 0] /H /I /C [1 0 0] /Rect [140.121 684.633 172.962 696.586] /A << /S /GoTo /D <74686D2E442E322E322E32> >> >>}
\MusicPDFWriteObject{2428}{<< /Type /Annot /Subtype /Link /Border [0 0 0] /H /I /C [0 1 0] /Rect [87.492 598.184 105.691 607.052] /A << /S /GoTo /D <636974652E4A6172> >> >>}
\MusicPDFWriteObject{2429}{<< /Type /Annot /Subtype /Link /Border [0 0 0] /H /I /C [0 1 0] /Rect [524.96 558.732 543.159 567.6] /A << /S /GoTo /D <636974652E4A6172> >> >>}
\MusicPDFWriteObject{2430}{<< /Type /Annot /Subtype /Link /Border [0 0 0] /H /I /C [1 0 0] /Rect [145.619 454.948 177.844 466.652] /A << /S /GoTo /D <74686D2E442E312E322E31> >> >>}
\MusicPDFWriteObject{2431}{<< /Type /Annot /Subtype /Link /Border [0 0 0] /H /I /C [1 0 0] /Rect [439.902 297.93 457.249 311.59] /A << /S /GoTo /D <73656374696F6E2E442E31> >> >>}
\MusicPDFWriteObject{2432}{<< /Type /Annot /Subtype /Link /Border [0 0 0] /H /I /C [1 0 0] /Rect [387.316 272.588 417.529 284.541] /A << /S /GoTo /D <6571756174696F6E2E352E332E342E38> >> >>}
\MusicPDFWriteObject{2433}{<< /Type /Annot /Subtype /Link /Border [0 0 0] /H /I /C [1 0 0] /Rect [428.297 171.062 461 183.015] /A << /S /GoTo /D <6571756174696F6E2E412E322E342E35> >> >>}
\MusicPDFWriteObject{2434}{<< /Type /Annot /Subtype /Link /Border [0 0 0] /H /I /C [1 0 0] /Rect [453.115 656.756 485.818 668.709] /A << /S /GoTo /D <6571756174696F6E2E412E322E342E38> >> >>}
\MusicPDFWriteObject{2435}{<< /Type /Annot /Subtype /Link /Border [0 0 0] /H /I /C [1 0 0] /Rect [142.195 643.152 177.388 655.104] /A << /S /GoTo /D <6571756174696F6E2E352E332E342E3130> >> >>}
\MusicPDFWriteObject{2436}{<< /Type /Annot /Subtype /Link /Border [0 0 0] /H /I /C [1 0 0] /Rect [183.606 614.878 213.819 626.831] /A << /S /GoTo /D <6571756174696F6E2E352E332E342E39> >> >>}
\MusicPDFWriteObject{2437}{<< /Type /Annot /Subtype /Link /Border [0 0 0] /H /I /C [1 0 0] /Rect [511.045 614.878 543.27 626.831] /A << /S /GoTo /D <74686D2E442E312E322E31> >> >>}
\MusicPDFWriteObject{2438}{<< /Type /Annot /Subtype /Link /Border [0 0 0] /H /I /C [1 0 0] /Rect [244.087 520.67 276.79 532.983] /A << /S /GoTo /D <74686D2E412E322E322E35> >> >>}
\MusicPDFWriteObject{2439}{<< /Type /Annot /Subtype /Link /Border [0 0 0] /H /I /C [1 0 0] /Rect [405.755 492.514 436.307 506.436] /A << /S /GoTo /D <74686D2E352E332E332E33> >> >>}
\MusicPDFWriteObject{2440}{<< /Type /Annot /Subtype /Link /Border [0 0 0] /H /I /C [1 0 0] /Rect [107.048 479.364 137.684 493.285] /A << /S /GoTo /D <74686D2E352E332E332E33> >> >>}
\MusicPDFWriteObject{2441}{<< /Type /Annot /Subtype /Link /Border [0 0 0] /H /I /C [0 1 0] /Rect [524.199 482.548 543.173 491.416] /A << /S /GoTo /D <636974652E547265> >> >>}
\MusicPDFWriteObject{2442}{<< /Type /Annot /Subtype /Link /Border [0 0 0] /H /I /C [1 0 0] /Rect [176.325 150.348 211.519 162.301] /A << /S /GoTo /D <6571756174696F6E2E352E332E342E3131> >> >>}
\MusicPDFWriteObject{2443}{<< /Type /Annot /Subtype /Link /Border [0 0 0] /H /I /C [1 0 0] /Rect [488.908 150.348 519.121 162.301] /A << /S /GoTo /D <6571756174696F6E2E352E332E342E36> >> >>}
\MusicPDFWriteObject{2444}{<< /Type /Annot /Subtype /Link /Border [0 0 0] /H /I /C [1 0 0] /Rect [141.288 618.755 171.501 630.708] /A << /S /GoTo /D <6571756174696F6E2E332E332E312E32> >> >>}
\MusicPDFWriteObject{2445}{<< /Type /Annot /Subtype /Link /Border [0 0 0] /H /I /C [1 0 0] /Rect [232.629 618.755 263.407 630.71] /A << /S /GoTo /D <74686D2E332E332E312E31> >> >>}
\MusicPDFWriteObject{2446}{<< /Type /Annot /Subtype /Link /Border [0 0 0] /H /I /C [1 0 0] /Rect [70.755 602.589 100.545 617.559] /A << /S /GoTo /D <74686D2E332E332E312E35> >> >>}
\MusicPDFWriteObject{2447}{<< /Type /Annot /Subtype /Link /Border [0 0 0] /H /I /C [1 0 0] /Rect [483.048 591.27 515.942 603.585] /A << /S /GoTo /D <74686D2E422E332E322E31> >> >>}
\MusicPDFWriteObject{2448}{<< /Type /Annot /Subtype /Link /Border [0 0 0] /H /I /C [1 0 0] /Rect [333.597 559.048 365.821 576.46] /A << /S /GoTo /D <74686D2E442E332E322E35> >> >>}
\MusicPDFWriteObject{2449}{<< /Type /Annot /Subtype /Link /Border [0 0 0] /H /I /C [1 0 0] /Rect [494.726 548.092 526.982 560.045] /A << /S /GoTo /D <74686D2E442E332E332E37> >> >>}
\MusicPDFWriteObject{2450}{<< /Type /Annot /Subtype /Link /Border [0 0 0] /H /I /C [1 0 0] /Rect [216.493 521.79 246.875 533.743] /A << /S /GoTo /D <74686D2E352E332E342E33> >> >>}
\MusicPDFWriteObject{2451}{<< /Type /Annot /Subtype /Link /Border [0 0 0] /H /I /C [1 0 0] /Rect [351.547 521.79 381.93 533.743] /A << /S /GoTo /D <74686D2E352E332E332E33> >> >>}
\MusicPDFWriteObject{2452}{<< /Type /Annot /Subtype /Link /Border [0 0 0] /H /I /C [1 0 0] /Rect [211.585 435.24 241.234 447.555] /A << /S /GoTo /D <74686D2E352E332E342E32> >> >>}
\MusicPDFWriteObject{2453}{<< /Type /Annot /Subtype /Link /Border [0 0 0] /H /I /C [1 0 0] /Rect [209.897 662.12 239.545 674.495] /A << /S /GoTo /D <74686D2E352E332E342E35> >> >>}
\MusicPDFWriteObject{2454}{<< /Type /Annot /Subtype /Link /Border [0 0 0] /H /I /C [1 0 0] /Rect [223.058 601.062 258.251 613.014] /A << /S /GoTo /D <6571756174696F6E2E352E332E342E3132> >> >>}
\MusicPDFWriteObject{2455}{<< /Type /Annot /Subtype /Link /Border [0 0 0] /H /I /C [1 0 0] /Rect [345.45 601.062 375.663 613.014] /A << /S /GoTo /D <74686D2E332E342E332E33> >> >>}
\MusicPDFWriteObject{2456}{<< /Type /Annot /Subtype /Link /Border [0 0 0] /H /I /C [1 0 0] /Rect [412.443 601.062 447.636 613.014] /A << /S /GoTo /D <6571756174696F6E2E332E342E332E3133> >> >>}
\MusicPDFWriteObject{2457}{<< /Type /Annot /Subtype /Link /Border [0 0 0] /H /I /C [1 0 0] /Rect [270.118 587.911 300.331 599.864] /A << /S /GoTo /D <6571756174696F6E2E332E342E312E36> >> >>}
\MusicPDFWriteObject{2458}{<< /Type /Annot /Subtype /Link /Border [0 0 0] /H /I /C [1 0 0] /Rect [369.233 587.911 398.881 599.864] /A << /S /GoTo /D <74686D2E332E342E312E31> >> >>}
\MusicPDFWriteObject{2459}{<< /Type /Annot /Subtype /Link /Border [0 0 0] /H /I /C [1 0 0] /Rect [142.109 561.609 172.322 573.562] /A << /S /GoTo /D <6571756174696F6E2E332E332E312E32> >> >>}
\MusicPDFWriteObject{2460}{<< /Type /Annot /Subtype /Link /Border [0 0 0] /H /I /C [1 0 0] /Rect [243.149 561.609 273.362 573.562] /A << /S /GoTo /D <74686D2E332E332E312E31> >> >>}
\MusicPDFWriteObject{2461}{<< /Type /Annot /Subtype /Link /Border [0 0 0] /H /I /C [1 0 0] /Rect [386.996 321.192 417.209 333.145] /A << /S /GoTo /D <74686D2E352E342E312E31> >> >>}
\MusicPDFWriteObject{2462}{<< /Type /Annot /Subtype /Link /Border [0 0 0] /H /I /C [1 0 0] /Rect [336.102 297.152 366.315 309.105] /A << /S /GoTo /D <74686D2E352E312E312E31> >> >>}
\MusicPDFWriteObject{2463}{<< /Type /Annot /Subtype /Link /Border [0 0 0] /H /I /C [1 0 0] /Rect [429.202 297.152 459.415 309.105] /A << /S /GoTo /D <6571756174696F6E2E352E312E312E34> >> >>}
\MusicPDFWriteObject{2464}{<< /Type /Annot /Subtype /Link /Border [0 0 0] /H /I /C [1 0 0] /Rect [381.639 233.715 411.852 245.667] /A << /S /GoTo /D <6571756174696F6E2E332E332E312E33> >> >>}
\MusicPDFWriteObject{2465}{<< /Type /Annot /Subtype /Link /Border [0 0 0] /H /I /C [1 0 0] /Rect [470.479 233.715 500.692 245.667] /A << /S /GoTo /D <74686D2E332E332E312E31> >> >>}
\MusicPDFWriteObject{2466}{<< /Type /Annot /Subtype /Link /Border [0 0 0] /H /I /C [1 0 0] /Rect [235.283 671.482 265.497 683.435] /A << /S /GoTo /D <6571756174696F6E2E332E312E322E35> >> >>}
\MusicPDFWriteObject{2467}{<< /Type /Annot /Subtype /Link /Border [0 0 0] /H /I /C [1 0 0] /Rect [496.942 595.603 527.155 607.555] /A << /S /GoTo /D <6571756174696F6E2E332E332E312E33> >> >>}
\MusicPDFWriteObject{2468}{<< /Type /Annot /Subtype /Link /Border [0 0 0] /H /I /C [1 0 0] /Rect [468.863 582.452 499.077 594.405] /A << /S /GoTo /D <6571756174696F6E2E332E312E322E35> >> >>}
\MusicPDFWriteObject{2469}{<< /Type /Annot /Subtype /Link /Border [0 0 0] /H /I /C [1 0 0] /Rect [272.841 362.655 303.054 373.446] /A << /S /GoTo /D <74686D2E342E342E312E31> >> >>}
\MusicPDFWriteObject{2470}{<< /Type /Annot /Subtype /Link /Border [0 0 0] /H /I /C [1 0 0] /Rect [366.688 362.655 373.661 373.446] /A << /S /GoTo /D <636861707465722E35> >> >>}
\MusicPDFWriteObject{2471}{<< /Type /Annot /Subtype /Link /Border [0 0 0] /H /I /C [1 0 0] /Rect [263.666 332.555 293.879 344.507] /A << /S /GoTo /D <6571756174696F6E2E352E342E312E32> >> >>}
\MusicPDFWriteObject{2472}{<< /Type /Annot /Subtype /Link /Border [0 0 0] /H /I /C [1 0 0] /Rect [234.568 275.816 264.781 287.768] /A << /S /GoTo /D <6571756174696F6E2E352E342E312E32> >> >>}
\MusicPDFWriteObject{2473}{<< /Type /Annot /Subtype /Link /Border [0 0 0] /H /I /C [1 0 0] /Rect [125.986 219.607 156.764 231.922] /A << /S /GoTo /D <74686D2E322E332E312E34> >> >>}
\MusicPDFWriteObject{2474}{<< /Type /Annot /Subtype /Link /Border [0 0 0] /H /I /C [1 0 0] /Rect [73.712 144.087 103.925 156.04] /A << /S /GoTo /D <6571756174696F6E2E322E342E332E35> >> >>}
\MusicPDFWriteObject{2475}{<< /Type /Annot /Subtype /Link /Border [0 0 0] /H /I /C [1 0 0] /Rect [177.144 144.087 207.555 156.04] /A << /S /GoTo /D <74686D2E322E342E332E32> >> >>}
\MusicPDFWriteObject{2476}{<< /Type /Annot /Subtype /Link /Border [0 0 0] /H /I /C [1 0 0] /Rect [462.38 144.087 492.593 156.04] /A << /S /GoTo /D <6571756174696F6E2E322E342E332E35> >> >>}
\MusicPDFWriteObject{2477}{<< /Type /Annot /Subtype /Link /Border [0 0 0] /H /I /C [1 0 0] /Rect [187.656 684.633 217.869 696.586] /A << /S /GoTo /D <6571756174696F6E2E352E342E312E34> >> >>}
\MusicPDFWriteObject{2478}{<< /Type /Annot /Subtype /Link /Border [0 0 0] /H /I /C [1 0 0] /Rect [166.97 611.354 197.183 623.307] /A << /S /GoTo /D <6571756174696F6E2E352E342E312E35> >> >>}
\MusicPDFWriteObject{2479}{<< /Type /Annot /Subtype /Link /Border [0 0 0] /H /I /C [1 0 0] /Rect [225.4 611.354 255.613 623.307] /A << /S /GoTo /D <6571756174696F6E2E352E342E312E36> >> >>}
\MusicPDFWriteObject{2480}{<< /Type /Annot /Subtype /Link /Border [0 0 0] /H /I /C [1 0 0] /Rect [180.339 597.843 209.988 610.158] /A << /S /GoTo /D <74686D2E322E312E342E31> >> >>}
\MusicPDFWriteObject{2481}{<< /Type /Annot /Subtype /Link /Border [0 0 0] /H /I /C [0 1 0] /Rect [433.565 600.694 458.625 609.561] /A << /S /GoTo /D <636974652E53636861> >> >>}
\MusicPDFWriteObject{2482}{<< /Type /Annot /Subtype /Link /Border [0 0 0] /H /I /C [0 1 0] /Rect [107.445 587.543 126.55 596.411] /A << /S /GoTo /D <636974652E4E42> >> >>}
\MusicPDFWriteObject{2483}{<< /Type /Annot /Subtype /Link /Border [0 0 0] /H /I /C [1 0 0] /Rect [175.06 571.902 207.902 583.855] /A << /S /GoTo /D <74686D2E442E332E312E32> >> >>}
\MusicPDFWriteObject{2484}{<< /Type /Annot /Subtype /Link /Border [0 0 0] /H /I /C [0 1 0] /Rect [346.846 574.392 365.951 583.26] /A << /S /GoTo /D <636974652E4E42> >> >>}
\MusicPDFWriteObject{2485}{<< /Type /Annot /Subtype /Link /Border [0 0 0] /H /I /C [0 1 0] /Rect [107.326 561.241 126.43 570.109] /A << /S /GoTo /D <636974652E4E42> >> >>}
\MusicPDFWriteObject{2486}{<< /Type /Annot /Subtype /Link /Border [0 0 0] /H /I /C [1 0 0] /Rect [73.712 532.45 103.925 544.402] /A << /S /GoTo /D <6571756174696F6E2E352E312E322E34> >> >>}
\MusicPDFWriteObject{2487}{<< /Type /Annot /Subtype /Link /Border [0 0 0] /H /I /C [1 0 0] /Rect [168.223 532.089 198.436 544.404] /A << /S /GoTo /D <74686D2E352E312E322E32> >> >>}
\MusicPDFWriteObject{2488}{<< /Type /Annot /Subtype /Link /Border [0 0 0] /H /I /C [1 0 0] /Rect [203.7 508.366 233.913 519.157] /A << /S /GoTo /D <74686D2E352E342E312E31> >> >>}
\MusicPDFWriteObject{2489}{<< /Type /Annot /Subtype /Link /Border [0 0 0] /H /I /C [1 0 0] /Rect [333.098 416.8 363.311 428.753] /A << /S /GoTo /D <6571756174696F6E2E352E342E312E32> >> >>}
\MusicPDFWriteObject{2490}{<< /Type /Annot /Subtype /Link /Border [0 0 0] /H /I /C [1 0 0] /Rect [218.031 366.702 248.245 378.655] /A << /S /GoTo /D <6571756174696F6E2E332E312E322E35> >> >>}
\MusicPDFWriteObject{2491}{<< /Type /Annot /Subtype /Link /Border [0 0 0] /H /I /C [1 0 0] /Rect [477.259 290.822 507.472 302.775] /A << /S /GoTo /D <6571756174696F6E2E332E332E312E33> >> >>}
\MusicPDFWriteObject{2492}{<< /Type /Annot /Subtype /Link /Border [0 0 0] /H /I /C [1 0 0] /Rect [429.62 277.671 459.833 289.624] /A << /S /GoTo /D <6571756174696F6E2E332E312E322E35> >> >>}
\MusicPDFWriteObject{2493}{<< /Type /Annot /Subtype /Link /Border [0 0 0] /H /I /C [1 0 0] /Rect [195.234 242.01 224.882 254.325] /A << /S /GoTo /D <74686D2E352E332E342E34> >> >>}
\MusicPDFWriteObject{2494}{<< /Type /Annot /Subtype /Link /Border [0 0 0] /H /I /C [1 0 0] /Rect [113.072 150.081 143.539 162.036] /A << /S /GoTo /D <74686D2E332E332E312E35> >> >>}
\MusicPDFWriteObject{2495}{<< /Type /Annot /Subtype /Link /Border [0 0 0] /H /I /C [0 1 0] /Rect [336.888 634.52 361.948 643.388] /A << /S /GoTo /D <636974652E53636861> >> >>}
\MusicPDFWriteObject{2496}{<< /Type /Annot /Subtype /Link /Border [0 0 0] /H /I /C [1 0 0] /Rect [347.768 618.879 379.844 630.832] /A << /S /GoTo /D <74686D2E422E322E322E34> >> >>}
\MusicPDFWriteObject{2497}{<< /Type /Annot /Subtype /Link /Border [0 0 0] /H /I /C [1 0 0] /Rect [509.29 618.879 541.916 630.832] /A << /S /GoTo /D <74686D2E442E332E312E32> >> >>}
\MusicPDFWriteObject{2498}{<< /Type /Annot /Subtype /Link /Border [0 0 0] /H /I /C [1 0 0] /Rect [474.621 592.578 504.834 604.531] /A << /S /GoTo /D <6571756174696F6E2E352E342E312E38> >> >>}
\MusicPDFWriteObject{2499}{<< /Type /Annot /Subtype /Link /Border [0 0 0] /H /I /C [1 0 0] /Rect [400.815 539.454 433.656 551.409] /A << /S /GoTo /D <74686D2E442E332E322E35> >> >>}
\MusicPDFWriteObject{2500}{<< /Type /Annot /Subtype /Link /Border [0 0 0] /H /I /C [1 0 0] /Rect [140.131 525.944 170.344 538.259] /A << /S /GoTo /D <74686D2E352E332E342E33> >> >>}
\MusicPDFWriteObject{2501}{<< /Type /Annot /Subtype /Link /Border [0 0 0] /H /I /C [1 0 0] /Rect [104.414 457.47 134.627 469.423] /A << /S /GoTo /D <6571756174696F6E2E332E332E312E33> >> >>}
\MusicPDFWriteObject{2502}{<< /Type /Annot /Subtype /Link /Border [0 0 0] /H /I /C [1 0 0] /Rect [169.229 457.47 199.443 469.423] /A << /S /GoTo /D <6571756174696F6E2E332E312E322E35> >> >>}
\MusicPDFWriteObject{2503}{<< /Type /Annot /Subtype /Link /Border [0 0 0] /H /I /C [1 0 0] /Rect [343.102 441.537 373.315 456.856] /A << /S /GoTo /D <74686D2E352E332E342E32> >> >>}
\MusicPDFWriteObject{2504}{<< /Type /Annot /Subtype /Link /Border [0 0 0] /H /I /C [1 0 0] /Rect [511.675 388.197 541.888 400.149] /A << /S /GoTo /D <6571756174696F6E2E332E312E322E35> >> >>}
\MusicPDFWriteObject{2505}{<< /Type /Annot /Subtype /Link /Border [0 0 0] /H /I /C [1 0 0] /Rect [152.29 334.772 182.503 346.725] /A << /S /GoTo /D <74686D2E352E332E342E34> >> >>}
\MusicPDFWriteObject{2506}{<< /Type /Annot /Subtype /Link /Border [0 0 0] /H /I /C [1 0 0] /Rect [236.569 321.621 266.782 333.574] /A << /S /GoTo /D <6571756174696F6E2E332E332E312E32> >> >>}
\MusicPDFWriteObject{2507}{<< /Type /Annot /Subtype /Link /Border [0 0 0] /H /I /C [1 0 0] /Rect [318.879 321.621 349.092 333.574] /A << /S /GoTo /D <74686D2E332E332E312E31> >> >>}
\MusicPDFWriteObject{2508}{<< /Type /Annot /Subtype /Link /Border [0 0 0] /H /I /C [1 0 0] /Rect [410.784 285.034 440.997 296.987] /A << /S /GoTo /D <6571756174696F6E2E332E312E322E35> >> >>}
\MusicPDFWriteObject{2509}{<< /Type /Annot /Subtype /Link /Border [0 0 0] /H /I /C [1 0 0] /Rect [489.572 205.336 519.418 217.651] /A << /S /GoTo /D <74686D2E352E312E312E33> >> >>}
\MusicPDFWriteObject{2510}{<< /Type /Annot /Subtype /Link /Border [0 0 0] /H /I /C [1 0 0] /Rect [501.148 139.582 530.797 151.897] /A << /S /GoTo /D <74686D2E322E312E342E31> >> >>}
\MusicPDFWriteObject{2511}{<< /Type /Annot /Subtype /Link /Border [0 0 0] /H /I /C [1 0 0] /Rect [141.585 113.281 172.363 125.595] /A << /S /GoTo /D <74686D2E352E342E312E33> >> >>}
\MusicPDFWriteObject{2512}{<< /Type /Annot /Subtype /Link /Border [0 0 0] /H /I /C [1 0 0] /Rect [234.784 113.641 264.997 125.594] /A << /S /GoTo /D <6571756174696F6E2E352E342E312E37> >> >>}
\MusicPDFWriteObject{2513}{<< /Type /Annot /Subtype /Link /Border [0 0 0] /H /I /C [1 0 0] /Rect [422.99 100.49 453.203 112.443] /A << /S /GoTo /D <6571756174696F6E2E352E342E312E32> >> >>}
\MusicPDFWriteObject{2514}{<< /Type /Annot /Subtype /Link /Border [0 0 0] /H /I /C [1 0 0] /Rect [138.829 684.273 168.477 696.588] /A << /S /GoTo /D <74686D2E352E342E312E33> >> >>}
\MusicPDFWriteObject{2515}{<< /Type /Annot /Subtype /Link /Border [0 0 0] /H /I /C [0 1 0] /Rect [514.509 673.973 543.27 682.84] /A << /S /GoTo /D <636974652E47726F2D66> >> >>}
\MusicPDFWriteObject{2516}{<< /Type /Annot /Subtype /Link /Border [0 0 0] /H /I /C [0 1 0] /Rect [204.914 660.822 222.234 669.69] /A << /S /GoTo /D <636974652E534E> >> >>}
\MusicPDFWriteObject{2517}{<< /Type /Annot /Subtype /Link /Border [0 0 0] /H /I /C [1 0 0] /Rect [187.794 598.225 221.355 607.079] /A << /S /GoTo /D <74686D2E352E342E312E31> >> >>}
\MusicPDFWriteObject{2518}{<< /Type /Annot /Subtype /Link /Border [0 0 0] /H /I /C [1 0 0] /Rect [179.57 583.137 209.783 593.928] /A << /S /GoTo /D <74686D2E352E342E312E34> >> >>}
\MusicPDFWriteObject{2519}{<< /Type /Annot /Subtype /Link /Border [0 0 0] /H /I /C [1 0 0] /Rect [219.239 569.073 249.452 580.777] /A << /S /GoTo /D <74686D2E352E342E312E33> >> >>}
\MusicPDFWriteObject{2520}{<< /Type /Annot /Subtype /Link /Border [0 0 0] /H /I /C [1 0 0] /Rect [219.989 556.282 250.202 568.235] /A << /S /GoTo /D <6571756174696F6E2E352E342E312E32> >> >>}
\MusicPDFWriteObject{2521}{<< /Type /Annot /Subtype /Link /Border [0 0 0] /H /I /C [1 0 0] /Rect [164.997 530.534 187.054 541.382] /A << /S /GoTo /D <73756273656374696F6E2E352E332E33> >> >>}
\MusicPDFWriteObject{2522}{<< /Type /Annot /Subtype /Link /Border [0 0 0] /H /I /C [1 0 0] /Rect [270.877 530.534 300.525 541.382] /A << /S /GoTo /D <74686D2E352E332E342E35> >> >>}
\MusicPDFWriteObject{2523}{<< /Type /Annot /Subtype /Link /Border [0 0 0] /H /I /C [1 0 0] /Rect [71.004 516.47 100.652 528.785] /A << /S /GoTo /D <74686D2E352E332E332E34> >> >>}
\MusicPDFWriteObject{2524}{<< /Type /Annot /Subtype /Link /Border [0 0 0] /H /I /C [1 0 0] /Rect [174.45 516.47 204.099 528.785] /A << /S /GoTo /D <74686D2E352E332E342E34> >> >>}
\MusicPDFWriteObject{2525}{<< /Type /Annot /Subtype /Link /Border [0 0 0] /H /I /C [1 0 0] /Rect [274.972 464.26 305.185 475.051] /A << /S /GoTo /D <74686D2E352E342E312E34> >> >>}
\MusicPDFWriteObject{2526}{<< /Type /Annot /Subtype /Link /Border [0 0 0] /H /I /C [1 0 0] /Rect [470.69 464.26 500.903 475.051] /A << /S /GoTo /D <74686D2E352E342E312E31> >> >>}
\MusicPDFWriteObject{2527}{<< /Type /Annot /Subtype /Link /Border [0 0 0] /H /I /C [1 0 0] /Rect [243.48 321.586 275.828 338.97] /A << /S /GoTo /D <74686D2E422E322E312E31> >> >>}
\MusicPDFWriteObject{2528}{<< /Type /Annot /Subtype /Link /Border [0 0 0] /H /I /C [1 0 0] /Rect [354.205 321.586 387.108 338.97] /A << /S /GoTo /D <74686D2E442E332E332E37> >> >>}
\MusicPDFWriteObject{2529}{<< /Type /Annot /Subtype /Link /Border [0 0 0] /H /I /C [1 0 0] /Rect [288.936 310.229 326.342 322.181] /A << /S /GoTo /D <6571756174696F6E2E432E322E332E3233> >> >>}
\MusicPDFWriteObject{2530}{<< /Type /Annot /Subtype /Link /Border [0 0 0] /H /I /C [1 0 0] /Rect [166.373 297.078 199.214 309.031] /A << /S /GoTo /D <74686D2E442E312E382E31> >> >>}
\MusicPDFWriteObject{2531}{<< /Type /Annot /Subtype /Link /Border [0 0 0] /H /I /C [1 0 0] /Rect [201.018 284.48 231.231 295.271] /A << /S /GoTo /D <74686D2E322E332E312E32> >> >>}
\MusicPDFWriteObject{2532}{<< /Type /Annot /Subtype /Link /Border [0 0 0] /H /I /C [1 0 0] /Rect [295.972 284.48 326.185 295.271] /A << /S /GoTo /D <74686D2E322E342E332E32> >> >>}
\MusicPDFWriteObject{2533}{<< /Type /Annot /Subtype /Link /Border [0 0 0] /H /I /C [1 0 0] /Rect [187.859 245.028 219.723 255.819] /A << /S /GoTo /D <74686D2E422E322E312E35> >> >>}
\MusicPDFWriteObject{2534}{<< /Type /Annot /Subtype /Link /Border [0 0 0] /H /I /C [1 0 0] /Rect [249.814 148.889 279.83 160.842] /A << /S /GoTo /D <74686D2E352E332E342E35> >> >>}
\MusicPDFWriteObject{2535}{<< /Type /Annot /Subtype /Link /Border [0 0 0] /H /I /C [1 0 0] /Rect [298.801 94.717 328.45 107.701] /A << /S /GoTo /D <74686D2E352E332E312E34> >> >>}
\MusicPDFWriteObject{2536}{<< /Type /Annot /Subtype /Link /Border [0 0 0] /H /I /C [0 1 0] /Rect [438.356 405.955 467.234 414.823] /A << /S /GoTo /D <636974652E5363682D61> >> >>}
\MusicPDFWriteObject{2537}{<< /Type /Annot /Subtype /Link /Border [0 0 0] /H /I /C [1 0 0] /Rect [157.637 389.284 187.596 402.269] /A << /S /GoTo /D <74686D2E352E332E342E35> >> >>}
\MusicPDFWriteObject{2538}{<< /Type /Annot /Subtype /Link /Border [0 0 0] /H /I /C [1 0 0] /Rect [248.642 333.708 279.42 346.023] /A << /S /GoTo /D <74686D2E352E342E312E34> >> >>}
\MusicPDFWriteObject{2539}{<< /Type /Annot /Subtype /Link /Border [0 0 0] /H /I /C [1 0 0] /Rect [294.776 291.765 324.99 303.718] /A << /S /GoTo /D <74686D2E352E332E342E35> >> >>}
\MusicPDFWriteObject{2540}{<< /Type /Annot /Subtype /Link /Border [0 0 0] /H /I /C [1 0 0] /Rect [451.596 229.036 486.789 240.989] /A << /S /GoTo /D <6571756174696F6E2E352E342E312E3133> >> >>}
\MusicPDFWriteObject{2541}{<< /Type /Annot /Subtype /Link /Border [0 0 0] /H /I /C [1 0 0] /Rect [157.27 561.482 187.709 575.795] /A << /S /GoTo /D <74686D2E352E312E322E32> >> >>}
\MusicPDFWriteObject{2542}{<< /Type /Annot /Subtype /Link /Border [0 0 0] /H /I /C [0 1 0] /Rect [91.561 155.75 111.758 164.617] /A << /S /GoTo /D <636974652E41686C> >> >>}
\MusicPDFWriteObject{2543}{<< /Type /Annot /Subtype /Link /Border [0 0 0] /H /I /C [1 0 0] /Rect [156.432 518.035 186.081 530.35] /A << /S /GoTo /D <74686D2E352E332E342E34> >> >>}
\MusicPDFWriteObject{2544}{<< /Type /Annot /Subtype /Link /Border [0 0 0] /H /I /C [1 0 0] /Rect [241.323 505.245 276.516 517.198] /A << /S /GoTo /D <6571756174696F6E2E352E342E312E3136> >> >>}
\MusicPDFWriteObject{2545}{<< /Type /Annot /Subtype /Link /Border [0 0 0] /H /I /C [1 0 0] /Rect [340.068 424.667 375.261 436.619] /A << /S /GoTo /D <6571756174696F6E2E352E342E312E3134> >> >>}
\MusicPDFWriteObject{2546}{<< /Type /Annot /Subtype /Link /Border [0 0 0] /H /I /C [1 0 0] /Rect [320.063 227.304 350.841 240.98] /A << /S /GoTo /D <74686D2E352E332E342E34> >> >>}
\MusicPDFWriteObject{2547}{<< /Type /Annot /Subtype /Link /Border [0 0 0] /H /I /C [1 0 0] /Rect [125.725 211.668 156.502 228.3] /A << /S /GoTo /D <74686D2E352E342E312E34> >> >>}
\MusicPDFWriteObject{2548}{<< /Type /Annot /Subtype /Link /Border [0 0 0] /H /I /C [0 1 0] /Rect [366.478 185.81 385.583 194.678] /A << /S /GoTo /D <636974652E4E42> >> >>}
\MusicPDFWriteObject{2549}{<< /Type /Annot /Subtype /Link /Border [0 0 0] /H /I /C [1 0 0] /Rect [143.215 104.79 178.408 116.742] /A << /S /GoTo /D <6571756174696F6E2E352E342E312E3136> >> >>}
\MusicPDFWriteObject{2550}{<< /Type /Annot /Subtype /Link /Border [0 0 0] /H /I /C [0 1 0] /Rect [367.442 641.949 385.64 650.816] /A << /S /GoTo /D <636974652E4A6172> >> >>}
\MusicPDFWriteObject{2551}{<< /Type /Annot /Subtype /Link /Border [0 0 0] /H /I /C [0 1 0] /Rect [274.996 602.496 303.757 611.364] /A << /S /GoTo /D <636974652E47726F2D66> >> >>}
\MusicPDFWriteObject{2552}{<< /Type /Annot /Subtype /Link /Border [0 0 0] /H /I /C [0 1 0] /Rect [442.28 602.496 459.6 611.364] /A << /S /GoTo /D <636974652E534E> >> >>}
\MusicPDFWriteObject{2553}{<< /Type /Annot /Subtype /Link /Border [0 0 0] /H /I /C [1 0 0] /Rect [184.429 400.737 219.623 412.69] /A << /S /GoTo /D <6571756174696F6E2E352E342E312E3231> >> >>}
\MusicPDFWriteObject{2554}{<< /Type /Annot /Subtype /Link /Border [0 0 0] /H /I /C [1 0 0] /Rect [159.391 361.838 190.14 372.629] /A << /S /GoTo /D <74686D2E352E342E312E34> >> >>}
\MusicPDFWriteObject{2555}{<< /Type /Annot /Subtype /Link /Border [0 0 0] /H /I /C [0 1 0] /Rect [73.826 207.449 91.298 216.317] /A << /S /GoTo /D <636974652E4346> >> >>}
\MusicPDFWriteObject{2556}{<< /Type /Annot /Subtype /Link /Border [0 0 0] /H /I /C [0 1 0] /Rect [227.945 207.449 245.915 216.317] /A << /S /GoTo /D <636974652E5256> >> >>}
\MusicPDFWriteObject{2557}{<< /Type /Annot /Subtype /Link /Border [0 0 0] /H /I /C [0 1 0] /Rect [248.73 194.298 266.203 203.166] /A << /S /GoTo /D <636974652E4346> >> >>}
\MusicPDFWriteObject{2558}{<< /Type /Annot /Subtype /Link /Border [0 0 0] /H /I /C [1 0 0] /Rect [233.576 549.559 263.79 561.512] /A << /S /GoTo /D <74686D2E352E332E332E34> >> >>}
\MusicPDFWriteObject{2559}{<< /Type /Annot /Subtype /Link /Border [0 0 0] /H /I /C [1 0 0] /Rect [160.068 523.258 195.262 535.211] /A << /S /GoTo /D <6571756174696F6E2E352E342E312E3231> >> >>}
\MusicPDFWriteObject{2560}{<< /Type /Annot /Subtype /Link /Border [0 0 0] /H /I /C [1 0 0] /Rect [160.575 510.107 190.788 522.06] /A << /S /GoTo /D <6571756174696F6E2E352E342E312E32> >> >>}
\MusicPDFWriteObject{2561}{<< /Type /Annot /Subtype /Link /Border [0 0 0] /H /I /C [1 0 0] /Rect [490.148 510.107 520.361 522.06] /A << /S /GoTo /D <74686D2E352E342E312E31> >> >>}
\MusicPDFWriteObject{2562}{<< /Type /Annot /Subtype /Link /Border [0 0 0] /H /I /C [1 0 0] /Rect [150.741 471.208 180.954 481.999] /A << /S /GoTo /D <74686D2E352E322E332E31> >> >>}
\MusicPDFWriteObject{2563}{<< /Type /Annot /Subtype /Link /Border [0 0 0] /H /I /C [1 0 0] /Rect [300.144 471.208 332.985 481.999] /A << /S /GoTo /D <74686D2E442E332E312E32> >> >>}
\MusicPDFWriteObject{2564}{<< /Type /Annot /Subtype /Link /Border [0 0 0] /H /I /C [1 0 0] /Rect [424.114 434.738 454.186 446.691] /A << /S /GoTo /D <74686D2E352E322E312E31> >> >>}
\MusicPDFWriteObject{2565}{<< /Type /Annot /Subtype /Link /Border [0 0 0] /H /I /C [1 0 0] /Rect [347.418 408.076 377.321 420.391] /A << /S /GoTo /D <74686D2E352E342E312E31> >> >>}
\MusicPDFWriteObject{2566}{<< /Type /Annot /Subtype /Link /Border [0 0 0] /H /I /C [1 0 0] /Rect [71.004 382.688 101.781 393.479] /A << /S /GoTo /D <74686D2E352E322E312E31> >> >>}
\MusicPDFWriteObject{2567}{<< /Type /Annot /Subtype /Link /Border [0 0 0] /H /I /C [1 0 0] /Rect [453.189 319.916 483.402 331.869] /A << /S /GoTo /D <6571756174696F6E2E352E322E342E31> >> >>}
\MusicPDFWriteObject{2568}{<< /Type /Annot /Subtype /Link /Border [0 0 0] /H /I /C [1 0 0] /Rect [295.306 249.167 325.52 261.482] /A << /S /GoTo /D <74686D2E332E332E312E35> >> >>}
\MusicPDFWriteObject{2569}{<< /Type /Annot /Subtype /Link /Border [0 0 0] /H /I /C [1 0 0] /Rect [89.175 212.579 119.388 224.532] /A << /S /GoTo /D <6571756174696F6E2E332E332E312E32> >> >>}
\MusicPDFWriteObject{2570}{<< /Type /Annot /Subtype /Link /Border [0 0 0] /H /I /C [1 0 0] /Rect [191.725 212.219 222.503 224.534] /A << /S /GoTo /D <74686D2E332E332E312E31> >> >>}
\MusicPDFWriteObject{2571}{<< /Type /Annot /Subtype /Link /Border [0 0 0] /H /I /C [1 0 0] /Rect [388.365 212.219 419.142 224.534] /A << /S /GoTo /D <74686D2E332E332E312E35> >> >>}
\MusicPDFWriteObject{2572}{<< /Type /Annot /Subtype /Link /Border [0 0 0] /H /I /C [1 0 0] /Rect [436.825 95.747 467.039 107.699] /A << /S /GoTo /D <6571756174696F6E2E352E342E312E31> >> >>}
\MusicPDFWriteObject{2573}{<< /Type /Annot /Subtype /Link /Border [0 0 0] /H /I /C [1 0 0] /Rect [185.85 82.596 216.064 94.549] /A << /S /GoTo /D <6571756174696F6E2E332E332E312E32> >> >>}
\MusicPDFWriteObject{2574}{<< /Type /Annot /Subtype /Link /Border [0 0 0] /H /I /C [1 0 0] /Rect [274.69 82.596 304.903 94.549] /A << /S /GoTo /D <74686D2E332E332E312E31> >> >>}
\MusicPDFWriteObject{2575}{<< /Type /Annot /Subtype /Link /Border [0 0 0] /H /I /C [1 0 0] /Rect [112.317 671.122 142.22 683.435] /A << /S /GoTo /D <74686D2E352E342E312E31> >> >>}
\MusicPDFWriteObject{2576}{<< /Type /Annot /Subtype /Link /Border [0 0 0] /H /I /C [1 0 0] /Rect [71.004 658.332 103.845 670.284] /A << /S /GoTo /D <74686D2E442E332E332E31> >> >>}
\MusicPDFWriteObject{2577}{<< /Type /Annot /Subtype /Link /Border [0 0 0] /H /I /C [1 0 0] /Rect [166.622 658.332 199.463 670.284] /A << /S /GoTo /D <74686D2E442E332E332E37> >> >>}
\MusicPDFWriteObject{2578}{<< /Type /Annot /Subtype /Link /Border [0 0 0] /H /I /C [1 0 0] /Rect [182.445 458.326 212.658 470.279] /A << /S /GoTo /D <6571756174696F6E2E352E342E322E33> >> >>}
\MusicPDFWriteObject{2579}{<< /Type /Annot /Subtype /Link /Border [0 0 0] /H /I /C [1 0 0] /Rect [311.215 352.005 344.673 367.804] /A << /S /GoTo /D <74686D2E442E332E332E37> >> >>}
\MusicPDFWriteObject{2580}{<< /Type /Annot /Subtype /Link /Border [0 0 0] /H /I /C [1 0 0] /Rect [442.967 352.005 476.426 367.804] /A << /S /GoTo /D <74686D2E442E332E322E35> >> >>}
\MusicPDFWriteObject{2581}{<< /Type /Annot /Subtype /Link /Border [0 0 0] /H /I /C [1 0 0] /Rect [71.004 303.871 103.845 312.725] /A << /S /GoTo /D <74686D2E442E332E362E31> >> >>}
\MusicPDFWriteObject{2582}{<< /Type /Annot /Subtype /Link /Border [0 0 0] /H /I /C [1 0 0] /Rect [243.005 222.817 275.847 233.608] /A << /S /GoTo /D <74686D2E442E332E322E35> >> >>}
\MusicPDFWriteObject{2583}{<< /Type /Annot /Subtype /Link /Border [0 0 0] /H /I /C [1 0 0] /Rect [266.653 209.113 296.867 221.066] /A << /S /GoTo /D <6571756174696F6E2E332E332E312E32> >> >>}
\MusicPDFWriteObject{2584}{<< /Type /Annot /Subtype /Link /Border [0 0 0] /H /I /C [1 0 0] /Rect [354.343 209.113 383.992 221.066] /A << /S /GoTo /D <74686D2E332E332E312E31> >> >>}
\MusicPDFWriteObject{2585}{<< /Type /Annot /Subtype /Link /Border [0 0 0] /H /I /C [1 0 0] /Rect [143.433 82.596 173.646 94.549] /A << /S /GoTo /D <74686D2E332E332E312E35> >> >>}
\MusicPDFWriteObject{2586}{<< /Type /Annot /Subtype /Link /Border [0 0 0] /H /I /C [1 0 0] /Rect [488.522 630.721 518.735 642.674] /A << /S /GoTo /D <6571756174696F6E2E352E342E322E33> >> >>}
\MusicPDFWriteObject{2587}{<< /Type /Annot /Subtype /Link /Border [0 0 0] /H /I /C [1 0 0] /Rect [384.423 419.616 414.636 431.568] /A << /S /GoTo /D <74686D2E352E342E312E31> >> >>}
\MusicPDFWriteObject{2588}{<< /Type /Annot /Subtype /Link /Border [0 0 0] /H /I /C [1 0 0] /Rect [466.864 208.913 501.451 223.883] /A << /S /GoTo /D <74686D2E442E332E332E37> >> >>}
\MusicPDFWriteObject{2589}{<< /Type /Annot /Subtype /Link /Border [0 0 0] /H /I /C [1 0 0] /Rect [142.215 194.11 176.803 209.909] /A << /S /GoTo /D <74686D2E442E332E322E35> >> >>}
\MusicPDFWriteObject{2590}{<< /Type /Annot /Subtype /Link /Border [0 0 0] /H /I /C [1 0 0] /Rect [273.54 590.019 303.753 601.972] /A << /S /GoTo /D <6571756174696F6E2E352E342E322E36> >> >>}
\MusicPDFWriteObject{2591}{<< /Type /Annot /Subtype /Link /Border [0 0 0] /H /I /C [1 0 0] /Rect [132.894 515.054 163.671 525.845] /A << /S /GoTo /D <74686D2E352E342E322E32> >> >>}
\MusicPDFWriteObject{2592}{<< /Type /Annot /Subtype /Link /Border [0 0 0] /H /I /C [1 0 0] /Rect [315.87 515.054 346.648 525.845] /A << /S /GoTo /D <74686D2E352E342E312E31> >> >>}
\MusicPDFWriteObject{2593}{<< /Type /Annot /Subtype /Link /Border [0 0 0] /H /I /C [1 0 0] /Rect [197.92 390.17 228.359 400.961] /A << /S /GoTo /D <74686D2E352E342E322E31> >> >>}
\MusicPDFWriteObject{2594}{<< /Type /Annot /Subtype /Link /Border [0 0 0] /H /I /C [1 0 0] /Rect [291.547 390.17 329.655 400.961] /A << /S /GoTo /D <74686D2E442E332E332E3133> >> >>}
\MusicPDFWriteObject{2595}{<< /Type /Annot /Subtype /Link /Border [0 0 0] /H /I /C [1 0 0] /Rect [174.785 376.466 204.998 388.418] /A << /S /GoTo /D <74686D2E352E342E312E31> >> >>}
\MusicPDFWriteObject{2596}{<< /Type /Annot /Subtype /Link /Border [0 0 0] /H /I /C [0 1 0] /Rect [300.448 249.876 334.327 258.744] /A << /S /GoTo /D <636974652E5765696C2D61> >> >>}
\MusicPDFWriteObject{2597}{<< /Type /Annot /Subtype /Link /Border [0 0 0] /H /I /C [0 1 0] /Rect [188.829 582.478 222.708 591.346] /A << /S /GoTo /D <636974652E5765696C2D61> >> >>}
\MusicPDFWriteObject{2598}{<< /Type /Annot /Subtype /Link /Border [0 0 0] /H /I /C [0 1 0] /Rect [303.181 489.952 337.06 498.82] /A << /S /GoTo /D <636974652E5765696C2D61> >> >>}
\MusicPDFWriteObject{2599}{<< /Type /Annot /Subtype /Link /Border [0 0 0] /H /I /C [1 0 0] /Rect [512.381 264.44 543.159 278.1] /A << /S /GoTo /D <74686D2E332E342E332E33> >> >>}
\MusicPDFWriteObject{2600}{<< /Type /Annot /Subtype /Link /Border [0 0 0] /H /I /C [1 0 0] /Rect [112.107 225.8 141.868 238.133] /A << /S /GoTo /D <74686D2E332E342E332E33> >> >>}
\MusicPDFWriteObject{2601}{<< /Type /Annot /Subtype /Link /Border [0 0 0] /H /I /C [1 0 0] /Rect [223.418 225.8 255.358 238.133] /A << /S /GoTo /D <74686D2E432E352E322E36> >> >>}
\MusicPDFWriteObject{2602}{<< /Type /Annot /Subtype /Link /Border [0 0 0] /H /I /C [1 0 0] /Rect [317.257 225.8 349.197 238.133] /A << /S /GoTo /D <74686D2E432E352E322E34> >> >>}
\MusicPDFWriteObject{2603}{<< /Type /Annot /Subtype /Link /Border [0 0 0] /H /I /C [0 1 0] /Rect [524.054 215.5 543.159 224.368] /A << /S /GoTo /D <636974652E4E42> >> >>}
\MusicPDFWriteObject{2604}{<< /Type /Annot /Subtype /Link /Border [0 0 0] /H /I /C [0 1 0] /Rect [204.165 201.641 223.269 210.509] /A << /S /GoTo /D <636974652E4E42> >> >>}
\MusicPDFWriteObject{2605}{<< /Type /Annot /Subtype /Link /Border [0 0 0] /H /I /C [1 0 0] /Rect [73.712 158.99 103.925 170.943] /A << /S /GoTo /D <6571756174696F6E2E352E352E312E35> >> >>}
\MusicPDFWriteObject{2606}{<< /Type /Annot /Subtype /Link /Border [0 0 0] /H /I /C [1 0 0] /Rect [164.506 82.596 194.719 94.549] /A << /S /GoTo /D <6571756174696F6E2E332E332E312E33> >> >>}
\MusicPDFWriteObject{2607}{<< /Type /Annot /Subtype /Link /Border [0 0 0] /H /I /C [1 0 0] /Rect [505.121 82.596 535.335 94.549] /A << /S /GoTo /D <6571756174696F6E2E332E332E312E32> >> >>}
\MusicPDFWriteObject{2608}{<< /Type /Annot /Subtype /Link /Border [0 0 0] /H /I /C [1 0 0] /Rect [233.412 684.633 268.605 696.586] /A << /S /GoTo /D <6571756174696F6E2E332E342E332E3133> >> >>}
\MusicPDFWriteObject{2609}{<< /Type /Annot /Subtype /Link /Border [0 0 0] /H /I /C [1 0 0] /Rect [331.795 684.273 362.177 696.606] /A << /S /GoTo /D <74686D2E352E312E322E32> >> >>}
\MusicPDFWriteObject{2610}{<< /Type /Annot /Subtype /Link /Border [0 0 0] /H /I /C [0 1 0] /Rect [524.082 522.612 543.187 531.48] /A << /S /GoTo /D <636974652E4E42> >> >>}
\MusicPDFWriteObject{2611}{<< /Type /Annot /Subtype /Link /Border [0 0 0] /H /I /C [0 1 0] /Rect [334.729 509.462 353.834 518.329] /A << /S /GoTo /D <636974652E4E42> >> >>}
\MusicPDFWriteObject{2612}{<< /Type /Annot /Subtype /Link /Border [0 0 0] /H /I /C [0 1 0] /Rect [184.628 496.311 203.733 505.179] /A << /S /GoTo /D <636974652E4E42> >> >>}
\MusicPDFWriteObject{2613}{<< /Type /Annot /Subtype /Link /Border [0 0 0] /H /I /C [1 0 0] /Rect [148.271 465.909 179.048 479.492] /A << /S /GoTo /D <74686D2E332E342E332E33> >> >>}
\MusicPDFWriteObject{2614}{<< /Type /Annot /Subtype /Link /Border [0 0 0] /H /I /C [1 0 0] /Rect [247.437 467.519 282.63 479.472] /A << /S /GoTo /D <6571756174696F6E2E352E352E312E3133> >> >>}
\MusicPDFWriteObject{2615}{<< /Type /Annot /Subtype /Link /Border [0 0 0] /H /I /C [1 0 0] /Rect [294.445 454.368 329.638 466.321] /A << /S /GoTo /D <6571756174696F6E2E352E352E312E3133> >> >>}
\MusicPDFWriteObject{2616}{<< /Type /Annot /Subtype /Link /Border [0 0 0] /H /I /C [1 0 0] /Rect [73.712 441.218 103.925 453.17] /A << /S /GoTo /D <6571756174696F6E2E352E352E312E39> >> >>}
\MusicPDFWriteObject{2617}{<< /Type /Annot /Subtype /Link /Border [0 0 0] /H /I /C [1 0 0] /Rect [105.422 413.306 137.239 426.888] /A << /S /GoTo /D <74686D2E432E342E332E38> >> >>}
\MusicPDFWriteObject{2618}{<< /Type /Annot /Subtype /Link /Border [0 0 0] /H /I /C [0 1 0] /Rect [352.751 417.406 371.856 426.274] /A << /S /GoTo /D <636974652E4E42> >> >>}
\MusicPDFWriteObject{2619}{<< /Type /Annot /Subtype /Link /Border [0 0 0] /H /I /C [1 0 0] /Rect [71.004 401.405 103.43 413.72] /A << /S /GoTo /D <74686D2E432E352E322E36> >> >>}
\MusicPDFWriteObject{2620}{<< /Type /Annot /Subtype /Link /Border [0 0 0] /H /I /C [1 0 0] /Rect [166.207 401.405 198.633 413.72] /A << /S /GoTo /D <74686D2E432E352E322E34> >> >>}
\MusicPDFWriteObject{2621}{<< /Type /Annot /Subtype /Link /Border [0 0 0] /H /I /C [1 0 0] /Rect [236.265 372.433 266.478 383.224] /A << /S /GoTo /D <74686D2E352E352E312E32> >> >>}
\MusicPDFWriteObject{2622}{<< /Type /Annot /Subtype /Link /Border [0 0 0] /H /I /C [0 1 0] /Rect [303.181 284.44 337.06 293.308] /A << /S /GoTo /D <636974652E5765696C2D61> >> >>}
\MusicPDFWriteObject{2623}{<< /Type /Annot /Subtype /Link /Border [0 0 0] /H /I /C [0 1 0] /Rect [146.629 687.123 165.734 695.991] /A << /S /GoTo /D <636974652E4E42> >> >>}
\MusicPDFWriteObject{2624}{<< /Type /Annot /Subtype /Link /Border [0 0 0] /H /I /C [0 1 0] /Rect [453.909 687.123 473.014 695.991] /A << /S /GoTo /D <636974652E4E42> >> >>}
\MusicPDFWriteObject{2625}{<< /Type /Annot /Subtype /Link /Border [0 0 0] /H /I /C [0 1 0] /Rect [316.542 673.973 335.646 682.84] /A << /S /GoTo /D <636974652E4E42> >> >>}
\MusicPDFWriteObject{2626}{<< /Type /Annot /Subtype /Link /Border [0 0 0] /H /I /C [1 0 0] /Rect [169.027 627.289 198.676 641.391] /A << /S /GoTo /D <74686D2E352E322E312E31> >> >>}
\MusicPDFWriteObject{2627}{<< /Type /Annot /Subtype /Link /Border [0 0 0] /H /I /C [1 0 0] /Rect [257.036 627.289 286.685 641.391] /A << /S /GoTo /D <74686D2E322E312E342E31> >> >>}
\MusicPDFWriteObject{2628}{<< /Type /Annot /Subtype /Link /Border [0 0 0] /H /I /C [1 0 0] /Rect [377.926 627.289 409.744 641.391] /A << /S /GoTo /D <74686D2E432E312E362E31> >> >>}
\MusicPDFWriteObject{2629}{<< /Type /Annot /Subtype /Link /Border [0 0 0] /H /I /C [1 0 0] /Rect [494.527 587.73 524.176 602.586] /A << /S /GoTo /D <74686D2E352E352E312E32> >> >>}
\MusicPDFWriteObject{2630}{<< /Type /Annot /Subtype /Link /Border [0 0 0] /H /I /C [1 0 0] /Rect [270.302 570.228 300.317 582.18] /A << /S /GoTo /D <74686D2E352E352E312E32> >> >>}
\MusicPDFWriteObject{2631}{<< /Type /Annot /Subtype /Link /Border [0 0 0] /H /I /C [1 0 0] /Rect [71.004 501.447 100.709 516.303] /A << /S /GoTo /D <74686D2E352E352E332E32> >> >>}
\MusicPDFWriteObject{2632}{<< /Type /Annot /Subtype /Link /Border [0 0 0] /H /I /C [1 0 0] /Rect [112.872 490.09 143.198 502.043] /A << /S /GoTo /D <74686D2E352E342E312E31> >> >>}
\MusicPDFWriteObject{2633}{<< /Type /Annot /Subtype /Link /Border [0 0 0] /H /I /C [0 1 0] /Rect [122.604 479.43 156.573 488.297] /A << /S /GoTo /D <636974652E436E772D61> >> >>}
\MusicPDFWriteObject{2634}{<< /Type /Annot /Subtype /Link /Border [0 0 0] /H /I /C [1 0 0] /Rect [388.051 476.94 417.7 488.892] /A << /S /GoTo /D <74686D2E352E352E312E33> >> >>}
\MusicPDFWriteObject{2635}{<< /Type /Annot /Subtype /Link /Border [0 0 0] /H /I /C [1 0 0] /Rect [146.374 405.04 181.567 416.993] /A << /S /GoTo /D <6571756174696F6E2E322E322E342E3233> >> >>}
\MusicPDFWriteObject{2636}{<< /Type /Annot /Subtype /Link /Border [0 0 0] /H /I /C [1 0 0] /Rect [211.807 212.868 242.02 223.658] /A << /S /GoTo /D <74686D2E352E352E312E33> >> >>}
\MusicPDFWriteObject{2637}{<< /Type /Annot /Subtype /Link /Border [0 0 0] /H /I /C [1 0 0] /Rect [338.382 148.349 369.159 160.304] /A << /S /GoTo /D <74686D2E332E322E322E31> >> >>}
\MusicPDFWriteObject{2638}{<< /Type /Annot /Subtype /Link /Border [0 0 0] /H /I /C [1 0 0] /Rect [229.126 135.198 259.537 147.153] /A << /S /GoTo /D <74686D2E352E312E322E31> >> >>}
\MusicPDFWriteObject{2639}{<< /Type /Annot /Subtype /Link /Border [0 0 0] /H /I /C [1 0 0] /Rect [364.11 645.18 393.758 657.136] /A << /S /GoTo /D <74686D2E352E352E322E31> >> >>}
\MusicPDFWriteObject{2640}{<< /Type /Annot /Subtype /Link /Border [0 0 0] /H /I /C [1 0 0] /Rect [287.926 631.869 318.139 646.177] /A << /S /GoTo /D <74686D2E352E312E322E31> >> >>}
\MusicPDFWriteObject{2641}{<< /Type /Annot /Subtype /Link /Border [0 0 0] /H /I /C [1 0 0] /Rect [285.293 618.718 315.224 630.674] /A << /S /GoTo /D <74686D2E322E322E332E34> >> >>}
\MusicPDFWriteObject{2642}{<< /Type /Annot /Subtype /Link /Border [0 0 0] /H /I /C [1 0 0] /Rect [73.712 605.568 103.925 617.521] /A << /S /GoTo /D <6571756174696F6E2E322E322E332E35> >> >>}
\MusicPDFWriteObject{2643}{<< /Type /Annot /Subtype /Link /Border [0 0 0] /H /I /C [1 0 0] /Rect [174.972 605.568 205.298 617.521] /A << /S /GoTo /D <74686D2E322E322E332E34> >> >>}
\MusicPDFWriteObject{2644}{<< /Type /Annot /Subtype /Link /Border [0 0 0] /H /I /C [1 0 0] /Rect [449.125 530.562 479.902 542.515] /A << /S /GoTo /D <74686D2E352E352E332E32> >> >>}
\MusicPDFWriteObject{2645}{<< /Type /Annot /Subtype /Link /Border [0 0 0] /H /I /C [1 0 0] /Rect [365.986 477.959 395.72 489.912] /A << /S /GoTo /D <74686D2E352E352E312E32> >> >>}
\MusicPDFWriteObject{2646}{<< /Type /Annot /Subtype /Link /Border [0 0 0] /H /I /C [1 0 0] /Rect [291.749 464.808 321.962 476.761] /A << /S /GoTo /D <74686D2E352E352E332E32> >> >>}
\MusicPDFWriteObject{2647}{<< /Type /Annot /Subtype /Link /Border [0 0 0] /H /I /C [0 1 0] /Rect [180.422 337.016 199.526 345.884] /A << /S /GoTo /D <636974652E4E42> >> >>}
\MusicPDFWriteObject{2648}{<< /Type /Annot /Subtype /Link /Border [0 0 0] /H /I /C [0 1 0] /Rect [217.138 323.866 236.243 332.733] /A << /S /GoTo /D <636974652E4E42> >> >>}
\MusicPDFWriteObject{2649}{<< /Type /Annot /Subtype /Link /Border [0 0 0] /H /I /C [0 1 0] /Rect [342.392 297.564 360.59 306.432] /A << /S /GoTo /D <636974652E4A6172> >> >>}
\MusicPDFWriteObject{2650}{<< /Type /Annot /Subtype /Link /Border [0 0 0] /H /I /C [1 0 0] /Rect [151.816 180.628 184.851 192.583] /A << /S /GoTo /D <74686D2E432E342E332E35> >> >>}
\MusicPDFWriteObject{2651}{<< /Type /Annot /Subtype /Link /Border [0 0 0] /H /I /C [1 0 0] /Rect [141.52 154.881 179.635 165.671] /A << /S /GoTo /D <74686D2E432E342E332E3132> >> >>}
\MusicPDFWriteObject{2652}{<< /Type /Annot /Subtype /Link /Border [0 0 0] /H /I /C [1 0 0] /Rect [220.708 93.953 251.034 108.808] /A << /S /GoTo /D <74686D2E352E352E332E31> >> >>}
\MusicPDFWriteObject{2653}{<< /Type /Annot /Subtype /Link /Border [0 0 0] /H /I /C [1 0 0] /Rect [313.965 82.596 343.614 94.549] /A << /S /GoTo /D <74686D2E352E342E312E31> >> >>}
\MusicPDFWriteObject{2654}{<< /Type /Annot /Subtype /Link /Border [0 0 0] /H /I /C [1 0 0] /Rect [71.004 684.633 100.652 696.586] /A << /S /GoTo /D <74686D2E352E322E312E31> >> >>}
\MusicPDFWriteObject{2655}{<< /Type /Annot /Subtype /Link /Border [0 0 0] /H /I /C [0 1 0] /Rect [210.063 490.207 228.262 499.075] /A << /S /GoTo /D <636974652E4A6172> >> >>}
\MusicPDFWriteObject{2656}{<< /Type /Annot /Subtype /Link /Border [0 0 0] /H /I /C [0 1 0] /Rect [411.267 429.684 429.466 438.552] /A << /S /GoTo /D <636974652E4A6172> >> >>}
\MusicPDFWriteObject{2657}{<< /Type /Annot /Subtype /Link /Border [0 0 0] /H /I /C [0 1 0] /Rect [342.862 390.232 361.967 399.099] /A << /S /GoTo /D <636974652E4E42> >> >>}
\MusicPDFWriteObject{2658}{<< /Type /Annot /Subtype /Link /Border [0 0 0] /H /I /C [1 0 0] /Rect [158.682 359.996 188.895 371.949] /A << /S /GoTo /D <74686D2E352E342E312E31> >> >>}
\MusicPDFWriteObject{2659}{<< /Type /Annot /Subtype /Link /Border [0 0 0] /H /I /C [1 0 0] /Rect [149.468 181.22 180.245 193.535] /A << /S /GoTo /D <74686D2E352E322E312E31> >> >>}
\MusicPDFWriteObject{2660}{<< /Type /Annot /Subtype /Link /Border [0 0 0] /H /I /C [0 1 0] /Rect [499.735 640.724 517.934 649.592] /A << /S /GoTo /D <636974652E4A6172> >> >>}
\MusicPDFWriteObject{2661}{<< /Type /Annot /Subtype /Link /Border [0 0 0] /H /I /C [0 1 0] /Rect [110.295 601.272 144.264 610.14] /A << /S /GoTo /D <636974652E436E772D61> >> >>}
\MusicPDFWriteObject{2662}{<< /Type /Annot /Subtype /Link /Border [0 0 0] /H /I /C [0 1 0] /Rect [509.19 271.007 543.159 279.875] /A << /S /GoTo /D <636974652E436E772D61> >> >>}
\MusicPDFWriteObject{2663}{<< /Type /Annot /Subtype /Link /Border [0 0 0] /H /I /C [0 1 0] /Rect [284.659 257.856 318.628 266.724] /A << /S /GoTo /D <636974652E436E772D61> >> >>}
\MusicPDFWriteObject{2664}{<< /Type /Annot /Subtype /Link /Border [0 0 0] /H /I /C [0 1 0] /Rect [291.442 165.388 325.411 174.256] /A << /S /GoTo /D <636974652E436E772D61> >> >>}
\MusicPDFWriteObject{2665}{<< /Type /Annot /Subtype /Link /Border [0 0 0] /H /I /C [1 0 0] /Rect [206.122 632.03 236.335 643.983] /A << /S /GoTo /D <6571756174696F6E2E352E352E332E36> >> >>}
\MusicPDFWriteObject{2666}{<< /Type /Annot /Subtype /Link /Border [0 0 0] /H /I /C [1 0 0] /Rect [505.641 447.739 535.854 459.692] /A << /S /GoTo /D <6571756174696F6E2E352E352E332E33> >> >>}
\MusicPDFWriteObject{2667}{<< /Type /Annot /Subtype /Link /Border [0 0 0] /H /I /C [1 0 0] /Rect [99.269 408.287 129.482 420.24] /A << /S /GoTo /D <6571756174696F6E2E332E312E322E35> >> >>}
\MusicPDFWriteObject{2668}{<< /Type /Annot /Subtype /Link /Border [0 0 0] /H /I /C [1 0 0] /Rect [502.7 381.985 532.914 393.938] /A << /S /GoTo /D <6571756174696F6E2E352E352E332E33> >> >>}
\MusicPDFWriteObject{2669}{<< /Type /Annot /Subtype /Link /Border [0 0 0] /H /I /C [1 0 0] /Rect [371.33 329.327 401.544 341.28] /A << /S /GoTo /D <6571756174696F6E2E352E352E312E36> >> >>}
\MusicPDFWriteObject{2670}{<< /Type /Annot /Subtype /Link /Border [0 0 0] /H /I /C [1 0 0] /Rect [460.17 328.967 490.383 341.282] /A << /S /GoTo /D <74686D2E352E352E312E32> >> >>}
\MusicPDFWriteObject{2671}{<< /Type /Annot /Subtype /Link /Border [0 0 0] /H /I /C [1 0 0] /Rect [222.557 303.026 252.77 314.978] /A << /S /GoTo /D <6571756174696F6E2E352E322E332E39> >> >>}
\MusicPDFWriteObject{2672}{<< /Type /Annot /Subtype /Link /Border [0 0 0] /H /I /C [0 1 0] /Rect [429.26 95.07 448.365 103.938] /A << /S /GoTo /D <636974652E4E42> >> >>}
\MusicPDFWriteObject{2673}{<< /Type /Annot /Subtype /Link /Border [0 0 0] /H /I /C [1 0 0] /Rect [308.404 668.98 339.182 683.836] /A << /S /GoTo /D <74686D2E352E322E312E31> >> >>}
\MusicPDFWriteObject{2674}{<< /Type /Annot /Subtype /Link /Border [0 0 0] /H /I /C [0 1 0] /Rect [225.966 645.546 244.164 654.414] /A << /S /GoTo /D <636974652E4A6172> >> >>}
\MusicPDFWriteObject{2675}{<< /Type /Annot /Subtype /Link /Border [0 0 0] /H /I /C [1 0 0] /Rect [331.582 623.883 361.231 634.674] /A << /S /GoTo /D <74686D2E352E352E332E32> >> >>}
\MusicPDFWriteObject{2676}{<< /Type /Annot /Subtype /Link /Border [0 0 0] /H /I /C [1 0 0] /Rect [362.115 597.028 392.328 608.981] /A << /S /GoTo /D <6571756174696F6E2E352E312E312E32> >> >>}
\MusicPDFWriteObject{2677}{<< /Type /Annot /Subtype /Link /Border [0 0 0] /H /I /C [0 1 0] /Rect [412.602 533.056 431.706 541.924] /A << /S /GoTo /D <636974652E4E42> >> >>}
\MusicPDFWriteObject{2678}{<< /Type /Annot /Subtype /Link /Border [0 0 0] /H /I /C [0 1 0] /Rect [107.178 519.197 125.376 528.065] /A << /S /GoTo /D <636974652E4A6172> >> >>}
\MusicPDFWriteObject{2679}{<< /Type /Annot /Subtype /Link /Border [0 0 0] /H /I /C [1 0 0] /Rect [500.711 514.913 530.642 529.769] /A << /S /GoTo /D <74686D2E352E322E312E31> >> >>}
\MusicPDFWriteObject{2680}{<< /Type /Annot /Subtype /Link /Border [0 0 0] /H /I /C [1 0 0] /Rect [131.018 484.384 161.231 495.174] /A << /S /GoTo /D <74686D2E352E352E332E32> >> >>}
\MusicPDFWriteObject{2681}{<< /Type /Annot /Subtype /Link /Border [0 0 0] /H /I /C [1 0 0] /Rect [298.948 348.533 329.725 361.379] /A << /S /GoTo /D <74686D2E342E332E332E31> >> >>}
\MusicPDFWriteObject{2682}{<< /Type /Annot /Subtype /Link /Border [0 0 0] /H /I /C [1 0 0] /Rect [364.354 333.948 395.132 348.051] /A << /S /GoTo /D <74686D2E352E322E332E31> >> >>}
\MusicPDFWriteObject{2683}{<< /Type /Annot /Subtype /Link /Border [0 0 0] /H /I /C [0 1 0] /Rect [518.099 338.126 543.159 346.994] /A << /S /GoTo /D <636974652E53636861> >> >>}
\MusicPDFWriteObject{2684}{<< /Type /Annot /Subtype /Link /Border [0 0 0] /H /I /C [1 0 0] /Rect [514.169 322.099 543.817 334.945] /A << /S /GoTo /D <74686D2E352E312E332E31> >> >>}
\MusicPDFWriteObject{2685}{<< /Type /Annot /Subtype /Link /Border [0 0 0] /H /I /C [1 0 0] /Rect [169.027 307.514 198.676 321.617] /A << /S /GoTo /D <74686D2E352E322E312E31> >> >>}
\MusicPDFWriteObject{2686}{<< /Type /Annot /Subtype /Link /Border [0 0 0] /H /I /C [1 0 0] /Rect [257.036 307.514 286.685 321.617] /A << /S /GoTo /D <74686D2E322E312E342E31> >> >>}
\MusicPDFWriteObject{2687}{<< /Type /Annot /Subtype /Link /Border [0 0 0] /H /I /C [1 0 0] /Rect [377.926 307.514 409.744 321.617] /A << /S /GoTo /D <74686D2E432E312E362E31> >> >>}
\MusicPDFWriteObject{2688}{<< /Type /Annot /Subtype /Link /Border [0 0 0] /H /I /C [1 0 0] /Rect [494.527 267.956 524.176 282.811] /A << /S /GoTo /D <74686D2E352E352E312E32> >> >>}
\MusicPDFWriteObject{2689}{<< /Type /Annot /Subtype /Link /Border [0 0 0] /H /I /C [1 0 0] /Rect [82.959 558.462 294.475 569.252] /A << /S /GoTo /D <73656374696F6E2E362E31> >> >>}
\MusicPDFWriteObject{2690}{<< /Type /Annot /Subtype /Link /Border [0 0 0] /H /I /C [1 0 0] /Rect [82.959 544.397 238.77 556.102] /A << /S /GoTo /D <73756273656374696F6E2E362E312E31> >> >>}
\MusicPDFWriteObject{2691}{<< /Type /Annot /Subtype /Link /Border [0 0 0] /H /I /C [1 0 0] /Rect [82.959 531.247 344.956 542.951] /A << /S /GoTo /D <73756273656374696F6E2E362E312E32> >> >>}
\MusicPDFWriteObject{2692}{<< /Type /Annot /Subtype /Link /Border [0 0 0] /H /I /C [1 0 0] /Rect [82.959 519.009 310.467 529.8] /A << /S /GoTo /D <73656374696F6E2E362E32> >> >>}
\MusicPDFWriteObject{2693}{<< /Type /Annot /Subtype /Link /Border [0 0 0] /H /I /C [1 0 0] /Rect [82.959 504.945 279.717 516.649] /A << /S /GoTo /D <73756273656374696F6E2E362E322E31> >> >>}
\MusicPDFWriteObject{2694}{<< /Type /Annot /Subtype /Link /Border [0 0 0] /H /I /C [1 0 0] /Rect [82.959 490.361 353.842 505.216] /A << /S /GoTo /D <73756273656374696F6E2E362E322E32> >> >>}
\MusicPDFWriteObject{2695}{<< /Type /Annot /Subtype /Link /Border [0 0 0] /H /I /C [1 0 0] /Rect [82.959 481.494 274.016 490.348] /A << /S /GoTo /D <73656374696F6E2E362E33> >> >>}
\MusicPDFWriteObject{2696}{<< /Type /Annot /Subtype /Link /Border [0 0 0] /H /I /C [1 0 0] /Rect [82.959 468.343 358.112 477.197] /A << /S /GoTo /D <73756273656374696F6E2E362E332E31> >> >>}
\MusicPDFWriteObject{2697}{<< /Type /Annot /Subtype /Link /Border [0 0 0] /H /I /C [1 0 0] /Rect [82.959 455.192 203.257 464.046] /A << /S /GoTo /D <73756273656374696F6E2E362E332E32> >> >>}
\MusicPDFWriteObject{2698}{<< /Type /Annot /Subtype /Link /Border [0 0 0] /H /I /C [1 0 0] /Rect [460.912 336.673 491.69 348.377] /A << /S /GoTo /D <74686D2E352E322E332E31> >> >>}
\MusicPDFWriteObject{2699}{<< /Type /Annot /Subtype /Link /Border [0 0 0] /H /I /C [1 0 0] /Rect [499.044 324.435 506.016 335.226] /A << /S /GoTo /D <636861707465722E35> >> >>}
\MusicPDFWriteObject{2700}{<< /Type /Annot /Subtype /Link /Border [0 0 0] /H /I /C [1 0 0] /Rect [71.004 641.785 101.217 656.17] /A << /S /GoTo /D <74686D2E352E322E332E31> >> >>}
\MusicPDFWriteObject{2701}{<< /Type /Annot /Subtype /Link /Border [0 0 0] /H /I /C [1 0 0] /Rect [402.799 592.608 433.012 604.938] /A << /S /GoTo /D <74686D2E322E322E322E31> >> >>}
\MusicPDFWriteObject{2702}{<< /Type /Annot /Subtype /Link /Border [0 0 0] /H /I /C [1 0 0] /Rect [315.648 546.625 345.861 558.578] /A << /S /GoTo /D <6571756174696F6E2E332E312E322E31> >> >>}
\MusicPDFWriteObject{2703}{<< /Type /Annot /Subtype /Link /Border [0 0 0] /H /I /C [1 0 0] /Rect [345.516 353.957 375.729 365.909] /A << /S /GoTo /D <6571756174696F6E2E352E322E312E31> >> >>}
\MusicPDFWriteObject{2704}{<< /Type /Annot /Subtype /Link /Border [0 0 0] /H /I /C [1 0 0] /Rect [91.42 313.907 121.633 325.86] /A << /S /GoTo /D <6571756174696F6E2E352E322E312E32> >> >>}
\MusicPDFWriteObject{2705}{<< /Type /Annot /Subtype /Link /Border [0 0 0] /H /I /C [1 0 0] /Rect [175.777 621.904 205.99 633.857] /A << /S /GoTo /D <6571756174696F6E2E362E312E322E35> >> >>}
\MusicPDFWriteObject{2706}{<< /Type /Annot /Subtype /Link /Border [0 0 0] /H /I /C [1 0 0] /Rect [270.518 497.272 300.731 509.865] /A << /S /GoTo /D <74686D2E362E312E322E32> >> >>}
\MusicPDFWriteObject{2707}{<< /Type /Annot /Subtype /Link /Border [0 0 0] /H /I /C [0 1 0] /Rect [390.106 137.689 409.079 146.557] /A << /S /GoTo /D <636974652E547265> >> >>}
\MusicPDFWriteObject{2708}{<< /Type /Annot /Subtype /Link /Border [0 0 0] /H /I /C [0 1 0] /Rect [73.77 124.538 98.831 133.406] /A << /S /GoTo /D <636974652E53636861> >> >>}
\MusicPDFWriteObject{2709}{<< /Type /Annot /Subtype /Link /Border [0 0 0] /H /I /C [0 1 0] /Rect [236.53 124.538 261.59 133.406] /A << /S /GoTo /D <636974652E53636861> >> >>}
\MusicPDFWriteObject{2710}{<< /Type /Annot /Subtype /Link /Border [0 0 0] /H /I /C [0 1 0] /Rect [107.145 111.387 126.118 120.255] /A << /S /GoTo /D <636974652E547265> >> >>}
\MusicPDFWriteObject{2711}{<< /Type /Annot /Subtype /Link /Border [0 0 0] /H /I /C [0 1 0] /Rect [120.86 85.086 163.883 93.954] /A << /S /GoTo /D <636974652E426F752D4754> >> >>}
\MusicPDFWriteObject{2712}{<< /Type /Annot /Subtype /Link /Border [0 0 0] /H /I /C [1 0 0] /Rect [197.665 684.633 227.879 696.586] /A << /S /GoTo /D <74686D2E322E312E342E32> >> >>}
\MusicPDFWriteObject{2713}{<< /Type /Annot /Subtype /Link /Border [0 0 0] /H /I /C [1 0 0] /Rect [364.828 540.152 395.041 552.105] /A << /S /GoTo /D <6571756174696F6E2E362E312E322E36> >> >>}
\MusicPDFWriteObject{2714}{<< /Type /Annot /Subtype /Link /Border [0 0 0] /H /I /C [1 0 0] /Rect [169.049 444.386 198.698 456.7] /A << /S /GoTo /D <74686D2E322E312E342E31> >> >>}
\MusicPDFWriteObject{2715}{<< /Type /Annot /Subtype /Link /Border [0 0 0] /H /I /C [0 1 0] /Rect [336.916 420.934 355.114 429.802] /A << /S /GoTo /D <636974652E4A6172> >> >>}
\MusicPDFWriteObject{2716}{<< /Type /Annot /Subtype /Link /Border [0 0 0] /H /I /C [1 0 0] /Rect [221.921 392.143 252.134 404.096] /A << /S /GoTo /D <6571756174696F6E2E362E312E322E35> >> >>}
\MusicPDFWriteObject{2717}{<< /Type /Annot /Subtype /Link /Border [0 0 0] /H /I /C [1 0 0] /Rect [329.178 365.841 359.391 377.794] /A << /S /GoTo /D <6571756174696F6E2E362E312E322E37> >> >>}
\MusicPDFWriteObject{2718}{<< /Type /Annot /Subtype /Link /Border [0 0 0] /H /I /C [1 0 0] /Rect [509.103 326.389 539.316 338.342] /A << /S /GoTo /D <6571756174696F6E2E362E312E322E37> >> >>}
\MusicPDFWriteObject{2719}{<< /Type /Annot /Subtype /Link /Border [0 0 0] /H /I /C [0 1 0] /Rect [248.71 315.728 266.909 324.596] /A << /S /GoTo /D <636974652E4A6172> >> >>}
\MusicPDFWriteObject{2720}{<< /Type /Annot /Subtype /Link /Border [0 0 0] /H /I /C [1 0 0] /Rect [263.193 124.504 298.387 136.456] /A << /S /GoTo /D <6571756174696F6E2E362E312E322E3130> >> >>}
\MusicPDFWriteObject{2721}{<< /Type /Annot /Subtype /Link /Border [0 0 0] /H /I /C [0 1 0] /Rect [73.77 113.843 92.073 122.711] /A << /S /GoTo /D <636974652E4445> >> >>}
\MusicPDFWriteObject{2722}{<< /Type /Annot /Subtype /Link /Border [0 0 0] /H /I /C [1 0 0] /Rect [183.092 577.384 216.409 593.562] /A << /S /GoTo /D <74686D2E412E332E322E32> >> >>}
\MusicPDFWriteObject{2723}{<< /Type /Annot /Subtype /Link /Border [0 0 0] /H /I /C [1 0 0] /Rect [275.416 565.319 310.609 577.272] /A << /S /GoTo /D <6571756174696F6E2E362E312E322E3131> >> >>}
\MusicPDFWriteObject{2724}{<< /Type /Annot /Subtype /Link /Border [0 0 0] /H /I /C [0 1 0] /Rect [461.799 493.539 480.101 502.407] /A << /S /GoTo /D <636974652E4445> >> >>}
\MusicPDFWriteObject{2725}{<< /Type /Annot /Subtype /Link /Border [0 0 0] /H /I /C [0 1 0] /Rect [300.063 433.277 318.365 442.145] /A << /S /GoTo /D <636974652E4445> >> >>}
\MusicPDFWriteObject{2726}{<< /Type /Annot /Subtype /Link /Border [0 0 0] /H /I /C [1 0 0] /Rect [318.741 197.539 348.503 208.329] /A << /S /GoTo /D <74686D2E322E312E342E32> >> >>}
\MusicPDFWriteObject{2727}{<< /Type /Annot /Subtype /Link /Border [0 0 0] /H /I /C [1 0 0] /Rect [258.257 184.388 288.611 195.179] /A << /S /GoTo /D <74686D2E322E312E342E32> >> >>}
\MusicPDFWriteObject{2728}{<< /Type /Annot /Subtype /Link /Border [0 0 0] /H /I /C [0 1 0] /Rect [216.219 373.54 234.522 382.407] /A << /S /GoTo /D <636974652E4445> >> >>}
\MusicPDFWriteObject{2729}{<< /Type /Annot /Subtype /Link /Border [0 0 0] /H /I /C [1 0 0] /Rect [511.593 236.568 543.817 248.523] /A << /S /GoTo /D <74686D2E442E312E322E31> >> >>}
\MusicPDFWriteObject{2730}{<< /Type /Annot /Subtype /Link /Border [0 0 0] /H /I /C [1 0 0] /Rect [148.692 124.537 179.47 137.479] /A << /S /GoTo /D <74686D2E332E342E332E33> >> >>}
\MusicPDFWriteObject{2731}{<< /Type /Annot /Subtype /Link /Border [0 0 0] /H /I /C [1 0 0] /Rect [464.744 684.633 495.521 697.574] /A << /S /GoTo /D <74686D2E352E352E332E32> >> >>}
\MusicPDFWriteObject{2732}{<< /Type /Annot /Subtype /Link /Border [0 0 0] /H /I /C [0 1 0] /Rect [457.654 660.822 476.759 669.69] /A << /S /GoTo /D <636974652E4E42> >> >>}
\MusicPDFWriteObject{2733}{<< /Type /Annot /Subtype /Link /Border [0 0 0] /H /I /C [0 1 0] /Rect [259.1 647.671 278.204 656.539] /A << /S /GoTo /D <636974652E4E42> >> >>}
\MusicPDFWriteObject{2734}{<< /Type /Annot /Subtype /Link /Border [0 0 0] /H /I /C [1 0 0] /Rect [218.868 502.396 249.646 515.697] /A << /S /GoTo /D <74686D2E352E352E332E32> >> >>}
\MusicPDFWriteObject{2735}{<< /Type /Annot /Subtype /Link /Border [0 0 0] /H /I /C [0 1 0] /Rect [438.924 462.8 458.029 471.667] /A << /S /GoTo /D <636974652E4E42> >> >>}
\MusicPDFWriteObject{2736}{<< /Type /Annot /Subtype /Link /Border [0 0 0] /H /I /C [0 1 0] /Rect [169.994 436.498 189.099 445.366] /A << /S /GoTo /D <636974652E4E42> >> >>}
\MusicPDFWriteObject{2737}{<< /Type /Annot /Subtype /Link /Border [0 0 0] /H /I /C [0 1 0] /Rect [332.4 384.52 364.024 393.388] /A << /S /GoTo /D <636974652E47726F2D64> >> >>}
\MusicPDFWriteObject{2738}{<< /Type /Annot /Subtype /Link /Border [0 0 0] /H /I /C [1 0 0] /Rect [187.794 295.206 221.355 304.06] /A << /S /GoTo /D <74686D2E362E312E322E31> >> >>}
\MusicPDFWriteObject{2739}{<< /Type /Annot /Subtype /Link /Border [0 0 0] /H /I /C [1 0 0] /Rect [140.121 208.995 170.334 223.85] /A << /S /GoTo /D <74686D2E362E312E322E32> >> >>}
\MusicPDFWriteObject{2740}{<< /Type /Annot /Subtype /Link /Border [0 0 0] /H /I /C [1 0 0] /Rect [324.243 105.584 359.436 117.537] /A << /S /GoTo /D <6571756174696F6E2E362E312E322E3138> >> >>}
\MusicPDFWriteObject{2741}{<< /Type /Annot /Subtype /Link /Border [0 0 0] /H /I /C [1 0 0] /Rect [122.026 684.273 152.126 696.586] /A << /S /GoTo /D <74686D2E362E312E322E34> >> >>}
\MusicPDFWriteObject{2742}{<< /Type /Annot /Subtype /Link /Border [0 0 0] /H /I /C [1 0 0] /Rect [277.033 684.273 307.134 696.586] /A << /S /GoTo /D <74686D2E322E312E342E32> >> >>}
\MusicPDFWriteObject{2743}{<< /Type /Annot /Subtype /Link /Border [0 0 0] /H /I /C [1 0 0] /Rect [287.279 634.457 322.472 646.41] /A << /S /GoTo /D <6571756174696F6E2E362E312E322E3136> >> >>}
\MusicPDFWriteObject{2744}{<< /Type /Annot /Subtype /Link /Border [0 0 0] /H /I /C [1 0 0] /Rect [374.569 634.457 404.782 646.41] /A << /S /GoTo /D <74686D2E362E312E322E34> >> >>}
\MusicPDFWriteObject{2745}{<< /Type /Annot /Subtype /Link /Border [0 0 0] /H /I /C [1 0 0] /Rect [312.275 478.374 347.468 490.326] /A << /S /GoTo /D <6571756174696F6E2E362E312E322E3230> >> >>}
\MusicPDFWriteObject{2746}{<< /Type /Annot /Subtype /Link /Border [0 0 0] /H /I /C [1 0 0] /Rect [258.116 465.223 293.309 477.176] /A << /S /GoTo /D <6571756174696F6E2E362E312E322E3139> >> >>}
\MusicPDFWriteObject{2747}{<< /Type /Annot /Subtype /Link /Border [0 0 0] /H /I /C [1 0 0] /Rect [408.404 464.863 438.617 477.176] /A << /S /GoTo /D <74686D2E322E312E342E32> >> >>}
\MusicPDFWriteObject{2748}{<< /Type /Annot /Subtype /Link /Border [0 0 0] /H /I /C [1 0 0] /Rect [73.712 349.34 108.905 361.292] /A << /S /GoTo /D <6571756174696F6E2E362E312E322E3231> >> >>}
\MusicPDFWriteObject{2749}{<< /Type /Annot /Subtype /Link /Border [0 0 0] /H /I /C [1 0 0] /Rect [411.188 211.432 440.836 227.61] /A << /S /GoTo /D <74686D2E342E332E332E31> >> >>}
\MusicPDFWriteObject{2750}{<< /Type /Annot /Subtype /Link /Border [0 0 0] /H /I /C [1 0 0] /Rect [335.448 197.573 365.351 212.428] /A << /S /GoTo /D <74686D2E352E322E312E31> >> >>}
\MusicPDFWriteObject{2751}{<< /Type /Annot /Subtype /Link /Border [0 0 0] /H /I /C [1 0 0] /Rect [406.06 143.404 412.023 157.59] /A << /S /GoTo /D <48666F6F746E6F74652E3232> >> >>}
\MusicPDFWriteObject{2752}{<< /Type /Annot /Subtype /Link /Border [0 0 0] /H /I /C [0 1 0] /Rect [73.715 687.123 98.775 695.991] /A << /S /GoTo /D <636974652E53636861> >> >>}
\MusicPDFWriteObject{2753}{<< /Type /Annot /Subtype /Link /Border [0 0 0] /H /I /C [1 0 0] /Rect [452.363 670.114 458.326 684.3] /A << /S /GoTo /D <48666F6F746E6F74652E3233> >> >>}
\MusicPDFWriteObject{2754}{<< /Type /Annot /Subtype /Link /Border [0 0 0] /H /I /C [1 0 0] /Rect [526.24 560.858 532.203 575.354] /A << /S /GoTo /D <48666F6F746E6F74652E3234> >> >>}
\MusicPDFWriteObject{2755}{<< /Type /Annot /Subtype /Link /Border [0 0 0] /H /I /C [1 0 0] /Rect [257.503 551.632 292.696 560.486] /A << /S /GoTo /D <74686D2E322E312E322E3130> >> >>}
\MusicPDFWriteObject{2756}{<< /Type /Annot /Subtype /Link /Border [0 0 0] /H /I /C [1 0 0] /Rect [212.555 536.544 242.768 547.335] /A << /S /GoTo /D <74686D2E342E332E332E31> >> >>}
\MusicPDFWriteObject{2757}{<< /Type /Annot /Subtype /Link /Border [0 0 0] /H /I /C [1 0 0] /Rect [143.101 521.046 173.879 537.224] /A << /S /GoTo /D <74686D2E342E332E332E31> >> >>}
\MusicPDFWriteObject{2758}{<< /Type /Annot /Subtype /Link /Border [0 0 0] /H /I /C [1 0 0] /Rect [465.152 356.734 494.801 368.687] /A << /S /GoTo /D <74686D2E352E322E332E31> >> >>}
\MusicPDFWriteObject{2759}{<< /Type /Annot /Subtype /Link /Border [0 0 0] /H /I /C [1 0 0] /Rect [206.223 250.996 236.437 262.949] /A << /S /GoTo /D <6571756174696F6E2E352E322E332E31> >> >>}
\MusicPDFWriteObject{2760}{<< /Type /Annot /Subtype /Link /Border [0 0 0] /H /I /C [1 0 0] /Rect [390.912 224.694 421.125 236.647] /A << /S /GoTo /D <6571756174696F6E2E352E322E332E32> >> >>}
\MusicPDFWriteObject{2761}{<< /Type /Annot /Subtype /Link /Border [0 0 0] /H /I /C [1 0 0] /Rect [190.033 139.163 220.81 155.683] /A << /S /GoTo /D <74686D2E362E322E312E31> >> >>}
\MusicPDFWriteObject{2762}{<< /Type /Annot /Subtype /Link /Border [0 0 0] /H /I /C [1 0 0] /Rect [402.26 681.617 432.049 699.017] /A << /S /GoTo /D <74686D2E352E342E312E31> >> >>}
\MusicPDFWriteObject{2763}{<< /Type /Annot /Subtype /Link /Border [0 0 0] /H /I /C [0 1 0] /Rect [188.741 266.593 207.845 275.461] /A << /S /GoTo /D <636974652E4E42> >> >>}
\MusicPDFWriteObject{2764}{<< /Type /Annot /Subtype /Link /Border [0 0 0] /H /I /C [1 0 0] /Rect [187.794 165.885 221.355 174.739] /A << /S /GoTo /D <74686D2E362E322E312E31> >> >>}
\MusicPDFWriteObject{2765}{<< /Type /Annot /Subtype /Link /Border [0 0 0] /H /I /C [1 0 0] /Rect [341.085 161.601 371.863 176.457] /A << /S /GoTo /D <74686D2E342E332E332E31> >> >>}
\MusicPDFWriteObject{2766}{<< /Type /Annot /Subtype /Link /Border [0 0 0] /H /I /C [1 0 0] /Rect [129.634 136.494 160.412 148.807] /A << /S /GoTo /D <74686D2E352E322E312E31> >> >>}
\MusicPDFWriteObject{2767}{<< /Type /Annot /Subtype /Link /Border [0 0 0] /H /I /C [1 0 0] /Rect [226.066 121.91 256.279 136.765] /A << /S /GoTo /D <74686D2E362E322E312E33> >> >>}
\MusicPDFWriteObject{2768}{<< /Type /Annot /Subtype /Link /Border [0 0 0] /H /I /C [1 0 0] /Rect [244.791 684.273 274.44 696.588] /A << /S /GoTo /D <74686D2E342E332E332E31> >> >>}
\MusicPDFWriteObject{2769}{<< /Type /Annot /Subtype /Link /Border [0 0 0] /H /I /C [1 0 0] /Rect [427.691 564.608 458.468 576.921] /A << /S /GoTo /D <74686D2E362E312E322E31> >> >>}
\MusicPDFWriteObject{2770}{<< /Type /Annot /Subtype /Link /Border [0 0 0] /H /I /C [1 0 0] /Rect [126.691 385.674 156.904 397.627] /A << /S /GoTo /D <6571756174696F6E2E362E322E312E31> >> >>}
\MusicPDFWriteObject{2771}{<< /Type /Annot /Subtype /Link /Border [0 0 0] /H /I /C [1 0 0] /Rect [199.265 275.045 229.309 287.377] /A << /S /GoTo /D <74686D2E352E322E312E31> >> >>}
\MusicPDFWriteObject{2772}{<< /Type /Annot /Subtype /Link /Border [0 0 0] /H /I /C [1 0 0] /Rect [508.978 150.578 539.756 162.531] /A << /S /GoTo /D <74686D2E352E322E332E31> >> >>}
\MusicPDFWriteObject{2773}{<< /Type /Annot /Subtype /Link /Border [0 0 0] /H /I /C [1 0 0] /Rect [91.42 645.181 121.633 657.134] /A << /S /GoTo /D <6571756174696F6E2E362E322E312E32> >> >>}
\MusicPDFWriteObject{2774}{<< /Type /Annot /Subtype /Link /Border [0 0 0] /H /I /C [0 1 0] /Rect [403.096 634.52 420.568 643.388] /A << /S /GoTo /D <636974652E4346> >> >>}
\MusicPDFWriteObject{2775}{<< /Type /Annot /Subtype /Link /Border [0 0 0] /H /I /C [1 0 0] /Rect [188.578 590.076 218.735 604.931] /A << /S /GoTo /D <74686D2E352E322E332E32> >> >>}
\MusicPDFWriteObject{2776}{<< /Type /Annot /Subtype /Link /Border [0 0 0] /H /I /C [1 0 0] /Rect [424.095 563.702 459.288 575.655] /A << /S /GoTo /D <6571756174696F6E2E352E322E332E3132> >> >>}
\MusicPDFWriteObject{2777}{<< /Type /Annot /Subtype /Link /Border [0 0 0] /H /I /C [1 0 0] /Rect [488.441 563.702 523.634 575.655] /A << /S /GoTo /D <6571756174696F6E2E352E322E332E3133> >> >>}
\MusicPDFWriteObject{2778}{<< /Type /Annot /Subtype /Link /Border [0 0 0] /H /I /C [1 0 0] /Rect [203.789 520.304 234.002 532.256] /A << /S /GoTo /D <6571756174696F6E2E352E322E322E33> >> >>}
\MusicPDFWriteObject{2779}{<< /Type /Annot /Subtype /Link /Border [0 0 0] /H /I /C [1 0 0] /Rect [279.871 507.153 310.649 519.106] /A << /S /GoTo /D <74686D2E352E322E312E31> >> >>}
\MusicPDFWriteObject{2780}{<< /Type /Annot /Subtype /Link /Border [0 0 0] /H /I /C [1 0 0] /Rect [111.812 455.869 141.461 468.182] /A << /S /GoTo /D <74686D2E362E322E312E31> >> >>}
\MusicPDFWriteObject{2781}{<< /Type /Annot /Subtype /Link /Border [0 0 0] /H /I /C [1 0 0] /Rect [413.235 608.922 443.449 620.874] /A << /S /GoTo /D <6571756174696F6E2E362E322E322E31> >> >>}
\MusicPDFWriteObject{2782}{<< /Type /Annot /Subtype /Link /Border [0 0 0] /H /I /C [1 0 0] /Rect [317.986 524.829 348.594 537.144] /A << /S /GoTo /D <74686D2E362E322E322E31> >> >>}
\MusicPDFWriteObject{2783}{<< /Type /Annot /Subtype /Link /Border [0 0 0] /H /I /C [1 0 0] /Rect [445.921 511.142 476.134 523.993] /A << /S /GoTo /D <74686D2E362E322E312E31> >> >>}
\MusicPDFWriteObject{2784}{<< /Type /Annot /Subtype /Link /Border [0 0 0] /H /I /C [1 0 0] /Rect [157.182 462.63 186.944 477.109] /A << /S /GoTo /D <74686D2E322E312E342E31> >> >>}
\MusicPDFWriteObject{2785}{<< /Type /Annot /Subtype /Link /Border [0 0 0] /H /I /C [1 0 0] /Rect [481.255 451.313 511.976 463.626] /A << /S /GoTo /D <74686D2E322E312E342E31> >> >>}
\MusicPDFWriteObject{2786}{<< /Type /Annot /Subtype /Link /Border [0 0 0] /H /I /C [1 0 0] /Rect [328.994 438.163 358.671 450.478] /A << /S /GoTo /D <74686D2E322E312E312E36> >> >>}
\MusicPDFWriteObject{2787}{<< /Type /Annot /Subtype /Link /Border [0 0 0] /H /I /C [0 1 0] /Rect [524.151 414.711 543.256 423.579] /A << /S /GoTo /D <636974652E4E42> >> >>}
\MusicPDFWriteObject{2788}{<< /Type /Annot /Subtype /Link /Border [0 0 0] /H /I /C [1 0 0] /Rect [412.637 396.214 445.459 410.693] /A << /S /GoTo /D <74686D2E432E322E352E31> >> >>}
\MusicPDFWriteObject{2789}{<< /Type /Annot /Subtype /Link /Border [0 0 0] /H /I /C [0 1 0] /Rect [107.031 374.927 125.229 383.795] /A << /S /GoTo /D <636974652E4A6172> >> >>}
\MusicPDFWriteObject{2790}{<< /Type /Annot /Subtype /Link /Border [0 0 0] /H /I /C [0 1 0] /Rect [363.501 335.475 381.7 344.343] /A << /S /GoTo /D <636974652E4A6172> >> >>}
\MusicPDFWriteObject{2791}{<< /Type /Annot /Subtype /Link /Border [0 0 0] /H /I /C [0 1 0] /Rect [107.293 309.173 136.054 318.041] /A << /S /GoTo /D <636974652E47726F2D66> >> >>}
\MusicPDFWriteObject{2792}{<< /Type /Annot /Subtype /Link /Border [0 0 0] /H /I /C [0 1 0] /Rect [291.47 309.173 308.79 318.041] /A << /S /GoTo /D <636974652E534E> >> >>}
\MusicPDFWriteObject{2793}{<< /Type /Annot /Subtype /Link /Border [0 0 0] /H /I /C [1 0 0] /Rect [245.695 293.533 276.472 305.485] /A << /S /GoTo /D <74686D2E352E322E312E31> >> >>}
\MusicPDFWriteObject{2794}{<< /Type /Annot /Subtype /Link /Border [0 0 0] /H /I /C [1 0 0] /Rect [353.992 293.533 384.77 305.485] /A << /S /GoTo /D <74686D2E352E342E312E31> >> >>}
\MusicPDFWriteObject{2795}{<< /Type /Annot /Subtype /Link /Border [0 0 0] /H /I /C [1 0 0] /Rect [268.52 280.382 298.733 292.334] /A << /S /GoTo /D <6571756174696F6E2E362E322E322E31> >> >>}
\MusicPDFWriteObject{2796}{<< /Type /Annot /Subtype /Link /Border [0 0 0] /H /I /C [1 0 0] /Rect [402.258 254.633 432.725 265.481] /A << /S /GoTo /D <74686D2E322E322E322E31> >> >>}
\MusicPDFWriteObject{2797}{<< /Type /Annot /Subtype /Link /Border [0 0 0] /H /I /C [1 0 0] /Rect [221.534 215.181 253.623 225.972] /A << /S /GoTo /D <74686D2E412E322E342E34> >> >>}
\MusicPDFWriteObject{2798}{<< /Type /Annot /Subtype /Link /Border [0 0 0] /H /I /C [1 0 0] /Rect [98.089 162.025 128.302 173.978] /A << /S /GoTo /D <6571756174696F6E2E322E322E322E33> >> >>}
\MusicPDFWriteObject{2799}{<< /Type /Annot /Subtype /Link /Border [0 0 0] /H /I /C [1 0 0] /Rect [179.725 162.025 209.487 173.978] /A << /S /GoTo /D <74686D2E322E322E322E31> >> >>}
\MusicPDFWriteObject{2800}{<< /Type /Annot /Subtype /Link /Border [0 0 0] /H /I /C [1 0 0] /Rect [254.077 162.025 284.29 173.978] /A << /S /GoTo /D <6571756174696F6E2E322E332E312E34> >> >>}
\MusicPDFWriteObject{2801}{<< /Type /Annot /Subtype /Link /Border [0 0 0] /H /I /C [1 0 0] /Rect [419.296 162.025 449.058 173.978] /A << /S /GoTo /D <74686D2E322E332E312E32> >> >>}
\MusicPDFWriteObject{2802}{<< /Type /Annot /Subtype /Link /Border [0 0 0] /H /I /C [1 0 0] /Rect [359.006 94.137 388.655 107.701] /A << /S /GoTo /D <74686D2E322E312E322E34> >> >>}
\MusicPDFWriteObject{2803}{<< /Type /Annot /Subtype /Link /Border [0 0 0] /H /I /C [1 0 0] /Rect [262.139 516.323 294.734 531.262] /A << /S /GoTo /D <74686D2E442E312E322E31> >> >>}
\MusicPDFWriteObject{2804}{<< /Type /Annot /Subtype /Link /Border [0 0 0] /H /I /C [1 0 0] /Rect [302.874 491.061 335.577 503.836] /A << /S /GoTo /D <74686D2E412E332E322E32> >> >>}
\MusicPDFWriteObject{2805}{<< /Type /Annot /Subtype /Link /Border [0 0 0] /H /I /C [1 0 0] /Rect [319.402 450.023 349.051 462.799] /A << /S /GoTo /D <74686D2E322E312E342E31> >> >>}
\MusicPDFWriteObject{2806}{<< /Type /Annot /Subtype /Link /Border [0 0 0] /H /I /C [1 0 0] /Rect [268.945 318.236 299.159 330.189] /A << /S /GoTo /D <6571756174696F6E2E362E322E322E37> >> >>}
\MusicPDFWriteObject{2807}{<< /Type /Annot /Subtype /Link /Border [0 0 0] /H /I /C [1 0 0] /Rect [384.541 237.251 414.754 249.203] /A << /S /GoTo /D <6571756174696F6E2E362E322E322E37> >> >>}
\MusicPDFWriteObject{2808}{<< /Type /Annot /Subtype /Link /Border [0 0 0] /H /I /C [1 0 0] /Rect [271.209 138.505 301.422 150.457] /A << /S /GoTo /D <6571756174696F6E2E362E322E322E36> >> >>}
\MusicPDFWriteObject{2809}{<< /Type /Annot /Subtype /Link /Border [0 0 0] /H /I /C [0 1 0] /Rect [351.85 109.647 370.152 118.515] /A << /S /GoTo /D <636974652E4445> >> >>}
\MusicPDFWriteObject{2810}{<< /Type /Annot /Subtype /Link /Border [0 0 0] /H /I /C [0 1 0] /Rect [524.915 109.647 543.217 118.515] /A << /S /GoTo /D <636974652E4445> >> >>}
\MusicPDFWriteObject{2811}{<< /Type /Annot /Subtype /Link /Border [0 0 0] /H /I /C [1 0 0] /Rect [422.821 643.879 453.034 655.832] /A << /S /GoTo /D <6571756174696F6E2E362E322E322E36> >> >>}
\MusicPDFWriteObject{2812}{<< /Type /Annot /Subtype /Link /Border [0 0 0] /H /I /C [1 0 0] /Rect [176.77 630.728 206.983 642.681] /A << /S /GoTo /D <6571756174696F6E2E362E322E322E39> >> >>}
\MusicPDFWriteObject{2813}{<< /Type /Annot /Subtype /Link /Border [0 0 0] /H /I /C [1 0 0] /Rect [120.805 533.681 151.018 545.634] /A << /S /GoTo /D <6571756174696F6E2E362E322E322E34> >> >>}
\MusicPDFWriteObject{2814}{<< /Type /Annot /Subtype /Link /Border [0 0 0] /H /I /C [1 0 0] /Rect [167.354 370.697 197.68 384.142] /A << /S /GoTo /D <74686D2E352E332E312E31> >> >>}
\MusicPDFWriteObject{2815}{<< /Type /Annot /Subtype /Link /Border [0 0 0] /H /I /C [1 0 0] /Rect [274.53 357.546 304.997 370.549] /A << /S /GoTo /D <74686D2E352E332E312E33> >> >>}
\MusicPDFWriteObject{2816}{<< /Type /Annot /Subtype /Link /Border [0 0 0] /H /I /C [1 0 0] /Rect [211.119 345.426 241.333 357.378] /A << /S /GoTo /D <74686D2E352E332E312E32> >> >>}
\MusicPDFWriteObject{2817}{<< /Type /Annot /Subtype /Link /Border [0 0 0] /H /I /C [1 0 0] /Rect [308.37 345.426 338.583 357.378] /A << /S /GoTo /D <74686D2E352E332E312E34> >> >>}
\MusicPDFWriteObject{2818}{<< /Type /Annot /Subtype /Link /Border [0 0 0] /H /I /C [1 0 0] /Rect [415.414 345.426 445.628 357.378] /A << /S /GoTo /D <74686D2E352E332E312E33> >> >>}
\MusicPDFWriteObject{2819}{<< /Type /Annot /Subtype /Link /Border [0 0 0] /H /I /C [1 0 0] /Rect [186.984 629.853 217.197 641.805] /A << /S /GoTo /D <6571756174696F6E2E362E322E322E37> >> >>}
\MusicPDFWriteObject{2820}{<< /Type /Annot /Subtype /Link /Border [0 0 0] /H /I /C [1 0 0] /Rect [414.688 629.853 444.901 641.805] /A << /S /GoTo /D <6571756174696F6E2E362E322E322E35> >> >>}
\MusicPDFWriteObject{2821}{<< /Type /Annot /Subtype /Link /Border [0 0 0] /H /I /C [1 0 0] /Rect [491.685 601.941 522.237 615.506] /A << /S /GoTo /D <74686D2E322E312E312E36> >> >>}
\MusicPDFWriteObject{2822}{<< /Type /Annot /Subtype /Link /Border [0 0 0] /H /I /C [1 0 0] /Rect [298.959 590.04 329.172 602.355] /A << /S /GoTo /D <74686D2E352E342E312E31> >> >>}
\MusicPDFWriteObject{2823}{<< /Type /Annot /Subtype /Link /Border [0 0 0] /H /I /C [0 1 0] /Rect [106.663 579.74 125.768 588.607] /A << /S /GoTo /D <636974652E4E42> >> >>}
\MusicPDFWriteObject{2824}{<< /Type /Annot /Subtype /Link /Border [0 0 0] /H /I /C [1 0 0] /Rect [251.644 550.948 281.857 562.901] /A << /S /GoTo /D <6571756174696F6E2E322E342E322E31> >> >>}
\MusicPDFWriteObject{2825}{<< /Type /Annot /Subtype /Link /Border [0 0 0] /H /I /C [1 0 0] /Rect [352.685 550.948 382.898 562.901] /A << /S /GoTo /D <74686D2E322E342E322E31> >> >>}
\MusicPDFWriteObject{2826}{<< /Type /Annot /Subtype /Link /Border [0 0 0] /H /I /C [1 0 0] /Rect [264.404 398.785 297.721 409.576] /A << /S /GoTo /D <74686D2E412E342E312E31> >> >>}
\MusicPDFWriteObject{2827}{<< /Type /Annot /Subtype /Link /Border [0 0 0] /H /I /C [1 0 0] /Rect [144.409 348.653 179.603 360.606] /A << /S /GoTo /D <6571756174696F6E2E362E322E322E3132> >> >>}
\MusicPDFWriteObject{2828}{<< /Type /Annot /Subtype /Link /Border [0 0 0] /H /I /C [1 0 0] /Rect [490.644 262.491 525.837 274.444] /A << /S /GoTo /D <6571756174696F6E2E352E322E332E3134> >> >>}
\MusicPDFWriteObject{2829}{<< /Type /Annot /Subtype /Link /Border [0 0 0] /H /I /C [0 1 0] /Rect [518.701 251.83 543.159 260.698] /A << /S /GoTo /D <636974652E435053> >> >>}
\MusicPDFWriteObject{2830}{<< /Type /Annot /Subtype /Link /Border [0 0 0] /H /I /C [1 0 0] /Rect [484.961 137.406 515.174 149.359] /A << /S /GoTo /D <6571756174696F6E2E362E332E312E31> >> >>}
\MusicPDFWriteObject{2831}{<< /Type /Annot /Subtype /Link /Border [0 0 0] /H /I /C [1 0 0] /Rect [162.712 684.633 192.925 696.586] /A << /S /GoTo /D <74686D2E332E322E312E32> >> >>}
\MusicPDFWriteObject{2832}{<< /Type /Annot /Subtype /Link /Border [0 0 0] /H /I /C [1 0 0] /Rect [73.712 671.482 103.925 683.435] /A << /S /GoTo /D <6571756174696F6E2E332E322E312E32> >> >>}
\MusicPDFWriteObject{2833}{<< /Type /Annot /Subtype /Link /Border [0 0 0] /H /I /C [1 0 0] /Rect [168.223 671.482 198.436 683.435] /A << /S /GoTo /D <74686D2E332E322E312E31> >> >>}
\MusicPDFWriteObject{2834}{<< /Type /Annot /Subtype /Link /Border [0 0 0] /H /I /C [1 0 0] /Rect [336.75 658.078 366.963 670.286] /A << /S /GoTo /D <74686D2E332E322E312E31> >> >>}
\MusicPDFWriteObject{2835}{<< /Type /Annot /Subtype /Link /Border [0 0 0] /H /I /C [1 0 0] /Rect [161.835 500.426 192.048 512.379] /A << /S /GoTo /D <6571756174696F6E2E362E332E312E32> >> >>}
\MusicPDFWriteObject{2836}{<< /Type /Annot /Subtype /Link /Border [0 0 0] /H /I /C [1 0 0] /Rect [313.697 421.522 343.911 433.474] /A << /S /GoTo /D <6571756174696F6E2E362E332E312E31> >> >>}
\MusicPDFWriteObject{2837}{<< /Type /Annot /Subtype /Link /Border [0 0 0] /H /I /C [1 0 0] /Rect [177.646 352.586 207.859 364.538] /A << /S /GoTo /D <6571756174696F6E2E332E322E312E34> >> >>}
\MusicPDFWriteObject{2838}{<< /Type /Annot /Subtype /Link /Border [0 0 0] /H /I /C [1 0 0] /Rect [236.301 352.586 266.514 364.538] /A << /S /GoTo /D <6571756174696F6E2E332E322E312E35> >> >>}
\MusicPDFWriteObject{2839}{<< /Type /Annot /Subtype /Link /Border [0 0 0] /H /I /C [1 0 0] /Rect [73.712 202.377 103.925 214.33] /A << /S /GoTo /D <6571756174696F6E2E332E332E312E33> >> >>}
\MusicPDFWriteObject{2840}{<< /Type /Annot /Subtype /Link /Border [0 0 0] /H /I /C [1 0 0] /Rect [285.653 132.655 320.847 144.608] /A << /S /GoTo /D <6571756174696F6E2E322E322E342E3233> >> >>}
\MusicPDFWriteObject{2841}{<< /Type /Annot /Subtype /Link /Border [0 0 0] /H /I /C [1 0 0] /Rect [349.288 132.655 379.501 144.608] /A << /S /GoTo /D <6571756174696F6E2E352E352E322E31> >> >>}
\MusicPDFWriteObject{2842}{<< /Type /Annot /Subtype /Link /Border [0 0 0] /H /I /C [1 0 0] /Rect [153.943 656.323 184.72 672.318] /A << /S /GoTo /D <74686D2E352E352E322E31> >> >>}
\MusicPDFWriteObject{2843}{<< /Type /Annot /Subtype /Link /Border [0 0 0] /H /I /C [1 0 0] /Rect [213.467 567.409 243.68 579.361] /A << /S /GoTo /D <74686D2E352E352E312E33> >> >>}
\MusicPDFWriteObject{2844}{<< /Type /Annot /Subtype /Link /Border [0 0 0] /H /I /C [1 0 0] /Rect [268.857 567.409 299.07 579.361] /A << /S /GoTo /D <6571756174696F6E2E352E352E322E33> >> >>}
\MusicPDFWriteObject{2845}{<< /Type /Annot /Subtype /Link /Border [0 0 0] /H /I /C [1 0 0] /Rect [341.234 514.358 375.822 528.36] /A << /S /GoTo /D <74686D2E442E332E312E33> >> >>}
\MusicPDFWriteObject{2846}{<< /Type /Annot /Subtype /Link /Border [0 0 0] /H /I /C [1 0 0] /Rect [246.485 158.636 276.698 170.589] /A << /S /GoTo /D <6571756174696F6E2E362E332E322E37> >> >>}
\MusicPDFWriteObject{2847}{<< /Type /Annot /Subtype /Link /Border [0 0 0] /H /I /C [1 0 0] /Rect [256.056 82.596 286.269 94.549] /A << /S /GoTo /D <6571756174696F6E2E362E332E322E38> >> >>}
\MusicPDFWriteObject{2848}{<< /Type /Annot /Subtype /Link /Border [0 0 0] /H /I /C [1 0 0] /Rect [375.814 636.246 406.591 648.199] /A << /S /GoTo /D <74686D2E352E322E332E31> >> >>}
\MusicPDFWriteObject{2849}{<< /Type /Annot /Subtype /Link /Border [0 0 0] /H /I /C [1 0 0] /Rect [264.237 623.095 294.45 635.048] /A << /S /GoTo /D <6571756174696F6E2E352E322E332E32> >> >>}
\MusicPDFWriteObject{2850}{<< /Type /Annot /Subtype /Link /Border [0 0 0] /H /I /C [1 0 0] /Rect [297.243 546.723 327.766 558.676] /A << /S /GoTo /D <74686D2E362E322E312E31> >> >>}
\MusicPDFWriteObject{2851}{<< /Type /Annot /Subtype /Link /Border [0 0 0] /H /I /C [1 0 0] /Rect [163.904 520.421 194.117 532.374] /A << /S /GoTo /D <6571756174696F6E2E362E332E322E39> >> >>}
\MusicPDFWriteObject{2852}{<< /Type /Annot /Subtype /Link /Border [0 0 0] /H /I /C [1 0 0] /Rect [118.838 400.441 149.051 412.394] /A << /S /GoTo /D <6571756174696F6E2E352E342E312E32> >> >>}
\MusicPDFWriteObject{2853}{<< /Type /Annot /Subtype /Link /Border [0 0 0] /H /I /C [1 0 0] /Rect [131.98 384.996 162.249 396.949] /A << /S /GoTo /D <74686D2E352E342E312E31> >> >>}
\MusicPDFWriteObject{2854}{<< /Type /Annot /Subtype /Link /Border [0 0 0] /H /I /C [1 0 0] /Rect [260.912 371.845 291.125 383.798] /A << /S /GoTo /D <74686D2E352E322E332E31> >> >>}
\MusicPDFWriteObject{2855}{<< /Type /Annot /Subtype /Link /Border [0 0 0] /H /I /C [1 0 0] /Rect [130.82 180.267 161.033 192.219] /A << /S /GoTo /D <6571756174696F6E2E362E332E312E36> >> >>}
\MusicPDFWriteObject{2856}{<< /Type /Annot /Subtype /Link /Border [0 0 0] /H /I /C [1 0 0] /Rect [235.153 180.267 265.367 192.219] /A << /S /GoTo /D <7465726D696E6F6C6F67792E362E332E312E31> >> >>}
\MusicPDFWriteObject{2857}{<< /Type /Annot /Subtype /Link /Border [0 0 0] /H /I /C [1 0 0] /Rect [479.95 153.964 509.599 165.92] /A << /S /GoTo /D <74686D2E332E322E312E33> >> >>}
\MusicPDFWriteObject{2858}{<< /Type /Annot /Subtype /Link /Border [0 0 0] /H /I /C [1 0 0] /Rect [127.331 127.663 159.619 139.616] /A << /S /GoTo /D <74686D2E422E312E312E34> >> >>}
\MusicPDFWriteObject{2859}{<< /Type /Annot /Subtype /Link /Border [0 0 0] /H /I /C [1 0 0] /Rect [233.527 658.332 263.74 670.284] /A << /S /GoTo /D <6571756174696F6E2E362E332E322E33> >> >>}
\MusicPDFWriteObject{2860}{<< /Type /Annot /Subtype /Link /Border [0 0 0] /H /I /C [1 0 0] /Rect [294.685 658.332 324.898 670.284] /A << /S /GoTo /D <6571756174696F6E2E362E332E322E34> >> >>}
\MusicPDFWriteObject{2861}{<< /Type /Annot /Subtype /Link /Border [0 0 0] /H /I /C [1 0 0] /Rect [315.413 590.551 345.626 602.503] /A << /S /GoTo /D <6571756174696F6E2E332E332E312E32> >> >>}
\MusicPDFWriteObject{2862}{<< /Type /Annot /Subtype /Link /Border [0 0 0] /H /I /C [1 0 0] /Rect [402.124 590.551 431.772 602.503] /A << /S /GoTo /D <74686D2E332E332E312E31> >> >>}
\MusicPDFWriteObject{2863}{<< /Type /Annot /Subtype /Link /Border [0 0 0] /H /I /C [1 0 0] /Rect [135.524 577.4 165.737 589.353] /A << /S /GoTo /D <6571756174696F6E2E332E322E322E32> >> >>}
\MusicPDFWriteObject{2864}{<< /Type /Annot /Subtype /Link /Border [0 0 0] /H /I /C [1 0 0] /Rect [224.363 577.039 254.576 589.353] /A << /S /GoTo /D <74686D2E332E322E322E31> >> >>}
\MusicPDFWriteObject{2865}{<< /Type /Annot /Subtype /Link /Border [0 0 0] /H /I /C [1 0 0] /Rect [211.213 418.046 243.229 430.007] /A << /S /GoTo /D <74686D2E422E332E322E31> >> >>}
\MusicPDFWriteObject{2866}{<< /Type /Annot /Subtype /Link /Border [0 0 0] /H /I /C [1 0 0] /Rect [279.936 390.699 312.224 403.704] /A << /S /GoTo /D <74686D2E422E362E312E38> >> >>}
\MusicPDFWriteObject{2867}{<< /Type /Annot /Subtype /Link /Border [0 0 0] /H /I /C [1 0 0] /Rect [265.762 339.148 298.05 351.101] /A << /S /GoTo /D <74686D2E422E332E322E31> >> >>}
\MusicPDFWriteObject{2868}{<< /Type /Annot /Subtype /Link /Border [0 0 0] /H /I /C [1 0 0] /Rect [358.171 339.148 391.012 351.101] /A << /S /GoTo /D <74686D2E442E332E312E31> >> >>}
\MusicPDFWriteObject{2869}{<< /Type /Annot /Subtype /Link /Border [0 0 0] /H /I /C [1 0 0] /Rect [341.27 272.937 374.482 285.252] /A << /S /GoTo /D <74686D2E442E332E312E32> >> >>}
\MusicPDFWriteObject{2870}{<< /Type /Annot /Subtype /Link /Border [0 0 0] /H /I /C [1 0 0] /Rect [98.422 112.957 133.615 124.909] /A << /S /GoTo /D <6571756174696F6E2E362E332E322E3132> >> >>}
\MusicPDFWriteObject{2871}{<< /Type /Annot /Subtype /Link /Border [0 0 0] /H /I /C [1 0 0] /Rect [194.102 593.131 224.203 603.922] /A << /S /GoTo /D <74686D2E362E322E312E31> >> >>}
\MusicPDFWriteObject{2872}{<< /Type /Annot /Subtype /Link /Border [0 0 0] /H /I /C [1 0 0] /Rect [355.438 593.131 385.538 603.922] /A << /S /GoTo /D <74686D2E362E322E322E31> >> >>}
\MusicPDFWriteObject{2873}{<< /Type /Annot /Subtype /Link /Border [0 0 0] /H /I /C [1 0 0] /Rect [298.439 465.54 328.652 477.492] /A << /S /GoTo /D <6571756174696F6E2E362E322E322E32> >> >>}
\MusicPDFWriteObject{2874}{<< /Type /Annot /Subtype /Link /Border [0 0 0] /H /I /C [1 0 0] /Rect [191.455 411.635 221.978 422.426] /A << /S /GoTo /D <74686D2E362E322E322E31> >> >>}
\MusicPDFWriteObject{2875}{<< /Type /Annot /Subtype /Link /Border [0 0 0] /H /I /C [1 0 0] /Rect [176.681 372.183 206.641 382.974] /A << /S /GoTo /D <74686D2E322E312E342E31> >> >>}
\MusicPDFWriteObject{2876}{<< /Type /Annot /Subtype /Link /Border [0 0 0] /H /I /C [1 0 0] /Rect [268.919 372.183 298.879 382.974] /A << /S /GoTo /D <74686D2E362E322E322E31> >> >>}
\MusicPDFWriteObject{2877}{<< /Type /Annot /Subtype /Link /Border [0 0 0] /H /I /C [1 0 0] /Rect [137.52 250.097 172.713 262.05] /A << /S /GoTo /D <6571756174696F6E2E362E332E322E3135> >> >>}
\MusicPDFWriteObject{2878}{<< /Type /Annot /Subtype /Link /Border [0 0 0] /H /I /C [1 0 0] /Rect [178.591 237.5 209.368 248.29] /A << /S /GoTo /D <74686D2E322E342E322E31> >> >>}
\MusicPDFWriteObject{2879}{<< /Type /Annot /Subtype /Link /Border [0 0 0] /H /I /C [1 0 0] /Rect [303.038 237.5 333.815 248.29] /A << /S /GoTo /D <74686D2E322E342E332E32> >> >>}
\MusicPDFWriteObject{2880}{<< /Type /Annot /Subtype /Link /Border [0 0 0] /H /I /C [1 0 0] /Rect [514.843 208.851 545.592 223.706] /A << /S /GoTo /D <74686D2E322E342E322E31> >> >>}
\MusicPDFWriteObject{2881}{<< /Type /Annot /Subtype /Link /Border [0 0 0] /H /I /C [1 0 0] /Rect [211.609 197.962 241.681 208.895] /A << /S /GoTo /D <74686D2E322E342E332E32> >> >>}
\MusicPDFWriteObject{2882}{<< /Type /Annot /Subtype /Link /Border [0 0 0] /H /I /C [1 0 0] /Rect [150.464 184.343 180.677 196.296] /A << /S /GoTo /D <74686D2E362E322E322E31> >> >>}
\MusicPDFWriteObject{2883}{<< /Type /Annot /Subtype /Link /Border [0 0 0] /H /I /C [1 0 0] /Rect [136.157 539.249 168.998 551.201] /A << /S /GoTo /D <74686D2E442E332E332E31> >> >>}
\MusicPDFWriteObject{2884}{<< /Type /Annot /Subtype /Link /Border [0 0 0] /H /I /C [1 0 0] /Rect [291.681 539.249 324.522 551.201] /A << /S /GoTo /D <74686D2E442E332E322E35> >> >>}
\MusicPDFWriteObject{2885}{<< /Type /Annot /Subtype /Link /Border [0 0 0] /H /I /C [1 0 0] /Rect [160.276 499.436 191.053 511.751] /A << /S /GoTo /D <74686D2E352E332E342E34> >> >>}
\MusicPDFWriteObject{2886}{<< /Type /Annot /Subtype /Link /Border [0 0 0] /H /I /C [1 0 0] /Rect [486.14 499.796 516.353 511.749] /A << /S /GoTo /D <6571756174696F6E2E352E332E342E31> >> >>}
\MusicPDFWriteObject{2887}{<< /Type /Annot /Subtype /Link /Border [0 0 0] /H /I /C [1 0 0] /Rect [119.178 486.285 149.956 498.6] /A << /S /GoTo /D <74686D2E352E332E342E31> >> >>}
\MusicPDFWriteObject{2888}{<< /Type /Annot /Subtype /Link /Border [0 0 0] /H /I /C [1 0 0] /Rect [207.684 447.747 239.972 458.537] /A << /S /GoTo /D <74686D2E422E362E312E38> >> >>}
\MusicPDFWriteObject{2889}{<< /Type /Annot /Subtype /Link /Border [0 0 0] /H /I /C [1 0 0] /Rect [154.559 359.149 186.847 371.101] /A << /S /GoTo /D <74686D2E422E312E312E34> >> >>}
\MusicPDFWriteObject{2890}{<< /Type /Annot /Subtype /Link /Border [0 0 0] /H /I /C [1 0 0] /Rect [208.383 273.365 238.031 285.069] /A << /S /GoTo /D <7465726D696E6F6C6F67792E362E332E312E31> >> >>}
\MusicPDFWriteObject{2891}{<< /Type /Annot /Subtype /Link /Border [0 0 0] /H /I /C [1 0 0] /Rect [265.699 179.695 295.912 191.648] /A << /S /GoTo /D <6571756174696F6E2E332E332E312E32> >> >>}
\MusicPDFWriteObject{2892}{<< /Type /Annot /Subtype /Link /Border [0 0 0] /H /I /C [1 0 0] /Rect [366.739 178.085 396.953 191.65] /A << /S /GoTo /D <74686D2E332E332E312E31> >> >>}
\MusicPDFWriteObject{2893}{<< /Type /Annot /Subtype /Link /Border [0 0 0] /H /I /C [1 0 0] /Rect [296.923 684.633 327.137 696.586] /A << /S /GoTo /D <6571756174696F6E2E332E332E312E32> >> >>}
\MusicPDFWriteObject{2894}{<< /Type /Annot /Subtype /Link /Border [0 0 0] /H /I /C [1 0 0] /Rect [384.603 682.839 414.252 698.94] /A << /S /GoTo /D <74686D2E332E332E312E31> >> >>}
\MusicPDFWriteObject{2895}{<< /Type /Annot /Subtype /Link /Border [0 0 0] /H /I /C [1 0 0] /Rect [466.811 574.985 499.099 586.938] /A << /S /GoTo /D <6571756174696F6E2E422E312E312E39> >> >>}
\MusicPDFWriteObject{2896}{<< /Type /Annot /Subtype /Link /Border [0 0 0] /H /I /C [1 0 0] /Rect [71.004 535.527 103.019 547.487] /A << /S /GoTo /D <74686D2E422E362E312E38> >> >>}
\MusicPDFWriteObject{2897}{<< /Type /Annot /Subtype /Link /Border [0 0 0] /H /I /C [1 0 0] /Rect [286.037 509.231 318.325 521.184] /A << /S /GoTo /D <6571756174696F6E2E422E312E312E39> >> >>}
\MusicPDFWriteObject{2898}{<< /Type /Annot /Subtype /Link /Border [0 0 0] /H /I /C [1 0 0] /Rect [511.675 509.231 541.888 521.184] /A << /S /GoTo /D <6571756174696F6E2E332E332E312E32> >> >>}
\MusicPDFWriteObject{2899}{<< /Type /Annot /Subtype /Link /Border [0 0 0] /H /I /C [1 0 0] /Rect [136.63 496.634 166.844 507.424] /A << /S /GoTo /D <74686D2E332E332E312E31> >> >>}
\MusicPDFWriteObject{2900}{<< /Type /Annot /Subtype /Link /Border [0 0 0] /H /I /C [1 0 0] /Rect [183.652 385.664 213.3 403.075] /A << /S /GoTo /D <74686D2E352E342E312E31> >> >>}
\MusicPDFWriteObject{2901}{<< /Type /Annot /Subtype /Link /Border [0 0 0] /H /I /C [1 0 0] /Rect [509.183 391.121 539.396 403.074] /A << /S /GoTo /D <6571756174696F6E2E352E342E312E32> >> >>}
\MusicPDFWriteObject{2902}{<< /Type /Annot /Subtype /Link /Border [0 0 0] /H /I /C [1 0 0] /Rect [210.702 282.777 240.379 294.73] /A << /S /GoTo /D <74686D2E352E352E322E31> >> >>}
\MusicPDFWriteObject{2903}{<< /Type /Annot /Subtype /Link /Border [0 0 0] /H /I /C [1 0 0] /Rect [371.716 282.777 401.929 294.73] /A << /S /GoTo /D <6571756174696F6E2E362E332E322E35> >> >>}
\MusicPDFWriteObject{2904}{<< /Type /Annot /Subtype /Link /Border [0 0 0] /H /I /C [1 0 0] /Rect [148.859 269.627 179.072 281.579] /A << /S /GoTo /D <6571756174696F6E2E362E332E322E37> >> >>}
\MusicPDFWriteObject{2905}{<< /Type /Annot /Subtype /Link /Border [0 0 0] /H /I /C [1 0 0] /Rect [207.514 269.627 237.727 281.579] /A << /S /GoTo /D <6571756174696F6E2E362E332E322E38> >> >>}
\MusicPDFWriteObject{2906}{<< /Type /Annot /Subtype /Link /Border [0 0 0] /H /I /C [1 0 0] /Rect [253.523 237.602 283.172 248.393] /A << /S /GoTo /D <74686D2E362E332E322E31> >> >>}
\MusicPDFWriteObject{2907}{<< /Type /Annot /Subtype /Link /Border [0 0 0] /H /I /C [1 0 0] /Rect [108.51 223.898 143.703 235.851] /A << /S /GoTo /D <6571756174696F6E2E352E322E332E3134> >> >>}
\MusicPDFWriteObject{2908}{<< /Type /Annot /Subtype /Link /Border [0 0 0] /H /I /C [0 1 0] /Rect [366.771 200.086 391.229 208.954] /A << /S /GoTo /D <636974652E435053> >> >>}
\MusicPDFWriteObject{2909}{<< /Type /Annot /Subtype /Link /Border [0 0 0] /H /I /C [1 0 0] /Rect [434.874 158.475 468.95 170.43] /A << /S /GoTo /D <74686D2E442E332E312E33> >> >>}
\MusicPDFWriteObject{2910}{<< /Type /Annot /Subtype /Link /Border [0 0 0] /H /I /C [1 0 0] /Rect [248.743 316.16 283.936 328.112] /A << /S /GoTo /D <6571756174696F6E2E362E332E322E3233> >> >>}
\MusicPDFWriteObject{2911}{<< /Type /Annot /Subtype /Link /Border [0 0 0] /H /I /C [1 0 0] /Rect [316.259 180.782 351.452 192.734] /A << /S /GoTo /D <6571756174696F6E2E362E332E322E3234> >> >>}
\MusicPDFWriteObject{2912}{<< /Type /Annot /Subtype /Link /Border [0 0 0] /H /I /C [1 0 0] /Rect [377.499 134.754 407.712 146.707] /A << /S /GoTo /D <74686D2E352E332E312E35> >> >>}
\MusicPDFWriteObject{2913}{<< /Type /Annot /Subtype /Link /Border [0 0 0] /H /I /C [1 0 0] /Rect [470.488 134.754 500.701 146.707] /A << /S /GoTo /D <74686D2E362E332E322E31> >> >>}
\MusicPDFWriteObject{2914}{<< /Type /Annot /Subtype /Link /Border [0 0 0] /H /I /C [1 0 0] /Rect [82.959 558.462 194.182 569.252] /A << /S /GoTo /D <73656374696F6E2E372E31> >> >>}
\MusicPDFWriteObject{2915}{<< /Type /Annot /Subtype /Link /Border [0 0 0] /H /I /C [1 0 0] /Rect [82.959 545.311 230.619 556.102] /A << /S /GoTo /D <73756273656374696F6E2E372E312E31> >> >>}
\MusicPDFWriteObject{2916}{<< /Type /Annot /Subtype /Link /Border [0 0 0] /H /I /C [1 0 0] /Rect [82.959 531.247 343.332 544.022] /A << /S /GoTo /D <73756273656374696F6E2E372E312E32> >> >>}
\MusicPDFWriteObject{2917}{<< /Type /Annot /Subtype /Link /Border [0 0 0] /H /I /C [1 0 0] /Rect [82.959 519.009 201.929 529.8] /A << /S /GoTo /D <73756273656374696F6E2E372E312E33> >> >>}
\MusicPDFWriteObject{2918}{<< /Type /Annot /Subtype /Link /Border [0 0 0] /H /I /C [1 0 0] /Rect [82.959 505.859 260.528 516.649] /A << /S /GoTo /D <73756273656374696F6E2E372E312E34> >> >>}
\MusicPDFWriteObject{2919}{<< /Type /Annot /Subtype /Link /Border [0 0 0] /H /I /C [1 0 0] /Rect [82.959 494.645 263.144 503.499] /A << /S /GoTo /D <73656374696F6E2E372E32> >> >>}
\MusicPDFWriteObject{2920}{<< /Type /Annot /Subtype /Link /Border [0 0 0] /H /I /C [1 0 0] /Rect [82.959 479.557 183.142 490.348] /A << /S /GoTo /D <73756273656374696F6E2E372E322E31> >> >>}
\MusicPDFWriteObject{2921}{<< /Type /Annot /Subtype /Link /Border [0 0 0] /H /I /C [1 0 0] /Rect [82.959 466.406 240.718 477.197] /A << /S /GoTo /D <73756273656374696F6E2E372E322E32> >> >>}
\MusicPDFWriteObject{2922}{<< /Type /Annot /Subtype /Link /Border [0 0 0] /H /I /C [1 0 0] /Rect [82.959 453.256 369.564 464.046] /A << /S /GoTo /D <73656374696F6E2E372E33> >> >>}
\MusicPDFWriteObject{2923}{<< /Type /Annot /Subtype /Link /Border [0 0 0] /H /I /C [1 0 0] /Rect [82.959 442.042 220.272 450.896] /A << /S /GoTo /D <73756273656374696F6E2E372E332E31> >> >>}
\MusicPDFWriteObject{2924}{<< /Type /Annot /Subtype /Link /Border [0 0 0] /H /I /C [1 0 0] /Rect [82.959 426.954 206.77 437.745] /A << /S /GoTo /D <73756273656374696F6E2E372E332E32> >> >>}
\MusicPDFWriteObject{2925}{<< /Type /Annot /Subtype /Link /Border [0 0 0] /H /I /C [1 0 0] /Rect [82.959 413.803 343.197 424.594] /A << /S /GoTo /D <73756273656374696F6E2E372E332E33> >> >>}
\MusicPDFWriteObject{2926}{<< /Type /Annot /Subtype /Link /Border [0 0 0] /H /I /C [1 0 0] /Rect [82.959 402.589 211.254 411.443] /A << /S /GoTo /D <73756273656374696F6E2E372E332E34> >> >>}
\MusicPDFWriteObject{2927}{<< /Type /Annot /Subtype /Link /Border [0 0 0] /H /I /C [1 0 0] /Rect [82.959 387.502 155.531 398.293] /A << /S /GoTo /D <73656374696F6E2E372E34> >> >>}
\MusicPDFWriteObject{2928}{<< /Type /Annot /Subtype /Link /Border [0 0 0] /H /I /C [1 0 0] /Rect [82.959 374.351 202.316 385.142] /A << /S /GoTo /D <73756273656374696F6E2E372E342E31> >> >>}
\MusicPDFWriteObject{2929}{<< /Type /Annot /Subtype /Link /Border [0 0 0] /H /I /C [1 0 0] /Rect [82.959 361.2 313.898 371.991] /A << /S /GoTo /D <73756273656374696F6E2E372E342E32> >> >>}
\MusicPDFWriteObject{2930}{<< /Type /Annot /Subtype /Link /Border [0 0 0] /H /I /C [1 0 0] /Rect [82.959 348.05 294.531 358.84] /A << /S /GoTo /D <73756273656374696F6E2E372E342E33> >> >>}
\MusicPDFWriteObject{2931}{<< /Type /Annot /Subtype /Link /Border [0 0 0] /H /I /C [1 0 0] /Rect [82.959 334.899 199.134 345.69] /A << /S /GoTo /D <73656374696F6E2E372E35> >> >>}
\MusicPDFWriteObject{2932}{<< /Type /Annot /Subtype /Link /Border [0 0 0] /H /I /C [0 1 0] /Rect [157.596 202.584 175.795 211.451] /A << /S /GoTo /D <636974652E4A6172> >> >>}
\MusicPDFWriteObject{2933}{<< /Type /Annot /Subtype /Link /Border [0 0 0] /H /I /C [0 1 0] /Rect [358.123 98.237 376.322 107.105] /A << /S /GoTo /D <636974652E4A6172> >> >>}
\MusicPDFWriteObject{2934}{<< /Type /Annot /Subtype /Link /Border [0 0 0] /H /I /C [0 1 0] /Rect [187.389 673.973 205.587 682.84] /A << /S /GoTo /D <636974652E4A6172> >> >>}
\MusicPDFWriteObject{2935}{<< /Type /Annot /Subtype /Link /Border [0 0 0] /H /I /C [0 1 0] /Rect [218.505 647.671 236.704 656.539] /A << /S /GoTo /D <636974652E4A6172> >> >>}
\MusicPDFWriteObject{2936}{<< /Type /Annot /Subtype /Link /Border [0 0 0] /H /I /C [0 1 0] /Rect [220.42 621.369 238.618 630.237] /A << /S /GoTo /D <636974652E4A6172> >> >>}
\MusicPDFWriteObject{2937}{<< /Type /Annot /Subtype /Link /Border [0 0 0] /H /I /C [0 1 0] /Rect [132.213 608.219 151.933 617.087] /A << /S /GoTo /D <636974652E4B616C> >> >>}
\MusicPDFWriteObject{2938}{<< /Type /Annot /Subtype /Link /Border [0 0 0] /H /I /C [0 1 0] /Rect [161.163 581.917 180.883 590.785] /A << /S /GoTo /D <636974652E4B616C> >> >>}
\MusicPDFWriteObject{2939}{<< /Type /Annot /Subtype /Link /Border [0 0 0] /H /I /C [0 1 0] /Rect [192.517 542.465 212.238 551.333] /A << /S /GoTo /D <636974652E4B616C> >> >>}
\MusicPDFWriteObject{2940}{<< /Type /Annot /Subtype /Link /Border [0 0 0] /H /I /C [0 1 0] /Rect [151.077 516.163 170.797 525.031] /A << /S /GoTo /D <636974652E4B616C> >> >>}
\MusicPDFWriteObject{2941}{<< /Type /Annot /Subtype /Link /Border [0 0 0] /H /I /C [1 0 0] /Rect [220.677 459.383 250.326 473.043] /A << /S /GoTo /D <74686D2E342E322E312E32> >> >>}
\MusicPDFWriteObject{2942}{<< /Type /Annot /Subtype /Link /Border [0 0 0] /H /I /C [0 1 0] /Rect [92.24 317.647 131.984 326.515] /A << /S /GoTo /D <636974652E4977616E696563> >> >>}
\MusicPDFWriteObject{2943}{<< /Type /Annot /Subtype /Link /Border [0 0 0] /H /I /C [0 1 0] /Rect [85.944 291.345 122.562 300.213] /A << /S /GoTo /D <636974652E5361726E616B> >> >>}
\MusicPDFWriteObject{2944}{<< /Type /Annot /Subtype /Link /Border [0 0 0] /H /I /C [0 1 0] /Rect [433.637 263.266 466.313 272.134] /A << /S /GoTo /D <636974652E426F687231> >> >>}
\MusicPDFWriteObject{2945}{<< /Type /Annot /Subtype /Link /Border [0 0 0] /H /I /C [0 1 0] /Rect [471.372 263.266 504.048 272.134] /A << /S /GoTo /D <636974652E426F687232> >> >>}
\MusicPDFWriteObject{2946}{<< /Type /Annot /Subtype /Link /Border [0 0 0] /H /I /C [0 1 0] /Rect [509.107 263.266 541.783 272.134] /A << /S /GoTo /D <636974652E426F687233> >> >>}
\MusicPDFWriteObject{2947}{<< /Type /Annot /Subtype /Link /Border [0 0 0] /H /I /C [1 0 0] /Rect [322.743 221.324 354.425 233.276] /A << /S /GoTo /D <74686D2E422E352E332E31> >> >>}
\MusicPDFWriteObject{2948}{<< /Type /Annot /Subtype /Link /Border [0 0 0] /H /I /C [1 0 0] /Rect [511.675 671.482 541.888 683.435] /A << /S /GoTo /D <6571756174696F6E2E372E312E312E31> >> >>}
\MusicPDFWriteObject{2949}{<< /Type /Annot /Subtype /Link /Border [0 0 0] /H /I /C [1 0 0] /Rect [172.999 618.879 203.212 630.832] /A << /S /GoTo /D <6571756174696F6E2E372E312E312E31> >> >>}
\MusicPDFWriteObject{2950}{<< /Type /Annot /Subtype /Link /Border [0 0 0] /H /I /C [1 0 0] /Rect [150.527 519.971 180.74 531.924] /A << /S /GoTo /D <6571756174696F6E2E372E312E312E32> >> >>}
\MusicPDFWriteObject{2951}{<< /Type /Annot /Subtype /Link /Border [0 0 0] /H /I /C [0 1 0] /Rect [159.39 509.311 186.338 518.179] /A << /S /GoTo /D <636974652E426F6872> >> >>}
\MusicPDFWriteObject{2952}{<< /Type /Annot /Subtype /Link /Border [0 0 0] /H /I /C [1 0 0] /Rect [440.847 506.821 471.06 518.773] /A << /S /GoTo /D <6571756174696F6E2E372E312E312E31> >> >>}
\MusicPDFWriteObject{2953}{<< /Type /Annot /Subtype /Link /Border [0 0 0] /H /I /C [1 0 0] /Rect [106.039 467.368 138.843 479.321] /A << /S /GoTo /D <73756273756273656374696F6E2E422E352E332E31> >> >>}
\MusicPDFWriteObject{2954}{<< /Type /Annot /Subtype /Link /Border [0 0 0] /H /I /C [0 1 0] /Rect [521.748 456.708 543.184 465.576] /A << /S /GoTo /D <636974652E426131> >> >>}
\MusicPDFWriteObject{2955}{<< /Type /Annot /Subtype /Link /Border [0 0 0] /H /I /C [0 1 0] /Rect [173.422 430.406 194.858 439.274] /A << /S /GoTo /D <636974652E426131> >> >>}
\MusicPDFWriteObject{2956}{<< /Type /Annot /Subtype /Link /Border [0 0 0] /H /I /C [0 1 0] /Rect [309.901 430.406 331.337 439.274] /A << /S /GoTo /D <636974652E426131> >> >>}
\MusicPDFWriteObject{2957}{<< /Type /Annot /Subtype /Link /Border [0 0 0] /H /I /C [0 1 0] /Rect [491.036 417.255 512.472 426.123] /A << /S /GoTo /D <636974652E426132> >> >>}
\MusicPDFWriteObject{2958}{<< /Type /Annot /Subtype /Link /Border [0 0 0] /H /I /C [0 1 0] /Rect [162.362 390.954 195.037 399.822] /A << /S /GoTo /D <636974652E426F687234> >> >>}
\MusicPDFWriteObject{2959}{<< /Type /Annot /Subtype /Link /Border [0 0 0] /H /I /C [0 1 0] /Rect [361.473 390.954 388.422 399.822] /A << /S /GoTo /D <636974652E426F6872> >> >>}
\MusicPDFWriteObject{2960}{<< /Type /Annot /Subtype /Link /Border [0 0 0] /H /I /C [0 1 0] /Rect [521.784 390.954 543.22 399.822] /A << /S /GoTo /D <636974652E426131> >> >>}
\MusicPDFWriteObject{2961}{<< /Type /Annot /Subtype /Link /Border [0 0 0] /H /I /C [0 1 0] /Rect [230.074 377.803 262.749 386.671] /A << /S /GoTo /D <636974652E426F687234> >> >>}
\MusicPDFWriteObject{2962}{<< /Type /Annot /Subtype /Link /Border [0 0 0] /H /I /C [0 1 0] /Rect [379.353 364.652 400.789 373.52] /A << /S /GoTo /D <636974652E426131> >> >>}
\MusicPDFWriteObject{2963}{<< /Type /Annot /Subtype /Link /Border [0 0 0] /H /I /C [1 0 0] /Rect [311.64 336.414 343.322 347.205] /A << /S /GoTo /D <74686D2E422E352E332E33> >> >>}
\MusicPDFWriteObject{2964}{<< /Type /Annot /Subtype /Link /Border [0 0 0] /H /I /C [1 0 0] /Rect [361.012 322.71 393.3 334.663] /A << /S /GoTo /D <74686D2E422E352E322E31> >> >>}
\MusicPDFWriteObject{2965}{<< /Type /Annot /Subtype /Link /Border [0 0 0] /H /I /C [1 0 0] /Rect [373.803 672.036 403.452 682.827] /A << /S /GoTo /D <74686D2E372E312E312E31> >> >>}
\MusicPDFWriteObject{2966}{<< /Type /Annot /Subtype /Link /Border [0 0 0] /H /I /C [1 0 0] /Rect [203.949 588.484 233.852 603.721] /A << /S /GoTo /D <74686D2E372E312E312E31> >> >>}
\MusicPDFWriteObject{2967}{<< /Type /Annot /Subtype /Link /Border [0 0 0] /H /I /C [1 0 0] /Rect [176.323 491.23 209.884 500.084] /A << /S /GoTo /D <74686D2E372E312E312E31> >> >>}
\MusicPDFWriteObject{2968}{<< /Type /Annot /Subtype /Link /Border [0 0 0] /H /I /C [0 1 0] /Rect [238.638 464.928 260.073 473.796] /A << /S /GoTo /D <636974652E426131> >> >>}
\MusicPDFWriteObject{2969}{<< /Type /Annot /Subtype /Link /Border [0 0 0] /H /I /C [0 1 0] /Rect [336.565 464.928 358.001 473.796] /A << /S /GoTo /D <636974652E426131> >> >>}
\MusicPDFWriteObject{2970}{<< /Type /Annot /Subtype /Link /Border [0 0 0] /H /I /C [1 0 0] /Rect [267.574 436.136 297.787 448.089] /A << /S /GoTo /D <6571756174696F6E2E372E312E312E31> >> >>}
\MusicPDFWriteObject{2971}{<< /Type /Annot /Subtype /Link /Border [0 0 0] /H /I /C [0 1 0] /Rect [344.826 425.476 366.262 434.344] /A << /S /GoTo /D <636974652E426131> >> >>}
\MusicPDFWriteObject{2972}{<< /Type /Annot /Subtype /Link /Border [0 0 0] /H /I /C [1 0 0] /Rect [271.075 184.646 301.288 196.598] /A << /S /GoTo /D <6571756174696F6E2E372E312E312E35> >> >>}
\MusicPDFWriteObject{2973}{<< /Type /Annot /Subtype /Link /Border [0 0 0] /H /I /C [1 0 0] /Rect [226.866 684.633 257.079 696.586] /A << /S /GoTo /D <6571756174696F6E2E372E312E312E37> >> >>}
\MusicPDFWriteObject{2974}{<< /Type /Annot /Subtype /Link /Border [0 0 0] /H /I /C [1 0 0] /Rect [248.86 648.206 279.073 660.158] /A << /S /GoTo /D <6571756174696F6E2E372E312E312E36> >> >>}
\MusicPDFWriteObject{2975}{<< /Type /Annot /Subtype /Link /Border [0 0 0] /H /I /C [1 0 0] /Rect [449.405 147.786 482.299 158.577] /A << /S /GoTo /D <74686D2E422E332E312E31> >> >>}
\MusicPDFWriteObject{2976}{<< /Type /Annot /Subtype /Link /Border [0 0 0] /H /I /C [1 0 0] /Rect [364.277 134.082 396.565 146.035] /A << /S /GoTo /D <74686D2E422E332E322E31> >> >>}
\MusicPDFWriteObject{2977}{<< /Type /Annot /Subtype /Link /Border [0 0 0] /H /I /C [1 0 0] /Rect [70.755 96.3 100.798 107.091] /A << /S /GoTo /D <74686D2E322E342E312E31> >> >>}
\MusicPDFWriteObject{2978}{<< /Type /Annot /Subtype /Link /Border [0 0 0] /H /I /C [1 0 0] /Rect [256.1 96.3 286.144 107.091] /A << /S /GoTo /D <74686D2E322E342E322E31> >> >>}
\MusicPDFWriteObject{2979}{<< /Type /Annot /Subtype /Link /Border [0 0 0] /H /I /C [1 0 0] /Rect [172.902 85.086 205.19 93.94] /A << /S /GoTo /D <74686D2E422E322E312E37> >> >>}
\MusicPDFWriteObject{2980}{<< /Type /Annot /Subtype /Link /Border [0 0 0] /H /I /C [1 0 0] /Rect [209.624 685.187 240.401 695.977] /A << /S /GoTo /D <74686D2E372E312E322E33> >> >>}
\MusicPDFWriteObject{2981}{<< /Type /Annot /Subtype /Link /Border [0 0 0] /H /I /C [0 1 0] /Rect [372.023 611.031 408.641 619.898] /A << /S /GoTo /D <636974652E5361726E616B> >> >>}
\MusicPDFWriteObject{2982}{<< /Type /Annot /Subtype /Link /Border [0 0 0] /H /I /C [0 1 0] /Rect [503.94 584.729 540.558 593.597] /A << /S /GoTo /D <636974652E5361726E616B> >> >>}
\MusicPDFWriteObject{2983}{<< /Type /Annot /Subtype /Link /Border [0 0 0] /H /I /C [1 0 0] /Rect [514.576 569.642 544.845 580.432] /A << /S /GoTo /D <74686D2E372E312E322E33> >> >>}
\MusicPDFWriteObject{2984}{<< /Type /Annot /Subtype /Link /Border [0 0 0] /H /I /C [1 0 0] /Rect [146.553 322.901 168.855 338.984] /A << /S /GoTo /D <73756273656374696F6E2E332E312E31> >> >>}
\MusicPDFWriteObject{2985}{<< /Type /Annot /Subtype /Link /Border [0 0 0] /H /I /C [1 0 0] /Rect [99.26 285.059 129.473 297.012] /A << /S /GoTo /D <6571756174696F6E2E372E312E322E31> >> >>}
\MusicPDFWriteObject{2986}{<< /Type /Annot /Subtype /Link /Border [0 0 0] /H /I /C [1 0 0] /Rect [158.424 285.059 188.638 297.012] /A << /S /GoTo /D <6571756174696F6E2E372E312E322E32> >> >>}
\MusicPDFWriteObject{2987}{<< /Type /Annot /Subtype /Link /Border [0 0 0] /H /I /C [0 1 0] /Rect [461.135 196.905 478.608 205.773] /A << /S /GoTo /D <636974652E4346> >> >>}
\MusicPDFWriteObject{2988}{<< /Type /Annot /Subtype /Link /Border [0 0 0] /H /I /C [1 0 0] /Rect [408.793 181.265 439.006 193.217] /A << /S /GoTo /D <6571756174696F6E2E332E332E312E33> >> >>}
\MusicPDFWriteObject{2989}{<< /Type /Annot /Subtype /Link /Border [0 0 0] /H /I /C [1 0 0] /Rect [503.304 181.265 533.518 193.217] /A << /S /GoTo /D <74686D2E332E332E312E31> >> >>}
\MusicPDFWriteObject{2990}{<< /Type /Annot /Subtype /Link /Border [0 0 0] /H /I /C [1 0 0] /Rect [511.085 496.557 533.551 507.348] /A << /S /GoTo /D <73756273656374696F6E2E372E312E33> >> >>}
\MusicPDFWriteObject{2991}{<< /Type /Annot /Subtype /Link /Border [0 0 0] /H /I /C [0 1 0] /Rect [428.536 444.792 446.008 453.66] /A << /S /GoTo /D <636974652E4346> >> >>}
\MusicPDFWriteObject{2992}{<< /Type /Annot /Subtype /Link /Border [0 0 0] /H /I /C [0 1 0] /Rect [374.446 246.432 411.064 255.3] /A << /S /GoTo /D <636974652E5361726E616B> >> >>}
\MusicPDFWriteObject{2993}{<< /Type /Annot /Subtype /Link /Border [0 0 0] /H /I /C [0 1 0] /Rect [187.671 206.98 224.289 215.847] /A << /S /GoTo /D <636974652E5361726E616B> >> >>}
\MusicPDFWriteObject{2994}{<< /Type /Annot /Subtype /Link /Border [0 0 0] /H /I /C [0 1 0] /Rect [118.624 193.829 155.242 202.697] /A << /S /GoTo /D <636974652E5361726E616B> >> >>}
\MusicPDFWriteObject{2995}{<< /Type /Annot /Subtype /Link /Border [0 0 0] /H /I /C [1 0 0] /Rect [261.813 541.791 292.026 553.744] /A << /S /GoTo /D <6571756174696F6E2E332E332E312E33> >> >>}
\MusicPDFWriteObject{2996}{<< /Type /Annot /Subtype /Link /Border [0 0 0] /H /I /C [1 0 0] /Rect [476.263 541.79 507.782 553.746] /A << /S /GoTo /D <74686D2E372E312E322E32> >> >>}
\MusicPDFWriteObject{2997}{<< /Type /Annot /Subtype /Link /Border [0 0 0] /H /I /C [1 0 0] /Rect [421.524 179.614 456.718 191.567] /A << /S /GoTo /D <74686D2E31342E312E322E33> >> >>}
\MusicPDFWriteObject{2998}{<< /Type /Annot /Subtype /Link /Border [0 0 0] /H /I /C [1 0 0] /Rect [203.115 162.054 233.469 172.844] /A << /S /GoTo /D <74686D2E372E312E322E33> >> >>}
\MusicPDFWriteObject{2999}{<< /Type /Annot /Subtype /Link /Border [0 0 0] /H /I /C [0 1 0] /Rect [141.97 150.84 178.588 159.708] /A << /S /GoTo /D <636974652E5361726E616B> >> >>}
\MusicPDFWriteObject{3000}{<< /Type /Annot /Subtype /Link /Border [0 0 0] /H /I /C [0 1 0] /Rect [197.298 647.671 233.916 656.539] /A << /S /GoTo /D <636974652E5361726E616B> >> >>}
\MusicPDFWriteObject{3001}{<< /Type /Annot /Subtype /Link /Border [0 0 0] /H /I /C [0 1 0] /Rect [167.483 555.616 204.102 564.484] /A << /S /GoTo /D <636974652E5361726E616B> >> >>}
\MusicPDFWriteObject{3002}{<< /Type /Annot /Subtype /Link /Border [0 0 0] /H /I /C [1 0 0] /Rect [255.069 417.295 285.847 432.15] /A << /S /GoTo /D <74686D2E352E352E312E31> >> >>}
\MusicPDFWriteObject{3003}{<< /Type /Annot /Subtype /Link /Border [0 0 0] /H /I /C [1 0 0] /Rect [351.532 417.295 382.31 432.15] /A << /S /GoTo /D <74686D2E352E352E312E32> >> >>}
\MusicPDFWriteObject{3004}{<< /Type /Annot /Subtype /Link /Border [0 0 0] /H /I /C [1 0 0] /Rect [265.824 405.938 296.037 417.891] /A << /S /GoTo /D <74686D2E372E322E312E34> >> >>}
\MusicPDFWriteObject{3005}{<< /Type /Annot /Subtype /Link /Border [0 0 0] /H /I /C [0 1 0] /Rect [149.489 382.127 186.107 390.994] /A << /S /GoTo /D <636974652E5361726E616B> >> >>}
\MusicPDFWriteObject{3006}{<< /Type /Annot /Subtype /Link /Border [0 0 0] /H /I /C [0 1 0] /Rect [363.767 316.373 400.385 325.241] /A << /S /GoTo /D <636974652E5361726E616B> >> >>}
\MusicPDFWriteObject{3007}{<< /Type /Annot /Subtype /Link /Border [0 0 0] /H /I /C [0 1 0] /Rect [439.719 316.373 476.337 325.241] /A << /S /GoTo /D <636974652E5361726E616B> >> >>}
\MusicPDFWriteObject{3008}{<< /Type /Annot /Subtype /Link /Border [0 0 0] /H /I /C [0 1 0] /Rect [85.534 290.071 122.152 298.939] /A << /S /GoTo /D <636974652E5361726E616B> >> >>}
\MusicPDFWriteObject{3009}{<< /Type /Annot /Subtype /Link /Border [0 0 0] /H /I /C [0 1 0] /Rect [380.931 276.921 417.549 285.788] /A << /S /GoTo /D <636974652E5361726E616B> >> >>}
\MusicPDFWriteObject{3010}{<< /Type /Annot /Subtype /Link /Border [0 0 0] /H /I /C [0 1 0] /Rect [506.64 276.921 543.258 285.788] /A << /S /GoTo /D <636974652E5361726E616B> >> >>}
\MusicPDFWriteObject{3011}{<< /Type /Annot /Subtype /Link /Border [0 0 0] /H /I /C [0 1 0] /Rect [350.697 263.77 387.316 272.638] /A << /S /GoTo /D <636974652E5361726E616B> >> >>}
\MusicPDFWriteObject{3012}{<< /Type /Annot /Subtype /Link /Border [0 0 0] /H /I /C [0 1 0] /Rect [212.293 250.619 248.911 259.487] /A << /S /GoTo /D <636974652E5361726E616B> >> >>}
\MusicPDFWriteObject{3013}{<< /Type /Annot /Subtype /Link /Border [0 0 0] /H /I /C [0 1 0] /Rect [422.227 250.619 461.972 259.487] /A << /S /GoTo /D <636974652E4977616E696563> >> >>}
\MusicPDFWriteObject{3014}{<< /Type /Annot /Subtype /Link /Border [0 0 0] /H /I /C [0 1 0] /Rect [73.77 237.468 113.536 246.336] /A << /S /GoTo /D <636974652E4B75626F7461> >> >>}
\MusicPDFWriteObject{3015}{<< /Type /Annot /Subtype /Link /Border [0 0 0] /H /I /C [0 1 0] /Rect [347.518 235.532 380.955 246.336] /A << /S /GoTo /D <636974652E5A6167696572> >> >>}
\MusicPDFWriteObject{3016}{<< /Type /Annot /Subtype /Link /Border [0 0 0] /H /I /C [1 0 0] /Rect [338.231 146.501 368.642 157.292] /A << /S /GoTo /D <74686D2E372E312E322E33> >> >>}
\MusicPDFWriteObject{3017}{<< /Type /Annot /Subtype /Link /Border [0 0 0] /H /I /C [0 1 0] /Rect [193.341 135.287 229.959 144.155] /A << /S /GoTo /D <636974652E5361726E616B> >> >>}
\MusicPDFWriteObject{3018}{<< /Type /Annot /Subtype /Link /Border [0 0 0] /H /I /C [0 1 0] /Rect [366.056 121.395 402.674 130.262] /A << /S /GoTo /D <636974652E5361726E616B> >> >>}
\MusicPDFWriteObject{3019}{<< /Type /Annot /Subtype /Link /Border [0 0 0] /H /I /C [1 0 0] /Rect [286.514 623.701 316.727 635.654] /A << /S /GoTo /D <6571756174696F6E2E372E312E322E34> >> >>}
\MusicPDFWriteObject{3020}{<< /Type /Annot /Subtype /Link /Border [0 0 0] /H /I /C [1 0 0] /Rect [322.445 610.551 352.658 622.503] /A << /S /GoTo /D <6571756174696F6E2E372E312E322E34> >> >>}
\MusicPDFWriteObject{3021}{<< /Type /Annot /Subtype /Link /Border [0 0 0] /H /I /C [1 0 0] /Rect [469.333 541.355 499.546 553.307] /A << /S /GoTo /D <6571756174696F6E2E372E312E322E34> >> >>}
\MusicPDFWriteObject{3022}{<< /Type /Annot /Subtype /Link /Border [0 0 0] /H /I /C [1 0 0] /Rect [498.426 498.392 528.639 510.344] /A << /S /GoTo /D <6571756174696F6E2E372E312E322E33> >> >>}
\MusicPDFWriteObject{3023}{<< /Type /Annot /Subtype /Link /Border [0 0 0] /H /I /C [1 0 0] /Rect [111.625 484.819 141.273 497.594] /A << /S /GoTo /D <74686D2E372E312E322E33> >> >>}
\MusicPDFWriteObject{3024}{<< /Type /Annot /Subtype /Link /Border [0 0 0] /H /I /C [1 0 0] /Rect [418.993 484.819 448.641 497.594] /A << /S /GoTo /D <74686D2E372E312E322E33> >> >>}
\MusicPDFWriteObject{3025}{<< /Type /Annot /Subtype /Link /Border [0 0 0] /H /I /C [0 1 0] /Rect [145.871 474.518 182.489 483.386] /A << /S /GoTo /D <636974652E5361726E616B> >> >>}
\MusicPDFWriteObject{3026}{<< /Type /Annot /Subtype /Link /Border [0 0 0] /H /I /C [0 1 0] /Rect [85.574 447.476 122.192 456.344] /A << /S /GoTo /D <636974652E5361726E616B> >> >>}
\MusicPDFWriteObject{3027}{<< /Type /Annot /Subtype /Link /Border [0 0 0] /H /I /C [1 0 0] /Rect [202.78 444.986 232.993 456.938] /A << /S /GoTo /D <6571756174696F6E2E372E312E322E34> >> >>}
\MusicPDFWriteObject{3028}{<< /Type /Annot /Subtype /Link /Border [0 0 0] /H /I /C [0 1 0] /Rect [176.223 368.571 210.897 377.439] /A << /S /GoTo /D <636974652E4174436865> >> >>}
\MusicPDFWriteObject{3029}{<< /Type /Annot /Subtype /Link /Border [0 0 0] /H /I /C [0 1 0] /Rect [224.478 355.42 261.096 364.288] /A << /S /GoTo /D <636974652E5361726E616B> >> >>}
\MusicPDFWriteObject{3030}{<< /Type /Annot /Subtype /Link /Border [0 0 0] /H /I /C [0 1 0] /Rect [335.519 355.42 370.193 364.288] /A << /S /GoTo /D <636974652E4174436865> >> >>}
\MusicPDFWriteObject{3031}{<< /Type /Annot /Subtype /Link /Border [0 0 0] /H /I /C [0 1 0] /Rect [514.93 342.27 540.342 351.138] /A << /S /GoTo /D <636974652E42435A> >> >>}
\MusicPDFWriteObject{3032}{<< /Type /Annot /Subtype /Link /Border [0 0 0] /H /I /C [0 1 0] /Rect [215.618 302.817 246.891 311.685] /A << /S /GoTo /D <636974652E536572696573> >> >>}
\MusicPDFWriteObject{3033}{<< /Type /Annot /Subtype /Link /Border [0 0 0] /H /I /C [0 1 0] /Rect [311.396 302.817 342.668 311.685] /A << /S /GoTo /D <636974652E536572696573> >> >>}
\MusicPDFWriteObject{3034}{<< /Type /Annot /Subtype /Link /Border [0 0 0] /H /I /C [0 1 0] /Rect [365.738 289.667 400.412 298.535] /A << /S /GoTo /D <636974652E4174436865> >> >>}
\MusicPDFWriteObject{3035}{<< /Type /Annot /Subtype /Link /Border [0 0 0] /H /I /C [1 0 0] /Rect [375.175 261.428 405.585 272.219] /A << /S /GoTo /D <74686D2E372E312E322E33> >> >>}
\MusicPDFWriteObject{3036}{<< /Type /Annot /Subtype /Link /Border [0 0 0] /H /I /C [1 0 0] /Rect [112.727 234.574 142.94 246.526] /A << /S /GoTo /D <74686D2E322E342E322E31> >> >>}
\MusicPDFWriteObject{3037}{<< /Type /Annot /Subtype /Link /Border [0 0 0] /H /I /C [0 1 0] /Rect [156.754 210.762 193.372 219.63] /A << /S /GoTo /D <636974652E5361726E616B> >> >>}
\MusicPDFWriteObject{3038}{<< /Type /Annot /Subtype /Link /Border [0 0 0] /H /I /C [1 0 0] /Rect [375.811 208.272 406.024 220.225] /A << /S /GoTo /D <74686D2E372E312E322E33> >> >>}
\MusicPDFWriteObject{3039}{<< /Type /Annot /Subtype /Link /Border [0 0 0] /H /I /C [1 0 0] /Rect [255.488 155.669 285.701 167.622] /A << /S /GoTo /D <6571756174696F6E2E352E312E312E32> >> >>}
\MusicPDFWriteObject{3040}{<< /Type /Annot /Subtype /Link /Border [0 0 0] /H /I /C [0 1 0] /Rect [451.437 673.973 488.056 682.84] /A << /S /GoTo /D <636974652E5361726E616B> >> >>}
\MusicPDFWriteObject{3041}{<< /Type /Annot /Subtype /Link /Border [0 0 0] /H /I /C [1 0 0] /Rect [214.306 628.52 244.434 640.473] /A << /S /GoTo /D <74686D2E372E312E322E33> >> >>}
\MusicPDFWriteObject{3042}{<< /Type /Annot /Subtype /Link /Border [0 0 0] /H /I /C [0 1 0] /Rect [506.604 631.01 543.223 639.878] /A << /S /GoTo /D <636974652E5361726E616B> >> >>}
\MusicPDFWriteObject{3043}{<< /Type /Annot /Subtype /Link /Border [0 0 0] /H /I /C [1 0 0] /Rect [325.574 615.369 355.787 627.322] /A << /S /GoTo /D <6571756174696F6E2E332E332E312E33> >> >>}
\MusicPDFWriteObject{3044}{<< /Type /Annot /Subtype /Link /Border [0 0 0] /H /I /C [1 0 0] /Rect [418.485 615.369 448.134 627.322] /A << /S /GoTo /D <74686D2E332E332E312E31> >> >>}
\MusicPDFWriteObject{3045}{<< /Type /Annot /Subtype /Link /Border [0 0 0] /H /I /C [0 1 0] /Rect [279.541 604.709 316.159 613.576] /A << /S /GoTo /D <636974652E5361726E616B> >> >>}
\MusicPDFWriteObject{3046}{<< /Type /Annot /Subtype /Link /Border [0 0 0] /H /I /C [0 1 0] /Rect [420.076 604.709 456.694 613.576] /A << /S /GoTo /D <636974652E5361726E616B> >> >>}
\MusicPDFWriteObject{3047}{<< /Type /Annot /Subtype /Link /Border [0 0 0] /H /I /C [1 0 0] /Rect [414.612 589.068 446.386 601.02] /A << /S /GoTo /D <74686D2E422E332E322E31> >> >>}
\MusicPDFWriteObject{3048}{<< /Type /Annot /Subtype /Link /Border [0 0 0] /H /I /C [1 0 0] /Rect [508.228 589.068 540.001 601.02] /A << /S /GoTo /D <74686D2E422E322E312E37> >> >>}
\MusicPDFWriteObject{3049}{<< /Type /Annot /Subtype /Link /Border [0 0 0] /H /I /C [1 0 0] /Rect [258.827 575.917 288.927 587.87] /A << /S /GoTo /D <74686D2E372E312E322E33> >> >>}
\MusicPDFWriteObject{3050}{<< /Type /Annot /Subtype /Link /Border [0 0 0] /H /I /C [0 1 0] /Rect [312.322 578.407 348.94 587.275] /A << /S /GoTo /D <636974652E5361726E616B> >> >>}
\MusicPDFWriteObject{3051}{<< /Type /Annot /Subtype /Link /Border [0 0 0] /H /I /C [1 0 0] /Rect [224.487 562.766 254.7 574.719] /A << /S /GoTo /D <6571756174696F6E2E372E312E322E33> >> >>}
\MusicPDFWriteObject{3052}{<< /Type /Annot /Subtype /Link /Border [0 0 0] /H /I /C [1 0 0] /Rect [315.798 537.018 346.35 547.809] /A << /S /GoTo /D <74686D2E372E312E312E31> >> >>}
\MusicPDFWriteObject{3053}{<< /Type /Annot /Subtype /Link /Border [0 0 0] /H /I /C [1 0 0] /Rect [131.896 523.314 164.184 535.267] /A << /S /GoTo /D <74686D2E422E352E322E31> >> >>}
\MusicPDFWriteObject{3054}{<< /Type /Annot /Subtype /Link /Border [0 0 0] /H /I /C [1 0 0] /Rect [239.605 523.314 271.893 535.267] /A << /S /GoTo /D <74686D2E422E352E332E33> >> >>}
\MusicPDFWriteObject{3055}{<< /Type /Annot /Subtype /Link /Border [0 0 0] /H /I /C [1 0 0] /Rect [195.787 510.717 226 521.507] /A << /S /GoTo /D <74686D2E372E312E322E33> >> >>}
\MusicPDFWriteObject{3056}{<< /Type /Annot /Subtype /Link /Border [0 0 0] /H /I /C [1 0 0] /Rect [187.794 488.694 221.355 497.548] /A << /S /GoTo /D <74686D2E372E312E322E33> >> >>}
\MusicPDFWriteObject{3057}{<< /Type /Annot /Subtype /Link /Border [0 0 0] /H /I /C [1 0 0] /Rect [116.265 221.679 147.947 236.916] /A << /S /GoTo /D <74686D2E422E332E312E31> >> >>}
\MusicPDFWriteObject{3058}{<< /Type /Annot /Subtype /Link /Border [0 0 0] /H /I /C [1 0 0] /Rect [145.622 209.818 175.835 222.676] /A << /S /GoTo /D <74686D2E372E312E312E31> >> >>}
\MusicPDFWriteObject{3059}{<< /Type /Annot /Subtype /Link /Border [0 0 0] /H /I /C [1 0 0] /Rect [341.98 209.818 347.942 222.676] /A << /S /GoTo /D <48666F6F746E6F74652E3235> >> >>}
\MusicPDFWriteObject{3060}{<< /Type /Annot /Subtype /Link /Border [0 0 0] /H /I /C [1 0 0] /Rect [125.141 193.652 155.918 208.622] /A << /S /GoTo /D <74686D2E372E312E322E32> >> >>}
\MusicPDFWriteObject{3061}{<< /Type /Annot /Subtype /Link /Border [0 0 0] /H /I /C [1 0 0] /Rect [147.091 93.056 175.346 104.644] /A << /S /GoTo /D <74686D2E422E332E312E31> >> >>}
\MusicPDFWriteObject{3062}{<< /Type /Annot /Subtype /Link /Border [0 0 0] /H /I /C [1 0 0] /Rect [440.363 93.056 466.823 104.644] /A << /S /GoTo /D <74686D2E372E312E312E31> >> >>}
\MusicPDFWriteObject{3063}{<< /Type /Annot /Subtype /Link /Border [0 0 0] /H /I /C [1 0 0] /Rect [314.209 684.633 346.103 696.588] /A << /S /GoTo /D <74686D2E422E332E322E31> >> >>}
\MusicPDFWriteObject{3064}{<< /Type /Annot /Subtype /Link /Border [0 0 0] /H /I /C [1 0 0] /Rect [156.922 671.482 187.135 683.435] /A << /S /GoTo /D <6571756174696F6E2E372E312E322E33> >> >>}
\MusicPDFWriteObject{3065}{<< /Type /Annot /Subtype /Link /Border [0 0 0] /H /I /C [1 0 0] /Rect [229.239 565.199 259.452 577.151] /A << /S /GoTo /D <6571756174696F6E2E372E312E322E33> >> >>}
\MusicPDFWriteObject{3066}{<< /Type /Annot /Subtype /Link /Border [0 0 0] /H /I /C [1 0 0] /Rect [70.755 525.746 100.403 538.162] /A << /S /GoTo /D <74686D2E372E312E312E31> >> >>}
\MusicPDFWriteObject{3067}{<< /Type /Annot /Subtype /Link /Border [0 0 0] /H /I /C [1 0 0] /Rect [228.431 499.445 258.644 511.398] /A << /S /GoTo /D <6571756174696F6E2E372E312E322E33> >> >>}
\MusicPDFWriteObject{3068}{<< /Type /Annot /Subtype /Link /Border [0 0 0] /H /I /C [1 0 0] /Rect [71.004 432.446 103.262 444.398] /A << /S /GoTo /D <74686D2E422E332E322E31> >> >>}
\MusicPDFWriteObject{3069}{<< /Type /Annot /Subtype /Link /Border [0 0 0] /H /I /C [1 0 0] /Rect [274.445 432.446 304.658 444.398] /A << /S /GoTo /D <6571756174696F6E2E372E312E322E33> >> >>}
\MusicPDFWriteObject{3070}{<< /Type /Annot /Subtype /Link /Border [0 0 0] /H /I /C [1 0 0] /Rect [208.938 331.97 239.151 343.923] /A << /S /GoTo /D <6571756174696F6E2E372E312E322E33> >> >>}
\MusicPDFWriteObject{3071}{<< /Type /Annot /Subtype /Link /Border [0 0 0] /H /I /C [1 0 0] /Rect [403.201 317.025 432.934 331.881] /A << /S /GoTo /D <74686D2E372E322E312E34> >> >>}
\MusicPDFWriteObject{3072}{<< /Type /Annot /Subtype /Link /Border [0 0 0] /H /I /C [1 0 0] /Rect [131.018 281.857 161.231 290.711] /A << /S /GoTo /D <74686D2E372E312E322E33> >> >>}
\MusicPDFWriteObject{3073}{<< /Type /Annot /Subtype /Link /Border [0 0 0] /H /I /C [1 0 0] /Rect [248.881 266.77 278.811 277.56] /A << /S /GoTo /D <74686D2E322E342E322E31> >> >>}
\MusicPDFWriteObject{3074}{<< /Type /Annot /Subtype /Link /Border [0 0 0] /H /I /C [1 0 0] /Rect [352.486 266.77 382.417 277.56] /A << /S /GoTo /D <74686D2E322E332E312E34> >> >>}
\MusicPDFWriteObject{3075}{<< /Type /Annot /Subtype /Link /Border [0 0 0] /H /I /C [1 0 0] /Rect [294.878 253.065 325.091 265.018] /A << /S /GoTo /D <6571756174696F6E2E372E312E322E33> >> >>}
\MusicPDFWriteObject{3076}{<< /Type /Annot /Subtype /Link /Border [0 0 0] /H /I /C [1 0 0] /Rect [221.989 627.857 251.638 639.809] /A << /S /GoTo /D <74686D2E372E312E322E33> >> >>}
\MusicPDFWriteObject{3077}{<< /Type /Annot /Subtype /Link /Border [0 0 0] /H /I /C [0 1 0] /Rect [342.651 630.347 379.269 639.215] /A << /S /GoTo /D <636974652E5361726E616B> >> >>}
\MusicPDFWriteObject{3078}{<< /Type /Annot /Subtype /Link /Border [0 0 0] /H /I /C [1 0 0] /Rect [233.397 562.103 265.823 574.056] /A << /S /GoTo /D <74686D2E432E352E332E32> >> >>}
\MusicPDFWriteObject{3079}{<< /Type /Annot /Subtype /Link /Border [0 0 0] /H /I /C [1 0 0] /Rect [476.841 562.103 509.267 574.056] /A << /S /GoTo /D <6571756174696F6E2E432E352E332E33> >> >>}
\MusicPDFWriteObject{3080}{<< /Type /Annot /Subtype /Link /Border [0 0 0] /H /I /C [1 0 0] /Rect [384.833 547.342 406.889 560.907] /A << /S /GoTo /D <73756273656374696F6E2E332E312E31> >> >>}
\MusicPDFWriteObject{3081}{<< /Type /Annot /Subtype /Link /Border [0 0 0] /H /I /C [1 0 0] /Rect [256.264 658.332 286.477 670.284] /A << /S /GoTo /D <6571756174696F6E2E372E312E332E31> >> >>}
\MusicPDFWriteObject{3082}{<< /Type /Annot /Subtype /Link /Border [0 0 0] /H /I /C [1 0 0] /Rect [226.898 572.987 246.597 583.835] /A << /S /GoTo /D <73656374696F6E2E31342E32> >> >>}
\MusicPDFWriteObject{3083}{<< /Type /Annot /Subtype /Link /Border [0 0 0] /H /I /C [1 0 0] /Rect [265.947 559.837 296.16 570.627] /A << /S /GoTo /D <74686D2E372E312E332E32> >> >>}
\MusicPDFWriteObject{3084}{<< /Type /Annot /Subtype /Link /Border [0 0 0] /H /I /C [1 0 0] /Rect [475.331 486.176 505.551 498.491] /A << /S /GoTo /D <74686D2E372E312E322E32> >> >>}
\MusicPDFWriteObject{3085}{<< /Type /Annot /Subtype /Link /Border [0 0 0] /H /I /C [1 0 0] /Rect [157.192 203.807 187.405 214.598] /A << /S /GoTo /D <74686D2E372E312E332E32> >> >>}
\MusicPDFWriteObject{3086}{<< /Type /Annot /Subtype /Link /Border [0 0 0] /H /I /C [1 0 0] /Rect [250.181 203.807 280.394 214.598] /A << /S /GoTo /D <74686D2E372E312E322E33> >> >>}
\MusicPDFWriteObject{3087}{<< /Type /Annot /Subtype /Link /Border [0 0 0] /H /I /C [1 0 0] /Rect [511.675 110.418 541.888 122.371] /A << /S /GoTo /D <6571756174696F6E2E372E312E332E32> >> >>}
\MusicPDFWriteObject{3088}{<< /Type /Annot /Subtype /Link /Border [0 0 0] /H /I /C [1 0 0] /Rect [433.505 97.268 463.718 109.22] /A << /S /GoTo /D <6571756174696F6E2E332E332E312E33> >> >>}
\MusicPDFWriteObject{3089}{<< /Type /Annot /Subtype /Link /Border [0 0 0] /H /I /C [1 0 0] /Rect [196.444 654.37 226.657 666.323] /A << /S /GoTo /D <74686D2E372E312E322E33> >> >>}
\MusicPDFWriteObject{3090}{<< /Type /Annot /Subtype /Link /Border [0 0 0] /H /I /C [1 0 0] /Rect [392.841 654.37 423.055 666.323] /A << /S /GoTo /D <6571756174696F6E2E372E312E322E33> >> >>}
\MusicPDFWriteObject{3091}{<< /Type /Annot /Subtype /Link /Border [0 0 0] /H /I /C [1 0 0] /Rect [164.441 589.25 194.655 601.203] /A << /S /GoTo /D <6571756174696F6E2E372E312E332E32> >> >>}
\MusicPDFWriteObject{3092}{<< /Type /Annot /Subtype /Link /Border [0 0 0] /H /I /C [1 0 0] /Rect [259.976 589.25 290.19 601.203] /A << /S /GoTo /D <6571756174696F6E2E372E312E322E33> >> >>}
\MusicPDFWriteObject{3093}{<< /Type /Annot /Subtype /Link /Border [0 0 0] /H /I /C [1 0 0] /Rect [237.327 533.694 268.105 545.647] /A << /S /GoTo /D <74686D2E322E342E322E31> >> >>}
\MusicPDFWriteObject{3094}{<< /Type /Annot /Subtype /Link /Border [0 0 0] /H /I /C [1 0 0] /Rect [447.11 533.694 477.323 545.647] /A << /S /GoTo /D <6571756174696F6E2E372E312E332E32> >> >>}
\MusicPDFWriteObject{3095}{<< /Type /Annot /Subtype /Link /Border [0 0 0] /H /I /C [0 1 0] /Rect [232.791 509.883 248.783 518.751] /A << /S /GoTo /D <636974652E436F> >> >>}
\MusicPDFWriteObject{3096}{<< /Type /Annot /Subtype /Link /Border [0 0 0] /H /I /C [1 0 0] /Rect [354.538 494.795 384.751 505.586] /A << /S /GoTo /D <74686D2E322E342E322E31> >> >>}
\MusicPDFWriteObject{3097}{<< /Type /Annot /Subtype /Link /Border [0 0 0] /H /I /C [1 0 0] /Rect [357.735 475.667 387.384 486.472] /A << /S /GoTo /D <74686D2E362E332E322E31> >> >>}
\MusicPDFWriteObject{3098}{<< /Type /Annot /Subtype /Link /Border [0 0 0] /H /I /C [1 0 0] /Rect [338.897 278.687 370.416 289.478] /A << /S /GoTo /D <74686D2E372E312E342E31> >> >>}
\MusicPDFWriteObject{3099}{<< /Type /Annot /Subtype /Link /Border [0 0 0] /H /I /C [1 0 0] /Rect [329.185 671.482 359.398 683.435] /A << /S /GoTo /D <6571756174696F6E2E372E312E342E31> >> >>}
\MusicPDFWriteObject{3100}{<< /Type /Annot /Subtype /Link /Border [0 0 0] /H /I /C [1 0 0] /Rect [197.678 658.332 227.891 670.284] /A << /S /GoTo /D <6571756174696F6E2E372E312E342E35> >> >>}
\MusicPDFWriteObject{3101}{<< /Type /Annot /Subtype /Link /Border [0 0 0] /H /I /C [1 0 0] /Rect [266.967 468.779 297.18 480.732] /A << /S /GoTo /D <6571756174696F6E2E372E312E342E36> >> >>}
\MusicPDFWriteObject{3102}{<< /Type /Annot /Subtype /Link /Border [0 0 0] /H /I /C [1 0 0] /Rect [353.445 466.732 384.374 481.72] /A << /S /GoTo /D <74686D2E372E312E342E32> >> >>}
\MusicPDFWriteObject{3103}{<< /Type /Annot /Subtype /Link /Border [0 0 0] /H /I /C [0 1 0] /Rect [168.663 394.341 189.089 403.208] /A << /S /GoTo /D <636974652E426172> >> >>}
\MusicPDFWriteObject{3104}{<< /Type /Annot /Subtype /Link /Border [0 0 0] /H /I /C [0 1 0] /Rect [235.555 251.422 255.981 260.29] /A << /S /GoTo /D <636974652E426172> >> >>}
\MusicPDFWriteObject{3105}{<< /Type /Annot /Subtype /Link /Border [0 0 0] /H /I /C [0 1 0] /Rect [167.12 223.695 187.546 232.563] /A << /S /GoTo /D <636974652E426172> >> >>}
\MusicPDFWriteObject{3106}{<< /Type /Annot /Subtype /Link /Border [0 0 0] /H /I /C [0 1 0] /Rect [312.151 160.966 337.212 169.834] /A << /S /GoTo /D <636974652E53636861> >> >>}
\MusicPDFWriteObject{3107}{<< /Type /Annot /Subtype /Link /Border [0 0 0] /H /I /C [0 1 0] /Rect [406.748 147.815 425.853 156.683] /A << /S /GoTo /D <636974652E4E42> >> >>}
\MusicPDFWriteObject{3108}{<< /Type /Annot /Subtype /Link /Border [0 0 0] /H /I /C [0 1 0] /Rect [524.054 147.815 543.159 156.683] /A << /S /GoTo /D <636974652E4E42> >> >>}
\MusicPDFWriteObject{3109}{<< /Type /Annot /Subtype /Link /Border [0 0 0] /H /I /C [0 1 0] /Rect [227.234 624.394 261.999 633.262] /A << /S /GoTo /D <636974652E436E772D62> >> >>}
\MusicPDFWriteObject{3110}{<< /Type /Annot /Subtype /Link /Border [0 0 0] /H /I /C [0 1 0] /Rect [196.475 382.57 215.579 391.437] /A << /S /GoTo /D <636974652E4E42> >> >>}
\MusicPDFWriteObject{3111}{<< /Type /Annot /Subtype /Link /Border [0 0 0] /H /I /C [0 1 0] /Rect [397.143 356.268 416.117 365.136] /A << /S /GoTo /D <636974652E547265> >> >>}
\MusicPDFWriteObject{3112}{<< /Type /Annot /Subtype /Link /Border [0 0 0] /H /I /C [1 0 0] /Rect [202.494 172.569 234.782 184.524] /A << /S /GoTo /D <74686D2E422E322E312E31> >> >>}
\MusicPDFWriteObject{3113}{<< /Type /Annot /Subtype /Link /Border [0 0 0] /H /I /C [1 0 0] /Rect [119.74 559.262 149.953 571.215] /A << /S /GoTo /D <6571756174696F6E2E372E312E342E36> >> >>}
\MusicPDFWriteObject{3114}{<< /Type /Annot /Subtype /Link /Border [0 0 0] /H /I /C [1 0 0] /Rect [501.148 559.262 530.797 571.215] /A << /S /GoTo /D <74686D2E362E332E322E31> >> >>}
\MusicPDFWriteObject{3115}{<< /Type /Annot /Subtype /Link /Border [0 0 0] /H /I /C [1 0 0] /Rect [415.699 383.409 448.593 395.722] /A << /S /GoTo /D <74686D2E422E332E322E31> >> >>}
\MusicPDFWriteObject{3116}{<< /Type /Annot /Subtype /Link /Border [0 0 0] /H /I /C [1 0 0] /Rect [485.137 370.618 518.031 383.947] /A << /S /GoTo /D <74686D2E422E332E312E31> >> >>}
\MusicPDFWriteObject{3117}{<< /Type /Annot /Subtype /Link /Border [0 0 0] /H /I /C [1 0 0] /Rect [148.857 181.414 184.05 193.366] /A << /S /GoTo /D <6571756174696F6E2E372E312E342E3137> >> >>}
\MusicPDFWriteObject{3118}{<< /Type /Annot /Subtype /Link /Border [0 0 0] /H /I /C [0 1 0] /Rect [179.238 547.773 198.343 556.641] /A << /S /GoTo /D <636974652E4E42> >> >>}
\MusicPDFWriteObject{3119}{<< /Type /Annot /Subtype /Link /Border [0 0 0] /H /I /C [1 0 0] /Rect [506.695 444.391 541.888 456.343] /A << /S /GoTo /D <6571756174696F6E2E372E312E342E3230> >> >>}
\MusicPDFWriteObject{3120}{<< /Type /Annot /Subtype /Link /Border [0 0 0] /H /I /C [0 1 0] /Rect [375.245 326.86 393.444 335.728] /A << /S /GoTo /D <636974652E4A6172> >> >>}
\MusicPDFWriteObject{3121}{<< /Type /Annot /Subtype /Link /Border [0 0 0] /H /I /C [1 0 0] /Rect [71.004 298.069 103.845 310.022] /A << /S /GoTo /D <74686D2E442E332E332E31> >> >>}
\MusicPDFWriteObject{3122}{<< /Type /Annot /Subtype /Link /Border [0 0 0] /H /I /C [1 0 0] /Rect [357.437 221.102 389.725 233.054] /A << /S /GoTo /D <74686D2E422E332E322E31> >> >>}
\MusicPDFWriteObject{3123}{<< /Type /Annot /Subtype /Link /Border [0 0 0] /H /I /C [1 0 0] /Rect [262.148 173.403 297.341 185.356] /A << /S /GoTo /D <6571756174696F6E2E372E312E342E3231> >> >>}
\MusicPDFWriteObject{3124}{<< /Type /Annot /Subtype /Link /Border [0 0 0] /H /I /C [0 1 0] /Rect [518.099 687.123 543.159 695.991] /A << /S /GoTo /D <636974652E53636861> >> >>}
\MusicPDFWriteObject{3125}{<< /Type /Annot /Subtype /Link /Border [0 0 0] /H /I /C [0 1 0] /Rect [107.342 660.822 132.403 669.69] /A << /S /GoTo /D <636974652E53636861> >> >>}
\MusicPDFWriteObject{3126}{<< /Type /Annot /Subtype /Link /Border [0 0 0] /H /I /C [1 0 0] /Rect [245.491 613.44 275.14 624.23] /A << /S /GoTo /D <74686D2E372E312E342E32> >> >>}
\MusicPDFWriteObject{3127}{<< /Type /Annot /Subtype /Link /Border [0 0 0] /H /I /C [1 0 0] /Rect [333.747 417.534 363.961 429.487] /A << /S /GoTo /D <6571756174696F6E2E372E312E342E33> >> >>}
\MusicPDFWriteObject{3128}{<< /Type /Annot /Subtype /Link /Border [0 0 0] /H /I /C [1 0 0] /Rect [282.72 226.276 317.914 238.229] /A << /S /GoTo /D <6571756174696F6E2E372E312E342E3236> >> >>}
\MusicPDFWriteObject{3129}{<< /Type /Annot /Subtype /Link /Border [0 0 0] /H /I /C [1 0 0] /Rect [70.755 671.482 100.968 683.437] /A << /S /GoTo /D <74686D2E372E312E342E32> >> >>}
\MusicPDFWriteObject{3130}{<< /Type /Annot /Subtype /Link /Border [0 0 0] /H /I /C [0 1 0] /Rect [127.805 634.249 145.277 643.117] /A << /S /GoTo /D <636974652E4346> >> >>}
\MusicPDFWriteObject{3131}{<< /Type /Annot /Subtype /Link /Border [0 0 0] /H /I /C [1 0 0] /Rect [390 618.608 419.733 630.563] /A << /S /GoTo /D <74686D2E372E312E342E32> >> >>}
\MusicPDFWriteObject{3132}{<< /Type /Annot /Subtype /Link /Border [0 0 0] /H /I /C [1 0 0] /Rect [98.175 574.922 133.368 586.875] /A << /S /GoTo /D <6571756174696F6E2E372E312E342E3132> >> >>}
\MusicPDFWriteObject{3133}{<< /Type /Annot /Subtype /Link /Border [0 0 0] /H /I /C [0 1 0] /Rect [148.22 563.52 166.418 572.388] /A << /S /GoTo /D <636974652E4A6172> >> >>}
\MusicPDFWriteObject{3134}{<< /Type /Annot /Subtype /Link /Border [0 0 0] /H /I /C [1 0 0] /Rect [230.661 321.053 260.874 333.006] /A << /S /GoTo /D <6571756174696F6E2E372E322E302E31> >> >>}
\MusicPDFWriteObject{3135}{<< /Type /Annot /Subtype /Link /Border [0 0 0] /H /I /C [1 0 0] /Rect [301.85 321.053 332.063 333.006] /A << /S /GoTo /D <6571756174696F6E2E372E322E302E32> >> >>}
\MusicPDFWriteObject{3136}{<< /Type /Annot /Subtype /Link /Border [0 0 0] /H /I /C [0 1 0] /Rect [190.619 178.885 209.143 187.753] /A << /S /GoTo /D <636974652E4745> >> >>}
\MusicPDFWriteObject{3137}{<< /Type /Annot /Subtype /Link /Border [0 0 0] /H /I /C [0 1 0] /Rect [297.995 178.885 325.739 187.753] /A << /S /GoTo /D <636974652E4A6F722D61> >> >>}
\MusicPDFWriteObject{3138}{<< /Type /Annot /Subtype /Link /Border [0 0 0] /H /I /C [0 1 0] /Rect [73.77 165.734 101.514 174.602] /A << /S /GoTo /D <636974652E4A6F722D61> >> >>}
\MusicPDFWriteObject{3139}{<< /Type /Annot /Subtype /Link /Border [0 0 0] /H /I /C [1 0 0] /Rect [198.305 101.363 228.518 112.154] /A << /S /GoTo /D <74686D2E332E332E312E31> >> >>}
\MusicPDFWriteObject{3140}{<< /Type /Annot /Subtype /Link /Border [0 0 0] /H /I /C [1 0 0] /Rect [99.062 684.633 129.275 696.586] /A << /S /GoTo /D <6571756174696F6E2E352E312E312E32> >> >>}
\MusicPDFWriteObject{3141}{<< /Type /Annot /Subtype /Link /Border [0 0 0] /H /I /C [1 0 0] /Rect [216.118 634.535 246.332 646.487] /A << /S /GoTo /D <6571756174696F6E2E332E332E312E32> >> >>}
\MusicPDFWriteObject{3142}{<< /Type /Annot /Subtype /Link /Border [0 0 0] /H /I /C [1 0 0] /Rect [452.923 608.233 483.136 620.186] /A << /S /GoTo /D <6571756174696F6E2E332E332E312E33> >> >>}
\MusicPDFWriteObject{3143}{<< /Type /Annot /Subtype /Link /Border [0 0 0] /H /I /C [1 0 0] /Rect [393.349 581.411 423.563 593.364] /A << /S /GoTo /D <6571756174696F6E2E332E332E312E33> >> >>}
\MusicPDFWriteObject{3144}{<< /Type /Annot /Subtype /Link /Border [0 0 0] /H /I /C [0 1 0] /Rect [334.508 537.11 351.981 545.978] /A << /S /GoTo /D <636974652E4346> >> >>}
\MusicPDFWriteObject{3145}{<< /Type /Annot /Subtype /Link /Border [0 0 0] /H /I /C [1 0 0] /Rect [188.792 508.158 219.006 520.11] /A << /S /GoTo /D <6571756174696F6E2E372E322E312E31> >> >>}
\MusicPDFWriteObject{3146}{<< /Type /Annot /Subtype /Link /Border [0 0 0] /H /I /C [1 0 0] /Rect [394.512 481.856 424.725 493.811] /A << /S /GoTo /D <74686D2E372E312E322E32> >> >>}
\MusicPDFWriteObject{3147}{<< /Type /Annot /Subtype /Link /Border [0 0 0] /H /I /C [1 0 0] /Rect [487.502 481.856 517.715 493.811] /A << /S /GoTo /D <74686D2E372E312E332E32> >> >>}
\MusicPDFWriteObject{3148}{<< /Type /Annot /Subtype /Link /Border [0 0 0] /H /I /C [1 0 0] /Rect [172.732 455.194 204.414 467.509] /A << /S /GoTo /D <74686D2E422E332E322E31> >> >>}
\MusicPDFWriteObject{3149}{<< /Type /Annot /Subtype /Link /Border [0 0 0] /H /I /C [1 0 0] /Rect [236.196 258.201 266.409 270.516] /A << /S /GoTo /D <74686D2E372E312E322E32> >> >>}
\MusicPDFWriteObject{3150}{<< /Type /Annot /Subtype /Link /Border [0 0 0] /H /I /C [1 0 0] /Rect [379.589 202.694 409.83 214.647] /A << /S /GoTo /D <74686D2E352E352E312E31> >> >>}
\MusicPDFWriteObject{3151}{<< /Type /Annot /Subtype /Link /Border [0 0 0] /H /I /C [1 0 0] /Rect [475.446 202.694 505.688 214.647] /A << /S /GoTo /D <74686D2E352E352E312E32> >> >>}
\MusicPDFWriteObject{3152}{<< /Type /Annot /Subtype /Link /Border [0 0 0] /H /I /C [1 0 0] /Rect [280.362 189.543 310.575 201.496] /A << /S /GoTo /D <74686D2E372E322E312E31> >> >>}
\MusicPDFWriteObject{3153}{<< /Type /Annot /Subtype /Link /Border [0 0 0] /H /I /C [1 0 0] /Rect [149.219 122.145 179.432 134.097] /A << /S /GoTo /D <74686D2E342E322E312E33> >> >>}
\MusicPDFWriteObject{3154}{<< /Type /Annot /Subtype /Link /Border [0 0 0] /H /I /C [1 0 0] /Rect [327.848 467.191 358.061 479.144] /A << /S /GoTo /D <6571756174696F6E2E332E332E312E32> >> >>}
\MusicPDFWriteObject{3155}{<< /Type /Annot /Subtype /Link /Border [0 0 0] /H /I /C [1 0 0] /Rect [428.889 467.191 459.102 479.144] /A << /S /GoTo /D <74686D2E332E332E312E31> >> >>}
\MusicPDFWriteObject{3156}{<< /Type /Annot /Subtype /Link /Border [0 0 0] /H /I /C [1 0 0] /Rect [171.293 417.092 200.941 429.407] /A << /S /GoTo /D <74686D2E352E312E332E31> >> >>}
\MusicPDFWriteObject{3157}{<< /Type /Annot /Subtype /Link /Border [0 0 0] /H /I /C [1 0 0] /Rect [311.987 250.591 342.2 262.543] /A << /S /GoTo /D <6571756174696F6E2E372E322E312E38> >> >>}
\MusicPDFWriteObject{3158}{<< /Type /Annot /Subtype /Link /Border [0 0 0] /H /I /C [1 0 0] /Rect [227.811 237.44 258.025 249.392] /A << /S /GoTo /D <6571756174696F6E2E372E322E312E38> >> >>}
\MusicPDFWriteObject{3159}{<< /Type /Annot /Subtype /Link /Border [0 0 0] /H /I /C [1 0 0] /Rect [196.131 221.273 225.78 236.244] /A << /S /GoTo /D <74686D2E352E312E322E32> >> >>}
\MusicPDFWriteObject{3160}{<< /Type /Annot /Subtype /Link /Border [0 0 0] /H /I /C [1 0 0] /Rect [191.951 182.083 222.164 194.036] /A << /S /GoTo /D <6571756174696F6E2E352E312E332E32> >> >>}
\MusicPDFWriteObject{3161}{<< /Type /Annot /Subtype /Link /Border [0 0 0] /H /I /C [1 0 0] /Rect [286.922 182.082 317.389 194.037] /A << /S /GoTo /D <74686D2E352E312E332E31> >> >>}
\MusicPDFWriteObject{3162}{<< /Type /Annot /Subtype /Link /Border [0 0 0] /H /I /C [1 0 0] /Rect [368.608 182.083 398.822 194.036] /A << /S /GoTo /D <6571756174696F6E2E372E322E312E38> >> >>}
\MusicPDFWriteObject{3163}{<< /Type /Annot /Subtype /Link /Border [0 0 0] /H /I /C [0 1 0] /Rect [324.151 647.19 343.256 656.058] /A << /S /GoTo /D <636974652E4E42> >> >>}
\MusicPDFWriteObject{3164}{<< /Type /Annot /Subtype /Link /Border [0 0 0] /H /I /C [1 0 0] /Rect [318.71 592.097 348.923 604.049] /A << /S /GoTo /D <6571756174696F6E2E372E322E312E39> >> >>}
\MusicPDFWriteObject{3165}{<< /Type /Annot /Subtype /Link /Border [0 0 0] /H /I /C [0 1 0] /Rect [435.327 594.587 454.432 603.454] /A << /S /GoTo /D <636974652E4E42> >> >>}
\MusicPDFWriteObject{3166}{<< /Type /Annot /Subtype /Link /Border [0 0 0] /H /I /C [1 0 0] /Rect [261.533 359.973 291.633 372.288] /A << /S /GoTo /D <74686D2E372E332E312E32> >> >>}
\MusicPDFWriteObject{3167}{<< /Type /Annot /Subtype /Link /Border [0 0 0] /H /I /C [1 0 0] /Rect [251.177 204.298 281.39 215.088] /A << /S /GoTo /D <74686D2E352E332E342E34> >> >>}
\MusicPDFWriteObject{3168}{<< /Type /Annot /Subtype /Link /Border [0 0 0] /H /I /C [1 0 0] /Rect [458.281 113.96 488.494 125.913] /A << /S /GoTo /D <6571756174696F6E2E372E322E312E32> >> >>}
\MusicPDFWriteObject{3169}{<< /Type /Annot /Subtype /Link /Border [0 0 0] /H /I /C [0 1 0] /Rect [149.831 583.547 167.303 592.414] /A << /S /GoTo /D <636974652E4346> >> >>}
\MusicPDFWriteObject{3170}{<< /Type /Annot /Subtype /Link /Border [0 0 0] /H /I /C [0 1 0] /Rect [270.558 583.547 288.529 592.414] /A << /S /GoTo /D <636974652E5256> >> >>}
\MusicPDFWriteObject{3171}{<< /Type /Annot /Subtype /Link /Border [0 0 0] /H /I /C [0 1 0] /Rect [229.006 570.396 246.479 579.264] /A << /S /GoTo /D <636974652E4346> >> >>}
\MusicPDFWriteObject{3172}{<< /Type /Annot /Subtype /Link /Border [0 0 0] /H /I /C [0 1 0] /Rect [348.639 570.396 366.61 579.264] /A << /S /GoTo /D <636974652E5256> >> >>}
\MusicPDFWriteObject{3173}{<< /Type /Annot /Subtype /Link /Border [0 0 0] /H /I /C [0 1 0] /Rect [405.255 481.552 422.727 490.42] /A << /S /GoTo /D <636974652E4346> >> >>}
\MusicPDFWriteObject{3174}{<< /Type /Annot /Subtype /Link /Border [0 0 0] /H /I /C [0 1 0] /Rect [525.266 481.552 543.236 490.42] /A << /S /GoTo /D <636974652E5256> >> >>}
\MusicPDFWriteObject{3175}{<< /Type /Annot /Subtype /Link /Border [0 0 0] /H /I /C [1 0 0] /Rect [264.54 442.099 294.754 450.953] /A << /S /GoTo /D <74686D2E332E332E312E31> >> >>}
\MusicPDFWriteObject{3176}{<< /Type /Annot /Subtype /Link /Border [0 0 0] /H /I /C [1 0 0] /Rect [241.722 125.328 276.915 137.281] /A << /S /GoTo /D <6571756174696F6E2E332E342E332E3134> >> >>}
\MusicPDFWriteObject{3177}{<< /Type /Annot /Subtype /Link /Border [0 0 0] /H /I /C [1 0 0] /Rect [512.049 123.641 542.978 137.301] /A << /S /GoTo /D <74686D2E352E352E312E32> >> >>}
\MusicPDFWriteObject{3178}{<< /Type /Annot /Subtype /Link /Border [0 0 0] /H /I /C [1 0 0] /Rect [511.675 590.523 541.888 602.475] /A << /S /GoTo /D <6571756174696F6E2E332E332E312E33> >> >>}
\MusicPDFWriteObject{3179}{<< /Type /Annot /Subtype /Link /Border [0 0 0] /H /I /C [1 0 0] /Rect [70.755 515.805 100.968 526.596] /A << /S /GoTo /D <74686D2E332E332E322E31> >> >>}
\MusicPDFWriteObject{3180}{<< /Type /Annot /Subtype /Link /Border [0 0 0] /H /I /C [1 0 0] /Rect [246.842 515.805 277.055 526.596] /A << /S /GoTo /D <74686D2E332E332E322E32> >> >>}
\MusicPDFWriteObject{3181}{<< /Type /Annot /Subtype /Link /Border [0 0 0] /H /I /C [1 0 0] /Rect [170.378 489.922 205.572 501.875] /A << /S /GoTo /D <6571756174696F6E2E332E342E332E3133> >> >>}
\MusicPDFWriteObject{3182}{<< /Type /Annot /Subtype /Link /Border [0 0 0] /H /I /C [1 0 0] /Rect [277.399 489.922 308.092 501.875] /A << /S /GoTo /D <74686D2E332E342E332E33> >> >>}
\MusicPDFWriteObject{3183}{<< /Type /Annot /Subtype /Link /Border [0 0 0] /H /I /C [1 0 0] /Rect [256.567 427 291.76 438.953] /A << /S /GoTo /D <6571756174696F6E2E332E342E332E3131> >> >>}
\MusicPDFWriteObject{3184}{<< /Type /Annot /Subtype /Link /Border [0 0 0] /H /I /C [1 0 0] /Rect [361.925 426.999 391.828 438.972] /A << /S /GoTo /D <74686D2E332E342E332E33> >> >>}
\MusicPDFWriteObject{3185}{<< /Type /Annot /Subtype /Link /Border [0 0 0] /H /I /C [1 0 0] /Rect [249.517 367.004 279.73 378.956] /A << /S /GoTo /D <6571756174696F6E2E332E332E312E32> >> >>}
\MusicPDFWriteObject{3186}{<< /Type /Annot /Subtype /Link /Border [0 0 0] /H /I /C [1 0 0] /Rect [346.154 367.004 375.803 378.956] /A << /S /GoTo /D <74686D2E332E332E312E31> >> >>}
\MusicPDFWriteObject{3187}{<< /Type /Annot /Subtype /Link /Border [0 0 0] /H /I /C [1 0 0] /Rect [320.846 276.426 356.04 288.378] /A << /S /GoTo /D <6571756174696F6E2E372E322E312E3137> >> >>}
\MusicPDFWriteObject{3188}{<< /Type /Annot /Subtype /Link /Border [0 0 0] /H /I /C [1 0 0] /Rect [238.307 239.949 273.5 251.902] /A << /S /GoTo /D <6571756174696F6E2E372E322E312E3138> >> >>}
\MusicPDFWriteObject{3189}{<< /Type /Annot /Subtype /Link /Border [0 0 0] /H /I /C [1 0 0] /Rect [178.112 226.798 208.833 238.771] /A << /S /GoTo /D <74686D2E332E342E332E33> >> >>}
\MusicPDFWriteObject{3190}{<< /Type /Annot /Subtype /Link /Border [0 0 0] /H /I /C [1 0 0] /Rect [128.789 118.411 164.049 132.071] /A << /S /GoTo /D <6571756174696F6E2E372E322E312E3138> >> >>}
\MusicPDFWriteObject{3191}{<< /Type /Annot /Subtype /Link /Border [0 0 0] /H /I /C [1 0 0] /Rect [193.691 631.767 228.884 643.72] /A << /S /GoTo /D <6571756174696F6E2E372E322E312E3231> >> >>}
\MusicPDFWriteObject{3192}{<< /Type /Annot /Subtype /Link /Border [0 0 0] /H /I /C [1 0 0] /Rect [398.033 618.616 427.681 630.571] /A << /S /GoTo /D <74686D2E332E332E322E31> >> >>}
\MusicPDFWriteObject{3193}{<< /Type /Annot /Subtype /Link /Border [0 0 0] /H /I /C [1 0 0] /Rect [482.147 618.616 511.796 630.571] /A << /S /GoTo /D <74686D2E332E332E322E32> >> >>}
\MusicPDFWriteObject{3194}{<< /Type /Annot /Subtype /Link /Border [0 0 0] /H /I /C [1 0 0] /Rect [101.261 516.116 136.454 528.069] /A << /S /GoTo /D <6571756174696F6E2E372E322E312E3232> >> >>}
\MusicPDFWriteObject{3195}{<< /Type /Annot /Subtype /Link /Border [0 0 0] /H /I /C [1 0 0] /Rect [218.294 516.116 253.487 528.069] /A << /S /GoTo /D <6571756174696F6E2E372E322E312E3133> >> >>}
\MusicPDFWriteObject{3196}{<< /Type /Annot /Subtype /Link /Border [0 0 0] /H /I /C [1 0 0] /Rect [133.297 502.966 168.49 514.918] /A << /S /GoTo /D <6571756174696F6E2E372E322E312E3133> >> >>}
\MusicPDFWriteObject{3197}{<< /Type /Annot /Subtype /Link /Border [0 0 0] /H /I /C [1 0 0] /Rect [340.224 476.664 375.417 488.617] /A << /S /GoTo /D <6571756174696F6E2E372E322E312E3139> >> >>}
\MusicPDFWriteObject{3198}{<< /Type /Annot /Subtype /Link /Border [0 0 0] /H /I /C [1 0 0] /Rect [164.497 313.533 199.69 325.485] /A << /S /GoTo /D <6571756174696F6E2E372E322E312E3236> >> >>}
\MusicPDFWriteObject{3199}{<< /Type /Annot /Subtype /Link /Border [0 0 0] /H /I /C [1 0 0] /Rect [210.806 259.054 246 271.007] /A << /S /GoTo /D <6571756174696F6E2E372E322E312E3237> >> >>}
\MusicPDFWriteObject{3200}{<< /Type /Annot /Subtype /Link /Border [0 0 0] /H /I /C [1 0 0] /Rect [164.237 245.904 199.43 257.856] /A << /S /GoTo /D <6571756174696F6E2E372E322E312E3134> >> >>}
\MusicPDFWriteObject{3201}{<< /Type /Annot /Subtype /Link /Border [0 0 0] /H /I /C [1 0 0] /Rect [137.631 140.712 172.824 152.665] /A << /S /GoTo /D <6571756174696F6E2E372E322E312E3131> >> >>}
\MusicPDFWriteObject{3202}{<< /Type /Annot /Subtype /Link /Border [0 0 0] /H /I /C [1 0 0] /Rect [265.384 671.482 295.315 683.437] /A << /S /GoTo /D <74686D2E332E332E322E31> >> >>}
\MusicPDFWriteObject{3203}{<< /Type /Annot /Subtype /Link /Border [0 0 0] /H /I /C [1 0 0] /Rect [351.092 671.482 381.023 683.437] /A << /S /GoTo /D <74686D2E332E332E322E32> >> >>}
\MusicPDFWriteObject{3204}{<< /Type /Annot /Subtype /Link /Border [0 0 0] /H /I /C [1 0 0] /Rect [288.889 497.037 319.102 508.99] /A << /S /GoTo /D <6571756174696F6E2E352E322E332E31> >> >>}
\MusicPDFWriteObject{3205}{<< /Type /Annot /Subtype /Link /Border [0 0 0] /H /I /C [1 0 0] /Rect [381.882 496.677 413.992 508.99] /A << /S /GoTo /D <74686D2E352E322E332E31> >> >>}
\MusicPDFWriteObject{3206}{<< /Type /Annot /Subtype /Link /Border [0 0 0] /H /I /C [1 0 0] /Rect [257.357 376.492 287.457 390.8] /A << /S /GoTo /D <74686D2E372E322E312E34> >> >>}
\MusicPDFWriteObject{3207}{<< /Type /Annot /Subtype /Link /Border [0 0 0] /H /I /C [1 0 0] /Rect [298.181 326.053 328.394 338.006] /A << /S /GoTo /D <6571756174696F6E2E332E322E322E32> >> >>}
\MusicPDFWriteObject{3208}{<< /Type /Annot /Subtype /Link /Border [0 0 0] /H /I /C [1 0 0] /Rect [387.02 326.053 417.233 338.006] /A << /S /GoTo /D <74686D2E332E322E322E31> >> >>}
\MusicPDFWriteObject{3209}{<< /Type /Annot /Subtype /Link /Border [0 0 0] /H /I /C [1 0 0] /Rect [154.52 212.267 186.202 223.058] /A << /S /GoTo /D <74686D2E422E332E322E31> >> >>}
\MusicPDFWriteObject{3210}{<< /Type /Annot /Subtype /Link /Border [0 0 0] /H /I /C [1 0 0] /Rect [198.077 172.261 228.29 184.216] /A << /S /GoTo /D <74686D2E362E332E322E31> >> >>}
\MusicPDFWriteObject{3211}{<< /Type /Annot /Subtype /Link /Border [0 0 0] /H /I /C [1 0 0] /Rect [140.121 129.793 170.334 144.764] /A << /S /GoTo /D <74686D2E352E332E312E35> >> >>}
\MusicPDFWriteObject{3212}{<< /Type /Annot /Subtype /Link /Border [0 0 0] /H /I /C [1 0 0] /Rect [215.024 118.817 244.673 130.79] /A << /S /GoTo /D <74686D2E372E322E312E31> >> >>}
\MusicPDFWriteObject{3213}{<< /Type /Annot /Subtype /Link /Border [0 0 0] /H /I /C [1 0 0] /Rect [277.944 593.497 308.157 605.45] /A << /S /GoTo /D <6571756174696F6E2E352E322E332E31> >> >>}
\MusicPDFWriteObject{3214}{<< /Type /Annot /Subtype /Link /Border [0 0 0] /H /I /C [1 0 0] /Rect [366.673 593.137 398.192 605.45] /A << /S /GoTo /D <74686D2E352E322E332E31> >> >>}
\MusicPDFWriteObject{3215}{<< /Type /Annot /Subtype /Link /Border [0 0 0] /H /I /C [1 0 0] /Rect [323.99 466.763 354.768 478.718] /A << /S /GoTo /D <74686D2E372E322E322E31> >> >>}
\MusicPDFWriteObject{3216}{<< /Type /Annot /Subtype /Link /Border [0 0 0] /H /I /C [1 0 0] /Rect [301.275 398.772 330.924 409.562] /A << /S /GoTo /D <74686D2E372E312E342E32> >> >>}
\MusicPDFWriteObject{3217}{<< /Type /Annot /Subtype /Link /Border [0 0 0] /H /I /C [1 0 0] /Rect [211.154 196.709 241.932 209.023] /A << /S /GoTo /D <74686D2E342E342E322E33> >> >>}
\MusicPDFWriteObject{3218}{<< /Type /Annot /Subtype /Link /Border [0 0 0] /H /I /C [1 0 0] /Rect [360.319 95.747 390.222 107.699] /A << /S /GoTo /D <74686D2E372E312E322E31> >> >>}
\MusicPDFWriteObject{3219}{<< /Type /Annot /Subtype /Link /Border [0 0 0] /H /I /C [1 0 0] /Rect [291.062 645.734 322.581 656.525] /A << /S /GoTo /D <74686D2E372E332E312E31> >> >>}
\MusicPDFWriteObject{3220}{<< /Type /Annot /Subtype /Link /Border [0 0 0] /H /I /C [1 0 0] /Rect [241.4 405.651 271.613 417.605] /A << /S /GoTo /D <74686D2E372E322E312E31> >> >>}
\MusicPDFWriteObject{3221}{<< /Type /Annot /Subtype /Link /Border [0 0 0] /H /I /C [1 0 0] /Rect [212.26 392.501 243.038 404.456] /A << /S /GoTo /D <74686D2E372E322E312E31> >> >>}
\MusicPDFWriteObject{3222}{<< /Type /Annot /Subtype /Link /Border [0 0 0] /H /I /C [1 0 0] /Rect [220.633 167.778 250.846 179.731] /A << /S /GoTo /D <6571756174696F6E2E372E332E312E32> >> >>}
\MusicPDFWriteObject{3223}{<< /Type /Annot /Subtype /Link /Border [0 0 0] /H /I /C [1 0 0] /Rect [448.593 167.778 478.806 179.731] /A << /S /GoTo /D <6571756174696F6E2E372E332E312E32> >> >>}
\MusicPDFWriteObject{3224}{<< /Type /Annot /Subtype /Link /Border [0 0 0] /H /I /C [0 1 0] /Rect [107.748 586.168 125.946 595.036] /A << /S /GoTo /D <636974652E4A6172> >> >>}
\MusicPDFWriteObject{3225}{<< /Type /Annot /Subtype /Link /Border [0 0 0] /H /I /C [1 0 0] /Rect [106.556 224.638 136.769 236.953] /A << /S /GoTo /D <74686D2E352E332E342E33> >> >>}
\MusicPDFWriteObject{3226}{<< /Type /Annot /Subtype /Link /Border [0 0 0] /H /I /C [1 0 0] /Rect [167.366 198.337 197.579 210.652] /A << /S /GoTo /D <74686D2E352E332E342E32> >> >>}
\MusicPDFWriteObject{3227}{<< /Type /Annot /Subtype /Link /Border [0 0 0] /H /I /C [1 0 0] /Rect [334.424 172.396 364.637 184.348] /A << /S /GoTo /D <74686D2E352E332E312E35> >> >>}
\MusicPDFWriteObject{3228}{<< /Type /Annot /Subtype /Link /Border [0 0 0] /H /I /C [1 0 0] /Rect [290.52 637.979 320.733 646.833] /A << /S /GoTo /D <74686D2E372E332E312E31> >> >>}
\MusicPDFWriteObject{3229}{<< /Type /Annot /Subtype /Link /Border [0 0 0] /H /I /C [1 0 0] /Rect [350.213 548.882 380.426 560.835] /A << /S /GoTo /D <6571756174696F6E2E372E322E322E32> >> >>}
\MusicPDFWriteObject{3230}{<< /Type /Annot /Subtype /Link /Border [0 0 0] /H /I /C [1 0 0] /Rect [395.918 535.731 426.131 547.684] /A << /S /GoTo /D <6571756174696F6E2E372E322E322E31> >> >>}
\MusicPDFWriteObject{3231}{<< /Type /Annot /Subtype /Link /Border [0 0 0] /H /I /C [0 1 0] /Rect [223.735 498.769 241.208 507.637] /A << /S /GoTo /D <636974652E4346> >> >>}
\MusicPDFWriteObject{3232}{<< /Type /Annot /Subtype /Link /Border [0 0 0] /H /I /C [1 0 0] /Rect [270.889 483.128 300.538 495.081] /A << /S /GoTo /D <74686D2E332E332E312E35> >> >>}
\MusicPDFWriteObject{3233}{<< /Type /Annot /Subtype /Link /Border [0 0 0] /H /I /C [1 0 0] /Rect [290.52 423.323 320.733 432.177] /A << /S /GoTo /D <74686D2E372E332E312E31> >> >>}
\MusicPDFWriteObject{3234}{<< /Type /Annot /Subtype /Link /Border [0 0 0] /H /I /C [1 0 0] /Rect [242.186 615.371 272.963 627.624] /A << /S /GoTo /D <74686D2E342E312E322E34> >> >>}
\MusicPDFWriteObject{3235}{<< /Type /Annot /Subtype /Link /Border [0 0 0] /H /I /C [0 1 0] /Rect [285.01 591.858 310.07 600.725] /A << /S /GoTo /D <636974652E53636861> >> >>}
\MusicPDFWriteObject{3236}{<< /Type /Annot /Subtype /Link /Border [0 0 0] /H /I /C [0 1 0] /Rect [161.921 578.707 181.025 587.575] /A << /S /GoTo /D <636974652E4E42> >> >>}
\MusicPDFWriteObject{3237}{<< /Type /Annot /Subtype /Link /Border [0 0 0] /H /I /C [0 1 0] /Rect [97.546 565.556 115.745 574.424] /A << /S /GoTo /D <636974652E4A6172> >> >>}
\MusicPDFWriteObject{3238}{<< /Type /Annot /Subtype /Link /Border [0 0 0] /H /I /C [0 1 0] /Rect [524.974 552.405 543.173 561.273] /A << /S /GoTo /D <636974652E4A6172> >> >>}
\MusicPDFWriteObject{3239}{<< /Type /Annot /Subtype /Link /Border [0 0 0] /H /I /C [1 0 0] /Rect [348.149 510.463 378.363 522.416] /A << /S /GoTo /D <6571756174696F6E2E372E332E312E39> >> >>}
\MusicPDFWriteObject{3240}{<< /Type /Annot /Subtype /Link /Border [0 0 0] /H /I /C [1 0 0] /Rect [512.598 510.463 543.178 522.416] /A << /S /GoTo /D <74686D2E342E312E322E34> >> >>}
\MusicPDFWriteObject{3241}{<< /Type /Annot /Subtype /Link /Border [0 0 0] /H /I /C [1 0 0] /Rect [358.991 478.183 388.64 490.138] /A << /S /GoTo /D <74686D2E372E312E322E32> >> >>}
\MusicPDFWriteObject{3242}{<< /Type /Annot /Subtype /Link /Border [0 0 0] /H /I /C [1 0 0] /Rect [324.389 463.423 354.602 476.988] /A << /S /GoTo /D <74686D2E372E332E312E31> >> >>}
\MusicPDFWriteObject{3243}{<< /Type /Annot /Subtype /Link /Border [0 0 0] /H /I /C [1 0 0] /Rect [291.062 271.767 322.581 282.557] /A << /S /GoTo /D <74686D2E372E332E312E31> >> >>}
\MusicPDFWriteObject{3244}{<< /Type /Annot /Subtype /Link /Border [0 0 0] /H /I /C [1 0 0] /Rect [167.899 670.988 203.092 682.941] /A << /S /GoTo /D <6571756174696F6E2E372E332E312E3132> >> >>}
\MusicPDFWriteObject{3245}{<< /Type /Annot /Subtype /Link /Border [0 0 0] /H /I /C [1 0 0] /Rect [268.399 670.988 303.592 682.941] /A << /S /GoTo /D <6571756174696F6E2E372E332E312E3133> >> >>}
\MusicPDFWriteObject{3246}{<< /Type /Annot /Subtype /Link /Border [0 0 0] /H /I /C [1 0 0] /Rect [159.704 274.4 190.481 291.812] /A << /S /GoTo /D <74686D2E322E342E332E31> >> >>}
\MusicPDFWriteObject{3247}{<< /Type /Annot /Subtype /Link /Border [0 0 0] /H /I /C [1 0 0] /Rect [437.194 176.174 470.036 191.145] /A << /S /GoTo /D <74686D2E442E332E362E31> >> >>}
\MusicPDFWriteObject{3248}{<< /Type /Annot /Subtype /Link /Border [0 0 0] /H /I /C [1 0 0] /Rect [200.727 162.2 234.185 177.171] /A << /S /GoTo /D <74686D2E442E332E332E36> >> >>}
\MusicPDFWriteObject{3249}{<< /Type /Annot /Subtype /Link /Border [0 0 0] /H /I /C [1 0 0] /Rect [158.188 133.245 188.401 150.046] /A << /S /GoTo /D <74686D2E352E332E342E34> >> >>}
\MusicPDFWriteObject{3250}{<< /Type /Annot /Subtype /Link /Border [0 0 0] /H /I /C [1 0 0] /Rect [207.459 95.747 237.673 107.699] /A << /S /GoTo /D <6571756174696F6E2E352E332E342E31> >> >>}
\MusicPDFWriteObject{3251}{<< /Type /Annot /Subtype /Link /Border [0 0 0] /H /I /C [1 0 0] /Rect [316.279 95.747 347.057 107.699] /A << /S /GoTo /D <74686D2E352E332E342E31> >> >>}
\MusicPDFWriteObject{3252}{<< /Type /Annot /Subtype /Link /Border [0 0 0] /H /I /C [1 0 0] /Rect [226.054 684.633 261.248 696.586] /A << /S /GoTo /D <6571756174696F6E2E372E332E312E3138> >> >>}
\MusicPDFWriteObject{3253}{<< /Type /Annot /Subtype /Link /Border [0 0 0] /H /I /C [1 0 0] /Rect [449.8 684.633 480.577 696.586] /A << /S /GoTo /D <74686D2E352E332E342E33> >> >>}
\MusicPDFWriteObject{3254}{<< /Type /Annot /Subtype /Link /Border [0 0 0] /H /I /C [1 0 0] /Rect [317.851 658.332 348.064 670.284] /A << /S /GoTo /D <74686D2E352E332E342E32> >> >>}
\MusicPDFWriteObject{3255}{<< /Type /Annot /Subtype /Link /Border [0 0 0] /H /I /C [1 0 0] /Rect [296.569 645.734 326.782 656.525] /A << /S /GoTo /D <74686D2E352E332E342E34> >> >>}
\MusicPDFWriteObject{3256}{<< /Type /Annot /Subtype /Link /Border [0 0 0] /H /I /C [1 0 0] /Rect [169.361 602.8 199.574 614.752] /A << /S /GoTo /D <74686D2E352E332E342E32> >> >>}
\MusicPDFWriteObject{3257}{<< /Type /Annot /Subtype /Link /Border [0 0 0] /H /I /C [1 0 0] /Rect [270.512 602.8 303.354 614.752] /A << /S /GoTo /D <74686D2E442E332E332E36> >> >>}
\MusicPDFWriteObject{3258}{<< /Type /Annot /Subtype /Link /Border [0 0 0] /H /I /C [1 0 0] /Rect [98.422 521.109 128.635 533.062] /A << /S /GoTo /D <6571756174696F6E2E352E322E342E31> >> >>}
\MusicPDFWriteObject{3259}{<< /Type /Annot /Subtype /Link /Border [0 0 0] /H /I /C [1 0 0] /Rect [187.261 521.109 217.475 533.062] /A << /S /GoTo /D <74686D2E352E322E342E31> >> >>}
\MusicPDFWriteObject{3260}{<< /Type /Annot /Subtype /Link /Border [0 0 0] /H /I /C [1 0 0] /Rect [422.35 521.109 452.563 533.062] /A << /S /GoTo /D <74686D2E352E322E342E31> >> >>}
\MusicPDFWriteObject{3261}{<< /Type /Annot /Subtype /Link /Border [0 0 0] /H /I /C [1 0 0] /Rect [515.351 454.215 545.592 469.186] /A << /S /GoTo /D <74686D2E342E362E312E31> >> >>}
\MusicPDFWriteObject{3262}{<< /Type /Annot /Subtype /Link /Border [0 0 0] /H /I /C [1 0 0] /Rect [347.511 442.897 377.894 455.211] /A << /S /GoTo /D <74686D2E342E362E312E35> >> >>}
\MusicPDFWriteObject{3263}{<< /Type /Annot /Subtype /Link /Border [0 0 0] /H /I /C [1 0 0] /Rect [230.566 416.955 260.779 428.908] /A << /S /GoTo /D <6571756174696F6E2E342E362E312E33> >> >>}
\MusicPDFWriteObject{3264}{<< /Type /Annot /Subtype /Link /Border [0 0 0] /H /I /C [1 0 0] /Rect [325.077 416.595 355.29 428.908] /A << /S /GoTo /D <74686D2E342E362E312E33> >> >>}
\MusicPDFWriteObject{3265}{<< /Type /Annot /Subtype /Link /Border [0 0 0] /H /I /C [1 0 0] /Rect [348.04 404.358 378.253 415.149] /A << /S /GoTo /D <74686D2E342E362E332E31> >> >>}
\MusicPDFWriteObject{3266}{<< /Type /Annot /Subtype /Link /Border [0 0 0] /H /I /C [1 0 0] /Rect [319.907 271.213 350.12 283.165] /A << /S /GoTo /D <6571756174696F6E2E352E312E312E32> >> >>}
\MusicPDFWriteObject{3267}{<< /Type /Annot /Subtype /Link /Border [0 0 0] /H /I /C [1 0 0] /Rect [203.739 258.062 233.952 270.014] /A << /S /GoTo /D <6571756174696F6E2E352E332E312E36> >> >>}
\MusicPDFWriteObject{3268}{<< /Type /Annot /Subtype /Link /Border [0 0 0] /H /I /C [1 0 0] /Rect [296.508 257.093 326.523 270.014] /A << /S /GoTo /D <74686D2E352E332E312E34> >> >>}
\MusicPDFWriteObject{3269}{<< /Type /Annot /Subtype /Link /Border [0 0 0] /H /I /C [1 0 0] /Rect [346.68 244.911 381.873 256.864] /A << /S /GoTo /D <6571756174696F6E2E372E332E312E3139> >> >>}
\MusicPDFWriteObject{3270}{<< /Type /Annot /Subtype /Link /Border [0 0 0] /H /I /C [1 0 0] /Rect [384.052 215.736 413.701 229.741] /A << /S /GoTo /D <74686D2E342E362E332E31> >> >>}
\MusicPDFWriteObject{3271}{<< /Type /Annot /Subtype /Link /Border [0 0 0] /H /I /C [1 0 0] /Rect [213.011 188.966 248.204 200.919] /A << /S /GoTo /D <6571756174696F6E2E372E332E312E3139> >> >>}
\MusicPDFWriteObject{3272}{<< /Type /Annot /Subtype /Link /Border [0 0 0] /H /I /C [1 0 0] /Rect [495.213 188.966 530.406 200.919] /A << /S /GoTo /D <6571756174696F6E2E372E332E312E3139> >> >>}
\MusicPDFWriteObject{3273}{<< /Type /Annot /Subtype /Link /Border [0 0 0] /H /I /C [0 1 0] /Rect [274.303 178.305 293.407 187.173] /A << /S /GoTo /D <636974652E4E42> >> >>}
\MusicPDFWriteObject{3274}{<< /Type /Annot /Subtype /Link /Border [0 0 0] /H /I /C [1 0 0] /Rect [131.585 162.665 162.362 174.617] /A << /S /GoTo /D <74686D2E342E362E312E35> >> >>}
\MusicPDFWriteObject{3275}{<< /Type /Annot /Subtype /Link /Border [0 0 0] /H /I /C [1 0 0] /Rect [220.023 162.665 250.801 174.617] /A << /S /GoTo /D <74686D2E372E332E312E35> >> >>}
\MusicPDFWriteObject{3276}{<< /Type /Annot /Subtype /Link /Border [0 0 0] /H /I /C [1 0 0] /Rect [159.098 149.514 189.847 161.467] /A << /S /GoTo /D <74686D2E352E322E332E31> >> >>}
\MusicPDFWriteObject{3277}{<< /Type /Annot /Subtype /Link /Border [0 0 0] /H /I /C [1 0 0] /Rect [333.566 149.514 366.994 161.467] /A << /S /GoTo /D <74686D2E442E332E332E37> >> >>}
\MusicPDFWriteObject{3278}{<< /Type /Annot /Subtype /Link /Border [0 0 0] /H /I /C [1 0 0] /Rect [337.878 136.363 373.071 148.316] /A << /S /GoTo /D <6571756174696F6E2E372E332E312E3133> >> >>}
\MusicPDFWriteObject{3279}{<< /Type /Annot /Subtype /Link /Border [0 0 0] /H /I /C [1 0 0] /Rect [506.695 136.363 541.888 148.316] /A << /S /GoTo /D <6571756174696F6E2E372E332E312E3137> >> >>}
\MusicPDFWriteObject{3280}{<< /Type /Annot /Subtype /Link /Border [0 0 0] /H /I /C [1 0 0] /Rect [93.688 122.391 128.881 134.344] /A << /S /GoTo /D <6571756174696F6E2E342E362E312E3137> >> >>}
\MusicPDFWriteObject{3281}{<< /Type /Annot /Subtype /Link /Border [0 0 0] /H /I /C [1 0 0] /Rect [442.646 109.24 477.839 121.193] /A << /S /GoTo /D <6571756174696F6E2E372E332E312E3133> >> >>}
\MusicPDFWriteObject{3282}{<< /Type /Annot /Subtype /Link /Border [0 0 0] /H /I /C [1 0 0] /Rect [245.939 684.633 281.132 696.586] /A << /S /GoTo /D <6571756174696F6E2E342E362E312E3134> >> >>}
\MusicPDFWriteObject{3283}{<< /Type /Annot /Subtype /Link /Border [0 0 0] /H /I /C [1 0 0] /Rect [112.018 672.036 141.723 682.827] /A << /S /GoTo /D <74686D2E342E362E312E35> >> >>}
\MusicPDFWriteObject{3284}{<< /Type /Annot /Subtype /Link /Border [0 0 0] /H /I /C [1 0 0] /Rect [513.559 672.036 543.264 682.827] /A << /S /GoTo /D <74686D2E342E312E322E34> >> >>}
\MusicPDFWriteObject{3285}{<< /Type /Annot /Subtype /Link /Border [0 0 0] /H /I /C [1 0 0] /Rect [388.312 237.724 418.525 249.677] /A << /S /GoTo /D <6571756174696F6E2E352E322E342E31> >> >>}
\MusicPDFWriteObject{3286}{<< /Type /Annot /Subtype /Link /Border [0 0 0] /H /I /C [1 0 0] /Rect [318.239 95.747 353.432 107.699] /A << /S /GoTo /D <6571756174696F6E2E372E332E312E3230> >> >>}
\MusicPDFWriteObject{3287}{<< /Type /Annot /Subtype /Link /Border [0 0 0] /H /I /C [1 0 0] /Rect [460.432 95.747 490.279 107.699] /A << /S /GoTo /D <74686D2E352E332E342E34> >> >>}
\MusicPDFWriteObject{3288}{<< /Type /Annot /Subtype /Link /Border [0 0 0] /H /I /C [1 0 0] /Rect [88.712 82.596 118.925 94.549] /A << /S /GoTo /D <74686D2E352E332E342E32> >> >>}
\MusicPDFWriteObject{3289}{<< /Type /Annot /Subtype /Link /Border [0 0 0] /H /I /C [1 0 0] /Rect [204.97 670.633 240.163 682.586] /A << /S /GoTo /D <6571756174696F6E2E372E332E312E3230> >> >>}
\MusicPDFWriteObject{3290}{<< /Type /Annot /Subtype /Link /Border [0 0 0] /H /I /C [1 0 0] /Rect [348.314 670.633 378.527 682.586] /A << /S /GoTo /D <74686D2E352E332E342E34> >> >>}
\MusicPDFWriteObject{3291}{<< /Type /Annot /Subtype /Link /Border [0 0 0] /H /I /C [1 0 0] /Rect [434.774 670.633 464.987 682.586] /A << /S /GoTo /D <74686D2E352E332E342E32> >> >>}
\MusicPDFWriteObject{3292}{<< /Type /Annot /Subtype /Link /Border [0 0 0] /H /I /C [1 0 0] /Rect [316.285 557.156 346.498 569.109] /A << /S /GoTo /D <6571756174696F6E2E332E332E312E32> >> >>}
\MusicPDFWriteObject{3293}{<< /Type /Annot /Subtype /Link /Border [0 0 0] /H /I /C [1 0 0] /Rect [404.224 557.156 433.873 569.109] /A << /S /GoTo /D <74686D2E332E332E312E31> >> >>}
\MusicPDFWriteObject{3294}{<< /Type /Annot /Subtype /Link /Border [0 0 0] /H /I /C [1 0 0] /Rect [91.42 544.005 121.633 555.958] /A << /S /GoTo /D <6571756174696F6E2E332E342E312E36> >> >>}
\MusicPDFWriteObject{3295}{<< /Type /Annot /Subtype /Link /Border [0 0 0] /H /I /C [1 0 0] /Rect [180.26 544.005 210.473 555.958] /A << /S /GoTo /D <74686D2E332E342E312E31> >> >>}
\MusicPDFWriteObject{3296}{<< /Type /Annot /Subtype /Link /Border [0 0 0] /H /I /C [1 0 0] /Rect [414.988 544.005 450.182 555.958] /A << /S /GoTo /D <6571756174696F6E2E372E332E312E3233> >> >>}
\MusicPDFWriteObject{3297}{<< /Type /Annot /Subtype /Link /Border [0 0 0] /H /I /C [1 0 0] /Rect [143.37 462.477 173.019 475.253] /A << /S /GoTo /D <74686D2E332E332E312E31> >> >>}
\MusicPDFWriteObject{3298}{<< /Type /Annot /Subtype /Link /Border [0 0 0] /H /I /C [1 0 0] /Rect [437.005 433.52 467.783 448.491] /A << /S /GoTo /D <74686D2E352E332E342E34> >> >>}
\MusicPDFWriteObject{3299}{<< /Type /Annot /Subtype /Link /Border [0 0 0] /H /I /C [1 0 0] /Rect [215.078 419.546 248.074 434.517] /A << /S /GoTo /D <74686D2E442E332E332E36> >> >>}
\MusicPDFWriteObject{3300}{<< /Type /Annot /Subtype /Link /Border [0 0 0] /H /I /C [1 0 0] /Rect [304.576 419.546 334.93 434.517] /A << /S /GoTo /D <74686D2E352E332E342E32> >> >>}
\MusicPDFWriteObject{3301}{<< /Type /Annot /Subtype /Link /Border [0 0 0] /H /I /C [1 0 0] /Rect [511.387 390.314 544.845 406.568] /A << /S /GoTo /D <74686D2E442E332E332E36> >> >>}
\MusicPDFWriteObject{3302}{<< /Type /Annot /Subtype /Link /Border [0 0 0] /H /I /C [1 0 0] /Rect [191.496 379.357 221.709 391.31] /A << /S /GoTo /D <74686D2E352E332E342E32> >> >>}
\MusicPDFWriteObject{3303}{<< /Type /Annot /Subtype /Link /Border [0 0 0] /H /I /C [1 0 0] /Rect [450.647 285.331 485.84 297.284] /A << /S /GoTo /D <6571756174696F6E2E372E332E312E3230> >> >>}
\MusicPDFWriteObject{3304}{<< /Type /Annot /Subtype /Link /Border [0 0 0] /H /I /C [1 0 0] /Rect [320.457 258.208 351.235 270.161] /A << /S /GoTo /D <74686D2E352E342E312E31> >> >>}
\MusicPDFWriteObject{3305}{<< /Type /Annot /Subtype /Link /Border [0 0 0] /H /I /C [1 0 0] /Rect [326.918 231.906 357.131 243.859] /A << /S /GoTo /D <6571756174696F6E2E352E312E312E32> >> >>}
\MusicPDFWriteObject{3306}{<< /Type /Annot /Subtype /Link /Border [0 0 0] /H /I /C [1 0 0] /Rect [170.358 150.41 205.551 162.362] /A << /S /GoTo /D <6571756174696F6E2E372E332E312E3234> >> >>}
\MusicPDFWriteObject{3307}{<< /Type /Annot /Subtype /Link /Border [0 0 0] /H /I /C [1 0 0] /Rect [473.734 102.829 503.382 116.393] /A << /S /GoTo /D <74686D2E332E332E312E31> >> >>}
\MusicPDFWriteObject{3308}{<< /Type /Annot /Subtype /Link /Border [0 0 0] /H /I /C [1 0 0] /Rect [164.359 643.051 194.572 655.364] /A << /S /GoTo /D <74686D2E372E332E312E32> >> >>}
\MusicPDFWriteObject{3309}{<< /Type /Annot /Subtype /Link /Border [0 0 0] /H /I /C [1 0 0] /Rect [258.448 603.721 293.641 615.673] /A << /S /GoTo /D <6571756174696F6E2E372E332E312E3230> >> >>}
\MusicPDFWriteObject{3310}{<< /Type /Annot /Subtype /Link /Border [0 0 0] /H /I /C [1 0 0] /Rect [513.621 577.419 543.27 589.372] /A << /S /GoTo /D <74686D2E352E342E312E31> >> >>}
\MusicPDFWriteObject{3311}{<< /Type /Annot /Subtype /Link /Border [0 0 0] /H /I /C [1 0 0] /Rect [91.42 520.817 121.633 532.77] /A << /S /GoTo /D <6571756174696F6E2E352E312E312E32> >> >>}
\MusicPDFWriteObject{3312}{<< /Type /Annot /Subtype /Link /Border [0 0 0] /H /I /C [1 0 0] /Rect [306.118 480.023 336.331 491.976] /A << /S /GoTo /D <6571756174696F6E2E352E342E322E36> >> >>}
\MusicPDFWriteObject{3313}{<< /Type /Annot /Subtype /Link /Border [0 0 0] /H /I /C [1 0 0] /Rect [447.704 480.023 478.481 491.976] /A << /S /GoTo /D <74686D2E352E342E322E32> >> >>}
\MusicPDFWriteObject{3314}{<< /Type /Annot /Subtype /Link /Border [0 0 0] /H /I /C [1 0 0] /Rect [174.024 264.969 207.112 279.939] /A << /S /GoTo /D <74686D2E442E332E332E37> >> >>}
\MusicPDFWriteObject{3315}{<< /Type /Annot /Subtype /Link /Border [0 0 0] /H /I /C [0 1 0] /Rect [524.993 270.475 543.192 279.343] /A << /S /GoTo /D <636974652E4A6172> >> >>}
\MusicPDFWriteObject{3316}{<< /Type /Annot /Subtype /Link /Border [0 0 0] /H /I /C [1 0 0] /Rect [411.91 253.652 441.982 265.965] /A << /S /GoTo /D <74686D2E342E362E312E33> >> >>}
\MusicPDFWriteObject{3317}{<< /Type /Annot /Subtype /Link /Border [0 0 0] /H /I /C [1 0 0] /Rect [308.33 137.585 337.979 152.555] /A << /S /GoTo /D <74686D2E342E362E332E31> >> >>}
\MusicPDFWriteObject{3318}{<< /Type /Annot /Subtype /Link /Border [0 0 0] /H /I /C [1 0 0] /Rect [427.563 126.628 462.756 138.581] /A << /S /GoTo /D <6571756174696F6E2E372E332E312E3236> >> >>}
\MusicPDFWriteObject{3319}{<< /Type /Annot /Subtype /Link /Border [0 0 0] /H /I /C [1 0 0] /Rect [393.865 99.503 429.058 111.456] /A << /S /GoTo /D <6571756174696F6E2E372E332E312E3132> >> >>}
\MusicPDFWriteObject{3320}{<< /Type /Annot /Subtype /Link /Border [0 0 0] /H /I /C [1 0 0] /Rect [148.002 82.595 178.215 97.484] /A << /S /GoTo /D <74686D2E372E322E312E34> >> >>}
\MusicPDFWriteObject{3321}{<< /Type /Annot /Subtype /Link /Border [0 0 0] /H /I /C [1 0 0] /Rect [297.514 658.885 328.291 669.676] /A << /S /GoTo /D <74686D2E372E322E312E34> >> >>}
\MusicPDFWriteObject{3322}{<< /Type /Annot /Subtype /Link /Border [0 0 0] /H /I /C [1 0 0] /Rect [194.831 618.879 230.024 630.832] /A << /S /GoTo /D <6571756174696F6E2E372E322E312E3132> >> >>}
\MusicPDFWriteObject{3323}{<< /Type /Annot /Subtype /Link /Border [0 0 0] /H /I /C [1 0 0] /Rect [492.142 562.822 524.067 573.613] /A << /S /GoTo /D <74686D2E422E352E332E33> >> >>}
\MusicPDFWriteObject{3324}{<< /Type /Annot /Subtype /Link /Border [0 0 0] /H /I /C [1 0 0] /Rect [487.49 531.705 517.138 542.496] /A << /S /GoTo /D <74686D2E372E332E312E31> >> >>}
\MusicPDFWriteObject{3325}{<< /Type /Annot /Subtype /Link /Border [0 0 0] /H /I /C [1 0 0] /Rect [472.962 475.533 503.175 490.504] /A << /S /GoTo /D <74686D2E372E332E312E36> >> >>}
\MusicPDFWriteObject{3326}{<< /Type /Annot /Subtype /Link /Border [0 0 0] /H /I /C [1 0 0] /Rect [380.97 462.465 416.163 474.418] /A << /S /GoTo /D <6571756174696F6E2E372E332E312E3132> >> >>}
\MusicPDFWriteObject{3327}{<< /Type /Annot /Subtype /Link /Border [0 0 0] /H /I /C [1 0 0] /Rect [384.44 435.95 419.634 447.903] /A << /S /GoTo /D <6571756174696F6E2E372E332E312E3133> >> >>}
\MusicPDFWriteObject{3328}{<< /Type /Annot /Subtype /Link /Border [0 0 0] /H /I /C [1 0 0] /Rect [307.871 294.599 338.829 306.551] /A << /S /GoTo /D <74686D2E372E332E312E32> >> >>}
\MusicPDFWriteObject{3329}{<< /Type /Annot /Subtype /Link /Border [0 0 0] /H /I /C [0 1 0] /Rect [206.79 237.578 224.988 246.446] /A << /S /GoTo /D <636974652E4A6172> >> >>}
\MusicPDFWriteObject{3330}{<< /Type /Annot /Subtype /Link /Border [0 0 0] /H /I /C [1 0 0] /Rect [493.882 235.088 524.659 247.041] /A << /S /GoTo /D <74686D2E352E342E312E31> >> >>}
\MusicPDFWriteObject{3331}{<< /Type /Annot /Subtype /Link /Border [0 0 0] /H /I /C [1 0 0] /Rect [127.253 218.921 157.664 233.892] /A << /S /GoTo /D <74686D2E362E322E312E31> >> >>}
\MusicPDFWriteObject{3332}{<< /Type /Annot /Subtype /Link /Border [0 0 0] /H /I /C [1 0 0] /Rect [243.482 218.921 273.892 233.892] /A << /S /GoTo /D <74686D2E362E322E312E32> >> >>}
\MusicPDFWriteObject{3333}{<< /Type /Annot /Subtype /Link /Border [0 0 0] /H /I /C [0 1 0] /Rect [87.919 210.455 107.023 219.323] /A << /S /GoTo /D <636974652E4E42> >> >>}
\MusicPDFWriteObject{3334}{<< /Type /Annot /Subtype /Link /Border [0 0 0] /H /I /C [1 0 0] /Rect [499.948 207.965 530.725 219.918] /A << /S /GoTo /D <74686D2E352E322E332E31> >> >>}
\MusicPDFWriteObject{3335}{<< /Type /Annot /Subtype /Link /Border [0 0 0] /H /I /C [1 0 0] /Rect [363.216 178.647 393.034 193.618] /A << /S /GoTo /D <74686D2E352E322E332E31> >> >>}
\MusicPDFWriteObject{3336}{<< /Type /Annot /Subtype /Link /Border [0 0 0] /H /I /C [0 1 0] /Rect [525.054 184.153 543.253 193.021] /A << /S /GoTo /D <636974652E4A6172> >> >>}
\MusicPDFWriteObject{3337}{<< /Type /Annot /Subtype /Link /Border [0 0 0] /H /I /C [1 0 0] /Rect [257.426 631.817 287.639 643.77] /A << /S /GoTo /D <74686D2E372E332E312E36> >> >>}
\MusicPDFWriteObject{3338}{<< /Type /Annot /Subtype /Link /Border [0 0 0] /H /I /C [1 0 0] /Rect [386.437 419.724 416.651 431.677] /A << /S /GoTo /D <6571756174696F6E2E322E312E342E32> >> >>}
\MusicPDFWriteObject{3339}{<< /Type /Annot /Subtype /Link /Border [0 0 0] /H /I /C [1 0 0] /Rect [490.328 419.363 521.106 431.678] /A << /S /GoTo /D <74686D2E322E312E342E31> >> >>}
\MusicPDFWriteObject{3340}{<< /Type /Annot /Subtype /Link /Border [0 0 0] /H /I /C [1 0 0] /Rect [383.41 208.734 413.623 226.144] /A << /S /GoTo /D <74686D2E352E322E312E31> >> >>}
\MusicPDFWriteObject{3341}{<< /Type /Annot /Subtype /Link /Border [0 0 0] /H /I /C [0 1 0] /Rect [107.335 637.545 136.213 646.413] /A << /S /GoTo /D <636974652E5363682D61> >> >>}
\MusicPDFWriteObject{3342}{<< /Type /Annot /Subtype /Link /Border [0 0 0] /H /I /C [1 0 0] /Rect [324.06 621.904 359.253 633.857] /A << /S /GoTo /D <6571756174696F6E2E372E332E322E3130> >> >>}
\MusicPDFWriteObject{3343}{<< /Type /Annot /Subtype /Link /Border [0 0 0] /H /I /C [0 1 0] /Rect [87.718 611.243 105.917 620.111] /A << /S /GoTo /D <636974652E4A6172> >> >>}
\MusicPDFWriteObject{3344}{<< /Type /Annot /Subtype /Link /Border [0 0 0] /H /I /C [1 0 0] /Rect [204.14 554.787 234.353 566.74] /A << /S /GoTo /D <74686D2E322E322E332E32> >> >>}
\MusicPDFWriteObject{3345}{<< /Type /Annot /Subtype /Link /Border [0 0 0] /H /I /C [1 0 0] /Rect [225.058 472.893 255.271 484.846] /A << /S /GoTo /D <6571756174696F6E2E372E332E322E33> >> >>}
\MusicPDFWriteObject{3346}{<< /Type /Annot /Subtype /Link /Border [0 0 0] /H /I /C [1 0 0] /Rect [504.205 406.619 541.888 418.572] /A << /S /GoTo /D <6571756174696F6E2E412E322E332E3134> >> >>}
\MusicPDFWriteObject{3347}{<< /Type /Annot /Subtype /Link /Border [0 0 0] /H /I /C [1 0 0] /Rect [128.69 395.959 161.393 404.813] /A << /S /GoTo /D <74686D2E412E322E332E36> >> >>}
\MusicPDFWriteObject{3348}{<< /Type /Annot /Subtype /Link /Border [0 0 0] /H /I /C [1 0 0] /Rect [259.969 291.287 295.162 303.24] /A << /S /GoTo /D <6571756174696F6E2E372E332E322E3131> >> >>}
\MusicPDFWriteObject{3349}{<< /Type /Annot /Subtype /Link /Border [0 0 0] /H /I /C [1 0 0] /Rect [429.676 136.503 464.869 148.456] /A << /S /GoTo /D <6571756174696F6E2E372E332E322E3131> >> >>}
\MusicPDFWriteObject{3350}{<< /Type /Annot /Subtype /Link /Border [0 0 0] /H /I /C [1 0 0] /Rect [165.977 298.236 196.19 310.189] /A << /S /GoTo /D <6571756174696F6E2E372E332E322E35> >> >>}
\MusicPDFWriteObject{3351}{<< /Type /Annot /Subtype /Link /Border [0 0 0] /H /I /C [1 0 0] /Rect [152.656 271.935 183.434 286.242] /A << /S /GoTo /D <74686D2E362E322E312E31> >> >>}
\MusicPDFWriteObject{3352}{<< /Type /Annot /Subtype /Link /Border [0 0 0] /H /I /C [1 0 0] /Rect [141.115 258.784 171.328 270.737] /A << /S /GoTo /D <6571756174696F6E2E372E332E322E33> >> >>}
\MusicPDFWriteObject{3353}{<< /Type /Annot /Subtype /Link /Border [0 0 0] /H /I /C [1 0 0] /Rect [386.437 569.374 416.651 581.327] /A << /S /GoTo /D <6571756174696F6E2E322E312E342E32> >> >>}
\MusicPDFWriteObject{3354}{<< /Type /Annot /Subtype /Link /Border [0 0 0] /H /I /C [1 0 0] /Rect [490.328 569.014 521.106 581.329] /A << /S /GoTo /D <74686D2E322E312E342E31> >> >>}
\MusicPDFWriteObject{3355}{<< /Type /Annot /Subtype /Link /Border [0 0 0] /H /I /C [1 0 0] /Rect [271.817 265.18 302.03 275.971] /A << /S /GoTo /D <74686D2E372E332E322E32> >> >>}
\MusicPDFWriteObject{3356}{<< /Type /Annot /Subtype /Link /Border [0 0 0] /H /I /C [1 0 0] /Rect [209.678 478.269 244.871 490.221] /A << /S /GoTo /D <6571756174696F6E2E372E332E322E3138> >> >>}
\MusicPDFWriteObject{3357}{<< /Type /Annot /Subtype /Link /Border [0 0 0] /H /I /C [1 0 0] /Rect [273.055 478.269 308.249 490.221] /A << /S /GoTo /D <6571756174696F6E2E372E332E322E3139> >> >>}
\MusicPDFWriteObject{3358}{<< /Type /Annot /Subtype /Link /Border [0 0 0] /H /I /C [1 0 0] /Rect [275.465 465.118 305.678 477.071] /A << /S /GoTo /D <74686D2E372E332E322E32> >> >>}
\MusicPDFWriteObject{3359}{<< /Type /Annot /Subtype /Link /Border [0 0 0] /H /I /C [1 0 0] /Rect [206.518 451.37 237.295 466.114] /A << /S /GoTo /D <74686D2E322E322E332E32> >> >>}
\MusicPDFWriteObject{3360}{<< /Type /Annot /Subtype /Link /Border [0 0 0] /H /I /C [1 0 0] /Rect [123.322 348.867 161.006 360.82] /A << /S /GoTo /D <6571756174696F6E2E412E322E332E3134> >> >>}
\MusicPDFWriteObject{3361}{<< /Type /Annot /Subtype /Link /Border [0 0 0] /H /I /C [1 0 0] /Rect [226.846 113.184 257.06 125.137] /A << /S /GoTo /D <6571756174696F6E2E352E312E312E32> >> >>}
\MusicPDFWriteObject{3362}{<< /Type /Annot /Subtype /Link /Border [0 0 0] /H /I /C [1 0 0] /Rect [355.496 113.184 385.709 125.137] /A << /S /GoTo /D <6571756174696F6E2E352E312E312E31> >> >>}
\MusicPDFWriteObject{3363}{<< /Type /Annot /Subtype /Link /Border [0 0 0] /H /I /C [1 0 0] /Rect [202.447 600.149 237.64 612.102] /A << /S /GoTo /D <6571756174696F6E2E372E332E322E3230> >> >>}
\MusicPDFWriteObject{3364}{<< /Type /Annot /Subtype /Link /Border [0 0 0] /H /I /C [1 0 0] /Rect [328.701 518.709 358.688 530.47] /A << /S /GoTo /D <74686D2E352E332E312E34> >> >>}
\MusicPDFWriteObject{3365}{<< /Type /Annot /Subtype /Link /Border [0 0 0] /H /I /C [1 0 0] /Rect [375.735 518.709 398.037 530.47] /A << /S /GoTo /D <73756273656374696F6E2E352E332E31> >> >>}
\MusicPDFWriteObject{3366}{<< /Type /Annot /Subtype /Link /Border [0 0 0] /H /I /C [1 0 0] /Rect [278.98 467.019 301.446 477.867] /A << /S /GoTo /D <73756273656374696F6E2E362E332E32> >> >>}
\MusicPDFWriteObject{3367}{<< /Type /Annot /Subtype /Link /Border [0 0 0] /H /I /C [1 0 0] /Rect [195.518 438.836 235.195 452.119] /A << /S /GoTo /D <6571756174696F6E2E31302E342E332E3237> >> >>}
\MusicPDFWriteObject{3368}{<< /Type /Annot /Subtype /Link /Border [0 0 0] /H /I /C [1 0 0] /Rect [98.469 427.567 125.406 438.415] /A << /S /GoTo /D <73756273656374696F6E2E31322E322E31> >> >>}
\MusicPDFWriteObject{3369}{<< /Type /Annot /Subtype /Link /Border [0 0 0] /H /I /C [1 0 0] /Rect [199.384 401.265 229.597 412.056] /A << /S /GoTo /D <74686D2E372E332E312E31> >> >>}
\MusicPDFWriteObject{3370}{<< /Type /Annot /Subtype /Link /Border [0 0 0] /H /I /C [1 0 0] /Rect [311.888 298.127 341.536 310.08] /A << /S /GoTo /D <74686D2E372E332E312E31> >> >>}
\MusicPDFWriteObject{3371}{<< /Type /Annot /Subtype /Link /Border [0 0 0] /H /I /C [1 0 0] /Rect [319.188 130.016 350.708 140.806] /A << /S /GoTo /D <74686D2E372E332E332E31> >> >>}
\MusicPDFWriteObject{3372}{<< /Type /Annot /Subtype /Link /Border [0 0 0] /H /I /C [1 0 0] /Rect [99.138 684.606 129.352 696.558] /A << /S /GoTo /D <6571756174696F6E2E352E332E312E39> >> >>}
\MusicPDFWriteObject{3373}{<< /Type /Annot /Subtype /Link /Border [0 0 0] /H /I /C [1 0 0] /Rect [181.218 600.927 211.431 612.88] /A << /S /GoTo /D <6571756174696F6E2E372E332E332E33> >> >>}
\MusicPDFWriteObject{3374}{<< /Type /Annot /Subtype /Link /Border [0 0 0] /H /I /C [1 0 0] /Rect [160.278 471.665 191.055 484.44] /A << /S /GoTo /D <74686D2E372E332E322E34> >> >>}
\MusicPDFWriteObject{3375}{<< /Type /Annot /Subtype /Link /Border [0 0 0] /H /I /C [1 0 0] /Rect [172.155 457.961 202.368 469.913] /A << /S /GoTo /D <6571756174696F6E2E372E332E332E31> >> >>}
\MusicPDFWriteObject{3376}{<< /Type /Annot /Subtype /Link /Border [0 0 0] /H /I /C [1 0 0] /Rect [347.358 444.81 377.571 456.763] /A << /S /GoTo /D <6571756174696F6E2E372E332E332E31> >> >>}
\MusicPDFWriteObject{3377}{<< /Type /Annot /Subtype /Link /Border [0 0 0] /H /I /C [1 0 0] /Rect [290.52 338.693 320.733 347.547] /A << /S /GoTo /D <74686D2E372E332E332E31> >> >>}
\MusicPDFWriteObject{3378}{<< /Type /Annot /Subtype /Link /Border [0 0 0] /H /I /C [1 0 0] /Rect [369.618 227.112 399.831 237.903] /A << /S /GoTo /D <74686D2E372E332E332E32> >> >>}
\MusicPDFWriteObject{3379}{<< /Type /Annot /Subtype /Link /Border [0 0 0] /H /I /C [1 0 0] /Rect [369.618 144.806 399.831 155.596] /A << /S /GoTo /D <74686D2E372E332E332E32> >> >>}
\MusicPDFWriteObject{3380}{<< /Type /Annot /Subtype /Link /Border [0 0 0] /H /I /C [1 0 0] /Rect [237.09 634.233 267.303 646.186] /A << /S /GoTo /D <74686D2E372E332E312E33> >> >>}
\MusicPDFWriteObject{3381}{<< /Type /Annot /Subtype /Link /Border [0 0 0] /H /I /C [1 0 0] /Rect [330.08 634.233 360.293 646.186] /A << /S /GoTo /D <74686D2E372E332E312E36> >> >>}
\MusicPDFWriteObject{3382}{<< /Type /Annot /Subtype /Link /Border [0 0 0] /H /I /C [1 0 0] /Rect [319.188 571.995 350.708 582.786] /A << /S /GoTo /D <74686D2E372E332E332E31> >> >>}
\MusicPDFWriteObject{3383}{<< /Type /Annot /Subtype /Link /Border [0 0 0] /H /I /C [1 0 0] /Rect [167.899 346.976 198.112 358.929] /A << /S /GoTo /D <6571756174696F6E2E372E332E332E37> >> >>}
\MusicPDFWriteObject{3384}{<< /Type /Annot /Subtype /Link /Border [0 0 0] /H /I /C [1 0 0] /Rect [263.419 346.976 293.632 358.929] /A << /S /GoTo /D <6571756174696F6E2E372E332E332E38> >> >>}
\MusicPDFWriteObject{3385}{<< /Type /Annot /Subtype /Link /Border [0 0 0] /H /I /C [1 0 0] /Rect [512.049 81.101 542.978 96.902] /A << /S /GoTo /D <74686D2E372E332E332E32> >> >>}
\MusicPDFWriteObject{3386}{<< /Type /Annot /Subtype /Link /Border [0 0 0] /H /I /C [1 0 0] /Rect [153.619 622.181 185.138 634.137] /A << /S /GoTo /D <74686D2E372E332E312E36> >> >>}
\MusicPDFWriteObject{3387}{<< /Type /Annot /Subtype /Link /Border [0 0 0] /H /I /C [1 0 0] /Rect [467.028 604.889 497.806 615.68] /A << /S /GoTo /D <74686D2E372E332E312E36> >> >>}
\MusicPDFWriteObject{3388}{<< /Type /Annot /Subtype /Link /Border [0 0 0] /H /I /C [1 0 0] /Rect [291.737 591.738 322.515 602.529] /A << /S /GoTo /D <74686D2E342E362E312E35> >> >>}
\MusicPDFWriteObject{3389}{<< /Type /Annot /Subtype /Link /Border [0 0 0] /H /I /C [1 0 0] /Rect [268.206 578.587 298.815 589.378] /A << /S /GoTo /D <74686D2E342E362E322E31> >> >>}
\MusicPDFWriteObject{3390}{<< /Type /Annot /Subtype /Link /Border [0 0 0] /H /I /C [1 0 0] /Rect [359.352 552.286 389.565 563.077] /A << /S /GoTo /D <74686D2E372E332E332E32> >> >>}
\MusicPDFWriteObject{3391}{<< /Type /Annot /Subtype /Link /Border [0 0 0] /H /I /C [1 0 0] /Rect [279.982 303.744 310.759 315.699] /A << /S /GoTo /D <74686D2E342E362E312E31> >> >>}
\MusicPDFWriteObject{3392}{<< /Type /Annot /Subtype /Link /Border [0 0 0] /H /I /C [1 0 0] /Rect [337.61 303.745 367.823 315.697] /A << /S /GoTo /D <6571756174696F6E2E342E352E332E33> >> >>}
\MusicPDFWriteObject{3393}{<< /Type /Annot /Subtype /Link /Border [0 0 0] /H /I /C [1 0 0] /Rect [97.389 684.633 127.602 696.586] /A << /S /GoTo /D <6571756174696F6E2E352E322E342E31> >> >>}
\MusicPDFWriteObject{3394}{<< /Type /Annot /Subtype /Link /Border [0 0 0] /H /I /C [1 0 0] /Rect [183.406 683.664 213.055 696.586] /A << /S /GoTo /D <74686D2E352E322E342E31> >> >>}
\MusicPDFWriteObject{3395}{<< /Type /Annot /Subtype /Link /Border [0 0 0] /H /I /C [1 0 0] /Rect [445.541 644.439 476.318 655.23] /A << /S /GoTo /D <74686D2E372E332E312E36> >> >>}
\MusicPDFWriteObject{3396}{<< /Type /Annot /Subtype /Link /Border [0 0 0] /H /I /C [1 0 0] /Rect [452.69 627.719 483.467 643.974] /A << /S /GoTo /D <74686D2E342E362E312E35> >> >>}
\MusicPDFWriteObject{3397}{<< /Type /Annot /Subtype /Link /Border [0 0 0] /H /I /C [1 0 0] /Rect [125.918 584.513 156.696 599.483] /A << /S /GoTo /D <74686D2E352E332E342E34> >> >>}
\MusicPDFWriteObject{3398}{<< /Type /Annot /Subtype /Link /Border [0 0 0] /H /I /C [1 0 0] /Rect [258.509 584.513 289.286 599.483] /A << /S /GoTo /D <74686D2E352E332E342E32> >> >>}
\MusicPDFWriteObject{3399}{<< /Type /Annot /Subtype /Link /Border [0 0 0] /H /I /C [1 0 0] /Rect [349.387 584.513 382.845 599.483] /A << /S /GoTo /D <74686D2E442E332E332E36> >> >>}
\MusicPDFWriteObject{3400}{<< /Type /Annot /Subtype /Link /Border [0 0 0] /H /I /C [1 0 0] /Rect [113.45 560.404 144.2 572.356] /A << /S /GoTo /D <74686D2E352E342E312E31> >> >>}
\MusicPDFWriteObject{3401}{<< /Type /Annot /Subtype /Link /Border [0 0 0] /H /I /C [1 0 0] /Rect [98.422 547.253 133.615 559.206] /A << /S /GoTo /D <6571756174696F6E2E372E332E312E3230> >> >>}
\MusicPDFWriteObject{3402}{<< /Type /Annot /Subtype /Link /Border [0 0 0] /H /I /C [1 0 0] /Rect [173.954 547.253 209.147 559.206] /A << /S /GoTo /D <6571756174696F6E2E372E332E332E3135> >> >>}
\MusicPDFWriteObject{3403}{<< /Type /Annot /Subtype /Link /Border [0 0 0] /H /I /C [1 0 0] /Rect [71.004 318.639 101.217 336.881] /A << /S /GoTo /D <74686D2E352E332E342E34> >> >>}
\MusicPDFWriteObject{3404}{<< /Type /Annot /Subtype /Link /Border [0 0 0] /H /I /C [1 0 0] /Rect [191.068 233.082 221.281 245.035] /A << /S /GoTo /D <6571756174696F6E2E352E332E342E31> >> >>}
\MusicPDFWriteObject{3405}{<< /Type /Annot /Subtype /Link /Border [0 0 0] /H /I /C [1 0 0] /Rect [295.445 232.722 325.235 245.037] /A << /S /GoTo /D <74686D2E352E332E342E31> >> >>}
\MusicPDFWriteObject{3406}{<< /Type /Annot /Subtype /Link /Border [0 0 0] /H /I /C [1 0 0] /Rect [413.358 232.722 443.148 245.037] /A << /S /GoTo /D <74686D2E352E332E342E34> >> >>}
\MusicPDFWriteObject{3407}{<< /Type /Annot /Subtype /Link /Border [0 0 0] /H /I /C [1 0 0] /Rect [392.365 630.225 427.558 642.178] /A << /S /GoTo /D <6571756174696F6E2E372E332E332E3235> >> >>}
\MusicPDFWriteObject{3408}{<< /Type /Annot /Subtype /Link /Border [0 0 0] /H /I /C [1 0 0] /Rect [120.694 361.358 153.536 373.311] /A << /S /GoTo /D <74686D2E442E332E332E36> >> >>}
\MusicPDFWriteObject{3409}{<< /Type /Annot /Subtype /Link /Border [0 0 0] /H /I /C [1 0 0] /Rect [301.509 253.508 331.525 264.299] /A << /S /GoTo /D <74686D2E372E332E312E36> >> >>}
\MusicPDFWriteObject{3410}{<< /Type /Annot /Subtype /Link /Border [0 0 0] /H /I /C [1 0 0] /Rect [112.256 240.357 142.131 251.148] /A << /S /GoTo /D <74686D2E342E362E322E31> >> >>}
\MusicPDFWriteObject{3411}{<< /Type /Annot /Subtype /Link /Border [0 0 0] /H /I /C [1 0 0] /Rect [206.173 240.357 236.047 251.148] /A << /S /GoTo /D <74686D2E342E362E332E31> >> >>}
\MusicPDFWriteObject{3412}{<< /Type /Annot /Subtype /Link /Border [0 0 0] /H /I /C [0 1 0] /Rect [465.19 229.143 484.295 238.011] /A << /S /GoTo /D <636974652E4E42> >> >>}
\MusicPDFWriteObject{3413}{<< /Type /Annot /Subtype /Link /Border [0 0 0] /H /I /C [1 0 0] /Rect [98.422 187.201 133.615 199.154] /A << /S /GoTo /D <6571756174696F6E2E372E332E332E3138> >> >>}
\MusicPDFWriteObject{3414}{<< /Type /Annot /Subtype /Link /Border [0 0 0] /H /I /C [1 0 0] /Rect [199.905 187.201 235.099 199.154] /A << /S /GoTo /D <6571756174696F6E2E372E332E332E3139> >> >>}
\MusicPDFWriteObject{3415}{<< /Type /Annot /Subtype /Link /Border [0 0 0] /H /I /C [1 0 0] /Rect [292.597 161.127 322.726 175.046] /A << /S /GoTo /D <74686D2E372E332E312E36> >> >>}
\MusicPDFWriteObject{3416}{<< /Type /Annot /Subtype /Link /Border [0 0 0] /H /I /C [1 0 0] /Rect [193.796 644.656 228.99 656.609] /A << /S /GoTo /D <6571756174696F6E2E372E332E332E3330> >> >>}
\MusicPDFWriteObject{3417}{<< /Type /Annot /Subtype /Link /Border [0 0 0] /H /I /C [1 0 0] /Rect [140.121 632.059 170.334 642.85] /A << /S /GoTo /D <74686D2E372E332E312E35> >> >>}
\MusicPDFWriteObject{3418}{<< /Type /Annot /Subtype /Link /Border [0 0 0] /H /I /C [1 0 0] /Rect [378.059 632.059 408.272 642.85] /A << /S /GoTo /D <74686D2E342E362E312E35> >> >>}
\MusicPDFWriteObject{3419}{<< /Type /Annot /Subtype /Link /Border [0 0 0] /H /I /C [1 0 0] /Rect [380.987 603.802 410.721 614.606] /A << /S /GoTo /D <74686D2E372E332E312E32> >> >>}
\MusicPDFWriteObject{3420}{<< /Type /Annot /Subtype /Link /Border [0 0 0] /H /I /C [1 0 0] /Rect [158.387 299.588 189.164 313.901] /A << /S /GoTo /D <74686D2E352E312E322E32> >> >>}
\MusicPDFWriteObject{3421}{<< /Type /Annot /Subtype /Link /Border [0 0 0] /H /I /C [1 0 0] /Rect [163.073 472.657 193.286 484.61] /A << /S /GoTo /D <6571756174696F6E2E372E332E342E36> >> >>}
\MusicPDFWriteObject{3422}{<< /Type /Annot /Subtype /Link /Border [0 0 0] /H /I /C [1 0 0] /Rect [498.525 349.22 529.303 361.173] /A << /S /GoTo /D <74686D2E352E342E312E31> >> >>}
\MusicPDFWriteObject{3423}{<< /Type /Annot /Subtype /Link /Border [0 0 0] /H /I /C [1 0 0] /Rect [197.775 322.918 227.988 334.871] /A << /S /GoTo /D <6571756174696F6E2E372E332E342E34> >> >>}
\MusicPDFWriteObject{3424}{<< /Type /Annot /Subtype /Link /Border [0 0 0] /H /I /C [0 1 0] /Rect [73.826 259.606 91.298 268.474] /A << /S /GoTo /D <636974652E4346> >> >>}
\MusicPDFWriteObject{3425}{<< /Type /Annot /Subtype /Link /Border [0 0 0] /H /I /C [0 1 0] /Rect [227.945 259.606 245.915 268.474] /A << /S /GoTo /D <636974652E5256> >> >>}
\MusicPDFWriteObject{3426}{<< /Type /Annot /Subtype /Link /Border [0 0 0] /H /I /C [0 1 0] /Rect [248.73 246.456 266.203 255.324] /A << /S /GoTo /D <636974652E4346> >> >>}
\MusicPDFWriteObject{3427}{<< /Type /Annot /Subtype /Link /Border [0 0 0] /H /I /C [1 0 0] /Rect [233.576 82.596 263.79 94.549] /A << /S /GoTo /D <74686D2E352E332E332E34> >> >>}
\MusicPDFWriteObject{3428}{<< /Type /Annot /Subtype /Link /Border [0 0 0] /H /I /C [1 0 0] /Rect [469.737 684.633 499.95 696.586] /A << /S /GoTo /D <6571756174696F6E2E372E332E342E38> >> >>}
\MusicPDFWriteObject{3429}{<< /Type /Annot /Subtype /Link /Border [0 0 0] /H /I /C [1 0 0] /Rect [443.663 644.968 473.877 656.921] /A << /S /GoTo /D <6571756174696F6E2E372E332E342E38> >> >>}
\MusicPDFWriteObject{3430}{<< /Type /Annot /Subtype /Link /Border [0 0 0] /H /I /C [0 1 0] /Rect [107.335 621.157 125.534 630.024] /A << /S /GoTo /D <636974652E4A6172> >> >>}
\MusicPDFWriteObject{3431}{<< /Type /Annot /Subtype /Link /Border [0 0 0] /H /I /C [1 0 0] /Rect [112.727 450.863 147.92 461.654] /A << /S /GoTo /D <74686D2E31302E332E312E31> >> >>}
\MusicPDFWriteObject{3432}{<< /Type /Annot /Subtype /Link /Border [0 0 0] /H /I /C [1 0 0] /Rect [229.791 344.173 261.9 358.175] /A << /S /GoTo /D <74686D2E372E312E342E32> >> >>}
\MusicPDFWriteObject{3433}{<< /Type /Annot /Subtype /Link /Border [0 0 0] /H /I /C [1 0 0] /Rect [146.298 204.274 176.314 216.589] /A << /S /GoTo /D <74686D2E372E312E342E32> >> >>}
\MusicPDFWriteObject{3434}{<< /Type /Annot /Subtype /Link /Border [0 0 0] /H /I /C [0 1 0] /Rect [331.671 193.974 362.5 202.841] /A << /S /GoTo /D <636974652E47726F2D61> >> >>}
\MusicPDFWriteObject{3435}{<< /Type /Annot /Subtype /Link /Border [0 0 0] /H /I /C [1 0 0] /Rect [166.51 151.851 196.271 163.804] /A << /S /GoTo /D <74686D2E372E312E342E32> >> >>}
\MusicPDFWriteObject{3436}{<< /Type /Annot /Subtype /Link /Border [0 0 0] /H /I /C [0 1 0] /Rect [380.66 154.341 411.489 163.209] /A << /S /GoTo /D <636974652E47726F2D61> >> >>}
\MusicPDFWriteObject{3437}{<< /Type /Annot /Subtype /Link /Border [0 0 0] /H /I /C [0 1 0] /Rect [134.251 141.191 165.08 150.058] /A << /S /GoTo /D <636974652E47726F2D61> >> >>}
\MusicPDFWriteObject{3438}{<< /Type /Annot /Subtype /Link /Border [0 0 0] /H /I /C [1 0 0] /Rect [235.421 669.025 265.634 680.978] /A << /S /GoTo /D <6571756174696F6E2E372E312E342E36> >> >>}
\MusicPDFWriteObject{3439}{<< /Type /Annot /Subtype /Link /Border [0 0 0] /H /I /C [1 0 0] /Rect [189.946 543.428 225.14 555.381] /A << /S /GoTo /D <74686D2E31302E332E312E31> >> >>}
\MusicPDFWriteObject{3440}{<< /Type /Annot /Subtype /Link /Border [0 0 0] /H /I /C [1 0 0] /Rect [329.703 490.825 359.916 502.778] /A << /S /GoTo /D <6571756174696F6E2E372E342E312E32> >> >>}
\MusicPDFWriteObject{3441}{<< /Type /Annot /Subtype /Link /Border [0 0 0] /H /I /C [0 1 0] /Rect [420 427.291 449.342 436.159] /A << /S /GoTo /D <636974652E4C69742D4949> >> >>}
\MusicPDFWriteObject{3442}{<< /Type /Annot /Subtype /Link /Border [0 0 0] /H /I /C [1 0 0] /Rect [209.976 411.65 240.189 423.603] /A << /S /GoTo /D <6571756174696F6E2E372E342E312E31> >> >>}
\MusicPDFWriteObject{3443}{<< /Type /Annot /Subtype /Link /Border [0 0 0] /H /I /C [1 0 0] /Rect [290.893 184.299 321.67 195.089] /A << /S /GoTo /D <74686D2E372E342E312E31> >> >>}
\MusicPDFWriteObject{3444}{<< /Type /Annot /Subtype /Link /Border [0 0 0] /H /I /C [1 0 0] /Rect [421.304 184.299 452.081 195.089] /A << /S /GoTo /D <74686D2E372E312E342E33> >> >>}
\MusicPDFWriteObject{3445}{<< /Type /Annot /Subtype /Link /Border [0 0 0] /H /I /C [1 0 0] /Rect [70.755 170.234 100.403 182.549] /A << /S /GoTo /D <74686D2E372E312E342E32> >> >>}
\MusicPDFWriteObject{3446}{<< /Type /Annot /Subtype /Link /Border [0 0 0] /H /I /C [1 0 0] /Rect [172.67 170.234 202.319 182.549] /A << /S /GoTo /D <74686D2E372E312E342E33> >> >>}
\MusicPDFWriteObject{3447}{<< /Type /Annot /Subtype /Link /Border [0 0 0] /H /I /C [1 0 0] /Rect [291.393 170.234 321.042 182.549] /A << /S /GoTo /D <74686D2E372E312E342E32> >> >>}
\MusicPDFWriteObject{3448}{<< /Type /Annot /Subtype /Link /Border [0 0 0] /H /I /C [1 0 0] /Rect [271.469 157.443 302.246 169.398] /A << /S /GoTo /D <74686D2E372E312E342E33> >> >>}
\MusicPDFWriteObject{3449}{<< /Type /Annot /Subtype /Link /Border [0 0 0] /H /I /C [0 1 0] /Rect [365.138 146.783 394.48 155.651] /A << /S /GoTo /D <636974652E4C69742D4949> >> >>}
\MusicPDFWriteObject{3450}{<< /Type /Annot /Subtype /Link /Border [0 0 0] /H /I /C [1 0 0] /Rect [154.033 131.142 184.246 143.095] /A << /S /GoTo /D <6571756174696F6E2E372E342E322E31> >> >>}
\MusicPDFWriteObject{3451}{<< /Type /Annot /Subtype /Link /Border [0 0 0] /H /I /C [1 0 0] /Rect [353.569 131.142 393.743 143.095] /A << /S /GoTo /D <6571756174696F6E2E31312E322E322E3132> >> >>}
\MusicPDFWriteObject{3452}{<< /Type /Annot /Subtype /Link /Border [0 0 0] /H /I /C [1 0 0] /Rect [452.369 131.142 487.562 143.095] /A << /S /GoTo /D <74686D2E31312E322E322E31> >> >>}
\MusicPDFWriteObject{3453}{<< /Type /Annot /Subtype /Link /Border [0 0 0] /H /I /C [1 0 0] /Rect [514.068 109.451 544.845 120.255] /A << /S /GoTo /D <74686D2E372E312E342E33> >> >>}
\MusicPDFWriteObject{3454}{<< /Type /Annot /Subtype /Link /Border [0 0 0] /H /I /C [1 0 0] /Rect [236.233 85.086 266.446 93.94] /A << /S /GoTo /D <74686D2E372E312E342E33> >> >>}
\MusicPDFWriteObject{3455}{<< /Type /Annot /Subtype /Link /Border [0 0 0] /H /I /C [1 0 0] /Rect [162.11 685.187 191.956 695.977] /A << /S /GoTo /D <74686D2E372E332E342E31> >> >>}
\MusicPDFWriteObject{3456}{<< /Type /Annot /Subtype /Link /Border [0 0 0] /H /I /C [1 0 0] /Rect [344.121 672.036 374.475 682.827] /A << /S /GoTo /D <74686D2E372E332E342E31> >> >>}
\MusicPDFWriteObject{3457}{<< /Type /Annot /Subtype /Link /Border [0 0 0] /H /I /C [1 0 0] /Rect [294.142 489.616 329.335 501.569] /A << /S /GoTo /D <6571756174696F6E2E352E332E342E3133> >> >>}
\MusicPDFWriteObject{3458}{<< /Type /Annot /Subtype /Link /Border [0 0 0] /H /I /C [1 0 0] /Rect [432.121 416.414 467.315 428.367] /A << /S /GoTo /D <6571756174696F6E2E352E332E342E3134> >> >>}
\MusicPDFWriteObject{3459}{<< /Type /Annot /Subtype /Link /Border [0 0 0] /H /I /C [1 0 0] /Rect [117.345 204.385 147.558 216.337] /A << /S /GoTo /D <6571756174696F6E2E372E342E332E34> >> >>}
\MusicPDFWriteObject{3460}{<< /Type /Annot /Subtype /Link /Border [0 0 0] /H /I /C [1 0 0] /Rect [203.785 150.94 233.998 162.893] /A << /S /GoTo /D <6571756174696F6E2E352E312E312E31> >> >>}
\MusicPDFWriteObject{3461}{<< /Type /Annot /Subtype /Link /Border [0 0 0] /H /I /C [1 0 0] /Rect [189.392 522.077 219.605 534.03] /A << /S /GoTo /D <74686D2E352E332E332E33> >> >>}
\MusicPDFWriteObject{3462}{<< /Type /Annot /Subtype /Link /Border [0 0 0] /H /I /C [1 0 0] /Rect [144.996 241.164 174.644 255.166] /A << /S /GoTo /D <74686D2E362E322E322E31> >> >>}
\MusicPDFWriteObject{3463}{<< /Type /Annot /Subtype /Link /Border [0 0 0] /H /I /C [1 0 0] /Rect [320.75 117.836 350.963 130.763] /A << /S /GoTo /D <74686D2E362E322E322E31> >> >>}
\MusicPDFWriteObject{3464}{<< /Type /Annot /Subtype /Link /Border [0 0 0] /H /I /C [0 1 0] /Rect [312.151 344.913 337.212 353.781] /A << /S /GoTo /D <636974652E53636861> >> >>}
\MusicPDFWriteObject{3465}{<< /Type /Annot /Subtype /Link /Border [0 0 0] /H /I /C [1 0 0] /Rect [183.332 266.542 213.037 278.498] /A << /S /GoTo /D <74686D2E372E312E342E32> >> >>}
\MusicPDFWriteObject{3466}{<< /Type /Annot /Subtype /Link /Border [0 0 0] /H /I /C [0 1 0] /Rect [73.77 255.882 108.535 264.75] /A << /S /GoTo /D <636974652E436E772D62> >> >>}
\MusicPDFWriteObject{3467}{<< /Type /Annot /Subtype /Link /Border [0 0 0] /H /I /C [1 0 0] /Rect [272.915 194.333 303.693 206.648] /A << /S /GoTo /D <74686D2E372E312E342E32> >> >>}
\MusicPDFWriteObject{3468}{<< /Type /Annot /Subtype /Link /Border [0 0 0] /H /I /C [0 1 0] /Rect [87.604 605.025 106.709 613.893] /A << /S /GoTo /D <636974652E4E42> >> >>}
\MusicPDFWriteObject{3469}{<< /Type /Annot /Subtype /Link /Border [0 0 0] /H /I /C [0 1 0] /Rect [360.956 438.315 380.061 447.183] /A << /S /GoTo /D <636974652E4E42> >> >>}
\MusicPDFWriteObject{3470}{<< /Type /Annot /Subtype /Link /Border [0 0 0] /H /I /C [1 0 0] /Rect [291.928 422.673 325.386 434.629] /A << /S /GoTo /D <74686D2E442E332E312E32> >> >>}
\MusicPDFWriteObject{3471}{<< /Type /Annot /Subtype /Link /Border [0 0 0] /H /I /C [1 0 0] /Rect [437.457 422.673 470.915 434.629] /A << /S /GoTo /D <74686D2E442E332E322E35> >> >>}
\MusicPDFWriteObject{3472}{<< /Type /Annot /Subtype /Link /Border [0 0 0] /H /I /C [0 1 0] /Rect [499.275 412.013 517.473 420.881] /A << /S /GoTo /D <636974652E4A6172> >> >>}
\MusicPDFWriteObject{3473}{<< /Type /Annot /Subtype /Link /Border [0 0 0] /H /I /C [1 0 0] /Rect [425.553 396.372 460.746 408.325] /A << /S /GoTo /D <6571756174696F6E2E372E342E332E3131> >> >>}
\MusicPDFWriteObject{3474}{<< /Type /Annot /Subtype /Link /Border [0 0 0] /H /I /C [1 0 0] /Rect [308.324 383.222 338.537 395.174] /A << /S /GoTo /D <6571756174696F6E2E372E342E332E34> >> >>}
\MusicPDFWriteObject{3475}{<< /Type /Annot /Subtype /Link /Border [0 0 0] /H /I /C [1 0 0] /Rect [239.221 343.769 269.434 355.722] /A << /S /GoTo /D <6571756174696F6E2E372E342E332E35> >> >>}
\MusicPDFWriteObject{3476}{<< /Type /Annot /Subtype /Link /Border [0 0 0] /H /I /C [1 0 0] /Rect [327.193 330.619 357.406 342.571] /A << /S /GoTo /D <6571756174696F6E2E332E322E322E32> >> >>}
\MusicPDFWriteObject{3477}{<< /Type /Annot /Subtype /Link /Border [0 0 0] /H /I /C [1 0 0] /Rect [416.405 330.618 446.844 342.591] /A << /S /GoTo /D <74686D2E332E322E322E31> >> >>}
\MusicPDFWriteObject{3478}{<< /Type /Annot /Subtype /Link /Border [0 0 0] /H /I /C [1 0 0] /Rect [329.485 290.646 359.698 302.961] /A << /S /GoTo /D <74686D2E362E322E312E31> >> >>}
\MusicPDFWriteObject{3479}{<< /Type /Annot /Subtype /Link /Border [0 0 0] /H /I /C [1 0 0] /Rect [387.399 205.969 417.048 221.671] /A << /S /GoTo /D <74686D2E352E332E312E35> >> >>}
\MusicPDFWriteObject{3480}{<< /Type /Annot /Subtype /Link /Border [0 0 0] /H /I /C [1 0 0] /Rect [402.304 147.961 437.498 159.914] /A << /S /GoTo /D <6571756174696F6E2E372E342E332E3132> >> >>}
\MusicPDFWriteObject{3481}{<< /Type /Annot /Subtype /Link /Border [0 0 0] /H /I /C [1 0 0] /Rect [126.919 98.383 157.132 110.335] /A << /S /GoTo /D <74686D2E372E332E342E31> >> >>}
\MusicPDFWriteObject{3482}{<< /Type /Annot /Subtype /Link /Border [0 0 0] /H /I /C [1 0 0] /Rect [314.337 684.633 349.53 696.586] /A << /S /GoTo /D <6571756174696F6E2E382E312E302E3130> >> >>}
\MusicPDFWriteObject{3483}{<< /Type /Annot /Subtype /Link /Border [0 0 0] /H /I /C [1 0 0] /Rect [476.8 684.633 507.521 696.586] /A << /S /GoTo /D <74686D2E372E352E302E31> >> >>}
\MusicPDFWriteObject{3484}{<< /Type /Annot /Subtype /Link /Border [0 0 0] /H /I /C [1 0 0] /Rect [512.948 672.036 543.161 682.827] /A << /S /GoTo /D <74686D2E372E322E312E34> >> >>}
\MusicPDFWriteObject{3485}{<< /Type /Annot /Subtype /Link /Border [0 0 0] /H /I /C [1 0 0] /Rect [296.237 589.626 325.885 604.182] /A << /S /GoTo /D <74686D2E372E352E302E31> >> >>}
\MusicPDFWriteObject{3486}{<< /Type /Annot /Subtype /Link /Border [0 0 0] /H /I /C [1 0 0] /Rect [349.798 589.986 384.992 601.939] /A << /S /GoTo /D <6571756174696F6E2E382E312E302E3130> >> >>}
\MusicPDFWriteObject{3487}{<< /Type /Annot /Subtype /Link /Border [0 0 0] /H /I /C [1 0 0] /Rect [73.712 563.685 103.925 575.638] /A << /S /GoTo /D <6571756174696F6E2E372E342E332E34> >> >>}
\MusicPDFWriteObject{3488}{<< /Type /Annot /Subtype /Link /Border [0 0 0] /H /I /C [1 0 0] /Rect [290.589 563.685 320.746 575.638] /A << /S /GoTo /D <74686D2E372E332E342E31> >> >>}
\MusicPDFWriteObject{3489}{<< /Type /Annot /Subtype /Link /Border [0 0 0] /H /I /C [0 1 0] /Rect [73.793 539.873 91.991 548.741] /A << /S /GoTo /D <636974652E4A6172> >> >>}
\MusicPDFWriteObject{3490}{<< /Type /Annot /Subtype /Link /Border [0 0 0] /H /I /C [1 0 0] /Rect [426.19 497.931 461.383 509.884] /A << /S /GoTo /D <6571756174696F6E2E372E342E332E3131> >> >>}
\MusicPDFWriteObject{3491}{<< /Type /Annot /Subtype /Link /Border [0 0 0] /H /I /C [1 0 0] /Rect [236.093 484.78 266.306 496.733] /A << /S /GoTo /D <6571756174696F6E2E372E342E332E34> >> >>}
\MusicPDFWriteObject{3492}{<< /Type /Annot /Subtype /Link /Border [0 0 0] /H /I /C [0 1 0] /Rect [367.781 487.27 386.886 496.138] /A << /S /GoTo /D <636974652E4E42> >> >>}
\MusicPDFWriteObject{3493}{<< /Type /Annot /Subtype /Link /Border [0 0 0] /H /I /C [1 0 0] /Rect [278.965 458.479 309.178 470.431] /A << /S /GoTo /D <6571756174696F6E2E372E342E332E35> >> >>}
\MusicPDFWriteObject{3494}{<< /Type /Annot /Subtype /Link /Border [0 0 0] /H /I /C [0 1 0] /Rect [106.226 421.516 124.425 430.384] /A << /S /GoTo /D <636974652E4A6172> >> >>}
\MusicPDFWriteObject{3495}{<< /Type /Annot /Subtype /Link /Border [0 0 0] /H /I /C [1 0 0] /Rect [96.911 405.876 127.124 417.828] /A << /S /GoTo /D <6571756174696F6E2E372E342E332E35> >> >>}
\MusicPDFWriteObject{3496}{<< /Type /Annot /Subtype /Link /Border [0 0 0] /H /I /C [1 0 0] /Rect [382.933 403.828 412.582 417.83] /A << /S /GoTo /D <74686D2E372E312E342E33> >> >>}
\MusicPDFWriteObject{3497}{<< /Type /Annot /Subtype /Link /Border [0 0 0] /H /I /C [1 0 0] /Rect [86.497 311.261 116.71 323.213] /A << /S /GoTo /D <6571756174696F6E2E372E342E332E31> >> >>}
\MusicPDFWriteObject{3498}{<< /Type /Annot /Subtype /Link /Border [0 0 0] /H /I /C [1 0 0] /Rect [186.127 311.261 216.34 323.213] /A << /S /GoTo /D <74686D2E372E342E332E31> >> >>}
\MusicPDFWriteObject{3499}{<< /Type /Annot /Subtype /Link /Border [0 0 0] /H /I /C [1 0 0] /Rect [196.52 274.833 226.225 286.29] /A << /S /GoTo /D <74686D2E372E342E332E31> >> >>}
\MusicPDFWriteObject{3500}{<< /Type /Annot /Subtype /Link /Border [0 0 0] /H /I /C [1 0 0] /Rect [320.622 209.633 350.835 220.423] /A << /S /GoTo /D <74686D2E372E342E322E31> >> >>}
\MusicPDFWriteObject{3501}{<< /Type /Annot /Subtype /Link /Border [0 0 0] /H /I /C [1 0 0] /Rect [138.403 122.048 168.616 134.001] /A << /S /GoTo /D <6571756174696F6E2E372E342E332E31> >> >>}
\MusicPDFWriteObject{3502}{<< /Type /Annot /Subtype /Link /Border [0 0 0] /H /I /C [1 0 0] /Rect [128.506 684.633 160.025 696.586] /A << /S /GoTo /D <74686D2E372E342E332E31> >> >>}
\MusicPDFWriteObject{3503}{<< /Type /Annot /Subtype /Link /Border [0 0 0] /H /I /C [1 0 0] /Rect [126.919 537.461 157.132 549.413] /A << /S /GoTo /D <74686D2E372E342E332E31> >> >>}
\MusicPDFWriteObject{3504}{<< /Type /Annot /Subtype /Link /Border [0 0 0] /H /I /C [1 0 0] /Rect [346.201 518.045 381.394 529.998] /A << /S /GoTo /D <6571756174696F6E2E372E342E332E3134> >> >>}
\MusicPDFWriteObject{3505}{<< /Type /Annot /Subtype /Link /Border [0 0 0] /H /I /C [1 0 0] /Rect [428.856 504.894 464.05 516.847] /A << /S /GoTo /D <6571756174696F6E2E31302E332E312E38> >> >>}
\MusicPDFWriteObject{3506}{<< /Type /Annot /Subtype /Link /Border [0 0 0] /H /I /C [1 0 0] /Rect [73.712 491.743 108.905 503.696] /A << /S /GoTo /D <6571756174696F6E2E31302E332E312E33> >> >>}
\MusicPDFWriteObject{3507}{<< /Type /Annot /Subtype /Link /Border [0 0 0] /H /I /C [1 0 0] /Rect [167.531 491.743 202.725 503.696] /A << /S /GoTo /D <74686D2E31302E332E312E31> >> >>}
\MusicPDFWriteObject{3508}{<< /Type /Annot /Subtype /Link /Border [0 0 0] /H /I /C [1 0 0] /Rect [364.698 465.442 396.986 477.395] /A << /S /GoTo /D <74686D2E422E322E312E39> >> >>}
\MusicPDFWriteObject{3509}{<< /Type /Annot /Subtype /Link /Border [0 0 0] /H /I /C [1 0 0] /Rect [363.638 419.721 387.729 431.676] /A << /S /GoTo /D <73756273656374696F6E2E422E322E31> >> >>}
\MusicPDFWriteObject{3510}{<< /Type /Annot /Subtype /Link /Border [0 0 0] /H /I /C [1 0 0] /Rect [174.785 405.076 207.073 418.526] /A << /S /GoTo /D <74686D2E422E332E322E31> >> >>}
\MusicPDFWriteObject{3511}{<< /Type /Annot /Subtype /Link /Border [0 0 0] /H /I /C [1 0 0] /Rect [289.368 338.157 324.562 350.109] /A << /S /GoTo /D <6571756174696F6E2E372E342E332E3135> >> >>}
\MusicPDFWriteObject{3512}{<< /Type /Annot /Subtype /Link /Border [0 0 0] /H /I /C [1 0 0] /Rect [325.07 311.855 360.263 323.808] /A << /S /GoTo /D <6571756174696F6E2E372E342E332E3135> >> >>}
\MusicPDFWriteObject{3513}{<< /Type /Annot /Subtype /Link /Border [0 0 0] /H /I /C [1 0 0] /Rect [463.625 311.855 498.818 323.808] /A << /S /GoTo /D <6571756174696F6E2E372E342E332E3134> >> >>}
\MusicPDFWriteObject{3514}{<< /Type /Annot /Subtype /Link /Border [0 0 0] /H /I /C [0 1 0] /Rect [388.59 288.044 410.98 296.911] /A << /S /GoTo /D <636974652E44656E> >> >>}
\MusicPDFWriteObject{3515}{<< /Type /Annot /Subtype /Link /Border [0 0 0] /H /I /C [1 0 0] /Rect [159.843 211.616 195.036 223.569] /A << /S /GoTo /D <6571756174696F6E2E372E342E332E3136> >> >>}
\MusicPDFWriteObject{3516}{<< /Type /Annot /Subtype /Link /Border [0 0 0] /H /I /C [1 0 0] /Rect [108.58 198.465 145.848 210.418] /A << /S /GoTo /D <6571756174696F6E2E422E342E312E3232> >> >>}
\MusicPDFWriteObject{3517}{<< /Type /Annot /Subtype /Link /Border [0 0 0] /H /I /C [1 0 0] /Rect [210.147 198.465 242.435 210.418] /A << /S /GoTo /D <74686D2E422E342E312E39> >> >>}
\MusicPDFWriteObject{3518}{<< /Type /Annot /Subtype /Link /Border [0 0 0] /H /I /C [0 1 0] /Rect [206.785 150.221 229.175 159.088] /A << /S /GoTo /D <636974652E44656E> >> >>}
\MusicPDFWriteObject{3519}{<< /Type /Annot /Subtype /Link /Border [0 0 0] /H /I /C [1 0 0] /Rect [181.501 684.633 214.395 696.586] /A << /S /GoTo /D <74686D2E422E322E322E34> >> >>}
\MusicPDFWriteObject{3520}{<< /Type /Annot /Subtype /Link /Border [0 0 0] /H /I /C [0 1 0] /Rect [259.658 687.123 281.515 695.991] /A << /S /GoTo /D <636974652E447572> >> >>}
\MusicPDFWriteObject{3521}{<< /Type /Annot /Subtype /Link /Border [0 0 0] /H /I /C [1 0 0] /Rect [449.104 684.633 484.297 696.586] /A << /S /GoTo /D <6571756174696F6E2E372E342E332E3137> >> >>}
\MusicPDFWriteObject{3522}{<< /Type /Annot /Subtype /Link /Border [0 0 0] /H /I /C [1 0 0] /Rect [201.79 671.482 236.984 683.435] /A << /S /GoTo /D <6571756174696F6E2E372E342E332E3135> >> >>}
\MusicPDFWriteObject{3523}{<< /Type /Annot /Subtype /Link /Border [0 0 0] /H /I /C [1 0 0] /Rect [73.712 658.332 108.905 670.284] /A << /S /GoTo /D <6571756174696F6E2E372E342E332E3134> >> >>}
\MusicPDFWriteObject{3524}{<< /Type /Annot /Subtype /Link /Border [0 0 0] /H /I /C [1 0 0] /Rect [468.32 303.79 500.002 315.743] /A << /S /GoTo /D <74686D2E422E362E312E36> >> >>}
\MusicPDFWriteObject{3525}{<< /Type /Annot /Subtype /Link /Border [0 0 0] /H /I /C [1 0 0] /Rect [485.296 291.193 515.65 301.983] /A << /S /GoTo /D <74686D2E372E332E312E36> >> >>}
\MusicPDFWriteObject{3526}{<< /Type /Annot /Subtype /Link /Border [0 0 0] /H /I /C [1 0 0] /Rect [71.004 82.595 104.462 94.55] /A << /S /GoTo /D <74686D2E442E322E312E31> >> >>}
\MusicPDFWriteObject{3527}{<< /Type /Annot /Subtype /Link /Border [0 0 0] /H /I /C [1 0 0] /Rect [331.135 82.596 361.348 94.549] /A << /S /GoTo /D <6571756174696F6E2E342E352E312E33> >> >>}
\MusicPDFWriteObject{3528}{<< /Type /Annot /Subtype /Link /Border [0 0 0] /H /I /C [1 0 0] /Rect [445.401 82.595 476.179 94.55] /A << /S /GoTo /D <74686D2E342E352E312E31> >> >>}
\MusicPDFWriteObject{3529}{<< /Type /Annot /Subtype /Link /Border [0 0 0] /H /I /C [1 0 0] /Rect [255.713 617.849 285.475 632.598] /A << /S /GoTo /D <74686D2E342E362E312E31> >> >>}
\MusicPDFWriteObject{3530}{<< /Type /Annot /Subtype /Link /Border [0 0 0] /H /I /C [1 0 0] /Rect [148.776 507.096 178.989 519.049] /A << /S /GoTo /D <6571756174696F6E2E372E352E302E32> >> >>}
\MusicPDFWriteObject{3531}{<< /Type /Annot /Subtype /Link /Border [0 0 0] /H /I /C [1 0 0] /Rect [108.585 320.093 138.798 332.045] /A << /S /GoTo /D <6571756174696F6E2E342E362E312E37> >> >>}
\MusicPDFWriteObject{3532}{<< /Type /Annot /Subtype /Link /Border [0 0 0] /H /I /C [1 0 0] /Rect [203.109 320.093 233.322 332.045] /A << /S /GoTo /D <74686D2E342E362E312E33> >> >>}
\MusicPDFWriteObject{3533}{<< /Type /Annot /Subtype /Link /Border [0 0 0] /H /I /C [1 0 0] /Rect [380.013 320.093 410.226 332.045] /A << /S /GoTo /D <6571756174696F6E2E382E312E302E35> >> >>}
\MusicPDFWriteObject{3534}{<< /Type /Annot /Subtype /Link /Border [0 0 0] /H /I /C [1 0 0] /Rect [285.509 281.194 317.433 291.984] /A << /S /GoTo /D <74686D2E422E332E312E31> >> >>}
\MusicPDFWriteObject{3535}{<< /Type /Annot /Subtype /Link /Border [0 0 0] /H /I /C [1 0 0] /Rect [131.896 267.49 164.184 279.442] /A << /S /GoTo /D <74686D2E422E332E322E31> >> >>}
\MusicPDFWriteObject{3536}{<< /Type /Annot /Subtype /Link /Border [0 0 0] /H /I /C [1 0 0] /Rect [239.691 254.339 269.904 266.292] /A << /S /GoTo /D <6571756174696F6E2E382E312E302E35> >> >>}
\MusicPDFWriteObject{3537}{<< /Type /Annot /Subtype /Link /Border [0 0 0] /H /I /C [0 1 0] /Rect [87.923 243.678 107.028 252.546] /A << /S /GoTo /D <636974652E4E42> >> >>}
\MusicPDFWriteObject{3538}{<< /Type /Annot /Subtype /Link /Border [0 0 0] /H /I /C [1 0 0] /Rect [73.712 228.037 103.925 239.99] /A << /S /GoTo /D <6571756174696F6E2E382E312E302E35> >> >>}
\MusicPDFWriteObject{3539}{<< /Type /Annot /Subtype /Link /Border [0 0 0] /H /I /C [1 0 0] /Rect [383.703 214.887 413.917 226.839] /A << /S /GoTo /D <6571756174696F6E2E382E312E302E34> >> >>}
\MusicPDFWriteObject{3540}{<< /Type /Annot /Subtype /Link /Border [0 0 0] /H /I /C [1 0 0] /Rect [106.244 188.585 136.175 200.538] /A << /S /GoTo /D <74686D2E342E362E312E33> >> >>}
\MusicPDFWriteObject{3541}{<< /Type /Annot /Subtype /Link /Border [0 0 0] /H /I /C [1 0 0] /Rect [299.817 188.585 330.03 200.538] /A << /S /GoTo /D <6571756174696F6E2E382E312E302E34> >> >>}
\MusicPDFWriteObject{3542}{<< /Type /Annot /Subtype /Link /Border [0 0 0] /H /I /C [1 0 0] /Rect [240.281 175.434 270.494 187.387] /A << /S /GoTo /D <6571756174696F6E2E372E352E302E31> >> >>}
\MusicPDFWriteObject{3543}{<< /Type /Annot /Subtype /Link /Border [0 0 0] /H /I /C [1 0 0] /Rect [159.19 149.686 189.967 160.477] /A << /S /GoTo /D <74686D2E372E312E342E33> >> >>}
\MusicPDFWriteObject{3544}{<< /Type /Annot /Subtype /Link /Border [0 0 0] /H /I /C [1 0 0] /Rect [297.439 149.686 328.216 160.477] /A << /S /GoTo /D <74686D2E342E342E322E36> >> >>}
\MusicPDFWriteObject{3545}{<< /Type /Annot /Subtype /Link /Border [0 0 0] /H /I /C [1 0 0] /Rect [73.712 135.982 103.925 147.935] /A << /S /GoTo /D <6571756174696F6E2E382E312E302E34> >> >>}
\MusicPDFWriteObject{3546}{<< /Type /Annot /Subtype /Link /Border [0 0 0] /H /I /C [0 1 0] /Rect [260.112 138.472 279.217 147.34] /A << /S /GoTo /D <636974652E4E42> >> >>}
\MusicPDFWriteObject{3547}{<< /Type /Annot /Subtype /Link /Border [0 0 0] /H /I /C [1 0 0] /Rect [334.605 122.831 365.326 134.784] /A << /S /GoTo /D <74686D2E362E322E312E32> >> >>}
\MusicPDFWriteObject{3548}{<< /Type /Annot /Subtype /Link /Border [0 0 0] /H /I /C [1 0 0] /Rect [478.562 122.831 511.958 134.784] /A << /S /GoTo /D <74686D2E442E332E332E37> >> >>}
\MusicPDFWriteObject{3549}{<< /Type /Annot /Subtype /Link /Border [0 0 0] /H /I /C [0 1 0] /Rect [123.098 112.171 148.158 121.038] /A << /S /GoTo /D <636974652E53636861> >> >>}
\MusicPDFWriteObject{3550}{<< /Type /Annot /Subtype /Link /Border [0 0 0] /H /I /C [0 1 0] /Rect [403.967 112.171 422.94 121.038] /A << /S /GoTo /D <636974652E547265> >> >>}
\MusicPDFWriteObject{3551}{<< /Type /Annot /Subtype /Link /Border [0 0 0] /H /I /C [1 0 0] /Rect [515.197 94.482 544.845 108.484] /A << /S /GoTo /D <74686D2E342E342E322E36> >> >>}
\MusicPDFWriteObject{3552}{<< /Type /Annot /Subtype /Link /Border [0 0 0] /H /I /C [1 0 0] /Rect [219.026 671.482 249.239 683.435] /A << /S /GoTo /D <6571756174696F6E2E382E312E302E34> >> >>}
\MusicPDFWriteObject{3553}{<< /Type /Annot /Subtype /Link /Border [0 0 0] /H /I /C [0 1 0] /Rect [73.77 647.671 91.969 656.539] /A << /S /GoTo /D <636974652E4A6172> >> >>}
\MusicPDFWriteObject{3554}{<< /Type /Annot /Subtype /Link /Border [0 0 0] /H /I /C [1 0 0] /Rect [390.772 632.03 420.985 643.983] /A << /S /GoTo /D <6571756174696F6E2E382E312E302E34> >> >>}
\MusicPDFWriteObject{3555}{<< /Type /Annot /Subtype /Link /Border [0 0 0] /H /I /C [0 1 0] /Rect [193.806 621.369 212.91 630.237] /A << /S /GoTo /D <636974652E4E42> >> >>}
\MusicPDFWriteObject{3556}{<< /Type /Annot /Subtype /Link /Border [0 0 0] /H /I /C [1 0 0] /Rect [221.32 605.729 251.533 617.681] /A << /S /GoTo /D <6571756174696F6E2E382E312E302E34> >> >>}
\MusicPDFWriteObject{3557}{<< /Type /Annot /Subtype /Link /Border [0 0 0] /H /I /C [1 0 0] /Rect [144.524 529.849 174.173 541.802] /A << /S /GoTo /D <74686D2E372E312E342E33> >> >>}
\MusicPDFWriteObject{3558}{<< /Type /Annot /Subtype /Link /Border [0 0 0] /H /I /C [1 0 0] /Rect [322.776 529.849 352.989 541.802] /A << /S /GoTo /D <6571756174696F6E2E382E312E302E37> >> >>}
\MusicPDFWriteObject{3559}{<< /Type /Annot /Subtype /Link /Border [0 0 0] /H /I /C [1 0 0] /Rect [109.531 489.394 139.744 501.346] /A << /S /GoTo /D <6571756174696F6E2E352E312E312E31> >> >>}
\MusicPDFWriteObject{3560}{<< /Type /Annot /Subtype /Link /Border [0 0 0] /H /I /C [1 0 0] /Rect [296.346 476.243 326.559 488.196] /A << /S /GoTo /D <6571756174696F6E2E382E312E302E37> >> >>}
\MusicPDFWriteObject{3561}{<< /Type /Annot /Subtype /Link /Border [0 0 0] /H /I /C [0 1 0] /Rect [339.402 452.431 357.601 461.299] /A << /S /GoTo /D <636974652E4A6172> >> >>}
\MusicPDFWriteObject{3562}{<< /Type /Annot /Subtype /Link /Border [0 0 0] /H /I /C [1 0 0] /Rect [380.937 436.43 411.038 448.135] /A << /S /GoTo /D <74686D2E372E352E302E31> >> >>}
\MusicPDFWriteObject{3563}{<< /Type /Annot /Subtype /Link /Border [0 0 0] /H /I /C [1 0 0] /Rect [70.755 423.64 101.109 435.593] /A << /S /GoTo /D <74686D2E372E312E342E33> >> >>}
\MusicPDFWriteObject{3564}{<< /Type /Annot /Subtype /Link /Border [0 0 0] /H /I /C [0 1 0] /Rect [124.639 426.13 143.744 434.998] /A << /S /GoTo /D <636974652E4E42> >> >>}
\MusicPDFWriteObject{3565}{<< /Type /Annot /Subtype /Link /Border [0 0 0] /H /I /C [1 0 0] /Rect [221.375 410.489 251.588 422.442] /A << /S /GoTo /D <6571756174696F6E2E372E352E302E31> >> >>}
\MusicPDFWriteObject{3566}{<< /Type /Annot /Subtype /Link /Border [0 0 0] /H /I /C [0 1 0] /Rect [444.654 399.828 463.759 408.696] /A << /S /GoTo /D <636974652E4E42> >> >>}
\MusicPDFWriteObject{3567}{<< /Type /Annot /Subtype /Link /Border [0 0 0] /H /I /C [1 0 0] /Rect [406.272 383.827 435.92 395.532] /A << /S /GoTo /D <74686D2E372E332E312E36> >> >>}
\MusicPDFWriteObject{3568}{<< /Type /Annot /Subtype /Link /Border [0 0 0] /H /I /C [1 0 0] /Rect [330.767 371.037 360.981 382.99] /A << /S /GoTo /D <6571756174696F6E2E372E352E302E31> >> >>}
\MusicPDFWriteObject{3569}{<< /Type /Annot /Subtype /Link /Border [0 0 0] /H /I /C [1 0 0] /Rect [134.175 294.335 163.824 306.288] /A << /S /GoTo /D <74686D2E352E322E312E31> >> >>}
\MusicPDFWriteObject{3570}{<< /Type /Annot /Subtype /Link /Border [0 0 0] /H /I /C [1 0 0] /Rect [206.855 254.646 237.068 266.599] /A << /S /GoTo /D <6571756174696F6E2E372E332E312E35> >> >>}
\MusicPDFWriteObject{3571}{<< /Type /Annot /Subtype /Link /Border [0 0 0] /H /I /C [1 0 0] /Rect [306.615 251.63 336.263 266.6] /A << /S /GoTo /D <74686D2E372E332E312E33> >> >>}
\MusicPDFWriteObject{3572}{<< /Type /Annot /Subtype /Link /Border [0 0 0] /H /I /C [1 0 0] /Rect [224.855 186.148 260.048 198.1] /A << /S /GoTo /D <6571756174696F6E2E372E332E312E3137> >> >>}
\MusicPDFWriteObject{3573}{<< /Type /Annot /Subtype /Link /Border [0 0 0] /H /I /C [1 0 0] /Rect [329.971 186.148 359.761 198.1] /A << /S /GoTo /D <74686D2E372E332E312E36> >> >>}
\MusicPDFWriteObject{3574}{<< /Type /Annot /Subtype /Link /Border [0 0 0] /H /I /C [1 0 0] /Rect [191.356 107.243 223.644 119.196] /A << /S /GoTo /D <74686D2E422E312E312E34> >> >>}
\MusicPDFWriteObject{3575}{<< /Type /Annot /Subtype /Link /Border [0 0 0] /H /I /C [1 0 0] /Rect [286.421 107.243 318.709 119.196] /A << /S /GoTo /D <74686D2E422E362E312E36> >> >>}
\MusicPDFWriteObject{3576}{<< /Type /Annot /Subtype /Link /Border [0 0 0] /H /I /C [1 0 0] /Rect [247.661 684.633 277.874 696.586] /A << /S /GoTo /D <6571756174696F6E2E372E352E302E31> >> >>}
\MusicPDFWriteObject{3577}{<< /Type /Annot /Subtype /Link /Border [0 0 0] /H /I /C [1 0 0] /Rect [140.121 658.885 170.334 669.676] /A << /S /GoTo /D <74686D2E352E332E312E35> >> >>}
\MusicPDFWriteObject{3578}{<< /Type /Annot /Subtype /Link /Border [0 0 0] /H /I /C [1 0 0] /Rect [146.65 615.526 176.863 627.479] /A << /S /GoTo /D <74686D2E362E322E312E31> >> >>}
\MusicPDFWriteObject{3579}{<< /Type /Annot /Subtype /Link /Border [0 0 0] /H /I /C [1 0 0] /Rect [97.951 229.151 133.144 241.103] /A << /S /GoTo /D <6571756174696F6E2E372E332E312E3133> >> >>}
\MusicPDFWriteObject{3580}{<< /Type /Annot /Subtype /Link /Border [0 0 0] /H /I /C [1 0 0] /Rect [190.635 229.151 220.284 241.103] /A << /S /GoTo /D <74686D2E372E332E312E36> >> >>}
\MusicPDFWriteObject{3581}{<< /Type /Annot /Subtype /Link /Border [0 0 0] /H /I /C [1 0 0] /Rect [143.265 202.027 173.478 213.98] /A << /S /GoTo /D <6571756174696F6E2E382E312E302E38> >> >>}
\MusicPDFWriteObject{3582}{<< /Type /Annot /Subtype /Link /Border [0 0 0] /H /I /C [0 1 0] /Rect [314.098 559.114 331.571 567.982] /A << /S /GoTo /D <636974652E4346> >> >>}
\MusicPDFWriteObject{3583}{<< /Type /Annot /Subtype /Link /Border [0 0 0] /H /I /C [1 0 0] /Rect [437.211 542.731 466.86 557.619] /A << /S /GoTo /D <74686D2E372E312E342E33> >> >>}
\MusicPDFWriteObject{3584}{<< /Type /Annot /Subtype /Link /Border [0 0 0] /H /I /C [1 0 0] /Rect [184.404 528.839 214.617 540.792] /A << /S /GoTo /D <6571756174696F6E2E332E322E322E32> >> >>}
\MusicPDFWriteObject{3585}{<< /Type /Annot /Subtype /Link /Border [0 0 0] /H /I /C [1 0 0] /Rect [272.557 528.838 302.375 543.727] /A << /S /GoTo /D <74686D2E332E322E322E31> >> >>}
\MusicPDFWriteObject{3586}{<< /Type /Annot /Subtype /Link /Border [0 0 0] /H /I /C [1 0 0] /Rect [213.833 515.688 249.026 527.641] /A << /S /GoTo /D <6571756174696F6E2E372E312E342E3237> >> >>}
\MusicPDFWriteObject{3587}{<< /Type /Annot /Subtype /Link /Border [0 0 0] /H /I /C [1 0 0] /Rect [307.653 513.641 337.866 527.661] /A << /S /GoTo /D <74686D2E372E312E342E33> >> >>}
\MusicPDFWriteObject{3588}{<< /Type /Annot /Subtype /Link /Border [0 0 0] /H /I /C [1 0 0] /Rect [147.572 502.537 178.293 514.492] /A << /S /GoTo /D <74686D2E342E362E312E35> >> >>}
\MusicPDFWriteObject{3589}{<< /Type /Annot /Subtype /Link /Border [0 0 0] /H /I /C [1 0 0] /Rect [197.925 472.25 228.702 489.414] /A << /S /GoTo /D <74686D2E342E362E312E31> >> >>}
\MusicPDFWriteObject{3590}{<< /Type /Annot /Subtype /Link /Border [0 0 0] /H /I /C [1 0 0] /Rect [190.613 461.294 220.826 473.247] /A << /S /GoTo /D <6571756174696F6E2E372E352E302E31> >> >>}
\MusicPDFWriteObject{3591}{<< /Type /Annot /Subtype /Link /Border [0 0 0] /H /I /C [1 0 0] /Rect [73.712 448.143 108.905 460.096] /A << /S /GoTo /D <6571756174696F6E2E342E362E312E3134> >> >>}
\MusicPDFWriteObject{3592}{<< /Type /Annot /Subtype /Link /Border [0 0 0] /H /I /C [1 0 0] /Rect [335.08 325.119 367.369 337.327] /A << /S /GoTo /D <74686D2E422E362E312E36> >> >>}
\MusicPDFWriteObject{3593}{<< /Type /Annot /Subtype /Link /Border [0 0 0] /H /I /C [1 0 0] /Rect [425.58 325.119 457.868 337.327] /A << /S /GoTo /D <74686D2E422E312E312E34> >> >>}
\MusicPDFWriteObject{3594}{<< /Type /Annot /Subtype /Link /Border [0 0 0] /H /I /C [1 0 0] /Rect [193.121 312.222 228.314 324.175] /A << /S /GoTo /D <6571756174696F6E2E342E362E312E3137> >> >>}
\MusicPDFWriteObject{3595}{<< /Type /Annot /Subtype /Link /Border [0 0 0] /H /I /C [1 0 0] /Rect [287.02 312.222 317.29 324.175] /A << /S /GoTo /D <74686D2E342E362E312E35> >> >>}
\MusicPDFWriteObject{3596}{<< /Type /Annot /Subtype /Link /Border [0 0 0] /H /I /C [1 0 0] /Rect [365.461 312.222 400.655 324.175] /A << /S /GoTo /D <6571756174696F6E2E372E332E312E3137> >> >>}
\MusicPDFWriteObject{3597}{<< /Type /Annot /Subtype /Link /Border [0 0 0] /H /I /C [1 0 0] /Rect [459.361 312.222 489.63 324.175] /A << /S /GoTo /D <74686D2E372E332E312E36> >> >>}
\MusicPDFWriteObject{3598}{<< /Type /Annot /Subtype /Link /Border [0 0 0] /H /I /C [1 0 0] /Rect [432.4 265.502 462.049 284.168] /A << /S /GoTo /D <74686D2E372E312E322E32> >> >>}
\MusicPDFWriteObject{3599}{<< /Type /Annot /Subtype /Link /Border [0 0 0] /H /I /C [1 0 0] /Rect [206.9 225.268 237.113 237.221] /A << /S /GoTo /D <6571756174696F6E2E372E352E302E31> >> >>}
\MusicPDFWriteObject{3600}{<< /Type /Annot /Subtype /Link /Border [0 0 0] /H /I /C [1 0 0] /Rect [210.563 196.267 240.776 207.058] /A << /S /GoTo /D <74686D2E372E352E302E31> >> >>}
\MusicPDFWriteObject{3601}{<< /Type /Annot /Subtype /Link /Border [0 0 0] /H /I /C [1 0 0] /Rect [179.42 146.689 209.435 157.48] /A << /S /GoTo /D <74686D2E362E322E312E31> >> >>}
\MusicPDFWriteObject{3602}{<< /Type /Annot /Subtype /Link /Border [0 0 0] /H /I /C [1 0 0] /Rect [455.11 82.596 490.303 94.549] /A << /S /GoTo /D <6571756174696F6E2E372E352E302E3136> >> >>}
\MusicPDFWriteObject{3603}{<< /Type /Annot /Subtype /Link /Border [0 0 0] /H /I /C [1 0 0] /Rect [436.71 658.332 466.979 670.284] /A << /S /GoTo /D <74686D2E362E322E312E31> >> >>}
\MusicPDFWriteObject{3604}{<< /Type /Annot /Subtype /Link /Border [0 0 0] /H /I /C [1 0 0] /Rect [399.715 644.212 429.928 657.134] /A << /S /GoTo /D <74686D2E352E332E312E35> >> >>}
\MusicPDFWriteObject{3605}{<< /Type /Annot /Subtype /Link /Border [0 0 0] /H /I /C [1 0 0] /Rect [258.68 618.879 293.874 630.832] /A << /S /GoTo /D <6571756174696F6E2E372E352E302E3135> >> >>}
\MusicPDFWriteObject{3606}{<< /Type /Annot /Subtype /Link /Border [0 0 0] /H /I /C [0 1 0] /Rect [398.746 535.578 416.218 544.445] /A << /S /GoTo /D <636974652E4346> >> >>}
\MusicPDFWriteObject{3607}{<< /Type /Annot /Subtype /Link /Border [0 0 0] /H /I /C [1 0 0] /Rect [333.538 457.761 363.751 468.552] /A << /S /GoTo /D <74686D2E372E352E302E31> >> >>}
\MusicPDFWriteObject{3608}{<< /Type /Annot /Subtype /Link /Border [0 0 0] /H /I /C [1 0 0] /Rect [340.177 291.833 370.39 302.624] /A << /S /GoTo /D <74686D2E372E352E302E31> >> >>}
\MusicPDFWriteObject{3609}{<< /Type /Annot /Subtype /Link /Border [0 0 0] /H /I /C [1 0 0] /Rect [436.486 291.833 466.7 302.624] /A << /S /GoTo /D <74686D2E362E332E322E32> >> >>}
\MusicPDFWriteObject{3610}{<< /Type /Annot /Subtype /Link /Border [0 0 0] /H /I /C [1 0 0] /Rect [82.959 558.462 214.545 569.252] /A << /S /GoTo /D <73656374696F6E2E382E31> >> >>}
\MusicPDFWriteObject{3611}{<< /Type /Annot /Subtype /Link /Border [0 0 0] /H /I /C [1 0 0] /Rect [82.959 547.248 323.498 556.102] /A << /S /GoTo /D <73656374696F6E2E382E32> >> >>}
\MusicPDFWriteObject{3612}{<< /Type /Annot /Subtype /Link /Border [0 0 0] /H /I /C [1 0 0] /Rect [82.959 532.16 218.418 542.951] /A << /S /GoTo /D <73656374696F6E2E382E33> >> >>}
\MusicPDFWriteObject{3613}{<< /Type /Annot /Subtype /Link /Border [0 0 0] /H /I /C [1 0 0] /Rect [82.959 520.946 173.321 529.8] /A << /S /GoTo /D <73656374696F6E2E382E34> >> >>}
\MusicPDFWriteObject{3614}{<< /Type /Annot /Subtype /Link /Border [0 0 0] /H /I /C [1 0 0] /Rect [82.959 505.859 177.941 516.649] /A << /S /GoTo /D <73656374696F6E2E382E35> >> >>}
\MusicPDFWriteObject{3615}{<< /Type /Annot /Subtype /Link /Border [0 0 0] /H /I /C [1 0 0] /Rect [70.755 453.256 100.911 464.046] /A << /S /GoTo /D <74686D2E372E312E322E33> >> >>}
\MusicPDFWriteObject{3616}{<< /Type /Annot /Subtype /Link /Border [0 0 0] /H /I /C [1 0 0] /Rect [260.384 244.526 291.162 255.317] /A << /S /GoTo /D <74686D2E372E352E302E31> >> >>}
\MusicPDFWriteObject{3617}{<< /Type /Annot /Subtype /Link /Border [0 0 0] /H /I /C [1 0 0] /Rect [127.071 231.375 157.848 242.166] /A << /S /GoTo /D <74686D2E372E332E312E36> >> >>}
\MusicPDFWriteObject{3618}{<< /Type /Annot /Subtype /Link /Border [0 0 0] /H /I /C [1 0 0] /Rect [121.457 188.015 151.67 199.968] /A << /S /GoTo /D <6571756174696F6E2E372E332E312E35> >> >>}
\MusicPDFWriteObject{3619}{<< /Type /Annot /Subtype /Link /Border [0 0 0] /H /I /C [1 0 0] /Rect [180.369 188.015 210.582 199.968] /A << /S /GoTo /D <6571756174696F6E2E372E332E342E36> >> >>}
\MusicPDFWriteObject{3620}{<< /Type /Annot /Subtype /Link /Border [0 0 0] /H /I /C [1 0 0] /Rect [112.727 671.482 142.94 683.435] /A << /S /GoTo /D <74686D2E372E332E332E34> >> >>}
\MusicPDFWriteObject{3621}{<< /Type /Annot /Subtype /Link /Border [0 0 0] /H /I /C [1 0 0] /Rect [98.422 578.056 128.635 590.008] /A << /S /GoTo /D <6571756174696F6E2E372E332E332E35> >> >>}
\MusicPDFWriteObject{3622}{<< /Type /Annot /Subtype /Link /Border [0 0 0] /H /I /C [1 0 0] /Rect [157.077 578.056 187.29 590.008] /A << /S /GoTo /D <6571756174696F6E2E372E332E332E36> >> >>}
\MusicPDFWriteObject{3623}{<< /Type /Annot /Subtype /Link /Border [0 0 0] /H /I /C [1 0 0] /Rect [332.537 251.072 362.75 263.025] /A << /S /GoTo /D <74686D2E362E322E312E31> >> >>}
\MusicPDFWriteObject{3624}{<< /Type /Annot /Subtype /Link /Border [0 0 0] /H /I /C [1 0 0] /Rect [289.964 213.823 320.544 225.776] /A << /S /GoTo /D <74686D2E352E342E312E31> >> >>}
\MusicPDFWriteObject{3625}{<< /Type /Annot /Subtype /Link /Border [0 0 0] /H /I /C [1 0 0] /Rect [438.446 213.823 468.659 225.776] /A << /S /GoTo /D <6571756174696F6E2E352E312E312E32> >> >>}
\MusicPDFWriteObject{3626}{<< /Type /Annot /Subtype /Link /Border [0 0 0] /H /I /C [1 0 0] /Rect [179.252 684.633 209.465 696.586] /A << /S /GoTo /D <6571756174696F6E2E372E322E322E32> >> >>}
\MusicPDFWriteObject{3627}{<< /Type /Annot /Subtype /Link /Border [0 0 0] /H /I /C [1 0 0] /Rect [237.906 684.633 268.12 696.586] /A << /S /GoTo /D <6571756174696F6E2E332E332E312E33> >> >>}
\MusicPDFWriteObject{3628}{<< /Type /Annot /Subtype /Link /Border [0 0 0] /H /I /C [1 0 0] /Rect [291.709 590.985 321.922 602.937] /A << /S /GoTo /D <6571756174696F6E2E332E322E312E37> >> >>}
\MusicPDFWriteObject{3629}{<< /Type /Annot /Subtype /Link /Border [0 0 0] /H /I /C [1 0 0] /Rect [396.833 590.624 426.962 602.939] /A << /S /GoTo /D <74686D2E332E322E312E34> >> >>}
\MusicPDFWriteObject{3630}{<< /Type /Annot /Subtype /Link /Border [0 0 0] /H /I /C [1 0 0] /Rect [167.239 402.435 197.96 420.615] /A << /S /GoTo /D <74686D2E332E332E312E35> >> >>}
\MusicPDFWriteObject{3631}{<< /Type /Annot /Subtype /Link /Border [0 0 0] /H /I /C [1 0 0] /Rect [146.65 345.054 176.863 355.845] /A << /S /GoTo /D <74686D2E372E352E302E31> >> >>}
\MusicPDFWriteObject{3632}{<< /Type /Annot /Subtype /Link /Border [0 0 0] /H /I /C [1 0 0] /Rect [255.075 297.646 285.288 309.599] /A << /S /GoTo /D <6571756174696F6E2E382E312E302E38> >> >>}
\MusicPDFWriteObject{3633}{<< /Type /Annot /Subtype /Link /Border [0 0 0] /H /I /C [1 0 0] /Rect [140.37 219.177 170.583 231.129] /A << /S /GoTo /D <6571756174696F6E2E352E312E312E34> >> >>}
\MusicPDFWriteObject{3634}{<< /Type /Annot /Subtype /Link /Border [0 0 0] /H /I /C [1 0 0] /Rect [126.919 602.697 157.132 613.488] /A << /S /GoTo /D <74686D2E372E332E342E31> >> >>}
\MusicPDFWriteObject{3635}{<< /Type /Annot /Subtype /Link /Border [0 0 0] /H /I /C [1 0 0] /Rect [301.273 588.993 331.486 600.946] /A << /S /GoTo /D <6571756174696F6E2E382E312E302E31> >> >>}
\MusicPDFWriteObject{3636}{<< /Type /Annot /Subtype /Link /Border [0 0 0] /H /I /C [0 1 0] /Rect [455.361 591.483 474.465 600.351] /A << /S /GoTo /D <636974652E4E42> >> >>}
\MusicPDFWriteObject{3637}{<< /Type /Annot /Subtype /Link /Border [0 0 0] /H /I /C [1 0 0] /Rect [226.633 563.245 256.846 574.035] /A << /S /GoTo /D <74686D2E372E312E312E31> >> >>}
\MusicPDFWriteObject{3638}{<< /Type /Annot /Subtype /Link /Border [0 0 0] /H /I /C [1 0 0] /Rect [430.101 525.44 460.314 537.395] /A << /S /GoTo /D <74686D2E372E352E302E31> >> >>}
\MusicPDFWriteObject{3639}{<< /Type /Annot /Subtype /Link /Border [0 0 0] /H /I /C [1 0 0] /Rect [120.805 447.038 151.018 458.991] /A << /S /GoTo /D <6571756174696F6E2E372E332E312E35> >> >>}
\MusicPDFWriteObject{3640}{<< /Type /Annot /Subtype /Link /Border [0 0 0] /H /I /C [1 0 0] /Rect [226.106 447.038 256.319 458.991] /A << /S /GoTo /D <74686D2E372E332E312E33> >> >>}
\MusicPDFWriteObject{3641}{<< /Type /Annot /Subtype /Link /Border [0 0 0] /H /I /C [1 0 0] /Rect [280.887 447.038 311.101 458.991] /A << /S /GoTo /D <6571756174696F6E2E372E332E342E33> >> >>}
\MusicPDFWriteObject{3642}{<< /Type /Annot /Subtype /Link /Border [0 0 0] /H /I /C [1 0 0] /Rect [381.624 447.038 411.837 458.991] /A << /S /GoTo /D <6E6F746174696F6E2E372E332E342E31> >> >>}
\MusicPDFWriteObject{3643}{<< /Type /Annot /Subtype /Link /Border [0 0 0] /H /I /C [1 0 0] /Rect [304.035 434.441 333.684 445.231] /A << /S /GoTo /D <74686D2E382E332E302E31> >> >>}
\MusicPDFWriteObject{3644}{<< /Type /Annot /Subtype /Link /Border [0 0 0] /H /I /C [1 0 0] /Rect [187.955 208.054 218.168 220.007] /A << /S /GoTo /D <6571756174696F6E2E382E322E302E31> >> >>}
\MusicPDFWriteObject{3645}{<< /Type /Annot /Subtype /Link /Border [0 0 0] /H /I /C [1 0 0] /Rect [347.011 194.903 377.224 206.856] /A << /S /GoTo /D <6571756174696F6E2E382E322E302E32> >> >>}
\MusicPDFWriteObject{3646}{<< /Type /Annot /Subtype /Link /Border [0 0 0] /H /I /C [1 0 0] /Rect [173.238 608.082 203.451 620.035] /A << /S /GoTo /D <6571756174696F6E2E382E322E302E36> >> >>}
\MusicPDFWriteObject{3647}{<< /Type /Annot /Subtype /Link /Border [0 0 0] /H /I /C [1 0 0] /Rect [215.648 608.082 245.861 620.035] /A << /S /GoTo /D <6571756174696F6E2E382E322E302E37> >> >>}
\MusicPDFWriteObject{3648}{<< /Type /Annot /Subtype /Link /Border [0 0 0] /H /I /C [1 0 0] /Rect [275.103 608.082 305.316 620.035] /A << /S /GoTo /D <6571756174696F6E2E382E322E302E38> >> >>}
\MusicPDFWriteObject{3649}{<< /Type /Annot /Subtype /Link /Border [0 0 0] /H /I /C [1 0 0] /Rect [369.944 608.082 400.721 620.035] /A << /S /GoTo /D <74686D2E382E322E302E31> >> >>}
\MusicPDFWriteObject{3650}{<< /Type /Annot /Subtype /Link /Border [0 0 0] /H /I /C [0 1 0] /Rect [107.11 111.443 150.133 120.31] /A << /S /GoTo /D <636974652E426F752D4754> >> >>}
\MusicPDFWriteObject{3651}{<< /Type /Annot /Subtype /Link /Border [0 0 0] /H /I /C [0 1 0] /Rect [297.583 647.671 315.056 656.539] /A << /S /GoTo /D <636974652E4346> >> >>}
\MusicPDFWriteObject{3652}{<< /Type /Annot /Subtype /Link /Border [0 0 0] /H /I /C [0 1 0] /Rect [422.984 550.268 466.007 559.136] /A << /S /GoTo /D <636974652E426F752D4754> >> >>}
\MusicPDFWriteObject{3653}{<< /Type /Annot /Subtype /Link /Border [0 0 0] /H /I /C [1 0 0] /Rect [236.395 498.763 270.227 510.718] /A << /S /GoTo /D <74686D2E412E322E322E35> >> >>}
\MusicPDFWriteObject{3654}{<< /Type /Annot /Subtype /Link /Border [0 0 0] /H /I /C [1 0 0] /Rect [277.374 423.317 307.587 434.108] /A << /S /GoTo /D <74686D2E382E322E302E33> >> >>}
\MusicPDFWriteObject{3655}{<< /Type /Annot /Subtype /Link /Border [0 0 0] /H /I /C [1 0 0] /Rect [435.821 311.081 466.034 323.034] /A << /S /GoTo /D <6571756174696F6E2E372E332E312E35> >> >>}
\MusicPDFWriteObject{3656}{<< /Type /Annot /Subtype /Link /Border [0 0 0] /H /I /C [1 0 0] /Rect [494.15 311.081 524.363 323.034] /A << /S /GoTo /D <6571756174696F6E2E372E332E342E33> >> >>}
\MusicPDFWriteObject{3657}{<< /Type /Annot /Subtype /Link /Border [0 0 0] /H /I /C [1 0 0] /Rect [99.186 157.962 129.399 169.915] /A << /S /GoTo /D <6571756174696F6E2E372E332E312E35> >> >>}
\MusicPDFWriteObject{3658}{<< /Type /Annot /Subtype /Link /Border [0 0 0] /H /I /C [1 0 0] /Rect [194.706 157.962 224.919 169.915] /A << /S /GoTo /D <6571756174696F6E2E372E332E342E33> >> >>}
\MusicPDFWriteObject{3659}{<< /Type /Annot /Subtype /Link /Border [0 0 0] /H /I /C [1 0 0] /Rect [216.073 293.078 246.85 308.066] /A << /S /GoTo /D <74686D2E382E312E302E31> >> >>}
\MusicPDFWriteObject{3660}{<< /Type /Annot /Subtype /Link /Border [0 0 0] /H /I /C [1 0 0] /Rect [509.731 268.146 539.944 280.098] /A << /S /GoTo /D <6571756174696F6E2E382E312E302E34> >> >>}
\MusicPDFWriteObject{3661}{<< /Type /Annot /Subtype /Link /Border [0 0 0] /H /I /C [1 0 0] /Rect [71.004 213.535 101.217 225.24] /A << /S /GoTo /D <74686D2E382E322E302E34> >> >>}
\MusicPDFWriteObject{3662}{<< /Type /Annot /Subtype /Link /Border [0 0 0] /H /I /C [1 0 0] /Rect [191.906 166.129 222.119 178.082] /A << /S /GoTo /D <6571756174696F6E2E382E322E302E36> >> >>}
\MusicPDFWriteObject{3663}{<< /Type /Annot /Subtype /Link /Border [0 0 0] /H /I /C [1 0 0] /Rect [297.207 165.769 327.42 178.082] /A << /S /GoTo /D <74686D2E382E322E302E31> >> >>}
\MusicPDFWriteObject{3664}{<< /Type /Annot /Subtype /Link /Border [0 0 0] /H /I /C [1 0 0] /Rect [124.623 111.519 154.836 123.223] /A << /S /GoTo /D <74686D2E382E322E302E35> >> >>}
\MusicPDFWriteObject{3665}{<< /Type /Annot /Subtype /Link /Border [0 0 0] /H /I /C [1 0 0] /Rect [182.776 684.633 212.989 696.586] /A << /S /GoTo /D <6571756174696F6E2E382E322E302E36> >> >>}
\MusicPDFWriteObject{3666}{<< /Type /Annot /Subtype /Link /Border [0 0 0] /H /I /C [1 0 0] /Rect [288.077 684.273 318.29 696.586] /A << /S /GoTo /D <74686D2E382E322E302E31> >> >>}
\MusicPDFWriteObject{3667}{<< /Type /Annot /Subtype /Link /Border [0 0 0] /H /I /C [1 0 0] /Rect [190.835 489.176 221.048 499.967] /A << /S /GoTo /D <74686D2E372E312E322E33> >> >>}
\MusicPDFWriteObject{3668}{<< /Type /Annot /Subtype /Link /Border [0 0 0] /H /I /C [1 0 0] /Rect [260.293 462.321 290.845 474.274] /A << /S /GoTo /D <74686D2E372E312E312E31> >> >>}
\MusicPDFWriteObject{3669}{<< /Type /Annot /Subtype /Link /Border [0 0 0] /H /I /C [1 0 0] /Rect [365.965 462.321 398.617 474.274] /A << /S /GoTo /D <74686D2E422E352E332E33> >> >>}
\MusicPDFWriteObject{3670}{<< /Type /Annot /Subtype /Link /Border [0 0 0] /H /I /C [1 0 0] /Rect [241.258 449.171 272.94 461.123] /A << /S /GoTo /D <74686D2E422E352E322E31> >> >>}
\MusicPDFWriteObject{3671}{<< /Type /Annot /Subtype /Link /Border [0 0 0] /H /I /C [1 0 0] /Rect [284.794 305.556 315.571 316.347] /A << /S /GoTo /D <74686D2E372E332E312E36> >> >>}
\MusicPDFWriteObject{3672}{<< /Type /Annot /Subtype /Link /Border [0 0 0] /H /I /C [1 0 0] /Rect [220.992 278.701 251.064 290.654] /A << /S /GoTo /D <74686D2E382E332E302E31> >> >>}
\MusicPDFWriteObject{3673}{<< /Type /Annot /Subtype /Link /Border [0 0 0] /H /I /C [1 0 0] /Rect [313.55 278.701 343.622 290.654] /A << /S /GoTo /D <74686D2E372E332E342E31> >> >>}
\MusicPDFWriteObject{3674}{<< /Type /Annot /Subtype /Link /Border [0 0 0] /H /I /C [1 0 0] /Rect [212.431 110.972 243.209 121.762] /A << /S /GoTo /D <74686D2E382E332E302E31> >> >>}
\MusicPDFWriteObject{3675}{<< /Type /Annot /Subtype /Link /Border [0 0 0] /H /I /C [1 0 0] /Rect [347.885 110.972 378.662 121.762] /A << /S /GoTo /D <74686D2E382E312E302E31> >> >>}
\MusicPDFWriteObject{3676}{<< /Type /Annot /Subtype /Link /Border [0 0 0] /H /I /C [1 0 0] /Rect [165.41 511.866 195.623 523.818] /A << /S /GoTo /D <6571756174696F6E2E382E342E302E36> >> >>}
\MusicPDFWriteObject{3677}{<< /Type /Annot /Subtype /Link /Border [0 0 0] /H /I /C [1 0 0] /Rect [282.99 499.268 313.203 510.059] /A << /S /GoTo /D <74686D2E352E332E342E34> >> >>}
\MusicPDFWriteObject{3678}{<< /Type /Annot /Subtype /Link /Border [0 0 0] /H /I /C [1 0 0] /Rect [387.517 499.268 417.73 510.059] /A << /S /GoTo /D <74686D2E352E332E342E35> >> >>}
\MusicPDFWriteObject{3679}{<< /Type /Annot /Subtype /Link /Border [0 0 0] /H /I /C [1 0 0] /Rect [387.013 405.838 417.226 417.791] /A << /S /GoTo /D <6571756174696F6E2E352E312E312E32> >> >>}
\MusicPDFWriteObject{3680}{<< /Type /Annot /Subtype /Link /Border [0 0 0] /H /I /C [1 0 0] /Rect [509.731 405.838 539.944 417.791] /A << /S /GoTo /D <6571756174696F6E2E382E342E302E36> >> >>}
\MusicPDFWriteObject{3681}{<< /Type /Annot /Subtype /Link /Border [0 0 0] /H /I /C [1 0 0] /Rect [306.515 228.504 341.708 240.456] /A << /S /GoTo /D <6571756174696F6E2E352E332E342E3133> >> >>}
\MusicPDFWriteObject{3682}{<< /Type /Annot /Subtype /Link /Border [0 0 0] /H /I /C [1 0 0] /Rect [420.213 227.534 449.862 240.456] /A << /S /GoTo /D <74686D2E352E332E342E35> >> >>}
\MusicPDFWriteObject{3683}{<< /Type /Annot /Subtype /Link /Border [0 0 0] /H /I /C [1 0 0] /Rect [509.731 228.504 539.944 240.456] /A << /S /GoTo /D <6571756174696F6E2E372E322E322E31> >> >>}
\MusicPDFWriteObject{3684}{<< /Type /Annot /Subtype /Link /Border [0 0 0] /H /I /C [1 0 0] /Rect [297.966 215.353 328.743 227.305] /A << /S /GoTo /D <74686D2E352E342E312E31> >> >>}
\MusicPDFWriteObject{3685}{<< /Type /Annot /Subtype /Link /Border [0 0 0] /H /I /C [1 0 0] /Rect [220.097 202.202 250.311 214.155] /A << /S /GoTo /D <6571756174696F6E2E382E342E302E37> >> >>}
\MusicPDFWriteObject{3686}{<< /Type /Annot /Subtype /Link /Border [0 0 0] /H /I /C [1 0 0] /Rect [338.229 175.688 368.443 187.64] /A << /S /GoTo /D <6571756174696F6E2E332E332E312E32> >> >>}
\MusicPDFWriteObject{3687}{<< /Type /Annot /Subtype /Link /Border [0 0 0] /H /I /C [1 0 0] /Rect [427.069 175.688 457.282 187.64] /A << /S /GoTo /D <74686D2E332E332E312E31> >> >>}
\MusicPDFWriteObject{3688}{<< /Type /Annot /Subtype /Link /Border [0 0 0] /H /I /C [1 0 0] /Rect [73.712 149.386 103.925 161.339] /A << /S /GoTo /D <6571756174696F6E2E382E342E302E37> >> >>}
\MusicPDFWriteObject{3689}{<< /Type /Annot /Subtype /Link /Border [0 0 0] /H /I /C [1 0 0] /Rect [73.712 82.596 103.925 94.549] /A << /S /GoTo /D <6571756174696F6E2E382E342E302E36> >> >>}
\MusicPDFWriteObject{3690}{<< /Type /Annot /Subtype /Link /Border [0 0 0] /H /I /C [1 0 0] /Rect [132.366 82.596 162.579 94.549] /A << /S /GoTo /D <6571756174696F6E2E382E342E302E37> >> >>}
\MusicPDFWriteObject{3691}{<< /Type /Annot /Subtype /Link /Border [0 0 0] /H /I /C [1 0 0] /Rect [338.716 524.301 368.929 536.254] /A << /S /GoTo /D <6571756174696F6E2E382E342E302E38> >> >>}
\MusicPDFWriteObject{3692}{<< /Type /Annot /Subtype /Link /Border [0 0 0] /H /I /C [1 0 0] /Rect [185.553 410.967 215.766 422.92] /A << /S /GoTo /D <6571756174696F6E2E382E342E302E36> >> >>}
\MusicPDFWriteObject{3693}{<< /Type /Annot /Subtype /Link /Border [0 0 0] /H /I /C [1 0 0] /Rect [228.119 410.967 258.332 422.92] /A << /S /GoTo /D <6571756174696F6E2E382E342E302E37> >> >>}
\MusicPDFWriteObject{3694}{<< /Type /Annot /Subtype /Link /Border [0 0 0] /H /I /C [1 0 0] /Rect [270.685 410.967 300.899 422.92] /A << /S /GoTo /D <6571756174696F6E2E382E342E302E38> >> >>}
\MusicPDFWriteObject{3695}{<< /Type /Annot /Subtype /Link /Border [0 0 0] /H /I /C [1 0 0] /Rect [330.391 410.967 360.605 422.92] /A << /S /GoTo /D <6571756174696F6E2E382E342E302E39> >> >>}
\MusicPDFWriteObject{3696}{<< /Type /Annot /Subtype /Link /Border [0 0 0] /H /I /C [1 0 0] /Rect [192.033 291.065 222.246 303.017] /A << /S /GoTo /D <6571756174696F6E2E382E342E302E36> >> >>}
\MusicPDFWriteObject{3697}{<< /Type /Annot /Subtype /Link /Border [0 0 0] /H /I /C [1 0 0] /Rect [234.085 291.065 264.298 303.017] /A << /S /GoTo /D <6571756174696F6E2E382E342E302E37> >> >>}
\MusicPDFWriteObject{3698}{<< /Type /Annot /Subtype /Link /Border [0 0 0] /H /I /C [1 0 0] /Rect [276.137 291.065 306.35 303.017] /A << /S /GoTo /D <6571756174696F6E2E382E342E302E38> >> >>}
\MusicPDFWriteObject{3699}{<< /Type /Annot /Subtype /Link /Border [0 0 0] /H /I /C [1 0 0] /Rect [372.49 291.065 402.703 303.017] /A << /S /GoTo /D <6571756174696F6E2E382E342E302E39> >> >>}
\MusicPDFWriteObject{3700}{<< /Type /Annot /Subtype /Link /Border [0 0 0] /H /I /C [1 0 0] /Rect [276.908 277.914 306.557 289.867] /A << /S /GoTo /D <74686D2E382E312E302E31> >> >>}
\MusicPDFWriteObject{3701}{<< /Type /Annot /Subtype /Link /Border [0 0 0] /H /I /C [1 0 0] /Rect [365.744 277.914 395.392 289.867] /A << /S /GoTo /D <74686D2E372E332E342E31> >> >>}
\MusicPDFWriteObject{3702}{<< /Type /Annot /Subtype /Link /Border [0 0 0] /H /I /C [1 0 0] /Rect [385.987 198.186 416.201 210.139] /A << /S /GoTo /D <6571756174696F6E2E382E342E302E36> >> >>}
\MusicPDFWriteObject{3703}{<< /Type /Annot /Subtype /Link /Border [0 0 0] /H /I /C [1 0 0] /Rect [383.498 158.734 413.711 170.686] /A << /S /GoTo /D <6571756174696F6E2E382E342E302E36> >> >>}
\MusicPDFWriteObject{3704}{<< /Type /Annot /Subtype /Link /Border [0 0 0] /H /I /C [1 0 0] /Rect [126.919 552.228 157.132 563.019] /A << /S /GoTo /D <74686D2E372E342E322E31> >> >>}
\MusicPDFWriteObject{3705}{<< /Type /Annot /Subtype /Link /Border [0 0 0] /H /I /C [1 0 0] /Rect [214.059 524.552 244.272 536.505] /A << /S /GoTo /D <74686D2E372E342E332E32> >> >>}
\MusicPDFWriteObject{3706}{<< /Type /Annot /Subtype /Link /Border [0 0 0] /H /I /C [1 0 0] /Rect [452.04 320.568 482.563 331.358] /A << /S /GoTo /D <74686D2E372E342E322E31> >> >>}
\MusicPDFWriteObject{3707}{<< /Type /Annot /Subtype /Link /Border [0 0 0] /H /I /C [1 0 0] /Rect [70.755 307.417 100.968 318.208] /A << /S /GoTo /D <74686D2E372E342E332E32> >> >>}
\MusicPDFWriteObject{3708}{<< /Type /Annot /Subtype /Link /Border [0 0 0] /H /I /C [1 0 0] /Rect [124.581 172.81 156.839 184.763] /A << /S /GoTo /D <74686D2E422E342E312E32> >> >>}
\MusicPDFWriteObject{3709}{<< /Type /Annot /Subtype /Link /Border [0 0 0] /H /I /C [1 0 0] /Rect [393.551 172.81 425.809 184.763] /A << /S /GoTo /D <74686D2E422E332E322E31> >> >>}
\MusicPDFWriteObject{3710}{<< /Type /Annot /Subtype /Link /Border [0 0 0] /H /I /C [1 0 0] /Rect [488.555 172.81 520.813 184.763] /A << /S /GoTo /D <74686D2E422E342E322E32> >> >>}
\MusicPDFWriteObject{3711}{<< /Type /Annot /Subtype /Link /Border [0 0 0] /H /I /C [1 0 0] /Rect [229.012 684.19 264.206 696.143] /A << /S /GoTo /D <6571756174696F6E2E382E342E302E3130> >> >>}
\MusicPDFWriteObject{3712}{<< /Type /Annot /Subtype /Link /Border [0 0 0] /H /I /C [0 1 0] /Rect [225.331 673.53 247.722 682.398] /A << /S /GoTo /D <636974652E44656E> >> >>}
\MusicPDFWriteObject{3713}{<< /Type /Annot /Subtype /Link /Border [0 0 0] /H /I /C [1 0 0] /Rect [194.012 644.738 229.206 656.691] /A << /S /GoTo /D <6571756174696F6E2E382E342E302E3130> >> >>}
\MusicPDFWriteObject{3714}{<< /Type /Annot /Subtype /Link /Border [0 0 0] /H /I /C [1 0 0] /Rect [282.967 471.451 313.18 483.404] /A << /S /GoTo /D <6571756174696F6E2E382E342E302E36> >> >>}
\MusicPDFWriteObject{3715}{<< /Type /Annot /Subtype /Link /Border [0 0 0] /H /I /C [1 0 0] /Rect [198.25 276.875 228.463 287.666] /A << /S /GoTo /D <74686D2E372E332E332E34> >> >>}
\MusicPDFWriteObject{3716}{<< /Type /Annot /Subtype /Link /Border [0 0 0] /H /I /C [1 0 0] /Rect [456.316 165.874 486.727 177.826] /A << /S /GoTo /D <74686D2E372E332E322E33> >> >>}
\MusicPDFWriteObject{3717}{<< /Type /Annot /Subtype /Link /Border [0 0 0] /H /I /C [1 0 0] /Rect [430.443 410.603 460.656 424.605] /A << /S /GoTo /D <74686D2E372E342E312E31> >> >>}
\MusicPDFWriteObject{3718}{<< /Type /Annot /Subtype /Link /Border [0 0 0] /H /I /C [1 0 0] /Rect [298.24 333.746 328.453 345.698] /A << /S /GoTo /D <6571756174696F6E2E372E312E342E36> >> >>}
\MusicPDFWriteObject{3719}{<< /Type /Annot /Subtype /Link /Border [0 0 0] /H /I /C [1 0 0] /Rect [387.058 333.746 417.271 345.698] /A << /S /GoTo /D <74686D2E372E312E342E32> >> >>}
\MusicPDFWriteObject{3720}{<< /Type /Annot /Subtype /Link /Border [0 0 0] /H /I /C [1 0 0] /Rect [82.959 557.908 214.994 570.501] /A << /S /GoTo /D <73656374696F6E2E392E31> >> >>}
\MusicPDFWriteObject{3721}{<< /Type /Annot /Subtype /Link /Border [0 0 0] /H /I /C [1 0 0] /Rect [82.959 545.311 377.394 556.102] /A << /S /GoTo /D <73756273656374696F6E2E392E312E31> >> >>}
\MusicPDFWriteObject{3722}{<< /Type /Annot /Subtype /Link /Border [0 0 0] /H /I /C [1 0 0] /Rect [82.959 532.075 322.813 544.199] /A << /S /GoTo /D <73756273656374696F6E2E392E312E32> >> >>}
\MusicPDFWriteObject{3723}{<< /Type /Annot /Subtype /Link /Border [0 0 0] /H /I /C [1 0 0] /Rect [82.959 518.202 367.993 530.411] /A << /S /GoTo /D <73656374696F6E2E392E32> >> >>}
\MusicPDFWriteObject{3724}{<< /Type /Annot /Subtype /Link /Border [0 0 0] /H /I /C [1 0 0] /Rect [82.959 505.305 476.493 517.26] /A << /S /GoTo /D <73756273656374696F6E2E392E322E31> >> >>}
\MusicPDFWriteObject{3725}{<< /Type /Annot /Subtype /Link /Border [0 0 0] /H /I /C [1 0 0] /Rect [82.959 492.708 158.408 503.499] /A << /S /GoTo /D <73756273756273656374696F6E2E392E322E312E31> >> >>}
\MusicPDFWriteObject{3726}{<< /Type /Annot /Subtype /Link /Border [0 0 0] /H /I /C [1 0 0] /Rect [82.959 478.75 367.177 490.958] /A << /S /GoTo /D <73756273656374696F6E2E392E322E32> >> >>}
\MusicPDFWriteObject{3727}{<< /Type /Annot /Subtype /Link /Border [0 0 0] /H /I /C [1 0 0] /Rect [82.959 465.599 286.542 477.808] /A << /S /GoTo /D <73756273756273656374696F6E2E392E322E322E31> >> >>}
\MusicPDFWriteObject{3728}{<< /Type /Annot /Subtype /Link /Border [0 0 0] /H /I /C [1 0 0] /Rect [82.959 452.448 275.086 464.657] /A << /S /GoTo /D <73756273756273656374696F6E2E392E322E322E32> >> >>}
\MusicPDFWriteObject{3729}{<< /Type /Annot /Subtype /Link /Border [0 0 0] /H /I /C [1 0 0] /Rect [82.959 439.551 340.105 451.506] /A << /S /GoTo /D <73656374696F6E2E392E33> >> >>}
\MusicPDFWriteObject{3730}{<< /Type /Annot /Subtype /Link /Border [0 0 0] /H /I /C [1 0 0] /Rect [82.959 426.4 242.025 438.355] /A << /S /GoTo /D <73756273656374696F6E2E392E332E31> >> >>}
\MusicPDFWriteObject{3731}{<< /Type /Annot /Subtype /Link /Border [0 0 0] /H /I /C [1 0 0] /Rect [82.959 413.803 199.853 424.594] /A << /S /GoTo /D <73756273656374696F6E2E392E332E32> >> >>}
\MusicPDFWriteObject{3732}{<< /Type /Annot /Subtype /Link /Border [0 0 0] /H /I /C [1 0 0] /Rect [82.959 400.099 200.138 412.054] /A << /S /GoTo /D <73756273756273656374696F6E2E392E332E322E31> >> >>}
\MusicPDFWriteObject{3733}{<< /Type /Annot /Subtype /Link /Border [0 0 0] /H /I /C [1 0 0] /Rect [82.959 386.695 330.448 398.903] /A << /S /GoTo /D <73756273756273656374696F6E2E392E332E322E32> >> >>}
\MusicPDFWriteObject{3734}{<< /Type /Annot /Subtype /Link /Border [0 0 0] /H /I /C [1 0 0] /Rect [82.959 373.544 317.472 385.752] /A << /S /GoTo /D <73756273656374696F6E2E392E332E33> >> >>}
\MusicPDFWriteObject{3735}{<< /Type /Annot /Subtype /Link /Border [0 0 0] /H /I /C [1 0 0] /Rect [82.959 363.137 157.218 371.991] /A << /S /GoTo /D <73756273756273656374696F6E2E392E332E332E31> >> >>}
\MusicPDFWriteObject{3736}{<< /Type /Annot /Subtype /Link /Border [0 0 0] /H /I /C [1 0 0] /Rect [82.959 347.136 289.233 359.451] /A << /S /GoTo /D <73756273756273656374696F6E2E392E332E332E32> >> >>}
\MusicPDFWriteObject{3737}{<< /Type /Annot /Subtype /Link /Border [0 0 0] /H /I /C [0 1 0] /Rect [115.825 610.119 135.414 618.987] /A << /S /GoTo /D <636974652E4752> >> >>}
\MusicPDFWriteObject{3738}{<< /Type /Annot /Subtype /Link /Border [0 0 0] /H /I /C [0 1 0] /Rect [235.791 610.119 251.783 618.987] /A << /S /GoTo /D <636974652E436F> >> >>}
\MusicPDFWriteObject{3739}{<< /Type /Annot /Subtype /Link /Border [0 0 0] /H /I /C [0 1 0] /Rect [107.221 334.799 124.376 343.667] /A << /S /GoTo /D <636974652E466F6C> >> >>}
\MusicPDFWriteObject{3740}{<< /Type /Annot /Subtype /Link /Border [0 0 0] /H /I /C [0 1 0] /Rect [87.359 307.57 104.513 316.438] /A << /S /GoTo /D <636974652E466F6C> >> >>}
\MusicPDFWriteObject{3741}{<< /Type /Annot /Subtype /Link /Border [0 0 0] /H /I /C [0 1 0] /Rect [401.758 307.57 424.114 316.438] /A << /S /GoTo /D <636974652E436F68> >> >>}
\MusicPDFWriteObject{3742}{<< /Type /Annot /Subtype /Link /Border [0 0 0] /H /I /C [0 1 0] /Rect [194.783 271.677 214.371 280.545] /A << /S /GoTo /D <636974652E4752> >> >>}
\MusicPDFWriteObject{3743}{<< /Type /Annot /Subtype /Link /Border [0 0 0] /H /I /C [0 1 0] /Rect [238.937 258.526 256.824 267.394] /A << /S /GoTo /D <636974652E4647> >> >>}
\MusicPDFWriteObject{3744}{<< /Type /Annot /Subtype /Link /Border [0 0 0] /H /I /C [1 0 0] /Rect [515.944 256.036 545.592 267.989] /A << /S /GoTo /D <74686D2E392E312E312E31> >> >>}
\MusicPDFWriteObject{3745}{<< /Type /Annot /Subtype /Link /Border [0 0 0] /H /I /C [0 1 0] /Rect [140.651 245.375 160.24 254.243] /A << /S /GoTo /D <636974652E4752> >> >>}
\MusicPDFWriteObject{3746}{<< /Type /Annot /Subtype /Link /Border [0 0 0] /H /I /C [0 1 0] /Rect [353.057 245.375 372.646 254.243] /A << /S /GoTo /D <636974652E4752> >> >>}
\MusicPDFWriteObject{3747}{<< /Type /Annot /Subtype /Link /Border [0 0 0] /H /I /C [0 1 0] /Rect [523.57 232.225 543.159 241.092] /A << /S /GoTo /D <636974652E4752> >> >>}
\MusicPDFWriteObject{3748}{<< /Type /Annot /Subtype /Link /Border [0 0 0] /H /I /C [0 1 0] /Rect [209.749 542.642 228.051 551.51] /A << /S /GoTo /D <636974652E4445> >> >>}
\MusicPDFWriteObject{3749}{<< /Type /Annot /Subtype /Link /Border [0 0 0] /H /I /C [1 0 0] /Rect [511.675 408.669 541.888 420.622] /A << /S /GoTo /D <6571756174696F6E2E392E312E312E36> >> >>}
\MusicPDFWriteObject{3750}{<< /Type /Annot /Subtype /Link /Border [0 0 0] /H /I /C [0 1 0] /Rect [401.933 370.52 418.88 379.388] /A << /S /GoTo /D <636974652E5275> >> >>}
\MusicPDFWriteObject{3751}{<< /Type /Annot /Subtype /Link /Border [0 0 0] /H /I /C [0 1 0] /Rect [111.149 357.369 129.451 366.237] /A << /S /GoTo /D <636974652E4445> >> >>}
\MusicPDFWriteObject{3752}{<< /Type /Annot /Subtype /Link /Border [0 0 0] /H /I /C [0 1 0] /Rect [200.5 357.369 218.803 366.237] /A << /S /GoTo /D <636974652E4445> >> >>}
\MusicPDFWriteObject{3753}{<< /Type /Annot /Subtype /Link /Border [0 0 0] /H /I /C [0 1 0] /Rect [363.083 236.685 381.385 245.553] /A << /S /GoTo /D <636974652E4445> >> >>}
\MusicPDFWriteObject{3754}{<< /Type /Annot /Subtype /Link /Border [0 0 0] /H /I /C [0 1 0] /Rect [129.767 208.715 148.069 217.583] /A << /S /GoTo /D <636974652E4445> >> >>}
\MusicPDFWriteObject{3755}{<< /Type /Annot /Subtype /Link /Border [0 0 0] /H /I /C [1 0 0] /Rect [96.84 190.34 127.053 202.293] /A << /S /GoTo /D <6571756174696F6E2E392E312E312E38> >> >>}
\MusicPDFWriteObject{3756}{<< /Type /Annot /Subtype /Link /Border [0 0 0] /H /I /C [1 0 0] /Rect [436.663 190.34 466.876 202.293] /A << /S /GoTo /D <6571756174696F6E2E392E312E312E38> >> >>}
\MusicPDFWriteObject{3757}{<< /Type /Annot /Subtype /Link /Border [0 0 0] /H /I /C [0 1 0] /Rect [411.239 179.679 429.542 188.547] /A << /S /GoTo /D <636974652E4445> >> >>}
\MusicPDFWriteObject{3758}{<< /Type /Annot /Subtype /Link /Border [0 0 0] /H /I /C [0 1 0] /Rect [145.982 166.528 164.284 175.396] /A << /S /GoTo /D <636974652E4445> >> >>}
\MusicPDFWriteObject{3759}{<< /Type /Annot /Subtype /Link /Border [0 0 0] /H /I /C [0 1 0] /Rect [304.77 141.489 324.359 150.357] /A << /S /GoTo /D <636974652E4752> >> >>}
\MusicPDFWriteObject{3760}{<< /Type /Annot /Subtype /Link /Border [0 0 0] /H /I /C [1 0 0] /Rect [503.172 138.999 533.385 150.952] /A << /S /GoTo /D <74686D2E392E312E312E33> >> >>}
\MusicPDFWriteObject{3761}{<< /Type /Annot /Subtype /Link /Border [0 0 0] /H /I /C [1 0 0] /Rect [329.046 555.073 359.259 567.026] /A << /S /GoTo /D <6571756174696F6E2E392E312E312E33> >> >>}
\MusicPDFWriteObject{3762}{<< /Type /Annot /Subtype /Link /Border [0 0 0] /H /I /C [0 1 0] /Rect [146.588 527.433 163.742 536.301] /A << /S /GoTo /D <636974652E466F6C> >> >>}
\MusicPDFWriteObject{3763}{<< /Type /Annot /Subtype /Link /Border [0 0 0] /H /I /C [0 1 0] /Rect [211.222 415.585 228.377 424.453] /A << /S /GoTo /D <636974652E466F6C> >> >>}
\MusicPDFWriteObject{3764}{<< /Type /Annot /Subtype /Link /Border [0 0 0] /H /I /C [1 0 0] /Rect [355.597 277.484 390.79 289.437] /A << /S /GoTo /D <6571756174696F6E2E392E312E312E3130> >> >>}
\MusicPDFWriteObject{3765}{<< /Type /Annot /Subtype /Link /Border [0 0 0] /H /I /C [1 0 0] /Rect [223.921 214.296 259.114 226.249] /A << /S /GoTo /D <6571756174696F6E2E392E312E312E3136> >> >>}
\MusicPDFWriteObject{3766}{<< /Type /Annot /Subtype /Link /Border [0 0 0] /H /I /C [1 0 0] /Rect [158.913 200.785 189.464 213.1] /A << /S /GoTo /D <74686D2E392E312E312E33> >> >>}
\MusicPDFWriteObject{3767}{<< /Type /Annot /Subtype /Link /Border [0 0 0] /H /I /C [1 0 0] /Rect [394.022 148.542 424.179 160.497] /A << /S /GoTo /D <74686D2E392E312E312E32> >> >>}
\MusicPDFWriteObject{3768}{<< /Type /Annot /Subtype /Link /Border [0 0 0] /H /I /C [0 1 0] /Rect [221.396 124.731 270.436 133.599] /A << /S /GoTo /D <636974652E426F752D545653> >> >>}
\MusicPDFWriteObject{3769}{<< /Type /Annot /Subtype /Link /Border [0 0 0] /H /I /C [1 0 0] /Rect [148.168 109.643 185.574 120.434] /A << /S /GoTo /D <74686D2E432E322E332E3131> >> >>}
\MusicPDFWriteObject{3770}{<< /Type /Annot /Subtype /Link /Border [0 0 0] /H /I /C [0 1 0] /Rect [322.593 687.123 339.747 695.991] /A << /S /GoTo /D <636974652E466F6C> >> >>}
\MusicPDFWriteObject{3771}{<< /Type /Annot /Subtype /Link /Border [0 0 0] /H /I /C [0 1 0] /Rect [229.343 515.786 246.498 524.654] /A << /S /GoTo /D <636974652E466F6C> >> >>}
\MusicPDFWriteObject{3772}{<< /Type /Annot /Subtype /Link /Border [0 0 0] /H /I /C [1 0 0] /Rect [165.826 500.145 201.019 512.098] /A << /S /GoTo /D <6571756174696F6E2E392E312E312E3137> >> >>}
\MusicPDFWriteObject{3773}{<< /Type /Annot /Subtype /Link /Border [0 0 0] /H /I /C [1 0 0] /Rect [413.403 500.145 448.597 512.098] /A << /S /GoTo /D <6571756174696F6E2E392E312E312E3132> >> >>}
\MusicPDFWriteObject{3774}{<< /Type /Annot /Subtype /Link /Border [0 0 0] /H /I /C [1 0 0] /Rect [164.069 402.666 194.564 416.585] /A << /S /GoTo /D <74686D2E392E312E312E31> >> >>}
\MusicPDFWriteObject{3775}{<< /Type /Annot /Subtype /Link /Border [0 0 0] /H /I /C [1 0 0] /Rect [324.921 402.666 355.416 416.585] /A << /S /GoTo /D <74686D2E392E312E312E33> >> >>}
\MusicPDFWriteObject{3776}{<< /Type /Annot /Subtype /Link /Border [0 0 0] /H /I /C [1 0 0] /Rect [174.578 339.383 204.791 351.336] /A << /S /GoTo /D <6571756174696F6E2E392E312E312E32> >> >>}
\MusicPDFWriteObject{3777}{<< /Type /Annot /Subtype /Link /Border [0 0 0] /H /I /C [1 0 0] /Rect [184.125 273.303 216.967 285.256] /A << /S /GoTo /D <74686D2E442E322E312E31> >> >>}
\MusicPDFWriteObject{3778}{<< /Type /Annot /Subtype /Link /Border [0 0 0] /H /I /C [1 0 0] /Rect [262.674 273.303 287.768 285.256] /A << /S /GoTo /D <73756273656374696F6E2E442E332E33> >> >>}
\MusicPDFWriteObject{3779}{<< /Type /Annot /Subtype /Link /Border [0 0 0] /H /I /C [1 0 0] /Rect [230.962 220.7 263.772 232.653] /A << /S /GoTo /D <74686D2E442E332E332E37> >> >>}
\MusicPDFWriteObject{3780}{<< /Type /Annot /Subtype /Link /Border [0 0 0] /H /I /C [1 0 0] /Rect [251.103 114.198 281.316 126.151] /A << /S /GoTo /D <6571756174696F6E2E392E312E322E31> >> >>}
\MusicPDFWriteObject{3781}{<< /Type /Annot /Subtype /Link /Border [0 0 0] /H /I /C [1 0 0] /Rect [71.004 639.069 104.308 651.022] /A << /S /GoTo /D <74686D2E442E322E312E32> >> >>}
\MusicPDFWriteObject{3782}{<< /Type /Annot /Subtype /Link /Border [0 0 0] /H /I /C [0 1 0] /Rect [404.888 602.107 423.19 610.975] /A << /S /GoTo /D <636974652E4445> >> >>}
\MusicPDFWriteObject{3783}{<< /Type /Annot /Subtype /Link /Border [0 0 0] /H /I /C [1 0 0] /Rect [269.985 466.946 303.043 477.795] /A << /S /GoTo /D <74686D2E442E312E362E35> >> >>}
\MusicPDFWriteObject{3784}{<< /Type /Annot /Subtype /Link /Border [0 0 0] /H /I /C [1 0 0] /Rect [370.529 466.946 403.586 477.795] /A << /S /GoTo /D <74686D2E442E312E362E38> >> >>}
\MusicPDFWriteObject{3785}{<< /Type /Annot /Subtype /Link /Border [0 0 0] /H /I /C [1 0 0] /Rect [420.765 466.946 446.021 477.795] /A << /S /GoTo /D <73756273656374696F6E2E442E312E36> >> >>}
\MusicPDFWriteObject{3786}{<< /Type /Annot /Subtype /Link /Border [0 0 0] /H /I /C [0 1 0] /Rect [141.489 204.502 160.594 213.37] /A << /S /GoTo /D <636974652E4E42> >> >>}
\MusicPDFWriteObject{3787}{<< /Type /Annot /Subtype /Link /Border [0 0 0] /H /I /C [0 1 0] /Rect [318.822 687.123 367.863 695.991] /A << /S /GoTo /D <636974652E426F752D545653> >> >>}
\MusicPDFWriteObject{3788}{<< /Type /Annot /Subtype /Link /Border [0 0 0] /H /I /C [0 1 0] /Rect [107.335 548.512 125.638 557.38] /A << /S /GoTo /D <636974652E4445> >> >>}
\MusicPDFWriteObject{3789}{<< /Type /Annot /Subtype /Link /Border [0 0 0] /H /I /C [1 0 0] /Rect [371.113 342.122 406.306 354.075] /A << /S /GoTo /D <6571756174696F6E2E392E312E322E3131> >> >>}
\MusicPDFWriteObject{3790}{<< /Type /Annot /Subtype /Link /Border [0 0 0] /H /I /C [0 1 0] /Rect [526.005 344.612 543.159 353.48] /A << /S /GoTo /D <636974652E466F6C> >> >>}
\MusicPDFWriteObject{3791}{<< /Type /Annot /Subtype /Link /Border [0 0 0] /H /I /C [1 0 0] /Rect [257.503 284.969 295.325 293.823] /A << /S /GoTo /D <74686D2E442E312E362E3139> >> >>}
\MusicPDFWriteObject{3792}{<< /Type /Annot /Subtype /Link /Border [0 0 0] /H /I /C [1 0 0] /Rect [215.036 269.328 250.229 281.281] /A << /S /GoTo /D <6571756174696F6E2E392E312E322E3130> >> >>}
\MusicPDFWriteObject{3793}{<< /Type /Annot /Subtype /Link /Border [0 0 0] /H /I /C [1 0 0] /Rect [204.775 184.82 235.185 195.611] /A << /S /GoTo /D <74686D2E392E312E322E32> >> >>}
\MusicPDFWriteObject{3794}{<< /Type /Annot /Subtype /Link /Border [0 0 0] /H /I /C [1 0 0] /Rect [318.855 157.966 349.068 169.918] /A << /S /GoTo /D <6571756174696F6E2E392E312E322E35> >> >>}
\MusicPDFWriteObject{3795}{<< /Type /Annot /Subtype /Link /Border [0 0 0] /H /I /C [1 0 0] /Rect [495.495 157.966 525.708 169.918] /A << /S /GoTo /D <6571756174696F6E2E392E312E322E35> >> >>}
\MusicPDFWriteObject{3796}{<< /Type /Annot /Subtype /Link /Border [0 0 0] /H /I /C [1 0 0] /Rect [450.998 392.556 481.211 404.509] /A << /S /GoTo /D <6571756174696F6E2E392E312E322E33> >> >>}
\MusicPDFWriteObject{3797}{<< /Type /Annot /Subtype /Link /Border [0 0 0] /H /I /C [1 0 0] /Rect [179.998 272.242 210.606 284.195] /A << /S /GoTo /D <74686D2E392E312E312E35> >> >>}
\MusicPDFWriteObject{3798}{<< /Type /Annot /Subtype /Link /Border [0 0 0] /H /I /C [1 0 0] /Rect [348.354 272.242 383.548 284.195] /A << /S /GoTo /D <6571756174696F6E2E392E312E322E3134> >> >>}
\MusicPDFWriteObject{3799}{<< /Type /Annot /Subtype /Link /Border [0 0 0] /H /I /C [1 0 0] /Rect [506.695 272.242 541.888 284.195] /A << /S /GoTo /D <6571756174696F6E2E392E312E312E3132> >> >>}
\MusicPDFWriteObject{3800}{<< /Type /Annot /Subtype /Link /Border [0 0 0] /H /I /C [1 0 0] /Rect [379.4 232.741 414.593 244.693] /A << /S /GoTo /D <6571756174696F6E2E392E312E322E3134> >> >>}
\MusicPDFWriteObject{3801}{<< /Type /Annot /Subtype /Link /Border [0 0 0] /H /I /C [1 0 0] /Rect [296.596 219.59 331.79 231.543] /A << /S /GoTo /D <6571756174696F6E2E392E312E312E3131> >> >>}
\MusicPDFWriteObject{3802}{<< /Type /Annot /Subtype /Link /Border [0 0 0] /H /I /C [1 0 0] /Rect [455.936 219.59 491.129 231.543] /A << /S /GoTo /D <6571756174696F6E2E392E312E322E3134> >> >>}
\MusicPDFWriteObject{3803}{<< /Type /Annot /Subtype /Link /Border [0 0 0] /H /I /C [1 0 0] /Rect [195.872 166.987 231.065 178.94] /A << /S /GoTo /D <6571756174696F6E2E392E312E322E3135> >> >>}
\MusicPDFWriteObject{3804}{<< /Type /Annot /Subtype /Link /Border [0 0 0] /H /I /C [1 0 0] /Rect [445.293 508.13 476.07 521.459] /A << /S /GoTo /D <74686D2E392E312E322E34> >> >>}
\MusicPDFWriteObject{3805}{<< /Type /Annot /Subtype /Link /Border [0 0 0] /H /I /C [1 0 0] /Rect [294.298 494.238 329.491 506.191] /A << /S /GoTo /D <6571756174696F6E2E392E312E322E3137> >> >>}
\MusicPDFWriteObject{3806}{<< /Type /Annot /Subtype /Link /Border [0 0 0] /H /I /C [1 0 0] /Rect [203.033 439.902 228.404 450.693] /A << /S /GoTo /D <73756273656374696F6E2E442E332E32> >> >>}
\MusicPDFWriteObject{3807}{<< /Type /Annot /Subtype /Link /Border [0 0 0] /H /I /C [1 0 0] /Rect [92.672 413.047 118.229 425.002] /A << /S /GoTo /D <73756273656374696F6E2E442E332E33> >> >>}
\MusicPDFWriteObject{3808}{<< /Type /Annot /Subtype /Link /Border [0 0 0] /H /I /C [1 0 0] /Rect [98.625 363.469 131.466 375.422] /A << /S /GoTo /D <6571756174696F6E2E442E332E322E32> >> >>}
\MusicPDFWriteObject{3809}{<< /Type /Annot /Subtype /Link /Border [0 0 0] /H /I /C [1 0 0] /Rect [163.531 297.715 201.46 309.668] /A << /S /GoTo /D <73756273756273656374696F6E2E442E312E31302E31> >> >>}
\MusicPDFWriteObject{3810}{<< /Type /Annot /Subtype /Link /Border [0 0 0] /H /I /C [0 1 0] /Rect [227.765 300.205 247.423 309.073] /A << /S /GoTo /D <636974652E4B47> >> >>}
\MusicPDFWriteObject{3811}{<< /Type /Annot /Subtype /Link /Border [0 0 0] /H /I /C [1 0 0] /Rect [492.242 297.715 525.176 309.668] /A << /S /GoTo /D <74686D2E442E332E312E33> >> >>}
\MusicPDFWriteObject{3812}{<< /Type /Annot /Subtype /Link /Border [0 0 0] /H /I /C [1 0 0] /Rect [106.556 284.564 139.398 296.517] /A << /S /GoTo /D <74686D2E442E332E312E32> >> >>}
\MusicPDFWriteObject{3813}{<< /Type /Annot /Subtype /Link /Border [0 0 0] /H /I /C [1 0 0] /Rect [407.77 207.6 441.136 219.556] /A << /S /GoTo /D <74686D2E442E312E322E31> >> >>}
\MusicPDFWriteObject{3814}{<< /Type /Annot /Subtype /Link /Border [0 0 0] /H /I /C [1 0 0] /Rect [305.741 181.299 338.582 193.252] /A << /S /GoTo /D <74686D2E442E332E322E35> >> >>}
\MusicPDFWriteObject{3815}{<< /Type /Annot /Subtype /Link /Border [0 0 0] /H /I /C [1 0 0] /Rect [433.998 154.107 471.82 166.062] /A << /S /GoTo /D <74686D2E442E332E332E3133> >> >>}
\MusicPDFWriteObject{3816}{<< /Type /Annot /Subtype /Link /Border [0 0 0] /H /I /C [1 0 0] /Rect [108.679 121.794 141.675 134.64] /A << /S /GoTo /D <74686D2E442E322E312E34> >> >>}
\MusicPDFWriteObject{3817}{<< /Type /Annot /Subtype /Link /Border [0 0 0] /H /I /C [0 1 0] /Rect [165.205 124.538 198.766 133.406] /A << /S /GoTo /D <636974652E5765696C2D65> >> >>}
\MusicPDFWriteObject{3818}{<< /Type /Annot /Subtype /Link /Border [0 0 0] /H /I /C [0 1 0] /Rect [305.166 98.237 322.639 107.105] /A << /S /GoTo /D <636974652E4346> >> >>}
\MusicPDFWriteObject{3819}{<< /Type /Annot /Subtype /Link /Border [0 0 0] /H /I /C [1 0 0] /Rect [73.712 82.596 103.925 94.549] /A << /S /GoTo /D <6571756174696F6E2E352E312E312E31> >> >>}
\MusicPDFWriteObject{3820}{<< /Type /Annot /Subtype /Link /Border [0 0 0] /H /I /C [1 0 0] /Rect [242.545 649.125 275.387 659.916] /A << /S /GoTo /D <74686D2E442E332E322E35> >> >>}
\MusicPDFWriteObject{3821}{<< /Type /Annot /Subtype /Link /Border [0 0 0] /H /I /C [1 0 0] /Rect [98.507 593.009 128.72 604.962] /A << /S /GoTo /D <6571756174696F6E2E392E322E302E31> >> >>}
\MusicPDFWriteObject{3822}{<< /Type /Annot /Subtype /Link /Border [0 0 0] /H /I /C [1 0 0] /Rect [406.228 593.009 436.554 604.962] /A << /S /GoTo /D <74686D2E392E312E312E31> >> >>}
\MusicPDFWriteObject{3823}{<< /Type /Annot /Subtype /Link /Border [0 0 0] /H /I /C [1 0 0] /Rect [196.287 328.264 218.733 342.815] /A << /S /GoTo /D <73756273656374696F6E2E392E322E32> >> >>}
\MusicPDFWriteObject{3824}{<< /Type /Annot /Subtype /Link /Border [0 0 0] /H /I /C [1 0 0] /Rect [442.437 588.109 472.65 600.061] /A << /S /GoTo /D <6571756174696F6E2E392E322E312E36> >> >>}
\MusicPDFWriteObject{3825}{<< /Type /Annot /Subtype /Link /Border [0 0 0] /H /I /C [1 0 0] /Rect [106.556 363.479 136.769 372.333] /A << /S /GoTo /D <74686D2E392E312E322E32> >> >>}
\MusicPDFWriteObject{3826}{<< /Type /Annot /Subtype /Link /Border [0 0 0] /H /I /C [0 1 0] /Rect [91.862 229.744 109.016 238.612] /A << /S /GoTo /D <636974652E466F6C> >> >>}
\MusicPDFWriteObject{3827}{<< /Type /Annot /Subtype /Link /Border [0 0 0] /H /I /C [0 1 0] /Rect [311.595 229.744 329.897 238.612] /A << /S /GoTo /D <636974652E4445> >> >>}
\MusicPDFWriteObject{3828}{<< /Type /Annot /Subtype /Link /Border [0 0 0] /H /I /C [0 1 0] /Rect [404.183 229.744 421.338 238.612] /A << /S /GoTo /D <636974652E466F6C> >> >>}
\MusicPDFWriteObject{3829}{<< /Type /Annot /Subtype /Link /Border [0 0 0] /H /I /C [0 1 0] /Rect [182.807 203.443 199.961 212.311] /A << /S /GoTo /D <636974652E466F6C> >> >>}
\MusicPDFWriteObject{3830}{<< /Type /Annot /Subtype /Link /Border [0 0 0] /H /I /C [0 1 0] /Rect [350.03 203.443 367.184 212.311] /A << /S /GoTo /D <636974652E466F6C> >> >>}
\MusicPDFWriteObject{3831}{<< /Type /Annot /Subtype /Link /Border [0 0 0] /H /I /C [0 1 0] /Rect [87.166 163.99 121.135 172.858] /A << /S /GoTo /D <636974652E436E772D61> >> >>}
\MusicPDFWriteObject{3832}{<< /Type /Annot /Subtype /Link /Border [0 0 0] /H /I /C [0 1 0] /Rect [332.296 111.387 349.45 120.255] /A << /S /GoTo /D <636974652E466F6C> >> >>}
\MusicPDFWriteObject{3833}{<< /Type /Annot /Subtype /Link /Border [0 0 0] /H /I /C [0 1 0] /Rect [426.963 111.387 444.118 120.255] /A << /S /GoTo /D <636974652E466F6C> >> >>}
\MusicPDFWriteObject{3834}{<< /Type /Annot /Subtype /Link /Border [0 0 0] /H /I /C [1 0 0] /Rect [189.586 82.596 219.799 94.549] /A << /S /GoTo /D <6571756174696F6E2E392E322E322E32> >> >>}
\MusicPDFWriteObject{3835}{<< /Type /Annot /Subtype /Link /Border [0 0 0] /H /I /C [1 0 0] /Rect [86.497 542.142 116.71 554.095] /A << /S /GoTo /D <6571756174696F6E2E392E322E312E36> >> >>}
\MusicPDFWriteObject{3836}{<< /Type /Annot /Subtype /Link /Border [0 0 0] /H /I /C [0 1 0] /Rect [317.472 399.517 339.33 408.385] /A << /S /GoTo /D <636974652E447572> >> >>}
\MusicPDFWriteObject{3837}{<< /Type /Annot /Subtype /Link /Border [0 0 0] /H /I /C [0 1 0] /Rect [168.067 307.3 189.925 316.168] /A << /S /GoTo /D <636974652E447572> >> >>}
\MusicPDFWriteObject{3838}{<< /Type /Annot /Subtype /Link /Border [0 0 0] /H /I /C [0 1 0] /Rect [115.824 267.848 137.682 276.716] /A << /S /GoTo /D <636974652E447572> >> >>}
\MusicPDFWriteObject{3839}{<< /Type /Annot /Subtype /Link /Border [0 0 0] /H /I /C [0 1 0] /Rect [523.216 240.605 543.159 249.473] /A << /S /GoTo /D <636974652E596F73> >> >>}
\MusicPDFWriteObject{3840}{<< /Type /Annot /Subtype /Link /Border [0 0 0] /H /I /C [1 0 0] /Rect [170.894 185.109 202.576 197.061] /A << /S /GoTo /D <74686D2E422E332E322E31> >> >>}
\MusicPDFWriteObject{3841}{<< /Type /Annot /Subtype /Link /Border [0 0 0] /H /I /C [1 0 0] /Rect [422.777 185.109 452.99 197.061] /A << /S /GoTo /D <6571756174696F6E2E392E322E312E36> >> >>}
\MusicPDFWriteObject{3842}{<< /Type /Annot /Subtype /Link /Border [0 0 0] /H /I /C [1 0 0] /Rect [492.59 108.786 524.109 123.337] /A << /S /GoTo /D <74686D2E392E322E312E32> >> >>}
\MusicPDFWriteObject{3843}{<< /Type /Annot /Subtype /Link /Border [0 0 0] /H /I /C [0 1 0] /Rect [524.934 554.737 543.236 563.605] /A << /S /GoTo /D <636974652E4445> >> >>}
\MusicPDFWriteObject{3844}{<< /Type /Annot /Subtype /Link /Border [0 0 0] /H /I /C [0 1 0] /Rect [143.133 541.586 160.288 550.454] /A << /S /GoTo /D <636974652E466F6C> >> >>}
\MusicPDFWriteObject{3845}{<< /Type /Annot /Subtype /Link /Border [0 0 0] /H /I /C [1 0 0] /Rect [230.233 490.239 260.446 502.192] /A << /S /GoTo /D <74686D2E392E322E312E32> >> >>}
\MusicPDFWriteObject{3846}{<< /Type /Annot /Subtype /Link /Border [0 0 0] /H /I /C [0 1 0] /Rect [104.072 446.396 121.226 455.264] /A << /S /GoTo /D <636974652E466F6C> >> >>}
\MusicPDFWriteObject{3847}{<< /Type /Annot /Subtype /Link /Border [0 0 0] /H /I /C [1 0 0] /Rect [147.951 368.528 178.729 381.535] /A << /S /GoTo /D <74686D2E392E322E312E32> >> >>}
\MusicPDFWriteObject{3848}{<< /Type /Annot /Subtype /Link /Border [0 0 0] /H /I /C [1 0 0] /Rect [328.942 170.088 364.135 182.04] /A << /S /GoTo /D <6571756174696F6E2E392E322E312E3136> >> >>}
\MusicPDFWriteObject{3849}{<< /Type /Annot /Subtype /Link /Border [0 0 0] /H /I /C [1 0 0] /Rect [235.596 132.277 265.809 143.068] /A << /S /GoTo /D <74686D2E392E322E312E32> >> >>}
\MusicPDFWriteObject{3850}{<< /Type /Annot /Subtype /Link /Border [0 0 0] /H /I /C [1 0 0] /Rect [330.55 132.277 360.763 143.068] /A << /S /GoTo /D <74686D2E392E322E312E34> >> >>}
\MusicPDFWriteObject{3851}{<< /Type /Annot /Subtype /Link /Border [0 0 0] /H /I /C [1 0 0] /Rect [438.337 232.692 473.53 244.645] /A << /S /GoTo /D <6571756174696F6E2E392E322E312E3232> >> >>}
\MusicPDFWriteObject{3852}{<< /Type /Annot /Subtype /Link /Border [0 0 0] /H /I /C [1 0 0] /Rect [226.451 173.066 259.909 187.617] /A << /S /GoTo /D <74686D2E442E312E322E31> >> >>}
\MusicPDFWriteObject{3853}{<< /Type /Annot /Subtype /Link /Border [0 0 0] /H /I /C [1 0 0] /Rect [181.454 147.28 216.647 159.233] /A << /S /GoTo /D <6571756174696F6E2E392E322E322E3436> >> >>}
\MusicPDFWriteObject{3854}{<< /Type /Annot /Subtype /Link /Border [0 0 0] /H /I /C [1 0 0] /Rect [496.27 147.28 529.512 159.233] /A << /S /GoTo /D <74686D2E442E312E322E31> >> >>}
\MusicPDFWriteObject{3855}{<< /Type /Annot /Subtype /Link /Border [0 0 0] /H /I /C [1 0 0] /Rect [301.782 121.532 331.995 132.323] /A << /S /GoTo /D <74686D2E392E322E312E32> >> >>}
\MusicPDFWriteObject{3856}{<< /Type /Annot /Subtype /Link /Border [0 0 0] /H /I /C [1 0 0] /Rect [293.84 102.403 324.053 113.194] /A << /S /GoTo /D <74686D2E392E322E312E35> >> >>}
\MusicPDFWriteObject{3857}{<< /Type /Annot /Subtype /Link /Border [0 0 0] /H /I /C [1 0 0] /Rect [388.794 102.403 419.007 113.194] /A << /S /GoTo /D <74686D2E392E322E312E33> >> >>}
\MusicPDFWriteObject{3858}{<< /Type /Annot /Subtype /Link /Border [0 0 0] /H /I /C [1 0 0] /Rect [70.21 671.482 101.729 683.437] /A << /S /GoTo /D <74686D2E392E322E312E35> >> >>}
\MusicPDFWriteObject{3859}{<< /Type /Annot /Subtype /Link /Border [0 0 0] /H /I /C [1 0 0] /Rect [262.244 504.073 292.457 514.864] /A << /S /GoTo /D <74686D2E392E322E312E34> >> >>}
\MusicPDFWriteObject{3860}{<< /Type /Annot /Subtype /Link /Border [0 0 0] /H /I /C [1 0 0] /Rect [278.534 337.16 308.747 347.951] /A << /S /GoTo /D <74686D2E392E322E312E35> >> >>}
\MusicPDFWriteObject{3861}{<< /Type /Annot /Subtype /Link /Border [0 0 0] /H /I /C [1 0 0] /Rect [373.488 337.16 403.701 347.951] /A << /S /GoTo /D <74686D2E392E322E312E34> >> >>}
\MusicPDFWriteObject{3862}{<< /Type /Annot /Subtype /Link /Border [0 0 0] /H /I /C [1 0 0] /Rect [139.266 611.318 174.459 623.271] /A << /S /GoTo /D <6571756174696F6E2E392E322E312E3238> >> >>}
\MusicPDFWriteObject{3863}{<< /Type /Annot /Subtype /Link /Border [0 0 0] /H /I /C [0 1 0] /Rect [146.588 587.948 163.742 596.816] /A << /S /GoTo /D <636974652E466F6C> >> >>}
\MusicPDFWriteObject{3864}{<< /Type /Annot /Subtype /Link /Border [0 0 0] /H /I /C [0 1 0] /Rect [171.74 478.031 188.894 486.898] /A << /S /GoTo /D <636974652E466F6C> >> >>}
\MusicPDFWriteObject{3865}{<< /Type /Annot /Subtype /Link /Border [0 0 0] /H /I /C [1 0 0] /Rect [376.2 327.595 411.393 339.548] /A << /S /GoTo /D <6571756174696F6E2E392E322E312E3237> >> >>}
\MusicPDFWriteObject{3866}{<< /Type /Annot /Subtype /Link /Border [0 0 0] /H /I /C [1 0 0] /Rect [392.102 224.274 422.682 236.483] /A << /S /GoTo /D <74686D2E392E322E322E32> >> >>}
\MusicPDFWriteObject{3867}{<< /Type /Annot /Subtype /Link /Border [0 0 0] /H /I /C [0 1 0] /Rect [221.396 200.717 270.436 209.585] /A << /S /GoTo /D <636974652E426F752D545653> >> >>}
\MusicPDFWriteObject{3868}{<< /Type /Annot /Subtype /Link /Border [0 0 0] /H /I /C [1 0 0] /Rect [148.168 185.629 185.574 196.42] /A << /S /GoTo /D <74686D2E432E322E332E3131> >> >>}
\MusicPDFWriteObject{3869}{<< /Type /Annot /Subtype /Link /Border [0 0 0] /H /I /C [0 1 0] /Rect [322.527 148.114 339.681 156.982] /A << /S /GoTo /D <636974652E466F6C> >> >>}
\MusicPDFWriteObject{3870}{<< /Type /Annot /Subtype /Link /Border [0 0 0] /H /I /C [0 1 0] /Rect [229.343 594.059 246.498 602.927] /A << /S /GoTo /D <636974652E466F6C> >> >>}
\MusicPDFWriteObject{3871}{<< /Type /Annot /Subtype /Link /Border [0 0 0] /H /I /C [1 0 0] /Rect [165.826 578.418 201.019 590.371] /A << /S /GoTo /D <6571756174696F6E2E392E322E312E3336> >> >>}
\MusicPDFWriteObject{3872}{<< /Type /Annot /Subtype /Link /Border [0 0 0] /H /I /C [1 0 0] /Rect [413.403 578.418 448.597 590.371] /A << /S /GoTo /D <6571756174696F6E2E392E322E312E3331> >> >>}
\MusicPDFWriteObject{3873}{<< /Type /Annot /Subtype /Link /Border [0 0 0] /H /I /C [1 0 0] /Rect [155.499 253.556 185.713 265.508] /A << /S /GoTo /D <6571756174696F6E2E392E322E322E31> >> >>}
\MusicPDFWriteObject{3874}{<< /Type /Annot /Subtype /Link /Border [0 0 0] /H /I /C [0 1 0] /Rect [107.639 242.895 129.496 251.763] /A << /S /GoTo /D <636974652E447572> >> >>}
\MusicPDFWriteObject{3875}{<< /Type /Annot /Subtype /Link /Border [0 0 0] /H /I /C [0 1 0] /Rect [241.861 242.895 263.718 251.763] /A << /S /GoTo /D <636974652E447572> >> >>}
\MusicPDFWriteObject{3876}{<< /Type /Annot /Subtype /Link /Border [0 0 0] /H /I /C [1 0 0] /Rect [169.756 227.254 201.438 240.118] /A << /S /GoTo /D <74686D2E422E322E322E34> >> >>}
\MusicPDFWriteObject{3877}{<< /Type /Annot /Subtype /Link /Border [0 0 0] /H /I /C [1 0 0] /Rect [366.668 214.103 396.882 226.056] /A << /S /GoTo /D <6571756174696F6E2E392E322E322E32> >> >>}
\MusicPDFWriteObject{3878}{<< /Type /Annot /Subtype /Link /Border [0 0 0] /H /I /C [0 1 0] /Rect [73.77 177.141 95.628 186.009] /A << /S /GoTo /D <636974652E447572> >> >>}
\MusicPDFWriteObject{3879}{<< /Type /Annot /Subtype /Link /Border [0 0 0] /H /I /C [0 1 0] /Rect [322.946 150.84 342.051 159.708] /A << /S /GoTo /D <636974652E4E42> >> >>}
\MusicPDFWriteObject{3880}{<< /Type /Annot /Subtype /Link /Border [0 0 0] /H /I /C [1 0 0] /Rect [185.11 108.897 215.323 120.85] /A << /S /GoTo /D <6571756174696F6E2E392E322E322E31> >> >>}
\MusicPDFWriteObject{3881}{<< /Type /Annot /Subtype /Link /Border [0 0 0] /H /I /C [1 0 0] /Rect [432.881 684.633 466.308 696.588] /A << /S /GoTo /D <74686D2E442E332E312E32> >> >>}
\MusicPDFWriteObject{3882}{<< /Type /Annot /Subtype /Link /Border [0 0 0] /H /I /C [1 0 0] /Rect [175.807 658.332 208.095 670.284] /A << /S /GoTo /D <74686D2E422E322E322E34> >> >>}
\MusicPDFWriteObject{3883}{<< /Type /Annot /Subtype /Link /Border [0 0 0] /H /I /C [1 0 0] /Rect [193.135 478.209 228.328 490.162] /A << /S /GoTo /D <6571756174696F6E2E392E322E322E3130> >> >>}
\MusicPDFWriteObject{3884}{<< /Type /Annot /Subtype /Link /Border [0 0 0] /H /I /C [1 0 0] /Rect [158.959 282.25 194.152 294.203] /A << /S /GoTo /D <6571756174696F6E2E392E322E322E3134> >> >>}
\MusicPDFWriteObject{3885}{<< /Type /Annot /Subtype /Link /Border [0 0 0] /H /I /C [1 0 0] /Rect [222.594 282.25 257.787 294.203] /A << /S /GoTo /D <6571756174696F6E2E392E322E322E3135> >> >>}
\MusicPDFWriteObject{3886}{<< /Type /Annot /Subtype /Link /Border [0 0 0] /H /I /C [1 0 0] /Rect [192.187 245.6 222.4 257.553] /A << /S /GoTo /D <6571756174696F6E2E392E322E322E34> >> >>}
\MusicPDFWriteObject{3887}{<< /Type /Annot /Subtype /Link /Border [0 0 0] /H /I /C [0 1 0] /Rect [274.694 234.939 291.849 243.807] /A << /S /GoTo /D <636974652E466F6C> >> >>}
\MusicPDFWriteObject{3888}{<< /Type /Annot /Subtype /Link /Border [0 0 0] /H /I /C [1 0 0] /Rect [124.787 182.578 155.565 194.534] /A << /S /GoTo /D <74686D2E392E312E322E32> >> >>}
\MusicPDFWriteObject{3889}{<< /Type /Annot /Subtype /Link /Border [0 0 0] /H /I /C [1 0 0] /Rect [130.322 95.747 165.516 107.699] /A << /S /GoTo /D <6571756174696F6E2E392E322E322E3136> >> >>}
\MusicPDFWriteObject{3890}{<< /Type /Annot /Subtype /Link /Border [0 0 0] /H /I /C [1 0 0] /Rect [207.214 685.187 237.427 695.977] /A << /S /GoTo /D <74686D2E392E322E322E31> >> >>}
\MusicPDFWriteObject{3891}{<< /Type /Annot /Subtype /Link /Border [0 0 0] /H /I /C [1 0 0] /Rect [346.242 685.187 376.455 695.977] /A << /S /GoTo /D <74686D2E392E322E322E32> >> >>}
\MusicPDFWriteObject{3892}{<< /Type /Annot /Subtype /Link /Border [0 0 0] /H /I /C [1 0 0] /Rect [291.096 573.833 326.289 585.786] /A << /S /GoTo /D <6571756174696F6E2E392E322E322E3139> >> >>}
\MusicPDFWriteObject{3893}{<< /Type /Annot /Subtype /Link /Border [0 0 0] /H /I /C [1 0 0] /Rect [454.713 510.491 489.906 522.444] /A << /S /GoTo /D <6571756174696F6E2E392E322E322E3139> >> >>}
\MusicPDFWriteObject{3894}{<< /Type /Annot /Subtype /Link /Border [0 0 0] /H /I /C [1 0 0] /Rect [116.894 379.447 149.736 391.4] /A << /S /GoTo /D <6571756174696F6E2E442E332E322E32> >> >>}
\MusicPDFWriteObject{3895}{<< /Type /Annot /Subtype /Link /Border [0 0 0] /H /I /C [1 0 0] /Rect [180.509 340.907 210.157 351.697] /A << /S /GoTo /D <74686D2E392E322E322E34> >> >>}
\MusicPDFWriteObject{3896}{<< /Type /Annot /Subtype /Link /Border [0 0 0] /H /I /C [1 0 0] /Rect [300.446 340.907 330.095 351.697] /A << /S /GoTo /D <74686D2E392E322E322E32> >> >>}
\MusicPDFWriteObject{3897}{<< /Type /Annot /Subtype /Link /Border [0 0 0] /H /I /C [1 0 0] /Rect [325.242 327.756 355.398 338.546] /A << /S /GoTo /D <74686D2E392E322E322E32> >> >>}
\MusicPDFWriteObject{3898}{<< /Type /Annot /Subtype /Link /Border [0 0 0] /H /I /C [1 0 0] /Rect [172.238 314.052 204.526 326.909] /A << /S /GoTo /D <74686D2E422E322E322E34> >> >>}
\MusicPDFWriteObject{3899}{<< /Type /Annot /Subtype /Link /Border [0 0 0] /H /I /C [1 0 0] /Rect [209.174 314.052 215.137 326.909] /A << /S /GoTo /D <48666F6F746E6F74652E3236> >> >>}
\MusicPDFWriteObject{3900}{<< /Type /Annot /Subtype /Link /Border [0 0 0] /H /I /C [1 0 0] /Rect [233.928 130.57 264.706 141.361] /A << /S /GoTo /D <74686D2E392E322E322E35> >> >>}
\MusicPDFWriteObject{3901}{<< /Type /Annot /Subtype /Link /Border [0 0 0] /H /I /C [1 0 0] /Rect [428.492 130.57 459.269 141.361] /A << /S /GoTo /D <74686D2E392E322E312E34> >> >>}
\MusicPDFWriteObject{3902}{<< /Type /Annot /Subtype /Link /Border [0 0 0] /H /I /C [1 0 0] /Rect [71.004 119.356 101.217 129.723] /A << /S /GoTo /D <74686D2E392E322E322E34> >> >>}
\MusicPDFWriteObject{3903}{<< /Type /Annot /Subtype /Link /Border [0 0 0] /H /I /C [1 0 0] /Rect [101.991 119.356 107.954 129.723] /A << /S /GoTo /D <48666F6F746E6F74652E3237> >> >>}
\MusicPDFWriteObject{3904}{<< /Type /Annot /Subtype /Link /Border [0 0 0] /H /I /C [1 0 0] /Rect [449.59 399.104 484.783 411.056] /A << /S /GoTo /D <6571756174696F6E2E392E322E322E3237> >> >>}
\MusicPDFWriteObject{3905}{<< /Type /Annot /Subtype /Link /Border [0 0 0] /H /I /C [1 0 0] /Rect [252.991 255.836 283.769 267.791] /A << /S /GoTo /D <74686D2E392E322E322E36> >> >>}
\MusicPDFWriteObject{3906}{<< /Type /Annot /Subtype /Link /Border [0 0 0] /H /I /C [1 0 0] /Rect [187.794 223.97 221.355 232.824] /A << /S /GoTo /D <74686D2E392E322E322E36> >> >>}
\MusicPDFWriteObject{3907}{<< /Type /Annot /Subtype /Link /Border [0 0 0] /H /I /C [1 0 0] /Rect [326.16 222.034 356.063 232.824] /A << /S /GoTo /D <74686D2E392E322E322E36> >> >>}
\MusicPDFWriteObject{3908}{<< /Type /Annot /Subtype /Link /Border [0 0 0] /H /I /C [1 0 0] /Rect [71.004 210.82 101.217 221.187] /A << /S /GoTo /D <74686D2E392E322E312E32> >> >>}
\MusicPDFWriteObject{3909}{<< /Type /Annot /Subtype /Link /Border [0 0 0] /H /I /C [1 0 0] /Rect [101.991 210.82 107.954 221.187] /A << /S /GoTo /D <48666F6F746E6F74652E3238> >> >>}
\MusicPDFWriteObject{3910}{<< /Type /Annot /Subtype /Link /Border [0 0 0] /H /I /C [1 0 0] /Rect [280.498 603.536 310.712 615.84] /A << /S /GoTo /D <74686D2E392E322E322E35> >> >>}
\MusicPDFWriteObject{3911}{<< /Type /Annot /Subtype /Link /Border [0 0 0] /H /I /C [1 0 0] /Rect [311.486 603.536 317.449 615.84] /A << /S /GoTo /D <48666F6F746E6F74652E3239> >> >>}
\MusicPDFWriteObject{3912}{<< /Type /Annot /Subtype /Link /Border [0 0 0] /H /I /C [1 0 0] /Rect [426.183 258.078 456.396 268.954] /A << /S /GoTo /D <74686D2E392E322E322E32> >> >>}
\MusicPDFWriteObject{3913}{<< /Type /Annot /Subtype /Link /Border [0 0 0] /H /I /C [1 0 0] /Rect [399.9 421.685 429.803 433.637] /A << /S /GoTo /D <74686D2E392E322E322E32> >> >>}
\MusicPDFWriteObject{3914}{<< /Type /Annot /Subtype /Link /Border [0 0 0] /H /I /C [1 0 0] /Rect [308.078 408.174 340.92 421.391] /A << /S /GoTo /D <74686D2E442E312E322E31> >> >>}
\MusicPDFWriteObject{3915}{<< /Type /Annot /Subtype /Link /Border [0 0 0] /H /I /C [1 0 0] /Rect [341.694 408.174 347.657 421.391] /A << /S /GoTo /D <48666F6F746E6F74652E3330> >> >>}
\MusicPDFWriteObject{3916}{<< /Type /Annot /Subtype /Link /Border [0 0 0] /H /I /C [1 0 0] /Rect [139.266 109.528 174.459 121.48] /A << /S /GoTo /D <6571756174696F6E2E392E322E322E3339> >> >>}
\MusicPDFWriteObject{3917}{<< /Type /Annot /Subtype /Link /Border [0 0 0] /H /I /C [1 0 0] /Rect [424.938 685.187 454.897 695.977] /A << /S /GoTo /D <74686D2E392E322E312E38> >> >>}
\MusicPDFWriteObject{3918}{<< /Type /Annot /Subtype /Link /Border [0 0 0] /H /I /C [1 0 0] /Rect [71.004 671.482 106.197 683.435] /A << /S /GoTo /D <74686D2E392E322E322E3130> >> >>}
\MusicPDFWriteObject{3919}{<< /Type /Annot /Subtype /Link /Border [0 0 0] /H /I /C [1 0 0] /Rect [284.151 467.214 319.344 478.005] /A << /S /GoTo /D <74686D2E392E322E322E3130> >> >>}
\MusicPDFWriteObject{3920}{<< /Type /Annot /Subtype /Link /Border [0 0 0] /H /I /C [1 0 0] /Rect [431.949 467.214 462.162 478.005] /A << /S /GoTo /D <74686D2E392E322E322E35> >> >>}
\MusicPDFWriteObject{3921}{<< /Type /Annot /Subtype /Link /Border [0 0 0] /H /I /C [1 0 0] /Rect [449.59 684.633 484.783 696.586] /A << /S /GoTo /D <6571756174696F6E2E392E322E322E3437> >> >>}
\MusicPDFWriteObject{3922}{<< /Type /Annot /Subtype /Link /Border [0 0 0] /H /I /C [1 0 0] /Rect [283.986 637.413 319.444 649.365] /A << /S /GoTo /D <74686D2E392E322E322E3130> >> >>}
\MusicPDFWriteObject{3923}{<< /Type /Annot /Subtype /Link /Border [0 0 0] /H /I /C [1 0 0] /Rect [71.004 626.752 101.217 635.469] /A << /S /GoTo /D <74686D2E392E322E322E36> >> >>}
\MusicPDFWriteObject{3924}{<< /Type /Annot /Subtype /Link /Border [0 0 0] /H /I /C [0 1 0] /Rect [305.745 492.004 326.967 500.872] /A << /S /GoTo /D <636974652E427275> >> >>}
\MusicPDFWriteObject{3925}{<< /Type /Annot /Subtype /Link /Border [0 0 0] /H /I /C [0 1 0] /Rect [249.571 478.854 270.793 487.722] /A << /S /GoTo /D <636974652E427275> >> >>}
\MusicPDFWriteObject{3926}{<< /Type /Annot /Subtype /Link /Border [0 0 0] /H /I /C [1 0 0] /Rect [355.4 476.364 365.001 488.316] /A << /S /GoTo /D <617070656E6469782E44> >> >>}
\MusicPDFWriteObject{3927}{<< /Type /Annot /Subtype /Link /Border [0 0 0] /H /I /C [1 0 0] /Rect [355.055 156.118 388.09 168.071] /A << /S /GoTo /D <74686D2E432E342E332E35> >> >>}
\MusicPDFWriteObject{3928}{<< /Type /Annot /Subtype /Link /Border [0 0 0] /H /I /C [1 0 0] /Rect [71.004 142.967 103.845 154.922] /A << /S /GoTo /D <74686D2E442E332E332E37> >> >>}
\MusicPDFWriteObject{3929}{<< /Type /Annot /Subtype /Link /Border [0 0 0] /H /I /C [1 0 0] /Rect [366.573 108.897 396.786 120.85] /A << /S /GoTo /D <6571756174696F6E2E392E332E312E31> >> >>}
\MusicPDFWriteObject{3930}{<< /Type /Annot /Subtype /Link /Border [0 0 0] /H /I /C [1 0 0] /Rect [180.993 503.159 218.313 515.115] /A << /S /GoTo /D <74686D2E442E332E332E3133> >> >>}
\MusicPDFWriteObject{3931}{<< /Type /Annot /Subtype /Link /Border [0 0 0] /H /I /C [1 0 0] /Rect [140.121 449.815 177.527 461.77] /A << /S /GoTo /D <74686D2E432E342E332E3132> >> >>}
\MusicPDFWriteObject{3932}{<< /Type /Annot /Subtype /Link /Border [0 0 0] /H /I /C [1 0 0] /Rect [421.406 156.022 451.28 166.812] /A << /S /GoTo /D <74686D2E392E332E312E33> >> >>}
\MusicPDFWriteObject{3933}{<< /Type /Annot /Subtype /Link /Border [0 0 0] /H /I /C [1 0 0] /Rect [280.632 142.317 310.309 154.27] /A << /S /GoTo /D <74686D2E352E342E312E31> >> >>}
\MusicPDFWriteObject{3934}{<< /Type /Annot /Subtype /Link /Border [0 0 0] /H /I /C [1 0 0] /Rect [366.011 339.501 396.224 351.454] /A << /S /GoTo /D <6571756174696F6E2E392E312E312E34> >> >>}
\MusicPDFWriteObject{3935}{<< /Type /Annot /Subtype /Link /Border [0 0 0] /H /I /C [1 0 0] /Rect [477.837 339.501 507.74 351.473] /A << /S /GoTo /D <74686D2E392E312E312E33> >> >>}
\MusicPDFWriteObject{3936}{<< /Type /Annot /Subtype /Link /Border [0 0 0] /H /I /C [1 0 0] /Rect [108.155 313.199 138.368 325.152] /A << /S /GoTo /D <6571756174696F6E2E392E312E312E34> >> >>}
\MusicPDFWriteObject{3937}{<< /Type /Annot /Subtype /Link /Border [0 0 0] /H /I /C [1 0 0] /Rect [218.791 313.199 248.44 325.152] /A << /S /GoTo /D <74686D2E392E312E312E33> >> >>}
\MusicPDFWriteObject{3938}{<< /Type /Annot /Subtype /Link /Border [0 0 0] /H /I /C [1 0 0] /Rect [183.193 257.019 213.406 268.972] /A << /S /GoTo /D <6571756174696F6E2E392E312E312E34> >> >>}
\MusicPDFWriteObject{3939}{<< /Type /Annot /Subtype /Link /Border [0 0 0] /H /I /C [1 0 0] /Rect [295.771 257.019 325.984 268.972] /A << /S /GoTo /D <74686D2E392E312E312E33> >> >>}
\MusicPDFWriteObject{3940}{<< /Type /Annot /Subtype /Link /Border [0 0 0] /H /I /C [0 1 0] /Rect [149.385 85.086 175.877 93.954] /A << /S /GoTo /D <636974652E52752D32> >> >>}
\MusicPDFWriteObject{3941}{<< /Type /Annot /Subtype /Link /Border [0 0 0] /H /I /C [0 1 0] /Rect [327.488 647.671 353.98 656.539] /A << /S /GoTo /D <636974652E52752D32> >> >>}
\MusicPDFWriteObject{3942}{<< /Type /Annot /Subtype /Link /Border [0 0 0] /H /I /C [1 0 0] /Rect [151.986 582.55 187.179 594.503] /A << /S /GoTo /D <74686D2E392E322E322E3130> >> >>}
\MusicPDFWriteObject{3943}{<< /Type /Annot /Subtype /Link /Border [0 0 0] /H /I /C [0 1 0] /Rect [394.092 500.158 410.084 509.026] /A << /S /GoTo /D <636974652E436F> >> >>}
\MusicPDFWriteObject{3944}{<< /Type /Annot /Subtype /Link /Border [0 0 0] /H /I /C [1 0 0] /Rect [110.222 451.766 140.435 463.718] /A << /S /GoTo /D <6571756174696F6E2E392E332E312E31> >> >>}
\MusicPDFWriteObject{3945}{<< /Type /Annot /Subtype /Link /Border [0 0 0] /H /I /C [1 0 0] /Rect [364.716 239.848 396.625 254.265] /A << /S /GoTo /D <74686D2E432E322E312E38> >> >>}
\MusicPDFWriteObject{3946}{<< /Type /Annot /Subtype /Link /Border [0 0 0] /H /I /C [1 0 0] /Rect [509.596 242.311 539.809 254.264] /A << /S /GoTo /D <6571756174696F6E2E392E332E312E31> >> >>}
\MusicPDFWriteObject{3947}{<< /Type /Annot /Subtype /Link /Border [0 0 0] /H /I /C [0 1 0] /Rect [93.24 185.748 118.3 194.616] /A << /S /GoTo /D <636974652E53636861> >> >>}
\MusicPDFWriteObject{3948}{<< /Type /Annot /Subtype /Link /Border [0 0 0] /H /I /C [1 0 0] /Rect [235.236 658.061 267.115 670.654] /A << /S /GoTo /D <74686D2E432E342E322E31> >> >>}
\MusicPDFWriteObject{3949}{<< /Type /Annot /Subtype /Link /Border [0 0 0] /H /I /C [1 0 0] /Rect [204.104 618.609 233.894 630.564] /A << /S /GoTo /D <74686D2E392E332E312E35> >> >>}
\MusicPDFWriteObject{3950}{<< /Type /Annot /Subtype /Link /Border [0 0 0] /H /I /C [1 0 0] /Rect [150.661 579.157 188.068 591.749] /A << /S /GoTo /D <74686D2E432E342E332E3131> >> >>}
\MusicPDFWriteObject{3951}{<< /Type /Annot /Subtype /Link /Border [0 0 0] /H /I /C [1 0 0] /Rect [335.513 566.007 365.726 577.959] /A << /S /GoTo /D <6571756174696F6E2E392E332E322E32> >> >>}
\MusicPDFWriteObject{3952}{<< /Type /Annot /Subtype /Link /Border [0 0 0] /H /I /C [1 0 0] /Rect [254.707 539.705 284.92 551.658] /A << /S /GoTo /D <6571756174696F6E2E392E332E322E32> >> >>}
\MusicPDFWriteObject{3953}{<< /Type /Annot /Subtype /Link /Border [0 0 0] /H /I /C [1 0 0] /Rect [511.675 539.705 541.888 551.658] /A << /S /GoTo /D <6571756174696F6E2E392E332E312E31> >> >>}
\MusicPDFWriteObject{3954}{<< /Type /Annot /Subtype /Link /Border [0 0 0] /H /I /C [1 0 0] /Rect [430.089 526.554 460.302 538.507] /A << /S /GoTo /D <6571756174696F6E2E392E332E312E31> >> >>}
\MusicPDFWriteObject{3955}{<< /Type /Annot /Subtype /Link /Border [0 0 0] /H /I /C [1 0 0] /Rect [255.744 487.101 288.566 499.694] /A << /S /GoTo /D <74686D2E432E342E332E39> >> >>}
\MusicPDFWriteObject{3956}{<< /Type /Annot /Subtype /Link /Border [0 0 0] /H /I /C [1 0 0] /Rect [261.695 415.528 293.604 427.483] /A << /S /GoTo /D <74686D2E432E342E332E39> >> >>}
\MusicPDFWriteObject{3957}{<< /Type /Annot /Subtype /Link /Border [0 0 0] /H /I /C [0 1 0] /Rect [152.901 365.415 183.848 374.283] /A << /S /GoTo /D <636974652E486F726D> >> >>}
\MusicPDFWriteObject{3958}{<< /Type /Annot /Subtype /Link /Border [0 0 0] /H /I /C [1 0 0] /Rect [278.985 128.499 312.443 140.538] /A << /S /GoTo /D <74686D2E442E312E322E31> >> >>}
\MusicPDFWriteObject{3959}{<< /Type /Annot /Subtype /Link /Border [0 0 0] /H /I /C [1 0 0] /Rect [172.267 462.021 202.48 473.974] /A << /S /GoTo /D <6571756174696F6E2E392E332E322E36> >> >>}
\MusicPDFWriteObject{3960}{<< /Type /Annot /Subtype /Link /Border [0 0 0] /H /I /C [0 1 0] /Rect [237.064 464.511 266.156 473.379] /A << /S /GoTo /D <636974652E48522D31> >> >>}
\MusicPDFWriteObject{3961}{<< /Type /Annot /Subtype /Link /Border [0 0 0] /H /I /C [1 0 0] /Rect [209.881 425.593 240.094 437.546] /A << /S /GoTo /D <6571756174696F6E2E392E332E322E35> >> >>}
\MusicPDFWriteObject{3962}{<< /Type /Annot /Subtype /Link /Border [0 0 0] /H /I /C [1 0 0] /Rect [164.582 285.819 194.231 300.448] /A << /S /GoTo /D <74686D2E342E332E322E32> >> >>}
\MusicPDFWriteObject{3963}{<< /Type /Annot /Subtype /Link /Border [0 0 0] /H /I /C [1 0 0] /Rect [298.903 89.704 329.116 100.495] /A << /S /GoTo /D <74686D2E392E332E322E31> >> >>}
\MusicPDFWriteObject{3964}{<< /Type /Annot /Subtype /Link /Border [0 0 0] /H /I /C [0 1 0] /Rect [214.508 580.048 232.81 588.916] /A << /S /GoTo /D <636974652E4445> >> >>}
\MusicPDFWriteObject{3965}{<< /Type /Annot /Subtype /Link /Border [0 0 0] /H /I /C [0 1 0] /Rect [419.292 527.445 450.238 536.313] /A << /S /GoTo /D <636974652E486F726D> >> >>}
\MusicPDFWriteObject{3966}{<< /Type /Annot /Subtype /Link /Border [0 0 0] /H /I /C [1 0 0] /Rect [345.901 484.761 375.662 497.354] /A << /S /GoTo /D <74686D2E392E332E322E31> >> >>}
\MusicPDFWriteObject{3967}{<< /Type /Annot /Subtype /Link /Border [0 0 0] /H /I /C [0 1 0] /Rect [88.919 234.446 125.641 245.25] /A << /S /GoTo /D <636974652E4C616E672D62> >> >>}
\MusicPDFWriteObject{3968}{<< /Type /Annot /Subtype /Link /Border [0 0 0] /H /I /C [1 0 0] /Rect [159.629 597.258 189.278 609.213] /A << /S /GoTo /D <74686D2E392E332E322E34> >> >>}
\MusicPDFWriteObject{3969}{<< /Type /Annot /Subtype /Link /Border [0 0 0] /H /I /C [1 0 0] /Rect [250.676 597.258 287.781 609.213] /A << /S /GoTo /D <74686D2E442E332E332E3133> >> >>}
\MusicPDFWriteObject{3970}{<< /Type /Annot /Subtype /Link /Border [0 0 0] /H /I /C [1 0 0] /Rect [161.365 584.107 194.206 596.965] /A << /S /GoTo /D <74686D2E442E332E312E32> >> >>}
\MusicPDFWriteObject{3971}{<< /Type /Annot /Subtype /Link /Border [0 0 0] /H /I /C [1 0 0] /Rect [257.121 584.107 287.334 596.965] /A << /S /GoTo /D <74686D2E342E312E322E36> >> >>}
\MusicPDFWriteObject{3972}{<< /Type /Annot /Subtype /Link /Border [0 0 0] /H /I /C [1 0 0] /Rect [504.439 584.107 510.402 596.965] /A << /S /GoTo /D <48666F6F746E6F74652E3331> >> >>}
\MusicPDFWriteObject{3973}{<< /Type /Annot /Subtype /Link /Border [0 0 0] /H /I /C [1 0 0] /Rect [338.28 441.229 369.001 454.21] /A << /S /GoTo /D <74686D2E392E332E322E35> >> >>}
\MusicPDFWriteObject{3974}{<< /Type /Annot /Subtype /Link /Border [0 0 0] /H /I /C [1 0 0] /Rect [382.145 252.365 413.963 264.958] /A << /S /GoTo /D <74686D2E432E342E332E35> >> >>}
\MusicPDFWriteObject{3975}{<< /Type /Annot /Subtype /Link /Border [0 0 0] /H /I /C [1 0 0] /Rect [105.892 226.064 135.541 238.656] /A << /S /GoTo /D <74686D2E342E332E332E31> >> >>}
\MusicPDFWriteObject{3976}{<< /Type /Annot /Subtype /Link /Border [0 0 0] /H /I /C [1 0 0] /Rect [297.38 226.064 329.604 238.656] /A << /S /GoTo /D <74686D2E442E322E322E31> >> >>}
\MusicPDFWriteObject{3977}{<< /Type /Annot /Subtype /Link /Border [0 0 0] /H /I /C [1 0 0] /Rect [71.004 107.707 103.845 119.66] /A << /S /GoTo /D <74686D2E442E332E332E31> >> >>}
\MusicPDFWriteObject{3978}{<< /Type /Annot /Subtype /Link /Border [0 0 0] /H /I /C [1 0 0] /Rect [133.145 107.707 150.493 119.66] /A << /S /GoTo /D <73656374696F6E2E442E33> >> >>}
\MusicPDFWriteObject{3979}{<< /Type /Annot /Subtype /Link /Border [0 0 0] /H /I /C [1 0 0] /Rect [235.058 658.8 267.838 669.676] /A << /S /GoTo /D <74686D2E442E322E312E31> >> >>}
\MusicPDFWriteObject{3980}{<< /Type /Annot /Subtype /Link /Border [0 0 0] /H /I /C [0 1 0] /Rect [294.124 687.123 328.093 695.991] /A << /S /GoTo /D <636974652E436E772D61> >> >>}
\MusicPDFWriteObject{3981}{<< /Type /Annot /Subtype /Link /Border [0 0 0] /H /I /C [0 1 0] /Rect [309.835 673.973 343.804 682.84] /A << /S /GoTo /D <636974652E436E772D61> >> >>}
\MusicPDFWriteObject{3982}{<< /Type /Annot /Subtype /Link /Border [0 0 0] /H /I /C [0 1 0] /Rect [509.195 673.973 543.164 682.84] /A << /S /GoTo /D <636974652E436E772D61> >> >>}
\MusicPDFWriteObject{3983}{<< /Type /Annot /Subtype /Link /Border [0 0 0] /H /I /C [0 1 0] /Rect [402.955 647.671 436.924 656.539] /A << /S /GoTo /D <636974652E436E772D61> >> >>}
\MusicPDFWriteObject{3984}{<< /Type /Annot /Subtype /Link /Border [0 0 0] /H /I /C [1 0 0] /Rect [230.592 334.72 260.41 347.034] /A << /S /GoTo /D <74686D2E392E332E332E32> >> >>}
\MusicPDFWriteObject{3985}{<< /Type /Annot /Subtype /Link /Border [0 0 0] /H /I /C [0 1 0] /Rect [243.297 265.912 277.266 274.78] /A << /S /GoTo /D <636974652E436E772D61> >> >>}
\MusicPDFWriteObject{3986}{<< /Type /Annot /Subtype /Link /Border [0 0 0] /H /I /C [1 0 0] /Rect [71.004 250.271 101.245 262.224] /A << /S /GoTo /D <74686D2E392E332E332E32> >> >>}
\MusicPDFWriteObject{3987}{<< /Type /Annot /Subtype /Link /Border [0 0 0] /H /I /C [0 1 0] /Rect [365.62 252.761 399.589 261.629] /A << /S /GoTo /D <636974652E436E772D61> >> >>}
\MusicPDFWriteObject{3988}{<< /Type /Annot /Subtype /Link /Border [0 0 0] /H /I /C [0 1 0] /Rect [87.604 239.61 121.573 248.478] /A << /S /GoTo /D <636974652E436E772D61> >> >>}
\MusicPDFWriteObject{3989}{<< /Type /Annot /Subtype /Link /Border [0 0 0] /H /I /C [0 1 0] /Rect [107.479 213.309 141.448 222.177] /A << /S /GoTo /D <636974652E436E772D61> >> >>}
\MusicPDFWriteObject{3990}{<< /Type /Annot /Subtype /Link /Border [0 0 0] /H /I /C [0 1 0] /Rect [413.687 187.007 447.656 195.875] /A << /S /GoTo /D <636974652E436E772D61> >> >>}
\MusicPDFWriteObject{3991}{<< /Type /Annot /Subtype /Link /Border [0 0 0] /H /I /C [0 1 0] /Rect [239.822 173.857 273.791 182.725] /A << /S /GoTo /D <636974652E436E772D61> >> >>}
\MusicPDFWriteObject{3992}{<< /Type /Annot /Subtype /Link /Border [0 0 0] /H /I /C [1 0 0] /Rect [250.058 135.198 283.393 147.791] /A << /S /GoTo /D <74686D2E442E322E312E32> >> >>}
\MusicPDFWriteObject{3993}{<< /Type /Annot /Subtype /Link /Border [0 0 0] /H /I /C [1 0 0] /Rect [132.393 96.3 165.71 107.091] /A << /S /GoTo /D <74686D2E412E352E342E31> >> >>}
\MusicPDFWriteObject{3994}{<< /Type /Annot /Subtype /Link /Border [0 0 0] /H /I /C [1 0 0] /Rect [492.383 82.596 525.178 94.549] /A << /S /GoTo /D <74686D2E412E352E332E36> >> >>}
\MusicPDFWriteObject{3995}{<< /Type /Annot /Subtype /Link /Border [0 0 0] /H /I /C [1 0 0] /Rect [182.328 684.633 214.939 696.586] /A << /S /GoTo /D <74686D2E412E352E332E37> >> >>}
\MusicPDFWriteObject{3996}{<< /Type /Annot /Subtype /Link /Border [0 0 0] /H /I /C [1 0 0] /Rect [374.707 671.482 407.41 683.435] /A << /S /GoTo /D <7465726D696E6F6C6F67792E412E352E342E31> >> >>}
\MusicPDFWriteObject{3997}{<< /Type /Annot /Subtype /Link /Border [0 0 0] /H /I /C [1 0 0] /Rect [191.615 658.332 216.112 670.284] /A << /S /GoTo /D <73756273656374696F6E2E412E352E33> >> >>}
\MusicPDFWriteObject{3998}{<< /Type /Annot /Subtype /Link /Border [0 0 0] /H /I /C [1 0 0] /Rect [224.571 658.332 249.068 670.284] /A << /S /GoTo /D <73756273656374696F6E2E412E352E34> >> >>}
\MusicPDFWriteObject{3999}{<< /Type /Annot /Subtype /Link /Border [0 0 0] /H /I /C [0 1 0] /Rect [286.83 660.822 318.774 669.69] /A << /S /GoTo /D <636974652E427269732D61> >> >>}
\MusicPDFWriteObject{4000}{<< /Type /Annot /Subtype /Link /Border [0 0 0] /H /I /C [0 1 0] /Rect [376.674 634.52 408.618 643.388] /A << /S /GoTo /D <636974652E427269732D61> >> >>}
\MusicPDFWriteObject{4001}{<< /Type /Annot /Subtype /Link /Border [0 0 0] /H /I /C [0 1 0] /Rect [272.311 621.369 305.05 630.237] /A << /S /GoTo /D <636974652E427269732D62> >> >>}
\MusicPDFWriteObject{4002}{<< /Type /Annot /Subtype /Link /Border [0 0 0] /H /I /C [1 0 0] /Rect [144.91 403.594 175.123 415.547] /A << /S /GoTo /D <6571756174696F6E2E392E332E332E38> >> >>}
\MusicPDFWriteObject{4003}{<< /Type /Annot /Subtype /Link /Border [0 0 0] /H /I /C [1 0 0] /Rect [305.38 364.142 335.593 376.095] /A << /S /GoTo /D <6571756174696F6E2E392E332E332E38> >> >>}
\MusicPDFWriteObject{4004}{<< /Type /Annot /Subtype /Link /Border [0 0 0] /H /I /C [1 0 0] /Rect [349.748 253.192 380.526 265.507] /A << /S /GoTo /D <74686D2E392E312E312E33> >> >>}
\MusicPDFWriteObject{4005}{<< /Type /Annot /Subtype /Link /Border [0 0 0] /H /I /C [0 1 0] /Rect [107.335 197.62 132.395 206.487] /A << /S /GoTo /D <636974652E53636861> >> >>}
\MusicPDFWriteObject{4006}{<< /Type /Annot /Subtype /Link /Border [0 0 0] /H /I /C [1 0 0] /Rect [278.342 150.523 310.769 161.314] /A << /S /GoTo /D <74686D2E432E332E322E32> >> >>}
\MusicPDFWriteObject{4007}{<< /Type /Annot /Subtype /Link /Border [0 0 0] /H /I /C [1 0 0] /Rect [349.513 559.175 384.707 571.128] /A << /S /GoTo /D <6571756174696F6E2E392E332E332E3133> >> >>}
\MusicPDFWriteObject{4008}{<< /Type /Annot /Subtype /Link /Border [0 0 0] /H /I /C [1 0 0] /Rect [249.704 419.142 284.897 431.095] /A << /S /GoTo /D <6571756174696F6E2E392E332E332E3132> >> >>}
\MusicPDFWriteObject{4009}{<< /Type /Annot /Subtype /Link /Border [0 0 0] /H /I /C [1 0 0] /Rect [483.114 405.992 518.307 417.944] /A << /S /GoTo /D <6571756174696F6E2E392E332E332E3132> >> >>}
\MusicPDFWriteObject{4010}{<< /Type /Annot /Subtype /Link /Border [0 0 0] /H /I /C [0 1 0] /Rect [237.521 355.879 256.495 364.747] /A << /S /GoTo /D <636974652E547265> >> >>}
\MusicPDFWriteObject{4011}{<< /Type /Annot /Subtype /Link /Border [0 0 0] /H /I /C [0 1 0] /Rect [421.624 355.879 440.597 364.747] /A << /S /GoTo /D <636974652E547265> >> >>}
\MusicPDFWriteObject{4012}{<< /Type /Annot /Subtype /Link /Border [0 0 0] /H /I /C [1 0 0] /Rect [131.28 339.878 167.978 352.83] /A << /S /GoTo /D <74686D2E432E322E342E3139> >> >>}
\MusicPDFWriteObject{4013}{<< /Type /Annot /Subtype /Link /Border [0 0 0] /H /I /C [1 0 0] /Rect [154.758 327.087 189.951 339.04] /A << /S /GoTo /D <6571756174696F6E2E392E332E332E3133> >> >>}
\MusicPDFWriteObject{4014}{<< /Type /Annot /Subtype /Link /Border [0 0 0] /H /I /C [1 0 0] /Rect [208.957 214.159 241.66 226.112] /A << /S /GoTo /D <74686D2E412E352E342E31> >> >>}
\MusicPDFWriteObject{4015}{<< /Type /Annot /Subtype /Link /Border [0 0 0] /H /I /C [1 0 0] /Rect [299.872 214.159 332.575 226.112] /A << /S /GoTo /D <74686D2E412E352E342E32> >> >>}
\MusicPDFWriteObject{4016}{<< /Type /Annot /Subtype /Link /Border [0 0 0] /H /I /C [1 0 0] /Rect [455.367 200.599 487.917 215.155] /A << /S /GoTo /D <74686D2E412E352E342E31> >> >>}
\MusicPDFWriteObject{4017}{<< /Type /Annot /Subtype /Link /Border [0 0 0] /H /I /C [0 1 0] /Rect [511.284 203.449 543.228 212.317] /A << /S /GoTo /D <636974652E427269732D61> >> >>}
\MusicPDFWriteObject{4018}{<< /Type /Annot /Subtype /Link /Border [0 0 0] /H /I /C [0 1 0] /Rect [129.595 190.298 161.539 199.166] /A << /S /GoTo /D <636974652E427269732D61> >> >>}
\MusicPDFWriteObject{4019}{<< /Type /Annot /Subtype /Link /Border [0 0 0] /H /I /C [1 0 0] /Rect [181.847 148.356 217.041 160.309] /A << /S /GoTo /D <6571756174696F6E2E392E332E332E3133> >> >>}
\MusicPDFWriteObject{4020}{<< /Type /Annot /Subtype /Link /Border [0 0 0] /H /I /C [0 1 0] /Rect [235.55 687.123 253.749 695.991] /A << /S /GoTo /D <636974652E4A6172> >> >>}
\MusicPDFWriteObject{4021}{<< /Type /Annot /Subtype /Link /Border [0 0 0] /H /I /C [1 0 0] /Rect [333.266 684.633 365.084 697.833] /A << /S /GoTo /D <74686D2E432E332E372E31> >> >>}
\MusicPDFWriteObject{4022}{<< /Type /Annot /Subtype /Link /Border [0 0 0] /H /I /C [1 0 0] /Rect [501.472 684.633 533.289 697.833] /A << /S /GoTo /D <74686D2E432E332E372E31> >> >>}
\MusicPDFWriteObject{4023}{<< /Type /Annot /Subtype /Link /Border [0 0 0] /H /I /C [0 1 0] /Rect [112.557 673.973 144.5 682.84] /A << /S /GoTo /D <636974652E427269732D61> >> >>}
\MusicPDFWriteObject{4024}{<< /Type /Annot /Subtype /Link /Border [0 0 0] /H /I /C [1 0 0] /Rect [229.47 658.332 264.663 670.284] /A << /S /GoTo /D <6571756174696F6E2E392E332E332E3131> >> >>}
\MusicPDFWriteObject{4025}{<< /Type /Annot /Subtype /Link /Border [0 0 0] /H /I /C [1 0 0] /Rect [247.742 632.03 282.935 643.983] /A << /S /GoTo /D <6571756174696F6E2E392E332E332E3133> >> >>}
\MusicPDFWriteObject{4026}{<< /Type /Annot /Subtype /Link /Border [0 0 0] /H /I /C [1 0 0] /Rect [500.393 524.474 532.971 536.427] /A << /S /GoTo /D <74686D2E432E332E372E31> >> >>}
\MusicPDFWriteObject{4027}{<< /Type /Annot /Subtype /Link /Border [0 0 0] /H /I /C [1 0 0] /Rect [285.694 511.323 320.888 523.276] /A << /S /GoTo /D <6571756174696F6E2E392E332E332E3134> >> >>}
\MusicPDFWriteObject{4028}{<< /Type /Annot /Subtype /Link /Border [0 0 0] /H /I /C [0 1 0] /Rect [289.105 500.662 321.049 509.53] /A << /S /GoTo /D <636974652E427269732D61> >> >>}
\MusicPDFWriteObject{4029}{<< /Type /Annot /Subtype /Link /Border [0 0 0] /H /I /C [1 0 0] /Rect [213.925 106.206 244.138 118.158] /A << /S /GoTo /D <6571756174696F6E2E392E332E332E39> >> >>}
\MusicPDFWriteObject{4030}{<< /Type /Annot /Subtype /Link /Border [0 0 0] /H /I /C [1 0 0] /Rect [82.959 556.115 447.467 570.97] /A << /S /GoTo /D <73656374696F6E2E31302E31> >> >>}
\MusicPDFWriteObject{4031}{<< /Type /Annot /Subtype /Link /Border [0 0 0] /H /I /C [1 0 0] /Rect [82.959 542.255 465.217 557.111] /A << /S /GoTo /D <73756273656374696F6E2E31302E312E31> >> >>}
\MusicPDFWriteObject{4032}{<< /Type /Annot /Subtype /Link /Border [0 0 0] /H /I /C [1 0 0] /Rect [82.959 533.389 200.075 542.243] /A << /S /GoTo /D <73756273656374696F6E2E31302E312E32> >> >>}
\MusicPDFWriteObject{4033}{<< /Type /Annot /Subtype /Link /Border [0 0 0] /H /I /C [1 0 0] /Rect [82.959 518.301 242.461 529.092] /A << /S /GoTo /D <73656374696F6E2E31302E32> >> >>}
\MusicPDFWriteObject{4034}{<< /Type /Annot /Subtype /Link /Border [0 0 0] /H /I /C [1 0 0] /Rect [82.959 505.15 296.422 515.941] /A << /S /GoTo /D <73756273656374696F6E2E31302E322E31> >> >>}
\MusicPDFWriteObject{4035}{<< /Type /Annot /Subtype /Link /Border [0 0 0] /H /I /C [1 0 0] /Rect [82.959 493.936 199.891 502.79] /A << /S /GoTo /D <73756273756273656374696F6E2E31302E322E312E31> >> >>}
\MusicPDFWriteObject{4036}{<< /Type /Annot /Subtype /Link /Border [0 0 0] /H /I /C [1 0 0] /Rect [82.959 478.849 281.26 489.64] /A << /S /GoTo /D <73756273756273656374696F6E2E31302E322E312E32> >> >>}
\MusicPDFWriteObject{4037}{<< /Type /Annot /Subtype /Link /Border [0 0 0] /H /I /C [1 0 0] /Rect [82.959 465.144 463.249 477.099] /A << /S /GoTo /D <73756273656374696F6E2E31302E322E32> >> >>}
\MusicPDFWriteObject{4038}{<< /Type /Annot /Subtype /Link /Border [0 0 0] /H /I /C [1 0 0] /Rect [82.959 452.547 326.723 463.338] /A << /S /GoTo /D <73756273756273656374696F6E2E31302E322E322E31> >> >>}
\MusicPDFWriteObject{4039}{<< /Type /Annot /Subtype /Link /Border [0 0 0] /H /I /C [1 0 0] /Rect [82.959 439.397 448.652 450.187] /A << /S /GoTo /D <73756273756273656374696F6E2E31302E322E322E32> >> >>}
\MusicPDFWriteObject{4040}{<< /Type /Annot /Subtype /Link /Border [0 0 0] /H /I /C [1 0 0] /Rect [82.959 426.246 199.549 437.037] /A << /S /GoTo /D <73656374696F6E2E31302E33> >> >>}
\MusicPDFWriteObject{4041}{<< /Type /Annot /Subtype /Link /Border [0 0 0] /H /I /C [1 0 0] /Rect [82.959 412.541 283.901 424.496] /A << /S /GoTo /D <73756273656374696F6E2E31302E332E31> >> >>}
\MusicPDFWriteObject{4042}{<< /Type /Annot /Subtype /Link /Border [0 0 0] /H /I /C [1 0 0] /Rect [82.959 399.39 424.931 411.346] /A << /S /GoTo /D <73756273656374696F6E2E31302E332E32> >> >>}
\MusicPDFWriteObject{4043}{<< /Type /Annot /Subtype /Link /Border [0 0 0] /H /I /C [1 0 0] /Rect [82.959 386.794 240.524 397.584] /A << /S /GoTo /D <73656374696F6E2E31302E34> >> >>}
\MusicPDFWriteObject{4044}{<< /Type /Annot /Subtype /Link /Border [0 0 0] /H /I /C [1 0 0] /Rect [82.959 375.58 174.953 384.434] /A << /S /GoTo /D <73756273656374696F6E2E31302E342E31> >> >>}
\MusicPDFWriteObject{4045}{<< /Type /Annot /Subtype /Link /Border [0 0 0] /H /I /C [1 0 0] /Rect [82.959 360.492 180.431 371.283] /A << /S /GoTo /D <73756273656374696F6E2E31302E342E32> >> >>}
\MusicPDFWriteObject{4046}{<< /Type /Annot /Subtype /Link /Border [0 0 0] /H /I /C [1 0 0] /Rect [82.959 347.341 229.151 358.132] /A << /S /GoTo /D <73756273756273656374696F6E2E31302E342E322E31> >> >>}
\MusicPDFWriteObject{4047}{<< /Type /Annot /Subtype /Link /Border [0 0 0] /H /I /C [1 0 0] /Rect [82.959 334.191 316.333 344.981] /A << /S /GoTo /D <73756273756273656374696F6E2E31302E342E322E32> >> >>}
\MusicPDFWriteObject{4048}{<< /Type /Annot /Subtype /Link /Border [0 0 0] /H /I /C [1 0 0] /Rect [82.959 321.04 194.901 331.83] /A << /S /GoTo /D <73756273656374696F6E2E31302E342E33> >> >>}
\MusicPDFWriteObject{4049}{<< /Type /Annot /Subtype /Link /Border [0 0 0] /H /I /C [1 0 0] /Rect [82.959 307.335 278.699 319.29] /A << /S /GoTo /D <73756273756273656374696F6E2E31302E342E332E31> >> >>}
\MusicPDFWriteObject{4050}{<< /Type /Annot /Subtype /Link /Border [0 0 0] /H /I /C [1 0 0] /Rect [82.959 294.184 419.73 306.139] /A << /S /GoTo /D <73756273756273656374696F6E2E31302E342E332E32> >> >>}
\MusicPDFWriteObject{4051}{<< /Type /Annot /Subtype /Link /Border [0 0 0] /H /I /C [1 0 0] /Rect [86.497 626.053 116.71 638.005] /A << /S /GoTo /D <6571756174696F6E2E392E312E322E33> >> >>}
\MusicPDFWriteObject{4052}{<< /Type /Annot /Subtype /Link /Border [0 0 0] /H /I /C [1 0 0] /Rect [179.598 625.692 209.811 638.005] /A << /S /GoTo /D <74686D2E392E312E322E31> >> >>}
\MusicPDFWriteObject{4053}{<< /Type /Annot /Subtype /Link /Border [0 0 0] /H /I /C [1 0 0] /Rect [146.65 580.25 176.863 591.041] /A << /S /GoTo /D <74686D2E322E322E332E34> >> >>}
\MusicPDFWriteObject{4054}{<< /Type /Annot /Subtype /Link /Border [0 0 0] /H /I /C [1 0 0] /Rect [233.11 580.25 263.324 591.041] /A << /S /GoTo /D <74686D2E322E322E342E33> >> >>}
\MusicPDFWriteObject{4055}{<< /Type /Annot /Subtype /Link /Border [0 0 0] /H /I /C [1 0 0] /Rect [211.124 491.22 240.773 502.01] /A << /S /GoTo /D <74686D2E322E322E342E33> >> >>}
\MusicPDFWriteObject{4056}{<< /Type /Annot /Subtype /Link /Border [0 0 0] /H /I /C [1 0 0] /Rect [85.83 451.214 121.024 463.167] /A << /S /GoTo /D <6571756174696F6E2E31302E312E312E31> >> >>}
\MusicPDFWriteObject{4057}{<< /Type /Annot /Subtype /Link /Border [0 0 0] /H /I /C [1 0 0] /Rect [228.982 438.063 264.175 450.016] /A << /S /GoTo /D <6571756174696F6E2E31302E312E312E31> >> >>}
\MusicPDFWriteObject{4058}{<< /Type /Annot /Subtype /Link /Border [0 0 0] /H /I /C [1 0 0] /Rect [226.549 424.913 256.762 436.865] /A << /S /GoTo /D <74686D2E352E322E312E31> >> >>}
\MusicPDFWriteObject{4059}{<< /Type /Annot /Subtype /Link /Border [0 0 0] /H /I /C [1 0 0] /Rect [283.321 411.762 318.514 423.715] /A << /S /GoTo /D <6571756174696F6E2E31302E312E312E31> >> >>}
\MusicPDFWriteObject{4060}{<< /Type /Annot /Subtype /Link /Border [0 0 0] /H /I /C [1 0 0] /Rect [325.679 109.012 356.457 120.965] /A << /S /GoTo /D <74686D2E322E342E332E31> >> >>}
\MusicPDFWriteObject{4061}{<< /Type /Annot /Subtype /Link /Border [0 0 0] /H /I /C [1 0 0] /Rect [434.783 109.012 469.976 120.965] /A << /S /GoTo /D <6571756174696F6E2E31302E312E312E35> >> >>}
\MusicPDFWriteObject{4062}{<< /Type /Annot /Subtype /Link /Border [0 0 0] /H /I /C [1 0 0] /Rect [294.211 655.83 300.174 670.685] /A << /S /GoTo /D <48666F6F746E6F74652E3332> >> >>}
\MusicPDFWriteObject{4063}{<< /Type /Annot /Subtype /Link /Border [0 0 0] /H /I /C [1 0 0] /Rect [414.016 609.059 444.285 621.835] /A << /S /GoTo /D <74686D2E352E322E342E31> >> >>}
\MusicPDFWriteObject{4064}{<< /Type /Annot /Subtype /Link /Border [0 0 0] /H /I /C [1 0 0] /Rect [263.789 583.118 298.982 595.073] /A << /S /GoTo /D <74686D2E31302E312E312E31> >> >>}
\MusicPDFWriteObject{4065}{<< /Type /Annot /Subtype /Link /Border [0 0 0] /H /I /C [1 0 0] /Rect [472.958 381.749 503.087 394.86] /A << /S /GoTo /D <74686D2E392E312E312E32> >> >>}
\MusicPDFWriteObject{4066}{<< /Type /Annot /Subtype /Link /Border [0 0 0] /H /I /C [1 0 0] /Rect [501.094 381.749 507.057 394.86] /A << /S /GoTo /D <48666F6F746E6F74652E3333> >> >>}
\MusicPDFWriteObject{4067}{<< /Type /Annot /Subtype /Link /Border [0 0 0] /H /I /C [1 0 0] /Rect [470.042 355.121 505.235 367.074] /A << /S /GoTo /D <6571756174696F6E2E31302E312E312E36> >> >>}
\MusicPDFWriteObject{4068}{<< /Type /Annot /Subtype /Link /Border [0 0 0] /H /I /C [1 0 0] /Rect [419.278 255.525 451.981 267.478] /A << /S /GoTo /D <6571756174696F6E2E412E332E342E32> >> >>}
\MusicPDFWriteObject{4069}{<< /Type /Annot /Subtype /Link /Border [0 0 0] /H /I /C [1 0 0] /Rect [243.061 240.58 273.839 255.436] /A << /S /GoTo /D <74686D2E322E342E322E31> >> >>}
\MusicPDFWriteObject{4070}{<< /Type /Annot /Subtype /Link /Border [0 0 0] /H /I /C [1 0 0] /Rect [350.636 240.58 381.414 255.436] /A << /S /GoTo /D <74686D2E322E332E312E32> >> >>}
\MusicPDFWriteObject{4071}{<< /Type /Annot /Subtype /Link /Border [0 0 0] /H /I /C [0 1 0] /Rect [260.638 504.41 277.792 513.278] /A << /S /GoTo /D <636974652E466F6C> >> >>}
\MusicPDFWriteObject{4072}{<< /Type /Annot /Subtype /Link /Border [0 0 0] /H /I /C [1 0 0] /Rect [485.316 474.91 520.509 486.863] /A << /S /GoTo /D <6571756174696F6E2E31302E312E312E35> >> >>}
\MusicPDFWriteObject{4073}{<< /Type /Annot /Subtype /Link /Border [0 0 0] /H /I /C [1 0 0] /Rect [298.81 460.895 328.459 474.113] /A << /S /GoTo /D <74686D2E392E312E312E32> >> >>}
\MusicPDFWriteObject{4074}{<< /Type /Annot /Subtype /Link /Border [0 0 0] /H /I /C [1 0 0] /Rect [326.467 460.895 332.429 474.113] /A << /S /GoTo /D <48666F6F746E6F74652E3334> >> >>}
\MusicPDFWriteObject{4075}{<< /Type /Annot /Subtype /Link /Border [0 0 0] /H /I /C [1 0 0] /Rect [142.155 481.056 177.083 491.847] /A << /S /GoTo /D <74686D2E31302E312E312E35> >> >>}
\MusicPDFWriteObject{4076}{<< /Type /Annot /Subtype /Link /Border [0 0 0] /H /I /C [1 0 0] /Rect [486.86 481.056 516.848 491.847] /A << /S /GoTo /D <74686D2E352E322E332E32> >> >>}
\MusicPDFWriteObject{4077}{<< /Type /Annot /Subtype /Link /Border [0 0 0] /H /I /C [1 0 0] /Rect [307.477 454.201 329.943 467.059] /A << /S /GoTo /D <73756273656374696F6E2E342E332E32> >> >>}
\MusicPDFWriteObject{4078}{<< /Type /Annot /Subtype /Link /Border [0 0 0] /H /I /C [1 0 0] /Rect [491.132 454.201 521.345 467.059] /A << /S /GoTo /D <74686D2E342E332E322E32> >> >>}
\MusicPDFWriteObject{4079}{<< /Type /Annot /Subtype /Link /Border [0 0 0] /H /I /C [1 0 0] /Rect [525.993 454.201 531.955 467.059] /A << /S /GoTo /D <48666F6F746E6F74652E3335> >> >>}
\MusicPDFWriteObject{4080}{<< /Type /Annot /Subtype /Link /Border [0 0 0] /H /I /C [1 0 0] /Rect [345.841 300.553 376.054 312.506] /A << /S /GoTo /D <6571756174696F6E2E352E322E342E33> >> >>}
\MusicPDFWriteObject{4081}{<< /Type /Annot /Subtype /Link /Border [0 0 0] /H /I /C [1 0 0] /Rect [280.913 645.734 311.691 656.525] /A << /S /GoTo /D <74686D2E322E322E332E32> >> >>}
\MusicPDFWriteObject{4082}{<< /Type /Annot /Subtype /Link /Border [0 0 0] /H /I /C [1 0 0] /Rect [403.883 492.738 409.846 505.85] /A << /S /GoTo /D <48666F6F746E6F74652E3336> >> >>}
\MusicPDFWriteObject{4083}{<< /Type /Annot /Subtype /Link /Border [0 0 0] /H /I /C [1 0 0] /Rect [485.281 492.992 520.475 504.945] /A << /S /GoTo /D <6571756174696F6E2E31302E312E322E34> >> >>}
\MusicPDFWriteObject{4084}{<< /Type /Annot /Subtype /Link /Border [0 0 0] /H /I /C [1 0 0] /Rect [86.995 466.691 122.189 478.643] /A << /S /GoTo /D <6571756174696F6E2E31302E312E322E35> >> >>}
\MusicPDFWriteObject{4085}{<< /Type /Annot /Subtype /Link /Border [0 0 0] /H /I /C [1 0 0] /Rect [200.708 466.691 235.901 478.643] /A << /S /GoTo /D <6571756174696F6E2E31302E312E322E32> >> >>}
\MusicPDFWriteObject{4086}{<< /Type /Annot /Subtype /Link /Border [0 0 0] /H /I /C [1 0 0] /Rect [294.308 440.389 329.501 452.342] /A << /S /GoTo /D <6571756174696F6E2E31302E312E322E32> >> >>}
\MusicPDFWriteObject{4087}{<< /Type /Annot /Subtype /Link /Border [0 0 0] /H /I /C [1 0 0] /Rect [237.79 408.663 267.439 419.454] /A << /S /GoTo /D <74686D2E352E342E322E31> >> >>}
\MusicPDFWriteObject{4088}{<< /Type /Annot /Subtype /Link /Border [0 0 0] /H /I /C [1 0 0] /Rect [182.157 341.938 212.652 353.893] /A << /S /GoTo /D <74686D2E332E332E312E35> >> >>}
\MusicPDFWriteObject{4089}{<< /Type /Annot /Subtype /Link /Border [0 0 0] /H /I /C [1 0 0] /Rect [306.525 341.939 336.738 353.891] /A << /S /GoTo /D <6571756174696F6E2E332E332E312E33> >> >>}
\MusicPDFWriteObject{4090}{<< /Type /Annot /Subtype /Link /Border [0 0 0] /H /I /C [0 1 0] /Rect [193.659 264.329 211.858 273.197] /A << /S /GoTo /D <636974652E4A6172> >> >>}
\MusicPDFWriteObject{4091}{<< /Type /Annot /Subtype /Link /Border [0 0 0] /H /I /C [1 0 0] /Rect [325.874 671.823 361.068 682.614] /A << /S /GoTo /D <74686D2E31302E312E312E31> >> >>}
\MusicPDFWriteObject{4092}{<< /Type /Annot /Subtype /Link /Border [0 0 0] /H /I /C [1 0 0] /Rect [226.332 426.855 256.545 438.808] /A << /S /GoTo /D <6571756174696F6E2E372E312E342E37> >> >>}
\MusicPDFWriteObject{4093}{<< /Type /Annot /Subtype /Link /Border [0 0 0] /H /I /C [1 0 0] /Rect [281.417 374.252 311.63 386.205] /A << /S /GoTo /D <6571756174696F6E2E392E332E312E32> >> >>}
\MusicPDFWriteObject{4094}{<< /Type /Annot /Subtype /Link /Border [0 0 0] /H /I /C [0 1 0] /Rect [347.763 363.592 363.755 372.459] /A << /S /GoTo /D <636974652E436F> >> >>}
\MusicPDFWriteObject{4095}{<< /Type /Annot /Subtype /Link /Border [0 0 0] /H /I /C [0 1 0] /Rect [445.921 363.592 461.913 372.459] /A << /S /GoTo /D <636974652E436F> >> >>}
\MusicPDFWriteObject{4096}{<< /Type /Annot /Subtype /Link /Border [0 0 0] /H /I /C [0 1 0] /Rect [194.757 310.989 210.749 319.856] /A << /S /GoTo /D <636974652E436F> >> >>}
\MusicPDFWriteObject{4097}{<< /Type /Annot /Subtype /Link /Border [0 0 0] /H /I /C [0 1 0] /Rect [311.44 310.989 327.432 319.856] /A << /S /GoTo /D <636974652E436F> >> >>}
\MusicPDFWriteObject{4098}{<< /Type /Annot /Subtype /Link /Border [0 0 0] /H /I /C [0 1 0] /Rect [104.693 258.386 120.685 267.253] /A << /S /GoTo /D <636974652E436F> >> >>}
\MusicPDFWriteObject{4099}{<< /Type /Annot /Subtype /Link /Border [0 0 0] /H /I /C [0 1 0] /Rect [260.216 232.084 276.209 240.952] /A << /S /GoTo /D <636974652E436F> >> >>}
\MusicPDFWriteObject{4100}{<< /Type /Annot /Subtype /Link /Border [0 0 0] /H /I /C [0 1 0] /Rect [269.478 218.933 285.47 227.801] /A << /S /GoTo /D <636974652E436F> >> >>}
\MusicPDFWriteObject{4101}{<< /Type /Annot /Subtype /Link /Border [0 0 0] /H /I /C [1 0 0] /Rect [348.119 113.96 378.332 125.913] /A << /S /GoTo /D <74686D2E342E332E332E31> >> >>}
\MusicPDFWriteObject{4102}{<< /Type /Annot /Subtype /Link /Border [0 0 0] /H /I /C [1 0 0] /Rect [318.511 100.81 353.704 112.762] /A << /S /GoTo /D <6571756174696F6E2E392E312E312E3131> >> >>}
\MusicPDFWriteObject{4103}{<< /Type /Annot /Subtype /Link /Border [0 0 0] /H /I /C [0 1 0] /Rect [300.65 648.62 319.754 657.488] /A << /S /GoTo /D <636974652E4E42> >> >>}
\MusicPDFWriteObject{4104}{<< /Type /Annot /Subtype /Link /Border [0 0 0] /H /I /C [0 1 0] /Rect [293.838 635.414 312.943 644.282] /A << /S /GoTo /D <636974652E4E42> >> >>}
\MusicPDFWriteObject{4105}{<< /Type /Annot /Subtype /Link /Border [0 0 0] /H /I /C [1 0 0] /Rect [503.789 252.201 538.982 264.153] /A << /S /GoTo /D <6571756174696F6E2E392E312E322E3134> >> >>}
\MusicPDFWriteObject{4106}{<< /Type /Annot /Subtype /Link /Border [0 0 0] /H /I /C [1 0 0] /Rect [298.797 239.05 333.991 251.003] /A << /S /GoTo /D <6571756174696F6E2E31302E322E312E33> >> >>}
\MusicPDFWriteObject{4107}{<< /Type /Annot /Subtype /Link /Border [0 0 0] /H /I /C [1 0 0] /Rect [337.625 182.061 372.818 194.014] /A << /S /GoTo /D <6571756174696F6E2E31302E322E312E35> >> >>}
\MusicPDFWriteObject{4108}{<< /Type /Annot /Subtype /Link /Border [0 0 0] /H /I /C [0 1 0] /Rect [351.475 150.84 369.674 159.708] /A << /S /GoTo /D <636974652E4A6172> >> >>}
\MusicPDFWriteObject{4109}{<< /Type /Annot /Subtype /Link /Border [0 0 0] /H /I /C [0 1 0] /Rect [292.302 85.086 310.501 93.954] /A << /S /GoTo /D <636974652E4A6172> >> >>}
\MusicPDFWriteObject{4110}{<< /Type /Annot /Subtype /Link /Border [0 0 0] /H /I /C [1 0 0] /Rect [187.794 604.554 226.861 613.408] /A << /S /GoTo /D <74686D2E31302E322E312E31> >> >>}
\MusicPDFWriteObject{4111}{<< /Type /Annot /Subtype /Link /Border [0 0 0] /H /I /C [1 0 0] /Rect [113.149 468.959 143.673 481.274] /A << /S /GoTo /D <74686D2E362E332E322E31> >> >>}
\MusicPDFWriteObject{4112}{<< /Type /Annot /Subtype /Link /Border [0 0 0] /H /I /C [1 0 0] /Rect [71.004 443.572 104.462 454.362] /A << /S /GoTo /D <74686D2E442E332E332E31> >> >>}
\MusicPDFWriteObject{4113}{<< /Type /Annot /Subtype /Link /Border [0 0 0] /H /I /C [1 0 0] /Rect [146.042 443.572 179.5 454.362] /A << /S /GoTo /D <74686D2E442E332E322E37> >> >>}
\MusicPDFWriteObject{4114}{<< /Type /Annot /Subtype /Link /Border [0 0 0] /H /I /C [1 0 0] /Rect [245.017 443.572 278.475 454.362] /A << /S /GoTo /D <74686D2E442E332E332E37> >> >>}
\MusicPDFWriteObject{4115}{<< /Type /Annot /Subtype /Link /Border [0 0 0] /H /I /C [1 0 0] /Rect [238.04 360.943 273.233 372.896] /A << /S /GoTo /D <6571756174696F6E2E31302E322E312E37> >> >>}
\MusicPDFWriteObject{4116}{<< /Type /Annot /Subtype /Link /Border [0 0 0] /H /I /C [1 0 0] /Rect [98.422 347.793 133.615 359.745] /A << /S /GoTo /D <6571756174696F6E2E31302E322E312E33> >> >>}
\MusicPDFWriteObject{4117}{<< /Type /Annot /Subtype /Link /Border [0 0 0] /H /I /C [1 0 0] /Rect [262.087 285.063 297.28 297.016] /A << /S /GoTo /D <6571756174696F6E2E31302E322E312E37> >> >>}
\MusicPDFWriteObject{4118}{<< /Type /Annot /Subtype /Link /Border [0 0 0] /H /I /C [1 0 0] /Rect [511.675 195.553 541.888 207.506] /A << /S /GoTo /D <6571756174696F6E2E352E312E312E34> >> >>}
\MusicPDFWriteObject{4119}{<< /Type /Annot /Subtype /Link /Border [0 0 0] /H /I /C [1 0 0] /Rect [119.1 166.878 148.918 181.351] /A << /S /GoTo /D <74686D2E342E342E312E31> >> >>}
\MusicPDFWriteObject{4120}{<< /Type /Annot /Subtype /Link /Border [0 0 0] /H /I /C [1 0 0] /Rect [153.957 95.939 184.735 109.24] /A << /S /GoTo /D <74686D2E362E332E322E31> >> >>}
\MusicPDFWriteObject{4121}{<< /Type /Annot /Subtype /Link /Border [0 0 0] /H /I /C [1 0 0] /Rect [312.437 645.374 342.65 657.327] /A << /S /GoTo /D <6571756174696F6E2E372E312E342E33> >> >>}
\MusicPDFWriteObject{4122}{<< /Type /Annot /Subtype /Link /Border [0 0 0] /H /I /C [1 0 0] /Rect [273.052 632.224 308.245 644.176] /A << /S /GoTo /D <6571756174696F6E2E31302E322E312E38> >> >>}
\MusicPDFWriteObject{4123}{<< /Type /Annot /Subtype /Link /Border [0 0 0] /H /I /C [1 0 0] /Rect [285.156 504.025 315.369 516.34] /A << /S /GoTo /D <74686D2E392E312E322E35> >> >>}
\MusicPDFWriteObject{4124}{<< /Type /Annot /Subtype /Link /Border [0 0 0] /H /I /C [1 0 0] /Rect [356.336 478.084 391.529 490.037] /A << /S /GoTo /D <6571756174696F6E2E31302E322E312E38> >> >>}
\MusicPDFWriteObject{4125}{<< /Type /Annot /Subtype /Link /Border [0 0 0] /H /I /C [1 0 0] /Rect [440.802 451.422 471.579 463.737] /A << /S /GoTo /D <74686D2E372E312E342E32> >> >>}
\MusicPDFWriteObject{4126}{<< /Type /Annot /Subtype /Link /Border [0 0 0] /H /I /C [1 0 0] /Rect [500.274 437.5 530.487 452.419] /A << /S /GoTo /D <74686D2E362E332E322E31> >> >>}
\MusicPDFWriteObject{4127}{<< /Type /Annot /Subtype /Link /Border [0 0 0] /H /I /C [1 0 0] /Rect [339.729 424.38 374.922 436.333] /A << /S /GoTo /D <6571756174696F6E2E31302E322E312E39> >> >>}
\MusicPDFWriteObject{4128}{<< /Type /Annot /Subtype /Link /Border [0 0 0] /H /I /C [1 0 0] /Rect [267.089 398.078 297.867 410.033] /A << /S /GoTo /D <74686D2E362E332E322E31> >> >>}
\MusicPDFWriteObject{4129}{<< /Type /Annot /Subtype /Link /Border [0 0 0] /H /I /C [1 0 0] /Rect [142.012 370.394 172.789 384.314] /A << /S /GoTo /D <74686D2E392E332E322E31> >> >>}
\MusicPDFWriteObject{4130}{<< /Type /Annot /Subtype /Link /Border [0 0 0] /H /I /C [1 0 0] /Rect [181.952 352.415 217.809 368.645] /A << /S /GoTo /D <6E6F746174696F6E2E31302E322E312E31> >> >>}
\MusicPDFWriteObject{4131}{<< /Type /Annot /Subtype /Link /Border [0 0 0] /H /I /C [1 0 0] /Rect [515.415 201.74 545.431 214.333] /A << /S /GoTo /D <74686D2E392E332E322E31> >> >>}
\MusicPDFWriteObject{4132}{<< /Type /Annot /Subtype /Link /Border [0 0 0] /H /I /C [1 0 0] /Rect [421.255 669.435 450.904 684.075] /A << /S /GoTo /D <74686D2E392E312E322E36> >> >>}
\MusicPDFWriteObject{4133}{<< /Type /Annot /Subtype /Link /Border [0 0 0] /H /I /C [1 0 0] /Rect [392.987 599.072 422.635 613.712] /A << /S /GoTo /D <74686D2E392E332E322E31> >> >>}
\MusicPDFWriteObject{4134}{<< /Type /Annot /Subtype /Link /Border [0 0 0] /H /I /C [1 0 0] /Rect [174.843 417.072 210.036 429.025] /A << /S /GoTo /D <6571756174696F6E2E31302E322E312E35> >> >>}
\MusicPDFWriteObject{4135}{<< /Type /Annot /Subtype /Link /Border [0 0 0] /H /I /C [1 0 0] /Rect [268.663 416.712 303.856 429.026] /A << /S /GoTo /D <74686D2E31302E322E312E31> >> >>}
\MusicPDFWriteObject{4136}{<< /Type /Annot /Subtype /Link /Border [0 0 0] /H /I /C [1 0 0] /Rect [314.007 380.644 349.2 392.597] /A << /S /GoTo /D <6571756174696F6E2E31302E322E312E33> >> >>}
\MusicPDFWriteObject{4137}{<< /Type /Annot /Subtype /Link /Border [0 0 0] /H /I /C [0 1 0] /Rect [335.31 343.682 354.415 352.55] /A << /S /GoTo /D <636974652E4E42> >> >>}
\MusicPDFWriteObject{4138}{<< /Type /Annot /Subtype /Link /Border [0 0 0] /H /I /C [0 1 0] /Rect [452.891 343.682 471.995 352.55] /A << /S /GoTo /D <636974652E4E42> >> >>}
\MusicPDFWriteObject{4139}{<< /Type /Annot /Subtype /Link /Border [0 0 0] /H /I /C [1 0 0] /Rect [488.442 234.121 518.655 246.074] /A << /S /GoTo /D <6571756174696F6E2E362E332E312E31> >> >>}
\MusicPDFWriteObject{4140}{<< /Type /Annot /Subtype /Link /Border [0 0 0] /H /I /C [1 0 0] /Rect [391.203 207.819 426.396 219.772] /A << /S /GoTo /D <6571756174696F6E2E31302E322E312E36> >> >>}
\MusicPDFWriteObject{4141}{<< /Type /Annot /Subtype /Link /Border [0 0 0] /H /I /C [1 0 0] /Rect [248.914 151.996 289.087 163.949] /A << /S /GoTo /D <6571756174696F6E2E31302E322E312E3131> >> >>}
\MusicPDFWriteObject{4142}{<< /Type /Annot /Subtype /Link /Border [0 0 0] /H /I /C [1 0 0] /Rect [211.393 114.453 246.587 125.243] /A << /S /GoTo /D <74686D2E31312E322E302E34> >> >>}
\MusicPDFWriteObject{4143}{<< /Type /Annot /Subtype /Link /Border [0 0 0] /H /I /C [1 0 0] /Rect [138.22 684.273 168.32 696.588] /A << /S /GoTo /D <74686D2E392E312E322E35> >> >>}
\MusicPDFWriteObject{4144}{<< /Type /Annot /Subtype /Link /Border [0 0 0] /H /I /C [1 0 0] /Rect [446.834 638.727 482.559 650.68] /A << /S /GoTo /D <74686D2E31312E322E302E35> >> >>}
\MusicPDFWriteObject{4145}{<< /Type /Annot /Subtype /Link /Border [0 0 0] /H /I /C [0 1 0] /Rect [194.455 628.066 214.175 636.934] /A << /S /GoTo /D <636974652E4B616C> >> >>}
\MusicPDFWriteObject{4146}{<< /Type /Annot /Subtype /Link /Border [0 0 0] /H /I /C [0 1 0] /Rect [161.163 614.916 180.883 623.783] /A << /S /GoTo /D <636974652E4B616C> >> >>}
\MusicPDFWriteObject{4147}{<< /Type /Annot /Subtype /Link /Border [0 0 0] /H /I /C [0 1 0] /Rect [192.517 575.463 212.238 584.331] /A << /S /GoTo /D <636974652E4B616C> >> >>}
\MusicPDFWriteObject{4148}{<< /Type /Annot /Subtype /Link /Border [0 0 0] /H /I /C [0 1 0] /Rect [151.077 549.162 170.797 558.03] /A << /S /GoTo /D <636974652E4B616C> >> >>}
\MusicPDFWriteObject{4149}{<< /Type /Annot /Subtype /Link /Border [0 0 0] /H /I /C [1 0 0] /Rect [488.442 361.645 518.655 373.597] /A << /S /GoTo /D <6571756174696F6E2E362E332E312E31> >> >>}
\MusicPDFWriteObject{4150}{<< /Type /Annot /Subtype /Link /Border [0 0 0] /H /I /C [1 0 0] /Rect [394.375 335.343 429.568 347.296] /A << /S /GoTo /D <6571756174696F6E2E31302E322E312E36> >> >>}
\MusicPDFWriteObject{4151}{<< /Type /Annot /Subtype /Link /Border [0 0 0] /H /I /C [1 0 0] /Rect [468.646 320.1 503.839 332.053] /A << /S /GoTo /D <6571756174696F6E2E392E312E322E3134> >> >>}
\MusicPDFWriteObject{4152}{<< /Type /Annot /Subtype /Link /Border [0 0 0] /H /I /C [1 0 0] /Rect [498.847 306.468 534.041 318.421] /A << /S /GoTo /D <6571756174696F6E2E31302E322E312E33> >> >>}
\MusicPDFWriteObject{4153}{<< /Type /Annot /Subtype /Link /Border [0 0 0] /H /I /C [1 0 0] /Rect [295.544 273.713 335.717 285.665] /A << /S /GoTo /D <6571756174696F6E2E31302E322E312E3132> >> >>}
\MusicPDFWriteObject{4154}{<< /Type /Annot /Subtype /Link /Border [0 0 0] /H /I /C [1 0 0] /Rect [405.881 273.352 441.074 285.667] /A << /S /GoTo /D <74686D2E31302E322E312E33> >> >>}
\MusicPDFWriteObject{4155}{<< /Type /Annot /Subtype /Link /Border [0 0 0] /H /I /C [1 0 0] /Rect [238.2 260.202 274.057 272.516] /A << /S /GoTo /D <74686D2E31302E322E312E33> >> >>}
\MusicPDFWriteObject{4156}{<< /Type /Annot /Subtype /Link /Border [0 0 0] /H /I /C [1 0 0] /Rect [164.305 245.724 199.498 259.383] /A << /S /GoTo /D <74686D2E31302E322E312E31> >> >>}
\MusicPDFWriteObject{4157}{<< /Type /Annot /Subtype /Link /Border [0 0 0] /H /I /C [1 0 0] /Rect [166.76 215.209 201.953 226] /A << /S /GoTo /D <7465726D696E6F6C6F67792E31302E322E312E32> >> >>}
\MusicPDFWriteObject{4158}{<< /Type /Annot /Subtype /Link /Border [0 0 0] /H /I /C [1 0 0] /Rect [293.443 186.667 329.3 200.327] /A << /S /GoTo /D <74686D2E31302E322E312E31> >> >>}
\MusicPDFWriteObject{4159}{<< /Type /Annot /Subtype /Link /Border [0 0 0] /H /I /C [1 0 0] /Rect [488.072 186.667 523.929 200.327] /A << /S /GoTo /D <74686D2E31302E322E312E33> >> >>}
\MusicPDFWriteObject{4160}{<< /Type /Annot /Subtype /Link /Border [0 0 0] /H /I /C [1 0 0] /Rect [115.05 175.757 150.243 186.548] /A << /S /GoTo /D <74686D2E31302E322E312E34> >> >>}
\MusicPDFWriteObject{4161}{<< /Type /Annot /Subtype /Link /Border [0 0 0] /H /I /C [1 0 0] /Rect [310.433 134.45 345.029 149.034] /A << /S /GoTo /D <74686D2E31302E322E312E31> >> >>}
\MusicPDFWriteObject{4162}{<< /Type /Annot /Subtype /Link /Border [0 0 0] /H /I /C [1 0 0] /Rect [303.034 81.101 333.247 95.103] /A << /S /GoTo /D <74686D2E392E332E322E31> >> >>}
\MusicPDFWriteObject{4163}{<< /Type /Annot /Subtype /Link /Border [0 0 0] /H /I /C [1 0 0] /Rect [287.856 501.016 318.069 512.969] /A << /S /GoTo /D <6571756174696F6E2E392E332E312E32> >> >>}
\MusicPDFWriteObject{4164}{<< /Type /Annot /Subtype /Link /Border [0 0 0] /H /I /C [1 0 0] /Rect [348.119 399.086 378.332 411.039] /A << /S /GoTo /D <74686D2E342E332E332E31> >> >>}
\MusicPDFWriteObject{4165}{<< /Type /Annot /Subtype /Link /Border [0 0 0] /H /I /C [1 0 0] /Rect [316.632 385.936 351.826 397.888] /A << /S /GoTo /D <6571756174696F6E2E392E322E312E3238> >> >>}
\MusicPDFWriteObject{4166}{<< /Type /Annot /Subtype /Link /Border [0 0 0] /H /I /C [0 1 0] /Rect [305.908 313.495 325.013 322.362] /A << /S /GoTo /D <636974652E4E42> >> >>}
\MusicPDFWriteObject{4167}{<< /Type /Annot /Subtype /Link /Border [0 0 0] /H /I /C [0 1 0] /Rect [317.431 300.344 336.536 309.212] /A << /S /GoTo /D <636974652E4E42> >> >>}
\MusicPDFWriteObject{4168}{<< /Type /Annot /Subtype /Link /Border [0 0 0] /H /I /C [1 0 0] /Rect [70.506 212.104 105.699 220.488] /A << /S /GoTo /D <6E6F746174696F6E2E31302E322E312E31> >> >>}
\MusicPDFWriteObject{4169}{<< /Type /Annot /Subtype /Link /Border [0 0 0] /H /I /C [1 0 0] /Rect [503.877 498.949 539.071 510.902] /A << /S /GoTo /D <6571756174696F6E2E392E322E312E3331> >> >>}
\MusicPDFWriteObject{4170}{<< /Type /Annot /Subtype /Link /Border [0 0 0] /H /I /C [1 0 0] /Rect [298.899 485.799 334.092 497.751] /A << /S /GoTo /D <6571756174696F6E2E31302E322E322E33> >> >>}
\MusicPDFWriteObject{4171}{<< /Type /Annot /Subtype /Link /Border [0 0 0] /H /I /C [1 0 0] /Rect [341.916 425.707 377.109 437.66] /A << /S /GoTo /D <6571756174696F6E2E31302E322E322E34> >> >>}
\MusicPDFWriteObject{4172}{<< /Type /Annot /Subtype /Link /Border [0 0 0] /H /I /C [1 0 0] /Rect [229.646 395.246 265.504 406.037] /A << /S /GoTo /D <74686D2E31302E322E322E31> >> >>}
\MusicPDFWriteObject{4173}{<< /Type /Annot /Subtype /Link /Border [0 0 0] /H /I /C [1 0 0] /Rect [507.302 395.246 543.159 406.037] /A << /S /GoTo /D <74686D2E31302E322E312E31> >> >>}
\MusicPDFWriteObject{4174}{<< /Type /Annot /Subtype /Link /Border [0 0 0] /H /I /C [1 0 0] /Rect [135.414 382.095 165.627 394.399] /A << /S /GoTo /D <74686D2E392E332E322E36> >> >>}
\MusicPDFWriteObject{4175}{<< /Type /Annot /Subtype /Link /Border [0 0 0] /H /I /C [1 0 0] /Rect [245.142 382.095 275.356 394.399] /A << /S /GoTo /D <74686D2E392E332E322E31> >> >>}
\MusicPDFWriteObject{4176}{<< /Type /Annot /Subtype /Link /Border [0 0 0] /H /I /C [1 0 0] /Rect [276.13 382.095 282.093 394.399] /A << /S /GoTo /D <48666F6F746E6F74652E3337> >> >>}
\MusicPDFWriteObject{4177}{<< /Type /Annot /Subtype /Link /Border [0 0 0] /H /I /C [1 0 0] /Rect [336.434 336.112 371.627 348.065] /A << /S /GoTo /D <6571756174696F6E2E31302E322E322E34> >> >>}
\MusicPDFWriteObject{4178}{<< /Type /Annot /Subtype /Link /Border [0 0 0] /H /I /C [1 0 0] /Rect [427.899 335.752 462.428 348.066] /A << /S /GoTo /D <74686D2E31302E322E322E31> >> >>}
\MusicPDFWriteObject{4179}{<< /Type /Annot /Subtype /Link /Border [0 0 0] /H /I /C [1 0 0] /Rect [307.771 299.441 342.964 311.394] /A << /S /GoTo /D <6571756174696F6E2E31302E322E322E33> >> >>}
\MusicPDFWriteObject{4180}{<< /Type /Annot /Subtype /Link /Border [0 0 0] /H /I /C [1 0 0] /Rect [488.442 164.153 518.655 176.105] /A << /S /GoTo /D <6571756174696F6E2E362E332E312E31> >> >>}
\MusicPDFWriteObject{4181}{<< /Type /Annot /Subtype /Link /Border [0 0 0] /H /I /C [1 0 0] /Rect [413.123 137.851 448.316 149.804] /A << /S /GoTo /D <6571756174696F6E2E31302E322E322E35> >> >>}
\MusicPDFWriteObject{4182}{<< /Type /Annot /Subtype /Link /Border [0 0 0] /H /I /C [1 0 0] /Rect [478.008 137.851 513.201 149.804] /A << /S /GoTo /D <6571756174696F6E2E31302E322E322E35> >> >>}
\MusicPDFWriteObject{4183}{<< /Type /Annot /Subtype /Link /Border [0 0 0] /H /I /C [1 0 0] /Rect [251.083 684.633 286.276 696.586] /A << /S /GoTo /D <6571756174696F6E2E31302E322E322E37> >> >>}
\MusicPDFWriteObject{4184}{<< /Type /Annot /Subtype /Link /Border [0 0 0] /H /I /C [1 0 0] /Rect [138.577 601.219 173.836 613.534] /A << /S /GoTo /D <74686D2E392E322E322E3131> >> >>}
\MusicPDFWriteObject{4185}{<< /Type /Annot /Subtype /Link /Border [0 0 0] /H /I /C [1 0 0] /Rect [429.797 556.485 464.99 567.276] /A << /S /GoTo /D <7465726D696E6F6C6F67792E31302E322E312E32> >> >>}
\MusicPDFWriteObject{4186}{<< /Type /Annot /Subtype /Link /Border [0 0 0] /H /I /C [1 0 0] /Rect [488.442 458.062 518.655 470.015] /A << /S /GoTo /D <6571756174696F6E2E362E332E312E31> >> >>}
\MusicPDFWriteObject{4187}{<< /Type /Annot /Subtype /Link /Border [0 0 0] /H /I /C [1 0 0] /Rect [418.992 431.76 454.186 443.713] /A << /S /GoTo /D <6571756174696F6E2E31302E322E322E35> >> >>}
\MusicPDFWriteObject{4188}{<< /Type /Annot /Subtype /Link /Border [0 0 0] /H /I /C [1 0 0] /Rect [483.878 431.76 519.071 443.713] /A << /S /GoTo /D <6571756174696F6E2E31302E322E322E36> >> >>}
\MusicPDFWriteObject{4189}{<< /Type /Annot /Subtype /Link /Border [0 0 0] /H /I /C [1 0 0] /Rect [469.108 416.441 504.301 428.394] /A << /S /GoTo /D <6571756174696F6E2E392E322E312E3331> >> >>}
\MusicPDFWriteObject{4190}{<< /Type /Annot /Subtype /Link /Border [0 0 0] /H /I /C [1 0 0] /Rect [498.228 402.809 533.422 414.762] /A << /S /GoTo /D <6571756174696F6E2E392E322E312E3238> >> >>}
\MusicPDFWriteObject{4191}{<< /Type /Annot /Subtype /Link /Border [0 0 0] /H /I /C [1 0 0] /Rect [307.283 369.312 342.476 381.264] /A << /S /GoTo /D <6571756174696F6E2E31302E322E322E38> >> >>}
\MusicPDFWriteObject{4192}{<< /Type /Annot /Subtype /Link /Border [0 0 0] /H /I /C [1 0 0] /Rect [412.64 368.951 447.833 381.266] /A << /S /GoTo /D <74686D2E31302E322E322E32> >> >>}
\MusicPDFWriteObject{4193}{<< /Type /Annot /Subtype /Link /Border [0 0 0] /H /I /C [1 0 0] /Rect [233.777 355.801 269.07 368.116] /A << /S /GoTo /D <74686D2E31302E322E322E32> >> >>}
\MusicPDFWriteObject{4194}{<< /Type /Annot /Subtype /Link /Border [0 0 0] /H /I /C [1 0 0] /Rect [170.744 341.323 205.937 354.983] /A << /S /GoTo /D <74686D2E31302E322E322E31> >> >>}
\MusicPDFWriteObject{4195}{<< /Type /Annot /Subtype /Link /Border [0 0 0] /H /I /C [1 0 0] /Rect [212.466 300.317 248.323 313.383] /A << /S /GoTo /D <74686D2E31302E322E312E31> >> >>}
\MusicPDFWriteObject{4196}{<< /Type /Annot /Subtype /Link /Border [0 0 0] /H /I /C [0 1 0] /Rect [257.718 278.115 273.71 286.983] /A << /S /GoTo /D <636974652E436F> >> >>}
\MusicPDFWriteObject{4197}{<< /Type /Annot /Subtype /Link /Border [0 0 0] /H /I /C [0 1 0] /Rect [428.477 278.115 444.469 286.983] /A << /S /GoTo /D <636974652E436F> >> >>}
\MusicPDFWriteObject{4198}{<< /Type /Annot /Subtype /Link /Border [0 0 0] /H /I /C [1 0 0] /Rect [331.971 249.324 365.006 261.277] /A << /S /GoTo /D <74686D2E432E332E362E32> >> >>}
\MusicPDFWriteObject{4199}{<< /Type /Annot /Subtype /Link /Border [0 0 0] /H /I /C [0 1 0] /Rect [160.953 225.512 201.562 234.38] /A << /S /GoTo /D <636974652E4B6F7468652D49> >> >>}
\MusicPDFWriteObject{4200}{<< /Type /Annot /Subtype /Link /Border [0 0 0] /H /I /C [0 1 0] /Rect [206.785 687.123 225.758 695.991] /A << /S /GoTo /D <636974652E547265> >> >>}
\MusicPDFWriteObject{4201}{<< /Type /Annot /Subtype /Link /Border [0 0 0] /H /I /C [0 1 0] /Rect [502.66 673.973 543.27 682.84] /A << /S /GoTo /D <636974652E4B6F7468652D49> >> >>}
\MusicPDFWriteObject{4202}{<< /Type /Annot /Subtype /Link /Border [0 0 0] /H /I /C [1 0 0] /Rect [193.354 543.552 210.287 554.401] /A << /S /GoTo /D <73656374696F6E2E432E33> >> >>}
\MusicPDFWriteObject{4203}{<< /Type /Annot /Subtype /Link /Border [0 0 0] /H /I /C [1 0 0] /Rect [235.409 543.552 260.089 554.401] /A << /S /GoTo /D <73756273656374696F6E2E432E332E34> >> >>}
\MusicPDFWriteObject{4204}{<< /Type /Annot /Subtype /Link /Border [0 0 0] /H /I /C [1 0 0] /Rect [283.166 530.402 315.166 541.193] /A << /S /GoTo /D <74686D2E432E332E342E32> >> >>}
\MusicPDFWriteObject{4205}{<< /Type /Annot /Subtype /Link /Border [0 0 0] /H /I /C [1 0 0] /Rect [99.822 393.296 135.016 405.248] /A << /S /GoTo /D <6571756174696F6E2E31302E322E312E34> >> >>}
\MusicPDFWriteObject{4206}{<< /Type /Annot /Subtype /Link /Border [0 0 0] /H /I /C [1 0 0] /Rect [201.317 393.296 237.927 405.248] /A << /S /GoTo /D <6E6F746174696F6E2E31302E322E312E31> >> >>}
\MusicPDFWriteObject{4207}{<< /Type /Annot /Subtype /Link /Border [0 0 0] /H /I /C [1 0 0] /Rect [338.768 320.045 375.378 332.997] /A << /S /GoTo /D <74686D2E31302E322E312E31> >> >>}
\MusicPDFWriteObject{4208}{<< /Type /Annot /Subtype /Link /Border [0 0 0] /H /I /C [1 0 0] /Rect [337.377 592.138 372.57 604.091] /A << /S /GoTo /D <6571756174696F6E2E31302E322E312E34> >> >>}
\MusicPDFWriteObject{4209}{<< /Type /Annot /Subtype /Link /Border [0 0 0] /H /I /C [1 0 0] /Rect [483.12 566.675 518.313 578.628] /A << /S /GoTo /D <6571756174696F6E2E31302E322E312E35> >> >>}
\MusicPDFWriteObject{4210}{<< /Type /Annot /Subtype /Link /Border [0 0 0] /H /I /C [1 0 0] /Rect [147.123 515.634 182.615 528.227] /A << /S /GoTo /D <74686D2E31302E322E312E31> >> >>}
\MusicPDFWriteObject{4211}{<< /Type /Annot /Subtype /Link /Border [0 0 0] /H /I /C [1 0 0] /Rect [383.011 515.635 418.204 527.588] /A << /S /GoTo /D <6571756174696F6E2E31302E332E312E36> >> >>}
\MusicPDFWriteObject{4212}{<< /Type /Annot /Subtype /Link /Border [0 0 0] /H /I /C [1 0 0] /Rect [387.468 502.484 422.661 514.437] /A << /S /GoTo /D <6571756174696F6E2E392E312E312E3131> >> >>}
\MusicPDFWriteObject{4213}{<< /Type /Annot /Subtype /Link /Border [0 0 0] /H /I /C [1 0 0] /Rect [73.712 476.183 108.905 488.135] /A << /S /GoTo /D <6571756174696F6E2E31302E332E312E31> >> >>}
\MusicPDFWriteObject{4214}{<< /Type /Annot /Subtype /Link /Border [0 0 0] /H /I /C [1 0 0] /Rect [73.712 499.306 113.885 511.259] /A << /S /GoTo /D <6571756174696F6E2E31302E332E312E3135> >> >>}
\MusicPDFWriteObject{4215}{<< /Type /Annot /Subtype /Link /Border [0 0 0] /H /I /C [0 1 0] /Rect [87.604 391.202 124.326 402.007] /A << /S /GoTo /D <636974652E4C616E672D62> >> >>}
\MusicPDFWriteObject{4216}{<< /Type /Annot /Subtype /Link /Border [0 0 0] /H /I /C [1 0 0] /Rect [229.618 182.326 267.439 194.279] /A << /S /GoTo /D <6571756174696F6E2E442E332E332E3131> >> >>}
\MusicPDFWriteObject{4217}{<< /Type /Annot /Subtype /Link /Border [0 0 0] /H /I /C [1 0 0] /Rect [327.754 182.326 361.212 194.279] /A << /S /GoTo /D <74686D2E442E332E332E37> >> >>}
\MusicPDFWriteObject{4218}{<< /Type /Annot /Subtype /Link /Border [0 0 0] /H /I /C [1 0 0] /Rect [433.734 182.326 459.29 194.279] /A << /S /GoTo /D <73756273656374696F6E2E442E332E32> >> >>}
\MusicPDFWriteObject{4219}{<< /Type /Annot /Subtype /Link /Border [0 0 0] /H /I /C [1 0 0] /Rect [70.655 169.175 100.727 181.13] /A << /S /GoTo /D <74686D2E342E342E312E31> >> >>}
\MusicPDFWriteObject{4220}{<< /Type /Annot /Subtype /Link /Border [0 0 0] /H /I /C [1 0 0] /Rect [200.791 605.005 240.964 616.958] /A << /S /GoTo /D <6571756174696F6E2E31302E332E312E3137> >> >>}
\MusicPDFWriteObject{4221}{<< /Type /Annot /Subtype /Link /Border [0 0 0] /H /I /C [1 0 0] /Rect [333.173 468.313 368.133 482.896] /A << /S /GoTo /D <7465726D696E6F6C6F67792E31302E322E312E32> >> >>}
\MusicPDFWriteObject{4222}{<< /Type /Annot /Subtype /Link /Border [0 0 0] /H /I /C [1 0 0] /Rect [496.472 428.605 531.665 440.558] /A << /S /GoTo /D <74686D2E31302E322E312E33> >> >>}
\MusicPDFWriteObject{4223}{<< /Type /Annot /Subtype /Link /Border [0 0 0] /H /I /C [1 0 0] /Rect [165.722 415.455 205.895 427.407] /A << /S /GoTo /D <6571756174696F6E2E31302E332E312E3137> >> >>}
\MusicPDFWriteObject{4224}{<< /Type /Annot /Subtype /Link /Border [0 0 0] /H /I /C [0 1 0] /Rect [233.081 417.945 251.28 426.813] /A << /S /GoTo /D <636974652E4A6172> >> >>}
\MusicPDFWriteObject{4225}{<< /Type /Annot /Subtype /Link /Border [0 0 0] /H /I /C [1 0 0] /Rect [471.975 415.455 504.188 427.407] /A << /S /GoTo /D <74686D2E432E332E362E33> >> >>}
\MusicPDFWriteObject{4226}{<< /Type /Annot /Subtype /Link /Border [0 0 0] /H /I /C [1 0 0] /Rect [288.136 402.304 328.31 414.257] /A << /S /GoTo /D <6571756174696F6E2E31302E332E312E3135> >> >>}
\MusicPDFWriteObject{4227}{<< /Type /Annot /Subtype /Link /Border [0 0 0] /H /I /C [1 0 0] /Rect [332.355 339.537 372.528 351.49] /A << /S /GoTo /D <6571756174696F6E2E31302E332E312E3139> >> >>}
\MusicPDFWriteObject{4228}{<< /Type /Annot /Subtype /Link /Border [0 0 0] /H /I /C [1 0 0] /Rect [243.928 214.551 284.102 226.504] /A << /S /GoTo /D <6571756174696F6E2E31302E332E312E3230> >> >>}
\MusicPDFWriteObject{4229}{<< /Type /Annot /Subtype /Link /Border [0 0 0] /H /I /C [1 0 0] /Rect [478.178 214.551 518.352 226.504] /A << /S /GoTo /D <6571756174696F6E2E31302E332E312E3230> >> >>}
\MusicPDFWriteObject{4230}{<< /Type /Annot /Subtype /Link /Border [0 0 0] /H /I /C [1 0 0] /Rect [166.201 188.25 206.374 200.202] /A << /S /GoTo /D <6571756174696F6E2E31302E332E312E3139> >> >>}
\MusicPDFWriteObject{4231}{<< /Type /Annot /Subtype /Link /Border [0 0 0] /H /I /C [1 0 0] /Rect [453.77 125.483 488.963 137.435] /A << /S /GoTo /D <6571756174696F6E2E31302E322E312E36> >> >>}
\MusicPDFWriteObject{4232}{<< /Type /Annot /Subtype /Link /Border [0 0 0] /H /I /C [1 0 0] /Rect [480.848 112.332 521.022 124.284] /A << /S /GoTo /D <6571756174696F6E2E31302E332E312E3139> >> >>}
\MusicPDFWriteObject{4233}{<< /Type /Annot /Subtype /Link /Border [0 0 0] /H /I /C [1 0 0] /Rect [262.668 625.069 302.841 637.022] /A << /S /GoTo /D <6571756174696F6E2E31302E332E312E3232> >> >>}
\MusicPDFWriteObject{4234}{<< /Type /Annot /Subtype /Link /Border [0 0 0] /H /I /C [1 0 0] /Rect [504.203 508.071 539.396 520.024] /A << /S /GoTo /D <6571756174696F6E2E31302E322E312E33> >> >>}
\MusicPDFWriteObject{4235}{<< /Type /Annot /Subtype /Link /Border [0 0 0] /H /I /C [1 0 0] /Rect [313.999 494.921 354.172 506.873] /A << /S /GoTo /D <6571756174696F6E2E31302E332E312E3136> >> >>}
\MusicPDFWriteObject{4236}{<< /Type /Annot /Subtype /Link /Border [0 0 0] /H /I /C [1 0 0] /Rect [410.134 467.009 440.912 480.572] /A << /S /GoTo /D <74686D2E392E332E322E33> >> >>}
\MusicPDFWriteObject{4237}{<< /Type /Annot /Subtype /Link /Border [0 0 0] /H /I /C [1 0 0] /Rect [450.705 455.468 483.131 467.421] /A << /S /GoTo /D <74686D2E432E332E352E32> >> >>}
\MusicPDFWriteObject{4238}{<< /Type /Annot /Subtype /Link /Border [0 0 0] /H /I /C [1 0 0] /Rect [499.165 442.317 539.338 454.27] /A << /S /GoTo /D <6571756174696F6E2E31302E332E312E3233> >> >>}
\MusicPDFWriteObject{4239}{<< /Type /Annot /Subtype /Link /Border [0 0 0] /H /I /C [1 0 0] /Rect [318.565 403.419 348.778 414.209] /A << /S /GoTo /D <74686D2E372E342E312E31> >> >>}
\MusicPDFWriteObject{4240}{<< /Type /Annot /Subtype /Link /Border [0 0 0] /H /I /C [0 1 0] /Rect [141.881 379.054 160.854 387.922] /A << /S /GoTo /D <636974652E547265> >> >>}
\MusicPDFWriteObject{4241}{<< /Type /Annot /Subtype /Link /Border [0 0 0] /H /I /C [1 0 0] /Rect [391.712 350.262 426.905 362.215] /A << /S /GoTo /D <6571756174696F6E2E31302E332E312E36> >> >>}
\MusicPDFWriteObject{4242}{<< /Type /Annot /Subtype /Link /Border [0 0 0] /H /I /C [1 0 0] /Rect [358.573 337.111 393.767 349.064] /A << /S /GoTo /D <6571756174696F6E2E31302E322E312E33> >> >>}
\MusicPDFWriteObject{4243}{<< /Type /Annot /Subtype /Link /Border [0 0 0] /H /I /C [1 0 0] /Rect [457.301 335.424 492.262 349.704] /A << /S /GoTo /D <74686D2E31302E322E312E31> >> >>}
\MusicPDFWriteObject{4244}{<< /Type /Annot /Subtype /Link /Border [0 0 0] /H /I /C [1 0 0] /Rect [159.872 323.961 195.065 335.913] /A << /S /GoTo /D <6571756174696F6E2E31302E322E312E34> >> >>}
\MusicPDFWriteObject{4245}{<< /Type /Annot /Subtype /Link /Border [0 0 0] /H /I /C [1 0 0] /Rect [266.99 283.064 302.183 295.017] /A << /S /GoTo /D <6571756174696F6E2E31302E322E312E36> >> >>}
\MusicPDFWriteObject{4246}{<< /Type /Annot /Subtype /Link /Border [0 0 0] /H /I /C [1 0 0] /Rect [373.011 282.704 408.204 295.018] /A << /S /GoTo /D <74686D2E31302E322E312E31> >> >>}
\MusicPDFWriteObject{4247}{<< /Type /Annot /Subtype /Link /Border [0 0 0] /H /I /C [1 0 0] /Rect [494.648 243.612 529.841 255.564] /A << /S /GoTo /D <6571756174696F6E2E31302E332E312E33> >> >>}
\MusicPDFWriteObject{4248}{<< /Type /Annot /Subtype /Link /Border [0 0 0] /H /I /C [1 0 0] /Rect [417.049 230.461 448.731 242.414] /A << /S /GoTo /D <74686D2E422E342E342E38> >> >>}
\MusicPDFWriteObject{4249}{<< /Type /Annot /Subtype /Link /Border [0 0 0] /H /I /C [1 0 0] /Rect [313.922 217.31 346.764 229.263] /A << /S /GoTo /D <74686D2E442E332E342E32> >> >>}
\MusicPDFWriteObject{4250}{<< /Type /Annot /Subtype /Link /Border [0 0 0] /H /I /C [1 0 0] /Rect [490.882 217.31 523.17 229.263] /A << /S /GoTo /D <74686D2E422E322E322E33> >> >>}
\MusicPDFWriteObject{4251}{<< /Type /Annot /Subtype /Link /Border [0 0 0] /H /I /C [1 0 0] /Rect [147.436 204.159 180.74 216.112] /A << /S /GoTo /D <74686D2E442E312E322E31> >> >>}
\MusicPDFWriteObject{4252}{<< /Type /Annot /Subtype /Link /Border [0 0 0] /H /I /C [1 0 0] /Rect [313.781 204.159 348.974 216.112] /A << /S /GoTo /D <6571756174696F6E2E31302E332E312E33> >> >>}
\MusicPDFWriteObject{4253}{<< /Type /Annot /Subtype /Link /Border [0 0 0] /H /I /C [1 0 0] /Rect [294.181 493.361 329.375 505.314] /A << /S /GoTo /D <6571756174696F6E2E31302E332E312E33> >> >>}
\MusicPDFWriteObject{4254}{<< /Type /Annot /Subtype /Link /Border [0 0 0] /H /I /C [1 0 0] /Rect [356.992 493.361 392.185 505.314] /A << /S /GoTo /D <6571756174696F6E2E31302E332E312E35> >> >>}
\MusicPDFWriteObject{4255}{<< /Type /Annot /Subtype /Link /Border [0 0 0] /H /I /C [0 1 0] /Rect [450.742 495.851 480.084 504.719] /A << /S /GoTo /D <636974652E4C69742D4949> >> >>}
\MusicPDFWriteObject{4256}{<< /Type /Annot /Subtype /Link /Border [0 0 0] /H /I /C [1 0 0] /Rect [86.48 467.06 121.673 479.012] /A << /S /GoTo /D <6571756174696F6E2E31302E332E312E33> >> >>}
\MusicPDFWriteObject{4257}{<< /Type /Annot /Subtype /Link /Border [0 0 0] /H /I /C [1 0 0] /Rect [251.245 440.758 283.945 452.711] /A << /S /GoTo /D <74686D2E432E332E362E35> >> >>}
\MusicPDFWriteObject{4258}{<< /Type /Annot /Subtype /Link /Border [0 0 0] /H /I /C [1 0 0] /Rect [509.153 440.758 541.854 452.711] /A << /S /GoTo /D <74686D2E432E332E342E34> >> >>}
\MusicPDFWriteObject{4259}{<< /Type /Annot /Subtype /Link /Border [0 0 0] /H /I /C [1 0 0] /Rect [87.971 350.283 128.145 362.236] /A << /S /GoTo /D <6571756174696F6E2E31302E332E312E3137> >> >>}
\MusicPDFWriteObject{4260}{<< /Type /Annot /Subtype /Link /Border [0 0 0] /H /I /C [1 0 0] /Rect [162.97 288.691 198.163 300.647] /A << /S /GoTo /D <74686D2E31302E332E312E31> >> >>}
\MusicPDFWriteObject{4261}{<< /Type /Annot /Subtype /Link /Border [0 0 0] /H /I /C [1 0 0] /Rect [99.822 202.662 135.016 214.615] /A << /S /GoTo /D <6571756174696F6E2E31302E322E312E34> >> >>}
\MusicPDFWriteObject{4262}{<< /Type /Annot /Subtype /Link /Border [0 0 0] /H /I /C [1 0 0] /Rect [201.317 202.662 237.927 214.615] /A << /S /GoTo /D <6E6F746174696F6E2E31302E322E312E31> >> >>}
\MusicPDFWriteObject{4263}{<< /Type /Annot /Subtype /Link /Border [0 0 0] /H /I /C [1 0 0] /Rect [348.63 132.388 384.547 145.887] /A << /S /GoTo /D <74686D2E31302E322E322E31> >> >>}
\MusicPDFWriteObject{4264}{<< /Type /Annot /Subtype /Link /Border [0 0 0] /H /I /C [1 0 0] /Rect [337.377 403.093 372.57 415.046] /A << /S /GoTo /D <6571756174696F6E2E31302E322E312E34> >> >>}
\MusicPDFWriteObject{4265}{<< /Type /Annot /Subtype /Link /Border [0 0 0] /H /I /C [1 0 0] /Rect [247.917 318.093 282.447 328.884] /A << /S /GoTo /D <74686D2E31302E332E322E31> >> >>}
\MusicPDFWriteObject{4266}{<< /Type /Annot /Subtype /Link /Border [0 0 0] /H /I /C [1 0 0] /Rect [511.063 318.093 545.592 328.884] /A << /S /GoTo /D <74686D2E31302E332E312E31> >> >>}
\MusicPDFWriteObject{4267}{<< /Type /Annot /Subtype /Link /Border [0 0 0] /H /I /C [1 0 0] /Rect [212.498 304.389 247.691 316.342] /A << /S /GoTo /D <6571756174696F6E2E31302E332E312E34> >> >>}
\MusicPDFWriteObject{4268}{<< /Type /Annot /Subtype /Link /Border [0 0 0] /H /I /C [1 0 0] /Rect [506.695 304.389 541.888 316.342] /A << /S /GoTo /D <6571756174696F6E2E31302E322E322E35> >> >>}
\MusicPDFWriteObject{4269}{<< /Type /Annot /Subtype /Link /Border [0 0 0] /H /I /C [1 0 0] /Rect [94.644 291.239 129.838 303.191] /A << /S /GoTo /D <6571756174696F6E2E31302E322E322E36> >> >>}
\MusicPDFWriteObject{4270}{<< /Type /Annot /Subtype /Link /Border [0 0 0] /H /I /C [1 0 0] /Rect [370.891 221.622 406.051 232.414] /A << /S /GoTo /D <74686D2E31302E332E312E31> >> >>}
\MusicPDFWriteObject{4271}{<< /Type /Annot /Subtype /Link /Border [0 0 0] /H /I /C [1 0 0] /Rect [73.712 207.919 108.905 219.872] /A << /S /GoTo /D <6571756174696F6E2E31302E332E312E31> >> >>}
\MusicPDFWriteObject{4272}{<< /Type /Annot /Subtype /Link /Border [0 0 0] /H /I /C [1 0 0] /Rect [373.205 194.768 408.399 206.721] /A << /S /GoTo /D <6571756174696F6E2E31302E332E322E31> >> >>}
\MusicPDFWriteObject{4273}{<< /Type /Annot /Subtype /Link /Border [0 0 0] /H /I /C [1 0 0] /Rect [193.298 168.467 228.492 180.419] /A << /S /GoTo /D <6571756174696F6E2E31302E332E322E36> >> >>}
\MusicPDFWriteObject{4274}{<< /Type /Annot /Subtype /Link /Border [0 0 0] /H /I /C [1 0 0] /Rect [169.056 497.516 204.25 509.468] /A << /S /GoTo /D <6571756174696F6E2E31302E332E322E38> >> >>}
\MusicPDFWriteObject{4275}{<< /Type /Annot /Subtype /Link /Border [0 0 0] /H /I /C [0 1 0] /Rect [219.492 451.024 255.101 461.829] /A << /S /GoTo /D <636974652E4C616E672D65> >> >>}
\MusicPDFWriteObject{4276}{<< /Type /Annot /Subtype /Link /Border [0 0 0] /H /I /C [1 0 0] /Rect [283.319 411.019 318.512 422.971] /A << /S /GoTo /D <6571756174696F6E2E31302E332E322E37> >> >>}
\MusicPDFWriteObject{4277}{<< /Type /Annot /Subtype /Link /Border [0 0 0] /H /I /C [1 0 0] /Rect [111.979 398.421 146.542 409.212] /A << /S /GoTo /D <74686D2E31302E332E312E31> >> >>}
\MusicPDFWriteObject{4278}{<< /Type /Annot /Subtype /Link /Border [0 0 0] /H /I /C [1 0 0] /Rect [400.526 398.421 435.088 409.212] /A << /S /GoTo /D <74686D2E31302E332E312E31> >> >>}
\MusicPDFWriteObject{4279}{<< /Type /Annot /Subtype /Link /Border [0 0 0] /H /I /C [1 0 0] /Rect [70.506 387.207 105.699 395.924] /A << /S /GoTo /D <74686D2E31302E332E322E31> >> >>}
\MusicPDFWriteObject{4280}{<< /Type /Annot /Subtype /Link /Border [0 0 0] /H /I /C [1 0 0] /Rect [197.375 288.059 227.024 300.014] /A << /S /GoTo /D <74686D2E392E312E322E36> >> >>}
\MusicPDFWriteObject{4281}{<< /Type /Annot /Subtype /Link /Border [0 0 0] /H /I /C [1 0 0] /Rect [214.982 95.493 250.839 107.701] /A << /S /GoTo /D <74686D2E31302E342E322E31> >> >>}
\MusicPDFWriteObject{4282}{<< /Type /Annot /Subtype /Link /Border [0 0 0] /H /I /C [1 0 0] /Rect [224.308 83.149 259.501 93.94] /A << /S /GoTo /D <74686D2E31302E342E322E34> >> >>}
\MusicPDFWriteObject{4283}{<< /Type /Annot /Subtype /Link /Border [0 0 0] /H /I /C [1 0 0] /Rect [256.005 685.187 291.862 695.977] /A << /S /GoTo /D <74686D2E31302E342E322E32> >> >>}
\MusicPDFWriteObject{4284}{<< /Type /Annot /Subtype /Link /Border [0 0 0] /H /I /C [1 0 0] /Rect [292.206 658.331 322.419 670.286] /A << /S /GoTo /D <74686D2E392E332E332E31> >> >>}
\MusicPDFWriteObject{4285}{<< /Type /Annot /Subtype /Link /Border [0 0 0] /H /I /C [1 0 0] /Rect [389.456 658.331 419.669 670.286] /A << /S /GoTo /D <74686D2E392E332E332E33> >> >>}
\MusicPDFWriteObject{4286}{<< /Type /Annot /Subtype /Link /Border [0 0 0] /H /I /C [1 0 0] /Rect [210.569 570.25 240.472 584.275] /A << /S /GoTo /D <74686D2E392E332E312E31> >> >>}
\MusicPDFWriteObject{4287}{<< /Type /Annot /Subtype /Link /Border [0 0 0] /H /I /C [1 0 0] /Rect [483.244 135.199 518.437 147.152] /A << /S /GoTo /D <6571756174696F6E2E31302E342E312E34> >> >>}
\MusicPDFWriteObject{4288}{<< /Type /Annot /Subtype /Link /Border [0 0 0] /H /I /C [0 1 0] /Rect [73.77 660.822 92.875 669.69] /A << /S /GoTo /D <636974652E4E42> >> >>}
\MusicPDFWriteObject{4289}{<< /Type /Annot /Subtype /Link /Border [0 0 0] /H /I /C [1 0 0] /Rect [240.995 620.783 271.772 632.738] /A << /S /GoTo /D <74686D2E392E332E332E31> >> >>}
\MusicPDFWriteObject{4290}{<< /Type /Annot /Subtype /Link /Border [0 0 0] /H /I /C [1 0 0] /Rect [341.581 620.783 372.359 632.738] /A << /S /GoTo /D <74686D2E392E332E332E33> >> >>}
\MusicPDFWriteObject{4291}{<< /Type /Annot /Subtype /Link /Border [0 0 0] /H /I /C [1 0 0] /Rect [172.219 258.365 202.997 270.573] /A << /S /GoTo /D <74686D2E392E332E312E35> >> >>}
\MusicPDFWriteObject{4292}{<< /Type /Annot /Subtype /Link /Border [0 0 0] /H /I /C [1 0 0] /Rect [311.671 219.166 346.865 231.119] /A << /S /GoTo /D <6571756174696F6E2E392E322E312E3238> >> >>}
\MusicPDFWriteObject{4293}{<< /Type /Annot /Subtype /Link /Border [0 0 0] /H /I /C [1 0 0] /Rect [244.726 179.714 274.939 191.667] /A << /S /GoTo /D <74686D2E392E312E322E32> >> >>}
\MusicPDFWriteObject{4294}{<< /Type /Annot /Subtype /Link /Border [0 0 0] /H /I /C [0 1 0] /Rect [416.429 182.204 433.583 191.072] /A << /S /GoTo /D <636974652E466F6C> >> >>}
\MusicPDFWriteObject{4295}{<< /Type /Annot /Subtype /Link /Border [0 0 0] /H /I /C [1 0 0] /Rect [448.832 153.413 478.508 165.365] /A << /S /GoTo /D <74686D2E392E312E322E32> >> >>}
\MusicPDFWriteObject{4296}{<< /Type /Annot /Subtype /Link /Border [0 0 0] /H /I /C [1 0 0] /Rect [107.437 140.815 138.215 151.606] /A << /S /GoTo /D <74686D2E392E332E312E35> >> >>}
\MusicPDFWriteObject{4297}{<< /Type /Annot /Subtype /Link /Border [0 0 0] /H /I /C [1 0 0] /Rect [189.613 127.111 219.826 139.064] /A << /S /GoTo /D <74686D2E392E332E312E35> >> >>}
\MusicPDFWriteObject{4298}{<< /Type /Annot /Subtype /Link /Border [0 0 0] /H /I /C [1 0 0] /Rect [119.58 608.753 149.793 620.706] /A << /S /GoTo /D <6571756174696F6E2E392E312E322E33> >> >>}
\MusicPDFWriteObject{4299}{<< /Type /Annot /Subtype /Link /Border [0 0 0] /H /I /C [1 0 0] /Rect [345.855 582.452 386.028 594.405] /A << /S /GoTo /D <6571756174696F6E2E31302E342E312E3130> >> >>}
\MusicPDFWriteObject{4300}{<< /Type /Annot /Subtype /Link /Border [0 0 0] /H /I /C [1 0 0] /Rect [209.202 267.517 244.395 278.307] /A << /S /GoTo /D <74686D2E31302E332E322E31> >> >>}
\MusicPDFWriteObject{4301}{<< /Type /Annot /Subtype /Link /Border [0 0 0] /H /I /C [1 0 0] /Rect [161.443 234.356 196.636 246.308] /A << /S /GoTo /D <6571756174696F6E2E31302E342E312E39> >> >>}
\MusicPDFWriteObject{4302}{<< /Type /Annot /Subtype /Link /Border [0 0 0] /H /I /C [1 0 0] /Rect [401.778 234.355 431.455 246.31] /A << /S /GoTo /D <74686D2E392E332E332E31> >> >>}
\MusicPDFWriteObject{4303}{<< /Type /Annot /Subtype /Link /Border [0 0 0] /H /I /C [1 0 0] /Rect [497.382 234.355 527.058 246.31] /A << /S /GoTo /D <74686D2E392E332E332E33> >> >>}
\MusicPDFWriteObject{4304}{<< /Type /Annot /Subtype /Link /Border [0 0 0] /H /I /C [1 0 0] /Rect [331.426 194.543 361.64 206.858] /A << /S /GoTo /D <74686D2E392E332E332E35> >> >>}
\MusicPDFWriteObject{4305}{<< /Type /Annot /Subtype /Link /Border [0 0 0] /H /I /C [1 0 0] /Rect [486.176 645.181 521.37 657.134] /A << /S /GoTo /D <6571756174696F6E2E31302E342E312E37> >> >>}
\MusicPDFWriteObject{4306}{<< /Type /Annot /Subtype /Link /Border [0 0 0] /H /I /C [1 0 0] /Rect [157.03 632.03 192.223 643.983] /A << /S /GoTo /D <6571756174696F6E2E31302E342E312E39> >> >>}
\MusicPDFWriteObject{4307}{<< /Type /Annot /Subtype /Link /Border [0 0 0] /H /I /C [1 0 0] /Rect [245.116 576.879 280.309 588.832] /A << /S /GoTo /D <6571756174696F6E2E31302E342E312E38> >> >>}
\MusicPDFWriteObject{4308}{<< /Type /Annot /Subtype /Link /Border [0 0 0] /H /I /C [1 0 0] /Rect [486.346 527.301 521.539 539.254] /A << /S /GoTo /D <6571756174696F6E2E31302E322E312E31> >> >>}
\MusicPDFWriteObject{4309}{<< /Type /Annot /Subtype /Link /Border [0 0 0] /H /I /C [1 0 0] /Rect [112.727 514.15 147.92 526.103] /A << /S /GoTo /D <74686D2E31302E332E312E31> >> >>}
\MusicPDFWriteObject{4310}{<< /Type /Annot /Subtype /Link /Border [0 0 0] /H /I /C [1 0 0] /Rect [225.543 451.421 260.736 463.374] /A << /S /GoTo /D <6571756174696F6E2E31302E332E312E32> >> >>}
\MusicPDFWriteObject{4311}{<< /Type /Annot /Subtype /Link /Border [0 0 0] /H /I /C [1 0 0] /Rect [319.362 451.42 354.556 463.376] /A << /S /GoTo /D <74686D2E31302E332E312E31> >> >>}
\MusicPDFWriteObject{4312}{<< /Type /Annot /Subtype /Link /Border [0 0 0] /H /I /C [1 0 0] /Rect [155.56 410.038 190.754 421.991] /A << /S /GoTo /D <6571756174696F6E2E31302E342E322E35> >> >>}
\MusicPDFWriteObject{4313}{<< /Type /Annot /Subtype /Link /Border [0 0 0] /H /I /C [1 0 0] /Rect [229.361 334.103 264.554 346.056] /A << /S /GoTo /D <6571756174696F6E2E31302E342E322E36> >> >>}
\MusicPDFWriteObject{4314}{<< /Type /Annot /Subtype /Link /Border [0 0 0] /H /I /C [1 0 0] /Rect [203.94 247.937 239.133 259.89] /A << /S /GoTo /D <6571756174696F6E2E31302E342E322E35> >> >>}
\MusicPDFWriteObject{4315}{<< /Type /Annot /Subtype /Link /Border [0 0 0] /H /I /C [1 0 0] /Rect [88.849 221.222 124.042 233.174] /A << /S /GoTo /D <6571756174696F6E2E31302E342E322E37> >> >>}
\MusicPDFWriteObject{4316}{<< /Type /Annot /Subtype /Link /Border [0 0 0] /H /I /C [1 0 0] /Rect [149.166 208.07 185.023 220.025] /A << /S /GoTo /D <74686D2E31302E332E312E31> >> >>}
\MusicPDFWriteObject{4317}{<< /Type /Annot /Subtype /Link /Border [0 0 0] /H /I /C [0 1 0] /Rect [319.371 210.561 354.135 219.429] /A << /S /GoTo /D <636974652E436E772D62> >> >>}
\MusicPDFWriteObject{4318}{<< /Type /Annot /Subtype /Link /Border [0 0 0] /H /I /C [1 0 0] /Rect [71.004 193.591 103.43 206.875] /A << /S /GoTo /D <74686D2E432E332E342E34> >> >>}
\MusicPDFWriteObject{4319}{<< /Type /Annot /Subtype /Link /Border [0 0 0] /H /I /C [1 0 0] /Rect [128.02 167.29 158.797 180.573] /A << /S /GoTo /D <74686D2E392E332E332E35> >> >>}
\MusicPDFWriteObject{4320}{<< /Type /Annot /Subtype /Link /Border [0 0 0] /H /I /C [1 0 0] /Rect [71.004 155.468 103.43 167.421] /A << /S /GoTo /D <74686D2E432E332E352E31> >> >>}
\MusicPDFWriteObject{4321}{<< /Type /Annot /Subtype /Link /Border [0 0 0] /H /I /C [1 0 0] /Rect [130.212 155.468 165.405 167.421] /A << /S /GoTo /D <6571756174696F6E2E31302E342E322E33> >> >>}
\MusicPDFWriteObject{4322}{<< /Type /Annot /Subtype /Link /Border [0 0 0] /H /I /C [0 1 0] /Rect [309.204 647.356 328.308 656.223] /A << /S /GoTo /D <636974652E4E42> >> >>}
\MusicPDFWriteObject{4323}{<< /Type /Annot /Subtype /Link /Border [0 0 0] /H /I /C [0 1 0] /Rect [317.732 634.205 336.837 643.073] /A << /S /GoTo /D <636974652E4E42> >> >>}
\MusicPDFWriteObject{4324}{<< /Type /Annot /Subtype /Link /Border [0 0 0] /H /I /C [1 0 0] /Rect [174.621 618.509 209.815 630.462] /A << /S /GoTo /D <6571756174696F6E2E31302E322E312E33> >> >>}
\MusicPDFWriteObject{4325}{<< /Type /Annot /Subtype /Link /Border [0 0 0] /H /I /C [1 0 0] /Rect [503.53 579.057 538.723 591.009] /A << /S /GoTo /D <74686D2E31302E342E322E38> >> >>}
\MusicPDFWriteObject{4326}{<< /Type /Annot /Subtype /Link /Border [0 0 0] /H /I /C [1 0 0] /Rect [328.085 493.543 363.278 505.496] /A << /S /GoTo /D <6571756174696F6E2E31302E332E312E32> >> >>}
\MusicPDFWriteObject{4327}{<< /Type /Annot /Subtype /Link /Border [0 0 0] /H /I /C [1 0 0] /Rect [87.218 480.392 122.411 492.345] /A << /S /GoTo /D <6571756174696F6E2E31302E332E312E33> >> >>}
\MusicPDFWriteObject{4328}{<< /Type /Annot /Subtype /Link /Border [0 0 0] /H /I /C [1 0 0] /Rect [180.926 480.392 217.536 492.345] /A << /S /GoTo /D <74686D2E31302E332E312E31> >> >>}
\MusicPDFWriteObject{4329}{<< /Type /Annot /Subtype /Link /Border [0 0 0] /H /I /C [1 0 0] /Rect [220.868 458.161 256.725 468.952] /A << /S /GoTo /D <74686D2E31302E342E322E32> >> >>}
\MusicPDFWriteObject{4330}{<< /Type /Annot /Subtype /Link /Border [0 0 0] /H /I /C [1 0 0] /Rect [222.013 431.307 257.206 443.259] /A << /S /GoTo /D <6571756174696F6E2E31302E342E322E33> >> >>}
\MusicPDFWriteObject{4331}{<< /Type /Annot /Subtype /Link /Border [0 0 0] /H /I /C [1 0 0] /Rect [318.5 431.307 354.357 443.259] /A << /S /GoTo /D <74686D2E31302E342E322E31> >> >>}
\MusicPDFWriteObject{4332}{<< /Type /Annot /Subtype /Link /Border [0 0 0] /H /I /C [1 0 0] /Rect [392.373 342.276 432.546 354.229] /A << /S /GoTo /D <6571756174696F6E2E31302E342E322E3130> >> >>}
\MusicPDFWriteObject{4333}{<< /Type /Annot /Subtype /Link /Border [0 0 0] /H /I /C [1 0 0] /Rect [370.041 329.679 405.234 340.469] /A << /S /GoTo /D <74686D2E31302E342E322E32> >> >>}
\MusicPDFWriteObject{4334}{<< /Type /Annot /Subtype /Link /Border [0 0 0] /H /I /C [1 0 0] /Rect [468.011 329.679 503.204 340.469] /A << /S /GoTo /D <74686D2E31302E332E312E31> >> >>}
\MusicPDFWriteObject{4335}{<< /Type /Annot /Subtype /Link /Border [0 0 0] /H /I /C [1 0 0] /Rect [221.961 306.894 257.154 317.685] /A << /S /GoTo /D <74686D2E31302E342E322E32> >> >>}
\MusicPDFWriteObject{4336}{<< /Type /Annot /Subtype /Link /Border [0 0 0] /H /I /C [1 0 0] /Rect [222.931 242.494 259.541 256.059] /A << /S /GoTo /D <74686D2E31302E332E312E31> >> >>}
\MusicPDFWriteObject{4337}{<< /Type /Annot /Subtype /Link /Border [0 0 0] /H /I /C [1 0 0] /Rect [147.549 184.892 187.723 196.845] /A << /S /GoTo /D <6571756174696F6E2E31302E342E312E3131> >> >>}
\MusicPDFWriteObject{4338}{<< /Type /Annot /Subtype /Link /Border [0 0 0] /H /I /C [1 0 0] /Rect [252.785 184.892 287.314 196.847] /A << /S /GoTo /D <74686D2E31302E342E312E34> >> >>}
\MusicPDFWriteObject{4339}{<< /Type /Annot /Subtype /Link /Border [0 0 0] /H /I /C [1 0 0] /Rect [124.18 121.909 157.215 134.118] /A << /S /GoTo /D <74686D2E432E332E322E32> >> >>}
\MusicPDFWriteObject{4340}{<< /Type /Annot /Subtype /Link /Border [0 0 0] /H /I /C [1 0 0] /Rect [420.331 95.862 455.524 107.814] /A << /S /GoTo /D <6571756174696F6E2E31302E342E312E33> >> >>}
\MusicPDFWriteObject{4341}{<< /Type /Annot /Subtype /Link /Border [0 0 0] /H /I /C [1 0 0] /Rect [313.202 82.711 348.395 94.664] /A << /S /GoTo /D <6571756174696F6E2E31302E342E312E35> >> >>}
\MusicPDFWriteObject{4342}{<< /Type /Annot /Subtype /Link /Border [0 0 0] /H /I /C [1 0 0] /Rect [400.492 81.101 435.685 94.664] /A << /S /GoTo /D <74686D2E31302E342E312E32> >> >>}
\MusicPDFWriteObject{4343}{<< /Type /Annot /Subtype /Link /Border [0 0 0] /H /I /C [1 0 0] /Rect [121.911 562.2 162.085 574.152] /A << /S /GoTo /D <6571756174696F6E2E31302E342E322E3131> >> >>}
\MusicPDFWriteObject{4344}{<< /Type /Annot /Subtype /Link /Border [0 0 0] /H /I /C [1 0 0] /Rect [187.794 542.773 226.861 551.627] /A << /S /GoTo /D <74686D2E31302E342E322E32> >> >>}
\MusicPDFWriteObject{4345}{<< /Type /Annot /Subtype /Link /Border [0 0 0] /H /I /C [1 0 0] /Rect [449.54 540.283 484.301 552.238] /A << /S /GoTo /D <74686D2E31302E342E322E31> >> >>}
\MusicPDFWriteObject{4346}{<< /Type /Annot /Subtype /Link /Border [0 0 0] /H /I /C [1 0 0] /Rect [460.403 525.523 495.596 539.087] /A << /S /GoTo /D <74686D2E31302E332E322E31> >> >>}
\MusicPDFWriteObject{4347}{<< /Type /Annot /Subtype /Link /Border [0 0 0] /H /I /C [1 0 0] /Rect [383.808 474.169 414.021 485.874] /A << /S /GoTo /D <74686D2E392E332E332E32> >> >>}
\MusicPDFWriteObject{4348}{<< /Type /Annot /Subtype /Link /Border [0 0 0] /H /I /C [1 0 0] /Rect [240.678 460.05 275.24 472.723] /A << /S /GoTo /D <74686D2E31302E342E312E33> >> >>}
\MusicPDFWriteObject{4349}{<< /Type /Annot /Subtype /Link /Border [0 0 0] /H /I /C [1 0 0] /Rect [511.419 460.05 543.267 472.723] /A << /S /GoTo /D <74686D2E432E332E352E35> >> >>}
\MusicPDFWriteObject{4350}{<< /Type /Annot /Subtype /Link /Border [0 0 0] /H /I /C [1 0 0] /Rect [506.695 448.228 541.888 460.181] /A << /S /GoTo /D <6571756174696F6E2E31302E342E322E39> >> >>}
\MusicPDFWriteObject{4351}{<< /Type /Annot /Subtype /Link /Border [0 0 0] /H /I /C [1 0 0] /Rect [437.594 631.732 472.123 643.985] /A << /S /GoTo /D <74686D2E31302E342E322E35> >> >>}
\MusicPDFWriteObject{4352}{<< /Type /Annot /Subtype /Link /Border [0 0 0] /H /I /C [1 0 0] /Rect [176.323 599.046 215.39 607.9] /A << /S /GoTo /D <74686D2E31302E342E322E35> >> >>}
\MusicPDFWriteObject{4353}{<< /Type /Annot /Subtype /Link /Border [0 0 0] /H /I /C [1 0 0] /Rect [396.269 599.046 402.232 609.413] /A << /S /GoTo /D <48666F6F746E6F74652E3338> >> >>}
\MusicPDFWriteObject{4354}{<< /Type /Annot /Subtype /Link /Border [0 0 0] /H /I /C [0 1 0] /Rect [73.782 533.292 92.886 542.16] /A << /S /GoTo /D <636974652E4E42> >> >>}
\MusicPDFWriteObject{4355}{<< /Type /Annot /Subtype /Link /Border [0 0 0] /H /I /C [0 1 0] /Rect [197.598 533.292 216.703 542.16] /A << /S /GoTo /D <636974652E4E42> >> >>}
\MusicPDFWriteObject{4356}{<< /Type /Annot /Subtype /Link /Border [0 0 0] /H /I /C [0 1 0] /Rect [297.631 493.839 316.735 502.707] /A << /S /GoTo /D <636974652E4E42> >> >>}
\MusicPDFWriteObject{4357}{<< /Type /Annot /Subtype /Link /Border [0 0 0] /H /I /C [0 1 0] /Rect [445.51 440.525 465.23 449.393] /A << /S /GoTo /D <636974652E4B616C> >> >>}
\MusicPDFWriteObject{4358}{<< /Type /Annot /Subtype /Link /Border [0 0 0] /H /I /C [1 0 0] /Rect [73.712 398.583 113.885 410.536] /A << /S /GoTo /D <6571756174696F6E2E31302E342E322E3137> >> >>}
\MusicPDFWriteObject{4359}{<< /Type /Annot /Subtype /Link /Border [0 0 0] /H /I /C [1 0 0] /Rect [73.712 359.131 113.885 371.083] /A << /S /GoTo /D <6571756174696F6E2E31302E342E322E3137> >> >>}
\MusicPDFWriteObject{4360}{<< /Type /Annot /Subtype /Link /Border [0 0 0] /H /I /C [0 1 0] /Rect [498.938 206.836 517.136 215.704] /A << /S /GoTo /D <636974652E4A6172> >> >>}
\MusicPDFWriteObject{4361}{<< /Type /Annot /Subtype /Link /Border [0 0 0] /H /I /C [0 1 0] /Rect [326.955 193.686 345.153 202.553] /A << /S /GoTo /D <636974652E4A6172> >> >>}
\MusicPDFWriteObject{4362}{<< /Type /Annot /Subtype /Link /Border [0 0 0] /H /I /C [0 1 0] /Rect [157.418 139.146 181.26 149.95] /A << /S /GoTo /D <636974652E4D6567> >> >>}
\MusicPDFWriteObject{4363}{<< /Type /Annot /Subtype /Link /Border [0 0 0] /H /I /C [0 1 0] /Rect [350.808 139.146 374.651 149.95] /A << /S /GoTo /D <636974652E4D6567> >> >>}
\MusicPDFWriteObject{4364}{<< /Type /Annot /Subtype /Link /Border [0 0 0] /H /I /C [1 0 0] /Rect [397.725 684.633 433.582 696.588] /A << /S /GoTo /D <74686D2E31302E342E322E35> >> >>}
\MusicPDFWriteObject{4365}{<< /Type /Annot /Subtype /Link /Border [0 0 0] /H /I /C [0 1 0] /Rect [362.9 660.822 381.099 669.69] /A << /S /GoTo /D <636974652E4A6172> >> >>}
\MusicPDFWriteObject{4366}{<< /Type /Annot /Subtype /Link /Border [0 0 0] /H /I /C [1 0 0] /Rect [328.085 343.471 363.278 355.424] /A << /S /GoTo /D <6571756174696F6E2E31302E332E322E32> >> >>}
\MusicPDFWriteObject{4367}{<< /Type /Annot /Subtype /Link /Border [0 0 0] /H /I /C [1 0 0] /Rect [87.218 330.321 122.411 342.273] /A << /S /GoTo /D <6571756174696F6E2E31302E332E322E33> >> >>}
\MusicPDFWriteObject{4368}{<< /Type /Annot /Subtype /Link /Border [0 0 0] /H /I /C [1 0 0] /Rect [180.926 330.321 217.536 342.273] /A << /S /GoTo /D <74686D2E31302E332E322E31> >> >>}
\MusicPDFWriteObject{4369}{<< /Type /Annot /Subtype /Link /Border [0 0 0] /H /I /C [1 0 0] /Rect [181.278 257.743 216.471 269.696] /A << /S /GoTo /D <6571756174696F6E2E31302E322E322E34> >> >>}
\MusicPDFWriteObject{4370}{<< /Type /Annot /Subtype /Link /Border [0 0 0] /H /I /C [1 0 0] /Rect [318.757 257.743 353.286 269.696] /A << /S /GoTo /D <74686D2E31302E322E322E31> >> >>}
\MusicPDFWriteObject{4371}{<< /Type /Annot /Subtype /Link /Border [0 0 0] /H /I /C [1 0 0] /Rect [486.865 145.435 522.722 157.39] /A << /S /GoTo /D <74686D2E31302E322E322E32> >> >>}
\MusicPDFWriteObject{4372}{<< /Type /Annot /Subtype /Link /Border [0 0 0] /H /I /C [1 0 0] /Rect [280.649 82.596 320.822 94.549] /A << /S /GoTo /D <6571756174696F6E2E31302E342E322E3233> >> >>}
\MusicPDFWriteObject{4373}{<< /Type /Annot /Subtype /Link /Border [0 0 0] /H /I /C [0 1 0] /Rect [185.564 650.696 220.329 659.564] /A << /S /GoTo /D <636974652E436E772D62> >> >>}
\MusicPDFWriteObject{4374}{<< /Type /Annot /Subtype /Link /Border [0 0 0] /H /I /C [1 0 0] /Rect [379 584.259 414.193 598.926] /A << /S /GoTo /D <74686D2E31302E342E322E35> >> >>}
\MusicPDFWriteObject{4375}{<< /Type /Annot /Subtype /Link /Border [0 0 0] /H /I /C [1 0 0] /Rect [235.148 571.468 275.321 583.421] /A << /S /GoTo /D <6571756174696F6E2E31302E342E322E3236> >> >>}
\MusicPDFWriteObject{4376}{<< /Type /Annot /Subtype /Link /Border [0 0 0] /H /I /C [1 0 0] /Rect [415.013 571.468 455.186 583.421] /A << /S /GoTo /D <6571756174696F6E2E31302E342E322E3236> >> >>}
\MusicPDFWriteObject{4377}{<< /Type /Annot /Subtype /Link /Border [0 0 0] /H /I /C [0 1 0] /Rect [192.444 520.151 209.598 529.018] /A << /S /GoTo /D <636974652E466F6C> >> >>}
\MusicPDFWriteObject{4378}{<< /Type /Annot /Subtype /Link /Border [0 0 0] /H /I /C [1 0 0] /Rect [192.491 504.51 232.664 516.463] /A << /S /GoTo /D <6571756174696F6E2E31302E342E322E3237> >> >>}
\MusicPDFWriteObject{4379}{<< /Type /Annot /Subtype /Link /Border [0 0 0] /H /I /C [1 0 0] /Rect [170.166 298.474 200.379 310.427] /A << /S /GoTo /D <6571756174696F6E2E392E332E312E36> >> >>}
\MusicPDFWriteObject{4380}{<< /Type /Annot /Subtype /Link /Border [0 0 0] /H /I /C [1 0 0] /Rect [263.813 298.474 293.546 310.427] /A << /S /GoTo /D <74686D2E392E332E312E35> >> >>}
\MusicPDFWriteObject{4381}{<< /Type /Annot /Subtype /Link /Border [0 0 0] /H /I /C [1 0 0] /Rect [367.693 298.474 407.867 310.427] /A << /S /GoTo /D <6571756174696F6E2E31302E342E322E3239> >> >>}
\MusicPDFWriteObject{4382}{<< /Type /Annot /Subtype /Link /Border [0 0 0] /H /I /C [1 0 0] /Rect [328.217 272.173 368.391 284.126] /A << /S /GoTo /D <6571756174696F6E2E31302E342E322E3235> >> >>}
\MusicPDFWriteObject{4383}{<< /Type /Annot /Subtype /Link /Border [0 0 0] /H /I /C [1 0 0] /Rect [247.112 211.665 282.206 222.456] /A << /S /GoTo /D <74686D2E31302E332E322E31> >> >>}
\MusicPDFWriteObject{4384}{<< /Type /Annot /Subtype /Link /Border [0 0 0] /H /I /C [1 0 0] /Rect [190.238 197.961 230.411 209.914] /A << /S /GoTo /D <6571756174696F6E2E31302E342E322E3235> >> >>}
\MusicPDFWriteObject{4385}{<< /Type /Annot /Subtype /Link /Border [0 0 0] /H /I /C [1 0 0] /Rect [499.843 184.81 505.806 197.668] /A << /S /GoTo /D <48666F6F746E6F74652E3339> >> >>}
\MusicPDFWriteObject{4386}{<< /Type /Annot /Subtype /Link /Border [0 0 0] /H /I /C [1 0 0] /Rect [219.108 459.599 249.886 470.39] /A << /S /GoTo /D <74686D2E362E332E322E31> >> >>}
\MusicPDFWriteObject{4387}{<< /Type /Annot /Subtype /Link /Border [0 0 0] /H /I /C [1 0 0] /Rect [425.26 445.895 460.453 457.848] /A << /S /GoTo /D <6571756174696F6E2E31302E322E312E31> >> >>}
\MusicPDFWriteObject{4388}{<< /Type /Annot /Subtype /Link /Border [0 0 0] /H /I /C [1 0 0] /Rect [506.695 418.431 541.888 430.384] /A << /S /GoTo /D <6571756174696F6E2E352E322E332E3134> >> >>}
\MusicPDFWriteObject{4389}{<< /Type /Annot /Subtype /Link /Border [0 0 0] /H /I /C [1 0 0] /Rect [70.655 379.532 100.304 390.323] /A << /S /GoTo /D <74686D2E342E342E312E31> >> >>}
\MusicPDFWriteObject{4390}{<< /Type /Annot /Subtype /Link /Border [0 0 0] /H /I /C [1 0 0] /Rect [159.895 298.027 189.996 312.029] /A << /S /GoTo /D <74686D2E362E332E322E31> >> >>}
\MusicPDFWriteObject{4391}{<< /Type /Annot /Subtype /Link /Border [0 0 0] /H /I /C [1 0 0] /Rect [449.363 273.772 482.821 285.728] /A << /S /GoTo /D <74686D2E442E332E332E37> >> >>}
\MusicPDFWriteObject{4392}{<< /Type /Annot /Subtype /Link /Border [0 0 0] /H /I /C [1 0 0] /Rect [71.004 221.17 103.845 233.123] /A << /S /GoTo /D <74686D2E442E332E312E32> >> >>}
\MusicPDFWriteObject{4393}{<< /Type /Annot /Subtype /Link /Border [0 0 0] /H /I /C [1 0 0] /Rect [166.622 221.17 204.443 233.123] /A << /S /GoTo /D <74686D2E442E332E332E3133> >> >>}
\MusicPDFWriteObject{4394}{<< /Type /Annot /Subtype /Link /Border [0 0 0] /H /I /C [1 0 0] /Rect [144.736 208.019 174.385 220.611] /A << /S /GoTo /D <74686D2E362E332E322E31> >> >>}
\MusicPDFWriteObject{4395}{<< /Type /Annot /Subtype /Link /Border [0 0 0] /H /I /C [1 0 0] /Rect [352.061 208.019 387.254 219.972] /A << /S /GoTo /D <6571756174696F6E2E31302E342E332E32> >> >>}
\MusicPDFWriteObject{4396}{<< /Type /Annot /Subtype /Link /Border [0 0 0] /H /I /C [1 0 0] /Rect [278.181 194.869 313.375 206.821] /A << /S /GoTo /D <6571756174696F6E2E31302E342E332E32> >> >>}
\MusicPDFWriteObject{4397}{<< /Type /Annot /Subtype /Link /Border [0 0 0] /H /I /C [1 0 0] /Rect [211.372 180.03 241.021 194.98] /A << /S /GoTo /D <74686D2E372E312E342E32> >> >>}
\MusicPDFWriteObject{4398}{<< /Type /Annot /Subtype /Link /Border [0 0 0] /H /I /C [1 0 0] /Rect [240.149 658.332 275.343 670.284] /A << /S /GoTo /D <6571756174696F6E2E31302E342E332E33> >> >>}
\MusicPDFWriteObject{4399}{<< /Type /Annot /Subtype /Link /Border [0 0 0] /H /I /C [1 0 0] /Rect [450.93 522.86 481.143 534.622] /A << /S /GoTo /D <74686D2E372E312E342E32> >> >>}
\MusicPDFWriteObject{4400}{<< /Type /Annot /Subtype /Link /Border [0 0 0] /H /I /C [1 0 0] /Rect [372.83 409.206 408.023 419.996] /A << /S /GoTo /D <74686D2E31302E332E312E31> >> >>}
\MusicPDFWriteObject{4401}{<< /Type /Annot /Subtype /Link /Border [0 0 0] /H /I /C [1 0 0] /Rect [324.468 369.754 359.661 380.544] /A << /S /GoTo /D <74686D2E31302E322E312E31> >> >>}
\MusicPDFWriteObject{4402}{<< /Type /Annot /Subtype /Link /Border [0 0 0] /H /I /C [1 0 0] /Rect [422.437 369.754 457.631 380.544] /A << /S /GoTo /D <74686D2E31302E332E312E31> >> >>}
\MusicPDFWriteObject{4403}{<< /Type /Annot /Subtype /Link /Border [0 0 0] /H /I /C [0 1 0] /Rect [225.225 225.401 243.424 234.268] /A << /S /GoTo /D <636974652E4A6172> >> >>}
\MusicPDFWriteObject{4404}{<< /Type /Annot /Subtype /Link /Border [0 0 0] /H /I /C [1 0 0] /Rect [385.468 222.911 417.895 234.863] /A << /S /GoTo /D <6571756174696F6E2E432E332E362E35> >> >>}
\MusicPDFWriteObject{4405}{<< /Type /Annot /Subtype /Link /Border [0 0 0] /H /I /C [1 0 0] /Rect [508.255 222.911 540.682 234.863] /A << /S /GoTo /D <74686D2E432E332E362E34> >> >>}
\MusicPDFWriteObject{4406}{<< /Type /Annot /Subtype /Link /Border [0 0 0] /H /I /C [1 0 0] /Rect [410.235 135.199 445.428 147.152] /A << /S /GoTo /D <6571756174696F6E2E31302E322E312E34> >> >>}
\MusicPDFWriteObject{4407}{<< /Type /Annot /Subtype /Link /Border [0 0 0] /H /I /C [1 0 0] /Rect [505.553 135.198 542.855 147.153] /A << /S /GoTo /D <6E6F746174696F6E2E31302E322E312E31> >> >>}
\MusicPDFWriteObject{4408}{<< /Type /Annot /Subtype /Link /Border [0 0 0] /H /I /C [1 0 0] /Rect [213.946 671.482 249.139 683.435] /A << /S /GoTo /D <6571756174696F6E2E31302E322E312E36> >> >>}
\MusicPDFWriteObject{4409}{<< /Type /Annot /Subtype /Link /Border [0 0 0] /H /I /C [1 0 0] /Rect [375.141 519.526 410.335 531.478] /A << /S /GoTo /D <6571756174696F6E2E392E312E322E3134> >> >>}
\MusicPDFWriteObject{4410}{<< /Type /Annot /Subtype /Link /Border [0 0 0] /H /I /C [1 0 0] /Rect [509.239 517.838 545.156 532.118] /A << /S /GoTo /D <74686D2E31302E322E312E31> >> >>}
\MusicPDFWriteObject{4411}{<< /Type /Annot /Subtype /Link /Border [0 0 0] /H /I /C [1 0 0] /Rect [120.694 493.091 155.888 505.044] /A << /S /GoTo /D <6571756174696F6E2E31302E322E312E31> >> >>}
\MusicPDFWriteObject{4412}{<< /Type /Annot /Subtype /Link /Border [0 0 0] /H /I /C [1 0 0] /Rect [437.733 401.447 472.926 413.4] /A << /S /GoTo /D <6571756174696F6E2E31302E342E312E35> >> >>}
\MusicPDFWriteObject{4413}{<< /Type /Annot /Subtype /Link /Border [0 0 0] /H /I /C [1 0 0] /Rect [509.88 141.372 542.962 154.544] /A << /S /GoTo /D <74686D2E432E332E362E34> >> >>}
\MusicPDFWriteObject{4414}{<< /Type /Annot /Subtype /Link /Border [0 0 0] /H /I /C [1 0 0] /Rect [150.122 100.81 190.295 112.762] /A << /S /GoTo /D <6571756174696F6E2E31302E342E332E3134> >> >>}
\MusicPDFWriteObject{4415}{<< /Type /Annot /Subtype /Link /Border [0 0 0] /H /I /C [1 0 0] /Rect [274.178 684.633 309.371 696.586] /A << /S /GoTo /D <6571756174696F6E2E31302E332E312E31> >> >>}
\MusicPDFWriteObject{4416}{<< /Type /Annot /Subtype /Link /Border [0 0 0] /H /I /C [1 0 0] /Rect [150.464 672.036 185.657 682.827] /A << /S /GoTo /D <74686D2E31302E332E312E31> >> >>}
\MusicPDFWriteObject{4417}{<< /Type /Annot /Subtype /Link /Border [0 0 0] /H /I /C [1 0 0] /Rect [234.871 637.083 270.064 649.036] /A << /S /GoTo /D <74686D2E31302E342E312E32> >> >>}
\MusicPDFWriteObject{4418}{<< /Type /Annot /Subtype /Link /Border [0 0 0] /H /I /C [1 0 0] /Rect [432.255 623.932 467.448 635.885] /A << /S /GoTo /D <6571756174696F6E2E31302E342E312E33> >> >>}
\MusicPDFWriteObject{4419}{<< /Type /Annot /Subtype /Link /Border [0 0 0] /H /I /C [0 1 0] /Rect [404.673 600.066 423.646 608.934] /A << /S /GoTo /D <636974652E547265> >> >>}
\MusicPDFWriteObject{4420}{<< /Type /Annot /Subtype /Link /Border [0 0 0] /H /I /C [1 0 0] /Rect [371.059 494.879 411.233 506.832] /A << /S /GoTo /D <6571756174696F6E2E31302E342E332E3137> >> >>}
\MusicPDFWriteObject{4421}{<< /Type /Annot /Subtype /Link /Border [0 0 0] /H /I /C [1 0 0] /Rect [106.473 458.292 146.646 470.245] /A << /S /GoTo /D <6571756174696F6E2E31302E342E332E3132> >> >>}
\MusicPDFWriteObject{4422}{<< /Type /Annot /Subtype /Link /Border [0 0 0] /H /I /C [1 0 0] /Rect [146.65 406.187 181.844 416.978] /A << /S /GoTo /D <74686D2E31302E322E312E31> >> >>}
\MusicPDFWriteObject{4423}{<< /Type /Annot /Subtype /Link /Border [0 0 0] /H /I /C [1 0 0] /Rect [420.612 375.338 455.805 387.291] /A << /S /GoTo /D <6571756174696F6E2E392E312E322E3134> >> >>}
\MusicPDFWriteObject{4424}{<< /Type /Annot /Subtype /Link /Border [0 0 0] /H /I /C [1 0 0] /Rect [98.726 220.09 138.9 232.043] /A << /S /GoTo /D <6571756174696F6E2E31302E342E332E3131> >> >>}
\MusicPDFWriteObject{4425}{<< /Type /Annot /Subtype /Link /Border [0 0 0] /H /I /C [1 0 0] /Rect [145.961 140.997 180.491 152.95] /A << /S /GoTo /D <74686D2E31302E322E312E34> >> >>}
\MusicPDFWriteObject{4426}{<< /Type /Annot /Subtype /Link /Border [0 0 0] /H /I /C [1 0 0] /Rect [70.506 127.847 105.699 139.799] /A << /S /GoTo /D <7465726D696E6F6C6F67792E31302E322E312E32> >> >>}
\MusicPDFWriteObject{4427}{<< /Type /Annot /Subtype /Link /Border [0 0 0] /H /I /C [1 0 0] /Rect [362.174 127.847 397.367 139.799] /A << /S /GoTo /D <74686D2E31302E322E312E34> >> >>}
\MusicPDFWriteObject{4428}{<< /Type /Annot /Subtype /Link /Border [0 0 0] /H /I /C [1 0 0] /Rect [462.108 127.847 497.301 139.799] /A << /S /GoTo /D <74686D2E31312E322E302E35> >> >>}
\MusicPDFWriteObject{4429}{<< /Type /Annot /Subtype /Link /Border [0 0 0] /H /I /C [1 0 0] /Rect [270.822 114.696 310.995 126.649] /A << /S /GoTo /D <6571756174696F6E2E31302E342E332E3133> >> >>}
\MusicPDFWriteObject{4430}{<< /Type /Annot /Subtype /Link /Border [0 0 0] /H /I /C [1 0 0] /Rect [365.743 671.482 400.937 683.435] /A << /S /GoTo /D <6571756174696F6E2E31302E332E312E31> >> >>}
\MusicPDFWriteObject{4431}{<< /Type /Annot /Subtype /Link /Border [0 0 0] /H /I /C [1 0 0] /Rect [153.516 658.332 188.71 670.284] /A << /S /GoTo /D <74686D2E31302E332E312E31> >> >>}
\MusicPDFWriteObject{4432}{<< /Type /Annot /Subtype /Link /Border [0 0 0] /H /I /C [1 0 0] /Rect [332.905 645.181 373.079 657.134] /A << /S /GoTo /D <6571756174696F6E2E31302E342E332E3133> >> >>}
\MusicPDFWriteObject{4433}{<< /Type /Annot /Subtype /Link /Border [0 0 0] /H /I /C [1 0 0] /Rect [434.344 631.29 466.771 643.243] /A << /S /GoTo /D <6571756174696F6E2E432E332E362E37> >> >>}
\MusicPDFWriteObject{4434}{<< /Type /Annot /Subtype /Link /Border [0 0 0] /H /I /C [1 0 0] /Rect [71.004 620.629 103.43 629.483] /A << /S /GoTo /D <74686D2E432E332E362E34> >> >>}
\MusicPDFWriteObject{4435}{<< /Type /Annot /Subtype /Link /Border [0 0 0] /H /I /C [1 0 0] /Rect [266.702 578.554 301.895 590.507] /A << /S /GoTo /D <6571756174696F6E2E392E312E322E3134> >> >>}
\MusicPDFWriteObject{4436}{<< /Type /Annot /Subtype /Link /Border [0 0 0] /H /I /C [1 0 0] /Rect [372.059 578.554 402.272 591.146] /A << /S /GoTo /D <74686D2E392E312E322E35> >> >>}
\MusicPDFWriteObject{4437}{<< /Type /Annot /Subtype /Link /Border [0 0 0] /H /I /C [0 1 0] /Rect [88.492 389.644 106.691 398.512] /A << /S /GoTo /D <636974652E4A6172> >> >>}
\MusicPDFWriteObject{4438}{<< /Type /Annot /Subtype /Link /Border [0 0 0] /H /I /C [0 1 0] /Rect [340.434 389.644 358.633 398.512] /A << /S /GoTo /D <636974652E4A6172> >> >>}
\MusicPDFWriteObject{4439}{<< /Type /Annot /Subtype /Link /Border [0 0 0] /H /I /C [0 1 0] /Rect [336.424 350.192 354.623 359.06] /A << /S /GoTo /D <636974652E4A6172> >> >>}
\MusicPDFWriteObject{4440}{<< /Type /Annot /Subtype /Link /Border [0 0 0] /H /I /C [0 1 0] /Rect [73.715 321.551 91.914 330.419] /A << /S /GoTo /D <636974652E4A6172> >> >>}
\MusicPDFWriteObject{4441}{<< /Type /Annot /Subtype /Link /Border [0 0 0] /H /I /C [0 1 0] /Rect [251.916 308.401 270.115 317.268] /A << /S /GoTo /D <636974652E4A6172> >> >>}
\MusicPDFWriteObject{4442}{<< /Type /Annot /Subtype /Link /Border [0 0 0] /H /I /C [0 1 0] /Rect [206.785 282.099 225.758 290.967] /A << /S /GoTo /D <636974652E547265> >> >>}
\MusicPDFWriteObject{4443}{<< /Type /Annot /Subtype /Link /Border [0 0 0] /H /I /C [1 0 0] /Rect [324.37 266.458 364.544 278.411] /A << /S /GoTo /D <6571756174696F6E2E31302E342E332E3139> >> >>}
\MusicPDFWriteObject{4444}{<< /Type /Annot /Subtype /Link /Border [0 0 0] /H /I /C [0 1 0] /Rect [73.77 208.346 91.969 217.214] /A << /S /GoTo /D <636974652E4A6172> >> >>}
\MusicPDFWriteObject{4445}{<< /Type /Annot /Subtype /Link /Border [0 0 0] /H /I /C [1 0 0] /Rect [299.908 192.705 340.082 204.658] /A << /S /GoTo /D <6571756174696F6E2E31302E342E332E3139> >> >>}
\MusicPDFWriteObject{4446}{<< /Type /Annot /Subtype /Link /Border [0 0 0] /H /I /C [1 0 0] /Rect [70.506 180.108 106.363 190.898] /A << /S /GoTo /D <74686D2E31302E322E312E31> >> >>}
\MusicPDFWriteObject{4447}{<< /Type /Annot /Subtype /Link /Border [0 0 0] /H /I /C [1 0 0] /Rect [159.139 180.108 194.997 190.898] /A << /S /GoTo /D <74686D2E31302E322E312E34> >> >>}
\MusicPDFWriteObject{4448}{<< /Type /Annot /Subtype /Link /Border [0 0 0] /H /I /C [1 0 0] /Rect [263.891 180.108 299.748 190.898] /A << /S /GoTo /D <74686D2E31302E332E312E31> >> >>}
\MusicPDFWriteObject{4449}{<< /Type /Annot /Subtype /Link /Border [0 0 0] /H /I /C [1 0 0] /Rect [198.283 166.404 238.456 178.356] /A << /S /GoTo /D <6571756174696F6E2E31302E342E332E3231> >> >>}
\MusicPDFWriteObject{4450}{<< /Type /Annot /Subtype /Link /Border [0 0 0] /H /I /C [1 0 0] /Rect [398.391 671.482 433.585 683.435] /A << /S /GoTo /D <6571756174696F6E2E31302E332E312E38> >> >>}
\MusicPDFWriteObject{4451}{<< /Type /Annot /Subtype /Link /Border [0 0 0] /H /I /C [1 0 0] /Rect [496.472 671.482 531.665 683.435] /A << /S /GoTo /D <74686D2E31302E342E312E32> >> >>}
\MusicPDFWriteObject{4452}{<< /Type /Annot /Subtype /Link /Border [0 0 0] /H /I /C [1 0 0] /Rect [376.345 658.332 408.771 670.284] /A << /S /GoTo /D <74686D2E432E332E362E34> >> >>}
\MusicPDFWriteObject{4453}{<< /Type /Annot /Subtype /Link /Border [0 0 0] /H /I /C [1 0 0] /Rect [98.422 598.463 138.595 610.416] /A << /S /GoTo /D <6571756174696F6E2E31302E342E332E3135> >> >>}
\MusicPDFWriteObject{4454}{<< /Type /Annot /Subtype /Link /Border [0 0 0] /H /I /C [1 0 0] /Rect [379.355 562.589 411.782 573.38] /A << /S /GoTo /D <74686D2E432E332E362E34> >> >>}
\MusicPDFWriteObject{4455}{<< /Type /Annot /Subtype /Link /Border [0 0 0] /H /I /C [1 0 0] /Rect [460.082 512.457 500.256 524.41] /A << /S /GoTo /D <6571756174696F6E2E31302E342E332E3233> >> >>}
\MusicPDFWriteObject{4456}{<< /Type /Annot /Subtype /Link /Border [0 0 0] /H /I /C [1 0 0] /Rect [216.727 499.307 256.9 511.259] /A << /S /GoTo /D <6571756174696F6E2E31302E342E332E3234> >> >>}
\MusicPDFWriteObject{4457}{<< /Type /Annot /Subtype /Link /Border [0 0 0] /H /I /C [1 0 0] /Rect [271.996 198.415 312.169 210.368] /A << /S /GoTo /D <6571756174696F6E2E31302E342E332E3235> >> >>}
\MusicPDFWriteObject{4458}{<< /Type /Annot /Subtype /Link /Border [0 0 0] /H /I /C [1 0 0] /Rect [168.993 472.104 209.166 484.057] /A << /S /GoTo /D <6571756174696F6E2E31302E342E332E3238> >> >>}
\MusicPDFWriteObject{4459}{<< /Type /Annot /Subtype /Link /Border [0 0 0] /H /I /C [1 0 0] /Rect [279.343 470.05 315.201 484.612] /A << /S /GoTo /D <74686D2E31302E342E332E33> >> >>}
\MusicPDFWriteObject{4460}{<< /Type /Annot /Subtype /Link /Border [0 0 0] /H /I /C [1 0 0] /Rect [501.715 445.803 541.888 457.755] /A << /S /GoTo /D <6571756174696F6E2E31302E342E332E3239> >> >>}
\MusicPDFWriteObject{4461}{<< /Type /Annot /Subtype /Link /Border [0 0 0] /H /I /C [1 0 0] /Rect [255.376 109.451 289.939 120.241] /A << /S /GoTo /D <74686D2E31302E342E322E38> >> >>}
\MusicPDFWriteObject{4462}{<< /Type /Annot /Subtype /Link /Border [0 0 0] /H /I /C [1 0 0] /Rect [508.704 109.451 543.267 120.241] /A << /S /GoTo /D <74686D2E31302E342E332E31> >> >>}
\MusicPDFWriteObject{4463}{<< /Type /Annot /Subtype /Link /Border [0 0 0] /H /I /C [1 0 0] /Rect [70.506 82.596 105.699 94.549] /A << /S /GoTo /D <74686D2E31302E342E312E35> >> >>}
\MusicPDFWriteObject{4464}{<< /Type /Annot /Subtype /Link /Border [0 0 0] /H /I /C [1 0 0] /Rect [410.235 607.379 445.428 619.332] /A << /S /GoTo /D <6571756174696F6E2E31302E322E312E34> >> >>}
\MusicPDFWriteObject{4465}{<< /Type /Annot /Subtype /Link /Border [0 0 0] /H /I /C [1 0 0] /Rect [505.553 607.378 542.855 619.333] /A << /S /GoTo /D <6E6F746174696F6E2E31302E322E312E31> >> >>}
\MusicPDFWriteObject{4466}{<< /Type /Annot /Subtype /Link /Border [0 0 0] /H /I /C [1 0 0] /Rect [247.869 528.474 283.063 540.427] /A << /S /GoTo /D <6571756174696F6E2E31302E322E322E35> >> >>}
\MusicPDFWriteObject{4467}{<< /Type /Annot /Subtype /Link /Border [0 0 0] /H /I /C [1 0 0] /Rect [312.755 528.474 347.948 540.427] /A << /S /GoTo /D <6571756174696F6E2E31302E322E322E36> >> >>}
\MusicPDFWriteObject{4468}{<< /Type /Annot /Subtype /Link /Border [0 0 0] /H /I /C [1 0 0] /Rect [365.431 373.347 400.624 385.3] /A << /S /GoTo /D <6571756174696F6E2E392E322E322E3432> >> >>}
\MusicPDFWriteObject{4469}{<< /Type /Annot /Subtype /Link /Border [0 0 0] /H /I /C [1 0 0] /Rect [509.239 371.659 545.156 386.408] /A << /S /GoTo /D <74686D2E392E322E322E3131> >> >>}
\MusicPDFWriteObject{4470}{<< /Type /Annot /Subtype /Link /Border [0 0 0] /H /I /C [1 0 0] /Rect [120.694 347.045 155.888 358.998] /A << /S /GoTo /D <6571756174696F6E2E31302E322E322E31> >> >>}
\MusicPDFWriteObject{4471}{<< /Type /Annot /Subtype /Link /Border [0 0 0] /H /I /C [1 0 0] /Rect [452.748 259.665 487.941 271.618] /A << /S /GoTo /D <6571756174696F6E2E31302E342E312E35> >> >>}
\MusicPDFWriteObject{4472}{<< /Type /Annot /Subtype /Link /Border [0 0 0] /H /I /C [1 0 0] /Rect [509.88 620.11 542.962 635.143] /A << /S /GoTo /D <74686D2E432E332E362E34> >> >>}
\MusicPDFWriteObject{4473}{<< /Type /Annot /Subtype /Link /Border [0 0 0] /H /I /C [1 0 0] /Rect [150.122 581.607 190.295 593.56] /A << /S /GoTo /D <6571756174696F6E2E31302E342E332E3432> >> >>}
\MusicPDFWriteObject{4474}{<< /Type /Annot /Subtype /Link /Border [0 0 0] /H /I /C [1 0 0] /Rect [275.21 495.939 310.403 507.892] /A << /S /GoTo /D <6571756174696F6E2E392E322E322E3431> >> >>}
\MusicPDFWriteObject{4475}{<< /Type /Annot /Subtype /Link /Border [0 0 0] /H /I /C [1 0 0] /Rect [340.242 263.655 376.099 274.446] /A << /S /GoTo /D <74686D2E31302E342E332E32> >> >>}
\MusicPDFWriteObject{4476}{<< /Type /Annot /Subtype /Link /Border [0 0 0] /H /I /C [1 0 0] /Rect [285.827 150.727 321.02 162.679] /A << /S /GoTo /D <6571756174696F6E2E392E322E322E3431> >> >>}
\MusicPDFWriteObject{4477}{<< /Type /Annot /Subtype /Link /Border [0 0 0] /H /I /C [1 0 0] /Rect [82.959 560.398 201.209 569.252] /A << /S /GoTo /D <73656374696F6E2E31312E31> >> >>}
\MusicPDFWriteObject{4478}{<< /Type /Annot /Subtype /Link /Border [0 0 0] /H /I /C [1 0 0] /Rect [82.959 545.311 246.431 556.102] /A << /S /GoTo /D <73756273656374696F6E2E31312E312E31> >> >>}
\MusicPDFWriteObject{4479}{<< /Type /Annot /Subtype /Link /Border [0 0 0] /H /I /C [1 0 0] /Rect [82.959 532.16 192.426 542.951] /A << /S /GoTo /D <73756273656374696F6E2E31312E312E32> >> >>}
\MusicPDFWriteObject{4480}{<< /Type /Annot /Subtype /Link /Border [0 0 0] /H /I /C [1 0 0] /Rect [82.959 519.009 458.266 529.8] /A << /S /GoTo /D <73656374696F6E2E31312E32> >> >>}
\MusicPDFWriteObject{4481}{<< /Type /Annot /Subtype /Link /Border [0 0 0] /H /I /C [1 0 0] /Rect [82.959 505.859 211.75 516.649] /A << /S /GoTo /D <73756273656374696F6E2E31312E322E31> >> >>}
\MusicPDFWriteObject{4482}{<< /Type /Annot /Subtype /Link /Border [0 0 0] /H /I /C [1 0 0] /Rect [82.959 492.708 368.983 503.499] /A << /S /GoTo /D <73756273656374696F6E2E31312E322E32> >> >>}
\MusicPDFWriteObject{4483}{<< /Type /Annot /Subtype /Link /Border [0 0 0] /H /I /C [1 0 0] /Rect [82.959 481.494 268.607 490.348] /A << /S /GoTo /D <73656374696F6E2E31312E33> >> >>}
\MusicPDFWriteObject{4484}{<< /Type /Annot /Subtype /Link /Border [0 0 0] /H /I /C [1 0 0] /Rect [82.959 466.406 544.596 477.197] /A << /S /GoTo /D <73756273656374696F6E2E31312E332E31> >> >>}
\MusicPDFWriteObject{4485}{<< /Type /Annot /Subtype /Link /Border [0 0 0] /H /I /C [1 0 0] /Rect [71.004 455.192 182.133 464.046] /A << /S /GoTo /D <73756273656374696F6E2E31312E332E31> >> >>}
\MusicPDFWriteObject{4486}{<< /Type /Annot /Subtype /Link /Border [0 0 0] /H /I /C [1 0 0] /Rect [82.959 442.042 287.005 450.896] /A << /S /GoTo /D <73756273756273656374696F6E2E31312E332E312E31> >> >>}
\MusicPDFWriteObject{4487}{<< /Type /Annot /Subtype /Link /Border [0 0 0] /H /I /C [1 0 0] /Rect [82.959 428.891 267.666 437.745] /A << /S /GoTo /D <73756273656374696F6E2E31312E332E32> >> >>}
\MusicPDFWriteObject{4488}{<< /Type /Annot /Subtype /Link /Border [0 0 0] /H /I /C [1 0 0] /Rect [97.048 320.507 103.01 332.811] /A << /S /GoTo /D <48666F6F746E6F74652E3430> >> >>}
\MusicPDFWriteObject{4489}{<< /Type /Annot /Subtype /Link /Border [0 0 0] /H /I /C [1 0 0] /Rect [70.805 290.165 76.768 300.532] /A << /S /GoTo /D <48666F6F746E6F74652E3431> >> >>}
\MusicPDFWriteObject{4490}{<< /Type /Annot /Subtype /Link /Border [0 0 0] /H /I /C [0 1 0] /Rect [347.168 647.671 373.66 656.539] /A << /S /GoTo /D <636974652E52752D32> >> >>}
\MusicPDFWriteObject{4491}{<< /Type /Annot /Subtype /Link /Border [0 0 0] /H /I /C [1 0 0] /Rect [470.59 399.776 505.783 411.729] /A << /S /GoTo /D <6571756174696F6E2E31312E312E312E33> >> >>}
\MusicPDFWriteObject{4492}{<< /Type /Annot /Subtype /Link /Border [0 0 0] /H /I /C [1 0 0] /Rect [344.619 297.028 380.476 309.444] /A << /S /GoTo /D <74686D2E31312E312E312E32> >> >>}
\MusicPDFWriteObject{4493}{<< /Type /Annot /Subtype /Link /Border [0 0 0] /H /I /C [0 1 0] /Rect [154.061 249.996 171.215 258.864] /A << /S /GoTo /D <636974652E466F6C> >> >>}
\MusicPDFWriteObject{4494}{<< /Type /Annot /Subtype /Link /Border [0 0 0] /H /I /C [0 1 0] /Rect [219.183 137.689 236.337 146.557] /A << /S /GoTo /D <636974652E466F6C> >> >>}
\MusicPDFWriteObject{4495}{<< /Type /Annot /Subtype /Link /Border [0 0 0] /H /I /C [0 1 0] /Rect [181.21 111.387 216.265 120.255] /A << /S /GoTo /D <636974652E426F752D4949> >> >>}
\MusicPDFWriteObject{4496}{<< /Type /Annot /Subtype /Link /Border [0 0 0] /H /I /C [0 1 0] /Rect [212.982 85.086 248.037 93.954] /A << /S /GoTo /D <636974652E426F752D4949> >> >>}
\MusicPDFWriteObject{4497}{<< /Type /Annot /Subtype /Link /Border [0 0 0] /H /I /C [1 0 0] /Rect [349.061 671.122 379.274 683.898] /A << /S /GoTo /D <74686D2E322E322E322E31> >> >>}
\MusicPDFWriteObject{4498}{<< /Type /Annot /Subtype /Link /Border [0 0 0] /H /I /C [0 1 0] /Rect [197.918 660.822 215.072 669.69] /A << /S /GoTo /D <636974652E466F6C> >> >>}
\MusicPDFWriteObject{4499}{<< /Type /Annot /Subtype /Link /Border [0 0 0] /H /I /C [1 0 0] /Rect [272.123 607.885 307.316 619.838] /A << /S /GoTo /D <6571756174696F6E2E31312E312E312E35> >> >>}
\MusicPDFWriteObject{4500}{<< /Type /Annot /Subtype /Link /Border [0 0 0] /H /I /C [1 0 0] /Rect [113.196 594.734 148.389 606.687] /A << /S /GoTo /D <6571756174696F6E2E31312E312E312E33> >> >>}
\MusicPDFWriteObject{4501}{<< /Type /Annot /Subtype /Link /Border [0 0 0] /H /I /C [0 1 0] /Rect [384.249 531.471 401.403 540.338] /A << /S /GoTo /D <636974652E466F6C> >> >>}
\MusicPDFWriteObject{4502}{<< /Type /Annot /Subtype /Link /Border [0 0 0] /H /I /C [1 0 0] /Rect [506.695 403.065 541.888 415.017] /A << /S /GoTo /D <6571756174696F6E2E31312E312E312E36> >> >>}
\MusicPDFWriteObject{4503}{<< /Type /Annot /Subtype /Link /Border [0 0 0] /H /I /C [0 1 0] /Rect [279.084 313.149 297.387 322.017] /A << /S /GoTo /D <636974652E4445> >> >>}
\MusicPDFWriteObject{4504}{<< /Type /Annot /Subtype /Link /Border [0 0 0] /H /I /C [1 0 0] /Rect [70.506 166 106.164 177.955] /A << /S /GoTo /D <74686D2E31312E312E312E31> >> >>}
\MusicPDFWriteObject{4505}{<< /Type /Annot /Subtype /Link /Border [0 0 0] /H /I /C [0 1 0] /Rect [150.737 168.49 177.229 177.358] /A << /S /GoTo /D <636974652E52752D32> >> >>}
\MusicPDFWriteObject{4506}{<< /Type /Annot /Subtype /Link /Border [0 0 0] /H /I /C [1 0 0] /Rect [266.82 139.699 306.994 151.652] /A << /S /GoTo /D <6571756174696F6E2E31312E312E312E3131> >> >>}
\MusicPDFWriteObject{4507}{<< /Type /Annot /Subtype /Link /Border [0 0 0] /H /I /C [1 0 0] /Rect [109.286 126.548 149.459 138.501] /A << /S /GoTo /D <6571756174696F6E2E31312E312E312E3131> >> >>}
\MusicPDFWriteObject{4508}{<< /Type /Annot /Subtype /Link /Border [0 0 0] /H /I /C [1 0 0] /Rect [107.788 684.633 147.961 696.586] /A << /S /GoTo /D <6571756174696F6E2E31312E312E312E3131> >> >>}
\MusicPDFWriteObject{4509}{<< /Type /Annot /Subtype /Link /Border [0 0 0] /H /I /C [1 0 0] /Rect [314.119 622.312 346.823 634.267] /A << /S /GoTo /D <74686D2E412E322E322E39> >> >>}
\MusicPDFWriteObject{4510}{<< /Type /Annot /Subtype /Link /Border [0 0 0] /H /I /C [1 0 0] /Rect [297.713 360.494 332.907 372.447] /A << /S /GoTo /D <6571756174696F6E2E31312E312E312E38> >> >>}
\MusicPDFWriteObject{4511}{<< /Type /Annot /Subtype /Link /Border [0 0 0] /H /I /C [1 0 0] /Rect [405.699 360.494 440.892 372.447] /A << /S /GoTo /D <74686D2E31312E312E312E34> >> >>}
\MusicPDFWriteObject{4512}{<< /Type /Annot /Subtype /Link /Border [0 0 0] /H /I /C [1 0 0] /Rect [227.697 336.573 263.554 349.58] /A << /S /GoTo /D <74686D2E31312E312E312E32> >> >>}
\MusicPDFWriteObject{4513}{<< /Type /Annot /Subtype /Link /Border [0 0 0] /H /I /C [1 0 0] /Rect [455.485 324.474 490.28 336.429] /A << /S /GoTo /D <74686D2E31312E312E312E32> >> >>}
\MusicPDFWriteObject{4514}{<< /Type /Annot /Subtype /Link /Border [0 0 0] /H /I /C [1 0 0] /Rect [228.111 311.324 268.284 323.277] /A << /S /GoTo /D <6571756174696F6E2E31312E312E312E3132> >> >>}
\MusicPDFWriteObject{4515}{<< /Type /Annot /Subtype /Link /Border [0 0 0] /H /I /C [1 0 0] /Rect [406.002 311.324 446.175 323.277] /A << /S /GoTo /D <6571756174696F6E2E31312E312E312E3133> >> >>}
\MusicPDFWriteObject{4516}{<< /Type /Annot /Subtype /Link /Border [0 0 0] /H /I /C [1 0 0] /Rect [86.497 271.872 126.671 283.824] /A << /S /GoTo /D <6571756174696F6E2E31312E312E312E3133> >> >>}
\MusicPDFWriteObject{4517}{<< /Type /Annot /Subtype /Link /Border [0 0 0] /H /I /C [1 0 0] /Rect [98.422 186.274 136.105 198.226] /A << /S /GoTo /D <6571756174696F6E2E412E322E322E3132> >> >>}
\MusicPDFWriteObject{4518}{<< /Type /Annot /Subtype /Link /Border [0 0 0] /H /I /C [1 0 0] /Rect [73.712 159.972 111.395 171.925] /A << /S /GoTo /D <6571756174696F6E2E412E322E332E3134> >> >>}
\MusicPDFWriteObject{4519}{<< /Type /Annot /Subtype /Link /Border [0 0 0] /H /I /C [1 0 0] /Rect [255.646 133.67 293.33 145.625] /A << /S /GoTo /D <74686D2E412E322E332E3130> >> >>}
\MusicPDFWriteObject{4520}{<< /Type /Annot /Subtype /Link /Border [0 0 0] /H /I /C [1 0 0] /Rect [228.345 443.332 264.202 456.178] /A << /S /GoTo /D <74686D2E31312E312E312E32> >> >>}
\MusicPDFWriteObject{4521}{<< /Type /Annot /Subtype /Link /Border [0 0 0] /H /I /C [1 0 0] /Rect [277.305 430.435 317.478 442.388] /A << /S /GoTo /D <6571756174696F6E2E31312E312E312E3132> >> >>}
\MusicPDFWriteObject{4522}{<< /Type /Annot /Subtype /Link /Border [0 0 0] /H /I /C [1 0 0] /Rect [136.611 417.285 176.785 429.237] /A << /S /GoTo /D <6571756174696F6E2E31312E312E312E3132> >> >>}
\MusicPDFWriteObject{4523}{<< /Type /Annot /Subtype /Link /Border [0 0 0] /H /I /C [1 0 0] /Rect [356.649 416.924 391.842 429.239] /A << /S /GoTo /D <74686D2E31312E312E312E32> >> >>}
\MusicPDFWriteObject{4524}{<< /Type /Annot /Subtype /Link /Border [0 0 0] /H /I /C [1 0 0] /Rect [429.613 343.107 468.01 357.554] /A << /S /GoTo /D <74686D2E412E322E332E3130> >> >>}
\MusicPDFWriteObject{4525}{<< /Type /Annot /Subtype /Link /Border [0 0 0] /H /I /C [1 0 0] /Rect [237.477 318.406 273.334 331.252] /A << /S /GoTo /D <74686D2E31312E312E312E32> >> >>}
\MusicPDFWriteObject{4526}{<< /Type /Annot /Subtype /Link /Border [0 0 0] /H /I /C [1 0 0] /Rect [217.197 305.51 257.371 317.462] /A << /S /GoTo /D <6571756174696F6E2E31312E312E312E3134> >> >>}
\MusicPDFWriteObject{4527}{<< /Type /Annot /Subtype /Link /Border [0 0 0] /H /I /C [1 0 0] /Rect [222.791 292.359 255.632 304.312] /A << /S /GoTo /D <74686D2E442E312E322E34> >> >>}
\MusicPDFWriteObject{4528}{<< /Type /Annot /Subtype /Link /Border [0 0 0] /H /I /C [0 1 0] /Rect [259.726 294.849 294.781 303.717] /A << /S /GoTo /D <636974652E426F752D4949> >> >>}
\MusicPDFWriteObject{4529}{<< /Type /Annot /Subtype /Link /Border [0 0 0] /H /I /C [1 0 0] /Rect [147.31 161.14 182.503 173.455] /A << /S /GoTo /D <74686D2E31312E312E312E32> >> >>}
\MusicPDFWriteObject{4530}{<< /Type /Annot /Subtype /Link /Border [0 0 0] /H /I /C [1 0 0] /Rect [501.715 135.199 541.888 147.152] /A << /S /GoTo /D <6571756174696F6E2E31312E312E312E3135> >> >>}
\MusicPDFWriteObject{4531}{<< /Type /Annot /Subtype /Link /Border [0 0 0] /H /I /C [1 0 0] /Rect [388.882 108.897 429.056 120.85] /A << /S /GoTo /D <6571756174696F6E2E31312E312E312E3133> >> >>}
\MusicPDFWriteObject{4532}{<< /Type /Annot /Subtype /Link /Border [0 0 0] /H /I /C [1 0 0] /Rect [380.699 561.956 386.662 576.429] /A << /S /GoTo /D <48666F6F746E6F74652E3432> >> >>}
\MusicPDFWriteObject{4533}{<< /Type /Annot /Subtype /Link /Border [0 0 0] /H /I /C [1 0 0] /Rect [189.28 253.517 224.473 265.469] /A << /S /GoTo /D <6571756174696F6E2E31312E312E322E35> >> >>}
\MusicPDFWriteObject{4534}{<< /Type /Annot /Subtype /Link /Border [0 0 0] /H /I /C [1 0 0] /Rect [278.437 227.215 313.63 239.168] /A << /S /GoTo /D <6571756174696F6E2E31312E312E322E36> >> >>}
\MusicPDFWriteObject{4535}{<< /Type /Annot /Subtype /Link /Border [0 0 0] /H /I /C [1 0 0] /Rect [233.938 214.064 269.131 226.017] /A << /S /GoTo /D <6571756174696F6E2E31312E312E322E34> >> >>}
\MusicPDFWriteObject{4536}{<< /Type /Annot /Subtype /Link /Border [0 0 0] /H /I /C [1 0 0] /Rect [229.001 684.633 264.858 696.586] /A << /S /GoTo /D <74686D2E31312E322E302E31> >> >>}
\MusicPDFWriteObject{4537}{<< /Type /Annot /Subtype /Link /Border [0 0 0] /H /I /C [1 0 0] /Rect [485.935 514.549 521.128 526.502] /A << /S /GoTo /D <6571756174696F6E2E31312E312E322E38> >> >>}
\MusicPDFWriteObject{4538}{<< /Type /Annot /Subtype /Link /Border [0 0 0] /H /I /C [1 0 0] /Rect [250.878 421.679 286.071 433.632] /A << /S /GoTo /D <6571756174696F6E2E31312E312E322E39> >> >>}
\MusicPDFWriteObject{4539}{<< /Type /Annot /Subtype /Link /Border [0 0 0] /H /I /C [1 0 0] /Rect [412.079 394.199 447.272 406.152] /A << /S /GoTo /D <6571756174696F6E2E31312E312E322E39> >> >>}
\MusicPDFWriteObject{4540}{<< /Type /Annot /Subtype /Link /Border [0 0 0] /H /I /C [1 0 0] /Rect [378.417 381.048 413.61 393.001] /A << /S /GoTo /D <6571756174696F6E2E31312E312E322E38> >> >>}
\MusicPDFWriteObject{4541}{<< /Type /Annot /Subtype /Link /Border [0 0 0] /H /I /C [1 0 0] /Rect [168.106 291.525 208.279 303.478] /A << /S /GoTo /D <6571756174696F6E2E31312E312E322E3130> >> >>}
\MusicPDFWriteObject{4542}{<< /Type /Annot /Subtype /Link /Border [0 0 0] /H /I /C [0 1 0] /Rect [509.226 280.865 543.195 289.732] /A << /S /GoTo /D <636974652E436E772D61> >> >>}
\MusicPDFWriteObject{4543}{<< /Type /Annot /Subtype /Link /Border [0 0 0] /H /I /C [1 0 0] /Rect [112.727 166.892 142.94 175.746] /A << /S /GoTo /D <74686D2E362E332E322E32> >> >>}
\MusicPDFWriteObject{4544}{<< /Type /Annot /Subtype /Link /Border [0 0 0] /H /I /C [1 0 0] /Rect [167.982 658.332 208.155 670.284] /A << /S /GoTo /D <6571756174696F6E2E31312E312E322E3132> >> >>}
\MusicPDFWriteObject{4545}{<< /Type /Annot /Subtype /Link /Border [0 0 0] /H /I /C [1 0 0] /Rect [111.012 525.392 116.974 538.25] /A << /S /GoTo /D <48666F6F746E6F74652E3433> >> >>}
\MusicPDFWriteObject{4546}{<< /Type /Annot /Subtype /Link /Border [0 0 0] /H /I /C [1 0 0] /Rect [256.972 512.242 292.166 524.194] /A << /S /GoTo /D <6571756174696F6E2E31312E312E322E35> >> >>}
\MusicPDFWriteObject{4547}{<< /Type /Annot /Subtype /Link /Border [0 0 0] /H /I /C [1 0 0] /Rect [402.015 512.242 437.872 524.194] /A << /S /GoTo /D <74686D2E31312E322E302E32> >> >>}
\MusicPDFWriteObject{4548}{<< /Type /Annot /Subtype /Link /Border [0 0 0] /H /I /C [1 0 0] /Rect [468.267 432.691 504.124 443.482] /A << /S /GoTo /D <74686D2E31312E322E302E32> >> >>}
\MusicPDFWriteObject{4549}{<< /Type /Annot /Subtype /Link /Border [0 0 0] /H /I /C [1 0 0] /Rect [442.461 392.685 477.654 404.638] /A << /S /GoTo /D <6571756174696F6E2E31312E312E322E36> >> >>}
\MusicPDFWriteObject{4550}{<< /Type /Annot /Subtype /Link /Border [0 0 0] /H /I /C [1 0 0] /Rect [88.936 332.446 94.899 342.814] /A << /S /GoTo /D <48666F6F746E6F74652E3434> >> >>}
\MusicPDFWriteObject{4551}{<< /Type /Annot /Subtype /Link /Border [0 0 0] /H /I /C [0 1 0] /Rect [378.263 305.652 412.232 314.52] /A << /S /GoTo /D <636974652E436E772D61> >> >>}
\MusicPDFWriteObject{4552}{<< /Type /Annot /Subtype /Link /Border [0 0 0] /H /I /C [1 0 0] /Rect [145.207 290.012 180.4 301.964] /A << /S /GoTo /D <6571756174696F6E2E31312E312E322E35> >> >>}
\MusicPDFWriteObject{4553}{<< /Type /Annot /Subtype /Link /Border [0 0 0] /H /I /C [1 0 0] /Rect [262.346 276.861 297.871 288.814] /A << /S /GoTo /D <74686D2E31312E312E312E32> >> >>}
\MusicPDFWriteObject{4554}{<< /Type /Annot /Subtype /Link /Border [0 0 0] /H /I /C [1 0 0] /Rect [275.537 250.559 313.221 262.512] /A << /S /GoTo /D <6571756174696F6E2E412E322E332E3134> >> >>}
\MusicPDFWriteObject{4555}{<< /Type /Annot /Subtype /Link /Border [0 0 0] /H /I /C [1 0 0] /Rect [112.727 226.748 150.41 235.602] /A << /S /GoTo /D <74686D2E412E322E332E3130> >> >>}
\MusicPDFWriteObject{4556}{<< /Type /Annot /Subtype /Link /Border [0 0 0] /H /I /C [1 0 0] /Rect [211 198.51 246.193 209.3] /A << /S /GoTo /D <74686D2E31312E312E312E32> >> >>}
\MusicPDFWriteObject{4557}{<< /Type /Annot /Subtype /Link /Border [0 0 0] /H /I /C [1 0 0] /Rect [219.466 117.281 249.679 129.233] /A << /S /GoTo /D <74686D2E362E322E312E31> >> >>}
\MusicPDFWriteObject{4558}{<< /Type /Annot /Subtype /Link /Border [0 0 0] /H /I /C [1 0 0] /Rect [174.979 643.924 209.509 654.772] /A << /S /GoTo /D <74686D2E31312E312E312E32> >> >>}
\MusicPDFWriteObject{4559}{<< /Type /Annot /Subtype /Link /Border [0 0 0] /H /I /C [1 0 0] /Rect [123.292 556.547 157.821 569.437] /A << /S /GoTo /D <74686D2E31312E312E312E32> >> >>}
\MusicPDFWriteObject{4560}{<< /Type /Annot /Subtype /Link /Border [0 0 0] /H /I /C [1 0 0] /Rect [88.936 226.212 94.899 236.579] /A << /S /GoTo /D <48666F6F746E6F74652E3435> >> >>}
\MusicPDFWriteObject{4561}{<< /Type /Annot /Subtype /Link /Border [0 0 0] /H /I /C [1 0 0] /Rect [158.132 176.484 193.325 188.437] /A << /S /GoTo /D <6571756174696F6E2E31302E322E312E33> >> >>}
\MusicPDFWriteObject{4562}{<< /Type /Annot /Subtype /Link /Border [0 0 0] /H /I /C [1 0 0] /Rect [384.359 176.484 419.552 188.437] /A << /S /GoTo /D <6571756174696F6E2E392E312E312E3131> >> >>}
\MusicPDFWriteObject{4563}{<< /Type /Annot /Subtype /Link /Border [0 0 0] /H /I /C [1 0 0] /Rect [216.118 140.056 251.312 152.009] /A << /S /GoTo /D <6571756174696F6E2E31302E322E312E33> >> >>}
\MusicPDFWriteObject{4564}{<< /Type /Annot /Subtype /Link /Border [0 0 0] /H /I /C [1 0 0] /Rect [216.118 648.206 251.312 660.158] /A << /S /GoTo /D <6571756174696F6E2E352E332E312E3130> >> >>}
\MusicPDFWriteObject{4565}{<< /Type /Annot /Subtype /Link /Border [0 0 0] /H /I /C [1 0 0] /Rect [135.65 635.054 165.299 646.399] /A << /S /GoTo /D <74686D2E352E332E312E35> >> >>}
\MusicPDFWriteObject{4566}{<< /Type /Annot /Subtype /Link /Border [0 0 0] /H /I /C [1 0 0] /Rect [71.004 622.457 101.217 633.248] /A << /S /GoTo /D <74686D2E362E332E322E32> >> >>}
\MusicPDFWriteObject{4567}{<< /Type /Annot /Subtype /Link /Border [0 0 0] /H /I /C [1 0 0] /Rect [332.718 471.256 367.911 483.209] /A << /S /GoTo /D <6571756174696F6E2E352E332E312E3130> >> >>}
\MusicPDFWriteObject{4568}{<< /Type /Annot /Subtype /Link /Border [0 0 0] /H /I /C [1 0 0] /Rect [91.42 440.455 126.613 452.407] /A << /S /GoTo /D <6571756174696F6E2E31302E312E312E35> >> >>}
\MusicPDFWriteObject{4569}{<< /Type /Annot /Subtype /Link /Border [0 0 0] /H /I /C [1 0 0] /Rect [150.124 342.611 185.982 354.564] /A << /S /GoTo /D <74686D2E31302E322E312E31> >> >>}
\MusicPDFWriteObject{4570}{<< /Type /Annot /Subtype /Link /Border [0 0 0] /H /I /C [0 1 0] /Rect [324.115 345.101 341.588 353.969] /A << /S /GoTo /D <636974652E4346> >> >>}
\MusicPDFWriteObject{4571}{<< /Type /Annot /Subtype /Link /Border [0 0 0] /H /I /C [0 1 0] /Rect [455.849 331.95 474.954 340.818] /A << /S /GoTo /D <636974652E4E42> >> >>}
\MusicPDFWriteObject{4572}{<< /Type /Annot /Subtype /Link /Border [0 0 0] /H /I /C [0 1 0] /Rect [89.402 305.649 107.601 314.517] /A << /S /GoTo /D <636974652E4A6172> >> >>}
\MusicPDFWriteObject{4573}{<< /Type /Annot /Subtype /Link /Border [0 0 0] /H /I /C [1 0 0] /Rect [190.344 290.008 220.557 301.961] /A << /S /GoTo /D <6571756174696F6E2E352E332E312E39> >> >>}
\MusicPDFWriteObject{4574}{<< /Type /Annot /Subtype /Link /Border [0 0 0] /H /I /C [1 0 0] /Rect [445.19 290.008 475.403 301.961] /A << /S /GoTo /D <6571756174696F6E2E352E332E312E31> >> >>}
\MusicPDFWriteObject{4575}{<< /Type /Annot /Subtype /Link /Border [0 0 0] /H /I /C [1 0 0] /Rect [106.556 263.706 136.769 275.661] /A << /S /GoTo /D <74686D2E352E332E312E35> >> >>}
\MusicPDFWriteObject{4576}{<< /Type /Annot /Subtype /Link /Border [0 0 0] /H /I /C [1 0 0] /Rect [179.104 250.556 209.317 262.509] /A << /S /GoTo /D <6571756174696F6E2E352E332E312E31> >> >>}
\MusicPDFWriteObject{4577}{<< /Type /Annot /Subtype /Link /Border [0 0 0] /H /I /C [1 0 0] /Rect [452.054 95.747 482.267 107.699] /A << /S /GoTo /D <6571756174696F6E2E352E332E312E31> >> >>}
\MusicPDFWriteObject{4578}{<< /Type /Annot /Subtype /Link /Border [0 0 0] /H /I /C [1 0 0] /Rect [299.427 573.075 334.62 585.028] /A << /S /GoTo /D <6571756174696F6E2E31302E322E312E39> >> >>}
\MusicPDFWriteObject{4579}{<< /Type /Annot /Subtype /Link /Border [0 0 0] /H /I /C [1 0 0] /Rect [461.987 573.075 497.844 585.028] /A << /S /GoTo /D <74686D2E31302E322E312E31> >> >>}
\MusicPDFWriteObject{4580}{<< /Type /Annot /Subtype /Link /Border [0 0 0] /H /I /C [1 0 0] /Rect [247.391 560.478 277.604 571.268] /A << /S /GoTo /D <74686D2E392E332E322E31> >> >>}
\MusicPDFWriteObject{4581}{<< /Type /Annot /Subtype /Link /Border [0 0 0] /H /I /C [1 0 0] /Rect [349.362 534.39 383.891 545.181] /A << /S /GoTo /D <74686D2E31302E322E312E31> >> >>}
\MusicPDFWriteObject{4582}{<< /Type /Annot /Subtype /Link /Border [0 0 0] /H /I /C [1 0 0] /Rect [380.765 484.259 415.958 496.211] /A << /S /GoTo /D <6571756174696F6E2E31312E322E302E33> >> >>}
\MusicPDFWriteObject{4583}{<< /Type /Annot /Subtype /Link /Border [0 0 0] /H /I /C [1 0 0] /Rect [283.354 471.108 318.548 483.061] /A << /S /GoTo /D <6571756174696F6E2E392E312E312E3131> >> >>}
\MusicPDFWriteObject{4584}{<< /Type /Annot /Subtype /Link /Border [0 0 0] /H /I /C [1 0 0] /Rect [466.871 471.108 502.064 483.061] /A << /S /GoTo /D <6571756174696F6E2E31302E322E312E33> >> >>}
\MusicPDFWriteObject{4585}{<< /Type /Annot /Subtype /Link /Border [0 0 0] /H /I /C [1 0 0] /Rect [314.662 221.272 349.855 233.225] /A << /S /GoTo /D <6571756174696F6E2E31302E322E312E33> >> >>}
\MusicPDFWriteObject{4586}{<< /Type /Annot /Subtype /Link /Border [0 0 0] /H /I /C [1 0 0] /Rect [332.668 185.333 363.446 197.286] /A << /S /GoTo /D <74686D2E342E322E312E33> >> >>}
\MusicPDFWriteObject{4587}{<< /Type /Annot /Subtype /Link /Border [0 0 0] /H /I /C [1 0 0] /Rect [98.884 159.032 134.077 170.985] /A << /S /GoTo /D <6571756174696F6E2E31302E322E312E33> >> >>}
\MusicPDFWriteObject{4588}{<< /Type /Annot /Subtype /Link /Border [0 0 0] /H /I /C [1 0 0] /Rect [168.604 94.744 203.963 110.556] /A << /S /GoTo /D <74686D2E31312E322E302E32> >> >>}
\MusicPDFWriteObject{4589}{<< /Type /Annot /Subtype /Link /Border [0 0 0] /H /I /C [1 0 0] /Rect [112.727 563.041 142.94 575.356] /A << /S /GoTo /D <74686D2E362E332E322E32> >> >>}
\MusicPDFWriteObject{4590}{<< /Type /Annot /Subtype /Link /Border [0 0 0] /H /I /C [1 0 0] /Rect [185.098 550.251 220.291 562.204] /A << /S /GoTo /D <6571756174696F6E2E31312E322E302E37> >> >>}
\MusicPDFWriteObject{4591}{<< /Type /Annot /Subtype /Link /Border [0 0 0] /H /I /C [1 0 0] /Rect [508.74 550.251 543.27 562.204] /A << /S /GoTo /D <74686D2E31302E322E312E31> >> >>}
\MusicPDFWriteObject{4592}{<< /Type /Annot /Subtype /Link /Border [0 0 0] /H /I /C [1 0 0] /Rect [183.61 537.653 214.387 548.444] /A << /S /GoTo /D <74686D2E362E332E322E31> >> >>}
\MusicPDFWriteObject{4593}{<< /Type /Annot /Subtype /Link /Border [0 0 0] /H /I /C [1 0 0] /Rect [282.365 537.653 313.142 548.444] /A << /S /GoTo /D <74686D2E362E332E322E32> >> >>}
\MusicPDFWriteObject{4594}{<< /Type /Annot /Subtype /Link /Border [0 0 0] /H /I /C [1 0 0] /Rect [322.967 523.949 358.16 535.902] /A << /S /GoTo /D <6571756174696F6E2E31302E322E312E33> >> >>}
\MusicPDFWriteObject{4595}{<< /Type /Annot /Subtype /Link /Border [0 0 0] /H /I /C [1 0 0] /Rect [73.712 510.799 108.905 522.751] /A << /S /GoTo /D <6571756174696F6E2E31312E322E302E35> >> >>}
\MusicPDFWriteObject{4596}{<< /Type /Annot /Subtype /Link /Border [0 0 0] /H /I /C [1 0 0] /Rect [475.431 510.799 505.08 522.751] /A << /S /GoTo /D <74686D2E392E332E322E31> >> >>}
\MusicPDFWriteObject{4597}{<< /Type /Annot /Subtype /Link /Border [0 0 0] /H /I /C [1 0 0] /Rect [146.597 498.201 181.126 508.992] /A << /S /GoTo /D <74686D2E31302E322E312E31> >> >>}
\MusicPDFWriteObject{4598}{<< /Type /Annot /Subtype /Link /Border [0 0 0] /H /I /C [1 0 0] /Rect [197.028 484.497 232.221 496.45] /A << /S /GoTo /D <6571756174696F6E2E31312E322E302E35> >> >>}
\MusicPDFWriteObject{4599}{<< /Type /Annot /Subtype /Link /Border [0 0 0] /H /I /C [1 0 0] /Rect [212.017 465.72 246.547 476.511] /A << /S /GoTo /D <74686D2E31302E322E312E33> >> >>}
\MusicPDFWriteObject{4600}{<< /Type /Annot /Subtype /Link /Border [0 0 0] /H /I /C [1 0 0] /Rect [309.977 465.72 344.506 476.511] /A << /S /GoTo /D <74686D2E31302E322E312E34> >> >>}
\MusicPDFWriteObject{4601}{<< /Type /Annot /Subtype /Link /Border [0 0 0] /H /I /C [1 0 0] /Rect [349.31 452.569 384.503 463.36] /A << /S /GoTo /D <74686D2E31312E322E302E32> >> >>}
\MusicPDFWriteObject{4602}{<< /Type /Annot /Subtype /Link /Border [0 0 0] /H /I /C [1 0 0] /Rect [314.007 415.096 349.2 427.049] /A << /S /GoTo /D <6571756174696F6E2E31302E322E312E33> >> >>}
\MusicPDFWriteObject{4603}{<< /Type /Annot /Subtype /Link /Border [0 0 0] /H /I /C [1 0 0] /Rect [281.088 402.498 316.48 413.289] /A << /S /GoTo /D <7465726D696E6F6C6F67792E31302E322E312E32> >> >>}
\MusicPDFWriteObject{4604}{<< /Type /Annot /Subtype /Link /Border [0 0 0] /H /I /C [1 0 0] /Rect [289.501 293.05 324.695 305.002] /A << /S /GoTo /D <6571756174696F6E2E31312E322E302E36> >> >>}
\MusicPDFWriteObject{4605}{<< /Type /Annot /Subtype /Link /Border [0 0 0] /H /I /C [1 0 0] /Rect [251.499 256.13 286.692 268.082] /A << /S /GoTo /D <6571756174696F6E2E31312E322E302E38> >> >>}
\MusicPDFWriteObject{4606}{<< /Type /Annot /Subtype /Link /Border [0 0 0] /H /I /C [1 0 0] /Rect [138.096 173.984 168.14 186.299] /A << /S /GoTo /D <74686D2E392E312E322E35> >> >>}
\MusicPDFWriteObject{4607}{<< /Type /Annot /Subtype /Link /Border [0 0 0] /H /I /C [1 0 0] /Rect [204.918 129.266 240.111 140.056] /A << /S /GoTo /D <74686D2E31302E322E312E34> >> >>}
\MusicPDFWriteObject{4608}{<< /Type /Annot /Subtype /Link /Border [0 0 0] /H /I /C [1 0 0] /Rect [312.43 635.005 347.623 646.958] /A << /S /GoTo /D <6571756174696F6E2E31312E322E302E36> >> >>}
\MusicPDFWriteObject{4609}{<< /Type /Annot /Subtype /Link /Border [0 0 0] /H /I /C [1 0 0] /Rect [143.141 542.298 178.334 554.251] /A << /S /GoTo /D <6571756174696F6E2E31312E322E302E39> >> >>}
\MusicPDFWriteObject{4610}{<< /Type /Annot /Subtype /Link /Border [0 0 0] /H /I /C [1 0 0] /Rect [248.498 542.298 283.691 554.251] /A << /S /GoTo /D <74686D2E31312E322E302E34> >> >>}
\MusicPDFWriteObject{4611}{<< /Type /Annot /Subtype /Link /Border [0 0 0] /H /I /C [1 0 0] /Rect [226.67 528.178 261.199 541.102] /A << /S /GoTo /D <74686D2E31312E322E302E34> >> >>}
\MusicPDFWriteObject{4612}{<< /Type /Annot /Subtype /Link /Border [0 0 0] /H /I /C [1 0 0] /Rect [366.327 528.178 400.856 541.102] /A << /S /GoTo /D <74686D2E31312E322E302E32> >> >>}
\MusicPDFWriteObject{4613}{<< /Type /Annot /Subtype /Link /Border [0 0 0] /H /I /C [1 0 0] /Rect [226.681 502.259 261.875 514.212] /A << /S /GoTo /D <6571756174696F6E2E31302E322E312E33> >> >>}
\MusicPDFWriteObject{4614}{<< /Type /Annot /Subtype /Link /Border [0 0 0] /H /I /C [1 0 0] /Rect [169.73 475.958 204.923 487.911] /A << /S /GoTo /D <6571756174696F6E2E392E312E322E3134> >> >>}
\MusicPDFWriteObject{4615}{<< /Type /Annot /Subtype /Link /Border [0 0 0] /H /I /C [1 0 0] /Rect [290.969 361.138 326.827 371.929] /A << /S /GoTo /D <74686D2E31312E322E312E32> >> >>}
\MusicPDFWriteObject{4616}{<< /Type /Annot /Subtype /Link /Border [0 0 0] /H /I /C [1 0 0] /Rect [222.145 371.831 252.923 386.305] /A << /S /GoTo /D <74686D2E372E332E322E32> >> >>}
\MusicPDFWriteObject{4617}{<< /Type /Annot /Subtype /Link /Border [0 0 0] /H /I /C [1 0 0] /Rect [332.425 371.831 365.46 386.305] /A << /S /GoTo /D <74686D2E432E332E352E35> >> >>}
\MusicPDFWriteObject{4618}{<< /Type /Annot /Subtype /Link /Border [0 0 0] /H /I /C [1 0 0] /Rect [501.39 327.796 531.237 343.597] /A << /S /GoTo /D <74686D2E372E332E322E32> >> >>}
\MusicPDFWriteObject{4619}{<< /Type /Annot /Subtype /Link /Border [0 0 0] /H /I /C [1 0 0] /Rect [277.707 316.84 310.133 328.793] /A << /S /GoTo /D <74686D2E432E332E352E35> >> >>}
\MusicPDFWriteObject{4620}{<< /Type /Annot /Subtype /Link /Border [0 0 0] /H /I /C [1 0 0] /Rect [264.54 292.64 299.734 303.431] /A << /S /GoTo /D <74686D2E31312E322E312E31> >> >>}
\MusicPDFWriteObject{4621}{<< /Type /Annot /Subtype /Link /Border [0 0 0] /H /I /C [1 0 0] /Rect [260.28 555.631 295.473 566.422] /A << /S /GoTo /D <74686D2E31312E322E312E32> >> >>}
\MusicPDFWriteObject{4622}{<< /Type /Annot /Subtype /Link /Border [0 0 0] /H /I /C [1 0 0] /Rect [176.821 355.761 212.015 367.714] /A << /S /GoTo /D <6571756174696F6E2E372E332E322E3230> >> >>}
\MusicPDFWriteObject{4623}{<< /Type /Annot /Subtype /Link /Border [0 0 0] /H /I /C [1 0 0] /Rect [297.502 350.727 328.279 368.625] /A << /S /GoTo /D <74686D2E372E332E322E34> >> >>}
\MusicPDFWriteObject{4624}{<< /Type /Annot /Subtype /Link /Border [0 0 0] /H /I /C [1 0 0] /Rect [165.404 342.869 200.598 351.723] /A << /S /GoTo /D <74686D2E31312E322E312E32> >> >>}
\MusicPDFWriteObject{4625}{<< /Type /Annot /Subtype /Link /Border [0 0 0] /H /I /C [1 0 0] /Rect [249.148 262.663 283.677 277.303] /A << /S /GoTo /D <74686D2E31302E332E312E31> >> >>}
\MusicPDFWriteObject{4626}{<< /Type /Annot /Subtype /Link /Border [0 0 0] /H /I /C [1 0 0] /Rect [160.996 187.861 190.645 200.175] /A << /S /GoTo /D <74686D2E392E332E312E31> >> >>}
\MusicPDFWriteObject{4627}{<< /Type /Annot /Subtype /Link /Border [0 0 0] /H /I /C [1 0 0] /Rect [165.636 175.68 188.102 187.632] /A << /S /GoTo /D <73756273656374696F6E2E352E332E31> >> >>}
\MusicPDFWriteObject{4628}{<< /Type /Annot /Subtype /Link /Border [0 0 0] /H /I /C [1 0 0] /Rect [308.316 175.68 338.529 187.632] /A << /S /GoTo /D <74686D2E352E332E312E32> >> >>}
\MusicPDFWriteObject{4629}{<< /Type /Annot /Subtype /Link /Border [0 0 0] /H /I /C [1 0 0] /Rect [405.566 175.68 435.779 187.632] /A << /S /GoTo /D <74686D2E352E332E312E34> >> >>}
\MusicPDFWriteObject{4630}{<< /Type /Annot /Subtype /Link /Border [0 0 0] /H /I /C [1 0 0] /Rect [414.561 147.713 449.755 159.666] /A << /S /GoTo /D <6571756174696F6E2E31312E322E302E33> >> >>}
\MusicPDFWriteObject{4631}{<< /Type /Annot /Subtype /Link /Border [0 0 0] /H /I /C [1 0 0] /Rect [462.865 132.898 498.059 144.851] /A << /S /GoTo /D <6571756174696F6E2E31312E322E302E32> >> >>}
\MusicPDFWriteObject{4632}{<< /Type /Annot /Subtype /Link /Border [0 0 0] /H /I /C [1 0 0] /Rect [227.322 684.633 257.535 696.586] /A << /S /GoTo /D <6571756174696F6E2E372E342E302E32> >> >>}
\MusicPDFWriteObject{4633}{<< /Type /Annot /Subtype /Link /Border [0 0 0] /H /I /C [1 0 0] /Rect [495.398 667.599 525.612 683.411] /A << /S /GoTo /D <74686D2E372E342E322E31> >> >>}
\MusicPDFWriteObject{4634}{<< /Type /Annot /Subtype /Link /Border [0 0 0] /H /I /C [1 0 0] /Rect [271.411 445.183 306.604 457.136] /A << /S /GoTo /D <6571756174696F6E2E31312E322E302E36> >> >>}
\MusicPDFWriteObject{4635}{<< /Type /Annot /Subtype /Link /Border [0 0 0] /H /I /C [0 1 0] /Rect [111.436 353.054 130.409 361.922] /A << /S /GoTo /D <636974652E547265> >> >>}
\MusicPDFWriteObject{4636}{<< /Type /Annot /Subtype /Link /Border [0 0 0] /H /I /C [1 0 0] /Rect [154.966 285.616 189.495 299.807] /A << /S /GoTo /D <74686D2E31312E322E302E32> >> >>}
\MusicPDFWriteObject{4637}{<< /Type /Annot /Subtype /Link /Border [0 0 0] /H /I /C [1 0 0] /Rect [427.698 285.616 462.227 299.807] /A << /S /GoTo /D <74686D2E31312E322E302E35> >> >>}
\MusicPDFWriteObject{4638}{<< /Type /Annot /Subtype /Link /Border [0 0 0] /H /I /C [1 0 0] /Rect [208.676 82.596 248.849 94.549] /A << /S /GoTo /D <6571756174696F6E2E31302E332E312E3137> >> >>}
\MusicPDFWriteObject{4639}{<< /Type /Annot /Subtype /Link /Border [0 0 0] /H /I /C [1 0 0] /Rect [374.209 82.596 409.402 94.549] /A << /S /GoTo /D <74686D2E31302E332E312E31> >> >>}
\MusicPDFWriteObject{4640}{<< /Type /Annot /Subtype /Link /Border [0 0 0] /H /I /C [1 0 0] /Rect [70.506 579.524 106.363 593.697] /A << /S /GoTo /D <74686D2E31302E332E312E31> >> >>}
\MusicPDFWriteObject{4641}{<< /Type /Annot /Subtype /Link /Border [0 0 0] /H /I /C [1 0 0] /Rect [255.055 581.742 290.249 593.695] /A << /S /GoTo /D <6571756174696F6E2E31312E322E322E34> >> >>}
\MusicPDFWriteObject{4642}{<< /Type /Annot /Subtype /Link /Border [0 0 0] /H /I /C [1 0 0] /Rect [278.008 568.567 318.182 580.52] /A << /S /GoTo /D <6571756174696F6E2E31312E322E322E3131> >> >>}
\MusicPDFWriteObject{4643}{<< /Type /Annot /Subtype /Link /Border [0 0 0] /H /I /C [1 0 0] /Rect [101.576 555.416 141.749 567.369] /A << /S /GoTo /D <6571756174696F6E2E31312E322E322E3132> >> >>}
\MusicPDFWriteObject{4644}{<< /Type /Annot /Subtype /Link /Border [0 0 0] /H /I /C [1 0 0] /Rect [326.695 467.419 361.789 479.734] /A << /S /GoTo /D <74686D2E31312E322E322E31> >> >>}
\MusicPDFWriteObject{4645}{<< /Type /Annot /Subtype /Link /Border [0 0 0] /H /I /C [1 0 0] /Rect [73.712 454.629 113.885 466.582] /A << /S /GoTo /D <6571756174696F6E2E31312E322E322E3133> >> >>}
\MusicPDFWriteObject{4646}{<< /Type /Annot /Subtype /Link /Border [0 0 0] /H /I /C [1 0 0] /Rect [170.978 498.401 206.171 510.353] /A << /S /GoTo /D <6571756174696F6E2E31312E332E312E31> >> >>}
\MusicPDFWriteObject{4647}{<< /Type /Annot /Subtype /Link /Border [0 0 0] /H /I /C [1 0 0] /Rect [70.506 585.477 100.719 597.429] /A << /S /GoTo /D <74686D2E312E382E312E31> >> >>}
\MusicPDFWriteObject{4648}{<< /Type /Annot /Subtype /Link /Border [0 0 0] /H /I /C [1 0 0] /Rect [354.945 472.736 390.138 484.689] /A << /S /GoTo /D <6571756174696F6E2E31312E332E312E33> >> >>}
\MusicPDFWriteObject{4649}{<< /Type /Annot /Subtype /Link /Border [0 0 0] /H /I /C [1 0 0] /Rect [139.539 393.09 174.732 405.043] /A << /S /GoTo /D <6571756174696F6E2E31312E332E312E33> >> >>}
\MusicPDFWriteObject{4650}{<< /Type /Annot /Subtype /Link /Border [0 0 0] /H /I /C [1 0 0] /Rect [191.651 238.142 226.844 250.094] /A << /S /GoTo /D <6571756174696F6E2E31312E332E322E31> >> >>}
\MusicPDFWriteObject{4651}{<< /Type /Annot /Subtype /Link /Border [0 0 0] /H /I /C [1 0 0] /Rect [112.727 199.243 147.92 210.033] /A << /S /GoTo /D <74686D2E31312E312E322E31> >> >>}
\MusicPDFWriteObject{4652}{<< /Type /Annot /Subtype /Link /Border [0 0 0] /H /I /C [1 0 0] /Rect [238.834 199.243 274.028 210.033] /A << /S /GoTo /D <74686D2E31322E322E322E32> >> >>}
\MusicPDFWriteObject{4653}{<< /Type /Annot /Subtype /Link /Border [0 0 0] /H /I /C [1 0 0] /Rect [170.761 630.841 201.538 642.26] /A << /S /GoTo /D <74686D2E322E322E332E32> >> >>}
\MusicPDFWriteObject{4654}{<< /Type /Annot /Subtype /Link /Border [0 0 0] /H /I /C [1 0 0] /Rect [201.335 604.539 237.192 615.33] /A << /S /GoTo /D <74686D2E31302E332E322E31> >> >>}
\MusicPDFWriteObject{4655}{<< /Type /Annot /Subtype /Link /Border [0 0 0] /H /I /C [1 0 0] /Rect [73.712 590.835 113.885 602.788] /A << /S /GoTo /D <6571756174696F6E2E31312E322E322E3132> >> >>}
\MusicPDFWriteObject{4656}{<< /Type /Annot /Subtype /Link /Border [0 0 0] /H /I /C [1 0 0] /Rect [239.245 590.835 274.438 602.788] /A << /S /GoTo /D <74686D2E31312E322E322E31> >> >>}
\MusicPDFWriteObject{4657}{<< /Type /Annot /Subtype /Link /Border [0 0 0] /H /I /C [1 0 0] /Rect [185.916 468.977 216.13 480.929] /A << /S /GoTo /D <6571756174696F6E2E372E342E322E31> >> >>}
\MusicPDFWriteObject{4658}{<< /Type /Annot /Subtype /Link /Border [0 0 0] /H /I /C [1 0 0] /Rect [286.661 468.977 316.733 480.929] /A << /S /GoTo /D <74686D2E372E342E322E31> >> >>}
\MusicPDFWriteObject{4659}{<< /Type /Annot /Subtype /Link /Border [0 0 0] /H /I /C [1 0 0] /Rect [183.775 455.826 216.064 467.779] /A << /S /GoTo /D <74686D2E422E342E312E32> >> >>}
\MusicPDFWriteObject{4660}{<< /Type /Annot /Subtype /Link /Border [0 0 0] /H /I /C [1 0 0] /Rect [82.959 558.462 259.587 569.252] /A << /S /GoTo /D <73656374696F6E2E31322E31> >> >>}
\MusicPDFWriteObject{4661}{<< /Type /Annot /Subtype /Link /Border [0 0 0] /H /I /C [1 0 0] /Rect [82.959 545.311 288.997 556.102] /A << /S /GoTo /D <73756273656374696F6E2E31322E312E31> >> >>}
\MusicPDFWriteObject{4662}{<< /Type /Annot /Subtype /Link /Border [0 0 0] /H /I /C [1 0 0] /Rect [82.959 532.16 304.989 542.951] /A << /S /GoTo /D <73656374696F6E2E31322E32> >> >>}
\MusicPDFWriteObject{4663}{<< /Type /Annot /Subtype /Link /Border [0 0 0] /H /I /C [1 0 0] /Rect [82.959 520.946 177.415 529.8] /A << /S /GoTo /D <73756273656374696F6E2E31322E322E31> >> >>}
\MusicPDFWriteObject{4664}{<< /Type /Annot /Subtype /Link /Border [0 0 0] /H /I /C [1 0 0] /Rect [82.959 505.859 333.957 516.649] /A << /S /GoTo /D <73756273656374696F6E2E31322E322E32> >> >>}
\MusicPDFWriteObject{4665}{<< /Type /Annot /Subtype /Link /Border [0 0 0] /H /I /C [1 0 0] /Rect [82.959 492.708 483.209 503.499] /A << /S /GoTo /D <73656374696F6E2E31322E33> >> >>}
\MusicPDFWriteObject{4666}{<< /Type /Annot /Subtype /Link /Border [0 0 0] /H /I /C [1 0 0] /Rect [82.959 481.494 220.868 490.348] /A << /S /GoTo /D <73656374696F6E2E31322E34> >> >>}
\MusicPDFWriteObject{4667}{<< /Type /Annot /Subtype /Link /Border [0 0 0] /H /I /C [1 0 0] /Rect [82.959 466.406 301.186 477.197] /A << /S /GoTo /D <73756273656374696F6E2E31322E342E31> >> >>}
\MusicPDFWriteObject{4668}{<< /Type /Annot /Subtype /Link /Border [0 0 0] /H /I /C [1 0 0] /Rect [82.959 455.192 320.012 464.046] /A << /S /GoTo /D <73656374696F6E2E31322E35> >> >>}
\MusicPDFWriteObject{4669}{<< /Type /Annot /Subtype /Link /Border [0 0 0] /H /I /C [1 0 0] /Rect [82.959 442.042 253.694 450.896] /A << /S /GoTo /D <73756273656374696F6E2E31322E352E31> >> >>}
\MusicPDFWriteObject{4670}{<< /Type /Annot /Subtype /Link /Border [0 0 0] /H /I /C [1 0 0] /Rect [116.452 637.874 146.101 651.158] /A << /S /GoTo /D <74686D2E392E332E332E31> >> >>}
\MusicPDFWriteObject{4671}{<< /Type /Annot /Subtype /Link /Border [0 0 0] /H /I /C [1 0 0] /Rect [272.338 533.997 305.041 545.95] /A << /S /GoTo /D <74686D2E412E352E332E37> >> >>}
\MusicPDFWriteObject{4672}{<< /Type /Annot /Subtype /Link /Border [0 0 0] /H /I /C [1 0 0] /Rect [367.818 533.997 400.521 545.95] /A << /S /GoTo /D <74686D2E412E352E342E31> >> >>}
\MusicPDFWriteObject{4673}{<< /Type /Annot /Subtype /Link /Border [0 0 0] /H /I /C [1 0 0] /Rect [345.135 448.966 377.838 461.236] /A << /S /GoTo /D <74686D2E412E352E322E31> >> >>}
\MusicPDFWriteObject{4674}{<< /Type /Annot /Subtype /Link /Border [0 0 0] /H /I /C [1 0 0] /Rect [257.731 377.742 292.924 389.694] /A << /S /GoTo /D <6571756174696F6E2E31322E312E312E31> >> >>}
\MusicPDFWriteObject{4675}{<< /Type /Annot /Subtype /Link /Border [0 0 0] /H /I /C [1 0 0] /Rect [284.794 364.591 319.987 376.544] /A << /S /GoTo /D <6571756174696F6E2E31322E312E312E31> >> >>}
\MusicPDFWriteObject{4676}{<< /Type /Annot /Subtype /Link /Border [0 0 0] /H /I /C [1 0 0] /Rect [500.196 148.35 525.496 160.322] /A << /S /GoTo /D <73756273656374696F6E2E412E352E32> >> >>}
\MusicPDFWriteObject{4677}{<< /Type /Annot /Subtype /Link /Border [0 0 0] /H /I /C [1 0 0] /Rect [117.346 137.689 150.05 146.543] /A << /S /GoTo /D <74686D2E412E352E322E31> >> >>}
\MusicPDFWriteObject{4678}{<< /Type /Annot /Subtype /Link /Border [0 0 0] /H /I /C [1 0 0] /Rect [217.087 137.689 254.77 146.543] /A << /S /GoTo /D <74686D2E412E352E322E3138> >> >>}
\MusicPDFWriteObject{4679}{<< /Type /Annot /Subtype /Link /Border [0 0 0] /H /I /C [0 1 0] /Rect [385.08 98.237 414.172 107.105] /A << /S /GoTo /D <636974652E48522D32> >> >>}
\MusicPDFWriteObject{4680}{<< /Type /Annot /Subtype /Link /Border [0 0 0] /H /I /C [1 0 0] /Rect [70.625 85.086 103.328 93.94] /A << /S /GoTo /D <74686D2E412E352E332E37> >> >>}
\MusicPDFWriteObject{4681}{<< /Type /Annot /Subtype /Link /Border [0 0 0] /H /I /C [1 0 0] /Rect [166.105 85.086 198.808 93.94] /A << /S /GoTo /D <74686D2E412E352E342E31> >> >>}
\MusicPDFWriteObject{4682}{<< /Type /Annot /Subtype /Link /Border [0 0 0] /H /I /C [1 0 0] /Rect [387.883 644.883 419.972 657.153] /A << /S /GoTo /D <74686D2E412E352E322E35> >> >>}
\MusicPDFWriteObject{4683}{<< /Type /Annot /Subtype /Link /Border [0 0 0] /H /I /C [1 0 0] /Rect [70.26 559.333 104.092 568.187] /A << /S /GoTo /D <74686D2E412E352E322E31> >> >>}
\MusicPDFWriteObject{4684}{<< /Type /Annot /Subtype /Link /Border [0 0 0] /H /I /C [0 1 0] /Rect [89.402 673.973 123.371 682.84] /A << /S /GoTo /D <636974652E436E772D61> >> >>}
\MusicPDFWriteObject{4685}{<< /Type /Annot /Subtype /Link /Border [0 0 0] /H /I /C [0 1 0] /Rect [107.784 660.822 141.753 669.69] /A << /S /GoTo /D <636974652E436E772D61> >> >>}
\MusicPDFWriteObject{4686}{<< /Type /Annot /Subtype /Link /Border [0 0 0] /H /I /C [1 0 0] /Rect [348.175 644.943 380.328 657.536] /A << /S /GoTo /D <74686D2E432E332E362E31> >> >>}
\MusicPDFWriteObject{4687}{<< /Type /Annot /Subtype /Link /Border [0 0 0] /H /I /C [0 1 0] /Rect [168.561 594.831 202.53 603.698] /A << /S /GoTo /D <636974652E436E772D61> >> >>}
\MusicPDFWriteObject{4688}{<< /Type /Annot /Subtype /Link /Border [0 0 0] /H /I /C [0 1 0] /Rect [338.701 594.831 372.67 603.698] /A << /S /GoTo /D <636974652E436E772D61> >> >>}
\MusicPDFWriteObject{4689}{<< /Type /Annot /Subtype /Link /Border [0 0 0] /H /I /C [1 0 0] /Rect [124.623 578.829 157.049 591.782] /A << /S /GoTo /D <74686D2E432E332E362E35> >> >>}
\MusicPDFWriteObject{4690}{<< /Type /Annot /Subtype /Link /Border [0 0 0] /H /I /C [1 0 0] /Rect [242.375 578.829 274.801 591.782] /A << /S /GoTo /D <74686D2E432E332E352E35> >> >>}
\MusicPDFWriteObject{4691}{<< /Type /Annot /Subtype /Link /Border [0 0 0] /H /I /C [1 0 0] /Rect [297.714 199.708 337.888 211.661] /A << /S /GoTo /D <6571756174696F6E2E31322E312E312E3131> >> >>}
\MusicPDFWriteObject{4692}{<< /Type /Annot /Subtype /Link /Border [0 0 0] /H /I /C [1 0 0] /Rect [319.087 186.197 352.405 198.512] /A << /S /GoTo /D <74686D2E412E352E342E33> >> >>}
\MusicPDFWriteObject{4693}{<< /Type /Annot /Subtype /Link /Border [0 0 0] /H /I /C [0 1 0] /Rect [246.328 162.746 278.272 171.614] /A << /S /GoTo /D <636974652E427269732D61> >> >>}
\MusicPDFWriteObject{4694}{<< /Type /Annot /Subtype /Link /Border [0 0 0] /H /I /C [1 0 0] /Rect [180.476 147.105 212.903 159.06] /A << /S /GoTo /D <74686D2E432E332E362E35> >> >>}
\MusicPDFWriteObject{4695}{<< /Type /Annot /Subtype /Link /Border [0 0 0] /H /I /C [1 0 0] /Rect [70.625 87.812 103.943 99.768] /A << /S /GoTo /D <74686D2E412E352E332E37> >> >>}
\MusicPDFWriteObject{4696}{<< /Type /Annot /Subtype /Link /Border [0 0 0] /H /I /C [1 0 0] /Rect [313.218 87.813 350.901 99.766] /A << /S /GoTo /D <6571756174696F6E2E412E352E332E3136> >> >>}
\MusicPDFWriteObject{4697}{<< /Type /Annot /Subtype /Link /Border [0 0 0] /H /I /C [1 0 0] /Rect [246.296 684.633 283.979 696.586] /A << /S /GoTo /D <6571756174696F6E2E412E352E332E3138> >> >>}
\MusicPDFWriteObject{4698}{<< /Type /Annot /Subtype /Link /Border [0 0 0] /H /I /C [0 1 0] /Rect [73.732 673.973 105.675 682.84] /A << /S /GoTo /D <636974652E427269732D61> >> >>}
\MusicPDFWriteObject{4699}{<< /Type /Annot /Subtype /Link /Border [0 0 0] /H /I /C [0 1 0] /Rect [189.822 673.973 221.766 682.84] /A << /S /GoTo /D <636974652E427269732D61> >> >>}
\MusicPDFWriteObject{4700}{<< /Type /Annot /Subtype /Link /Border [0 0 0] /H /I /C [1 0 0] /Rect [112.144 658.331 144.386 670.286] /A << /S /GoTo /D <74686D2E412E352E342E31> >> >>}
\MusicPDFWriteObject{4701}{<< /Type /Annot /Subtype /Link /Border [0 0 0] /H /I /C [1 0 0] /Rect [73.712 468.252 108.905 480.205] /A << /S /GoTo /D <6571756174696F6E2E31322E312E312E33> >> >>}
\MusicPDFWriteObject{4702}{<< /Type /Annot /Subtype /Link /Border [0 0 0] /H /I /C [1 0 0] /Rect [174.227 468.252 209.42 480.205] /A << /S /GoTo /D <6571756174696F6E2E31322E312E312E32> >> >>}
\MusicPDFWriteObject{4703}{<< /Type /Annot /Subtype /Link /Border [0 0 0] /H /I /C [1 0 0] /Rect [135.022 408.711 168.339 420.981] /A << /S /GoTo /D <74686D2E412E352E322E35> >> >>}
\MusicPDFWriteObject{4704}{<< /Type /Annot /Subtype /Link /Border [0 0 0] /H /I /C [1 0 0] /Rect [70.625 395.858 103.943 407.811] /A << /S /GoTo /D <74686D2E412E352E342E33> >> >>}
\MusicPDFWriteObject{4705}{<< /Type /Annot /Subtype /Link /Border [0 0 0] /H /I /C [1 0 0] /Rect [165.894 395.858 206.068 407.811] /A << /S /GoTo /D <6571756174696F6E2E31322E312E312E3133> >> >>}
\MusicPDFWriteObject{4706}{<< /Type /Annot /Subtype /Link /Border [0 0 0] /H /I /C [0 1 0] /Rect [272.399 371.671 289.553 380.538] /A << /S /GoTo /D <636974652E466F6C> >> >>}
\MusicPDFWriteObject{4707}{<< /Type /Annot /Subtype /Link /Border [0 0 0] /H /I /C [0 1 0] /Rect [196.455 332.218 230.424 341.086] /A << /S /GoTo /D <636974652E436E772D61> >> >>}
\MusicPDFWriteObject{4708}{<< /Type /Annot /Subtype /Link /Border [0 0 0] /H /I /C [1 0 0] /Rect [171.641 317.13 204.067 327.922] /A << /S /GoTo /D <74686D2E432E332E342E34> >> >>}
\MusicPDFWriteObject{4709}{<< /Type /Annot /Subtype /Link /Border [0 0 0] /H /I /C [1 0 0] /Rect [215.113 262.812 247.54 276.406] /A << /S /GoTo /D <74686D2E432E332E342E34> >> >>}
\MusicPDFWriteObject{4710}{<< /Type /Annot /Subtype /Link /Border [0 0 0] /H /I /C [1 0 0] /Rect [164.237 94.838 199.43 106.791] /A << /S /GoTo /D <6571756174696F6E2E31322E312E312E33> >> >>}
\MusicPDFWriteObject{4711}{<< /Type /Annot /Subtype /Link /Border [0 0 0] /H /I /C [1 0 0] /Rect [138.875 605.048 170.693 618.001] /A << /S /GoTo /D <74686D2E432E332E342E34> >> >>}
\MusicPDFWriteObject{4712}{<< /Type /Annot /Subtype /Link /Border [0 0 0] /H /I /C [1 0 0] /Rect [506.695 605.408 541.888 617.361] /A << /S /GoTo /D <6571756174696F6E2E31322E312E312E33> >> >>}
\MusicPDFWriteObject{4713}{<< /Type /Annot /Subtype /Link /Border [0 0 0] /H /I /C [1 0 0] /Rect [250.068 384.9 290.242 396.853] /A << /S /GoTo /D <6571756174696F6E2E31322E312E312E3138> >> >>}
\MusicPDFWriteObject{4714}{<< /Type /Annot /Subtype /Link /Border [0 0 0] /H /I /C [1 0 0] /Rect [453.645 359.152 486.68 370] /A << /S /GoTo /D <74686D2E432E332E362E35> >> >>}
\MusicPDFWriteObject{4715}{<< /Type /Annot /Subtype /Link /Border [0 0 0] /H /I /C [1 0 0] /Rect [149.196 313.025 182.514 327.692] /A << /S /GoTo /D <74686D2E412E352E342E31> >> >>}
\MusicPDFWriteObject{4716}{<< /Type /Annot /Subtype /Link /Border [0 0 0] /H /I /C [1 0 0] /Rect [328.715 299.874 361.08 312.187] /A << /S /GoTo /D <74686D2E412E352E342E32> >> >>}
\MusicPDFWriteObject{4717}{<< /Type /Annot /Subtype /Link /Border [0 0 0] /H /I /C [1 0 0] /Rect [427.244 299.874 459.609 312.187] /A << /S /GoTo /D <74686D2E412E352E342E31> >> >>}
\MusicPDFWriteObject{4718}{<< /Type /Annot /Subtype /Link /Border [0 0 0] /H /I /C [1 0 0] /Rect [282.756 260.782 317.949 272.735] /A << /S /GoTo /D <6571756174696F6E2E31322E312E312E31> >> >>}
\MusicPDFWriteObject{4719}{<< /Type /Annot /Subtype /Link /Border [0 0 0] /H /I /C [1 0 0] /Rect [181.036 195.582 213.74 206.372] /A << /S /GoTo /D <74686D2E412E352E342E31> >> >>}
\MusicPDFWriteObject{4720}{<< /Type /Annot /Subtype /Link /Border [0 0 0] /H /I /C [1 0 0] /Rect [209.011 590.348 244.735 602.301] /A << /S /GoTo /D <74686D2E31322E312E312E32> >> >>}
\MusicPDFWriteObject{4721}{<< /Type /Annot /Subtype /Link /Border [0 0 0] /H /I /C [1 0 0] /Rect [291.559 590.348 327.283 602.301] /A << /S /GoTo /D <74686D2E31302E342E332E34> >> >>}
\MusicPDFWriteObject{4722}{<< /Type /Annot /Subtype /Link /Border [0 0 0] /H /I /C [1 0 0] /Rect [390.909 590.348 426.633 602.301] /A << /S /GoTo /D <74686D2E31302E342E322E31> >> >>}
\MusicPDFWriteObject{4723}{<< /Type /Annot /Subtype /Link /Border [0 0 0] /H /I /C [1 0 0] /Rect [70.625 564.047 102.714 575.999] /A << /S /GoTo /D <74686D2E412E352E342E31> >> >>}
\MusicPDFWriteObject{4724}{<< /Type /Annot /Subtype /Link /Border [0 0 0] /H /I /C [1 0 0] /Rect [170.975 534.928 206.832 550.563] /A << /S /GoTo /D <74686D2E31322E312E312E32> >> >>}
\MusicPDFWriteObject{4725}{<< /Type /Annot /Subtype /Link /Border [0 0 0] /H /I /C [1 0 0] /Rect [284.237 534.928 320.094 550.563] /A << /S /GoTo /D <74686D2E31302E342E332E34> >> >>}
\MusicPDFWriteObject{4726}{<< /Type /Annot /Subtype /Link /Border [0 0 0] /H /I /C [1 0 0] /Rect [507.302 509.954 543.159 522.547] /A << /S /GoTo /D <74686D2E31302E342E332E34> >> >>}
\MusicPDFWriteObject{4727}{<< /Type /Annot /Subtype /Link /Border [0 0 0] /H /I /C [1 0 0] /Rect [108.58 470.342 143.773 482.295] /A << /S /GoTo /D <6571756174696F6E2E31322E312E312E32> >> >>}
\MusicPDFWriteObject{4728}{<< /Type /Annot /Subtype /Link /Border [0 0 0] /H /I /C [1 0 0] /Rect [214.601 470.342 249.794 482.295] /A << /S /GoTo /D <74686D2E31322E312E312E32> >> >>}
\MusicPDFWriteObject{4729}{<< /Type /Annot /Subtype /Link /Border [0 0 0] /H /I /C [1 0 0] /Rect [257.762 470.342 297.935 482.295] /A << /S /GoTo /D <6571756174696F6E2E31322E312E312E3230> >> >>}
\MusicPDFWriteObject{4730}{<< /Type /Annot /Subtype /Link /Border [0 0 0] /H /I /C [1 0 0] /Rect [326.377 470.342 366.55 482.295] /A << /S /GoTo /D <6571756174696F6E2E31322E312E312E3231> >> >>}
\MusicPDFWriteObject{4731}{<< /Type /Annot /Subtype /Link /Border [0 0 0] /H /I /C [1 0 0] /Rect [146.286 407.293 181.48 419.246] /A << /S /GoTo /D <74686D2E31302E342E332E31> >> >>}
\MusicPDFWriteObject{4732}{<< /Type /Annot /Subtype /Link /Border [0 0 0] /H /I /C [1 0 0] /Rect [296.602 407.293 336.776 419.246] /A << /S /GoTo /D <6571756174696F6E2E31322E312E312E3233> >> >>}
\MusicPDFWriteObject{4733}{<< /Type /Annot /Subtype /Link /Border [0 0 0] /H /I /C [1 0 0] /Rect [108.58 319.144 148.754 331.097] /A << /S /GoTo /D <6571756174696F6E2E31302E342E332E3332> >> >>}
\MusicPDFWriteObject{4734}{<< /Type /Annot /Subtype /Link /Border [0 0 0] /H /I /C [1 0 0] /Rect [215.403 319.144 250.597 331.097] /A << /S /GoTo /D <74686D2E31322E312E312E32> >> >>}
\MusicPDFWriteObject{4735}{<< /Type /Annot /Subtype /Link /Border [0 0 0] /H /I /C [1 0 0] /Rect [359.052 319.144 399.225 331.097] /A << /S /GoTo /D <6571756174696F6E2E31322E312E312E3234> >> >>}
\MusicPDFWriteObject{4736}{<< /Type /Annot /Subtype /Link /Border [0 0 0] /H /I /C [1 0 0] /Rect [106.473 277.771 146.646 289.723] /A << /S /GoTo /D <6571756174696F6E2E31302E342E332E3331> >> >>}
\MusicPDFWriteObject{4737}{<< /Type /Annot /Subtype /Link /Border [0 0 0] /H /I /C [1 0 0] /Rect [382.323 256.412 418.181 272.048] /A << /S /GoTo /D <74686D2E31302E342E332E34> >> >>}
\MusicPDFWriteObject{4738}{<< /Type /Annot /Subtype /Link /Border [0 0 0] /H /I /C [1 0 0] /Rect [71.004 94.418 100.652 110.054] /A << /S /GoTo /D <74686D2E392E332E332E31> >> >>}
\MusicPDFWriteObject{4739}{<< /Type /Annot /Subtype /Link /Border [0 0 0] /H /I /C [0 1 0] /Rect [247.742 595.068 280.481 603.936] /A << /S /GoTo /D <636974652E427269732D62> >> >>}
\MusicPDFWriteObject{4740}{<< /Type /Annot /Subtype /Link /Border [0 0 0] /H /I /C [0 1 0] /Rect [345.068 444.91 362.222 453.778] /A << /S /GoTo /D <636974652E466F6C> >> >>}
\MusicPDFWriteObject{4741}{<< /Type /Annot /Subtype /Link /Border [0 0 0] /H /I /C [0 1 0] /Rect [137.673 366.006 154.827 374.874] /A << /S /GoTo /D <636974652E466F6C> >> >>}
\MusicPDFWriteObject{4742}{<< /Type /Annot /Subtype /Link /Border [0 0 0] /H /I /C [0 1 0] /Rect [134.344 324.982 151.498 333.85] /A << /S /GoTo /D <636974652E466F6C> >> >>}
\MusicPDFWriteObject{4743}{<< /Type /Annot /Subtype /Link /Border [0 0 0] /H /I /C [0 1 0] /Rect [235.313 324.982 265.706 333.85] /A << /S /GoTo /D <636974652E426F752D49> >> >>}
\MusicPDFWriteObject{4744}{<< /Type /Annot /Subtype /Link /Border [0 0 0] /H /I /C [0 1 0] /Rect [423.631 324.982 454.024 333.85] /A << /S /GoTo /D <636974652E426F752D49> >> >>}
\MusicPDFWriteObject{4745}{<< /Type /Annot /Subtype /Link /Border [0 0 0] /H /I /C [0 1 0] /Rect [126.883 311.831 144.038 320.699] /A << /S /GoTo /D <636974652E466F6C> >> >>}
\MusicPDFWriteObject{4746}{<< /Type /Annot /Subtype /Link /Border [0 0 0] /H /I /C [0 1 0] /Rect [183.155 298.68 213.548 307.548] /A << /S /GoTo /D <636974652E426F752D49> >> >>}
\MusicPDFWriteObject{4747}{<< /Type /Annot /Subtype /Link /Border [0 0 0] /H /I /C [0 1 0] /Rect [218.616 298.68 253.671 307.548] /A << /S /GoTo /D <636974652E426F752D4949> >> >>}
\MusicPDFWriteObject{4748}{<< /Type /Annot /Subtype /Link /Border [0 0 0] /H /I /C [0 1 0] /Rect [434.941 298.68 452.095 307.548] /A << /S /GoTo /D <636974652E466F6C> >> >>}
\MusicPDFWriteObject{4749}{<< /Type /Annot /Subtype /Link /Border [0 0 0] /H /I /C [0 1 0] /Rect [526.005 298.68 543.159 307.548] /A << /S /GoTo /D <636974652E466F6C> >> >>}
\MusicPDFWriteObject{4750}{<< /Type /Annot /Subtype /Link /Border [0 0 0] /H /I /C [0 1 0] /Rect [284.949 187.509 320.004 196.377] /A << /S /GoTo /D <636974652E426F752D4949> >> >>}
\MusicPDFWriteObject{4751}{<< /Type /Annot /Subtype /Link /Border [0 0 0] /H /I /C [0 1 0] /Rect [265.219 477.181 283.521 486.048] /A << /S /GoTo /D <636974652E4445> >> >>}
\MusicPDFWriteObject{4752}{<< /Type /Annot /Subtype /Link /Border [0 0 0] /H /I /C [0 1 0] /Rect [111.149 450.879 129.451 459.747] /A << /S /GoTo /D <636974652E4445> >> >>}
\MusicPDFWriteObject{4753}{<< /Type /Annot /Subtype /Link /Border [0 0 0] /H /I /C [0 1 0] /Rect [73.826 248.03 108.881 256.897] /A << /S /GoTo /D <636974652E426F752D4949> >> >>}
\MusicPDFWriteObject{4754}{<< /Type /Annot /Subtype /Link /Border [0 0 0] /H /I /C [1 0 0] /Rect [86.901 179.786 122.095 191.739] /A << /S /GoTo /D <6571756174696F6E2E31322E322E312E34> >> >>}
\MusicPDFWriteObject{4755}{<< /Type /Annot /Subtype /Link /Border [0 0 0] /H /I /C [0 1 0] /Rect [360.588 182.276 395.643 191.144] /A << /S /GoTo /D <636974652E426F752D4949> >> >>}
\MusicPDFWriteObject{4756}{<< /Type /Annot /Subtype /Link /Border [0 0 0] /H /I /C [0 1 0] /Rect [279.249 155.974 314.304 164.842] /A << /S /GoTo /D <636974652E426F752D4949> >> >>}
\MusicPDFWriteObject{4757}{<< /Type /Annot /Subtype /Link /Border [0 0 0] /H /I /C [1 0 0] /Rect [221.33 510.842 256.523 522.795] /A << /S /GoTo /D <6571756174696F6E2E31322E322E312E35> >> >>}
\MusicPDFWriteObject{4758}{<< /Type /Annot /Subtype /Link /Border [0 0 0] /H /I /C [1 0 0] /Rect [294.556 437.986 330.474 450.579] /A << /S /GoTo /D <74686D2E31322E322E312E31> >> >>}
\MusicPDFWriteObject{4759}{<< /Type /Annot /Subtype /Link /Border [0 0 0] /H /I /C [0 1 0] /Rect [187.455 407.44 222.51 416.308] /A << /S /GoTo /D <636974652E426F752D4949> >> >>}
\MusicPDFWriteObject{4760}{<< /Type /Annot /Subtype /Link /Border [0 0 0] /H /I /C [0 1 0] /Rect [266.725 394.289 301.78 403.157] /A << /S /GoTo /D <636974652E426F752D4949> >> >>}
\MusicPDFWriteObject{4761}{<< /Type /Annot /Subtype /Link /Border [0 0 0] /H /I /C [0 1 0] /Rect [422.166 328.535 439.32 337.403] /A << /S /GoTo /D <636974652E466F6C> >> >>}
\MusicPDFWriteObject{4762}{<< /Type /Annot /Subtype /Link /Border [0 0 0] /H /I /C [0 1 0] /Rect [73.729 302.234 104.122 311.101] /A << /S /GoTo /D <636974652E426F752D49> >> >>}
\MusicPDFWriteObject{4763}{<< /Type /Annot /Subtype /Link /Border [0 0 0] /H /I /C [0 1 0] /Rect [268.71 302.234 299.103 311.101] /A << /S /GoTo /D <636974652E426F752D49> >> >>}
\MusicPDFWriteObject{4764}{<< /Type /Annot /Subtype /Link /Border [0 0 0] /H /I /C [0 1 0] /Rect [468.08 302.234 485.234 311.101] /A << /S /GoTo /D <636974652E466F6C> >> >>}
\MusicPDFWriteObject{4765}{<< /Type /Annot /Subtype /Link /Border [0 0 0] /H /I /C [0 1 0] /Rect [477.394 289.083 494.549 297.951] /A << /S /GoTo /D <636974652E466F6C> >> >>}
\MusicPDFWriteObject{4766}{<< /Type /Annot /Subtype /Link /Border [0 0 0] /H /I /C [0 1 0] /Rect [93.138 275.932 110.292 284.8] /A << /S /GoTo /D <636974652E466F6C> >> >>}
\MusicPDFWriteObject{4767}{<< /Type /Annot /Subtype /Link /Border [0 0 0] /H /I /C [0 1 0] /Rect [107.666 262.781 138.059 271.649] /A << /S /GoTo /D <636974652E426F752D49> >> >>}
\MusicPDFWriteObject{4768}{<< /Type /Annot /Subtype /Link /Border [0 0 0] /H /I /C [1 0 0] /Rect [131.526 220.839 166.72 232.792] /A << /S /GoTo /D <6571756174696F6E2E31322E322E312E36> >> >>}
\MusicPDFWriteObject{4769}{<< /Type /Annot /Subtype /Link /Border [0 0 0] /H /I /C [1 0 0] /Rect [322.31 194.538 357.504 206.49] /A << /S /GoTo /D <6571756174696F6E2E31322E322E312E37> >> >>}
\MusicPDFWriteObject{4770}{<< /Type /Annot /Subtype /Link /Border [0 0 0] /H /I /C [1 0 0] /Rect [303.333 168.236 338.526 180.189] /A << /S /GoTo /D <6571756174696F6E2E31322E322E312E36> >> >>}
\MusicPDFWriteObject{4771}{<< /Type /Annot /Subtype /Link /Border [0 0 0] /H /I /C [1 0 0] /Rect [399.951 168.236 435.808 180.828] /A << /S /GoTo /D <74686D2E31322E322E312E31> >> >>}
\MusicPDFWriteObject{4772}{<< /Type /Annot /Subtype /Link /Border [0 0 0] /H /I /C [0 1 0] /Rect [88.122 144.425 105.277 153.292] /A << /S /GoTo /D <636974652E466F6C> >> >>}
\MusicPDFWriteObject{4773}{<< /Type /Annot /Subtype /Link /Border [0 0 0] /H /I /C [1 0 0] /Rect [163.329 129.337 198.523 140.128] /A << /S /GoTo /D <74686D2E31322E322E312E31> >> >>}
\MusicPDFWriteObject{4774}{<< /Type /Annot /Subtype /Link /Border [0 0 0] /H /I /C [1 0 0] /Rect [381.435 82.298 416.628 93.94] /A << /S /GoTo /D <74686D2E31322E322E312E31> >> >>}
\MusicPDFWriteObject{4775}{<< /Type /Annot /Subtype /Link /Border [0 0 0] /H /I /C [1 0 0] /Rect [257.458 672.036 292.652 682.827] /A << /S /GoTo /D <74686D2E31322E322E312E32> >> >>}
\MusicPDFWriteObject{4776}{<< /Type /Annot /Subtype /Link /Border [0 0 0] /H /I /C [0 1 0] /Rect [398.225 583.03 441.248 591.897] /A << /S /GoTo /D <636974652E426F752D4754> >> >>}
\MusicPDFWriteObject{4777}{<< /Type /Annot /Subtype /Link /Border [0 0 0] /H /I /C [0 1 0] /Rect [262.021 541.046 305.044 549.914] /A << /S /GoTo /D <636974652E426F752D4754> >> >>}
\MusicPDFWriteObject{4778}{<< /Type /Annot /Subtype /Link /Border [0 0 0] /H /I /C [0 1 0] /Rect [107.335 527.895 150.358 536.763] /A << /S /GoTo /D <636974652E426F752D4754> >> >>}
\MusicPDFWriteObject{4779}{<< /Type /Annot /Subtype /Link /Border [0 0 0] /H /I /C [0 1 0] /Rect [351.878 514.744 394.901 523.612] /A << /S /GoTo /D <636974652E426F752D4754> >> >>}
\MusicPDFWriteObject{4780}{<< /Type /Annot /Subtype /Link /Border [0 0 0] /H /I /C [0 1 0] /Rect [73.77 501.593 116.794 510.461] /A << /S /GoTo /D <636974652E426F752D4754> >> >>}
\MusicPDFWriteObject{4781}{<< /Type /Annot /Subtype /Link /Border [0 0 0] /H /I /C [1 0 0] /Rect [262.034 337.588 295.493 349.541] /A << /S /GoTo /D <74686D2E442E312E322E31> >> >>}
\MusicPDFWriteObject{4782}{<< /Type /Annot /Subtype /Link /Border [0 0 0] /H /I /C [0 1 0] /Rect [508.104 340.078 543.159 348.946] /A << /S /GoTo /D <636974652E426F752D4949> >> >>}
\MusicPDFWriteObject{4783}{<< /Type /Annot /Subtype /Link /Border [0 0 0] /H /I /C [0 1 0] /Rect [87.604 85.086 105.906 93.954] /A << /S /GoTo /D <636974652E4445> >> >>}
\MusicPDFWriteObject{4784}{<< /Type /Annot /Subtype /Link /Border [0 0 0] /H /I /C [1 0 0] /Rect [115.05 671.482 150.243 683.435] /A << /S /GoTo /D <74686D2E31322E322E312E32> >> >>}
\MusicPDFWriteObject{4785}{<< /Type /Annot /Subtype /Link /Border [0 0 0] /H /I /C [1 0 0] /Rect [115.05 553.126 150.243 565.078] /A << /S /GoTo /D <74686D2E31322E322E312E32> >> >>}
\MusicPDFWriteObject{4786}{<< /Type /Annot /Subtype /Link /Border [0 0 0] /H /I /C [1 0 0] /Rect [449.375 504.973 483.323 517.98] /A << /S /GoTo /D <74686D2E442E312E322E31> >> >>}
\MusicPDFWriteObject{4787}{<< /Type /Annot /Subtype /Link /Border [0 0 0] /H /I /C [1 0 0] /Rect [321.832 493.427 358.442 504.218] /A << /S /GoTo /D <74686D2E31322E322E322E31> >> >>}
\MusicPDFWriteObject{4788}{<< /Type /Annot /Subtype /Link /Border [0 0 0] /H /I /C [0 1 0] /Rect [184.267 580.987 201.421 589.855] /A << /S /GoTo /D <636974652E466F6C> >> >>}
\MusicPDFWriteObject{4789}{<< /Type /Annot /Subtype /Link /Border [0 0 0] /H /I /C [0 1 0] /Rect [276.918 580.987 307.311 589.855] /A << /S /GoTo /D <636974652E426F752D49> >> >>}
\MusicPDFWriteObject{4790}{<< /Type /Annot /Subtype /Link /Border [0 0 0] /H /I /C [0 1 0] /Rect [142.018 567.836 172.411 576.704] /A << /S /GoTo /D <636974652E426F752D49> >> >>}
\MusicPDFWriteObject{4791}{<< /Type /Annot /Subtype /Link /Border [0 0 0] /H /I /C [0 1 0] /Rect [407.877 554.685 438.27 563.553] /A << /S /GoTo /D <636974652E426F752D49> >> >>}
\MusicPDFWriteObject{4792}{<< /Type /Annot /Subtype /Link /Border [0 0 0] /H /I /C [0 1 0] /Rect [512.766 541.535 543.159 550.402] /A << /S /GoTo /D <636974652E426F752D49> >> >>}
\MusicPDFWriteObject{4793}{<< /Type /Annot /Subtype /Link /Border [0 0 0] /H /I /C [1 0 0] /Rect [264.78 493.042 299.973 505.635] /A << /S /GoTo /D <74686D2E31322E322E312E31> >> >>}
\MusicPDFWriteObject{4794}{<< /Type /Annot /Subtype /Link /Border [0 0 0] /H /I /C [1 0 0] /Rect [287.666 368.888 327.839 380.841] /A << /S /GoTo /D <6571756174696F6E2E31322E322E322E3130> >> >>}
\MusicPDFWriteObject{4795}{<< /Type /Annot /Subtype /Link /Border [0 0 0] /H /I /C [1 0 0] /Rect [347.077 332.46 387.251 344.413] /A << /S /GoTo /D <6571756174696F6E2E31322E322E322E3130> >> >>}
\MusicPDFWriteObject{4796}{<< /Type /Annot /Subtype /Link /Border [0 0 0] /H /I /C [1 0 0] /Rect [160.85 319.31 196.044 331.262] /A << /S /GoTo /D <6571756174696F6E2E31322E322E322E31> >> >>}
\MusicPDFWriteObject{4797}{<< /Type /Annot /Subtype /Link /Border [0 0 0] /H /I /C [0 1 0] /Rect [304.35 321.8 339.404 330.667] /A << /S /GoTo /D <636974652E426F752D4949> >> >>}
\MusicPDFWriteObject{4798}{<< /Type /Annot /Subtype /Link /Border [0 0 0] /H /I /C [0 1 0] /Rect [382.224 308.649 417.279 317.517] /A << /S /GoTo /D <636974652E426F752D4949> >> >>}
\MusicPDFWriteObject{4799}{<< /Type /Annot /Subtype /Link /Border [0 0 0] /H /I /C [0 1 0] /Rect [382.412 256.046 417.467 264.914] /A << /S /GoTo /D <636974652E426F752D4949> >> >>}
\MusicPDFWriteObject{4800}{<< /Type /Annot /Subtype /Link /Border [0 0 0] /H /I /C [0 1 0] /Rect [106.354 229.744 141.409 238.612] /A << /S /GoTo /D <636974652E426F752D4949> >> >>}
\MusicPDFWriteObject{4801}{<< /Type /Annot /Subtype /Link /Border [0 0 0] /H /I /C [1 0 0] /Rect [130.512 186.207 163.829 199.756] /A << /S /GoTo /D <74686D2E412E352E342E31> >> >>}
\MusicPDFWriteObject{4802}{<< /Type /Annot /Subtype /Link /Border [0 0 0] /H /I /C [1 0 0] /Rect [186.226 174.651 218.93 186.606] /A << /S /GoTo /D <74686D2E412E352E342E31> >> >>}
\MusicPDFWriteObject{4803}{<< /Type /Annot /Subtype /Link /Border [0 0 0] /H /I /C [0 1 0] /Rect [508.104 137.689 543.159 146.557] /A << /S /GoTo /D <636974652E426F752D4949> >> >>}
\MusicPDFWriteObject{4804}{<< /Type /Annot /Subtype /Link /Border [0 0 0] /H /I /C [0 1 0] /Rect [293.489 687.123 328.544 695.991] /A << /S /GoTo /D <636974652E426F752D4949> >> >>}
\MusicPDFWriteObject{4805}{<< /Type /Annot /Subtype /Link /Border [0 0 0] /H /I /C [0 1 0] /Rect [73.79 673.973 104.183 682.84] /A << /S /GoTo /D <636974652E426F752D49> >> >>}
\MusicPDFWriteObject{4806}{<< /Type /Annot /Subtype /Link /Border [0 0 0] /H /I /C [1 0 0] /Rect [186.823 645.181 219.527 657.134] /A << /S /GoTo /D <74686D2E412E352E332E33> >> >>}
\MusicPDFWriteObject{4807}{<< /Type /Annot /Subtype /Link /Border [0 0 0] /H /I /C [1 0 0] /Rect [441.971 579.427 477.164 591.38] /A << /S /GoTo /D <6571756174696F6E2E31322E322E322E33> >> >>}
\MusicPDFWriteObject{4808}{<< /Type /Annot /Subtype /Link /Border [0 0 0] /H /I /C [1 0 0] /Rect [506.695 579.427 541.888 591.38] /A << /S /GoTo /D <6571756174696F6E2E31322E322E322E34> >> >>}
\MusicPDFWriteObject{4809}{<< /Type /Annot /Subtype /Link /Border [0 0 0] /H /I /C [1 0 0] /Rect [178.651 553.126 213.845 565.078] /A << /S /GoTo /D <6571756174696F6E2E31322E322E322E33> >> >>}
\MusicPDFWriteObject{4810}{<< /Type /Annot /Subtype /Link /Border [0 0 0] /H /I /C [1 0 0] /Rect [268.12 402.182 301.579 414.137] /A << /S /GoTo /D <74686D2E442E312E322E31> >> >>}
\MusicPDFWriteObject{4811}{<< /Type /Annot /Subtype /Link /Border [0 0 0] /H /I /C [1 0 0] /Rect [207.381 349.579 247.554 361.532] /A << /S /GoTo /D <6571756174696F6E2E31322E322E322E3134> >> >>}
\MusicPDFWriteObject{4812}{<< /Type /Annot /Subtype /Link /Border [0 0 0] /H /I /C [1 0 0] /Rect [307.916 242.229 343.74 254.182] /A << /S /GoTo /D <74686D2E31322E322E312E31> >> >>}
\MusicPDFWriteObject{4813}{<< /Type /Annot /Subtype /Link /Border [0 0 0] /H /I /C [1 0 0] /Rect [439.828 242.229 480.001 254.182] /A << /S /GoTo /D <6571756174696F6E2E31322E322E322E3136> >> >>}
\MusicPDFWriteObject{4814}{<< /Type /Annot /Subtype /Link /Border [0 0 0] /H /I /C [1 0 0] /Rect [289.984 202.777 330.158 214.73] /A << /S /GoTo /D <6571756174696F6E2E31322E322E322E3134> >> >>}
\MusicPDFWriteObject{4815}{<< /Type /Annot /Subtype /Link /Border [0 0 0] /H /I /C [0 1 0] /Rect [89.402 159.995 124.457 168.863] /A << /S /GoTo /D <636974652E426F752D4949> >> >>}
\MusicPDFWriteObject{4816}{<< /Type /Annot /Subtype /Link /Border [0 0 0] /H /I /C [1 0 0] /Rect [360.875 157.505 401.048 169.458] /A << /S /GoTo /D <6571756174696F6E2E31322E322E322E3137> >> >>}
\MusicPDFWriteObject{4817}{<< /Type /Annot /Subtype /Link /Border [0 0 0] /H /I /C [0 1 0] /Rect [262.6 133.09 297.655 141.958] /A << /S /GoTo /D <636974652E426F752D4949> >> >>}
\MusicPDFWriteObject{4818}{<< /Type /Annot /Subtype /Link /Border [0 0 0] /H /I /C [0 1 0] /Rect [87.604 119.94 122.659 128.808] /A << /S /GoTo /D <636974652E426F752D4949> >> >>}
\MusicPDFWriteObject{4819}{<< /Type /Annot /Subtype /Link /Border [0 0 0] /H /I /C [1 0 0] /Rect [178.712 684.633 218.886 696.586] /A << /S /GoTo /D <6571756174696F6E2E31322E322E322E3138> >> >>}
\MusicPDFWriteObject{4820}{<< /Type /Annot /Subtype /Link /Border [0 0 0] /H /I /C [1 0 0] /Rect [335.669 684.633 370.862 696.586] /A << /S /GoTo /D <6571756174696F6E2E31322E322E322E34> >> >>}
\MusicPDFWriteObject{4821}{<< /Type /Annot /Subtype /Link /Border [0 0 0] /H /I /C [1 0 0] /Rect [194.57 587.515 229.763 599.468] /A << /S /GoTo /D <6571756174696F6E2E31322E322E322E34> >> >>}
\MusicPDFWriteObject{4822}{<< /Type /Annot /Subtype /Link /Border [0 0 0] /H /I /C [0 1 0] /Rect [178.455 521.396 195.609 530.264] /A << /S /GoTo /D <636974652E466F6C> >> >>}
\MusicPDFWriteObject{4823}{<< /Type /Annot /Subtype /Link /Border [0 0 0] /H /I /C [0 1 0] /Rect [443.102 521.396 460.256 530.264] /A << /S /GoTo /D <636974652E466F6C> >> >>}
\MusicPDFWriteObject{4824}{<< /Type /Annot /Subtype /Link /Border [0 0 0] /H /I /C [1 0 0] /Rect [86.497 409.598 121.691 421.551] /A << /S /GoTo /D <6571756174696F6E2E31322E322E322E34> >> >>}
\MusicPDFWriteObject{4825}{<< /Type /Annot /Subtype /Link /Border [0 0 0] /H /I /C [1 0 0] /Rect [183.084 409.598 223.257 421.551] /A << /S /GoTo /D <6571756174696F6E2E31322E322E322E3231> >> >>}
\MusicPDFWriteObject{4826}{<< /Type /Annot /Subtype /Link /Border [0 0 0] /H /I /C [1 0 0] /Rect [70.625 293.999 103.328 305.952] /A << /S /GoTo /D <74686D2E412E352E342E33> >> >>}
\MusicPDFWriteObject{4827}{<< /Type /Annot /Subtype /Link /Border [0 0 0] /H /I /C [1 0 0] /Rect [233.355 247.927 266.058 259.879] /A << /S /GoTo /D <74686D2E412E352E342E33> >> >>}
\MusicPDFWriteObject{4828}{<< /Type /Annot /Subtype /Link /Border [0 0 0] /H /I /C [1 0 0] /Rect [437.263 247.927 477.436 259.879] /A << /S /GoTo /D <6571756174696F6E2E31322E322E322E3232> >> >>}
\MusicPDFWriteObject{4829}{<< /Type /Annot /Subtype /Link /Border [0 0 0] /H /I /C [1 0 0] /Rect [159.29 201.854 199.463 213.807] /A << /S /GoTo /D <6571756174696F6E2E31322E322E322E3233> >> >>}
\MusicPDFWriteObject{4830}{<< /Type /Annot /Subtype /Link /Border [0 0 0] /H /I /C [0 1 0] /Rect [375.046 149.797 410.101 158.665] /A << /S /GoTo /D <636974652E426F752D4949> >> >>}
\MusicPDFWriteObject{4831}{<< /Type /Annot /Subtype /Link /Border [0 0 0] /H /I /C [0 1 0] /Rect [133.162 660.822 168.217 669.69] /A << /S /GoTo /D <636974652E426F752D4949> >> >>}
\MusicPDFWriteObject{4832}{<< /Type /Annot /Subtype /Link /Border [0 0 0] /H /I /C [0 1 0] /Rect [387.57 660.822 417.963 669.69] /A << /S /GoTo /D <636974652E426F752D49> >> >>}
\MusicPDFWriteObject{4833}{<< /Type /Annot /Subtype /Link /Border [0 0 0] /H /I /C [1 0 0] /Rect [244.046 645.181 276.718 657.134] /A << /S /GoTo /D <74686D2E412E352E332E33> >> >>}
\MusicPDFWriteObject{4834}{<< /Type /Annot /Subtype /Link /Border [0 0 0] /H /I /C [1 0 0] /Rect [479.496 645.181 514.69 657.134] /A << /S /GoTo /D <6571756174696F6E2E31322E322E322E39> >> >>}
\MusicPDFWriteObject{4835}{<< /Type /Annot /Subtype /Link /Border [0 0 0] /H /I /C [1 0 0] /Rect [399.022 584.81 432.97 596.765] /A << /S /GoTo /D <74686D2E442E312E322E31> >> >>}
\MusicPDFWriteObject{4836}{<< /Type /Annot /Subtype /Link /Border [0 0 0] /H /I /C [1 0 0] /Rect [321.832 572.212 358.442 583.004] /A << /S /GoTo /D <74686D2E31322E322E322E31> >> >>}
\MusicPDFWriteObject{4837}{<< /Type /Annot /Subtype /Link /Border [0 0 0] /H /I /C [0 1 0] /Rect [388.805 302.599 407.107 311.467] /A << /S /GoTo /D <636974652E4445> >> >>}
\MusicPDFWriteObject{4838}{<< /Type /Annot /Subtype /Link /Border [0 0 0] /H /I /C [1 0 0] /Rect [86.497 171.626 126.671 183.579] /A << /S /GoTo /D <6571756174696F6E2E31322E322E322E3238> >> >>}
\MusicPDFWriteObject{4839}{<< /Type /Annot /Subtype /Link /Border [0 0 0] /H /I /C [1 0 0] /Rect [73.712 145.325 113.885 157.278] /A << /S /GoTo /D <6571756174696F6E2E31322E322E322E3238> >> >>}
\MusicPDFWriteObject{4840}{<< /Type /Annot /Subtype /Link /Border [0 0 0] /H /I /C [0 1 0] /Rect [202.815 85.086 221.117 93.954] /A << /S /GoTo /D <636974652E4445> >> >>}
\MusicPDFWriteObject{4841}{<< /Type /Annot /Subtype /Link /Border [0 0 0] /H /I /C [1 0 0] /Rect [425.511 684.633 465.684 696.586] /A << /S /GoTo /D <6571756174696F6E2E31322E322E322E3239> >> >>}
\MusicPDFWriteObject{4842}{<< /Type /Annot /Subtype /Link /Border [0 0 0] /H /I /C [0 1 0] /Rect [391.373 624.394 409.675 633.262] /A << /S /GoTo /D <636974652E4445> >> >>}
\MusicPDFWriteObject{4843}{<< /Type /Annot /Subtype /Link /Border [0 0 0] /H /I /C [0 1 0] /Rect [524.857 624.394 543.159 633.262] /A << /S /GoTo /D <636974652E4445> >> >>}
\MusicPDFWriteObject{4844}{<< /Type /Annot /Subtype /Link /Border [0 0 0] /H /I /C [1 0 0] /Rect [186.957 272.239 227.13 284.192] /A << /S /GoTo /D <6571756174696F6E2E31322E322E322E3331> >> >>}
\MusicPDFWriteObject{4845}{<< /Type /Annot /Subtype /Link /Border [0 0 0] /H /I /C [1 0 0] /Rect [70.506 219.636 105.699 231.589] /A << /S /GoTo /D <74686D2E31322E322E322E32> >> >>}
\MusicPDFWriteObject{4846}{<< /Type /Annot /Subtype /Link /Border [0 0 0] /H /I /C [1 0 0] /Rect [417.808 87.891 457.982 99.844] /A << /S /GoTo /D <6571756174696F6E2E31322E322E322E3236> >> >>}
\MusicPDFWriteObject{4847}{<< /Type /Annot /Subtype /Link /Border [0 0 0] /H /I /C [1 0 0] /Rect [88.936 667.995 94.899 678.362] /A << /S /GoTo /D <48666F6F746E6F74652E3436> >> >>}
\MusicPDFWriteObject{4848}{<< /Type /Annot /Subtype /Link /Border [0 0 0] /H /I /C [0 1 0] /Rect [444.351 612.977 460.343 621.845] /A << /S /GoTo /D <636974652E436F> >> >>}
\MusicPDFWriteObject{4849}{<< /Type /Annot /Subtype /Link /Border [0 0 0] /H /I /C [0 1 0] /Rect [209.059 599.826 232.902 608.694] /A << /S /GoTo /D <636974652E4D6531> >> >>}
\MusicPDFWriteObject{4850}{<< /Type /Annot /Subtype /Link /Border [0 0 0] /H /I /C [0 1 0] /Rect [264.315 547.223 280.307 556.091] /A << /S /GoTo /D <636974652E436F> >> >>}
\MusicPDFWriteObject{4851}{<< /Type /Annot /Subtype /Link /Border [0 0 0] /H /I /C [1 0 0] /Rect [73.712 268.28 108.905 280.233] /A << /S /GoTo /D <6571756174696F6E2E31322E342E312E34> >> >>}
\MusicPDFWriteObject{4852}{<< /Type /Annot /Subtype /Link /Border [0 0 0] /H /I /C [1 0 0] /Rect [177.08 244.469 207.293 253.323] /A << /S /GoTo /D <74686D2E372E342E322E31> >> >>}
\MusicPDFWriteObject{4853}{<< /Type /Annot /Subtype /Link /Border [0 0 0] /H /I /C [1 0 0] /Rect [270.07 244.469 300.283 253.323] /A << /S /GoTo /D <74686D2E372E342E332E32> >> >>}
\MusicPDFWriteObject{4854}{<< /Type /Annot /Subtype /Link /Border [0 0 0] /H /I /C [1 0 0] /Rect [370.132 211.886 402.42 224.301] /A << /S /GoTo /D <74686D2E422E342E342E32> >> >>}
\MusicPDFWriteObject{4855}{<< /Type /Annot /Subtype /Link /Border [0 0 0] /H /I /C [1 0 0] /Rect [117.346 168.337 147.559 178.704] /A << /S /GoTo /D <74686D2E392E332E332E31> >> >>}
\MusicPDFWriteObject{4856}{<< /Type /Annot /Subtype /Link /Border [0 0 0] /H /I /C [1 0 0] /Rect [148.334 168.337 154.296 178.704] /A << /S /GoTo /D <48666F6F746E6F74652E3437> >> >>}
\MusicPDFWriteObject{4857}{<< /Type /Annot /Subtype /Link /Border [0 0 0] /H /I /C [1 0 0] /Rect [143.299 106.612 169.135 118.2] /A << /S /GoTo /D <74686D2E392E332E332E31> >> >>}
\MusicPDFWriteObject{4858}{<< /Type /Annot /Subtype /Link /Border [0 0 0] /H /I /C [1 0 0] /Rect [114.988 83.094 149.434 93.055] /A << /S /GoTo /D <6571756174696F6E2E31322E342E312E3233> >> >>}
\MusicPDFWriteObject{4859}{<< /Type /Annot /Subtype /Link /Border [0 0 0] /H /I /C [1 0 0] /Rect [504.547 604.017 510.51 617.766] /A << /S /GoTo /D <48666F6F746E6F74652E3438> >> >>}
\MusicPDFWriteObject{4860}{<< /Type /Annot /Subtype /Link /Border [0 0 0] /H /I /C [1 0 0] /Rect [503.104 496.445 533.317 513.127] /A << /S /GoTo /D <74686D2E352E332E342E34> >> >>}
\MusicPDFWriteObject{4861}{<< /Type /Annot /Subtype /Link /Border [0 0 0] /H /I /C [1 0 0] /Rect [531.324 496.445 537.287 513.127] /A << /S /GoTo /D <48666F6F746E6F74652E3439> >> >>}
\MusicPDFWriteObject{4862}{<< /Type /Annot /Subtype /Link /Border [0 0 0] /H /I /C [1 0 0] /Rect [140.121 484.728 172.547 497.043] /A << /S /GoTo /D <74686D2E432E332E322E32> >> >>}
\MusicPDFWriteObject{4863}{<< /Type /Annot /Subtype /Link /Border [0 0 0] /H /I /C [0 1 0] /Rect [376.937 450.049 392.929 458.917] /A << /S /GoTo /D <636974652E436F> >> >>}
\MusicPDFWriteObject{4864}{<< /Type /Annot /Subtype /Link /Border [0 0 0] /H /I /C [0 1 0] /Rect [73.77 436.899 97.613 445.766] /A << /S /GoTo /D <636974652E4D6531> >> >>}
\MusicPDFWriteObject{4865}{<< /Type /Annot /Subtype /Link /Border [0 0 0] /H /I /C [1 0 0] /Rect [218.11 341.499 248.324 353.452] /A << /S /GoTo /D <74686D2E342E322E312E35> >> >>}
\MusicPDFWriteObject{4866}{<< /Type /Annot /Subtype /Link /Border [0 0 0] /H /I /C [1 0 0] /Rect [124.128 250.605 158.989 261.453] /A << /S /GoTo /D <74686D2E31322E322E322E32> >> >>}
\MusicPDFWriteObject{4867}{<< /Type /Annot /Subtype /Link /Border [0 0 0] /H /I /C [1 0 0] /Rect [188.097 250.605 215.289 261.453] /A << /S /GoTo /D <73756273656374696F6E2E31322E322E32> >> >>}
\MusicPDFWriteObject{4868}{<< /Type /Annot /Subtype /Link /Border [0 0 0] /H /I /C [1 0 0] /Rect [239.732 684.273 274.527 696.588] /A << /S /GoTo /D <74686D2E31322E322E322E32> >> >>}
\MusicPDFWriteObject{4869}{<< /Type /Annot /Subtype /Link /Border [0 0 0] /H /I /C [1 0 0] /Rect [303.594 684.273 330.735 696.588] /A << /S /GoTo /D <73756273656374696F6E2E31322E322E32> >> >>}
\MusicPDFWriteObject{4870}{<< /Type /Annot /Subtype /Link /Border [0 0 0] /H /I /C [1 0 0] /Rect [162.639 671.482 197.832 683.435] /A << /S /GoTo /D <6571756174696F6E2E31322E342E312E34> >> >>}
\MusicPDFWriteObject{4871}{<< /Type /Annot /Subtype /Link /Border [0 0 0] /H /I /C [1 0 0] /Rect [458.383 495.769 490.962 507.722] /A << /S /GoTo /D <74686D2E432E342E332E35> >> >>}
\MusicPDFWriteObject{4872}{<< /Type /Annot /Subtype /Link /Border [0 0 0] /H /I /C [1 0 0] /Rect [178.382 430.015 213.576 441.968] /A << /S /GoTo /D <6571756174696F6E2E31322E342E312E35> >> >>}
\MusicPDFWriteObject{4873}{<< /Type /Annot /Subtype /Link /Border [0 0 0] /H /I /C [1 0 0] /Rect [452.478 415.177 489.884 428.837] /A << /S /GoTo /D <74686D2E432E322E342E3139> >> >>}
\MusicPDFWriteObject{4874}{<< /Type /Annot /Subtype /Link /Border [0 0 0] /H /I /C [1 0 0] /Rect [400.004 352.809 432.093 367.359] /A << /S /GoTo /D <74686D2E412E352E332E37> >> >>}
\MusicPDFWriteObject{4875}{<< /Type /Annot /Subtype /Link /Border [0 0 0] /H /I /C [1 0 0] /Rect [173.1 523.115 213.274 535.068] /A << /S /GoTo /D <6571756174696F6E2E31322E342E312E3136> >> >>}
\MusicPDFWriteObject{4876}{<< /Type /Annot /Subtype /Link /Border [0 0 0] /H /I /C [1 0 0] /Rect [378.329 473.493 418.503 485.446] /A << /S /GoTo /D <6571756174696F6E2E31322E342E312E3138> >> >>}
\MusicPDFWriteObject{4877}{<< /Type /Annot /Subtype /Link /Border [0 0 0] /H /I /C [1 0 0] /Rect [128.054 380.21 168.227 392.162] /A << /S /GoTo /D <6571756174696F6E2E31322E342E312E3139> >> >>}
\MusicPDFWriteObject{4878}{<< /Type /Annot /Subtype /Link /Border [0 0 0] /H /I /C [1 0 0] /Rect [140.342 333.523 180.516 345.476] /A << /S /GoTo /D <6571756174696F6E2E31322E342E312E3132> >> >>}
\MusicPDFWriteObject{4879}{<< /Type /Annot /Subtype /Link /Border [0 0 0] /H /I /C [1 0 0] /Rect [269.327 333.523 309.5 345.476] /A << /S /GoTo /D <6571756174696F6E2E31322E342E312E3131> >> >>}
\MusicPDFWriteObject{4880}{<< /Type /Annot /Subtype /Link /Border [0 0 0] /H /I /C [1 0 0] /Rect [148.999 242.023 189.172 253.976] /A << /S /GoTo /D <6571756174696F6E2E31322E342E312E3232> >> >>}
\MusicPDFWriteObject{4881}{<< /Type /Annot /Subtype /Link /Border [0 0 0] /H /I /C [0 1 0] /Rect [185.257 687.123 201.25 695.991] /A << /S /GoTo /D <636974652E436F> >> >>}
\MusicPDFWriteObject{4882}{<< /Type /Annot /Subtype /Link /Border [0 0 0] /H /I /C [1 0 0] /Rect [422.043 684.633 454.058 696.586] /A << /S /GoTo /D <74686D2E422E342E342E33> >> >>}
\MusicPDFWriteObject{4883}{<< /Type /Annot /Subtype /Link /Border [0 0 0] /H /I /C [1 0 0] /Rect [493.37 684.633 525.385 696.586] /A << /S /GoTo /D <74686D2E422E342E342E35> >> >>}
\MusicPDFWriteObject{4884}{<< /Type /Annot /Subtype /Link /Border [0 0 0] /H /I /C [1 0 0] /Rect [106.556 669.689 138.844 684.544] /A << /S /GoTo /D <74686D2E422E342E332E32> >> >>}
\MusicPDFWriteObject{4885}{<< /Type /Annot /Subtype /Link /Border [0 0 0] /H /I /C [1 0 0] /Rect [383.16 671.482 423.333 683.435] /A << /S /GoTo /D <6571756174696F6E2E31322E342E312E3233> >> >>}
\MusicPDFWriteObject{4886}{<< /Type /Annot /Subtype /Link /Border [0 0 0] /H /I /C [1 0 0] /Rect [274.705 514.178 314.879 526.131] /A << /S /GoTo /D <6571756174696F6E2E31322E342E312E3235> >> >>}
\MusicPDFWriteObject{4887}{<< /Type /Annot /Subtype /Link /Border [0 0 0] /H /I /C [1 0 0] /Rect [240.324 501.028 275.517 512.98] /A << /S /GoTo /D <6571756174696F6E2E31322E342E312E34> >> >>}
\MusicPDFWriteObject{4888}{<< /Type /Annot /Subtype /Link /Border [0 0 0] /H /I /C [1 0 0] /Rect [335.996 404.43 366.209 416.385] /A << /S /GoTo /D <74686D2E392E332E332E31> >> >>}
\MusicPDFWriteObject{4889}{<< /Type /Annot /Subtype /Link /Border [0 0 0] /H /I /C [1 0 0] /Rect [313.607 586.003 348.8 597.955] /A << /S /GoTo /D <6571756174696F6E2E31312E332E312E34> >> >>}
\MusicPDFWriteObject{4890}{<< /Type /Annot /Subtype /Link /Border [0 0 0] /H /I /C [1 0 0] /Rect [145.746 529.927 180.674 544.489] /A << /S /GoTo /D <74686D2E31312E332E322E31> >> >>}
\MusicPDFWriteObject{4891}{<< /Type /Annot /Subtype /Link /Border [0 0 0] /H /I /C [1 0 0] /Rect [149.409 503.625 185.266 518.187] /A << /S /GoTo /D <74686D2E31322E342E312E31> >> >>}
\MusicPDFWriteObject{4892}{<< /Type /Annot /Subtype /Link /Border [0 0 0] /H /I /C [1 0 0] /Rect [479.052 505.679 514.246 517.632] /A << /S /GoTo /D <6571756174696F6E2E31322E352E312E31> >> >>}
\MusicPDFWriteObject{4893}{<< /Type /Annot /Subtype /Link /Border [0 0 0] /H /I /C [1 0 0] /Rect [314.226 426.685 349.42 438.638] /A << /S /GoTo /D <74686D2E31302E342E332E37> >> >>}
\MusicPDFWriteObject{4894}{<< /Type /Annot /Subtype /Link /Border [0 0 0] /H /I /C [1 0 0] /Rect [429.903 426.685 465.096 438.638] /A << /S /GoTo /D <74686D2E31302E342E322E38> >> >>}
\MusicPDFWriteObject{4895}{<< /Type /Annot /Subtype /Link /Border [0 0 0] /H /I /C [1 0 0] /Rect [167.926 413.534 203.12 425.487] /A << /S /GoTo /D <6571756174696F6E2E31322E352E312E34> >> >>}
\MusicPDFWriteObject{4896}{<< /Type /Annot /Subtype /Link /Border [0 0 0] /H /I /C [1 0 0] /Rect [432.73 413.534 467.923 425.487] /A << /S /GoTo /D <6571756174696F6E2E31322E352E312E31> >> >>}
\MusicPDFWriteObject{4897}{<< /Type /Annot /Subtype /Link /Border [0 0 0] /H /I /C [1 0 0] /Rect [82.959 558.462 283.602 569.252] /A << /S /GoTo /D <73656374696F6E2E31332E31> >> >>}
\MusicPDFWriteObject{4898}{<< /Type /Annot /Subtype /Link /Border [0 0 0] /H /I /C [1 0 0] /Rect [82.959 545.311 281.029 556.102] /A << /S /GoTo /D <73656374696F6E2E31332E32> >> >>}
\MusicPDFWriteObject{4899}{<< /Type /Annot /Subtype /Link /Border [0 0 0] /H /I /C [1 0 0] /Rect [82.959 534.097 306.179 542.951] /A << /S /GoTo /D <73656374696F6E2E31332E33> >> >>}
\MusicPDFWriteObject{4900}{<< /Type /Annot /Subtype /Link /Border [0 0 0] /H /I /C [1 0 0] /Rect [82.959 518.455 337.138 530.411] /A << /S /GoTo /D <73756273656374696F6E2E31332E332E31> >> >>}
\MusicPDFWriteObject{4901}{<< /Type /Annot /Subtype /Link /Border [0 0 0] /H /I /C [1 0 0] /Rect [336.127 672.036 371.32 682.827] /A << /S /GoTo /D <74686D2E31322E352E312E31> >> >>}
\MusicPDFWriteObject{4902}{<< /Type /Annot /Subtype /Link /Border [0 0 0] /H /I /C [1 0 0] /Rect [397.011 579.981 432.005 590.771] /A << /S /GoTo /D <74686D2E31322E352E312E31> >> >>}
\MusicPDFWriteObject{4903}{<< /Type /Annot /Subtype /Link /Border [0 0 0] /H /I /C [1 0 0] /Rect [262.24 534.551 297.433 545.342] /A << /S /GoTo /D <74686D2E31322E352E312E31> >> >>}
\MusicPDFWriteObject{4904}{<< /Type /Annot /Subtype /Link /Border [0 0 0] /H /I /C [1 0 0] /Rect [359.461 514.271 365.424 527.129] /A << /S /GoTo /D <48666F6F746E6F74652E3530> >> >>}
\MusicPDFWriteObject{4905}{<< /Type /Annot /Subtype /Link /Border [0 0 0] /H /I /C [1 0 0] /Rect [396.067 202.619 431.261 214.572] /A << /S /GoTo /D <6571756174696F6E2E31332E332E312E32> >> >>}
\MusicPDFWriteObject{4906}{<< /Type /Annot /Subtype /Link /Border [0 0 0] /H /I /C [1 0 0] /Rect [205.997 170.295 241.19 181.1] /A << /S /GoTo /D <74686D2E31332E332E312E31> >> >>}
\MusicPDFWriteObject{4907}{<< /Type /Annot /Subtype /Link /Border [0 0 0] /H /I /C [1 0 0] /Rect [266.12 150.015 272.083 162.873] /A << /S /GoTo /D <48666F6F746E6F74652E3531> >> >>}
\MusicPDFWriteObject{4908}{<< /Type /Annot /Subtype /Link /Border [0 0 0] /H /I /C [1 0 0] /Rect [88.936 132.78 94.899 143.147] /A << /S /GoTo /D <48666F6F746E6F74652E3532> >> >>}
\MusicPDFWriteObject{4909}{<< /Type /Annot /Subtype /Link /Border [0 0 0] /H /I /C [1 0 0] /Rect [82.959 558.462 258.84 569.252] /A << /S /GoTo /D <73656374696F6E2E31342E31> >> >>}
\MusicPDFWriteObject{4910}{<< /Type /Annot /Subtype /Link /Border [0 0 0] /H /I /C [1 0 0] /Rect [82.959 547.248 222.404 556.102] /A << /S /GoTo /D <73756273656374696F6E2E31342E312E31> >> >>}
\MusicPDFWriteObject{4911}{<< /Type /Annot /Subtype /Link /Border [0 0 0] /H /I /C [1 0 0] /Rect [82.959 532.16 250.125 542.951] /A << /S /GoTo /D <73756273656374696F6E2E31342E312E32> >> >>}
\MusicPDFWriteObject{4912}{<< /Type /Annot /Subtype /Link /Border [0 0 0] /H /I /C [1 0 0] /Rect [82.959 519.009 368.264 529.8] /A << /S /GoTo /D <73656374696F6E2E31342E32> >> >>}
\MusicPDFWriteObject{4913}{<< /Type /Annot /Subtype /Link /Border [0 0 0] /H /I /C [1 0 0] /Rect [82.959 505.859 207.075 516.649] /A << /S /GoTo /D <73656374696F6E2E31342E33> >> >>}
\MusicPDFWriteObject{4914}{<< /Type /Annot /Subtype /Link /Border [0 0 0] /H /I /C [1 0 0] /Rect [248.51 652.907 283.039 663.698] /A << /S /GoTo /D <74686D2E31332E332E312E32> >> >>}
\MusicPDFWriteObject{4915}{<< /Type /Annot /Subtype /Link /Border [0 0 0] /H /I /C [1 0 0] /Rect [426.76 197.928 461.953 209.881] /A << /S /GoTo /D <6571756174696F6E2E31342E312E312E32> >> >>}
\MusicPDFWriteObject{4916}{<< /Type /Annot /Subtype /Link /Border [0 0 0] /H /I /C [0 1 0] /Rect [73.723 160.966 98.002 169.834] /A << /S /GoTo /D <636974652E54617465> >> >>}
\MusicPDFWriteObject{4917}{<< /Type /Annot /Subtype /Link /Border [0 0 0] /H /I /C [0 1 0] /Rect [194.8 687.123 219.079 695.991] /A << /S /GoTo /D <636974652E54617465> >> >>}
\MusicPDFWriteObject{4918}{<< /Type /Annot /Subtype /Link /Border [0 0 0] /H /I /C [0 1 0] /Rect [525.493 687.123 543.159 695.991] /A << /S /GoTo /D <636974652E4154> >> >>}
\MusicPDFWriteObject{4919}{<< /Type /Annot /Subtype /Link /Border [0 0 0] /H /I /C [0 1 0] /Rect [191.503 673.973 209.169 682.84] /A << /S /GoTo /D <636974652E4154> >> >>}
\MusicPDFWriteObject{4920}{<< /Type /Annot /Subtype /Link /Border [0 0 0] /H /I /C [0 1 0] /Rect [93.441 660.822 127.002 669.69] /A << /S /GoTo /D <636974652E5765696C2D65> >> >>}
\MusicPDFWriteObject{4921}{<< /Type /Annot /Subtype /Link /Border [0 0 0] /H /I /C [0 1 0] /Rect [372.858 660.822 404.669 669.69] /A << /S /GoTo /D <636974652E5765696C2D66> >> >>}
\MusicPDFWriteObject{4922}{<< /Type /Annot /Subtype /Link /Border [0 0 0] /H /I /C [0 1 0] /Rect [421.064 619.433 473.599 630.237] /A << /S /GoTo /D <636974652E4C616E676C616E6473> >> >>}
\MusicPDFWriteObject{4923}{<< /Type /Annot /Subtype /Link /Border [0 0 0] /H /I /C [0 1 0] /Rect [343.035 595.068 360.701 603.936] /A << /S /GoTo /D <636974652E4154> >> >>}
\MusicPDFWriteObject{4924}{<< /Type /Annot /Subtype /Link /Border [0 0 0] /H /I /C [0 1 0] /Rect [195.261 555.616 212.927 564.484] /A << /S /GoTo /D <636974652E4154> >> >>}
\MusicPDFWriteObject{4925}{<< /Type /Annot /Subtype /Link /Border [0 0 0] /H /I /C [0 1 0] /Rect [256.347 516.163 274.013 525.031] /A << /S /GoTo /D <636974652E4154> >> >>}
\MusicPDFWriteObject{4926}{<< /Type /Annot /Subtype /Link /Border [0 0 0] /H /I /C [0 1 0] /Rect [261.653 687.123 279.319 695.991] /A << /S /GoTo /D <636974652E4154> >> >>}
\MusicPDFWriteObject{4927}{<< /Type /Annot /Subtype /Link /Border [0 0 0] /H /I /C [0 1 0] /Rect [399.615 687.123 417.281 695.991] /A << /S /GoTo /D <636974652E4154> >> >>}
\MusicPDFWriteObject{4928}{<< /Type /Annot /Subtype /Link /Border [0 0 0] /H /I /C [0 1 0] /Rect [222.4 673.973 240.066 682.84] /A << /S /GoTo /D <636974652E4154> >> >>}
\MusicPDFWriteObject{4929}{<< /Type /Annot /Subtype /Link /Border [0 0 0] /H /I /C [0 1 0] /Rect [360.363 673.973 378.029 682.84] /A << /S /GoTo /D <636974652E4154> >> >>}
\MusicPDFWriteObject{4930}{<< /Type /Annot /Subtype /Link /Border [0 0 0] /H /I /C [0 1 0] /Rect [146.536 660.822 170.815 669.69] /A << /S /GoTo /D <636974652E54617465> >> >>}
\MusicPDFWriteObject{4931}{<< /Type /Annot /Subtype /Link /Border [0 0 0] /H /I /C [0 1 0] /Rect [123.418 581.729 156.979 590.597] /A << /S /GoTo /D <636974652E5765696C2D65> >> >>}
\MusicPDFWriteObject{4932}{<< /Type /Annot /Subtype /Link /Border [0 0 0] /H /I /C [0 1 0] /Rect [158.46 529.126 182.738 537.994] /A << /S /GoTo /D <636974652E54617465> >> >>}
\MusicPDFWriteObject{4933}{<< /Type /Annot /Subtype /Link /Border [0 0 0] /H /I /C [0 1 0] /Rect [229.533 458.657 282.068 469.462] /A << /S /GoTo /D <636974652E4C616E676C616E6473> >> >>}
\MusicPDFWriteObject{4934}{<< /Type /Annot /Subtype /Link /Border [0 0 0] /H /I /C [1 0 0] /Rect [438.807 406.083 474.665 416.874] /A << /S /GoTo /D <74686D2E31332E332E312E31> >> >>}
\MusicPDFWriteObject{4935}{<< /Type /Annot /Subtype /Link /Border [0 0 0] /H /I /C [1 0 0] /Rect [169.531 392.379 204.06 404.332] /A << /S /GoTo /D <74686D2E31332E332E312E32> >> >>}
\MusicPDFWriteObject{4936}{<< /Type /Annot /Subtype /Link /Border [0 0 0] /H /I /C [0 1 0] /Rect [201.666 287.726 254.201 298.531] /A << /S /GoTo /D <636974652E4C616E676C616E6473> >> >>}
\MusicPDFWriteObject{4937}{<< /Type /Annot /Subtype /Link /Border [0 0 0] /H /I /C [0 1 0] /Rect [427.211 276.512 440.816 285.38] /A << /S /GoTo /D <636974652E5365> >> >>}
\MusicPDFWriteObject{4938}{<< /Type /Annot /Subtype /Link /Border [0 0 0] /H /I /C [0 1 0] /Rect [196.635 261.425 249.17 272.23] /A << /S /GoTo /D <636974652E4C616E676C616E6473> >> >>}
\MusicPDFWriteObject{4939}{<< /Type /Annot /Subtype /Link /Border [0 0 0] /H /I /C [0 1 0] /Rect [282.362 247.259 334.897 258.063] /A << /S /GoTo /D <636974652E4C616E676C616E6473> >> >>}
\MusicPDFWriteObject{4940}{<< /Type /Annot /Subtype /Link /Border [0 0 0] /H /I /C [0 1 0] /Rect [416.947 247.259 469.482 258.063] /A << /S /GoTo /D <636974652E4C616E676C616E6473> >> >>}
\MusicPDFWriteObject{4941}{<< /Type /Annot /Subtype /Link /Border [0 0 0] /H /I /C [0 1 0] /Rect [93.761 234.108 146.296 244.913] /A << /S /GoTo /D <636974652E4C616E676C616E6473> >> >>}
\MusicPDFWriteObject{4942}{<< /Type /Annot /Subtype /Link /Border [0 0 0] /H /I /C [0 1 0] /Rect [73.77 206.791 126.305 217.596] /A << /S /GoTo /D <636974652E4C616E676C616E6473> >> >>}
\MusicPDFWriteObject{4943}{<< /Type /Annot /Subtype /Link /Border [0 0 0] /H /I /C [0 1 0] /Rect [389.309 122.601 441.845 133.406] /A << /S /GoTo /D <636974652E4C616E676C616E6473> >> >>}
\MusicPDFWriteObject{4944}{<< /Type /Annot /Subtype /Link /Border [0 0 0] /H /I /C [0 1 0] /Rect [401.335 109.451 453.87 120.255] /A << /S /GoTo /D <636974652E4C616E676C616E6473> >> >>}
\MusicPDFWriteObject{4945}{<< /Type /Annot /Subtype /Link /Border [0 0 0] /H /I /C [0 1 0] /Rect [511.348 111.387 543.159 120.255] /A << /S /GoTo /D <636974652E5765696C2D66> >> >>}
\MusicPDFWriteObject{4946}{<< /Type /Annot /Subtype /Link /Border [0 0 0] /H /I /C [1 0 0] /Rect [107.663 630.375 142.856 642.328] /A << /S /GoTo /D <6571756174696F6E2E31342E312E312E32> >> >>}
\MusicPDFWriteObject{4947}{<< /Type /Annot /Subtype /Link /Border [0 0 0] /H /I /C [0 1 0] /Rect [312.179 461.532 364.714 472.337] /A << /S /GoTo /D <636974652E4C616E676C616E6473> >> >>}
\MusicPDFWriteObject{4948}{<< /Type /Annot /Subtype /Link /Border [0 0 0] /H /I /C [1 0 0] /Rect [391.346 421.527 426.539 433.479] /A << /S /GoTo /D <6571756174696F6E2E31342E312E312E36> >> >>}
\MusicPDFWriteObject{4949}{<< /Type /Annot /Subtype /Link /Border [0 0 0] /H /I /C [1 0 0] /Rect [496.231 419.727 531.159 433.479] /A << /S /GoTo /D <74686D2E31342E312E312E31> >> >>}
\MusicPDFWriteObject{4950}{<< /Type /Annot /Subtype /Link /Border [0 0 0] /H /I /C [0 1 0] /Rect [250.509 362.954 303.044 373.758] /A << /S /GoTo /D <636974652E4C616E676C616E6473> >> >>}
\MusicPDFWriteObject{4951}{<< /Type /Annot /Subtype /Link /Border [0 0 0] /H /I /C [1 0 0] /Rect [125.061 336.098 160.553 349.382] /A << /S /GoTo /D <74686D2E31342E312E312E31> >> >>}
\MusicPDFWriteObject{4952}{<< /Type /Annot /Subtype /Link /Border [0 0 0] /H /I /C [0 1 0] /Rect [169.226 211.194 221.762 221.999] /A << /S /GoTo /D <636974652E4C616E676C616E6473> >> >>}
\MusicPDFWriteObject{4953}{<< /Type /Annot /Subtype /Link /Border [0 0 0] /H /I /C [0 1 0] /Rect [280.703 211.194 333.238 221.999] /A << /S /GoTo /D <636974652E4C616E676C616E6473> >> >>}
\MusicPDFWriteObject{4954}{<< /Type /Annot /Subtype /Link /Border [0 0 0] /H /I /C [1 0 0] /Rect [490.458 210.641 530.631 222.594] /A << /S /GoTo /D <6571756174696F6E2E31342E312E312E3130> >> >>}
\MusicPDFWriteObject{4955}{<< /Type /Annot /Subtype /Link /Border [0 0 0] /H /I /C [0 1 0] /Rect [291.731 199.98 316.009 208.848] /A << /S /GoTo /D <636974652E54617465> >> >>}
\MusicPDFWriteObject{4956}{<< /Type /Annot /Subtype /Link /Border [0 0 0] /H /I /C [0 1 0] /Rect [93.138 186.829 117.416 195.697] /A << /S /GoTo /D <636974652E54617465> >> >>}
\MusicPDFWriteObject{4957}{<< /Type /Annot /Subtype /Link /Border [0 0 0] /H /I /C [0 1 0] /Rect [121.939 93.993 174.474 104.797] /A << /S /GoTo /D <636974652E4C616E676C616E6473> >> >>}
\MusicPDFWriteObject{4958}{<< /Type /Annot /Subtype /Link /Border [0 0 0] /H /I /C [1 0 0] /Rect [216.473 684.633 256.647 696.586] /A << /S /GoTo /D <6571756174696F6E2E31342E312E312E3130> >> >>}
\MusicPDFWriteObject{4959}{<< /Type /Annot /Subtype /Link /Border [0 0 0] /H /I /C [1 0 0] /Rect [383.78 463.837 423.954 475.79] /A << /S /GoTo /D <6571756174696F6E2E31342E312E312E3133> >> >>}
\MusicPDFWriteObject{4960}{<< /Type /Annot /Subtype /Link /Border [0 0 0] /H /I /C [0 1 0] /Rect [191.693 453.176 223.504 462.044] /A << /S /GoTo /D <636974652E5765696C2D66> >> >>}
\MusicPDFWriteObject{4961}{<< /Type /Annot /Subtype /Link /Border [0 0 0] /H /I /C [1 0 0] /Rect [499.154 411.234 539.327 423.186] /A << /S /GoTo /D <6571756174696F6E2E31342E312E312E3131> >> >>}
\MusicPDFWriteObject{4962}{<< /Type /Annot /Subtype /Link /Border [0 0 0] /H /I /C [1 0 0] /Rect [73.712 398.083 113.885 410.036] /A << /S /GoTo /D <6571756174696F6E2E31342E312E312E3132> >> >>}
\MusicPDFWriteObject{4963}{<< /Type /Annot /Subtype /Link /Border [0 0 0] /H /I /C [1 0 0] /Rect [144.557 398.083 184.73 410.036] /A << /S /GoTo /D <6571756174696F6E2E31342E312E312E3133> >> >>}
\MusicPDFWriteObject{4964}{<< /Type /Annot /Subtype /Link /Border [0 0 0] /H /I /C [1 0 0] /Rect [376.976 398.083 417.149 410.036] /A << /S /GoTo /D <6571756174696F6E2E31342E312E312E3130> >> >>}
\MusicPDFWriteObject{4965}{<< /Type /Annot /Subtype /Link /Border [0 0 0] /H /I /C [0 1 0] /Rect [290.648 387.422 320.017 396.29] /A << /S /GoTo /D <636974652E4172742D32> >> >>}
\MusicPDFWriteObject{4966}{<< /Type /Annot /Subtype /Link /Border [0 0 0] /H /I /C [0 1 0] /Rect [360.377 385.486 394.505 396.29] /A << /S /GoTo /D <636974652E4D75727479> >> >>}
\MusicPDFWriteObject{4967}{<< /Type /Annot /Subtype /Link /Border [0 0 0] /H /I /C [0 1 0] /Rect [235.663 361.121 259.942 369.989] /A << /S /GoTo /D <636974652E54617465> >> >>}
\MusicPDFWriteObject{4968}{<< /Type /Annot /Subtype /Link /Border [0 0 0] /H /I /C [0 1 0] /Rect [518.88 361.121 543.159 369.989] /A << /S /GoTo /D <636974652E54617465> >> >>}
\MusicPDFWriteObject{4969}{<< /Type /Annot /Subtype /Link /Border [0 0 0] /H /I /C [0 1 0] /Rect [298.969 334.819 323.248 343.687] /A << /S /GoTo /D <636974652E54617465> >> >>}
\MusicPDFWriteObject{4970}{<< /Type /Annot /Subtype /Link /Border [0 0 0] /H /I /C [0 1 0] /Rect [195.455 321.668 214.664 330.536] /A << /S /GoTo /D <636974652E44656C> >> >>}
\MusicPDFWriteObject{4971}{<< /Type /Annot /Subtype /Link /Border [0 0 0] /H /I /C [0 1 0] /Rect [296.298 295.367 328.109 304.235] /A << /S /GoTo /D <636974652E5765696C2D66> >> >>}
\MusicPDFWriteObject{4972}{<< /Type /Annot /Subtype /Link /Border [0 0 0] /H /I /C [1 0 0] /Rect [529.779 267.889 535.741 281.809] /A << /S /GoTo /D <48666F6F746E6F74652E3533> >> >>}
\MusicPDFWriteObject{4973}{<< /Type /Annot /Subtype /Link /Border [0 0 0] /H /I /C [1 0 0] /Rect [222.517 254.186 257.877 266.138] /A << /S /GoTo /D <74686D2E31332E332E312E31> >> >>}
\MusicPDFWriteObject{4974}{<< /Type /Annot /Subtype /Link /Border [0 0 0] /H /I /C [1 0 0] /Rect [388.586 241.034 423.214 252.99] /A << /S /GoTo /D <74686D2E31332E332E312E32> >> >>}
\MusicPDFWriteObject{4975}{<< /Type /Annot /Subtype /Link /Border [0 0 0] /H /I /C [0 1 0] /Rect [126.919 256.046 165.979 264.914] /A << /S /GoTo /D <636974652E5275642D5361> >> >>}
\MusicPDFWriteObject{4976}{<< /Type /Annot /Subtype /Link /Border [0 0 0] /H /I /C [0 1 0] /Rect [192.071 256.046 224.768 264.914] /A << /S /GoTo /D <636974652E52752D5361> >> >>}
\MusicPDFWriteObject{4977}{<< /Type /Annot /Subtype /Link /Border [0 0 0] /H /I /C [0 1 0] /Rect [170.891 203.443 203.587 212.311] /A << /S /GoTo /D <636974652E52752D5361> >> >>}
\MusicPDFWriteObject{4978}{<< /Type /Annot /Subtype /Link /Border [0 0 0] /H /I /C [0 1 0] /Rect [424.539 163.99 449.454 172.858] /A << /S /GoTo /D <636974652E41726D> >> >>}
\MusicPDFWriteObject{4979}{<< /Type /Annot /Subtype /Link /Border [0 0 0] /H /I /C [0 1 0] /Rect [92.581 83.149 109.265 93.954] /A << /S /GoTo /D <636974652E4E67> >> >>}
\MusicPDFWriteObject{4980}{<< /Type /Annot /Subtype /Link /Border [0 0 0] /H /I /C [1 0 0] /Rect [112.727 672.036 147.92 682.827] /A << /S /GoTo /D <74686D2E31332E332E312E32> >> >>}
\MusicPDFWriteObject{4981}{<< /Type /Annot /Subtype /Link /Border [0 0 0] /H /I /C [1 0 0] /Rect [363.143 672.036 398.336 682.827] /A << /S /GoTo /D <74686D2E31332E332E312E32> >> >>}
\MusicPDFWriteObject{4982}{<< /Type /Annot /Subtype /Link /Border [0 0 0] /H /I /C [0 1 0] /Rect [116.433 621.369 142.195 630.237] /A << /S /GoTo /D <636974652E476F7373> >> >>}
\MusicPDFWriteObject{4983}{<< /Type /Annot /Subtype /Link /Border [0 0 0] /H /I /C [0 1 0] /Rect [396.818 608.219 427.245 617.087] /A << /S /GoTo /D <636974652E4F6D6172> >> >>}
\MusicPDFWriteObject{4984}{<< /Type /Annot /Subtype /Link /Border [0 0 0] /H /I /C [1 0 0] /Rect [491.206 566.83 525.735 577.62] /A << /S /GoTo /D <74686D2E31332E332E312E32> >> >>}
\MusicPDFWriteObject{4985}{<< /Type /Annot /Subtype /Link /Border [0 0 0] /H /I /C [1 0 0] /Rect [358.964 540.528 394.158 551.319] /A << /S /GoTo /D <74686D2E31342E312E322E31> >> >>}
\MusicPDFWriteObject{4986}{<< /Type /Annot /Subtype /Link /Border [0 0 0] /H /I /C [0 1 0] /Rect [462.464 354.083 514.999 364.888] /A << /S /GoTo /D <636974652E4C616E676C616E6473> >> >>}
\MusicPDFWriteObject{4987}{<< /Type /Annot /Subtype /Link /Border [0 0 0] /H /I /C [0 1 0] /Rect [186.537 340.933 239.072 351.737] /A << /S /GoTo /D <636974652E4C616E676C616E6473> >> >>}
\MusicPDFWriteObject{4988}{<< /Type /Annot /Subtype /Link /Border [0 0 0] /H /I /C [1 0 0] /Rect [160.968 213.009 195.497 223.8] /A << /S /GoTo /D <74686D2E31342E312E322E31> >> >>}
\MusicPDFWriteObject{4989}{<< /Type /Annot /Subtype /Link /Border [0 0 0] /H /I /C [1 0 0] /Rect [262.562 213.009 297.091 223.8] /A << /S /GoTo /D <74686D2E31342E312E322E32> >> >>}
\MusicPDFWriteObject{4990}{<< /Type /Annot /Subtype /Link /Border [0 0 0] /H /I /C [1 0 0] /Rect [304.943 186.708 340.8 197.499] /A << /S /GoTo /D <74686D2E31332E332E312E32> >> >>}
\MusicPDFWriteObject{4991}{<< /Type /Annot /Subtype /Link /Border [0 0 0] /H /I /C [1 0 0] /Rect [463.496 160.406 498.689 171.197] /A << /S /GoTo /D <74686D2E31332E332E312E32> >> >>}
\MusicPDFWriteObject{4992}{<< /Type /Annot /Subtype /Link /Border [0 0 0] /H /I /C [1 0 0] /Rect [183.289 147.256 217.818 158.046] /A << /S /GoTo /D <74686D2E31342E312E322E31> >> >>}
\MusicPDFWriteObject{4993}{<< /Type /Annot /Subtype /Link /Border [0 0 0] /H /I /C [1 0 0] /Rect [287.495 147.256 322.024 158.046] /A << /S /GoTo /D <74686D2E31342E312E322E32> >> >>}
\MusicPDFWriteObject{4994}{<< /Type /Annot /Subtype /Link /Border [0 0 0] /H /I /C [1 0 0] /Rect [124.073 134.105 158.835 144.895] /A << /S /GoTo /D <74686D2E31332E332E312E32> >> >>}
\MusicPDFWriteObject{4995}{<< /Type /Annot /Subtype /Link /Border [0 0 0] /H /I /C [1 0 0] /Rect [201.756 658.005 239.058 669.96] /A << /S /GoTo /D <74686D2E31342E312E322E31> >> >>}
\MusicPDFWriteObject{4996}{<< /Type /Annot /Subtype /Link /Border [0 0 0] /H /I /C [1 0 0] /Rect [314.553 658.005 351.855 669.96] /A << /S /GoTo /D <74686D2E31342E312E322E32> >> >>}
\MusicPDFWriteObject{4997}{<< /Type /Annot /Subtype /Link /Border [0 0 0] /H /I /C [1 0 0] /Rect [226.345 583.846 261.339 594.637] /A << /S /GoTo /D <74686D2E31342E312E322E33> >> >>}
\MusicPDFWriteObject{4998}{<< /Type /Annot /Subtype /Link /Border [0 0 0] /H /I /C [1 0 0] /Rect [509.851 583.846 544.845 594.637] /A << /S /GoTo /D <74686D2E31332E332E312E32> >> >>}
\MusicPDFWriteObject{4999}{<< /Type /Annot /Subtype /Link /Border [0 0 0] /H /I /C [0 1 0] /Rect [434.505 500.14 466.316 509.008] /A << /S /GoTo /D <636974652E5765696C2D66> >> >>}
\MusicPDFWriteObject{5000}{<< /Type /Annot /Subtype /Link /Border [0 0 0] /H /I /C [1 0 0] /Rect [149.158 445.047 189.331 457] /A << /S /GoTo /D <6571756174696F6E2E31342E312E312E3133> >> >>}
\MusicPDFWriteObject{5001}{<< /Type /Annot /Subtype /Link /Border [0 0 0] /H /I /C [1 0 0] /Rect [140.255 405.688 175.448 417.641] /A << /S /GoTo /D <74686D2E31342E312E312E31> >> >>}
\MusicPDFWriteObject{5002}{<< /Type /Annot /Subtype /Link /Border [0 0 0] /H /I /C [1 0 0] /Rect [337.605 405.688 377.779 417.641] /A << /S /GoTo /D <6571756174696F6E2E31342E312E312E3132> >> >>}
\MusicPDFWriteObject{5003}{<< /Type /Annot /Subtype /Link /Border [0 0 0] /H /I /C [1 0 0] /Rect [275.523 341.894 305.737 353.847] /A << /S /GoTo /D <74686D2E372E312E322E34> >> >>}
\MusicPDFWriteObject{5004}{<< /Type /Annot /Subtype /Link /Border [0 0 0] /H /I /C [1 0 0] /Rect [157.225 329.296 193.083 340.144] /A << /S /GoTo /D <74686D2E31342E312E322E32> >> >>}
\MusicPDFWriteObject{5005}{<< /Type /Annot /Subtype /Link /Border [0 0 0] /H /I /C [1 0 0] /Rect [157.25 289.29 193.108 301.245] /A << /S /GoTo /D <74686D2E31342E312E322E31> >> >>}
\MusicPDFWriteObject{5006}{<< /Type /Annot /Subtype /Link /Border [0 0 0] /H /I /C [1 0 0] /Rect [221.479 211.013 256.672 222.966] /A << /S /GoTo /D <6571756174696F6E2E31342E312E312E31> >> >>}
\MusicPDFWriteObject{5007}{<< /Type /Annot /Subtype /Link /Border [0 0 0] /H /I /C [1 0 0] /Rect [261.489 568.848 296.683 580.801] /A << /S /GoTo /D <6571756174696F6E2E31342E312E322E34> >> >>}
\MusicPDFWriteObject{5008}{<< /Type /Annot /Subtype /Link /Border [0 0 0] /H /I /C [1 0 0] /Rect [246.747 556.251 281.276 567.041] /A << /S /GoTo /D <74686D2E31342E312E322E35> >> >>}
\MusicPDFWriteObject{5009}{<< /Type /Annot /Subtype /Link /Border [0 0 0] /H /I /C [1 0 0] /Rect [144.626 542.547 179.819 554.499] /A << /S /GoTo /D <6571756174696F6E2E31342E312E322E34> >> >>}
\MusicPDFWriteObject{5010}{<< /Type /Annot /Subtype /Link /Border [0 0 0] /H /I /C [0 1 0] /Rect [412.764 543.1 439.692 553.904] /A << /S /GoTo /D <636974652E51756572> >> >>}
\MusicPDFWriteObject{5011}{<< /Type /Annot /Subtype /Link /Border [0 0 0] /H /I /C [0 1 0] /Rect [107.295 473.036 134.223 483.841] /A << /S /GoTo /D <636974652E51756572> >> >>}
\MusicPDFWriteObject{5012}{<< /Type /Annot /Subtype /Link /Border [0 0 0] /H /I /C [1 0 0] /Rect [459.853 472.483 495.046 484.436] /A << /S /GoTo /D <6571756174696F6E2E31342E312E322E34> >> >>}
\MusicPDFWriteObject{5013}{<< /Type /Annot /Subtype /Link /Border [0 0 0] /H /I /C [0 1 0] /Rect [267.158 446.735 294.086 457.539] /A << /S /GoTo /D <636974652E51756572> >> >>}
\MusicPDFWriteObject{5014}{<< /Type /Annot /Subtype /Link /Border [0 0 0] /H /I /C [0 1 0] /Rect [491.268 422.37 543.159 431.238] /A << /S /GoTo /D <636974652E53696C7665726D616E> >> >>}
\MusicPDFWriteObject{5015}{<< /Type /Annot /Subtype /Link /Border [0 0 0] /H /I /C [0 1 0] /Rect [350.925 238.597 365.728 247.465] /A << /S /GoTo /D <636974652E4A4C> >> >>}
\MusicPDFWriteObject{5016}{<< /Type /Annot /Subtype /Link /Border [0 0 0] /H /I /C [0 1 0] /Rect [147.219 225.446 162.022 234.314] /A << /S /GoTo /D <636974652E4A4C> >> >>}
\MusicPDFWriteObject{5017}{<< /Type /Annot /Subtype /Link /Border [0 0 0] /H /I /C [1 0 0] /Rect [358.009 168.835 393.202 180.788] /A << /S /GoTo /D <6571756174696F6E2E31342E312E322E34> >> >>}
\MusicPDFWriteObject{5018}{<< /Type /Annot /Subtype /Link /Border [0 0 0] /H /I /C [1 0 0] /Rect [317.893 156.237 352.422 167.028] /A << /S /GoTo /D <74686D2E31342E312E322E35> >> >>}
\MusicPDFWriteObject{5019}{<< /Type /Annot /Subtype /Link /Border [0 0 0] /H /I /C [0 1 0] /Rect [361.959 145.024 387.72 153.891] /A << /S /GoTo /D <636974652E476F7373> >> >>}
\MusicPDFWriteObject{5020}{<< /Type /Annot /Subtype /Link /Border [0 0 0] /H /I /C [1 0 0] /Rect [369.013 96.3 403.841 107.091] /A << /S /GoTo /D <74686D2E31332E332E312E32> >> >>}
\MusicPDFWriteObject{5021}{<< /Type /Annot /Subtype /Link /Border [0 0 0] /H /I /C [1 0 0] /Rect [464.256 684.633 498.785 696.588] /A << /S /GoTo /D <74686D2E31332E332E312E31> >> >>}
\MusicPDFWriteObject{5022}{<< /Type /Annot /Subtype /Link /Border [0 0 0] /H /I /C [1 0 0] /Rect [125.31 672.036 160.536 682.84] /A << /S /GoTo /D <74686D2E31332E332E312E32> >> >>}
\MusicPDFWriteObject{5023}{<< /Type /Annot /Subtype /Link /Border [0 0 0] /H /I /C [1 0 0] /Rect [343.827 569.506 374.04 581.461] /A << /S /GoTo /D <74686D2E372E312E332E32> >> >>}
\MusicPDFWriteObject{5024}{<< /Type /Annot /Subtype /Link /Border [0 0 0] /H /I /C [0 1 0] /Rect [155.93 545.695 190.058 554.563] /A << /S /GoTo /D <636974652E4D6F726F7A> >> >>}
\MusicPDFWriteObject{5025}{<< /Type /Annot /Subtype /Link /Border [0 0 0] /H /I /C [0 1 0] /Rect [324.382 545.695 360.385 554.563] /A << /S /GoTo /D <636974652E4B6C654D6F> >> >>}
\MusicPDFWriteObject{5026}{<< /Type /Annot /Subtype /Link /Border [0 0 0] /H /I /C [1 0 0] /Rect [384.908 351.462 420.101 363.417] /A << /S /GoTo /D <74686D2E31332E332E312E31> >> >>}
\MusicPDFWriteObject{5027}{<< /Type /Annot /Subtype /Link /Border [0 0 0] /H /I /C [1 0 0] /Rect [257.52 338.312 293.377 350.264] /A << /S /GoTo /D <74686D2E31342E312E322E33> >> >>}
\MusicPDFWriteObject{5028}{<< /Type /Annot /Subtype /Link /Border [0 0 0] /H /I /C [1 0 0] /Rect [434.197 338.312 470.054 350.264] /A << /S /GoTo /D <74686D2E31342E312E312E31> >> >>}
\MusicPDFWriteObject{5029}{<< /Type /Annot /Subtype /Link /Border [0 0 0] /H /I /C [0 1 0] /Rect [444.959 270.786 477.655 279.654] /A << /S /GoTo /D <636974652E52752D5361> >> >>}
\MusicPDFWriteObject{5030}{<< /Type /Annot /Subtype /Link /Border [0 0 0] /H /I /C [1 0 0] /Rect [352.143 187.566 386.872 198.357] /A << /S /GoTo /D <74686D2E31332E332E312E31> >> >>}
\MusicPDFWriteObject{5031}{<< /Type /Annot /Subtype /Link /Border [0 0 0] /H /I /C [1 0 0] /Rect [482.21 666.058 518.067 676.849] /A << /S /GoTo /D <74686D2E31332E332E312E31> >> >>}
\MusicPDFWriteObject{5032}{<< /Type /Annot /Subtype /Link /Border [0 0 0] /H /I /C [0 1 0] /Rect [73.726 536.487 91.199 545.355] /A << /S /GoTo /D <636974652E4346> >> >>}
\MusicPDFWriteObject{5033}{<< /Type /Annot /Subtype /Link /Border [0 0 0] /H /I /C [0 1 0] /Rect [293.205 536.487 311.175 545.355] /A << /S /GoTo /D <636974652E5256> >> >>}
\MusicPDFWriteObject{5034}{<< /Type /Annot /Subtype /Link /Border [0 0 0] /H /I /C [1 0 0] /Rect [508.165 195.391 543.225 206.181] /A << /S /GoTo /D <74686D2E31312E322E322E31> >> >>}
\MusicPDFWriteObject{5035}{<< /Type /Annot /Subtype /Link /Border [0 0 0] /H /I /C [1 0 0] /Rect [357.188 142.234 387.966 154.187] /A << /S /GoTo /D <74686D2E372E312E332E32> >> >>}
\MusicPDFWriteObject{5036}{<< /Type /Annot /Subtype /Link /Border [0 0 0] /H /I /C [0 1 0] /Rect [151.861 647.343 174.045 656.211] /A << /S /GoTo /D <636974652E547356> >> >>}
\MusicPDFWriteObject{5037}{<< /Type /Annot /Subtype /Link /Border [0 0 0] /H /I /C [1 0 0] /Rect [324.441 605.954 358.971 616.745] /A << /S /GoTo /D <74686D2E31332E332E312E32> >> >>}
\MusicPDFWriteObject{5038}{<< /Type /Annot /Subtype /Link /Border [0 0 0] /H /I /C [1 0 0] /Rect [252.531 539.647 282.744 551.599] /A << /S /GoTo /D <6571756174696F6E2E352E332E312E35> >> >>}
\MusicPDFWriteObject{5039}{<< /Type /Annot /Subtype /Link /Border [0 0 0] /H /I /C [0 1 0] /Rect [233.139 436.931 251.255 445.799] /A << /S /GoTo /D <636974652E4D65> >> >>}
\MusicPDFWriteObject{5040}{<< /Type /Annot /Subtype /Link /Border [0 0 0] /H /I /C [0 1 0] /Rect [296.631 358.026 314.746 366.894] /A << /S /GoTo /D <636974652E4D65> >> >>}
\MusicPDFWriteObject{5041}{<< /Type /Annot /Subtype /Link /Border [0 0 0] /H /I /C [0 1 0] /Rect [383.575 344.875 401.691 353.743] /A << /S /GoTo /D <636974652E4D65> >> >>}
\MusicPDFWriteObject{5042}{<< /Type /Annot /Subtype /Link /Border [0 0 0] /H /I /C [1 0 0] /Rect [197.609 303.486 232.802 314.277] /A << /S /GoTo /D <74686D2E31312E322E322E31> >> >>}
\MusicPDFWriteObject{5043}{<< /Type /Annot /Subtype /Link /Border [0 0 0] /H /I /C [0 1 0] /Rect [390.651 120.404 416.507 129.272] /A << /S /GoTo /D <636974652E464754> >> >>}
\MusicPDFWriteObject{5044}{<< /Type /Annot /Subtype /Link /Border [0 0 0] /H /I /C [0 1 0] /Rect [419.747 687.123 445.508 695.991] /A << /S /GoTo /D <636974652E476F7373> >> >>}
\MusicPDFWriteObject{5045}{<< /Type /Annot /Subtype /Link /Border [0 0 0] /H /I /C [1 0 0] /Rect [88.936 472.939 94.899 483.306] /A << /S /GoTo /D <48666F6F746E6F74652E3534> >> >>}
\MusicPDFWriteObject{5046}{<< /Type /Annot /Subtype /Link /Border [0 0 0] /H /I /C [0 1 0] /Rect [494.754 473.976 512.227 482.844] /A << /S /GoTo /D <636974652E4346> >> >>}
\MusicPDFWriteObject{5047}{<< /Type /Annot /Subtype /Link /Border [0 0 0] /H /I /C [1 0 0] /Rect [82.959 397.854 431.249 410.169] /A << /S /GoTo /D <73656374696F6E2E412E31> >> >>}
\MusicPDFWriteObject{5048}{<< /Type /Annot /Subtype /Link /Border [0 0 0] /H /I /C [1 0 0] /Rect [82.959 385.617 277.445 396.407] /A << /S /GoTo /D <73756273656374696F6E2E412E312E31> >> >>}
\MusicPDFWriteObject{5049}{<< /Type /Annot /Subtype /Link /Border [0 0 0] /H /I /C [1 0 0] /Rect [82.959 372.466 287.31 383.257] /A << /S /GoTo /D <73756273756273656374696F6E2E412E312E312E31> >> >>}
\MusicPDFWriteObject{5050}{<< /Type /Annot /Subtype /Link /Border [0 0 0] /H /I /C [1 0 0] /Rect [82.959 359.315 306.704 370.106] /A << /S /GoTo /D <73756273656374696F6E2E412E312E32> >> >>}
\MusicPDFWriteObject{5051}{<< /Type /Annot /Subtype /Link /Border [0 0 0] /H /I /C [1 0 0] /Rect [82.959 346.607 238.96 356.955] /A << /S /GoTo /D <73756273756273656374696F6E2E412E312E322E31> >> >>}
\MusicPDFWriteObject{5052}{<< /Type /Annot /Subtype /Link /Border [0 0 0] /H /I /C [1 0 0] /Rect [82.959 332.46 234.152 344.415] /A << /S /GoTo /D <73756273656374696F6E2E412E312E33> >> >>}
\MusicPDFWriteObject{5053}{<< /Type /Annot /Subtype /Link /Border [0 0 0] /H /I /C [1 0 0] /Rect [82.959 319.309 273.966 331.264] /A << /S /GoTo /D <73756273756273656374696F6E2E412E312E332E31> >> >>}
\MusicPDFWriteObject{5054}{<< /Type /Annot /Subtype /Link /Border [0 0 0] /H /I /C [1 0 0] /Rect [82.959 306.158 319.247 318.574] /A << /S /GoTo /D <73756273756273656374696F6E2E412E312E332E32> >> >>}
\MusicPDFWriteObject{5055}{<< /Type /Annot /Subtype /Link /Border [0 0 0] /H /I /C [1 0 0] /Rect [82.959 293.561 229.699 304.352] /A << /S /GoTo /D <73756273656374696F6E2E412E312E34> >> >>}
\MusicPDFWriteObject{5056}{<< /Type /Annot /Subtype /Link /Border [0 0 0] /H /I /C [1 0 0] /Rect [82.959 280.411 312.348 291.201] /A << /S /GoTo /D <73756273756273656374696F6E2E412E312E342E31> >> >>}
\MusicPDFWriteObject{5057}{<< /Type /Annot /Subtype /Link /Border [0 0 0] /H /I /C [1 0 0] /Rect [82.959 267.26 295.364 278.051] /A << /S /GoTo /D <73756273656374696F6E2E412E312E35> >> >>}
\MusicPDFWriteObject{5058}{<< /Type /Annot /Subtype /Link /Border [0 0 0] /H /I /C [1 0 0] /Rect [82.959 256.046 223.761 264.9] /A << /S /GoTo /D <73756273756273656374696F6E2E412E312E352E31> >> >>}
\MusicPDFWriteObject{5059}{<< /Type /Annot /Subtype /Link /Border [0 0 0] /H /I /C [1 0 0] /Rect [82.959 240.045 380.605 252.36] /A << /S /GoTo /D <73756273756273656374696F6E2E412E312E352E32> >> >>}
\MusicPDFWriteObject{5060}{<< /Type /Annot /Subtype /Link /Border [0 0 0] /H /I /C [1 0 0] /Rect [82.959 226.894 436.055 239.209] /A << /S /GoTo /D <73756273756273656374696F6E2E412E312E352E33> >> >>}
\MusicPDFWriteObject{5061}{<< /Type /Annot /Subtype /Link /Border [0 0 0] /H /I /C [1 0 0] /Rect [82.959 216.594 280.863 225.448] /A << /S /GoTo /D <73656374696F6E2E412E32> >> >>}
\MusicPDFWriteObject{5062}{<< /Type /Annot /Subtype /Link /Border [0 0 0] /H /I /C [1 0 0] /Rect [82.959 201.948 403.458 212.297] /A << /S /GoTo /D <73756273656374696F6E2E412E322E31> >> >>}
\MusicPDFWriteObject{5063}{<< /Type /Annot /Subtype /Link /Border [0 0 0] /H /I /C [1 0 0] /Rect [82.959 190.292 280.602 199.146] /A << /S /GoTo /D <73756273756273656374696F6E2E412E322E312E31> >> >>}
\MusicPDFWriteObject{5064}{<< /Type /Annot /Subtype /Link /Border [0 0 0] /H /I /C [1 0 0] /Rect [82.959 175.205 351.055 185.995] /A << /S /GoTo /D <73756273656374696F6E2E412E322E32> >> >>}
\MusicPDFWriteObject{5065}{<< /Type /Annot /Subtype /Link /Border [0 0 0] /H /I /C [1 0 0] /Rect [82.959 163.99 225.28 172.844] /A << /S /GoTo /D <73756273756273656374696F6E2E412E322E322E31> >> >>}
\MusicPDFWriteObject{5066}{<< /Type /Annot /Subtype /Link /Border [0 0 0] /H /I /C [1 0 0] /Rect [82.959 147.381 391.898 160.765] /A << /S /GoTo /D <73756273756273656374696F6E2E412E322E322E32> >> >>}
\MusicPDFWriteObject{5067}{<< /Type /Annot /Subtype /Link /Border [0 0 0] /H /I /C [1 0 0] /Rect [82.959 134.839 300.15 147.153] /A << /S /GoTo /D <73756273756273656374696F6E2E412E322E322E33> >> >>}
\MusicPDFWriteObject{5068}{<< /Type /Annot /Subtype /Link /Border [0 0 0] /H /I /C [1 0 0] /Rect [82.959 122.048 465.538 134.64] /A << /S /GoTo /D <73756273656374696F6E2E412E322E33> >> >>}
\MusicPDFWriteObject{5069}{<< /Type /Annot /Subtype /Link /Border [0 0 0] /H /I /C [1 0 0] /Rect [82.959 109.451 264.069 120.241] /A << /S /GoTo /D <73756273756273656374696F6E2E412E322E332E31> >> >>}
\MusicPDFWriteObject{5070}{<< /Type /Annot /Subtype /Link /Border [0 0 0] /H /I /C [1 0 0] /Rect [82.959 95.493 284.576 108.339] /A << /S /GoTo /D <73756273756273656374696F6E2E412E322E332E32> >> >>}
\MusicPDFWriteObject{5071}{<< /Type /Annot /Subtype /Link /Border [0 0 0] /H /I /C [1 0 0] /Rect [82.959 83.149 250.162 93.94] /A << /S /GoTo /D <73756273756273656374696F6E2E412E322E332E33> >> >>}
\MusicPDFWriteObject{5072}{<< /Type /Annot /Subtype /Link /Border [0 0 0] /H /I /C [1 0 0] /Rect [82.959 684.633 267.201 696.588] /A << /S /GoTo /D <73756273756273656374696F6E2E412E322E332E34> >> >>}
\MusicPDFWriteObject{5073}{<< /Type /Annot /Subtype /Link /Border [0 0 0] /H /I /C [1 0 0] /Rect [82.959 671.228 252.316 684.075] /A << /S /GoTo /D <73756273756273656374696F6E2E412E322E332E35> >> >>}
\MusicPDFWriteObject{5074}{<< /Type /Annot /Subtype /Link /Border [0 0 0] /H /I /C [1 0 0] /Rect [82.959 658.331 379.539 670.747] /A << /S /GoTo /D <73756273756273656374696F6E2E412E322E332E36> >> >>}
\MusicPDFWriteObject{5075}{<< /Type /Annot /Subtype /Link /Border [0 0 0] /H /I /C [1 0 0] /Rect [82.959 645.734 303.91 656.525] /A << /S /GoTo /D <73756273756273656374696F6E2E412E322E332E37> >> >>}
\MusicPDFWriteObject{5076}{<< /Type /Annot /Subtype /Link /Border [0 0 0] /H /I /C [1 0 0] /Rect [82.959 631.67 322.192 643.985] /A << /S /GoTo /D <73756273656374696F6E2E412E322E34> >> >>}
\MusicPDFWriteObject{5077}{<< /Type /Annot /Subtype /Link /Border [0 0 0] /H /I /C [1 0 0] /Rect [82.959 618.519 386.88 630.358] /A << /S /GoTo /D <73756273756273656374696F6E2E412E322E342E31> >> >>}
\MusicPDFWriteObject{5078}{<< /Type /Annot /Subtype /Link /Border [0 0 0] /H /I /C [1 0 0] /Rect [82.959 603.935 442.236 618.79] /A << /S /GoTo /D <73756273756273656374696F6E2E412E322E342E32> >> >>}
\MusicPDFWriteObject{5079}{<< /Type /Annot /Subtype /Link /Border [0 0 0] /H /I /C [1 0 0] /Rect [82.959 595.068 176.973 603.922] /A << /S /GoTo /D <73656374696F6E2E412E33> >> >>}
\MusicPDFWriteObject{5080}{<< /Type /Annot /Subtype /Link /Border [0 0 0] /H /I /C [1 0 0] /Rect [82.959 579.981 238.102 590.771] /A << /S /GoTo /D <73756273656374696F6E2E412E332E31> >> >>}
\MusicPDFWriteObject{5081}{<< /Type /Annot /Subtype /Link /Border [0 0 0] /H /I /C [1 0 0] /Rect [82.959 566.83 345.992 577.62] /A << /S /GoTo /D <73756273756273656374696F6E2E412E332E312E31> >> >>}
\MusicPDFWriteObject{5082}{<< /Type /Annot /Subtype /Link /Border [0 0 0] /H /I /C [1 0 0] /Rect [82.959 553.125 264.365 565.08] /A << /S /GoTo /D <73756273656374696F6E2E412E332E32> >> >>}
\MusicPDFWriteObject{5083}{<< /Type /Annot /Subtype /Link /Border [0 0 0] /H /I /C [1 0 0] /Rect [82.959 539.721 252.048 552.567] /A << /S /GoTo /D <73756273656374696F6E2E412E332E33> >> >>}
\MusicPDFWriteObject{5084}{<< /Type /Annot /Subtype /Link /Border [0 0 0] /H /I /C [1 0 0] /Rect [82.959 527.377 456.163 538.168] /A << /S /GoTo /D <73756273756273656374696F6E2E412E332E332E31> >> >>}
\MusicPDFWriteObject{5085}{<< /Type /Annot /Subtype /Link /Border [0 0 0] /H /I /C [1 0 0] /Rect [82.959 514.227 243.113 525.017] /A << /S /GoTo /D <73756273656374696F6E2E412E332E34> >> >>}
\MusicPDFWriteObject{5086}{<< /Type /Annot /Subtype /Link /Border [0 0 0] /H /I /C [1 0 0] /Rect [82.959 501.076 389.141 511.867] /A << /S /GoTo /D <73656374696F6E2E412E34> >> >>}
\MusicPDFWriteObject{5087}{<< /Type /Annot /Subtype /Link /Border [0 0 0] /H /I /C [1 0 0] /Rect [82.959 487.925 360.118 498.716] /A << /S /GoTo /D <73756273656374696F6E2E412E342E31> >> >>}
\MusicPDFWriteObject{5088}{<< /Type /Annot /Subtype /Link /Border [0 0 0] /H /I /C [1 0 0] /Rect [82.959 474.774 429.409 485.565] /A << /S /GoTo /D <73656374696F6E2E412E35> >> >>}
\MusicPDFWriteObject{5089}{<< /Type /Annot /Subtype /Link /Border [0 0 0] /H /I /C [1 0 0] /Rect [82.959 461.624 350.778 472.414] /A << /S /GoTo /D <73756273656374696F6E2E412E352E31> >> >>}
\MusicPDFWriteObject{5090}{<< /Type /Annot /Subtype /Link /Border [0 0 0] /H /I /C [1 0 0] /Rect [82.959 448.473 355.234 459.264] /A << /S /GoTo /D <73756273756273656374696F6E2E412E352E312E31> >> >>}
\MusicPDFWriteObject{5091}{<< /Type /Annot /Subtype /Link /Border [0 0 0] /H /I /C [1 0 0] /Rect [82.959 435.322 342.424 446.113] /A << /S /GoTo /D <73756273756273656374696F6E2E412E352E312E32> >> >>}
\MusicPDFWriteObject{5092}{<< /Type /Annot /Subtype /Link /Border [0 0 0] /H /I /C [1 0 0] /Rect [82.959 422.171 497.123 433.59] /A << /S /GoTo /D <73756273656374696F6E2E412E352E32> >> >>}
\MusicPDFWriteObject{5093}{<< /Type /Annot /Subtype /Link /Border [0 0 0] /H /I /C [1 0 0] /Rect [82.959 409.021 459.325 420.44] /A << /S /GoTo /D <73756273756273656374696F6E2E412E352E322E31> >> >>}
\MusicPDFWriteObject{5094}{<< /Type /Annot /Subtype /Link /Border [0 0 0] /H /I /C [1 0 0] /Rect [82.959 395.87 407.954 406.661] /A << /S /GoTo /D <73756273756273656374696F6E2E412E352E322E32> >> >>}
\MusicPDFWriteObject{5095}{<< /Type /Annot /Subtype /Link /Border [0 0 0] /H /I /C [1 0 0] /Rect [82.959 382.719 479.016 393.51] /A << /S /GoTo /D <73756273656374696F6E2E412E352E33> >> >>}
\MusicPDFWriteObject{5096}{<< /Type /Annot /Subtype /Link /Border [0 0 0] /H /I /C [1 0 0] /Rect [82.959 369.568 493.458 380.359] /A << /S /GoTo /D <73756273756273656374696F6E2E412E352E332E31> >> >>}
\MusicPDFWriteObject{5097}{<< /Type /Annot /Subtype /Link /Border [0 0 0] /H /I /C [1 0 0] /Rect [82.959 356.418 424.124 367.208] /A << /S /GoTo /D <73756273656374696F6E2E412E352E34> >> >>}
\MusicPDFWriteObject{5098}{<< /Type /Annot /Subtype /Link /Border [0 0 0] /H /I /C [0 1 0] /Rect [491.775 654.844 509.248 663.712] /A << /S /GoTo /D <636974652E4346> >> >>}
\MusicPDFWriteObject{5099}{<< /Type /Annot /Subtype /Link /Border [0 0 0] /H /I /C [0 1 0] /Rect [422.99 615.392 440.462 624.26] /A << /S /GoTo /D <636974652E4346> >> >>}
\MusicPDFWriteObject{5100}{<< /Type /Annot /Subtype /Link /Border [0 0 0] /H /I /C [0 1 0] /Rect [466.468 615.392 484.438 624.26] /A << /S /GoTo /D <636974652E5256> >> >>}
\MusicPDFWriteObject{5101}{<< /Type /Annot /Subtype /Link /Border [0 0 0] /H /I /C [0 1 0] /Rect [188 186.515 205.97 195.383] /A << /S /GoTo /D <636974652E5256> >> >>}
\MusicPDFWriteObject{5102}{<< /Type /Annot /Subtype /Link /Border [0 0 0] /H /I /C [0 1 0] /Rect [265.182 186.515 283.485 195.383] /A << /S /GoTo /D <636974652E4445> >> >>}
\MusicPDFWriteObject{5103}{<< /Type /Annot /Subtype /Link /Border [0 0 0] /H /I /C [0 1 0] /Rect [324.958 186.515 342.431 195.383] /A << /S /GoTo /D <636974652E4346> >> >>}
\MusicPDFWriteObject{5104}{<< /Type /Annot /Subtype /Link /Border [0 0 0] /H /I /C [0 1 0] /Rect [488.593 186.515 506.065 195.383] /A << /S /GoTo /D <636974652E4346> >> >>}
\MusicPDFWriteObject{5105}{<< /Type /Annot /Subtype /Link /Border [0 0 0] /H /I /C [1 0 0] /Rect [407.881 346.2 441.199 358.153] /A << /S /GoTo /D <74686D2E412E312E312E32> >> >>}
\MusicPDFWriteObject{5106}{<< /Type /Annot /Subtype /Link /Border [0 0 0] /H /I /C [1 0 0] /Rect [204.239 621.666 236.942 633.619] /A << /S /GoTo /D <6571756174696F6E2E412E312E322E31> >> >>}
\MusicPDFWriteObject{5107}{<< /Type /Annot /Subtype /Link /Border [0 0 0] /H /I /C [1 0 0] /Rect [364.912 621.666 397.616 633.619] /A << /S /GoTo /D <6571756174696F6E2E412E312E322E32> >> >>}
\MusicPDFWriteObject{5108}{<< /Type /Annot /Subtype /Link /Border [0 0 0] /H /I /C [0 1 0] /Rect [190.918 431.164 208.391 440.032] /A << /S /GoTo /D <636974652E4346> >> >>}
\MusicPDFWriteObject{5109}{<< /Type /Annot /Subtype /Link /Border [0 0 0] /H /I /C [0 1 0] /Rect [297.37 431.164 315.672 440.032] /A << /S /GoTo /D <636974652E4445> >> >>}
\MusicPDFWriteObject{5110}{<< /Type /Annot /Subtype /Link /Border [0 0 0] /H /I /C [0 1 0] /Rect [195.72 387.743 229.122 396.611] /A << /S /GoTo /D <636974652E5765696C2D63> >> >>}
\MusicPDFWriteObject{5111}{<< /Type /Annot /Subtype /Link /Border [0 0 0] /H /I /C [0 1 0] /Rect [308.487 387.743 341.889 396.611] /A << /S /GoTo /D <636974652E5765696C2D63> >> >>}
\MusicPDFWriteObject{5112}{<< /Type /Annot /Subtype /Link /Border [0 0 0] /H /I /C [0 1 0] /Rect [504.265 374.592 543.27 383.46] /A << /S /GoTo /D <636974652E436F6E726164> >> >>}
\MusicPDFWriteObject{5113}{<< /Type /Annot /Subtype /Link /Border [0 0 0] /H /I /C [1 0 0] /Rect [208.13 220.304 240.833 232.257] /A << /S /GoTo /D <6571756174696F6E2E412E312E332E33> >> >>}
\MusicPDFWriteObject{5114}{<< /Type /Annot /Subtype /Link /Border [0 0 0] /H /I /C [1 0 0] /Rect [198.819 684.633 231.522 696.586] /A << /S /GoTo /D <6571756174696F6E2E412E312E332E32> >> >>}
\MusicPDFWriteObject{5115}{<< /Type /Annot /Subtype /Link /Border [0 0 0] /H /I /C [1 0 0] /Rect [400.385 542.676 433.089 554.629] /A << /S /GoTo /D <6571756174696F6E2E412E312E332E35> >> >>}
\MusicPDFWriteObject{5116}{<< /Type /Annot /Subtype /Link /Border [0 0 0] /H /I /C [1 0 0] /Rect [394.134 418.387 419.09 430.339] /A << /S /GoTo /D <73756273656374696F6E2E412E322E32> >> >>}
\MusicPDFWriteObject{5117}{<< /Type /Annot /Subtype /Link /Border [0 0 0] /H /I /C [1 0 0] /Rect [153.354 283.088 185.443 295.043] /A << /S /GoTo /D <74686D2E412E312E332E34> >> >>}
\MusicPDFWriteObject{5118}{<< /Type /Annot /Subtype /Link /Border [0 0 0] /H /I /C [1 0 0] /Rect [382.199 233.256 414.902 246.103] /A << /S /GoTo /D <74686D2E412E312E332E34> >> >>}
\MusicPDFWriteObject{5119}{<< /Type /Annot /Subtype /Link /Border [0 0 0] /H /I /C [1 0 0] /Rect [282.022 200.472 314.725 212.425] /A << /S /GoTo /D <6571756174696F6E2E412E312E332E35> >> >>}
\MusicPDFWriteObject{5120}{<< /Type /Annot /Subtype /Link /Border [0 0 0] /H /I /C [1 0 0] /Rect [119.144 337.309 152.461 349.013] /A << /S /GoTo /D <74686D2E412E312E322E31> >> >>}
\MusicPDFWriteObject{5121}{<< /Type /Annot /Subtype /Link /Border [0 0 0] /H /I /C [1 0 0] /Rect [244.022 671.482 276.726 683.435] /A << /S /GoTo /D <6571756174696F6E2E412E312E332E32> >> >>}
\MusicPDFWriteObject{5122}{<< /Type /Annot /Subtype /Link /Border [0 0 0] /H /I /C [1 0 0] /Rect [373.429 516.181 406.746 528.494] /A << /S /GoTo /D <74686D2E412E312E352E31> >> >>}
\MusicPDFWriteObject{5123}{<< /Type /Annot /Subtype /Link /Border [0 0 0] /H /I /C [1 0 0] /Rect [273.228 400.53 305.931 411.321] /A << /S /GoTo /D <74686D2E412E312E342E33> >> >>}
\MusicPDFWriteObject{5124}{<< /Type /Annot /Subtype /Link /Border [0 0 0] /H /I /C [1 0 0] /Rect [357.56 131.804 390.264 143.757] /A << /S /GoTo /D <6571756174696F6E2E412E312E352E39> >> >>}
\MusicPDFWriteObject{5125}{<< /Type /Annot /Subtype /Link /Border [0 0 0] /H /I /C [1 0 0] /Rect [176.777 614.091 214.46 626.044] /A << /S /GoTo /D <6571756174696F6E2E412E312E352E3130> >> >>}
\MusicPDFWriteObject{5126}{<< /Type /Annot /Subtype /Link /Border [0 0 0] /H /I /C [0 1 0] /Rect [525.194 518.03 543.164 526.898] /A << /S /GoTo /D <636974652E5256> >> >>}
\MusicPDFWriteObject{5127}{<< /Type /Annot /Subtype /Link /Border [0 0 0] /H /I /C [0 1 0] /Rect [131.233 504.88 148.706 513.747] /A << /S /GoTo /D <636974652E4346> >> >>}
\MusicPDFWriteObject{5128}{<< /Type /Annot /Subtype /Link /Border [0 0 0] /H /I /C [1 0 0] /Rect [70.625 457.513 102.714 468.304] /A << /S /GoTo /D <74686D2E412E312E322E31> >> >>}
\MusicPDFWriteObject{5129}{<< /Type /Annot /Subtype /Link /Border [0 0 0] /H /I /C [1 0 0] /Rect [203.268 194.39 235.971 206.343] /A << /S /GoTo /D <6571756174696F6E2E412E322E312E33> >> >>}
\MusicPDFWriteObject{5130}{<< /Type /Annot /Subtype /Link /Border [0 0 0] /H /I /C [1 0 0] /Rect [227.981 181.24 260.684 193.192] /A << /S /GoTo /D <6571756174696F6E2E412E322E312E32> >> >>}
\MusicPDFWriteObject{5131}{<< /Type /Annot /Subtype /Link /Border [0 0 0] /H /I /C [1 0 0] /Rect [176.666 684.633 209.369 696.586] /A << /S /GoTo /D <6571756174696F6E2E412E322E312E35> >> >>}
\MusicPDFWriteObject{5132}{<< /Type /Annot /Subtype /Link /Border [0 0 0] /H /I /C [1 0 0] /Rect [497.912 647.768 530.615 659.721] /A << /S /GoTo /D <6571756174696F6E2E412E322E312E35> >> >>}
\MusicPDFWriteObject{5133}{<< /Type /Annot /Subtype /Link /Border [0 0 0] /H /I /C [1 0 0] /Rect [491.948 609.901 525.02 620.705] /A << /S /GoTo /D <74686D2E412E312E332E33> >> >>}
\MusicPDFWriteObject{5134}{<< /Type /Annot /Subtype /Link /Border [0 0 0] /H /I /C [1 0 0] /Rect [106.556 598.687 139.259 607.541] /A << /S /GoTo /D <74686D2E412E312E332E34> >> >>}
\MusicPDFWriteObject{5135}{<< /Type /Annot /Subtype /Link /Border [0 0 0] /H /I /C [1 0 0] /Rect [154.21 346.881 186.913 358.834] /A << /S /GoTo /D <6571756174696F6E2E412E322E322E31> >> >>}
\MusicPDFWriteObject{5136}{<< /Type /Annot /Subtype /Link /Border [0 0 0] /H /I /C [1 0 0] /Rect [364.964 647.972 397.667 659.925] /A << /S /GoTo /D <6571756174696F6E2E412E322E322E37> >> >>}
\MusicPDFWriteObject{5137}{<< /Type /Annot /Subtype /Link /Border [0 0 0] /H /I /C [1 0 0] /Rect [376.833 538.073 409.536 550.026] /A << /S /GoTo /D <6571756174696F6E2E412E322E322E38> >> >>}
\MusicPDFWriteObject{5138}{<< /Type /Annot /Subtype /Link /Border [0 0 0] /H /I /C [1 0 0] /Rect [164.043 480.099 196.746 492.051] /A << /S /GoTo /D <6571756174696F6E2E412E322E322E33> >> >>}
\MusicPDFWriteObject{5139}{<< /Type /Annot /Subtype /Link /Border [0 0 0] /H /I /C [1 0 0] /Rect [172.027 466.588 204.73 478.902] /A << /S /GoTo /D <74686D2E412E322E322E31> >> >>}
\MusicPDFWriteObject{5140}{<< /Type /Annot /Subtype /Link /Border [0 0 0] /H /I /C [1 0 0] /Rect [370.8 354.024 408.484 365.977] /A << /S /GoTo /D <6571756174696F6E2E412E322E322E3130> >> >>}
\MusicPDFWriteObject{5141}{<< /Type /Annot /Subtype /Link /Border [0 0 0] /H /I /C [1 0 0] /Rect [151.038 340.729 183.741 352.681] /A << /S /GoTo /D <6571756174696F6E2E412E322E332E33> >> >>}
\MusicPDFWriteObject{5142}{<< /Type /Annot /Subtype /Link /Border [0 0 0] /H /I /C [1 0 0] /Rect [123.702 221.365 156.405 233.318] /A << /S /GoTo /D <6571756174696F6E2E412E322E332E33> >> >>}
\MusicPDFWriteObject{5143}{<< /Type /Annot /Subtype /Link /Border [0 0 0] /H /I /C [1 0 0] /Rect [266.705 220.03 304.389 231.983] /A << /S /GoTo /D <6571756174696F6E2E412E322E332E3134> >> >>}
\MusicPDFWriteObject{5144}{<< /Type /Annot /Subtype /Link /Border [0 0 0] /H /I /C [0 1 0] /Rect [422.638 222.52 440.11 231.388] /A << /S /GoTo /D <636974652E4346> >> >>}
\MusicPDFWriteObject{5145}{<< /Type /Annot /Subtype /Link /Border [0 0 0] /H /I /C [1 0 0] /Rect [216.03 198.256 253.713 210.208] /A << /S /GoTo /D <6571756174696F6E2E412E322E332E3132> >> >>}
\MusicPDFWriteObject{5146}{<< /Type /Annot /Subtype /Link /Border [0 0 0] /H /I /C [1 0 0] /Rect [243.971 119.515 281.654 131.467] /A << /S /GoTo /D <6571756174696F6E2E412E322E332E3134> >> >>}
\MusicPDFWriteObject{5147}{<< /Type /Annot /Subtype /Link /Border [0 0 0] /H /I /C [1 0 0] /Rect [398.187 684.633 435.87 696.586] /A << /S /GoTo /D <6571756174696F6E2E412E322E332E3133> >> >>}
\MusicPDFWriteObject{5148}{<< /Type /Annot /Subtype /Link /Border [0 0 0] /H /I /C [1 0 0] /Rect [307.543 481.652 340.246 493.605] /A << /S /GoTo /D <6571756174696F6E2E412E322E332E31> >> >>}
\MusicPDFWriteObject{5149}{<< /Type /Annot /Subtype /Link /Border [0 0 0] /H /I /C [1 0 0] /Rect [98.422 389.652 131.125 401.605] /A << /S /GoTo /D <6571756174696F6E2E412E322E332E33> >> >>}
\MusicPDFWriteObject{5150}{<< /Type /Annot /Subtype /Link /Border [0 0 0] /H /I /C [1 0 0] /Rect [206.891 376.501 239.594 388.454] /A << /S /GoTo /D <6571756174696F6E2E412E322E332E39> >> >>}
\MusicPDFWriteObject{5151}{<< /Type /Annot /Subtype /Link /Border [0 0 0] /H /I /C [1 0 0] /Rect [161.415 363.35 193.995 375.943] /A << /S /GoTo /D <74686D2E412E322E332E35> >> >>}
\MusicPDFWriteObject{5152}{<< /Type /Annot /Subtype /Link /Border [0 0 0] /H /I /C [1 0 0] /Rect [379.295 263.833 411.999 276.148] /A << /S /GoTo /D <74686D2E412E322E322E35> >> >>}
\MusicPDFWriteObject{5153}{<< /Type /Annot /Subtype /Link /Border [0 0 0] /H /I /C [1 0 0] /Rect [219.83 251.043 257.514 262.996] /A << /S /GoTo /D <6571756174696F6E2E412E322E332E3139> >> >>}
\MusicPDFWriteObject{5154}{<< /Type /Annot /Subtype /Link /Border [0 0 0] /H /I /C [1 0 0] /Rect [117.723 224.741 155.727 237.605] /A << /S /GoTo /D <6571756174696F6E2E412E322E322E3130> >> >>}
\MusicPDFWriteObject{5155}{<< /Type /Annot /Subtype /Link /Border [0 0 0] /H /I /C [1 0 0] /Rect [159.189 162.012 196.872 173.965] /A << /S /GoTo /D <6571756174696F6E2E412E322E332E3133> >> >>}
\MusicPDFWriteObject{5156}{<< /Type /Annot /Subtype /Link /Border [0 0 0] /H /I /C [1 0 0] /Rect [259.759 162.012 292.463 174.605] /A << /S /GoTo /D <74686D2E412E322E332E36> >> >>}
\MusicPDFWriteObject{5157}{<< /Type /Annot /Subtype /Link /Border [0 0 0] /H /I /C [1 0 0] /Rect [138.886 508.132 170.975 521.631] /A << /S /GoTo /D <74686D2E412E322E332E35> >> >>}
\MusicPDFWriteObject{5158}{<< /Type /Annot /Subtype /Link /Border [0 0 0] /H /I /C [1 0 0] /Rect [139.068 481.908 171.156 494.86] /A << /S /GoTo /D <74686D2E412E322E332E35> >> >>}
\MusicPDFWriteObject{5159}{<< /Type /Annot /Subtype /Link /Border [0 0 0] /H /I /C [1 0 0] /Rect [122.807 455.967 160.491 467.919] /A << /S /GoTo /D <6571756174696F6E2E412E322E332E3230> >> >>}
\MusicPDFWriteObject{5160}{<< /Type /Annot /Subtype /Link /Border [0 0 0] /H /I /C [1 0 0] /Rect [189.102 455.967 226.785 467.919] /A << /S /GoTo /D <6571756174696F6E2E412E322E332E3231> >> >>}
\MusicPDFWriteObject{5161}{<< /Type /Annot /Subtype /Link /Border [0 0 0] /H /I /C [1 0 0] /Rect [141.376 416.514 174.694 429.107] /A << /S /GoTo /D <74686D2E412E322E332E32> >> >>}
\MusicPDFWriteObject{5162}{<< /Type /Annot /Subtype /Link /Border [0 0 0] /H /I /C [1 0 0] /Rect [498.108 376.624 530.596 390.124] /A << /S /GoTo /D <74686D2E412E322E332E35> >> >>}
\MusicPDFWriteObject{5163}{<< /Type /Annot /Subtype /Link /Border [0 0 0] /H /I /C [1 0 0] /Rect [382.142 363.474 415.459 376.973] /A << /S /GoTo /D <74686D2E412E322E332E34> >> >>}
\MusicPDFWriteObject{5164}{<< /Type /Annot /Subtype /Link /Border [0 0 0] /H /I /C [1 0 0] /Rect [123.615 314.332 156.932 326.925] /A << /S /GoTo /D <74686D2E412E322E332E35> >> >>}
\MusicPDFWriteObject{5165}{<< /Type /Annot /Subtype /Link /Border [0 0 0] /H /I /C [1 0 0] /Rect [127.801 553.147 165.484 565.1] /A << /S /GoTo /D <6571756174696F6E2E412E322E332E3233> >> >>}
\MusicPDFWriteObject{5166}{<< /Type /Annot /Subtype /Link /Border [0 0 0] /H /I /C [1 0 0] /Rect [387.129 273.922 419.833 286.874] /A << /S /GoTo /D <74686D2E412E322E332E38> >> >>}
\MusicPDFWriteObject{5167}{<< /Type /Annot /Subtype /Link /Border [0 0 0] /H /I /C [1 0 0] /Rect [178.044 261.132 210.133 273.084] /A << /S /GoTo /D <74686D2E412E322E332E39> >> >>}
\MusicPDFWriteObject{5168}{<< /Type /Annot /Subtype /Link /Border [0 0 0] /H /I /C [1 0 0] /Rect [171.768 211.553 209.451 223.506] /A << /S /GoTo /D <6571756174696F6E2E412E322E332E3235> >> >>}
\MusicPDFWriteObject{5169}{<< /Type /Annot /Subtype /Link /Border [0 0 0] /H /I /C [1 0 0] /Rect [470.856 174.451 504.174 187.95] /A << /S /GoTo /D <74686D2E412E322E332E38> >> >>}
\MusicPDFWriteObject{5170}{<< /Type /Annot /Subtype /Link /Border [0 0 0] /H /I /C [1 0 0] /Rect [308.941 161.738 346.625 173.69] /A << /S /GoTo /D <6571756174696F6E2E412E322E332E3236> >> >>}
\MusicPDFWriteObject{5171}{<< /Type /Annot /Subtype /Link /Border [0 0 0] /H /I /C [1 0 0] /Rect [411.435 161.738 449.118 173.69] /A << /S /GoTo /D <6571756174696F6E2E412E322E332E3137> >> >>}
\MusicPDFWriteObject{5172}{<< /Type /Annot /Subtype /Link /Border [0 0 0] /H /I /C [1 0 0] /Rect [500.81 161.738 533.175 174.153] /A << /S /GoTo /D <74686D2E412E322E332E38> >> >>}
\MusicPDFWriteObject{5173}{<< /Type /Annot /Subtype /Link /Border [0 0 0] /H /I /C [1 0 0] /Rect [113.64 108.897 146.957 122.18] /A << /S /GoTo /D <74686D2E412E312E342E33> >> >>}
\MusicPDFWriteObject{5174}{<< /Type /Annot /Subtype /Link /Border [0 0 0] /H /I /C [1 0 0] /Rect [314.386 108.897 352.069 120.85] /A << /S /GoTo /D <6571756174696F6E2E412E322E332E3136> >> >>}
\MusicPDFWriteObject{5175}{<< /Type /Annot /Subtype /Link /Border [0 0 0] /H /I /C [1 0 0] /Rect [106.556 82.596 139.259 94.549] /A << /S /GoTo /D <74686D2E412E322E332E38> >> >>}
\MusicPDFWriteObject{5176}{<< /Type /Annot /Subtype /Link /Border [0 0 0] /H /I /C [1 0 0] /Rect [287.338 590.934 325.022 602.887] /A << /S /GoTo /D <74686D2E412E322E332E3130> >> >>}
\MusicPDFWriteObject{5177}{<< /Type /Annot /Subtype /Link /Border [0 0 0] /H /I /C [1 0 0] /Rect [366.713 333.202 400.03 346.486] /A << /S /GoTo /D <74686D2E412E322E332E38> >> >>}
\MusicPDFWriteObject{5178}{<< /Type /Annot /Subtype /Link /Border [0 0 0] /H /I /C [1 0 0] /Rect [106.556 267.717 139.259 282.495] /A << /S /GoTo /D <74686D2E412E322E332E38> >> >>}
\MusicPDFWriteObject{5179}{<< /Type /Annot /Subtype /Link /Border [0 0 0] /H /I /C [1 0 0] /Rect [267.983 108.897 300.686 120.85] /A << /S /GoTo /D <6571756174696F6E2E412E322E312E32> >> >>}
\MusicPDFWriteObject{5180}{<< /Type /Annot /Subtype /Link /Border [0 0 0] /H /I /C [1 0 0] /Rect [98.422 601.071 131.125 613.024] /A << /S /GoTo /D <6571756174696F6E2E412E322E312E31> >> >>}
\MusicPDFWriteObject{5181}{<< /Type /Annot /Subtype /Link /Border [0 0 0] /H /I /C [1 0 0] /Rect [197.548 458.072 230.252 470.025] /A << /S /GoTo /D <6571756174696F6E2E412E322E342E32> >> >>}
\MusicPDFWriteObject{5182}{<< /Type /Annot /Subtype /Link /Border [0 0 0] /H /I /C [1 0 0] /Rect [234.763 444.921 267.466 456.874] /A << /S /GoTo /D <6571756174696F6E2E412E322E342E31> >> >>}
\MusicPDFWriteObject{5183}{<< /Type /Annot /Subtype /Link /Border [0 0 0] /H /I /C [1 0 0] /Rect [228.959 369.087 262.092 382.73] /A << /S /GoTo /D <74686D2E412E322E312E32> >> >>}
\MusicPDFWriteObject{5184}{<< /Type /Annot /Subtype /Link /Border [0 0 0] /H /I /C [1 0 0] /Rect [125.057 356.296 157.76 368.249] /A << /S /GoTo /D <6571756174696F6E2E412E322E312E31> >> >>}
\MusicPDFWriteObject{5185}{<< /Type /Annot /Subtype /Link /Border [0 0 0] /H /I /C [1 0 0] /Rect [271.817 565.544 309.5 576.335] /A << /S /GoTo /D <74686D2E412E322E332E3130> >> >>}
\MusicPDFWriteObject{5186}{<< /Type /Annot /Subtype /Link /Border [0 0 0] /H /I /C [1 0 0] /Rect [458.734 226.327 483.828 237.176] /A << /S /GoTo /D <73756273656374696F6E2E442E312E31> >> >>}
\MusicPDFWriteObject{5187}{<< /Type /Annot /Subtype /Link /Border [0 0 0] /H /I /C [0 1 0] /Rect [318.528 582.917 336.499 591.785] /A << /S /GoTo /D <636974652E5256> >> >>}
\MusicPDFWriteObject{5188}{<< /Type /Annot /Subtype /Link /Border [0 0 0] /H /I /C [1 0 0] /Rect [266.318 335.156 272.281 348.374] /A << /S /GoTo /D <48666F6F746E6F74652E3535> >> >>}
\MusicPDFWriteObject{5189}{<< /Type /Annot /Subtype /Link /Border [0 0 0] /H /I /C [1 0 0] /Rect [265.242 311.615 297.945 323.568] /A << /S /GoTo /D <6571756174696F6E2E412E322E312E32> >> >>}
\MusicPDFWriteObject{5190}{<< /Type /Annot /Subtype /Link /Border [0 0 0] /H /I /C [0 1 0] /Rect [322.58 237.244 340.053 246.112] /A << /S /GoTo /D <636974652E4346> >> >>}
\MusicPDFWriteObject{5191}{<< /Type /Annot /Subtype /Link /Border [0 0 0] /H /I /C [1 0 0] /Rect [233.814 221.603 266.241 233.558] /A << /S /GoTo /D <74686D2E432E342E342E33> >> >>}
\MusicPDFWriteObject{5192}{<< /Type /Annot /Subtype /Link /Border [0 0 0] /H /I /C [0 1 0] /Rect [163.355 146.122 180.828 154.99] /A << /S /GoTo /D <636974652E4346> >> >>}
\MusicPDFWriteObject{5193}{<< /Type /Annot /Subtype /Link /Border [0 0 0] /H /I /C [1 0 0] /Rect [278.984 672.036 312.087 682.827] /A << /S /GoTo /D <74686D2E412E332E312E32> >> >>}
\MusicPDFWriteObject{5194}{<< /Type /Annot /Subtype /Link /Border [0 0 0] /H /I /C [1 0 0] /Rect [368.841 632.584 402.299 643.374] /A << /S /GoTo /D <74686D2E442E312E312E32> >> >>}
\MusicPDFWriteObject{5195}{<< /Type /Annot /Subtype /Link /Border [0 0 0] /H /I /C [1 0 0] /Rect [504.067 605.729 541.888 617.681] /A << /S /GoTo /D <6571756174696F6E2E442E312E312E3131> >> >>}
\MusicPDFWriteObject{5196}{<< /Type /Annot /Subtype /Link /Border [0 0 0] /H /I /C [1 0 0] /Rect [70.625 549.072 103.328 557.926] /A << /S /GoTo /D <74686D2E412E332E312E32> >> >>}
\MusicPDFWriteObject{5197}{<< /Type /Annot /Subtype /Link /Border [0 0 0] /H /I /C [1 0 0] /Rect [70.625 399.855 103.328 411.808] /A << /S /GoTo /D <74686D2E412E332E312E32> >> >>}
\MusicPDFWriteObject{5198}{<< /Type /Annot /Subtype /Link /Border [0 0 0] /H /I /C [0 1 0] /Rect [188.04 364.788 205.512 373.656] /A << /S /GoTo /D <636974652E4346> >> >>}
\MusicPDFWriteObject{5199}{<< /Type /Annot /Subtype /Link /Border [0 0 0] /H /I /C [1 0 0] /Rect [195.832 624.045 228.535 636.969] /A << /S /GoTo /D <74686D2E412E332E312E36> >> >>}
\MusicPDFWriteObject{5200}{<< /Type /Annot /Subtype /Link /Border [0 0 0] /H /I /C [1 0 0] /Rect [141.208 608.808 174.666 623.816] /A << /S /GoTo /D <74686D2E442E312E312E32> >> >>}
\MusicPDFWriteObject{5201}{<< /Type /Annot /Subtype /Link /Border [0 0 0] /H /I /C [1 0 0] /Rect [342.85 597.491 375.553 609.804] /A << /S /GoTo /D <74686D2E412E332E312E32> >> >>}
\MusicPDFWriteObject{5202}{<< /Type /Annot /Subtype /Link /Border [0 0 0] /H /I /C [1 0 0] /Rect [346.895 518.873 379.168 530.634] /A << /S /GoTo /D <74686D2E412E332E312E36> >> >>}
\MusicPDFWriteObject{5203}{<< /Type /Annot /Subtype /Link /Border [0 0 0] /H /I /C [1 0 0] /Rect [175.736 505.094 208.9 518.037] /A << /S /GoTo /D <74686D2E412E332E312E35> >> >>}
\MusicPDFWriteObject{5204}{<< /Type /Annot /Subtype /Link /Border [0 0 0] /H /I /C [0 1 0] /Rect [256.959 495.422 274.431 504.29] /A << /S /GoTo /D <636974652E4346> >> >>}
\MusicPDFWriteObject{5205}{<< /Type /Annot /Subtype /Link /Border [0 0 0] /H /I /C [1 0 0] /Rect [511.59 451.963 544.845 463.667] /A << /S /GoTo /D <74686D2E412E332E312E36> >> >>}
\MusicPDFWriteObject{5206}{<< /Type /Annot /Subtype /Link /Border [0 0 0] /H /I /C [0 1 0] /Rect [139.932 441.663 157.405 450.53] /A << /S /GoTo /D <636974652E4346> >> >>}
\MusicPDFWriteObject{5207}{<< /Type /Annot /Subtype /Link /Border [0 0 0] /H /I /C [0 1 0] /Rect [322.622 441.663 340.094 450.53] /A << /S /GoTo /D <636974652E4346> >> >>}
\MusicPDFWriteObject{5208}{<< /Type /Annot /Subtype /Link /Border [0 0 0] /H /I /C [1 0 0] /Rect [325.685 423.968 358.388 439.18] /A << /S /GoTo /D <74686D2E412E332E312E31> >> >>}
\MusicPDFWriteObject{5209}{<< /Type /Annot /Subtype /Link /Border [0 0 0] /H /I /C [1 0 0] /Rect [398.588 348.947 430.677 359.738] /A << /S /GoTo /D <74686D2E412E332E312E31> >> >>}
\MusicPDFWriteObject{5210}{<< /Type /Annot /Subtype /Link /Border [0 0 0] /H /I /C [0 1 0] /Rect [212.992 337.733 230.464 346.601] /A << /S /GoTo /D <636974652E4346> >> >>}
\MusicPDFWriteObject{5211}{<< /Type /Annot /Subtype /Link /Border [0 0 0] /H /I /C [1 0 0] /Rect [298.068 134.839 330.771 147.152] /A << /S /GoTo /D <74686D2E412E312E352E31> >> >>}
\MusicPDFWriteObject{5212}{<< /Type /Annot /Subtype /Link /Border [0 0 0] /H /I /C [1 0 0] /Rect [112.67 108.897 145.619 120.85] /A << /S /GoTo /D <74686D2E412E332E312E38> >> >>}
\MusicPDFWriteObject{5213}{<< /Type /Annot /Subtype /Link /Border [0 0 0] /H /I /C [1 0 0] /Rect [208.106 108.897 241.055 120.85] /A << /S /GoTo /D <74686D2E412E332E312E39> >> >>}
\MusicPDFWriteObject{5214}{<< /Type /Annot /Subtype /Link /Border [0 0 0] /H /I /C [1 0 0] /Rect [347.357 108.897 380.306 120.85] /A << /S /GoTo /D <74686D2E412E332E312E33> >> >>}
\MusicPDFWriteObject{5215}{<< /Type /Annot /Subtype /Link /Border [0 0 0] /H /I /C [1 0 0] /Rect [390.452 83.149 423.77 93.94] /A << /S /GoTo /D <74686D2E412E332E312E31> >> >>}
\MusicPDFWriteObject{5216}{<< /Type /Annot /Subtype /Link /Border [0 0 0] /H /I /C [1 0 0] /Rect [185.491 534.004 191.454 549.926] /A << /S /GoTo /D <48666F6F746E6F74652E3536> >> >>}
\MusicPDFWriteObject{5217}{<< /Type /Annot /Subtype /Link /Border [0 0 0] /H /I /C [1 0 0] /Rect [388.642 534.004 421.114 549.926] /A << /S /GoTo /D <74686D2E442E312E342E36> >> >>}
\MusicPDFWriteObject{5218}{<< /Type /Annot /Subtype /Link /Border [0 0 0] /H /I /C [1 0 0] /Rect [437.782 508.063 475.465 520.016] /A << /S /GoTo /D <6571756174696F6E2E412E332E312E3130> >> >>}
\MusicPDFWriteObject{5219}{<< /Type /Annot /Subtype /Link /Border [0 0 0] /H /I /C [0 1 0] /Rect [378.221 405.901 395.693 414.769] /A << /S /GoTo /D <636974652E4346> >> >>}
\MusicPDFWriteObject{5220}{<< /Type /Annot /Subtype /Link /Border [0 0 0] /H /I /C [1 0 0] /Rect [504.205 377.109 541.888 389.062] /A << /S /GoTo /D <6571756174696F6E2E412E332E312E3132> >> >>}
\MusicPDFWriteObject{5221}{<< /Type /Annot /Subtype /Link /Border [0 0 0] /H /I /C [1 0 0] /Rect [70.625 246.976 103.943 259.291] /A << /S /GoTo /D <74686D2E412E312E352E31> >> >>}
\MusicPDFWriteObject{5222}{<< /Type /Annot /Subtype /Link /Border [0 0 0] /H /I /C [1 0 0] /Rect [205.824 253.313 238.527 265.266] /A << /S /GoTo /D <6571756174696F6E2E412E332E322E33> >> >>}
\MusicPDFWriteObject{5223}{<< /Type /Annot /Subtype /Link /Border [0 0 0] /H /I /C [0 1 0] /Rect [123.131 146.25 141.101 155.118] /A << /S /GoTo /D <636974652E5256> >> >>}
\MusicPDFWriteObject{5224}{<< /Type /Annot /Subtype /Link /Border [0 0 0] /H /I /C [0 1 0] /Rect [156.819 85.086 174.291 93.954] /A << /S /GoTo /D <636974652E4346> >> >>}
\MusicPDFWriteObject{5225}{<< /Type /Annot /Subtype /Link /Border [0 0 0] /H /I /C [0 1 0] /Rect [410.202 687.123 426.194 695.991] /A << /S /GoTo /D <636974652E436F> >> >>}
\MusicPDFWriteObject{5226}{<< /Type /Annot /Subtype /Link /Border [0 0 0] /H /I /C [1 0 0] /Rect [283.682 658.332 313.895 670.284] /A << /S /GoTo /D <74686D2E322E322E332E32> >> >>}
\MusicPDFWriteObject{5227}{<< /Type /Annot /Subtype /Link /Border [0 0 0] /H /I /C [0 1 0] /Rect [362.644 660.822 380.116 669.69] /A << /S /GoTo /D <636974652E4346> >> >>}
\MusicPDFWriteObject{5228}{<< /Type /Annot /Subtype /Link /Border [0 0 0] /H /I /C [0 1 0] /Rect [294.131 634.52 314.496 643.388] /A << /S /GoTo /D <636974652E4F6573> >> >>}
\MusicPDFWriteObject{5229}{<< /Type /Annot /Subtype /Link /Border [0 0 0] /H /I /C [1 0 0] /Rect [468.708 509.855 501.411 521.811] /A << /S /GoTo /D <74686D2E412E322E322E39> >> >>}
\MusicPDFWriteObject{5230}{<< /Type /Annot /Subtype /Link /Border [0 0 0] /H /I /C [1 0 0] /Rect [188.461 423.596 221.164 436.442] /A << /S /GoTo /D <6571756174696F6E2E412E332E322E31> >> >>}
\MusicPDFWriteObject{5231}{<< /Type /Annot /Subtype /Link /Border [0 0 0] /H /I /C [1 0 0] /Rect [225.71 257.947 258.413 268.295] /A << /S /GoTo /D <74686D2E412E332E312E32> >> >>}
\MusicPDFWriteObject{5232}{<< /Type /Annot /Subtype /Link /Border [0 0 0] /H /I /C [1 0 0] /Rect [278.251 184.346 310.954 196.734] /A << /S /GoTo /D <74686D2E412E332E322E32> >> >>}
\MusicPDFWriteObject{5233}{<< /Type /Annot /Subtype /Link /Border [0 0 0] /H /I /C [1 0 0] /Rect [450.049 157.657 482.274 170.609] /A << /S /GoTo /D <74686D2E442E312E322E31> >> >>}
\MusicPDFWriteObject{5234}{<< /Type /Annot /Subtype /Link /Border [0 0 0] /H /I /C [1 0 0] /Rect [419.289 126.735 451.822 139.048] /A << /S /GoTo /D <74686D2E442E312E322E31> >> >>}
\MusicPDFWriteObject{5235}{<< /Type /Annot /Subtype /Link /Border [0 0 0] /H /I /C [1 0 0] /Rect [427.189 98.945 459.969 113.147] /A << /S /GoTo /D <74686D2E442E312E322E31> >> >>}
\MusicPDFWriteObject{5236}{<< /Type /Annot /Subtype /Link /Border [0 0 0] /H /I /C [1 0 0] /Rect [233.575 645.181 266.094 657.134] /A << /S /GoTo /D <74686D2E412E332E312E38> >> >>}
\MusicPDFWriteObject{5237}{<< /Type /Annot /Subtype /Link /Border [0 0 0] /H /I /C [1 0 0] /Rect [327.857 645.181 360.376 657.134] /A << /S /GoTo /D <74686D2E412E332E312E39> >> >>}
\MusicPDFWriteObject{5238}{<< /Type /Annot /Subtype /Link /Border [0 0 0] /H /I /C [1 0 0] /Rect [438.806 645.181 471.325 657.134] /A << /S /GoTo /D <74686D2E412E332E322E35> >> >>}
\MusicPDFWriteObject{5239}{<< /Type /Annot /Subtype /Link /Border [0 0 0] /H /I /C [1 0 0] /Rect [157.851 523.833 195.534 536.179] /A << /S /GoTo /D <74686D2E412E332E312E3130> >> >>}
\MusicPDFWriteObject{5240}{<< /Type /Annot /Subtype /Link /Border [0 0 0] /H /I /C [1 0 0] /Rect [497.788 511.324 530.722 523.639] /A << /S /GoTo /D <74686D2E442E312E322E31> >> >>}
\MusicPDFWriteObject{5241}{<< /Type /Annot /Subtype /Link /Border [0 0 0] /H /I /C [1 0 0] /Rect [475.928 485.022 508.294 498.711] /A << /S /GoTo /D <74686D2E412E332E312E33> >> >>}
\MusicPDFWriteObject{5242}{<< /Type /Annot /Subtype /Link /Border [0 0 0] /H /I /C [1 0 0] /Rect [495.512 410.549 528.829 422.502] /A << /S /GoTo /D <74686D2E412E312E352E32> >> >>}
\MusicPDFWriteObject{5243}{<< /Type /Annot /Subtype /Link /Border [0 0 0] /H /I /C [0 1 0] /Rect [73.776 384.801 138.269 395.606] /A << /S /GoTo /D <636974652E446176656E706F72742D31> >> >>}
\MusicPDFWriteObject{5244}{<< /Type /Annot /Subtype /Link /Border [0 0 0] /H /I /C [1 0 0] /Rect [415.256 358.5 445.469 369.348] /A << /S /GoTo /D <74686D2E322E342E322E31> >> >>}
\MusicPDFWriteObject{5245}{<< /Type /Annot /Subtype /Link /Border [0 0 0] /H /I /C [1 0 0] /Rect [448.822 333.853 478.979 344.657] /A << /S /GoTo /D <74686D2E322E342E322E31> >> >>}
\MusicPDFWriteObject{5246}{<< /Type /Annot /Subtype /Link /Border [0 0 0] /H /I /C [1 0 0] /Rect [112.174 309.488 118.136 319.855] /A << /S /GoTo /D <48666F6F746E6F74652E3537> >> >>}
\MusicPDFWriteObject{5247}{<< /Type /Annot /Subtype /Link /Border [0 0 0] /H /I /C [1 0 0] /Rect [185.209 684.633 217.298 696.586] /A << /S /GoTo /D <74686D2E412E342E312E31> >> >>}
\MusicPDFWriteObject{5248}{<< /Type /Annot /Subtype /Link /Border [0 0 0] /H /I /C [0 1 0] /Rect [478.327 687.123 507.364 695.991] /A << /S /GoTo /D <636974652E5363686D> >> >>}
\MusicPDFWriteObject{5249}{<< /Type /Annot /Subtype /Link /Border [0 0 0] /H /I /C [0 1 0] /Rect [278.702 672.036 343.195 682.84] /A << /S /GoTo /D <636974652E446176656E706F72742D31> >> >>}
\MusicPDFWriteObject{5250}{<< /Type /Annot /Subtype /Link /Border [0 0 0] /H /I /C [0 1 0] /Rect [140.476 632.584 204.969 643.388] /A << /S /GoTo /D <636974652E446176656E706F72742D31> >> >>}
\MusicPDFWriteObject{5251}{<< /Type /Annot /Subtype /Link /Border [0 0 0] /H /I /C [0 1 0] /Rect [409.68 634.52 438.718 643.388] /A << /S /GoTo /D <636974652E5363686D> >> >>}
\MusicPDFWriteObject{5252}{<< /Type /Annot /Subtype /Link /Border [0 0 0] /H /I /C [1 0 0] /Rect [118.59 375.169 156.273 387.122] /A << /S /GoTo /D <6571756174696F6E2E412E322E322E3131> >> >>}
\MusicPDFWriteObject{5253}{<< /Type /Annot /Subtype /Link /Border [0 0 0] /H /I /C [1 0 0] /Rect [271.989 335.717 304.692 347.67] /A << /S /GoTo /D <6571756174696F6E2E412E322E342E35> >> >>}
\MusicPDFWriteObject{5254}{<< /Type /Annot /Subtype /Link /Border [0 0 0] /H /I /C [1 0 0] /Rect [231.06 306.205 263.149 321.06] /A << /S /GoTo /D <6571756174696F6E2E412E342E312E33> >> >>}
\MusicPDFWriteObject{5255}{<< /Type /Annot /Subtype /Link /Border [0 0 0] /H /I /C [1 0 0] /Rect [200.241 243.315 233.558 255.271] /A << /S /GoTo /D <6571756174696F6E2E412E342E312E33> >> >>}
\MusicPDFWriteObject{5256}{<< /Type /Annot /Subtype /Link /Border [0 0 0] /H /I /C [1 0 0] /Rect [138.722 203.626 171.425 215.579] /A << /S /GoTo /D <6571756174696F6E2E412E342E312E36> >> >>}
\MusicPDFWriteObject{5257}{<< /Type /Annot /Subtype /Link /Border [0 0 0] /H /I /C [1 0 0] /Rect [247.804 134.839 280.09 147.153] /A << /S /GoTo /D <74686D2E442E312E322E31> >> >>}
\MusicPDFWriteObject{5258}{<< /Type /Annot /Subtype /Link /Border [0 0 0] /H /I /C [0 1 0] /Rect [398.55 569.929 427.587 578.797] /A << /S /GoTo /D <636974652E5363686D> >> >>}
\MusicPDFWriteObject{5259}{<< /Type /Annot /Subtype /Link /Border [0 0 0] /H /I /C [0 1 0] /Rect [133.726 556.778 162.764 565.646] /A << /S /GoTo /D <636974652E5363686D> >> >>}
\MusicPDFWriteObject{5260}{<< /Type /Annot /Subtype /Link /Border [0 0 0] /H /I /C [0 1 0] /Rect [73.815 468.235 102.852 477.103] /A << /S /GoTo /D <636974652E5363686D> >> >>}
\MusicPDFWriteObject{5261}{<< /Type /Annot /Subtype /Link /Border [0 0 0] /H /I /C [0 1 0] /Rect [322.333 441.934 351.37 450.801] /A << /S /GoTo /D <636974652E5363686D> >> >>}
\MusicPDFWriteObject{5262}{<< /Type /Annot /Subtype /Link /Border [0 0 0] /H /I /C [1 0 0] /Rect [455.143 377.513 487.417 389.828] /A << /S /GoTo /D <74686D2E412E342E312E32> >> >>}
\MusicPDFWriteObject{5263}{<< /Type /Annot /Subtype /Link /Border [0 0 0] /H /I /C [1 0 0] /Rect [172.478 297.169 208.639 306.023] /A << /S /GoTo /D <74686D2E412E342E312E32> >> >>}
\MusicPDFWriteObject{5264}{<< /Type /Annot /Subtype /Link /Border [0 0 0] /H /I /C [1 0 0] /Rect [325.232 280.93 357.936 292.883] /A << /S /GoTo /D <6571756174696F6E2E412E342E312E37> >> >>}
\MusicPDFWriteObject{5265}{<< /Type /Annot /Subtype /Link /Border [0 0 0] /H /I /C [1 0 0] /Rect [458.985 280.93 491.688 292.883] /A << /S /GoTo /D <6571756174696F6E2E412E342E312E34> >> >>}
\MusicPDFWriteObject{5266}{<< /Type /Annot /Subtype /Link /Border [0 0 0] /H /I /C [1 0 0] /Rect [206.918 319.436 239.007 331.391] /A << /S /GoTo /D <6571756174696F6E2E412E342E312E39> >> >>}
\MusicPDFWriteObject{5267}{<< /Type /Annot /Subtype /Link /Border [0 0 0] /H /I /C [1 0 0] /Rect [232.882 201.895 265.585 213.848] /A << /S /GoTo /D <6571756174696F6E2E412E342E312E38> >> >>}
\MusicPDFWriteObject{5268}{<< /Type /Annot /Subtype /Link /Border [0 0 0] /H /I /C [1 0 0] /Rect [229.731 188.744 262.434 200.697] /A << /S /GoTo /D <6571756174696F6E2E412E342E312E33> >> >>}
\MusicPDFWriteObject{5269}{<< /Type /Annot /Subtype /Link /Border [0 0 0] /H /I /C [1 0 0] /Rect [133.896 176.147 166.599 186.938] /A << /S /GoTo /D <74686D2E412E342E312E32> >> >>}
\MusicPDFWriteObject{5270}{<< /Type /Annot /Subtype /Link /Border [0 0 0] /H /I /C [1 0 0] /Rect [468.508 123.813 501.211 135.765] /A << /S /GoTo /D <6571756174696F6E2E412E342E312E33> >> >>}
\MusicPDFWriteObject{5271}{<< /Type /Annot /Subtype /Link /Border [0 0 0] /H /I /C [1 0 0] /Rect [449.284 95.747 481.988 107.699] /A << /S /GoTo /D <6571756174696F6E2E412E342E312E33> >> >>}
\MusicPDFWriteObject{5272}{<< /Type /Annot /Subtype /Link /Border [0 0 0] /H /I /C [0 1 0] /Rect [93.114 654.844 110.269 663.712] /A << /S /GoTo /D <636974652E466F6C> >> >>}
\MusicPDFWriteObject{5273}{<< /Type /Annot /Subtype /Link /Border [0 0 0] /H /I /C [0 1 0] /Rect [408.363 654.844 425.517 663.712] /A << /S /GoTo /D <636974652E466F6C> >> >>}
\MusicPDFWriteObject{5274}{<< /Type /Annot /Subtype /Link /Border [0 0 0] /H /I /C [0 1 0] /Rect [83.786 628.543 112.879 637.411] /A << /S /GoTo /D <636974652E48522D32> >> >>}
\MusicPDFWriteObject{5275}{<< /Type /Annot /Subtype /Link /Border [0 0 0] /H /I /C [1 0 0] /Rect [121.524 527.222 154.227 539.175] /A << /S /GoTo /D <6571756174696F6E2E412E352E312E31> >> >>}
\MusicPDFWriteObject{5276}{<< /Type /Annot /Subtype /Link /Border [0 0 0] /H /I /C [0 1 0] /Rect [313.923 516.562 331.077 525.43] /A << /S /GoTo /D <636974652E466F6C> >> >>}
\MusicPDFWriteObject{5277}{<< /Type /Annot /Subtype /Link /Border [0 0 0] /H /I /C [0 1 0] /Rect [404.893 516.562 422.047 525.43] /A << /S /GoTo /D <636974652E466F6C> >> >>}
\MusicPDFWriteObject{5278}{<< /Type /Annot /Subtype /Link /Border [0 0 0] /H /I /C [1 0 0] /Rect [114.746 435.72 120.708 448.025] /A << /S /GoTo /D <48666F6F746E6F74652E3538> >> >>}
\MusicPDFWriteObject{5279}{<< /Type /Annot /Subtype /Link /Border [0 0 0] /H /I /C [0 1 0] /Rect [302.02 115.061 331.112 123.929] /A << /S /GoTo /D <636974652E48522D32> >> >>}
\MusicPDFWriteObject{5280}{<< /Type /Annot /Subtype /Link /Border [0 0 0] /H /I /C [0 1 0] /Rect [139.62 558.587 168.712 567.455] /A << /S /GoTo /D <636974652E48522D32> >> >>}
\MusicPDFWriteObject{5281}{<< /Type /Annot /Subtype /Link /Border [0 0 0] /H /I /C [0 1 0] /Rect [159.493 545.436 176.648 554.304] /A << /S /GoTo /D <636974652E466F6C> >> >>}
\MusicPDFWriteObject{5282}{<< /Type /Annot /Subtype /Link /Border [0 0 0] /H /I /C [0 1 0] /Rect [394.658 428.393 443.249 437.261] /A << /S /GoTo /D <636974652E53745368612D494949> >> >>}
\MusicPDFWriteObject{5283}{<< /Type /Annot /Subtype /Link /Border [0 0 0] /H /I /C [1 0 0] /Rect [175.31 412.753 208.013 424.705] /A << /S /GoTo /D <6571756174696F6E2E412E352E312E33> >> >>}
\MusicPDFWriteObject{5284}{<< /Type /Annot /Subtype /Link /Border [0 0 0] /H /I /C [0 1 0] /Rect [381.491 375.79 398.645 384.658] /A << /S /GoTo /D <636974652E466F6C> >> >>}
\MusicPDFWriteObject{5285}{<< /Type /Annot /Subtype /Link /Border [0 0 0] /H /I /C [1 0 0] /Rect [138.512 217.921 171.215 229.874] /A << /S /GoTo /D <6571756174696F6E2E412E352E312E39> >> >>}
\MusicPDFWriteObject{5286}{<< /Type /Annot /Subtype /Link /Border [0 0 0] /H /I /C [1 0 0] /Rect [199.656 217.921 232.36 229.874] /A << /S /GoTo /D <6571756174696F6E2E412E352E312E38> >> >>}
\MusicPDFWriteObject{5287}{<< /Type /Annot /Subtype /Link /Border [0 0 0] /H /I /C [1 0 0] /Rect [149.994 119.146 187.677 131.099] /A << /S /GoTo /D <6571756174696F6E2E412E352E312E3131> >> >>}
\MusicPDFWriteObject{5288}{<< /Type /Annot /Subtype /Link /Border [0 0 0] /H /I /C [1 0 0] /Rect [216.118 119.146 253.802 131.099] /A << /S /GoTo /D <6571756174696F6E2E412E352E312E3130> >> >>}
\MusicPDFWriteObject{5289}{<< /Type /Annot /Subtype /Link /Border [0 0 0] /H /I /C [1 0 0] /Rect [332.321 119.146 365.024 131.099] /A << /S /GoTo /D <6571756174696F6E2E412E352E312E34> >> >>}
\MusicPDFWriteObject{5290}{<< /Type /Annot /Subtype /Link /Border [0 0 0] /H /I /C [1 0 0] /Rect [279.84 685.187 285.803 697.491] /A << /S /GoTo /D <48666F6F746E6F74652E3539> >> >>}
\MusicPDFWriteObject{5291}{<< /Type /Annot /Subtype /Link /Border [0 0 0] /H /I /C [0 1 0] /Rect [268.157 673.973 301.614 682.84] /A << /S /GoTo /D <636974652E48755765> >> >>}
\MusicPDFWriteObject{5292}{<< /Type /Annot /Subtype /Link /Border [0 0 0] /H /I /C [1 0 0] /Rect [404.09 671.482 436.793 683.435] /A << /S /GoTo /D <6571756174696F6E2E412E352E312E34> >> >>}
\MusicPDFWriteObject{5293}{<< /Type /Annot /Subtype /Link /Border [0 0 0] /H /I /C [0 1 0] /Rect [192.538 647.671 241.129 656.539] /A << /S /GoTo /D <636974652E53745368612D494949> >> >>}
\MusicPDFWriteObject{5294}{<< /Type /Annot /Subtype /Link /Border [0 0 0] /H /I /C [0 1 0] /Rect [314.878 634.52 332.032 643.388] /A << /S /GoTo /D <636974652E466F6C> >> >>}
\MusicPDFWriteObject{5295}{<< /Type /Annot /Subtype /Link /Border [0 0 0] /H /I /C [1 0 0] /Rect [459.421 618.879 492.124 630.832] /A << /S /GoTo /D <6571756174696F6E2E412E352E312E33> >> >>}
\MusicPDFWriteObject{5296}{<< /Type /Annot /Subtype /Link /Border [0 0 0] /H /I /C [0 1 0] /Rect [230.256 608.219 247.411 617.087] /A << /S /GoTo /D <636974652E466F6C> >> >>}
\MusicPDFWriteObject{5297}{<< /Type /Annot /Subtype /Link /Border [0 0 0] /H /I /C [0 1 0] /Rect [353.569 608.219 402.16 617.087] /A << /S /GoTo /D <636974652E53745368612D494949> >> >>}
\MusicPDFWriteObject{5298}{<< /Type /Annot /Subtype /Link /Border [0 0 0] /H /I /C [1 0 0] /Rect [273.331 592.578 305.85 604.531] /A << /S /GoTo /D <74686D2E412E352E312E31> >> >>}
\MusicPDFWriteObject{5299}{<< /Type /Annot /Subtype /Link /Border [0 0 0] /H /I /C [1 0 0] /Rect [118.142 566.276 150.846 578.229] /A << /S /GoTo /D <6571756174696F6E2E412E352E312E34> >> >>}
\MusicPDFWriteObject{5300}{<< /Type /Annot /Subtype /Link /Border [0 0 0] /H /I /C [0 1 0] /Rect [293.724 568.766 322.07 577.634] /A << /S /GoTo /D <636974652E5374537472> >> >>}
\MusicPDFWriteObject{5301}{<< /Type /Annot /Subtype /Link /Border [0 0 0] /H /I /C [1 0 0] /Rect [171.108 539.975 203.811 551.928] /A << /S /GoTo /D <6571756174696F6E2E412E352E312E33> >> >>}
\MusicPDFWriteObject{5302}{<< /Type /Annot /Subtype /Link /Border [0 0 0] /H /I /C [1 0 0] /Rect [347.387 539.975 380.09 551.928] /A << /S /GoTo /D <6571756174696F6E2E412E352E312E34> >> >>}
\MusicPDFWriteObject{5303}{<< /Type /Annot /Subtype /Link /Border [0 0 0] /H /I /C [1 0 0] /Rect [509.185 506.23 541.888 518.183] /A << /S /GoTo /D <6571756174696F6E2E412E352E312E33> >> >>}
\MusicPDFWriteObject{5304}{<< /Type /Annot /Subtype /Link /Border [0 0 0] /H /I /C [0 1 0] /Rect [107.3 320.928 124.454 329.795] /A << /S /GoTo /D <636974652E466F6C> >> >>}
\MusicPDFWriteObject{5305}{<< /Type /Annot /Subtype /Link /Border [0 0 0] /H /I /C [1 0 0] /Rect [183.838 278.985 216.541 290.938] /A << /S /GoTo /D <6571756174696F6E2E412E352E312E33> >> >>}
\MusicPDFWriteObject{5306}{<< /Type /Annot /Subtype /Link /Border [0 0 0] /H /I /C [1 0 0] /Rect [287.368 278.625 320.072 290.938] /A << /S /GoTo /D <74686D2E412E352E312E31> >> >>}
\MusicPDFWriteObject{5307}{<< /Type /Annot /Subtype /Link /Border [0 0 0] /H /I /C [1 0 0] /Rect [538.832 206.893 544.795 220.803] /A << /S /GoTo /D <48666F6F746E6F74652E3630> >> >>}
\MusicPDFWriteObject{5308}{<< /Type /Annot /Subtype /Link /Border [0 0 0] /H /I /C [1 0 0] /Rect [378.934 141.976 416.618 153.929] /A << /S /GoTo /D <6571756174696F6E2E412E352E312E3132> >> >>}
\MusicPDFWriteObject{5309}{<< /Type /Annot /Subtype /Link /Border [0 0 0] /H /I /C [1 0 0] /Rect [477.608 128.825 515.291 140.778] /A << /S /GoTo /D <6571756174696F6E2E412E352E312E3132> >> >>}
\MusicPDFWriteObject{5310}{<< /Type /Annot /Subtype /Link /Border [0 0 0] /H /I /C [1 0 0] /Rect [452.328 672.036 484.417 682.827] /A << /S /GoTo /D <74686D2E412E352E312E31> >> >>}
\MusicPDFWriteObject{5311}{<< /Type /Annot /Subtype /Link /Border [0 0 0] /H /I /C [1 0 0] /Rect [70.625 658.885 103.943 669.676] /A << /S /GoTo /D <74686D2E412E352E312E33> >> >>}
\MusicPDFWriteObject{5312}{<< /Type /Annot /Subtype /Link /Border [0 0 0] /H /I /C [0 1 0] /Rect [218.018 634.52 247.111 643.388] /A << /S /GoTo /D <636974652E48522D32> >> >>}
\MusicPDFWriteObject{5313}{<< /Type /Annot /Subtype /Link /Border [0 0 0] /H /I /C [1 0 0] /Rect [222.156 619.433 254.245 630.223] /A << /S /GoTo /D <74686D2E412E352E312E31> >> >>}
\MusicPDFWriteObject{5314}{<< /Type /Annot /Subtype /Link /Border [0 0 0] /H /I /C [1 0 0] /Rect [314.556 619.433 346.645 630.223] /A << /S /GoTo /D <74686D2E412E352E312E33> >> >>}
\MusicPDFWriteObject{5315}{<< /Type /Annot /Subtype /Link /Border [0 0 0] /H /I /C [0 1 0] /Rect [191.993 489.862 223.937 498.73] /A << /S /GoTo /D <636974652E427269732D61> >> >>}
\MusicPDFWriteObject{5316}{<< /Type /Annot /Subtype /Link /Border [0 0 0] /H /I /C [0 1 0] /Rect [228.946 489.862 261.685 498.73] /A << /S /GoTo /D <636974652E427269732D62> >> >>}
\MusicPDFWriteObject{5317}{<< /Type /Annot /Subtype /Link /Border [0 0 0] /H /I /C [0 1 0] /Rect [168.504 588.713 186.806 597.581] /A << /S /GoTo /D <636974652E4445> >> >>}
\MusicPDFWriteObject{5318}{<< /Type /Annot /Subtype /Link /Border [0 0 0] /H /I /C [0 1 0] /Rect [264.699 588.713 283.001 597.581] /A << /S /GoTo /D <636974652E4445> >> >>}
\MusicPDFWriteObject{5319}{<< /Type /Annot /Subtype /Link /Border [0 0 0] /H /I /C [0 1 0] /Rect [334.637 374.039 352.94 382.907] /A << /S /GoTo /D <636974652E4445> >> >>}
\MusicPDFWriteObject{5320}{<< /Type /Annot /Subtype /Link /Border [0 0 0] /H /I /C [1 0 0] /Rect [509.185 318.946 541.888 330.899] /A << /S /GoTo /D <6571756174696F6E2E412E352E322E32> >> >>}
\MusicPDFWriteObject{5321}{<< /Type /Annot /Subtype /Link /Border [0 0 0] /H /I /C [1 0 0] /Rect [110.723 122.048 144.607 134.003] /A << /S /GoTo /D <74686D2E442E312E322E31> >> >>}
\MusicPDFWriteObject{5322}{<< /Type /Annot /Subtype /Link /Border [0 0 0] /H /I /C [1 0 0] /Rect [252.765 603.56 290.306 614.35] /A << /S /GoTo /D <74686D2E412E352E322E3130> >> >>}
\MusicPDFWriteObject{5323}{<< /Type /Annot /Subtype /Link /Border [0 0 0] /H /I /C [0 1 0] /Rect [413.028 547.146 440.384 556.014] /A << /S /GoTo /D <636974652E4B656368> >> >>}
\MusicPDFWriteObject{5324}{<< /Type /Annot /Subtype /Link /Border [0 0 0] /H /I /C [1 0 0] /Rect [114.514 500.01 120.477 512.315] /A << /S /GoTo /D <48666F6F746E6F74652E3631> >> >>}
\MusicPDFWriteObject{5325}{<< /Type /Annot /Subtype /Link /Border [0 0 0] /H /I /C [1 0 0] /Rect [436.926 427.956 469.783 439.911] /A << /S /GoTo /D <74686D2E412E352E322E38> >> >>}
\MusicPDFWriteObject{5326}{<< /Type /Annot /Subtype /Link /Border [0 0 0] /H /I /C [0 1 0] /Rect [401.863 385.247 419.017 394.115] /A << /S /GoTo /D <636974652E466F6C> >> >>}
\MusicPDFWriteObject{5327}{<< /Type /Annot /Subtype /Link /Border [0 0 0] /H /I /C [0 1 0] /Rect [381.69 319.493 424.713 328.361] /A << /S /GoTo /D <636974652E426F752D4754> >> >>}
\MusicPDFWriteObject{5328}{<< /Type /Annot /Subtype /Link /Border [0 0 0] /H /I /C [0 1 0] /Rect [90.061 306.343 133.084 315.21] /A << /S /GoTo /D <636974652E426F752D4754> >> >>}
\MusicPDFWriteObject{5329}{<< /Type /Annot /Subtype /Link /Border [0 0 0] /H /I /C [0 1 0] /Rect [355.091 293.192 397.948 302.06] /A << /S /GoTo /D <636974652E4269726B686F6666> >> >>}
\MusicPDFWriteObject{5330}{<< /Type /Annot /Subtype /Link /Border [0 0 0] /H /I /C [0 1 0] /Rect [73.77 266.89 112.457 275.758] /A << /S /GoTo /D <636974652E53747275626C65> >> >>}
\MusicPDFWriteObject{5331}{<< /Type /Annot /Subtype /Link /Border [0 0 0] /H /I /C [0 1 0] /Rect [168.865 227.438 207.551 236.306] /A << /S /GoTo /D <636974652E53747275626C65> >> >>}
\MusicPDFWriteObject{5332}{<< /Type /Annot /Subtype /Link /Border [0 0 0] /H /I /C [0 1 0] /Rect [468.405 201.137 511.428 210.004] /A << /S /GoTo /D <636974652E426F752D4754> >> >>}
\MusicPDFWriteObject{5333}{<< /Type /Annot /Subtype /Link /Border [0 0 0] /H /I /C [0 1 0] /Rect [504.473 148.534 543.159 157.401] /A << /S /GoTo /D <636974652E53747275626C65> >> >>}
\MusicPDFWriteObject{5334}{<< /Type /Annot /Subtype /Link /Border [0 0 0] /H /I /C [1 0 0] /Rect [459.899 114.548 465.862 126.852] /A << /S /GoTo /D <48666F6F746E6F74652E3632> >> >>}
\MusicPDFWriteObject{5335}{<< /Type /Annot /Subtype /Link /Border [0 0 0] /H /I /C [0 1 0] /Rect [107.377 562.884 128.972 573.689] /A << /S /GoTo /D <636974652E426F67> >> >>}
\MusicPDFWriteObject{5336}{<< /Type /Annot /Subtype /Link /Border [0 0 0] /H /I /C [0 1 0] /Rect [160.037 551.67 203.06 560.538] /A << /S /GoTo /D <636974652E426F752D4754> >> >>}
\MusicPDFWriteObject{5337}{<< /Type /Annot /Subtype /Link /Border [0 0 0] /H /I /C [1 0 0] /Rect [203.197 289.002 239.358 297.856] /A << /S /GoTo /D <74686D2E412E352E322E36> >> >>}
\MusicPDFWriteObject{5338}{<< /Type /Annot /Subtype /Link /Border [0 0 0] /H /I /C [0 1 0] /Rect [107.564 262.701 132.216 271.569] /A << /S /GoTo /D <636974652E4B6C6565> >> >>}
\MusicPDFWriteObject{5339}{<< /Type /Annot /Subtype /Link /Border [0 0 0] /H /I /C [0 1 0] /Rect [106.454 223.248 131.106 232.116] /A << /S /GoTo /D <636974652E4B6C6565> >> >>}
\MusicPDFWriteObject{5340}{<< /Type /Annot /Subtype /Link /Border [0 0 0] /H /I /C [1 0 0] /Rect [281.052 181.86 287.015 194.164] /A << /S /GoTo /D <48666F6F746E6F74652E3633> >> >>}
\MusicPDFWriteObject{5341}{<< /Type /Annot /Subtype /Link /Border [0 0 0] /H /I /C [1 0 0] /Rect [218.141 155.004 250.23 166.959] /A << /S /GoTo /D <74686D2E412E352E322E38> >> >>}
\MusicPDFWriteObject{5342}{<< /Type /Annot /Subtype /Link /Border [0 0 0] /H /I /C [1 0 0] /Rect [145.615 115.552 177.704 127.507] /A << /S /GoTo /D <74686D2E412E352E322E38> >> >>}
\MusicPDFWriteObject{5343}{<< /Type /Annot /Subtype /Link /Border [0 0 0] /H /I /C [1 0 0] /Rect [228.074 684.633 265.044 696.588] /A << /S /GoTo /D <74686D2E412E352E322E3130> >> >>}
\MusicPDFWriteObject{5344}{<< /Type /Annot /Subtype /Link /Border [0 0 0] /H /I /C [0 1 0] /Rect [374.6 660.822 409.655 669.69] /A << /S /GoTo /D <636974652E426F752D4949> >> >>}
\MusicPDFWriteObject{5345}{<< /Type /Annot /Subtype /Link /Border [0 0 0] /H /I /C [1 0 0] /Rect [411.457 598.831 449.034 609.622] /A << /S /GoTo /D <74686D2E412E352E322E3130> >> >>}
\MusicPDFWriteObject{5346}{<< /Type /Annot /Subtype /Link /Border [0 0 0] /H /I /C [1 0 0] /Rect [191.959 585.127 224.662 597.08] /A << /S /GoTo /D <6571756174696F6E2E412E352E322E34> >> >>}
\MusicPDFWriteObject{5347}{<< /Type /Annot /Subtype /Link /Border [0 0 0] /H /I /C [1 0 0] /Rect [285.419 583.992 318.661 597.081] /A << /S /GoTo /D <74686D2E442E312E322E33> >> >>}
\MusicPDFWriteObject{5348}{<< /Type /Annot /Subtype /Link /Border [0 0 0] /H /I /C [0 1 0] /Rect [101.438 574.466 136.493 583.334] /A << /S /GoTo /D <636974652E426F752D4949> >> >>}
\MusicPDFWriteObject{5349}{<< /Type /Annot /Subtype /Link /Border [0 0 0] /H /I /C [0 1 0] /Rect [338.02 561.315 373.075 570.183] /A << /S /GoTo /D <636974652E426F752D4949> >> >>}
\MusicPDFWriteObject{5350}{<< /Type /Annot /Subtype /Link /Border [0 0 0] /H /I /C [1 0 0] /Rect [382.01 446.518 414.714 458.471] /A << /S /GoTo /D <6571756174696F6E2E412E352E322E36> >> >>}
\MusicPDFWriteObject{5351}{<< /Type /Annot /Subtype /Link /Border [0 0 0] /H /I /C [1 0 0] /Rect [457.14 406.768 490.458 419.658] /A << /S /GoTo /D <74686D2E412E352E322E38> >> >>}
\MusicPDFWriteObject{5352}{<< /Type /Annot /Subtype /Link /Border [0 0 0] /H /I /C [1 0 0] /Rect [337.143 393.915 369.846 405.868] /A << /S /GoTo /D <6571756174696F6E2E412E352E322E36> >> >>}
\MusicPDFWriteObject{5353}{<< /Type /Annot /Subtype /Link /Border [0 0 0] /H /I /C [1 0 0] /Rect [394.881 144.05 427.584 156.003] /A << /S /GoTo /D <6571756174696F6E2E412E352E322E36> >> >>}
\MusicPDFWriteObject{5354}{<< /Type /Annot /Subtype /Link /Border [0 0 0] /H /I /C [1 0 0] /Rect [509.185 144.05 541.888 156.003] /A << /S /GoTo /D <6571756174696F6E2E412E352E322E32> >> >>}
\MusicPDFWriteObject{5355}{<< /Type /Annot /Subtype /Link /Border [0 0 0] /H /I /C [0 1 0] /Rect [267.994 575.571 292.646 584.439] /A << /S /GoTo /D <636974652E4B6C6565> >> >>}
\MusicPDFWriteObject{5356}{<< /Type /Annot /Subtype /Link /Border [0 0 0] /H /I /C [1 0 0] /Rect [162.521 461.529 200.204 473.482] /A << /S /GoTo /D <6571756174696F6E2E412E322E332E3134> >> >>}
\MusicPDFWriteObject{5357}{<< /Type /Annot /Subtype /Link /Border [0 0 0] /H /I /C [1 0 0] /Rect [98.232 306.101 130.936 318.054] /A << /S /GoTo /D <6571756174696F6E2E412E322E332E39> >> >>}
\MusicPDFWriteObject{5358}{<< /Type /Annot /Subtype /Link /Border [0 0 0] /H /I /C [1 0 0] /Rect [200.343 178.833 238.027 190.786] /A << /S /GoTo /D <6571756174696F6E2E412E322E332E3134> >> >>}
\MusicPDFWriteObject{5359}{<< /Type /Annot /Subtype /Link /Border [0 0 0] /H /I /C [1 0 0] /Rect [288.92 599.181 322.115 609.971] /A << /S /GoTo /D <74686D2E412E322E342E33> >> >>}
\MusicPDFWriteObject{5360}{<< /Type /Annot /Subtype /Link /Border [0 0 0] /H /I /C [1 0 0] /Rect [90.739 122.584 128.422 134.537] /A << /S /GoTo /D <6571756174696F6E2E412E322E332E3234> >> >>}
\MusicPDFWriteObject{5361}{<< /Type /Annot /Subtype /Link /Border [0 0 0] /H /I /C [1 0 0] /Rect [235.305 618.587 272.988 630.539] /A << /S /GoTo /D <6571756174696F6E2E412E352E322E3134> >> >>}
\MusicPDFWriteObject{5362}{<< /Type /Annot /Subtype /Link /Border [0 0 0] /H /I /C [1 0 0] /Rect [244.837 497.096 282.52 509.048] /A << /S /GoTo /D <6571756174696F6E2E412E352E322E3132> >> >>}
\MusicPDFWriteObject{5363}{<< /Type /Annot /Subtype /Link /Border [0 0 0] /H /I /C [1 0 0] /Rect [464.049 483.585 473.982 496.802] /A << /S /GoTo /D <48666F6F746E6F74652E3634> >> >>}
\MusicPDFWriteObject{5364}{<< /Type /Annot /Subtype /Link /Border [0 0 0] /H /I /C [1 0 0] /Rect [233.798 448.031 243.731 460.335] /A << /S /GoTo /D <48666F6F746E6F74652E3635> >> >>}
\MusicPDFWriteObject{5365}{<< /Type /Annot /Subtype /Link /Border [0 0 0] /H /I /C [1 0 0] /Rect [157.937 381.426 167.87 394.581] /A << /S /GoTo /D <48666F6F746E6F74652E3636> >> >>}
\MusicPDFWriteObject{5366}{<< /Type /Annot /Subtype /Link /Border [0 0 0] /H /I /C [1 0 0] /Rect [218.519 192.307 251.223 204.26] /A << /S /GoTo /D <74686D2E412E352E322E33> >> >>}
\MusicPDFWriteObject{5367}{<< /Type /Annot /Subtype /Link /Border [0 0 0] /H /I /C [1 0 0] /Rect [370.828 192.307 403.531 204.26] /A << /S /GoTo /D <74686D2E412E352E322E33> >> >>}
\MusicPDFWriteObject{5368}{<< /Type /Annot /Subtype /Link /Border [0 0 0] /H /I /C [1 0 0] /Rect [472.505 192.307 505.208 204.26] /A << /S /GoTo /D <74686D2E412E352E322E35> >> >>}
\MusicPDFWriteObject{5369}{<< /Type /Annot /Subtype /Link /Border [0 0 0] /H /I /C [1 0 0] /Rect [421.087 164.157 453.79 178.359] /A << /S /GoTo /D <74686D2E412E352E322E36> >> >>}
\MusicPDFWriteObject{5370}{<< /Type /Annot /Subtype /Link /Border [0 0 0] /H /I /C [1 0 0] /Rect [290.308 669.451 327.992 681.404] /A << /S /GoTo /D <6571756174696F6E2E412E352E322E3231> >> >>}
\MusicPDFWriteObject{5371}{<< /Type /Annot /Subtype /Link /Border [0 0 0] /H /I /C [1 0 0] /Rect [327.915 573.594 361.11 584.384] /A << /S /GoTo /D <74686D2E412E352E322E35> >> >>}
\MusicPDFWriteObject{5372}{<< /Type /Annot /Subtype /Link /Border [0 0 0] /H /I /C [1 0 0] /Rect [421.462 573.594 459.716 584.384] /A << /S /GoTo /D <74686D2E412E352E322E3134> >> >>}
\MusicPDFWriteObject{5373}{<< /Type /Annot /Subtype /Link /Border [0 0 0] /H /I /C [1 0 0] /Rect [200.048 559.592 237.731 573.172] /A << /S /GoTo /D <74686D2E412E352E322E3132> >> >>}
\MusicPDFWriteObject{5374}{<< /Type /Annot /Subtype /Link /Border [0 0 0] /H /I /C [1 0 0] /Rect [324.771 559.592 357.475 573.172] /A << /S /GoTo /D <74686D2E412E352E322E36> >> >>}
\MusicPDFWriteObject{5375}{<< /Type /Annot /Subtype /Link /Border [0 0 0] /H /I /C [1 0 0] /Rect [218.519 371.285 251.223 383.238] /A << /S /GoTo /D <74686D2E412E352E322E33> >> >>}
\MusicPDFWriteObject{5376}{<< /Type /Annot /Subtype /Link /Border [0 0 0] /H /I /C [1 0 0] /Rect [370.828 371.285 403.531 383.238] /A << /S /GoTo /D <74686D2E412E352E322E33> >> >>}
\MusicPDFWriteObject{5377}{<< /Type /Annot /Subtype /Link /Border [0 0 0] /H /I /C [1 0 0] /Rect [472.505 371.285 505.208 383.238] /A << /S /GoTo /D <74686D2E412E352E322E35> >> >>}
\MusicPDFWriteObject{5378}{<< /Type /Annot /Subtype /Link /Border [0 0 0] /H /I /C [1 0 0] /Rect [222.664 355.665 259.74 370.727] /A << /S /GoTo /D <74686D2E412E352E322E3132> >> >>}
\MusicPDFWriteObject{5379}{<< /Type /Annot /Subtype /Link /Border [0 0 0] /H /I /C [1 0 0] /Rect [417.368 317.407 450.071 329.721] /A << /S /GoTo /D <74686D2E412E352E322E36> >> >>}
\MusicPDFWriteObject{5380}{<< /Type /Annot /Subtype /Link /Border [0 0 0] /H /I /C [1 0 0] /Rect [278.6 235.475 316.283 247.428] /A << /S /GoTo /D <6571756174696F6E2E412E352E322E3233> >> >>}
\MusicPDFWriteObject{5381}{<< /Type /Annot /Subtype /Link /Border [0 0 0] /H /I /C [1 0 0] /Rect [200.915 95.746 233.342 107.719] /A << /S /GoTo /D <74686D2E412E352E322E35> >> >>}
\MusicPDFWriteObject{5382}{<< /Type /Annot /Subtype /Link /Border [0 0 0] /H /I /C [1 0 0] /Rect [310.803 95.747 348.486 107.699] /A << /S /GoTo /D <6571756174696F6E2E412E352E322E3234> >> >>}
\MusicPDFWriteObject{5383}{<< /Type /Annot /Subtype /Link /Border [0 0 0] /H /I /C [1 0 0] /Rect [271.1 671.482 308.783 684.075] /A << /S /GoTo /D <74686D2E412E352E322E3136> >> >>}
\MusicPDFWriteObject{5384}{<< /Type /Annot /Subtype /Link /Border [0 0 0] /H /I /C [1 0 0] /Rect [410.51 656.54 442.599 671.615] /A << /S /GoTo /D <74686D2E412E352E322E36> >> >>}
\MusicPDFWriteObject{5385}{<< /Type /Annot /Subtype /Link /Border [0 0 0] /H /I /C [1 0 0] /Rect [120.211 644.646 157.752 657.536] /A << /S /GoTo /D <74686D2E412E352E322E3132> >> >>}
\MusicPDFWriteObject{5386}{<< /Type /Annot /Subtype /Link /Border [0 0 0] /H /I /C [1 0 0] /Rect [156.217 551.493 193.9 563.446] /A << /S /GoTo /D <6571756174696F6E2E412E352E322E3234> >> >>}
\MusicPDFWriteObject{5387}{<< /Type /Annot /Subtype /Link /Border [0 0 0] /H /I /C [1 0 0] /Rect [188.982 537.982 226.665 551.625] /A << /S /GoTo /D <74686D2E412E352E322E3137> >> >>}
\MusicPDFWriteObject{5388}{<< /Type /Annot /Subtype /Link /Border [0 0 0] /H /I /C [1 0 0] /Rect [320.623 537.982 358.306 551.625] /A << /S /GoTo /D <74686D2E412E352E322E3135> >> >>}
\MusicPDFWriteObject{5389}{<< /Type /Annot /Subtype /Link /Border [0 0 0] /H /I /C [1 0 0] /Rect [163.85 524.594 201.533 536.547] /A << /S /GoTo /D <6571756174696F6E2E412E352E322E3233> >> >>}
\MusicPDFWriteObject{5390}{<< /Type /Annot /Subtype /Link /Border [0 0 0] /H /I /C [1 0 0] /Rect [196.615 509.833 234.298 524.726] /A << /S /GoTo /D <74686D2E412E352E322E3136> >> >>}
\MusicPDFWriteObject{5391}{<< /Type /Annot /Subtype /Link /Border [0 0 0] /H /I /C [1 0 0] /Rect [401.308 509.833 438.991 524.726] /A << /S /GoTo /D <74686D2E412E352E322E3136> >> >>}
\MusicPDFWriteObject{5392}{<< /Type /Annot /Subtype /Link /Border [0 0 0] /H /I /C [1 0 0] /Rect [156.217 498.293 193.9 510.245] /A << /S /GoTo /D <6571756174696F6E2E412E352E322E3232> >> >>}
\MusicPDFWriteObject{5393}{<< /Type /Annot /Subtype /Link /Border [0 0 0] /H /I /C [1 0 0] /Rect [188.982 484.782 226.665 498.425] /A << /S /GoTo /D <74686D2E412E352E322E3135> >> >>}
\MusicPDFWriteObject{5394}{<< /Type /Annot /Subtype /Link /Border [0 0 0] /H /I /C [1 0 0] /Rect [163.427 471.991 201.11 483.944] /A << /S /GoTo /D <6571756174696F6E2E412E352E322E3231> >> >>}
\MusicPDFWriteObject{5395}{<< /Type /Annot /Subtype /Link /Border [0 0 0] /H /I /C [1 0 0] /Rect [196.192 457.23 233.876 472.123] /A << /S /GoTo /D <74686D2E412E352E322E3134> >> >>}
\MusicPDFWriteObject{5396}{<< /Type /Annot /Subtype /Link /Border [0 0 0] /H /I /C [1 0 0] /Rect [195.786 328.615 228.489 340.568] /A << /S /GoTo /D <6571756174696F6E2E412E352E322E31> >> >>}
\MusicPDFWriteObject{5397}{<< /Type /Annot /Subtype /Link /Border [0 0 0] /H /I /C [1 0 0] /Rect [294.06 328.615 303.993 341.473] /A << /S /GoTo /D <48666F6F746E6F74652E3637> >> >>}
\MusicPDFWriteObject{5398}{<< /Type /Annot /Subtype /Link /Border [0 0 0] /H /I /C [0 1 0] /Rect [268.884 164.973 297.977 173.841] /A << /S /GoTo /D <636974652E48522D32> >> >>}
\MusicPDFWriteObject{5399}{<< /Type /Annot /Subtype /Link /Border [0 0 0] /H /I /C [1 0 0] /Rect [175.141 684.633 185.074 697.491] /A << /S /GoTo /D <48666F6F746E6F74652E3638> >> >>}
\MusicPDFWriteObject{5400}{<< /Type /Annot /Subtype /Link /Border [0 0 0] /H /I /C [1 0 0] /Rect [268.319 593.14 301.022 605.093] /A << /S /GoTo /D <6571756174696F6E2E412E352E322E31> >> >>}
\MusicPDFWriteObject{5401}{<< /Type /Annot /Subtype /Link /Border [0 0 0] /H /I /C [0 1 0] /Rect [247.956 576.663 277.049 585.531] /A << /S /GoTo /D <636974652E48522D32> >> >>}
\MusicPDFWriteObject{5402}{<< /Type /Annot /Subtype /Link /Border [0 0 0] /H /I /C [0 1 0] /Rect [198.669 563.513 247.26 572.381] /A << /S /GoTo /D <636974652E53745368612D494949> >> >>}
\MusicPDFWriteObject{5403}{<< /Type /Annot /Subtype /Link /Border [0 0 0] /H /I /C [0 1 0] /Rect [164.12 550.362 193.212 559.23] /A << /S /GoTo /D <636974652E48522D32> >> >>}
\MusicPDFWriteObject{5404}{<< /Type /Annot /Subtype /Link /Border [0 0 0] /H /I /C [0 1 0] /Rect [479.286 537.211 509.679 546.079] /A << /S /GoTo /D <636974652E426F752D49> >> >>}
\MusicPDFWriteObject{5405}{<< /Type /Annot /Subtype /Link /Border [0 0 0] /H /I /C [0 1 0] /Rect [409.437 505.094 438.53 513.962] /A << /S /GoTo /D <636974652E48522D32> >> >>}
\MusicPDFWriteObject{5406}{<< /Type /Annot /Subtype /Link /Border [0 0 0] /H /I /C [0 1 0] /Rect [87.604 491.943 116.697 500.811] /A << /S /GoTo /D <636974652E48522D32> >> >>}
\MusicPDFWriteObject{5407}{<< /Type /Annot /Subtype /Link /Border [0 0 0] /H /I /C [1 0 0] /Rect [296.589 437.403 306.522 449.708] /A << /S /GoTo /D <48666F6F746E6F74652E3639> >> >>}
\MusicPDFWriteObject{5408}{<< /Type /Annot /Subtype /Link /Border [0 0 0] /H /I /C [0 1 0] /Rect [187.102 413.039 217.495 421.907] /A << /S /GoTo /D <636974652E426F752D49> >> >>}
\MusicPDFWriteObject{5409}{<< /Type /Annot /Subtype /Link /Border [0 0 0] /H /I /C [0 1 0] /Rect [156.938 290.224 186.031 299.092] /A << /S /GoTo /D <636974652E48522D32> >> >>}
\MusicPDFWriteObject{5410}{<< /Type /Annot /Subtype /Link /Border [0 0 0] /H /I /C [1 0 0] /Rect [205.405 229.868 215.338 242.172] /A << /S /GoTo /D <48666F6F746E6F74652E3730> >> >>}
\MusicPDFWriteObject{5411}{<< /Type /Annot /Subtype /Link /Border [0 0 0] /H /I /C [0 1 0] /Rect [243.104 665.892 272.197 674.76] /A << /S /GoTo /D <636974652E48522D32> >> >>}
\MusicPDFWriteObject{5412}{<< /Type /Annot /Subtype /Link /Border [0 0 0] /H /I /C [1 0 0] /Rect [483.329 623.949 515.603 635.904] /A << /S /GoTo /D <74686D2E412E352E332E31> >> >>}
\MusicPDFWriteObject{5413}{<< /Type /Annot /Subtype /Link /Border [0 0 0] /H /I /C [1 0 0] /Rect [198.313 597.648 231.016 609.601] /A << /S /GoTo /D <6571756174696F6E2E412E352E332E31> >> >>}
\MusicPDFWriteObject{5414}{<< /Type /Annot /Subtype /Link /Border [0 0 0] /H /I /C [1 0 0] /Rect [302.795 597.648 335.99 610.506] /A << /S /GoTo /D <74686D2E412E352E332E31> >> >>}
\MusicPDFWriteObject{5415}{<< /Type /Annot /Subtype /Link /Border [0 0 0] /H /I /C [1 0 0] /Rect [369.506 597.648 379.439 610.506] /A << /S /GoTo /D <48666F6F746E6F74652E3731> >> >>}
\MusicPDFWriteObject{5416}{<< /Type /Annot /Subtype /Link /Border [0 0 0] /H /I /C [1 0 0] /Rect [218.867 584.497 251.57 596.45] /A << /S /GoTo /D <6571756174696F6E2E412E352E332E33> >> >>}
\MusicPDFWriteObject{5417}{<< /Type /Annot /Subtype /Link /Border [0 0 0] /H /I /C [1 0 0] /Rect [443.005 558.195 476.323 570.151] /A << /S /GoTo /D <74686D2E412E352E332E33> >> >>}
\MusicPDFWriteObject{5418}{<< /Type /Annot /Subtype /Link /Border [0 0 0] /H /I /C [1 0 0] /Rect [73.712 441.894 106.415 453.847] /A << /S /GoTo /D <6571756174696F6E2E412E352E332E34> >> >>}
\MusicPDFWriteObject{5419}{<< /Type /Annot /Subtype /Link /Border [0 0 0] /H /I /C [0 1 0] /Rect [73.77 408.684 102.863 417.552] /A << /S /GoTo /D <636974652E48522D32> >> >>}
\MusicPDFWriteObject{5420}{<< /Type /Annot /Subtype /Link /Border [0 0 0] /H /I /C [1 0 0] /Rect [205.405 346.064 215.338 358.368] /A << /S /GoTo /D <48666F6F746E6F74652E3732> >> >>}
\MusicPDFWriteObject{5421}{<< /Type /Annot /Subtype /Link /Border [0 0 0] /H /I /C [0 1 0] /Rect [428.385 257.218 457.478 266.086] /A << /S /GoTo /D <636974652E48522D32> >> >>}
\MusicPDFWriteObject{5422}{<< /Type /Annot /Subtype /Link /Border [0 0 0] /H /I /C [1 0 0] /Rect [222.561 188.973 255.264 200.928] /A << /S /GoTo /D <74686D2E412E352E332E34> >> >>}
\MusicPDFWriteObject{5423}{<< /Type /Annot /Subtype /Link /Border [0 0 0] /H /I /C [0 1 0] /Rect [184.245 687.123 201.399 695.991] /A << /S /GoTo /D <636974652E466F6C> >> >>}
\MusicPDFWriteObject{5424}{<< /Type /Annot /Subtype /Link /Border [0 0 0] /H /I /C [0 1 0] /Rect [293.199 687.123 310.353 695.991] /A << /S /GoTo /D <636974652E466F6C> >> >>}
\MusicPDFWriteObject{5425}{<< /Type /Annot /Subtype /Link /Border [0 0 0] /H /I /C [1 0 0] /Rect [409.372 485.431 419.305 498.288] /A << /S /GoTo /D <48666F6F746E6F74652E3733> >> >>}
\MusicPDFWriteObject{5426}{<< /Type /Annot /Subtype /Link /Border [0 0 0] /H /I /C [1 0 0] /Rect [332.797 376.417 364.886 388.369] /A << /S /GoTo /D <74686D2E412E352E312E31> >> >>}
\MusicPDFWriteObject{5427}{<< /Type /Annot /Subtype /Link /Border [0 0 0] /H /I /C [0 1 0] /Rect [514.177 378.907 543.27 387.774] /A << /S /GoTo /D <636974652E48522D32> >> >>}
\MusicPDFWriteObject{5428}{<< /Type /Annot /Subtype /Link /Border [0 0 0] /H /I /C [1 0 0] /Rect [346.962 349.817 356.895 363.398] /A << /S /GoTo /D <48666F6F746E6F74652E3734> >> >>}
\MusicPDFWriteObject{5429}{<< /Type /Annot /Subtype /Link /Border [0 0 0] /H /I /C [0 1 0] /Rect [487.243 352.605 513.735 361.473] /A << /S /GoTo /D <636974652E52752D32> >> >>}
\MusicPDFWriteObject{5430}{<< /Type /Annot /Subtype /Link /Border [0 0 0] /H /I /C [0 1 0] /Rect [162.632 326.304 189.124 335.171] /A << /S /GoTo /D <636974652E52752D32> >> >>}
\MusicPDFWriteObject{5431}{<< /Type /Annot /Subtype /Link /Border [0 0 0] /H /I /C [0 1 0] /Rect [425.455 286.851 454.548 295.719] /A << /S /GoTo /D <636974652E48522D32> >> >>}
\MusicPDFWriteObject{5432}{<< /Type /Annot /Subtype /Link /Border [0 0 0] /H /I /C [0 1 0] /Rect [283.765 273.701 300.919 282.568] /A << /S /GoTo /D <636974652E466F6C> >> >>}
\MusicPDFWriteObject{5433}{<< /Type /Annot /Subtype /Link /Border [0 0 0] /H /I /C [1 0 0] /Rect [171.297 199.195 204.001 211.147] /A << /S /GoTo /D <74686D2E412E352E332E33> >> >>}
\MusicPDFWriteObject{5434}{<< /Type /Annot /Subtype /Link /Border [0 0 0] /H /I /C [1 0 0] /Rect [195.611 635.055 228.315 647.008] /A << /S /GoTo /D <74686D2E412E352E332E34> >> >>}
\MusicPDFWriteObject{5435}{<< /Type /Annot /Subtype /Link /Border [0 0 0] /H /I /C [1 0 0] /Rect [311.842 635.055 349.525 647.008] /A << /S /GoTo /D <6571756174696F6E2E412E352E332E3134> >> >>}
\MusicPDFWriteObject{5436}{<< /Type /Annot /Subtype /Link /Border [0 0 0] /H /I /C [1 0 0] /Rect [377.966 635.055 410.669 647.008] /A << /S /GoTo /D <6571756174696F6E2E412E352E332E33> >> >>}
\MusicPDFWriteObject{5437}{<< /Type /Annot /Subtype /Link /Border [0 0 0] /H /I /C [1 0 0] /Rect [129.95 542.286 162.653 554.238] /A << /S /GoTo /D <6571756174696F6E2E412E352E332E39> >> >>}
\MusicPDFWriteObject{5438}{<< /Type /Annot /Subtype /Link /Border [0 0 0] /H /I /C [1 0 0] /Rect [267.665 542.286 305.349 554.238] /A << /S /GoTo /D <6571756174696F6E2E412E352E332E3130> >> >>}
\MusicPDFWriteObject{5439}{<< /Type /Annot /Subtype /Link /Border [0 0 0] /H /I /C [0 1 0] /Rect [164.173 531.625 193.265 540.493] /A << /S /GoTo /D <636974652E48522D32> >> >>}
\MusicPDFWriteObject{5440}{<< /Type /Annot /Subtype /Link /Border [0 0 0] /H /I /C [0 1 0] /Rect [526.005 531.625 543.159 540.493] /A << /S /GoTo /D <636974652E466F6C> >> >>}
\MusicPDFWriteObject{5441}{<< /Type /Annot /Subtype /Link /Border [0 0 0] /H /I /C [1 0 0] /Rect [261.11 491.504 293.813 502.294] /A << /S /GoTo /D <74686D2E412E352E332E36> >> >>}
\MusicPDFWriteObject{5442}{<< /Type /Annot /Subtype /Link /Border [0 0 0] /H /I /C [1 0 0] /Rect [184.894 285.207 216.983 297.163] /A << /S /GoTo /D <74686D2E412E352E332E36> >> >>}
\MusicPDFWriteObject{5443}{<< /Type /Annot /Subtype /Link /Border [0 0 0] /H /I /C [1 0 0] /Rect [127.117 605.67 159.82 617.622] /A << /S /GoTo /D <6571756174696F6E2E412E352E342E31> >> >>}
\MusicPDFWriteObject{5444}{<< /Type /Annot /Subtype /Link /Border [0 0 0] /H /I /C [1 0 0] /Rect [498.431 605.67 530.704 617.622] /A << /S /GoTo /D <74686D2E412E352E342E31> >> >>}
\MusicPDFWriteObject{5445}{<< /Type /Annot /Subtype /Link /Border [0 0 0] /H /I /C [1 0 0] /Rect [371.658 545.594 381.591 560.67] /A << /S /GoTo /D <48666F6F746E6F74652E3735> >> >>}
\MusicPDFWriteObject{5446}{<< /Type /Annot /Subtype /Link /Border [0 0 0] /H /I /C [0 1 0] /Rect [299.315 211.099 342.338 219.967] /A << /S /GoTo /D <636974652E426F752D4754> >> >>}
\MusicPDFWriteObject{5447}{<< /Type /Annot /Subtype /Link /Border [0 0 0] /H /I /C [1 0 0] /Rect [478.788 195.459 511.492 207.411] /A << /S /GoTo /D <6571756174696F6E2E412E352E342E35> >> >>}
\MusicPDFWriteObject{5448}{<< /Type /Annot /Subtype /Link /Border [0 0 0] /H /I /C [1 0 0] /Rect [290.903 182.308 322.992 194.261] /A << /S /GoTo /D <74686D2E412E352E312E33> >> >>}
\MusicPDFWriteObject{5449}{<< /Type /Annot /Subtype /Link /Border [0 0 0] /H /I /C [0 1 0] /Rect [253.389 158.496 270.543 167.364] /A << /S /GoTo /D <636974652E466F6C> >> >>}
\MusicPDFWriteObject{5450}{<< /Type /Annot /Subtype /Link /Border [0 0 0] /H /I /C [1 0 0] /Rect [402.009 588.816 434.713 600.769] /A << /S /GoTo /D <6571756174696F6E2E412E352E342E38> >> >>}
\MusicPDFWriteObject{5451}{<< /Type /Annot /Subtype /Link /Border [0 0 0] /H /I /C [1 0 0] /Rect [73.712 575.665 106.415 587.618] /A << /S /GoTo /D <6571756174696F6E2E412E352E342E37> >> >>}
\MusicPDFWriteObject{5452}{<< /Type /Annot /Subtype /Link /Border [0 0 0] /H /I /C [1 0 0] /Rect [149.025 515.426 158.958 525.794] /A << /S /GoTo /D <48666F6F746E6F74652E3736> >> >>}
\MusicPDFWriteObject{5453}{<< /Type /Annot /Subtype /Link /Border [0 0 0] /H /I /C [1 0 0] /Rect [206.684 413.779 239.387 425.732] /A << /S /GoTo /D <74686D2E412E352E332E37> >> >>}
\MusicPDFWriteObject{5454}{<< /Type /Annot /Subtype /Link /Border [0 0 0] /H /I /C [1 0 0] /Rect [117.789 371.319 155.472 383.272] /A << /S /GoTo /D <6571756174696F6E2E412E352E342E3131> >> >>}
\MusicPDFWriteObject{5455}{<< /Type /Annot /Subtype /Link /Border [0 0 0] /H /I /C [0 1 0] /Rect [307.995 189.283 325.149 198.15] /A << /S /GoTo /D <636974652E466F6C> >> >>}
\MusicPDFWriteObject{5456}{<< /Type /Annot /Subtype /Link /Border [0 0 0] /H /I /C [1 0 0] /Rect [443.974 186.793 453.907 199.65] /A << /S /GoTo /D <48666F6F746E6F74652E3737> >> >>}
\MusicPDFWriteObject{5457}{<< /Type /Annot /Subtype /Link /Border [0 0 0] /H /I /C [1 0 0] /Rect [177.472 592.578 210.175 604.531] /A << /S /GoTo /D <6571756174696F6E2E412E352E342E32> >> >>}
\MusicPDFWriteObject{5458}{<< /Type /Annot /Subtype /Link /Border [0 0 0] /H /I /C [1 0 0] /Rect [446.816 592.578 456.749 605.435] /A << /S /GoTo /D <48666F6F746E6F74652E3738> >> >>}
\MusicPDFWriteObject{5459}{<< /Type /Annot /Subtype /Link /Border [0 0 0] /H /I /C [0 1 0] /Rect [321.028 521.549 338.183 530.417] /A << /S /GoTo /D <636974652E466F6C> >> >>}
\MusicPDFWriteObject{5460}{<< /Type /Annot /Subtype /Link /Border [0 0 0] /H /I /C [1 0 0] /Rect [307.295 492.757 317.228 505.615] /A << /S /GoTo /D <48666F6F746E6F74652E3739> >> >>}
\MusicPDFWriteObject{5461}{<< /Type /Annot /Subtype /Link /Border [0 0 0] /H /I /C [0 1 0] /Rect [334.99 495.248 364.082 504.116] /A << /S /GoTo /D <636974652E48522D31> >> >>}
\MusicPDFWriteObject{5462}{<< /Type /Annot /Subtype /Link /Border [0 0 0] /H /I /C [1 0 0] /Rect [308.31 384.509 341.013 395.3] /A << /S /GoTo /D <74686D2E412E352E342E31> >> >>}
\MusicPDFWriteObject{5463}{<< /Type /Annot /Subtype /Link /Border [0 0 0] /H /I /C [1 0 0] /Rect [385.799 645.181 423.483 657.134] /A << /S /GoTo /D <6571756174696F6E2E412E352E342E3137> >> >>}
\MusicPDFWriteObject{5464}{<< /Type /Annot /Subtype /Link /Border [0 0 0] /H /I /C [1 0 0] /Rect [268.796 595.603 306.479 607.555] /A << /S /GoTo /D <6571756174696F6E2E412E352E342E3137> >> >>}
\MusicPDFWriteObject{5465}{<< /Type /Annot /Subtype /Link /Border [0 0 0] /H /I /C [1 0 0] /Rect [211.391 532.16 217.354 544.464] /A << /S /GoTo /D <48666F6F746E6F74652E3830> >> >>}
\MusicPDFWriteObject{5466}{<< /Type /Annot /Subtype /Link /Border [0 0 0] /H /I /C [1 0 0] /Rect [82.959 465.493 239.084 477.808] /A << /S /GoTo /D <73656374696F6E2E422E31> >> >>}
\MusicPDFWriteObject{5467}{<< /Type /Annot /Subtype /Link /Border [0 0 0] /H /I /C [1 0 0] /Rect [82.959 452.342 324.036 464.657] /A << /S /GoTo /D <73756273656374696F6E2E422E312E31> >> >>}
\MusicPDFWriteObject{5468}{<< /Type /Annot /Subtype /Link /Border [0 0 0] /H /I /C [1 0 0] /Rect [82.959 439.191 380.881 451.506] /A << /S /GoTo /D <73756273756273656374696F6E2E422E312E312E31> >> >>}
\MusicPDFWriteObject{5469}{<< /Type /Annot /Subtype /Link /Border [0 0 0] /H /I /C [1 0 0] /Rect [82.959 426.954 274.859 437.745] /A << /S /GoTo /D <73656374696F6E2E422E32> >> >>}
\MusicPDFWriteObject{5470}{<< /Type /Annot /Subtype /Link /Border [0 0 0] /H /I /C [1 0 0] /Rect [82.959 413.803 289.163 424.594] /A << /S /GoTo /D <73756273656374696F6E2E422E322E31> >> >>}
\MusicPDFWriteObject{5471}{<< /Type /Annot /Subtype /Link /Border [0 0 0] /H /I /C [1 0 0] /Rect [82.959 400.653 411.148 411.443] /A << /S /GoTo /D <73756273756273656374696F6E2E422E322E312E31> >> >>}
\MusicPDFWriteObject{5472}{<< /Type /Annot /Subtype /Link /Border [0 0 0] /H /I /C [1 0 0] /Rect [82.959 387.502 326.459 398.293] /A << /S /GoTo /D <73756273756273656374696F6E2E422E322E312E32> >> >>}
\MusicPDFWriteObject{5473}{<< /Type /Annot /Subtype /Link /Border [0 0 0] /H /I /C [1 0 0] /Rect [82.959 374.351 282.606 385.142] /A << /S /GoTo /D <73756273656374696F6E2E422E322E32> >> >>}
\MusicPDFWriteObject{5474}{<< /Type /Annot /Subtype /Link /Border [0 0 0] /H /I /C [1 0 0] /Rect [82.959 361.2 204.142 371.991] /A << /S /GoTo /D <73756273756273656374696F6E2E422E322E322E31> >> >>}
\MusicPDFWriteObject{5475}{<< /Type /Annot /Subtype /Link /Border [0 0 0] /H /I /C [1 0 0] /Rect [82.959 348.05 242.24 358.84] /A << /S /GoTo /D <73656374696F6E2E422E33> >> >>}
\MusicPDFWriteObject{5476}{<< /Type /Annot /Subtype /Link /Border [0 0 0] /H /I /C [1 0 0] /Rect [82.959 334.899 238.477 345.69] /A << /S /GoTo /D <73756273656374696F6E2E422E332E31> >> >>}
\MusicPDFWriteObject{5477}{<< /Type /Annot /Subtype /Link /Border [0 0 0] /H /I /C [1 0 0] /Rect [82.959 321.748 284.017 332.539] /A << /S /GoTo /D <73756273756273656374696F6E2E422E332E312E31> >> >>}
\MusicPDFWriteObject{5478}{<< /Type /Annot /Subtype /Link /Border [0 0 0] /H /I /C [1 0 0] /Rect [82.959 308.597 270.46 319.388] /A << /S /GoTo /D <73756273656374696F6E2E422E332E32> >> >>}
\MusicPDFWriteObject{5479}{<< /Type /Annot /Subtype /Link /Border [0 0 0] /H /I /C [1 0 0] /Rect [82.959 295.447 281.389 306.237] /A << /S /GoTo /D <73756273756273656374696F6E2E422E332E322E31> >> >>}
\MusicPDFWriteObject{5480}{<< /Type /Annot /Subtype /Link /Border [0 0 0] /H /I /C [1 0 0] /Rect [82.959 282.296 286.813 293.087] /A << /S /GoTo /D <73656374696F6E2E422E34> >> >>}
\MusicPDFWriteObject{5481}{<< /Type /Annot /Subtype /Link /Border [0 0 0] /H /I /C [1 0 0] /Rect [82.959 269.145 378.197 279.936] /A << /S /GoTo /D <73756273656374696F6E2E422E342E31> >> >>}
\MusicPDFWriteObject{5482}{<< /Type /Annot /Subtype /Link /Border [0 0 0] /H /I /C [1 0 0] /Rect [82.959 255.994 373.659 266.785] /A << /S /GoTo /D <73756273756273656374696F6E2E422E342E312E31> >> >>}
\MusicPDFWriteObject{5483}{<< /Type /Annot /Subtype /Link /Border [0 0 0] /H /I /C [1 0 0] /Rect [82.959 242.29 427.905 254.263] /A << /S /GoTo /D <73756273756273656374696F6E2E422E342E312E32> >> >>}
\MusicPDFWriteObject{5484}{<< /Type /Annot /Subtype /Link /Border [0 0 0] /H /I /C [1 0 0] /Rect [82.959 229.693 368.651 240.483] /A << /S /GoTo /D <73756273656374696F6E2E422E342E32> >> >>}
\MusicPDFWriteObject{5485}{<< /Type /Annot /Subtype /Link /Border [0 0 0] /H /I /C [1 0 0] /Rect [82.959 215.988 487.96 227.943] /A << /S /GoTo /D <73756273756273656374696F6E2E422E342E322E31> >> >>}
\MusicPDFWriteObject{5486}{<< /Type /Annot /Subtype /Link /Border [0 0 0] /H /I /C [1 0 0] /Rect [82.959 205.328 262.271 214.182] /A << /S /GoTo /D <73756273756273656374696F6E2E422E342E322E32> >> >>}
\MusicPDFWriteObject{5487}{<< /Type /Annot /Subtype /Link /Border [0 0 0] /H /I /C [1 0 0] /Rect [82.959 192.177 247.054 201.031] /A << /S /GoTo /D <73756273756273656374696F6E2E422E342E322E33> >> >>}
\MusicPDFWriteObject{5488}{<< /Type /Annot /Subtype /Link /Border [0 0 0] /H /I /C [1 0 0] /Rect [82.959 177.09 302.307 187.88] /A << /S /GoTo /D <73756273756273656374696F6E2E422E342E322E34> >> >>}
\MusicPDFWriteObject{5489}{<< /Type /Annot /Subtype /Link /Border [0 0 0] /H /I /C [1 0 0] /Rect [82.959 163.385 242.165 175.34] /A << /S /GoTo /D <73756273656374696F6E2E422E342E33> >> >>}
\MusicPDFWriteObject{5490}{<< /Type /Annot /Subtype /Link /Border [0 0 0] /H /I /C [1 0 0] /Rect [82.959 150.234 393.228 162.189] /A << /S /GoTo /D <73756273756273656374696F6E2E422E342E332E31> >> >>}
\MusicPDFWriteObject{5491}{<< /Type /Annot /Subtype /Link /Border [0 0 0] /H /I /C [1 0 0] /Rect [82.959 137.637 224.339 148.428] /A << /S /GoTo /D <73756273656374696F6E2E422E342E34> >> >>}
\MusicPDFWriteObject{5492}{<< /Type /Annot /Subtype /Link /Border [0 0 0] /H /I /C [1 0 0] /Rect [82.959 124.487 507.984 135.277] /A << /S /GoTo /D <73756273756273656374696F6E2E422E342E342E31> >> >>}
\MusicPDFWriteObject{5493}{<< /Type /Annot /Subtype /Link /Border [0 0 0] /H /I /C [1 0 0] /Rect [82.959 111.336 279.922 122.127] /A << /S /GoTo /D <73756273756273656374696F6E2E422E342E342E32> >> >>}
\MusicPDFWriteObject{5494}{<< /Type /Annot /Subtype /Link /Border [0 0 0] /H /I /C [1 0 0] /Rect [82.959 685.187 295.637 695.977] /A << /S /GoTo /D <73756273756273656374696F6E2E422E342E342E33> >> >>}
\MusicPDFWriteObject{5495}{<< /Type /Annot /Subtype /Link /Border [0 0 0] /H /I /C [1 0 0] /Rect [82.959 672.036 311.435 682.827] /A << /S /GoTo /D <73756273756273656374696F6E2E422E342E342E34> >> >>}
\MusicPDFWriteObject{5496}{<< /Type /Annot /Subtype /Link /Border [0 0 0] /H /I /C [1 0 0] /Rect [82.959 658.885 252.477 669.676] /A << /S /GoTo /D <73656374696F6E2E422E35> >> >>}
\MusicPDFWriteObject{5497}{<< /Type /Annot /Subtype /Link /Border [0 0 0] /H /I /C [1 0 0] /Rect [82.959 645.734 175.811 656.525] /A << /S /GoTo /D <73756273656374696F6E2E422E352E31> >> >>}
\MusicPDFWriteObject{5498}{<< /Type /Annot /Subtype /Link /Border [0 0 0] /H /I /C [1 0 0] /Rect [82.959 632.584 189.451 643.374] /A << /S /GoTo /D <73756273656374696F6E2E422E352E32> >> >>}
\MusicPDFWriteObject{5499}{<< /Type /Annot /Subtype /Link /Border [0 0 0] /H /I /C [1 0 0] /Rect [82.959 619.433 298.045 630.223] /A << /S /GoTo /D <73756273756273656374696F6E2E422E352E322E31> >> >>}
\MusicPDFWriteObject{5500}{<< /Type /Annot /Subtype /Link /Border [0 0 0] /H /I /C [1 0 0] /Rect [82.959 606.282 347.486 617.073] /A << /S /GoTo /D <73756273656374696F6E2E422E352E33> >> >>}
\MusicPDFWriteObject{5501}{<< /Type /Annot /Subtype /Link /Border [0 0 0] /H /I /C [1 0 0] /Rect [82.959 593.131 205 603.922] /A << /S /GoTo /D <73756273756273656374696F6E2E422E352E332E31> >> >>}
\MusicPDFWriteObject{5502}{<< /Type /Annot /Subtype /Link /Border [0 0 0] /H /I /C [1 0 0] /Rect [82.959 579.981 491.124 590.771] /A << /S /GoTo /D <73756273756273656374696F6E2E422E352E332E32> >> >>}
\MusicPDFWriteObject{5503}{<< /Type /Annot /Subtype /Link /Border [0 0 0] /H /I /C [1 0 0] /Rect [82.959 566.83 245.615 577.62] /A << /S /GoTo /D <73756273756273656374696F6E2E422E352E332E33> >> >>}
\MusicPDFWriteObject{5504}{<< /Type /Annot /Subtype /Link /Border [0 0 0] /H /I /C [1 0 0] /Rect [82.959 553.679 355.537 564.47] /A << /S /GoTo /D <73756273756273656374696F6E2E422E352E332E34> >> >>}
\MusicPDFWriteObject{5505}{<< /Type /Annot /Subtype /Link /Border [0 0 0] /H /I /C [1 0 0] /Rect [82.959 540.528 321.287 551.319] /A << /S /GoTo /D <73656374696F6E2E422E36> >> >>}
\MusicPDFWriteObject{5506}{<< /Type /Annot /Subtype /Link /Border [0 0 0] /H /I /C [1 0 0] /Rect [82.959 525.137 445.012 538.796] /A << /S /GoTo /D <73756273656374696F6E2E422E362E31> >> >>}
\MusicPDFWriteObject{5507}{<< /Type /Annot /Subtype /Link /Border [0 0 0] /H /I /C [1 0 0] /Rect [82.959 513.313 292.15 525.628] /A << /S /GoTo /D <73756273756273656374696F6E2E422E362E312E31> >> >>}
\MusicPDFWriteObject{5508}{<< /Type /Annot /Subtype /Link /Border [0 0 0] /H /I /C [1 0 0] /Rect [82.959 498.835 267.193 512.495] /A << /S /GoTo /D <73756273756273656374696F6E2E422E362E312E32> >> >>}
\MusicPDFWriteObject{5509}{<< /Type /Annot /Subtype /Link /Border [0 0 0] /H /I /C [1 0 0] /Rect [82.959 487.925 445.957 498.716] /A << /S /GoTo /D <73756273756273656374696F6E2E422E362E312E33> >> >>}
\MusicPDFWriteObject{5510}{<< /Type /Annot /Subtype /Link /Border [0 0 0] /H /I /C [1 0 0] /Rect [82.959 474.774 340.819 485.565] /A << /S /GoTo /D <73656374696F6E2E422E37> >> >>}
\MusicPDFWriteObject{5511}{<< /Type /Annot /Subtype /Link /Border [0 0 0] /H /I /C [1 0 0] /Rect [82.959 459.383 381.681 473.043] /A << /S /GoTo /D <73756273656374696F6E2E422E372E31> >> >>}
\MusicPDFWriteObject{5512}{<< /Type /Annot /Subtype /Link /Border [0 0 0] /H /I /C [1 0 0] /Rect [82.959 448.473 351.693 459.264] /A << /S /GoTo /D <73756273756273656374696F6E2E422E372E312E31> >> >>}
\MusicPDFWriteObject{5513}{<< /Type /Annot /Subtype /Link /Border [0 0 0] /H /I /C [1 0 0] /Rect [82.959 437.259 336.087 446.113] /A << /S /GoTo /D <73756273756273656374696F6E2E422E372E312E32> >> >>}
\MusicPDFWriteObject{5514}{<< /Type /Annot /Subtype /Link /Border [0 0 0] /H /I /C [1 0 0] /Rect [82.959 419.931 463.658 433.59] /A << /S /GoTo /D <73756273756273656374696F6E2E422E372E312E33> >> >>}
\MusicPDFWriteObject{5515}{<< /Type /Annot /Subtype /Link /Border [0 0 0] /H /I /C [1 0 0] /Rect [82.959 409.021 245.81 419.811] /A << /S /GoTo /D <73756273656374696F6E2E422E372E32> >> >>}
\MusicPDFWriteObject{5516}{<< /Type /Annot /Subtype /Link /Border [0 0 0] /H /I /C [1 0 0] /Rect [82.959 395.87 434.169 406.661] /A << /S /GoTo /D <73756273756273656374696F6E2E422E372E322E31> >> >>}
\MusicPDFWriteObject{5517}{<< /Type /Annot /Subtype /Link /Border [0 0 0] /H /I /C [1 0 0] /Rect [82.959 382.719 333.984 393.51] /A << /S /GoTo /D <73756273756273656374696F6E2E422E372E322E32> >> >>}
\MusicPDFWriteObject{5518}{<< /Type /Annot /Subtype /Link /Border [0 0 0] /H /I /C [1 0 0] /Rect [82.959 369.568 283.88 380.359] /A << /S /GoTo /D <73756273756273656374696F6E2E422E372E322E33> >> >>}
\MusicPDFWriteObject{5519}{<< /Type /Annot /Subtype /Link /Border [0 0 0] /H /I /C [1 0 0] /Rect [82.959 356.418 467.037 367.208] /A << /S /GoTo /D <73756273756273656374696F6E2E422E372E322E34> >> >>}
\MusicPDFWriteObject{5520}{<< /Type /Annot /Subtype /Link /Border [0 0 0] /H /I /C [0 1 0] /Rect [116.423 611.216 142.244 620.084] /A << /S /GoTo /D <636974652E4D5665> >> >>}
\MusicPDFWriteObject{5521}{<< /Type /Annot /Subtype /Link /Border [0 0 0] /H /I /C [0 1 0] /Rect [366.711 598.065 382.399 606.933] /A << /S /GoTo /D <636974652E5432> >> >>}
\MusicPDFWriteObject{5522}{<< /Type /Annot /Subtype /Link /Border [0 0 0] /H /I /C [0 1 0] /Rect [158.574 584.914 174.262 593.782] /A << /S /GoTo /D <636974652E5431> >> >>}
\MusicPDFWriteObject{5523}{<< /Type /Annot /Subtype /Link /Border [0 0 0] /H /I /C [0 1 0] /Rect [196.983 558.613 212.671 567.481] /A << /S /GoTo /D <636974652E5431> >> >>}
\MusicPDFWriteObject{5524}{<< /Type /Annot /Subtype /Link /Border [0 0 0] /H /I /C [0 1 0] /Rect [158.242 450.875 173.93 459.743] /A << /S /GoTo /D <636974652E5431> >> >>}
\MusicPDFWriteObject{5525}{<< /Type /Annot /Subtype /Link /Border [0 0 0] /H /I /C [0 1 0] /Rect [320.181 437.724 335.869 446.592] /A << /S /GoTo /D <636974652E5431> >> >>}
\MusicPDFWriteObject{5526}{<< /Type /Annot /Subtype /Link /Border [0 0 0] /H /I /C [0 1 0] /Rect [89.469 650.696 105.157 659.564] /A << /S /GoTo /D <636974652E5431> >> >>}
\MusicPDFWriteObject{5527}{<< /Type /Annot /Subtype /Link /Border [0 0 0] /H /I /C [1 0 0] /Rect [307.645 648.205 340.055 661.534] /A << /S /GoTo /D <6571756174696F6E2E422E312E312E31> >> >>}
\MusicPDFWriteObject{5528}{<< /Type /Annot /Subtype /Link /Border [0 0 0] /H /I /C [1 0 0] /Rect [251.468 465.108 283.757 477.06] /A << /S /GoTo /D <6571756174696F6E2E422E312E312E36> >> >>}
\MusicPDFWriteObject{5529}{<< /Type /Annot /Subtype /Link /Border [0 0 0] /H /I /C [1 0 0] /Rect [348.962 465.108 381.251 477.06] /A << /S /GoTo /D <6571756174696F6E2E422E312E312E33> >> >>}
\MusicPDFWriteObject{5530}{<< /Type /Annot /Subtype /Link /Border [0 0 0] /H /I /C [1 0 0] /Rect [73.712 451.957 106 463.91] /A << /S /GoTo /D <6571756174696F6E2E422E312E312E33> >> >>}
\MusicPDFWriteObject{5531}{<< /Type /Annot /Subtype /Link /Border [0 0 0] /H /I /C [0 1 0] /Rect [127.801 454.447 143.489 463.315] /A << /S /GoTo /D <636974652E5431> >> >>}
\MusicPDFWriteObject{5532}{<< /Type /Annot /Subtype /Link /Border [0 0 0] /H /I /C [0 1 0] /Rect [199.799 435.319 215.113 444.186] /A << /S /GoTo /D <636974652E494B> >> >>}
\MusicPDFWriteObject{5533}{<< /Type /Annot /Subtype /Link /Border [0 0 0] /H /I /C [1 0 0] /Rect [173.559 384.652 206.453 395.443] /A << /S /GoTo /D <74686D2E422E312E312E32> >> >>}
\MusicPDFWriteObject{5534}{<< /Type /Annot /Subtype /Link /Border [0 0 0] /H /I /C [1 0 0] /Rect [422.017 357.797 454.305 369.752] /A << /S /GoTo /D <74686D2E422E312E312E31> >> >>}
\MusicPDFWriteObject{5535}{<< /Type /Annot /Subtype /Link /Border [0 0 0] /H /I /C [1 0 0] /Rect [278.613 344.647 308.826 356.599] /A << /S /GoTo /D <6571756174696F6E2E332E322E312E33> >> >>}
\MusicPDFWriteObject{5536}{<< /Type /Annot /Subtype /Link /Border [0 0 0] /H /I /C [1 0 0] /Rect [190.204 226.179 220.7 238.494] /A << /S /GoTo /D <74686D2E332E322E312E33> >> >>}
\MusicPDFWriteObject{5537}{<< /Type /Annot /Subtype /Link /Border [0 0 0] /H /I /C [1 0 0] /Rect [245.477 226.179 275.972 238.494] /A << /S /GoTo /D <6571756174696F6E2E332E322E312E38> >> >>}
\MusicPDFWriteObject{5538}{<< /Type /Annot /Subtype /Link /Border [0 0 0] /H /I /C [1 0 0] /Rect [285.157 226.179 320.682 238.494] /A << /S /GoTo /D <6571756174696F6E2E332E322E312E3130> >> >>}
\MusicPDFWriteObject{5539}{<< /Type /Annot /Subtype /Link /Border [0 0 0] /H /I /C [0 1 0] /Rect [341.426 129.055 363.996 137.923] /A << /S /GoTo /D <636974652E4E6575> >> >>}
\MusicPDFWriteObject{5540}{<< /Type /Annot /Subtype /Link /Border [0 0 0] /H /I /C [0 1 0] /Rect [234.93 687.123 257.5 695.991] /A << /S /GoTo /D <636974652E4E6575> >> >>}
\MusicPDFWriteObject{5541}{<< /Type /Annot /Subtype /Link /Border [0 0 0] /H /I /C [1 0 0] /Rect [428.601 684.633 465.87 696.586] /A << /S /GoTo /D <6571756174696F6E2E422E312E312E3130> >> >>}
\MusicPDFWriteObject{5542}{<< /Type /Annot /Subtype /Link /Border [0 0 0] /H /I /C [1 0 0] /Rect [373.645 638.059 379.607 650.917] /A << /S /GoTo /D <48666F6F746E6F74652E3831> >> >>}
\MusicPDFWriteObject{5543}{<< /Type /Annot /Subtype /Link /Border [0 0 0] /H /I /C [1 0 0] /Rect [444.418 567.357 474.631 579.31] /A << /S /GoTo /D <6571756174696F6E2E332E322E312E37> >> >>}
\MusicPDFWriteObject{5544}{<< /Type /Annot /Subtype /Link /Border [0 0 0] /H /I /C [1 0 0] /Rect [505.68 554.206 537.968 566.159] /A << /S /GoTo /D <6571756174696F6E2E422E312E312E39> >> >>}
\MusicPDFWriteObject{5545}{<< /Type /Annot /Subtype /Link /Border [0 0 0] /H /I /C [1 0 0] /Rect [237.81 541.055 275.078 553.008] /A << /S /GoTo /D <6571756174696F6E2E422E312E312E3130> >> >>}
\MusicPDFWriteObject{5546}{<< /Type /Annot /Subtype /Link /Border [0 0 0] /H /I /C [1 0 0] /Rect [191.466 420.914 223.754 432.867] /A << /S /GoTo /D <74686D2E422E312E312E31> >> >>}
\MusicPDFWriteObject{5547}{<< /Type /Annot /Subtype /Link /Border [0 0 0] /H /I /C [1 0 0] /Rect [358.71 376.087 391.604 386.878] /A << /S /GoTo /D <74686D2E422E312E312E31> >> >>}
\MusicPDFWriteObject{5548}{<< /Type /Annot /Subtype /Link /Border [0 0 0] /H /I /C [1 0 0] /Rect [73.712 362.383 110.98 374.336] /A << /S /GoTo /D <6571756174696F6E2E422E312E312E3132> >> >>}
\MusicPDFWriteObject{5549}{<< /Type /Annot /Subtype /Link /Border [0 0 0] /H /I /C [1 0 0] /Rect [98.619 300.127 130.908 312.08] /A << /S /GoTo /D <6571756174696F6E2E422E312E312E39> >> >>}
\MusicPDFWriteObject{5550}{<< /Type /Annot /Subtype /Link /Border [0 0 0] /H /I /C [1 0 0] /Rect [116.408 289.466 122.37 299.834] /A << /S /GoTo /D <48666F6F746E6F74652E3832> >> >>}
\MusicPDFWriteObject{5551}{<< /Type /Annot /Subtype /Link /Border [0 0 0] /H /I /C [1 0 0] /Rect [428.203 209.443 465.471 221.395] /A << /S /GoTo /D <6571756174696F6E2E422E312E312E3135> >> >>}
\MusicPDFWriteObject{5552}{<< /Type /Annot /Subtype /Link /Border [0 0 0] /H /I /C [1 0 0] /Rect [156.877 183.695 162.84 195.999] /A << /S /GoTo /D <48666F6F746E6F74652E3833> >> >>}
\MusicPDFWriteObject{5553}{<< /Type /Annot /Subtype /Link /Border [0 0 0] /H /I /C [1 0 0] /Rect [430.104 183.695 459.753 195.999] /A << /S /GoTo /D <74686D2E332E332E312E35> >> >>}
\MusicPDFWriteObject{5554}{<< /Type /Annot /Subtype /Link /Border [0 0 0] /H /I /C [1 0 0] /Rect [242.237 103.281 269.977 115.865] /A << /S /GoTo /D <74686D2E422E362E312E37> >> >>}
\MusicPDFWriteObject{5555}{<< /Type /Annot /Subtype /Link /Border [0 0 0] /H /I /C [1 0 0] /Rect [538.832 589.082 544.795 602.991] /A << /S /GoTo /D <48666F6F746E6F74652E3834> >> >>}
\MusicPDFWriteObject{5556}{<< /Type /Annot /Subtype /Link /Border [0 0 0] /H /I /C [0 1 0] /Rect [97.315 579.473 117.499 588.341] /A << /S /GoTo /D <636974652E4E696B> >> >>}
\MusicPDFWriteObject{5557}{<< /Type /Annot /Subtype /Link /Border [0 0 0] /H /I /C [0 1 0] /Rect [117.876 566.322 139.734 575.19] /A << /S /GoTo /D <636974652E447572> >> >>}
\MusicPDFWriteObject{5558}{<< /Type /Annot /Subtype /Link /Border [0 0 0] /H /I /C [0 1 0] /Rect [106.564 513.456 128.421 522.324] /A << /S /GoTo /D <636974652E447572> >> >>}
\MusicPDFWriteObject{5559}{<< /Type /Annot /Subtype /Link /Border [0 0 0] /H /I /C [1 0 0] /Rect [512.603 468.983 518.565 481.84] /A << /S /GoTo /D <48666F6F746E6F74652E3835> >> >>}
\MusicPDFWriteObject{5560}{<< /Type /Annot /Subtype /Link /Border [0 0 0] /H /I /C [0 1 0] /Rect [522.975 471.473 543.159 480.341] /A << /S /GoTo /D <636974652E4E696B> >> >>}
\MusicPDFWriteObject{5561}{<< /Type /Annot /Subtype /Link /Border [0 0 0] /H /I /C [0 1 0] /Rect [168.778 341.86 190.636 350.727] /A << /S /GoTo /D <636974652E447572> >> >>}
\MusicPDFWriteObject{5562}{<< /Type /Annot /Subtype /Link /Border [0 0 0] /H /I /C [1 0 0] /Rect [95.415 300.385 101.377 312.775] /A << /S /GoTo /D <48666F6F746E6F74652E3836> >> >>}
\MusicPDFWriteObject{5563}{<< /Type /Annot /Subtype /Link /Border [0 0 0] /H /I /C [1 0 0] /Rect [147.448 363.499 153.411 375.804] /A << /S /GoTo /D <48666F6F746E6F74652E3837> >> >>}
\MusicPDFWriteObject{5564}{<< /Type /Annot /Subtype /Link /Border [0 0 0] /H /I /C [1 0 0] /Rect [200.463 169.476 232.752 181.429] /A << /S /GoTo /D <6571756174696F6E2E422E322E312E36> >> >>}
\MusicPDFWriteObject{5565}{<< /Type /Annot /Subtype /Link /Border [0 0 0] /H /I /C [1 0 0] /Rect [203.197 151.232 238.929 160.086] /A << /S /GoTo /D <74686D2E422E322E312E35> >> >>}
\MusicPDFWriteObject{5566}{<< /Type /Annot /Subtype /Link /Border [0 0 0] /H /I /C [1 0 0] /Rect [171.741 684.633 204.029 696.586] /A << /S /GoTo /D <6571756174696F6E2E422E322E312E36> >> >>}
\MusicPDFWriteObject{5567}{<< /Type /Annot /Subtype /Link /Border [0 0 0] /H /I /C [1 0 0] /Rect [268.328 684.633 300.616 696.586] /A << /S /GoTo /D <74686D2E422E322E312E36> >> >>}
\MusicPDFWriteObject{5568}{<< /Type /Annot /Subtype /Link /Border [0 0 0] /H /I /C [1 0 0] /Rect [206.678 588.207 238.966 600.16] /A << /S /GoTo /D <6571756174696F6E2E422E322E312E37> >> >>}
\MusicPDFWriteObject{5569}{<< /Type /Annot /Subtype /Link /Border [0 0 0] /H /I /C [1 0 0] /Rect [293.866 369.157 331.134 381.11] /A << /S /GoTo /D <6571756174696F6E2E422E322E312E3131> >> >>}
\MusicPDFWriteObject{5570}{<< /Type /Annot /Subtype /Link /Border [0 0 0] /H /I /C [1 0 0] /Rect [418.206 369.157 455.475 381.11] /A << /S /GoTo /D <6571756174696F6E2E422E322E312E3131> >> >>}
\MusicPDFWriteObject{5571}{<< /Type /Annot /Subtype /Link /Border [0 0 0] /H /I /C [1 0 0] /Rect [485.429 369.157 522.697 381.11] /A << /S /GoTo /D <6571756174696F6E2E422E322E312E3130> >> >>}
\MusicPDFWriteObject{5572}{<< /Type /Annot /Subtype /Link /Border [0 0 0] /H /I /C [1 0 0] /Rect [227.207 311.609 264.475 323.561] /A << /S /GoTo /D <6571756174696F6E2E422E322E312E3132> >> >>}
\MusicPDFWriteObject{5573}{<< /Type /Annot /Subtype /Link /Border [0 0 0] /H /I /C [1 0 0] /Rect [429.316 211.995 435.279 224.853] /A << /S /GoTo /D <48666F6F746E6F74652E3838> >> >>}
\MusicPDFWriteObject{5574}{<< /Type /Annot /Subtype /Link /Border [0 0 0] /H /I /C [1 0 0] /Rect [174.56 112.851 207.09 125.858] /A << /S /GoTo /D <6571756174696F6E2E422E322E312E32> >> >>}
\MusicPDFWriteObject{5575}{<< /Type /Annot /Subtype /Link /Border [0 0 0] /H /I /C [0 1 0] /Rect [229.563 82.435 245.019 89.928] /A << /S /GoTo /D <636974652E484C> >> >>}
\MusicPDFWriteObject{5576}{<< /Type /Annot /Subtype /Link /Border [0 0 0] /H /I /C [1 0 0] /Rect [372.012 344.161 379.085 356.114] /A << /S /GoTo /D <6669677572652E422E31> >> >>}
\MusicPDFWriteObject{5577}{<< /Type /Annot /Subtype /Link /Border [0 0 0] /H /I /C [1 0 0] /Rect [120.17 125.044 157.438 136.997] /A << /S /GoTo /D <6571756174696F6E2E422E322E312E3134> >> >>}
\MusicPDFWriteObject{5578}{<< /Type /Annot /Subtype /Link /Border [0 0 0] /H /I /C [1 0 0] /Rect [321.997 684.633 359.265 696.586] /A << /S /GoTo /D <6571756174696F6E2E422E322E312E3135> >> >>}
\MusicPDFWriteObject{5579}{<< /Type /Annot /Subtype /Link /Border [0 0 0] /H /I /C [1 0 0] /Rect [486.21 552.815 523.478 564.768] /A << /S /GoTo /D <6571756174696F6E2E422E322E312E3138> >> >>}
\MusicPDFWriteObject{5580}{<< /Type /Annot /Subtype /Link /Border [0 0 0] /H /I /C [1 0 0] /Rect [387.297 414.318 419.585 425.109] /A << /S /GoTo /D <74686D2E422E322E312E37> >> >>}
\MusicPDFWriteObject{5581}{<< /Type /Annot /Subtype /Link /Border [0 0 0] /H /I /C [0 1 0] /Rect [471.916 250.117 489.071 258.985] /A << /S /GoTo /D <636974652E466F6C> >> >>}
\MusicPDFWriteObject{5582}{<< /Type /Annot /Subtype /Link /Border [0 0 0] /H /I /C [1 0 0] /Rect [71.004 672.035 103.722 682.884] /A << /S /GoTo /D <74686D2E442E332E312E31> >> >>}
\MusicPDFWriteObject{5583}{<< /Type /Annot /Subtype /Link /Border [0 0 0] /H /I /C [1 0 0] /Rect [196.76 386.859 229.602 397.707] /A << /S /GoTo /D <74686D2E442E332E312E31> >> >>}
\MusicPDFWriteObject{5584}{<< /Type /Annot /Subtype /Link /Border [0 0 0] /H /I /C [1 0 0] /Rect [249.733 339.161 287.001 351.113] /A << /S /GoTo /D <6571756174696F6E2E422E322E312E3231> >> >>}
\MusicPDFWriteObject{5585}{<< /Type /Annot /Subtype /Link /Border [0 0 0] /H /I /C [1 0 0] /Rect [118.122 326.01 155.391 337.963] /A << /S /GoTo /D <6571756174696F6E2E422E322E312E3231> >> >>}
\MusicPDFWriteObject{5586}{<< /Type /Annot /Subtype /Link /Border [0 0 0] /H /I /C [1 0 0] /Rect [293.248 629.899 330.516 641.852] /A << /S /GoTo /D <6571756174696F6E2E422E322E312E3236> >> >>}
\MusicPDFWriteObject{5587}{<< /Type /Annot /Subtype /Link /Border [0 0 0] /H /I /C [1 0 0] /Rect [370.266 486.86 407.534 498.813] /A << /S /GoTo /D <6571756174696F6E2E422E322E312E3330> >> >>}
\MusicPDFWriteObject{5588}{<< /Type /Annot /Subtype /Link /Border [0 0 0] /H /I /C [1 0 0] /Rect [504.62 473.709 541.888 485.662] /A << /S /GoTo /D <6571756174696F6E2E422E322E312E3239> >> >>}
\MusicPDFWriteObject{5589}{<< /Type /Annot /Subtype /Link /Border [0 0 0] /H /I /C [1 0 0] /Rect [283.579 434.257 315.443 446.21] /A << /S /GoTo /D <74686D2E422E322E312E35> >> >>}
\MusicPDFWriteObject{5590}{<< /Type /Annot /Subtype /Link /Border [0 0 0] /H /I /C [1 0 0] /Rect [352.933 434.257 385.221 446.21] /A << /S /GoTo /D <6571756174696F6E2E422E322E312E37> >> >>}
\MusicPDFWriteObject{5591}{<< /Type /Annot /Subtype /Link /Border [0 0 0] /H /I /C [1 0 0] /Rect [224.696 237.622 256.985 249.575] /A << /S /GoTo /D <74686D2E422E322E312E34> >> >>}
\MusicPDFWriteObject{5592}{<< /Type /Annot /Subtype /Link /Border [0 0 0] /H /I /C [1 0 0] /Rect [145.456 164.767 182.725 176.719] /A << /S /GoTo /D <6571756174696F6E2E422E322E312E3336> >> >>}
\MusicPDFWriteObject{5593}{<< /Type /Annot /Subtype /Link /Border [0 0 0] /H /I /C [1 0 0] /Rect [211.166 164.767 248.434 176.719] /A << /S /GoTo /D <6571756174696F6E2E422E322E312E3334> >> >>}
\MusicPDFWriteObject{5594}{<< /Type /Annot /Subtype /Link /Border [0 0 0] /H /I /C [1 0 0] /Rect [166.201 119.816 203.469 131.769] /A << /S /GoTo /D <6571756174696F6E2E422E322E312E3337> >> >>}
\MusicPDFWriteObject{5595}{<< /Type /Annot /Subtype /Link /Border [0 0 0] /H /I /C [1 0 0] /Rect [149.635 684.633 186.903 696.586] /A << /S /GoTo /D <6571756174696F6E2E422E322E312E3238> >> >>}
\MusicPDFWriteObject{5596}{<< /Type /Annot /Subtype /Link /Border [0 0 0] /H /I /C [1 0 0] /Rect [91.42 598.377 128.688 610.33] /A << /S /GoTo /D <6571756174696F6E2E422E322E312E3238> >> >>}
\MusicPDFWriteObject{5597}{<< /Type /Annot /Subtype /Link /Border [0 0 0] /H /I /C [1 0 0] /Rect [383.861 522.231 421.129 534.184] /A << /S /GoTo /D <6571756174696F6E2E422E322E312E3236> >> >>}
\MusicPDFWriteObject{5598}{<< /Type /Annot /Subtype /Link /Border [0 0 0] /H /I /C [1 0 0] /Rect [124.623 490.67 156.911 502.623] /A << /S /GoTo /D <74686D2E422E322E312E39> >> >>}
\MusicPDFWriteObject{5599}{<< /Type /Annot /Subtype /Link /Border [0 0 0] /H /I /C [1 0 0] /Rect [164.526 451.218 201.794 463.171] /A << /S /GoTo /D <6571756174696F6E2E422E322E312E3236> >> >>}
\MusicPDFWriteObject{5600}{<< /Type /Annot /Subtype /Link /Border [0 0 0] /H /I /C [1 0 0] /Rect [287.021 451.218 324.29 463.171] /A << /S /GoTo /D <6571756174696F6E2E422E322E312E3237> >> >>}
\MusicPDFWriteObject{5601}{<< /Type /Annot /Subtype /Link /Border [0 0 0] /H /I /C [1 0 0] /Rect [459.718 451.218 496.986 463.171] /A << /S /GoTo /D <6571756174696F6E2E422E322E312E3339> >> >>}
\MusicPDFWriteObject{5602}{<< /Type /Annot /Subtype /Link /Border [0 0 0] /H /I /C [1 0 0] /Rect [118.591 198.32 150.88 210.273] /A << /S /GoTo /D <6571756174696F6E2E422E322E322E31> >> >>}
\MusicPDFWriteObject{5603}{<< /Type /Annot /Subtype /Link /Border [0 0 0] /H /I /C [1 0 0] /Rect [347.184 178.581 380.642 192.623] /A << /S /GoTo /D <74686D2E442E332E312E31> >> >>}
\MusicPDFWriteObject{5604}{<< /Type /Annot /Subtype /Link /Border [0 0 0] /H /I /C [1 0 0] /Rect [260.79 108.897 293.685 120.85] /A << /S /GoTo /D <74686D2E422E322E312E39> >> >>}
\MusicPDFWriteObject{5605}{<< /Type /Annot /Subtype /Link /Border [0 0 0] /H /I /C [1 0 0] /Rect [426.848 94.694 458.53 107.699] /A << /S /GoTo /D <74686D2E422E322E312E39> >> >>}
\MusicPDFWriteObject{5606}{<< /Type /Annot /Subtype /Link /Border [0 0 0] /H /I /C [1 0 0] /Rect [73.712 82.596 106 94.549] /A << /S /GoTo /D <6571756174696F6E2E422E322E322E31> >> >>}
\MusicPDFWriteObject{5607}{<< /Type /Annot /Subtype /Link /Border [0 0 0] /H /I /C [1 0 0] /Rect [158.219 684.633 190.538 696.586] /A << /S /GoTo /D <74686D2E422E322E312E35> >> >>}
\MusicPDFWriteObject{5608}{<< /Type /Annot /Subtype /Link /Border [0 0 0] /H /I /C [1 0 0] /Rect [265.963 349.088 298.797 361.04] /A << /S /GoTo /D <74686D2E422E322E312E35> >> >>}
\MusicPDFWriteObject{5609}{<< /Type /Annot /Subtype /Link /Border [0 0 0] /H /I /C [1 0 0] /Rect [337.22 349.088 369.508 361.04] /A << /S /GoTo /D <6571756174696F6E2E422E322E312E37> >> >>}
\MusicPDFWriteObject{5610}{<< /Type /Annot /Subtype /Link /Border [0 0 0] /H /I /C [1 0 0] /Rect [259.817 288.536 292.105 300.488] /A << /S /GoTo /D <6571756174696F6E2E422E322E322E33> >> >>}
\MusicPDFWriteObject{5611}{<< /Type /Annot /Subtype /Link /Border [0 0 0] /H /I /C [1 0 0] /Rect [399.537 288.536 432.431 300.488] /A << /S /GoTo /D <74686D2E422E322E312E36> >> >>}
\MusicPDFWriteObject{5612}{<< /Type /Annot /Subtype /Link /Border [0 0 0] /H /I /C [1 0 0] /Rect [243.85 222.782 276.139 234.735] /A << /S /GoTo /D <6571756174696F6E2E422E322E322E33> >> >>}
\MusicPDFWriteObject{5613}{<< /Type /Annot /Subtype /Link /Border [0 0 0] /H /I /C [1 0 0] /Rect [190.582 209.631 222.87 221.584] /A << /S /GoTo /D <6571756174696F6E2E422E322E322E32> >> >>}
\MusicPDFWriteObject{5614}{<< /Type /Annot /Subtype /Link /Border [0 0 0] /H /I /C [1 0 0] /Rect [271.979 209.631 304.267 221.584] /A << /S /GoTo /D <6571756174696F6E2E422E322E322E33> >> >>}
\MusicPDFWriteObject{5615}{<< /Type /Annot /Subtype /Link /Border [0 0 0] /H /I /C [1 0 0] /Rect [248.37 188.488 255.343 200.443] /A << /S /GoTo /D <6669677572652E422E32> >> >>}
\MusicPDFWriteObject{5616}{<< /Type /Annot /Subtype /Link /Border [0 0 0] /H /I /C [1 0 0] /Rect [115.743 660.822 125.676 671.189] /A << /S /GoTo /D <48666F6F746E6F74652E3839> >> >>}
\MusicPDFWriteObject{5617}{<< /Type /Annot /Subtype /Link /Border [0 0 0] /H /I /C [1 0 0] /Rect [206.687 567.065 243.955 579.018] /A << /S /GoTo /D <6571756174696F6E2E422E322E322E3132> >> >>}
\MusicPDFWriteObject{5618}{<< /Type /Annot /Subtype /Link /Border [0 0 0] /H /I /C [1 0 0] /Rect [359.722 567.065 396.99 579.018] /A << /S /GoTo /D <6571756174696F6E2E422E322E322E3132> >> >>}
\MusicPDFWriteObject{5619}{<< /Type /Annot /Subtype /Link /Border [0 0 0] /H /I /C [1 0 0] /Rect [87.106 506.216 119.394 518.169] /A << /S /GoTo /D <6571756174696F6E2E422E322E322E34> >> >>}
\MusicPDFWriteObject{5620}{<< /Type /Annot /Subtype /Link /Border [0 0 0] /H /I /C [1 0 0] /Rect [147.836 506.216 180.124 518.169] /A << /S /GoTo /D <6571756174696F6E2E422E322E322E35> >> >>}
\MusicPDFWriteObject{5621}{<< /Type /Annot /Subtype /Link /Border [0 0 0] /H /I /C [1 0 0] /Rect [332.97 438.27 370.239 450.223] /A << /S /GoTo /D <6571756174696F6E2E422E322E312E3237> >> >>}
\MusicPDFWriteObject{5622}{<< /Type /Annot /Subtype /Link /Border [0 0 0] /H /I /C [1 0 0] /Rect [246.86 345.999 284.129 357.952] /A << /S /GoTo /D <6571756174696F6E2E422E322E312E3237> >> >>}
\MusicPDFWriteObject{5623}{<< /Type /Annot /Subtype /Link /Border [0 0 0] /H /I /C [0 1 0] /Rect [512.226 335.339 543.173 344.207] /A << /S /GoTo /D <636974652E486F726D> >> >>}
\MusicPDFWriteObject{5624}{<< /Type /Annot /Subtype /Link /Border [0 0 0] /H /I /C [1 0 0] /Rect [325.21 319.698 335.143 332.555] /A << /S /GoTo /D <48666F6F746E6F74652E3930> >> >>}
\MusicPDFWriteObject{5625}{<< /Type /Annot /Subtype /Link /Border [0 0 0] /H /I /C [1 0 0] /Rect [209.588 136.462 241.876 148.415] /A << /S /GoTo /D <74686D2E422E322E322E34> >> >>}
\MusicPDFWriteObject{5626}{<< /Type /Annot /Subtype /Link /Border [0 0 0] /H /I /C [1 0 0] /Rect [185.188 671.482 217.264 683.435] /A << /S /GoTo /D <74686D2E422E322E322E34> >> >>}
\MusicPDFWriteObject{5627}{<< /Type /Annot /Subtype /Link /Border [0 0 0] /H /I /C [1 0 0] /Rect [465.529 645.181 502.127 657.134] /A << /S /GoTo /D <6571756174696F6E2E422E322E322E3134> >> >>}
\MusicPDFWriteObject{5628}{<< /Type /Annot /Subtype /Link /Border [0 0 0] /H /I /C [1 0 0] /Rect [260.34 445.773 293.234 456.563] /A << /S /GoTo /D <74686D2E422E322E322E31> >> >>}
\MusicPDFWriteObject{5629}{<< /Type /Annot /Subtype /Link /Border [0 0 0] /H /I /C [1 0 0] /Rect [334.021 469.952 366.915 482.957] /A << /S /GoTo /D <74686D2E422E332E312E31> >> >>}
\MusicPDFWriteObject{5630}{<< /Type /Annot /Subtype /Link /Border [0 0 0] /H /I /C [1 0 0] /Rect [537.923 192.914 544.845 205.229] /A << /S /GoTo /D <6669677572652E422E33> >> >>}
\MusicPDFWriteObject{5631}{<< /Type /Annot /Subtype /Link /Border [0 0 0] /H /I /C [1 0 0] /Rect [174.074 636.445 206.362 648.398] /A << /S /GoTo /D <6571756174696F6E2E422E332E322E31> >> >>}
\MusicPDFWriteObject{5632}{<< /Type /Annot /Subtype /Link /Border [0 0 0] /H /I /C [1 0 0] /Rect [187.794 551.811 223.527 560.665] /A << /S /GoTo /D <74686D2E422E332E322E31> >> >>}
\MusicPDFWriteObject{5633}{<< /Type /Annot /Subtype /Link /Border [0 0 0] /H /I /C [1 0 0] /Rect [208.886 536.17 241.78 548.125] /A << /S /GoTo /D <74686D2E422E332E312E32> >> >>}
\MusicPDFWriteObject{5634}{<< /Type /Annot /Subtype /Link /Border [0 0 0] /H /I /C [1 0 0] /Rect [162.449 445.072 194.738 457.024] /A << /S /GoTo /D <6571756174696F6E2E422E332E322E33> >> >>}
\MusicPDFWriteObject{5635}{<< /Type /Annot /Subtype /Link /Border [0 0 0] /H /I /C [1 0 0] /Rect [308.29 370.616 339.972 382.571] /A << /S /GoTo /D <74686D2E422E332E312E32> >> >>}
\MusicPDFWriteObject{5636}{<< /Type /Annot /Subtype /Link /Border [0 0 0] /H /I /C [1 0 0] /Rect [486.137 370.617 518.425 382.569] /A << /S /GoTo /D <6571756174696F6E2E422E332E322E34> >> >>}
\MusicPDFWriteObject{5637}{<< /Type /Annot /Subtype /Link /Border [0 0 0] /H /I /C [1 0 0] /Rect [280.338 357.465 312.626 369.42] /A << /S /GoTo /D <6571756174696F6E2E422E332E322E33> >> >>}
\MusicPDFWriteObject{5638}{<< /Type /Annot /Subtype /Link /Border [0 0 0] /H /I /C [1 0 0] /Rect [193.7 309.373 225.988 321.325] /A << /S /GoTo /D <6571756174696F6E2E422E332E322E31> >> >>}
\MusicPDFWriteObject{5639}{<< /Type /Annot /Subtype /Link /Border [0 0 0] /H /I /C [1 0 0] /Rect [319.885 309.373 352.173 321.325] /A << /S /GoTo /D <6571756174696F6E2E422E332E322E32> >> >>}
\MusicPDFWriteObject{5640}{<< /Type /Annot /Subtype /Link /Border [0 0 0] /H /I /C [0 1 0] /Rect [73.776 239.507 108.45 248.374] /A << /S /GoTo /D <636974652E5765696C2D62> >> >>}
\MusicPDFWriteObject{5641}{<< /Type /Annot /Subtype /Link /Border [0 0 0] /H /I /C [0 1 0] /Rect [234.469 239.507 254.895 248.374] /A << /S /GoTo /D <636974652E426172> >> >>}
\MusicPDFWriteObject{5642}{<< /Type /Annot /Subtype /Link /Border [0 0 0] /H /I /C [0 1 0] /Rect [320.859 198.118 357.581 208.922] /A << /S /GoTo /D <636974652E4C616E672D62> >> >>}
\MusicPDFWriteObject{5643}{<< /Type /Annot /Subtype /Link /Border [0 0 0] /H /I /C [0 1 0] /Rect [340.068 186.904 360.494 195.771] /A << /S /GoTo /D <636974652E426172> >> >>}
\MusicPDFWriteObject{5644}{<< /Type /Annot /Subtype /Link /Border [0 0 0] /H /I /C [0 1 0] /Rect [307.324 171.816 343.25 182.621] /A << /S /GoTo /D <636974652E4C616E672D61> >> >>}
\MusicPDFWriteObject{5645}{<< /Type /Annot /Subtype /Link /Border [0 0 0] /H /I /C [0 1 0] /Rect [472.625 171.816 509.347 182.621] /A << /S /GoTo /D <636974652E4C616E672D62> >> >>}
\MusicPDFWriteObject{5646}{<< /Type /Annot /Subtype /Link /Border [0 0 0] /H /I /C [0 1 0] /Rect [247.741 160.602 282.415 169.47] /A << /S /GoTo /D <636974652E5765696C2D62> >> >>}
\MusicPDFWriteObject{5647}{<< /Type /Annot /Subtype /Link /Border [0 0 0] /H /I /C [0 1 0] /Rect [73.715 542.465 100.414 551.333] /A << /S /GoTo /D <636974652E466F6C2D32> >> >>}
\MusicPDFWriteObject{5648}{<< /Type /Annot /Subtype /Link /Border [0 0 0] /H /I /C [1 0 0] /Rect [184.222 177.182 217.116 189.137] /A << /S /GoTo /D <6571756174696F6E2E422E342E312E32> >> >>}
\MusicPDFWriteObject{5649}{<< /Type /Annot /Subtype /Link /Border [0 0 0] /H /I /C [1 0 0] /Rect [237.714 122.048 269.487 134.001] /A << /S /GoTo /D <74686D2E422E322E312E39> >> >>}
\MusicPDFWriteObject{5650}{<< /Type /Annot /Subtype /Link /Border [0 0 0] /H /I /C [1 0 0] /Rect [164.237 540.114 196.525 552.067] /A << /S /GoTo /D <6571756174696F6E2E422E342E312E33> >> >>}
\MusicPDFWriteObject{5651}{<< /Type /Annot /Subtype /Link /Border [0 0 0] /H /I /C [1 0 0] /Rect [277.85 507.362 310.138 518.295] /A << /S /GoTo /D <74686D2E422E342E312E31> >> >>}
\MusicPDFWriteObject{5652}{<< /Type /Annot /Subtype /Link /Border [0 0 0] /H /I /C [1 0 0] /Rect [341.474 461.046 374.369 473.001] /A << /S /GoTo /D <74686D2E422E342E312E31> >> >>}
\MusicPDFWriteObject{5653}{<< /Type /Annot /Subtype /Link /Border [0 0 0] /H /I /C [0 1 0] /Rect [191.735 365.27 227.711 376.075] /A << /S /GoTo /D <636974652E50696E736B79> >> >>}
\MusicPDFWriteObject{5654}{<< /Type /Annot /Subtype /Link /Border [0 0 0] /H /I /C [1 0 0] /Rect [191.771 352.119 224.059 362.91] /A << /S /GoTo /D <74686D2E422E342E312E32> >> >>}
\MusicPDFWriteObject{5655}{<< /Type /Annot /Subtype /Link /Border [0 0 0] /H /I /C [1 0 0] /Rect [309.892 273.148 342.18 285.101] /A << /S /GoTo /D <6571756174696F6E2E422E342E312E37> >> >>}
\MusicPDFWriteObject{5656}{<< /Type /Annot /Subtype /Link /Border [0 0 0] /H /I /C [1 0 0] /Rect [200.499 259.997 210.432 272.855] /A << /S /GoTo /D <48666F6F746E6F74652E3931> >> >>}
\MusicPDFWriteObject{5657}{<< /Type /Annot /Subtype /Link /Border [0 0 0] /H /I /C [1 0 0] /Rect [210.517 136.271 242.805 148.224] /A << /S /GoTo /D <6571756174696F6E2E422E342E312E37> >> >>}
\MusicPDFWriteObject{5658}{<< /Type /Annot /Subtype /Link /Border [0 0 0] /H /I /C [1 0 0] /Rect [203.197 687.123 238.929 695.977] /A << /S /GoTo /D <74686D2E422E342E312E32> >> >>}
\MusicPDFWriteObject{5659}{<< /Type /Annot /Subtype /Link /Border [0 0 0] /H /I /C [1 0 0] /Rect [185.425 625.493 218.077 637.446] /A << /S /GoTo /D <6571756174696F6E2E422E342E312E37> >> >>}
\MusicPDFWriteObject{5660}{<< /Type /Annot /Subtype /Link /Border [0 0 0] /H /I /C [1 0 0] /Rect [337.894 625.493 370.545 637.446] /A << /S /GoTo /D <6571756174696F6E2E422E342E312E37> >> >>}
\MusicPDFWriteObject{5661}{<< /Type /Annot /Subtype /Link /Border [0 0 0] /H /I /C [0 1 0] /Rect [339.009 516.239 365.708 525.107] /A << /S /GoTo /D <636974652E466F6C2D32> >> >>}
\MusicPDFWriteObject{5662}{<< /Type /Annot /Subtype /Link /Border [0 0 0] /H /I /C [1 0 0] /Rect [191.356 407.363 228.625 419.316] /A << /S /GoTo /D <6571756174696F6E2E422E342E312E3131> >> >>}
\MusicPDFWriteObject{5663}{<< /Type /Annot /Subtype /Link /Border [0 0 0] /H /I /C [1 0 0] /Rect [272.092 350.165 309.36 362.118] /A << /S /GoTo /D <6571756174696F6E2E422E342E312E3131> >> >>}
\MusicPDFWriteObject{5664}{<< /Type /Annot /Subtype /Link /Border [0 0 0] /H /I /C [1 0 0] /Rect [219.888 578.246 257.156 590.198] /A << /S /GoTo /D <6571756174696F6E2E422E342E312E3134> >> >>}
\MusicPDFWriteObject{5665}{<< /Type /Annot /Subtype /Link /Border [0 0 0] /H /I /C [1 0 0] /Rect [167.926 480.82 205.195 492.773] /A << /S /GoTo /D <6571756174696F6E2E422E342E312E3130> >> >>}
\MusicPDFWriteObject{5666}{<< /Type /Annot /Subtype /Link /Border [0 0 0] /H /I /C [1 0 0] /Rect [233.636 480.82 270.904 492.773] /A << /S /GoTo /D <6571756174696F6E2E422E342E312E3136> >> >>}
\MusicPDFWriteObject{5667}{<< /Type /Annot /Subtype /Link /Border [0 0 0] /H /I /C [1 0 0] /Rect [294.686 415.346 331.955 427.298] /A << /S /GoTo /D <6571756174696F6E2E422E342E312E3137> >> >>}
\MusicPDFWriteObject{5668}{<< /Type /Annot /Subtype /Link /Border [0 0 0] /H /I /C [1 0 0] /Rect [254.31 337.444 286.598 349.397] /A << /S /GoTo /D <6571756174696F6E2E422E342E312E37> >> >>}
\MusicPDFWriteObject{5669}{<< /Type /Annot /Subtype /Link /Border [0 0 0] /H /I /C [0 1 0] /Rect [474.891 339.934 505.838 348.802] /A << /S /GoTo /D <636974652E486F726D> >> >>}
\MusicPDFWriteObject{5670}{<< /Type /Annot /Subtype /Link /Border [0 0 0] /H /I /C [0 1 0] /Rect [184.12 326.783 215.066 335.651] /A << /S /GoTo /D <636974652E486F726D> >> >>}
\MusicPDFWriteObject{5671}{<< /Type /Annot /Subtype /Link /Border [0 0 0] /H /I /C [0 1 0] /Rect [425.856 313.633 443.011 322.5] /A << /S /GoTo /D <636974652E466F6C> >> >>}
\MusicPDFWriteObject{5672}{<< /Type /Annot /Subtype /Link /Border [0 0 0] /H /I /C [0 1 0] /Rect [377.001 300.482 407.947 309.35] /A << /S /GoTo /D <636974652E486F726D> >> >>}
\MusicPDFWriteObject{5673}{<< /Type /Annot /Subtype /Link /Border [0 0 0] /H /I /C [1 0 0] /Rect [224.613 179.952 256.902 191.905] /A << /S /GoTo /D <74686D2E422E322E312E31> >> >>}
\MusicPDFWriteObject{5674}{<< /Type /Annot /Subtype /Link /Border [0 0 0] /H /I /C [1 0 0] /Rect [259.504 560.764 284.025 571.555] /A << /S /GoTo /D <73756273656374696F6E2E432E322E31> >> >>}
\MusicPDFWriteObject{5675}{<< /Type /Annot /Subtype /Link /Border [0 0 0] /H /I /C [1 0 0] /Rect [291.373 560.764 315.894 571.555] /A << /S /GoTo /D <73756273656374696F6E2E432E322E32> >> >>}
\MusicPDFWriteObject{5676}{<< /Type /Annot /Subtype /Link /Border [0 0 0] /H /I /C [0 1 0] /Rect [214.799 523.249 239.859 532.116] /A << /S /GoTo /D <636974652E53636861> >> >>}
\MusicPDFWriteObject{5677}{<< /Type /Annot /Subtype /Link /Border [0 0 0] /H /I /C [1 0 0] /Rect [236.794 507.608 269.018 519.561] /A << /S /GoTo /D <74686D2E442E312E332E37> >> >>}
\MusicPDFWriteObject{5678}{<< /Type /Annot /Subtype /Link /Border [0 0 0] /H /I /C [0 1 0] /Rect [464.916 483.796 489.976 492.664] /A << /S /GoTo /D <636974652E53636861> >> >>}
\MusicPDFWriteObject{5679}{<< /Type /Annot /Subtype /Link /Border [0 0 0] /H /I /C [0 1 0] /Rect [241.911 457.495 266.971 466.363] /A << /S /GoTo /D <636974652E53636861> >> >>}
\MusicPDFWriteObject{5680}{<< /Type /Annot /Subtype /Link /Border [0 0 0] /H /I /C [1 0 0] /Rect [444.165 407.914 476.756 418.704] /A << /S /GoTo /D <74686D2E422E342E312E32> >> >>}
\MusicPDFWriteObject{5681}{<< /Type /Annot /Subtype /Link /Border [0 0 0] /H /I /C [1 0 0] /Rect [286.006 394.21 318.294 406.162] /A << /S /GoTo /D <74686D2E422E332E322E31> >> >>}
\MusicPDFWriteObject{5682}{<< /Type /Annot /Subtype /Link /Border [0 0 0] /H /I /C [1 0 0] /Rect [245.285 346.239 277.573 358.194] /A << /S /GoTo /D <74686D2E422E322E312E31> >> >>}
\MusicPDFWriteObject{5683}{<< /Type /Annot /Subtype /Link /Border [0 0 0] /H /I /C [1 0 0] /Rect [222.763 196.547 255.051 207.338] /A << /S /GoTo /D <74686D2E422E342E312E37> >> >>}
\MusicPDFWriteObject{5684}{<< /Type /Annot /Subtype /Link /Border [0 0 0] /H /I /C [0 1 0] /Rect [356.406 177.141 378.796 186.009] /A << /S /GoTo /D <636974652E44656E> >> >>}
\MusicPDFWriteObject{5685}{<< /Type /Annot /Subtype /Link /Border [0 0 0] /H /I /C [0 1 0] /Rect [440.277 177.141 462.667 186.009] /A << /S /GoTo /D <636974652E44656E> >> >>}
\MusicPDFWriteObject{5686}{<< /Type /Annot /Subtype /Link /Border [0 0 0] /H /I /C [0 1 0] /Rect [520.802 177.141 543.192 186.009] /A << /S /GoTo /D <636974652E44656E> >> >>}
\MusicPDFWriteObject{5687}{<< /Type /Annot /Subtype /Link /Border [0 0 0] /H /I /C [0 1 0] /Rect [291.146 163.99 313.536 172.858] /A << /S /GoTo /D <636974652E44656E> >> >>}
\MusicPDFWriteObject{5688}{<< /Type /Annot /Subtype /Link /Border [0 0 0] /H /I /C [1 0 0] /Rect [498.197 135.199 530.849 147.152] /A << /S /GoTo /D <74686D2E422E322E312E39> >> >>}
\MusicPDFWriteObject{5689}{<< /Type /Annot /Subtype /Link /Border [0 0 0] /H /I /C [0 1 0] /Rect [252.489 124.538 274.879 133.406] /A << /S /GoTo /D <636974652E44656E> >> >>}
\MusicPDFWriteObject{5690}{<< /Type /Annot /Subtype /Link /Border [0 0 0] /H /I /C [0 1 0] /Rect [520.769 124.538 543.159 133.406] /A << /S /GoTo /D <636974652E44656E> >> >>}
\MusicPDFWriteObject{5691}{<< /Type /Annot /Subtype /Link /Border [0 0 0] /H /I /C [1 0 0] /Rect [345.077 108.897 377.365 120.85] /A << /S /GoTo /D <74686D2E422E322E312E39> >> >>}
\MusicPDFWriteObject{5692}{<< /Type /Annot /Subtype /Link /Border [0 0 0] /H /I /C [1 0 0] /Rect [160.309 95.747 197.577 107.699] /A << /S /GoTo /D <6571756174696F6E2E422E342E312E3230> >> >>}
\MusicPDFWriteObject{5693}{<< /Type /Annot /Subtype /Link /Border [0 0 0] /H /I /C [0 1 0] /Rect [285.508 98.237 307.898 107.105] /A << /S /GoTo /D <636974652E44656E> >> >>}
\MusicPDFWriteObject{5694}{<< /Type /Annot /Subtype /Link /Border [0 0 0] /H /I /C [0 1 0] /Rect [85.165 85.086 107.555 93.954] /A << /S /GoTo /D <636974652E44656E> >> >>}
\MusicPDFWriteObject{5695}{<< /Type /Annot /Subtype /Link /Border [0 0 0] /H /I /C [0 1 0] /Rect [423.47 85.086 445.86 93.954] /A << /S /GoTo /D <636974652E44656E> >> >>}
\MusicPDFWriteObject{5696}{<< /Type /Annot /Subtype /Link /Border [0 0 0] /H /I /C [0 1 0] /Rect [235.337 687.123 257.727 695.991] /A << /S /GoTo /D <636974652E44656E> >> >>}
\MusicPDFWriteObject{5697}{<< /Type /Annot /Subtype /Link /Border [0 0 0] /H /I /C [0 1 0] /Rect [461.396 685.187 498.118 695.991] /A << /S /GoTo /D <636974652E4C616E672D62> >> >>}
\MusicPDFWriteObject{5698}{<< /Type /Annot /Subtype /Link /Border [0 0 0] /H /I /C [1 0 0] /Rect [251.918 585.539 289.186 597.492] /A << /S /GoTo /D <6571756174696F6E2E422E342E312E3231> >> >>}
\MusicPDFWriteObject{5699}{<< /Type /Annot /Subtype /Link /Border [0 0 0] /H /I /C [1 0 0] /Rect [335.218 379.479 372.487 391.432] /A << /S /GoTo /D <6571756174696F6E2E422E342E312E3231> >> >>}
\MusicPDFWriteObject{5700}{<< /Type /Annot /Subtype /Link /Border [0 0 0] /H /I /C [1 0 0] /Rect [108.611 117.652 146.268 137.577] /A << /S /GoTo /D <6571756174696F6E2E422E342E312E3234> >> >>}
\MusicPDFWriteObject{5701}{<< /Type /Annot /Subtype /Link /Border [0 0 0] /H /I /C [1 0 0] /Rect [158.217 117.652 195.874 137.577] /A << /S /GoTo /D <6571756174696F6E2E422E342E312E3235> >> >>}
\MusicPDFWriteObject{5702}{<< /Type /Annot /Subtype /Link /Border [0 0 0] /H /I /C [1 0 0] /Rect [227.359 117.652 265.015 137.577] /A << /S /GoTo /D <6571756174696F6E2E422E342E312E3236> >> >>}
\MusicPDFWriteObject{5703}{<< /Type /Annot /Subtype /Link /Border [0 0 0] /H /I /C [1 0 0] /Rect [343.412 117.652 381.069 137.577] /A << /S /GoTo /D <6571756174696F6E2E422E342E312E3237> >> >>}
\MusicPDFWriteObject{5704}{<< /Type /Annot /Subtype /Link /Border [0 0 0] /H /I /C [1 0 0] /Rect [141.913 684.633 179.181 696.586] /A << /S /GoTo /D <6571756174696F6E2E422E342E312E3231> >> >>}
\MusicPDFWriteObject{5705}{<< /Type /Annot /Subtype /Link /Border [0 0 0] /H /I /C [1 0 0] /Rect [299.076 684.633 336.344 696.586] /A << /S /GoTo /D <6571756174696F6E2E422E342E312E3237> >> >>}
\MusicPDFWriteObject{5706}{<< /Type /Annot /Subtype /Link /Border [0 0 0] /H /I /C [1 0 0] /Rect [365.411 684.633 403.385 696.586] /A << /S /GoTo /D <6571756174696F6E2E422E342E312E3238> >> >>}
\MusicPDFWriteObject{5707}{<< /Type /Annot /Subtype /Link /Border [0 0 0] /H /I /C [1 0 0] /Rect [113.031 614.091 150.3 626.043] /A << /S /GoTo /D <6571756174696F6E2E422E342E312E3232> >> >>}
\MusicPDFWriteObject{5708}{<< /Type /Annot /Subtype /Link /Border [0 0 0] /H /I /C [1 0 0] /Rect [138.102 600.94 175.371 612.893] /A << /S /GoTo /D <6571756174696F6E2E422E342E312E3233> >> >>}
\MusicPDFWriteObject{5709}{<< /Type /Annot /Subtype /Link /Border [0 0 0] /H /I /C [1 0 0] /Rect [388.76 539.281 426.029 551.233] /A << /S /GoTo /D <6571756174696F6E2E422E342E312E3330> >> >>}
\MusicPDFWriteObject{5710}{<< /Type /Annot /Subtype /Link /Border [0 0 0] /H /I /C [1 0 0] /Rect [219.115 498.59 256.384 510.543] /A << /S /GoTo /D <6571756174696F6E2E422E342E312E3233> >> >>}
\MusicPDFWriteObject{5711}{<< /Type /Annot /Subtype /Link /Border [0 0 0] /H /I /C [1 0 0] /Rect [176.323 453.767 212.056 462.621] /A << /S /GoTo /D <74686D2E422E342E312E37> >> >>}
\MusicPDFWriteObject{5712}{<< /Type /Annot /Subtype /Link /Border [0 0 0] /H /I /C [0 1 0] /Rect [298.313 453.767 320.704 462.635] /A << /S /GoTo /D <636974652E44656E> >> >>}
\MusicPDFWriteObject{5713}{<< /Type /Annot /Subtype /Link /Border [0 0 0] /H /I /C [1 0 0] /Rect [86.497 409.203 123.766 421.155] /A << /S /GoTo /D <6571756174696F6E2E422E342E312E3231> >> >>}
\MusicPDFWriteObject{5714}{<< /Type /Annot /Subtype /Link /Border [0 0 0] /H /I /C [1 0 0] /Rect [210.223 305.933 247.491 317.886] /A << /S /GoTo /D <6571756174696F6E2E422E342E312E3331> >> >>}
\MusicPDFWriteObject{5715}{<< /Type /Annot /Subtype /Link /Border [0 0 0] /H /I /C [1 0 0] /Rect [413.525 305.933 445.844 317.886] /A << /S /GoTo /D <74686D2E422E342E312E39> >> >>}
\MusicPDFWriteObject{5716}{<< /Type /Annot /Subtype /Link /Border [0 0 0] /H /I /C [1 0 0] /Rect [106.903 292.783 144.172 304.735] /A << /S /GoTo /D <6571756174696F6E2E422E342E312E3332> >> >>}
\MusicPDFWriteObject{5717}{<< /Type /Annot /Subtype /Link /Border [0 0 0] /H /I /C [1 0 0] /Rect [219.438 279.632 256.707 291.585] /A << /S /GoTo /D <6571756174696F6E2E422E342E312E3331> >> >>}
\MusicPDFWriteObject{5718}{<< /Type /Annot /Subtype /Link /Border [0 0 0] /H /I /C [1 0 0] /Rect [228.569 266.481 265.837 278.434] /A << /S /GoTo /D <6571756174696F6E2E422E342E312E3233> >> >>}
\MusicPDFWriteObject{5719}{<< /Type /Annot /Subtype /Link /Border [0 0 0] /H /I /C [1 0 0] /Rect [330.135 266.481 362.423 278.434] /A << /S /GoTo /D <74686D2E422E342E312E39> >> >>}
\MusicPDFWriteObject{5720}{<< /Type /Annot /Subtype /Link /Border [0 0 0] /H /I /C [0 1 0] /Rect [99.731 163.642 122.121 172.51] /A << /S /GoTo /D <636974652E44656E> >> >>}
\MusicPDFWriteObject{5721}{<< /Type /Annot /Subtype /Link /Border [0 0 0] /H /I /C [1 0 0] /Rect [292.032 671.122 330.006 683.437] /A << /S /GoTo /D <6571756174696F6E2E422E342E312E3333> >> >>}
\MusicPDFWriteObject{5722}{<< /Type /Annot /Subtype /Link /Border [0 0 0] /H /I /C [1 0 0] /Rect [279.7 658.885 311.988 669.676] /A << /S /GoTo /D <74686D2E422E322E312E39> >> >>}
\MusicPDFWriteObject{5723}{<< /Type /Annot /Subtype /Link /Border [0 0 0] /H /I /C [1 0 0] /Rect [356.325 643.298 388.613 657.689] /A << /S /GoTo /D <74686D2E422E322E322E34> >> >>}
\MusicPDFWriteObject{5724}{<< /Type /Annot /Subtype /Link /Border [0 0 0] /H /I /C [1 0 0] /Rect [115.675 591.361 147.963 603.313] /A << /S /GoTo /D <6571756174696F6E2E422E322E322E36> >> >>}
\MusicPDFWriteObject{5725}{<< /Type /Annot /Subtype /Link /Border [0 0 0] /H /I /C [1 0 0] /Rect [216.729 591.36 248.412 603.315] /A << /S /GoTo /D <74686D2E422E322E322E34> >> >>}
\MusicPDFWriteObject{5726}{<< /Type /Annot /Subtype /Link /Border [0 0 0] /H /I /C [1 0 0] /Rect [140.545 538.188 172.985 552.073] /A << /S /GoTo /D <74686D2E422E342E312E32> >> >>}
\MusicPDFWriteObject{5727}{<< /Type /Annot /Subtype /Link /Border [0 0 0] /H /I /C [1 0 0] /Rect [112.56 526.33 149.828 538.283] /A << /S /GoTo /D <6571756174696F6E2E422E342E312E3230> >> >>}
\MusicPDFWriteObject{5728}{<< /Type /Annot /Subtype /Link /Border [0 0 0] /H /I /C [0 1 0] /Rect [506.492 462.096 543.214 472.901] /A << /S /GoTo /D <636974652E4C616E672D62> >> >>}
\MusicPDFWriteObject{5729}{<< /Type /Annot /Subtype /Link /Border [0 0 0] /H /I /C [1 0 0] /Rect [328.869 401.707 361.763 414.712] /A << /S /GoTo /D <74686D2E422E332E312E31> >> >>}
\MusicPDFWriteObject{5730}{<< /Type /Annot /Subtype /Link /Border [0 0 0] /H /I /C [1 0 0] /Rect [193.266 363.307 225.555 375.26] /A << /S /GoTo /D <74686D2E422E332E322E31> >> >>}
\MusicPDFWriteObject{5731}{<< /Type /Annot /Subtype /Link /Border [0 0 0] /H /I /C [1 0 0] /Rect [307.529 189.042 337.968 201.9] /A << /S /GoTo /D <74686D2E332E322E312E33> >> >>}
\MusicPDFWriteObject{5732}{<< /Type /Annot /Subtype /Link /Border [0 0 0] /H /I /C [1 0 0] /Rect [420.54 189.042 430.473 201.9] /A << /S /GoTo /D <48666F6F746E6F74652E3932> >> >>}
\MusicPDFWriteObject{5733}{<< /Type /Annot /Subtype /Link /Border [0 0 0] /H /I /C [1 0 0] /Rect [452.194 531.116 482.407 543.069] /A << /S /GoTo /D <6571756174696F6E2E332E322E312E31> >> >>}
\MusicPDFWriteObject{5734}{<< /Type /Annot /Subtype /Link /Border [0 0 0] /H /I /C [1 0 0] /Rect [509.731 531.116 539.944 543.069] /A << /S /GoTo /D <6571756174696F6E2E332E322E312E33> >> >>}
\MusicPDFWriteObject{5735}{<< /Type /Annot /Subtype /Link /Border [0 0 0] /H /I /C [0 1 0] /Rect [462.102 413.213 496.777 422.081] /A << /S /GoTo /D <636974652E5765696C2D62> >> >>}
\MusicPDFWriteObject{5736}{<< /Type /Annot /Subtype /Link /Border [0 0 0] /H /I /C [0 1 0] /Rect [506.47 411.276 543.192 422.081] /A << /S /GoTo /D <636974652E4C616E672D62> >> >>}
\MusicPDFWriteObject{5737}{<< /Type /Annot /Subtype /Link /Border [0 0 0] /H /I /C [0 1 0] /Rect [185.823 400.062 206.249 408.93] /A << /S /GoTo /D <636974652E426172> >> >>}
\MusicPDFWriteObject{5738}{<< /Type /Annot /Subtype /Link /Border [0 0 0] /H /I /C [1 0 0] /Rect [241.438 657.06 273.726 669.013] /A << /S /GoTo /D <6571756174696F6E2E422E342E322E34> >> >>}
\MusicPDFWriteObject{5739}{<< /Type /Annot /Subtype /Link /Border [0 0 0] /H /I /C [1 0 0] /Rect [128.189 467.453 138.122 480.31] /A << /S /GoTo /D <48666F6F746E6F74652E3933> >> >>}
\MusicPDFWriteObject{5740}{<< /Type /Annot /Subtype /Link /Border [0 0 0] /H /I /C [0 1 0] /Rect [306.351 469.943 326.777 478.811] /A << /S /GoTo /D <636974652E426172> >> >>}
\MusicPDFWriteObject{5741}{<< /Type /Annot /Subtype /Link /Border [0 0 0] /H /I /C [1 0 0] /Rect [447.624 454.301 479.306 466.256] /A << /S /GoTo /D <74686D2E422E332E322E31> >> >>}
\MusicPDFWriteObject{5742}{<< /Type /Annot /Subtype /Link /Border [0 0 0] /H /I /C [1 0 0] /Rect [179.631 428 216.899 439.953] /A << /S /GoTo /D <6571756174696F6E2E422E342E322E3134> >> >>}
\MusicPDFWriteObject{5743}{<< /Type /Annot /Subtype /Link /Border [0 0 0] /H /I /C [1 0 0] /Rect [186.925 378.773 224.193 390.726] /A << /S /GoTo /D <6571756174696F6E2E422E342E322E3136> >> >>}
\MusicPDFWriteObject{5744}{<< /Type /Annot /Subtype /Link /Border [0 0 0] /H /I /C [0 1 0] /Rect [302.223 381.263 322.649 390.131] /A << /S /GoTo /D <636974652E426172> >> >>}
\MusicPDFWriteObject{5745}{<< /Type /Annot /Subtype /Link /Border [0 0 0] /H /I /C [1 0 0] /Rect [184.017 365.622 221.285 377.575] /A << /S /GoTo /D <6571756174696F6E2E422E342E322E3135> >> >>}
\MusicPDFWriteObject{5746}{<< /Type /Annot /Subtype /Link /Border [0 0 0] /H /I /C [0 1 0] /Rect [383.561 368.112 403.987 376.98] /A << /S /GoTo /D <636974652E426172> >> >>}
\MusicPDFWriteObject{5747}{<< /Type /Annot /Subtype /Link /Border [0 0 0] /H /I /C [1 0 0] /Rect [248.433 352.472 285.702 364.424] /A << /S /GoTo /D <6571756174696F6E2E422E342E322E3135> >> >>}
\MusicPDFWriteObject{5748}{<< /Type /Annot /Subtype /Link /Border [0 0 0] /H /I /C [1 0 0] /Rect [313.458 352.472 350.726 364.424] /A << /S /GoTo /D <6571756174696F6E2E422E342E322E3136> >> >>}
\MusicPDFWriteObject{5749}{<< /Type /Annot /Subtype /Link /Border [0 0 0] /H /I /C [1 0 0] /Rect [229.786 289.538 239.719 301.842] /A << /S /GoTo /D <48666F6F746E6F74652E3934> >> >>}
\MusicPDFWriteObject{5750}{<< /Type /Annot /Subtype /Link /Border [0 0 0] /H /I /C [0 1 0] /Rect [267.195 95.678 284.894 103.17] /A << /S /GoTo /D <636974652E426172> >> >>}
\MusicPDFWriteObject{5751}{<< /Type /Annot /Subtype /Link /Border [0 0 0] /H /I /C [1 0 0] /Rect [199.082 684.633 236.35 696.586] /A << /S /GoTo /D <6571756174696F6E2E422E342E322E3137> >> >>}
\MusicPDFWriteObject{5752}{<< /Type /Annot /Subtype /Link /Border [0 0 0] /H /I /C [1 0 0] /Rect [355.686 684.633 387.974 696.586] /A << /S /GoTo /D <6571756174696F6E2E422E342E322E32> >> >>}
\MusicPDFWriteObject{5753}{<< /Type /Annot /Subtype /Link /Border [0 0 0] /H /I /C [1 0 0] /Rect [419.133 684.633 451.421 696.586] /A << /S /GoTo /D <6571756174696F6E2E422E342E322E35> >> >>}
\MusicPDFWriteObject{5754}{<< /Type /Annot /Subtype /Link /Border [0 0 0] /H /I /C [1 0 0] /Rect [463.12 684.633 495.409 696.586] /A << /S /GoTo /D <6571756174696F6E2E422E342E322E36> >> >>}
\MusicPDFWriteObject{5755}{<< /Type /Annot /Subtype /Link /Border [0 0 0] /H /I /C [1 0 0] /Rect [507.108 684.633 539.396 696.586] /A << /S /GoTo /D <6571756174696F6E2E422E342E322E37> >> >>}
\MusicPDFWriteObject{5756}{<< /Type /Annot /Subtype /Link /Border [0 0 0] /H /I /C [1 0 0] /Rect [73.712 671.482 110.98 683.435] /A << /S /GoTo /D <6571756174696F6E2E422E342E322E3130> >> >>}
\MusicPDFWriteObject{5757}{<< /Type /Annot /Subtype /Link /Border [0 0 0] /H /I /C [1 0 0] /Rect [139.422 671.482 176.69 683.435] /A << /S /GoTo /D <6571756174696F6E2E422E342E322E3131> >> >>}
\MusicPDFWriteObject{5758}{<< /Type /Annot /Subtype /Link /Border [0 0 0] /H /I /C [1 0 0] /Rect [300.279 671.482 332.567 683.435] /A << /S /GoTo /D <74686D2E422E342E322E31> >> >>}
\MusicPDFWriteObject{5759}{<< /Type /Annot /Subtype /Link /Border [0 0 0] /H /I /C [0 1 0] /Rect [375.861 651.623 410.535 660.491] /A << /S /GoTo /D <636974652E5765696C2D62> >> >>}
\MusicPDFWriteObject{5760}{<< /Type /Annot /Subtype /Link /Border [0 0 0] /H /I /C [0 1 0] /Rect [474.599 649.687 510.525 660.491] /A << /S /GoTo /D <636974652E4C616E672D61> >> >>}
\MusicPDFWriteObject{5761}{<< /Type /Annot /Subtype /Link /Border [0 0 0] /H /I /C [0 1 0] /Rect [104.706 638.473 125.132 647.34] /A << /S /GoTo /D <636974652E426172> >> >>}
\MusicPDFWriteObject{5762}{<< /Type /Annot /Subtype /Link /Border [0 0 0] /H /I /C [0 1 0] /Rect [224.287 636.536 261.009 647.34] /A << /S /GoTo /D <636974652E4C616E672D62> >> >>}
\MusicPDFWriteObject{5763}{<< /Type /Annot /Subtype /Link /Border [0 0 0] /H /I /C [1 0 0] /Rect [209.506 596.53 241.794 608.483] /A << /S /GoTo /D <74686D2E422E332E322E31> >> >>}
\MusicPDFWriteObject{5764}{<< /Type /Annot /Subtype /Link /Border [0 0 0] /H /I /C [0 1 0] /Rect [323.4 562.232 339.392 571.1] /A << /S /GoTo /D <636974652E436F> >> >>}
\MusicPDFWriteObject{5765}{<< /Type /Annot /Subtype /Link /Border [0 0 0] /H /I /C [1 0 0] /Rect [189.644 533.993 204.109 544.841] /A << /S /GoTo /D <73656374696F6E2E322E31> >> >>}
\MusicPDFWriteObject{5766}{<< /Type /Annot /Subtype /Link /Border [0 0 0] /H /I /C [0 1 0] /Rect [138.884 522.78 160.106 531.647] /A << /S /GoTo /D <636974652E427275> >> >>}
\MusicPDFWriteObject{5767}{<< /Type /Annot /Subtype /Link /Border [0 0 0] /H /I /C [1 0 0] /Rect [71.004 429.825 100.652 444.32] /A << /S /GoTo /D <74686D2E362E312E322E32> >> >>}
\MusicPDFWriteObject{5768}{<< /Type /Annot /Subtype /Link /Border [0 0 0] /H /I /C [1 0 0] /Rect [70.755 416.421 100.403 430.08] /A << /S /GoTo /D <74686D2E332E342E332E33> >> >>}
\MusicPDFWriteObject{5769}{<< /Type /Annot /Subtype /Link /Border [0 0 0] /H /I /C [1 0 0] /Rect [422.168 416.421 451.817 430.08] /A << /S /GoTo /D <74686D2E352E352E332E31> >> >>}
\MusicPDFWriteObject{5770}{<< /Type /Annot /Subtype /Link /Border [0 0 0] /H /I /C [1 0 0] /Rect [211.548 379.077 244.866 389.867] /A << /S /GoTo /D <74686D2E412E332E312E38> >> >>}
\MusicPDFWriteObject{5771}{<< /Type /Annot /Subtype /Link /Border [0 0 0] /H /I /C [1 0 0] /Rect [308.039 379.077 341.357 389.867] /A << /S /GoTo /D <74686D2E412E332E312E39> >> >>}
\MusicPDFWriteObject{5772}{<< /Type /Annot /Subtype /Link /Border [0 0 0] /H /I /C [1 0 0] /Rect [432.514 684.66 466.346 697.048] /A << /S /GoTo /D <74686D2E412E332E322E32> >> >>}
\MusicPDFWriteObject{5773}{<< /Type /Annot /Subtype /Link /Border [0 0 0] /H /I /C [1 0 0] /Rect [444.332 669.168 476.62 681.123] /A << /S /GoTo /D <6571756174696F6E2E422E342E332E35> >> >>}
\MusicPDFWriteObject{5774}{<< /Type /Annot /Subtype /Link /Border [0 0 0] /H /I /C [1 0 0] /Rect [260.397 649.164 292.686 661.137] /A << /S /GoTo /D <6571756174696F6E2E422E342E332E33> >> >>}
\MusicPDFWriteObject{5775}{<< /Type /Annot /Subtype /Link /Border [0 0 0] /H /I /C [1 0 0] /Rect [394.214 584.822 427.531 595.67] /A << /S /GoTo /D <74686D2E412E332E312E32> >> >>}
\MusicPDFWriteObject{5776}{<< /Type /Annot /Subtype /Link /Border [0 0 0] /H /I /C [0 1 0] /Rect [211.166 573.608 228.638 582.476] /A << /S /GoTo /D <636974652E4346> >> >>}
\MusicPDFWriteObject{5777}{<< /Type /Annot /Subtype /Link /Border [0 0 0] /H /I /C [1 0 0] /Rect [219.975 478.362 252.263 490.777] /A << /S /GoTo /D <6571756174696F6E2E422E342E332E37> >> >>}
\MusicPDFWriteObject{5778}{<< /Type /Annot /Subtype /Link /Border [0 0 0] /H /I /C [1 0 0] /Rect [139.979 432.747 172.267 444.7] /A << /S /GoTo /D <6571756174696F6E2E422E342E332E33> >> >>}
\MusicPDFWriteObject{5779}{<< /Type /Annot /Subtype /Link /Border [0 0 0] /H /I /C [1 0 0] /Rect [284.922 419.547 317.21 431.5] /A << /S /GoTo /D <6571756174696F6E2E422E342E332E33> >> >>}
\MusicPDFWriteObject{5780}{<< /Type /Annot /Subtype /Link /Border [0 0 0] /H /I /C [1 0 0] /Rect [300.825 269.114 333.114 281.067] /A << /S /GoTo /D <6571756174696F6E2E422E342E332E34> >> >>}
\MusicPDFWriteObject{5781}{<< /Type /Annot /Subtype /Link /Border [0 0 0] /H /I /C [1 0 0] /Rect [170.744 182.251 208.012 194.204] /A << /S /GoTo /D <6571756174696F6E2E422E342E332E3130> >> >>}
\MusicPDFWriteObject{5782}{<< /Type /Annot /Subtype /Link /Border [0 0 0] /H /I /C [1 0 0] /Rect [241.612 182.251 273.9 194.204] /A << /S /GoTo /D <6571756174696F6E2E422E342E332E39> >> >>}
\MusicPDFWriteObject{5783}{<< /Type /Annot /Subtype /Link /Border [0 0 0] /H /I /C [1 0 0] /Rect [509.6 182.251 541.888 194.204] /A << /S /GoTo /D <6571756174696F6E2E422E342E332E35> >> >>}
\MusicPDFWriteObject{5784}{<< /Type /Annot /Subtype /Link /Border [0 0 0] /H /I /C [0 1 0] /Rect [363.961 445.157 379.953 454.025] /A << /S /GoTo /D <636974652E436F> >> >>}
\MusicPDFWriteObject{5785}{<< /Type /Annot /Subtype /Link /Border [0 0 0] /H /I /C [1 0 0] /Rect [515.301 442.667 525.234 455.524] /A << /S /GoTo /D <48666F6F746E6F74652E3935> >> >>}
\MusicPDFWriteObject{5786}{<< /Type /Annot /Subtype /Link /Border [0 0 0] /H /I /C [0 1 0] /Rect [138.843 432.006 154.835 440.874] /A << /S /GoTo /D <636974652E436F> >> >>}
\MusicPDFWriteObject{5787}{<< /Type /Annot /Subtype /Link /Border [0 0 0] /H /I /C [0 1 0] /Rect [288.14 432.006 304.132 440.874] /A << /S /GoTo /D <636974652E436F> >> >>}
\MusicPDFWriteObject{5788}{<< /Type /Annot /Subtype /Link /Border [0 0 0] /H /I /C [1 0 0] /Rect [243.843 416.365 276.132 428.318] /A << /S /GoTo /D <74686D2E422E342E332E32> >> >>}
\MusicPDFWriteObject{5789}{<< /Type /Annot /Subtype /Link /Border [0 0 0] /H /I /C [0 1 0] /Rect [308.582 379.403 324.575 388.271] /A << /S /GoTo /D <636974652E436F> >> >>}
\MusicPDFWriteObject{5790}{<< /Type /Annot /Subtype /Link /Border [0 0 0] /H /I /C [1 0 0] /Rect [432.145 376.913 442.078 389.771] /A << /S /GoTo /D <48666F6F746E6F74652E3936> >> >>}
\MusicPDFWriteObject{5791}{<< /Type /Annot /Subtype /Link /Border [0 0 0] /H /I /C [1 0 0] /Rect [336.585 283.166 346.518 296.384] /A << /S /GoTo /D <48666F6F746E6F74652E3937> >> >>}
\MusicPDFWriteObject{5792}{<< /Type /Annot /Subtype /Link /Border [0 0 0] /H /I /C [1 0 0] /Rect [178.276 267.916 187.573 280.774] /A << /S /GoTo /D <48666F6F746E6F74652E3938> >> >>}
\MusicPDFWriteObject{5793}{<< /Type /Annot /Subtype /Link /Border [0 0 0] /H /I /C [1 0 0] /Rect [329.779 267.916 339.712 280.774] /A << /S /GoTo /D <48666F6F746E6F74652E3939> >> >>}
\MusicPDFWriteObject{5794}{<< /Type /Annot /Subtype /Link /Border [0 0 0] /H /I /C [1 0 0] /Rect [516.604 196.202 526.537 210.869] /A << /S /GoTo /D <48666F6F746E6F74652E313030> >> >>}
\MusicPDFWriteObject{5795}{<< /Type /Annot /Subtype /Link /Border [0 0 0] /H /I /C [1 0 0] /Rect [178.276 684.633 188.209 697.491] /A << /S /GoTo /D <48666F6F746E6F74652E313031> >> >>}
\MusicPDFWriteObject{5796}{<< /Type /Annot /Subtype /Link /Border [0 0 0] /H /I /C [1 0 0] /Rect [257.329 633.904 289.617 645.857] /A << /S /GoTo /D <6571756174696F6E2E422E342E342E32> >> >>}
\MusicPDFWriteObject{5797}{<< /Type /Annot /Subtype /Link /Border [0 0 0] /H /I /C [1 0 0] /Rect [339.78 615.545 372.068 634.54] /A << /S /GoTo /D <6571756174696F6E2E422E342E342E35> >> >>}
\MusicPDFWriteObject{5798}{<< /Type /Annot /Subtype /Link /Border [0 0 0] /H /I /C [1 0 0] /Rect [257.327 572.28 289.552 584.235] /A << /S /GoTo /D <74686D2E442E322E312E31> >> >>}
\MusicPDFWriteObject{5799}{<< /Type /Annot /Subtype /Link /Border [0 0 0] /H /I /C [1 0 0] /Rect [484.14 423.224 516.982 435.176] /A << /S /GoTo /D <6571756174696F6E2E442E322E312E34> >> >>}
\MusicPDFWriteObject{5800}{<< /Type /Annot /Subtype /Link /Border [0 0 0] /H /I /C [1 0 0] /Rect [117.346 410.073 150.188 422.026] /A << /S /GoTo /D <74686D2E442E322E312E31> >> >>}
\MusicPDFWriteObject{5801}{<< /Type /Annot /Subtype /Link /Border [0 0 0] /H /I /C [1 0 0] /Rect [98.699 383.799 108.632 396.629] /A << /S /GoTo /D <48666F6F746E6F74652E313032> >> >>}
\MusicPDFWriteObject{5802}{<< /Type /Annot /Subtype /Link /Border [0 0 0] /H /I /C [1 0 0] /Rect [345.028 82.097 366.183 94.681] /A << /S /GoTo /D <73756273656374696F6E2E422E342E33> >> >>}
\MusicPDFWriteObject{5803}{<< /Type /Annot /Subtype /Link /Border [0 0 0] /H /I /C [1 0 0] /Rect [179.374 685.187 212.077 695.977] /A << /S /GoTo /D <74686D2E412E332E322E33> >> >>}
\MusicPDFWriteObject{5804}{<< /Type /Annot /Subtype /Link /Border [0 0 0] /H /I /C [1 0 0] /Rect [509.6 578.341 541.888 590.294] /A << /S /GoTo /D <6571756174696F6E2E422E342E332E35> >> >>}
\MusicPDFWriteObject{5805}{<< /Type /Annot /Subtype /Link /Border [0 0 0] /H /I /C [1 0 0] /Rect [148.168 565.191 180.456 577.143] /A << /S /GoTo /D <74686D2E422E342E332E31> >> >>}
\MusicPDFWriteObject{5806}{<< /Type /Annot /Subtype /Link /Border [0 0 0] /H /I /C [1 0 0] /Rect [205.024 565.191 237.312 577.143] /A << /S /GoTo /D <6571756174696F6E2E422E342E342E31> >> >>}
\MusicPDFWriteObject{5807}{<< /Type /Annot /Subtype /Link /Border [0 0 0] /H /I /C [1 0 0] /Rect [305.594 565.191 337.882 577.143] /A << /S /GoTo /D <6571756174696F6E2E422E342E342E32> >> >>}
\MusicPDFWriteObject{5808}{<< /Type /Annot /Subtype /Link /Border [0 0 0] /H /I /C [1 0 0] /Rect [146.304 394.146 178.592 406.099] /A << /S /GoTo /D <6571756174696F6E2E422E342E342E35> >> >>}
\MusicPDFWriteObject{5809}{<< /Type /Annot /Subtype /Link /Border [0 0 0] /H /I /C [1 0 0] /Rect [173.784 287.347 206.073 299.3] /A << /S /GoTo /D <6571756174696F6E2E422E342E342E35> >> >>}
\MusicPDFWriteObject{5810}{<< /Type /Annot /Subtype /Link /Border [0 0 0] /H /I /C [0 1 0] /Rect [195.562 120.353 211.554 129.221] /A << /S /GoTo /D <636974652E436F> >> >>}
\MusicPDFWriteObject{5811}{<< /Type /Annot /Subtype /Link /Border [0 0 0] /H /I /C [1 0 0] /Rect [275.099 539.093 307.387 551.046] /A << /S /GoTo /D <6571756174696F6E2E422E342E332E34> >> >>}
\MusicPDFWriteObject{5812}{<< /Type /Annot /Subtype /Link /Border [0 0 0] /H /I /C [1 0 0] /Rect [393.897 539.093 426.791 551.046] /A << /S /GoTo /D <74686D2E422E342E332E31> >> >>}
\MusicPDFWriteObject{5813}{<< /Type /Annot /Subtype /Link /Border [0 0 0] /H /I /C [0 1 0] /Rect [525.687 527.692 543.159 536.56] /A << /S /GoTo /D <636974652E4346> >> >>}
\MusicPDFWriteObject{5814}{<< /Type /Annot /Subtype /Link /Border [0 0 0] /H /I /C [1 0 0] /Rect [114.631 375.403 146.798 387.995] /A << /S /GoTo /D <6571756174696F6E2E422E342E332E34> >> >>}
\MusicPDFWriteObject{5815}{<< /Type /Annot /Subtype /Link /Border [0 0 0] /H /I /C [1 0 0] /Rect [235.851 375.403 268.018 387.995] /A << /S /GoTo /D <6571756174696F6E2E422E342E342E31> >> >>}
\MusicPDFWriteObject{5816}{<< /Type /Annot /Subtype /Link /Border [0 0 0] /H /I /C [1 0 0] /Rect [207.466 362.252 239.754 374.205] /A << /S /GoTo /D <6571756174696F6E2E422E342E332E34> >> >>}
\MusicPDFWriteObject{5817}{<< /Type /Annot /Subtype /Link /Border [0 0 0] /H /I /C [1 0 0] /Rect [106.473 273.724 143.741 285.677] /A << /S /GoTo /D <6571756174696F6E2E422E342E342E3131> >> >>}
\MusicPDFWriteObject{5818}{<< /Type /Annot /Subtype /Link /Border [0 0 0] /H /I /C [1 0 0] /Rect [209.063 273.724 246.332 285.677] /A << /S /GoTo /D <6571756174696F6E2E422E342E342E3130> >> >>}
\MusicPDFWriteObject{5819}{<< /Type /Annot /Subtype /Link /Border [0 0 0] /H /I /C [0 1 0] /Rect [221.278 119.226 254.839 128.094] /A << /S /GoTo /D <636974652E5765696C2D65> >> >>}
\MusicPDFWriteObject{5820}{<< /Type /Annot /Subtype /Link /Border [0 0 0] /H /I /C [0 1 0] /Rect [212.411 659.032 228.403 667.899] /A << /S /GoTo /D <636974652E436F> >> >>}
\MusicPDFWriteObject{5821}{<< /Type /Annot /Subtype /Link /Border [0 0 0] /H /I /C [1 0 0] /Rect [327.242 635.736 364.511 647.689] /A << /S /GoTo /D <6571756174696F6E2E422E342E342E3134> >> >>}
\MusicPDFWriteObject{5822}{<< /Type /Annot /Subtype /Link /Border [0 0 0] /H /I /C [1 0 0] /Rect [413.325 635.736 445.492 647.689] /A << /S /GoTo /D <6571756174696F6E2E422E342E342E35> >> >>}
\MusicPDFWriteObject{5823}{<< /Type /Annot /Subtype /Link /Border [0 0 0] /H /I /C [1 0 0] /Rect [473.833 635.736 510.961 647.689] /A << /S /GoTo /D <6571756174696F6E2E422E342E342E3131> >> >>}
\MusicPDFWriteObject{5824}{<< /Type /Annot /Subtype /Link /Border [0 0 0] /H /I /C [1 0 0] /Rect [73.712 622.585 110.98 634.538] /A << /S /GoTo /D <6571756174696F6E2E422E342E342E3134> >> >>}
\MusicPDFWriteObject{5825}{<< /Type /Annot /Subtype /Link /Border [0 0 0] /H /I /C [1 0 0] /Rect [238.65 563.95 270.938 580.112] /A << /S /GoTo /D <74686D2E422E342E342E34> >> >>}
\MusicPDFWriteObject{5826}{<< /Type /Annot /Subtype /Link /Border [0 0 0] /H /I /C [1 0 0] /Rect [269.365 503.443 302.259 525.272] /A << /S /GoTo /D <74686D2E422E342E342E32> >> >>}
\MusicPDFWriteObject{5827}{<< /Type /Annot /Subtype /Link /Border [0 0 0] /H /I /C [1 0 0] /Rect [386.259 475.264 396.192 488.122] /A << /S /GoTo /D <48666F6F746E6F74652E313033> >> >>}
\MusicPDFWriteObject{5828}{<< /Type /Annot /Subtype /Link /Border [0 0 0] /H /I /C [1 0 0] /Rect [144.433 424.096 181.701 436.049] /A << /S /GoTo /D <6571756174696F6E2E422E342E342E3134> >> >>}
\MusicPDFWriteObject{5829}{<< /Type /Annot /Subtype /Link /Border [0 0 0] /H /I /C [0 1 0] /Rect [269.963 144.899 291.184 153.767] /A << /S /GoTo /D <636974652E427275> >> >>}
\MusicPDFWriteObject{5830}{<< /Type /Annot /Subtype /Link /Border [0 0 0] /H /I /C [1 0 0] /Rect [458.893 142.409 490.576 154.362] /A << /S /GoTo /D <74686D2E422E342E342E37> >> >>}
\MusicPDFWriteObject{5831}{<< /Type /Annot /Subtype /Link /Border [0 0 0] /H /I /C [1 0 0] /Rect [194.181 129.258 226.406 141.211] /A << /S /GoTo /D <74686D2E442E322E312E31> >> >>}
\MusicPDFWriteObject{5832}{<< /Type /Annot /Subtype /Link /Border [0 0 0] /H /I /C [1 0 0] /Rect [534.862 658.331 544.795 671.189] /A << /S /GoTo /D <48666F6F746E6F74652E313034> >> >>}
\MusicPDFWriteObject{5833}{<< /Type /Annot /Subtype /Link /Border [0 0 0] /H /I /C [1 0 0] /Rect [217.183 542.622 254.451 554.578] /A << /S /GoTo /D <6571756174696F6E2E422E342E342E3137> >> >>}
\MusicPDFWriteObject{5834}{<< /Type /Annot /Subtype /Link /Border [0 0 0] /H /I /C [1 0 0] /Rect [356.974 439.909 394.066 453.07] /A << /S /GoTo /D <6571756174696F6E2E422E342E342E3138> >> >>}
\MusicPDFWriteObject{5835}{<< /Type /Annot /Subtype /Link /Border [0 0 0] /H /I /C [1 0 0] /Rect [490.225 305.993 522.449 317.945] /A << /S /GoTo /D <74686D2E442E332E342E32> >> >>}
\MusicPDFWriteObject{5836}{<< /Type /Annot /Subtype /Link /Border [0 0 0] /H /I /C [1 0 0] /Rect [475.094 292.842 512.362 304.795] /A << /S /GoTo /D <6571756174696F6E2E422E342E342E3137> >> >>}
\MusicPDFWriteObject{5837}{<< /Type /Annot /Subtype /Link /Border [0 0 0] /H /I /C [1 0 0] /Rect [516.954 291.05 526.888 305.705] /A << /S /GoTo /D <48666F6F746E6F74652E313035> >> >>}
\MusicPDFWriteObject{5838}{<< /Type /Annot /Subtype /Link /Border [0 0 0] /H /I /C [1 0 0] /Rect [401.052 613.474 433.276 625.429] /A << /S /GoTo /D <74686D2E442E312E322E31> >> >>}
\MusicPDFWriteObject{5839}{<< /Type /Annot /Subtype /Link /Border [0 0 0] /H /I /C [1 0 0] /Rect [452.569 613.475 489.837 625.427] /A << /S /GoTo /D <6571756174696F6E2E422E342E342E3139> >> >>}
\MusicPDFWriteObject{5840}{<< /Type /Annot /Subtype /Link /Border [0 0 0] /H /I /C [1 0 0] /Rect [185.491 602.814 217.779 611.668] /A << /S /GoTo /D <74686D2E422E342E342E38> >> >>}
\MusicPDFWriteObject{5841}{<< /Type /Annot /Subtype /Link /Border [0 0 0] /H /I /C [1 0 0] /Rect [187.794 582.41 223.527 591.264] /A << /S /GoTo /D <74686D2E422E342E342E37> >> >>}
\MusicPDFWriteObject{5842}{<< /Type /Annot /Subtype /Link /Border [0 0 0] /H /I /C [1 0 0] /Rect [300.782 580.473 333.676 591.264] /A << /S /GoTo /D <74686D2E422E342E332E32> >> >>}
\MusicPDFWriteObject{5843}{<< /Type /Annot /Subtype /Link /Border [0 0 0] /H /I /C [1 0 0] /Rect [391.768 580.473 424.662 591.264] /A << /S /GoTo /D <74686D2E422E342E342E35> >> >>}
\MusicPDFWriteObject{5844}{<< /Type /Annot /Subtype /Link /Border [0 0 0] /H /I /C [1 0 0] /Rect [73.712 503.349 110.98 515.304] /A << /S /GoTo /D <6571756174696F6E2E422E342E332E3133> >> >>}
\MusicPDFWriteObject{5845}{<< /Type /Annot /Subtype /Link /Border [0 0 0] /H /I /C [1 0 0] /Rect [139.422 503.349 176.69 515.304] /A << /S /GoTo /D <6571756174696F6E2E422E342E332E3134> >> >>}
\MusicPDFWriteObject{5846}{<< /Type /Annot /Subtype /Link /Border [0 0 0] /H /I /C [1 0 0] /Rect [304.622 427.025 341.89 438.978] /A << /S /GoTo /D <6571756174696F6E2E422E342E342E3136> >> >>}
\MusicPDFWriteObject{5847}{<< /Type /Annot /Subtype /Link /Border [0 0 0] /H /I /C [1 0 0] /Rect [104.481 414.428 114.414 426.732] /A << /S /GoTo /D <48666F6F746E6F74652E313036> >> >>}
\MusicPDFWriteObject{5848}{<< /Type /Annot /Subtype /Link /Border [0 0 0] /H /I /C [0 1 0] /Rect [176.748 378.069 210.309 386.937] /A << /S /GoTo /D <636974652E5765696C2D65> >> >>}
\MusicPDFWriteObject{5849}{<< /Type /Annot /Subtype /Link /Border [0 0 0] /H /I /C [1 0 0] /Rect [436.938 325.358 470.396 338.133] /A << /S /GoTo /D <74686D2E442E322E312E31> >> >>}
\MusicPDFWriteObject{5850}{<< /Type /Annot /Subtype /Link /Border [0 0 0] /H /I /C [1 0 0] /Rect [298.144 265.413 335.412 277.368] /A << /S /GoTo /D <6571756174696F6E2E422E342E342E3136> >> >>}
\MusicPDFWriteObject{5851}{<< /Type /Annot /Subtype /Link /Border [0 0 0] /H /I /C [1 0 0] /Rect [491.067 265.413 523.355 277.368] /A << /S /GoTo /D <74686D2E422E342E342E39> >> >>}
\MusicPDFWriteObject{5852}{<< /Type /Annot /Subtype /Link /Border [0 0 0] /H /I /C [1 0 0] /Rect [366.532 252.263 403.8 264.215] /A << /S /GoTo /D <6571756174696F6E2E422E342E342E3136> >> >>}
\MusicPDFWriteObject{5853}{<< /Type /Annot /Subtype /Link /Border [0 0 0] /H /I /C [1 0 0] /Rect [325.625 183.006 358.004 198.945] /A << /S /GoTo /D <74686D2E442E312E322E31> >> >>}
\MusicPDFWriteObject{5854}{<< /Type /Annot /Subtype /Link /Border [0 0 0] /H /I /C [1 0 0] /Rect [504.62 158.899 541.888 170.852] /A << /S /GoTo /D <6571756174696F6E2E422E342E342E3232> >> >>}
\MusicPDFWriteObject{5855}{<< /Type /Annot /Subtype /Link /Border [0 0 0] /H /I /C [1 0 0] /Rect [217.295 132.597 254.564 144.55] /A << /S /GoTo /D <6571756174696F6E2E422E342E342E3136> >> >>}
\MusicPDFWriteObject{5856}{<< /Type /Annot /Subtype /Link /Border [0 0 0] /H /I /C [1 0 0] /Rect [428.448 117.759 460.131 131.419] /A << /S /GoTo /D <74686D2E422E342E312E36> >> >>}
\MusicPDFWriteObject{5857}{<< /Type /Annot /Subtype /Link /Border [0 0 0] /H /I /C [1 0 0] /Rect [338.724 82.097 366.463 94.681] /A << /S /GoTo /D <74686D2E422E342E322E32> >> >>}
\MusicPDFWriteObject{5858}{<< /Type /Annot /Subtype /Link /Border [0 0 0] /H /I /C [0 1 0] /Rect [334.599 660.822 355.82 669.69] /A << /S /GoTo /D <636974652E427275> >> >>}
\MusicPDFWriteObject{5859}{<< /Type /Annot /Subtype /Link /Border [0 0 0] /H /I /C [1 0 0] /Rect [511.554 658.332 543.267 670.747] /A << /S /GoTo /D <74686D2E422E342E342E31> >> >>}
\MusicPDFWriteObject{5860}{<< /Type /Annot /Subtype /Link /Border [0 0 0] /H /I /C [0 1 0] /Rect [434.62 621.369 455.841 630.237] /A << /S /GoTo /D <636974652E427275> >> >>}
\MusicPDFWriteObject{5861}{<< /Type /Annot /Subtype /Link /Border [0 0 0] /H /I /C [0 1 0] /Rect [488.436 608.219 513.496 617.087] /A << /S /GoTo /D <636974652E53636861> >> >>}
\MusicPDFWriteObject{5862}{<< /Type /Annot /Subtype /Link /Border [0 0 0] /H /I /C [0 1 0] /Rect [431.088 595.068 452.309 603.936] /A << /S /GoTo /D <636974652E427275> >> >>}
\MusicPDFWriteObject{5863}{<< /Type /Annot /Subtype /Link /Border [0 0 0] /H /I /C [0 1 0] /Rect [512.96 410.546 543.159 419.414] /A << /S /GoTo /D <636974652E43465353> >> >>}
\MusicPDFWriteObject{5864}{<< /Type /Annot /Subtype /Link /Border [0 0 0] /H /I /C [0 1 0] /Rect [497.222 344.792 517.157 353.66] /A << /S /GoTo /D <636974652E484B> >> >>}
\MusicPDFWriteObject{5865}{<< /Type /Annot /Subtype /Link /Border [0 0 0] /H /I /C [0 1 0] /Rect [475.979 331.641 499.026 340.509] /A << /S /GoTo /D <636974652E487531> >> >>}
\MusicPDFWriteObject{5866}{<< /Type /Annot /Subtype /Link /Border [0 0 0] /H /I /C [0 1 0] /Rect [73.757 318.491 96.804 327.358] /A << /S /GoTo /D <636974652E487532> >> >>}
\MusicPDFWriteObject{5867}{<< /Type /Annot /Subtype /Link /Border [0 0 0] /H /I /C [1 0 0] /Rect [112.46 572.381 144.536 585.189] /A << /S /GoTo /D <74686D2E422E352E322E31> >> >>}
\MusicPDFWriteObject{5868}{<< /Type /Annot /Subtype /Link /Border [0 0 0] /H /I /C [0 1 0] /Rect [487.161 528.271 507.096 537.139] /A << /S /GoTo /D <636974652E484B> >> >>}
\MusicPDFWriteObject{5869}{<< /Type /Annot /Subtype /Link /Border [0 0 0] /H /I /C [0 1 0] /Rect [375.53 515.12 395.464 523.988] /A << /S /GoTo /D <636974652E484B> >> >>}
\MusicPDFWriteObject{5870}{<< /Type /Annot /Subtype /Link /Border [0 0 0] /H /I /C [0 1 0] /Rect [183.51 475.668 206.557 484.536] /A << /S /GoTo /D <636974652E487531> >> >>}
\MusicPDFWriteObject{5871}{<< /Type /Annot /Subtype /Link /Border [0 0 0] /H /I /C [0 1 0] /Rect [380.458 638.979 397.612 647.847] /A << /S /GoTo /D <636974652E466F6C> >> >>}
\MusicPDFWriteObject{5872}{<< /Type /Annot /Subtype /Link /Border [0 0 0] /H /I /C [0 1 0] /Rect [76.611 612.677 93.766 621.545] /A << /S /GoTo /D <636974652E466F6C> >> >>}
\MusicPDFWriteObject{5873}{<< /Type /Annot /Subtype /Link /Border [0 0 0] /H /I /C [1 0 0] /Rect [377.673 312.139 409.9 322.93] /A << /S /GoTo /D <74686D2E422E352E322E32> >> >>}
\MusicPDFWriteObject{5874}{<< /Type /Annot /Subtype /Link /Border [0 0 0] /H /I /C [0 1 0] /Rect [385.778 278.73 408.826 287.597] /A << /S /GoTo /D <636974652E487531> >> >>}
\MusicPDFWriteObject{5875}{<< /Type /Annot /Subtype /Link /Border [0 0 0] /H /I /C [1 0 0] /Rect [187.794 253.989 223.527 262.843] /A << /S /GoTo /D <74686D2E422E352E322E31> >> >>}
\MusicPDFWriteObject{5876}{<< /Type /Annot /Subtype /Link /Border [0 0 0] /H /I /C [0 1 0] /Rect [524.857 645.503 543.159 654.37] /A << /S /GoTo /D <636974652E4445> >> >>}
\MusicPDFWriteObject{5877}{<< /Type /Annot /Subtype /Link /Border [0 0 0] /H /I /C [1 0 0] /Rect [208.221 629.862 240.445 641.815] /A << /S /GoTo /D <74686D2E442E312E342E36> >> >>}
\MusicPDFWriteObject{5878}{<< /Type /Annot /Subtype /Link /Border [0 0 0] /H /I /C [1 0 0] /Rect [175.488 590.41 207.776 602.362] /A << /S /GoTo /D <6571756174696F6E2E422E352E322E31> >> >>}
\MusicPDFWriteObject{5879}{<< /Type /Annot /Subtype /Link /Border [0 0 0] /H /I /C [1 0 0] /Rect [140.069 511.405 172.327 524.27] /A << /S /GoTo /D <74686D2E422E352E322E32> >> >>}
\MusicPDFWriteObject{5880}{<< /Type /Annot /Subtype /Link /Border [0 0 0] /H /I /C [1 0 0] /Rect [108.58 415.307 140.868 427.259] /A << /S /GoTo /D <6571756174696F6E2E422E352E322E31> >> >>}
\MusicPDFWriteObject{5881}{<< /Type /Annot /Subtype /Link /Border [0 0 0] /H /I /C [1 0 0] /Rect [509.6 281.089 541.888 293.042] /A << /S /GoTo /D <6571756174696F6E2E422E352E322E31> >> >>}
\MusicPDFWriteObject{5882}{<< /Type /Annot /Subtype /Link /Border [0 0 0] /H /I /C [1 0 0] /Rect [210.623 254.739 242.911 266.692] /A << /S /GoTo /D <6571756174696F6E2E422E352E322E31> >> >>}
\MusicPDFWriteObject{5883}{<< /Type /Annot /Subtype /Link /Border [0 0 0] /H /I /C [1 0 0] /Rect [186.899 215.287 219.187 227.239] /A << /S /GoTo /D <6571756174696F6E2E422E352E322E31> >> >>}
\MusicPDFWriteObject{5884}{<< /Type /Annot /Subtype /Link /Border [0 0 0] /H /I /C [1 0 0] /Rect [329.654 645.181 361.942 657.134] /A << /S /GoTo /D <6571756174696F6E2E422E352E322E31> >> >>}
\MusicPDFWriteObject{5885}{<< /Type /Annot /Subtype /Link /Border [0 0 0] /H /I /C [1 0 0] /Rect [269.422 618.879 301.71 630.832] /A << /S /GoTo /D <6571756174696F6E2E422E352E322E31> >> >>}
\MusicPDFWriteObject{5886}{<< /Type /Annot /Subtype /Link /Border [0 0 0] /H /I /C [1 0 0] /Rect [335.077 295.631 367.365 307.584] /A << /S /GoTo /D <6571756174696F6E2E422E352E322E31> >> >>}
\MusicPDFWriteObject{5887}{<< /Type /Annot /Subtype /Link /Border [0 0 0] /H /I /C [1 0 0] /Rect [244.89 229.877 277.602 241.832] /A << /S /GoTo /D <74686D2E422E352E322E32> >> >>}
\MusicPDFWriteObject{5888}{<< /Type /Annot /Subtype /Link /Border [0 0 0] /H /I /C [1 0 0] /Rect [509.6 229.877 541.888 241.83] /A << /S /GoTo /D <6571756174696F6E2E422E352E322E31> >> >>}
\MusicPDFWriteObject{5889}{<< /Type /Annot /Subtype /Link /Border [0 0 0] /H /I /C [1 0 0] /Rect [173.741 381.353 211.009 393.306] /A << /S /GoTo /D <6571756174696F6E2E422E352E322E3133> >> >>}
\MusicPDFWriteObject{5890}{<< /Type /Annot /Subtype /Link /Border [0 0 0] /H /I /C [1 0 0] /Rect [471.313 381.353 508.582 393.306] /A << /S /GoTo /D <6571756174696F6E2E422E352E322E3133> >> >>}
\MusicPDFWriteObject{5891}{<< /Type /Annot /Subtype /Link /Border [0 0 0] /H /I /C [1 0 0] /Rect [288.053 267.378 320.895 283.241] /A << /S /GoTo /D <74686D2E442E312E322E37> >> >>}
\MusicPDFWriteObject{5892}{<< /Type /Annot /Subtype /Link /Border [0 0 0] /H /I /C [1 0 0] /Rect [73.712 130.616 110.98 142.569] /A << /S /GoTo /D <6571756174696F6E2E422E352E322E3135> >> >>}
\MusicPDFWriteObject{5893}{<< /Type /Annot /Subtype /Link /Border [0 0 0] /H /I /C [1 0 0] /Rect [233.296 130.616 265.585 142.569] /A << /S /GoTo /D <6571756174696F6E2E422E352E322E31> >> >>}
\MusicPDFWriteObject{5894}{<< /Type /Annot /Subtype /Link /Border [0 0 0] /H /I /C [1 0 0] /Rect [336.412 130.616 368.7 142.569] /A << /S /GoTo /D <74686D2E422E352E322E31> >> >>}
\MusicPDFWriteObject{5895}{<< /Type /Annot /Subtype /Link /Border [0 0 0] /H /I /C [1 0 0] /Rect [387.024 109.451 419.222 120.241] /A << /S /GoTo /D <74686D2E422E352E322E31> >> >>}
\MusicPDFWriteObject{5896}{<< /Type /Annot /Subtype /Link /Border [0 0 0] /H /I /C [1 0 0] /Rect [483.814 109.451 516.011 120.241] /A << /S /GoTo /D <74686D2E422E352E322E34> >> >>}
\MusicPDFWriteObject{5897}{<< /Type /Annot /Subtype /Link /Border [0 0 0] /H /I /C [0 1 0] /Rect [483.147 98.237 506.195 107.105] /A << /S /GoTo /D <636974652E487532> >> >>}
\MusicPDFWriteObject{5898}{<< /Type /Annot /Subtype /Link /Border [0 0 0] /H /I /C [0 1 0] /Rect [391.471 371.929 412.907 380.796] /A << /S /GoTo /D <636974652E426131> >> >>}
\MusicPDFWriteObject{5899}{<< /Type /Annot /Subtype /Link /Border [0 0 0] /H /I /C [0 1 0] /Rect [493.867 358.778 515.303 367.646] /A << /S /GoTo /D <636974652E426131> >> >>}
\MusicPDFWriteObject{5900}{<< /Type /Annot /Subtype /Link /Border [0 0 0] /H /I /C [0 1 0] /Rect [155.196 345.627 176.632 354.495] /A << /S /GoTo /D <636974652E426131> >> >>}
\MusicPDFWriteObject{5901}{<< /Type /Annot /Subtype /Link /Border [0 0 0] /H /I /C [1 0 0] /Rect [371.682 155.584 404.273 166.375] /A << /S /GoTo /D <74686D2E422E352E332E32> >> >>}
\MusicPDFWriteObject{5902}{<< /Type /Annot /Subtype /Link /Border [0 0 0] /H /I /C [1 0 0] /Rect [136.243 128.729 166.456 140.682] /A << /S /GoTo /D <74686D2E372E312E312E31> >> >>}
\MusicPDFWriteObject{5903}{<< /Type /Annot /Subtype /Link /Border [0 0 0] /H /I /C [0 1 0] /Rect [256.789 131.22 278.225 140.087] /A << /S /GoTo /D <636974652E426131> >> >>}
\MusicPDFWriteObject{5904}{<< /Type /Annot /Subtype /Link /Border [0 0 0] /H /I /C [1 0 0] /Rect [401.046 527.013 410.979 539.317] /A << /S /GoTo /D <48666F6F746E6F74652E313037> >> >>}
\MusicPDFWriteObject{5905}{<< /Type /Annot /Subtype /Link /Border [0 0 0] /H /I /C [1 0 0] /Rect [111.536 475.928 143.824 487.881] /A << /S /GoTo /D <6571756174696F6E2E422E352E332E39> >> >>}
\MusicPDFWriteObject{5906}{<< /Type /Annot /Subtype /Link /Border [0 0 0] /H /I /C [1 0 0] /Rect [118.817 439.5 156.085 451.453] /A << /S /GoTo /D <6571756174696F6E2E422E352E332E3130> >> >>}
\MusicPDFWriteObject{5907}{<< /Type /Annot /Subtype /Link /Border [0 0 0] /H /I /C [1 0 0] /Rect [198.969 439.5 236.237 451.453] /A << /S /GoTo /D <6571756174696F6E2E422E352E332E3131> >> >>}
\MusicPDFWriteObject{5908}{<< /Type /Annot /Subtype /Link /Border [0 0 0] /H /I /C [1 0 0] /Rect [495.285 420.925 527.573 431.716] /A << /S /GoTo /D <74686D2E422E352E332E33> >> >>}
\MusicPDFWriteObject{5909}{<< /Type /Annot /Subtype /Link /Border [0 0 0] /H /I /C [1 0 0] /Rect [203.197 403.734 238.929 412.588] /A << /S /GoTo /D <74686D2E422E352E332E33> >> >>}
\MusicPDFWriteObject{5910}{<< /Type /Annot /Subtype /Link /Border [0 0 0] /H /I /C [1 0 0] /Rect [439.252 109.624 471.54 121.58] /A << /S /GoTo /D <74686D2E422E352E322E34> >> >>}
\MusicPDFWriteObject{5911}{<< /Type /Annot /Subtype /Link /Border [0 0 0] /H /I /C [1 0 0] /Rect [362.154 621.028 399.422 632.98] /A << /S /GoTo /D <6571756174696F6E2E422E352E332E3134> >> >>}
\MusicPDFWriteObject{5912}{<< /Type /Annot /Subtype /Link /Border [0 0 0] /H /I /C [1 0 0] /Rect [147.508 606.85 179.19 620.469] /A << /S /GoTo /D <74686D2E422E352E332E35> >> >>}
\MusicPDFWriteObject{5913}{<< /Type /Annot /Subtype /Link /Border [0 0 0] /H /I /C [1 0 0] /Rect [122.88 568.425 160.148 580.377] /A << /S /GoTo /D <6571756174696F6E2E422E352E332E3135> >> >>}
\MusicPDFWriteObject{5914}{<< /Type /Annot /Subtype /Link /Border [0 0 0] /H /I /C [1 0 0] /Rect [191.91 568.425 229.178 580.377] /A << /S /GoTo /D <6571756174696F6E2E422E352E332E3134> >> >>}
\MusicPDFWriteObject{5915}{<< /Type /Annot /Subtype /Link /Border [0 0 0] /H /I /C [1 0 0] /Rect [71.004 542.123 102.686 554.076] /A << /S /GoTo /D <74686D2E422E352E322E31> >> >>}
\MusicPDFWriteObject{5916}{<< /Type /Annot /Subtype /Link /Border [0 0 0] /H /I /C [1 0 0] /Rect [202.405 542.123 234.088 554.076] /A << /S /GoTo /D <74686D2E422E352E322E34> >> >>}
\MusicPDFWriteObject{5917}{<< /Type /Annot /Subtype /Link /Border [0 0 0] /H /I /C [0 1 0] /Rect [207.73 296.198 228.232 305.066] /A << /S /GoTo /D <636974652E4D6974> >> >>}
\MusicPDFWriteObject{5918}{<< /Type /Annot /Subtype /Link /Border [0 0 0] /H /I /C [1 0 0] /Rect [141.794 201.506 173.628 212.297] /A << /S /GoTo /D <74686D2E422E352E332E32> >> >>}
\MusicPDFWriteObject{5919}{<< /Type /Annot /Subtype /Link /Border [0 0 0] /H /I /C [1 0 0] /Rect [283.278 201.506 315.112 212.297] /A << /S /GoTo /D <74686D2E422E352E332E33> >> >>}
\MusicPDFWriteObject{5920}{<< /Type /Annot /Subtype /Link /Border [0 0 0] /H /I /C [1 0 0] /Rect [409.203 175.205 442.097 185.995] /A << /S /GoTo /D <74686D2E422E352E332E33> >> >>}
\MusicPDFWriteObject{5921}{<< /Type /Annot /Subtype /Link /Border [0 0 0] /H /I /C [1 0 0] /Rect [124.332 135.198 156.438 147.153] /A << /S /GoTo /D <74686D2E422E352E332E33> >> >>}
\MusicPDFWriteObject{5922}{<< /Type /Annot /Subtype /Link /Border [0 0 0] /H /I /C [1 0 0] /Rect [124.623 83.149 156.911 93.94] /A << /S /GoTo /D <74686D2E422E352E332E32> >> >>}
\MusicPDFWriteObject{5923}{<< /Type /Annot /Subtype /Link /Border [0 0 0] /H /I /C [0 1 0] /Rect [73.815 414.618 89.129 423.486] /A << /S /GoTo /D <636974652E494B> >> >>}
\MusicPDFWriteObject{5924}{<< /Type /Annot /Subtype /Link /Border [0 0 0] /H /I /C [1 0 0] /Rect [230.672 412.128 263.445 424.081] /A << /S /GoTo /D <74686D2E422E362E312E31> >> >>}
\MusicPDFWriteObject{5925}{<< /Type /Annot /Subtype /Link /Border [0 0 0] /H /I /C [0 1 0] /Rect [362.946 414.618 378.26 423.486] /A << /S /GoTo /D <636974652E494B> >> >>}
\MusicPDFWriteObject{5926}{<< /Type /Annot /Subtype /Link /Border [0 0 0] /H /I /C [0 1 0] /Rect [73.715 388.316 89.029 397.184] /A << /S /GoTo /D <636974652E494B> >> >>}
\MusicPDFWriteObject{5927}{<< /Type /Annot /Subtype /Link /Border [0 0 0] /H /I /C [0 1 0] /Rect [161.443 151.581 178.915 160.449] /A << /S /GoTo /D <636974652E4346> >> >>}
\MusicPDFWriteObject{5928}{<< /Type /Annot /Subtype /Link /Border [0 0 0] /H /I /C [1 0 0] /Rect [452.434 82.596 484.722 94.549] /A << /S /GoTo /D <6571756174696F6E2E422E362E312E32> >> >>}
\MusicPDFWriteObject{5929}{<< /Type /Annot /Subtype /Link /Border [0 0 0] /H /I /C [0 1 0] /Rect [241.665 396.009 267.804 404.877] /A << /S /GoTo /D <636974652E4D5661> >> >>}
\MusicPDFWriteObject{5930}{<< /Type /Annot /Subtype /Link /Border [0 0 0] /H /I /C [1 0 0] /Rect [509.6 393.519 541.888 405.472] /A << /S /GoTo /D <6571756174696F6E2E422E362E312E37> >> >>}
\MusicPDFWriteObject{5931}{<< /Type /Annot /Subtype /Link /Border [0 0 0] /H /I /C [0 1 0] /Rect [427.798 369.708 453.937 378.575] /A << /S /GoTo /D <636974652E4D5661> >> >>}
\MusicPDFWriteObject{5932}{<< /Type /Annot /Subtype /Link /Border [0 0 0] /H /I /C [1 0 0] /Rect [126.865 354.067 159.153 366.019] /A << /S /GoTo /D <6571756174696F6E2E422E362E312E38> >> >>}
\MusicPDFWriteObject{5933}{<< /Type /Annot /Subtype /Link /Border [0 0 0] /H /I /C [0 1 0] /Rect [365.456 337.071 385.654 345.939] /A << /S /GoTo /D <636974652E5357> >> >>}
\MusicPDFWriteObject{5934}{<< /Type /Annot /Subtype /Link /Border [0 0 0] /H /I /C [0 1 0] /Rect [527.582 337.071 543.27 345.939] /A << /S /GoTo /D <636974652E5432> >> >>}
\MusicPDFWriteObject{5935}{<< /Type /Annot /Subtype /Link /Border [0 0 0] /H /I /C [0 1 0] /Rect [469.968 321.983 500.485 332.788] /A << /S /GoTo /D <636974652E496E676833> >> >>}
\MusicPDFWriteObject{5936}{<< /Type /Annot /Subtype /Link /Border [0 0 0] /H /I /C [1 0 0] /Rect [227.497 460.571 264.06 472.527] /A << /S /GoTo /D <6571756174696F6E2E422E362E312E3134> >> >>}
\MusicPDFWriteObject{5937}{<< /Type /Annot /Subtype /Link /Border [0 0 0] /H /I /C [1 0 0] /Rect [240.116 330.529 272.798 342.482] /A << /S /GoTo /D <74686D2E422E312E312E33> >> >>}
\MusicPDFWriteObject{5938}{<< /Type /Annot /Subtype /Link /Border [0 0 0] /H /I /C [0 1 0] /Rect [399.766 331.083 430.284 341.887] /A << /S /GoTo /D <636974652E496E676833> >> >>}
\MusicPDFWriteObject{5939}{<< /Type /Annot /Subtype /Link /Border [0 0 0] /H /I /C [1 0 0] /Rect [73.712 291.077 110.98 303.03] /A << /S /GoTo /D <6571756174696F6E2E422E362E312E3132> >> >>}
\MusicPDFWriteObject{5940}{<< /Type /Annot /Subtype /Link /Border [0 0 0] /H /I /C [0 1 0] /Rect [512.675 641.685 543.192 652.49] /A << /S /GoTo /D <636974652E496E676833> >> >>}
\MusicPDFWriteObject{5941}{<< /Type /Annot /Subtype /Link /Border [0 0 0] /H /I /C [0 1 0] /Rect [160.757 628.535 191.275 639.339] /A << /S /GoTo /D <636974652E496E676833> >> >>}
\MusicPDFWriteObject{5942}{<< /Type /Annot /Subtype /Link /Border [0 0 0] /H /I /C [1 0 0] /Rect [481.353 627.981 518.621 639.934] /A << /S /GoTo /D <6571756174696F6E2E422E362E312E3137> >> >>}
\MusicPDFWriteObject{5943}{<< /Type /Annot /Subtype /Link /Border [0 0 0] /H /I /C [0 1 0] /Rect [197.331 589.082 227.848 599.887] /A << /S /GoTo /D <636974652E496E676833> >> >>}
\MusicPDFWriteObject{5944}{<< /Type /Annot /Subtype /Link /Border [0 0 0] /H /I /C [1 0 0] /Rect [173.377 549.076 206.21 561.031] /A << /S /GoTo /D <74686D2E422E312E312E33> >> >>}
\MusicPDFWriteObject{5945}{<< /Type /Annot /Subtype /Link /Border [0 0 0] /H /I /C [0 1 0] /Rect [107.089 379.801 137.607 390.606] /A << /S /GoTo /D <636974652E496E676833> >> >>}
\MusicPDFWriteObject{5946}{<< /Type /Annot /Subtype /Link /Border [0 0 0] /H /I /C [1 0 0] /Rect [342.402 331.568 379.671 343.521] /A << /S /GoTo /D <6571756174696F6E2E422E362E312E3230> >> >>}
\MusicPDFWriteObject{5947}{<< /Type /Annot /Subtype /Link /Border [0 0 0] /H /I /C [1 0 0] /Rect [230.771 292.116 268.039 304.068] /A << /S /GoTo /D <6571756174696F6E2E422E362E312E3230> >> >>}
\MusicPDFWriteObject{5948}{<< /Type /Annot /Subtype /Link /Border [0 0 0] /H /I /C [1 0 0] /Rect [204.947 278.219 242.216 290.172] /A << /S /GoTo /D <6571756174696F6E2E422E362E312E3130> >> >>}
\MusicPDFWriteObject{5949}{<< /Type /Annot /Subtype /Link /Border [0 0 0] /H /I /C [1 0 0] /Rect [266.338 107.21 299.232 121.49] /A << /S /GoTo /D <74686D2E422E362E312E31> >> >>}
\MusicPDFWriteObject{5950}{<< /Type /Annot /Subtype /Link /Border [0 0 0] /H /I /C [1 0 0] /Rect [364.022 107.21 396.916 121.49] /A << /S /GoTo /D <74686D2E422E362E312E33> >> >>}
\MusicPDFWriteObject{5951}{<< /Type /Annot /Subtype /Link /Border [0 0 0] /H /I /C [1 0 0] /Rect [73.712 82.596 110.98 94.549] /A << /S /GoTo /D <6571756174696F6E2E422E362E312E3231> >> >>}
\MusicPDFWriteObject{5952}{<< /Type /Annot /Subtype /Link /Border [0 0 0] /H /I /C [0 1 0] /Rect [482.166 522.072 499.639 530.94] /A << /S /GoTo /D <636974652E4346> >> >>}
\MusicPDFWriteObject{5953}{<< /Type /Annot /Subtype /Link /Border [0 0 0] /H /I /C [1 0 0] /Rect [361.926 470.557 393.608 481.348] /A << /S /GoTo /D <74686D2E422E312E312E33> >> >>}
\MusicPDFWriteObject{5954}{<< /Type /Annot /Subtype /Link /Border [0 0 0] /H /I /C [1 0 0] /Rect [149.989 406.976 182.883 417.766] /A << /S /GoTo /D <74686D2E422E362E312E35> >> >>}
\MusicPDFWriteObject{5955}{<< /Type /Annot /Subtype /Link /Border [0 0 0] /H /I /C [1 0 0] /Rect [131.817 332.487 169.402 344.738] /A << /S /GoTo /D <6571756174696F6E2E422E362E312E3130> >> >>}
\MusicPDFWriteObject{5956}{<< /Type /Annot /Subtype /Link /Border [0 0 0] /H /I /C [0 1 0] /Rect [453.672 333.338 489.121 344.143] /A << /S /GoTo /D <636974652E4C616E672D63> >> >>}
\MusicPDFWriteObject{5957}{<< /Type /Annot /Subtype /Link /Border [0 0 0] /H /I /C [1 0 0] /Rect [108.339 233.474 145.607 245.427] /A << /S /GoTo /D <6571756174696F6E2E422E362E312E3234> >> >>}
\MusicPDFWriteObject{5958}{<< /Type /Annot /Subtype /Link /Border [0 0 0] /H /I /C [1 0 0] /Rect [173.737 233.474 211.005 245.427] /A << /S /GoTo /D <6571756174696F6E2E422E362E312E3236> >> >>}
\MusicPDFWriteObject{5959}{<< /Type /Annot /Subtype /Link /Border [0 0 0] /H /I /C [1 0 0] /Rect [109.112 182.621 146.381 194.573] /A << /S /GoTo /D <6571756174696F6E2E422E362E312E3233> >> >>}
\MusicPDFWriteObject{5960}{<< /Type /Annot /Subtype /Link /Border [0 0 0] /H /I /C [1 0 0] /Rect [175.715 182.621 212.983 194.573] /A << /S /GoTo /D <6571756174696F6E2E422E362E312E3234> >> >>}
\MusicPDFWriteObject{5961}{<< /Type /Annot /Subtype /Link /Border [0 0 0] /H /I /C [1 0 0] /Rect [212.456 639.374 249.724 651.329] /A << /S /GoTo /D <6571756174696F6E2E422E362E312E3239> >> >>}
\MusicPDFWriteObject{5962}{<< /Type /Annot /Subtype /Link /Border [0 0 0] /H /I /C [1 0 0] /Rect [278.166 639.374 315.434 651.329] /A << /S /GoTo /D <6571756174696F6E2E422E362E312E3330> >> >>}
\MusicPDFWriteObject{5963}{<< /Type /Annot /Subtype /Link /Border [0 0 0] /H /I /C [1 0 0] /Rect [215.887 440.924 253.155 452.877] /A << /S /GoTo /D <6571756174696F6E2E422E362E312E3232> >> >>}
\MusicPDFWriteObject{5964}{<< /Type /Annot /Subtype /Link /Border [0 0 0] /H /I /C [0 1 0] /Rect [506.437 392.899 543.159 403.704] /A << /S /GoTo /D <636974652E4C616E672D62> >> >>}
\MusicPDFWriteObject{5965}{<< /Type /Annot /Subtype /Link /Border [0 0 0] /H /I /C [0 1 0] /Rect [307.15 355.384 322.838 364.252] /A << /S /GoTo /D <636974652E5432> >> >>}
\MusicPDFWriteObject{5966}{<< /Type /Annot /Subtype /Link /Border [0 0 0] /H /I /C [1 0 0] /Rect [265.247 339.417 295.46 351.37] /A << /S /GoTo /D <6571756174696F6E2E362E332E312E36> >> >>}
\MusicPDFWriteObject{5967}{<< /Type /Annot /Subtype /Link /Border [0 0 0] /H /I /C [1 0 0] /Rect [370.276 339.417 400.8 353.89] /A << /S /GoTo /D <7465726D696E6F6C6F67792E362E332E312E31> >> >>}
\MusicPDFWriteObject{5968}{<< /Type /Annot /Subtype /Link /Border [0 0 0] /H /I /C [1 0 0] /Rect [124.623 608.293 156.911 619.998] /A << /S /GoTo /D <74686D2E422E312E312E33> >> >>}
\MusicPDFWriteObject{5969}{<< /Type /Annot /Subtype /Link /Border [0 0 0] /H /I /C [1 0 0] /Rect [276.275 559.075 314.249 571.03] /A << /S /GoTo /D <6571756174696F6E2E422E362E312E3337> >> >>}
\MusicPDFWriteObject{5970}{<< /Type /Annot /Subtype /Link /Border [0 0 0] /H /I /C [0 1 0] /Rect [344.133 561.566 359.821 570.433] /A << /S /GoTo /D <636974652E5432> >> >>}
\MusicPDFWriteObject{5971}{<< /Type /Annot /Subtype /Link /Border [0 0 0] /H /I /C [0 1 0] /Rect [292.426 443.262 308.114 452.13] /A << /S /GoTo /D <636974652E5432> >> >>}
\MusicPDFWriteObject{5972}{<< /Type /Annot /Subtype /Link /Border [0 0 0] /H /I /C [0 1 0] /Rect [399.641 430.111 415.329 438.979] /A << /S /GoTo /D <636974652E5432> >> >>}
\MusicPDFWriteObject{5973}{<< /Type /Annot /Subtype /Link /Border [0 0 0] /H /I /C [0 1 0] /Rect [359.267 264.285 395.989 275.09] /A << /S /GoTo /D <636974652E4C616E672D62> >> >>}
\MusicPDFWriteObject{5974}{<< /Type /Annot /Subtype /Link /Border [0 0 0] /H /I /C [1 0 0] /Rect [373.642 141.578 410.91 153.531] /A << /S /GoTo /D <6571756174696F6E2E422E362E312E3432> >> >>}
\MusicPDFWriteObject{5975}{<< /Type /Annot /Subtype /Link /Border [0 0 0] /H /I /C [1 0 0] /Rect [170.486 684.633 207.754 696.586] /A << /S /GoTo /D <6571756174696F6E2E422E362E312E3432> >> >>}
\MusicPDFWriteObject{5976}{<< /Type /Annot /Subtype /Link /Border [0 0 0] /H /I /C [1 0 0] /Rect [264.366 513.275 301.635 525.227] /A << /S /GoTo /D <6571756174696F6E2E422E362E312E3434> >> >>}
\MusicPDFWriteObject{5977}{<< /Type /Annot /Subtype /Link /Border [0 0 0] /H /I /C [0 1 0] /Rect [386.735 515.765 402.423 524.633] /A << /S /GoTo /D <636974652E5432> >> >>}
\MusicPDFWriteObject{5978}{<< /Type /Annot /Subtype /Link /Border [0 0 0] /H /I /C [1 0 0] /Rect [119.975 500.124 157.243 512.077] /A << /S /GoTo /D <6571756174696F6E2E422E362E312E3335> >> >>}
\MusicPDFWriteObject{5979}{<< /Type /Annot /Subtype /Link /Border [0 0 0] /H /I /C [1 0 0] /Rect [185.684 500.124 222.953 512.077] /A << /S /GoTo /D <6571756174696F6E2E422E362E312E3336> >> >>}
\MusicPDFWriteObject{5980}{<< /Type /Annot /Subtype /Link /Border [0 0 0] /H /I /C [1 0 0] /Rect [284.628 368.189 317.522 380.144] /A << /S /GoTo /D <74686D2E422E362E312E37> >> >>}
\MusicPDFWriteObject{5981}{<< /Type /Annot /Subtype /Link /Border [0 0 0] /H /I /C [1 0 0] /Rect [93.86 355.039 103.793 367.896] /A << /S /GoTo /D <48666F6F746E6F74652E313038> >> >>}
\MusicPDFWriteObject{5982}{<< /Type /Annot /Subtype /Link /Border [0 0 0] /H /I /C [1 0 0] /Rect [285.459 131.954 322.022 143.909] /A << /S /GoTo /D <6571756174696F6E2E422E362E312E3133> >> >>}
\MusicPDFWriteObject{5983}{<< /Type /Annot /Subtype /Link /Border [0 0 0] /H /I /C [0 1 0] /Rect [491.119 687.123 506.807 695.991] /A << /S /GoTo /D <636974652E5432> >> >>}
\MusicPDFWriteObject{5984}{<< /Type /Annot /Subtype /Link /Border [0 0 0] /H /I /C [1 0 0] /Rect [127.03 357.999 159.925 369.952] /A << /S /GoTo /D <74686D2E422E372E312E31> >> >>}
\MusicPDFWriteObject{5985}{<< /Type /Annot /Subtype /Link /Border [0 0 0] /H /I /C [0 1 0] /Rect [178.017 360.489 193.705 369.357] /A << /S /GoTo /D <636974652E5432> >> >>}
\MusicPDFWriteObject{5986}{<< /Type /Annot /Subtype /Link /Border [0 0 0] /H /I /C [0 1 0] /Rect [110.802 347.338 126.49 356.206] /A << /S /GoTo /D <636974652E5431> >> >>}
\MusicPDFWriteObject{5987}{<< /Type /Annot /Subtype /Link /Border [0 0 0] /H /I /C [1 0 0] /Rect [176.323 314.396 212.056 323.25] /A << /S /GoTo /D <74686D2E422E372E312E31> >> >>}
\MusicPDFWriteObject{5988}{<< /Type /Annot /Subtype /Link /Border [0 0 0] /H /I /C [1 0 0] /Rect [157.584 215.06 189.873 227.013] /A << /S /GoTo /D <6571756174696F6E2E422E372E312E36> >> >>}
\MusicPDFWriteObject{5989}{<< /Type /Annot /Subtype /Link /Border [0 0 0] /H /I /C [0 1 0] /Rect [250.705 215.614 269.906 226.418] /A << /S /GoTo /D <636974652E5151> >> >>}
\MusicPDFWriteObject{5990}{<< /Type /Annot /Subtype /Link /Border [0 0 0] /H /I /C [1 0 0] /Rect [408.195 188.759 440.483 200.712] /A << /S /GoTo /D <6571756174696F6E2E422E372E312E36> >> >>}
\MusicPDFWriteObject{5991}{<< /Type /Annot /Subtype /Link /Border [0 0 0] /H /I /C [1 0 0] /Rect [115.247 175.608 147.535 187.561] /A << /S /GoTo /D <6571756174696F6E2E422E372E312E36> >> >>}
\MusicPDFWriteObject{5992}{<< /Type /Annot /Subtype /Link /Border [0 0 0] /H /I /C [0 1 0] /Rect [165.738 176.162 184.939 186.966] /A << /S /GoTo /D <636974652E5151> >> >>}
\MusicPDFWriteObject{5993}{<< /Type /Annot /Subtype /Link /Border [0 0 0] /H /I /C [1 0 0] /Rect [476.849 135.199 509.137 147.152] /A << /S /GoTo /D <6571756174696F6E2E422E372E312E33> >> >>}
\MusicPDFWriteObject{5994}{<< /Type /Annot /Subtype /Link /Border [0 0 0] /H /I /C [1 0 0] /Rect [86.497 108.897 118.786 120.85] /A << /S /GoTo /D <6571756174696F6E2E422E372E312E34> >> >>}
\MusicPDFWriteObject{5995}{<< /Type /Annot /Subtype /Link /Border [0 0 0] /H /I /C [0 1 0] /Rect [439.078 98.237 454.766 107.105] /A << /S /GoTo /D <636974652E5432> >> >>}
\MusicPDFWriteObject{5996}{<< /Type /Annot /Subtype /Link /Border [0 0 0] /H /I /C [1 0 0] /Rect [444.25 684.633 476.659 696.588] /A << /S /GoTo /D <6571756174696F6E2E422E372E312E34> >> >>}
\MusicPDFWriteObject{5997}{<< /Type /Annot /Subtype /Link /Border [0 0 0] /H /I /C [1 0 0] /Rect [86.497 473.687 118.786 485.64] /A << /S /GoTo /D <6571756174696F6E2E422E372E312E34> >> >>}
\MusicPDFWriteObject{5998}{<< /Type /Annot /Subtype /Link /Border [0 0 0] /H /I /C [1 0 0] /Rect [379.818 460.537 412.106 472.489] /A << /S /GoTo /D <6571756174696F6E2E422E372E312E34> >> >>}
\MusicPDFWriteObject{5999}{<< /Type /Annot /Subtype /Link /Border [0 0 0] /H /I /C [0 1 0] /Rect [208.433 449.876 224.12 458.744] /A << /S /GoTo /D <636974652E5432> >> >>}
\MusicPDFWriteObject{6000}{<< /Type /Annot /Subtype /Link /Border [0 0 0] /H /I /C [1 0 0] /Rect [354.854 373.787 387.143 385.739] /A << /S /GoTo /D <6571756174696F6E2E422E372E312E34> >> >>}
\MusicPDFWriteObject{6001}{<< /Type /Annot /Subtype /Link /Border [0 0 0] /H /I /C [1 0 0] /Rect [421.534 373.787 453.822 385.739] /A << /S /GoTo /D <6571756174696F6E2E422E372E312E37> >> >>}
\MusicPDFWriteObject{6002}{<< /Type /Annot /Subtype /Link /Border [0 0 0] /H /I /C [1 0 0] /Rect [493.633 373.787 525.921 385.739] /A << /S /GoTo /D <6571756174696F6E2E422E372E312E35> >> >>}
\MusicPDFWriteObject{6003}{<< /Type /Annot /Subtype /Link /Border [0 0 0] /H /I /C [0 1 0] /Rect [138.788 258.822 154.476 267.69] /A << /S /GoTo /D <636974652E5432> >> >>}
\MusicPDFWriteObject{6004}{<< /Type /Annot /Subtype /Link /Border [0 0 0] /H /I /C [1 0 0] /Rect [375.18 595.924 408.074 610.121] /A << /S /GoTo /D <74686D2E422E372E312E31> >> >>}
\MusicPDFWriteObject{6005}{<< /Type /Annot /Subtype /Link /Border [0 0 0] /H /I /C [1 0 0] /Rect [484.833 598.166 517.122 610.119] /A << /S /GoTo /D <6571756174696F6E2E422E372E312E33> >> >>}
\MusicPDFWriteObject{6006}{<< /Type /Annot /Subtype /Link /Border [0 0 0] /H /I /C [1 0 0] /Rect [193.21 584.965 225.498 596.92] /A << /S /GoTo /D <74686D2E422E362E312E34> >> >>}
\MusicPDFWriteObject{6007}{<< /Type /Annot /Subtype /Link /Border [0 0 0] /H /I /C [1 0 0] /Rect [374.218 518.931 411.486 530.884] /A << /S /GoTo /D <6571756174696F6E2E422E372E312E3131> >> >>}
\MusicPDFWriteObject{6008}{<< /Type /Annot /Subtype /Link /Border [0 0 0] /H /I /C [1 0 0] /Rect [106.556 508.271 138.844 517.125] /A << /S /GoTo /D <74686D2E422E372E312E31> >> >>}
\MusicPDFWriteObject{6009}{<< /Type /Annot /Subtype /Link /Border [0 0 0] /H /I /C [1 0 0] /Rect [508.438 479.225 515.385 491.452] /A << /S /GoTo /D <6669677572652E422E34> >> >>}
\MusicPDFWriteObject{6010}{<< /Type /Annot /Subtype /Link /Border [0 0 0] /H /I /C [1 0 0] /Rect [396.118 639.576 428.558 651.829] /A << /S /GoTo /D <74686D2E422E362E312E31> >> >>}
\MusicPDFWriteObject{6011}{<< /Type /Annot /Subtype /Link /Border [0 0 0] /H /I /C [1 0 0] /Rect [491.599 639.576 524.039 651.829] /A << /S /GoTo /D <74686D2E422E362E312E33> >> >>}
\MusicPDFWriteObject{6012}{<< /Type /Annot /Subtype /Link /Border [0 0 0] /H /I /C [1 0 0] /Rect [170.441 554.786 177.414 566.741] /A << /S /GoTo /D <6669677572652E422E34> >> >>}
\MusicPDFWriteObject{6013}{<< /Type /Annot /Subtype /Link /Border [0 0 0] /H /I /C [1 0 0] /Rect [167.479 365.417 174.452 376.208] /A << /S /GoTo /D <6669677572652E422E34> >> >>}
\MusicPDFWriteObject{6014}{<< /Type /Annot /Subtype /Link /Border [0 0 0] /H /I /C [1 0 0] /Rect [244.063 671.482 281.332 683.435] /A << /S /GoTo /D <6571756174696F6E2E422E372E312E3131> >> >>}
\MusicPDFWriteObject{6015}{<< /Type /Annot /Subtype /Link /Border [0 0 0] /H /I /C [1 0 0] /Rect [294.361 533.775 331.629 545.727] /A << /S /GoTo /D <6571756174696F6E2E422E372E312E3131> >> >>}
\MusicPDFWriteObject{6016}{<< /Type /Annot /Subtype /Link /Border [0 0 0] /H /I /C [1 0 0] /Rect [343.503 533.775 380.772 545.727] /A << /S /GoTo /D <6571756174696F6E2E422E372E312E3133> >> >>}
\MusicPDFWriteObject{6017}{<< /Type /Annot /Subtype /Link /Border [0 0 0] /H /I /C [1 0 0] /Rect [412.219 533.775 449.487 545.727] /A << /S /GoTo /D <6571756174696F6E2E422E372E312E3134> >> >>}
\MusicPDFWriteObject{6018}{<< /Type /Annot /Subtype /Link /Border [0 0 0] /H /I /C [0 1 0] /Rect [527.49 536.265 543.178 545.133] /A << /S /GoTo /D <636974652E5432> >> >>}
\MusicPDFWriteObject{6019}{<< /Type /Annot /Subtype /Link /Border [0 0 0] /H /I /C [1 0 0] /Rect [117.208 494.322 154.476 506.275] /A << /S /GoTo /D <6571756174696F6E2E422E372E312E3134> >> >>}
\MusicPDFWriteObject{6020}{<< /Type /Annot /Subtype /Link /Border [0 0 0] /H /I /C [1 0 0] /Rect [153.315 275.228 185.603 287.18] /A << /S /GoTo /D <6571756174696F6E2E422E372E312E38> >> >>}
\MusicPDFWriteObject{6021}{<< /Type /Annot /Subtype /Link /Border [0 0 0] /H /I /C [1 0 0] /Rect [233.106 275.228 243.039 288.085] /A << /S /GoTo /D <48666F6F746E6F74652E313039> >> >>}
\MusicPDFWriteObject{6022}{<< /Type /Annot /Subtype /Link /Border [0 0 0] /H /I /C [1 0 0] /Rect [248.246 275.228 280.534 287.18] /A << /S /GoTo /D <6571756174696F6E2E422E372E312E39> >> >>}
\MusicPDFWriteObject{6023}{<< /Type /Annot /Subtype /Link /Border [0 0 0] /H /I /C [1 0 0] /Rect [498.421 275.228 530.709 287.18] /A << /S /GoTo /D <6571756174696F6E2E422E372E312E39> >> >>}
\MusicPDFWriteObject{6024}{<< /Type /Annot /Subtype /Link /Border [0 0 0] /H /I /C [1 0 0] /Rect [299.758 262.077 332.046 274.03] /A << /S /GoTo /D <6571756174696F6E2E422E372E312E38> >> >>}
\MusicPDFWriteObject{6025}{<< /Type /Annot /Subtype /Link /Border [0 0 0] /H /I /C [1 0 0] /Rect [164.021 248.926 201.289 260.879] /A << /S /GoTo /D <6571756174696F6E2E422E372E312E3136> >> >>}
\MusicPDFWriteObject{6026}{<< /Type /Annot /Subtype /Link /Border [0 0 0] /H /I /C [1 0 0] /Rect [290.728 235.775 322.411 247.728] /A << /S /GoTo /D <74686D2E422E362E312E31> >> >>}
\MusicPDFWriteObject{6027}{<< /Type /Annot /Subtype /Link /Border [0 0 0] /H /I /C [1 0 0] /Rect [353.974 235.775 391.243 247.728] /A << /S /GoTo /D <6571756174696F6E2E422E372E312E3130> >> >>}
\MusicPDFWriteObject{6028}{<< /Type /Annot /Subtype /Link /Border [0 0 0] /H /I /C [1 0 0] /Rect [288.128 472.145 319.81 484.1] /A << /S /GoTo /D <74686D2E422E372E322E32> >> >>}
\MusicPDFWriteObject{6029}{<< /Type /Annot /Subtype /Link /Border [0 0 0] /H /I /C [1 0 0] /Rect [296.815 323.756 329.709 335.711] /A << /S /GoTo /D <74686D2E422E372E322E32> >> >>}
\MusicPDFWriteObject{6030}{<< /Type /Annot /Subtype /Link /Border [0 0 0] /H /I /C [1 0 0] /Rect [186.88 275.789 219.168 287.742] /A << /S /GoTo /D <6571756174696F6E2E422E372E322E31> >> >>}
\MusicPDFWriteObject{6031}{<< /Type /Annot /Subtype /Link /Border [0 0 0] /H /I /C [1 0 0] /Rect [247.609 275.789 279.897 287.742] /A << /S /GoTo /D <6571756174696F6E2E422E372E322E33> >> >>}
\MusicPDFWriteObject{6032}{<< /Type /Annot /Subtype /Link /Border [0 0 0] /H /I /C [1 0 0] /Rect [73.712 671.482 106 683.435] /A << /S /GoTo /D <6571756174696F6E2E422E372E322E34> >> >>}
\MusicPDFWriteObject{6033}{<< /Type /Annot /Subtype /Link /Border [0 0 0] /H /I /C [1 0 0] /Rect [178.764 591.127 210.87 603.08] /A << /S /GoTo /D <74686D2E422E372E312E32> >> >>}
\MusicPDFWriteObject{6034}{<< /Type /Annot /Subtype /Link /Border [0 0 0] /H /I /C [1 0 0] /Rect [254.113 577.976 286.401 589.929] /A << /S /GoTo /D <6571756174696F6E2E422E372E312E38> >> >>}
\MusicPDFWriteObject{6035}{<< /Type /Annot /Subtype /Link /Border [0 0 0] /H /I /C [1 0 0] /Rect [165.939 499.514 197.621 511.469] /A << /S /GoTo /D <74686D2E422E372E312E32> >> >>}
\MusicPDFWriteObject{6036}{<< /Type /Annot /Subtype /Link /Border [0 0 0] /H /I /C [1 0 0] /Rect [151.785 486.364 184.073 498.316] /A << /S /GoTo /D <6571756174696F6E2E422E372E322E36> >> >>}
\MusicPDFWriteObject{6037}{<< /Type /Annot /Subtype /Link /Border [0 0 0] /H /I /C [1 0 0] /Rect [292.078 213.08 324.366 225.035] /A << /S /GoTo /D <74686D2E422E372E312E32> >> >>}
\MusicPDFWriteObject{6038}{<< /Type /Annot /Subtype /Link /Border [0 0 0] /H /I /C [1 0 0] /Rect [248.948 477.098 286.216 489.05] /A << /S /GoTo /D <6571756174696F6E2E422E372E322E3133> >> >>}
\MusicPDFWriteObject{6039}{<< /Type /Annot /Subtype /Link /Border [0 0 0] /H /I /C [1 0 0] /Rect [368.944 477.098 401.232 489.05] /A << /S /GoTo /D <6571756174696F6E2E422E372E322E39> >> >>}
\MusicPDFWriteObject{6040}{<< /Type /Annot /Subtype /Link /Border [0 0 0] /H /I /C [1 0 0] /Rect [243.071 391.997 274.754 403.952] /A << /S /GoTo /D <74686D2E422E372E322E32> >> >>}
\MusicPDFWriteObject{6041}{<< /Type /Annot /Subtype /Link /Border [0 0 0] /H /I /C [1 0 0] /Rect [310.24 352.545 342.529 364.498] /A << /S /GoTo /D <6571756174696F6E2E422E372E322E33> >> >>}
\MusicPDFWriteObject{6042}{<< /Type /Annot /Subtype /Link /Border [0 0 0] /H /I /C [1 0 0] /Rect [453.184 352.545 484.867 364.498] /A << /S /GoTo /D <74686D2E422E372E312E32> >> >>}
\MusicPDFWriteObject{6043}{<< /Type /Annot /Subtype /Link /Border [0 0 0] /H /I /C [1 0 0] /Rect [112.727 339.394 145.015 351.349] /A << /S /GoTo /D <74686D2E422E372E312E32> >> >>}
\MusicPDFWriteObject{6044}{<< /Type /Annot /Subtype /Link /Border [0 0 0] /H /I /C [1 0 0] /Rect [120.073 684.633 157.341 696.586] /A << /S /GoTo /D <6571756174696F6E2E422E372E322E3136> >> >>}
\MusicPDFWriteObject{6045}{<< /Type /Annot /Subtype /Link /Border [0 0 0] /H /I /C [1 0 0] /Rect [226.935 684.633 259.829 696.586] /A << /S /GoTo /D <74686D2E422E372E322E31> >> >>}
\MusicPDFWriteObject{6046}{<< /Type /Annot /Subtype /Link /Border [0 0 0] /H /I /C [1 0 0] /Rect [430.933 684.633 468.201 696.586] /A << /S /GoTo /D <6571756174696F6E2E422E372E322E3135> >> >>}
\MusicPDFWriteObject{6047}{<< /Type /Annot /Subtype /Link /Border [0 0 0] /H /I /C [1 0 0] /Rect [232.499 625.281 264.787 634.135] /A << /S /GoTo /D <74686D2E422E372E322E32> >> >>}
\MusicPDFWriteObject{6048}{<< /Type /Annot /Subtype /Link /Border [0 0 0] /H /I /C [1 0 0] /Rect [298.792 360.095 336.06 372.048] /A << /S /GoTo /D <6571756174696F6E2E422E372E322E3139> >> >>}
\MusicPDFWriteObject{6049}{<< /Type /Annot /Subtype /Link /Border [0 0 0] /H /I /C [1 0 0] /Rect [349.642 685.187 381.93 695.977] /A << /S /GoTo /D <74686D2E422E372E322E31> >> >>}
\MusicPDFWriteObject{6050}{<< /Type /Annot /Subtype /Link /Border [0 0 0] /H /I /C [1 0 0] /Rect [160.563 619.063 197.831 631.016] /A << /S /GoTo /D <6571756174696F6E2E422E372E322E3233> >> >>}
\MusicPDFWriteObject{6051}{<< /Type /Annot /Subtype /Link /Border [0 0 0] /H /I /C [1 0 0] /Rect [149.765 605.912 187.034 617.865] /A << /S /GoTo /D <6571756174696F6E2E422E372E322E3233> >> >>}
\MusicPDFWriteObject{6052}{<< /Type /Annot /Subtype /Link /Border [0 0 0] /H /I /C [1 0 0] /Rect [272.024 605.912 304.918 617.865] /A << /S /GoTo /D <74686D2E422E372E322E31> >> >>}
\MusicPDFWriteObject{6053}{<< /Type /Annot /Subtype /Link /Border [0 0 0] /H /I /C [1 0 0] /Rect [73.712 579.611 110.98 591.563] /A << /S /GoTo /D <6571756174696F6E2E422E372E322E3233> >> >>}
\MusicPDFWriteObject{6054}{<< /Type /Annot /Subtype /Link /Border [0 0 0] /H /I /C [1 0 0] /Rect [192.289 579.611 229.557 591.563] /A << /S /GoTo /D <6571756174696F6E2E422E372E322E3233> >> >>}
\MusicPDFWriteObject{6055}{<< /Type /Annot /Subtype /Link /Border [0 0 0] /H /I /C [1 0 0] /Rect [297.624 419.578 334.892 431.53] /A << /S /GoTo /D <6571756174696F6E2E422E372E322E3234> >> >>}
\MusicPDFWriteObject{6056}{<< /Type /Annot /Subtype /Link /Border [0 0 0] /H /I /C [1 0 0] /Rect [502.017 419.578 539.285 431.53] /A << /S /GoTo /D <6571756174696F6E2E422E372E322E3137> >> >>}
\MusicPDFWriteObject{6057}{<< /Type /Annot /Subtype /Link /Border [0 0 0] /H /I /C [1 0 0] /Rect [312.962 406.067 345.856 418.381] /A << /S /GoTo /D <74686D2E422E372E322E34> >> >>}
\MusicPDFWriteObject{6058}{<< /Type /Annot /Subtype /Link /Border [0 0 0] /H /I /C [1 0 0] /Rect [163.627 298.486 200.895 310.439] /A << /S /GoTo /D <6571756174696F6E2E422E372E322E3235> >> >>}
\MusicPDFWriteObject{6059}{<< /Type /Annot /Subtype /Link /Border [0 0 0] /H /I /C [1 0 0] /Rect [265.274 298.486 297.562 310.439] /A << /S /GoTo /D <6571756174696F6E2E422E372E312E39> >> >>}
\MusicPDFWriteObject{6060}{<< /Type /Annot /Subtype /Link /Border [0 0 0] /H /I /C [1 0 0] /Rect [355.305 298.486 387.017 310.439] /A << /S /GoTo /D <74686D2E422E372E312E32> >> >>}
\MusicPDFWriteObject{6061}{<< /Type /Annot /Subtype /Link /Border [0 0 0] /H /I /C [1 0 0] /Rect [504.62 298.486 541.888 310.439] /A << /S /GoTo /D <6571756174696F6E2E422E372E322E3235> >> >>}
\MusicPDFWriteObject{6062}{<< /Type /Annot /Subtype /Link /Border [0 0 0] /H /I /C [1 0 0] /Rect [98.422 285.335 130.71 297.288] /A << /S /GoTo /D <6571756174696F6E2E422E372E312E39> >> >>}
\MusicPDFWriteObject{6063}{<< /Type /Annot /Subtype /Link /Border [0 0 0] /H /I /C [0 1 0] /Rect [269.184 287.825 291.982 296.693] /A << /S /GoTo /D <636974652E4857> >> >>}
\MusicPDFWriteObject{6064}{<< /Type /Annot /Subtype /Link /Border [0 0 0] /H /I /C [0 1 0] /Rect [415.157 287.825 433.045 296.693] /A << /S /GoTo /D <636974652E474C> >> >>}
\MusicPDFWriteObject{6065}{<< /Type /Annot /Subtype /Link /Border [0 0 0] /H /I /C [1 0 0] /Rect [82.959 558.462 288.029 569.252] /A << /S /GoTo /D <73656374696F6E2E432E31> >> >>}
\MusicPDFWriteObject{6066}{<< /Type /Annot /Subtype /Link /Border [0 0 0] /H /I /C [1 0 0] /Rect [82.959 545.311 245.477 556.102] /A << /S /GoTo /D <73756273656374696F6E2E432E312E31> >> >>}
\MusicPDFWriteObject{6067}{<< /Type /Annot /Subtype /Link /Border [0 0 0] /H /I /C [1 0 0] /Rect [82.959 534.097 212.276 542.951] /A << /S /GoTo /D <73756273756273656374696F6E2E432E312E312E31> >> >>}
\MusicPDFWriteObject{6068}{<< /Type /Annot /Subtype /Link /Border [0 0 0] /H /I /C [1 0 0] /Rect [82.959 519.009 378.363 529.8] /A << /S /GoTo /D <73756273656374696F6E2E432E312E32> >> >>}
\MusicPDFWriteObject{6069}{<< /Type /Annot /Subtype /Link /Border [0 0 0] /H /I /C [1 0 0] /Rect [82.959 505.859 419.836 516.649] /A << /S /GoTo /D <73756273756273656374696F6E2E432E312E322E31> >> >>}
\MusicPDFWriteObject{6070}{<< /Type /Annot /Subtype /Link /Border [0 0 0] /H /I /C [1 0 0] /Rect [82.959 492.708 254.137 503.499] /A << /S /GoTo /D <73756273656374696F6E2E432E312E33> >> >>}
\MusicPDFWriteObject{6071}{<< /Type /Annot /Subtype /Link /Border [0 0 0] /H /I /C [1 0 0] /Rect [82.959 479.557 281.804 490.348] /A << /S /GoTo /D <73756273656374696F6E2E432E312E34> >> >>}
\MusicPDFWriteObject{6072}{<< /Type /Annot /Subtype /Link /Border [0 0 0] /H /I /C [1 0 0] /Rect [82.959 466.406 331.965 477.197] /A << /S /GoTo /D <73756273656374696F6E2E432E312E35> >> >>}
\MusicPDFWriteObject{6073}{<< /Type /Annot /Subtype /Link /Border [0 0 0] /H /I /C [1 0 0] /Rect [82.959 453.256 317.633 464.046] /A << /S /GoTo /D <73756273756273656374696F6E2E432E312E352E31> >> >>}
\MusicPDFWriteObject{6074}{<< /Type /Annot /Subtype /Link /Border [0 0 0] /H /I /C [1 0 0] /Rect [82.959 440.105 243.125 450.896] /A << /S /GoTo /D <73756273656374696F6E2E432E312E36> >> >>}
\MusicPDFWriteObject{6075}{<< /Type /Annot /Subtype /Link /Border [0 0 0] /H /I /C [1 0 0] /Rect [82.959 426.954 449.965 437.745] /A << /S /GoTo /D <73656374696F6E2E432E32> >> >>}
\MusicPDFWriteObject{6076}{<< /Type /Annot /Subtype /Link /Border [0 0 0] /H /I /C [1 0 0] /Rect [82.959 413.803 189.368 424.594] /A << /S /GoTo /D <73756273656374696F6E2E432E322E31> >> >>}
\MusicPDFWriteObject{6077}{<< /Type /Annot /Subtype /Link /Border [0 0 0] /H /I /C [1 0 0] /Rect [82.959 400.653 386.912 411.443] /A << /S /GoTo /D <73756273656374696F6E2E432E322E32> >> >>}
\MusicPDFWriteObject{6078}{<< /Type /Annot /Subtype /Link /Border [0 0 0] /H /I /C [1 0 0] /Rect [82.959 387.502 436.09 398.293] /A << /S /GoTo /D <73756273756273656374696F6E2E432E322E322E31> >> >>}
\MusicPDFWriteObject{6079}{<< /Type /Annot /Subtype /Link /Border [0 0 0] /H /I /C [1 0 0] /Rect [82.959 374.351 220.687 385.142] /A << /S /GoTo /D <73756273756273656374696F6E2E432E322E322E32> >> >>}
\MusicPDFWriteObject{6080}{<< /Type /Annot /Subtype /Link /Border [0 0 0] /H /I /C [1 0 0] /Rect [82.959 361.2 523.77 372.619] /A << /S /GoTo /D <73756273756273656374696F6E2E432E322E322E33> >> >>}
\MusicPDFWriteObject{6081}{<< /Type /Annot /Subtype /Link /Border [0 0 0] /H /I /C [1 0 0] /Rect [82.959 348.05 241.437 358.84] /A << /S /GoTo /D <73756273656374696F6E2E432E322E33> >> >>}
\MusicPDFWriteObject{6082}{<< /Type /Annot /Subtype /Link /Border [0 0 0] /H /I /C [1 0 0] /Rect [82.959 334.899 384.837 345.69] /A << /S /GoTo /D <73756273756273656374696F6E2E432E322E332E31> >> >>}
\MusicPDFWriteObject{6083}{<< /Type /Annot /Subtype /Link /Border [0 0 0] /H /I /C [1 0 0] /Rect [82.959 321.748 243.54 332.539] /A << /S /GoTo /D <73756273756273656374696F6E2E432E322E332E32> >> >>}
\MusicPDFWriteObject{6084}{<< /Type /Annot /Subtype /Link /Border [0 0 0] /H /I /C [1 0 0] /Rect [82.959 308.597 438.871 319.388] /A << /S /GoTo /D <73756273756273656374696F6E2E432E322E332E33> >> >>}
\MusicPDFWriteObject{6085}{<< /Type /Annot /Subtype /Link /Border [0 0 0] /H /I /C [1 0 0] /Rect [82.959 295.447 240.054 306.237] /A << /S /GoTo /D <73756273756273656374696F6E2E432E322E332E34> >> >>}
\MusicPDFWriteObject{6086}{<< /Type /Annot /Subtype /Link /Border [0 0 0] /H /I /C [1 0 0] /Rect [82.959 282.296 517.335 293.087] /A << /S /GoTo /D <73756273756273656374696F6E2E432E322E332E35> >> >>}
\MusicPDFWriteObject{6087}{<< /Type /Annot /Subtype /Link /Border [0 0 0] /H /I /C [1 0 0] /Rect [82.959 269.145 184.415 279.936] /A << /S /GoTo /D <73756273656374696F6E2E432E322E34> >> >>}
\MusicPDFWriteObject{6088}{<< /Type /Annot /Subtype /Link /Border [0 0 0] /H /I /C [1 0 0] /Rect [82.959 255.994 262.16 266.785] /A << /S /GoTo /D <73756273756273656374696F6E2E432E322E342E31> >> >>}
\MusicPDFWriteObject{6089}{<< /Type /Annot /Subtype /Link /Border [0 0 0] /H /I /C [1 0 0] /Rect [82.959 243.286 246.37 253.634] /A << /S /GoTo /D <73756273756273656374696F6E2E432E322E342E32> >> >>}
\MusicPDFWriteObject{6090}{<< /Type /Annot /Subtype /Link /Border [0 0 0] /H /I /C [1 0 0] /Rect [82.959 229.693 360.462 240.483] /A << /S /GoTo /D <73756273756273656374696F6E2E432E322E342E33> >> >>}
\MusicPDFWriteObject{6091}{<< /Type /Annot /Subtype /Link /Border [0 0 0] /H /I /C [1 0 0] /Rect [82.959 216.542 227.438 227.333] /A << /S /GoTo /D <73756273756273656374696F6E2E432E322E342E34> >> >>}
\MusicPDFWriteObject{6092}{<< /Type /Annot /Subtype /Link /Border [0 0 0] /H /I /C [1 0 0] /Rect [82.959 203.391 235.683 214.182] /A << /S /GoTo /D <73756273756273656374696F6E2E432E322E342E35> >> >>}
\MusicPDFWriteObject{6093}{<< /Type /Annot /Subtype /Link /Border [0 0 0] /H /I /C [1 0 0] /Rect [82.959 192.177 202.925 201.031] /A << /S /GoTo /D <73756273756273656374696F6E2E432E322E342E36> >> >>}
\MusicPDFWriteObject{6094}{<< /Type /Annot /Subtype /Link /Border [0 0 0] /H /I /C [1 0 0] /Rect [82.959 177.09 342.091 187.88] /A << /S /GoTo /D <73756273656374696F6E2E432E322E35> >> >>}
\MusicPDFWriteObject{6095}{<< /Type /Annot /Subtype /Link /Border [0 0 0] /H /I /C [1 0 0] /Rect [82.959 163.939 190.225 174.73] /A << /S /GoTo /D <73656374696F6E2E432E33> >> >>}
\MusicPDFWriteObject{6096}{<< /Type /Annot /Subtype /Link /Border [0 0 0] /H /I /C [1 0 0] /Rect [82.959 150.788 175.949 161.579] /A << /S /GoTo /D <73756273656374696F6E2E432E332E31> >> >>}
\MusicPDFWriteObject{6097}{<< /Type /Annot /Subtype /Link /Border [0 0 0] /H /I /C [1 0 0] /Rect [82.959 137.637 403.484 148.428] /A << /S /GoTo /D <73756273656374696F6E2E432E332E32> >> >>}
\MusicPDFWriteObject{6098}{<< /Type /Annot /Subtype /Link /Border [0 0 0] /H /I /C [1 0 0] /Rect [82.959 124.487 379.857 135.277] /A << /S /GoTo /D <73756273656374696F6E2E432E332E33> >> >>}
\MusicPDFWriteObject{6099}{<< /Type /Annot /Subtype /Link /Border [0 0 0] /H /I /C [1 0 0] /Rect [82.959 111.336 290.63 122.127] /A << /S /GoTo /D <73756273656374696F6E2E432E332E34> >> >>}
\MusicPDFWriteObject{6100}{<< /Type /Annot /Subtype /Link /Border [0 0 0] /H /I /C [1 0 0] /Rect [82.959 98.185 362.924 108.976] /A << /S /GoTo /D <73756273756273656374696F6E2E432E332E342E31> >> >>}
\MusicPDFWriteObject{6101}{<< /Type /Annot /Subtype /Link /Border [0 0 0] /H /I /C [1 0 0] /Rect [82.959 85.034 390.536 95.825] /A << /S /GoTo /D <73756273656374696F6E2E432E332E35> >> >>}
\MusicPDFWriteObject{6102}{<< /Type /Annot /Subtype /Link /Border [0 0 0] /H /I /C [1 0 0] /Rect [82.959 685.187 245.477 695.977] /A << /S /GoTo /D <73756273756273656374696F6E2E432E332E352E31> >> >>}
\MusicPDFWriteObject{6103}{<< /Type /Annot /Subtype /Link /Border [0 0 0] /H /I /C [1 0 0] /Rect [82.959 672.035 328.075 682.827] /A << /S /GoTo /D <73756273656374696F6E2E432E332E36> >> >>}
\MusicPDFWriteObject{6104}{<< /Type /Annot /Subtype /Link /Border [0 0 0] /H /I /C [1 0 0] /Rect [82.959 658.885 475.17 669.676] /A << /S /GoTo /D <73756273756273656374696F6E2E432E332E362E31> >> >>}
\MusicPDFWriteObject{6105}{<< /Type /Annot /Subtype /Link /Border [0 0 0] /H /I /C [1 0 0] /Rect [82.959 645.734 355.897 656.525] /A << /S /GoTo /D <73756273756273656374696F6E2E432E332E362E32> >> >>}
\MusicPDFWriteObject{6106}{<< /Type /Annot /Subtype /Link /Border [0 0 0] /H /I /C [1 0 0] /Rect [82.959 632.584 187.431 643.374] /A << /S /GoTo /D <73756273656374696F6E2E432E332E37> >> >>}
\MusicPDFWriteObject{6107}{<< /Type /Annot /Subtype /Link /Border [0 0 0] /H /I /C [1 0 0] /Rect [82.959 619.433 255.16 630.223] /A << /S /GoTo /D <73656374696F6E2E432E34> >> >>}
\MusicPDFWriteObject{6108}{<< /Type /Annot /Subtype /Link /Border [0 0 0] /H /I /C [1 0 0] /Rect [82.959 606.282 195.344 617.073] /A << /S /GoTo /D <73756273656374696F6E2E432E342E31> >> >>}
\MusicPDFWriteObject{6109}{<< /Type /Annot /Subtype /Link /Border [0 0 0] /H /I /C [1 0 0] /Rect [82.959 593.131 340.514 603.922] /A << /S /GoTo /D <73756273656374696F6E2E432E342E32> >> >>}
\MusicPDFWriteObject{6110}{<< /Type /Annot /Subtype /Link /Border [0 0 0] /H /I /C [1 0 0] /Rect [82.959 579.981 329.751 590.771] /A << /S /GoTo /D <73756273656374696F6E2E432E342E33> >> >>}
\MusicPDFWriteObject{6111}{<< /Type /Annot /Subtype /Link /Border [0 0 0] /H /I /C [1 0 0] /Rect [82.959 568.766 288.693 577.62] /A << /S /GoTo /D <73756273756273656374696F6E2E432E342E332E31> >> >>}
\MusicPDFWriteObject{6112}{<< /Type /Annot /Subtype /Link /Border [0 0 0] /H /I /C [1 0 0] /Rect [82.959 553.679 434.693 564.47] /A << /S /GoTo /D <73756273756273656374696F6E2E432E342E332E32> >> >>}
\MusicPDFWriteObject{6113}{<< /Type /Annot /Subtype /Link /Border [0 0 0] /H /I /C [1 0 0] /Rect [82.959 540.528 214.489 551.319] /A << /S /GoTo /D <73756273756273656374696F6E2E432E342E332E33> >> >>}
\MusicPDFWriteObject{6114}{<< /Type /Annot /Subtype /Link /Border [0 0 0] /H /I /C [1 0 0] /Rect [82.959 527.377 321.977 538.168] /A << /S /GoTo /D <73756273656374696F6E2E432E342E34> >> >>}
\MusicPDFWriteObject{6115}{<< /Type /Annot /Subtype /Link /Border [0 0 0] /H /I /C [1 0 0] /Rect [82.959 514.227 413.03 525.017] /A << /S /GoTo /D <73656374696F6E2E432E35> >> >>}
\MusicPDFWriteObject{6116}{<< /Type /Annot /Subtype /Link /Border [0 0 0] /H /I /C [1 0 0] /Rect [82.959 503.013 351.83 511.867] /A << /S /GoTo /D <73756273656374696F6E2E432E352E31> >> >>}
\MusicPDFWriteObject{6117}{<< /Type /Annot /Subtype /Link /Border [0 0 0] /H /I /C [1 0 0] /Rect [82.959 487.925 473.305 498.716] /A << /S /GoTo /D <73756273756273656374696F6E2E432E352E312E31> >> >>}
\MusicPDFWriteObject{6118}{<< /Type /Annot /Subtype /Link /Border [0 0 0] /H /I /C [1 0 0] /Rect [82.959 474.774 257.374 485.565] /A << /S /GoTo /D <73756273756273656374696F6E2E432E352E312E32> >> >>}
\MusicPDFWriteObject{6119}{<< /Type /Annot /Subtype /Link /Border [0 0 0] /H /I /C [1 0 0] /Rect [82.959 461.624 366.936 472.414] /A << /S /GoTo /D <73756273756273656374696F6E2E432E352E312E33> >> >>}
\MusicPDFWriteObject{6120}{<< /Type /Annot /Subtype /Link /Border [0 0 0] /H /I /C [1 0 0] /Rect [82.959 448.473 300.037 459.264] /A << /S /GoTo /D <73756273756273656374696F6E2E432E352E312E34> >> >>}
\MusicPDFWriteObject{6121}{<< /Type /Annot /Subtype /Link /Border [0 0 0] /H /I /C [1 0 0] /Rect [82.959 435.322 310.356 446.113] /A << /S /GoTo /D <73756273756273656374696F6E2E432E352E312E35> >> >>}
\MusicPDFWriteObject{6122}{<< /Type /Annot /Subtype /Link /Border [0 0 0] /H /I /C [1 0 0] /Rect [82.959 422.171 280.116 432.962] /A << /S /GoTo /D <73756273656374696F6E2E432E352E32> >> >>}
\MusicPDFWriteObject{6123}{<< /Type /Annot /Subtype /Link /Border [0 0 0] /H /I /C [1 0 0] /Rect [82.959 409.021 302.955 419.811] /A << /S /GoTo /D <73756273756273656374696F6E2E432E352E322E31> >> >>}
\MusicPDFWriteObject{6124}{<< /Type /Annot /Subtype /Link /Border [0 0 0] /H /I /C [1 0 0] /Rect [82.959 395.87 282.357 406.661] /A << /S /GoTo /D <73756273656374696F6E2E432E352E33> >> >>}
\MusicPDFWriteObject{6125}{<< /Type /Annot /Subtype /Link /Border [0 0 0] /H /I /C [1 0 0] /Rect [82.959 384.656 225.39 393.51] /A << /S /GoTo /D <73756273756273656374696F6E2E432E352E332E31> >> >>}
\MusicPDFWriteObject{6126}{<< /Type /Annot /Subtype /Link /Border [0 0 0] /H /I /C [1 0 0] /Rect [82.959 371.505 259.227 380.359] /A << /S /GoTo /D <73756273656374696F6E2E432E352E34> >> >>}
\MusicPDFWriteObject{6127}{<< /Type /Annot /Subtype /Link /Border [0 0 0] /H /I /C [1 0 0] /Rect [82.959 356.418 267.334 367.208] /A << /S /GoTo /D <73756273756273656374696F6E2E432E352E342E31> >> >>}
\MusicPDFWriteObject{6128}{<< /Type /Annot /Subtype /Link /Border [0 0 0] /H /I /C [1 0 0] /Rect [365.854 547.148 375.455 559.101] /A << /S /GoTo /D <617070656E6469782E44> >> >>}
\MusicPDFWriteObject{6129}{<< /Type /Annot /Subtype /Link /Border [0 0 0] /H /I /C [0 1 0] /Rect [199.068 464.756 248.109 473.624] /A << /S /GoTo /D <636974652E426F752D545653> >> >>}
\MusicPDFWriteObject{6130}{<< /Type /Annot /Subtype /Link /Border [0 0 0] /H /I /C [0 1 0] /Rect [285.654 464.756 310.714 473.624] /A << /S /GoTo /D <636974652E53636861> >> >>}
\MusicPDFWriteObject{6131}{<< /Type /Annot /Subtype /Link /Border [0 0 0] /H /I /C [0 1 0] /Rect [348.259 464.756 366.458 473.624] /A << /S /GoTo /D <636974652E4A6172> >> >>}
\MusicPDFWriteObject{6132}{<< /Type /Annot /Subtype /Link /Border [0 0 0] /H /I /C [0 1 0] /Rect [405.467 464.756 446.076 473.624] /A << /S /GoTo /D <636974652E4B6F7468652D49> >> >>}
\MusicPDFWriteObject{6133}{<< /Type /Annot /Subtype /Link /Border [0 0 0] /H /I /C [0 1 0] /Rect [524.054 464.756 543.159 473.624] /A << /S /GoTo /D <636974652E4E42> >> >>}
\MusicPDFWriteObject{6134}{<< /Type /Annot /Subtype /Link /Border [0 0 0] /H /I /C [0 1 0] /Rect [270.047 438.454 319.087 447.322] /A << /S /GoTo /D <636974652E426F752D545653> >> >>}
\MusicPDFWriteObject{6135}{<< /Type /Annot /Subtype /Link /Border [0 0 0] /H /I /C [0 1 0] /Rect [280.696 425.304 329.736 434.172] /A << /S /GoTo /D <636974652E426F752D545653> >> >>}
\MusicPDFWriteObject{6136}{<< /Type /Annot /Subtype /Link /Border [0 0 0] /H /I /C [0 1 0] /Rect [180.93 383.759 205.99 392.627] /A << /S /GoTo /D <636974652E53636861> >> >>}
\MusicPDFWriteObject{6137}{<< /Type /Annot /Subtype /Link /Border [0 0 0] /H /I /C [0 1 0] /Rect [268.401 252.252 293.461 261.12] /A << /S /GoTo /D <636974652E53636861> >> >>}
\MusicPDFWriteObject{6138}{<< /Type /Annot /Subtype /Link /Border [0 0 0] /H /I /C [0 1 0] /Rect [345.711 252.252 386.32 261.12] /A << /S /GoTo /D <636974652E4B6F7468652D49> >> >>}
\MusicPDFWriteObject{6139}{<< /Type /Annot /Subtype /Link /Border [0 0 0] /H /I /C [0 1 0] /Rect [209.833 225.95 228.938 234.818] /A << /S /GoTo /D <636974652E4E42> >> >>}
\MusicPDFWriteObject{6140}{<< /Type /Annot /Subtype /Link /Border [0 0 0] /H /I /C [0 1 0] /Rect [204.404 212.799 222.603 221.667] /A << /S /GoTo /D <636974652E4A6172> >> >>}
\MusicPDFWriteObject{6141}{<< /Type /Annot /Subtype /Link /Border [0 0 0] /H /I /C [0 1 0] /Rect [356.448 212.799 375.553 221.667] /A << /S /GoTo /D <636974652E4E42> >> >>}
\MusicPDFWriteObject{6142}{<< /Type /Annot /Subtype /Link /Border [0 0 0] /H /I /C [0 1 0] /Rect [286.204 173.347 311.264 182.215] /A << /S /GoTo /D <636974652E53636861> >> >>}
\MusicPDFWriteObject{6143}{<< /Type /Annot /Subtype /Link /Border [0 0 0] /H /I /C [0 1 0] /Rect [363.5 173.347 382.604 182.215] /A << /S /GoTo /D <636974652E4E42> >> >>}
\MusicPDFWriteObject{6144}{<< /Type /Annot /Subtype /Link /Border [0 0 0] /H /I /C [0 1 0] /Rect [73.748 142.367 92.722 151.235] /A << /S /GoTo /D <636974652E547265> >> >>}
\MusicPDFWriteObject{6145}{<< /Type /Annot /Subtype /Link /Border [0 0 0] /H /I /C [0 1 0] /Rect [156.814 142.367 175.919 151.235] /A << /S /GoTo /D <636974652E4E42> >> >>}
\MusicPDFWriteObject{6146}{<< /Type /Annot /Subtype /Link /Border [0 0 0] /H /I /C [0 1 0] /Rect [518.209 687.123 543.27 695.991] /A << /S /GoTo /D <636974652E53636861> >> >>}
\MusicPDFWriteObject{6147}{<< /Type /Annot /Subtype /Link /Border [0 0 0] /H /I /C [0 1 0] /Rect [465.007 549.256 483.206 558.124] /A << /S /GoTo /D <636974652E4A6172> >> >>}
\MusicPDFWriteObject{6148}{<< /Type /Annot /Subtype /Link /Border [0 0 0] /H /I /C [0 1 0] /Rect [147.332 466.084 172.392 474.952] /A << /S /GoTo /D <636974652E53636861> >> >>}
\MusicPDFWriteObject{6149}{<< /Type /Annot /Subtype /Link /Border [0 0 0] /H /I /C [0 1 0] /Rect [218.833 422.365 261.856 431.233] /A << /S /GoTo /D <636974652E426F752D4754> >> >>}
\MusicPDFWriteObject{6150}{<< /Type /Annot /Subtype /Link /Border [0 0 0] /H /I /C [0 1 0] /Rect [265.771 409.214 290.831 418.082] /A << /S /GoTo /D <636974652E53636861> >> >>}
\MusicPDFWriteObject{6151}{<< /Type /Annot /Subtype /Link /Border [0 0 0] /H /I /C [0 1 0] /Rect [458.277 396.063 498.886 404.931] /A << /S /GoTo /D <636974652E4B6F7468652D49> >> >>}
\MusicPDFWriteObject{6152}{<< /Type /Annot /Subtype /Link /Border [0 0 0] /H /I /C [0 1 0] /Rect [73.795 382.912 122.836 391.78] /A << /S /GoTo /D <636974652E426F752D545653> >> >>}
\MusicPDFWriteObject{6153}{<< /Type /Annot /Subtype /Link /Border [0 0 0] /H /I /C [0 1 0] /Rect [73.77 299.74 98.831 308.608] /A << /S /GoTo /D <636974652E53636861> >> >>}
\MusicPDFWriteObject{6154}{<< /Type /Annot /Subtype /Link /Border [0 0 0] /H /I /C [0 1 0] /Rect [166.839 299.74 215.88 308.608] /A << /S /GoTo /D <636974652E426F752D545653> >> >>}
\MusicPDFWriteObject{6155}{<< /Type /Annot /Subtype /Link /Border [0 0 0] /H /I /C [0 1 0] /Rect [283.889 299.74 302.994 308.608] /A << /S /GoTo /D <636974652E4E42> >> >>}
\MusicPDFWriteObject{6156}{<< /Type /Annot /Subtype /Link /Border [0 0 0] /H /I /C [0 1 0] /Rect [73.721 232.697 98.781 241.565] /A << /S /GoTo /D <636974652E53636861> >> >>}
\MusicPDFWriteObject{6157}{<< /Type /Annot /Subtype /Link /Border [0 0 0] /H /I /C [0 1 0] /Rect [439.304 219.546 464.364 228.414] /A << /S /GoTo /D <636974652E53636861> >> >>}
\MusicPDFWriteObject{6158}{<< /Type /Annot /Subtype /Link /Border [0 0 0] /H /I /C [0 1 0] /Rect [518.204 85.086 543.264 93.954] /A << /S /GoTo /D <636974652E53636861> >> >>}
\MusicPDFWriteObject{6159}{<< /Type /Annot /Subtype /Link /Border [0 0 0] /H /I /C [0 1 0] /Rect [105.317 621.369 148.34 630.237] /A << /S /GoTo /D <636974652E426F752D4754> >> >>}
\MusicPDFWriteObject{6160}{<< /Type /Annot /Subtype /Link /Border [0 0 0] /H /I /C [0 1 0] /Rect [157.117 595.068 200.14 603.936] /A << /S /GoTo /D <636974652E426F752D4754> >> >>}
\MusicPDFWriteObject{6161}{<< /Type /Annot /Subtype /Link /Border [0 0 0] /H /I /C [0 1 0] /Rect [117.946 568.766 135.847 577.634] /A << /S /GoTo /D <636974652E536D> >> >>}
\MusicPDFWriteObject{6162}{<< /Type /Annot /Subtype /Link /Border [0 0 0] /H /I /C [0 1 0] /Rect [175.963 568.766 189.409 577.634] /A << /S /GoTo /D <636974652E436C> >> >>}
\MusicPDFWriteObject{6163}{<< /Type /Annot /Subtype /Link /Border [0 0 0] /H /I /C [0 1 0] /Rect [308.776 568.766 326.677 577.634] /A << /S /GoTo /D <636974652E536D> >> >>}
\MusicPDFWriteObject{6164}{<< /Type /Annot /Subtype /Link /Border [0 0 0] /H /I /C [0 1 0] /Rect [471.583 568.766 487.292 577.634] /A << /S /GoTo /D <636974652E4261> >> >>}
\MusicPDFWriteObject{6165}{<< /Type /Annot /Subtype /Link /Border [0 0 0] /H /I /C [0 1 0] /Rect [343.381 555.616 356.828 564.484] /A << /S /GoTo /D <636974652E436C> >> >>}
\MusicPDFWriteObject{6166}{<< /Type /Annot /Subtype /Link /Border [0 0 0] /H /I /C [0 1 0] /Rect [462.514 555.616 481.619 564.484] /A << /S /GoTo /D <636974652E4E42> >> >>}
\MusicPDFWriteObject{6167}{<< /Type /Annot /Subtype /Link /Border [0 0 0] /H /I /C [0 1 0] /Rect [152.651 542.465 171.756 551.333] /A << /S /GoTo /D <636974652E4E42> >> >>}
\MusicPDFWriteObject{6168}{<< /Type /Annot /Subtype /Link /Border [0 0 0] /H /I /C [0 1 0] /Rect [133.087 316.089 158.147 324.957] /A << /S /GoTo /D <636974652E53636861> >> >>}
\MusicPDFWriteObject{6169}{<< /Type /Annot /Subtype /Link /Border [0 0 0] /H /I /C [0 1 0] /Rect [225.835 589.217 244.808 598.085] /A << /S /GoTo /D <636974652E547265> >> >>}
\MusicPDFWriteObject{6170}{<< /Type /Annot /Subtype /Link /Border [0 0 0] /H /I /C [0 1 0] /Rect [138.176 562.915 157.281 571.783] /A << /S /GoTo /D <636974652E4E42> >> >>}
\MusicPDFWriteObject{6171}{<< /Type /Annot /Subtype /Link /Border [0 0 0] /H /I /C [0 1 0] /Rect [122.618 549.765 141.723 558.633] /A << /S /GoTo /D <636974652E4E42> >> >>}
\MusicPDFWriteObject{6172}{<< /Type /Annot /Subtype /Link /Border [0 0 0] /H /I /C [0 1 0] /Rect [456.88 536.614 475.985 545.482] /A << /S /GoTo /D <636974652E4E42> >> >>}
\MusicPDFWriteObject{6173}{<< /Type /Annot /Subtype /Link /Border [0 0 0] /H /I /C [0 1 0] /Rect [73.734 425.557 98.795 434.425] /A << /S /GoTo /D <636974652E53636861> >> >>}
\MusicPDFWriteObject{6174}{<< /Type /Annot /Subtype /Link /Border [0 0 0] /H /I /C [1 0 0] /Rect [362.297 396.765 371.898 408.72] /A << /S /GoTo /D <617070656E6469782E44> >> >>}
\MusicPDFWriteObject{6175}{<< /Type /Annot /Subtype /Link /Border [0 0 0] /H /I /C [1 0 0] /Rect [318.037 384.168 342.399 394.959] /A << /S /GoTo /D <73756273656374696F6E2E432E312E35> >> >>}
\MusicPDFWriteObject{6176}{<< /Type /Annot /Subtype /Link /Border [0 0 0] /H /I /C [0 1 0] /Rect [224.907 320.351 249.968 329.219] /A << /S /GoTo /D <636974652E53636861> >> >>}
\MusicPDFWriteObject{6177}{<< /Type /Annot /Subtype /Link /Border [0 0 0] /H /I /C [1 0 0] /Rect [512.884 304.71 544.945 316.663] /A << /S /GoTo /D <74686D2E432E312E352E34> >> >>}
\MusicPDFWriteObject{6178}{<< /Type /Annot /Subtype /Link /Border [0 0 0] /H /I /C [0 1 0] /Rect [259.508 222.445 284.568 231.312] /A << /S /GoTo /D <636974652E53636861> >> >>}
\MusicPDFWriteObject{6179}{<< /Type /Annot /Subtype /Link /Border [0 0 0] /H /I /C [0 1 0] /Rect [255.959 196.143 281.019 205.011] /A << /S /GoTo /D <636974652E53636861> >> >>}
\MusicPDFWriteObject{6180}{<< /Type /Annot /Subtype /Link /Border [0 0 0] /H /I /C [0 1 0] /Rect [471.574 631.266 496.634 640.134] /A << /S /GoTo /D <636974652E53636861> >> >>}
\MusicPDFWriteObject{6181}{<< /Type /Annot /Subtype /Link /Border [0 0 0] /H /I /C [0 1 0] /Rect [220.471 614.862 245.531 623.73] /A << /S /GoTo /D <636974652E53636861> >> >>}
\MusicPDFWriteObject{6182}{<< /Type /Annot /Subtype /Link /Border [0 0 0] /H /I /C [0 1 0] /Rect [117.728 498.65 136.702 507.518] /A << /S /GoTo /D <636974652E547265> >> >>}
\MusicPDFWriteObject{6183}{<< /Type /Annot /Subtype /Link /Border [0 0 0] /H /I /C [0 1 0] /Rect [460.044 385.408 479.149 394.276] /A << /S /GoTo /D <636974652E4E42> >> >>}
\MusicPDFWriteObject{6184}{<< /Type /Annot /Subtype /Link /Border [0 0 0] /H /I /C [0 1 0] /Rect [495.316 372.208 520.377 381.076] /A << /S /GoTo /D <636974652E53636861> >> >>}
\MusicPDFWriteObject{6185}{<< /Type /Annot /Subtype /Link /Border [0 0 0] /H /I /C [0 1 0] /Rect [202.411 345.907 227.471 354.775] /A << /S /GoTo /D <636974652E53636861> >> >>}
\MusicPDFWriteObject{6186}{<< /Type /Annot /Subtype /Link /Border [0 0 0] /H /I /C [0 1 0] /Rect [437.949 332.756 463.009 341.624] /A << /S /GoTo /D <636974652E53636861> >> >>}
\MusicPDFWriteObject{6187}{<< /Type /Annot /Subtype /Link /Border [0 0 0] /H /I /C [0 1 0] /Rect [238.214 306.455 257.188 315.322] /A << /S /GoTo /D <636974652E547265> >> >>}
\MusicPDFWriteObject{6188}{<< /Type /Annot /Subtype /Link /Border [0 0 0] /H /I /C [0 1 0] /Rect [147.514 198.328 172.574 207.196] /A << /S /GoTo /D <636974652E53636861> >> >>}
\MusicPDFWriteObject{6189}{<< /Type /Annot /Subtype /Link /Border [0 0 0] /H /I /C [1 0 0] /Rect [330.496 108.897 363.532 120.852] /A << /S /GoTo /D <74686D2E432E312E342E34> >> >>}
\MusicPDFWriteObject{6190}{<< /Type /Annot /Subtype /Link /Border [0 0 0] /H /I /C [1 0 0] /Rect [299.012 82.595 331.439 94.55] /A << /S /GoTo /D <74686D2E432E342E332E34> >> >>}
\MusicPDFWriteObject{6191}{<< /Type /Annot /Subtype /Link /Border [0 0 0] /H /I /C [1 0 0] /Rect [234.791 630.978 266.609 643.983] /A << /S /GoTo /D <74686D2E432E312E342E34> >> >>}
\MusicPDFWriteObject{6192}{<< /Type /Annot /Subtype /Link /Border [0 0 0] /H /I /C [1 0 0] /Rect [463.855 630.978 495.673 643.983] /A << /S /GoTo /D <74686D2E432E312E342E31> >> >>}
\MusicPDFWriteObject{6193}{<< /Type /Annot /Subtype /Link /Border [0 0 0] /H /I /C [0 1 0] /Rect [318.833 608.219 343.893 617.087] /A << /S /GoTo /D <636974652E53636861> >> >>}
\MusicPDFWriteObject{6194}{<< /Type /Annot /Subtype /Link /Border [0 0 0] /H /I /C [1 0 0] /Rect [171.504 592.577 204.539 604.532] /A << /S /GoTo /D <74686D2E432E312E342E35> >> >>}
\MusicPDFWriteObject{6195}{<< /Type /Annot /Subtype /Link /Border [0 0 0] /H /I /C [1 0 0] /Rect [106.556 579.427 138.983 591.38] /A << /S /GoTo /D <74686D2E432E322E312E37> >> >>}
\MusicPDFWriteObject{6196}{<< /Type /Annot /Subtype /Link /Border [0 0 0] /H /I /C [0 1 0] /Rect [524.186 529.314 543.159 538.182] /A << /S /GoTo /D <636974652E547265> >> >>}
\MusicPDFWriteObject{6197}{<< /Type /Annot /Subtype /Link /Border [0 0 0] /H /I /C [0 1 0] /Rect [98.722 503.013 123.782 511.881] /A << /S /GoTo /D <636974652E53636861> >> >>}
\MusicPDFWriteObject{6198}{<< /Type /Annot /Subtype /Link /Border [0 0 0] /H /I /C [1 0 0] /Rect [162.985 474.774 180.037 485.565] /A << /S /GoTo /D <73656374696F6E2E432E34> >> >>}
\MusicPDFWriteObject{6199}{<< /Type /Annot /Subtype /Link /Border [0 0 0] /H /I /C [1 0 0] /Rect [365.208 474.774 382.26 485.565] /A << /S /GoTo /D <73656374696F6E2E432E35> >> >>}
\MusicPDFWriteObject{6200}{<< /Type /Annot /Subtype /Link /Border [0 0 0] /H /I /C [1 0 0] /Rect [393.223 333.014 426.014 343.804] /A << /S /GoTo /D <74686D2E432E312E332E33> >> >>}
\MusicPDFWriteObject{6201}{<< /Type /Annot /Subtype /Link /Border [0 0 0] /H /I /C [1 0 0] /Rect [437.696 306.158 462.83 318.113] /A << /S /GoTo /D <73756273656374696F6E2E432E322E34> >> >>}
\MusicPDFWriteObject{6202}{<< /Type /Annot /Subtype /Link /Border [0 0 0] /H /I /C [0 1 0] /Rect [302.826 282.347 327.886 291.215] /A << /S /GoTo /D <636974652E53636861> >> >>}
\MusicPDFWriteObject{6203}{<< /Type /Annot /Subtype /Link /Border [0 0 0] /H /I /C [1 0 0] /Rect [290.462 267.26 323.303 278.051] /A << /S /GoTo /D <74686D2E442E312E352E37> >> >>}
\MusicPDFWriteObject{6204}{<< /Type /Annot /Subtype /Link /Border [0 0 0] /H /I /C [1 0 0] /Rect [248.934 253.555 274.067 265.51] /A << /S /GoTo /D <73756273656374696F6E2E432E312E35> >> >>}
\MusicPDFWriteObject{6205}{<< /Type /Annot /Subtype /Link /Border [0 0 0] /H /I /C [0 1 0] /Rect [171.35 229.744 196.41 238.612] /A << /S /GoTo /D <636974652E53636861> >> >>}
\MusicPDFWriteObject{6206}{<< /Type /Annot /Subtype /Link /Border [0 0 0] /H /I /C [1 0 0] /Rect [140.228 187.802 172.655 199.755] /A << /S /GoTo /D <6571756174696F6E2E432E312E352E31> >> >>}
\MusicPDFWriteObject{6207}{<< /Type /Annot /Subtype /Link /Border [0 0 0] /H /I /C [0 1 0] /Rect [340.762 163.99 369.523 172.858] /A << /S /GoTo /D <636974652E47726F2D66> >> >>}
\MusicPDFWriteObject{6208}{<< /Type /Annot /Subtype /Link /Border [0 0 0] /H /I /C [0 1 0] /Rect [305.266 124.538 330.327 133.406] /A << /S /GoTo /D <636974652E53636861> >> >>}
\MusicPDFWriteObject{6209}{<< /Type /Annot /Subtype /Link /Border [0 0 0] /H /I /C [0 1 0] /Rect [247.733 111.387 272.794 120.255] /A << /S /GoTo /D <636974652E53636861> >> >>}
\MusicPDFWriteObject{6210}{<< /Type /Annot /Subtype /Link /Border [0 0 0] /H /I /C [0 1 0] /Rect [442.342 111.387 461.316 120.255] /A << /S /GoTo /D <636974652E547265> >> >>}
\MusicPDFWriteObject{6211}{<< /Type /Annot /Subtype /Link /Border [0 0 0] /H /I /C [0 1 0] /Rect [487.766 111.387 506.74 120.255] /A << /S /GoTo /D <636974652E547265> >> >>}
\MusicPDFWriteObject{6212}{<< /Type /Annot /Subtype /Link /Border [0 0 0] /H /I /C [1 0 0] /Rect [129.29 82.596 161.108 94.549] /A << /S /GoTo /D <74686D2E432E312E352E35> >> >>}
\MusicPDFWriteObject{6213}{<< /Type /Annot /Subtype /Link /Border [0 0 0] /H /I /C [0 1 0] /Rect [268.189 634.52 285.51 643.388] /A << /S /GoTo /D <636974652E534E> >> >>}
\MusicPDFWriteObject{6214}{<< /Type /Annot /Subtype /Link /Border [0 0 0] /H /I /C [0 1 0] /Rect [73.715 573.701 98.775 582.568] /A << /S /GoTo /D <636974652E53636861> >> >>}
\MusicPDFWriteObject{6215}{<< /Type /Annot /Subtype /Link /Border [0 0 0] /H /I /C [0 1 0] /Rect [298.419 573.701 317.392 582.568] /A << /S /GoTo /D <636974652E547265> >> >>}
\MusicPDFWriteObject{6216}{<< /Type /Annot /Subtype /Link /Border [0 0 0] /H /I /C [0 1 0] /Rect [211.923 547.399 236.983 556.267] /A << /S /GoTo /D <636974652E53636861> >> >>}
\MusicPDFWriteObject{6217}{<< /Type /Annot /Subtype /Link /Border [0 0 0] /H /I /C [0 1 0] /Rect [315.03 499.73 340.09 508.598] /A << /S /GoTo /D <636974652E53636861> >> >>}
\MusicPDFWriteObject{6218}{<< /Type /Annot /Subtype /Link /Border [0 0 0] /H /I /C [0 1 0] /Rect [189.014 486.58 207.988 495.447] /A << /S /GoTo /D <636974652E547265> >> >>}
\MusicPDFWriteObject{6219}{<< /Type /Annot /Subtype /Link /Border [0 0 0] /H /I /C [0 1 0] /Rect [121.01 473.429 152.634 482.297] /A << /S /GoTo /D <636974652E47726F2D62> >> >>}
\MusicPDFWriteObject{6220}{<< /Type /Annot /Subtype /Link /Border [0 0 0] /H /I /C [0 1 0] /Rect [126.209 447.127 157.833 455.995] /A << /S /GoTo /D <636974652E47726F2D62> >> >>}
\MusicPDFWriteObject{6221}{<< /Type /Annot /Subtype /Link /Border [0 0 0] /H /I /C [0 1 0] /Rect [324.283 447.127 373.324 455.995] /A << /S /GoTo /D <636974652E426F752D545653> >> >>}
\MusicPDFWriteObject{6222}{<< /Type /Annot /Subtype /Link /Border [0 0 0] /H /I /C [0 1 0] /Rect [431.662 447.127 449.86 455.995] /A << /S /GoTo /D <636974652E4A6172> >> >>}
\MusicPDFWriteObject{6223}{<< /Type /Annot /Subtype /Link /Border [0 0 0] /H /I /C [0 1 0] /Rect [509.19 447.127 543.159 455.995] /A << /S /GoTo /D <636974652E436E772D61> >> >>}
\MusicPDFWriteObject{6224}{<< /Type /Annot /Subtype /Link /Border [0 0 0] /H /I /C [0 1 0] /Rect [130.212 433.976 149.317 442.844] /A << /S /GoTo /D <636974652E4E42> >> >>}
\MusicPDFWriteObject{6225}{<< /Type /Annot /Subtype /Link /Border [0 0 0] /H /I /C [1 0 0] /Rect [304.014 637.488 336.775 649.44] /A << /S /GoTo /D <6571756174696F6E2E432E312E352E31> >> >>}
\MusicPDFWriteObject{6226}{<< /Type /Annot /Subtype /Link /Border [0 0 0] /H /I /C [1 0 0] /Rect [390.17 323.781 423.205 335.734] /A << /S /GoTo /D <74686D2E432E312E342E36> >> >>}
\MusicPDFWriteObject{6227}{<< /Type /Annot /Subtype /Link /Border [0 0 0] /H /I /C [0 1 0] /Rect [498.843 313.12 523.903 321.988] /A << /S /GoTo /D <636974652E53636861> >> >>}
\MusicPDFWriteObject{6228}{<< /Type /Annot /Subtype /Link /Border [0 0 0] /H /I /C [0 1 0] /Rect [343.799 286.819 372.56 295.686] /A << /S /GoTo /D <636974652E47726F2D66> >> >>}
\MusicPDFWriteObject{6229}{<< /Type /Annot /Subtype /Link /Border [0 0 0] /H /I /C [1 0 0] /Rect [432.975 203.188 466.011 215.141] /A << /S /GoTo /D <74686D2E432E312E332E34> >> >>}
\MusicPDFWriteObject{6230}{<< /Type /Annot /Subtype /Link /Border [0 0 0] /H /I /C [1 0 0] /Rect [226.043 190.59 259.079 201.439] /A << /S /GoTo /D <74686D2E432E312E352E32> >> >>}
\MusicPDFWriteObject{6231}{<< /Type /Annot /Subtype /Link /Border [0 0 0] /H /I /C [0 1 0] /Rect [203.364 166.226 228.424 175.094] /A << /S /GoTo /D <636974652E53636861> >> >>}
\MusicPDFWriteObject{6232}{<< /Type /Annot /Subtype /Link /Border [0 0 0] /H /I /C [0 1 0] /Rect [190.169 638.707 218.93 647.575] /A << /S /GoTo /D <636974652E47726F2D66> >> >>}
\MusicPDFWriteObject{6233}{<< /Type /Annot /Subtype /Link /Border [0 0 0] /H /I /C [0 1 0] /Rect [396.593 533.501 425.353 542.369] /A << /S /GoTo /D <636974652E47726F2D66> >> >>}
\MusicPDFWriteObject{6234}{<< /Type /Annot /Subtype /Link /Border [0 0 0] /H /I /C [0 1 0] /Rect [475.526 415.144 495.184 424.012] /A << /S /GoTo /D <636974652E4B47> >> >>}
\MusicPDFWriteObject{6235}{<< /Type /Annot /Subtype /Link /Border [0 0 0] /H /I /C [0 1 0] /Rect [307.85 401.993 332.91 410.861] /A << /S /GoTo /D <636974652E53636861> >> >>}
\MusicPDFWriteObject{6236}{<< /Type /Annot /Subtype /Link /Border [0 0 0] /H /I /C [0 1 0] /Rect [137.118 214.849 165.879 223.717] /A << /S /GoTo /D <636974652E47726F2D66> >> >>}
\MusicPDFWriteObject{6237}{<< /Type /Annot /Subtype /Link /Border [0 0 0] /H /I /C [0 1 0] /Rect [171.309 201.699 200.07 210.566] /A << /S /GoTo /D <636974652E47726F2D66> >> >>}
\MusicPDFWriteObject{6238}{<< /Type /Annot /Subtype /Link /Border [0 0 0] /H /I /C [0 1 0] /Rect [397.934 201.699 417.039 210.566] /A << /S /GoTo /D <636974652E4E42> >> >>}
\MusicPDFWriteObject{6239}{<< /Type /Annot /Subtype /Link /Border [0 0 0] /H /I /C [1 0 0] /Rect [124.706 175.397 157.132 184.251] /A << /S /GoTo /D <74686D2E432E312E352E35> >> >>}
\MusicPDFWriteObject{6240}{<< /Type /Annot /Subtype /Link /Border [0 0 0] /H /I /C [1 0 0] /Rect [412.752 615.243 445.787 626.034] /A << /S /GoTo /D <74686D2E432E312E352E34> >> >>}
\MusicPDFWriteObject{6241}{<< /Type /Annot /Subtype /Link /Border [0 0 0] /H /I /C [1 0 0] /Rect [88.936 563.582 94.899 573.95] /A << /S /GoTo /D <48666F6F746E6F74652E313130> >> >>}
\MusicPDFWriteObject{6242}{<< /Type /Annot /Subtype /Link /Border [0 0 0] /H /I /C [1 0 0] /Rect [512.336 436.886 544.945 449.19] /A << /S /GoTo /D <74686D2E432E312E352E34> >> >>}
\MusicPDFWriteObject{6243}{<< /Type /Annot /Subtype /Link /Border [0 0 0] /H /I /C [0 1 0] /Rect [278.395 360.803 293.468 369.671] /A << /S /GoTo /D <636974652E5061> >> >>}
\MusicPDFWriteObject{6244}{<< /Type /Annot /Subtype /Link /Border [0 0 0] /H /I /C [0 1 0] /Rect [73.826 347.652 94.113 356.52] /A << /S /GoTo /D <636974652E526574> >> >>}
\MusicPDFWriteObject{6245}{<< /Type /Annot /Subtype /Link /Border [0 0 0] /H /I /C [0 1 0] /Rect [252.932 345.716 288.084 356.52] /A << /S /GoTo /D <636974652E566F67742D31> >> >>}
\MusicPDFWriteObject{6246}{<< /Type /Annot /Subtype /Link /Border [0 0 0] /H /I /C [0 1 0] /Rect [385.227 345.716 420.379 356.52] /A << /S /GoTo /D <636974652E566F67742D31> >> >>}
\MusicPDFWriteObject{6247}{<< /Type /Annot /Subtype /Link /Border [0 0 0] /H /I /C [1 0 0] /Rect [476.184 319.414 513.767 330.205] /A << /S /GoTo /D <74686D2E432E312E352E3131> >> >>}
\MusicPDFWriteObject{6248}{<< /Type /Annot /Subtype /Link /Border [0 0 0] /H /I /C [0 1 0] /Rect [73.715 242.446 104.226 251.314] /A << /S /GoTo /D <636974652E47726F2D65> >> >>}
\MusicPDFWriteObject{6249}{<< /Type /Annot /Subtype /Link /Border [0 0 0] /H /I /C [0 1 0] /Rect [282.907 167.358 300.047 176.226] /A << /S /GoTo /D <636974652E4453> >> >>}
\MusicPDFWriteObject{6250}{<< /Type /Annot /Subtype /Link /Border [0 0 0] /H /I /C [0 1 0] /Rect [319.722 141.057 344.782 149.924] /A << /S /GoTo /D <636974652E53636861> >> >>}
\MusicPDFWriteObject{6251}{<< /Type /Annot /Subtype /Link /Border [0 0 0] /H /I /C [1 0 0] /Rect [353.954 93.499 382.333 104.644] /A << /S /GoTo /D <74686D2E432E312E352E35> >> >>}
\MusicPDFWriteObject{6252}{<< /Type /Annot /Subtype /Link /Border [0 0 0] /H /I /C [1 0 0] /Rect [94.556 685.187 100.519 698.986] /A << /S /GoTo /D <48666F6F746E6F74652E313131> >> >>}
\MusicPDFWriteObject{6253}{<< /Type /Annot /Subtype /Link /Border [0 0 0] /H /I /C [1 0 0] /Rect [292.55 685.187 298.194 698.986] /A << /S /GoTo /D <48666F6F746E6F74652E313132> >> >>}
\MusicPDFWriteObject{6254}{<< /Type /Annot /Subtype /Link /Border [0 0 0] /H /I /C [1 0 0] /Rect [262.172 672.036 268.135 684.34] /A << /S /GoTo /D <48666F6F746E6F74652E313133> >> >>}
\MusicPDFWriteObject{6255}{<< /Type /Annot /Subtype /Link /Border [0 0 0] /H /I /C [1 0 0] /Rect [205.474 511.964 238.509 525.711] /A << /S /GoTo /D <74686D2E432E312E352E38> >> >>}
\MusicPDFWriteObject{6256}{<< /Type /Annot /Subtype /Link /Border [0 0 0] /H /I /C [0 1 0] /Rect [107.631 388.088 142.782 398.893] /A << /S /GoTo /D <636974652E566F67742D31> >> >>}
\MusicPDFWriteObject{6257}{<< /Type /Annot /Subtype /Link /Border [0 0 0] /H /I /C [1 0 0] /Rect [470.539 387.535 503.543 399.488] /A << /S /GoTo /D <74686D2E432E312E352E34> >> >>}
\MusicPDFWriteObject{6258}{<< /Type /Annot /Subtype /Link /Border [0 0 0] /H /I /C [1 0 0] /Rect [287.458 374.384 319.581 386.337] /A << /S /GoTo /D <74686D2E432E312E352E35> >> >>}
\MusicPDFWriteObject{6259}{<< /Type /Annot /Subtype /Link /Border [0 0 0] /H /I /C [1 0 0] /Rect [373.379 374.384 405.501 386.337] /A << /S /GoTo /D <74686D2E432E312E352E38> >> >>}
\MusicPDFWriteObject{6260}{<< /Type /Annot /Subtype /Link /Border [0 0 0] /H /I /C [0 1 0] /Rect [420.742 348.636 455.894 359.44] /A << /S /GoTo /D <636974652E566F67742D31> >> >>}
\MusicPDFWriteObject{6261}{<< /Type /Annot /Subtype /Link /Border [0 0 0] /H /I /C [1 0 0] /Rect [187.372 335.485 193.335 347.789] /A << /S /GoTo /D <48666F6F746E6F74652E313134> >> >>}
\MusicPDFWriteObject{6262}{<< /Type /Annot /Subtype /Link /Border [0 0 0] /H /I /C [1 0 0] /Rect [507.478 322.334 545.592 333.125] /A << /S /GoTo /D <74686D2E432E312E352E3131> >> >>}
\MusicPDFWriteObject{6263}{<< /Type /Annot /Subtype /Link /Border [0 0 0] /H /I /C [0 1 0] /Rect [419.817 209.779 450.328 218.647] /A << /S /GoTo /D <636974652E47726F2D65> >> >>}
\MusicPDFWriteObject{6264}{<< /Type /Annot /Subtype /Link /Border [0 0 0] /H /I /C [0 1 0] /Rect [376.407 196.629 406.918 205.496] /A << /S /GoTo /D <636974652E47726F2D65> >> >>}
\MusicPDFWriteObject{6265}{<< /Type /Annot /Subtype /Link /Border [0 0 0] /H /I /C [0 1 0] /Rect [325.367 93.499 355.585 102.541] /A << /S /GoTo /D <636974652E566F67742D31> >> >>}
\MusicPDFWriteObject{6266}{<< /Type /Annot /Subtype /Link /Border [0 0 0] /H /I /C [0 1 0] /Rect [89.219 628.34 114.279 637.207] /A << /S /GoTo /D <636974652E53636861> >> >>}
\MusicPDFWriteObject{6267}{<< /Type /Annot /Subtype /Link /Border [0 0 0] /H /I /C [0 1 0] /Rect [73.826 588.561 98.886 597.429] /A << /S /GoTo /D <636974652E53636861> >> >>}
\MusicPDFWriteObject{6268}{<< /Type /Annot /Subtype /Link /Border [0 0 0] /H /I /C [1 0 0] /Rect [71.004 514.462 107.702 526.417] /A << /S /GoTo /D <74686D2E432E312E352E3133> >> >>}
\MusicPDFWriteObject{6269}{<< /Type /Annot /Subtype /Link /Border [0 0 0] /H /I /C [0 1 0] /Rect [430.661 358.493 461.172 367.36] /A << /S /GoTo /D <636974652E47726F2D65> >> >>}
\MusicPDFWriteObject{6270}{<< /Type /Annot /Subtype /Link /Border [0 0 0] /H /I /C [1 0 0] /Rect [303.768 330.254 341.068 341.045] /A << /S /GoTo /D <74686D2E432E312E352E3131> >> >>}
\MusicPDFWriteObject{6271}{<< /Type /Annot /Subtype /Link /Border [0 0 0] /H /I /C [0 1 0] /Rect [84.329 317.104 123.562 327.908] /A << /S /GoTo /D <636974652E576E67722D32> >> >>}
\MusicPDFWriteObject{6272}{<< /Type /Annot /Subtype /Link /Border [0 0 0] /H /I /C [1 0 0] /Rect [172.432 303.399 209.838 315.352] /A << /S /GoTo /D <74686D2E432E312E352E3130> >> >>}
\MusicPDFWriteObject{6273}{<< /Type /Annot /Subtype /Link /Border [0 0 0] /H /I /C [1 0 0] /Rect [71.004 271.242 107.702 283.215] /A << /S /GoTo /D <74686D2E432E312E352E3131> >> >>}
\MusicPDFWriteObject{6274}{<< /Type /Annot /Subtype /Link /Border [0 0 0] /H /I /C [0 1 0] /Rect [294.156 247.432 319.217 256.299] /A << /S /GoTo /D <636974652E53636861> >> >>}
\MusicPDFWriteObject{6275}{<< /Type /Annot /Subtype /Link /Border [0 0 0] /H /I /C [0 1 0] /Rect [149.555 156.817 172.81 165.685] /A << /S /GoTo /D <636974652E466C4D> >> >>}
\MusicPDFWriteObject{6276}{<< /Type /Annot /Subtype /Link /Border [0 0 0] /H /I /C [1 0 0] /Rect [71.004 122.048 109.119 134.003] /A << /S /GoTo /D <74686D2E432E312E352E3131> >> >>}
\MusicPDFWriteObject{6277}{<< /Type /Annot /Subtype /Link /Border [0 0 0] /H /I /C [1 0 0] /Rect [345.544 618.879 377.361 630.832] /A << /S /GoTo /D <74686D2E432E312E352E35> >> >>}
\MusicPDFWriteObject{6278}{<< /Type /Annot /Subtype /Link /Border [0 0 0] /H /I /C [1 0 0] /Rect [106.814 591.526 139.484 604.532] /A << /S /GoTo /D <74686D2E432E312E352E34> >> >>}
\MusicPDFWriteObject{6279}{<< /Type /Annot /Subtype /Link /Border [0 0 0] /H /I /C [0 1 0] /Rect [252.2 542.465 277.26 551.333] /A << /S /GoTo /D <636974652E53636861> >> >>}
\MusicPDFWriteObject{6280}{<< /Type /Annot /Subtype /Link /Border [0 0 0] /H /I /C [1 0 0] /Rect [71.004 512.621 108.198 525.626] /A << /S /GoTo /D <74686D2E432E312E352E3130> >> >>}
\MusicPDFWriteObject{6281}{<< /Type /Annot /Subtype /Link /Border [0 0 0] /H /I /C [0 1 0] /Rect [246.93 516.163 277.441 525.031] /A << /S /GoTo /D <636974652E47726F2D65> >> >>}
\MusicPDFWriteObject{6282}{<< /Type /Annot /Subtype /Link /Border [0 0 0] /H /I /C [0 1 0] /Rect [268.25 238.222 293.31 247.09] /A << /S /GoTo /D <636974652E53636861> >> >>}
\MusicPDFWriteObject{6283}{<< /Type /Annot /Subtype /Link /Border [0 0 0] /H /I /C [1 0 0] /Rect [226.905 586.003 259.088 597.956] /A << /S /GoTo /D <74686D2E432E312E362E31> >> >>}
\MusicPDFWriteObject{6284}{<< /Type /Annot /Subtype /Link /Border [0 0 0] /H /I /C [0 1 0] /Rect [238.296 549.041 282.356 557.908] /A << /S /GoTo /D <636974652E424B4D634B> >> >>}
\MusicPDFWriteObject{6285}{<< /Type /Annot /Subtype /Link /Border [0 0 0] /H /I /C [0 1 0] /Rect [119.38 372.537 144.44 381.405] /A << /S /GoTo /D <636974652E53636861> >> >>}
\MusicPDFWriteObject{6286}{<< /Type /Annot /Subtype /Link /Border [0 0 0] /H /I /C [0 1 0] /Rect [144.275 322.959 169.336 331.827] /A << /S /GoTo /D <636974652E53636861> >> >>}
\MusicPDFWriteObject{6287}{<< /Type /Annot /Subtype /Link /Border [0 0 0] /H /I /C [0 1 0] /Rect [518.113 247.079 543.173 255.947] /A << /S /GoTo /D <636974652E53636861> >> >>}
\MusicPDFWriteObject{6288}{<< /Type /Annot /Subtype /Link /Border [0 0 0] /H /I /C [0 1 0] /Rect [486.134 220.777 511.195 229.645] /A << /S /GoTo /D <636974652E53636861> >> >>}
\MusicPDFWriteObject{6289}{<< /Type /Annot /Subtype /Link /Border [0 0 0] /H /I /C [0 1 0] /Rect [349.037 207.627 374.097 216.494] /A << /S /GoTo /D <636974652E53636861> >> >>}
\MusicPDFWriteObject{6290}{<< /Type /Annot /Subtype /Link /Border [0 0 0] /H /I /C [0 1 0] /Rect [138.014 194.476 163.074 203.344] /A << /S /GoTo /D <636974652E53636861> >> >>}
\MusicPDFWriteObject{6291}{<< /Type /Annot /Subtype /Link /Border [0 0 0] /H /I /C [0 1 0] /Rect [217.755 85.086 242.815 93.954] /A << /S /GoTo /D <636974652E53636861> >> >>}
\MusicPDFWriteObject{6292}{<< /Type /Annot /Subtype /Link /Border [0 0 0] /H /I /C [0 1 0] /Rect [339.231 536.421 364.291 545.289] /A << /S /GoTo /D <636974652E53636861> >> >>}
\MusicPDFWriteObject{6293}{<< /Type /Annot /Subtype /Link /Border [0 0 0] /H /I /C [0 1 0] /Rect [73.77 322.365 98.831 331.233] /A << /S /GoTo /D <636974652E53636861> >> >>}
\MusicPDFWriteObject{6294}{<< /Type /Annot /Subtype /Link /Border [0 0 0] /H /I /C [1 0 0] /Rect [87.051 159.652 119.477 170.443] /A << /S /GoTo /D <73756273756273656374696F6E2E432E322E342E35> >> >>}
\MusicPDFWriteObject{6295}{<< /Type /Annot /Subtype /Link /Border [0 0 0] /H /I /C [1 0 0] /Rect [295.401 684.633 327.828 696.586] /A << /S /GoTo /D <6571756174696F6E2E432E322E312E34> >> >>}
\MusicPDFWriteObject{6296}{<< /Type /Annot /Subtype /Link /Border [0 0 0] /H /I /C [0 1 0] /Rect [291.259 561.983 316.319 570.851] /A << /S /GoTo /D <636974652E53636861> >> >>}
\MusicPDFWriteObject{6297}{<< /Type /Annot /Subtype /Link /Border [0 0 0] /H /I /C [0 1 0] /Rect [518.157 561.983 543.217 570.851] /A << /S /GoTo /D <636974652E53636861> >> >>}
\MusicPDFWriteObject{6298}{<< /Type /Annot /Subtype /Link /Border [0 0 0] /H /I /C [1 0 0] /Rect [268.217 546.342 274.18 559.2] /A << /S /GoTo /D <48666F6F746E6F74652E313135> >> >>}
\MusicPDFWriteObject{6299}{<< /Type /Annot /Subtype /Link /Border [0 0 0] /H /I /C [1 0 0] /Rect [377.425 412.747 410.46 423.537] /A << /S /GoTo /D <74686D2E432E322E312E37> >> >>}
\MusicPDFWriteObject{6300}{<< /Type /Annot /Subtype /Link /Border [0 0 0] /H /I /C [0 1 0] /Rect [158.644 388.382 183.704 397.25] /A << /S /GoTo /D <636974652E53636861> >> >>}
\MusicPDFWriteObject{6301}{<< /Type /Annot /Subtype /Link /Border [0 0 0] /H /I /C [0 1 0] /Rect [261.297 233.915 286.357 242.783] /A << /S /GoTo /D <636974652E53636861> >> >>}
\MusicPDFWriteObject{6302}{<< /Type /Annot /Subtype /Link /Border [0 0 0] /H /I /C [0 1 0] /Rect [73.77 207.614 98.831 216.482] /A << /S /GoTo /D <636974652E53636861> >> >>}
\MusicPDFWriteObject{6303}{<< /Type /Annot /Subtype /Link /Border [0 0 0] /H /I /C [0 1 0] /Rect [464.22 595.068 489.28 603.936] /A << /S /GoTo /D <636974652E53636861> >> >>}
\MusicPDFWriteObject{6304}{<< /Type /Annot /Subtype /Link /Border [0 0 0] /H /I /C [1 0 0] /Rect [297.318 566.83 329.744 577.677] /A << /S /GoTo /D <74686D2E432E322E312E35> >> >>}
\MusicPDFWriteObject{6305}{<< /Type /Annot /Subtype /Link /Border [0 0 0] /H /I /C [0 1 0] /Rect [118.042 555.616 143.103 564.484] /A << /S /GoTo /D <636974652E53636861> >> >>}
\MusicPDFWriteObject{6306}{<< /Type /Annot /Subtype /Link /Border [0 0 0] /H /I /C [0 1 0] /Rect [243.852 377.554 268.912 386.422] /A << /S /GoTo /D <636974652E53636861> >> >>}
\MusicPDFWriteObject{6307}{<< /Type /Annot /Subtype /Link /Border [0 0 0] /H /I /C [0 1 0] /Rect [109.121 364.404 158.162 373.272] /A << /S /GoTo /D <636974652E426F752D545653> >> >>}
\MusicPDFWriteObject{6308}{<< /Type /Annot /Subtype /Link /Border [0 0 0] /H /I /C [1 0 0] /Rect [149.235 260.286 158.684 271.076] /A << /S /GoTo /D <617070656E6469782E44> >> >>}
\MusicPDFWriteObject{6309}{<< /Type /Annot /Subtype /Link /Border [0 0 0] /H /I /C [1 0 0] /Rect [353.487 173.182 385.914 185.135] /A << /S /GoTo /D <6571756174696F6E2E432E322E322E34> >> >>}
\MusicPDFWriteObject{6310}{<< /Type /Annot /Subtype /Link /Border [0 0 0] /H /I /C [1 0 0] /Rect [176.879 684.633 209.306 696.586] /A << /S /GoTo /D <6571756174696F6E2E432E322E322E35> >> >>}
\MusicPDFWriteObject{6311}{<< /Type /Annot /Subtype /Link /Border [0 0 0] /H /I /C [1 0 0] /Rect [71.004 624.175 103.43 636.127] /A << /S /GoTo /D <74686D2E432E322E322E31> >> >>}
\MusicPDFWriteObject{6312}{<< /Type /Annot /Subtype /Link /Border [0 0 0] /H /I /C [0 1 0] /Rect [195.7 626.665 220.76 635.532] /A << /S /GoTo /D <636974652E53636861> >> >>}
\MusicPDFWriteObject{6313}{<< /Type /Annot /Subtype /Link /Border [0 0 0] /H /I /C [0 1 0] /Rect [368.229 613.514 393.289 622.382] /A << /S /GoTo /D <636974652E53636861> >> >>}
\MusicPDFWriteObject{6314}{<< /Type /Annot /Subtype /Link /Border [0 0 0] /H /I /C [0 1 0] /Rect [468.7 613.514 487.804 622.382] /A << /S /GoTo /D <636974652E4E42> >> >>}
\MusicPDFWriteObject{6315}{<< /Type /Annot /Subtype /Link /Border [0 0 0] /H /I /C [0 1 0] /Rect [212.496 256.046 237.556 264.914] /A << /S /GoTo /D <636974652E53636861> >> >>}
\MusicPDFWriteObject{6316}{<< /Type /Annot /Subtype /Link /Border [0 0 0] /H /I /C [0 1 0] /Rect [518.135 229.744 543.195 238.612] /A << /S /GoTo /D <636974652E53636861> >> >>}
\MusicPDFWriteObject{6317}{<< /Type /Annot /Subtype /Link /Border [0 0 0] /H /I /C [0 1 0] /Rect [73.715 203.443 98.775 212.311] /A << /S /GoTo /D <636974652E53636861> >> >>}
\MusicPDFWriteObject{6318}{<< /Type /Annot /Subtype /Link /Border [0 0 0] /H /I /C [0 1 0] /Rect [73.77 163.99 91.969 172.858] /A << /S /GoTo /D <636974652E4A6172> >> >>}
\MusicPDFWriteObject{6319}{<< /Type /Annot /Subtype /Link /Border [0 0 0] /H /I /C [0 1 0] /Rect [182.627 135.752 211.415 146.557] /A << /S /GoTo /D <636974652E566F696774> >> >>}
\MusicPDFWriteObject{6320}{<< /Type /Annot /Subtype /Link /Border [0 0 0] /H /I /C [0 1 0] /Rect [89.288 111.387 114.348 120.255] /A << /S /GoTo /D <636974652E53636861> >> >>}
\MusicPDFWriteObject{6321}{<< /Type /Annot /Subtype /Link /Border [0 0 0] /H /I /C [0 1 0] /Rect [245.546 98.237 270.606 107.105] /A << /S /GoTo /D <636974652E53636861> >> >>}
\MusicPDFWriteObject{6322}{<< /Type /Annot /Subtype /Link /Border [0 0 0] /H /I /C [0 1 0] /Rect [518.209 673.973 543.27 682.84] /A << /S /GoTo /D <636974652E53636861> >> >>}
\MusicPDFWriteObject{6323}{<< /Type /Annot /Subtype /Link /Border [0 0 0] /H /I /C [1 0 0] /Rect [533.821 606.282 543.27 617.073] /A << /S /GoTo /D <617070656E6469782E44> >> >>}
\MusicPDFWriteObject{6324}{<< /Type /Annot /Subtype /Link /Border [0 0 0] /H /I /C [0 1 0] /Rect [346.764 595.068 371.824 603.936] /A << /S /GoTo /D <636974652E53636861> >> >>}
\MusicPDFWriteObject{6325}{<< /Type /Annot /Subtype /Link /Border [0 0 0] /H /I /C [1 0 0] /Rect [218.685 474.981 250.503 486.4] /A << /S /GoTo /D <74686D2E432E322E322E32> >> >>}
\MusicPDFWriteObject{6326}{<< /Type /Annot /Subtype /Link /Border [0 0 0] /H /I /C [1 0 0] /Rect [355.178 461.277 387.969 473.23] /A << /S /GoTo /D <6571756174696F6E2E432E322E322E36> >> >>}
\MusicPDFWriteObject{6327}{<< /Type /Annot /Subtype /Link /Border [0 0 0] /H /I /C [0 1 0] /Rect [205.121 450.617 230.181 459.485] /A << /S /GoTo /D <636974652E53636861> >> >>}
\MusicPDFWriteObject{6328}{<< /Type /Annot /Subtype /Link /Border [0 0 0] /H /I /C [1 0 0] /Rect [493.586 422.378 526.408 433.304] /A << /S /GoTo /D <73756273756273656374696F6E2E432E322E332E31> >> >>}
\MusicPDFWriteObject{6329}{<< /Type /Annot /Subtype /Link /Border [0 0 0] /H /I /C [1 0 0] /Rect [123.988 409.228 157.023 420.018] /A << /S /GoTo /D <74686D2E432E322E322E32> >> >>}
\MusicPDFWriteObject{6330}{<< /Type /Annot /Subtype /Link /Border [0 0 0] /H /I /C [1 0 0] /Rect [322.194 342.921 355.229 354.873] /A << /S /GoTo /D <73756273756273656374696F6E2E432E322E332E31> >> >>}
\MusicPDFWriteObject{6331}{<< /Type /Annot /Subtype /Link /Border [0 0 0] /H /I /C [1 0 0] /Rect [380.784 342.921 413.82 354.873] /A << /S /GoTo /D <73756273756273656374696F6E2E432E322E332E32> >> >>}
\MusicPDFWriteObject{6332}{<< /Type /Annot /Subtype /Link /Border [0 0 0] /H /I /C [1 0 0] /Rect [394.649 264.569 426.467 275.36] /A << /S /GoTo /D <74686D2E432E322E322E32> >> >>}
\MusicPDFWriteObject{6333}{<< /Type /Annot /Subtype /Link /Border [0 0 0] /H /I /C [1 0 0] /Rect [272.612 171.961 297.745 183.913] /A << /S /GoTo /D <73756273656374696F6E2E432E312E32> >> >>}
\MusicPDFWriteObject{6334}{<< /Type /Annot /Subtype /Link /Border [0 0 0] /H /I /C [0 1 0] /Rect [456.396 621.369 475.37 630.237] /A << /S /GoTo /D <636974652E547265> >> >>}
\MusicPDFWriteObject{6335}{<< /Type /Annot /Subtype /Link /Border [0 0 0] /H /I /C [0 1 0] /Rect [454.594 581.917 473.567 590.785] /A << /S /GoTo /D <636974652E547265> >> >>}
\MusicPDFWriteObject{6336}{<< /Type /Annot /Subtype /Link /Border [0 0 0] /H /I /C [0 1 0] /Rect [151.746 379.165 176.806 388.033] /A << /S /GoTo /D <636974652E53636861> >> >>}
\MusicPDFWriteObject{6337}{<< /Type /Annot /Subtype /Link /Border [0 0 0] /H /I /C [1 0 0] /Rect [399.697 376.675 409.146 388.628] /A << /S /GoTo /D <617070656E6469782E44> >> >>}
\MusicPDFWriteObject{6338}{<< /Type /Annot /Subtype /Link /Border [0 0 0] /H /I /C [0 1 0] /Rect [278.71 352.864 297.684 361.732] /A << /S /GoTo /D <636974652E547265> >> >>}
\MusicPDFWriteObject{6339}{<< /Type /Annot /Subtype /Link /Border [0 0 0] /H /I /C [0 1 0] /Rect [413.042 352.864 438.102 361.732] /A << /S /GoTo /D <636974652E53636861> >> >>}
\MusicPDFWriteObject{6340}{<< /Type /Annot /Subtype /Link /Border [0 0 0] /H /I /C [1 0 0] /Rect [160.038 337.223 184.718 349.176] /A << /S /GoTo /D <73756273656374696F6E2E432E322E34> >> >>}
\MusicPDFWriteObject{6341}{<< /Type /Annot /Subtype /Link /Border [0 0 0] /H /I /C [1 0 0] /Rect [371.692 226.526 404.727 238.499] /A << /S /GoTo /D <6571756174696F6E2E432E322E332E33> >> >>}
\MusicPDFWriteObject{6342}{<< /Type /Annot /Subtype /Link /Border [0 0 0] /H /I /C [0 1 0] /Rect [396.974 215.866 415.947 224.734] /A << /S /GoTo /D <636974652E547265> >> >>}
\MusicPDFWriteObject{6343}{<< /Type /Annot /Subtype /Link /Border [0 0 0] /H /I /C [1 0 0] /Rect [364.389 187.074 396.816 199.027] /A << /S /GoTo /D <6571756174696F6E2E432E322E332E33> >> >>}
\MusicPDFWriteObject{6344}{<< /Type /Annot /Subtype /Link /Border [0 0 0] /H /I /C [0 1 0] /Rect [136.386 163.263 155.359 172.131] /A << /S /GoTo /D <636974652E547265> >> >>}
\MusicPDFWriteObject{6345}{<< /Type /Annot /Subtype /Link /Border [0 0 0] /H /I /C [0 1 0] /Rect [387.755 112.267 412.815 121.135] /A << /S /GoTo /D <636974652E53636861> >> >>}
\MusicPDFWriteObject{6346}{<< /Type /Annot /Subtype /Link /Border [0 0 0] /H /I /C [1 0 0] /Rect [172.721 684.633 178.684 697.491] /A << /S /GoTo /D <48666F6F746E6F74652E313136> >> >>}
\MusicPDFWriteObject{6347}{<< /Type /Annot /Subtype /Link /Border [0 0 0] /H /I /C [0 1 0] /Rect [495.813 687.123 514.011 695.991] /A << /S /GoTo /D <636974652E4A6172> >> >>}
\MusicPDFWriteObject{6348}{<< /Type /Annot /Subtype /Link /Border [0 0 0] /H /I /C [1 0 0] /Rect [312.848 671.482 318.811 684.34] /A << /S /GoTo /D <48666F6F746E6F74652E313137> >> >>}
\MusicPDFWriteObject{6349}{<< /Type /Annot /Subtype /Link /Border [0 0 0] /H /I /C [0 1 0] /Rect [492.794 673.973 510.992 682.84] /A << /S /GoTo /D <636974652E4A6172> >> >>}
\MusicPDFWriteObject{6350}{<< /Type /Annot /Subtype /Link /Border [0 0 0] /H /I /C [0 1 0] /Rect [236.575 634.471 261.635 643.339] /A << /S /GoTo /D <636974652E53636861> >> >>}
\MusicPDFWriteObject{6351}{<< /Type /Annot /Subtype /Link /Border [0 0 0] /H /I /C [1 0 0] /Rect [71.004 579.378 103.43 591.331] /A << /S /GoTo /D <74686D2E432E322E322E32> >> >>}
\MusicPDFWriteObject{6352}{<< /Type /Annot /Subtype /Link /Border [0 0 0] /H /I /C [0 1 0] /Rect [306.183 482.711 324.382 491.579] /A << /S /GoTo /D <636974652E4A6172> >> >>}
\MusicPDFWriteObject{6353}{<< /Type /Annot /Subtype /Link /Border [0 0 0] /H /I /C [1 0 0] /Rect [500.775 461.972 533.81 473.925] /A << /S /GoTo /D <74686D2E432E322E332E32> >> >>}
\MusicPDFWriteObject{6354}{<< /Type /Annot /Subtype /Link /Border [0 0 0] /H /I /C [0 1 0] /Rect [483.495 425.01 508.555 433.878] /A << /S /GoTo /D <636974652E53636861> >> >>}
\MusicPDFWriteObject{6355}{<< /Type /Annot /Subtype /Link /Border [0 0 0] /H /I /C [0 1 0] /Rect [329.025 411.859 354.085 420.727] /A << /S /GoTo /D <636974652E53636861> >> >>}
\MusicPDFWriteObject{6356}{<< /Type /Annot /Subtype /Link /Border [0 0 0] /H /I /C [0 1 0] /Rect [413.388 411.859 431.587 420.727] /A << /S /GoTo /D <636974652E4A6172> >> >>}
\MusicPDFWriteObject{6357}{<< /Type /Annot /Subtype /Link /Border [0 0 0] /H /I /C [0 1 0] /Rect [256.85 398.709 281.91 407.576] /A << /S /GoTo /D <636974652E53636861> >> >>}
\MusicPDFWriteObject{6358}{<< /Type /Annot /Subtype /Link /Border [0 0 0] /H /I /C [1 0 0] /Rect [496.818 338.517 528.758 350.47] /A << /S /GoTo /D <6571756174696F6E2E432E322E332E31> >> >>}
\MusicPDFWriteObject{6359}{<< /Type /Annot /Subtype /Link /Border [0 0 0] /H /I /C [1 0 0] /Rect [103.231 194.809 135.78 206.782] /A << /S /GoTo /D <6571756174696F6E2E432E322E332E38> >> >>}
\MusicPDFWriteObject{6360}{<< /Type /Annot /Subtype /Link /Border [0 0 0] /H /I /C [1 0 0] /Rect [464.899 150.259 497.934 162.212] /A << /S /GoTo /D <6571756174696F6E2E432E322E332E37> >> >>}
\MusicPDFWriteObject{6361}{<< /Type /Annot /Subtype /Link /Border [0 0 0] /H /I /C [0 1 0] /Rect [524.186 152.749 543.159 161.617] /A << /S /GoTo /D <636974652E547265> >> >>}
\MusicPDFWriteObject{6362}{<< /Type /Annot /Subtype /Link /Border [0 0 0] /H /I /C [0 1 0] /Rect [450.225 459.621 469.198 468.489] /A << /S /GoTo /D <636974652E547265> >> >>}
\MusicPDFWriteObject{6363}{<< /Type /Annot /Subtype /Link /Border [0 0 0] /H /I /C [1 0 0] /Rect [317.381 443.98 349.807 455.933] /A << /S /GoTo /D <6571756174696F6E2E432E322E332E37> >> >>}
\MusicPDFWriteObject{6364}{<< /Type /Annot /Subtype /Link /Border [0 0 0] /H /I /C [0 1 0] /Rect [201.357 380.717 220.331 389.585] /A << /S /GoTo /D <636974652E547265> >> >>}
\MusicPDFWriteObject{6365}{<< /Type /Annot /Subtype /Link /Border [0 0 0] /H /I /C [0 1 0] /Rect [439.656 129.612 457.855 138.48] /A << /S /GoTo /D <636974652E4A6172> >> >>}
\MusicPDFWriteObject{6366}{<< /Type /Annot /Subtype /Link /Border [0 0 0] /H /I /C [1 0 0] /Rect [520.925 127.122 526.888 139.979] /A << /S /GoTo /D <48666F6F746E6F74652E313138> >> >>}
\MusicPDFWriteObject{6367}{<< /Type /Annot /Subtype /Link /Border [0 0 0] /H /I /C [0 1 0] /Rect [374.044 85.086 389.845 92.579] /A << /S /GoTo /D <636974652E4A6172> >> >>}
\MusicPDFWriteObject{6368}{<< /Type /Annot /Subtype /Link /Border [0 0 0] /H /I /C [0 1 0] /Rect [108.3 545.49 127.273 554.357] /A << /S /GoTo /D <636974652E547265> >> >>}
\MusicPDFWriteObject{6369}{<< /Type /Annot /Subtype /Link /Border [0 0 0] /H /I /C [0 1 0] /Rect [186.494 545.49 211.554 554.357] /A << /S /GoTo /D <636974652E53636861> >> >>}
\MusicPDFWriteObject{6370}{<< /Type /Annot /Subtype /Link /Border [0 0 0] /H /I /C [0 1 0] /Rect [73.826 443.308 92.799 452.176] /A << /S /GoTo /D <636974652E547265> >> >>}
\MusicPDFWriteObject{6371}{<< /Type /Annot /Subtype /Link /Border [0 0 0] /H /I /C [0 1 0] /Rect [336.955 430.158 355.928 439.025] /A << /S /GoTo /D <636974652E547265> >> >>}
\MusicPDFWriteObject{6372}{<< /Type /Annot /Subtype /Link /Border [0 0 0] /H /I /C [1 0 0] /Rect [187.847 401.366 220.273 414.223] /A << /S /GoTo /D <74686D2E432E322E332E35> >> >>}
\MusicPDFWriteObject{6373}{<< /Type /Annot /Subtype /Link /Border [0 0 0] /H /I /C [0 1 0] /Rect [397.011 403.856 415.21 412.724] /A << /S /GoTo /D <636974652E4A6172> >> >>}
\MusicPDFWriteObject{6374}{<< /Type /Annot /Subtype /Link /Border [0 0 0] /H /I /C [1 0 0] /Rect [515.035 401.366 524.968 414.223] /A << /S /GoTo /D <48666F6F746E6F74652E313139> >> >>}
\MusicPDFWriteObject{6375}{<< /Type /Annot /Subtype /Link /Border [0 0 0] /H /I /C [0 1 0] /Rect [107.413 390.705 132.473 399.573] /A << /S /GoTo /D <636974652E53636861> >> >>}
\MusicPDFWriteObject{6376}{<< /Type /Annot /Subtype /Link /Border [0 0 0] /H /I /C [1 0 0] /Rect [116.276 289.905 125.877 300.696] /A << /S /GoTo /D <617070656E6469782E44> >> >>}
\MusicPDFWriteObject{6377}{<< /Type /Annot /Subtype /Link /Border [0 0 0] /H /I /C [1 0 0] /Rect [190.277 249.899 227.684 261.852] /A << /S /GoTo /D <6571756174696F6E2E432E322E332E3135> >> >>}
\MusicPDFWriteObject{6378}{<< /Type /Annot /Subtype /Link /Border [0 0 0] /H /I /C [0 1 0] /Rect [466.894 252.389 485.868 261.257] /A << /S /GoTo /D <636974652E547265> >> >>}
\MusicPDFWriteObject{6379}{<< /Type /Annot /Subtype /Link /Border [0 0 0] /H /I /C [0 1 0] /Rect [273.74 163.861 292.713 172.729] /A << /S /GoTo /D <636974652E547265> >> >>}
\MusicPDFWriteObject{6380}{<< /Type /Annot /Subtype /Link /Border [0 0 0] /H /I /C [1 0 0] /Rect [131.951 148.221 157.084 160.173] /A << /S /GoTo /D <73756273656374696F6E2E432E322E34> >> >>}
\MusicPDFWriteObject{6381}{<< /Type /Annot /Subtype /Link /Border [0 0 0] /H /I /C [0 1 0] /Rect [182.816 106.447 201.789 115.315] /A << /S /GoTo /D <636974652E547265> >> >>}
\MusicPDFWriteObject{6382}{<< /Type /Annot /Subtype /Link /Border [0 0 0] /H /I /C [1 0 0] /Rect [128.824 618.879 138.508 630.832] /A << /S /GoTo /D <617070656E6469782E44> >> >>}
\MusicPDFWriteObject{6383}{<< /Type /Annot /Subtype /Link /Border [0 0 0] /H /I /C [1 0 0] /Rect [198.215 618.879 230.977 630.832] /A << /S /GoTo /D <74686D2E432E322E322E32> >> >>}
\MusicPDFWriteObject{6384}{<< /Type /Annot /Subtype /Link /Border [0 0 0] /H /I /C [0 1 0] /Rect [518.168 581.917 543.228 590.785] /A << /S /GoTo /D <636974652E53636861> >> >>}
\MusicPDFWriteObject{6385}{<< /Type /Annot /Subtype /Link /Border [0 0 0] /H /I /C [1 0 0] /Rect [186.607 553.126 219.572 565.078] /A << /S /GoTo /D <74686D2E442E312E332E37> >> >>}
\MusicPDFWriteObject{6386}{<< /Type /Annot /Subtype /Link /Border [0 0 0] /H /I /C [1 0 0] /Rect [259.36 553.126 292.325 565.078] /A << /S /GoTo /D <74686D2E442E322E322E32> >> >>}
\MusicPDFWriteObject{6387}{<< /Type /Annot /Subtype /Link /Border [0 0 0] /H /I /C [0 1 0] /Rect [318.64 555.616 337.614 564.484] /A << /S /GoTo /D <636974652E547265> >> >>}
\MusicPDFWriteObject{6388}{<< /Type /Annot /Subtype /Link /Border [0 0 0] /H /I /C [1 0 0] /Rect [516.917 526.824 541.846 538.777] /A << /S /GoTo /D <73756273656374696F6E2E432E322E34> >> >>}
\MusicPDFWriteObject{6389}{<< /Type /Annot /Subtype /Link /Border [0 0 0] /H /I /C [0 1 0] /Rect [451.569 482.414 483.194 491.282] /A << /S /GoTo /D <636974652E47726F2D62> >> >>}
\MusicPDFWriteObject{6390}{<< /Type /Annot /Subtype /Link /Border [0 0 0] /H /I /C [0 1 0] /Rect [511.535 482.414 543.159 491.282] /A << /S /GoTo /D <636974652E47726F2D62> >> >>}
\MusicPDFWriteObject{6391}{<< /Type /Annot /Subtype /Link /Border [0 0 0] /H /I /C [0 1 0] /Rect [257.613 203.443 282.673 212.311] /A << /S /GoTo /D <636974652E53636861> >> >>}
\MusicPDFWriteObject{6392}{<< /Type /Annot /Subtype /Link /Border [0 0 0] /H /I /C [0 1 0] /Rect [189.677 177.141 207.876 186.009] /A << /S /GoTo /D <636974652E4A6172> >> >>}
\MusicPDFWriteObject{6393}{<< /Type /Annot /Subtype /Link /Border [0 0 0] /H /I /C [0 1 0] /Rect [363.586 177.141 395.21 186.009] /A << /S /GoTo /D <636974652E47726F2D62> >> >>}
\MusicPDFWriteObject{6394}{<< /Type /Annot /Subtype /Link /Border [0 0 0] /H /I /C [0 1 0] /Rect [87.932 163.99 119.556 172.858] /A << /S /GoTo /D <636974652E47726F2D62> >> >>}
\MusicPDFWriteObject{6395}{<< /Type /Annot /Subtype /Link /Border [0 0 0] /H /I /C [0 1 0] /Rect [487.224 137.689 512.284 146.557] /A << /S /GoTo /D <636974652E53636861> >> >>}
\MusicPDFWriteObject{6396}{<< /Type /Annot /Subtype /Link /Border [0 0 0] /H /I /C [0 1 0] /Rect [161.373 673.973 180.346 682.84] /A << /S /GoTo /D <636974652E547265> >> >>}
\MusicPDFWriteObject{6397}{<< /Type /Annot /Subtype /Link /Border [0 0 0] /H /I /C [1 0 0] /Rect [493.566 658.332 530.972 670.284] /A << /S /GoTo /D <6571756174696F6E2E432E322E332E3139> >> >>}
\MusicPDFWriteObject{6398}{<< /Type /Annot /Subtype /Link /Border [0 0 0] /H /I /C [1 0 0] /Rect [141.334 608.219 151.267 618.586] /A << /S /GoTo /D <48666F6F746E6F74652E313230> >> >>}
\MusicPDFWriteObject{6399}{<< /Type /Annot /Subtype /Link /Border [0 0 0] /H /I /C [1 0 0] /Rect [380.684 473.713 413.111 484.504] /A << /S /GoTo /D <74686D2E432E312E352E31> >> >>}
\MusicPDFWriteObject{6400}{<< /Type /Annot /Subtype /Link /Border [0 0 0] /H /I /C [1 0 0] /Rect [203.197 456.522 244.587 465.376] /A << /S /GoTo /D <74686D2E432E322E332E3130> >> >>}
\MusicPDFWriteObject{6401}{<< /Type /Annot /Subtype /Link /Border [0 0 0] /H /I /C [1 0 0] /Rect [396.951 316.88 429.986 327.727] /A << /S /GoTo /D <74686D2E432E312E352E31> >> >>}
\MusicPDFWriteObject{6402}{<< /Type /Annot /Subtype /Link /Border [0 0 0] /H /I /C [0 1 0] /Rect [271.753 168.18 320.793 177.048] /A << /S /GoTo /D <636974652E426F752D545653> >> >>}
\MusicPDFWriteObject{6403}{<< /Type /Annot /Subtype /Link /Border [0 0 0] /H /I /C [0 1 0] /Rect [122.695 601.341 147.755 610.208] /A << /S /GoTo /D <636974652E53636861> >> >>}
\MusicPDFWriteObject{6404}{<< /Type /Annot /Subtype /Link /Border [0 0 0] /H /I /C [0 1 0] /Rect [323.611 502.619 348.671 511.487] /A << /S /GoTo /D <636974652E53636861> >> >>}
\MusicPDFWriteObject{6405}{<< /Type /Annot /Subtype /Link /Border [0 0 0] /H /I /C [0 1 0] /Rect [438.724 502.619 470.348 511.487] /A << /S /GoTo /D <636974652E47726F2D62> >> >>}
\MusicPDFWriteObject{6406}{<< /Type /Annot /Subtype /Link /Border [0 0 0] /H /I /C [0 1 0] /Rect [439.221 408.033 470.845 416.9] /A << /S /GoTo /D <636974652E47726F2D62> >> >>}
\MusicPDFWriteObject{6407}{<< /Type /Annot /Subtype /Link /Border [0 0 0] /H /I /C [1 0 0] /Rect [138.125 392.392 148.058 405.249] /A << /S /GoTo /D <48666F6F746E6F74652E313231> >> >>}
\MusicPDFWriteObject{6408}{<< /Type /Annot /Subtype /Link /Border [0 0 0] /H /I /C [1 0 0] /Rect [229.114 241.185 265.813 251.976] /A << /S /GoTo /D <74686D2E432E322E332E3131> >> >>}
\MusicPDFWriteObject{6409}{<< /Type /Annot /Subtype /Link /Border [0 0 0] /H /I /C [1 0 0] /Rect [374.589 578.575 384.522 590.879] /A << /S /GoTo /D <48666F6F746E6F74652E313232> >> >>}
\MusicPDFWriteObject{6410}{<< /Type /Annot /Subtype /Link /Border [0 0 0] /H /I /C [1 0 0] /Rect [305.558 447.179 342.964 459.134] /A << /S /GoTo /D <6571756174696F6E2E432E322E332E3235> >> >>}
\MusicPDFWriteObject{6411}{<< /Type /Annot /Subtype /Link /Border [0 0 0] /H /I /C [1 0 0] /Rect [197.69 420.878 235.096 432.831] /A << /S /GoTo /D <6571756174696F6E2E432E322E332E3237> >> >>}
\MusicPDFWriteObject{6412}{<< /Type /Annot /Subtype /Link /Border [0 0 0] /H /I /C [1 0 0] /Rect [229.126 373.878 253.806 384.669] /A << /S /GoTo /D <73756273656374696F6E2E432E322E34> >> >>}
\MusicPDFWriteObject{6413}{<< /Type /Annot /Subtype /Link /Border [0 0 0] /H /I /C [0 1 0] /Rect [372.188 353.218 403.812 362.086] /A << /S /GoTo /D <636974652E47726F2D62> >> >>}
\MusicPDFWriteObject{6414}{<< /Type /Annot /Subtype /Link /Border [0 0 0] /H /I /C [1 0 0] /Rect [418.535 337.577 455.517 349.53] /A << /S /GoTo /D <6571756174696F6E2E432E322E332E3237> >> >>}
\MusicPDFWriteObject{6415}{<< /Type /Annot /Subtype /Link /Border [0 0 0] /H /I /C [0 1 0] /Rect [417.552 313.766 442.612 322.633] /A << /S /GoTo /D <636974652E53636861> >> >>}
\MusicPDFWriteObject{6416}{<< /Type /Annot /Subtype /Link /Border [0 0 0] /H /I /C [1 0 0] /Rect [71.004 284.614 109.083 298.174] /A << /S /GoTo /D <74686D2E432E322E332E3132> >> >>}
\MusicPDFWriteObject{6417}{<< /Type /Annot /Subtype /Link /Border [0 0 0] /H /I /C [0 1 0] /Rect [73.754 274.313 92.727 283.181] /A << /S /GoTo /D <636974652E547265> >> >>}
\MusicPDFWriteObject{6418}{<< /Type /Annot /Subtype /Link /Border [0 0 0] /H /I /C [1 0 0] /Rect [71.004 258.673 103.43 270.625] /A << /S /GoTo /D <74686D2E432E322E332E35> >> >>}
\MusicPDFWriteObject{6419}{<< /Type /Annot /Subtype /Link /Border [0 0 0] /H /I /C [0 1 0] /Rect [356.425 248.012 388.049 256.88] /A << /S /GoTo /D <636974652E47726F2D62> >> >>}
\MusicPDFWriteObject{6420}{<< /Type /Annot /Subtype /Link /Border [0 0 0] /H /I /C [1 0 0] /Rect [491.747 245.522 501.68 258.379] /A << /S /GoTo /D <48666F6F746E6F74652E313233> >> >>}
\MusicPDFWriteObject{6421}{<< /Type /Annot /Subtype /Link /Border [0 0 0] /H /I /C [1 0 0] /Rect [498.658 183.473 530.689 195.425] /A << /S /GoTo /D <74686D2E432E322E332E35> >> >>}
\MusicPDFWriteObject{6422}{<< /Type /Annot /Subtype /Link /Border [0 0 0] /H /I /C [0 1 0] /Rect [231.367 159.661 250.34 168.529] /A << /S /GoTo /D <636974652E547265> >> >>}
\MusicPDFWriteObject{6423}{<< /Type /Annot /Subtype /Link /Border [0 0 0] /H /I /C [1 0 0] /Rect [518.687 157.171 543.231 170.371] /A << /S /GoTo /D <73756273656374696F6E2E432E322E34> >> >>}
\MusicPDFWriteObject{6424}{<< /Type /Annot /Subtype /Link /Border [0 0 0] /H /I /C [0 1 0] /Rect [218.359 133.36 243.42 142.228] /A << /S /GoTo /D <636974652E53636861> >> >>}
\MusicPDFWriteObject{6425}{<< /Type /Annot /Subtype /Link /Border [0 0 0] /H /I /C [0 1 0] /Rect [470.012 95.049 497.262 102.541] /A << /S /GoTo /D <636974652E47726F2D62> >> >>}
\MusicPDFWriteObject{6426}{<< /Type /Annot /Subtype /Link /Border [0 0 0] /H /I /C [1 0 0] /Rect [376.268 555.149 414.135 567.102] /A << /S /GoTo /D <74686D2E432E322E332E3133> >> >>}
\MusicPDFWriteObject{6427}{<< /Type /Annot /Subtype /Link /Border [0 0 0] /H /I /C [1 0 0] /Rect [408.258 516.25 440.411 527.041] /A << /S /GoTo /D <74686D2E432E312E352E37> >> >>}
\MusicPDFWriteObject{6428}{<< /Type /Annot /Subtype /Link /Border [0 0 0] /H /I /C [1 0 0] /Rect [277.623 452.546 314.427 464.833] /A << /S /GoTo /D <74686D2E432E322E332E3132> >> >>}
\MusicPDFWriteObject{6429}{<< /Type /Annot /Subtype /Link /Border [0 0 0] /H /I /C [0 1 0] /Rect [384.706 455.036 403.679 463.904] /A << /S /GoTo /D <636974652E547265> >> >>}
\MusicPDFWriteObject{6430}{<< /Type /Annot /Subtype /Link /Border [0 0 0] /H /I /C [0 1 0] /Rect [465.534 455.036 490.594 463.904] /A << /S /GoTo /D <636974652E53636861> >> >>}
\MusicPDFWriteObject{6431}{<< /Type /Annot /Subtype /Link /Border [0 0 0] /H /I /C [1 0 0] /Rect [373.312 439.395 411.073 451.35] /A << /S /GoTo /D <74686D2E432E322E332E3133> >> >>}
\MusicPDFWriteObject{6432}{<< /Type /Annot /Subtype /Link /Border [0 0 0] /H /I /C [1 0 0] /Rect [224.236 363.516 262.068 375.468] /A << /S /GoTo /D <74686D2E432E322E332E3132> >> >>}
\MusicPDFWriteObject{6433}{<< /Type /Annot /Subtype /Link /Border [0 0 0] /H /I /C [1 0 0] /Rect [106.777 350.365 139.416 362.318] /A << /S /GoTo /D <74686D2E432E322E332E35> >> >>}
\MusicPDFWriteObject{6434}{<< /Type /Annot /Subtype /Link /Border [0 0 0] /H /I /C [0 1 0] /Rect [102.903 339.704 127.963 348.572] /A << /S /GoTo /D <636974652E53636861> >> >>}
\MusicPDFWriteObject{6435}{<< /Type /Annot /Subtype /Link /Border [0 0 0] /H /I /C [1 0 0] /Rect [237.348 243.539 269.166 256.739] /A << /S /GoTo /D <74686D2E432E322E332E35> >> >>}
\MusicPDFWriteObject{6436}{<< /Type /Annot /Subtype /Link /Border [0 0 0] /H /I /C [1 0 0] /Rect [443.394 217.237 476.429 229.192] /A << /S /GoTo /D <74686D2E432E322E332E37> >> >>}
\MusicPDFWriteObject{6437}{<< /Type /Annot /Subtype /Link /Border [0 0 0] /H /I /C [0 1 0] /Rect [240.471 206.576 265.531 215.444] /A << /S /GoTo /D <636974652E53636861> >> >>}
\MusicPDFWriteObject{6438}{<< /Type /Annot /Subtype /Link /Border [0 0 0] /H /I /C [1 0 0] /Rect [338.577 684.633 371.125 696.586] /A << /S /GoTo /D <6571756174696F6E2E432E312E352E31> >> >>}
\MusicPDFWriteObject{6439}{<< /Type /Annot /Subtype /Link /Border [0 0 0] /H /I /C [1 0 0] /Rect [446.603 605.729 484.541 617.681] /A << /S /GoTo /D <74686D2E432E312E352E3133> >> >>}
\MusicPDFWriteObject{6440}{<< /Type /Annot /Subtype /Link /Border [0 0 0] /H /I /C [1 0 0] /Rect [203.829 564.115 241.944 579.476] /A << /S /GoTo /D <74686D2E432E312E352E3133> >> >>}
\MusicPDFWriteObject{6441}{<< /Type /Annot /Subtype /Link /Border [0 0 0] /H /I /C [1 0 0] /Rect [435.854 552.825 473.261 564.778] /A << /S /GoTo /D <6571756174696F6E2E432E322E332E3238> >> >>}
\MusicPDFWriteObject{6442}{<< /Type /Annot /Subtype /Link /Border [0 0 0] /H /I /C [0 1 0] /Rect [411.643 459.216 430.616 468.083] /A << /S /GoTo /D <636974652E547265> >> >>}
\MusicPDFWriteObject{6443}{<< /Type /Annot /Subtype /Link /Border [0 0 0] /H /I /C [1 0 0] /Rect [164.133 427.95 196.559 439.903] /A << /S /GoTo /D <6571756174696F6E2E432E322E342E31> >> >>}
\MusicPDFWriteObject{6444}{<< /Type /Annot /Subtype /Link /Border [0 0 0] /H /I /C [0 1 0] /Rect [187.593 375.363 219.218 384.231] /A << /S /GoTo /D <636974652E47726F2D62> >> >>}
\MusicPDFWriteObject{6445}{<< /Type /Annot /Subtype /Link /Border [0 0 0] /H /I /C [1 0 0] /Rect [73.712 280.818 106.138 292.77] /A << /S /GoTo /D <6571756174696F6E2E432E322E342E31> >> >>}
\MusicPDFWriteObject{6446}{<< /Type /Annot /Subtype /Link /Border [0 0 0] /H /I /C [0 1 0] /Rect [334.733 257.006 366.358 265.874] /A << /S /GoTo /D <636974652E47726F2D62> >> >>}
\MusicPDFWriteObject{6447}{<< /Type /Annot /Subtype /Link /Border [0 0 0] /H /I /C [0 1 0] /Rect [157.996 163.99 183.056 172.858] /A << /S /GoTo /D <636974652E53636861> >> >>}
\MusicPDFWriteObject{6448}{<< /Type /Annot /Subtype /Link /Border [0 0 0] /H /I /C [0 1 0] /Rect [518.14 163.99 543.2 172.858] /A << /S /GoTo /D <636974652E53636861> >> >>}
\MusicPDFWriteObject{6449}{<< /Type /Annot /Subtype /Link /Border [0 0 0] /H /I /C [1 0 0] /Rect [314.725 148.35 347.151 160.302] /A << /S /GoTo /D <74686D2E432E322E342E31> >> >>}
\MusicPDFWriteObject{6450}{<< /Type /Annot /Subtype /Link /Border [0 0 0] /H /I /C [0 1 0] /Rect [388.873 150.84 407.846 159.708] /A << /S /GoTo /D <636974652E547265> >> >>}
\MusicPDFWriteObject{6451}{<< /Type /Annot /Subtype /Link /Border [0 0 0] /H /I /C [1 0 0] /Rect [197.557 135.752 230.592 146.543] /A << /S /GoTo /D <74686D2E432E322E342E31> >> >>}
\MusicPDFWriteObject{6452}{<< /Type /Annot /Subtype /Link /Border [0 0 0] /H /I /C [1 0 0] /Rect [283.385 109.45 300.616 120.298] /A << /S /GoTo /D <73656374696F6E2E432E33> >> >>}
\MusicPDFWriteObject{6453}{<< /Type /Annot /Subtype /Link /Border [0 0 0] /H /I /C [0 1 0] /Rect [255.608 98.237 280.668 107.105] /A << /S /GoTo /D <636974652E53636861> >> >>}
\MusicPDFWriteObject{6454}{<< /Type /Annot /Subtype /Link /Border [0 0 0] /H /I /C [0 1 0] /Rect [199.324 602.264 224.384 611.132] /A << /S /GoTo /D <636974652E53636861> >> >>}
\MusicPDFWriteObject{6455}{<< /Type /Annot /Subtype /Link /Border [0 0 0] /H /I /C [0 1 0] /Rect [339.427 602.264 358.401 611.132] /A << /S /GoTo /D <636974652E547265> >> >>}
\MusicPDFWriteObject{6456}{<< /Type /Annot /Subtype /Link /Border [0 0 0] /H /I /C [1 0 0] /Rect [279.854 510.644 312.281 522.596] /A << /S /GoTo /D <6571756174696F6E2E432E312E322E31> >> >>}
\MusicPDFWriteObject{6457}{<< /Type /Annot /Subtype /Link /Border [0 0 0] /H /I /C [0 1 0] /Rect [180.76 400.512 199.733 409.38] /A << /S /GoTo /D <636974652E547265> >> >>}
\MusicPDFWriteObject{6458}{<< /Type /Annot /Subtype /Link /Border [0 0 0] /H /I /C [0 1 0] /Rect [85.305 361.059 104.278 369.927] /A << /S /GoTo /D <636974652E547265> >> >>}
\MusicPDFWriteObject{6459}{<< /Type /Annot /Subtype /Link /Border [0 0 0] /H /I /C [0 1 0] /Rect [524.186 347.909 543.159 356.776] /A << /S /GoTo /D <636974652E547265> >> >>}
\MusicPDFWriteObject{6460}{<< /Type /Annot /Subtype /Link /Border [0 0 0] /H /I /C [0 1 0] /Rect [460.934 282.155 479.133 291.023] /A << /S /GoTo /D <636974652E4A6172> >> >>}
\MusicPDFWriteObject{6461}{<< /Type /Annot /Subtype /Link /Border [0 0 0] /H /I /C [0 1 0] /Rect [443.737 255.853 462.71 264.721] /A << /S /GoTo /D <636974652E547265> >> >>}
\MusicPDFWriteObject{6462}{<< /Type /Annot /Subtype /Link /Border [0 0 0] /H /I /C [0 1 0] /Rect [491.451 242.703 510.425 251.57] /A << /S /GoTo /D <636974652E547265> >> >>}
\MusicPDFWriteObject{6463}{<< /Type /Annot /Subtype /Link /Border [0 0 0] /H /I /C [1 0 0] /Rect [417.375 199.734 450.41 212.715] /A << /S /GoTo /D <6571756174696F6E2E432E312E322E31> >> >>}
\MusicPDFWriteObject{6464}{<< /Type /Annot /Subtype /Link /Border [0 0 0] /H /I /C [0 1 0] /Rect [222.786 175.823 247.846 184.691] /A << /S /GoTo /D <636974652E53636861> >> >>}
\MusicPDFWriteObject{6465}{<< /Type /Annot /Subtype /Link /Border [0 0 0] /H /I /C [0 1 0] /Rect [370.454 155.208 395.514 164.075] /A << /S /GoTo /D <636974652E53636861> >> >>}
\MusicPDFWriteObject{6466}{<< /Type /Annot /Subtype /Link /Border [0 0 0] /H /I /C [0 1 0] /Rect [524.971 155.208 543.17 164.075] /A << /S /GoTo /D <636974652E4A6172> >> >>}
\MusicPDFWriteObject{6467}{<< /Type /Annot /Subtype /Link /Border [0 0 0] /H /I /C [1 0 0] /Rect [360.117 506.444 392.787 518.389] /A << /S /GoTo /D <74686D2E432E322E312E37> >> >>}
\MusicPDFWriteObject{6468}{<< /Type /Annot /Subtype /Link /Border [0 0 0] /H /I /C [0 1 0] /Rect [99.972 456.304 125.032 465.171] /A << /S /GoTo /D <636974652E53636861> >> >>}
\MusicPDFWriteObject{6469}{<< /Type /Annot /Subtype /Link /Border [0 0 0] /H /I /C [0 1 0] /Rect [233.543 429.953 258.603 438.821] /A << /S /GoTo /D <636974652E53636861> >> >>}
\MusicPDFWriteObject{6470}{<< /Type /Annot /Subtype /Link /Border [0 0 0] /H /I /C [0 1 0] /Rect [197.332 358.564 227.843 367.432] /A << /S /GoTo /D <636974652E47726F2D65> >> >>}
\MusicPDFWriteObject{6471}{<< /Type /Annot /Subtype /Link /Border [0 0 0] /H /I /C [0 1 0] /Rect [337.184 345.414 362.244 354.281] /A << /S /GoTo /D <636974652E53636861> >> >>}
\MusicPDFWriteObject{6472}{<< /Type /Annot /Subtype /Link /Border [0 0 0] /H /I /C [1 0 0] /Rect [305.033 316.649 314.966 329.479] /A << /S /GoTo /D <48666F6F746E6F74652E313234> >> >>}
\MusicPDFWriteObject{6473}{<< /Type /Annot /Subtype /Link /Border [0 0 0] /H /I /C [0 1 0] /Rect [410.27 300.376 429.243 309.243] /A << /S /GoTo /D <636974652E547265> >> >>}
\MusicPDFWriteObject{6474}{<< /Type /Annot /Subtype /Link /Border [0 0 0] /H /I /C [1 0 0] /Rect [350.994 149.939 368.061 160.788] /A << /S /GoTo /D <73656374696F6E2E432E34> >> >>}
\MusicPDFWriteObject{6475}{<< /Type /Annot /Subtype /Link /Border [0 0 0] /H /I /C [1 0 0] /Rect [512.145 149.939 544.845 160.788] /A << /S /GoTo /D <74686D2E432E322E342E36> >> >>}
\MusicPDFWriteObject{6476}{<< /Type /Annot /Subtype /Link /Border [0 0 0] /H /I /C [1 0 0] /Rect [334.01 136.789 350.942 147.637] /A << /S /GoTo /D <73656374696F6E2E432E34> >> >>}
\MusicPDFWriteObject{6477}{<< /Type /Annot /Subtype /Link /Border [0 0 0] /H /I /C [0 1 0] /Rect [518.124 106.447 543.184 115.315] /A << /S /GoTo /D <636974652E53636861> >> >>}
\MusicPDFWriteObject{6478}{<< /Type /Annot /Subtype /Link /Border [0 0 0] /H /I /C [0 1 0] /Rect [253.749 85.086 270.222 92.579] /A << /S /GoTo /D <636974652E547265> >> >>}
\MusicPDFWriteObject{6479}{<< /Type /Annot /Subtype /Link /Border [0 0 0] /H /I /C [1 0 0] /Rect [308.353 671.482 341.389 683.455] /A << /S /GoTo /D <74686D2E432E342E332E32> >> >>}
\MusicPDFWriteObject{6480}{<< /Type /Annot /Subtype /Link /Border [0 0 0] /H /I /C [0 1 0] /Rect [306.581 617.424 325.554 626.292] /A << /S /GoTo /D <636974652E547265> >> >>}
\MusicPDFWriteObject{6481}{<< /Type /Annot /Subtype /Link /Border [0 0 0] /H /I /C [1 0 0] /Rect [251.356 601.783 283.357 615.177] /A << /S /GoTo /D <74686D2E432E322E342E36> >> >>}
\MusicPDFWriteObject{6482}{<< /Type /Annot /Subtype /Link /Border [0 0 0] /H /I /C [1 0 0] /Rect [363.109 589.186 395.079 599.977] /A << /S /GoTo /D <74686D2E432E322E342E36> >> >>}
\MusicPDFWriteObject{6483}{<< /Type /Annot /Subtype /Link /Border [0 0 0] /H /I /C [0 1 0] /Rect [280.279 577.972 299.252 586.84] /A << /S /GoTo /D <636974652E547265> >> >>}
\MusicPDFWriteObject{6484}{<< /Type /Annot /Subtype /Link /Border [0 0 0] /H /I /C [1 0 0] /Rect [458.517 509.678 490.487 521.651] /A << /S /GoTo /D <74686D2E432E322E342E34> >> >>}
\MusicPDFWriteObject{6485}{<< /Type /Annot /Subtype /Link /Border [0 0 0] /H /I /C [0 1 0] /Rect [328.064 433.264 347.038 442.132] /A << /S /GoTo /D <636974652E547265> >> >>}
\MusicPDFWriteObject{6486}{<< /Type /Annot /Subtype /Link /Border [0 0 0] /H /I /C [1 0 0] /Rect [512.784 430.774 544.845 442.727] /A << /S /GoTo /D <74686D2E432E322E342E36> >> >>}
\MusicPDFWriteObject{6487}{<< /Type /Annot /Subtype /Link /Border [0 0 0] /H /I /C [0 1 0] /Rect [221.928 420.113 240.901 428.981] /A << /S /GoTo /D <636974652E547265> >> >>}
\MusicPDFWriteObject{6488}{<< /Type /Annot /Subtype /Link /Border [0 0 0] /H /I /C [0 1 0] /Rect [189.904 403.017 208.877 411.885] /A << /S /GoTo /D <636974652E547265> >> >>}
\MusicPDFWriteObject{6489}{<< /Type /Annot /Subtype /Link /Border [0 0 0] /H /I /C [1 0 0] /Rect [132.619 387.376 165.045 399.329] /A << /S /GoTo /D <74686D2E432E322E342E36> >> >>}
\MusicPDFWriteObject{6490}{<< /Type /Annot /Subtype /Link /Border [0 0 0] /H /I /C [0 1 0] /Rect [267.967 389.866 286.94 398.734] /A << /S /GoTo /D <636974652E547265> >> >>}
\MusicPDFWriteObject{6491}{<< /Type /Annot /Subtype /Link /Border [0 0 0] /H /I /C [1 0 0] /Rect [208.678 252.169 241.104 265.027] /A << /S /GoTo /D <74686D2E432E322E342E32> >> >>}
\MusicPDFWriteObject{6492}{<< /Type /Annot /Subtype /Link /Border [0 0 0] /H /I /C [0 1 0] /Rect [317.549 254.659 336.522 263.527] /A << /S /GoTo /D <636974652E547265> >> >>}
\MusicPDFWriteObject{6493}{<< /Type /Annot /Subtype /Link /Border [0 0 0] /H /I /C [0 1 0] /Rect [420.492 254.659 439.465 263.527] /A << /S /GoTo /D <636974652E547265> >> >>}
\MusicPDFWriteObject{6494}{<< /Type /Annot /Subtype /Link /Border [0 0 0] /H /I /C [1 0 0] /Rect [490.179 252.169 500.112 265.027] /A << /S /GoTo /D <48666F6F746E6F74652E313235> >> >>}
\MusicPDFWriteObject{6495}{<< /Type /Annot /Subtype /Link /Border [0 0 0] /H /I /C [1 0 0] /Rect [288.955 150.636 321.746 161.426] /A << /S /GoTo /D <74686D2E432E322E342E39> >> >>}
\MusicPDFWriteObject{6496}{<< /Type /Annot /Subtype /Link /Border [0 0 0] /H /I /C [0 1 0] /Rect [319.598 139.422 338.571 148.289] /A << /S /GoTo /D <636974652E547265> >> >>}
\MusicPDFWriteObject{6497}{<< /Type /Annot /Subtype /Link /Border [0 0 0] /H /I /C [0 1 0] /Rect [362.838 126.271 381.811 135.139] /A << /S /GoTo /D <636974652E547265> >> >>}
\MusicPDFWriteObject{6498}{<< /Type /Annot /Subtype /Link /Border [0 0 0] /H /I /C [1 0 0] /Rect [257.329 92.7 284.931 103.845] /A << /S /GoTo /D <73756273756273656374696F6E2E432E322E342E33> >> >>}
\MusicPDFWriteObject{6499}{<< /Type /Annot /Subtype /Link /Border [0 0 0] /H /I /C [0 1 0] /Rect [97.108 84.287 113.581 91.779] /A << /S /GoTo /D <636974652E547265> >> >>}
\MusicPDFWriteObject{6500}{<< /Type /Annot /Subtype /Link /Border [0 0 0] /H /I /C [0 1 0] /Rect [73.77 673.973 98.831 682.84] /A << /S /GoTo /D <636974652E53636861> >> >>}
\MusicPDFWriteObject{6501}{<< /Type /Annot /Subtype /Link /Border [0 0 0] /H /I /C [0 1 0] /Rect [337.859 567.163 356.832 576.031] /A << /S /GoTo /D <636974652E547265> >> >>}
\MusicPDFWriteObject{6502}{<< /Type /Annot /Subtype /Link /Border [0 0 0] /H /I /C [1 0 0] /Rect [409.787 564.673 419.72 577.531] /A << /S /GoTo /D <48666F6F746E6F74652E313236> >> >>}
\MusicPDFWriteObject{6503}{<< /Type /Annot /Subtype /Link /Border [0 0 0] /H /I /C [1 0 0] /Rect [412.726 512.559 450.133 524.512] /A << /S /GoTo /D <74686D2E432E322E342E3130> >> >>}
\MusicPDFWriteObject{6504}{<< /Type /Annot /Subtype /Link /Border [0 0 0] /H /I /C [1 0 0] /Rect [382.168 405.75 392.101 418.607] /A << /S /GoTo /D <48666F6F746E6F74652E313237> >> >>}
\MusicPDFWriteObject{6505}{<< /Type /Annot /Subtype /Link /Border [0 0 0] /H /I /C [0 1 0] /Rect [407.658 408.24 439.282 417.108] /A << /S /GoTo /D <636974652E47726F2D62> >> >>}
\MusicPDFWriteObject{6506}{<< /Type /Annot /Subtype /Link /Border [0 0 0] /H /I /C [0 1 0] /Rect [310.802 336.155 335.862 345.023] /A << /S /GoTo /D <636974652E53636861> >> >>}
\MusicPDFWriteObject{6507}{<< /Type /Annot /Subtype /Link /Border [0 0 0] /H /I /C [0 1 0] /Rect [275.548 309.854 300.608 318.721] /A << /S /GoTo /D <636974652E53636861> >> >>}
\MusicPDFWriteObject{6508}{<< /Type /Annot /Subtype /Link /Border [0 0 0] /H /I /C [0 1 0] /Rect [191.907 204.366 216.967 213.234] /A << /S /GoTo /D <636974652E53636861> >> >>}
\MusicPDFWriteObject{6509}{<< /Type /Annot /Subtype /Link /Border [0 0 0] /H /I /C [1 0 0] /Rect [494.117 201.516 526.514 213.829] /A << /S /GoTo /D <74686D2E432E322E312E35> >> >>}
\MusicPDFWriteObject{6510}{<< /Type /Annot /Subtype /Link /Border [0 0 0] /H /I /C [1 0 0] /Rect [465.664 154.001 498.09 165.953] /A << /S /GoTo /D <6571756174696F6E2E432E322E342E36> >> >>}
\MusicPDFWriteObject{6511}{<< /Type /Annot /Subtype /Link /Border [0 0 0] /H /I /C [1 0 0] /Rect [306.196 140.85 338.622 152.803] /A << /S /GoTo /D <6571756174696F6E2E432E322E342E36> >> >>}
\MusicPDFWriteObject{6512}{<< /Type /Annot /Subtype /Link /Border [0 0 0] /H /I /C [1 0 0] /Rect [126.269 114.548 158.696 126.501] /A << /S /GoTo /D <6571756174696F6E2E432E322E342E36> >> >>}
\MusicPDFWriteObject{6513}{<< /Type /Annot /Subtype /Link /Border [0 0 0] /H /I /C [1 0 0] /Rect [126.269 671.482 158.696 683.435] /A << /S /GoTo /D <6571756174696F6E2E432E322E342E36> >> >>}
\MusicPDFWriteObject{6514}{<< /Type /Annot /Subtype /Link /Border [0 0 0] /H /I /C [1 0 0] /Rect [180.11 640.741 211.927 651.531] /A << /S /GoTo /D <74686D2E432E322E342E31> >> >>}
\MusicPDFWriteObject{6515}{<< /Type /Annot /Subtype /Link /Border [0 0 0] /H /I /C [0 1 0] /Rect [288.745 616.376 307.718 625.244] /A << /S /GoTo /D <636974652E547265> >> >>}
\MusicPDFWriteObject{6516}{<< /Type /Annot /Subtype /Link /Border [0 0 0] /H /I /C [0 1 0] /Rect [127.974 590.074 153.034 598.942] /A << /S /GoTo /D <636974652E53636861> >> >>}
\MusicPDFWriteObject{6517}{<< /Type /Annot /Subtype /Link /Border [0 0 0] /H /I /C [1 0 0] /Rect [175.456 574.434 212.154 586.386] /A << /S /GoTo /D <74686D2E432E322E342E3135> >> >>}
\MusicPDFWriteObject{6518}{<< /Type /Annot /Subtype /Link /Border [0 0 0] /H /I /C [1 0 0] /Rect [436.191 574.434 452.825 586.386] /A << /S /GoTo /D <73656374696F6E2E432E33> >> >>}
\MusicPDFWriteObject{6519}{<< /Type /Annot /Subtype /Link /Border [0 0 0] /H /I /C [1 0 0] /Rect [256.476 488.443 288.902 500.396] /A << /S /GoTo /D <6571756174696F6E2E432E322E342E36> >> >>}
\MusicPDFWriteObject{6520}{<< /Type /Annot /Subtype /Link /Border [0 0 0] /H /I /C [1 0 0] /Rect [91.42 475.292 128.827 487.247] /A << /S /GoTo /D <74686D2E432E322E342E3135> >> >>}
\MusicPDFWriteObject{6521}{<< /Type /Annot /Subtype /Link /Border [0 0 0] /H /I /C [0 1 0] /Rect [179.595 459.638 204.655 468.506] /A << /S /GoTo /D <636974652E53636861> >> >>}
\MusicPDFWriteObject{6522}{<< /Type /Annot /Subtype /Link /Border [0 0 0] /H /I /C [0 1 0] /Rect [372.375 459.638 397.435 468.506] /A << /S /GoTo /D <636974652E53636861> >> >>}
\MusicPDFWriteObject{6523}{<< /Type /Annot /Subtype /Link /Border [0 0 0] /H /I /C [1 0 0] /Rect [294.433 439.557 331.131 450.348] /A << /S /GoTo /D <74686D2E432E322E342E3136> >> >>}
\MusicPDFWriteObject{6524}{<< /Type /Annot /Subtype /Link /Border [0 0 0] /H /I /C [0 1 0] /Rect [342.601 375.74 361.574 384.608] /A << /S /GoTo /D <636974652E547265> >> >>}
\MusicPDFWriteObject{6525}{<< /Type /Annot /Subtype /Link /Border [0 0 0] /H /I /C [0 1 0] /Rect [483.306 331.294 508.366 340.162] /A << /S /GoTo /D <636974652E53636861> >> >>}
\MusicPDFWriteObject{6526}{<< /Type /Annot /Subtype /Link /Border [0 0 0] /H /I /C [0 1 0] /Rect [188.38 182.204 213.44 191.072] /A << /S /GoTo /D <636974652E53636861> >> >>}
\MusicPDFWriteObject{6527}{<< /Type /Annot /Subtype /Link /Border [0 0 0] /H /I /C [0 1 0] /Rect [109.625 142.752 128.598 151.62] /A << /S /GoTo /D <636974652E547265> >> >>}
\MusicPDFWriteObject{6528}{<< /Type /Annot /Subtype /Link /Border [0 0 0] /H /I /C [1 0 0] /Rect [192.899 684.335 230.448 696.606] /A << /S /GoTo /D <74686D2E432E322E342E3138> >> >>}
\MusicPDFWriteObject{6529}{<< /Type /Annot /Subtype /Link /Border [0 0 0] /H /I /C [1 0 0] /Rect [440.4 624.435 477.099 636.388] /A << /S /GoTo /D <74686D2E432E322E342E3138> >> >>}
\MusicPDFWriteObject{6530}{<< /Type /Annot /Subtype /Link /Border [0 0 0] /H /I /C [0 1 0] /Rect [123.383 613.774 142.356 622.642] /A << /S /GoTo /D <636974652E547265> >> >>}
\MusicPDFWriteObject{6531}{<< /Type /Annot /Subtype /Link /Border [0 0 0] /H /I /C [1 0 0] /Rect [205.439 600.624 215.372 610.991] /A << /S /GoTo /D <48666F6F746E6F74652E313238> >> >>}
\MusicPDFWriteObject{6532}{<< /Type /Annot /Subtype /Link /Border [0 0 0] /H /I /C [1 0 0] /Rect [235.266 486.485 244.905 498.441] /A << /S /GoTo /D <617070656E6469782E44> >> >>}
\MusicPDFWriteObject{6533}{<< /Type /Annot /Subtype /Link /Border [0 0 0] /H /I /C [1 0 0] /Rect [71.004 473.889 103.845 484.679] /A << /S /GoTo /D <74686D2E442E312E382E31> >> >>}
\MusicPDFWriteObject{6534}{<< /Type /Annot /Subtype /Link /Border [0 0 0] /H /I /C [1 0 0] /Rect [188.838 454.206 222.297 466.161] /A << /S /GoTo /D <74686D2E442E312E352E37> >> >>}
\MusicPDFWriteObject{6535}{<< /Type /Annot /Subtype /Link /Border [0 0 0] /H /I /C [0 1 0] /Rect [253.372 430.395 278.432 439.263] /A << /S /GoTo /D <636974652E53636861> >> >>}
\MusicPDFWriteObject{6536}{<< /Type /Annot /Subtype /Link /Border [0 0 0] /H /I /C [0 1 0] /Rect [461.855 430.395 486.915 439.263] /A << /S /GoTo /D <636974652E53636861> >> >>}
\MusicPDFWriteObject{6537}{<< /Type /Annot /Subtype /Link /Border [0 0 0] /H /I /C [0 1 0] /Rect [146.314 417.245 165.287 426.112] /A << /S /GoTo /D <636974652E547265> >> >>}
\MusicPDFWriteObject{6538}{<< /Type /Annot /Subtype /Link /Border [0 0 0] /H /I /C [0 1 0] /Rect [332.297 335.387 357.357 344.255] /A << /S /GoTo /D <636974652E53636861> >> >>}
\MusicPDFWriteObject{6539}{<< /Type /Annot /Subtype /Link /Border [0 0 0] /H /I /C [0 1 0] /Rect [132.177 95.049 148.65 102.541] /A << /S /GoTo /D <636974652E547265> >> >>}
\MusicPDFWriteObject{6540}{<< /Type /Annot /Subtype /Link /Border [0 0 0] /H /I /C [0 1 0] /Rect [440.816 520.681 465.876 529.549] /A << /S /GoTo /D <636974652E53636861> >> >>}
\MusicPDFWriteObject{6541}{<< /Type /Annot /Subtype /Link /Border [0 0 0] /H /I /C [0 1 0] /Rect [192.927 494.379 217.987 503.247] /A << /S /GoTo /D <636974652E53636861> >> >>}
\MusicPDFWriteObject{6542}{<< /Type /Annot /Subtype /Link /Border [0 0 0] /H /I /C [0 1 0] /Rect [329.511 441.776 354.571 450.644] /A << /S /GoTo /D <636974652E53636861> >> >>}
\MusicPDFWriteObject{6543}{<< /Type /Annot /Subtype /Link /Border [0 0 0] /H /I /C [0 1 0] /Rect [279.961 428.626 305.021 437.493] /A << /S /GoTo /D <636974652E53636861> >> >>}
\MusicPDFWriteObject{6544}{<< /Type /Annot /Subtype /Link /Border [0 0 0] /H /I /C [0 1 0] /Rect [304.697 415.475 323.67 424.343] /A << /S /GoTo /D <636974652E547265> >> >>}
\MusicPDFWriteObject{6545}{<< /Type /Annot /Subtype /Link /Border [0 0 0] /H /I /C [0 1 0] /Rect [223.856 402.324 248.916 411.192] /A << /S /GoTo /D <636974652E53636861> >> >>}
\MusicPDFWriteObject{6546}{<< /Type /Annot /Subtype /Link /Border [0 0 0] /H /I /C [0 1 0] /Rect [286.725 389.173 311.785 398.041] /A << /S /GoTo /D <636974652E53636861> >> >>}
\MusicPDFWriteObject{6547}{<< /Type /Annot /Subtype /Link /Border [0 0 0] /H /I /C [1 0 0] /Rect [106.922 347.231 139.683 359.184] /A << /S /GoTo /D <74686D2E432E312E362E31> >> >>}
\MusicPDFWriteObject{6548}{<< /Type /Annot /Subtype /Link /Border [0 0 0] /H /I /C [1 0 0] /Rect [71.004 200.953 103.704 212.905] /A << /S /GoTo /D <74686D2E432E322E352E31> >> >>}
\MusicPDFWriteObject{6549}{<< /Type /Annot /Subtype /Link /Border [0 0 0] /H /I /C [1 0 0] /Rect [315.582 187.802 348.617 199.755] /A << /S /GoTo /D <74686D2E432E322E342E32> >> >>}
\MusicPDFWriteObject{6550}{<< /Type /Annot /Subtype /Link /Border [0 0 0] /H /I /C [1 0 0] /Rect [109.475 135.199 147.59 147.152] /A << /S /GoTo /D <6571756174696F6E2E432E322E332E3236> >> >>}
\MusicPDFWriteObject{6551}{<< /Type /Annot /Subtype /Link /Border [0 0 0] /H /I /C [1 0 0] /Rect [290.662 108.897 322.723 120.85] /A << /S /GoTo /D <74686D2E432E322E332E39> >> >>}
\MusicPDFWriteObject{6552}{<< /Type /Annot /Subtype /Link /Border [0 0 0] /H /I /C [0 1 0] /Rect [238.891 615.392 262.733 624.26] /A << /S /GoTo /D <636974652E4D6531> >> >>}
\MusicPDFWriteObject{6553}{<< /Type /Annot /Subtype /Link /Border [0 0 0] /H /I /C [1 0 0] /Rect [508.776 344.477 541.811 356.43] /A << /S /GoTo /D <6571756174696F6E2E432E332E312E31> >> >>}
\MusicPDFWriteObject{6554}{<< /Type /Annot /Subtype /Link /Border [0 0 0] /H /I /C [1 0 0] /Rect [457.632 294.899 490.059 306.852] /A << /S /GoTo /D <6571756174696F6E2E432E332E312E32> >> >>}
\MusicPDFWriteObject{6555}{<< /Type /Annot /Subtype /Link /Border [0 0 0] /H /I /C [1 0 0] /Rect [198.688 281.748 231.115 293.701] /A << /S /GoTo /D <6571756174696F6E2E432E322E342E32> >> >>}
\MusicPDFWriteObject{6556}{<< /Type /Annot /Subtype /Link /Border [0 0 0] /H /I /C [1 0 0] /Rect [147.48 246.412 180.515 260.072] /A << /S /GoTo /D <6571756174696F6E2E432E332E312E31> >> >>}
\MusicPDFWriteObject{6557}{<< /Type /Annot /Subtype /Link /Border [0 0 0] /H /I /C [0 1 0] /Rect [438.678 237.439 457.651 246.306] /A << /S /GoTo /D <636974652E547265> >> >>}
\MusicPDFWriteObject{6558}{<< /Type /Annot /Subtype /Link /Border [0 0 0] /H /I /C [0 1 0] /Rect [273.218 224.288 292.191 233.156] /A << /S /GoTo /D <636974652E547265> >> >>}
\MusicPDFWriteObject{6559}{<< /Type /Annot /Subtype /Link /Border [0 0 0] /H /I /C [1 0 0] /Rect [455.195 220.11 488.017 233.77] /A << /S /GoTo /D <6571756174696F6E2E432E332E312E31> >> >>}
\MusicPDFWriteObject{6560}{<< /Type /Annot /Subtype /Link /Border [0 0 0] /H /I /C [1 0 0] /Rect [285.646 208.647 318.073 220.6] /A << /S /GoTo /D <6571756174696F6E2E432E332E312E32> >> >>}
\MusicPDFWriteObject{6561}{<< /Type /Annot /Subtype /Link /Border [0 0 0] /H /I /C [1 0 0] /Rect [234.465 148.35 266.892 160.302] /A << /S /GoTo /D <6571756174696F6E2E432E312E322E31> >> >>}
\MusicPDFWriteObject{6562}{<< /Type /Annot /Subtype /Link /Border [0 0 0] /H /I /C [1 0 0] /Rect [213.877 671.482 246.304 683.435] /A << /S /GoTo /D <6571756174696F6E2E432E322E342E32> >> >>}
\MusicPDFWriteObject{6563}{<< /Type /Annot /Subtype /Link /Border [0 0 0] /H /I /C [1 0 0] /Rect [419.816 645.181 452.791 657.134] /A << /S /GoTo /D <6571756174696F6E2E432E332E312E31> >> >>}
\MusicPDFWriteObject{6564}{<< /Type /Annot /Subtype /Link /Border [0 0 0] /H /I /C [1 0 0] /Rect [161.481 618.879 194.09 630.832] /A << /S /GoTo /D <6571756174696F6E2E432E322E342E32> >> >>}
\MusicPDFWriteObject{6565}{<< /Type /Annot /Subtype /Link /Border [0 0 0] /H /I /C [0 1 0] /Rect [73.77 608.219 98.831 617.087] /A << /S /GoTo /D <636974652E53636861> >> >>}
\MusicPDFWriteObject{6566}{<< /Type /Annot /Subtype /Link /Border [0 0 0] /H /I /C [1 0 0] /Rect [219.897 433.446 252.324 445.399] /A << /S /GoTo /D <6571756174696F6E2E432E332E312E33> >> >>}
\MusicPDFWriteObject{6567}{<< /Type /Annot /Subtype /Link /Border [0 0 0] /H /I /C [1 0 0] /Rect [234.976 323.717 272.383 335.672] /A << /S /GoTo /D <74686D2E432E322E342E3135> >> >>}
\MusicPDFWriteObject{6568}{<< /Type /Annot /Subtype /Link /Border [0 0 0] /H /I /C [1 0 0] /Rect [211.772 276.967 244.198 288.92] /A << /S /GoTo /D <6571756174696F6E2E432E332E312E33> >> >>}
\MusicPDFWriteObject{6569}{<< /Type /Annot /Subtype /Link /Border [0 0 0] /H /I /C [0 1 0] /Rect [162.336 167.222 187.016 176.09] /A << /S /GoTo /D <636974652E4C69742D49> >> >>}
\MusicPDFWriteObject{6570}{<< /Type /Annot /Subtype /Link /Border [0 0 0] /H /I /C [1 0 0] /Rect [86.568 151.581 118.994 163.534] /A << /S /GoTo /D <6571756174696F6E2E432E332E322E33> >> >>}
\MusicPDFWriteObject{6571}{<< /Type /Annot /Subtype /Link /Border [0 0 0] /H /I /C [1 0 0] /Rect [315.756 684.633 348.182 696.586] /A << /S /GoTo /D <6571756174696F6E2E432E332E322E32> >> >>}
\MusicPDFWriteObject{6572}{<< /Type /Annot /Subtype /Link /Border [0 0 0] /H /I /C [1 0 0] /Rect [245.415 531.617 277.842 543.569] /A << /S /GoTo /D <6571756174696F6E2E432E332E322E32> >> >>}
\MusicPDFWriteObject{6573}{<< /Type /Annot /Subtype /Link /Border [0 0 0] /H /I /C [0 1 0] /Rect [98.208 347.222 122.05 356.09] /A << /S /GoTo /D <636974652E4D6531> >> >>}
\MusicPDFWriteObject{6574}{<< /Type /Annot /Subtype /Link /Border [0 0 0] /H /I /C [0 1 0] /Rect [518.099 242.718 543.159 251.586] /A << /S /GoTo /D <636974652E53636861> >> >>}
\MusicPDFWriteObject{6575}{<< /Type /Annot /Subtype /Link /Border [0 0 0] /H /I /C [1 0 0] /Rect [214.902 563.848 247.329 575.801] /A << /S /GoTo /D <6571756174696F6E2E432E332E322E35> >> >>}
\MusicPDFWriteObject{6576}{<< /Type /Annot /Subtype /Link /Border [0 0 0] /H /I /C [1 0 0] /Rect [335.663 550.697 373.777 562.65] /A << /S /GoTo /D <74686D2E432E322E342E3135> >> >>}
\MusicPDFWriteObject{6577}{<< /Type /Annot /Subtype /Link /Border [0 0 0] /H /I /C [1 0 0] /Rect [71.004 511.245 103.43 523.198] /A << /S /GoTo /D <74686D2E432E322E312E35> >> >>}
\MusicPDFWriteObject{6578}{<< /Type /Annot /Subtype /Link /Border [0 0 0] /H /I /C [1 0 0] /Rect [364.394 458.642 396.82 470.595] /A << /S /GoTo /D <6571756174696F6E2E432E332E322E35> >> >>}
\MusicPDFWriteObject{6579}{<< /Type /Annot /Subtype /Link /Border [0 0 0] /H /I /C [0 1 0] /Rect [111.518 447.981 136.198 456.849] /A << /S /GoTo /D <636974652E4C69742D49> >> >>}
\MusicPDFWriteObject{6580}{<< /Type /Annot /Subtype /Link /Border [0 0 0] /H /I /C [0 1 0] /Rect [483.226 434.831 508.286 443.698] /A << /S /GoTo /D <636974652E53636861> >> >>}
\MusicPDFWriteObject{6581}{<< /Type /Annot /Subtype /Link /Border [0 0 0] /H /I /C [0 1 0] /Rect [186.827 327.131 211.887 335.999] /A << /S /GoTo /D <636974652E53636861> >> >>}
\MusicPDFWriteObject{6582}{<< /Type /Annot /Subtype /Link /Border [0 0 0] /H /I /C [0 1 0] /Rect [271.787 313.98 290.892 322.848] /A << /S /GoTo /D <636974652E4E42> >> >>}
\MusicPDFWriteObject{6583}{<< /Type /Annot /Subtype /Link /Border [0 0 0] /H /I /C [0 1 0] /Rect [470.318 313.98 489.423 322.848] /A << /S /GoTo /D <636974652E4E42> >> >>}
\MusicPDFWriteObject{6584}{<< /Type /Annot /Subtype /Link /Border [0 0 0] /H /I /C [0 1 0] /Rect [327.247 235.075 352.307 243.943] /A << /S /GoTo /D <636974652E53636861> >> >>}
\MusicPDFWriteObject{6585}{<< /Type /Annot /Subtype /Link /Border [0 0 0] /H /I /C [1 0 0] /Rect [226.828 216.972 259.437 231.407] /A << /S /GoTo /D <74686D2E432E322E312E37> >> >>}
\MusicPDFWriteObject{6586}{<< /Type /Annot /Subtype /Link /Border [0 0 0] /H /I /C [1 0 0] /Rect [218.628 206.016 251.45 217.968] /A << /S /GoTo /D <74686D2E432E342E312E32> >> >>}
\MusicPDFWriteObject{6587}{<< /Type /Annot /Subtype /Link /Border [0 0 0] /H /I /C [1 0 0] /Rect [71.004 166.563 104.039 178.516] /A << /S /GoTo /D <74686D2E432E312E332E31> >> >>}
\MusicPDFWriteObject{6588}{<< /Type /Annot /Subtype /Link /Border [0 0 0] /H /I /C [0 1 0] /Rect [186.522 561.626 220.491 570.494] /A << /S /GoTo /D <636974652E436E772D61> >> >>}
\MusicPDFWriteObject{6589}{<< /Type /Annot /Subtype /Link /Border [0 0 0] /H /I /C [0 1 0] /Rect [328.842 548.475 353.903 557.343] /A << /S /GoTo /D <636974652E53636861> >> >>}
\MusicPDFWriteObject{6590}{<< /Type /Annot /Subtype /Link /Border [0 0 0] /H /I /C [0 1 0] /Rect [175.633 482.721 200.693 491.589] /A << /S /GoTo /D <636974652E53636861> >> >>}
\MusicPDFWriteObject{6591}{<< /Type /Annot /Subtype /Link /Border [0 0 0] /H /I /C [1 0 0] /Rect [264.411 433.835 295.189 445.79] /A << /S /GoTo /D <74686D2E342E332E312E36> >> >>}
\MusicPDFWriteObject{6592}{<< /Type /Annot /Subtype /Link /Border [0 0 0] /H /I /C [0 1 0] /Rect [73.765 363.242 98.825 372.11] /A << /S /GoTo /D <636974652E53636861> >> >>}
\MusicPDFWriteObject{6593}{<< /Type /Annot /Subtype /Link /Border [0 0 0] /H /I /C [0 1 0] /Rect [92.581 350.091 116.278 358.959] /A << /S /GoTo /D <636974652E4B7562> >> >>}
\MusicPDFWriteObject{6594}{<< /Type /Annot /Subtype /Link /Border [0 0 0] /H /I /C [0 1 0] /Rect [187.283 189.002 220.775 197.87] /A << /S /GoTo /D <636974652E436E772D63> >> >>}
\MusicPDFWriteObject{6595}{<< /Type /Annot /Subtype /Link /Border [0 0 0] /H /I /C [1 0 0] /Rect [400.597 186.512 432.415 198.465] /A << /S /GoTo /D <74686D2E432E332E342E31> >> >>}
\MusicPDFWriteObject{6596}{<< /Type /Annot /Subtype /Link /Border [0 0 0] /H /I /C [0 1 0] /Rect [509.778 189.002 543.27 197.87] /A << /S /GoTo /D <636974652E436E772D63> >> >>}
\MusicPDFWriteObject{6597}{<< /Type /Annot /Subtype /Link /Border [0 0 0] /H /I /C [1 0 0] /Rect [123.932 685.187 155.841 695.977] /A << /S /GoTo /D <74686D2E432E332E342E31> >> >>}
\MusicPDFWriteObject{6598}{<< /Type /Annot /Subtype /Link /Border [0 0 0] /H /I /C [1 0 0] /Rect [359.266 685.187 391.175 695.977] /A << /S /GoTo /D <74686D2E432E332E342E32> >> >>}
\MusicPDFWriteObject{6599}{<< /Type /Annot /Subtype /Link /Border [0 0 0] /H /I /C [0 1 0] /Rect [147.167 610.192 181.136 619.06] /A << /S /GoTo /D <636974652E436E772D61> >> >>}
\MusicPDFWriteObject{6600}{<< /Type /Annot /Subtype /Link /Border [0 0 0] /H /I /C [0 1 0] /Rect [174.341 137.689 198.038 146.557] /A << /S /GoTo /D <636974652E4B7562> >> >>}
\MusicPDFWriteObject{6601}{<< /Type /Annot /Subtype /Link /Border [0 0 0] /H /I /C [0 1 0] /Rect [146.536 124.538 180.028 133.406] /A << /S /GoTo /D <636974652E436E772D63> >> >>}
\MusicPDFWriteObject{6602}{<< /Type /Annot /Subtype /Link /Border [0 0 0] /H /I /C [0 1 0] /Rect [146.536 111.387 180.028 120.255] /A << /S /GoTo /D <636974652E436E772D63> >> >>}
\MusicPDFWriteObject{6603}{<< /Type /Annot /Subtype /Link /Border [0 0 0] /H /I /C [0 1 0] /Rect [146.536 98.237 180.028 107.105] /A << /S /GoTo /D <636974652E436E772D63> >> >>}
\MusicPDFWriteObject{6604}{<< /Type /Annot /Subtype /Link /Border [0 0 0] /H /I /C [0 1 0] /Rect [146.536 85.086 180.028 93.954] /A << /S /GoTo /D <636974652E436E772D63> >> >>}
\MusicPDFWriteObject{6605}{<< /Type /Annot /Subtype /Link /Border [0 0 0] /H /I /C [0 1 0] /Rect [199.199 203.674 233.168 212.542] /A << /S /GoTo /D <636974652E436E772D61> >> >>}
\MusicPDFWriteObject{6606}{<< /Type /Annot /Subtype /Link /Border [0 0 0] /H /I /C [0 1 0] /Rect [116.727 177.372 150.696 186.24] /A << /S /GoTo /D <636974652E436E772D61> >> >>}
\MusicPDFWriteObject{6607}{<< /Type /Annot /Subtype /Link /Border [0 0 0] /H /I /C [1 0 0] /Rect [368.587 129.292 400.404 140.083] /A << /S /GoTo /D <74686D2E432E332E352E31> >> >>}
\MusicPDFWriteObject{6608}{<< /Type /Annot /Subtype /Link /Border [0 0 0] /H /I /C [0 1 0] /Rect [509.259 687.123 543.228 695.991] /A << /S /GoTo /D <636974652E436E772D61> >> >>}
\MusicPDFWriteObject{6609}{<< /Type /Annot /Subtype /Link /Border [0 0 0] /H /I /C [0 1 0] /Rect [430.396 518.828 464.365 527.696] /A << /S /GoTo /D <636974652E436E772D61> >> >>}
\MusicPDFWriteObject{6610}{<< /Type /Annot /Subtype /Link /Border [0 0 0] /H /I /C [1 0 0] /Rect [221.57 430.931 253.996 442.884] /A << /S /GoTo /D <6571756174696F6E2E432E332E352E35> >> >>}
\MusicPDFWriteObject{6611}{<< /Type /Annot /Subtype /Link /Border [0 0 0] /H /I /C [1 0 0] /Rect [167.59 308.286 200.017 320.239] /A << /S /GoTo /D <6571756174696F6E2E432E332E352E36> >> >>}
\MusicPDFWriteObject{6612}{<< /Type /Annot /Subtype /Link /Border [0 0 0] /H /I /C [0 1 0] /Rect [107.038 203.647 141.007 212.515] /A << /S /GoTo /D <636974652E436E772D61> >> >>}
\MusicPDFWriteObject{6613}{<< /Type /Annot /Subtype /Link /Border [0 0 0] /H /I /C [1 0 0] /Rect [485.977 188.006 518.403 199.959] /A << /S /GoTo /D <6571756174696F6E2E432E332E352E37> >> >>}
\MusicPDFWriteObject{6614}{<< /Type /Annot /Subtype /Link /Border [0 0 0] /H /I /C [1 0 0] /Rect [429.271 119.606 461.698 131.559] /A << /S /GoTo /D <6571756174696F6E2E432E332E352E37> >> >>}
\MusicPDFWriteObject{6615}{<< /Type /Annot /Subtype /Link /Border [0 0 0] /H /I /C [1 0 0] /Rect [304.488 670.946 336.914 682.898] /A << /S /GoTo /D <6571756174696F6E2E432E332E352E34> >> >>}
\MusicPDFWriteObject{6616}{<< /Type /Annot /Subtype /Link /Border [0 0 0] /H /I /C [1 0 0] /Rect [206.333 657.795 238.333 669.748] /A << /S /GoTo /D <6571756174696F6E2E432E332E352E34> >> >>}
\MusicPDFWriteObject{6617}{<< /Type /Annot /Subtype /Link /Border [0 0 0] /H /I /C [0 1 0] /Rect [216.986 607.682 250.955 616.55] /A << /S /GoTo /D <636974652E436E772D61> >> >>}
\MusicPDFWriteObject{6618}{<< /Type /Annot /Subtype /Link /Border [0 0 0] /H /I /C [0 1 0] /Rect [73.77 568.23 107.739 577.097] /A << /S /GoTo /D <636974652E436E772D61> >> >>}
\MusicPDFWriteObject{6619}{<< /Type /Annot /Subtype /Link /Border [0 0 0] /H /I /C [0 1 0] /Rect [244.056 360.031 278.025 368.899] /A << /S /GoTo /D <636974652E436E772D61> >> >>}
\MusicPDFWriteObject{6620}{<< /Type /Annot /Subtype /Link /Border [0 0 0] /H /I /C [1 0 0] /Rect [277.558 344.39 309.985 356.343] /A << /S /GoTo /D <6571756174696F6E2E432E332E352E34> >> >>}
\MusicPDFWriteObject{6621}{<< /Type /Annot /Subtype /Link /Border [0 0 0] /H /I /C [1 0 0] /Rect [377.261 344.39 410.296 356.343] /A << /S /GoTo /D <74686D2E432E332E352E33> >> >>}
\MusicPDFWriteObject{6622}{<< /Type /Annot /Subtype /Link /Border [0 0 0] /H /I /C [1 0 0] /Rect [212.855 188.099 250.261 200.052] /A << /S /GoTo /D <6571756174696F6E2E432E332E352E3134> >> >>}
\MusicPDFWriteObject{6623}{<< /Type /Annot /Subtype /Link /Border [0 0 0] /H /I /C [1 0 0] /Rect [226.278 151.671 263.684 163.624] /A << /S /GoTo /D <6571756174696F6E2E432E332E352E3135> >> >>}
\MusicPDFWriteObject{6624}{<< /Type /Annot /Subtype /Link /Border [0 0 0] /H /I /C [1 0 0] /Rect [321.647 151.671 354.073 163.626] /A << /S /GoTo /D <74686D2E432E332E352E33> >> >>}
\MusicPDFWriteObject{6625}{<< /Type /Annot /Subtype /Link /Border [0 0 0] /H /I /C [1 0 0] /Rect [222.258 130.284 254.076 142.237] /A << /S /GoTo /D <74686D2E432E332E352E31> >> >>}
\MusicPDFWriteObject{6626}{<< /Type /Annot /Subtype /Link /Border [0 0 0] /H /I /C [0 1 0] /Rect [354.645 132.774 378.342 141.642] /A << /S /GoTo /D <636974652E4B7562> >> >>}
\MusicPDFWriteObject{6627}{<< /Type /Annot /Subtype /Link /Border [0 0 0] /H /I /C [1 0 0] /Rect [73.712 632.03 111.118 643.983] /A << /S /GoTo /D <6571756174696F6E2E432E332E352E3131> >> >>}
\MusicPDFWriteObject{6628}{<< /Type /Annot /Subtype /Link /Border [0 0 0] /H /I /C [1 0 0] /Rect [175.417 632.03 207.843 643.983] /A << /S /GoTo /D <74686D2E432E332E352E34> >> >>}
\MusicPDFWriteObject{6629}{<< /Type /Annot /Subtype /Link /Border [0 0 0] /H /I /C [0 1 0] /Rect [232.056 598.093 265.547 606.961] /A << /S /GoTo /D <636974652E436E772D63> >> >>}
\MusicPDFWriteObject{6630}{<< /Type /Annot /Subtype /Link /Border [0 0 0] /H /I /C [1 0 0] /Rect [123.533 535.638 156.568 547.654] /A << /S /GoTo /D <74686D2E432E332E352E33> >> >>}
\MusicPDFWriteObject{6631}{<< /Type /Annot /Subtype /Link /Border [0 0 0] /H /I /C [1 0 0] /Rect [496.438 535.638 529.473 547.654] /A << /S /GoTo /D <74686D2E432E332E342E33> >> >>}
\MusicPDFWriteObject{6632}{<< /Type /Annot /Subtype /Link /Border [0 0 0] /H /I /C [0 1 0] /Rect [155.112 511.888 180.172 520.756] /A << /S /GoTo /D <636974652E53636861> >> >>}
\MusicPDFWriteObject{6633}{<< /Type /Annot /Subtype /Link /Border [0 0 0] /H /I /C [1 0 0] /Rect [108.58 496.248 145.987 508.2] /A << /S /GoTo /D <6571756174696F6E2E432E332E352E3136> >> >>}
\MusicPDFWriteObject{6634}{<< /Type /Annot /Subtype /Link /Border [0 0 0] /H /I /C [0 1 0] /Rect [130.018 412.534 163.51 421.402] /A << /S /GoTo /D <636974652E436E772D63> >> >>}
\MusicPDFWriteObject{6635}{<< /Type /Annot /Subtype /Link /Border [0 0 0] /H /I /C [0 1 0] /Rect [313.576 634.52 337.273 643.388] /A << /S /GoTo /D <636974652E4B7562> >> >>}
\MusicPDFWriteObject{6636}{<< /Type /Annot /Subtype /Link /Border [0 0 0] /H /I /C [0 1 0] /Rect [376.288 621.369 411.052 630.237] /A << /S /GoTo /D <636974652E436E772D62> >> >>}
\MusicPDFWriteObject{6637}{<< /Type /Annot /Subtype /Link /Border [0 0 0] /H /I /C [1 0 0] /Rect [160.399 557.795 193.434 571.542] /A << /S /GoTo /D <74686D2E432E332E352E31> >> >>}
\MusicPDFWriteObject{6638}{<< /Type /Annot /Subtype /Link /Border [0 0 0] /H /I /C [1 0 0] /Rect [450.569 546.436 483.513 558.391] /A << /S /GoTo /D <74686D2E432E332E352E31> >> >>}
\MusicPDFWriteObject{6639}{<< /Type /Annot /Subtype /Link /Border [0 0 0] /H /I /C [1 0 0] /Rect [504.482 546.436 541.888 558.389] /A << /S /GoTo /D <6571756174696F6E2E432E332E352E3138> >> >>}
\MusicPDFWriteObject{6640}{<< /Type /Annot /Subtype /Link /Border [0 0 0] /H /I /C [1 0 0] /Rect [121.234 474.182 154.269 486.134] /A << /S /GoTo /D <74686D2E432E332E352E33> >> >>}
\MusicPDFWriteObject{6641}{<< /Type /Annot /Subtype /Link /Border [0 0 0] /H /I /C [0 1 0] /Rect [324.644 476.672 358.613 485.54] /A << /S /GoTo /D <636974652E436E772D61> >> >>}
\MusicPDFWriteObject{6642}{<< /Type /Annot /Subtype /Link /Border [0 0 0] /H /I /C [1 0 0] /Rect [336.609 447.88 374.016 459.833] /A << /S /GoTo /D <6571756174696F6E2E432E332E352E3233> >> >>}
\MusicPDFWriteObject{6643}{<< /Type /Annot /Subtype /Link /Border [0 0 0] /H /I /C [1 0 0] /Rect [412.187 387.918 444.857 399.871] /A << /S /GoTo /D <74686D2E432E332E352E33> >> >>}
\MusicPDFWriteObject{6644}{<< /Type /Annot /Subtype /Link /Border [0 0 0] /H /I /C [0 1 0] /Rect [160.77 270.57 194.739 279.438] /A << /S /GoTo /D <636974652E436E772D61> >> >>}
\MusicPDFWriteObject{6645}{<< /Type /Annot /Subtype /Link /Border [0 0 0] /H /I /C [0 1 0] /Rect [352.933 257.419 371.132 266.287] /A << /S /GoTo /D <636974652E4A6172> >> >>}
\MusicPDFWriteObject{6646}{<< /Type /Annot /Subtype /Link /Border [0 0 0] /H /I /C [1 0 0] /Rect [226.787 95.747 259.214 107.699] /A << /S /GoTo /D <6571756174696F6E2E432E332E362E31> >> >>}
\MusicPDFWriteObject{6647}{<< /Type /Annot /Subtype /Link /Border [0 0 0] /H /I /C [1 0 0] /Rect [415.127 623.655 447.553 635.61] /A << /S /GoTo /D <74686D2E432E332E352E33> >> >>}
\MusicPDFWriteObject{6648}{<< /Type /Annot /Subtype /Link /Border [0 0 0] /H /I /C [1 0 0] /Rect [512 487.69 543.817 499.942] /A << /S /GoTo /D <74686D2E432E332E362E31> >> >>}
\MusicPDFWriteObject{6649}{<< /Type /Annot /Subtype /Link /Border [0 0 0] /H /I /C [0 1 0] /Rect [169.253 477.327 187.451 486.194] /A << /S /GoTo /D <636974652E4A6172> >> >>}
\MusicPDFWriteObject{6650}{<< /Type /Annot /Subtype /Link /Border [0 0 0] /H /I /C [0 1 0] /Rect [112.468 451.025 130.667 459.893] /A << /S /GoTo /D <636974652E4A6172> >> >>}
\MusicPDFWriteObject{6651}{<< /Type /Annot /Subtype /Link /Border [0 0 0] /H /I /C [1 0 0] /Rect [282.53 220.209 315.565 231] /A << /S /GoTo /D <74686D2E432E332E352E34> >> >>}
\MusicPDFWriteObject{6652}{<< /Type /Annot /Subtype /Link /Border [0 0 0] /H /I /C [0 1 0] /Rect [163.355 108.983 181.554 117.851] /A << /S /GoTo /D <636974652E4A6172> >> >>}
\MusicPDFWriteObject{6653}{<< /Type /Annot /Subtype /Link /Border [0 0 0] /H /I /C [1 0 0] /Rect [206.858 353.307 239.285 365.26] /A << /S /GoTo /D <6571756174696F6E2E432E332E362E38> >> >>}
\MusicPDFWriteObject{6654}{<< /Type /Annot /Subtype /Link /Border [0 0 0] /H /I /C [1 0 0] /Rect [238.468 340.156 270.894 352.109] /A << /S /GoTo /D <6571756174696F6E2E432E332E362E37> >> >>}
\MusicPDFWriteObject{6655}{<< /Type /Annot /Subtype /Link /Border [0 0 0] /H /I /C [1 0 0] /Rect [248.557 254.274 280.984 266.227] /A << /S /GoTo /D <6571756174696F6E2E432E332E362E38> >> >>}
\MusicPDFWriteObject{6656}{<< /Type /Annot /Subtype /Link /Border [0 0 0] /H /I /C [0 1 0] /Rect [514.92 254.827 543.162 265.632] /A << /S /GoTo /D <636974652E5279616E> >> >>}
\MusicPDFWriteObject{6657}{<< /Type /Annot /Subtype /Link /Border [0 0 0] /H /I /C [0 1 0] /Rect [73.74 228.526 101.982 239.33] /A << /S /GoTo /D <636974652E5279616E> >> >>}
\MusicPDFWriteObject{6658}{<< /Type /Annot /Subtype /Link /Border [0 0 0] /H /I /C [0 1 0] /Rect [73.77 202.224 102.012 213.029] /A << /S /GoTo /D <636974652E5279616E> >> >>}
\MusicPDFWriteObject{6659}{<< /Type /Annot /Subtype /Link /Border [0 0 0] /H /I /C [0 1 0] /Rect [342.182 189.074 370.424 199.878] /A << /S /GoTo /D <636974652E5279616E> >> >>}
\MusicPDFWriteObject{6660}{<< /Type /Annot /Subtype /Link /Border [0 0 0] /H /I /C [0 1 0] /Rect [514.917 189.074 543.159 199.878] /A << /S /GoTo /D <636974652E5279616E> >> >>}
\MusicPDFWriteObject{6661}{<< /Type /Annot /Subtype /Link /Border [0 0 0] /H /I /C [1 0 0] /Rect [71.004 128.665 103.43 140.98] /A << /S /GoTo /D <74686D2E432E332E352E31> >> >>}
\MusicPDFWriteObject{6662}{<< /Type /Annot /Subtype /Link /Border [0 0 0] /H /I /C [1 0 0] /Rect [168.171 128.665 200.598 140.98] /A << /S /GoTo /D <74686D2E432E332E352E35> >> >>}
\MusicPDFWriteObject{6663}{<< /Type /Annot /Subtype /Link /Border [0 0 0] /H /I /C [1 0 0] /Rect [422.008 448.893 454.434 460.846] /A << /S /GoTo /D <6571756174696F6E2E432E332E362E34> >> >>}
\MusicPDFWriteObject{6664}{<< /Type /Annot /Subtype /Link /Border [0 0 0] /H /I /C [1 0 0] /Rect [451.823 400.204 484.858 413.127] /A << /S /GoTo /D <74686D2E432E332E362E33> >> >>}
\MusicPDFWriteObject{6665}{<< /Type /Annot /Subtype /Link /Border [0 0 0] /H /I /C [0 1 0] /Rect [351.286 390.512 369.484 399.38] /A << /S /GoTo /D <636974652E4A6172> >> >>}
\MusicPDFWriteObject{6666}{<< /Type /Annot /Subtype /Link /Border [0 0 0] /H /I /C [1 0 0] /Rect [158.959 361.721 196.366 373.673] /A << /S /GoTo /D <6571756174696F6E2E432E332E362E3133> >> >>}
\MusicPDFWriteObject{6667}{<< /Type /Annot /Subtype /Link /Border [0 0 0] /H /I /C [1 0 0] /Rect [404.115 297.1 440.813 310.024] /A << /S /GoTo /D <6571756174696F6E2E432E332E352E3131> >> >>}
\MusicPDFWriteObject{6668}{<< /Type /Annot /Subtype /Link /Border [0 0 0] /H /I /C [1 0 0] /Rect [499.527 297.1 531.345 310.024] /A << /S /GoTo /D <74686D2E432E332E362E33> >> >>}
\MusicPDFWriteObject{6669}{<< /Type /Annot /Subtype /Link /Border [0 0 0] /H /I /C [0 1 0] /Rect [251.81 234.806 270.009 243.673] /A << /S /GoTo /D <636974652E4A6172> >> >>}
\MusicPDFWriteObject{6670}{<< /Type /Annot /Subtype /Link /Border [0 0 0] /H /I /C [1 0 0] /Rect [166.201 156.436 203.608 168.388] /A << /S /GoTo /D <6571756174696F6E2E432E332E362E3135> >> >>}
\MusicPDFWriteObject{6671}{<< /Type /Annot /Subtype /Link /Border [0 0 0] /H /I /C [1 0 0] /Rect [345.25 684.633 382.657 696.586] /A << /S /GoTo /D <6571756174696F6E2E432E332E362E3137> >> >>}
\MusicPDFWriteObject{6672}{<< /Type /Annot /Subtype /Link /Border [0 0 0] /H /I /C [1 0 0] /Rect [264.836 564.808 302.242 576.761] /A << /S /GoTo /D <6571756174696F6E2E432E332E362E3139> >> >>}
\MusicPDFWriteObject{6673}{<< /Type /Annot /Subtype /Link /Border [0 0 0] /H /I /C [1 0 0] /Rect [280.125 543.731 313.16 554.522] /A << /S /GoTo /D <74686D2E432E332E362E35> >> >>}
\MusicPDFWriteObject{6674}{<< /Type /Annot /Subtype /Link /Border [0 0 0] /H /I /C [1 0 0] /Rect [410.896 543.731 443.932 554.522] /A << /S /GoTo /D <74686D2E432E332E352E31> >> >>}
\MusicPDFWriteObject{6675}{<< /Type /Annot /Subtype /Link /Border [0 0 0] /H /I /C [1 0 0] /Rect [511.81 543.731 544.845 554.522] /A << /S /GoTo /D <74686D2E432E332E352E35> >> >>}
\MusicPDFWriteObject{6676}{<< /Type /Annot /Subtype /Link /Border [0 0 0] /H /I /C [1 0 0] /Rect [440.071 529.73 473.106 541.982] /A << /S /GoTo /D <74686D2E432E332E362E35> >> >>}
\MusicPDFWriteObject{6677}{<< /Type /Annot /Subtype /Link /Border [0 0 0] /H /I /C [1 0 0] /Rect [295.977 515.321 328.403 528.831] /A << /S /GoTo /D <74686D2E432E332E352E31> >> >>}
\MusicPDFWriteObject{6678}{<< /Type /Annot /Subtype /Link /Border [0 0 0] /H /I /C [1 0 0] /Rect [406.895 515.321 439.321 528.831] /A << /S /GoTo /D <74686D2E432E332E362E35> >> >>}
\MusicPDFWriteObject{6679}{<< /Type /Annot /Subtype /Link /Border [0 0 0] /H /I /C [1 0 0] /Rect [133.356 503.664 166.056 515.07] /A << /S /GoTo /D <74686D2E432E332E352E35> >> >>}
\MusicPDFWriteObject{6680}{<< /Type /Annot /Subtype /Link /Border [0 0 0] /H /I /C [1 0 0] /Rect [296.574 503.664 329.274 515.07] /A << /S /GoTo /D <74686D2E432E332E362E35> >> >>}
\MusicPDFWriteObject{6681}{<< /Type /Annot /Subtype /Link /Border [0 0 0] /H /I /C [0 1 0] /Rect [426.272 418.832 459.764 427.7] /A << /S /GoTo /D <636974652E436E772D63> >> >>}
\MusicPDFWriteObject{6682}{<< /Type /Annot /Subtype /Link /Border [0 0 0] /H /I /C [1 0 0] /Rect [509.839 402.831 541.931 416.391] /A << /S /GoTo /D <74686D2E432E332E362E34> >> >>}
\MusicPDFWriteObject{6683}{<< /Type /Annot /Subtype /Link /Border [0 0 0] /H /I /C [1 0 0] /Rect [467.743 293.329 492.86 305.282] /A << /S /GoTo /D <73756273656374696F6E2E412E352E31> >> >>}
\MusicPDFWriteObject{6684}{<< /Type /Annot /Subtype /Link /Border [0 0 0] /H /I /C [1 0 0] /Rect [518.078 293.329 543.195 305.282] /A << /S /GoTo /D <73756273656374696F6E2E412E352E32> >> >>}
\MusicPDFWriteObject{6685}{<< /Type /Annot /Subtype /Link /Border [0 0 0] /H /I /C [1 0 0] /Rect [75.432 280.731 100.686 291.579] /A << /S /GoTo /D <73756273656374696F6E2E412E352E33> >> >>}
\MusicPDFWriteObject{6686}{<< /Type /Annot /Subtype /Link /Border [0 0 0] /H /I /C [1 0 0] /Rect [126.011 280.731 151.266 291.579] /A << /S /GoTo /D <73756273656374696F6E2E412E352E34> >> >>}
\MusicPDFWriteObject{6687}{<< /Type /Annot /Subtype /Link /Border [0 0 0] /H /I /C [1 0 0] /Rect [210.229 280.731 227.636 291.579] /A << /S /GoTo /D <73656374696F6E2E412E35> >> >>}
\MusicPDFWriteObject{6688}{<< /Type /Annot /Subtype /Link /Border [0 0 0] /H /I /C [0 1 0] /Rect [447.048 243.216 478.992 252.084] /A << /S /GoTo /D <636974652E427269732D61> >> >>}
\MusicPDFWriteObject{6689}{<< /Type /Annot /Subtype /Link /Border [0 0 0] /H /I /C [0 1 0] /Rect [103.819 230.065 136.559 238.933] /A << /S /GoTo /D <636974652E427269732D62> >> >>}
\MusicPDFWriteObject{6690}{<< /Type /Annot /Subtype /Link /Border [0 0 0] /H /I /C [0 1 0] /Rect [136.07 190.292 168.014 199.16] /A << /S /GoTo /D <636974652E427269732D61> >> >>}
\MusicPDFWriteObject{6691}{<< /Type /Annot /Subtype /Link /Border [0 0 0] /H /I /C [0 1 0] /Rect [130.309 150.84 163.048 159.708] /A << /S /GoTo /D <636974652E427269732D62> >> >>}
\MusicPDFWriteObject{6692}{<< /Type /Annot /Subtype /Link /Border [0 0 0] /H /I /C [0 1 0] /Rect [73.77 673.973 105.714 682.84] /A << /S /GoTo /D <636974652E427269732D61> >> >>}
\MusicPDFWriteObject{6693}{<< /Type /Annot /Subtype /Link /Border [0 0 0] /H /I /C [0 1 0] /Rect [266.591 660.822 303.541 669.69] /A << /S /GoTo /D <636974652E417274687572> >> >>}
\MusicPDFWriteObject{6694}{<< /Type /Annot /Subtype /Link /Border [0 0 0] /H /I /C [0 1 0] /Rect [312.546 568.766 345.285 577.634] /A << /S /GoTo /D <636974652E427269732D62> >> >>}
\MusicPDFWriteObject{6695}{<< /Type /Annot /Subtype /Link /Border [0 0 0] /H /I /C [1 0 0] /Rect [163.631 317.936 196.057 329.889] /A << /S /GoTo /D <6571756174696F6E2E432E332E372E32> >> >>}
\MusicPDFWriteObject{6696}{<< /Type /Annot /Subtype /Link /Border [0 0 0] /H /I /C [0 1 0] /Rect [146.898 297.606 178.842 306.474] /A << /S /GoTo /D <636974652E427269732D61> >> >>}
\MusicPDFWriteObject{6697}{<< /Type /Annot /Subtype /Link /Border [0 0 0] /H /I /C [1 0 0] /Rect [370.339 295.116 403.374 307.069] /A << /S /GoTo /D <74686D2E432E332E372E31> >> >>}
\MusicPDFWriteObject{6698}{<< /Type /Annot /Subtype /Link /Border [0 0 0] /H /I /C [0 1 0] /Rect [465.358 297.606 497.301 306.474] /A << /S /GoTo /D <636974652E427269732D61> >> >>}
\MusicPDFWriteObject{6699}{<< /Type /Annot /Subtype /Link /Border [0 0 0] /H /I /C [1 0 0] /Rect [193.653 281.966 226.079 293.918] /A << /S /GoTo /D <74686D2E432E332E372E31> >> >>}
\MusicPDFWriteObject{6700}{<< /Type /Annot /Subtype /Link /Border [0 0 0] /H /I /C [1 0 0] /Rect [385.383 243.067 402.303 253.857] /A << /S /GoTo /D <73656374696F6E2E412E35> >> >>}
\MusicPDFWriteObject{6701}{<< /Type /Annot /Subtype /Link /Border [0 0 0] /H /I /C [0 1 0] /Rect [322.471 231.853 354.415 240.72] /A << /S /GoTo /D <636974652E427269732D61> >> >>}
\MusicPDFWriteObject{6702}{<< /Type /Annot /Subtype /Link /Border [0 0 0] /H /I /C [0 1 0] /Rect [181.23 598.093 206.29 606.961] /A << /S /GoTo /D <636974652E53636861> >> >>}
\MusicPDFWriteObject{6703}{<< /Type /Annot /Subtype /Link /Border [0 0 0] /H /I /C [0 1 0] /Rect [386.925 522.213 435.966 531.081] /A << /S /GoTo /D <636974652E426F752D545653> >> >>}
\MusicPDFWriteObject{6704}{<< /Type /Annot /Subtype /Link /Border [0 0 0] /H /I /C [0 1 0] /Rect [289.121 430.158 314.181 439.025] /A << /S /GoTo /D <636974652E53636861> >> >>}
\MusicPDFWriteObject{6705}{<< /Type /Annot /Subtype /Link /Border [0 0 0] /H /I /C [0 1 0] /Rect [248.793 354.278 273.854 363.146] /A << /S /GoTo /D <636974652E53636861> >> >>}
\MusicPDFWriteObject{6706}{<< /Type /Annot /Subtype /Link /Border [0 0 0] /H /I /C [1 0 0] /Rect [73.712 305.266 106.138 317.218] /A << /S /GoTo /D <6571756174696F6E2E432E342E312E31> >> >>}
\MusicPDFWriteObject{6707}{<< /Type /Annot /Subtype /Link /Border [0 0 0] /H /I /C [1 0 0] /Rect [382.659 629.937 415.695 641.91] /A << /S /GoTo /D <6571756174696F6E2E432E342E312E31> >> >>}
\MusicPDFWriteObject{6708}{<< /Type /Annot /Subtype /Link /Border [0 0 0] /H /I /C [0 1 0] /Rect [297.417 572.394 322.477 581.262] /A << /S /GoTo /D <636974652E53636861> >> >>}
\MusicPDFWriteObject{6709}{<< /Type /Annot /Subtype /Link /Border [0 0 0] /H /I /C [0 1 0] /Rect [389.578 553.266 414.638 562.133] /A << /S /GoTo /D <636974652E53636861> >> >>}
\MusicPDFWriteObject{6710}{<< /Type /Annot /Subtype /Link /Border [0 0 0] /H /I /C [0 1 0] /Rect [433.778 540.115 452.883 548.983] /A << /S /GoTo /D <636974652E4E42> >> >>}
\MusicPDFWriteObject{6711}{<< /Type /Annot /Subtype /Link /Border [0 0 0] /H /I /C [0 1 0] /Rect [272.184 526.964 297.244 535.832] /A << /S /GoTo /D <636974652E53636861> >> >>}
\MusicPDFWriteObject{6712}{<< /Type /Annot /Subtype /Link /Border [0 0 0] /H /I /C [0 1 0] /Rect [494.229 500.663 543.27 509.53] /A << /S /GoTo /D <636974652E426F752D545653> >> >>}
\MusicPDFWriteObject{6713}{<< /Type /Annot /Subtype /Link /Border [0 0 0] /H /I /C [0 1 0] /Rect [199.382 487.512 248.422 496.38] /A << /S /GoTo /D <636974652E426F752D545653> >> >>}
\MusicPDFWriteObject{6714}{<< /Type /Annot /Subtype /Link /Border [0 0 0] /H /I /C [0 1 0] /Rect [270.262 455.516 301.091 464.384] /A << /S /GoTo /D <636974652E47726F2D61> >> >>}
\MusicPDFWriteObject{6715}{<< /Type /Annot /Subtype /Link /Border [0 0 0] /H /I /C [0 1 0] /Rect [207.006 410.369 224.161 419.237] /A << /S /GoTo /D <636974652E466F6C> >> >>}
\MusicPDFWriteObject{6716}{<< /Type /Annot /Subtype /Link /Border [0 0 0] /H /I /C [0 1 0] /Rect [281.717 384.068 298.872 392.935] /A << /S /GoTo /D <636974652E466F6C> >> >>}
\MusicPDFWriteObject{6717}{<< /Type /Annot /Subtype /Link /Border [0 0 0] /H /I /C [1 0 0] /Rect [410.134 219.798 442.561 231.751] /A << /S /GoTo /D <6571756174696F6E2E432E342E322E31> >> >>}
\MusicPDFWriteObject{6718}{<< /Type /Annot /Subtype /Link /Border [0 0 0] /H /I /C [1 0 0] /Rect [250.637 193.496 283.064 205.449] /A << /S /GoTo /D <6571756174696F6E2E432E342E322E31> >> >>}
\MusicPDFWriteObject{6719}{<< /Type /Annot /Subtype /Link /Border [0 0 0] /H /I /C [0 1 0] /Rect [447.025 162.054 483.746 172.858] /A << /S /GoTo /D <636974652E4C616E672D64> >> >>}
\MusicPDFWriteObject{6720}{<< /Type /Annot /Subtype /Link /Border [0 0 0] /H /I /C [1 0 0] /Rect [516.678 108.897 541.811 120.85] /A << /S /GoTo /D <73756273656374696F6E2E432E312E34> >> >>}
\MusicPDFWriteObject{6721}{<< /Type /Annot /Subtype /Link /Border [0 0 0] /H /I /C [0 1 0] /Rect [117.987 85.086 136.96 93.954] /A << /S /GoTo /D <636974652E547265> >> >>}
\MusicPDFWriteObject{6722}{<< /Type /Annot /Subtype /Link /Border [0 0 0] /H /I /C [0 1 0] /Rect [414.081 687.123 428.642 695.991] /A << /S /GoTo /D <636974652E5A61> >> >>}
\MusicPDFWriteObject{6723}{<< /Type /Annot /Subtype /Link /Border [0 0 0] /H /I /C [0 1 0] /Rect [399.386 673.973 418.975 682.84] /A << /S /GoTo /D <636974652E4455> >> >>}
\MusicPDFWriteObject{6724}{<< /Type /Annot /Subtype /Link /Border [0 0 0] /H /I /C [0 1 0] /Rect [244.064 575.593 261.219 584.461] /A << /S /GoTo /D <636974652E466F6C> >> >>}
\MusicPDFWriteObject{6725}{<< /Type /Annot /Subtype /Link /Border [0 0 0] /H /I /C [0 1 0] /Rect [375.337 562.443 400.397 571.31] /A << /S /GoTo /D <636974652E53636861> >> >>}
\MusicPDFWriteObject{6726}{<< /Type /Annot /Subtype /Link /Border [0 0 0] /H /I /C [0 1 0] /Rect [134.793 522.99 151.947 531.858] /A << /S /GoTo /D <636974652E466F6C> >> >>}
\MusicPDFWriteObject{6727}{<< /Type /Annot /Subtype /Link /Border [0 0 0] /H /I /C [0 1 0] /Rect [193.381 306.929 223.774 315.797] /A << /S /GoTo /D <636974652E426F752D49> >> >>}
\MusicPDFWriteObject{6728}{<< /Type /Annot /Subtype /Link /Border [0 0 0] /H /I /C [1 0 0] /Rect [512.588 304.439 545.592 316.392] /A << /S /GoTo /D <74686D2E432E342E332E31> >> >>}
\MusicPDFWriteObject{6729}{<< /Type /Annot /Subtype /Link /Border [0 0 0] /H /I /C [0 1 0] /Rect [455.298 229.194 504.338 238.062] /A << /S /GoTo /D <636974652E426F752D545653> >> >>}
\MusicPDFWriteObject{6730}{<< /Type /Annot /Subtype /Link /Border [0 0 0] /H /I /C [0 1 0] /Rect [433.287 189.742 450.427 198.609] /A << /S /GoTo /D <636974652E4453> >> >>}
\MusicPDFWriteObject{6731}{<< /Type /Annot /Subtype /Link /Border [0 0 0] /H /I /C [1 0 0] /Rect [227.938 610.23 260.364 622.183] /A << /S /GoTo /D <6571756174696F6E2E432E342E332E32> >> >>}
\MusicPDFWriteObject{6732}{<< /Type /Annot /Subtype /Link /Border [0 0 0] /H /I /C [1 0 0] /Rect [192.55 597.079 224.976 609.032] /A << /S /GoTo /D <6571756174696F6E2E432E342E332E32> >> >>}
\MusicPDFWriteObject{6733}{<< /Type /Annot /Subtype /Link /Border [0 0 0] /H /I /C [0 1 0] /Rect [185.789 553.145 214.577 563.949] /A << /S /GoTo /D <636974652E566F696774> >> >>}
\MusicPDFWriteObject{6734}{<< /Type /Annot /Subtype /Link /Border [0 0 0] /H /I /C [0 1 0] /Rect [371.093 555.081 401.486 563.949] /A << /S /GoTo /D <636974652E426F752D49> >> >>}
\MusicPDFWriteObject{6735}{<< /Type /Annot /Subtype /Link /Border [0 0 0] /H /I /C [1 0 0] /Rect [194.734 513.138 227.16 525.093] /A << /S /GoTo /D <74686D2E432E342E332E31> >> >>}
\MusicPDFWriteObject{6736}{<< /Type /Annot /Subtype /Link /Border [0 0 0] /H /I /C [0 1 0] /Rect [241.31 376.359 271.702 385.227] /A << /S /GoTo /D <636974652E426F752D49> >> >>}
\MusicPDFWriteObject{6737}{<< /Type /Annot /Subtype /Link /Border [0 0 0] /H /I /C [1 0 0] /Rect [280.208 329.381 312.726 341.336] /A << /S /GoTo /D <74686D2E432E342E332E32> >> >>}
\MusicPDFWriteObject{6738}{<< /Type /Annot /Subtype /Link /Border [0 0 0] /H /I /C [1 0 0] /Rect [196.307 232.843 229.342 243.634] /A << /S /GoTo /D <74686D2E432E342E332E31> >> >>}
\MusicPDFWriteObject{6739}{<< /Type /Annot /Subtype /Link /Border [0 0 0] /H /I /C [1 0 0] /Rect [176.001 219.139 208.427 231.092] /A << /S /GoTo /D <74686D2E432E342E332E35> >> >>}
\MusicPDFWriteObject{6740}{<< /Type /Annot /Subtype /Link /Border [0 0 0] /H /I /C [1 0 0] /Rect [106.111 192.838 138.111 204.791] /A << /S /GoTo /D <74686D2E432E342E332E31> >> >>}
\MusicPDFWriteObject{6741}{<< /Type /Annot /Subtype /Link /Border [0 0 0] /H /I /C [1 0 0] /Rect [337.233 179.687 370.269 191.64] /A << /S /GoTo /D <74686D2E432E342E322E31> >> >>}
\MusicPDFWriteObject{6742}{<< /Type /Annot /Subtype /Link /Border [0 0 0] /H /I /C [0 1 0] /Rect [154.911 129.574 179.971 138.442] /A << /S /GoTo /D <636974652E53636861> >> >>}
\MusicPDFWriteObject{6743}{<< /Type /Annot /Subtype /Link /Border [0 0 0] /H /I /C [0 1 0] /Rect [357.451 489.862 382.511 498.73] /A << /S /GoTo /D <636974652E53636861> >> >>}
\MusicPDFWriteObject{6744}{<< /Type /Annot /Subtype /Link /Border [0 0 0] /H /I /C [0 1 0] /Rect [254.88 476.711 273.854 485.579] /A << /S /GoTo /D <636974652E547265> >> >>}
\MusicPDFWriteObject{6745}{<< /Type /Annot /Subtype /Link /Border [0 0 0] /H /I /C [1 0 0] /Rect [182.567 408.467 214.384 420.42] /A << /S /GoTo /D <74686D2E432E332E332E31> >> >>}
\MusicPDFWriteObject{6746}{<< /Type /Annot /Subtype /Link /Border [0 0 0] /H /I /C [0 1 0] /Rect [465.675 332.053 490.735 340.921] /A << /S /GoTo /D <636974652E53636861> >> >>}
\MusicPDFWriteObject{6747}{<< /Type /Annot /Subtype /Link /Border [0 0 0] /H /I /C [0 1 0] /Rect [124.853 111.387 155.682 120.255] /A << /S /GoTo /D <636974652E47726F2D61> >> >>}
\MusicPDFWriteObject{6748}{<< /Type /Annot /Subtype /Link /Border [0 0 0] /H /I /C [0 1 0] /Rect [430.431 85.086 455.491 93.954] /A << /S /GoTo /D <636974652E53636861> >> >>}
\MusicPDFWriteObject{6749}{<< /Type /Annot /Subtype /Link /Border [0 0 0] /H /I /C [1 0 0] /Rect [261.797 671.482 294.224 683.437] /A << /S /GoTo /D <74686D2E432E342E332E31> >> >>}
\MusicPDFWriteObject{6750}{<< /Type /Annot /Subtype /Link /Border [0 0 0] /H /I /C [1 0 0] /Rect [423.6 599.023 456.392 610.976] /A << /S /GoTo /D <74686D2E432E342E312E32> >> >>}
\MusicPDFWriteObject{6751}{<< /Type /Annot /Subtype /Link /Border [0 0 0] /H /I /C [0 1 0] /Rect [268.233 567.832 298.626 576.7] /A << /S /GoTo /D <636974652E426F752D49> >> >>}
\MusicPDFWriteObject{6752}{<< /Type /Annot /Subtype /Link /Border [0 0 0] /H /I /C [1 0 0] /Rect [251.307 466.46 283.277 477.251] /A << /S /GoTo /D <74686D2E432E342E332E35> >> >>}
\MusicPDFWriteObject{6753}{<< /Type /Annot /Subtype /Link /Border [0 0 0] /H /I /C [0 1 0] /Rect [442.771 387.883 462.36 396.751] /A << /S /GoTo /D <636974652E4455> >> >>}
\MusicPDFWriteObject{6754}{<< /Type /Annot /Subtype /Link /Border [0 0 0] /H /I /C [1 0 0] /Rect [275.436 280.802 307.376 291.592] /A << /S /GoTo /D <74686D2E432E342E332E37> >> >>}
\MusicPDFWriteObject{6755}{<< /Type /Annot /Subtype /Link /Border [0 0 0] /H /I /C [1 0 0] /Rect [212.4 108.897 244.827 120.85] /A << /S /GoTo /D <6571756174696F6E2E432E342E332E35> >> >>}
\MusicPDFWriteObject{6756}{<< /Type /Annot /Subtype /Link /Border [0 0 0] /H /I /C [1 0 0] /Rect [231.778 82.596 264.205 94.549] /A << /S /GoTo /D <6571756174696F6E2E432E342E332E35> >> >>}
\MusicPDFWriteObject{6757}{<< /Type /Annot /Subtype /Link /Border [0 0 0] /H /I /C [1 0 0] /Rect [295.966 82.596 328.393 94.549] /A << /S /GoTo /D <6571756174696F6E2E432E342E312E32> >> >>}
\MusicPDFWriteObject{6758}{<< /Type /Annot /Subtype /Link /Border [0 0 0] /H /I /C [1 0 0] /Rect [397.505 82.596 429.931 94.549] /A << /S /GoTo /D <6571756174696F6E2E432E342E332E36> >> >>}
\MusicPDFWriteObject{6759}{<< /Type /Annot /Subtype /Link /Border [0 0 0] /H /I /C [0 1 0] /Rect [207.231 638.471 256.272 647.338] /A << /S /GoTo /D <636974652E426F752D545653> >> >>}
\MusicPDFWriteObject{6760}{<< /Type /Annot /Subtype /Link /Border [0 0 0] /H /I /C [1 0 0] /Rect [307.572 622.83 339.999 634.782] /A << /S /GoTo /D <6571756174696F6E2E432E342E332E37> >> >>}
\MusicPDFWriteObject{6761}{<< /Type /Annot /Subtype /Link /Border [0 0 0] /H /I /C [0 1 0] /Rect [379.847 550.348 428.888 559.216] /A << /S /GoTo /D <636974652E426F752D545653> >> >>}
\MusicPDFWriteObject{6762}{<< /Type /Annot /Subtype /Link /Border [0 0 0] /H /I /C [0 1 0] /Rect [323.576 537.198 372.616 546.065] /A << /S /GoTo /D <636974652E426F752D545653> >> >>}
\MusicPDFWriteObject{6763}{<< /Type /Annot /Subtype /Link /Border [0 0 0] /H /I /C [1 0 0] /Rect [217.446 521.557 249.873 533.509] /A << /S /GoTo /D <6571756174696F6E2E432E342E332E38> >> >>}
\MusicPDFWriteObject{6764}{<< /Type /Annot /Subtype /Link /Border [0 0 0] /H /I /C [1 0 0] /Rect [216.595 472.887 249.021 484.839] /A << /S /GoTo /D <6571756174696F6E2E432E342E332E35> >> >>}
\MusicPDFWriteObject{6765}{<< /Type /Annot /Subtype /Link /Border [0 0 0] /H /I /C [1 0 0] /Rect [281.283 458.684 314.227 472.244] /A << /S /GoTo /D <74686D2E432E342E312E31> >> >>}
\MusicPDFWriteObject{6766}{<< /Type /Annot /Subtype /Link /Border [0 0 0] /H /I /C [1 0 0] /Rect [471.554 446.585 503.98 458.538] /A << /S /GoTo /D <6571756174696F6E2E432E342E332E36> >> >>}
\MusicPDFWriteObject{6767}{<< /Type /Annot /Subtype /Link /Border [0 0 0] /H /I /C [0 1 0] /Rect [183.918 341.5 214.311 350.368] /A << /S /GoTo /D <636974652E426F752D49> >> >>}
\MusicPDFWriteObject{6768}{<< /Type /Annot /Subtype /Link /Border [0 0 0] /H /I /C [1 0 0] /Rect [129.882 305.685 162.308 317.638] /A << /S /GoTo /D <6571756174696F6E2E432E342E332E39> >> >>}
\MusicPDFWriteObject{6769}{<< /Type /Annot /Subtype /Link /Border [0 0 0] /H /I /C [1 0 0] /Rect [468.168 305.685 500.595 317.638] /A << /S /GoTo /D <6571756174696F6E2E432E342E332E39> >> >>}
\MusicPDFWriteObject{6770}{<< /Type /Annot /Subtype /Link /Border [0 0 0] /H /I /C [1 0 0] /Rect [482.375 604.492 519.782 616.445] /A << /S /GoTo /D <6571756174696F6E2E432E342E332E3132> >> >>}
\MusicPDFWriteObject{6771}{<< /Type /Annot /Subtype /Link /Border [0 0 0] /H /I /C [1 0 0] /Rect [299.974 549.113 337.381 561.066] /A << /S /GoTo /D <6571756174696F6E2E432E342E332E3133> >> >>}
\MusicPDFWriteObject{6772}{<< /Type /Annot /Subtype /Link /Border [0 0 0] /H /I /C [1 0 0] /Rect [364.896 536.515 397.931 547.306] /A << /S /GoTo /D <74686D2E432E342E332E37> >> >>}
\MusicPDFWriteObject{6773}{<< /Type /Annot /Subtype /Link /Border [0 0 0] /H /I /C [1 0 0] /Rect [476.438 522.811 509.29 534.766] /A << /S /GoTo /D <74686D2E432E342E332E39> >> >>}
\MusicPDFWriteObject{6774}{<< /Type /Annot /Subtype /Link /Border [0 0 0] /H /I /C [1 0 0] /Rect [359.641 249.6 391.459 263.16] /A << /S /GoTo /D <74686D2E432E342E332E35> >> >>}
\MusicPDFWriteObject{6775}{<< /Type /Annot /Subtype /Link /Border [0 0 0] /H /I /C [1 0 0] /Rect [246.457 238.055 278.884 248.845] /A << /S /GoTo /D <74686D2E432E342E332E39> >> >>}
\MusicPDFWriteObject{6776}{<< /Type /Annot /Subtype /Link /Border [0 0 0] /H /I /C [1 0 0] /Rect [504.482 192.063 541.888 204.016] /A << /S /GoTo /D <6571756174696F6E2E432E342E332E3136> >> >>}
\MusicPDFWriteObject{6777}{<< /Type /Annot /Subtype /Link /Border [0 0 0] /H /I /C [1 0 0] /Rect [284.955 83.149 317.381 93.94] /A << /S /GoTo /D <74686D2E432E332E332E31> >> >>}
\MusicPDFWriteObject{6778}{<< /Type /Annot /Subtype /Link /Border [0 0 0] /H /I /C [1 0 0] /Rect [364.867 637.51 374.8 649.815] /A << /S /GoTo /D <48666F6F746E6F74652E313239> >> >>}
\MusicPDFWriteObject{6779}{<< /Type /Annot /Subtype /Link /Border [0 0 0] /H /I /C [1 0 0] /Rect [278.857 247.056 311.283 259.009] /A << /S /GoTo /D <6571756174696F6E2E432E342E342E33> >> >>}
\MusicPDFWriteObject{6780}{<< /Type /Annot /Subtype /Link /Border [0 0 0] /H /I /C [1 0 0] /Rect [441.321 247.056 473.747 259.009] /A << /S /GoTo /D <6571756174696F6E2E432E342E342E31> >> >>}
\MusicPDFWriteObject{6781}{<< /Type /Annot /Subtype /Link /Border [0 0 0] /H /I /C [1 0 0] /Rect [313.333 213.084 345.151 223.932] /A << /S /GoTo /D <74686D2E432E342E342E31> >> >>}
\MusicPDFWriteObject{6782}{<< /Type /Annot /Subtype /Link /Border [0 0 0] /H /I /C [0 1 0] /Rect [367.684 201.871 398.077 210.739] /A << /S /GoTo /D <636974652E426F752D49> >> >>}
\MusicPDFWriteObject{6783}{<< /Type /Annot /Subtype /Link /Border [0 0 0] /H /I /C [0 1 0] /Rect [73.79 162.418 98.85 171.286] /A << /S /GoTo /D <636974652E53636861> >> >>}
\MusicPDFWriteObject{6784}{<< /Type /Annot /Subtype /Link /Border [0 0 0] /H /I /C [0 1 0] /Rect [512.802 162.418 543.195 171.286] /A << /S /GoTo /D <636974652E426F752D49> >> >>}
\MusicPDFWriteObject{6785}{<< /Type /Annot /Subtype /Link /Border [0 0 0] /H /I /C [1 0 0] /Rect [212.594 132.331 244.959 145.891] /A << /S /GoTo /D <74686D2E432E342E332E35> >> >>}
\MusicPDFWriteObject{6786}{<< /Type /Annot /Subtype /Link /Border [0 0 0] /H /I /C [0 1 0] /Rect [173.091 687.123 203.484 695.991] /A << /S /GoTo /D <636974652E426F752D49> >> >>}
\MusicPDFWriteObject{6787}{<< /Type /Annot /Subtype /Link /Border [0 0 0] /H /I /C [0 1 0] /Rect [319.218 687.123 349.611 695.991] /A << /S /GoTo /D <636974652E426F752D49> >> >>}
\MusicPDFWriteObject{6788}{<< /Type /Annot /Subtype /Link /Border [0 0 0] /H /I /C [0 1 0] /Rect [460.34 673.48 490.733 682.348] /A << /S /GoTo /D <636974652E426F752D49> >> >>}
\MusicPDFWriteObject{6789}{<< /Type /Annot /Subtype /Link /Border [0 0 0] /H /I /C [0 1 0] /Rect [460.616 634.028 491.009 642.896] /A << /S /GoTo /D <636974652E426F752D49> >> >>}
\MusicPDFWriteObject{6790}{<< /Type /Annot /Subtype /Link /Border [0 0 0] /H /I /C [1 0 0] /Rect [473.372 605.236 505.799 617.189] /A << /S /GoTo /D <6571756174696F6E2E432E342E332E33> >> >>}
\MusicPDFWriteObject{6791}{<< /Type /Annot /Subtype /Link /Border [0 0 0] /H /I /C [1 0 0] /Rect [368.254 592.085 400.68 604.038] /A << /S /GoTo /D <6571756174696F6E2E432E342E342E31> >> >>}
\MusicPDFWriteObject{6792}{<< /Type /Annot /Subtype /Link /Border [0 0 0] /H /I /C [1 0 0] /Rect [71.004 577.882 104.039 591.443] /A << /S /GoTo /D <74686D2E432E342E332E35> >> >>}
\MusicPDFWriteObject{6793}{<< /Type /Annot /Subtype /Link /Border [0 0 0] /H /I /C [1 0 0] /Rect [177.528 565.783 209.955 577.756] /A << /S /GoTo /D <74686D2E432E342E342E31> >> >>}
\MusicPDFWriteObject{6794}{<< /Type /Annot /Subtype /Link /Border [0 0 0] /H /I /C [0 1 0] /Rect [524.926 548.048 543.228 556.916] /A << /S /GoTo /D <636974652E4445> >> >>}
\MusicPDFWriteObject{6795}{<< /Type /Annot /Subtype /Link /Border [0 0 0] /H /I /C [1 0 0] /Rect [420.39 412.478 452.816 424.43] /A << /S /GoTo /D <6571756174696F6E2E432E342E342E35> >> >>}
\MusicPDFWriteObject{6796}{<< /Type /Annot /Subtype /Link /Border [0 0 0] /H /I /C [1 0 0] /Rect [370.967 379.101 403.393 391.053] /A << /S /GoTo /D <6571756174696F6E2E432E342E342E35> >> >>}
\MusicPDFWriteObject{6797}{<< /Type /Annot /Subtype /Link /Border [0 0 0] /H /I /C [0 1 0] /Rect [524.857 381.591 543.159 390.459] /A << /S /GoTo /D <636974652E4445> >> >>}
\MusicPDFWriteObject{6798}{<< /Type /Annot /Subtype /Link /Border [0 0 0] /H /I /C [1 0 0] /Rect [164.674 319.975 197.1 330.823] /A << /S /GoTo /D <74686D2E432E342E342E33> >> >>}
\MusicPDFWriteObject{6799}{<< /Type /Annot /Subtype /Link /Border [0 0 0] /H /I /C [1 0 0] /Rect [305.102 273.515 337.529 285.468] /A << /S /GoTo /D <6571756174696F6E2E432E342E342E35> >> >>}
\MusicPDFWriteObject{6800}{<< /Type /Annot /Subtype /Link /Border [0 0 0] /H /I /C [1 0 0] /Rect [303.344 253.289 336.379 265.242] /A << /S /GoTo /D <6571756174696F6E2E432E342E342E35> >> >>}
\MusicPDFWriteObject{6801}{<< /Type /Annot /Subtype /Link /Border [0 0 0] /H /I /C [1 0 0] /Rect [253.388 240.692 263.321 252.996] /A << /S /GoTo /D <48666F6F746E6F74652E313330> >> >>}
\MusicPDFWriteObject{6802}{<< /Type /Annot /Subtype /Link /Border [0 0 0] /H /I /C [1 0 0] /Rect [384.467 193.61 418.653 205.565] /A << /S /GoTo /D <74686D2E432E342E342E33> >> >>}
\MusicPDFWriteObject{6803}{<< /Type /Annot /Subtype /Link /Border [0 0 0] /H /I /C [1 0 0] /Rect [504.482 514.908 541.888 526.861] /A << /S /GoTo /D <6571756174696F6E2E432E342E342E3130> >> >>}
\MusicPDFWriteObject{6804}{<< /Type /Annot /Subtype /Link /Border [0 0 0] /H /I /C [1 0 0] /Rect [223.722 409.121 256.757 422.345] /A << /S /GoTo /D <74686D2E432E342E342E33> >> >>}
\MusicPDFWriteObject{6805}{<< /Type /Annot /Subtype /Link /Border [0 0 0] /H /I /C [1 0 0] /Rect [280.301 382.82 313.336 395.152] /A << /S /GoTo /D <74686D2E432E342E342E34> >> >>}
\MusicPDFWriteObject{6806}{<< /Type /Annot /Subtype /Link /Border [0 0 0] /H /I /C [1 0 0] /Rect [146.286 370.583 178.713 381.373] /A << /S /GoTo /D <74686D2E432E342E342E33> >> >>}
\MusicPDFWriteObject{6807}{<< /Type /Annot /Subtype /Link /Border [0 0 0] /H /I /C [0 1 0] /Rect [357.201 216.594 388.03 225.461] /A << /S /GoTo /D <636974652E47726F2D61> >> >>}
\MusicPDFWriteObject{6808}{<< /Type /Annot /Subtype /Link /Border [0 0 0] /H /I /C [0 1 0] /Rect [511.535 216.594 543.159 225.461] /A << /S /GoTo /D <636974652E47726F2D62> >> >>}
\MusicPDFWriteObject{6809}{<< /Type /Annot /Subtype /Link /Border [0 0 0] /H /I /C [0 1 0] /Rect [448.592 163.99 480.411 172.858] /A << /S /GoTo /D <636974652E4B72757365> >> >>}
\MusicPDFWriteObject{6810}{<< /Type /Annot /Subtype /Link /Border [0 0 0] /H /I /C [0 1 0] /Rect [515.415 163.99 543.159 172.858] /A << /S /GoTo /D <636974652E4A6F722D61> >> >>}
\MusicPDFWriteObject{6811}{<< /Type /Annot /Subtype /Link /Border [0 0 0] /H /I /C [0 1 0] /Rect [112.693 150.84 139.627 159.708] /A << /S /GoTo /D <636974652E424A4D> >> >>}
\MusicPDFWriteObject{6812}{<< /Type /Annot /Subtype /Link /Border [0 0 0] /H /I /C [0 1 0] /Rect [491.148 150.84 509.672 159.708] /A << /S /GoTo /D <636974652E4745> >> >>}
\MusicPDFWriteObject{6813}{<< /Type /Annot /Subtype /Link /Border [0 0 0] /H /I /C [0 1 0] /Rect [229.788 250.116 254.848 258.984] /A << /S /GoTo /D <636974652E53636861> >> >>}
\MusicPDFWriteObject{6814}{<< /Type /Annot /Subtype /Link /Border [0 0 0] /H /I /C [0 1 0] /Rect [201.664 184.362 232.057 193.23] /A << /S /GoTo /D <636974652E426F752D49> >> >>}
\MusicPDFWriteObject{6815}{<< /Type /Annot /Subtype /Link /Border [0 0 0] /H /I /C [0 1 0] /Rect [73.77 608.219 104.163 617.087] /A << /S /GoTo /D <636974652E426F752D49> >> >>}
\MusicPDFWriteObject{6816}{<< /Type /Annot /Subtype /Link /Border [0 0 0] /H /I /C [1 0 0] /Rect [98.422 459.873 130.848 471.825] /A << /S /GoTo /D <6571756174696F6E2E432E352E312E32> >> >>}
\MusicPDFWriteObject{6817}{<< /Type /Annot /Subtype /Link /Border [0 0 0] /H /I /C [1 0 0] /Rect [233.637 443.053 266.064 453.844] /A << /S /GoTo /D <74686D2E432E352E312E31> >> >>}
\MusicPDFWriteObject{6818}{<< /Type /Annot /Subtype /Link /Border [0 0 0] /H /I /C [0 1 0] /Rect [390.05 416.596 415.11 425.464] /A << /S /GoTo /D <636974652E53636861> >> >>}
\MusicPDFWriteObject{6819}{<< /Type /Annot /Subtype /Link /Border [0 0 0] /H /I /C [1 0 0] /Rect [511.336 401.508 543.762 412.299] /A << /S /GoTo /D <74686D2E432E352E312E35> >> >>}
\MusicPDFWriteObject{6820}{<< /Type /Annot /Subtype /Link /Border [0 0 0] /H /I /C [1 0 0] /Rect [511.81 96.3 544.845 107.091] /A << /S /GoTo /D <74686D2E432E352E312E33> >> >>}
\MusicPDFWriteObject{6821}{<< /Type /Annot /Subtype /Link /Border [0 0 0] /H /I /C [0 1 0] /Rect [424.513 85.086 467.536 93.954] /A << /S /GoTo /D <636974652E426F752D4754> >> >>}
\MusicPDFWriteObject{6822}{<< /Type /Annot /Subtype /Link /Border [0 0 0] /H /I /C [0 1 0] /Rect [224.808 637.361 255.201 646.229] /A << /S /GoTo /D <636974652E426F752D49> >> >>}
\MusicPDFWriteObject{6823}{<< /Type /Annot /Subtype /Link /Border [0 0 0] /H /I /C [0 1 0] /Rect [377.791 637.361 408.184 646.229] /A << /S /GoTo /D <636974652E426F752D49> >> >>}
\MusicPDFWriteObject{6824}{<< /Type /Annot /Subtype /Link /Border [0 0 0] /H /I /C [0 1 0] /Rect [324.423 624.21 354.816 633.078] /A << /S /GoTo /D <636974652E426F752D49> >> >>}
\MusicPDFWriteObject{6825}{<< /Type /Annot /Subtype /Link /Border [0 0 0] /H /I /C [0 1 0] /Rect [73.77 611.059 104.163 619.927] /A << /S /GoTo /D <636974652E426F752D49> >> >>}
\MusicPDFWriteObject{6826}{<< /Type /Annot /Subtype /Link /Border [0 0 0] /H /I /C [0 1 0] /Rect [280.875 611.059 311.268 619.927] /A << /S /GoTo /D <636974652E426F752D49> >> >>}
\MusicPDFWriteObject{6827}{<< /Type /Annot /Subtype /Link /Border [0 0 0] /H /I /C [1 0 0] /Rect [99.443 451.605 131.869 463.558] /A << /S /GoTo /D <6571756174696F6E2E432E352E312E35> >> >>}
\MusicPDFWriteObject{6828}{<< /Type /Annot /Subtype /Link /Border [0 0 0] /H /I /C [0 1 0] /Rect [133.535 440.944 158.595 449.812] /A << /S /GoTo /D <636974652E53636861> >> >>}
\MusicPDFWriteObject{6829}{<< /Type /Annot /Subtype /Link /Border [0 0 0] /H /I /C [0 1 0] /Rect [399.746 406.161 424.806 415.029] /A << /S /GoTo /D <636974652E53636861> >> >>}
\MusicPDFWriteObject{6830}{<< /Type /Annot /Subtype /Link /Border [0 0 0] /H /I /C [1 0 0] /Rect [361.696 355.736 393.514 367.692] /A << /S /GoTo /D <74686D2E432E352E312E33> >> >>}
\MusicPDFWriteObject{6831}{<< /Type /Annot /Subtype /Link /Border [0 0 0] /H /I /C [0 1 0] /Rect [390.641 345.076 415.701 353.944] /A << /S /GoTo /D <636974652E53636861> >> >>}
\MusicPDFWriteObject{6832}{<< /Type /Annot /Subtype /Link /Border [0 0 0] /H /I /C [1 0 0] /Rect [511.9 639.448 543.809 651.403] /A << /S /GoTo /D <74686D2E432E352E312E31> >> >>}
\MusicPDFWriteObject{6833}{<< /Type /Annot /Subtype /Link /Border [0 0 0] /H /I /C [0 1 0] /Rect [511.42 576.185 543.239 585.053] /A << /S /GoTo /D <636974652E4B72757365> >> >>}
\MusicPDFWriteObject{6834}{<< /Type /Annot /Subtype /Link /Border [0 0 0] /H /I /C [1 0 0] /Rect [164.509 429.037 189.166 440.989] /A << /S /GoTo /D <73756273656374696F6E2E432E342E33> >> >>}
\MusicPDFWriteObject{6835}{<< /Type /Annot /Subtype /Link /Border [0 0 0] /H /I /C [1 0 0] /Rect [354.48 429.037 386.876 440.989] /A << /S /GoTo /D <74686D2E432E342E332E31> >> >>}
\MusicPDFWriteObject{6836}{<< /Type /Annot /Subtype /Link /Border [0 0 0] /H /I /C [1 0 0] /Rect [469.07 429.037 501.467 440.989] /A << /S /GoTo /D <74686D2E432E342E332E32> >> >>}
\MusicPDFWriteObject{6837}{<< /Type /Annot /Subtype /Link /Border [0 0 0] /H /I /C [1 0 0] /Rect [224.114 416.439 256.541 427.23] /A << /S /GoTo /D <74686D2E432E342E332E35> >> >>}
\MusicPDFWriteObject{6838}{<< /Type /Annot /Subtype /Link /Border [0 0 0] /H /I /C [1 0 0] /Rect [192.462 321.123 224.889 333.076] /A << /S /GoTo /D <74686D2E432E352E312E37> >> >>}
\MusicPDFWriteObject{6839}{<< /Type /Annot /Subtype /Link /Border [0 0 0] /H /I /C [1 0 0] /Rect [73.712 127.78 106.138 139.733] /A << /S /GoTo /D <6571756174696F6E2E432E352E312E38> >> >>}
\MusicPDFWriteObject{6840}{<< /Type /Annot /Subtype /Link /Border [0 0 0] /H /I /C [1 0 0] /Rect [304.554 400.869 341.961 412.822] /A << /S /GoTo /D <6571756174696F6E2E432E352E312E3130> >> >>}
\MusicPDFWriteObject{6841}{<< /Type /Annot /Subtype /Link /Border [0 0 0] /H /I /C [1 0 0] /Rect [269.839 229.276 307.245 241.228] /A << /S /GoTo /D <6571756174696F6E2E432E352E312E3133> >> >>}
\MusicPDFWriteObject{6842}{<< /Type /Annot /Subtype /Link /Border [0 0 0] /H /I /C [1 0 0] /Rect [262.81 201.605 294.628 213.992] /A << /S /GoTo /D <74686D2E432E352E312E38> >> >>}
\MusicPDFWriteObject{6843}{<< /Type /Annot /Subtype /Link /Border [0 0 0] /H /I /C [1 0 0] /Rect [135.329 161.599 167.785 173.552] /A << /S /GoTo /D <74686D2E432E352E312E38> >> >>}
\MusicPDFWriteObject{6844}{<< /Type /Annot /Subtype /Link /Border [0 0 0] /H /I /C [0 1 0] /Rect [238.698 164.089 270.518 172.957] /A << /S /GoTo /D <636974652E4B72757365> >> >>}
\MusicPDFWriteObject{6845}{<< /Type /Annot /Subtype /Link /Border [0 0 0] /H /I /C [1 0 0] /Rect [248.19 148.448 281.225 160.421] /A << /S /GoTo /D <74686D2E432E352E312E38> >> >>}
\MusicPDFWriteObject{6846}{<< /Type /Annot /Subtype /Link /Border [0 0 0] /H /I /C [0 1 0] /Rect [223.366 124.637 254.195 133.505] /A << /S /GoTo /D <636974652E47726F2D61> >> >>}
\MusicPDFWriteObject{6847}{<< /Type /Annot /Subtype /Link /Border [0 0 0] /H /I /C [1 0 0] /Rect [397.704 122.147 430.131 134.1] /A << /S /GoTo /D <74686D2E432E342E312E32> >> >>}
\MusicPDFWriteObject{6848}{<< /Type /Annot /Subtype /Link /Border [0 0 0] /H /I /C [1 0 0] /Rect [488.342 122.147 520.769 134.1] /A << /S /GoTo /D <74686D2E432E342E312E33> >> >>}
\MusicPDFWriteObject{6849}{<< /Type /Annot /Subtype /Link /Border [0 0 0] /H /I /C [0 1 0] /Rect [168.046 111.437 199.865 120.305] /A << /S /GoTo /D <636974652E4B72757365> >> >>}
\MusicPDFWriteObject{6850}{<< /Type /Annot /Subtype /Link /Border [0 0 0] /H /I /C [1 0 0] /Rect [75.432 632.583 107.858 643.431] /A << /S /GoTo /D <74686D2E432E352E322E31> >> >>}
\MusicPDFWriteObject{6851}{<< /Type /Annot /Subtype /Link /Border [0 0 0] /H /I /C [1 0 0] /Rect [128.967 600.161 161.393 609.015] /A << /S /GoTo /D <74686D2E432E352E312E38> >> >>}
\MusicPDFWriteObject{6852}{<< /Type /Annot /Subtype /Link /Border [0 0 0] /H /I /C [0 1 0] /Rect [299.374 422.72 331.193 431.588] /A << /S /GoTo /D <636974652E4B72757365> >> >>}
\MusicPDFWriteObject{6853}{<< /Type /Annot /Subtype /Link /Border [0 0 0] /H /I /C [1 0 0] /Rect [501.879 164.829 539.285 176.781] /A << /S /GoTo /D <6571756174696F6E2E432E352E312E3137> >> >>}
\MusicPDFWriteObject{6854}{<< /Type /Annot /Subtype /Link /Border [0 0 0] /H /I /C [1 0 0] /Rect [309.424 95.747 346.831 107.699] /A << /S /GoTo /D <6571756174696F6E2E432E352E312E3137> >> >>}
\MusicPDFWriteObject{6855}{<< /Type /Annot /Subtype /Link /Border [0 0 0] /H /I /C [1 0 0] /Rect [121.911 648.206 159.318 660.158] /A << /S /GoTo /D <6571756174696F6E2E432E352E312E3134> >> >>}
\MusicPDFWriteObject{6856}{<< /Type /Annot /Subtype /Link /Border [0 0 0] /H /I /C [1 0 0] /Rect [463.189 635.055 500.596 647.008] /A << /S /GoTo /D <6571756174696F6E2E432E352E312E3137> >> >>}
\MusicPDFWriteObject{6857}{<< /Type /Annot /Subtype /Link /Border [0 0 0] /H /I /C [1 0 0] /Rect [118.591 621.904 151.018 633.857] /A << /S /GoTo /D <6571756174696F6E2E432E352E312E39> >> >>}
\MusicPDFWriteObject{6858}{<< /Type /Annot /Subtype /Link /Border [0 0 0] /H /I /C [1 0 0] /Rect [271.315 621.904 308.721 633.857] /A << /S /GoTo /D <6571756174696F6E2E432E352E312E3130> >> >>}
\MusicPDFWriteObject{6859}{<< /Type /Annot /Subtype /Link /Border [0 0 0] /H /I /C [1 0 0] /Rect [378.808 571.402 415.719 584.342] /A << /S /GoTo /D <74686D2E432E352E312E3133> >> >>}
\MusicPDFWriteObject{6860}{<< /Type /Annot /Subtype /Link /Border [0 0 0] /H /I /C [0 1 0] /Rect [511.434 573.892 543.253 582.759] /A << /S /GoTo /D <636974652E4B72757365> >> >>}
\MusicPDFWriteObject{6861}{<< /Type /Annot /Subtype /Link /Border [0 0 0] /H /I /C [0 1 0] /Rect [378.041 560.741 408.87 569.609] /A << /S /GoTo /D <636974652E47726F2D61> >> >>}
\MusicPDFWriteObject{6862}{<< /Type /Annot /Subtype /Link /Border [0 0 0] /H /I /C [1 0 0] /Rect [262.638 249.334 295.673 261.286] /A << /S /GoTo /D <74686D2E432E342E332E35> >> >>}
\MusicPDFWriteObject{6863}{<< /Type /Annot /Subtype /Link /Border [0 0 0] /H /I /C [1 0 0] /Rect [210.852 95.747 248.259 107.699] /A << /S /GoTo /D <6571756174696F6E2E432E352E312E3231> >> >>}
\MusicPDFWriteObject{6864}{<< /Type /Annot /Subtype /Link /Border [0 0 0] /H /I /C [0 1 0] /Rect [412.815 687.123 437.875 695.991] /A << /S /GoTo /D <636974652E53636861> >> >>}
\MusicPDFWriteObject{6865}{<< /Type /Annot /Subtype /Link /Border [0 0 0] /H /I /C [1 0 0] /Rect [146.65 658.885 179.077 669.676] /A << /S /GoTo /D <74686D2E432E342E332E31> >> >>}
\MusicPDFWriteObject{6866}{<< /Type /Annot /Subtype /Link /Border [0 0 0] /H /I /C [0 1 0] /Rect [320.355 545.233 345.415 554.101] /A << /S /GoTo /D <636974652E53636861> >> >>}
\MusicPDFWriteObject{6867}{<< /Type /Annot /Subtype /Link /Border [0 0 0] /H /I /C [1 0 0] /Rect [422.996 399.132 460.402 411.085] /A << /S /GoTo /D <6571756174696F6E2E432E352E312E3230> >> >>}
\MusicPDFWriteObject{6868}{<< /Type /Annot /Subtype /Link /Border [0 0 0] /H /I /C [1 0 0] /Rect [292.207 385.981 328.905 397.954] /A << /S /GoTo /D <74686D2E432E352E312E3132> >> >>}
\MusicPDFWriteObject{6869}{<< /Type /Annot /Subtype /Link /Border [0 0 0] /H /I /C [1 0 0] /Rect [146.728 307.076 179.215 319.049] /A << /S /GoTo /D <74686D2E432E342E332E35> >> >>}
\MusicPDFWriteObject{6870}{<< /Type /Annot /Subtype /Link /Border [0 0 0] /H /I /C [1 0 0] /Rect [166.201 233.644 203.608 245.597] /A << /S /GoTo /D <6571756174696F6E2E432E352E312E3236> >> >>}
\MusicPDFWriteObject{6871}{<< /Type /Annot /Subtype /Link /Border [0 0 0] /H /I /C [0 1 0] /Rect [346.155 189.023 371.215 197.891] /A << /S /GoTo /D <636974652E53636861> >> >>}
\MusicPDFWriteObject{6872}{<< /Type /Annot /Subtype /Link /Border [0 0 0] /H /I /C [1 0 0] /Rect [194.772 160.231 232.179 172.184] /A << /S /GoTo /D <6571756174696F6E2E432E352E312E3237> >> >>}
\MusicPDFWriteObject{6873}{<< /Type /Annot /Subtype /Link /Border [0 0 0] /H /I /C [1 0 0] /Rect [412.311 147.081 449.717 160.021] /A << /S /GoTo /D <6571756174696F6E2E432E352E312E3237> >> >>}
\MusicPDFWriteObject{6874}{<< /Type /Annot /Subtype /Link /Border [0 0 0] /H /I /C [0 1 0] /Rect [511.414 673.973 543.234 682.84] /A << /S /GoTo /D <636974652E4B72757365> >> >>}
\MusicPDFWriteObject{6875}{<< /Type /Annot /Subtype /Link /Border [0 0 0] /H /I /C [1 0 0] /Rect [176.128 485.795 213.535 497.748] /A << /S /GoTo /D <6571756174696F6E2E432E352E312E3238> >> >>}
\MusicPDFWriteObject{6876}{<< /Type /Annot /Subtype /Link /Border [0 0 0] /H /I /C [1 0 0] /Rect [167.464 335.503 204.871 347.455] /A << /S /GoTo /D <6571756174696F6E2E432E352E312E3130> >> >>}
\MusicPDFWriteObject{6877}{<< /Type /Annot /Subtype /Link /Border [0 0 0] /H /I /C [1 0 0] /Rect [281.512 335.503 314.547 348.444] /A << /S /GoTo /D <74686D2E432E352E312E38> >> >>}
\MusicPDFWriteObject{6878}{<< /Type /Annot /Subtype /Link /Border [0 0 0] /H /I /C [1 0 0] /Rect [263.629 322.352 301.036 334.305] /A << /S /GoTo /D <6571756174696F6E2E432E352E312E3130> >> >>}
\MusicPDFWriteObject{6879}{<< /Type /Annot /Subtype /Link /Border [0 0 0] /H /I /C [1 0 0] /Rect [378.431 296.05 415.837 308.003] /A << /S /GoTo /D <6571756174696F6E2E432E352E312E3332> >> >>}
\MusicPDFWriteObject{6880}{<< /Type /Annot /Subtype /Link /Border [0 0 0] /H /I /C [1 0 0] /Rect [443.826 95.747 481.232 107.699] /A << /S /GoTo /D <6571756174696F6E2E432E352E312E3331> >> >>}
\MusicPDFWriteObject{6881}{<< /Type /Annot /Subtype /Link /Border [0 0 0] /H /I /C [1 0 0] /Rect [73.712 564.773 111.118 576.725] /A << /S /GoTo /D <6571756174696F6E2E432E352E312E3333> >> >>}
\MusicPDFWriteObject{6882}{<< /Type /Annot /Subtype /Link /Border [0 0 0] /H /I /C [1 0 0] /Rect [308.014 564.772 344.925 576.745] /A << /S /GoTo /D <6571756174696F6E2E432E352E312E3239> >> >>}
\MusicPDFWriteObject{6883}{<< /Type /Annot /Subtype /Link /Border [0 0 0] /H /I /C [1 0 0] /Rect [261.299 143.381 298.706 155.334] /A << /S /GoTo /D <6571756174696F6E2E432E352E312E3335> >> >>}
\MusicPDFWriteObject{6884}{<< /Type /Annot /Subtype /Link /Border [0 0 0] /H /I /C [1 0 0] /Rect [96.458 117.08 133.864 129.033] /A << /S /GoTo /D <6571756174696F6E2E432E352E312E3336> >> >>}
\MusicPDFWriteObject{6885}{<< /Type /Annot /Subtype /Link /Border [0 0 0] /H /I /C [1 0 0] /Rect [401.103 543.836 438.262 554.626] /A << /S /GoTo /D <74686D2E432E352E312E3139> >> >>}
\MusicPDFWriteObject{6886}{<< /Type /Annot /Subtype /Link /Border [0 0 0] /H /I /C [1 0 0] /Rect [248.793 530.685 281.22 541.476] /A << /S /GoTo /D <74686D2E432E342E322E31> >> >>}
\MusicPDFWriteObject{6887}{<< /Type /Annot /Subtype /Link /Border [0 0 0] /H /I /C [1 0 0] /Rect [324.223 446.905 361.629 458.857] /A << /S /GoTo /D <6571756174696F6E2E432E352E312E3134> >> >>}
\MusicPDFWriteObject{6888}{<< /Type /Annot /Subtype /Link /Border [0 0 0] /H /I /C [1 0 0] /Rect [424.872 446.904 461.57 458.859] /A << /S /GoTo /D <74686D2E432E352E312E3132> >> >>}
\MusicPDFWriteObject{6889}{<< /Type /Annot /Subtype /Link /Border [0 0 0] /H /I /C [1 0 0] /Rect [213.03 378.619 244.909 390.592] /A << /S /GoTo /D <74686D2E432E342E332E35> >> >>}
\MusicPDFWriteObject{6890}{<< /Type /Annot /Subtype /Link /Border [0 0 0] /H /I /C [1 0 0] /Rect [457.263 237.554 494.67 249.507] /A << /S /GoTo /D <6571756174696F6E2E432E352E312E3337> >> >>}
\MusicPDFWriteObject{6891}{<< /Type /Annot /Subtype /Link /Border [0 0 0] /H /I /C [1 0 0] /Rect [108.169 192.678 146.284 203.469] /A << /S /GoTo /D <74686D2E432E352E312E3133> >> >>}
\MusicPDFWriteObject{6892}{<< /Type /Annot /Subtype /Link /Border [0 0 0] /H /I /C [1 0 0] /Rect [432.393 166.376 469.8 177.167] /A << /S /GoTo /D <74686D2E432E352E312E3234> >> >>}
\MusicPDFWriteObject{6893}{<< /Type /Annot /Subtype /Link /Border [0 0 0] /H /I /C [1 0 0] /Rect [327.737 153.226 365.852 164.016] /A << /S /GoTo /D <74686D2E432E352E312E3132> >> >>}
\MusicPDFWriteObject{6894}{<< /Type /Annot /Subtype /Link /Border [0 0 0] /H /I /C [1 0 0] /Rect [407.071 153.226 445.186 164.016] /A << /S /GoTo /D <74686D2E432E352E312E3133> >> >>}
\MusicPDFWriteObject{6895}{<< /Type /Annot /Subtype /Link /Border [0 0 0] /H /I /C [1 0 0] /Rect [486.406 153.226 524.521 164.016] /A << /S /GoTo /D <74686D2E432E352E312E3135> >> >>}
\MusicPDFWriteObject{6896}{<< /Type /Annot /Subtype /Link /Border [0 0 0] /H /I /C [1 0 0] /Rect [106.556 142.012 143.963 150.866] /A << /S /GoTo /D <74686D2E432E352E312E3136> >> >>}
\MusicPDFWriteObject{6897}{<< /Type /Annot /Subtype /Link /Border [0 0 0] /H /I /C [1 0 0] /Rect [353.565 536.755 390.263 548.708] /A << /S /GoTo /D <74686D2E432E352E312E3132> >> >>}
\MusicPDFWriteObject{6898}{<< /Type /Annot /Subtype /Link /Border [0 0 0] /H /I /C [1 0 0] /Rect [403.282 497.303 440.228 509.255] /A << /S /GoTo /D <74686D2E432E352E312E3135> >> >>}
\MusicPDFWriteObject{6899}{<< /Type /Annot /Subtype /Link /Border [0 0 0] /H /I /C [1 0 0] /Rect [481.156 497.303 518.562 509.255] /A << /S /GoTo /D <6571756174696F6E2E432E352E312E3134> >> >>}
\MusicPDFWriteObject{6900}{<< /Type /Annot /Subtype /Link /Border [0 0 0] /H /I /C [1 0 0] /Rect [120.39 370.999 157.796 382.952] /A << /S /GoTo /D <74686D2E432E352E312E3133> >> >>}
\MusicPDFWriteObject{6901}{<< /Type /Annot /Subtype /Link /Border [0 0 0] /H /I /C [1 0 0] /Rect [208.095 357.849 244.794 369.801] /A << /S /GoTo /D <74686D2E432E352E312E3136> >> >>}
\MusicPDFWriteObject{6902}{<< /Type /Annot /Subtype /Link /Border [0 0 0] /H /I /C [1 0 0] /Rect [429.894 357.849 466.592 369.801] /A << /S /GoTo /D <74686D2E432E352E312E3133> >> >>}
\MusicPDFWriteObject{6903}{<< /Type /Annot /Subtype /Link /Border [0 0 0] /H /I /C [0 1 0] /Rect [192.022 320.985 222.851 329.853] /A << /S /GoTo /D <636974652E47726F2D61> >> >>}
\MusicPDFWriteObject{6904}{<< /Type /Annot /Subtype /Link /Border [0 0 0] /H /I /C [0 1 0] /Rect [126.524 307.835 156.197 316.702] /A << /S /GoTo /D <636974652E5363682D62> >> >>}
\MusicPDFWriteObject{6905}{<< /Type /Annot /Subtype /Link /Border [0 0 0] /H /I /C [0 1 0] /Rect [262.572 281.533 288.683 290.401] /A << /S /GoTo /D <636974652E47612D32> >> >>}
\MusicPDFWriteObject{6906}{<< /Type /Annot /Subtype /Link /Border [0 0 0] /H /I /C [1 0 0] /Rect [177.234 95.746 213.932 107.701] /A << /S /GoTo /D <74686D2E432E352E312E3133> >> >>}
\MusicPDFWriteObject{6907}{<< /Type /Annot /Subtype /Link /Border [0 0 0] /H /I /C [0 1 0] /Rect [484.761 443.149 509.822 452.017] /A << /S /GoTo /D <636974652E53636861> >> >>}
\MusicPDFWriteObject{6908}{<< /Type /Annot /Subtype /Link /Border [0 0 0] /H /I /C [1 0 0] /Rect [240.22 401.207 277.626 413.159] /A << /S /GoTo /D <6571756174696F6E2E432E352E312E3431> >> >>}
\MusicPDFWriteObject{6909}{<< /Type /Annot /Subtype /Link /Border [0 0 0] /H /I /C [1 0 0] /Rect [482.276 364.779 519.683 376.732] /A << /S /GoTo /D <6571756174696F6E2E432E352E312E3432> >> >>}
\MusicPDFWriteObject{6910}{<< /Type /Annot /Subtype /Link /Border [0 0 0] /H /I /C [1 0 0] /Rect [220.808 325.327 258.215 337.279] /A << /S /GoTo /D <6571756174696F6E2E432E352E312E3432> >> >>}
\MusicPDFWriteObject{6911}{<< /Type /Annot /Subtype /Link /Border [0 0 0] /H /I /C [0 1 0] /Rect [378.682 614.22 408.355 623.088] /A << /S /GoTo /D <636974652E5363682D62> >> >>}
\MusicPDFWriteObject{6912}{<< /Type /Annot /Subtype /Link /Border [0 0 0] /H /I /C [1 0 0] /Rect [244.554 535.862 281.253 547.815] /A << /S /GoTo /D <74686D2E432E352E312E3133> >> >>}
\MusicPDFWriteObject{6913}{<< /Type /Annot /Subtype /Link /Border [0 0 0] /H /I /C [1 0 0] /Rect [475.57 384.118 512.268 396.711] /A << /S /GoTo /D <74686D2E432E352E312E3232> >> >>}
\MusicPDFWriteObject{6914}{<< /Type /Annot /Subtype /Link /Border [0 0 0] /H /I /C [1 0 0] /Rect [296.871 357.817 333.888 369.772] /A << /S /GoTo /D <74686D2E432E352E312E3136> >> >>}
\MusicPDFWriteObject{6915}{<< /Type /Annot /Subtype /Link /Border [0 0 0] /H /I /C [1 0 0] /Rect [374.128 357.817 411.145 369.772] /A << /S /GoTo /D <6571756174696F6E2E432E352E312E3332> >> >>}
\MusicPDFWriteObject{6916}{<< /Type /Annot /Subtype /Link /Border [0 0 0] /H /I /C [1 0 0] /Rect [315.109 331.155 351.807 343.47] /A << /S /GoTo /D <74686D2E432E352E312E3133> >> >>}
\MusicPDFWriteObject{6917}{<< /Type /Annot /Subtype /Link /Border [0 0 0] /H /I /C [1 0 0] /Rect [71.004 318.365 108.41 330.957] /A << /S /GoTo /D <74686D2E432E352E312E3136> >> >>}
\MusicPDFWriteObject{6918}{<< /Type /Annot /Subtype /Link /Border [0 0 0] /H /I /C [1 0 0] /Rect [358.401 108.897 390.827 120.85] /A << /S /GoTo /D <74686D2E432E352E322E31> >> >>}
\MusicPDFWriteObject{6919}{<< /Type /Annot /Subtype /Link /Border [0 0 0] /H /I /C [0 1 0] /Rect [376.612 98.237 395.585 107.105] /A << /S /GoTo /D <636974652E547265> >> >>}
\MusicPDFWriteObject{6920}{<< /Type /Annot /Subtype /Link /Border [0 0 0] /H /I /C [1 0 0] /Rect [130.54 82.596 168.655 94.549] /A << /S /GoTo /D <74686D2E432E352E312E3233> >> >>}
\MusicPDFWriteObject{6921}{<< /Type /Annot /Subtype /Link /Border [0 0 0] /H /I /C [0 1 0] /Rect [307.616 85.086 337.289 93.954] /A << /S /GoTo /D <636974652E5363682D62> >> >>}
\MusicPDFWriteObject{6922}{<< /Type /Annot /Subtype /Link /Border [0 0 0] /H /I /C [1 0 0] /Rect [492.063 82.596 530.177 94.549] /A << /S /GoTo /D <74686D2E432E352E312E3235> >> >>}
\MusicPDFWriteObject{6923}{<< /Type /Annot /Subtype /Link /Border [0 0 0] /H /I /C [1 0 0] /Rect [376.2 685.187 414.315 695.977] /A << /S /GoTo /D <74686D2E432E352E312E3235> >> >>}
\MusicPDFWriteObject{6924}{<< /Type /Annot /Subtype /Link /Border [0 0 0] /H /I /C [1 0 0] /Rect [365.578 671.482 402.276 683.435] /A << /S /GoTo /D <74686D2E432E352E312E3136> >> >>}
\MusicPDFWriteObject{6925}{<< /Type /Annot /Subtype /Link /Border [0 0 0] /H /I /C [1 0 0] /Rect [336.755 658.332 374.303 670.284] /A << /S /GoTo /D <74686D2E432E352E312E3235> >> >>}
\MusicPDFWriteObject{6926}{<< /Type /Annot /Subtype /Link /Border [0 0 0] /H /I /C [1 0 0] /Rect [338.491 645.734 375.897 656.525] /A << /S /GoTo /D <74686D2E432E352E312E3233> >> >>}
\MusicPDFWriteObject{6927}{<< /Type /Annot /Subtype /Link /Border [0 0 0] /H /I /C [0 1 0] /Rect [400.531 383.752 425.591 392.62] /A << /S /GoTo /D <636974652E53636861> >> >>}
\MusicPDFWriteObject{6928}{<< /Type /Annot /Subtype /Link /Border [0 0 0] /H /I /C [1 0 0] /Rect [221.406 305.9 258.635 316.69] /A << /S /GoTo /D <74686D2E432E352E312E3236> >> >>}
\MusicPDFWriteObject{6929}{<< /Type /Annot /Subtype /Link /Border [0 0 0] /H /I /C [0 1 0] /Rect [369.32 241.673 398.993 250.541] /A << /S /GoTo /D <636974652E5363682D62> >> >>}
\MusicPDFWriteObject{6930}{<< /Type /Annot /Subtype /Link /Border [0 0 0] /H /I /C [0 1 0] /Rect [164.031 228.522 193.705 237.39] /A << /S /GoTo /D <636974652E5363682D62> >> >>}
\MusicPDFWriteObject{6931}{<< /Type /Annot /Subtype /Link /Border [0 0 0] /H /I /C [1 0 0] /Rect [490.796 226.032 528.026 238.973] /A << /S /GoTo /D <74686D2E432E352E312E3236> >> >>}
\MusicPDFWriteObject{6932}{<< /Type /Annot /Subtype /Link /Border [0 0 0] /H /I /C [1 0 0] /Rect [201.206 212.882 238.612 224.834] /A << /S /GoTo /D <74686D2E432E352E312E3235> >> >>}
\MusicPDFWriteObject{6933}{<< /Type /Annot /Subtype /Link /Border [0 0 0] /H /I /C [0 1 0] /Rect [264.84 215.372 294.514 224.239] /A << /S /GoTo /D <636974652E5363682D62> >> >>}
\MusicPDFWriteObject{6934}{<< /Type /Annot /Subtype /Link /Border [0 0 0] /H /I /C [1 0 0] /Rect [71.004 177.004 104.039 189.945] /A << /S /GoTo /D <74686D2E432E322E322E32> >> >>}
\MusicPDFWriteObject{6935}{<< /Type /Annot /Subtype /Link /Border [0 0 0] /H /I /C [1 0 0] /Rect [217.299 651.118 253.997 661.909] /A << /S /GoTo /D <74686D2E432E352E312E3236> >> >>}
\MusicPDFWriteObject{6936}{<< /Type /Annot /Subtype /Link /Border [0 0 0] /H /I /C [0 1 0] /Rect [369.32 586.892 398.993 595.76] /A << /S /GoTo /D <636974652E5363682D62> >> >>}
\MusicPDFWriteObject{6937}{<< /Type /Annot /Subtype /Link /Border [0 0 0] /H /I /C [0 1 0] /Rect [165.067 573.741 194.741 582.609] /A << /S /GoTo /D <636974652E5363682D62> >> >>}
\MusicPDFWriteObject{6938}{<< /Type /Annot /Subtype /Link /Border [0 0 0] /H /I /C [1 0 0] /Rect [489.887 571.25 527.789 584.192] /A << /S /GoTo /D <74686D2E432E352E312E3236> >> >>}
\MusicPDFWriteObject{6939}{<< /Type /Annot /Subtype /Link /Border [0 0 0] /H /I /C [1 0 0] /Rect [201.206 558.1 238.612 570.053] /A << /S /GoTo /D <74686D2E432E352E312E3235> >> >>}
\MusicPDFWriteObject{6940}{<< /Type /Annot /Subtype /Link /Border [0 0 0] /H /I /C [0 1 0] /Rect [264.84 560.59 294.514 569.458] /A << /S /GoTo /D <636974652E5363682D62> >> >>}
\MusicPDFWriteObject{6941}{<< /Type /Annot /Subtype /Link /Border [0 0 0] /H /I /C [0 1 0] /Rect [73.77 526.522 103.444 535.39] /A << /S /GoTo /D <636974652E5363682D62> >> >>}
\MusicPDFWriteObject{6942}{<< /Type /Annot /Subtype /Link /Border [0 0 0] /H /I /C [1 0 0] /Rect [356.189 306.158 393.879 319.1] /A << /S /GoTo /D <74686D2E432E352E312E3238> >> >>}
\MusicPDFWriteObject{6943}{<< /Type /Annot /Subtype /Link /Border [0 0 0] /H /I /C [1 0 0] /Rect [149.282 254.109 187.149 266.497] /A << /S /GoTo /D <74686D2E432E352E312E3235> >> >>}
\MusicPDFWriteObject{6944}{<< /Type /Annot /Subtype /Link /Border [0 0 0] /H /I /C [0 1 0] /Rect [308.475 229.744 333.535 238.612] /A << /S /GoTo /D <636974652E53636861> >> >>}
\MusicPDFWriteObject{6945}{<< /Type /Annot /Subtype /Link /Border [0 0 0] /H /I /C [1 0 0] /Rect [73.712 95.747 111.118 107.699] /A << /S /GoTo /D <6571756174696F6E2E432E352E312E3435> >> >>}
\MusicPDFWriteObject{6946}{<< /Type /Annot /Subtype /Link /Border [0 0 0] /H /I /C [1 0 0] /Rect [141.439 95.747 178.845 107.699] /A << /S /GoTo /D <6571756174696F6E2E432E352E312E3436> >> >>}
\MusicPDFWriteObject{6947}{<< /Type /Annot /Subtype /Link /Border [0 0 0] /H /I /C [1 0 0] /Rect [185.617 637.883 223.732 649.838] /A << /S /GoTo /D <6571756174696F6E2E432E352E312E3437> >> >>}
\MusicPDFWriteObject{6948}{<< /Type /Annot /Subtype /Link /Border [0 0 0] /H /I /C [1 0 0] /Rect [380.384 624.733 417.79 636.686] /A << /S /GoTo /D <6571756174696F6E2E432E352E312E3438> >> >>}
\MusicPDFWriteObject{6949}{<< /Type /Annot /Subtype /Link /Border [0 0 0] /H /I /C [1 0 0] /Rect [505.268 598.431 541.966 610.384] /A << /S /GoTo /D <6571756174696F6E2E432E352E312E3437> >> >>}
\MusicPDFWriteObject{6950}{<< /Type /Annot /Subtype /Link /Border [0 0 0] /H /I /C [1 0 0] /Rect [207.001 566.152 244.549 578.105] /A << /S /GoTo /D <74686D2E432E352E312E3330> >> >>}
\MusicPDFWriteObject{6951}{<< /Type /Annot /Subtype /Link /Border [0 0 0] /H /I /C [1 0 0] /Rect [316.121 566.152 325.752 578.105] /A << /S /GoTo /D <617070656E6469782E44> >> >>}
\MusicPDFWriteObject{6952}{<< /Type /Annot /Subtype /Link /Border [0 0 0] /H /I /C [1 0 0] /Rect [437.5 566.152 470.465 578.105] /A << /S /GoTo /D <74686D2E442E312E332E37> >> >>}
\MusicPDFWriteObject{6953}{<< /Type /Annot /Subtype /Link /Border [0 0 0] /H /I /C [1 0 0] /Rect [107.44 553.002 140.898 564.954] /A << /S /GoTo /D <74686D2E442E322E322E32> >> >>}
\MusicPDFWriteObject{6954}{<< /Type /Annot /Subtype /Link /Border [0 0 0] /H /I /C [0 1 0] /Rect [256.889 529.19 275.862 538.058] /A << /S /GoTo /D <636974652E547265> >> >>}
\MusicPDFWriteObject{6955}{<< /Type /Annot /Subtype /Link /Border [0 0 0] /H /I /C [1 0 0] /Rect [71.004 513.549 107.702 525.502] /A << /S /GoTo /D <74686D2E432E352E312E3330> >> >>}
\MusicPDFWriteObject{6956}{<< /Type /Annot /Subtype /Link /Border [0 0 0] /H /I /C [1 0 0] /Rect [288.35 513.549 320.575 525.502] /A << /S /GoTo /D <74686D2E442E312E332E37> >> >>}
\MusicPDFWriteObject{6957}{<< /Type /Annot /Subtype /Link /Border [0 0 0] /H /I /C [1 0 0] /Rect [375.704 513.549 407.929 525.502] /A << /S /GoTo /D <74686D2E442E322E322E32> >> >>}
\MusicPDFWriteObject{6958}{<< /Type /Annot /Subtype /Link /Border [0 0 0] /H /I /C [1 0 0] /Rect [477.57 513.549 487.019 525.502] /A << /S /GoTo /D <617070656E6469782E44> >> >>}
\MusicPDFWriteObject{6959}{<< /Type /Annot /Subtype /Link /Border [0 0 0] /H /I /C [0 1 0] /Rect [173.037 502.889 192.011 511.756] /A << /S /GoTo /D <636974652E547265> >> >>}
\MusicPDFWriteObject{6960}{<< /Type /Annot /Subtype /Link /Border [0 0 0] /H /I /C [1 0 0] /Rect [189.112 389.215 221.539 401.168] /A << /S /GoTo /D <6571756174696F6E2E432E342E312E31> >> >>}
\MusicPDFWriteObject{6961}{<< /Type /Annot /Subtype /Link /Border [0 0 0] /H /I /C [1 0 0] /Rect [328.486 389.214 361.004 401.187] /A << /S /GoTo /D <74686D2E432E342E312E31> >> >>}
\MusicPDFWriteObject{6962}{<< /Type /Annot /Subtype /Link /Border [0 0 0] /H /I /C [1 0 0] /Rect [71.004 349.763 103.43 361.715] /A << /S /GoTo /D <74686D2E432E342E322E31> >> >>}
\MusicPDFWriteObject{6963}{<< /Type /Annot /Subtype /Link /Border [0 0 0] /H /I /C [1 0 0] /Rect [349.853 153.59 381.671 164.381] /A << /S /GoTo /D <74686D2E432E352E312E38> >> >>}
\MusicPDFWriteObject{6964}{<< /Type /Annot /Subtype /Link /Border [0 0 0] /H /I /C [0 1 0] /Rect [105.28 124.538 136.109 133.406] /A << /S /GoTo /D <636974652E47726F2D61> >> >>}
\MusicPDFWriteObject{6965}{<< /Type /Annot /Subtype /Link /Border [0 0 0] /H /I /C [1 0 0] /Rect [118.559 108.897 151.594 120.85] /A << /S /GoTo /D <74686D2E432E352E322E31> >> >>}
\MusicPDFWriteObject{6966}{<< /Type /Annot /Subtype /Link /Border [0 0 0] /H /I /C [0 1 0] /Rect [186.632 687.123 217.46 695.991] /A << /S /GoTo /D <636974652E47726F2D61> >> >>}
\MusicPDFWriteObject{6967}{<< /Type /Annot /Subtype /Link /Border [0 0 0] /H /I /C [0 1 0] /Rect [273.526 651.987 304.355 660.854] /A << /S /GoTo /D <636974652E47726F2D61> >> >>}
\MusicPDFWriteObject{6968}{<< /Type /Annot /Subtype /Link /Border [0 0 0] /H /I /C [1 0 0] /Rect [366.36 544.716 399.335 556.669] /A << /S /GoTo /D <74686D2E432E352E322E31> >> >>}
\MusicPDFWriteObject{6969}{<< /Type /Annot /Subtype /Link /Border [0 0 0] /H /I /C [1 0 0] /Rect [232.561 394.261 264.379 406.214] /A << /S /GoTo /D <6571756174696F6E2E432E352E322E32> >> >>}
\MusicPDFWriteObject{6970}{<< /Type /Annot /Subtype /Link /Border [0 0 0] /H /I /C [1 0 0] /Rect [292.316 394.261 324.134 406.214] /A << /S /GoTo /D <6571756174696F6E2E432E352E322E33> >> >>}
\MusicPDFWriteObject{6971}{<< /Type /Annot /Subtype /Link /Border [0 0 0] /H /I /C [1 0 0] /Rect [229.377 83.149 261.803 93.94] /A << /S /GoTo /D <74686D2E432E352E322E36> >> >>}
\MusicPDFWriteObject{6972}{<< /Type /Annot /Subtype /Link /Border [0 0 0] /H /I /C [0 1 0] /Rect [172.056 687.123 202.884 695.991] /A << /S /GoTo /D <636974652E47726F2D61> >> >>}
\MusicPDFWriteObject{6973}{<< /Type /Annot /Subtype /Link /Border [0 0 0] /H /I /C [1 0 0] /Rect [275.972 612.835 307.79 624.788] /A << /S /GoTo /D <74686D2E432E352E322E36> >> >>}
\MusicPDFWriteObject{6974}{<< /Type /Annot /Subtype /Link /Border [0 0 0] /H /I /C [0 1 0] /Rect [333.916 615.325 364.745 624.193] /A << /S /GoTo /D <636974652E47726F2D61> >> >>}
\MusicPDFWriteObject{6975}{<< /Type /Annot /Subtype /Link /Border [0 0 0] /H /I /C [1 0 0] /Rect [511.81 586.533 544.845 598.486] /A << /S /GoTo /D <74686D2E432E342E332E35> >> >>}
\MusicPDFWriteObject{6976}{<< /Type /Annot /Subtype /Link /Border [0 0 0] /H /I /C [1 0 0] /Rect [315.868 560.232 348.295 572.185] /A << /S /GoTo /D <6571756174696F6E2E432E352E322E34> >> >>}
\MusicPDFWriteObject{6977}{<< /Type /Annot /Subtype /Link /Border [0 0 0] /H /I /C [1 0 0] /Rect [124.264 520.779 156.691 532.732] /A << /S /GoTo /D <74686D2E432E342E332E31> >> >>}
\MusicPDFWriteObject{6978}{<< /Type /Annot /Subtype /Link /Border [0 0 0] /H /I /C [0 1 0] /Rect [377.807 523.27 408.636 532.137] /A << /S /GoTo /D <636974652E47726F2D61> >> >>}
\MusicPDFWriteObject{6979}{<< /Type /Annot /Subtype /Link /Border [0 0 0] /H /I /C [0 1 0] /Rect [237.565 510.119 268.394 518.987] /A << /S /GoTo /D <636974652E47726F2D61> >> >>}
\MusicPDFWriteObject{6980}{<< /Type /Annot /Subtype /Link /Border [0 0 0] /H /I /C [1 0 0] /Rect [71.004 458.186 104.039 470.139] /A << /S /GoTo /D <74686D2E432E352E322E34> >> >>}
\MusicPDFWriteObject{6981}{<< /Type /Annot /Subtype /Link /Border [0 0 0] /H /I /C [0 1 0] /Rect [334.218 460.676 378.03 469.544] /A << /S /GoTo /D <636974652E44756E666F7264> >> >>}
\MusicPDFWriteObject{6982}{<< /Type /Annot /Subtype /Link /Border [0 0 0] /H /I /C [1 0 0] /Rect [136.438 186.29 168.864 198.242] /A << /S /GoTo /D <6571756174696F6E2E432E352E332E31> >> >>}
\MusicPDFWriteObject{6983}{<< /Type /Annot /Subtype /Link /Border [0 0 0] /H /I /C [1 0 0] /Rect [325.851 132.592 358.278 143.383] /A << /S /GoTo /D <74686D2E432E352E322E34> >> >>}
\MusicPDFWriteObject{6984}{<< /Type /Annot /Subtype /Link /Border [0 0 0] /H /I /C [1 0 0] /Rect [216.168 619.276 249.203 631.229] /A << /S /GoTo /D <6571756174696F6E2E432E352E322E34> >> >>}
\MusicPDFWriteObject{6985}{<< /Type /Annot /Subtype /Link /Border [0 0 0] /H /I /C [1 0 0] /Rect [157.996 606.125 190.423 618.078] /A << /S /GoTo /D <6571756174696F6E2E432E352E332E32> >> >>}
\MusicPDFWriteObject{6986}{<< /Type /Annot /Subtype /Link /Border [0 0 0] /H /I /C [1 0 0] /Rect [265.254 354.555 297.68 366.508] /A << /S /GoTo /D <6571756174696F6E2E432E352E332E35> >> >>}
\MusicPDFWriteObject{6987}{<< /Type /Annot /Subtype /Link /Border [0 0 0] /H /I /C [1 0 0] /Rect [361.101 140.206 393.527 152.158] /A << /S /GoTo /D <74686D2E432E342E332E38> >> >>}
\MusicPDFWriteObject{6988}{<< /Type /Annot /Subtype /Link /Border [0 0 0] /H /I /C [1 0 0] /Rect [239.65 574.554 271.894 586.526] /A << /S /GoTo /D <6571756174696F6E2E432E352E342E31> >> >>}
\MusicPDFWriteObject{6989}{<< /Type /Annot /Subtype /Link /Border [0 0 0] /H /I /C [1 0 0] /Rect [186.666 539.47 218.697 550.318] /A << /S /GoTo /D <74686D2E432E352E342E31> >> >>}
\MusicPDFWriteObject{6990}{<< /Type /Annot /Subtype /Link /Border [0 0 0] /H /I /C [1 0 0] /Rect [112.727 515.106 145.154 523.96] /A << /S /GoTo /D <74686D2E432E342E332E35> >> >>}
\MusicPDFWriteObject{6991}{<< /Type /Annot /Subtype /Link /Border [0 0 0] /H /I /C [0 1 0] /Rect [73.77 488.804 105.59 497.672] /A << /S /GoTo /D <636974652E4B72757365> >> >>}
\MusicPDFWriteObject{6992}{<< /Type /Annot /Subtype /Link /Border [0 0 0] /H /I /C [0 1 0] /Rect [311.072 393.794 339.611 402.662] /A << /S /GoTo /D <636974652E4A6F722D62> >> >>}
\MusicPDFWriteObject{6993}{<< /Type /Annot /Subtype /Link /Border [0 0 0] /H /I /C [0 1 0] /Rect [226.577 673.973 254.32 682.84] /A << /S /GoTo /D <636974652E4A6F722D61> >> >>}
\MusicPDFWriteObject{6994}{<< /Type /Annot /Subtype /Link /Border [0 0 0] /H /I /C [0 1 0] /Rect [259.328 673.973 286.262 682.84] /A << /S /GoTo /D <636974652E424A4D> >> >>}
\MusicPDFWriteObject{6995}{<< /Type /Annot /Subtype /Link /Border [0 0 0] /H /I /C [0 1 0] /Rect [366.995 673.973 385.519 682.84] /A << /S /GoTo /D <636974652E4745> >> >>}
\MusicPDFWriteObject{6996}{<< /Type /Annot /Subtype /Link /Border [0 0 0] /H /I /C [1 0 0] /Rect [284.062 645.734 293.511 656.525] /A << /S /GoTo /D <617070656E6469782E44> >> >>}
\MusicPDFWriteObject{6997}{<< /Type /Annot /Subtype /Link /Border [0 0 0] /H /I /C [1 0 0] /Rect [82.959 571.612 412.2 582.403] /A << /S /GoTo /D <73656374696F6E2E442E31> >> >>}
\MusicPDFWriteObject{6998}{<< /Type /Annot /Subtype /Link /Border [0 0 0] /H /I /C [1 0 0] /Rect [82.959 560.398 200.213 569.252] /A << /S /GoTo /D <73756273656374696F6E2E442E312E31> >> >>}
\MusicPDFWriteObject{6999}{<< /Type /Annot /Subtype /Link /Border [0 0 0] /H /I /C [1 0 0] /Rect [82.959 545.311 280.171 556.102] /A << /S /GoTo /D <73756273656374696F6E2E442E312E32> >> >>}
\MusicPDFWriteObject{7000}{<< /Type /Annot /Subtype /Link /Border [0 0 0] /H /I /C [1 0 0] /Rect [82.959 532.16 297.63 542.951] /A << /S /GoTo /D <73756273656374696F6E2E442E312E33> >> >>}
\MusicPDFWriteObject{7001}{<< /Type /Annot /Subtype /Link /Border [0 0 0] /H /I /C [1 0 0] /Rect [82.959 519.009 433.752 529.8] /A << /S /GoTo /D <73756273756273656374696F6E2E442E312E332E31> >> >>}
\MusicPDFWriteObject{7002}{<< /Type /Annot /Subtype /Link /Border [0 0 0] /H /I /C [1 0 0] /Rect [82.959 507.795 335.036 516.649] /A << /S /GoTo /D <73756273756273656374696F6E2E442E312E332E32> >> >>}
\MusicPDFWriteObject{7003}{<< /Type /Annot /Subtype /Link /Border [0 0 0] /H /I /C [1 0 0] /Rect [82.959 492.708 463.771 503.499] /A << /S /GoTo /D <73756273756273656374696F6E2E442E312E332E33> >> >>}
\MusicPDFWriteObject{7004}{<< /Type /Annot /Subtype /Link /Border [0 0 0] /H /I /C [1 0 0] /Rect [82.959 479.557 388.018 490.348] /A << /S /GoTo /D <73756273656374696F6E2E442E312E34> >> >>}
\MusicPDFWriteObject{7005}{<< /Type /Annot /Subtype /Link /Border [0 0 0] /H /I /C [1 0 0] /Rect [82.959 466.406 525.331 477.197] /A << /S /GoTo /D <73756273756273656374696F6E2E442E312E342E31> >> >>}
\MusicPDFWriteObject{7006}{<< /Type /Annot /Subtype /Link /Border [0 0 0] /H /I /C [1 0 0] /Rect [82.959 453.256 311.85 464.046] /A << /S /GoTo /D <73756273756273656374696F6E2E442E312E342E32> >> >>}
\MusicPDFWriteObject{7007}{<< /Type /Annot /Subtype /Link /Border [0 0 0] /H /I /C [1 0 0] /Rect [82.959 442.042 242.862 450.896] /A << /S /GoTo /D <73756273756273656374696F6E2E442E312E342E33> >> >>}
\MusicPDFWriteObject{7008}{<< /Type /Annot /Subtype /Link /Border [0 0 0] /H /I /C [1 0 0] /Rect [82.959 426.954 419.946 437.745] /A << /S /GoTo /D <73756273656374696F6E2E442E312E35> >> >>}
\MusicPDFWriteObject{7009}{<< /Type /Annot /Subtype /Link /Border [0 0 0] /H /I /C [1 0 0] /Rect [82.959 413.803 292.165 424.594] /A << /S /GoTo /D <73756273756273656374696F6E2E442E312E352E31> >> >>}
\MusicPDFWriteObject{7010}{<< /Type /Annot /Subtype /Link /Border [0 0 0] /H /I /C [1 0 0] /Rect [82.959 400.099 275.921 412.054] /A << /S /GoTo /D <73756273756273656374696F6E2E442E312E352E32> >> >>}
\MusicPDFWriteObject{7011}{<< /Type /Annot /Subtype /Link /Border [0 0 0] /H /I /C [1 0 0] /Rect [82.959 387.502 364.972 398.293] /A << /S /GoTo /D <73756273756273656374696F6E2E442E312E352E33> >> >>}
\MusicPDFWriteObject{7012}{<< /Type /Annot /Subtype /Link /Border [0 0 0] /H /I /C [1 0 0] /Rect [82.959 374.351 410.899 385.142] /A << /S /GoTo /D <73756273656374696F6E2E442E312E36> >> >>}
\MusicPDFWriteObject{7013}{<< /Type /Annot /Subtype /Link /Border [0 0 0] /H /I /C [1 0 0] /Rect [82.959 361.2 469.666 371.991] /A << /S /GoTo /D <73756273756273656374696F6E2E442E312E362E31> >> >>}
\MusicPDFWriteObject{7014}{<< /Type /Annot /Subtype /Link /Border [0 0 0] /H /I /C [1 0 0] /Rect [82.959 348.05 259.075 358.84] /A << /S /GoTo /D <73756273756273656374696F6E2E442E312E362E32> >> >>}
\MusicPDFWriteObject{7015}{<< /Type /Annot /Subtype /Link /Border [0 0 0] /H /I /C [1 0 0] /Rect [82.959 334.899 381.544 345.69] /A << /S /GoTo /D <73756273756273656374696F6E2E442E312E362E33> >> >>}
\MusicPDFWriteObject{7016}{<< /Type /Annot /Subtype /Link /Border [0 0 0] /H /I /C [1 0 0] /Rect [82.959 323.685 270.654 332.539] /A << /S /GoTo /D <73756273656374696F6E2E442E312E37> >> >>}
\MusicPDFWriteObject{7017}{<< /Type /Annot /Subtype /Link /Border [0 0 0] /H /I /C [1 0 0] /Rect [82.959 308.597 466.981 319.388] /A << /S /GoTo /D <73756273656374696F6E2E442E312E38> >> >>}
\MusicPDFWriteObject{7018}{<< /Type /Annot /Subtype /Link /Border [0 0 0] /H /I /C [1 0 0] /Rect [82.959 297.383 206.106 306.237] /A << /S /GoTo /D <73756273756273656374696F6E2E442E312E382E31> >> >>}
\MusicPDFWriteObject{7019}{<< /Type /Annot /Subtype /Link /Border [0 0 0] /H /I /C [1 0 0] /Rect [82.959 282.296 170.388 293.087] /A << /S /GoTo /D <73756273756273656374696F6E2E442E312E382E32> >> >>}
\MusicPDFWriteObject{7020}{<< /Type /Annot /Subtype /Link /Border [0 0 0] /H /I /C [1 0 0] /Rect [82.959 269.145 349.976 279.936] /A << /S /GoTo /D <73756273656374696F6E2E442E312E39> >> >>}
\MusicPDFWriteObject{7021}{<< /Type /Annot /Subtype /Link /Border [0 0 0] /H /I /C [1 0 0] /Rect [82.959 255.994 308.767 266.785] /A << /S /GoTo /D <73756273656374696F6E2E442E312E3130> >> >>}
\MusicPDFWriteObject{7022}{<< /Type /Annot /Subtype /Link /Border [0 0 0] /H /I /C [1 0 0] /Rect [82.959 242.844 449.55 253.634] /A << /S /GoTo /D <73756273756273656374696F6E2E442E312E31302E31> >> >>}
\MusicPDFWriteObject{7023}{<< /Type /Annot /Subtype /Link /Border [0 0 0] /H /I /C [1 0 0] /Rect [82.959 229.139 326.514 241.094] /A << /S /GoTo /D <73756273756273656374696F6E2E442E312E31302E32> >> >>}
\MusicPDFWriteObject{7024}{<< /Type /Annot /Subtype /Link /Border [0 0 0] /H /I /C [1 0 0] /Rect [82.959 215.628 333.296 227.943] /A << /S /GoTo /D <73756273756273656374696F6E2E442E312E31302E33> >> >>}
\MusicPDFWriteObject{7025}{<< /Type /Annot /Subtype /Link /Border [0 0 0] /H /I /C [1 0 0] /Rect [82.959 203.391 324.743 214.182] /A << /S /GoTo /D <73656374696F6E2E442E32> >> >>}
\MusicPDFWriteObject{7026}{<< /Type /Annot /Subtype /Link /Border [0 0 0] /H /I /C [1 0 0] /Rect [82.959 190.241 271.899 201.031] /A << /S /GoTo /D <73756273656374696F6E2E442E322E31> >> >>}
\MusicPDFWriteObject{7027}{<< /Type /Annot /Subtype /Link /Border [0 0 0] /H /I /C [1 0 0] /Rect [82.959 177.09 360.545 187.88] /A << /S /GoTo /D <73756273756273656374696F6E2E442E322E312E31> >> >>}
\MusicPDFWriteObject{7028}{<< /Type /Annot /Subtype /Link /Border [0 0 0] /H /I /C [1 0 0] /Rect [82.959 163.385 408.137 175.978] /A << /S /GoTo /D <73756273656374696F6E2E442E322E32> >> >>}
\MusicPDFWriteObject{7029}{<< /Type /Annot /Subtype /Link /Border [0 0 0] /H /I /C [1 0 0] /Rect [82.959 150.703 329.756 161.579] /A << /S /GoTo /D <73756273756273656374696F6E2E442E322E322E31> >> >>}
\MusicPDFWriteObject{7030}{<< /Type /Annot /Subtype /Link /Border [0 0 0] /H /I /C [1 0 0] /Rect [82.959 137.637 341.442 148.428] /A << /S /GoTo /D <73656374696F6E2E442E33> >> >>}
\MusicPDFWriteObject{7031}{<< /Type /Annot /Subtype /Link /Border [0 0 0] /H /I /C [1 0 0] /Rect [82.959 124.487 232.861 135.277] /A << /S /GoTo /D <73756273656374696F6E2E442E332E31> >> >>}
\MusicPDFWriteObject{7032}{<< /Type /Annot /Subtype /Link /Border [0 0 0] /H /I /C [1 0 0] /Rect [82.959 110.782 282.596 122.737] /A << /S /GoTo /D <73756273656374696F6E2E442E332E32> >> >>}
\MusicPDFWriteObject{7033}{<< /Type /Annot /Subtype /Link /Border [0 0 0] /H /I /C [1 0 0] /Rect [82.959 97.631 252.992 109.586] /A << /S /GoTo /D <73756273756273656374696F6E2E442E332E322E31> >> >>}
\MusicPDFWriteObject{7034}{<< /Type /Annot /Subtype /Link /Border [0 0 0] /H /I /C [1 0 0] /Rect [82.959 84.48 293.11 96.436] /A << /S /GoTo /D <73756273756273656374696F6E2E442E332E322E32> >> >>}
\MusicPDFWriteObject{7035}{<< /Type /Annot /Subtype /Link /Border [0 0 0] /H /I /C [1 0 0] /Rect [82.959 684.633 390.48 696.588] /A << /S /GoTo /D <73756273756273656374696F6E2E442E332E322E33> >> >>}
\MusicPDFWriteObject{7036}{<< /Type /Annot /Subtype /Link /Border [0 0 0] /H /I /C [1 0 0] /Rect [82.959 671.95 467.333 682.827] /A << /S /GoTo /D <73756273656374696F6E2E442E332E33> >> >>}
\MusicPDFWriteObject{7037}{<< /Type /Annot /Subtype /Link /Border [0 0 0] /H /I /C [1 0 0] /Rect [82.959 658.885 504.485 669.676] /A << /S /GoTo /D <73756273756273656374696F6E2E442E332E332E31> >> >>}
\MusicPDFWriteObject{7038}{<< /Type /Annot /Subtype /Link /Border [0 0 0] /H /I /C [1 0 0] /Rect [82.959 645.18 284.421 657.136] /A << /S /GoTo /D <73756273756273656374696F6E2E442E332E332E32> >> >>}
\MusicPDFWriteObject{7039}{<< /Type /Annot /Subtype /Link /Border [0 0 0] /H /I /C [1 0 0] /Rect [82.959 632.498 395.742 643.374] /A << /S /GoTo /D <73756273656374696F6E2E442E332E34> >> >>}
\MusicPDFWriteObject{7040}{<< /Type /Annot /Subtype /Link /Border [0 0 0] /H /I /C [1 0 0] /Rect [82.959 619.347 459.764 630.223] /A << /S /GoTo /D <73756273656374696F6E2E442E332E35> >> >>}
\MusicPDFWriteObject{7041}{<< /Type /Annot /Subtype /Link /Border [0 0 0] /H /I /C [1 0 0] /Rect [82.959 606.282 373.05 617.073] /A << /S /GoTo /D <73756273656374696F6E2E442E332E36> >> >>}
\MusicPDFWriteObject{7042}{<< /Type /Annot /Subtype /Link /Border [0 0 0] /H /I /C [1 0 0] /Rect [466.944 495.098 483.877 505.889] /A << /S /GoTo /D <73656374696F6E2E432E31> >> >>}
\MusicPDFWriteObject{7043}{<< /Type /Annot /Subtype /Link /Border [0 0 0] /H /I /C [0 1 0] /Rect [419.896 470.734 444.956 479.601] /A << /S /GoTo /D <636974652E53636861> >> >>}
\MusicPDFWriteObject{7044}{<< /Type /Annot /Subtype /Link /Border [0 0 0] /H /I /C [0 1 0] /Rect [301.664 457.583 326.724 466.451] /A << /S /GoTo /D <636974652E53636861> >> >>}
\MusicPDFWriteObject{7045}{<< /Type /Annot /Subtype /Link /Border [0 0 0] /H /I /C [0 1 0] /Rect [415.681 444.432 440.741 453.3] /A << /S /GoTo /D <636974652E53636861> >> >>}
\MusicPDFWriteObject{7046}{<< /Type /Annot /Subtype /Link /Border [0 0 0] /H /I /C [0 1 0] /Rect [320.479 431.281 345.539 440.149] /A << /S /GoTo /D <636974652E53636861> >> >>}
\MusicPDFWriteObject{7047}{<< /Type /Annot /Subtype /Link /Border [0 0 0] /H /I /C [0 1 0] /Rect [73.826 391.829 98.886 400.697] /A << /S /GoTo /D <636974652E53636861> >> >>}
\MusicPDFWriteObject{7048}{<< /Type /Annot /Subtype /Link /Border [0 0 0] /H /I /C [0 1 0] /Rect [225.235 378.678 244.208 387.546] /A << /S /GoTo /D <636974652E547265> >> >>}
\MusicPDFWriteObject{7049}{<< /Type /Annot /Subtype /Link /Border [0 0 0] /H /I /C [0 1 0] /Rect [73.77 352.377 98.831 361.245] /A << /S /GoTo /D <636974652E53636861> >> >>}
\MusicPDFWriteObject{7050}{<< /Type /Annot /Subtype /Link /Border [0 0 0] /H /I /C [0 1 0] /Rect [100.717 299.774 125.777 308.642] /A << /S /GoTo /D <636974652E53636861> >> >>}
\MusicPDFWriteObject{7051}{<< /Type /Annot /Subtype /Link /Border [0 0 0] /H /I /C [0 1 0] /Rect [73.826 286.623 92.799 295.491] /A << /S /GoTo /D <636974652E547265> >> >>}
\MusicPDFWriteObject{7052}{<< /Type /Annot /Subtype /Link /Border [0 0 0] /H /I /C [0 1 0] /Rect [372.885 260.321 394.106 269.189] /A << /S /GoTo /D <636974652E427275> >> >>}
\MusicPDFWriteObject{7053}{<< /Type /Annot /Subtype /Link /Border [0 0 0] /H /I /C [0 1 0] /Rect [521.938 260.321 543.159 269.189] /A << /S /GoTo /D <636974652E427275> >> >>}
\MusicPDFWriteObject{7054}{<< /Type /Annot /Subtype /Link /Border [0 0 0] /H /I /C [0 1 0] /Rect [520.553 247.171 541.774 256.039] /A << /S /GoTo /D <636974652E427275> >> >>}
\MusicPDFWriteObject{7055}{<< /Type /Annot /Subtype /Link /Border [0 0 0] /H /I /C [1 0 0] /Rect [231.984 205.782 265.302 216.572] /A << /S /GoTo /D <74686D2E412E352E322E39> >> >>}
\MusicPDFWriteObject{7056}{<< /Type /Annot /Subtype /Link /Border [0 0 0] /H /I /C [0 1 0] /Rect [330.943 168.266 352.164 177.134] /A << /S /GoTo /D <636974652E427275> >> >>}
\MusicPDFWriteObject{7057}{<< /Type /Annot /Subtype /Link /Border [0 0 0] /H /I /C [1 0 0] /Rect [236.27 153.179 269.235 163.969] /A << /S /GoTo /D <74686D2E442E312E352E38> >> >>}
\MusicPDFWriteObject{7058}{<< /Type /Annot /Subtype /Link /Border [0 0 0] /H /I /C [0 1 0] /Rect [165.428 115.663 183.731 124.531] /A << /S /GoTo /D <636974652E4445> >> >>}
\MusicPDFWriteObject{7059}{<< /Type /Annot /Subtype /Link /Border [0 0 0] /H /I /C [0 1 0] /Rect [328.13 115.663 345.077 124.531] /A << /S /GoTo /D <636974652E5275> >> >>}
\MusicPDFWriteObject{7060}{<< /Type /Annot /Subtype /Link /Border [0 0 0] /H /I /C [0 1 0] /Rect [410.405 115.663 428.708 124.531] /A << /S /GoTo /D <636974652E4445> >> >>}
\MusicPDFWriteObject{7061}{<< /Type /Annot /Subtype /Link /Border [0 0 0] /H /I /C [0 1 0] /Rect [225.462 89.362 246.684 98.229] /A << /S /GoTo /D <636974652E427275> >> >>}
\MusicPDFWriteObject{7062}{<< /Type /Annot /Subtype /Link /Border [0 0 0] /H /I /C [0 1 0] /Rect [271.968 89.362 290.416 98.229] /A << /S /GoTo /D <636974652E5761> >> >>}
\MusicPDFWriteObject{7063}{<< /Type /Annot /Subtype /Link /Border [0 0 0] /H /I /C [1 0 0] /Rect [133.227 610.324 166.685 622.277] /A << /S /GoTo /D <74686D2E442E312E322E31> >> >>}
\MusicPDFWriteObject{7064}{<< /Type /Annot /Subtype /Link /Border [0 0 0] /H /I /C [1 0 0] /Rect [417.929 610.324 451.388 622.277] /A << /S /GoTo /D <74686D2E442E312E322E37> >> >>}
\MusicPDFWriteObject{7065}{<< /Type /Annot /Subtype /Link /Border [0 0 0] /H /I /C [0 1 0] /Rect [254.89 391.419 271.837 400.287] /A << /S /GoTo /D <636974652E5275> >> >>}
\MusicPDFWriteObject{7066}{<< /Type /Annot /Subtype /Link /Border [0 0 0] /H /I /C [0 1 0] /Rect [447.015 391.419 463.961 400.287] /A << /S /GoTo /D <636974652E5275> >> >>}
\MusicPDFWriteObject{7067}{<< /Type /Annot /Subtype /Link /Border [0 0 0] /H /I /C [0 1 0] /Rect [315.942 378.268 332.889 387.136] /A << /S /GoTo /D <636974652E5275> >> >>}
\MusicPDFWriteObject{7068}{<< /Type /Annot /Subtype /Link /Border [0 0 0] /H /I /C [0 1 0] /Rect [439.961 378.268 456.907 387.136] /A << /S /GoTo /D <636974652E5275> >> >>}
\MusicPDFWriteObject{7069}{<< /Type /Annot /Subtype /Link /Border [0 0 0] /H /I /C [0 1 0] /Rect [87.604 192.536 104.551 201.404] /A << /S /GoTo /D <636974652E5275> >> >>}
\MusicPDFWriteObject{7070}{<< /Type /Annot /Subtype /Link /Border [0 0 0] /H /I /C [0 1 0] /Rect [197.88 446.947 214.826 455.815] /A << /S /GoTo /D <636974652E5275> >> >>}
\MusicPDFWriteObject{7071}{<< /Type /Annot /Subtype /Link /Border [0 0 0] /H /I /C [0 1 0] /Rect [343.811 446.947 362.114 455.815] /A << /S /GoTo /D <636974652E4445> >> >>}
\MusicPDFWriteObject{7072}{<< /Type /Annot /Subtype /Link /Border [0 0 0] /H /I /C [0 1 0] /Rect [146.536 433.796 163.483 442.664] /A << /S /GoTo /D <636974652E5275> >> >>}
\MusicPDFWriteObject{7073}{<< /Type /Annot /Subtype /Link /Border [0 0 0] /H /I /C [1 0 0] /Rect [486.027 431.306 518.869 443.259] /A << /S /GoTo /D <6571756174696F6E2E442E312E312E38> >> >>}
\MusicPDFWriteObject{7074}{<< /Type /Annot /Subtype /Link /Border [0 0 0] /H /I /C [0 1 0] /Rect [146.536 420.645 164.839 429.513] /A << /S /GoTo /D <636974652E4445> >> >>}
\MusicPDFWriteObject{7075}{<< /Type /Annot /Subtype /Link /Border [0 0 0] /H /I /C [0 1 0] /Rect [146.536 407.495 164.839 416.362] /A << /S /GoTo /D <636974652E4445> >> >>}
\MusicPDFWriteObject{7076}{<< /Type /Annot /Subtype /Link /Border [0 0 0] /H /I /C [0 1 0] /Rect [146.536 394.344 163.483 403.212] /A << /S /GoTo /D <636974652E5275> >> >>}
\MusicPDFWriteObject{7077}{<< /Type /Annot /Subtype /Link /Border [0 0 0] /H /I /C [1 0 0] /Rect [73.712 216.011 106.553 227.963] /A << /S /GoTo /D <6571756174696F6E2E442E312E322E31> >> >>}
\MusicPDFWriteObject{7078}{<< /Type /Annot /Subtype /Link /Border [0 0 0] /H /I /C [0 1 0] /Rect [264.202 199.275 300.924 210.08] /A << /S /GoTo /D <636974652E4C616E672D64> >> >>}
\MusicPDFWriteObject{7079}{<< /Type /Annot /Subtype /Link /Border [0 0 0] /H /I /C [0 1 0] /Rect [106.722 188.061 125.025 196.929] /A << /S /GoTo /D <636974652E4445> >> >>}
\MusicPDFWriteObject{7080}{<< /Type /Annot /Subtype /Link /Border [0 0 0] /H /I /C [1 0 0] /Rect [212.085 168.836 244.618 179.627] /A << /S /GoTo /D <74686D2E442E312E322E31> >> >>}
\MusicPDFWriteObject{7081}{<< /Type /Annot /Subtype /Link /Border [0 0 0] /H /I /C [1 0 0] /Rect [190.143 111.541 222.984 123.496] /A << /S /GoTo /D <74686D2E442E312E322E31> >> >>}
\MusicPDFWriteObject{7082}{<< /Type /Annot /Subtype /Link /Border [0 0 0] /H /I /C [1 0 0] /Rect [275.576 684.633 308.417 696.586] /A << /S /GoTo /D <6571756174696F6E2E442E312E322E31> >> >>}
\MusicPDFWriteObject{7083}{<< /Type /Annot /Subtype /Link /Border [0 0 0] /H /I /C [0 1 0] /Rect [148.663 667.313 166.965 676.181] /A << /S /GoTo /D <636974652E4445> >> >>}
\MusicPDFWriteObject{7084}{<< /Type /Annot /Subtype /Link /Border [0 0 0] /H /I /C [1 0 0] /Rect [300.318 376.574 333.159 388.527] /A << /S /GoTo /D <6571756174696F6E2E442E312E322E33> >> >>}
\MusicPDFWriteObject{7085}{<< /Type /Annot /Subtype /Link /Border [0 0 0] /H /I /C [0 1 0] /Rect [399.655 379.064 434.71 387.932] /A << /S /GoTo /D <636974652E426F752D4949> >> >>}
\MusicPDFWriteObject{7086}{<< /Type /Annot /Subtype /Link /Border [0 0 0] /H /I /C [1 0 0] /Rect [434.966 337.122 467.808 349.074] /A << /S /GoTo /D <6571756174696F6E2E442E312E322E35> >> >>}
\MusicPDFWriteObject{7087}{<< /Type /Annot /Subtype /Link /Border [0 0 0] /H /I /C [0 1 0] /Rect [85.518 326.461 120.573 335.329] /A << /S /GoTo /D <636974652E426F752D4949> >> >>}
\MusicPDFWriteObject{7088}{<< /Type /Annot /Subtype /Link /Border [0 0 0] /H /I /C [1 0 0] /Rect [216.316 291.502 248.664 303.454] /A << /S /GoTo /D <6571756174696F6E2E442E312E322E35> >> >>}
\MusicPDFWriteObject{7089}{<< /Type /Annot /Subtype /Link /Border [0 0 0] /H /I /C [1 0 0] /Rect [390.721 128.707 423.871 140.66] /A << /S /GoTo /D <74686D2E442E312E322E31> >> >>}
\MusicPDFWriteObject{7090}{<< /Type /Annot /Subtype /Link /Border [0 0 0] /H /I /C [1 0 0] /Rect [427.154 81.718 460.089 94.55] /A << /S /GoTo /D <74686D2E442E312E322E31> >> >>}
\MusicPDFWriteObject{7091}{<< /Type /Annot /Subtype /Link /Border [0 0 0] /H /I /C [1 0 0] /Rect [511.055 662.918 543.773 682.054] /A << /S /GoTo /D <74686D2E442E312E342E36> >> >>}
\MusicPDFWriteObject{7092}{<< /Type /Annot /Subtype /Link /Border [0 0 0] /H /I /C [0 1 0] /Rect [107.335 632.554 125.638 641.422] /A << /S /GoTo /D <636974652E4445> >> >>}
\MusicPDFWriteObject{7093}{<< /Type /Annot /Subtype /Link /Border [0 0 0] /H /I /C [1 0 0] /Rect [254.316 596.558 287.157 607.349] /A << /S /GoTo /D <74686D2E442E312E322E36> >> >>}
\MusicPDFWriteObject{7094}{<< /Type /Annot /Subtype /Link /Border [0 0 0] /H /I /C [1 0 0] /Rect [279.955 422.453 312.796 434.406] /A << /S /GoTo /D <6571756174696F6E2E442E312E322E38> >> >>}
\MusicPDFWriteObject{7095}{<< /Type /Annot /Subtype /Link /Border [0 0 0] /H /I /C [1 0 0] /Rect [117.744 409.303 150.586 421.255] /A << /S /GoTo /D <6571756174696F6E2E442E312E322E38> >> >>}
\MusicPDFWriteObject{7096}{<< /Type /Annot /Subtype /Link /Border [0 0 0] /H /I /C [1 0 0] /Rect [259.551 388.394 292.393 400.346] /A << /S /GoTo /D <6571756174696F6E2E442E312E322E38> >> >>}
\MusicPDFWriteObject{7097}{<< /Type /Annot /Subtype /Link /Border [0 0 0] /H /I /C [1 0 0] /Rect [221.707 375.083 254.549 389.39] /A << /S /GoTo /D <6571756174696F6E2E442E312E322E38> >> >>}
\MusicPDFWriteObject{7098}{<< /Type /Annot /Subtype /Link /Border [0 0 0] /H /I /C [1 0 0] /Rect [181.977 513.119 207.072 523.91] /A << /S /GoTo /D <73756273656374696F6E2E442E312E33> >> >>}
\MusicPDFWriteObject{7099}{<< /Type /Annot /Subtype /Link /Border [0 0 0] /H /I /C [1 0 0] /Rect [123.368 671.482 160.545 683.435] /A << /S /GoTo /D <6571756174696F6E2E442E312E332E3130> >> >>}
\MusicPDFWriteObject{7100}{<< /Type /Annot /Subtype /Link /Border [0 0 0] /H /I /C [1 0 0] /Rect [166.041 632.03 203.862 643.983] /A << /S /GoTo /D <6571756174696F6E2E442E312E332E3130> >> >>}
\MusicPDFWriteObject{7101}{<< /Type /Annot /Subtype /Link /Border [0 0 0] /H /I /C [1 0 0] /Rect [248.946 383.349 281.788 395.301] /A << /S /GoTo /D <6571756174696F6E2E442E312E332E37> >> >>}
\MusicPDFWriteObject{7102}{<< /Type /Annot /Subtype /Link /Border [0 0 0] /H /I /C [1 0 0] /Rect [437.894 684.633 470.736 696.586] /A << /S /GoTo /D <6571756174696F6E2E442E312E332E36> >> >>}
\MusicPDFWriteObject{7103}{<< /Type /Annot /Subtype /Link /Border [0 0 0] /H /I /C [1 0 0] /Rect [349.791 566.214 382.632 578.167] /A << /S /GoTo /D <6571756174696F6E2E442E312E332E36> >> >>}
\MusicPDFWriteObject{7104}{<< /Type /Annot /Subtype /Link /Border [0 0 0] /H /I /C [1 0 0] /Rect [469.712 500.46 502.831 512.415] /A << /S /GoTo /D <6571756174696F6E2E442E312E332E38> >> >>}
\MusicPDFWriteObject{7105}{<< /Type /Annot /Subtype /Link /Border [0 0 0] /H /I /C [1 0 0] /Rect [133.172 461.562 166.014 472.487] /A << /S /GoTo /D <74686D2E442E312E332E34> >> >>}
\MusicPDFWriteObject{7106}{<< /Type /Annot /Subtype /Link /Border [0 0 0] /H /I /C [1 0 0] /Rect [110.291 385.128 148.113 397.081] /A << /S /GoTo /D <6571756174696F6E2E442E312E332E3135> >> >>}
\MusicPDFWriteObject{7107}{<< /Type /Annot /Subtype /Link /Border [0 0 0] /H /I /C [1 0 0] /Rect [344.832 335.327 382.653 347.28] /A << /S /GoTo /D <6571756174696F6E2E442E312E332E3136> >> >>}
\MusicPDFWriteObject{7108}{<< /Type /Annot /Subtype /Link /Border [0 0 0] /H /I /C [1 0 0] /Rect [162.179 295.875 195.02 307.828] /A << /S /GoTo /D <6571756174696F6E2E442E312E332E38> >> >>}
\MusicPDFWriteObject{7109}{<< /Type /Annot /Subtype /Link /Border [0 0 0] /H /I /C [1 0 0] /Rect [483.942 200.001 516.167 211.954] /A << /S /GoTo /D <74686D2E442E312E332E37> >> >>}
\MusicPDFWriteObject{7110}{<< /Type /Annot /Subtype /Link /Border [0 0 0] /H /I /C [0 1 0] /Rect [248.221 139.862 267.195 148.73] /A << /S /GoTo /D <636974652E547265> >> >>}
\MusicPDFWriteObject{7111}{<< /Type /Annot /Subtype /Link /Border [0 0 0] /H /I /C [1 0 0] /Rect [333.405 671.122 371.944 684.075] /A << /S /GoTo /D <6571756174696F6E2E442E312E332E3137> >> >>}
\MusicPDFWriteObject{7112}{<< /Type /Annot /Subtype /Link /Border [0 0 0] /H /I /C [0 1 0] /Rect [208.059 647.671 227.033 656.539] /A << /S /GoTo /D <636974652E547265> >> >>}
\MusicPDFWriteObject{7113}{<< /Type /Annot /Subtype /Link /Border [0 0 0] /H /I /C [1 0 0] /Rect [207.044 505.524 244.865 517.477] /A << /S /GoTo /D <6571756174696F6E2E442E312E332E3138> >> >>}
\MusicPDFWriteObject{7114}{<< /Type /Annot /Subtype /Link /Border [0 0 0] /H /I /C [1 0 0] /Rect [477.142 505.524 483.105 518.381] /A << /S /GoTo /D <48666F6F746E6F74652E313331> >> >>}
\MusicPDFWriteObject{7115}{<< /Type /Annot /Subtype /Link /Border [0 0 0] /H /I /C [1 0 0] /Rect [421.064 433.692 437.698 444.483] /A << /S /GoTo /D <73656374696F6E2E432E31> >> >>}
\MusicPDFWriteObject{7116}{<< /Type /Annot /Subtype /Link /Border [0 0 0] /H /I /C [1 0 0] /Rect [107.63 420.541 140.665 431.332] /A << /S /GoTo /D <74686D2E432E322E322E32> >> >>}
\MusicPDFWriteObject{7117}{<< /Type /Annot /Subtype /Link /Border [0 0 0] /H /I /C [1 0 0] /Rect [216.686 407.39 249.722 418.181] /A << /S /GoTo /D <74686D2E432E322E332E38> >> >>}
\MusicPDFWriteObject{7118}{<< /Type /Annot /Subtype /Link /Border [0 0 0] /H /I /C [1 0 0] /Rect [297.901 393.686 329.719 405.639] /A << /S /GoTo /D <74686D2E432E322E342E38> >> >>}
\MusicPDFWriteObject{7119}{<< /Type /Annot /Subtype /Link /Border [0 0 0] /H /I /C [1 0 0] /Rect [502.266 302.184 534.601 312.975] /A << /S /GoTo /D <74686D2E432E322E322E32> >> >>}
\MusicPDFWriteObject{7120}{<< /Type /Annot /Subtype /Link /Border [0 0 0] /H /I /C [1 0 0] /Rect [117.493 239.925 150.026 250.716] /A << /S /GoTo /D <74686D2E442E312E332E37> >> >>}
\MusicPDFWriteObject{7121}{<< /Type /Annot /Subtype /Link /Border [0 0 0] /H /I /C [1 0 0] /Rect [351.441 239.925 368.224 250.716] /A << /S /GoTo /D <73656374696F6E2E432E31> >> >>}
\MusicPDFWriteObject{7122}{<< /Type /Annot /Subtype /Link /Border [0 0 0] /H /I /C [1 0 0] /Rect [200.493 199.919 237.191 211.872] /A << /S /GoTo /D <74686D2E432E352E312E3330> >> >>}
\MusicPDFWriteObject{7123}{<< /Type /Annot /Subtype /Link /Border [0 0 0] /H /I /C [1 0 0] /Rect [98.422 137.19 135.829 149.143] /A << /S /GoTo /D <6571756174696F6E2E432E352E312E3438> >> >>}
\MusicPDFWriteObject{7124}{<< /Type /Annot /Subtype /Link /Border [0 0 0] /H /I /C [0 1 0] /Rect [445.942 687.123 464.916 695.991] /A << /S /GoTo /D <636974652E547265> >> >>}
\MusicPDFWriteObject{7125}{<< /Type /Annot /Subtype /Link /Border [0 0 0] /H /I /C [1 0 0] /Rect [230.497 635.055 268.319 647.008] /A << /S /GoTo /D <6571756174696F6E2E442E312E332E3232> >> >>}
\MusicPDFWriteObject{7126}{<< /Type /Annot /Subtype /Link /Border [0 0 0] /H /I /C [0 1 0] /Rect [159.044 624.394 178.017 633.262] /A << /S /GoTo /D <636974652E547265> >> >>}
\MusicPDFWriteObject{7127}{<< /Type /Annot /Subtype /Link /Border [0 0 0] /H /I /C [0 1 0] /Rect [89.402 574.592 108.376 583.46] /A << /S /GoTo /D <636974652E547265> >> >>}
\MusicPDFWriteObject{7128}{<< /Type /Annot /Subtype /Link /Border [0 0 0] /H /I /C [1 0 0] /Rect [201.123 456.608 238.944 468.56] /A << /S /GoTo /D <6571756174696F6E2E442E312E332E3236> >> >>}
\MusicPDFWriteObject{7129}{<< /Type /Annot /Subtype /Link /Border [0 0 0] /H /I /C [1 0 0] /Rect [176.323 420.956 212.637 429.81] /A << /S /GoTo /D <74686D2E442E312E332E37> >> >>}
\MusicPDFWriteObject{7130}{<< /Type /Annot /Subtype /Link /Border [0 0 0] /H /I /C [1 0 0] /Rect [296.954 418.466 335.492 430.421] /A << /S /GoTo /D <6571756174696F6E2E442E312E332E3138> >> >>}
\MusicPDFWriteObject{7131}{<< /Type /Annot /Subtype /Link /Border [0 0 0] /H /I /C [0 1 0] /Rect [236.51 394.655 255.484 403.523] /A << /S /GoTo /D <636974652E547265> >> >>}
\MusicPDFWriteObject{7132}{<< /Type /Annot /Subtype /Link /Border [0 0 0] /H /I /C [1 0 0] /Rect [434.341 379.014 467.182 390.967] /A << /S /GoTo /D <6571756174696F6E2E442E312E332E37> >> >>}
\MusicPDFWriteObject{7133}{<< /Type /Annot /Subtype /Link /Border [0 0 0] /H /I /C [1 0 0] /Rect [129.812 326.411 167.634 338.364] /A << /S /GoTo /D <6571756174696F6E2E442E312E332E3138> >> >>}
\MusicPDFWriteObject{7134}{<< /Type /Annot /Subtype /Link /Border [0 0 0] /H /I /C [0 1 0] /Rect [518.14 279.323 543.2 288.19] /A << /S /GoTo /D <636974652E53636861> >> >>}
\MusicPDFWriteObject{7135}{<< /Type /Annot /Subtype /Link /Border [0 0 0] /H /I /C [0 1 0] /Rect [370.008 266.172 395.068 275.04] /A << /S /GoTo /D <636974652E53636861> >> >>}
\MusicPDFWriteObject{7136}{<< /Type /Annot /Subtype /Link /Border [0 0 0] /H /I /C [0 1 0] /Rect [518.209 177.141 543.27 186.009] /A << /S /GoTo /D <636974652E53636861> >> >>}
\MusicPDFWriteObject{7137}{<< /Type /Annot /Subtype /Link /Border [0 0 0] /H /I /C [0 1 0] /Rect [306.401 137.689 331.461 146.557] /A << /S /GoTo /D <636974652E53636861> >> >>}
\MusicPDFWriteObject{7138}{<< /Type /Annot /Subtype /Link /Border [0 0 0] /H /I /C [1 0 0] /Rect [219.374 122.601 251.192 133.392] /A << /S /GoTo /D <74686D2E432E322E322E32> >> >>}
\MusicPDFWriteObject{7139}{<< /Type /Annot /Subtype /Link /Border [0 0 0] /H /I /C [1 0 0] /Rect [347.136 122.601 378.954 133.392] /A << /S /GoTo /D <74686D2E432E322E332E38> >> >>}
\MusicPDFWriteObject{7140}{<< /Type /Annot /Subtype /Link /Border [0 0 0] /H /I /C [1 0 0] /Rect [132.065 648.759 164.907 659.55] /A << /S /GoTo /D <74686D2E442E312E332E38> >> >>}
\MusicPDFWriteObject{7141}{<< /Type /Annot /Subtype /Link /Border [0 0 0] /H /I /C [1 0 0] /Rect [285.328 529.861 318.786 541.816] /A << /S /GoTo /D <74686D2E442E312E332E35> >> >>}
\MusicPDFWriteObject{7142}{<< /Type /Annot /Subtype /Link /Border [0 0 0] /H /I /C [0 1 0] /Rect [518.099 532.351 543.159 541.219] /A << /S /GoTo /D <636974652E53636861> >> >>}
\MusicPDFWriteObject{7143}{<< /Type /Annot /Subtype /Link /Border [0 0 0] /H /I /C [0 1 0] /Rect [435.162 519.201 460.222 528.068] /A << /S /GoTo /D <636974652E53636861> >> >>}
\MusicPDFWriteObject{7144}{<< /Type /Annot /Subtype /Link /Border [0 0 0] /H /I /C [1 0 0] /Rect [506.953 371.731 539.38 383.684] /A << /S /GoTo /D <6571756174696F6E2E432E322E332E37> >> >>}
\MusicPDFWriteObject{7145}{<< /Type /Annot /Subtype /Link /Border [0 0 0] /H /I /C [1 0 0] /Rect [504.742 358.58 541.954 370.535] /A << /S /GoTo /D <6571756174696F6E2E442E312E332E3238> >> >>}
\MusicPDFWriteObject{7146}{<< /Type /Annot /Subtype /Link /Border [0 0 0] /H /I /C [1 0 0] /Rect [352.511 319.128 385.546 331.083] /A << /S /GoTo /D <74686D2E432E322E332E38> >> >>}
\MusicPDFWriteObject{7147}{<< /Type /Annot /Subtype /Link /Border [0 0 0] /H /I /C [1 0 0] /Rect [105.332 292.826 137.15 304.781] /A << /S /GoTo /D <74686D2E432E322E322E32> >> >>}
\MusicPDFWriteObject{7148}{<< /Type /Annot /Subtype /Link /Border [0 0 0] /H /I /C [1 0 0] /Rect [131.11 279.676 169.29 291.629] /A << /S /GoTo /D <6571756174696F6E2E442E312E332E3238> >> >>}
\MusicPDFWriteObject{7149}{<< /Type /Annot /Subtype /Link /Border [0 0 0] /H /I /C [1 0 0] /Rect [393.974 279.676 432.154 291.629] /A << /S /GoTo /D <6571756174696F6E2E442E312E332E3238> >> >>}
\MusicPDFWriteObject{7150}{<< /Type /Annot /Subtype /Link /Border [0 0 0] /H /I /C [1 0 0] /Rect [276.189 253.374 314.011 265.327] /A << /S /GoTo /D <6571756174696F6E2E442E312E332E3238> >> >>}
\MusicPDFWriteObject{7151}{<< /Type /Annot /Subtype /Link /Border [0 0 0] /H /I /C [1 0 0] /Rect [457.963 253.374 495.785 265.327] /A << /S /GoTo /D <6571756174696F6E2E442E312E332E3138> >> >>}
\MusicPDFWriteObject{7152}{<< /Type /Annot /Subtype /Link /Border [0 0 0] /H /I /C [1 0 0] /Rect [445.354 216.946 478.01 228.901] /A << /S /GoTo /D <74686D2E442E312E332E35> >> >>}
\MusicPDFWriteObject{7153}{<< /Type /Annot /Subtype /Link /Border [0 0 0] /H /I /C [0 1 0] /Rect [276.679 206.286 301.739 215.154] /A << /S /GoTo /D <636974652E53636861> >> >>}
\MusicPDFWriteObject{7154}{<< /Type /Annot /Subtype /Link /Border [0 0 0] /H /I /C [0 1 0] /Rect [272.203 130.406 291.177 139.274] /A << /S /GoTo /D <636974652E547265> >> >>}
\MusicPDFWriteObject{7155}{<< /Type /Annot /Subtype /Link /Border [0 0 0] /H /I /C [1 0 0] /Rect [186.376 592.578 224.198 604.531] /A << /S /GoTo /D <6571756174696F6E2E442E312E332E3138> >> >>}
\MusicPDFWriteObject{7156}{<< /Type /Annot /Subtype /Link /Border [0 0 0] /H /I /C [1 0 0] /Rect [335.752 579.427 369.21 591.382] /A << /S /GoTo /D <74686D2E442E312E332E35> >> >>}
\MusicPDFWriteObject{7157}{<< /Type /Annot /Subtype /Link /Border [0 0 0] /H /I /C [1 0 0] /Rect [221.28 566.276 259.101 578.229] /A << /S /GoTo /D <6571756174696F6E2E442E312E332E3138> >> >>}
\MusicPDFWriteObject{7158}{<< /Type /Annot /Subtype /Link /Border [0 0 0] /H /I /C [1 0 0] /Rect [343.522 500.523 376.363 512.475] /A << /S /GoTo /D <6571756174696F6E2E442E312E332E37> >> >>}
\MusicPDFWriteObject{7159}{<< /Type /Annot /Subtype /Link /Border [0 0 0] /H /I /C [1 0 0] /Rect [157.405 474.774 189.722 486.193] /A << /S /GoTo /D <74686D2E442E312E332E34> >> >>}
\MusicPDFWriteObject{7160}{<< /Type /Annot /Subtype /Link /Border [0 0 0] /H /I /C [1 0 0] /Rect [119.052 253.937 156.873 265.89] /A << /S /GoTo /D <6571756174696F6E2E442E312E332E3333> >> >>}
\MusicPDFWriteObject{7161}{<< /Type /Annot /Subtype /Link /Border [0 0 0] /H /I /C [1 0 0] /Rect [149.01 142.38 186.832 154.333] /A << /S /GoTo /D <6571756174696F6E2E442E312E332E3335> >> >>}
\MusicPDFWriteObject{7162}{<< /Type /Annot /Subtype /Link /Border [0 0 0] /H /I /C [1 0 0] /Rect [324.05 129.23 361.872 141.182] /A << /S /GoTo /D <6571756174696F6E2E442E312E332E3335> >> >>}
\MusicPDFWriteObject{7163}{<< /Type /Annot /Subtype /Link /Border [0 0 0] /H /I /C [1 0 0] /Rect [73.712 102.928 111.533 114.881] /A << /S /GoTo /D <6571756174696F6E2E442E312E332E3331> >> >>}
\MusicPDFWriteObject{7164}{<< /Type /Annot /Subtype /Link /Border [0 0 0] /H /I /C [1 0 0] /Rect [304.022 684.633 341.843 696.586] /A << /S /GoTo /D <6571756174696F6E2E442E312E332E3138> >> >>}
\MusicPDFWriteObject{7165}{<< /Type /Annot /Subtype /Link /Border [0 0 0] /H /I /C [0 1 0] /Rect [300.042 567.3 319.016 576.168] /A << /S /GoTo /D <636974652E547265> >> >>}
\MusicPDFWriteObject{7166}{<< /Type /Annot /Subtype /Link /Border [0 0 0] /H /I /C [1 0 0] /Rect [161.679 499.789 199.501 511.741] /A << /S /GoTo /D <6571756174696F6E2E442E312E332E3138> >> >>}
\MusicPDFWriteObject{7167}{<< /Type /Annot /Subtype /Link /Border [0 0 0] /H /I /C [1 0 0] /Rect [321.607 499.789 327.57 512.646] /A << /S /GoTo /D <48666F6F746E6F74652E313332> >> >>}
\MusicPDFWriteObject{7168}{<< /Type /Annot /Subtype /Link /Border [0 0 0] /H /I /C [1 0 0] /Rect [163.132 460.336 196.436 472.291] /A << /S /GoTo /D <74686D2E442E312E332E35> >> >>}
\MusicPDFWriteObject{7169}{<< /Type /Annot /Subtype /Link /Border [0 0 0] /H /I /C [0 1 0] /Rect [443.642 449.676 462.615 458.543] /A << /S /GoTo /D <636974652E547265> >> >>}
\MusicPDFWriteObject{7170}{<< /Type /Annot /Subtype /Link /Border [0 0 0] /H /I /C [0 1 0] /Rect [142.671 436.525 161.644 445.393] /A << /S /GoTo /D <636974652E547265> >> >>}
\MusicPDFWriteObject{7171}{<< /Type /Annot /Subtype /Link /Border [0 0 0] /H /I /C [0 1 0] /Rect [430.779 436.525 455.839 445.393] /A << /S /GoTo /D <636974652E53636861> >> >>}
\MusicPDFWriteObject{7172}{<< /Type /Annot /Subtype /Link /Border [0 0 0] /H /I /C [1 0 0] /Rect [226.74 419.197 260.045 432.856] /A << /S /GoTo /D <74686D2E442E312E332E35> >> >>}
\MusicPDFWriteObject{7173}{<< /Type /Annot /Subtype /Link /Border [0 0 0] /H /I /C [0 1 0] /Rect [150.763 410.223 175.823 419.091] /A << /S /GoTo /D <636974652E53636861> >> >>}
\MusicPDFWriteObject{7174}{<< /Type /Annot /Subtype /Link /Border [0 0 0] /H /I /C [0 1 0] /Rect [173.866 397.073 192.839 405.94] /A << /S /GoTo /D <636974652E547265> >> >>}
\MusicPDFWriteObject{7175}{<< /Type /Annot /Subtype /Link /Border [0 0 0] /H /I /C [0 1 0] /Rect [186.135 383.922 205.108 392.79] /A << /S /GoTo /D <636974652E547265> >> >>}
\MusicPDFWriteObject{7176}{<< /Type /Annot /Subtype /Link /Border [0 0 0] /H /I /C [0 1 0] /Rect [202.595 370.771 227.655 379.639] /A << /S /GoTo /D <636974652E53636861> >> >>}
\MusicPDFWriteObject{7177}{<< /Type /Annot /Subtype /Link /Border [0 0 0] /H /I /C [0 1 0] /Rect [437.235 370.771 462.295 379.639] /A << /S /GoTo /D <636974652E53636861> >> >>}
\MusicPDFWriteObject{7178}{<< /Type /Annot /Subtype /Link /Border [0 0 0] /H /I /C [0 1 0] /Rect [125.402 209.325 143.705 218.193] /A << /S /GoTo /D <636974652E4445> >> >>}
\MusicPDFWriteObject{7179}{<< /Type /Annot /Subtype /Link /Border [0 0 0] /H /I /C [0 1 0] /Rect [167.414 687.123 185.717 695.991] /A << /S /GoTo /D <636974652E4445> >> >>}
\MusicPDFWriteObject{7180}{<< /Type /Annot /Subtype /Link /Border [0 0 0] /H /I /C [1 0 0] /Rect [237.301 658.331 270.759 670.286] /A << /S /GoTo /D <74686D2E442E312E342E32> >> >>}
\MusicPDFWriteObject{7181}{<< /Type /Annot /Subtype /Link /Border [0 0 0] /H /I /C [1 0 0] /Rect [462.747 658.331 496.206 670.286] /A << /S /GoTo /D <6571756174696F6E2E442E312E342E31> >> >>}
\MusicPDFWriteObject{7182}{<< /Type /Annot /Subtype /Link /Border [0 0 0] /H /I /C [1 0 0] /Rect [202.728 551.202 235.416 563.155] /A << /S /GoTo /D <74686D2E442E312E342E32> >> >>}
\MusicPDFWriteObject{7183}{<< /Type /Annot /Subtype /Link /Border [0 0 0] /H /I /C [1 0 0] /Rect [297.93 551.202 330.617 563.155] /A << /S /GoTo /D <74686D2E442E312E342E33> >> >>}
\MusicPDFWriteObject{7184}{<< /Type /Annot /Subtype /Link /Border [0 0 0] /H /I /C [1 0 0] /Rect [509.22 551.202 541.908 563.155] /A << /S /GoTo /D <6571756174696F6E2E442E312E342E33> >> >>}
\MusicPDFWriteObject{7185}{<< /Type /Annot /Subtype /Link /Border [0 0 0] /H /I /C [0 1 0] /Rect [325.191 405.794 346.412 414.662] /A << /S /GoTo /D <636974652E427275> >> >>}
\MusicPDFWriteObject{7186}{<< /Type /Annot /Subtype /Link /Border [0 0 0] /H /I /C [1 0 0] /Rect [127.694 377.003 160.535 388.956] /A << /S /GoTo /D <6571756174696F6E2E442E312E342E34> >> >>}
\MusicPDFWriteObject{7187}{<< /Type /Annot /Subtype /Link /Border [0 0 0] /H /I /C [0 1 0] /Rect [196.106 117.659 214.409 126.527] /A << /S /GoTo /D <636974652E4445> >> >>}
\MusicPDFWriteObject{7188}{<< /Type /Annot /Subtype /Link /Border [0 0 0] /H /I /C [0 1 0] /Rect [397.419 117.659 415.721 126.527] /A << /S /GoTo /D <636974652E4445> >> >>}
\MusicPDFWriteObject{7189}{<< /Type /Annot /Subtype /Link /Border [0 0 0] /H /I /C [1 0 0] /Rect [308.314 685.187 341.309 695.977] /A << /S /GoTo /D <74686D2E442E312E342E34> >> >>}
\MusicPDFWriteObject{7190}{<< /Type /Annot /Subtype /Link /Border [0 0 0] /H /I /C [1 0 0] /Rect [134.555 616.551 167.397 627.342] /A << /S /GoTo /D <74686D2E442E312E342E36> >> >>}
\MusicPDFWriteObject{7191}{<< /Type /Annot /Subtype /Link /Border [0 0 0] /H /I /C [1 0 0] /Rect [436.789 542.816 470.247 557.289] /A << /S /GoTo /D <6571756174696F6E2E442E312E342E38> >> >>}
\MusicPDFWriteObject{7192}{<< /Type /Annot /Subtype /Link /Border [0 0 0] /H /I /C [1 0 0] /Rect [197.5 359.757 230.342 371.71] /A << /S /GoTo /D <6571756174696F6E2E442E312E342E39> >> >>}
\MusicPDFWriteObject{7193}{<< /Type /Annot /Subtype /Link /Border [0 0 0] /H /I /C [1 0 0] /Rect [269.675 314.826 301.899 327.139] /A << /S /GoTo /D <6571756174696F6E2E442E312E342E38> >> >>}
\MusicPDFWriteObject{7194}{<< /Type /Annot /Subtype /Link /Border [0 0 0] /H /I /C [1 0 0] /Rect [333.503 302.035 366.345 313.988] /A << /S /GoTo /D <74686D2E442E312E342E33> >> >>}
\MusicPDFWriteObject{7195}{<< /Type /Annot /Subtype /Link /Border [0 0 0] /H /I /C [1 0 0] /Rect [138.886 288.885 171.111 300.837] /A << /S /GoTo /D <74686D2E442E312E342E33> >> >>}
\MusicPDFWriteObject{7196}{<< /Type /Annot /Subtype /Link /Border [0 0 0] /H /I /C [0 1 0] /Rect [133.089 128.595 151.537 137.462] /A << /S /GoTo /D <636974652E5761> >> >>}
\MusicPDFWriteObject{7197}{<< /Type /Annot /Subtype /Link /Border [0 0 0] /H /I /C [1 0 0] /Rect [293.449 500.398 326.907 511.189] /A << /S /GoTo /D <74686D2E442E312E342E37> >> >>}
\MusicPDFWriteObject{7198}{<< /Type /Annot /Subtype /Link /Border [0 0 0] /H /I /C [0 1 0] /Rect [199.181 332.95 220.402 341.817] /A << /S /GoTo /D <636974652E427275> >> >>}
\MusicPDFWriteObject{7199}{<< /Type /Annot /Subtype /Link /Border [0 0 0] /H /I /C [0 1 0] /Rect [484.088 293.497 505.309 302.365] /A << /S /GoTo /D <636974652E427275> >> >>}
\MusicPDFWriteObject{7200}{<< /Type /Annot /Subtype /Link /Border [0 0 0] /H /I /C [1 0 0] /Rect [326.183 552.389 359.024 564.342] /A << /S /GoTo /D <6571756174696F6E2E442E312E352E31> >> >>}
\MusicPDFWriteObject{7201}{<< /Type /Annot /Subtype /Link /Border [0 0 0] /H /I /C [1 0 0] /Rect [263.463 489.66 296.675 501.613] /A << /S /GoTo /D <74686D2E442E312E342E33> >> >>}
\MusicPDFWriteObject{7202}{<< /Type /Annot /Subtype /Link /Border [0 0 0] /H /I /C [1 0 0] /Rect [225.422 448.226 258.572 463.363] /A << /S /GoTo /D <74686D2E442E312E342E33> >> >>}
\MusicPDFWriteObject{7203}{<< /Type /Annot /Subtype /Link /Border [0 0 0] /H /I /C [1 0 0] /Rect [348.008 421.925 380.85 433.878] /A << /S /GoTo /D <74686D2E442E312E342E33> >> >>}
\MusicPDFWriteObject{7204}{<< /Type /Annot /Subtype /Link /Border [0 0 0] /H /I /C [1 0 0] /Rect [490.44 381.385 522.664 396.57] /A << /S /GoTo /D <74686D2E442E312E312E32> >> >>}
\MusicPDFWriteObject{7205}{<< /Type /Annot /Subtype /Link /Border [0 0 0] /H /I /C [1 0 0] /Rect [352.521 368.074 390.343 380.027] /A << /S /GoTo /D <6571756174696F6E2E442E312E312E3131> >> >>}
\MusicPDFWriteObject{7206}{<< /Type /Annot /Subtype /Link /Border [0 0 0] /H /I /C [1 0 0] /Rect [483.694 108.947 521.515 120.899] /A << /S /GoTo /D <6571756174696F6E2E442E312E342E3130> >> >>}
\MusicPDFWriteObject{7207}{<< /Type /Annot /Subtype /Link /Border [0 0 0] /H /I /C [1 0 0] /Rect [106.334 95.746 138.96 109.943] /A << /S /GoTo /D <74686D2E442E312E342E37> >> >>}
\MusicPDFWriteObject{7208}{<< /Type /Annot /Subtype /Link /Border [0 0 0] /H /I /C [1 0 0] /Rect [464.862 82.596 497.704 94.549] /A << /S /GoTo /D <6571756174696F6E2E442E312E342E39> >> >>}
\MusicPDFWriteObject{7209}{<< /Type /Annot /Subtype /Link /Border [0 0 0] /H /I /C [1 0 0] /Rect [326.455 487.794 359.297 499.749] /A << /S /GoTo /D <74686D2E442E312E332E31> >> >>}
\MusicPDFWriteObject{7210}{<< /Type /Annot /Subtype /Link /Border [0 0 0] /H /I /C [1 0 0] /Rect [506.544 435.192 539.385 447.144] /A << /S /GoTo /D <6571756174696F6E2E442E312E352E33> >> >>}
\MusicPDFWriteObject{7211}{<< /Type /Annot /Subtype /Link /Border [0 0 0] /H /I /C [1 0 0] /Rect [222.04 422.594 254.882 433.385] /A << /S /GoTo /D <74686D2E442E312E332E33> >> >>}
\MusicPDFWriteObject{7212}{<< /Type /Annot /Subtype /Link /Border [0 0 0] /H /I /C [1 0 0] /Rect [417.479 174.651 454.763 186.606] /A << /S /GoTo /D <6571756174696F6E2E442E312E332E3230> >> >>}
\MusicPDFWriteObject{7213}{<< /Type /Annot /Subtype /Link /Border [0 0 0] /H /I /C [1 0 0] /Rect [394.465 671.482 433.003 683.455] /A << /S /GoTo /D <6571756174696F6E2E442E312E332E3230> >> >>}
\MusicPDFWriteObject{7214}{<< /Type /Annot /Subtype /Link /Border [0 0 0] /H /I /C [1 0 0] /Rect [283.106 592.578 315.947 604.55] /A << /S /GoTo /D <6571756174696F6E2E442E312E332E39> >> >>}
\MusicPDFWriteObject{7215}{<< /Type /Annot /Subtype /Link /Border [0 0 0] /H /I /C [1 0 0] /Rect [124.623 452.785 157.464 463.576] /A << /S /GoTo /D <74686D2E442E312E392E32> >> >>}
\MusicPDFWriteObject{7216}{<< /Type /Annot /Subtype /Link /Border [0 0 0] /H /I /C [1 0 0] /Rect [71.004 379.171 103.228 391.126] /A << /S /GoTo /D <74686D2E442E312E342E36> >> >>}
\MusicPDFWriteObject{7217}{<< /Type /Annot /Subtype /Link /Border [0 0 0] /H /I /C [0 1 0] /Rect [122.913 152.471 147.973 161.339] /A << /S /GoTo /D <636974652E53636861> >> >>}
\MusicPDFWriteObject{7218}{<< /Type /Annot /Subtype /Link /Border [0 0 0] /H /I /C [0 1 0] /Rect [518.209 152.471 543.27 161.339] /A << /S /GoTo /D <636974652E53636861> >> >>}
\MusicPDFWriteObject{7219}{<< /Type /Annot /Subtype /Link /Border [0 0 0] /H /I /C [0 1 0] /Rect [408.355 118.644 429.576 127.512] /A << /S /GoTo /D <636974652E427275> >> >>}
\MusicPDFWriteObject{7220}{<< /Type /Annot /Subtype /Link /Border [0 0 0] /H /I /C [1 0 0] /Rect [385.855 650.273 419.313 662.228] /A << /S /GoTo /D <74686D2E442E312E332E37> >> >>}
\MusicPDFWriteObject{7221}{<< /Type /Annot /Subtype /Link /Border [0 0 0] /H /I /C [0 1 0] /Rect [190.357 639.613 215.417 648.481] /A << /S /GoTo /D <636974652E53636861> >> >>}
\MusicPDFWriteObject{7222}{<< /Type /Annot /Subtype /Link /Border [0 0 0] /H /I /C [0 1 0] /Rect [224.695 626.462 249.755 635.33] /A << /S /GoTo /D <636974652E53636861> >> >>}
\MusicPDFWriteObject{7223}{<< /Type /Annot /Subtype /Link /Border [0 0 0] /H /I /C [1 0 0] /Rect [392.594 610.821 426.053 622.776] /A << /S /GoTo /D <74686D2E442E312E332E37> >> >>}
\MusicPDFWriteObject{7224}{<< /Type /Annot /Subtype /Link /Border [0 0 0] /H /I /C [0 1 0] /Rect [218.888 587.01 243.948 595.878] /A << /S /GoTo /D <636974652E53636861> >> >>}
\MusicPDFWriteObject{7225}{<< /Type /Annot /Subtype /Link /Border [0 0 0] /H /I /C [0 1 0] /Rect [374.407 587.01 399.468 595.878] /A << /S /GoTo /D <636974652E53636861> >> >>}
\MusicPDFWriteObject{7226}{<< /Type /Annot /Subtype /Link /Border [0 0 0] /H /I /C [0 1 0] /Rect [92.963 573.859 118.023 582.727] /A << /S /GoTo /D <636974652E53636861> >> >>}
\MusicPDFWriteObject{7227}{<< /Type /Annot /Subtype /Link /Border [0 0 0] /H /I /C [0 1 0] /Rect [257.306 508.105 282.366 516.973] /A << /S /GoTo /D <636974652E53636861> >> >>}
\MusicPDFWriteObject{7228}{<< /Type /Annot /Subtype /Link /Border [0 0 0] /H /I /C [0 1 0] /Rect [330.454 494.954 355.514 503.822] /A << /S /GoTo /D <636974652E53636861> >> >>}
\MusicPDFWriteObject{7229}{<< /Type /Annot /Subtype /Link /Border [0 0 0] /H /I /C [0 1 0] /Rect [357.996 468.653 376.969 477.521] /A << /S /GoTo /D <636974652E547265> >> >>}
\MusicPDFWriteObject{7230}{<< /Type /Annot /Subtype /Link /Border [0 0 0] /H /I /C [0 1 0] /Rect [73.737 429.201 98.797 438.068] /A << /S /GoTo /D <636974652E53636861> >> >>}
\MusicPDFWriteObject{7231}{<< /Type /Annot /Subtype /Link /Border [0 0 0] /H /I /C [0 1 0] /Rect [372.249 402.833 397.309 411.701] /A << /S /GoTo /D <636974652E53636861> >> >>}
\MusicPDFWriteObject{7232}{<< /Type /Annot /Subtype /Link /Border [0 0 0] /H /I /C [0 1 0] /Rect [341.247 389.682 366.307 398.55] /A << /S /GoTo /D <636974652E53636861> >> >>}
\MusicPDFWriteObject{7233}{<< /Type /Annot /Subtype /Link /Border [0 0 0] /H /I /C [0 1 0] /Rect [409.317 368.473 430.538 377.341] /A << /S /GoTo /D <636974652E427275> >> >>}
\MusicPDFWriteObject{7234}{<< /Type /Annot /Subtype /Link /Border [0 0 0] /H /I /C [0 1 0] /Rect [435.251 355.322 456.472 364.19] /A << /S /GoTo /D <636974652E427275> >> >>}
\MusicPDFWriteObject{7235}{<< /Type /Annot /Subtype /Link /Border [0 0 0] /H /I /C [0 1 0] /Rect [218.441 315.87 243.501 324.738] /A << /S /GoTo /D <636974652E53636861> >> >>}
\MusicPDFWriteObject{7236}{<< /Type /Annot /Subtype /Link /Border [0 0 0] /H /I /C [0 1 0] /Rect [339.401 302.719 358.374 311.587] /A << /S /GoTo /D <636974652E547265> >> >>}
\MusicPDFWriteObject{7237}{<< /Type /Annot /Subtype /Link /Border [0 0 0] /H /I /C [0 1 0] /Rect [131.816 276.418 153.037 285.286] /A << /S /GoTo /D <636974652E427275> >> >>}
\MusicPDFWriteObject{7238}{<< /Type /Annot /Subtype /Link /Border [0 0 0] /H /I /C [0 1 0] /Rect [518.168 276.418 543.228 285.286] /A << /S /GoTo /D <636974652E53636861> >> >>}
\MusicPDFWriteObject{7239}{<< /Type /Annot /Subtype /Link /Border [0 0 0] /H /I /C [1 0 0] /Rect [425.326 226.417 449.892 238.37] /A << /S /GoTo /D <73756273656374696F6E2E432E312E35> >> >>}
\MusicPDFWriteObject{7240}{<< /Type /Annot /Subtype /Link /Border [0 0 0] /H /I /C [1 0 0] /Rect [355.672 200.116 388.514 212.068] /A << /S /GoTo /D <74686D2E442E332E332E37> >> >>}
\MusicPDFWriteObject{7241}{<< /Type /Annot /Subtype /Link /Border [0 0 0] /H /I /C [0 1 0] /Rect [155.628 133.886 180.688 142.754] /A << /S /GoTo /D <636974652E53636861> >> >>}
\MusicPDFWriteObject{7242}{<< /Type /Annot /Subtype /Link /Border [0 0 0] /H /I /C [0 1 0] /Rect [87.568 118.799 126.8 129.603] /A << /S /GoTo /D <636974652E576E67722D32> >> >>}
\MusicPDFWriteObject{7243}{<< /Type /Annot /Subtype /Link /Border [0 0 0] /H /I /C [0 1 0] /Rect [518.16 120.735 543.22 129.603] /A << /S /GoTo /D <636974652E53636861> >> >>}
\MusicPDFWriteObject{7244}{<< /Type /Annot /Subtype /Link /Border [0 0 0] /H /I /C [0 1 0] /Rect [144.599 107.585 163.573 116.452] /A << /S /GoTo /D <636974652E547265> >> >>}
\MusicPDFWriteObject{7245}{<< /Type /Annot /Subtype /Link /Border [0 0 0] /H /I /C [1 0 0] /Rect [333.045 83.149 365.886 93.94] /A << /S /GoTo /D <74686D2E442E312E352E39> >> >>}
\MusicPDFWriteObject{7246}{<< /Type /Annot /Subtype /Link /Border [0 0 0] /H /I /C [1 0 0] /Rect [163.206 464.25 196.048 476.203] /A << /S /GoTo /D <6571756174696F6E2E442E312E352E36> >> >>}
\MusicPDFWriteObject{7247}{<< /Type /Annot /Subtype /Link /Border [0 0 0] /H /I /C [1 0 0] /Rect [160.596 437.351 194.055 449.603] /A << /S /GoTo /D <74686D2E442E312E352E39> >> >>}
\MusicPDFWriteObject{7248}{<< /Type /Annot /Subtype /Link /Border [0 0 0] /H /I /C [0 1 0] /Rect [281.776 425.051 316.928 435.856] /A << /S /GoTo /D <636974652E566F67742D31> >> >>}
\MusicPDFWriteObject{7249}{<< /Type /Annot /Subtype /Link /Border [0 0 0] /H /I /C [1 0 0] /Rect [191.437 388.102 215.663 398.951] /A << /S /GoTo /D <73756273656374696F6E2E432E312E35> >> >>}
\MusicPDFWriteObject{7250}{<< /Type /Annot /Subtype /Link /Border [0 0 0] /H /I /C [1 0 0] /Rect [71.004 150.941 104.462 161.732] /A << /S /GoTo /D <74686D2E442E312E362E31> >> >>}
\MusicPDFWriteObject{7251}{<< /Type /Annot /Subtype /Link /Border [0 0 0] /H /I /C [1 0 0] /Rect [229.73 585.367 261.955 596.158] /A << /S /GoTo /D <74686D2E442E312E342E37> >> >>}
\MusicPDFWriteObject{7252}{<< /Type /Annot /Subtype /Link /Border [0 0 0] /H /I /C [0 1 0] /Rect [347.059 547.851 364.075 556.719] /A << /S /GoTo /D <636974652E4175> >> >>}
\MusicPDFWriteObject{7253}{<< /Type /Annot /Subtype /Link /Border [0 0 0] /H /I /C [1 0 0] /Rect [418.874 461.64 451.715 473.592] /A << /S /GoTo /D <6571756174696F6E2E442E312E362E34> >> >>}
\MusicPDFWriteObject{7254}{<< /Type /Annot /Subtype /Link /Border [0 0 0] /H /I /C [1 0 0] /Rect [331.341 230.597 361.554 241.388] /A << /S /GoTo /D <74686D2E322E312E312E34> >> >>}
\MusicPDFWriteObject{7255}{<< /Type /Annot /Subtype /Link /Border [0 0 0] /H /I /C [1 0 0] /Rect [171.799 583.936 204.641 595.889] /A << /S /GoTo /D <6571756174696F6E2E442E312E352E32> >> >>}
\MusicPDFWriteObject{7256}{<< /Type /Annot /Subtype /Link /Border [0 0 0] /H /I /C [1 0 0] /Rect [71.004 570.785 104.462 582.738] /A << /S /GoTo /D <74686D2E442E312E352E33> >> >>}
\MusicPDFWriteObject{7257}{<< /Type /Annot /Subtype /Link /Border [0 0 0] /H /I /C [1 0 0] /Rect [205.98 518.182 238.822 530.775] /A << /S /GoTo /D <6571756174696F6E2E442E312E362E37> >> >>}
\MusicPDFWriteObject{7258}{<< /Type /Annot /Subtype /Link /Border [0 0 0] /H /I /C [0 1 0] /Rect [164.247 370.86 182.694 379.727] /A << /S /GoTo /D <636974652E5761> >> >>}
\MusicPDFWriteObject{7259}{<< /Type /Annot /Subtype /Link /Border [0 0 0] /H /I /C [1 0 0] /Rect [223.984 241.866 262.272 252.657] /A << /S /GoTo /D <74686D2E442E312E362E3130> >> >>}
\MusicPDFWriteObject{7260}{<< /Type /Annot /Subtype /Link /Border [0 0 0] /H /I /C [1 0 0] /Rect [468.8 228.715 501.024 239.506] /A << /S /GoTo /D <74686D2E442E312E352E35> >> >>}
\MusicPDFWriteObject{7261}{<< /Type /Annot /Subtype /Link /Border [0 0 0] /H /I /C [1 0 0] /Rect [194.701 615.637 233.24 627.59] /A << /S /GoTo /D <74686D2E442E312E362E3139> >> >>}
\MusicPDFWriteObject{7262}{<< /Type /Annot /Subtype /Link /Border [0 0 0] /H /I /C [1 0 0] /Rect [188.595 459.409 226.166 471.362] /A << /S /GoTo /D <6571756174696F6E2E442E312E362E3133> >> >>}
\MusicPDFWriteObject{7263}{<< /Type /Annot /Subtype /Link /Border [0 0 0] /H /I /C [1 0 0] /Rect [410.704 361.424 449.242 373.377] /A << /S /GoTo /D <6571756174696F6E2E442E312E362E3133> >> >>}
\MusicPDFWriteObject{7264}{<< /Type /Annot /Subtype /Link /Border [0 0 0] /H /I /C [1 0 0] /Rect [420.723 685.187 452.947 695.977] /A << /S /GoTo /D <74686D2E442E312E362E34> >> >>}
\MusicPDFWriteObject{7265}{<< /Type /Annot /Subtype /Link /Border [0 0 0] /H /I /C [1 0 0] /Rect [122.839 658.332 156.297 670.284] /A << /S /GoTo /D <74686D2E442E312E362E37> >> >>}
\MusicPDFWriteObject{7266}{<< /Type /Annot /Subtype /Link /Border [0 0 0] /H /I /C [1 0 0] /Rect [199.048 579.295 236.869 591.247] /A << /S /GoTo /D <6571756174696F6E2E442E312E362E3135> >> >>}
\MusicPDFWriteObject{7267}{<< /Type /Annot /Subtype /Link /Border [0 0 0] /H /I /C [1 0 0] /Rect [392.592 529.677 425.403 541.629] /A << /S /GoTo /D <74686D2E442E312E312E31> >> >>}
\MusicPDFWriteObject{7268}{<< /Type /Annot /Subtype /Link /Border [0 0 0] /H /I /C [1 0 0] /Rect [187.258 479.689 225.079 491.642] /A << /S /GoTo /D <6571756174696F6E2E442E312E362E3135> >> >>}
\MusicPDFWriteObject{7269}{<< /Type /Annot /Subtype /Link /Border [0 0 0] /H /I /C [1 0 0] /Rect [165.404 466.538 203.226 478.491] /A << /S /GoTo /D <6571756174696F6E2E442E312E362E3135> >> >>}
\MusicPDFWriteObject{7270}{<< /Type /Annot /Subtype /Link /Border [0 0 0] /H /I /C [1 0 0] /Rect [230.724 378.968 268.546 389.758] /A << /S /GoTo /D <74686D2E442E312E362E3133> >> >>}
\MusicPDFWriteObject{7271}{<< /Type /Annot /Subtype /Link /Border [0 0 0] /H /I /C [1 0 0] /Rect [423.413 335.524 461.951 347.479] /A << /S /GoTo /D <6571756174696F6E2E442E312E362E3136> >> >>}
\MusicPDFWriteObject{7272}{<< /Type /Annot /Subtype /Link /Border [0 0 0] /H /I /C [1 0 0] /Rect [158.901 135.199 191.743 147.152] /A << /S /GoTo /D <6571756174696F6E2E442E312E352E32> >> >>}
\MusicPDFWriteObject{7273}{<< /Type /Annot /Subtype /Link /Border [0 0 0] /H /I /C [1 0 0] /Rect [191.031 122.048 224.058 134.003] /A << /S /GoTo /D <74686D2E442E312E332E33> >> >>}
\MusicPDFWriteObject{7274}{<< /Type /Annot /Subtype /Link /Border [0 0 0] /H /I /C [1 0 0] /Rect [157.176 82.595 194.281 94.55] /A << /S /GoTo /D <74686D2E442E312E362E3133> >> >>}
\MusicPDFWriteObject{7275}{<< /Type /Annot /Subtype /Link /Border [0 0 0] /H /I /C [1 0 0] /Rect [271.404 547.522 309.226 559.477] /A << /S /GoTo /D <6571756174696F6E2E442E312E362E3132> >> >>}
\MusicPDFWriteObject{7276}{<< /Type /Annot /Subtype /Link /Border [0 0 0] /H /I /C [1 0 0] /Rect [423.519 547.522 461.341 559.475] /A << /S /GoTo /D <6571756174696F6E2E442E312E362E3138> >> >>}
\MusicPDFWriteObject{7277}{<< /Type /Annot /Subtype /Link /Border [0 0 0] /H /I /C [1 0 0] /Rect [354.771 402.808 386.995 414.763] /A << /S /GoTo /D <6571756174696F6E2E442E312E332E39> >> >>}
\MusicPDFWriteObject{7278}{<< /Type /Annot /Subtype /Link /Border [0 0 0] /H /I /C [1 0 0] /Rect [216.821 389.658 249.662 401.613] /A << /S /GoTo /D <74686D2E442E312E332E34> >> >>}
\MusicPDFWriteObject{7279}{<< /Type /Annot /Subtype /Link /Border [0 0 0] /H /I /C [1 0 0] /Rect [316.807 308.203 349.649 323.277] /A << /S /GoTo /D <74686D2E442E312E312E32> >> >>}
\MusicPDFWriteObject{7280}{<< /Type /Annot /Subtype /Link /Border [0 0 0] /H /I /C [1 0 0] /Rect [304.531 269.628 336.756 281.583] /A << /S /GoTo /D <74686D2E442E312E322E31> >> >>}
\MusicPDFWriteObject{7281}{<< /Type /Annot /Subtype /Link /Border [0 0 0] /H /I /C [1 0 0] /Rect [251.289 253.148 284.624 270.624] /A << /S /GoTo /D <74686D2E442E312E312E32> >> >>}
\MusicPDFWriteObject{7282}{<< /Type /Annot /Subtype /Link /Border [0 0 0] /H /I /C [1 0 0] /Rect [151.082 223.08 157.045 236.723] /A << /S /GoTo /D <48666F6F746E6F74652E313333> >> >>}
\MusicPDFWriteObject{7283}{<< /Type /Annot /Subtype /Link /Border [0 0 0] /H /I /C [1 0 0] /Rect [182.261 458.954 219.438 470.907] /A << /S /GoTo /D <6571756174696F6E2E442E312E362E3231> >> >>}
\MusicPDFWriteObject{7284}{<< /Type /Annot /Subtype /Link /Border [0 0 0] /H /I /C [1 0 0] /Rect [318.68 396.225 356.502 408.197] /A << /S /GoTo /D <6571756174696F6E2E442E312E332E3230> >> >>}
\MusicPDFWriteObject{7285}{<< /Type /Annot /Subtype /Link /Border [0 0 0] /H /I /C [1 0 0] /Rect [262.355 346.646 295.165 358.601] /A << /S /GoTo /D <74686D2E442E312E332E34> >> >>}
\MusicPDFWriteObject{7286}{<< /Type /Annot /Subtype /Link /Border [0 0 0] /H /I /C [1 0 0] /Rect [209.252 307.194 247.074 319.149] /A << /S /GoTo /D <6571756174696F6E2E442E312E332E3130> >> >>}
\MusicPDFWriteObject{7287}{<< /Type /Annot /Subtype /Link /Border [0 0 0] /H /I /C [1 0 0] /Rect [469.317 267.742 507.174 279.695] /A << /S /GoTo /D <6571756174696F6E2E442E312E362E3231> >> >>}
\MusicPDFWriteObject{7288}{<< /Type /Annot /Subtype /Link /Border [0 0 0] /H /I /C [1 0 0] /Rect [139.296 228.184 177.763 241.384] /A << /S /GoTo /D <6571756174696F6E2E442E312E332E3230> >> >>}
\MusicPDFWriteObject{7289}{<< /Type /Annot /Subtype /Link /Border [0 0 0] /H /I /C [0 1 0] /Rect [377.704 217.524 402.764 226.392] /A << /S /GoTo /D <636974652E53636861> >> >>}
\MusicPDFWriteObject{7290}{<< /Type /Annot /Subtype /Link /Border [0 0 0] /H /I /C [1 0 0] /Rect [349.846 435.887 355.809 448.744] /A << /S /GoTo /D <48666F6F746E6F74652E313334> >> >>}
\MusicPDFWriteObject{7291}{<< /Type /Annot /Subtype /Link /Border [0 0 0] /H /I /C [1 0 0] /Rect [201.341 421.684 238.446 434.689] /A << /S /GoTo /D <74686D2E442E312E362E3139> >> >>}
\MusicPDFWriteObject{7292}{<< /Type /Annot /Subtype /Link /Border [0 0 0] /H /I /C [1 0 0] /Rect [172.238 409.585 205.08 421.538] /A << /S /GoTo /D <74686D2E442E312E322E37> >> >>}
\MusicPDFWriteObject{7293}{<< /Type /Annot /Subtype /Link /Border [0 0 0] /H /I /C [0 1 0] /Rect [221.375 412.075 239.678 420.943] /A << /S /GoTo /D <636974652E4445> >> >>}
\MusicPDFWriteObject{7294}{<< /Type /Annot /Subtype /Link /Border [0 0 0] /H /I /C [0 1 0] /Rect [344.825 141.589 363.273 150.457] /A << /S /GoTo /D <636974652E5761> >> >>}
\MusicPDFWriteObject{7295}{<< /Type /Annot /Subtype /Link /Border [0 0 0] /H /I /C [0 1 0] /Rect [154.351 128.438 172.798 137.306] /A << /S /GoTo /D <636974652E5761> >> >>}
\MusicPDFWriteObject{7296}{<< /Type /Annot /Subtype /Link /Border [0 0 0] /H /I /C [1 0 0] /Rect [449.748 125.948 482.589 137.901] /A << /S /GoTo /D <6571756174696F6E2E442E312E312E34> >> >>}
\MusicPDFWriteObject{7297}{<< /Type /Annot /Subtype /Link /Border [0 0 0] /H /I /C [1 0 0] /Rect [266.107 383.999 304.645 395.954] /A << /S /GoTo /D <6571756174696F6E2E442E312E362E3133> >> >>}
\MusicPDFWriteObject{7298}{<< /Type /Annot /Subtype /Link /Border [0 0 0] /H /I /C [1 0 0] /Rect [175.709 331.026 207.934 345.333] /A << /S /GoTo /D <74686D2E442E312E352E32> >> >>}
\MusicPDFWriteObject{7299}{<< /Type /Annot /Subtype /Link /Border [0 0 0] /H /I /C [1 0 0] /Rect [106.676 304.564 139.641 318.871] /A << /S /GoTo /D <74686D2E442E312E352E32> >> >>}
\MusicPDFWriteObject{7300}{<< /Type /Annot /Subtype /Link /Border [0 0 0] /H /I /C [1 0 0] /Rect [387.003 189.93 392.965 205.289] /A << /S /GoTo /D <48666F6F746E6F74652E313335> >> >>}
\MusicPDFWriteObject{7301}{<< /Type /Annot /Subtype /Link /Border [0 0 0] /H /I /C [1 0 0] /Rect [241.13 162.638 247.093 175.496] /A << /S /GoTo /D <48666F6F746E6F74652E313336> >> >>}
\MusicPDFWriteObject{7302}{<< /Type /Annot /Subtype /Link /Border [0 0 0] /H /I /C [1 0 0] /Rect [435.858 672.036 473.106 682.827] /A << /S /GoTo /D <74686D2E442E312E362E3136> >> >>}
\MusicPDFWriteObject{7303}{<< /Type /Annot /Subtype /Link /Border [0 0 0] /H /I /C [1 0 0] /Rect [106.072 605.358 138.45 617.311] /A << /S /GoTo /D <74686D2E442E312E312E32> >> >>}
\MusicPDFWriteObject{7304}{<< /Type /Annot /Subtype /Link /Border [0 0 0] /H /I /C [1 0 0] /Rect [221.533 579.007 259.426 590.96] /A << /S /GoTo /D <6571756174696F6E2E442E312E332E3230> >> >>}
\MusicPDFWriteObject{7305}{<< /Type /Annot /Subtype /Link /Border [0 0 0] /H /I /C [0 1 0] /Rect [303.255 528.894 328.315 537.762] /A << /S /GoTo /D <636974652E53636861> >> >>}
\MusicPDFWriteObject{7306}{<< /Type /Annot /Subtype /Link /Border [0 0 0] /H /I /C [0 1 0] /Rect [422.457 450.776 442.115 459.644] /A << /S /GoTo /D <636974652E4B47> >> >>}
\MusicPDFWriteObject{7307}{<< /Type /Annot /Subtype /Link /Border [0 0 0] /H /I /C [0 1 0] /Rect [121.939 164.912 140.912 173.78] /A << /S /GoTo /D <636974652E547265> >> >>}
\MusicPDFWriteObject{7308}{<< /Type /Annot /Subtype /Link /Border [0 0 0] /H /I /C [1 0 0] /Rect [174.928 149.272 212.75 161.224] /A << /S /GoTo /D <6571756174696F6E2E442E312E362E3235> >> >>}
\MusicPDFWriteObject{7309}{<< /Type /Annot /Subtype /Link /Border [0 0 0] /H /I /C [1 0 0] /Rect [291.238 149.272 297.201 162.129] /A << /S /GoTo /D <48666F6F746E6F74652E313337> >> >>}
\MusicPDFWriteObject{7310}{<< /Type /Annot /Subtype /Link /Border [0 0 0] /H /I /C [0 1 0] /Rect [524.857 151.762 543.159 160.629] /A << /S /GoTo /D <636974652E4445> >> >>}
\MusicPDFWriteObject{7311}{<< /Type /Annot /Subtype /Link /Border [0 0 0] /H /I /C [1 0 0] /Rect [120.446 136.121 126.409 148.978] /A << /S /GoTo /D <48666F6F746E6F74652E313338> >> >>}
\MusicPDFWriteObject{7312}{<< /Type /Annot /Subtype /Link /Border [0 0 0] /H /I /C [1 0 0] /Rect [295.919 82.097 324.129 94.681] /A << /S /GoTo /D <6571756174696F6E2E442E312E322E38> >> >>}
\MusicPDFWriteObject{7313}{<< /Type /Annot /Subtype /Link /Border [0 0 0] /H /I /C [1 0 0] /Rect [295.772 684.633 334.31 698.829] /A << /S /GoTo /D <74686D2E442E312E362E3139> >> >>}
\MusicPDFWriteObject{7314}{<< /Type /Annot /Subtype /Link /Border [0 0 0] /H /I /C [1 0 0] /Rect [162.972 645.131 201.51 657.086] /A << /S /GoTo /D <74686D2E442E312E362E3132> >> >>}
\MusicPDFWriteObject{7315}{<< /Type /Annot /Subtype /Link /Border [0 0 0] /H /I /C [1 0 0] /Rect [73.712 631.544 111.533 643.497] /A << /S /GoTo /D <6571756174696F6E2E442E312E362E3233> >> >>}
\MusicPDFWriteObject{7316}{<< /Type /Annot /Subtype /Link /Border [0 0 0] /H /I /C [1 0 0] /Rect [268.512 631.544 306.334 646.127] /A << /S /GoTo /D <74686D2E442E312E362E3135> >> >>}
\MusicPDFWriteObject{7317}{<< /Type /Annot /Subtype /Link /Border [0 0 0] /H /I /C [0 1 0] /Rect [351.708 634.034 368.862 642.902] /A << /S /GoTo /D <636974652E466F6C> >> >>}
\MusicPDFWriteObject{7318}{<< /Type /Annot /Subtype /Link /Border [0 0 0] /H /I /C [1 0 0] /Rect [229.259 591.23 262.717 603.182] /A << /S /GoTo /D <74686D2E442E312E322E37> >> >>}
\MusicPDFWriteObject{7319}{<< /Type /Annot /Subtype /Link /Border [0 0 0] /H /I /C [1 0 0] /Rect [249.795 578.632 280.403 589.423] /A << /S /GoTo /D <74686D2E342E362E322E31> >> >>}
\MusicPDFWriteObject{7320}{<< /Type /Annot /Subtype /Link /Border [0 0 0] /H /I /C [1 0 0] /Rect [327.156 578.632 357.765 589.423] /A << /S /GoTo /D <74686D2E342E362E312E35> >> >>}
\MusicPDFWriteObject{7321}{<< /Type /Annot /Subtype /Link /Border [0 0 0] /H /I /C [1 0 0] /Rect [414.681 578.632 445.289 589.423] /A << /S /GoTo /D <74686D2E342E342E322E36> >> >>}
\MusicPDFWriteObject{7322}{<< /Type /Annot /Subtype /Link /Border [0 0 0] /H /I /C [1 0 0] /Rect [193.736 567.418 223.949 576.272] /A << /S /GoTo /D <74686D2E372E332E312E36> >> >>}
\MusicPDFWriteObject{7323}{<< /Type /Annot /Subtype /Link /Border [0 0 0] /H /I /C [1 0 0] /Rect [70.444 500.802 104.392 509.656] /A << /S /GoTo /D <74686D2E442E312E322E31> >> >>}
\MusicPDFWriteObject{7324}{<< /Type /Annot /Subtype /Link /Border [0 0 0] /H /I /C [1 0 0] /Rect [120.694 448.734 153.536 460.687] /A << /S /GoTo /D <6571756174696F6E2E442E312E322E37> >> >>}
\MusicPDFWriteObject{7325}{<< /Type /Annot /Subtype /Link /Border [0 0 0] /H /I /C [1 0 0] /Rect [105.768 255.742 142.873 267.698] /A << /S /GoTo /D <74686D2E442E312E362E3138> >> >>}
\MusicPDFWriteObject{7326}{<< /Type /Annot /Subtype /Link /Border [0 0 0] /H /I /C [1 0 0] /Rect [428.589 200.953 461.369 215.305] /A << /S /GoTo /D <74686D2E442E312E342E33> >> >>}
\MusicPDFWriteObject{7327}{<< /Type /Annot /Subtype /Link /Border [0 0 0] /H /I /C [0 1 0] /Rect [514.445 660.822 543.206 669.69] /A << /S /GoTo /D <636974652E47726F2D66> >> >>}
\MusicPDFWriteObject{7328}{<< /Type /Annot /Subtype /Link /Border [0 0 0] /H /I /C [1 0 0] /Rect [182.227 632.03 215.685 643.985] /A << /S /GoTo /D <6571756174696F6E2E442E312E372E33> >> >>}
\MusicPDFWriteObject{7329}{<< /Type /Annot /Subtype /Link /Border [0 0 0] /H /I /C [1 0 0] /Rect [207.651 603.686 240.493 616.693] /A << /S /GoTo /D <74686D2E442E312E312E32> >> >>}
\MusicPDFWriteObject{7330}{<< /Type /Annot /Subtype /Link /Border [0 0 0] /H /I /C [1 0 0] /Rect [474.663 356.658 499.295 368.61] /A << /S /GoTo /D <73756273656374696F6E2E442E312E35> >> >>}
\MusicPDFWriteObject{7331}{<< /Type /Annot /Subtype /Link /Border [0 0 0] /H /I /C [0 1 0] /Rect [524.967 359.148 543.27 368.016] /A << /S /GoTo /D <636974652E4445> >> >>}
\MusicPDFWriteObject{7332}{<< /Type /Annot /Subtype /Link /Border [0 0 0] /H /I /C [1 0 0] /Rect [479.918 343.507 512.482 355.46] /A << /S /GoTo /D <74686D2E442E312E352E37> >> >>}
\MusicPDFWriteObject{7333}{<< /Type /Annot /Subtype /Link /Border [0 0 0] /H /I /C [0 1 0] /Rect [83.847 332.846 105.068 341.714] /A << /S /GoTo /D <636974652E427275> >> >>}
\MusicPDFWriteObject{7334}{<< /Type /Annot /Subtype /Link /Border [0 0 0] /H /I /C [1 0 0] /Rect [152.674 264.602 189.922 276.557] /A << /S /GoTo /D <6571756174696F6E2E442E312E332E3230> >> >>}
\MusicPDFWriteObject{7335}{<< /Type /Annot /Subtype /Link /Border [0 0 0] /H /I /C [0 1 0] /Rect [256.109 253.942 275.767 262.81] /A << /S /GoTo /D <636974652E4B47> >> >>}
\MusicPDFWriteObject{7336}{<< /Type /Annot /Subtype /Link /Border [0 0 0] /H /I /C [1 0 0] /Rect [104.394 523.91 137.236 535.863] /A << /S /GoTo /D <6571756174696F6E2E442E312E382E32> >> >>}
\MusicPDFWriteObject{7337}{<< /Type /Annot /Subtype /Link /Border [0 0 0] /H /I /C [1 0 0] /Rect [261.311 522.223 298.009 537.11] /A << /S /GoTo /D <74686D2E432E322E332E3132> >> >>}
\MusicPDFWriteObject{7338}{<< /Type /Annot /Subtype /Link /Border [0 0 0] /H /I /C [1 0 0] /Rect [71.004 510.76 103.845 522.712] /A << /S /GoTo /D <74686D2E442E312E352E37> >> >>}
\MusicPDFWriteObject{7339}{<< /Type /Annot /Subtype /Link /Border [0 0 0] /H /I /C [0 1 0] /Rect [334.064 513.25 359.864 522.117] /A << /S /GoTo /D <636974652E416C2D42> >> >>}
\MusicPDFWriteObject{7340}{<< /Type /Annot /Subtype /Link /Border [0 0 0] /H /I /C [0 1 0] /Rect [237.321 467.445 262.381 476.313] /A << /S /GoTo /D <636974652E53636861> >> >>}
\MusicPDFWriteObject{7341}{<< /Type /Annot /Subtype /Link /Border [0 0 0] /H /I /C [1 0 0] /Rect [130.101 451.805 162.942 463.757] /A << /S /GoTo /D <74686D2E442E312E332E37> >> >>}
\MusicPDFWriteObject{7342}{<< /Type /Annot /Subtype /Link /Border [0 0 0] /H /I /C [1 0 0] /Rect [137.094 399.202 170.429 411.154] /A << /S /GoTo /D <74686D2E442E312E352E31> >> >>}
\MusicPDFWriteObject{7343}{<< /Type /Annot /Subtype /Link /Border [0 0 0] /H /I /C [1 0 0] /Rect [273.992 359.749 312.107 371.702] /A << /S /GoTo /D <74686D2E432E322E332E3132> >> >>}
\MusicPDFWriteObject{7344}{<< /Type /Annot /Subtype /Link /Border [0 0 0] /H /I /C [1 0 0] /Rect [445.431 359.749 483.546 371.702] /A << /S /GoTo /D <74686D2E432E322E332E3133> >> >>}
\MusicPDFWriteObject{7345}{<< /Type /Annot /Subtype /Link /Border [0 0 0] /H /I /C [1 0 0] /Rect [163.302 320.297 200.284 332.25] /A << /S /GoTo /D <74686D2E432E322E332E3132> >> >>}
\MusicPDFWriteObject{7346}{<< /Type /Annot /Subtype /Link /Border [0 0 0] /H /I /C [0 1 0] /Rect [73.732 296.486 92.705 305.353] /A << /S /GoTo /D <636974652E547265> >> >>}
\MusicPDFWriteObject{7347}{<< /Type /Annot /Subtype /Link /Border [0 0 0] /H /I /C [0 1 0] /Rect [524.28 296.486 543.253 305.353] /A << /S /GoTo /D <636974652E547265> >> >>}
\MusicPDFWriteObject{7348}{<< /Type /Annot /Subtype /Link /Border [0 0 0] /H /I /C [1 0 0] /Rect [404.548 241.032 437.583 253.347] /A << /S /GoTo /D <74686D2E432E312E362E31> >> >>}
\MusicPDFWriteObject{7349}{<< /Type /Annot /Subtype /Link /Border [0 0 0] /H /I /C [0 1 0] /Rect [107.806 111.387 126.005 120.255] /A << /S /GoTo /D <636974652E4A6172> >> >>}
\MusicPDFWriteObject{7350}{<< /Type /Annot /Subtype /Link /Border [0 0 0] /H /I /C [1 0 0] /Rect [336.328 82.596 369.786 94.549] /A << /S /GoTo /D <6571756174696F6E2E442E312E382E33> >> >>}
\MusicPDFWriteObject{7351}{<< /Type /Annot /Subtype /Link /Border [0 0 0] /H /I /C [1 0 0] /Rect [71.004 671.482 109.119 683.435] /A << /S /GoTo /D <74686D2E432E322E332E3132> >> >>}
\MusicPDFWriteObject{7352}{<< /Type /Annot /Subtype /Link /Border [0 0 0] /H /I /C [0 1 0] /Rect [524.186 673.973 543.159 682.84] /A << /S /GoTo /D <636974652E547265> >> >>}
\MusicPDFWriteObject{7353}{<< /Type /Annot /Subtype /Link /Border [0 0 0] /H /I /C [1 0 0] /Rect [86.286 645.181 119.282 657.134] /A << /S /GoTo /D <6571756174696F6E2E442E312E382E33> >> >>}
\MusicPDFWriteObject{7354}{<< /Type /Annot /Subtype /Link /Border [0 0 0] /H /I /C [1 0 0] /Rect [258.798 645.181 291.793 657.134] /A << /S /GoTo /D <6571756174696F6E2E442E312E382E33> >> >>}
\MusicPDFWriteObject{7355}{<< /Type /Annot /Subtype /Link /Border [0 0 0] /H /I /C [0 1 0] /Rect [73.77 634.52 98.831 643.388] /A << /S /GoTo /D <636974652E53636861> >> >>}
\MusicPDFWriteObject{7356}{<< /Type /Annot /Subtype /Link /Border [0 0 0] /H /I /C [0 1 0] /Rect [188.464 504.011 219.41 512.879] /A << /S /GoTo /D <636974652E486F726D> >> >>}
\MusicPDFWriteObject{7357}{<< /Type /Annot /Subtype /Link /Border [0 0 0] /H /I /C [0 1 0] /Rect [411.687 151.2 431.345 160.067] /A << /S /GoTo /D <636974652E4B47> >> >>}
\MusicPDFWriteObject{7358}{<< /Type /Annot /Subtype /Link /Border [0 0 0] /H /I /C [1 0 0] /Rect [73.789 109.257 112.327 121.212] /A << /S /GoTo /D <6571756174696F6E2E442E312E332E3230> >> >>}
\MusicPDFWriteObject{7359}{<< /Type /Annot /Subtype /Link /Border [0 0 0] /H /I /C [1 0 0] /Rect [419.862 531.724 452.087 545.717] /A << /S /GoTo /D <74686D2E442E312E352E31> >> >>}
\MusicPDFWriteObject{7360}{<< /Type /Annot /Subtype /Link /Border [0 0 0] /H /I /C [1 0 0] /Rect [470.209 519.625 503.05 531.58] /A << /S /GoTo /D <74686D2E442E312E352E33> >> >>}
\MusicPDFWriteObject{7361}{<< /Type /Annot /Subtype /Link /Border [0 0 0] /H /I /C [1 0 0] /Rect [203.451 342.092 236.91 352.94] /A << /S /GoTo /D <74686D2E442E312E342E37> >> >>}
\MusicPDFWriteObject{7362}{<< /Type /Annot /Subtype /Link /Border [0 0 0] /H /I /C [1 0 0] /Rect [515.944 281.862 545.592 292.667] /A << /S /GoTo /D <74686D2E342E332E332E31> >> >>}
\MusicPDFWriteObject{7363}{<< /Type /Annot /Subtype /Link /Border [0 0 0] /H /I /C [0 1 0] /Rect [257.434 137.689 282.494 146.557] /A << /S /GoTo /D <636974652E53636861> >> >>}
\MusicPDFWriteObject{7364}{<< /Type /Annot /Subtype /Link /Border [0 0 0] /H /I /C [0 1 0] /Rect [326.271 660.822 351.331 669.69] /A << /S /GoTo /D <636974652E53636861> >> >>}
\MusicPDFWriteObject{7365}{<< /Type /Annot /Subtype /Link /Border [0 0 0] /H /I /C [1 0 0] /Rect [335.532 450.436 368.99 462.392] /A << /S /GoTo /D <74686D2E442E312E352E31> >> >>}
\MusicPDFWriteObject{7366}{<< /Type /Annot /Subtype /Link /Border [0 0 0] /H /I /C [1 0 0] /Rect [149.018 424.135 181.243 436.09] /A << /S /GoTo /D <74686D2E442E312E392E31> >> >>}
\MusicPDFWriteObject{7367}{<< /Type /Annot /Subtype /Link /Border [0 0 0] /H /I /C [0 1 0] /Rect [234.405 400.324 259.465 409.192] /A << /S /GoTo /D <636974652E53636861> >> >>}
\MusicPDFWriteObject{7368}{<< /Type /Annot /Subtype /Link /Border [0 0 0] /H /I /C [1 0 0] /Rect [502.369 371.532 539.474 383.485] /A << /S /GoTo /D <6571756174696F6E2E442E312E332E3230> >> >>}
\MusicPDFWriteObject{7369}{<< /Type /Annot /Subtype /Link /Border [0 0 0] /H /I /C [1 0 0] /Rect [270.097 158.233 303.463 170.188] /A << /S /GoTo /D <6571756174696F6E2E442E312E332E38> >> >>}
\MusicPDFWriteObject{7370}{<< /Type /Annot /Subtype /Link /Border [0 0 0] /H /I /C [1 0 0] /Rect [295.537 419.335 328.378 431.291] /A << /S /GoTo /D <6571756174696F6E2E442E312E392E32> >> >>}
\MusicPDFWriteObject{7371}{<< /Type /Annot /Subtype /Link /Border [0 0 0] /H /I /C [1 0 0] /Rect [224.533 650.335 262.355 662.288] /A << /S /GoTo /D <6571756174696F6E2E442E312E31302E31> >> >>}
\MusicPDFWriteObject{7372}{<< /Type /Annot /Subtype /Link /Border [0 0 0] /H /I /C [1 0 0] /Rect [235.566 451.869 265.497 463.822] /A << /S /GoTo /D <74686D2E322E312E312E34> >> >>}
\MusicPDFWriteObject{7373}{<< /Type /Annot /Subtype /Link /Border [0 0 0] /H /I /C [0 1 0] /Rect [465.212 241.488 490.272 250.356] /A << /S /GoTo /D <636974652E53636861> >> >>}
\MusicPDFWriteObject{7374}{<< /Type /Annot /Subtype /Link /Border [0 0 0] /H /I /C [0 1 0] /Rect [446.296 202.036 463.45 210.904] /A << /S /GoTo /D <636974652E466F6C> >> >>}
\MusicPDFWriteObject{7375}{<< /Type /Annot /Subtype /Link /Border [0 0 0] /H /I /C [1 0 0] /Rect [502.52 82.595 532.733 94.55] /A << /S /GoTo /D <74686D2E322E312E312E34> >> >>}
\MusicPDFWriteObject{7376}{<< /Type /Annot /Subtype /Link /Border [0 0 0] /H /I /C [1 0 0] /Rect [450.319 368.079 482.667 380.034] /A << /S /GoTo /D <74686D2E442E312E332E31> >> >>}
\MusicPDFWriteObject{7377}{<< /Type /Annot /Subtype /Link /Border [0 0 0] /H /I /C [1 0 0] /Rect [511.675 313.277 541.888 325.23] /A << /S /GoTo /D <6571756174696F6E2E322E312E312E33> >> >>}
\MusicPDFWriteObject{7378}{<< /Type /Annot /Subtype /Link /Border [0 0 0] /H /I /C [1 0 0] /Rect [253.204 300.127 285.737 312.079] /A << /S /GoTo /D <6571756174696F6E2E442E312E332E38> >> >>}
\MusicPDFWriteObject{7379}{<< /Type /Annot /Subtype /Link /Border [0 0 0] /H /I /C [0 1 0] /Rect [303.436 289.466 323.094 298.334] /A << /S /GoTo /D <636974652E4B47> >> >>}
\MusicPDFWriteObject{7380}{<< /Type /Annot /Subtype /Link /Border [0 0 0] /H /I /C [0 1 0] /Rect [233.67 201.368 253.329 210.236] /A << /S /GoTo /D <636974652E4B47> >> >>}
\MusicPDFWriteObject{7381}{<< /Type /Annot /Subtype /Link /Border [0 0 0] /H /I /C [1 0 0] /Rect [117.346 155.058 150.188 167.011] /A << /S /GoTo /D <74686D2E442E312E332E31> >> >>}
\MusicPDFWriteObject{7382}{<< /Type /Annot /Subtype /Link /Border [0 0 0] /H /I /C [0 1 0] /Rect [280.146 598.093 304.107 606.961] /A << /S /GoTo /D <636974652E57726E> >> >>}
\MusicPDFWriteObject{7383}{<< /Type /Annot /Subtype /Link /Border [0 0 0] /H /I /C [1 0 0] /Rect [430.828 304.842 472.814 316.797] /A << /S /GoTo /D <6571756174696F6E2E442E312E31302E3131> >> >>}
\MusicPDFWriteObject{7384}{<< /Type /Annot /Subtype /Link /Border [0 0 0] /H /I /C [1 0 0] /Rect [445.745 259.17 478.093 271.126] /A << /S /GoTo /D <74686D2E442E312E332E31> >> >>}
\MusicPDFWriteObject{7385}{<< /Type /Annot /Subtype /Link /Border [0 0 0] /H /I /C [1 0 0] /Rect [302.541 120.435 339.897 131.226] /A << /S /GoTo /D <74686D2E442E312E31302E33> >> >>}
\MusicPDFWriteObject{7386}{<< /Type /Annot /Subtype /Link /Border [0 0 0] /H /I /C [1 0 0] /Rect [71.004 619.365 103.845 631.318] /A << /S /GoTo /D <74686D2E442E312E332E34> >> >>}
\MusicPDFWriteObject{7387}{<< /Type /Annot /Subtype /Link /Border [0 0 0] /H /I /C [0 1 0] /Rect [400.939 621.855 419.387 630.723] /A << /S /GoTo /D <636974652E5761> >> >>}
\MusicPDFWriteObject{7388}{<< /Type /Annot /Subtype /Link /Border [0 0 0] /H /I /C [1 0 0] /Rect [71.004 567.247 108.825 579.202] /A << /S /GoTo /D <74686D2E442E312E362E3139> >> >>}
\MusicPDFWriteObject{7389}{<< /Type /Annot /Subtype /Link /Border [0 0 0] /H /I /C [1 0 0] /Rect [142.114 453.223 180.652 465.178] /A << /S /GoTo /D <74686D2E442E312E362E3136> >> >>}
\MusicPDFWriteObject{7390}{<< /Type /Annot /Subtype /Link /Border [0 0 0] /H /I /C [1 0 0] /Rect [71.004 440.072 108.825 452.045] /A << /S /GoTo /D <74686D2E442E312E362E3230> >> >>}
\MusicPDFWriteObject{7391}{<< /Type /Annot /Subtype /Link /Border [0 0 0] /H /I /C [1 0 0] /Rect [144.917 262.889 182.022 274.844] /A << /S /GoTo /D <74686D2E442E312E362E3132> >> >>}
\MusicPDFWriteObject{7392}{<< /Type /Annot /Subtype /Link /Border [0 0 0] /H /I /C [1 0 0] /Rect [237.293 262.889 274.398 274.844] /A << /S /GoTo /D <74686D2E442E312E362E3133> >> >>}
\MusicPDFWriteObject{7393}{<< /Type /Annot /Subtype /Link /Border [0 0 0] /H /I /C [1 0 0] /Rect [353.414 213.105 386.873 225.06] /A << /S /GoTo /D <6571756174696F6E2E442E312E332E39> >> >>}
\MusicPDFWriteObject{7394}{<< /Type /Annot /Subtype /Link /Border [0 0 0] /H /I /C [1 0 0] /Rect [281.076 158.524 318.897 169.315] /A << /S /GoTo /D <74686D2E442E312E31302E33> >> >>}
\MusicPDFWriteObject{7395}{<< /Type /Annot /Subtype /Link /Border [0 0 0] /H /I /C [1 0 0] /Rect [321.885 158.524 359.706 169.315] /A << /S /GoTo /D <74686D2E442E312E31302E35> >> >>}
\MusicPDFWriteObject{7396}{<< /Type /Annot /Subtype /Link /Border [0 0 0] /H /I /C [1 0 0] /Rect [469.57 684.633 507.464 697.574] /A << /S /GoTo /D <74686D2E442E312E31302E33> >> >>}
\MusicPDFWriteObject{7397}{<< /Type /Annot /Subtype /Link /Border [0 0 0] /H /I /C [1 0 0] /Rect [343.805 562.893 382.164 576.554] /A << /S /GoTo /D <74686D2E442E312E31302E34> >> >>}
\MusicPDFWriteObject{7398}{<< /Type /Annot /Subtype /Link /Border [0 0 0] /H /I /C [1 0 0] /Rect [441.574 551.071 484.376 563.384] /A << /S /GoTo /D <6571756174696F6E2E442E312E31302E3133> >> >>}
\MusicPDFWriteObject{7399}{<< /Type /Annot /Subtype /Link /Border [0 0 0] /H /I /C [1 0 0] /Rect [167.815 396.934 206.174 408.352] /A << /S /GoTo /D <74686D2E442E312E31302E35> >> >>}
\MusicPDFWriteObject{7400}{<< /Type /Annot /Subtype /Link /Border [0 0 0] /H /I /C [1 0 0] /Rect [291.661 335.164 334.463 347.117] /A << /S /GoTo /D <6571756174696F6E2E442E312E31302E3133> >> >>}
\MusicPDFWriteObject{7401}{<< /Type /Annot /Subtype /Link /Border [0 0 0] /H /I /C [1 0 0] /Rect [130.361 124.29 163.357 136.243] /A << /S /GoTo /D <74686D2E442E312E332E37> >> >>}
\MusicPDFWriteObject{7402}{<< /Type /Annot /Subtype /Link /Border [0 0 0] /H /I /C [1 0 0] /Rect [242.555 124.29 280.556 136.243] /A << /S /GoTo /D <74686D2E442E312E31302E33> >> >>}
\MusicPDFWriteObject{7403}{<< /Type /Annot /Subtype /Link /Border [0 0 0] /H /I /C [1 0 0] /Rect [283.568 124.29 321.569 136.243] /A << /S /GoTo /D <74686D2E442E312E31302E35> >> >>}
\MusicPDFWriteObject{7404}{<< /Type /Annot /Subtype /Link /Border [0 0 0] /H /I /C [0 1 0] /Rect [253.381 620.783 278.441 629.651] /A << /S /GoTo /D <636974652E53636861> >> >>}
\MusicPDFWriteObject{7405}{<< /Type /Annot /Subtype /Link /Border [0 0 0] /H /I /C [1 0 0] /Rect [330.844 425.42 363.547 439.623] /A << /S /GoTo /D <74686D2E412E312E342E33> >> >>}
\MusicPDFWriteObject{7406}{<< /Type /Annot /Subtype /Link /Border [0 0 0] /H /I /C [1 0 0] /Rect [409.734 256.9 447.985 269.853] /A << /S /GoTo /D <74686D2E442E312E31302E37> >> >>}
\MusicPDFWriteObject{7407}{<< /Type /Annot /Subtype /Link /Border [0 0 0] /H /I /C [0 1 0] /Rect [109.199 160.966 126.865 169.834] /A << /S /GoTo /D <636974652E4154> >> >>}
\MusicPDFWriteObject{7408}{<< /Type /Annot /Subtype /Link /Border [0 0 0] /H /I /C [1 0 0] /Rect [248.518 132.174 280.866 144.129] /A << /S /GoTo /D <74686D2E442E312E382E31> >> >>}
\MusicPDFWriteObject{7409}{<< /Type /Annot /Subtype /Link /Border [0 0 0] /H /I /C [1 0 0] /Rect [448.461 684.633 486.576 696.92] /A << /S /GoTo /D <74686D2E432E322E332E3133> >> >>}
\MusicPDFWriteObject{7410}{<< /Type /Annot /Subtype /Link /Border [0 0 0] /H /I /C [1 0 0] /Rect [495.461 513.673 533.576 526.266] /A << /S /GoTo /D <74686D2E432E322E332E3132> >> >>}
\MusicPDFWriteObject{7411}{<< /Type /Annot /Subtype /Link /Border [0 0 0] /H /I /C [1 0 0] /Rect [135.911 316.82 173.733 327.61] /A << /S /GoTo /D <74686D2E442E312E31302E33> >> >>}
\MusicPDFWriteObject{7412}{<< /Type /Annot /Subtype /Link /Border [0 0 0] /H /I /C [1 0 0] /Rect [456.369 273.909 493.725 286.755] /A << /S /GoTo /D <74686D2E442E312E31302E37> >> >>}
\MusicPDFWriteObject{7413}{<< /Type /Annot /Subtype /Link /Border [0 0 0] /H /I /C [0 1 0] /Rect [167.119 660.822 184.786 669.69] /A << /S /GoTo /D <636974652E4154> >> >>}
\MusicPDFWriteObject{7414}{<< /Type /Annot /Subtype /Link /Border [0 0 0] /H /I /C [1 0 0] /Rect [71.004 417.541 109.542 430.134] /A << /S /GoTo /D <74686D2E442E312E31302E36> >> >>}
\MusicPDFWriteObject{7415}{<< /Type /Annot /Subtype /Link /Border [0 0 0] /H /I /C [1 0 0] /Rect [418.618 338.994 461.419 350.946] /A << /S /GoTo /D <6571756174696F6E2E442E312E31302E3137> >> >>}
\MusicPDFWriteObject{7416}{<< /Type /Annot /Subtype /Link /Border [0 0 0] /H /I /C [1 0 0] /Rect [233.643 325.843 276.444 337.796] /A << /S /GoTo /D <6571756174696F6E2E442E312E31302E3137> >> >>}
\MusicPDFWriteObject{7417}{<< /Type /Annot /Subtype /Link /Border [0 0 0] /H /I /C [1 0 0] /Rect [71.004 262.56 108.825 274.264] /A << /S /GoTo /D <74686D2E442E312E31302E39> >> >>}
\MusicPDFWriteObject{7418}{<< /Type /Annot /Subtype /Link /Border [0 0 0] /H /I /C [1 0 0] /Rect [477.471 213.372 514.612 227.575] /A << /S /GoTo /D <74686D2E442E312E31302E37> >> >>}
\MusicPDFWriteObject{7419}{<< /Type /Annot /Subtype /Link /Border [0 0 0] /H /I /C [1 0 0] /Rect [206.449 650.742 250.067 661.532] /A << /S /GoTo /D <74686D2E442E312E31302E3130> >> >>}
\MusicPDFWriteObject{7420}{<< /Type /Annot /Subtype /Link /Border [0 0 0] /H /I /C [1 0 0] /Rect [380.511 650.742 419.05 661.532] /A << /S /GoTo /D <74686D2E442E312E31302E39> >> >>}
\MusicPDFWriteObject{7421}{<< /Type /Annot /Subtype /Link /Border [0 0 0] /H /I /C [1 0 0] /Rect [211.955 635.427 250.493 649.63] /A << /S /GoTo /D <74686D2E442E312E31302E31> >> >>}
\MusicPDFWriteObject{7422}{<< /Type /Annot /Subtype /Link /Border [0 0 0] /H /I /C [1 0 0] /Rect [448.326 637.038 491.128 648.99] /A << /S /GoTo /D <6571756174696F6E2E442E312E31302E3230> >> >>}
\MusicPDFWriteObject{7423}{<< /Type /Annot /Subtype /Link /Border [0 0 0] /H /I /C [1 0 0] /Rect [154.434 623.887 197.236 635.839] /A << /S /GoTo /D <6571756174696F6E2E442E312E31302E3137> >> >>}
\MusicPDFWriteObject{7424}{<< /Type /Annot /Subtype /Link /Border [0 0 0] /H /I /C [1 0 0] /Rect [71.004 610.736 108.825 622.689] /A << /S /GoTo /D <74686D2E442E312E31302E31> >> >>}
\MusicPDFWriteObject{7425}{<< /Type /Annot /Subtype /Link /Border [0 0 0] /H /I /C [1 0 0] /Rect [71.004 471.479 104.462 485.681] /A << /S /GoTo /D <74686D2E442E312E392E32> >> >>}
\MusicPDFWriteObject{7426}{<< /Type /Annot /Subtype /Link /Border [0 0 0] /H /I /C [0 1 0] /Rect [333.537 449.278 367.098 458.145] /A << /S /GoTo /D <636974652E5765696C2D65> >> >>}
\MusicPDFWriteObject{7427}{<< /Type /Annot /Subtype /Link /Border [0 0 0] /H /I /C [1 0 0] /Rect [183.278 304.262 215.656 316.217] /A << /S /GoTo /D <74686D2E442E312E352E37> >> >>}
\MusicPDFWriteObject{7428}{<< /Type /Annot /Subtype /Link /Border [0 0 0] /H /I /C [1 0 0] /Rect [157.807 291.112 190.649 303.067] /A << /S /GoTo /D <74686D2E442E312E392E32> >> >>}
\MusicPDFWriteObject{7429}{<< /Type /Annot /Subtype /Link /Border [0 0 0] /H /I /C [1 0 0] /Rect [223.409 143.881 253.368 156.473] /A << /S /GoTo /D <74686D2E342E332E332E31> >> >>}
\MusicPDFWriteObject{7430}{<< /Type /Annot /Subtype /Link /Border [0 0 0] /H /I /C [1 0 0] /Rect [307.249 684.273 340.492 696.588] /A << /S /GoTo /D <6571756174696F6E2E442E322E312E32> >> >>}
\MusicPDFWriteObject{7431}{<< /Type /Annot /Subtype /Link /Border [0 0 0] /H /I /C [1 0 0] /Rect [345.845 671.482 376.058 684.075] /A << /S /GoTo /D <74686D2E342E332E332E31> >> >>}
\MusicPDFWriteObject{7432}{<< /Type /Annot /Subtype /Link /Border [0 0 0] /H /I /C [1 0 0] /Rect [155.562 486.293 187.787 498.246] /A << /S /GoTo /D <6571756174696F6E2E442E322E312E34> >> >>}
\MusicPDFWriteObject{7433}{<< /Type /Annot /Subtype /Link /Border [0 0 0] /H /I /C [0 1 0] /Rect [156.765 388.351 173.712 397.219] /A << /S /GoTo /D <636974652E5275> >> >>}
\MusicPDFWriteObject{7434}{<< /Type /Annot /Subtype /Link /Border [0 0 0] /H /I /C [0 1 0] /Rect [242.806 362.049 276.367 370.917] /A << /S /GoTo /D <636974652E5765696C2D65> >> >>}
\MusicPDFWriteObject{7435}{<< /Type /Annot /Subtype /Link /Border [0 0 0] /H /I /C [1 0 0] /Rect [327.899 172.523 360.895 185.407] /A << /S /GoTo /D <6571756174696F6E2E442E322E312E35> >> >>}
\MusicPDFWriteObject{7436}{<< /Type /Annot /Subtype /Link /Border [0 0 0] /H /I /C [1 0 0] /Rect [368.024 159.372 401.421 171.328] /A << /S /GoTo /D <6571756174696F6E2E442E322E312E35> >> >>}
\MusicPDFWriteObject{7437}{<< /Type /Annot /Subtype /Link /Border [0 0 0] /H /I /C [1 0 0] /Rect [255.347 96.415 288.188 107.206] /A << /S /GoTo /D <74686D2E442E322E312E33> >> >>}
\MusicPDFWriteObject{7438}{<< /Type /Annot /Subtype /Link /Border [0 0 0] /H /I /C [0 1 0] /Rect [128.69 544.019 146.163 552.887] /A << /S /GoTo /D <636974652E4346> >> >>}
\MusicPDFWriteObject{7439}{<< /Type /Annot /Subtype /Link /Border [0 0 0] /H /I /C [0 1 0] /Rect [92.581 530.868 110.053 539.736] /A << /S /GoTo /D <636974652E4346> >> >>}
\MusicPDFWriteObject{7440}{<< /Type /Annot /Subtype /Link /Border [0 0 0] /H /I /C [1 0 0] /Rect [503.518 506.518 533.251 519.111] /A << /S /GoTo /D <74686D2E342E332E332E31> >> >>}
\MusicPDFWriteObject{7441}{<< /Type /Annot /Subtype /Link /Border [0 0 0] /H /I /C [1 0 0] /Rect [70.655 417.488 101.376 430.08] /A << /S /GoTo /D <74686D2E342E332E332E31> >> >>}
\MusicPDFWriteObject{7442}{<< /Type /Annot /Subtype /Link /Border [0 0 0] /H /I /C [1 0 0] /Rect [512.621 338.583 544.845 350.538] /A << /S /GoTo /D <74686D2E442E312E392E32> >> >>}
\MusicPDFWriteObject{7443}{<< /Type /Annot /Subtype /Link /Border [0 0 0] /H /I /C [1 0 0] /Rect [158.576 95.746 188.796 108.339] /A << /S /GoTo /D <74686D2E342E332E332E31> >> >>}
\MusicPDFWriteObject{7444}{<< /Type /Annot /Subtype /Link /Border [0 0 0] /H /I /C [1 0 0] /Rect [364.293 621.904 397.135 633.857] /A << /S /GoTo /D <6571756174696F6E2E442E322E322E35> >> >>}
\MusicPDFWriteObject{7445}{<< /Type /Annot /Subtype /Link /Border [0 0 0] /H /I /C [1 0 0] /Rect [189.875 510.15 195.838 522.88] /A << /S /GoTo /D <48666F6F746E6F74652E313339> >> >>}
\MusicPDFWriteObject{7446}{<< /Type /Annot /Subtype /Link /Border [0 0 0] /H /I /C [1 0 0] /Rect [192.878 281.007 223.091 293.599] /A << /S /GoTo /D <6571756174696F6E2E342E332E332E32> >> >>}
\MusicPDFWriteObject{7447}{<< /Type /Annot /Subtype /Link /Border [0 0 0] /H /I /C [1 0 0] /Rect [231.059 233.956 263.912 245.909] /A << /S /GoTo /D <74686D2E432E342E332E38> >> >>}
\MusicPDFWriteObject{7448}{<< /Type /Annot /Subtype /Link /Border [0 0 0] /H /I /C [1 0 0] /Rect [243.005 205.718 273.726 220.192] /A << /S /GoTo /D <74686D2E342E332E332E31> >> >>}
\MusicPDFWriteObject{7449}{<< /Type /Annot /Subtype /Link /Border [0 0 0] /H /I /C [1 0 0] /Rect [176.627 120.079 210.086 132.034] /A << /S /GoTo /D <74686D2E442E322E322E31> >> >>}
\MusicPDFWriteObject{7450}{<< /Type /Annot /Subtype /Link /Border [0 0 0] /H /I /C [1 0 0] /Rect [346.592 631.139 379.434 643.092] /A << /S /GoTo /D <6571756174696F6E2E442E322E322E36> >> >>}
\MusicPDFWriteObject{7451}{<< /Type /Annot /Subtype /Link /Border [0 0 0] /H /I /C [1 0 0] /Rect [205.742 565.939 237.682 576.729] /A << /S /GoTo /D <73756273756273656374696F6E2E432E322E332E34> >> >>}
\MusicPDFWriteObject{7452}{<< /Type /Annot /Subtype /Link /Border [0 0 0] /H /I /C [1 0 0] /Rect [338.347 565.939 375.187 576.729] /A << /S /GoTo /D <74686D2E432E322E332E3132> >> >>}
\MusicPDFWriteObject{7453}{<< /Type /Annot /Subtype /Link /Border [0 0 0] /H /I /C [1 0 0] /Rect [448.467 565.939 485.307 576.729] /A << /S /GoTo /D <74686D2E432E322E332E3133> >> >>}
\MusicPDFWriteObject{7454}{<< /Type /Annot /Subtype /Link /Border [0 0 0] /H /I /C [1 0 0] /Rect [261.921 367.796 292.699 382.27] /A << /S /GoTo /D <74686D2E342E332E332E31> >> >>}
\MusicPDFWriteObject{7455}{<< /Type /Annot /Subtype /Link /Border [0 0 0] /H /I /C [1 0 0] /Rect [312.075 148.35 345.534 160.302] /A << /S /GoTo /D <74686D2E442E312E382E31> >> >>}
\MusicPDFWriteObject{7456}{<< /Type /Annot /Subtype /Link /Border [0 0 0] /H /I /C [1 0 0] /Rect [375.925 95.386 406.703 107.701] /A << /S /GoTo /D <6571756174696F6E2E342E332E332E32> >> >>}
\MusicPDFWriteObject{7457}{<< /Type /Annot /Subtype /Link /Border [0 0 0] /H /I /C [1 0 0] /Rect [347.06 671.122 380.518 684.34] /A << /S /GoTo /D <74686D2E442E322E322E31> >> >>}
\MusicPDFWriteObject{7458}{<< /Type /Annot /Subtype /Link /Border [0 0 0] /H /I /C [1 0 0] /Rect [454.828 671.122 464.761 684.34] /A << /S /GoTo /D <48666F6F746E6F74652E313430> >> >>}
\MusicPDFWriteObject{7459}{<< /Type /Annot /Subtype /Link /Border [0 0 0] /H /I /C [1 0 0] /Rect [93.423 618.519 123.269 631.472] /A << /S /GoTo /D <6571756174696F6E2E342E332E332E32> >> >>}
\MusicPDFWriteObject{7460}{<< /Type /Annot /Subtype /Link /Border [0 0 0] /H /I /C [1 0 0] /Rect [372.782 532.114 406.241 542.99] /A << /S /GoTo /D <74686D2E442E322E322E32> >> >>}
\MusicPDFWriteObject{7461}{<< /Type /Annot /Subtype /Link /Border [0 0 0] /H /I /C [1 0 0] /Rect [185.868 485.052 195.801 497.442] /A << /S /GoTo /D <48666F6F746E6F74652E313431> >> >>}
\MusicPDFWriteObject{7462}{<< /Type /Annot /Subtype /Link /Border [0 0 0] /H /I /C [1 0 0] /Rect [196.141 354.8 228.983 365.677] /A << /S /GoTo /D <74686D2E442E322E322E32> >> >>}
\MusicPDFWriteObject{7463}{<< /Type /Annot /Subtype /Link /Border [0 0 0] /H /I /C [1 0 0] /Rect [174.785 247.111 207.073 259.064] /A << /S /GoTo /D <74686D2E422E332E322E31> >> >>}
\MusicPDFWriteObject{7464}{<< /Type /Annot /Subtype /Link /Border [0 0 0] /H /I /C [1 0 0] /Rect [176.666 530.733 209.508 542.686] /A << /S /GoTo /D <6571756174696F6E2E442E332E312E32> >> >>}
\MusicPDFWriteObject{7465}{<< /Type /Annot /Subtype /Link /Border [0 0 0] /H /I /C [1 0 0] /Rect [377.201 298.82 384.174 309.611] /A << /S /GoTo /D <6669677572652E442E31> >> >>}
\MusicPDFWriteObject{7466}{<< /Type /Annot /Subtype /Link /Border [0 0 0] /H /I /C [1 0 0] /Rect [140.342 242.253 173.184 254.206] /A << /S /GoTo /D <6571756174696F6E2E442E332E312E32> >> >>}
\MusicPDFWriteObject{7467}{<< /Type /Annot /Subtype /Link /Border [0 0 0] /H /I /C [1 0 0] /Rect [73.712 671.482 106.553 683.435] /A << /S /GoTo /D <6571756174696F6E2E442E332E312E35> >> >>}
\MusicPDFWriteObject{7468}{<< /Type /Annot /Subtype /Link /Border [0 0 0] /H /I /C [0 1 0] /Rect [516.307 512.316 543.242 521.184] /A << /S /GoTo /D <636974652E484B5A> >> >>}
\MusicPDFWriteObject{7469}{<< /Type /Annot /Subtype /Link /Border [0 0 0] /H /I /C [1 0 0] /Rect [451.437 158.726 483.725 171.869] /A << /S /GoTo /D <74686D2E422E322E322E34> >> >>}
\MusicPDFWriteObject{7470}{<< /Type /Annot /Subtype /Link /Border [0 0 0] /H /I /C [1 0 0] /Rect [495.648 81.101 528.49 95.629] /A << /S /GoTo /D <6571756174696F6E2E442E332E312E38> >> >>}
\MusicPDFWriteObject{7471}{<< /Type /Annot /Subtype /Link /Border [0 0 0] /H /I /C [1 0 0] /Rect [296.966 610.915 329.345 622.868] /A << /S /GoTo /D <6571756174696F6E2E442E332E312E39> >> >>}
\MusicPDFWriteObject{7472}{<< /Type /Annot /Subtype /Link /Border [0 0 0] /H /I /C [1 0 0] /Rect [314.039 312.567 346.881 324.52] /A << /S /GoTo /D <6571756174696F6E2E442E332E312E37> >> >>}
\MusicPDFWriteObject{7473}{<< /Type /Annot /Subtype /Link /Border [0 0 0] /H /I /C [1 0 0] /Rect [120.758 286.265 153.538 299.129] /A << /S /GoTo /D <6571756174696F6E2E442E332E312E38> >> >>}
\MusicPDFWriteObject{7474}{<< /Type /Annot /Subtype /Link /Border [0 0 0] /H /I /C [1 0 0] /Rect [464.938 199.238 497.533 212.179] /A << /S /GoTo /D <74686D2E442E332E312E31> >> >>}
\MusicPDFWriteObject{7475}{<< /Type /Annot /Subtype /Link /Border [0 0 0] /H /I /C [1 0 0] /Rect [260.22 133.484 298.042 144.829] /A << /S /GoTo /D <74686D2E442E312E31302E33> >> >>}
\MusicPDFWriteObject{7476}{<< /Type /Annot /Subtype /Link /Border [0 0 0] /H /I /C [1 0 0] /Rect [431.087 647.908 463.928 660.16] /A << /S /GoTo /D <74686D2E442E332E312E32> >> >>}
\MusicPDFWriteObject{7477}{<< /Type /Annot /Subtype /Link /Border [0 0 0] /H /I /C [1 0 0] /Rect [205.203 628.014 237.613 639.969] /A << /S /GoTo /D <74686D2E442E332E312E32> >> >>}
\MusicPDFWriteObject{7478}{<< /Type /Annot /Subtype /Link /Border [0 0 0] /H /I /C [0 1 0] /Rect [325.922 498.067 350.982 506.935] /A << /S /GoTo /D <636974652E53636861> >> >>}
\MusicPDFWriteObject{7479}{<< /Type /Annot /Subtype /Link /Border [0 0 0] /H /I /C [0 1 0] /Rect [434.072 184.176 453.045 193.044] /A << /S /GoTo /D <636974652E547265> >> >>}
\MusicPDFWriteObject{7480}{<< /Type /Annot /Subtype /Link /Border [0 0 0] /H /I /C [0 1 0] /Rect [238.747 171.025 263.807 179.893] /A << /S /GoTo /D <636974652E53636861> >> >>}
\MusicPDFWriteObject{7481}{<< /Type /Annot /Subtype /Link /Border [0 0 0] /H /I /C [0 1 0] /Rect [350.139 124.538 375.199 133.406] /A << /S /GoTo /D <636974652E53636861> >> >>}
\MusicPDFWriteObject{7482}{<< /Type /Annot /Subtype /Link /Border [0 0 0] /H /I /C [0 1 0] /Rect [129.438 463.56 154.498 472.428] /A << /S /GoTo /D <636974652E53636861> >> >>}
\MusicPDFWriteObject{7483}{<< /Type /Annot /Subtype /Link /Border [0 0 0] /H /I /C [0 1 0] /Rect [379.345 463.56 398.319 472.428] /A << /S /GoTo /D <636974652E547265> >> >>}
\MusicPDFWriteObject{7484}{<< /Type /Annot /Subtype /Link /Border [0 0 0] /H /I /C [0 1 0] /Rect [123.433 450.41 148.493 459.277] /A << /S /GoTo /D <636974652E53636861> >> >>}
\MusicPDFWriteObject{7485}{<< /Type /Annot /Subtype /Link /Border [0 0 0] /H /I /C [1 0 0] /Rect [431.051 408.467 463.339 421.325] /A << /S /GoTo /D <74686D2E422E322E312E31> >> >>}
\MusicPDFWriteObject{7486}{<< /Type /Annot /Subtype /Link /Border [0 0 0] /H /I /C [1 0 0] /Rect [467.987 408.467 477.92 421.325] /A << /S /GoTo /D <48666F6F746E6F74652E313432> >> >>}
\MusicPDFWriteObject{7487}{<< /Type /Annot /Subtype /Link /Border [0 0 0] /H /I /C [1 0 0] /Rect [281.946 378.22 314.17 391.103] /A << /S /GoTo /D <74686D2E442E332E322E33> >> >>}
\MusicPDFWriteObject{7488}{<< /Type /Annot /Subtype /Link /Border [0 0 0] /H /I /C [0 1 0] /Rect [185.473 354.409 204.446 363.277] /A << /S /GoTo /D <636974652E547265> >> >>}
\MusicPDFWriteObject{7489}{<< /Type /Annot /Subtype /Link /Border [0 0 0] /H /I /C [0 1 0] /Rect [280.666 341.258 299.639 350.126] /A << /S /GoTo /D <636974652E547265> >> >>}
\MusicPDFWriteObject{7490}{<< /Type /Annot /Subtype /Link /Border [0 0 0] /H /I /C [0 1 0] /Rect [456.707 341.258 481.767 350.126] /A << /S /GoTo /D <636974652E53636861> >> >>}
\MusicPDFWriteObject{7491}{<< /Type /Annot /Subtype /Link /Border [0 0 0] /H /I /C [1 0 0] /Rect [197.138 229.813 229.274 241.769] /A << /S /GoTo /D <74686D2E422E322E322E33> >> >>}
\MusicPDFWriteObject{7492}{<< /Type /Annot /Subtype /Link /Border [0 0 0] /H /I /C [1 0 0] /Rect [371.123 203.513 403.964 215.465] /A << /S /GoTo /D <6571756174696F6E2E442E332E322E32> >> >>}
\MusicPDFWriteObject{7493}{<< /Type /Annot /Subtype /Link /Border [0 0 0] /H /I /C [1 0 0] /Rect [187.367 177.21 219.564 189.166] /A << /S /GoTo /D <74686D2E422E322E312E38> >> >>}
\MusicPDFWriteObject{7494}{<< /Type /Annot /Subtype /Link /Border [0 0 0] /H /I /C [1 0 0] /Rect [143.67 164.614 177.129 176.918] /A << /S /GoTo /D <74686D2E442E332E312E31> >> >>}
\MusicPDFWriteObject{7495}{<< /Type /Annot /Subtype /Link /Border [0 0 0] /H /I /C [1 0 0] /Rect [534.862 164.614 544.795 176.918] /A << /S /GoTo /D <48666F6F746E6F74652E313433> >> >>}
\MusicPDFWriteObject{7496}{<< /Type /Annot /Subtype /Link /Border [0 0 0] /H /I /C [1 0 0] /Rect [143.651 671.482 177.11 683.437] /A << /S /GoTo /D <74686D2E442E332E312E33> >> >>}
\MusicPDFWriteObject{7497}{<< /Type /Annot /Subtype /Link /Border [0 0 0] /H /I /C [1 0 0] /Rect [446.994 671.482 480.452 683.437] /A << /S /GoTo /D <74686D2E442E332E322E33> >> >>}
\MusicPDFWriteObject{7498}{<< /Type /Annot /Subtype /Link /Border [0 0 0] /H /I /C [1 0 0] /Rect [507.055 606.057 539.897 618.009] /A << /S /GoTo /D <6571756174696F6E2E442E332E322E32> >> >>}
\MusicPDFWriteObject{7499}{<< /Type /Annot /Subtype /Link /Border [0 0 0] /H /I /C [1 0 0] /Rect [482.591 553.781 516.049 565.736] /A << /S /GoTo /D <74686D2E442E332E322E33> >> >>}
\MusicPDFWriteObject{7500}{<< /Type /Annot /Subtype /Link /Border [0 0 0] /H /I /C [1 0 0] /Rect [123.714 622.168 157.172 634.12] /A << /S /GoTo /D <74686D2E442E332E312E32> >> >>}
\MusicPDFWriteObject{7501}{<< /Type /Annot /Subtype /Link /Border [0 0 0] /H /I /C [1 0 0] /Rect [259.757 622.168 293.215 634.12] /A << /S /GoTo /D <74686D2E442E332E312E34> >> >>}
\MusicPDFWriteObject{7502}{<< /Type /Annot /Subtype /Link /Border [0 0 0] /H /I /C [1 0 0] /Rect [369.802 622.168 403.26 634.12] /A << /S /GoTo /D <74686D2E442E332E322E35> >> >>}
\MusicPDFWriteObject{7503}{<< /Type /Annot /Subtype /Link /Border [0 0 0] /H /I /C [1 0 0] /Rect [71.004 330.118 103.845 342.711] /A << /S /GoTo /D <74686D2E442E332E322E32> >> >>}
\MusicPDFWriteObject{7504}{<< /Type /Annot /Subtype /Link /Border [0 0 0] /H /I /C [1 0 0] /Rect [369.415 228.382 402.843 240.337] /A << /S /GoTo /D <74686D2E442E332E322E31> >> >>}
\MusicPDFWriteObject{7505}{<< /Type /Annot /Subtype /Link /Border [0 0 0] /H /I /C [1 0 0] /Rect [320.001 214.18 352.626 227.187] /A << /S /GoTo /D <74686D2E442E332E322E33> >> >>}
\MusicPDFWriteObject{7506}{<< /Type /Annot /Subtype /Link /Border [0 0 0] /H /I /C [1 0 0] /Rect [310.513 671.482 343.972 683.435] /A << /S /GoTo /D <74686D2E442E332E322E33> >> >>}
\MusicPDFWriteObject{7507}{<< /Type /Annot /Subtype /Link /Border [0 0 0] /H /I /C [1 0 0] /Rect [309.76 658.332 342.386 670.284] /A << /S /GoTo /D <6571756174696F6E2E442E312E332E38> >> >>}
\MusicPDFWriteObject{7508}{<< /Type /Annot /Subtype /Link /Border [0 0 0] /H /I /C [1 0 0] /Rect [428.573 381.544 462.032 393.499] /A << /S /GoTo /D <74686D2E442E332E322E31> >> >>}
\MusicPDFWriteObject{7509}{<< /Type /Annot /Subtype /Link /Border [0 0 0] /H /I /C [1 0 0] /Rect [107.895 282.824 141.353 293.615] /A << /S /GoTo /D <74686D2E442E332E332E32> >> >>}
\MusicPDFWriteObject{7510}{<< /Type /Annot /Subtype /Link /Border [0 0 0] /H /I /C [1 0 0] /Rect [71.004 255.969 103.845 267.922] /A << /S /GoTo /D <74686D2E442E332E322E33> >> >>}
\MusicPDFWriteObject{7511}{<< /Type /Annot /Subtype /Link /Border [0 0 0] /H /I /C [0 1 0] /Rect [356.255 219.007 381.316 227.875] /A << /S /GoTo /D <636974652E53636861> >> >>}
\MusicPDFWriteObject{7512}{<< /Type /Annot /Subtype /Link /Border [0 0 0] /H /I /C [1 0 0] /Rect [362.913 190.215 401.451 202.17] /A << /S /GoTo /D <6571756174696F6E2E442E312E332E3230> >> >>}
\MusicPDFWriteObject{7513}{<< /Type /Annot /Subtype /Link /Border [0 0 0] /H /I /C [1 0 0] /Rect [237.051 177.064 246.984 189.922] /A << /S /GoTo /D <48666F6F746E6F74652E313434> >> >>}
\MusicPDFWriteObject{7514}{<< /Type /Annot /Subtype /Link /Border [0 0 0] /H /I /C [1 0 0] /Rect [105.71 104.584 137.935 116.539] /A << /S /GoTo /D <74686D2E442E312E332E37> >> >>}
\MusicPDFWriteObject{7515}{<< /Type /Annot /Subtype /Link /Border [0 0 0] /H /I /C [0 1 0] /Rect [436.084 687.123 461.144 695.991] /A << /S /GoTo /D <636974652E53636861> >> >>}
\MusicPDFWriteObject{7516}{<< /Type /Annot /Subtype /Link /Border [0 0 0] /H /I /C [1 0 0] /Rect [490.464 626.273 522.751 638.228] /A << /S /GoTo /D <74686D2E442E332E332E31> >> >>}
\MusicPDFWriteObject{7517}{<< /Type /Annot /Subtype /Link /Border [0 0 0] /H /I /C [1 0 0] /Rect [305.394 298.423 315.327 310.727] /A << /S /GoTo /D <48666F6F746E6F74652E313435> >> >>}
\MusicPDFWriteObject{7518}{<< /Type /Annot /Subtype /Link /Border [0 0 0] /H /I /C [1 0 0] /Rect [121.226 156.56 154.068 168.513] /A << /S /GoTo /D <6571756174696F6E2E442E332E332E33> >> >>}
\MusicPDFWriteObject{7519}{<< /Type /Annot /Subtype /Link /Border [0 0 0] /H /I /C [1 0 0] /Rect [204.013 143.409 241.419 155.364] /A << /S /GoTo /D <74686D2E432E322E332E3132> >> >>}
\MusicPDFWriteObject{7520}{<< /Type /Annot /Subtype /Link /Border [0 0 0] /H /I /C [1 0 0] /Rect [210.945 585.989 248.351 596.78] /A << /S /GoTo /D <74686D2E432E322E332E3132> >> >>}
\MusicPDFWriteObject{7521}{<< /Type /Annot /Subtype /Link /Border [0 0 0] /H /I /C [1 0 0] /Rect [322.665 585.989 360.072 596.78] /A << /S /GoTo /D <74686D2E432E322E332E3133> >> >>}
\MusicPDFWriteObject{7522}{<< /Type /Annot /Subtype /Link /Border [0 0 0] /H /I /C [1 0 0] /Rect [418.906 565.481 443.132 577.433] /A << /S /GoTo /D <73756273656374696F6E2E432E312E35> >> >>}
\MusicPDFWriteObject{7523}{<< /Type /Annot /Subtype /Link /Border [0 0 0] /H /I /C [1 0 0] /Rect [445.701 526.028 483.108 537.981] /A << /S /GoTo /D <74686D2E432E312E352E3133> >> >>}
\MusicPDFWriteObject{7524}{<< /Type /Annot /Subtype /Link /Border [0 0 0] /H /I /C [1 0 0] /Rect [154.008 513.431 186.325 524.222] /A << /S /GoTo /D <74686D2E442E312E352E39> >> >>}
\MusicPDFWriteObject{7525}{<< /Type /Annot /Subtype /Link /Border [0 0 0] /H /I /C [1 0 0] /Rect [497.27 473.425 530.728 486.017] /A << /S /GoTo /D <74686D2E442E332E332E37> >> >>}
\MusicPDFWriteObject{7526}{<< /Type /Annot /Subtype /Link /Border [0 0 0] /H /I /C [0 1 0] /Rect [248.764 449.614 277.525 458.482] /A << /S /GoTo /D <636974652E47726F2D66> >> >>}
\MusicPDFWriteObject{7527}{<< /Type /Annot /Subtype /Link /Border [0 0 0] /H /I /C [1 0 0] /Rect [73.712 433.973 106.553 445.926] /A << /S /GoTo /D <6571756174696F6E2E442E332E332E33> >> >>}
\MusicPDFWriteObject{7528}{<< /Type /Annot /Subtype /Link /Border [0 0 0] /H /I /C [1 0 0] /Rect [206.241 414.017 239.699 425.972] /A << /S /GoTo /D <74686D2E442E332E332E37> >> >>}
\MusicPDFWriteObject{7529}{<< /Type /Annot /Subtype /Link /Border [0 0 0] /H /I /C [1 0 0] /Rect [529.26 387.716 539.193 400.573] /A << /S /GoTo /D <48666F6F746E6F74652E313436> >> >>}
\MusicPDFWriteObject{7530}{<< /Type /Annot /Subtype /Link /Border [0 0 0] /H /I /C [1 0 0] /Rect [416.418 154.318 449.259 166.273] /A << /S /GoTo /D <74686D2E442E332E332E35> >> >>}
\MusicPDFWriteObject{7531}{<< /Type /Annot /Subtype /Link /Border [0 0 0] /H /I /C [0 1 0] /Rect [122.938 687.123 147.998 695.991] /A << /S /GoTo /D <636974652E53636861> >> >>}
\MusicPDFWriteObject{7532}{<< /Type /Annot /Subtype /Link /Border [0 0 0] /H /I /C [0 1 0] /Rect [220.592 687.123 245.652 695.991] /A << /S /GoTo /D <636974652E53636861> >> >>}
\MusicPDFWriteObject{7533}{<< /Type /Annot /Subtype /Link /Border [0 0 0] /H /I /C [1 0 0] /Rect [420.858 632.03 454.317 643.983] /A << /S /GoTo /D <6571756174696F6E2E442E332E332E36> >> >>}
\MusicPDFWriteObject{7534}{<< /Type /Annot /Subtype /Link /Border [0 0 0] /H /I /C [0 1 0] /Rect [301.471 621.369 333.095 630.237] /A << /S /GoTo /D <636974652E47726F2D62> >> >>}
\MusicPDFWriteObject{7535}{<< /Type /Annot /Subtype /Link /Border [0 0 0] /H /I /C [1 0 0] /Rect [228.126 579.427 265.533 591.38] /A << /S /GoTo /D <74686D2E432E322E332E3133> >> >>}
\MusicPDFWriteObject{7536}{<< /Type /Annot /Subtype /Link /Border [0 0 0] /H /I /C [1 0 0] /Rect [258.437 543.553 295.171 554.344] /A << /S /GoTo /D <74686D2E432E322E332E3132> >> >>}
\MusicPDFWriteObject{7537}{<< /Type /Annot /Subtype /Link /Border [0 0 0] /H /I /C [1 0 0] /Rect [215.618 503.547 249.077 515.502] /A << /S /GoTo /D <6571756174696F6E2E442E332E332E36> >> >>}
\MusicPDFWriteObject{7538}{<< /Type /Annot /Subtype /Link /Border [0 0 0] /H /I /C [1 0 0] /Rect [350.1 503.547 383.558 515.502] /A << /S /GoTo /D <74686D2E442E332E322E33> >> >>}
\MusicPDFWriteObject{7539}{<< /Type /Annot /Subtype /Link /Border [0 0 0] /H /I /C [0 1 0] /Rect [391.562 479.736 416.622 488.604] /A << /S /GoTo /D <636974652E53636861> >> >>}
\MusicPDFWriteObject{7540}{<< /Type /Annot /Subtype /Link /Border [0 0 0] /H /I /C [0 1 0] /Rect [445.595 466.585 464.568 475.453] /A << /S /GoTo /D <636974652E547265> >> >>}
\MusicPDFWriteObject{7541}{<< /Type /Annot /Subtype /Link /Border [0 0 0] /H /I /C [0 1 0] /Rect [139.665 277.171 158.77 286.038] /A << /S /GoTo /D <636974652E4E42> >> >>}
\MusicPDFWriteObject{7542}{<< /Type /Annot /Subtype /Link /Border [0 0 0] /H /I /C [0 1 0] /Rect [203.47 237.718 228.53 246.586] /A << /S /GoTo /D <636974652E53636861> >> >>}
\MusicPDFWriteObject{7543}{<< /Type /Annot /Subtype /Link /Border [0 0 0] /H /I /C [0 1 0] /Rect [319.147 224.567 344.207 233.435] /A << /S /GoTo /D <636974652E53636861> >> >>}
\MusicPDFWriteObject{7544}{<< /Type /Annot /Subtype /Link /Border [0 0 0] /H /I /C [0 1 0] /Rect [305.934 184.847 330.995 193.715] /A << /S /GoTo /D <636974652E53636861> >> >>}
\MusicPDFWriteObject{7545}{<< /Type /Annot /Subtype /Link /Border [0 0 0] /H /I /C [0 1 0] /Rect [73.826 171.696 98.886 180.564] /A << /S /GoTo /D <636974652E53636861> >> >>}
\MusicPDFWriteObject{7546}{<< /Type /Annot /Subtype /Link /Border [0 0 0] /H /I /C [0 1 0] /Rect [259.102 171.696 278.207 180.564] /A << /S /GoTo /D <636974652E4E42> >> >>}
\MusicPDFWriteObject{7547}{<< /Type /Annot /Subtype /Link /Border [0 0 0] /H /I /C [0 1 0] /Rect [208.96 92.527 234.02 101.394] /A << /S /GoTo /D <636974652E53636861> >> >>}
\MusicPDFWriteObject{7548}{<< /Type /Annot /Subtype /Link /Border [0 0 0] /H /I /C [0 1 0] /Rect [232.757 673.973 257.817 682.84] /A << /S /GoTo /D <636974652E53636861> >> >>}
\MusicPDFWriteObject{7549}{<< /Type /Annot /Subtype /Link /Border [0 0 0] /H /I /C [1 0 0] /Rect [208.013 618.879 240.855 630.832] /A << /S /GoTo /D <6571756174696F6E2E442E332E332E36> >> >>}
\MusicPDFWriteObject{7550}{<< /Type /Annot /Subtype /Link /Border [0 0 0] /H /I /C [1 0 0] /Rect [272.531 605.729 305.373 617.681] /A << /S /GoTo /D <6571756174696F6E2E442E332E332E36> >> >>}
\MusicPDFWriteObject{7551}{<< /Type /Annot /Subtype /Link /Border [0 0 0] /H /I /C [1 0 0] /Rect [268.625 579.427 301.467 591.38] /A << /S /GoTo /D <6571756174696F6E2E442E332E332E36> >> >>}
\MusicPDFWriteObject{7552}{<< /Type /Annot /Subtype /Link /Border [0 0 0] /H /I /C [1 0 0] /Rect [362.075 566.276 394.917 578.229] /A << /S /GoTo /D <6571756174696F6E2E442E332E332E36> >> >>}
\MusicPDFWriteObject{7553}{<< /Type /Annot /Subtype /Link /Border [0 0 0] /H /I /C [1 0 0] /Rect [299.721 553.125 309.654 566.326] /A << /S /GoTo /D <48666F6F746E6F74652E313437> >> >>}
\MusicPDFWriteObject{7554}{<< /Type /Annot /Subtype /Link /Border [0 0 0] /H /I /C [0 1 0] /Rect [333.185 529.314 351.384 538.182] /A << /S /GoTo /D <636974652E4A6172> >> >>}
\MusicPDFWriteObject{7555}{<< /Type /Annot /Subtype /Link /Border [0 0 0] /H /I /C [1 0 0] /Rect [268.133 514.227 278.066 526.531] /A << /S /GoTo /D <48666F6F746E6F74652E313438> >> >>}
\MusicPDFWriteObject{7556}{<< /Type /Annot /Subtype /Link /Border [0 0 0] /H /I /C [0 1 0] /Rect [318.874 503.013 347.634 511.881] /A << /S /GoTo /D <636974652E47726F2D66> >> >>}
\MusicPDFWriteObject{7557}{<< /Type /Annot /Subtype /Link /Border [0 0 0] /H /I /C [0 1 0] /Rect [328.171 461.029 346.37 469.897] /A << /S /GoTo /D <636974652E4A6172> >> >>}
\MusicPDFWriteObject{7558}{<< /Type /Annot /Subtype /Link /Border [0 0 0] /H /I /C [1 0 0] /Rect [385.393 445.387 395.326 458.329] /A << /S /GoTo /D <48666F6F746E6F74652E313439> >> >>}
\MusicPDFWriteObject{7559}{<< /Type /Annot /Subtype /Link /Border [0 0 0] /H /I /C [1 0 0] /Rect [73.712 405.936 106.553 417.889] /A << /S /GoTo /D <6571756174696F6E2E442E332E332E36> >> >>}
\MusicPDFWriteObject{7560}{<< /Type /Annot /Subtype /Link /Border [0 0 0] /H /I /C [0 1 0] /Rect [524.109 395.275 543.214 404.143] /A << /S /GoTo /D <636974652E4E42> >> >>}
\MusicPDFWriteObject{7561}{<< /Type /Annot /Subtype /Link /Border [0 0 0] /H /I /C [0 1 0] /Rect [245.024 382.124 270.084 390.992] /A << /S /GoTo /D <636974652E53636861> >> >>}
\MusicPDFWriteObject{7562}{<< /Type /Annot /Subtype /Link /Border [0 0 0] /H /I /C [0 1 0] /Rect [182.055 355.823 201.16 364.691] /A << /S /GoTo /D <636974652E4E42> >> >>}
\MusicPDFWriteObject{7563}{<< /Type /Annot /Subtype /Link /Border [0 0 0] /H /I /C [0 1 0] /Rect [171.294 303.22 189.493 312.088] /A << /S /GoTo /D <636974652E4A6172> >> >>}
\MusicPDFWriteObject{7564}{<< /Type /Annot /Subtype /Link /Border [0 0 0] /H /I /C [1 0 0] /Rect [395.311 300.73 405.244 313.587] /A << /S /GoTo /D <48666F6F746E6F74652E313530> >> >>}
\MusicPDFWriteObject{7565}{<< /Type /Annot /Subtype /Link /Border [0 0 0] /H /I /C [0 1 0] /Rect [169.195 173.952 194.255 182.82] /A << /S /GoTo /D <636974652E53636861> >> >>}
\MusicPDFWriteObject{7566}{<< /Type /Annot /Subtype /Link /Border [0 0 0] /H /I /C [1 0 0] /Rect [390.577 684.633 423.419 696.586] /A << /S /GoTo /D <6571756174696F6E2E442E332E332E36> >> >>}
\MusicPDFWriteObject{7567}{<< /Type /Annot /Subtype /Link /Border [0 0 0] /H /I /C [0 1 0] /Rect [107.002 647.671 125.201 656.539] /A << /S /GoTo /D <636974652E4A6172> >> >>}
\MusicPDFWriteObject{7568}{<< /Type /Annot /Subtype /Link /Border [0 0 0] /H /I /C [0 1 0] /Rect [184.173 563.412 203.147 572.28] /A << /S /GoTo /D <636974652E547265> >> >>}
\MusicPDFWriteObject{7569}{<< /Type /Annot /Subtype /Link /Border [0 0 0] /H /I /C [1 0 0] /Rect [104.924 523.959 114.857 534.327] /A << /S /GoTo /D <48666F6F746E6F74652E313531> >> >>}
\MusicPDFWriteObject{7570}{<< /Type /Annot /Subtype /Link /Border [0 0 0] /H /I /C [0 1 0] /Rect [455.968 484.507 481.028 493.375] /A << /S /GoTo /D <636974652E53636861> >> >>}
\MusicPDFWriteObject{7571}{<< /Type /Annot /Subtype /Link /Border [0 0 0] /H /I /C [1 0 0] /Rect [187.794 459.774 224.108 468.628] /A << /S /GoTo /D <74686D2E442E332E332E37> >> >>}
\MusicPDFWriteObject{7572}{<< /Type /Annot /Subtype /Link /Border [0 0 0] /H /I /C [1 0 0] /Rect [382.656 417.831 414.88 430.424] /A << /S /GoTo /D <6571756174696F6E2E442E332E332E37> >> >>}
\MusicPDFWriteObject{7573}{<< /Type /Annot /Subtype /Link /Border [0 0 0] /H /I /C [1 0 0] /Rect [245.387 404.68 283.501 417.881] /A << /S /GoTo /D <74686D2E432E322E332E3132> >> >>}
\MusicPDFWriteObject{7574}{<< /Type /Annot /Subtype /Link /Border [0 0 0] /H /I /C [1 0 0] /Rect [102.137 378.38 134.454 390.332] /A << /S /GoTo /D <6571756174696F6E2E442E332E332E33> >> >>}
\MusicPDFWriteObject{7575}{<< /Type /Annot /Subtype /Link /Border [0 0 0] /H /I /C [0 1 0] /Rect [334.866 367.719 359.926 376.587] /A << /S /GoTo /D <636974652E53636861> >> >>}
\MusicPDFWriteObject{7576}{<< /Type /Annot /Subtype /Link /Border [0 0 0] /H /I /C [1 0 0] /Rect [165.62 352.077 198.863 364.033] /A << /S /GoTo /D <74686D2E442E332E332E35> >> >>}
\MusicPDFWriteObject{7577}{<< /Type /Annot /Subtype /Link /Border [0 0 0] /H /I /C [0 1 0] /Rect [151.243 341.417 176.303 350.285] /A << /S /GoTo /D <636974652E53636861> >> >>}
\MusicPDFWriteObject{7578}{<< /Type /Annot /Subtype /Link /Border [0 0 0] /H /I /C [0 1 0] /Rect [362.231 341.417 381.204 350.285] /A << /S /GoTo /D <636974652E547265> >> >>}
\MusicPDFWriteObject{7579}{<< /Type /Annot /Subtype /Link /Border [0 0 0] /H /I /C [0 1 0] /Rect [73.77 301.965 98.831 310.833] /A << /S /GoTo /D <636974652E53636861> >> >>}
\MusicPDFWriteObject{7580}{<< /Type /Annot /Subtype /Link /Border [0 0 0] /H /I /C [0 1 0] /Rect [376.877 234.984 401.937 243.852] /A << /S /GoTo /D <636974652E53636861> >> >>}
\MusicPDFWriteObject{7581}{<< /Type /Annot /Subtype /Link /Border [0 0 0] /H /I /C [0 1 0] /Rect [232.627 221.833 257.687 230.701] /A << /S /GoTo /D <636974652E53636861> >> >>}
\MusicPDFWriteObject{7582}{<< /Type /Annot /Subtype /Link /Border [0 0 0] /H /I /C [1 0 0] /Rect [274.275 204.505 306.499 218.165] /A << /S /GoTo /D <74686D2E442E312E352E37> >> >>}
\MusicPDFWriteObject{7583}{<< /Type /Annot /Subtype /Link /Border [0 0 0] /H /I /C [0 1 0] /Rect [470.008 195.532 488.982 204.399] /A << /S /GoTo /D <636974652E547265> >> >>}
\MusicPDFWriteObject{7584}{<< /Type /Annot /Subtype /Link /Border [0 0 0] /H /I /C [1 0 0] /Rect [236.877 82.097 265.087 93.687] /A << /S /GoTo /D <6571756174696F6E2E442E332E332E37> >> >>}
\MusicPDFWriteObject{7585}{<< /Type /Annot /Subtype /Link /Border [0 0 0] /H /I /C [1 0 0] /Rect [108.724 645.181 142.09 657.134] /A << /S /GoTo /D <6571756174696F6E2E442E332E332E33> >> >>}
\MusicPDFWriteObject{7586}{<< /Type /Annot /Subtype /Link /Border [0 0 0] /H /I /C [1 0 0] /Rect [394.26 645.181 427.626 657.134] /A << /S /GoTo /D <6571756174696F6E2E442E312E382E31> >> >>}
\MusicPDFWriteObject{7587}{<< /Type /Annot /Subtype /Link /Border [0 0 0] /H /I /C [1 0 0] /Rect [348.838 595.603 381.68 607.555] /A << /S /GoTo /D <74686D2E442E332E322E33> >> >>}
\MusicPDFWriteObject{7588}{<< /Type /Annot /Subtype /Link /Border [0 0 0] /H /I /C [1 0 0] /Rect [344.47 556.15 382.585 568.103] /A << /S /GoTo /D <74686D2E432E322E332E3134> >> >>}
\MusicPDFWriteObject{7589}{<< /Type /Annot /Subtype /Link /Border [0 0 0] /H /I /C [1 0 0] /Rect [497.777 511.697 530.927 524.897] /A << /S /GoTo /D <6571756174696F6E2E442E332E332E36> >> >>}
\MusicPDFWriteObject{7590}{<< /Type /Annot /Subtype /Link /Border [0 0 0] /H /I /C [1 0 0] /Rect [73.789 448.968 107.248 460.921] /A << /S /GoTo /D <6571756174696F6E2E442E332E332E38> >> >>}
\MusicPDFWriteObject{7591}{<< /Type /Annot /Subtype /Link /Border [0 0 0] /H /I /C [0 1 0] /Rect [142.161 425.157 167.221 434.025] /A << /S /GoTo /D <636974652E53636861> >> >>}
\MusicPDFWriteObject{7592}{<< /Type /Annot /Subtype /Link /Border [0 0 0] /H /I /C [1 0 0] /Rect [187.123 383.768 224.53 394.558] /A << /S /GoTo /D <74686D2E432E322E332E3132> >> >>}
\MusicPDFWriteObject{7593}{<< /Type /Annot /Subtype /Link /Border [0 0 0] /H /I /C [0 1 0] /Rect [434.525 240.699 453.499 249.567] /A << /S /GoTo /D <636974652E547265> >> >>}
\MusicPDFWriteObject{7594}{<< /Type /Annot /Subtype /Link /Border [0 0 0] /H /I /C [0 1 0] /Rect [252.598 227.548 284.222 236.416] /A << /S /GoTo /D <636974652E47726F2D62> >> >>}
\MusicPDFWriteObject{7595}{<< /Type /Annot /Subtype /Link /Border [0 0 0] /H /I /C [1 0 0] /Rect [309.563 211.908 342.404 223.86] /A << /S /GoTo /D <6571756174696F6E2E442E332E332E38> >> >>}
\MusicPDFWriteObject{7596}{<< /Type /Annot /Subtype /Link /Border [0 0 0] /H /I /C [1 0 0] /Rect [106.556 172.455 139.398 184.408] /A << /S /GoTo /D <74686D2E442E312E382E31> >> >>}
\MusicPDFWriteObject{7597}{<< /Type /Annot /Subtype /Link /Border [0 0 0] /H /I /C [1 0 0] /Rect [320.347 591.348 353.189 603.3] /A << /S /GoTo /D <6571756174696F6E2E442E322E312E31> >> >>}
\MusicPDFWriteObject{7598}{<< /Type /Annot /Subtype /Link /Border [0 0 0] /H /I /C [1 0 0] /Rect [203.197 526.08 245.016 534.934] /A << /S /GoTo /D <74686D2E442E332E332E3132> >> >>}
\MusicPDFWriteObject{7599}{<< /Type /Annot /Subtype /Link /Border [0 0 0] /H /I /C [1 0 0] /Rect [73.712 381.957 111.533 393.909] /A << /S /GoTo /D <6571756174696F6E2E442E332E332E3130> >> >>}
\MusicPDFWriteObject{7600}{<< /Type /Annot /Subtype /Link /Border [0 0 0] /H /I /C [1 0 0] /Rect [263.817 295.328 296.658 306.204] /A << /S /GoTo /D <74686D2E442E332E322E35> >> >>}
\MusicPDFWriteObject{7601}{<< /Type /Annot /Subtype /Link /Border [0 0 0] /H /I /C [1 0 0] /Rect [298.051 171.979 330.893 183.931] /A << /S /GoTo /D <6571756174696F6E2E442E332E332E33> >> >>}
\MusicPDFWriteObject{7602}{<< /Type /Annot /Subtype /Link /Border [0 0 0] /H /I /C [1 0 0] /Rect [100.011 684.633 133.469 696.588] /A << /S /GoTo /D <6571756174696F6E2E442E332E332E33> >> >>}
\MusicPDFWriteObject{7603}{<< /Type /Annot /Subtype /Link /Border [0 0 0] /H /I /C [1 0 0] /Rect [202.476 608.753 235.873 620.706] /A << /S /GoTo /D <6571756174696F6E2E442E332E332E35> >> >>}
\MusicPDFWriteObject{7604}{<< /Type /Annot /Subtype /Link /Border [0 0 0] /H /I /C [0 1 0] /Rect [217.723 598.093 242.783 606.961] /A << /S /GoTo /D <636974652E53636861> >> >>}
\MusicPDFWriteObject{7605}{<< /Type /Annot /Subtype /Link /Border [0 0 0] /H /I /C [1 0 0] /Rect [403.581 318.869 441.402 331.461] /A << /S /GoTo /D <6571756174696F6E2E442E312E31302E38> >> >>}
\MusicPDFWriteObject{7606}{<< /Type /Annot /Subtype /Link /Border [0 0 0] /H /I /C [1 0 0] /Rect [370.909 215.176 408.014 228.923] /A << /S /GoTo /D <6571756174696F6E2E442E332E332E3133> >> >>}
\MusicPDFWriteObject{7607}{<< /Type /Annot /Subtype /Link /Border [0 0 0] /H /I /C [1 0 0] /Rect [295.441 117.002 328.9 128.957] /A << /S /GoTo /D <6571756174696F6E2E442E332E322E32> >> >>}
\MusicPDFWriteObject{7608}{<< /Type /Annot /Subtype /Link /Border [0 0 0] /H /I /C [1 0 0] /Rect [443.054 103.851 476.512 115.806] /A << /S /GoTo /D <6571756174696F6E2E442E312E312E34> >> >>}
\MusicPDFWriteObject{7609}{<< /Type /Annot /Subtype /Link /Border [0 0 0] /H /I /C [1 0 0] /Rect [126.186 432.992 162.885 444.947] /A << /S /GoTo /D <74686D2E432E322E332E3133> >> >>}
\MusicPDFWriteObject{7610}{<< /Type /Annot /Subtype /Link /Border [0 0 0] /H /I /C [1 0 0] /Rect [363.77 392.168 402.308 404.76] /A << /S /GoTo /D <6571756174696F6E2E442E332E332E3134> >> >>}
\MusicPDFWriteObject{7611}{<< /Type /Annot /Subtype /Link /Border [0 0 0] /H /I /C [1 0 0] /Rect [415.879 148.452 454.417 163.49] /A << /S /GoTo /D <74686D2E442E312E362E3139> >> >>}
\MusicPDFWriteObject{7612}{<< /Type /Annot /Subtype /Link /Border [0 0 0] /H /I /C [1 0 0] /Rect [198.328 135.301 231.169 147.256] /A << /S /GoTo /D <6571756174696F6E2E442E332E322E32> >> >>}
\MusicPDFWriteObject{7613}{<< /Type /Annot /Subtype /Link /Border [0 0 0] /H /I /C [1 0 0] /Rect [182.015 671.482 218.713 684.075] /A << /S /GoTo /D <74686D2E432E322E332E3132> >> >>}
\MusicPDFWriteObject{7614}{<< /Type /Annot /Subtype /Link /Border [0 0 0] /H /I /C [0 1 0] /Rect [235.462 521.565 260.522 530.433] /A << /S /GoTo /D <636974652E53636861> >> >>}
\MusicPDFWriteObject{7615}{<< /Type /Annot /Subtype /Link /Border [0 0 0] /H /I /C [1 0 0] /Rect [71.004 341.584 103.292 352.432] /A << /S /GoTo /D <74686D2E422E322E322E33> >> >>}
\MusicPDFWriteObject{7616}{<< /Type /Annot /Subtype /Link /Border [0 0 0] /H /I /C [1 0 0] /Rect [212.904 341.584 237.999 352.432] /A << /S /GoTo /D <73756273656374696F6E2E442E312E36> >> >>}
\MusicPDFWriteObject{7617}{<< /Type /Annot /Subtype /Link /Border [0 0 0] /H /I /C [1 0 0] /Rect [473.258 327.88 511.079 339.833] /A << /S /GoTo /D <74686D2E442E312E362E3139> >> >>}
\MusicPDFWriteObject{7618}{<< /Type /Annot /Subtype /Link /Border [0 0 0] /H /I /C [1 0 0] /Rect [245.695 191.921 277.983 203.874] /A << /S /GoTo /D <74686D2E422E322E322E33> >> >>}
\MusicPDFWriteObject{7619}{<< /Type /Annot /Subtype /Link /Border [0 0 0] /H /I /C [1 0 0] /Rect [500.518 621.854 539.056 636.051] /A << /S /GoTo /D <74686D2E442E312E362E3130> >> >>}
\MusicPDFWriteObject{7620}{<< /Type /Annot /Subtype /Link /Border [0 0 0] /H /I /C [0 1 0] /Rect [320.799 276.042 345.859 284.91] /A << /S /GoTo /D <636974652E53636861> >> >>}
\MusicPDFWriteObject{7621}{<< /Type /Annot /Subtype /Link /Border [0 0 0] /H /I /C [1 0 0] /Rect [121.341 232.147 153.566 244.639] /A << /S /GoTo /D <74686D2E442E332E342E31> >> >>}
\MusicPDFWriteObject{7622}{<< /Type /Annot /Subtype /Link /Border [0 0 0] /H /I /C [1 0 0] /Rect [393.158 160.963 426.616 174.894] /A << /S /GoTo /D <74686D2E442E312E352E34> >> >>}
\MusicPDFWriteObject{7623}{<< /Type /Annot /Subtype /Link /Border [0 0 0] /H /I /C [1 0 0] /Rect [337.709 96.3 370.55 107.091] /A << /S /GoTo /D <74686D2E442E312E322E31> >> >>}
\MusicPDFWriteObject{7624}{<< /Type /Annot /Subtype /Link /Border [0 0 0] /H /I /C [1 0 0] /Rect [268.607 82.595 307.145 94.55] /A << /S /GoTo /D <74686D2E442E332E332E3133> >> >>}
\MusicPDFWriteObject{7625}{<< /Type /Annot /Subtype /Link /Border [0 0 0] /H /I /C [1 0 0] /Rect [191.873 502.318 225.332 514.81] /A << /S /GoTo /D <74686D2E442E332E332E37> >> >>}
\MusicPDFWriteObject{7626}{<< /Type /Annot /Subtype /Link /Border [0 0 0] /H /I /C [0 1 0] /Rect [415.995 505.346 441.055 514.213] /A << /S /GoTo /D <636974652E53636861> >> >>}
\MusicPDFWriteObject{7627}{<< /Type /Annot /Subtype /Link /Border [0 0 0] /H /I /C [0 1 0] /Rect [114.653 492.195 133.627 501.063] /A << /S /GoTo /D <636974652E547265> >> >>}
\MusicPDFWriteObject{7628}{<< /Type /Annot /Subtype /Link /Border [0 0 0] /H /I /C [1 0 0] /Rect [371.718 210.985 404.612 222.94] /A << /S /GoTo /D <74686D2E422E322E322E33> >> >>}
\MusicPDFWriteObject{7629}{<< /Type /Annot /Subtype /Link /Border [0 0 0] /H /I /C [1 0 0] /Rect [73.712 108.804 111.533 120.757] /A << /S /GoTo /D <6571756174696F6E2E442E312E31302E37> >> >>}
\MusicPDFWriteObject{7630}{<< /Type /Annot /Subtype /Link /Border [0 0 0] /H /I /C [1 0 0] /Rect [342.385 584.985 384.534 596.941] /A << /S /GoTo /D <6571756174696F6E2E442E312E31302E3130> >> >>}
\MusicPDFWriteObject{7631}{<< /Type /Annot /Subtype /Link /Border [0 0 0] /H /I /C [1 0 0] /Rect [133.856 523.655 166.698 535.607] /A << /S /GoTo /D <6571756174696F6E2E442E312E322E38> >> >>}
\MusicPDFWriteObject{7632}{<< /Type /Annot /Subtype /Link /Border [0 0 0] /H /I /C [1 0 0] /Rect [233.404 510.504 265.968 522.457] /A << /S /GoTo /D <74686D2E442E312E322E37> >> >>}
\MusicPDFWriteObject{7633}{<< /Type /Annot /Subtype /Link /Border [0 0 0] /H /I /C [1 0 0] /Rect [443.461 444.75 476.303 457.342] /A << /S /GoTo /D <74686D2E442E312E322E31> >> >>}
\MusicPDFWriteObject{7634}{<< /Type /Annot /Subtype /Link /Border [0 0 0] /H /I /C [1 0 0] /Rect [203.197 569.683 239.51 578.537] /A << /S /GoTo /D <74686D2E442E332E352E31> >> >>}
\MusicPDFWriteObject{7635}{<< /Type /Annot /Subtype /Link /Border [0 0 0] /H /I /C [1 0 0] /Rect [184.275 554.042 217.116 565.995] /A << /S /GoTo /D <6571756174696F6E2E442E332E352E33> >> >>}
\MusicPDFWriteObject{7636}{<< /Type /Annot /Subtype /Link /Border [0 0 0] /H /I /C [1 0 0] /Rect [155.389 540.891 188.23 552.844] /A << /S /GoTo /D <6571756174696F6E2E442E332E352E33> >> >>}
\MusicPDFWriteObject{7637}{<< /Type /Annot /Subtype /Link /Border [0 0 0] /H /I /C [1 0 0] /Rect [335.766 514.589 368.607 526.544] /A << /S /GoTo /D <6571756174696F6E2E442E332E352E31> >> >>}
\MusicPDFWriteObject{7638}{<< /Type /Annot /Subtype /Link /Border [0 0 0] /H /I /C [1 0 0] /Rect [191.689 468.274 224.531 480.227] /A << /S /GoTo /D <6571756174696F6E2E442E332E352E35> >> >>}
\MusicPDFWriteObject{7639}{<< /Type /Annot /Subtype /Link /Border [0 0 0] /H /I /C [1 0 0] /Rect [431.549 468.274 464.391 480.227] /A << /S /GoTo /D <6571756174696F6E2E442E332E352E33> >> >>}
\MusicPDFWriteObject{7640}{<< /Type /Annot /Subtype /Link /Border [0 0 0] /H /I /C [1 0 0] /Rect [150.616 455.123 183.51 467.078] /A << /S /GoTo /D <74686D2E422E342E312E32> >> >>}
\MusicPDFWriteObject{7641}{<< /Type /Annot /Subtype /Link /Border [0 0 0] /H /I /C [1 0 0] /Rect [412.083 224.229 444.832 236.184] /A << /S /GoTo /D <74686D2E442E332E332E37> >> >>}
\MusicPDFWriteObject{7642}{<< /Type /Annot /Subtype /Link /Border [0 0 0] /H /I /C [0 1 0] /Rect [165.12 187.267 186.341 196.135] /A << /S /GoTo /D <636974652E427275> >> >>}
\MusicPDFWriteObject{7643}{<< /Type /Annot /Subtype /Link /Border [0 0 0] /H /I /C [1 0 0] /Rect [133.322 108.897 165.67 120.85] /A << /S /GoTo /D <6571756174696F6E2E442E332E362E31> >> >>}
\MusicPDFWriteObject{7644}{<< /Type /Annot /Subtype /Link /Border [0 0 0] /H /I /C [1 0 0] /Rect [242.56 95.747 275.586 107.699] /A << /S /GoTo /D <6571756174696F6E2E442E332E362E31> >> >>}
\MusicPDFWriteObject{7645}{<< /Type /Annot /Subtype /Link /Border [0 0 0] /H /I /C [1 0 0] /Rect [462.512 95.747 495.539 107.699] /A << /S /GoTo /D <6571756174696F6E2E442E332E362E31> >> >>}
\MusicPDFWriteObject{7646}{<< /Type /Annot /Subtype /Link /Border [0 0 0] /H /I /C [1 0 0] /Rect [353.668 684.633 386.51 696.586] /A << /S /GoTo /D <6571756174696F6E2E442E332E362E34> >> >>}
\MusicPDFWriteObject{7647}{<< /Type /Annot /Subtype /Link /Border [0 0 0] /H /I /C [1 0 0] /Rect [88.936 573.549 94.318 582.611] /A << /S /GoTo /D <48666F6F746E6F74652E313532> >> >>}
\MusicPDFWriteObject{7648}{<< /Type /Annot /Subtype /Link /Border [0 0 0] /H /I /C [1 0 0] /Rect [82.959 545.311 235.461 556.102] /A << /S /GoTo /D <73656374696F6E2E452E31> >> >>}
\MusicPDFWriteObject{7649}{<< /Type /Annot /Subtype /Link /Border [0 0 0] /H /I /C [1 0 0] /Rect [82.959 532.16 313.179 542.951] /A << /S /GoTo /D <73756273656374696F6E2E452E312E31> >> >>}
\MusicPDFWriteObject{7650}{<< /Type /Annot /Subtype /Link /Border [0 0 0] /H /I /C [1 0 0] /Rect [82.959 519.009 419.31 529.8] /A << /S /GoTo /D <73756273656374696F6E2E452E312E32> >> >>}
\MusicPDFWriteObject{7651}{<< /Type /Annot /Subtype /Link /Border [0 0 0] /H /I /C [1 0 0] /Rect [82.959 505.859 313.843 516.649] /A << /S /GoTo /D <73656374696F6E2E452E32> >> >>}
\MusicPDFWriteObject{7652}{<< /Type /Annot /Subtype /Link /Border [0 0 0] /H /I /C [1 0 0] /Rect [82.959 492.708 391.56 503.499] /A << /S /GoTo /D <73756273656374696F6E2E452E322E31> >> >>}
\MusicPDFWriteObject{7653}{<< /Type /Annot /Subtype /Link /Border [0 0 0] /H /I /C [1 0 0] /Rect [82.959 479.557 497.691 490.348] /A << /S /GoTo /D <73756273656374696F6E2E452E322E32> >> >>}
\MusicPDFWriteObject{7654}{<< /Type /Annot /Subtype /Link /Border [0 0 0] /H /I /C [1 0 0] /Rect [82.959 466.406 432.812 477.197] /A << /S /GoTo /D <73756273656374696F6E2E452E322E33> >> >>}
\MusicPDFWriteObject{7655}{<< /Type /Annot /Subtype /Link /Border [0 0 0] /H /I /C [1 0 0] /Rect [82.959 453.256 415.686 464.046] /A << /S /GoTo /D <73756273656374696F6E2E452E322E34> >> >>}
\MusicPDFWriteObject{7656}{<< /Type /Annot /Subtype /Link /Border [0 0 0] /H /I /C [1 0 0] /Rect [82.959 440.105 338.743 450.896] /A << /S /GoTo /D <73656374696F6E2E452E33> >> >>}
\MusicPDFWriteObject{7657}{<< /Type /Annot /Subtype /Link /Border [0 0 0] /H /I /C [1 0 0] /Rect [82.959 426.954 449.55 437.745] /A << /S /GoTo /D <73756273656374696F6E2E452E332E31> >> >>}
\MusicPDFWriteObject{7658}{<< /Type /Annot /Subtype /Link /Border [0 0 0] /H /I /C [1 0 0] /Rect [82.959 413.803 373.982 424.594] /A << /S /GoTo /D <73756273656374696F6E2E452E332E32> >> >>}
\MusicPDFWriteObject{7659}{<< /Type /Annot /Subtype /Link /Border [0 0 0] /H /I /C [1 0 0] /Rect [82.959 402.589 202.178 411.443] /A << /S /GoTo /D <73756273656374696F6E2E452E332E33> >> >>}
\MusicPDFWriteObject{7660}{<< /Type /Annot /Subtype /Link /Border [0 0 0] /H /I /C [1 0 0] /Rect [82.959 387.502 257.983 398.293] /A << /S /GoTo /D <73756273656374696F6E2E452E332E34> >> >>}
\MusicPDFWriteObject{7661}{<< /Type /Annot /Subtype /Link /Border [0 0 0] /H /I /C [1 0 0] /Rect [82.959 374.351 395.018 385.142] /A << /S /GoTo /D <73756273656374696F6E2E452E332E35> >> >>}
\MusicPDFWriteObject{7662}{<< /Type /Annot /Subtype /Link /Border [0 0 0] /H /I /C [1 0 0] /Rect [82.959 361.2 207.739 371.991] /A << /S /GoTo /D <73656374696F6E2E452E34> >> >>}
\MusicPDFWriteObject{7663}{<< /Type /Annot /Subtype /Link /Border [0 0 0] /H /I /C [1 0 0] /Rect [82.959 348.05 298.224 358.84] /A << /S /GoTo /D <73756273656374696F6E2E452E342E31> >> >>}
\MusicPDFWriteObject{7664}{<< /Type /Annot /Subtype /Link /Border [0 0 0] /H /I /C [1 0 0] /Rect [82.959 334.899 255.612 345.69] /A << /S /GoTo /D <73756273656374696F6E2E452E342E32> >> >>}
\MusicPDFWriteObject{7665}{<< /Type /Annot /Subtype /Link /Border [0 0 0] /H /I /C [1 0 0] /Rect [82.959 321.748 224.672 332.539] /A << /S /GoTo /D <73656374696F6E2E452E35> >> >>}
\MusicPDFWriteObject{7666}{<< /Type /Annot /Subtype /Link /Border [0 0 0] /H /I /C [1 0 0] /Rect [82.959 308.597 274.5 319.388] /A << /S /GoTo /D <73756273656374696F6E2E452E352E31> >> >>}
\MusicPDFWriteObject{7667}{<< /Type /Annot /Subtype /Link /Border [0 0 0] /H /I /C [1 0 0] /Rect [82.959 295.447 170.942 306.237] /A << /S /GoTo /D <73756273656374696F6E2E452E352E32> >> >>}
\MusicPDFWriteObject{7668}{<< /Type /Annot /Subtype /Link /Border [0 0 0] /H /I /C [1 0 0] /Rect [82.959 284.233 189.672 293.087] /A << /S /GoTo /D <73756273656374696F6E2E452E352E33> >> >>}
\MusicPDFWriteObject{7669}{<< /Type /Annot /Subtype /Link /Border [0 0 0] /H /I /C [1 0 0] /Rect [82.959 269.145 235.305 280.692] /A << /S /GoTo /D <73756273656374696F6E2E452E352E34> >> >>}
\MusicPDFWriteObject{7670}{<< /Type /Annot /Subtype /Link /Border [0 0 0] /H /I /C [1 0 0] /Rect [82.959 255.994 286.897 266.785] /A << /S /GoTo /D <73756273656374696F6E2E452E352E35> >> >>}
\MusicPDFWriteObject{7671}{<< /Type /Annot /Subtype /Link /Border [0 0 0] /H /I /C [1 0 0] /Rect [82.959 242.29 266.395 254.245] /A << /S /GoTo /D <73756273656374696F6E2E452E352E36> >> >>}
\MusicPDFWriteObject{7672}{<< /Type /Annot /Subtype /Link /Border [0 0 0] /H /I /C [1 0 0] /Rect [82.959 229.693 280.644 240.483] /A << /S /GoTo /D <73756273656374696F6E2E452E352E37> >> >>}
\MusicPDFWriteObject{7673}{<< /Type /Annot /Subtype /Link /Border [0 0 0] /H /I /C [1 0 0] /Rect [82.959 215.628 207.697 227.333] /A << /S /GoTo /D <73656374696F6E2E452E36> >> >>}
\MusicPDFWriteObject{7674}{<< /Type /Annot /Subtype /Link /Border [0 0 0] /H /I /C [1 0 0] /Rect [82.959 203.391 183.557 214.182] /A << /S /GoTo /D <73756273656374696F6E2E452E362E31> >> >>}
\MusicPDFWriteObject{7675}{<< /Type /Annot /Subtype /Link /Border [0 0 0] /H /I /C [1 0 0] /Rect [82.959 190.241 379.599 201.031] /A << /S /GoTo /D <73756273656374696F6E2E452E362E32> >> >>}
\MusicPDFWriteObject{7676}{<< /Type /Annot /Subtype /Link /Border [0 0 0] /H /I /C [1 0 0] /Rect [82.959 177.09 250.485 187.88] /A << /S /GoTo /D <73756273656374696F6E2E452E362E33> >> >>}
\MusicPDFWriteObject{7677}{<< /Type /Annot /Subtype /Link /Border [0 0 0] /H /I /C [1 0 0] /Rect [82.959 163.939 175.479 174.73] /A << /S /GoTo /D <73756273656374696F6E2E452E362E34> >> >>}
\MusicPDFWriteObject{7678}{<< /Type /Annot /Subtype /Link /Border [0 0 0] /H /I /C [0 1 0] /Rect [419.746 301.872 436.886 310.74] /A << /S /GoTo /D <636974652E4475> >> >>}
\MusicPDFWriteObject{7679}{<< /Type /Annot /Subtype /Link /Border [0 0 0] /H /I /C [0 1 0] /Rect [447.862 301.872 490.885 310.74] /A << /S /GoTo /D <636974652E426F752D4754> >> >>}
\MusicPDFWriteObject{7680}{<< /Type /Annot /Subtype /Link /Border [0 0 0] /H /I /C [0 1 0] /Rect [501.862 301.872 518.082 310.74] /A << /S /GoTo /D <636974652E4B65> >> >>}
\MusicPDFWriteObject{7681}{<< /Type /Annot /Subtype /Link /Border [0 0 0] /H /I /C [0 1 0] /Rect [73.826 288.721 90.012 297.589] /A << /S /GoTo /D <636974652E5350> >> >>}
\MusicPDFWriteObject{7682}{<< /Type /Annot /Subtype /Link /Border [0 0 0] /H /I /C [0 1 0] /Rect [223.331 288.721 266.354 297.589] /A << /S /GoTo /D <636974652E426F752D4754> >> >>}
\MusicPDFWriteObject{7683}{<< /Type /Annot /Subtype /Link /Border [0 0 0] /H /I /C [0 1 0] /Rect [340.172 288.721 357.312 297.589] /A << /S /GoTo /D <636974652E4475> >> >>}
\MusicPDFWriteObject{7684}{<< /Type /Annot /Subtype /Link /Border [0 0 0] /H /I /C [0 1 0] /Rect [425.689 288.721 441.875 297.589] /A << /S /GoTo /D <636974652E5350> >> >>}
\MusicPDFWriteObject{7685}{<< /Type /Annot /Subtype /Link /Border [0 0 0] /H /I /C [0 1 0] /Rect [275.014 275.571 318.037 284.438] /A << /S /GoTo /D <636974652E426F752D4754> >> >>}
\MusicPDFWriteObject{7686}{<< /Type /Annot /Subtype /Link /Border [0 0 0] /H /I /C [0 1 0] /Rect [369.556 275.571 386.696 284.438] /A << /S /GoTo /D <636974652E4475> >> >>}
\MusicPDFWriteObject{7687}{<< /Type /Annot /Subtype /Link /Border [0 0 0] /H /I /C [0 1 0] /Rect [476.805 275.571 492.991 284.438] /A << /S /GoTo /D <636974652E5350> >> >>}
\MusicPDFWriteObject{7688}{<< /Type /Annot /Subtype /Link /Border [0 0 0] /H /I /C [0 1 0] /Rect [275.084 262.42 318.107 271.288] /A << /S /GoTo /D <636974652E426F752D4754> >> >>}
\MusicPDFWriteObject{7689}{<< /Type /Annot /Subtype /Link /Border [0 0 0] /H /I /C [0 1 0] /Rect [368.632 262.42 385.772 271.288] /A << /S /GoTo /D <636974652E4475> >> >>}
\MusicPDFWriteObject{7690}{<< /Type /Annot /Subtype /Link /Border [0 0 0] /H /I /C [0 1 0] /Rect [455.277 262.42 471.497 271.288] /A << /S /GoTo /D <636974652E4B65> >> >>}
\MusicPDFWriteObject{7691}{<< /Type /Annot /Subtype /Link /Border [0 0 0] /H /I /C [0 1 0] /Rect [527.084 262.42 543.27 271.288] /A << /S /GoTo /D <636974652E5350> >> >>}
\MusicPDFWriteObject{7692}{<< /Type /Annot /Subtype /Link /Border [0 0 0] /H /I /C [0 1 0] /Rect [374.381 249.269 391.521 258.137] /A << /S /GoTo /D <636974652E4475> >> >>}
\MusicPDFWriteObject{7693}{<< /Type /Annot /Subtype /Link /Border [0 0 0] /H /I /C [0 1 0] /Rect [482.392 249.269 498.612 258.137] /A << /S /GoTo /D <636974652E4B65> >> >>}
\MusicPDFWriteObject{7694}{<< /Type /Annot /Subtype /Link /Border [0 0 0] /H /I /C [0 1 0] /Rect [359.753 236.118 375.939 244.986] /A << /S /GoTo /D <636974652E5350> >> >>}
\MusicPDFWriteObject{7695}{<< /Type /Annot /Subtype /Link /Border [0 0 0] /H /I /C [0 1 0] /Rect [527.017 236.118 543.203 244.986] /A << /S /GoTo /D <636974652E5350> >> >>}
\MusicPDFWriteObject{7696}{<< /Type /Annot /Subtype /Link /Border [0 0 0] /H /I /C [0 1 0] /Rect [204.072 196.666 221.212 205.534] /A << /S /GoTo /D <636974652E4475> >> >>}
\MusicPDFWriteObject{7697}{<< /Type /Annot /Subtype /Link /Border [0 0 0] /H /I /C [0 1 0] /Rect [227.208 196.666 243.428 205.534] /A << /S /GoTo /D <636974652E4B65> >> >>}
\MusicPDFWriteObject{7698}{<< /Type /Annot /Subtype /Link /Border [0 0 0] /H /I /C [0 1 0] /Rect [289.367 196.666 332.39 205.534] /A << /S /GoTo /D <636974652E426F752D4754> >> >>}
\MusicPDFWriteObject{7699}{<< /Type /Annot /Subtype /Link /Border [0 0 0] /H /I /C [0 1 0] /Rect [338.386 196.666 354.572 205.534] /A << /S /GoTo /D <636974652E5350> >> >>}
\MusicPDFWriteObject{7700}{<< /Type /Annot /Subtype /Link /Border [0 0 0] /H /I /C [0 1 0] /Rect [73.82 157.214 98.597 166.082] /A << /S /GoTo /D <636974652E4C474C> >> >>}
\MusicPDFWriteObject{7701}{<< /Type /Annot /Subtype /Link /Border [0 0 0] /H /I /C [0 1 0] /Rect [103.653 157.214 130.623 166.082] /A << /S /GoTo /D <636974652E52657A6B> >> >>}
\MusicPDFWriteObject{7702}{<< /Type /Annot /Subtype /Link /Border [0 0 0] /H /I /C [0 1 0] /Rect [135.679 157.214 153.739 166.082] /A << /S /GoTo /D <636974652E537465> >> >>}
\MusicPDFWriteObject{7703}{<< /Type /Annot /Subtype /Link /Border [0 0 0] /H /I /C [0 1 0] /Rect [117.704 673.973 133.89 682.84] /A << /S /GoTo /D <636974652E5350> >> >>}
\MusicPDFWriteObject{7704}{<< /Type /Annot /Subtype /Link /Border [0 0 0] /H /I /C [0 1 0] /Rect [433.675 595.068 449.861 603.936] /A << /S /GoTo /D <636974652E5350> >> >>}
\MusicPDFWriteObject{7705}{<< /Type /Annot /Subtype /Link /Border [0 0 0] /H /I /C [0 1 0] /Rect [136.969 494.235 163.938 503.102] /A << /S /GoTo /D <636974652E52657A6B> >> >>}
\MusicPDFWriteObject{7706}{<< /Type /Annot /Subtype /Link /Border [0 0 0] /H /I /C [0 1 0] /Rect [168.913 494.235 186.44 503.102] /A << /S /GoTo /D <636974652E537472> >> >>}
\MusicPDFWriteObject{7707}{<< /Type /Annot /Subtype /Link /Border [0 0 0] /H /I /C [0 1 0] /Rect [336.248 481.084 356.833 489.952] /A << /S /GoTo /D <636974652E427234> >> >>}
\MusicPDFWriteObject{7708}{<< /Type /Annot /Subtype /Link /Border [0 0 0] /H /I /C [0 1 0] /Rect [479.39 481.084 496.917 489.952] /A << /S /GoTo /D <636974652E537472> >> >>}
\MusicPDFWriteObject{7709}{<< /Type /Annot /Subtype /Link /Border [0 0 0] /H /I /C [0 1 0] /Rect [289.517 467.933 316.486 476.801] /A << /S /GoTo /D <636974652E52657A6B> >> >>}
\MusicPDFWriteObject{7710}{<< /Type /Annot /Subtype /Link /Border [0 0 0] /H /I /C [0 1 0] /Rect [365.86 467.933 383.92 476.801] /A << /S /GoTo /D <636974652E537465> >> >>}
\MusicPDFWriteObject{7711}{<< /Type /Annot /Subtype /Link /Border [0 0 0] /H /I /C [1 0 0] /Rect [227.474 372.466 243.745 383.257] /A << /S /GoTo /D <73656374696F6E2E452E32> >> >>}
\MusicPDFWriteObject{7712}{<< /Type /Annot /Subtype /Link /Border [0 0 0] /H /I /C [0 1 0] /Rect [307.026 269.197 326.172 278.064] /A << /S /GoTo /D <636974652E4247> >> >>}
\MusicPDFWriteObject{7713}{<< /Type /Annot /Subtype /Link /Border [0 0 0] /H /I /C [0 1 0] /Rect [522.277 269.197 540.337 278.064] /A << /S /GoTo /D <636974652E537465> >> >>}
\MusicPDFWriteObject{7714}{<< /Type /Annot /Subtype /Link /Border [0 0 0] /H /I /C [0 1 0] /Rect [278.072 256.046 297.218 264.914] /A << /S /GoTo /D <636974652E4247> >> >>}
\MusicPDFWriteObject{7715}{<< /Type /Annot /Subtype /Link /Border [0 0 0] /H /I /C [0 1 0] /Rect [436.244 256.046 462.362 264.914] /A << /S /GoTo /D <636974652E4C4D53> >> >>}
\MusicPDFWriteObject{7716}{<< /Type /Annot /Subtype /Link /Border [0 0 0] /H /I /C [0 1 0] /Rect [474.01 242.895 500.128 251.763] /A << /S /GoTo /D <636974652E4C4D53> >> >>}
\MusicPDFWriteObject{7717}{<< /Type /Annot /Subtype /Link /Border [0 0 0] /H /I /C [0 1 0] /Rect [324.446 227.808 348.607 238.612] /A << /S /GoTo /D <636974652E4D6179> >> >>}
\MusicPDFWriteObject{7718}{<< /Type /Annot /Subtype /Link /Border [0 0 0] /H /I /C [0 1 0] /Rect [226.539 216.594 251.316 225.461] /A << /S /GoTo /D <636974652E4C474C> >> >>}
\MusicPDFWriteObject{7719}{<< /Type /Annot /Subtype /Link /Border [0 0 0] /H /I /C [0 1 0] /Rect [334.615 216.594 359.392 225.461] /A << /S /GoTo /D <636974652E4C474C> >> >>}
\MusicPDFWriteObject{7720}{<< /Type /Annot /Subtype /Link /Border [0 0 0] /H /I /C [0 1 0] /Rect [101.498 150.84 122.083 159.708] /A << /S /GoTo /D <636974652E427231> >> >>}
\MusicPDFWriteObject{7721}{<< /Type /Annot /Subtype /Link /Border [0 0 0] /H /I /C [0 1 0] /Rect [256.713 150.84 277.298 159.708] /A << /S /GoTo /D <636974652E427231> >> >>}
\MusicPDFWriteObject{7722}{<< /Type /Annot /Subtype /Link /Border [0 0 0] /H /I /C [0 1 0] /Rect [73.721 137.689 94.306 146.557] /A << /S /GoTo /D <636974652E427231> >> >>}
\MusicPDFWriteObject{7723}{<< /Type /Annot /Subtype /Link /Border [0 0 0] /H /I /C [0 1 0] /Rect [523.825 124.538 541.885 133.406] /A << /S /GoTo /D <636974652E537465> >> >>}
\MusicPDFWriteObject{7724}{<< /Type /Annot /Subtype /Link /Border [0 0 0] /H /I /C [0 1 0] /Rect [501.797 98.237 522.382 107.105] /A << /S /GoTo /D <636974652E427232> >> >>}
\MusicPDFWriteObject{7725}{<< /Type /Annot /Subtype /Link /Border [0 0 0] /H /I /C [0 1 0] /Rect [73.726 85.086 94.311 93.954] /A << /S /GoTo /D <636974652E427233> >> >>}
\MusicPDFWriteObject{7726}{<< /Type /Annot /Subtype /Link /Border [0 0 0] /H /I /C [0 1 0] /Rect [166.511 673.973 184.571 682.84] /A << /S /GoTo /D <636974652E537465> >> >>}
\MusicPDFWriteObject{7727}{<< /Type /Annot /Subtype /Link /Border [0 0 0] /H /I /C [0 1 0] /Rect [210.799 672.036 236.406 682.84] /A << /S /GoTo /D <636974652E566F6774> >> >>}
\MusicPDFWriteObject{7728}{<< /Type /Annot /Subtype /Link /Border [0 0 0] /H /I /C [0 1 0] /Rect [409.788 634.52 426.354 643.388] /A << /S /GoTo /D <636974652E4761> >> >>}
\MusicPDFWriteObject{7729}{<< /Type /Annot /Subtype /Link /Border [0 0 0] /H /I /C [0 1 0] /Rect [304.725 621.369 320.945 630.237] /A << /S /GoTo /D <636974652E4B65> >> >>}
\MusicPDFWriteObject{7730}{<< /Type /Annot /Subtype /Link /Border [0 0 0] /H /I /C [0 1 0] /Rect [501.191 621.369 521.776 630.237] /A << /S /GoTo /D <636974652E427231> >> >>}
\MusicPDFWriteObject{7731}{<< /Type /Annot /Subtype /Link /Border [0 0 0] /H /I /C [0 1 0] /Rect [137.971 608.219 158.556 617.087] /A << /S /GoTo /D <636974652E427232> >> >>}
\MusicPDFWriteObject{7732}{<< /Type /Annot /Subtype /Link /Border [0 0 0] /H /I /C [0 1 0] /Rect [407.911 608.219 424.131 617.087] /A << /S /GoTo /D <636974652E4B65> >> >>}
\MusicPDFWriteObject{7733}{<< /Type /Annot /Subtype /Link /Border [0 0 0] /H /I /C [0 1 0] /Rect [171.739 595.068 192.324 603.936] /A << /S /GoTo /D <636974652E427232> >> >>}
\MusicPDFWriteObject{7734}{<< /Type /Annot /Subtype /Link /Border [0 0 0] /H /I /C [0 1 0] /Rect [498.969 581.917 519.554 590.785] /A << /S /GoTo /D <636974652E427232> >> >>}
\MusicPDFWriteObject{7735}{<< /Type /Annot /Subtype /Link /Border [0 0 0] /H /I /C [0 1 0] /Rect [73.715 566.83 99.322 577.634] /A << /S /GoTo /D <636974652E566F6774> >> >>}
\MusicPDFWriteObject{7736}{<< /Type /Annot /Subtype /Link /Border [0 0 0] /H /I /C [0 1 0] /Rect [510.366 568.766 543.27 577.634] /A << /S /GoTo /D <636974652E6E4C6162> >> >>}
\MusicPDFWriteObject{7737}{<< /Type /Annot /Subtype /Link /Border [0 0 0] /H /I /C [0 1 0] /Rect [525.207 484.735 543.267 493.603] /A << /S /GoTo /D <636974652E537465> >> >>}
\MusicPDFWriteObject{7738}{<< /Type /Annot /Subtype /Link /Border [0 0 0] /H /I /C [0 1 0] /Rect [298.706 456.497 324.312 467.302] /A << /S /GoTo /D <636974652E566F6774> >> >>}
\MusicPDFWriteObject{7739}{<< /Type /Annot /Subtype /Link /Border [0 0 0] /H /I /C [0 1 0] /Rect [384.984 443.346 410.591 454.151] /A << /S /GoTo /D <636974652E566F6774> >> >>}
\MusicPDFWriteObject{7740}{<< /Type /Annot /Subtype /Link /Border [0 0 0] /H /I /C [0 1 0] /Rect [352.934 430.196 378.541 441] /A << /S /GoTo /D <636974652E566F6774> >> >>}
\MusicPDFWriteObject{7741}{<< /Type /Annot /Subtype /Link /Border [0 0 0] /H /I /C [0 1 0] /Rect [198.139 417.045 223.746 427.849] /A << /S /GoTo /D <636974652E566F6774> >> >>}
\MusicPDFWriteObject{7742}{<< /Type /Annot /Subtype /Link /Border [0 0 0] /H /I /C [0 1 0] /Rect [424.853 418.981 445.437 427.849] /A << /S /GoTo /D <636974652E427233> >> >>}
\MusicPDFWriteObject{7743}{<< /Type /Annot /Subtype /Link /Border [0 0 0] /H /I /C [0 1 0] /Rect [151.482 405.831 169.542 414.699] /A << /S /GoTo /D <636974652E537465> >> >>}
\MusicPDFWriteObject{7744}{<< /Type /Annot /Subtype /Link /Border [0 0 0] /H /I /C [0 1 0] /Rect [473.128 405.831 503.542 414.699] /A << /S /GoTo /D <636974652E4D61634C> >> >>}
\MusicPDFWriteObject{7745}{<< /Type /Annot /Subtype /Link /Border [0 0 0] /H /I /C [0 1 0] /Rect [245.821 374.403 268.592 383.27] /A << /S /GoTo /D <636974652E454C53> >> >>}
\MusicPDFWriteObject{7746}{<< /Type /Annot /Subtype /Link /Border [0 0 0] /H /I /C [0 1 0] /Rect [344.011 374.403 366.782 383.27] /A << /S /GoTo /D <636974652E454C53> >> >>}
\MusicPDFWriteObject{7747}{<< /Type /Annot /Subtype /Link /Border [0 0 0] /H /I /C [0 1 0] /Rect [438.68 374.403 461.451 383.27] /A << /S /GoTo /D <636974652E454C53> >> >>}
\MusicPDFWriteObject{7748}{<< /Type /Annot /Subtype /Link /Border [0 0 0] /H /I /C [0 1 0] /Rect [136.615 334.95 169.519 343.818] /A << /S /GoTo /D <636974652E6E4C6162> >> >>}
\MusicPDFWriteObject{7749}{<< /Type /Annot /Subtype /Link /Border [0 0 0] /H /I /C [0 1 0] /Rect [316.011 334.95 348.915 343.818] /A << /S /GoTo /D <636974652E6E4C6162> >> >>}
\MusicPDFWriteObject{7750}{<< /Type /Annot /Subtype /Link /Border [0 0 0] /H /I /C [0 1 0] /Rect [141.844 308.649 164.615 317.517] /A << /S /GoTo /D <636974652E454C53> >> >>}
\MusicPDFWriteObject{7751}{<< /Type /Annot /Subtype /Link /Border [0 0 0] /H /I /C [0 1 0] /Rect [466.508 282.347 492.737 291.215] /A << /S /GoTo /D <636974652E4D6343> >> >>}
\MusicPDFWriteObject{7752}{<< /Type /Annot /Subtype /Link /Border [0 0 0] /H /I /C [0 1 0] /Rect [225.02 269.197 249.797 278.064] /A << /S /GoTo /D <636974652E4C474C> >> >>}
\MusicPDFWriteObject{7753}{<< /Type /Annot /Subtype /Link /Border [0 0 0] /H /I /C [0 1 0] /Rect [485.576 269.197 503.104 278.064] /A << /S /GoTo /D <636974652E537472> >> >>}
\MusicPDFWriteObject{7754}{<< /Type /Annot /Subtype /Link /Border [0 0 0] /H /I /C [0 1 0] /Rect [93.657 256.046 120.627 264.914] /A << /S /GoTo /D <636974652E52657A6B> >> >>}
\MusicPDFWriteObject{7755}{<< /Type /Annot /Subtype /Link /Border [0 0 0] /H /I /C [0 1 0] /Rect [237.56 242.895 253.78 251.763] /A << /S /GoTo /D <636974652E4B65> >> >>}
\MusicPDFWriteObject{7756}{<< /Type /Annot /Subtype /Link /Border [0 0 0] /H /I /C [1 0 0] /Rect [345.018 214.103 369.728 226.056] /A << /S /GoTo /D <73756273656374696F6E2E452E322E34> >> >>}
\MusicPDFWriteObject{7757}{<< /Type /Annot /Subtype /Link /Border [0 0 0] /H /I /C [1 0 0] /Rect [475.8 214.103 513.491 226.056] /A << /S /GoTo /D <74686D2E452E322E342E3131> >> >>}
\MusicPDFWriteObject{7758}{<< /Type /Annot /Subtype /Link /Border [0 0 0] /H /I /C [0 1 0] /Rect [221.634 190.292 241.389 199.16] /A << /S /GoTo /D <636974652E4C4D> >> >>}
\MusicPDFWriteObject{7759}{<< /Type /Annot /Subtype /Link /Border [0 0 0] /H /I /C [0 1 0] /Rect [354.137 177.141 373.892 186.009] /A << /S /GoTo /D <636974652E4C4D> >> >>}
\MusicPDFWriteObject{7760}{<< /Type /Annot /Subtype /Link /Border [0 0 0] /H /I /C [0 1 0] /Rect [431.706 177.141 447.927 186.009] /A << /S /GoTo /D <636974652E4B65> >> >>}
\MusicPDFWriteObject{7761}{<< /Type /Annot /Subtype /Link /Border [0 0 0] /H /I /C [0 1 0] /Rect [85.505 162.054 111.111 172.858] /A << /S /GoTo /D <636974652E566F6774> >> >>}
\MusicPDFWriteObject{7762}{<< /Type /Annot /Subtype /Link /Border [0 0 0] /H /I /C [0 1 0] /Rect [290.767 135.752 316.374 146.557] /A << /S /GoTo /D <636974652E566F6774> >> >>}
\MusicPDFWriteObject{7763}{<< /Type /Annot /Subtype /Link /Border [0 0 0] /H /I /C [0 1 0] /Rect [497.836 137.689 514.056 146.557] /A << /S /GoTo /D <636974652E4B65> >> >>}
\MusicPDFWriteObject{7764}{<< /Type /Annot /Subtype /Link /Border [0 0 0] /H /I /C [1 0 0] /Rect [154.748 122.048 187.12 134.001] /A << /S /GoTo /D <74686D2E452E322E342E32> >> >>}
\MusicPDFWriteObject{7765}{<< /Type /Annot /Subtype /Link /Border [0 0 0] /H /I /C [0 1 0] /Rect [305.253 124.538 323.313 133.406] /A << /S /GoTo /D <636974652E537465> >> >>}
\MusicPDFWriteObject{7766}{<< /Type /Annot /Subtype /Link /Border [0 0 0] /H /I /C [0 1 0] /Rect [498.082 98.237 524.311 107.105] /A << /S /GoTo /D <636974652E4D6343> >> >>}
\MusicPDFWriteObject{7767}{<< /Type /Annot /Subtype /Link /Border [0 0 0] /H /I /C [0 1 0] /Rect [73.768 85.086 98.544 93.954] /A << /S /GoTo /D <636974652E4C474C> >> >>}
\MusicPDFWriteObject{7768}{<< /Type /Annot /Subtype /Link /Border [0 0 0] /H /I /C [0 1 0] /Rect [445.239 687.123 463.299 695.991] /A << /S /GoTo /D <636974652E537465> >> >>}
\MusicPDFWriteObject{7769}{<< /Type /Annot /Subtype /Link /Border [0 0 0] /H /I /C [1 0 0] /Rect [71.004 605.729 108.695 617.681] /A << /S /GoTo /D <74686D2E452E322E332E3134> >> >>}
\MusicPDFWriteObject{7770}{<< /Type /Annot /Subtype /Link /Border [0 0 0] /H /I /C [0 1 0] /Rect [304.655 555.616 330.884 564.484] /A << /S /GoTo /D <636974652E4D6343> >> >>}
\MusicPDFWriteObject{7771}{<< /Type /Annot /Subtype /Link /Border [0 0 0] /H /I /C [0 1 0] /Rect [130.765 529.314 156.994 538.182] /A << /S /GoTo /D <636974652E4D6343> >> >>}
\MusicPDFWriteObject{7772}{<< /Type /Annot /Subtype /Link /Border [0 0 0] /H /I /C [0 1 0] /Rect [499.681 529.314 517.741 538.182] /A << /S /GoTo /D <636974652E537465> >> >>}
\MusicPDFWriteObject{7773}{<< /Type /Annot /Subtype /Link /Border [0 0 0] /H /I /C [1 0 0] /Rect [184.376 434.769 208.44 446.722] /A << /S /GoTo /D <73756273656374696F6E2E452E322E34> >> >>}
\MusicPDFWriteObject{7774}{<< /Type /Annot /Subtype /Link /Border [0 0 0] /H /I /C [1 0 0] /Rect [308.786 434.769 345.463 446.722] /A << /S /GoTo /D <74686D2E452E322E342E3132> >> >>}
\MusicPDFWriteObject{7775}{<< /Type /Annot /Subtype /Link /Border [0 0 0] /H /I /C [0 1 0] /Rect [509.545 410.957 527.606 419.825] /A << /S /GoTo /D <636974652E537465> >> >>}
\MusicPDFWriteObject{7776}{<< /Type /Annot /Subtype /Link /Border [0 0 0] /H /I /C [0 1 0] /Rect [256.048 397.807 274.108 406.674] /A << /S /GoTo /D <636974652E537465> >> >>}
\MusicPDFWriteObject{7777}{<< /Type /Annot /Subtype /Link /Border [0 0 0] /H /I /C [0 1 0] /Rect [368.014 384.656 392.79 393.524] /A << /S /GoTo /D <636974652E4C474C> >> >>}
\MusicPDFWriteObject{7778}{<< /Type /Annot /Subtype /Link /Border [0 0 0] /H /I /C [0 1 0] /Rect [459.853 384.656 486.082 393.524] /A << /S /GoTo /D <636974652E4D6343> >> >>}
\MusicPDFWriteObject{7779}{<< /Type /Annot /Subtype /Link /Border [0 0 0] /H /I /C [0 1 0] /Rect [510.562 384.656 528.09 393.524] /A << /S /GoTo /D <636974652E537472> >> >>}
\MusicPDFWriteObject{7780}{<< /Type /Annot /Subtype /Link /Border [0 0 0] /H /I /C [0 1 0] /Rect [109.575 371.505 136.544 380.373] /A << /S /GoTo /D <636974652E52657A6B> >> >>}
\MusicPDFWriteObject{7781}{<< /Type /Annot /Subtype /Link /Border [0 0 0] /H /I /C [0 1 0] /Rect [361.867 345.204 388.096 354.071] /A << /S /GoTo /D <636974652E4D6343> >> >>}
\MusicPDFWriteObject{7782}{<< /Type /Annot /Subtype /Link /Border [0 0 0] /H /I /C [0 1 0] /Rect [429.839 345.204 456.808 354.071] /A << /S /GoTo /D <636974652E52657A6B> >> >>}
\MusicPDFWriteObject{7783}{<< /Type /Annot /Subtype /Link /Border [0 0 0] /H /I /C [0 1 0] /Rect [195.411 318.902 222.38 327.77] /A << /S /GoTo /D <636974652E52657A6B> >> >>}
\MusicPDFWriteObject{7784}{<< /Type /Annot /Subtype /Link /Border [0 0 0] /H /I /C [0 1 0] /Rect [172.161 305.751 199.131 314.619] /A << /S /GoTo /D <636974652E52657A6B> >> >>}
\MusicPDFWriteObject{7785}{<< /Type /Annot /Subtype /Link /Border [0 0 0] /H /I /C [0 1 0] /Rect [290.123 239.998 314.9 248.865] /A << /S /GoTo /D <636974652E4C474C> >> >>}
\MusicPDFWriteObject{7786}{<< /Type /Annot /Subtype /Link /Border [0 0 0] /H /I /C [0 1 0] /Rect [275.917 226.847 300.694 235.715] /A << /S /GoTo /D <636974652E4C474C> >> >>}
\MusicPDFWriteObject{7787}{<< /Type /Annot /Subtype /Link /Border [0 0 0] /H /I /C [1 0 0] /Rect [214.046 211.206 251.177 223.159] /A << /S /GoTo /D <74686D2E452E312E322E3138> >> >>}
\MusicPDFWriteObject{7788}{<< /Type /Annot /Subtype /Link /Border [0 0 0] /H /I /C [1 0 0] /Rect [279.272 198.055 316.964 210.008] /A << /S /GoTo /D <74686D2E452E322E322E3130> >> >>}
\MusicPDFWriteObject{7789}{<< /Type /Annot /Subtype /Link /Border [0 0 0] /H /I /C [1 0 0] /Rect [413.447 198.055 451.139 210.008] /A << /S /GoTo /D <74686D2E452E312E322E3133> >> >>}
\MusicPDFWriteObject{7790}{<< /Type /Annot /Subtype /Link /Border [0 0 0] /H /I /C [1 0 0] /Rect [458.304 184.904 489.715 196.857] /A << /S /GoTo /D <74686D2E452E322E332E38> >> >>}
\MusicPDFWriteObject{7791}{<< /Type /Annot /Subtype /Link /Border [0 0 0] /H /I /C [1 0 0] /Rect [106.128 171.754 142.665 183.706] /A << /S /GoTo /D <74686D2E452E312E322E3134> >> >>}
\MusicPDFWriteObject{7792}{<< /Type /Annot /Subtype /Link /Border [0 0 0] /H /I /C [1 0 0] /Rect [212.243 171.754 248.779 183.706] /A << /S /GoTo /D <74686D2E452E322E332E3134> >> >>}
\MusicPDFWriteObject{7793}{<< /Type /Annot /Subtype /Link /Border [0 0 0] /H /I /C [1 0 0] /Rect [508.344 642.186 544.845 654.139] /A << /S /GoTo /D <74686D2E452E312E322E3235> >> >>}
\MusicPDFWriteObject{7794}{<< /Type /Annot /Subtype /Link /Border [0 0 0] /H /I /C [1 0 0] /Rect [176.029 629.035 208.04 640.988] /A << /S /GoTo /D <74686D2E452E312E322E32> >> >>}
\MusicPDFWriteObject{7795}{<< /Type /Annot /Subtype /Link /Border [0 0 0] /H /I /C [1 0 0] /Rect [411.311 629.035 448.302 640.988] /A << /S /GoTo /D <74686D2E452E312E322E3136> >> >>}
\MusicPDFWriteObject{7796}{<< /Type /Annot /Subtype /Link /Border [0 0 0] /H /I /C [0 1 0] /Rect [73.715 605.224 91.775 614.092] /A << /S /GoTo /D <636974652E537465> >> >>}
\MusicPDFWriteObject{7797}{<< /Type /Annot /Subtype /Link /Border [0 0 0] /H /I /C [0 1 0] /Rect [335.087 605.224 353.147 614.092] /A << /S /GoTo /D <636974652E537465> >> >>}
\MusicPDFWriteObject{7798}{<< /Type /Annot /Subtype /Link /Border [0 0 0] /H /I /C [1 0 0] /Rect [469.85 602.734 501.262 614.687] /A << /S /GoTo /D <74686D2E452E322E322E35> >> >>}
\MusicPDFWriteObject{7799}{<< /Type /Annot /Subtype /Link /Border [0 0 0] /H /I /C [1 0 0] /Rect [71.004 576.432 103.045 588.385] /A << /S /GoTo /D <74686D2E452E322E332E36> >> >>}
\MusicPDFWriteObject{7800}{<< /Type /Annot /Subtype /Link /Border [0 0 0] /H /I /C [0 1 0] /Rect [357.203 578.922 375.264 587.79] /A << /S /GoTo /D <636974652E537465> >> >>}
\MusicPDFWriteObject{7801}{<< /Type /Annot /Subtype /Link /Border [0 0 0] /H /I /C [1 0 0] /Rect [147.224 563.282 184.215 575.234] /A << /S /GoTo /D <74686D2E452E322E332E3132> >> >>}
\MusicPDFWriteObject{7802}{<< /Type /Annot /Subtype /Link /Border [0 0 0] /H /I /C [1 0 0] /Rect [490.311 537.533 506.771 548.324] /A << /S /GoTo /D <73656374696F6E2E452E33> >> >>}
\MusicPDFWriteObject{7803}{<< /Type /Annot /Subtype /Link /Border [0 0 0] /H /I /C [0 1 0] /Rect [482.242 513.169 507.019 522.036] /A << /S /GoTo /D <636974652E4C474C> >> >>}
\MusicPDFWriteObject{7804}{<< /Type /Annot /Subtype /Link /Border [0 0 0] /H /I /C [0 1 0] /Rect [73.826 500.018 100.055 508.886] /A << /S /GoTo /D <636974652E4D6343> >> >>}
\MusicPDFWriteObject{7805}{<< /Type /Annot /Subtype /Link /Border [0 0 0] /H /I /C [0 1 0] /Rect [125.901 500.018 143.428 508.886] /A << /S /GoTo /D <636974652E537472> >> >>}
\MusicPDFWriteObject{7806}{<< /Type /Annot /Subtype /Link /Border [0 0 0] /H /I /C [0 1 0] /Rect [169.274 500.018 196.243 508.886] /A << /S /GoTo /D <636974652E52657A6B> >> >>}
\MusicPDFWriteObject{7807}{<< /Type /Annot /Subtype /Link /Border [0 0 0] /H /I /C [0 1 0] /Rect [319.455 500.018 342.225 508.886] /A << /S /GoTo /D <636974652E454C53> >> >>}
\MusicPDFWriteObject{7808}{<< /Type /Annot /Subtype /Link /Border [0 0 0] /H /I /C [0 1 0] /Rect [522.809 500.018 540.337 508.886] /A << /S /GoTo /D <636974652E537472> >> >>}
\MusicPDFWriteObject{7809}{<< /Type /Annot /Subtype /Link /Border [0 0 0] /H /I /C [1 0 0] /Rect [400.934 424.632 406.897 437.489] /A << /S /GoTo /D <48666F6F746E6F74652E313533> >> >>}
\MusicPDFWriteObject{7810}{<< /Type /Annot /Subtype /Link /Border [0 0 0] /H /I /C [1 0 0] /Rect [233.222 351.721 265.234 362.512] /A << /S /GoTo /D <74686D2E452E312E312E31> >> >>}
\MusicPDFWriteObject{7811}{<< /Type /Annot /Subtype /Link /Border [0 0 0] /H /I /C [1 0 0] /Rect [367.446 282.419 399.968 294.372] /A << /S /GoTo /D <74686D2E452E312E312E31> >> >>}
\MusicPDFWriteObject{7812}{<< /Type /Annot /Subtype /Link /Border [0 0 0] /H /I /C [1 0 0] /Rect [215.431 625.213 247.443 636.004] /A << /S /GoTo /D <74686D2E452E312E312E33> >> >>}
\MusicPDFWriteObject{7813}{<< /Type /Annot /Subtype /Link /Border [0 0 0] /H /I /C [1 0 0] /Rect [428.518 575.666 460.53 586.457] /A << /S /GoTo /D <74686D2E452E312E312E33> >> >>}
\MusicPDFWriteObject{7814}{<< /Type /Annot /Subtype /Link /Border [0 0 0] /H /I /C [1 0 0] /Rect [273.657 431.885 305.669 443.84] /A << /S /GoTo /D <74686D2E452E312E312E33> >> >>}
\MusicPDFWriteObject{7815}{<< /Type /Annot /Subtype /Link /Border [0 0 0] /H /I /C [1 0 0] /Rect [484.694 345.611 516.495 357.564] /A << /S /GoTo /D <74686D2E452E312E312E33> >> >>}
\MusicPDFWriteObject{7816}{<< /Type /Annot /Subtype /Link /Border [0 0 0] /H /I /C [1 0 0] /Rect [488.332 214.103 520.343 226.056] /A << /S /GoTo /D <6571756174696F6E2E452E312E312E31> >> >>}
\MusicPDFWriteObject{7817}{<< /Type /Annot /Subtype /Link /Border [0 0 0] /H /I /C [1 0 0] /Rect [464.814 135.198 497.426 147.153] /A << /S /GoTo /D <74686D2E452E322E312E31> >> >>}
\MusicPDFWriteObject{7818}{<< /Type /Annot /Subtype /Link /Border [0 0 0] /H /I /C [1 0 0] /Rect [283.186 122.048 299.994 134.001] /A << /S /GoTo /D <73656374696F6E2E452E32> >> >>}
\MusicPDFWriteObject{7819}{<< /Type /Annot /Subtype /Link /Border [0 0 0] /H /I /C [1 0 0] /Rect [358.576 122.048 391.188 134.001] /A << /S /GoTo /D <74686D2E452E322E312E32> >> >>}
\MusicPDFWriteObject{7820}{<< /Type /Annot /Subtype /Link /Border [0 0 0] /H /I /C [1 0 0] /Rect [217.549 109.451 248.961 120.241] /A << /S /GoTo /D <74686D2E452E322E312E34> >> >>}
\MusicPDFWriteObject{7821}{<< /Type /Annot /Subtype /Link /Border [0 0 0] /H /I /C [1 0 0] /Rect [382.58 421.618 415.192 433.571] /A << /S /GoTo /D <74686D2E452E312E312E33> >> >>}
\MusicPDFWriteObject{7822}{<< /Type /Annot /Subtype /Link /Border [0 0 0] /H /I /C [1 0 0] /Rect [512.603 363.888 544.945 375.843] /A << /S /GoTo /D <74686D2E452E312E312E34> >> >>}
\MusicPDFWriteObject{7823}{<< /Type /Annot /Subtype /Link /Border [0 0 0] /H /I /C [1 0 0] /Rect [176.323 321.8 211.766 330.654] /A << /S /GoTo /D <74686D2E452E312E312E33> >> >>}
\MusicPDFWriteObject{7824}{<< /Type /Annot /Subtype /Link /Border [0 0 0] /H /I /C [1 0 0] /Rect [183.974 266.706 216.076 279.299] /A << /S /GoTo /D <74686D2E452E312E312E32> >> >>}
\MusicPDFWriteObject{7825}{<< /Type /Annot /Subtype /Link /Border [0 0 0] /H /I /C [1 0 0] /Rect [309.844 227.808 341.255 238.598] /A << /S /GoTo /D <74686D2E452E312E312E33> >> >>}
\MusicPDFWriteObject{7826}{<< /Type /Annot /Subtype /Link /Border [0 0 0] /H /I /C [1 0 0] /Rect [457.801 227.808 489.212 238.598] /A << /S /GoTo /D <74686D2E452E312E312E34> >> >>}
\MusicPDFWriteObject{7827}{<< /Type /Annot /Subtype /Link /Border [0 0 0] /H /I /C [1 0 0] /Rect [367.282 200.953 398.693 212.905] /A << /S /GoTo /D <74686D2E452E312E312E34> >> >>}
\MusicPDFWriteObject{7828}{<< /Type /Annot /Subtype /Link /Border [0 0 0] /H /I /C [1 0 0] /Rect [71.004 188.355 102.415 199.146] /A << /S /GoTo /D <74686D2E452E312E312E33> >> >>}
\MusicPDFWriteObject{7829}{<< /Type /Annot /Subtype /Link /Border [0 0 0] /H /I /C [0 1 0] /Rect [472.384 660.822 496.684 669.69] /A << /S /GoTo /D <636974652E4A656368> >> >>}
\MusicPDFWriteObject{7830}{<< /Type /Annot /Subtype /Link /Border [0 0 0] /H /I /C [1 0 0] /Rect [399.402 645.181 431.174 657.134] /A << /S /GoTo /D <74686D2E452E312E312E36> >> >>}
\MusicPDFWriteObject{7831}{<< /Type /Annot /Subtype /Link /Border [0 0 0] /H /I /C [1 0 0] /Rect [151.53 575.568 184.142 587.521] /A << /S /GoTo /D <6571756174696F6E2E452E312E312E32> >> >>}
\MusicPDFWriteObject{7832}{<< /Type /Annot /Subtype /Link /Border [0 0 0] /H /I /C [1 0 0] /Rect [293.052 288.098 325.664 300.691] /A << /S /GoTo /D <6571756174696F6E2E452E312E312E35> >> >>}
\MusicPDFWriteObject{7833}{<< /Type /Annot /Subtype /Link /Border [0 0 0] /H /I /C [1 0 0] /Rect [97.805 247.594 129.216 260.619] /A << /S /GoTo /D <6571756174696F6E2E452E312E312E32> >> >>}
\MusicPDFWriteObject{7834}{<< /Type /Annot /Subtype /Link /Border [0 0 0] /H /I /C [0 1 0] /Rect [426.688 172.232 461.998 181.1] /A << /S /GoTo /D <636974652E4B756E656E> >> >>}
\MusicPDFWriteObject{7835}{<< /Type /Annot /Subtype /Link /Border [0 0 0] /H /I /C [0 1 0] /Rect [73.806 145.93 98.106 154.798] /A << /S /GoTo /D <636974652E4A656368> >> >>}
\MusicPDFWriteObject{7836}{<< /Type /Annot /Subtype /Link /Border [0 0 0] /H /I /C [0 1 0] /Rect [199.661 145.93 234.972 154.798] /A << /S /GoTo /D <636974652E4B756E656E> >> >>}
\MusicPDFWriteObject{7837}{<< /Type /Annot /Subtype /Link /Border [0 0 0] /H /I /C [0 1 0] /Rect [182.685 549.652 225.708 558.52] /A << /S /GoTo /D <636974652E426F752D4754> >> >>}
\MusicPDFWriteObject{7838}{<< /Type /Annot /Subtype /Link /Border [0 0 0] /H /I /C [0 1 0] /Rect [378.049 549.652 394.235 558.52] /A << /S /GoTo /D <636974652E5350> >> >>}
\MusicPDFWriteObject{7839}{<< /Type /Annot /Subtype /Link /Border [0 0 0] /H /I /C [0 1 0] /Rect [278.888 396.904 303.187 405.772] /A << /S /GoTo /D <636974652E4A656368> >> >>}
\MusicPDFWriteObject{7840}{<< /Type /Annot /Subtype /Link /Border [0 0 0] /H /I /C [0 1 0] /Rect [518.937 396.904 543.236 405.772] /A << /S /GoTo /D <636974652E4A656368> >> >>}
\MusicPDFWriteObject{7841}{<< /Type /Annot /Subtype /Link /Border [0 0 0] /H /I /C [1 0 0] /Rect [71.004 355.515 102.655 366.306] /A << /S /GoTo /D <74686D2E452E312E312E33> >> >>}
\MusicPDFWriteObject{7842}{<< /Type /Annot /Subtype /Link /Border [0 0 0] /H /I /C [1 0 0] /Rect [135.608 293.464 167.619 302.318] /A << /S /GoTo /D <74686D2E452E312E312E33> >> >>}
\MusicPDFWriteObject{7843}{<< /Type /Annot /Subtype /Link /Border [0 0 0] /H /I /C [1 0 0] /Rect [371.391 243.432 403.853 255.387] /A << /S /GoTo /D <74686D2E452E312E312E31> >> >>}
\MusicPDFWriteObject{7844}{<< /Type /Annot /Subtype /Link /Border [0 0 0] /H /I /C [1 0 0] /Rect [230.796 230.835 262.207 241.626] /A << /S /GoTo /D <74686D2E452E312E312E31> >> >>}
\MusicPDFWriteObject{7845}{<< /Type /Annot /Subtype /Link /Border [0 0 0] /H /I /C [1 0 0] /Rect [360.139 230.835 391.55 241.626] /A << /S /GoTo /D <74686D2E452E312E312E37> >> >>}
\MusicPDFWriteObject{7846}{<< /Type /Annot /Subtype /Link /Border [0 0 0] /H /I /C [0 1 0] /Rect [409.07 660.822 433.494 669.69] /A << /S /GoTo /D <636974652E53745365> >> >>}
\MusicPDFWriteObject{7847}{<< /Type /Annot /Subtype /Link /Border [0 0 0] /H /I /C [1 0 0] /Rect [146.472 638.575 171.115 649.424] /A << /S /GoTo /D <73756273656374696F6E2E452E322E34> >> >>}
\MusicPDFWriteObject{7848}{<< /Type /Annot /Subtype /Link /Border [0 0 0] /H /I /C [1 0 0] /Rect [374.655 625.425 398.474 636.273] /A << /S /GoTo /D <73756273656374696F6E2E452E322E33> >> >>}
\MusicPDFWriteObject{7849}{<< /Type /Annot /Subtype /Link /Border [0 0 0] /H /I /C [1 0 0] /Rect [340.373 612.274 364.638 623.122] /A << /S /GoTo /D <73756273656374696F6E2E452E322E33> >> >>}
\MusicPDFWriteObject{7850}{<< /Type /Annot /Subtype /Link /Border [0 0 0] /H /I /C [0 1 0] /Rect [487.593 566.926 512.37 575.794] /A << /S /GoTo /D <636974652E4C474C> >> >>}
\MusicPDFWriteObject{7851}{<< /Type /Annot /Subtype /Link /Border [0 0 0] /H /I /C [0 1 0] /Rect [92.143 553.775 118.372 562.643] /A << /S /GoTo /D <636974652E4D6343> >> >>}
\MusicPDFWriteObject{7852}{<< /Type /Annot /Subtype /Link /Border [0 0 0] /H /I /C [0 1 0] /Rect [143.622 553.775 161.149 562.643] /A << /S /GoTo /D <636974652E537472> >> >>}
\MusicPDFWriteObject{7853}{<< /Type /Annot /Subtype /Link /Border [0 0 0] /H /I /C [0 1 0] /Rect [201.955 553.775 228.924 562.643] /A << /S /GoTo /D <636974652E52657A6B> >> >>}
\MusicPDFWriteObject{7854}{<< /Type /Annot /Subtype /Link /Border [0 0 0] /H /I /C [0 1 0] /Rect [331.454 553.775 354.224 562.643] /A << /S /GoTo /D <636974652E454C53> >> >>}
\MusicPDFWriteObject{7855}{<< /Type /Annot /Subtype /Link /Border [0 0 0] /H /I /C [1 0 0] /Rect [161.78 498.682 193.791 510.635] /A << /S /GoTo /D <6E6F746174696F6E2E452E302E302E31> >> >>}
\MusicPDFWriteObject{7856}{<< /Type /Annot /Subtype /Link /Border [0 0 0] /H /I /C [1 0 0] /Rect [499.934 95.747 531.345 107.699] /A << /S /GoTo /D <74686D2E452E312E322E31> >> >>}
\MusicPDFWriteObject{7857}{<< /Type /Annot /Subtype /Link /Border [0 0 0] /H /I /C [1 0 0] /Rect [106.098 656.477 137.689 670.286] /A << /S /GoTo /D <74686D2E452E312E312E33> >> >>}
\MusicPDFWriteObject{7858}{<< /Type /Annot /Subtype /Link /Border [0 0 0] /H /I /C [1 0 0] /Rect [435.539 606.031 468.151 618.484] /A << /S /GoTo /D <74686D2E452E312E312E33> >> >>}
\MusicPDFWriteObject{7859}{<< /Type /Annot /Subtype /Link /Border [0 0 0] /H /I /C [1 0 0] /Rect [460.562 593.379 492.483 605.332] /A << /S /GoTo /D <74686D2E452E312E312E33> >> >>}
\MusicPDFWriteObject{7860}{<< /Type /Annot /Subtype /Link /Border [0 0 0] /H /I /C [1 0 0] /Rect [122.658 538.513 154.67 547.367] /A << /S /GoTo /D <74686D2E452E322E342E39> >> >>}
\MusicPDFWriteObject{7861}{<< /Type /Annot /Subtype /Link /Border [0 0 0] /H /I /C [0 1 0] /Rect [475.201 520.61 491.387 529.478] /A << /S /GoTo /D <636974652E5350> >> >>}
\MusicPDFWriteObject{7862}{<< /Type /Annot /Subtype /Link /Border [0 0 0] /H /I /C [1 0 0] /Rect [142.861 504.969 175.473 516.922] /A << /S /GoTo /D <74686D2E452E312E322E32> >> >>}
\MusicPDFWriteObject{7863}{<< /Type /Annot /Subtype /Link /Border [0 0 0] /H /I /C [1 0 0] /Rect [497.91 492.372 530.131 503.163] /A << /S /GoTo /D <74686D2E452E312E322E33> >> >>}
\MusicPDFWriteObject{7864}{<< /Type /Annot /Subtype /Link /Border [0 0 0] /H /I /C [1 0 0] /Rect [233.222 248.114 265.234 258.905] /A << /S /GoTo /D <74686D2E452E312E322E34> >> >>}
\MusicPDFWriteObject{7865}{<< /Type /Annot /Subtype /Link /Border [0 0 0] /H /I /C [1 0 0] /Rect [190.246 232.318 222.857 244.27] /A << /S /GoTo /D <74686D2E452E312E322E34> >> >>}
\MusicPDFWriteObject{7866}{<< /Type /Annot /Subtype /Link /Border [0 0 0] /H /I /C [0 1 0] /Rect [262.082 234.808 279.61 243.676] /A << /S /GoTo /D <636974652E537472> >> >>}
\MusicPDFWriteObject{7867}{<< /Type /Annot /Subtype /Link /Border [0 0 0] /H /I /C [1 0 0] /Rect [147.46 219.167 179.922 231.12] /A << /S /GoTo /D <74686D2E452E312E322E34> >> >>}
\MusicPDFWriteObject{7868}{<< /Type /Annot /Subtype /Link /Border [0 0 0] /H /I /C [1 0 0] /Rect [376.03 219.167 408.492 231.12] /A << /S /GoTo /D <74686D2E452E312E322E32> >> >>}
\MusicPDFWriteObject{7869}{<< /Type /Annot /Subtype /Link /Border [0 0 0] /H /I /C [1 0 0] /Rect [409.199 179.714 441.631 191.669] /A << /S /GoTo /D <74686D2E452E312E322E32> >> >>}
\MusicPDFWriteObject{7870}{<< /Type /Annot /Subtype /Link /Border [0 0 0] /H /I /C [1 0 0] /Rect [138.098 113.65 170.109 125.603] /A << /S /GoTo /D <74686D2E452E312E322E32> >> >>}
\MusicPDFWriteObject{7871}{<< /Type /Annot /Subtype /Link /Border [0 0 0] /H /I /C [0 1 0] /Rect [343.86 85.086 363.263 93.954] /A << /S /GoTo /D <636974652E4B656C> >> >>}
\MusicPDFWriteObject{7872}{<< /Type /Annot /Subtype /Link /Border [0 0 0] /H /I /C [0 1 0] /Rect [419.189 85.086 449.603 93.954] /A << /S /GoTo /D <636974652E4D61634C> >> >>}
\MusicPDFWriteObject{7873}{<< /Type /Annot /Subtype /Link /Border [0 0 0] /H /I /C [0 1 0] /Rect [73.759 673.973 104.173 682.84] /A << /S /GoTo /D <636974652E4D61634C> >> >>}
\MusicPDFWriteObject{7874}{<< /Type /Annot /Subtype /Link /Border [0 0 0] /H /I /C [0 1 0] /Rect [399.11 634.52 429.523 643.388] /A << /S /GoTo /D <636974652E4D61634C> >> >>}
\MusicPDFWriteObject{7875}{<< /Type /Annot /Subtype /Link /Border [0 0 0] /H /I /C [0 1 0] /Rect [173.57 621.369 203.984 630.237] /A << /S /GoTo /D <636974652E4D61634C> >> >>}
\MusicPDFWriteObject{7876}{<< /Type /Annot /Subtype /Link /Border [0 0 0] /H /I /C [0 1 0] /Rect [289.469 621.369 307.965 630.237] /A << /S /GoTo /D <636974652E454B> >> >>}
\MusicPDFWriteObject{7877}{<< /Type /Annot /Subtype /Link /Border [0 0 0] /H /I /C [0 1 0] /Rect [184.669 608.219 215.083 617.087] /A << /S /GoTo /D <636974652E4D61634C> >> >>}
\MusicPDFWriteObject{7878}{<< /Type /Annot /Subtype /Link /Border [0 0 0] /H /I /C [0 1 0] /Rect [120.556 555.616 138.616 564.484] /A << /S /GoTo /D <636974652E537465> >> >>}
\MusicPDFWriteObject{7879}{<< /Type /Annot /Subtype /Link /Border [0 0 0] /H /I /C [0 1 0] /Rect [180.339 555.616 210.753 564.484] /A << /S /GoTo /D <636974652E4D61634C> >> >>}
\MusicPDFWriteObject{7880}{<< /Type /Annot /Subtype /Link /Border [0 0 0] /H /I /C [0 1 0] /Rect [127.243 529.314 150.014 538.182] /A << /S /GoTo /D <636974652E454C53> >> >>}
\MusicPDFWriteObject{7881}{<< /Type /Annot /Subtype /Link /Border [0 0 0] /H /I /C [0 1 0] /Rect [397.779 529.314 433.92 538.182] /A << /S /GoTo /D <636974652E4D61634C32> >> >>}
\MusicPDFWriteObject{7882}{<< /Type /Annot /Subtype /Link /Border [0 0 0] /H /I /C [0 1 0] /Rect [341.36 503.013 371.773 511.881] /A << /S /GoTo /D <636974652E4D61634C> >> >>}
\MusicPDFWriteObject{7883}{<< /Type /Annot /Subtype /Link /Border [0 0 0] /H /I /C [0 1 0] /Rect [288.301 489.862 318.714 498.73] /A << /S /GoTo /D <636974652E4D61634C> >> >>}
\MusicPDFWriteObject{7884}{<< /Type /Annot /Subtype /Link /Border [0 0 0] /H /I /C [0 1 0] /Rect [438.472 489.862 468.886 498.73] /A << /S /GoTo /D <636974652E4D61634C> >> >>}
\MusicPDFWriteObject{7885}{<< /Type /Annot /Subtype /Link /Border [0 0 0] /H /I /C [0 1 0] /Rect [131.092 476.711 167.233 485.579] /A << /S /GoTo /D <636974652E4D61634C32> >> >>}
\MusicPDFWriteObject{7886}{<< /Type /Annot /Subtype /Link /Border [0 0 0] /H /I /C [0 1 0] /Rect [300.958 476.711 337.099 485.579] /A << /S /GoTo /D <636974652E4D61634C32> >> >>}
\MusicPDFWriteObject{7887}{<< /Type /Annot /Subtype /Link /Border [0 0 0] /H /I /C [0 1 0] /Rect [507.018 463.56 543.159 472.428] /A << /S /GoTo /D <636974652E4D61634C32> >> >>}
\MusicPDFWriteObject{7888}{<< /Type /Annot /Subtype /Link /Border [0 0 0] /H /I /C [0 1 0] /Rect [174.744 450.41 210.885 459.277] /A << /S /GoTo /D <636974652E4D61634C32> >> >>}
\MusicPDFWriteObject{7889}{<< /Type /Annot /Subtype /Link /Border [0 0 0] /H /I /C [1 0 0] /Rect [118.414 422.171 151.026 432.962] /A << /S /GoTo /D <74686D2E452E312E312E37> >> >>}
\MusicPDFWriteObject{7890}{<< /Type /Annot /Subtype /Link /Border [0 0 0] /H /I /C [0 1 0] /Rect [339.858 410.957 364.635 419.825] /A << /S /GoTo /D <636974652E4C474C> >> >>}
\MusicPDFWriteObject{7891}{<< /Type /Annot /Subtype /Link /Border [0 0 0] /H /I /C [0 1 0] /Rect [389.516 410.957 412.287 419.825] /A << /S /GoTo /D <636974652E454C53> >> >>}
\MusicPDFWriteObject{7892}{<< /Type /Annot /Subtype /Link /Border [0 0 0] /H /I /C [0 1 0] /Rect [437.169 410.957 454.696 419.825] /A << /S /GoTo /D <636974652E537472> >> >>}
\MusicPDFWriteObject{7893}{<< /Type /Annot /Subtype /Link /Border [0 0 0] /H /I /C [0 1 0] /Rect [498.705 410.957 525.674 419.825] /A << /S /GoTo /D <636974652E52657A6B> >> >>}
\MusicPDFWriteObject{7894}{<< /Type /Annot /Subtype /Link /Border [0 0 0] /H /I /C [0 1 0] /Rect [91.423 261.639 107.609 270.507] /A << /S /GoTo /D <636974652E5350> >> >>}
\MusicPDFWriteObject{7895}{<< /Type /Annot /Subtype /Link /Border [0 0 0] /H /I /C [0 1 0] /Rect [177.858 261.639 194.043 270.507] /A << /S /GoTo /D <636974652E5350> >> >>}
\MusicPDFWriteObject{7896}{<< /Type /Annot /Subtype /Link /Border [0 0 0] /H /I /C [0 1 0] /Rect [293.335 261.639 310.475 270.507] /A << /S /GoTo /D <636974652E4475> >> >>}
\MusicPDFWriteObject{7897}{<< /Type /Annot /Subtype /Link /Border [0 0 0] /H /I /C [0 1 0] /Rect [499.321 261.639 515.541 270.507] /A << /S /GoTo /D <636974652E4B65> >> >>}
\MusicPDFWriteObject{7898}{<< /Type /Annot /Subtype /Link /Border [0 0 0] /H /I /C [1 0 0] /Rect [513.254 232.847 544.845 244.802] /A << /S /GoTo /D <74686D2E452E312E322E38> >> >>}
\MusicPDFWriteObject{7899}{<< /Type /Annot /Subtype /Link /Border [0 0 0] /H /I /C [1 0 0] /Rect [381.815 135.199 413.226 147.152] /A << /S /GoTo /D <74686D2E452E312E322E36> >> >>}
\MusicPDFWriteObject{7900}{<< /Type /Annot /Subtype /Link /Border [0 0 0] /H /I /C [1 0 0] /Rect [223.062 122.048 254.804 134.003] /A << /S /GoTo /D <74686D2E452E312E322E36> >> >>}
\MusicPDFWriteObject{7901}{<< /Type /Annot /Subtype /Link /Border [0 0 0] /H /I /C [0 1 0] /Rect [342.506 98.237 372.919 107.105] /A << /S /GoTo /D <636974652E4D61634C> >> >>}
\MusicPDFWriteObject{7902}{<< /Type /Annot /Subtype /Link /Border [0 0 0] /H /I /C [0 1 0] /Rect [194.097 85.086 224.511 93.954] /A << /S /GoTo /D <636974652E4D61634C> >> >>}
\MusicPDFWriteObject{7903}{<< /Type /Annot /Subtype /Link /Border [0 0 0] /H /I /C [1 0 0] /Rect [106.556 671.482 143.548 683.435] /A << /S /GoTo /D <74686D2E452E312E322E3132> >> >>}
\MusicPDFWriteObject{7904}{<< /Type /Annot /Subtype /Link /Border [0 0 0] /H /I /C [1 0 0] /Rect [217.861 671.482 254.853 683.435] /A << /S /GoTo /D <74686D2E452E312E322E3133> >> >>}
\MusicPDFWriteObject{7905}{<< /Type /Annot /Subtype /Link /Border [0 0 0] /H /I /C [0 1 0] /Rect [520.413 634.52 543.184 643.388] /A << /S /GoTo /D <636974652E454C53> >> >>}
\MusicPDFWriteObject{7906}{<< /Type /Annot /Subtype /Link /Border [0 0 0] /H /I /C [0 1 0] /Rect [217.393 608.219 240.164 617.087] /A << /S /GoTo /D <636974652E454C53> >> >>}
\MusicPDFWriteObject{7907}{<< /Type /Annot /Subtype /Link /Border [0 0 0] /H /I /C [0 1 0] /Rect [368.598 608.219 393.375 617.087] /A << /S /GoTo /D <636974652E4C474C> >> >>}
\MusicPDFWriteObject{7908}{<< /Type /Annot /Subtype /Link /Border [0 0 0] /H /I /C [1 0 0] /Rect [236.537 579.427 273.528 591.38] /A << /S /GoTo /D <74686D2E452E322E322E3130> >> >>}
\MusicPDFWriteObject{7909}{<< /Type /Annot /Subtype /Link /Border [0 0 0] /H /I /C [1 0 0] /Rect [404.225 465.676 435.636 477.631] /A << /S /GoTo /D <6571756174696F6E2E452E312E322E34> >> >>}
\MusicPDFWriteObject{7910}{<< /Type /Annot /Subtype /Link /Border [0 0 0] /H /I /C [0 1 0] /Rect [89.402 386.73 113.702 395.598] /A << /S /GoTo /D <636974652E4A656368> >> >>}
\MusicPDFWriteObject{7911}{<< /Type /Annot /Subtype /Link /Border [0 0 0] /H /I /C [0 1 0] /Rect [335.151 386.73 359.45 395.598] /A << /S /GoTo /D <636974652E4A656368> >> >>}
\MusicPDFWriteObject{7912}{<< /Type /Annot /Subtype /Link /Border [0 0 0] /H /I /C [1 0 0] /Rect [176.817 334.661 208.948 346.617] /A << /S /GoTo /D <74686D2E452E312E322E36> >> >>}
\MusicPDFWriteObject{7913}{<< /Type /Annot /Subtype /Link /Border [0 0 0] /H /I /C [1 0 0] /Rect [420.983 334.661 453.114 346.617] /A << /S /GoTo /D <6571756174696F6E2E452E312E322E36> >> >>}
\MusicPDFWriteObject{7914}{<< /Type /Annot /Subtype /Link /Border [0 0 0] /H /I /C [1 0 0] /Rect [250.702 239.112 282.714 249.903] /A << /S /GoTo /D <74686D2E452E312E322E36> >> >>}
\MusicPDFWriteObject{7915}{<< /Type /Annot /Subtype /Link /Border [0 0 0] /H /I /C [0 1 0] /Rect [378.382 203.9 403.158 212.768] /A << /S /GoTo /D <636974652E4C474C> >> >>}
\MusicPDFWriteObject{7916}{<< /Type /Annot /Subtype /Link /Border [0 0 0] /H /I /C [0 1 0] /Rect [143.1 190.749 167.877 199.617] /A << /S /GoTo /D <636974652E4C474C> >> >>}
\MusicPDFWriteObject{7917}{<< /Type /Annot /Subtype /Link /Border [0 0 0] /H /I /C [0 1 0] /Rect [246.536 190.749 276.95 199.617] /A << /S /GoTo /D <636974652E4D61634C> >> >>}
\MusicPDFWriteObject{7918}{<< /Type /Annot /Subtype /Link /Border [0 0 0] /H /I /C [0 1 0] /Rect [364.161 190.749 381.688 199.617] /A << /S /GoTo /D <636974652E537472> >> >>}
\MusicPDFWriteObject{7919}{<< /Type /Annot /Subtype /Link /Border [0 0 0] /H /I /C [0 1 0] /Rect [466.297 190.749 493.266 199.617] /A << /S /GoTo /D <636974652E52657A6B> >> >>}
\MusicPDFWriteObject{7920}{<< /Type /Annot /Subtype /Link /Border [0 0 0] /H /I /C [1 0 0] /Rect [204.503 454.821 241.494 465.282] /A << /S /GoTo /D <74686D2E452E312E322E3130> >> >>}
\MusicPDFWriteObject{7921}{<< /Type /Annot /Subtype /Link /Border [0 0 0] /H /I /C [0 1 0] /Rect [271.951 406.737 302.365 415.605] /A << /S /GoTo /D <636974652E4D61634C> >> >>}
\MusicPDFWriteObject{7922}{<< /Type /Annot /Subtype /Link /Border [0 0 0] /H /I /C [0 1 0] /Rect [401.797 406.737 419.857 415.605] /A << /S /GoTo /D <636974652E537465> >> >>}
\MusicPDFWriteObject{7923}{<< /Type /Annot /Subtype /Link /Border [0 0 0] /H /I /C [0 1 0] /Rect [73.762 393.586 98.539 402.454] /A << /S /GoTo /D <636974652E4C474C> >> >>}
\MusicPDFWriteObject{7924}{<< /Type /Annot /Subtype /Link /Border [0 0 0] /H /I /C [0 1 0] /Rect [152.312 393.586 169.84 402.454] /A << /S /GoTo /D <636974652E537472> >> >>}
\MusicPDFWriteObject{7925}{<< /Type /Annot /Subtype /Link /Border [0 0 0] /H /I /C [0 1 0] /Rect [248.072 393.586 275.041 402.454] /A << /S /GoTo /D <636974652E52657A6B> >> >>}
\MusicPDFWriteObject{7926}{<< /Type /Annot /Subtype /Link /Border [0 0 0] /H /I /C [0 1 0] /Rect [169.5 367.285 194.277 376.152] /A << /S /GoTo /D <636974652E4C474C> >> >>}
\MusicPDFWriteObject{7927}{<< /Type /Annot /Subtype /Link /Border [0 0 0] /H /I /C [0 1 0] /Rect [219.399 367.285 236.927 376.152] /A << /S /GoTo /D <636974652E537472> >> >>}
\MusicPDFWriteObject{7928}{<< /Type /Annot /Subtype /Link /Border [0 0 0] /H /I /C [0 1 0] /Rect [278.65 367.285 305.619 376.152] /A << /S /GoTo /D <636974652E52657A6B> >> >>}
\MusicPDFWriteObject{7929}{<< /Type /Annot /Subtype /Link /Border [0 0 0] /H /I /C [0 1 0] /Rect [410.268 333.416 435.045 342.284] /A << /S /GoTo /D <636974652E4C474C> >> >>}
\MusicPDFWriteObject{7930}{<< /Type /Annot /Subtype /Link /Border [0 0 0] /H /I /C [1 0 0] /Rect [420.968 317.775 453.34 329.728] /A << /S /GoTo /D <74686D2E452E312E322E36> >> >>}
\MusicPDFWriteObject{7931}{<< /Type /Annot /Subtype /Link /Border [0 0 0] /H /I /C [0 1 0] /Rect [502.807 320.265 527.583 329.133] /A << /S /GoTo /D <636974652E4C474C> >> >>}
\MusicPDFWriteObject{7932}{<< /Type /Annot /Subtype /Link /Border [0 0 0] /H /I /C [0 1 0] /Rect [254.112 307.114 271.64 315.982] /A << /S /GoTo /D <636974652E537472> >> >>}
\MusicPDFWriteObject{7933}{<< /Type /Annot /Subtype /Link /Border [0 0 0] /H /I /C [0 1 0] /Rect [169.229 293.964 196.198 302.831] /A << /S /GoTo /D <636974652E52657A6B> >> >>}
\MusicPDFWriteObject{7934}{<< /Type /Annot /Subtype /Link /Border [0 0 0] /H /I /C [0 1 0] /Rect [73.826 280.813 98.602 289.681] /A << /S /GoTo /D <636974652E4C474C> >> >>}
\MusicPDFWriteObject{7935}{<< /Type /Annot /Subtype /Link /Border [0 0 0] /H /I /C [0 1 0] /Rect [173.151 280.813 190.679 289.681] /A << /S /GoTo /D <636974652E537472> >> >>}
\MusicPDFWriteObject{7936}{<< /Type /Annot /Subtype /Link /Border [0 0 0] /H /I /C [0 1 0] /Rect [478.405 280.813 505.374 289.681] /A << /S /GoTo /D <636974652E52657A6B> >> >>}
\MusicPDFWriteObject{7937}{<< /Type /Annot /Subtype /Link /Border [0 0 0] /H /I /C [0 1 0] /Rect [319.601 267.662 342.372 276.53] /A << /S /GoTo /D <636974652E454C53> >> >>}
\MusicPDFWriteObject{7938}{<< /Type /Annot /Subtype /Link /Border [0 0 0] /H /I /C [0 1 0] /Rect [85.226 241.361 112.195 250.228] /A << /S /GoTo /D <636974652E52657A6B> >> >>}
\MusicPDFWriteObject{7939}{<< /Type /Annot /Subtype /Link /Border [0 0 0] /H /I /C [0 1 0] /Rect [503.099 241.361 527.875 250.228] /A << /S /GoTo /D <636974652E4C474C> >> >>}
\MusicPDFWriteObject{7940}{<< /Type /Annot /Subtype /Link /Border [0 0 0] /H /I /C [0 1 0] /Rect [107.525 228.21 125.053 237.078] /A << /S /GoTo /D <636974652E537472> >> >>}
\MusicPDFWriteObject{7941}{<< /Type /Annot /Subtype /Link /Border [0 0 0] /H /I /C [1 0 0] /Rect [279.281 441.805 316.973 453.757] /A << /S /GoTo /D <6571756174696F6E2E452E312E322E3130> >> >>}
\MusicPDFWriteObject{7942}{<< /Type /Annot /Subtype /Link /Border [0 0 0] /H /I /C [1 0 0] /Rect [151.42 428.654 182.831 440.607] /A << /S /GoTo /D <6571756174696F6E2E452E312E322E39> >> >>}
\MusicPDFWriteObject{7943}{<< /Type /Annot /Subtype /Link /Border [0 0 0] /H /I /C [1 0 0] /Rect [316.922 428.654 353.214 440.607] /A << /S /GoTo /D <6571756174696F6E2E452E312E322E3131> >> >>}
\MusicPDFWriteObject{7944}{<< /Type /Annot /Subtype /Link /Border [0 0 0] /H /I /C [1 0 0] /Rect [88.711 684.633 125.702 696.586] /A << /S /GoTo /D <6571756174696F6E2E452E312E322E3135> >> >>}
\MusicPDFWriteObject{7945}{<< /Type /Annot /Subtype /Link /Border [0 0 0] /H /I /C [1 0 0] /Rect [314.475 684.633 351.467 696.586] /A << /S /GoTo /D <6571756174696F6E2E452E312E322E3137> >> >>}
\MusicPDFWriteObject{7946}{<< /Type /Annot /Subtype /Link /Border [0 0 0] /H /I /C [1 0 0] /Rect [449.635 658.165 486.767 671.57] /A << /S /GoTo /D <6571756174696F6E2E452E312E322E3136> >> >>}
\MusicPDFWriteObject{7947}{<< /Type /Annot /Subtype /Link /Border [0 0 0] /H /I /C [1 0 0] /Rect [71.004 308.259 103.616 320.211] /A << /S /GoTo /D <74686D2E452E312E322E31> >> >>}
\MusicPDFWriteObject{7948}{<< /Type /Annot /Subtype /Link /Border [0 0 0] /H /I /C [1 0 0] /Rect [491.883 281.957 528.245 293.912] /A << /S /GoTo /D <74686D2E452E312E322E3130> >> >>}
\MusicPDFWriteObject{7949}{<< /Type /Annot /Subtype /Link /Border [0 0 0] /H /I /C [1 0 0] /Rect [127.159 167.118 164.851 179.506] /A << /S /GoTo /D <6571756174696F6E2E452E312E322E3231> >> >>}
\MusicPDFWriteObject{7950}{<< /Type /Annot /Subtype /Link /Border [0 0 0] /H /I /C [1 0 0] /Rect [239.287 568.903 271.899 580.875] /A << /S /GoTo /D <74686D2E452E312E312E33> >> >>}
\MusicPDFWriteObject{7951}{<< /Type /Annot /Subtype /Link /Border [0 0 0] /H /I /C [0 1 0] /Rect [505.019 571.393 521.205 580.26] /A << /S /GoTo /D <636974652E5350> >> >>}
\MusicPDFWriteObject{7952}{<< /Type /Annot /Subtype /Link /Border [0 0 0] /H /I /C [1 0 0] /Rect [249.318 377.691 287.01 389.643] /A << /S /GoTo /D <6571756174696F6E2E452E312E322E3130> >> >>}
\MusicPDFWriteObject{7953}{<< /Type /Annot /Subtype /Link /Border [0 0 0] /H /I /C [1 0 0] /Rect [461.581 377.691 494.193 389.643] /A << /S /GoTo /D <6571756174696F6E2E452E312E322E39> >> >>}
\MusicPDFWriteObject{7954}{<< /Type /Annot /Subtype /Link /Border [0 0 0] /H /I /C [1 0 0] /Rect [339.979 328.112 377.671 340.065] /A << /S /GoTo /D <6571756174696F6E2E452E312E322E3232> >> >>}
\MusicPDFWriteObject{7955}{<< /Type /Annot /Subtype /Link /Border [0 0 0] /H /I /C [1 0 0] /Rect [372.178 314.961 409.87 326.934] /A << /S /GoTo /D <6571756174696F6E2E452E312E322E3233> >> >>}
\MusicPDFWriteObject{7956}{<< /Type /Annot /Subtype /Link /Border [0 0 0] /H /I /C [1 0 0] /Rect [73.669 163.202 110.276 175.155] /A << /S /GoTo /D <6571756174696F6E2E452E312E322E3237> >> >>}
\MusicPDFWriteObject{7957}{<< /Type /Annot /Subtype /Link /Border [0 0 0] /H /I /C [1 0 0] /Rect [450.631 163.202 487.238 175.155] /A << /S /GoTo /D <6571756174696F6E2E452E312E322E3237> >> >>}
\MusicPDFWriteObject{7958}{<< /Type /Annot /Subtype /Link /Border [0 0 0] /H /I /C [1 0 0] /Rect [348.777 150.05 385.069 162.006] /A << /S /GoTo /D <6571756174696F6E2E452E312E322E3238> >> >>}
\MusicPDFWriteObject{7959}{<< /Type /Annot /Subtype /Link /Border [0 0 0] /H /I /C [1 0 0] /Rect [233.529 123.749 239.492 136.607] /A << /S /GoTo /D <48666F6F746E6F74652E313534> >> >>}
\MusicPDFWriteObject{7960}{<< /Type /Annot /Subtype /Link /Border [0 0 0] /H /I /C [1 0 0] /Rect [143.579 93.057 174.751 104.644] /A << /S /GoTo /D <6571756174696F6E2E452E312E322E3238> >> >>}
\MusicPDFWriteObject{7961}{<< /Type /Annot /Subtype /Link /Border [0 0 0] /H /I /C [0 1 0] /Rect [222.015 687.123 252.429 695.991] /A << /S /GoTo /D <636974652E4D61634C> >> >>}
\MusicPDFWriteObject{7962}{<< /Type /Annot /Subtype /Link /Border [0 0 0] /H /I /C [0 1 0] /Rect [218.322 560.368 235.849 569.236] /A << /S /GoTo /D <636974652E537472> >> >>}
\MusicPDFWriteObject{7963}{<< /Type /Annot /Subtype /Link /Border [0 0 0] /H /I /C [0 1 0] /Rect [340.237 560.368 357.764 569.236] /A << /S /GoTo /D <636974652E537472> >> >>}
\MusicPDFWriteObject{7964}{<< /Type /Annot /Subtype /Link /Border [0 0 0] /H /I /C [0 1 0] /Rect [218.853 547.218 236.381 556.086] /A << /S /GoTo /D <636974652E537472> >> >>}
\MusicPDFWriteObject{7965}{<< /Type /Annot /Subtype /Link /Border [0 0 0] /H /I /C [1 0 0] /Rect [469.694 544.728 507.386 556.68] /A << /S /GoTo /D <74686D2E452E312E322E3132> >> >>}
\MusicPDFWriteObject{7966}{<< /Type /Annot /Subtype /Link /Border [0 0 0] /H /I /C [1 0 0] /Rect [179.374 513.002 216.365 523.793] /A << /S /GoTo /D <74686D2E452E312E322E3132> >> >>}
\MusicPDFWriteObject{7967}{<< /Type /Annot /Subtype /Link /Border [0 0 0] /H /I /C [0 1 0] /Rect [314.252 478.511 344.665 487.379] /A << /S /GoTo /D <636974652E4D61634C> >> >>}
\MusicPDFWriteObject{7968}{<< /Type /Annot /Subtype /Link /Border [0 0 0] /H /I /C [0 1 0] /Rect [423.848 478.511 442.344 487.379] /A << /S /GoTo /D <636974652E454B> >> >>}
\MusicPDFWriteObject{7969}{<< /Type /Annot /Subtype /Link /Border [0 0 0] /H /I /C [0 1 0] /Rect [376.328 465.36 406.741 474.228] /A << /S /GoTo /D <636974652E4D61634C> >> >>}
\MusicPDFWriteObject{7970}{<< /Type /Annot /Subtype /Link /Border [0 0 0] /H /I /C [1 0 0] /Rect [257.542 431.144 294.743 441.935] /A << /S /GoTo /D <74686D2E452E312E322E3133> >> >>}
\MusicPDFWriteObject{7971}{<< /Type /Annot /Subtype /Link /Border [0 0 0] /H /I /C [0 1 0] /Rect [73.77 406.78 98.547 415.647] /A << /S /GoTo /D <636974652E4C474C> >> >>}
\MusicPDFWriteObject{7972}{<< /Type /Annot /Subtype /Link /Border [0 0 0] /H /I /C [1 0 0] /Rect [498.021 272.418 530.332 284.371] /A << /S /GoTo /D <74686D2E452E312E322E32> >> >>}
\MusicPDFWriteObject{7973}{<< /Type /Annot /Subtype /Link /Border [0 0 0] /H /I /C [1 0 0] /Rect [117.683 206.663 149.545 218.636] /A << /S /GoTo /D <74686D2E452E312E322E32> >> >>}
\MusicPDFWriteObject{7974}{<< /Type /Annot /Subtype /Link /Border [0 0 0] /H /I /C [1 0 0] /Rect [474.789 178.675 506.621 192.955] /A << /S /GoTo /D <74686D2E452E312E322E32> >> >>}
\MusicPDFWriteObject{7975}{<< /Type /Annot /Subtype /Link /Border [0 0 0] /H /I /C [1 0 0] /Rect [299.58 122.601 337.272 133.392] /A << /S /GoTo /D <74686D2E452E312E322E3134> >> >>}
\MusicPDFWriteObject{7976}{<< /Type /Annot /Subtype /Link /Border [0 0 0] /H /I /C [0 1 0] /Rect [308.849 98.237 335.161 107.105] /A << /S /GoTo /D <636974652E426F7232> >> >>}
\MusicPDFWriteObject{7977}{<< /Type /Annot /Subtype /Link /Border [0 0 0] /H /I /C [0 1 0] /Rect [210.858 85.086 237.17 93.954] /A << /S /GoTo /D <636974652E426F7232> >> >>}
\MusicPDFWriteObject{7978}{<< /Type /Annot /Subtype /Link /Border [0 0 0] /H /I /C [1 0 0] /Rect [507.469 672.036 544.845 682.827] /A << /S /GoTo /D <74686D2E452E312E322E3136> >> >>}
\MusicPDFWriteObject{7979}{<< /Type /Annot /Subtype /Link /Border [0 0 0] /H /I /C [1 0 0] /Rect [249.591 645.734 265.818 656.582] /A << /S /GoTo /D <73656374696F6E2E452E32> >> >>}
\MusicPDFWriteObject{7980}{<< /Type /Annot /Subtype /Link /Border [0 0 0] /H /I /C [1 0 0] /Rect [361.676 645.734 397.968 656.582] /A << /S /GoTo /D <74686D2E452E312E322E3137> >> >>}
\MusicPDFWriteObject{7981}{<< /Type /Annot /Subtype /Link /Border [0 0 0] /H /I /C [0 1 0] /Rect [402.199 634.52 426.975 643.388] /A << /S /GoTo /D <636974652E4C474C> >> >>}
\MusicPDFWriteObject{7982}{<< /Type /Annot /Subtype /Link /Border [0 0 0] /H /I /C [0 1 0] /Rect [199.757 615.484 226.069 624.352] /A << /S /GoTo /D <636974652E426F7232> >> >>}
\MusicPDFWriteObject{7983}{<< /Type /Annot /Subtype /Link /Border [0 0 0] /H /I /C [0 1 0] /Rect [268.328 504.393 285.856 513.261] /A << /S /GoTo /D <636974652E537472> >> >>}
\MusicPDFWriteObject{7984}{<< /Type /Annot /Subtype /Link /Border [0 0 0] /H /I /C [1 0 0] /Rect [223.558 488.752 256.17 500.705] /A << /S /GoTo /D <74686D2E452E312E322E32> >> >>}
\MusicPDFWriteObject{7985}{<< /Type /Annot /Subtype /Link /Border [0 0 0] /H /I /C [1 0 0] /Rect [390.608 363.912 428.3 375.331] /A << /S /GoTo /D <74686D2E452E312E322E3134> >> >>}
\MusicPDFWriteObject{7986}{<< /Type /Annot /Subtype /Link /Border [0 0 0] /H /I /C [1 0 0] /Rect [321.69 349.848 359.381 362.163] /A << /S /GoTo /D <74686D2E452E312E322E3132> >> >>}
\MusicPDFWriteObject{7987}{<< /Type /Annot /Subtype /Link /Border [0 0 0] /H /I /C [1 0 0] /Rect [181.092 323.907 213.103 335.859] /A << /S /GoTo /D <74686D2E452E312E322E32> >> >>}
\MusicPDFWriteObject{7988}{<< /Type /Annot /Subtype /Link /Border [0 0 0] /H /I /C [1 0 0] /Rect [174.024 305.331 211.4 316.122] /A << /S /GoTo /D <74686D2E452E312E322E3137> >> >>}
\MusicPDFWriteObject{7989}{<< /Type /Annot /Subtype /Link /Border [0 0 0] /H /I /C [1 0 0] /Rect [238.551 292.181 275.682 302.972] /A << /S /GoTo /D <74686D2E452E312E322E3138> >> >>}
\MusicPDFWriteObject{7990}{<< /Type /Annot /Subtype /Link /Border [0 0 0] /H /I /C [1 0 0] /Rect [130.379 265.326 146.897 277.279] /A << /S /GoTo /D <73656374696F6E2E452E32> >> >>}
\MusicPDFWriteObject{7991}{<< /Type /Annot /Subtype /Link /Border [0 0 0] /H /I /C [0 1 0] /Rect [393.958 209.328 414.543 218.196] /A << /S /GoTo /D <636974652E427234> >> >>}
\MusicPDFWriteObject{7992}{<< /Type /Annot /Subtype /Link /Border [0 0 0] /H /I /C [1 0 0] /Rect [223.927 148.903 256.539 159.694] /A << /S /GoTo /D <74686D2E452E312E322E31> >> >>}
\MusicPDFWriteObject{7993}{<< /Type /Annot /Subtype /Link /Border [0 0 0] /H /I /C [1 0 0] /Rect [298.433 148.903 331.044 159.694] /A << /S /GoTo /D <74686D2E452E312E322E32> >> >>}
\MusicPDFWriteObject{7994}{<< /Type /Annot /Subtype /Link /Border [0 0 0] /H /I /C [1 0 0] /Rect [372.938 148.903 410.63 159.694] /A << /S /GoTo /D <74686D2E452E312E322E3130> >> >>}
\MusicPDFWriteObject{7995}{<< /Type /Annot /Subtype /Link /Border [0 0 0] /H /I /C [1 0 0] /Rect [488.327 148.903 526.019 159.694] /A << /S /GoTo /D <74686D2E452E312E322E3133> >> >>}
\MusicPDFWriteObject{7996}{<< /Type /Annot /Subtype /Link /Border [0 0 0] /H /I /C [0 1 0] /Rect [191.693 627.482 209.22 636.35] /A << /S /GoTo /D <636974652E537472> >> >>}
\MusicPDFWriteObject{7997}{<< /Type /Annot /Subtype /Link /Border [0 0 0] /H /I /C [1 0 0] /Rect [347.445 494.055 379.817 506.007] /A << /S /GoTo /D <74686D2E452E312E312E33> >> >>}
\MusicPDFWriteObject{7998}{<< /Type /Annot /Subtype /Link /Border [0 0 0] /H /I /C [1 0 0] /Rect [336.2 480.904 342.162 493.761] /A << /S /GoTo /D <48666F6F746E6F74652E313535> >> >>}
\MusicPDFWriteObject{7999}{<< /Type /Annot /Subtype /Link /Border [0 0 0] /H /I /C [1 0 0] /Rect [383.448 480.904 421.139 493.761] /A << /S /GoTo /D <6571756174696F6E2E452E312E322E3239> >> >>}
\MusicPDFWriteObject{8000}{<< /Type /Annot /Subtype /Link /Border [0 0 0] /H /I /C [1 0 0] /Rect [276.829 450.21 314.52 461] /A << /S /GoTo /D <74686D2E452E312E322E3139> >> >>}
\MusicPDFWriteObject{8001}{<< /Type /Annot /Subtype /Link /Border [0 0 0] /H /I /C [1 0 0] /Rect [510.547 436.505 543.159 448.458] /A << /S /GoTo /D <74686D2E452E322E342E39> >> >>}
\MusicPDFWriteObject{8002}{<< /Type /Annot /Subtype /Link /Border [0 0 0] /H /I /C [0 1 0] /Rect [424.832 412.694 449.255 421.562] /A << /S /GoTo /D <636974652E53745365> >> >>}
\MusicPDFWriteObject{8003}{<< /Type /Annot /Subtype /Link /Border [0 0 0] /H /I /C [1 0 0] /Rect [71.004 384.455 102.865 395.303] /A << /S /GoTo /D <74686D2E452E322E342E39> >> >>}
\MusicPDFWriteObject{8004}{<< /Type /Annot /Subtype /Link /Border [0 0 0] /H /I /C [1 0 0] /Rect [250.229 384.455 282.091 395.303] /A << /S /GoTo /D <74686D2E452E322E312E31> >> >>}
\MusicPDFWriteObject{8005}{<< /Type /Annot /Subtype /Link /Border [0 0 0] /H /I /C [1 0 0] /Rect [513.157 384.455 529.602 395.303] /A << /S /GoTo /D <73656374696F6E2E452E32> >> >>}
\MusicPDFWriteObject{8006}{<< /Type /Annot /Subtype /Link /Border [0 0 0] /H /I /C [0 1 0] /Rect [73.826 315.692 86.788 324.56] /A << /S /GoTo /D <636974652E4672> >> >>}
\MusicPDFWriteObject{8007}{<< /Type /Annot /Subtype /Link /Border [0 0 0] /H /I /C [0 1 0] /Rect [237.6 315.692 250.563 324.56] /A << /S /GoTo /D <636974652E4672> >> >>}
\MusicPDFWriteObject{8008}{<< /Type /Annot /Subtype /Link /Border [0 0 0] /H /I /C [0 1 0] /Rect [336.832 315.692 349.794 324.56] /A << /S /GoTo /D <636974652E4672> >> >>}
\MusicPDFWriteObject{8009}{<< /Type /Annot /Subtype /Link /Border [0 0 0] /H /I /C [0 1 0] /Rect [341.672 289.391 354.635 298.259] /A << /S /GoTo /D <636974652E4672> >> >>}
\MusicPDFWriteObject{8010}{<< /Type /Annot /Subtype /Link /Border [0 0 0] /H /I /C [1 0 0] /Rect [340.316 103.957 371.997 115.909] /A << /S /GoTo /D <74686D2E452E312E312E33> >> >>}
\MusicPDFWriteObject{8011}{<< /Type /Annot /Subtype /Link /Border [0 0 0] /H /I /C [0 1 0] /Rect [358.706 647.671 371.668 656.539] /A << /S /GoTo /D <636974652E4672> >> >>}
\MusicPDFWriteObject{8012}{<< /Type /Annot /Subtype /Link /Border [0 0 0] /H /I /C [1 0 0] /Rect [381.098 618.879 413.109 630.832] /A << /S /GoTo /D <74686D2E452E312E322E32> >> >>}
\MusicPDFWriteObject{8013}{<< /Type /Annot /Subtype /Link /Border [0 0 0] /H /I /C [1 0 0] /Rect [177.633 481.19 214.065 493.783] /A << /S /GoTo /D <74686D2E452E312E322E3231> >> >>}
\MusicPDFWriteObject{8014}{<< /Type /Annot /Subtype /Link /Border [0 0 0] /H /I /C [1 0 0] /Rect [209.928 434.481 246.36 445.271] /A << /S /GoTo /D <74686D2E452E312E322E3234> >> >>}
\MusicPDFWriteObject{8015}{<< /Type /Annot /Subtype /Link /Border [0 0 0] /H /I /C [0 1 0] /Rect [146.326 423.267 176 432.135] /A << /S /GoTo /D <636974652E41424D> >> >>}
\MusicPDFWriteObject{8016}{<< /Type /Annot /Subtype /Link /Border [0 0 0] /H /I /C [0 1 0] /Rect [275.036 423.267 304.709 432.135] /A << /S /GoTo /D <636974652E41424D> >> >>}
\MusicPDFWriteObject{8017}{<< /Type /Annot /Subtype /Link /Border [0 0 0] /H /I /C [1 0 0] /Rect [237.92 407.626 274.912 419.579] /A << /S /GoTo /D <74686D2E452E312E322E3231> >> >>}
\MusicPDFWriteObject{8018}{<< /Type /Annot /Subtype /Link /Border [0 0 0] /H /I /C [1 0 0] /Rect [251.895 387.217 289.587 398.008] /A << /S /GoTo /D <74686D2E452E312E322E3138> >> >>}
\MusicPDFWriteObject{8019}{<< /Type /Annot /Subtype /Link /Border [0 0 0] /H /I /C [0 1 0] /Rect [348.314 333.477 361.276 342.344] /A << /S /GoTo /D <636974652E4672> >> >>}
\MusicPDFWriteObject{8020}{<< /Type /Annot /Subtype /Link /Border [0 0 0] /H /I /C [0 1 0] /Rect [444.528 333.477 457.491 342.344] /A << /S /GoTo /D <636974652E4672> >> >>}
\MusicPDFWriteObject{8021}{<< /Type /Annot /Subtype /Link /Border [0 0 0] /H /I /C [1 0 0] /Rect [352.054 315.314 388.556 326.105] /A << /S /GoTo /D <74686D2E452E312E322E3231> >> >>}
\MusicPDFWriteObject{8022}{<< /Type /Annot /Subtype /Link /Border [0 0 0] /H /I /C [1 0 0] /Rect [474.722 315.314 511.224 326.105] /A << /S /GoTo /D <74686D2E452E312E322E3137> >> >>}
\MusicPDFWriteObject{8023}{<< /Type /Annot /Subtype /Link /Border [0 0 0] /H /I /C [1 0 0] /Rect [293.338 289.013 330.329 299.804] /A << /S /GoTo /D <74686D2E452E312E322E3138> >> >>}
\MusicPDFWriteObject{8024}{<< /Type /Annot /Subtype /Link /Border [0 0 0] /H /I /C [0 1 0] /Rect [221.739 143.011 239.267 151.879] /A << /S /GoTo /D <636974652E537472> >> >>}
\MusicPDFWriteObject{8025}{<< /Type /Annot /Subtype /Link /Border [0 0 0] /H /I /C [1 0 0] /Rect [178.54 119.138 209.981 129.929] /A << /S /GoTo /D <74686D2E452E312E312E33> >> >>}
\MusicPDFWriteObject{8026}{<< /Type /Annot /Subtype /Link /Border [0 0 0] /H /I /C [1 0 0] /Rect [323.445 632.03 361.137 643.983] /A << /S /GoTo /D <74686D2E452E312E322E3139> >> >>}
\MusicPDFWriteObject{8027}{<< /Type /Annot /Subtype /Link /Border [0 0 0] /H /I /C [1 0 0] /Rect [430.328 618.879 462.94 631.472] /A << /S /GoTo /D <74686D2E452E312E312E33> >> >>}
\MusicPDFWriteObject{8028}{<< /Type /Annot /Subtype /Link /Border [0 0 0] /H /I /C [1 0 0] /Rect [247.076 591.443 279.688 605.17] /A << /S /GoTo /D <74686D2E452E312E312E33> >> >>}
\MusicPDFWriteObject{8029}{<< /Type /Annot /Subtype /Link /Border [0 0 0] /H /I /C [1 0 0] /Rect [245.473 547.701 282.465 558.492] /A << /S /GoTo /D <74686D2E452E312E322E3137> >> >>}
\MusicPDFWriteObject{8030}{<< /Type /Annot /Subtype /Link /Border [0 0 0] /H /I /C [1 0 0] /Rect [384.086 547.701 421.078 558.492] /A << /S /GoTo /D <74686D2E452E312E322E3235> >> >>}
\MusicPDFWriteObject{8031}{<< /Type /Annot /Subtype /Link /Border [0 0 0] /H /I /C [1 0 0] /Rect [173.033 426.821 205.014 438.774] /A << /S /GoTo /D <74686D2E452E312E312E38> >> >>}
\MusicPDFWriteObject{8032}{<< /Type /Annot /Subtype /Link /Border [0 0 0] /H /I /C [1 0 0] /Rect [320.113 413.67 356.579 425.643] /A << /S /GoTo /D <74686D2E452E312E322E3235> >> >>}
\MusicPDFWriteObject{8033}{<< /Type /Annot /Subtype /Link /Border [0 0 0] /H /I /C [1 0 0] /Rect [216.785 216.627 254.091 227.418] /A << /S /GoTo /D <74686D2E452E312E322E3237> >> >>}
\MusicPDFWriteObject{8034}{<< /Type /Annot /Subtype /Link /Border [0 0 0] /H /I /C [1 0 0] /Rect [281.987 109.45 318.278 120.298] /A << /S /GoTo /D <74686D2E452E312E322E3132> >> >>}
\MusicPDFWriteObject{8035}{<< /Type /Annot /Subtype /Link /Border [0 0 0] /H /I /C [1 0 0] /Rect [106.556 83.149 143.548 93.94] /A << /S /GoTo /D <74686D2E452E312E322E3237> >> >>}
\MusicPDFWriteObject{8036}{<< /Type /Annot /Subtype /Link /Border [0 0 0] /H /I /C [0 1 0] /Rect [285.609 556.091 310.385 564.959] /A << /S /GoTo /D <636974652E4C474C> >> >>}
\MusicPDFWriteObject{8037}{<< /Type /Annot /Subtype /Link /Border [0 0 0] /H /I /C [0 1 0] /Rect [335.508 556.091 361.737 564.959] /A << /S /GoTo /D <636974652E4D6343> >> >>}
\MusicPDFWriteObject{8038}{<< /Type /Annot /Subtype /Link /Border [0 0 0] /H /I /C [0 1 0] /Rect [386.861 556.091 413.83 564.959] /A << /S /GoTo /D <636974652E52657A6B> >> >>}
\MusicPDFWriteObject{8039}{<< /Type /Annot /Subtype /Link /Border [0 0 0] /H /I /C [0 1 0] /Rect [455.553 556.091 473.08 564.959] /A << /S /GoTo /D <636974652E537472> >> >>}
\MusicPDFWriteObject{8040}{<< /Type /Annot /Subtype /Link /Border [0 0 0] /H /I /C [1 0 0] /Rect [301.583 422.97 333.594 433.761] /A << /S /GoTo /D <74686D2E452E322E342E39> >> >>}
\MusicPDFWriteObject{8041}{<< /Type /Annot /Subtype /Link /Border [0 0 0] /H /I /C [0 1 0] /Rect [191.802 687.123 209.33 695.991] /A << /S /GoTo /D <636974652E537472> >> >>}
\MusicPDFWriteObject{8042}{<< /Type /Annot /Subtype /Link /Border [0 0 0] /H /I /C [0 1 0] /Rect [297.24 687.123 324.209 695.991] /A << /S /GoTo /D <636974652E52657A6B> >> >>}
\MusicPDFWriteObject{8043}{<< /Type /Annot /Subtype /Link /Border [0 0 0] /H /I /C [1 0 0] /Rect [71.004 672.036 103.616 682.827] /A << /S /GoTo /D <74686D2E452E322E342E39> >> >>}
\MusicPDFWriteObject{8044}{<< /Type /Annot /Subtype /Link /Border [0 0 0] /H /I /C [1 0 0] /Rect [71.004 486.439 102.415 499.032] /A << /S /GoTo /D <74686D2E452E322E312E34> >> >>}
\MusicPDFWriteObject{8045}{<< /Type /Annot /Subtype /Link /Border [0 0 0] /H /I /C [1 0 0] /Rect [394.08 486.439 425.491 499.032] /A << /S /GoTo /D <6571756174696F6E2E452E322E312E32> >> >>}
\MusicPDFWriteObject{8046}{<< /Type /Annot /Subtype /Link /Border [0 0 0] /H /I /C [1 0 0] /Rect [98.885 473.289 131.467 485.241] /A << /S /GoTo /D <6571756174696F6E2E452E322E312E31> >> >>}
\MusicPDFWriteObject{8047}{<< /Type /Annot /Subtype /Link /Border [0 0 0] /H /I /C [0 1 0] /Rect [450.173 144.706 466.358 153.574] /A << /S /GoTo /D <636974652E5350> >> >>}
\MusicPDFWriteObject{8048}{<< /Type /Annot /Subtype /Link /Border [0 0 0] /H /I /C [1 0 0] /Rect [474.021 129.065 506.033 141.02] /A << /S /GoTo /D <74686D2E452E312E312E38> >> >>}
\MusicPDFWriteObject{8049}{<< /Type /Annot /Subtype /Link /Border [0 0 0] /H /I /C [1 0 0] /Rect [314.933 108.897 346.825 120.85] /A << /S /GoTo /D <74686D2E452E322E312E36> >> >>}
\MusicPDFWriteObject{8050}{<< /Type /Annot /Subtype /Link /Border [0 0 0] /H /I /C [1 0 0] /Rect [309.04 82.596 346.102 94.549] /A << /S /GoTo /D <74686D2E452E322E312E3130> >> >>}
\MusicPDFWriteObject{8051}{<< /Type /Annot /Subtype /Link /Border [0 0 0] /H /I /C [1 0 0] /Rect [338.916 684.633 371.528 696.586] /A << /S /GoTo /D <74686D2E452E322E312E36> >> >>}
\MusicPDFWriteObject{8052}{<< /Type /Annot /Subtype /Link /Border [0 0 0] /H /I /C [0 1 0] /Rect [483.303 687.123 510.272 695.991] /A << /S /GoTo /D <636974652E52657A6B> >> >>}
\MusicPDFWriteObject{8053}{<< /Type /Annot /Subtype /Link /Border [0 0 0] /H /I /C [1 0 0] /Rect [320.29 618.879 352.902 630.832] /A << /S /GoTo /D <74686D2E452E322E312E38> >> >>}
\MusicPDFWriteObject{8054}{<< /Type /Annot /Subtype /Link /Border [0 0 0] /H /I /C [1 0 0] /Rect [512.234 618.879 544.845 630.832] /A << /S /GoTo /D <74686D2E452E322E312E39> >> >>}
\MusicPDFWriteObject{8055}{<< /Type /Annot /Subtype /Link /Border [0 0 0] /H /I /C [1 0 0] /Rect [380.269 605.729 412.49 617.681] /A << /S /GoTo /D <74686D2E452E322E342E35> >> >>}
\MusicPDFWriteObject{8056}{<< /Type /Annot /Subtype /Link /Border [0 0 0] /H /I /C [1 0 0] /Rect [234.712 560.189 266.723 572.142] /A << /S /GoTo /D <74686D2E452E322E312E36> >> >>}
\MusicPDFWriteObject{8057}{<< /Type /Annot /Subtype /Link /Border [0 0 0] /H /I /C [0 1 0] /Rect [136.283 517.139 160.706 526.007] /A << /S /GoTo /D <636974652E53745365> >> >>}
\MusicPDFWriteObject{8058}{<< /Type /Annot /Subtype /Link /Border [0 0 0] /H /I /C [1 0 0] /Rect [320.997 514.649 353.309 526.602] /A << /S /GoTo /D <74686D2E452E322E342E34> >> >>}
\MusicPDFWriteObject{8059}{<< /Type /Annot /Subtype /Link /Border [0 0 0] /H /I /C [0 1 0] /Rect [310.014 503.989 334.438 512.856] /A << /S /GoTo /D <636974652E53745365> >> >>}
\MusicPDFWriteObject{8060}{<< /Type /Annot /Subtype /Link /Border [0 0 0] /H /I /C [0 1 0] /Rect [495.439 484.75 519.863 493.618] /A << /S /GoTo /D <636974652E53745365> >> >>}
\MusicPDFWriteObject{8061}{<< /Type /Annot /Subtype /Link /Border [0 0 0] /H /I /C [0 1 0] /Rect [252.475 445.298 276.899 454.166] /A << /S /GoTo /D <636974652E53745365> >> >>}
\MusicPDFWriteObject{8062}{<< /Type /Annot /Subtype /Link /Border [0 0 0] /H /I /C [0 1 0] /Rect [368.951 432.147 386.921 441.015] /A << /S /GoTo /D <636974652E5042> >> >>}
\MusicPDFWriteObject{8063}{<< /Type /Annot /Subtype /Link /Border [0 0 0] /H /I /C [0 1 0] /Rect [191.693 332.889 209.22 341.756] /A << /S /GoTo /D <636974652E537472> >> >>}
\MusicPDFWriteObject{8064}{<< /Type /Annot /Subtype /Link /Border [0 0 0] /H /I /C [0 1 0] /Rect [337.043 332.889 354.571 341.756] /A << /S /GoTo /D <636974652E537472> >> >>}
\MusicPDFWriteObject{8065}{<< /Type /Annot /Subtype /Link /Border [0 0 0] /H /I /C [1 0 0] /Rect [424.264 298.119 461.781 310.074] /A << /S /GoTo /D <74686D2E452E312E322E3137> >> >>}
\MusicPDFWriteObject{8066}{<< /Type /Annot /Subtype /Link /Border [0 0 0] /H /I /C [1 0 0] /Rect [155.555 245.516 187.566 257.469] /A << /S /GoTo /D <74686D2E452E322E312E36> >> >>}
\MusicPDFWriteObject{8067}{<< /Type /Annot /Subtype /Link /Border [0 0 0] /H /I /C [1 0 0] /Rect [126.632 173.675 163.169 185.63] /A << /S /GoTo /D <74686D2E452E312E322E3139> >> >>}
\MusicPDFWriteObject{8068}{<< /Type /Annot /Subtype /Link /Border [0 0 0] /H /I /C [1 0 0] /Rect [233.466 590.182 266.077 602.135] /A << /S /GoTo /D <74686D2E452E312E312E33> >> >>}
\MusicPDFWriteObject{8069}{<< /Type /Annot /Subtype /Link /Border [0 0 0] /H /I /C [0 1 0] /Rect [296.963 440.78 339.986 449.647] /A << /S /GoTo /D <636974652E426F752D4754> >> >>}
\MusicPDFWriteObject{8070}{<< /Type /Annot /Subtype /Link /Border [0 0 0] /H /I /C [1 0 0] /Rect [286.041 335.923 322.332 350.125] /A << /S /GoTo /D <74686D2E452E312E322E3139> >> >>}
\MusicPDFWriteObject{8071}{<< /Type /Annot /Subtype /Link /Border [0 0 0] /H /I /C [1 0 0] /Rect [378.321 324.382 416.012 336.335] /A << /S /GoTo /D <74686D2E452E312E322E3235> >> >>}
\MusicPDFWriteObject{8072}{<< /Type /Annot /Subtype /Link /Border [0 0 0] /H /I /C [1 0 0] /Rect [486.587 324.382 524.279 336.335] /A << /S /GoTo /D <74686D2E452E322E312E3131> >> >>}
\MusicPDFWriteObject{8073}{<< /Type /Annot /Subtype /Link /Border [0 0 0] /H /I /C [1 0 0] /Rect [71.004 298.08 107.995 310.033] /A << /S /GoTo /D <74686D2E452E312E322E3134> >> >>}
\MusicPDFWriteObject{8074}{<< /Type /Annot /Subtype /Link /Border [0 0 0] /H /I /C [1 0 0] /Rect [509.877 226.606 541.888 238.559] /A << /S /GoTo /D <6571756174696F6E2E452E322E312E34> >> >>}
\MusicPDFWriteObject{8075}{<< /Type /Annot /Subtype /Link /Border [0 0 0] /H /I /C [1 0 0] /Rect [215.651 213.455 247.212 225.408] /A << /S /GoTo /D <74686D2E452E312E312E33> >> >>}
\MusicPDFWriteObject{8076}{<< /Type /Annot /Subtype /Link /Border [0 0 0] /H /I /C [1 0 0] /Rect [512.406 142.289 543.817 153.08] /A << /S /GoTo /D <74686D2E452E322E312E34> >> >>}
\MusicPDFWriteObject{8077}{<< /Type /Annot /Subtype /Link /Border [0 0 0] /H /I /C [1 0 0] /Rect [488.969 82.596 520.38 94.568] /A << /S /GoTo /D <6571756174696F6E2E452E322E312E34> >> >>}
\MusicPDFWriteObject{8078}{<< /Type /Annot /Subtype /Link /Border [0 0 0] /H /I /C [1 0 0] /Rect [192.444 685.187 230.136 695.977] /A << /S /GoTo /D <74686D2E452E322E312E3133> >> >>}
\MusicPDFWriteObject{8079}{<< /Type /Annot /Subtype /Link /Border [0 0 0] /H /I /C [1 0 0] /Rect [333.216 645.734 365.828 656.525] /A << /S /GoTo /D <74686D2E452E312E312E33> >> >>}
\MusicPDFWriteObject{8080}{<< /Type /Annot /Subtype /Link /Border [0 0 0] /H /I /C [0 1 0] /Rect [211.668 496.36 238.174 505.228] /A << /S /GoTo /D <636974652E564C48> >> >>}
\MusicPDFWriteObject{8081}{<< /Type /Annot /Subtype /Link /Border [0 0 0] /H /I /C [1 0 0] /Rect [128.189 473.179 134.152 486.037] /A << /S /GoTo /D <48666F6F746E6F74652E313536> >> >>}
\MusicPDFWriteObject{8082}{<< /Type /Annot /Subtype /Link /Border [0 0 0] /H /I /C [1 0 0] /Rect [230.875 446.877 236.838 459.735] /A << /S /GoTo /D <48666F6F746E6F74652E313537> >> >>}
\MusicPDFWriteObject{8083}{<< /Type /Annot /Subtype /Link /Border [0 0 0] /H /I /C [1 0 0] /Rect [222.598 394.274 228.561 407.132] /A << /S /GoTo /D <48666F6F746E6F74652E313538> >> >>}
\MusicPDFWriteObject{8084}{<< /Type /Annot /Subtype /Link /Border [0 0 0] /H /I /C [1 0 0] /Rect [110.92 136.659 116.883 148.963] /A << /S /GoTo /D <48666F6F746E6F74652E313539> >> >>}
\MusicPDFWriteObject{8085}{<< /Type /Annot /Subtype /Link /Border [0 0 0] /H /I /C [1 0 0] /Rect [195.622 684.335 231.913 696.588] /A << /S /GoTo /D <74686D2E452E322E312E3133> >> >>}
\MusicPDFWriteObject{8086}{<< /Type /Annot /Subtype /Link /Border [0 0 0] /H /I /C [1 0 0] /Rect [259.529 671.185 297.221 682.884] /A << /S /GoTo /D <74686D2E452E322E312E3135> >> >>}
\MusicPDFWriteObject{8087}{<< /Type /Annot /Subtype /Link /Border [0 0 0] /H /I /C [1 0 0] /Rect [449.224 671.185 486.915 682.884] /A << /S /GoTo /D <74686D2E452E322E312E3136> >> >>}
\MusicPDFWriteObject{8088}{<< /Type /Annot /Subtype /Link /Border [0 0 0] /H /I /C [0 1 0] /Rect [470.518 227.757 488.046 236.625] /A << /S /GoTo /D <636974652E537472> >> >>}
\MusicPDFWriteObject{8089}{<< /Type /Annot /Subtype /Link /Border [0 0 0] /H /I /C [1 0 0] /Rect [73.712 671.482 105.723 683.437] /A << /S /GoTo /D <6571756174696F6E2E452E322E322E33> >> >>}
\MusicPDFWriteObject{8090}{<< /Type /Annot /Subtype /Link /Border [0 0 0] /H /I /C [1 0 0] /Rect [492.431 671.482 529.423 683.437] /A << /S /GoTo /D <74686D2E452E322E312E3131> >> >>}
\MusicPDFWriteObject{8091}{<< /Type /Annot /Subtype /Link /Border [0 0 0] /H /I /C [1 0 0] /Rect [276.046 650.997 307.638 662.952] /A << /S /GoTo /D <6571756174696F6E2E452E322E322E33> >> >>}
\MusicPDFWriteObject{8092}{<< /Type /Annot /Subtype /Link /Border [0 0 0] /H /I /C [1 0 0] /Rect [256.586 614.57 288.597 626.522] /A << /S /GoTo /D <74686D2E452E312E322E31> >> >>}
\MusicPDFWriteObject{8093}{<< /Type /Annot /Subtype /Link /Border [0 0 0] /H /I /C [0 1 0] /Rect [311.56 617.06 329.088 625.928] /A << /S /GoTo /D <636974652E537472> >> >>}
\MusicPDFWriteObject{8094}{<< /Type /Annot /Subtype /Link /Border [0 0 0] /H /I /C [0 1 0] /Rect [404.73 617.06 422.258 625.928] /A << /S /GoTo /D <636974652E537472> >> >>}
\MusicPDFWriteObject{8095}{<< /Type /Annot /Subtype /Link /Border [0 0 0] /H /I /C [1 0 0] /Rect [328.205 325.424 360.817 338.017] /A << /S /GoTo /D <74686D2E452E322E312E36> >> >>}
\MusicPDFWriteObject{8096}{<< /Type /Annot /Subtype /Link /Border [0 0 0] /H /I /C [1 0 0] /Rect [309.313 312.273 340.725 324.866] /A << /S /GoTo /D <74686D2E452E312E322E31> >> >>}
\MusicPDFWriteObject{8097}{<< /Type /Annot /Subtype /Link /Border [0 0 0] /H /I /C [1 0 0] /Rect [287.732 299.122 319.743 311.078] /A << /S /GoTo /D <6571756174696F6E2E452E322E322E38> >> >>}
\MusicPDFWriteObject{8098}{<< /Type /Annot /Subtype /Link /Border [0 0 0] /H /I /C [1 0 0] /Rect [73.712 137.237 110.703 149.19] /A << /S /GoTo /D <6571756174696F6E2E452E322E322E3130> >> >>}
\MusicPDFWriteObject{8099}{<< /Type /Annot /Subtype /Link /Border [0 0 0] /H /I /C [1 0 0] /Rect [121.635 100.81 158.626 112.782] /A << /S /GoTo /D <6571756174696F6E2E452E322E322E3131> >> >>}
\MusicPDFWriteObject{8100}{<< /Type /Annot /Subtype /Link /Border [0 0 0] /H /I /C [1 0 0] /Rect [284.47 100.81 316.481 112.782] /A << /S /GoTo /D <6571756174696F6E2E452E322E322E39> >> >>}
\MusicPDFWriteObject{8101}{<< /Type /Annot /Subtype /Link /Border [0 0 0] /H /I /C [1 0 0] /Rect [216.947 672.035 253.238 682.884] /A << /S /GoTo /D <74686D2E452E312E322E3134> >> >>}
\MusicPDFWriteObject{8102}{<< /Type /Annot /Subtype /Link /Border [0 0 0] /H /I /C [1 0 0] /Rect [150.735 658.332 188.427 670.304] /A << /S /GoTo /D <6571756174696F6E2E452E322E322E3130> >> >>}
\MusicPDFWriteObject{8103}{<< /Type /Annot /Subtype /Link /Border [0 0 0] /H /I /C [1 0 0] /Rect [507.144 645.18 539.756 657.136] /A << /S /GoTo /D <6571756174696F6E2E452E322E322E35> >> >>}
\MusicPDFWriteObject{8104}{<< /Type /Annot /Subtype /Link /Border [0 0 0] /H /I /C [1 0 0] /Rect [457.106 479.788 488.727 490.579] /A << /S /GoTo /D <74686D2E452E322E312E36> >> >>}
\MusicPDFWriteObject{8105}{<< /Type /Annot /Subtype /Link /Border [0 0 0] /H /I /C [1 0 0] /Rect [267.542 402.202 300.154 414.157] /A << /S /GoTo /D <74686D2E452E322E312E36> >> >>}
\MusicPDFWriteObject{8106}{<< /Type /Annot /Subtype /Link /Border [0 0 0] /H /I /C [1 0 0] /Rect [139.757 376.454 171.768 387.245] /A << /S /GoTo /D <74686D2E452E322E312E36> >> >>}
\MusicPDFWriteObject{8107}{<< /Type /Annot /Subtype /Link /Border [0 0 0] /H /I /C [1 0 0] /Rect [349.779 357.326 381.79 368.117] /A << /S /GoTo /D <74686D2E452E312E322E33> >> >>}
\MusicPDFWriteObject{8108}{<< /Type /Annot /Subtype /Link /Border [0 0 0] /H /I /C [1 0 0] /Rect [356.803 297.543 389.145 309.496] /A << /S /GoTo /D <74686D2E452E322E322E35> >> >>}
\MusicPDFWriteObject{8109}{<< /Type /Annot /Subtype /Link /Border [0 0 0] /H /I /C [1 0 0] /Rect [475.54 297.543 507.882 309.496] /A << /S /GoTo /D <74686D2E452E312E322E32> >> >>}
\MusicPDFWriteObject{8110}{<< /Type /Annot /Subtype /Link /Border [0 0 0] /H /I /C [1 0 0] /Rect [272.505 206.041 304.517 216.832] /A << /S /GoTo /D <74686D2E452E322E322E37> >> >>}
\MusicPDFWriteObject{8111}{<< /Type /Annot /Subtype /Link /Border [0 0 0] /H /I /C [0 1 0] /Rect [428.066 98.237 445.207 107.105] /A << /S /GoTo /D <636974652E4475> >> >>}
\MusicPDFWriteObject{8112}{<< /Type /Annot /Subtype /Link /Border [0 0 0] /H /I /C [1 0 0] /Rect [227.478 685.187 259.099 695.977] /A << /S /GoTo /D <74686D2E452E312E322E34> >> >>}
\MusicPDFWriteObject{8113}{<< /Type /Annot /Subtype /Link /Border [0 0 0] /H /I /C [1 0 0] /Rect [398.719 658.332 430.521 670.284] /A << /S /GoTo /D <74686D2E452E312E322E32> >> >>}
\MusicPDFWriteObject{8114}{<< /Type /Annot /Subtype /Link /Border [0 0 0] /H /I /C [1 0 0] /Rect [379.688 632.03 411.34 643.983] /A << /S /GoTo /D <74686D2E452E312E322E32> >> >>}
\MusicPDFWriteObject{8115}{<< /Type /Annot /Subtype /Link /Border [0 0 0] /H /I /C [1 0 0] /Rect [284.021 606.282 316.483 617.073] /A << /S /GoTo /D <74686D2E452E322E322E37> >> >>}
\MusicPDFWriteObject{8116}{<< /Type /Annot /Subtype /Link /Border [0 0 0] /H /I /C [1 0 0] /Rect [290.438 482.127 322.93 494.58] /A << /S /GoTo /D <74686D2E452E322E322E35> >> >>}
\MusicPDFWriteObject{8117}{<< /Type /Annot /Subtype /Link /Border [0 0 0] /H /I /C [1 0 0] /Rect [372.051 468.976 404.663 481.43] /A << /S /GoTo /D <74686D2E452E312E322E32> >> >>}
\MusicPDFWriteObject{8118}{<< /Type /Annot /Subtype /Link /Border [0 0 0] /H /I /C [1 0 0] /Rect [203.696 456.324 235.707 468.277] /A << /S /GoTo /D <74686D2E452E312E322E32> >> >>}
\MusicPDFWriteObject{8119}{<< /Type /Annot /Subtype /Link /Border [0 0 0] /H /I /C [1 0 0] /Rect [334.849 410.881 366.86 422.836] /A << /S /GoTo /D <74686D2E452E322E312E36> >> >>}
\MusicPDFWriteObject{8120}{<< /Type /Annot /Subtype /Link /Border [0 0 0] /H /I /C [0 1 0] /Rect [525.684 355.905 543.211 364.773] /A << /S /GoTo /D <636974652E537472> >> >>}
\MusicPDFWriteObject{8121}{<< /Type /Annot /Subtype /Link /Border [0 0 0] /H /I /C [1 0 0] /Rect [71.004 154.314 102.625 166.907] /A << /S /GoTo /D <74686D2E452E322E312E36> >> >>}
\MusicPDFWriteObject{8122}{<< /Type /Annot /Subtype /Link /Border [0 0 0] /H /I /C [1 0 0] /Rect [308.279 154.314 339.9 166.907] /A << /S /GoTo /D <74686D2E452E312E322E31> >> >>}
\MusicPDFWriteObject{8123}{<< /Type /Annot /Subtype /Link /Border [0 0 0] /H /I /C [1 0 0] /Rect [97.126 140.138 133.418 153.756] /A << /S /GoTo /D <6571756174696F6E2E452E322E322E3134> >> >>}
\MusicPDFWriteObject{8124}{<< /Type /Annot /Subtype /Link /Border [0 0 0] /H /I /C [1 0 0] /Rect [132.619 82.595 164.63 94.55] /A << /S /GoTo /D <74686D2E452E322E312E36> >> >>}
\MusicPDFWriteObject{8125}{<< /Type /Annot /Subtype /Link /Border [0 0 0] /H /I /C [1 0 0] /Rect [180.037 574.873 217.554 586.826] /A << /S /GoTo /D <74686D2E452E312E322E3139> >> >>}
\MusicPDFWriteObject{8126}{<< /Type /Annot /Subtype /Link /Border [0 0 0] /H /I /C [1 0 0] /Rect [123.826 561.722 161.517 573.675] /A << /S /GoTo /D <74686D2E452E312E322E3139> >> >>}
\MusicPDFWriteObject{8127}{<< /Type /Annot /Subtype /Link /Border [0 0 0] /H /I /C [1 0 0] /Rect [177.313 548.571 209.925 560.524] /A << /S /GoTo /D <74686D2E452E322E322E35> >> >>}
\MusicPDFWriteObject{8128}{<< /Type /Annot /Subtype /Link /Border [0 0 0] /H /I /C [1 0 0] /Rect [267.806 515.367 299.818 526.158] /A << /S /GoTo /D <74686D2E452E322E332E31> >> >>}
\MusicPDFWriteObject{8129}{<< /Type /Annot /Subtype /Link /Border [0 0 0] /H /I /C [1 0 0] /Rect [124.264 433.67 156.276 444.46] /A << /S /GoTo /D <74686D2E452E322E342E39> >> >>}
\MusicPDFWriteObject{8130}{<< /Type /Annot /Subtype /Link /Border [0 0 0] /H /I /C [1 0 0] /Rect [316.684 386.761 353.326 397.552] /A << /S /GoTo /D <74686D2E452E312E322E3237> >> >>}
\MusicPDFWriteObject{8131}{<< /Type /Annot /Subtype /Link /Border [0 0 0] /H /I /C [0 1 0] /Rect [250.126 278.938 267.654 287.805] /A << /S /GoTo /D <636974652E537472> >> >>}
\MusicPDFWriteObject{8132}{<< /Type /Annot /Subtype /Link /Border [0 0 0] /H /I /C [1 0 0] /Rect [179.476 255.87 216.538 266.718] /A << /S /GoTo /D <74686D2E452E312E322E3237> >> >>}
\MusicPDFWriteObject{8133}{<< /Type /Annot /Subtype /Link /Border [0 0 0] /H /I /C [1 0 0] /Rect [105.686 202.714 137.097 214.669] /A << /S /GoTo /D <74686D2E452E322E322E34> >> >>}
\MusicPDFWriteObject{8134}{<< /Type /Annot /Subtype /Link /Border [0 0 0] /H /I /C [1 0 0] /Rect [71.004 189.563 102.805 201.518] /A << /S /GoTo /D <74686D2E452E322E322E34> >> >>}
\MusicPDFWriteObject{8135}{<< /Type /Annot /Subtype /Link /Border [0 0 0] /H /I /C [1 0 0] /Rect [388.396 684.633 421.008 696.588] /A << /S /GoTo /D <74686D2E452E322E312E36> >> >>}
\MusicPDFWriteObject{8136}{<< /Type /Annot /Subtype /Link /Border [0 0 0] /H /I /C [1 0 0] /Rect [175.117 515.755 207.128 526.546] /A << /S /GoTo /D <74686D2E452E322E332E34> >> >>}
\MusicPDFWriteObject{8137}{<< /Type /Annot /Subtype /Link /Border [0 0 0] /H /I /C [0 1 0] /Rect [302.717 483.819 320.245 492.687] /A << /S /GoTo /D <636974652E537472> >> >>}
\MusicPDFWriteObject{8138}{<< /Type /Annot /Subtype /Link /Border [0 0 0] /H /I /C [1 0 0] /Rect [205.323 345.77 242.42 357.723] /A << /S /GoTo /D <74686D2E452E312E322E3136> >> >>}
\MusicPDFWriteObject{8139}{<< /Type /Annot /Subtype /Link /Border [0 0 0] /H /I /C [1 0 0] /Rect [466.475 298.747 498.306 310.7] /A << /S /GoTo /D <74686D2E452E322E322E35> >> >>}
\MusicPDFWriteObject{8140}{<< /Type /Annot /Subtype /Link /Border [0 0 0] /H /I /C [1 0 0] /Rect [389.014 285.596 421.626 297.549] /A << /S /GoTo /D <74686D2E452E312E322E32> >> >>}
\MusicPDFWriteObject{8141}{<< /Type /Annot /Subtype /Link /Border [0 0 0] /H /I /C [1 0 0] /Rect [106.556 645.734 138.568 656.582] /A << /S /GoTo /D <74686D2E452E322E332E33> >> >>}
\MusicPDFWriteObject{8142}{<< /Type /Annot /Subtype /Link /Border [0 0 0] /H /I /C [1 0 0] /Rect [136.223 595.603 168.834 607.555] /A << /S /GoTo /D <6571756174696F6E2E452E322E332E32> >> >>}
\MusicPDFWriteObject{8143}{<< /Type /Annot /Subtype /Link /Border [0 0 0] /H /I /C [1 0 0] /Rect [196.124 556.15 228.046 568.103] /A << /S /GoTo /D <6571756174696F6E2E452E322E332E31> >> >>}
\MusicPDFWriteObject{8144}{<< /Type /Annot /Subtype /Link /Border [0 0 0] /H /I /C [1 0 0] /Rect [506.336 556.15 543.223 568.103] /A << /S /GoTo /D <74686D2E452E312E322E3137> >> >>}
\MusicPDFWriteObject{8145}{<< /Type /Annot /Subtype /Link /Border [0 0 0] /H /I /C [1 0 0] /Rect [222.634 543 254.615 554.952] /A << /S /GoTo /D <74686D2E452E322E322E35> >> >>}
\MusicPDFWriteObject{8146}{<< /Type /Annot /Subtype /Link /Border [0 0 0] /H /I /C [1 0 0] /Rect [269.841 431.472 302.453 443.428] /A << /S /GoTo /D <74686D2E452E322E332E34> >> >>}
\MusicPDFWriteObject{8147}{<< /Type /Annot /Subtype /Link /Border [0 0 0] /H /I /C [1 0 0] /Rect [218.956 418.322 251.568 430.295] /A << /S /GoTo /D <6571756174696F6E2E452E322E332E33> >> >>}
\MusicPDFWriteObject{8148}{<< /Type /Annot /Subtype /Link /Border [0 0 0] /H /I /C [1 0 0] /Rect [487.236 418.322 519.848 430.295] /A << /S /GoTo /D <6571756174696F6E2E452E322E332E33> >> >>}
\MusicPDFWriteObject{8149}{<< /Type /Annot /Subtype /Link /Border [0 0 0] /H /I /C [1 0 0] /Rect [73.789 405.171 106.401 417.764] /A << /S /GoTo /D <6571756174696F6E2E452E322E332E31> >> >>}
\MusicPDFWriteObject{8150}{<< /Type /Annot /Subtype /Link /Border [0 0 0] /H /I /C [1 0 0] /Rect [511.858 351.965 543.27 362.755] /A << /S /GoTo /D <74686D2E452E322E332E36> >> >>}
\MusicPDFWriteObject{8151}{<< /Type /Annot /Subtype /Link /Border [0 0 0] /H /I /C [1 0 0] /Rect [287.267 108.897 318.678 120.85] /A << /S /GoTo /D <74686D2E452E322E322E35> >> >>}
\MusicPDFWriteObject{8152}{<< /Type /Annot /Subtype /Link /Border [0 0 0] /H /I /C [1 0 0] /Rect [292.18 95.746 323.591 107.719] /A << /S /GoTo /D <74686D2E452E312E322E32> >> >>}
\MusicPDFWriteObject{8153}{<< /Type /Annot /Subtype /Link /Border [0 0 0] /H /I /C [1 0 0] /Rect [365.731 536.216 398.343 548.189] /A << /S /GoTo /D <74686D2E452E322E332E36> >> >>}
\MusicPDFWriteObject{8154}{<< /Type /Annot /Subtype /Link /Border [0 0 0] /H /I /C [1 0 0] /Rect [284.746 497.317 316.758 508.165] /A << /S /GoTo /D <74686D2E452E322E332E33> >> >>}
\MusicPDFWriteObject{8155}{<< /Type /Annot /Subtype /Link /Border [0 0 0] /H /I /C [1 0 0] /Rect [171.436 431.01 203.447 442.963] /A << /S /GoTo /D <74686D2E452E322E332E36> >> >>}
\MusicPDFWriteObject{8156}{<< /Type /Annot /Subtype /Link /Border [0 0 0] /H /I /C [1 0 0] /Rect [106.556 394.453 143.548 403.307] /A << /S /GoTo /D <74686D2E452E312E322E3134> >> >>}
\MusicPDFWriteObject{8157}{<< /Type /Annot /Subtype /Link /Border [0 0 0] /H /I /C [1 0 0] /Rect [375.523 263.188 408.134 275.143] /A << /S /GoTo /D <74686D2E452E322E322E35> >> >>}
\MusicPDFWriteObject{8158}{<< /Type /Annot /Subtype /Link /Border [0 0 0] /H /I /C [1 0 0] /Rect [98.422 158.476 130.433 170.428] /A << /S /GoTo /D <6571756174696F6E2E452E322E322E35> >> >>}
\MusicPDFWriteObject{8159}{<< /Type /Annot /Subtype /Link /Border [0 0 0] /H /I /C [1 0 0] /Rect [398.049 122.048 434.34 134.02] /A << /S /GoTo /D <74686D2E452E312E322E3134> >> >>}
\MusicPDFWriteObject{8160}{<< /Type /Annot /Subtype /Link /Border [0 0 0] /H /I /C [1 0 0] /Rect [440.918 671.482 472.839 683.437] /A << /S /GoTo /D <74686D2E452E322E312E36> >> >>}
\MusicPDFWriteObject{8161}{<< /Type /Annot /Subtype /Link /Border [0 0 0] /H /I /C [1 0 0] /Rect [349.124 645.18 386.815 657.136] /A << /S /GoTo /D <74686D2E452E312E322E3134> >> >>}
\MusicPDFWriteObject{8162}{<< /Type /Annot /Subtype /Link /Border [0 0 0] /H /I /C [1 0 0] /Rect [71.004 618.519 102.415 630.223] /A << /S /GoTo /D <74686D2E452E322E322E35> >> >>}
\MusicPDFWriteObject{8163}{<< /Type /Annot /Subtype /Link /Border [0 0 0] /H /I /C [1 0 0] /Rect [341.151 618.519 377.442 630.223] /A << /S /GoTo /D <74686D2E452E322E312E3131> >> >>}
\MusicPDFWriteObject{8164}{<< /Type /Annot /Subtype /Link /Border [0 0 0] /H /I /C [1 0 0] /Rect [204 389.964 236.101 402.989] /A << /S /GoTo /D <74686D2E452E322E322E35> >> >>}
\MusicPDFWriteObject{8165}{<< /Type /Annot /Subtype /Link /Border [0 0 0] /H /I /C [1 0 0] /Rect [190.346 160.069 222.958 170.86] /A << /S /GoTo /D <74686D2E452E322E322E39> >> >>}
\MusicPDFWriteObject{8166}{<< /Type /Annot /Subtype /Link /Border [0 0 0] /H /I /C [1 0 0] /Rect [192.847 485.241 225.459 497.214] /A << /S /GoTo /D <74686D2E452E322E312E36> >> >>}
\MusicPDFWriteObject{8167}{<< /Type /Annot /Subtype /Link /Border [0 0 0] /H /I /C [1 0 0] /Rect [336.859 458.939 374.55 470.912] /A << /S /GoTo /D <6571756174696F6E2E452E322E332E3130> >> >>}
\MusicPDFWriteObject{8168}{<< /Type /Annot /Subtype /Link /Border [0 0 0] /H /I /C [1 0 0] /Rect [116.542 445.789 154.234 457.761] /A << /S /GoTo /D <6571756174696F6E2E452E322E332E3130> >> >>}
\MusicPDFWriteObject{8169}{<< /Type /Annot /Subtype /Link /Border [0 0 0] /H /I /C [1 0 0] /Rect [408.221 445.789 445.912 457.761] /A << /S /GoTo /D <6571756174696F6E2E452E322E332E3131> >> >>}
\MusicPDFWriteObject{8170}{<< /Type /Annot /Subtype /Link /Border [0 0 0] /H /I /C [1 0 0] /Rect [176.108 288.533 208.33 299.324] /A << /S /GoTo /D <74686D2E452E322E322E39> >> >>}
\MusicPDFWriteObject{8171}{<< /Type /Annot /Subtype /Link /Border [0 0 0] /H /I /C [1 0 0] /Rect [257.628 275.382 289.039 286.173] /A << /S /GoTo /D <74686D2E452E312E322E32> >> >>}
\MusicPDFWriteObject{8172}{<< /Type /Annot /Subtype /Link /Border [0 0 0] /H /I /C [1 0 0] /Rect [137.236 261.678 169.248 273.633] /A << /S /GoTo /D <74686D2E452E312E322E31> >> >>}
\MusicPDFWriteObject{8173}{<< /Type /Annot /Subtype /Link /Border [0 0 0] /H /I /C [1 0 0] /Rect [262.991 210.488 299.283 222.461] /A << /S /GoTo /D <74686D2E452E322E312E3131> >> >>}
\MusicPDFWriteObject{8174}{<< /Type /Annot /Subtype /Link /Border [0 0 0] /H /I /C [1 0 0] /Rect [119.664 105.282 151.165 117.255] /A << /S /GoTo /D <74686D2E452E322E332E39> >> >>}
\MusicPDFWriteObject{8175}{<< /Type /Annot /Subtype /Link /Border [0 0 0] /H /I /C [0 1 0] /Rect [209.268 653.813 234.044 662.681] /A << /S /GoTo /D <636974652E4C474C> >> >>}
\MusicPDFWriteObject{8176}{<< /Type /Annot /Subtype /Link /Border [0 0 0] /H /I /C [1 0 0] /Rect [71.004 526.49 107.995 537.282] /A << /S /GoTo /D <74686D2E452E322E332E3130> >> >>}
\MusicPDFWriteObject{8177}{<< /Type /Annot /Subtype /Link /Border [0 0 0] /H /I /C [1 0 0] /Rect [231.524 505.758 263.835 517.711] /A << /S /GoTo /D <74686D2E452E322E322E37> >> >>}
\MusicPDFWriteObject{8178}{<< /Type /Annot /Subtype /Link /Border [0 0 0] /H /I /C [0 1 0] /Rect [310.71 508.248 327.85 517.116] /A << /S /GoTo /D <636974652E4475> >> >>}
\MusicPDFWriteObject{8179}{<< /Type /Annot /Subtype /Link /Border [0 0 0] /H /I /C [0 1 0] /Rect [372.344 455.645 390.405 464.513] /A << /S /GoTo /D <636974652E537465> >> >>}
\MusicPDFWriteObject{8180}{<< /Type /Annot /Subtype /Link /Border [0 0 0] /H /I /C [0 1 0] /Rect [85.578 409.164 110.002 418.032] /A << /S /GoTo /D <636974652E53745365> >> >>}
\MusicPDFWriteObject{8181}{<< /Type /Annot /Subtype /Link /Border [0 0 0] /H /I /C [1 0 0] /Rect [71.004 378.517 108.24 392.327] /A << /S /GoTo /D <74686D2E452E322E332E3130> >> >>}
\MusicPDFWriteObject{8182}{<< /Type /Annot /Subtype /Link /Border [0 0 0] /H /I /C [0 1 0] /Rect [434.504 317.108 452.564 325.976] /A << /S /GoTo /D <636974652E537465> >> >>}
\MusicPDFWriteObject{8183}{<< /Type /Annot /Subtype /Link /Border [0 0 0] /H /I /C [1 0 0] /Rect [445.837 288.87 483.528 299.661] /A << /S /GoTo /D <74686D2E452E322E332E3131> >> >>}
\MusicPDFWriteObject{8184}{<< /Type /Annot /Subtype /Link /Border [0 0 0] /H /I /C [0 1 0] /Rect [206.518 264.505 224.578 273.373] /A << /S /GoTo /D <636974652E537465> >> >>}
\MusicPDFWriteObject{8185}{<< /Type /Annot /Subtype /Link /Border [0 0 0] /H /I /C [1 0 0] /Rect [367.259 228.673 404.95 240.977] /A << /S /GoTo /D <74686D2E452E322E332E3132> >> >>}
\MusicPDFWriteObject{8186}{<< /Type /Annot /Subtype /Link /Border [0 0 0] /H /I /C [1 0 0] /Rect [464.027 228.673 469.99 240.977] /A << /S /GoTo /D <48666F6F746E6F74652E313630> >> >>}
\MusicPDFWriteObject{8187}{<< /Type /Annot /Subtype /Link /Border [0 0 0] /H /I /C [1 0 0] /Rect [506.978 189.221 543.27 200.012] /A << /S /GoTo /D <74686D2E452E322E312E3131> >> >>}
\MusicPDFWriteObject{8188}{<< /Type /Annot /Subtype /Link /Border [0 0 0] /H /I /C [1 0 0] /Rect [239.553 162.919 249.486 175.223] /A << /S /GoTo /D <48666F6F746E6F74652E313631> >> >>}
\MusicPDFWriteObject{8189}{<< /Type /Annot /Subtype /Link /Border [0 0 0] /H /I /C [1 0 0] /Rect [351.115 129.589 387.406 140.38] /A << /S /GoTo /D <74686D2E452E312E322E3136> >> >>}
\MusicPDFWriteObject{8190}{<< /Type /Annot /Subtype /Link /Border [0 0 0] /H /I /C [1 0 0] /Rect [352.351 606.282 362.284 618.586] /A << /S /GoTo /D <48666F6F746E6F74652E313632> >> >>}
\MusicPDFWriteObject{8191}{<< /Type /Annot /Subtype /Link /Border [0 0 0] /H /I /C [0 1 0] /Rect [298.262 595.068 324.575 603.936] /A << /S /GoTo /D <636974652E426F7232> >> >>}
\MusicPDFWriteObject{8192}{<< /Type /Annot /Subtype /Link /Border [0 0 0] /H /I /C [0 1 0] /Rect [379.954 327.475 406.266 336.343] /A << /S /GoTo /D <636974652E426F7232> >> >>}
\MusicPDFWriteObject{8193}{<< /Type /Annot /Subtype /Link /Border [0 0 0] /H /I /C [1 0 0] /Rect [125.423 259.784 158.035 270.575] /A << /S /GoTo /D <74686D2E452E322E332E38> >> >>}
\MusicPDFWriteObject{8194}{<< /Type /Annot /Subtype /Link /Border [0 0 0] /H /I /C [1 0 0] /Rect [105.984 246.08 137.455 258.033] /A << /S /GoTo /D <74686D2E452E322E332E36> >> >>}
\MusicPDFWriteObject{8195}{<< /Type /Annot /Subtype /Link /Border [0 0 0] /H /I /C [1 0 0] /Rect [268.768 206.628 306.46 218.581] /A << /S /GoTo /D <74686D2E452E322E332E3135> >> >>}
\MusicPDFWriteObject{8196}{<< /Type /Annot /Subtype /Link /Border [0 0 0] /H /I /C [1 0 0] /Rect [172.189 194.031 208.551 204.821] /A << /S /GoTo /D <74686D2E452E312E322E3136> >> >>}
\MusicPDFWriteObject{8197}{<< /Type /Annot /Subtype /Link /Border [0 0 0] /H /I /C [0 1 0] /Rect [437.537 169.666 455.065 178.534] /A << /S /GoTo /D <636974652E537472> >> >>}
\MusicPDFWriteObject{8198}{<< /Type /Annot /Subtype /Link /Border [0 0 0] /H /I /C [0 1 0] /Rect [444.669 156.515 487.692 165.383] /A << /S /GoTo /D <636974652E426F752D4754> >> >>}
\MusicPDFWriteObject{8199}{<< /Type /Annot /Subtype /Link /Border [0 0 0] /H /I /C [1 0 0] /Rect [106.556 143.364 143.548 152.218] /A << /S /GoTo /D <74686D2E452E322E332E3130> >> >>}
\MusicPDFWriteObject{8200}{<< /Type /Annot /Subtype /Link /Border [0 0 0] /H /I /C [1 0 0] /Rect [199.795 143.364 231.806 152.218] /A << /S /GoTo /D <74686D2E452E322E322E35> >> >>}
\MusicPDFWriteObject{8201}{<< /Type /Annot /Subtype /Link /Border [0 0 0] /H /I /C [1 0 0] /Rect [213.256 128.277 249.653 139.068] /A << /S /GoTo /D <74686D2E452E322E332E3135> >> >>}
\MusicPDFWriteObject{8202}{<< /Type /Annot /Subtype /Link /Border [0 0 0] /H /I /C [1 0 0] /Rect [495.408 477.441 532.399 489.396] /A << /S /GoTo /D <74686D2E452E322E332E3135> >> >>}
\MusicPDFWriteObject{8203}{<< /Type /Annot /Subtype /Link /Border [0 0 0] /H /I /C [1 0 0] /Rect [161.467 443.36 198.459 455.315] /A << /S /GoTo /D <6571756174696F6E2E452E322E332E3132> >> >>}
\MusicPDFWriteObject{8204}{<< /Type /Annot /Subtype /Link /Border [0 0 0] /H /I /C [1 0 0] /Rect [281.18 443.36 318.171 455.315] /A << /S /GoTo /D <6571756174696F6E2E452E322E332E3133> >> >>}
\MusicPDFWriteObject{8205}{<< /Type /Annot /Subtype /Link /Border [0 0 0] /H /I /C [1 0 0] /Rect [417.39 443.36 454.382 455.315] /A << /S /GoTo /D <6571756174696F6E2E452E322E332E3134> >> >>}
\MusicPDFWriteObject{8206}{<< /Type /Annot /Subtype /Link /Border [0 0 0] /H /I /C [1 0 0] /Rect [237.781 414.008 274.772 425.963] /A << /S /GoTo /D <6571756174696F6E2E452E322E332E3132> >> >>}
\MusicPDFWriteObject{8207}{<< /Type /Annot /Subtype /Link /Border [0 0 0] /H /I /C [1 0 0] /Rect [140.847 122.048 178.399 134.003] /A << /S /GoTo /D <74686D2E452E322E332E3130> >> >>}
\MusicPDFWriteObject{8208}{<< /Type /Annot /Subtype /Link /Border [0 0 0] /H /I /C [1 0 0] /Rect [318.302 95.746 355.084 107.701] /A << /S /GoTo /D <74686D2E452E322E332E3130> >> >>}
\MusicPDFWriteObject{8209}{<< /Type /Annot /Subtype /Link /Border [0 0 0] /H /I /C [1 0 0] /Rect [252.743 618.879 262.676 631.737] /A << /S /GoTo /D <48666F6F746E6F74652E313633> >> >>}
\MusicPDFWriteObject{8210}{<< /Type /Annot /Subtype /Link /Border [0 0 0] /H /I /C [1 0 0] /Rect [266.2 605.729 302.492 617.681] /A << /S /GoTo /D <74686D2E452E322E332E3130> >> >>}
\MusicPDFWriteObject{8211}{<< /Type /Annot /Subtype /Link /Border [0 0 0] /H /I /C [1 0 0] /Rect [443.823 605.729 480.115 617.681] /A << /S /GoTo /D <6571756174696F6E2E452E322E332E3135> >> >>}
\MusicPDFWriteObject{8212}{<< /Type /Annot /Subtype /Link /Border [0 0 0] /H /I /C [1 0 0] /Rect [480.771 592.578 517.062 604.531] /A << /S /GoTo /D <6571756174696F6E2E452E322E332E3136> >> >>}
\MusicPDFWriteObject{8213}{<< /Type /Annot /Subtype /Link /Border [0 0 0] /H /I /C [1 0 0] /Rect [340.141 566.276 376.888 578.869] /A << /S /GoTo /D <6571756174696F6E2E452E322E332E3134> >> >>}
\MusicPDFWriteObject{8214}{<< /Type /Annot /Subtype /Link /Border [0 0 0] /H /I /C [1 0 0] /Rect [207.325 526.837 244.316 537.628] /A << /S /GoTo /D <74686D2E452E322E332E3135> >> >>}
\MusicPDFWriteObject{8215}{<< /Type /Annot /Subtype /Link /Border [0 0 0] /H /I /C [1 0 0] /Rect [203.197 501.932 244.145 510.786] /A << /S /GoTo /D <74686D2E452E322E332E3135> >> >>}
\MusicPDFWriteObject{8216}{<< /Type /Annot /Subtype /Link /Border [0 0 0] /H /I /C [1 0 0] /Rect [351.592 325.95 388.688 337.905] /A << /S /GoTo /D <74686D2E452E322E332E3136> >> >>}
\MusicPDFWriteObject{8217}{<< /Type /Annot /Subtype /Link /Border [0 0 0] /H /I /C [1 0 0] /Rect [422.967 312.8 459.328 325.392] /A << /S /GoTo /D <74686D2E452E322E332E3137> >> >>}
\MusicPDFWriteObject{8218}{<< /Type /Annot /Subtype /Link /Border [0 0 0] /H /I /C [1 0 0] /Rect [155.953 273.348 193.645 285.301] /A << /S /GoTo /D <74686D2E452E322E332E3136> >> >>}
\MusicPDFWriteObject{8219}{<< /Type /Annot /Subtype /Link /Border [0 0 0] /H /I /C [1 0 0] /Rect [397.893 247.046 429.934 258.999] /A << /S /GoTo /D <74686D2E452E322E332E38> >> >>}
\MusicPDFWriteObject{8220}{<< /Type /Annot /Subtype /Link /Border [0 0 0] /H /I /C [1 0 0] /Rect [380.848 220.745 413.279 232.698] /A << /S /GoTo /D <74686D2E452E322E332E38> >> >>}
\MusicPDFWriteObject{8221}{<< /Type /Annot /Subtype /Link /Border [0 0 0] /H /I /C [1 0 0] /Rect [106.556 207.594 138.568 219.547] /A << /S /GoTo /D <74686D2E452E322E332E36> >> >>}
\MusicPDFWriteObject{8222}{<< /Type /Annot /Subtype /Link /Border [0 0 0] /H /I /C [1 0 0] /Rect [381.605 168.142 391.538 180.999] /A << /S /GoTo /D <48666F6F746E6F74652E313634> >> >>}
\MusicPDFWriteObject{8223}{<< /Type /Annot /Subtype /Link /Border [0 0 0] /H /I /C [1 0 0] /Rect [313.591 154.991 345.332 166.944] /A << /S /GoTo /D <74686D2E452E322E332E36> >> >>}
\MusicPDFWriteObject{8224}{<< /Type /Annot /Subtype /Link /Border [0 0 0] /H /I /C [1 0 0] /Rect [358.501 141.84 390.513 153.793] /A << /S /GoTo /D <74686D2E452E322E322E35> >> >>}
\MusicPDFWriteObject{8225}{<< /Type /Annot /Subtype /Link /Border [0 0 0] /H /I /C [1 0 0] /Rect [71.004 128.69 102.415 140.642] /A << /S /GoTo /D <74686D2E452E322E322E35> >> >>}
\MusicPDFWriteObject{8226}{<< /Type /Annot /Subtype /Link /Border [0 0 0] /H /I /C [1 0 0] /Rect [508.61 128.69 540.021 140.642] /A << /S /GoTo /D <74686D2E452E322E332E36> >> >>}
\MusicPDFWriteObject{8227}{<< /Type /Annot /Subtype /Link /Border [0 0 0] /H /I /C [0 1 0] /Rect [358.426 687.123 375.068 695.991] /A << /S /GoTo /D <636974652E764F> >> >>}
\MusicPDFWriteObject{8228}{<< /Type /Annot /Subtype /Link /Border [0 0 0] /H /I /C [0 1 0] /Rect [437.838 687.123 464.15 695.991] /A << /S /GoTo /D <636974652E426F7232> >> >>}
\MusicPDFWriteObject{8229}{<< /Type /Annot /Subtype /Link /Border [0 0 0] /H /I /C [0 1 0] /Rect [99.833 673.973 116.475 682.84] /A << /S /GoTo /D <636974652E764F> >> >>}
\MusicPDFWriteObject{8230}{<< /Type /Annot /Subtype /Link /Border [0 0 0] /H /I /C [1 0 0] /Rect [161.683 671.482 171.616 684.34] /A << /S /GoTo /D <48666F6F746E6F74652E313635> >> >>}
\MusicPDFWriteObject{8231}{<< /Type /Annot /Subtype /Link /Border [0 0 0] /H /I /C [1 0 0] /Rect [397.256 598.8 429.267 610.753] /A << /S /GoTo /D <74686D2E452E312E312E37> >> >>}
\MusicPDFWriteObject{8232}{<< /Type /Annot /Subtype /Link /Border [0 0 0] /H /I /C [1 0 0] /Rect [436.166 504.037 468.178 514.828] /A << /S /GoTo /D <74686D2E452E312E312E33> >> >>}
\MusicPDFWriteObject{8233}{<< /Type /Annot /Subtype /Link /Border [0 0 0] /H /I /C [0 1 0] /Rect [500.216 121.34 543.239 130.208] /A << /S /GoTo /D <636974652E426F752D4754> >> >>}
\MusicPDFWriteObject{8234}{<< /Type /Annot /Subtype /Link /Border [0 0 0] /H /I /C [1 0 0] /Rect [514.181 639.088 545.592 649.878] /A << /S /GoTo /D <74686D2E452E322E342E31> >> >>}
\MusicPDFWriteObject{8235}{<< /Type /Annot /Subtype /Link /Border [0 0 0] /H /I /C [1 0 0] /Rect [71.004 612.786 103.015 623.577] /A << /S /GoTo /D <74686D2E452E312E312E37> >> >>}
\MusicPDFWriteObject{8236}{<< /Type /Annot /Subtype /Link /Border [0 0 0] /H /I /C [1 0 0] /Rect [118.43 490.961 151.042 502.913] /A << /S /GoTo /D <74686D2E452E312E312E37> >> >>}
\MusicPDFWriteObject{8237}{<< /Type /Annot /Subtype /Link /Border [0 0 0] /H /I /C [1 0 0] /Rect [485.836 490.961 518.448 502.913] /A << /S /GoTo /D <74686D2E452E312E312E33> >> >>}
\MusicPDFWriteObject{8238}{<< /Type /Annot /Subtype /Link /Border [0 0 0] /H /I /C [1 0 0] /Rect [275.963 464.659 307.974 476.612] /A << /S /GoTo /D <74686D2E452E322E342E32> >> >>}
\MusicPDFWriteObject{8239}{<< /Type /Annot /Subtype /Link /Border [0 0 0] /H /I /C [0 1 0] /Rect [277.934 374.728 302.358 383.596] /A << /S /GoTo /D <636974652E53745365> >> >>}
\MusicPDFWriteObject{8240}{<< /Type /Annot /Subtype /Link /Border [0 0 0] /H /I /C [1 0 0] /Rect [378.997 339.62 411.609 350.411] /A << /S /GoTo /D <74686D2E452E322E312E34> >> >>}
\MusicPDFWriteObject{8241}{<< /Type /Annot /Subtype /Link /Border [0 0 0] /H /I /C [1 0 0] /Rect [357.361 260.127 389.973 270.917] /A << /S /GoTo /D <74686D2E452E322E342E34> >> >>}
\MusicPDFWriteObject{8242}{<< /Type /Annot /Subtype /Link /Border [0 0 0] /H /I /C [1 0 0] /Rect [304.819 246.423 341.985 258.375] /A << /S /GoTo /D <74686D2E452E312E322E3231> >> >>}
\MusicPDFWriteObject{8243}{<< /Type /Annot /Subtype /Link /Border [0 0 0] /H /I /C [0 1 0] /Rect [460.252 248.913 484.676 257.781] /A << /S /GoTo /D <636974652E53745365> >> >>}
\MusicPDFWriteObject{8244}{<< /Type /Annot /Subtype /Link /Border [0 0 0] /H /I /C [1 0 0] /Rect [251.215 122.601 283.647 133.392] /A << /S /GoTo /D <74686D2E452E322E342E36> >> >>}
\MusicPDFWriteObject{8245}{<< /Type /Annot /Subtype /Link /Border [0 0 0] /H /I /C [1 0 0] /Rect [439.553 122.601 477.035 133.392] /A << /S /GoTo /D <74686D2E452E312E322E3231> >> >>}
\MusicPDFWriteObject{8246}{<< /Type /Annot /Subtype /Link /Border [0 0 0] /H /I /C [1 0 0] /Rect [71.004 559.469 103.165 572.45] /A << /S /GoTo /D <74686D2E452E322E342E32> >> >>}
\MusicPDFWriteObject{8247}{<< /Type /Annot /Subtype /Link /Border [0 0 0] /H /I /C [1 0 0] /Rect [396.932 521.043 429.543 533.015] /A << /S /GoTo /D <74686D2E452E312E312E33> >> >>}
\MusicPDFWriteObject{8248}{<< /Type /Annot /Subtype /Link /Border [0 0 0] /H /I /C [1 0 0] /Rect [213.866 507.892 245.277 519.845] /A << /S /GoTo /D <74686D2E452E312E312E33> >> >>}
\MusicPDFWriteObject{8249}{<< /Type /Annot /Subtype /Link /Border [0 0 0] /H /I /C [1 0 0] /Rect [473.958 375.101 505.369 385.892] /A << /S /GoTo /D <74686D2E452E322E312E34> >> >>}
\MusicPDFWriteObject{8250}{<< /Type /Annot /Subtype /Link /Border [0 0 0] /H /I /C [1 0 0] /Rect [458.448 264.114 491.06 275.975] /A << /S /GoTo /D <74686D2E452E322E342E34> >> >>}
\MusicPDFWriteObject{8251}{<< /Type /Annot /Subtype /Link /Border [0 0 0] /H /I /C [0 1 0] /Rect [447.955 252.9 472.379 261.768] /A << /S /GoTo /D <636974652E53745365> >> >>}
\MusicPDFWriteObject{8252}{<< /Type /Annot /Subtype /Link /Border [0 0 0] /H /I /C [1 0 0] /Rect [176.189 210.957 208.201 222.91] /A << /S /GoTo /D <74686D2E452E322E312E36> >> >>}
\MusicPDFWriteObject{8253}{<< /Type /Annot /Subtype /Link /Border [0 0 0] /H /I /C [0 1 0] /Rect [73.77 85.086 98.194 93.954] /A << /S /GoTo /D <636974652E53745365> >> >>}
\MusicPDFWriteObject{8254}{<< /Type /Annot /Subtype /Link /Border [0 0 0] /H /I /C [1 0 0] /Rect [272.032 651.617 309.374 662.408] /A << /S /GoTo /D <74686D2E452E322E342E3130> >> >>}
\MusicPDFWriteObject{8255}{<< /Type /Annot /Subtype /Link /Border [0 0 0] /H /I /C [1 0 0] /Rect [193.336 624.762 231.028 636.715] /A << /S /GoTo /D <74686D2E452E312E322E3235> >> >>}
\MusicPDFWriteObject{8256}{<< /Type /Annot /Subtype /Link /Border [0 0 0] /H /I /C [1 0 0] /Rect [331.316 542.382 368.307 553.173] /A << /S /GoTo /D <74686D2E452E322E342E3130> >> >>}
\MusicPDFWriteObject{8257}{<< /Type /Annot /Subtype /Link /Border [0 0 0] /H /I /C [1 0 0] /Rect [211.451 529.231 248.092 540.022] /A << /S /GoTo /D <74686D2E452E322E342E3131> >> >>}
\MusicPDFWriteObject{8258}{<< /Type /Annot /Subtype /Link /Border [0 0 0] /H /I /C [1 0 0] /Rect [329.167 515.527 366.158 528.385] /A << /S /GoTo /D <74686D2E452E322E342E3131> >> >>}
\MusicPDFWriteObject{8259}{<< /Type /Annot /Subtype /Link /Border [0 0 0] /H /I /C [1 0 0] /Rect [370.806 515.527 380.739 528.385] /A << /S /GoTo /D <48666F6F746E6F74652E313636> >> >>}
\MusicPDFWriteObject{8260}{<< /Type /Annot /Subtype /Link /Border [0 0 0] /H /I /C [1 0 0] /Rect [374.57 495.378 412.261 506.168] /A << /S /GoTo /D <74686D2E452E312E322E3235> >> >>}
\MusicPDFWriteObject{8261}{<< /Type /Annot /Subtype /Link /Border [0 0 0] /H /I /C [1 0 0] /Rect [507.154 495.378 544.845 506.168] /A << /S /GoTo /D <74686D2E452E322E342E3132> >> >>}
\MusicPDFWriteObject{8262}{<< /Type /Annot /Subtype /Link /Border [0 0 0] /H /I /C [0 1 0] /Rect [97.662 427.216 113.848 436.084] /A << /S /GoTo /D <636974652E5350> >> >>}
\MusicPDFWriteObject{8263}{<< /Type /Annot /Subtype /Link /Border [0 0 0] /H /I /C [0 1 0] /Rect [384.041 172.597 419.352 181.465] /A << /S /GoTo /D <636974652E4B756E656E> >> >>}
\MusicPDFWriteObject{8264}{<< /Type /Annot /Subtype /Link /Border [0 0 0] /H /I /C [0 1 0] /Rect [487.567 172.597 511.866 181.465] /A << /S /GoTo /D <636974652E4A656368> >> >>}
\MusicPDFWriteObject{8265}{<< /Type /Annot /Subtype /Link /Border [0 0 0] /H /I /C [0 1 0] /Rect [219.917 159.446 236.102 168.314] /A << /S /GoTo /D <636974652E5350> >> >>}
\MusicPDFWriteObject{8266}{<< /Type /Annot /Subtype /Link /Border [0 0 0] /H /I /C [0 1 0] /Rect [367.354 146.296 383.54 155.163] /A << /S /GoTo /D <636974652E5350> >> >>}
\MusicPDFWriteObject{8267}{<< /Type /Annot /Subtype /Link /Border [0 0 0] /H /I /C [0 1 0] /Rect [526.99 133.145 543.175 142.013] /A << /S /GoTo /D <636974652E5350> >> >>}
\MusicPDFWriteObject{8268}{<< /Type /Annot /Subtype /Link /Border [0 0 0] /H /I /C [0 1 0] /Rect [351.028 538.62 386.339 547.488] /A << /S /GoTo /D <636974652E4B756E656E> >> >>}
\MusicPDFWriteObject{8269}{<< /Type /Annot /Subtype /Link /Border [0 0 0] /H /I /C [0 1 0] /Rect [103.682 512.319 127.981 521.187] /A << /S /GoTo /D <636974652E4A656368> >> >>}
\MusicPDFWriteObject{8270}{<< /Type /Annot /Subtype /Link /Border [0 0 0] /H /I /C [0 1 0] /Rect [293.101 512.319 309.287 521.187] /A << /S /GoTo /D <636974652E5350> >> >>}
\MusicPDFWriteObject{8271}{<< /Type /Annot /Subtype /Link /Border [0 0 0] /H /I /C [0 1 0] /Rect [386.467 486.017 421.778 494.885] /A << /S /GoTo /D <636974652E4B756E656E> >> >>}
\MusicPDFWriteObject{8272}{<< /Type /Annot /Subtype /Link /Border [0 0 0] /H /I /C [0 1 0] /Rect [111.144 416.055 146.455 424.923] /A << /S /GoTo /D <636974652E4B756E656E> >> >>}
\MusicPDFWriteObject{8273}{<< /Type /Annot /Subtype /Link /Border [0 0 0] /H /I /C [1 0 0] /Rect [262.44 252.785 272.373 266.479] /A << /S /GoTo /D <48666F6F746E6F74652E313637> >> >>}
\MusicPDFWriteObject{8274}{<< /Type /Annot /Subtype /Link /Border [0 0 0] /H /I /C [0 1 0] /Rect [87.498 242.961 122.808 251.829] /A << /S /GoTo /D <636974652E4B756E656E> >> >>}
\MusicPDFWriteObject{8275}{<< /Type /Annot /Subtype /Link /Border [0 0 0] /H /I /C [1 0 0] /Rect [178.395 162.09 210.346 174.394] /A << /S /GoTo /D <74686D2E452E332E312E32> >> >>}
\MusicPDFWriteObject{8276}{<< /Type /Annot /Subtype /Link /Border [0 0 0] /H /I /C [1 0 0] /Rect [429.703 162.09 439.636 174.394] /A << /S /GoTo /D <48666F6F746E6F74652E313638> >> >>}
\MusicPDFWriteObject{8277}{<< /Type /Annot /Subtype /Link /Border [0 0 0] /H /I /C [1 0 0] /Rect [332.439 595.859 365.051 607.501] /A << /S /GoTo /D <74686D2E452E332E312E32> >> >>}
\MusicPDFWriteObject{8278}{<< /Type /Annot /Subtype /Link /Border [0 0 0] /H /I /C [1 0 0] /Rect [283.623 497 315.635 508.953] /A << /S /GoTo /D <6571756174696F6E2E452E342E312E31> >> >>}
\MusicPDFWriteObject{8279}{<< /Type /Annot /Subtype /Link /Border [0 0 0] /H /I /C [1 0 0] /Rect [278.97 187.572 311.581 199.214] /A << /S /GoTo /D <74686D2E452E332E312E32> >> >>}
\MusicPDFWriteObject{8280}{<< /Type /Annot /Subtype /Link /Border [0 0 0] /H /I /C [1 0 0] /Rect [203.82 134.901 236.432 147.171] /A << /S /GoTo /D <74686D2E452E332E312E32> >> >>}
\MusicPDFWriteObject{8281}{<< /Type /Annot /Subtype /Link /Border [0 0 0] /H /I /C [1 0 0] /Rect [204.235 108.897 235.646 120.87] /A << /S /GoTo /D <6571756174696F6E2E452E332E312E32> >> >>}
\MusicPDFWriteObject{8282}{<< /Type /Annot /Subtype /Link /Border [0 0 0] /H /I /C [1 0 0] /Rect [508.062 108.897 539.474 120.87] /A << /S /GoTo /D <6571756174696F6E2E452E332E312E33> >> >>}
\MusicPDFWriteObject{8283}{<< /Type /Annot /Subtype /Link /Border [0 0 0] /H /I /C [1 0 0] /Rect [152.116 95.746 184.728 107.719] /A << /S /GoTo /D <6571756174696F6E2E452E332E312E34> >> >>}
\MusicPDFWriteObject{8284}{<< /Type /Annot /Subtype /Link /Border [0 0 0] /H /I /C [1 0 0] /Rect [291.428 95.746 324.04 107.719] /A << /S /GoTo /D <6571756174696F6E2E452E332E312E35> >> >>}
\MusicPDFWriteObject{8285}{<< /Type /Annot /Subtype /Link /Border [0 0 0] /H /I /C [1 0 0] /Rect [117.417 495.424 149.429 507.379] /A << /S /GoTo /D <74686D2E452E332E312E31> >> >>}
\MusicPDFWriteObject{8286}{<< /Type /Annot /Subtype /Link /Border [0 0 0] /H /I /C [1 0 0] /Rect [339.473 95.746 370.914 107.701] /A << /S /GoTo /D <74686D2E452E332E312E31> >> >>}
\MusicPDFWriteObject{8287}{<< /Type /Annot /Subtype /Link /Border [0 0 0] /H /I /C [1 0 0] /Rect [420.622 608.753 452.994 620.708] /A << /S /GoTo /D <74686D2E452E332E312E35> >> >>}
\MusicPDFWriteObject{8288}{<< /Type /Annot /Subtype /Link /Border [0 0 0] /H /I /C [1 0 0] /Rect [196.686 563.877 228.338 574.668] /A << /S /GoTo /D <74686D2E452E332E312E32> >> >>}
\MusicPDFWriteObject{8289}{<< /Type /Annot /Subtype /Link /Border [0 0 0] /H /I /C [1 0 0] /Rect [283.972 563.877 315.624 574.668] /A << /S /GoTo /D <74686D2E452E332E312E36> >> >>}
\MusicPDFWriteObject{8290}{<< /Type /Annot /Subtype /Link /Border [0 0 0] /H /I /C [1 0 0] /Rect [412.675 563.877 449.246 574.668] /A << /S /GoTo /D <74686D2E452E322E312E3133> >> >>}
\MusicPDFWriteObject{8291}{<< /Type /Annot /Subtype /Link /Border [0 0 0] /H /I /C [0 1 0] /Rect [264.093 500.06 280.279 508.928] /A << /S /GoTo /D <636974652E5350> >> >>}
\MusicPDFWriteObject{8292}{<< /Type /Annot /Subtype /Link /Border [0 0 0] /H /I /C [1 0 0] /Rect [106.331 418.49 138.132 430.189] /A << /S /GoTo /D <74686D2E452E332E312E32> >> >>}
\MusicPDFWriteObject{8293}{<< /Type /Annot /Subtype /Link /Border [0 0 0] /H /I /C [1 0 0] /Rect [483.877 418.49 515.679 430.189] /A << /S /GoTo /D <74686D2E452E332E312E36> >> >>}
\MusicPDFWriteObject{8294}{<< /Type /Annot /Subtype /Link /Border [0 0 0] /H /I /C [1 0 0] /Rect [192.463 318.271 229.455 329.062] /A << /S /GoTo /D <74686D2E452E322E312E3131> >> >>}
\MusicPDFWriteObject{8295}{<< /Type /Annot /Subtype /Link /Border [0 0 0] /H /I /C [1 0 0] /Rect [296.346 255.542 333.128 266.333] /A << /S /GoTo /D <74686D2E452E322E312E3133> >> >>}
\MusicPDFWriteObject{8296}{<< /Type /Annot /Subtype /Link /Border [0 0 0] /H /I /C [1 0 0] /Rect [71.004 215.239 103.105 226.88] /A << /S /GoTo /D <74686D2E452E332E312E36> >> >>}
\MusicPDFWriteObject{8297}{<< /Type /Annot /Subtype /Link /Border [0 0 0] /H /I /C [1 0 0] /Rect [234.493 612.007 266.805 622.797] /A << /S /GoTo /D <74686D2E452E332E322E31> >> >>}
\MusicPDFWriteObject{8298}{<< /Type /Annot /Subtype /Link /Border [0 0 0] /H /I /C [0 1 0] /Rect [376.076 600.793 411.387 609.66] /A << /S /GoTo /D <636974652E4B756E656E> >> >>}
\MusicPDFWriteObject{8299}{<< /Type /Annot /Subtype /Link /Border [0 0 0] /H /I /C [0 1 0] /Rect [73.723 561.34 98.023 570.208] /A << /S /GoTo /D <636974652E4A656368> >> >>}
\MusicPDFWriteObject{8300}{<< /Type /Annot /Subtype /Link /Border [0 0 0] /H /I /C [0 1 0] /Rect [253.903 527.613 270.089 536.48] /A << /S /GoTo /D <636974652E5350> >> >>}
\MusicPDFWriteObject{8301}{<< /Type /Annot /Subtype /Link /Border [0 0 0] /H /I /C [1 0 0] /Rect [111.585 412.746 143.596 424.719] /A << /S /GoTo /D <6571756174696F6E2E452E332E312E31> >> >>}
\MusicPDFWriteObject{8302}{<< /Type /Annot /Subtype /Link /Border [0 0 0] /H /I /C [0 1 0] /Rect [478.333 368.358 494.519 377.226] /A << /S /GoTo /D <636974652E5350> >> >>}
\MusicPDFWriteObject{8303}{<< /Type /Annot /Subtype /Link /Border [0 0 0] /H /I /C [0 1 0] /Rect [261.621 236.552 277.807 245.42] /A << /S /GoTo /D <636974652E5350> >> >>}
\MusicPDFWriteObject{8304}{<< /Type /Annot /Subtype /Link /Border [0 0 0] /H /I /C [1 0 0] /Rect [389.781 208.314 422.393 219.105] /A << /S /GoTo /D <74686D2E452E332E312E37> >> >>}
\MusicPDFWriteObject{8305}{<< /Type /Annot /Subtype /Link /Border [0 0 0] /H /I /C [0 1 0] /Rect [345.67 197.1 361.856 205.968] /A << /S /GoTo /D <636974652E5350> >> >>}
\MusicPDFWriteObject{8306}{<< /Type /Annot /Subtype /Link /Border [0 0 0] /H /I /C [1 0 0] /Rect [122.672 168.308 155.284 180.263] /A << /S /GoTo /D <74686D2E452E332E312E34> >> >>}
\MusicPDFWriteObject{8307}{<< /Type /Annot /Subtype /Link /Border [0 0 0] /H /I /C [0 1 0] /Rect [370.823 631.989 387.009 640.857] /A << /S /GoTo /D <636974652E5350> >> >>}
\MusicPDFWriteObject{8308}{<< /Type /Annot /Subtype /Link /Border [0 0 0] /H /I /C [0 1 0] /Rect [384.41 450.501 400.596 459.368] /A << /S /GoTo /D <636974652E5350> >> >>}
\MusicPDFWriteObject{8309}{<< /Type /Annot /Subtype /Link /Border [0 0 0] /H /I /C [0 1 0] /Rect [438.049 333.245 454.235 342.113] /A << /S /GoTo /D <636974652E5350> >> >>}
\MusicPDFWriteObject{8310}{<< /Type /Annot /Subtype /Link /Border [0 0 0] /H /I /C [1 0 0] /Rect [152.986 267.746 184.998 279.699] /A << /S /GoTo /D <74686D2E452E332E312E31> >> >>}
\MusicPDFWriteObject{8311}{<< /Type /Annot /Subtype /Link /Border [0 0 0] /H /I /C [1 0 0] /Rect [465.491 177.382 501.782 190.838] /A << /S /GoTo /D <74686D2E452E322E312E3133> >> >>}
\MusicPDFWriteObject{8312}{<< /Type /Annot /Subtype /Link /Border [0 0 0] /H /I /C [1 0 0] /Rect [167.77 153.135 205.461 163.926] /A << /S /GoTo /D <74686D2E452E312E322E3235> >> >>}
\MusicPDFWriteObject{8313}{<< /Type /Annot /Subtype /Link /Border [0 0 0] /H /I /C [1 0 0] /Rect [454.243 153.135 491.934 163.926] /A << /S /GoTo /D <74686D2E452E322E312E3133> >> >>}
\MusicPDFWriteObject{8314}{<< /Type /Annot /Subtype /Link /Border [0 0 0] /H /I /C [1 0 0] /Rect [338.12 597.312 370.012 608.103] /A << /S /GoTo /D <74686D2E452E322E332E31> >> >>}
\MusicPDFWriteObject{8315}{<< /Type /Annot /Subtype /Link /Border [0 0 0] /H /I /C [1 0 0] /Rect [500.382 581.92 532.844 596.2] /A << /S /GoTo /D <74686D2E452E322E322E35> >> >>}
\MusicPDFWriteObject{8316}{<< /Type /Annot /Subtype /Link /Border [0 0 0] /H /I /C [1 0 0] /Rect [125.923 570.457 157.935 582.41] /A << /S /GoTo /D <74686D2E452E312E322E32> >> >>}
\MusicPDFWriteObject{8317}{<< /Type /Annot /Subtype /Link /Border [0 0 0] /H /I /C [1 0 0] /Rect [227.047 570.457 259.058 582.41] /A << /S /GoTo /D <74686D2E452E332E332E32> >> >>}
\MusicPDFWriteObject{8318}{<< /Type /Annot /Subtype /Link /Border [0 0 0] /H /I /C [1 0 0] /Rect [139.757 528.308 171.768 539.099] /A << /S /GoTo /D <74686D2E452E332E332E31> >> >>}
\MusicPDFWriteObject{8319}{<< /Type /Annot /Subtype /Link /Border [0 0 0] /H /I /C [0 1 0] /Rect [89.402 323.777 124.713 332.645] /A << /S /GoTo /D <636974652E4B756E656E> >> >>}
\MusicPDFWriteObject{8320}{<< /Type /Annot /Subtype /Link /Border [0 0 0] /H /I /C [0 1 0] /Rect [89.781 173.739 125.092 182.607] /A << /S /GoTo /D <636974652E4B756E656E> >> >>}
\MusicPDFWriteObject{8321}{<< /Type /Annot /Subtype /Link /Border [0 0 0] /H /I /C [1 0 0] /Rect [344.83 131.797 376.842 143.749] /A << /S /GoTo /D <6571756174696F6E2E452E332E332E31> >> >>}
\MusicPDFWriteObject{8322}{<< /Type /Annot /Subtype /Link /Border [0 0 0] /H /I /C [1 0 0] /Rect [286.342 628.474 317.994 639.264] /A << /S /GoTo /D <74686D2E452E332E322E33> >> >>}
\MusicPDFWriteObject{8323}{<< /Type /Annot /Subtype /Link /Border [0 0 0] /H /I /C [1 0 0] /Rect [309.717 537.36 341.729 548.151] /A << /S /GoTo /D <74686D2E452E332E332E32> >> >>}
\MusicPDFWriteObject{8324}{<< /Type /Annot /Subtype /Link /Border [0 0 0] /H /I /C [1 0 0] /Rect [122.188 495.157 154.2 507.13] /A << /S /GoTo /D <74686D2E452E332E322E33> >> >>}
\MusicPDFWriteObject{8325}{<< /Type /Annot /Subtype /Link /Border [0 0 0] /H /I /C [1 0 0] /Rect [348.05 428.724 381.858 439.515] /A << /S /GoTo /D <74686D2E452E332E322E33> >> >>}
\MusicPDFWriteObject{8326}{<< /Type /Annot /Subtype /Link /Border [0 0 0] /H /I /C [1 0 0] /Rect [379.311 307.277 416.478 319.25] /A << /S /GoTo /D <74686D2E452E312E322E3235> >> >>}
\MusicPDFWriteObject{8327}{<< /Type /Annot /Subtype /Link /Border [0 0 0] /H /I /C [1 0 0] /Rect [181.443 273.841 214.055 285.794] /A << /S /GoTo /D <74686D2E452E332E322E33> >> >>}
\MusicPDFWriteObject{8328}{<< /Type /Annot /Subtype /Link /Border [0 0 0] /H /I /C [1 0 0] /Rect [71.004 208.087 108.135 220.04] /A << /S /GoTo /D <74686D2E452E322E312E3133> >> >>}
\MusicPDFWriteObject{8329}{<< /Type /Annot /Subtype /Link /Border [0 0 0] /H /I /C [1 0 0] /Rect [164.683 208.087 196.815 220.04] /A << /S /GoTo /D <74686D2E452E332E332E34> >> >>}
\MusicPDFWriteObject{8330}{<< /Type /Annot /Subtype /Link /Border [0 0 0] /H /I /C [1 0 0] /Rect [245.127 82.596 277.739 94.549] /A << /S /GoTo /D <74686D2E452E332E342E31> >> >>}
\MusicPDFWriteObject{8331}{<< /Type /Annot /Subtype /Link /Border [0 0 0] /H /I /C [1 0 0] /Rect [482.925 672.036 515.537 682.827] /A << /S /GoTo /D <74686D2E452E332E342E33> >> >>}
\MusicPDFWriteObject{8332}{<< /Type /Annot /Subtype /Link /Border [0 0 0] /H /I /C [1 0 0] /Rect [158.04 658.885 190.412 669.676] /A << /S /GoTo /D <74686D2E452E332E342E31> >> >>}
\MusicPDFWriteObject{8333}{<< /Type /Annot /Subtype /Link /Border [0 0 0] /H /I /C [1 0 0] /Rect [124.623 618.879 156.634 630.832] /A << /S /GoTo /D <74686D2E452E332E342E31> >> >>}
\MusicPDFWriteObject{8334}{<< /Type /Annot /Subtype /Link /Border [0 0 0] /H /I /C [1 0 0] /Rect [391.713 588.62 423.124 600.576] /A << /S /GoTo /D <74686D2E452E322E312E32> >> >>}
\MusicPDFWriteObject{8335}{<< /Type /Annot /Subtype /Link /Border [0 0 0] /H /I /C [1 0 0] /Rect [122.957 519.975 159.249 532.288] /A << /S /GoTo /D <74686D2E452E322E322E3130> >> >>}
\MusicPDFWriteObject{8336}{<< /Type /Annot /Subtype /Link /Border [0 0 0] /H /I /C [0 1 0] /Rect [411.64 522.826 446.95 531.694] /A << /S /GoTo /D <636974652E4B756E656E> >> >>}
\MusicPDFWriteObject{8337}{<< /Type /Annot /Subtype /Link /Border [0 0 0] /H /I /C [0 1 0] /Rect [422.396 509.675 457.707 518.543] /A << /S /GoTo /D <636974652E4B756E656E> >> >>}
\MusicPDFWriteObject{8338}{<< /Type /Annot /Subtype /Link /Border [0 0 0] /H /I /C [1 0 0] /Rect [210.524 401.6 242.536 412.391] /A << /S /GoTo /D <74686D2E452E332E342E31> >> >>}
\MusicPDFWriteObject{8339}{<< /Type /Annot /Subtype /Link /Border [0 0 0] /H /I /C [1 0 0] /Rect [179.21 375.299 216.202 386.089] /A << /S /GoTo /D <74686D2E452E322E332E3135> >> >>}
\MusicPDFWriteObject{8340}{<< /Type /Annot /Subtype /Link /Border [0 0 0] /H /I /C [0 1 0] /Rect [288.478 143.667 305.12 152.534] /A << /S /GoTo /D <636974652E764F> >> >>}
\MusicPDFWriteObject{8341}{<< /Type /Annot /Subtype /Link /Border [0 0 0] /H /I /C [1 0 0] /Rect [140.575 684.633 172.407 696.586] /A << /S /GoTo /D <74686D2E452E332E342E31> >> >>}
\MusicPDFWriteObject{8342}{<< /Type /Annot /Subtype /Link /Border [0 0 0] /H /I /C [1 0 0] /Rect [126.016 632.583 163.707 643.431] /A << /S /GoTo /D <74686D2E452E322E332E3135> >> >>}
\MusicPDFWriteObject{8343}{<< /Type /Annot /Subtype /Link /Border [0 0 0] /H /I /C [1 0 0] /Rect [186.68 605.729 223.672 617.681] /A << /S /GoTo /D <74686D2E452E322E332E3135> >> >>}
\MusicPDFWriteObject{8344}{<< /Type /Annot /Subtype /Link /Border [0 0 0] /H /I /C [1 0 0] /Rect [308.281 592.578 340.233 604.531] /A << /S /GoTo /D <74686D2E452E332E342E31> >> >>}
\MusicPDFWriteObject{8345}{<< /Type /Annot /Subtype /Link /Border [0 0 0] /H /I /C [1 0 0] /Rect [459.581 566.276 492.193 578.229] /A << /S /GoTo /D <74686D2E452E332E342E31> >> >>}
\MusicPDFWriteObject{8346}{<< /Type /Annot /Subtype /Link /Border [0 0 0] /H /I /C [1 0 0] /Rect [506.997 539.975 543.779 551.928] /A << /S /GoTo /D <74686D2E452E322E332E3135> >> >>}
\MusicPDFWriteObject{8347}{<< /Type /Annot /Subtype /Link /Border [0 0 0] /H /I /C [1 0 0] /Rect [168.441 500.522 200.122 512.477] /A << /S /GoTo /D <74686D2E452E332E342E31> >> >>}
\MusicPDFWriteObject{8348}{<< /Type /Annot /Subtype /Link /Border [0 0 0] /H /I /C [1 0 0] /Rect [372.89 447.92 405.442 459.872] /A << /S /GoTo /D <74686D2E452E322E332E36> >> >>}
\MusicPDFWriteObject{8349}{<< /Type /Annot /Subtype /Link /Border [0 0 0] /H /I /C [1 0 0] /Rect [310.431 395.317 347.353 407.269] /A << /S /GoTo /D <74686D2E452E322E332E3135> >> >>}
\MusicPDFWriteObject{8350}{<< /Type /Annot /Subtype /Link /Border [0 0 0] /H /I /C [1 0 0] /Rect [473.895 382.166 510.432 394.119] /A << /S /GoTo /D <74686D2E452E322E332E3135> >> >>}
\MusicPDFWriteObject{8351}{<< /Type /Annot /Subtype /Link /Border [0 0 0] /H /I /C [1 0 0] /Rect [106.937 355.864 139.308 367.817] /A << /S /GoTo /D <74686D2E452E332E342E31> >> >>}
\MusicPDFWriteObject{8352}{<< /Type /Annot /Subtype /Link /Border [0 0 0] /H /I /C [0 1 0] /Rect [347.336 345.204 363.979 354.071] /A << /S /GoTo /D <636974652E764F> >> >>}
\MusicPDFWriteObject{8353}{<< /Type /Annot /Subtype /Link /Border [0 0 0] /H /I /C [0 1 0] /Rect [447.912 345.204 464.554 354.071] /A << /S /GoTo /D <636974652E764F> >> >>}
\MusicPDFWriteObject{8354}{<< /Type /Annot /Subtype /Link /Border [0 0 0] /H /I /C [1 0 0] /Rect [189.678 310.177 221.179 320.968] /A << /S /GoTo /D <74686D2E452E332E342E34> >> >>}
\MusicPDFWriteObject{8355}{<< /Type /Annot /Subtype /Link /Border [0 0 0] /H /I /C [1 0 0] /Rect [351.402 205.704 384.014 217.659] /A << /S /GoTo /D <74686D2E452E332E322E33> >> >>}
\MusicPDFWriteObject{8356}{<< /Type /Annot /Subtype /Link /Border [0 0 0] /H /I /C [1 0 0] /Rect [221.668 135.752 253.079 146.543] /A << /S /GoTo /D <74686D2E452E332E342E31> >> >>}
\MusicPDFWriteObject{8357}{<< /Type /Annot /Subtype /Link /Border [0 0 0] /H /I /C [1 0 0] /Rect [513.534 135.752 544.945 146.543] /A << /S /GoTo /D <74686D2E452E332E342E34> >> >>}
\MusicPDFWriteObject{8358}{<< /Type /Annot /Subtype /Link /Border [0 0 0] /H /I /C [1 0 0] /Rect [256.602 96.3 289.214 107.091] /A << /S /GoTo /D <74686D2E452E332E342E34> >> >>}
\MusicPDFWriteObject{8359}{<< /Type /Annot /Subtype /Link /Border [0 0 0] /H /I /C [1 0 0] /Rect [256.741 82.596 289.173 94.549] /A << /S /GoTo /D <74686D2E452E332E342E34> >> >>}
\MusicPDFWriteObject{8360}{<< /Type /Annot /Subtype /Link /Border [0 0 0] /H /I /C [1 0 0] /Rect [507.364 82.596 539.796 94.549] /A << /S /GoTo /D <74686D2E452E332E352E37> >> >>}
\MusicPDFWriteObject{8361}{<< /Type /Annot /Subtype /Link /Border [0 0 0] /H /I /C [0 1 0] /Rect [286.739 687.123 302.925 695.991] /A << /S /GoTo /D <636974652E5350> >> >>}
\MusicPDFWriteObject{8362}{<< /Type /Annot /Subtype /Link /Border [0 0 0] /H /I /C [0 1 0] /Rect [381.282 687.123 398.616 695.991] /A << /S /GoTo /D <636974652E4B53> >> >>}
\MusicPDFWriteObject{8363}{<< /Type /Annot /Subtype /Link /Border [0 0 0] /H /I /C [0 1 0] /Rect [525.933 687.123 543.267 695.991] /A << /S /GoTo /D <636974652E4B53> >> >>}
\MusicPDFWriteObject{8364}{<< /Type /Annot /Subtype /Link /Border [0 0 0] /H /I /C [0 1 0] /Rect [154.56 673.973 171.894 682.84] /A << /S /GoTo /D <636974652E4B53> >> >>}
\MusicPDFWriteObject{8365}{<< /Type /Annot /Subtype /Link /Border [0 0 0] /H /I /C [1 0 0] /Rect [265.361 443.939 297.373 455.892] /A << /S /GoTo /D <6571756174696F6E2E452E332E312E31> >> >>}
\MusicPDFWriteObject{8366}{<< /Type /Annot /Subtype /Link /Border [0 0 0] /H /I /C [1 0 0] /Rect [371.649 441.831 403.271 455.894] /A << /S /GoTo /D <74686D2E452E332E312E31> >> >>}
\MusicPDFWriteObject{8367}{<< /Type /Annot /Subtype /Link /Border [0 0 0] /H /I /C [1 0 0] /Rect [204.036 417.34 235.447 429.61] /A << /S /GoTo /D <74686D2E452E332E312E36> >> >>}
\MusicPDFWriteObject{8368}{<< /Type /Annot /Subtype /Link /Border [0 0 0] /H /I /C [1 0 0] /Rect [259.381 405.04 291.393 415.831] /A << /S /GoTo /D <74686D2E452E332E322E33> >> >>}
\MusicPDFWriteObject{8369}{<< /Type /Annot /Subtype /Link /Border [0 0 0] /H /I /C [1 0 0] /Rect [194.655 382.004 219.276 393.702] /A << /S /GoTo /D <73756273656374696F6E2E452E332E31> >> >>}
\MusicPDFWriteObject{8370}{<< /Type /Annot /Subtype /Link /Border [0 0 0] /H /I /C [1 0 0] /Rect [135.883 317.101 167.895 327.892] /A << /S /GoTo /D <74686D2E452E332E312E32> >> >>}
\MusicPDFWriteObject{8371}{<< /Type /Annot /Subtype /Link /Border [0 0 0] /H /I /C [1 0 0] /Rect [207.542 317.101 239.553 327.892] /A << /S /GoTo /D <74686D2E452E332E312E34> >> >>}
\MusicPDFWriteObject{8372}{<< /Type /Annot /Subtype /Link /Border [0 0 0] /H /I /C [1 0 0] /Rect [298.567 317.101 330.578 327.892] /A << /S /GoTo /D <74686D2E452E332E312E36> >> >>}
\MusicPDFWriteObject{8373}{<< /Type /Annot /Subtype /Link /Border [0 0 0] /H /I /C [1 0 0] /Rect [71.004 290.245 103.015 302.201] /A << /S /GoTo /D <74686D2E452E332E352E31> >> >>}
\MusicPDFWriteObject{8374}{<< /Type /Annot /Subtype /Link /Border [0 0 0] /H /I /C [1 0 0] /Rect [372.501 290.245 404.512 302.201] /A << /S /GoTo /D <74686D2E452E332E352E31> >> >>}
\MusicPDFWriteObject{8375}{<< /Type /Annot /Subtype /Link /Border [0 0 0] /H /I /C [1 0 0] /Rect [431.455 237.413 463.016 249.665] /A << /S /GoTo /D <74686D2E452E332E352E31> >> >>}
\MusicPDFWriteObject{8376}{<< /Type /Annot /Subtype /Link /Border [0 0 0] /H /I /C [0 1 0] /Rect [169.531 461.356 185.717 470.223] /A << /S /GoTo /D <636974652E5350> >> >>}
\MusicPDFWriteObject{8377}{<< /Type /Annot /Subtype /Link /Border [0 0 0] /H /I /C [0 1 0] /Rect [378.945 461.356 395.13 470.223] /A << /S /GoTo /D <636974652E5350> >> >>}
\MusicPDFWriteObject{8378}{<< /Type /Annot /Subtype /Link /Border [0 0 0] /H /I /C [0 1 0] /Rect [474.756 461.356 492.09 470.223] /A << /S /GoTo /D <636974652E4B53> >> >>}
\MusicPDFWriteObject{8379}{<< /Type /Annot /Subtype /Link /Border [0 0 0] /H /I /C [0 1 0] /Rect [222.799 443.718 238.985 452.586] /A << /S /GoTo /D <636974652E5350> >> >>}
\MusicPDFWriteObject{8380}{<< /Type /Annot /Subtype /Link /Border [0 0 0] /H /I /C [0 1 0] /Rect [527.045 326.294 543.231 335.161] /A << /S /GoTo /D <636974652E5350> >> >>}
\MusicPDFWriteObject{8381}{<< /Type /Annot /Subtype /Link /Border [0 0 0] /H /I /C [1 0 0] /Rect [294.121 266.714 326.132 278.666] /A << /S /GoTo /D <74686D2E452E332E352E34> >> >>}
\MusicPDFWriteObject{8382}{<< /Type /Annot /Subtype /Link /Border [0 0 0] /H /I /C [0 1 0] /Rect [427.726 269.204 443.912 278.072] /A << /S /GoTo /D <636974652E5350> >> >>}
\MusicPDFWriteObject{8383}{<< /Type /Annot /Subtype /Link /Border [0 0 0] /H /I /C [1 0 0] /Rect [233.222 177.297 265.234 188.087] /A << /S /GoTo /D <74686D2E452E332E352E35> >> >>}
\MusicPDFWriteObject{8384}{<< /Type /Annot /Subtype /Link /Border [0 0 0] /H /I /C [0 1 0] /Rect [385.91 163.99 402.095 172.858] /A << /S /GoTo /D <636974652E5350> >> >>}
\MusicPDFWriteObject{8385}{<< /Type /Annot /Subtype /Link /Border [0 0 0] /H /I /C [0 1 0] /Rect [91.423 150.84 108.757 159.708] /A << /S /GoTo /D <636974652E4B53> >> >>}
\MusicPDFWriteObject{8386}{<< /Type /Annot /Subtype /Link /Border [0 0 0] /H /I /C [0 1 0] /Rect [331.715 150.84 349.049 159.708] /A << /S /GoTo /D <636974652E4B53> >> >>}
\MusicPDFWriteObject{8387}{<< /Type /Annot /Subtype /Link /Border [0 0 0] /H /I /C [0 1 0] /Rect [440.786 150.84 458.12 159.708] /A << /S /GoTo /D <636974652E4B53> >> >>}
\MusicPDFWriteObject{8388}{<< /Type /Annot /Subtype /Link /Border [0 0 0] /H /I /C [0 1 0] /Rect [111.328 137.689 127.514 146.557] /A << /S /GoTo /D <636974652E5350> >> >>}
\MusicPDFWriteObject{8389}{<< /Type /Annot /Subtype /Link /Border [0 0 0] /H /I /C [0 1 0] /Rect [182.328 111.387 199.662 120.255] /A << /S /GoTo /D <636974652E4B53> >> >>}
\MusicPDFWriteObject{8390}{<< /Type /Annot /Subtype /Link /Border [0 0 0] /H /I /C [0 1 0] /Rect [399.694 111.387 417.028 120.255] /A << /S /GoTo /D <636974652E4B53> >> >>}
\MusicPDFWriteObject{8391}{<< /Type /Annot /Subtype /Link /Border [0 0 0] /H /I /C [0 1 0] /Rect [160.03 98.237 176.216 107.105] /A << /S /GoTo /D <636974652E5350> >> >>}
\MusicPDFWriteObject{8392}{<< /Type /Annot /Subtype /Link /Border [0 0 0] /H /I /C [0 1 0] /Rect [179.405 85.086 196.739 93.954] /A << /S /GoTo /D <636974652E4B53> >> >>}
\MusicPDFWriteObject{8393}{<< /Type /Annot /Subtype /Link /Border [0 0 0] /H /I /C [0 1 0] /Rect [295.09 687.123 311.275 695.991] /A << /S /GoTo /D <636974652E5350> >> >>}
\MusicPDFWriteObject{8394}{<< /Type /Annot /Subtype /Link /Border [0 0 0] /H /I /C [0 1 0] /Rect [433.752 687.123 451.086 695.991] /A << /S /GoTo /D <636974652E4B53> >> >>}
\MusicPDFWriteObject{8395}{<< /Type /Annot /Subtype /Link /Border [0 0 0] /H /I /C [0 1 0] /Rect [338.549 673.973 355.883 682.84] /A << /S /GoTo /D <636974652E4B53> >> >>}
\MusicPDFWriteObject{8396}{<< /Type /Annot /Subtype /Link /Border [0 0 0] /H /I /C [0 1 0] /Rect [362.494 621.369 378.68 630.237] /A << /S /GoTo /D <636974652E5350> >> >>}
\MusicPDFWriteObject{8397}{<< /Type /Annot /Subtype /Link /Border [0 0 0] /H /I /C [0 1 0] /Rect [91.479 608.219 108.813 617.087] /A << /S /GoTo /D <636974652E4B53> >> >>}
\MusicPDFWriteObject{8398}{<< /Type /Annot /Subtype /Link /Border [0 0 0] /H /I /C [0 1 0] /Rect [218.902 608.219 236.236 617.087] /A << /S /GoTo /D <636974652E4B53> >> >>}
\MusicPDFWriteObject{8399}{<< /Type /Annot /Subtype /Link /Border [0 0 0] /H /I /C [0 1 0] /Rect [368.828 595.068 385.014 603.936] /A << /S /GoTo /D <636974652E5350> >> >>}
\MusicPDFWriteObject{8400}{<< /Type /Annot /Subtype /Link /Border [0 0 0] /H /I /C [0 1 0] /Rect [315.593 535.946 331.778 544.814] /A << /S /GoTo /D <636974652E5350> >> >>}
\MusicPDFWriteObject{8401}{<< /Type /Annot /Subtype /Link /Border [0 0 0] /H /I /C [0 1 0] /Rect [91.534 470.192 107.72 479.06] /A << /S /GoTo /D <636974652E5350> >> >>}
\MusicPDFWriteObject{8402}{<< /Type /Annot /Subtype /Link /Border [0 0 0] /H /I /C [0 1 0] /Rect [238.035 457.041 254.221 465.909] /A << /S /GoTo /D <636974652E5350> >> >>}
\MusicPDFWriteObject{8403}{<< /Type /Annot /Subtype /Link /Border [0 0 0] /H /I /C [0 1 0] /Rect [163.107 404.438 179.293 413.306] /A << /S /GoTo /D <636974652E5350> >> >>}
\MusicPDFWriteObject{8404}{<< /Type /Annot /Subtype /Link /Border [0 0 0] /H /I /C [1 0 0] /Rect [224.422 375.646 256.704 387.601] /A << /S /GoTo /D <74686D2E452E332E352E32> >> >>}
\MusicPDFWriteObject{8405}{<< /Type /Annot /Subtype /Link /Border [0 0 0] /H /I /C [1 0 0] /Rect [310.175 302.384 341.586 314.337] /A << /S /GoTo /D <74686D2E452E332E352E36> >> >>}
\MusicPDFWriteObject{8406}{<< /Type /Annot /Subtype /Link /Border [0 0 0] /H /I /C [0 1 0] /Rect [198.777 278.572 214.963 287.44] /A << /S /GoTo /D <636974652E5350> >> >>}
\MusicPDFWriteObject{8407}{<< /Type /Annot /Subtype /Link /Border [0 0 0] /H /I /C [0 1 0] /Rect [314.546 278.572 330.732 287.44] /A << /S /GoTo /D <636974652E5350> >> >>}
\MusicPDFWriteObject{8408}{<< /Type /Annot /Subtype /Link /Border [0 0 0] /H /I /C [1 0 0] /Rect [318.122 221.612 350.104 233.565] /A << /S /GoTo /D <74686D2E452E332E342E34> >> >>}
\MusicPDFWriteObject{8409}{<< /Type /Annot /Subtype /Link /Border [0 0 0] /H /I /C [0 1 0] /Rect [511.457 210.952 543.206 219.819] /A << /S /GoTo /D <636974652E53474133> >> >>}
\MusicPDFWriteObject{8410}{<< /Type /Annot /Subtype /Link /Border [0 0 0] /H /I /C [0 1 0] /Rect [293.789 177.141 312.438 186.009] /A << /S /GoTo /D <636974652E4C6573> >> >>}
\MusicPDFWriteObject{8411}{<< /Type /Annot /Subtype /Link /Border [0 0 0] /H /I /C [1 0 0] /Rect [384.582 162.054 417.164 172.844] /A << /S /GoTo /D <74686D2E452E332E352E37> >> >>}
\MusicPDFWriteObject{8412}{<< /Type /Annot /Subtype /Link /Border [0 0 0] /H /I /C [0 1 0] /Rect [408.369 150.84 425.703 159.708] /A << /S /GoTo /D <636974652E4B53> >> >>}
\MusicPDFWriteObject{8413}{<< /Type /Annot /Subtype /Link /Border [0 0 0] /H /I /C [0 1 0] /Rect [137.36 124.538 169.109 133.406] /A << /S /GoTo /D <636974652E53474134> >> >>}
\MusicPDFWriteObject{8414}{<< /Type /Annot /Subtype /Link /Border [0 0 0] /H /I /C [0 1 0] /Rect [422.485 124.538 454.234 133.406] /A << /S /GoTo /D <636974652E53474134> >> >>}
\MusicPDFWriteObject{8415}{<< /Type /Annot /Subtype /Link /Border [0 0 0] /H /I /C [1 0 0] /Rect [374.545 671.482 405.986 683.435] /A << /S /GoTo /D <74686D2E452E322E332E36> >> >>}
\MusicPDFWriteObject{8416}{<< /Type /Annot /Subtype /Link /Border [0 0 0] /H /I /C [1 0 0] /Rect [478.36 658.332 514.652 670.284] /A << /S /GoTo /D <74686D2E452E312E322E3136> >> >>}
\MusicPDFWriteObject{8417}{<< /Type /Annot /Subtype /Link /Border [0 0 0] /H /I /C [1 0 0] /Rect [424.595 605.92 456.576 616.71] /A << /S /GoTo /D <74686D2E452E332E352E31> >> >>}
\MusicPDFWriteObject{8418}{<< /Type /Annot /Subtype /Link /Border [0 0 0] /H /I /C [1 0 0] /Rect [314.065 592.769 346.076 603.617] /A << /S /GoTo /D <74686D2E452E332E352E34> >> >>}
\MusicPDFWriteObject{8419}{<< /Type /Annot /Subtype /Link /Border [0 0 0] /H /I /C [1 0 0] /Rect [87.051 434.863 103.569 445.654] /A << /S /GoTo /D <73656374696F6E2E452E35> >> >>}
\MusicPDFWriteObject{8420}{<< /Type /Annot /Subtype /Link /Border [0 0 0] /H /I /C [1 0 0] /Rect [109.323 434.863 125.841 445.654] /A << /S /GoTo /D <73656374696F6E2E452E36> >> >>}
\MusicPDFWriteObject{8421}{<< /Type /Annot /Subtype /Link /Border [0 0 0] /H /I /C [0 1 0] /Rect [434.932 402.008 454.798 412.812] /A << /S /GoTo /D <636974652E486967> >> >>}
\MusicPDFWriteObject{8422}{<< /Type /Annot /Subtype /Link /Border [0 0 0] /H /I /C [0 1 0] /Rect [465.289 403.944 500.122 412.812] /A << /S /GoTo /D <636974652E43645357> >> >>}
\MusicPDFWriteObject{8423}{<< /Type /Annot /Subtype /Link /Border [0 0 0] /H /I /C [0 1 0] /Rect [93.71 390.793 104.667 399.661] /A << /S /GoTo /D <636974652E4E> >> >>}
\MusicPDFWriteObject{8424}{<< /Type /Annot /Subtype /Link /Border [0 0 0] /H /I /C [0 1 0] /Rect [393.47 390.793 409.656 399.661] /A << /S /GoTo /D <636974652E5350> >> >>}
\MusicPDFWriteObject{8425}{<< /Type /Annot /Subtype /Link /Border [0 0 0] /H /I /C [0 1 0] /Rect [96.245 377.643 112.431 386.511] /A << /S /GoTo /D <636974652E5350> >> >>}
\MusicPDFWriteObject{8426}{<< /Type /Annot /Subtype /Link /Border [0 0 0] /H /I /C [1 0 0] /Rect [131.564 623.736 163.365 635.692] /A << /S /GoTo /D <6571756174696F6E2E452E342E312E37> >> >>}
\MusicPDFWriteObject{8427}{<< /Type /Annot /Subtype /Link /Border [0 0 0] /H /I /C [1 0 0] /Rect [233.222 589.696 265.234 600.486] /A << /S /GoTo /D <74686D2E452E342E312E31> >> >>}
\MusicPDFWriteObject{8428}{<< /Type /Annot /Subtype /Link /Border [0 0 0] /H /I /C [1 0 0] /Rect [156.541 559.382 194.232 571.335] /A << /S /GoTo /D <74686D2E452E352E312E3135> >> >>}
\MusicPDFWriteObject{8429}{<< /Type /Annot /Subtype /Link /Border [0 0 0] /H /I /C [1 0 0] /Rect [284.755 519.929 316.256 532.522] /A << /S /GoTo /D <74686D2E452E342E312E31> >> >>}
\MusicPDFWriteObject{8430}{<< /Type /Annot /Subtype /Link /Border [0 0 0] /H /I /C [1 0 0] /Rect [105.006 467.327 136.957 479.28] /A << /S /GoTo /D <6571756174696F6E2E452E342E312E38> >> >>}
\MusicPDFWriteObject{8431}{<< /Type /Annot /Subtype /Link /Border [0 0 0] /H /I /C [1 0 0] /Rect [182.647 454.175 214.058 466.131] /A << /S /GoTo /D <6571756174696F6E2E452E342E312E39> >> >>}
\MusicPDFWriteObject{8432}{<< /Type /Annot /Subtype /Link /Border [0 0 0] /H /I /C [1 0 0] /Rect [213.269 545.781 245.28 556.572] /A << /S /GoTo /D <74686D2E452E322E322E31> >> >>}
\MusicPDFWriteObject{8433}{<< /Type /Annot /Subtype /Link /Border [0 0 0] /H /I /C [1 0 0] /Rect [416.569 499.289 448.58 510.214] /A << /S /GoTo /D <74686D2E452E342E312E31> >> >>}
\MusicPDFWriteObject{8434}{<< /Type /Annot /Subtype /Link /Border [0 0 0] /H /I /C [1 0 0] /Rect [131.65 459.283 163.662 471.236] /A << /S /GoTo /D <6571756174696F6E2E452E342E312E38> >> >>}
\MusicPDFWriteObject{8435}{<< /Type /Annot /Subtype /Link /Border [0 0 0] /H /I /C [1 0 0] /Rect [174.396 459.283 211.388 471.236] /A << /S /GoTo /D <6571756174696F6E2E452E342E312E3131> >> >>}
\MusicPDFWriteObject{8436}{<< /Type /Annot /Subtype /Link /Border [0 0 0] /H /I /C [0 1 0] /Rect [281.107 428.432 312.856 437.3] /A << /S /GoTo /D <636974652E53474134> >> >>}
\MusicPDFWriteObject{8437}{<< /Type /Annot /Subtype /Link /Border [0 0 0] /H /I /C [0 1 0] /Rect [424.479 428.432 440.665 437.3] /A << /S /GoTo /D <636974652E5350> >> >>}
\MusicPDFWriteObject{8438}{<< /Type /Annot /Subtype /Link /Border [0 0 0] /H /I /C [0 1 0] /Rect [221.403 386.448 237.589 395.316] /A << /S /GoTo /D <636974652E5350> >> >>}
\MusicPDFWriteObject{8439}{<< /Type /Annot /Subtype /Link /Border [0 0 0] /H /I /C [1 0 0] /Rect [117.346 308.631 149.358 319.557] /A << /S /GoTo /D <74686D2E452E342E312E31> >> >>}
\MusicPDFWriteObject{8440}{<< /Type /Annot /Subtype /Link /Border [0 0 0] /H /I /C [1 0 0] /Rect [273.416 235.777 305.278 246.68] /A << /S /GoTo /D <74686D2E452E342E312E31> >> >>}
\MusicPDFWriteObject{8441}{<< /Type /Annot /Subtype /Link /Border [0 0 0] /H /I /C [1 0 0] /Rect [202.506 209.475 234.518 220.378] /A << /S /GoTo /D <74686D2E452E342E312E31> >> >>}
\MusicPDFWriteObject{8442}{<< /Type /Annot /Subtype /Link /Border [0 0 0] /H /I /C [1 0 0] /Rect [128.69 176.133 160.701 186.981] /A << /S /GoTo /D <74686D2E452E342E322E31> >> >>}
\MusicPDFWriteObject{8443}{<< /Type /Annot /Subtype /Link /Border [0 0 0] /H /I /C [1 0 0] /Rect [425.472 129.088 457.483 141.041] /A << /S /GoTo /D <74686D2E452E322E322E35> >> >>}
\MusicPDFWriteObject{8444}{<< /Type /Annot /Subtype /Link /Border [0 0 0] /H /I /C [1 0 0] /Rect [124.041 671.482 161.733 683.435] /A << /S /GoTo /D <74686D2E452E312E322E3235> >> >>}
\MusicPDFWriteObject{8445}{<< /Type /Annot /Subtype /Link /Border [0 0 0] /H /I /C [1 0 0] /Rect [492.499 658.332 529.49 670.284] /A << /S /GoTo /D <74686D2E452E312E322E3235> >> >>}
\MusicPDFWriteObject{8446}{<< /Type /Annot /Subtype /Link /Border [0 0 0] /H /I /C [1 0 0] /Rect [289.137 567.448 321.749 578.239] /A << /S /GoTo /D <74686D2E452E322E332E34> >> >>}
\MusicPDFWriteObject{8447}{<< /Type /Annot /Subtype /Link /Border [0 0 0] /H /I /C [1 0 0] /Rect [268.864 499.105 300.755 511.058] /A << /S /GoTo /D <74686D2E452E322E322E35> >> >>}
\MusicPDFWriteObject{8448}{<< /Type /Annot /Subtype /Link /Border [0 0 0] /H /I /C [1 0 0] /Rect [335.461 449.526 366.872 461.481] /A << /S /GoTo /D <74686D2E452E322E312E36> >> >>}
\MusicPDFWriteObject{8449}{<< /Type /Annot /Subtype /Link /Border [0 0 0] /H /I /C [1 0 0] /Rect [71.004 408.354 107.89 422.556] /A << /S /GoTo /D <74686D2E452E312E322E3235> >> >>}
\MusicPDFWriteObject{8450}{<< /Type /Annot /Subtype /Link /Border [0 0 0] /H /I /C [1 0 0] /Rect [265.741 305.746 302.208 317.699] /A << /S /GoTo /D <74686D2E452E312E322E3235> >> >>}
\MusicPDFWriteObject{8451}{<< /Type /Annot /Subtype /Link /Border [0 0 0] /H /I /C [1 0 0] /Rect [175.689 566.276 208.301 578.229] /A << /S /GoTo /D <7465726D696E6F6C6F67792E452E302E302E32> >> >>}
\MusicPDFWriteObject{8452}{<< /Type /Annot /Subtype /Link /Border [0 0 0] /H /I /C [1 0 0] /Rect [260.385 684.633 292.396 696.586] /A << /S /GoTo /D <6571756174696F6E2E452E342E322E34> >> >>}
\MusicPDFWriteObject{8453}{<< /Type /Annot /Subtype /Link /Border [0 0 0] /H /I /C [1 0 0] /Rect [457.226 589.055 489.238 599.845] /A << /S /GoTo /D <74686D2E452E342E322E33> >> >>}
\MusicPDFWriteObject{8454}{<< /Type /Annot /Subtype /Link /Border [0 0 0] /H /I /C [0 1 0] /Rect [420.028 492.921 435.419 501.789] /A << /S /GoTo /D <636974652E4265> >> >>}
\MusicPDFWriteObject{8455}{<< /Type /Annot /Subtype /Link /Border [0 0 0] /H /I /C [0 1 0] /Rect [126.064 479.771 149.215 488.638] /A << /S /GoTo /D <636974652E4D6172> >> >>}
\MusicPDFWriteObject{8456}{<< /Type /Annot /Subtype /Link /Border [0 0 0] /H /I /C [0 1 0] /Rect [124.497 466.62 150.823 475.488] /A << /S /GoTo /D <636974652E45687231> >> >>}
\MusicPDFWriteObject{8457}{<< /Type /Annot /Subtype /Link /Border [0 0 0] /H /I /C [0 1 0] /Rect [177.427 466.62 203.753 475.488] /A << /S /GoTo /D <636974652E45687232> >> >>}
\MusicPDFWriteObject{8458}{<< /Type /Annot /Subtype /Link /Border [0 0 0] /H /I /C [0 1 0] /Rect [343.95 466.62 370.262 475.488] /A << /S /GoTo /D <636974652E426F7231> >> >>}
\MusicPDFWriteObject{8459}{<< /Type /Annot /Subtype /Link /Border [0 0 0] /H /I /C [0 1 0] /Rect [408.658 466.62 428.247 475.488] /A << /S /GoTo /D <636974652E4752> >> >>}
\MusicPDFWriteObject{8460}{<< /Type /Annot /Subtype /Link /Border [0 0 0] /H /I /C [0 1 0] /Rect [100.582 453.469 117.107 462.337] /A << /S /GoTo /D <636974652E4861> >> >>}
\MusicPDFWriteObject{8461}{<< /Type /Annot /Subtype /Link /Border [0 0 0] /H /I /C [0 1 0] /Rect [195.646 453.469 211.832 462.337] /A << /S /GoTo /D <636974652E5350> >> >>}
\MusicPDFWriteObject{8462}{<< /Type /Annot /Subtype /Link /Border [0 0 0] /H /I /C [0 1 0] /Rect [283.658 453.469 299.843 462.337] /A << /S /GoTo /D <636974652E5350> >> >>}
\MusicPDFWriteObject{8463}{<< /Type /Annot /Subtype /Link /Border [0 0 0] /H /I /C [1 0 0] /Rect [139.804 398.375 163.69 410.33] /A << /S /GoTo /D <73756273656374696F6E2E452E332E35> >> >>}
\MusicPDFWriteObject{8464}{<< /Type /Annot /Subtype /Link /Border [0 0 0] /H /I /C [0 1 0] /Rect [90.398 224.5 106.584 233.368] /A << /S /GoTo /D <636974652E5350> >> >>}
\MusicPDFWriteObject{8465}{<< /Type /Annot /Subtype /Link /Border [0 0 0] /H /I /C [0 1 0] /Rect [181.092 224.5 207.404 233.368] /A << /S /GoTo /D <636974652E426F7231> >> >>}
\MusicPDFWriteObject{8466}{<< /Type /Annot /Subtype /Link /Border [0 0 0] /H /I /C [1 0 0] /Rect [144.397 572.326 176.408 584.279] /A << /S /GoTo /D <6571756174696F6E2E452E352E312E32> >> >>}
\MusicPDFWriteObject{8467}{<< /Type /Annot /Subtype /Link /Border [0 0 0] /H /I /C [1 0 0] /Rect [288.5 321.383 320.512 333.336] /A << /S /GoTo /D <6571756174696F6E2E452E352E312E32> >> >>}
\MusicPDFWriteObject{8468}{<< /Type /Annot /Subtype /Link /Border [0 0 0] /H /I /C [1 0 0] /Rect [381.789 321.383 414.401 333.336] /A << /S /GoTo /D <74686D2E452E352E312E32> >> >>}
\MusicPDFWriteObject{8469}{<< /Type /Annot /Subtype /Link /Border [0 0 0] /H /I /C [1 0 0] /Rect [265.504 205.573 297.515 216.363] /A << /S /GoTo /D <74686D2E452E352E312E33> >> >>}
\MusicPDFWriteObject{8470}{<< /Type /Annot /Subtype /Link /Border [0 0 0] /H /I /C [1 0 0] /Rect [444.747 647.132 476.669 658.551] /A << /S /GoTo /D <74686D2E452E352E312E33> >> >>}
\MusicPDFWriteObject{8471}{<< /Type /Annot /Subtype /Link /Border [0 0 0] /H /I /C [1 0 0] /Rect [209.426 448.319 241.438 460.271] /A << /S /GoTo /D <6571756174696F6E2E452E352E312E38> >> >>}
\MusicPDFWriteObject{8472}{<< /Type /Annot /Subtype /Link /Border [0 0 0] /H /I /C [1 0 0] /Rect [164.499 202.147 195.91 212.938] /A << /S /GoTo /D <74686D2E452E352E312E33> >> >>}
\MusicPDFWriteObject{8473}{<< /Type /Annot /Subtype /Link /Border [0 0 0] /H /I /C [1 0 0] /Rect [337.679 299.514 370.291 310.305] /A << /S /GoTo /D <74686D2E452E352E312E36> >> >>}
\MusicPDFWriteObject{8474}{<< /Type /Annot /Subtype /Link /Border [0 0 0] /H /I /C [1 0 0] /Rect [429.807 299.514 462.419 310.305] /A << /S /GoTo /D <74686D2E452E352E312E38> >> >>}
\MusicPDFWriteObject{8475}{<< /Type /Annot /Subtype /Link /Border [0 0 0] /H /I /C [1 0 0] /Rect [301.243 172.5 333.855 183.291] /A << /S /GoTo /D <74686D2E452E352E312E36> >> >>}
\MusicPDFWriteObject{8476}{<< /Type /Annot /Subtype /Link /Border [0 0 0] /H /I /C [1 0 0] /Rect [393.432 172.5 426.044 183.291] /A << /S /GoTo /D <74686D2E452E352E312E38> >> >>}
\MusicPDFWriteObject{8477}{<< /Type /Annot /Subtype /Link /Border [0 0 0] /H /I /C [0 1 0] /Rect [116.433 134.985 132.619 143.853] /A << /S /GoTo /D <636974652E5350> >> >>}
\MusicPDFWriteObject{8478}{<< /Type /Annot /Subtype /Link /Border [0 0 0] /H /I /C [1 0 0] /Rect [177.449 685.187 215.141 695.977] /A << /S /GoTo /D <74686D2E452E352E312E3130> >> >>}
\MusicPDFWriteObject{8479}{<< /Type /Annot /Subtype /Link /Border [0 0 0] /H /I /C [1 0 0] /Rect [349.72 671.482 381.732 683.435] /A << /S /GoTo /D <6571756174696F6E2E452E352E312E35> >> >>}
\MusicPDFWriteObject{8480}{<< /Type /Annot /Subtype /Link /Border [0 0 0] /H /I /C [1 0 0] /Rect [298.179 377.956 329.59 389.915] /A << /S /GoTo /D <74686D2E452E352E312E32> >> >>}
\MusicPDFWriteObject{8481}{<< /Type /Annot /Subtype /Link /Border [0 0 0] /H /I /C [1 0 0] /Rect [481.473 377.956 517.765 389.915] /A << /S /GoTo /D <6571756174696F6E2E452E352E312E3131> >> >>}
\MusicPDFWriteObject{8482}{<< /Type /Annot /Subtype /Link /Border [0 0 0] /H /I /C [1 0 0] /Rect [106.556 364.812 138.568 376.764] /A << /S /GoTo /D <74686D2E452E352E312E36> >> >>}
\MusicPDFWriteObject{8483}{<< /Type /Annot /Subtype /Link /Border [0 0 0] /H /I /C [1 0 0] /Rect [163.136 364.812 200.127 376.764] /A << /S /GoTo /D <6571756174696F6E2E452E352E312E3135> >> >>}
\MusicPDFWriteObject{8484}{<< /Type /Annot /Subtype /Link /Border [0 0 0] /H /I /C [1 0 0] /Rect [264.425 364.812 296.437 376.764] /A << /S /GoTo /D <74686D2E452E352E312E38> >> >>}
\MusicPDFWriteObject{8485}{<< /Type /Annot /Subtype /Link /Border [0 0 0] /H /I /C [1 0 0] /Rect [139.757 175.517 171.768 186.308] /A << /S /GoTo /D <74686D2E452E352E312E32> >> >>}
\MusicPDFWriteObject{8486}{<< /Type /Annot /Subtype /Link /Border [0 0 0] /H /I /C [1 0 0] /Rect [270.597 684.633 302.609 696.586] /A << /S /GoTo /D <6571756174696F6E2E452E352E312E31> >> >>}
\MusicPDFWriteObject{8487}{<< /Type /Annot /Subtype /Link /Border [0 0 0] /H /I /C [1 0 0] /Rect [170.457 671.482 207.448 683.435] /A << /S /GoTo /D <6571756174696F6E2E452E352E312E3230> >> >>}
\MusicPDFWriteObject{8488}{<< /Type /Annot /Subtype /Link /Border [0 0 0] /H /I /C [0 1 0] /Rect [73.77 535.673 89.956 544.54] /A << /S /GoTo /D <636974652E5350> >> >>}
\MusicPDFWriteObject{8489}{<< /Type /Annot /Subtype /Link /Border [0 0 0] /H /I /C [0 1 0] /Rect [285.397 535.673 301.583 544.54] /A << /S /GoTo /D <636974652E5350> >> >>}
\MusicPDFWriteObject{8490}{<< /Type /Annot /Subtype /Link /Border [0 0 0] /H /I /C [1 0 0] /Rect [310.054 514.607 347.745 525.398] /A << /S /GoTo /D <74686D2E452E352E312E3132> >> >>}
\MusicPDFWriteObject{8491}{<< /Type /Annot /Subtype /Link /Border [0 0 0] /H /I /C [1 0 0] /Rect [506.015 514.607 543.707 525.398] /A << /S /GoTo /D <74686D2E452E352E312E3133> >> >>}
\MusicPDFWriteObject{8492}{<< /Type /Annot /Subtype /Link /Border [0 0 0] /H /I /C [1 0 0] /Rect [336.968 300.058 373.26 312.031] /A << /S /GoTo /D <74686D2E452E352E312E3133> >> >>}
\MusicPDFWriteObject{8493}{<< /Type /Annot /Subtype /Link /Border [0 0 0] /H /I /C [1 0 0] /Rect [246.526 286.908 282.817 298.861] /A << /S /GoTo /D <6571756174696F6E2E452E352E312E3131> >> >>}
\MusicPDFWriteObject{8494}{<< /Type /Annot /Subtype /Link /Border [0 0 0] /H /I /C [1 0 0] /Rect [345.87 286.908 377.281 298.861] /A << /S /GoTo /D <74686D2E452E352E312E36> >> >>}
\MusicPDFWriteObject{8495}{<< /Type /Annot /Subtype /Link /Border [0 0 0] /H /I /C [1 0 0] /Rect [412.514 286.908 448.806 298.861] /A << /S /GoTo /D <6571756174696F6E2E452E352E312E3135> >> >>}
\MusicPDFWriteObject{8496}{<< /Type /Annot /Subtype /Link /Border [0 0 0] /H /I /C [1 0 0] /Rect [511.858 286.908 543.27 298.861] /A << /S /GoTo /D <74686D2E452E352E312E38> >> >>}
\MusicPDFWriteObject{8497}{<< /Type /Annot /Subtype /Link /Border [0 0 0] /H /I /C [1 0 0] /Rect [220.694 273.757 257.651 285.71] /A << /S /GoTo /D <74686D2E452E352E312E3132> >> >>}
\MusicPDFWriteObject{8498}{<< /Type /Annot /Subtype /Link /Border [0 0 0] /H /I /C [1 0 0] /Rect [160.244 197.325 192.316 209.278] /A << /S /GoTo /D <74686D2E452E342E312E31> >> >>}
\MusicPDFWriteObject{8499}{<< /Type /Annot /Subtype /Link /Border [0 0 0] /H /I /C [1 0 0] /Rect [173.338 184.175 183.271 197.032] /A << /S /GoTo /D <48666F6F746E6F74652E313639> >> >>}
\MusicPDFWriteObject{8500}{<< /Type /Annot /Subtype /Link /Border [0 0 0] /H /I /C [1 0 0] /Rect [509.674 184.175 541.865 197.032] /A << /S /GoTo /D <6571756174696F6E2E452E342E312E37> >> >>}
\MusicPDFWriteObject{8501}{<< /Type /Annot /Subtype /Link /Border [0 0 0] /H /I /C [1 0 0] /Rect [160.741 140.575 192.752 151.423] /A << /S /GoTo /D <73756273756273656374696F6E2E452E352E312E31> >> >>}
\MusicPDFWriteObject{8502}{<< /Type /Annot /Subtype /Link /Border [0 0 0] /H /I /C [1 0 0] /Rect [149.936 224.703 186.927 236.656] /A << /S /GoTo /D <6571756174696F6E2E452E352E312E3331> >> >>}
\MusicPDFWriteObject{8503}{<< /Type /Annot /Subtype /Link /Border [0 0 0] /H /I /C [1 0 0] /Rect [88.711 155.122 125.702 167.075] /A << /S /GoTo /D <6571756174696F6E2E452E352E312E3234> >> >>}
\MusicPDFWriteObject{8504}{<< /Type /Annot /Subtype /Link /Border [0 0 0] /H /I /C [1 0 0] /Rect [156.523 155.122 193.515 167.075] /A << /S /GoTo /D <6571756174696F6E2E452E352E312E3332> >> >>}
\MusicPDFWriteObject{8505}{<< /Type /Annot /Subtype /Link /Border [0 0 0] /H /I /C [1 0 0] /Rect [191.688 684.633 228.68 696.586] /A << /S /GoTo /D <6571756174696F6E2E452E352E312E3332> >> >>}
\MusicPDFWriteObject{8506}{<< /Type /Annot /Subtype /Link /Border [0 0 0] /H /I /C [1 0 0] /Rect [415.655 684.633 452.646 696.586] /A << /S /GoTo /D <6571756174696F6E2E452E352E312E3333> >> >>}
\MusicPDFWriteObject{8507}{<< /Type /Annot /Subtype /Link /Border [0 0 0] /H /I /C [1 0 0] /Rect [168.338 514.484 205.33 526.437] /A << /S /GoTo /D <6571756174696F6E2E452E352E312E3332> >> >>}
\MusicPDFWriteObject{8508}{<< /Type /Annot /Subtype /Link /Border [0 0 0] /H /I /C [1 0 0] /Rect [403.897 514.484 440.889 526.437] /A << /S /GoTo /D <6571756174696F6E2E452E352E312E3333> >> >>}
\MusicPDFWriteObject{8509}{<< /Type /Annot /Subtype /Link /Border [0 0 0] /H /I /C [1 0 0] /Rect [376.183 478.056 413.175 490.011] /A << /S /GoTo /D <6571756174696F6E2E452E352E312E3334> >> >>}
\MusicPDFWriteObject{8510}{<< /Type /Annot /Subtype /Link /Border [0 0 0] /H /I /C [1 0 0] /Rect [92.75 358.33 129.742 370.282] /A << /S /GoTo /D <6571756174696F6E2E452E352E312E3335> >> >>}
\MusicPDFWriteObject{8511}{<< /Type /Annot /Subtype /Link /Border [0 0 0] /H /I /C [1 0 0] /Rect [225.534 221.61 261.825 233.583] /A << /S /GoTo /D <6571756174696F6E2E452E352E312E3234> >> >>}
\MusicPDFWriteObject{8512}{<< /Type /Annot /Subtype /Link /Border [0 0 0] /H /I /C [1 0 0] /Rect [92.75 208.46 129.742 220.413] /A << /S /GoTo /D <6571756174696F6E2E452E352E312E3336> >> >>}
\MusicPDFWriteObject{8513}{<< /Type /Annot /Subtype /Link /Border [0 0 0] /H /I /C [1 0 0] /Rect [372.604 208.46 409.596 220.413] /A << /S /GoTo /D <74686D2E452E352E312E3133> >> >>}
\MusicPDFWriteObject{8514}{<< /Type /Annot /Subtype /Link /Border [0 0 0] /H /I /C [1 0 0] /Rect [266.699 180.055 303.69 190.846] /A << /S /GoTo /D <74686D2E452E352E312E3134> >> >>}
\MusicPDFWriteObject{8515}{<< /Type /Annot /Subtype /Link /Border [0 0 0] /H /I /C [1 0 0] /Rect [99.186 606.222 136.178 618.175] /A << /S /GoTo /D <6571756174696F6E2E452E352E312E3330> >> >>}
\MusicPDFWriteObject{8516}{<< /Type /Annot /Subtype /Link /Border [0 0 0] /H /I /C [1 0 0] /Rect [195.559 504.796 233.251 515.586] /A << /S /GoTo /D <74686D2E452E352E312E3138> >> >>}
\MusicPDFWriteObject{8517}{<< /Type /Annot /Subtype /Link /Border [0 0 0] /H /I /C [1 0 0] /Rect [396.833 504.796 429.444 515.586] /A << /S /GoTo /D <74686D2E452E352E312E32> >> >>}
\MusicPDFWriteObject{8518}{<< /Type /Annot /Subtype /Link /Border [0 0 0] /H /I /C [1 0 0] /Rect [471.832 504.796 504.444 515.586] /A << /S /GoTo /D <74686D2E452E352E312E36> >> >>}
\MusicPDFWriteObject{8519}{<< /Type /Annot /Subtype /Link /Border [0 0 0] /H /I /C [1 0 0] /Rect [71.004 491.645 103.616 502.436] /A << /S /GoTo /D <74686D2E452E352E312E38> >> >>}
\MusicPDFWriteObject{8520}{<< /Type /Annot /Subtype /Link /Border [0 0 0] /H /I /C [1 0 0] /Rect [144.738 491.645 177.35 502.436] /A << /S /GoTo /D <74686D2E452E352E312E39> >> >>}
\MusicPDFWriteObject{8521}{<< /Type /Annot /Subtype /Link /Border [0 0 0] /H /I /C [1 0 0] /Rect [218.473 491.645 256.165 502.436] /A << /S /GoTo /D <74686D2E452E352E312E3130> >> >>}
\MusicPDFWriteObject{8522}{<< /Type /Annot /Subtype /Link /Border [0 0 0] /H /I /C [1 0 0] /Rect [297.287 491.645 334.979 502.436] /A << /S /GoTo /D <74686D2E452E352E312E3132> >> >>}
\MusicPDFWriteObject{8523}{<< /Type /Annot /Subtype /Link /Border [0 0 0] /H /I /C [1 0 0] /Rect [393.226 491.645 430.918 502.436] /A << /S /GoTo /D <74686D2E452E352E312E3133> >> >>}
\MusicPDFWriteObject{8524}{<< /Type /Annot /Subtype /Link /Border [0 0 0] /H /I /C [1 0 0] /Rect [310.132 478.494 347.123 489.285] /A << /S /GoTo /D <74686D2E452E352E312E3134> >> >>}
\MusicPDFWriteObject{8525}{<< /Type /Annot /Subtype /Link /Border [0 0 0] /H /I /C [0 1 0] /Rect [460.895 252.819 487.207 261.687] /A << /S /GoTo /D <636974652E426F7231> >> >>}
\MusicPDFWriteObject{8526}{<< /Type /Annot /Subtype /Link /Border [0 0 0] /H /I /C [0 1 0] /Rect [523.606 252.819 543.195 261.687] /A << /S /GoTo /D <636974652E4752> >> >>}
\MusicPDFWriteObject{8527}{<< /Type /Annot /Subtype /Link /Border [0 0 0] /H /I /C [0 1 0] /Rect [103.089 239.668 119.614 248.536] /A << /S /GoTo /D <636974652E4861> >> >>}
\MusicPDFWriteObject{8528}{<< /Type /Annot /Subtype /Link /Border [0 0 0] /H /I /C [0 1 0] /Rect [185.933 239.668 202.119 248.536] /A << /S /GoTo /D <636974652E5350> >> >>}
\MusicPDFWriteObject{8529}{<< /Type /Annot /Subtype /Link /Border [0 0 0] /H /I /C [1 0 0] /Rect [401.263 171.425 433.185 183.377] /A << /S /GoTo /D <74686D2E452E342E312E31> >> >>}
\MusicPDFWriteObject{8530}{<< /Type /Annot /Subtype /Link /Border [0 0 0] /H /I /C [0 1 0] /Rect [506.062 173.915 522.248 182.783] /A << /S /GoTo /D <636974652E5350> >> >>}
\MusicPDFWriteObject{8531}{<< /Type /Annot /Subtype /Link /Border [0 0 0] /H /I /C [1 0 0] /Rect [301.857 542.999 311.79 555.857] /A << /S /GoTo /D <48666F6F746E6F74652E313730> >> >>}
\MusicPDFWriteObject{8532}{<< /Type /Annot /Subtype /Link /Border [0 0 0] /H /I /C [1 0 0] /Rect [345.039 506.572 382.031 518.525] /A << /S /GoTo /D <74686D2E452E352E312E3230> >> >>}
\MusicPDFWriteObject{8533}{<< /Type /Annot /Subtype /Link /Border [0 0 0] /H /I /C [1 0 0] /Rect [507.385 384.148 539.396 396.101] /A << /S /GoTo /D <6571756174696F6E2E452E352E322E31> >> >>}
\MusicPDFWriteObject{8534}{<< /Type /Annot /Subtype /Link /Border [0 0 0] /H /I /C [1 0 0] /Rect [186.463 270.591 218.474 282.543] /A << /S /GoTo /D <6571756174696F6E2E452E352E322E32> >> >>}
\MusicPDFWriteObject{8535}{<< /Type /Annot /Subtype /Link /Border [0 0 0] /H /I /C [1 0 0] /Rect [406.696 220.096 439.008 232.051] /A << /S /GoTo /D <6571756174696F6E2E452E352E322E31> >> >>}
\MusicPDFWriteObject{8536}{<< /Type /Annot /Subtype /Link /Border [0 0 0] /H /I /C [0 1 0] /Rect [367.344 169.983 386.933 178.851] /A << /S /GoTo /D <636974652E4752> >> >>}
\MusicPDFWriteObject{8537}{<< /Type /Annot /Subtype /Link /Border [0 0 0] /H /I /C [1 0 0] /Rect [465.127 154.343 497.738 166.295] /A << /S /GoTo /D <74686D2E452E352E322E31> >> >>}
\MusicPDFWriteObject{8538}{<< /Type /Annot /Subtype /Link /Border [0 0 0] /H /I /C [1 0 0] /Rect [158.351 603.197 190.633 615.169] /A << /S /GoTo /D <74686D2E452E352E322E31> >> >>}
\MusicPDFWriteObject{8539}{<< /Type /Annot /Subtype /Link /Border [0 0 0] /H /I /C [1 0 0] /Rect [156.823 517.191 188.234 529.144] /A << /S /GoTo /D <74686D2E452E352E322E31> >> >>}
\MusicPDFWriteObject{8540}{<< /Type /Annot /Subtype /Link /Border [0 0 0] /H /I /C [0 1 0] /Rect [229.021 519.681 255.333 528.549] /A << /S /GoTo /D <636974652E426F7231> >> >>}
\MusicPDFWriteObject{8541}{<< /Type /Annot /Subtype /Link /Border [0 0 0] /H /I /C [0 1 0] /Rect [496.152 519.681 515.74 528.549] /A << /S /GoTo /D <636974652E4752> >> >>}
\MusicPDFWriteObject{8542}{<< /Type /Annot /Subtype /Link /Border [0 0 0] /H /I /C [0 1 0] /Rect [286.993 467.078 306.582 475.946] /A << /S /GoTo /D <636974652E4752> >> >>}
\MusicPDFWriteObject{8543}{<< /Type /Annot /Subtype /Link /Border [0 0 0] /H /I /C [1 0 0] /Rect [514.121 425.136 545.592 437.089] /A << /S /GoTo /D <74686D2E452E352E322E31> >> >>}
\MusicPDFWriteObject{8544}{<< /Type /Annot /Subtype /Link /Border [0 0 0] /H /I /C [0 1 0] /Rect [275.582 414.475 292.107 423.343] /A << /S /GoTo /D <636974652E4861> >> >>}
\MusicPDFWriteObject{8545}{<< /Type /Annot /Subtype /Link /Border [0 0 0] /H /I /C [0 1 0] /Rect [443.699 388.174 470.011 397.042] /A << /S /GoTo /D <636974652E426F7231> >> >>}
\MusicPDFWriteObject{8546}{<< /Type /Annot /Subtype /Link /Border [0 0 0] /H /I /C [0 1 0] /Rect [523.578 388.174 543.167 397.042] /A << /S /GoTo /D <636974652E4752> >> >>}
\MusicPDFWriteObject{8547}{<< /Type /Annot /Subtype /Link /Border [0 0 0] /H /I /C [0 1 0] /Rect [241.933 361.872 268.245 370.74] /A << /S /GoTo /D <636974652E426F7231> >> >>}
\MusicPDFWriteObject{8548}{<< /Type /Annot /Subtype /Link /Border [0 0 0] /H /I /C [0 1 0] /Rect [324.876 361.872 344.465 370.74] /A << /S /GoTo /D <636974652E4752> >> >>}
\MusicPDFWriteObject{8549}{<< /Type /Annot /Subtype /Link /Border [0 0 0] /H /I /C [0 1 0] /Rect [147.8 348.721 164.325 357.589] /A << /S /GoTo /D <636974652E4861> >> >>}
\MusicPDFWriteObject{8550}{<< /Type /Annot /Subtype /Link /Border [0 0 0] /H /I /C [0 1 0] /Rect [445.59 348.721 460.981 357.589] /A << /S /GoTo /D <636974652E4265> >> >>}
\MusicPDFWriteObject{8551}{<< /Type /Annot /Subtype /Link /Border [0 0 0] /H /I /C [0 1 0] /Rect [193.022 335.571 209.208 344.439] /A << /S /GoTo /D <636974652E5350> >> >>}
\MusicPDFWriteObject{8552}{<< /Type /Annot /Subtype /Link /Border [0 0 0] /H /I /C [0 1 0] /Rect [347.682 335.571 363.868 344.439] /A << /S /GoTo /D <636974652E5350> >> >>}
\MusicPDFWriteObject{8553}{<< /Type /Annot /Subtype /Link /Border [0 0 0] /H /I /C [0 1 0] /Rect [439.359 286.094 465.671 294.962] /A << /S /GoTo /D <636974652E426F7231> >> >>}
\MusicPDFWriteObject{8554}{<< /Type /Annot /Subtype /Link /Border [0 0 0] /H /I /C [0 1 0] /Rect [318.906 272.944 335.431 281.811] /A << /S /GoTo /D <636974652E4861> >> >>}
\MusicPDFWriteObject{8555}{<< /Type /Annot /Subtype /Link /Border [0 0 0] /H /I /C [0 1 0] /Rect [402.669 272.944 418.855 281.811] /A << /S /GoTo /D <636974652E5350> >> >>}
\MusicPDFWriteObject{8556}{<< /Type /Annot /Subtype /Link /Border [0 0 0] /H /I /C [0 1 0] /Rect [496.433 272.944 516.022 281.811] /A << /S /GoTo /D <636974652E4752> >> >>}
\MusicPDFWriteObject{8557}{<< /Type /Annot /Subtype /Link /Border [0 0 0] /H /I /C [1 0 0] /Rect [148.624 424.026 180.996 435.979] /A << /S /GoTo /D <6571756174696F6E2E452E352E322E33> >> >>}
\MusicPDFWriteObject{8558}{<< /Type /Annot /Subtype /Link /Border [0 0 0] /H /I /C [1 0 0] /Rect [201.595 410.876 233.607 422.828] /A << /S /GoTo /D <6571756174696F6E2E452E352E322E33> >> >>}
\MusicPDFWriteObject{8559}{<< /Type /Annot /Subtype /Link /Border [0 0 0] /H /I /C [1 0 0] /Rect [152.122 397.724 184.734 409.68] /A << /S /GoTo /D <6571756174696F6E2E452E352E322E33> >> >>}
\MusicPDFWriteObject{8560}{<< /Type /Annot /Subtype /Link /Border [0 0 0] /H /I /C [1 0 0] /Rect [125.858 371.423 157.329 383.376] /A << /S /GoTo /D <6571756174696F6E2E452E352E322E33> >> >>}
\MusicPDFWriteObject{8561}{<< /Type /Annot /Subtype /Link /Border [0 0 0] /H /I /C [0 1 0] /Rect [232.991 261.987 252.58 270.855] /A << /S /GoTo /D <636974652E4752> >> >>}
\MusicPDFWriteObject{8562}{<< /Type /Annot /Subtype /Link /Border [0 0 0] /H /I /C [1 0 0] /Rect [444.227 259.497 476.028 271.45] /A << /S /GoTo /D <74686D2E452E352E322E33> >> >>}
\MusicPDFWriteObject{8563}{<< /Type /Annot /Subtype /Link /Border [0 0 0] /H /I /C [1 0 0] /Rect [487.862 227.254 519.453 239.207] /A << /S /GoTo /D <74686D2E452E352E322E33> >> >>}
\MusicPDFWriteObject{8564}{<< /Type /Annot /Subtype /Link /Border [0 0 0] /H /I /C [1 0 0] /Rect [309.487 214.103 341.348 226.056] /A << /S /GoTo /D <74686D2E452E352E322E33> >> >>}
\MusicPDFWriteObject{8565}{<< /Type /Annot /Subtype /Link /Border [0 0 0] /H /I /C [0 1 0] /Rect [523.639 216.594 543.228 225.461] /A << /S /GoTo /D <636974652E4752> >> >>}
\MusicPDFWriteObject{8566}{<< /Type /Annot /Subtype /Link /Border [0 0 0] /H /I /C [0 1 0] /Rect [289.084 177.141 308.673 186.009] /A << /S /GoTo /D <636974652E4752> >> >>}
\MusicPDFWriteObject{8567}{<< /Type /Annot /Subtype /Link /Border [0 0 0] /H /I /C [1 0 0] /Rect [157.243 135.752 188.744 146.543] /A << /S /GoTo /D <74686D2E452E352E322E33> >> >>}
\MusicPDFWriteObject{8568}{<< /Type /Annot /Subtype /Link /Border [0 0 0] /H /I /C [0 1 0] /Rect [176.698 124.538 196.287 133.406] /A << /S /GoTo /D <636974652E4752> >> >>}
\MusicPDFWriteObject{8569}{<< /Type /Annot /Subtype /Link /Border [0 0 0] /H /I /C [0 1 0] /Rect [413.205 85.086 432.794 93.954] /A << /S /GoTo /D <636974652E4752> >> >>}
\MusicPDFWriteObject{8570}{<< /Type /Annot /Subtype /Link /Border [0 0 0] /H /I /C [0 1 0] /Rect [199.883 660.822 219.472 669.69] /A << /S /GoTo /D <636974652E4752> >> >>}
\MusicPDFWriteObject{8571}{<< /Type /Annot /Subtype /Link /Border [0 0 0] /H /I /C [1 0 0] /Rect [117.346 562.428 149.358 571.282] /A << /S /GoTo /D <74686D2E452E352E322E31> >> >>}
\MusicPDFWriteObject{8572}{<< /Type /Annot /Subtype /Link /Border [0 0 0] /H /I /C [1 0 0] /Rect [466.631 684.633 503.763 696.588] /A << /S /GoTo /D <74686D2E452E352E312E3230> >> >>}
\MusicPDFWriteObject{8573}{<< /Type /Annot /Subtype /Link /Border [0 0 0] /H /I /C [1 0 0] /Rect [176.592 671.476 209.204 683.437] /A << /S /GoTo /D <74686D2E452E352E322E32> >> >>}
\MusicPDFWriteObject{8574}{<< /Type /Annot /Subtype /Link /Border [0 0 0] /H /I /C [1 0 0] /Rect [448.761 630.268 481.373 641.059] /A << /S /GoTo /D <74686D2E452E352E322E31> >> >>}
\MusicPDFWriteObject{8575}{<< /Type /Annot /Subtype /Link /Border [0 0 0] /H /I /C [1 0 0] /Rect [225.065 616.564 257.166 628.519] /A << /S /GoTo /D <74686D2E452E352E322E31> >> >>}
\MusicPDFWriteObject{8576}{<< /Type /Annot /Subtype /Link /Border [0 0 0] /H /I /C [1 0 0] /Rect [221.677 548.279 253.088 560.232] /A << /S /GoTo /D <74686D2E452E352E322E31> >> >>}
\MusicPDFWriteObject{8577}{<< /Type /Annot /Subtype /Link /Border [0 0 0] /H /I /C [1 0 0] /Rect [414.638 535.128 446.65 547.083] /A << /S /GoTo /D <74686D2E452E352E322E31> >> >>}
\MusicPDFWriteObject{8578}{<< /Type /Annot /Subtype /Link /Border [0 0 0] /H /I /C [1 0 0] /Rect [184.957 459.137 209.667 471.237] /A << /S /GoTo /D <73756273656374696F6E2E452E352E31> >> >>}
\MusicPDFWriteObject{8579}{<< /Type /Annot /Subtype /Link /Border [0 0 0] /H /I /C [1 0 0] /Rect [117.346 446.54 154.338 457.331] /A << /S /GoTo /D <74686D2E452E352E312E3135> >> >>}
\MusicPDFWriteObject{8580}{<< /Type /Annot /Subtype /Link /Border [0 0 0] /H /I /C [0 1 0] /Rect [448.191 435.326 474.503 444.194] /A << /S /GoTo /D <636974652E426F7231> >> >>}
\MusicPDFWriteObject{8581}{<< /Type /Annot /Subtype /Link /Border [0 0 0] /H /I /C [0 1 0] /Rect [169.706 422.175 189.295 431.043] /A << /S /GoTo /D <636974652E4752> >> >>}
\MusicPDFWriteObject{8582}{<< /Type /Annot /Subtype /Link /Border [0 0 0] /H /I /C [0 1 0] /Rect [489.664 422.175 515.976 431.043] /A << /S /GoTo /D <636974652E426F7231> >> >>}
\MusicPDFWriteObject{8583}{<< /Type /Annot /Subtype /Link /Border [0 0 0] /H /I /C [0 1 0] /Rect [93.138 409.024 112.726 417.892] /A << /S /GoTo /D <636974652E4752> >> >>}
\MusicPDFWriteObject{8584}{<< /Type /Annot /Subtype /Link /Border [0 0 0] /H /I /C [1 0 0] /Rect [242.527 180.415 274.508 192.368] /A << /S /GoTo /D <6571756174696F6E2E452E352E332E33> >> >>}
\MusicPDFWriteObject{8585}{<< /Type /Annot /Subtype /Link /Border [0 0 0] /H /I /C [1 0 0] /Rect [349.292 356.728 380.703 367.576] /A << /S /GoTo /D <74686D2E452E352E332E31> >> >>}
\MusicPDFWriteObject{8586}{<< /Type /Annot /Subtype /Link /Border [0 0 0] /H /I /C [1 0 0] /Rect [297.131 306.237 328.542 318.552] /A << /S /GoTo /D <74686D2E452E352E322E34> >> >>}
\MusicPDFWriteObject{8587}{<< /Type /Annot /Subtype /Link /Border [0 0 0] /H /I /C [1 0 0] /Rect [286.548 218.92 317.96 230.339] /A << /S /GoTo /D <74686D2E452E352E322E36> >> >>}
\MusicPDFWriteObject{8588}{<< /Type /Annot /Subtype /Link /Border [0 0 0] /H /I /C [1 0 0] /Rect [446.594 192.065 478.245 204.018] /A << /S /GoTo /D <74686D2E452E352E332E31> >> >>}
\MusicPDFWriteObject{8589}{<< /Type /Annot /Subtype /Link /Border [0 0 0] /H /I /C [1 0 0] /Rect [320.445 153.413 351.856 165.365] /A << /S /GoTo /D <74686D2E452E352E322E37> >> >>}
\MusicPDFWriteObject{8590}{<< /Type /Annot /Subtype /Link /Border [0 0 0] /H /I /C [1 0 0] /Rect [147.344 114.514 179.506 125.304] /A << /S /GoTo /D <74686D2E452E352E332E31> >> >>}
\MusicPDFWriteObject{8591}{<< /Type /Annot /Subtype /Link /Border [0 0 0] /H /I /C [1 0 0] /Rect [238.725 639.312 271.337 651.627] /A << /S /GoTo /D <74686D2E452E352E332E31> >> >>}
\MusicPDFWriteObject{8592}{<< /Type /Annot /Subtype /Link /Border [0 0 0] /H /I /C [1 0 0] /Rect [86.882 560.768 119.493 572.721] /A << /S /GoTo /D <6571756174696F6E2E452E352E332E33> >> >>}
\MusicPDFWriteObject{8593}{<< /Type /Annot /Subtype /Link /Border [0 0 0] /H /I /C [1 0 0] /Rect [320.48 547.617 352.492 560.475] /A << /S /GoTo /D <74686D2E452E352E322E31> >> >>}
\MusicPDFWriteObject{8594}{<< /Type /Annot /Subtype /Link /Border [0 0 0] /H /I /C [1 0 0] /Rect [353.266 547.617 363.199 560.475] /A << /S /GoTo /D <48666F6F746E6F74652E313731> >> >>}
\MusicPDFWriteObject{8595}{<< /Type /Annot /Subtype /Link /Border [0 0 0] /H /I /C [1 0 0] /Rect [71.004 493.962 103.015 506.969] /A << /S /GoTo /D <74686D2E452E352E322E31> >> >>}
\MusicPDFWriteObject{8596}{<< /Type /Annot /Subtype /Link /Border [0 0 0] /H /I /C [1 0 0] /Rect [232.94 468.713 264.952 480.666] /A << /S /GoTo /D <6571756174696F6E2E452E352E332E33> >> >>}
\MusicPDFWriteObject{8597}{<< /Type /Annot /Subtype /Link /Border [0 0 0] /H /I /C [1 0 0] /Rect [293.393 468.713 325.405 480.666] /A << /S /GoTo /D <6571756174696F6E2E452E352E332E34> >> >>}
\MusicPDFWriteObject{8598}{<< /Type /Annot /Subtype /Link /Border [0 0 0] /H /I /C [1 0 0] /Rect [199.601 454.51 232.212 467.517] /A << /S /GoTo /D <74686D2E452E352E332E31> >> >>}
\MusicPDFWriteObject{8599}{<< /Type /Annot /Subtype /Link /Border [0 0 0] /H /I /C [1 0 0] /Rect [396.198 272.047 428.81 282.838] /A << /S /GoTo /D <74686D2E452E352E332E31> >> >>}
\MusicPDFWriteObject{8600}{<< /Type /Annot /Subtype /Link /Border [0 0 0] /H /I /C [1 0 0] /Rect [223.821 244.832 255.232 257.147] /A << /S /GoTo /D <74686D2E452E352E332E31> >> >>}
\MusicPDFWriteObject{8601}{<< /Type /Annot /Subtype /Link /Border [0 0 0] /H /I /C [0 1 0] /Rect [182.088 234.532 198.612 243.399] /A << /S /GoTo /D <636974652E4861> >> >>}
\MusicPDFWriteObject{8602}{<< /Type /Annot /Subtype /Link /Border [0 0 0] /H /I /C [0 1 0] /Rect [482.744 206.138 499.268 215.006] /A << /S /GoTo /D <636974652E4861> >> >>}
\MusicPDFWriteObject{8603}{<< /Type /Annot /Subtype /Link /Border [0 0 0] /H /I /C [0 1 0] /Rect [146.896 130.258 163.42 139.126] /A << /S /GoTo /D <636974652E4861> >> >>}
\MusicPDFWriteObject{8604}{<< /Type /Annot /Subtype /Link /Border [0 0 0] /H /I /C [0 1 0] /Rect [322.733 117.107 342.322 125.975] /A << /S /GoTo /D <636974652E4752> >> >>}
\MusicPDFWriteObject{8605}{<< /Type /Annot /Subtype /Link /Border [0 0 0] /H /I /C [0 1 0] /Rect [484.846 117.107 501.371 125.975] /A << /S /GoTo /D <636974652E4861> >> >>}
\MusicPDFWriteObject{8606}{<< /Type /Annot /Subtype /Link /Border [0 0 0] /H /I /C [1 0 0] /Rect [406.83 671.482 438.811 683.437] /A << /S /GoTo /D <74686D2E452E352E322E37> >> >>}
\MusicPDFWriteObject{8607}{<< /Type /Annot /Subtype /Link /Border [0 0 0] /H /I /C [0 1 0] /Rect [144.791 647.671 160.977 656.539] /A << /S /GoTo /D <636974652E5350> >> >>}
\MusicPDFWriteObject{8608}{<< /Type /Annot /Subtype /Link /Border [0 0 0] /H /I /C [0 1 0] /Rect [379.148 634.52 398.737 643.388] /A << /S /GoTo /D <636974652E4752> >> >>}
\MusicPDFWriteObject{8609}{<< /Type /Annot /Subtype /Link /Border [0 0 0] /H /I /C [0 1 0] /Rect [220.924 621.369 237.11 630.237] /A << /S /GoTo /D <636974652E5350> >> >>}
\MusicPDFWriteObject{8610}{<< /Type /Annot /Subtype /Link /Border [0 0 0] /H /I /C [1 0 0] /Rect [340.832 605.728 372.844 617.683] /A << /S /GoTo /D <74686D2E452E352E332E31> >> >>}
\MusicPDFWriteObject{8611}{<< /Type /Annot /Subtype /Link /Border [0 0 0] /H /I /C [1 0 0] /Rect [134.804 522.635 166.816 533.425] /A << /S /GoTo /D <74686D2E452E352E332E31> >> >>}
\MusicPDFWriteObject{8612}{<< /Type /Annot /Subtype /Link /Border [0 0 0] /H /I /C [0 1 0] /Rect [277.506 375.503 293.692 384.371] /A << /S /GoTo /D <636974652E5350> >> >>}
\MusicPDFWriteObject{8613}{<< /Type /Annot /Subtype /Link /Border [0 0 0] /H /I /C [1 0 0] /Rect [484.707 378.837 522.399 390.792] /A << /S /GoTo /D <74686D2E452E352E312E3135> >> >>}
\MusicPDFWriteObject{8614}{<< /Type /Annot /Subtype /Link /Border [0 0 0] /H /I /C [1 0 0] /Rect [408.725 175.45 440.737 186.298] /A << /S /GoTo /D <74686D2E452E352E342E31> >> >>}
\MusicPDFWriteObject{8615}{<< /Type /Annot /Subtype /Link /Border [0 0 0] /H /I /C [1 0 0] /Rect [176.001 648.759 209.185 659.55] /A << /S /GoTo /D <74686D2E452E352E342E31> >> >>}
\MusicPDFWriteObject{8616}{<< /Type /Annot /Subtype /Link /Border [0 0 0] /H /I /C [1 0 0] /Rect [468.882 382.512 506.118 394.467] /A << /S /GoTo /D <74686D2E452E352E312E3135> >> >>}
\MusicPDFWriteObject{8617}{<< /Type /Annot /Subtype /Link /Border [0 0 0] /H /I /C [1 0 0] /Rect [73.712 343.06 105.723 355.013] /A << /S /GoTo /D <6571756174696F6E2E452E352E342E32> >> >>}
\MusicPDFWriteObject{8618}{<< /Type /Annot /Subtype /Link /Border [0 0 0] /H /I /C [1 0 0] /Rect [321.554 244.652 358.545 256.605] /A << /S /GoTo /D <6571756174696F6E2E452E352E312E3234> >> >>}
\MusicPDFWriteObject{8619}{<< /Type /Annot /Subtype /Link /Border [0 0 0] /H /I /C [1 0 0] /Rect [431.515 244.652 467.807 256.607] /A << /S /GoTo /D <74686D2E452E352E312E3136> >> >>}
\MusicPDFWriteObject{8620}{<< /Type /Annot /Subtype /Link /Border [0 0 0] /H /I /C [1 0 0] /Rect [383.398 196.141 422.399 207.845] /A << /S /GoTo /D <74686D2E452E352E312E3134> >> >>}
\MusicPDFWriteObject{8621}{<< /Type /Annot /Subtype /Link /Border [0 0 0] /H /I /C [1 0 0] /Rect [480.151 196.141 519.152 207.845] /A << /S /GoTo /D <74686D2E452E352E312E3139> >> >>}
\MusicPDFWriteObject{8622}{<< /Type /Annot /Subtype /Link /Border [0 0 0] /H /I /C [1 0 0] /Rect [326.074 135.199 362.366 147.152] /A << /S /GoTo /D <74686D2E452E352E312E3135> >> >>}
\MusicPDFWriteObject{8623}{<< /Type /Annot /Subtype /Link /Border [0 0 0] /H /I /C [1 0 0] /Rect [439.452 135.199 470.863 147.152] /A << /S /GoTo /D <74686D2E452E352E342E34> >> >>}
\MusicPDFWriteObject{8624}{<< /Type /Annot /Subtype /Link /Border [0 0 0] /H /I /C [1 0 0] /Rect [239.543 109.451 271.795 120.241] /A << /S /GoTo /D <74686D2E452E352E342E33> >> >>}
\MusicPDFWriteObject{8625}{<< /Type /Annot /Subtype /Link /Border [0 0 0] /H /I /C [1 0 0] /Rect [117.34 605.729 154.332 617.681] /A << /S /GoTo /D <74686D2E452E352E312E3136> >> >>}
\MusicPDFWriteObject{8626}{<< /Type /Annot /Subtype /Link /Border [0 0 0] /H /I /C [1 0 0] /Rect [450.262 593.131 482.874 603.979] /A << /S /GoTo /D <74686D2E452E352E342E33> >> >>}
\MusicPDFWriteObject{8627}{<< /Type /Annot /Subtype /Link /Border [0 0 0] /H /I /C [1 0 0] /Rect [247.347 559.544 278.758 570.334] /A << /S /GoTo /D <74686D2E452E352E342E35> >> >>}
\MusicPDFWriteObject{8628}{<< /Type /Annot /Subtype /Link /Border [0 0 0] /H /I /C [1 0 0] /Rect [106.556 506.94 138.568 518.487] /A << /S /GoTo /D <74686D2E452E352E342E35> >> >>}
\MusicPDFWriteObject{8629}{<< /Type /Annot /Subtype /Link /Border [0 0 0] /H /I /C [1 0 0] /Rect [366.93 358.579 399.541 370.534] /A << /S /GoTo /D <74686D2E452E342E312E31> >> >>}
\MusicPDFWriteObject{8630}{<< /Type /Annot /Subtype /Link /Border [0 0 0] /H /I /C [1 0 0] /Rect [299.159 275.053 338.16 287.514] /A << /S /GoTo /D <74686D2E452E352E312E3134> >> >>}
\MusicPDFWriteObject{8631}{<< /Type /Annot /Subtype /Link /Border [0 0 0] /H /I /C [1 0 0] /Rect [73.712 262.263 105.723 274.216] /A << /S /GoTo /D <6571756174696F6E2E452E352E342E33> >> >>}
\MusicPDFWriteObject{8632}{<< /Type /Annot /Subtype /Link /Border [0 0 0] /H /I /C [1 0 0] /Rect [159.484 202.927 191.495 213.718] /A << /S /GoTo /D <74686D2E452E352E342E37> >> >>}
\MusicPDFWriteObject{8633}{<< /Type /Annot /Subtype /Link /Border [0 0 0] /H /I /C [1 0 0] /Rect [128.967 116.736 160.978 127.585] /A << /S /GoTo /D <74686D2E452E342E312E31> >> >>}
\MusicPDFWriteObject{8634}{<< /Type /Annot /Subtype /Link /Border [0 0 0] /H /I /C [1 0 0] /Rect [71.004 596.014 103.015 607.967] /A << /S /GoTo /D <74686D2E452E342E312E31> >> >>}
\MusicPDFWriteObject{8635}{<< /Type /Annot /Subtype /Link /Border [0 0 0] /H /I /C [1 0 0] /Rect [243.842 544.002 275.434 554.793] /A << /S /GoTo /D <74686D2E452E352E322E34> >> >>}
\MusicPDFWriteObject{8636}{<< /Type /Annot /Subtype /Link /Border [0 0 0] /H /I /C [1 0 0] /Rect [223.59 530.851 247.854 541.699] /A << /S /GoTo /D <73756273656374696F6E2E452E352E37> >> >>}
\MusicPDFWriteObject{8637}{<< /Type /Annot /Subtype /Link /Border [0 0 0] /H /I /C [1 0 0] /Rect [289.598 312.149 321.01 322.939] /A << /S /GoTo /D <74686D2E452E352E322E34> >> >>}
\MusicPDFWriteObject{8638}{<< /Type /Annot /Subtype /Link /Border [0 0 0] /H /I /C [1 0 0] /Rect [147.034 425.512 179.046 437.464] /A << /S /GoTo /D <6571756174696F6E2E452E352E352E32> >> >>}
\MusicPDFWriteObject{8639}{<< /Type /Annot /Subtype /Link /Border [0 0 0] /H /I /C [1 0 0] /Rect [490.55 684.633 523.162 696.586] /A << /S /GoTo /D <74686D2E452E352E322E31> >> >>}
\MusicPDFWriteObject{8640}{<< /Type /Annot /Subtype /Link /Border [0 0 0] /H /I /C [1 0 0] /Rect [346.517 297.612 378.529 309.567] /A << /S /GoTo /D <74686D2E452E352E352E31> >> >>}
\MusicPDFWriteObject{8641}{<< /Type /Annot /Subtype /Link /Border [0 0 0] /H /I /C [1 0 0] /Rect [235.3 250.525 268.484 259.379] /A << /S /GoTo /D <74686D2E452E352E352E32> >> >>}
\MusicPDFWriteObject{8642}{<< /Type /Annot /Subtype /Link /Border [0 0 0] /H /I /C [1 0 0] /Rect [271.214 234.883 303.225 246.838] /A << /S /GoTo /D <6571756174696F6E2E452E352E352E36> >> >>}
\MusicPDFWriteObject{8643}{<< /Type /Annot /Subtype /Link /Border [0 0 0] /H /I /C [1 0 0] /Rect [331.888 234.883 363.899 246.838] /A << /S /GoTo /D <6571756174696F6E2E452E352E352E38> >> >>}
\MusicPDFWriteObject{8644}{<< /Type /Annot /Subtype /Link /Border [0 0 0] /H /I /C [1 0 0] /Rect [386.27 95.747 417.711 107.699] /A << /S /GoTo /D <74686D2E452E352E322E31> >> >>}
\MusicPDFWriteObject{8645}{<< /Type /Annot /Subtype /Link /Border [0 0 0] /H /I /C [1 0 0] /Rect [380.602 619.373 412.013 631.325] /A << /S /GoTo /D <74686D2E452E352E322E31> >> >>}
\MusicPDFWriteObject{8646}{<< /Type /Annot /Subtype /Link /Border [0 0 0] /H /I /C [1 0 0] /Rect [376.645 606.221 409.257 618.176] /A << /S /GoTo /D <74686D2E452E352E322E31> >> >>}
\MusicPDFWriteObject{8647}{<< /Type /Annot /Subtype /Link /Border [0 0 0] /H /I /C [1 0 0] /Rect [507.144 529.716 539.756 541.671] /A << /S /GoTo /D <74686D2E452E352E322E36> >> >>}
\MusicPDFWriteObject{8648}{<< /Type /Annot /Subtype /Link /Border [0 0 0] /H /I /C [0 1 0] /Rect [305.247 505.905 321.786 514.773] /A << /S /GoTo /D <636974652E4B61> >> >>}
\MusicPDFWriteObject{8649}{<< /Type /Annot /Subtype /Link /Border [0 0 0] /H /I /C [0 1 0] /Rect [476.975 505.905 493.513 514.773] /A << /S /GoTo /D <636974652E4B61> >> >>}
\MusicPDFWriteObject{8650}{<< /Type /Annot /Subtype /Link /Border [0 0 0] /H /I /C [0 1 0] /Rect [523.639 405.885 543.228 414.753] /A << /S /GoTo /D <636974652E4752> >> >>}
\MusicPDFWriteObject{8651}{<< /Type /Annot /Subtype /Link /Border [0 0 0] /H /I /C [0 1 0] /Rect [157.388 379.583 176.977 388.451] /A << /S /GoTo /D <636974652E4752> >> >>}
\MusicPDFWriteObject{8652}{<< /Type /Annot /Subtype /Link /Border [0 0 0] /H /I /C [0 1 0] /Rect [296.15 379.583 315.739 388.451] /A << /S /GoTo /D <636974652E4752> >> >>}
\MusicPDFWriteObject{8653}{<< /Type /Annot /Subtype /Link /Border [0 0 0] /H /I /C [1 0 0] /Rect [314.489 351.345 345.901 362.136] /A << /S /GoTo /D <74686D2E452E352E362E31> >> >>}
\MusicPDFWriteObject{8654}{<< /Type /Annot /Subtype /Link /Border [0 0 0] /H /I /C [1 0 0] /Rect [447.815 337.641 479.947 349.596] /A << /S /GoTo /D <74686D2E452E352E352E34> >> >>}
\MusicPDFWriteObject{8655}{<< /Type /Annot /Subtype /Link /Border [0 0 0] /H /I /C [0 1 0] /Rect [477.447 313.83 497.035 322.698] /A << /S /GoTo /D <636974652E4752> >> >>}
\MusicPDFWriteObject{8656}{<< /Type /Annot /Subtype /Link /Border [0 0 0] /H /I /C [1 0 0] /Rect [183.719 228.787 215.641 239.635] /A << /S /GoTo /D <74686D2E452E352E362E32> >> >>}
\MusicPDFWriteObject{8657}{<< /Type /Annot /Subtype /Link /Border [0 0 0] /H /I /C [0 1 0] /Rect [270.273 560.343 286.459 569.211] /A << /S /GoTo /D <636974652E5350> >> >>}
\MusicPDFWriteObject{8658}{<< /Type /Annot /Subtype /Link /Border [0 0 0] /H /I /C [0 1 0] /Rect [174.195 459.865 194.268 468.733] /A << /S /GoTo /D <636974652E464D> >> >>}
\MusicPDFWriteObject{8659}{<< /Type /Annot /Subtype /Link /Border [0 0 0] /H /I /C [0 1 0] /Rect [261.647 446.714 277.832 455.582] /A << /S /GoTo /D <636974652E5350> >> >>}
\MusicPDFWriteObject{8660}{<< /Type /Annot /Subtype /Link /Border [0 0 0] /H /I /C [1 0 0] /Rect [415.697 241.591 439.516 253.564] /A << /S /GoTo /D <73756273656374696F6E2E452E352E31> >> >>}
\MusicPDFWriteObject{8661}{<< /Type /Annot /Subtype /Link /Border [0 0 0] /H /I /C [1 0 0] /Rect [505.674 241.591 541.966 253.564] /A << /S /GoTo /D <74686D2E452E352E312E3131> >> >>}
\MusicPDFWriteObject{8662}{<< /Type /Annot /Subtype /Link /Border [0 0 0] /H /I /C [1 0 0] /Rect [256.057 192.567 287.558 203.358] /A << /S /GoTo /D <74686D2E452E352E322E31> >> >>}
\MusicPDFWriteObject{8663}{<< /Type /Annot /Subtype /Link /Border [0 0 0] /H /I /C [1 0 0] /Rect [337.555 192.567 373.952 203.358] /A << /S /GoTo /D <74686D2E452E352E312E3134> >> >>}
\MusicPDFWriteObject{8664}{<< /Type /Annot /Subtype /Link /Border [0 0 0] /H /I /C [1 0 0] /Rect [393.348 165.712 430.06 177.665] /A << /S /GoTo /D <6571756174696F6E2E452E352E312E3230> >> >>}
\MusicPDFWriteObject{8665}{<< /Type /Annot /Subtype /Link /Border [0 0 0] /H /I /C [1 0 0] /Rect [485.906 165.712 522.618 177.665] /A << /S /GoTo /D <74686D2E452E352E312E3134> >> >>}
\MusicPDFWriteObject{8666}{<< /Type /Annot /Subtype /Link /Border [0 0 0] /H /I /C [1 0 0] /Rect [309.296 135.198 340.827 147.153] /A << /S /GoTo /D <74686D2E452E352E372E32> >> >>}
\MusicPDFWriteObject{8667}{<< /Type /Annot /Subtype /Link /Border [0 0 0] /H /I /C [0 1 0] /Rect [290.323 111.387 306.509 120.255] /A << /S /GoTo /D <636974652E5350> >> >>}
\MusicPDFWriteObject{8668}{<< /Type /Annot /Subtype /Link /Border [0 0 0] /H /I /C [0 1 0] /Rect [434.055 111.387 450.24 120.255] /A << /S /GoTo /D <636974652E5350> >> >>}
\MusicPDFWriteObject{8669}{<< /Type /Annot /Subtype /Link /Border [0 0 0] /H /I /C [1 0 0] /Rect [452.147 95.746 484.099 108.535] /A << /S /GoTo /D <74686D2E452E352E372E32> >> >>}
\MusicPDFWriteObject{8670}{<< /Type /Annot /Subtype /Link /Border [0 0 0] /H /I /C [1 0 0] /Rect [316.26 375.434 349.445 388.416] /A << /S /GoTo /D <74686D2E452E352E372E33> >> >>}
\MusicPDFWriteObject{8671}{<< /Type /Annot /Subtype /Link /Border [0 0 0] /H /I /C [0 1 0] /Rect [163.355 358.695 179.541 367.563] /A << /S /GoTo /D <636974652E5350> >> >>}
\MusicPDFWriteObject{8672}{<< /Type /Annot /Subtype /Link /Border [0 0 0] /H /I /C [0 1 0] /Rect [381.622 358.695 397.047 367.563] /A << /S /GoTo /D <636974652E4163> >> >>}
\MusicPDFWriteObject{8673}{<< /Type /Annot /Subtype /Link /Border [0 0 0] /H /I /C [0 1 0] /Rect [404.592 339.567 420.778 348.435] /A << /S /GoTo /D <636974652E5350> >> >>}
\MusicPDFWriteObject{8674}{<< /Type /Annot /Subtype /Link /Border [0 0 0] /H /I /C [0 1 0] /Rect [230.437 104.214 246.623 113.082] /A << /S /GoTo /D <636974652E5350> >> >>}
\MusicPDFWriteObject{8675}{<< /Type /Annot /Subtype /Link /Border [0 0 0] /H /I /C [0 1 0] /Rect [322.237 104.214 337.662 113.082] /A << /S /GoTo /D <636974652E4163> >> >>}
\MusicPDFWriteObject{8676}{<< /Type /Annot /Subtype /Link /Border [0 0 0] /H /I /C [0 1 0] /Rect [136.386 85.086 152.572 93.954] /A << /S /GoTo /D <636974652E5350> >> >>}
\MusicPDFWriteObject{8677}{<< /Type /Annot /Subtype /Link /Border [0 0 0] /H /I /C [1 0 0] /Rect [169.145 494.205 193.856 506.161] /A << /S /GoTo /D <73756273656374696F6E2E452E352E31> >> >>}
\MusicPDFWriteObject{8678}{<< /Type /Annot /Subtype /Link /Border [0 0 0] /H /I /C [0 1 0] /Rect [189.365 465.237 205.551 474.105] /A << /S /GoTo /D <636974652E5350> >> >>}
\MusicPDFWriteObject{8679}{<< /Type /Annot /Subtype /Link /Border [0 0 0] /H /I /C [0 1 0] /Rect [119.338 438.936 135.524 447.804] /A << /S /GoTo /D <636974652E5350> >> >>}
\MusicPDFWriteObject{8680}{<< /Type /Annot /Subtype /Link /Border [0 0 0] /H /I /C [1 0 0] /Rect [463.606 404.987 496.167 417.776] /A << /S /GoTo /D <74686D2E452E352E372E34> >> >>}
\MusicPDFWriteObject{8681}{<< /Type /Annot /Subtype /Link /Border [0 0 0] /H /I /C [0 1 0] /Rect [163.955 326.44 180.14 335.308] /A << /S /GoTo /D <636974652E5350> >> >>}
\MusicPDFWriteObject{8682}{<< /Type /Annot /Subtype /Link /Border [0 0 0] /H /I /C [0 1 0] /Rect [258.379 326.44 274.564 335.308] /A << /S /GoTo /D <636974652E5350> >> >>}
\MusicPDFWriteObject{8683}{<< /Type /Annot /Subtype /Link /Border [0 0 0] /H /I /C [0 1 0] /Rect [467.845 326.44 483.27 335.308] /A << /S /GoTo /D <636974652E4163> >> >>}
\MusicPDFWriteObject{8684}{<< /Type /Annot /Subtype /Link /Border [0 0 0] /H /I /C [0 1 0] /Rect [131.374 313.29 146.799 322.158] /A << /S /GoTo /D <636974652E4163> >> >>}
\MusicPDFWriteObject{8685}{<< /Type /Annot /Subtype /Link /Border [0 0 0] /H /I /C [0 1 0] /Rect [163.355 205.763 179.541 214.631] /A << /S /GoTo /D <636974652E5350> >> >>}
\MusicPDFWriteObject{8686}{<< /Type /Annot /Subtype /Link /Border [0 0 0] /H /I /C [0 1 0] /Rect [369.449 205.763 384.874 214.631] /A << /S /GoTo /D <636974652E4163> >> >>}
\MusicPDFWriteObject{8687}{<< /Type /Annot /Subtype /Link /Border [0 0 0] /H /I /C [1 0 0] /Rect [225.929 184.965 263.621 196.918] /A << /S /GoTo /D <74686D2E452E352E372E3131> >> >>}
\MusicPDFWriteObject{8688}{<< /Type /Annot /Subtype /Link /Border [0 0 0] /H /I /C [1 0 0] /Rect [109.842 171.814 147.534 183.769] /A << /S /GoTo /D <74686D2E452E352E312E3230> >> >>}
\MusicPDFWriteObject{8689}{<< /Type /Annot /Subtype /Link /Border [0 0 0] /H /I /C [1 0 0] /Rect [456.714 159.217 488.725 170.008] /A << /S /GoTo /D <74686D2E452E352E372E33> >> >>}
\MusicPDFWriteObject{8690}{<< /Type /Annot /Subtype /Link /Border [0 0 0] /H /I /C [0 1 0] /Rect [526.973 124.538 543.159 133.406] /A << /S /GoTo /D <636974652E5350> >> >>}
\MusicPDFWriteObject{8691}{<< /Type /Annot /Subtype /Link /Border [0 0 0] /H /I /C [0 1 0] /Rect [163.355 687.123 179.541 695.991] /A << /S /GoTo /D <636974652E5350> >> >>}
\MusicPDFWriteObject{8692}{<< /Type /Annot /Subtype /Link /Border [0 0 0] /H /I /C [1 0 0] /Rect [507.58 613.455 539.861 624.246] /A << /S /GoTo /D <74686D2E452E352E372E34> >> >>}
\MusicPDFWriteObject{8693}{<< /Type /Annot /Subtype /Link /Border [0 0 0] /H /I /C [0 1 0] /Rect [314.879 589.09 330.304 597.958] /A << /S /GoTo /D <636974652E4163> >> >>}
\MusicPDFWriteObject{8694}{<< /Type /Annot /Subtype /Link /Border [0 0 0] /H /I /C [0 1 0] /Rect [500.975 589.09 516.4 597.958] /A << /S /GoTo /D <636974652E4163> >> >>}
\MusicPDFWriteObject{8695}{<< /Type /Annot /Subtype /Link /Border [0 0 0] /H /I /C [0 1 0] /Rect [273.833 575.94 289.258 584.807] /A << /S /GoTo /D <636974652E4163> >> >>}
\MusicPDFWriteObject{8696}{<< /Type /Annot /Subtype /Link /Border [0 0 0] /H /I /C [1 0 0] /Rect [313.235 527.759 345.247 539.715] /A << /S /GoTo /D <74686D2E452E352E322E36> >> >>}
\MusicPDFWriteObject{8697}{<< /Type /Annot /Subtype /Link /Border [0 0 0] /H /I /C [0 1 0] /Rect [233.544 365.75 248.969 374.618] /A << /S /GoTo /D <636974652E4163> >> >>}
\MusicPDFWriteObject{8698}{<< /Type /Annot /Subtype /Link /Border [0 0 0] /H /I /C [1 0 0] /Rect [157.519 350.109 194.511 362.062] /A << /S /GoTo /D <74686D2E452E352E312E3230> >> >>}
\MusicPDFWriteObject{8699}{<< /Type /Annot /Subtype /Link /Border [0 0 0] /H /I /C [0 1 0] /Rect [290.84 215.658 306.265 224.526] /A << /S /GoTo /D <636974652E4163> >> >>}
\MusicPDFWriteObject{8700}{<< /Type /Annot /Subtype /Link /Border [0 0 0] /H /I /C [0 1 0] /Rect [85.662 170.228 101.087 179.096] /A << /S /GoTo /D <636974652E4163> >> >>}
\MusicPDFWriteObject{8701}{<< /Type /Annot /Subtype /Link /Border [0 0 0] /H /I /C [0 1 0] /Rect [246.537 170.228 261.962 179.096] /A << /S /GoTo /D <636974652E4163> >> >>}
\MusicPDFWriteObject{8702}{<< /Type /Annot /Subtype /Link /Border [0 0 0] /H /I /C [0 1 0] /Rect [73.77 117.625 89.195 126.493] /A << /S /GoTo /D <636974652E4163> >> >>}
\MusicPDFWriteObject{8703}{<< /Type /Annot /Subtype /Link /Border [0 0 0] /H /I /C [0 1 0] /Rect [241.76 613.019 257.185 621.887] /A << /S /GoTo /D <636974652E4163> >> >>}
\MusicPDFWriteObject{8704}{<< /Type /Annot /Subtype /Link /Border [0 0 0] /H /I /C [0 1 0] /Rect [400.908 599.869 416.333 608.736] /A << /S /GoTo /D <636974652E4163> >> >>}
\MusicPDFWriteObject{8705}{<< /Type /Annot /Subtype /Link /Border [0 0 0] /H /I /C [0 1 0] /Rect [232.307 521.731 247.732 530.599] /A << /S /GoTo /D <636974652E4163> >> >>}
\MusicPDFWriteObject{8706}{<< /Type /Annot /Subtype /Link /Border [0 0 0] /H /I /C [1 0 0] /Rect [225.929 447.362 263.621 459.315] /A << /S /GoTo /D <74686D2E452E352E372E3138> >> >>}
\MusicPDFWriteObject{8707}{<< /Type /Annot /Subtype /Link /Border [0 0 0] /H /I /C [1 0 0] /Rect [109.267 434.211 146.958 446.166] /A << /S /GoTo /D <74686D2E452E352E312E3230> >> >>}
\MusicPDFWriteObject{8708}{<< /Type /Annot /Subtype /Link /Border [0 0 0] /H /I /C [1 0 0] /Rect [290.765 421.614 327.756 432.405] /A << /S /GoTo /D <74686D2E452E352E372E3134> >> >>}
\MusicPDFWriteObject{8709}{<< /Type /Annot /Subtype /Link /Border [0 0 0] /H /I /C [0 1 0] /Rect [230.437 274.599 245.862 283.467] /A << /S /GoTo /D <636974652E4163> >> >>}
\MusicPDFWriteObject{8710}{<< /Type /Annot /Subtype /Link /Border [0 0 0] /H /I /C [1 0 0] /Rect [235.087 252.981 266.529 264.934] /A << /S /GoTo /D <6571756174696F6E2E452E352E372E32> >> >>}
\MusicPDFWriteObject{8711}{<< /Type /Annot /Subtype /Link /Border [0 0 0] /H /I /C [1 0 0] /Rect [318.086 252.981 354.413 264.934] /A << /S /GoTo /D <74686D2E452E352E372E3230> >> >>}
\MusicPDFWriteObject{8712}{<< /Type /Annot /Subtype /Link /Border [0 0 0] /H /I /C [0 1 0] /Rect [73.77 242.32 89.195 251.188] /A << /S /GoTo /D <636974652E4163> >> >>}
\MusicPDFWriteObject{8713}{<< /Type /Annot /Subtype /Link /Border [0 0 0] /H /I /C [1 0 0] /Rect [120.751 95.938 158.023 107.894] /A << /S /GoTo /D <74686D2E452E352E372E3138> >> >>}
\MusicPDFWriteObject{8714}{<< /Type /Annot /Subtype /Link /Border [0 0 0] /H /I /C [1 0 0] /Rect [225.929 626.287 263.621 638.24] /A << /S /GoTo /D <74686D2E452E352E372E3231> >> >>}
\MusicPDFWriteObject{8715}{<< /Type /Annot /Subtype /Link /Border [0 0 0] /H /I /C [1 0 0] /Rect [108.254 613.136 144.965 625.091] /A << /S /GoTo /D <74686D2E452E352E312E3230> >> >>}
\MusicPDFWriteObject{8716}{<< /Type /Annot /Subtype /Link /Border [0 0 0] /H /I /C [1 0 0] /Rect [503.424 599.985 540.415 611.94] /A << /S /GoTo /D <74686D2E452E352E372E3139> >> >>}
\MusicPDFWriteObject{8717}{<< /Type /Annot /Subtype /Link /Border [0 0 0] /H /I /C [1 0 0] /Rect [198.142 360.05 230.754 372.003] /A << /S /GoTo /D <6571756174696F6E2E452E352E372E36> >> >>}
\MusicPDFWriteObject{8718}{<< /Type /Annot /Subtype /Link /Border [0 0 0] /H /I /C [0 1 0] /Rect [159.751 336.239 175.936 345.106] /A << /S /GoTo /D <636974652E5350> >> >>}
\MusicPDFWriteObject{8719}{<< /Type /Annot /Subtype /Link /Border [0 0 0] /H /I /C [0 1 0] /Rect [134.835 323.088 151.021 331.956] /A << /S /GoTo /D <636974652E5350> >> >>}
\MusicPDFWriteObject{8720}{<< /Type /Annot /Subtype /Link /Border [0 0 0] /H /I /C [0 1 0] /Rect [146.245 183.119 161.67 191.987] /A << /S /GoTo /D <636974652E4163> >> >>}
\MusicPDFWriteObject{8721}{<< /Type /Annot /Subtype /Link /Border [0 0 0] /H /I /C [1 0 0] /Rect [71.004 154.327 107.995 166.282] /A << /S /GoTo /D <74686D2E452E352E372E3231> >> >>}
\MusicPDFWriteObject{8722}{<< /Type /Annot /Subtype /Link /Border [0 0 0] /H /I /C [0 1 0] /Rect [460.669 124.538 476.855 133.406] /A << /S /GoTo /D <636974652E5350> >> >>}
\MusicPDFWriteObject{8723}{<< /Type /Annot /Subtype /Link /Border [0 0 0] /H /I /C [0 1 0] /Rect [233.382 595.193 248.807 604.061] /A << /S /GoTo /D <636974652E4163> >> >>}
\MusicPDFWriteObject{8724}{<< /Type /Annot /Subtype /Link /Border [0 0 0] /H /I /C [1 0 0] /Rect [216.197 559.259 253.679 571.563] /A << /S /GoTo /D <74686D2E452E352E372E3234> >> >>}
\MusicPDFWriteObject{8725}{<< /Type /Annot /Subtype /Link /Border [0 0 0] /H /I /C [1 0 0] /Rect [310.588 559.259 348.069 571.563] /A << /S /GoTo /D <74686D2E452E352E372E3235> >> >>}
\MusicPDFWriteObject{8726}{<< /Type /Annot /Subtype /Link /Border [0 0 0] /H /I /C [1 0 0] /Rect [383.347 559.259 393.28 571.563] /A << /S /GoTo /D <48666F6F746E6F74652E313732> >> >>}
\MusicPDFWriteObject{8727}{<< /Type /Annot /Subtype /Link /Border [0 0 0] /H /I /C [1 0 0] /Rect [440.6 344.537 478.292 355.328] /A << /S /GoTo /D <74686D2E452E352E372E3230> >> >>}
\MusicPDFWriteObject{8728}{<< /Type /Annot /Subtype /Link /Border [0 0 0] /H /I /C [1 0 0] /Rect [71.004 307.022 107.995 315.876] /A << /S /GoTo /D <74686D2E452E352E372E3138> >> >>}
\MusicPDFWriteObject{8729}{<< /Type /Annot /Subtype /Link /Border [0 0 0] /H /I /C [0 1 0] /Rect [505.167 262.936 521.353 271.804] /A << /S /GoTo /D <636974652E5350> >> >>}
\MusicPDFWriteObject{8730}{<< /Type /Annot /Subtype /Link /Border [0 0 0] /H /I /C [1 0 0] /Rect [472.708 226.314 510.4 237.119] /A << /S /GoTo /D <74686D2E452E352E372E3233> >> >>}
\MusicPDFWriteObject{8731}{<< /Type /Annot /Subtype /Link /Border [0 0 0] /H /I /C [1 0 0] /Rect [221.602 178.976 253.013 190.931] /A << /S /GoTo /D <74686D2E452E352E322E37> >> >>}
\MusicPDFWriteObject{8732}{<< /Type /Annot /Subtype /Link /Border [0 0 0] /H /I /C [1 0 0] /Rect [446.23 178.976 477.641 190.931] /A << /S /GoTo /D <74686D2E452E352E332E31> >> >>}
\MusicPDFWriteObject{8733}{<< /Type /Annot /Subtype /Link /Border [0 0 0] /H /I /C [0 1 0] /Rect [445.078 460.546 460.503 469.414] /A << /S /GoTo /D <636974652E4163> >> >>}
\MusicPDFWriteObject{8734}{<< /Type /Annot /Subtype /Link /Border [0 0 0] /H /I /C [1 0 0] /Rect [261.634 444.905 294.245 456.86] /A << /S /GoTo /D <74686D2E452E362E312E31> >> >>}
\MusicPDFWriteObject{8735}{<< /Type /Annot /Subtype /Link /Border [0 0 0] /H /I /C [1 0 0] /Rect [238.767 418.604 271.379 430.557] /A << /S /GoTo /D <6571756174696F6E2E452E362E312E33> >> >>}
\MusicPDFWriteObject{8736}{<< /Type /Annot /Subtype /Link /Border [0 0 0] /H /I /C [1 0 0] /Rect [351.901 418.604 384.513 430.557] /A << /S /GoTo /D <74686D2E452E362E312E31> >> >>}
\MusicPDFWriteObject{8737}{<< /Type /Annot /Subtype /Link /Border [0 0 0] /H /I /C [0 1 0] /Rect [366.408 407.943 381.833 416.811] /A << /S /GoTo /D <636974652E4163> >> >>}
\MusicPDFWriteObject{8738}{<< /Type /Annot /Subtype /Link /Border [0 0 0] /H /I /C [0 1 0] /Rect [345.261 316.381 361.447 325.249] /A << /S /GoTo /D <636974652E5350> >> >>}
\MusicPDFWriteObject{8739}{<< /Type /Annot /Subtype /Link /Border [0 0 0] /H /I /C [0 1 0] /Rect [435.816 316.381 452.002 325.249] /A << /S /GoTo /D <636974652E5350> >> >>}
\MusicPDFWriteObject{8740}{<< /Type /Annot /Subtype /Link /Border [0 0 0] /H /I /C [1 0 0] /Rect [357.317 213.479 389.929 225.435] /A << /S /GoTo /D <74686D2E452E362E312E31> >> >>}
\MusicPDFWriteObject{8741}{<< /Type /Annot /Subtype /Link /Border [0 0 0] /H /I /C [1 0 0] /Rect [447.604 213.479 480.216 225.435] /A << /S /GoTo /D <74686D2E452E362E312E32> >> >>}
\MusicPDFWriteObject{8742}{<< /Type /Annot /Subtype /Link /Border [0 0 0] /H /I /C [0 1 0] /Rect [325.895 189.669 342.081 198.536] /A << /S /GoTo /D <636974652E5350> >> >>}
\MusicPDFWriteObject{8743}{<< /Type /Annot /Subtype /Link /Border [0 0 0] /H /I /C [1 0 0] /Rect [211.849 160.876 244.461 172.832] /A << /S /GoTo /D <74686D2E452E352E322E33> >> >>}
\MusicPDFWriteObject{8744}{<< /Type /Annot /Subtype /Link /Border [0 0 0] /H /I /C [0 1 0] /Rect [308.611 137.066 324.797 145.933] /A << /S /GoTo /D <636974652E5350> >> >>}
\MusicPDFWriteObject{8745}{<< /Type /Annot /Subtype /Link /Border [0 0 0] /H /I /C [0 1 0] /Rect [360.546 117.651 380.135 126.519] /A << /S /GoTo /D <636974652E4752> >> >>}
\MusicPDFWriteObject{8746}{<< /Type /Annot /Subtype /Link /Border [0 0 0] /H /I /C [0 1 0] /Rect [199.803 624.674 219.392 633.542] /A << /S /GoTo /D <636974652E4752> >> >>}
\MusicPDFWriteObject{8747}{<< /Type /Annot /Subtype /Link /Border [0 0 0] /H /I /C [0 1 0] /Rect [324.683 524.419 344.272 533.287] /A << /S /GoTo /D <636974652E4752> >> >>}
\MusicPDFWriteObject{8748}{<< /Type /Annot /Subtype /Link /Border [0 0 0] /H /I /C [1 0 0] /Rect [71.004 508.777 102.985 520.732] /A << /S /GoTo /D <74686D2E452E362E312E34> >> >>}
\MusicPDFWriteObject{8749}{<< /Type /Annot /Subtype /Link /Border [0 0 0] /H /I /C [1 0 0] /Rect [488.026 495.627 520.038 507.58] /A << /S /GoTo /D <6571756174696F6E2E452E362E312E35> >> >>}
\MusicPDFWriteObject{8750}{<< /Type /Annot /Subtype /Link /Border [0 0 0] /H /I /C [1 0 0] /Rect [150.27 484.966 182.282 493.82] /A << /S /GoTo /D <74686D2E452E362E312E34> >> >>}
\MusicPDFWriteObject{8751}{<< /Type /Annot /Subtype /Link /Border [0 0 0] /H /I /C [0 1 0] /Rect [505.473 415.053 521.659 423.921] /A << /S /GoTo /D <636974652E5350> >> >>}
\MusicPDFWriteObject{8752}{<< /Type /Annot /Subtype /Link /Border [0 0 0] /H /I /C [0 1 0] /Rect [388.535 368.286 404.72 377.154] /A << /S /GoTo /D <636974652E5350> >> >>}
\MusicPDFWriteObject{8753}{<< /Type /Annot /Subtype /Link /Border [0 0 0] /H /I /C [0 1 0] /Rect [128.04 355.136 147.629 364.003] /A << /S /GoTo /D <636974652E4752> >> >>}
\MusicPDFWriteObject{8754}{<< /Type /Annot /Subtype /Link /Border [0 0 0] /H /I /C [1 0 0] /Rect [467.116 298.563 499.727 310.518] /A << /S /GoTo /D <74686D2E452E342E322E33> >> >>}
\MusicPDFWriteObject{8755}{<< /Type /Annot /Subtype /Link /Border [0 0 0] /H /I /C [1 0 0] /Rect [173.115 272.262 205.126 284.215] /A << /S /GoTo /D <74686D2E452E332E322E31> >> >>}
\MusicPDFWriteObject{8756}{<< /Type /Annot /Subtype /Link /Border [0 0 0] /H /I /C [0 1 0] /Rect [102.019 134.664 118.205 143.532] /A << /S /GoTo /D <636974652E5350> >> >>}
\MusicPDFWriteObject{8757}{<< /Type /Annot /Subtype /Link /Border [0 0 0] /H /I /C [1 0 0] /Rect [407.674 95.747 439.685 107.699] /A << /S /GoTo /D <6571756174696F6E2E452E342E322E32> >> >>}
\MusicPDFWriteObject{8758}{<< /Type /Annot /Subtype /Link /Border [0 0 0] /H /I /C [1 0 0] /Rect [73.712 671.482 105.723 683.435] /A << /S /GoTo /D <6571756174696F6E2E452E362E322E31> >> >>}
\MusicPDFWriteObject{8759}{<< /Type /Annot /Subtype /Link /Border [0 0 0] /H /I /C [1 0 0] /Rect [164.493 671.482 196.505 683.435] /A << /S /GoTo /D <6571756174696F6E2E452E362E322E32> >> >>}
\MusicPDFWriteObject{8760}{<< /Type /Annot /Subtype /Link /Border [0 0 0] /H /I /C [1 0 0] /Rect [287.301 579.951 319.312 591.904] /A << /S /GoTo /D <6571756174696F6E2E452E362E312E34> >> >>}
\MusicPDFWriteObject{8761}{<< /Type /Annot /Subtype /Link /Border [0 0 0] /H /I /C [1 0 0] /Rect [422.92 579.951 454.932 591.904] /A << /S /GoTo /D <6571756174696F6E2E452E362E322E31> >> >>}
\MusicPDFWriteObject{8762}{<< /Type /Annot /Subtype /Link /Border [0 0 0] /H /I /C [1 0 0] /Rect [507.932 579.951 539.944 591.904] /A << /S /GoTo /D <6571756174696F6E2E452E362E322E32> >> >>}
\MusicPDFWriteObject{8763}{<< /Type /Annot /Subtype /Link /Border [0 0 0] /H /I /C [1 0 0] /Rect [281.689 454.493 313.701 466.446] /A << /S /GoTo /D <6571756174696F6E2E452E362E312E34> >> >>}
\MusicPDFWriteObject{8764}{<< /Type /Annot /Subtype /Link /Border [0 0 0] /H /I /C [1 0 0] /Rect [414.643 454.493 446.655 466.446] /A << /S /GoTo /D <6571756174696F6E2E452E362E322E32> >> >>}
\MusicPDFWriteObject{8765}{<< /Type /Annot /Subtype /Link /Border [0 0 0] /H /I /C [1 0 0] /Rect [431.289 441.342 468.28 453.295] /A << /S /GoTo /D <74686D2E452E342E322E3132> >> >>}
\MusicPDFWriteObject{8766}{<< /Type /Annot /Subtype /Link /Border [0 0 0] /H /I /C [1 0 0] /Rect [281.689 378.613 313.701 390.566] /A << /S /GoTo /D <6571756174696F6E2E452E362E312E34> >> >>}
\MusicPDFWriteObject{8767}{<< /Type /Annot /Subtype /Link /Border [0 0 0] /H /I /C [1 0 0] /Rect [414.643 378.613 446.655 390.566] /A << /S /GoTo /D <6571756174696F6E2E452E362E322E32> >> >>}
\MusicPDFWriteObject{8768}{<< /Type /Annot /Subtype /Link /Border [0 0 0] /H /I /C [1 0 0] /Rect [288.134 365.462 320.145 377.415] /A << /S /GoTo /D <6571756174696F6E2E452E342E322E32> >> >>}
\MusicPDFWriteObject{8769}{<< /Type /Annot /Subtype /Link /Border [0 0 0] /H /I /C [1 0 0] /Rect [293.456 302.733 325.467 314.686] /A << /S /GoTo /D <6571756174696F6E2E452E362E312E34> >> >>}
\MusicPDFWriteObject{8770}{<< /Type /Annot /Subtype /Link /Border [0 0 0] /H /I /C [1 0 0] /Rect [416.573 302.733 448.584 314.686] /A << /S /GoTo /D <6571756174696F6E2E452E362E322E32> >> >>}
\MusicPDFWriteObject{8771}{<< /Type /Annot /Subtype /Link /Border [0 0 0] /H /I /C [1 0 0] /Rect [263.061 247.63 295.072 259.582] /A << /S /GoTo /D <6571756174696F6E2E452E362E322E37> >> >>}
\MusicPDFWriteObject{8772}{<< /Type /Annot /Subtype /Link /Border [0 0 0] /H /I /C [1 0 0] /Rect [262.084 234.479 294.096 246.432] /A << /S /GoTo /D <6571756174696F6E2E452E362E322E38> >> >>}
\MusicPDFWriteObject{8773}{<< /Type /Annot /Subtype /Link /Border [0 0 0] /H /I /C [0 1 0] /Rect [277.343 205.481 293.529 214.349] /A << /S /GoTo /D <636974652E5350> >> >>}
\MusicPDFWriteObject{8774}{<< /Type /Annot /Subtype /Link /Border [0 0 0] /H /I /C [0 1 0] /Rect [452.127 205.481 468.313 214.349] /A << /S /GoTo /D <636974652E5350> >> >>}
\MusicPDFWriteObject{8775}{<< /Type /Annot /Subtype /Link /Border [0 0 0] /H /I /C [0 1 0] /Rect [213.301 103.3 229.487 112.168] /A << /S /GoTo /D <636974652E5350> >> >>}
\MusicPDFWriteObject{8776}{<< /Type /Annot /Subtype /Link /Border [0 0 0] /H /I /C [1 0 0] /Rect [359.927 645.181 396.918 657.134] /A << /S /GoTo /D <6571756174696F6E2E452E362E322E3130> >> >>}
\MusicPDFWriteObject{8777}{<< /Type /Annot /Subtype /Link /Border [0 0 0] /H /I /C [1 0 0] /Rect [281.667 592.578 313.678 604.531] /A << /S /GoTo /D <6571756174696F6E2E452E362E322E39> >> >>}
\MusicPDFWriteObject{8778}{<< /Type /Annot /Subtype /Link /Border [0 0 0] /H /I /C [1 0 0] /Rect [241.277 519.723 278.269 531.676] /A << /S /GoTo /D <6571756174696F6E2E452E362E322E3131> >> >>}
\MusicPDFWriteObject{8779}{<< /Type /Annot /Subtype /Link /Border [0 0 0] /H /I /C [0 1 0] /Rect [222.413 509.062 238.598 517.93] /A << /S /GoTo /D <636974652E5350> >> >>}
\MusicPDFWriteObject{8780}{<< /Type /Annot /Subtype /Link /Border [0 0 0] /H /I /C [1 0 0] /Rect [144.156 409.44 176.167 421.393] /A << /S /GoTo /D <6571756174696F6E2E452E362E322E33> >> >>}
\MusicPDFWriteObject{8781}{<< /Type /Annot /Subtype /Link /Border [0 0 0] /H /I /C [0 1 0] /Rect [322.031 364.376 338.217 373.244] /A << /S /GoTo /D <636974652E5350> >> >>}
\MusicPDFWriteObject{8782}{<< /Type /Annot /Subtype /Link /Border [0 0 0] /H /I /C [0 1 0] /Rect [465.156 364.376 481.342 373.244] /A << /S /GoTo /D <636974652E5350> >> >>}
\MusicPDFWriteObject{8783}{<< /Type /Annot /Subtype /Link /Border [0 0 0] /H /I /C [0 1 0] /Rect [464.361 351.225 480.547 360.093] /A << /S /GoTo /D <636974652E5350> >> >>}
\MusicPDFWriteObject{8784}{<< /Type /Annot /Subtype /Link /Border [0 0 0] /H /I /C [0 1 0] /Rect [250.785 155.917 266.971 164.784] /A << /S /GoTo /D <636974652E5350> >> >>}
\MusicPDFWriteObject{8785}{<< /Type /Annot /Subtype /Link /Border [0 0 0] /H /I /C [0 1 0] /Rect [263.353 687.123 279.539 695.991] /A << /S /GoTo /D <636974652E5350> >> >>}
\MusicPDFWriteObject{8786}{<< /Type /Annot /Subtype /Link /Border [0 0 0] /H /I /C [0 1 0] /Rect [350.241 687.123 366.427 695.991] /A << /S /GoTo /D <636974652E5350> >> >>}
\MusicPDFWriteObject{8787}{<< /Type /Annot /Subtype /Link /Border [0 0 0] /H /I /C [1 0 0] /Rect [291.021 280.315 322.432 292.845] /A << /S /GoTo /D <74686D2E452E322E332E39> >> >>}
\MusicPDFWriteObject{8788}{<< /Type /Annot /Subtype /Link /Border [0 0 0] /H /I /C [1 0 0] /Rect [207.281 268.769 238.993 279.56] /A << /S /GoTo /D <74686D2E452E322E312E36> >> >>}
\MusicPDFWriteObject{8789}{<< /Type /Annot /Subtype /Link /Border [0 0 0] /H /I /C [1 0 0] /Rect [144.079 254.013 176.571 266.544] /A << /S /GoTo /D <74686D2E452E322E332E31> >> >>}
\MusicPDFWriteObject{8790}{<< /Type /Annot /Subtype /Link /Border [0 0 0] /H /I /C [1 0 0] /Rect [228.585 242.467 260.596 253.259] /A << /S /GoTo /D <74686D2E452E322E332E39> >> >>}
\MusicPDFWriteObject{8791}{<< /Type /Annot /Subtype /Link /Border [0 0 0] /H /I /C [0 1 0] /Rect [172.424 637.545 188.61 646.413] /A << /S /GoTo /D <636974652E5350> >> >>}
\MusicPDFWriteObject{8792}{<< /Type /Annot /Subtype /Link /Border [0 0 0] /H /I /C [0 1 0] /Rect [465.015 637.545 481.201 646.413] /A << /S /GoTo /D <636974652E5350> >> >>}
\MusicPDFWriteObject{8793}{<< /Type /Annot /Subtype /Link /Border [0 0 0] /H /I /C [0 1 0] /Rect [171.048 453.011 187.234 461.879] /A << /S /GoTo /D <636974652E5350> >> >>}
\MusicPDFWriteObject{8794}{<< /Type /Annot /Subtype /Link /Border [0 0 0] /H /I /C [1 0 0] /Rect [301.01 224.782 333.021 233.636] /A << /S /GoTo /D <74686D2E452E362E332E33> >> >>}
\MusicPDFWriteObject{8795}{<< /Type /Annot /Subtype /Link /Border [0 0 0] /H /I /C [1 0 0] /Rect [244.725 171.661 254.658 185.571] /A << /S /GoTo /D <48666F6F746E6F74652E313733> >> >>}
\MusicPDFWriteObject{8796}{<< /Type /Annot /Subtype /Link /Border [0 0 0] /H /I /C [1 0 0] /Rect [331.8 126.663 341.733 138.967] /A << /S /GoTo /D <48666F6F746E6F74652E313734> >> >>}
\MusicPDFWriteObject{8797}{<< /Type /Annot /Subtype /Link /Border [0 0 0] /H /I /C [1 0 0] /Rect [203.197 522.636 238.639 531.49] /A << /S /GoTo /D <74686D2E452E362E332E38> >> >>}
\MusicPDFWriteObject{8798}{<< /Type /Annot /Subtype /Link /Border [0 0 0] /H /I /C [1 0 0] /Rect [220.125 506.995 251.537 518.95] /A << /S /GoTo /D <74686D2E452E362E332E37> >> >>}
\MusicPDFWriteObject{8799}{<< /Type /Annot /Subtype /Link /Border [0 0 0] /H /I /C [1 0 0] /Rect [193.018 441.242 224.579 453.194] /A << /S /GoTo /D <74686D2E452E362E332E37> >> >>}
\MusicPDFWriteObject{8800}{<< /Type /Annot /Subtype /Link /Border [0 0 0] /H /I /C [1 0 0] /Rect [511.17 427.039 543.211 439.569] /A << /S /GoTo /D <74686D2E452E322E332E35> >> >>}
\MusicPDFWriteObject{8801}{<< /Type /Annot /Subtype /Link /Border [0 0 0] /H /I /C [1 0 0] /Rect [486.919 94.694 519.05 107.701] /A << /S /GoTo /D <74686D2E452E362E332E34> >> >>}
\MusicPDFWriteObject{8802}{<< /Type /Annot /Subtype /Link /Border [0 0 0] /H /I /C [1 0 0] /Rect [73.634 463.721 105.045 475.694] /A << /S /GoTo /D <6571756174696F6E2E452E362E332E31> >> >>}
\MusicPDFWriteObject{8803}{<< /Type /Annot /Subtype /Link /Border [0 0 0] /H /I /C [1 0 0] /Rect [268.336 354.006 300.648 365.425] /A << /S /GoTo /D <74686D2E452E362E332E36> >> >>}
\MusicPDFWriteObject{8804}{<< /Type /Annot /Subtype /Link /Border [0 0 0] /H /I /C [0 1 0] /Rect [343.451 638.948 359.637 647.816] /A << /S /GoTo /D <636974652E5350> >> >>}
\MusicPDFWriteObject{8805}{<< /Type /Annot /Subtype /Link /Border [0 0 0] /H /I /C [0 1 0] /Rect [457.608 625.798 473.793 634.666] /A << /S /GoTo /D <636974652E5350> >> >>}
\MusicPDFWriteObject{8806}{<< /Type /Annot /Subtype /Link /Border [0 0 0] /H /I /C [0 1 0] /Rect [452.765 372.843 468.95 381.711] /A << /S /GoTo /D <636974652E5350> >> >>}
\MusicPDFWriteObject{8807}{<< /Type /Annot /Subtype /Link /Border [0 0 0] /H /I /C [0 1 0] /Rect [94.958 359.692 111.144 368.56] /A << /S /GoTo /D <636974652E5350> >> >>}
\MusicPDFWriteObject{8808}{<< /Type /Annot /Subtype /Link /Border [0 0 0] /H /I /C [0 1 0] /Rect [157.74 346.541 173.926 355.409] /A << /S /GoTo /D <636974652E5350> >> >>}
\MusicPDFWriteObject{8809}{<< /Type /Annot /Subtype /Link /Border [0 0 0] /H /I /C [1 0 0] /Rect [71.004 330.901 103.105 342.853] /A << /S /GoTo /D <74686D2E452E362E342E31> >> >>}
\MusicPDFWriteObject{8810}{<< /Type /Annot /Subtype /Link /Border [0 0 0] /H /I /C [0 1 0] /Rect [245.932 333.391 262.117 342.259] /A << /S /GoTo /D <636974652E5350> >> >>}
\MusicPDFWriteObject{8811}{<< /Type /Annot /Subtype /Link /Border [0 0 0] /H /I /C [1 0 0] /Rect [495.59 662.308 527.001 674.64] /A << /S /GoTo /D <74686D2E452E322E332E36> >> >>}
\MusicPDFWriteObject{8812}{<< /Type /Annot /Subtype /Link /Border [0 0 0] /H /I /C [1 0 0] /Rect [479.965 649.517 511.976 661.489] /A << /S /GoTo /D <74686D2E452E362E342E32> >> >>}
\MusicPDFWriteObject{8813}{<< /Type /Annot /Subtype /Link /Border [0 0 0] /H /I /C [1 0 0] /Rect [403.937 612.303 413.87 625.161] /A << /S /GoTo /D <48666F6F746E6F74652E313735> >> >>}
\MusicPDFWriteObject{8814}{<< /Type /Annot /Subtype /Link /Border [0 0 0] /H /I /C [1 0 0] /Rect [102.787 599.153 112.72 612.011] /A << /S /GoTo /D <48666F6F746E6F74652E313736> >> >>}
\MusicPDFWriteObject{8815}{<< /Type /Annot /Subtype /Link /Border [0 0 0] /H /I /C [1 0 0] /Rect [113.335 588.493 123.268 598.86] /A << /S /GoTo /D <48666F6F746E6F74652E313737> >> >>}
\MusicPDFWriteObject{8816}{<< /Type /Annot /Subtype /Link /Border [0 0 0] /H /I /C [1 0 0] /Rect [517.561 526.783 527.494 542.773] /A << /S /GoTo /D <48666F6F746E6F74652E313738> >> >>}
\MusicPDFWriteObject{8817}{<< /Title \MusicPDFRef{8820} /A \MusicPDFRef{8821} /Parent \MusicPDFRef{1} /Next \MusicPDFRef{8822} >>}
\MusicPDFWriteObject{8818}{<< /Title \MusicPDFRef{8823} /A \MusicPDFRef{8824} /Parent \MusicPDFRef{1} /Prev \MusicPDFRef{8825} >>}
\MusicPDFWriteObject{8819}{<< /Kids [\MusicPDFRef{8826} \MusicPDFRef{8827} \MusicPDFRef{8828} \MusicPDFRef{8829} \MusicPDFRef{8830} \MusicPDFRef{8831}] /Limits [<446F632D5374617274> <74686D2E452E362E342E38>] >>}
\MusicPDFWriteObject{8820}{<FEFF0049006E0064006500780020006F0066002000530079006D0062006F006C0073>}
\MusicPDFWriteObject{8821}{<< /S /GoTo /D <636861707465722A2E32> >>}
\MusicPDFWriteObject{8822}{<< /Title \MusicPDFRef{8832} /A \MusicPDFRef{8833} /Parent \MusicPDFRef{1} /Prev \MusicPDFRef{8817} /Next \MusicPDFRef{8834} /First \MusicPDFRef{8835} /Last \MusicPDFRef{8836} /Count -8 >>}
\MusicPDFWriteObject{8823}{<FEFF004200690062006C0069006F006700720061007000680079>}
\MusicPDFWriteObject{8824}{<< /S /GoTo /D <617070656E6469782A2E33> >>}
\MusicPDFWriteObject{8825}{<< /Title \MusicPDFRef{8837} /A \MusicPDFRef{8838} /Parent \MusicPDFRef{1} /Prev \MusicPDFRef{8839} /Next \MusicPDFRef{8818} /First \MusicPDFRef{8840} /Last \MusicPDFRef{8841} /Count -6 >>}
\MusicPDFWriteObject{8826}{<< /Kids [\MusicPDFRef{8842} \MusicPDFRef{8843} \MusicPDFRef{8844} \MusicPDFRef{8845} \MusicPDFRef{8846} \MusicPDFRef{8847}] /Limits [<446F632D5374617274> <4974656D2E353632>] >>}
\MusicPDFWriteObject{8827}{<< /Kids [\MusicPDFRef{8848} \MusicPDFRef{8849} \MusicPDFRef{8850} \MusicPDFRef{8851} \MusicPDFRef{8852} \MusicPDFRef{8853}] /Limits [<4974656D2E353633> <6571756174696F6E2E322E322E332E38>] >>}
\MusicPDFWriteObject{8828}{<< /Kids [\MusicPDFRef{8854} \MusicPDFRef{8855} \MusicPDFRef{8856} \MusicPDFRef{8857} \MusicPDFRef{8858} \MusicPDFRef{8859}] /Limits [<6571756174696F6E2E322E322E332E39> <6571756174696F6E2E422E332E322E34>] >>}
\MusicPDFWriteObject{8829}{<< /Kids [\MusicPDFRef{8860} \MusicPDFRef{8861} \MusicPDFRef{8862} \MusicPDFRef{8863} \MusicPDFRef{8864} \MusicPDFRef{8865}] /Limits [<6571756174696F6E2E422E332E322E35> <706167652E3532>] >>}
\MusicPDFWriteObject{8830}{<< /Kids [\MusicPDFRef{8866} \MusicPDFRef{8867} \MusicPDFRef{8868} \MusicPDFRef{8869} \MusicPDFRef{8870} \MusicPDFRef{8871}] /Limits [<706167652E353230> <74686D2E352E352E312E32>] >>}
\MusicPDFWriteObject{8831}{<< /Kids [\MusicPDFRef{8872} \MusicPDFRef{8873} \MusicPDFRef{8874} \MusicPDFRef{8875}] /Limits [<74686D2E352E352E312E33> <74686D2E452E362E342E38>] >>}
\MusicPDFWriteObject{8832}{<FEFF004300680061007000740065007200200031002E00200049006E00740072006F00640075006300740069006F006E>}
\MusicPDFWriteObject{8833}{<< /S /GoTo /D <636861707465722E31> >>}
\MusicPDFWriteObject{8834}{<< /Title \MusicPDFRef{8876} /A \MusicPDFRef{8877} /Parent \MusicPDFRef{1} /Prev \MusicPDFRef{8822} /Next \MusicPDFRef{8878} /First \MusicPDFRef{8879} /Last \MusicPDFRef{8880} /Count -4 >>}
\MusicPDFWriteObject{8835}{<< /Title \MusicPDFRef{8881} /A \MusicPDFRef{8882} /Parent \MusicPDFRef{8822} /Next \MusicPDFRef{8883} >>}
\MusicPDFWriteObject{8836}{<< /Title \MusicPDFRef{8884} /A \MusicPDFRef{8885} /Parent \MusicPDFRef{8822} /Prev \MusicPDFRef{8886} >>}
\MusicPDFWriteObject{8837}{<FEFF0041007000700065006E00640069007800200045002E00200054006F0070006F006C006F0067006900630061006C00200073007400610063006B0073>}
\MusicPDFWriteObject{8838}{<< /S /GoTo /D <617070656E6469782E45> >>}
\MusicPDFWriteObject{8839}{<< /Title \MusicPDFRef{8887} /A \MusicPDFRef{8888} /Parent \MusicPDFRef{1} /Prev \MusicPDFRef{8889} /Next \MusicPDFRef{8825} /First \MusicPDFRef{8890} /Last \MusicPDFRef{8891} /Count -3 >>}
\MusicPDFWriteObject{8840}{<< /Title \MusicPDFRef{8892} /A \MusicPDFRef{8893} /Parent \MusicPDFRef{8825} /Next \MusicPDFRef{8894} >>}
\MusicPDFWriteObject{8841}{<< /Title \MusicPDFRef{8895} /A \MusicPDFRef{8896} /Parent \MusicPDFRef{8825} /Prev \MusicPDFRef{8897} >>}
\MusicPDFWriteObject{8842}{<< /Kids [\MusicPDFRef{8898} \MusicPDFRef{8899} \MusicPDFRef{8900} \MusicPDFRef{8901} \MusicPDFRef{8902} \MusicPDFRef{8903}] /Limits [<446F632D5374617274> <4974656D2E31303330>] >>}
\MusicPDFWriteObject{8843}{<< /Kids [\MusicPDFRef{8904} \MusicPDFRef{8905} \MusicPDFRef{8906} \MusicPDFRef{8907} \MusicPDFRef{8908} \MusicPDFRef{8909}] /Limits [<4974656D2E31303331> <4974656D2E31323235>] >>}
\MusicPDFWriteObject{8844}{<< /Kids [\MusicPDFRef{8910} \MusicPDFRef{8911} \MusicPDFRef{8912} \MusicPDFRef{8913} \MusicPDFRef{8914} \MusicPDFRef{8915}] /Limits [<4974656D2E31323236> <4974656D2E313432>] >>}
\MusicPDFWriteObject{8845}{<< /Kids [\MusicPDFRef{8916} \MusicPDFRef{8917} \MusicPDFRef{8918} \MusicPDFRef{8919} \MusicPDFRef{8920} \MusicPDFRef{8921}] /Limits [<4974656D2E31343230> <4974656D2E313733>] >>}
\MusicPDFWriteObject{8846}{<< /Kids [\MusicPDFRef{8922} \MusicPDFRef{8923} \MusicPDFRef{8924} \MusicPDFRef{8925} \MusicPDFRef{8926} \MusicPDFRef{8927}] /Limits [<4974656D2E313734> <4974656D2E333638>] >>}
\MusicPDFWriteObject{8847}{<< /Kids [\MusicPDFRef{8928} \MusicPDFRef{8929} \MusicPDFRef{8930} \MusicPDFRef{8931} \MusicPDFRef{8932} \MusicPDFRef{8933}] /Limits [<4974656D2E333639> <4974656D2E353632>] >>}
\MusicPDFWriteObject{8848}{<< /Kids [\MusicPDFRef{8934} \MusicPDFRef{8935} \MusicPDFRef{8936} \MusicPDFRef{8937} \MusicPDFRef{8938} \MusicPDFRef{8939}] /Limits [<4974656D2E353633> <4974656D2E373537>] >>}
\MusicPDFWriteObject{8849}{<< /Kids [\MusicPDFRef{8940} \MusicPDFRef{8941} \MusicPDFRef{8942} \MusicPDFRef{8943} \MusicPDFRef{8944} \MusicPDFRef{8945}] /Limits [<4974656D2E373538> <4974656D2E393531>] >>}
\MusicPDFWriteObject{8850}{<< /Kids [\MusicPDFRef{8946} \MusicPDFRef{8947} \MusicPDFRef{8948} \MusicPDFRef{8949} \MusicPDFRef{8950} \MusicPDFRef{8951}] /Limits [<4974656D2E393532> <636974652E4850>] >>}
\MusicPDFWriteObject{8851}{<< /Kids [\MusicPDFRef{8952} \MusicPDFRef{8953} \MusicPDFRef{8954} \MusicPDFRef{8955} \MusicPDFRef{8956} \MusicPDFRef{8957}] /Limits [<636974652E48522D31> <6571756174696F6E2E31302E312E312E3131>] >>}
\MusicPDFWriteObject{8852}{<< /Kids [\MusicPDFRef{8958} \MusicPDFRef{8959} \MusicPDFRef{8960} \MusicPDFRef{8961} \MusicPDFRef{8962} \MusicPDFRef{8963}] /Limits [<6571756174696F6E2E31302E312E312E3132> <6571756174696F6E2E31312E322E312E31>] >>}
\MusicPDFWriteObject{8853}{<< /Kids [\MusicPDFRef{8964} \MusicPDFRef{8965} \MusicPDFRef{8966} \MusicPDFRef{8967} \MusicPDFRef{8968} \MusicPDFRef{8969}] /Limits [<6571756174696F6E2E31312E322E312E3130> <6571756174696F6E2E322E322E332E38>] >>}
\MusicPDFWriteObject{8854}{<< /Kids [\MusicPDFRef{8970} \MusicPDFRef{8971} \MusicPDFRef{8972} \MusicPDFRef{8973} \MusicPDFRef{8974} \MusicPDFRef{8975}] /Limits [<6571756174696F6E2E322E322E332E39> <6571756174696F6E2E332E342E322E34>] >>}
\MusicPDFWriteObject{8855}{<< /Kids [\MusicPDFRef{8976} \MusicPDFRef{8977} \MusicPDFRef{8978} \MusicPDFRef{8979} \MusicPDFRef{8980} \MusicPDFRef{8981}] /Limits [<6571756174696F6E2E332E342E322E35> <6571756174696F6E2E352E332E312E39>] >>}
\MusicPDFWriteObject{8856}{<< /Kids [\MusicPDFRef{8982} \MusicPDFRef{8983} \MusicPDFRef{8984} \MusicPDFRef{8985} \MusicPDFRef{8986} \MusicPDFRef{8987}] /Limits [<6571756174696F6E2E352E332E322E31> <6571756174696F6E2E372E312E332E32>] >>}
\MusicPDFWriteObject{8857}{<< /Kids [\MusicPDFRef{8988} \MusicPDFRef{8989} \MusicPDFRef{8990} \MusicPDFRef{8991} \MusicPDFRef{8992} \MusicPDFRef{8993}] /Limits [<6571756174696F6E2E372E312E342E31> <6571756174696F6E2E382E342E302E3130>] >>}
\MusicPDFWriteObject{8858}{<< /Kids [\MusicPDFRef{8994} \MusicPDFRef{8995} \MusicPDFRef{8996} \MusicPDFRef{8997} \MusicPDFRef{8998} \MusicPDFRef{8999}] /Limits [<6571756174696F6E2E382E342E302E3131> <6571756174696F6E2E412E322E332E31>] >>}
\MusicPDFWriteObject{8859}{<< /Kids [\MusicPDFRef{9000} \MusicPDFRef{9001} \MusicPDFRef{9002} \MusicPDFRef{9003} \MusicPDFRef{9004} \MusicPDFRef{9005}] /Limits [<6571756174696F6E2E412E322E332E3130> <6571756174696F6E2E422E332E322E34>] >>}
\MusicPDFWriteObject{8860}{<< /Kids [\MusicPDFRef{9006} \MusicPDFRef{9007} \MusicPDFRef{9008} \MusicPDFRef{9009} \MusicPDFRef{9010} \MusicPDFRef{9011}] /Limits [<6571756174696F6E2E422E332E322E35> <6571756174696F6E2E432E322E312E32>] >>}
\MusicPDFWriteObject{8861}{<< /Kids [\MusicPDFRef{9012} \MusicPDFRef{9013} \MusicPDFRef{9014} \MusicPDFRef{9015} \MusicPDFRef{9016} \MusicPDFRef{9017}] /Limits [<6571756174696F6E2E432E322E312E33> <6571756174696F6E2E442E312E31302E31>] >>}
\MusicPDFWriteObject{8862}{<< /Kids [\MusicPDFRef{9018} \MusicPDFRef{9019} \MusicPDFRef{9020} \MusicPDFRef{9021} \MusicPDFRef{9022} \MusicPDFRef{9023}] /Limits [<6571756174696F6E2E442E312E31302E3130> <6571756174696F6E2E452E322E312E32>] >>}
\MusicPDFWriteObject{8863}{<< /Kids [\MusicPDFRef{9024} \MusicPDFRef{9025} \MusicPDFRef{9026} \MusicPDFRef{9027} \MusicPDFRef{9028} \MusicPDFRef{9029}] /Limits [<6571756174696F6E2E452E322E312E33> <706167652E313330>] >>}
\MusicPDFWriteObject{8864}{<< /Kids [\MusicPDFRef{9030} \MusicPDFRef{9031} \MusicPDFRef{9032} \MusicPDFRef{9033} \MusicPDFRef{9034} \MusicPDFRef{9035}] /Limits [<706167652E313331> <706167652E333235>] >>}
\MusicPDFWriteObject{8865}{<< /Kids [\MusicPDFRef{9036} \MusicPDFRef{9037} \MusicPDFRef{9038} \MusicPDFRef{9039} \MusicPDFRef{9040} \MusicPDFRef{9041}] /Limits [<706167652E333236> <706167652E3532>] >>}
\MusicPDFWriteObject{8866}{<< /Kids [\MusicPDFRef{9042} \MusicPDFRef{9043} \MusicPDFRef{9044} \MusicPDFRef{9045} \MusicPDFRef{9046} \MusicPDFRef{9047}] /Limits [<706167652E353230> <706167652E373134>] >>}
\MusicPDFWriteObject{8867}{<< /Kids [\MusicPDFRef{9048} \MusicPDFRef{9049} \MusicPDFRef{9050} \MusicPDFRef{9051} \MusicPDFRef{9052} \MusicPDFRef{9053}] /Limits [<706167652E373135> <706167652E393039>] >>}
\MusicPDFWriteObject{8868}{<< /Kids [\MusicPDFRef{9054} \MusicPDFRef{9055} \MusicPDFRef{9056} \MusicPDFRef{9057} \MusicPDFRef{9058} \MusicPDFRef{9059}] /Limits [<706167652E3931> <73756273656374696F6E2E332E342E33>] >>}
\MusicPDFWriteObject{8869}{<< /Kids [\MusicPDFRef{9060} \MusicPDFRef{9061} \MusicPDFRef{9062} \MusicPDFRef{9063} \MusicPDFRef{9064} \MusicPDFRef{9065}] /Limits [<73756273656374696F6E2E342E312E31> <73756273756273656374696F6E2E392E322E322E32>] >>}
\MusicPDFWriteObject{8870}{<< /Kids [\MusicPDFRef{9066} \MusicPDFRef{9067} \MusicPDFRef{9068} \MusicPDFRef{9069} \MusicPDFRef{9070} \MusicPDFRef{9071}] /Limits [<73756273756273656374696F6E2E392E332E322E31> <74686D2E31322E322E322E32>] >>}
\MusicPDFWriteObject{8871}{<< /Kids [\MusicPDFRef{9072} \MusicPDFRef{9073} \MusicPDFRef{9074} \MusicPDFRef{9075} \MusicPDFRef{9076} \MusicPDFRef{9077}] /Limits [<74686D2E31322E322E322E33> <74686D2E352E352E312E32>] >>}
\MusicPDFWriteObject{8872}{<< /Kids [\MusicPDFRef{9078} \MusicPDFRef{9079} \MusicPDFRef{9080} \MusicPDFRef{9081} \MusicPDFRef{9082} \MusicPDFRef{9083}] /Limits [<74686D2E352E352E312E33> <74686D2E412E352E342E32>] >>}
\MusicPDFWriteObject{8873}{<< /Kids [\MusicPDFRef{9084} \MusicPDFRef{9085} \MusicPDFRef{9086} \MusicPDFRef{9087} \MusicPDFRef{9088} \MusicPDFRef{9089}] /Limits [<74686D2E412E352E342E33> <74686D2E432E352E312E33>] >>}
\MusicPDFWriteObject{8874}{<< /Kids [\MusicPDFRef{9090} \MusicPDFRef{9091} \MusicPDFRef{9092} \MusicPDFRef{9093} \MusicPDFRef{9094} \MusicPDFRef{9095}] /Limits [<74686D2E432E352E312E3330> <74686D2E452E322E342E32>] >>}
\MusicPDFWriteObject{8875}{<< /Kids [\MusicPDFRef{9096} \MusicPDFRef{9097} \MusicPDFRef{9098} \MusicPDFRef{9099} \MusicPDFRef{9100}] /Limits [<74686D2E452E322E342E33> <74686D2E452E362E342E38>] >>}
\MusicPDFWriteObject{8876}{<FEFF004300680061007000740065007200200032002E0020004100640065006C006900630020002D00660075006E006300740069006F006E0073>}
\MusicPDFWriteObject{8877}{<< /S /GoTo /D <636861707465722E32> >>}
\MusicPDFWriteObject{8878}{<< /Title \MusicPDFRef{9101} /A \MusicPDFRef{9102} /Parent \MusicPDFRef{1} /Prev \MusicPDFRef{8834} /Next \MusicPDFRef{9103} /First \MusicPDFRef{9104} /Last \MusicPDFRef{9105} /Count -4 >>}
\MusicPDFWriteObject{8879}{<< /Title \MusicPDFRef{9106} /A \MusicPDFRef{9107} /Parent \MusicPDFRef{8834} /Next \MusicPDFRef{9108} >>}
\MusicPDFWriteObject{8880}{<< /Title \MusicPDFRef{9109} /A \MusicPDFRef{9110} /Parent \MusicPDFRef{8834} /Prev \MusicPDFRef{9111} >>}
\MusicPDFWriteObject{8881}{<FEFF0031002E0031002E002000540068006500200070006100720061006400690067006D0020006F0066002000660061006D0069006C006900650073>}
\MusicPDFWriteObject{8882}{<< /S /GoTo /D <73656374696F6E2E312E31> >>}
\MusicPDFWriteObject{8883}{<< /Title \MusicPDFRef{9112} /A \MusicPDFRef{9113} /Parent \MusicPDFRef{8822} /Prev \MusicPDFRef{8835} /Next \MusicPDFRef{9114} >>}
\MusicPDFWriteObject{8884}{<FEFF0031002E0038002E0020004100200076006900650077002000660072006F006D002000740068006500200073006B0079003A002000660072006F006D002000610020006C0065007400740065007200200074006F0020005000650074006500720020005300610072006E0061006B>}
\MusicPDFWriteObject{8885}{<< /S /GoTo /D <73656374696F6E2E312E38> >>}
\MusicPDFWriteObject{8886}{<< /Title \MusicPDFRef{9115} /A \MusicPDFRef{9116} /Parent \MusicPDFRef{8822} /Prev \MusicPDFRef{9117} /Next \MusicPDFRef{8836} >>}
\MusicPDFWriteObject{8887}{<FEFF0041007000700065006E00640069007800200044002E00200053006300680077006100720074007A002D004200720075006800610074002000660075006E006300740069006F006E0073>}
\MusicPDFWriteObject{8888}{<< /S /GoTo /D <617070656E6469782E44> >>}
\MusicPDFWriteObject{8889}{<< /Title \MusicPDFRef{9118} /A \MusicPDFRef{9119} /Parent \MusicPDFRef{1} /Prev \MusicPDFRef{9120} /Next \MusicPDFRef{8839} /First \MusicPDFRef{9121} /Last \MusicPDFRef{9122} /Count -5 >>}
\MusicPDFWriteObject{8890}{<< /Title \MusicPDFRef{9123} /A \MusicPDFRef{9124} /Parent \MusicPDFRef{8839} /Next \MusicPDFRef{9125} >>}
\MusicPDFWriteObject{8891}{<< /Title \MusicPDFRef{9126} /A \MusicPDFRef{9127} /Parent \MusicPDFRef{8839} /Prev \MusicPDFRef{9125} >>}
\MusicPDFWriteObject{8892}{<FEFF0045002E0031002E00200043006F006D0070006100630074006C0079002000670065006E0065007200610074006500640020007300700061006300650073>}
\MusicPDFWriteObject{8893}{<< /S /GoTo /D <73656374696F6E2E452E31> >>}
\MusicPDFWriteObject{8894}{<< /Title \MusicPDFRef{9128} /A \MusicPDFRef{9129} /Parent \MusicPDFRef{8825} /Prev \MusicPDFRef{8840} /Next \MusicPDFRef{9130} >>}
\MusicPDFWriteObject{8895}{<FEFF0045002E0036002E00200053007400610063006B00730020006F006E00200043004700570048002C>}
\MusicPDFWriteObject{8896}{<< /S /GoTo /D <73656374696F6E2E452E36> >>}
\MusicPDFWriteObject{8897}{<< /Title \MusicPDFRef{9131} /A \MusicPDFRef{9132} /Parent \MusicPDFRef{8825} /Prev \MusicPDFRef{9133} /Next \MusicPDFRef{8841} >>}
\MusicPDFWriteObject{8898}{<< /Kids [\MusicPDFRef{9134} \MusicPDFRef{9135} \MusicPDFRef{9136} \MusicPDFRef{9137} \MusicPDFRef{9138} \MusicPDFRef{9139}] /Limits [<446F632D5374617274> <48666F6F746E6F74652E3133>] >>}
\MusicPDFWriteObject{8899}{<< /Kids [\MusicPDFRef{9140} \MusicPDFRef{9141} \MusicPDFRef{9142} \MusicPDFRef{9143} \MusicPDFRef{9144} \MusicPDFRef{9145}] /Limits [<48666F6F746E6F74652E313330> <48666F6F746E6F74652E313632>] >>}
\MusicPDFWriteObject{8900}{<< /Kids [\MusicPDFRef{9146} \MusicPDFRef{9147} \MusicPDFRef{9148} \MusicPDFRef{9149} \MusicPDFRef{9150} \MusicPDFRef{9151}] /Limits [<48666F6F746E6F74652E313633> <48666F6F746E6F74652E3334>] >>}
\MusicPDFWriteObject{8901}{<< /Kids [\MusicPDFRef{9152} \MusicPDFRef{9153} \MusicPDFRef{9154} \MusicPDFRef{9155} \MusicPDFRef{9156} \MusicPDFRef{9157}] /Limits [<48666F6F746E6F74652E3335> <48666F6F746E6F74652E3637>] >>}
\MusicPDFWriteObject{8902}{<< /Kids [\MusicPDFRef{9158} \MusicPDFRef{9159} \MusicPDFRef{9160} \MusicPDFRef{9161} \MusicPDFRef{9162} \MusicPDFRef{9163}] /Limits [<48666F6F746E6F74652E3638> <4974656D2E31>] >>}
\MusicPDFWriteObject{8903}{<< /Kids [\MusicPDFRef{9164} \MusicPDFRef{9165} \MusicPDFRef{9166} \MusicPDFRef{9167} \MusicPDFRef{9168} \MusicPDFRef{9169}] /Limits [<4974656D2E3130> <4974656D2E31303330>] >>}
\MusicPDFWriteObject{8904}{<< /Kids [\MusicPDFRef{9170} \MusicPDFRef{9171} \MusicPDFRef{9172} \MusicPDFRef{9173} \MusicPDFRef{9174} \MusicPDFRef{9175}] /Limits [<4974656D2E31303331> <4974656D2E31303633>] >>}
\MusicPDFWriteObject{8905}{<< /Kids [\MusicPDFRef{9176} \MusicPDFRef{9177} \MusicPDFRef{9178} \MusicPDFRef{9179} \MusicPDFRef{9180} \MusicPDFRef{9181}] /Limits [<4974656D2E31303634> <4974656D2E31303936>] >>}
\MusicPDFWriteObject{8906}{<< /Kids [\MusicPDFRef{9182} \MusicPDFRef{9183} \MusicPDFRef{9184} \MusicPDFRef{9185} \MusicPDFRef{9186} \MusicPDFRef{9187}] /Limits [<4974656D2E31303937> <4974656D2E31313238>] >>}
\MusicPDFWriteObject{8907}{<< /Kids [\MusicPDFRef{9188} \MusicPDFRef{9189} \MusicPDFRef{9190} \MusicPDFRef{9191} \MusicPDFRef{9192} \MusicPDFRef{9193}] /Limits [<4974656D2E31313239> <4974656D2E31313630>] >>}
\MusicPDFWriteObject{8908}{<< /Kids [\MusicPDFRef{9194} \MusicPDFRef{9195} \MusicPDFRef{9196} \MusicPDFRef{9197} \MusicPDFRef{9198} \MusicPDFRef{9199}] /Limits [<4974656D2E31313631> <4974656D2E31313933>] >>}
\MusicPDFWriteObject{8909}{<< /Kids [\MusicPDFRef{9200} \MusicPDFRef{9201} \MusicPDFRef{9202} \MusicPDFRef{9203} \MusicPDFRef{9204} \MusicPDFRef{9205}] /Limits [<4974656D2E31313934> <4974656D2E31323235>] >>}
\MusicPDFWriteObject{8910}{<< /Kids [\MusicPDFRef{9206} \MusicPDFRef{9207} \MusicPDFRef{9208} \MusicPDFRef{9209} \MusicPDFRef{9210} \MusicPDFRef{9211}] /Limits [<4974656D2E31323236> <4974656D2E31323538>] >>}
\MusicPDFWriteObject{8911}{<< /Kids [\MusicPDFRef{9212} \MusicPDFRef{9213} \MusicPDFRef{9214} \MusicPDFRef{9215} \MusicPDFRef{9216} \MusicPDFRef{9217}] /Limits [<4974656D2E31323539> <4974656D2E31323930>] >>}
\MusicPDFWriteObject{8912}{<< /Kids [\MusicPDFRef{9218} \MusicPDFRef{9219} \MusicPDFRef{9220} \MusicPDFRef{9221} \MusicPDFRef{9222} \MusicPDFRef{9223}] /Limits [<4974656D2E31323931> <4974656D2E31333232>] >>}
\MusicPDFWriteObject{8913}{<< /Kids [\MusicPDFRef{9224} \MusicPDFRef{9225} \MusicPDFRef{9226} \MusicPDFRef{9227} \MusicPDFRef{9228} \MusicPDFRef{9229}] /Limits [<4974656D2E31333233> <4974656D2E31333535>] >>}
\MusicPDFWriteObject{8914}{<< /Kids [\MusicPDFRef{9230} \MusicPDFRef{9231} \MusicPDFRef{9232} \MusicPDFRef{9233} \MusicPDFRef{9234} \MusicPDFRef{9235}] /Limits [<4974656D2E31333536> <4974656D2E31333838>] >>}
\MusicPDFWriteObject{8915}{<< /Kids [\MusicPDFRef{9236} \MusicPDFRef{9237} \MusicPDFRef{9238} \MusicPDFRef{9239} \MusicPDFRef{9240} \MusicPDFRef{9241}] /Limits [<4974656D2E31333839> <4974656D2E313432>] >>}
\MusicPDFWriteObject{8916}{<< /Kids [\MusicPDFRef{9242} \MusicPDFRef{9243} \MusicPDFRef{9244} \MusicPDFRef{9245} \MusicPDFRef{9246} \MusicPDFRef{9247}] /Limits [<4974656D2E31343230> <4974656D2E31343532>] >>}
\MusicPDFWriteObject{8917}{<< /Kids [\MusicPDFRef{9248} \MusicPDFRef{9249} \MusicPDFRef{9250} \MusicPDFRef{9251} \MusicPDFRef{9252} \MusicPDFRef{9253}] /Limits [<4974656D2E31343533> <4974656D2E31343835>] >>}
\MusicPDFWriteObject{8918}{<< /Kids [\MusicPDFRef{9254} \MusicPDFRef{9255} \MusicPDFRef{9256} \MusicPDFRef{9257} \MusicPDFRef{9258} \MusicPDFRef{9259}] /Limits [<4974656D2E31343836> <4974656D2E31353137>] >>}
\MusicPDFWriteObject{8919}{<< /Kids [\MusicPDFRef{9260} \MusicPDFRef{9261} \MusicPDFRef{9262} \MusicPDFRef{9263} \MusicPDFRef{9264} \MusicPDFRef{9265}] /Limits [<4974656D2E31353138> <4974656D2E313535>] >>}
\MusicPDFWriteObject{8920}{<< /Kids [\MusicPDFRef{9266} \MusicPDFRef{9267} \MusicPDFRef{9268} \MusicPDFRef{9269} \MusicPDFRef{9270} \MusicPDFRef{9271}] /Limits [<4974656D2E31353530> <4974656D2E31353832>] >>}
\MusicPDFWriteObject{8921}{<< /Kids [\MusicPDFRef{9272} \MusicPDFRef{9273} \MusicPDFRef{9274} \MusicPDFRef{9275} \MusicPDFRef{9276} \MusicPDFRef{9277}] /Limits [<4974656D2E31353833> <4974656D2E313733>] >>}
\MusicPDFWriteObject{8922}{<< /Kids [\MusicPDFRef{9278} \MusicPDFRef{9279} \MusicPDFRef{9280} \MusicPDFRef{9281} \MusicPDFRef{9282} \MusicPDFRef{9283}] /Limits [<4974656D2E313734> <4974656D2E323035>] >>}
\MusicPDFWriteObject{8923}{<< /Kids [\MusicPDFRef{9284} \MusicPDFRef{9285} \MusicPDFRef{9286} \MusicPDFRef{9287} \MusicPDFRef{9288} \MusicPDFRef{9289}] /Limits [<4974656D2E323036> <4974656D2E323338>] >>}
\MusicPDFWriteObject{8924}{<< /Kids [\MusicPDFRef{9290} \MusicPDFRef{9291} \MusicPDFRef{9292} \MusicPDFRef{9293} \MusicPDFRef{9294} \MusicPDFRef{9295}] /Limits [<4974656D2E323339> <4974656D2E323730>] >>}
\MusicPDFWriteObject{8925}{<< /Kids [\MusicPDFRef{9296} \MusicPDFRef{9297} \MusicPDFRef{9298} \MusicPDFRef{9299} \MusicPDFRef{9300} \MusicPDFRef{9301}] /Limits [<4974656D2E323731> <4974656D2E333032>] >>}
\MusicPDFWriteObject{8926}{<< /Kids [\MusicPDFRef{9302} \MusicPDFRef{9303} \MusicPDFRef{9304} \MusicPDFRef{9305} \MusicPDFRef{9306} \MusicPDFRef{9307}] /Limits [<4974656D2E333033> <4974656D2E333335>] >>}
\MusicPDFWriteObject{8927}{<< /Kids [\MusicPDFRef{9308} \MusicPDFRef{9309} \MusicPDFRef{9310} \MusicPDFRef{9311} \MusicPDFRef{9312} \MusicPDFRef{9313}] /Limits [<4974656D2E333336> <4974656D2E333638>] >>}
\MusicPDFWriteObject{8928}{<< /Kids [\MusicPDFRef{9314} \MusicPDFRef{9315} \MusicPDFRef{9316} \MusicPDFRef{9317} \MusicPDFRef{9318} \MusicPDFRef{9319}] /Limits [<4974656D2E333639> <4974656D2E3430>] >>}
\MusicPDFWriteObject{8929}{<< /Kids [\MusicPDFRef{9320} \MusicPDFRef{9321} \MusicPDFRef{9322} \MusicPDFRef{9323} \MusicPDFRef{9324} \MusicPDFRef{9325}] /Limits [<4974656D2E343030> <4974656D2E343332>] >>}
\MusicPDFWriteObject{8930}{<< /Kids [\MusicPDFRef{9326} \MusicPDFRef{9327} \MusicPDFRef{9328} \MusicPDFRef{9329} \MusicPDFRef{9330} \MusicPDFRef{9331}] /Limits [<4974656D2E343333> <4974656D2E343635>] >>}
\MusicPDFWriteObject{8931}{<< /Kids [\MusicPDFRef{9332} \MusicPDFRef{9333} \MusicPDFRef{9334} \MusicPDFRef{9335} \MusicPDFRef{9336} \MusicPDFRef{9337}] /Limits [<4974656D2E343636> <4974656D2E343938>] >>}
\MusicPDFWriteObject{8932}{<< /Kids [\MusicPDFRef{9338} \MusicPDFRef{9339} \MusicPDFRef{9340} \MusicPDFRef{9341} \MusicPDFRef{9342} \MusicPDFRef{9343}] /Limits [<4974656D2E343939> <4974656D2E3533>] >>}
\MusicPDFWriteObject{8933}{<< /Kids [\MusicPDFRef{9344} \MusicPDFRef{9345} \MusicPDFRef{9346} \MusicPDFRef{9347} \MusicPDFRef{9348} \MusicPDFRef{9349}] /Limits [<4974656D2E353330> <4974656D2E353632>] >>}
\MusicPDFWriteObject{8934}{<< /Kids [\MusicPDFRef{9350} \MusicPDFRef{9351} \MusicPDFRef{9352} \MusicPDFRef{9353} \MusicPDFRef{9354} \MusicPDFRef{9355}] /Limits [<4974656D2E353633> <4974656D2E353935>] >>}
\MusicPDFWriteObject{8935}{<< /Kids [\MusicPDFRef{9356} \MusicPDFRef{9357} \MusicPDFRef{9358} \MusicPDFRef{9359} \MusicPDFRef{9360} \MusicPDFRef{9361}] /Limits [<4974656D2E353936> <4974656D2E363237>] >>}
\MusicPDFWriteObject{8936}{<< /Kids [\MusicPDFRef{9362} \MusicPDFRef{9363} \MusicPDFRef{9364} \MusicPDFRef{9365} \MusicPDFRef{9366} \MusicPDFRef{9367}] /Limits [<4974656D2E363238> <4974656D2E3636>] >>}
\MusicPDFWriteObject{8937}{<< /Kids [\MusicPDFRef{9368} \MusicPDFRef{9369} \MusicPDFRef{9370} \MusicPDFRef{9371} \MusicPDFRef{9372} \MusicPDFRef{9373}] /Limits [<4974656D2E363630> <4974656D2E363932>] >>}
\MusicPDFWriteObject{8938}{<< /Kids [\MusicPDFRef{9374} \MusicPDFRef{9375} \MusicPDFRef{9376} \MusicPDFRef{9377} \MusicPDFRef{9378} \MusicPDFRef{9379}] /Limits [<4974656D2E363933> <4974656D2E373234>] >>}
\MusicPDFWriteObject{8939}{<< /Kids [\MusicPDFRef{9380} \MusicPDFRef{9381} \MusicPDFRef{9382} \MusicPDFRef{9383} \MusicPDFRef{9384} \MusicPDFRef{9385}] /Limits [<4974656D2E373235> <4974656D2E373537>] >>}
\MusicPDFWriteObject{8940}{<< /Kids [\MusicPDFRef{9386} \MusicPDFRef{9387} \MusicPDFRef{9388} \MusicPDFRef{9389} \MusicPDFRef{9390} \MusicPDFRef{9391}] /Limits [<4974656D2E373538> <4974656D2E3739>] >>}
\MusicPDFWriteObject{8941}{<< /Kids [\MusicPDFRef{9392} \MusicPDFRef{9393} \MusicPDFRef{9394} \MusicPDFRef{9395} \MusicPDFRef{9396} \MusicPDFRef{9397}] /Limits [<4974656D2E373930> <4974656D2E383231>] >>}
\MusicPDFWriteObject{8942}{<< /Kids [\MusicPDFRef{9398} \MusicPDFRef{9399} \MusicPDFRef{9400} \MusicPDFRef{9401} \MusicPDFRef{9402} \MusicPDFRef{9403}] /Limits [<4974656D2E383232> <4974656D2E383534>] >>}
\MusicPDFWriteObject{8943}{<< /Kids [\MusicPDFRef{9404} \MusicPDFRef{9405} \MusicPDFRef{9406} \MusicPDFRef{9407} \MusicPDFRef{9408} \MusicPDFRef{9409}] /Limits [<4974656D2E383535> <4974656D2E383837>] >>}
\MusicPDFWriteObject{8944}{<< /Kids [\MusicPDFRef{9410} \MusicPDFRef{9411} \MusicPDFRef{9412} \MusicPDFRef{9413} \MusicPDFRef{9414} \MusicPDFRef{9415}] /Limits [<4974656D2E383838> <4974656D2E393139>] >>}
\MusicPDFWriteObject{8945}{<< /Kids [\MusicPDFRef{9416} \MusicPDFRef{9417} \MusicPDFRef{9418} \MusicPDFRef{9419} \MusicPDFRef{9420} \MusicPDFRef{9421}] /Limits [<4974656D2E3932> <4974656D2E393531>] >>}
\MusicPDFWriteObject{8946}{<< /Kids [\MusicPDFRef{9422} \MusicPDFRef{9423} \MusicPDFRef{9424} \MusicPDFRef{9425} \MusicPDFRef{9426} \MusicPDFRef{9427}] /Limits [<4974656D2E393532> <4974656D2E393834>] >>}
\MusicPDFWriteObject{8947}{<< /Kids [\MusicPDFRef{9428} \MusicPDFRef{9429} \MusicPDFRef{9430} \MusicPDFRef{9431} \MusicPDFRef{9432} \MusicPDFRef{9433}] /Limits [<4974656D2E393835> <636861707465722E38>] >>}
\MusicPDFWriteObject{8948}{<< /Kids [\MusicPDFRef{9434} \MusicPDFRef{9435} \MusicPDFRef{9436} \MusicPDFRef{9437} \MusicPDFRef{9438} \MusicPDFRef{9439}] /Limits [<636861707465722E39> <636974652E4269726B686F6666>] >>}
\MusicPDFWriteObject{8949}{<< /Kids [\MusicPDFRef{9440} \MusicPDFRef{9441} \MusicPDFRef{9442} \MusicPDFRef{9443} \MusicPDFRef{9444} \MusicPDFRef{9445}] /Limits [<636974652E426C> <636974652E436F>] >>}
\MusicPDFWriteObject{8950}{<< /Kids [\MusicPDFRef{9446} \MusicPDFRef{9447} \MusicPDFRef{9448} \MusicPDFRef{9449} \MusicPDFRef{9450} \MusicPDFRef{9451}] /Limits [<636974652E436F68> <636974652E464B>] >>}
\MusicPDFWriteObject{8951}{<< /Kids [\MusicPDFRef{9452} \MusicPDFRef{9453} \MusicPDFRef{9454} \MusicPDFRef{9455} \MusicPDFRef{9456} \MusicPDFRef{9457}] /Limits [<636974652E464D> <636974652E4850>] >>}
\MusicPDFWriteObject{8952}{<< /Kids [\MusicPDFRef{9458} \MusicPDFRef{9459} \MusicPDFRef{9460} \MusicPDFRef{9461} \MusicPDFRef{9462} \MusicPDFRef{9463}] /Limits [<636974652E48522D31> <636974652E4B656368>] >>}
\MusicPDFWriteObject{8953}{<< /Kids [\MusicPDFRef{9464} \MusicPDFRef{9465} \MusicPDFRef{9466} \MusicPDFRef{9467} \MusicPDFRef{9468} \MusicPDFRef{9469}] /Limits [<636974652E4B656C> <636974652E4D5665>] >>}
\MusicPDFWriteObject{8954}{<< /Kids [\MusicPDFRef{9470} \MusicPDFRef{9471} \MusicPDFRef{9472} \MusicPDFRef{9473} \MusicPDFRef{9474} \MusicPDFRef{9475}] /Limits [<636974652E4D5A> <636974652E52657A6B>] >>}
\MusicPDFWriteObject{8955}{<< /Kids [\MusicPDFRef{9476} \MusicPDFRef{9477} \MusicPDFRef{9478} \MusicPDFRef{9479} \MusicPDFRef{9480} \MusicPDFRef{9481}] /Limits [<636974652E5269656D616E6E> <636974652E54617465>] >>}
\MusicPDFWriteObject{8956}{<< /Kids [\MusicPDFRef{9482} \MusicPDFRef{9483} \MusicPDFRef{9484} \MusicPDFRef{9485} \MusicPDFRef{9486} \MusicPDFRef{9487}] /Limits [<636974652E547265> <6571756174696F6E2E312E332E312E3132>] >>}
\MusicPDFWriteObject{8957}{<< /Kids [\MusicPDFRef{9488} \MusicPDFRef{9489} \MusicPDFRef{9490} \MusicPDFRef{9491} \MusicPDFRef{9492} \MusicPDFRef{9493}] /Limits [<6571756174696F6E2E312E332E312E32> <6571756174696F6E2E31302E312E312E3131>] >>}
\MusicPDFWriteObject{8958}{<< /Kids [\MusicPDFRef{9494} \MusicPDFRef{9495} \MusicPDFRef{9496} \MusicPDFRef{9497} \MusicPDFRef{9498} \MusicPDFRef{9499}] /Limits [<6571756174696F6E2E31302E312E312E3132> <6571756174696F6E2E31302E322E322E35>] >>}
\MusicPDFWriteObject{8959}{<< /Kids [\MusicPDFRef{9500} \MusicPDFRef{9501} \MusicPDFRef{9502} \MusicPDFRef{9503} \MusicPDFRef{9504} \MusicPDFRef{9505}] /Limits [<6571756174696F6E2E31302E322E322E36> <6571756174696F6E2E31302E332E322E34>] >>}
\MusicPDFWriteObject{8960}{<< /Kids [\MusicPDFRef{9506} \MusicPDFRef{9507} \MusicPDFRef{9508} \MusicPDFRef{9509} \MusicPDFRef{9510} \MusicPDFRef{9511}] /Limits [<6571756174696F6E2E31302E332E322E35> <6571756174696F6E2E31302E342E322E3235>] >>}
\MusicPDFWriteObject{8961}{<< /Kids [\MusicPDFRef{9512} \MusicPDFRef{9513} \MusicPDFRef{9514} \MusicPDFRef{9515} \MusicPDFRef{9516} \MusicPDFRef{9517}] /Limits [<6571756174696F6E2E31302E342E322E3236> <6571756174696F6E2E31302E342E332E3237>] >>}
\MusicPDFWriteObject{8962}{<< /Kids [\MusicPDFRef{9518} \MusicPDFRef{9519} \MusicPDFRef{9520} \MusicPDFRef{9521} \MusicPDFRef{9522} \MusicPDFRef{9523}] /Limits [<6571756174696F6E2E31302E342E332E3238> <6571756174696F6E2E31312E312E312E3134>] >>}
\MusicPDFWriteObject{8963}{<< /Kids [\MusicPDFRef{9524} \MusicPDFRef{9525} \MusicPDFRef{9526} \MusicPDFRef{9527} \MusicPDFRef{9528} \MusicPDFRef{9529}] /Limits [<6571756174696F6E2E31312E312E312E3135> <6571756174696F6E2E31312E322E312E31>] >>}
\MusicPDFWriteObject{8964}{<< /Kids [\MusicPDFRef{9530} \MusicPDFRef{9531} \MusicPDFRef{9532} \MusicPDFRef{9533} \MusicPDFRef{9534} \MusicPDFRef{9535}] /Limits [<6571756174696F6E2E31312E322E312E3130> <6571756174696F6E2E31322E312E312E3133>] >>}
\MusicPDFWriteObject{8965}{<< /Kids [\MusicPDFRef{9536} \MusicPDFRef{9537} \MusicPDFRef{9538} \MusicPDFRef{9539} \MusicPDFRef{9540} \MusicPDFRef{9541}] /Limits [<6571756174696F6E2E31322E312E312E3134> <6571756174696F6E2E31322E322E322E3132>] >>}
\MusicPDFWriteObject{8966}{<< /Kids [\MusicPDFRef{9542} \MusicPDFRef{9543} \MusicPDFRef{9544} \MusicPDFRef{9545} \MusicPDFRef{9546} \MusicPDFRef{9547}] /Limits [<6571756174696F6E2E31322E322E322E3133> <6571756174696F6E2E31322E342E312E3137>] >>}
\MusicPDFWriteObject{8967}{<< /Kids [\MusicPDFRef{9548} \MusicPDFRef{9549} \MusicPDFRef{9550} \MusicPDFRef{9551} \MusicPDFRef{9552} \MusicPDFRef{9553}] /Limits [<6571756174696F6E2E31322E342E312E3138> <6571756174696F6E2E31342E312E312E34>] >>}
\MusicPDFWriteObject{8968}{<< /Kids [\MusicPDFRef{9554} \MusicPDFRef{9555} \MusicPDFRef{9556} \MusicPDFRef{9557} \MusicPDFRef{9558} \MusicPDFRef{9559}] /Limits [<6571756174696F6E2E31342E312E312E35> <6571756174696F6E2E322E312E342E36>] >>}
\MusicPDFWriteObject{8969}{<< /Kids [\MusicPDFRef{9560} \MusicPDFRef{9561} \MusicPDFRef{9562} \MusicPDFRef{9563} \MusicPDFRef{9564} \MusicPDFRef{9565}] /Limits [<6571756174696F6E2E322E312E342E37> <6571756174696F6E2E322E322E332E38>] >>}
\MusicPDFWriteObject{8970}{<< /Kids [\MusicPDFRef{9566} \MusicPDFRef{9567} \MusicPDFRef{9568} \MusicPDFRef{9569} \MusicPDFRef{9570} \MusicPDFRef{9571}] /Limits [<6571756174696F6E2E322E322E332E39> <6571756174696F6E2E322E322E342E3430>] >>}
\MusicPDFWriteObject{8971}{<< /Kids [\MusicPDFRef{9572} \MusicPDFRef{9573} \MusicPDFRef{9574} \MusicPDFRef{9575} \MusicPDFRef{9576} \MusicPDFRef{9577}] /Limits [<6571756174696F6E2E322E322E342E3431> <6571756174696F6E2E322E332E312E39>] >>}
\MusicPDFWriteObject{8972}{<< /Kids [\MusicPDFRef{9578} \MusicPDFRef{9579} \MusicPDFRef{9580} \MusicPDFRef{9581} \MusicPDFRef{9582} \MusicPDFRef{9583}] /Limits [<6571756174696F6E2E322E332E322E31> <6571756174696F6E2E332E312E322E31>] >>}
\MusicPDFWriteObject{8973}{<< /Kids [\MusicPDFRef{9584} \MusicPDFRef{9585} \MusicPDFRef{9586} \MusicPDFRef{9587} \MusicPDFRef{9588} \MusicPDFRef{9589}] /Limits [<6571756174696F6E2E332E312E322E32> <6571756174696F6E2E332E332E312E3135>] >>}
\MusicPDFWriteObject{8974}{<< /Kids [\MusicPDFRef{9590} \MusicPDFRef{9591} \MusicPDFRef{9592} \MusicPDFRef{9593} \MusicPDFRef{9594} \MusicPDFRef{9595}] /Limits [<6571756174696F6E2E332E332E312E3136> <6571756174696F6E2E332E332E322E3137>] >>}
\MusicPDFWriteObject{8975}{<< /Kids [\MusicPDFRef{9596} \MusicPDFRef{9597} \MusicPDFRef{9598} \MusicPDFRef{9599} \MusicPDFRef{9600} \MusicPDFRef{9601}] /Limits [<6571756174696F6E2E332E332E322E3138> <6571756174696F6E2E332E342E322E34>] >>}
\MusicPDFWriteObject{8976}{<< /Kids [\MusicPDFRef{9602} \MusicPDFRef{9603} \MusicPDFRef{9604} \MusicPDFRef{9605} \MusicPDFRef{9606} \MusicPDFRef{9607}] /Limits [<6571756174696F6E2E332E342E322E35> <6571756174696F6E2E342E322E312E3137>] >>}
\MusicPDFWriteObject{8977}{<< /Kids [\MusicPDFRef{9608} \MusicPDFRef{9609} \MusicPDFRef{9610} \MusicPDFRef{9611} \MusicPDFRef{9612} \MusicPDFRef{9613}] /Limits [<6571756174696F6E2E342E322E312E3138> <6571756174696F6E2E342E332E322E38>] >>}
\MusicPDFWriteObject{8978}{<< /Kids [\MusicPDFRef{9614} \MusicPDFRef{9615} \MusicPDFRef{9616} \MusicPDFRef{9617} \MusicPDFRef{9618} \MusicPDFRef{9619}] /Limits [<6571756174696F6E2E342E332E332E31> <6571756174696F6E2E342E352E312E3131>] >>}
\MusicPDFWriteObject{8979}{<< /Kids [\MusicPDFRef{9620} \MusicPDFRef{9621} \MusicPDFRef{9622} \MusicPDFRef{9623} \MusicPDFRef{9624} \MusicPDFRef{9625}] /Limits [<6571756174696F6E2E342E352E312E3132> <6571756174696F6E2E342E362E312E3236>] >>}
\MusicPDFWriteObject{8980}{<< /Kids [\MusicPDFRef{9626} \MusicPDFRef{9627} \MusicPDFRef{9628} \MusicPDFRef{9629} \MusicPDFRef{9630} \MusicPDFRef{9631}] /Limits [<6571756174696F6E2E342E362E312E3237> <6571756174696F6E2E352E322E312E31>] >>}
\MusicPDFWriteObject{8981}{<< /Kids [\MusicPDFRef{9632} \MusicPDFRef{9633} \MusicPDFRef{9634} \MusicPDFRef{9635} \MusicPDFRef{9636} \MusicPDFRef{9637}] /Limits [<6571756174696F6E2E352E322E312E32> <6571756174696F6E2E352E332E312E39>] >>}
\MusicPDFWriteObject{8982}{<< /Kids [\MusicPDFRef{9638} \MusicPDFRef{9639} \MusicPDFRef{9640} \MusicPDFRef{9641} \MusicPDFRef{9642} \MusicPDFRef{9643}] /Limits [<6571756174696F6E2E352E332E322E31> <6571756174696F6E2E352E332E332E3234>] >>}
\MusicPDFWriteObject{8983}{<< /Kids [\MusicPDFRef{9644} \MusicPDFRef{9645} \MusicPDFRef{9646} \MusicPDFRef{9647} \MusicPDFRef{9648} \MusicPDFRef{9649}] /Limits [<6571756174696F6E2E352E332E332E3235> <6571756174696F6E2E352E332E342E37>] >>}
\MusicPDFWriteObject{8984}{<< /Kids [\MusicPDFRef{9650} \MusicPDFRef{9651} \MusicPDFRef{9652} \MusicPDFRef{9653} \MusicPDFRef{9654} \MusicPDFRef{9655}] /Limits [<6571756174696F6E2E352E332E342E38> <6571756174696F6E2E352E352E312E3132>] >>}
\MusicPDFWriteObject{8985}{<< /Kids [\MusicPDFRef{9656} \MusicPDFRef{9657} \MusicPDFRef{9658} \MusicPDFRef{9659} \MusicPDFRef{9660} \MusicPDFRef{9661}] /Limits [<6571756174696F6E2E352E352E312E3133> <6571756174696F6E2E362E312E322E3139>] >>}
\MusicPDFWriteObject{8986}{<< /Kids [\MusicPDFRef{9662} \MusicPDFRef{9663} \MusicPDFRef{9664} \MusicPDFRef{9665} \MusicPDFRef{9666} \MusicPDFRef{9667}] /Limits [<6571756174696F6E2E362E312E322E32> <6571756174696F6E2E362E332E322E3133>] >>}
\MusicPDFWriteObject{8987}{<< /Kids [\MusicPDFRef{9668} \MusicPDFRef{9669} \MusicPDFRef{9670} \MusicPDFRef{9671} \MusicPDFRef{9672} \MusicPDFRef{9673}] /Limits [<6571756174696F6E2E362E332E322E3134> <6571756174696F6E2E372E312E332E32>] >>}
\MusicPDFWriteObject{8988}{<< /Kids [\MusicPDFRef{9674} \MusicPDFRef{9675} \MusicPDFRef{9676} \MusicPDFRef{9677} \MusicPDFRef{9678} \MusicPDFRef{9679}] /Limits [<6571756174696F6E2E372E312E342E31> <6571756174696F6E2E372E322E312E3134>] >>}
\MusicPDFWriteObject{8989}{<< /Kids [\MusicPDFRef{9680} \MusicPDFRef{9681} \MusicPDFRef{9682} \MusicPDFRef{9683} \MusicPDFRef{9684} \MusicPDFRef{9685}] /Limits [<6571756174696F6E2E372E322E312E3135> <6571756174696F6E2E372E332E312E32>] >>}
\MusicPDFWriteObject{8990}{<< /Kids [\MusicPDFRef{9686} \MusicPDFRef{9687} \MusicPDFRef{9688} \MusicPDFRef{9689} \MusicPDFRef{9690} \MusicPDFRef{9691}] /Limits [<6571756174696F6E2E372E332E312E3230> <6571756174696F6E2E372E332E322E39>] >>}
\MusicPDFWriteObject{8991}{<< /Kids [\MusicPDFRef{9692} \MusicPDFRef{9693} \MusicPDFRef{9694} \MusicPDFRef{9695} \MusicPDFRef{9696} \MusicPDFRef{9697}] /Limits [<6571756174696F6E2E372E332E332E31> <6571756174696F6E2E372E332E342E36>] >>}
\MusicPDFWriteObject{8992}{<< /Kids [\MusicPDFRef{9698} \MusicPDFRef{9699} \MusicPDFRef{9700} \MusicPDFRef{9701} \MusicPDFRef{9702} \MusicPDFRef{9703}] /Limits [<6571756174696F6E2E372E332E342E37> <6571756174696F6E2E372E352E302E33>] >>}
\MusicPDFWriteObject{8993}{<< /Kids [\MusicPDFRef{9704} \MusicPDFRef{9705} \MusicPDFRef{9706} \MusicPDFRef{9707} \MusicPDFRef{9708} \MusicPDFRef{9709}] /Limits [<6571756174696F6E2E372E352E302E34> <6571756174696F6E2E382E342E302E3130>] >>}
\MusicPDFWriteObject{8994}{<< /Kids [\MusicPDFRef{9710} \MusicPDFRef{9711} \MusicPDFRef{9712} \MusicPDFRef{9713} \MusicPDFRef{9714} \MusicPDFRef{9715}] /Limits [<6571756174696F6E2E382E342E302E3131> <6571756174696F6E2E392E312E322E3135>] >>}
\MusicPDFWriteObject{8995}{<< /Kids [\MusicPDFRef{9716} \MusicPDFRef{9717} \MusicPDFRef{9718} \MusicPDFRef{9719} \MusicPDFRef{9720} \MusicPDFRef{9721}] /Limits [<6571756174696F6E2E392E312E322E3136> <6571756174696F6E2E392E322E312E3331>] >>}
\MusicPDFWriteObject{8996}{<< /Kids [\MusicPDFRef{9722} \MusicPDFRef{9723} \MusicPDFRef{9724} \MusicPDFRef{9725} \MusicPDFRef{9726} \MusicPDFRef{9727}] /Limits [<6571756174696F6E2E392E322E312E3332> <6571756174696F6E2E392E322E322E3331>] >>}
\MusicPDFWriteObject{8997}{<< /Kids [\MusicPDFRef{9728} \MusicPDFRef{9729} \MusicPDFRef{9730} \MusicPDFRef{9731} \MusicPDFRef{9732} \MusicPDFRef{9733}] /Limits [<6571756174696F6E2E392E322E322E3332> <6571756174696F6E2E392E332E322E36>] >>}
\MusicPDFWriteObject{8998}{<< /Kids [\MusicPDFRef{9734} \MusicPDFRef{9735} \MusicPDFRef{9736} \MusicPDFRef{9737} \MusicPDFRef{9738} \MusicPDFRef{9739}] /Limits [<6571756174696F6E2E392E332E322E37> <6571756174696F6E2E412E312E332E32>] >>}
\MusicPDFWriteObject{8999}{<< /Kids [\MusicPDFRef{9740} \MusicPDFRef{9741} \MusicPDFRef{9742} \MusicPDFRef{9743} \MusicPDFRef{9744} \MusicPDFRef{9745}] /Limits [<6571756174696F6E2E412E312E332E33> <6571756174696F6E2E412E322E332E31>] >>}
\MusicPDFWriteObject{9000}{<< /Kids [\MusicPDFRef{9746} \MusicPDFRef{9747} \MusicPDFRef{9748} \MusicPDFRef{9749} \MusicPDFRef{9750} \MusicPDFRef{9751}] /Limits [<6571756174696F6E2E412E322E332E3130> <6571756174696F6E2E412E322E342E36>] >>}
\MusicPDFWriteObject{9001}{<< /Kids [\MusicPDFRef{9752} \MusicPDFRef{9753} \MusicPDFRef{9754} \MusicPDFRef{9755} \MusicPDFRef{9756} \MusicPDFRef{9757}] /Limits [<6571756174696F6E2E412E322E342E37> <6571756174696F6E2E412E352E312E3132>] >>}
\MusicPDFWriteObject{9002}{<< /Kids [\MusicPDFRef{9758} \MusicPDFRef{9759} \MusicPDFRef{9760} \MusicPDFRef{9761} \MusicPDFRef{9762} \MusicPDFRef{9763}] /Limits [<6571756174696F6E2E412E352E312E3133> <6571756174696F6E2E412E352E332E3130>] >>}
\MusicPDFWriteObject{9003}{<< /Kids [\MusicPDFRef{9764} \MusicPDFRef{9765} \MusicPDFRef{9766} \MusicPDFRef{9767} \MusicPDFRef{9768} \MusicPDFRef{9769}] /Limits [<6571756174696F6E2E412E352E332E3131> <6571756174696F6E2E422E312E312E3131>] >>}
\MusicPDFWriteObject{9004}{<< /Kids [\MusicPDFRef{9770} \MusicPDFRef{9771} \MusicPDFRef{9772} \MusicPDFRef{9773} \MusicPDFRef{9774} \MusicPDFRef{9775}] /Limits [<6571756174696F6E2E422E312E312E3132> <6571756174696F6E2E422E322E312E3330>] >>}
\MusicPDFWriteObject{9005}{<< /Kids [\MusicPDFRef{9776} \MusicPDFRef{9777} \MusicPDFRef{9778} \MusicPDFRef{9779} \MusicPDFRef{9780} \MusicPDFRef{9781}] /Limits [<6571756174696F6E2E422E322E312E3331> <6571756174696F6E2E422E332E322E34>] >>}
\MusicPDFWriteObject{9006}{<< /Kids [\MusicPDFRef{9782} \MusicPDFRef{9783} \MusicPDFRef{9784} \MusicPDFRef{9785} \MusicPDFRef{9786} \MusicPDFRef{9787}] /Limits [<6571756174696F6E2E422E332E322E35> <6571756174696F6E2E422E342E312E39>] >>}
\MusicPDFWriteObject{9007}{<< /Kids [\MusicPDFRef{9788} \MusicPDFRef{9789} \MusicPDFRef{9790} \MusicPDFRef{9791} \MusicPDFRef{9792} \MusicPDFRef{9793}] /Limits [<6571756174696F6E2E422E342E322E31> <6571756174696F6E2E422E342E342E3133>] >>}
\MusicPDFWriteObject{9008}{<< /Kids [\MusicPDFRef{9794} \MusicPDFRef{9795} \MusicPDFRef{9796} \MusicPDFRef{9797} \MusicPDFRef{9798} \MusicPDFRef{9799}] /Limits [<6571756174696F6E2E422E342E342E3134> <6571756174696F6E2E422E352E332E31>] >>}
\MusicPDFWriteObject{9009}{<< /Kids [\MusicPDFRef{9800} \MusicPDFRef{9801} \MusicPDFRef{9802} \MusicPDFRef{9803} \MusicPDFRef{9804} \MusicPDFRef{9805}] /Limits [<6571756174696F6E2E422E352E332E3130> <6571756174696F6E2E422E362E312E3239>] >>}
\MusicPDFWriteObject{9010}{<< /Kids [\MusicPDFRef{9806} \MusicPDFRef{9807} \MusicPDFRef{9808} \MusicPDFRef{9809} \MusicPDFRef{9810} \MusicPDFRef{9811}] /Limits [<6571756174696F6E2E422E362E312E33> <6571756174696F6E2E422E372E312E35>] >>}
\MusicPDFWriteObject{9011}{<< /Kids [\MusicPDFRef{9812} \MusicPDFRef{9813} \MusicPDFRef{9814} \MusicPDFRef{9815} \MusicPDFRef{9816} \MusicPDFRef{9817}] /Limits [<6571756174696F6E2E422E372E312E36> <6571756174696F6E2E432E322E312E32>] >>}
\MusicPDFWriteObject{9012}{<< /Kids [\MusicPDFRef{9818} \MusicPDFRef{9819} \MusicPDFRef{9820} \MusicPDFRef{9821} \MusicPDFRef{9822} \MusicPDFRef{9823}] /Limits [<6571756174696F6E2E432E322E312E33> <6571756174696F6E2E432E322E332E38>] >>}
\MusicPDFWriteObject{9013}{<< /Kids [\MusicPDFRef{9824} \MusicPDFRef{9825} \MusicPDFRef{9826} \MusicPDFRef{9827} \MusicPDFRef{9828} \MusicPDFRef{9829}] /Limits [<6571756174696F6E2E432E322E332E39> <6571756174696F6E2E432E332E352E3138>] >>}
\MusicPDFWriteObject{9014}{<< /Kids [\MusicPDFRef{9830} \MusicPDFRef{9831} \MusicPDFRef{9832} \MusicPDFRef{9833} \MusicPDFRef{9834} \MusicPDFRef{9835}] /Limits [<6571756174696F6E2E432E332E352E3139> <6571756174696F6E2E432E342E312E31>] >>}
\MusicPDFWriteObject{9015}{<< /Kids [\MusicPDFRef{9836} \MusicPDFRef{9837} \MusicPDFRef{9838} \MusicPDFRef{9839} \MusicPDFRef{9840} \MusicPDFRef{9841}] /Limits [<6571756174696F6E2E432E342E312E32> <6571756174696F6E2E432E352E312E3134>] >>}
\MusicPDFWriteObject{9016}{<< /Kids [\MusicPDFRef{9842} \MusicPDFRef{9843} \MusicPDFRef{9844} \MusicPDFRef{9845} \MusicPDFRef{9846} \MusicPDFRef{9847}] /Limits [<6571756174696F6E2E432E352E312E3135> <6571756174696F6E2E432E352E312E3437>] >>}
\MusicPDFWriteObject{9017}{<< /Kids [\MusicPDFRef{9848} \MusicPDFRef{9849} \MusicPDFRef{9850} \MusicPDFRef{9851} \MusicPDFRef{9852} \MusicPDFRef{9853}] /Limits [<6571756174696F6E2E432E352E312E3438> <6571756174696F6E2E442E312E31302E31>] >>}
\MusicPDFWriteObject{9018}{<< /Kids [\MusicPDFRef{9854} \MusicPDFRef{9855} \MusicPDFRef{9856} \MusicPDFRef{9857} \MusicPDFRef{9858} \MusicPDFRef{9859}] /Limits [<6571756174696F6E2E442E312E31302E3130> <6571756174696F6E2E442E312E332E3137>] >>}
\MusicPDFWriteObject{9019}{<< /Kids [\MusicPDFRef{9860} \MusicPDFRef{9861} \MusicPDFRef{9862} \MusicPDFRef{9863} \MusicPDFRef{9864} \MusicPDFRef{9865}] /Limits [<6571756174696F6E2E442E312E332E3138> <6571756174696F6E2E442E312E342E33>] >>}
\MusicPDFWriteObject{9020}{<< /Kids [\MusicPDFRef{9866} \MusicPDFRef{9867} \MusicPDFRef{9868} \MusicPDFRef{9869} \MusicPDFRef{9870} \MusicPDFRef{9871}] /Limits [<6571756174696F6E2E442E312E342E34> <6571756174696F6E2E442E312E362E37>] >>}
\MusicPDFWriteObject{9021}{<< /Kids [\MusicPDFRef{9872} \MusicPDFRef{9873} \MusicPDFRef{9874} \MusicPDFRef{9875} \MusicPDFRef{9876} \MusicPDFRef{9877}] /Limits [<6571756174696F6E2E442E312E362E38> <6571756174696F6E2E442E332E312E35>] >>}
\MusicPDFWriteObject{9022}{<< /Kids [\MusicPDFRef{9878} \MusicPDFRef{9879} \MusicPDFRef{9880} \MusicPDFRef{9881} \MusicPDFRef{9882} \MusicPDFRef{9883}] /Limits [<6571756174696F6E2E442E332E312E36> <6571756174696F6E2E442E332E362E34>] >>}
\MusicPDFWriteObject{9023}{<< /Kids [\MusicPDFRef{9884} \MusicPDFRef{9885} \MusicPDFRef{9886} \MusicPDFRef{9887} \MusicPDFRef{9888} \MusicPDFRef{9889}] /Limits [<6571756174696F6E2E452E312E312E31> <6571756174696F6E2E452E322E312E32>] >>}
\MusicPDFWriteObject{9024}{<< /Kids [\MusicPDFRef{9890} \MusicPDFRef{9891} \MusicPDFRef{9892} \MusicPDFRef{9893} \MusicPDFRef{9894} \MusicPDFRef{9895}] /Limits [<6571756174696F6E2E452E322E312E33> <6571756174696F6E2E452E332E312E32>] >>}
\MusicPDFWriteObject{9025}{<< /Kids [\MusicPDFRef{9896} \MusicPDFRef{9897} \MusicPDFRef{9898} \MusicPDFRef{9899} \MusicPDFRef{9900} \MusicPDFRef{9901}] /Limits [<6571756174696F6E2E452E332E312E33> <6571756174696F6E2E452E352E312E3138>] >>}
\MusicPDFWriteObject{9026}{<< /Kids [\MusicPDFRef{9902} \MusicPDFRef{9903} \MusicPDFRef{9904} \MusicPDFRef{9905} \MusicPDFRef{9906} \MusicPDFRef{9907}] /Limits [<6571756174696F6E2E452E352E312E3139> <6571756174696F6E2E452E352E332E33>] >>}
\MusicPDFWriteObject{9027}{<< /Kids [\MusicPDFRef{9908} \MusicPDFRef{9909} \MusicPDFRef{9910} \MusicPDFRef{9911} \MusicPDFRef{9912} \MusicPDFRef{9913}] /Limits [<6571756174696F6E2E452E352E332E34> <6571756174696F6E2E452E362E322E34>] >>}
\MusicPDFWriteObject{9028}{<< /Kids [\MusicPDFRef{9914} \MusicPDFRef{9915} \MusicPDFRef{9916} \MusicPDFRef{9917} \MusicPDFRef{9918} \MusicPDFRef{9919}] /Limits [<6571756174696F6E2E452E362E322E35> <6E6F746174696F6E2E452E362E332E31>] >>}
\MusicPDFWriteObject{9029}{<< /Kids [\MusicPDFRef{9920} \MusicPDFRef{9921} \MusicPDFRef{9922} \MusicPDFRef{9923} \MusicPDFRef{9924} \MusicPDFRef{9925}] /Limits [<706167652E31> <706167652E313330>] >>}
\MusicPDFWriteObject{9030}{<< /Kids [\MusicPDFRef{9926} \MusicPDFRef{9927} \MusicPDFRef{9928} \MusicPDFRef{9929} \MusicPDFRef{9930} \MusicPDFRef{9931}] /Limits [<706167652E313331> <706167652E313633>] >>}
\MusicPDFWriteObject{9031}{<< /Kids [\MusicPDFRef{9932} \MusicPDFRef{9933} \MusicPDFRef{9934} \MusicPDFRef{9935} \MusicPDFRef{9936} \MusicPDFRef{9937}] /Limits [<706167652E313634> <706167652E313936>] >>}
\MusicPDFWriteObject{9032}{<< /Kids [\MusicPDFRef{9938} \MusicPDFRef{9939} \MusicPDFRef{9940} \MusicPDFRef{9941} \MusicPDFRef{9942} \MusicPDFRef{9943}] /Limits [<706167652E313937> <706167652E323238>] >>}
\MusicPDFWriteObject{9033}{<< /Kids [\MusicPDFRef{9944} \MusicPDFRef{9945} \MusicPDFRef{9946} \MusicPDFRef{9947} \MusicPDFRef{9948} \MusicPDFRef{9949}] /Limits [<706167652E323239> <706167652E323630>] >>}
\MusicPDFWriteObject{9034}{<< /Kids [\MusicPDFRef{9950} \MusicPDFRef{9951} \MusicPDFRef{9952} \MusicPDFRef{9953} \MusicPDFRef{9954} \MusicPDFRef{9955}] /Limits [<706167652E323631> <706167652E323933>] >>}
\MusicPDFWriteObject{9035}{<< /Kids [\MusicPDFRef{9956} \MusicPDFRef{9957} \MusicPDFRef{9958} \MusicPDFRef{9959} \MusicPDFRef{9960} \MusicPDFRef{9961}] /Limits [<706167652E323934> <706167652E333235>] >>}
\MusicPDFWriteObject{9036}{<< /Kids [\MusicPDFRef{9962} \MusicPDFRef{9963} \MusicPDFRef{9964} \MusicPDFRef{9965} \MusicPDFRef{9966} \MusicPDFRef{9967}] /Limits [<706167652E333236> <706167652E333538>] >>}
\MusicPDFWriteObject{9037}{<< /Kids [\MusicPDFRef{9968} \MusicPDFRef{9969} \MusicPDFRef{9970} \MusicPDFRef{9971} \MusicPDFRef{9972} \MusicPDFRef{9973}] /Limits [<706167652E333539> <706167652E333930>] >>}
\MusicPDFWriteObject{9038}{<< /Kids [\MusicPDFRef{9974} \MusicPDFRef{9975} \MusicPDFRef{9976} \MusicPDFRef{9977} \MusicPDFRef{9978} \MusicPDFRef{9979}] /Limits [<706167652E333931> <706167652E343232>] >>}
\MusicPDFWriteObject{9039}{<< /Kids [\MusicPDFRef{9980} \MusicPDFRef{9981} \MusicPDFRef{9982} \MusicPDFRef{9983} \MusicPDFRef{9984} \MusicPDFRef{9985}] /Limits [<706167652E343233> <706167652E343535>] >>}
\MusicPDFWriteObject{9040}{<< /Kids [\MusicPDFRef{9986} \MusicPDFRef{9987} \MusicPDFRef{9988} \MusicPDFRef{9989} \MusicPDFRef{9990} \MusicPDFRef{9991}] /Limits [<706167652E343536> <706167652E343838>] >>}
\MusicPDFWriteObject{9041}{<< /Kids [\MusicPDFRef{9992} \MusicPDFRef{9993} \MusicPDFRef{9994} \MusicPDFRef{9995} \MusicPDFRef{9996} \MusicPDFRef{9997}] /Limits [<706167652E343839> <706167652E3532>] >>}
\MusicPDFWriteObject{9042}{<< /Kids [\MusicPDFRef{9998} \MusicPDFRef{9999} \MusicPDFRef{10000} \MusicPDFRef{10001} \MusicPDFRef{10002} \MusicPDFRef{10003}] /Limits [<706167652E353230> <706167652E353532>] >>}
\MusicPDFWriteObject{9043}{<< /Kids [\MusicPDFRef{10004} \MusicPDFRef{10005} \MusicPDFRef{10006} \MusicPDFRef{10007} \MusicPDFRef{10008} \MusicPDFRef{10009}] /Limits [<706167652E353533> <706167652E353835>] >>}
\MusicPDFWriteObject{9044}{<< /Kids [\MusicPDFRef{10010} \MusicPDFRef{10011} \MusicPDFRef{10012} \MusicPDFRef{10013} \MusicPDFRef{10014} \MusicPDFRef{10015}] /Limits [<706167652E353836> <706167652E363137>] >>}
\MusicPDFWriteObject{9045}{<< /Kids [\MusicPDFRef{10016} \MusicPDFRef{10017} \MusicPDFRef{10018} \MusicPDFRef{10019} \MusicPDFRef{10020} \MusicPDFRef{10021}] /Limits [<706167652E363138> <706167652E3635>] >>}
\MusicPDFWriteObject{9046}{<< /Kids [\MusicPDFRef{10022} \MusicPDFRef{10023} \MusicPDFRef{10024} \MusicPDFRef{10025} \MusicPDFRef{10026} \MusicPDFRef{10027}] /Limits [<706167652E363530> <706167652E363832>] >>}
\MusicPDFWriteObject{9047}{<< /Kids [\MusicPDFRef{10028} \MusicPDFRef{10029} \MusicPDFRef{10030} \MusicPDFRef{10031} \MusicPDFRef{10032} \MusicPDFRef{10033}] /Limits [<706167652E363833> <706167652E373134>] >>}
\MusicPDFWriteObject{9048}{<< /Kids [\MusicPDFRef{10034} \MusicPDFRef{10035} \MusicPDFRef{10036} \MusicPDFRef{10037} \MusicPDFRef{10038} \MusicPDFRef{10039}] /Limits [<706167652E373135> <706167652E373437>] >>}
\MusicPDFWriteObject{9049}{<< /Kids [\MusicPDFRef{10040} \MusicPDFRef{10041} \MusicPDFRef{10042} \MusicPDFRef{10043} \MusicPDFRef{10044} \MusicPDFRef{10045}] /Limits [<706167652E373438> <706167652E3738>] >>}
\MusicPDFWriteObject{9050}{<< /Kids [\MusicPDFRef{10046} \MusicPDFRef{10047} \MusicPDFRef{10048} \MusicPDFRef{10049} \MusicPDFRef{10050} \MusicPDFRef{10051}] /Limits [<706167652E373830> <706167652E383131>] >>}
\MusicPDFWriteObject{9051}{<< /Kids [\MusicPDFRef{10052} \MusicPDFRef{10053} \MusicPDFRef{10054} \MusicPDFRef{10055} \MusicPDFRef{10056} \MusicPDFRef{10057}] /Limits [<706167652E383132> <706167652E383434>] >>}
\MusicPDFWriteObject{9052}{<< /Kids [\MusicPDFRef{10058} \MusicPDFRef{10059} \MusicPDFRef{10060} \MusicPDFRef{10061} \MusicPDFRef{10062} \MusicPDFRef{10063}] /Limits [<706167652E383435> <706167652E383737>] >>}
\MusicPDFWriteObject{9053}{<< /Kids [\MusicPDFRef{10064} \MusicPDFRef{10065} \MusicPDFRef{10066} \MusicPDFRef{10067} \MusicPDFRef{10068} \MusicPDFRef{10069}] /Limits [<706167652E383738> <706167652E393039>] >>}
\MusicPDFWriteObject{9054}{<< /Kids [\MusicPDFRef{10070} \MusicPDFRef{10071} \MusicPDFRef{10072} \MusicPDFRef{10073} \MusicPDFRef{10074} \MusicPDFRef{10075}] /Limits [<706167652E3931> <706167652E3935>] >>}
\MusicPDFWriteObject{9055}{<< /Kids [\MusicPDFRef{10076} \MusicPDFRef{10077} \MusicPDFRef{10078} \MusicPDFRef{10079} \MusicPDFRef{10080} \MusicPDFRef{10081}] /Limits [<706167652E3936> <73656374696F6E2E332E31>] >>}
\MusicPDFWriteObject{9056}{<< /Kids [\MusicPDFRef{10082} \MusicPDFRef{10083} \MusicPDFRef{10084} \MusicPDFRef{10085} \MusicPDFRef{10086} \MusicPDFRef{10087}] /Limits [<73656374696F6E2E332E32> <73656374696F6E2E422E31>] >>}
\MusicPDFWriteObject{9057}{<< /Kids [\MusicPDFRef{10088} \MusicPDFRef{10089} \MusicPDFRef{10090} \MusicPDFRef{10091} \MusicPDFRef{10092} \MusicPDFRef{10093}] /Limits [<73656374696F6E2E422E32> <73756273656374696F6E2E312E332E31>] >>}
\MusicPDFWriteObject{9058}{<< /Kids [\MusicPDFRef{10094} \MusicPDFRef{10095} \MusicPDFRef{10096} \MusicPDFRef{10097} \MusicPDFRef{10098} \MusicPDFRef{10099}] /Limits [<73756273656374696F6E2E312E332E32> <73756273656374696F6E2E31312E312E31>] >>}
\MusicPDFWriteObject{9059}{<< /Kids [\MusicPDFRef{10100} \MusicPDFRef{10101} \MusicPDFRef{10102} \MusicPDFRef{10103} \MusicPDFRef{10104} \MusicPDFRef{10105}] /Limits [<73756273656374696F6E2E31312E312E32> <73756273656374696F6E2E332E342E33>] >>}
\MusicPDFWriteObject{9060}{<< /Kids [\MusicPDFRef{10106} \MusicPDFRef{10107} \MusicPDFRef{10108} \MusicPDFRef{10109} \MusicPDFRef{10110} \MusicPDFRef{10111}] /Limits [<73756273656374696F6E2E342E312E31> <73756273656374696F6E2E362E332E31>] >>}
\MusicPDFWriteObject{9061}{<< /Kids [\MusicPDFRef{10112} \MusicPDFRef{10113} \MusicPDFRef{10114} \MusicPDFRef{10115} \MusicPDFRef{10116} \MusicPDFRef{10117}] /Limits [<73756273656374696F6E2E362E332E32> <73756273656374696F6E2E412E352E31>] >>}
\MusicPDFWriteObject{9062}{<< /Kids [\MusicPDFRef{10118} \MusicPDFRef{10119} \MusicPDFRef{10120} \MusicPDFRef{10121} \MusicPDFRef{10122} \MusicPDFRef{10123}] /Limits [<73756273656374696F6E2E412E352E32> <73756273656374696F6E2E432E332E37>] >>}
\MusicPDFWriteObject{9063}{<< /Kids [\MusicPDFRef{10124} \MusicPDFRef{10125} \MusicPDFRef{10126} \MusicPDFRef{10127} \MusicPDFRef{10128} \MusicPDFRef{10129}] /Limits [<73756273656374696F6E2E432E342E31> <73756273656374696F6E2E452E332E34>] >>}
\MusicPDFWriteObject{9064}{<< /Kids [\MusicPDFRef{10130} \MusicPDFRef{10131} \MusicPDFRef{10132} \MusicPDFRef{10133} \MusicPDFRef{10134} \MusicPDFRef{10135}] /Limits [<73756273656374696F6E2E452E332E35> <73756273756273656374696F6E2E322E312E312E33>] >>}
\MusicPDFWriteObject{9065}{<< /Kids [\MusicPDFRef{10136} \MusicPDFRef{10137} \MusicPDFRef{10138} \MusicPDFRef{10139} \MusicPDFRef{10140} \MusicPDFRef{10141}] /Limits [<73756273756273656374696F6E2E322E312E322E31> <73756273756273656374696F6E2E392E322E322E32>] >>}
\MusicPDFWriteObject{9066}{<< /Kids [\MusicPDFRef{10142} \MusicPDFRef{10143} \MusicPDFRef{10144} \MusicPDFRef{10145} \MusicPDFRef{10146} \MusicPDFRef{10147}] /Limits [<73756273756273656374696F6E2E392E332E322E31> <73756273756273656374696F6E2E422E322E322E31>] >>}
\MusicPDFWriteObject{9067}{<< /Kids [\MusicPDFRef{10148} \MusicPDFRef{10149} \MusicPDFRef{10150} \MusicPDFRef{10151} \MusicPDFRef{10152} \MusicPDFRef{10153}] /Limits [<73756273756273656374696F6E2E422E332E312E31> <73756273756273656374696F6E2E432E322E332E32>] >>}
\MusicPDFWriteObject{9068}{<< /Kids [\MusicPDFRef{10154} \MusicPDFRef{10155} \MusicPDFRef{10156} \MusicPDFRef{10157} \MusicPDFRef{10158} \MusicPDFRef{10159}] /Limits [<73756273756273656374696F6E2E432E322E332E33> <73756273756273656374696F6E2E442E312E352E33>] >>}
\MusicPDFWriteObject{9069}{<< /Kids [\MusicPDFRef{10160} \MusicPDFRef{10161} \MusicPDFRef{10162} \MusicPDFRef{10163} \MusicPDFRef{10164} \MusicPDFRef{10165}] /Limits [<73756273756273656374696F6E2E442E312E362E31> <7465726D696E6F6C6F67792E412E352E342E31>] >>}
\MusicPDFWriteObject{9070}{<< /Kids [\MusicPDFRef{10166} \MusicPDFRef{10167} \MusicPDFRef{10168} \MusicPDFRef{10169} \MusicPDFRef{10170} \MusicPDFRef{10171}] /Limits [<7465726D696E6F6C6F67792E442E332E332E31> <74686D2E31302E342E322E38>] >>}
\MusicPDFWriteObject{9071}{<< /Kids [\MusicPDFRef{10172} \MusicPDFRef{10173} \MusicPDFRef{10174} \MusicPDFRef{10175} \MusicPDFRef{10176} \MusicPDFRef{10177}] /Limits [<74686D2E31302E342E332E31> <74686D2E31322E322E322E32>] >>}
\MusicPDFWriteObject{9072}{<< /Kids [\MusicPDFRef{10178} \MusicPDFRef{10179} \MusicPDFRef{10180} \MusicPDFRef{10181} \MusicPDFRef{10182} \MusicPDFRef{10183}] /Limits [<74686D2E31322E322E322E33> <74686D2E322E312E332E32>] >>}
\MusicPDFWriteObject{9073}{<< /Kids [\MusicPDFRef{10184} \MusicPDFRef{10185} \MusicPDFRef{10186} \MusicPDFRef{10187} \MusicPDFRef{10188} \MusicPDFRef{10189}] /Limits [<74686D2E322E312E332E33> <74686D2E322E342E312E32>] >>}
\MusicPDFWriteObject{9074}{<< /Kids [\MusicPDFRef{10190} \MusicPDFRef{10191} \MusicPDFRef{10192} \MusicPDFRef{10193} \MusicPDFRef{10194} \MusicPDFRef{10195}] /Limits [<74686D2E322E342E322E31> <74686D2E342E312E312E31>] >>}
\MusicPDFWriteObject{9075}{<< /Kids [\MusicPDFRef{10196} \MusicPDFRef{10197} \MusicPDFRef{10198} \MusicPDFRef{10199} \MusicPDFRef{10200} \MusicPDFRef{10201}] /Limits [<74686D2E342E312E322E31> <74686D2E342E332E332E37>] >>}
\MusicPDFWriteObject{9076}{<< /Kids [\MusicPDFRef{10202} \MusicPDFRef{10203} \MusicPDFRef{10204} \MusicPDFRef{10205} \MusicPDFRef{10206} \MusicPDFRef{10207}] /Limits [<74686D2E342E332E332E38> <74686D2E352E322E312E31>] >>}
\MusicPDFWriteObject{9077}{<< /Kids [\MusicPDFRef{10208} \MusicPDFRef{10209} \MusicPDFRef{10210} \MusicPDFRef{10211} \MusicPDFRef{10212} \MusicPDFRef{10213}] /Limits [<74686D2E352E322E322E31> <74686D2E352E352E312E32>] >>}
\MusicPDFWriteObject{9078}{<< /Kids [\MusicPDFRef{10214} \MusicPDFRef{10215} \MusicPDFRef{10216} \MusicPDFRef{10217} \MusicPDFRef{10218} \MusicPDFRef{10219}] /Limits [<74686D2E352E352E312E33> <74686D2E372E322E312E34>] >>}
\MusicPDFWriteObject{9079}{<< /Kids [\MusicPDFRef{10220} \MusicPDFRef{10221} \MusicPDFRef{10222} \MusicPDFRef{10223} \MusicPDFRef{10224} \MusicPDFRef{10225}] /Limits [<74686D2E372E322E322E31> <74686D2E382E342E302E31>] >>}
\MusicPDFWriteObject{9080}{<< /Kids [\MusicPDFRef{10226} \MusicPDFRef{10227} \MusicPDFRef{10228} \MusicPDFRef{10229} \MusicPDFRef{10230} \MusicPDFRef{10231}] /Limits [<74686D2E392E312E312E31> <74686D2E392E332E312E33>] >>}
\MusicPDFWriteObject{9081}{<< /Kids [\MusicPDFRef{10232} \MusicPDFRef{10233} \MusicPDFRef{10234} \MusicPDFRef{10235} \MusicPDFRef{10236} \MusicPDFRef{10237}] /Limits [<74686D2E392E332E312E34> <74686D2E412E322E322E36>] >>}
\MusicPDFWriteObject{9082}{<< /Kids [\MusicPDFRef{10238} \MusicPDFRef{10239} \MusicPDFRef{10240} \MusicPDFRef{10241} \MusicPDFRef{10242} \MusicPDFRef{10243}] /Limits [<74686D2E412E322E322E37> <74686D2E412E332E322E34>] >>}
\MusicPDFWriteObject{9083}{<< /Kids [\MusicPDFRef{10244} \MusicPDFRef{10245} \MusicPDFRef{10246} \MusicPDFRef{10247} \MusicPDFRef{10248} \MusicPDFRef{10249}] /Limits [<74686D2E412E332E322E35> <74686D2E412E352E342E32>] >>}
\MusicPDFWriteObject{9084}{<< /Kids [\MusicPDFRef{10250} \MusicPDFRef{10251} \MusicPDFRef{10252} \MusicPDFRef{10253} \MusicPDFRef{10254} \MusicPDFRef{10255}] /Limits [<74686D2E412E352E342E33> <74686D2E422E342E332E31>] >>}
\MusicPDFWriteObject{9085}{<< /Kids [\MusicPDFRef{10256} \MusicPDFRef{10257} \MusicPDFRef{10258} \MusicPDFRef{10259} \MusicPDFRef{10260} \MusicPDFRef{10261}] /Limits [<74686D2E422E342E332E32> <74686D2E422E372E322E35>] >>}
\MusicPDFWriteObject{9086}{<< /Kids [\MusicPDFRef{10262} \MusicPDFRef{10263} \MusicPDFRef{10264} \MusicPDFRef{10265} \MusicPDFRef{10266} \MusicPDFRef{10267}] /Limits [<74686D2E422E372E322E36> <74686D2E432E322E312E35>] >>}
\MusicPDFWriteObject{9087}{<< /Kids [\MusicPDFRef{10268} \MusicPDFRef{10269} \MusicPDFRef{10270} \MusicPDFRef{10271} \MusicPDFRef{10272} \MusicPDFRef{10273}] /Limits [<74686D2E432E322E312E36> <74686D2E432E322E342E36>] >>}
\MusicPDFWriteObject{9088}{<< /Kids [\MusicPDFRef{10274} \MusicPDFRef{10275} \MusicPDFRef{10276} \MusicPDFRef{10277} \MusicPDFRef{10278} \MusicPDFRef{10279}] /Limits [<74686D2E432E322E342E37> <74686D2E432E342E332E3133>] >>}
\MusicPDFWriteObject{9089}{<< /Kids [\MusicPDFRef{10280} \MusicPDFRef{10281} \MusicPDFRef{10282} \MusicPDFRef{10283} \MusicPDFRef{10284} \MusicPDFRef{10285}] /Limits [<74686D2E432E342E332E32> <74686D2E432E352E312E33>] >>}
\MusicPDFWriteObject{9090}{<< /Kids [\MusicPDFRef{10286} \MusicPDFRef{10287} \MusicPDFRef{10288} \MusicPDFRef{10289} \MusicPDFRef{10290} \MusicPDFRef{10291}] /Limits [<74686D2E432E352E312E3330> <74686D2E442E312E322E34>] >>}
\MusicPDFWriteObject{9091}{<< /Kids [\MusicPDFRef{10292} \MusicPDFRef{10293} \MusicPDFRef{10294} \MusicPDFRef{10295} \MusicPDFRef{10296} \MusicPDFRef{10297}] /Limits [<74686D2E442E312E322E35> <74686D2E442E312E362E3135>] >>}
\MusicPDFWriteObject{9092}{<< /Kids [\MusicPDFRef{10298} \MusicPDFRef{10299} \MusicPDFRef{10300} \MusicPDFRef{10301} \MusicPDFRef{10302} \MusicPDFRef{10303}] /Limits [<74686D2E442E312E362E3136> <74686D2E442E332E322E33>] >>}
\MusicPDFWriteObject{9093}{<< /Kids [\MusicPDFRef{10304} \MusicPDFRef{10305} \MusicPDFRef{10306} \MusicPDFRef{10307} \MusicPDFRef{10308} \MusicPDFRef{10309}] /Limits [<74686D2E442E332E322E34> <74686D2E452E312E322E3133>] >>}
\MusicPDFWriteObject{9094}{<< /Kids [\MusicPDFRef{10310} \MusicPDFRef{10311} \MusicPDFRef{10312} \MusicPDFRef{10313} \MusicPDFRef{10314} \MusicPDFRef{10315}] /Limits [<74686D2E452E312E322E3134> <74686D2E452E322E312E35>] >>}
\MusicPDFWriteObject{9095}{<< /Kids [\MusicPDFRef{10316} \MusicPDFRef{10317} \MusicPDFRef{10318} \MusicPDFRef{10319} \MusicPDFRef{10320} \MusicPDFRef{10321}] /Limits [<74686D2E452E322E312E36> <74686D2E452E322E342E32>] >>}
\MusicPDFWriteObject{9096}{<< /Kids [\MusicPDFRef{10322} \MusicPDFRef{10323} \MusicPDFRef{10324} \MusicPDFRef{10325} \MusicPDFRef{10326} \MusicPDFRef{10327}] /Limits [<74686D2E452E322E342E33> <74686D2E452E342E312E32>] >>}
\MusicPDFWriteObject{9097}{<< /Kids [\MusicPDFRef{10328} \MusicPDFRef{10329} \MusicPDFRef{10330} \MusicPDFRef{10331} \MusicPDFRef{10332} \MusicPDFRef{10333}] /Limits [<74686D2E452E342E312E33> <74686D2E452E352E322E32>] >>}
\MusicPDFWriteObject{9098}{<< /Kids [\MusicPDFRef{10334} \MusicPDFRef{10335} \MusicPDFRef{10336} \MusicPDFRef{10337} \MusicPDFRef{10338} \MusicPDFRef{10339}] /Limits [<74686D2E452E352E322E33> <74686D2E452E352E372E3136>] >>}
\MusicPDFWriteObject{9099}{<< /Kids [\MusicPDFRef{10340} \MusicPDFRef{10341} \MusicPDFRef{10342} \MusicPDFRef{10343} \MusicPDFRef{10344} \MusicPDFRef{10345}] /Limits [<74686D2E452E352E372E3137> <74686D2E452E362E332E32>] >>}
\MusicPDFWriteObject{9100}{<< /Kids [\MusicPDFRef{10346} \MusicPDFRef{10347} \MusicPDFRef{10348}] /Limits [<74686D2E452E362E332E33> <74686D2E452E362E342E38>] >>}
\MusicPDFWriteObject{9101}{<FEFF004300680061007000740065007200200033002E0020004900640065006C006900630020005A006500740061002000460075006E006300740069006F006E0073>}
\MusicPDFWriteObject{9102}{<< /S /GoTo /D <636861707465722E33> >>}
\MusicPDFWriteObject{9103}{<< /Title \MusicPDFRef{10349} /A \MusicPDFRef{10350} /Parent \MusicPDFRef{1} /Prev \MusicPDFRef{8878} /Next \MusicPDFRef{10351} /First \MusicPDFRef{10352} /Last \MusicPDFRef{10353} /Count -6 >>}
\MusicPDFWriteObject{9104}{<< /Title \MusicPDFRef{10354} /A \MusicPDFRef{10355} /Parent \MusicPDFRef{8878} /Next \MusicPDFRef{10356} >>}
\MusicPDFWriteObject{9105}{<< /Title \MusicPDFRef{10357} /A \MusicPDFRef{10358} /Parent \MusicPDFRef{8878} /Prev \MusicPDFRef{10359} >>}
\MusicPDFWriteObject{9106}{<FEFF0032002E0031002E00200053006300680077006100720074007A002D004200720075006800610074002000660075006E006300740069006F006E0073>}
\MusicPDFWriteObject{9107}{<< /S /GoTo /D <73656374696F6E2E322E31> >>}
\MusicPDFWriteObject{9108}{<< /Title \MusicPDFRef{10360} /A \MusicPDFRef{10361} /Parent \MusicPDFRef{8834} /Prev \MusicPDFRef{8879} /Next \MusicPDFRef{9111} >>}
\MusicPDFWriteObject{9109}{<FEFF0032002E0034002E0020004100730079006D00700074006F0074006900630020006200650068006100760069006F00720020006100730020007C0067007C0046002C005300200030002B0020006200790020006C00610074007400690063006500200070006F0069006E0074002D0063006F0075006E00740069006E0067>}
\MusicPDFWriteObject{9110}{<< /S /GoTo /D <73656374696F6E2E322E34> >>}
\MusicPDFWriteObject{9111}{<< /Title \MusicPDFRef{10362} /A \MusicPDFRef{10363} /Parent \MusicPDFRef{8834} /Prev \MusicPDFRef{9108} /Next \MusicPDFRef{8880} >>}
\MusicPDFWriteObject{9112}{<FEFF0031002E0032002E00200049006E0074006F00200074006800650020006D006500740068006F0064003A0020006F00720062006900740061006C002000530063006800610075006400650072002000620061007300650073>}
\MusicPDFWriteObject{9113}{<< /S /GoTo /D <73656374696F6E2E312E32> >>}
\MusicPDFWriteObject{9114}{<< /Title \MusicPDFRef{10364} /A \MusicPDFRef{10365} /Parent \MusicPDFRef{8822} /Prev \MusicPDFRef{8883} /Next \MusicPDFRef{10366} >>}
\MusicPDFWriteObject{9115}{<FEFF0031002E0037002E002000450078006300650072007000740073002000660072006F006D002000610020006C0065007400740065007200200074006F0020005000650074006500720020005300610072006E0061006B>}
\MusicPDFWriteObject{9116}{<< /S /GoTo /D <73656374696F6E2E312E37> >>}
\MusicPDFWriteObject{9117}{<< /Title \MusicPDFRef{10367} /A \MusicPDFRef{10368} /Parent \MusicPDFRef{8822} /Prev \MusicPDFRef{10369} /Next \MusicPDFRef{8886} >>}
\MusicPDFWriteObject{9118}{<FEFF0041007000700065006E00640069007800200043002E0020005000720065007000610072006100740069006F006E00730020006F006E0020006C006F00630061006C006C007900200063006F006E00760065007800200074006F0070006F006C006F0067006900630061006C00200076006500630074006F00720020007300700061006300650073>}
\MusicPDFWriteObject{9119}{<< /S /GoTo /D <617070656E6469782E43> >>}
\MusicPDFWriteObject{9120}{<< /Title \MusicPDFRef{10370} /A \MusicPDFRef{10371} /Parent \MusicPDFRef{1} /Prev \MusicPDFRef{10372} /Next \MusicPDFRef{8889} /First \MusicPDFRef{10373} /Last \MusicPDFRef{10374} /Count -7 >>}
\MusicPDFWriteObject{9121}{<< /Title \MusicPDFRef{10375} /A \MusicPDFRef{10376} /Parent \MusicPDFRef{8889} /Next \MusicPDFRef{10377} >>}
\MusicPDFWriteObject{9122}{<< /Title \MusicPDFRef{10378} /A \MusicPDFRef{10379} /Parent \MusicPDFRef{8889} /Prev \MusicPDFRef{10380} >>}
\MusicPDFWriteObject{9123}{<FEFF0044002E0031002E00200053006300680077006100720074007A002D0042007200750068006100740020007300700061006300650020006F006600200061002000480061007500730064006F0072006600660020006C006F00630061006C006C007900200063006F006D00700061006300740020006100620065006C00690061006E002000670072006F00750070>}
\MusicPDFWriteObject{9124}{<< /S /GoTo /D <73656374696F6E2E442E31> >>}
\MusicPDFWriteObject{9125}{<< /Title \MusicPDFRef{10381} /A \MusicPDFRef{10382} /Parent \MusicPDFRef{8839} /Prev \MusicPDFRef{8890} /Next \MusicPDFRef{8891} >>}
\MusicPDFWriteObject{9126}{<FEFF0044002E0033002E00200053006300680077006100720074007A002D004200720075006800610074002000660075006E006300740069006F006E00730020007700690074006800200049002D006500780070006F006E0065006E007400690061006C002000640065006300610079>}
\MusicPDFWriteObject{9127}{<< /S /GoTo /D <73656374696F6E2E442E33> >>}
\MusicPDFWriteObject{9128}{<FEFF0045002E0032002E00200043006F006D0070006100630074006C0079002000670065006E0065007200610074006500640020007700650061006B006C0079002000480061007500730064006F0072006600660020007300700061006300650073>}
\MusicPDFWriteObject{9129}{<< /S /GoTo /D <73656374696F6E2E452E32> >>}
\MusicPDFWriteObject{9130}{<< /Title \MusicPDFRef{10383} /A \MusicPDFRef{10384} /Parent \MusicPDFRef{8825} /Prev \MusicPDFRef{8894} /Next \MusicPDFRef{9133} >>}
\MusicPDFWriteObject{9131}{<FEFF0045002E0035002E00200032002D00630061007400650067006F0072006900630061006C0020007000720065007000610072006100740069006F006E0073>}
\MusicPDFWriteObject{9132}{<< /S /GoTo /D <73656374696F6E2E452E35> >>}
\MusicPDFWriteObject{9133}{<< /Title \MusicPDFRef{10385} /A \MusicPDFRef{10386} /Parent \MusicPDFRef{8825} /Prev \MusicPDFRef{9130} /Next \MusicPDFRef{8897} >>}
\MusicPDFWriteObject{9134}{<< /Names [<446F632D5374617274> \MusicPDFRef{10387} <48666F6F746E6F74652E31> \MusicPDFRef{10388} <48666F6F746E6F74652E3130> \MusicPDFRef{10389} <48666F6F746E6F74652E313030> \MusicPDFRef{10390} <48666F6F746E6F74652E313031> \MusicPDFRef{10391} <48666F6F746E6F74652E313032> \MusicPDFRef{10392}] /Limits [<446F632D5374617274> <48666F6F746E6F74652E313032>] >>}
\MusicPDFWriteObject{9135}{<< /Names [<48666F6F746E6F74652E313033> \MusicPDFRef{10393} <48666F6F746E6F74652E313034> \MusicPDFRef{10394} <48666F6F746E6F74652E313035> \MusicPDFRef{10395} <48666F6F746E6F74652E313036> \MusicPDFRef{10396} <48666F6F746E6F74652E313037> \MusicPDFRef{10397} <48666F6F746E6F74652E313038> \MusicPDFRef{10398}] /Limits [<48666F6F746E6F74652E313033> <48666F6F746E6F74652E313038>] >>}
\MusicPDFWriteObject{9136}{<< /Names [<48666F6F746E6F74652E313039> \MusicPDFRef{10399} <48666F6F746E6F74652E3131> \MusicPDFRef{10400} <48666F6F746E6F74652E313130> \MusicPDFRef{10401} <48666F6F746E6F74652E313131> \MusicPDFRef{10402} <48666F6F746E6F74652E313132> \MusicPDFRef{10403} <48666F6F746E6F74652E313133> \MusicPDFRef{10404}] /Limits [<48666F6F746E6F74652E313039> <48666F6F746E6F74652E313133>] >>}
\MusicPDFWriteObject{9137}{<< /Names [<48666F6F746E6F74652E313134> \MusicPDFRef{10405} <48666F6F746E6F74652E313135> \MusicPDFRef{10406} <48666F6F746E6F74652E313136> \MusicPDFRef{10407} <48666F6F746E6F74652E313137> \MusicPDFRef{10408} <48666F6F746E6F74652E313138> \MusicPDFRef{10409} <48666F6F746E6F74652E313139> \MusicPDFRef{10410}] /Limits [<48666F6F746E6F74652E313134> <48666F6F746E6F74652E313139>] >>}
\MusicPDFWriteObject{9138}{<< /Names [<48666F6F746E6F74652E3132> \MusicPDFRef{10411} <48666F6F746E6F74652E313230> \MusicPDFRef{10412} <48666F6F746E6F74652E313231> \MusicPDFRef{10413} <48666F6F746E6F74652E313232> \MusicPDFRef{10414} <48666F6F746E6F74652E313233> \MusicPDFRef{10415} <48666F6F746E6F74652E313234> \MusicPDFRef{10416}] /Limits [<48666F6F746E6F74652E3132> <48666F6F746E6F74652E313234>] >>}
\MusicPDFWriteObject{9139}{<< /Names [<48666F6F746E6F74652E313235> \MusicPDFRef{10417} <48666F6F746E6F74652E313236> \MusicPDFRef{10418} <48666F6F746E6F74652E313237> \MusicPDFRef{10419} <48666F6F746E6F74652E313238> \MusicPDFRef{10420} <48666F6F746E6F74652E313239> \MusicPDFRef{10421} <48666F6F746E6F74652E3133> \MusicPDFRef{10422}] /Limits [<48666F6F746E6F74652E313235> <48666F6F746E6F74652E3133>] >>}
\MusicPDFWriteObject{9140}{<< /Names [<48666F6F746E6F74652E313330> \MusicPDFRef{10423} <48666F6F746E6F74652E313331> \MusicPDFRef{10424} <48666F6F746E6F74652E313332> \MusicPDFRef{10425} <48666F6F746E6F74652E313333> \MusicPDFRef{10426} <48666F6F746E6F74652E313334> \MusicPDFRef{10427} <48666F6F746E6F74652E313335> \MusicPDFRef{10428}] /Limits [<48666F6F746E6F74652E313330> <48666F6F746E6F74652E313335>] >>}
\MusicPDFWriteObject{9141}{<< /Names [<48666F6F746E6F74652E313336> \MusicPDFRef{10429} <48666F6F746E6F74652E313337> \MusicPDFRef{10430} <48666F6F746E6F74652E313338> \MusicPDFRef{10431} <48666F6F746E6F74652E313339> \MusicPDFRef{10432} <48666F6F746E6F74652E3134> \MusicPDFRef{10433} <48666F6F746E6F74652E313430> \MusicPDFRef{10434}] /Limits [<48666F6F746E6F74652E313336> <48666F6F746E6F74652E313430>] >>}
\MusicPDFWriteObject{9142}{<< /Names [<48666F6F746E6F74652E313431> \MusicPDFRef{10435} <48666F6F746E6F74652E313432> \MusicPDFRef{10436} <48666F6F746E6F74652E313433> \MusicPDFRef{10437} <48666F6F746E6F74652E313434> \MusicPDFRef{10438} <48666F6F746E6F74652E313435> \MusicPDFRef{10439} <48666F6F746E6F74652E313436> \MusicPDFRef{10440}] /Limits [<48666F6F746E6F74652E313431> <48666F6F746E6F74652E313436>] >>}
\MusicPDFWriteObject{9143}{<< /Names [<48666F6F746E6F74652E313437> \MusicPDFRef{10441} <48666F6F746E6F74652E313438> \MusicPDFRef{10442} <48666F6F746E6F74652E313439> \MusicPDFRef{10443} <48666F6F746E6F74652E3135> \MusicPDFRef{10444} <48666F6F746E6F74652E313530> \MusicPDFRef{10445} <48666F6F746E6F74652E313531> \MusicPDFRef{10446}] /Limits [<48666F6F746E6F74652E313437> <48666F6F746E6F74652E313531>] >>}
\MusicPDFWriteObject{9144}{<< /Names [<48666F6F746E6F74652E313532> \MusicPDFRef{10447} <48666F6F746E6F74652E313533> \MusicPDFRef{10448} <48666F6F746E6F74652E313534> \MusicPDFRef{10449} <48666F6F746E6F74652E313535> \MusicPDFRef{10450} <48666F6F746E6F74652E313536> \MusicPDFRef{10451} <48666F6F746E6F74652E313537> \MusicPDFRef{10452}] /Limits [<48666F6F746E6F74652E313532> <48666F6F746E6F74652E313537>] >>}
\MusicPDFWriteObject{9145}{<< /Names [<48666F6F746E6F74652E313538> \MusicPDFRef{10453} <48666F6F746E6F74652E313539> \MusicPDFRef{10454} <48666F6F746E6F74652E3136> \MusicPDFRef{10455} <48666F6F746E6F74652E313630> \MusicPDFRef{10456} <48666F6F746E6F74652E313631> \MusicPDFRef{10457} <48666F6F746E6F74652E313632> \MusicPDFRef{10458}] /Limits [<48666F6F746E6F74652E313538> <48666F6F746E6F74652E313632>] >>}
\MusicPDFWriteObject{9146}{<< /Names [<48666F6F746E6F74652E313633> \MusicPDFRef{10459} <48666F6F746E6F74652E313634> \MusicPDFRef{10460} <48666F6F746E6F74652E313635> \MusicPDFRef{10461} <48666F6F746E6F74652E313636> \MusicPDFRef{10462} <48666F6F746E6F74652E313637> \MusicPDFRef{10463} <48666F6F746E6F74652E313638> \MusicPDFRef{10464}] /Limits [<48666F6F746E6F74652E313633> <48666F6F746E6F74652E313638>] >>}
\MusicPDFWriteObject{9147}{<< /Names [<48666F6F746E6F74652E313639> \MusicPDFRef{10465} <48666F6F746E6F74652E3137> \MusicPDFRef{10466} <48666F6F746E6F74652E313730> \MusicPDFRef{10467} <48666F6F746E6F74652E313731> \MusicPDFRef{10468} <48666F6F746E6F74652E313732> \MusicPDFRef{10469} <48666F6F746E6F74652E313733> \MusicPDFRef{10470}] /Limits [<48666F6F746E6F74652E313639> <48666F6F746E6F74652E313733>] >>}
\MusicPDFWriteObject{9148}{<< /Names [<48666F6F746E6F74652E313734> \MusicPDFRef{10471} <48666F6F746E6F74652E313735> \MusicPDFRef{10472} <48666F6F746E6F74652E313736> \MusicPDFRef{10473} <48666F6F746E6F74652E313737> \MusicPDFRef{10474} <48666F6F746E6F74652E313738> \MusicPDFRef{10475} <48666F6F746E6F74652E3138> \MusicPDFRef{10476}] /Limits [<48666F6F746E6F74652E313734> <48666F6F746E6F74652E3138>] >>}
\MusicPDFWriteObject{9149}{<< /Names [<48666F6F746E6F74652E3139> \MusicPDFRef{10477} <48666F6F746E6F74652E32> \MusicPDFRef{10478} <48666F6F746E6F74652E3230> \MusicPDFRef{10479} <48666F6F746E6F74652E3231> \MusicPDFRef{10480} <48666F6F746E6F74652E3232> \MusicPDFRef{10481} <48666F6F746E6F74652E3233> \MusicPDFRef{10482}] /Limits [<48666F6F746E6F74652E3139> <48666F6F746E6F74652E3233>] >>}
\MusicPDFWriteObject{9150}{<< /Names [<48666F6F746E6F74652E3234> \MusicPDFRef{10483} <48666F6F746E6F74652E3235> \MusicPDFRef{10484} <48666F6F746E6F74652E3236> \MusicPDFRef{10485} <48666F6F746E6F74652E3237> \MusicPDFRef{10486} <48666F6F746E6F74652E3238> \MusicPDFRef{10487} <48666F6F746E6F74652E3239> \MusicPDFRef{10488}] /Limits [<48666F6F746E6F74652E3234> <48666F6F746E6F74652E3239>] >>}
\MusicPDFWriteObject{9151}{<< /Names [<48666F6F746E6F74652E33> \MusicPDFRef{10489} <48666F6F746E6F74652E3330> \MusicPDFRef{10490} <48666F6F746E6F74652E3331> \MusicPDFRef{10491} <48666F6F746E6F74652E3332> \MusicPDFRef{10492} <48666F6F746E6F74652E3333> \MusicPDFRef{10493} <48666F6F746E6F74652E3334> \MusicPDFRef{10494}] /Limits [<48666F6F746E6F74652E33> <48666F6F746E6F74652E3334>] >>}
\MusicPDFWriteObject{9152}{<< /Names [<48666F6F746E6F74652E3335> \MusicPDFRef{10495} <48666F6F746E6F74652E3336> \MusicPDFRef{10496} <48666F6F746E6F74652E3337> \MusicPDFRef{10497} <48666F6F746E6F74652E3338> \MusicPDFRef{10498} <48666F6F746E6F74652E3339> \MusicPDFRef{10499} <48666F6F746E6F74652E34> \MusicPDFRef{10500}] /Limits [<48666F6F746E6F74652E3335> <48666F6F746E6F74652E34>] >>}
\MusicPDFWriteObject{9153}{<< /Names [<48666F6F746E6F74652E3430> \MusicPDFRef{10501} <48666F6F746E6F74652E3431> \MusicPDFRef{10502} <48666F6F746E6F74652E3432> \MusicPDFRef{10503} <48666F6F746E6F74652E3433> \MusicPDFRef{10504} <48666F6F746E6F74652E3434> \MusicPDFRef{10505} <48666F6F746E6F74652E3435> \MusicPDFRef{10506}] /Limits [<48666F6F746E6F74652E3430> <48666F6F746E6F74652E3435>] >>}
\MusicPDFWriteObject{9154}{<< /Names [<48666F6F746E6F74652E3436> \MusicPDFRef{10507} <48666F6F746E6F74652E3437> \MusicPDFRef{10508} <48666F6F746E6F74652E3438> \MusicPDFRef{10509} <48666F6F746E6F74652E3439> \MusicPDFRef{10510} <48666F6F746E6F74652E35> \MusicPDFRef{10511} <48666F6F746E6F74652E3530> \MusicPDFRef{10512}] /Limits [<48666F6F746E6F74652E3436> <48666F6F746E6F74652E3530>] >>}
\MusicPDFWriteObject{9155}{<< /Names [<48666F6F746E6F74652E3531> \MusicPDFRef{10513} <48666F6F746E6F74652E3532> \MusicPDFRef{10514} <48666F6F746E6F74652E3533> \MusicPDFRef{10515} <48666F6F746E6F74652E3534> \MusicPDFRef{10516} <48666F6F746E6F74652E3535> \MusicPDFRef{10517} <48666F6F746E6F74652E3536> \MusicPDFRef{10518}] /Limits [<48666F6F746E6F74652E3531> <48666F6F746E6F74652E3536>] >>}
\MusicPDFWriteObject{9156}{<< /Names [<48666F6F746E6F74652E3537> \MusicPDFRef{10519} <48666F6F746E6F74652E3538> \MusicPDFRef{10520} <48666F6F746E6F74652E3539> \MusicPDFRef{10521} <48666F6F746E6F74652E36> \MusicPDFRef{10522} <48666F6F746E6F74652E3630> \MusicPDFRef{10523} <48666F6F746E6F74652E3631> \MusicPDFRef{10524}] /Limits [<48666F6F746E6F74652E3537> <48666F6F746E6F74652E3631>] >>}
\MusicPDFWriteObject{9157}{<< /Names [<48666F6F746E6F74652E3632> \MusicPDFRef{10525} <48666F6F746E6F74652E3633> \MusicPDFRef{10526} <48666F6F746E6F74652E3634> \MusicPDFRef{10527} <48666F6F746E6F74652E3635> \MusicPDFRef{10528} <48666F6F746E6F74652E3636> \MusicPDFRef{10529} <48666F6F746E6F74652E3637> \MusicPDFRef{10530}] /Limits [<48666F6F746E6F74652E3632> <48666F6F746E6F74652E3637>] >>}
\MusicPDFWriteObject{9158}{<< /Names [<48666F6F746E6F74652E3638> \MusicPDFRef{10531} <48666F6F746E6F74652E3639> \MusicPDFRef{10532} <48666F6F746E6F74652E37> \MusicPDFRef{10533} <48666F6F746E6F74652E3730> \MusicPDFRef{10534} <48666F6F746E6F74652E3731> \MusicPDFRef{10535} <48666F6F746E6F74652E3732> \MusicPDFRef{10536}] /Limits [<48666F6F746E6F74652E3638> <48666F6F746E6F74652E3732>] >>}
\MusicPDFWriteObject{9159}{<< /Names [<48666F6F746E6F74652E3733> \MusicPDFRef{10537} <48666F6F746E6F74652E3734> \MusicPDFRef{10538} <48666F6F746E6F74652E3735> \MusicPDFRef{10539} <48666F6F746E6F74652E3736> \MusicPDFRef{10540} <48666F6F746E6F74652E3737> \MusicPDFRef{10541} <48666F6F746E6F74652E3738> \MusicPDFRef{10542}] /Limits [<48666F6F746E6F74652E3733> <48666F6F746E6F74652E3738>] >>}
\MusicPDFWriteObject{9160}{<< /Names [<48666F6F746E6F74652E3739> \MusicPDFRef{10543} <48666F6F746E6F74652E38> \MusicPDFRef{10544} <48666F6F746E6F74652E3830> \MusicPDFRef{10545} <48666F6F746E6F74652E3831> \MusicPDFRef{10546} <48666F6F746E6F74652E3832> \MusicPDFRef{10547} <48666F6F746E6F74652E3833> \MusicPDFRef{10548}] /Limits [<48666F6F746E6F74652E3739> <48666F6F746E6F74652E3833>] >>}
\MusicPDFWriteObject{9161}{<< /Names [<48666F6F746E6F74652E3834> \MusicPDFRef{10549} <48666F6F746E6F74652E3835> \MusicPDFRef{10550} <48666F6F746E6F74652E3836> \MusicPDFRef{10551} <48666F6F746E6F74652E3837> \MusicPDFRef{10552} <48666F6F746E6F74652E3838> \MusicPDFRef{10553} <48666F6F746E6F74652E3839> \MusicPDFRef{10554}] /Limits [<48666F6F746E6F74652E3834> <48666F6F746E6F74652E3839>] >>}
\MusicPDFWriteObject{9162}{<< /Names [<48666F6F746E6F74652E39> \MusicPDFRef{10555} <48666F6F746E6F74652E3930> \MusicPDFRef{10556} <48666F6F746E6F74652E3931> \MusicPDFRef{10557} <48666F6F746E6F74652E3932> \MusicPDFRef{10558} <48666F6F746E6F74652E3933> \MusicPDFRef{10559} <48666F6F746E6F74652E3934> \MusicPDFRef{10560}] /Limits [<48666F6F746E6F74652E39> <48666F6F746E6F74652E3934>] >>}
\MusicPDFWriteObject{9163}{<< /Names [<48666F6F746E6F74652E3935> \MusicPDFRef{10561} <48666F6F746E6F74652E3936> \MusicPDFRef{10562} <48666F6F746E6F74652E3937> \MusicPDFRef{10563} <48666F6F746E6F74652E3938> \MusicPDFRef{10564} <48666F6F746E6F74652E3939> \MusicPDFRef{10565} <4974656D2E31> \MusicPDFRef{10566}] /Limits [<48666F6F746E6F74652E3935> <4974656D2E31>] >>}
\MusicPDFWriteObject{9164}{<< /Names [<4974656D2E3130> \MusicPDFRef{10567} <4974656D2E313030> \MusicPDFRef{10568} <4974656D2E31303030> \MusicPDFRef{10569} <4974656D2E31303031> \MusicPDFRef{10570} <4974656D2E31303032> \MusicPDFRef{10571} <4974656D2E31303033> \MusicPDFRef{10572}] /Limits [<4974656D2E3130> <4974656D2E31303033>] >>}
\MusicPDFWriteObject{9165}{<< /Names [<4974656D2E31303034> \MusicPDFRef{10573} <4974656D2E31303035> \MusicPDFRef{10574} <4974656D2E31303036> \MusicPDFRef{10575} <4974656D2E31303037> \MusicPDFRef{10576} <4974656D2E31303038> \MusicPDFRef{10577} <4974656D2E31303039> \MusicPDFRef{10578}] /Limits [<4974656D2E31303034> <4974656D2E31303039>] >>}
\MusicPDFWriteObject{9166}{<< /Names [<4974656D2E313031> \MusicPDFRef{10579} <4974656D2E31303130> \MusicPDFRef{10580} <4974656D2E31303131> \MusicPDFRef{10581} <4974656D2E31303132> \MusicPDFRef{10582} <4974656D2E31303133> \MusicPDFRef{10583} <4974656D2E31303134> \MusicPDFRef{10584}] /Limits [<4974656D2E313031> <4974656D2E31303134>] >>}
\MusicPDFWriteObject{9167}{<< /Names [<4974656D2E31303135> \MusicPDFRef{10585} <4974656D2E31303136> \MusicPDFRef{10586} <4974656D2E31303137> \MusicPDFRef{10587} <4974656D2E31303138> \MusicPDFRef{10588} <4974656D2E31303139> \MusicPDFRef{10589} <4974656D2E313032> \MusicPDFRef{10590}] /Limits [<4974656D2E31303135> <4974656D2E313032>] >>}
\MusicPDFWriteObject{9168}{<< /Names [<4974656D2E31303230> \MusicPDFRef{10591} <4974656D2E31303231> \MusicPDFRef{10592} <4974656D2E31303232> \MusicPDFRef{10593} <4974656D2E31303233> \MusicPDFRef{10594} <4974656D2E31303234> \MusicPDFRef{10595} <4974656D2E31303235> \MusicPDFRef{10596}] /Limits [<4974656D2E31303230> <4974656D2E31303235>] >>}
\MusicPDFWriteObject{9169}{<< /Names [<4974656D2E31303236> \MusicPDFRef{10597} <4974656D2E31303237> \MusicPDFRef{10598} <4974656D2E31303238> \MusicPDFRef{10599} <4974656D2E31303239> \MusicPDFRef{10600} <4974656D2E313033> \MusicPDFRef{10601} <4974656D2E31303330> \MusicPDFRef{10602}] /Limits [<4974656D2E31303236> <4974656D2E31303330>] >>}
\MusicPDFWriteObject{9170}{<< /Names [<4974656D2E31303331> \MusicPDFRef{10603} <4974656D2E31303332> \MusicPDFRef{10604} <4974656D2E31303333> \MusicPDFRef{10605} <4974656D2E31303334> \MusicPDFRef{10606} <4974656D2E31303335> \MusicPDFRef{10607} <4974656D2E31303336> \MusicPDFRef{10608}] /Limits [<4974656D2E31303331> <4974656D2E31303336>] >>}
\MusicPDFWriteObject{9171}{<< /Names [<4974656D2E31303337> \MusicPDFRef{10609} <4974656D2E31303338> \MusicPDFRef{10610} <4974656D2E31303339> \MusicPDFRef{10611} <4974656D2E313034> \MusicPDFRef{10612} <4974656D2E31303430> \MusicPDFRef{10613} <4974656D2E31303431> \MusicPDFRef{10614}] /Limits [<4974656D2E31303337> <4974656D2E31303431>] >>}
\MusicPDFWriteObject{9172}{<< /Names [<4974656D2E31303432> \MusicPDFRef{10615} <4974656D2E31303433> \MusicPDFRef{10616} <4974656D2E31303434> \MusicPDFRef{10617} <4974656D2E31303435> \MusicPDFRef{10618} <4974656D2E31303436> \MusicPDFRef{10619} <4974656D2E31303437> \MusicPDFRef{10620}] /Limits [<4974656D2E31303432> <4974656D2E31303437>] >>}
\MusicPDFWriteObject{9173}{<< /Names [<4974656D2E31303438> \MusicPDFRef{10621} <4974656D2E31303439> \MusicPDFRef{10622} <4974656D2E313035> \MusicPDFRef{10623} <4974656D2E31303530> \MusicPDFRef{10624} <4974656D2E31303531> \MusicPDFRef{10625} <4974656D2E31303532> \MusicPDFRef{10626}] /Limits [<4974656D2E31303438> <4974656D2E31303532>] >>}
\MusicPDFWriteObject{9174}{<< /Names [<4974656D2E31303533> \MusicPDFRef{10627} <4974656D2E31303534> \MusicPDFRef{10628} <4974656D2E31303535> \MusicPDFRef{10629} <4974656D2E31303536> \MusicPDFRef{10630} <4974656D2E31303537> \MusicPDFRef{10631} <4974656D2E31303538> \MusicPDFRef{10632}] /Limits [<4974656D2E31303533> <4974656D2E31303538>] >>}
\MusicPDFWriteObject{9175}{<< /Names [<4974656D2E31303539> \MusicPDFRef{10633} <4974656D2E313036> \MusicPDFRef{10634} <4974656D2E31303630> \MusicPDFRef{10635} <4974656D2E31303631> \MusicPDFRef{10636} <4974656D2E31303632> \MusicPDFRef{10637} <4974656D2E31303633> \MusicPDFRef{10638}] /Limits [<4974656D2E31303539> <4974656D2E31303633>] >>}
\MusicPDFWriteObject{9176}{<< /Names [<4974656D2E31303634> \MusicPDFRef{10639} <4974656D2E31303635> \MusicPDFRef{10640} <4974656D2E31303636> \MusicPDFRef{10641} <4974656D2E31303637> \MusicPDFRef{10642} <4974656D2E31303638> \MusicPDFRef{10643} <4974656D2E31303639> \MusicPDFRef{10644}] /Limits [<4974656D2E31303634> <4974656D2E31303639>] >>}
\MusicPDFWriteObject{9177}{<< /Names [<4974656D2E313037> \MusicPDFRef{10645} <4974656D2E31303730> \MusicPDFRef{10646} <4974656D2E31303731> \MusicPDFRef{10647} <4974656D2E31303732> \MusicPDFRef{10648} <4974656D2E31303733> \MusicPDFRef{10649} <4974656D2E31303734> \MusicPDFRef{10650}] /Limits [<4974656D2E313037> <4974656D2E31303734>] >>}
\MusicPDFWriteObject{9178}{<< /Names [<4974656D2E31303735> \MusicPDFRef{10651} <4974656D2E31303736> \MusicPDFRef{10652} <4974656D2E31303737> \MusicPDFRef{10653} <4974656D2E31303738> \MusicPDFRef{10654} <4974656D2E31303739> \MusicPDFRef{10655} <4974656D2E313038> \MusicPDFRef{10656}] /Limits [<4974656D2E31303735> <4974656D2E313038>] >>}
\MusicPDFWriteObject{9179}{<< /Names [<4974656D2E31303830> \MusicPDFRef{10657} <4974656D2E31303831> \MusicPDFRef{10658} <4974656D2E31303832> \MusicPDFRef{10659} <4974656D2E31303833> \MusicPDFRef{10660} <4974656D2E31303834> \MusicPDFRef{10661} <4974656D2E31303835> \MusicPDFRef{10662}] /Limits [<4974656D2E31303830> <4974656D2E31303835>] >>}
\MusicPDFWriteObject{9180}{<< /Names [<4974656D2E31303836> \MusicPDFRef{10663} <4974656D2E31303837> \MusicPDFRef{10664} <4974656D2E31303838> \MusicPDFRef{10665} <4974656D2E31303839> \MusicPDFRef{10666} <4974656D2E313039> \MusicPDFRef{10667} <4974656D2E31303930> \MusicPDFRef{10668}] /Limits [<4974656D2E31303836> <4974656D2E31303930>] >>}
\MusicPDFWriteObject{9181}{<< /Names [<4974656D2E31303931> \MusicPDFRef{10669} <4974656D2E31303932> \MusicPDFRef{10670} <4974656D2E31303933> \MusicPDFRef{10671} <4974656D2E31303934> \MusicPDFRef{10672} <4974656D2E31303935> \MusicPDFRef{10673} <4974656D2E31303936> \MusicPDFRef{10674}] /Limits [<4974656D2E31303931> <4974656D2E31303936>] >>}
\MusicPDFWriteObject{9182}{<< /Names [<4974656D2E31303937> \MusicPDFRef{10675} <4974656D2E31303938> \MusicPDFRef{10676} <4974656D2E31303939> \MusicPDFRef{10677} <4974656D2E3131> \MusicPDFRef{10678} <4974656D2E313130> \MusicPDFRef{10679} <4974656D2E31313030> \MusicPDFRef{10680}] /Limits [<4974656D2E31303937> <4974656D2E31313030>] >>}
\MusicPDFWriteObject{9183}{<< /Names [<4974656D2E31313031> \MusicPDFRef{10681} <4974656D2E31313032> \MusicPDFRef{10682} <4974656D2E31313033> \MusicPDFRef{10683} <4974656D2E31313034> \MusicPDFRef{10684} <4974656D2E31313035> \MusicPDFRef{10685} <4974656D2E31313036> \MusicPDFRef{10686}] /Limits [<4974656D2E31313031> <4974656D2E31313036>] >>}
\MusicPDFWriteObject{9184}{<< /Names [<4974656D2E31313037> \MusicPDFRef{10687} <4974656D2E31313038> \MusicPDFRef{10688} <4974656D2E31313039> \MusicPDFRef{10689} <4974656D2E313131> \MusicPDFRef{10690} <4974656D2E31313130> \MusicPDFRef{10691} <4974656D2E31313131> \MusicPDFRef{10692}] /Limits [<4974656D2E31313037> <4974656D2E31313131>] >>}
\MusicPDFWriteObject{9185}{<< /Names [<4974656D2E31313132> \MusicPDFRef{10693} <4974656D2E31313133> \MusicPDFRef{10694} <4974656D2E31313134> \MusicPDFRef{10695} <4974656D2E31313135> \MusicPDFRef{10696} <4974656D2E31313136> \MusicPDFRef{10697} <4974656D2E31313137> \MusicPDFRef{10698}] /Limits [<4974656D2E31313132> <4974656D2E31313137>] >>}
\MusicPDFWriteObject{9186}{<< /Names [<4974656D2E31313138> \MusicPDFRef{10699} <4974656D2E31313139> \MusicPDFRef{10700} <4974656D2E313132> \MusicPDFRef{10701} <4974656D2E31313230> \MusicPDFRef{10702} <4974656D2E31313231> \MusicPDFRef{10703} <4974656D2E31313232> \MusicPDFRef{10704}] /Limits [<4974656D2E31313138> <4974656D2E31313232>] >>}
\MusicPDFWriteObject{9187}{<< /Names [<4974656D2E31313233> \MusicPDFRef{10705} <4974656D2E31313234> \MusicPDFRef{10706} <4974656D2E31313235> \MusicPDFRef{10707} <4974656D2E31313236> \MusicPDFRef{10708} <4974656D2E31313237> \MusicPDFRef{10709} <4974656D2E31313238> \MusicPDFRef{10710}] /Limits [<4974656D2E31313233> <4974656D2E31313238>] >>}
\MusicPDFWriteObject{9188}{<< /Names [<4974656D2E31313239> \MusicPDFRef{10711} <4974656D2E313133> \MusicPDFRef{10712} <4974656D2E31313330> \MusicPDFRef{10713} <4974656D2E31313331> \MusicPDFRef{10714} <4974656D2E31313332> \MusicPDFRef{10715} <4974656D2E31313333> \MusicPDFRef{10716}] /Limits [<4974656D2E31313239> <4974656D2E31313333>] >>}
\MusicPDFWriteObject{9189}{<< /Names [<4974656D2E31313334> \MusicPDFRef{10717} <4974656D2E31313335> \MusicPDFRef{10718} <4974656D2E31313336> \MusicPDFRef{10719} <4974656D2E31313337> \MusicPDFRef{10720} <4974656D2E31313338> \MusicPDFRef{10721} <4974656D2E31313339> \MusicPDFRef{10722}] /Limits [<4974656D2E31313334> <4974656D2E31313339>] >>}
\MusicPDFWriteObject{9190}{<< /Names [<4974656D2E313134> \MusicPDFRef{10723} <4974656D2E31313430> \MusicPDFRef{10724} <4974656D2E31313431> \MusicPDFRef{10725} <4974656D2E31313432> \MusicPDFRef{10726} <4974656D2E31313433> \MusicPDFRef{10727} <4974656D2E31313434> \MusicPDFRef{10728}] /Limits [<4974656D2E313134> <4974656D2E31313434>] >>}
\MusicPDFWriteObject{9191}{<< /Names [<4974656D2E31313435> \MusicPDFRef{10729} <4974656D2E31313436> \MusicPDFRef{10730} <4974656D2E31313437> \MusicPDFRef{10731} <4974656D2E31313438> \MusicPDFRef{10732} <4974656D2E31313439> \MusicPDFRef{10733} <4974656D2E313135> \MusicPDFRef{10734}] /Limits [<4974656D2E31313435> <4974656D2E313135>] >>}
\MusicPDFWriteObject{9192}{<< /Names [<4974656D2E31313530> \MusicPDFRef{10735} <4974656D2E31313531> \MusicPDFRef{10736} <4974656D2E31313532> \MusicPDFRef{10737} <4974656D2E31313533> \MusicPDFRef{10738} <4974656D2E31313534> \MusicPDFRef{10739} <4974656D2E31313535> \MusicPDFRef{10740}] /Limits [<4974656D2E31313530> <4974656D2E31313535>] >>}
\MusicPDFWriteObject{9193}{<< /Names [<4974656D2E31313536> \MusicPDFRef{10741} <4974656D2E31313537> \MusicPDFRef{10742} <4974656D2E31313538> \MusicPDFRef{10743} <4974656D2E31313539> \MusicPDFRef{10744} <4974656D2E313136> \MusicPDFRef{10745} <4974656D2E31313630> \MusicPDFRef{10746}] /Limits [<4974656D2E31313536> <4974656D2E31313630>] >>}
\MusicPDFWriteObject{9194}{<< /Names [<4974656D2E31313631> \MusicPDFRef{10747} <4974656D2E31313632> \MusicPDFRef{10748} <4974656D2E31313633> \MusicPDFRef{10749} <4974656D2E31313634> \MusicPDFRef{10750} <4974656D2E31313635> \MusicPDFRef{10751} <4974656D2E31313636> \MusicPDFRef{10752}] /Limits [<4974656D2E31313631> <4974656D2E31313636>] >>}
\MusicPDFWriteObject{9195}{<< /Names [<4974656D2E31313637> \MusicPDFRef{10753} <4974656D2E31313638> \MusicPDFRef{10754} <4974656D2E31313639> \MusicPDFRef{10755} <4974656D2E313137> \MusicPDFRef{10756} <4974656D2E31313730> \MusicPDFRef{10757} <4974656D2E31313731> \MusicPDFRef{10758}] /Limits [<4974656D2E31313637> <4974656D2E31313731>] >>}
\MusicPDFWriteObject{9196}{<< /Names [<4974656D2E31313732> \MusicPDFRef{10759} <4974656D2E31313733> \MusicPDFRef{10760} <4974656D2E31313734> \MusicPDFRef{10761} <4974656D2E31313735> \MusicPDFRef{10762} <4974656D2E31313736> \MusicPDFRef{10763} <4974656D2E31313737> \MusicPDFRef{10764}] /Limits [<4974656D2E31313732> <4974656D2E31313737>] >>}
\MusicPDFWriteObject{9197}{<< /Names [<4974656D2E31313738> \MusicPDFRef{10765} <4974656D2E31313739> \MusicPDFRef{10766} <4974656D2E313138> \MusicPDFRef{10767} <4974656D2E31313830> \MusicPDFRef{10768} <4974656D2E31313831> \MusicPDFRef{10769} <4974656D2E31313832> \MusicPDFRef{10770}] /Limits [<4974656D2E31313738> <4974656D2E31313832>] >>}
\MusicPDFWriteObject{9198}{<< /Names [<4974656D2E31313833> \MusicPDFRef{10771} <4974656D2E31313834> \MusicPDFRef{10772} <4974656D2E31313835> \MusicPDFRef{10773} <4974656D2E31313836> \MusicPDFRef{10774} <4974656D2E31313837> \MusicPDFRef{10775} <4974656D2E31313838> \MusicPDFRef{10776}] /Limits [<4974656D2E31313833> <4974656D2E31313838>] >>}
\MusicPDFWriteObject{9199}{<< /Names [<4974656D2E31313839> \MusicPDFRef{10777} <4974656D2E313139> \MusicPDFRef{10778} <4974656D2E31313930> \MusicPDFRef{10779} <4974656D2E31313931> \MusicPDFRef{10780} <4974656D2E31313932> \MusicPDFRef{10781} <4974656D2E31313933> \MusicPDFRef{10782}] /Limits [<4974656D2E31313839> <4974656D2E31313933>] >>}
\MusicPDFWriteObject{9200}{<< /Names [<4974656D2E31313934> \MusicPDFRef{10783} <4974656D2E31313935> \MusicPDFRef{10784} <4974656D2E31313936> \MusicPDFRef{10785} <4974656D2E31313937> \MusicPDFRef{10786} <4974656D2E31313938> \MusicPDFRef{10787} <4974656D2E31313939> \MusicPDFRef{10788}] /Limits [<4974656D2E31313934> <4974656D2E31313939>] >>}
\MusicPDFWriteObject{9201}{<< /Names [<4974656D2E3132> \MusicPDFRef{10789} <4974656D2E313230> \MusicPDFRef{10790} <4974656D2E31323030> \MusicPDFRef{10791} <4974656D2E31323031> \MusicPDFRef{10792} <4974656D2E31323032> \MusicPDFRef{10793} <4974656D2E31323033> \MusicPDFRef{10794}] /Limits [<4974656D2E3132> <4974656D2E31323033>] >>}
\MusicPDFWriteObject{9202}{<< /Names [<4974656D2E31323034> \MusicPDFRef{10795} <4974656D2E31323035> \MusicPDFRef{10796} <4974656D2E31323036> \MusicPDFRef{10797} <4974656D2E31323037> \MusicPDFRef{10798} <4974656D2E31323038> \MusicPDFRef{10799} <4974656D2E31323039> \MusicPDFRef{10800}] /Limits [<4974656D2E31323034> <4974656D2E31323039>] >>}
\MusicPDFWriteObject{9203}{<< /Names [<4974656D2E313231> \MusicPDFRef{10801} <4974656D2E31323130> \MusicPDFRef{10802} <4974656D2E31323131> \MusicPDFRef{10803} <4974656D2E31323132> \MusicPDFRef{10804} <4974656D2E31323133> \MusicPDFRef{10805} <4974656D2E31323134> \MusicPDFRef{10806}] /Limits [<4974656D2E313231> <4974656D2E31323134>] >>}
\MusicPDFWriteObject{9204}{<< /Names [<4974656D2E31323135> \MusicPDFRef{10807} <4974656D2E31323136> \MusicPDFRef{10808} <4974656D2E31323137> \MusicPDFRef{10809} <4974656D2E31323138> \MusicPDFRef{10810} <4974656D2E31323139> \MusicPDFRef{10811} <4974656D2E313232> \MusicPDFRef{10812}] /Limits [<4974656D2E31323135> <4974656D2E313232>] >>}
\MusicPDFWriteObject{9205}{<< /Names [<4974656D2E31323230> \MusicPDFRef{10813} <4974656D2E31323231> \MusicPDFRef{10814} <4974656D2E31323232> \MusicPDFRef{10815} <4974656D2E31323233> \MusicPDFRef{10816} <4974656D2E31323234> \MusicPDFRef{10817} <4974656D2E31323235> \MusicPDFRef{10818}] /Limits [<4974656D2E31323230> <4974656D2E31323235>] >>}
\MusicPDFWriteObject{9206}{<< /Names [<4974656D2E31323236> \MusicPDFRef{10819} <4974656D2E31323237> \MusicPDFRef{10820} <4974656D2E31323238> \MusicPDFRef{10821} <4974656D2E31323239> \MusicPDFRef{10822} <4974656D2E313233> \MusicPDFRef{10823} <4974656D2E31323330> \MusicPDFRef{10824}] /Limits [<4974656D2E31323236> <4974656D2E31323330>] >>}
\MusicPDFWriteObject{9207}{<< /Names [<4974656D2E31323331> \MusicPDFRef{10825} <4974656D2E31323332> \MusicPDFRef{10826} <4974656D2E31323333> \MusicPDFRef{10827} <4974656D2E31323334> \MusicPDFRef{10828} <4974656D2E31323335> \MusicPDFRef{10829} <4974656D2E31323336> \MusicPDFRef{10830}] /Limits [<4974656D2E31323331> <4974656D2E31323336>] >>}
\MusicPDFWriteObject{9208}{<< /Names [<4974656D2E31323337> \MusicPDFRef{10831} <4974656D2E31323338> \MusicPDFRef{10832} <4974656D2E31323339> \MusicPDFRef{10833} <4974656D2E313234> \MusicPDFRef{10834} <4974656D2E31323430> \MusicPDFRef{10835} <4974656D2E31323431> \MusicPDFRef{10836}] /Limits [<4974656D2E31323337> <4974656D2E31323431>] >>}
\MusicPDFWriteObject{9209}{<< /Names [<4974656D2E31323432> \MusicPDFRef{10837} <4974656D2E31323433> \MusicPDFRef{10838} <4974656D2E31323434> \MusicPDFRef{10839} <4974656D2E31323435> \MusicPDFRef{10840} <4974656D2E31323436> \MusicPDFRef{10841} <4974656D2E31323437> \MusicPDFRef{10842}] /Limits [<4974656D2E31323432> <4974656D2E31323437>] >>}
\MusicPDFWriteObject{9210}{<< /Names [<4974656D2E31323438> \MusicPDFRef{10843} <4974656D2E31323439> \MusicPDFRef{10844} <4974656D2E313235> \MusicPDFRef{10845} <4974656D2E31323530> \MusicPDFRef{10846} <4974656D2E31323531> \MusicPDFRef{10847} <4974656D2E31323532> \MusicPDFRef{10848}] /Limits [<4974656D2E31323438> <4974656D2E31323532>] >>}
\MusicPDFWriteObject{9211}{<< /Names [<4974656D2E31323533> \MusicPDFRef{10849} <4974656D2E31323534> \MusicPDFRef{10850} <4974656D2E31323535> \MusicPDFRef{10851} <4974656D2E31323536> \MusicPDFRef{10852} <4974656D2E31323537> \MusicPDFRef{10853} <4974656D2E31323538> \MusicPDFRef{10854}] /Limits [<4974656D2E31323533> <4974656D2E31323538>] >>}
\MusicPDFWriteObject{9212}{<< /Names [<4974656D2E31323539> \MusicPDFRef{10855} <4974656D2E313236> \MusicPDFRef{10856} <4974656D2E31323630> \MusicPDFRef{10857} <4974656D2E31323631> \MusicPDFRef{10858} <4974656D2E31323632> \MusicPDFRef{10859} <4974656D2E31323633> \MusicPDFRef{10860}] /Limits [<4974656D2E31323539> <4974656D2E31323633>] >>}
\MusicPDFWriteObject{9213}{<< /Names [<4974656D2E31323634> \MusicPDFRef{10861} <4974656D2E31323635> \MusicPDFRef{10862} <4974656D2E31323636> \MusicPDFRef{10863} <4974656D2E31323637> \MusicPDFRef{10864} <4974656D2E31323638> \MusicPDFRef{10865} <4974656D2E31323639> \MusicPDFRef{10866}] /Limits [<4974656D2E31323634> <4974656D2E31323639>] >>}
\MusicPDFWriteObject{9214}{<< /Names [<4974656D2E313237> \MusicPDFRef{10867} <4974656D2E31323730> \MusicPDFRef{10868} <4974656D2E31323731> \MusicPDFRef{10869} <4974656D2E31323732> \MusicPDFRef{10870} <4974656D2E31323733> \MusicPDFRef{10871} <4974656D2E31323734> \MusicPDFRef{10872}] /Limits [<4974656D2E313237> <4974656D2E31323734>] >>}
\MusicPDFWriteObject{9215}{<< /Names [<4974656D2E31323735> \MusicPDFRef{10873} <4974656D2E31323736> \MusicPDFRef{10874} <4974656D2E31323737> \MusicPDFRef{10875} <4974656D2E31323738> \MusicPDFRef{10876} <4974656D2E31323739> \MusicPDFRef{10877} <4974656D2E313238> \MusicPDFRef{10878}] /Limits [<4974656D2E31323735> <4974656D2E313238>] >>}
\MusicPDFWriteObject{9216}{<< /Names [<4974656D2E31323830> \MusicPDFRef{10879} <4974656D2E31323831> \MusicPDFRef{10880} <4974656D2E31323832> \MusicPDFRef{10881} <4974656D2E31323833> \MusicPDFRef{10882} <4974656D2E31323834> \MusicPDFRef{10883} <4974656D2E31323835> \MusicPDFRef{10884}] /Limits [<4974656D2E31323830> <4974656D2E31323835>] >>}
\MusicPDFWriteObject{9217}{<< /Names [<4974656D2E31323836> \MusicPDFRef{10885} <4974656D2E31323837> \MusicPDFRef{10886} <4974656D2E31323838> \MusicPDFRef{10887} <4974656D2E31323839> \MusicPDFRef{10888} <4974656D2E313239> \MusicPDFRef{10889} <4974656D2E31323930> \MusicPDFRef{10890}] /Limits [<4974656D2E31323836> <4974656D2E31323930>] >>}
\MusicPDFWriteObject{9218}{<< /Names [<4974656D2E31323931> \MusicPDFRef{10891} <4974656D2E31323932> \MusicPDFRef{10892} <4974656D2E31323933> \MusicPDFRef{10893} <4974656D2E31323934> \MusicPDFRef{10894} <4974656D2E31323935> \MusicPDFRef{10895} <4974656D2E31323936> \MusicPDFRef{10896}] /Limits [<4974656D2E31323931> <4974656D2E31323936>] >>}
\MusicPDFWriteObject{9219}{<< /Names [<4974656D2E31323937> \MusicPDFRef{10897} <4974656D2E31323938> \MusicPDFRef{10898} <4974656D2E31323939> \MusicPDFRef{10899} <4974656D2E3133> \MusicPDFRef{10900} <4974656D2E313330> \MusicPDFRef{10901} <4974656D2E31333030> \MusicPDFRef{10902}] /Limits [<4974656D2E31323937> <4974656D2E31333030>] >>}
\MusicPDFWriteObject{9220}{<< /Names [<4974656D2E31333031> \MusicPDFRef{10903} <4974656D2E31333032> \MusicPDFRef{10904} <4974656D2E31333033> \MusicPDFRef{10905} <4974656D2E31333034> \MusicPDFRef{10906} <4974656D2E31333035> \MusicPDFRef{10907} <4974656D2E31333036> \MusicPDFRef{10908}] /Limits [<4974656D2E31333031> <4974656D2E31333036>] >>}
\MusicPDFWriteObject{9221}{<< /Names [<4974656D2E31333037> \MusicPDFRef{10909} <4974656D2E31333038> \MusicPDFRef{10910} <4974656D2E31333039> \MusicPDFRef{10911} <4974656D2E313331> \MusicPDFRef{10912} <4974656D2E31333130> \MusicPDFRef{10913} <4974656D2E31333131> \MusicPDFRef{10914}] /Limits [<4974656D2E31333037> <4974656D2E31333131>] >>}
\MusicPDFWriteObject{9222}{<< /Names [<4974656D2E31333132> \MusicPDFRef{10915} <4974656D2E31333133> \MusicPDFRef{10916} <4974656D2E31333134> \MusicPDFRef{10917} <4974656D2E31333135> \MusicPDFRef{10918} <4974656D2E31333136> \MusicPDFRef{10919} <4974656D2E31333137> \MusicPDFRef{10920}] /Limits [<4974656D2E31333132> <4974656D2E31333137>] >>}
\MusicPDFWriteObject{9223}{<< /Names [<4974656D2E31333138> \MusicPDFRef{10921} <4974656D2E31333139> \MusicPDFRef{10922} <4974656D2E313332> \MusicPDFRef{10923} <4974656D2E31333230> \MusicPDFRef{10924} <4974656D2E31333231> \MusicPDFRef{10925} <4974656D2E31333232> \MusicPDFRef{10926}] /Limits [<4974656D2E31333138> <4974656D2E31333232>] >>}
\MusicPDFWriteObject{9224}{<< /Names [<4974656D2E31333233> \MusicPDFRef{10927} <4974656D2E31333234> \MusicPDFRef{10928} <4974656D2E31333235> \MusicPDFRef{10929} <4974656D2E31333236> \MusicPDFRef{10930} <4974656D2E31333237> \MusicPDFRef{10931} <4974656D2E31333238> \MusicPDFRef{10932}] /Limits [<4974656D2E31333233> <4974656D2E31333238>] >>}
\MusicPDFWriteObject{9225}{<< /Names [<4974656D2E31333239> \MusicPDFRef{10933} <4974656D2E313333> \MusicPDFRef{10934} <4974656D2E31333330> \MusicPDFRef{10935} <4974656D2E31333331> \MusicPDFRef{10936} <4974656D2E31333332> \MusicPDFRef{10937} <4974656D2E31333333> \MusicPDFRef{10938}] /Limits [<4974656D2E31333239> <4974656D2E31333333>] >>}
\MusicPDFWriteObject{9226}{<< /Names [<4974656D2E31333334> \MusicPDFRef{10939} <4974656D2E31333335> \MusicPDFRef{10940} <4974656D2E31333336> \MusicPDFRef{10941} <4974656D2E31333337> \MusicPDFRef{10942} <4974656D2E31333338> \MusicPDFRef{10943} <4974656D2E31333339> \MusicPDFRef{10944}] /Limits [<4974656D2E31333334> <4974656D2E31333339>] >>}
\MusicPDFWriteObject{9227}{<< /Names [<4974656D2E313334> \MusicPDFRef{10945} <4974656D2E31333430> \MusicPDFRef{10946} <4974656D2E31333431> \MusicPDFRef{10947} <4974656D2E31333432> \MusicPDFRef{10948} <4974656D2E31333433> \MusicPDFRef{10949} <4974656D2E31333434> \MusicPDFRef{10950}] /Limits [<4974656D2E313334> <4974656D2E31333434>] >>}
\MusicPDFWriteObject{9228}{<< /Names [<4974656D2E31333435> \MusicPDFRef{10951} <4974656D2E31333436> \MusicPDFRef{10952} <4974656D2E31333437> \MusicPDFRef{10953} <4974656D2E31333438> \MusicPDFRef{10954} <4974656D2E31333439> \MusicPDFRef{10955} <4974656D2E313335> \MusicPDFRef{10956}] /Limits [<4974656D2E31333435> <4974656D2E313335>] >>}
\MusicPDFWriteObject{9229}{<< /Names [<4974656D2E31333530> \MusicPDFRef{10957} <4974656D2E31333531> \MusicPDFRef{10958} <4974656D2E31333532> \MusicPDFRef{10959} <4974656D2E31333533> \MusicPDFRef{10960} <4974656D2E31333534> \MusicPDFRef{10961} <4974656D2E31333535> \MusicPDFRef{10962}] /Limits [<4974656D2E31333530> <4974656D2E31333535>] >>}
\MusicPDFWriteObject{9230}{<< /Names [<4974656D2E31333536> \MusicPDFRef{10963} <4974656D2E31333537> \MusicPDFRef{10964} <4974656D2E31333538> \MusicPDFRef{10965} <4974656D2E31333539> \MusicPDFRef{10966} <4974656D2E313336> \MusicPDFRef{10967} <4974656D2E31333630> \MusicPDFRef{10968}] /Limits [<4974656D2E31333536> <4974656D2E31333630>] >>}
\MusicPDFWriteObject{9231}{<< /Names [<4974656D2E31333631> \MusicPDFRef{10969} <4974656D2E31333632> \MusicPDFRef{10970} <4974656D2E31333633> \MusicPDFRef{10971} <4974656D2E31333634> \MusicPDFRef{10972} <4974656D2E31333635> \MusicPDFRef{10973} <4974656D2E31333636> \MusicPDFRef{10974}] /Limits [<4974656D2E31333631> <4974656D2E31333636>] >>}
\MusicPDFWriteObject{9232}{<< /Names [<4974656D2E31333637> \MusicPDFRef{10975} <4974656D2E31333638> \MusicPDFRef{10976} <4974656D2E31333639> \MusicPDFRef{10977} <4974656D2E313337> \MusicPDFRef{10978} <4974656D2E31333730> \MusicPDFRef{10979} <4974656D2E31333731> \MusicPDFRef{10980}] /Limits [<4974656D2E31333637> <4974656D2E31333731>] >>}
\MusicPDFWriteObject{9233}{<< /Names [<4974656D2E31333732> \MusicPDFRef{10981} <4974656D2E31333733> \MusicPDFRef{10982} <4974656D2E31333734> \MusicPDFRef{10983} <4974656D2E31333735> \MusicPDFRef{10984} <4974656D2E31333736> \MusicPDFRef{10985} <4974656D2E31333737> \MusicPDFRef{10986}] /Limits [<4974656D2E31333732> <4974656D2E31333737>] >>}
\MusicPDFWriteObject{9234}{<< /Names [<4974656D2E31333738> \MusicPDFRef{10987} <4974656D2E31333739> \MusicPDFRef{10988} <4974656D2E313338> \MusicPDFRef{10989} <4974656D2E31333830> \MusicPDFRef{10990} <4974656D2E31333831> \MusicPDFRef{10991} <4974656D2E31333832> \MusicPDFRef{10992}] /Limits [<4974656D2E31333738> <4974656D2E31333832>] >>}
\MusicPDFWriteObject{9235}{<< /Names [<4974656D2E31333833> \MusicPDFRef{10993} <4974656D2E31333834> \MusicPDFRef{10994} <4974656D2E31333835> \MusicPDFRef{10995} <4974656D2E31333836> \MusicPDFRef{10996} <4974656D2E31333837> \MusicPDFRef{10997} <4974656D2E31333838> \MusicPDFRef{10998}] /Limits [<4974656D2E31333833> <4974656D2E31333838>] >>}
\MusicPDFWriteObject{9236}{<< /Names [<4974656D2E31333839> \MusicPDFRef{10999} <4974656D2E313339> \MusicPDFRef{11000} <4974656D2E31333930> \MusicPDFRef{11001} <4974656D2E31333931> \MusicPDFRef{11002} <4974656D2E31333932> \MusicPDFRef{11003} <4974656D2E31333933> \MusicPDFRef{11004}] /Limits [<4974656D2E31333839> <4974656D2E31333933>] >>}
\MusicPDFWriteObject{9237}{<< /Names [<4974656D2E31333934> \MusicPDFRef{11005} <4974656D2E31333935> \MusicPDFRef{11006} <4974656D2E31333936> \MusicPDFRef{11007} <4974656D2E31333937> \MusicPDFRef{11008} <4974656D2E31333938> \MusicPDFRef{11009} <4974656D2E31333939> \MusicPDFRef{11010}] /Limits [<4974656D2E31333934> <4974656D2E31333939>] >>}
\MusicPDFWriteObject{9238}{<< /Names [<4974656D2E3134> \MusicPDFRef{11011} <4974656D2E313430> \MusicPDFRef{11012} <4974656D2E31343030> \MusicPDFRef{11013} <4974656D2E31343031> \MusicPDFRef{11014} <4974656D2E31343032> \MusicPDFRef{11015} <4974656D2E31343033> \MusicPDFRef{11016}] /Limits [<4974656D2E3134> <4974656D2E31343033>] >>}
\MusicPDFWriteObject{9239}{<< /Names [<4974656D2E31343034> \MusicPDFRef{11017} <4974656D2E31343035> \MusicPDFRef{11018} <4974656D2E31343036> \MusicPDFRef{11019} <4974656D2E31343037> \MusicPDFRef{11020} <4974656D2E31343038> \MusicPDFRef{11021} <4974656D2E31343039> \MusicPDFRef{11022}] /Limits [<4974656D2E31343034> <4974656D2E31343039>] >>}
\MusicPDFWriteObject{9240}{<< /Names [<4974656D2E313431> \MusicPDFRef{11023} <4974656D2E31343130> \MusicPDFRef{11024} <4974656D2E31343131> \MusicPDFRef{11025} <4974656D2E31343132> \MusicPDFRef{11026} <4974656D2E31343133> \MusicPDFRef{11027} <4974656D2E31343134> \MusicPDFRef{11028}] /Limits [<4974656D2E313431> <4974656D2E31343134>] >>}
\MusicPDFWriteObject{9241}{<< /Names [<4974656D2E31343135> \MusicPDFRef{11029} <4974656D2E31343136> \MusicPDFRef{11030} <4974656D2E31343137> \MusicPDFRef{11031} <4974656D2E31343138> \MusicPDFRef{11032} <4974656D2E31343139> \MusicPDFRef{11033} <4974656D2E313432> \MusicPDFRef{11034}] /Limits [<4974656D2E31343135> <4974656D2E313432>] >>}
\MusicPDFWriteObject{9242}{<< /Names [<4974656D2E31343230> \MusicPDFRef{11035} <4974656D2E31343231> \MusicPDFRef{11036} <4974656D2E31343232> \MusicPDFRef{11037} <4974656D2E31343233> \MusicPDFRef{11038} <4974656D2E31343234> \MusicPDFRef{11039} <4974656D2E31343235> \MusicPDFRef{11040}] /Limits [<4974656D2E31343230> <4974656D2E31343235>] >>}
\MusicPDFWriteObject{9243}{<< /Names [<4974656D2E31343236> \MusicPDFRef{11041} <4974656D2E31343237> \MusicPDFRef{11042} <4974656D2E31343238> \MusicPDFRef{11043} <4974656D2E31343239> \MusicPDFRef{11044} <4974656D2E313433> \MusicPDFRef{11045} <4974656D2E31343330> \MusicPDFRef{11046}] /Limits [<4974656D2E31343236> <4974656D2E31343330>] >>}
\MusicPDFWriteObject{9244}{<< /Names [<4974656D2E31343331> \MusicPDFRef{11047} <4974656D2E31343332> \MusicPDFRef{11048} <4974656D2E31343333> \MusicPDFRef{11049} <4974656D2E31343334> \MusicPDFRef{11050} <4974656D2E31343335> \MusicPDFRef{11051} <4974656D2E31343336> \MusicPDFRef{11052}] /Limits [<4974656D2E31343331> <4974656D2E31343336>] >>}
\MusicPDFWriteObject{9245}{<< /Names [<4974656D2E31343337> \MusicPDFRef{11053} <4974656D2E31343338> \MusicPDFRef{11054} <4974656D2E31343339> \MusicPDFRef{11055} <4974656D2E313434> \MusicPDFRef{11056} <4974656D2E31343430> \MusicPDFRef{11057} <4974656D2E31343431> \MusicPDFRef{11058}] /Limits [<4974656D2E31343337> <4974656D2E31343431>] >>}
\MusicPDFWriteObject{9246}{<< /Names [<4974656D2E31343432> \MusicPDFRef{11059} <4974656D2E31343433> \MusicPDFRef{11060} <4974656D2E31343434> \MusicPDFRef{11061} <4974656D2E31343435> \MusicPDFRef{11062} <4974656D2E31343436> \MusicPDFRef{11063} <4974656D2E31343437> \MusicPDFRef{11064}] /Limits [<4974656D2E31343432> <4974656D2E31343437>] >>}
\MusicPDFWriteObject{9247}{<< /Names [<4974656D2E31343438> \MusicPDFRef{11065} <4974656D2E31343439> \MusicPDFRef{11066} <4974656D2E313435> \MusicPDFRef{11067} <4974656D2E31343530> \MusicPDFRef{11068} <4974656D2E31343531> \MusicPDFRef{11069} <4974656D2E31343532> \MusicPDFRef{11070}] /Limits [<4974656D2E31343438> <4974656D2E31343532>] >>}
\MusicPDFWriteObject{9248}{<< /Names [<4974656D2E31343533> \MusicPDFRef{11071} <4974656D2E31343534> \MusicPDFRef{11072} <4974656D2E31343535> \MusicPDFRef{11073} <4974656D2E31343536> \MusicPDFRef{11074} <4974656D2E31343537> \MusicPDFRef{11075} <4974656D2E31343538> \MusicPDFRef{11076}] /Limits [<4974656D2E31343533> <4974656D2E31343538>] >>}
\MusicPDFWriteObject{9249}{<< /Names [<4974656D2E31343539> \MusicPDFRef{11077} <4974656D2E313436> \MusicPDFRef{11078} <4974656D2E31343630> \MusicPDFRef{11079} <4974656D2E31343631> \MusicPDFRef{11080} <4974656D2E31343632> \MusicPDFRef{11081} <4974656D2E31343633> \MusicPDFRef{11082}] /Limits [<4974656D2E31343539> <4974656D2E31343633>] >>}
\MusicPDFWriteObject{9250}{<< /Names [<4974656D2E31343634> \MusicPDFRef{11083} <4974656D2E31343635> \MusicPDFRef{11084} <4974656D2E31343636> \MusicPDFRef{11085} <4974656D2E31343637> \MusicPDFRef{11086} <4974656D2E31343638> \MusicPDFRef{11087} <4974656D2E31343639> \MusicPDFRef{11088}] /Limits [<4974656D2E31343634> <4974656D2E31343639>] >>}
\MusicPDFWriteObject{9251}{<< /Names [<4974656D2E313437> \MusicPDFRef{11089} <4974656D2E31343730> \MusicPDFRef{11090} <4974656D2E31343731> \MusicPDFRef{11091} <4974656D2E31343732> \MusicPDFRef{11092} <4974656D2E31343733> \MusicPDFRef{11093} <4974656D2E31343734> \MusicPDFRef{11094}] /Limits [<4974656D2E313437> <4974656D2E31343734>] >>}
\MusicPDFWriteObject{9252}{<< /Names [<4974656D2E31343735> \MusicPDFRef{11095} <4974656D2E31343736> \MusicPDFRef{11096} <4974656D2E31343737> \MusicPDFRef{11097} <4974656D2E31343738> \MusicPDFRef{11098} <4974656D2E31343739> \MusicPDFRef{11099} <4974656D2E313438> \MusicPDFRef{11100}] /Limits [<4974656D2E31343735> <4974656D2E313438>] >>}
\MusicPDFWriteObject{9253}{<< /Names [<4974656D2E31343830> \MusicPDFRef{11101} <4974656D2E31343831> \MusicPDFRef{11102} <4974656D2E31343832> \MusicPDFRef{11103} <4974656D2E31343833> \MusicPDFRef{11104} <4974656D2E31343834> \MusicPDFRef{11105} <4974656D2E31343835> \MusicPDFRef{11106}] /Limits [<4974656D2E31343830> <4974656D2E31343835>] >>}
\MusicPDFWriteObject{9254}{<< /Names [<4974656D2E31343836> \MusicPDFRef{11107} <4974656D2E31343837> \MusicPDFRef{11108} <4974656D2E31343838> \MusicPDFRef{11109} <4974656D2E31343839> \MusicPDFRef{11110} <4974656D2E313439> \MusicPDFRef{11111} <4974656D2E31343930> \MusicPDFRef{11112}] /Limits [<4974656D2E31343836> <4974656D2E31343930>] >>}
\MusicPDFWriteObject{9255}{<< /Names [<4974656D2E31343931> \MusicPDFRef{11113} <4974656D2E31343932> \MusicPDFRef{11114} <4974656D2E31343933> \MusicPDFRef{11115} <4974656D2E31343934> \MusicPDFRef{11116} <4974656D2E31343935> \MusicPDFRef{11117} <4974656D2E31343936> \MusicPDFRef{11118}] /Limits [<4974656D2E31343931> <4974656D2E31343936>] >>}
\MusicPDFWriteObject{9256}{<< /Names [<4974656D2E31343937> \MusicPDFRef{11119} <4974656D2E31343938> \MusicPDFRef{11120} <4974656D2E31343939> \MusicPDFRef{11121} <4974656D2E3135> \MusicPDFRef{11122} <4974656D2E313530> \MusicPDFRef{11123} <4974656D2E31353030> \MusicPDFRef{11124}] /Limits [<4974656D2E31343937> <4974656D2E31353030>] >>}
\MusicPDFWriteObject{9257}{<< /Names [<4974656D2E31353031> \MusicPDFRef{11125} <4974656D2E31353032> \MusicPDFRef{11126} <4974656D2E31353033> \MusicPDFRef{11127} <4974656D2E31353034> \MusicPDFRef{11128} <4974656D2E31353035> \MusicPDFRef{11129} <4974656D2E31353036> \MusicPDFRef{11130}] /Limits [<4974656D2E31353031> <4974656D2E31353036>] >>}
\MusicPDFWriteObject{9258}{<< /Names [<4974656D2E31353037> \MusicPDFRef{11131} <4974656D2E31353038> \MusicPDFRef{11132} <4974656D2E31353039> \MusicPDFRef{11133} <4974656D2E313531> \MusicPDFRef{11134} <4974656D2E31353130> \MusicPDFRef{11135} <4974656D2E31353131> \MusicPDFRef{11136}] /Limits [<4974656D2E31353037> <4974656D2E31353131>] >>}
\MusicPDFWriteObject{9259}{<< /Names [<4974656D2E31353132> \MusicPDFRef{11137} <4974656D2E31353133> \MusicPDFRef{11138} <4974656D2E31353134> \MusicPDFRef{11139} <4974656D2E31353135> \MusicPDFRef{11140} <4974656D2E31353136> \MusicPDFRef{11141} <4974656D2E31353137> \MusicPDFRef{11142}] /Limits [<4974656D2E31353132> <4974656D2E31353137>] >>}
\MusicPDFWriteObject{9260}{<< /Names [<4974656D2E31353138> \MusicPDFRef{11143} <4974656D2E31353139> \MusicPDFRef{11144} <4974656D2E313532> \MusicPDFRef{11145} <4974656D2E31353230> \MusicPDFRef{11146} <4974656D2E31353231> \MusicPDFRef{11147} <4974656D2E31353232> \MusicPDFRef{11148}] /Limits [<4974656D2E31353138> <4974656D2E31353232>] >>}
\MusicPDFWriteObject{9261}{<< /Names [<4974656D2E31353233> \MusicPDFRef{11149} <4974656D2E31353234> \MusicPDFRef{11150} <4974656D2E31353235> \MusicPDFRef{11151} <4974656D2E31353236> \MusicPDFRef{11152} <4974656D2E31353237> \MusicPDFRef{11153} <4974656D2E31353238> \MusicPDFRef{11154}] /Limits [<4974656D2E31353233> <4974656D2E31353238>] >>}
\MusicPDFWriteObject{9262}{<< /Names [<4974656D2E31353239> \MusicPDFRef{11155} <4974656D2E313533> \MusicPDFRef{11156} <4974656D2E31353330> \MusicPDFRef{11157} <4974656D2E31353331> \MusicPDFRef{11158} <4974656D2E31353332> \MusicPDFRef{11159} <4974656D2E31353333> \MusicPDFRef{11160}] /Limits [<4974656D2E31353239> <4974656D2E31353333>] >>}
\MusicPDFWriteObject{9263}{<< /Names [<4974656D2E31353334> \MusicPDFRef{11161} <4974656D2E31353335> \MusicPDFRef{11162} <4974656D2E31353336> \MusicPDFRef{11163} <4974656D2E31353337> \MusicPDFRef{11164} <4974656D2E31353338> \MusicPDFRef{11165} <4974656D2E31353339> \MusicPDFRef{11166}] /Limits [<4974656D2E31353334> <4974656D2E31353339>] >>}
\MusicPDFWriteObject{9264}{<< /Names [<4974656D2E313534> \MusicPDFRef{11167} <4974656D2E31353430> \MusicPDFRef{11168} <4974656D2E31353431> \MusicPDFRef{11169} <4974656D2E31353432> \MusicPDFRef{11170} <4974656D2E31353433> \MusicPDFRef{11171} <4974656D2E31353434> \MusicPDFRef{11172}] /Limits [<4974656D2E313534> <4974656D2E31353434>] >>}
\MusicPDFWriteObject{9265}{<< /Names [<4974656D2E31353435> \MusicPDFRef{11173} <4974656D2E31353436> \MusicPDFRef{11174} <4974656D2E31353437> \MusicPDFRef{11175} <4974656D2E31353438> \MusicPDFRef{11176} <4974656D2E31353439> \MusicPDFRef{11177} <4974656D2E313535> \MusicPDFRef{11178}] /Limits [<4974656D2E31353435> <4974656D2E313535>] >>}
\MusicPDFWriteObject{9266}{<< /Names [<4974656D2E31353530> \MusicPDFRef{11179} <4974656D2E31353531> \MusicPDFRef{11180} <4974656D2E31353532> \MusicPDFRef{11181} <4974656D2E31353533> \MusicPDFRef{11182} <4974656D2E31353534> \MusicPDFRef{11183} <4974656D2E31353535> \MusicPDFRef{11184}] /Limits [<4974656D2E31353530> <4974656D2E31353535>] >>}
\MusicPDFWriteObject{9267}{<< /Names [<4974656D2E31353536> \MusicPDFRef{11185} <4974656D2E31353537> \MusicPDFRef{11186} <4974656D2E31353538> \MusicPDFRef{11187} <4974656D2E31353539> \MusicPDFRef{11188} <4974656D2E313536> \MusicPDFRef{11189} <4974656D2E31353630> \MusicPDFRef{11190}] /Limits [<4974656D2E31353536> <4974656D2E31353630>] >>}
\MusicPDFWriteObject{9268}{<< /Names [<4974656D2E31353631> \MusicPDFRef{11191} <4974656D2E31353632> \MusicPDFRef{11192} <4974656D2E31353633> \MusicPDFRef{11193} <4974656D2E31353634> \MusicPDFRef{11194} <4974656D2E31353635> \MusicPDFRef{11195} <4974656D2E31353636> \MusicPDFRef{11196}] /Limits [<4974656D2E31353631> <4974656D2E31353636>] >>}
\MusicPDFWriteObject{9269}{<< /Names [<4974656D2E31353637> \MusicPDFRef{11197} <4974656D2E31353638> \MusicPDFRef{11198} <4974656D2E31353639> \MusicPDFRef{11199} <4974656D2E313537> \MusicPDFRef{11200} <4974656D2E31353730> \MusicPDFRef{11201} <4974656D2E31353731> \MusicPDFRef{11202}] /Limits [<4974656D2E31353637> <4974656D2E31353731>] >>}
\MusicPDFWriteObject{9270}{<< /Names [<4974656D2E31353732> \MusicPDFRef{11203} <4974656D2E31353733> \MusicPDFRef{11204} <4974656D2E31353734> \MusicPDFRef{11205} <4974656D2E31353735> \MusicPDFRef{11206} <4974656D2E31353736> \MusicPDFRef{11207} <4974656D2E31353737> \MusicPDFRef{11208}] /Limits [<4974656D2E31353732> <4974656D2E31353737>] >>}
\MusicPDFWriteObject{9271}{<< /Names [<4974656D2E31353738> \MusicPDFRef{11209} <4974656D2E31353739> \MusicPDFRef{11210} <4974656D2E313538> \MusicPDFRef{11211} <4974656D2E31353830> \MusicPDFRef{11212} <4974656D2E31353831> \MusicPDFRef{11213} <4974656D2E31353832> \MusicPDFRef{11214}] /Limits [<4974656D2E31353738> <4974656D2E31353832>] >>}
\MusicPDFWriteObject{9272}{<< /Names [<4974656D2E31353833> \MusicPDFRef{11215} <4974656D2E31353834> \MusicPDFRef{11216} <4974656D2E31353835> \MusicPDFRef{11217} <4974656D2E31353836> \MusicPDFRef{11218} <4974656D2E31353837> \MusicPDFRef{11219} <4974656D2E31353838> \MusicPDFRef{11220}] /Limits [<4974656D2E31353833> <4974656D2E31353838>] >>}
\MusicPDFWriteObject{9273}{<< /Names [<4974656D2E31353839> \MusicPDFRef{11221} <4974656D2E313539> \MusicPDFRef{11222} <4974656D2E31353930> \MusicPDFRef{11223} <4974656D2E31353931> \MusicPDFRef{11224} <4974656D2E31353932> \MusicPDFRef{11225} <4974656D2E31353933> \MusicPDFRef{11226}] /Limits [<4974656D2E31353839> <4974656D2E31353933>] >>}
\MusicPDFWriteObject{9274}{<< /Names [<4974656D2E31353934> \MusicPDFRef{11227} <4974656D2E31353935> \MusicPDFRef{11228} <4974656D2E31353936> \MusicPDFRef{11229} <4974656D2E31353937> \MusicPDFRef{11230} <4974656D2E31353938> \MusicPDFRef{11231} <4974656D2E31353939> \MusicPDFRef{11232}] /Limits [<4974656D2E31353934> <4974656D2E31353939>] >>}
\MusicPDFWriteObject{9275}{<< /Names [<4974656D2E3136> \MusicPDFRef{11233} <4974656D2E313630> \MusicPDFRef{11234} <4974656D2E31363030> \MusicPDFRef{11235} <4974656D2E31363031> \MusicPDFRef{11236} <4974656D2E313631> \MusicPDFRef{11237} <4974656D2E313632> \MusicPDFRef{11238}] /Limits [<4974656D2E3136> <4974656D2E313632>] >>}
\MusicPDFWriteObject{9276}{<< /Names [<4974656D2E313633> \MusicPDFRef{11239} <4974656D2E313634> \MusicPDFRef{11240} <4974656D2E313635> \MusicPDFRef{11241} <4974656D2E313636> \MusicPDFRef{11242} <4974656D2E313637> \MusicPDFRef{11243} <4974656D2E313638> \MusicPDFRef{11244}] /Limits [<4974656D2E313633> <4974656D2E313638>] >>}
\MusicPDFWriteObject{9277}{<< /Names [<4974656D2E313639> \MusicPDFRef{11245} <4974656D2E3137> \MusicPDFRef{11246} <4974656D2E313730> \MusicPDFRef{11247} <4974656D2E313731> \MusicPDFRef{11248} <4974656D2E313732> \MusicPDFRef{11249} <4974656D2E313733> \MusicPDFRef{11250}] /Limits [<4974656D2E313639> <4974656D2E313733>] >>}
\MusicPDFWriteObject{9278}{<< /Names [<4974656D2E313734> \MusicPDFRef{11251} <4974656D2E313735> \MusicPDFRef{11252} <4974656D2E313736> \MusicPDFRef{11253} <4974656D2E313737> \MusicPDFRef{11254} <4974656D2E313738> \MusicPDFRef{11255} <4974656D2E313739> \MusicPDFRef{11256}] /Limits [<4974656D2E313734> <4974656D2E313739>] >>}
\MusicPDFWriteObject{9279}{<< /Names [<4974656D2E3138> \MusicPDFRef{11257} <4974656D2E313830> \MusicPDFRef{11258} <4974656D2E313831> \MusicPDFRef{11259} <4974656D2E313832> \MusicPDFRef{11260} <4974656D2E313833> \MusicPDFRef{11261} <4974656D2E313834> \MusicPDFRef{11262}] /Limits [<4974656D2E3138> <4974656D2E313834>] >>}
\MusicPDFWriteObject{9280}{<< /Names [<4974656D2E313835> \MusicPDFRef{11263} <4974656D2E313836> \MusicPDFRef{11264} <4974656D2E313837> \MusicPDFRef{11265} <4974656D2E313838> \MusicPDFRef{11266} <4974656D2E313839> \MusicPDFRef{11267} <4974656D2E3139> \MusicPDFRef{11268}] /Limits [<4974656D2E313835> <4974656D2E3139>] >>}
\MusicPDFWriteObject{9281}{<< /Names [<4974656D2E313930> \MusicPDFRef{11269} <4974656D2E313931> \MusicPDFRef{11270} <4974656D2E313932> \MusicPDFRef{11271} <4974656D2E313933> \MusicPDFRef{11272} <4974656D2E313934> \MusicPDFRef{11273} <4974656D2E313935> \MusicPDFRef{11274}] /Limits [<4974656D2E313930> <4974656D2E313935>] >>}
\MusicPDFWriteObject{9282}{<< /Names [<4974656D2E313936> \MusicPDFRef{11275} <4974656D2E313937> \MusicPDFRef{11276} <4974656D2E313938> \MusicPDFRef{11277} <4974656D2E313939> \MusicPDFRef{11278} <4974656D2E32> \MusicPDFRef{11279} <4974656D2E3230> \MusicPDFRef{11280}] /Limits [<4974656D2E313936> <4974656D2E3230>] >>}
\MusicPDFWriteObject{9283}{<< /Names [<4974656D2E323030> \MusicPDFRef{11281} <4974656D2E323031> \MusicPDFRef{11282} <4974656D2E323032> \MusicPDFRef{11283} <4974656D2E323033> \MusicPDFRef{11284} <4974656D2E323034> \MusicPDFRef{11285} <4974656D2E323035> \MusicPDFRef{11286}] /Limits [<4974656D2E323030> <4974656D2E323035>] >>}
\MusicPDFWriteObject{9284}{<< /Names [<4974656D2E323036> \MusicPDFRef{11287} <4974656D2E323037> \MusicPDFRef{11288} <4974656D2E323038> \MusicPDFRef{11289} <4974656D2E323039> \MusicPDFRef{11290} <4974656D2E3231> \MusicPDFRef{11291} <4974656D2E323130> \MusicPDFRef{11292}] /Limits [<4974656D2E323036> <4974656D2E323130>] >>}
\MusicPDFWriteObject{9285}{<< /Names [<4974656D2E323131> \MusicPDFRef{11293} <4974656D2E323132> \MusicPDFRef{11294} <4974656D2E323133> \MusicPDFRef{11295} <4974656D2E323134> \MusicPDFRef{11296} <4974656D2E323135> \MusicPDFRef{11297} <4974656D2E323136> \MusicPDFRef{11298}] /Limits [<4974656D2E323131> <4974656D2E323136>] >>}
\MusicPDFWriteObject{9286}{<< /Names [<4974656D2E323137> \MusicPDFRef{11299} <4974656D2E323138> \MusicPDFRef{11300} <4974656D2E323139> \MusicPDFRef{11301} <4974656D2E3232> \MusicPDFRef{11302} <4974656D2E323230> \MusicPDFRef{11303} <4974656D2E323231> \MusicPDFRef{11304}] /Limits [<4974656D2E323137> <4974656D2E323231>] >>}
\MusicPDFWriteObject{9287}{<< /Names [<4974656D2E323232> \MusicPDFRef{11305} <4974656D2E323233> \MusicPDFRef{11306} <4974656D2E323234> \MusicPDFRef{11307} <4974656D2E323235> \MusicPDFRef{11308} <4974656D2E323236> \MusicPDFRef{11309} <4974656D2E323237> \MusicPDFRef{11310}] /Limits [<4974656D2E323232> <4974656D2E323237>] >>}
\MusicPDFWriteObject{9288}{<< /Names [<4974656D2E323238> \MusicPDFRef{11311} <4974656D2E323239> \MusicPDFRef{11312} <4974656D2E3233> \MusicPDFRef{11313} <4974656D2E323330> \MusicPDFRef{11314} <4974656D2E323331> \MusicPDFRef{11315} <4974656D2E323332> \MusicPDFRef{11316}] /Limits [<4974656D2E323238> <4974656D2E323332>] >>}
\MusicPDFWriteObject{9289}{<< /Names [<4974656D2E323333> \MusicPDFRef{11317} <4974656D2E323334> \MusicPDFRef{11318} <4974656D2E323335> \MusicPDFRef{11319} <4974656D2E323336> \MusicPDFRef{11320} <4974656D2E323337> \MusicPDFRef{11321} <4974656D2E323338> \MusicPDFRef{11322}] /Limits [<4974656D2E323333> <4974656D2E323338>] >>}
\MusicPDFWriteObject{9290}{<< /Names [<4974656D2E323339> \MusicPDFRef{11323} <4974656D2E3234> \MusicPDFRef{11324} <4974656D2E323430> \MusicPDFRef{11325} <4974656D2E323431> \MusicPDFRef{11326} <4974656D2E323432> \MusicPDFRef{11327} <4974656D2E323433> \MusicPDFRef{11328}] /Limits [<4974656D2E323339> <4974656D2E323433>] >>}
\MusicPDFWriteObject{9291}{<< /Names [<4974656D2E323434> \MusicPDFRef{11329} <4974656D2E323435> \MusicPDFRef{11330} <4974656D2E323436> \MusicPDFRef{11331} <4974656D2E323437> \MusicPDFRef{11332} <4974656D2E323438> \MusicPDFRef{11333} <4974656D2E323439> \MusicPDFRef{11334}] /Limits [<4974656D2E323434> <4974656D2E323439>] >>}
\MusicPDFWriteObject{9292}{<< /Names [<4974656D2E3235> \MusicPDFRef{11335} <4974656D2E323530> \MusicPDFRef{11336} <4974656D2E323531> \MusicPDFRef{11337} <4974656D2E323532> \MusicPDFRef{11338} <4974656D2E323533> \MusicPDFRef{11339} <4974656D2E323534> \MusicPDFRef{11340}] /Limits [<4974656D2E3235> <4974656D2E323534>] >>}
\MusicPDFWriteObject{9293}{<< /Names [<4974656D2E323535> \MusicPDFRef{11341} <4974656D2E323536> \MusicPDFRef{11342} <4974656D2E323537> \MusicPDFRef{11343} <4974656D2E323538> \MusicPDFRef{11344} <4974656D2E323539> \MusicPDFRef{11345} <4974656D2E3236> \MusicPDFRef{11346}] /Limits [<4974656D2E323535> <4974656D2E3236>] >>}
\MusicPDFWriteObject{9294}{<< /Names [<4974656D2E323630> \MusicPDFRef{11347} <4974656D2E323631> \MusicPDFRef{11348} <4974656D2E323632> \MusicPDFRef{11349} <4974656D2E323633> \MusicPDFRef{11350} <4974656D2E323634> \MusicPDFRef{11351} <4974656D2E323635> \MusicPDFRef{11352}] /Limits [<4974656D2E323630> <4974656D2E323635>] >>}
\MusicPDFWriteObject{9295}{<< /Names [<4974656D2E323636> \MusicPDFRef{11353} <4974656D2E323637> \MusicPDFRef{11354} <4974656D2E323638> \MusicPDFRef{11355} <4974656D2E323639> \MusicPDFRef{11356} <4974656D2E3237> \MusicPDFRef{11357} <4974656D2E323730> \MusicPDFRef{11358}] /Limits [<4974656D2E323636> <4974656D2E323730>] >>}
\MusicPDFWriteObject{9296}{<< /Names [<4974656D2E323731> \MusicPDFRef{11359} <4974656D2E323732> \MusicPDFRef{11360} <4974656D2E323733> \MusicPDFRef{11361} <4974656D2E323734> \MusicPDFRef{11362} <4974656D2E323735> \MusicPDFRef{11363} <4974656D2E323736> \MusicPDFRef{11364}] /Limits [<4974656D2E323731> <4974656D2E323736>] >>}
\MusicPDFWriteObject{9297}{<< /Names [<4974656D2E323737> \MusicPDFRef{11365} <4974656D2E323738> \MusicPDFRef{11366} <4974656D2E323739> \MusicPDFRef{11367} <4974656D2E3238> \MusicPDFRef{11368} <4974656D2E323830> \MusicPDFRef{11369} <4974656D2E323831> \MusicPDFRef{11370}] /Limits [<4974656D2E323737> <4974656D2E323831>] >>}
\MusicPDFWriteObject{9298}{<< /Names [<4974656D2E323832> \MusicPDFRef{11371} <4974656D2E323833> \MusicPDFRef{11372} <4974656D2E323834> \MusicPDFRef{11373} <4974656D2E323835> \MusicPDFRef{11374} <4974656D2E323836> \MusicPDFRef{11375} <4974656D2E323837> \MusicPDFRef{11376}] /Limits [<4974656D2E323832> <4974656D2E323837>] >>}
\MusicPDFWriteObject{9299}{<< /Names [<4974656D2E323838> \MusicPDFRef{11377} <4974656D2E323839> \MusicPDFRef{11378} <4974656D2E3239> \MusicPDFRef{11379} <4974656D2E323930> \MusicPDFRef{11380} <4974656D2E323931> \MusicPDFRef{11381} <4974656D2E323932> \MusicPDFRef{11382}] /Limits [<4974656D2E323838> <4974656D2E323932>] >>}
\MusicPDFWriteObject{9300}{<< /Names [<4974656D2E323933> \MusicPDFRef{11383} <4974656D2E323934> \MusicPDFRef{11384} <4974656D2E323935> \MusicPDFRef{11385} <4974656D2E323936> \MusicPDFRef{11386} <4974656D2E323937> \MusicPDFRef{11387} <4974656D2E323938> \MusicPDFRef{11388}] /Limits [<4974656D2E323933> <4974656D2E323938>] >>}
\MusicPDFWriteObject{9301}{<< /Names [<4974656D2E323939> \MusicPDFRef{11389} <4974656D2E33> \MusicPDFRef{11390} <4974656D2E3330> \MusicPDFRef{11391} <4974656D2E333030> \MusicPDFRef{11392} <4974656D2E333031> \MusicPDFRef{11393} <4974656D2E333032> \MusicPDFRef{11394}] /Limits [<4974656D2E323939> <4974656D2E333032>] >>}
\MusicPDFWriteObject{9302}{<< /Names [<4974656D2E333033> \MusicPDFRef{11395} <4974656D2E333034> \MusicPDFRef{11396} <4974656D2E333035> \MusicPDFRef{11397} <4974656D2E333036> \MusicPDFRef{11398} <4974656D2E333037> \MusicPDFRef{11399} <4974656D2E333038> \MusicPDFRef{11400}] /Limits [<4974656D2E333033> <4974656D2E333038>] >>}
\MusicPDFWriteObject{9303}{<< /Names [<4974656D2E333039> \MusicPDFRef{11401} <4974656D2E3331> \MusicPDFRef{11402} <4974656D2E333130> \MusicPDFRef{11403} <4974656D2E333131> \MusicPDFRef{11404} <4974656D2E333132> \MusicPDFRef{11405} <4974656D2E333133> \MusicPDFRef{11406}] /Limits [<4974656D2E333039> <4974656D2E333133>] >>}
\MusicPDFWriteObject{9304}{<< /Names [<4974656D2E333134> \MusicPDFRef{11407} <4974656D2E333135> \MusicPDFRef{11408} <4974656D2E333136> \MusicPDFRef{11409} <4974656D2E333137> \MusicPDFRef{11410} <4974656D2E333138> \MusicPDFRef{11411} <4974656D2E333139> \MusicPDFRef{11412}] /Limits [<4974656D2E333134> <4974656D2E333139>] >>}
\MusicPDFWriteObject{9305}{<< /Names [<4974656D2E3332> \MusicPDFRef{11413} <4974656D2E333230> \MusicPDFRef{11414} <4974656D2E333231> \MusicPDFRef{11415} <4974656D2E333232> \MusicPDFRef{11416} <4974656D2E333233> \MusicPDFRef{11417} <4974656D2E333234> \MusicPDFRef{11418}] /Limits [<4974656D2E3332> <4974656D2E333234>] >>}
\MusicPDFWriteObject{9306}{<< /Names [<4974656D2E333235> \MusicPDFRef{11419} <4974656D2E333236> \MusicPDFRef{11420} <4974656D2E333237> \MusicPDFRef{11421} <4974656D2E333238> \MusicPDFRef{11422} <4974656D2E333239> \MusicPDFRef{11423} <4974656D2E3333> \MusicPDFRef{11424}] /Limits [<4974656D2E333235> <4974656D2E3333>] >>}
\MusicPDFWriteObject{9307}{<< /Names [<4974656D2E333330> \MusicPDFRef{11425} <4974656D2E333331> \MusicPDFRef{11426} <4974656D2E333332> \MusicPDFRef{11427} <4974656D2E333333> \MusicPDFRef{11428} <4974656D2E333334> \MusicPDFRef{11429} <4974656D2E333335> \MusicPDFRef{11430}] /Limits [<4974656D2E333330> <4974656D2E333335>] >>}
\MusicPDFWriteObject{9308}{<< /Names [<4974656D2E333336> \MusicPDFRef{11431} <4974656D2E333337> \MusicPDFRef{11432} <4974656D2E333338> \MusicPDFRef{11433} <4974656D2E333339> \MusicPDFRef{11434} <4974656D2E3334> \MusicPDFRef{11435} <4974656D2E333430> \MusicPDFRef{11436}] /Limits [<4974656D2E333336> <4974656D2E333430>] >>}
\MusicPDFWriteObject{9309}{<< /Names [<4974656D2E333431> \MusicPDFRef{11437} <4974656D2E333432> \MusicPDFRef{11438} <4974656D2E333433> \MusicPDFRef{11439} <4974656D2E333434> \MusicPDFRef{11440} <4974656D2E333435> \MusicPDFRef{11441} <4974656D2E333436> \MusicPDFRef{11442}] /Limits [<4974656D2E333431> <4974656D2E333436>] >>}
\MusicPDFWriteObject{9310}{<< /Names [<4974656D2E333437> \MusicPDFRef{11443} <4974656D2E333438> \MusicPDFRef{11444} <4974656D2E333439> \MusicPDFRef{11445} <4974656D2E3335> \MusicPDFRef{11446} <4974656D2E333530> \MusicPDFRef{11447} <4974656D2E333531> \MusicPDFRef{11448}] /Limits [<4974656D2E333437> <4974656D2E333531>] >>}
\MusicPDFWriteObject{9311}{<< /Names [<4974656D2E333532> \MusicPDFRef{11449} <4974656D2E333533> \MusicPDFRef{11450} <4974656D2E333534> \MusicPDFRef{11451} <4974656D2E333535> \MusicPDFRef{11452} <4974656D2E333536> \MusicPDFRef{11453} <4974656D2E333537> \MusicPDFRef{11454}] /Limits [<4974656D2E333532> <4974656D2E333537>] >>}
\MusicPDFWriteObject{9312}{<< /Names [<4974656D2E333538> \MusicPDFRef{11455} <4974656D2E333539> \MusicPDFRef{11456} <4974656D2E3336> \MusicPDFRef{11457} <4974656D2E333630> \MusicPDFRef{11458} <4974656D2E333631> \MusicPDFRef{11459} <4974656D2E333632> \MusicPDFRef{11460}] /Limits [<4974656D2E333538> <4974656D2E333632>] >>}
\MusicPDFWriteObject{9313}{<< /Names [<4974656D2E333633> \MusicPDFRef{11461} <4974656D2E333634> \MusicPDFRef{11462} <4974656D2E333635> \MusicPDFRef{11463} <4974656D2E333636> \MusicPDFRef{11464} <4974656D2E333637> \MusicPDFRef{11465} <4974656D2E333638> \MusicPDFRef{11466}] /Limits [<4974656D2E333633> <4974656D2E333638>] >>}
\MusicPDFWriteObject{9314}{<< /Names [<4974656D2E333639> \MusicPDFRef{11467} <4974656D2E3337> \MusicPDFRef{11468} <4974656D2E333730> \MusicPDFRef{11469} <4974656D2E333731> \MusicPDFRef{11470} <4974656D2E333732> \MusicPDFRef{11471} <4974656D2E333733> \MusicPDFRef{11472}] /Limits [<4974656D2E333639> <4974656D2E333733>] >>}
\MusicPDFWriteObject{9315}{<< /Names [<4974656D2E333734> \MusicPDFRef{11473} <4974656D2E333735> \MusicPDFRef{11474} <4974656D2E333736> \MusicPDFRef{11475} <4974656D2E333737> \MusicPDFRef{11476} <4974656D2E333738> \MusicPDFRef{11477} <4974656D2E333739> \MusicPDFRef{11478}] /Limits [<4974656D2E333734> <4974656D2E333739>] >>}
\MusicPDFWriteObject{9316}{<< /Names [<4974656D2E3338> \MusicPDFRef{11479} <4974656D2E333830> \MusicPDFRef{11480} <4974656D2E333831> \MusicPDFRef{11481} <4974656D2E333832> \MusicPDFRef{11482} <4974656D2E333833> \MusicPDFRef{11483} <4974656D2E333834> \MusicPDFRef{11484}] /Limits [<4974656D2E3338> <4974656D2E333834>] >>}
\MusicPDFWriteObject{9317}{<< /Names [<4974656D2E333835> \MusicPDFRef{11485} <4974656D2E333836> \MusicPDFRef{11486} <4974656D2E333837> \MusicPDFRef{11487} <4974656D2E333838> \MusicPDFRef{11488} <4974656D2E333839> \MusicPDFRef{11489} <4974656D2E3339> \MusicPDFRef{11490}] /Limits [<4974656D2E333835> <4974656D2E3339>] >>}
\MusicPDFWriteObject{9318}{<< /Names [<4974656D2E333930> \MusicPDFRef{11491} <4974656D2E333931> \MusicPDFRef{11492} <4974656D2E333932> \MusicPDFRef{11493} <4974656D2E333933> \MusicPDFRef{11494} <4974656D2E333934> \MusicPDFRef{11495} <4974656D2E333935> \MusicPDFRef{11496}] /Limits [<4974656D2E333930> <4974656D2E333935>] >>}
\MusicPDFWriteObject{9319}{<< /Names [<4974656D2E333936> \MusicPDFRef{11497} <4974656D2E333937> \MusicPDFRef{11498} <4974656D2E333938> \MusicPDFRef{11499} <4974656D2E333939> \MusicPDFRef{11500} <4974656D2E34> \MusicPDFRef{11501} <4974656D2E3430> \MusicPDFRef{11502}] /Limits [<4974656D2E333936> <4974656D2E3430>] >>}
\MusicPDFWriteObject{9320}{<< /Names [<4974656D2E343030> \MusicPDFRef{11503} <4974656D2E343031> \MusicPDFRef{11504} <4974656D2E343032> \MusicPDFRef{11505} <4974656D2E343033> \MusicPDFRef{11506} <4974656D2E343034> \MusicPDFRef{11507} <4974656D2E343035> \MusicPDFRef{11508}] /Limits [<4974656D2E343030> <4974656D2E343035>] >>}
\MusicPDFWriteObject{9321}{<< /Names [<4974656D2E343036> \MusicPDFRef{11509} <4974656D2E343037> \MusicPDFRef{11510} <4974656D2E343038> \MusicPDFRef{11511} <4974656D2E343039> \MusicPDFRef{11512} <4974656D2E3431> \MusicPDFRef{11513} <4974656D2E343130> \MusicPDFRef{11514}] /Limits [<4974656D2E343036> <4974656D2E343130>] >>}
\MusicPDFWriteObject{9322}{<< /Names [<4974656D2E343131> \MusicPDFRef{11515} <4974656D2E343132> \MusicPDFRef{11516} <4974656D2E343133> \MusicPDFRef{11517} <4974656D2E343134> \MusicPDFRef{11518} <4974656D2E343135> \MusicPDFRef{11519} <4974656D2E343136> \MusicPDFRef{11520}] /Limits [<4974656D2E343131> <4974656D2E343136>] >>}
\MusicPDFWriteObject{9323}{<< /Names [<4974656D2E343137> \MusicPDFRef{11521} <4974656D2E343138> \MusicPDFRef{11522} <4974656D2E343139> \MusicPDFRef{11523} <4974656D2E3432> \MusicPDFRef{11524} <4974656D2E343230> \MusicPDFRef{11525} <4974656D2E343231> \MusicPDFRef{11526}] /Limits [<4974656D2E343137> <4974656D2E343231>] >>}
\MusicPDFWriteObject{9324}{<< /Names [<4974656D2E343232> \MusicPDFRef{11527} <4974656D2E343233> \MusicPDFRef{11528} <4974656D2E343234> \MusicPDFRef{11529} <4974656D2E343235> \MusicPDFRef{11530} <4974656D2E343236> \MusicPDFRef{11531} <4974656D2E343237> \MusicPDFRef{11532}] /Limits [<4974656D2E343232> <4974656D2E343237>] >>}
\MusicPDFWriteObject{9325}{<< /Names [<4974656D2E343238> \MusicPDFRef{11533} <4974656D2E343239> \MusicPDFRef{11534} <4974656D2E3433> \MusicPDFRef{11535} <4974656D2E343330> \MusicPDFRef{11536} <4974656D2E343331> \MusicPDFRef{11537} <4974656D2E343332> \MusicPDFRef{11538}] /Limits [<4974656D2E343238> <4974656D2E343332>] >>}
\MusicPDFWriteObject{9326}{<< /Names [<4974656D2E343333> \MusicPDFRef{11539} <4974656D2E343334> \MusicPDFRef{11540} <4974656D2E343335> \MusicPDFRef{11541} <4974656D2E343336> \MusicPDFRef{11542} <4974656D2E343337> \MusicPDFRef{11543} <4974656D2E343338> \MusicPDFRef{11544}] /Limits [<4974656D2E343333> <4974656D2E343338>] >>}
\MusicPDFWriteObject{9327}{<< /Names [<4974656D2E343339> \MusicPDFRef{11545} <4974656D2E3434> \MusicPDFRef{11546} <4974656D2E343430> \MusicPDFRef{11547} <4974656D2E343431> \MusicPDFRef{11548} <4974656D2E343432> \MusicPDFRef{11549} <4974656D2E343433> \MusicPDFRef{11550}] /Limits [<4974656D2E343339> <4974656D2E343433>] >>}
\MusicPDFWriteObject{9328}{<< /Names [<4974656D2E343434> \MusicPDFRef{11551} <4974656D2E343435> \MusicPDFRef{11552} <4974656D2E343436> \MusicPDFRef{11553} <4974656D2E343437> \MusicPDFRef{11554} <4974656D2E343438> \MusicPDFRef{11555} <4974656D2E343439> \MusicPDFRef{11556}] /Limits [<4974656D2E343434> <4974656D2E343439>] >>}
\MusicPDFWriteObject{9329}{<< /Names [<4974656D2E3435> \MusicPDFRef{11557} <4974656D2E343530> \MusicPDFRef{11558} <4974656D2E343531> \MusicPDFRef{11559} <4974656D2E343532> \MusicPDFRef{11560} <4974656D2E343533> \MusicPDFRef{11561} <4974656D2E343534> \MusicPDFRef{11562}] /Limits [<4974656D2E3435> <4974656D2E343534>] >>}
\MusicPDFWriteObject{9330}{<< /Names [<4974656D2E343535> \MusicPDFRef{11563} <4974656D2E343536> \MusicPDFRef{11564} <4974656D2E343537> \MusicPDFRef{11565} <4974656D2E343538> \MusicPDFRef{11566} <4974656D2E343539> \MusicPDFRef{11567} <4974656D2E3436> \MusicPDFRef{11568}] /Limits [<4974656D2E343535> <4974656D2E3436>] >>}
\MusicPDFWriteObject{9331}{<< /Names [<4974656D2E343630> \MusicPDFRef{11569} <4974656D2E343631> \MusicPDFRef{11570} <4974656D2E343632> \MusicPDFRef{11571} <4974656D2E343633> \MusicPDFRef{11572} <4974656D2E343634> \MusicPDFRef{11573} <4974656D2E343635> \MusicPDFRef{11574}] /Limits [<4974656D2E343630> <4974656D2E343635>] >>}
\MusicPDFWriteObject{9332}{<< /Names [<4974656D2E343636> \MusicPDFRef{11575} <4974656D2E343637> \MusicPDFRef{11576} <4974656D2E343638> \MusicPDFRef{11577} <4974656D2E343639> \MusicPDFRef{11578} <4974656D2E3437> \MusicPDFRef{11579} <4974656D2E343730> \MusicPDFRef{11580}] /Limits [<4974656D2E343636> <4974656D2E343730>] >>}
\MusicPDFWriteObject{9333}{<< /Names [<4974656D2E343731> \MusicPDFRef{11581} <4974656D2E343732> \MusicPDFRef{11582} <4974656D2E343733> \MusicPDFRef{11583} <4974656D2E343734> \MusicPDFRef{11584} <4974656D2E343735> \MusicPDFRef{11585} <4974656D2E343736> \MusicPDFRef{11586}] /Limits [<4974656D2E343731> <4974656D2E343736>] >>}
\MusicPDFWriteObject{9334}{<< /Names [<4974656D2E343737> \MusicPDFRef{11587} <4974656D2E343738> \MusicPDFRef{11588} <4974656D2E343739> \MusicPDFRef{11589} <4974656D2E3438> \MusicPDFRef{11590} <4974656D2E343830> \MusicPDFRef{11591} <4974656D2E343831> \MusicPDFRef{11592}] /Limits [<4974656D2E343737> <4974656D2E343831>] >>}
\MusicPDFWriteObject{9335}{<< /Names [<4974656D2E343832> \MusicPDFRef{11593} <4974656D2E343833> \MusicPDFRef{11594} <4974656D2E343834> \MusicPDFRef{11595} <4974656D2E343835> \MusicPDFRef{11596} <4974656D2E343836> \MusicPDFRef{11597} <4974656D2E343837> \MusicPDFRef{11598}] /Limits [<4974656D2E343832> <4974656D2E343837>] >>}
\MusicPDFWriteObject{9336}{<< /Names [<4974656D2E343838> \MusicPDFRef{11599} <4974656D2E343839> \MusicPDFRef{11600} <4974656D2E3439> \MusicPDFRef{11601} <4974656D2E343930> \MusicPDFRef{11602} <4974656D2E343931> \MusicPDFRef{11603} <4974656D2E343932> \MusicPDFRef{11604}] /Limits [<4974656D2E343838> <4974656D2E343932>] >>}
\MusicPDFWriteObject{9337}{<< /Names [<4974656D2E343933> \MusicPDFRef{11605} <4974656D2E343934> \MusicPDFRef{11606} <4974656D2E343935> \MusicPDFRef{11607} <4974656D2E343936> \MusicPDFRef{11608} <4974656D2E343937> \MusicPDFRef{11609} <4974656D2E343938> \MusicPDFRef{11610}] /Limits [<4974656D2E343933> <4974656D2E343938>] >>}
\MusicPDFWriteObject{9338}{<< /Names [<4974656D2E343939> \MusicPDFRef{11611} <4974656D2E35> \MusicPDFRef{11612} <4974656D2E3530> \MusicPDFRef{11613} <4974656D2E353030> \MusicPDFRef{11614} <4974656D2E353031> \MusicPDFRef{11615} <4974656D2E353032> \MusicPDFRef{11616}] /Limits [<4974656D2E343939> <4974656D2E353032>] >>}
\MusicPDFWriteObject{9339}{<< /Names [<4974656D2E353033> \MusicPDFRef{11617} <4974656D2E353034> \MusicPDFRef{11618} <4974656D2E353035> \MusicPDFRef{11619} <4974656D2E353036> \MusicPDFRef{11620} <4974656D2E353037> \MusicPDFRef{11621} <4974656D2E353038> \MusicPDFRef{11622}] /Limits [<4974656D2E353033> <4974656D2E353038>] >>}
\MusicPDFWriteObject{9340}{<< /Names [<4974656D2E353039> \MusicPDFRef{11623} <4974656D2E3531> \MusicPDFRef{11624} <4974656D2E353130> \MusicPDFRef{11625} <4974656D2E353131> \MusicPDFRef{11626} <4974656D2E353132> \MusicPDFRef{11627} <4974656D2E353133> \MusicPDFRef{11628}] /Limits [<4974656D2E353039> <4974656D2E353133>] >>}
\MusicPDFWriteObject{9341}{<< /Names [<4974656D2E353134> \MusicPDFRef{11629} <4974656D2E353135> \MusicPDFRef{11630} <4974656D2E353136> \MusicPDFRef{11631} <4974656D2E353137> \MusicPDFRef{11632} <4974656D2E353138> \MusicPDFRef{11633} <4974656D2E353139> \MusicPDFRef{11634}] /Limits [<4974656D2E353134> <4974656D2E353139>] >>}
\MusicPDFWriteObject{9342}{<< /Names [<4974656D2E3532> \MusicPDFRef{11635} <4974656D2E353230> \MusicPDFRef{11636} <4974656D2E353231> \MusicPDFRef{11637} <4974656D2E353232> \MusicPDFRef{11638} <4974656D2E353233> \MusicPDFRef{11639} <4974656D2E353234> \MusicPDFRef{11640}] /Limits [<4974656D2E3532> <4974656D2E353234>] >>}
\MusicPDFWriteObject{9343}{<< /Names [<4974656D2E353235> \MusicPDFRef{11641} <4974656D2E353236> \MusicPDFRef{11642} <4974656D2E353237> \MusicPDFRef{11643} <4974656D2E353238> \MusicPDFRef{11644} <4974656D2E353239> \MusicPDFRef{11645} <4974656D2E3533> \MusicPDFRef{11646}] /Limits [<4974656D2E353235> <4974656D2E3533>] >>}
\MusicPDFWriteObject{9344}{<< /Names [<4974656D2E353330> \MusicPDFRef{11647} <4974656D2E353331> \MusicPDFRef{11648} <4974656D2E353332> \MusicPDFRef{11649} <4974656D2E353333> \MusicPDFRef{11650} <4974656D2E353334> \MusicPDFRef{11651} <4974656D2E353335> \MusicPDFRef{11652}] /Limits [<4974656D2E353330> <4974656D2E353335>] >>}
\MusicPDFWriteObject{9345}{<< /Names [<4974656D2E353336> \MusicPDFRef{11653} <4974656D2E353337> \MusicPDFRef{11654} <4974656D2E353338> \MusicPDFRef{11655} <4974656D2E353339> \MusicPDFRef{11656} <4974656D2E3534> \MusicPDFRef{11657} <4974656D2E353430> \MusicPDFRef{11658}] /Limits [<4974656D2E353336> <4974656D2E353430>] >>}
\MusicPDFWriteObject{9346}{<< /Names [<4974656D2E353431> \MusicPDFRef{11659} <4974656D2E353432> \MusicPDFRef{11660} <4974656D2E353433> \MusicPDFRef{11661} <4974656D2E353434> \MusicPDFRef{11662} <4974656D2E353435> \MusicPDFRef{11663} <4974656D2E353436> \MusicPDFRef{11664}] /Limits [<4974656D2E353431> <4974656D2E353436>] >>}
\MusicPDFWriteObject{9347}{<< /Names [<4974656D2E353437> \MusicPDFRef{11665} <4974656D2E353438> \MusicPDFRef{11666} <4974656D2E353439> \MusicPDFRef{11667} <4974656D2E3535> \MusicPDFRef{11668} <4974656D2E353530> \MusicPDFRef{11669} <4974656D2E353531> \MusicPDFRef{11670}] /Limits [<4974656D2E353437> <4974656D2E353531>] >>}
\MusicPDFWriteObject{9348}{<< /Names [<4974656D2E353532> \MusicPDFRef{11671} <4974656D2E353533> \MusicPDFRef{11672} <4974656D2E353534> \MusicPDFRef{11673} <4974656D2E353535> \MusicPDFRef{11674} <4974656D2E353536> \MusicPDFRef{11675} <4974656D2E353537> \MusicPDFRef{11676}] /Limits [<4974656D2E353532> <4974656D2E353537>] >>}
\MusicPDFWriteObject{9349}{<< /Names [<4974656D2E353538> \MusicPDFRef{11677} <4974656D2E353539> \MusicPDFRef{11678} <4974656D2E3536> \MusicPDFRef{11679} <4974656D2E353630> \MusicPDFRef{11680} <4974656D2E353631> \MusicPDFRef{11681} <4974656D2E353632> \MusicPDFRef{11682}] /Limits [<4974656D2E353538> <4974656D2E353632>] >>}
\MusicPDFWriteObject{9350}{<< /Names [<4974656D2E353633> \MusicPDFRef{11683} <4974656D2E353634> \MusicPDFRef{11684} <4974656D2E353635> \MusicPDFRef{11685} <4974656D2E353636> \MusicPDFRef{11686} <4974656D2E353637> \MusicPDFRef{11687} <4974656D2E353638> \MusicPDFRef{11688}] /Limits [<4974656D2E353633> <4974656D2E353638>] >>}
\MusicPDFWriteObject{9351}{<< /Names [<4974656D2E353639> \MusicPDFRef{11689} <4974656D2E3537> \MusicPDFRef{11690} <4974656D2E353730> \MusicPDFRef{11691} <4974656D2E353731> \MusicPDFRef{11692} <4974656D2E353732> \MusicPDFRef{11693} <4974656D2E353733> \MusicPDFRef{11694}] /Limits [<4974656D2E353639> <4974656D2E353733>] >>}
\MusicPDFWriteObject{9352}{<< /Names [<4974656D2E353734> \MusicPDFRef{11695} <4974656D2E353735> \MusicPDFRef{11696} <4974656D2E353736> \MusicPDFRef{11697} <4974656D2E353737> \MusicPDFRef{11698} <4974656D2E353738> \MusicPDFRef{11699} <4974656D2E353739> \MusicPDFRef{11700}] /Limits [<4974656D2E353734> <4974656D2E353739>] >>}
\MusicPDFWriteObject{9353}{<< /Names [<4974656D2E3538> \MusicPDFRef{11701} <4974656D2E353830> \MusicPDFRef{11702} <4974656D2E353831> \MusicPDFRef{11703} <4974656D2E353832> \MusicPDFRef{11704} <4974656D2E353833> \MusicPDFRef{11705} <4974656D2E353834> \MusicPDFRef{11706}] /Limits [<4974656D2E3538> <4974656D2E353834>] >>}
\MusicPDFWriteObject{9354}{<< /Names [<4974656D2E353835> \MusicPDFRef{11707} <4974656D2E353836> \MusicPDFRef{11708} <4974656D2E353837> \MusicPDFRef{11709} <4974656D2E353838> \MusicPDFRef{11710} <4974656D2E353839> \MusicPDFRef{11711} <4974656D2E3539> \MusicPDFRef{11712}] /Limits [<4974656D2E353835> <4974656D2E3539>] >>}
\MusicPDFWriteObject{9355}{<< /Names [<4974656D2E353930> \MusicPDFRef{11713} <4974656D2E353931> \MusicPDFRef{11714} <4974656D2E353932> \MusicPDFRef{11715} <4974656D2E353933> \MusicPDFRef{11716} <4974656D2E353934> \MusicPDFRef{11717} <4974656D2E353935> \MusicPDFRef{11718}] /Limits [<4974656D2E353930> <4974656D2E353935>] >>}
\MusicPDFWriteObject{9356}{<< /Names [<4974656D2E353936> \MusicPDFRef{11719} <4974656D2E353937> \MusicPDFRef{11720} <4974656D2E353938> \MusicPDFRef{11721} <4974656D2E353939> \MusicPDFRef{11722} <4974656D2E36> \MusicPDFRef{11723} <4974656D2E3630> \MusicPDFRef{11724}] /Limits [<4974656D2E353936> <4974656D2E3630>] >>}
\MusicPDFWriteObject{9357}{<< /Names [<4974656D2E363030> \MusicPDFRef{11725} <4974656D2E363031> \MusicPDFRef{11726} <4974656D2E363032> \MusicPDFRef{11727} <4974656D2E363033> \MusicPDFRef{11728} <4974656D2E363034> \MusicPDFRef{11729} <4974656D2E363035> \MusicPDFRef{11730}] /Limits [<4974656D2E363030> <4974656D2E363035>] >>}
\MusicPDFWriteObject{9358}{<< /Names [<4974656D2E363036> \MusicPDFRef{11731} <4974656D2E363037> \MusicPDFRef{11732} <4974656D2E363038> \MusicPDFRef{11733} <4974656D2E363039> \MusicPDFRef{11734} <4974656D2E3631> \MusicPDFRef{11735} <4974656D2E363130> \MusicPDFRef{11736}] /Limits [<4974656D2E363036> <4974656D2E363130>] >>}
\MusicPDFWriteObject{9359}{<< /Names [<4974656D2E363131> \MusicPDFRef{11737} <4974656D2E363132> \MusicPDFRef{11738} <4974656D2E363133> \MusicPDFRef{11739} <4974656D2E363134> \MusicPDFRef{11740} <4974656D2E363135> \MusicPDFRef{11741} <4974656D2E363136> \MusicPDFRef{11742}] /Limits [<4974656D2E363131> <4974656D2E363136>] >>}
\MusicPDFWriteObject{9360}{<< /Names [<4974656D2E363137> \MusicPDFRef{11743} <4974656D2E363138> \MusicPDFRef{11744} <4974656D2E363139> \MusicPDFRef{11745} <4974656D2E3632> \MusicPDFRef{11746} <4974656D2E363230> \MusicPDFRef{11747} <4974656D2E363231> \MusicPDFRef{11748}] /Limits [<4974656D2E363137> <4974656D2E363231>] >>}
\MusicPDFWriteObject{9361}{<< /Names [<4974656D2E363232> \MusicPDFRef{11749} <4974656D2E363233> \MusicPDFRef{11750} <4974656D2E363234> \MusicPDFRef{11751} <4974656D2E363235> \MusicPDFRef{11752} <4974656D2E363236> \MusicPDFRef{11753} <4974656D2E363237> \MusicPDFRef{11754}] /Limits [<4974656D2E363232> <4974656D2E363237>] >>}
\MusicPDFWriteObject{9362}{<< /Names [<4974656D2E363238> \MusicPDFRef{11755} <4974656D2E363239> \MusicPDFRef{11756} <4974656D2E3633> \MusicPDFRef{11757} <4974656D2E363330> \MusicPDFRef{11758} <4974656D2E363331> \MusicPDFRef{11759} <4974656D2E363332> \MusicPDFRef{11760}] /Limits [<4974656D2E363238> <4974656D2E363332>] >>}
\MusicPDFWriteObject{9363}{<< /Names [<4974656D2E363333> \MusicPDFRef{11761} <4974656D2E363334> \MusicPDFRef{11762} <4974656D2E363335> \MusicPDFRef{11763} <4974656D2E363336> \MusicPDFRef{11764} <4974656D2E363337> \MusicPDFRef{11765} <4974656D2E363338> \MusicPDFRef{11766}] /Limits [<4974656D2E363333> <4974656D2E363338>] >>}
\MusicPDFWriteObject{9364}{<< /Names [<4974656D2E363339> \MusicPDFRef{11767} <4974656D2E3634> \MusicPDFRef{11768} <4974656D2E363430> \MusicPDFRef{11769} <4974656D2E363431> \MusicPDFRef{11770} <4974656D2E363432> \MusicPDFRef{11771} <4974656D2E363433> \MusicPDFRef{11772}] /Limits [<4974656D2E363339> <4974656D2E363433>] >>}
\MusicPDFWriteObject{9365}{<< /Names [<4974656D2E363434> \MusicPDFRef{11773} <4974656D2E363435> \MusicPDFRef{11774} <4974656D2E363436> \MusicPDFRef{11775} <4974656D2E363437> \MusicPDFRef{11776} <4974656D2E363438> \MusicPDFRef{11777} <4974656D2E363439> \MusicPDFRef{11778}] /Limits [<4974656D2E363434> <4974656D2E363439>] >>}
\MusicPDFWriteObject{9366}{<< /Names [<4974656D2E3635> \MusicPDFRef{11779} <4974656D2E363530> \MusicPDFRef{11780} <4974656D2E363531> \MusicPDFRef{11781} <4974656D2E363532> \MusicPDFRef{11782} <4974656D2E363533> \MusicPDFRef{11783} <4974656D2E363534> \MusicPDFRef{11784}] /Limits [<4974656D2E3635> <4974656D2E363534>] >>}
\MusicPDFWriteObject{9367}{<< /Names [<4974656D2E363535> \MusicPDFRef{11785} <4974656D2E363536> \MusicPDFRef{11786} <4974656D2E363537> \MusicPDFRef{11787} <4974656D2E363538> \MusicPDFRef{11788} <4974656D2E363539> \MusicPDFRef{11789} <4974656D2E3636> \MusicPDFRef{11790}] /Limits [<4974656D2E363535> <4974656D2E3636>] >>}
\MusicPDFWriteObject{9368}{<< /Names [<4974656D2E363630> \MusicPDFRef{11791} <4974656D2E363631> \MusicPDFRef{11792} <4974656D2E363632> \MusicPDFRef{11793} <4974656D2E363633> \MusicPDFRef{11794} <4974656D2E363634> \MusicPDFRef{11795} <4974656D2E363635> \MusicPDFRef{11796}] /Limits [<4974656D2E363630> <4974656D2E363635>] >>}
\MusicPDFWriteObject{9369}{<< /Names [<4974656D2E363636> \MusicPDFRef{11797} <4974656D2E363637> \MusicPDFRef{11798} <4974656D2E363638> \MusicPDFRef{11799} <4974656D2E363639> \MusicPDFRef{11800} <4974656D2E3637> \MusicPDFRef{11801} <4974656D2E363730> \MusicPDFRef{11802}] /Limits [<4974656D2E363636> <4974656D2E363730>] >>}
\MusicPDFWriteObject{9370}{<< /Names [<4974656D2E363731> \MusicPDFRef{11803} <4974656D2E363732> \MusicPDFRef{11804} <4974656D2E363733> \MusicPDFRef{11805} <4974656D2E363734> \MusicPDFRef{11806} <4974656D2E363735> \MusicPDFRef{11807} <4974656D2E363736> \MusicPDFRef{11808}] /Limits [<4974656D2E363731> <4974656D2E363736>] >>}
\MusicPDFWriteObject{9371}{<< /Names [<4974656D2E363737> \MusicPDFRef{11809} <4974656D2E363738> \MusicPDFRef{11810} <4974656D2E363739> \MusicPDFRef{11811} <4974656D2E3638> \MusicPDFRef{11812} <4974656D2E363830> \MusicPDFRef{11813} <4974656D2E363831> \MusicPDFRef{11814}] /Limits [<4974656D2E363737> <4974656D2E363831>] >>}
\MusicPDFWriteObject{9372}{<< /Names [<4974656D2E363832> \MusicPDFRef{11815} <4974656D2E363833> \MusicPDFRef{11816} <4974656D2E363834> \MusicPDFRef{11817} <4974656D2E363835> \MusicPDFRef{11818} <4974656D2E363836> \MusicPDFRef{11819} <4974656D2E363837> \MusicPDFRef{11820}] /Limits [<4974656D2E363832> <4974656D2E363837>] >>}
\MusicPDFWriteObject{9373}{<< /Names [<4974656D2E363838> \MusicPDFRef{11821} <4974656D2E363839> \MusicPDFRef{11822} <4974656D2E3639> \MusicPDFRef{11823} <4974656D2E363930> \MusicPDFRef{11824} <4974656D2E363931> \MusicPDFRef{11825} <4974656D2E363932> \MusicPDFRef{11826}] /Limits [<4974656D2E363838> <4974656D2E363932>] >>}
\MusicPDFWriteObject{9374}{<< /Names [<4974656D2E363933> \MusicPDFRef{11827} <4974656D2E363934> \MusicPDFRef{11828} <4974656D2E363935> \MusicPDFRef{11829} <4974656D2E363936> \MusicPDFRef{11830} <4974656D2E363937> \MusicPDFRef{11831} <4974656D2E363938> \MusicPDFRef{11832}] /Limits [<4974656D2E363933> <4974656D2E363938>] >>}
\MusicPDFWriteObject{9375}{<< /Names [<4974656D2E363939> \MusicPDFRef{11833} <4974656D2E37> \MusicPDFRef{11834} <4974656D2E3730> \MusicPDFRef{11835} <4974656D2E373030> \MusicPDFRef{11836} <4974656D2E373031> \MusicPDFRef{11837} <4974656D2E373032> \MusicPDFRef{11838}] /Limits [<4974656D2E363939> <4974656D2E373032>] >>}
\MusicPDFWriteObject{9376}{<< /Names [<4974656D2E373033> \MusicPDFRef{11839} <4974656D2E373034> \MusicPDFRef{11840} <4974656D2E373035> \MusicPDFRef{11841} <4974656D2E373036> \MusicPDFRef{11842} <4974656D2E373037> \MusicPDFRef{11843} <4974656D2E373038> \MusicPDFRef{11844}] /Limits [<4974656D2E373033> <4974656D2E373038>] >>}
\MusicPDFWriteObject{9377}{<< /Names [<4974656D2E373039> \MusicPDFRef{11845} <4974656D2E3731> \MusicPDFRef{11846} <4974656D2E373130> \MusicPDFRef{11847} <4974656D2E373131> \MusicPDFRef{11848} <4974656D2E373132> \MusicPDFRef{11849} <4974656D2E373133> \MusicPDFRef{11850}] /Limits [<4974656D2E373039> <4974656D2E373133>] >>}
\MusicPDFWriteObject{9378}{<< /Names [<4974656D2E373134> \MusicPDFRef{11851} <4974656D2E373135> \MusicPDFRef{11852} <4974656D2E373136> \MusicPDFRef{11853} <4974656D2E373137> \MusicPDFRef{11854} <4974656D2E373138> \MusicPDFRef{11855} <4974656D2E373139> \MusicPDFRef{11856}] /Limits [<4974656D2E373134> <4974656D2E373139>] >>}
\MusicPDFWriteObject{9379}{<< /Names [<4974656D2E3732> \MusicPDFRef{11857} <4974656D2E373230> \MusicPDFRef{11858} <4974656D2E373231> \MusicPDFRef{11859} <4974656D2E373232> \MusicPDFRef{11860} <4974656D2E373233> \MusicPDFRef{11861} <4974656D2E373234> \MusicPDFRef{11862}] /Limits [<4974656D2E3732> <4974656D2E373234>] >>}
\MusicPDFWriteObject{9380}{<< /Names [<4974656D2E373235> \MusicPDFRef{11863} <4974656D2E373236> \MusicPDFRef{11864} <4974656D2E373237> \MusicPDFRef{11865} <4974656D2E373238> \MusicPDFRef{11866} <4974656D2E373239> \MusicPDFRef{11867} <4974656D2E3733> \MusicPDFRef{11868}] /Limits [<4974656D2E373235> <4974656D2E3733>] >>}
\MusicPDFWriteObject{9381}{<< /Names [<4974656D2E373330> \MusicPDFRef{11869} <4974656D2E373331> \MusicPDFRef{11870} <4974656D2E373332> \MusicPDFRef{11871} <4974656D2E373333> \MusicPDFRef{11872} <4974656D2E373334> \MusicPDFRef{11873} <4974656D2E373335> \MusicPDFRef{11874}] /Limits [<4974656D2E373330> <4974656D2E373335>] >>}
\MusicPDFWriteObject{9382}{<< /Names [<4974656D2E373336> \MusicPDFRef{11875} <4974656D2E373337> \MusicPDFRef{11876} <4974656D2E373338> \MusicPDFRef{11877} <4974656D2E373339> \MusicPDFRef{11878} <4974656D2E3734> \MusicPDFRef{11879} <4974656D2E373430> \MusicPDFRef{11880}] /Limits [<4974656D2E373336> <4974656D2E373430>] >>}
\MusicPDFWriteObject{9383}{<< /Names [<4974656D2E373431> \MusicPDFRef{11881} <4974656D2E373432> \MusicPDFRef{11882} <4974656D2E373433> \MusicPDFRef{11883} <4974656D2E373434> \MusicPDFRef{11884} <4974656D2E373435> \MusicPDFRef{11885} <4974656D2E373436> \MusicPDFRef{11886}] /Limits [<4974656D2E373431> <4974656D2E373436>] >>}
\MusicPDFWriteObject{9384}{<< /Names [<4974656D2E373437> \MusicPDFRef{11887} <4974656D2E373438> \MusicPDFRef{11888} <4974656D2E373439> \MusicPDFRef{11889} <4974656D2E3735> \MusicPDFRef{11890} <4974656D2E373530> \MusicPDFRef{11891} <4974656D2E373531> \MusicPDFRef{11892}] /Limits [<4974656D2E373437> <4974656D2E373531>] >>}
\MusicPDFWriteObject{9385}{<< /Names [<4974656D2E373532> \MusicPDFRef{11893} <4974656D2E373533> \MusicPDFRef{11894} <4974656D2E373534> \MusicPDFRef{11895} <4974656D2E373535> \MusicPDFRef{11896} <4974656D2E373536> \MusicPDFRef{11897} <4974656D2E373537> \MusicPDFRef{11898}] /Limits [<4974656D2E373532> <4974656D2E373537>] >>}
\MusicPDFWriteObject{9386}{<< /Names [<4974656D2E373538> \MusicPDFRef{11899} <4974656D2E373539> \MusicPDFRef{11900} <4974656D2E3736> \MusicPDFRef{11901} <4974656D2E373630> \MusicPDFRef{11902} <4974656D2E373631> \MusicPDFRef{11903} <4974656D2E373632> \MusicPDFRef{11904}] /Limits [<4974656D2E373538> <4974656D2E373632>] >>}
\MusicPDFWriteObject{9387}{<< /Names [<4974656D2E373633> \MusicPDFRef{11905} <4974656D2E373634> \MusicPDFRef{11906} <4974656D2E373635> \MusicPDFRef{11907} <4974656D2E373636> \MusicPDFRef{11908} <4974656D2E373637> \MusicPDFRef{11909} <4974656D2E373638> \MusicPDFRef{11910}] /Limits [<4974656D2E373633> <4974656D2E373638>] >>}
\MusicPDFWriteObject{9388}{<< /Names [<4974656D2E373639> \MusicPDFRef{11911} <4974656D2E3737> \MusicPDFRef{11912} <4974656D2E373730> \MusicPDFRef{11913} <4974656D2E373731> \MusicPDFRef{11914} <4974656D2E373732> \MusicPDFRef{11915} <4974656D2E373733> \MusicPDFRef{11916}] /Limits [<4974656D2E373639> <4974656D2E373733>] >>}
\MusicPDFWriteObject{9389}{<< /Names [<4974656D2E373734> \MusicPDFRef{11917} <4974656D2E373735> \MusicPDFRef{11918} <4974656D2E373736> \MusicPDFRef{11919} <4974656D2E373737> \MusicPDFRef{11920} <4974656D2E373738> \MusicPDFRef{11921} <4974656D2E373739> \MusicPDFRef{11922}] /Limits [<4974656D2E373734> <4974656D2E373739>] >>}
\MusicPDFWriteObject{9390}{<< /Names [<4974656D2E3738> \MusicPDFRef{11923} <4974656D2E373830> \MusicPDFRef{11924} <4974656D2E373831> \MusicPDFRef{11925} <4974656D2E373832> \MusicPDFRef{11926} <4974656D2E373833> \MusicPDFRef{11927} <4974656D2E373834> \MusicPDFRef{11928}] /Limits [<4974656D2E3738> <4974656D2E373834>] >>}
\MusicPDFWriteObject{9391}{<< /Names [<4974656D2E373835> \MusicPDFRef{11929} <4974656D2E373836> \MusicPDFRef{11930} <4974656D2E373837> \MusicPDFRef{11931} <4974656D2E373838> \MusicPDFRef{11932} <4974656D2E373839> \MusicPDFRef{11933} <4974656D2E3739> \MusicPDFRef{11934}] /Limits [<4974656D2E373835> <4974656D2E3739>] >>}
\MusicPDFWriteObject{9392}{<< /Names [<4974656D2E373930> \MusicPDFRef{11935} <4974656D2E373931> \MusicPDFRef{11936} <4974656D2E373932> \MusicPDFRef{11937} <4974656D2E373933> \MusicPDFRef{11938} <4974656D2E373934> \MusicPDFRef{11939} <4974656D2E373935> \MusicPDFRef{11940}] /Limits [<4974656D2E373930> <4974656D2E373935>] >>}
\MusicPDFWriteObject{9393}{<< /Names [<4974656D2E373936> \MusicPDFRef{11941} <4974656D2E373937> \MusicPDFRef{11942} <4974656D2E373938> \MusicPDFRef{11943} <4974656D2E373939> \MusicPDFRef{11944} <4974656D2E38> \MusicPDFRef{11945} <4974656D2E3830> \MusicPDFRef{11946}] /Limits [<4974656D2E373936> <4974656D2E3830>] >>}
\MusicPDFWriteObject{9394}{<< /Names [<4974656D2E383030> \MusicPDFRef{11947} <4974656D2E383031> \MusicPDFRef{11948} <4974656D2E383032> \MusicPDFRef{11949} <4974656D2E383033> \MusicPDFRef{11950} <4974656D2E383034> \MusicPDFRef{11951} <4974656D2E383035> \MusicPDFRef{11952}] /Limits [<4974656D2E383030> <4974656D2E383035>] >>}
\MusicPDFWriteObject{9395}{<< /Names [<4974656D2E383036> \MusicPDFRef{11953} <4974656D2E383037> \MusicPDFRef{11954} <4974656D2E383038> \MusicPDFRef{11955} <4974656D2E383039> \MusicPDFRef{11956} <4974656D2E3831> \MusicPDFRef{11957} <4974656D2E383130> \MusicPDFRef{11958}] /Limits [<4974656D2E383036> <4974656D2E383130>] >>}
\MusicPDFWriteObject{9396}{<< /Names [<4974656D2E383131> \MusicPDFRef{11959} <4974656D2E383132> \MusicPDFRef{11960} <4974656D2E383133> \MusicPDFRef{11961} <4974656D2E383134> \MusicPDFRef{11962} <4974656D2E383135> \MusicPDFRef{11963} <4974656D2E383136> \MusicPDFRef{11964}] /Limits [<4974656D2E383131> <4974656D2E383136>] >>}
\MusicPDFWriteObject{9397}{<< /Names [<4974656D2E383137> \MusicPDFRef{11965} <4974656D2E383138> \MusicPDFRef{11966} <4974656D2E383139> \MusicPDFRef{11967} <4974656D2E3832> \MusicPDFRef{11968} <4974656D2E383230> \MusicPDFRef{11969} <4974656D2E383231> \MusicPDFRef{11970}] /Limits [<4974656D2E383137> <4974656D2E383231>] >>}
\MusicPDFWriteObject{9398}{<< /Names [<4974656D2E383232> \MusicPDFRef{11971} <4974656D2E383233> \MusicPDFRef{11972} <4974656D2E383234> \MusicPDFRef{11973} <4974656D2E383235> \MusicPDFRef{11974} <4974656D2E383236> \MusicPDFRef{11975} <4974656D2E383237> \MusicPDFRef{11976}] /Limits [<4974656D2E383232> <4974656D2E383237>] >>}
\MusicPDFWriteObject{9399}{<< /Names [<4974656D2E383238> \MusicPDFRef{11977} <4974656D2E383239> \MusicPDFRef{11978} <4974656D2E3833> \MusicPDFRef{11979} <4974656D2E383330> \MusicPDFRef{11980} <4974656D2E383331> \MusicPDFRef{11981} <4974656D2E383332> \MusicPDFRef{11982}] /Limits [<4974656D2E383238> <4974656D2E383332>] >>}
\MusicPDFWriteObject{9400}{<< /Names [<4974656D2E383333> \MusicPDFRef{11983} <4974656D2E383334> \MusicPDFRef{11984} <4974656D2E383335> \MusicPDFRef{11985} <4974656D2E383336> \MusicPDFRef{11986} <4974656D2E383337> \MusicPDFRef{11987} <4974656D2E383338> \MusicPDFRef{11988}] /Limits [<4974656D2E383333> <4974656D2E383338>] >>}
\MusicPDFWriteObject{9401}{<< /Names [<4974656D2E383339> \MusicPDFRef{11989} <4974656D2E3834> \MusicPDFRef{11990} <4974656D2E383430> \MusicPDFRef{11991} <4974656D2E383431> \MusicPDFRef{11992} <4974656D2E383432> \MusicPDFRef{11993} <4974656D2E383433> \MusicPDFRef{11994}] /Limits [<4974656D2E383339> <4974656D2E383433>] >>}
\MusicPDFWriteObject{9402}{<< /Names [<4974656D2E383434> \MusicPDFRef{11995} <4974656D2E383435> \MusicPDFRef{11996} <4974656D2E383436> \MusicPDFRef{11997} <4974656D2E383437> \MusicPDFRef{11998} <4974656D2E383438> \MusicPDFRef{11999} <4974656D2E383439> \MusicPDFRef{12000}] /Limits [<4974656D2E383434> <4974656D2E383439>] >>}
\MusicPDFWriteObject{9403}{<< /Names [<4974656D2E3835> \MusicPDFRef{12001} <4974656D2E383530> \MusicPDFRef{12002} <4974656D2E383531> \MusicPDFRef{12003} <4974656D2E383532> \MusicPDFRef{12004} <4974656D2E383533> \MusicPDFRef{12005} <4974656D2E383534> \MusicPDFRef{12006}] /Limits [<4974656D2E3835> <4974656D2E383534>] >>}
\MusicPDFWriteObject{9404}{<< /Names [<4974656D2E383535> \MusicPDFRef{12007} <4974656D2E383536> \MusicPDFRef{12008} <4974656D2E383537> \MusicPDFRef{12009} <4974656D2E383538> \MusicPDFRef{12010} <4974656D2E383539> \MusicPDFRef{12011} <4974656D2E3836> \MusicPDFRef{12012}] /Limits [<4974656D2E383535> <4974656D2E3836>] >>}
\MusicPDFWriteObject{9405}{<< /Names [<4974656D2E383630> \MusicPDFRef{12013} <4974656D2E383631> \MusicPDFRef{12014} <4974656D2E383632> \MusicPDFRef{12015} <4974656D2E383633> \MusicPDFRef{12016} <4974656D2E383634> \MusicPDFRef{12017} <4974656D2E383635> \MusicPDFRef{12018}] /Limits [<4974656D2E383630> <4974656D2E383635>] >>}
\MusicPDFWriteObject{9406}{<< /Names [<4974656D2E383636> \MusicPDFRef{12019} <4974656D2E383637> \MusicPDFRef{12020} <4974656D2E383638> \MusicPDFRef{12021} <4974656D2E383639> \MusicPDFRef{12022} <4974656D2E3837> \MusicPDFRef{12023} <4974656D2E383730> \MusicPDFRef{12024}] /Limits [<4974656D2E383636> <4974656D2E383730>] >>}
\MusicPDFWriteObject{9407}{<< /Names [<4974656D2E383731> \MusicPDFRef{12025} <4974656D2E383732> \MusicPDFRef{12026} <4974656D2E383733> \MusicPDFRef{12027} <4974656D2E383734> \MusicPDFRef{12028} <4974656D2E383735> \MusicPDFRef{12029} <4974656D2E383736> \MusicPDFRef{12030}] /Limits [<4974656D2E383731> <4974656D2E383736>] >>}
\MusicPDFWriteObject{9408}{<< /Names [<4974656D2E383737> \MusicPDFRef{12031} <4974656D2E383738> \MusicPDFRef{12032} <4974656D2E383739> \MusicPDFRef{12033} <4974656D2E3838> \MusicPDFRef{12034} <4974656D2E383830> \MusicPDFRef{12035} <4974656D2E383831> \MusicPDFRef{12036}] /Limits [<4974656D2E383737> <4974656D2E383831>] >>}
\MusicPDFWriteObject{9409}{<< /Names [<4974656D2E383832> \MusicPDFRef{12037} <4974656D2E383833> \MusicPDFRef{12038} <4974656D2E383834> \MusicPDFRef{12039} <4974656D2E383835> \MusicPDFRef{12040} <4974656D2E383836> \MusicPDFRef{12041} <4974656D2E383837> \MusicPDFRef{12042}] /Limits [<4974656D2E383832> <4974656D2E383837>] >>}
\MusicPDFWriteObject{9410}{<< /Names [<4974656D2E383838> \MusicPDFRef{12043} <4974656D2E383839> \MusicPDFRef{12044} <4974656D2E3839> \MusicPDFRef{12045} <4974656D2E383930> \MusicPDFRef{12046} <4974656D2E383931> \MusicPDFRef{12047} <4974656D2E383932> \MusicPDFRef{12048}] /Limits [<4974656D2E383838> <4974656D2E383932>] >>}
\MusicPDFWriteObject{9411}{<< /Names [<4974656D2E383933> \MusicPDFRef{12049} <4974656D2E383934> \MusicPDFRef{12050} <4974656D2E383935> \MusicPDFRef{12051} <4974656D2E383936> \MusicPDFRef{12052} <4974656D2E383937> \MusicPDFRef{12053} <4974656D2E383938> \MusicPDFRef{12054}] /Limits [<4974656D2E383933> <4974656D2E383938>] >>}
\MusicPDFWriteObject{9412}{<< /Names [<4974656D2E383939> \MusicPDFRef{12055} <4974656D2E39> \MusicPDFRef{12056} <4974656D2E3930> \MusicPDFRef{12057} <4974656D2E393030> \MusicPDFRef{12058} <4974656D2E393031> \MusicPDFRef{12059} <4974656D2E393032> \MusicPDFRef{12060}] /Limits [<4974656D2E383939> <4974656D2E393032>] >>}
\MusicPDFWriteObject{9413}{<< /Names [<4974656D2E393033> \MusicPDFRef{12061} <4974656D2E393034> \MusicPDFRef{12062} <4974656D2E393035> \MusicPDFRef{12063} <4974656D2E393036> \MusicPDFRef{12064} <4974656D2E393037> \MusicPDFRef{12065} <4974656D2E393038> \MusicPDFRef{12066}] /Limits [<4974656D2E393033> <4974656D2E393038>] >>}
\MusicPDFWriteObject{9414}{<< /Names [<4974656D2E393039> \MusicPDFRef{12067} <4974656D2E3931> \MusicPDFRef{12068} <4974656D2E393130> \MusicPDFRef{12069} <4974656D2E393131> \MusicPDFRef{12070} <4974656D2E393132> \MusicPDFRef{12071} <4974656D2E393133> \MusicPDFRef{12072}] /Limits [<4974656D2E393039> <4974656D2E393133>] >>}
\MusicPDFWriteObject{9415}{<< /Names [<4974656D2E393134> \MusicPDFRef{12073} <4974656D2E393135> \MusicPDFRef{12074} <4974656D2E393136> \MusicPDFRef{12075} <4974656D2E393137> \MusicPDFRef{12076} <4974656D2E393138> \MusicPDFRef{12077} <4974656D2E393139> \MusicPDFRef{12078}] /Limits [<4974656D2E393134> <4974656D2E393139>] >>}
\MusicPDFWriteObject{9416}{<< /Names [<4974656D2E3932> \MusicPDFRef{12079} <4974656D2E393230> \MusicPDFRef{12080} <4974656D2E393231> \MusicPDFRef{12081} <4974656D2E393232> \MusicPDFRef{12082} <4974656D2E393233> \MusicPDFRef{12083} <4974656D2E393234> \MusicPDFRef{12084}] /Limits [<4974656D2E3932> <4974656D2E393234>] >>}
\MusicPDFWriteObject{9417}{<< /Names [<4974656D2E393235> \MusicPDFRef{12085} <4974656D2E393236> \MusicPDFRef{12086} <4974656D2E393237> \MusicPDFRef{12087} <4974656D2E393238> \MusicPDFRef{12088} <4974656D2E393239> \MusicPDFRef{12089} <4974656D2E3933> \MusicPDFRef{12090}] /Limits [<4974656D2E393235> <4974656D2E3933>] >>}
\MusicPDFWriteObject{9418}{<< /Names [<4974656D2E393330> \MusicPDFRef{12091} <4974656D2E393331> \MusicPDFRef{12092} <4974656D2E393332> \MusicPDFRef{12093} <4974656D2E393333> \MusicPDFRef{12094} <4974656D2E393334> \MusicPDFRef{12095} <4974656D2E393335> \MusicPDFRef{12096}] /Limits [<4974656D2E393330> <4974656D2E393335>] >>}
\MusicPDFWriteObject{9419}{<< /Names [<4974656D2E393336> \MusicPDFRef{12097} <4974656D2E393337> \MusicPDFRef{12098} <4974656D2E393338> \MusicPDFRef{12099} <4974656D2E393339> \MusicPDFRef{12100} <4974656D2E3934> \MusicPDFRef{12101} <4974656D2E393430> \MusicPDFRef{12102}] /Limits [<4974656D2E393336> <4974656D2E393430>] >>}
\MusicPDFWriteObject{9420}{<< /Names [<4974656D2E393431> \MusicPDFRef{12103} <4974656D2E393432> \MusicPDFRef{12104} <4974656D2E393433> \MusicPDFRef{12105} <4974656D2E393434> \MusicPDFRef{12106} <4974656D2E393435> \MusicPDFRef{12107} <4974656D2E393436> \MusicPDFRef{12108}] /Limits [<4974656D2E393431> <4974656D2E393436>] >>}
\MusicPDFWriteObject{9421}{<< /Names [<4974656D2E393437> \MusicPDFRef{12109} <4974656D2E393438> \MusicPDFRef{12110} <4974656D2E393439> \MusicPDFRef{12111} <4974656D2E3935> \MusicPDFRef{12112} <4974656D2E393530> \MusicPDFRef{12113} <4974656D2E393531> \MusicPDFRef{12114}] /Limits [<4974656D2E393437> <4974656D2E393531>] >>}
\MusicPDFWriteObject{9422}{<< /Names [<4974656D2E393532> \MusicPDFRef{12115} <4974656D2E393533> \MusicPDFRef{12116} <4974656D2E393534> \MusicPDFRef{12117} <4974656D2E393535> \MusicPDFRef{12118} <4974656D2E393536> \MusicPDFRef{12119} <4974656D2E393537> \MusicPDFRef{12120}] /Limits [<4974656D2E393532> <4974656D2E393537>] >>}
\MusicPDFWriteObject{9423}{<< /Names [<4974656D2E393538> \MusicPDFRef{12121} <4974656D2E393539> \MusicPDFRef{12122} <4974656D2E3936> \MusicPDFRef{12123} <4974656D2E393630> \MusicPDFRef{12124} <4974656D2E393631> \MusicPDFRef{12125} <4974656D2E393632> \MusicPDFRef{12126}] /Limits [<4974656D2E393538> <4974656D2E393632>] >>}
\MusicPDFWriteObject{9424}{<< /Names [<4974656D2E393633> \MusicPDFRef{12127} <4974656D2E393634> \MusicPDFRef{12128} <4974656D2E393635> \MusicPDFRef{12129} <4974656D2E393636> \MusicPDFRef{12130} <4974656D2E393637> \MusicPDFRef{12131} <4974656D2E393638> \MusicPDFRef{12132}] /Limits [<4974656D2E393633> <4974656D2E393638>] >>}
\MusicPDFWriteObject{9425}{<< /Names [<4974656D2E393639> \MusicPDFRef{12133} <4974656D2E3937> \MusicPDFRef{12134} <4974656D2E393730> \MusicPDFRef{12135} <4974656D2E393731> \MusicPDFRef{12136} <4974656D2E393732> \MusicPDFRef{12137} <4974656D2E393733> \MusicPDFRef{12138}] /Limits [<4974656D2E393639> <4974656D2E393733>] >>}
\MusicPDFWriteObject{9426}{<< /Names [<4974656D2E393734> \MusicPDFRef{12139} <4974656D2E393735> \MusicPDFRef{12140} <4974656D2E393736> \MusicPDFRef{12141} <4974656D2E393737> \MusicPDFRef{12142} <4974656D2E393738> \MusicPDFRef{12143} <4974656D2E393739> \MusicPDFRef{12144}] /Limits [<4974656D2E393734> <4974656D2E393739>] >>}
\MusicPDFWriteObject{9427}{<< /Names [<4974656D2E3938> \MusicPDFRef{12145} <4974656D2E393830> \MusicPDFRef{12146} <4974656D2E393831> \MusicPDFRef{12147} <4974656D2E393832> \MusicPDFRef{12148} <4974656D2E393833> \MusicPDFRef{12149} <4974656D2E393834> \MusicPDFRef{12150}] /Limits [<4974656D2E3938> <4974656D2E393834>] >>}
\MusicPDFWriteObject{9428}{<< /Names [<4974656D2E393835> \MusicPDFRef{12151} <4974656D2E393836> \MusicPDFRef{12152} <4974656D2E393837> \MusicPDFRef{12153} <4974656D2E393838> \MusicPDFRef{12154} <4974656D2E393839> \MusicPDFRef{12155} <4974656D2E3939> \MusicPDFRef{12156}] /Limits [<4974656D2E393835> <4974656D2E3939>] >>}
\MusicPDFWriteObject{9429}{<< /Names [<4974656D2E393930> \MusicPDFRef{12157} <4974656D2E393931> \MusicPDFRef{12158} <4974656D2E393932> \MusicPDFRef{12159} <4974656D2E393933> \MusicPDFRef{12160} <4974656D2E393934> \MusicPDFRef{12161} <4974656D2E393935> \MusicPDFRef{12162}] /Limits [<4974656D2E393930> <4974656D2E393935>] >>}
\MusicPDFWriteObject{9430}{<< /Names [<4974656D2E393936> \MusicPDFRef{12163} <4974656D2E393937> \MusicPDFRef{12164} <4974656D2E393938> \MusicPDFRef{12165} <4974656D2E393939> \MusicPDFRef{12166} <617070656E6469782A2E33> \MusicPDFRef{12167} <617070656E6469782E41> \MusicPDFRef{12168}] /Limits [<4974656D2E393936> <617070656E6469782E41>] >>}
\MusicPDFWriteObject{9431}{<< /Names [<617070656E6469782E42> \MusicPDFRef{12169} <617070656E6469782E43> \MusicPDFRef{12170} <617070656E6469782E44> \MusicPDFRef{12171} <617070656E6469782E45> \MusicPDFRef{12172} <636861707465722A2E32> \MusicPDFRef{12173} <636861707465722E31> \MusicPDFRef{12174}] /Limits [<617070656E6469782E42> <636861707465722E31>] >>}
\MusicPDFWriteObject{9432}{<< /Names [<636861707465722E3130> \MusicPDFRef{12175} <636861707465722E3131> \MusicPDFRef{12176} <636861707465722E3132> \MusicPDFRef{12177} <636861707465722E3133> \MusicPDFRef{12178} <636861707465722E3134> \MusicPDFRef{12179} <636861707465722E32> \MusicPDFRef{12180}] /Limits [<636861707465722E3130> <636861707465722E32>] >>}
\MusicPDFWriteObject{9433}{<< /Names [<636861707465722E33> \MusicPDFRef{12181} <636861707465722E34> \MusicPDFRef{12182} <636861707465722E35> \MusicPDFRef{12183} <636861707465722E36> \MusicPDFRef{12184} <636861707465722E37> \MusicPDFRef{12185} <636861707465722E38> \MusicPDFRef{12186}] /Limits [<636861707465722E33> <636861707465722E38>] >>}
\MusicPDFWriteObject{9434}{<< /Names [<636861707465722E39> \MusicPDFRef{12187} <636974652E41> \MusicPDFRef{12188} <636974652E4142> \MusicPDFRef{12189} <636974652E41424D> \MusicPDFRef{12190} <636974652E414252> \MusicPDFRef{12191} <636974652E4153> \MusicPDFRef{12192}] /Limits [<636861707465722E39> <636974652E4153>] >>}
\MusicPDFWriteObject{9435}{<< /Names [<636974652E4154> \MusicPDFRef{12193} <636974652E4163> \MusicPDFRef{12194} <636974652E41686C> \MusicPDFRef{12195} <636974652E416C2D42> \MusicPDFRef{12196} <636974652E4170696F6C61> \MusicPDFRef{12197} <636974652E417261> \MusicPDFRef{12198}] /Limits [<636974652E4154> <636974652E417261>] >>}
\MusicPDFWriteObject{9436}{<< /Names [<636974652E41726D> \MusicPDFRef{12199} <636974652E4172742D31> \MusicPDFRef{12200} <636974652E4172742D32> \MusicPDFRef{12201} <636974652E417274687572> \MusicPDFRef{12202} <636974652E4174436865> \MusicPDFRef{12203} <636974652E4175> \MusicPDFRef{12204}] /Limits [<636974652E41726D> <636974652E4175>] >>}
\MusicPDFWriteObject{9437}{<< /Names [<636974652E417562> \MusicPDFRef{12205} <636974652E4242> \MusicPDFRef{12206} <636974652E424352> \MusicPDFRef{12207} <636974652E42435A> \MusicPDFRef{12208} <636974652E4247> \MusicPDFRef{12209} <636974652E424756> \MusicPDFRef{12210}] /Limits [<636974652E417562> <636974652E424756>] >>}
\MusicPDFWriteObject{9438}{<< /Names [<636974652E424A> \MusicPDFRef{12211} <636974652E424A4C> \MusicPDFRef{12212} <636974652E424A4D> \MusicPDFRef{12213} <636974652E424B4D634B> \MusicPDFRef{12214} <636974652E424C> \MusicPDFRef{12215} <636974652E4254> \MusicPDFRef{12216}] /Limits [<636974652E424A> <636974652E4254>] >>}
\MusicPDFWriteObject{9439}{<< /Names [<636974652E4261> \MusicPDFRef{12217} <636974652E426131> \MusicPDFRef{12218} <636974652E426132> \MusicPDFRef{12219} <636974652E426172> \MusicPDFRef{12220} <636974652E4265> \MusicPDFRef{12221} <636974652E4269726B686F6666> \MusicPDFRef{12222}] /Limits [<636974652E4261> <636974652E4269726B686F6666>] >>}
\MusicPDFWriteObject{9440}{<< /Names [<636974652E426C> \MusicPDFRef{12223} <636974652E426F> \MusicPDFRef{12224} <636974652E426F67> \MusicPDFRef{12225} <636974652E426F6872> \MusicPDFRef{12226} <636974652E426F687231> \MusicPDFRef{12227} <636974652E426F687232> \MusicPDFRef{12228}] /Limits [<636974652E426C> <636974652E426F687232>] >>}
\MusicPDFWriteObject{9441}{<< /Names [<636974652E426F687233> \MusicPDFRef{12229} <636974652E426F687234> \MusicPDFRef{12230} <636974652E426F7231> \MusicPDFRef{12231} <636974652E426F7232> \MusicPDFRef{12232} <636974652E426F75> \MusicPDFRef{12233} <636974652E426F752D4754> \MusicPDFRef{12234}] /Limits [<636974652E426F687233> <636974652E426F752D4754>] >>}
\MusicPDFWriteObject{9442}{<< /Names [<636974652E426F752D49> \MusicPDFRef{12235} <636974652E426F752D4949> \MusicPDFRef{12236} <636974652E426F752D545653> \MusicPDFRef{12237} <636974652E427231> \MusicPDFRef{12238} <636974652E427232> \MusicPDFRef{12239} <636974652E427233> \MusicPDFRef{12240}] /Limits [<636974652E426F752D49> <636974652E427233>] >>}
\MusicPDFWriteObject{9443}{<< /Names [<636974652E427234> \MusicPDFRef{12241} <636974652E427269732D61> \MusicPDFRef{12242} <636974652E427269732D62> \MusicPDFRef{12243} <636974652E427275> \MusicPDFRef{12244} <636974652E4346> \MusicPDFRef{12245} <636974652E43465353> \MusicPDFRef{12246}] /Limits [<636974652E427234> <636974652E43465353>] >>}
\MusicPDFWriteObject{9444}{<< /Names [<636974652E434852> \MusicPDFRef{12247} <636974652E435053> \MusicPDFRef{12248} <636974652E436172> \MusicPDFRef{12249} <636974652E4361726C65736F6E> \MusicPDFRef{12250} <636974652E43645357> \MusicPDFRef{12251} <636974652E4368> \MusicPDFRef{12252}] /Limits [<636974652E434852> <636974652E4368>] >>}
\MusicPDFWriteObject{9445}{<< /Names [<636974652E436861> \MusicPDFRef{12253} <636974652E436C> \MusicPDFRef{12254} <636974652E436E772D61> \MusicPDFRef{12255} <636974652E436E772D62> \MusicPDFRef{12256} <636974652E436E772D63> \MusicPDFRef{12257} <636974652E436F> \MusicPDFRef{12258}] /Limits [<636974652E436861> <636974652E436F>] >>}
\MusicPDFWriteObject{9446}{<< /Names [<636974652E436F68> \MusicPDFRef{12259} <636974652E436F6E726164> \MusicPDFRef{12260} <636974652E4445> \MusicPDFRef{12261} <636974652E4453> \MusicPDFRef{12262} <636974652E4455> \MusicPDFRef{12263} <636974652E446176656E706F72742D31> \MusicPDFRef{12264}] /Limits [<636974652E436F68> <636974652E446176656E706F72742D31>] >>}
\MusicPDFWriteObject{9447}{<< /Names [<636974652E446176656E706F72742D32> \MusicPDFRef{12265} <636974652E4465> \MusicPDFRef{12266} <636974652E446532> \MusicPDFRef{12267} <636974652E446533> \MusicPDFRef{12268} <636974652E446534> \MusicPDFRef{12269} <636974652E44656C> \MusicPDFRef{12270}] /Limits [<636974652E446176656E706F72742D32> <636974652E44656C>] >>}
\MusicPDFWriteObject{9448}{<< /Names [<636974652E44656E> \MusicPDFRef{12271} <636974652E44694D> \MusicPDFRef{12272} <636974652E446978> \MusicPDFRef{12273} <636974652E4475> \MusicPDFRef{12274} <636974652E44756E666F7264> \MusicPDFRef{12275} <636974652E447572> \MusicPDFRef{12276}] /Limits [<636974652E44656E> <636974652E447572>] >>}
\MusicPDFWriteObject{9449}{<< /Names [<636974652E44794D> \MusicPDFRef{12277} <636974652E45> \MusicPDFRef{12278} <636974652E45474149> \MusicPDFRef{12279} <636974652E4547414949> \MusicPDFRef{12280} <636974652E454741494949245F3124> \MusicPDFRef{12281} <636974652E454741494949245F3224> \MusicPDFRef{12282}] /Limits [<636974652E44794D> <636974652E454741494949245F3224>] >>}
\MusicPDFWriteObject{9450}{<< /Names [<636974652E4547414956245F3124> \MusicPDFRef{12283} <636974652E4547414956245F3224> \MusicPDFRef{12284} <636974652E4547414956245F3324> \MusicPDFRef{12285} <636974652E4547414956245F3424> \MusicPDFRef{12286} <636974652E454B> \MusicPDFRef{12287} <636974652E454C53> \MusicPDFRef{12288}] /Limits [<636974652E4547414956245F3124> <636974652E454C53>] >>}
\MusicPDFWriteObject{9451}{<< /Names [<636974652E4555> \MusicPDFRef{12289} <636974652E45687231> \MusicPDFRef{12290} <636974652E45687232> \MusicPDFRef{12291} <636974652E4647> \MusicPDFRef{12292} <636974652E464754> \MusicPDFRef{12293} <636974652E464B> \MusicPDFRef{12294}] /Limits [<636974652E4555> <636974652E464B>] >>}
\MusicPDFWriteObject{9452}{<< /Names [<636974652E464D> \MusicPDFRef{12295} <636974652E465374> \MusicPDFRef{12296} <636974652E466C4D> \MusicPDFRef{12297} <636974652E466C6F> \MusicPDFRef{12298} <636974652E466F6C> \MusicPDFRef{12299} <636974652E466F6C2D32> \MusicPDFRef{12300}] /Limits [<636974652E464D> <636974652E466F6C2D32>] >>}
\MusicPDFWriteObject{9453}{<< /Names [<636974652E466F72> \MusicPDFRef{12301} <636974652E4672> \MusicPDFRef{12302} <636974652E46756A4B61> \MusicPDFRef{12303} <636974652E4745> \MusicPDFRef{12304} <636974652E474C> \MusicPDFRef{12305} <636974652E474E6172> \MusicPDFRef{12306}] /Limits [<636974652E466F72> <636974652E474E6172>] >>}
\MusicPDFWriteObject{9454}{<< /Names [<636974652E4752> \MusicPDFRef{12307} <636974652E475261> \MusicPDFRef{12308} <636974652E4753> \MusicPDFRef{12309} <636974652E4761> \MusicPDFRef{12310} <636974652E47612D31> \MusicPDFRef{12311} <636974652E47612D32> \MusicPDFRef{12312}] /Limits [<636974652E4752> <636974652E47612D32>] >>}
\MusicPDFWriteObject{9455}{<< /Names [<636974652E47616F> \MusicPDFRef{12313} <636974652E47646D> \MusicPDFRef{12314} <636974652E476C65> \MusicPDFRef{12315} <636974652E476F4B7265> \MusicPDFRef{12316} <636974652E476F7373> \MusicPDFRef{12317} <636974652E47724D> \MusicPDFRef{12318}] /Limits [<636974652E47616F> <636974652E47724D>] >>}
\MusicPDFWriteObject{9456}{<< /Names [<636974652E47726F2D61> \MusicPDFRef{12319} <636974652E47726F2D62> \MusicPDFRef{12320} <636974652E47726F2D63> \MusicPDFRef{12321} <636974652E47726F2D64> \MusicPDFRef{12322} <636974652E47726F2D65> \MusicPDFRef{12323} <636974652E47726F2D66> \MusicPDFRef{12324}] /Limits [<636974652E47726F2D61> <636974652E47726F2D66>] >>}
\MusicPDFWriteObject{9457}{<< /Names [<636974652E48> \MusicPDFRef{12325} <636974652E484B> \MusicPDFRef{12326} <636974652E484B5A> \MusicPDFRef{12327} <636974652E484C> \MusicPDFRef{12328} <636974652E484E> \MusicPDFRef{12329} <636974652E4850> \MusicPDFRef{12330}] /Limits [<636974652E48> <636974652E4850>] >>}
\MusicPDFWriteObject{9458}{<< /Names [<636974652E48522D31> \MusicPDFRef{12331} <636974652E48522D32> \MusicPDFRef{12332} <636974652E4853> \MusicPDFRef{12333} <636974652E485353> \MusicPDFRef{12334} <636974652E485374> \MusicPDFRef{12335} <636974652E4857> \MusicPDFRef{12336}] /Limits [<636974652E48522D31> <636974652E4857>] >>}
\MusicPDFWriteObject{9459}{<< /Names [<636974652E4861> \MusicPDFRef{12337} <636974652E486967> \MusicPDFRef{12338} <636974652E486F72> \MusicPDFRef{12339} <636974652E486F726D> \MusicPDFRef{12340} <636974652E487531> \MusicPDFRef{12341} <636974652E487532> \MusicPDFRef{12342}] /Limits [<636974652E4861> <636974652E487532>] >>}
\MusicPDFWriteObject{9460}{<< /Names [<636974652E48755765> \MusicPDFRef{12343} <636974652E48756E74> \MusicPDFRef{12344} <636974652E494B> \MusicPDFRef{12345} <636974652E496E676831> \MusicPDFRef{12346} <636974652E496E676832> \MusicPDFRef{12347} <636974652E496E676833> \MusicPDFRef{12348}] /Limits [<636974652E48755765> <636974652E496E676833>] >>}
\MusicPDFWriteObject{9461}{<< /Names [<636974652E4977> \MusicPDFRef{12349} <636974652E4977616E696563> \MusicPDFRef{12350} <636974652E4A4C> \MusicPDFRef{12351} <636974652E4A5050> \MusicPDFRef{12352} <636974652E4A6172> \MusicPDFRef{12353} <636974652E4A656368> \MusicPDFRef{12354}] /Limits [<636974652E4977> <636974652E4A656368>] >>}
\MusicPDFWriteObject{9462}{<< /Names [<636974652E4A6F722D61> \MusicPDFRef{12355} <636974652E4A6F722D62> \MusicPDFRef{12356} <636974652E4B47> \MusicPDFRef{12357} <636974652E4B53> \MusicPDFRef{12358} <636974652E4B61> \MusicPDFRef{12359} <636974652E4B616B7574616E69> \MusicPDFRef{12360}] /Limits [<636974652E4A6F722D61> <636974652E4B616B7574616E69>] >>}
\MusicPDFWriteObject{9463}{<< /Names [<636974652E4B616C> \MusicPDFRef{12361} <636974652E4B61747A6E> \MusicPDFRef{12362} <636974652E4B65> \MusicPDFRef{12363} <636974652E4B652D61> \MusicPDFRef{12364} <636974652E4B652D62> \MusicPDFRef{12365} <636974652E4B656368> \MusicPDFRef{12366}] /Limits [<636974652E4B616C> <636974652E4B656368>] >>}
\MusicPDFWriteObject{9464}{<< /Names [<636974652E4B656C> \MusicPDFRef{12367} <636974652E4B6C654D6F> \MusicPDFRef{12368} <636974652E4B6C6565> \MusicPDFRef{12369} <636974652E4B6E> \MusicPDFRef{12370} <636974652E4B6F62> \MusicPDFRef{12371} <636974652E4B6F6C> \MusicPDFRef{12372}] /Limits [<636974652E4B656C> <636974652E4B6F6C>] >>}
\MusicPDFWriteObject{9465}{<< /Names [<636974652E4B6F7468652D49> \MusicPDFRef{12373} <636974652E4B6F7468652D4949> \MusicPDFRef{12374} <636974652E4B72757365> \MusicPDFRef{12375} <636974652E4B7562> \MusicPDFRef{12376} <636974652E4B75626F7461> \MusicPDFRef{12377} <636974652E4B75646C61> \MusicPDFRef{12378}] /Limits [<636974652E4B6F7468652D49> <636974652E4B75646C61>] >>}
\MusicPDFWriteObject{9466}{<< /Names [<636974652E4B756E656E> \MusicPDFRef{12379} <636974652E4C474C> \MusicPDFRef{12380} <636974652E4C4D> \MusicPDFRef{12381} <636974652E4C4D53> \MusicPDFRef{12382} <636974652E4C54> \MusicPDFRef{12383} <636974652E4C61> \MusicPDFRef{12384}] /Limits [<636974652E4B756E656E> <636974652E4C61>] >>}
\MusicPDFWriteObject{9467}{<< /Names [<636974652E4C616E672D61> \MusicPDFRef{12385} <636974652E4C616E672D62> \MusicPDFRef{12386} <636974652E4C616E672D63> \MusicPDFRef{12387} <636974652E4C616E672D64> \MusicPDFRef{12388} <636974652E4C616E672D65> \MusicPDFRef{12389} <636974652E4C616E672D66> \MusicPDFRef{12390}] /Limits [<636974652E4C616E672D61> <636974652E4C616E672D66>] >>}
\MusicPDFWriteObject{9468}{<< /Names [<636974652E4C616E676C616E6473> \MusicPDFRef{12391} <636974652E4C6565> \MusicPDFRef{12392} <636974652E4C6573> \MusicPDFRef{12393} <636974652E4C69742D49> \MusicPDFRef{12394} <636974652E4C69742D4949> \MusicPDFRef{12395} <636974652E4D> \MusicPDFRef{12396}] /Limits [<636974652E4C616E676C616E6473> <636974652E4D>] >>}
\MusicPDFWriteObject{9469}{<< /Names [<636974652E4D4C6569> \MusicPDFRef{12397} <636974652E4D53> \MusicPDFRef{12398} <636974652E4D56> \MusicPDFRef{12399} <636974652E4D562D32> \MusicPDFRef{12400} <636974652E4D5661> \MusicPDFRef{12401} <636974652E4D5665> \MusicPDFRef{12402}] /Limits [<636974652E4D4C6569> <636974652E4D5665>] >>}
\MusicPDFWriteObject{9470}{<< /Names [<636974652E4D5A> \MusicPDFRef{12403} <636974652E4D61634C> \MusicPDFRef{12404} <636974652E4D61634C32> \MusicPDFRef{12405} <636974652E4D6172> \MusicPDFRef{12406} <636974652E4D61742D49> \MusicPDFRef{12407} <636974652E4D61742D4949> \MusicPDFRef{12408}] /Limits [<636974652E4D5A> <636974652E4D61742D4949>] >>}
\MusicPDFWriteObject{9471}{<< /Names [<636974652E4D6179> \MusicPDFRef{12409} <636974652E4D634172> \MusicPDFRef{12410} <636974652E4D6343> \MusicPDFRef{12411} <636974652E4D65> \MusicPDFRef{12412} <636974652E4D6531> \MusicPDFRef{12413} <636974652E4D6532> \MusicPDFRef{12414}] /Limits [<636974652E4D6179> <636974652E4D6532>] >>}
\MusicPDFWriteObject{9472}{<< /Names [<636974652E4D6533> \MusicPDFRef{12415} <636974652E4D6567> \MusicPDFRef{12416} <636974652E4D69> \MusicPDFRef{12417} <636974652E4D6974> \MusicPDFRef{12418} <636974652E4D6F726F7A> \MusicPDFRef{12419} <636974652E4D6F73> \MusicPDFRef{12420}] /Limits [<636974652E4D6533> <636974652E4D6F73>] >>}
\MusicPDFWriteObject{9473}{<< /Names [<636974652E4D75> \MusicPDFRef{12421} <636974652E4D75727479> \MusicPDFRef{12422} <636974652E4E> \MusicPDFRef{12423} <636974652E4E42> \MusicPDFRef{12424} <636974652E4E6575> \MusicPDFRef{12425} <636974652E4E67> \MusicPDFRef{12426}] /Limits [<636974652E4D75> <636974652E4E67>] >>}
\MusicPDFWriteObject{9474}{<< /Names [<636974652E4E696B> \MusicPDFRef{12427} <636974652E4F6573> \MusicPDFRef{12428} <636974652E4F6D6172> \MusicPDFRef{12429} <636974652E5042> \MusicPDFRef{12430} <636974652E5061> \MusicPDFRef{12431} <636974652E50696E736B79> \MusicPDFRef{12432}] /Limits [<636974652E4E696B> <636974652E50696E736B79>] >>}
\MusicPDFWriteObject{9475}{<< /Names [<636974652E5151> \MusicPDFRef{12433} <636974652E51756572> \MusicPDFRef{12434} <636974652E5253> \MusicPDFRef{12435} <636974652E5256> \MusicPDFRef{12436} <636974652E526574> \MusicPDFRef{12437} <636974652E52657A6B> \MusicPDFRef{12438}] /Limits [<636974652E5151> <636974652E52657A6B>] >>}
\MusicPDFWriteObject{9476}{<< /Names [<636974652E5269656D616E6E> \MusicPDFRef{12439} <636974652E526F> \MusicPDFRef{12440} <636974652E5275> \MusicPDFRef{12441} <636974652E52752D32> \MusicPDFRef{12442} <636974652E52752D33> \MusicPDFRef{12443} <636974652E52752D5361> \MusicPDFRef{12444}] /Limits [<636974652E5269656D616E6E> <636974652E52752D5361>] >>}
\MusicPDFWriteObject{9477}{<< /Names [<636974652E5275642D5361> \MusicPDFRef{12445} <636974652E5279616E> \MusicPDFRef{12446} <636974652E53474131> \MusicPDFRef{12447} <636974652E53474133> \MusicPDFRef{12448} <636974652E53474134> \MusicPDFRef{12449} <636974652E534E> \MusicPDFRef{12450}] /Limits [<636974652E5275642D5361> <636974652E534E>] >>}
\MusicPDFWriteObject{9478}{<< /Names [<636974652E534F> \MusicPDFRef{12451} <636974652E5350> \MusicPDFRef{12452} <636974652E5357> \MusicPDFRef{12453} <636974652E5361726E616B> \MusicPDFRef{12454} <636974652E5363682D61> \MusicPDFRef{12455} <636974652E5363682D62> \MusicPDFRef{12456}] /Limits [<636974652E534F> <636974652E5363682D62>] >>}
\MusicPDFWriteObject{9479}{<< /Names [<636974652E53636861> \MusicPDFRef{12457} <636974652E5363686D> \MusicPDFRef{12458} <636974652E5365> \MusicPDFRef{12459} <636974652E536572696573> \MusicPDFRef{12460} <636974652E536869> \MusicPDFRef{12461} <636974652E53696C7665726D616E> \MusicPDFRef{12462}] /Limits [<636974652E53636861> <636974652E53696C7665726D616E>] >>}
\MusicPDFWriteObject{9480}{<< /Names [<636974652E53696D6F6E> \MusicPDFRef{12463} <636974652E536D> \MusicPDFRef{12464} <636974652E53745365> \MusicPDFRef{12465} <636974652E53745368612D494949> \MusicPDFRef{12466} <636974652E5374537472> \MusicPDFRef{12467} <636974652E537465> \MusicPDFRef{12468}] /Limits [<636974652E53696D6F6E> <636974652E537465>] >>}
\MusicPDFWriteObject{9481}{<< /Names [<636974652E537472> \MusicPDFRef{12469} <636974652E53747275626C65> \MusicPDFRef{12470} <636974652E5431> \MusicPDFRef{12471} <636974652E5432> \MusicPDFRef{12472} <636974652E544B594B> \MusicPDFRef{12473} <636974652E54617465> \MusicPDFRef{12474}] /Limits [<636974652E537472> <636974652E54617465>] >>}
\MusicPDFWriteObject{9482}{<< /Names [<636974652E547265> \MusicPDFRef{12475} <636974652E547356> \MusicPDFRef{12476} <636974652E564B> \MusicPDFRef{12477} <636974652E564C48> \MusicPDFRef{12478} <636974652E56616C2D61> \MusicPDFRef{12479} <636974652E56616C2D62> \MusicPDFRef{12480}] /Limits [<636974652E547265> <636974652E56616C2D62>] >>}
\MusicPDFWriteObject{9483}{<< /Names [<636974652E566F6774> \MusicPDFRef{12481} <636974652E566F67742D31> \MusicPDFRef{12482} <636974652E566F67742D32> \MusicPDFRef{12483} <636974652E566F696774> \MusicPDFRef{12484} <636974652E566F72> \MusicPDFRef{12485} <636974652E57> \MusicPDFRef{12486}] /Limits [<636974652E566F6774> <636974652E57>] >>}
\MusicPDFWriteObject{9484}{<< /Names [<636974652E5761> \MusicPDFRef{12487} <636974652E5765696C2D61> \MusicPDFRef{12488} <636974652E5765696C2D62> \MusicPDFRef{12489} <636974652E5765696C2D63> \MusicPDFRef{12490} <636974652E5765696C2D64> \MusicPDFRef{12491} <636974652E5765696C2D65> \MusicPDFRef{12492}] /Limits [<636974652E5761> <636974652E5765696C2D65>] >>}
\MusicPDFWriteObject{9485}{<< /Names [<636974652E5765696C2D66> \MusicPDFRef{12493} <636974652E5769> \MusicPDFRef{12494} <636974652E57696C> \MusicPDFRef{12495} <636974652E576E67722D31> \MusicPDFRef{12496} <636974652E576E67722D32> \MusicPDFRef{12497} <636974652E57726E> \MusicPDFRef{12498}] /Limits [<636974652E5765696C2D66> <636974652E57726E>] >>}
\MusicPDFWriteObject{9486}{<< /Names [<636974652E596F73> \MusicPDFRef{12499} <636974652E596F732D32> \MusicPDFRef{12500} <636974652E5A61> \MusicPDFRef{12501} <636974652E5A6167696572> \MusicPDFRef{12502} <636974652E6E4C6162> \MusicPDFRef{12503} <636974652E764F> \MusicPDFRef{12504}] /Limits [<636974652E596F73> <636974652E764F>] >>}
\MusicPDFWriteObject{9487}{<< /Names [<6571756174696F6E2E312E312E322E31> \MusicPDFRef{12505} <6571756174696F6E2E312E322E312E31> \MusicPDFRef{12506} <6571756174696F6E2E312E332E312E31> \MusicPDFRef{12507} <6571756174696F6E2E312E332E312E3130> \MusicPDFRef{12508} <6571756174696F6E2E312E332E312E3131> \MusicPDFRef{12509} <6571756174696F6E2E312E332E312E3132> \MusicPDFRef{12510}] /Limits [<6571756174696F6E2E312E312E322E31> <6571756174696F6E2E312E332E312E3132>] >>}
\MusicPDFWriteObject{9488}{<< /Names [<6571756174696F6E2E312E332E312E32> \MusicPDFRef{12511} <6571756174696F6E2E312E332E312E33> \MusicPDFRef{12512} <6571756174696F6E2E312E332E312E34> \MusicPDFRef{12513} <6571756174696F6E2E312E332E312E35> \MusicPDFRef{12514} <6571756174696F6E2E312E332E312E36> \MusicPDFRef{12515} <6571756174696F6E2E312E332E312E37> \MusicPDFRef{12516}] /Limits [<6571756174696F6E2E312E332E312E32> <6571756174696F6E2E312E332E312E37>] >>}
\MusicPDFWriteObject{9489}{<< /Names [<6571756174696F6E2E312E332E312E38> \MusicPDFRef{12517} <6571756174696F6E2E312E332E312E39> \MusicPDFRef{12518} <6571756174696F6E2E312E332E322E31> \MusicPDFRef{12519} <6571756174696F6E2E312E332E322E32> \MusicPDFRef{12520} <6571756174696F6E2E312E372E312E31> \MusicPDFRef{12521} <6571756174696F6E2E312E372E312E32> \MusicPDFRef{12522}] /Limits [<6571756174696F6E2E312E332E312E38> <6571756174696F6E2E312E372E312E32>] >>}
\MusicPDFWriteObject{9490}{<< /Names [<6571756174696F6E2E312E372E322E31> \MusicPDFRef{12523} <6571756174696F6E2E312E372E322E32> \MusicPDFRef{12524} <6571756174696F6E2E312E372E332E31> \MusicPDFRef{12525} <6571756174696F6E2E312E372E332E32> \MusicPDFRef{12526} <6571756174696F6E2E312E382E312E31> \MusicPDFRef{12527} <6571756174696F6E2E312E382E312E32> \MusicPDFRef{12528}] /Limits [<6571756174696F6E2E312E372E322E31> <6571756174696F6E2E312E382E312E32>] >>}
\MusicPDFWriteObject{9491}{<< /Names [<6571756174696F6E2E312E382E312E33> \MusicPDFRef{12529} <6571756174696F6E2E312E382E312E34> \MusicPDFRef{12530} <6571756174696F6E2E312E382E322E31> \MusicPDFRef{12531} <6571756174696F6E2E312E382E322E32> \MusicPDFRef{12532} <6571756174696F6E2E312E382E322E33> \MusicPDFRef{12533} <6571756174696F6E2E312E382E332E31> \MusicPDFRef{12534}] /Limits [<6571756174696F6E2E312E382E312E33> <6571756174696F6E2E312E382E332E31>] >>}
\MusicPDFWriteObject{9492}{<< /Names [<6571756174696F6E2E312E382E332E32> \MusicPDFRef{12535} <6571756174696F6E2E312E382E332E33> \MusicPDFRef{12536} <6571756174696F6E2E312E382E332E34> \MusicPDFRef{12537} <6571756174696F6E2E312E382E332E35> \MusicPDFRef{12538} <6571756174696F6E2E312E382E332E36> \MusicPDFRef{12539} <6571756174696F6E2E312E382E332E37> \MusicPDFRef{12540}] /Limits [<6571756174696F6E2E312E382E332E32> <6571756174696F6E2E312E382E332E37>] >>}
\MusicPDFWriteObject{9493}{<< /Names [<6571756174696F6E2E312E382E332E38> \MusicPDFRef{12541} <6571756174696F6E2E312E382E332E39> \MusicPDFRef{12542} <6571756174696F6E2E312E382E352E31> \MusicPDFRef{12543} <6571756174696F6E2E31302E312E312E31> \MusicPDFRef{12544} <6571756174696F6E2E31302E312E312E3130> \MusicPDFRef{12545} <6571756174696F6E2E31302E312E312E3131> \MusicPDFRef{12546}] /Limits [<6571756174696F6E2E312E382E332E38> <6571756174696F6E2E31302E312E312E3131>] >>}
\MusicPDFWriteObject{9494}{<< /Names [<6571756174696F6E2E31302E312E312E3132> \MusicPDFRef{12547} <6571756174696F6E2E31302E312E312E3133> \MusicPDFRef{12548} <6571756174696F6E2E31302E312E312E3134> \MusicPDFRef{12549} <6571756174696F6E2E31302E312E312E3135> \MusicPDFRef{12550} <6571756174696F6E2E31302E312E312E32> \MusicPDFRef{12551} <6571756174696F6E2E31302E312E312E33> \MusicPDFRef{12552}] /Limits [<6571756174696F6E2E31302E312E312E3132> <6571756174696F6E2E31302E312E312E33>] >>}
\MusicPDFWriteObject{9495}{<< /Names [<6571756174696F6E2E31302E312E312E34> \MusicPDFRef{12553} <6571756174696F6E2E31302E312E312E35> \MusicPDFRef{12554} <6571756174696F6E2E31302E312E312E36> \MusicPDFRef{12555} <6571756174696F6E2E31302E312E312E37> \MusicPDFRef{12556} <6571756174696F6E2E31302E312E312E38> \MusicPDFRef{12557} <6571756174696F6E2E31302E312E312E39> \MusicPDFRef{12558}] /Limits [<6571756174696F6E2E31302E312E312E34> <6571756174696F6E2E31302E312E312E39>] >>}
\MusicPDFWriteObject{9496}{<< /Names [<6571756174696F6E2E31302E312E322E31> \MusicPDFRef{12559} <6571756174696F6E2E31302E312E322E32> \MusicPDFRef{12560} <6571756174696F6E2E31302E312E322E33> \MusicPDFRef{12561} <6571756174696F6E2E31302E312E322E34> \MusicPDFRef{12562} <6571756174696F6E2E31302E312E322E35> \MusicPDFRef{12563} <6571756174696F6E2E31302E312E322E36> \MusicPDFRef{12564}] /Limits [<6571756174696F6E2E31302E312E322E31> <6571756174696F6E2E31302E312E322E36>] >>}
\MusicPDFWriteObject{9497}{<< /Names [<6571756174696F6E2E31302E322E312E31> \MusicPDFRef{12565} <6571756174696F6E2E31302E322E312E3130> \MusicPDFRef{12566} <6571756174696F6E2E31302E322E312E3131> \MusicPDFRef{12567} <6571756174696F6E2E31302E322E312E3132> \MusicPDFRef{12568} <6571756174696F6E2E31302E322E312E3133> \MusicPDFRef{12569} <6571756174696F6E2E31302E322E312E32> \MusicPDFRef{12570}] /Limits [<6571756174696F6E2E31302E322E312E31> <6571756174696F6E2E31302E322E312E32>] >>}
\MusicPDFWriteObject{9498}{<< /Names [<6571756174696F6E2E31302E322E312E33> \MusicPDFRef{12571} <6571756174696F6E2E31302E322E312E34> \MusicPDFRef{12572} <6571756174696F6E2E31302E322E312E35> \MusicPDFRef{12573} <6571756174696F6E2E31302E322E312E36> \MusicPDFRef{12574} <6571756174696F6E2E31302E322E312E37> \MusicPDFRef{12575} <6571756174696F6E2E31302E322E312E38> \MusicPDFRef{12576}] /Limits [<6571756174696F6E2E31302E322E312E33> <6571756174696F6E2E31302E322E312E38>] >>}
\MusicPDFWriteObject{9499}{<< /Names [<6571756174696F6E2E31302E322E312E39> \MusicPDFRef{12577} <6571756174696F6E2E31302E322E322E31> \MusicPDFRef{12578} <6571756174696F6E2E31302E322E322E32> \MusicPDFRef{12579} <6571756174696F6E2E31302E322E322E33> \MusicPDFRef{12580} <6571756174696F6E2E31302E322E322E34> \MusicPDFRef{12581} <6571756174696F6E2E31302E322E322E35> \MusicPDFRef{12582}] /Limits [<6571756174696F6E2E31302E322E312E39> <6571756174696F6E2E31302E322E322E35>] >>}
\MusicPDFWriteObject{9500}{<< /Names [<6571756174696F6E2E31302E322E322E36> \MusicPDFRef{12583} <6571756174696F6E2E31302E322E322E37> \MusicPDFRef{12584} <6571756174696F6E2E31302E322E322E38> \MusicPDFRef{12585} <6571756174696F6E2E31302E322E322E39> \MusicPDFRef{12586} <6571756174696F6E2E31302E332E312E31> \MusicPDFRef{12587} <6571756174696F6E2E31302E332E312E3130> \MusicPDFRef{12588}] /Limits [<6571756174696F6E2E31302E322E322E36> <6571756174696F6E2E31302E332E312E3130>] >>}
\MusicPDFWriteObject{9501}{<< /Names [<6571756174696F6E2E31302E332E312E3131> \MusicPDFRef{12589} <6571756174696F6E2E31302E332E312E3132> \MusicPDFRef{12590} <6571756174696F6E2E31302E332E312E3133> \MusicPDFRef{12591} <6571756174696F6E2E31302E332E312E3134> \MusicPDFRef{12592} <6571756174696F6E2E31302E332E312E3135> \MusicPDFRef{12593} <6571756174696F6E2E31302E332E312E3136> \MusicPDFRef{12594}] /Limits [<6571756174696F6E2E31302E332E312E3131> <6571756174696F6E2E31302E332E312E3136>] >>}
\MusicPDFWriteObject{9502}{<< /Names [<6571756174696F6E2E31302E332E312E3137> \MusicPDFRef{12595} <6571756174696F6E2E31302E332E312E3138> \MusicPDFRef{12596} <6571756174696F6E2E31302E332E312E3139> \MusicPDFRef{12597} <6571756174696F6E2E31302E332E312E32> \MusicPDFRef{12598} <6571756174696F6E2E31302E332E312E3230> \MusicPDFRef{12599} <6571756174696F6E2E31302E332E312E3231> \MusicPDFRef{12600}] /Limits [<6571756174696F6E2E31302E332E312E3137> <6571756174696F6E2E31302E332E312E3231>] >>}
\MusicPDFWriteObject{9503}{<< /Names [<6571756174696F6E2E31302E332E312E3232> \MusicPDFRef{12601} <6571756174696F6E2E31302E332E312E3233> \MusicPDFRef{12602} <6571756174696F6E2E31302E332E312E3234> \MusicPDFRef{12603} <6571756174696F6E2E31302E332E312E3235> \MusicPDFRef{12604} <6571756174696F6E2E31302E332E312E3236> \MusicPDFRef{12605} <6571756174696F6E2E31302E332E312E33> \MusicPDFRef{12606}] /Limits [<6571756174696F6E2E31302E332E312E3232> <6571756174696F6E2E31302E332E312E33>] >>}
\MusicPDFWriteObject{9504}{<< /Names [<6571756174696F6E2E31302E332E312E34> \MusicPDFRef{12607} <6571756174696F6E2E31302E332E312E35> \MusicPDFRef{12608} <6571756174696F6E2E31302E332E312E36> \MusicPDFRef{12609} <6571756174696F6E2E31302E332E312E37> \MusicPDFRef{12610} <6571756174696F6E2E31302E332E312E38> \MusicPDFRef{12611} <6571756174696F6E2E31302E332E312E39> \MusicPDFRef{12612}] /Limits [<6571756174696F6E2E31302E332E312E34> <6571756174696F6E2E31302E332E312E39>] >>}
\MusicPDFWriteObject{9505}{<< /Names [<6571756174696F6E2E31302E332E322E31> \MusicPDFRef{12613} <6571756174696F6E2E31302E332E322E3130> \MusicPDFRef{12614} <6571756174696F6E2E31302E332E322E3131> \MusicPDFRef{12615} <6571756174696F6E2E31302E332E322E32> \MusicPDFRef{12616} <6571756174696F6E2E31302E332E322E33> \MusicPDFRef{12617} <6571756174696F6E2E31302E332E322E34> \MusicPDFRef{12618}] /Limits [<6571756174696F6E2E31302E332E322E31> <6571756174696F6E2E31302E332E322E34>] >>}
\MusicPDFWriteObject{9506}{<< /Names [<6571756174696F6E2E31302E332E322E35> \MusicPDFRef{12619} <6571756174696F6E2E31302E332E322E36> \MusicPDFRef{12620} <6571756174696F6E2E31302E332E322E37> \MusicPDFRef{12621} <6571756174696F6E2E31302E332E322E38> \MusicPDFRef{12622} <6571756174696F6E2E31302E332E322E39> \MusicPDFRef{12623} <6571756174696F6E2E31302E342E312E31> \MusicPDFRef{12624}] /Limits [<6571756174696F6E2E31302E332E322E35> <6571756174696F6E2E31302E342E312E31>] >>}
\MusicPDFWriteObject{9507}{<< /Names [<6571756174696F6E2E31302E342E312E3130> \MusicPDFRef{12625} <6571756174696F6E2E31302E342E312E3131> \MusicPDFRef{12626} <6571756174696F6E2E31302E342E312E3132> \MusicPDFRef{12627} <6571756174696F6E2E31302E342E312E3133> \MusicPDFRef{12628} <6571756174696F6E2E31302E342E312E32> \MusicPDFRef{12629} <6571756174696F6E2E31302E342E312E33> \MusicPDFRef{12630}] /Limits [<6571756174696F6E2E31302E342E312E3130> <6571756174696F6E2E31302E342E312E33>] >>}
\MusicPDFWriteObject{9508}{<< /Names [<6571756174696F6E2E31302E342E312E34> \MusicPDFRef{12631} <6571756174696F6E2E31302E342E312E35> \MusicPDFRef{12632} <6571756174696F6E2E31302E342E312E36> \MusicPDFRef{12633} <6571756174696F6E2E31302E342E312E37> \MusicPDFRef{12634} <6571756174696F6E2E31302E342E312E38> \MusicPDFRef{12635} <6571756174696F6E2E31302E342E312E39> \MusicPDFRef{12636}] /Limits [<6571756174696F6E2E31302E342E312E34> <6571756174696F6E2E31302E342E312E39>] >>}
\MusicPDFWriteObject{9509}{<< /Names [<6571756174696F6E2E31302E342E322E31> \MusicPDFRef{12637} <6571756174696F6E2E31302E342E322E3130> \MusicPDFRef{12638} <6571756174696F6E2E31302E342E322E3131> \MusicPDFRef{12639} <6571756174696F6E2E31302E342E322E3132> \MusicPDFRef{12640} <6571756174696F6E2E31302E342E322E3133> \MusicPDFRef{12641} <6571756174696F6E2E31302E342E322E3134> \MusicPDFRef{12642}] /Limits [<6571756174696F6E2E31302E342E322E31> <6571756174696F6E2E31302E342E322E3134>] >>}
\MusicPDFWriteObject{9510}{<< /Names [<6571756174696F6E2E31302E342E322E3135> \MusicPDFRef{12643} <6571756174696F6E2E31302E342E322E3136> \MusicPDFRef{12644} <6571756174696F6E2E31302E342E322E3137> \MusicPDFRef{12645} <6571756174696F6E2E31302E342E322E3138> \MusicPDFRef{12646} <6571756174696F6E2E31302E342E322E3139> \MusicPDFRef{12647} <6571756174696F6E2E31302E342E322E32> \MusicPDFRef{12648}] /Limits [<6571756174696F6E2E31302E342E322E3135> <6571756174696F6E2E31302E342E322E32>] >>}
\MusicPDFWriteObject{9511}{<< /Names [<6571756174696F6E2E31302E342E322E3230> \MusicPDFRef{12649} <6571756174696F6E2E31302E342E322E3231> \MusicPDFRef{12650} <6571756174696F6E2E31302E342E322E3232> \MusicPDFRef{12651} <6571756174696F6E2E31302E342E322E3233> \MusicPDFRef{12652} <6571756174696F6E2E31302E342E322E3234> \MusicPDFRef{12653} <6571756174696F6E2E31302E342E322E3235> \MusicPDFRef{12654}] /Limits [<6571756174696F6E2E31302E342E322E3230> <6571756174696F6E2E31302E342E322E3235>] >>}
\MusicPDFWriteObject{9512}{<< /Names [<6571756174696F6E2E31302E342E322E3236> \MusicPDFRef{12655} <6571756174696F6E2E31302E342E322E3237> \MusicPDFRef{12656} <6571756174696F6E2E31302E342E322E3238> \MusicPDFRef{12657} <6571756174696F6E2E31302E342E322E3239> \MusicPDFRef{12658} <6571756174696F6E2E31302E342E322E33> \MusicPDFRef{12659} <6571756174696F6E2E31302E342E322E3330> \MusicPDFRef{12660}] /Limits [<6571756174696F6E2E31302E342E322E3236> <6571756174696F6E2E31302E342E322E3330>] >>}
\MusicPDFWriteObject{9513}{<< /Names [<6571756174696F6E2E31302E342E322E3331> \MusicPDFRef{12661} <6571756174696F6E2E31302E342E322E3332> \MusicPDFRef{12662} <6571756174696F6E2E31302E342E322E3333> \MusicPDFRef{12663} <6571756174696F6E2E31302E342E322E3334> \MusicPDFRef{12664} <6571756174696F6E2E31302E342E322E34> \MusicPDFRef{12665} <6571756174696F6E2E31302E342E322E35> \MusicPDFRef{12666}] /Limits [<6571756174696F6E2E31302E342E322E3331> <6571756174696F6E2E31302E342E322E35>] >>}
\MusicPDFWriteObject{9514}{<< /Names [<6571756174696F6E2E31302E342E322E36> \MusicPDFRef{12667} <6571756174696F6E2E31302E342E322E37> \MusicPDFRef{12668} <6571756174696F6E2E31302E342E322E38> \MusicPDFRef{12669} <6571756174696F6E2E31302E342E322E39> \MusicPDFRef{12670} <6571756174696F6E2E31302E342E332E31> \MusicPDFRef{12671} <6571756174696F6E2E31302E342E332E3130> \MusicPDFRef{12672}] /Limits [<6571756174696F6E2E31302E342E322E36> <6571756174696F6E2E31302E342E332E3130>] >>}
\MusicPDFWriteObject{9515}{<< /Names [<6571756174696F6E2E31302E342E332E3131> \MusicPDFRef{12673} <6571756174696F6E2E31302E342E332E3132> \MusicPDFRef{12674} <6571756174696F6E2E31302E342E332E3133> \MusicPDFRef{12675} <6571756174696F6E2E31302E342E332E3134> \MusicPDFRef{12676} <6571756174696F6E2E31302E342E332E3135> \MusicPDFRef{12677} <6571756174696F6E2E31302E342E332E3136> \MusicPDFRef{12678}] /Limits [<6571756174696F6E2E31302E342E332E3131> <6571756174696F6E2E31302E342E332E3136>] >>}
\MusicPDFWriteObject{9516}{<< /Names [<6571756174696F6E2E31302E342E332E3137> \MusicPDFRef{12679} <6571756174696F6E2E31302E342E332E3138> \MusicPDFRef{12680} <6571756174696F6E2E31302E342E332E3139> \MusicPDFRef{12681} <6571756174696F6E2E31302E342E332E32> \MusicPDFRef{12682} <6571756174696F6E2E31302E342E332E3230> \MusicPDFRef{12683} <6571756174696F6E2E31302E342E332E3231> \MusicPDFRef{12684}] /Limits [<6571756174696F6E2E31302E342E332E3137> <6571756174696F6E2E31302E342E332E3231>] >>}
\MusicPDFWriteObject{9517}{<< /Names [<6571756174696F6E2E31302E342E332E3232> \MusicPDFRef{12685} <6571756174696F6E2E31302E342E332E3233> \MusicPDFRef{12686} <6571756174696F6E2E31302E342E332E3234> \MusicPDFRef{12687} <6571756174696F6E2E31302E342E332E3235> \MusicPDFRef{12688} <6571756174696F6E2E31302E342E332E3236> \MusicPDFRef{12689} <6571756174696F6E2E31302E342E332E3237> \MusicPDFRef{12690}] /Limits [<6571756174696F6E2E31302E342E332E3232> <6571756174696F6E2E31302E342E332E3237>] >>}
\MusicPDFWriteObject{9518}{<< /Names [<6571756174696F6E2E31302E342E332E3238> \MusicPDFRef{12691} <6571756174696F6E2E31302E342E332E3239> \MusicPDFRef{12692} <6571756174696F6E2E31302E342E332E33> \MusicPDFRef{12693} <6571756174696F6E2E31302E342E332E3330> \MusicPDFRef{12694} <6571756174696F6E2E31302E342E332E3331> \MusicPDFRef{12695} <6571756174696F6E2E31302E342E332E3332> \MusicPDFRef{12696}] /Limits [<6571756174696F6E2E31302E342E332E3238> <6571756174696F6E2E31302E342E332E3332>] >>}
\MusicPDFWriteObject{9519}{<< /Names [<6571756174696F6E2E31302E342E332E3333> \MusicPDFRef{12697} <6571756174696F6E2E31302E342E332E3334> \MusicPDFRef{12698} <6571756174696F6E2E31302E342E332E3335> \MusicPDFRef{12699} <6571756174696F6E2E31302E342E332E3336> \MusicPDFRef{12700} <6571756174696F6E2E31302E342E332E3337> \MusicPDFRef{12701} <6571756174696F6E2E31302E342E332E3338> \MusicPDFRef{12702}] /Limits [<6571756174696F6E2E31302E342E332E3333> <6571756174696F6E2E31302E342E332E3338>] >>}
\MusicPDFWriteObject{9520}{<< /Names [<6571756174696F6E2E31302E342E332E3339> \MusicPDFRef{12703} <6571756174696F6E2E31302E342E332E34> \MusicPDFRef{12704} <6571756174696F6E2E31302E342E332E3430> \MusicPDFRef{12705} <6571756174696F6E2E31302E342E332E3431> \MusicPDFRef{12706} <6571756174696F6E2E31302E342E332E3432> \MusicPDFRef{12707} <6571756174696F6E2E31302E342E332E3433> \MusicPDFRef{12708}] /Limits [<6571756174696F6E2E31302E342E332E3339> <6571756174696F6E2E31302E342E332E3433>] >>}
\MusicPDFWriteObject{9521}{<< /Names [<6571756174696F6E2E31302E342E332E3434> \MusicPDFRef{12709} <6571756174696F6E2E31302E342E332E3435> \MusicPDFRef{12710} <6571756174696F6E2E31302E342E332E3436> \MusicPDFRef{12711} <6571756174696F6E2E31302E342E332E3437> \MusicPDFRef{12712} <6571756174696F6E2E31302E342E332E3438> \MusicPDFRef{12713} <6571756174696F6E2E31302E342E332E3439> \MusicPDFRef{12714}] /Limits [<6571756174696F6E2E31302E342E332E3434> <6571756174696F6E2E31302E342E332E3439>] >>}
\MusicPDFWriteObject{9522}{<< /Names [<6571756174696F6E2E31302E342E332E35> \MusicPDFRef{12715} <6571756174696F6E2E31302E342E332E3530> \MusicPDFRef{12716} <6571756174696F6E2E31302E342E332E36> \MusicPDFRef{12717} <6571756174696F6E2E31302E342E332E37> \MusicPDFRef{12718} <6571756174696F6E2E31302E342E332E38> \MusicPDFRef{12719} <6571756174696F6E2E31302E342E332E39> \MusicPDFRef{12720}] /Limits [<6571756174696F6E2E31302E342E332E35> <6571756174696F6E2E31302E342E332E39>] >>}
\MusicPDFWriteObject{9523}{<< /Names [<6571756174696F6E2E31312E312E312E31> \MusicPDFRef{12721} <6571756174696F6E2E31312E312E312E3130> \MusicPDFRef{12722} <6571756174696F6E2E31312E312E312E3131> \MusicPDFRef{12723} <6571756174696F6E2E31312E312E312E3132> \MusicPDFRef{12724} <6571756174696F6E2E31312E312E312E3133> \MusicPDFRef{12725} <6571756174696F6E2E31312E312E312E3134> \MusicPDFRef{12726}] /Limits [<6571756174696F6E2E31312E312E312E31> <6571756174696F6E2E31312E312E312E3134>] >>}
\MusicPDFWriteObject{9524}{<< /Names [<6571756174696F6E2E31312E312E312E3135> \MusicPDFRef{12727} <6571756174696F6E2E31312E312E312E32> \MusicPDFRef{12728} <6571756174696F6E2E31312E312E312E33> \MusicPDFRef{12729} <6571756174696F6E2E31312E312E312E34> \MusicPDFRef{12730} <6571756174696F6E2E31312E312E312E35> \MusicPDFRef{12731} <6571756174696F6E2E31312E312E312E36> \MusicPDFRef{12732}] /Limits [<6571756174696F6E2E31312E312E312E3135> <6571756174696F6E2E31312E312E312E36>] >>}
\MusicPDFWriteObject{9525}{<< /Names [<6571756174696F6E2E31312E312E312E37> \MusicPDFRef{12733} <6571756174696F6E2E31312E312E312E38> \MusicPDFRef{12734} <6571756174696F6E2E31312E312E312E39> \MusicPDFRef{12735} <6571756174696F6E2E31312E312E322E31> \MusicPDFRef{12736} <6571756174696F6E2E31312E312E322E3130> \MusicPDFRef{12737} <6571756174696F6E2E31312E312E322E3131> \MusicPDFRef{12738}] /Limits [<6571756174696F6E2E31312E312E312E37> <6571756174696F6E2E31312E312E322E3131>] >>}
\MusicPDFWriteObject{9526}{<< /Names [<6571756174696F6E2E31312E312E322E3132> \MusicPDFRef{12739} <6571756174696F6E2E31312E312E322E3133> \MusicPDFRef{12740} <6571756174696F6E2E31312E312E322E3134> \MusicPDFRef{12741} <6571756174696F6E2E31312E312E322E3135> \MusicPDFRef{12742} <6571756174696F6E2E31312E312E322E3136> \MusicPDFRef{12743} <6571756174696F6E2E31312E312E322E32> \MusicPDFRef{12744}] /Limits [<6571756174696F6E2E31312E312E322E3132> <6571756174696F6E2E31312E312E322E32>] >>}
\MusicPDFWriteObject{9527}{<< /Names [<6571756174696F6E2E31312E312E322E33> \MusicPDFRef{12745} <6571756174696F6E2E31312E312E322E34> \MusicPDFRef{12746} <6571756174696F6E2E31312E312E322E35> \MusicPDFRef{12747} <6571756174696F6E2E31312E312E322E36> \MusicPDFRef{12748} <6571756174696F6E2E31312E312E322E37> \MusicPDFRef{12749} <6571756174696F6E2E31312E312E322E38> \MusicPDFRef{12750}] /Limits [<6571756174696F6E2E31312E312E322E33> <6571756174696F6E2E31312E312E322E38>] >>}
\MusicPDFWriteObject{9528}{<< /Names [<6571756174696F6E2E31312E312E322E39> \MusicPDFRef{12751} <6571756174696F6E2E31312E322E302E31> \MusicPDFRef{12752} <6571756174696F6E2E31312E322E302E3130> \MusicPDFRef{12753} <6571756174696F6E2E31312E322E302E32> \MusicPDFRef{12754} <6571756174696F6E2E31312E322E302E33> \MusicPDFRef{12755} <6571756174696F6E2E31312E322E302E34> \MusicPDFRef{12756}] /Limits [<6571756174696F6E2E31312E312E322E39> <6571756174696F6E2E31312E322E302E34>] >>}
\MusicPDFWriteObject{9529}{<< /Names [<6571756174696F6E2E31312E322E302E35> \MusicPDFRef{12757} <6571756174696F6E2E31312E322E302E36> \MusicPDFRef{12758} <6571756174696F6E2E31312E322E302E37> \MusicPDFRef{12759} <6571756174696F6E2E31312E322E302E38> \MusicPDFRef{12760} <6571756174696F6E2E31312E322E302E39> \MusicPDFRef{12761} <6571756174696F6E2E31312E322E312E31> \MusicPDFRef{12762}] /Limits [<6571756174696F6E2E31312E322E302E35> <6571756174696F6E2E31312E322E312E31>] >>}
\MusicPDFWriteObject{9530}{<< /Names [<6571756174696F6E2E31312E322E312E3130> \MusicPDFRef{12763} <6571756174696F6E2E31312E322E312E32> \MusicPDFRef{12764} <6571756174696F6E2E31312E322E312E33> \MusicPDFRef{12765} <6571756174696F6E2E31312E322E312E34> \MusicPDFRef{12766} <6571756174696F6E2E31312E322E312E35> \MusicPDFRef{12767} <6571756174696F6E2E31312E322E312E36> \MusicPDFRef{12768}] /Limits [<6571756174696F6E2E31312E322E312E3130> <6571756174696F6E2E31312E322E312E36>] >>}
\MusicPDFWriteObject{9531}{<< /Names [<6571756174696F6E2E31312E322E312E37> \MusicPDFRef{12769} <6571756174696F6E2E31312E322E312E38> \MusicPDFRef{12770} <6571756174696F6E2E31312E322E312E39> \MusicPDFRef{12771} <6571756174696F6E2E31312E322E322E31> \MusicPDFRef{12772} <6571756174696F6E2E31312E322E322E3130> \MusicPDFRef{12773} <6571756174696F6E2E31312E322E322E3131> \MusicPDFRef{12774}] /Limits [<6571756174696F6E2E31312E322E312E37> <6571756174696F6E2E31312E322E322E3131>] >>}
\MusicPDFWriteObject{9532}{<< /Names [<6571756174696F6E2E31312E322E322E3132> \MusicPDFRef{12775} <6571756174696F6E2E31312E322E322E3133> \MusicPDFRef{12776} <6571756174696F6E2E31312E322E322E32> \MusicPDFRef{12777} <6571756174696F6E2E31312E322E322E33> \MusicPDFRef{12778} <6571756174696F6E2E31312E322E322E34> \MusicPDFRef{12779} <6571756174696F6E2E31312E322E322E35> \MusicPDFRef{12780}] /Limits [<6571756174696F6E2E31312E322E322E3132> <6571756174696F6E2E31312E322E322E35>] >>}
\MusicPDFWriteObject{9533}{<< /Names [<6571756174696F6E2E31312E322E322E36> \MusicPDFRef{12781} <6571756174696F6E2E31312E322E322E37> \MusicPDFRef{12782} <6571756174696F6E2E31312E322E322E38> \MusicPDFRef{12783} <6571756174696F6E2E31312E322E322E39> \MusicPDFRef{12784} <6571756174696F6E2E31312E332E312E31> \MusicPDFRef{12785} <6571756174696F6E2E31312E332E312E32> \MusicPDFRef{12786}] /Limits [<6571756174696F6E2E31312E322E322E36> <6571756174696F6E2E31312E332E312E32>] >>}
\MusicPDFWriteObject{9534}{<< /Names [<6571756174696F6E2E31312E332E312E33> \MusicPDFRef{12787} <6571756174696F6E2E31312E332E312E34> \MusicPDFRef{12788} <6571756174696F6E2E31312E332E322E31> \MusicPDFRef{12789} <6571756174696F6E2E31312E332E322E32> \MusicPDFRef{12790} <6571756174696F6E2E31312E332E322E33> \MusicPDFRef{12791} <6571756174696F6E2E31312E332E322E34> \MusicPDFRef{12792}] /Limits [<6571756174696F6E2E31312E332E312E33> <6571756174696F6E2E31312E332E322E34>] >>}
\MusicPDFWriteObject{9535}{<< /Names [<6571756174696F6E2E31312E332E322E35> \MusicPDFRef{12793} <6571756174696F6E2E31322E312E312E31> \MusicPDFRef{12794} <6571756174696F6E2E31322E312E312E3130> \MusicPDFRef{12795} <6571756174696F6E2E31322E312E312E3131> \MusicPDFRef{12796} <6571756174696F6E2E31322E312E312E3132> \MusicPDFRef{12797} <6571756174696F6E2E31322E312E312E3133> \MusicPDFRef{12798}] /Limits [<6571756174696F6E2E31312E332E322E35> <6571756174696F6E2E31322E312E312E3133>] >>}
\MusicPDFWriteObject{9536}{<< /Names [<6571756174696F6E2E31322E312E312E3134> \MusicPDFRef{12799} <6571756174696F6E2E31322E312E312E3135> \MusicPDFRef{12800} <6571756174696F6E2E31322E312E312E3136> \MusicPDFRef{12801} <6571756174696F6E2E31322E312E312E3137> \MusicPDFRef{12802} <6571756174696F6E2E31322E312E312E3138> \MusicPDFRef{12803} <6571756174696F6E2E31322E312E312E3139> \MusicPDFRef{12804}] /Limits [<6571756174696F6E2E31322E312E312E3134> <6571756174696F6E2E31322E312E312E3139>] >>}
\MusicPDFWriteObject{9537}{<< /Names [<6571756174696F6E2E31322E312E312E32> \MusicPDFRef{12805} <6571756174696F6E2E31322E312E312E3230> \MusicPDFRef{12806} <6571756174696F6E2E31322E312E312E3231> \MusicPDFRef{12807} <6571756174696F6E2E31322E312E312E3232> \MusicPDFRef{12808} <6571756174696F6E2E31322E312E312E3233> \MusicPDFRef{12809} <6571756174696F6E2E31322E312E312E3234> \MusicPDFRef{12810}] /Limits [<6571756174696F6E2E31322E312E312E32> <6571756174696F6E2E31322E312E312E3234>] >>}
\MusicPDFWriteObject{9538}{<< /Names [<6571756174696F6E2E31322E312E312E33> \MusicPDFRef{12811} <6571756174696F6E2E31322E312E312E34> \MusicPDFRef{12812} <6571756174696F6E2E31322E312E312E35> \MusicPDFRef{12813} <6571756174696F6E2E31322E312E312E36> \MusicPDFRef{12814} <6571756174696F6E2E31322E312E312E37> \MusicPDFRef{12815} <6571756174696F6E2E31322E312E312E38> \MusicPDFRef{12816}] /Limits [<6571756174696F6E2E31322E312E312E33> <6571756174696F6E2E31322E312E312E38>] >>}
\MusicPDFWriteObject{9539}{<< /Names [<6571756174696F6E2E31322E312E312E39> \MusicPDFRef{12817} <6571756174696F6E2E31322E322E312E31> \MusicPDFRef{12818} <6571756174696F6E2E31322E322E312E3130> \MusicPDFRef{12819} <6571756174696F6E2E31322E322E312E3131> \MusicPDFRef{12820} <6571756174696F6E2E31322E322E312E3132> \MusicPDFRef{12821} <6571756174696F6E2E31322E322E312E3133> \MusicPDFRef{12822}] /Limits [<6571756174696F6E2E31322E312E312E39> <6571756174696F6E2E31322E322E312E3133>] >>}
\MusicPDFWriteObject{9540}{<< /Names [<6571756174696F6E2E31322E322E312E32> \MusicPDFRef{12823} <6571756174696F6E2E31322E322E312E33> \MusicPDFRef{12824} <6571756174696F6E2E31322E322E312E34> \MusicPDFRef{12825} <6571756174696F6E2E31322E322E312E35> \MusicPDFRef{12826} <6571756174696F6E2E31322E322E312E36> \MusicPDFRef{12827} <6571756174696F6E2E31322E322E312E37> \MusicPDFRef{12828}] /Limits [<6571756174696F6E2E31322E322E312E32> <6571756174696F6E2E31322E322E312E37>] >>}
\MusicPDFWriteObject{9541}{<< /Names [<6571756174696F6E2E31322E322E312E38> \MusicPDFRef{12829} <6571756174696F6E2E31322E322E312E39> \MusicPDFRef{12830} <6571756174696F6E2E31322E322E322E31> \MusicPDFRef{12831} <6571756174696F6E2E31322E322E322E3130> \MusicPDFRef{12832} <6571756174696F6E2E31322E322E322E3131> \MusicPDFRef{12833} <6571756174696F6E2E31322E322E322E3132> \MusicPDFRef{12834}] /Limits [<6571756174696F6E2E31322E322E312E38> <6571756174696F6E2E31322E322E322E3132>] >>}
\MusicPDFWriteObject{9542}{<< /Names [<6571756174696F6E2E31322E322E322E3133> \MusicPDFRef{12835} <6571756174696F6E2E31322E322E322E3134> \MusicPDFRef{12836} <6571756174696F6E2E31322E322E322E3135> \MusicPDFRef{12837} <6571756174696F6E2E31322E322E322E3136> \MusicPDFRef{12838} <6571756174696F6E2E31322E322E322E3137> \MusicPDFRef{12839} <6571756174696F6E2E31322E322E322E3138> \MusicPDFRef{12840}] /Limits [<6571756174696F6E2E31322E322E322E3133> <6571756174696F6E2E31322E322E322E3138>] >>}
\MusicPDFWriteObject{9543}{<< /Names [<6571756174696F6E2E31322E322E322E3139> \MusicPDFRef{12841} <6571756174696F6E2E31322E322E322E32> \MusicPDFRef{12842} <6571756174696F6E2E31322E322E322E3230> \MusicPDFRef{12843} <6571756174696F6E2E31322E322E322E3231> \MusicPDFRef{12844} <6571756174696F6E2E31322E322E322E3232> \MusicPDFRef{12845} <6571756174696F6E2E31322E322E322E3233> \MusicPDFRef{12846}] /Limits [<6571756174696F6E2E31322E322E322E3139> <6571756174696F6E2E31322E322E322E3233>] >>}
\MusicPDFWriteObject{9544}{<< /Names [<6571756174696F6E2E31322E322E322E3234> \MusicPDFRef{12847} <6571756174696F6E2E31322E322E322E3235> \MusicPDFRef{12848} <6571756174696F6E2E31322E322E322E3236> \MusicPDFRef{12849} <6571756174696F6E2E31322E322E322E3237> \MusicPDFRef{12850} <6571756174696F6E2E31322E322E322E3238> \MusicPDFRef{12851} <6571756174696F6E2E31322E322E322E3239> \MusicPDFRef{12852}] /Limits [<6571756174696F6E2E31322E322E322E3234> <6571756174696F6E2E31322E322E322E3239>] >>}
\MusicPDFWriteObject{9545}{<< /Names [<6571756174696F6E2E31322E322E322E33> \MusicPDFRef{12853} <6571756174696F6E2E31322E322E322E3330> \MusicPDFRef{12854} <6571756174696F6E2E31322E322E322E3331> \MusicPDFRef{12855} <6571756174696F6E2E31322E322E322E34> \MusicPDFRef{12856} <6571756174696F6E2E31322E322E322E35> \MusicPDFRef{12857} <6571756174696F6E2E31322E322E322E36> \MusicPDFRef{12858}] /Limits [<6571756174696F6E2E31322E322E322E33> <6571756174696F6E2E31322E322E322E36>] >>}
\MusicPDFWriteObject{9546}{<< /Names [<6571756174696F6E2E31322E322E322E37> \MusicPDFRef{12859} <6571756174696F6E2E31322E322E322E38> \MusicPDFRef{12860} <6571756174696F6E2E31322E322E322E39> \MusicPDFRef{12861} <6571756174696F6E2E31322E342E312E31> \MusicPDFRef{12862} <6571756174696F6E2E31322E342E312E3130> \MusicPDFRef{12863} <6571756174696F6E2E31322E342E312E3131> \MusicPDFRef{12864}] /Limits [<6571756174696F6E2E31322E322E322E37> <6571756174696F6E2E31322E342E312E3131>] >>}
\MusicPDFWriteObject{9547}{<< /Names [<6571756174696F6E2E31322E342E312E3132> \MusicPDFRef{12865} <6571756174696F6E2E31322E342E312E3133> \MusicPDFRef{12866} <6571756174696F6E2E31322E342E312E3134> \MusicPDFRef{12867} <6571756174696F6E2E31322E342E312E3135> \MusicPDFRef{12868} <6571756174696F6E2E31322E342E312E3136> \MusicPDFRef{12869} <6571756174696F6E2E31322E342E312E3137> \MusicPDFRef{12870}] /Limits [<6571756174696F6E2E31322E342E312E3132> <6571756174696F6E2E31322E342E312E3137>] >>}
\MusicPDFWriteObject{9548}{<< /Names [<6571756174696F6E2E31322E342E312E3138> \MusicPDFRef{12871} <6571756174696F6E2E31322E342E312E3139> \MusicPDFRef{12872} <6571756174696F6E2E31322E342E312E32> \MusicPDFRef{12873} <6571756174696F6E2E31322E342E312E3230> \MusicPDFRef{12874} <6571756174696F6E2E31322E342E312E3231> \MusicPDFRef{12875} <6571756174696F6E2E31322E342E312E3232> \MusicPDFRef{12876}] /Limits [<6571756174696F6E2E31322E342E312E3138> <6571756174696F6E2E31322E342E312E3232>] >>}
\MusicPDFWriteObject{9549}{<< /Names [<6571756174696F6E2E31322E342E312E3233> \MusicPDFRef{12877} <6571756174696F6E2E31322E342E312E3234> \MusicPDFRef{12878} <6571756174696F6E2E31322E342E312E3235> \MusicPDFRef{12879} <6571756174696F6E2E31322E342E312E3236> \MusicPDFRef{12880} <6571756174696F6E2E31322E342E312E3237> \MusicPDFRef{12881} <6571756174696F6E2E31322E342E312E3238> \MusicPDFRef{12882}] /Limits [<6571756174696F6E2E31322E342E312E3233> <6571756174696F6E2E31322E342E312E3238>] >>}
\MusicPDFWriteObject{9550}{<< /Names [<6571756174696F6E2E31322E342E312E33> \MusicPDFRef{12883} <6571756174696F6E2E31322E342E312E34> \MusicPDFRef{12884} <6571756174696F6E2E31322E342E312E35> \MusicPDFRef{12885} <6571756174696F6E2E31322E342E312E36> \MusicPDFRef{12886} <6571756174696F6E2E31322E342E312E37> \MusicPDFRef{12887} <6571756174696F6E2E31322E342E312E38> \MusicPDFRef{12888}] /Limits [<6571756174696F6E2E31322E342E312E33> <6571756174696F6E2E31322E342E312E38>] >>}
\MusicPDFWriteObject{9551}{<< /Names [<6571756174696F6E2E31322E342E312E39> \MusicPDFRef{12889} <6571756174696F6E2E31322E352E312E31> \MusicPDFRef{12890} <6571756174696F6E2E31322E352E312E32> \MusicPDFRef{12891} <6571756174696F6E2E31322E352E312E33> \MusicPDFRef{12892} <6571756174696F6E2E31322E352E312E34> \MusicPDFRef{12893} <6571756174696F6E2E31332E332E312E31> \MusicPDFRef{12894}] /Limits [<6571756174696F6E2E31322E342E312E39> <6571756174696F6E2E31332E332E312E31>] >>}
\MusicPDFWriteObject{9552}{<< /Names [<6571756174696F6E2E31332E332E312E32> \MusicPDFRef{12895} <6571756174696F6E2E31332E332E312E33> \MusicPDFRef{12896} <6571756174696F6E2E31332E332E312E34> \MusicPDFRef{12897} <6571756174696F6E2E31332E332E312E35> \MusicPDFRef{12898} <6571756174696F6E2E31342E312E312E31> \MusicPDFRef{12899} <6571756174696F6E2E31342E312E312E3130> \MusicPDFRef{12900}] /Limits [<6571756174696F6E2E31332E332E312E32> <6571756174696F6E2E31342E312E312E3130>] >>}
\MusicPDFWriteObject{9553}{<< /Names [<6571756174696F6E2E31342E312E312E3131> \MusicPDFRef{12901} <6571756174696F6E2E31342E312E312E3132> \MusicPDFRef{12902} <6571756174696F6E2E31342E312E312E3133> \MusicPDFRef{12903} <6571756174696F6E2E31342E312E312E32> \MusicPDFRef{12904} <6571756174696F6E2E31342E312E312E33> \MusicPDFRef{12905} <6571756174696F6E2E31342E312E312E34> \MusicPDFRef{12906}] /Limits [<6571756174696F6E2E31342E312E312E3131> <6571756174696F6E2E31342E312E312E34>] >>}
\MusicPDFWriteObject{9554}{<< /Names [<6571756174696F6E2E31342E312E312E35> \MusicPDFRef{12907} <6571756174696F6E2E31342E312E312E36> \MusicPDFRef{12908} <6571756174696F6E2E31342E312E312E37> \MusicPDFRef{12909} <6571756174696F6E2E31342E312E312E38> \MusicPDFRef{12910} <6571756174696F6E2E31342E312E312E39> \MusicPDFRef{12911} <6571756174696F6E2E31342E312E322E31> \MusicPDFRef{12912}] /Limits [<6571756174696F6E2E31342E312E312E35> <6571756174696F6E2E31342E312E322E31>] >>}
\MusicPDFWriteObject{9555}{<< /Names [<6571756174696F6E2E31342E312E322E32> \MusicPDFRef{12913} <6571756174696F6E2E31342E312E322E33> \MusicPDFRef{12914} <6571756174696F6E2E31342E312E322E34> \MusicPDFRef{12915} <6571756174696F6E2E31342E312E322E35> \MusicPDFRef{12916} <6571756174696F6E2E31342E312E322E36> \MusicPDFRef{12917} <6571756174696F6E2E31342E312E322E37> \MusicPDFRef{12918}] /Limits [<6571756174696F6E2E31342E312E322E32> <6571756174696F6E2E31342E312E322E37>] >>}
\MusicPDFWriteObject{9556}{<< /Names [<6571756174696F6E2E31342E322E302E31> \MusicPDFRef{12919} <6571756174696F6E2E322E312E312E31> \MusicPDFRef{12920} <6571756174696F6E2E322E312E312E32> \MusicPDFRef{12921} <6571756174696F6E2E322E312E312E33> \MusicPDFRef{12922} <6571756174696F6E2E322E312E312E34> \MusicPDFRef{12923} <6571756174696F6E2E322E312E322E31> \MusicPDFRef{12924}] /Limits [<6571756174696F6E2E31342E322E302E31> <6571756174696F6E2E322E312E322E31>] >>}
\MusicPDFWriteObject{9557}{<< /Names [<6571756174696F6E2E322E312E322E32> \MusicPDFRef{12925} <6571756174696F6E2E322E312E322E33> \MusicPDFRef{12926} <6571756174696F6E2E322E312E322E34> \MusicPDFRef{12927} <6571756174696F6E2E322E312E322E35> \MusicPDFRef{12928} <6571756174696F6E2E322E312E332E31> \MusicPDFRef{12929} <6571756174696F6E2E322E312E332E32> \MusicPDFRef{12930}] /Limits [<6571756174696F6E2E322E312E322E32> <6571756174696F6E2E322E312E332E32>] >>}
\MusicPDFWriteObject{9558}{<< /Names [<6571756174696F6E2E322E312E332E33> \MusicPDFRef{12931} <6571756174696F6E2E322E312E342E31> \MusicPDFRef{12932} <6571756174696F6E2E322E312E342E3130> \MusicPDFRef{12933} <6571756174696F6E2E322E312E342E3131> \MusicPDFRef{12934} <6571756174696F6E2E322E312E342E3132> \MusicPDFRef{12935} <6571756174696F6E2E322E312E342E3133> \MusicPDFRef{12936}] /Limits [<6571756174696F6E2E322E312E332E33> <6571756174696F6E2E322E312E342E3133>] >>}
\MusicPDFWriteObject{9559}{<< /Names [<6571756174696F6E2E322E312E342E3134> \MusicPDFRef{12937} <6571756174696F6E2E322E312E342E32> \MusicPDFRef{12938} <6571756174696F6E2E322E312E342E33> \MusicPDFRef{12939} <6571756174696F6E2E322E312E342E34> \MusicPDFRef{12940} <6571756174696F6E2E322E312E342E35> \MusicPDFRef{12941} <6571756174696F6E2E322E312E342E36> \MusicPDFRef{12942}] /Limits [<6571756174696F6E2E322E312E342E3134> <6571756174696F6E2E322E312E342E36>] >>}
\MusicPDFWriteObject{9560}{<< /Names [<6571756174696F6E2E322E312E342E37> \MusicPDFRef{12943} <6571756174696F6E2E322E312E342E38> \MusicPDFRef{12944} <6571756174696F6E2E322E312E342E39> \MusicPDFRef{12945} <6571756174696F6E2E322E322E312E31> \MusicPDFRef{12946} <6571756174696F6E2E322E322E312E32> \MusicPDFRef{12947} <6571756174696F6E2E322E322E312E33> \MusicPDFRef{12948}] /Limits [<6571756174696F6E2E322E312E342E37> <6571756174696F6E2E322E322E312E33>] >>}
\MusicPDFWriteObject{9561}{<< /Names [<6571756174696F6E2E322E322E312E34> \MusicPDFRef{12949} <6571756174696F6E2E322E322E312E35> \MusicPDFRef{12950} <6571756174696F6E2E322E322E322E31> \MusicPDFRef{12951} <6571756174696F6E2E322E322E322E3130> \MusicPDFRef{12952} <6571756174696F6E2E322E322E322E3131> \MusicPDFRef{12953} <6571756174696F6E2E322E322E322E3132> \MusicPDFRef{12954}] /Limits [<6571756174696F6E2E322E322E312E34> <6571756174696F6E2E322E322E322E3132>] >>}
\MusicPDFWriteObject{9562}{<< /Names [<6571756174696F6E2E322E322E322E3133> \MusicPDFRef{12955} <6571756174696F6E2E322E322E322E3134> \MusicPDFRef{12956} <6571756174696F6E2E322E322E322E32> \MusicPDFRef{12957} <6571756174696F6E2E322E322E322E33> \MusicPDFRef{12958} <6571756174696F6E2E322E322E322E34> \MusicPDFRef{12959} <6571756174696F6E2E322E322E322E35> \MusicPDFRef{12960}] /Limits [<6571756174696F6E2E322E322E322E3133> <6571756174696F6E2E322E322E322E35>] >>}
\MusicPDFWriteObject{9563}{<< /Names [<6571756174696F6E2E322E322E322E36> \MusicPDFRef{12961} <6571756174696F6E2E322E322E322E37> \MusicPDFRef{12962} <6571756174696F6E2E322E322E322E38> \MusicPDFRef{12963} <6571756174696F6E2E322E322E322E39> \MusicPDFRef{12964} <6571756174696F6E2E322E322E332E31> \MusicPDFRef{12965} <6571756174696F6E2E322E322E332E3130> \MusicPDFRef{12966}] /Limits [<6571756174696F6E2E322E322E322E36> <6571756174696F6E2E322E322E332E3130>] >>}
\MusicPDFWriteObject{9564}{<< /Names [<6571756174696F6E2E322E322E332E3131> \MusicPDFRef{12967} <6571756174696F6E2E322E322E332E3132> \MusicPDFRef{12968} <6571756174696F6E2E322E322E332E3133> \MusicPDFRef{12969} <6571756174696F6E2E322E322E332E3134> \MusicPDFRef{12970} <6571756174696F6E2E322E322E332E3135> \MusicPDFRef{12971} <6571756174696F6E2E322E322E332E32> \MusicPDFRef{12972}] /Limits [<6571756174696F6E2E322E322E332E3131> <6571756174696F6E2E322E322E332E32>] >>}
\MusicPDFWriteObject{9565}{<< /Names [<6571756174696F6E2E322E322E332E33> \MusicPDFRef{12973} <6571756174696F6E2E322E322E332E34> \MusicPDFRef{12974} <6571756174696F6E2E322E322E332E35> \MusicPDFRef{12975} <6571756174696F6E2E322E322E332E36> \MusicPDFRef{12976} <6571756174696F6E2E322E322E332E37> \MusicPDFRef{12977} <6571756174696F6E2E322E322E332E38> \MusicPDFRef{12978}] /Limits [<6571756174696F6E2E322E322E332E33> <6571756174696F6E2E322E322E332E38>] >>}
\MusicPDFWriteObject{9566}{<< /Names [<6571756174696F6E2E322E322E332E39> \MusicPDFRef{12979} <6571756174696F6E2E322E322E342E31> \MusicPDFRef{12980} <6571756174696F6E2E322E322E342E3130> \MusicPDFRef{12981} <6571756174696F6E2E322E322E342E3131> \MusicPDFRef{12982} <6571756174696F6E2E322E322E342E3132> \MusicPDFRef{12983} <6571756174696F6E2E322E322E342E3133> \MusicPDFRef{12984}] /Limits [<6571756174696F6E2E322E322E332E39> <6571756174696F6E2E322E322E342E3133>] >>}
\MusicPDFWriteObject{9567}{<< /Names [<6571756174696F6E2E322E322E342E3134> \MusicPDFRef{12985} <6571756174696F6E2E322E322E342E3135> \MusicPDFRef{12986} <6571756174696F6E2E322E322E342E3136> \MusicPDFRef{12987} <6571756174696F6E2E322E322E342E3137> \MusicPDFRef{12988} <6571756174696F6E2E322E322E342E3138> \MusicPDFRef{12989} <6571756174696F6E2E322E322E342E3139> \MusicPDFRef{12990}] /Limits [<6571756174696F6E2E322E322E342E3134> <6571756174696F6E2E322E322E342E3139>] >>}
\MusicPDFWriteObject{9568}{<< /Names [<6571756174696F6E2E322E322E342E32> \MusicPDFRef{12991} <6571756174696F6E2E322E322E342E3230> \MusicPDFRef{12992} <6571756174696F6E2E322E322E342E3231> \MusicPDFRef{12993} <6571756174696F6E2E322E322E342E3232> \MusicPDFRef{12994} <6571756174696F6E2E322E322E342E3233> \MusicPDFRef{12995} <6571756174696F6E2E322E322E342E3234> \MusicPDFRef{12996}] /Limits [<6571756174696F6E2E322E322E342E32> <6571756174696F6E2E322E322E342E3234>] >>}
\MusicPDFWriteObject{9569}{<< /Names [<6571756174696F6E2E322E322E342E3235> \MusicPDFRef{12997} <6571756174696F6E2E322E322E342E3236> \MusicPDFRef{12998} <6571756174696F6E2E322E322E342E3237> \MusicPDFRef{12999} <6571756174696F6E2E322E322E342E3238> \MusicPDFRef{13000} <6571756174696F6E2E322E322E342E3239> \MusicPDFRef{13001} <6571756174696F6E2E322E322E342E33> \MusicPDFRef{13002}] /Limits [<6571756174696F6E2E322E322E342E3235> <6571756174696F6E2E322E322E342E33>] >>}
\MusicPDFWriteObject{9570}{<< /Names [<6571756174696F6E2E322E322E342E3330> \MusicPDFRef{13003} <6571756174696F6E2E322E322E342E3331> \MusicPDFRef{13004} <6571756174696F6E2E322E322E342E3332> \MusicPDFRef{13005} <6571756174696F6E2E322E322E342E3333> \MusicPDFRef{13006} <6571756174696F6E2E322E322E342E3334> \MusicPDFRef{13007} <6571756174696F6E2E322E322E342E3335> \MusicPDFRef{13008}] /Limits [<6571756174696F6E2E322E322E342E3330> <6571756174696F6E2E322E322E342E3335>] >>}
\MusicPDFWriteObject{9571}{<< /Names [<6571756174696F6E2E322E322E342E3336> \MusicPDFRef{13009} <6571756174696F6E2E322E322E342E3337> \MusicPDFRef{13010} <6571756174696F6E2E322E322E342E3338> \MusicPDFRef{13011} <6571756174696F6E2E322E322E342E3339> \MusicPDFRef{13012} <6571756174696F6E2E322E322E342E34> \MusicPDFRef{13013} <6571756174696F6E2E322E322E342E3430> \MusicPDFRef{13014}] /Limits [<6571756174696F6E2E322E322E342E3336> <6571756174696F6E2E322E322E342E3430>] >>}
\MusicPDFWriteObject{9572}{<< /Names [<6571756174696F6E2E322E322E342E3431> \MusicPDFRef{13015} <6571756174696F6E2E322E322E342E3432> \MusicPDFRef{13016} <6571756174696F6E2E322E322E342E35> \MusicPDFRef{13017} <6571756174696F6E2E322E322E342E36> \MusicPDFRef{13018} <6571756174696F6E2E322E322E342E37> \MusicPDFRef{13019} <6571756174696F6E2E322E322E342E38> \MusicPDFRef{13020}] /Limits [<6571756174696F6E2E322E322E342E3431> <6571756174696F6E2E322E322E342E38>] >>}
\MusicPDFWriteObject{9573}{<< /Names [<6571756174696F6E2E322E322E342E39> \MusicPDFRef{13021} <6571756174696F6E2E322E322E352E31> \MusicPDFRef{13022} <6571756174696F6E2E322E322E352E3130> \MusicPDFRef{13023} <6571756174696F6E2E322E322E352E3131> \MusicPDFRef{13024} <6571756174696F6E2E322E322E352E3132> \MusicPDFRef{13025} <6571756174696F6E2E322E322E352E3133> \MusicPDFRef{13026}] /Limits [<6571756174696F6E2E322E322E342E39> <6571756174696F6E2E322E322E352E3133>] >>}
\MusicPDFWriteObject{9574}{<< /Names [<6571756174696F6E2E322E322E352E3134> \MusicPDFRef{13027} <6571756174696F6E2E322E322E352E3135> \MusicPDFRef{13028} <6571756174696F6E2E322E322E352E32> \MusicPDFRef{13029} <6571756174696F6E2E322E322E352E33> \MusicPDFRef{13030} <6571756174696F6E2E322E322E352E34> \MusicPDFRef{13031} <6571756174696F6E2E322E322E352E35> \MusicPDFRef{13032}] /Limits [<6571756174696F6E2E322E322E352E3134> <6571756174696F6E2E322E322E352E35>] >>}
\MusicPDFWriteObject{9575}{<< /Names [<6571756174696F6E2E322E322E352E36> \MusicPDFRef{13033} <6571756174696F6E2E322E322E352E37> \MusicPDFRef{13034} <6571756174696F6E2E322E322E352E38> \MusicPDFRef{13035} <6571756174696F6E2E322E322E352E39> \MusicPDFRef{13036} <6571756174696F6E2E322E332E312E31> \MusicPDFRef{13037} <6571756174696F6E2E322E332E312E3130> \MusicPDFRef{13038}] /Limits [<6571756174696F6E2E322E322E352E36> <6571756174696F6E2E322E332E312E3130>] >>}
\MusicPDFWriteObject{9576}{<< /Names [<6571756174696F6E2E322E332E312E3131> \MusicPDFRef{13039} <6571756174696F6E2E322E332E312E3132> \MusicPDFRef{13040} <6571756174696F6E2E322E332E312E3133> \MusicPDFRef{13041} <6571756174696F6E2E322E332E312E3134> \MusicPDFRef{13042} <6571756174696F6E2E322E332E312E32> \MusicPDFRef{13043} <6571756174696F6E2E322E332E312E33> \MusicPDFRef{13044}] /Limits [<6571756174696F6E2E322E332E312E3131> <6571756174696F6E2E322E332E312E33>] >>}
\MusicPDFWriteObject{9577}{<< /Names [<6571756174696F6E2E322E332E312E34> \MusicPDFRef{13045} <6571756174696F6E2E322E332E312E35> \MusicPDFRef{13046} <6571756174696F6E2E322E332E312E36> \MusicPDFRef{13047} <6571756174696F6E2E322E332E312E37> \MusicPDFRef{13048} <6571756174696F6E2E322E332E312E38> \MusicPDFRef{13049} <6571756174696F6E2E322E332E312E39> \MusicPDFRef{13050}] /Limits [<6571756174696F6E2E322E332E312E34> <6571756174696F6E2E322E332E312E39>] >>}
\MusicPDFWriteObject{9578}{<< /Names [<6571756174696F6E2E322E332E322E31> \MusicPDFRef{13051} <6571756174696F6E2E322E332E322E32> \MusicPDFRef{13052} <6571756174696F6E2E322E332E322E33> \MusicPDFRef{13053} <6571756174696F6E2E322E332E322E34> \MusicPDFRef{13054} <6571756174696F6E2E322E342E312E31> \MusicPDFRef{13055} <6571756174696F6E2E322E342E312E3130> \MusicPDFRef{13056}] /Limits [<6571756174696F6E2E322E332E322E31> <6571756174696F6E2E322E342E312E3130>] >>}
\MusicPDFWriteObject{9579}{<< /Names [<6571756174696F6E2E322E342E312E3131> \MusicPDFRef{13057} <6571756174696F6E2E322E342E312E3132> \MusicPDFRef{13058} <6571756174696F6E2E322E342E312E3133> \MusicPDFRef{13059} <6571756174696F6E2E322E342E312E3134> \MusicPDFRef{13060} <6571756174696F6E2E322E342E312E3135> \MusicPDFRef{13061} <6571756174696F6E2E322E342E312E3136> \MusicPDFRef{13062}] /Limits [<6571756174696F6E2E322E342E312E3131> <6571756174696F6E2E322E342E312E3136>] >>}
\MusicPDFWriteObject{9580}{<< /Names [<6571756174696F6E2E322E342E312E32> \MusicPDFRef{13063} <6571756174696F6E2E322E342E312E33> \MusicPDFRef{13064} <6571756174696F6E2E322E342E312E34> \MusicPDFRef{13065} <6571756174696F6E2E322E342E312E35> \MusicPDFRef{13066} <6571756174696F6E2E322E342E312E36> \MusicPDFRef{13067} <6571756174696F6E2E322E342E312E37> \MusicPDFRef{13068}] /Limits [<6571756174696F6E2E322E342E312E32> <6571756174696F6E2E322E342E312E37>] >>}
\MusicPDFWriteObject{9581}{<< /Names [<6571756174696F6E2E322E342E312E38> \MusicPDFRef{13069} <6571756174696F6E2E322E342E312E39> \MusicPDFRef{13070} <6571756174696F6E2E322E342E322E31> \MusicPDFRef{13071} <6571756174696F6E2E322E342E322E32> \MusicPDFRef{13072} <6571756174696F6E2E322E342E322E33> \MusicPDFRef{13073} <6571756174696F6E2E322E342E322E34> \MusicPDFRef{13074}] /Limits [<6571756174696F6E2E322E342E312E38> <6571756174696F6E2E322E342E322E34>] >>}
\MusicPDFWriteObject{9582}{<< /Names [<6571756174696F6E2E322E342E322E35> \MusicPDFRef{13075} <6571756174696F6E2E322E342E332E31> \MusicPDFRef{13076} <6571756174696F6E2E322E342E332E32> \MusicPDFRef{13077} <6571756174696F6E2E322E342E332E33> \MusicPDFRef{13078} <6571756174696F6E2E322E342E332E34> \MusicPDFRef{13079} <6571756174696F6E2E322E342E332E35> \MusicPDFRef{13080}] /Limits [<6571756174696F6E2E322E342E322E35> <6571756174696F6E2E322E342E332E35>] >>}
\MusicPDFWriteObject{9583}{<< /Names [<6571756174696F6E2E322E342E332E36> \MusicPDFRef{13081} <6571756174696F6E2E332E312E312E31> \MusicPDFRef{13082} <6571756174696F6E2E332E312E312E32> \MusicPDFRef{13083} <6571756174696F6E2E332E312E312E33> \MusicPDFRef{13084} <6571756174696F6E2E332E312E312E34> \MusicPDFRef{13085} <6571756174696F6E2E332E312E322E31> \MusicPDFRef{13086}] /Limits [<6571756174696F6E2E322E342E332E36> <6571756174696F6E2E332E312E322E31>] >>}
\MusicPDFWriteObject{9584}{<< /Names [<6571756174696F6E2E332E312E322E32> \MusicPDFRef{13087} <6571756174696F6E2E332E312E322E33> \MusicPDFRef{13088} <6571756174696F6E2E332E312E322E34> \MusicPDFRef{13089} <6571756174696F6E2E332E312E322E35> \MusicPDFRef{13090} <6571756174696F6E2E332E312E322E36> \MusicPDFRef{13091} <6571756174696F6E2E332E312E322E37> \MusicPDFRef{13092}] /Limits [<6571756174696F6E2E332E312E322E32> <6571756174696F6E2E332E312E322E37>] >>}
\MusicPDFWriteObject{9585}{<< /Names [<6571756174696F6E2E332E312E322E38> \MusicPDFRef{13093} <6571756174696F6E2E332E312E322E39> \MusicPDFRef{13094} <6571756174696F6E2E332E322E312E31> \MusicPDFRef{13095} <6571756174696F6E2E332E322E312E3130> \MusicPDFRef{13096} <6571756174696F6E2E332E322E312E32> \MusicPDFRef{13097} <6571756174696F6E2E332E322E312E33> \MusicPDFRef{13098}] /Limits [<6571756174696F6E2E332E312E322E38> <6571756174696F6E2E332E322E312E33>] >>}
\MusicPDFWriteObject{9586}{<< /Names [<6571756174696F6E2E332E322E312E34> \MusicPDFRef{13099} <6571756174696F6E2E332E322E312E35> \MusicPDFRef{13100} <6571756174696F6E2E332E322E312E36> \MusicPDFRef{13101} <6571756174696F6E2E332E322E312E37> \MusicPDFRef{13102} <6571756174696F6E2E332E322E312E38> \MusicPDFRef{13103} <6571756174696F6E2E332E322E312E39> \MusicPDFRef{13104}] /Limits [<6571756174696F6E2E332E322E312E34> <6571756174696F6E2E332E322E312E39>] >>}
\MusicPDFWriteObject{9587}{<< /Names [<6571756174696F6E2E332E322E322E31> \MusicPDFRef{13105} <6571756174696F6E2E332E322E322E3130> \MusicPDFRef{13106} <6571756174696F6E2E332E322E322E3131> \MusicPDFRef{13107} <6571756174696F6E2E332E322E322E32> \MusicPDFRef{13108} <6571756174696F6E2E332E322E322E33> \MusicPDFRef{13109} <6571756174696F6E2E332E322E322E34> \MusicPDFRef{13110}] /Limits [<6571756174696F6E2E332E322E322E31> <6571756174696F6E2E332E322E322E34>] >>}
\MusicPDFWriteObject{9588}{<< /Names [<6571756174696F6E2E332E322E322E35> \MusicPDFRef{13111} <6571756174696F6E2E332E322E322E36> \MusicPDFRef{13112} <6571756174696F6E2E332E322E322E37> \MusicPDFRef{13113} <6571756174696F6E2E332E322E322E38> \MusicPDFRef{13114} <6571756174696F6E2E332E322E322E39> \MusicPDFRef{13115} <6571756174696F6E2E332E332E312E31> \MusicPDFRef{13116}] /Limits [<6571756174696F6E2E332E322E322E35> <6571756174696F6E2E332E332E312E31>] >>}
\MusicPDFWriteObject{9589}{<< /Names [<6571756174696F6E2E332E332E312E3130> \MusicPDFRef{13117} <6571756174696F6E2E332E332E312E3131> \MusicPDFRef{13118} <6571756174696F6E2E332E332E312E3132> \MusicPDFRef{13119} <6571756174696F6E2E332E332E312E3133> \MusicPDFRef{13120} <6571756174696F6E2E332E332E312E3134> \MusicPDFRef{13121} <6571756174696F6E2E332E332E312E3135> \MusicPDFRef{13122}] /Limits [<6571756174696F6E2E332E332E312E3130> <6571756174696F6E2E332E332E312E3135>] >>}
\MusicPDFWriteObject{9590}{<< /Names [<6571756174696F6E2E332E332E312E3136> \MusicPDFRef{13123} <6571756174696F6E2E332E332E312E3137> \MusicPDFRef{13124} <6571756174696F6E2E332E332E312E3138> \MusicPDFRef{13125} <6571756174696F6E2E332E332E312E3139> \MusicPDFRef{13126} <6571756174696F6E2E332E332E312E32> \MusicPDFRef{13127} <6571756174696F6E2E332E332E312E3230> \MusicPDFRef{13128}] /Limits [<6571756174696F6E2E332E332E312E3136> <6571756174696F6E2E332E332E312E3230>] >>}
\MusicPDFWriteObject{9591}{<< /Names [<6571756174696F6E2E332E332E312E3231> \MusicPDFRef{13129} <6571756174696F6E2E332E332E312E3232> \MusicPDFRef{13130} <6571756174696F6E2E332E332E312E3233> \MusicPDFRef{13131} <6571756174696F6E2E332E332E312E3234> \MusicPDFRef{13132} <6571756174696F6E2E332E332E312E3235> \MusicPDFRef{13133} <6571756174696F6E2E332E332E312E3236> \MusicPDFRef{13134}] /Limits [<6571756174696F6E2E332E332E312E3231> <6571756174696F6E2E332E332E312E3236>] >>}
\MusicPDFWriteObject{9592}{<< /Names [<6571756174696F6E2E332E332E312E3237> \MusicPDFRef{13135} <6571756174696F6E2E332E332E312E3238> \MusicPDFRef{13136} <6571756174696F6E2E332E332E312E3239> \MusicPDFRef{13137} <6571756174696F6E2E332E332E312E33> \MusicPDFRef{13138} <6571756174696F6E2E332E332E312E3330> \MusicPDFRef{13139} <6571756174696F6E2E332E332E312E3331> \MusicPDFRef{13140}] /Limits [<6571756174696F6E2E332E332E312E3237> <6571756174696F6E2E332E332E312E3331>] >>}
\MusicPDFWriteObject{9593}{<< /Names [<6571756174696F6E2E332E332E312E3332> \MusicPDFRef{13141} <6571756174696F6E2E332E332E312E3333> \MusicPDFRef{13142} <6571756174696F6E2E332E332E312E3334> \MusicPDFRef{13143} <6571756174696F6E2E332E332E312E34> \MusicPDFRef{13144} <6571756174696F6E2E332E332E312E35> \MusicPDFRef{13145} <6571756174696F6E2E332E332E312E36> \MusicPDFRef{13146}] /Limits [<6571756174696F6E2E332E332E312E3332> <6571756174696F6E2E332E332E312E36>] >>}
\MusicPDFWriteObject{9594}{<< /Names [<6571756174696F6E2E332E332E312E37> \MusicPDFRef{13147} <6571756174696F6E2E332E332E312E38> \MusicPDFRef{13148} <6571756174696F6E2E332E332E312E39> \MusicPDFRef{13149} <6571756174696F6E2E332E332E322E31> \MusicPDFRef{13150} <6571756174696F6E2E332E332E322E3130> \MusicPDFRef{13151} <6571756174696F6E2E332E332E322E3131> \MusicPDFRef{13152}] /Limits [<6571756174696F6E2E332E332E312E37> <6571756174696F6E2E332E332E322E3131>] >>}
\MusicPDFWriteObject{9595}{<< /Names [<6571756174696F6E2E332E332E322E3132> \MusicPDFRef{13153} <6571756174696F6E2E332E332E322E3133> \MusicPDFRef{13154} <6571756174696F6E2E332E332E322E3134> \MusicPDFRef{13155} <6571756174696F6E2E332E332E322E3135> \MusicPDFRef{13156} <6571756174696F6E2E332E332E322E3136> \MusicPDFRef{13157} <6571756174696F6E2E332E332E322E3137> \MusicPDFRef{13158}] /Limits [<6571756174696F6E2E332E332E322E3132> <6571756174696F6E2E332E332E322E3137>] >>}
\MusicPDFWriteObject{9596}{<< /Names [<6571756174696F6E2E332E332E322E3138> \MusicPDFRef{13159} <6571756174696F6E2E332E332E322E3139> \MusicPDFRef{13160} <6571756174696F6E2E332E332E322E32> \MusicPDFRef{13161} <6571756174696F6E2E332E332E322E3230> \MusicPDFRef{13162} <6571756174696F6E2E332E332E322E3231> \MusicPDFRef{13163} <6571756174696F6E2E332E332E322E3232> \MusicPDFRef{13164}] /Limits [<6571756174696F6E2E332E332E322E3138> <6571756174696F6E2E332E332E322E3232>] >>}
\MusicPDFWriteObject{9597}{<< /Names [<6571756174696F6E2E332E332E322E3233> \MusicPDFRef{13165} <6571756174696F6E2E332E332E322E3234> \MusicPDFRef{13166} <6571756174696F6E2E332E332E322E3235> \MusicPDFRef{13167} <6571756174696F6E2E332E332E322E33> \MusicPDFRef{13168} <6571756174696F6E2E332E332E322E34> \MusicPDFRef{13169} <6571756174696F6E2E332E332E322E35> \MusicPDFRef{13170}] /Limits [<6571756174696F6E2E332E332E322E3233> <6571756174696F6E2E332E332E322E35>] >>}
\MusicPDFWriteObject{9598}{<< /Names [<6571756174696F6E2E332E332E322E36> \MusicPDFRef{13171} <6571756174696F6E2E332E332E322E37> \MusicPDFRef{13172} <6571756174696F6E2E332E332E322E38> \MusicPDFRef{13173} <6571756174696F6E2E332E332E322E39> \MusicPDFRef{13174} <6571756174696F6E2E332E342E312E31> \MusicPDFRef{13175} <6571756174696F6E2E332E342E312E32> \MusicPDFRef{13176}] /Limits [<6571756174696F6E2E332E332E322E36> <6571756174696F6E2E332E342E312E32>] >>}
\MusicPDFWriteObject{9599}{<< /Names [<6571756174696F6E2E332E342E312E33> \MusicPDFRef{13177} <6571756174696F6E2E332E342E312E34> \MusicPDFRef{13178} <6571756174696F6E2E332E342E312E35> \MusicPDFRef{13179} <6571756174696F6E2E332E342E312E36> \MusicPDFRef{13180} <6571756174696F6E2E332E342E312E37> \MusicPDFRef{13181} <6571756174696F6E2E332E342E322E31> \MusicPDFRef{13182}] /Limits [<6571756174696F6E2E332E342E312E33> <6571756174696F6E2E332E342E322E31>] >>}
\MusicPDFWriteObject{9600}{<< /Names [<6571756174696F6E2E332E342E322E3130> \MusicPDFRef{13183} <6571756174696F6E2E332E342E322E3131> \MusicPDFRef{13184} <6571756174696F6E2E332E342E322E3132> \MusicPDFRef{13185} <6571756174696F6E2E332E342E322E3133> \MusicPDFRef{13186} <6571756174696F6E2E332E342E322E3134> \MusicPDFRef{13187} <6571756174696F6E2E332E342E322E3135> \MusicPDFRef{13188}] /Limits [<6571756174696F6E2E332E342E322E3130> <6571756174696F6E2E332E342E322E3135>] >>}
\MusicPDFWriteObject{9601}{<< /Names [<6571756174696F6E2E332E342E322E3136> \MusicPDFRef{13189} <6571756174696F6E2E332E342E322E3137> \MusicPDFRef{13190} <6571756174696F6E2E332E342E322E3138> \MusicPDFRef{13191} <6571756174696F6E2E332E342E322E32> \MusicPDFRef{13192} <6571756174696F6E2E332E342E322E33> \MusicPDFRef{13193} <6571756174696F6E2E332E342E322E34> \MusicPDFRef{13194}] /Limits [<6571756174696F6E2E332E342E322E3136> <6571756174696F6E2E332E342E322E34>] >>}
\MusicPDFWriteObject{9602}{<< /Names [<6571756174696F6E2E332E342E322E35> \MusicPDFRef{13195} <6571756174696F6E2E332E342E322E36> \MusicPDFRef{13196} <6571756174696F6E2E332E342E322E37> \MusicPDFRef{13197} <6571756174696F6E2E332E342E322E38> \MusicPDFRef{13198} <6571756174696F6E2E332E342E322E39> \MusicPDFRef{13199} <6571756174696F6E2E332E342E332E31> \MusicPDFRef{13200}] /Limits [<6571756174696F6E2E332E342E322E35> <6571756174696F6E2E332E342E332E31>] >>}
\MusicPDFWriteObject{9603}{<< /Names [<6571756174696F6E2E332E342E332E3130> \MusicPDFRef{13201} <6571756174696F6E2E332E342E332E3131> \MusicPDFRef{13202} <6571756174696F6E2E332E342E332E3132> \MusicPDFRef{13203} <6571756174696F6E2E332E342E332E3133> \MusicPDFRef{13204} <6571756174696F6E2E332E342E332E3134> \MusicPDFRef{13205} <6571756174696F6E2E332E342E332E3135> \MusicPDFRef{13206}] /Limits [<6571756174696F6E2E332E342E332E3130> <6571756174696F6E2E332E342E332E3135>] >>}
\MusicPDFWriteObject{9604}{<< /Names [<6571756174696F6E2E332E342E332E32> \MusicPDFRef{13207} <6571756174696F6E2E332E342E332E33> \MusicPDFRef{13208} <6571756174696F6E2E332E342E332E34> \MusicPDFRef{13209} <6571756174696F6E2E332E342E332E35> \MusicPDFRef{13210} <6571756174696F6E2E332E342E332E36> \MusicPDFRef{13211} <6571756174696F6E2E332E342E332E37> \MusicPDFRef{13212}] /Limits [<6571756174696F6E2E332E342E332E32> <6571756174696F6E2E332E342E332E37>] >>}
\MusicPDFWriteObject{9605}{<< /Names [<6571756174696F6E2E332E342E332E38> \MusicPDFRef{13213} <6571756174696F6E2E332E342E332E39> \MusicPDFRef{13214} <6571756174696F6E2E342E312E312E31> \MusicPDFRef{13215} <6571756174696F6E2E342E312E312E32> \MusicPDFRef{13216} <6571756174696F6E2E342E312E322E31> \MusicPDFRef{13217} <6571756174696F6E2E342E312E322E32> \MusicPDFRef{13218}] /Limits [<6571756174696F6E2E332E342E332E38> <6571756174696F6E2E342E312E322E32>] >>}
\MusicPDFWriteObject{9606}{<< /Names [<6571756174696F6E2E342E312E322E33> \MusicPDFRef{13219} <6571756174696F6E2E342E312E322E34> \MusicPDFRef{13220} <6571756174696F6E2E342E312E322E35> \MusicPDFRef{13221} <6571756174696F6E2E342E322E312E31> \MusicPDFRef{13222} <6571756174696F6E2E342E322E312E3130> \MusicPDFRef{13223} <6571756174696F6E2E342E322E312E3131> \MusicPDFRef{13224}] /Limits [<6571756174696F6E2E342E312E322E33> <6571756174696F6E2E342E322E312E3131>] >>}
\MusicPDFWriteObject{9607}{<< /Names [<6571756174696F6E2E342E322E312E3132> \MusicPDFRef{13225} <6571756174696F6E2E342E322E312E3133> \MusicPDFRef{13226} <6571756174696F6E2E342E322E312E3134> \MusicPDFRef{13227} <6571756174696F6E2E342E322E312E3135> \MusicPDFRef{13228} <6571756174696F6E2E342E322E312E3136> \MusicPDFRef{13229} <6571756174696F6E2E342E322E312E3137> \MusicPDFRef{13230}] /Limits [<6571756174696F6E2E342E322E312E3132> <6571756174696F6E2E342E322E312E3137>] >>}
\MusicPDFWriteObject{9608}{<< /Names [<6571756174696F6E2E342E322E312E3138> \MusicPDFRef{13231} <6571756174696F6E2E342E322E312E3139> \MusicPDFRef{13232} <6571756174696F6E2E342E322E312E32> \MusicPDFRef{13233} <6571756174696F6E2E342E322E312E33> \MusicPDFRef{13234} <6571756174696F6E2E342E322E312E34> \MusicPDFRef{13235} <6571756174696F6E2E342E322E312E35> \MusicPDFRef{13236}] /Limits [<6571756174696F6E2E342E322E312E3138> <6571756174696F6E2E342E322E312E35>] >>}
\MusicPDFWriteObject{9609}{<< /Names [<6571756174696F6E2E342E322E312E36> \MusicPDFRef{13237} <6571756174696F6E2E342E322E312E37> \MusicPDFRef{13238} <6571756174696F6E2E342E322E312E38> \MusicPDFRef{13239} <6571756174696F6E2E342E322E312E39> \MusicPDFRef{13240} <6571756174696F6E2E342E322E322E31> \MusicPDFRef{13241} <6571756174696F6E2E342E322E322E32> \MusicPDFRef{13242}] /Limits [<6571756174696F6E2E342E322E312E36> <6571756174696F6E2E342E322E322E32>] >>}
\MusicPDFWriteObject{9610}{<< /Names [<6571756174696F6E2E342E322E322E33> \MusicPDFRef{13243} <6571756174696F6E2E342E322E322E34> \MusicPDFRef{13244} <6571756174696F6E2E342E322E322E35> \MusicPDFRef{13245} <6571756174696F6E2E342E322E322E36> \MusicPDFRef{13246} <6571756174696F6E2E342E322E322E37> \MusicPDFRef{13247} <6571756174696F6E2E342E332E312E31> \MusicPDFRef{13248}] /Limits [<6571756174696F6E2E342E322E322E33> <6571756174696F6E2E342E332E312E31>] >>}
\MusicPDFWriteObject{9611}{<< /Names [<6571756174696F6E2E342E332E312E3130> \MusicPDFRef{13249} <6571756174696F6E2E342E332E312E3131> \MusicPDFRef{13250} <6571756174696F6E2E342E332E312E32> \MusicPDFRef{13251} <6571756174696F6E2E342E332E312E33> \MusicPDFRef{13252} <6571756174696F6E2E342E332E312E34> \MusicPDFRef{13253} <6571756174696F6E2E342E332E312E35> \MusicPDFRef{13254}] /Limits [<6571756174696F6E2E342E332E312E3130> <6571756174696F6E2E342E332E312E35>] >>}
\MusicPDFWriteObject{9612}{<< /Names [<6571756174696F6E2E342E332E312E36> \MusicPDFRef{13255} <6571756174696F6E2E342E332E312E37> \MusicPDFRef{13256} <6571756174696F6E2E342E332E312E38> \MusicPDFRef{13257} <6571756174696F6E2E342E332E312E39> \MusicPDFRef{13258} <6571756174696F6E2E342E332E322E31> \MusicPDFRef{13259} <6571756174696F6E2E342E332E322E32> \MusicPDFRef{13260}] /Limits [<6571756174696F6E2E342E332E312E36> <6571756174696F6E2E342E332E322E32>] >>}
\MusicPDFWriteObject{9613}{<< /Names [<6571756174696F6E2E342E332E322E33> \MusicPDFRef{13261} <6571756174696F6E2E342E332E322E34> \MusicPDFRef{13262} <6571756174696F6E2E342E332E322E35> \MusicPDFRef{13263} <6571756174696F6E2E342E332E322E36> \MusicPDFRef{13264} <6571756174696F6E2E342E332E322E37> \MusicPDFRef{13265} <6571756174696F6E2E342E332E322E38> \MusicPDFRef{13266}] /Limits [<6571756174696F6E2E342E332E322E33> <6571756174696F6E2E342E332E322E38>] >>}
\MusicPDFWriteObject{9614}{<< /Names [<6571756174696F6E2E342E332E332E31> \MusicPDFRef{13267} <6571756174696F6E2E342E332E332E3130> \MusicPDFRef{13268} <6571756174696F6E2E342E332E332E3131> \MusicPDFRef{13269} <6571756174696F6E2E342E332E332E3132> \MusicPDFRef{13270} <6571756174696F6E2E342E332E332E3133> \MusicPDFRef{13271} <6571756174696F6E2E342E332E332E3134> \MusicPDFRef{13272}] /Limits [<6571756174696F6E2E342E332E332E31> <6571756174696F6E2E342E332E332E3134>] >>}
\MusicPDFWriteObject{9615}{<< /Names [<6571756174696F6E2E342E332E332E3135> \MusicPDFRef{13273} <6571756174696F6E2E342E332E332E3136> \MusicPDFRef{13274} <6571756174696F6E2E342E332E332E3137> \MusicPDFRef{13275} <6571756174696F6E2E342E332E332E32> \MusicPDFRef{13276} <6571756174696F6E2E342E332E332E33> \MusicPDFRef{13277} <6571756174696F6E2E342E332E332E34> \MusicPDFRef{13278}] /Limits [<6571756174696F6E2E342E332E332E3135> <6571756174696F6E2E342E332E332E34>] >>}
\MusicPDFWriteObject{9616}{<< /Names [<6571756174696F6E2E342E332E332E35> \MusicPDFRef{13279} <6571756174696F6E2E342E332E332E36> \MusicPDFRef{13280} <6571756174696F6E2E342E332E332E37> \MusicPDFRef{13281} <6571756174696F6E2E342E332E332E38> \MusicPDFRef{13282} <6571756174696F6E2E342E332E332E39> \MusicPDFRef{13283} <6571756174696F6E2E342E342E312E31> \MusicPDFRef{13284}] /Limits [<6571756174696F6E2E342E332E332E35> <6571756174696F6E2E342E342E312E31>] >>}
\MusicPDFWriteObject{9617}{<< /Names [<6571756174696F6E2E342E342E312E32> \MusicPDFRef{13285} <6571756174696F6E2E342E342E312E33> \MusicPDFRef{13286} <6571756174696F6E2E342E342E312E34> \MusicPDFRef{13287} <6571756174696F6E2E342E342E322E31> \MusicPDFRef{13288} <6571756174696F6E2E342E342E322E3130> \MusicPDFRef{13289} <6571756174696F6E2E342E342E322E3131> \MusicPDFRef{13290}] /Limits [<6571756174696F6E2E342E342E312E32> <6571756174696F6E2E342E342E322E3131>] >>}
\MusicPDFWriteObject{9618}{<< /Names [<6571756174696F6E2E342E342E322E3132> \MusicPDFRef{13291} <6571756174696F6E2E342E342E322E32> \MusicPDFRef{13292} <6571756174696F6E2E342E342E322E33> \MusicPDFRef{13293} <6571756174696F6E2E342E342E322E34> \MusicPDFRef{13294} <6571756174696F6E2E342E342E322E35> \MusicPDFRef{13295} <6571756174696F6E2E342E342E322E36> \MusicPDFRef{13296}] /Limits [<6571756174696F6E2E342E342E322E3132> <6571756174696F6E2E342E342E322E36>] >>}
\MusicPDFWriteObject{9619}{<< /Names [<6571756174696F6E2E342E342E322E37> \MusicPDFRef{13297} <6571756174696F6E2E342E342E322E38> \MusicPDFRef{13298} <6571756174696F6E2E342E342E322E39> \MusicPDFRef{13299} <6571756174696F6E2E342E352E312E31> \MusicPDFRef{13300} <6571756174696F6E2E342E352E312E3130> \MusicPDFRef{13301} <6571756174696F6E2E342E352E312E3131> \MusicPDFRef{13302}] /Limits [<6571756174696F6E2E342E342E322E37> <6571756174696F6E2E342E352E312E3131>] >>}
\MusicPDFWriteObject{9620}{<< /Names [<6571756174696F6E2E342E352E312E3132> \MusicPDFRef{13303} <6571756174696F6E2E342E352E312E32> \MusicPDFRef{13304} <6571756174696F6E2E342E352E312E33> \MusicPDFRef{13305} <6571756174696F6E2E342E352E312E34> \MusicPDFRef{13306} <6571756174696F6E2E342E352E312E35> \MusicPDFRef{13307} <6571756174696F6E2E342E352E312E36> \MusicPDFRef{13308}] /Limits [<6571756174696F6E2E342E352E312E3132> <6571756174696F6E2E342E352E312E36>] >>}
\MusicPDFWriteObject{9621}{<< /Names [<6571756174696F6E2E342E352E312E37> \MusicPDFRef{13309} <6571756174696F6E2E342E352E312E38> \MusicPDFRef{13310} <6571756174696F6E2E342E352E312E39> \MusicPDFRef{13311} <6571756174696F6E2E342E352E322E31> \MusicPDFRef{13312} <6571756174696F6E2E342E352E322E32> \MusicPDFRef{13313} <6571756174696F6E2E342E352E332E31> \MusicPDFRef{13314}] /Limits [<6571756174696F6E2E342E352E312E37> <6571756174696F6E2E342E352E332E31>] >>}
\MusicPDFWriteObject{9622}{<< /Names [<6571756174696F6E2E342E352E332E32> \MusicPDFRef{13315} <6571756174696F6E2E342E352E332E33> \MusicPDFRef{13316} <6571756174696F6E2E342E352E332E34> \MusicPDFRef{13317} <6571756174696F6E2E342E352E332E35> \MusicPDFRef{13318} <6571756174696F6E2E342E352E332E36> \MusicPDFRef{13319} <6571756174696F6E2E342E362E312E31> \MusicPDFRef{13320}] /Limits [<6571756174696F6E2E342E352E332E32> <6571756174696F6E2E342E362E312E31>] >>}
\MusicPDFWriteObject{9623}{<< /Names [<6571756174696F6E2E342E362E312E3130> \MusicPDFRef{13321} <6571756174696F6E2E342E362E312E3131> \MusicPDFRef{13322} <6571756174696F6E2E342E362E312E3132> \MusicPDFRef{13323} <6571756174696F6E2E342E362E312E3133> \MusicPDFRef{13324} <6571756174696F6E2E342E362E312E3134> \MusicPDFRef{13325} <6571756174696F6E2E342E362E312E3135> \MusicPDFRef{13326}] /Limits [<6571756174696F6E2E342E362E312E3130> <6571756174696F6E2E342E362E312E3135>] >>}
\MusicPDFWriteObject{9624}{<< /Names [<6571756174696F6E2E342E362E312E3136> \MusicPDFRef{13327} <6571756174696F6E2E342E362E312E3137> \MusicPDFRef{13328} <6571756174696F6E2E342E362E312E3138> \MusicPDFRef{13329} <6571756174696F6E2E342E362E312E3139> \MusicPDFRef{13330} <6571756174696F6E2E342E362E312E32> \MusicPDFRef{13331} <6571756174696F6E2E342E362E312E3230> \MusicPDFRef{13332}] /Limits [<6571756174696F6E2E342E362E312E3136> <6571756174696F6E2E342E362E312E3230>] >>}
\MusicPDFWriteObject{9625}{<< /Names [<6571756174696F6E2E342E362E312E3231> \MusicPDFRef{13333} <6571756174696F6E2E342E362E312E3232> \MusicPDFRef{13334} <6571756174696F6E2E342E362E312E3233> \MusicPDFRef{13335} <6571756174696F6E2E342E362E312E3234> \MusicPDFRef{13336} <6571756174696F6E2E342E362E312E3235> \MusicPDFRef{13337} <6571756174696F6E2E342E362E312E3236> \MusicPDFRef{13338}] /Limits [<6571756174696F6E2E342E362E312E3231> <6571756174696F6E2E342E362E312E3236>] >>}
\MusicPDFWriteObject{9626}{<< /Names [<6571756174696F6E2E342E362E312E3237> \MusicPDFRef{13339} <6571756174696F6E2E342E362E312E3238> \MusicPDFRef{13340} <6571756174696F6E2E342E362E312E3239> \MusicPDFRef{13341} <6571756174696F6E2E342E362E312E33> \MusicPDFRef{13342} <6571756174696F6E2E342E362E312E3330> \MusicPDFRef{13343} <6571756174696F6E2E342E362E312E34> \MusicPDFRef{13344}] /Limits [<6571756174696F6E2E342E362E312E3237> <6571756174696F6E2E342E362E312E34>] >>}
\MusicPDFWriteObject{9627}{<< /Names [<6571756174696F6E2E342E362E312E35> \MusicPDFRef{13345} <6571756174696F6E2E342E362E312E36> \MusicPDFRef{13346} <6571756174696F6E2E342E362E312E37> \MusicPDFRef{13347} <6571756174696F6E2E342E362E312E38> \MusicPDFRef{13348} <6571756174696F6E2E342E362E312E39> \MusicPDFRef{13349} <6571756174696F6E2E342E362E322E31> \MusicPDFRef{13350}] /Limits [<6571756174696F6E2E342E362E312E35> <6571756174696F6E2E342E362E322E31>] >>}
\MusicPDFWriteObject{9628}{<< /Names [<6571756174696F6E2E342E362E322E32> \MusicPDFRef{13351} <6571756174696F6E2E342E362E322E33> \MusicPDFRef{13352} <6571756174696F6E2E342E362E322E34> \MusicPDFRef{13353} <6571756174696F6E2E342E362E322E35> \MusicPDFRef{13354} <6571756174696F6E2E342E362E322E36> \MusicPDFRef{13355} <6571756174696F6E2E342E362E332E31> \MusicPDFRef{13356}] /Limits [<6571756174696F6E2E342E362E322E32> <6571756174696F6E2E342E362E332E31>] >>}
\MusicPDFWriteObject{9629}{<< /Names [<6571756174696F6E2E342E362E332E32> \MusicPDFRef{13357} <6571756174696F6E2E352E312E312E31> \MusicPDFRef{13358} <6571756174696F6E2E352E312E312E32> \MusicPDFRef{13359} <6571756174696F6E2E352E312E312E33> \MusicPDFRef{13360} <6571756174696F6E2E352E312E312E34> \MusicPDFRef{13361} <6571756174696F6E2E352E312E312E35> \MusicPDFRef{13362}] /Limits [<6571756174696F6E2E342E362E332E32> <6571756174696F6E2E352E312E312E35>] >>}
\MusicPDFWriteObject{9630}{<< /Names [<6571756174696F6E2E352E312E322E31> \MusicPDFRef{13363} <6571756174696F6E2E352E312E322E32> \MusicPDFRef{13364} <6571756174696F6E2E352E312E322E33> \MusicPDFRef{13365} <6571756174696F6E2E352E312E322E34> \MusicPDFRef{13366} <6571756174696F6E2E352E312E332E31> \MusicPDFRef{13367} <6571756174696F6E2E352E312E332E32> \MusicPDFRef{13368}] /Limits [<6571756174696F6E2E352E312E322E31> <6571756174696F6E2E352E312E332E32>] >>}
\MusicPDFWriteObject{9631}{<< /Names [<6571756174696F6E2E352E312E332E33> \MusicPDFRef{13369} <6571756174696F6E2E352E312E332E34> \MusicPDFRef{13370} <6571756174696F6E2E352E312E332E35> \MusicPDFRef{13371} <6571756174696F6E2E352E312E332E36> \MusicPDFRef{13372} <6571756174696F6E2E352E312E332E37> \MusicPDFRef{13373} <6571756174696F6E2E352E322E312E31> \MusicPDFRef{13374}] /Limits [<6571756174696F6E2E352E312E332E33> <6571756174696F6E2E352E322E312E31>] >>}
\MusicPDFWriteObject{9632}{<< /Names [<6571756174696F6E2E352E322E312E32> \MusicPDFRef{13375} <6571756174696F6E2E352E322E312E33> \MusicPDFRef{13376} <6571756174696F6E2E352E322E312E34> \MusicPDFRef{13377} <6571756174696F6E2E352E322E322E31> \MusicPDFRef{13378} <6571756174696F6E2E352E322E322E32> \MusicPDFRef{13379} <6571756174696F6E2E352E322E322E33> \MusicPDFRef{13380}] /Limits [<6571756174696F6E2E352E322E312E32> <6571756174696F6E2E352E322E322E33>] >>}
\MusicPDFWriteObject{9633}{<< /Names [<6571756174696F6E2E352E322E322E34> \MusicPDFRef{13381} <6571756174696F6E2E352E322E332E31> \MusicPDFRef{13382} <6571756174696F6E2E352E322E332E3130> \MusicPDFRef{13383} <6571756174696F6E2E352E322E332E3131> \MusicPDFRef{13384} <6571756174696F6E2E352E322E332E3132> \MusicPDFRef{13385} <6571756174696F6E2E352E322E332E3133> \MusicPDFRef{13386}] /Limits [<6571756174696F6E2E352E322E322E34> <6571756174696F6E2E352E322E332E3133>] >>}
\MusicPDFWriteObject{9634}{<< /Names [<6571756174696F6E2E352E322E332E3134> \MusicPDFRef{13387} <6571756174696F6E2E352E322E332E3135> \MusicPDFRef{13388} <6571756174696F6E2E352E322E332E32> \MusicPDFRef{13389} <6571756174696F6E2E352E322E332E33> \MusicPDFRef{13390} <6571756174696F6E2E352E322E332E34> \MusicPDFRef{13391} <6571756174696F6E2E352E322E332E35> \MusicPDFRef{13392}] /Limits [<6571756174696F6E2E352E322E332E3134> <6571756174696F6E2E352E322E332E35>] >>}
\MusicPDFWriteObject{9635}{<< /Names [<6571756174696F6E2E352E322E332E36> \MusicPDFRef{13393} <6571756174696F6E2E352E322E332E37> \MusicPDFRef{13394} <6571756174696F6E2E352E322E332E38> \MusicPDFRef{13395} <6571756174696F6E2E352E322E332E39> \MusicPDFRef{13396} <6571756174696F6E2E352E322E342E31> \MusicPDFRef{13397} <6571756174696F6E2E352E322E342E32> \MusicPDFRef{13398}] /Limits [<6571756174696F6E2E352E322E332E36> <6571756174696F6E2E352E322E342E32>] >>}
\MusicPDFWriteObject{9636}{<< /Names [<6571756174696F6E2E352E322E342E33> \MusicPDFRef{13399} <6571756174696F6E2E352E322E342E34> \MusicPDFRef{13400} <6571756174696F6E2E352E332E312E31> \MusicPDFRef{13401} <6571756174696F6E2E352E332E312E3130> \MusicPDFRef{13402} <6571756174696F6E2E352E332E312E32> \MusicPDFRef{13403} <6571756174696F6E2E352E332E312E33> \MusicPDFRef{13404}] /Limits [<6571756174696F6E2E352E322E342E33> <6571756174696F6E2E352E332E312E33>] >>}
\MusicPDFWriteObject{9637}{<< /Names [<6571756174696F6E2E352E332E312E34> \MusicPDFRef{13405} <6571756174696F6E2E352E332E312E35> \MusicPDFRef{13406} <6571756174696F6E2E352E332E312E36> \MusicPDFRef{13407} <6571756174696F6E2E352E332E312E37> \MusicPDFRef{13408} <6571756174696F6E2E352E332E312E38> \MusicPDFRef{13409} <6571756174696F6E2E352E332E312E39> \MusicPDFRef{13410}] /Limits [<6571756174696F6E2E352E332E312E34> <6571756174696F6E2E352E332E312E39>] >>}
\MusicPDFWriteObject{9638}{<< /Names [<6571756174696F6E2E352E332E322E31> \MusicPDFRef{13411} <6571756174696F6E2E352E332E322E3130> \MusicPDFRef{13412} <6571756174696F6E2E352E332E322E3131> \MusicPDFRef{13413} <6571756174696F6E2E352E332E322E3132> \MusicPDFRef{13414} <6571756174696F6E2E352E332E322E3133> \MusicPDFRef{13415} <6571756174696F6E2E352E332E322E3134> \MusicPDFRef{13416}] /Limits [<6571756174696F6E2E352E332E322E31> <6571756174696F6E2E352E332E322E3134>] >>}
\MusicPDFWriteObject{9639}{<< /Names [<6571756174696F6E2E352E332E322E3135> \MusicPDFRef{13417} <6571756174696F6E2E352E332E322E3136> \MusicPDFRef{13418} <6571756174696F6E2E352E332E322E3137> \MusicPDFRef{13419} <6571756174696F6E2E352E332E322E3138> \MusicPDFRef{13420} <6571756174696F6E2E352E332E322E3139> \MusicPDFRef{13421} <6571756174696F6E2E352E332E322E32> \MusicPDFRef{13422}] /Limits [<6571756174696F6E2E352E332E322E3135> <6571756174696F6E2E352E332E322E32>] >>}
\MusicPDFWriteObject{9640}{<< /Names [<6571756174696F6E2E352E332E322E33> \MusicPDFRef{13423} <6571756174696F6E2E352E332E322E34> \MusicPDFRef{13424} <6571756174696F6E2E352E332E322E35> \MusicPDFRef{13425} <6571756174696F6E2E352E332E322E36> \MusicPDFRef{13426} <6571756174696F6E2E352E332E322E37> \MusicPDFRef{13427} <6571756174696F6E2E352E332E322E38> \MusicPDFRef{13428}] /Limits [<6571756174696F6E2E352E332E322E33> <6571756174696F6E2E352E332E322E38>] >>}
\MusicPDFWriteObject{9641}{<< /Names [<6571756174696F6E2E352E332E322E39> \MusicPDFRef{13429} <6571756174696F6E2E352E332E332E31> \MusicPDFRef{13430} <6571756174696F6E2E352E332E332E3130> \MusicPDFRef{13431} <6571756174696F6E2E352E332E332E3131> \MusicPDFRef{13432} <6571756174696F6E2E352E332E332E3132> \MusicPDFRef{13433} <6571756174696F6E2E352E332E332E3133> \MusicPDFRef{13434}] /Limits [<6571756174696F6E2E352E332E322E39> <6571756174696F6E2E352E332E332E3133>] >>}
\MusicPDFWriteObject{9642}{<< /Names [<6571756174696F6E2E352E332E332E3134> \MusicPDFRef{13435} <6571756174696F6E2E352E332E332E3135> \MusicPDFRef{13436} <6571756174696F6E2E352E332E332E3136> \MusicPDFRef{13437} <6571756174696F6E2E352E332E332E3137> \MusicPDFRef{13438} <6571756174696F6E2E352E332E332E3138> \MusicPDFRef{13439} <6571756174696F6E2E352E332E332E3139> \MusicPDFRef{13440}] /Limits [<6571756174696F6E2E352E332E332E3134> <6571756174696F6E2E352E332E332E3139>] >>}
\MusicPDFWriteObject{9643}{<< /Names [<6571756174696F6E2E352E332E332E32> \MusicPDFRef{13441} <6571756174696F6E2E352E332E332E3230> \MusicPDFRef{13442} <6571756174696F6E2E352E332E332E3231> \MusicPDFRef{13443} <6571756174696F6E2E352E332E332E3232> \MusicPDFRef{13444} <6571756174696F6E2E352E332E332E3233> \MusicPDFRef{13445} <6571756174696F6E2E352E332E332E3234> \MusicPDFRef{13446}] /Limits [<6571756174696F6E2E352E332E332E32> <6571756174696F6E2E352E332E332E3234>] >>}
\MusicPDFWriteObject{9644}{<< /Names [<6571756174696F6E2E352E332E332E3235> \MusicPDFRef{13447} <6571756174696F6E2E352E332E332E3236> \MusicPDFRef{13448} <6571756174696F6E2E352E332E332E3237> \MusicPDFRef{13449} <6571756174696F6E2E352E332E332E3238> \MusicPDFRef{13450} <6571756174696F6E2E352E332E332E3239> \MusicPDFRef{13451} <6571756174696F6E2E352E332E332E33> \MusicPDFRef{13452}] /Limits [<6571756174696F6E2E352E332E332E3235> <6571756174696F6E2E352E332E332E33>] >>}
\MusicPDFWriteObject{9645}{<< /Names [<6571756174696F6E2E352E332E332E3330> \MusicPDFRef{13453} <6571756174696F6E2E352E332E332E3331> \MusicPDFRef{13454} <6571756174696F6E2E352E332E332E3332> \MusicPDFRef{13455} <6571756174696F6E2E352E332E332E3333> \MusicPDFRef{13456} <6571756174696F6E2E352E332E332E3334> \MusicPDFRef{13457} <6571756174696F6E2E352E332E332E3335> \MusicPDFRef{13458}] /Limits [<6571756174696F6E2E352E332E332E3330> <6571756174696F6E2E352E332E332E3335>] >>}
\MusicPDFWriteObject{9646}{<< /Names [<6571756174696F6E2E352E332E332E3336> \MusicPDFRef{13459} <6571756174696F6E2E352E332E332E3337> \MusicPDFRef{13460} <6571756174696F6E2E352E332E332E3338> \MusicPDFRef{13461} <6571756174696F6E2E352E332E332E3339> \MusicPDFRef{13462} <6571756174696F6E2E352E332E332E34> \MusicPDFRef{13463} <6571756174696F6E2E352E332E332E3430> \MusicPDFRef{13464}] /Limits [<6571756174696F6E2E352E332E332E3336> <6571756174696F6E2E352E332E332E3430>] >>}
\MusicPDFWriteObject{9647}{<< /Names [<6571756174696F6E2E352E332E332E35> \MusicPDFRef{13465} <6571756174696F6E2E352E332E332E36> \MusicPDFRef{13466} <6571756174696F6E2E352E332E332E37> \MusicPDFRef{13467} <6571756174696F6E2E352E332E332E38> \MusicPDFRef{13468} <6571756174696F6E2E352E332E332E39> \MusicPDFRef{13469} <6571756174696F6E2E352E332E342E31> \MusicPDFRef{13470}] /Limits [<6571756174696F6E2E352E332E332E35> <6571756174696F6E2E352E332E342E31>] >>}
\MusicPDFWriteObject{9648}{<< /Names [<6571756174696F6E2E352E332E342E3130> \MusicPDFRef{13471} <6571756174696F6E2E352E332E342E3131> \MusicPDFRef{13472} <6571756174696F6E2E352E332E342E3132> \MusicPDFRef{13473} <6571756174696F6E2E352E332E342E3133> \MusicPDFRef{13474} <6571756174696F6E2E352E332E342E3134> \MusicPDFRef{13475} <6571756174696F6E2E352E332E342E3135> \MusicPDFRef{13476}] /Limits [<6571756174696F6E2E352E332E342E3130> <6571756174696F6E2E352E332E342E3135>] >>}
\MusicPDFWriteObject{9649}{<< /Names [<6571756174696F6E2E352E332E342E32> \MusicPDFRef{13477} <6571756174696F6E2E352E332E342E33> \MusicPDFRef{13478} <6571756174696F6E2E352E332E342E34> \MusicPDFRef{13479} <6571756174696F6E2E352E332E342E35> \MusicPDFRef{13480} <6571756174696F6E2E352E332E342E36> \MusicPDFRef{13481} <6571756174696F6E2E352E332E342E37> \MusicPDFRef{13482}] /Limits [<6571756174696F6E2E352E332E342E32> <6571756174696F6E2E352E332E342E37>] >>}
\MusicPDFWriteObject{9650}{<< /Names [<6571756174696F6E2E352E332E342E38> \MusicPDFRef{13483} <6571756174696F6E2E352E332E342E39> \MusicPDFRef{13484} <6571756174696F6E2E352E342E312E31> \MusicPDFRef{13485} <6571756174696F6E2E352E342E312E3130> \MusicPDFRef{13486} <6571756174696F6E2E352E342E312E3131> \MusicPDFRef{13487} <6571756174696F6E2E352E342E312E3132> \MusicPDFRef{13488}] /Limits [<6571756174696F6E2E352E332E342E38> <6571756174696F6E2E352E342E312E3132>] >>}
\MusicPDFWriteObject{9651}{<< /Names [<6571756174696F6E2E352E342E312E3133> \MusicPDFRef{13489} <6571756174696F6E2E352E342E312E3134> \MusicPDFRef{13490} <6571756174696F6E2E352E342E312E3135> \MusicPDFRef{13491} <6571756174696F6E2E352E342E312E3136> \MusicPDFRef{13492} <6571756174696F6E2E352E342E312E3137> \MusicPDFRef{13493} <6571756174696F6E2E352E342E312E3138> \MusicPDFRef{13494}] /Limits [<6571756174696F6E2E352E342E312E3133> <6571756174696F6E2E352E342E312E3138>] >>}
\MusicPDFWriteObject{9652}{<< /Names [<6571756174696F6E2E352E342E312E3139> \MusicPDFRef{13495} <6571756174696F6E2E352E342E312E32> \MusicPDFRef{13496} <6571756174696F6E2E352E342E312E3230> \MusicPDFRef{13497} <6571756174696F6E2E352E342E312E3231> \MusicPDFRef{13498} <6571756174696F6E2E352E342E312E3232> \MusicPDFRef{13499} <6571756174696F6E2E352E342E312E3233> \MusicPDFRef{13500}] /Limits [<6571756174696F6E2E352E342E312E3139> <6571756174696F6E2E352E342E312E3233>] >>}
\MusicPDFWriteObject{9653}{<< /Names [<6571756174696F6E2E352E342E312E33> \MusicPDFRef{13501} <6571756174696F6E2E352E342E312E34> \MusicPDFRef{13502} <6571756174696F6E2E352E342E312E35> \MusicPDFRef{13503} <6571756174696F6E2E352E342E312E36> \MusicPDFRef{13504} <6571756174696F6E2E352E342E312E37> \MusicPDFRef{13505} <6571756174696F6E2E352E342E312E38> \MusicPDFRef{13506}] /Limits [<6571756174696F6E2E352E342E312E33> <6571756174696F6E2E352E342E312E38>] >>}
\MusicPDFWriteObject{9654}{<< /Names [<6571756174696F6E2E352E342E312E39> \MusicPDFRef{13507} <6571756174696F6E2E352E342E322E31> \MusicPDFRef{13508} <6571756174696F6E2E352E342E322E32> \MusicPDFRef{13509} <6571756174696F6E2E352E342E322E33> \MusicPDFRef{13510} <6571756174696F6E2E352E342E322E34> \MusicPDFRef{13511} <6571756174696F6E2E352E342E322E35> \MusicPDFRef{13512}] /Limits [<6571756174696F6E2E352E342E312E39> <6571756174696F6E2E352E342E322E35>] >>}
\MusicPDFWriteObject{9655}{<< /Names [<6571756174696F6E2E352E342E322E36> \MusicPDFRef{13513} <6571756174696F6E2E352E342E322E37> \MusicPDFRef{13514} <6571756174696F6E2E352E352E312E31> \MusicPDFRef{13515} <6571756174696F6E2E352E352E312E3130> \MusicPDFRef{13516} <6571756174696F6E2E352E352E312E3131> \MusicPDFRef{13517} <6571756174696F6E2E352E352E312E3132> \MusicPDFRef{13518}] /Limits [<6571756174696F6E2E352E342E322E36> <6571756174696F6E2E352E352E312E3132>] >>}
\MusicPDFWriteObject{9656}{<< /Names [<6571756174696F6E2E352E352E312E3133> \MusicPDFRef{13519} <6571756174696F6E2E352E352E312E3134> \MusicPDFRef{13520} <6571756174696F6E2E352E352E312E3135> \MusicPDFRef{13521} <6571756174696F6E2E352E352E312E3136> \MusicPDFRef{13522} <6571756174696F6E2E352E352E312E3137> \MusicPDFRef{13523} <6571756174696F6E2E352E352E312E3138> \MusicPDFRef{13524}] /Limits [<6571756174696F6E2E352E352E312E3133> <6571756174696F6E2E352E352E312E3138>] >>}
\MusicPDFWriteObject{9657}{<< /Names [<6571756174696F6E2E352E352E312E32> \MusicPDFRef{13525} <6571756174696F6E2E352E352E312E33> \MusicPDFRef{13526} <6571756174696F6E2E352E352E312E34> \MusicPDFRef{13527} <6571756174696F6E2E352E352E312E35> \MusicPDFRef{13528} <6571756174696F6E2E352E352E312E36> \MusicPDFRef{13529} <6571756174696F6E2E352E352E312E37> \MusicPDFRef{13530}] /Limits [<6571756174696F6E2E352E352E312E32> <6571756174696F6E2E352E352E312E37>] >>}
\MusicPDFWriteObject{9658}{<< /Names [<6571756174696F6E2E352E352E312E38> \MusicPDFRef{13531} <6571756174696F6E2E352E352E312E39> \MusicPDFRef{13532} <6571756174696F6E2E352E352E322E31> \MusicPDFRef{13533} <6571756174696F6E2E352E352E322E32> \MusicPDFRef{13534} <6571756174696F6E2E352E352E322E33> \MusicPDFRef{13535} <6571756174696F6E2E352E352E332E31> \MusicPDFRef{13536}] /Limits [<6571756174696F6E2E352E352E312E38> <6571756174696F6E2E352E352E332E31>] >>}
\MusicPDFWriteObject{9659}{<< /Names [<6571756174696F6E2E352E352E332E32> \MusicPDFRef{13537} <6571756174696F6E2E352E352E332E33> \MusicPDFRef{13538} <6571756174696F6E2E352E352E332E34> \MusicPDFRef{13539} <6571756174696F6E2E352E352E332E35> \MusicPDFRef{13540} <6571756174696F6E2E352E352E332E36> \MusicPDFRef{13541} <6571756174696F6E2E362E312E312E31> \MusicPDFRef{13542}] /Limits [<6571756174696F6E2E352E352E332E32> <6571756174696F6E2E362E312E312E31>] >>}
\MusicPDFWriteObject{9660}{<< /Names [<6571756174696F6E2E362E312E312E32> \MusicPDFRef{13543} <6571756174696F6E2E362E312E322E31> \MusicPDFRef{13544} <6571756174696F6E2E362E312E322E3130> \MusicPDFRef{13545} <6571756174696F6E2E362E312E322E3131> \MusicPDFRef{13546} <6571756174696F6E2E362E312E322E3132> \MusicPDFRef{13547} <6571756174696F6E2E362E312E322E3133> \MusicPDFRef{13548}] /Limits [<6571756174696F6E2E362E312E312E32> <6571756174696F6E2E362E312E322E3133>] >>}
\MusicPDFWriteObject{9661}{<< /Names [<6571756174696F6E2E362E312E322E3134> \MusicPDFRef{13549} <6571756174696F6E2E362E312E322E3135> \MusicPDFRef{13550} <6571756174696F6E2E362E312E322E3136> \MusicPDFRef{13551} <6571756174696F6E2E362E312E322E3137> \MusicPDFRef{13552} <6571756174696F6E2E362E312E322E3138> \MusicPDFRef{13553} <6571756174696F6E2E362E312E322E3139> \MusicPDFRef{13554}] /Limits [<6571756174696F6E2E362E312E322E3134> <6571756174696F6E2E362E312E322E3139>] >>}
\MusicPDFWriteObject{9662}{<< /Names [<6571756174696F6E2E362E312E322E32> \MusicPDFRef{13555} <6571756174696F6E2E362E312E322E3230> \MusicPDFRef{13556} <6571756174696F6E2E362E312E322E3231> \MusicPDFRef{13557} <6571756174696F6E2E362E312E322E3232> \MusicPDFRef{13558} <6571756174696F6E2E362E312E322E33> \MusicPDFRef{13559} <6571756174696F6E2E362E312E322E34> \MusicPDFRef{13560}] /Limits [<6571756174696F6E2E362E312E322E32> <6571756174696F6E2E362E312E322E34>] >>}
\MusicPDFWriteObject{9663}{<< /Names [<6571756174696F6E2E362E312E322E35> \MusicPDFRef{13561} <6571756174696F6E2E362E312E322E36> \MusicPDFRef{13562} <6571756174696F6E2E362E312E322E37> \MusicPDFRef{13563} <6571756174696F6E2E362E312E322E38> \MusicPDFRef{13564} <6571756174696F6E2E362E312E322E39> \MusicPDFRef{13565} <6571756174696F6E2E362E322E312E31> \MusicPDFRef{13566}] /Limits [<6571756174696F6E2E362E312E322E35> <6571756174696F6E2E362E322E312E31>] >>}
\MusicPDFWriteObject{9664}{<< /Names [<6571756174696F6E2E362E322E312E32> \MusicPDFRef{13567} <6571756174696F6E2E362E322E322E31> \MusicPDFRef{13568} <6571756174696F6E2E362E322E322E3130> \MusicPDFRef{13569} <6571756174696F6E2E362E322E322E3131> \MusicPDFRef{13570} <6571756174696F6E2E362E322E322E3132> \MusicPDFRef{13571} <6571756174696F6E2E362E322E322E32> \MusicPDFRef{13572}] /Limits [<6571756174696F6E2E362E322E312E32> <6571756174696F6E2E362E322E322E32>] >>}
\MusicPDFWriteObject{9665}{<< /Names [<6571756174696F6E2E362E322E322E33> \MusicPDFRef{13573} <6571756174696F6E2E362E322E322E34> \MusicPDFRef{13574} <6571756174696F6E2E362E322E322E35> \MusicPDFRef{13575} <6571756174696F6E2E362E322E322E36> \MusicPDFRef{13576} <6571756174696F6E2E362E322E322E37> \MusicPDFRef{13577} <6571756174696F6E2E362E322E322E38> \MusicPDFRef{13578}] /Limits [<6571756174696F6E2E362E322E322E33> <6571756174696F6E2E362E322E322E38>] >>}
\MusicPDFWriteObject{9666}{<< /Names [<6571756174696F6E2E362E322E322E39> \MusicPDFRef{13579} <6571756174696F6E2E362E332E312E31> \MusicPDFRef{13580} <6571756174696F6E2E362E332E312E32> \MusicPDFRef{13581} <6571756174696F6E2E362E332E312E33> \MusicPDFRef{13582} <6571756174696F6E2E362E332E312E34> \MusicPDFRef{13583} <6571756174696F6E2E362E332E312E35> \MusicPDFRef{13584}] /Limits [<6571756174696F6E2E362E322E322E39> <6571756174696F6E2E362E332E312E35>] >>}
\MusicPDFWriteObject{9667}{<< /Names [<6571756174696F6E2E362E332E312E36> \MusicPDFRef{13585} <6571756174696F6E2E362E332E322E31> \MusicPDFRef{13586} <6571756174696F6E2E362E332E322E3130> \MusicPDFRef{13587} <6571756174696F6E2E362E332E322E3131> \MusicPDFRef{13588} <6571756174696F6E2E362E332E322E3132> \MusicPDFRef{13589} <6571756174696F6E2E362E332E322E3133> \MusicPDFRef{13590}] /Limits [<6571756174696F6E2E362E332E312E36> <6571756174696F6E2E362E332E322E3133>] >>}
\MusicPDFWriteObject{9668}{<< /Names [<6571756174696F6E2E362E332E322E3134> \MusicPDFRef{13591} <6571756174696F6E2E362E332E322E3135> \MusicPDFRef{13592} <6571756174696F6E2E362E332E322E3136> \MusicPDFRef{13593} <6571756174696F6E2E362E332E322E3137> \MusicPDFRef{13594} <6571756174696F6E2E362E332E322E3138> \MusicPDFRef{13595} <6571756174696F6E2E362E332E322E3139> \MusicPDFRef{13596}] /Limits [<6571756174696F6E2E362E332E322E3134> <6571756174696F6E2E362E332E322E3139>] >>}
\MusicPDFWriteObject{9669}{<< /Names [<6571756174696F6E2E362E332E322E32> \MusicPDFRef{13597} <6571756174696F6E2E362E332E322E3230> \MusicPDFRef{13598} <6571756174696F6E2E362E332E322E3231> \MusicPDFRef{13599} <6571756174696F6E2E362E332E322E3232> \MusicPDFRef{13600} <6571756174696F6E2E362E332E322E3233> \MusicPDFRef{13601} <6571756174696F6E2E362E332E322E3234> \MusicPDFRef{13602}] /Limits [<6571756174696F6E2E362E332E322E32> <6571756174696F6E2E362E332E322E3234>] >>}
\MusicPDFWriteObject{9670}{<< /Names [<6571756174696F6E2E362E332E322E3235> \MusicPDFRef{13603} <6571756174696F6E2E362E332E322E3236> \MusicPDFRef{13604} <6571756174696F6E2E362E332E322E33> \MusicPDFRef{13605} <6571756174696F6E2E362E332E322E34> \MusicPDFRef{13606} <6571756174696F6E2E362E332E322E35> \MusicPDFRef{13607} <6571756174696F6E2E362E332E322E36> \MusicPDFRef{13608}] /Limits [<6571756174696F6E2E362E332E322E3235> <6571756174696F6E2E362E332E322E36>] >>}
\MusicPDFWriteObject{9671}{<< /Names [<6571756174696F6E2E362E332E322E37> \MusicPDFRef{13609} <6571756174696F6E2E362E332E322E38> \MusicPDFRef{13610} <6571756174696F6E2E362E332E322E39> \MusicPDFRef{13611} <6571756174696F6E2E372E312E312E31> \MusicPDFRef{13612} <6571756174696F6E2E372E312E312E32> \MusicPDFRef{13613} <6571756174696F6E2E372E312E312E33> \MusicPDFRef{13614}] /Limits [<6571756174696F6E2E362E332E322E37> <6571756174696F6E2E372E312E312E33>] >>}
\MusicPDFWriteObject{9672}{<< /Names [<6571756174696F6E2E372E312E312E34> \MusicPDFRef{13615} <6571756174696F6E2E372E312E312E35> \MusicPDFRef{13616} <6571756174696F6E2E372E312E312E36> \MusicPDFRef{13617} <6571756174696F6E2E372E312E312E37> \MusicPDFRef{13618} <6571756174696F6E2E372E312E312E38> \MusicPDFRef{13619} <6571756174696F6E2E372E312E322E31> \MusicPDFRef{13620}] /Limits [<6571756174696F6E2E372E312E312E34> <6571756174696F6E2E372E312E322E31>] >>}
\MusicPDFWriteObject{9673}{<< /Names [<6571756174696F6E2E372E312E322E32> \MusicPDFRef{13621} <6571756174696F6E2E372E312E322E33> \MusicPDFRef{13622} <6571756174696F6E2E372E312E322E34> \MusicPDFRef{13623} <6571756174696F6E2E372E312E322E35> \MusicPDFRef{13624} <6571756174696F6E2E372E312E332E31> \MusicPDFRef{13625} <6571756174696F6E2E372E312E332E32> \MusicPDFRef{13626}] /Limits [<6571756174696F6E2E372E312E322E32> <6571756174696F6E2E372E312E332E32>] >>}
\MusicPDFWriteObject{9674}{<< /Names [<6571756174696F6E2E372E312E342E31> \MusicPDFRef{13627} <6571756174696F6E2E372E312E342E3130> \MusicPDFRef{13628} <6571756174696F6E2E372E312E342E3131> \MusicPDFRef{13629} <6571756174696F6E2E372E312E342E3132> \MusicPDFRef{13630} <6571756174696F6E2E372E312E342E3133> \MusicPDFRef{13631} <6571756174696F6E2E372E312E342E3134> \MusicPDFRef{13632}] /Limits [<6571756174696F6E2E372E312E342E31> <6571756174696F6E2E372E312E342E3134>] >>}
\MusicPDFWriteObject{9675}{<< /Names [<6571756174696F6E2E372E312E342E3135> \MusicPDFRef{13633} <6571756174696F6E2E372E312E342E3136> \MusicPDFRef{13634} <6571756174696F6E2E372E312E342E3137> \MusicPDFRef{13635} <6571756174696F6E2E372E312E342E3138> \MusicPDFRef{13636} <6571756174696F6E2E372E312E342E3139> \MusicPDFRef{13637} <6571756174696F6E2E372E312E342E32> \MusicPDFRef{13638}] /Limits [<6571756174696F6E2E372E312E342E3135> <6571756174696F6E2E372E312E342E32>] >>}
\MusicPDFWriteObject{9676}{<< /Names [<6571756174696F6E2E372E312E342E3230> \MusicPDFRef{13639} <6571756174696F6E2E372E312E342E3231> \MusicPDFRef{13640} <6571756174696F6E2E372E312E342E3232> \MusicPDFRef{13641} <6571756174696F6E2E372E312E342E3233> \MusicPDFRef{13642} <6571756174696F6E2E372E312E342E3234> \MusicPDFRef{13643} <6571756174696F6E2E372E312E342E3235> \MusicPDFRef{13644}] /Limits [<6571756174696F6E2E372E312E342E3230> <6571756174696F6E2E372E312E342E3235>] >>}
\MusicPDFWriteObject{9677}{<< /Names [<6571756174696F6E2E372E312E342E3236> \MusicPDFRef{13645} <6571756174696F6E2E372E312E342E3237> \MusicPDFRef{13646} <6571756174696F6E2E372E312E342E3238> \MusicPDFRef{13647} <6571756174696F6E2E372E312E342E33> \MusicPDFRef{13648} <6571756174696F6E2E372E312E342E34> \MusicPDFRef{13649} <6571756174696F6E2E372E312E342E35> \MusicPDFRef{13650}] /Limits [<6571756174696F6E2E372E312E342E3236> <6571756174696F6E2E372E312E342E35>] >>}
\MusicPDFWriteObject{9678}{<< /Names [<6571756174696F6E2E372E312E342E36> \MusicPDFRef{13651} <6571756174696F6E2E372E312E342E37> \MusicPDFRef{13652} <6571756174696F6E2E372E312E342E38> \MusicPDFRef{13653} <6571756174696F6E2E372E312E342E39> \MusicPDFRef{13654} <6571756174696F6E2E372E322E302E31> \MusicPDFRef{13655} <6571756174696F6E2E372E322E302E32> \MusicPDFRef{13656}] /Limits [<6571756174696F6E2E372E312E342E36> <6571756174696F6E2E372E322E302E32>] >>}
\MusicPDFWriteObject{9679}{<< /Names [<6571756174696F6E2E372E322E312E31> \MusicPDFRef{13657} <6571756174696F6E2E372E322E312E3130> \MusicPDFRef{13658} <6571756174696F6E2E372E322E312E3131> \MusicPDFRef{13659} <6571756174696F6E2E372E322E312E3132> \MusicPDFRef{13660} <6571756174696F6E2E372E322E312E3133> \MusicPDFRef{13661} <6571756174696F6E2E372E322E312E3134> \MusicPDFRef{13662}] /Limits [<6571756174696F6E2E372E322E312E31> <6571756174696F6E2E372E322E312E3134>] >>}
\MusicPDFWriteObject{9680}{<< /Names [<6571756174696F6E2E372E322E312E3135> \MusicPDFRef{13663} <6571756174696F6E2E372E322E312E3136> \MusicPDFRef{13664} <6571756174696F6E2E372E322E312E3137> \MusicPDFRef{13665} <6571756174696F6E2E372E322E312E3138> \MusicPDFRef{13666} <6571756174696F6E2E372E322E312E3139> \MusicPDFRef{13667} <6571756174696F6E2E372E322E312E32> \MusicPDFRef{13668}] /Limits [<6571756174696F6E2E372E322E312E3135> <6571756174696F6E2E372E322E312E32>] >>}
\MusicPDFWriteObject{9681}{<< /Names [<6571756174696F6E2E372E322E312E3230> \MusicPDFRef{13669} <6571756174696F6E2E372E322E312E3231> \MusicPDFRef{13670} <6571756174696F6E2E372E322E312E3232> \MusicPDFRef{13671} <6571756174696F6E2E372E322E312E3233> \MusicPDFRef{13672} <6571756174696F6E2E372E322E312E3234> \MusicPDFRef{13673} <6571756174696F6E2E372E322E312E3235> \MusicPDFRef{13674}] /Limits [<6571756174696F6E2E372E322E312E3230> <6571756174696F6E2E372E322E312E3235>] >>}
\MusicPDFWriteObject{9682}{<< /Names [<6571756174696F6E2E372E322E312E3236> \MusicPDFRef{13675} <6571756174696F6E2E372E322E312E3237> \MusicPDFRef{13676} <6571756174696F6E2E372E322E312E33> \MusicPDFRef{13677} <6571756174696F6E2E372E322E312E34> \MusicPDFRef{13678} <6571756174696F6E2E372E322E312E35> \MusicPDFRef{13679} <6571756174696F6E2E372E322E312E36> \MusicPDFRef{13680}] /Limits [<6571756174696F6E2E372E322E312E3236> <6571756174696F6E2E372E322E312E36>] >>}
\MusicPDFWriteObject{9683}{<< /Names [<6571756174696F6E2E372E322E312E37> \MusicPDFRef{13681} <6571756174696F6E2E372E322E312E38> \MusicPDFRef{13682} <6571756174696F6E2E372E322E312E39> \MusicPDFRef{13683} <6571756174696F6E2E372E322E322E31> \MusicPDFRef{13684} <6571756174696F6E2E372E322E322E32> \MusicPDFRef{13685} <6571756174696F6E2E372E322E322E33> \MusicPDFRef{13686}] /Limits [<6571756174696F6E2E372E322E312E37> <6571756174696F6E2E372E322E322E33>] >>}
\MusicPDFWriteObject{9684}{<< /Names [<6571756174696F6E2E372E332E312E31> \MusicPDFRef{13687} <6571756174696F6E2E372E332E312E3130> \MusicPDFRef{13688} <6571756174696F6E2E372E332E312E3131> \MusicPDFRef{13689} <6571756174696F6E2E372E332E312E3132> \MusicPDFRef{13690} <6571756174696F6E2E372E332E312E3133> \MusicPDFRef{13691} <6571756174696F6E2E372E332E312E3134> \MusicPDFRef{13692}] /Limits [<6571756174696F6E2E372E332E312E31> <6571756174696F6E2E372E332E312E3134>] >>}
\MusicPDFWriteObject{9685}{<< /Names [<6571756174696F6E2E372E332E312E3135> \MusicPDFRef{13693} <6571756174696F6E2E372E332E312E3136> \MusicPDFRef{13694} <6571756174696F6E2E372E332E312E3137> \MusicPDFRef{13695} <6571756174696F6E2E372E332E312E3138> \MusicPDFRef{13696} <6571756174696F6E2E372E332E312E3139> \MusicPDFRef{13697} <6571756174696F6E2E372E332E312E32> \MusicPDFRef{13698}] /Limits [<6571756174696F6E2E372E332E312E3135> <6571756174696F6E2E372E332E312E32>] >>}
\MusicPDFWriteObject{9686}{<< /Names [<6571756174696F6E2E372E332E312E3230> \MusicPDFRef{13699} <6571756174696F6E2E372E332E312E3231> \MusicPDFRef{13700} <6571756174696F6E2E372E332E312E3232> \MusicPDFRef{13701} <6571756174696F6E2E372E332E312E3233> \MusicPDFRef{13702} <6571756174696F6E2E372E332E312E3234> \MusicPDFRef{13703} <6571756174696F6E2E372E332E312E3235> \MusicPDFRef{13704}] /Limits [<6571756174696F6E2E372E332E312E3230> <6571756174696F6E2E372E332E312E3235>] >>}
\MusicPDFWriteObject{9687}{<< /Names [<6571756174696F6E2E372E332E312E3236> \MusicPDFRef{13705} <6571756174696F6E2E372E332E312E33> \MusicPDFRef{13706} <6571756174696F6E2E372E332E312E34> \MusicPDFRef{13707} <6571756174696F6E2E372E332E312E35> \MusicPDFRef{13708} <6571756174696F6E2E372E332E312E36> \MusicPDFRef{13709} <6571756174696F6E2E372E332E312E37> \MusicPDFRef{13710}] /Limits [<6571756174696F6E2E372E332E312E3236> <6571756174696F6E2E372E332E312E37>] >>}
\MusicPDFWriteObject{9688}{<< /Names [<6571756174696F6E2E372E332E312E38> \MusicPDFRef{13711} <6571756174696F6E2E372E332E312E39> \MusicPDFRef{13712} <6571756174696F6E2E372E332E322E31> \MusicPDFRef{13713} <6571756174696F6E2E372E332E322E3130> \MusicPDFRef{13714} <6571756174696F6E2E372E332E322E3131> \MusicPDFRef{13715} <6571756174696F6E2E372E332E322E3132> \MusicPDFRef{13716}] /Limits [<6571756174696F6E2E372E332E312E38> <6571756174696F6E2E372E332E322E3132>] >>}
\MusicPDFWriteObject{9689}{<< /Names [<6571756174696F6E2E372E332E322E3133> \MusicPDFRef{13717} <6571756174696F6E2E372E332E322E3134> \MusicPDFRef{13718} <6571756174696F6E2E372E332E322E3135> \MusicPDFRef{13719} <6571756174696F6E2E372E332E322E3136> \MusicPDFRef{13720} <6571756174696F6E2E372E332E322E3137> \MusicPDFRef{13721} <6571756174696F6E2E372E332E322E3138> \MusicPDFRef{13722}] /Limits [<6571756174696F6E2E372E332E322E3133> <6571756174696F6E2E372E332E322E3138>] >>}
\MusicPDFWriteObject{9690}{<< /Names [<6571756174696F6E2E372E332E322E3139> \MusicPDFRef{13723} <6571756174696F6E2E372E332E322E32> \MusicPDFRef{13724} <6571756174696F6E2E372E332E322E3230> \MusicPDFRef{13725} <6571756174696F6E2E372E332E322E3231> \MusicPDFRef{13726} <6571756174696F6E2E372E332E322E3232> \MusicPDFRef{13727} <6571756174696F6E2E372E332E322E33> \MusicPDFRef{13728}] /Limits [<6571756174696F6E2E372E332E322E3139> <6571756174696F6E2E372E332E322E33>] >>}
\MusicPDFWriteObject{9691}{<< /Names [<6571756174696F6E2E372E332E322E34> \MusicPDFRef{13729} <6571756174696F6E2E372E332E322E35> \MusicPDFRef{13730} <6571756174696F6E2E372E332E322E36> \MusicPDFRef{13731} <6571756174696F6E2E372E332E322E37> \MusicPDFRef{13732} <6571756174696F6E2E372E332E322E38> \MusicPDFRef{13733} <6571756174696F6E2E372E332E322E39> \MusicPDFRef{13734}] /Limits [<6571756174696F6E2E372E332E322E34> <6571756174696F6E2E372E332E322E39>] >>}
\MusicPDFWriteObject{9692}{<< /Names [<6571756174696F6E2E372E332E332E31> \MusicPDFRef{13735} <6571756174696F6E2E372E332E332E3130> \MusicPDFRef{13736} <6571756174696F6E2E372E332E332E3131> \MusicPDFRef{13737} <6571756174696F6E2E372E332E332E3132> \MusicPDFRef{13738} <6571756174696F6E2E372E332E332E3133> \MusicPDFRef{13739} <6571756174696F6E2E372E332E332E3134> \MusicPDFRef{13740}] /Limits [<6571756174696F6E2E372E332E332E31> <6571756174696F6E2E372E332E332E3134>] >>}
\MusicPDFWriteObject{9693}{<< /Names [<6571756174696F6E2E372E332E332E3135> \MusicPDFRef{13741} <6571756174696F6E2E372E332E332E3136> \MusicPDFRef{13742} <6571756174696F6E2E372E332E332E3137> \MusicPDFRef{13743} <6571756174696F6E2E372E332E332E3138> \MusicPDFRef{13744} <6571756174696F6E2E372E332E332E3139> \MusicPDFRef{13745} <6571756174696F6E2E372E332E332E32> \MusicPDFRef{13746}] /Limits [<6571756174696F6E2E372E332E332E3135> <6571756174696F6E2E372E332E332E32>] >>}
\MusicPDFWriteObject{9694}{<< /Names [<6571756174696F6E2E372E332E332E3230> \MusicPDFRef{13747} <6571756174696F6E2E372E332E332E3231> \MusicPDFRef{13748} <6571756174696F6E2E372E332E332E3232> \MusicPDFRef{13749} <6571756174696F6E2E372E332E332E3233> \MusicPDFRef{13750} <6571756174696F6E2E372E332E332E3234> \MusicPDFRef{13751} <6571756174696F6E2E372E332E332E3235> \MusicPDFRef{13752}] /Limits [<6571756174696F6E2E372E332E332E3230> <6571756174696F6E2E372E332E332E3235>] >>}
\MusicPDFWriteObject{9695}{<< /Names [<6571756174696F6E2E372E332E332E3236> \MusicPDFRef{13753} <6571756174696F6E2E372E332E332E3237> \MusicPDFRef{13754} <6571756174696F6E2E372E332E332E3238> \MusicPDFRef{13755} <6571756174696F6E2E372E332E332E3239> \MusicPDFRef{13756} <6571756174696F6E2E372E332E332E33> \MusicPDFRef{13757} <6571756174696F6E2E372E332E332E3330> \MusicPDFRef{13758}] /Limits [<6571756174696F6E2E372E332E332E3236> <6571756174696F6E2E372E332E332E3330>] >>}
\MusicPDFWriteObject{9696}{<< /Names [<6571756174696F6E2E372E332E332E34> \MusicPDFRef{13759} <6571756174696F6E2E372E332E332E35> \MusicPDFRef{13760} <6571756174696F6E2E372E332E332E36> \MusicPDFRef{13761} <6571756174696F6E2E372E332E332E37> \MusicPDFRef{13762} <6571756174696F6E2E372E332E332E38> \MusicPDFRef{13763} <6571756174696F6E2E372E332E332E39> \MusicPDFRef{13764}] /Limits [<6571756174696F6E2E372E332E332E34> <6571756174696F6E2E372E332E332E39>] >>}
\MusicPDFWriteObject{9697}{<< /Names [<6571756174696F6E2E372E332E342E31> \MusicPDFRef{13765} <6571756174696F6E2E372E332E342E32> \MusicPDFRef{13766} <6571756174696F6E2E372E332E342E33> \MusicPDFRef{13767} <6571756174696F6E2E372E332E342E34> \MusicPDFRef{13768} <6571756174696F6E2E372E332E342E35> \MusicPDFRef{13769} <6571756174696F6E2E372E332E342E36> \MusicPDFRef{13770}] /Limits [<6571756174696F6E2E372E332E342E31> <6571756174696F6E2E372E332E342E36>] >>}
\MusicPDFWriteObject{9698}{<< /Names [<6571756174696F6E2E372E332E342E37> \MusicPDFRef{13771} <6571756174696F6E2E372E332E342E38> \MusicPDFRef{13772} <6571756174696F6E2E372E332E342E39> \MusicPDFRef{13773} <6571756174696F6E2E372E342E302E31> \MusicPDFRef{13774} <6571756174696F6E2E372E342E302E32> \MusicPDFRef{13775} <6571756174696F6E2E372E342E312E31> \MusicPDFRef{13776}] /Limits [<6571756174696F6E2E372E332E342E37> <6571756174696F6E2E372E342E312E31>] >>}
\MusicPDFWriteObject{9699}{<< /Names [<6571756174696F6E2E372E342E312E32> \MusicPDFRef{13777} <6571756174696F6E2E372E342E322E31> \MusicPDFRef{13778} <6571756174696F6E2E372E342E332E31> \MusicPDFRef{13779} <6571756174696F6E2E372E342E332E3130> \MusicPDFRef{13780} <6571756174696F6E2E372E342E332E3131> \MusicPDFRef{13781} <6571756174696F6E2E372E342E332E3132> \MusicPDFRef{13782}] /Limits [<6571756174696F6E2E372E342E312E32> <6571756174696F6E2E372E342E332E3132>] >>}
\MusicPDFWriteObject{9700}{<< /Names [<6571756174696F6E2E372E342E332E3133> \MusicPDFRef{13783} <6571756174696F6E2E372E342E332E3134> \MusicPDFRef{13784} <6571756174696F6E2E372E342E332E3135> \MusicPDFRef{13785} <6571756174696F6E2E372E342E332E3136> \MusicPDFRef{13786} <6571756174696F6E2E372E342E332E3137> \MusicPDFRef{13787} <6571756174696F6E2E372E342E332E32> \MusicPDFRef{13788}] /Limits [<6571756174696F6E2E372E342E332E3133> <6571756174696F6E2E372E342E332E32>] >>}
\MusicPDFWriteObject{9701}{<< /Names [<6571756174696F6E2E372E342E332E33> \MusicPDFRef{13789} <6571756174696F6E2E372E342E332E34> \MusicPDFRef{13790} <6571756174696F6E2E372E342E332E35> \MusicPDFRef{13791} <6571756174696F6E2E372E342E332E36> \MusicPDFRef{13792} <6571756174696F6E2E372E342E332E37> \MusicPDFRef{13793} <6571756174696F6E2E372E342E332E38> \MusicPDFRef{13794}] /Limits [<6571756174696F6E2E372E342E332E33> <6571756174696F6E2E372E342E332E38>] >>}
\MusicPDFWriteObject{9702}{<< /Names [<6571756174696F6E2E372E342E332E39> \MusicPDFRef{13795} <6571756174696F6E2E372E352E302E31> \MusicPDFRef{13796} <6571756174696F6E2E372E352E302E3130> \MusicPDFRef{13797} <6571756174696F6E2E372E352E302E3131> \MusicPDFRef{13798} <6571756174696F6E2E372E352E302E3132> \MusicPDFRef{13799} <6571756174696F6E2E372E352E302E3133> \MusicPDFRef{13800}] /Limits [<6571756174696F6E2E372E342E332E39> <6571756174696F6E2E372E352E302E3133>] >>}
\MusicPDFWriteObject{9703}{<< /Names [<6571756174696F6E2E372E352E302E3134> \MusicPDFRef{13801} <6571756174696F6E2E372E352E302E3135> \MusicPDFRef{13802} <6571756174696F6E2E372E352E302E3136> \MusicPDFRef{13803} <6571756174696F6E2E372E352E302E3137> \MusicPDFRef{13804} <6571756174696F6E2E372E352E302E32> \MusicPDFRef{13805} <6571756174696F6E2E372E352E302E33> \MusicPDFRef{13806}] /Limits [<6571756174696F6E2E372E352E302E3134> <6571756174696F6E2E372E352E302E33>] >>}
\MusicPDFWriteObject{9704}{<< /Names [<6571756174696F6E2E372E352E302E34> \MusicPDFRef{13807} <6571756174696F6E2E372E352E302E35> \MusicPDFRef{13808} <6571756174696F6E2E372E352E302E36> \MusicPDFRef{13809} <6571756174696F6E2E372E352E302E37> \MusicPDFRef{13810} <6571756174696F6E2E372E352E302E38> \MusicPDFRef{13811} <6571756174696F6E2E372E352E302E39> \MusicPDFRef{13812}] /Limits [<6571756174696F6E2E372E352E302E34> <6571756174696F6E2E372E352E302E39>] >>}
\MusicPDFWriteObject{9705}{<< /Names [<6571756174696F6E2E382E312E302E31> \MusicPDFRef{13813} <6571756174696F6E2E382E312E302E3130> \MusicPDFRef{13814} <6571756174696F6E2E382E312E302E32> \MusicPDFRef{13815} <6571756174696F6E2E382E312E302E33> \MusicPDFRef{13816} <6571756174696F6E2E382E312E302E34> \MusicPDFRef{13817} <6571756174696F6E2E382E312E302E35> \MusicPDFRef{13818}] /Limits [<6571756174696F6E2E382E312E302E31> <6571756174696F6E2E382E312E302E35>] >>}
\MusicPDFWriteObject{9706}{<< /Names [<6571756174696F6E2E382E312E302E36> \MusicPDFRef{13819} <6571756174696F6E2E382E312E302E37> \MusicPDFRef{13820} <6571756174696F6E2E382E312E302E38> \MusicPDFRef{13821} <6571756174696F6E2E382E312E302E39> \MusicPDFRef{13822} <6571756174696F6E2E382E322E302E31> \MusicPDFRef{13823} <6571756174696F6E2E382E322E302E3130> \MusicPDFRef{13824}] /Limits [<6571756174696F6E2E382E312E302E36> <6571756174696F6E2E382E322E302E3130>] >>}
\MusicPDFWriteObject{9707}{<< /Names [<6571756174696F6E2E382E322E302E32> \MusicPDFRef{13825} <6571756174696F6E2E382E322E302E33> \MusicPDFRef{13826} <6571756174696F6E2E382E322E302E34> \MusicPDFRef{13827} <6571756174696F6E2E382E322E302E35> \MusicPDFRef{13828} <6571756174696F6E2E382E322E302E36> \MusicPDFRef{13829} <6571756174696F6E2E382E322E302E37> \MusicPDFRef{13830}] /Limits [<6571756174696F6E2E382E322E302E32> <6571756174696F6E2E382E322E302E37>] >>}
\MusicPDFWriteObject{9708}{<< /Names [<6571756174696F6E2E382E322E302E38> \MusicPDFRef{13831} <6571756174696F6E2E382E322E302E39> \MusicPDFRef{13832} <6571756174696F6E2E382E332E302E31> \MusicPDFRef{13833} <6571756174696F6E2E382E332E302E32> \MusicPDFRef{13834} <6571756174696F6E2E382E332E302E33> \MusicPDFRef{13835} <6571756174696F6E2E382E332E302E34> \MusicPDFRef{13836}] /Limits [<6571756174696F6E2E382E322E302E38> <6571756174696F6E2E382E332E302E34>] >>}
\MusicPDFWriteObject{9709}{<< /Names [<6571756174696F6E2E382E332E302E35> \MusicPDFRef{13837} <6571756174696F6E2E382E332E302E36> \MusicPDFRef{13838} <6571756174696F6E2E382E332E302E37> \MusicPDFRef{13839} <6571756174696F6E2E382E332E302E38> \MusicPDFRef{13840} <6571756174696F6E2E382E342E302E31> \MusicPDFRef{13841} <6571756174696F6E2E382E342E302E3130> \MusicPDFRef{13842}] /Limits [<6571756174696F6E2E382E332E302E35> <6571756174696F6E2E382E342E302E3130>] >>}
\MusicPDFWriteObject{9710}{<< /Names [<6571756174696F6E2E382E342E302E3131> \MusicPDFRef{13843} <6571756174696F6E2E382E342E302E3132> \MusicPDFRef{13844} <6571756174696F6E2E382E342E302E3133> \MusicPDFRef{13845} <6571756174696F6E2E382E342E302E3134> \MusicPDFRef{13846} <6571756174696F6E2E382E342E302E32> \MusicPDFRef{13847} <6571756174696F6E2E382E342E302E33> \MusicPDFRef{13848}] /Limits [<6571756174696F6E2E382E342E302E3131> <6571756174696F6E2E382E342E302E33>] >>}
\MusicPDFWriteObject{9711}{<< /Names [<6571756174696F6E2E382E342E302E34> \MusicPDFRef{13849} <6571756174696F6E2E382E342E302E35> \MusicPDFRef{13850} <6571756174696F6E2E382E342E302E36> \MusicPDFRef{13851} <6571756174696F6E2E382E342E302E37> \MusicPDFRef{13852} <6571756174696F6E2E382E342E302E38> \MusicPDFRef{13853} <6571756174696F6E2E382E342E302E39> \MusicPDFRef{13854}] /Limits [<6571756174696F6E2E382E342E302E34> <6571756174696F6E2E382E342E302E39>] >>}
\MusicPDFWriteObject{9712}{<< /Names [<6571756174696F6E2E392E312E312E31> \MusicPDFRef{13855} <6571756174696F6E2E392E312E312E3130> \MusicPDFRef{13856} <6571756174696F6E2E392E312E312E3131> \MusicPDFRef{13857} <6571756174696F6E2E392E312E312E3132> \MusicPDFRef{13858} <6571756174696F6E2E392E312E312E3133> \MusicPDFRef{13859} <6571756174696F6E2E392E312E312E3134> \MusicPDFRef{13860}] /Limits [<6571756174696F6E2E392E312E312E31> <6571756174696F6E2E392E312E312E3134>] >>}
\MusicPDFWriteObject{9713}{<< /Names [<6571756174696F6E2E392E312E312E3135> \MusicPDFRef{13861} <6571756174696F6E2E392E312E312E3136> \MusicPDFRef{13862} <6571756174696F6E2E392E312E312E3137> \MusicPDFRef{13863} <6571756174696F6E2E392E312E312E32> \MusicPDFRef{13864} <6571756174696F6E2E392E312E312E33> \MusicPDFRef{13865} <6571756174696F6E2E392E312E312E34> \MusicPDFRef{13866}] /Limits [<6571756174696F6E2E392E312E312E3135> <6571756174696F6E2E392E312E312E34>] >>}
\MusicPDFWriteObject{9714}{<< /Names [<6571756174696F6E2E392E312E312E35> \MusicPDFRef{13867} <6571756174696F6E2E392E312E312E36> \MusicPDFRef{13868} <6571756174696F6E2E392E312E312E37> \MusicPDFRef{13869} <6571756174696F6E2E392E312E312E38> \MusicPDFRef{13870} <6571756174696F6E2E392E312E312E39> \MusicPDFRef{13871} <6571756174696F6E2E392E312E322E31> \MusicPDFRef{13872}] /Limits [<6571756174696F6E2E392E312E312E35> <6571756174696F6E2E392E312E322E31>] >>}
\MusicPDFWriteObject{9715}{<< /Names [<6571756174696F6E2E392E312E322E3130> \MusicPDFRef{13873} <6571756174696F6E2E392E312E322E3131> \MusicPDFRef{13874} <6571756174696F6E2E392E312E322E3132> \MusicPDFRef{13875} <6571756174696F6E2E392E312E322E3133> \MusicPDFRef{13876} <6571756174696F6E2E392E312E322E3134> \MusicPDFRef{13877} <6571756174696F6E2E392E312E322E3135> \MusicPDFRef{13878}] /Limits [<6571756174696F6E2E392E312E322E3130> <6571756174696F6E2E392E312E322E3135>] >>}
\MusicPDFWriteObject{9716}{<< /Names [<6571756174696F6E2E392E312E322E3136> \MusicPDFRef{13879} <6571756174696F6E2E392E312E322E3137> \MusicPDFRef{13880} <6571756174696F6E2E392E312E322E32> \MusicPDFRef{13881} <6571756174696F6E2E392E312E322E33> \MusicPDFRef{13882} <6571756174696F6E2E392E312E322E34> \MusicPDFRef{13883} <6571756174696F6E2E392E312E322E35> \MusicPDFRef{13884}] /Limits [<6571756174696F6E2E392E312E322E3136> <6571756174696F6E2E392E312E322E35>] >>}
\MusicPDFWriteObject{9717}{<< /Names [<6571756174696F6E2E392E312E322E36> \MusicPDFRef{13885} <6571756174696F6E2E392E312E322E37> \MusicPDFRef{13886} <6571756174696F6E2E392E312E322E38> \MusicPDFRef{13887} <6571756174696F6E2E392E312E322E39> \MusicPDFRef{13888} <6571756174696F6E2E392E322E302E31> \MusicPDFRef{13889} <6571756174696F6E2E392E322E312E31> \MusicPDFRef{13890}] /Limits [<6571756174696F6E2E392E312E322E36> <6571756174696F6E2E392E322E312E31>] >>}
\MusicPDFWriteObject{9718}{<< /Names [<6571756174696F6E2E392E322E312E3130> \MusicPDFRef{13891} <6571756174696F6E2E392E322E312E3131> \MusicPDFRef{13892} <6571756174696F6E2E392E322E312E3132> \MusicPDFRef{13893} <6571756174696F6E2E392E322E312E3133> \MusicPDFRef{13894} <6571756174696F6E2E392E322E312E3134> \MusicPDFRef{13895} <6571756174696F6E2E392E322E312E3135> \MusicPDFRef{13896}] /Limits [<6571756174696F6E2E392E322E312E3130> <6571756174696F6E2E392E322E312E3135>] >>}
\MusicPDFWriteObject{9719}{<< /Names [<6571756174696F6E2E392E322E312E3136> \MusicPDFRef{13897} <6571756174696F6E2E392E322E312E3137> \MusicPDFRef{13898} <6571756174696F6E2E392E322E312E3138> \MusicPDFRef{13899} <6571756174696F6E2E392E322E312E3139> \MusicPDFRef{13900} <6571756174696F6E2E392E322E312E32> \MusicPDFRef{13901} <6571756174696F6E2E392E322E312E3230> \MusicPDFRef{13902}] /Limits [<6571756174696F6E2E392E322E312E3136> <6571756174696F6E2E392E322E312E3230>] >>}
\MusicPDFWriteObject{9720}{<< /Names [<6571756174696F6E2E392E322E312E3231> \MusicPDFRef{13903} <6571756174696F6E2E392E322E312E3232> \MusicPDFRef{13904} <6571756174696F6E2E392E322E312E3233> \MusicPDFRef{13905} <6571756174696F6E2E392E322E312E3234> \MusicPDFRef{13906} <6571756174696F6E2E392E322E312E3235> \MusicPDFRef{13907} <6571756174696F6E2E392E322E312E3236> \MusicPDFRef{13908}] /Limits [<6571756174696F6E2E392E322E312E3231> <6571756174696F6E2E392E322E312E3236>] >>}
\MusicPDFWriteObject{9721}{<< /Names [<6571756174696F6E2E392E322E312E3237> \MusicPDFRef{13909} <6571756174696F6E2E392E322E312E3238> \MusicPDFRef{13910} <6571756174696F6E2E392E322E312E3239> \MusicPDFRef{13911} <6571756174696F6E2E392E322E312E33> \MusicPDFRef{13912} <6571756174696F6E2E392E322E312E3330> \MusicPDFRef{13913} <6571756174696F6E2E392E322E312E3331> \MusicPDFRef{13914}] /Limits [<6571756174696F6E2E392E322E312E3237> <6571756174696F6E2E392E322E312E3331>] >>}
\MusicPDFWriteObject{9722}{<< /Names [<6571756174696F6E2E392E322E312E3332> \MusicPDFRef{13915} <6571756174696F6E2E392E322E312E3333> \MusicPDFRef{13916} <6571756174696F6E2E392E322E312E3334> \MusicPDFRef{13917} <6571756174696F6E2E392E322E312E3335> \MusicPDFRef{13918} <6571756174696F6E2E392E322E312E3336> \MusicPDFRef{13919} <6571756174696F6E2E392E322E312E34> \MusicPDFRef{13920}] /Limits [<6571756174696F6E2E392E322E312E3332> <6571756174696F6E2E392E322E312E34>] >>}
\MusicPDFWriteObject{9723}{<< /Names [<6571756174696F6E2E392E322E312E35> \MusicPDFRef{13921} <6571756174696F6E2E392E322E312E36> \MusicPDFRef{13922} <6571756174696F6E2E392E322E312E37> \MusicPDFRef{13923} <6571756174696F6E2E392E322E312E38> \MusicPDFRef{13924} <6571756174696F6E2E392E322E312E39> \MusicPDFRef{13925} <6571756174696F6E2E392E322E322E31> \MusicPDFRef{13926}] /Limits [<6571756174696F6E2E392E322E312E35> <6571756174696F6E2E392E322E322E31>] >>}
\MusicPDFWriteObject{9724}{<< /Names [<6571756174696F6E2E392E322E322E3130> \MusicPDFRef{13927} <6571756174696F6E2E392E322E322E3131> \MusicPDFRef{13928} <6571756174696F6E2E392E322E322E3132> \MusicPDFRef{13929} <6571756174696F6E2E392E322E322E3133> \MusicPDFRef{13930} <6571756174696F6E2E392E322E322E3134> \MusicPDFRef{13931} <6571756174696F6E2E392E322E322E3135> \MusicPDFRef{13932}] /Limits [<6571756174696F6E2E392E322E322E3130> <6571756174696F6E2E392E322E322E3135>] >>}
\MusicPDFWriteObject{9725}{<< /Names [<6571756174696F6E2E392E322E322E3136> \MusicPDFRef{13933} <6571756174696F6E2E392E322E322E3137> \MusicPDFRef{13934} <6571756174696F6E2E392E322E322E3138> \MusicPDFRef{13935} <6571756174696F6E2E392E322E322E3139> \MusicPDFRef{13936} <6571756174696F6E2E392E322E322E32> \MusicPDFRef{13937} <6571756174696F6E2E392E322E322E3230> \MusicPDFRef{13938}] /Limits [<6571756174696F6E2E392E322E322E3136> <6571756174696F6E2E392E322E322E3230>] >>}
\MusicPDFWriteObject{9726}{<< /Names [<6571756174696F6E2E392E322E322E3231> \MusicPDFRef{13939} <6571756174696F6E2E392E322E322E3232> \MusicPDFRef{13940} <6571756174696F6E2E392E322E322E3233> \MusicPDFRef{13941} <6571756174696F6E2E392E322E322E3234> \MusicPDFRef{13942} <6571756174696F6E2E392E322E322E3235> \MusicPDFRef{13943} <6571756174696F6E2E392E322E322E3236> \MusicPDFRef{13944}] /Limits [<6571756174696F6E2E392E322E322E3231> <6571756174696F6E2E392E322E322E3236>] >>}
\MusicPDFWriteObject{9727}{<< /Names [<6571756174696F6E2E392E322E322E3237> \MusicPDFRef{13945} <6571756174696F6E2E392E322E322E3238> \MusicPDFRef{13946} <6571756174696F6E2E392E322E322E3239> \MusicPDFRef{13947} <6571756174696F6E2E392E322E322E33> \MusicPDFRef{13948} <6571756174696F6E2E392E322E322E3330> \MusicPDFRef{13949} <6571756174696F6E2E392E322E322E3331> \MusicPDFRef{13950}] /Limits [<6571756174696F6E2E392E322E322E3237> <6571756174696F6E2E392E322E322E3331>] >>}
\MusicPDFWriteObject{9728}{<< /Names [<6571756174696F6E2E392E322E322E3332> \MusicPDFRef{13951} <6571756174696F6E2E392E322E322E3333> \MusicPDFRef{13952} <6571756174696F6E2E392E322E322E3334> \MusicPDFRef{13953} <6571756174696F6E2E392E322E322E3335> \MusicPDFRef{13954} <6571756174696F6E2E392E322E322E3336> \MusicPDFRef{13955} <6571756174696F6E2E392E322E322E3337> \MusicPDFRef{13956}] /Limits [<6571756174696F6E2E392E322E322E3332> <6571756174696F6E2E392E322E322E3337>] >>}
\MusicPDFWriteObject{9729}{<< /Names [<6571756174696F6E2E392E322E322E3338> \MusicPDFRef{13957} <6571756174696F6E2E392E322E322E3339> \MusicPDFRef{13958} <6571756174696F6E2E392E322E322E34> \MusicPDFRef{13959} <6571756174696F6E2E392E322E322E3430> \MusicPDFRef{13960} <6571756174696F6E2E392E322E322E3431> \MusicPDFRef{13961} <6571756174696F6E2E392E322E322E3432> \MusicPDFRef{13962}] /Limits [<6571756174696F6E2E392E322E322E3338> <6571756174696F6E2E392E322E322E3432>] >>}
\MusicPDFWriteObject{9730}{<< /Names [<6571756174696F6E2E392E322E322E3433> \MusicPDFRef{13963} <6571756174696F6E2E392E322E322E3434> \MusicPDFRef{13964} <6571756174696F6E2E392E322E322E3435> \MusicPDFRef{13965} <6571756174696F6E2E392E322E322E3436> \MusicPDFRef{13966} <6571756174696F6E2E392E322E322E3437> \MusicPDFRef{13967} <6571756174696F6E2E392E322E322E35> \MusicPDFRef{13968}] /Limits [<6571756174696F6E2E392E322E322E3433> <6571756174696F6E2E392E322E322E35>] >>}
\MusicPDFWriteObject{9731}{<< /Names [<6571756174696F6E2E392E322E322E36> \MusicPDFRef{13969} <6571756174696F6E2E392E322E322E37> \MusicPDFRef{13970} <6571756174696F6E2E392E322E322E38> \MusicPDFRef{13971} <6571756174696F6E2E392E322E322E39> \MusicPDFRef{13972} <6571756174696F6E2E392E332E312E31> \MusicPDFRef{13973} <6571756174696F6E2E392E332E312E32> \MusicPDFRef{13974}] /Limits [<6571756174696F6E2E392E322E322E36> <6571756174696F6E2E392E332E312E32>] >>}
\MusicPDFWriteObject{9732}{<< /Names [<6571756174696F6E2E392E332E312E33> \MusicPDFRef{13975} <6571756174696F6E2E392E332E312E34> \MusicPDFRef{13976} <6571756174696F6E2E392E332E312E35> \MusicPDFRef{13977} <6571756174696F6E2E392E332E312E36> \MusicPDFRef{13978} <6571756174696F6E2E392E332E312E37> \MusicPDFRef{13979} <6571756174696F6E2E392E332E322E31> \MusicPDFRef{13980}] /Limits [<6571756174696F6E2E392E332E312E33> <6571756174696F6E2E392E332E322E31>] >>}
\MusicPDFWriteObject{9733}{<< /Names [<6571756174696F6E2E392E332E322E3130> \MusicPDFRef{13981} <6571756174696F6E2E392E332E322E32> \MusicPDFRef{13982} <6571756174696F6E2E392E332E322E33> \MusicPDFRef{13983} <6571756174696F6E2E392E332E322E34> \MusicPDFRef{13984} <6571756174696F6E2E392E332E322E35> \MusicPDFRef{13985} <6571756174696F6E2E392E332E322E36> \MusicPDFRef{13986}] /Limits [<6571756174696F6E2E392E332E322E3130> <6571756174696F6E2E392E332E322E36>] >>}
\MusicPDFWriteObject{9734}{<< /Names [<6571756174696F6E2E392E332E322E37> \MusicPDFRef{13987} <6571756174696F6E2E392E332E322E38> \MusicPDFRef{13988} <6571756174696F6E2E392E332E322E39> \MusicPDFRef{13989} <6571756174696F6E2E392E332E332E31> \MusicPDFRef{13990} <6571756174696F6E2E392E332E332E3130> \MusicPDFRef{13991} <6571756174696F6E2E392E332E332E3131> \MusicPDFRef{13992}] /Limits [<6571756174696F6E2E392E332E322E37> <6571756174696F6E2E392E332E332E3131>] >>}
\MusicPDFWriteObject{9735}{<< /Names [<6571756174696F6E2E392E332E332E3132> \MusicPDFRef{13993} <6571756174696F6E2E392E332E332E3133> \MusicPDFRef{13994} <6571756174696F6E2E392E332E332E3134> \MusicPDFRef{13995} <6571756174696F6E2E392E332E332E3135> \MusicPDFRef{13996} <6571756174696F6E2E392E332E332E3136> \MusicPDFRef{13997} <6571756174696F6E2E392E332E332E3137> \MusicPDFRef{13998}] /Limits [<6571756174696F6E2E392E332E332E3132> <6571756174696F6E2E392E332E332E3137>] >>}
\MusicPDFWriteObject{9736}{<< /Names [<6571756174696F6E2E392E332E332E3138> \MusicPDFRef{13999} <6571756174696F6E2E392E332E332E3139> \MusicPDFRef{14000} <6571756174696F6E2E392E332E332E32> \MusicPDFRef{14001} <6571756174696F6E2E392E332E332E3230> \MusicPDFRef{14002} <6571756174696F6E2E392E332E332E3231> \MusicPDFRef{14003} <6571756174696F6E2E392E332E332E3232> \MusicPDFRef{14004}] /Limits [<6571756174696F6E2E392E332E332E3138> <6571756174696F6E2E392E332E332E3232>] >>}
\MusicPDFWriteObject{9737}{<< /Names [<6571756174696F6E2E392E332E332E3233> \MusicPDFRef{14005} <6571756174696F6E2E392E332E332E33> \MusicPDFRef{14006} <6571756174696F6E2E392E332E332E34> \MusicPDFRef{14007} <6571756174696F6E2E392E332E332E35> \MusicPDFRef{14008} <6571756174696F6E2E392E332E332E36> \MusicPDFRef{14009} <6571756174696F6E2E392E332E332E37> \MusicPDFRef{14010}] /Limits [<6571756174696F6E2E392E332E332E3233> <6571756174696F6E2E392E332E332E37>] >>}
\MusicPDFWriteObject{9738}{<< /Names [<6571756174696F6E2E392E332E332E38> \MusicPDFRef{14011} <6571756174696F6E2E392E332E332E39> \MusicPDFRef{14012} <6571756174696F6E2E412E312E312E31> \MusicPDFRef{14013} <6571756174696F6E2E412E312E312E32> \MusicPDFRef{14014} <6571756174696F6E2E412E312E312E33> \MusicPDFRef{14015} <6571756174696F6E2E412E312E312E34> \MusicPDFRef{14016}] /Limits [<6571756174696F6E2E392E332E332E38> <6571756174696F6E2E412E312E312E34>] >>}
\MusicPDFWriteObject{9739}{<< /Names [<6571756174696F6E2E412E312E312E35> \MusicPDFRef{14017} <6571756174696F6E2E412E312E312E36> \MusicPDFRef{14018} <6571756174696F6E2E412E312E322E31> \MusicPDFRef{14019} <6571756174696F6E2E412E312E322E32> \MusicPDFRef{14020} <6571756174696F6E2E412E312E332E31> \MusicPDFRef{14021} <6571756174696F6E2E412E312E332E32> \MusicPDFRef{14022}] /Limits [<6571756174696F6E2E412E312E312E35> <6571756174696F6E2E412E312E332E32>] >>}
\MusicPDFWriteObject{9740}{<< /Names [<6571756174696F6E2E412E312E332E33> \MusicPDFRef{14023} <6571756174696F6E2E412E312E332E34> \MusicPDFRef{14024} <6571756174696F6E2E412E312E332E35> \MusicPDFRef{14025} <6571756174696F6E2E412E312E332E36> \MusicPDFRef{14026} <6571756174696F6E2E412E312E342E31> \MusicPDFRef{14027} <6571756174696F6E2E412E312E342E32> \MusicPDFRef{14028}] /Limits [<6571756174696F6E2E412E312E332E33> <6571756174696F6E2E412E312E342E32>] >>}
\MusicPDFWriteObject{9741}{<< /Names [<6571756174696F6E2E412E312E342E33> \MusicPDFRef{14029} <6571756174696F6E2E412E312E352E31> \MusicPDFRef{14030} <6571756174696F6E2E412E312E352E3130> \MusicPDFRef{14031} <6571756174696F6E2E412E312E352E32> \MusicPDFRef{14032} <6571756174696F6E2E412E312E352E33> \MusicPDFRef{14033} <6571756174696F6E2E412E312E352E34> \MusicPDFRef{14034}] /Limits [<6571756174696F6E2E412E312E342E33> <6571756174696F6E2E412E312E352E34>] >>}
\MusicPDFWriteObject{9742}{<< /Names [<6571756174696F6E2E412E312E352E35> \MusicPDFRef{14035} <6571756174696F6E2E412E312E352E36> \MusicPDFRef{14036} <6571756174696F6E2E412E312E352E37> \MusicPDFRef{14037} <6571756174696F6E2E412E312E352E38> \MusicPDFRef{14038} <6571756174696F6E2E412E312E352E39> \MusicPDFRef{14039} <6571756174696F6E2E412E322E312E31> \MusicPDFRef{14040}] /Limits [<6571756174696F6E2E412E312E352E35> <6571756174696F6E2E412E322E312E31>] >>}
\MusicPDFWriteObject{9743}{<< /Names [<6571756174696F6E2E412E322E312E32> \MusicPDFRef{14041} <6571756174696F6E2E412E322E312E33> \MusicPDFRef{14042} <6571756174696F6E2E412E322E312E34> \MusicPDFRef{14043} <6571756174696F6E2E412E322E312E35> \MusicPDFRef{14044} <6571756174696F6E2E412E322E312E36> \MusicPDFRef{14045} <6571756174696F6E2E412E322E322E31> \MusicPDFRef{14046}] /Limits [<6571756174696F6E2E412E322E312E32> <6571756174696F6E2E412E322E322E31>] >>}
\MusicPDFWriteObject{9744}{<< /Names [<6571756174696F6E2E412E322E322E3130> \MusicPDFRef{14047} <6571756174696F6E2E412E322E322E3131> \MusicPDFRef{14048} <6571756174696F6E2E412E322E322E3132> \MusicPDFRef{14049} <6571756174696F6E2E412E322E322E32> \MusicPDFRef{14050} <6571756174696F6E2E412E322E322E33> \MusicPDFRef{14051} <6571756174696F6E2E412E322E322E34> \MusicPDFRef{14052}] /Limits [<6571756174696F6E2E412E322E322E3130> <6571756174696F6E2E412E322E322E34>] >>}
\MusicPDFWriteObject{9745}{<< /Names [<6571756174696F6E2E412E322E322E35> \MusicPDFRef{14053} <6571756174696F6E2E412E322E322E36> \MusicPDFRef{14054} <6571756174696F6E2E412E322E322E37> \MusicPDFRef{14055} <6571756174696F6E2E412E322E322E38> \MusicPDFRef{14056} <6571756174696F6E2E412E322E322E39> \MusicPDFRef{14057} <6571756174696F6E2E412E322E332E31> \MusicPDFRef{14058}] /Limits [<6571756174696F6E2E412E322E322E35> <6571756174696F6E2E412E322E332E31>] >>}
\MusicPDFWriteObject{9746}{<< /Names [<6571756174696F6E2E412E322E332E3130> \MusicPDFRef{14059} <6571756174696F6E2E412E322E332E3131> \MusicPDFRef{14060} <6571756174696F6E2E412E322E332E3132> \MusicPDFRef{14061} <6571756174696F6E2E412E322E332E3133> \MusicPDFRef{14062} <6571756174696F6E2E412E322E332E3134> \MusicPDFRef{14063} <6571756174696F6E2E412E322E332E3135> \MusicPDFRef{14064}] /Limits [<6571756174696F6E2E412E322E332E3130> <6571756174696F6E2E412E322E332E3135>] >>}
\MusicPDFWriteObject{9747}{<< /Names [<6571756174696F6E2E412E322E332E3136> \MusicPDFRef{14065} <6571756174696F6E2E412E322E332E3137> \MusicPDFRef{14066} <6571756174696F6E2E412E322E332E3138> \MusicPDFRef{14067} <6571756174696F6E2E412E322E332E3139> \MusicPDFRef{14068} <6571756174696F6E2E412E322E332E32> \MusicPDFRef{14069} <6571756174696F6E2E412E322E332E3230> \MusicPDFRef{14070}] /Limits [<6571756174696F6E2E412E322E332E3136> <6571756174696F6E2E412E322E332E3230>] >>}
\MusicPDFWriteObject{9748}{<< /Names [<6571756174696F6E2E412E322E332E3231> \MusicPDFRef{14071} <6571756174696F6E2E412E322E332E3232> \MusicPDFRef{14072} <6571756174696F6E2E412E322E332E3233> \MusicPDFRef{14073} <6571756174696F6E2E412E322E332E3234> \MusicPDFRef{14074} <6571756174696F6E2E412E322E332E3235> \MusicPDFRef{14075} <6571756174696F6E2E412E322E332E3236> \MusicPDFRef{14076}] /Limits [<6571756174696F6E2E412E322E332E3231> <6571756174696F6E2E412E322E332E3236>] >>}
\MusicPDFWriteObject{9749}{<< /Names [<6571756174696F6E2E412E322E332E3237> \MusicPDFRef{14077} <6571756174696F6E2E412E322E332E3238> \MusicPDFRef{14078} <6571756174696F6E2E412E322E332E3239> \MusicPDFRef{14079} <6571756174696F6E2E412E322E332E33> \MusicPDFRef{14080} <6571756174696F6E2E412E322E332E34> \MusicPDFRef{14081} <6571756174696F6E2E412E322E332E35> \MusicPDFRef{14082}] /Limits [<6571756174696F6E2E412E322E332E3237> <6571756174696F6E2E412E322E332E35>] >>}
\MusicPDFWriteObject{9750}{<< /Names [<6571756174696F6E2E412E322E332E36> \MusicPDFRef{14083} <6571756174696F6E2E412E322E332E37> \MusicPDFRef{14084} <6571756174696F6E2E412E322E332E38> \MusicPDFRef{14085} <6571756174696F6E2E412E322E332E39> \MusicPDFRef{14086} <6571756174696F6E2E412E322E342E31> \MusicPDFRef{14087} <6571756174696F6E2E412E322E342E3130> \MusicPDFRef{14088}] /Limits [<6571756174696F6E2E412E322E332E36> <6571756174696F6E2E412E322E342E3130>] >>}
\MusicPDFWriteObject{9751}{<< /Names [<6571756174696F6E2E412E322E342E3131> \MusicPDFRef{14089} <6571756174696F6E2E412E322E342E32> \MusicPDFRef{14090} <6571756174696F6E2E412E322E342E33> \MusicPDFRef{14091} <6571756174696F6E2E412E322E342E34> \MusicPDFRef{14092} <6571756174696F6E2E412E322E342E35> \MusicPDFRef{14093} <6571756174696F6E2E412E322E342E36> \MusicPDFRef{14094}] /Limits [<6571756174696F6E2E412E322E342E3131> <6571756174696F6E2E412E322E342E36>] >>}
\MusicPDFWriteObject{9752}{<< /Names [<6571756174696F6E2E412E322E342E37> \MusicPDFRef{14095} <6571756174696F6E2E412E322E342E38> \MusicPDFRef{14096} <6571756174696F6E2E412E322E342E39> \MusicPDFRef{14097} <6571756174696F6E2E412E332E312E31> \MusicPDFRef{14098} <6571756174696F6E2E412E332E312E3130> \MusicPDFRef{14099} <6571756174696F6E2E412E332E312E3131> \MusicPDFRef{14100}] /Limits [<6571756174696F6E2E412E322E342E37> <6571756174696F6E2E412E332E312E3131>] >>}
\MusicPDFWriteObject{9753}{<< /Names [<6571756174696F6E2E412E332E312E3132> \MusicPDFRef{14101} <6571756174696F6E2E412E332E312E3133> \MusicPDFRef{14102} <6571756174696F6E2E412E332E312E32> \MusicPDFRef{14103} <6571756174696F6E2E412E332E312E33> \MusicPDFRef{14104} <6571756174696F6E2E412E332E312E34> \MusicPDFRef{14105} <6571756174696F6E2E412E332E312E35> \MusicPDFRef{14106}] /Limits [<6571756174696F6E2E412E332E312E3132> <6571756174696F6E2E412E332E312E35>] >>}
\MusicPDFWriteObject{9754}{<< /Names [<6571756174696F6E2E412E332E312E36> \MusicPDFRef{14107} <6571756174696F6E2E412E332E312E37> \MusicPDFRef{14108} <6571756174696F6E2E412E332E312E38> \MusicPDFRef{14109} <6571756174696F6E2E412E332E312E39> \MusicPDFRef{14110} <6571756174696F6E2E412E332E322E31> \MusicPDFRef{14111} <6571756174696F6E2E412E332E322E32> \MusicPDFRef{14112}] /Limits [<6571756174696F6E2E412E332E312E36> <6571756174696F6E2E412E332E322E32>] >>}
\MusicPDFWriteObject{9755}{<< /Names [<6571756174696F6E2E412E332E322E33> \MusicPDFRef{14113} <6571756174696F6E2E412E332E322E34> \MusicPDFRef{14114} <6571756174696F6E2E412E332E322E35> \MusicPDFRef{14115} <6571756174696F6E2E412E332E342E31> \MusicPDFRef{14116} <6571756174696F6E2E412E332E342E32> \MusicPDFRef{14117} <6571756174696F6E2E412E342E312E31> \MusicPDFRef{14118}] /Limits [<6571756174696F6E2E412E332E322E33> <6571756174696F6E2E412E342E312E31>] >>}
\MusicPDFWriteObject{9756}{<< /Names [<6571756174696F6E2E412E342E312E32> \MusicPDFRef{14119} <6571756174696F6E2E412E342E312E33> \MusicPDFRef{14120} <6571756174696F6E2E412E342E312E34> \MusicPDFRef{14121} <6571756174696F6E2E412E342E312E35> \MusicPDFRef{14122} <6571756174696F6E2E412E342E312E36> \MusicPDFRef{14123} <6571756174696F6E2E412E342E312E37> \MusicPDFRef{14124}] /Limits [<6571756174696F6E2E412E342E312E32> <6571756174696F6E2E412E342E312E37>] >>}
\MusicPDFWriteObject{9757}{<< /Names [<6571756174696F6E2E412E342E312E38> \MusicPDFRef{14125} <6571756174696F6E2E412E342E312E39> \MusicPDFRef{14126} <6571756174696F6E2E412E352E312E31> \MusicPDFRef{14127} <6571756174696F6E2E412E352E312E3130> \MusicPDFRef{14128} <6571756174696F6E2E412E352E312E3131> \MusicPDFRef{14129} <6571756174696F6E2E412E352E312E3132> \MusicPDFRef{14130}] /Limits [<6571756174696F6E2E412E342E312E38> <6571756174696F6E2E412E352E312E3132>] >>}
\MusicPDFWriteObject{9758}{<< /Names [<6571756174696F6E2E412E352E312E3133> \MusicPDFRef{14131} <6571756174696F6E2E412E352E312E3134> \MusicPDFRef{14132} <6571756174696F6E2E412E352E312E32> \MusicPDFRef{14133} <6571756174696F6E2E412E352E312E33> \MusicPDFRef{14134} <6571756174696F6E2E412E352E312E34> \MusicPDFRef{14135} <6571756174696F6E2E412E352E312E35> \MusicPDFRef{14136}] /Limits [<6571756174696F6E2E412E352E312E3133> <6571756174696F6E2E412E352E312E35>] >>}
\MusicPDFWriteObject{9759}{<< /Names [<6571756174696F6E2E412E352E312E36> \MusicPDFRef{14137} <6571756174696F6E2E412E352E312E37> \MusicPDFRef{14138} <6571756174696F6E2E412E352E312E38> \MusicPDFRef{14139} <6571756174696F6E2E412E352E312E39> \MusicPDFRef{14140} <6571756174696F6E2E412E352E322E31> \MusicPDFRef{14141} <6571756174696F6E2E412E352E322E3130> \MusicPDFRef{14142}] /Limits [<6571756174696F6E2E412E352E312E36> <6571756174696F6E2E412E352E322E3130>] >>}
\MusicPDFWriteObject{9760}{<< /Names [<6571756174696F6E2E412E352E322E3131> \MusicPDFRef{14143} <6571756174696F6E2E412E352E322E3132> \MusicPDFRef{14144} <6571756174696F6E2E412E352E322E3133> \MusicPDFRef{14145} <6571756174696F6E2E412E352E322E3134> \MusicPDFRef{14146} <6571756174696F6E2E412E352E322E3135> \MusicPDFRef{14147} <6571756174696F6E2E412E352E322E3136> \MusicPDFRef{14148}] /Limits [<6571756174696F6E2E412E352E322E3131> <6571756174696F6E2E412E352E322E3136>] >>}
\MusicPDFWriteObject{9761}{<< /Names [<6571756174696F6E2E412E352E322E3137> \MusicPDFRef{14149} <6571756174696F6E2E412E352E322E3138> \MusicPDFRef{14150} <6571756174696F6E2E412E352E322E3139> \MusicPDFRef{14151} <6571756174696F6E2E412E352E322E32> \MusicPDFRef{14152} <6571756174696F6E2E412E352E322E3230> \MusicPDFRef{14153} <6571756174696F6E2E412E352E322E3231> \MusicPDFRef{14154}] /Limits [<6571756174696F6E2E412E352E322E3137> <6571756174696F6E2E412E352E322E3231>] >>}
\MusicPDFWriteObject{9762}{<< /Names [<6571756174696F6E2E412E352E322E3232> \MusicPDFRef{14155} <6571756174696F6E2E412E352E322E3233> \MusicPDFRef{14156} <6571756174696F6E2E412E352E322E3234> \MusicPDFRef{14157} <6571756174696F6E2E412E352E322E33> \MusicPDFRef{14158} <6571756174696F6E2E412E352E322E34> \MusicPDFRef{14159} <6571756174696F6E2E412E352E322E35> \MusicPDFRef{14160}] /Limits [<6571756174696F6E2E412E352E322E3232> <6571756174696F6E2E412E352E322E35>] >>}
\MusicPDFWriteObject{9763}{<< /Names [<6571756174696F6E2E412E352E322E36> \MusicPDFRef{14161} <6571756174696F6E2E412E352E322E37> \MusicPDFRef{14162} <6571756174696F6E2E412E352E322E38> \MusicPDFRef{14163} <6571756174696F6E2E412E352E322E39> \MusicPDFRef{14164} <6571756174696F6E2E412E352E332E31> \MusicPDFRef{14165} <6571756174696F6E2E412E352E332E3130> \MusicPDFRef{14166}] /Limits [<6571756174696F6E2E412E352E322E36> <6571756174696F6E2E412E352E332E3130>] >>}
\MusicPDFWriteObject{9764}{<< /Names [<6571756174696F6E2E412E352E332E3131> \MusicPDFRef{14167} <6571756174696F6E2E412E352E332E3132> \MusicPDFRef{14168} <6571756174696F6E2E412E352E332E3133> \MusicPDFRef{14169} <6571756174696F6E2E412E352E332E3134> \MusicPDFRef{14170} <6571756174696F6E2E412E352E332E3135> \MusicPDFRef{14171} <6571756174696F6E2E412E352E332E3136> \MusicPDFRef{14172}] /Limits [<6571756174696F6E2E412E352E332E3131> <6571756174696F6E2E412E352E332E3136>] >>}
\MusicPDFWriteObject{9765}{<< /Names [<6571756174696F6E2E412E352E332E3137> \MusicPDFRef{14173} <6571756174696F6E2E412E352E332E3138> \MusicPDFRef{14174} <6571756174696F6E2E412E352E332E32> \MusicPDFRef{14175} <6571756174696F6E2E412E352E332E33> \MusicPDFRef{14176} <6571756174696F6E2E412E352E332E34> \MusicPDFRef{14177} <6571756174696F6E2E412E352E332E35> \MusicPDFRef{14178}] /Limits [<6571756174696F6E2E412E352E332E3137> <6571756174696F6E2E412E352E332E35>] >>}
\MusicPDFWriteObject{9766}{<< /Names [<6571756174696F6E2E412E352E332E36> \MusicPDFRef{14179} <6571756174696F6E2E412E352E332E37> \MusicPDFRef{14180} <6571756174696F6E2E412E352E332E38> \MusicPDFRef{14181} <6571756174696F6E2E412E352E332E39> \MusicPDFRef{14182} <6571756174696F6E2E412E352E342E31> \MusicPDFRef{14183} <6571756174696F6E2E412E352E342E3130> \MusicPDFRef{14184}] /Limits [<6571756174696F6E2E412E352E332E36> <6571756174696F6E2E412E352E342E3130>] >>}
\MusicPDFWriteObject{9767}{<< /Names [<6571756174696F6E2E412E352E342E3131> \MusicPDFRef{14185} <6571756174696F6E2E412E352E342E3132> \MusicPDFRef{14186} <6571756174696F6E2E412E352E342E3133> \MusicPDFRef{14187} <6571756174696F6E2E412E352E342E3134> \MusicPDFRef{14188} <6571756174696F6E2E412E352E342E3135> \MusicPDFRef{14189} <6571756174696F6E2E412E352E342E3136> \MusicPDFRef{14190}] /Limits [<6571756174696F6E2E412E352E342E3131> <6571756174696F6E2E412E352E342E3136>] >>}
\MusicPDFWriteObject{9768}{<< /Names [<6571756174696F6E2E412E352E342E3137> \MusicPDFRef{14191} <6571756174696F6E2E412E352E342E32> \MusicPDFRef{14192} <6571756174696F6E2E412E352E342E33> \MusicPDFRef{14193} <6571756174696F6E2E412E352E342E34> \MusicPDFRef{14194} <6571756174696F6E2E412E352E342E35> \MusicPDFRef{14195} <6571756174696F6E2E412E352E342E36> \MusicPDFRef{14196}] /Limits [<6571756174696F6E2E412E352E342E3137> <6571756174696F6E2E412E352E342E36>] >>}
\MusicPDFWriteObject{9769}{<< /Names [<6571756174696F6E2E412E352E342E37> \MusicPDFRef{14197} <6571756174696F6E2E412E352E342E38> \MusicPDFRef{14198} <6571756174696F6E2E412E352E342E39> \MusicPDFRef{14199} <6571756174696F6E2E422E312E312E31> \MusicPDFRef{14200} <6571756174696F6E2E422E312E312E3130> \MusicPDFRef{14201} <6571756174696F6E2E422E312E312E3131> \MusicPDFRef{14202}] /Limits [<6571756174696F6E2E412E352E342E37> <6571756174696F6E2E422E312E312E3131>] >>}
\MusicPDFWriteObject{9770}{<< /Names [<6571756174696F6E2E422E312E312E3132> \MusicPDFRef{14203} <6571756174696F6E2E422E312E312E3133> \MusicPDFRef{14204} <6571756174696F6E2E422E312E312E3134> \MusicPDFRef{14205} <6571756174696F6E2E422E312E312E3135> \MusicPDFRef{14206} <6571756174696F6E2E422E312E312E32> \MusicPDFRef{14207} <6571756174696F6E2E422E312E312E33> \MusicPDFRef{14208}] /Limits [<6571756174696F6E2E422E312E312E3132> <6571756174696F6E2E422E312E312E33>] >>}
\MusicPDFWriteObject{9771}{<< /Names [<6571756174696F6E2E422E312E312E34> \MusicPDFRef{14209} <6571756174696F6E2E422E312E312E35> \MusicPDFRef{14210} <6571756174696F6E2E422E312E312E36> \MusicPDFRef{14211} <6571756174696F6E2E422E312E312E37> \MusicPDFRef{14212} <6571756174696F6E2E422E312E312E38> \MusicPDFRef{14213} <6571756174696F6E2E422E312E312E39> \MusicPDFRef{14214}] /Limits [<6571756174696F6E2E422E312E312E34> <6571756174696F6E2E422E312E312E39>] >>}
\MusicPDFWriteObject{9772}{<< /Names [<6571756174696F6E2E422E322E312E31> \MusicPDFRef{14215} <6571756174696F6E2E422E322E312E3130> \MusicPDFRef{14216} <6571756174696F6E2E422E322E312E3131> \MusicPDFRef{14217} <6571756174696F6E2E422E322E312E3132> \MusicPDFRef{14218} <6571756174696F6E2E422E322E312E3133> \MusicPDFRef{14219} <6571756174696F6E2E422E322E312E3134> \MusicPDFRef{14220}] /Limits [<6571756174696F6E2E422E322E312E31> <6571756174696F6E2E422E322E312E3134>] >>}
\MusicPDFWriteObject{9773}{<< /Names [<6571756174696F6E2E422E322E312E3135> \MusicPDFRef{14221} <6571756174696F6E2E422E322E312E3136> \MusicPDFRef{14222} <6571756174696F6E2E422E322E312E3137> \MusicPDFRef{14223} <6571756174696F6E2E422E322E312E3138> \MusicPDFRef{14224} <6571756174696F6E2E422E322E312E3139> \MusicPDFRef{14225} <6571756174696F6E2E422E322E312E32> \MusicPDFRef{14226}] /Limits [<6571756174696F6E2E422E322E312E3135> <6571756174696F6E2E422E322E312E32>] >>}
\MusicPDFWriteObject{9774}{<< /Names [<6571756174696F6E2E422E322E312E3230> \MusicPDFRef{14227} <6571756174696F6E2E422E322E312E3231> \MusicPDFRef{14228} <6571756174696F6E2E422E322E312E3232> \MusicPDFRef{14229} <6571756174696F6E2E422E322E312E3233> \MusicPDFRef{14230} <6571756174696F6E2E422E322E312E3234> \MusicPDFRef{14231} <6571756174696F6E2E422E322E312E3235> \MusicPDFRef{14232}] /Limits [<6571756174696F6E2E422E322E312E3230> <6571756174696F6E2E422E322E312E3235>] >>}
\MusicPDFWriteObject{9775}{<< /Names [<6571756174696F6E2E422E322E312E3236> \MusicPDFRef{14233} <6571756174696F6E2E422E322E312E3237> \MusicPDFRef{14234} <6571756174696F6E2E422E322E312E3238> \MusicPDFRef{14235} <6571756174696F6E2E422E322E312E3239> \MusicPDFRef{14236} <6571756174696F6E2E422E322E312E33> \MusicPDFRef{14237} <6571756174696F6E2E422E322E312E3330> \MusicPDFRef{14238}] /Limits [<6571756174696F6E2E422E322E312E3236> <6571756174696F6E2E422E322E312E3330>] >>}
\MusicPDFWriteObject{9776}{<< /Names [<6571756174696F6E2E422E322E312E3331> \MusicPDFRef{14239} <6571756174696F6E2E422E322E312E3332> \MusicPDFRef{14240} <6571756174696F6E2E422E322E312E3333> \MusicPDFRef{14241} <6571756174696F6E2E422E322E312E3334> \MusicPDFRef{14242} <6571756174696F6E2E422E322E312E3335> \MusicPDFRef{14243} <6571756174696F6E2E422E322E312E3336> \MusicPDFRef{14244}] /Limits [<6571756174696F6E2E422E322E312E3331> <6571756174696F6E2E422E322E312E3336>] >>}
\MusicPDFWriteObject{9777}{<< /Names [<6571756174696F6E2E422E322E312E3337> \MusicPDFRef{14245} <6571756174696F6E2E422E322E312E3338> \MusicPDFRef{14246} <6571756174696F6E2E422E322E312E3339> \MusicPDFRef{14247} <6571756174696F6E2E422E322E312E34> \MusicPDFRef{14248} <6571756174696F6E2E422E322E312E35> \MusicPDFRef{14249} <6571756174696F6E2E422E322E312E36> \MusicPDFRef{14250}] /Limits [<6571756174696F6E2E422E322E312E3337> <6571756174696F6E2E422E322E312E36>] >>}
\MusicPDFWriteObject{9778}{<< /Names [<6571756174696F6E2E422E322E312E37> \MusicPDFRef{14251} <6571756174696F6E2E422E322E312E38> \MusicPDFRef{14252} <6571756174696F6E2E422E322E312E39> \MusicPDFRef{14253} <6571756174696F6E2E422E322E322E31> \MusicPDFRef{14254} <6571756174696F6E2E422E322E322E3130> \MusicPDFRef{14255} <6571756174696F6E2E422E322E322E3131> \MusicPDFRef{14256}] /Limits [<6571756174696F6E2E422E322E312E37> <6571756174696F6E2E422E322E322E3131>] >>}
\MusicPDFWriteObject{9779}{<< /Names [<6571756174696F6E2E422E322E322E3132> \MusicPDFRef{14257} <6571756174696F6E2E422E322E322E3133> \MusicPDFRef{14258} <6571756174696F6E2E422E322E322E3134> \MusicPDFRef{14259} <6571756174696F6E2E422E322E322E32> \MusicPDFRef{14260} <6571756174696F6E2E422E322E322E33> \MusicPDFRef{14261} <6571756174696F6E2E422E322E322E34> \MusicPDFRef{14262}] /Limits [<6571756174696F6E2E422E322E322E3132> <6571756174696F6E2E422E322E322E34>] >>}
\MusicPDFWriteObject{9780}{<< /Names [<6571756174696F6E2E422E322E322E35> \MusicPDFRef{14263} <6571756174696F6E2E422E322E322E36> \MusicPDFRef{14264} <6571756174696F6E2E422E322E322E37> \MusicPDFRef{14265} <6571756174696F6E2E422E322E322E38> \MusicPDFRef{14266} <6571756174696F6E2E422E322E322E39> \MusicPDFRef{14267} <6571756174696F6E2E422E332E312E31> \MusicPDFRef{14268}] /Limits [<6571756174696F6E2E422E322E322E35> <6571756174696F6E2E422E332E312E31>] >>}
\MusicPDFWriteObject{9781}{<< /Names [<6571756174696F6E2E422E332E312E32> \MusicPDFRef{14269} <6571756174696F6E2E422E332E312E33> \MusicPDFRef{14270} <6571756174696F6E2E422E332E322E31> \MusicPDFRef{14271} <6571756174696F6E2E422E332E322E32> \MusicPDFRef{14272} <6571756174696F6E2E422E332E322E33> \MusicPDFRef{14273} <6571756174696F6E2E422E332E322E34> \MusicPDFRef{14274}] /Limits [<6571756174696F6E2E422E332E312E32> <6571756174696F6E2E422E332E322E34>] >>}
\MusicPDFWriteObject{9782}{<< /Names [<6571756174696F6E2E422E332E322E35> \MusicPDFRef{14275} <6571756174696F6E2E422E342E312E31> \MusicPDFRef{14276} <6571756174696F6E2E422E342E312E3130> \MusicPDFRef{14277} <6571756174696F6E2E422E342E312E3131> \MusicPDFRef{14278} <6571756174696F6E2E422E342E312E3132> \MusicPDFRef{14279} <6571756174696F6E2E422E342E312E3133> \MusicPDFRef{14280}] /Limits [<6571756174696F6E2E422E332E322E35> <6571756174696F6E2E422E342E312E3133>] >>}
\MusicPDFWriteObject{9783}{<< /Names [<6571756174696F6E2E422E342E312E3134> \MusicPDFRef{14281} <6571756174696F6E2E422E342E312E3135> \MusicPDFRef{14282} <6571756174696F6E2E422E342E312E3136> \MusicPDFRef{14283} <6571756174696F6E2E422E342E312E3137> \MusicPDFRef{14284} <6571756174696F6E2E422E342E312E3138> \MusicPDFRef{14285} <6571756174696F6E2E422E342E312E3139> \MusicPDFRef{14286}] /Limits [<6571756174696F6E2E422E342E312E3134> <6571756174696F6E2E422E342E312E3139>] >>}
\MusicPDFWriteObject{9784}{<< /Names [<6571756174696F6E2E422E342E312E32> \MusicPDFRef{14287} <6571756174696F6E2E422E342E312E3230> \MusicPDFRef{14288} <6571756174696F6E2E422E342E312E3231> \MusicPDFRef{14289} <6571756174696F6E2E422E342E312E3232> \MusicPDFRef{14290} <6571756174696F6E2E422E342E312E3233> \MusicPDFRef{14291} <6571756174696F6E2E422E342E312E3234> \MusicPDFRef{14292}] /Limits [<6571756174696F6E2E422E342E312E32> <6571756174696F6E2E422E342E312E3234>] >>}
\MusicPDFWriteObject{9785}{<< /Names [<6571756174696F6E2E422E342E312E3235> \MusicPDFRef{14293} <6571756174696F6E2E422E342E312E3236> \MusicPDFRef{14294} <6571756174696F6E2E422E342E312E3237> \MusicPDFRef{14295} <6571756174696F6E2E422E342E312E3238> \MusicPDFRef{14296} <6571756174696F6E2E422E342E312E3239> \MusicPDFRef{14297} <6571756174696F6E2E422E342E312E33> \MusicPDFRef{14298}] /Limits [<6571756174696F6E2E422E342E312E3235> <6571756174696F6E2E422E342E312E33>] >>}
\MusicPDFWriteObject{9786}{<< /Names [<6571756174696F6E2E422E342E312E3330> \MusicPDFRef{14299} <6571756174696F6E2E422E342E312E3331> \MusicPDFRef{14300} <6571756174696F6E2E422E342E312E3332> \MusicPDFRef{14301} <6571756174696F6E2E422E342E312E3333> \MusicPDFRef{14302} <6571756174696F6E2E422E342E312E3334> \MusicPDFRef{14303} <6571756174696F6E2E422E342E312E3335> \MusicPDFRef{14304}] /Limits [<6571756174696F6E2E422E342E312E3330> <6571756174696F6E2E422E342E312E3335>] >>}
\MusicPDFWriteObject{9787}{<< /Names [<6571756174696F6E2E422E342E312E34> \MusicPDFRef{14305} <6571756174696F6E2E422E342E312E35> \MusicPDFRef{14306} <6571756174696F6E2E422E342E312E36> \MusicPDFRef{14307} <6571756174696F6E2E422E342E312E37> \MusicPDFRef{14308} <6571756174696F6E2E422E342E312E38> \MusicPDFRef{14309} <6571756174696F6E2E422E342E312E39> \MusicPDFRef{14310}] /Limits [<6571756174696F6E2E422E342E312E34> <6571756174696F6E2E422E342E312E39>] >>}
\MusicPDFWriteObject{9788}{<< /Names [<6571756174696F6E2E422E342E322E31> \MusicPDFRef{14311} <6571756174696F6E2E422E342E322E3130> \MusicPDFRef{14312} <6571756174696F6E2E422E342E322E3131> \MusicPDFRef{14313} <6571756174696F6E2E422E342E322E3132> \MusicPDFRef{14314} <6571756174696F6E2E422E342E322E3133> \MusicPDFRef{14315} <6571756174696F6E2E422E342E322E3134> \MusicPDFRef{14316}] /Limits [<6571756174696F6E2E422E342E322E31> <6571756174696F6E2E422E342E322E3134>] >>}
\MusicPDFWriteObject{9789}{<< /Names [<6571756174696F6E2E422E342E322E3135> \MusicPDFRef{14317} <6571756174696F6E2E422E342E322E3136> \MusicPDFRef{14318} <6571756174696F6E2E422E342E322E3137> \MusicPDFRef{14319} <6571756174696F6E2E422E342E322E32> \MusicPDFRef{14320} <6571756174696F6E2E422E342E322E33> \MusicPDFRef{14321} <6571756174696F6E2E422E342E322E34> \MusicPDFRef{14322}] /Limits [<6571756174696F6E2E422E342E322E3135> <6571756174696F6E2E422E342E322E34>] >>}
\MusicPDFWriteObject{9790}{<< /Names [<6571756174696F6E2E422E342E322E35> \MusicPDFRef{14323} <6571756174696F6E2E422E342E322E36> \MusicPDFRef{14324} <6571756174696F6E2E422E342E322E37> \MusicPDFRef{14325} <6571756174696F6E2E422E342E322E38> \MusicPDFRef{14326} <6571756174696F6E2E422E342E322E39> \MusicPDFRef{14327} <6571756174696F6E2E422E342E332E31> \MusicPDFRef{14328}] /Limits [<6571756174696F6E2E422E342E322E35> <6571756174696F6E2E422E342E332E31>] >>}
\MusicPDFWriteObject{9791}{<< /Names [<6571756174696F6E2E422E342E332E3130> \MusicPDFRef{14329} <6571756174696F6E2E422E342E332E3131> \MusicPDFRef{14330} <6571756174696F6E2E422E342E332E3132> \MusicPDFRef{14331} <6571756174696F6E2E422E342E332E3133> \MusicPDFRef{14332} <6571756174696F6E2E422E342E332E3134> \MusicPDFRef{14333} <6571756174696F6E2E422E342E332E32> \MusicPDFRef{14334}] /Limits [<6571756174696F6E2E422E342E332E3130> <6571756174696F6E2E422E342E332E32>] >>}
\MusicPDFWriteObject{9792}{<< /Names [<6571756174696F6E2E422E342E332E33> \MusicPDFRef{14335} <6571756174696F6E2E422E342E332E34> \MusicPDFRef{14336} <6571756174696F6E2E422E342E332E35> \MusicPDFRef{14337} <6571756174696F6E2E422E342E332E36> \MusicPDFRef{14338} <6571756174696F6E2E422E342E332E37> \MusicPDFRef{14339} <6571756174696F6E2E422E342E332E38> \MusicPDFRef{14340}] /Limits [<6571756174696F6E2E422E342E332E33> <6571756174696F6E2E422E342E332E38>] >>}
\MusicPDFWriteObject{9793}{<< /Names [<6571756174696F6E2E422E342E332E39> \MusicPDFRef{14341} <6571756174696F6E2E422E342E342E31> \MusicPDFRef{14342} <6571756174696F6E2E422E342E342E3130> \MusicPDFRef{14343} <6571756174696F6E2E422E342E342E3131> \MusicPDFRef{14344} <6571756174696F6E2E422E342E342E3132> \MusicPDFRef{14345} <6571756174696F6E2E422E342E342E3133> \MusicPDFRef{14346}] /Limits [<6571756174696F6E2E422E342E332E39> <6571756174696F6E2E422E342E342E3133>] >>}
\MusicPDFWriteObject{9794}{<< /Names [<6571756174696F6E2E422E342E342E3134> \MusicPDFRef{14347} <6571756174696F6E2E422E342E342E3135> \MusicPDFRef{14348} <6571756174696F6E2E422E342E342E3136> \MusicPDFRef{14349} <6571756174696F6E2E422E342E342E3137> \MusicPDFRef{14350} <6571756174696F6E2E422E342E342E3138> \MusicPDFRef{14351} <6571756174696F6E2E422E342E342E3139> \MusicPDFRef{14352}] /Limits [<6571756174696F6E2E422E342E342E3134> <6571756174696F6E2E422E342E342E3139>] >>}
\MusicPDFWriteObject{9795}{<< /Names [<6571756174696F6E2E422E342E342E32> \MusicPDFRef{14353} <6571756174696F6E2E422E342E342E3230> \MusicPDFRef{14354} <6571756174696F6E2E422E342E342E3231> \MusicPDFRef{14355} <6571756174696F6E2E422E342E342E3232> \MusicPDFRef{14356} <6571756174696F6E2E422E342E342E33> \MusicPDFRef{14357} <6571756174696F6E2E422E342E342E34> \MusicPDFRef{14358}] /Limits [<6571756174696F6E2E422E342E342E32> <6571756174696F6E2E422E342E342E34>] >>}
\MusicPDFWriteObject{9796}{<< /Names [<6571756174696F6E2E422E342E342E35> \MusicPDFRef{14359} <6571756174696F6E2E422E342E342E36> \MusicPDFRef{14360} <6571756174696F6E2E422E342E342E37> \MusicPDFRef{14361} <6571756174696F6E2E422E342E342E38> \MusicPDFRef{14362} <6571756174696F6E2E422E342E342E39> \MusicPDFRef{14363} <6571756174696F6E2E422E352E312E31> \MusicPDFRef{14364}] /Limits [<6571756174696F6E2E422E342E342E35> <6571756174696F6E2E422E352E312E31>] >>}
\MusicPDFWriteObject{9797}{<< /Names [<6571756174696F6E2E422E352E312E32> \MusicPDFRef{14365} <6571756174696F6E2E422E352E312E33> \MusicPDFRef{14366} <6571756174696F6E2E422E352E322E31> \MusicPDFRef{14367} <6571756174696F6E2E422E352E322E3130> \MusicPDFRef{14368} <6571756174696F6E2E422E352E322E3131> \MusicPDFRef{14369} <6571756174696F6E2E422E352E322E3132> \MusicPDFRef{14370}] /Limits [<6571756174696F6E2E422E352E312E32> <6571756174696F6E2E422E352E322E3132>] >>}
\MusicPDFWriteObject{9798}{<< /Names [<6571756174696F6E2E422E352E322E3133> \MusicPDFRef{14371} <6571756174696F6E2E422E352E322E3134> \MusicPDFRef{14372} <6571756174696F6E2E422E352E322E3135> \MusicPDFRef{14373} <6571756174696F6E2E422E352E322E32> \MusicPDFRef{14374} <6571756174696F6E2E422E352E322E33> \MusicPDFRef{14375} <6571756174696F6E2E422E352E322E34> \MusicPDFRef{14376}] /Limits [<6571756174696F6E2E422E352E322E3133> <6571756174696F6E2E422E352E322E34>] >>}
\MusicPDFWriteObject{9799}{<< /Names [<6571756174696F6E2E422E352E322E35> \MusicPDFRef{14377} <6571756174696F6E2E422E352E322E36> \MusicPDFRef{14378} <6571756174696F6E2E422E352E322E37> \MusicPDFRef{14379} <6571756174696F6E2E422E352E322E38> \MusicPDFRef{14380} <6571756174696F6E2E422E352E322E39> \MusicPDFRef{14381} <6571756174696F6E2E422E352E332E31> \MusicPDFRef{14382}] /Limits [<6571756174696F6E2E422E352E322E35> <6571756174696F6E2E422E352E332E31>] >>}
\MusicPDFWriteObject{9800}{<< /Names [<6571756174696F6E2E422E352E332E3130> \MusicPDFRef{14383} <6571756174696F6E2E422E352E332E3131> \MusicPDFRef{14384} <6571756174696F6E2E422E352E332E3132> \MusicPDFRef{14385} <6571756174696F6E2E422E352E332E3133> \MusicPDFRef{14386} <6571756174696F6E2E422E352E332E3134> \MusicPDFRef{14387} <6571756174696F6E2E422E352E332E3135> \MusicPDFRef{14388}] /Limits [<6571756174696F6E2E422E352E332E3130> <6571756174696F6E2E422E352E332E3135>] >>}
\MusicPDFWriteObject{9801}{<< /Names [<6571756174696F6E2E422E352E332E32> \MusicPDFRef{14389} <6571756174696F6E2E422E352E332E33> \MusicPDFRef{14390} <6571756174696F6E2E422E352E332E34> \MusicPDFRef{14391} <6571756174696F6E2E422E352E332E35> \MusicPDFRef{14392} <6571756174696F6E2E422E352E332E36> \MusicPDFRef{14393} <6571756174696F6E2E422E352E332E37> \MusicPDFRef{14394}] /Limits [<6571756174696F6E2E422E352E332E32> <6571756174696F6E2E422E352E332E37>] >>}
\MusicPDFWriteObject{9802}{<< /Names [<6571756174696F6E2E422E352E332E38> \MusicPDFRef{14395} <6571756174696F6E2E422E352E332E39> \MusicPDFRef{14396} <6571756174696F6E2E422E362E312E31> \MusicPDFRef{14397} <6571756174696F6E2E422E362E312E3130> \MusicPDFRef{14398} <6571756174696F6E2E422E362E312E3131> \MusicPDFRef{14399} <6571756174696F6E2E422E362E312E3132> \MusicPDFRef{14400}] /Limits [<6571756174696F6E2E422E352E332E38> <6571756174696F6E2E422E362E312E3132>] >>}
\MusicPDFWriteObject{9803}{<< /Names [<6571756174696F6E2E422E362E312E3133> \MusicPDFRef{14401} <6571756174696F6E2E422E362E312E3134> \MusicPDFRef{14402} <6571756174696F6E2E422E362E312E3135> \MusicPDFRef{14403} <6571756174696F6E2E422E362E312E3136> \MusicPDFRef{14404} <6571756174696F6E2E422E362E312E3137> \MusicPDFRef{14405} <6571756174696F6E2E422E362E312E3138> \MusicPDFRef{14406}] /Limits [<6571756174696F6E2E422E362E312E3133> <6571756174696F6E2E422E362E312E3138>] >>}
\MusicPDFWriteObject{9804}{<< /Names [<6571756174696F6E2E422E362E312E3139> \MusicPDFRef{14407} <6571756174696F6E2E422E362E312E32> \MusicPDFRef{14408} <6571756174696F6E2E422E362E312E3230> \MusicPDFRef{14409} <6571756174696F6E2E422E362E312E3231> \MusicPDFRef{14410} <6571756174696F6E2E422E362E312E3232> \MusicPDFRef{14411} <6571756174696F6E2E422E362E312E3233> \MusicPDFRef{14412}] /Limits [<6571756174696F6E2E422E362E312E3139> <6571756174696F6E2E422E362E312E3233>] >>}
\MusicPDFWriteObject{9805}{<< /Names [<6571756174696F6E2E422E362E312E3234> \MusicPDFRef{14413} <6571756174696F6E2E422E362E312E3235> \MusicPDFRef{14414} <6571756174696F6E2E422E362E312E3236> \MusicPDFRef{14415} <6571756174696F6E2E422E362E312E3237> \MusicPDFRef{14416} <6571756174696F6E2E422E362E312E3238> \MusicPDFRef{14417} <6571756174696F6E2E422E362E312E3239> \MusicPDFRef{14418}] /Limits [<6571756174696F6E2E422E362E312E3234> <6571756174696F6E2E422E362E312E3239>] >>}
\MusicPDFWriteObject{9806}{<< /Names [<6571756174696F6E2E422E362E312E33> \MusicPDFRef{14419} <6571756174696F6E2E422E362E312E3330> \MusicPDFRef{14420} <6571756174696F6E2E422E362E312E3331> \MusicPDFRef{14421} <6571756174696F6E2E422E362E312E3332> \MusicPDFRef{14422} <6571756174696F6E2E422E362E312E3333> \MusicPDFRef{14423} <6571756174696F6E2E422E362E312E3334> \MusicPDFRef{14424}] /Limits [<6571756174696F6E2E422E362E312E33> <6571756174696F6E2E422E362E312E3334>] >>}
\MusicPDFWriteObject{9807}{<< /Names [<6571756174696F6E2E422E362E312E3335> \MusicPDFRef{14425} <6571756174696F6E2E422E362E312E3336> \MusicPDFRef{14426} <6571756174696F6E2E422E362E312E3337> \MusicPDFRef{14427} <6571756174696F6E2E422E362E312E3338> \MusicPDFRef{14428} <6571756174696F6E2E422E362E312E3339> \MusicPDFRef{14429} <6571756174696F6E2E422E362E312E34> \MusicPDFRef{14430}] /Limits [<6571756174696F6E2E422E362E312E3335> <6571756174696F6E2E422E362E312E34>] >>}
\MusicPDFWriteObject{9808}{<< /Names [<6571756174696F6E2E422E362E312E3430> \MusicPDFRef{14431} <6571756174696F6E2E422E362E312E3431> \MusicPDFRef{14432} <6571756174696F6E2E422E362E312E3432> \MusicPDFRef{14433} <6571756174696F6E2E422E362E312E3433> \MusicPDFRef{14434} <6571756174696F6E2E422E362E312E3434> \MusicPDFRef{14435} <6571756174696F6E2E422E362E312E3435> \MusicPDFRef{14436}] /Limits [<6571756174696F6E2E422E362E312E3430> <6571756174696F6E2E422E362E312E3435>] >>}
\MusicPDFWriteObject{9809}{<< /Names [<6571756174696F6E2E422E362E312E35> \MusicPDFRef{14437} <6571756174696F6E2E422E362E312E36> \MusicPDFRef{14438} <6571756174696F6E2E422E362E312E37> \MusicPDFRef{14439} <6571756174696F6E2E422E362E312E38> \MusicPDFRef{14440} <6571756174696F6E2E422E362E312E39> \MusicPDFRef{14441} <6571756174696F6E2E422E372E312E31> \MusicPDFRef{14442}] /Limits [<6571756174696F6E2E422E362E312E35> <6571756174696F6E2E422E372E312E31>] >>}
\MusicPDFWriteObject{9810}{<< /Names [<6571756174696F6E2E422E372E312E3130> \MusicPDFRef{14443} <6571756174696F6E2E422E372E312E3131> \MusicPDFRef{14444} <6571756174696F6E2E422E372E312E3132> \MusicPDFRef{14445} <6571756174696F6E2E422E372E312E3133> \MusicPDFRef{14446} <6571756174696F6E2E422E372E312E3134> \MusicPDFRef{14447} <6571756174696F6E2E422E372E312E3135> \MusicPDFRef{14448}] /Limits [<6571756174696F6E2E422E372E312E3130> <6571756174696F6E2E422E372E312E3135>] >>}
\MusicPDFWriteObject{9811}{<< /Names [<6571756174696F6E2E422E372E312E3136> \MusicPDFRef{14449} <6571756174696F6E2E422E372E312E3137> \MusicPDFRef{14450} <6571756174696F6E2E422E372E312E32> \MusicPDFRef{14451} <6571756174696F6E2E422E372E312E33> \MusicPDFRef{14452} <6571756174696F6E2E422E372E312E34> \MusicPDFRef{14453} <6571756174696F6E2E422E372E312E35> \MusicPDFRef{14454}] /Limits [<6571756174696F6E2E422E372E312E3136> <6571756174696F6E2E422E372E312E35>] >>}
\MusicPDFWriteObject{9812}{<< /Names [<6571756174696F6E2E422E372E312E36> \MusicPDFRef{14455} <6571756174696F6E2E422E372E312E37> \MusicPDFRef{14456} <6571756174696F6E2E422E372E312E38> \MusicPDFRef{14457} <6571756174696F6E2E422E372E312E39> \MusicPDFRef{14458} <6571756174696F6E2E422E372E322E31> \MusicPDFRef{14459} <6571756174696F6E2E422E372E322E3130> \MusicPDFRef{14460}] /Limits [<6571756174696F6E2E422E372E312E36> <6571756174696F6E2E422E372E322E3130>] >>}
\MusicPDFWriteObject{9813}{<< /Names [<6571756174696F6E2E422E372E322E3131> \MusicPDFRef{14461} <6571756174696F6E2E422E372E322E3132> \MusicPDFRef{14462} <6571756174696F6E2E422E372E322E3133> \MusicPDFRef{14463} <6571756174696F6E2E422E372E322E3134> \MusicPDFRef{14464} <6571756174696F6E2E422E372E322E3135> \MusicPDFRef{14465} <6571756174696F6E2E422E372E322E3136> \MusicPDFRef{14466}] /Limits [<6571756174696F6E2E422E372E322E3131> <6571756174696F6E2E422E372E322E3136>] >>}
\MusicPDFWriteObject{9814}{<< /Names [<6571756174696F6E2E422E372E322E3137> \MusicPDFRef{14467} <6571756174696F6E2E422E372E322E3138> \MusicPDFRef{14468} <6571756174696F6E2E422E372E322E3139> \MusicPDFRef{14469} <6571756174696F6E2E422E372E322E32> \MusicPDFRef{14470} <6571756174696F6E2E422E372E322E3230> \MusicPDFRef{14471} <6571756174696F6E2E422E372E322E3231> \MusicPDFRef{14472}] /Limits [<6571756174696F6E2E422E372E322E3137> <6571756174696F6E2E422E372E322E3231>] >>}
\MusicPDFWriteObject{9815}{<< /Names [<6571756174696F6E2E422E372E322E3232> \MusicPDFRef{14473} <6571756174696F6E2E422E372E322E3233> \MusicPDFRef{14474} <6571756174696F6E2E422E372E322E3234> \MusicPDFRef{14475} <6571756174696F6E2E422E372E322E3235> \MusicPDFRef{14476} <6571756174696F6E2E422E372E322E33> \MusicPDFRef{14477} <6571756174696F6E2E422E372E322E34> \MusicPDFRef{14478}] /Limits [<6571756174696F6E2E422E372E322E3232> <6571756174696F6E2E422E372E322E34>] >>}
\MusicPDFWriteObject{9816}{<< /Names [<6571756174696F6E2E422E372E322E35> \MusicPDFRef{14479} <6571756174696F6E2E422E372E322E36> \MusicPDFRef{14480} <6571756174696F6E2E422E372E322E37> \MusicPDFRef{14481} <6571756174696F6E2E422E372E322E38> \MusicPDFRef{14482} <6571756174696F6E2E422E372E322E39> \MusicPDFRef{14483} <6571756174696F6E2E432E312E322E31> \MusicPDFRef{14484}] /Limits [<6571756174696F6E2E422E372E322E35> <6571756174696F6E2E432E312E322E31>] >>}
\MusicPDFWriteObject{9817}{<< /Names [<6571756174696F6E2E432E312E352E31> \MusicPDFRef{14485} <6571756174696F6E2E432E312E352E32> \MusicPDFRef{14486} <6571756174696F6E2E432E312E352E33> \MusicPDFRef{14487} <6571756174696F6E2E432E312E352E34> \MusicPDFRef{14488} <6571756174696F6E2E432E322E312E31> \MusicPDFRef{14489} <6571756174696F6E2E432E322E312E32> \MusicPDFRef{14490}] /Limits [<6571756174696F6E2E432E312E352E31> <6571756174696F6E2E432E322E312E32>] >>}
\MusicPDFWriteObject{9818}{<< /Names [<6571756174696F6E2E432E322E312E33> \MusicPDFRef{14491} <6571756174696F6E2E432E322E312E34> \MusicPDFRef{14492} <6571756174696F6E2E432E322E322E31> \MusicPDFRef{14493} <6571756174696F6E2E432E322E322E32> \MusicPDFRef{14494} <6571756174696F6E2E432E322E322E33> \MusicPDFRef{14495} <6571756174696F6E2E432E322E322E34> \MusicPDFRef{14496}] /Limits [<6571756174696F6E2E432E322E312E33> <6571756174696F6E2E432E322E322E34>] >>}
\MusicPDFWriteObject{9819}{<< /Names [<6571756174696F6E2E432E322E322E35> \MusicPDFRef{14497} <6571756174696F6E2E432E322E322E36> \MusicPDFRef{14498} <6571756174696F6E2E432E322E332E31> \MusicPDFRef{14499} <6571756174696F6E2E432E322E332E3130> \MusicPDFRef{14500} <6571756174696F6E2E432E322E332E3131> \MusicPDFRef{14501} <6571756174696F6E2E432E322E332E3132> \MusicPDFRef{14502}] /Limits [<6571756174696F6E2E432E322E322E35> <6571756174696F6E2E432E322E332E3132>] >>}
\MusicPDFWriteObject{9820}{<< /Names [<6571756174696F6E2E432E322E332E3133> \MusicPDFRef{14503} <6571756174696F6E2E432E322E332E3134> \MusicPDFRef{14504} <6571756174696F6E2E432E322E332E3135> \MusicPDFRef{14505} <6571756174696F6E2E432E322E332E3136> \MusicPDFRef{14506} <6571756174696F6E2E432E322E332E3137> \MusicPDFRef{14507} <6571756174696F6E2E432E322E332E3138> \MusicPDFRef{14508}] /Limits [<6571756174696F6E2E432E322E332E3133> <6571756174696F6E2E432E322E332E3138>] >>}
\MusicPDFWriteObject{9821}{<< /Names [<6571756174696F6E2E432E322E332E3139> \MusicPDFRef{14509} <6571756174696F6E2E432E322E332E32> \MusicPDFRef{14510} <6571756174696F6E2E432E322E332E3230> \MusicPDFRef{14511} <6571756174696F6E2E432E322E332E3231> \MusicPDFRef{14512} <6571756174696F6E2E432E322E332E3232> \MusicPDFRef{14513} <6571756174696F6E2E432E322E332E3233> \MusicPDFRef{14514}] /Limits [<6571756174696F6E2E432E322E332E3139> <6571756174696F6E2E432E322E332E3233>] >>}
\MusicPDFWriteObject{9822}{<< /Names [<6571756174696F6E2E432E322E332E3234> \MusicPDFRef{14515} <6571756174696F6E2E432E322E332E3235> \MusicPDFRef{14516} <6571756174696F6E2E432E322E332E3236> \MusicPDFRef{14517} <6571756174696F6E2E432E322E332E3237> \MusicPDFRef{14518} <6571756174696F6E2E432E322E332E3238> \MusicPDFRef{14519} <6571756174696F6E2E432E322E332E3239> \MusicPDFRef{14520}] /Limits [<6571756174696F6E2E432E322E332E3234> <6571756174696F6E2E432E322E332E3239>] >>}
\MusicPDFWriteObject{9823}{<< /Names [<6571756174696F6E2E432E322E332E33> \MusicPDFRef{14521} <6571756174696F6E2E432E322E332E34> \MusicPDFRef{14522} <6571756174696F6E2E432E322E332E35> \MusicPDFRef{14523} <6571756174696F6E2E432E322E332E36> \MusicPDFRef{14524} <6571756174696F6E2E432E322E332E37> \MusicPDFRef{14525} <6571756174696F6E2E432E322E332E38> \MusicPDFRef{14526}] /Limits [<6571756174696F6E2E432E322E332E33> <6571756174696F6E2E432E322E332E38>] >>}
\MusicPDFWriteObject{9824}{<< /Names [<6571756174696F6E2E432E322E332E39> \MusicPDFRef{14527} <6571756174696F6E2E432E322E342E31> \MusicPDFRef{14528} <6571756174696F6E2E432E322E342E3130> \MusicPDFRef{14529} <6571756174696F6E2E432E322E342E32> \MusicPDFRef{14530} <6571756174696F6E2E432E322E342E33> \MusicPDFRef{14531} <6571756174696F6E2E432E322E342E34> \MusicPDFRef{14532}] /Limits [<6571756174696F6E2E432E322E332E39> <6571756174696F6E2E432E322E342E34>] >>}
\MusicPDFWriteObject{9825}{<< /Names [<6571756174696F6E2E432E322E342E35> \MusicPDFRef{14533} <6571756174696F6E2E432E322E342E36> \MusicPDFRef{14534} <6571756174696F6E2E432E322E342E37> \MusicPDFRef{14535} <6571756174696F6E2E432E322E342E38> \MusicPDFRef{14536} <6571756174696F6E2E432E322E342E39> \MusicPDFRef{14537} <6571756174696F6E2E432E332E312E31> \MusicPDFRef{14538}] /Limits [<6571756174696F6E2E432E322E342E35> <6571756174696F6E2E432E332E312E31>] >>}
\MusicPDFWriteObject{9826}{<< /Names [<6571756174696F6E2E432E332E312E32> \MusicPDFRef{14539} <6571756174696F6E2E432E332E312E33> \MusicPDFRef{14540} <6571756174696F6E2E432E332E322E31> \MusicPDFRef{14541} <6571756174696F6E2E432E332E322E32> \MusicPDFRef{14542} <6571756174696F6E2E432E332E322E33> \MusicPDFRef{14543} <6571756174696F6E2E432E332E322E34> \MusicPDFRef{14544}] /Limits [<6571756174696F6E2E432E332E312E32> <6571756174696F6E2E432E332E322E34>] >>}
\MusicPDFWriteObject{9827}{<< /Names [<6571756174696F6E2E432E332E322E35> \MusicPDFRef{14545} <6571756174696F6E2E432E332E322E36> \MusicPDFRef{14546} <6571756174696F6E2E432E332E322E37> \MusicPDFRef{14547} <6571756174696F6E2E432E332E332E31> \MusicPDFRef{14548} <6571756174696F6E2E432E332E332E32> \MusicPDFRef{14549} <6571756174696F6E2E432E332E342E31> \MusicPDFRef{14550}] /Limits [<6571756174696F6E2E432E332E322E35> <6571756174696F6E2E432E332E342E31>] >>}
\MusicPDFWriteObject{9828}{<< /Names [<6571756174696F6E2E432E332E342E32> \MusicPDFRef{14551} <6571756174696F6E2E432E332E342E33> \MusicPDFRef{14552} <6571756174696F6E2E432E332E352E31> \MusicPDFRef{14553} <6571756174696F6E2E432E332E352E3130> \MusicPDFRef{14554} <6571756174696F6E2E432E332E352E3131> \MusicPDFRef{14555} <6571756174696F6E2E432E332E352E3132> \MusicPDFRef{14556}] /Limits [<6571756174696F6E2E432E332E342E32> <6571756174696F6E2E432E332E352E3132>] >>}
\MusicPDFWriteObject{9829}{<< /Names [<6571756174696F6E2E432E332E352E3133> \MusicPDFRef{14557} <6571756174696F6E2E432E332E352E3134> \MusicPDFRef{14558} <6571756174696F6E2E432E332E352E3135> \MusicPDFRef{14559} <6571756174696F6E2E432E332E352E3136> \MusicPDFRef{14560} <6571756174696F6E2E432E332E352E3137> \MusicPDFRef{14561} <6571756174696F6E2E432E332E352E3138> \MusicPDFRef{14562}] /Limits [<6571756174696F6E2E432E332E352E3133> <6571756174696F6E2E432E332E352E3138>] >>}
\MusicPDFWriteObject{9830}{<< /Names [<6571756174696F6E2E432E332E352E3139> \MusicPDFRef{14563} <6571756174696F6E2E432E332E352E32> \MusicPDFRef{14564} <6571756174696F6E2E432E332E352E3230> \MusicPDFRef{14565} <6571756174696F6E2E432E332E352E3231> \MusicPDFRef{14566} <6571756174696F6E2E432E332E352E3232> \MusicPDFRef{14567} <6571756174696F6E2E432E332E352E3233> \MusicPDFRef{14568}] /Limits [<6571756174696F6E2E432E332E352E3139> <6571756174696F6E2E432E332E352E3233>] >>}
\MusicPDFWriteObject{9831}{<< /Names [<6571756174696F6E2E432E332E352E3234> \MusicPDFRef{14569} <6571756174696F6E2E432E332E352E33> \MusicPDFRef{14570} <6571756174696F6E2E432E332E352E34> \MusicPDFRef{14571} <6571756174696F6E2E432E332E352E35> \MusicPDFRef{14572} <6571756174696F6E2E432E332E352E36> \MusicPDFRef{14573} <6571756174696F6E2E432E332E352E37> \MusicPDFRef{14574}] /Limits [<6571756174696F6E2E432E332E352E3234> <6571756174696F6E2E432E332E352E37>] >>}
\MusicPDFWriteObject{9832}{<< /Names [<6571756174696F6E2E432E332E352E38> \MusicPDFRef{14575} <6571756174696F6E2E432E332E352E39> \MusicPDFRef{14576} <6571756174696F6E2E432E332E362E31> \MusicPDFRef{14577} <6571756174696F6E2E432E332E362E3130> \MusicPDFRef{14578} <6571756174696F6E2E432E332E362E3131> \MusicPDFRef{14579} <6571756174696F6E2E432E332E362E3132> \MusicPDFRef{14580}] /Limits [<6571756174696F6E2E432E332E352E38> <6571756174696F6E2E432E332E362E3132>] >>}
\MusicPDFWriteObject{9833}{<< /Names [<6571756174696F6E2E432E332E362E3133> \MusicPDFRef{14581} <6571756174696F6E2E432E332E362E3134> \MusicPDFRef{14582} <6571756174696F6E2E432E332E362E3135> \MusicPDFRef{14583} <6571756174696F6E2E432E332E362E3136> \MusicPDFRef{14584} <6571756174696F6E2E432E332E362E3137> \MusicPDFRef{14585} <6571756174696F6E2E432E332E362E3138> \MusicPDFRef{14586}] /Limits [<6571756174696F6E2E432E332E362E3133> <6571756174696F6E2E432E332E362E3138>] >>}
\MusicPDFWriteObject{9834}{<< /Names [<6571756174696F6E2E432E332E362E3139> \MusicPDFRef{14587} <6571756174696F6E2E432E332E362E32> \MusicPDFRef{14588} <6571756174696F6E2E432E332E362E33> \MusicPDFRef{14589} <6571756174696F6E2E432E332E362E34> \MusicPDFRef{14590} <6571756174696F6E2E432E332E362E35> \MusicPDFRef{14591} <6571756174696F6E2E432E332E362E36> \MusicPDFRef{14592}] /Limits [<6571756174696F6E2E432E332E362E3139> <6571756174696F6E2E432E332E362E36>] >>}
\MusicPDFWriteObject{9835}{<< /Names [<6571756174696F6E2E432E332E362E37> \MusicPDFRef{14593} <6571756174696F6E2E432E332E362E38> \MusicPDFRef{14594} <6571756174696F6E2E432E332E362E39> \MusicPDFRef{14595} <6571756174696F6E2E432E332E372E31> \MusicPDFRef{14596} <6571756174696F6E2E432E332E372E32> \MusicPDFRef{14597} <6571756174696F6E2E432E342E312E31> \MusicPDFRef{14598}] /Limits [<6571756174696F6E2E432E332E362E37> <6571756174696F6E2E432E342E312E31>] >>}
\MusicPDFWriteObject{9836}{<< /Names [<6571756174696F6E2E432E342E312E32> \MusicPDFRef{14599} <6571756174696F6E2E432E342E322E31> \MusicPDFRef{14600} <6571756174696F6E2E432E342E332E31> \MusicPDFRef{14601} <6571756174696F6E2E432E342E332E3130> \MusicPDFRef{14602} <6571756174696F6E2E432E342E332E3131> \MusicPDFRef{14603} <6571756174696F6E2E432E342E332E3132> \MusicPDFRef{14604}] /Limits [<6571756174696F6E2E432E342E312E32> <6571756174696F6E2E432E342E332E3132>] >>}
\MusicPDFWriteObject{9837}{<< /Names [<6571756174696F6E2E432E342E332E3133> \MusicPDFRef{14605} <6571756174696F6E2E432E342E332E3134> \MusicPDFRef{14606} <6571756174696F6E2E432E342E332E3135> \MusicPDFRef{14607} <6571756174696F6E2E432E342E332E3136> \MusicPDFRef{14608} <6571756174696F6E2E432E342E332E32> \MusicPDFRef{14609} <6571756174696F6E2E432E342E332E33> \MusicPDFRef{14610}] /Limits [<6571756174696F6E2E432E342E332E3133> <6571756174696F6E2E432E342E332E33>] >>}
\MusicPDFWriteObject{9838}{<< /Names [<6571756174696F6E2E432E342E332E34> \MusicPDFRef{14611} <6571756174696F6E2E432E342E332E35> \MusicPDFRef{14612} <6571756174696F6E2E432E342E332E36> \MusicPDFRef{14613} <6571756174696F6E2E432E342E332E37> \MusicPDFRef{14614} <6571756174696F6E2E432E342E332E38> \MusicPDFRef{14615} <6571756174696F6E2E432E342E332E39> \MusicPDFRef{14616}] /Limits [<6571756174696F6E2E432E342E332E34> <6571756174696F6E2E432E342E332E39>] >>}
\MusicPDFWriteObject{9839}{<< /Names [<6571756174696F6E2E432E342E342E31> \MusicPDFRef{14617} <6571756174696F6E2E432E342E342E3130> \MusicPDFRef{14618} <6571756174696F6E2E432E342E342E3131> \MusicPDFRef{14619} <6571756174696F6E2E432E342E342E3132> \MusicPDFRef{14620} <6571756174696F6E2E432E342E342E32> \MusicPDFRef{14621} <6571756174696F6E2E432E342E342E33> \MusicPDFRef{14622}] /Limits [<6571756174696F6E2E432E342E342E31> <6571756174696F6E2E432E342E342E33>] >>}
\MusicPDFWriteObject{9840}{<< /Names [<6571756174696F6E2E432E342E342E34> \MusicPDFRef{14623} <6571756174696F6E2E432E342E342E35> \MusicPDFRef{14624} <6571756174696F6E2E432E342E342E36> \MusicPDFRef{14625} <6571756174696F6E2E432E342E342E37> \MusicPDFRef{14626} <6571756174696F6E2E432E342E342E38> \MusicPDFRef{14627} <6571756174696F6E2E432E342E342E39> \MusicPDFRef{14628}] /Limits [<6571756174696F6E2E432E342E342E34> <6571756174696F6E2E432E342E342E39>] >>}
\MusicPDFWriteObject{9841}{<< /Names [<6571756174696F6E2E432E352E312E31> \MusicPDFRef{14629} <6571756174696F6E2E432E352E312E3130> \MusicPDFRef{14630} <6571756174696F6E2E432E352E312E3131> \MusicPDFRef{14631} <6571756174696F6E2E432E352E312E3132> \MusicPDFRef{14632} <6571756174696F6E2E432E352E312E3133> \MusicPDFRef{14633} <6571756174696F6E2E432E352E312E3134> \MusicPDFRef{14634}] /Limits [<6571756174696F6E2E432E352E312E31> <6571756174696F6E2E432E352E312E3134>] >>}
\MusicPDFWriteObject{9842}{<< /Names [<6571756174696F6E2E432E352E312E3135> \MusicPDFRef{14635} <6571756174696F6E2E432E352E312E3136> \MusicPDFRef{14636} <6571756174696F6E2E432E352E312E3137> \MusicPDFRef{14637} <6571756174696F6E2E432E352E312E3138> \MusicPDFRef{14638} <6571756174696F6E2E432E352E312E3139> \MusicPDFRef{14639} <6571756174696F6E2E432E352E312E32> \MusicPDFRef{14640}] /Limits [<6571756174696F6E2E432E352E312E3135> <6571756174696F6E2E432E352E312E32>] >>}
\MusicPDFWriteObject{9843}{<< /Names [<6571756174696F6E2E432E352E312E3230> \MusicPDFRef{14641} <6571756174696F6E2E432E352E312E3231> \MusicPDFRef{14642} <6571756174696F6E2E432E352E312E3232> \MusicPDFRef{14643} <6571756174696F6E2E432E352E312E3233> \MusicPDFRef{14644} <6571756174696F6E2E432E352E312E3234> \MusicPDFRef{14645} <6571756174696F6E2E432E352E312E3235> \MusicPDFRef{14646}] /Limits [<6571756174696F6E2E432E352E312E3230> <6571756174696F6E2E432E352E312E3235>] >>}
\MusicPDFWriteObject{9844}{<< /Names [<6571756174696F6E2E432E352E312E3236> \MusicPDFRef{14647} <6571756174696F6E2E432E352E312E3237> \MusicPDFRef{14648} <6571756174696F6E2E432E352E312E3238> \MusicPDFRef{14649} <6571756174696F6E2E432E352E312E3239> \MusicPDFRef{14650} <6571756174696F6E2E432E352E312E33> \MusicPDFRef{14651} <6571756174696F6E2E432E352E312E3330> \MusicPDFRef{14652}] /Limits [<6571756174696F6E2E432E352E312E3236> <6571756174696F6E2E432E352E312E3330>] >>}
\MusicPDFWriteObject{9845}{<< /Names [<6571756174696F6E2E432E352E312E3331> \MusicPDFRef{14653} <6571756174696F6E2E432E352E312E3332> \MusicPDFRef{14654} <6571756174696F6E2E432E352E312E3333> \MusicPDFRef{14655} <6571756174696F6E2E432E352E312E3334> \MusicPDFRef{14656} <6571756174696F6E2E432E352E312E3335> \MusicPDFRef{14657} <6571756174696F6E2E432E352E312E3336> \MusicPDFRef{14658}] /Limits [<6571756174696F6E2E432E352E312E3331> <6571756174696F6E2E432E352E312E3336>] >>}
\MusicPDFWriteObject{9846}{<< /Names [<6571756174696F6E2E432E352E312E3337> \MusicPDFRef{14659} <6571756174696F6E2E432E352E312E3338> \MusicPDFRef{14660} <6571756174696F6E2E432E352E312E3339> \MusicPDFRef{14661} <6571756174696F6E2E432E352E312E34> \MusicPDFRef{14662} <6571756174696F6E2E432E352E312E3430> \MusicPDFRef{14663} <6571756174696F6E2E432E352E312E3431> \MusicPDFRef{14664}] /Limits [<6571756174696F6E2E432E352E312E3337> <6571756174696F6E2E432E352E312E3431>] >>}
\MusicPDFWriteObject{9847}{<< /Names [<6571756174696F6E2E432E352E312E3432> \MusicPDFRef{14665} <6571756174696F6E2E432E352E312E3433> \MusicPDFRef{14666} <6571756174696F6E2E432E352E312E3434> \MusicPDFRef{14667} <6571756174696F6E2E432E352E312E3435> \MusicPDFRef{14668} <6571756174696F6E2E432E352E312E3436> \MusicPDFRef{14669} <6571756174696F6E2E432E352E312E3437> \MusicPDFRef{14670}] /Limits [<6571756174696F6E2E432E352E312E3432> <6571756174696F6E2E432E352E312E3437>] >>}
\MusicPDFWriteObject{9848}{<< /Names [<6571756174696F6E2E432E352E312E3438> \MusicPDFRef{14671} <6571756174696F6E2E432E352E312E35> \MusicPDFRef{14672} <6571756174696F6E2E432E352E312E36> \MusicPDFRef{14673} <6571756174696F6E2E432E352E312E37> \MusicPDFRef{14674} <6571756174696F6E2E432E352E312E38> \MusicPDFRef{14675} <6571756174696F6E2E432E352E312E39> \MusicPDFRef{14676}] /Limits [<6571756174696F6E2E432E352E312E3438> <6571756174696F6E2E432E352E312E39>] >>}
\MusicPDFWriteObject{9849}{<< /Names [<6571756174696F6E2E432E352E322E31> \MusicPDFRef{14677} <6571756174696F6E2E432E352E322E32> \MusicPDFRef{14678} <6571756174696F6E2E432E352E322E33> \MusicPDFRef{14679} <6571756174696F6E2E432E352E322E34> \MusicPDFRef{14680} <6571756174696F6E2E432E352E322E35> \MusicPDFRef{14681} <6571756174696F6E2E432E352E322E36> \MusicPDFRef{14682}] /Limits [<6571756174696F6E2E432E352E322E31> <6571756174696F6E2E432E352E322E36>] >>}
\MusicPDFWriteObject{9850}{<< /Names [<6571756174696F6E2E432E352E332E31> \MusicPDFRef{14683} <6571756174696F6E2E432E352E332E32> \MusicPDFRef{14684} <6571756174696F6E2E432E352E332E33> \MusicPDFRef{14685} <6571756174696F6E2E432E352E332E34> \MusicPDFRef{14686} <6571756174696F6E2E432E352E332E35> \MusicPDFRef{14687} <6571756174696F6E2E432E352E332E36> \MusicPDFRef{14688}] /Limits [<6571756174696F6E2E432E352E332E31> <6571756174696F6E2E432E352E332E36>] >>}
\MusicPDFWriteObject{9851}{<< /Names [<6571756174696F6E2E432E352E332E37> \MusicPDFRef{14689} <6571756174696F6E2E432E352E332E38> \MusicPDFRef{14690} <6571756174696F6E2E432E352E332E39> \MusicPDFRef{14691} <6571756174696F6E2E432E352E342E31> \MusicPDFRef{14692} <6571756174696F6E2E432E352E342E32> \MusicPDFRef{14693} <6571756174696F6E2E442E312E312E31> \MusicPDFRef{14694}] /Limits [<6571756174696F6E2E432E352E332E37> <6571756174696F6E2E442E312E312E31>] >>}
\MusicPDFWriteObject{9852}{<< /Names [<6571756174696F6E2E442E312E312E3130> \MusicPDFRef{14695} <6571756174696F6E2E442E312E312E3131> \MusicPDFRef{14696} <6571756174696F6E2E442E312E312E3132> \MusicPDFRef{14697} <6571756174696F6E2E442E312E312E32> \MusicPDFRef{14698} <6571756174696F6E2E442E312E312E33> \MusicPDFRef{14699} <6571756174696F6E2E442E312E312E34> \MusicPDFRef{14700}] /Limits [<6571756174696F6E2E442E312E312E3130> <6571756174696F6E2E442E312E312E34>] >>}
\MusicPDFWriteObject{9853}{<< /Names [<6571756174696F6E2E442E312E312E35> \MusicPDFRef{14701} <6571756174696F6E2E442E312E312E36> \MusicPDFRef{14702} <6571756174696F6E2E442E312E312E37> \MusicPDFRef{14703} <6571756174696F6E2E442E312E312E38> \MusicPDFRef{14704} <6571756174696F6E2E442E312E312E39> \MusicPDFRef{14705} <6571756174696F6E2E442E312E31302E31> \MusicPDFRef{14706}] /Limits [<6571756174696F6E2E442E312E312E35> <6571756174696F6E2E442E312E31302E31>] >>}
\MusicPDFWriteObject{9854}{<< /Names [<6571756174696F6E2E442E312E31302E3130> \MusicPDFRef{14707} <6571756174696F6E2E442E312E31302E3131> \MusicPDFRef{14708} <6571756174696F6E2E442E312E31302E3132> \MusicPDFRef{14709} <6571756174696F6E2E442E312E31302E3133> \MusicPDFRef{14710} <6571756174696F6E2E442E312E31302E3134> \MusicPDFRef{14711} <6571756174696F6E2E442E312E31302E3135> \MusicPDFRef{14712}] /Limits [<6571756174696F6E2E442E312E31302E3130> <6571756174696F6E2E442E312E31302E3135>] >>}
\MusicPDFWriteObject{9855}{<< /Names [<6571756174696F6E2E442E312E31302E3136> \MusicPDFRef{14713} <6571756174696F6E2E442E312E31302E3137> \MusicPDFRef{14714} <6571756174696F6E2E442E312E31302E3138> \MusicPDFRef{14715} <6571756174696F6E2E442E312E31302E3139> \MusicPDFRef{14716} <6571756174696F6E2E442E312E31302E32> \MusicPDFRef{14717} <6571756174696F6E2E442E312E31302E3230> \MusicPDFRef{14718}] /Limits [<6571756174696F6E2E442E312E31302E3136> <6571756174696F6E2E442E312E31302E3230>] >>}
\MusicPDFWriteObject{9856}{<< /Names [<6571756174696F6E2E442E312E31302E33> \MusicPDFRef{14719} <6571756174696F6E2E442E312E31302E34> \MusicPDFRef{14720} <6571756174696F6E2E442E312E31302E35> \MusicPDFRef{14721} <6571756174696F6E2E442E312E31302E36> \MusicPDFRef{14722} <6571756174696F6E2E442E312E31302E37> \MusicPDFRef{14723} <6571756174696F6E2E442E312E31302E38> \MusicPDFRef{14724}] /Limits [<6571756174696F6E2E442E312E31302E33> <6571756174696F6E2E442E312E31302E38>] >>}
\MusicPDFWriteObject{9857}{<< /Names [<6571756174696F6E2E442E312E31302E39> \MusicPDFRef{14725} <6571756174696F6E2E442E312E322E31> \MusicPDFRef{14726} <6571756174696F6E2E442E312E322E32> \MusicPDFRef{14727} <6571756174696F6E2E442E312E322E33> \MusicPDFRef{14728} <6571756174696F6E2E442E312E322E34> \MusicPDFRef{14729} <6571756174696F6E2E442E312E322E35> \MusicPDFRef{14730}] /Limits [<6571756174696F6E2E442E312E31302E39> <6571756174696F6E2E442E312E322E35>] >>}
\MusicPDFWriteObject{9858}{<< /Names [<6571756174696F6E2E442E312E322E36> \MusicPDFRef{14731} <6571756174696F6E2E442E312E322E37> \MusicPDFRef{14732} <6571756174696F6E2E442E312E322E38> \MusicPDFRef{14733} <6571756174696F6E2E442E312E332E31> \MusicPDFRef{14734} <6571756174696F6E2E442E312E332E3130> \MusicPDFRef{14735} <6571756174696F6E2E442E312E332E3131> \MusicPDFRef{14736}] /Limits [<6571756174696F6E2E442E312E322E36> <6571756174696F6E2E442E312E332E3131>] >>}
\MusicPDFWriteObject{9859}{<< /Names [<6571756174696F6E2E442E312E332E3132> \MusicPDFRef{14737} <6571756174696F6E2E442E312E332E3133> \MusicPDFRef{14738} <6571756174696F6E2E442E312E332E3134> \MusicPDFRef{14739} <6571756174696F6E2E442E312E332E3135> \MusicPDFRef{14740} <6571756174696F6E2E442E312E332E3136> \MusicPDFRef{14741} <6571756174696F6E2E442E312E332E3137> \MusicPDFRef{14742}] /Limits [<6571756174696F6E2E442E312E332E3132> <6571756174696F6E2E442E312E332E3137>] >>}
\MusicPDFWriteObject{9860}{<< /Names [<6571756174696F6E2E442E312E332E3138> \MusicPDFRef{14743} <6571756174696F6E2E442E312E332E3139> \MusicPDFRef{14744} <6571756174696F6E2E442E312E332E32> \MusicPDFRef{14745} <6571756174696F6E2E442E312E332E3230> \MusicPDFRef{14746} <6571756174696F6E2E442E312E332E3231> \MusicPDFRef{14747} <6571756174696F6E2E442E312E332E3232> \MusicPDFRef{14748}] /Limits [<6571756174696F6E2E442E312E332E3138> <6571756174696F6E2E442E312E332E3232>] >>}
\MusicPDFWriteObject{9861}{<< /Names [<6571756174696F6E2E442E312E332E3233> \MusicPDFRef{14749} <6571756174696F6E2E442E312E332E3234> \MusicPDFRef{14750} <6571756174696F6E2E442E312E332E3235> \MusicPDFRef{14751} <6571756174696F6E2E442E312E332E3236> \MusicPDFRef{14752} <6571756174696F6E2E442E312E332E3237> \MusicPDFRef{14753} <6571756174696F6E2E442E312E332E3238> \MusicPDFRef{14754}] /Limits [<6571756174696F6E2E442E312E332E3233> <6571756174696F6E2E442E312E332E3238>] >>}
\MusicPDFWriteObject{9862}{<< /Names [<6571756174696F6E2E442E312E332E3239> \MusicPDFRef{14755} <6571756174696F6E2E442E312E332E33> \MusicPDFRef{14756} <6571756174696F6E2E442E312E332E3330> \MusicPDFRef{14757} <6571756174696F6E2E442E312E332E3331> \MusicPDFRef{14758} <6571756174696F6E2E442E312E332E3332> \MusicPDFRef{14759} <6571756174696F6E2E442E312E332E3333> \MusicPDFRef{14760}] /Limits [<6571756174696F6E2E442E312E332E3239> <6571756174696F6E2E442E312E332E3333>] >>}
\MusicPDFWriteObject{9863}{<< /Names [<6571756174696F6E2E442E312E332E3334> \MusicPDFRef{14761} <6571756174696F6E2E442E312E332E3335> \MusicPDFRef{14762} <6571756174696F6E2E442E312E332E34> \MusicPDFRef{14763} <6571756174696F6E2E442E312E332E35> \MusicPDFRef{14764} <6571756174696F6E2E442E312E332E36> \MusicPDFRef{14765} <6571756174696F6E2E442E312E332E37> \MusicPDFRef{14766}] /Limits [<6571756174696F6E2E442E312E332E3334> <6571756174696F6E2E442E312E332E37>] >>}
\MusicPDFWriteObject{9864}{<< /Names [<6571756174696F6E2E442E312E332E38> \MusicPDFRef{14767} <6571756174696F6E2E442E312E332E39> \MusicPDFRef{14768} <6571756174696F6E2E442E312E342E31> \MusicPDFRef{14769} <6571756174696F6E2E442E312E342E3130> \MusicPDFRef{14770} <6571756174696F6E2E442E312E342E3131> \MusicPDFRef{14771} <6571756174696F6E2E442E312E342E3132> \MusicPDFRef{14772}] /Limits [<6571756174696F6E2E442E312E332E38> <6571756174696F6E2E442E312E342E3132>] >>}
\MusicPDFWriteObject{9865}{<< /Names [<6571756174696F6E2E442E312E342E3133> \MusicPDFRef{14773} <6571756174696F6E2E442E312E342E3134> \MusicPDFRef{14774} <6571756174696F6E2E442E312E342E3135> \MusicPDFRef{14775} <6571756174696F6E2E442E312E342E3136> \MusicPDFRef{14776} <6571756174696F6E2E442E312E342E32> \MusicPDFRef{14777} <6571756174696F6E2E442E312E342E33> \MusicPDFRef{14778}] /Limits [<6571756174696F6E2E442E312E342E3133> <6571756174696F6E2E442E312E342E33>] >>}
\MusicPDFWriteObject{9866}{<< /Names [<6571756174696F6E2E442E312E342E34> \MusicPDFRef{14779} <6571756174696F6E2E442E312E342E35> \MusicPDFRef{14780} <6571756174696F6E2E442E312E342E36> \MusicPDFRef{14781} <6571756174696F6E2E442E312E342E37> \MusicPDFRef{14782} <6571756174696F6E2E442E312E342E38> \MusicPDFRef{14783} <6571756174696F6E2E442E312E342E39> \MusicPDFRef{14784}] /Limits [<6571756174696F6E2E442E312E342E34> <6571756174696F6E2E442E312E342E39>] >>}
\MusicPDFWriteObject{9867}{<< /Names [<6571756174696F6E2E442E312E352E31> \MusicPDFRef{14785} <6571756174696F6E2E442E312E352E32> \MusicPDFRef{14786} <6571756174696F6E2E442E312E352E33> \MusicPDFRef{14787} <6571756174696F6E2E442E312E352E34> \MusicPDFRef{14788} <6571756174696F6E2E442E312E352E35> \MusicPDFRef{14789} <6571756174696F6E2E442E312E352E36> \MusicPDFRef{14790}] /Limits [<6571756174696F6E2E442E312E352E31> <6571756174696F6E2E442E312E352E36>] >>}
\MusicPDFWriteObject{9868}{<< /Names [<6571756174696F6E2E442E312E362E31> \MusicPDFRef{14791} <6571756174696F6E2E442E312E362E3130> \MusicPDFRef{14792} <6571756174696F6E2E442E312E362E3131> \MusicPDFRef{14793} <6571756174696F6E2E442E312E362E3132> \MusicPDFRef{14794} <6571756174696F6E2E442E312E362E3133> \MusicPDFRef{14795} <6571756174696F6E2E442E312E362E3134> \MusicPDFRef{14796}] /Limits [<6571756174696F6E2E442E312E362E31> <6571756174696F6E2E442E312E362E3134>] >>}
\MusicPDFWriteObject{9869}{<< /Names [<6571756174696F6E2E442E312E362E3135> \MusicPDFRef{14797} <6571756174696F6E2E442E312E362E3136> \MusicPDFRef{14798} <6571756174696F6E2E442E312E362E3137> \MusicPDFRef{14799} <6571756174696F6E2E442E312E362E3138> \MusicPDFRef{14800} <6571756174696F6E2E442E312E362E3139> \MusicPDFRef{14801} <6571756174696F6E2E442E312E362E32> \MusicPDFRef{14802}] /Limits [<6571756174696F6E2E442E312E362E3135> <6571756174696F6E2E442E312E362E32>] >>}
\MusicPDFWriteObject{9870}{<< /Names [<6571756174696F6E2E442E312E362E3230> \MusicPDFRef{14803} <6571756174696F6E2E442E312E362E3231> \MusicPDFRef{14804} <6571756174696F6E2E442E312E362E3232> \MusicPDFRef{14805} <6571756174696F6E2E442E312E362E3233> \MusicPDFRef{14806} <6571756174696F6E2E442E312E362E3234> \MusicPDFRef{14807} <6571756174696F6E2E442E312E362E3235> \MusicPDFRef{14808}] /Limits [<6571756174696F6E2E442E312E362E3230> <6571756174696F6E2E442E312E362E3235>] >>}
\MusicPDFWriteObject{9871}{<< /Names [<6571756174696F6E2E442E312E362E3236> \MusicPDFRef{14809} <6571756174696F6E2E442E312E362E33> \MusicPDFRef{14810} <6571756174696F6E2E442E312E362E34> \MusicPDFRef{14811} <6571756174696F6E2E442E312E362E35> \MusicPDFRef{14812} <6571756174696F6E2E442E312E362E36> \MusicPDFRef{14813} <6571756174696F6E2E442E312E362E37> \MusicPDFRef{14814}] /Limits [<6571756174696F6E2E442E312E362E3236> <6571756174696F6E2E442E312E362E37>] >>}
\MusicPDFWriteObject{9872}{<< /Names [<6571756174696F6E2E442E312E362E38> \MusicPDFRef{14815} <6571756174696F6E2E442E312E362E39> \MusicPDFRef{14816} <6571756174696F6E2E442E312E372E31> \MusicPDFRef{14817} <6571756174696F6E2E442E312E372E32> \MusicPDFRef{14818} <6571756174696F6E2E442E312E372E33> \MusicPDFRef{14819} <6571756174696F6E2E442E312E382E31> \MusicPDFRef{14820}] /Limits [<6571756174696F6E2E442E312E362E38> <6571756174696F6E2E442E312E382E31>] >>}
\MusicPDFWriteObject{9873}{<< /Names [<6571756174696F6E2E442E312E382E32> \MusicPDFRef{14821} <6571756174696F6E2E442E312E382E33> \MusicPDFRef{14822} <6571756174696F6E2E442E312E382E34> \MusicPDFRef{14823} <6571756174696F6E2E442E312E392E31> \MusicPDFRef{14824} <6571756174696F6E2E442E312E392E32> \MusicPDFRef{14825} <6571756174696F6E2E442E322E312E31> \MusicPDFRef{14826}] /Limits [<6571756174696F6E2E442E312E382E32> <6571756174696F6E2E442E322E312E31>] >>}
\MusicPDFWriteObject{9874}{<< /Names [<6571756174696F6E2E442E322E312E32> \MusicPDFRef{14827} <6571756174696F6E2E442E322E312E33> \MusicPDFRef{14828} <6571756174696F6E2E442E322E312E34> \MusicPDFRef{14829} <6571756174696F6E2E442E322E312E35> \MusicPDFRef{14830} <6571756174696F6E2E442E322E312E36> \MusicPDFRef{14831} <6571756174696F6E2E442E322E312E37> \MusicPDFRef{14832}] /Limits [<6571756174696F6E2E442E322E312E32> <6571756174696F6E2E442E322E312E37>] >>}
\MusicPDFWriteObject{9875}{<< /Names [<6571756174696F6E2E442E322E322E31> \MusicPDFRef{14833} <6571756174696F6E2E442E322E322E3130> \MusicPDFRef{14834} <6571756174696F6E2E442E322E322E32> \MusicPDFRef{14835} <6571756174696F6E2E442E322E322E33> \MusicPDFRef{14836} <6571756174696F6E2E442E322E322E34> \MusicPDFRef{14837} <6571756174696F6E2E442E322E322E35> \MusicPDFRef{14838}] /Limits [<6571756174696F6E2E442E322E322E31> <6571756174696F6E2E442E322E322E35>] >>}
\MusicPDFWriteObject{9876}{<< /Names [<6571756174696F6E2E442E322E322E36> \MusicPDFRef{14839} <6571756174696F6E2E442E322E322E37> \MusicPDFRef{14840} <6571756174696F6E2E442E322E322E38> \MusicPDFRef{14841} <6571756174696F6E2E442E322E322E39> \MusicPDFRef{14842} <6571756174696F6E2E442E332E312E31> \MusicPDFRef{14843} <6571756174696F6E2E442E332E312E3130> \MusicPDFRef{14844}] /Limits [<6571756174696F6E2E442E322E322E36> <6571756174696F6E2E442E332E312E3130>] >>}
\MusicPDFWriteObject{9877}{<< /Names [<6571756174696F6E2E442E332E312E3131> \MusicPDFRef{14845} <6571756174696F6E2E442E332E312E3132> \MusicPDFRef{14846} <6571756174696F6E2E442E332E312E32> \MusicPDFRef{14847} <6571756174696F6E2E442E332E312E33> \MusicPDFRef{14848} <6571756174696F6E2E442E332E312E34> \MusicPDFRef{14849} <6571756174696F6E2E442E332E312E35> \MusicPDFRef{14850}] /Limits [<6571756174696F6E2E442E332E312E3131> <6571756174696F6E2E442E332E312E35>] >>}
\MusicPDFWriteObject{9878}{<< /Names [<6571756174696F6E2E442E332E312E36> \MusicPDFRef{14851} <6571756174696F6E2E442E332E312E37> \MusicPDFRef{14852} <6571756174696F6E2E442E332E312E38> \MusicPDFRef{14853} <6571756174696F6E2E442E332E312E39> \MusicPDFRef{14854} <6571756174696F6E2E442E332E322E31> \MusicPDFRef{14855} <6571756174696F6E2E442E332E322E32> \MusicPDFRef{14856}] /Limits [<6571756174696F6E2E442E332E312E36> <6571756174696F6E2E442E332E322E32>] >>}
\MusicPDFWriteObject{9879}{<< /Names [<6571756174696F6E2E442E332E322E33> \MusicPDFRef{14857} <6571756174696F6E2E442E332E322E34> \MusicPDFRef{14858} <6571756174696F6E2E442E332E322E35> \MusicPDFRef{14859} <6571756174696F6E2E442E332E322E36> \MusicPDFRef{14860} <6571756174696F6E2E442E332E322E37> \MusicPDFRef{14861} <6571756174696F6E2E442E332E332E31> \MusicPDFRef{14862}] /Limits [<6571756174696F6E2E442E332E322E33> <6571756174696F6E2E442E332E332E31>] >>}
\MusicPDFWriteObject{9880}{<< /Names [<6571756174696F6E2E442E332E332E3130> \MusicPDFRef{14863} <6571756174696F6E2E442E332E332E3131> \MusicPDFRef{14864} <6571756174696F6E2E442E332E332E3132> \MusicPDFRef{14865} <6571756174696F6E2E442E332E332E3133> \MusicPDFRef{14866} <6571756174696F6E2E442E332E332E3134> \MusicPDFRef{14867} <6571756174696F6E2E442E332E332E32> \MusicPDFRef{14868}] /Limits [<6571756174696F6E2E442E332E332E3130> <6571756174696F6E2E442E332E332E32>] >>}
\MusicPDFWriteObject{9881}{<< /Names [<6571756174696F6E2E442E332E332E33> \MusicPDFRef{14869} <6571756174696F6E2E442E332E332E34> \MusicPDFRef{14870} <6571756174696F6E2E442E332E332E35> \MusicPDFRef{14871} <6571756174696F6E2E442E332E332E36> \MusicPDFRef{14872} <6571756174696F6E2E442E332E332E37> \MusicPDFRef{14873} <6571756174696F6E2E442E332E332E38> \MusicPDFRef{14874}] /Limits [<6571756174696F6E2E442E332E332E33> <6571756174696F6E2E442E332E332E38>] >>}
\MusicPDFWriteObject{9882}{<< /Names [<6571756174696F6E2E442E332E332E39> \MusicPDFRef{14875} <6571756174696F6E2E442E332E342E31> \MusicPDFRef{14876} <6571756174696F6E2E442E332E342E32> \MusicPDFRef{14877} <6571756174696F6E2E442E332E352E31> \MusicPDFRef{14878} <6571756174696F6E2E442E332E352E32> \MusicPDFRef{14879} <6571756174696F6E2E442E332E352E33> \MusicPDFRef{14880}] /Limits [<6571756174696F6E2E442E332E332E39> <6571756174696F6E2E442E332E352E33>] >>}
\MusicPDFWriteObject{9883}{<< /Names [<6571756174696F6E2E442E332E352E34> \MusicPDFRef{14881} <6571756174696F6E2E442E332E352E35> \MusicPDFRef{14882} <6571756174696F6E2E442E332E362E31> \MusicPDFRef{14883} <6571756174696F6E2E442E332E362E32> \MusicPDFRef{14884} <6571756174696F6E2E442E332E362E33> \MusicPDFRef{14885} <6571756174696F6E2E442E332E362E34> \MusicPDFRef{14886}] /Limits [<6571756174696F6E2E442E332E352E34> <6571756174696F6E2E442E332E362E34>] >>}
\MusicPDFWriteObject{9884}{<< /Names [<6571756174696F6E2E452E312E312E31> \MusicPDFRef{14887} <6571756174696F6E2E452E312E312E32> \MusicPDFRef{14888} <6571756174696F6E2E452E312E312E33> \MusicPDFRef{14889} <6571756174696F6E2E452E312E312E34> \MusicPDFRef{14890} <6571756174696F6E2E452E312E312E35> \MusicPDFRef{14891} <6571756174696F6E2E452E312E322E31> \MusicPDFRef{14892}] /Limits [<6571756174696F6E2E452E312E312E31> <6571756174696F6E2E452E312E322E31>] >>}
\MusicPDFWriteObject{9885}{<< /Names [<6571756174696F6E2E452E312E322E3130> \MusicPDFRef{14893} <6571756174696F6E2E452E312E322E3131> \MusicPDFRef{14894} <6571756174696F6E2E452E312E322E3132> \MusicPDFRef{14895} <6571756174696F6E2E452E312E322E3133> \MusicPDFRef{14896} <6571756174696F6E2E452E312E322E3134> \MusicPDFRef{14897} <6571756174696F6E2E452E312E322E3135> \MusicPDFRef{14898}] /Limits [<6571756174696F6E2E452E312E322E3130> <6571756174696F6E2E452E312E322E3135>] >>}
\MusicPDFWriteObject{9886}{<< /Names [<6571756174696F6E2E452E312E322E3136> \MusicPDFRef{14899} <6571756174696F6E2E452E312E322E3137> \MusicPDFRef{14900} <6571756174696F6E2E452E312E322E3138> \MusicPDFRef{14901} <6571756174696F6E2E452E312E322E3139> \MusicPDFRef{14902} <6571756174696F6E2E452E312E322E32> \MusicPDFRef{14903} <6571756174696F6E2E452E312E322E3230> \MusicPDFRef{14904}] /Limits [<6571756174696F6E2E452E312E322E3136> <6571756174696F6E2E452E312E322E3230>] >>}
\MusicPDFWriteObject{9887}{<< /Names [<6571756174696F6E2E452E312E322E3231> \MusicPDFRef{14905} <6571756174696F6E2E452E312E322E3232> \MusicPDFRef{14906} <6571756174696F6E2E452E312E322E3233> \MusicPDFRef{14907} <6571756174696F6E2E452E312E322E3234> \MusicPDFRef{14908} <6571756174696F6E2E452E312E322E3235> \MusicPDFRef{14909} <6571756174696F6E2E452E312E322E3236> \MusicPDFRef{14910}] /Limits [<6571756174696F6E2E452E312E322E3231> <6571756174696F6E2E452E312E322E3236>] >>}
\MusicPDFWriteObject{9888}{<< /Names [<6571756174696F6E2E452E312E322E3237> \MusicPDFRef{14911} <6571756174696F6E2E452E312E322E3238> \MusicPDFRef{14912} <6571756174696F6E2E452E312E322E3239> \MusicPDFRef{14913} <6571756174696F6E2E452E312E322E33> \MusicPDFRef{14914} <6571756174696F6E2E452E312E322E34> \MusicPDFRef{14915} <6571756174696F6E2E452E312E322E35> \MusicPDFRef{14916}] /Limits [<6571756174696F6E2E452E312E322E3237> <6571756174696F6E2E452E312E322E35>] >>}
\MusicPDFWriteObject{9889}{<< /Names [<6571756174696F6E2E452E312E322E36> \MusicPDFRef{14917} <6571756174696F6E2E452E312E322E37> \MusicPDFRef{14918} <6571756174696F6E2E452E312E322E38> \MusicPDFRef{14919} <6571756174696F6E2E452E312E322E39> \MusicPDFRef{14920} <6571756174696F6E2E452E322E312E31> \MusicPDFRef{14921} <6571756174696F6E2E452E322E312E32> \MusicPDFRef{14922}] /Limits [<6571756174696F6E2E452E312E322E36> <6571756174696F6E2E452E322E312E32>] >>}
\MusicPDFWriteObject{9890}{<< /Names [<6571756174696F6E2E452E322E312E33> \MusicPDFRef{14923} <6571756174696F6E2E452E322E312E34> \MusicPDFRef{14924} <6571756174696F6E2E452E322E312E35> \MusicPDFRef{14925} <6571756174696F6E2E452E322E322E31> \MusicPDFRef{14926} <6571756174696F6E2E452E322E322E3130> \MusicPDFRef{14927} <6571756174696F6E2E452E322E322E3131> \MusicPDFRef{14928}] /Limits [<6571756174696F6E2E452E322E312E33> <6571756174696F6E2E452E322E322E3131>] >>}
\MusicPDFWriteObject{9891}{<< /Names [<6571756174696F6E2E452E322E322E3132> \MusicPDFRef{14929} <6571756174696F6E2E452E322E322E3133> \MusicPDFRef{14930} <6571756174696F6E2E452E322E322E3134> \MusicPDFRef{14931} <6571756174696F6E2E452E322E322E32> \MusicPDFRef{14932} <6571756174696F6E2E452E322E322E33> \MusicPDFRef{14933} <6571756174696F6E2E452E322E322E34> \MusicPDFRef{14934}] /Limits [<6571756174696F6E2E452E322E322E3132> <6571756174696F6E2E452E322E322E34>] >>}
\MusicPDFWriteObject{9892}{<< /Names [<6571756174696F6E2E452E322E322E35> \MusicPDFRef{14935} <6571756174696F6E2E452E322E322E36> \MusicPDFRef{14936} <6571756174696F6E2E452E322E322E37> \MusicPDFRef{14937} <6571756174696F6E2E452E322E322E38> \MusicPDFRef{14938} <6571756174696F6E2E452E322E322E39> \MusicPDFRef{14939} <6571756174696F6E2E452E322E332E31> \MusicPDFRef{14940}] /Limits [<6571756174696F6E2E452E322E322E35> <6571756174696F6E2E452E322E332E31>] >>}
\MusicPDFWriteObject{9893}{<< /Names [<6571756174696F6E2E452E322E332E3130> \MusicPDFRef{14941} <6571756174696F6E2E452E322E332E3131> \MusicPDFRef{14942} <6571756174696F6E2E452E322E332E3132> \MusicPDFRef{14943} <6571756174696F6E2E452E322E332E3133> \MusicPDFRef{14944} <6571756174696F6E2E452E322E332E3134> \MusicPDFRef{14945} <6571756174696F6E2E452E322E332E3135> \MusicPDFRef{14946}] /Limits [<6571756174696F6E2E452E322E332E3130> <6571756174696F6E2E452E322E332E3135>] >>}
\MusicPDFWriteObject{9894}{<< /Names [<6571756174696F6E2E452E322E332E3136> \MusicPDFRef{14947} <6571756174696F6E2E452E322E332E3137> \MusicPDFRef{14948} <6571756174696F6E2E452E322E332E32> \MusicPDFRef{14949} <6571756174696F6E2E452E322E332E33> \MusicPDFRef{14950} <6571756174696F6E2E452E322E332E34> \MusicPDFRef{14951} <6571756174696F6E2E452E322E332E35> \MusicPDFRef{14952}] /Limits [<6571756174696F6E2E452E322E332E3136> <6571756174696F6E2E452E322E332E35>] >>}
\MusicPDFWriteObject{9895}{<< /Names [<6571756174696F6E2E452E322E332E36> \MusicPDFRef{14953} <6571756174696F6E2E452E322E332E37> \MusicPDFRef{14954} <6571756174696F6E2E452E322E332E38> \MusicPDFRef{14955} <6571756174696F6E2E452E322E332E39> \MusicPDFRef{14956} <6571756174696F6E2E452E332E312E31> \MusicPDFRef{14957} <6571756174696F6E2E452E332E312E32> \MusicPDFRef{14958}] /Limits [<6571756174696F6E2E452E322E332E36> <6571756174696F6E2E452E332E312E32>] >>}
\MusicPDFWriteObject{9896}{<< /Names [<6571756174696F6E2E452E332E312E33> \MusicPDFRef{14959} <6571756174696F6E2E452E332E312E34> \MusicPDFRef{14960} <6571756174696F6E2E452E332E312E35> \MusicPDFRef{14961} <6571756174696F6E2E452E332E332E31> \MusicPDFRef{14962} <6571756174696F6E2E452E332E342E31> \MusicPDFRef{14963} <6571756174696F6E2E452E332E352E31> \MusicPDFRef{14964}] /Limits [<6571756174696F6E2E452E332E312E33> <6571756174696F6E2E452E332E352E31>] >>}
\MusicPDFWriteObject{9897}{<< /Names [<6571756174696F6E2E452E342E312E31> \MusicPDFRef{14965} <6571756174696F6E2E452E342E312E3130> \MusicPDFRef{14966} <6571756174696F6E2E452E342E312E3131> \MusicPDFRef{14967} <6571756174696F6E2E452E342E312E32> \MusicPDFRef{14968} <6571756174696F6E2E452E342E312E33> \MusicPDFRef{14969} <6571756174696F6E2E452E342E312E34> \MusicPDFRef{14970}] /Limits [<6571756174696F6E2E452E342E312E31> <6571756174696F6E2E452E342E312E34>] >>}
\MusicPDFWriteObject{9898}{<< /Names [<6571756174696F6E2E452E342E312E35> \MusicPDFRef{14971} <6571756174696F6E2E452E342E312E36> \MusicPDFRef{14972} <6571756174696F6E2E452E342E312E37> \MusicPDFRef{14973} <6571756174696F6E2E452E342E312E38> \MusicPDFRef{14974} <6571756174696F6E2E452E342E312E39> \MusicPDFRef{14975} <6571756174696F6E2E452E342E322E31> \MusicPDFRef{14976}] /Limits [<6571756174696F6E2E452E342E312E35> <6571756174696F6E2E452E342E322E31>] >>}
\MusicPDFWriteObject{9899}{<< /Names [<6571756174696F6E2E452E342E322E32> \MusicPDFRef{14977} <6571756174696F6E2E452E342E322E33> \MusicPDFRef{14978} <6571756174696F6E2E452E342E322E34> \MusicPDFRef{14979} <6571756174696F6E2E452E342E322E35> \MusicPDFRef{14980} <6571756174696F6E2E452E342E322E36> \MusicPDFRef{14981} <6571756174696F6E2E452E342E322E37> \MusicPDFRef{14982}] /Limits [<6571756174696F6E2E452E342E322E32> <6571756174696F6E2E452E342E322E37>] >>}
\MusicPDFWriteObject{9900}{<< /Names [<6571756174696F6E2E452E342E322E38> \MusicPDFRef{14983} <6571756174696F6E2E452E342E322E39> \MusicPDFRef{14984} <6571756174696F6E2E452E352E312E31> \MusicPDFRef{14985} <6571756174696F6E2E452E352E312E3130> \MusicPDFRef{14986} <6571756174696F6E2E452E352E312E3131> \MusicPDFRef{14987} <6571756174696F6E2E452E352E312E3132> \MusicPDFRef{14988}] /Limits [<6571756174696F6E2E452E342E322E38> <6571756174696F6E2E452E352E312E3132>] >>}
\MusicPDFWriteObject{9901}{<< /Names [<6571756174696F6E2E452E352E312E3133> \MusicPDFRef{14989} <6571756174696F6E2E452E352E312E3134> \MusicPDFRef{14990} <6571756174696F6E2E452E352E312E3135> \MusicPDFRef{14991} <6571756174696F6E2E452E352E312E3136> \MusicPDFRef{14992} <6571756174696F6E2E452E352E312E3137> \MusicPDFRef{14993} <6571756174696F6E2E452E352E312E3138> \MusicPDFRef{14994}] /Limits [<6571756174696F6E2E452E352E312E3133> <6571756174696F6E2E452E352E312E3138>] >>}
\MusicPDFWriteObject{9902}{<< /Names [<6571756174696F6E2E452E352E312E3139> \MusicPDFRef{14995} <6571756174696F6E2E452E352E312E32> \MusicPDFRef{14996} <6571756174696F6E2E452E352E312E3230> \MusicPDFRef{14997} <6571756174696F6E2E452E352E312E3231> \MusicPDFRef{14998} <6571756174696F6E2E452E352E312E3232> \MusicPDFRef{14999} <6571756174696F6E2E452E352E312E3233> \MusicPDFRef{15000}] /Limits [<6571756174696F6E2E452E352E312E3139> <6571756174696F6E2E452E352E312E3233>] >>}
\MusicPDFWriteObject{9903}{<< /Names [<6571756174696F6E2E452E352E312E3234> \MusicPDFRef{15001} <6571756174696F6E2E452E352E312E3235> \MusicPDFRef{15002} <6571756174696F6E2E452E352E312E3236> \MusicPDFRef{15003} <6571756174696F6E2E452E352E312E3237> \MusicPDFRef{15004} <6571756174696F6E2E452E352E312E3238> \MusicPDFRef{15005} <6571756174696F6E2E452E352E312E3239> \MusicPDFRef{15006}] /Limits [<6571756174696F6E2E452E352E312E3234> <6571756174696F6E2E452E352E312E3239>] >>}
\MusicPDFWriteObject{9904}{<< /Names [<6571756174696F6E2E452E352E312E33> \MusicPDFRef{15007} <6571756174696F6E2E452E352E312E3330> \MusicPDFRef{15008} <6571756174696F6E2E452E352E312E3331> \MusicPDFRef{15009} <6571756174696F6E2E452E352E312E3332> \MusicPDFRef{15010} <6571756174696F6E2E452E352E312E3333> \MusicPDFRef{15011} <6571756174696F6E2E452E352E312E3334> \MusicPDFRef{15012}] /Limits [<6571756174696F6E2E452E352E312E33> <6571756174696F6E2E452E352E312E3334>] >>}
\MusicPDFWriteObject{9905}{<< /Names [<6571756174696F6E2E452E352E312E3335> \MusicPDFRef{15013} <6571756174696F6E2E452E352E312E3336> \MusicPDFRef{15014} <6571756174696F6E2E452E352E312E34> \MusicPDFRef{15015} <6571756174696F6E2E452E352E312E35> \MusicPDFRef{15016} <6571756174696F6E2E452E352E312E36> \MusicPDFRef{15017} <6571756174696F6E2E452E352E312E37> \MusicPDFRef{15018}] /Limits [<6571756174696F6E2E452E352E312E3335> <6571756174696F6E2E452E352E312E37>] >>}
\MusicPDFWriteObject{9906}{<< /Names [<6571756174696F6E2E452E352E312E38> \MusicPDFRef{15019} <6571756174696F6E2E452E352E312E39> \MusicPDFRef{15020} <6571756174696F6E2E452E352E322E31> \MusicPDFRef{15021} <6571756174696F6E2E452E352E322E32> \MusicPDFRef{15022} <6571756174696F6E2E452E352E322E33> \MusicPDFRef{15023} <6571756174696F6E2E452E352E322E34> \MusicPDFRef{15024}] /Limits [<6571756174696F6E2E452E352E312E38> <6571756174696F6E2E452E352E322E34>] >>}
\MusicPDFWriteObject{9907}{<< /Names [<6571756174696F6E2E452E352E322E35> \MusicPDFRef{15025} <6571756174696F6E2E452E352E322E36> \MusicPDFRef{15026} <6571756174696F6E2E452E352E322E37> \MusicPDFRef{15027} <6571756174696F6E2E452E352E332E31> \MusicPDFRef{15028} <6571756174696F6E2E452E352E332E32> \MusicPDFRef{15029} <6571756174696F6E2E452E352E332E33> \MusicPDFRef{15030}] /Limits [<6571756174696F6E2E452E352E322E35> <6571756174696F6E2E452E352E332E33>] >>}
\MusicPDFWriteObject{9908}{<< /Names [<6571756174696F6E2E452E352E332E34> \MusicPDFRef{15031} <6571756174696F6E2E452E352E332E35> \MusicPDFRef{15032} <6571756174696F6E2E452E352E342E31> \MusicPDFRef{15033} <6571756174696F6E2E452E352E342E32> \MusicPDFRef{15034} <6571756174696F6E2E452E352E342E33> \MusicPDFRef{15035} <6571756174696F6E2E452E352E352E31> \MusicPDFRef{15036}] /Limits [<6571756174696F6E2E452E352E332E34> <6571756174696F6E2E452E352E352E31>] >>}
\MusicPDFWriteObject{9909}{<< /Names [<6571756174696F6E2E452E352E352E3130> \MusicPDFRef{15037} <6571756174696F6E2E452E352E352E32> \MusicPDFRef{15038} <6571756174696F6E2E452E352E352E33> \MusicPDFRef{15039} <6571756174696F6E2E452E352E352E34> \MusicPDFRef{15040} <6571756174696F6E2E452E352E352E35> \MusicPDFRef{15041} <6571756174696F6E2E452E352E352E36> \MusicPDFRef{15042}] /Limits [<6571756174696F6E2E452E352E352E3130> <6571756174696F6E2E452E352E352E36>] >>}
\MusicPDFWriteObject{9910}{<< /Names [<6571756174696F6E2E452E352E352E37> \MusicPDFRef{15043} <6571756174696F6E2E452E352E352E38> \MusicPDFRef{15044} <6571756174696F6E2E452E352E352E39> \MusicPDFRef{15045} <6571756174696F6E2E452E352E372E31> \MusicPDFRef{15046} <6571756174696F6E2E452E352E372E32> \MusicPDFRef{15047} <6571756174696F6E2E452E352E372E33> \MusicPDFRef{15048}] /Limits [<6571756174696F6E2E452E352E352E37> <6571756174696F6E2E452E352E372E33>] >>}
\MusicPDFWriteObject{9911}{<< /Names [<6571756174696F6E2E452E352E372E34> \MusicPDFRef{15049} <6571756174696F6E2E452E352E372E35> \MusicPDFRef{15050} <6571756174696F6E2E452E352E372E36> \MusicPDFRef{15051} <6571756174696F6E2E452E352E372E37> \MusicPDFRef{15052} <6571756174696F6E2E452E352E372E38> \MusicPDFRef{15053} <6571756174696F6E2E452E362E312E31> \MusicPDFRef{15054}] /Limits [<6571756174696F6E2E452E352E372E34> <6571756174696F6E2E452E362E312E31>] >>}
\MusicPDFWriteObject{9912}{<< /Names [<6571756174696F6E2E452E362E312E32> \MusicPDFRef{15055} <6571756174696F6E2E452E362E312E33> \MusicPDFRef{15056} <6571756174696F6E2E452E362E312E34> \MusicPDFRef{15057} <6571756174696F6E2E452E362E312E35> \MusicPDFRef{15058} <6571756174696F6E2E452E362E322E31> \MusicPDFRef{15059} <6571756174696F6E2E452E362E322E3130> \MusicPDFRef{15060}] /Limits [<6571756174696F6E2E452E362E312E32> <6571756174696F6E2E452E362E322E3130>] >>}
\MusicPDFWriteObject{9913}{<< /Names [<6571756174696F6E2E452E362E322E3131> \MusicPDFRef{15061} <6571756174696F6E2E452E362E322E3132> \MusicPDFRef{15062} <6571756174696F6E2E452E362E322E3133> \MusicPDFRef{15063} <6571756174696F6E2E452E362E322E32> \MusicPDFRef{15064} <6571756174696F6E2E452E362E322E33> \MusicPDFRef{15065} <6571756174696F6E2E452E362E322E34> \MusicPDFRef{15066}] /Limits [<6571756174696F6E2E452E362E322E3131> <6571756174696F6E2E452E362E322E34>] >>}
\MusicPDFWriteObject{9914}{<< /Names [<6571756174696F6E2E452E362E322E35> \MusicPDFRef{15067} <6571756174696F6E2E452E362E322E36> \MusicPDFRef{15068} <6571756174696F6E2E452E362E322E37> \MusicPDFRef{15069} <6571756174696F6E2E452E362E322E38> \MusicPDFRef{15070} <6571756174696F6E2E452E362E322E39> \MusicPDFRef{15071} <6571756174696F6E2E452E362E332E31> \MusicPDFRef{15072}] /Limits [<6571756174696F6E2E452E362E322E35> <6571756174696F6E2E452E362E332E31>] >>}
\MusicPDFWriteObject{9915}{<< /Names [<6578706563746174696F6E2E312E372E31> \MusicPDFRef{15073} <6669677572652E422E31> \MusicPDFRef{15074} <6669677572652E422E32> \MusicPDFRef{15075} <6669677572652E422E33> \MusicPDFRef{15076} <6669677572652E422E34> \MusicPDFRef{15077} <6669677572652E442E31> \MusicPDFRef{15078}] /Limits [<6578706563746174696F6E2E312E372E31> <6669677572652E442E31>] >>}
\MusicPDFWriteObject{9916}{<< /Names [<676C6F3A4F465373746770> \MusicPDFRef{15079} <676C6F3A53696E746567657273> \MusicPDFRef{15080} <676C6F3A536C6F63616C6D656C6C696E> \MusicPDFRef{15081} <676C6F3A536C6F63616C7468657461> \MusicPDFRef{15082} <676C6F3A676C6F62616C6D656C6C696E> \MusicPDFRef{15083} <676C6F3A676C6F62616C7468657461> \MusicPDFRef{15084}] /Limits [<676C6F3A4F465373746770> <676C6F3A676C6F62616C7468657461>] >>}
\MusicPDFWriteObject{9917}{<< /Names [<676C6F3A6964636C6770> \MusicPDFRef{15085} <6E6F746174696F6E2E31302E322E312E31> \MusicPDFRef{15086} <6E6F746174696F6E2E31302E322E322E31> \MusicPDFRef{15087} <6E6F746174696F6E2E322E332E312E31> \MusicPDFRef{15088} <6E6F746174696F6E2E332E342E312E31> \MusicPDFRef{15089} <6E6F746174696F6E2E342E352E312E31> \MusicPDFRef{15090}] /Limits [<676C6F3A6964636C6770> <6E6F746174696F6E2E342E352E312E31>] >>}
\MusicPDFWriteObject{9918}{<< /Names [<6E6F746174696F6E2E342E362E312E31> \MusicPDFRef{15091} <6E6F746174696F6E2E372E332E342E31> \MusicPDFRef{15092} <6E6F746174696F6E2E452E302E302E31> \MusicPDFRef{15093} <6E6F746174696F6E2E452E322E322E31> \MusicPDFRef{15094} <6E6F746174696F6E2E452E332E312E31> \MusicPDFRef{15095} <6E6F746174696F6E2E452E342E312E31> \MusicPDFRef{15096}] /Limits [<6E6F746174696F6E2E342E362E312E31> <6E6F746174696F6E2E452E342E312E31>] >>}
\MusicPDFWriteObject{9919}{<< /Names [<6E6F746174696F6E2E452E352E312E31> \MusicPDFRef{15097} <6E6F746174696F6E2E452E352E332E31> \MusicPDFRef{15098} <6E6F746174696F6E2E452E352E352E31> \MusicPDFRef{15099} <6E6F746174696F6E2E452E352E352E32> \MusicPDFRef{15100} <6E6F746174696F6E2E452E352E362E31> \MusicPDFRef{15101} <6E6F746174696F6E2E452E362E332E31> \MusicPDFRef{15102}] /Limits [<6E6F746174696F6E2E452E352E312E31> <6E6F746174696F6E2E452E362E332E31>] >>}
\MusicPDFWriteObject{9920}{<< /Names [<706167652E31> \MusicPDFRef{15103} <706167652E3130> \MusicPDFRef{15104} <706167652E313030> \MusicPDFRef{15105} <706167652E313031> \MusicPDFRef{15106} <706167652E313032> \MusicPDFRef{15107} <706167652E313033> \MusicPDFRef{15108}] /Limits [<706167652E31> <706167652E313033>] >>}
\MusicPDFWriteObject{9921}{<< /Names [<706167652E313034> \MusicPDFRef{15109} <706167652E313035> \MusicPDFRef{15110} <706167652E313036> \MusicPDFRef{15111} <706167652E313037> \MusicPDFRef{15112} <706167652E313038> \MusicPDFRef{15113} <706167652E313039> \MusicPDFRef{15114}] /Limits [<706167652E313034> <706167652E313039>] >>}
\MusicPDFWriteObject{9922}{<< /Names [<706167652E3131> \MusicPDFRef{15115} <706167652E313130> \MusicPDFRef{15116} <706167652E313131> \MusicPDFRef{15117} <706167652E313132> \MusicPDFRef{15118} <706167652E313133> \MusicPDFRef{15119} <706167652E313134> \MusicPDFRef{15120}] /Limits [<706167652E3131> <706167652E313134>] >>}
\MusicPDFWriteObject{9923}{<< /Names [<706167652E313135> \MusicPDFRef{15121} <706167652E313136> \MusicPDFRef{15122} <706167652E313137> \MusicPDFRef{15123} <706167652E313138> \MusicPDFRef{15124} <706167652E313139> \MusicPDFRef{15125} <706167652E3132> \MusicPDFRef{15126}] /Limits [<706167652E313135> <706167652E3132>] >>}
\MusicPDFWriteObject{9924}{<< /Names [<706167652E313230> \MusicPDFRef{15127} <706167652E313231> \MusicPDFRef{15128} <706167652E313232> \MusicPDFRef{15129} <706167652E313233> \MusicPDFRef{15130} <706167652E313234> \MusicPDFRef{15131} <706167652E313235> \MusicPDFRef{15132}] /Limits [<706167652E313230> <706167652E313235>] >>}
\MusicPDFWriteObject{9925}{<< /Names [<706167652E313236> \MusicPDFRef{15133} <706167652E313237> \MusicPDFRef{15134} <706167652E313238> \MusicPDFRef{15135} <706167652E313239> \MusicPDFRef{15136} <706167652E3133> \MusicPDFRef{15137} <706167652E313330> \MusicPDFRef{15138}] /Limits [<706167652E313236> <706167652E313330>] >>}
\MusicPDFWriteObject{9926}{<< /Names [<706167652E313331> \MusicPDFRef{15139} <706167652E313332> \MusicPDFRef{15140} <706167652E313333> \MusicPDFRef{15141} <706167652E313334> \MusicPDFRef{15142} <706167652E313335> \MusicPDFRef{15143} <706167652E313336> \MusicPDFRef{15144}] /Limits [<706167652E313331> <706167652E313336>] >>}
\MusicPDFWriteObject{9927}{<< /Names [<706167652E313337> \MusicPDFRef{15145} <706167652E313338> \MusicPDFRef{15146} <706167652E313339> \MusicPDFRef{15147} <706167652E3134> \MusicPDFRef{15148} <706167652E313430> \MusicPDFRef{15149} <706167652E313431> \MusicPDFRef{15150}] /Limits [<706167652E313337> <706167652E313431>] >>}
\MusicPDFWriteObject{9928}{<< /Names [<706167652E313432> \MusicPDFRef{15151} <706167652E313433> \MusicPDFRef{15152} <706167652E313434> \MusicPDFRef{15153} <706167652E313435> \MusicPDFRef{15154} <706167652E313436> \MusicPDFRef{15155} <706167652E313437> \MusicPDFRef{15156}] /Limits [<706167652E313432> <706167652E313437>] >>}
\MusicPDFWriteObject{9929}{<< /Names [<706167652E313438> \MusicPDFRef{15157} <706167652E313439> \MusicPDFRef{15158} <706167652E3135> \MusicPDFRef{15159} <706167652E313530> \MusicPDFRef{15160} <706167652E313531> \MusicPDFRef{15161} <706167652E313532> \MusicPDFRef{15162}] /Limits [<706167652E313438> <706167652E313532>] >>}
\MusicPDFWriteObject{9930}{<< /Names [<706167652E313533> \MusicPDFRef{15163} <706167652E313534> \MusicPDFRef{15164} <706167652E313535> \MusicPDFRef{15165} <706167652E313536> \MusicPDFRef{15166} <706167652E313537> \MusicPDFRef{15167} <706167652E313538> \MusicPDFRef{15168}] /Limits [<706167652E313533> <706167652E313538>] >>}
\MusicPDFWriteObject{9931}{<< /Names [<706167652E313539> \MusicPDFRef{15169} <706167652E3136> \MusicPDFRef{15170} <706167652E313630> \MusicPDFRef{15171} <706167652E313631> \MusicPDFRef{15172} <706167652E313632> \MusicPDFRef{15173} <706167652E313633> \MusicPDFRef{15174}] /Limits [<706167652E313539> <706167652E313633>] >>}
\MusicPDFWriteObject{9932}{<< /Names [<706167652E313634> \MusicPDFRef{15175} <706167652E313635> \MusicPDFRef{15176} <706167652E313636> \MusicPDFRef{15177} <706167652E313637> \MusicPDFRef{15178} <706167652E313638> \MusicPDFRef{15179} <706167652E313639> \MusicPDFRef{15180}] /Limits [<706167652E313634> <706167652E313639>] >>}
\MusicPDFWriteObject{9933}{<< /Names [<706167652E3137> \MusicPDFRef{15181} <706167652E313730> \MusicPDFRef{15182} <706167652E313731> \MusicPDFRef{15183} <706167652E313732> \MusicPDFRef{15184} <706167652E313733> \MusicPDFRef{15185} <706167652E313734> \MusicPDFRef{15186}] /Limits [<706167652E3137> <706167652E313734>] >>}
\MusicPDFWriteObject{9934}{<< /Names [<706167652E313735> \MusicPDFRef{15187} <706167652E313736> \MusicPDFRef{15188} <706167652E313737> \MusicPDFRef{15189} <706167652E313738> \MusicPDFRef{15190} <706167652E313739> \MusicPDFRef{15191} <706167652E3138> \MusicPDFRef{15192}] /Limits [<706167652E313735> <706167652E3138>] >>}
\MusicPDFWriteObject{9935}{<< /Names [<706167652E313830> \MusicPDFRef{15193} <706167652E313831> \MusicPDFRef{15194} <706167652E313832> \MusicPDFRef{15195} <706167652E313833> \MusicPDFRef{15196} <706167652E313834> \MusicPDFRef{15197} <706167652E313835> \MusicPDFRef{15198}] /Limits [<706167652E313830> <706167652E313835>] >>}
\MusicPDFWriteObject{9936}{<< /Names [<706167652E313836> \MusicPDFRef{15199} <706167652E313837> \MusicPDFRef{15200} <706167652E313838> \MusicPDFRef{15201} <706167652E313839> \MusicPDFRef{15202} <706167652E3139> \MusicPDFRef{15203} <706167652E313930> \MusicPDFRef{15204}] /Limits [<706167652E313836> <706167652E313930>] >>}
\MusicPDFWriteObject{9937}{<< /Names [<706167652E313931> \MusicPDFRef{15205} <706167652E313932> \MusicPDFRef{15206} <706167652E313933> \MusicPDFRef{15207} <706167652E313934> \MusicPDFRef{15208} <706167652E313935> \MusicPDFRef{15209} <706167652E313936> \MusicPDFRef{15210}] /Limits [<706167652E313931> <706167652E313936>] >>}
\MusicPDFWriteObject{9938}{<< /Names [<706167652E313937> \MusicPDFRef{15211} <706167652E313938> \MusicPDFRef{15212} <706167652E313939> \MusicPDFRef{15213} <706167652E32> \MusicPDFRef{15214} <706167652E3230> \MusicPDFRef{15215} <706167652E323030> \MusicPDFRef{15216}] /Limits [<706167652E313937> <706167652E323030>] >>}
\MusicPDFWriteObject{9939}{<< /Names [<706167652E323031> \MusicPDFRef{15217} <706167652E323032> \MusicPDFRef{15218} <706167652E323033> \MusicPDFRef{15219} <706167652E323034> \MusicPDFRef{15220} <706167652E323035> \MusicPDFRef{15221} <706167652E323036> \MusicPDFRef{15222}] /Limits [<706167652E323031> <706167652E323036>] >>}
\MusicPDFWriteObject{9940}{<< /Names [<706167652E323037> \MusicPDFRef{15223} <706167652E323038> \MusicPDFRef{15224} <706167652E323039> \MusicPDFRef{15225} <706167652E3231> \MusicPDFRef{15226} <706167652E323130> \MusicPDFRef{15227} <706167652E323131> \MusicPDFRef{15228}] /Limits [<706167652E323037> <706167652E323131>] >>}
\MusicPDFWriteObject{9941}{<< /Names [<706167652E323132> \MusicPDFRef{15229} <706167652E323133> \MusicPDFRef{15230} <706167652E323134> \MusicPDFRef{15231} <706167652E323135> \MusicPDFRef{15232} <706167652E323136> \MusicPDFRef{15233} <706167652E323137> \MusicPDFRef{15234}] /Limits [<706167652E323132> <706167652E323137>] >>}
\MusicPDFWriteObject{9942}{<< /Names [<706167652E323138> \MusicPDFRef{15235} <706167652E323139> \MusicPDFRef{15236} <706167652E3232> \MusicPDFRef{15237} <706167652E323230> \MusicPDFRef{15238} <706167652E323231> \MusicPDFRef{15239} <706167652E323232> \MusicPDFRef{15240}] /Limits [<706167652E323138> <706167652E323232>] >>}
\MusicPDFWriteObject{9943}{<< /Names [<706167652E323233> \MusicPDFRef{15241} <706167652E323234> \MusicPDFRef{15242} <706167652E323235> \MusicPDFRef{15243} <706167652E323236> \MusicPDFRef{15244} <706167652E323237> \MusicPDFRef{15245} <706167652E323238> \MusicPDFRef{15246}] /Limits [<706167652E323233> <706167652E323238>] >>}
\MusicPDFWriteObject{9944}{<< /Names [<706167652E323239> \MusicPDFRef{15247} <706167652E3233> \MusicPDFRef{15248} <706167652E323330> \MusicPDFRef{15249} <706167652E323331> \MusicPDFRef{15250} <706167652E323332> \MusicPDFRef{15251} <706167652E323333> \MusicPDFRef{15252}] /Limits [<706167652E323239> <706167652E323333>] >>}
\MusicPDFWriteObject{9945}{<< /Names [<706167652E323334> \MusicPDFRef{15253} <706167652E323335> \MusicPDFRef{15254} <706167652E323336> \MusicPDFRef{15255} <706167652E323337> \MusicPDFRef{15256} <706167652E323338> \MusicPDFRef{15257} <706167652E323339> \MusicPDFRef{15258}] /Limits [<706167652E323334> <706167652E323339>] >>}
\MusicPDFWriteObject{9946}{<< /Names [<706167652E3234> \MusicPDFRef{15259} <706167652E323430> \MusicPDFRef{15260} <706167652E323431> \MusicPDFRef{15261} <706167652E323432> \MusicPDFRef{15262} <706167652E323433> \MusicPDFRef{15263} <706167652E323434> \MusicPDFRef{15264}] /Limits [<706167652E3234> <706167652E323434>] >>}
\MusicPDFWriteObject{9947}{<< /Names [<706167652E323435> \MusicPDFRef{15265} <706167652E323436> \MusicPDFRef{15266} <706167652E323437> \MusicPDFRef{15267} <706167652E323438> \MusicPDFRef{15268} <706167652E323439> \MusicPDFRef{15269} <706167652E3235> \MusicPDFRef{15270}] /Limits [<706167652E323435> <706167652E3235>] >>}
\MusicPDFWriteObject{9948}{<< /Names [<706167652E323530> \MusicPDFRef{15271} <706167652E323531> \MusicPDFRef{15272} <706167652E323532> \MusicPDFRef{15273} <706167652E323533> \MusicPDFRef{15274} <706167652E323534> \MusicPDFRef{15275} <706167652E323535> \MusicPDFRef{15276}] /Limits [<706167652E323530> <706167652E323535>] >>}
\MusicPDFWriteObject{9949}{<< /Names [<706167652E323536> \MusicPDFRef{15277} <706167652E323537> \MusicPDFRef{15278} <706167652E323538> \MusicPDFRef{15279} <706167652E323539> \MusicPDFRef{15280} <706167652E3236> \MusicPDFRef{15281} <706167652E323630> \MusicPDFRef{15282}] /Limits [<706167652E323536> <706167652E323630>] >>}
\MusicPDFWriteObject{9950}{<< /Names [<706167652E323631> \MusicPDFRef{15283} <706167652E323632> \MusicPDFRef{15284} <706167652E323633> \MusicPDFRef{15285} <706167652E323634> \MusicPDFRef{15286} <706167652E323635> \MusicPDFRef{15287} <706167652E323636> \MusicPDFRef{15288}] /Limits [<706167652E323631> <706167652E323636>] >>}
\MusicPDFWriteObject{9951}{<< /Names [<706167652E323637> \MusicPDFRef{15289} <706167652E323638> \MusicPDFRef{15290} <706167652E323639> \MusicPDFRef{15291} <706167652E3237> \MusicPDFRef{15292} <706167652E323730> \MusicPDFRef{15293} <706167652E323731> \MusicPDFRef{15294}] /Limits [<706167652E323637> <706167652E323731>] >>}
\MusicPDFWriteObject{9952}{<< /Names [<706167652E323732> \MusicPDFRef{15295} <706167652E323733> \MusicPDFRef{15296} <706167652E323734> \MusicPDFRef{15297} <706167652E323735> \MusicPDFRef{15298} <706167652E323736> \MusicPDFRef{15299} <706167652E323737> \MusicPDFRef{15300}] /Limits [<706167652E323732> <706167652E323737>] >>}
\MusicPDFWriteObject{9953}{<< /Names [<706167652E323738> \MusicPDFRef{15301} <706167652E323739> \MusicPDFRef{15302} <706167652E3238> \MusicPDFRef{15303} <706167652E323830> \MusicPDFRef{15304} <706167652E323831> \MusicPDFRef{15305} <706167652E323832> \MusicPDFRef{15306}] /Limits [<706167652E323738> <706167652E323832>] >>}
\MusicPDFWriteObject{9954}{<< /Names [<706167652E323833> \MusicPDFRef{15307} <706167652E323834> \MusicPDFRef{15308} <706167652E323835> \MusicPDFRef{15309} <706167652E323836> \MusicPDFRef{15310} <706167652E323837> \MusicPDFRef{15311} <706167652E323838> \MusicPDFRef{15312}] /Limits [<706167652E323833> <706167652E323838>] >>}
\MusicPDFWriteObject{9955}{<< /Names [<706167652E323839> \MusicPDFRef{15313} <706167652E3239> \MusicPDFRef{15314} <706167652E323930> \MusicPDFRef{15315} <706167652E323931> \MusicPDFRef{15316} <706167652E323932> \MusicPDFRef{15317} <706167652E323933> \MusicPDFRef{15318}] /Limits [<706167652E323839> <706167652E323933>] >>}
\MusicPDFWriteObject{9956}{<< /Names [<706167652E323934> \MusicPDFRef{15319} <706167652E323935> \MusicPDFRef{15320} <706167652E323936> \MusicPDFRef{15321} <706167652E323937> \MusicPDFRef{15322} <706167652E323938> \MusicPDFRef{15323} <706167652E323939> \MusicPDFRef{15324}] /Limits [<706167652E323934> <706167652E323939>] >>}
\MusicPDFWriteObject{9957}{<< /Names [<706167652E33> \MusicPDFRef{15325} <706167652E3330> \MusicPDFRef{15326} <706167652E333030> \MusicPDFRef{15327} <706167652E333031> \MusicPDFRef{15328} <706167652E333032> \MusicPDFRef{15329} <706167652E333033> \MusicPDFRef{15330}] /Limits [<706167652E33> <706167652E333033>] >>}
\MusicPDFWriteObject{9958}{<< /Names [<706167652E333034> \MusicPDFRef{15331} <706167652E333035> \MusicPDFRef{15332} <706167652E333036> \MusicPDFRef{15333} <706167652E333037> \MusicPDFRef{15334} <706167652E333038> \MusicPDFRef{15335} <706167652E333039> \MusicPDFRef{15336}] /Limits [<706167652E333034> <706167652E333039>] >>}
\MusicPDFWriteObject{9959}{<< /Names [<706167652E3331> \MusicPDFRef{15337} <706167652E333130> \MusicPDFRef{15338} <706167652E333131> \MusicPDFRef{15339} <706167652E333132> \MusicPDFRef{15340} <706167652E333133> \MusicPDFRef{15341} <706167652E333134> \MusicPDFRef{15342}] /Limits [<706167652E3331> <706167652E333134>] >>}
\MusicPDFWriteObject{9960}{<< /Names [<706167652E333135> \MusicPDFRef{15343} <706167652E333136> \MusicPDFRef{15344} <706167652E333137> \MusicPDFRef{15345} <706167652E333138> \MusicPDFRef{15346} <706167652E333139> \MusicPDFRef{15347} <706167652E3332> \MusicPDFRef{15348}] /Limits [<706167652E333135> <706167652E3332>] >>}
\MusicPDFWriteObject{9961}{<< /Names [<706167652E333230> \MusicPDFRef{15349} <706167652E333231> \MusicPDFRef{15350} <706167652E333232> \MusicPDFRef{15351} <706167652E333233> \MusicPDFRef{15352} <706167652E333234> \MusicPDFRef{15353} <706167652E333235> \MusicPDFRef{15354}] /Limits [<706167652E333230> <706167652E333235>] >>}
\MusicPDFWriteObject{9962}{<< /Names [<706167652E333236> \MusicPDFRef{15355} <706167652E333237> \MusicPDFRef{15356} <706167652E333238> \MusicPDFRef{15357} <706167652E333239> \MusicPDFRef{15358} <706167652E3333> \MusicPDFRef{15359} <706167652E333330> \MusicPDFRef{15360}] /Limits [<706167652E333236> <706167652E333330>] >>}
\MusicPDFWriteObject{9963}{<< /Names [<706167652E333331> \MusicPDFRef{15361} <706167652E333332> \MusicPDFRef{15362} <706167652E333333> \MusicPDFRef{15363} <706167652E333334> \MusicPDFRef{15364} <706167652E333335> \MusicPDFRef{15365} <706167652E333336> \MusicPDFRef{15366}] /Limits [<706167652E333331> <706167652E333336>] >>}
\MusicPDFWriteObject{9964}{<< /Names [<706167652E333337> \MusicPDFRef{15367} <706167652E333338> \MusicPDFRef{15368} <706167652E333339> \MusicPDFRef{15369} <706167652E3334> \MusicPDFRef{15370} <706167652E333430> \MusicPDFRef{15371} <706167652E333431> \MusicPDFRef{15372}] /Limits [<706167652E333337> <706167652E333431>] >>}
\MusicPDFWriteObject{9965}{<< /Names [<706167652E333432> \MusicPDFRef{15373} <706167652E333433> \MusicPDFRef{15374} <706167652E333434> \MusicPDFRef{15375} <706167652E333435> \MusicPDFRef{15376} <706167652E333436> \MusicPDFRef{15377} <706167652E333437> \MusicPDFRef{15378}] /Limits [<706167652E333432> <706167652E333437>] >>}
\MusicPDFWriteObject{9966}{<< /Names [<706167652E333438> \MusicPDFRef{15379} <706167652E333439> \MusicPDFRef{15380} <706167652E3335> \MusicPDFRef{15381} <706167652E333530> \MusicPDFRef{15382} <706167652E333531> \MusicPDFRef{15383} <706167652E333532> \MusicPDFRef{15384}] /Limits [<706167652E333438> <706167652E333532>] >>}
\MusicPDFWriteObject{9967}{<< /Names [<706167652E333533> \MusicPDFRef{15385} <706167652E333534> \MusicPDFRef{15386} <706167652E333535> \MusicPDFRef{15387} <706167652E333536> \MusicPDFRef{15388} <706167652E333537> \MusicPDFRef{15389} <706167652E333538> \MusicPDFRef{15390}] /Limits [<706167652E333533> <706167652E333538>] >>}
\MusicPDFWriteObject{9968}{<< /Names [<706167652E333539> \MusicPDFRef{15391} <706167652E3336> \MusicPDFRef{15392} <706167652E333630> \MusicPDFRef{15393} <706167652E333631> \MusicPDFRef{15394} <706167652E333632> \MusicPDFRef{15395} <706167652E333633> \MusicPDFRef{15396}] /Limits [<706167652E333539> <706167652E333633>] >>}
\MusicPDFWriteObject{9969}{<< /Names [<706167652E333634> \MusicPDFRef{15397} <706167652E333635> \MusicPDFRef{15398} <706167652E333636> \MusicPDFRef{15399} <706167652E333637> \MusicPDFRef{15400} <706167652E333638> \MusicPDFRef{15401} <706167652E333639> \MusicPDFRef{15402}] /Limits [<706167652E333634> <706167652E333639>] >>}
\MusicPDFWriteObject{9970}{<< /Names [<706167652E3337> \MusicPDFRef{15403} <706167652E333730> \MusicPDFRef{15404} <706167652E333731> \MusicPDFRef{15405} <706167652E333732> \MusicPDFRef{15406} <706167652E333733> \MusicPDFRef{15407} <706167652E333734> \MusicPDFRef{15408}] /Limits [<706167652E3337> <706167652E333734>] >>}
\MusicPDFWriteObject{9971}{<< /Names [<706167652E333735> \MusicPDFRef{15409} <706167652E333736> \MusicPDFRef{15410} <706167652E333737> \MusicPDFRef{15411} <706167652E333738> \MusicPDFRef{15412} <706167652E333739> \MusicPDFRef{15413} <706167652E3338> \MusicPDFRef{15414}] /Limits [<706167652E333735> <706167652E3338>] >>}
\MusicPDFWriteObject{9972}{<< /Names [<706167652E333830> \MusicPDFRef{15415} <706167652E333831> \MusicPDFRef{15416} <706167652E333832> \MusicPDFRef{15417} <706167652E333833> \MusicPDFRef{15418} <706167652E333834> \MusicPDFRef{15419} <706167652E333835> \MusicPDFRef{15420}] /Limits [<706167652E333830> <706167652E333835>] >>}
\MusicPDFWriteObject{9973}{<< /Names [<706167652E333836> \MusicPDFRef{15421} <706167652E333837> \MusicPDFRef{15422} <706167652E333838> \MusicPDFRef{15423} <706167652E333839> \MusicPDFRef{15424} <706167652E3339> \MusicPDFRef{15425} <706167652E333930> \MusicPDFRef{15426}] /Limits [<706167652E333836> <706167652E333930>] >>}
\MusicPDFWriteObject{9974}{<< /Names [<706167652E333931> \MusicPDFRef{15427} <706167652E333932> \MusicPDFRef{15428} <706167652E333933> \MusicPDFRef{15429} <706167652E333934> \MusicPDFRef{15430} <706167652E333935> \MusicPDFRef{15431} <706167652E333936> \MusicPDFRef{15432}] /Limits [<706167652E333931> <706167652E333936>] >>}
\MusicPDFWriteObject{9975}{<< /Names [<706167652E333937> \MusicPDFRef{15433} <706167652E333938> \MusicPDFRef{15434} <706167652E333939> \MusicPDFRef{15435} <706167652E34> \MusicPDFRef{15436} <706167652E3430> \MusicPDFRef{15437} <706167652E343030> \MusicPDFRef{15438}] /Limits [<706167652E333937> <706167652E343030>] >>}
\MusicPDFWriteObject{9976}{<< /Names [<706167652E343031> \MusicPDFRef{15439} <706167652E343032> \MusicPDFRef{15440} <706167652E343033> \MusicPDFRef{15441} <706167652E343034> \MusicPDFRef{15442} <706167652E343035> \MusicPDFRef{15443} <706167652E343036> \MusicPDFRef{15444}] /Limits [<706167652E343031> <706167652E343036>] >>}
\MusicPDFWriteObject{9977}{<< /Names [<706167652E343037> \MusicPDFRef{15445} <706167652E343038> \MusicPDFRef{15446} <706167652E343039> \MusicPDFRef{15447} <706167652E3431> \MusicPDFRef{15448} <706167652E343130> \MusicPDFRef{15449} <706167652E343131> \MusicPDFRef{15450}] /Limits [<706167652E343037> <706167652E343131>] >>}
\MusicPDFWriteObject{9978}{<< /Names [<706167652E343132> \MusicPDFRef{15451} <706167652E343133> \MusicPDFRef{15452} <706167652E343134> \MusicPDFRef{15453} <706167652E343135> \MusicPDFRef{15454} <706167652E343136> \MusicPDFRef{15455} <706167652E343137> \MusicPDFRef{15456}] /Limits [<706167652E343132> <706167652E343137>] >>}
\MusicPDFWriteObject{9979}{<< /Names [<706167652E343138> \MusicPDFRef{15457} <706167652E343139> \MusicPDFRef{15458} <706167652E3432> \MusicPDFRef{15459} <706167652E343230> \MusicPDFRef{15460} <706167652E343231> \MusicPDFRef{15461} <706167652E343232> \MusicPDFRef{15462}] /Limits [<706167652E343138> <706167652E343232>] >>}
\MusicPDFWriteObject{9980}{<< /Names [<706167652E343233> \MusicPDFRef{15463} <706167652E343234> \MusicPDFRef{15464} <706167652E343235> \MusicPDFRef{15465} <706167652E343236> \MusicPDFRef{15466} <706167652E343237> \MusicPDFRef{15467} <706167652E343238> \MusicPDFRef{15468}] /Limits [<706167652E343233> <706167652E343238>] >>}
\MusicPDFWriteObject{9981}{<< /Names [<706167652E343239> \MusicPDFRef{15469} <706167652E3433> \MusicPDFRef{15470} <706167652E343330> \MusicPDFRef{15471} <706167652E343331> \MusicPDFRef{15472} <706167652E343332> \MusicPDFRef{15473} <706167652E343333> \MusicPDFRef{15474}] /Limits [<706167652E343239> <706167652E343333>] >>}
\MusicPDFWriteObject{9982}{<< /Names [<706167652E343334> \MusicPDFRef{15475} <706167652E343335> \MusicPDFRef{15476} <706167652E343336> \MusicPDFRef{15477} <706167652E343337> \MusicPDFRef{15478} <706167652E343338> \MusicPDFRef{15479} <706167652E343339> \MusicPDFRef{15480}] /Limits [<706167652E343334> <706167652E343339>] >>}
\MusicPDFWriteObject{9983}{<< /Names [<706167652E3434> \MusicPDFRef{15481} <706167652E343430> \MusicPDFRef{15482} <706167652E343431> \MusicPDFRef{15483} <706167652E343432> \MusicPDFRef{15484} <706167652E343433> \MusicPDFRef{15485} <706167652E343434> \MusicPDFRef{15486}] /Limits [<706167652E3434> <706167652E343434>] >>}
\MusicPDFWriteObject{9984}{<< /Names [<706167652E343435> \MusicPDFRef{15487} <706167652E343436> \MusicPDFRef{15488} <706167652E343437> \MusicPDFRef{15489} <706167652E343438> \MusicPDFRef{15490} <706167652E343439> \MusicPDFRef{15491} <706167652E3435> \MusicPDFRef{15492}] /Limits [<706167652E343435> <706167652E3435>] >>}
\MusicPDFWriteObject{9985}{<< /Names [<706167652E343530> \MusicPDFRef{15493} <706167652E343531> \MusicPDFRef{15494} <706167652E343532> \MusicPDFRef{15495} <706167652E343533> \MusicPDFRef{15496} <706167652E343534> \MusicPDFRef{15497} <706167652E343535> \MusicPDFRef{15498}] /Limits [<706167652E343530> <706167652E343535>] >>}
\MusicPDFWriteObject{9986}{<< /Names [<706167652E343536> \MusicPDFRef{15499} <706167652E343537> \MusicPDFRef{15500} <706167652E343538> \MusicPDFRef{15501} <706167652E343539> \MusicPDFRef{15502} <706167652E3436> \MusicPDFRef{15503} <706167652E343630> \MusicPDFRef{15504}] /Limits [<706167652E343536> <706167652E343630>] >>}
\MusicPDFWriteObject{9987}{<< /Names [<706167652E343631> \MusicPDFRef{15505} <706167652E343632> \MusicPDFRef{15506} <706167652E343633> \MusicPDFRef{15507} <706167652E343634> \MusicPDFRef{15508} <706167652E343635> \MusicPDFRef{15509} <706167652E343636> \MusicPDFRef{15510}] /Limits [<706167652E343631> <706167652E343636>] >>}
\MusicPDFWriteObject{9988}{<< /Names [<706167652E343637> \MusicPDFRef{15511} <706167652E343638> \MusicPDFRef{15512} <706167652E343639> \MusicPDFRef{15513} <706167652E3437> \MusicPDFRef{15514} <706167652E343730> \MusicPDFRef{15515} <706167652E343731> \MusicPDFRef{15516}] /Limits [<706167652E343637> <706167652E343731>] >>}
\MusicPDFWriteObject{9989}{<< /Names [<706167652E343732> \MusicPDFRef{15517} <706167652E343733> \MusicPDFRef{15518} <706167652E343734> \MusicPDFRef{15519} <706167652E343735> \MusicPDFRef{15520} <706167652E343736> \MusicPDFRef{15521} <706167652E343737> \MusicPDFRef{15522}] /Limits [<706167652E343732> <706167652E343737>] >>}
\MusicPDFWriteObject{9990}{<< /Names [<706167652E343738> \MusicPDFRef{15523} <706167652E343739> \MusicPDFRef{15524} <706167652E3438> \MusicPDFRef{15525} <706167652E343830> \MusicPDFRef{15526} <706167652E343831> \MusicPDFRef{15527} <706167652E343832> \MusicPDFRef{15528}] /Limits [<706167652E343738> <706167652E343832>] >>}
\MusicPDFWriteObject{9991}{<< /Names [<706167652E343833> \MusicPDFRef{15529} <706167652E343834> \MusicPDFRef{15530} <706167652E343835> \MusicPDFRef{15531} <706167652E343836> \MusicPDFRef{15532} <706167652E343837> \MusicPDFRef{15533} <706167652E343838> \MusicPDFRef{15534}] /Limits [<706167652E343833> <706167652E343838>] >>}
\MusicPDFWriteObject{9992}{<< /Names [<706167652E343839> \MusicPDFRef{15535} <706167652E3439> \MusicPDFRef{15536} <706167652E343930> \MusicPDFRef{15537} <706167652E343931> \MusicPDFRef{15538} <706167652E343932> \MusicPDFRef{15539} <706167652E343933> \MusicPDFRef{15540}] /Limits [<706167652E343839> <706167652E343933>] >>}
\MusicPDFWriteObject{9993}{<< /Names [<706167652E343934> \MusicPDFRef{15541} <706167652E343935> \MusicPDFRef{15542} <706167652E343936> \MusicPDFRef{15543} <706167652E343937> \MusicPDFRef{15544} <706167652E343938> \MusicPDFRef{15545} <706167652E343939> \MusicPDFRef{15546}] /Limits [<706167652E343934> <706167652E343939>] >>}
\MusicPDFWriteObject{9994}{<< /Names [<706167652E35> \MusicPDFRef{15547} <706167652E3530> \MusicPDFRef{15548} <706167652E353030> \MusicPDFRef{15549} <706167652E353031> \MusicPDFRef{15550} <706167652E353032> \MusicPDFRef{15551} <706167652E353033> \MusicPDFRef{15552}] /Limits [<706167652E35> <706167652E353033>] >>}
\MusicPDFWriteObject{9995}{<< /Names [<706167652E353034> \MusicPDFRef{15553} <706167652E353035> \MusicPDFRef{15554} <706167652E353036> \MusicPDFRef{15555} <706167652E353037> \MusicPDFRef{15556} <706167652E353038> \MusicPDFRef{15557} <706167652E353039> \MusicPDFRef{15558}] /Limits [<706167652E353034> <706167652E353039>] >>}
\MusicPDFWriteObject{9996}{<< /Names [<706167652E3531> \MusicPDFRef{15559} <706167652E353130> \MusicPDFRef{15560} <706167652E353131> \MusicPDFRef{15561} <706167652E353132> \MusicPDFRef{15562} <706167652E353133> \MusicPDFRef{15563} <706167652E353134> \MusicPDFRef{15564}] /Limits [<706167652E3531> <706167652E353134>] >>}
\MusicPDFWriteObject{9997}{<< /Names [<706167652E353135> \MusicPDFRef{15565} <706167652E353136> \MusicPDFRef{15566} <706167652E353137> \MusicPDFRef{15567} <706167652E353138> \MusicPDFRef{15568} <706167652E353139> \MusicPDFRef{15569} <706167652E3532> \MusicPDFRef{15570}] /Limits [<706167652E353135> <706167652E3532>] >>}
\MusicPDFWriteObject{9998}{<< /Names [<706167652E353230> \MusicPDFRef{15571} <706167652E353231> \MusicPDFRef{15572} <706167652E353232> \MusicPDFRef{15573} <706167652E353233> \MusicPDFRef{15574} <706167652E353234> \MusicPDFRef{15575} <706167652E353235> \MusicPDFRef{15576}] /Limits [<706167652E353230> <706167652E353235>] >>}
\MusicPDFWriteObject{9999}{<< /Names [<706167652E353236> \MusicPDFRef{15577} <706167652E353237> \MusicPDFRef{15578} <706167652E353238> \MusicPDFRef{15579} <706167652E353239> \MusicPDFRef{15580} <706167652E3533> \MusicPDFRef{15581} <706167652E353330> \MusicPDFRef{15582}] /Limits [<706167652E353236> <706167652E353330>] >>}
\MusicPDFWriteObject{10000}{<< /Names [<706167652E353331> \MusicPDFRef{15583} <706167652E353332> \MusicPDFRef{15584} <706167652E353333> \MusicPDFRef{15585} <706167652E353334> \MusicPDFRef{15586} <706167652E353335> \MusicPDFRef{15587} <706167652E353336> \MusicPDFRef{15588}] /Limits [<706167652E353331> <706167652E353336>] >>}
\MusicPDFWriteObject{10001}{<< /Names [<706167652E353337> \MusicPDFRef{15589} <706167652E353338> \MusicPDFRef{15590} <706167652E353339> \MusicPDFRef{15591} <706167652E3534> \MusicPDFRef{15592} <706167652E353430> \MusicPDFRef{15593} <706167652E353431> \MusicPDFRef{15594}] /Limits [<706167652E353337> <706167652E353431>] >>}
\MusicPDFWriteObject{10002}{<< /Names [<706167652E353432> \MusicPDFRef{15595} <706167652E353433> \MusicPDFRef{15596} <706167652E353434> \MusicPDFRef{15597} <706167652E353435> \MusicPDFRef{15598} <706167652E353436> \MusicPDFRef{15599} <706167652E353437> \MusicPDFRef{15600}] /Limits [<706167652E353432> <706167652E353437>] >>}
\MusicPDFWriteObject{10003}{<< /Names [<706167652E353438> \MusicPDFRef{15601} <706167652E353439> \MusicPDFRef{15602} <706167652E3535> \MusicPDFRef{15603} <706167652E353530> \MusicPDFRef{15604} <706167652E353531> \MusicPDFRef{15605} <706167652E353532> \MusicPDFRef{15606}] /Limits [<706167652E353438> <706167652E353532>] >>}
\MusicPDFWriteObject{10004}{<< /Names [<706167652E353533> \MusicPDFRef{15607} <706167652E353534> \MusicPDFRef{15608} <706167652E353535> \MusicPDFRef{15609} <706167652E353536> \MusicPDFRef{15610} <706167652E353537> \MusicPDFRef{15611} <706167652E353538> \MusicPDFRef{15612}] /Limits [<706167652E353533> <706167652E353538>] >>}
\MusicPDFWriteObject{10005}{<< /Names [<706167652E353539> \MusicPDFRef{15613} <706167652E3536> \MusicPDFRef{15614} <706167652E353630> \MusicPDFRef{15615} <706167652E353631> \MusicPDFRef{15616} <706167652E353632> \MusicPDFRef{15617} <706167652E353633> \MusicPDFRef{15618}] /Limits [<706167652E353539> <706167652E353633>] >>}
\MusicPDFWriteObject{10006}{<< /Names [<706167652E353634> \MusicPDFRef{15619} <706167652E353635> \MusicPDFRef{15620} <706167652E353636> \MusicPDFRef{15621} <706167652E353637> \MusicPDFRef{15622} <706167652E353638> \MusicPDFRef{15623} <706167652E353639> \MusicPDFRef{15624}] /Limits [<706167652E353634> <706167652E353639>] >>}
\MusicPDFWriteObject{10007}{<< /Names [<706167652E3537> \MusicPDFRef{15625} <706167652E353730> \MusicPDFRef{15626} <706167652E353731> \MusicPDFRef{15627} <706167652E353732> \MusicPDFRef{15628} <706167652E353733> \MusicPDFRef{15629} <706167652E353734> \MusicPDFRef{15630}] /Limits [<706167652E3537> <706167652E353734>] >>}
\MusicPDFWriteObject{10008}{<< /Names [<706167652E353735> \MusicPDFRef{15631} <706167652E353736> \MusicPDFRef{15632} <706167652E353737> \MusicPDFRef{15633} <706167652E353738> \MusicPDFRef{15634} <706167652E353739> \MusicPDFRef{15635} <706167652E3538> \MusicPDFRef{15636}] /Limits [<706167652E353735> <706167652E3538>] >>}
\MusicPDFWriteObject{10009}{<< /Names [<706167652E353830> \MusicPDFRef{15637} <706167652E353831> \MusicPDFRef{15638} <706167652E353832> \MusicPDFRef{15639} <706167652E353833> \MusicPDFRef{15640} <706167652E353834> \MusicPDFRef{15641} <706167652E353835> \MusicPDFRef{15642}] /Limits [<706167652E353830> <706167652E353835>] >>}
\MusicPDFWriteObject{10010}{<< /Names [<706167652E353836> \MusicPDFRef{15643} <706167652E353837> \MusicPDFRef{15644} <706167652E353838> \MusicPDFRef{15645} <706167652E353839> \MusicPDFRef{15646} <706167652E3539> \MusicPDFRef{15647} <706167652E353930> \MusicPDFRef{15648}] /Limits [<706167652E353836> <706167652E353930>] >>}
\MusicPDFWriteObject{10011}{<< /Names [<706167652E353931> \MusicPDFRef{15649} <706167652E353932> \MusicPDFRef{15650} <706167652E353933> \MusicPDFRef{15651} <706167652E353934> \MusicPDFRef{15652} <706167652E353935> \MusicPDFRef{15653} <706167652E353936> \MusicPDFRef{15654}] /Limits [<706167652E353931> <706167652E353936>] >>}
\MusicPDFWriteObject{10012}{<< /Names [<706167652E353937> \MusicPDFRef{15655} <706167652E353938> \MusicPDFRef{15656} <706167652E353939> \MusicPDFRef{15657} <706167652E36> \MusicPDFRef{15658} <706167652E3630> \MusicPDFRef{15659} <706167652E363030> \MusicPDFRef{15660}] /Limits [<706167652E353937> <706167652E363030>] >>}
\MusicPDFWriteObject{10013}{<< /Names [<706167652E363031> \MusicPDFRef{15661} <706167652E363032> \MusicPDFRef{15662} <706167652E363033> \MusicPDFRef{15663} <706167652E363034> \MusicPDFRef{15664} <706167652E363035> \MusicPDFRef{15665} <706167652E363036> \MusicPDFRef{15666}] /Limits [<706167652E363031> <706167652E363036>] >>}
\MusicPDFWriteObject{10014}{<< /Names [<706167652E363037> \MusicPDFRef{15667} <706167652E363038> \MusicPDFRef{15668} <706167652E363039> \MusicPDFRef{15669} <706167652E3631> \MusicPDFRef{15670} <706167652E363130> \MusicPDFRef{15671} <706167652E363131> \MusicPDFRef{15672}] /Limits [<706167652E363037> <706167652E363131>] >>}
\MusicPDFWriteObject{10015}{<< /Names [<706167652E363132> \MusicPDFRef{15673} <706167652E363133> \MusicPDFRef{15674} <706167652E363134> \MusicPDFRef{15675} <706167652E363135> \MusicPDFRef{15676} <706167652E363136> \MusicPDFRef{15677} <706167652E363137> \MusicPDFRef{15678}] /Limits [<706167652E363132> <706167652E363137>] >>}
\MusicPDFWriteObject{10016}{<< /Names [<706167652E363138> \MusicPDFRef{15679} <706167652E363139> \MusicPDFRef{15680} <706167652E3632> \MusicPDFRef{15681} <706167652E363230> \MusicPDFRef{15682} <706167652E363231> \MusicPDFRef{15683} <706167652E363232> \MusicPDFRef{15684}] /Limits [<706167652E363138> <706167652E363232>] >>}
\MusicPDFWriteObject{10017}{<< /Names [<706167652E363233> \MusicPDFRef{15685} <706167652E363234> \MusicPDFRef{15686} <706167652E363235> \MusicPDFRef{15687} <706167652E363236> \MusicPDFRef{15688} <706167652E363237> \MusicPDFRef{15689} <706167652E363238> \MusicPDFRef{15690}] /Limits [<706167652E363233> <706167652E363238>] >>}
\MusicPDFWriteObject{10018}{<< /Names [<706167652E363239> \MusicPDFRef{15691} <706167652E3633> \MusicPDFRef{15692} <706167652E363330> \MusicPDFRef{15693} <706167652E363331> \MusicPDFRef{15694} <706167652E363332> \MusicPDFRef{15695} <706167652E363333> \MusicPDFRef{15696}] /Limits [<706167652E363239> <706167652E363333>] >>}
\MusicPDFWriteObject{10019}{<< /Names [<706167652E363334> \MusicPDFRef{15697} <706167652E363335> \MusicPDFRef{15698} <706167652E363336> \MusicPDFRef{15699} <706167652E363337> \MusicPDFRef{15700} <706167652E363338> \MusicPDFRef{15701} <706167652E363339> \MusicPDFRef{15702}] /Limits [<706167652E363334> <706167652E363339>] >>}
\MusicPDFWriteObject{10020}{<< /Names [<706167652E3634> \MusicPDFRef{15703} <706167652E363430> \MusicPDFRef{15704} <706167652E363431> \MusicPDFRef{15705} <706167652E363432> \MusicPDFRef{15706} <706167652E363433> \MusicPDFRef{15707} <706167652E363434> \MusicPDFRef{15708}] /Limits [<706167652E3634> <706167652E363434>] >>}
\MusicPDFWriteObject{10021}{<< /Names [<706167652E363435> \MusicPDFRef{15709} <706167652E363436> \MusicPDFRef{15710} <706167652E363437> \MusicPDFRef{15711} <706167652E363438> \MusicPDFRef{15712} <706167652E363439> \MusicPDFRef{15713} <706167652E3635> \MusicPDFRef{15714}] /Limits [<706167652E363435> <706167652E3635>] >>}
\MusicPDFWriteObject{10022}{<< /Names [<706167652E363530> \MusicPDFRef{15715} <706167652E363531> \MusicPDFRef{15716} <706167652E363532> \MusicPDFRef{15717} <706167652E363533> \MusicPDFRef{15718} <706167652E363534> \MusicPDFRef{15719} <706167652E363535> \MusicPDFRef{15720}] /Limits [<706167652E363530> <706167652E363535>] >>}
\MusicPDFWriteObject{10023}{<< /Names [<706167652E363536> \MusicPDFRef{15721} <706167652E363537> \MusicPDFRef{15722} <706167652E363538> \MusicPDFRef{15723} <706167652E363539> \MusicPDFRef{15724} <706167652E3636> \MusicPDFRef{15725} <706167652E363630> \MusicPDFRef{15726}] /Limits [<706167652E363536> <706167652E363630>] >>}
\MusicPDFWriteObject{10024}{<< /Names [<706167652E363631> \MusicPDFRef{15727} <706167652E363632> \MusicPDFRef{15728} <706167652E363633> \MusicPDFRef{15729} <706167652E363634> \MusicPDFRef{15730} <706167652E363635> \MusicPDFRef{15731} <706167652E363636> \MusicPDFRef{15732}] /Limits [<706167652E363631> <706167652E363636>] >>}
\MusicPDFWriteObject{10025}{<< /Names [<706167652E363637> \MusicPDFRef{15733} <706167652E363638> \MusicPDFRef{15734} <706167652E363639> \MusicPDFRef{15735} <706167652E3637> \MusicPDFRef{15736} <706167652E363730> \MusicPDFRef{15737} <706167652E363731> \MusicPDFRef{15738}] /Limits [<706167652E363637> <706167652E363731>] >>}
\MusicPDFWriteObject{10026}{<< /Names [<706167652E363732> \MusicPDFRef{15739} <706167652E363733> \MusicPDFRef{15740} <706167652E363734> \MusicPDFRef{15741} <706167652E363735> \MusicPDFRef{15742} <706167652E363736> \MusicPDFRef{15743} <706167652E363737> \MusicPDFRef{15744}] /Limits [<706167652E363732> <706167652E363737>] >>}
\MusicPDFWriteObject{10027}{<< /Names [<706167652E363738> \MusicPDFRef{15745} <706167652E363739> \MusicPDFRef{15746} <706167652E3638> \MusicPDFRef{15747} <706167652E363830> \MusicPDFRef{15748} <706167652E363831> \MusicPDFRef{15749} <706167652E363832> \MusicPDFRef{15750}] /Limits [<706167652E363738> <706167652E363832>] >>}
\MusicPDFWriteObject{10028}{<< /Names [<706167652E363833> \MusicPDFRef{15751} <706167652E363834> \MusicPDFRef{15752} <706167652E363835> \MusicPDFRef{15753} <706167652E363836> \MusicPDFRef{15754} <706167652E363837> \MusicPDFRef{15755} <706167652E363838> \MusicPDFRef{15756}] /Limits [<706167652E363833> <706167652E363838>] >>}
\MusicPDFWriteObject{10029}{<< /Names [<706167652E363839> \MusicPDFRef{15757} <706167652E3639> \MusicPDFRef{15758} <706167652E363930> \MusicPDFRef{15759} <706167652E363931> \MusicPDFRef{15760} <706167652E363932> \MusicPDFRef{15761} <706167652E363933> \MusicPDFRef{15762}] /Limits [<706167652E363839> <706167652E363933>] >>}
\MusicPDFWriteObject{10030}{<< /Names [<706167652E363934> \MusicPDFRef{15763} <706167652E363935> \MusicPDFRef{15764} <706167652E363936> \MusicPDFRef{15765} <706167652E363937> \MusicPDFRef{15766} <706167652E363938> \MusicPDFRef{15767} <706167652E363939> \MusicPDFRef{15768}] /Limits [<706167652E363934> <706167652E363939>] >>}
\MusicPDFWriteObject{10031}{<< /Names [<706167652E37> \MusicPDFRef{15769} <706167652E3730> \MusicPDFRef{15770} <706167652E373030> \MusicPDFRef{15771} <706167652E373031> \MusicPDFRef{15772} <706167652E373032> \MusicPDFRef{15773} <706167652E373033> \MusicPDFRef{15774}] /Limits [<706167652E37> <706167652E373033>] >>}
\MusicPDFWriteObject{10032}{<< /Names [<706167652E373034> \MusicPDFRef{15775} <706167652E373035> \MusicPDFRef{15776} <706167652E373036> \MusicPDFRef{15777} <706167652E373037> \MusicPDFRef{15778} <706167652E373038> \MusicPDFRef{15779} <706167652E373039> \MusicPDFRef{15780}] /Limits [<706167652E373034> <706167652E373039>] >>}
\MusicPDFWriteObject{10033}{<< /Names [<706167652E3731> \MusicPDFRef{15781} <706167652E373130> \MusicPDFRef{15782} <706167652E373131> \MusicPDFRef{15783} <706167652E373132> \MusicPDFRef{15784} <706167652E373133> \MusicPDFRef{15785} <706167652E373134> \MusicPDFRef{15786}] /Limits [<706167652E3731> <706167652E373134>] >>}
\MusicPDFWriteObject{10034}{<< /Names [<706167652E373135> \MusicPDFRef{15787} <706167652E373136> \MusicPDFRef{15788} <706167652E373137> \MusicPDFRef{15789} <706167652E373138> \MusicPDFRef{15790} <706167652E373139> \MusicPDFRef{15791} <706167652E3732> \MusicPDFRef{15792}] /Limits [<706167652E373135> <706167652E3732>] >>}
\MusicPDFWriteObject{10035}{<< /Names [<706167652E373230> \MusicPDFRef{15793} <706167652E373231> \MusicPDFRef{15794} <706167652E373232> \MusicPDFRef{15795} <706167652E373233> \MusicPDFRef{15796} <706167652E373234> \MusicPDFRef{15797} <706167652E373235> \MusicPDFRef{15798}] /Limits [<706167652E373230> <706167652E373235>] >>}
\MusicPDFWriteObject{10036}{<< /Names [<706167652E373236> \MusicPDFRef{15799} <706167652E373237> \MusicPDFRef{15800} <706167652E373238> \MusicPDFRef{15801} <706167652E373239> \MusicPDFRef{15802} <706167652E3733> \MusicPDFRef{15803} <706167652E373330> \MusicPDFRef{15804}] /Limits [<706167652E373236> <706167652E373330>] >>}
\MusicPDFWriteObject{10037}{<< /Names [<706167652E373331> \MusicPDFRef{15805} <706167652E373332> \MusicPDFRef{15806} <706167652E373333> \MusicPDFRef{15807} <706167652E373334> \MusicPDFRef{15808} <706167652E373335> \MusicPDFRef{15809} <706167652E373336> \MusicPDFRef{15810}] /Limits [<706167652E373331> <706167652E373336>] >>}
\MusicPDFWriteObject{10038}{<< /Names [<706167652E373337> \MusicPDFRef{15811} <706167652E373338> \MusicPDFRef{15812} <706167652E373339> \MusicPDFRef{15813} <706167652E3734> \MusicPDFRef{15814} <706167652E373430> \MusicPDFRef{15815} <706167652E373431> \MusicPDFRef{15816}] /Limits [<706167652E373337> <706167652E373431>] >>}
\MusicPDFWriteObject{10039}{<< /Names [<706167652E373432> \MusicPDFRef{15817} <706167652E373433> \MusicPDFRef{15818} <706167652E373434> \MusicPDFRef{15819} <706167652E373435> \MusicPDFRef{15820} <706167652E373436> \MusicPDFRef{15821} <706167652E373437> \MusicPDFRef{15822}] /Limits [<706167652E373432> <706167652E373437>] >>}
\MusicPDFWriteObject{10040}{<< /Names [<706167652E373438> \MusicPDFRef{15823} <706167652E373439> \MusicPDFRef{15824} <706167652E3735> \MusicPDFRef{15825} <706167652E373530> \MusicPDFRef{15826} <706167652E373531> \MusicPDFRef{15827} <706167652E373532> \MusicPDFRef{15828}] /Limits [<706167652E373438> <706167652E373532>] >>}
\MusicPDFWriteObject{10041}{<< /Names [<706167652E373533> \MusicPDFRef{15829} <706167652E373534> \MusicPDFRef{15830} <706167652E373535> \MusicPDFRef{15831} <706167652E373536> \MusicPDFRef{15832} <706167652E373537> \MusicPDFRef{15833} <706167652E373538> \MusicPDFRef{15834}] /Limits [<706167652E373533> <706167652E373538>] >>}
\MusicPDFWriteObject{10042}{<< /Names [<706167652E373539> \MusicPDFRef{15835} <706167652E3736> \MusicPDFRef{15836} <706167652E373630> \MusicPDFRef{15837} <706167652E373631> \MusicPDFRef{15838} <706167652E373632> \MusicPDFRef{15839} <706167652E373633> \MusicPDFRef{15840}] /Limits [<706167652E373539> <706167652E373633>] >>}
\MusicPDFWriteObject{10043}{<< /Names [<706167652E373634> \MusicPDFRef{15841} <706167652E373635> \MusicPDFRef{15842} <706167652E373636> \MusicPDFRef{15843} <706167652E373637> \MusicPDFRef{15844} <706167652E373638> \MusicPDFRef{15845} <706167652E373639> \MusicPDFRef{15846}] /Limits [<706167652E373634> <706167652E373639>] >>}
\MusicPDFWriteObject{10044}{<< /Names [<706167652E3737> \MusicPDFRef{15847} <706167652E373730> \MusicPDFRef{15848} <706167652E373731> \MusicPDFRef{15849} <706167652E373732> \MusicPDFRef{15850} <706167652E373733> \MusicPDFRef{15851} <706167652E373734> \MusicPDFRef{15852}] /Limits [<706167652E3737> <706167652E373734>] >>}
\MusicPDFWriteObject{10045}{<< /Names [<706167652E373735> \MusicPDFRef{15853} <706167652E373736> \MusicPDFRef{15854} <706167652E373737> \MusicPDFRef{15855} <706167652E373738> \MusicPDFRef{15856} <706167652E373739> \MusicPDFRef{15857} <706167652E3738> \MusicPDFRef{15858}] /Limits [<706167652E373735> <706167652E3738>] >>}
\MusicPDFWriteObject{10046}{<< /Names [<706167652E373830> \MusicPDFRef{15859} <706167652E373831> \MusicPDFRef{15860} <706167652E373832> \MusicPDFRef{15861} <706167652E373833> \MusicPDFRef{15862} <706167652E373834> \MusicPDFRef{15863} <706167652E373835> \MusicPDFRef{15864}] /Limits [<706167652E373830> <706167652E373835>] >>}
\MusicPDFWriteObject{10047}{<< /Names [<706167652E373836> \MusicPDFRef{15865} <706167652E373837> \MusicPDFRef{15866} <706167652E373838> \MusicPDFRef{15867} <706167652E373839> \MusicPDFRef{15868} <706167652E3739> \MusicPDFRef{15869} <706167652E373930> \MusicPDFRef{15870}] /Limits [<706167652E373836> <706167652E373930>] >>}
\MusicPDFWriteObject{10048}{<< /Names [<706167652E373931> \MusicPDFRef{15871} <706167652E373932> \MusicPDFRef{15872} <706167652E373933> \MusicPDFRef{15873} <706167652E373934> \MusicPDFRef{15874} <706167652E373935> \MusicPDFRef{15875} <706167652E373936> \MusicPDFRef{15876}] /Limits [<706167652E373931> <706167652E373936>] >>}
\MusicPDFWriteObject{10049}{<< /Names [<706167652E373937> \MusicPDFRef{15877} <706167652E373938> \MusicPDFRef{15878} <706167652E373939> \MusicPDFRef{15879} <706167652E38> \MusicPDFRef{15880} <706167652E3830> \MusicPDFRef{15881} <706167652E383030> \MusicPDFRef{15882}] /Limits [<706167652E373937> <706167652E383030>] >>}
\MusicPDFWriteObject{10050}{<< /Names [<706167652E383031> \MusicPDFRef{15883} <706167652E383032> \MusicPDFRef{15884} <706167652E383033> \MusicPDFRef{15885} <706167652E383034> \MusicPDFRef{15886} <706167652E383035> \MusicPDFRef{15887} <706167652E383036> \MusicPDFRef{15888}] /Limits [<706167652E383031> <706167652E383036>] >>}
\MusicPDFWriteObject{10051}{<< /Names [<706167652E383037> \MusicPDFRef{15889} <706167652E383038> \MusicPDFRef{15890} <706167652E383039> \MusicPDFRef{15891} <706167652E3831> \MusicPDFRef{15892} <706167652E383130> \MusicPDFRef{15893} <706167652E383131> \MusicPDFRef{15894}] /Limits [<706167652E383037> <706167652E383131>] >>}
\MusicPDFWriteObject{10052}{<< /Names [<706167652E383132> \MusicPDFRef{15895} <706167652E383133> \MusicPDFRef{15896} <706167652E383134> \MusicPDFRef{15897} <706167652E383135> \MusicPDFRef{15898} <706167652E383136> \MusicPDFRef{15899} <706167652E383137> \MusicPDFRef{15900}] /Limits [<706167652E383132> <706167652E383137>] >>}
\MusicPDFWriteObject{10053}{<< /Names [<706167652E383138> \MusicPDFRef{15901} <706167652E383139> \MusicPDFRef{15902} <706167652E3832> \MusicPDFRef{15903} <706167652E383230> \MusicPDFRef{15904} <706167652E383231> \MusicPDFRef{15905} <706167652E383232> \MusicPDFRef{15906}] /Limits [<706167652E383138> <706167652E383232>] >>}
\MusicPDFWriteObject{10054}{<< /Names [<706167652E383233> \MusicPDFRef{15907} <706167652E383234> \MusicPDFRef{15908} <706167652E383235> \MusicPDFRef{15909} <706167652E383236> \MusicPDFRef{15910} <706167652E383237> \MusicPDFRef{15911} <706167652E383238> \MusicPDFRef{15912}] /Limits [<706167652E383233> <706167652E383238>] >>}
\MusicPDFWriteObject{10055}{<< /Names [<706167652E383239> \MusicPDFRef{15913} <706167652E3833> \MusicPDFRef{15914} <706167652E383330> \MusicPDFRef{15915} <706167652E383331> \MusicPDFRef{15916} <706167652E383332> \MusicPDFRef{15917} <706167652E383333> \MusicPDFRef{15918}] /Limits [<706167652E383239> <706167652E383333>] >>}
\MusicPDFWriteObject{10056}{<< /Names [<706167652E383334> \MusicPDFRef{15919} <706167652E383335> \MusicPDFRef{15920} <706167652E383336> \MusicPDFRef{15921} <706167652E383337> \MusicPDFRef{15922} <706167652E383338> \MusicPDFRef{15923} <706167652E383339> \MusicPDFRef{15924}] /Limits [<706167652E383334> <706167652E383339>] >>}
\MusicPDFWriteObject{10057}{<< /Names [<706167652E3834> \MusicPDFRef{15925} <706167652E383430> \MusicPDFRef{15926} <706167652E383431> \MusicPDFRef{15927} <706167652E383432> \MusicPDFRef{15928} <706167652E383433> \MusicPDFRef{15929} <706167652E383434> \MusicPDFRef{15930}] /Limits [<706167652E3834> <706167652E383434>] >>}
\MusicPDFWriteObject{10058}{<< /Names [<706167652E383435> \MusicPDFRef{15931} <706167652E383436> \MusicPDFRef{15932} <706167652E383437> \MusicPDFRef{15933} <706167652E383438> \MusicPDFRef{15934} <706167652E383439> \MusicPDFRef{15935} <706167652E3835> \MusicPDFRef{15936}] /Limits [<706167652E383435> <706167652E3835>] >>}
\MusicPDFWriteObject{10059}{<< /Names [<706167652E383530> \MusicPDFRef{15937} <706167652E383531> \MusicPDFRef{15938} <706167652E383532> \MusicPDFRef{15939} <706167652E383533> \MusicPDFRef{15940} <706167652E383534> \MusicPDFRef{15941} <706167652E383535> \MusicPDFRef{15942}] /Limits [<706167652E383530> <706167652E383535>] >>}
\MusicPDFWriteObject{10060}{<< /Names [<706167652E383536> \MusicPDFRef{15943} <706167652E383537> \MusicPDFRef{15944} <706167652E383538> \MusicPDFRef{15945} <706167652E383539> \MusicPDFRef{15946} <706167652E3836> \MusicPDFRef{15947} <706167652E383630> \MusicPDFRef{15948}] /Limits [<706167652E383536> <706167652E383630>] >>}
\MusicPDFWriteObject{10061}{<< /Names [<706167652E383631> \MusicPDFRef{15949} <706167652E383632> \MusicPDFRef{15950} <706167652E383633> \MusicPDFRef{15951} <706167652E383634> \MusicPDFRef{15952} <706167652E383635> \MusicPDFRef{15953} <706167652E383636> \MusicPDFRef{15954}] /Limits [<706167652E383631> <706167652E383636>] >>}
\MusicPDFWriteObject{10062}{<< /Names [<706167652E383637> \MusicPDFRef{15955} <706167652E383638> \MusicPDFRef{15956} <706167652E383639> \MusicPDFRef{15957} <706167652E3837> \MusicPDFRef{15958} <706167652E383730> \MusicPDFRef{15959} <706167652E383731> \MusicPDFRef{15960}] /Limits [<706167652E383637> <706167652E383731>] >>}
\MusicPDFWriteObject{10063}{<< /Names [<706167652E383732> \MusicPDFRef{15961} <706167652E383733> \MusicPDFRef{15962} <706167652E383734> \MusicPDFRef{15963} <706167652E383735> \MusicPDFRef{15964} <706167652E383736> \MusicPDFRef{15965} <706167652E383737> \MusicPDFRef{15966}] /Limits [<706167652E383732> <706167652E383737>] >>}
\MusicPDFWriteObject{10064}{<< /Names [<706167652E383738> \MusicPDFRef{15967} <706167652E383739> \MusicPDFRef{15968} <706167652E3838> \MusicPDFRef{15969} <706167652E383830> \MusicPDFRef{15970} <706167652E383831> \MusicPDFRef{15971} <706167652E383832> \MusicPDFRef{15972}] /Limits [<706167652E383738> <706167652E383832>] >>}
\MusicPDFWriteObject{10065}{<< /Names [<706167652E383833> \MusicPDFRef{15973} <706167652E383834> \MusicPDFRef{15974} <706167652E383835> \MusicPDFRef{15975} <706167652E383836> \MusicPDFRef{15976} <706167652E383837> \MusicPDFRef{15977} <706167652E383838> \MusicPDFRef{15978}] /Limits [<706167652E383833> <706167652E383838>] >>}
\MusicPDFWriteObject{10066}{<< /Names [<706167652E383839> \MusicPDFRef{15979} <706167652E3839> \MusicPDFRef{15980} <706167652E383930> \MusicPDFRef{15981} <706167652E383931> \MusicPDFRef{15982} <706167652E383932> \MusicPDFRef{15983} <706167652E383933> \MusicPDFRef{15984}] /Limits [<706167652E383839> <706167652E383933>] >>}
\MusicPDFWriteObject{10067}{<< /Names [<706167652E383934> \MusicPDFRef{15985} <706167652E383935> \MusicPDFRef{15986} <706167652E383936> \MusicPDFRef{15987} <706167652E383937> \MusicPDFRef{15988} <706167652E383938> \MusicPDFRef{15989} <706167652E383939> \MusicPDFRef{15990}] /Limits [<706167652E383934> <706167652E383939>] >>}
\MusicPDFWriteObject{10068}{<< /Names [<706167652E39> \MusicPDFRef{15991} <706167652E3930> \MusicPDFRef{15992} <706167652E393030> \MusicPDFRef{15993} <706167652E393031> \MusicPDFRef{15994} <706167652E393032> \MusicPDFRef{15995} <706167652E393033> \MusicPDFRef{15996}] /Limits [<706167652E39> <706167652E393033>] >>}
\MusicPDFWriteObject{10069}{<< /Names [<706167652E393034> \MusicPDFRef{15997} <706167652E393035> \MusicPDFRef{15998} <706167652E393036> \MusicPDFRef{15999} <706167652E393037> \MusicPDFRef{16000} <706167652E393038> \MusicPDFRef{16001} <706167652E393039> \MusicPDFRef{16002}] /Limits [<706167652E393034> <706167652E393039>] >>}
\MusicPDFWriteObject{10070}{<< /Names [<706167652E3931> \MusicPDFRef{16003} <706167652E393130> \MusicPDFRef{16004} <706167652E393131> \MusicPDFRef{16005} <706167652E393132> \MusicPDFRef{16006} <706167652E393133> \MusicPDFRef{16007} <706167652E393134> \MusicPDFRef{16008}] /Limits [<706167652E3931> <706167652E393134>] >>}
\MusicPDFWriteObject{10071}{<< /Names [<706167652E393135> \MusicPDFRef{16009} <706167652E393136> \MusicPDFRef{16010} <706167652E393137> \MusicPDFRef{16011} <706167652E393138> \MusicPDFRef{16012} <706167652E393139> \MusicPDFRef{16013} <706167652E3932> \MusicPDFRef{16014}] /Limits [<706167652E393135> <706167652E3932>] >>}
\MusicPDFWriteObject{10072}{<< /Names [<706167652E393230> \MusicPDFRef{16015} <706167652E393231> \MusicPDFRef{16016} <706167652E393232> \MusicPDFRef{16017} <706167652E393233> \MusicPDFRef{16018} <706167652E393234> \MusicPDFRef{16019} <706167652E393235> \MusicPDFRef{16020}] /Limits [<706167652E393230> <706167652E393235>] >>}
\MusicPDFWriteObject{10073}{<< /Names [<706167652E393236> \MusicPDFRef{16021} <706167652E393237> \MusicPDFRef{16022} <706167652E393238> \MusicPDFRef{16023} <706167652E393239> \MusicPDFRef{16024} <706167652E3933> \MusicPDFRef{16025} <706167652E393330> \MusicPDFRef{16026}] /Limits [<706167652E393236> <706167652E393330>] >>}
\MusicPDFWriteObject{10074}{<< /Names [<706167652E393331> \MusicPDFRef{16027} <706167652E393332> \MusicPDFRef{16028} <706167652E393333> \MusicPDFRef{16029} <706167652E393334> \MusicPDFRef{16030} <706167652E393335> \MusicPDFRef{16031} <706167652E393336> \MusicPDFRef{16032}] /Limits [<706167652E393331> <706167652E393336>] >>}
\MusicPDFWriteObject{10075}{<< /Names [<706167652E393337> \MusicPDFRef{16033} <706167652E393338> \MusicPDFRef{16034} <706167652E393339> \MusicPDFRef{16035} <706167652E3934> \MusicPDFRef{16036} <706167652E393430> \MusicPDFRef{16037} <706167652E3935> \MusicPDFRef{16038}] /Limits [<706167652E393337> <706167652E3935>] >>}
\MusicPDFWriteObject{10076}{<< /Names [<706167652E3936> \MusicPDFRef{16039} <706167652E3937> \MusicPDFRef{16040} <706167652E3938> \MusicPDFRef{16041} <706167652E3939> \MusicPDFRef{16042} <73656374696F6E2A2E31> \MusicPDFRef{16043} <73656374696F6E2E312E31> \MusicPDFRef{16044}] /Limits [<706167652E3936> <73656374696F6E2E312E31>] >>}
\MusicPDFWriteObject{10077}{<< /Names [<73656374696F6E2E312E32> \MusicPDFRef{16045} <73656374696F6E2E312E33> \MusicPDFRef{16046} <73656374696F6E2E312E34> \MusicPDFRef{16047} <73656374696F6E2E312E35> \MusicPDFRef{16048} <73656374696F6E2E312E36> \MusicPDFRef{16049} <73656374696F6E2E312E37> \MusicPDFRef{16050}] /Limits [<73656374696F6E2E312E32> <73656374696F6E2E312E37>] >>}
\MusicPDFWriteObject{10078}{<< /Names [<73656374696F6E2E312E38> \MusicPDFRef{16051} <73656374696F6E2E31302E31> \MusicPDFRef{16052} <73656374696F6E2E31302E32> \MusicPDFRef{16053} <73656374696F6E2E31302E33> \MusicPDFRef{16054} <73656374696F6E2E31302E34> \MusicPDFRef{16055} <73656374696F6E2E31312E31> \MusicPDFRef{16056}] /Limits [<73656374696F6E2E312E38> <73656374696F6E2E31312E31>] >>}
\MusicPDFWriteObject{10079}{<< /Names [<73656374696F6E2E31312E32> \MusicPDFRef{16057} <73656374696F6E2E31312E33> \MusicPDFRef{16058} <73656374696F6E2E31322E31> \MusicPDFRef{16059} <73656374696F6E2E31322E32> \MusicPDFRef{16060} <73656374696F6E2E31322E33> \MusicPDFRef{16061} <73656374696F6E2E31322E34> \MusicPDFRef{16062}] /Limits [<73656374696F6E2E31312E32> <73656374696F6E2E31322E34>] >>}
\MusicPDFWriteObject{10080}{<< /Names [<73656374696F6E2E31322E35> \MusicPDFRef{16063} <73656374696F6E2E31332E31> \MusicPDFRef{16064} <73656374696F6E2E31332E32> \MusicPDFRef{16065} <73656374696F6E2E31332E33> \MusicPDFRef{16066} <73656374696F6E2E31342E31> \MusicPDFRef{16067} <73656374696F6E2E31342E32> \MusicPDFRef{16068}] /Limits [<73656374696F6E2E31322E35> <73656374696F6E2E31342E32>] >>}
\MusicPDFWriteObject{10081}{<< /Names [<73656374696F6E2E31342E33> \MusicPDFRef{16069} <73656374696F6E2E322E31> \MusicPDFRef{16070} <73656374696F6E2E322E32> \MusicPDFRef{16071} <73656374696F6E2E322E33> \MusicPDFRef{16072} <73656374696F6E2E322E34> \MusicPDFRef{16073} <73656374696F6E2E332E31> \MusicPDFRef{16074}] /Limits [<73656374696F6E2E31342E33> <73656374696F6E2E332E31>] >>}
\MusicPDFWriteObject{10082}{<< /Names [<73656374696F6E2E332E32> \MusicPDFRef{16075} <73656374696F6E2E332E33> \MusicPDFRef{16076} <73656374696F6E2E332E34> \MusicPDFRef{16077} <73656374696F6E2E342E31> \MusicPDFRef{16078} <73656374696F6E2E342E32> \MusicPDFRef{16079} <73656374696F6E2E342E33> \MusicPDFRef{16080}] /Limits [<73656374696F6E2E332E32> <73656374696F6E2E342E33>] >>}
\MusicPDFWriteObject{10083}{<< /Names [<73656374696F6E2E342E34> \MusicPDFRef{16081} <73656374696F6E2E342E35> \MusicPDFRef{16082} <73656374696F6E2E342E36> \MusicPDFRef{16083} <73656374696F6E2E352E31> \MusicPDFRef{16084} <73656374696F6E2E352E32> \MusicPDFRef{16085} <73656374696F6E2E352E33> \MusicPDFRef{16086}] /Limits [<73656374696F6E2E342E34> <73656374696F6E2E352E33>] >>}
\MusicPDFWriteObject{10084}{<< /Names [<73656374696F6E2E352E34> \MusicPDFRef{16087} <73656374696F6E2E352E35> \MusicPDFRef{16088} <73656374696F6E2E362E31> \MusicPDFRef{16089} <73656374696F6E2E362E32> \MusicPDFRef{16090} <73656374696F6E2E362E33> \MusicPDFRef{16091} <73656374696F6E2E372E31> \MusicPDFRef{16092}] /Limits [<73656374696F6E2E352E34> <73656374696F6E2E372E31>] >>}
\MusicPDFWriteObject{10085}{<< /Names [<73656374696F6E2E372E32> \MusicPDFRef{16093} <73656374696F6E2E372E33> \MusicPDFRef{16094} <73656374696F6E2E372E34> \MusicPDFRef{16095} <73656374696F6E2E372E35> \MusicPDFRef{16096} <73656374696F6E2E382E31> \MusicPDFRef{16097} <73656374696F6E2E382E32> \MusicPDFRef{16098}] /Limits [<73656374696F6E2E372E32> <73656374696F6E2E382E32>] >>}
\MusicPDFWriteObject{10086}{<< /Names [<73656374696F6E2E382E33> \MusicPDFRef{16099} <73656374696F6E2E382E34> \MusicPDFRef{16100} <73656374696F6E2E382E35> \MusicPDFRef{16101} <73656374696F6E2E392E31> \MusicPDFRef{16102} <73656374696F6E2E392E32> \MusicPDFRef{16103} <73656374696F6E2E392E33> \MusicPDFRef{16104}] /Limits [<73656374696F6E2E382E33> <73656374696F6E2E392E33>] >>}
\MusicPDFWriteObject{10087}{<< /Names [<73656374696F6E2E412E31> \MusicPDFRef{16105} <73656374696F6E2E412E32> \MusicPDFRef{16106} <73656374696F6E2E412E33> \MusicPDFRef{16107} <73656374696F6E2E412E34> \MusicPDFRef{16108} <73656374696F6E2E412E35> \MusicPDFRef{16109} <73656374696F6E2E422E31> \MusicPDFRef{16110}] /Limits [<73656374696F6E2E412E31> <73656374696F6E2E422E31>] >>}
\MusicPDFWriteObject{10088}{<< /Names [<73656374696F6E2E422E32> \MusicPDFRef{16111} <73656374696F6E2E422E33> \MusicPDFRef{16112} <73656374696F6E2E422E34> \MusicPDFRef{16113} <73656374696F6E2E422E35> \MusicPDFRef{16114} <73656374696F6E2E422E36> \MusicPDFRef{16115} <73656374696F6E2E422E37> \MusicPDFRef{16116}] /Limits [<73656374696F6E2E422E32> <73656374696F6E2E422E37>] >>}
\MusicPDFWriteObject{10089}{<< /Names [<73656374696F6E2E432E31> \MusicPDFRef{16117} <73656374696F6E2E432E32> \MusicPDFRef{16118} <73656374696F6E2E432E33> \MusicPDFRef{16119} <73656374696F6E2E432E34> \MusicPDFRef{16120} <73656374696F6E2E432E35> \MusicPDFRef{16121} <73656374696F6E2E442E31> \MusicPDFRef{16122}] /Limits [<73656374696F6E2E432E31> <73656374696F6E2E442E31>] >>}
\MusicPDFWriteObject{10090}{<< /Names [<73656374696F6E2E442E32> \MusicPDFRef{16123} <73656374696F6E2E442E33> \MusicPDFRef{16124} <73656374696F6E2E452E31> \MusicPDFRef{16125} <73656374696F6E2E452E32> \MusicPDFRef{16126} <73656374696F6E2E452E33> \MusicPDFRef{16127} <73656374696F6E2E452E34> \MusicPDFRef{16128}] /Limits [<73656374696F6E2E442E32> <73656374696F6E2E452E34>] >>}
\MusicPDFWriteObject{10091}{<< /Names [<73656374696F6E2E452E35> \MusicPDFRef{16129} <73656374696F6E2E452E36> \MusicPDFRef{16130} <73756273656374696F6E2E312E312E31> \MusicPDFRef{16131} <73756273656374696F6E2E312E312E32> \MusicPDFRef{16132} <73756273656374696F6E2E312E312E33> \MusicPDFRef{16133} <73756273656374696F6E2E312E312E34> \MusicPDFRef{16134}] /Limits [<73656374696F6E2E452E35> <73756273656374696F6E2E312E312E34>] >>}
\MusicPDFWriteObject{10092}{<< /Names [<73756273656374696F6E2E312E322E31> \MusicPDFRef{16135} <73756273656374696F6E2E312E322E3130> \MusicPDFRef{16136} <73756273656374696F6E2E312E322E3131> \MusicPDFRef{16137} <73756273656374696F6E2E312E322E32> \MusicPDFRef{16138} <73756273656374696F6E2E312E322E33> \MusicPDFRef{16139} <73756273656374696F6E2E312E322E34> \MusicPDFRef{16140}] /Limits [<73756273656374696F6E2E312E322E31> <73756273656374696F6E2E312E322E34>] >>}
\MusicPDFWriteObject{10093}{<< /Names [<73756273656374696F6E2E312E322E35> \MusicPDFRef{16141} <73756273656374696F6E2E312E322E36> \MusicPDFRef{16142} <73756273656374696F6E2E312E322E37> \MusicPDFRef{16143} <73756273656374696F6E2E312E322E38> \MusicPDFRef{16144} <73756273656374696F6E2E312E322E39> \MusicPDFRef{16145} <73756273656374696F6E2E312E332E31> \MusicPDFRef{16146}] /Limits [<73756273656374696F6E2E312E322E35> <73756273656374696F6E2E312E332E31>] >>}
\MusicPDFWriteObject{10094}{<< /Names [<73756273656374696F6E2E312E332E32> \MusicPDFRef{16147} <73756273656374696F6E2E312E332E33> \MusicPDFRef{16148} <73756273656374696F6E2E312E332E34> \MusicPDFRef{16149} <73756273656374696F6E2E312E342E31> \MusicPDFRef{16150} <73756273656374696F6E2E312E342E3130> \MusicPDFRef{16151} <73756273656374696F6E2E312E342E3131> \MusicPDFRef{16152}] /Limits [<73756273656374696F6E2E312E332E32> <73756273656374696F6E2E312E342E3131>] >>}
\MusicPDFWriteObject{10095}{<< /Names [<73756273656374696F6E2E312E342E32> \MusicPDFRef{16153} <73756273656374696F6E2E312E342E33> \MusicPDFRef{16154} <73756273656374696F6E2E312E342E34> \MusicPDFRef{16155} <73756273656374696F6E2E312E342E35> \MusicPDFRef{16156} <73756273656374696F6E2E312E342E36> \MusicPDFRef{16157} <73756273656374696F6E2E312E342E37> \MusicPDFRef{16158}] /Limits [<73756273656374696F6E2E312E342E32> <73756273656374696F6E2E312E342E37>] >>}
\MusicPDFWriteObject{10096}{<< /Names [<73756273656374696F6E2E312E342E38> \MusicPDFRef{16159} <73756273656374696F6E2E312E342E39> \MusicPDFRef{16160} <73756273656374696F6E2E312E352E31> \MusicPDFRef{16161} <73756273656374696F6E2E312E352E32> \MusicPDFRef{16162} <73756273656374696F6E2E312E362E31> \MusicPDFRef{16163} <73756273656374696F6E2E312E362E32> \MusicPDFRef{16164}] /Limits [<73756273656374696F6E2E312E342E38> <73756273656374696F6E2E312E362E32>] >>}
\MusicPDFWriteObject{10097}{<< /Names [<73756273656374696F6E2E312E372E31> \MusicPDFRef{16165} <73756273656374696F6E2E312E372E32> \MusicPDFRef{16166} <73756273656374696F6E2E312E372E33> \MusicPDFRef{16167} <73756273656374696F6E2E312E382E31> \MusicPDFRef{16168} <73756273656374696F6E2E312E382E32> \MusicPDFRef{16169} <73756273656374696F6E2E312E382E33> \MusicPDFRef{16170}] /Limits [<73756273656374696F6E2E312E372E31> <73756273656374696F6E2E312E382E33>] >>}
\MusicPDFWriteObject{10098}{<< /Names [<73756273656374696F6E2E312E382E34> \MusicPDFRef{16171} <73756273656374696F6E2E312E382E35> \MusicPDFRef{16172} <73756273656374696F6E2E31302E312E31> \MusicPDFRef{16173} <73756273656374696F6E2E31302E312E32> \MusicPDFRef{16174} <73756273656374696F6E2E31302E322E31> \MusicPDFRef{16175} <73756273656374696F6E2E31302E322E32> \MusicPDFRef{16176}] /Limits [<73756273656374696F6E2E312E382E34> <73756273656374696F6E2E31302E322E32>] >>}
\MusicPDFWriteObject{10099}{<< /Names [<73756273656374696F6E2E31302E332E31> \MusicPDFRef{16177} <73756273656374696F6E2E31302E332E32> \MusicPDFRef{16178} <73756273656374696F6E2E31302E342E31> \MusicPDFRef{16179} <73756273656374696F6E2E31302E342E32> \MusicPDFRef{16180} <73756273656374696F6E2E31302E342E33> \MusicPDFRef{16181} <73756273656374696F6E2E31312E312E31> \MusicPDFRef{16182}] /Limits [<73756273656374696F6E2E31302E332E31> <73756273656374696F6E2E31312E312E31>] >>}
\MusicPDFWriteObject{10100}{<< /Names [<73756273656374696F6E2E31312E312E32> \MusicPDFRef{16183} <73756273656374696F6E2E31312E322E31> \MusicPDFRef{16184} <73756273656374696F6E2E31312E322E32> \MusicPDFRef{16185} <73756273656374696F6E2E31312E332E31> \MusicPDFRef{16186} <73756273656374696F6E2E31312E332E32> \MusicPDFRef{16187} <73756273656374696F6E2E31322E312E31> \MusicPDFRef{16188}] /Limits [<73756273656374696F6E2E31312E312E32> <73756273656374696F6E2E31322E312E31>] >>}
\MusicPDFWriteObject{10101}{<< /Names [<73756273656374696F6E2E31322E322E31> \MusicPDFRef{16189} <73756273656374696F6E2E31322E322E32> \MusicPDFRef{16190} <73756273656374696F6E2E31322E342E31> \MusicPDFRef{16191} <73756273656374696F6E2E31322E352E31> \MusicPDFRef{16192} <73756273656374696F6E2E31332E332E31> \MusicPDFRef{16193} <73756273656374696F6E2E31342E312E31> \MusicPDFRef{16194}] /Limits [<73756273656374696F6E2E31322E322E31> <73756273656374696F6E2E31342E312E31>] >>}
\MusicPDFWriteObject{10102}{<< /Names [<73756273656374696F6E2E31342E312E32> \MusicPDFRef{16195} <73756273656374696F6E2E322E312E31> \MusicPDFRef{16196} <73756273656374696F6E2E322E312E32> \MusicPDFRef{16197} <73756273656374696F6E2E322E312E33> \MusicPDFRef{16198} <73756273656374696F6E2E322E312E34> \MusicPDFRef{16199} <73756273656374696F6E2E322E322E31> \MusicPDFRef{16200}] /Limits [<73756273656374696F6E2E31342E312E32> <73756273656374696F6E2E322E322E31>] >>}
\MusicPDFWriteObject{10103}{<< /Names [<73756273656374696F6E2E322E322E32> \MusicPDFRef{16201} <73756273656374696F6E2E322E322E33> \MusicPDFRef{16202} <73756273656374696F6E2E322E322E34> \MusicPDFRef{16203} <73756273656374696F6E2E322E322E35> \MusicPDFRef{16204} <73756273656374696F6E2E322E332E31> \MusicPDFRef{16205} <73756273656374696F6E2E322E332E32> \MusicPDFRef{16206}] /Limits [<73756273656374696F6E2E322E322E32> <73756273656374696F6E2E322E332E32>] >>}
\MusicPDFWriteObject{10104}{<< /Names [<73756273656374696F6E2E322E342E31> \MusicPDFRef{16207} <73756273656374696F6E2E322E342E32> \MusicPDFRef{16208} <73756273656374696F6E2E322E342E33> \MusicPDFRef{16209} <73756273656374696F6E2E332E312E31> \MusicPDFRef{16210} <73756273656374696F6E2E332E312E32> \MusicPDFRef{16211} <73756273656374696F6E2E332E322E31> \MusicPDFRef{16212}] /Limits [<73756273656374696F6E2E322E342E31> <73756273656374696F6E2E332E322E31>] >>}
\MusicPDFWriteObject{10105}{<< /Names [<73756273656374696F6E2E332E322E32> \MusicPDFRef{16213} <73756273656374696F6E2E332E332E31> \MusicPDFRef{16214} <73756273656374696F6E2E332E332E32> \MusicPDFRef{16215} <73756273656374696F6E2E332E342E31> \MusicPDFRef{16216} <73756273656374696F6E2E332E342E32> \MusicPDFRef{16217} <73756273656374696F6E2E332E342E33> \MusicPDFRef{16218}] /Limits [<73756273656374696F6E2E332E322E32> <73756273656374696F6E2E332E342E33>] >>}
\MusicPDFWriteObject{10106}{<< /Names [<73756273656374696F6E2E342E312E31> \MusicPDFRef{16219} <73756273656374696F6E2E342E312E32> \MusicPDFRef{16220} <73756273656374696F6E2E342E322E31> \MusicPDFRef{16221} <73756273656374696F6E2E342E322E32> \MusicPDFRef{16222} <73756273656374696F6E2E342E332E31> \MusicPDFRef{16223} <73756273656374696F6E2E342E332E32> \MusicPDFRef{16224}] /Limits [<73756273656374696F6E2E342E312E31> <73756273656374696F6E2E342E332E32>] >>}
\MusicPDFWriteObject{10107}{<< /Names [<73756273656374696F6E2E342E332E33> \MusicPDFRef{16225} <73756273656374696F6E2E342E342E31> \MusicPDFRef{16226} <73756273656374696F6E2E342E342E32> \MusicPDFRef{16227} <73756273656374696F6E2E342E352E31> \MusicPDFRef{16228} <73756273656374696F6E2E342E352E32> \MusicPDFRef{16229} <73756273656374696F6E2E342E352E33> \MusicPDFRef{16230}] /Limits [<73756273656374696F6E2E342E332E33> <73756273656374696F6E2E342E352E33>] >>}
\MusicPDFWriteObject{10108}{<< /Names [<73756273656374696F6E2E342E362E31> \MusicPDFRef{16231} <73756273656374696F6E2E342E362E32> \MusicPDFRef{16232} <73756273656374696F6E2E342E362E33> \MusicPDFRef{16233} <73756273656374696F6E2E352E312E31> \MusicPDFRef{16234} <73756273656374696F6E2E352E312E32> \MusicPDFRef{16235} <73756273656374696F6E2E352E312E33> \MusicPDFRef{16236}] /Limits [<73756273656374696F6E2E342E362E31> <73756273656374696F6E2E352E312E33>] >>}
\MusicPDFWriteObject{10109}{<< /Names [<73756273656374696F6E2E352E322E31> \MusicPDFRef{16237} <73756273656374696F6E2E352E322E32> \MusicPDFRef{16238} <73756273656374696F6E2E352E322E33> \MusicPDFRef{16239} <73756273656374696F6E2E352E322E34> \MusicPDFRef{16240} <73756273656374696F6E2E352E332E31> \MusicPDFRef{16241} <73756273656374696F6E2E352E332E32> \MusicPDFRef{16242}] /Limits [<73756273656374696F6E2E352E322E31> <73756273656374696F6E2E352E332E32>] >>}
\MusicPDFWriteObject{10110}{<< /Names [<73756273656374696F6E2E352E332E33> \MusicPDFRef{16243} <73756273656374696F6E2E352E332E34> \MusicPDFRef{16244} <73756273656374696F6E2E352E342E31> \MusicPDFRef{16245} <73756273656374696F6E2E352E342E32> \MusicPDFRef{16246} <73756273656374696F6E2E352E352E31> \MusicPDFRef{16247} <73756273656374696F6E2E352E352E32> \MusicPDFRef{16248}] /Limits [<73756273656374696F6E2E352E332E33> <73756273656374696F6E2E352E352E32>] >>}
\MusicPDFWriteObject{10111}{<< /Names [<73756273656374696F6E2E352E352E33> \MusicPDFRef{16249} <73756273656374696F6E2E362E312E31> \MusicPDFRef{16250} <73756273656374696F6E2E362E312E32> \MusicPDFRef{16251} <73756273656374696F6E2E362E322E31> \MusicPDFRef{16252} <73756273656374696F6E2E362E322E32> \MusicPDFRef{16253} <73756273656374696F6E2E362E332E31> \MusicPDFRef{16254}] /Limits [<73756273656374696F6E2E352E352E33> <73756273656374696F6E2E362E332E31>] >>}
\MusicPDFWriteObject{10112}{<< /Names [<73756273656374696F6E2E362E332E32> \MusicPDFRef{16255} <73756273656374696F6E2E372E312E31> \MusicPDFRef{16256} <73756273656374696F6E2E372E312E32> \MusicPDFRef{16257} <73756273656374696F6E2E372E312E33> \MusicPDFRef{16258} <73756273656374696F6E2E372E312E34> \MusicPDFRef{16259} <73756273656374696F6E2E372E322E31> \MusicPDFRef{16260}] /Limits [<73756273656374696F6E2E362E332E32> <73756273656374696F6E2E372E322E31>] >>}
\MusicPDFWriteObject{10113}{<< /Names [<73756273656374696F6E2E372E322E32> \MusicPDFRef{16261} <73756273656374696F6E2E372E332E31> \MusicPDFRef{16262} <73756273656374696F6E2E372E332E32> \MusicPDFRef{16263} <73756273656374696F6E2E372E332E33> \MusicPDFRef{16264} <73756273656374696F6E2E372E332E34> \MusicPDFRef{16265} <73756273656374696F6E2E372E342E31> \MusicPDFRef{16266}] /Limits [<73756273656374696F6E2E372E322E32> <73756273656374696F6E2E372E342E31>] >>}
\MusicPDFWriteObject{10114}{<< /Names [<73756273656374696F6E2E372E342E32> \MusicPDFRef{16267} <73756273656374696F6E2E372E342E33> \MusicPDFRef{16268} <73756273656374696F6E2E392E312E31> \MusicPDFRef{16269} <73756273656374696F6E2E392E312E32> \MusicPDFRef{16270} <73756273656374696F6E2E392E322E31> \MusicPDFRef{16271} <73756273656374696F6E2E392E322E32> \MusicPDFRef{16272}] /Limits [<73756273656374696F6E2E372E342E32> <73756273656374696F6E2E392E322E32>] >>}
\MusicPDFWriteObject{10115}{<< /Names [<73756273656374696F6E2E392E332E31> \MusicPDFRef{16273} <73756273656374696F6E2E392E332E32> \MusicPDFRef{16274} <73756273656374696F6E2E392E332E33> \MusicPDFRef{16275} <73756273656374696F6E2E412E312E31> \MusicPDFRef{16276} <73756273656374696F6E2E412E312E32> \MusicPDFRef{16277} <73756273656374696F6E2E412E312E33> \MusicPDFRef{16278}] /Limits [<73756273656374696F6E2E392E332E31> <73756273656374696F6E2E412E312E33>] >>}
\MusicPDFWriteObject{10116}{<< /Names [<73756273656374696F6E2E412E312E34> \MusicPDFRef{16279} <73756273656374696F6E2E412E312E35> \MusicPDFRef{16280} <73756273656374696F6E2E412E322E31> \MusicPDFRef{16281} <73756273656374696F6E2E412E322E32> \MusicPDFRef{16282} <73756273656374696F6E2E412E322E33> \MusicPDFRef{16283} <73756273656374696F6E2E412E322E34> \MusicPDFRef{16284}] /Limits [<73756273656374696F6E2E412E312E34> <73756273656374696F6E2E412E322E34>] >>}
\MusicPDFWriteObject{10117}{<< /Names [<73756273656374696F6E2E412E332E31> \MusicPDFRef{16285} <73756273656374696F6E2E412E332E32> \MusicPDFRef{16286} <73756273656374696F6E2E412E332E33> \MusicPDFRef{16287} <73756273656374696F6E2E412E332E34> \MusicPDFRef{16288} <73756273656374696F6E2E412E342E31> \MusicPDFRef{16289} <73756273656374696F6E2E412E352E31> \MusicPDFRef{16290}] /Limits [<73756273656374696F6E2E412E332E31> <73756273656374696F6E2E412E352E31>] >>}
\MusicPDFWriteObject{10118}{<< /Names [<73756273656374696F6E2E412E352E32> \MusicPDFRef{16291} <73756273656374696F6E2E412E352E33> \MusicPDFRef{16292} <73756273656374696F6E2E412E352E34> \MusicPDFRef{16293} <73756273656374696F6E2E422E312E31> \MusicPDFRef{16294} <73756273656374696F6E2E422E322E31> \MusicPDFRef{16295} <73756273656374696F6E2E422E322E32> \MusicPDFRef{16296}] /Limits [<73756273656374696F6E2E412E352E32> <73756273656374696F6E2E422E322E32>] >>}
\MusicPDFWriteObject{10119}{<< /Names [<73756273656374696F6E2E422E332E31> \MusicPDFRef{16297} <73756273656374696F6E2E422E332E32> \MusicPDFRef{16298} <73756273656374696F6E2E422E342E31> \MusicPDFRef{16299} <73756273656374696F6E2E422E342E32> \MusicPDFRef{16300} <73756273656374696F6E2E422E342E33> \MusicPDFRef{16301} <73756273656374696F6E2E422E342E34> \MusicPDFRef{16302}] /Limits [<73756273656374696F6E2E422E332E31> <73756273656374696F6E2E422E342E34>] >>}
\MusicPDFWriteObject{10120}{<< /Names [<73756273656374696F6E2E422E352E31> \MusicPDFRef{16303} <73756273656374696F6E2E422E352E32> \MusicPDFRef{16304} <73756273656374696F6E2E422E352E33> \MusicPDFRef{16305} <73756273656374696F6E2E422E362E31> \MusicPDFRef{16306} <73756273656374696F6E2E422E372E31> \MusicPDFRef{16307} <73756273656374696F6E2E422E372E32> \MusicPDFRef{16308}] /Limits [<73756273656374696F6E2E422E352E31> <73756273656374696F6E2E422E372E32>] >>}
\MusicPDFWriteObject{10121}{<< /Names [<73756273656374696F6E2E432E312E31> \MusicPDFRef{16309} <73756273656374696F6E2E432E312E32> \MusicPDFRef{16310} <73756273656374696F6E2E432E312E33> \MusicPDFRef{16311} <73756273656374696F6E2E432E312E34> \MusicPDFRef{16312} <73756273656374696F6E2E432E312E35> \MusicPDFRef{16313} <73756273656374696F6E2E432E312E36> \MusicPDFRef{16314}] /Limits [<73756273656374696F6E2E432E312E31> <73756273656374696F6E2E432E312E36>] >>}
\MusicPDFWriteObject{10122}{<< /Names [<73756273656374696F6E2E432E322E31> \MusicPDFRef{16315} <73756273656374696F6E2E432E322E32> \MusicPDFRef{16316} <73756273656374696F6E2E432E322E33> \MusicPDFRef{16317} <73756273656374696F6E2E432E322E34> \MusicPDFRef{16318} <73756273656374696F6E2E432E322E35> \MusicPDFRef{16319} <73756273656374696F6E2E432E332E31> \MusicPDFRef{16320}] /Limits [<73756273656374696F6E2E432E322E31> <73756273656374696F6E2E432E332E31>] >>}
\MusicPDFWriteObject{10123}{<< /Names [<73756273656374696F6E2E432E332E32> \MusicPDFRef{16321} <73756273656374696F6E2E432E332E33> \MusicPDFRef{16322} <73756273656374696F6E2E432E332E34> \MusicPDFRef{16323} <73756273656374696F6E2E432E332E35> \MusicPDFRef{16324} <73756273656374696F6E2E432E332E36> \MusicPDFRef{16325} <73756273656374696F6E2E432E332E37> \MusicPDFRef{16326}] /Limits [<73756273656374696F6E2E432E332E32> <73756273656374696F6E2E432E332E37>] >>}
\MusicPDFWriteObject{10124}{<< /Names [<73756273656374696F6E2E432E342E31> \MusicPDFRef{16327} <73756273656374696F6E2E432E342E32> \MusicPDFRef{16328} <73756273656374696F6E2E432E342E33> \MusicPDFRef{16329} <73756273656374696F6E2E432E342E34> \MusicPDFRef{16330} <73756273656374696F6E2E432E352E31> \MusicPDFRef{16331} <73756273656374696F6E2E432E352E32> \MusicPDFRef{16332}] /Limits [<73756273656374696F6E2E432E342E31> <73756273656374696F6E2E432E352E32>] >>}
\MusicPDFWriteObject{10125}{<< /Names [<73756273656374696F6E2E432E352E33> \MusicPDFRef{16333} <73756273656374696F6E2E432E352E34> \MusicPDFRef{16334} <73756273656374696F6E2E442E312E31> \MusicPDFRef{16335} <73756273656374696F6E2E442E312E3130> \MusicPDFRef{16336} <73756273656374696F6E2E442E312E32> \MusicPDFRef{16337} <73756273656374696F6E2E442E312E33> \MusicPDFRef{16338}] /Limits [<73756273656374696F6E2E432E352E33> <73756273656374696F6E2E442E312E33>] >>}
\MusicPDFWriteObject{10126}{<< /Names [<73756273656374696F6E2E442E312E34> \MusicPDFRef{16339} <73756273656374696F6E2E442E312E35> \MusicPDFRef{16340} <73756273656374696F6E2E442E312E36> \MusicPDFRef{16341} <73756273656374696F6E2E442E312E37> \MusicPDFRef{16342} <73756273656374696F6E2E442E312E38> \MusicPDFRef{16343} <73756273656374696F6E2E442E312E39> \MusicPDFRef{16344}] /Limits [<73756273656374696F6E2E442E312E34> <73756273656374696F6E2E442E312E39>] >>}
\MusicPDFWriteObject{10127}{<< /Names [<73756273656374696F6E2E442E322E31> \MusicPDFRef{16345} <73756273656374696F6E2E442E322E32> \MusicPDFRef{16346} <73756273656374696F6E2E442E332E31> \MusicPDFRef{16347} <73756273656374696F6E2E442E332E32> \MusicPDFRef{16348} <73756273656374696F6E2E442E332E33> \MusicPDFRef{16349} <73756273656374696F6E2E442E332E34> \MusicPDFRef{16350}] /Limits [<73756273656374696F6E2E442E322E31> <73756273656374696F6E2E442E332E34>] >>}
\MusicPDFWriteObject{10128}{<< /Names [<73756273656374696F6E2E442E332E35> \MusicPDFRef{16351} <73756273656374696F6E2E442E332E36> \MusicPDFRef{16352} <73756273656374696F6E2E452E312E31> \MusicPDFRef{16353} <73756273656374696F6E2E452E312E32> \MusicPDFRef{16354} <73756273656374696F6E2E452E322E31> \MusicPDFRef{16355} <73756273656374696F6E2E452E322E32> \MusicPDFRef{16356}] /Limits [<73756273656374696F6E2E442E332E35> <73756273656374696F6E2E452E322E32>] >>}
\MusicPDFWriteObject{10129}{<< /Names [<73756273656374696F6E2E452E322E33> \MusicPDFRef{16357} <73756273656374696F6E2E452E322E34> \MusicPDFRef{16358} <73756273656374696F6E2E452E332E31> \MusicPDFRef{16359} <73756273656374696F6E2E452E332E32> \MusicPDFRef{16360} <73756273656374696F6E2E452E332E33> \MusicPDFRef{16361} <73756273656374696F6E2E452E332E34> \MusicPDFRef{16362}] /Limits [<73756273656374696F6E2E452E322E33> <73756273656374696F6E2E452E332E34>] >>}
\MusicPDFWriteObject{10130}{<< /Names [<73756273656374696F6E2E452E332E35> \MusicPDFRef{16363} <73756273656374696F6E2E452E342E31> \MusicPDFRef{16364} <73756273656374696F6E2E452E342E32> \MusicPDFRef{16365} <73756273656374696F6E2E452E352E31> \MusicPDFRef{16366} <73756273656374696F6E2E452E352E32> \MusicPDFRef{16367} <73756273656374696F6E2E452E352E33> \MusicPDFRef{16368}] /Limits [<73756273656374696F6E2E452E332E35> <73756273656374696F6E2E452E352E33>] >>}
\MusicPDFWriteObject{10131}{<< /Names [<73756273656374696F6E2E452E352E34> \MusicPDFRef{16369} <73756273656374696F6E2E452E352E35> \MusicPDFRef{16370} <73756273656374696F6E2E452E352E36> \MusicPDFRef{16371} <73756273656374696F6E2E452E352E37> \MusicPDFRef{16372} <73756273656374696F6E2E452E362E31> \MusicPDFRef{16373} <73756273656374696F6E2E452E362E32> \MusicPDFRef{16374}] /Limits [<73756273656374696F6E2E452E352E34> <73756273656374696F6E2E452E362E32>] >>}
\MusicPDFWriteObject{10132}{<< /Names [<73756273656374696F6E2E452E362E33> \MusicPDFRef{16375} <73756273656374696F6E2E452E362E34> \MusicPDFRef{16376} <73756273756273656374696F6E2E312E382E322E31> \MusicPDFRef{16377} <73756273756273656374696F6E2E312E382E322E32> \MusicPDFRef{16378} <73756273756273656374696F6E2E312E382E322E33> \MusicPDFRef{16379} <73756273756273656374696F6E2E312E382E332E31> \MusicPDFRef{16380}] /Limits [<73756273656374696F6E2E452E362E33> <73756273756273656374696F6E2E312E382E332E31>] >>}
\MusicPDFWriteObject{10133}{<< /Names [<73756273756273656374696F6E2E312E382E332E32> \MusicPDFRef{16381} <73756273756273656374696F6E2E312E382E332E33> \MusicPDFRef{16382} <73756273756273656374696F6E2E312E382E342E31> \MusicPDFRef{16383} <73756273756273656374696F6E2E312E382E342E32> \MusicPDFRef{16384} <73756273756273656374696F6E2E312E382E342E33> \MusicPDFRef{16385} <73756273756273656374696F6E2E312E382E352E31> \MusicPDFRef{16386}] /Limits [<73756273756273656374696F6E2E312E382E332E32> <73756273756273656374696F6E2E312E382E352E31>] >>}
\MusicPDFWriteObject{10134}{<< /Names [<73756273756273656374696F6E2E31302E322E312E31> \MusicPDFRef{16387} <73756273756273656374696F6E2E31302E322E312E32> \MusicPDFRef{16388} <73756273756273656374696F6E2E31302E322E322E31> \MusicPDFRef{16389} <73756273756273656374696F6E2E31302E322E322E32> \MusicPDFRef{16390} <73756273756273656374696F6E2E31302E342E322E31> \MusicPDFRef{16391} <73756273756273656374696F6E2E31302E342E322E32> \MusicPDFRef{16392}] /Limits [<73756273756273656374696F6E2E31302E322E312E31> <73756273756273656374696F6E2E31302E342E322E32>] >>}
\MusicPDFWriteObject{10135}{<< /Names [<73756273756273656374696F6E2E31302E342E332E31> \MusicPDFRef{16393} <73756273756273656374696F6E2E31302E342E332E32> \MusicPDFRef{16394} <73756273756273656374696F6E2E31312E332E312E31> \MusicPDFRef{16395} <73756273756273656374696F6E2E322E312E312E31> \MusicPDFRef{16396} <73756273756273656374696F6E2E322E312E312E32> \MusicPDFRef{16397} <73756273756273656374696F6E2E322E312E312E33> \MusicPDFRef{16398}] /Limits [<73756273756273656374696F6E2E31302E342E332E31> <73756273756273656374696F6E2E322E312E312E33>] >>}
\MusicPDFWriteObject{10136}{<< /Names [<73756273756273656374696F6E2E322E312E322E31> \MusicPDFRef{16399} <73756273756273656374696F6E2E322E312E322E32> \MusicPDFRef{16400} <73756273756273656374696F6E2E322E312E322E33> \MusicPDFRef{16401} <73756273756273656374696F6E2E322E312E332E31> \MusicPDFRef{16402} <73756273756273656374696F6E2E322E322E322E31> \MusicPDFRef{16403} <73756273756273656374696F6E2E322E322E332E31> \MusicPDFRef{16404}] /Limits [<73756273756273656374696F6E2E322E312E322E31> <73756273756273656374696F6E2E322E322E332E31>] >>}
\MusicPDFWriteObject{10137}{<< /Names [<73756273756273656374696F6E2E322E322E342E31> \MusicPDFRef{16405} <73756273756273656374696F6E2E322E322E352E31> \MusicPDFRef{16406} <73756273756273656374696F6E2E322E322E352E32> \MusicPDFRef{16407} <73756273756273656374696F6E2E322E332E312E31> \MusicPDFRef{16408} <73756273756273656374696F6E2E322E332E312E32> \MusicPDFRef{16409} <73756273756273656374696F6E2E322E332E312E33> \MusicPDFRef{16410}] /Limits [<73756273756273656374696F6E2E322E322E342E31> <73756273756273656374696F6E2E322E332E312E33>] >>}
\MusicPDFWriteObject{10138}{<< /Names [<73756273756273656374696F6E2E332E312E322E31> \MusicPDFRef{16411} <73756273756273656374696F6E2E332E312E322E32> \MusicPDFRef{16412} <73756273756273656374696F6E2E332E312E322E33> \MusicPDFRef{16413} <73756273756273656374696F6E2E332E332E312E31> \MusicPDFRef{16414} <73756273756273656374696F6E2E332E332E312E32> \MusicPDFRef{16415} <73756273756273656374696F6E2E332E332E322E31> \MusicPDFRef{16416}] /Limits [<73756273756273656374696F6E2E332E312E322E31> <73756273756273656374696F6E2E332E332E322E31>] >>}
\MusicPDFWriteObject{10139}{<< /Names [<73756273756273656374696F6E2E332E332E322E32> \MusicPDFRef{16417} <73756273756273656374696F6E2E332E342E322E31> \MusicPDFRef{16418} <73756273756273656374696F6E2E332E342E322E32> \MusicPDFRef{16419} <73756273756273656374696F6E2E332E342E332E31> \MusicPDFRef{16420} <73756273756273656374696F6E2E332E342E332E32> \MusicPDFRef{16421} <73756273756273656374696F6E2E342E312E322E31> \MusicPDFRef{16422}] /Limits [<73756273756273656374696F6E2E332E332E322E32> <73756273756273656374696F6E2E342E312E322E31>] >>}
\MusicPDFWriteObject{10140}{<< /Names [<73756273756273656374696F6E2E342E322E312E31> \MusicPDFRef{16423} <73756273756273656374696F6E2E342E322E312E32> \MusicPDFRef{16424} <73756273756273656374696F6E2E342E332E312E31> \MusicPDFRef{16425} <73756273756273656374696F6E2E342E332E332E31> \MusicPDFRef{16426} <73756273756273656374696F6E2E342E332E332E32> \MusicPDFRef{16427} <73756273756273656374696F6E2E342E342E322E31> \MusicPDFRef{16428}] /Limits [<73756273756273656374696F6E2E342E322E312E31> <73756273756273656374696F6E2E342E342E322E31>] >>}
\MusicPDFWriteObject{10141}{<< /Names [<73756273756273656374696F6E2E342E352E312E31> \MusicPDFRef{16429} <73756273756273656374696F6E2E342E352E312E32> \MusicPDFRef{16430} <73756273756273656374696F6E2E342E362E312E31> \MusicPDFRef{16431} <73756273756273656374696F6E2E392E322E312E31> \MusicPDFRef{16432} <73756273756273656374696F6E2E392E322E322E31> \MusicPDFRef{16433} <73756273756273656374696F6E2E392E322E322E32> \MusicPDFRef{16434}] /Limits [<73756273756273656374696F6E2E342E352E312E31> <73756273756273656374696F6E2E392E322E322E32>] >>}
\MusicPDFWriteObject{10142}{<< /Names [<73756273756273656374696F6E2E392E332E322E31> \MusicPDFRef{16435} <73756273756273656374696F6E2E392E332E322E32> \MusicPDFRef{16436} <73756273756273656374696F6E2E392E332E332E31> \MusicPDFRef{16437} <73756273756273656374696F6E2E392E332E332E32> \MusicPDFRef{16438} <73756273756273656374696F6E2E412E312E312E31> \MusicPDFRef{16439} <73756273756273656374696F6E2E412E312E322E31> \MusicPDFRef{16440}] /Limits [<73756273756273656374696F6E2E392E332E322E31> <73756273756273656374696F6E2E412E312E322E31>] >>}
\MusicPDFWriteObject{10143}{<< /Names [<73756273756273656374696F6E2E412E312E332E31> \MusicPDFRef{16441} <73756273756273656374696F6E2E412E312E332E32> \MusicPDFRef{16442} <73756273756273656374696F6E2E412E312E342E31> \MusicPDFRef{16443} <73756273756273656374696F6E2E412E312E352E31> \MusicPDFRef{16444} <73756273756273656374696F6E2E412E312E352E32> \MusicPDFRef{16445} <73756273756273656374696F6E2E412E312E352E33> \MusicPDFRef{16446}] /Limits [<73756273756273656374696F6E2E412E312E332E31> <73756273756273656374696F6E2E412E312E352E33>] >>}
\MusicPDFWriteObject{10144}{<< /Names [<73756273756273656374696F6E2E412E322E312E31> \MusicPDFRef{16447} <73756273756273656374696F6E2E412E322E322E31> \MusicPDFRef{16448} <73756273756273656374696F6E2E412E322E322E32> \MusicPDFRef{16449} <73756273756273656374696F6E2E412E322E322E33> \MusicPDFRef{16450} <73756273756273656374696F6E2E412E322E332E31> \MusicPDFRef{16451} <73756273756273656374696F6E2E412E322E332E32> \MusicPDFRef{16452}] /Limits [<73756273756273656374696F6E2E412E322E312E31> <73756273756273656374696F6E2E412E322E332E32>] >>}
\MusicPDFWriteObject{10145}{<< /Names [<73756273756273656374696F6E2E412E322E332E33> \MusicPDFRef{16453} <73756273756273656374696F6E2E412E322E332E34> \MusicPDFRef{16454} <73756273756273656374696F6E2E412E322E332E35> \MusicPDFRef{16455} <73756273756273656374696F6E2E412E322E332E36> \MusicPDFRef{16456} <73756273756273656374696F6E2E412E322E332E37> \MusicPDFRef{16457} <73756273756273656374696F6E2E412E322E342E31> \MusicPDFRef{16458}] /Limits [<73756273756273656374696F6E2E412E322E332E33> <73756273756273656374696F6E2E412E322E342E31>] >>}
\MusicPDFWriteObject{10146}{<< /Names [<73756273756273656374696F6E2E412E322E342E32> \MusicPDFRef{16459} <73756273756273656374696F6E2E412E332E312E31> \MusicPDFRef{16460} <73756273756273656374696F6E2E412E332E332E31> \MusicPDFRef{16461} <73756273756273656374696F6E2E412E352E312E31> \MusicPDFRef{16462} <73756273756273656374696F6E2E412E352E312E32> \MusicPDFRef{16463} <73756273756273656374696F6E2E412E352E322E31> \MusicPDFRef{16464}] /Limits [<73756273756273656374696F6E2E412E322E342E32> <73756273756273656374696F6E2E412E352E322E31>] >>}
\MusicPDFWriteObject{10147}{<< /Names [<73756273756273656374696F6E2E412E352E322E32> \MusicPDFRef{16465} <73756273756273656374696F6E2E412E352E332E31> \MusicPDFRef{16466} <73756273756273656374696F6E2E422E312E312E31> \MusicPDFRef{16467} <73756273756273656374696F6E2E422E322E312E31> \MusicPDFRef{16468} <73756273756273656374696F6E2E422E322E312E32> \MusicPDFRef{16469} <73756273756273656374696F6E2E422E322E322E31> \MusicPDFRef{16470}] /Limits [<73756273756273656374696F6E2E412E352E322E32> <73756273756273656374696F6E2E422E322E322E31>] >>}
\MusicPDFWriteObject{10148}{<< /Names [<73756273756273656374696F6E2E422E332E312E31> \MusicPDFRef{16471} <73756273756273656374696F6E2E422E332E322E31> \MusicPDFRef{16472} <73756273756273656374696F6E2E422E342E312E31> \MusicPDFRef{16473} <73756273756273656374696F6E2E422E342E312E32> \MusicPDFRef{16474} <73756273756273656374696F6E2E422E342E322E31> \MusicPDFRef{16475} <73756273756273656374696F6E2E422E342E322E32> \MusicPDFRef{16476}] /Limits [<73756273756273656374696F6E2E422E332E312E31> <73756273756273656374696F6E2E422E342E322E32>] >>}
\MusicPDFWriteObject{10149}{<< /Names [<73756273756273656374696F6E2E422E342E322E33> \MusicPDFRef{16477} <73756273756273656374696F6E2E422E342E322E34> \MusicPDFRef{16478} <73756273756273656374696F6E2E422E342E332E31> \MusicPDFRef{16479} <73756273756273656374696F6E2E422E342E342E31> \MusicPDFRef{16480} <73756273756273656374696F6E2E422E342E342E32> \MusicPDFRef{16481} <73756273756273656374696F6E2E422E342E342E33> \MusicPDFRef{16482}] /Limits [<73756273756273656374696F6E2E422E342E322E33> <73756273756273656374696F6E2E422E342E342E33>] >>}
\MusicPDFWriteObject{10150}{<< /Names [<73756273756273656374696F6E2E422E342E342E34> \MusicPDFRef{16483} <73756273756273656374696F6E2E422E352E322E31> \MusicPDFRef{16484} <73756273756273656374696F6E2E422E352E332E31> \MusicPDFRef{16485} <73756273756273656374696F6E2E422E352E332E32> \MusicPDFRef{16486} <73756273756273656374696F6E2E422E352E332E33> \MusicPDFRef{16487} <73756273756273656374696F6E2E422E352E332E34> \MusicPDFRef{16488}] /Limits [<73756273756273656374696F6E2E422E342E342E34> <73756273756273656374696F6E2E422E352E332E34>] >>}
\MusicPDFWriteObject{10151}{<< /Names [<73756273756273656374696F6E2E422E362E312E31> \MusicPDFRef{16489} <73756273756273656374696F6E2E422E362E312E32> \MusicPDFRef{16490} <73756273756273656374696F6E2E422E362E312E33> \MusicPDFRef{16491} <73756273756273656374696F6E2E422E372E312E31> \MusicPDFRef{16492} <73756273756273656374696F6E2E422E372E312E32> \MusicPDFRef{16493} <73756273756273656374696F6E2E422E372E312E33> \MusicPDFRef{16494}] /Limits [<73756273756273656374696F6E2E422E362E312E31> <73756273756273656374696F6E2E422E372E312E33>] >>}
\MusicPDFWriteObject{10152}{<< /Names [<73756273756273656374696F6E2E422E372E322E31> \MusicPDFRef{16495} <73756273756273656374696F6E2E422E372E322E32> \MusicPDFRef{16496} <73756273756273656374696F6E2E422E372E322E33> \MusicPDFRef{16497} <73756273756273656374696F6E2E422E372E322E34> \MusicPDFRef{16498} <73756273756273656374696F6E2E432E312E312E31> \MusicPDFRef{16499} <73756273756273656374696F6E2E432E312E322E31> \MusicPDFRef{16500}] /Limits [<73756273756273656374696F6E2E422E372E322E31> <73756273756273656374696F6E2E432E312E322E31>] >>}
\MusicPDFWriteObject{10153}{<< /Names [<73756273756273656374696F6E2E432E312E352E31> \MusicPDFRef{16501} <73756273756273656374696F6E2E432E322E322E31> \MusicPDFRef{16502} <73756273756273656374696F6E2E432E322E322E32> \MusicPDFRef{16503} <73756273756273656374696F6E2E432E322E322E33> \MusicPDFRef{16504} <73756273756273656374696F6E2E432E322E332E31> \MusicPDFRef{16505} <73756273756273656374696F6E2E432E322E332E32> \MusicPDFRef{16506}] /Limits [<73756273756273656374696F6E2E432E312E352E31> <73756273756273656374696F6E2E432E322E332E32>] >>}
\MusicPDFWriteObject{10154}{<< /Names [<73756273756273656374696F6E2E432E322E332E33> \MusicPDFRef{16507} <73756273756273656374696F6E2E432E322E332E34> \MusicPDFRef{16508} <73756273756273656374696F6E2E432E322E332E35> \MusicPDFRef{16509} <73756273756273656374696F6E2E432E322E342E31> \MusicPDFRef{16510} <73756273756273656374696F6E2E432E322E342E32> \MusicPDFRef{16511} <73756273756273656374696F6E2E432E322E342E33> \MusicPDFRef{16512}] /Limits [<73756273756273656374696F6E2E432E322E332E33> <73756273756273656374696F6E2E432E322E342E33>] >>}
\MusicPDFWriteObject{10155}{<< /Names [<73756273756273656374696F6E2E432E322E342E34> \MusicPDFRef{16513} <73756273756273656374696F6E2E432E322E342E35> \MusicPDFRef{16514} <73756273756273656374696F6E2E432E322E342E36> \MusicPDFRef{16515} <73756273756273656374696F6E2E432E332E342E31> \MusicPDFRef{16516} <73756273756273656374696F6E2E432E332E352E31> \MusicPDFRef{16517} <73756273756273656374696F6E2E432E332E362E31> \MusicPDFRef{16518}] /Limits [<73756273756273656374696F6E2E432E322E342E34> <73756273756273656374696F6E2E432E332E362E31>] >>}
\MusicPDFWriteObject{10156}{<< /Names [<73756273756273656374696F6E2E432E332E362E32> \MusicPDFRef{16519} <73756273756273656374696F6E2E432E342E332E31> \MusicPDFRef{16520} <73756273756273656374696F6E2E432E342E332E32> \MusicPDFRef{16521} <73756273756273656374696F6E2E432E342E332E33> \MusicPDFRef{16522} <73756273756273656374696F6E2E432E352E312E31> \MusicPDFRef{16523} <73756273756273656374696F6E2E432E352E312E32> \MusicPDFRef{16524}] /Limits [<73756273756273656374696F6E2E432E332E362E32> <73756273756273656374696F6E2E432E352E312E32>] >>}
\MusicPDFWriteObject{10157}{<< /Names [<73756273756273656374696F6E2E432E352E312E33> \MusicPDFRef{16525} <73756273756273656374696F6E2E432E352E312E34> \MusicPDFRef{16526} <73756273756273656374696F6E2E432E352E312E35> \MusicPDFRef{16527} <73756273756273656374696F6E2E432E352E322E31> \MusicPDFRef{16528} <73756273756273656374696F6E2E432E352E332E31> \MusicPDFRef{16529} <73756273756273656374696F6E2E432E352E342E31> \MusicPDFRef{16530}] /Limits [<73756273756273656374696F6E2E432E352E312E33> <73756273756273656374696F6E2E432E352E342E31>] >>}
\MusicPDFWriteObject{10158}{<< /Names [<73756273756273656374696F6E2E442E312E31302E31> \MusicPDFRef{16531} <73756273756273656374696F6E2E442E312E31302E32> \MusicPDFRef{16532} <73756273756273656374696F6E2E442E312E31302E33> \MusicPDFRef{16533} <73756273756273656374696F6E2E442E312E332E31> \MusicPDFRef{16534} <73756273756273656374696F6E2E442E312E332E32> \MusicPDFRef{16535} <73756273756273656374696F6E2E442E312E332E33> \MusicPDFRef{16536}] /Limits [<73756273756273656374696F6E2E442E312E31302E31> <73756273756273656374696F6E2E442E312E332E33>] >>}
\MusicPDFWriteObject{10159}{<< /Names [<73756273756273656374696F6E2E442E312E342E31> \MusicPDFRef{16537} <73756273756273656374696F6E2E442E312E342E32> \MusicPDFRef{16538} <73756273756273656374696F6E2E442E312E342E33> \MusicPDFRef{16539} <73756273756273656374696F6E2E442E312E352E31> \MusicPDFRef{16540} <73756273756273656374696F6E2E442E312E352E32> \MusicPDFRef{16541} <73756273756273656374696F6E2E442E312E352E33> \MusicPDFRef{16542}] /Limits [<73756273756273656374696F6E2E442E312E342E31> <73756273756273656374696F6E2E442E312E352E33>] >>}
\MusicPDFWriteObject{10160}{<< /Names [<73756273756273656374696F6E2E442E312E362E31> \MusicPDFRef{16543} <73756273756273656374696F6E2E442E312E362E32> \MusicPDFRef{16544} <73756273756273656374696F6E2E442E312E362E33> \MusicPDFRef{16545} <73756273756273656374696F6E2E442E312E382E31> \MusicPDFRef{16546} <73756273756273656374696F6E2E442E312E382E32> \MusicPDFRef{16547} <73756273756273656374696F6E2E442E322E312E31> \MusicPDFRef{16548}] /Limits [<73756273756273656374696F6E2E442E312E362E31> <73756273756273656374696F6E2E442E322E312E31>] >>}
\MusicPDFWriteObject{10161}{<< /Names [<73756273756273656374696F6E2E442E322E322E31> \MusicPDFRef{16549} <73756273756273656374696F6E2E442E332E322E31> \MusicPDFRef{16550} <73756273756273656374696F6E2E442E332E322E32> \MusicPDFRef{16551} <73756273756273656374696F6E2E442E332E322E33> \MusicPDFRef{16552} <73756273756273656374696F6E2E442E332E332E31> \MusicPDFRef{16553} <73756273756273656374696F6E2E442E332E332E32> \MusicPDFRef{16554}] /Limits [<73756273756273656374696F6E2E442E322E322E31> <73756273756273656374696F6E2E442E332E332E32>] >>}
\MusicPDFWriteObject{10162}{<< /Names [<73756273756273656374696F6E2E452E312E322E31> \MusicPDFRef{16555} <73756273756273656374696F6E2E452E342E322E31> \MusicPDFRef{16556} <73756273756273656374696F6E2E452E342E322E32> \MusicPDFRef{16557} <73756273756273656374696F6E2E452E352E312E31> \MusicPDFRef{16558} <73756273756273656374696F6E2E452E352E312E32> \MusicPDFRef{16559} <73756273756273656374696F6E2E452E352E312E33> \MusicPDFRef{16560}] /Limits [<73756273756273656374696F6E2E452E312E322E31> <73756273756273656374696F6E2E452E352E312E33>] >>}
\MusicPDFWriteObject{10163}{<< /Names [<73756273756273656374696F6E2E452E352E322E31> \MusicPDFRef{16561} <73756273756273656374696F6E2E452E352E322E32> \MusicPDFRef{16562} <73756273756273656374696F6E2E452E352E322E33> \MusicPDFRef{16563} <73756273756273656374696F6E2E452E352E332E31> \MusicPDFRef{16564} <73756273756273656374696F6E2E452E352E332E32> \MusicPDFRef{16565} <73756273756273656374696F6E2E452E352E352E31> \MusicPDFRef{16566}] /Limits [<73756273756273656374696F6E2E452E352E322E31> <73756273756273656374696F6E2E452E352E352E31>] >>}
\MusicPDFWriteObject{10164}{<< /Names [<73756273756273656374696F6E2E452E352E372E31> \MusicPDFRef{16567} <73756273756273656374696F6E2E452E352E372E32> \MusicPDFRef{16568} <7465726D696E6F6C6F67792E31302E322E312E31> \MusicPDFRef{16569} <7465726D696E6F6C6F67792E31302E322E312E32> \MusicPDFRef{16570} <7465726D696E6F6C6F67792E31302E322E322E31> \MusicPDFRef{16571} <7465726D696E6F6C6F67792E31302E342E322E31> \MusicPDFRef{16572}] /Limits [<73756273756273656374696F6E2E452E352E372E31> <7465726D696E6F6C6F67792E31302E342E322E31>] >>}
\MusicPDFWriteObject{10165}{<< /Names [<7465726D696E6F6C6F67792E31312E322E302E31> \MusicPDFRef{16573} <7465726D696E6F6C6F67792E342E332E322E31> \MusicPDFRef{16574} <7465726D696E6F6C6F67792E342E332E332E31> \MusicPDFRef{16575} <7465726D696E6F6C6F67792E342E332E332E32> \MusicPDFRef{16576} <7465726D696E6F6C6F67792E362E332E312E31> \MusicPDFRef{16577} <7465726D696E6F6C6F67792E412E352E342E31> \MusicPDFRef{16578}] /Limits [<7465726D696E6F6C6F67792E31312E322E302E31> <7465726D696E6F6C6F67792E412E352E342E31>] >>}
\MusicPDFWriteObject{10166}{<< /Names [<7465726D696E6F6C6F67792E442E332E332E31> \MusicPDFRef{16579} <7465726D696E6F6C6F67792E452E302E302E31> \MusicPDFRef{16580} <7465726D696E6F6C6F67792E452E302E302E32> \MusicPDFRef{16581} <7465726D696E6F6C6F67792E452E352E372E31> \MusicPDFRef{16582} <7465726D696E6F6C6F67792E452E362E332E31> \MusicPDFRef{16583} <74686D2E312E372E312E31> \MusicPDFRef{16584}] /Limits [<7465726D696E6F6C6F67792E442E332E332E31> <74686D2E312E372E312E31>] >>}
\MusicPDFWriteObject{10167}{<< /Names [<74686D2E312E382E312E31> \MusicPDFRef{16585} <74686D2E31302E312E312E31> \MusicPDFRef{16586} <74686D2E31302E312E312E32> \MusicPDFRef{16587} <74686D2E31302E312E312E33> \MusicPDFRef{16588} <74686D2E31302E312E312E34> \MusicPDFRef{16589} <74686D2E31302E312E312E35> \MusicPDFRef{16590}] /Limits [<74686D2E312E382E312E31> <74686D2E31302E312E312E35>] >>}
\MusicPDFWriteObject{10168}{<< /Names [<74686D2E31302E312E322E31> \MusicPDFRef{16591} <74686D2E31302E322E312E31> \MusicPDFRef{16592} <74686D2E31302E322E312E32> \MusicPDFRef{16593} <74686D2E31302E322E312E33> \MusicPDFRef{16594} <74686D2E31302E322E312E34> \MusicPDFRef{16595} <74686D2E31302E322E312E35> \MusicPDFRef{16596}] /Limits [<74686D2E31302E312E322E31> <74686D2E31302E322E312E35>] >>}
\MusicPDFWriteObject{10169}{<< /Names [<74686D2E31302E322E322E31> \MusicPDFRef{16597} <74686D2E31302E322E322E32> \MusicPDFRef{16598} <74686D2E31302E322E322E33> \MusicPDFRef{16599} <74686D2E31302E332E312E31> \MusicPDFRef{16600} <74686D2E31302E332E322E31> \MusicPDFRef{16601} <74686D2E31302E342E312E31> \MusicPDFRef{16602}] /Limits [<74686D2E31302E322E322E31> <74686D2E31302E342E312E31>] >>}
\MusicPDFWriteObject{10170}{<< /Names [<74686D2E31302E342E312E32> \MusicPDFRef{16603} <74686D2E31302E342E312E33> \MusicPDFRef{16604} <74686D2E31302E342E312E34> \MusicPDFRef{16605} <74686D2E31302E342E312E35> \MusicPDFRef{16606} <74686D2E31302E342E322E31> \MusicPDFRef{16607} <74686D2E31302E342E322E32> \MusicPDFRef{16608}] /Limits [<74686D2E31302E342E312E32> <74686D2E31302E342E322E32>] >>}
\MusicPDFWriteObject{10171}{<< /Names [<74686D2E31302E342E322E33> \MusicPDFRef{16609} <74686D2E31302E342E322E34> \MusicPDFRef{16610} <74686D2E31302E342E322E35> \MusicPDFRef{16611} <74686D2E31302E342E322E36> \MusicPDFRef{16612} <74686D2E31302E342E322E37> \MusicPDFRef{16613} <74686D2E31302E342E322E38> \MusicPDFRef{16614}] /Limits [<74686D2E31302E342E322E33> <74686D2E31302E342E322E38>] >>}
\MusicPDFWriteObject{10172}{<< /Names [<74686D2E31302E342E332E31> \MusicPDFRef{16615} <74686D2E31302E342E332E32> \MusicPDFRef{16616} <74686D2E31302E342E332E33> \MusicPDFRef{16617} <74686D2E31302E342E332E34> \MusicPDFRef{16618} <74686D2E31302E342E332E35> \MusicPDFRef{16619} <74686D2E31302E342E332E36> \MusicPDFRef{16620}] /Limits [<74686D2E31302E342E332E31> <74686D2E31302E342E332E36>] >>}
\MusicPDFWriteObject{10173}{<< /Names [<74686D2E31302E342E332E37> \MusicPDFRef{16621} <74686D2E31312E312E312E31> \MusicPDFRef{16622} <74686D2E31312E312E312E32> \MusicPDFRef{16623} <74686D2E31312E312E312E33> \MusicPDFRef{16624} <74686D2E31312E312E312E34> \MusicPDFRef{16625} <74686D2E31312E312E312E35> \MusicPDFRef{16626}] /Limits [<74686D2E31302E342E332E37> <74686D2E31312E312E312E35>] >>}
\MusicPDFWriteObject{10174}{<< /Names [<74686D2E31312E312E322E31> \MusicPDFRef{16627} <74686D2E31312E312E322E32> \MusicPDFRef{16628} <74686D2E31312E322E302E31> \MusicPDFRef{16629} <74686D2E31312E322E302E32> \MusicPDFRef{16630} <74686D2E31312E322E302E33> \MusicPDFRef{16631} <74686D2E31312E322E302E34> \MusicPDFRef{16632}] /Limits [<74686D2E31312E312E322E31> <74686D2E31312E322E302E34>] >>}
\MusicPDFWriteObject{10175}{<< /Names [<74686D2E31312E322E302E35> \MusicPDFRef{16633} <74686D2E31312E322E312E31> \MusicPDFRef{16634} <74686D2E31312E322E312E32> \MusicPDFRef{16635} <74686D2E31312E322E312E33> \MusicPDFRef{16636} <74686D2E31312E322E312E34> \MusicPDFRef{16637} <74686D2E31312E322E322E31> \MusicPDFRef{16638}] /Limits [<74686D2E31312E322E302E35> <74686D2E31312E322E322E31>] >>}
\MusicPDFWriteObject{10176}{<< /Names [<74686D2E31312E322E322E32> \MusicPDFRef{16639} <74686D2E31312E332E322E31> \MusicPDFRef{16640} <74686D2E31322E312E312E31> \MusicPDFRef{16641} <74686D2E31322E312E312E32> \MusicPDFRef{16642} <74686D2E31322E312E312E33> \MusicPDFRef{16643} <74686D2E31322E312E312E34> \MusicPDFRef{16644}] /Limits [<74686D2E31312E322E322E32> <74686D2E31322E312E312E34>] >>}
\MusicPDFWriteObject{10177}{<< /Names [<74686D2E31322E322E312E31> \MusicPDFRef{16645} <74686D2E31322E322E312E32> \MusicPDFRef{16646} <74686D2E31322E322E312E33> \MusicPDFRef{16647} <74686D2E31322E322E312E34> \MusicPDFRef{16648} <74686D2E31322E322E322E31> \MusicPDFRef{16649} <74686D2E31322E322E322E32> \MusicPDFRef{16650}] /Limits [<74686D2E31322E322E312E31> <74686D2E31322E322E322E32>] >>}
\MusicPDFWriteObject{10178}{<< /Names [<74686D2E31322E322E322E33> \MusicPDFRef{16651} <74686D2E31322E322E322E34> \MusicPDFRef{16652} <74686D2E31322E342E312E31> \MusicPDFRef{16653} <74686D2E31322E352E312E31> \MusicPDFRef{16654} <74686D2E31332E332E312E31> \MusicPDFRef{16655} <74686D2E31332E332E312E32> \MusicPDFRef{16656}] /Limits [<74686D2E31322E322E322E33> <74686D2E31332E332E312E32>] >>}
\MusicPDFWriteObject{10179}{<< /Names [<74686D2E31342E312E312E31> \MusicPDFRef{16657} <74686D2E31342E312E322E31> \MusicPDFRef{16658} <74686D2E31342E312E322E32> \MusicPDFRef{16659} <74686D2E31342E312E322E33> \MusicPDFRef{16660} <74686D2E31342E312E322E34> \MusicPDFRef{16661} <74686D2E31342E312E322E35> \MusicPDFRef{16662}] /Limits [<74686D2E31342E312E312E31> <74686D2E31342E312E322E35>] >>}
\MusicPDFWriteObject{10180}{<< /Names [<74686D2E31342E312E322E36> \MusicPDFRef{16663} <74686D2E31342E312E322E37> \MusicPDFRef{16664} <74686D2E31342E312E322E38> \MusicPDFRef{16665} <74686D2E322E312E312E31> \MusicPDFRef{16666} <74686D2E322E312E312E32> \MusicPDFRef{16667} <74686D2E322E312E312E33> \MusicPDFRef{16668}] /Limits [<74686D2E31342E312E322E36> <74686D2E322E312E312E33>] >>}
\MusicPDFWriteObject{10181}{<< /Names [<74686D2E322E312E312E34> \MusicPDFRef{16669} <74686D2E322E312E312E35> \MusicPDFRef{16670} <74686D2E322E312E312E36> \MusicPDFRef{16671} <74686D2E322E312E312E37> \MusicPDFRef{16672} <74686D2E322E312E322E31> \MusicPDFRef{16673} <74686D2E322E312E322E3130> \MusicPDFRef{16674}] /Limits [<74686D2E322E312E312E34> <74686D2E322E312E322E3130>] >>}
\MusicPDFWriteObject{10182}{<< /Names [<74686D2E322E312E322E3131> \MusicPDFRef{16675} <74686D2E322E312E322E3132> \MusicPDFRef{16676} <74686D2E322E312E322E32> \MusicPDFRef{16677} <74686D2E322E312E322E33> \MusicPDFRef{16678} <74686D2E322E312E322E34> \MusicPDFRef{16679} <74686D2E322E312E322E35> \MusicPDFRef{16680}] /Limits [<74686D2E322E312E322E3131> <74686D2E322E312E322E35>] >>}
\MusicPDFWriteObject{10183}{<< /Names [<74686D2E322E312E322E36> \MusicPDFRef{16681} <74686D2E322E312E322E37> \MusicPDFRef{16682} <74686D2E322E312E322E38> \MusicPDFRef{16683} <74686D2E322E312E322E39> \MusicPDFRef{16684} <74686D2E322E312E332E31> \MusicPDFRef{16685} <74686D2E322E312E332E32> \MusicPDFRef{16686}] /Limits [<74686D2E322E312E322E36> <74686D2E322E312E332E32>] >>}
\MusicPDFWriteObject{10184}{<< /Names [<74686D2E322E312E332E33> \MusicPDFRef{16687} <74686D2E322E312E332E34> \MusicPDFRef{16688} <74686D2E322E312E332E35> \MusicPDFRef{16689} <74686D2E322E312E342E31> \MusicPDFRef{16690} <74686D2E322E312E342E32> \MusicPDFRef{16691} <74686D2E322E312E342E33> \MusicPDFRef{16692}] /Limits [<74686D2E322E312E332E33> <74686D2E322E312E342E33>] >>}
\MusicPDFWriteObject{10185}{<< /Names [<74686D2E322E312E342E34> \MusicPDFRef{16693} <74686D2E322E322E312E31> \MusicPDFRef{16694} <74686D2E322E322E312E32> \MusicPDFRef{16695} <74686D2E322E322E312E33> \MusicPDFRef{16696} <74686D2E322E322E322E31> \MusicPDFRef{16697} <74686D2E322E322E322E32> \MusicPDFRef{16698}] /Limits [<74686D2E322E312E342E34> <74686D2E322E322E322E32>] >>}
\MusicPDFWriteObject{10186}{<< /Names [<74686D2E322E322E322E33> \MusicPDFRef{16699} <74686D2E322E322E332E31> \MusicPDFRef{16700} <74686D2E322E322E332E32> \MusicPDFRef{16701} <74686D2E322E322E332E33> \MusicPDFRef{16702} <74686D2E322E322E332E34> \MusicPDFRef{16703} <74686D2E322E322E342E31> \MusicPDFRef{16704}] /Limits [<74686D2E322E322E322E33> <74686D2E322E322E342E31>] >>}
\MusicPDFWriteObject{10187}{<< /Names [<74686D2E322E322E342E32> \MusicPDFRef{16705} <74686D2E322E322E342E33> \MusicPDFRef{16706} <74686D2E322E322E342E34> \MusicPDFRef{16707} <74686D2E322E322E342E35> \MusicPDFRef{16708} <74686D2E322E322E342E36> \MusicPDFRef{16709} <74686D2E322E322E352E31> \MusicPDFRef{16710}] /Limits [<74686D2E322E322E342E32> <74686D2E322E322E352E31>] >>}
\MusicPDFWriteObject{10188}{<< /Names [<74686D2E322E322E352E32> \MusicPDFRef{16711} <74686D2E322E322E352E33> \MusicPDFRef{16712} <74686D2E322E322E352E34> \MusicPDFRef{16713} <74686D2E322E332E312E31> \MusicPDFRef{16714} <74686D2E322E332E312E32> \MusicPDFRef{16715} <74686D2E322E332E312E33> \MusicPDFRef{16716}] /Limits [<74686D2E322E322E352E32> <74686D2E322E332E312E33>] >>}
\MusicPDFWriteObject{10189}{<< /Names [<74686D2E322E332E312E34> \MusicPDFRef{16717} <74686D2E322E332E322E31> \MusicPDFRef{16718} <74686D2E322E332E322E32> \MusicPDFRef{16719} <74686D2E322E332E322E33> \MusicPDFRef{16720} <74686D2E322E342E312E31> \MusicPDFRef{16721} <74686D2E322E342E312E32> \MusicPDFRef{16722}] /Limits [<74686D2E322E332E312E34> <74686D2E322E342E312E32>] >>}
\MusicPDFWriteObject{10190}{<< /Names [<74686D2E322E342E322E31> \MusicPDFRef{16723} <74686D2E322E342E332E31> \MusicPDFRef{16724} <74686D2E322E342E332E32> \MusicPDFRef{16725} <74686D2E332E312E312E31> \MusicPDFRef{16726} <74686D2E332E312E312E32> \MusicPDFRef{16727} <74686D2E332E312E322E31> \MusicPDFRef{16728}] /Limits [<74686D2E322E342E322E31> <74686D2E332E312E322E31>] >>}
\MusicPDFWriteObject{10191}{<< /Names [<74686D2E332E312E322E32> \MusicPDFRef{16729} <74686D2E332E322E312E31> \MusicPDFRef{16730} <74686D2E332E322E312E32> \MusicPDFRef{16731} <74686D2E332E322E312E33> \MusicPDFRef{16732} <74686D2E332E322E312E34> \MusicPDFRef{16733} <74686D2E332E322E322E31> \MusicPDFRef{16734}] /Limits [<74686D2E332E312E322E32> <74686D2E332E322E322E31>] >>}
\MusicPDFWriteObject{10192}{<< /Names [<74686D2E332E322E322E32> \MusicPDFRef{16735} <74686D2E332E322E322E33> \MusicPDFRef{16736} <74686D2E332E322E322E34> \MusicPDFRef{16737} <74686D2E332E322E322E35> \MusicPDFRef{16738} <74686D2E332E332E312E31> \MusicPDFRef{16739} <74686D2E332E332E312E32> \MusicPDFRef{16740}] /Limits [<74686D2E332E322E322E32> <74686D2E332E332E312E32>] >>}
\MusicPDFWriteObject{10193}{<< /Names [<74686D2E332E332E312E33> \MusicPDFRef{16741} <74686D2E332E332E312E34> \MusicPDFRef{16742} <74686D2E332E332E312E35> \MusicPDFRef{16743} <74686D2E332E332E312E36> \MusicPDFRef{16744} <74686D2E332E332E312E37> \MusicPDFRef{16745} <74686D2E332E332E312E38> \MusicPDFRef{16746}] /Limits [<74686D2E332E332E312E33> <74686D2E332E332E312E38>] >>}
\MusicPDFWriteObject{10194}{<< /Names [<74686D2E332E332E322E31> \MusicPDFRef{16747} <74686D2E332E332E322E32> \MusicPDFRef{16748} <74686D2E332E342E312E31> \MusicPDFRef{16749} <74686D2E332E342E322E31> \MusicPDFRef{16750} <74686D2E332E342E322E32> \MusicPDFRef{16751} <74686D2E332E342E322E33> \MusicPDFRef{16752}] /Limits [<74686D2E332E332E322E31> <74686D2E332E342E322E33>] >>}
\MusicPDFWriteObject{10195}{<< /Names [<74686D2E332E342E322E34> \MusicPDFRef{16753} <74686D2E332E342E332E31> \MusicPDFRef{16754} <74686D2E332E342E332E32> \MusicPDFRef{16755} <74686D2E332E342E332E33> \MusicPDFRef{16756} <74686D2E332E342E332E34> \MusicPDFRef{16757} <74686D2E342E312E312E31> \MusicPDFRef{16758}] /Limits [<74686D2E332E342E322E34> <74686D2E342E312E312E31>] >>}
\MusicPDFWriteObject{10196}{<< /Names [<74686D2E342E312E322E31> \MusicPDFRef{16759} <74686D2E342E312E322E32> \MusicPDFRef{16760} <74686D2E342E312E322E33> \MusicPDFRef{16761} <74686D2E342E312E322E34> \MusicPDFRef{16762} <74686D2E342E312E322E35> \MusicPDFRef{16763} <74686D2E342E312E322E36> \MusicPDFRef{16764}] /Limits [<74686D2E342E312E322E31> <74686D2E342E312E322E36>] >>}
\MusicPDFWriteObject{10197}{<< /Names [<74686D2E342E322E312E31> \MusicPDFRef{16765} <74686D2E342E322E312E32> \MusicPDFRef{16766} <74686D2E342E322E312E33> \MusicPDFRef{16767} <74686D2E342E322E312E34> \MusicPDFRef{16768} <74686D2E342E322E312E35> \MusicPDFRef{16769} <74686D2E342E322E312E36> \MusicPDFRef{16770}] /Limits [<74686D2E342E322E312E31> <74686D2E342E322E312E36>] >>}
\MusicPDFWriteObject{10198}{<< /Names [<74686D2E342E322E312E37> \MusicPDFRef{16771} <74686D2E342E322E312E38> \MusicPDFRef{16772} <74686D2E342E322E322E31> \MusicPDFRef{16773} <74686D2E342E322E322E32> \MusicPDFRef{16774} <74686D2E342E322E322E33> \MusicPDFRef{16775} <74686D2E342E322E322E34> \MusicPDFRef{16776}] /Limits [<74686D2E342E322E312E37> <74686D2E342E322E322E34>] >>}
\MusicPDFWriteObject{10199}{<< /Names [<74686D2E342E322E322E35> \MusicPDFRef{16777} <74686D2E342E322E322E36> \MusicPDFRef{16778} <74686D2E342E332E302E31> \MusicPDFRef{16779} <74686D2E342E332E312E31> \MusicPDFRef{16780} <74686D2E342E332E312E32> \MusicPDFRef{16781} <74686D2E342E332E312E33> \MusicPDFRef{16782}] /Limits [<74686D2E342E322E322E35> <74686D2E342E332E312E33>] >>}
\MusicPDFWriteObject{10200}{<< /Names [<74686D2E342E332E312E34> \MusicPDFRef{16783} <74686D2E342E332E312E35> \MusicPDFRef{16784} <74686D2E342E332E312E36> \MusicPDFRef{16785} <74686D2E342E332E322E31> \MusicPDFRef{16786} <74686D2E342E332E322E32> \MusicPDFRef{16787} <74686D2E342E332E332E31> \MusicPDFRef{16788}] /Limits [<74686D2E342E332E312E34> <74686D2E342E332E332E31>] >>}
\MusicPDFWriteObject{10201}{<< /Names [<74686D2E342E332E332E32> \MusicPDFRef{16789} <74686D2E342E332E332E33> \MusicPDFRef{16790} <74686D2E342E332E332E34> \MusicPDFRef{16791} <74686D2E342E332E332E35> \MusicPDFRef{16792} <74686D2E342E332E332E36> \MusicPDFRef{16793} <74686D2E342E332E332E37> \MusicPDFRef{16794}] /Limits [<74686D2E342E332E332E32> <74686D2E342E332E332E37>] >>}
\MusicPDFWriteObject{10202}{<< /Names [<74686D2E342E332E332E38> \MusicPDFRef{16795} <74686D2E342E342E312E31> \MusicPDFRef{16796} <74686D2E342E342E322E31> \MusicPDFRef{16797} <74686D2E342E342E322E32> \MusicPDFRef{16798} <74686D2E342E342E322E33> \MusicPDFRef{16799} <74686D2E342E342E322E34> \MusicPDFRef{16800}] /Limits [<74686D2E342E332E332E38> <74686D2E342E342E322E34>] >>}
\MusicPDFWriteObject{10203}{<< /Names [<74686D2E342E342E322E35> \MusicPDFRef{16801} <74686D2E342E342E322E36> \MusicPDFRef{16802} <74686D2E342E342E322E37> \MusicPDFRef{16803} <74686D2E342E342E322E38> \MusicPDFRef{16804} <74686D2E342E342E322E39> \MusicPDFRef{16805} <74686D2E342E352E312E31> \MusicPDFRef{16806}] /Limits [<74686D2E342E342E322E35> <74686D2E342E352E312E31>] >>}
\MusicPDFWriteObject{10204}{<< /Names [<74686D2E342E352E312E32> \MusicPDFRef{16807} <74686D2E342E352E312E33> \MusicPDFRef{16808} <74686D2E342E352E312E34> \MusicPDFRef{16809} <74686D2E342E352E312E35> \MusicPDFRef{16810} <74686D2E342E352E312E36> \MusicPDFRef{16811} <74686D2E342E352E322E31> \MusicPDFRef{16812}] /Limits [<74686D2E342E352E312E32> <74686D2E342E352E322E31>] >>}
\MusicPDFWriteObject{10205}{<< /Names [<74686D2E342E352E322E32> \MusicPDFRef{16813} <74686D2E342E352E322E33> \MusicPDFRef{16814} <74686D2E342E352E332E31> \MusicPDFRef{16815} <74686D2E342E362E312E31> \MusicPDFRef{16816} <74686D2E342E362E312E32> \MusicPDFRef{16817} <74686D2E342E362E312E33> \MusicPDFRef{16818}] /Limits [<74686D2E342E352E322E32> <74686D2E342E362E312E33>] >>}
\MusicPDFWriteObject{10206}{<< /Names [<74686D2E342E362E312E34> \MusicPDFRef{16819} <74686D2E342E362E312E35> \MusicPDFRef{16820} <74686D2E342E362E312E36> \MusicPDFRef{16821} <74686D2E342E362E322E31> \MusicPDFRef{16822} <74686D2E342E362E332E31> \MusicPDFRef{16823} <74686D2E352E312E312E31> \MusicPDFRef{16824}] /Limits [<74686D2E342E362E312E34> <74686D2E352E312E312E31>] >>}
\MusicPDFWriteObject{10207}{<< /Names [<74686D2E352E312E312E32> \MusicPDFRef{16825} <74686D2E352E312E312E33> \MusicPDFRef{16826} <74686D2E352E312E322E31> \MusicPDFRef{16827} <74686D2E352E312E322E32> \MusicPDFRef{16828} <74686D2E352E312E332E31> \MusicPDFRef{16829} <74686D2E352E322E312E31> \MusicPDFRef{16830}] /Limits [<74686D2E352E312E312E32> <74686D2E352E322E312E31>] >>}
\MusicPDFWriteObject{10208}{<< /Names [<74686D2E352E322E322E31> \MusicPDFRef{16831} <74686D2E352E322E332E31> \MusicPDFRef{16832} <74686D2E352E322E332E32> \MusicPDFRef{16833} <74686D2E352E322E332E33> \MusicPDFRef{16834} <74686D2E352E322E332E34> \MusicPDFRef{16835} <74686D2E352E322E332E35> \MusicPDFRef{16836}] /Limits [<74686D2E352E322E322E31> <74686D2E352E322E332E35>] >>}
\MusicPDFWriteObject{10209}{<< /Names [<74686D2E352E322E342E31> \MusicPDFRef{16837} <74686D2E352E332E312E31> \MusicPDFRef{16838} <74686D2E352E332E312E32> \MusicPDFRef{16839} <74686D2E352E332E312E33> \MusicPDFRef{16840} <74686D2E352E332E312E34> \MusicPDFRef{16841} <74686D2E352E332E312E35> \MusicPDFRef{16842}] /Limits [<74686D2E352E322E342E31> <74686D2E352E332E312E35>] >>}
\MusicPDFWriteObject{10210}{<< /Names [<74686D2E352E332E322E31> \MusicPDFRef{16843} <74686D2E352E332E322E32> \MusicPDFRef{16844} <74686D2E352E332E322E33> \MusicPDFRef{16845} <74686D2E352E332E322E34> \MusicPDFRef{16846} <74686D2E352E332E322E35> \MusicPDFRef{16847} <74686D2E352E332E332E31> \MusicPDFRef{16848}] /Limits [<74686D2E352E332E322E31> <74686D2E352E332E332E31>] >>}
\MusicPDFWriteObject{10211}{<< /Names [<74686D2E352E332E332E32> \MusicPDFRef{16849} <74686D2E352E332E332E33> \MusicPDFRef{16850} <74686D2E352E332E332E34> \MusicPDFRef{16851} <74686D2E352E332E342E31> \MusicPDFRef{16852} <74686D2E352E332E342E32> \MusicPDFRef{16853} <74686D2E352E332E342E33> \MusicPDFRef{16854}] /Limits [<74686D2E352E332E332E32> <74686D2E352E332E342E33>] >>}
\MusicPDFWriteObject{10212}{<< /Names [<74686D2E352E332E342E34> \MusicPDFRef{16855} <74686D2E352E332E342E35> \MusicPDFRef{16856} <74686D2E352E342E312E31> \MusicPDFRef{16857} <74686D2E352E342E312E32> \MusicPDFRef{16858} <74686D2E352E342E312E33> \MusicPDFRef{16859} <74686D2E352E342E312E34> \MusicPDFRef{16860}] /Limits [<74686D2E352E332E342E34> <74686D2E352E342E312E34>] >>}
\MusicPDFWriteObject{10213}{<< /Names [<74686D2E352E342E312E35> \MusicPDFRef{16861} <74686D2E352E342E322E31> \MusicPDFRef{16862} <74686D2E352E342E322E32> \MusicPDFRef{16863} <74686D2E352E342E322E33> \MusicPDFRef{16864} <74686D2E352E352E312E31> \MusicPDFRef{16865} <74686D2E352E352E312E32> \MusicPDFRef{16866}] /Limits [<74686D2E352E342E312E35> <74686D2E352E352E312E32>] >>}
\MusicPDFWriteObject{10214}{<< /Names [<74686D2E352E352E312E33> \MusicPDFRef{16867} <74686D2E352E352E312E34> \MusicPDFRef{16868} <74686D2E352E352E322E31> \MusicPDFRef{16869} <74686D2E352E352E322E32> \MusicPDFRef{16870} <74686D2E352E352E332E31> \MusicPDFRef{16871} <74686D2E352E352E332E32> \MusicPDFRef{16872}] /Limits [<74686D2E352E352E312E33> <74686D2E352E352E332E32>] >>}
\MusicPDFWriteObject{10215}{<< /Names [<74686D2E352E352E332E33> \MusicPDFRef{16873} <74686D2E352E352E332E34> \MusicPDFRef{16874} <74686D2E362E312E312E31> \MusicPDFRef{16875} <74686D2E362E312E322E31> \MusicPDFRef{16876} <74686D2E362E312E322E32> \MusicPDFRef{16877} <74686D2E362E312E322E33> \MusicPDFRef{16878}] /Limits [<74686D2E352E352E332E33> <74686D2E362E312E322E33>] >>}
\MusicPDFWriteObject{10216}{<< /Names [<74686D2E362E312E322E34> \MusicPDFRef{16879} <74686D2E362E322E312E31> \MusicPDFRef{16880} <74686D2E362E322E312E32> \MusicPDFRef{16881} <74686D2E362E322E312E33> \MusicPDFRef{16882} <74686D2E362E322E312E34> \MusicPDFRef{16883} <74686D2E362E322E322E31> \MusicPDFRef{16884}] /Limits [<74686D2E362E312E322E34> <74686D2E362E322E322E31>] >>}
\MusicPDFWriteObject{10217}{<< /Names [<74686D2E362E322E322E32> \MusicPDFRef{16885} <74686D2E362E332E322E31> \MusicPDFRef{16886} <74686D2E362E332E322E32> \MusicPDFRef{16887} <74686D2E372E312E312E31> \MusicPDFRef{16888} <74686D2E372E312E322E31> \MusicPDFRef{16889} <74686D2E372E312E322E32> \MusicPDFRef{16890}] /Limits [<74686D2E362E322E322E32> <74686D2E372E312E322E32>] >>}
\MusicPDFWriteObject{10218}{<< /Names [<74686D2E372E312E322E33> \MusicPDFRef{16891} <74686D2E372E312E322E34> \MusicPDFRef{16892} <74686D2E372E312E332E31> \MusicPDFRef{16893} <74686D2E372E312E332E32> \MusicPDFRef{16894} <74686D2E372E312E332E33> \MusicPDFRef{16895} <74686D2E372E312E342E31> \MusicPDFRef{16896}] /Limits [<74686D2E372E312E322E33> <74686D2E372E312E342E31>] >>}
\MusicPDFWriteObject{10219}{<< /Names [<74686D2E372E312E342E32> \MusicPDFRef{16897} <74686D2E372E312E342E33> \MusicPDFRef{16898} <74686D2E372E322E312E31> \MusicPDFRef{16899} <74686D2E372E322E312E32> \MusicPDFRef{16900} <74686D2E372E322E312E33> \MusicPDFRef{16901} <74686D2E372E322E312E34> \MusicPDFRef{16902}] /Limits [<74686D2E372E312E342E32> <74686D2E372E322E312E34>] >>}
\MusicPDFWriteObject{10220}{<< /Names [<74686D2E372E322E322E31> \MusicPDFRef{16903} <74686D2E372E322E322E32> \MusicPDFRef{16904} <74686D2E372E332E312E31> \MusicPDFRef{16905} <74686D2E372E332E312E32> \MusicPDFRef{16906} <74686D2E372E332E312E33> \MusicPDFRef{16907} <74686D2E372E332E312E34> \MusicPDFRef{16908}] /Limits [<74686D2E372E322E322E31> <74686D2E372E332E312E34>] >>}
\MusicPDFWriteObject{10221}{<< /Names [<74686D2E372E332E312E35> \MusicPDFRef{16909} <74686D2E372E332E312E36> \MusicPDFRef{16910} <74686D2E372E332E312E37> \MusicPDFRef{16911} <74686D2E372E332E312E38> \MusicPDFRef{16912} <74686D2E372E332E322E31> \MusicPDFRef{16913} <74686D2E372E332E322E32> \MusicPDFRef{16914}] /Limits [<74686D2E372E332E312E35> <74686D2E372E332E322E32>] >>}
\MusicPDFWriteObject{10222}{<< /Names [<74686D2E372E332E322E33> \MusicPDFRef{16915} <74686D2E372E332E322E34> \MusicPDFRef{16916} <74686D2E372E332E332E31> \MusicPDFRef{16917} <74686D2E372E332E332E32> \MusicPDFRef{16918} <74686D2E372E332E332E33> \MusicPDFRef{16919} <74686D2E372E332E332E34> \MusicPDFRef{16920}] /Limits [<74686D2E372E332E322E33> <74686D2E372E332E332E34>] >>}
\MusicPDFWriteObject{10223}{<< /Names [<74686D2E372E332E342E31> \MusicPDFRef{16921} <74686D2E372E342E312E31> \MusicPDFRef{16922} <74686D2E372E342E322E31> \MusicPDFRef{16923} <74686D2E372E342E332E31> \MusicPDFRef{16924} <74686D2E372E342E332E32> \MusicPDFRef{16925} <74686D2E372E342E332E33> \MusicPDFRef{16926}] /Limits [<74686D2E372E332E342E31> <74686D2E372E342E332E33>] >>}
\MusicPDFWriteObject{10224}{<< /Names [<74686D2E372E352E302E31> \MusicPDFRef{16927} <74686D2E372E352E302E32> \MusicPDFRef{16928} <74686D2E372E352E302E33> \MusicPDFRef{16929} <74686D2E382E312E302E31> \MusicPDFRef{16930} <74686D2E382E322E302E31> \MusicPDFRef{16931} <74686D2E382E322E302E32> \MusicPDFRef{16932}] /Limits [<74686D2E372E352E302E31> <74686D2E382E322E302E32>] >>}
\MusicPDFWriteObject{10225}{<< /Names [<74686D2E382E322E302E33> \MusicPDFRef{16933} <74686D2E382E322E302E34> \MusicPDFRef{16934} <74686D2E382E322E302E35> \MusicPDFRef{16935} <74686D2E382E322E302E36> \MusicPDFRef{16936} <74686D2E382E332E302E31> \MusicPDFRef{16937} <74686D2E382E342E302E31> \MusicPDFRef{16938}] /Limits [<74686D2E382E322E302E33> <74686D2E382E342E302E31>] >>}
\MusicPDFWriteObject{10226}{<< /Names [<74686D2E392E312E312E31> \MusicPDFRef{16939} <74686D2E392E312E312E32> \MusicPDFRef{16940} <74686D2E392E312E312E33> \MusicPDFRef{16941} <74686D2E392E312E312E34> \MusicPDFRef{16942} <74686D2E392E312E312E35> \MusicPDFRef{16943} <74686D2E392E312E322E31> \MusicPDFRef{16944}] /Limits [<74686D2E392E312E312E31> <74686D2E392E312E322E31>] >>}
\MusicPDFWriteObject{10227}{<< /Names [<74686D2E392E312E322E32> \MusicPDFRef{16945} <74686D2E392E312E322E33> \MusicPDFRef{16946} <74686D2E392E312E322E34> \MusicPDFRef{16947} <74686D2E392E312E322E35> \MusicPDFRef{16948} <74686D2E392E312E322E36> \MusicPDFRef{16949} <74686D2E392E322E302E31> \MusicPDFRef{16950}] /Limits [<74686D2E392E312E322E32> <74686D2E392E322E302E31>] >>}
\MusicPDFWriteObject{10228}{<< /Names [<74686D2E392E322E312E31> \MusicPDFRef{16951} <74686D2E392E322E312E32> \MusicPDFRef{16952} <74686D2E392E322E312E33> \MusicPDFRef{16953} <74686D2E392E322E312E34> \MusicPDFRef{16954} <74686D2E392E322E312E35> \MusicPDFRef{16955} <74686D2E392E322E312E36> \MusicPDFRef{16956}] /Limits [<74686D2E392E322E312E31> <74686D2E392E322E312E36>] >>}
\MusicPDFWriteObject{10229}{<< /Names [<74686D2E392E322E312E37> \MusicPDFRef{16957} <74686D2E392E322E312E38> \MusicPDFRef{16958} <74686D2E392E322E322E31> \MusicPDFRef{16959} <74686D2E392E322E322E3130> \MusicPDFRef{16960} <74686D2E392E322E322E3131> \MusicPDFRef{16961} <74686D2E392E322E322E3132> \MusicPDFRef{16962}] /Limits [<74686D2E392E322E312E37> <74686D2E392E322E322E3132>] >>}
\MusicPDFWriteObject{10230}{<< /Names [<74686D2E392E322E322E3133> \MusicPDFRef{16963} <74686D2E392E322E322E32> \MusicPDFRef{16964} <74686D2E392E322E322E33> \MusicPDFRef{16965} <74686D2E392E322E322E34> \MusicPDFRef{16966} <74686D2E392E322E322E35> \MusicPDFRef{16967} <74686D2E392E322E322E36> \MusicPDFRef{16968}] /Limits [<74686D2E392E322E322E3133> <74686D2E392E322E322E36>] >>}
\MusicPDFWriteObject{10231}{<< /Names [<74686D2E392E322E322E37> \MusicPDFRef{16969} <74686D2E392E322E322E38> \MusicPDFRef{16970} <74686D2E392E322E322E39> \MusicPDFRef{16971} <74686D2E392E332E312E31> \MusicPDFRef{16972} <74686D2E392E332E312E32> \MusicPDFRef{16973} <74686D2E392E332E312E33> \MusicPDFRef{16974}] /Limits [<74686D2E392E322E322E37> <74686D2E392E332E312E33>] >>}
\MusicPDFWriteObject{10232}{<< /Names [<74686D2E392E332E312E34> \MusicPDFRef{16975} <74686D2E392E332E312E35> \MusicPDFRef{16976} <74686D2E392E332E322E31> \MusicPDFRef{16977} <74686D2E392E332E322E32> \MusicPDFRef{16978} <74686D2E392E332E322E33> \MusicPDFRef{16979} <74686D2E392E332E322E34> \MusicPDFRef{16980}] /Limits [<74686D2E392E332E312E34> <74686D2E392E332E322E34>] >>}
\MusicPDFWriteObject{10233}{<< /Names [<74686D2E392E332E322E35> \MusicPDFRef{16981} <74686D2E392E332E322E36> \MusicPDFRef{16982} <74686D2E392E332E332E31> \MusicPDFRef{16983} <74686D2E392E332E332E32> \MusicPDFRef{16984} <74686D2E392E332E332E33> \MusicPDFRef{16985} <74686D2E392E332E332E34> \MusicPDFRef{16986}] /Limits [<74686D2E392E332E322E35> <74686D2E392E332E332E34>] >>}
\MusicPDFWriteObject{10234}{<< /Names [<74686D2E392E332E332E35> \MusicPDFRef{16987} <74686D2E412E312E312E31> \MusicPDFRef{16988} <74686D2E412E312E312E32> \MusicPDFRef{16989} <74686D2E412E312E312E33> \MusicPDFRef{16990} <74686D2E412E312E322E31> \MusicPDFRef{16991} <74686D2E412E312E332E31> \MusicPDFRef{16992}] /Limits [<74686D2E392E332E332E35> <74686D2E412E312E332E31>] >>}
\MusicPDFWriteObject{10235}{<< /Names [<74686D2E412E312E332E32> \MusicPDFRef{16993} <74686D2E412E312E332E33> \MusicPDFRef{16994} <74686D2E412E312E332E34> \MusicPDFRef{16995} <74686D2E412E312E342E31> \MusicPDFRef{16996} <74686D2E412E312E342E32> \MusicPDFRef{16997} <74686D2E412E312E342E33> \MusicPDFRef{16998}] /Limits [<74686D2E412E312E332E32> <74686D2E412E312E342E33>] >>}
\MusicPDFWriteObject{10236}{<< /Names [<74686D2E412E312E352E31> \MusicPDFRef{16999} <74686D2E412E312E352E32> \MusicPDFRef{17000} <74686D2E412E312E352E33> \MusicPDFRef{17001} <74686D2E412E322E312E31> \MusicPDFRef{17002} <74686D2E412E322E312E32> \MusicPDFRef{17003} <74686D2E412E322E312E33> \MusicPDFRef{17004}] /Limits [<74686D2E412E312E352E31> <74686D2E412E322E312E33>] >>}
\MusicPDFWriteObject{10237}{<< /Names [<74686D2E412E322E322E31> \MusicPDFRef{17005} <74686D2E412E322E322E32> \MusicPDFRef{17006} <74686D2E412E322E322E33> \MusicPDFRef{17007} <74686D2E412E322E322E34> \MusicPDFRef{17008} <74686D2E412E322E322E35> \MusicPDFRef{17009} <74686D2E412E322E322E36> \MusicPDFRef{17010}] /Limits [<74686D2E412E322E322E31> <74686D2E412E322E322E36>] >>}
\MusicPDFWriteObject{10238}{<< /Names [<74686D2E412E322E322E37> \MusicPDFRef{17011} <74686D2E412E322E322E38> \MusicPDFRef{17012} <74686D2E412E322E322E39> \MusicPDFRef{17013} <74686D2E412E322E332E31> \MusicPDFRef{17014} <74686D2E412E322E332E3130> \MusicPDFRef{17015} <74686D2E412E322E332E3131> \MusicPDFRef{17016}] /Limits [<74686D2E412E322E322E37> <74686D2E412E322E332E3131>] >>}
\MusicPDFWriteObject{10239}{<< /Names [<74686D2E412E322E332E3132> \MusicPDFRef{17017} <74686D2E412E322E332E3133> \MusicPDFRef{17018} <74686D2E412E322E332E32> \MusicPDFRef{17019} <74686D2E412E322E332E33> \MusicPDFRef{17020} <74686D2E412E322E332E34> \MusicPDFRef{17021} <74686D2E412E322E332E35> \MusicPDFRef{17022}] /Limits [<74686D2E412E322E332E3132> <74686D2E412E322E332E35>] >>}
\MusicPDFWriteObject{10240}{<< /Names [<74686D2E412E322E332E36> \MusicPDFRef{17023} <74686D2E412E322E332E37> \MusicPDFRef{17024} <74686D2E412E322E332E38> \MusicPDFRef{17025} <74686D2E412E322E332E39> \MusicPDFRef{17026} <74686D2E412E322E342E31> \MusicPDFRef{17027} <74686D2E412E322E342E32> \MusicPDFRef{17028}] /Limits [<74686D2E412E322E332E36> <74686D2E412E322E342E32>] >>}
\MusicPDFWriteObject{10241}{<< /Names [<74686D2E412E322E342E33> \MusicPDFRef{17029} <74686D2E412E322E342E34> \MusicPDFRef{17030} <74686D2E412E322E342E35> \MusicPDFRef{17031} <74686D2E412E332E312E31> \MusicPDFRef{17032} <74686D2E412E332E312E3130> \MusicPDFRef{17033} <74686D2E412E332E312E3131> \MusicPDFRef{17034}] /Limits [<74686D2E412E322E342E33> <74686D2E412E332E312E3131>] >>}
\MusicPDFWriteObject{10242}{<< /Names [<74686D2E412E332E312E32> \MusicPDFRef{17035} <74686D2E412E332E312E33> \MusicPDFRef{17036} <74686D2E412E332E312E34> \MusicPDFRef{17037} <74686D2E412E332E312E35> \MusicPDFRef{17038} <74686D2E412E332E312E36> \MusicPDFRef{17039} <74686D2E412E332E312E37> \MusicPDFRef{17040}] /Limits [<74686D2E412E332E312E32> <74686D2E412E332E312E37>] >>}
\MusicPDFWriteObject{10243}{<< /Names [<74686D2E412E332E312E38> \MusicPDFRef{17041} <74686D2E412E332E312E39> \MusicPDFRef{17042} <74686D2E412E332E322E31> \MusicPDFRef{17043} <74686D2E412E332E322E32> \MusicPDFRef{17044} <74686D2E412E332E322E33> \MusicPDFRef{17045} <74686D2E412E332E322E34> \MusicPDFRef{17046}] /Limits [<74686D2E412E332E312E38> <74686D2E412E332E322E34>] >>}
\MusicPDFWriteObject{10244}{<< /Names [<74686D2E412E332E322E35> \MusicPDFRef{17047} <74686D2E412E332E342E31> \MusicPDFRef{17048} <74686D2E412E332E342E32> \MusicPDFRef{17049} <74686D2E412E342E312E31> \MusicPDFRef{17050} <74686D2E412E342E312E32> \MusicPDFRef{17051} <74686D2E412E352E312E31> \MusicPDFRef{17052}] /Limits [<74686D2E412E332E322E35> <74686D2E412E352E312E31>] >>}
\MusicPDFWriteObject{10245}{<< /Names [<74686D2E412E352E312E32> \MusicPDFRef{17053} <74686D2E412E352E312E33> \MusicPDFRef{17054} <74686D2E412E352E322E31> \MusicPDFRef{17055} <74686D2E412E352E322E3130> \MusicPDFRef{17056} <74686D2E412E352E322E3131> \MusicPDFRef{17057} <74686D2E412E352E322E3132> \MusicPDFRef{17058}] /Limits [<74686D2E412E352E312E32> <74686D2E412E352E322E3132>] >>}
\MusicPDFWriteObject{10246}{<< /Names [<74686D2E412E352E322E3133> \MusicPDFRef{17059} <74686D2E412E352E322E3134> \MusicPDFRef{17060} <74686D2E412E352E322E3135> \MusicPDFRef{17061} <74686D2E412E352E322E3136> \MusicPDFRef{17062} <74686D2E412E352E322E3137> \MusicPDFRef{17063} <74686D2E412E352E322E3138> \MusicPDFRef{17064}] /Limits [<74686D2E412E352E322E3133> <74686D2E412E352E322E3138>] >>}
\MusicPDFWriteObject{10247}{<< /Names [<74686D2E412E352E322E3139> \MusicPDFRef{17065} <74686D2E412E352E322E32> \MusicPDFRef{17066} <74686D2E412E352E322E33> \MusicPDFRef{17067} <74686D2E412E352E322E34> \MusicPDFRef{17068} <74686D2E412E352E322E35> \MusicPDFRef{17069} <74686D2E412E352E322E36> \MusicPDFRef{17070}] /Limits [<74686D2E412E352E322E3139> <74686D2E412E352E322E36>] >>}
\MusicPDFWriteObject{10248}{<< /Names [<74686D2E412E352E322E37> \MusicPDFRef{17071} <74686D2E412E352E322E38> \MusicPDFRef{17072} <74686D2E412E352E322E39> \MusicPDFRef{17073} <74686D2E412E352E332E31> \MusicPDFRef{17074} <74686D2E412E352E332E32> \MusicPDFRef{17075} <74686D2E412E352E332E33> \MusicPDFRef{17076}] /Limits [<74686D2E412E352E322E37> <74686D2E412E352E332E33>] >>}
\MusicPDFWriteObject{10249}{<< /Names [<74686D2E412E352E332E34> \MusicPDFRef{17077} <74686D2E412E352E332E35> \MusicPDFRef{17078} <74686D2E412E352E332E36> \MusicPDFRef{17079} <74686D2E412E352E332E37> \MusicPDFRef{17080} <74686D2E412E352E342E31> \MusicPDFRef{17081} <74686D2E412E352E342E32> \MusicPDFRef{17082}] /Limits [<74686D2E412E352E332E34> <74686D2E412E352E342E32>] >>}
\MusicPDFWriteObject{10250}{<< /Names [<74686D2E412E352E342E33> \MusicPDFRef{17083} <74686D2E422E312E312E31> \MusicPDFRef{17084} <74686D2E422E312E312E32> \MusicPDFRef{17085} <74686D2E422E312E312E33> \MusicPDFRef{17086} <74686D2E422E312E312E34> \MusicPDFRef{17087} <74686D2E422E322E312E31> \MusicPDFRef{17088}] /Limits [<74686D2E412E352E342E33> <74686D2E422E322E312E31>] >>}
\MusicPDFWriteObject{10251}{<< /Names [<74686D2E422E322E312E3130> \MusicPDFRef{17089} <74686D2E422E322E312E32> \MusicPDFRef{17090} <74686D2E422E322E312E33> \MusicPDFRef{17091} <74686D2E422E322E312E34> \MusicPDFRef{17092} <74686D2E422E322E312E35> \MusicPDFRef{17093} <74686D2E422E322E312E36> \MusicPDFRef{17094}] /Limits [<74686D2E422E322E312E3130> <74686D2E422E322E312E36>] >>}
\MusicPDFWriteObject{10252}{<< /Names [<74686D2E422E322E312E37> \MusicPDFRef{17095} <74686D2E422E322E312E38> \MusicPDFRef{17096} <74686D2E422E322E312E39> \MusicPDFRef{17097} <74686D2E422E322E322E31> \MusicPDFRef{17098} <74686D2E422E322E322E32> \MusicPDFRef{17099} <74686D2E422E322E322E33> \MusicPDFRef{17100}] /Limits [<74686D2E422E322E312E37> <74686D2E422E322E322E33>] >>}
\MusicPDFWriteObject{10253}{<< /Names [<74686D2E422E322E322E34> \MusicPDFRef{17101} <74686D2E422E322E322E35> \MusicPDFRef{17102} <74686D2E422E322E322E36> \MusicPDFRef{17103} <74686D2E422E332E312E31> \MusicPDFRef{17104} <74686D2E422E332E312E32> \MusicPDFRef{17105} <74686D2E422E332E322E31> \MusicPDFRef{17106}] /Limits [<74686D2E422E322E322E34> <74686D2E422E332E322E31>] >>}
\MusicPDFWriteObject{10254}{<< /Names [<74686D2E422E342E312E31> \MusicPDFRef{17107} <74686D2E422E342E312E32> \MusicPDFRef{17108} <74686D2E422E342E312E33> \MusicPDFRef{17109} <74686D2E422E342E312E34> \MusicPDFRef{17110} <74686D2E422E342E312E35> \MusicPDFRef{17111} <74686D2E422E342E312E36> \MusicPDFRef{17112}] /Limits [<74686D2E422E342E312E31> <74686D2E422E342E312E36>] >>}
\MusicPDFWriteObject{10255}{<< /Names [<74686D2E422E342E312E37> \MusicPDFRef{17113} <74686D2E422E342E312E38> \MusicPDFRef{17114} <74686D2E422E342E312E39> \MusicPDFRef{17115} <74686D2E422E342E322E31> \MusicPDFRef{17116} <74686D2E422E342E322E32> \MusicPDFRef{17117} <74686D2E422E342E332E31> \MusicPDFRef{17118}] /Limits [<74686D2E422E342E312E37> <74686D2E422E342E332E31>] >>}
\MusicPDFWriteObject{10256}{<< /Names [<74686D2E422E342E332E32> \MusicPDFRef{17119} <74686D2E422E342E342E31> \MusicPDFRef{17120} <74686D2E422E342E342E32> \MusicPDFRef{17121} <74686D2E422E342E342E33> \MusicPDFRef{17122} <74686D2E422E342E342E34> \MusicPDFRef{17123} <74686D2E422E342E342E35> \MusicPDFRef{17124}] /Limits [<74686D2E422E342E332E32> <74686D2E422E342E342E35>] >>}
\MusicPDFWriteObject{10257}{<< /Names [<74686D2E422E342E342E36> \MusicPDFRef{17125} <74686D2E422E342E342E37> \MusicPDFRef{17126} <74686D2E422E342E342E38> \MusicPDFRef{17127} <74686D2E422E342E342E39> \MusicPDFRef{17128} <74686D2E422E352E312E31> \MusicPDFRef{17129} <74686D2E422E352E322E31> \MusicPDFRef{17130}] /Limits [<74686D2E422E342E342E36> <74686D2E422E352E322E31>] >>}
\MusicPDFWriteObject{10258}{<< /Names [<74686D2E422E352E322E32> \MusicPDFRef{17131} <74686D2E422E352E322E33> \MusicPDFRef{17132} <74686D2E422E352E322E34> \MusicPDFRef{17133} <74686D2E422E352E322E35> \MusicPDFRef{17134} <74686D2E422E352E332E31> \MusicPDFRef{17135} <74686D2E422E352E332E32> \MusicPDFRef{17136}] /Limits [<74686D2E422E352E322E32> <74686D2E422E352E332E32>] >>}
\MusicPDFWriteObject{10259}{<< /Names [<74686D2E422E352E332E33> \MusicPDFRef{17137} <74686D2E422E352E332E34> \MusicPDFRef{17138} <74686D2E422E352E332E35> \MusicPDFRef{17139} <74686D2E422E362E312E31> \MusicPDFRef{17140} <74686D2E422E362E312E32> \MusicPDFRef{17141} <74686D2E422E362E312E33> \MusicPDFRef{17142}] /Limits [<74686D2E422E352E332E33> <74686D2E422E362E312E33>] >>}
\MusicPDFWriteObject{10260}{<< /Names [<74686D2E422E362E312E34> \MusicPDFRef{17143} <74686D2E422E362E312E35> \MusicPDFRef{17144} <74686D2E422E362E312E36> \MusicPDFRef{17145} <74686D2E422E362E312E37> \MusicPDFRef{17146} <74686D2E422E362E312E38> \MusicPDFRef{17147} <74686D2E422E372E312E31> \MusicPDFRef{17148}] /Limits [<74686D2E422E362E312E34> <74686D2E422E372E312E31>] >>}
\MusicPDFWriteObject{10261}{<< /Names [<74686D2E422E372E312E32> \MusicPDFRef{17149} <74686D2E422E372E322E31> \MusicPDFRef{17150} <74686D2E422E372E322E32> \MusicPDFRef{17151} <74686D2E422E372E322E33> \MusicPDFRef{17152} <74686D2E422E372E322E34> \MusicPDFRef{17153} <74686D2E422E372E322E35> \MusicPDFRef{17154}] /Limits [<74686D2E422E372E312E32> <74686D2E422E372E322E35>] >>}
\MusicPDFWriteObject{10262}{<< /Names [<74686D2E422E372E322E36> \MusicPDFRef{17155} <74686D2E422E372E322E37> \MusicPDFRef{17156} <74686D2E432E312E312E31> \MusicPDFRef{17157} <74686D2E432E312E312E32> \MusicPDFRef{17158} <74686D2E432E312E322E31> \MusicPDFRef{17159} <74686D2E432E312E332E31> \MusicPDFRef{17160}] /Limits [<74686D2E422E372E322E36> <74686D2E432E312E332E31>] >>}
\MusicPDFWriteObject{10263}{<< /Names [<74686D2E432E312E332E32> \MusicPDFRef{17161} <74686D2E432E312E332E33> \MusicPDFRef{17162} <74686D2E432E312E332E34> \MusicPDFRef{17163} <74686D2E432E312E342E31> \MusicPDFRef{17164} <74686D2E432E312E342E32> \MusicPDFRef{17165} <74686D2E432E312E342E33> \MusicPDFRef{17166}] /Limits [<74686D2E432E312E332E32> <74686D2E432E312E342E33>] >>}
\MusicPDFWriteObject{10264}{<< /Names [<74686D2E432E312E342E34> \MusicPDFRef{17167} <74686D2E432E312E342E35> \MusicPDFRef{17168} <74686D2E432E312E342E36> \MusicPDFRef{17169} <74686D2E432E312E342E37> \MusicPDFRef{17170} <74686D2E432E312E352E31> \MusicPDFRef{17171} <74686D2E432E312E352E3130> \MusicPDFRef{17172}] /Limits [<74686D2E432E312E342E34> <74686D2E432E312E352E3130>] >>}
\MusicPDFWriteObject{10265}{<< /Names [<74686D2E432E312E352E3131> \MusicPDFRef{17173} <74686D2E432E312E352E3132> \MusicPDFRef{17174} <74686D2E432E312E352E3133> \MusicPDFRef{17175} <74686D2E432E312E352E32> \MusicPDFRef{17176} <74686D2E432E312E352E33> \MusicPDFRef{17177} <74686D2E432E312E352E34> \MusicPDFRef{17178}] /Limits [<74686D2E432E312E352E3131> <74686D2E432E312E352E34>] >>}
\MusicPDFWriteObject{10266}{<< /Names [<74686D2E432E312E352E35> \MusicPDFRef{17179} <74686D2E432E312E352E36> \MusicPDFRef{17180} <74686D2E432E312E352E37> \MusicPDFRef{17181} <74686D2E432E312E352E38> \MusicPDFRef{17182} <74686D2E432E312E352E39> \MusicPDFRef{17183} <74686D2E432E312E362E31> \MusicPDFRef{17184}] /Limits [<74686D2E432E312E352E35> <74686D2E432E312E362E31>] >>}
\MusicPDFWriteObject{10267}{<< /Names [<74686D2E432E312E362E32> \MusicPDFRef{17185} <74686D2E432E322E312E31> \MusicPDFRef{17186} <74686D2E432E322E312E32> \MusicPDFRef{17187} <74686D2E432E322E312E33> \MusicPDFRef{17188} <74686D2E432E322E312E34> \MusicPDFRef{17189} <74686D2E432E322E312E35> \MusicPDFRef{17190}] /Limits [<74686D2E432E312E362E32> <74686D2E432E322E312E35>] >>}
\MusicPDFWriteObject{10268}{<< /Names [<74686D2E432E322E312E36> \MusicPDFRef{17191} <74686D2E432E322E312E37> \MusicPDFRef{17192} <74686D2E432E322E312E38> \MusicPDFRef{17193} <74686D2E432E322E312E39> \MusicPDFRef{17194} <74686D2E432E322E322E31> \MusicPDFRef{17195} <74686D2E432E322E322E32> \MusicPDFRef{17196}] /Limits [<74686D2E432E322E312E36> <74686D2E432E322E322E32>] >>}
\MusicPDFWriteObject{10269}{<< /Names [<74686D2E432E322E332E31> \MusicPDFRef{17197} <74686D2E432E322E332E3130> \MusicPDFRef{17198} <74686D2E432E322E332E3131> \MusicPDFRef{17199} <74686D2E432E322E332E3132> \MusicPDFRef{17200} <74686D2E432E322E332E3133> \MusicPDFRef{17201} <74686D2E432E322E332E3134> \MusicPDFRef{17202}] /Limits [<74686D2E432E322E332E31> <74686D2E432E322E332E3134>] >>}
\MusicPDFWriteObject{10270}{<< /Names [<74686D2E432E322E332E32> \MusicPDFRef{17203} <74686D2E432E322E332E33> \MusicPDFRef{17204} <74686D2E432E322E332E34> \MusicPDFRef{17205} <74686D2E432E322E332E35> \MusicPDFRef{17206} <74686D2E432E322E332E36> \MusicPDFRef{17207} <74686D2E432E322E332E37> \MusicPDFRef{17208}] /Limits [<74686D2E432E322E332E32> <74686D2E432E322E332E37>] >>}
\MusicPDFWriteObject{10271}{<< /Names [<74686D2E432E322E332E38> \MusicPDFRef{17209} <74686D2E432E322E332E39> \MusicPDFRef{17210} <74686D2E432E322E342E31> \MusicPDFRef{17211} <74686D2E432E322E342E3130> \MusicPDFRef{17212} <74686D2E432E322E342E3131> \MusicPDFRef{17213} <74686D2E432E322E342E3132> \MusicPDFRef{17214}] /Limits [<74686D2E432E322E332E38> <74686D2E432E322E342E3132>] >>}
\MusicPDFWriteObject{10272}{<< /Names [<74686D2E432E322E342E3133> \MusicPDFRef{17215} <74686D2E432E322E342E3134> \MusicPDFRef{17216} <74686D2E432E322E342E3135> \MusicPDFRef{17217} <74686D2E432E322E342E3136> \MusicPDFRef{17218} <74686D2E432E322E342E3137> \MusicPDFRef{17219} <74686D2E432E322E342E3138> \MusicPDFRef{17220}] /Limits [<74686D2E432E322E342E3133> <74686D2E432E322E342E3138>] >>}
\MusicPDFWriteObject{10273}{<< /Names [<74686D2E432E322E342E3139> \MusicPDFRef{17221} <74686D2E432E322E342E32> \MusicPDFRef{17222} <74686D2E432E322E342E33> \MusicPDFRef{17223} <74686D2E432E322E342E34> \MusicPDFRef{17224} <74686D2E432E322E342E35> \MusicPDFRef{17225} <74686D2E432E322E342E36> \MusicPDFRef{17226}] /Limits [<74686D2E432E322E342E3139> <74686D2E432E322E342E36>] >>}
\MusicPDFWriteObject{10274}{<< /Names [<74686D2E432E322E342E37> \MusicPDFRef{17227} <74686D2E432E322E342E38> \MusicPDFRef{17228} <74686D2E432E322E342E39> \MusicPDFRef{17229} <74686D2E432E322E352E31> \MusicPDFRef{17230} <74686D2E432E322E352E32> \MusicPDFRef{17231} <74686D2E432E332E312E31> \MusicPDFRef{17232}] /Limits [<74686D2E432E322E342E37> <74686D2E432E332E312E31>] >>}
\MusicPDFWriteObject{10275}{<< /Names [<74686D2E432E332E312E32> \MusicPDFRef{17233} <74686D2E432E332E322E31> \MusicPDFRef{17234} <74686D2E432E332E322E32> \MusicPDFRef{17235} <74686D2E432E332E332E31> \MusicPDFRef{17236} <74686D2E432E332E342E31> \MusicPDFRef{17237} <74686D2E432E332E342E32> \MusicPDFRef{17238}] /Limits [<74686D2E432E332E312E32> <74686D2E432E332E342E32>] >>}
\MusicPDFWriteObject{10276}{<< /Names [<74686D2E432E332E342E33> \MusicPDFRef{17239} <74686D2E432E332E342E34> \MusicPDFRef{17240} <74686D2E432E332E352E31> \MusicPDFRef{17241} <74686D2E432E332E352E32> \MusicPDFRef{17242} <74686D2E432E332E352E33> \MusicPDFRef{17243} <74686D2E432E332E352E34> \MusicPDFRef{17244}] /Limits [<74686D2E432E332E342E33> <74686D2E432E332E352E34>] >>}
\MusicPDFWriteObject{10277}{<< /Names [<74686D2E432E332E352E35> \MusicPDFRef{17245} <74686D2E432E332E362E31> \MusicPDFRef{17246} <74686D2E432E332E362E32> \MusicPDFRef{17247} <74686D2E432E332E362E33> \MusicPDFRef{17248} <74686D2E432E332E362E34> \MusicPDFRef{17249} <74686D2E432E332E362E35> \MusicPDFRef{17250}] /Limits [<74686D2E432E332E352E35> <74686D2E432E332E362E35>] >>}
\MusicPDFWriteObject{10278}{<< /Names [<74686D2E432E332E362E36> \MusicPDFRef{17251} <74686D2E432E332E372E31> \MusicPDFRef{17252} <74686D2E432E342E312E31> \MusicPDFRef{17253} <74686D2E432E342E312E32> \MusicPDFRef{17254} <74686D2E432E342E312E33> \MusicPDFRef{17255} <74686D2E432E342E322E31> \MusicPDFRef{17256}] /Limits [<74686D2E432E332E362E36> <74686D2E432E342E322E31>] >>}
\MusicPDFWriteObject{10279}{<< /Names [<74686D2E432E342E322E32> \MusicPDFRef{17257} <74686D2E432E342E332E31> \MusicPDFRef{17258} <74686D2E432E342E332E3130> \MusicPDFRef{17259} <74686D2E432E342E332E3131> \MusicPDFRef{17260} <74686D2E432E342E332E3132> \MusicPDFRef{17261} <74686D2E432E342E332E3133> \MusicPDFRef{17262}] /Limits [<74686D2E432E342E322E32> <74686D2E432E342E332E3133>] >>}
\MusicPDFWriteObject{10280}{<< /Names [<74686D2E432E342E332E32> \MusicPDFRef{17263} <74686D2E432E342E332E33> \MusicPDFRef{17264} <74686D2E432E342E332E34> \MusicPDFRef{17265} <74686D2E432E342E332E35> \MusicPDFRef{17266} <74686D2E432E342E332E36> \MusicPDFRef{17267} <74686D2E432E342E332E37> \MusicPDFRef{17268}] /Limits [<74686D2E432E342E332E32> <74686D2E432E342E332E37>] >>}
\MusicPDFWriteObject{10281}{<< /Names [<74686D2E432E342E332E38> \MusicPDFRef{17269} <74686D2E432E342E332E39> \MusicPDFRef{17270} <74686D2E432E342E342E31> \MusicPDFRef{17271} <74686D2E432E342E342E32> \MusicPDFRef{17272} <74686D2E432E342E342E33> \MusicPDFRef{17273} <74686D2E432E342E342E34> \MusicPDFRef{17274}] /Limits [<74686D2E432E342E332E38> <74686D2E432E342E342E34>] >>}
\MusicPDFWriteObject{10282}{<< /Names [<74686D2E432E342E342E35> \MusicPDFRef{17275} <74686D2E432E352E312E31> \MusicPDFRef{17276} <74686D2E432E352E312E3130> \MusicPDFRef{17277} <74686D2E432E352E312E3131> \MusicPDFRef{17278} <74686D2E432E352E312E3132> \MusicPDFRef{17279} <74686D2E432E352E312E3133> \MusicPDFRef{17280}] /Limits [<74686D2E432E342E342E35> <74686D2E432E352E312E3133>] >>}
\MusicPDFWriteObject{10283}{<< /Names [<74686D2E432E352E312E3134> \MusicPDFRef{17281} <74686D2E432E352E312E3135> \MusicPDFRef{17282} <74686D2E432E352E312E3136> \MusicPDFRef{17283} <74686D2E432E352E312E3137> \MusicPDFRef{17284} <74686D2E432E352E312E3138> \MusicPDFRef{17285} <74686D2E432E352E312E3139> \MusicPDFRef{17286}] /Limits [<74686D2E432E352E312E3134> <74686D2E432E352E312E3139>] >>}
\MusicPDFWriteObject{10284}{<< /Names [<74686D2E432E352E312E32> \MusicPDFRef{17287} <74686D2E432E352E312E3230> \MusicPDFRef{17288} <74686D2E432E352E312E3231> \MusicPDFRef{17289} <74686D2E432E352E312E3232> \MusicPDFRef{17290} <74686D2E432E352E312E3233> \MusicPDFRef{17291} <74686D2E432E352E312E3234> \MusicPDFRef{17292}] /Limits [<74686D2E432E352E312E32> <74686D2E432E352E312E3234>] >>}
\MusicPDFWriteObject{10285}{<< /Names [<74686D2E432E352E312E3235> \MusicPDFRef{17293} <74686D2E432E352E312E3236> \MusicPDFRef{17294} <74686D2E432E352E312E3237> \MusicPDFRef{17295} <74686D2E432E352E312E3238> \MusicPDFRef{17296} <74686D2E432E352E312E3239> \MusicPDFRef{17297} <74686D2E432E352E312E33> \MusicPDFRef{17298}] /Limits [<74686D2E432E352E312E3235> <74686D2E432E352E312E33>] >>}
\MusicPDFWriteObject{10286}{<< /Names [<74686D2E432E352E312E3330> \MusicPDFRef{17299} <74686D2E432E352E312E34> \MusicPDFRef{17300} <74686D2E432E352E312E35> \MusicPDFRef{17301} <74686D2E432E352E312E36> \MusicPDFRef{17302} <74686D2E432E352E312E37> \MusicPDFRef{17303} <74686D2E432E352E312E38> \MusicPDFRef{17304}] /Limits [<74686D2E432E352E312E3330> <74686D2E432E352E312E38>] >>}
\MusicPDFWriteObject{10287}{<< /Names [<74686D2E432E352E312E39> \MusicPDFRef{17305} <74686D2E432E352E322E31> \MusicPDFRef{17306} <74686D2E432E352E322E32> \MusicPDFRef{17307} <74686D2E432E352E322E33> \MusicPDFRef{17308} <74686D2E432E352E322E34> \MusicPDFRef{17309} <74686D2E432E352E322E35> \MusicPDFRef{17310}] /Limits [<74686D2E432E352E312E39> <74686D2E432E352E322E35>] >>}
\MusicPDFWriteObject{10288}{<< /Names [<74686D2E432E352E322E36> \MusicPDFRef{17311} <74686D2E432E352E322E37> \MusicPDFRef{17312} <74686D2E432E352E332E31> \MusicPDFRef{17313} <74686D2E432E352E332E32> \MusicPDFRef{17314} <74686D2E432E352E332E33> \MusicPDFRef{17315} <74686D2E432E352E342E31> \MusicPDFRef{17316}] /Limits [<74686D2E432E352E322E36> <74686D2E432E352E342E31>] >>}
\MusicPDFWriteObject{10289}{<< /Names [<74686D2E432E352E342E32> \MusicPDFRef{17317} <74686D2E432E352E342E33> \MusicPDFRef{17318} <74686D2E442E312E312E31> \MusicPDFRef{17319} <74686D2E442E312E312E32> \MusicPDFRef{17320} <74686D2E442E312E31302E31> \MusicPDFRef{17321} <74686D2E442E312E31302E3130> \MusicPDFRef{17322}] /Limits [<74686D2E432E352E342E32> <74686D2E442E312E31302E3130>] >>}
\MusicPDFWriteObject{10290}{<< /Names [<74686D2E442E312E31302E32> \MusicPDFRef{17323} <74686D2E442E312E31302E33> \MusicPDFRef{17324} <74686D2E442E312E31302E34> \MusicPDFRef{17325} <74686D2E442E312E31302E35> \MusicPDFRef{17326} <74686D2E442E312E31302E36> \MusicPDFRef{17327} <74686D2E442E312E31302E37> \MusicPDFRef{17328}] /Limits [<74686D2E442E312E31302E32> <74686D2E442E312E31302E37>] >>}
\MusicPDFWriteObject{10291}{<< /Names [<74686D2E442E312E31302E38> \MusicPDFRef{17329} <74686D2E442E312E31302E39> \MusicPDFRef{17330} <74686D2E442E312E322E31> \MusicPDFRef{17331} <74686D2E442E312E322E32> \MusicPDFRef{17332} <74686D2E442E312E322E33> \MusicPDFRef{17333} <74686D2E442E312E322E34> \MusicPDFRef{17334}] /Limits [<74686D2E442E312E31302E38> <74686D2E442E312E322E34>] >>}
\MusicPDFWriteObject{10292}{<< /Names [<74686D2E442E312E322E35> \MusicPDFRef{17335} <74686D2E442E312E322E36> \MusicPDFRef{17336} <74686D2E442E312E322E37> \MusicPDFRef{17337} <74686D2E442E312E322E38> \MusicPDFRef{17338} <74686D2E442E312E332E31> \MusicPDFRef{17339} <74686D2E442E312E332E32> \MusicPDFRef{17340}] /Limits [<74686D2E442E312E322E35> <74686D2E442E312E332E32>] >>}
\MusicPDFWriteObject{10293}{<< /Names [<74686D2E442E312E332E33> \MusicPDFRef{17341} <74686D2E442E312E332E34> \MusicPDFRef{17342} <74686D2E442E312E332E35> \MusicPDFRef{17343} <74686D2E442E312E332E36> \MusicPDFRef{17344} <74686D2E442E312E332E37> \MusicPDFRef{17345} <74686D2E442E312E332E38> \MusicPDFRef{17346}] /Limits [<74686D2E442E312E332E33> <74686D2E442E312E332E38>] >>}
\MusicPDFWriteObject{10294}{<< /Names [<74686D2E442E312E342E31> \MusicPDFRef{17347} <74686D2E442E312E342E32> \MusicPDFRef{17348} <74686D2E442E312E342E33> \MusicPDFRef{17349} <74686D2E442E312E342E34> \MusicPDFRef{17350} <74686D2E442E312E342E35> \MusicPDFRef{17351} <74686D2E442E312E342E36> \MusicPDFRef{17352}] /Limits [<74686D2E442E312E342E31> <74686D2E442E312E342E36>] >>}
\MusicPDFWriteObject{10295}{<< /Names [<74686D2E442E312E342E37> \MusicPDFRef{17353} <74686D2E442E312E352E31> \MusicPDFRef{17354} <74686D2E442E312E352E3130> \MusicPDFRef{17355} <74686D2E442E312E352E32> \MusicPDFRef{17356} <74686D2E442E312E352E33> \MusicPDFRef{17357} <74686D2E442E312E352E34> \MusicPDFRef{17358}] /Limits [<74686D2E442E312E342E37> <74686D2E442E312E352E34>] >>}
\MusicPDFWriteObject{10296}{<< /Names [<74686D2E442E312E352E35> \MusicPDFRef{17359} <74686D2E442E312E352E36> \MusicPDFRef{17360} <74686D2E442E312E352E37> \MusicPDFRef{17361} <74686D2E442E312E352E38> \MusicPDFRef{17362} <74686D2E442E312E352E39> \MusicPDFRef{17363} <74686D2E442E312E362E31> \MusicPDFRef{17364}] /Limits [<74686D2E442E312E352E35> <74686D2E442E312E362E31>] >>}
\MusicPDFWriteObject{10297}{<< /Names [<74686D2E442E312E362E3130> \MusicPDFRef{17365} <74686D2E442E312E362E3131> \MusicPDFRef{17366} <74686D2E442E312E362E3132> \MusicPDFRef{17367} <74686D2E442E312E362E3133> \MusicPDFRef{17368} <74686D2E442E312E362E3134> \MusicPDFRef{17369} <74686D2E442E312E362E3135> \MusicPDFRef{17370}] /Limits [<74686D2E442E312E362E3130> <74686D2E442E312E362E3135>] >>}
\MusicPDFWriteObject{10298}{<< /Names [<74686D2E442E312E362E3136> \MusicPDFRef{17371} <74686D2E442E312E362E3137> \MusicPDFRef{17372} <74686D2E442E312E362E3138> \MusicPDFRef{17373} <74686D2E442E312E362E3139> \MusicPDFRef{17374} <74686D2E442E312E362E32> \MusicPDFRef{17375} <74686D2E442E312E362E3230> \MusicPDFRef{17376}] /Limits [<74686D2E442E312E362E3136> <74686D2E442E312E362E3230>] >>}
\MusicPDFWriteObject{10299}{<< /Names [<74686D2E442E312E362E33> \MusicPDFRef{17377} <74686D2E442E312E362E34> \MusicPDFRef{17378} <74686D2E442E312E362E35> \MusicPDFRef{17379} <74686D2E442E312E362E36> \MusicPDFRef{17380} <74686D2E442E312E362E37> \MusicPDFRef{17381} <74686D2E442E312E362E38> \MusicPDFRef{17382}] /Limits [<74686D2E442E312E362E33> <74686D2E442E312E362E38>] >>}
\MusicPDFWriteObject{10300}{<< /Names [<74686D2E442E312E362E39> \MusicPDFRef{17383} <74686D2E442E312E372E31> \MusicPDFRef{17384} <74686D2E442E312E382E31> \MusicPDFRef{17385} <74686D2E442E312E382E32> \MusicPDFRef{17386} <74686D2E442E312E382E33> \MusicPDFRef{17387} <74686D2E442E312E382E34> \MusicPDFRef{17388}] /Limits [<74686D2E442E312E362E39> <74686D2E442E312E382E34>] >>}
\MusicPDFWriteObject{10301}{<< /Names [<74686D2E442E312E392E31> \MusicPDFRef{17389} <74686D2E442E312E392E32> \MusicPDFRef{17390} <74686D2E442E312E392E33> \MusicPDFRef{17391} <74686D2E442E312E392E34> \MusicPDFRef{17392} <74686D2E442E322E312E31> \MusicPDFRef{17393} <74686D2E442E322E312E32> \MusicPDFRef{17394}] /Limits [<74686D2E442E312E392E31> <74686D2E442E322E312E32>] >>}
\MusicPDFWriteObject{10302}{<< /Names [<74686D2E442E322E312E33> \MusicPDFRef{17395} <74686D2E442E322E312E34> \MusicPDFRef{17396} <74686D2E442E322E322E31> \MusicPDFRef{17397} <74686D2E442E322E322E32> \MusicPDFRef{17398} <74686D2E442E322E322E33> \MusicPDFRef{17399} <74686D2E442E332E312E31> \MusicPDFRef{17400}] /Limits [<74686D2E442E322E312E33> <74686D2E442E332E312E31>] >>}
\MusicPDFWriteObject{10303}{<< /Names [<74686D2E442E332E312E32> \MusicPDFRef{17401} <74686D2E442E332E312E33> \MusicPDFRef{17402} <74686D2E442E332E312E34> \MusicPDFRef{17403} <74686D2E442E332E322E31> \MusicPDFRef{17404} <74686D2E442E332E322E32> \MusicPDFRef{17405} <74686D2E442E332E322E33> \MusicPDFRef{17406}] /Limits [<74686D2E442E332E312E32> <74686D2E442E332E322E33>] >>}
\MusicPDFWriteObject{10304}{<< /Names [<74686D2E442E332E322E34> \MusicPDFRef{17407} <74686D2E442E332E322E35> \MusicPDFRef{17408} <74686D2E442E332E322E36> \MusicPDFRef{17409} <74686D2E442E332E322E37> \MusicPDFRef{17410} <74686D2E442E332E332E31> \MusicPDFRef{17411} <74686D2E442E332E332E3130> \MusicPDFRef{17412}] /Limits [<74686D2E442E332E322E34> <74686D2E442E332E332E3130>] >>}
\MusicPDFWriteObject{10305}{<< /Names [<74686D2E442E332E332E3131> \MusicPDFRef{17413} <74686D2E442E332E332E3132> \MusicPDFRef{17414} <74686D2E442E332E332E3133> \MusicPDFRef{17415} <74686D2E442E332E332E32> \MusicPDFRef{17416} <74686D2E442E332E332E33> \MusicPDFRef{17417} <74686D2E442E332E332E34> \MusicPDFRef{17418}] /Limits [<74686D2E442E332E332E3131> <74686D2E442E332E332E34>] >>}
\MusicPDFWriteObject{10306}{<< /Names [<74686D2E442E332E332E35> \MusicPDFRef{17419} <74686D2E442E332E332E36> \MusicPDFRef{17420} <74686D2E442E332E332E37> \MusicPDFRef{17421} <74686D2E442E332E332E38> \MusicPDFRef{17422} <74686D2E442E332E332E39> \MusicPDFRef{17423} <74686D2E442E332E342E31> \MusicPDFRef{17424}] /Limits [<74686D2E442E332E332E35> <74686D2E442E332E342E31>] >>}
\MusicPDFWriteObject{10307}{<< /Names [<74686D2E442E332E342E32> \MusicPDFRef{17425} <74686D2E442E332E352E31> \MusicPDFRef{17426} <74686D2E442E332E352E32> \MusicPDFRef{17427} <74686D2E442E332E362E31> \MusicPDFRef{17428} <74686D2E452E312E312E31> \MusicPDFRef{17429} <74686D2E452E312E312E32> \MusicPDFRef{17430}] /Limits [<74686D2E442E332E342E32> <74686D2E452E312E312E32>] >>}
\MusicPDFWriteObject{10308}{<< /Names [<74686D2E452E312E312E33> \MusicPDFRef{17431} <74686D2E452E312E312E34> \MusicPDFRef{17432} <74686D2E452E312E312E35> \MusicPDFRef{17433} <74686D2E452E312E312E36> \MusicPDFRef{17434} <74686D2E452E312E312E37> \MusicPDFRef{17435} <74686D2E452E312E312E38> \MusicPDFRef{17436}] /Limits [<74686D2E452E312E312E33> <74686D2E452E312E312E38>] >>}
\MusicPDFWriteObject{10309}{<< /Names [<74686D2E452E312E312E39> \MusicPDFRef{17437} <74686D2E452E312E322E31> \MusicPDFRef{17438} <74686D2E452E312E322E3130> \MusicPDFRef{17439} <74686D2E452E312E322E3131> \MusicPDFRef{17440} <74686D2E452E312E322E3132> \MusicPDFRef{17441} <74686D2E452E312E322E3133> \MusicPDFRef{17442}] /Limits [<74686D2E452E312E312E39> <74686D2E452E312E322E3133>] >>}
\MusicPDFWriteObject{10310}{<< /Names [<74686D2E452E312E322E3134> \MusicPDFRef{17443} <74686D2E452E312E322E3135> \MusicPDFRef{17444} <74686D2E452E312E322E3136> \MusicPDFRef{17445} <74686D2E452E312E322E3137> \MusicPDFRef{17446} <74686D2E452E312E322E3138> \MusicPDFRef{17447} <74686D2E452E312E322E3139> \MusicPDFRef{17448}] /Limits [<74686D2E452E312E322E3134> <74686D2E452E312E322E3139>] >>}
\MusicPDFWriteObject{10311}{<< /Names [<74686D2E452E312E322E32> \MusicPDFRef{17449} <74686D2E452E312E322E3230> \MusicPDFRef{17450} <74686D2E452E312E322E3231> \MusicPDFRef{17451} <74686D2E452E312E322E3232> \MusicPDFRef{17452} <74686D2E452E312E322E3233> \MusicPDFRef{17453} <74686D2E452E312E322E3234> \MusicPDFRef{17454}] /Limits [<74686D2E452E312E322E32> <74686D2E452E312E322E3234>] >>}
\MusicPDFWriteObject{10312}{<< /Names [<74686D2E452E312E322E3235> \MusicPDFRef{17455} <74686D2E452E312E322E3236> \MusicPDFRef{17456} <74686D2E452E312E322E3237> \MusicPDFRef{17457} <74686D2E452E312E322E3238> \MusicPDFRef{17458} <74686D2E452E312E322E33> \MusicPDFRef{17459} <74686D2E452E312E322E34> \MusicPDFRef{17460}] /Limits [<74686D2E452E312E322E3235> <74686D2E452E312E322E34>] >>}
\MusicPDFWriteObject{10313}{<< /Names [<74686D2E452E312E322E35> \MusicPDFRef{17461} <74686D2E452E312E322E36> \MusicPDFRef{17462} <74686D2E452E312E322E37> \MusicPDFRef{17463} <74686D2E452E312E322E38> \MusicPDFRef{17464} <74686D2E452E312E322E39> \MusicPDFRef{17465} <74686D2E452E322E312E31> \MusicPDFRef{17466}] /Limits [<74686D2E452E312E322E35> <74686D2E452E322E312E31>] >>}
\MusicPDFWriteObject{10314}{<< /Names [<74686D2E452E322E312E3130> \MusicPDFRef{17467} <74686D2E452E322E312E3131> \MusicPDFRef{17468} <74686D2E452E322E312E3132> \MusicPDFRef{17469} <74686D2E452E322E312E3133> \MusicPDFRef{17470} <74686D2E452E322E312E3134> \MusicPDFRef{17471} <74686D2E452E322E312E3135> \MusicPDFRef{17472}] /Limits [<74686D2E452E322E312E3130> <74686D2E452E322E312E3135>] >>}
\MusicPDFWriteObject{10315}{<< /Names [<74686D2E452E322E312E3136> \MusicPDFRef{17473} <74686D2E452E322E312E3137> \MusicPDFRef{17474} <74686D2E452E322E312E32> \MusicPDFRef{17475} <74686D2E452E322E312E33> \MusicPDFRef{17476} <74686D2E452E322E312E34> \MusicPDFRef{17477} <74686D2E452E322E312E35> \MusicPDFRef{17478}] /Limits [<74686D2E452E322E312E3136> <74686D2E452E322E312E35>] >>}
\MusicPDFWriteObject{10316}{<< /Names [<74686D2E452E322E312E36> \MusicPDFRef{17479} <74686D2E452E322E312E37> \MusicPDFRef{17480} <74686D2E452E322E312E38> \MusicPDFRef{17481} <74686D2E452E322E312E39> \MusicPDFRef{17482} <74686D2E452E322E322E31> \MusicPDFRef{17483} <74686D2E452E322E322E3130> \MusicPDFRef{17484}] /Limits [<74686D2E452E322E312E36> <74686D2E452E322E322E3130>] >>}
\MusicPDFWriteObject{10317}{<< /Names [<74686D2E452E322E322E32> \MusicPDFRef{17485} <74686D2E452E322E322E33> \MusicPDFRef{17486} <74686D2E452E322E322E34> \MusicPDFRef{17487} <74686D2E452E322E322E35> \MusicPDFRef{17488} <74686D2E452E322E322E36> \MusicPDFRef{17489} <74686D2E452E322E322E37> \MusicPDFRef{17490}] /Limits [<74686D2E452E322E322E32> <74686D2E452E322E322E37>] >>}
\MusicPDFWriteObject{10318}{<< /Names [<74686D2E452E322E322E38> \MusicPDFRef{17491} <74686D2E452E322E322E39> \MusicPDFRef{17492} <74686D2E452E322E332E31> \MusicPDFRef{17493} <74686D2E452E322E332E3130> \MusicPDFRef{17494} <74686D2E452E322E332E3131> \MusicPDFRef{17495} <74686D2E452E322E332E3132> \MusicPDFRef{17496}] /Limits [<74686D2E452E322E322E38> <74686D2E452E322E332E3132>] >>}
\MusicPDFWriteObject{10319}{<< /Names [<74686D2E452E322E332E3133> \MusicPDFRef{17497} <74686D2E452E322E332E3134> \MusicPDFRef{17498} <74686D2E452E322E332E3135> \MusicPDFRef{17499} <74686D2E452E322E332E3136> \MusicPDFRef{17500} <74686D2E452E322E332E3137> \MusicPDFRef{17501} <74686D2E452E322E332E32> \MusicPDFRef{17502}] /Limits [<74686D2E452E322E332E3133> <74686D2E452E322E332E32>] >>}
\MusicPDFWriteObject{10320}{<< /Names [<74686D2E452E322E332E33> \MusicPDFRef{17503} <74686D2E452E322E332E34> \MusicPDFRef{17504} <74686D2E452E322E332E35> \MusicPDFRef{17505} <74686D2E452E322E332E36> \MusicPDFRef{17506} <74686D2E452E322E332E37> \MusicPDFRef{17507} <74686D2E452E322E332E38> \MusicPDFRef{17508}] /Limits [<74686D2E452E322E332E33> <74686D2E452E322E332E38>] >>}
\MusicPDFWriteObject{10321}{<< /Names [<74686D2E452E322E332E39> \MusicPDFRef{17509} <74686D2E452E322E342E31> \MusicPDFRef{17510} <74686D2E452E322E342E3130> \MusicPDFRef{17511} <74686D2E452E322E342E3131> \MusicPDFRef{17512} <74686D2E452E322E342E3132> \MusicPDFRef{17513} <74686D2E452E322E342E32> \MusicPDFRef{17514}] /Limits [<74686D2E452E322E332E39> <74686D2E452E322E342E32>] >>}
\MusicPDFWriteObject{10322}{<< /Names [<74686D2E452E322E342E33> \MusicPDFRef{17515} <74686D2E452E322E342E34> \MusicPDFRef{17516} <74686D2E452E322E342E35> \MusicPDFRef{17517} <74686D2E452E322E342E36> \MusicPDFRef{17518} <74686D2E452E322E342E37> \MusicPDFRef{17519} <74686D2E452E322E342E38> \MusicPDFRef{17520}] /Limits [<74686D2E452E322E342E33> <74686D2E452E322E342E38>] >>}
\MusicPDFWriteObject{10323}{<< /Names [<74686D2E452E322E342E39> \MusicPDFRef{17521} <74686D2E452E332E312E31> \MusicPDFRef{17522} <74686D2E452E332E312E32> \MusicPDFRef{17523} <74686D2E452E332E312E33> \MusicPDFRef{17524} <74686D2E452E332E312E34> \MusicPDFRef{17525} <74686D2E452E332E312E35> \MusicPDFRef{17526}] /Limits [<74686D2E452E322E342E39> <74686D2E452E332E312E35>] >>}
\MusicPDFWriteObject{10324}{<< /Names [<74686D2E452E332E312E36> \MusicPDFRef{17527} <74686D2E452E332E312E37> \MusicPDFRef{17528} <74686D2E452E332E322E31> \MusicPDFRef{17529} <74686D2E452E332E322E32> \MusicPDFRef{17530} <74686D2E452E332E322E33> \MusicPDFRef{17531} <74686D2E452E332E332E31> \MusicPDFRef{17532}] /Limits [<74686D2E452E332E312E36> <74686D2E452E332E332E31>] >>}
\MusicPDFWriteObject{10325}{<< /Names [<74686D2E452E332E332E32> \MusicPDFRef{17533} <74686D2E452E332E332E33> \MusicPDFRef{17534} <74686D2E452E332E332E34> \MusicPDFRef{17535} <74686D2E452E332E342E31> \MusicPDFRef{17536} <74686D2E452E332E342E32> \MusicPDFRef{17537} <74686D2E452E332E342E33> \MusicPDFRef{17538}] /Limits [<74686D2E452E332E332E32> <74686D2E452E332E342E33>] >>}
\MusicPDFWriteObject{10326}{<< /Names [<74686D2E452E332E342E34> \MusicPDFRef{17539} <74686D2E452E332E352E31> \MusicPDFRef{17540} <74686D2E452E332E352E32> \MusicPDFRef{17541} <74686D2E452E332E352E33> \MusicPDFRef{17542} <74686D2E452E332E352E34> \MusicPDFRef{17543} <74686D2E452E332E352E35> \MusicPDFRef{17544}] /Limits [<74686D2E452E332E342E34> <74686D2E452E332E352E35>] >>}
\MusicPDFWriteObject{10327}{<< /Names [<74686D2E452E332E352E36> \MusicPDFRef{17545} <74686D2E452E332E352E37> \MusicPDFRef{17546} <74686D2E452E332E352E38> \MusicPDFRef{17547} <74686D2E452E332E352E39> \MusicPDFRef{17548} <74686D2E452E342E312E31> \MusicPDFRef{17549} <74686D2E452E342E312E32> \MusicPDFRef{17550}] /Limits [<74686D2E452E332E352E36> <74686D2E452E342E312E32>] >>}
\MusicPDFWriteObject{10328}{<< /Names [<74686D2E452E342E312E33> \MusicPDFRef{17551} <74686D2E452E342E312E34> \MusicPDFRef{17552} <74686D2E452E342E322E31> \MusicPDFRef{17553} <74686D2E452E342E322E3130> \MusicPDFRef{17554} <74686D2E452E342E322E3131> \MusicPDFRef{17555} <74686D2E452E342E322E3132> \MusicPDFRef{17556}] /Limits [<74686D2E452E342E312E33> <74686D2E452E342E322E3132>] >>}
\MusicPDFWriteObject{10329}{<< /Names [<74686D2E452E342E322E32> \MusicPDFRef{17557} <74686D2E452E342E322E33> \MusicPDFRef{17558} <74686D2E452E342E322E34> \MusicPDFRef{17559} <74686D2E452E342E322E35> \MusicPDFRef{17560} <74686D2E452E342E322E36> \MusicPDFRef{17561} <74686D2E452E342E322E37> \MusicPDFRef{17562}] /Limits [<74686D2E452E342E322E32> <74686D2E452E342E322E37>] >>}
\MusicPDFWriteObject{10330}{<< /Names [<74686D2E452E342E322E38> \MusicPDFRef{17563} <74686D2E452E342E322E39> \MusicPDFRef{17564} <74686D2E452E352E312E31> \MusicPDFRef{17565} <74686D2E452E352E312E3130> \MusicPDFRef{17566} <74686D2E452E352E312E3131> \MusicPDFRef{17567} <74686D2E452E352E312E3132> \MusicPDFRef{17568}] /Limits [<74686D2E452E342E322E38> <74686D2E452E352E312E3132>] >>}
\MusicPDFWriteObject{10331}{<< /Names [<74686D2E452E352E312E3133> \MusicPDFRef{17569} <74686D2E452E352E312E3134> \MusicPDFRef{17570} <74686D2E452E352E312E3135> \MusicPDFRef{17571} <74686D2E452E352E312E3136> \MusicPDFRef{17572} <74686D2E452E352E312E3137> \MusicPDFRef{17573} <74686D2E452E352E312E3138> \MusicPDFRef{17574}] /Limits [<74686D2E452E352E312E3133> <74686D2E452E352E312E3138>] >>}
\MusicPDFWriteObject{10332}{<< /Names [<74686D2E452E352E312E3139> \MusicPDFRef{17575} <74686D2E452E352E312E32> \MusicPDFRef{17576} <74686D2E452E352E312E3230> \MusicPDFRef{17577} <74686D2E452E352E312E33> \MusicPDFRef{17578} <74686D2E452E352E312E34> \MusicPDFRef{17579} <74686D2E452E352E312E35> \MusicPDFRef{17580}] /Limits [<74686D2E452E352E312E3139> <74686D2E452E352E312E35>] >>}
\MusicPDFWriteObject{10333}{<< /Names [<74686D2E452E352E312E36> \MusicPDFRef{17581} <74686D2E452E352E312E37> \MusicPDFRef{17582} <74686D2E452E352E312E38> \MusicPDFRef{17583} <74686D2E452E352E312E39> \MusicPDFRef{17584} <74686D2E452E352E322E31> \MusicPDFRef{17585} <74686D2E452E352E322E32> \MusicPDFRef{17586}] /Limits [<74686D2E452E352E312E36> <74686D2E452E352E322E32>] >>}
\MusicPDFWriteObject{10334}{<< /Names [<74686D2E452E352E322E33> \MusicPDFRef{17587} <74686D2E452E352E322E34> \MusicPDFRef{17588} <74686D2E452E352E322E35> \MusicPDFRef{17589} <74686D2E452E352E322E36> \MusicPDFRef{17590} <74686D2E452E352E322E37> \MusicPDFRef{17591} <74686D2E452E352E332E31> \MusicPDFRef{17592}] /Limits [<74686D2E452E352E322E33> <74686D2E452E352E332E31>] >>}
\MusicPDFWriteObject{10335}{<< /Names [<74686D2E452E352E332E32> \MusicPDFRef{17593} <74686D2E452E352E332E33> \MusicPDFRef{17594} <74686D2E452E352E332E34> \MusicPDFRef{17595} <74686D2E452E352E332E35> \MusicPDFRef{17596} <74686D2E452E352E332E36> \MusicPDFRef{17597} <74686D2E452E352E332E37> \MusicPDFRef{17598}] /Limits [<74686D2E452E352E332E32> <74686D2E452E352E332E37>] >>}
\MusicPDFWriteObject{10336}{<< /Names [<74686D2E452E352E342E31> \MusicPDFRef{17599} <74686D2E452E352E342E3130> \MusicPDFRef{17600} <74686D2E452E352E342E32> \MusicPDFRef{17601} <74686D2E452E352E342E33> \MusicPDFRef{17602} <74686D2E452E352E342E34> \MusicPDFRef{17603} <74686D2E452E352E342E35> \MusicPDFRef{17604}] /Limits [<74686D2E452E352E342E31> <74686D2E452E352E342E35>] >>}
\MusicPDFWriteObject{10337}{<< /Names [<74686D2E452E352E342E36> \MusicPDFRef{17605} <74686D2E452E352E342E37> \MusicPDFRef{17606} <74686D2E452E352E342E38> \MusicPDFRef{17607} <74686D2E452E352E342E39> \MusicPDFRef{17608} <74686D2E452E352E352E31> \MusicPDFRef{17609} <74686D2E452E352E352E32> \MusicPDFRef{17610}] /Limits [<74686D2E452E352E342E36> <74686D2E452E352E352E32>] >>}
\MusicPDFWriteObject{10338}{<< /Names [<74686D2E452E352E352E33> \MusicPDFRef{17611} <74686D2E452E352E352E34> \MusicPDFRef{17612} <74686D2E452E352E362E31> \MusicPDFRef{17613} <74686D2E452E352E362E32> \MusicPDFRef{17614} <74686D2E452E352E372E31> \MusicPDFRef{17615} <74686D2E452E352E372E3130> \MusicPDFRef{17616}] /Limits [<74686D2E452E352E352E33> <74686D2E452E352E372E3130>] >>}
\MusicPDFWriteObject{10339}{<< /Names [<74686D2E452E352E372E3131> \MusicPDFRef{17617} <74686D2E452E352E372E3132> \MusicPDFRef{17618} <74686D2E452E352E372E3133> \MusicPDFRef{17619} <74686D2E452E352E372E3134> \MusicPDFRef{17620} <74686D2E452E352E372E3135> \MusicPDFRef{17621} <74686D2E452E352E372E3136> \MusicPDFRef{17622}] /Limits [<74686D2E452E352E372E3131> <74686D2E452E352E372E3136>] >>}
\MusicPDFWriteObject{10340}{<< /Names [<74686D2E452E352E372E3137> \MusicPDFRef{17623} <74686D2E452E352E372E3138> \MusicPDFRef{17624} <74686D2E452E352E372E3139> \MusicPDFRef{17625} <74686D2E452E352E372E32> \MusicPDFRef{17626} <74686D2E452E352E372E3230> \MusicPDFRef{17627} <74686D2E452E352E372E3231> \MusicPDFRef{17628}] /Limits [<74686D2E452E352E372E3137> <74686D2E452E352E372E3231>] >>}
\MusicPDFWriteObject{10341}{<< /Names [<74686D2E452E352E372E3232> \MusicPDFRef{17629} <74686D2E452E352E372E3233> \MusicPDFRef{17630} <74686D2E452E352E372E3234> \MusicPDFRef{17631} <74686D2E452E352E372E3235> \MusicPDFRef{17632} <74686D2E452E352E372E3236> \MusicPDFRef{17633} <74686D2E452E352E372E33> \MusicPDFRef{17634}] /Limits [<74686D2E452E352E372E3232> <74686D2E452E352E372E33>] >>}
\MusicPDFWriteObject{10342}{<< /Names [<74686D2E452E352E372E34> \MusicPDFRef{17635} <74686D2E452E352E372E35> \MusicPDFRef{17636} <74686D2E452E352E372E36> \MusicPDFRef{17637} <74686D2E452E352E372E37> \MusicPDFRef{17638} <74686D2E452E352E372E38> \MusicPDFRef{17639} <74686D2E452E352E372E39> \MusicPDFRef{17640}] /Limits [<74686D2E452E352E372E34> <74686D2E452E352E372E39>] >>}
\MusicPDFWriteObject{10343}{<< /Names [<74686D2E452E362E312E31> \MusicPDFRef{17641} <74686D2E452E362E312E32> \MusicPDFRef{17642} <74686D2E452E362E312E33> \MusicPDFRef{17643} <74686D2E452E362E312E34> \MusicPDFRef{17644} <74686D2E452E362E312E35> \MusicPDFRef{17645} <74686D2E452E362E312E36> \MusicPDFRef{17646}] /Limits [<74686D2E452E362E312E31> <74686D2E452E362E312E36>] >>}
\MusicPDFWriteObject{10344}{<< /Names [<74686D2E452E362E312E37> \MusicPDFRef{17647} <74686D2E452E362E322E31> \MusicPDFRef{17648} <74686D2E452E362E322E32> \MusicPDFRef{17649} <74686D2E452E362E322E33> \MusicPDFRef{17650} <74686D2E452E362E322E34> \MusicPDFRef{17651} <74686D2E452E362E322E35> \MusicPDFRef{17652}] /Limits [<74686D2E452E362E312E37> <74686D2E452E362E322E35>] >>}
\MusicPDFWriteObject{10345}{<< /Names [<74686D2E452E362E322E36> \MusicPDFRef{17653} <74686D2E452E362E332E31> \MusicPDFRef{17654} <74686D2E452E362E332E3130> \MusicPDFRef{17655} <74686D2E452E362E332E3131> \MusicPDFRef{17656} <74686D2E452E362E332E3132> \MusicPDFRef{17657} <74686D2E452E362E332E32> \MusicPDFRef{17658}] /Limits [<74686D2E452E362E322E36> <74686D2E452E362E332E32>] >>}
\MusicPDFWriteObject{10346}{<< /Names [<74686D2E452E362E332E33> \MusicPDFRef{17659} <74686D2E452E362E332E34> \MusicPDFRef{17660} <74686D2E452E362E332E35> \MusicPDFRef{17661} <74686D2E452E362E332E36> \MusicPDFRef{17662} <74686D2E452E362E332E37> \MusicPDFRef{17663} <74686D2E452E362E332E38> \MusicPDFRef{17664}] /Limits [<74686D2E452E362E332E33> <74686D2E452E362E332E38>] >>}
\MusicPDFWriteObject{10347}{<< /Names [<74686D2E452E362E332E39> \MusicPDFRef{17665} <74686D2E452E362E342E31> \MusicPDFRef{17666} <74686D2E452E362E342E32> \MusicPDFRef{17667} <74686D2E452E362E342E33> \MusicPDFRef{17668} <74686D2E452E362E342E34> \MusicPDFRef{17669} <74686D2E452E362E342E35> \MusicPDFRef{17670}] /Limits [<74686D2E452E362E332E39> <74686D2E452E362E342E35>] >>}
\MusicPDFWriteObject{10348}{<< /Names [<74686D2E452E362E342E36> \MusicPDFRef{17671} <74686D2E452E362E342E37> \MusicPDFRef{17672} <74686D2E452E362E342E38> \MusicPDFRef{17673}] /Limits [<74686D2E452E362E342E36> <74686D2E452E362E342E38>] >>}
\MusicPDFWriteObject{10349}{<FEFF004300680061007000740065007200200034002E00200050006500720069006F0064006900630020005300630068006100750064006500720020004200610073006500730020006F006E0020006D006F00640075006C0061007400650064002000670072006F007500700073>}
\MusicPDFWriteObject{10350}{<< /S /GoTo /D <636861707465722E34> >>}
\MusicPDFWriteObject{10351}{<< /Title \MusicPDFRef{17674} /A \MusicPDFRef{17675} /Parent \MusicPDFRef{1} /Prev \MusicPDFRef{9103} /Next \MusicPDFRef{17676} /First \MusicPDFRef{17677} /Last \MusicPDFRef{17678} /Count -5 >>}
\MusicPDFWriteObject{10352}{<< /Title \MusicPDFRef{17679} /A \MusicPDFRef{17680} /Parent \MusicPDFRef{9103} /Next \MusicPDFRef{17681} >>}
\MusicPDFWriteObject{10353}{<< /Title \MusicPDFRef{17682} /A \MusicPDFRef{17683} /Parent \MusicPDFRef{9103} /Prev \MusicPDFRef{17684} >>}
\MusicPDFWriteObject{10354}{<FEFF0033002E0031002E0020005200690065006D0061006E006E0020007300750072006600610063006500730020006F006600200063006F006E00740069006E0075006F00750073002000710075006100730069002D0063006800610072006100630074006500720073>}
\MusicPDFWriteObject{10355}{<< /S /GoTo /D <73656374696F6E2E332E31> >>}
\MusicPDFWriteObject{10356}{<< /Title \MusicPDFRef{17685} /A \MusicPDFRef{17686} /Parent \MusicPDFRef{8878} /Prev \MusicPDFRef{9104} /Next \MusicPDFRef{10359} >>}
\MusicPDFWriteObject{10357}{<FEFF0033002E0034002E00200045006900670065006E0064006900730074007200690062007500740069006F006E007300200061006E00640020005700650069006C002700730020004900640065006E0074006900740079>}
\MusicPDFWriteObject{10358}{<< /S /GoTo /D <73656374696F6E2E332E34> >>}
\MusicPDFWriteObject{10359}{<< /Title \MusicPDFRef{17687} /A \MusicPDFRef{17688} /Parent \MusicPDFRef{8878} /Prev \MusicPDFRef{10356} /Next \MusicPDFRef{9105} >>}
\MusicPDFWriteObject{10360}{<FEFF0032002E0032002E00200047006C006F00620061006C0020006100640065006C006900630020002D00660075006E006300740069006F006E0073>}
\MusicPDFWriteObject{10361}{<< /S /GoTo /D <73656374696F6E2E322E32> >>}
\MusicPDFWriteObject{10362}{<FEFF0032002E0033002E00200053002D006C006F00630061006C0020006100640065006C006900630020002D00660075006E006300740069006F006E0073>}
\MusicPDFWriteObject{10363}{<< /S /GoTo /D <73656374696F6E2E322E33> >>}
\MusicPDFWriteObject{10364}{<FEFF0031002E0033002E0020004F00720062006900740061006C0020006400650063006F006D0070006F0073006900740069006F006E0020006F006600200074006800650020006C0069006600740065006400200063006F006E0076006F006C007500740069006F006E0020006B00650072006E0065006C002C00200061006E006400200074006800650020006F00720062006900740061006C00200070006900650063006500730020006F006600200074006800650020006B00650072006E0065006C>}
\MusicPDFWriteObject{10365}{<< /S /GoTo /D <73656374696F6E2E312E33> >>}
\MusicPDFWriteObject{10366}{<< /Title \MusicPDFRef{17689} /A \MusicPDFRef{17690} /Parent \MusicPDFRef{8822} /Prev \MusicPDFRef{9114} /Next \MusicPDFRef{10369} >>}
\MusicPDFWriteObject{10367}{<FEFF0031002E0036002E00200048006F00770020007400680069007300200077006F0072006B002000770061007300200063006F006E006400750063007400650064002C00200061006E00640020007700680079>}
\MusicPDFWriteObject{10368}{<< /S /GoTo /D <73656374696F6E2E312E36> >>}
\MusicPDFWriteObject{10369}{<< /Title \MusicPDFRef{17691} /A \MusicPDFRef{17692} /Parent \MusicPDFRef{8822} /Prev \MusicPDFRef{10366} /Next \MusicPDFRef{9117} >>}
\MusicPDFWriteObject{10370}{<FEFF0041007000700065006E00640069007800200042002E0020005000720065007000610072006100740069006F006E0073002000660072006F006D00200041006E0061006C00790074006900630020004E0075006D0062006500720020005400680065006F00720079>}
\MusicPDFWriteObject{10371}{<< /S /GoTo /D <617070656E6469782E42> >>}
\MusicPDFWriteObject{10372}{<< /Title \MusicPDFRef{17693} /A \MusicPDFRef{17694} /Parent \MusicPDFRef{1} /Prev \MusicPDFRef{17695} /Next \MusicPDFRef{9120} /First \MusicPDFRef{17696} /Last \MusicPDFRef{17697} /Count -5 >>}
\MusicPDFWriteObject{10373}{<< /Title \MusicPDFRef{17698} /A \MusicPDFRef{17699} /Parent \MusicPDFRef{9120} /Next \MusicPDFRef{17700} >>}
\MusicPDFWriteObject{10374}{<< /Title \MusicPDFRef{17701} /A \MusicPDFRef{17702} /Parent \MusicPDFRef{9120} /Prev \MusicPDFRef{17703} >>}
\MusicPDFWriteObject{10375}{<FEFF0043002E0031002E0020004C006F00630061006C006C007900200063006F006E00760065007800200074006F0070006F006C006F0067006900630061006C00200076006500630074006F00720020007300700061006300650073>}
\MusicPDFWriteObject{10376}{<< /S /GoTo /D <73656374696F6E2E432E31> >>}
\MusicPDFWriteObject{10377}{<< /Title \MusicPDFRef{17704} /A \MusicPDFRef{17705} /Parent \MusicPDFRef{8889} /Prev \MusicPDFRef{9121} /Next \MusicPDFRef{17706} >>}
\MusicPDFWriteObject{10378}{<FEFF0043002E0035002E0020004D00650072006F006D006F00720070006800690063002000660075006E006300740069006F006E0073002000760061006C00750065006400200069006E00200061002000710075006100730069002D0063006F006D0070006C0065007400650020006C006F00630061006C006C007900200063006F006E0076006500780020007400760073>}
\MusicPDFWriteObject{10379}{<< /S /GoTo /D <73656374696F6E2E432E35> >>}
\MusicPDFWriteObject{10380}{<< /Title \MusicPDFRef{17707} /A \MusicPDFRef{17708} /Parent \MusicPDFRef{8889} /Prev \MusicPDFRef{17706} /Next \MusicPDFRef{9122} >>}
\MusicPDFWriteObject{10381}{<FEFF0044002E0032002E00200053006300680077006100720074007A002D004200720075006800610074002000660075006E006300740069006F006E00730020006F006E0020006D006F00640075006C0061007400650064002000670072006F007500700073>}
\MusicPDFWriteObject{10382}{<< /S /GoTo /D <73656374696F6E2E442E32> >>}
\MusicPDFWriteObject{10383}{<FEFF0045002E0033002E00200053006D0061006C006C00200063006F006D0070006100630074006C0079002000670065006E0065007200610074006500640020007700650061006B006C0079002000480061007500730064006F0072006600660020007300700061006300650073>}
\MusicPDFWriteObject{10384}{<< /S /GoTo /D <73656374696F6E2E452E33> >>}
\MusicPDFWriteObject{10385}{<FEFF0045002E0034002E00200054006F0070006F006C006F0067006900630061006C002000670072006F00750070006F006900640073>}
\MusicPDFWriteObject{10386}{<< /S /GoTo /D <73656374696F6E2E452E34> >>}
\MusicPDFWriteObject{10387}{<< /D [\MusicPDFPageRef{1} /XYZ 72 698.082 null] >>}
\MusicPDFWriteObject{10388}{<< /D [\MusicPDFPageRef{9} /XYZ 75.573 96.045 null] >>}
\MusicPDFWriteObject{10389}{<< /D [\MusicPDFPageRef{137} /XYZ 75.772 93.296 null] >>}
\MusicPDFWriteObject{10390}{<< /D [\MusicPDFPageRef{619} /XYZ 79.742 96.045 null] >>}
\MusicPDFWriteObject{10391}{<< /D [\MusicPDFPageRef{620} /XYZ 79.742 106.637 null] >>}
\MusicPDFWriteObject{10392}{<< /D [\MusicPDFPageRef{620} /XYZ 79.742 96.045 null] >>}
\MusicPDFWriteObject{10393}{<< /D [\MusicPDFPageRef{623} /XYZ 79.742 96.045 null] >>}
\MusicPDFWriteObject{10394}{<< /D [\MusicPDFPageRef{624} /XYZ 79.742 106.637 null] >>}
\MusicPDFWriteObject{10395}{<< /D [\MusicPDFPageRef{624} /XYZ 79.742 96.045 null] >>}
\MusicPDFWriteObject{10396}{<< /D [\MusicPDFPageRef{625} /XYZ 79.742 96.045 null] >>}
\MusicPDFWriteObject{10397}{<< /D [\MusicPDFPageRef{633} /XYZ 79.742 93.377 null] >>}
\MusicPDFWriteObject{10398}{<< /D [\MusicPDFPageRef{643} /XYZ 79.742 96.045 null] >>}
\MusicPDFWriteObject{10399}{<< /D [\MusicPDFPageRef{648} /XYZ 79.742 96.045 null] >>}
\MusicPDFWriteObject{10400}{<< /D [\MusicPDFPageRef{144} /XYZ 75.772 94.451 null] >>}
\MusicPDFWriteObject{10401}{<< /D [\MusicPDFPageRef{666} /XYZ 75.573 106.007 null] >>}
\MusicPDFWriteObject{10402}{<< /D [\MusicPDFPageRef{667} /XYZ 75.772 147.328 null] >>}
\MusicPDFWriteObject{10403}{<< /D [\MusicPDFPageRef{667} /XYZ 75.469 137.375 null] >>}
\MusicPDFWriteObject{10404}{<< /D [\MusicPDFPageRef{667} /XYZ 75.692 116.599 null] >>}
\MusicPDFWriteObject{10405}{<< /D [\MusicPDFPageRef{667} /XYZ 75.97 106.007 null] >>}
\MusicPDFWriteObject{10406}{<< /D [\MusicPDFPageRef{672} /XYZ 75.97 106.007 null] >>}
\MusicPDFWriteObject{10407}{<< /D [\MusicPDFPageRef{677} /XYZ 75.772 116.599 null] >>}
\MusicPDFWriteObject{10408}{<< /D [\MusicPDFPageRef{677} /XYZ 75.97 96.045 null] >>}
\MusicPDFWriteObject{10409}{<< /D [\MusicPDFPageRef{678} /XYZ 75.97 96.045 null] >>}
\MusicPDFWriteObject{10410}{<< /D [\MusicPDFPageRef{679} /XYZ 79.543 96.045 null] >>}
\MusicPDFWriteObject{10411}{<< /D [\MusicPDFPageRef{199} /XYZ 75.573 96.045 null] >>}
\MusicPDFWriteObject{10412}{<< /D [\MusicPDFPageRef{681} /XYZ 79.543 96.045 null] >>}
\MusicPDFWriteObject{10413}{<< /D [\MusicPDFPageRef{682} /XYZ 79.543 96.045 null] >>}
\MusicPDFWriteObject{10414}{<< /D [\MusicPDFPageRef{683} /XYZ 79.543 116.599 null] >>}
\MusicPDFWriteObject{10415}{<< /D [\MusicPDFPageRef{683} /XYZ 79.543 106.007 null] >>}
\MusicPDFWriteObject{10416}{<< /D [\MusicPDFPageRef{687} /XYZ 79.543 96.045 null] >>}
\MusicPDFWriteObject{10417}{<< /D [\MusicPDFPageRef{688} /XYZ 79.543 105.208 null] >>}
\MusicPDFWriteObject{10418}{<< /D [\MusicPDFPageRef{689} /XYZ 79.543 106.637 null] >>}
\MusicPDFWriteObject{10419}{<< /D [\MusicPDFPageRef{689} /XYZ 79.543 96.045 null] >>}
\MusicPDFWriteObject{10420}{<< /D [\MusicPDFPageRef{691} /XYZ 79.543 106.007 null] >>}
\MusicPDFWriteObject{10421}{<< /D [\MusicPDFPageRef{717} /XYZ 79.742 106.007 null] >>}
\MusicPDFWriteObject{10422}{<< /D [\MusicPDFPageRef{200} /XYZ 75.772 96.045 null] >>}
\MusicPDFWriteObject{10423}{<< /D [\MusicPDFPageRef{718} /XYZ 79.742 96.045 null] >>}
\MusicPDFWriteObject{10424}{<< /D [\MusicPDFPageRef{753} /XYZ 75.573 96.045 null] >>}
\MusicPDFWriteObject{10425}{<< /D [\MusicPDFPageRef{757} /XYZ 75.772 96.045 null] >>}
\MusicPDFWriteObject{10426}{<< /D [\MusicPDFPageRef{771} /XYZ 75.772 115.97 null] >>}
\MusicPDFWriteObject{10427}{<< /D [\MusicPDFPageRef{773} /XYZ 75.692 105.439 null] >>}
\MusicPDFWriteObject{10428}{<< /D [\MusicPDFPageRef{774} /XYZ 75.97 115.827 null] >>}
\MusicPDFWriteObject{10429}{<< /D [\MusicPDFPageRef{774} /XYZ 75.97 95.273 null] >>}
\MusicPDFWriteObject{10430}{<< /D [\MusicPDFPageRef{775} /XYZ 75.772 116.599 null] >>}
\MusicPDFWriteObject{10431}{<< /D [\MusicPDFPageRef{775} /XYZ 75.97 96.045 null] >>}
\MusicPDFWriteObject{10432}{<< /D [\MusicPDFPageRef{795} /XYZ 75.97 96.045 null] >>}
\MusicPDFWriteObject{10433}{<< /D [\MusicPDFPageRef{230} /XYZ 75.573 96.045 null] >>}
\MusicPDFWriteObject{10434}{<< /D [\MusicPDFPageRef{797} /XYZ 79.543 126.562 null] >>}
\MusicPDFWriteObject{10435}{<< /D [\MusicPDFPageRef{797} /XYZ 79.543 96.045 null] >>}
\MusicPDFWriteObject{10436}{<< /D [\MusicPDFPageRef{802} /XYZ 79.543 116.398 null] >>}
\MusicPDFWriteObject{10437}{<< /D [\MusicPDFPageRef{802} /XYZ 79.543 95.844 null] >>}
\MusicPDFWriteObject{10438}{<< /D [\MusicPDFPageRef{805} /XYZ 79.543 96.045 null] >>}
\MusicPDFWriteObject{10439}{<< /D [\MusicPDFPageRef{806} /XYZ 79.543 96.045 null] >>}
\MusicPDFWriteObject{10440}{<< /D [\MusicPDFPageRef{807} /XYZ 79.543 96.045 null] >>}
\MusicPDFWriteObject{10441}{<< /D [\MusicPDFPageRef{809} /XYZ 79.543 157.708 null] >>}
\MusicPDFWriteObject{10442}{<< /D [\MusicPDFPageRef{809} /XYZ 79.543 137.154 null] >>}
\MusicPDFWriteObject{10443}{<< /D [\MusicPDFPageRef{809} /XYZ 79.543 116.599 null] >>}
\MusicPDFWriteObject{10444}{<< /D [\MusicPDFPageRef{231} /XYZ 75.772 96.045 null] >>}
\MusicPDFWriteObject{10445}{<< /D [\MusicPDFPageRef{809} /XYZ 79.742 106.007 null] >>}
\MusicPDFWriteObject{10446}{<< /D [\MusicPDFPageRef{810} /XYZ 79.742 116.136 null] >>}
\MusicPDFWriteObject{10447}{<< /D [\MusicPDFPageRef{823} /XYZ 75.05 96.045 null] >>}
\MusicPDFWriteObject{10448}{<< /D [\MusicPDFPageRef{828} /XYZ 75.772 96.045 null] >>}
\MusicPDFWriteObject{10449}{<< /D [\MusicPDFPageRef{840} /XYZ 75.772 106.007 null] >>}
\MusicPDFWriteObject{10450}{<< /D [\MusicPDFPageRef{843} /XYZ 75.692 96.045 null] >>}
\MusicPDFWriteObject{10451}{<< /D [\MusicPDFPageRef{850} /XYZ 75.97 127.82 null] >>}
\MusicPDFWriteObject{10452}{<< /D [\MusicPDFPageRef{850} /XYZ 75.97 117.228 null] >>}
\MusicPDFWriteObject{10453}{<< /D [\MusicPDFPageRef{850} /XYZ 75.772 106.637 null] >>}
\MusicPDFWriteObject{10454}{<< /D [\MusicPDFPageRef{850} /XYZ 75.97 96.045 null] >>}
\MusicPDFWriteObject{10455}{<< /D [\MusicPDFPageRef{232} /XYZ 75.772 96.045 null] >>}
\MusicPDFWriteObject{10456}{<< /D [\MusicPDFPageRef{861} /XYZ 75.97 106.637 null] >>}
\MusicPDFWriteObject{10457}{<< /D [\MusicPDFPageRef{861} /XYZ 79.543 96.045 null] >>}
\MusicPDFWriteObject{10458}{<< /D [\MusicPDFPageRef{862} /XYZ 79.543 96.045 null] >>}
\MusicPDFWriteObject{10459}{<< /D [\MusicPDFPageRef{864} /XYZ 79.543 106.637 null] >>}
\MusicPDFWriteObject{10460}{<< /D [\MusicPDFPageRef{864} /XYZ 79.543 96.045 null] >>}
\MusicPDFWriteObject{10461}{<< /D [\MusicPDFPageRef{865} /XYZ 79.543 96.045 null] >>}
\MusicPDFWriteObject{10462}{<< /D [\MusicPDFPageRef{868} /XYZ 79.543 96.045 null] >>}
\MusicPDFWriteObject{10463}{<< /D [\MusicPDFPageRef{869} /XYZ 79.543 116.599 null] >>}
\MusicPDFWriteObject{10464}{<< /D [\MusicPDFPageRef{869} /XYZ 79.543 106.007 null] >>}
\MusicPDFWriteObject{10465}{<< /D [\MusicPDFPageRef{894} /XYZ 79.543 96.045 null] >>}
\MusicPDFWriteObject{10466}{<< /D [\MusicPDFPageRef{233} /XYZ 75.692 96.045 null] >>}
\MusicPDFWriteObject{10467}{<< /D [\MusicPDFPageRef{899} /XYZ 79.543 96.045 null] >>}
\MusicPDFWriteObject{10468}{<< /D [\MusicPDFPageRef{905} /XYZ 79.742 96.045 null] >>}
\MusicPDFWriteObject{10469}{<< /D [\MusicPDFPageRef{920} /XYZ 79.742 96.045 null] >>}
\MusicPDFWriteObject{10470}{<< /D [\MusicPDFPageRef{926} /XYZ 79.742 116.599 null] >>}
\MusicPDFWriteObject{10471}{<< /D [\MusicPDFPageRef{926} /XYZ 79.742 106.007 null] >>}
\MusicPDFWriteObject{10472}{<< /D [\MusicPDFPageRef{931} /XYZ 79.742 137.783 null] >>}
\MusicPDFWriteObject{10473}{<< /D [\MusicPDFPageRef{931} /XYZ 79.742 127.191 null] >>}
\MusicPDFWriteObject{10474}{<< /D [\MusicPDFPageRef{931} /XYZ 79.742 116.599 null] >>}
\MusicPDFWriteObject{10475}{<< /D [\MusicPDFPageRef{931} /XYZ 79.742 106.007 null] >>}
\MusicPDFWriteObject{10476}{<< /D [\MusicPDFPageRef{234} /XYZ 75.97 96.045 null] >>}
\MusicPDFWriteObject{10477}{<< /D [\MusicPDFPageRef{235} /XYZ 75.97 96.045 null] >>}
\MusicPDFWriteObject{10478}{<< /D [\MusicPDFPageRef{16} /XYZ 75.573 106.007 null] >>}
\MusicPDFWriteObject{10479}{<< /D [\MusicPDFPageRef{236} /XYZ 75.772 96.045 null] >>}
\MusicPDFWriteObject{10480}{<< /D [\MusicPDFPageRef{250} /XYZ 75.97 96.045 null] >>}
\MusicPDFWriteObject{10481}{<< /D [\MusicPDFPageRef{294} /XYZ 75.573 96.045 null] >>}
\MusicPDFWriteObject{10482}{<< /D [\MusicPDFPageRef{295} /XYZ 75.772 106.637 null] >>}
\MusicPDFWriteObject{10483}{<< /D [\MusicPDFPageRef{295} /XYZ 75.772 96.045 null] >>}
\MusicPDFWriteObject{10484}{<< /D [\MusicPDFPageRef{321} /XYZ 75.573 106.007 null] >>}
\MusicPDFWriteObject{10485}{<< /D [\MusicPDFPageRef{414} /XYZ 75.573 106.637 null] >>}
\MusicPDFWriteObject{10486}{<< /D [\MusicPDFPageRef{414} /XYZ 75.772 96.045 null] >>}
\MusicPDFWriteObject{10487}{<< /D [\MusicPDFPageRef{415} /XYZ 75.772 96.045 null] >>}
\MusicPDFWriteObject{10488}{<< /D [\MusicPDFPageRef{416} /XYZ 75.692 95.181 null] >>}
\MusicPDFWriteObject{10489}{<< /D [\MusicPDFPageRef{20} /XYZ 75.772 106.007 null] >>}
\MusicPDFWriteObject{10490}{<< /D [\MusicPDFPageRef{417} /XYZ 75.97 96.045 null] >>}
\MusicPDFWriteObject{10491}{<< /D [\MusicPDFPageRef{426} /XYZ 75.97 96.045 null] >>}
\MusicPDFWriteObject{10492}{<< /D [\MusicPDFPageRef{435} /XYZ 75.573 106.637 null] >>}
\MusicPDFWriteObject{10493}{<< /D [\MusicPDFPageRef{435} /XYZ 75.772 96.045 null] >>}
\MusicPDFWriteObject{10494}{<< /D [\MusicPDFPageRef{436} /XYZ 75.772 96.045 null] >>}
\MusicPDFWriteObject{10495}{<< /D [\MusicPDFPageRef{437} /XYZ 75.692 94.373 null] >>}
\MusicPDFWriteObject{10496}{<< /D [\MusicPDFPageRef{438} /XYZ 75.97 94.07 null] >>}
\MusicPDFWriteObject{10497}{<< /D [\MusicPDFPageRef{446} /XYZ 75.97 96.045 null] >>}
\MusicPDFWriteObject{10498}{<< /D [\MusicPDFPageRef{462} /XYZ 75.772 96.045 null] >>}
\MusicPDFWriteObject{10499}{<< /D [\MusicPDFPageRef{464} /XYZ 75.97 96.045 null] >>}
\MusicPDFWriteObject{10500}{<< /D [\MusicPDFPageRef{21} /XYZ 75.772 106.007 null] >>}
\MusicPDFWriteObject{10501}{<< /D [\MusicPDFPageRef{475} /XYZ 75.573 106.637 null] >>}
\MusicPDFWriteObject{10502}{<< /D [\MusicPDFPageRef{475} /XYZ 75.772 96.045 null] >>}
\MusicPDFWriteObject{10503}{<< /D [\MusicPDFPageRef{480} /XYZ 75.772 96.045 null] >>}
\MusicPDFWriteObject{10504}{<< /D [\MusicPDFPageRef{482} /XYZ 75.692 106.637 null] >>}
\MusicPDFWriteObject{10505}{<< /D [\MusicPDFPageRef{482} /XYZ 75.97 96.045 null] >>}
\MusicPDFWriteObject{10506}{<< /D [\MusicPDFPageRef{483} /XYZ 75.97 96.045 null] >>}
\MusicPDFWriteObject{10507}{<< /D [\MusicPDFPageRef{514} /XYZ 75.573 130.155 null] >>}
\MusicPDFWriteObject{10508}{<< /D [\MusicPDFPageRef{514} /XYZ 75.772 119.563 null] >>}
\MusicPDFWriteObject{10509}{<< /D [\MusicPDFPageRef{515} /XYZ 75.772 106.637 null] >>}
\MusicPDFWriteObject{10510}{<< /D [\MusicPDFPageRef{515} /XYZ 75.692 96.045 null] >>}
\MusicPDFWriteObject{10511}{<< /D [\MusicPDFPageRef{22} /XYZ 75.692 106.007 null] >>}
\MusicPDFWriteObject{10512}{<< /D [\MusicPDFPageRef{522} /XYZ 75.573 117.228 null] >>}
\MusicPDFWriteObject{10513}{<< /D [\MusicPDFPageRef{522} /XYZ 75.772 106.637 null] >>}
\MusicPDFWriteObject{10514}{<< /D [\MusicPDFPageRef{522} /XYZ 75.772 96.045 null] >>}
\MusicPDFWriteObject{10515}{<< /D [\MusicPDFPageRef{528} /XYZ 75.573 96.045 null] >>}
\MusicPDFWriteObject{10516}{<< /D [\MusicPDFPageRef{536} /XYZ 75.772 96.045 null] >>}
\MusicPDFWriteObject{10517}{<< /D [\MusicPDFPageRef{556} /XYZ 75.573 96.045 null] >>}
\MusicPDFWriteObject{10518}{<< /D [\MusicPDFPageRef{559} /XYZ 75.772 106.007 null] >>}
\MusicPDFWriteObject{10519}{<< /D [\MusicPDFPageRef{562} /XYZ 75.772 96.045 null] >>}
\MusicPDFWriteObject{10520}{<< /D [\MusicPDFPageRef{566} /XYZ 75.692 96.045 null] >>}
\MusicPDFWriteObject{10521}{<< /D [\MusicPDFPageRef{568} /XYZ 75.97 106.637 null] >>}
\MusicPDFWriteObject{10522}{<< /D [\MusicPDFPageRef{27} /XYZ 75.97 106.637 null] >>}
\MusicPDFWriteObject{10523}{<< /D [\MusicPDFPageRef{568} /XYZ 75.97 96.045 null] >>}
\MusicPDFWriteObject{10524}{<< /D [\MusicPDFPageRef{571} /XYZ 75.772 106.637 null] >>}
\MusicPDFWriteObject{10525}{<< /D [\MusicPDFPageRef{571} /XYZ 75.97 96.045 null] >>}
\MusicPDFWriteObject{10526}{<< /D [\MusicPDFPageRef{572} /XYZ 75.97 94.489 null] >>}
\MusicPDFWriteObject{10527}{<< /D [\MusicPDFPageRef{576} /XYZ 79.543 117.228 null] >>}
\MusicPDFWriteObject{10528}{<< /D [\MusicPDFPageRef{576} /XYZ 79.543 106.637 null] >>}
\MusicPDFWriteObject{10529}{<< /D [\MusicPDFPageRef{576} /XYZ 79.543 96.045 null] >>}
\MusicPDFWriteObject{10530}{<< /D [\MusicPDFPageRef{578} /XYZ 79.543 96.045 null] >>}
\MusicPDFWriteObject{10531}{<< /D [\MusicPDFPageRef{579} /XYZ 79.543 117.228 null] >>}
\MusicPDFWriteObject{10532}{<< /D [\MusicPDFPageRef{579} /XYZ 79.543 106.637 null] >>}
\MusicPDFWriteObject{10533}{<< /D [\MusicPDFPageRef{27} /XYZ 75.97 96.045 null] >>}
\MusicPDFWriteObject{10534}{<< /D [\MusicPDFPageRef{579} /XYZ 79.543 96.045 null] >>}
\MusicPDFWriteObject{10535}{<< /D [\MusicPDFPageRef{580} /XYZ 79.543 106.637 null] >>}
\MusicPDFWriteObject{10536}{<< /D [\MusicPDFPageRef{580} /XYZ 79.543 96.045 null] >>}
\MusicPDFWriteObject{10537}{<< /D [\MusicPDFPageRef{581} /XYZ 79.543 106.637 null] >>}
\MusicPDFWriteObject{10538}{<< /D [\MusicPDFPageRef{581} /XYZ 79.742 96.045 null] >>}
\MusicPDFWriteObject{10539}{<< /D [\MusicPDFPageRef{583} /XYZ 79.742 96.045 null] >>}
\MusicPDFWriteObject{10540}{<< /D [\MusicPDFPageRef{584} /XYZ 79.742 106.637 null] >>}
\MusicPDFWriteObject{10541}{<< /D [\MusicPDFPageRef{584} /XYZ 79.742 96.045 null] >>}
\MusicPDFWriteObject{10542}{<< /D [\MusicPDFPageRef{585} /XYZ 79.742 106.637 null] >>}
\MusicPDFWriteObject{10543}{<< /D [\MusicPDFPageRef{585} /XYZ 79.742 96.045 null] >>}
\MusicPDFWriteObject{10544}{<< /D [\MusicPDFPageRef{52} /XYZ 75.772 106.007 null] >>}
\MusicPDFWriteObject{10545}{<< /D [\MusicPDFPageRef{587} /XYZ 75.573 96.045 null] >>}
\MusicPDFWriteObject{10546}{<< /D [\MusicPDFPageRef{591} /XYZ 75.772 117.228 null] >>}
\MusicPDFWriteObject{10547}{<< /D [\MusicPDFPageRef{591} /XYZ 75.772 106.637 null] >>}
\MusicPDFWriteObject{10548}{<< /D [\MusicPDFPageRef{591} /XYZ 75.692 96.045 null] >>}
\MusicPDFWriteObject{10549}{<< /D [\MusicPDFPageRef{592} /XYZ 75.97 117.228 null] >>}
\MusicPDFWriteObject{10550}{<< /D [\MusicPDFPageRef{592} /XYZ 75.97 106.637 null] >>}
\MusicPDFWriteObject{10551}{<< /D [\MusicPDFPageRef{592} /XYZ 75.772 96.045 null] >>}
\MusicPDFWriteObject{10552}{<< /D [\MusicPDFPageRef{593} /XYZ 75.97 96.045 null] >>}
\MusicPDFWriteObject{10553}{<< /D [\MusicPDFPageRef{595} /XYZ 75.97 93.394 null] >>}
\MusicPDFWriteObject{10554}{<< /D [\MusicPDFPageRef{603} /XYZ 79.543 126.562 null] >>}
\MusicPDFWriteObject{10555}{<< /D [\MusicPDFPageRef{131} /XYZ 75.573 96.045 null] >>}
\MusicPDFWriteObject{10556}{<< /D [\MusicPDFPageRef{603} /XYZ 79.543 106.007 null] >>}
\MusicPDFWriteObject{10557}{<< /D [\MusicPDFPageRef{608} /XYZ 79.543 96.045 null] >>}
\MusicPDFWriteObject{10558}{<< /D [\MusicPDFPageRef{614} /XYZ 79.543 96.045 null] >>}
\MusicPDFWriteObject{10559}{<< /D [\MusicPDFPageRef{616} /XYZ 79.543 106.637 null] >>}
\MusicPDFWriteObject{10560}{<< /D [\MusicPDFPageRef{616} /XYZ 79.543 96.045 null] >>}
\MusicPDFWriteObject{10561}{<< /D [\MusicPDFPageRef{619} /XYZ 79.543 171.427 null] >>}
\MusicPDFWriteObject{10562}{<< /D [\MusicPDFPageRef{619} /XYZ 79.543 159.715 null] >>}
\MusicPDFWriteObject{10563}{<< /D [\MusicPDFPageRef{619} /XYZ 79.543 147.233 null] >>}
\MusicPDFWriteObject{10564}{<< /D [\MusicPDFPageRef{619} /XYZ 78.939 127.359 null] >>}
\MusicPDFWriteObject{10565}{<< /D [\MusicPDFPageRef{619} /XYZ 79.742 116.767 null] >>}
\MusicPDFWriteObject{10566}{<< /D [\MusicPDFPageRef{13} /XYZ 72 321.089 null] >>}
\MusicPDFWriteObject{10567}{<< /D [\MusicPDFPageRef{21} /XYZ 72 152.178 null] >>}
\MusicPDFWriteObject{10568}{<< /D [\MusicPDFPageRef{100} /XYZ 72 571.594 null] >>}
\MusicPDFWriteObject{10569}{<< /D [\MusicPDFPageRef{710} /XYZ 72 335.335 null] >>}
\MusicPDFWriteObject{10570}{<< /D [\MusicPDFPageRef{710} /XYZ 72 246.043 null] >>}
\MusicPDFWriteObject{10571}{<< /D [\MusicPDFPageRef{712} /XYZ 72 670.386 null] >>}
\MusicPDFWriteObject{10572}{<< /D [\MusicPDFPageRef{712} /XYZ 72 598.075 null] >>}
\MusicPDFWriteObject{10573}{<< /D [\MusicPDFPageRef{715} /XYZ 72 152.694 null] >>}
\MusicPDFWriteObject{10574}{<< /D [\MusicPDFPageRef{715} /XYZ 72 139.544 null] >>}
\MusicPDFWriteObject{10575}{<< /D [\MusicPDFPageRef{717} /XYZ 72 621.399 null] >>}
\MusicPDFWriteObject{10576}{<< /D [\MusicPDFPageRef{717} /XYZ 72 608.249 null] >>}
\MusicPDFWriteObject{10577}{<< /D [\MusicPDFPageRef{717} /XYZ 72 456.279 null] >>}
\MusicPDFWriteObject{10578}{<< /D [\MusicPDFPageRef{717} /XYZ 72 440.278 null] >>}
\MusicPDFWriteObject{10579}{<< /D [\MusicPDFPageRef{100} /XYZ 72 529.794 null] >>}
\MusicPDFWriteObject{10580}{<< /D [\MusicPDFPageRef{717} /XYZ 72 427.127 null] >>}
\MusicPDFWriteObject{10581}{<< /D [\MusicPDFPageRef{717} /XYZ 72 401.74 null] >>}
\MusicPDFWriteObject{10582}{<< /D [\MusicPDFPageRef{718} /XYZ 72 167.916 null] >>}
\MusicPDFWriteObject{10583}{<< /D [\MusicPDFPageRef{719} /XYZ 72 698.082 null] >>}
\MusicPDFWriteObject{10584}{<< /D [\MusicPDFPageRef{721} /XYZ 72 441.957 null] >>}
\MusicPDFWriteObject{10585}{<< /D [\MusicPDFPageRef{721} /XYZ 72 319.222 null] >>}
\MusicPDFWriteObject{10586}{<< /D [\MusicPDFPageRef{721} /XYZ 72 277.279 null] >>}
\MusicPDFWriteObject{10587}{<< /D [\MusicPDFPageRef{721} /XYZ 72 206.725 null] >>}
\MusicPDFWriteObject{10588}{<< /D [\MusicPDFPageRef{721} /XYZ 72 193.02 null] >>}
\MusicPDFWriteObject{10589}{<< /D [\MusicPDFPageRef{721} /XYZ 72 166.718 null] >>}
\MusicPDFWriteObject{10590}{<< /D [\MusicPDFPageRef{100} /XYZ 72 481.65 null] >>}
\MusicPDFWriteObject{10591}{<< /D [\MusicPDFPageRef{721} /XYZ 72 142.908 null] >>}
\MusicPDFWriteObject{10592}{<< /D [\MusicPDFPageRef{724} /XYZ 72 651.029 null] >>}
\MusicPDFWriteObject{10593}{<< /D [\MusicPDFPageRef{724} /XYZ 72 637.325 null] >>}
\MusicPDFWriteObject{10594}{<< /D [\MusicPDFPageRef{724} /XYZ 72 525.936 null] >>}
\MusicPDFWriteObject{10595}{<< /D [\MusicPDFPageRef{724} /XYZ 72 499.634 null] >>}
\MusicPDFWriteObject{10596}{<< /D [\MusicPDFPageRef{724} /XYZ 72 356.283 null] >>}
\MusicPDFWriteObject{10597}{<< /D [\MusicPDFPageRef{724} /XYZ 72 342.579 null] >>}
\MusicPDFWriteObject{10598}{<< /D [\MusicPDFPageRef{724} /XYZ 72 329.428 null] >>}
\MusicPDFWriteObject{10599}{<< /D [\MusicPDFPageRef{725} /XYZ 72 535.289 null] >>}
\MusicPDFWriteObject{10600}{<< /D [\MusicPDFPageRef{725} /XYZ 72 508.434 null] >>}
\MusicPDFWriteObject{10601}{<< /D [\MusicPDFPageRef{100} /XYZ 72 442.197 null] >>}
\MusicPDFWriteObject{10602}{<< /D [\MusicPDFPageRef{726} /XYZ 72 507.673 null] >>}
\MusicPDFWriteObject{10603}{<< /D [\MusicPDFPageRef{726} /XYZ 72 423.196 null] >>}
\MusicPDFWriteObject{10604}{<< /D [\MusicPDFPageRef{728} /XYZ 72 411.037 null] >>}
\MusicPDFWriteObject{10605}{<< /D [\MusicPDFPageRef{728} /XYZ 72 384.181 null] >>}
\MusicPDFWriteObject{10606}{<< /D [\MusicPDFPageRef{728} /XYZ 72 371.031 null] >>}
\MusicPDFWriteObject{10607}{<< /D [\MusicPDFPageRef{728} /XYZ 72 357.88 null] >>}
\MusicPDFWriteObject{10608}{<< /D [\MusicPDFPageRef{729} /XYZ 72 698.082 null] >>}
\MusicPDFWriteObject{10609}{<< /D [\MusicPDFPageRef{729} /XYZ 72 673.032 null] >>}
\MusicPDFWriteObject{10610}{<< /D [\MusicPDFPageRef{729} /XYZ 72 550.463 null] >>}
\MusicPDFWriteObject{10611}{<< /D [\MusicPDFPageRef{730} /XYZ 72 457.612 null] >>}
\MusicPDFWriteObject{10612}{<< /D [\MusicPDFPageRef{105} /XYZ 72 466.789 null] >>}
\MusicPDFWriteObject{10613}{<< /D [\MusicPDFPageRef{730} /XYZ 72 441.971 null] >>}
\MusicPDFWriteObject{10614}{<< /D [\MusicPDFPageRef{731} /XYZ 72 633.424 null] >>}
\MusicPDFWriteObject{10615}{<< /D [\MusicPDFPageRef{731} /XYZ 72 617.783 null] >>}
\MusicPDFWriteObject{10616}{<< /D [\MusicPDFPageRef{732} /XYZ 72 698.082 null] >>}
\MusicPDFWriteObject{10617}{<< /D [\MusicPDFPageRef{732} /XYZ 72 685.629 null] >>}
\MusicPDFWriteObject{10618}{<< /D [\MusicPDFPageRef{732} /XYZ 72 659.881 null] >>}
\MusicPDFWriteObject{10619}{<< /D [\MusicPDFPageRef{732} /XYZ 72 633.58 null] >>}
\MusicPDFWriteObject{10620}{<< /D [\MusicPDFPageRef{732} /XYZ 72 226.33 null] >>}
\MusicPDFWriteObject{10621}{<< /D [\MusicPDFPageRef{732} /XYZ 72 213.18 null] >>}
\MusicPDFWriteObject{10622}{<< /D [\MusicPDFPageRef{733} /XYZ 72 123.597 null] >>}
\MusicPDFWriteObject{10623}{<< /D [\MusicPDFPageRef{105} /XYZ 72 428.586 null] >>}
\MusicPDFWriteObject{10624}{<< /D [\MusicPDFPageRef{733} /XYZ 72 109.894 null] >>}
\MusicPDFWriteObject{10625}{<< /D [\MusicPDFPageRef{734} /XYZ 72 698.082 null] >>}
\MusicPDFWriteObject{10626}{<< /D [\MusicPDFPageRef{734} /XYZ 72 181.258 null] >>}
\MusicPDFWriteObject{10627}{<< /D [\MusicPDFPageRef{734} /XYZ 72 167.554 null] >>}
\MusicPDFWriteObject{10628}{<< /D [\MusicPDFPageRef{735} /XYZ 72 501.168 null] >>}
\MusicPDFWriteObject{10629}{<< /D [\MusicPDFPageRef{735} /XYZ 72 422.263 null] >>}
\MusicPDFWriteObject{10630}{<< /D [\MusicPDFPageRef{736} /XYZ 72 488.473 null] >>}
\MusicPDFWriteObject{10631}{<< /D [\MusicPDFPageRef{736} /XYZ 72 474.769 null] >>}
\MusicPDFWriteObject{10632}{<< /D [\MusicPDFPageRef{736} /XYZ 72 461.618 null] >>}
\MusicPDFWriteObject{10633}{<< /D [\MusicPDFPageRef{736} /XYZ 72 448.467 null] >>}
\MusicPDFWriteObject{10634}{<< /D [\MusicPDFPageRef{105} /XYZ 72 379.008 null] >>}
\MusicPDFWriteObject{10635}{<< /D [\MusicPDFPageRef{737} /XYZ 72 305.08 null] >>}
\MusicPDFWriteObject{10636}{<< /D [\MusicPDFPageRef{737} /XYZ 72 278.225 null] >>}
\MusicPDFWriteObject{10637}{<< /D [\MusicPDFPageRef{737} /XYZ 72 251.564 null] >>}
\MusicPDFWriteObject{10638}{<< /D [\MusicPDFPageRef{737} /XYZ 72 226.176 null] >>}
\MusicPDFWriteObject{10639}{<< /D [\MusicPDFPageRef{738} /XYZ 72 572.567 null] >>}
\MusicPDFWriteObject{10640}{<< /D [\MusicPDFPageRef{738} /XYZ 72 558.863 null] >>}
\MusicPDFWriteObject{10641}{<< /D [\MusicPDFPageRef{738} /XYZ 72 488.378 null] >>}
\MusicPDFWriteObject{10642}{<< /D [\MusicPDFPageRef{738} /XYZ 72 308.496 null] >>}
\MusicPDFWriteObject{10643}{<< /D [\MusicPDFPageRef{738} /XYZ 72 294.792 null] >>}
\MusicPDFWriteObject{10644}{<< /D [\MusicPDFPageRef{738} /XYZ 72 281.642 null] >>}
\MusicPDFWriteObject{10645}{<< /D [\MusicPDFPageRef{105} /XYZ 72 352.707 null] >>}
\MusicPDFWriteObject{10646}{<< /D [\MusicPDFPageRef{738} /XYZ 72 197.107 null] >>}
\MusicPDFWriteObject{10647}{<< /D [\MusicPDFPageRef{739} /XYZ 72 698.082 null] >>}
\MusicPDFWriteObject{10648}{<< /D [\MusicPDFPageRef{739} /XYZ 72 521.776 null] >>}
\MusicPDFWriteObject{10649}{<< /D [\MusicPDFPageRef{739} /XYZ 72 216.077 null] >>}
\MusicPDFWriteObject{10650}{<< /D [\MusicPDFPageRef{739} /XYZ 72 187.286 null] >>}
\MusicPDFWriteObject{10651}{<< /D [\MusicPDFPageRef{740} /XYZ 72 405.683 null] >>}
\MusicPDFWriteObject{10652}{<< /D [\MusicPDFPageRef{740} /XYZ 72 355.191 null] >>}
\MusicPDFWriteObject{10653}{<< /D [\MusicPDFPageRef{740} /XYZ 72 285.085 null] >>}
\MusicPDFWriteObject{10654}{<< /D [\MusicPDFPageRef{740} /XYZ 72 224.777 null] >>}
\MusicPDFWriteObject{10655}{<< /D [\MusicPDFPageRef{746} /XYZ 72 561.746 null] >>}
\MusicPDFWriteObject{10656}{<< /D [\MusicPDFPageRef{107} /XYZ 72 532.662 null] >>}
\MusicPDFWriteObject{10657}{<< /D [\MusicPDFPageRef{746} /XYZ 72 535.948 null] >>}
\MusicPDFWriteObject{10658}{<< /D [\MusicPDFPageRef{746} /XYZ 72 423.661 null] >>}
\MusicPDFWriteObject{10659}{<< /D [\MusicPDFPageRef{746} /XYZ 72 128.871 null] >>}
\MusicPDFWriteObject{10660}{<< /D [\MusicPDFPageRef{746} /XYZ 72 115.166 null] >>}
\MusicPDFWriteObject{10661}{<< /D [\MusicPDFPageRef{747} /XYZ 72 685.629 null] >>}
\MusicPDFWriteObject{10662}{<< /D [\MusicPDFPageRef{747} /XYZ 72 649.201 null] >>}
\MusicPDFWriteObject{10663}{<< /D [\MusicPDFPageRef{747} /XYZ 72 562.108 null] >>}
\MusicPDFWriteObject{10664}{<< /D [\MusicPDFPageRef{747} /XYZ 72 512.923 null] >>}
\MusicPDFWriteObject{10665}{<< /D [\MusicPDFPageRef{748} /XYZ 72 613.482 null] >>}
\MusicPDFWriteObject{10666}{<< /D [\MusicPDFPageRef{748} /XYZ 72 585.08 null] >>}
\MusicPDFWriteObject{10667}{<< /D [\MusicPDFPageRef{107} /XYZ 72 483.083 null] >>}
\MusicPDFWriteObject{10668}{<< /D [\MusicPDFPageRef{748} /XYZ 72 571.929 null] >>}
\MusicPDFWriteObject{10669}{<< /D [\MusicPDFPageRef{748} /XYZ 72 244.781 null] >>}
\MusicPDFWriteObject{10670}{<< /D [\MusicPDFPageRef{748} /XYZ 72 231.077 null] >>}
\MusicPDFWriteObject{10671}{<< /D [\MusicPDFPageRef{748} /XYZ 72 217.926 null] >>}
\MusicPDFWriteObject{10672}{<< /D [\MusicPDFPageRef{749} /XYZ 72 535.046 null] >>}
\MusicPDFWriteObject{10673}{<< /D [\MusicPDFPageRef{749} /XYZ 72 520.481 null] >>}
\MusicPDFWriteObject{10674}{<< /D [\MusicPDFPageRef{749} /XYZ 72 463.615 null] >>}
\MusicPDFWriteObject{10675}{<< /D [\MusicPDFPageRef{749} /XYZ 72 449.911 null] >>}
\MusicPDFWriteObject{10676}{<< /D [\MusicPDFPageRef{749} /XYZ 72 436.76 null] >>}
\MusicPDFWriteObject{10677}{<< /D [\MusicPDFPageRef{750} /XYZ 72 479.814 null] >>}
\MusicPDFWriteObject{10678}{<< /D [\MusicPDFPageRef{22} /XYZ 72 698.082 null] >>}
\MusicPDFWriteObject{10679}{<< /D [\MusicPDFPageRef{113} /XYZ 72 623.116 null] >>}
\MusicPDFWriteObject{10680}{<< /D [\MusicPDFPageRef{752} /XYZ 72 698.082 null] >>}
\MusicPDFWriteObject{10681}{<< /D [\MusicPDFPageRef{752} /XYZ 72 635.785 null] >>}
\MusicPDFWriteObject{10682}{<< /D [\MusicPDFPageRef{753} /XYZ 72 618.827 null] >>}
\MusicPDFWriteObject{10683}{<< /D [\MusicPDFPageRef{753} /XYZ 72 520.224 null] >>}
\MusicPDFWriteObject{10684}{<< /D [\MusicPDFPageRef{753} /XYZ 72 483.243 null] >>}
\MusicPDFWriteObject{10685}{<< /D [\MusicPDFPageRef{757} /XYZ 72 173.833 null] >>}
\MusicPDFWriteObject{10686}{<< /D [\MusicPDFPageRef{757} /XYZ 72 146.978 null] >>}
\MusicPDFWriteObject{10687}{<< /D [\MusicPDFPageRef{758} /XYZ 72 190.641 null] >>}
\MusicPDFWriteObject{10688}{<< /D [\MusicPDFPageRef{758} /XYZ 72 176.937 null] >>}
\MusicPDFWriteObject{10689}{<< /D [\MusicPDFPageRef{758} /XYZ 72 162.372 null] >>}
\MusicPDFWriteObject{10690}{<< /D [\MusicPDFPageRef{113} /XYZ 72 556.232 null] >>}
\MusicPDFWriteObject{10691}{<< /D [\MusicPDFPageRef{758} /XYZ 72 149.172 null] >>}
\MusicPDFWriteObject{10692}{<< /D [\MusicPDFPageRef{760} /XYZ 72 271.759 null] >>}
\MusicPDFWriteObject{10693}{<< /D [\MusicPDFPageRef{760} /XYZ 72 245.408 null] >>}
\MusicPDFWriteObject{10694}{<< /D [\MusicPDFPageRef{761} /XYZ 72 668.664 null] >>}
\MusicPDFWriteObject{10695}{<< /D [\MusicPDFPageRef{761} /XYZ 72 654.959 null] >>}
\MusicPDFWriteObject{10696}{<< /D [\MusicPDFPageRef{761} /XYZ 72 641.759 null] >>}
\MusicPDFWriteObject{10697}{<< /D [\MusicPDFPageRef{762} /XYZ 72 670.43 null] >>}
\MusicPDFWriteObject{10698}{<< /D [\MusicPDFPageRef{762} /XYZ 72 590.972 null] >>}
\MusicPDFWriteObject{10699}{<< /D [\MusicPDFPageRef{762} /XYZ 72 577.822 null] >>}
\MusicPDFWriteObject{10700}{<< /D [\MusicPDFPageRef{762} /XYZ 72 564.67 null] >>}
\MusicPDFWriteObject{10701}{<< /D [\MusicPDFPageRef{113} /XYZ 72 516.78 null] >>}
\MusicPDFWriteObject{10702}{<< /D [\MusicPDFPageRef{763} /XYZ 72 284.116 null] >>}
\MusicPDFWriteObject{10703}{<< /D [\MusicPDFPageRef{763} /XYZ 72 257.814 null] >>}
\MusicPDFWriteObject{10704}{<< /D [\MusicPDFPageRef{767} /XYZ 72 399.632 null] >>}
\MusicPDFWriteObject{10705}{<< /D [\MusicPDFPageRef{767} /XYZ 72 373.33 null] >>}
\MusicPDFWriteObject{10706}{<< /D [\MusicPDFPageRef{768} /XYZ 72 670.304 null] >>}
\MusicPDFWriteObject{10707}{<< /D [\MusicPDFPageRef{768} /XYZ 72 615.42 null] >>}
\MusicPDFWriteObject{10708}{<< /D [\MusicPDFPageRef{770} /XYZ 72 257.153 null] >>}
\MusicPDFWriteObject{10709}{<< /D [\MusicPDFPageRef{770} /XYZ 72 243.449 null] >>}
\MusicPDFWriteObject{10710}{<< /D [\MusicPDFPageRef{773} /XYZ 72 364.164 null] >>}
\MusicPDFWriteObject{10711}{<< /D [\MusicPDFPageRef{773} /XYZ 72 291.539 null] >>}
\MusicPDFWriteObject{10712}{<< /D [\MusicPDFPageRef{119} /XYZ 72 247.202 null] >>}
\MusicPDFWriteObject{10713}{<< /D [\MusicPDFPageRef{773} /XYZ 72 265.138 null] >>}
\MusicPDFWriteObject{10714}{<< /D [\MusicPDFPageRef{776} /XYZ 72 384.406 null] >>}
\MusicPDFWriteObject{10715}{<< /D [\MusicPDFPageRef{776} /XYZ 72 309.152 null] >>}
\MusicPDFWriteObject{10716}{<< /D [\MusicPDFPageRef{778} /XYZ 72 670.809 null] >>}
\MusicPDFWriteObject{10717}{<< /D [\MusicPDFPageRef{778} /XYZ 72 608.08 null] >>}
\MusicPDFWriteObject{10718}{<< /D [\MusicPDFPageRef{779} /XYZ 72 438.859 null] >>}
\MusicPDFWriteObject{10719}{<< /D [\MusicPDFPageRef{779} /XYZ 72 410.066 null] >>}
\MusicPDFWriteObject{10720}{<< /D [\MusicPDFPageRef{784} /XYZ 72 665.035 null] >>}
\MusicPDFWriteObject{10721}{<< /D [\MusicPDFPageRef{784} /XYZ 72 636.387 null] >>}
\MusicPDFWriteObject{10722}{<< /D [\MusicPDFPageRef{784} /XYZ 72 611.518 null] >>}
\MusicPDFWriteObject{10723}{<< /D [\MusicPDFPageRef{119} /XYZ 72 198.177 null] >>}
\MusicPDFWriteObject{10724}{<< /D [\MusicPDFPageRef{785} /XYZ 72 285.88 null] >>}
\MusicPDFWriteObject{10725}{<< /D [\MusicPDFPageRef{785} /XYZ 72 248.124 null] >>}
\MusicPDFWriteObject{10726}{<< /D [\MusicPDFPageRef{790} /XYZ 72 226.275 null] >>}
\MusicPDFWriteObject{10727}{<< /D [\MusicPDFPageRef{790} /XYZ 72 168.018 null] >>}
\MusicPDFWriteObject{10728}{<< /D [\MusicPDFPageRef{790} /XYZ 72 157.059 null] >>}
\MusicPDFWriteObject{10729}{<< /D [\MusicPDFPageRef{791} /XYZ 72 172.357 null] >>}
\MusicPDFWriteObject{10730}{<< /D [\MusicPDFPageRef{791} /XYZ 72 94.55 null] >>}
\MusicPDFWriteObject{10731}{<< /D [\MusicPDFPageRef{792} /XYZ 72 698.082 null] >>}
\MusicPDFWriteObject{10732}{<< /D [\MusicPDFPageRef{793} /XYZ 72 597.847 null] >>}
\MusicPDFWriteObject{10733}{<< /D [\MusicPDFPageRef{793} /XYZ 72 584.696 null] >>}
\MusicPDFWriteObject{10734}{<< /D [\MusicPDFPageRef{119} /XYZ 72 184.473 null] >>}
\MusicPDFWriteObject{10735}{<< /D [\MusicPDFPageRef{795} /XYZ 72 698.082 null] >>}
\MusicPDFWriteObject{10736}{<< /D [\MusicPDFPageRef{795} /XYZ 72 647.873 null] >>}
\MusicPDFWriteObject{10737}{<< /D [\MusicPDFPageRef{795} /XYZ 72 636.051 null] >>}
\MusicPDFWriteObject{10738}{<< /D [\MusicPDFPageRef{795} /XYZ 72 622.9 null] >>}
\MusicPDFWriteObject{10739}{<< /D [\MusicPDFPageRef{796} /XYZ 72 485.734 null] >>}
\MusicPDFWriteObject{10740}{<< /D [\MusicPDFPageRef{796} /XYZ 72 423.559 null] >>}
\MusicPDFWriteObject{10741}{<< /D [\MusicPDFPageRef{797} /XYZ 72 455.581 null] >>}
\MusicPDFWriteObject{10742}{<< /D [\MusicPDFPageRef{797} /XYZ 72 442.43 null] >>}
\MusicPDFWriteObject{10743}{<< /D [\MusicPDFPageRef{801} /XYZ 72 318.759 null] >>}
\MusicPDFWriteObject{10744}{<< /D [\MusicPDFPageRef{801} /XYZ 72 305.054 null] >>}
\MusicPDFWriteObject{10745}{<< /D [\MusicPDFPageRef{123} /XYZ 72 345.121 null] >>}
\MusicPDFWriteObject{10746}{<< /D [\MusicPDFPageRef{801} /XYZ 72 290.216 null] >>}
\MusicPDFWriteObject{10747}{<< /D [\MusicPDFPageRef{801} /XYZ 72 244.815 null] >>}
\MusicPDFWriteObject{10748}{<< /D [\MusicPDFPageRef{801} /XYZ 72 229.174 null] >>}
\MusicPDFWriteObject{10749}{<< /D [\MusicPDFPageRef{801} /XYZ 72 216.023 null] >>}
\MusicPDFWriteObject{10750}{<< /D [\MusicPDFPageRef{803} /XYZ 72 398.081 null] >>}
\MusicPDFWriteObject{10751}{<< /D [\MusicPDFPageRef{803} /XYZ 72 384.93 null] >>}
\MusicPDFWriteObject{10752}{<< /D [\MusicPDFPageRef{803} /XYZ 72 346.392 null] >>}
\MusicPDFWriteObject{10753}{<< /D [\MusicPDFPageRef{804} /XYZ 72 328.573 null] >>}
\MusicPDFWriteObject{10754}{<< /D [\MusicPDFPageRef{804} /XYZ 72 302.826 null] >>}
\MusicPDFWriteObject{10755}{<< /D [\MusicPDFPageRef{804} /XYZ 72 276.524 null] >>}
\MusicPDFWriteObject{10756}{<< /D [\MusicPDFPageRef{123} /XYZ 72 329.37 null] >>}
\MusicPDFWriteObject{10757}{<< /D [\MusicPDFPageRef{806} /XYZ 72 362.018 null] >>}
\MusicPDFWriteObject{10758}{<< /D [\MusicPDFPageRef{806} /XYZ 72 144.405 null] >>}
\MusicPDFWriteObject{10759}{<< /D [\MusicPDFPageRef{807} /XYZ 72 698.082 null] >>}
\MusicPDFWriteObject{10760}{<< /D [\MusicPDFPageRef{807} /XYZ 72 683.942 null] >>}
\MusicPDFWriteObject{10761}{<< /D [\MusicPDFPageRef{807} /XYZ 72 254.174 null] >>}
\MusicPDFWriteObject{10762}{<< /D [\MusicPDFPageRef{807} /XYZ 72 241.023 null] >>}
\MusicPDFWriteObject{10763}{<< /D [\MusicPDFPageRef{807} /XYZ 72 201.318 null] >>}
\MusicPDFWriteObject{10764}{<< /D [\MusicPDFPageRef{818} /XYZ 72 369.655 null] >>}
\MusicPDFWriteObject{10765}{<< /D [\MusicPDFPageRef{818} /XYZ 72 319.523 null] >>}
\MusicPDFWriteObject{10766}{<< /D [\MusicPDFPageRef{818} /XYZ 72 244.197 null] >>}
\MusicPDFWriteObject{10767}{<< /D [\MusicPDFPageRef{123} /XYZ 72 302.638 null] >>}
\MusicPDFWriteObject{10768}{<< /D [\MusicPDFPageRef{824} /XYZ 72 576.539 null] >>}
\MusicPDFWriteObject{10769}{<< /D [\MusicPDFPageRef{824} /XYZ 72 560.899 null] >>}
\MusicPDFWriteObject{10770}{<< /D [\MusicPDFPageRef{824} /XYZ 72 535.15 null] >>}
\MusicPDFWriteObject{10771}{<< /D [\MusicPDFPageRef{824} /XYZ 72 521.446 null] >>}
\MusicPDFWriteObject{10772}{<< /D [\MusicPDFPageRef{824} /XYZ 72 491.261 null] >>}
\MusicPDFWriteObject{10773}{<< /D [\MusicPDFPageRef{824} /XYZ 72 449.872 null] >>}
\MusicPDFWriteObject{10774}{<< /D [\MusicPDFPageRef{824} /XYZ 72 423.57 null] >>}
\MusicPDFWriteObject{10775}{<< /D [\MusicPDFPageRef{824} /XYZ 72 399.206 null] >>}
\MusicPDFWriteObject{10776}{<< /D [\MusicPDFPageRef{824} /XYZ 72 357.263 null] >>}
\MusicPDFWriteObject{10777}{<< /D [\MusicPDFPageRef{824} /XYZ 72 344.112 null] >>}
\MusicPDFWriteObject{10778}{<< /D [\MusicPDFPageRef{123} /XYZ 72 240.96 null] >>}
\MusicPDFWriteObject{10779}{<< /D [\MusicPDFPageRef{824} /XYZ 72 110.447 null] >>}
\MusicPDFWriteObject{10780}{<< /D [\MusicPDFPageRef{825} /XYZ 72 672.479 null] >>}
\MusicPDFWriteObject{10781}{<< /D [\MusicPDFPageRef{825} /XYZ 72 646.731 null] >>}
\MusicPDFWriteObject{10782}{<< /D [\MusicPDFPageRef{825} /XYZ 72 620.429 null] >>}
\MusicPDFWriteObject{10783}{<< /D [\MusicPDFPageRef{825} /XYZ 72 606.725 null] >>}
\MusicPDFWriteObject{10784}{<< /D [\MusicPDFPageRef{825} /XYZ 72 593.574 null] >>}
\MusicPDFWriteObject{10785}{<< /D [\MusicPDFPageRef{826} /XYZ 72 523.247 null] >>}
\MusicPDFWriteObject{10786}{<< /D [\MusicPDFPageRef{826} /XYZ 72 496.392 null] >>}
\MusicPDFWriteObject{10787}{<< /D [\MusicPDFPageRef{826} /XYZ 72 470.091 null] >>}
\MusicPDFWriteObject{10788}{<< /D [\MusicPDFPageRef{826} /XYZ 72 456.94 null] >>}
\MusicPDFWriteObject{10789}{<< /D [\MusicPDFPageRef{23} /XYZ 72 624.238 null] >>}
\MusicPDFWriteObject{10790}{<< /D [\MusicPDFPageRef{124} /XYZ 72 317.888 null] >>}
\MusicPDFWriteObject{10791}{<< /D [\MusicPDFPageRef{826} /XYZ 72 430.638 null] >>}
\MusicPDFWriteObject{10792}{<< /D [\MusicPDFPageRef{827} /XYZ 72 136.749 null] >>}
\MusicPDFWriteObject{10793}{<< /D [\MusicPDFPageRef{827} /XYZ 72 123.044 null] >>}
\MusicPDFWriteObject{10794}{<< /D [\MusicPDFPageRef{827} /XYZ 72 109.894 null] >>}
\MusicPDFWriteObject{10795}{<< /D [\MusicPDFPageRef{827} /XYZ 72 96.743 null] >>}
\MusicPDFWriteObject{10796}{<< /D [\MusicPDFPageRef{828} /XYZ 72 698.082 null] >>}
\MusicPDFWriteObject{10797}{<< /D [\MusicPDFPageRef{828} /XYZ 72 439.332 null] >>}
\MusicPDFWriteObject{10798}{<< /D [\MusicPDFPageRef{828} /XYZ 72 412.476 null] >>}
\MusicPDFWriteObject{10799}{<< /D [\MusicPDFPageRef{828} /XYZ 72 349.722 null] >>}
\MusicPDFWriteObject{10800}{<< /D [\MusicPDFPageRef{828} /XYZ 72 296.565 null] >>}
\MusicPDFWriteObject{10801}{<< /D [\MusicPDFPageRef{124} /XYZ 72 245.124 null] >>}
\MusicPDFWriteObject{10802}{<< /D [\MusicPDFPageRef{828} /XYZ 72 226.101 null] >>}
\MusicPDFWriteObject{10803}{<< /D [\MusicPDFPageRef{828} /XYZ 72 212.397 null] >>}
\MusicPDFWriteObject{10804}{<< /D [\MusicPDFPageRef{828} /XYZ 72 199.246 null] >>}
\MusicPDFWriteObject{10805}{<< /D [\MusicPDFPageRef{828} /XYZ 72 175.435 null] >>}
\MusicPDFWriteObject{10806}{<< /D [\MusicPDFPageRef{828} /XYZ 72 147.196 null] >>}
\MusicPDFWriteObject{10807}{<< /D [\MusicPDFPageRef{828} /XYZ 72 133.493 null] >>}
\MusicPDFWriteObject{10808}{<< /D [\MusicPDFPageRef{829} /XYZ 72 698.082 null] >>}
\MusicPDFWriteObject{10809}{<< /D [\MusicPDFPageRef{829} /XYZ 72 672.478 null] >>}
\MusicPDFWriteObject{10810}{<< /D [\MusicPDFPageRef{829} /XYZ 72 602.411 null] >>}
\MusicPDFWriteObject{10811}{<< /D [\MusicPDFPageRef{829} /XYZ 72 576.662 null] >>}
\MusicPDFWriteObject{10812}{<< /D [\MusicPDFPageRef{124} /XYZ 72 218.409 null] >>}
\MusicPDFWriteObject{10813}{<< /D [\MusicPDFPageRef{829} /XYZ 72 562.958 null] >>}
\MusicPDFWriteObject{10814}{<< /D [\MusicPDFPageRef{829} /XYZ 72 549.807 null] >>}
\MusicPDFWriteObject{10815}{<< /D [\MusicPDFPageRef{829} /XYZ 72 536.657 null] >>}
\MusicPDFWriteObject{10816}{<< /D [\MusicPDFPageRef{829} /XYZ 72 485.485 null] >>}
\MusicPDFWriteObject{10817}{<< /D [\MusicPDFPageRef{829} /XYZ 72 459.183 null] >>}
\MusicPDFWriteObject{10818}{<< /D [\MusicPDFPageRef{829} /XYZ 72 446.033 null] >>}
\MusicPDFWriteObject{10819}{<< /D [\MusicPDFPageRef{832} /XYZ 72 320.762 null] >>}
\MusicPDFWriteObject{10820}{<< /D [\MusicPDFPageRef{832} /XYZ 72 294.46 null] >>}
\MusicPDFWriteObject{10821}{<< /D [\MusicPDFPageRef{833} /XYZ 72 404.511 null] >>}
\MusicPDFWriteObject{10822}{<< /D [\MusicPDFPageRef{833} /XYZ 72 377.656 null] >>}
\MusicPDFWriteObject{10823}{<< /D [\MusicPDFPageRef{124} /XYZ 72 182.075 null] >>}
\MusicPDFWriteObject{10824}{<< /D [\MusicPDFPageRef{833} /XYZ 72 196.225 null] >>}
\MusicPDFWriteObject{10825}{<< /D [\MusicPDFPageRef{833} /XYZ 72 142.993 null] >>}
\MusicPDFWriteObject{10826}{<< /D [\MusicPDFPageRef{834} /XYZ 72 410.644 null] >>}
\MusicPDFWriteObject{10827}{<< /D [\MusicPDFPageRef{834} /XYZ 72 396.94 null] >>}
\MusicPDFWriteObject{10828}{<< /D [\MusicPDFPageRef{834} /XYZ 72 371.191 null] >>}
\MusicPDFWriteObject{10829}{<< /D [\MusicPDFPageRef{834} /XYZ 72 292.938 null] >>}
\MusicPDFWriteObject{10830}{<< /D [\MusicPDFPageRef{834} /XYZ 72 247.018 null] >>}
\MusicPDFWriteObject{10831}{<< /D [\MusicPDFPageRef{834} /XYZ 72 193.363 null] >>}
\MusicPDFWriteObject{10832}{<< /D [\MusicPDFPageRef{835} /XYZ 72 389.181 null] >>}
\MusicPDFWriteObject{10833}{<< /D [\MusicPDFPageRef{835} /XYZ 72 286.999 null] >>}
\MusicPDFWriteObject{10834}{<< /D [\MusicPDFPageRef{124} /XYZ 72 110.054 null] >>}
\MusicPDFWriteObject{10835}{<< /D [\MusicPDFPageRef{841} /XYZ 72 371.734 null] >>}
\MusicPDFWriteObject{10836}{<< /D [\MusicPDFPageRef{841} /XYZ 72 344.879 null] >>}
\MusicPDFWriteObject{10837}{<< /D [\MusicPDFPageRef{841} /XYZ 72 331.368 null] >>}
\MusicPDFWriteObject{10838}{<< /D [\MusicPDFPageRef{843} /XYZ 72 686.027 null] >>}
\MusicPDFWriteObject{10839}{<< /D [\MusicPDFPageRef{843} /XYZ 72 670.387 null] >>}
\MusicPDFWriteObject{10840}{<< /D [\MusicPDFPageRef{843} /XYZ 72 257.047 null] >>}
\MusicPDFWriteObject{10841}{<< /D [\MusicPDFPageRef{843} /XYZ 72 241.407 null] >>}
\MusicPDFWriteObject{10842}{<< /D [\MusicPDFPageRef{844} /XYZ 72 358.284 null] >>}
\MusicPDFWriteObject{10843}{<< /D [\MusicPDFPageRef{844} /XYZ 72 345.134 null] >>}
\MusicPDFWriteObject{10844}{<< /D [\MusicPDFPageRef{844} /XYZ 72 218.71 null] >>}
\MusicPDFWriteObject{10845}{<< /D [\MusicPDFPageRef{125} /XYZ 72 698.082 null] >>}
\MusicPDFWriteObject{10846}{<< /D [\MusicPDFPageRef{844} /XYZ 72 205.006 null] >>}
\MusicPDFWriteObject{10847}{<< /D [\MusicPDFPageRef{844} /XYZ 72 191.855 null] >>}
\MusicPDFWriteObject{10848}{<< /D [\MusicPDFPageRef{845} /XYZ 72 348.252 null] >>}
\MusicPDFWriteObject{10849}{<< /D [\MusicPDFPageRef{845} /XYZ 72 334.548 null] >>}
\MusicPDFWriteObject{10850}{<< /D [\MusicPDFPageRef{846} /XYZ 72 152.027 null] >>}
\MusicPDFWriteObject{10851}{<< /D [\MusicPDFPageRef{846} /XYZ 72 138.323 null] >>}
\MusicPDFWriteObject{10852}{<< /D [\MusicPDFPageRef{846} /XYZ 72 125.172 null] >>}
\MusicPDFWriteObject{10853}{<< /D [\MusicPDFPageRef{848} /XYZ 72 377.488 null] >>}
\MusicPDFWriteObject{10854}{<< /D [\MusicPDFPageRef{848} /XYZ 72 363.783 null] >>}
\MusicPDFWriteObject{10855}{<< /D [\MusicPDFPageRef{849} /XYZ 72 698.082 null] >>}
\MusicPDFWriteObject{10856}{<< /D [\MusicPDFPageRef{125} /XYZ 72 646.177 null] >>}
\MusicPDFWriteObject{10857}{<< /D [\MusicPDFPageRef{849} /XYZ 72 628.225 null] >>}
\MusicPDFWriteObject{10858}{<< /D [\MusicPDFPageRef{850} /XYZ 72 570.224 null] >>}
\MusicPDFWriteObject{10859}{<< /D [\MusicPDFPageRef{850} /XYZ 72 557.371 null] >>}
\MusicPDFWriteObject{10860}{<< /D [\MusicPDFPageRef{850} /XYZ 72 325.78 null] >>}
\MusicPDFWriteObject{10861}{<< /D [\MusicPDFPageRef{850} /XYZ 72 312.076 null] >>}
\MusicPDFWriteObject{10862}{<< /D [\MusicPDFPageRef{852} /XYZ 72 581.724 null] >>}
\MusicPDFWriteObject{10863}{<< /D [\MusicPDFPageRef{852} /XYZ 72 566.083 null] >>}
\MusicPDFWriteObject{10864}{<< /D [\MusicPDFPageRef{852} /XYZ 72 488.286 null] >>}
\MusicPDFWriteObject{10865}{<< /D [\MusicPDFPageRef{852} /XYZ 72 475.388 null] >>}
\MusicPDFWriteObject{10866}{<< /D [\MusicPDFPageRef{852} /XYZ 72 462.238 null] >>}
\MusicPDFWriteObject{10867}{<< /D [\MusicPDFPageRef{125} /XYZ 72 572.115 null] >>}
\MusicPDFWriteObject{10868}{<< /D [\MusicPDFPageRef{852} /XYZ 72 449.087 null] >>}
\MusicPDFWriteObject{10869}{<< /D [\MusicPDFPageRef{853} /XYZ 72 587.501 null] >>}
\MusicPDFWriteObject{10870}{<< /D [\MusicPDFPageRef{853} /XYZ 72 563.136 null] >>}
\MusicPDFWriteObject{10871}{<< /D [\MusicPDFPageRef{853} /XYZ 72 337.991 null] >>}
\MusicPDFWriteObject{10872}{<< /D [\MusicPDFPageRef{853} /XYZ 72 311.689 null] >>}
\MusicPDFWriteObject{10873}{<< /D [\MusicPDFPageRef{853} /XYZ 72 285.942 null] >>}
\MusicPDFWriteObject{10874}{<< /D [\MusicPDFPageRef{853} /XYZ 72 246.489 null] >>}
\MusicPDFWriteObject{10875}{<< /D [\MusicPDFPageRef{854} /XYZ 72 532.635 null] >>}
\MusicPDFWriteObject{10876}{<< /D [\MusicPDFPageRef{854} /XYZ 72 518.432 null] >>}
\MusicPDFWriteObject{10877}{<< /D [\MusicPDFPageRef{855} /XYZ 72 638.418 null] >>}
\MusicPDFWriteObject{10878}{<< /D [\MusicPDFPageRef{125} /XYZ 72 545.653 null] >>}
\MusicPDFWriteObject{10879}{<< /D [\MusicPDFPageRef{855} /XYZ 72 622.778 null] >>}
\MusicPDFWriteObject{10880}{<< /D [\MusicPDFPageRef{855} /XYZ 72 513.57 null] >>}
\MusicPDFWriteObject{10881}{<< /D [\MusicPDFPageRef{855} /XYZ 72 462.904 null] >>}
\MusicPDFWriteObject{10882}{<< /D [\MusicPDFPageRef{855} /XYZ 72 338.056 null] >>}
\MusicPDFWriteObject{10883}{<< /D [\MusicPDFPageRef{855} /XYZ 72 324.352 null] >>}
\MusicPDFWriteObject{10884}{<< /D [\MusicPDFPageRef{856} /XYZ 72 616.083 null] >>}
\MusicPDFWriteObject{10885}{<< /D [\MusicPDFPageRef{856} /XYZ 72 602.379 null] >>}
\MusicPDFWriteObject{10886}{<< /D [\MusicPDFPageRef{856} /XYZ 72 589.228 null] >>}
\MusicPDFWriteObject{10887}{<< /D [\MusicPDFPageRef{856} /XYZ 72 446.651 null] >>}
\MusicPDFWriteObject{10888}{<< /D [\MusicPDFPageRef{856} /XYZ 72 432.947 null] >>}
\MusicPDFWriteObject{10889}{<< /D [\MusicPDFPageRef{129} /XYZ 72 667.635 null] >>}
\MusicPDFWriteObject{10890}{<< /D [\MusicPDFPageRef{856} /XYZ 72 419.796 null] >>}
\MusicPDFWriteObject{10891}{<< /D [\MusicPDFPageRef{856} /XYZ 72 406.646 null] >>}
\MusicPDFWriteObject{10892}{<< /D [\MusicPDFPageRef{856} /XYZ 72 393.495 null] >>}
\MusicPDFWriteObject{10893}{<< /D [\MusicPDFPageRef{856} /XYZ 72 380.344 null] >>}
\MusicPDFWriteObject{10894}{<< /D [\MusicPDFPageRef{857} /XYZ 72 212.742 null] >>}
\MusicPDFWriteObject{10895}{<< /D [\MusicPDFPageRef{857} /XYZ 72 199.039 null] >>}
\MusicPDFWriteObject{10896}{<< /D [\MusicPDFPageRef{857} /XYZ 72 185.888 null] >>}
\MusicPDFWriteObject{10897}{<< /D [\MusicPDFPageRef{857} /XYZ 72 172.737 null] >>}
\MusicPDFWriteObject{10898}{<< /D [\MusicPDFPageRef{857} /XYZ 72 159.587 null] >>}
\MusicPDFWriteObject{10899}{<< /D [\MusicPDFPageRef{857} /XYZ 72 146.436 null] >>}
\MusicPDFWriteObject{10900}{<< /D [\MusicPDFPageRef{23} /XYZ 72 540.265 null] >>}
\MusicPDFWriteObject{10901}{<< /D [\MusicPDFPageRef{129} /XYZ 72 595.558 null] >>}
\MusicPDFWriteObject{10902}{<< /D [\MusicPDFPageRef{858} /XYZ 72 353.196 null] >>}
\MusicPDFWriteObject{10903}{<< /D [\MusicPDFPageRef{858} /XYZ 72 326.341 null] >>}
\MusicPDFWriteObject{10904}{<< /D [\MusicPDFPageRef{858} /XYZ 72 300.593 null] >>}
\MusicPDFWriteObject{10905}{<< /D [\MusicPDFPageRef{859} /XYZ 72 233.357 null] >>}
\MusicPDFWriteObject{10906}{<< /D [\MusicPDFPageRef{859} /XYZ 72 219.293 null] >>}
\MusicPDFWriteObject{10907}{<< /D [\MusicPDFPageRef{862} /XYZ 72 572.649 null] >>}
\MusicPDFWriteObject{10908}{<< /D [\MusicPDFPageRef{862} /XYZ 72 557.009 null] >>}
\MusicPDFWriteObject{10909}{<< /D [\MusicPDFPageRef{862} /XYZ 72 543.858 null] >>}
\MusicPDFWriteObject{10910}{<< /D [\MusicPDFPageRef{862} /XYZ 72 530.707 null] >>}
\MusicPDFWriteObject{10911}{<< /D [\MusicPDFPageRef{862} /XYZ 72 517.556 null] >>}
\MusicPDFWriteObject{10912}{<< /D [\MusicPDFPageRef{129} /XYZ 72 484.632 null] >>}
\MusicPDFWriteObject{10913}{<< /D [\MusicPDFPageRef{862} /XYZ 72 504.406 null] >>}
\MusicPDFWriteObject{10914}{<< /D [\MusicPDFPageRef{862} /XYZ 72 491.255 null] >>}
\MusicPDFWriteObject{10915}{<< /D [\MusicPDFPageRef{862} /XYZ 72 478.104 null] >>}
\MusicPDFWriteObject{10916}{<< /D [\MusicPDFPageRef{862} /XYZ 72 464.953 null] >>}
\MusicPDFWriteObject{10917}{<< /D [\MusicPDFPageRef{862} /XYZ 72 451.803 null] >>}
\MusicPDFWriteObject{10918}{<< /D [\MusicPDFPageRef{862} /XYZ 72 438.652 null] >>}
\MusicPDFWriteObject{10919}{<< /D [\MusicPDFPageRef{862} /XYZ 72 425.501 null] >>}
\MusicPDFWriteObject{10920}{<< /D [\MusicPDFPageRef{862} /XYZ 72 412.35 null] >>}
\MusicPDFWriteObject{10921}{<< /D [\MusicPDFPageRef{862} /XYZ 72 386.602 null] >>}
\MusicPDFWriteObject{10922}{<< /D [\MusicPDFPageRef{863} /XYZ 72 441.306 null] >>}
\MusicPDFWriteObject{10923}{<< /D [\MusicPDFPageRef{129} /XYZ 72 458.884 null] >>}
\MusicPDFWriteObject{10924}{<< /D [\MusicPDFPageRef{863} /XYZ 72 428.155 null] >>}
\MusicPDFWriteObject{10925}{<< /D [\MusicPDFPageRef{863} /XYZ 72 364.385 null] >>}
\MusicPDFWriteObject{10926}{<< /D [\MusicPDFPageRef{863} /XYZ 72 303.593 null] >>}
\MusicPDFWriteObject{10927}{<< /D [\MusicPDFPageRef{865} /XYZ 72 577.781 null] >>}
\MusicPDFWriteObject{10928}{<< /D [\MusicPDFPageRef{865} /XYZ 72 551.479 null] >>}
\MusicPDFWriteObject{10929}{<< /D [\MusicPDFPageRef{865} /XYZ 72 469.313 null] >>}
\MusicPDFWriteObject{10930}{<< /D [\MusicPDFPageRef{865} /XYZ 72 455.609 null] >>}
\MusicPDFWriteObject{10931}{<< /D [\MusicPDFPageRef{865} /XYZ 72 442.458 null] >>}
\MusicPDFWriteObject{10932}{<< /D [\MusicPDFPageRef{865} /XYZ 72 429.308 null] >>}
\MusicPDFWriteObject{10933}{<< /D [\MusicPDFPageRef{865} /XYZ 72 416.157 null] >>}
\MusicPDFWriteObject{10934}{<< /D [\MusicPDFPageRef{129} /XYZ 72 384.981 null] >>}
\MusicPDFWriteObject{10935}{<< /D [\MusicPDFPageRef{865} /XYZ 72 389.855 null] >>}
\MusicPDFWriteObject{10936}{<< /D [\MusicPDFPageRef{866} /XYZ 72 591.435 null] >>}
\MusicPDFWriteObject{10937}{<< /D [\MusicPDFPageRef{866} /XYZ 72 577.731 null] >>}
\MusicPDFWriteObject{10938}{<< /D [\MusicPDFPageRef{867} /XYZ 72 635.694 null] >>}
\MusicPDFWriteObject{10939}{<< /D [\MusicPDFPageRef{867} /XYZ 72 621.99 null] >>}
\MusicPDFWriteObject{10940}{<< /D [\MusicPDFPageRef{867} /XYZ 72 608.839 null] >>}
\MusicPDFWriteObject{10941}{<< /D [\MusicPDFPageRef{867} /XYZ 72 325.608 null] >>}
\MusicPDFWriteObject{10942}{<< /D [\MusicPDFPageRef{867} /XYZ 72 311.903 null] >>}
\MusicPDFWriteObject{10943}{<< /D [\MusicPDFPageRef{867} /XYZ 72 166.192 null] >>}
\MusicPDFWriteObject{10944}{<< /D [\MusicPDFPageRef{867} /XYZ 72 152.488 null] >>}
\MusicPDFWriteObject{10945}{<< /D [\MusicPDFPageRef{132} /XYZ 72 471.963 null] >>}
\MusicPDFWriteObject{10946}{<< /D [\MusicPDFPageRef{867} /XYZ 72 139.337 null] >>}
\MusicPDFWriteObject{10947}{<< /D [\MusicPDFPageRef{868} /XYZ 72 590.099 null] >>}
\MusicPDFWriteObject{10948}{<< /D [\MusicPDFPageRef{868} /XYZ 72 576.395 null] >>}
\MusicPDFWriteObject{10949}{<< /D [\MusicPDFPageRef{868} /XYZ 72 317.65 null] >>}
\MusicPDFWriteObject{10950}{<< /D [\MusicPDFPageRef{868} /XYZ 72 280.226 null] >>}
\MusicPDFWriteObject{10951}{<< /D [\MusicPDFPageRef{869} /XYZ 72 482.419 null] >>}
\MusicPDFWriteObject{10952}{<< /D [\MusicPDFPageRef{869} /XYZ 72 469.269 null] >>}
\MusicPDFWriteObject{10953}{<< /D [\MusicPDFPageRef{869} /XYZ 72 456.118 null] >>}
\MusicPDFWriteObject{10954}{<< /D [\MusicPDFPageRef{869} /XYZ 72 326.665 null] >>}
\MusicPDFWriteObject{10955}{<< /D [\MusicPDFPageRef{869} /XYZ 72 313.514 null] >>}
\MusicPDFWriteObject{10956}{<< /D [\MusicPDFPageRef{132} /XYZ 72 435.735 null] >>}
\MusicPDFWriteObject{10957}{<< /D [\MusicPDFPageRef{871} /XYZ 72 141.067 null] >>}
\MusicPDFWriteObject{10958}{<< /D [\MusicPDFPageRef{871} /XYZ 72 127.917 null] >>}
\MusicPDFWriteObject{10959}{<< /D [\MusicPDFPageRef{872} /XYZ 72 123.598 null] >>}
\MusicPDFWriteObject{10960}{<< /D [\MusicPDFPageRef{873} /XYZ 72 698.082 null] >>}
\MusicPDFWriteObject{10961}{<< /D [\MusicPDFPageRef{873} /XYZ 72 673.032 null] >>}
\MusicPDFWriteObject{10962}{<< /D [\MusicPDFPageRef{873} /XYZ 72 477.025 null] >>}
\MusicPDFWriteObject{10963}{<< /D [\MusicPDFPageRef{873} /XYZ 72 463.321 null] >>}
\MusicPDFWriteObject{10964}{<< /D [\MusicPDFPageRef{873} /XYZ 72 413.742 null] >>}
\MusicPDFWriteObject{10965}{<< /D [\MusicPDFPageRef{873} /XYZ 72 316.731 null] >>}
\MusicPDFWriteObject{10966}{<< /D [\MusicPDFPageRef{873} /XYZ 72 303.58 null] >>}
\MusicPDFWriteObject{10967}{<< /D [\MusicPDFPageRef{132} /XYZ 72 382.576 null] >>}
\MusicPDFWriteObject{10968}{<< /D [\MusicPDFPageRef{873} /XYZ 72 276.919 null] >>}
\MusicPDFWriteObject{10969}{<< /D [\MusicPDFPageRef{877} /XYZ 72 599.052 null] >>}
\MusicPDFWriteObject{10970}{<< /D [\MusicPDFPageRef{877} /XYZ 72 586.927 null] >>}
\MusicPDFWriteObject{10971}{<< /D [\MusicPDFPageRef{877} /XYZ 72 414.033 null] >>}
\MusicPDFWriteObject{10972}{<< /D [\MusicPDFPageRef{877} /XYZ 72 389.668 null] >>}
\MusicPDFWriteObject{10973}{<< /D [\MusicPDFPageRef{877} /XYZ 72 374.027 null] >>}
\MusicPDFWriteObject{10974}{<< /D [\MusicPDFPageRef{877} /XYZ 72 360.877 null] >>}
\MusicPDFWriteObject{10975}{<< /D [\MusicPDFPageRef{877} /XYZ 72 335.128 null] >>}
\MusicPDFWriteObject{10976}{<< /D [\MusicPDFPageRef{878} /XYZ 72 359.032 null] >>}
\MusicPDFWriteObject{10977}{<< /D [\MusicPDFPageRef{878} /XYZ 72 343.391 null] >>}
\MusicPDFWriteObject{10978}{<< /D [\MusicPDFPageRef{135} /XYZ 72 401.965 null] >>}
\MusicPDFWriteObject{10979}{<< /D [\MusicPDFPageRef{878} /XYZ 72 330.24 null] >>}
\MusicPDFWriteObject{10980}{<< /D [\MusicPDFPageRef{878} /XYZ 72 317.089 null] >>}
\MusicPDFWriteObject{10981}{<< /D [\MusicPDFPageRef{878} /XYZ 72 303.939 null] >>}
\MusicPDFWriteObject{10982}{<< /D [\MusicPDFPageRef{878} /XYZ 72 290.788 null] >>}
\MusicPDFWriteObject{10983}{<< /D [\MusicPDFPageRef{878} /XYZ 72 277.637 null] >>}
\MusicPDFWriteObject{10984}{<< /D [\MusicPDFPageRef{878} /XYZ 72 264.486 null] >>}
\MusicPDFWriteObject{10985}{<< /D [\MusicPDFPageRef{878} /XYZ 72 251.336 null] >>}
\MusicPDFWriteObject{10986}{<< /D [\MusicPDFPageRef{878} /XYZ 72 238.185 null] >>}
\MusicPDFWriteObject{10987}{<< /D [\MusicPDFPageRef{878} /XYZ 72 225.034 null] >>}
\MusicPDFWriteObject{10988}{<< /D [\MusicPDFPageRef{878} /XYZ 72 211.883 null] >>}
\MusicPDFWriteObject{10989}{<< /D [\MusicPDFPageRef{135} /XYZ 72 364.122 null] >>}
\MusicPDFWriteObject{10990}{<< /D [\MusicPDFPageRef{878} /XYZ 72 198.733 null] >>}
\MusicPDFWriteObject{10991}{<< /D [\MusicPDFPageRef{878} /XYZ 72 172.984 null] >>}
\MusicPDFWriteObject{10992}{<< /D [\MusicPDFPageRef{879} /XYZ 72 249.082 null] >>}
\MusicPDFWriteObject{10993}{<< /D [\MusicPDFPageRef{879} /XYZ 72 235.378 null] >>}
\MusicPDFWriteObject{10994}{<< /D [\MusicPDFPageRef{879} /XYZ 72 191.727 null] >>}
\MusicPDFWriteObject{10995}{<< /D [\MusicPDFPageRef{879} /XYZ 72 178.023 null] >>}
\MusicPDFWriteObject{10996}{<< /D [\MusicPDFPageRef{880} /XYZ 72 620.497 null] >>}
\MusicPDFWriteObject{10997}{<< /D [\MusicPDFPageRef{880} /XYZ 72 606.793 null] >>}
\MusicPDFWriteObject{10998}{<< /D [\MusicPDFPageRef{880} /XYZ 72 593.642 null] >>}
\MusicPDFWriteObject{10999}{<< /D [\MusicPDFPageRef{880} /XYZ 72 580.491 null] >>}
\MusicPDFWriteObject{11000}{<< /D [\MusicPDFPageRef{135} /XYZ 72 227.97 null] >>}
\MusicPDFWriteObject{11001}{<< /D [\MusicPDFPageRef{880} /XYZ 72 566.012 null] >>}
\MusicPDFWriteObject{11002}{<< /D [\MusicPDFPageRef{880} /XYZ 72 554.19 null] >>}
\MusicPDFWriteObject{11003}{<< /D [\MusicPDFPageRef{880} /XYZ 72 540.679 null] >>}
\MusicPDFWriteObject{11004}{<< /D [\MusicPDFPageRef{880} /XYZ 72 526.075 null] >>}
\MusicPDFWriteObject{11005}{<< /D [\MusicPDFPageRef{880} /XYZ 72 514.737 null] >>}
\MusicPDFWriteObject{11006}{<< /D [\MusicPDFPageRef{880} /XYZ 72 501.587 null] >>}
\MusicPDFWriteObject{11007}{<< /D [\MusicPDFPageRef{880} /XYZ 72 488.436 null] >>}
\MusicPDFWriteObject{11008}{<< /D [\MusicPDFPageRef{880} /XYZ 72 267.499 null] >>}
\MusicPDFWriteObject{11009}{<< /D [\MusicPDFPageRef{880} /XYZ 72 225.556 null] >>}
\MusicPDFWriteObject{11010}{<< /D [\MusicPDFPageRef{880} /XYZ 72 155.001 null] >>}
\MusicPDFWriteObject{11011}{<< /D [\MusicPDFPageRef{23} /XYZ 72 276.334 null] >>}
\MusicPDFWriteObject{11012}{<< /D [\MusicPDFPageRef{135} /XYZ 72 213.991 null] >>}
\MusicPDFWriteObject{11013}{<< /D [\MusicPDFPageRef{881} /XYZ 72 377.402 null] >>}
\MusicPDFWriteObject{11014}{<< /D [\MusicPDFPageRef{881} /XYZ 72 363.892 null] >>}
\MusicPDFWriteObject{11015}{<< /D [\MusicPDFPageRef{881} /XYZ 72 350.741 null] >>}
\MusicPDFWriteObject{11016}{<< /D [\MusicPDFPageRef{881} /XYZ 72 337.59 null] >>}
\MusicPDFWriteObject{11017}{<< /D [\MusicPDFPageRef{881} /XYZ 72 311.649 null] >>}
\MusicPDFWriteObject{11018}{<< /D [\MusicPDFPageRef{881} /XYZ 72 250.47 null] >>}
\MusicPDFWriteObject{11019}{<< /D [\MusicPDFPageRef{881} /XYZ 72 222.232 null] >>}
\MusicPDFWriteObject{11020}{<< /D [\MusicPDFPageRef{881} /XYZ 72 176.201 null] >>}
\MusicPDFWriteObject{11021}{<< /D [\MusicPDFPageRef{882} /XYZ 72 698.082 null] >>}
\MusicPDFWriteObject{11022}{<< /D [\MusicPDFPageRef{882} /XYZ 72 606.725 null] >>}
\MusicPDFWriteObject{11023}{<< /D [\MusicPDFPageRef{135} /XYZ 72 154.408 null] >>}
\MusicPDFWriteObject{11024}{<< /D [\MusicPDFPageRef{882} /XYZ 72 495.553 null] >>}
\MusicPDFWriteObject{11025}{<< /D [\MusicPDFPageRef{882} /XYZ 72 428.886 null] >>}
\MusicPDFWriteObject{11026}{<< /D [\MusicPDFPageRef{883} /XYZ 72 591.529 null] >>}
\MusicPDFWriteObject{11027}{<< /D [\MusicPDFPageRef{883} /XYZ 72 577.465 null] >>}
\MusicPDFWriteObject{11028}{<< /D [\MusicPDFPageRef{883} /XYZ 72 314.553 null] >>}
\MusicPDFWriteObject{11029}{<< /D [\MusicPDFPageRef{883} /XYZ 72 301.402 null] >>}
\MusicPDFWriteObject{11030}{<< /D [\MusicPDFPageRef{883} /XYZ 72 288.252 null] >>}
\MusicPDFWriteObject{11031}{<< /D [\MusicPDFPageRef{883} /XYZ 72 275.101 null] >>}
\MusicPDFWriteObject{11032}{<< /D [\MusicPDFPageRef{883} /XYZ 72 233.67 null] >>}
\MusicPDFWriteObject{11033}{<< /D [\MusicPDFPageRef{883} /XYZ 72 219.966 null] >>}
\MusicPDFWriteObject{11034}{<< /D [\MusicPDFPageRef{135} /XYZ 72 141.258 null] >>}
\MusicPDFWriteObject{11035}{<< /D [\MusicPDFPageRef{883} /XYZ 72 189.746 null] >>}
\MusicPDFWriteObject{11036}{<< /D [\MusicPDFPageRef{883} /XYZ 72 176.042 null] >>}
\MusicPDFWriteObject{11037}{<< /D [\MusicPDFPageRef{884} /XYZ 72 698.082 null] >>}
\MusicPDFWriteObject{11038}{<< /D [\MusicPDFPageRef{884} /XYZ 72 661.992 null] >>}
\MusicPDFWriteObject{11039}{<< /D [\MusicPDFPageRef{884} /XYZ 72 649.202 null] >>}
\MusicPDFWriteObject{11040}{<< /D [\MusicPDFPageRef{884} /XYZ 72 612.774 null] >>}
\MusicPDFWriteObject{11041}{<< /D [\MusicPDFPageRef{884} /XYZ 72 576.346 null] >>}
\MusicPDFWriteObject{11042}{<< /D [\MusicPDFPageRef{884} /XYZ 72 563.196 null] >>}
\MusicPDFWriteObject{11043}{<< /D [\MusicPDFPageRef{884} /XYZ 72 526.768 null] >>}
\MusicPDFWriteObject{11044}{<< /D [\MusicPDFPageRef{884} /XYZ 72 410.519 null] >>}
\MusicPDFWriteObject{11045}{<< /D [\MusicPDFPageRef{135} /XYZ 72 115.51 null] >>}
\MusicPDFWriteObject{11046}{<< /D [\MusicPDFPageRef{884} /XYZ 72 334.25 null] >>}
\MusicPDFWriteObject{11047}{<< /D [\MusicPDFPageRef{884} /XYZ 72 247.044 null] >>}
\MusicPDFWriteObject{11048}{<< /D [\MusicPDFPageRef{884} /XYZ 72 162.039 null] >>}
\MusicPDFWriteObject{11049}{<< /D [\MusicPDFPageRef{885} /XYZ 72 587.233 null] >>}
\MusicPDFWriteObject{11050}{<< /D [\MusicPDFPageRef{885} /XYZ 72 481.473 null] >>}
\MusicPDFWriteObject{11051}{<< /D [\MusicPDFPageRef{885} /XYZ 72 301.641 null] >>}
\MusicPDFWriteObject{11052}{<< /D [\MusicPDFPageRef{885} /XYZ 72 287.937 null] >>}
\MusicPDFWriteObject{11053}{<< /D [\MusicPDFPageRef{885} /XYZ 72 274.787 null] >>}
\MusicPDFWriteObject{11054}{<< /D [\MusicPDFPageRef{885} /XYZ 72 261.636 null] >>}
\MusicPDFWriteObject{11055}{<< /D [\MusicPDFPageRef{885} /XYZ 72 248.485 null] >>}
\MusicPDFWriteObject{11056}{<< /D [\MusicPDFPageRef{136} /XYZ 72 607.673 null] >>}
\MusicPDFWriteObject{11057}{<< /D [\MusicPDFPageRef{885} /XYZ 72 139.653 null] >>}
\MusicPDFWriteObject{11058}{<< /D [\MusicPDFPageRef{885} /XYZ 72 126.143 null] >>}
\MusicPDFWriteObject{11059}{<< /D [\MusicPDFPageRef{887} /XYZ 72 614.937 null] >>}
\MusicPDFWriteObject{11060}{<< /D [\MusicPDFPageRef{887} /XYZ 72 601.786 null] >>}
\MusicPDFWriteObject{11061}{<< /D [\MusicPDFPageRef{887} /XYZ 72 222.754 null] >>}
\MusicPDFWriteObject{11062}{<< /D [\MusicPDFPageRef{887} /XYZ 72 112.384 null] >>}
\MusicPDFWriteObject{11063}{<< /D [\MusicPDFPageRef{887} /XYZ 72 96.743 null] >>}
\MusicPDFWriteObject{11064}{<< /D [\MusicPDFPageRef{888} /XYZ 72 673.032 null] >>}
\MusicPDFWriteObject{11065}{<< /D [\MusicPDFPageRef{890} /XYZ 72 551.13 null] >>}
\MusicPDFWriteObject{11066}{<< /D [\MusicPDFPageRef{890} /XYZ 72 537.427 null] >>}
\MusicPDFWriteObject{11067}{<< /D [\MusicPDFPageRef{136} /XYZ 72 595.595 null] >>}
\MusicPDFWriteObject{11068}{<< /D [\MusicPDFPageRef{894} /XYZ 72 670.387 null] >>}
\MusicPDFWriteObject{11069}{<< /D [\MusicPDFPageRef{894} /XYZ 72 646.575 null] >>}
\MusicPDFWriteObject{11070}{<< /D [\MusicPDFPageRef{894} /XYZ 72 630.574 null] >>}
\MusicPDFWriteObject{11071}{<< /D [\MusicPDFPageRef{894} /XYZ 72 617.423 null] >>}
\MusicPDFWriteObject{11072}{<< /D [\MusicPDFPageRef{894} /XYZ 72 592.035 null] >>}
\MusicPDFWriteObject{11073}{<< /D [\MusicPDFPageRef{894} /XYZ 72 577.971 null] >>}
\MusicPDFWriteObject{11074}{<< /D [\MusicPDFPageRef{895} /XYZ 72 474.907 null] >>}
\MusicPDFWriteObject{11075}{<< /D [\MusicPDFPageRef{895} /XYZ 72 424.775 null] >>}
\MusicPDFWriteObject{11076}{<< /D [\MusicPDFPageRef{895} /XYZ 72 318.292 null] >>}
\MusicPDFWriteObject{11077}{<< /D [\MusicPDFPageRef{895} /XYZ 72 188.365 null] >>}
\MusicPDFWriteObject{11078}{<< /D [\MusicPDFPageRef{136} /XYZ 72 534.115 null] >>}
\MusicPDFWriteObject{11079}{<< /D [\MusicPDFPageRef{896} /XYZ 72 649.755 null] >>}
\MusicPDFWriteObject{11080}{<< /D [\MusicPDFPageRef{898} /XYZ 72 604.569 null] >>}
\MusicPDFWriteObject{11081}{<< /D [\MusicPDFPageRef{898} /XYZ 72 591.418 null] >>}
\MusicPDFWriteObject{11082}{<< /D [\MusicPDFPageRef{898} /XYZ 72 577.215 null] >>}
\MusicPDFWriteObject{11083}{<< /D [\MusicPDFPageRef{898} /XYZ 72 563.533 null] >>}
\MusicPDFWriteObject{11084}{<< /D [\MusicPDFPageRef{898} /XYZ 72 551.966 null] >>}
\MusicPDFWriteObject{11085}{<< /D [\MusicPDFPageRef{898} /XYZ 72 538.455 null] >>}
\MusicPDFWriteObject{11086}{<< /D [\MusicPDFPageRef{898} /XYZ 72 456.415 null] >>}
\MusicPDFWriteObject{11087}{<< /D [\MusicPDFPageRef{898} /XYZ 72 403.812 null] >>}
\MusicPDFWriteObject{11088}{<< /D [\MusicPDFPageRef{898} /XYZ 72 110.447 null] >>}
\MusicPDFWriteObject{11089}{<< /D [\MusicPDFPageRef{138} /XYZ 72 447.904 null] >>}
\MusicPDFWriteObject{11090}{<< /D [\MusicPDFPageRef{898} /XYZ 72 96.742 null] >>}
\MusicPDFWriteObject{11091}{<< /D [\MusicPDFPageRef{899} /XYZ 72 698.082 null] >>}
\MusicPDFWriteObject{11092}{<< /D [\MusicPDFPageRef{899} /XYZ 72 685.629 null] >>}
\MusicPDFWriteObject{11093}{<< /D [\MusicPDFPageRef{899} /XYZ 72 659.881 null] >>}
\MusicPDFWriteObject{11094}{<< /D [\MusicPDFPageRef{899} /XYZ 72 633.58 null] >>}
\MusicPDFWriteObject{11095}{<< /D [\MusicPDFPageRef{899} /XYZ 72 584.001 null] >>}
\MusicPDFWriteObject{11096}{<< /D [\MusicPDFPageRef{899} /XYZ 72 569.937 null] >>}
\MusicPDFWriteObject{11097}{<< /D [\MusicPDFPageRef{899} /XYZ 72 556.786 null] >>}
\MusicPDFWriteObject{11098}{<< /D [\MusicPDFPageRef{899} /XYZ 72 487.769 null] >>}
\MusicPDFWriteObject{11099}{<< /D [\MusicPDFPageRef{899} /XYZ 72 372.547 null] >>}
\MusicPDFWriteObject{11100}{<< /D [\MusicPDFPageRef{138} /XYZ 72 421.603 null] >>}
\MusicPDFWriteObject{11101}{<< /D [\MusicPDFPageRef{899} /XYZ 72 234.797 null] >>}
\MusicPDFWriteObject{11102}{<< /D [\MusicPDFPageRef{899} /XYZ 72 194.79 null] >>}
\MusicPDFWriteObject{11103}{<< /D [\MusicPDFPageRef{899} /XYZ 72 168.49 null] >>}
\MusicPDFWriteObject{11104}{<< /D [\MusicPDFPageRef{900} /XYZ 72 630.495 null] >>}
\MusicPDFWriteObject{11105}{<< /D [\MusicPDFPageRef{900} /XYZ 72 530.978 null] >>}
\MusicPDFWriteObject{11106}{<< /D [\MusicPDFPageRef{900} /XYZ 72 439.282 null] >>}
\MusicPDFWriteObject{11107}{<< /D [\MusicPDFPageRef{900} /XYZ 72 412.981 null] >>}
\MusicPDFWriteObject{11108}{<< /D [\MusicPDFPageRef{900} /XYZ 72 230.288 null] >>}
\MusicPDFWriteObject{11109}{<< /D [\MusicPDFPageRef{900} /XYZ 72 181.263 null] >>}
\MusicPDFWriteObject{11110}{<< /D [\MusicPDFPageRef{900} /XYZ 72 130.597 null] >>}
\MusicPDFWriteObject{11111}{<< /D [\MusicPDFPageRef{138} /XYZ 72 347.834 null] >>}
\MusicPDFWriteObject{11112}{<< /D [\MusicPDFPageRef{901} /XYZ 72 685.629 null] >>}
\MusicPDFWriteObject{11113}{<< /D [\MusicPDFPageRef{901} /XYZ 72 511.886 null] >>}
\MusicPDFWriteObject{11114}{<< /D [\MusicPDFPageRef{901} /XYZ 72 438.173 null] >>}
\MusicPDFWriteObject{11115}{<< /D [\MusicPDFPageRef{901} /XYZ 72 411.872 null] >>}
\MusicPDFWriteObject{11116}{<< /D [\MusicPDFPageRef{901} /XYZ 72 385.57 null] >>}
\MusicPDFWriteObject{11117}{<< /D [\MusicPDFPageRef{901} /XYZ 72 241.955 null] >>}
\MusicPDFWriteObject{11118}{<< /D [\MusicPDFPageRef{901} /XYZ 72 201.948 null] >>}
\MusicPDFWriteObject{11119}{<< /D [\MusicPDFPageRef{901} /XYZ 72 175.647 null] >>}
\MusicPDFWriteObject{11120}{<< /D [\MusicPDFPageRef{901} /XYZ 72 123.044 null] >>}
\MusicPDFWriteObject{11121}{<< /D [\MusicPDFPageRef{902} /XYZ 72 698.082 null] >>}
\MusicPDFWriteObject{11122}{<< /D [\MusicPDFPageRef{23} /XYZ 72 228.948 null] >>}
\MusicPDFWriteObject{11123}{<< /D [\MusicPDFPageRef{138} /XYZ 72 274.979 null] >>}
\MusicPDFWriteObject{11124}{<< /D [\MusicPDFPageRef{902} /XYZ 72 471.632 null] >>}
\MusicPDFWriteObject{11125}{<< /D [\MusicPDFPageRef{902} /XYZ 72 421.5 null] >>}
\MusicPDFWriteObject{11126}{<< /D [\MusicPDFPageRef{902} /XYZ 72 407.298 null] >>}
\MusicPDFWriteObject{11127}{<< /D [\MusicPDFPageRef{902} /XYZ 72 345.62 null] >>}
\MusicPDFWriteObject{11128}{<< /D [\MusicPDFPageRef{902} /XYZ 72 267.27 null] >>}
\MusicPDFWriteObject{11129}{<< /D [\MusicPDFPageRef{902} /XYZ 72 219.628 null] >>}
\MusicPDFWriteObject{11130}{<< /D [\MusicPDFPageRef{902} /XYZ 72 168.113 null] >>}
\MusicPDFWriteObject{11131}{<< /D [\MusicPDFPageRef{902} /XYZ 72 114.596 null] >>}
\MusicPDFWriteObject{11132}{<< /D [\MusicPDFPageRef{903} /XYZ 72 356.079 null] >>}
\MusicPDFWriteObject{11133}{<< /D [\MusicPDFPageRef{903} /XYZ 72 342.374 null] >>}
\MusicPDFWriteObject{11134}{<< /D [\MusicPDFPageRef{138} /XYZ 72 262.188 null] >>}
\MusicPDFWriteObject{11135}{<< /D [\MusicPDFPageRef{903} /XYZ 72 305.587 null] >>}
\MusicPDFWriteObject{11136}{<< /D [\MusicPDFPageRef{903} /XYZ 72 205.521 null] >>}
\MusicPDFWriteObject{11137}{<< /D [\MusicPDFPageRef{904} /XYZ 72 698.082 null] >>}
\MusicPDFWriteObject{11138}{<< /D [\MusicPDFPageRef{904} /XYZ 72 499.942 null] >>}
\MusicPDFWriteObject{11139}{<< /D [\MusicPDFPageRef{905} /XYZ 72 344.305 null] >>}
\MusicPDFWriteObject{11140}{<< /D [\MusicPDFPageRef{905} /XYZ 72 303.939 null] >>}
\MusicPDFWriteObject{11141}{<< /D [\MusicPDFPageRef{905} /XYZ 72 230.946 null] >>}
\MusicPDFWriteObject{11142}{<< /D [\MusicPDFPageRef{905} /XYZ 72 191.493 null] >>}
\MusicPDFWriteObject{11143}{<< /D [\MusicPDFPageRef{906} /XYZ 72 604.353 null] >>}
\MusicPDFWriteObject{11144}{<< /D [\MusicPDFPageRef{906} /XYZ 72 591.202 null] >>}
\MusicPDFWriteObject{11145}{<< /D [\MusicPDFPageRef{138} /XYZ 72 158.779 null] >>}
\MusicPDFWriteObject{11146}{<< /D [\MusicPDFPageRef{906} /XYZ 72 578.051 null] >>}
\MusicPDFWriteObject{11147}{<< /D [\MusicPDFPageRef{906} /XYZ 72 563.849 null] >>}
\MusicPDFWriteObject{11148}{<< /D [\MusicPDFPageRef{906} /XYZ 72 521.26 null] >>}
\MusicPDFWriteObject{11149}{<< /D [\MusicPDFPageRef{906} /XYZ 72 507.195 null] >>}
\MusicPDFWriteObject{11150}{<< /D [\MusicPDFPageRef{906} /XYZ 72 441.243 null] >>}
\MusicPDFWriteObject{11151}{<< /D [\MusicPDFPageRef{906} /XYZ 72 332.037 null] >>}
\MusicPDFWriteObject{11152}{<< /D [\MusicPDFPageRef{906} /XYZ 72 284.396 null] >>}
\MusicPDFWriteObject{11153}{<< /D [\MusicPDFPageRef{906} /XYZ 72 232.88 null] >>}
\MusicPDFWriteObject{11154}{<< /D [\MusicPDFPageRef{906} /XYZ 72 183.302 null] >>}
\MusicPDFWriteObject{11155}{<< /D [\MusicPDFPageRef{906} /XYZ 72 133.724 null] >>}
\MusicPDFWriteObject{11156}{<< /D [\MusicPDFPageRef{146} /XYZ 72 491.354 null] >>}
\MusicPDFWriteObject{11157}{<< /D [\MusicPDFPageRef{907} /XYZ 72 698.082 null] >>}
\MusicPDFWriteObject{11158}{<< /D [\MusicPDFPageRef{907} /XYZ 72 672.478 null] >>}
\MusicPDFWriteObject{11159}{<< /D [\MusicPDFPageRef{907} /XYZ 72 646.177 null] >>}
\MusicPDFWriteObject{11160}{<< /D [\MusicPDFPageRef{907} /XYZ 72 633.026 null] >>}
\MusicPDFWriteObject{11161}{<< /D [\MusicPDFPageRef{907} /XYZ 72 565.013 null] >>}
\MusicPDFWriteObject{11162}{<< /D [\MusicPDFPageRef{907} /XYZ 72 551.309 null] >>}
\MusicPDFWriteObject{11163}{<< /D [\MusicPDFPageRef{907} /XYZ 72 538.159 null] >>}
\MusicPDFWriteObject{11164}{<< /D [\MusicPDFPageRef{907} /XYZ 72 525.008 null] >>}
\MusicPDFWriteObject{11165}{<< /D [\MusicPDFPageRef{907} /XYZ 72 498.346 null] >>}
\MusicPDFWriteObject{11166}{<< /D [\MusicPDFPageRef{907} /XYZ 72 485.555 null] >>}
\MusicPDFWriteObject{11167}{<< /D [\MusicPDFPageRef{146} /XYZ 72 428.071 null] >>}
\MusicPDFWriteObject{11168}{<< /D [\MusicPDFPageRef{907} /XYZ 72 472.405 null] >>}
\MusicPDFWriteObject{11169}{<< /D [\MusicPDFPageRef{907} /XYZ 72 448.593 null] >>}
\MusicPDFWriteObject{11170}{<< /D [\MusicPDFPageRef{907} /XYZ 72 363.334 null] >>}
\MusicPDFWriteObject{11171}{<< /D [\MusicPDFPageRef{907} /XYZ 72 350.183 null] >>}
\MusicPDFWriteObject{11172}{<< /D [\MusicPDFPageRef{907} /XYZ 72 336.501 null] >>}
\MusicPDFWriteObject{11173}{<< /D [\MusicPDFPageRef{907} /XYZ 72 307.188 null] >>}
\MusicPDFWriteObject{11174}{<< /D [\MusicPDFPageRef{907} /XYZ 72 292.986 null] >>}
\MusicPDFWriteObject{11175}{<< /D [\MusicPDFPageRef{910} /XYZ 72 154.408 null] >>}
\MusicPDFWriteObject{11176}{<< /D [\MusicPDFPageRef{910} /XYZ 72 141.258 null] >>}
\MusicPDFWriteObject{11177}{<< /D [\MusicPDFPageRef{910} /XYZ 72 128.107 null] >>}
\MusicPDFWriteObject{11178}{<< /D [\MusicPDFPageRef{151} /XYZ 72 268.353 null] >>}
\MusicPDFWriteObject{11179}{<< /D [\MusicPDFPageRef{911} /XYZ 72 698.082 null] >>}
\MusicPDFWriteObject{11180}{<< /D [\MusicPDFPageRef{912} /XYZ 72 312.313 null] >>}
\MusicPDFWriteObject{11181}{<< /D [\MusicPDFPageRef{912} /XYZ 72 298.608 null] >>}
\MusicPDFWriteObject{11182}{<< /D [\MusicPDFPageRef{912} /XYZ 72 251.521 null] >>}
\MusicPDFWriteObject{11183}{<< /D [\MusicPDFPageRef{913} /XYZ 72 479.374 null] >>}
\MusicPDFWriteObject{11184}{<< /D [\MusicPDFPageRef{913} /XYZ 72 465.67 null] >>}
\MusicPDFWriteObject{11185}{<< /D [\MusicPDFPageRef{914} /XYZ 72 373.716 null] >>}
\MusicPDFWriteObject{11186}{<< /D [\MusicPDFPageRef{914} /XYZ 72 360.012 null] >>}
\MusicPDFWriteObject{11187}{<< /D [\MusicPDFPageRef{914} /XYZ 72 346.861 null] >>}
\MusicPDFWriteObject{11188}{<< /D [\MusicPDFPageRef{914} /XYZ 72 333.711 null] >>}
\MusicPDFWriteObject{11189}{<< /D [\MusicPDFPageRef{151} /XYZ 72 190.221 null] >>}
\MusicPDFWriteObject{11190}{<< /D [\MusicPDFPageRef{915} /XYZ 72 670.026 null] >>}
\MusicPDFWriteObject{11191}{<< /D [\MusicPDFPageRef{915} /XYZ 72 644.638 null] >>}
\MusicPDFWriteObject{11192}{<< /D [\MusicPDFPageRef{915} /XYZ 72 618.336 null] >>}
\MusicPDFWriteObject{11193}{<< /D [\MusicPDFPageRef{915} /XYZ 72 604.632 null] >>}
\MusicPDFWriteObject{11194}{<< /D [\MusicPDFPageRef{915} /XYZ 72 175.495 null] >>}
\MusicPDFWriteObject{11195}{<< /D [\MusicPDFPageRef{915} /XYZ 72 161.792 null] >>}
\MusicPDFWriteObject{11196}{<< /D [\MusicPDFPageRef{915} /XYZ 72 148.64 null] >>}
\MusicPDFWriteObject{11197}{<< /D [\MusicPDFPageRef{916} /XYZ 72 625.777 null] >>}
\MusicPDFWriteObject{11198}{<< /D [\MusicPDFPageRef{916} /XYZ 72 612.626 null] >>}
\MusicPDFWriteObject{11199}{<< /D [\MusicPDFPageRef{916} /XYZ 72 586.878 null] >>}
\MusicPDFWriteObject{11200}{<< /D [\MusicPDFPageRef{152} /XYZ 72 606.854 null] >>}
\MusicPDFWriteObject{11201}{<< /D [\MusicPDFPageRef{918} /XYZ 72 670.94 null] >>}
\MusicPDFWriteObject{11202}{<< /D [\MusicPDFPageRef{918} /XYZ 72 655.548 null] >>}
\MusicPDFWriteObject{11203}{<< /D [\MusicPDFPageRef{918} /XYZ 72 334.324 null] >>}
\MusicPDFWriteObject{11204}{<< /D [\MusicPDFPageRef{918} /XYZ 72 318.683 null] >>}
\MusicPDFWriteObject{11205}{<< /D [\MusicPDFPageRef{921} /XYZ 72 530.71 null] >>}
\MusicPDFWriteObject{11206}{<< /D [\MusicPDFPageRef{921} /XYZ 72 505.322 null] >>}
\MusicPDFWriteObject{11207}{<< /D [\MusicPDFPageRef{921} /XYZ 72 293.934 null] >>}
\MusicPDFWriteObject{11208}{<< /D [\MusicPDFPageRef{921} /XYZ 72 226.298 null] >>}
\MusicPDFWriteObject{11209}{<< /D [\MusicPDFPageRef{922} /XYZ 72 426.35 null] >>}
\MusicPDFWriteObject{11210}{<< /D [\MusicPDFPageRef{922} /XYZ 72 400.049 null] >>}
\MusicPDFWriteObject{11211}{<< /D [\MusicPDFPageRef{152} /XYZ 72 594.063 null] >>}
\MusicPDFWriteObject{11212}{<< /D [\MusicPDFPageRef{923} /XYZ 72 630.525 null] >>}
\MusicPDFWriteObject{11213}{<< /D [\MusicPDFPageRef{923} /XYZ 72 579.895 null] >>}
\MusicPDFWriteObject{11214}{<< /D [\MusicPDFPageRef{923} /XYZ 72 530.317 null] >>}
\MusicPDFWriteObject{11215}{<< /D [\MusicPDFPageRef{923} /XYZ 72 442.339 null] >>}
\MusicPDFWriteObject{11216}{<< /D [\MusicPDFPageRef{923} /XYZ 72 366.459 null] >>}
\MusicPDFWriteObject{11217}{<< /D [\MusicPDFPageRef{923} /XYZ 72 261.416 null] >>}
\MusicPDFWriteObject{11218}{<< /D [\MusicPDFPageRef{923} /XYZ 72 248.626 null] >>}
\MusicPDFWriteObject{11219}{<< /D [\MusicPDFPageRef{924} /XYZ 72 299.425 null] >>}
\MusicPDFWriteObject{11220}{<< /D [\MusicPDFPageRef{924} /XYZ 72 285.72 null] >>}
\MusicPDFWriteObject{11221}{<< /D [\MusicPDFPageRef{925} /XYZ 72 171.188 null] >>}
\MusicPDFWriteObject{11222}{<< /D [\MusicPDFPageRef{152} /XYZ 72 556.61 null] >>}
\MusicPDFWriteObject{11223}{<< /D [\MusicPDFPageRef{925} /XYZ 72 156.432 null] >>}
\MusicPDFWriteObject{11224}{<< /D [\MusicPDFPageRef{925} /XYZ 72 143.281 null] >>}
\MusicPDFWriteObject{11225}{<< /D [\MusicPDFPageRef{927} /XYZ 72 683.537 null] >>}
\MusicPDFWriteObject{11226}{<< /D [\MusicPDFPageRef{927} /XYZ 72 670.386 null] >>}
\MusicPDFWriteObject{11227}{<< /D [\MusicPDFPageRef{928} /XYZ 72 97.296 null] >>}
\MusicPDFWriteObject{11228}{<< /D [\MusicPDFPageRef{929} /XYZ 72 698.082 null] >>}
\MusicPDFWriteObject{11229}{<< /D [\MusicPDFPageRef{929} /XYZ 72 504.267 null] >>}
\MusicPDFWriteObject{11230}{<< /D [\MusicPDFPageRef{929} /XYZ 72 490.563 null] >>}
\MusicPDFWriteObject{11231}{<< /D [\MusicPDFPageRef{929} /XYZ 72 477.412 null] >>}
\MusicPDFWriteObject{11232}{<< /D [\MusicPDFPageRef{929} /XYZ 72 435.981 null] >>}
\MusicPDFWriteObject{11233}{<< /D [\MusicPDFPageRef{25} /XYZ 72 600.189 null] >>}
\MusicPDFWriteObject{11234}{<< /D [\MusicPDFPageRef{152} /XYZ 72 519.516 null] >>}
\MusicPDFWriteObject{11235}{<< /D [\MusicPDFPageRef{929} /XYZ 72 211.859 null] >>}
\MusicPDFWriteObject{11236}{<< /D [\MusicPDFPageRef{929} /XYZ 72 198.156 null] >>}
\MusicPDFWriteObject{11237}{<< /D [\MusicPDFPageRef{152} /XYZ 72 506.366 null] >>}
\MusicPDFWriteObject{11238}{<< /D [\MusicPDFPageRef{152} /XYZ 72 462.459 null] >>}
\MusicPDFWriteObject{11239}{<< /D [\MusicPDFPageRef{152} /XYZ 72 420.386 null] >>}
\MusicPDFWriteObject{11240}{<< /D [\MusicPDFPageRef{152} /XYZ 72 109.921 null] >>}
\MusicPDFWriteObject{11241}{<< /D [\MusicPDFPageRef{153} /XYZ 72 698.082 null] >>}
\MusicPDFWriteObject{11242}{<< /D [\MusicPDFPageRef{153} /XYZ 72 659.663 null] >>}
\MusicPDFWriteObject{11243}{<< /D [\MusicPDFPageRef{153} /XYZ 72 621.446 null] >>}
\MusicPDFWriteObject{11244}{<< /D [\MusicPDFPageRef{153} /XYZ 72 524.068 null] >>}
\MusicPDFWriteObject{11245}{<< /D [\MusicPDFPageRef{157} /XYZ 72 451.63 null] >>}
\MusicPDFWriteObject{11246}{<< /D [\MusicPDFPageRef{25} /XYZ 72 586.485 null] >>}
\MusicPDFWriteObject{11247}{<< /D [\MusicPDFPageRef{157} /XYZ 72 392.955 null] >>}
\MusicPDFWriteObject{11248}{<< /D [\MusicPDFPageRef{158} /XYZ 72 282.163 null] >>}
\MusicPDFWriteObject{11249}{<< /D [\MusicPDFPageRef{158} /XYZ 72 255.862 null] >>}
\MusicPDFWriteObject{11250}{<< /D [\MusicPDFPageRef{158} /XYZ 72 242.711 null] >>}
\MusicPDFWriteObject{11251}{<< /D [\MusicPDFPageRef{158} /XYZ 72 181.467 null] >>}
\MusicPDFWriteObject{11252}{<< /D [\MusicPDFPageRef{158} /XYZ 72 155.566 null] >>}
\MusicPDFWriteObject{11253}{<< /D [\MusicPDFPageRef{158} /XYZ 72 142.52 null] >>}
\MusicPDFWriteObject{11254}{<< /D [\MusicPDFPageRef{161} /XYZ 72 570.611 null] >>}
\MusicPDFWriteObject{11255}{<< /D [\MusicPDFPageRef{161} /XYZ 72 557.46 null] >>}
\MusicPDFWriteObject{11256}{<< /D [\MusicPDFPageRef{161} /XYZ 72 544.309 null] >>}
\MusicPDFWriteObject{11257}{<< /D [\MusicPDFPageRef{25} /XYZ 72 573.334 null] >>}
\MusicPDFWriteObject{11258}{<< /D [\MusicPDFPageRef{162} /XYZ 72 342.07 null] >>}
\MusicPDFWriteObject{11259}{<< /D [\MusicPDFPageRef{162} /XYZ 72 328.365 null] >>}
\MusicPDFWriteObject{11260}{<< /D [\MusicPDFPageRef{162} /XYZ 72 215.739 null] >>}
\MusicPDFWriteObject{11261}{<< /D [\MusicPDFPageRef{163} /XYZ 72 123.264 null] >>}
\MusicPDFWriteObject{11262}{<< /D [\MusicPDFPageRef{164} /XYZ 72 698.082 null] >>}
\MusicPDFWriteObject{11263}{<< /D [\MusicPDFPageRef{165} /XYZ 72 345.868 null] >>}
\MusicPDFWriteObject{11264}{<< /D [\MusicPDFPageRef{165} /XYZ 72 329.867 null] >>}
\MusicPDFWriteObject{11265}{<< /D [\MusicPDFPageRef{165} /XYZ 72 316.716 null] >>}
\MusicPDFWriteObject{11266}{<< /D [\MusicPDFPageRef{170} /XYZ 72 144.537 null] >>}
\MusicPDFWriteObject{11267}{<< /D [\MusicPDFPageRef{171} /XYZ 72 698.082 null] >>}
\MusicPDFWriteObject{11268}{<< /D [\MusicPDFPageRef{26} /XYZ 72 584.003 null] >>}
\MusicPDFWriteObject{11269}{<< /D [\MusicPDFPageRef{171} /XYZ 72 540.819 null] >>}
\MusicPDFWriteObject{11270}{<< /D [\MusicPDFPageRef{171} /XYZ 72 464.47 null] >>}
\MusicPDFWriteObject{11271}{<< /D [\MusicPDFPageRef{173} /XYZ 72 471.417 null] >>}
\MusicPDFWriteObject{11272}{<< /D [\MusicPDFPageRef{173} /XYZ 72 414.122 null] >>}
\MusicPDFWriteObject{11273}{<< /D [\MusicPDFPageRef{174} /XYZ 72 448.55 null] >>}
\MusicPDFWriteObject{11274}{<< /D [\MusicPDFPageRef{174} /XYZ 72 434.485 null] >>}
\MusicPDFWriteObject{11275}{<< /D [\MusicPDFPageRef{174} /XYZ 72 421.334 null] >>}
\MusicPDFWriteObject{11276}{<< /D [\MusicPDFPageRef{176} /XYZ 72 652.331 null] >>}
\MusicPDFWriteObject{11277}{<< /D [\MusicPDFPageRef{176} /XYZ 72 639.18 null] >>}
\MusicPDFWriteObject{11278}{<< /D [\MusicPDFPageRef{176} /XYZ 72 588.55 null] >>}
\MusicPDFWriteObject{11279}{<< /D [\MusicPDFPageRef{13} /XYZ 72 229.934 null] >>}
\MusicPDFWriteObject{11280}{<< /D [\MusicPDFPageRef{26} /XYZ 72 570.853 null] >>}
\MusicPDFWriteObject{11281}{<< /D [\MusicPDFPageRef{180} /XYZ 72 110.447 null] >>}
\MusicPDFWriteObject{11282}{<< /D [\MusicPDFPageRef{181} /XYZ 72 698.082 null] >>}
\MusicPDFWriteObject{11283}{<< /D [\MusicPDFPageRef{181} /XYZ 72 657.402 null] >>}
\MusicPDFWriteObject{11284}{<< /D [\MusicPDFPageRef{181} /XYZ 72 544.975 null] >>}
\MusicPDFWriteObject{11285}{<< /D [\MusicPDFPageRef{181} /XYZ 72 455.944 null] >>}
\MusicPDFWriteObject{11286}{<< /D [\MusicPDFPageRef{183} /XYZ 72 287.061 null] >>}
\MusicPDFWriteObject{11287}{<< /D [\MusicPDFPageRef{183} /XYZ 72 213.174 null] >>}
\MusicPDFWriteObject{11288}{<< /D [\MusicPDFPageRef{183} /XYZ 72 125.057 null] >>}
\MusicPDFWriteObject{11289}{<< /D [\MusicPDFPageRef{186} /XYZ 72 577.301 null] >>}
\MusicPDFWriteObject{11290}{<< /D [\MusicPDFPageRef{186} /XYZ 72 537.849 null] >>}
\MusicPDFWriteObject{11291}{<< /D [\MusicPDFPageRef{26} /XYZ 72 557.702 null] >>}
\MusicPDFWriteObject{11292}{<< /D [\MusicPDFPageRef{186} /XYZ 72 488.92 null] >>}
\MusicPDFWriteObject{11293}{<< /D [\MusicPDFPageRef{186} /XYZ 72 376.494 null] >>}
\MusicPDFWriteObject{11294}{<< /D [\MusicPDFPageRef{186} /XYZ 72 349.638 null] >>}
\MusicPDFWriteObject{11295}{<< /D [\MusicPDFPageRef{188} /XYZ 72 526.103 null] >>}
\MusicPDFWriteObject{11296}{<< /D [\MusicPDFPageRef{188} /XYZ 72 452.215 null] >>}
\MusicPDFWriteObject{11297}{<< /D [\MusicPDFPageRef{188} /XYZ 72 425.913 null] >>}
\MusicPDFWriteObject{11298}{<< /D [\MusicPDFPageRef{191} /XYZ 72 571.127 null] >>}
\MusicPDFWriteObject{11299}{<< /D [\MusicPDFPageRef{191} /XYZ 72 557.975 null] >>}
\MusicPDFWriteObject{11300}{<< /D [\MusicPDFPageRef{191} /XYZ 72 531.237 null] >>}
\MusicPDFWriteObject{11301}{<< /D [\MusicPDFPageRef{191} /XYZ 72 453.378 null] >>}
\MusicPDFWriteObject{11302}{<< /D [\MusicPDFPageRef{28} /XYZ 72 452.744 null] >>}
\MusicPDFWriteObject{11303}{<< /D [\MusicPDFPageRef{191} /XYZ 72 411.435 null] >>}
\MusicPDFWriteObject{11304}{<< /D [\MusicPDFPageRef{192} /XYZ 72 441.045 null] >>}
\MusicPDFWriteObject{11305}{<< /D [\MusicPDFPageRef{192} /XYZ 72 427.425 null] >>}
\MusicPDFWriteObject{11306}{<< /D [\MusicPDFPageRef{192} /XYZ 72 388.527 null] >>}
\MusicPDFWriteObject{11307}{<< /D [\MusicPDFPageRef{192} /XYZ 72 374.822 null] >>}
\MusicPDFWriteObject{11308}{<< /D [\MusicPDFPageRef{192} /XYZ 72 347.292 null] >>}
\MusicPDFWriteObject{11309}{<< /D [\MusicPDFPageRef{193} /XYZ 72 247.548 null] >>}
\MusicPDFWriteObject{11310}{<< /D [\MusicPDFPageRef{193} /XYZ 72 97.296 null] >>}
\MusicPDFWriteObject{11311}{<< /D [\MusicPDFPageRef{194} /XYZ 72 698.082 null] >>}
\MusicPDFWriteObject{11312}{<< /D [\MusicPDFPageRef{194} /XYZ 72 644.699 null] >>}
\MusicPDFWriteObject{11313}{<< /D [\MusicPDFPageRef{28} /XYZ 72 439.04 null] >>}
\MusicPDFWriteObject{11314}{<< /D [\MusicPDFPageRef{197} /XYZ 72 163.639 null] >>}
\MusicPDFWriteObject{11315}{<< /D [\MusicPDFPageRef{198} /XYZ 72 659.327 null] >>}
\MusicPDFWriteObject{11316}{<< /D [\MusicPDFPageRef{198} /XYZ 72 645.817 null] >>}
\MusicPDFWriteObject{11317}{<< /D [\MusicPDFPageRef{198} /XYZ 72 632.666 null] >>}
\MusicPDFWriteObject{11318}{<< /D [\MusicPDFPageRef{198} /XYZ 72 606.724 null] >>}
\MusicPDFWriteObject{11319}{<< /D [\MusicPDFPageRef{198} /XYZ 72 580.423 null] >>}
\MusicPDFWriteObject{11320}{<< /D [\MusicPDFPageRef{205} /XYZ 72 308.065 null] >>}
\MusicPDFWriteObject{11321}{<< /D [\MusicPDFPageRef{205} /XYZ 72 241.758 null] >>}
\MusicPDFWriteObject{11322}{<< /D [\MusicPDFPageRef{205} /XYZ 72 152.39 null] >>}
\MusicPDFWriteObject{11323}{<< /D [\MusicPDFPageRef{205} /XYZ 72 125.534 null] >>}
\MusicPDFWriteObject{11324}{<< /D [\MusicPDFPageRef{35} /XYZ 72 585.856 null] >>}
\MusicPDFWriteObject{11325}{<< /D [\MusicPDFPageRef{205} /XYZ 72 99.233 null] >>}
\MusicPDFWriteObject{11326}{<< /D [\MusicPDFPageRef{207} /XYZ 72 349.738 null] >>}
\MusicPDFWriteObject{11327}{<< /D [\MusicPDFPageRef{207} /XYZ 72 264.599 null] >>}
\MusicPDFWriteObject{11328}{<< /D [\MusicPDFPageRef{208} /XYZ 72 286.591 null] >>}
\MusicPDFWriteObject{11329}{<< /D [\MusicPDFPageRef{208} /XYZ 72 247.139 null] >>}
\MusicPDFWriteObject{11330}{<< /D [\MusicPDFPageRef{209} /XYZ 72 477.302 null] >>}
\MusicPDFWriteObject{11331}{<< /D [\MusicPDFPageRef{209} /XYZ 72 463.683 null] >>}
\MusicPDFWriteObject{11332}{<< /D [\MusicPDFPageRef{209} /XYZ 72 424.784 null] >>}
\MusicPDFWriteObject{11333}{<< /D [\MusicPDFPageRef{209} /XYZ 72 385.332 null] >>}
\MusicPDFWriteObject{11334}{<< /D [\MusicPDFPageRef{209} /XYZ 72 357.802 null] >>}
\MusicPDFWriteObject{11335}{<< /D [\MusicPDFPageRef{35} /XYZ 72 558.585 null] >>}
\MusicPDFWriteObject{11336}{<< /D [\MusicPDFPageRef{210} /XYZ 72 448.887 null] >>}
\MusicPDFWriteObject{11337}{<< /D [\MusicPDFPageRef{210} /XYZ 72 349.73 null] >>}
\MusicPDFWriteObject{11338}{<< /D [\MusicPDFPageRef{210} /XYZ 72 300.706 null] >>}
\MusicPDFWriteObject{11339}{<< /D [\MusicPDFPageRef{211} /XYZ 72 353.493 null] >>}
\MusicPDFWriteObject{11340}{<< /D [\MusicPDFPageRef{211} /XYZ 72 340.343 null] >>}
\MusicPDFWriteObject{11341}{<< /D [\MusicPDFPageRef{212} /XYZ 72 199.6 null] >>}
\MusicPDFWriteObject{11342}{<< /D [\MusicPDFPageRef{212} /XYZ 72 135.958 null] >>}
\MusicPDFWriteObject{11343}{<< /D [\MusicPDFPageRef{218} /XYZ 72 338.126 null] >>}
\MusicPDFWriteObject{11344}{<< /D [\MusicPDFPageRef{218} /XYZ 72 324.422 null] >>}
\MusicPDFWriteObject{11345}{<< /D [\MusicPDFPageRef{218} /XYZ 72 311.271 null] >>}
\MusicPDFWriteObject{11346}{<< /D [\MusicPDFPageRef{35} /XYZ 72 520.102 null] >>}
\MusicPDFWriteObject{11347}{<< /D [\MusicPDFPageRef{223} /XYZ 72 598.506 null] >>}
\MusicPDFWriteObject{11348}{<< /D [\MusicPDFPageRef{223} /XYZ 72 495.77 null] >>}
\MusicPDFWriteObject{11349}{<< /D [\MusicPDFPageRef{223} /XYZ 72 367.841 null] >>}
\MusicPDFWriteObject{11350}{<< /D [\MusicPDFPageRef{223} /XYZ 72 242.383 null] >>}
\MusicPDFWriteObject{11351}{<< /D [\MusicPDFPageRef{224} /XYZ 72 698.082 null] >>}
\MusicPDFWriteObject{11352}{<< /D [\MusicPDFPageRef{224} /XYZ 72 672.118 null] >>}
\MusicPDFWriteObject{11353}{<< /D [\MusicPDFPageRef{224} /XYZ 72 644.372 null] >>}
\MusicPDFWriteObject{11354}{<< /D [\MusicPDFPageRef{224} /XYZ 72 297.988 null] >>}
\MusicPDFWriteObject{11355}{<< /D [\MusicPDFPageRef{224} /XYZ 72 259.07 null] >>}
\MusicPDFWriteObject{11356}{<< /D [\MusicPDFPageRef{225} /XYZ 72 379.947 null] >>}
\MusicPDFWriteObject{11357}{<< /D [\MusicPDFPageRef{38} /XYZ 72 612.019 null] >>}
\MusicPDFWriteObject{11358}{<< /D [\MusicPDFPageRef{225} /XYZ 72 343.52 null] >>}
\MusicPDFWriteObject{11359}{<< /D [\MusicPDFPageRef{227} /XYZ 72 301.659 null] >>}
\MusicPDFWriteObject{11360}{<< /D [\MusicPDFPageRef{227} /XYZ 72 198.924 null] >>}
\MusicPDFWriteObject{11361}{<< /D [\MusicPDFPageRef{230} /XYZ 72 444.519 null] >>}
\MusicPDFWriteObject{11362}{<< /D [\MusicPDFPageRef{230} /XYZ 72 381.79 null] >>}
\MusicPDFWriteObject{11363}{<< /D [\MusicPDFPageRef{230} /XYZ 72 239.136 null] >>}
\MusicPDFWriteObject{11364}{<< /D [\MusicPDFPageRef{232} /XYZ 72 280.034 null] >>}
\MusicPDFWriteObject{11365}{<< /D [\MusicPDFPageRef{232} /XYZ 72 227.606 null] >>}
\MusicPDFWriteObject{11366}{<< /D [\MusicPDFPageRef{233} /XYZ 72 602.149 null] >>}
\MusicPDFWriteObject{11367}{<< /D [\MusicPDFPageRef{233} /XYZ 72 449.987 null] >>}
\MusicPDFWriteObject{11368}{<< /D [\MusicPDFPageRef{38} /XYZ 72 584.616 null] >>}
\MusicPDFWriteObject{11369}{<< /D [\MusicPDFPageRef{233} /XYZ 72 399.321 null] >>}
\MusicPDFWriteObject{11370}{<< /D [\MusicPDFPageRef{236} /XYZ 72 668.777 null] >>}
\MusicPDFWriteObject{11371}{<< /D [\MusicPDFPageRef{236} /XYZ 72 503.789 null] >>}
\MusicPDFWriteObject{11372}{<< /D [\MusicPDFPageRef{236} /XYZ 72 433.666 null] >>}
\MusicPDFWriteObject{11373}{<< /D [\MusicPDFPageRef{237} /XYZ 72 300.139 null] >>}
\MusicPDFWriteObject{11374}{<< /D [\MusicPDFPageRef{237} /XYZ 72 259.773 null] >>}
\MusicPDFWriteObject{11375}{<< /D [\MusicPDFPageRef{238} /XYZ 72 115.809 null] >>}
\MusicPDFWriteObject{11376}{<< /D [\MusicPDFPageRef{239} /XYZ 72 569.481 null] >>}
\MusicPDFWriteObject{11377}{<< /D [\MusicPDFPageRef{241} /XYZ 72 503.15 null] >>}
\MusicPDFWriteObject{11378}{<< /D [\MusicPDFPageRef{241} /XYZ 72 491.433 null] >>}
\MusicPDFWriteObject{11379}{<< /D [\MusicPDFPageRef{46} /XYZ 72 539.513 null] >>}
\MusicPDFWriteObject{11380}{<< /D [\MusicPDFPageRef{241} /XYZ 72 478.282 null] >>}
\MusicPDFWriteObject{11381}{<< /D [\MusicPDFPageRef{241} /XYZ 72 117.447 null] >>}
\MusicPDFWriteObject{11382}{<< /D [\MusicPDFPageRef{242} /XYZ 72 611.362 null] >>}
\MusicPDFWriteObject{11383}{<< /D [\MusicPDFPageRef{243} /XYZ 72 404.557 null] >>}
\MusicPDFWriteObject{11384}{<< /D [\MusicPDFPageRef{243} /XYZ 72 266.905 null] >>}
\MusicPDFWriteObject{11385}{<< /D [\MusicPDFPageRef{245} /XYZ 72 459.532 null] >>}
\MusicPDFWriteObject{11386}{<< /D [\MusicPDFPageRef{245} /XYZ 72 430.735 null] >>}
\MusicPDFWriteObject{11387}{<< /D [\MusicPDFPageRef{245} /XYZ 72 367.746 null] >>}
\MusicPDFWriteObject{11388}{<< /D [\MusicPDFPageRef{245} /XYZ 72 305.57 null] >>}
\MusicPDFWriteObject{11389}{<< /D [\MusicPDFPageRef{245} /XYZ 72 254.857 null] >>}
\MusicPDFWriteObject{11390}{<< /D [\MusicPDFPageRef{13} /XYZ 72 216.229 null] >>}
\MusicPDFWriteObject{11391}{<< /D [\MusicPDFPageRef{46} /XYZ 72 464.829 null] >>}
\MusicPDFWriteObject{11392}{<< /D [\MusicPDFPageRef{246} /XYZ 72 319.295 null] >>}
\MusicPDFWriteObject{11393}{<< /D [\MusicPDFPageRef{246} /XYZ 72 178.134 null] >>}
\MusicPDFWriteObject{11394}{<< /D [\MusicPDFPageRef{249} /XYZ 72 176.201 null] >>}
\MusicPDFWriteObject{11395}{<< /D [\MusicPDFPageRef{249} /XYZ 72 135.143 null] >>}
\MusicPDFWriteObject{11396}{<< /D [\MusicPDFPageRef{249} /XYZ 72 123.044 null] >>}
\MusicPDFWriteObject{11397}{<< /D [\MusicPDFPageRef{250} /XYZ 72 673.032 null] >>}
\MusicPDFWriteObject{11398}{<< /D [\MusicPDFPageRef{252} /XYZ 72 250.993 null] >>}
\MusicPDFWriteObject{11399}{<< /D [\MusicPDFPageRef{252} /XYZ 72 235.351 null] >>}
\MusicPDFWriteObject{11400}{<< /D [\MusicPDFPageRef{252} /XYZ 72 175.113 null] >>}
\MusicPDFWriteObject{11401}{<< /D [\MusicPDFPageRef{255} /XYZ 72 146.875 null] >>}
\MusicPDFWriteObject{11402}{<< /D [\MusicPDFPageRef{46} /XYZ 72 382.036 null] >>}
\MusicPDFWriteObject{11403}{<< /D [\MusicPDFPageRef{255} /XYZ 72 133.17 null] >>}
\MusicPDFWriteObject{11404}{<< /D [\MusicPDFPageRef{258} /XYZ 72 445.074 null] >>}
\MusicPDFWriteObject{11405}{<< /D [\MusicPDFPageRef{258} /XYZ 72 371.505 null] >>}
\MusicPDFWriteObject{11406}{<< /D [\MusicPDFPageRef{260} /XYZ 72 698.082 null] >>}
\MusicPDFWriteObject{11407}{<< /D [\MusicPDFPageRef{260} /XYZ 72 659.881 null] >>}
\MusicPDFWriteObject{11408}{<< /D [\MusicPDFPageRef{260} /XYZ 72 645.817 null] >>}
\MusicPDFWriteObject{11409}{<< /D [\MusicPDFPageRef{262} /XYZ 72 384.432 null] >>}
\MusicPDFWriteObject{11410}{<< /D [\MusicPDFPageRef{262} /XYZ 72 340.226 null] >>}
\MusicPDFWriteObject{11411}{<< /D [\MusicPDFPageRef{264} /XYZ 72 698.082 null] >>}
\MusicPDFWriteObject{11412}{<< /D [\MusicPDFPageRef{264} /XYZ 72 682.614 null] >>}
\MusicPDFWriteObject{11413}{<< /D [\MusicPDFPageRef{47} /XYZ 72 532.386 null] >>}
\MusicPDFWriteObject{11414}{<< /D [\MusicPDFPageRef{264} /XYZ 72 668.639 null] >>}
\MusicPDFWriteObject{11415}{<< /D [\MusicPDFPageRef{264} /XYZ 72 338.55 null] >>}
\MusicPDFWriteObject{11416}{<< /D [\MusicPDFPageRef{264} /XYZ 72 277.927 null] >>}
\MusicPDFWriteObject{11417}{<< /D [\MusicPDFPageRef{264} /XYZ 72 155.424 null] >>}
\MusicPDFWriteObject{11418}{<< /D [\MusicPDFPageRef{265} /XYZ 72 116.401 null] >>}
\MusicPDFWriteObject{11419}{<< /D [\MusicPDFPageRef{266} /XYZ 72 583.088 null] >>}
\MusicPDFWriteObject{11420}{<< /D [\MusicPDFPageRef{275} /XYZ 72 588.917 null] >>}
\MusicPDFWriteObject{11421}{<< /D [\MusicPDFPageRef{275} /XYZ 72 533.833 null] >>}
\MusicPDFWriteObject{11422}{<< /D [\MusicPDFPageRef{276} /XYZ 72 226.629 null] >>}
\MusicPDFWriteObject{11423}{<< /D [\MusicPDFPageRef{276} /XYZ 72 181.132 null] >>}
\MusicPDFWriteObject{11424}{<< /D [\MusicPDFPageRef{47} /XYZ 72 492.097 null] >>}
\MusicPDFWriteObject{11425}{<< /D [\MusicPDFPageRef{276} /XYZ 72 137.934 null] >>}
\MusicPDFWriteObject{11426}{<< /D [\MusicPDFPageRef{276} /XYZ 72 109.924 null] >>}
\MusicPDFWriteObject{11427}{<< /D [\MusicPDFPageRef{277} /XYZ 72 698.082 null] >>}
\MusicPDFWriteObject{11428}{<< /D [\MusicPDFPageRef{277} /XYZ 72 656.706 null] >>}
\MusicPDFWriteObject{11429}{<< /D [\MusicPDFPageRef{277} /XYZ 72 210.101 null] >>}
\MusicPDFWriteObject{11430}{<< /D [\MusicPDFPageRef{278} /XYZ 72 698.082 null] >>}
\MusicPDFWriteObject{11431}{<< /D [\MusicPDFPageRef{278} /XYZ 72 548.152 null] >>}
\MusicPDFWriteObject{11432}{<< /D [\MusicPDFPageRef{278} /XYZ 72 451.671 null] >>}
\MusicPDFWriteObject{11433}{<< /D [\MusicPDFPageRef{279} /XYZ 72 343.113 null] >>}
\MusicPDFWriteObject{11434}{<< /D [\MusicPDFPageRef{279} /XYZ 72 246.518 null] >>}
\MusicPDFWriteObject{11435}{<< /D [\MusicPDFPageRef{47} /XYZ 72 445.009 null] >>}
\MusicPDFWriteObject{11436}{<< /D [\MusicPDFPageRef{282} /XYZ 72 640.081 null] >>}
\MusicPDFWriteObject{11437}{<< /D [\MusicPDFPageRef{282} /XYZ 72 585.674 null] >>}
\MusicPDFWriteObject{11438}{<< /D [\MusicPDFPageRef{285} /XYZ 72 463.008 null] >>}
\MusicPDFWriteObject{11439}{<< /D [\MusicPDFPageRef{285} /XYZ 72 408.961 null] >>}
\MusicPDFWriteObject{11440}{<< /D [\MusicPDFPageRef{288} /XYZ 72 367.743 null] >>}
\MusicPDFWriteObject{11441}{<< /D [\MusicPDFPageRef{288} /XYZ 72 339.411 null] >>}
\MusicPDFWriteObject{11442}{<< /D [\MusicPDFPageRef{288} /XYZ 72 262.804 null] >>}
\MusicPDFWriteObject{11443}{<< /D [\MusicPDFPageRef{288} /XYZ 72 212.672 null] >>}
\MusicPDFWriteObject{11444}{<< /D [\MusicPDFPageRef{288} /XYZ 72 185.369 null] >>}
\MusicPDFWriteObject{11445}{<< /D [\MusicPDFPageRef{288} /XYZ 72 147.472 null] >>}
\MusicPDFWriteObject{11446}{<< /D [\MusicPDFPageRef{47} /XYZ 72 403.303 null] >>}
\MusicPDFWriteObject{11447}{<< /D [\MusicPDFPageRef{289} /XYZ 72 698.082 null] >>}
\MusicPDFWriteObject{11448}{<< /D [\MusicPDFPageRef{289} /XYZ 72 634.257 null] >>}
\MusicPDFWriteObject{11449}{<< /D [\MusicPDFPageRef{290} /XYZ 72 601.124 null] >>}
\MusicPDFWriteObject{11450}{<< /D [\MusicPDFPageRef{290} /XYZ 72 525.967 null] >>}
\MusicPDFWriteObject{11451}{<< /D [\MusicPDFPageRef{291} /XYZ 72 110.447 null] >>}
\MusicPDFWriteObject{11452}{<< /D [\MusicPDFPageRef{292} /XYZ 72 698.082 null] >>}
\MusicPDFWriteObject{11453}{<< /D [\MusicPDFPageRef{292} /XYZ 72 538.606 null] >>}
\MusicPDFWriteObject{11454}{<< /D [\MusicPDFPageRef{295} /XYZ 72 277.933 null] >>}
\MusicPDFWriteObject{11455}{<< /D [\MusicPDFPageRef{295} /XYZ 72 251.632 null] >>}
\MusicPDFWriteObject{11456}{<< /D [\MusicPDFPageRef{296} /XYZ 72 604.068 null] >>}
\MusicPDFWriteObject{11457}{<< /D [\MusicPDFPageRef{47} /XYZ 72 287.811 null] >>}
\MusicPDFWriteObject{11458}{<< /D [\MusicPDFPageRef{296} /XYZ 72 574.711 null] >>}
\MusicPDFWriteObject{11459}{<< /D [\MusicPDFPageRef{297} /XYZ 72 229.662 null] >>}
\MusicPDFWriteObject{11460}{<< /D [\MusicPDFPageRef{298} /XYZ 72 620.429 null] >>}
\MusicPDFWriteObject{11461}{<< /D [\MusicPDFPageRef{298} /XYZ 72 319.853 null] >>}
\MusicPDFWriteObject{11462}{<< /D [\MusicPDFPageRef{298} /XYZ 72 277.502 null] >>}
\MusicPDFWriteObject{11463}{<< /D [\MusicPDFPageRef{298} /XYZ 72 260.859 null] >>}
\MusicPDFWriteObject{11464}{<< /D [\MusicPDFPageRef{299} /XYZ 72 625.559 null] >>}
\MusicPDFWriteObject{11465}{<< /D [\MusicPDFPageRef{299} /XYZ 72 558.396 null] >>}
\MusicPDFWriteObject{11466}{<< /D [\MusicPDFPageRef{304} /XYZ 72 488.568 null] >>}
\MusicPDFWriteObject{11467}{<< /D [\MusicPDFPageRef{304} /XYZ 72 368.107 null] >>}
\MusicPDFWriteObject{11468}{<< /D [\MusicPDFPageRef{48} /XYZ 72 698.082 null] >>}
\MusicPDFWriteObject{11469}{<< /D [\MusicPDFPageRef{304} /XYZ 72 323.066 null] >>}
\MusicPDFWriteObject{11470}{<< /D [\MusicPDFPageRef{304} /XYZ 72 294.827 null] >>}
\MusicPDFWriteObject{11471}{<< /D [\MusicPDFPageRef{304} /XYZ 72 266.84 null] >>}
\MusicPDFWriteObject{11472}{<< /D [\MusicPDFPageRef{307} /XYZ 72 566.794 null] >>}
\MusicPDFWriteObject{11473}{<< /D [\MusicPDFPageRef{307} /XYZ 72 463.52 null] >>}
\MusicPDFWriteObject{11474}{<< /D [\MusicPDFPageRef{307} /XYZ 72 452.561 null] >>}
\MusicPDFWriteObject{11475}{<< /D [\MusicPDFPageRef{310} /XYZ 72 698.082 null] >>}
\MusicPDFWriteObject{11476}{<< /D [\MusicPDFPageRef{310} /XYZ 72 571.286 null] >>}
\MusicPDFWriteObject{11477}{<< /D [\MusicPDFPageRef{310} /XYZ 72 524.558 null] >>}
\MusicPDFWriteObject{11478}{<< /D [\MusicPDFPageRef{310} /XYZ 72 496.32 null] >>}
\MusicPDFWriteObject{11479}{<< /D [\MusicPDFPageRef{48} /XYZ 72 171.141 null] >>}
\MusicPDFWriteObject{11480}{<< /D [\MusicPDFPageRef{310} /XYZ 72 468.333 null] >>}
\MusicPDFWriteObject{11481}{<< /D [\MusicPDFPageRef{311} /XYZ 72 184.989 null] >>}
\MusicPDFWriteObject{11482}{<< /D [\MusicPDFPageRef{311} /XYZ 72 126.25 null] >>}
\MusicPDFWriteObject{11483}{<< /D [\MusicPDFPageRef{313} /XYZ 72 219.955 null] >>}
\MusicPDFWriteObject{11484}{<< /D [\MusicPDFPageRef{313} /XYZ 72 206.805 null] >>}
\MusicPDFWriteObject{11485}{<< /D [\MusicPDFPageRef{313} /XYZ 72 192.602 null] >>}
\MusicPDFWriteObject{11486}{<< /D [\MusicPDFPageRef{317} /XYZ 72 698.082 null] >>}
\MusicPDFWriteObject{11487}{<< /D [\MusicPDFPageRef{317} /XYZ 72 617.268 null] >>}
\MusicPDFWriteObject{11488}{<< /D [\MusicPDFPageRef{324} /XYZ 72 484.748 null] >>}
\MusicPDFWriteObject{11489}{<< /D [\MusicPDFPageRef{324} /XYZ 72 457.837 null] >>}
\MusicPDFWriteObject{11490}{<< /D [\MusicPDFPageRef{48} /XYZ 72 157.437 null] >>}
\MusicPDFWriteObject{11491}{<< /D [\MusicPDFPageRef{324} /XYZ 72 432.145 null] >>}
\MusicPDFWriteObject{11492}{<< /D [\MusicPDFPageRef{324} /XYZ 72 393.052 null] >>}
\MusicPDFWriteObject{11493}{<< /D [\MusicPDFPageRef{324} /XYZ 72 366.751 null] >>}
\MusicPDFWriteObject{11494}{<< /D [\MusicPDFPageRef{326} /XYZ 72 277.591 null] >>}
\MusicPDFWriteObject{11495}{<< /D [\MusicPDFPageRef{326} /XYZ 72 114.18 null] >>}
\MusicPDFWriteObject{11496}{<< /D [\MusicPDFPageRef{327} /XYZ 72 608.766 null] >>}
\MusicPDFWriteObject{11497}{<< /D [\MusicPDFPageRef{327} /XYZ 72 532.699 null] >>}
\MusicPDFWriteObject{11498}{<< /D [\MusicPDFPageRef{331} /XYZ 72 542.764 null] >>}
\MusicPDFWriteObject{11499}{<< /D [\MusicPDFPageRef{331} /XYZ 72 301.215 null] >>}
\MusicPDFWriteObject{11500}{<< /D [\MusicPDFPageRef{331} /XYZ 72 174.248 null] >>}
\MusicPDFWriteObject{11501}{<< /D [\MusicPDFPageRef{13} /XYZ 72 202.641 null] >>}
\MusicPDFWriteObject{11502}{<< /D [\MusicPDFPageRef{48} /XYZ 72 144.286 null] >>}
\MusicPDFWriteObject{11503}{<< /D [\MusicPDFPageRef{334} /XYZ 72 698.082 null] >>}
\MusicPDFWriteObject{11504}{<< /D [\MusicPDFPageRef{334} /XYZ 72 685.269 null] >>}
\MusicPDFWriteObject{11505}{<< /D [\MusicPDFPageRef{334} /XYZ 72 656.312 null] >>}
\MusicPDFWriteObject{11506}{<< /D [\MusicPDFPageRef{334} /XYZ 72 644.993 null] >>}
\MusicPDFWriteObject{11507}{<< /D [\MusicPDFPageRef{336} /XYZ 72 141.966 null] >>}
\MusicPDFWriteObject{11508}{<< /D [\MusicPDFPageRef{337} /XYZ 72 659.167 null] >>}
\MusicPDFWriteObject{11509}{<< /D [\MusicPDFPageRef{339} /XYZ 72 481.751 null] >>}
\MusicPDFWriteObject{11510}{<< /D [\MusicPDFPageRef{339} /XYZ 72 465.945 null] >>}
\MusicPDFWriteObject{11511}{<< /D [\MusicPDFPageRef{339} /XYZ 72 427.954 null] >>}
\MusicPDFWriteObject{11512}{<< /D [\MusicPDFPageRef{339} /XYZ 72 411.234 null] >>}
\MusicPDFWriteObject{11513}{<< /D [\MusicPDFPageRef{49} /XYZ 72 136.195 null] >>}
\MusicPDFWriteObject{11514}{<< /D [\MusicPDFPageRef{340} /XYZ 72 554.842 null] >>}
\MusicPDFWriteObject{11515}{<< /D [\MusicPDFPageRef{340} /XYZ 72 528.077 null] >>}
\MusicPDFWriteObject{11516}{<< /D [\MusicPDFPageRef{340} /XYZ 72 512.879 null] >>}
\MusicPDFWriteObject{11517}{<< /D [\MusicPDFPageRef{340} /XYZ 72 498.76 null] >>}
\MusicPDFWriteObject{11518}{<< /D [\MusicPDFPageRef{340} /XYZ 72 277.523 null] >>}
\MusicPDFWriteObject{11519}{<< /D [\MusicPDFPageRef{340} /XYZ 72 238.513 null] >>}
\MusicPDFWriteObject{11520}{<< /D [\MusicPDFPageRef{340} /XYZ 72 213.706 null] >>}
\MusicPDFWriteObject{11521}{<< /D [\MusicPDFPageRef{341} /XYZ 72 597.582 null] >>}
\MusicPDFWriteObject{11522}{<< /D [\MusicPDFPageRef{341} /XYZ 72 567.966 null] >>}
\MusicPDFWriteObject{11523}{<< /D [\MusicPDFPageRef{343} /XYZ 72 248.265 null] >>}
\MusicPDFWriteObject{11524}{<< /D [\MusicPDFPageRef{49} /XYZ 72 109.894 null] >>}
\MusicPDFWriteObject{11525}{<< /D [\MusicPDFPageRef{343} /XYZ 72 234.561 null] >>}
\MusicPDFWriteObject{11526}{<< /D [\MusicPDFPageRef{343} /XYZ 72 221.41 null] >>}
\MusicPDFWriteObject{11527}{<< /D [\MusicPDFPageRef{343} /XYZ 72 207.962 null] >>}
\MusicPDFWriteObject{11528}{<< /D [\MusicPDFPageRef{343} /XYZ 72 125.727 null] >>}
\MusicPDFWriteObject{11529}{<< /D [\MusicPDFPageRef{344} /XYZ 72 698.082 null] >>}
\MusicPDFWriteObject{11530}{<< /D [\MusicPDFPageRef{345} /XYZ 72 669.892 null] >>}
\MusicPDFWriteObject{11531}{<< /D [\MusicPDFPageRef{345} /XYZ 72 629.109 null] >>}
\MusicPDFWriteObject{11532}{<< /D [\MusicPDFPageRef{345} /XYZ 72 568.7 null] >>}
\MusicPDFWriteObject{11533}{<< /D [\MusicPDFPageRef{345} /XYZ 72 518.909 null] >>}
\MusicPDFWriteObject{11534}{<< /D [\MusicPDFPageRef{345} /XYZ 72 385.861 null] >>}
\MusicPDFWriteObject{11535}{<< /D [\MusicPDFPageRef{49} /XYZ 72 96.743 null] >>}
\MusicPDFWriteObject{11536}{<< /D [\MusicPDFPageRef{347} /XYZ 72 360.654 null] >>}
\MusicPDFWriteObject{11537}{<< /D [\MusicPDFPageRef{347} /XYZ 72 328.023 null] >>}
\MusicPDFWriteObject{11538}{<< /D [\MusicPDFPageRef{347} /XYZ 72 317.067 null] >>}
\MusicPDFWriteObject{11539}{<< /D [\MusicPDFPageRef{347} /XYZ 72 300.9 null] >>}
\MusicPDFWriteObject{11540}{<< /D [\MusicPDFPageRef{347} /XYZ 72 289.943 null] >>}
\MusicPDFWriteObject{11541}{<< /D [\MusicPDFPageRef{347} /XYZ 72 225.209 null] >>}
\MusicPDFWriteObject{11542}{<< /D [\MusicPDFPageRef{347} /XYZ 72 212.419 null] >>}
\MusicPDFWriteObject{11543}{<< /D [\MusicPDFPageRef{347} /XYZ 72 182.278 null] >>}
\MusicPDFWriteObject{11544}{<< /D [\MusicPDFPageRef{347} /XYZ 72 171.321 null] >>}
\MusicPDFWriteObject{11545}{<< /D [\MusicPDFPageRef{347} /XYZ 72 155.155 null] >>}
\MusicPDFWriteObject{11546}{<< /D [\MusicPDFPageRef{52} /XYZ 72 344.604 null] >>}
\MusicPDFWriteObject{11547}{<< /D [\MusicPDFPageRef{347} /XYZ 72 126.356 null] >>}
\MusicPDFWriteObject{11548}{<< /D [\MusicPDFPageRef{348} /XYZ 72 698.082 null] >>}
\MusicPDFWriteObject{11549}{<< /D [\MusicPDFPageRef{348} /XYZ 72 671.629 null] >>}
\MusicPDFWriteObject{11550}{<< /D [\MusicPDFPageRef{348} /XYZ 72 365.908 null] >>}
\MusicPDFWriteObject{11551}{<< /D [\MusicPDFPageRef{348} /XYZ 72 350.71 null] >>}
\MusicPDFWriteObject{11552}{<< /D [\MusicPDFPageRef{348} /XYZ 72 339.754 null] >>}
\MusicPDFWriteObject{11553}{<< /D [\MusicPDFPageRef{348} /XYZ 72 326.603 null] >>}
\MusicPDFWriteObject{11554}{<< /D [\MusicPDFPageRef{349} /XYZ 72 592.12 null] >>}
\MusicPDFWriteObject{11555}{<< /D [\MusicPDFPageRef{349} /XYZ 72 442.12 null] >>}
\MusicPDFWriteObject{11556}{<< /D [\MusicPDFPageRef{349} /XYZ 72 402.115 null] >>}
\MusicPDFWriteObject{11557}{<< /D [\MusicPDFPageRef{52} /XYZ 72 132.506 null] >>}
\MusicPDFWriteObject{11558}{<< /D [\MusicPDFPageRef{349} /XYZ 72 388.964 null] >>}
\MusicPDFWriteObject{11559}{<< /D [\MusicPDFPageRef{349} /XYZ 72 375.813 null] >>}
\MusicPDFWriteObject{11560}{<< /D [\MusicPDFPageRef{350} /XYZ 72 474.437 null] >>}
\MusicPDFWriteObject{11561}{<< /D [\MusicPDFPageRef{350} /XYZ 72 450.651 null] >>}
\MusicPDFWriteObject{11562}{<< /D [\MusicPDFPageRef{350} /XYZ 72 280.905 null] >>}
\MusicPDFWriteObject{11563}{<< /D [\MusicPDFPageRef{350} /XYZ 72 267.201 null] >>}
\MusicPDFWriteObject{11564}{<< /D [\MusicPDFPageRef{351} /XYZ 72 519.588 null] >>}
\MusicPDFWriteObject{11565}{<< /D [\MusicPDFPageRef{351} /XYZ 72 299.153 null] >>}
\MusicPDFWriteObject{11566}{<< /D [\MusicPDFPageRef{352} /XYZ 72 594.529 null] >>}
\MusicPDFWriteObject{11567}{<< /D [\MusicPDFPageRef{352} /XYZ 72 470.271 null] >>}
\MusicPDFWriteObject{11568}{<< /D [\MusicPDFPageRef{53} /XYZ 72 389.314 null] >>}
\MusicPDFWriteObject{11569}{<< /D [\MusicPDFPageRef{352} /XYZ 72 414.831 null] >>}
\MusicPDFWriteObject{11570}{<< /D [\MusicPDFPageRef{352} /XYZ 72 379.547 null] >>}
\MusicPDFWriteObject{11571}{<< /D [\MusicPDFPageRef{355} /XYZ 72 669.476 null] >>}
\MusicPDFWriteObject{11572}{<< /D [\MusicPDFPageRef{355} /XYZ 72 448.726 null] >>}
\MusicPDFWriteObject{11573}{<< /D [\MusicPDFPageRef{355} /XYZ 72 157.205 null] >>}
\MusicPDFWriteObject{11574}{<< /D [\MusicPDFPageRef{356} /XYZ 72 649.629 null] >>}
\MusicPDFWriteObject{11575}{<< /D [\MusicPDFPageRef{356} /XYZ 72 619.852 null] >>}
\MusicPDFWriteObject{11576}{<< /D [\MusicPDFPageRef{356} /XYZ 72 587.744 null] >>}
\MusicPDFWriteObject{11577}{<< /D [\MusicPDFPageRef{357} /XYZ 72 325.978 null] >>}
\MusicPDFWriteObject{11578}{<< /D [\MusicPDFPageRef{357} /XYZ 72 312.274 null] >>}
\MusicPDFWriteObject{11579}{<< /D [\MusicPDFPageRef{53} /XYZ 72 375.61 null] >>}
\MusicPDFWriteObject{11580}{<< /D [\MusicPDFPageRef{357} /XYZ 72 244.728 null] >>}
\MusicPDFWriteObject{11581}{<< /D [\MusicPDFPageRef{359} /XYZ 72 297.173 null] >>}
\MusicPDFWriteObject{11582}{<< /D [\MusicPDFPageRef{359} /XYZ 72 211.388 null] >>}
\MusicPDFWriteObject{11583}{<< /D [\MusicPDFPageRef{360} /XYZ 72 343.265 null] >>}
\MusicPDFWriteObject{11584}{<< /D [\MusicPDFPageRef{360} /XYZ 72 150.269 null] >>}
\MusicPDFWriteObject{11585}{<< /D [\MusicPDFPageRef{361} /XYZ 72 666.232 null] >>}
\MusicPDFWriteObject{11586}{<< /D [\MusicPDFPageRef{364} /XYZ 72 138.078 null] >>}
\MusicPDFWriteObject{11587}{<< /D [\MusicPDFPageRef{364} /XYZ 72 119.864 null] >>}
\MusicPDFWriteObject{11588}{<< /D [\MusicPDFPageRef{365} /XYZ 72 621.86 null] >>}
\MusicPDFWriteObject{11589}{<< /D [\MusicPDFPageRef{365} /XYZ 72 568.636 null] >>}
\MusicPDFWriteObject{11590}{<< /D [\MusicPDFPageRef{57} /XYZ 72 313.642 null] >>}
\MusicPDFWriteObject{11591}{<< /D [\MusicPDFPageRef{365} /XYZ 72 429.084 null] >>}
\MusicPDFWriteObject{11592}{<< /D [\MusicPDFPageRef{366} /XYZ 72 400.373 null] >>}
\MusicPDFWriteObject{11593}{<< /D [\MusicPDFPageRef{366} /XYZ 72 384.622 null] >>}
\MusicPDFWriteObject{11594}{<< /D [\MusicPDFPageRef{367} /XYZ 72 353.523 null] >>}
\MusicPDFWriteObject{11595}{<< /D [\MusicPDFPageRef{367} /XYZ 72 310.986 null] >>}
\MusicPDFWriteObject{11596}{<< /D [\MusicPDFPageRef{368} /XYZ 72 521.7 null] >>}
\MusicPDFWriteObject{11597}{<< /D [\MusicPDFPageRef{368} /XYZ 72 501.571 null] >>}
\MusicPDFWriteObject{11598}{<< /D [\MusicPDFPageRef{368} /XYZ 72 231.244 null] >>}
\MusicPDFWriteObject{11599}{<< /D [\MusicPDFPageRef{368} /XYZ 72 178.791 null] >>}
\MusicPDFWriteObject{11600}{<< /D [\MusicPDFPageRef{368} /XYZ 72 113.462 null] >>}
\MusicPDFWriteObject{11601}{<< /D [\MusicPDFPageRef{57} /XYZ 72 273.636 null] >>}
\MusicPDFWriteObject{11602}{<< /D [\MusicPDFPageRef{372} /XYZ 72 165.582 null] >>}
\MusicPDFWriteObject{11603}{<< /D [\MusicPDFPageRef{372} /XYZ 72 123.044 null] >>}
\MusicPDFWriteObject{11604}{<< /D [\MusicPDFPageRef{382} /XYZ 72 516.997 null] >>}
\MusicPDFWriteObject{11605}{<< /D [\MusicPDFPageRef{382} /XYZ 72 412.202 null] >>}
\MusicPDFWriteObject{11606}{<< /D [\MusicPDFPageRef{384} /XYZ 72 297.546 null] >>}
\MusicPDFWriteObject{11607}{<< /D [\MusicPDFPageRef{384} /XYZ 72 170.151 null] >>}
\MusicPDFWriteObject{11608}{<< /D [\MusicPDFPageRef{385} /XYZ 72 698.082 null] >>}
\MusicPDFWriteObject{11609}{<< /D [\MusicPDFPageRef{385} /XYZ 72 526.157 null] >>}
\MusicPDFWriteObject{11610}{<< /D [\MusicPDFPageRef{385} /XYZ 72 438.214 null] >>}
\MusicPDFWriteObject{11611}{<< /D [\MusicPDFPageRef{385} /XYZ 72 330.791 null] >>}
\MusicPDFWriteObject{11612}{<< /D [\MusicPDFPageRef{15} /XYZ 72 630.135 null] >>}
\MusicPDFWriteObject{11613}{<< /D [\MusicPDFPageRef{57} /XYZ 72 247.888 null] >>}
\MusicPDFWriteObject{11614}{<< /D [\MusicPDFPageRef{385} /XYZ 72 300.334 null] >>}
\MusicPDFWriteObject{11615}{<< /D [\MusicPDFPageRef{386} /XYZ 72 296.702 null] >>}
\MusicPDFWriteObject{11616}{<< /D [\MusicPDFPageRef{386} /XYZ 72 269.846 null] >>}
\MusicPDFWriteObject{11617}{<< /D [\MusicPDFPageRef{386} /XYZ 72 109.893 null] >>}
\MusicPDFWriteObject{11618}{<< /D [\MusicPDFPageRef{387} /XYZ 72 698.082 null] >>}
\MusicPDFWriteObject{11619}{<< /D [\MusicPDFPageRef{387} /XYZ 72 496.968 null] >>}
\MusicPDFWriteObject{11620}{<< /D [\MusicPDFPageRef{387} /XYZ 72 483.457 null] >>}
\MusicPDFWriteObject{11621}{<< /D [\MusicPDFPageRef{388} /XYZ 72 682.845 null] >>}
\MusicPDFWriteObject{11622}{<< /D [\MusicPDFPageRef{388} /XYZ 72 475.496 null] >>}
\MusicPDFWriteObject{11623}{<< /D [\MusicPDFPageRef{396} /XYZ 72 151.938 null] >>}
\MusicPDFWriteObject{11624}{<< /D [\MusicPDFPageRef{57} /XYZ 72 207.522 null] >>}
\MusicPDFWriteObject{11625}{<< /D [\MusicPDFPageRef{397} /XYZ 72 698.082 null] >>}
\MusicPDFWriteObject{11626}{<< /D [\MusicPDFPageRef{397} /XYZ 72 647.179 null] >>}
\MusicPDFWriteObject{11627}{<< /D [\MusicPDFPageRef{404} /XYZ 72 541.767 null] >>}
\MusicPDFWriteObject{11628}{<< /D [\MusicPDFPageRef{404} /XYZ 72 437.785 null] >>}
\MusicPDFWriteObject{11629}{<< /D [\MusicPDFPageRef{404} /XYZ 72 175.517 null] >>}
\MusicPDFWriteObject{11630}{<< /D [\MusicPDFPageRef{404} /XYZ 72 132.118 null] >>}
\MusicPDFWriteObject{11631}{<< /D [\MusicPDFPageRef{405} /XYZ 72 698.082 null] >>}
\MusicPDFWriteObject{11632}{<< /D [\MusicPDFPageRef{405} /XYZ 72 563.357 null] >>}
\MusicPDFWriteObject{11633}{<< /D [\MusicPDFPageRef{405} /XYZ 72 549.652 null] >>}
\MusicPDFWriteObject{11634}{<< /D [\MusicPDFPageRef{408} /XYZ 72 510.061 null] >>}
\MusicPDFWriteObject{11635}{<< /D [\MusicPDFPageRef{64} /XYZ 72 408.739 null] >>}
\MusicPDFWriteObject{11636}{<< /D [\MusicPDFPageRef{408} /XYZ 72 453.667 null] >>}
\MusicPDFWriteObject{11637}{<< /D [\MusicPDFPageRef{408} /XYZ 72 351.893 null] >>}
\MusicPDFWriteObject{11638}{<< /D [\MusicPDFPageRef{412} /XYZ 72 420.712 null] >>}
\MusicPDFWriteObject{11639}{<< /D [\MusicPDFPageRef{412} /XYZ 72 381.513 null] >>}
\MusicPDFWriteObject{11640}{<< /D [\MusicPDFPageRef{412} /XYZ 72 266.386 null] >>}
\MusicPDFWriteObject{11641}{<< /D [\MusicPDFPageRef{412} /XYZ 72 255.429 null] >>}
\MusicPDFWriteObject{11642}{<< /D [\MusicPDFPageRef{412} /XYZ 72 127.059 null] >>}
\MusicPDFWriteObject{11643}{<< /D [\MusicPDFPageRef{414} /XYZ 72 527.128 null] >>}
\MusicPDFWriteObject{11644}{<< /D [\MusicPDFPageRef{414} /XYZ 72 498.039 null] >>}
\MusicPDFWriteObject{11645}{<< /D [\MusicPDFPageRef{414} /XYZ 72 429.506 null] >>}
\MusicPDFWriteObject{11646}{<< /D [\MusicPDFPageRef{64} /XYZ 72 382.437 null] >>}
\MusicPDFWriteObject{11647}{<< /D [\MusicPDFPageRef{414} /XYZ 72 278.561 null] >>}
\MusicPDFWriteObject{11648}{<< /D [\MusicPDFPageRef{414} /XYZ 72 230.575 null] >>}
\MusicPDFWriteObject{11649}{<< /D [\MusicPDFPageRef{415} /XYZ 72 582.669 null] >>}
\MusicPDFWriteObject{11650}{<< /D [\MusicPDFPageRef{415} /XYZ 72 516.398 null] >>}
\MusicPDFWriteObject{11651}{<< /D [\MusicPDFPageRef{415} /XYZ 72 374.352 null] >>}
\MusicPDFWriteObject{11652}{<< /D [\MusicPDFPageRef{415} /XYZ 72 360.647 null] >>}
\MusicPDFWriteObject{11653}{<< /D [\MusicPDFPageRef{415} /XYZ 72 157.498 null] >>}
\MusicPDFWriteObject{11654}{<< /D [\MusicPDFPageRef{416} /XYZ 72 698.082 null] >>}
\MusicPDFWriteObject{11655}{<< /D [\MusicPDFPageRef{416} /XYZ 72 226.821 null] >>}
\MusicPDFWriteObject{11656}{<< /D [\MusicPDFPageRef{417} /XYZ 72 599.147 null] >>}
\MusicPDFWriteObject{11657}{<< /D [\MusicPDFPageRef{64} /XYZ 72 356.136 null] >>}
\MusicPDFWriteObject{11658}{<< /D [\MusicPDFPageRef{417} /XYZ 72 549.192 null] >>}
\MusicPDFWriteObject{11659}{<< /D [\MusicPDFPageRef{417} /XYZ 72 538.236 null] >>}
\MusicPDFWriteObject{11660}{<< /D [\MusicPDFPageRef{417} /XYZ 72 486.081 null] >>}
\MusicPDFWriteObject{11661}{<< /D [\MusicPDFPageRef{418} /XYZ 72 618.364 null] >>}
\MusicPDFWriteObject{11662}{<< /D [\MusicPDFPageRef{418} /XYZ 72 518.159 null] >>}
\MusicPDFWriteObject{11663}{<< /D [\MusicPDFPageRef{418} /XYZ 72 340.622 null] >>}
\MusicPDFWriteObject{11664}{<< /D [\MusicPDFPageRef{418} /XYZ 72 297.627 null] >>}
\MusicPDFWriteObject{11665}{<< /D [\MusicPDFPageRef{418} /XYZ 72 181.892 null] >>}
\MusicPDFWriteObject{11666}{<< /D [\MusicPDFPageRef{422} /XYZ 72 358.994 null] >>}
\MusicPDFWriteObject{11667}{<< /D [\MusicPDFPageRef{422} /XYZ 72 343.382 null] >>}
\MusicPDFWriteObject{11668}{<< /D [\MusicPDFPageRef{65} /XYZ 72 498.843 null] >>}
\MusicPDFWriteObject{11669}{<< /D [\MusicPDFPageRef{422} /XYZ 72 328.98 null] >>}
\MusicPDFWriteObject{11670}{<< /D [\MusicPDFPageRef{422} /XYZ 72 302.376 null] >>}
\MusicPDFWriteObject{11671}{<< /D [\MusicPDFPageRef{425} /XYZ 72 670.011 null] >>}
\MusicPDFWriteObject{11672}{<< /D [\MusicPDFPageRef{425} /XYZ 72 628.677 null] >>}
\MusicPDFWriteObject{11673}{<< /D [\MusicPDFPageRef{427} /XYZ 72 474.232 null] >>}
\MusicPDFWriteObject{11674}{<< /D [\MusicPDFPageRef{427} /XYZ 72 460.528 null] >>}
\MusicPDFWriteObject{11675}{<< /D [\MusicPDFPageRef{427} /XYZ 72 447.377 null] >>}
\MusicPDFWriteObject{11676}{<< /D [\MusicPDFPageRef{427} /XYZ 72 408.478 null] >>}
\MusicPDFWriteObject{11677}{<< /D [\MusicPDFPageRef{427} /XYZ 72 394.774 null] >>}
\MusicPDFWriteObject{11678}{<< /D [\MusicPDFPageRef{427} /XYZ 72 381.624 null] >>}
\MusicPDFWriteObject{11679}{<< /D [\MusicPDFPageRef{65} /XYZ 72 474.59 null] >>}
\MusicPDFWriteObject{11680}{<< /D [\MusicPDFPageRef{429} /XYZ 72 509.262 null] >>}
\MusicPDFWriteObject{11681}{<< /D [\MusicPDFPageRef{429} /XYZ 72 431.445 null] >>}
\MusicPDFWriteObject{11682}{<< /D [\MusicPDFPageRef{429} /XYZ 72 391.993 null] >>}
\MusicPDFWriteObject{11683}{<< /D [\MusicPDFPageRef{434} /XYZ 72 315.091 null] >>}
\MusicPDFWriteObject{11684}{<< /D [\MusicPDFPageRef{434} /XYZ 72 251.448 null] >>}
\MusicPDFWriteObject{11685}{<< /D [\MusicPDFPageRef{435} /XYZ 72 459.694 null] >>}
\MusicPDFWriteObject{11686}{<< /D [\MusicPDFPageRef{435} /XYZ 72 442.443 null] >>}
\MusicPDFWriteObject{11687}{<< /D [\MusicPDFPageRef{435} /XYZ 72 178.931 null] >>}
\MusicPDFWriteObject{11688}{<< /D [\MusicPDFPageRef{435} /XYZ 72 151.801 null] >>}
\MusicPDFWriteObject{11689}{<< /D [\MusicPDFPageRef{436} /XYZ 72 640.058 null] >>}
\MusicPDFWriteObject{11690}{<< /D [\MusicPDFPageRef{65} /XYZ 72 448.289 null] >>}
\MusicPDFWriteObject{11691}{<< /D [\MusicPDFPageRef{436} /XYZ 72 626.199 null] >>}
\MusicPDFWriteObject{11692}{<< /D [\MusicPDFPageRef{436} /XYZ 72 272.297 null] >>}
\MusicPDFWriteObject{11693}{<< /D [\MusicPDFPageRef{437} /XYZ 72 624.131 null] >>}
\MusicPDFWriteObject{11694}{<< /D [\MusicPDFPageRef{437} /XYZ 72 314.699 null] >>}
\MusicPDFWriteObject{11695}{<< /D [\MusicPDFPageRef{437} /XYZ 72 157.667 null] >>}
\MusicPDFWriteObject{11696}{<< /D [\MusicPDFPageRef{440} /XYZ 72 355.185 null] >>}
\MusicPDFWriteObject{11697}{<< /D [\MusicPDFPageRef{440} /XYZ 72 339.106 null] >>}
\MusicPDFWriteObject{11698}{<< /D [\MusicPDFPageRef{442} /XYZ 72 337.189 null] >>}
\MusicPDFWriteObject{11699}{<< /D [\MusicPDFPageRef{442} /XYZ 72 323.484 null] >>}
\MusicPDFWriteObject{11700}{<< /D [\MusicPDFPageRef{442} /XYZ 72 308.452 null] >>}
\MusicPDFWriteObject{11701}{<< /D [\MusicPDFPageRef{66} /XYZ 72 547.538 null] >>}
\MusicPDFWriteObject{11702}{<< /D [\MusicPDFPageRef{442} /XYZ 72 281.57 null] >>}
\MusicPDFWriteObject{11703}{<< /D [\MusicPDFPageRef{443} /XYZ 72 250.758 null] >>}
\MusicPDFWriteObject{11704}{<< /D [\MusicPDFPageRef{443} /XYZ 72 235.117 null] >>}
\MusicPDFWriteObject{11705}{<< /D [\MusicPDFPageRef{444} /XYZ 72 517.587 null] >>}
\MusicPDFWriteObject{11706}{<< /D [\MusicPDFPageRef{444} /XYZ 72 505.826 null] >>}
\MusicPDFWriteObject{11707}{<< /D [\MusicPDFPageRef{444} /XYZ 72 482.312 null] >>}
\MusicPDFWriteObject{11708}{<< /D [\MusicPDFPageRef{444} /XYZ 72 378.282 null] >>}
\MusicPDFWriteObject{11709}{<< /D [\MusicPDFPageRef{444} /XYZ 72 362.64 null] >>}
\MusicPDFWriteObject{11710}{<< /D [\MusicPDFPageRef{446} /XYZ 72 636.449 null] >>}
\MusicPDFWriteObject{11711}{<< /D [\MusicPDFPageRef{446} /XYZ 72 620.37 null] >>}
\MusicPDFWriteObject{11712}{<< /D [\MusicPDFPageRef{66} /XYZ 72 485.397 null] >>}
\MusicPDFWriteObject{11713}{<< /D [\MusicPDFPageRef{446} /XYZ 72 180.79 null] >>}
\MusicPDFWriteObject{11714}{<< /D [\MusicPDFPageRef{446} /XYZ 72 165.148 null] >>}
\MusicPDFWriteObject{11715}{<< /D [\MusicPDFPageRef{447} /XYZ 72 474.699 null] >>}
\MusicPDFWriteObject{11716}{<< /D [\MusicPDFPageRef{447} /XYZ 72 459.057 null] >>}
\MusicPDFWriteObject{11717}{<< /D [\MusicPDFPageRef{448} /XYZ 72 347.702 null] >>}
\MusicPDFWriteObject{11718}{<< /D [\MusicPDFPageRef{448} /XYZ 72 269.815 null] >>}
\MusicPDFWriteObject{11719}{<< /D [\MusicPDFPageRef{449} /XYZ 72 698.082 null] >>}
\MusicPDFWriteObject{11720}{<< /D [\MusicPDFPageRef{453} /XYZ 72 160.123 null] >>}
\MusicPDFWriteObject{11721}{<< /D [\MusicPDFPageRef{454} /XYZ 72 698.082 null] >>}
\MusicPDFWriteObject{11722}{<< /D [\MusicPDFPageRef{454} /XYZ 72 510.034 null] >>}
\MusicPDFWriteObject{11723}{<< /D [\MusicPDFPageRef{15} /XYZ 72 602.919 null] >>}
\MusicPDFWriteObject{11724}{<< /D [\MusicPDFPageRef{66} /XYZ 72 425.121 null] >>}
\MusicPDFWriteObject{11725}{<< /D [\MusicPDFPageRef{456} /XYZ 72 511.243 null] >>}
\MusicPDFWriteObject{11726}{<< /D [\MusicPDFPageRef{456} /XYZ 72 436.505 null] >>}
\MusicPDFWriteObject{11727}{<< /D [\MusicPDFPageRef{457} /XYZ 72 530.044 null] >>}
\MusicPDFWriteObject{11728}{<< /D [\MusicPDFPageRef{457} /XYZ 72 429.559 null] >>}
\MusicPDFWriteObject{11729}{<< /D [\MusicPDFPageRef{458} /XYZ 72 502.162 null] >>}
\MusicPDFWriteObject{11730}{<< /D [\MusicPDFPageRef{458} /XYZ 72 439.073 null] >>}
\MusicPDFWriteObject{11731}{<< /D [\MusicPDFPageRef{461} /XYZ 72 235.054 null] >>}
\MusicPDFWriteObject{11732}{<< /D [\MusicPDFPageRef{461} /XYZ 72 198.564 null] >>}
\MusicPDFWriteObject{11733}{<< /D [\MusicPDFPageRef{466} /XYZ 72 162.496 null] >>}
\MusicPDFWriteObject{11734}{<< /D [\MusicPDFPageRef{466} /XYZ 72 122.747 null] >>}
\MusicPDFWriteObject{11735}{<< /D [\MusicPDFPageRef{68} /XYZ 72 465.13 null] >>}
\MusicPDFWriteObject{11736}{<< /D [\MusicPDFPageRef{467} /XYZ 72 698.082 null] >>}
\MusicPDFWriteObject{11737}{<< /D [\MusicPDFPageRef{467} /XYZ 72 626.621 null] >>}
\MusicPDFWriteObject{11738}{<< /D [\MusicPDFPageRef{467} /XYZ 72 418.084 null] >>}
\MusicPDFWriteObject{11739}{<< /D [\MusicPDFPageRef{467} /XYZ 72 335.759 null] >>}
\MusicPDFWriteObject{11740}{<< /D [\MusicPDFPageRef{467} /XYZ 72 192.472 null] >>}
\MusicPDFWriteObject{11741}{<< /D [\MusicPDFPageRef{470} /XYZ 72 369.729 null] >>}
\MusicPDFWriteObject{11742}{<< /D [\MusicPDFPageRef{470} /XYZ 72 316.331 null] >>}
\MusicPDFWriteObject{11743}{<< /D [\MusicPDFPageRef{470} /XYZ 72 161.035 null] >>}
\MusicPDFWriteObject{11744}{<< /D [\MusicPDFPageRef{470} /XYZ 72 112.762 null] >>}
\MusicPDFWriteObject{11745}{<< /D [\MusicPDFPageRef{471} /XYZ 72 657.999 null] >>}
\MusicPDFWriteObject{11746}{<< /D [\MusicPDFPageRef{68} /XYZ 72 423.991 null] >>}
\MusicPDFWriteObject{11747}{<< /D [\MusicPDFPageRef{472} /XYZ 72 634.676 null] >>}
\MusicPDFWriteObject{11748}{<< /D [\MusicPDFPageRef{472} /XYZ 72 594.927 null] >>}
\MusicPDFWriteObject{11749}{<< /D [\MusicPDFPageRef{472} /XYZ 72 555.412 null] >>}
\MusicPDFWriteObject{11750}{<< /D [\MusicPDFPageRef{472} /XYZ 72 476.556 null] >>}
\MusicPDFWriteObject{11751}{<< /D [\MusicPDFPageRef{472} /XYZ 72 276.302 null] >>}
\MusicPDFWriteObject{11752}{<< /D [\MusicPDFPageRef{472} /XYZ 72 195.641 null] >>}
\MusicPDFWriteObject{11753}{<< /D [\MusicPDFPageRef{473} /XYZ 72 673.032 null] >>}
\MusicPDFWriteObject{11754}{<< /D [\MusicPDFPageRef{473} /XYZ 72 460.19 null] >>}
\MusicPDFWriteObject{11755}{<< /D [\MusicPDFPageRef{473} /XYZ 72 408.177 null] >>}
\MusicPDFWriteObject{11756}{<< /D [\MusicPDFPageRef{473} /XYZ 72 333.892 null] >>}
\MusicPDFWriteObject{11757}{<< /D [\MusicPDFPageRef{68} /XYZ 72 363.502 null] >>}
\MusicPDFWriteObject{11758}{<< /D [\MusicPDFPageRef{473} /XYZ 72 203.791 null] >>}
\MusicPDFWriteObject{11759}{<< /D [\MusicPDFPageRef{473} /XYZ 72 113.462 null] >>}
\MusicPDFWriteObject{11760}{<< /D [\MusicPDFPageRef{476} /XYZ 72 565.184 null] >>}
\MusicPDFWriteObject{11761}{<< /D [\MusicPDFPageRef{476} /XYZ 72 526.029 null] >>}
\MusicPDFWriteObject{11762}{<< /D [\MusicPDFPageRef{476} /XYZ 72 463.76 null] >>}
\MusicPDFWriteObject{11763}{<< /D [\MusicPDFPageRef{476} /XYZ 72 387.621 null] >>}
\MusicPDFWriteObject{11764}{<< /D [\MusicPDFPageRef{478} /XYZ 72 587.874 null] >>}
\MusicPDFWriteObject{11765}{<< /D [\MusicPDFPageRef{478} /XYZ 72 572.233 null] >>}
\MusicPDFWriteObject{11766}{<< /D [\MusicPDFPageRef{478} /XYZ 72 548.422 null] >>}
\MusicPDFWriteObject{11767}{<< /D [\MusicPDFPageRef{478} /XYZ 72 530.318 null] >>}
\MusicPDFWriteObject{11768}{<< /D [\MusicPDFPageRef{68} /XYZ 72 349.798 null] >>}
\MusicPDFWriteObject{11769}{<< /D [\MusicPDFPageRef{478} /XYZ 72 99.233 null] >>}
\MusicPDFWriteObject{11770}{<< /D [\MusicPDFPageRef{479} /XYZ 72 698.082 null] >>}
\MusicPDFWriteObject{11771}{<< /D [\MusicPDFPageRef{479} /XYZ 72 674.969 null] >>}
\MusicPDFWriteObject{11772}{<< /D [\MusicPDFPageRef{479} /XYZ 72 633.28 null] >>}
\MusicPDFWriteObject{11773}{<< /D [\MusicPDFPageRef{480} /XYZ 72 436.788 null] >>}
\MusicPDFWriteObject{11774}{<< /D [\MusicPDFPageRef{480} /XYZ 72 367.374 null] >>}
\MusicPDFWriteObject{11775}{<< /D [\MusicPDFPageRef{480} /XYZ 72 228.211 null] >>}
\MusicPDFWriteObject{11776}{<< /D [\MusicPDFPageRef{484} /XYZ 72 468.02 null] >>}
\MusicPDFWriteObject{11777}{<< /D [\MusicPDFPageRef{484} /XYZ 72 441.451 null] >>}
\MusicPDFWriteObject{11778}{<< /D [\MusicPDFPageRef{485} /XYZ 72 384.223 null] >>}
\MusicPDFWriteObject{11779}{<< /D [\MusicPDFPageRef{68} /XYZ 72 336.647 null] >>}
\MusicPDFWriteObject{11780}{<< /D [\MusicPDFPageRef{485} /XYZ 72 339.712 null] >>}
\MusicPDFWriteObject{11781}{<< /D [\MusicPDFPageRef{487} /XYZ 72 191.555 null] >>}
\MusicPDFWriteObject{11782}{<< /D [\MusicPDFPageRef{488} /XYZ 72 667.172 null] >>}
\MusicPDFWriteObject{11783}{<< /D [\MusicPDFPageRef{488} /XYZ 72 601.972 null] >>}
\MusicPDFWriteObject{11784}{<< /D [\MusicPDFPageRef{488} /XYZ 72 485.978 null] >>}
\MusicPDFWriteObject{11785}{<< /D [\MusicPDFPageRef{488} /XYZ 72 468.727 null] >>}
\MusicPDFWriteObject{11786}{<< /D [\MusicPDFPageRef{488} /XYZ 72 189.81 null] >>}
\MusicPDFWriteObject{11787}{<< /D [\MusicPDFPageRef{489} /XYZ 72 668.002 null] >>}
\MusicPDFWriteObject{11788}{<< /D [\MusicPDFPageRef{489} /XYZ 72 603.675 null] >>}
\MusicPDFWriteObject{11789}{<< /D [\MusicPDFPageRef{489} /XYZ 72 467.556 null] >>}
\MusicPDFWriteObject{11790}{<< /D [\MusicPDFPageRef{68} /XYZ 72 296.835 null] >>}
\MusicPDFWriteObject{11791}{<< /D [\MusicPDFPageRef{489} /XYZ 72 450.305 null] >>}
\MusicPDFWriteObject{11792}{<< /D [\MusicPDFPageRef{491} /XYZ 72 163.91 null] >>}
\MusicPDFWriteObject{11793}{<< /D [\MusicPDFPageRef{491} /XYZ 72 150.206 null] >>}
\MusicPDFWriteObject{11794}{<< /D [\MusicPDFPageRef{491} /XYZ 72 137.056 null] >>}
\MusicPDFWriteObject{11795}{<< /D [\MusicPDFPageRef{492} /XYZ 72 136.195 null] >>}
\MusicPDFWriteObject{11796}{<< /D [\MusicPDFPageRef{492} /XYZ 72 109.894 null] >>}
\MusicPDFWriteObject{11797}{<< /D [\MusicPDFPageRef{492} /XYZ 72 96.743 null] >>}
\MusicPDFWriteObject{11798}{<< /D [\MusicPDFPageRef{494} /XYZ 72 562.763 null] >>}
\MusicPDFWriteObject{11799}{<< /D [\MusicPDFPageRef{494} /XYZ 72 341.484 null] >>}
\MusicPDFWriteObject{11800}{<< /D [\MusicPDFPageRef{499} /XYZ 72 487.719 null] >>}
\MusicPDFWriteObject{11801}{<< /D [\MusicPDFPageRef{71} /XYZ 72 657.347 null] >>}
\MusicPDFWriteObject{11802}{<< /D [\MusicPDFPageRef{499} /XYZ 72 427.337 null] >>}
\MusicPDFWriteObject{11803}{<< /D [\MusicPDFPageRef{499} /XYZ 72 341.488 null] >>}
\MusicPDFWriteObject{11804}{<< /D [\MusicPDFPageRef{499} /XYZ 72 290.439 null] >>}
\MusicPDFWriteObject{11805}{<< /D [\MusicPDFPageRef{502} /XYZ 72 151.248 null] >>}
\MusicPDFWriteObject{11806}{<< /D [\MusicPDFPageRef{502} /XYZ 72 124.947 null] >>}
\MusicPDFWriteObject{11807}{<< /D [\MusicPDFPageRef{503} /XYZ 72 653.358 null] >>}
\MusicPDFWriteObject{11808}{<< /D [\MusicPDFPageRef{504} /XYZ 72 123.598 null] >>}
\MusicPDFWriteObject{11809}{<< /D [\MusicPDFPageRef{504} /XYZ 72 109.894 null] >>}
\MusicPDFWriteObject{11810}{<< /D [\MusicPDFPageRef{505} /XYZ 72 698.082 null] >>}
\MusicPDFWriteObject{11811}{<< /D [\MusicPDFPageRef{506} /XYZ 72 584.618 null] >>}
\MusicPDFWriteObject{11812}{<< /D [\MusicPDFPageRef{71} /XYZ 72 584.547 null] >>}
\MusicPDFWriteObject{11813}{<< /D [\MusicPDFPageRef{506} /XYZ 72 524.988 null] >>}
\MusicPDFWriteObject{11814}{<< /D [\MusicPDFPageRef{508} /XYZ 72 465.177 null] >>}
\MusicPDFWriteObject{11815}{<< /D [\MusicPDFPageRef{508} /XYZ 72 393.248 null] >>}
\MusicPDFWriteObject{11816}{<< /D [\MusicPDFPageRef{508} /XYZ 72 130.978 null] >>}
\MusicPDFWriteObject{11817}{<< /D [\MusicPDFPageRef{509} /XYZ 72 698.082 null] >>}
\MusicPDFWriteObject{11818}{<< /D [\MusicPDFPageRef{512} /XYZ 72 519.502 null] >>}
\MusicPDFWriteObject{11819}{<< /D [\MusicPDFPageRef{512} /XYZ 72 503.861 null] >>}
\MusicPDFWriteObject{11820}{<< /D [\MusicPDFPageRef{512} /XYZ 72 490.711 null] >>}
\MusicPDFWriteObject{11821}{<< /D [\MusicPDFPageRef{512} /XYZ 72 477.56 null] >>}
\MusicPDFWriteObject{11822}{<< /D [\MusicPDFPageRef{515} /XYZ 72 183.958 null] >>}
\MusicPDFWriteObject{11823}{<< /D [\MusicPDFPageRef{71} /XYZ 72 558.799 null] >>}
\MusicPDFWriteObject{11824}{<< /D [\MusicPDFPageRef{515} /XYZ 72 156.632 null] >>}
\MusicPDFWriteObject{11825}{<< /D [\MusicPDFPageRef{516} /XYZ 72 698.082 null] >>}
\MusicPDFWriteObject{11826}{<< /D [\MusicPDFPageRef{525} /XYZ 72 455.105 null] >>}
\MusicPDFWriteObject{11827}{<< /D [\MusicPDFPageRef{525} /XYZ 72 392.376 null] >>}
\MusicPDFWriteObject{11828}{<< /D [\MusicPDFPageRef{525} /XYZ 72 319.771 null] >>}
\MusicPDFWriteObject{11829}{<< /D [\MusicPDFPageRef{525} /XYZ 72 243.891 null] >>}
\MusicPDFWriteObject{11830}{<< /D [\MusicPDFPageRef{525} /XYZ 72 202.833 null] >>}
\MusicPDFWriteObject{11831}{<< /D [\MusicPDFPageRef{525} /XYZ 72 180.074 null] >>}
\MusicPDFWriteObject{11832}{<< /D [\MusicPDFPageRef{525} /XYZ 72 151.282 null] >>}
\MusicPDFWriteObject{11833}{<< /D [\MusicPDFPageRef{525} /XYZ 72 125.534 null] >>}
\MusicPDFWriteObject{11834}{<< /D [\MusicPDFPageRef{16} /XYZ 72 568.131 null] >>}
\MusicPDFWriteObject{11835}{<< /D [\MusicPDFPageRef{71} /XYZ 72 446.482 null] >>}
\MusicPDFWriteObject{11836}{<< /D [\MusicPDFPageRef{529} /XYZ 72 333.373 null] >>}
\MusicPDFWriteObject{11837}{<< /D [\MusicPDFPageRef{529} /XYZ 72 317.731 null] >>}
\MusicPDFWriteObject{11838}{<< /D [\MusicPDFPageRef{530} /XYZ 72 437.368 null] >>}
\MusicPDFWriteObject{11839}{<< /D [\MusicPDFPageRef{530} /XYZ 72 421.727 null] >>}
\MusicPDFWriteObject{11840}{<< /D [\MusicPDFPageRef{531} /XYZ 72 632.969 null] >>}
\MusicPDFWriteObject{11841}{<< /D [\MusicPDFPageRef{531} /XYZ 72 617.328 null] >>}
\MusicPDFWriteObject{11842}{<< /D [\MusicPDFPageRef{532} /XYZ 72 683.487 null] >>}
\MusicPDFWriteObject{11843}{<< /D [\MusicPDFPageRef{532} /XYZ 72 669.783 null] >>}
\MusicPDFWriteObject{11844}{<< /D [\MusicPDFPageRef{537} /XYZ 72 486.186 null] >>}
\MusicPDFWriteObject{11845}{<< /D [\MusicPDFPageRef{537} /XYZ 72 459.331 null] >>}
\MusicPDFWriteObject{11846}{<< /D [\MusicPDFPageRef{71} /XYZ 72 433.331 null] >>}
\MusicPDFWriteObject{11847}{<< /D [\MusicPDFPageRef{537} /XYZ 72 446.181 null] >>}
\MusicPDFWriteObject{11848}{<< /D [\MusicPDFPageRef{539} /XYZ 72 97.456 null] >>}
\MusicPDFWriteObject{11849}{<< /D [\MusicPDFPageRef{540} /XYZ 72 698.082 null] >>}
\MusicPDFWriteObject{11850}{<< /D [\MusicPDFPageRef{540} /XYZ 72 642.594 null] >>}
\MusicPDFWriteObject{11851}{<< /D [\MusicPDFPageRef{540} /XYZ 72 601.402 null] >>}
\MusicPDFWriteObject{11852}{<< /D [\MusicPDFPageRef{540} /XYZ 72 540.459 null] >>}
\MusicPDFWriteObject{11853}{<< /D [\MusicPDFPageRef{540} /XYZ 72 458.451 null] >>}
\MusicPDFWriteObject{11854}{<< /D [\MusicPDFPageRef{540} /XYZ 72 430.213 null] >>}
\MusicPDFWriteObject{11855}{<< /D [\MusicPDFPageRef{540} /XYZ 72 110.447 null] >>}
\MusicPDFWriteObject{11856}{<< /D [\MusicPDFPageRef{540} /XYZ 72 96.743 null] >>}
\MusicPDFWriteObject{11857}{<< /D [\MusicPDFPageRef{73} /XYZ 72 664.347 null] >>}
\MusicPDFWriteObject{11858}{<< /D [\MusicPDFPageRef{541} /XYZ 72 698.082 null] >>}
\MusicPDFWriteObject{11859}{<< /D [\MusicPDFPageRef{542} /XYZ 72 528.301 null] >>}
\MusicPDFWriteObject{11860}{<< /D [\MusicPDFPageRef{542} /XYZ 72 515.15 null] >>}
\MusicPDFWriteObject{11861}{<< /D [\MusicPDFPageRef{542} /XYZ 72 501.999 null] >>}
\MusicPDFWriteObject{11862}{<< /D [\MusicPDFPageRef{542} /XYZ 72 97.296 null] >>}
\MusicPDFWriteObject{11863}{<< /D [\MusicPDFPageRef{543} /XYZ 72 698.082 null] >>}
\MusicPDFWriteObject{11864}{<< /D [\MusicPDFPageRef{543} /XYZ 72 685.192 null] >>}
\MusicPDFWriteObject{11865}{<< /D [\MusicPDFPageRef{543} /XYZ 72 214.426 null] >>}
\MusicPDFWriteObject{11866}{<< /D [\MusicPDFPageRef{543} /XYZ 72 201.635 null] >>}
\MusicPDFWriteObject{11867}{<< /D [\MusicPDFPageRef{543} /XYZ 72 188.124 null] >>}
\MusicPDFWriteObject{11868}{<< /D [\MusicPDFPageRef{73} /XYZ 72 600.81 null] >>}
\MusicPDFWriteObject{11869}{<< /D [\MusicPDFPageRef{544} /XYZ 72 344.668 null] >>}
\MusicPDFWriteObject{11870}{<< /D [\MusicPDFPageRef{544} /XYZ 72 329.354 null] >>}
\MusicPDFWriteObject{11871}{<< /D [\MusicPDFPageRef{544} /XYZ 72 315.495 null] >>}
\MusicPDFWriteObject{11872}{<< /D [\MusicPDFPageRef{545} /XYZ 72 245.518 null] >>}
\MusicPDFWriteObject{11873}{<< /D [\MusicPDFPageRef{545} /XYZ 72 182.236 null] >>}
\MusicPDFWriteObject{11874}{<< /D [\MusicPDFPageRef{548} /XYZ 72 335.86 null] >>}
\MusicPDFWriteObject{11875}{<< /D [\MusicPDFPageRef{548} /XYZ 72 322.349 null] >>}
\MusicPDFWriteObject{11876}{<< /D [\MusicPDFPageRef{548} /XYZ 72 296.407 null] >>}
\MusicPDFWriteObject{11877}{<< /D [\MusicPDFPageRef{548} /XYZ 72 192.582 null] >>}
\MusicPDFWriteObject{11878}{<< /D [\MusicPDFPageRef{548} /XYZ 72 178.44 null] >>}
\MusicPDFWriteObject{11879}{<< /D [\MusicPDFPageRef{73} /XYZ 72 465.12 null] >>}
\MusicPDFWriteObject{11880}{<< /D [\MusicPDFPageRef{548} /XYZ 72 153.129 null] >>}
\MusicPDFWriteObject{11881}{<< /D [\MusicPDFPageRef{549} /XYZ 72 509.245 null] >>}
\MusicPDFWriteObject{11882}{<< /D [\MusicPDFPageRef{549} /XYZ 72 496.094 null] >>}
\MusicPDFWriteObject{11883}{<< /D [\MusicPDFPageRef{550} /XYZ 72 660.347 null] >>}
\MusicPDFWriteObject{11884}{<< /D [\MusicPDFPageRef{550} /XYZ 72 634.959 null] >>}
\MusicPDFWriteObject{11885}{<< /D [\MusicPDFPageRef{550} /XYZ 72 584.389 null] >>}
\MusicPDFWriteObject{11886}{<< /D [\MusicPDFPageRef{552} /XYZ 72 466.211 null] >>}
\MusicPDFWriteObject{11887}{<< /D [\MusicPDFPageRef{552} /XYZ 72 452.807 null] >>}
\MusicPDFWriteObject{11888}{<< /D [\MusicPDFPageRef{552} /XYZ 72 380.206 null] >>}
\MusicPDFWriteObject{11889}{<< /D [\MusicPDFPageRef{553} /XYZ 72 409.133 null] >>}
\MusicPDFWriteObject{11890}{<< /D [\MusicPDFPageRef{77} /XYZ 72 215.726 null] >>}
\MusicPDFWriteObject{11891}{<< /D [\MusicPDFPageRef{553} /XYZ 72 382.831 null] >>}
\MusicPDFWriteObject{11892}{<< /D [\MusicPDFPageRef{553} /XYZ 72 368.546 null] >>}
\MusicPDFWriteObject{11893}{<< /D [\MusicPDFPageRef{554} /XYZ 72 511.137 null] >>}
\MusicPDFWriteObject{11894}{<< /D [\MusicPDFPageRef{554} /XYZ 72 445.557 null] >>}
\MusicPDFWriteObject{11895}{<< /D [\MusicPDFPageRef{555} /XYZ 72 532.588 null] >>}
\MusicPDFWriteObject{11896}{<< /D [\MusicPDFPageRef{555} /XYZ 72 516.587 null] >>}
\MusicPDFWriteObject{11897}{<< /D [\MusicPDFPageRef{555} /XYZ 72 443.023 null] >>}
\MusicPDFWriteObject{11898}{<< /D [\MusicPDFPageRef{555} /XYZ 72 392.515 null] >>}
\MusicPDFWriteObject{11899}{<< /D [\MusicPDFPageRef{556} /XYZ 72 553.182 null] >>}
\MusicPDFWriteObject{11900}{<< /D [\MusicPDFPageRef{556} /XYZ 72 527.876 null] >>}
\MusicPDFWriteObject{11901}{<< /D [\MusicPDFPageRef{77} /XYZ 72 152.997 null] >>}
\MusicPDFWriteObject{11902}{<< /D [\MusicPDFPageRef{556} /XYZ 72 501.133 null] >>}
\MusicPDFWriteObject{11903}{<< /D [\MusicPDFPageRef{560} /XYZ 72 587.462 null] >>}
\MusicPDFWriteObject{11904}{<< /D [\MusicPDFPageRef{560} /XYZ 72 572.555 null] >>}
\MusicPDFWriteObject{11905}{<< /D [\MusicPDFPageRef{561} /XYZ 72 330.523 null] >>}
\MusicPDFWriteObject{11906}{<< /D [\MusicPDFPageRef{561} /XYZ 72 258.943 null] >>}
\MusicPDFWriteObject{11907}{<< /D [\MusicPDFPageRef{561} /XYZ 72 185.342 null] >>}
\MusicPDFWriteObject{11908}{<< /D [\MusicPDFPageRef{561} /XYZ 72 172.164 null] >>}
\MusicPDFWriteObject{11909}{<< /D [\MusicPDFPageRef{561} /XYZ 72 140.882 null] >>}
\MusicPDFWriteObject{11910}{<< /D [\MusicPDFPageRef{561} /XYZ 72 113.147 null] >>}
\MusicPDFWriteObject{11911}{<< /D [\MusicPDFPageRef{562} /XYZ 72 257.313 null] >>}
\MusicPDFWriteObject{11912}{<< /D [\MusicPDFPageRef{77} /XYZ 72 126.695 null] >>}
\MusicPDFWriteObject{11913}{<< /D [\MusicPDFPageRef{562} /XYZ 72 219.972 null] >>}
\MusicPDFWriteObject{11914}{<< /D [\MusicPDFPageRef{566} /XYZ 72 272.835 null] >>}
\MusicPDFWriteObject{11915}{<< /D [\MusicPDFPageRef{566} /XYZ 72 246.533 null] >>}
\MusicPDFWriteObject{11916}{<< /D [\MusicPDFPageRef{566} /XYZ 72 201.712 null] >>}
\MusicPDFWriteObject{11917}{<< /D [\MusicPDFPageRef{569} /XYZ 72 356.451 null] >>}
\MusicPDFWriteObject{11918}{<< /D [\MusicPDFPageRef{569} /XYZ 72 343.597 null] >>}
\MusicPDFWriteObject{11919}{<< /D [\MusicPDFPageRef{569} /XYZ 72 330.447 null] >>}
\MusicPDFWriteObject{11920}{<< /D [\MusicPDFPageRef{570} /XYZ 72 672.877 null] >>}
\MusicPDFWriteObject{11921}{<< /D [\MusicPDFPageRef{570} /XYZ 72 643.57 null] >>}
\MusicPDFWriteObject{11922}{<< /D [\MusicPDFPageRef{575} /XYZ 72 365.443 null] >>}
\MusicPDFWriteObject{11923}{<< /D [\MusicPDFPageRef{79} /XYZ 72 601.761 null] >>}
\MusicPDFWriteObject{11924}{<< /D [\MusicPDFPageRef{575} /XYZ 72 336.796 null] >>}
\MusicPDFWriteObject{11925}{<< /D [\MusicPDFPageRef{575} /XYZ 72 322.937 null] >>}
\MusicPDFWriteObject{11926}{<< /D [\MusicPDFPageRef{576} /XYZ 72 633.287 null] >>}
\MusicPDFWriteObject{11927}{<< /D [\MusicPDFPageRef{576} /XYZ 72 619.222 null] >>}
\MusicPDFWriteObject{11928}{<< /D [\MusicPDFPageRef{576} /XYZ 72 319.933 null] >>}
\MusicPDFWriteObject{11929}{<< /D [\MusicPDFPageRef{576} /XYZ 72 240.87 null] >>}
\MusicPDFWriteObject{11930}{<< /D [\MusicPDFPageRef{577} /XYZ 72 525.024 null] >>}
\MusicPDFWriteObject{11931}{<< /D [\MusicPDFPageRef{577} /XYZ 72 460.109 null] >>}
\MusicPDFWriteObject{11932}{<< /D [\MusicPDFPageRef{577} /XYZ 72 433.807 null] >>}
\MusicPDFWriteObject{11933}{<< /D [\MusicPDFPageRef{577} /XYZ 72 220.588 null] >>}
\MusicPDFWriteObject{11934}{<< /D [\MusicPDFPageRef{79} /XYZ 72 546.308 null] >>}
\MusicPDFWriteObject{11935}{<< /D [\MusicPDFPageRef{577} /XYZ 72 157.561 null] >>}
\MusicPDFWriteObject{11936}{<< /D [\MusicPDFPageRef{578} /XYZ 72 566.194 null] >>}
\MusicPDFWriteObject{11937}{<< /D [\MusicPDFPageRef{578} /XYZ 72 552.129 null] >>}
\MusicPDFWriteObject{11938}{<< /D [\MusicPDFPageRef{578} /XYZ 72 538.978 null] >>}
\MusicPDFWriteObject{11939}{<< /D [\MusicPDFPageRef{578} /XYZ 72 525.23 null] >>}
\MusicPDFWriteObject{11940}{<< /D [\MusicPDFPageRef{578} /XYZ 72 510.83 null] >>}
\MusicPDFWriteObject{11941}{<< /D [\MusicPDFPageRef{578} /XYZ 72 498.929 null] >>}
\MusicPDFWriteObject{11942}{<< /D [\MusicPDFPageRef{578} /XYZ 72 485.778 null] >>}
\MusicPDFWriteObject{11943}{<< /D [\MusicPDFPageRef{578} /XYZ 72 472.627 null] >>}
\MusicPDFWriteObject{11944}{<< /D [\MusicPDFPageRef{579} /XYZ 72 698.082 null] >>}
\MusicPDFWriteObject{11945}{<< /D [\MusicPDFPageRef{16} /XYZ 72 358.522 null] >>}
\MusicPDFWriteObject{11946}{<< /D [\MusicPDFPageRef{79} /XYZ 72 197.209 null] >>}
\MusicPDFWriteObject{11947}{<< /D [\MusicPDFPageRef{579} /XYZ 72 685.629 null] >>}
\MusicPDFWriteObject{11948}{<< /D [\MusicPDFPageRef{579} /XYZ 72 181.749 null] >>}
\MusicPDFWriteObject{11949}{<< /D [\MusicPDFPageRef{580} /XYZ 72 698.082 null] >>}
\MusicPDFWriteObject{11950}{<< /D [\MusicPDFPageRef{581} /XYZ 72 499.28 null] >>}
\MusicPDFWriteObject{11951}{<< /D [\MusicPDFPageRef{581} /XYZ 72 437.772 null] >>}
\MusicPDFWriteObject{11952}{<< /D [\MusicPDFPageRef{582} /XYZ 72 394.674 null] >>}
\MusicPDFWriteObject{11953}{<< /D [\MusicPDFPageRef{582} /XYZ 72 347.528 null] >>}
\MusicPDFWriteObject{11954}{<< /D [\MusicPDFPageRef{582} /XYZ 72 168.602 null] >>}
\MusicPDFWriteObject{11955}{<< /D [\MusicPDFPageRef{582} /XYZ 72 111.171 null] >>}
\MusicPDFWriteObject{11956}{<< /D [\MusicPDFPageRef{582} /XYZ 72 97.169 null] >>}
\MusicPDFWriteObject{11957}{<< /D [\MusicPDFPageRef{79} /XYZ 72 183.145 null] >>}
\MusicPDFWriteObject{11958}{<< /D [\MusicPDFPageRef{589} /XYZ 72 390.753 null] >>}
\MusicPDFWriteObject{11959}{<< /D [\MusicPDFPageRef{589} /XYZ 72 377.135 null] >>}
\MusicPDFWriteObject{11960}{<< /D [\MusicPDFPageRef{593} /XYZ 72 555.688 null] >>}
\MusicPDFWriteObject{11961}{<< /D [\MusicPDFPageRef{593} /XYZ 72 489.935 null] >>}
\MusicPDFWriteObject{11962}{<< /D [\MusicPDFPageRef{593} /XYZ 72 231.95 null] >>}
\MusicPDFWriteObject{11963}{<< /D [\MusicPDFPageRef{593} /XYZ 72 218.799 null] >>}
\MusicPDFWriteObject{11964}{<< /D [\MusicPDFPageRef{593} /XYZ 72 192.498 null] >>}
\MusicPDFWriteObject{11965}{<< /D [\MusicPDFPageRef{593} /XYZ 72 179.347 null] >>}
\MusicPDFWriteObject{11966}{<< /D [\MusicPDFPageRef{594} /XYZ 72 419.358 null] >>}
\MusicPDFWriteObject{11967}{<< /D [\MusicPDFPageRef{594} /XYZ 72 405.739 null] >>}
\MusicPDFWriteObject{11968}{<< /D [\MusicPDFPageRef{83} /XYZ 72 698.082 null] >>}
\MusicPDFWriteObject{11969}{<< /D [\MusicPDFPageRef{594} /XYZ 72 329.859 null] >>}
\MusicPDFWriteObject{11970}{<< /D [\MusicPDFPageRef{595} /XYZ 72 140.149 null] >>}
\MusicPDFWriteObject{11971}{<< /D [\MusicPDFPageRef{595} /XYZ 72 128.05 null] >>}
\MusicPDFWriteObject{11972}{<< /D [\MusicPDFPageRef{596} /XYZ 72 698.082 null] >>}
\MusicPDFWriteObject{11973}{<< /D [\MusicPDFPageRef{596} /XYZ 72 684.577 null] >>}
\MusicPDFWriteObject{11974}{<< /D [\MusicPDFPageRef{596} /XYZ 72 623.112 null] >>}
\MusicPDFWriteObject{11975}{<< /D [\MusicPDFPageRef{598} /XYZ 72 173.176 null] >>}
\MusicPDFWriteObject{11976}{<< /D [\MusicPDFPageRef{598} /XYZ 72 159.472 null] >>}
\MusicPDFWriteObject{11977}{<< /D [\MusicPDFPageRef{599} /XYZ 72 698.082 null] >>}
\MusicPDFWriteObject{11978}{<< /D [\MusicPDFPageRef{599} /XYZ 72 641.852 null] >>}
\MusicPDFWriteObject{11979}{<< /D [\MusicPDFPageRef{83} /XYZ 72 649.201 null] >>}
\MusicPDFWriteObject{11980}{<< /D [\MusicPDFPageRef{600} /XYZ 72 625.675 null] >>}
\MusicPDFWriteObject{11981}{<< /D [\MusicPDFPageRef{600} /XYZ 72 598.321 null] >>}
\MusicPDFWriteObject{11982}{<< /D [\MusicPDFPageRef{601} /XYZ 72 539.257 null] >>}
\MusicPDFWriteObject{11983}{<< /D [\MusicPDFPageRef{601} /XYZ 72 509.413 null] >>}
\MusicPDFWriteObject{11984}{<< /D [\MusicPDFPageRef{601} /XYZ 72 127.708 null] >>}
\MusicPDFWriteObject{11985}{<< /D [\MusicPDFPageRef{602} /XYZ 72 635.597 null] >>}
\MusicPDFWriteObject{11986}{<< /D [\MusicPDFPageRef{604} /XYZ 72 393.577 null] >>}
\MusicPDFWriteObject{11987}{<< /D [\MusicPDFPageRef{604} /XYZ 72 367.276 null] >>}
\MusicPDFWriteObject{11988}{<< /D [\MusicPDFPageRef{604} /XYZ 72 352.519 null] >>}
\MusicPDFWriteObject{11989}{<< /D [\MusicPDFPageRef{604} /XYZ 72 339.369 null] >>}
\MusicPDFWriteObject{11990}{<< /D [\MusicPDFPageRef{83} /XYZ 72 599.463 null] >>}
\MusicPDFWriteObject{11991}{<< /D [\MusicPDFPageRef{604} /XYZ 72 301.522 null] >>}
\MusicPDFWriteObject{11992}{<< /D [\MusicPDFPageRef{604} /XYZ 72 251.03 null] >>}
\MusicPDFWriteObject{11993}{<< /D [\MusicPDFPageRef{606} /XYZ 72 698.082 null] >>}
\MusicPDFWriteObject{11994}{<< /D [\MusicPDFPageRef{606} /XYZ 72 637.081 null] >>}
\MusicPDFWriteObject{11995}{<< /D [\MusicPDFPageRef{607} /XYZ 72 472.85 null] >>}
\MusicPDFWriteObject{11996}{<< /D [\MusicPDFPageRef{607} /XYZ 72 430.935 null] >>}
\MusicPDFWriteObject{11997}{<< /D [\MusicPDFPageRef{617} /XYZ 72 231.138 null] >>}
\MusicPDFWriteObject{11998}{<< /D [\MusicPDFPageRef{617} /XYZ 72 126.302 null] >>}
\MusicPDFWriteObject{11999}{<< /D [\MusicPDFPageRef{619} /XYZ 72 594.643 null] >>}
\MusicPDFWriteObject{12000}{<< /D [\MusicPDFPageRef{619} /XYZ 72 528.97 null] >>}
\MusicPDFWriteObject{12001}{<< /D [\MusicPDFPageRef{83} /XYZ 72 536.574 null] >>}
\MusicPDFWriteObject{12002}{<< /D [\MusicPDFPageRef{619} /XYZ 72 281.703 null] >>}
\MusicPDFWriteObject{12003}{<< /D [\MusicPDFPageRef{620} /XYZ 72 698.082 null] >>}
\MusicPDFWriteObject{12004}{<< /D [\MusicPDFPageRef{622} /XYZ 72 657.368 null] >>}
\MusicPDFWriteObject{12005}{<< /D [\MusicPDFPageRef{622} /XYZ 72 606.217 null] >>}
\MusicPDFWriteObject{12006}{<< /D [\MusicPDFPageRef{626} /XYZ 72 214.748 null] >>}
\MusicPDFWriteObject{12007}{<< /D [\MusicPDFPageRef{626} /XYZ 72 163.251 null] >>}
\MusicPDFWriteObject{12008}{<< /D [\MusicPDFPageRef{627} /XYZ 72 265.71 null] >>}
\MusicPDFWriteObject{12009}{<< /D [\MusicPDFPageRef{627} /XYZ 72 239.408 null] >>}
\MusicPDFWriteObject{12010}{<< /D [\MusicPDFPageRef{628} /XYZ 72 431.038 null] >>}
\MusicPDFWriteObject{12011}{<< /D [\MusicPDFPageRef{628} /XYZ 72 339.436 null] >>}
\MusicPDFWriteObject{12012}{<< /D [\MusicPDFPageRef{83} /XYZ 72 496.762 null] >>}
\MusicPDFWriteObject{12013}{<< /D [\MusicPDFPageRef{632} /XYZ 72 522.451 null] >>}
\MusicPDFWriteObject{12014}{<< /D [\MusicPDFPageRef{632} /XYZ 72 473.777 null] >>}
\MusicPDFWriteObject{12015}{<< /D [\MusicPDFPageRef{632} /XYZ 72 243.63 null] >>}
\MusicPDFWriteObject{12016}{<< /D [\MusicPDFPageRef{632} /XYZ 72 230.479 null] >>}
\MusicPDFWriteObject{12017}{<< /D [\MusicPDFPageRef{632} /XYZ 72 216.277 null] >>}
\MusicPDFWriteObject{12018}{<< /D [\MusicPDFPageRef{633} /XYZ 72 495.67 null] >>}
\MusicPDFWriteObject{12019}{<< /D [\MusicPDFPageRef{633} /XYZ 72 470.916 null] >>}
\MusicPDFWriteObject{12020}{<< /D [\MusicPDFPageRef{633} /XYZ 72 246.021 null] >>}
\MusicPDFWriteObject{12021}{<< /D [\MusicPDFPageRef{633} /XYZ 72 232.317 null] >>}
\MusicPDFWriteObject{12022}{<< /D [\MusicPDFPageRef{637} /XYZ 72 518.244 null] >>}
\MusicPDFWriteObject{12023}{<< /D [\MusicPDFPageRef{87} /XYZ 72 647.663 null] >>}
\MusicPDFWriteObject{12024}{<< /D [\MusicPDFPageRef{637} /XYZ 72 467.038 null] >>}
\MusicPDFWriteObject{12025}{<< /D [\MusicPDFPageRef{637} /XYZ 72 250.975 null] >>}
\MusicPDFWriteObject{12026}{<< /D [\MusicPDFPageRef{637} /XYZ 72 236.973 null] >>}
\MusicPDFWriteObject{12027}{<< /D [\MusicPDFPageRef{637} /XYZ 72 224.119 null] >>}
\MusicPDFWriteObject{12028}{<< /D [\MusicPDFPageRef{644} /XYZ 72 617.993 null] >>}
\MusicPDFWriteObject{12029}{<< /D [\MusicPDFPageRef{644} /XYZ 72 604.374 null] >>}
\MusicPDFWriteObject{12030}{<< /D [\MusicPDFPageRef{645} /XYZ 72 222.242 null] >>}
\MusicPDFWriteObject{12031}{<< /D [\MusicPDFPageRef{646} /XYZ 72 698.082 null] >>}
\MusicPDFWriteObject{12032}{<< /D [\MusicPDFPageRef{649} /XYZ 72 596.577 null] >>}
\MusicPDFWriteObject{12033}{<< /D [\MusicPDFPageRef{649} /XYZ 72 414.252 null] >>}
\MusicPDFWriteObject{12034}{<< /D [\MusicPDFPageRef{87} /XYZ 72 620.808 null] >>}
\MusicPDFWriteObject{12035}{<< /D [\MusicPDFPageRef{652} /XYZ 72 488.486 null] >>}
\MusicPDFWriteObject{12036}{<< /D [\MusicPDFPageRef{652} /XYZ 72 238.458 null] >>}
\MusicPDFWriteObject{12037}{<< /D [\MusicPDFPageRef{657} /XYZ 72 409.12 null] >>}
\MusicPDFWriteObject{12038}{<< /D [\MusicPDFPageRef{657} /XYZ 72 341.021 null] >>}
\MusicPDFWriteObject{12039}{<< /D [\MusicPDFPageRef{657} /XYZ 72 316.511 null] >>}
\MusicPDFWriteObject{12040}{<< /D [\MusicPDFPageRef{657} /XYZ 72 211.306 null] >>}
\MusicPDFWriteObject{12041}{<< /D [\MusicPDFPageRef{658} /XYZ 72 618.337 null] >>}
\MusicPDFWriteObject{12042}{<< /D [\MusicPDFPageRef{658} /XYZ 72 604.632 null] >>}
\MusicPDFWriteObject{12043}{<< /D [\MusicPDFPageRef{658} /XYZ 72 591.481 null] >>}
\MusicPDFWriteObject{12044}{<< /D [\MusicPDFPageRef{659} /XYZ 72 519.243 null] >>}
\MusicPDFWriteObject{12045}{<< /D [\MusicPDFPageRef{91} /XYZ 72 282.806 null] >>}
\MusicPDFWriteObject{12046}{<< /D [\MusicPDFPageRef{659} /XYZ 72 505.179 null] >>}
\MusicPDFWriteObject{12047}{<< /D [\MusicPDFPageRef{659} /XYZ 72 478.877 null] >>}
\MusicPDFWriteObject{12048}{<< /D [\MusicPDFPageRef{659} /XYZ 72 413.484 null] >>}
\MusicPDFWriteObject{12049}{<< /D [\MusicPDFPageRef{659} /XYZ 72 389.672 null] >>}
\MusicPDFWriteObject{12050}{<< /D [\MusicPDFPageRef{661} /XYZ 72 671.715 null] >>}
\MusicPDFWriteObject{12051}{<< /D [\MusicPDFPageRef{661} /XYZ 72 642.923 null] >>}
\MusicPDFWriteObject{12052}{<< /D [\MusicPDFPageRef{662} /XYZ 72 426.065 null] >>}
\MusicPDFWriteObject{12053}{<< /D [\MusicPDFPageRef{662} /XYZ 72 399.21 null] >>}
\MusicPDFWriteObject{12054}{<< /D [\MusicPDFPageRef{662} /XYZ 72 373.462 null] >>}
\MusicPDFWriteObject{12055}{<< /D [\MusicPDFPageRef{662} /XYZ 72 268.256 null] >>}
\MusicPDFWriteObject{12056}{<< /D [\MusicPDFPageRef{21} /XYZ 72 165.882 null] >>}
\MusicPDFWriteObject{12057}{<< /D [\MusicPDFPageRef{91} /XYZ 72 230.737 null] >>}
\MusicPDFWriteObject{12058}{<< /D [\MusicPDFPageRef{668} /XYZ 72 468.613 null] >>}
\MusicPDFWriteObject{12059}{<< /D [\MusicPDFPageRef{668} /XYZ 72 428.31 null] >>}
\MusicPDFWriteObject{12060}{<< /D [\MusicPDFPageRef{669} /XYZ 72 350.53 null] >>}
\MusicPDFWriteObject{12061}{<< /D [\MusicPDFPageRef{669} /XYZ 72 310.711 null] >>}
\MusicPDFWriteObject{12062}{<< /D [\MusicPDFPageRef{670} /XYZ 72 158.567 null] >>}
\MusicPDFWriteObject{12063}{<< /D [\MusicPDFPageRef{670} /XYZ 72 144.862 null] >>}
\MusicPDFWriteObject{12064}{<< /D [\MusicPDFPageRef{671} /XYZ 72 272.189 null] >>}
\MusicPDFWriteObject{12065}{<< /D [\MusicPDFPageRef{671} /XYZ 72 245.888 null] >>}
\MusicPDFWriteObject{12066}{<< /D [\MusicPDFPageRef{671} /XYZ 72 220.14 null] >>}
\MusicPDFWriteObject{12067}{<< /D [\MusicPDFPageRef{672} /XYZ 72 507.119 null] >>}
\MusicPDFWriteObject{12068}{<< /D [\MusicPDFPageRef{92} /XYZ 72 698.082 null] >>}
\MusicPDFWriteObject{12069}{<< /D [\MusicPDFPageRef{672} /XYZ 72 493.415 null] >>}
\MusicPDFWriteObject{12070}{<< /D [\MusicPDFPageRef{672} /XYZ 72 480.264 null] >>}
\MusicPDFWriteObject{12071}{<< /D [\MusicPDFPageRef{672} /XYZ 72 165.9 null] >>}
\MusicPDFWriteObject{12072}{<< /D [\MusicPDFPageRef{672} /XYZ 72 139.599 null] >>}
\MusicPDFWriteObject{12073}{<< /D [\MusicPDFPageRef{678} /XYZ 72 581.602 null] >>}
\MusicPDFWriteObject{12074}{<< /D [\MusicPDFPageRef{678} /XYZ 72 555.3 null] >>}
\MusicPDFWriteObject{12075}{<< /D [\MusicPDFPageRef{678} /XYZ 72 492.571 null] >>}
\MusicPDFWriteObject{12076}{<< /D [\MusicPDFPageRef{678} /XYZ 72 303.535 null] >>}
\MusicPDFWriteObject{12077}{<< /D [\MusicPDFPageRef{678} /XYZ 72 233.048 null] >>}
\MusicPDFWriteObject{12078}{<< /D [\MusicPDFPageRef{680} /XYZ 72 405.485 null] >>}
\MusicPDFWriteObject{12079}{<< /D [\MusicPDFPageRef{95} /XYZ 72 249.562 null] >>}
\MusicPDFWriteObject{12080}{<< /D [\MusicPDFPageRef{680} /XYZ 72 298.013 null] >>}
\MusicPDFWriteObject{12081}{<< /D [\MusicPDFPageRef{683} /XYZ 72 698.082 null] >>}
\MusicPDFWriteObject{12082}{<< /D [\MusicPDFPageRef{683} /XYZ 72 621.251 null] >>}
\MusicPDFWriteObject{12083}{<< /D [\MusicPDFPageRef{683} /XYZ 72 421.875 null] >>}
\MusicPDFWriteObject{12084}{<< /D [\MusicPDFPageRef{684} /XYZ 72 669.898 null] >>}
\MusicPDFWriteObject{12085}{<< /D [\MusicPDFPageRef{684} /XYZ 72 656.194 null] >>}
\MusicPDFWriteObject{12086}{<< /D [\MusicPDFPageRef{685} /XYZ 72 442.65 null] >>}
\MusicPDFWriteObject{12087}{<< /D [\MusicPDFPageRef{685} /XYZ 72 428.946 null] >>}
\MusicPDFWriteObject{12088}{<< /D [\MusicPDFPageRef{686} /XYZ 72 564.845 null] >>}
\MusicPDFWriteObject{12089}{<< /D [\MusicPDFPageRef{686} /XYZ 72 551.141 null] >>}
\MusicPDFWriteObject{12090}{<< /D [\MusicPDFPageRef{95} /XYZ 72 177.053 null] >>}
\MusicPDFWriteObject{12091}{<< /D [\MusicPDFPageRef{686} /XYZ 72 511.64 null] >>}
\MusicPDFWriteObject{12092}{<< /D [\MusicPDFPageRef{686} /XYZ 72 485.892 null] >>}
\MusicPDFWriteObject{12093}{<< /D [\MusicPDFPageRef{688} /XYZ 72 221.38 null] >>}
\MusicPDFWriteObject{12094}{<< /D [\MusicPDFPageRef{688} /XYZ 72 207.676 null] >>}
\MusicPDFWriteObject{12095}{<< /D [\MusicPDFPageRef{689} /XYZ 72 638.308 null] >>}
\MusicPDFWriteObject{12096}{<< /D [\MusicPDFPageRef{689} /XYZ 72 624.603 null] >>}
\MusicPDFWriteObject{12097}{<< /D [\MusicPDFPageRef{689} /XYZ 72 611.453 null] >>}
\MusicPDFWriteObject{12098}{<< /D [\MusicPDFPageRef{689} /XYZ 72 479.384 null] >>}
\MusicPDFWriteObject{12099}{<< /D [\MusicPDFPageRef{689} /XYZ 72 453.083 null] >>}
\MusicPDFWriteObject{12100}{<< /D [\MusicPDFPageRef{689} /XYZ 72 168.701 null] >>}
\MusicPDFWriteObject{12101}{<< /D [\MusicPDFPageRef{97} /XYZ 72 239.977 null] >>}
\MusicPDFWriteObject{12102}{<< /D [\MusicPDFPageRef{689} /XYZ 72 154.637 null] >>}
\MusicPDFWriteObject{12103}{<< /D [\MusicPDFPageRef{689} /XYZ 72 141.486 null] >>}
\MusicPDFWriteObject{12104}{<< /D [\MusicPDFPageRef{690} /XYZ 72 698.082 null] >>}
\MusicPDFWriteObject{12105}{<< /D [\MusicPDFPageRef{690} /XYZ 72 529.445 null] >>}
\MusicPDFWriteObject{12106}{<< /D [\MusicPDFPageRef{690} /XYZ 72 502.23 null] >>}
\MusicPDFWriteObject{12107}{<< /D [\MusicPDFPageRef{691} /XYZ 72 194.374 null] >>}
\MusicPDFWriteObject{12108}{<< /D [\MusicPDFPageRef{691} /XYZ 72 168.073 null] >>}
\MusicPDFWriteObject{12109}{<< /D [\MusicPDFPageRef{691} /XYZ 72 141.217 null] >>}
\MusicPDFWriteObject{12110}{<< /D [\MusicPDFPageRef{692} /XYZ 72 698.082 null] >>}
\MusicPDFWriteObject{12111}{<< /D [\MusicPDFPageRef{692} /XYZ 72 673.032 null] >>}
\MusicPDFWriteObject{12112}{<< /D [\MusicPDFPageRef{97} /XYZ 72 226.904 null] >>}
\MusicPDFWriteObject{12113}{<< /D [\MusicPDFPageRef{692} /XYZ 72 646.731 null] >>}
\MusicPDFWriteObject{12114}{<< /D [\MusicPDFPageRef{694} /XYZ 72 569.557 null] >>}
\MusicPDFWriteObject{12115}{<< /D [\MusicPDFPageRef{694} /XYZ 72 554.614 null] >>}
\MusicPDFWriteObject{12116}{<< /D [\MusicPDFPageRef{694} /XYZ 72 542.958 null] >>}
\MusicPDFWriteObject{12117}{<< /D [\MusicPDFPageRef{695} /XYZ 72 475.132 null] >>}
\MusicPDFWriteObject{12118}{<< /D [\MusicPDFPageRef{695} /XYZ 72 461.428 null] >>}
\MusicPDFWriteObject{12119}{<< /D [\MusicPDFPageRef{695} /XYZ 72 448.277 null] >>}
\MusicPDFWriteObject{12120}{<< /D [\MusicPDFPageRef{695} /XYZ 72 435.127 null] >>}
\MusicPDFWriteObject{12121}{<< /D [\MusicPDFPageRef{697} /XYZ 72 606.585 null] >>}
\MusicPDFWriteObject{12122}{<< /D [\MusicPDFPageRef{697} /XYZ 72 593.433 null] >>}
\MusicPDFWriteObject{12123}{<< /D [\MusicPDFPageRef{97} /XYZ 72 114.703 null] >>}
\MusicPDFWriteObject{12124}{<< /D [\MusicPDFPageRef{697} /XYZ 72 260.262 null] >>}
\MusicPDFWriteObject{12125}{<< /D [\MusicPDFPageRef{697} /XYZ 72 234.514 null] >>}
\MusicPDFWriteObject{12126}{<< /D [\MusicPDFPageRef{697} /XYZ 72 218.955 null] >>}
\MusicPDFWriteObject{12127}{<< /D [\MusicPDFPageRef{697} /XYZ 72 151.776 null] >>}
\MusicPDFWriteObject{12128}{<< /D [\MusicPDFPageRef{697} /XYZ 72 127.965 null] >>}
\MusicPDFWriteObject{12129}{<< /D [\MusicPDFPageRef{698} /XYZ 72 323.423 null] >>}
\MusicPDFWriteObject{12130}{<< /D [\MusicPDFPageRef{698} /XYZ 72 297.122 null] >>}
\MusicPDFWriteObject{12131}{<< /D [\MusicPDFPageRef{698} /XYZ 72 260.14 null] >>}
\MusicPDFWriteObject{12132}{<< /D [\MusicPDFPageRef{698} /XYZ 72 233.839 null] >>}
\MusicPDFWriteObject{12133}{<< /D [\MusicPDFPageRef{698} /XYZ 72 208.091 null] >>}
\MusicPDFWriteObject{12134}{<< /D [\MusicPDFPageRef{98} /XYZ 72 674.969 null] >>}
\MusicPDFWriteObject{12135}{<< /D [\MusicPDFPageRef{699} /XYZ 72 435.799 null] >>}
\MusicPDFWriteObject{12136}{<< /D [\MusicPDFPageRef{699} /XYZ 72 363.726 null] >>}
\MusicPDFWriteObject{12137}{<< /D [\MusicPDFPageRef{700} /XYZ 72 698.082 null] >>}
\MusicPDFWriteObject{12138}{<< /D [\MusicPDFPageRef{700} /XYZ 72 672.478 null] >>}
\MusicPDFWriteObject{12139}{<< /D [\MusicPDFPageRef{700} /XYZ 72 659.327 null] >>}
\MusicPDFWriteObject{12140}{<< /D [\MusicPDFPageRef{701} /XYZ 72 455.691 null] >>}
\MusicPDFWriteObject{12141}{<< /D [\MusicPDFPageRef{701} /XYZ 72 441.986 null] >>}
\MusicPDFWriteObject{12142}{<< /D [\MusicPDFPageRef{702} /XYZ 72 353.394 null] >>}
\MusicPDFWriteObject{12143}{<< /D [\MusicPDFPageRef{702} /XYZ 72 337.753 null] >>}
\MusicPDFWriteObject{12144}{<< /D [\MusicPDFPageRef{702} /XYZ 72 311.451 null] >>}
\MusicPDFWriteObject{12145}{<< /D [\MusicPDFPageRef{98} /XYZ 72 597.152 null] >>}
\MusicPDFWriteObject{12146}{<< /D [\MusicPDFPageRef{703} /XYZ 72 200.968 null] >>}
\MusicPDFWriteObject{12147}{<< /D [\MusicPDFPageRef{703} /XYZ 72 187.818 null] >>}
\MusicPDFWriteObject{12148}{<< /D [\MusicPDFPageRef{704} /XYZ 72 381.921 null] >>}
\MusicPDFWriteObject{12149}{<< /D [\MusicPDFPageRef{704} /XYZ 72 266.586 null] >>}
\MusicPDFWriteObject{12150}{<< /D [\MusicPDFPageRef{704} /XYZ 72 255.627 null] >>}
\MusicPDFWriteObject{12151}{<< /D [\MusicPDFPageRef{704} /XYZ 72 169.319 null] >>}
\MusicPDFWriteObject{12152}{<< /D [\MusicPDFPageRef{704} /XYZ 72 155.614 null] >>}
\MusicPDFWriteObject{12153}{<< /D [\MusicPDFPageRef{704} /XYZ 72 141.129 null] >>}
\MusicPDFWriteObject{12154}{<< /D [\MusicPDFPageRef{705} /XYZ 72 516.742 null] >>}
\MusicPDFWriteObject{12155}{<< /D [\MusicPDFPageRef{705} /XYZ 72 466.25 null] >>}
\MusicPDFWriteObject{12156}{<< /D [\MusicPDFPageRef{98} /XYZ 72 556.709 null] >>}
\MusicPDFWriteObject{12157}{<< /D [\MusicPDFPageRef{705} /XYZ 72 354.303 null] >>}
\MusicPDFWriteObject{12158}{<< /D [\MusicPDFPageRef{708} /XYZ 72 474.699 null] >>}
\MusicPDFWriteObject{12159}{<< /D [\MusicPDFPageRef{708} /XYZ 72 415.224 null] >>}
\MusicPDFWriteObject{12160}{<< /D [\MusicPDFPageRef{708} /XYZ 72 403.126 null] >>}
\MusicPDFWriteObject{12161}{<< /D [\MusicPDFPageRef{708} /XYZ 72 345.554 null] >>}
\MusicPDFWriteObject{12162}{<< /D [\MusicPDFPageRef{709} /XYZ 72 291.179 null] >>}
\MusicPDFWriteObject{12163}{<< /D [\MusicPDFPageRef{709} /XYZ 72 264.324 null] >>}
\MusicPDFWriteObject{12164}{<< /D [\MusicPDFPageRef{710} /XYZ 72 670.386 null] >>}
\MusicPDFWriteObject{12165}{<< /D [\MusicPDFPageRef{710} /XYZ 72 644.084 null] >>}
\MusicPDFWriteObject{12166}{<< /D [\MusicPDFPageRef{710} /XYZ 72 349.039 null] >>}
\MusicPDFWriteObject{12167}{<< /D [\MusicPDFPageRef{933} /XYZ 72 698.082 null] >>}
\MusicPDFWriteObject{12168}{<< /D [\MusicPDFPageRef{537} /XYZ 72 698.082 null] >>}
\MusicPDFWriteObject{12169}{<< /D [\MusicPDFPageRef{587} /XYZ 72 698.082 null] >>}
\MusicPDFWriteObject{12170}{<< /D [\MusicPDFPageRef{655} /XYZ 72 698.082 null] >>}
\MusicPDFWriteObject{12171}{<< /D [\MusicPDFPageRef{743} /XYZ 72 698.082 null] >>}
\MusicPDFWriteObject{12172}{<< /D [\MusicPDFPageRef{823} /XYZ 72 698.082 null] >>}
\MusicPDFWriteObject{12173}{<< /D [\MusicPDFPageRef{9} /XYZ 72 698.082 null] >>}
\MusicPDFWriteObject{12174}{<< /D [\MusicPDFPageRef{11} /XYZ 72 698.082 null] >>}
\MusicPDFWriteObject{12175}{<< /D [\MusicPDFPageRef{433} /XYZ 72 698.082 null] >>}
\MusicPDFWriteObject{12176}{<< /D [\MusicPDFPageRef{475} /XYZ 72 698.082 null] >>}
\MusicPDFWriteObject{12177}{<< /D [\MusicPDFPageRef{497} /XYZ 72 698.082 null] >>}
\MusicPDFWriteObject{12178}{<< /D [\MusicPDFPageRef{521} /XYZ 72 698.082 null] >>}
\MusicPDFWriteObject{12179}{<< /D [\MusicPDFPageRef{523} /XYZ 72 698.082 null] >>}
\MusicPDFWriteObject{12180}{<< /D [\MusicPDFPageRef{57} /XYZ 72 698.082 null] >>}
\MusicPDFWriteObject{12181}{<< /D [\MusicPDFPageRef{117} /XYZ 72 698.082 null] >>}
\MusicPDFWriteObject{12182}{<< /D [\MusicPDFPageRef{155} /XYZ 72 698.082 null] >>}
\MusicPDFWriteObject{12183}{<< /D [\MusicPDFPageRef{221} /XYZ 72 698.082 null] >>}
\MusicPDFWriteObject{12184}{<< /D [\MusicPDFPageRef{287} /XYZ 72 698.082 null] >>}
\MusicPDFWriteObject{12185}{<< /D [\MusicPDFPageRef{311} /XYZ 72 698.082 null] >>}
\MusicPDFWriteObject{12186}{<< /D [\MusicPDFPageRef{381} /XYZ 72 698.082 null] >>}
\MusicPDFWriteObject{12187}{<< /D [\MusicPDFPageRef{395} /XYZ 72 698.082 null] >>}
\MusicPDFWriteObject{12188}{<< /D [\MusicPDFPageRef{938} /XYZ 72 567.571 null] >>}
\MusicPDFWriteObject{12189}{<< /D [\MusicPDFPageRef{933} /XYZ 72 587.696 null] >>}
\MusicPDFWriteObject{12190}{<< /D [\MusicPDFPageRef{933} /XYZ 72 576.737 null] >>}
\MusicPDFWriteObject{12191}{<< /D [\MusicPDFPageRef{933} /XYZ 72 565.778 null] >>}
\MusicPDFWriteObject{12192}{<< /D [\MusicPDFPageRef{933} /XYZ 72 445.23 null] >>}
\MusicPDFWriteObject{12193}{<< /D [\MusicPDFPageRef{933} /XYZ 72 434.271 null] >>}
\MusicPDFWriteObject{12194}{<< /D [\MusicPDFPageRef{933} /XYZ 72 554.819 null] >>}
\MusicPDFWriteObject{12195}{<< /D [\MusicPDFPageRef{933} /XYZ 72 543.86 null] >>}
\MusicPDFWriteObject{12196}{<< /D [\MusicPDFPageRef{933} /XYZ 72 532.901 null] >>}
\MusicPDFWriteObject{12197}{<< /D [\MusicPDFPageRef{933} /XYZ 72 521.942 null] >>}
\MusicPDFWriteObject{12198}{<< /D [\MusicPDFPageRef{933} /XYZ 72 510.983 null] >>}
\MusicPDFWriteObject{12199}{<< /D [\MusicPDFPageRef{933} /XYZ 72 489.066 null] >>}
\MusicPDFWriteObject{12200}{<< /D [\MusicPDFPageRef{933} /XYZ 72 478.107 null] >>}
\MusicPDFWriteObject{12201}{<< /D [\MusicPDFPageRef{933} /XYZ 72 467.148 null] >>}
\MusicPDFWriteObject{12202}{<< /D [\MusicPDFPageRef{933} /XYZ 72 456.189 null] >>}
\MusicPDFWriteObject{12203}{<< /D [\MusicPDFPageRef{933} /XYZ 72 423.312 null] >>}
\MusicPDFWriteObject{12204}{<< /D [\MusicPDFPageRef{933} /XYZ 72 401.394 null] >>}
\MusicPDFWriteObject{12205}{<< /D [\MusicPDFPageRef{933} /XYZ 72 390.435 null] >>}
\MusicPDFWriteObject{12206}{<< /D [\MusicPDFPageRef{933} /XYZ 72 313.722 null] >>}
\MusicPDFWriteObject{12207}{<< /D [\MusicPDFPageRef{933} /XYZ 72 291.804 null] >>}
\MusicPDFWriteObject{12208}{<< /D [\MusicPDFPageRef{933} /XYZ 72 280.845 null] >>}
\MusicPDFWriteObject{12209}{<< /D [\MusicPDFPageRef{933} /XYZ 72 247.968 null] >>}
\MusicPDFWriteObject{12210}{<< /D [\MusicPDFPageRef{933} /XYZ 72 226.05 null] >>}
\MusicPDFWriteObject{12211}{<< /D [\MusicPDFPageRef{933} /XYZ 72 204.132 null] >>}
\MusicPDFWriteObject{12212}{<< /D [\MusicPDFPageRef{933} /XYZ 72 193.173 null] >>}
\MusicPDFWriteObject{12213}{<< /D [\MusicPDFPageRef{933} /XYZ 72 182.215 null] >>}
\MusicPDFWriteObject{12214}{<< /D [\MusicPDFPageRef{933} /XYZ 72 171.256 null] >>}
\MusicPDFWriteObject{12215}{<< /D [\MusicPDFPageRef{933} /XYZ 72 127.42 null] >>}
\MusicPDFWriteObject{12216}{<< /D [\MusicPDFPageRef{934} /XYZ 72 370.31 null] >>}
\MusicPDFWriteObject{12217}{<< /D [\MusicPDFPageRef{933} /XYZ 72 368.517 null] >>}
\MusicPDFWriteObject{12218}{<< /D [\MusicPDFPageRef{933} /XYZ 72 357.558 null] >>}
\MusicPDFWriteObject{12219}{<< /D [\MusicPDFPageRef{933} /XYZ 72 335.64 null] >>}
\MusicPDFWriteObject{12220}{<< /D [\MusicPDFPageRef{933} /XYZ 72 324.681 null] >>}
\MusicPDFWriteObject{12221}{<< /D [\MusicPDFPageRef{933} /XYZ 72 258.927 null] >>}
\MusicPDFWriteObject{12222}{<< /D [\MusicPDFPageRef{933} /XYZ 72 215.091 null] >>}
\MusicPDFWriteObject{12223}{<< /D [\MusicPDFPageRef{933} /XYZ 72 149.338 null] >>}
\MusicPDFWriteObject{12224}{<< /D [\MusicPDFPageRef{934} /XYZ 72 699.079 null] >>}
\MusicPDFWriteObject{12225}{<< /D [\MusicPDFPageRef{934} /XYZ 72 666.202 null] >>}
\MusicPDFWriteObject{12226}{<< /D [\MusicPDFPageRef{934} /XYZ 72 655.243 null] >>}
\MusicPDFWriteObject{12227}{<< /D [\MusicPDFPageRef{934} /XYZ 72 644.284 null] >>}
\MusicPDFWriteObject{12228}{<< /D [\MusicPDFPageRef{934} /XYZ 72 633.325 null] >>}
\MusicPDFWriteObject{12229}{<< /D [\MusicPDFPageRef{934} /XYZ 72 622.366 null] >>}
\MusicPDFWriteObject{12230}{<< /D [\MusicPDFPageRef{934} /XYZ 72 611.407 null] >>}
\MusicPDFWriteObject{12231}{<< /D [\MusicPDFPageRef{934} /XYZ 72 600.448 null] >>}
\MusicPDFWriteObject{12232}{<< /D [\MusicPDFPageRef{934} /XYZ 72 578.53 null] >>}
\MusicPDFWriteObject{12233}{<< /D [\MusicPDFPageRef{934} /XYZ 72 556.612 null] >>}
\MusicPDFWriteObject{12234}{<< /D [\MusicPDFPageRef{934} /XYZ 72 523.735 null] >>}
\MusicPDFWriteObject{12235}{<< /D [\MusicPDFPageRef{934} /XYZ 72 545.653 null] >>}
\MusicPDFWriteObject{12236}{<< /D [\MusicPDFPageRef{934} /XYZ 72 534.694 null] >>}
\MusicPDFWriteObject{12237}{<< /D [\MusicPDFPageRef{934} /XYZ 72 512.776 null] >>}
\MusicPDFWriteObject{12238}{<< /D [\MusicPDFPageRef{934} /XYZ 72 501.817 null] >>}
\MusicPDFWriteObject{12239}{<< /D [\MusicPDFPageRef{934} /XYZ 72 479.899 null] >>}
\MusicPDFWriteObject{12240}{<< /D [\MusicPDFPageRef{934} /XYZ 72 468.94 null] >>}
\MusicPDFWriteObject{12241}{<< /D [\MusicPDFPageRef{934} /XYZ 72 457.981 null] >>}
\MusicPDFWriteObject{12242}{<< /D [\MusicPDFPageRef{934} /XYZ 72 436.063 null] >>}
\MusicPDFWriteObject{12243}{<< /D [\MusicPDFPageRef{934} /XYZ 72 414.145 null] >>}
\MusicPDFWriteObject{12244}{<< /D [\MusicPDFPageRef{934} /XYZ 72 392.227 null] >>}
\MusicPDFWriteObject{12245}{<< /D [\MusicPDFPageRef{934} /XYZ 72 293.597 null] >>}
\MusicPDFWriteObject{12246}{<< /D [\MusicPDFPageRef{934} /XYZ 72 282.638 null] >>}
\MusicPDFWriteObject{12247}{<< /D [\MusicPDFPageRef{934} /XYZ 72 238.802 null] >>}
\MusicPDFWriteObject{12248}{<< /D [\MusicPDFPageRef{934} /XYZ 72 129.212 null] >>}
\MusicPDFWriteObject{12249}{<< /D [\MusicPDFPageRef{934} /XYZ 72 348.392 null] >>}
\MusicPDFWriteObject{12250}{<< /D [\MusicPDFPageRef{934} /XYZ 72 326.474 null] >>}
\MusicPDFWriteObject{12251}{<< /D [\MusicPDFPageRef{934} /XYZ 72 315.515 null] >>}
\MusicPDFWriteObject{12252}{<< /D [\MusicPDFPageRef{934} /XYZ 72 271.679 null] >>}
\MusicPDFWriteObject{12253}{<< /D [\MusicPDFPageRef{934} /XYZ 72 249.761 null] >>}
\MusicPDFWriteObject{12254}{<< /D [\MusicPDFPageRef{934} /XYZ 72 216.884 null] >>}
\MusicPDFWriteObject{12255}{<< /D [\MusicPDFPageRef{934} /XYZ 72 205.925 null] >>}
\MusicPDFWriteObject{12256}{<< /D [\MusicPDFPageRef{934} /XYZ 72 194.966 null] >>}
\MusicPDFWriteObject{12257}{<< /D [\MusicPDFPageRef{934} /XYZ 72 184.007 null] >>}
\MusicPDFWriteObject{12258}{<< /D [\MusicPDFPageRef{934} /XYZ 72 173.048 null] >>}
\MusicPDFWriteObject{12259}{<< /D [\MusicPDFPageRef{934} /XYZ 72 151.13 null] >>}
\MusicPDFWriteObject{12260}{<< /D [\MusicPDFPageRef{934} /XYZ 72 140.171 null] >>}
\MusicPDFWriteObject{12261}{<< /D [\MusicPDFPageRef{935} /XYZ 72 655.243 null] >>}
\MusicPDFWriteObject{12262}{<< /D [\MusicPDFPageRef{935} /XYZ 72 501.817 null] >>}
\MusicPDFWriteObject{12263}{<< /D [\MusicPDFPageRef{935} /XYZ 72 490.858 null] >>}
\MusicPDFWriteObject{12264}{<< /D [\MusicPDFPageRef{935} /XYZ 72 699.079 null] >>}
\MusicPDFWriteObject{12265}{<< /D [\MusicPDFPageRef{935} /XYZ 72 677.161 null] >>}
\MusicPDFWriteObject{12266}{<< /D [\MusicPDFPageRef{935} /XYZ 72 622.366 null] >>}
\MusicPDFWriteObject{12267}{<< /D [\MusicPDFPageRef{935} /XYZ 72 600.448 null] >>}
\MusicPDFWriteObject{12268}{<< /D [\MusicPDFPageRef{935} /XYZ 72 578.53 null] >>}
\MusicPDFWriteObject{12269}{<< /D [\MusicPDFPageRef{935} /XYZ 72 567.571 null] >>}
\MusicPDFWriteObject{12270}{<< /D [\MusicPDFPageRef{935} /XYZ 72 644.284 null] >>}
\MusicPDFWriteObject{12271}{<< /D [\MusicPDFPageRef{935} /XYZ 72 545.653 null] >>}
\MusicPDFWriteObject{12272}{<< /D [\MusicPDFPageRef{935} /XYZ 72 523.735 null] >>}
\MusicPDFWriteObject{12273}{<< /D [\MusicPDFPageRef{935} /XYZ 72 512.776 null] >>}
\MusicPDFWriteObject{12274}{<< /D [\MusicPDFPageRef{935} /XYZ 72 479.899 null] >>}
\MusicPDFWriteObject{12275}{<< /D [\MusicPDFPageRef{935} /XYZ 72 468.94 null] >>}
\MusicPDFWriteObject{12276}{<< /D [\MusicPDFPageRef{935} /XYZ 72 457.981 null] >>}
\MusicPDFWriteObject{12277}{<< /D [\MusicPDFPageRef{935} /XYZ 72 447.022 null] >>}
\MusicPDFWriteObject{12278}{<< /D [\MusicPDFPageRef{935} /XYZ 72 425.104 null] >>}
\MusicPDFWriteObject{12279}{<< /D [\MusicPDFPageRef{935} /XYZ 72 414.145 null] >>}
\MusicPDFWriteObject{12280}{<< /D [\MusicPDFPageRef{935} /XYZ 72 392.227 null] >>}
\MusicPDFWriteObject{12281}{<< /D [\MusicPDFPageRef{935} /XYZ 72 370.31 null] >>}
\MusicPDFWriteObject{12282}{<< /D [\MusicPDFPageRef{935} /XYZ 72 348.392 null] >>}
\MusicPDFWriteObject{12283}{<< /D [\MusicPDFPageRef{935} /XYZ 72 326.474 null] >>}
\MusicPDFWriteObject{12284}{<< /D [\MusicPDFPageRef{935} /XYZ 72 304.556 null] >>}
\MusicPDFWriteObject{12285}{<< /D [\MusicPDFPageRef{935} /XYZ 72 282.638 null] >>}
\MusicPDFWriteObject{12286}{<< /D [\MusicPDFPageRef{935} /XYZ 72 260.72 null] >>}
\MusicPDFWriteObject{12287}{<< /D [\MusicPDFPageRef{935} /XYZ 72 216.884 null] >>}
\MusicPDFWriteObject{12288}{<< /D [\MusicPDFPageRef{935} /XYZ 72 194.966 null] >>}
\MusicPDFWriteObject{12289}{<< /D [\MusicPDFPageRef{935} /XYZ 72 173.048 null] >>}
\MusicPDFWriteObject{12290}{<< /D [\MusicPDFPageRef{935} /XYZ 72 238.802 null] >>}
\MusicPDFWriteObject{12291}{<< /D [\MusicPDFPageRef{935} /XYZ 72 227.843 null] >>}
\MusicPDFWriteObject{12292}{<< /D [\MusicPDFPageRef{935} /XYZ 72 162.089 null] >>}
\MusicPDFWriteObject{12293}{<< /D [\MusicPDFPageRef{935} /XYZ 72 140.171 null] >>}
\MusicPDFWriteObject{12294}{<< /D [\MusicPDFPageRef{935} /XYZ 72 129.212 null] >>}
\MusicPDFWriteObject{12295}{<< /D [\MusicPDFPageRef{936} /XYZ 72 677.161 null] >>}
\MusicPDFWriteObject{12296}{<< /D [\MusicPDFPageRef{936} /XYZ 72 611.407 null] >>}
\MusicPDFWriteObject{12297}{<< /D [\MusicPDFPageRef{935} /XYZ 72 118.253 null] >>}
\MusicPDFWriteObject{12298}{<< /D [\MusicPDFPageRef{936} /XYZ 72 699.079 null] >>}
\MusicPDFWriteObject{12299}{<< /D [\MusicPDFPageRef{936} /XYZ 72 655.243 null] >>}
\MusicPDFWriteObject{12300}{<< /D [\MusicPDFPageRef{936} /XYZ 72 644.284 null] >>}
\MusicPDFWriteObject{12301}{<< /D [\MusicPDFPageRef{936} /XYZ 72 633.325 null] >>}
\MusicPDFWriteObject{12302}{<< /D [\MusicPDFPageRef{936} /XYZ 72 622.366 null] >>}
\MusicPDFWriteObject{12303}{<< /D [\MusicPDFPageRef{936} /XYZ 72 589.489 null] >>}
\MusicPDFWriteObject{12304}{<< /D [\MusicPDFPageRef{936} /XYZ 72 523.735 null] >>}
\MusicPDFWriteObject{12305}{<< /D [\MusicPDFPageRef{936} /XYZ 72 501.817 null] >>}
\MusicPDFWriteObject{12306}{<< /D [\MusicPDFPageRef{936} /XYZ 72 468.94 null] >>}
\MusicPDFWriteObject{12307}{<< /D [\MusicPDFPageRef{936} /XYZ 72 425.104 null] >>}
\MusicPDFWriteObject{12308}{<< /D [\MusicPDFPageRef{936} /XYZ 72 414.145 null] >>}
\MusicPDFWriteObject{12309}{<< /D [\MusicPDFPageRef{936} /XYZ 72 304.556 null] >>}
\MusicPDFWriteObject{12310}{<< /D [\MusicPDFPageRef{936} /XYZ 72 578.53 null] >>}
\MusicPDFWriteObject{12311}{<< /D [\MusicPDFPageRef{936} /XYZ 72 567.571 null] >>}
\MusicPDFWriteObject{12312}{<< /D [\MusicPDFPageRef{936} /XYZ 72 556.612 null] >>}
\MusicPDFWriteObject{12313}{<< /D [\MusicPDFPageRef{936} /XYZ 72 545.653 null] >>}
\MusicPDFWriteObject{12314}{<< /D [\MusicPDFPageRef{936} /XYZ 72 534.694 null] >>}
\MusicPDFWriteObject{12315}{<< /D [\MusicPDFPageRef{936} /XYZ 72 479.899 null] >>}
\MusicPDFWriteObject{12316}{<< /D [\MusicPDFPageRef{936} /XYZ 72 457.981 null] >>}
\MusicPDFWriteObject{12317}{<< /D [\MusicPDFPageRef{936} /XYZ 72 436.063 null] >>}
\MusicPDFWriteObject{12318}{<< /D [\MusicPDFPageRef{936} /XYZ 72 403.186 null] >>}
\MusicPDFWriteObject{12319}{<< /D [\MusicPDFPageRef{936} /XYZ 72 392.227 null] >>}
\MusicPDFWriteObject{12320}{<< /D [\MusicPDFPageRef{936} /XYZ 72 370.31 null] >>}
\MusicPDFWriteObject{12321}{<< /D [\MusicPDFPageRef{936} /XYZ 72 359.351 null] >>}
\MusicPDFWriteObject{12322}{<< /D [\MusicPDFPageRef{936} /XYZ 72 337.433 null] >>}
\MusicPDFWriteObject{12323}{<< /D [\MusicPDFPageRef{936} /XYZ 72 326.474 null] >>}
\MusicPDFWriteObject{12324}{<< /D [\MusicPDFPageRef{936} /XYZ 72 315.515 null] >>}
\MusicPDFWriteObject{12325}{<< /D [\MusicPDFPageRef{936} /XYZ 72 260.72 null] >>}
\MusicPDFWriteObject{12326}{<< /D [\MusicPDFPageRef{936} /XYZ 72 238.802 null] >>}
\MusicPDFWriteObject{12327}{<< /D [\MusicPDFPageRef{936} /XYZ 72 216.884 null] >>}
\MusicPDFWriteObject{12328}{<< /D [\MusicPDFPageRef{936} /XYZ 72 205.925 null] >>}
\MusicPDFWriteObject{12329}{<< /D [\MusicPDFPageRef{936} /XYZ 72 184.007 null] >>}
\MusicPDFWriteObject{12330}{<< /D [\MusicPDFPageRef{936} /XYZ 72 140.171 null] >>}
\MusicPDFWriteObject{12331}{<< /D [\MusicPDFPageRef{936} /XYZ 72 129.212 null] >>}
\MusicPDFWriteObject{12332}{<< /D [\MusicPDFPageRef{936} /XYZ 72 118.253 null] >>}
\MusicPDFWriteObject{12333}{<< /D [\MusicPDFPageRef{936} /XYZ 72 107.294 null] >>}
\MusicPDFWriteObject{12334}{<< /D [\MusicPDFPageRef{937} /XYZ 72 699.079 null] >>}
\MusicPDFWriteObject{12335}{<< /D [\MusicPDFPageRef{937} /XYZ 72 677.161 null] >>}
\MusicPDFWriteObject{12336}{<< /D [\MusicPDFPageRef{937} /XYZ 72 578.53 null] >>}
\MusicPDFWriteObject{12337}{<< /D [\MusicPDFPageRef{936} /XYZ 72 282.638 null] >>}
\MusicPDFWriteObject{12338}{<< /D [\MusicPDFPageRef{936} /XYZ 72 249.761 null] >>}
\MusicPDFWriteObject{12339}{<< /D [\MusicPDFPageRef{936} /XYZ 72 151.13 null] >>}
\MusicPDFWriteObject{12340}{<< /D [\MusicPDFPageRef{936} /XYZ 72 173.048 null] >>}
\MusicPDFWriteObject{12341}{<< /D [\MusicPDFPageRef{937} /XYZ 72 666.202 null] >>}
\MusicPDFWriteObject{12342}{<< /D [\MusicPDFPageRef{937} /XYZ 72 644.284 null] >>}
\MusicPDFWriteObject{12343}{<< /D [\MusicPDFPageRef{937} /XYZ 72 600.448 null] >>}
\MusicPDFWriteObject{12344}{<< /D [\MusicPDFPageRef{937} /XYZ 72 622.366 null] >>}
\MusicPDFWriteObject{12345}{<< /D [\MusicPDFPageRef{937} /XYZ 72 556.612 null] >>}
\MusicPDFWriteObject{12346}{<< /D [\MusicPDFPageRef{937} /XYZ 72 545.653 null] >>}
\MusicPDFWriteObject{12347}{<< /D [\MusicPDFPageRef{937} /XYZ 72 534.694 null] >>}
\MusicPDFWriteObject{12348}{<< /D [\MusicPDFPageRef{937} /XYZ 72 523.735 null] >>}
\MusicPDFWriteObject{12349}{<< /D [\MusicPDFPageRef{937} /XYZ 72 512.776 null] >>}
\MusicPDFWriteObject{12350}{<< /D [\MusicPDFPageRef{937} /XYZ 72 501.817 null] >>}
\MusicPDFWriteObject{12351}{<< /D [\MusicPDFPageRef{937} /XYZ 72 457.981 null] >>}
\MusicPDFWriteObject{12352}{<< /D [\MusicPDFPageRef{937} /XYZ 72 403.186 null] >>}
\MusicPDFWriteObject{12353}{<< /D [\MusicPDFPageRef{937} /XYZ 72 479.899 null] >>}
\MusicPDFWriteObject{12354}{<< /D [\MusicPDFPageRef{937} /XYZ 72 468.94 null] >>}
\MusicPDFWriteObject{12355}{<< /D [\MusicPDFPageRef{937} /XYZ 72 425.104 null] >>}
\MusicPDFWriteObject{12356}{<< /D [\MusicPDFPageRef{937} /XYZ 72 414.145 null] >>}
\MusicPDFWriteObject{12357}{<< /D [\MusicPDFPageRef{937} /XYZ 72 271.679 null] >>}
\MusicPDFWriteObject{12358}{<< /D [\MusicPDFPageRef{937} /XYZ 72 129.212 null] >>}
\MusicPDFWriteObject{12359}{<< /D [\MusicPDFPageRef{937} /XYZ 72 392.227 null] >>}
\MusicPDFWriteObject{12360}{<< /D [\MusicPDFPageRef{937} /XYZ 72 381.269 null] >>}
\MusicPDFWriteObject{12361}{<< /D [\MusicPDFPageRef{937} /XYZ 72 370.31 null] >>}
\MusicPDFWriteObject{12362}{<< /D [\MusicPDFPageRef{937} /XYZ 72 359.351 null] >>}
\MusicPDFWriteObject{12363}{<< /D [\MusicPDFPageRef{937} /XYZ 72 293.597 null] >>}
\MusicPDFWriteObject{12364}{<< /D [\MusicPDFPageRef{937} /XYZ 72 315.515 null] >>}
\MusicPDFWriteObject{12365}{<< /D [\MusicPDFPageRef{937} /XYZ 72 304.556 null] >>}
\MusicPDFWriteObject{12366}{<< /D [\MusicPDFPageRef{937} /XYZ 72 326.474 null] >>}
\MusicPDFWriteObject{12367}{<< /D [\MusicPDFPageRef{937} /XYZ 72 282.638 null] >>}
\MusicPDFWriteObject{12368}{<< /D [\MusicPDFPageRef{937} /XYZ 72 238.802 null] >>}
\MusicPDFWriteObject{12369}{<< /D [\MusicPDFPageRef{937} /XYZ 72 249.761 null] >>}
\MusicPDFWriteObject{12370}{<< /D [\MusicPDFPageRef{937} /XYZ 72 216.884 null] >>}
\MusicPDFWriteObject{12371}{<< /D [\MusicPDFPageRef{937} /XYZ 72 205.925 null] >>}
\MusicPDFWriteObject{12372}{<< /D [\MusicPDFPageRef{937} /XYZ 72 184.007 null] >>}
\MusicPDFWriteObject{12373}{<< /D [\MusicPDFPageRef{937} /XYZ 72 162.089 null] >>}
\MusicPDFWriteObject{12374}{<< /D [\MusicPDFPageRef{937} /XYZ 72 151.13 null] >>}
\MusicPDFWriteObject{12375}{<< /D [\MusicPDFPageRef{937} /XYZ 72 140.171 null] >>}
\MusicPDFWriteObject{12376}{<< /D [\MusicPDFPageRef{937} /XYZ 72 107.294 null] >>}
\MusicPDFWriteObject{12377}{<< /D [\MusicPDFPageRef{937} /XYZ 72 118.253 null] >>}
\MusicPDFWriteObject{12378}{<< /D [\MusicPDFPageRef{938} /XYZ 72 699.079 null] >>}
\MusicPDFWriteObject{12379}{<< /D [\MusicPDFPageRef{938} /XYZ 72 677.161 null] >>}
\MusicPDFWriteObject{12380}{<< /D [\MusicPDFPageRef{938} /XYZ 72 523.735 null] >>}
\MusicPDFWriteObject{12381}{<< /D [\MusicPDFPageRef{938} /XYZ 72 468.94 null] >>}
\MusicPDFWriteObject{12382}{<< /D [\MusicPDFPageRef{938} /XYZ 72 457.981 null] >>}
\MusicPDFWriteObject{12383}{<< /D [\MusicPDFPageRef{938} /XYZ 72 436.063 null] >>}
\MusicPDFWriteObject{12384}{<< /D [\MusicPDFPageRef{938} /XYZ 72 666.202 null] >>}
\MusicPDFWriteObject{12385}{<< /D [\MusicPDFPageRef{938} /XYZ 72 633.325 null] >>}
\MusicPDFWriteObject{12386}{<< /D [\MusicPDFPageRef{938} /XYZ 72 622.366 null] >>}
\MusicPDFWriteObject{12387}{<< /D [\MusicPDFPageRef{938} /XYZ 72 611.407 null] >>}
\MusicPDFWriteObject{12388}{<< /D [\MusicPDFPageRef{938} /XYZ 72 600.448 null] >>}
\MusicPDFWriteObject{12389}{<< /D [\MusicPDFPageRef{938} /XYZ 72 589.489 null] >>}
\MusicPDFWriteObject{12390}{<< /D [\MusicPDFPageRef{938} /XYZ 72 578.53 null] >>}
\MusicPDFWriteObject{12391}{<< /D [\MusicPDFPageRef{938} /XYZ 72 644.284 null] >>}
\MusicPDFWriteObject{12392}{<< /D [\MusicPDFPageRef{938} /XYZ 72 545.653 null] >>}
\MusicPDFWriteObject{12393}{<< /D [\MusicPDFPageRef{938} /XYZ 72 534.694 null] >>}
\MusicPDFWriteObject{12394}{<< /D [\MusicPDFPageRef{938} /XYZ 72 512.776 null] >>}
\MusicPDFWriteObject{12395}{<< /D [\MusicPDFPageRef{938} /XYZ 72 490.858 null] >>}
\MusicPDFWriteObject{12396}{<< /D [\MusicPDFPageRef{938} /XYZ 72 414.145 null] >>}
\MusicPDFWriteObject{12397}{<< /D [\MusicPDFPageRef{938} /XYZ 72 194.966 null] >>}
\MusicPDFWriteObject{12398}{<< /D [\MusicPDFPageRef{938} /XYZ 72 140.171 null] >>}
\MusicPDFWriteObject{12399}{<< /D [\MusicPDFPageRef{939} /XYZ 72 699.079 null] >>}
\MusicPDFWriteObject{12400}{<< /D [\MusicPDFPageRef{939} /XYZ 72 688.12 null] >>}
\MusicPDFWriteObject{12401}{<< /D [\MusicPDFPageRef{939} /XYZ 72 677.161 null] >>}
\MusicPDFWriteObject{12402}{<< /D [\MusicPDFPageRef{939} /XYZ 72 666.202 null] >>}
\MusicPDFWriteObject{12403}{<< /D [\MusicPDFPageRef{939} /XYZ 72 644.284 null] >>}
\MusicPDFWriteObject{12404}{<< /D [\MusicPDFPageRef{938} /XYZ 72 403.186 null] >>}
\MusicPDFWriteObject{12405}{<< /D [\MusicPDFPageRef{938} /XYZ 72 392.227 null] >>}
\MusicPDFWriteObject{12406}{<< /D [\MusicPDFPageRef{938} /XYZ 72 381.269 null] >>}
\MusicPDFWriteObject{12407}{<< /D [\MusicPDFPageRef{938} /XYZ 72 359.351 null] >>}
\MusicPDFWriteObject{12408}{<< /D [\MusicPDFPageRef{938} /XYZ 72 348.392 null] >>}
\MusicPDFWriteObject{12409}{<< /D [\MusicPDFPageRef{938} /XYZ 72 370.31 null] >>}
\MusicPDFWriteObject{12410}{<< /D [\MusicPDFPageRef{938} /XYZ 72 326.474 null] >>}
\MusicPDFWriteObject{12411}{<< /D [\MusicPDFPageRef{938} /XYZ 72 304.556 null] >>}
\MusicPDFWriteObject{12412}{<< /D [\MusicPDFPageRef{938} /XYZ 72 282.638 null] >>}
\MusicPDFWriteObject{12413}{<< /D [\MusicPDFPageRef{938} /XYZ 72 271.679 null] >>}
\MusicPDFWriteObject{12414}{<< /D [\MusicPDFPageRef{938} /XYZ 72 249.761 null] >>}
\MusicPDFWriteObject{12415}{<< /D [\MusicPDFPageRef{938} /XYZ 72 238.802 null] >>}
\MusicPDFWriteObject{12416}{<< /D [\MusicPDFPageRef{938} /XYZ 72 227.843 null] >>}
\MusicPDFWriteObject{12417}{<< /D [\MusicPDFPageRef{938} /XYZ 72 216.884 null] >>}
\MusicPDFWriteObject{12418}{<< /D [\MusicPDFPageRef{938} /XYZ 72 205.925 null] >>}
\MusicPDFWriteObject{12419}{<< /D [\MusicPDFPageRef{938} /XYZ 72 184.007 null] >>}
\MusicPDFWriteObject{12420}{<< /D [\MusicPDFPageRef{938} /XYZ 72 162.089 null] >>}
\MusicPDFWriteObject{12421}{<< /D [\MusicPDFPageRef{938} /XYZ 72 129.212 null] >>}
\MusicPDFWriteObject{12422}{<< /D [\MusicPDFPageRef{938} /XYZ 72 118.253 null] >>}
\MusicPDFWriteObject{12423}{<< /D [\MusicPDFPageRef{939} /XYZ 72 622.366 null] >>}
\MusicPDFWriteObject{12424}{<< /D [\MusicPDFPageRef{939} /XYZ 72 611.407 null] >>}
\MusicPDFWriteObject{12425}{<< /D [\MusicPDFPageRef{939} /XYZ 72 589.489 null] >>}
\MusicPDFWriteObject{12426}{<< /D [\MusicPDFPageRef{939} /XYZ 72 600.448 null] >>}
\MusicPDFWriteObject{12427}{<< /D [\MusicPDFPageRef{939} /XYZ 72 578.53 null] >>}
\MusicPDFWriteObject{12428}{<< /D [\MusicPDFPageRef{939} /XYZ 72 556.612 null] >>}
\MusicPDFWriteObject{12429}{<< /D [\MusicPDFPageRef{939} /XYZ 72 545.653 null] >>}
\MusicPDFWriteObject{12430}{<< /D [\MusicPDFPageRef{939} /XYZ 72 523.735 null] >>}
\MusicPDFWriteObject{12431}{<< /D [\MusicPDFPageRef{939} /XYZ 72 534.694 null] >>}
\MusicPDFWriteObject{12432}{<< /D [\MusicPDFPageRef{939} /XYZ 72 512.776 null] >>}
\MusicPDFWriteObject{12433}{<< /D [\MusicPDFPageRef{939} /XYZ 72 479.899 null] >>}
\MusicPDFWriteObject{12434}{<< /D [\MusicPDFPageRef{939} /XYZ 72 501.817 null] >>}
\MusicPDFWriteObject{12435}{<< /D [\MusicPDFPageRef{939} /XYZ 72 414.145 null] >>}
\MusicPDFWriteObject{12436}{<< /D [\MusicPDFPageRef{939} /XYZ 72 315.515 null] >>}
\MusicPDFWriteObject{12437}{<< /D [\MusicPDFPageRef{939} /XYZ 72 468.94 null] >>}
\MusicPDFWriteObject{12438}{<< /D [\MusicPDFPageRef{939} /XYZ 72 457.981 null] >>}
\MusicPDFWriteObject{12439}{<< /D [\MusicPDFPageRef{939} /XYZ 72 447.022 null] >>}
\MusicPDFWriteObject{12440}{<< /D [\MusicPDFPageRef{939} /XYZ 72 425.104 null] >>}
\MusicPDFWriteObject{12441}{<< /D [\MusicPDFPageRef{939} /XYZ 72 370.31 null] >>}
\MusicPDFWriteObject{12442}{<< /D [\MusicPDFPageRef{939} /XYZ 72 359.351 null] >>}
\MusicPDFWriteObject{12443}{<< /D [\MusicPDFPageRef{939} /XYZ 72 348.392 null] >>}
\MusicPDFWriteObject{12444}{<< /D [\MusicPDFPageRef{939} /XYZ 72 381.269 null] >>}
\MusicPDFWriteObject{12445}{<< /D [\MusicPDFPageRef{939} /XYZ 72 337.433 null] >>}
\MusicPDFWriteObject{12446}{<< /D [\MusicPDFPageRef{939} /XYZ 72 392.227 null] >>}
\MusicPDFWriteObject{12447}{<< /D [\MusicPDFPageRef{939} /XYZ 72 216.884 null] >>}
\MusicPDFWriteObject{12448}{<< /D [\MusicPDFPageRef{939} /XYZ 72 194.966 null] >>}
\MusicPDFWriteObject{12449}{<< /D [\MusicPDFPageRef{939} /XYZ 72 173.048 null] >>}
\MusicPDFWriteObject{12450}{<< /D [\MusicPDFPageRef{940} /XYZ 72 699.079 null] >>}
\MusicPDFWriteObject{12451}{<< /D [\MusicPDFPageRef{940} /XYZ 72 677.161 null] >>}
\MusicPDFWriteObject{12452}{<< /D [\MusicPDFPageRef{940} /XYZ 72 655.243 null] >>}
\MusicPDFWriteObject{12453}{<< /D [\MusicPDFPageRef{940} /XYZ 72 556.612 null] >>}
\MusicPDFWriteObject{12454}{<< /D [\MusicPDFPageRef{939} /XYZ 72 304.556 null] >>}
\MusicPDFWriteObject{12455}{<< /D [\MusicPDFPageRef{939} /XYZ 72 282.638 null] >>}
\MusicPDFWriteObject{12456}{<< /D [\MusicPDFPageRef{939} /XYZ 72 271.679 null] >>}
\MusicPDFWriteObject{12457}{<< /D [\MusicPDFPageRef{939} /XYZ 72 260.72 null] >>}
\MusicPDFWriteObject{12458}{<< /D [\MusicPDFPageRef{939} /XYZ 72 249.761 null] >>}
\MusicPDFWriteObject{12459}{<< /D [\MusicPDFPageRef{939} /XYZ 72 227.843 null] >>}
\MusicPDFWriteObject{12460}{<< /D [\MusicPDFPageRef{939} /XYZ 72 238.802 null] >>}
\MusicPDFWriteObject{12461}{<< /D [\MusicPDFPageRef{939} /XYZ 72 151.13 null] >>}
\MusicPDFWriteObject{12462}{<< /D [\MusicPDFPageRef{939} /XYZ 72 140.171 null] >>}
\MusicPDFWriteObject{12463}{<< /D [\MusicPDFPageRef{939} /XYZ 72 118.253 null] >>}
\MusicPDFWriteObject{12464}{<< /D [\MusicPDFPageRef{939} /XYZ 72 107.294 null] >>}
\MusicPDFWriteObject{12465}{<< /D [\MusicPDFPageRef{940} /XYZ 72 611.407 null] >>}
\MusicPDFWriteObject{12466}{<< /D [\MusicPDFPageRef{940} /XYZ 72 633.325 null] >>}
\MusicPDFWriteObject{12467}{<< /D [\MusicPDFPageRef{940} /XYZ 72 600.448 null] >>}
\MusicPDFWriteObject{12468}{<< /D [\MusicPDFPageRef{940} /XYZ 72 578.53 null] >>}
\MusicPDFWriteObject{12469}{<< /D [\MusicPDFPageRef{940} /XYZ 72 567.571 null] >>}
\MusicPDFWriteObject{12470}{<< /D [\MusicPDFPageRef{940} /XYZ 72 644.284 null] >>}
\MusicPDFWriteObject{12471}{<< /D [\MusicPDFPageRef{940} /XYZ 72 545.653 null] >>}
\MusicPDFWriteObject{12472}{<< /D [\MusicPDFPageRef{940} /XYZ 72 534.694 null] >>}
\MusicPDFWriteObject{12473}{<< /D [\MusicPDFPageRef{937} /XYZ 72 173.048 null] >>}
\MusicPDFWriteObject{12474}{<< /D [\MusicPDFPageRef{940} /XYZ 72 523.735 null] >>}
\MusicPDFWriteObject{12475}{<< /D [\MusicPDFPageRef{940} /XYZ 72 512.776 null] >>}
\MusicPDFWriteObject{12476}{<< /D [\MusicPDFPageRef{940} /XYZ 72 501.817 null] >>}
\MusicPDFWriteObject{12477}{<< /D [\MusicPDFPageRef{940} /XYZ 72 457.981 null] >>}
\MusicPDFWriteObject{12478}{<< /D [\MusicPDFPageRef{940} /XYZ 72 468.94 null] >>}
\MusicPDFWriteObject{12479}{<< /D [\MusicPDFPageRef{940} /XYZ 72 490.858 null] >>}
\MusicPDFWriteObject{12480}{<< /D [\MusicPDFPageRef{940} /XYZ 72 479.899 null] >>}
\MusicPDFWriteObject{12481}{<< /D [\MusicPDFPageRef{940} /XYZ 72 392.227 null] >>}
\MusicPDFWriteObject{12482}{<< /D [\MusicPDFPageRef{940} /XYZ 72 447.022 null] >>}
\MusicPDFWriteObject{12483}{<< /D [\MusicPDFPageRef{940} /XYZ 72 425.104 null] >>}
\MusicPDFWriteObject{12484}{<< /D [\MusicPDFPageRef{940} /XYZ 72 403.186 null] >>}
\MusicPDFWriteObject{12485}{<< /D [\MusicPDFPageRef{940} /XYZ 72 370.31 null] >>}
\MusicPDFWriteObject{12486}{<< /D [\MusicPDFPageRef{940} /XYZ 72 249.761 null] >>}
\MusicPDFWriteObject{12487}{<< /D [\MusicPDFPageRef{940} /XYZ 72 359.351 null] >>}
\MusicPDFWriteObject{12488}{<< /D [\MusicPDFPageRef{940} /XYZ 72 337.433 null] >>}
\MusicPDFWriteObject{12489}{<< /D [\MusicPDFPageRef{940} /XYZ 72 326.474 null] >>}
\MusicPDFWriteObject{12490}{<< /D [\MusicPDFPageRef{940} /XYZ 72 304.556 null] >>}
\MusicPDFWriteObject{12491}{<< /D [\MusicPDFPageRef{940} /XYZ 72 293.597 null] >>}
\MusicPDFWriteObject{12492}{<< /D [\MusicPDFPageRef{940} /XYZ 72 282.638 null] >>}
\MusicPDFWriteObject{12493}{<< /D [\MusicPDFPageRef{940} /XYZ 72 260.72 null] >>}
\MusicPDFWriteObject{12494}{<< /D [\MusicPDFPageRef{940} /XYZ 72 238.802 null] >>}
\MusicPDFWriteObject{12495}{<< /D [\MusicPDFPageRef{940} /XYZ 72 227.843 null] >>}
\MusicPDFWriteObject{12496}{<< /D [\MusicPDFPageRef{940} /XYZ 72 216.884 null] >>}
\MusicPDFWriteObject{12497}{<< /D [\MusicPDFPageRef{940} /XYZ 72 205.925 null] >>}
\MusicPDFWriteObject{12498}{<< /D [\MusicPDFPageRef{940} /XYZ 72 194.966 null] >>}
\MusicPDFWriteObject{12499}{<< /D [\MusicPDFPageRef{940} /XYZ 72 184.007 null] >>}
\MusicPDFWriteObject{12500}{<< /D [\MusicPDFPageRef{940} /XYZ 72 173.048 null] >>}
\MusicPDFWriteObject{12501}{<< /D [\MusicPDFPageRef{940} /XYZ 72 140.171 null] >>}
\MusicPDFWriteObject{12502}{<< /D [\MusicPDFPageRef{940} /XYZ 72 162.089 null] >>}
\MusicPDFWriteObject{12503}{<< /D [\MusicPDFPageRef{939} /XYZ 72 567.571 null] >>}
\MusicPDFWriteObject{12504}{<< /D [\MusicPDFPageRef{940} /XYZ 72 381.269 null] >>}
\MusicPDFWriteObject{12505}{<< /D [\MusicPDFPageRef{11} /XYZ 251.39 393.995 null] >>}
\MusicPDFWriteObject{12506}{<< /D [\MusicPDFPageRef{14} /XYZ 125.425 244.64 null] >>}
\MusicPDFWriteObject{12507}{<< /D [\MusicPDFPageRef{22} /XYZ 263.277 679.161 null] >>}
\MusicPDFWriteObject{12508}{<< /D [\MusicPDFPageRef{23} /XYZ 210.825 262.803 null] >>}
\MusicPDFWriteObject{12509}{<< /D [\MusicPDFPageRef{23} /XYZ 164.563 209.328 null] >>}
\MusicPDFWriteObject{12510}{<< /D [\MusicPDFPageRef{24} /XYZ 180.935 264.522 null] >>}
\MusicPDFWriteObject{12511}{<< /D [\MusicPDFPageRef{22} /XYZ 242.186 587.596 null] >>}
\MusicPDFWriteObject{12512}{<< /D [\MusicPDFPageRef{23} /XYZ 243.477 604.065 null] >>}
\MusicPDFWriteObject{12513}{<< /D [\MusicPDFPageRef{23} /XYZ 247.226 522.615 null] >>}
\MusicPDFWriteObject{12514}{<< /D [\MusicPDFPageRef{23} /XYZ 152.028 429.688 null] >>}
\MusicPDFWriteObject{12515}{<< /D [\MusicPDFPageRef{23} /XYZ 72 389.421 null] >>}
\MusicPDFWriteObject{12516}{<< /D [\MusicPDFPageRef{23} /XYZ 72 360.756 null] >>}
\MusicPDFWriteObject{12517}{<< /D [\MusicPDFPageRef{23} /XYZ 72 333.834 null] >>}
\MusicPDFWriteObject{12518}{<< /D [\MusicPDFPageRef{23} /XYZ 72 306.913 null] >>}
\MusicPDFWriteObject{12519}{<< /D [\MusicPDFPageRef{25} /XYZ 268.126 501.788 null] >>}
\MusicPDFWriteObject{12520}{<< /D [\MusicPDFPageRef{25} /XYZ 279.214 151.942 null] >>}
\MusicPDFWriteObject{12521}{<< /D [\MusicPDFPageRef{33} /XYZ 278.188 130.597 null] >>}
\MusicPDFWriteObject{12522}{<< /D [\MusicPDFPageRef{34} /XYZ 271.112 525.247 null] >>}
\MusicPDFWriteObject{12523}{<< /D [\MusicPDFPageRef{35} /XYZ 283.279 341.01 null] >>}
\MusicPDFWriteObject{12524}{<< /D [\MusicPDFPageRef{36} /XYZ 159.224 610.307 null] >>}
\MusicPDFWriteObject{12525}{<< /D [\MusicPDFPageRef{37} /XYZ 297.452 448.886 null] >>}
\MusicPDFWriteObject{12526}{<< /D [\MusicPDFPageRef{37} /XYZ 266.459 352.511 null] >>}
\MusicPDFWriteObject{12527}{<< /D [\MusicPDFPageRef{39} /XYZ 258.552 352.489 null] >>}
\MusicPDFWriteObject{12528}{<< /D [\MusicPDFPageRef{39} /XYZ 276.289 250.943 null] >>}
\MusicPDFWriteObject{12529}{<< /D [\MusicPDFPageRef{41} /XYZ 215.043 447.829 null] >>}
\MusicPDFWriteObject{12530}{<< /D [\MusicPDFPageRef{41} /XYZ 192.623 411.401 null] >>}
\MusicPDFWriteObject{12531}{<< /D [\MusicPDFPageRef{42} /XYZ 251.39 288.337 null] >>}
\MusicPDFWriteObject{12532}{<< /D [\MusicPDFPageRef{44} /XYZ 161.183 421.954 null] >>}
\MusicPDFWriteObject{12533}{<< /D [\MusicPDFPageRef{44} /XYZ 209.673 196.841 null] >>}
\MusicPDFWriteObject{12534}{<< /D [\MusicPDFPageRef{45} /XYZ 157.661 579.693 null] >>}
\MusicPDFWriteObject{12535}{<< /D [\MusicPDFPageRef{45} /XYZ 139.369 397.053 null] >>}
\MusicPDFWriteObject{12536}{<< /D [\MusicPDFPageRef{46} /XYZ 248.568 352.36 null] >>}
\MusicPDFWriteObject{12537}{<< /D [\MusicPDFPageRef{47} /XYZ 233.726 681.264 null] >>}
\MusicPDFWriteObject{12538}{<< /D [\MusicPDFPageRef{47} /XYZ 125.682 119.906 null] >>}
\MusicPDFWriteObject{12539}{<< /D [\MusicPDFPageRef{48} /XYZ 135.01 603.834 null] >>}
\MusicPDFWriteObject{12540}{<< /D [\MusicPDFPageRef{49} /XYZ 268.126 683.057 null] >>}
\MusicPDFWriteObject{12541}{<< /D [\MusicPDFPageRef{49} /XYZ 279.214 333.21 null] >>}
\MusicPDFWriteObject{12542}{<< /D [\MusicPDFPageRef{51} /XYZ 181.641 591.494 null] >>}
\MusicPDFWriteObject{12543}{<< /D [\MusicPDFPageRef{55} /XYZ 117.071 663.657 null] >>}
\MusicPDFWriteObject{12544}{<< /D [\MusicPDFPageRef{434} /XYZ 225.747 578.12 null] >>}
\MusicPDFWriteObject{12545}{<< /D [\MusicPDFPageRef{436} /XYZ 172.274 164.754 null] >>}
\MusicPDFWriteObject{12546}{<< /D [\MusicPDFPageRef{436} /XYZ 178.565 128.327 null] >>}
\MusicPDFWriteObject{12547}{<< /D [\MusicPDFPageRef{437} /XYZ 172.082 683.057 null] >>}
\MusicPDFWriteObject{12548}{<< /D [\MusicPDFPageRef{437} /XYZ 155.191 646.629 null] >>}
\MusicPDFWriteObject{12549}{<< /D [\MusicPDFPageRef{437} /XYZ 174.947 609.405 null] >>}
\MusicPDFWriteObject{12550}{<< /D [\MusicPDFPageRef{437} /XYZ 160.222 559.609 null] >>}
\MusicPDFWriteObject{12551}{<< /D [\MusicPDFPageRef{434} /XYZ 226.915 397.034 null] >>}
\MusicPDFWriteObject{12552}{<< /D [\MusicPDFPageRef{434} /XYZ 227.365 298.814 null] >>}
\MusicPDFWriteObject{12553}{<< /D [\MusicPDFPageRef{434} /XYZ 248.153 236.085 null] >>}
\MusicPDFWriteObject{12554}{<< /D [\MusicPDFPageRef{434} /XYZ 149.35 172.758 null] >>}
\MusicPDFWriteObject{12555}{<< /D [\MusicPDFPageRef{435} /XYZ 135.596 499.428 null] >>}
\MusicPDFWriteObject{12556}{<< /D [\MusicPDFPageRef{436} /XYZ 118.863 670.725 null] >>}
\MusicPDFWriteObject{12557}{<< /D [\MusicPDFPageRef{436} /XYZ 239.446 254.288 null] >>}
\MusicPDFWriteObject{12558}{<< /D [\MusicPDFPageRef{436} /XYZ 226.443 208.291 null] >>}
\MusicPDFWriteObject{12559}{<< /D [\MusicPDFPageRef{437} /XYZ 174.587 363.797 null] >>}
\MusicPDFWriteObject{12560}{<< /D [\MusicPDFPageRef{437} /XYZ 189.689 278.078 null] >>}
\MusicPDFWriteObject{12561}{<< /D [\MusicPDFPageRef{437} /XYZ 259.166 193.131 null] >>}
\MusicPDFWriteObject{12562}{<< /D [\MusicPDFPageRef{438} /XYZ 253.277 630.453 null] >>}
\MusicPDFWriteObject{12563}{<< /D [\MusicPDFPageRef{438} /XYZ 184.59 588.442 null] >>}
\MusicPDFWriteObject{12564}{<< /D [\MusicPDFPageRef{438} /XYZ 181.859 392.645 null] >>}
\MusicPDFWriteObject{12565}{<< /D [\MusicPDFPageRef{439} /XYZ 207.524 188.565 null] >>}
\MusicPDFWriteObject{12566}{<< /D [\MusicPDFPageRef{442} /XYZ 213.72 134.083 null] >>}
\MusicPDFWriteObject{12567}{<< /D [\MusicPDFPageRef{443} /XYZ 239.823 288.367 null] >>}
\MusicPDFWriteObject{12568}{<< /D [\MusicPDFPageRef{443} /XYZ 169.966 186.847 null] >>}
\MusicPDFWriteObject{12569}{<< /D [\MusicPDFPageRef{444} /XYZ 239.823 411.738 null] >>}
\MusicPDFWriteObject{12570}{<< /D [\MusicPDFPageRef{439} /XYZ 198.578 148.811 null] >>}
\MusicPDFWriteObject{12571}{<< /D [\MusicPDFPageRef{439} /XYZ 225.527 99.233 null] >>}
\MusicPDFWriteObject{12572}{<< /D [\MusicPDFPageRef{440} /XYZ 276.988 518.374 null] >>}
\MusicPDFWriteObject{12573}{<< /D [\MusicPDFPageRef{440} /XYZ 239.823 389.306 null] >>}
\MusicPDFWriteObject{12574}{<< /D [\MusicPDFPageRef{440} /XYZ 221.965 305.622 null] >>}
\MusicPDFWriteObject{12575}{<< /D [\MusicPDFPageRef{441} /XYZ 168.058 560.475 null] >>}
\MusicPDFWriteObject{12576}{<< /D [\MusicPDFPageRef{441} /XYZ 231.54 283.487 null] >>}
\MusicPDFWriteObject{12577}{<< /D [\MusicPDFPageRef{442} /XYZ 154.351 544.633 null] >>}
\MusicPDFWriteObject{12578}{<< /D [\MusicPDFPageRef{445} /XYZ 207.833 473.138 null] >>}
\MusicPDFWriteObject{12579}{<< /D [\MusicPDFPageRef{445} /XYZ 198.888 433.937 null] >>}
\MusicPDFWriteObject{12580}{<< /D [\MusicPDFPageRef{445} /XYZ 231.814 384.359 null] >>}
\MusicPDFWriteObject{12581}{<< /D [\MusicPDFPageRef{446} /XYZ 232.227 669.906 null] >>}
\MusicPDFWriteObject{12582}{<< /D [\MusicPDFPageRef{446} /XYZ 192.191 587.55 null] >>}
\MusicPDFWriteObject{12583}{<< /D [\MusicPDFPageRef{446} /XYZ 229.047 546.951 null] >>}
\MusicPDFWriteObject{12584}{<< /D [\MusicPDFPageRef{446} /XYZ 232.227 214.246 null] >>}
\MusicPDFWriteObject{12585}{<< /D [\MusicPDFPageRef{446} /XYZ 162.998 121.031 null] >>}
\MusicPDFWriteObject{12586}{<< /D [\MusicPDFPageRef{447} /XYZ 232.227 508.232 null] >>}
\MusicPDFWriteObject{12587}{<< /D [\MusicPDFPageRef{448} /XYZ 155.426 317.672 null] >>}
\MusicPDFWriteObject{12588}{<< /D [\MusicPDFPageRef{449} /XYZ 72 341.015 null] >>}
\MusicPDFWriteObject{12589}{<< /D [\MusicPDFPageRef{449} /XYZ 72 307.343 null] >>}
\MusicPDFWriteObject{12590}{<< /D [\MusicPDFPageRef{449} /XYZ 72 273.671 null] >>}
\MusicPDFWriteObject{12591}{<< /D [\MusicPDFPageRef{449} /XYZ 96.543 211.437 null] >>}
\MusicPDFWriteObject{12592}{<< /D [\MusicPDFPageRef{449} /XYZ 224.93 168.468 null] >>}
\MusicPDFWriteObject{12593}{<< /D [\MusicPDFPageRef{450} /XYZ 171.055 628.511 null] >>}
\MusicPDFWriteObject{12594}{<< /D [\MusicPDFPageRef{450} /XYZ 178.924 470.936 null] >>}
\MusicPDFWriteObject{12595}{<< /D [\MusicPDFPageRef{450} /XYZ 72 274.615 null] >>}
\MusicPDFWriteObject{12596}{<< /D [\MusicPDFPageRef{451} /XYZ 168.537 573.598 null] >>}
\MusicPDFWriteObject{12597}{<< /D [\MusicPDFPageRef{451} /XYZ 135.565 399.931 null] >>}
\MusicPDFWriteObject{12598}{<< /D [\MusicPDFPageRef{448} /XYZ 255.313 255.701 null] >>}
\MusicPDFWriteObject{12599}{<< /D [\MusicPDFPageRef{451} /XYZ 167.153 337.158 null] >>}
\MusicPDFWriteObject{12600}{<< /D [\MusicPDFPageRef{451} /XYZ 178.686 261.793 null] >>}
\MusicPDFWriteObject{12601}{<< /D [\MusicPDFPageRef{451} /XYZ 172.267 185.877 null] >>}
\MusicPDFWriteObject{12602}{<< /D [\MusicPDFPageRef{452} /XYZ 177.224 607.935 null] >>}
\MusicPDFWriteObject{12603}{<< /D [\MusicPDFPageRef{452} /XYZ 153.851 181.408 null] >>}
\MusicPDFWriteObject{12604}{<< /D [\MusicPDFPageRef{452} /XYZ 176.078 129.883 null] >>}
\MusicPDFWriteObject{12605}{<< /D [\MusicPDFPageRef{453} /XYZ 232.143 601.9 null] >>}
\MusicPDFWriteObject{12606}{<< /D [\MusicPDFPageRef{448} /XYZ 169.837 175.924 null] >>}
\MusicPDFWriteObject{12607}{<< /D [\MusicPDFPageRef{448} /XYZ 221.965 122.229 null] >>}
\MusicPDFWriteObject{12608}{<< /D [\MusicPDFPageRef{449} /XYZ 181.539 654.203 null] >>}
\MusicPDFWriteObject{12609}{<< /D [\MusicPDFPageRef{449} /XYZ 210.125 551.948 null] >>}
\MusicPDFWriteObject{12610}{<< /D [\MusicPDFPageRef{449} /XYZ 155.426 473.809 null] >>}
\MusicPDFWriteObject{12611}{<< /D [\MusicPDFPageRef{449} /XYZ 72 417.633 null] >>}
\MusicPDFWriteObject{12612}{<< /D [\MusicPDFPageRef{449} /XYZ 72 369.824 null] >>}
\MusicPDFWriteObject{12613}{<< /D [\MusicPDFPageRef{453} /XYZ 148.408 130.015 null] >>}
\MusicPDFWriteObject{12614}{<< /D [\MusicPDFPageRef{455} /XYZ 72 581.207 null] >>}
\MusicPDFWriteObject{12615}{<< /D [\MusicPDFPageRef{455} /XYZ 213.21 489.964 null] >>}
\MusicPDFWriteObject{12616}{<< /D [\MusicPDFPageRef{454} /XYZ 252.094 683.057 null] >>}
\MusicPDFWriteObject{12617}{<< /D [\MusicPDFPageRef{454} /XYZ 156.999 603.28 null] >>}
\MusicPDFWriteObject{12618}{<< /D [\MusicPDFPageRef{454} /XYZ 128.751 548.671 null] >>}
\MusicPDFWriteObject{12619}{<< /D [\MusicPDFPageRef{454} /XYZ 149.393 465.159 null] >>}
\MusicPDFWriteObject{12620}{<< /D [\MusicPDFPageRef{454} /XYZ 209.842 267.702 null] >>}
\MusicPDFWriteObject{12621}{<< /D [\MusicPDFPageRef{454} /XYZ 183.375 155.591 null] >>}
\MusicPDFWriteObject{12622}{<< /D [\MusicPDFPageRef{455} /XYZ 72 677.609 null] >>}
\MusicPDFWriteObject{12623}{<< /D [\MusicPDFPageRef{455} /XYZ 72 630.286 null] >>}
\MusicPDFWriteObject{12624}{<< /D [\MusicPDFPageRef{456} /XYZ 246.203 494.966 null] >>}
\MusicPDFWriteObject{12625}{<< /D [\MusicPDFPageRef{457} /XYZ 225.845 99.233 null] >>}
\MusicPDFWriteObject{12626}{<< /D [\MusicPDFPageRef{458} /XYZ 230.607 683.057 null] >>}
\MusicPDFWriteObject{12627}{<< /D [\MusicPDFPageRef{458} /XYZ 227.371 486.439 null] >>}
\MusicPDFWriteObject{12628}{<< /D [\MusicPDFPageRef{458} /XYZ 228.958 423.71 null] >>}
\MusicPDFWriteObject{12629}{<< /D [\MusicPDFPageRef{456} /XYZ 247.79 422.11 null] >>}
\MusicPDFWriteObject{12630}{<< /D [\MusicPDFPageRef{456} /XYZ 172.698 352.306 null] >>}
\MusicPDFWriteObject{12631}{<< /D [\MusicPDFPageRef{456} /XYZ 234.857 289.577 null] >>}
\MusicPDFWriteObject{12632}{<< /D [\MusicPDFPageRef{456} /XYZ 244.94 211.514 null] >>}
\MusicPDFWriteObject{12633}{<< /D [\MusicPDFPageRef{457} /XYZ 246.203 514.321 null] >>}
\MusicPDFWriteObject{12634}{<< /D [\MusicPDFPageRef{457} /XYZ 152.822 451.592 null] >>}
\MusicPDFWriteObject{12635}{<< /D [\MusicPDFPageRef{457} /XYZ 247.79 415.164 null] >>}
\MusicPDFWriteObject{12636}{<< /D [\MusicPDFPageRef{457} /XYZ 193.246 365.586 null] >>}
\MusicPDFWriteObject{12637}{<< /D [\MusicPDFPageRef{458} /XYZ 203.598 193.327 null] >>}
\MusicPDFWriteObject{12638}{<< /D [\MusicPDFPageRef{460} /XYZ 211.873 416.579 null] >>}
\MusicPDFWriteObject{12639}{<< /D [\MusicPDFPageRef{460} /XYZ 221.959 242.527 null] >>}
\MusicPDFWriteObject{12640}{<< /D [\MusicPDFPageRef{460} /XYZ 168.63 157.014 null] >>}
\MusicPDFWriteObject{12641}{<< /D [\MusicPDFPageRef{461} /XYZ 257.789 683.057 null] >>}
\MusicPDFWriteObject{12642}{<< /D [\MusicPDFPageRef{461} /XYZ 225.106 646.629 null] >>}
\MusicPDFWriteObject{12643}{<< /D [\MusicPDFPageRef{461} /XYZ 217.341 597.051 null] >>}
\MusicPDFWriteObject{12644}{<< /D [\MusicPDFPageRef{461} /XYZ 214.887 324.877 null] >>}
\MusicPDFWriteObject{12645}{<< /D [\MusicPDFPageRef{461} /XYZ 273.651 219.628 null] >>}
\MusicPDFWriteObject{12646}{<< /D [\MusicPDFPageRef{461} /XYZ 213.087 183.201 null] >>}
\MusicPDFWriteObject{12647}{<< /D [\MusicPDFPageRef{462} /XYZ 256.599 489.773 null] >>}
\MusicPDFWriteObject{12648}{<< /D [\MusicPDFPageRef{458} /XYZ 166.211 156.899 null] >>}
\MusicPDFWriteObject{12649}{<< /D [\MusicPDFPageRef{463} /XYZ 195.238 378.322 null] >>}
\MusicPDFWriteObject{12650}{<< /D [\MusicPDFPageRef{463} /XYZ 163.902 292.594 null] >>}
\MusicPDFWriteObject{12651}{<< /D [\MusicPDFPageRef{463} /XYZ 122.615 229.865 null] >>}
\MusicPDFWriteObject{12652}{<< /D [\MusicPDFPageRef{463} /XYZ 163.902 193.437 null] >>}
\MusicPDFWriteObject{12653}{<< /D [\MusicPDFPageRef{463} /XYZ 229.42 130.708 null] >>}
\MusicPDFWriteObject{12654}{<< /D [\MusicPDFPageRef{464} /XYZ 207.134 683.057 null] >>}
\MusicPDFWriteObject{12655}{<< /D [\MusicPDFPageRef{464} /XYZ 149.871 627.219 null] >>}
\MusicPDFWriteObject{12656}{<< /D [\MusicPDFPageRef{464} /XYZ 72 557.908 null] >>}
\MusicPDFWriteObject{12657}{<< /D [\MusicPDFPageRef{464} /XYZ 72 490.949 null] >>}
\MusicPDFWriteObject{12658}{<< /D [\MusicPDFPageRef{464} /XYZ 72 453.024 null] >>}
\MusicPDFWriteObject{12659}{<< /D [\MusicPDFPageRef{459} /XYZ 166.211 617.303 null] >>}
\MusicPDFWriteObject{12660}{<< /D [\MusicPDFPageRef{464} /XYZ 72 435.192 null] >>}
\MusicPDFWriteObject{12661}{<< /D [\MusicPDFPageRef{464} /XYZ 72 412.018 null] >>}
\MusicPDFWriteObject{12662}{<< /D [\MusicPDFPageRef{464} /XYZ 72 379.078 null] >>}
\MusicPDFWriteObject{12663}{<< /D [\MusicPDFPageRef{464} /XYZ 72 346.139 null] >>}
\MusicPDFWriteObject{12664}{<< /D [\MusicPDFPageRef{464} /XYZ 129.731 257.806 null] >>}
\MusicPDFWriteObject{12665}{<< /D [\MusicPDFPageRef{459} /XYZ 266.172 575.302 null] >>}
\MusicPDFWriteObject{12666}{<< /D [\MusicPDFPageRef{459} /XYZ 234.251 499.423 null] >>}
\MusicPDFWriteObject{12667}{<< /D [\MusicPDFPageRef{459} /XYZ 259.427 368.954 null] >>}
\MusicPDFWriteObject{12668}{<< /D [\MusicPDFPageRef{459} /XYZ 214.153 332.527 null] >>}
\MusicPDFWriteObject{12669}{<< /D [\MusicPDFPageRef{459} /XYZ 156.918 282.948 null] >>}
\MusicPDFWriteObject{12670}{<< /D [\MusicPDFPageRef{460} /XYZ 207.091 528.394 null] >>}
\MusicPDFWriteObject{12671}{<< /D [\MusicPDFPageRef{465} /XYZ 236.41 678.385 null] >>}
\MusicPDFWriteObject{12672}{<< /D [\MusicPDFPageRef{467} /XYZ 250.595 554.376 null] >>}
\MusicPDFWriteObject{12673}{<< /D [\MusicPDFPageRef{467} /XYZ 165.181 478.364 null] >>}
\MusicPDFWriteObject{12674}{<< /D [\MusicPDFPageRef{467} /XYZ 248.218 356.823 null] >>}
\MusicPDFWriteObject{12675}{<< /D [\MusicPDFPageRef{467} /XYZ 72 311.512 null] >>}
\MusicPDFWriteObject{12676}{<< /D [\MusicPDFPageRef{467} /XYZ 75.02 252.562 null] >>}
\MusicPDFWriteObject{12677}{<< /D [\MusicPDFPageRef{467} /XYZ 204.431 176.195 null] >>}
\MusicPDFWriteObject{12678}{<< /D [\MusicPDFPageRef{467} /XYZ 250.37 99.233 null] >>}
\MusicPDFWriteObject{12679}{<< /D [\MusicPDFPageRef{468} /XYZ 168.089 579.309 null] >>}
\MusicPDFWriteObject{12680}{<< /D [\MusicPDFPageRef{468} /XYZ 249.602 493.143 null] >>}
\MusicPDFWriteObject{12681}{<< /D [\MusicPDFPageRef{468} /XYZ 171.966 112.323 null] >>}
\MusicPDFWriteObject{12682}{<< /D [\MusicPDFPageRef{465} /XYZ 124.852 609.84 null] >>}
\MusicPDFWriteObject{12683}{<< /D [\MusicPDFPageRef{469} /XYZ 186.017 461.457 null] >>}
\MusicPDFWriteObject{12684}{<< /D [\MusicPDFPageRef{469} /XYZ 128.748 264.085 null] >>}
\MusicPDFWriteObject{12685}{<< /D [\MusicPDFPageRef{469} /XYZ 189.72 149.551 null] >>}
\MusicPDFWriteObject{12686}{<< /D [\MusicPDFPageRef{470} /XYZ 191.607 596.887 null] >>}
\MusicPDFWriteObject{12687}{<< /D [\MusicPDFPageRef{470} /XYZ 139.785 547.308 null] >>}
\MusicPDFWriteObject{12688}{<< /D [\MusicPDFPageRef{470} /XYZ 168.312 457.406 null] >>}
\MusicPDFWriteObject{12689}{<< /D [\MusicPDFPageRef{470} /XYZ 218.8 351.515 null] >>}
\MusicPDFWriteObject{12690}{<< /D [\MusicPDFPageRef{470} /XYZ 228.742 301.936 null] >>}
\MusicPDFWriteObject{12691}{<< /D [\MusicPDFPageRef{470} /XYZ 218.489 238.66 null] >>}
\MusicPDFWriteObject{12692}{<< /D [\MusicPDFPageRef{470} /XYZ 166.861 141.055 null] >>}
\MusicPDFWriteObject{12693}{<< /D [\MusicPDFPageRef{465} /XYZ 207.515 108.804 null] >>}
\MusicPDFWriteObject{12694}{<< /D [\MusicPDFPageRef{470} /XYZ 159.667 99.233 null] >>}
\MusicPDFWriteObject{12695}{<< /D [\MusicPDFPageRef{471} /XYZ 231.414 643.604 null] >>}
\MusicPDFWriteObject{12696}{<< /D [\MusicPDFPageRef{471} /XYZ 72 588.478 null] >>}
\MusicPDFWriteObject{12697}{<< /D [\MusicPDFPageRef{471} /XYZ 72 560.555 null] >>}
\MusicPDFWriteObject{12698}{<< /D [\MusicPDFPageRef{471} /XYZ 72 532.631 null] >>}
\MusicPDFWriteObject{12699}{<< /D [\MusicPDFPageRef{471} /XYZ 72 510.557 null] >>}
\MusicPDFWriteObject{12700}{<< /D [\MusicPDFPageRef{471} /XYZ 224.917 266.078 null] >>}
\MusicPDFWriteObject{12701}{<< /D [\MusicPDFPageRef{471} /XYZ 170.97 202.839 null] >>}
\MusicPDFWriteObject{12702}{<< /D [\MusicPDFPageRef{472} /XYZ 72 514.028 null] >>}
\MusicPDFWriteObject{12703}{<< /D [\MusicPDFPageRef{472} /XYZ 242.999 408.198 null] >>}
\MusicPDFWriteObject{12704}{<< /D [\MusicPDFPageRef{466} /XYZ 196.288 652.499 null] >>}
\MusicPDFWriteObject{12705}{<< /D [\MusicPDFPageRef{472} /XYZ 165.181 332.318 null] >>}
\MusicPDFWriteObject{12706}{<< /D [\MusicPDFPageRef{472} /XYZ 242.348 216.706 null] >>}
\MusicPDFWriteObject{12707}{<< /D [\MusicPDFPageRef{472} /XYZ 72 171.394 null] >>}
\MusicPDFWriteObject{12708}{<< /D [\MusicPDFPageRef{472} /XYZ 72 112.384 null] >>}
\MusicPDFWriteObject{12709}{<< /D [\MusicPDFPageRef{473} /XYZ 213.937 656.755 null] >>}
\MusicPDFWriteObject{12710}{<< /D [\MusicPDFPageRef{473} /XYZ 247.151 580.03 null] >>}
\MusicPDFWriteObject{12711}{<< /D [\MusicPDFPageRef{473} /XYZ 163.786 436.47 null] >>}
\MusicPDFWriteObject{12712}{<< /D [\MusicPDFPageRef{473} /XYZ 195.429 394.648 null] >>}
\MusicPDFWriteObject{12713}{<< /D [\MusicPDFPageRef{473} /XYZ 249.44 318.529 null] >>}
\MusicPDFWriteObject{12714}{<< /D [\MusicPDFPageRef{473} /XYZ 237.235 185.577 null] >>}
\MusicPDFWriteObject{12715}{<< /D [\MusicPDFPageRef{466} /XYZ 226.141 571.222 null] >>}
\MusicPDFWriteObject{12716}{<< /D [\MusicPDFPageRef{473} /XYZ 222.447 99.532 null] >>}
\MusicPDFWriteObject{12717}{<< /D [\MusicPDFPageRef{466} /XYZ 195.762 507.696 null] >>}
\MusicPDFWriteObject{12718}{<< /D [\MusicPDFPageRef{466} /XYZ 215.811 341.322 null] >>}
\MusicPDFWriteObject{12719}{<< /D [\MusicPDFPageRef{466} /XYZ 239.517 274.696 null] >>}
\MusicPDFWriteObject{12720}{<< /D [\MusicPDFPageRef{467} /XYZ 207.748 664.094 null] >>}
\MusicPDFWriteObject{12721}{<< /D [\MusicPDFPageRef{475} /XYZ 265.02 150.557 null] >>}
\MusicPDFWriteObject{12722}{<< /D [\MusicPDFPageRef{478} /XYZ 257.694 419.7 null] >>}
\MusicPDFWriteObject{12723}{<< /D [\MusicPDFPageRef{479} /XYZ 253.626 616.206 null] >>}
\MusicPDFWriteObject{12724}{<< /D [\MusicPDFPageRef{479} /XYZ 212.29 547.744 null] >>}
\MusicPDFWriteObject{12725}{<< /D [\MusicPDFPageRef{479} /XYZ 250.917 500.555 null] >>}
\MusicPDFWriteObject{12726}{<< /D [\MusicPDFPageRef{479} /XYZ 72 403.724 null] >>}
\MusicPDFWriteObject{12727}{<< /D [\MusicPDFPageRef{479} /XYZ 171.632 196.351 null] >>}
\MusicPDFWriteObject{12728}{<< /D [\MusicPDFPageRef{476} /XYZ 252.242 509.509 null] >>}
\MusicPDFWriteObject{12729}{<< /D [\MusicPDFPageRef{476} /XYZ 284.532 444.75 null] >>}
\MusicPDFWriteObject{12730}{<< /D [\MusicPDFPageRef{476} /XYZ 257.694 371.898 null] >>}
\MusicPDFWriteObject{12731}{<< /D [\MusicPDFPageRef{477} /XYZ 170.397 652.859 null] >>}
\MusicPDFWriteObject{12732}{<< /D [\MusicPDFPageRef{477} /XYZ 252.609 474.801 null] >>}
\MusicPDFWriteObject{12733}{<< /D [\MusicPDFPageRef{477} /XYZ 258.336 380.34 null] >>}
\MusicPDFWriteObject{12734}{<< /D [\MusicPDFPageRef{478} /XYZ 242.048 515.098 null] >>}
\MusicPDFWriteObject{12735}{<< /D [\MusicPDFPageRef{478} /XYZ 284.532 466.251 null] >>}
\MusicPDFWriteObject{12736}{<< /D [\MusicPDFPageRef{480} /XYZ 216.794 630.453 null] >>}
\MusicPDFWriteObject{12737}{<< /D [\MusicPDFPageRef{481} /XYZ 176.386 339.527 null] >>}
\MusicPDFWriteObject{12738}{<< /D [\MusicPDFPageRef{481} /XYZ 175.002 248.832 null] >>}
\MusicPDFWriteObject{12739}{<< /D [\MusicPDFPageRef{481} /XYZ 164.386 212.404 null] >>}
\MusicPDFWriteObject{12740}{<< /D [\MusicPDFPageRef{481} /XYZ 229.358 148.581 null] >>}
\MusicPDFWriteObject{12741}{<< /D [\MusicPDFPageRef{482} /XYZ 162.227 655.959 null] >>}
\MusicPDFWriteObject{12742}{<< /D [\MusicPDFPageRef{482} /XYZ 162.227 596.772 null] >>}
\MusicPDFWriteObject{12743}{<< /D [\MusicPDFPageRef{482} /XYZ 169.042 377.958 null] >>}
\MusicPDFWriteObject{12744}{<< /D [\MusicPDFPageRef{480} /XYZ 279.002 520.927 null] >>}
\MusicPDFWriteObject{12745}{<< /D [\MusicPDFPageRef{480} /XYZ 241.256 484.5 null] >>}
\MusicPDFWriteObject{12746}{<< /D [\MusicPDFPageRef{480} /XYZ 190.12 418.574 null] >>}
\MusicPDFWriteObject{12747}{<< /D [\MusicPDFPageRef{480} /XYZ 168.383 324.795 null] >>}
\MusicPDFWriteObject{12748}{<< /D [\MusicPDFPageRef{480} /XYZ 166.999 288.368 null] >>}
\MusicPDFWriteObject{12749}{<< /D [\MusicPDFPageRef{480} /XYZ 247.008 199.093 null] >>}
\MusicPDFWriteObject{12750}{<< /D [\MusicPDFPageRef{481} /XYZ 90.172 619.036 null] >>}
\MusicPDFWriteObject{12751}{<< /D [\MusicPDFPageRef{481} /XYZ 186.223 572.624 null] >>}
\MusicPDFWriteObject{12752}{<< /D [\MusicPDFPageRef{483} /XYZ 225.527 174.907 null] >>}
\MusicPDFWriteObject{12753}{<< /D [\MusicPDFPageRef{487} /XYZ 239.446 683.057 null] >>}
\MusicPDFWriteObject{12754}{<< /D [\MusicPDFPageRef{484} /XYZ 195.351 683.057 null] >>}
\MusicPDFWriteObject{12755}{<< /D [\MusicPDFPageRef{484} /XYZ 164.342 554.372 null] >>}
\MusicPDFWriteObject{12756}{<< /D [\MusicPDFPageRef{484} /XYZ 172.302 425.727 null] >>}
\MusicPDFWriteObject{12757}{<< /D [\MusicPDFPageRef{485} /XYZ 244.164 420.116 null] >>}
\MusicPDFWriteObject{12758}{<< /D [\MusicPDFPageRef{485} /XYZ 233.322 304.019 null] >>}
\MusicPDFWriteObject{12759}{<< /D [\MusicPDFPageRef{486} /XYZ 156.223 654.557 null] >>}
\MusicPDFWriteObject{12760}{<< /D [\MusicPDFPageRef{486} /XYZ 239.446 341.101 null] >>}
\MusicPDFWriteObject{12761}{<< /D [\MusicPDFPageRef{486} /XYZ 152.776 291.473 null] >>}
\MusicPDFWriteObject{12762}{<< /D [\MusicPDFPageRef{487} /XYZ 169.605 173.342 null] >>}
\MusicPDFWriteObject{12763}{<< /D [\MusicPDFPageRef{489} /XYZ 171.795 398.502 null] >>}
\MusicPDFWriteObject{12764}{<< /D [\MusicPDFPageRef{487} /XYZ 168.162 136.701 null] >>}
\MusicPDFWriteObject{12765}{<< /D [\MusicPDFPageRef{488} /XYZ 171.147 653.642 null] >>}
\MusicPDFWriteObject{12766}{<< /D [\MusicPDFPageRef{488} /XYZ 222.826 454.613 null] >>}
\MusicPDFWriteObject{12767}{<< /D [\MusicPDFPageRef{488} /XYZ 172.535 416.924 null] >>}
\MusicPDFWriteObject{12768}{<< /D [\MusicPDFPageRef{488} /XYZ 168.908 171.596 null] >>}
\MusicPDFWriteObject{12769}{<< /D [\MusicPDFPageRef{488} /XYZ 180.469 134.955 null] >>}
\MusicPDFWriteObject{12770}{<< /D [\MusicPDFPageRef{489} /XYZ 182.715 654.472 null] >>}
\MusicPDFWriteObject{12771}{<< /D [\MusicPDFPageRef{489} /XYZ 222.086 436.191 null] >>}
\MusicPDFWriteObject{12772}{<< /D [\MusicPDFPageRef{489} /XYZ 161.389 118.171 null] >>}
\MusicPDFWriteObject{12773}{<< /D [\MusicPDFPageRef{490} /XYZ 72 156.895 null] >>}
\MusicPDFWriteObject{12774}{<< /D [\MusicPDFPageRef{491} /XYZ 239.603 681.651 null] >>}
\MusicPDFWriteObject{12775}{<< /D [\MusicPDFPageRef{491} /XYZ 221.108 640.437 null] >>}
\MusicPDFWriteObject{12776}{<< /D [\MusicPDFPageRef{491} /XYZ 279.832 515.781 null] >>}
\MusicPDFWriteObject{12777}{<< /D [\MusicPDFPageRef{490} /XYZ 279.832 632.3 null] >>}
\MusicPDFWriteObject{12778}{<< /D [\MusicPDFPageRef{490} /XYZ 241.292 569.571 null] >>}
\MusicPDFWriteObject{12779}{<< /D [\MusicPDFPageRef{490} /XYZ 197.203 489.768 null] >>}
\MusicPDFWriteObject{12780}{<< /D [\MusicPDFPageRef{490} /XYZ 228.717 438.825 null] >>}
\MusicPDFWriteObject{12781}{<< /D [\MusicPDFPageRef{490} /XYZ 72 263.2 null] >>}
\MusicPDFWriteObject{12782}{<< /D [\MusicPDFPageRef{490} /XYZ 72 246.074 null] >>}
\MusicPDFWriteObject{12783}{<< /D [\MusicPDFPageRef{490} /XYZ 72 222.9 null] >>}
\MusicPDFWriteObject{12784}{<< /D [\MusicPDFPageRef{490} /XYZ 72 189.898 null] >>}
\MusicPDFWriteObject{12785}{<< /D [\MusicPDFPageRef{492} /XYZ 268.126 683.057 null] >>}
\MusicPDFWriteObject{12786}{<< /D [\MusicPDFPageRef{492} /XYZ 279.214 333.21 null] >>}
\MusicPDFWriteObject{12787}{<< /D [\MusicPDFPageRef{494} /XYZ 181.641 633.478 null] >>}
\MusicPDFWriteObject{12788}{<< /D [\MusicPDFPageRef{494} /XYZ 224.428 520.737 null] >>}
\MusicPDFWriteObject{12789}{<< /D [\MusicPDFPageRef{495} /XYZ 149.521 476.074 null] >>}
\MusicPDFWriteObject{12790}{<< /D [\MusicPDFPageRef{495} /XYZ 197.644 413.583 null] >>}
\MusicPDFWriteObject{12791}{<< /D [\MusicPDFPageRef{495} /XYZ 156.495 193.672 null] >>}
\MusicPDFWriteObject{12792}{<< /D [\MusicPDFPageRef{496} /XYZ 85.28 564.026 null] >>}
\MusicPDFWriteObject{12793}{<< /D [\MusicPDFPageRef{496} /XYZ 216.343 513.561 null] >>}
\MusicPDFWriteObject{12794}{<< /D [\MusicPDFPageRef{498} /XYZ 221.044 418.042 null] >>}
\MusicPDFWriteObject{12795}{<< /D [\MusicPDFPageRef{500} /XYZ 236.692 288.649 null] >>}
\MusicPDFWriteObject{12796}{<< /D [\MusicPDFPageRef{500} /XYZ 86.925 240.238 null] >>}
\MusicPDFWriteObject{12797}{<< /D [\MusicPDFPageRef{500} /XYZ 237.679 144.732 null] >>}
\MusicPDFWriteObject{12798}{<< /D [\MusicPDFPageRef{501} /XYZ 150.92 449.628 null] >>}
\MusicPDFWriteObject{12799}{<< /D [\MusicPDFPageRef{501} /XYZ 218.683 193.141 null] >>}
\MusicPDFWriteObject{12800}{<< /D [\MusicPDFPageRef{501} /XYZ 208.264 146.624 null] >>}
\MusicPDFWriteObject{12801}{<< /D [\MusicPDFPageRef{502} /XYZ 149.696 587.338 null] >>}
\MusicPDFWriteObject{12802}{<< /D [\MusicPDFPageRef{502} /XYZ 253.734 530.216 null] >>}
\MusicPDFWriteObject{12803}{<< /D [\MusicPDFPageRef{502} /XYZ 233.568 481.529 null] >>}
\MusicPDFWriteObject{12804}{<< /D [\MusicPDFPageRef{502} /XYZ 90.518 429.276 null] >>}
\MusicPDFWriteObject{12805}{<< /D [\MusicPDFPageRef{499} /XYZ 245.54 454.628 null] >>}
\MusicPDFWriteObject{12806}{<< /D [\MusicPDFPageRef{502} /XYZ 199.715 105.327 null] >>}
\MusicPDFWriteObject{12807}{<< /D [\MusicPDFPageRef{503} /XYZ 209.36 681.651 null] >>}
\MusicPDFWriteObject{12808}{<< /D [\MusicPDFPageRef{503} /XYZ 177.158 632.249 null] >>}
\MusicPDFWriteObject{12809}{<< /D [\MusicPDFPageRef{503} /XYZ 249.836 442.144 null] >>}
\MusicPDFWriteObject{12810}{<< /D [\MusicPDFPageRef{503} /XYZ 166.683 359.39 null] >>}
\MusicPDFWriteObject{12811}{<< /D [\MusicPDFPageRef{499} /XYZ 224.391 368.779 null] >>}
\MusicPDFWriteObject{12812}{<< /D [\MusicPDFPageRef{499} /XYZ 234.203 261.565 null] >>}
\MusicPDFWriteObject{12813}{<< /D [\MusicPDFPageRef{499} /XYZ 239.846 221.241 null] >>}
\MusicPDFWriteObject{12814}{<< /D [\MusicPDFPageRef{499} /XYZ 228.749 177.018 null] >>}
\MusicPDFWriteObject{12815}{<< /D [\MusicPDFPageRef{500} /XYZ 216.441 564.219 null] >>}
\MusicPDFWriteObject{12816}{<< /D [\MusicPDFPageRef{500} /XYZ 263.006 465.424 null] >>}
\MusicPDFWriteObject{12817}{<< /D [\MusicPDFPageRef{500} /XYZ 215.709 402.695 null] >>}
\MusicPDFWriteObject{12818}{<< /D [\MusicPDFPageRef{504} /XYZ 230.276 259.322 null] >>}
\MusicPDFWriteObject{12819}{<< /D [\MusicPDFPageRef{507} /XYZ 277.855 497.527 null] >>}
\MusicPDFWriteObject{12820}{<< /D [\MusicPDFPageRef{507} /XYZ 180.876 404.268 null] >>}
\MusicPDFWriteObject{12821}{<< /D [\MusicPDFPageRef{507} /XYZ 231.621 309.467 null] >>}
\MusicPDFWriteObject{12822}{<< /D [\MusicPDFPageRef{507} /XYZ 282.621 250.867 null] >>}
\MusicPDFWriteObject{12823}{<< /D [\MusicPDFPageRef{505} /XYZ 263.277 679.161 null] >>}
\MusicPDFWriteObject{12824}{<< /D [\MusicPDFPageRef{505} /XYZ 242.186 602.604 null] >>}
\MusicPDFWriteObject{12825}{<< /D [\MusicPDFPageRef{505} /XYZ 242.505 364.565 null] >>}
\MusicPDFWriteObject{12826}{<< /D [\MusicPDFPageRef{505} /XYZ 254.728 99.233 null] >>}
\MusicPDFWriteObject{12827}{<< /D [\MusicPDFPageRef{506} /XYZ 224.034 564.999 null] >>}
\MusicPDFWriteObject{12828}{<< /D [\MusicPDFPageRef{506} /XYZ 241.913 509.265 null] >>}
\MusicPDFWriteObject{12829}{<< /D [\MusicPDFPageRef{506} /XYZ 261.493 472.838 null] >>}
\MusicPDFWriteObject{12830}{<< /D [\MusicPDFPageRef{507} /XYZ 274.383 615.39 null] >>}
\MusicPDFWriteObject{12831}{<< /D [\MusicPDFPageRef{508} /XYZ 243.477 445.004 null] >>}
\MusicPDFWriteObject{12832}{<< /D [\MusicPDFPageRef{509} /XYZ 148.313 412.561 null] >>}
\MusicPDFWriteObject{12833}{<< /D [\MusicPDFPageRef{509} /XYZ 158.753 367.311 null] >>}
\MusicPDFWriteObject{12834}{<< /D [\MusicPDFPageRef{510} /XYZ 72 546.129 null] >>}
\MusicPDFWriteObject{12835}{<< /D [\MusicPDFPageRef{510} /XYZ 72 500.208 null] >>}
\MusicPDFWriteObject{12836}{<< /D [\MusicPDFPageRef{510} /XYZ 72 472.758 null] >>}
\MusicPDFWriteObject{12837}{<< /D [\MusicPDFPageRef{510} /XYZ 162.481 344.106 null] >>}
\MusicPDFWriteObject{12838}{<< /D [\MusicPDFPageRef{510} /XYZ 149.219 297.007 null] >>}
\MusicPDFWriteObject{12839}{<< /D [\MusicPDFPageRef{510} /XYZ 254.116 197.304 null] >>}
\MusicPDFWriteObject{12840}{<< /D [\MusicPDFPageRef{510} /XYZ 199.465 111.977 null] >>}
\MusicPDFWriteObject{12841}{<< /D [\MusicPDFPageRef{511} /XYZ 72 661.818 null] >>}
\MusicPDFWriteObject{12842}{<< /D [\MusicPDFPageRef{508} /XYZ 247.226 376.704 null] >>}
\MusicPDFWriteObject{12843}{<< /D [\MusicPDFPageRef{511} /XYZ 72 573.955 null] >>}
\MusicPDFWriteObject{12844}{<< /D [\MusicPDFPageRef{511} /XYZ 155.9 487.685 null] >>}
\MusicPDFWriteObject{12845}{<< /D [\MusicPDFPageRef{511} /XYZ 171.008 288.526 null] >>}
\MusicPDFWriteObject{12846}{<< /D [\MusicPDFPageRef{511} /XYZ 180.706 242.454 null] >>}
\MusicPDFWriteObject{12847}{<< /D [\MusicPDFPageRef{511} /XYZ 74.486 188.294 null] >>}
\MusicPDFWriteObject{12848}{<< /D [\MusicPDFPageRef{511} /XYZ 164.077 127.631 null] >>}
\MusicPDFWriteObject{12849}{<< /D [\MusicPDFPageRef{512} /XYZ 210.668 379.891 null] >>}
\MusicPDFWriteObject{12850}{<< /D [\MusicPDFPageRef{512} /XYZ 209.57 342.058 null] >>}
\MusicPDFWriteObject{12851}{<< /D [\MusicPDFPageRef{512} /XYZ 259.065 272.231 null] >>}
\MusicPDFWriteObject{12852}{<< /D [\MusicPDFPageRef{512} /XYZ 273.649 130.597 null] >>}
\MusicPDFWriteObject{12853}{<< /D [\MusicPDFPageRef{508} /XYZ 154.275 284.331 null] >>}
\MusicPDFWriteObject{12854}{<< /D [\MusicPDFPageRef{513} /XYZ 295.025 656.755 null] >>}
\MusicPDFWriteObject{12855}{<< /D [\MusicPDFPageRef{513} /XYZ 156.422 366.592 null] >>}
\MusicPDFWriteObject{12856}{<< /D [\MusicPDFPageRef{508} /XYZ 72 244.065 null] >>}
\MusicPDFWriteObject{12857}{<< /D [\MusicPDFPageRef{508} /XYZ 72 215.4 null] >>}
\MusicPDFWriteObject{12858}{<< /D [\MusicPDFPageRef{508} /XYZ 72 188.478 null] >>}
\MusicPDFWriteObject{12859}{<< /D [\MusicPDFPageRef{508} /XYZ 72 161.557 null] >>}
\MusicPDFWriteObject{12860}{<< /D [\MusicPDFPageRef{508} /XYZ 210.825 117.447 null] >>}
\MusicPDFWriteObject{12861}{<< /D [\MusicPDFPageRef{509} /XYZ 164.563 679.16 null] >>}
\MusicPDFWriteObject{12862}{<< /D [\MusicPDFPageRef{514} /XYZ 72 457.711 null] >>}
\MusicPDFWriteObject{12863}{<< /D [\MusicPDFPageRef{517} /XYZ 72 589.318 null] >>}
\MusicPDFWriteObject{12864}{<< /D [\MusicPDFPageRef{517} /XYZ 248.495 319.922 null] >>}
\MusicPDFWriteObject{12865}{<< /D [\MusicPDFPageRef{517} /XYZ 74.21 265.756 null] >>}
\MusicPDFWriteObject{12866}{<< /D [\MusicPDFPageRef{517} /XYZ 72 231.305 null] >>}
\MusicPDFWriteObject{12867}{<< /D [\MusicPDFPageRef{517} /XYZ 72 183.413 null] >>}
\MusicPDFWriteObject{12868}{<< /D [\MusicPDFPageRef{517} /XYZ 206.903 136.538 null] >>}
\MusicPDFWriteObject{12869}{<< /D [\MusicPDFPageRef{518} /XYZ 86.251 671.073 null] >>}
\MusicPDFWriteObject{12870}{<< /D [\MusicPDFPageRef{518} /XYZ 259.971 581.999 null] >>}
\MusicPDFWriteObject{12871}{<< /D [\MusicPDFPageRef{518} /XYZ 173.545 517.642 null] >>}
\MusicPDFWriteObject{12872}{<< /D [\MusicPDFPageRef{518} /XYZ 195.215 464.938 null] >>}
\MusicPDFWriteObject{12873}{<< /D [\MusicPDFPageRef{514} /XYZ 72 399.362 null] >>}
\MusicPDFWriteObject{12874}{<< /D [\MusicPDFPageRef{518} /XYZ 255.3 420.061 null] >>}
\MusicPDFWriteObject{12875}{<< /D [\MusicPDFPageRef{518} /XYZ 199.954 374.737 null] >>}
\MusicPDFWriteObject{12876}{<< /D [\MusicPDFPageRef{518} /XYZ 198.528 328.05 null] >>}
\MusicPDFWriteObject{12877}{<< /D [\MusicPDFPageRef{518} /XYZ 72 245.688 null] >>}
\MusicPDFWriteObject{12878}{<< /D [\MusicPDFPageRef{518} /XYZ 72 196.116 null] >>}
\MusicPDFWriteObject{12879}{<< /D [\MusicPDFPageRef{519} /XYZ 72 663.919 null] >>}
\MusicPDFWriteObject{12880}{<< /D [\MusicPDFPageRef{519} /XYZ 72 628.155 null] >>}
\MusicPDFWriteObject{12881}{<< /D [\MusicPDFPageRef{519} /XYZ 72 592.39 null] >>}
\MusicPDFWriteObject{12882}{<< /D [\MusicPDFPageRef{519} /XYZ 72 565.468 null] >>}
\MusicPDFWriteObject{12883}{<< /D [\MusicPDFPageRef{514} /XYZ 72 363.597 null] >>}
\MusicPDFWriteObject{12884}{<< /D [\MusicPDFPageRef{515} /XYZ 263.648 403.94 null] >>}
\MusicPDFWriteObject{12885}{<< /D [\MusicPDFPageRef{516} /XYZ 214.533 536.983 null] >>}
\MusicPDFWriteObject{12886}{<< /D [\MusicPDFPageRef{516} /XYZ 225.112 414.548 null] >>}
\MusicPDFWriteObject{12887}{<< /D [\MusicPDFPageRef{516} /XYZ 72 352.986 null] >>}
\MusicPDFWriteObject{12888}{<< /D [\MusicPDFPageRef{517} /XYZ 72 659.498 null] >>}
\MusicPDFWriteObject{12889}{<< /D [\MusicPDFPageRef{517} /XYZ 72 618.21 null] >>}
\MusicPDFWriteObject{12890}{<< /D [\MusicPDFPageRef{519} /XYZ 72 300.116 null] >>}
\MusicPDFWriteObject{12891}{<< /D [\MusicPDFPageRef{519} /XYZ 72 241.766 null] >>}
\MusicPDFWriteObject{12892}{<< /D [\MusicPDFPageRef{519} /XYZ 72 211.098 null] >>}
\MusicPDFWriteObject{12893}{<< /D [\MusicPDFPageRef{520} /XYZ 203.27 471.109 null] >>}
\MusicPDFWriteObject{12894}{<< /D [\MusicPDFPageRef{522} /XYZ 240.045 499.544 null] >>}
\MusicPDFWriteObject{12895}{<< /D [\MusicPDFPageRef{522} /XYZ 161.767 429.044 null] >>}
\MusicPDFWriteObject{12896}{<< /D [\MusicPDFPageRef{522} /XYZ 72 359.578 null] >>}
\MusicPDFWriteObject{12897}{<< /D [\MusicPDFPageRef{522} /XYZ 72 300.454 null] >>}
\MusicPDFWriteObject{12898}{<< /D [\MusicPDFPageRef{522} /XYZ 72 264.153 null] >>}
\MusicPDFWriteObject{12899}{<< /D [\MusicPDFPageRef{524} /XYZ 235.038 348.111 null] >>}
\MusicPDFWriteObject{12900}{<< /D [\MusicPDFPageRef{527} /XYZ 245.634 245.492 null] >>}
\MusicPDFWriteObject{12901}{<< /D [\MusicPDFPageRef{528} /XYZ 265.357 643.604 null] >>}
\MusicPDFWriteObject{12902}{<< /D [\MusicPDFPageRef{528} /XYZ 252.693 594.026 null] >>}
\MusicPDFWriteObject{12903}{<< /D [\MusicPDFPageRef{528} /XYZ 248.519 512.551 null] >>}
\MusicPDFWriteObject{12904}{<< /D [\MusicPDFPageRef{524} /XYZ 273.531 272.231 null] >>}
\MusicPDFWriteObject{12905}{<< /D [\MusicPDFPageRef{524} /XYZ 268.025 117.447 null] >>}
\MusicPDFWriteObject{12906}{<< /D [\MusicPDFPageRef{525} /XYZ 229.718 439.382 null] >>}
\MusicPDFWriteObject{12907}{<< /D [\MusicPDFPageRef{525} /XYZ 269.812 376.653 null] >>}
\MusicPDFWriteObject{12908}{<< /D [\MusicPDFPageRef{525} /XYZ 273.562 277.193 null] >>}
\MusicPDFWriteObject{12909}{<< /D [\MusicPDFPageRef{527} /XYZ 221.748 572.011 null] >>}
\MusicPDFWriteObject{12910}{<< /D [\MusicPDFPageRef{527} /XYZ 254.903 406.471 null] >>}
\MusicPDFWriteObject{12911}{<< /D [\MusicPDFPageRef{527} /XYZ 267.632 321.372 null] >>}
\MusicPDFWriteObject{12912}{<< /D [\MusicPDFPageRef{530} /XYZ 271.938 388.381 null] >>}
\MusicPDFWriteObject{12913}{<< /D [\MusicPDFPageRef{531} /XYZ 221.526 441.86 null] >>}
\MusicPDFWriteObject{12914}{<< /D [\MusicPDFPageRef{531} /XYZ 259.26 403.315 null] >>}
\MusicPDFWriteObject{12915}{<< /D [\MusicPDFPageRef{531} /XYZ 242.962 130.597 null] >>}
\MusicPDFWriteObject{12916}{<< /D [\MusicPDFPageRef{532} /XYZ 250.366 637.335 null] >>}
\MusicPDFWriteObject{12917}{<< /D [\MusicPDFPageRef{532} /XYZ 253.333 507.334 null] >>}
\MusicPDFWriteObject{12918}{<< /D [\MusicPDFPageRef{532} /XYZ 230.63 404.826 null] >>}
\MusicPDFWriteObject{12919}{<< /D [\MusicPDFPageRef{534} /XYZ 247.014 490.026 null] >>}
\MusicPDFWriteObject{12920}{<< /D [\MusicPDFPageRef{60} /XYZ 207.797 209.27 null] >>}
\MusicPDFWriteObject{12921}{<< /D [\MusicPDFPageRef{60} /XYZ 278.158 160.208 null] >>}
\MusicPDFWriteObject{12922}{<< /D [\MusicPDFPageRef{60} /XYZ 245.55 99.233 null] >>}
\MusicPDFWriteObject{12923}{<< /D [\MusicPDFPageRef{61} /XYZ 237.592 683.057 null] >>}
\MusicPDFWriteObject{12924}{<< /D [\MusicPDFPageRef{65} /XYZ 262.299 623.481 null] >>}
\MusicPDFWriteObject{12925}{<< /D [\MusicPDFPageRef{65} /XYZ 253.324 143.748 null] >>}
\MusicPDFWriteObject{12926}{<< /D [\MusicPDFPageRef{66} /XYZ 239.438 518.11 null] >>}
\MusicPDFWriteObject{12927}{<< /D [\MusicPDFPageRef{66} /XYZ 230.181 458.715 null] >>}
\MusicPDFWriteObject{12928}{<< /D [\MusicPDFPageRef{66} /XYZ 213.456 397.428 null] >>}
\MusicPDFWriteObject{12929}{<< /D [\MusicPDFPageRef{67} /XYZ 254.031 153.573 null] >>}
\MusicPDFWriteObject{12930}{<< /D [\MusicPDFPageRef{68} /XYZ 242.75 409.954 null] >>}
\MusicPDFWriteObject{12931}{<< /D [\MusicPDFPageRef{69} /XYZ 209.24 156.899 null] >>}
\MusicPDFWriteObject{12932}{<< /D [\MusicPDFPageRef{71} /XYZ 231.106 641.187 null] >>}
\MusicPDFWriteObject{12933}{<< /D [\MusicPDFPageRef{73} /XYZ 251.154 585.34 null] >>}
\MusicPDFWriteObject{12934}{<< /D [\MusicPDFPageRef{73} /XYZ 269.802 548.913 null] >>}
\MusicPDFWriteObject{12935}{<< /D [\MusicPDFPageRef{73} /XYZ 246.416 512.485 null] >>}
\MusicPDFWriteObject{12936}{<< /D [\MusicPDFPageRef{73} /XYZ 258.869 449.756 null] >>}
\MusicPDFWriteObject{12937}{<< /D [\MusicPDFPageRef{74} /XYZ 247.854 527.851 null] >>}
\MusicPDFWriteObject{12938}{<< /D [\MusicPDFPageRef{71} /XYZ 248.66 542.522 null] >>}
\MusicPDFWriteObject{12939}{<< /D [\MusicPDFPageRef{71} /XYZ 240.891 505.852 null] >>}
\MusicPDFWriteObject{12940}{<< /D [\MusicPDFPageRef{71} /XYZ 248.604 417.968 null] >>}
\MusicPDFWriteObject{12941}{<< /D [\MusicPDFPageRef{71} /XYZ 258.35 378.409 null] >>}
\MusicPDFWriteObject{12942}{<< /D [\MusicPDFPageRef{72} /XYZ 216.94 682.813 null] >>}
\MusicPDFWriteObject{12943}{<< /D [\MusicPDFPageRef{72} /XYZ 166.295 260.775 null] >>}
\MusicPDFWriteObject{12944}{<< /D [\MusicPDFPageRef{72} /XYZ 245.062 167.142 null] >>}
\MusicPDFWriteObject{12945}{<< /D [\MusicPDFPageRef{73} /XYZ 260.845 648.07 null] >>}
\MusicPDFWriteObject{12946}{<< /D [\MusicPDFPageRef{75} /XYZ 217.384 566.319 null] >>}
\MusicPDFWriteObject{12947}{<< /D [\MusicPDFPageRef{75} /XYZ 257.049 529.892 null] >>}
\MusicPDFWriteObject{12948}{<< /D [\MusicPDFPageRef{76} /XYZ 267.025 501.994 null] >>}
\MusicPDFWriteObject{12949}{<< /D [\MusicPDFPageRef{76} /XYZ 209.599 393.384 null] >>}
\MusicPDFWriteObject{12950}{<< /D [\MusicPDFPageRef{76} /XYZ 153.843 347.914 null] >>}
\MusicPDFWriteObject{12951}{<< /D [\MusicPDFPageRef{77} /XYZ 180.311 276.756 null] >>}
\MusicPDFWriteObject{12952}{<< /D [\MusicPDFPageRef{79} /XYZ 151.465 250.863 null] >>}
\MusicPDFWriteObject{12953}{<< /D [\MusicPDFPageRef{80} /XYZ 115.598 675.286 null] >>}
\MusicPDFWriteObject{12954}{<< /D [\MusicPDFPageRef{80} /XYZ 245.997 135.369 null] >>}
\MusicPDFWriteObject{12955}{<< /D [\MusicPDFPageRef{81} /XYZ 248.532 682.813 null] >>}
\MusicPDFWriteObject{12956}{<< /D [\MusicPDFPageRef{81} /XYZ 259.022 584.125 null] >>}
\MusicPDFWriteObject{12957}{<< /D [\MusicPDFPageRef{77} /XYZ 270.948 110.175 null] >>}
\MusicPDFWriteObject{12958}{<< /D [\MusicPDFPageRef{78} /XYZ 234.178 656.755 null] >>}
\MusicPDFWriteObject{12959}{<< /D [\MusicPDFPageRef{78} /XYZ 176.827 384.777 null] >>}
\MusicPDFWriteObject{12960}{<< /D [\MusicPDFPageRef{78} /XYZ 221.883 99.233 null] >>}
\MusicPDFWriteObject{12961}{<< /D [\MusicPDFPageRef{79} /XYZ 195.207 683.057 null] >>}
\MusicPDFWriteObject{12962}{<< /D [\MusicPDFPageRef{79} /XYZ 228.429 517.793 null] >>}
\MusicPDFWriteObject{12963}{<< /D [\MusicPDFPageRef{79} /XYZ 264.735 385.933 null] >>}
\MusicPDFWriteObject{12964}{<< /D [\MusicPDFPageRef{79} /XYZ 250.893 322.8 null] >>}
\MusicPDFWriteObject{12965}{<< /D [\MusicPDFPageRef{81} /XYZ 221.713 110.054 null] >>}
\MusicPDFWriteObject{12966}{<< /D [\MusicPDFPageRef{84} /XYZ 240.382 580.875 null] >>}
\MusicPDFWriteObject{12967}{<< /D [\MusicPDFPageRef{84} /XYZ 263.969 478.694 null] >>}
\MusicPDFWriteObject{12968}{<< /D [\MusicPDFPageRef{84} /XYZ 169.805 415.233 null] >>}
\MusicPDFWriteObject{12969}{<< /D [\MusicPDFPageRef{84} /XYZ 202.992 324.99 null] >>}
\MusicPDFWriteObject{12970}{<< /D [\MusicPDFPageRef{85} /XYZ 271.183 288.534 null] >>}
\MusicPDFWriteObject{12971}{<< /D [\MusicPDFPageRef{86} /XYZ 223.131 643.361 null] >>}
\MusicPDFWriteObject{12972}{<< /D [\MusicPDFPageRef{82} /XYZ 201.855 395.13 null] >>}
\MusicPDFWriteObject{12973}{<< /D [\MusicPDFPageRef{82} /XYZ 198.465 219.002 null] >>}
\MusicPDFWriteObject{12974}{<< /D [\MusicPDFPageRef{83} /XYZ 258.651 683.057 null] >>}
\MusicPDFWriteObject{12975}{<< /D [\MusicPDFPageRef{83} /XYZ 219.408 633.318 null] >>}
\MusicPDFWriteObject{12976}{<< /D [\MusicPDFPageRef{83} /XYZ 246.179 570.429 null] >>}
\MusicPDFWriteObject{12977}{<< /D [\MusicPDFPageRef{83} /XYZ 200.195 404.753 null] >>}
\MusicPDFWriteObject{12978}{<< /D [\MusicPDFPageRef{83} /XYZ 234.178 222.19 null] >>}
\MusicPDFWriteObject{12979}{<< /D [\MusicPDFPageRef{84} /XYZ 239.144 630.453 null] >>}
\MusicPDFWriteObject{12980}{<< /D [\MusicPDFPageRef{87} /XYZ 176.71 683.057 null] >>}
\MusicPDFWriteObject{12981}{<< /D [\MusicPDFPageRef{89} /XYZ 251.063 494.008 null] >>}
\MusicPDFWriteObject{12982}{<< /D [\MusicPDFPageRef{89} /XYZ 118.262 410.367 null] >>}
\MusicPDFWriteObject{12983}{<< /D [\MusicPDFPageRef{89} /XYZ 257.533 204.138 null] >>}
\MusicPDFWriteObject{12984}{<< /D [\MusicPDFPageRef{90} /XYZ 267.032 669.906 null] >>}
\MusicPDFWriteObject{12985}{<< /D [\MusicPDFPageRef{90} /XYZ 230.493 501.326 null] >>}
\MusicPDFWriteObject{12986}{<< /D [\MusicPDFPageRef{90} /XYZ 270.794 458.046 null] >>}
\MusicPDFWriteObject{12987}{<< /D [\MusicPDFPageRef{90} /XYZ 231.218 421.618 null] >>}
\MusicPDFWriteObject{12988}{<< /D [\MusicPDFPageRef{90} /XYZ 144.008 314.655 null] >>}
\MusicPDFWriteObject{12989}{<< /D [\MusicPDFPageRef{90} /XYZ 230.441 240.59 null] >>}
\MusicPDFWriteObject{12990}{<< /D [\MusicPDFPageRef{90} /XYZ 149.729 162.651 null] >>}
\MusicPDFWriteObject{12991}{<< /D [\MusicPDFPageRef{87} /XYZ 237.506 433.75 null] >>}
\MusicPDFWriteObject{12992}{<< /D [\MusicPDFPageRef{90} /XYZ 161.016 117.824 null] >>}
\MusicPDFWriteObject{12993}{<< /D [\MusicPDFPageRef{91} /XYZ 205.561 665.124 null] >>}
\MusicPDFWriteObject{12994}{<< /D [\MusicPDFPageRef{91} /XYZ 243.321 608.068 null] >>}
\MusicPDFWriteObject{12995}{<< /D [\MusicPDFPageRef{91} /XYZ 247.831 396.575 null] >>}
\MusicPDFWriteObject{12996}{<< /D [\MusicPDFPageRef{91} /XYZ 262.266 264.593 null] >>}
\MusicPDFWriteObject{12997}{<< /D [\MusicPDFPageRef{91} /XYZ 244.775 215.014 null] >>}
\MusicPDFWriteObject{12998}{<< /D [\MusicPDFPageRef{91} /XYZ 224.381 178.426 null] >>}
\MusicPDFWriteObject{12999}{<< /D [\MusicPDFPageRef{91} /XYZ 249.573 136.575 null] >>}
\MusicPDFWriteObject{13000}{<< /D [\MusicPDFPageRef{92} /XYZ 240.962 679.56 null] >>}
\MusicPDFWriteObject{13001}{<< /D [\MusicPDFPageRef{92} /XYZ 224.761 636.017 null] >>}
\MusicPDFWriteObject{13002}{<< /D [\MusicPDFPageRef{87} /XYZ 268.242 99.233 null] >>}
\MusicPDFWriteObject{13003}{<< /D [\MusicPDFPageRef{92} /XYZ 206.594 592.711 null] >>}
\MusicPDFWriteObject{13004}{<< /D [\MusicPDFPageRef{92} /XYZ 230.827 429.006 null] >>}
\MusicPDFWriteObject{13005}{<< /D [\MusicPDFPageRef{92} /XYZ 72 140.842 null] >>}
\MusicPDFWriteObject{13006}{<< /D [\MusicPDFPageRef{93} /XYZ 209.432 653.412 null] >>}
\MusicPDFWriteObject{13007}{<< /D [\MusicPDFPageRef{93} /XYZ 84.908 502.881 null] >>}
\MusicPDFWriteObject{13008}{<< /D [\MusicPDFPageRef{93} /XYZ 72 441.38 null] >>}
\MusicPDFWriteObject{13009}{<< /D [\MusicPDFPageRef{93} /XYZ 199.048 379.09 null] >>}
\MusicPDFWriteObject{13010}{<< /D [\MusicPDFPageRef{93} /XYZ 157.842 134.795 null] >>}
\MusicPDFWriteObject{13011}{<< /D [\MusicPDFPageRef{94} /XYZ 244.463 663.846 null] >>}
\MusicPDFWriteObject{13012}{<< /D [\MusicPDFPageRef{94} /XYZ 143.816 594.014 null] >>}
\MusicPDFWriteObject{13013}{<< /D [\MusicPDFPageRef{88} /XYZ 189.796 667.848 null] >>}
\MusicPDFWriteObject{13014}{<< /D [\MusicPDFPageRef{94} /XYZ 261.108 545.626 null] >>}
\MusicPDFWriteObject{13015}{<< /D [\MusicPDFPageRef{94} /XYZ 191.459 336.359 null] >>}
\MusicPDFWriteObject{13016}{<< /D [\MusicPDFPageRef{94} /XYZ 212.031 276.813 null] >>}
\MusicPDFWriteObject{13017}{<< /D [\MusicPDFPageRef{88} /XYZ 72 586.794 null] >>}
\MusicPDFWriteObject{13018}{<< /D [\MusicPDFPageRef{88} /XYZ 239.329 412.231 null] >>}
\MusicPDFWriteObject{13019}{<< /D [\MusicPDFPageRef{88} /XYZ 263.24 272.237 null] >>}
\MusicPDFWriteObject{13020}{<< /D [\MusicPDFPageRef{89} /XYZ 232.794 681.9 null] >>}
\MusicPDFWriteObject{13021}{<< /D [\MusicPDFPageRef{89} /XYZ 269.484 536.858 null] >>}
\MusicPDFWriteObject{13022}{<< /D [\MusicPDFPageRef{95} /XYZ 256.723 468.205 null] >>}
\MusicPDFWriteObject{13023}{<< /D [\MusicPDFPageRef{97} /XYZ 171.762 148.811 null] >>}
\MusicPDFWriteObject{13024}{<< /D [\MusicPDFPageRef{97} /XYZ 255.079 99.233 null] >>}
\MusicPDFWriteObject{13025}{<< /D [\MusicPDFPageRef{98} /XYZ 242.79 643.604 null] >>}
\MusicPDFWriteObject{13026}{<< /D [\MusicPDFPageRef{98} /XYZ 206.6 437.089 null] >>}
\MusicPDFWriteObject{13027}{<< /D [\MusicPDFPageRef{98} /XYZ 162.47 352.958 null] >>}
\MusicPDFWriteObject{13028}{<< /D [\MusicPDFPageRef{98} /XYZ 150.423 308.752 null] >>}
\MusicPDFWriteObject{13029}{<< /D [\MusicPDFPageRef{95} /XYZ 217.993 428.435 null] >>}
\MusicPDFWriteObject{13030}{<< /D [\MusicPDFPageRef{95} /XYZ 247.314 230.197 null] >>}
\MusicPDFWriteObject{13031}{<< /D [\MusicPDFPageRef{95} /XYZ 215.712 157.433 null] >>}
\MusicPDFWriteObject{13032}{<< /D [\MusicPDFPageRef{96} /XYZ 232.78 682.813 null] >>}
\MusicPDFWriteObject{13033}{<< /D [\MusicPDFPageRef{96} /XYZ 151.501 595.732 null] >>}
\MusicPDFWriteObject{13034}{<< /D [\MusicPDFPageRef{96} /XYZ 255.558 528.747 null] >>}
\MusicPDFWriteObject{13035}{<< /D [\MusicPDFPageRef{96} /XYZ 254.093 438.21 null] >>}
\MusicPDFWriteObject{13036}{<< /D [\MusicPDFPageRef{97} /XYZ 255.263 198.39 null] >>}
\MusicPDFWriteObject{13037}{<< /D [\MusicPDFPageRef{99} /XYZ 236.037 150.847 null] >>}
\MusicPDFWriteObject{13038}{<< /D [\MusicPDFPageRef{104} /XYZ 235.542 497.647 null] >>}
\MusicPDFWriteObject{13039}{<< /D [\MusicPDFPageRef{104} /XYZ 259.682 293.458 null] >>}
\MusicPDFWriteObject{13040}{<< /D [\MusicPDFPageRef{105} /XYZ 247.134 683.057 null] >>}
\MusicPDFWriteObject{13041}{<< /D [\MusicPDFPageRef{105} /XYZ 204.991 642.373 null] >>}
\MusicPDFWriteObject{13042}{<< /D [\MusicPDFPageRef{105} /XYZ 251.657 413.223 null] >>}
\MusicPDFWriteObject{13043}{<< /D [\MusicPDFPageRef{99} /XYZ 250.041 113.263 null] >>}
\MusicPDFWriteObject{13044}{<< /D [\MusicPDFPageRef{100} /XYZ 255.261 515.865 null] >>}
\MusicPDFWriteObject{13045}{<< /D [\MusicPDFPageRef{101} /XYZ 209.917 674.787 null] >>}
\MusicPDFWriteObject{13046}{<< /D [\MusicPDFPageRef{101} /XYZ 202.42 496.668 null] >>}
\MusicPDFWriteObject{13047}{<< /D [\MusicPDFPageRef{101} /XYZ 232.544 341.554 null] >>}
\MusicPDFWriteObject{13048}{<< /D [\MusicPDFPageRef{101} /XYZ 234.167 189.731 null] >>}
\MusicPDFWriteObject{13049}{<< /D [\MusicPDFPageRef{101} /XYZ 253.302 126.482 null] >>}
\MusicPDFWriteObject{13050}{<< /D [\MusicPDFPageRef{102} /XYZ 186.565 239.776 null] >>}
\MusicPDFWriteObject{13051}{<< /D [\MusicPDFPageRef{105} /XYZ 237.506 133.026 null] >>}
\MusicPDFWriteObject{13052}{<< /D [\MusicPDFPageRef{106} /XYZ 199.849 559.32 null] >>}
\MusicPDFWriteObject{13053}{<< /D [\MusicPDFPageRef{106} /XYZ 189.039 518.464 null] >>}
\MusicPDFWriteObject{13054}{<< /D [\MusicPDFPageRef{106} /XYZ 193.915 479.055 null] >>}
\MusicPDFWriteObject{13055}{<< /D [\MusicPDFPageRef{107} /XYZ 234.778 620.349 null] >>}
\MusicPDFWriteObject{13056}{<< /D [\MusicPDFPageRef{108} /XYZ 72 525.489 null] >>}
\MusicPDFWriteObject{13057}{<< /D [\MusicPDFPageRef{108} /XYZ 184.791 475.387 null] >>}
\MusicPDFWriteObject{13058}{<< /D [\MusicPDFPageRef{108} /XYZ 151.344 424.885 null] >>}
\MusicPDFWriteObject{13059}{<< /D [\MusicPDFPageRef{108} /XYZ 239.819 196.703 null] >>}
\MusicPDFWriteObject{13060}{<< /D [\MusicPDFPageRef{109} /XYZ 72 637.908 null] >>}
\MusicPDFWriteObject{13061}{<< /D [\MusicPDFPageRef{109} /XYZ 278.329 494.862 null] >>}
\MusicPDFWriteObject{13062}{<< /D [\MusicPDFPageRef{109} /XYZ 152.92 161.346 null] >>}
\MusicPDFWriteObject{13063}{<< /D [\MusicPDFPageRef{107} /XYZ 285.833 516.385 null] >>}
\MusicPDFWriteObject{13064}{<< /D [\MusicPDFPageRef{107} /XYZ 213.748 447.631 null] >>}
\MusicPDFWriteObject{13065}{<< /D [\MusicPDFPageRef{107} /XYZ 80.973 278.568 null] >>}
\MusicPDFWriteObject{13066}{<< /D [\MusicPDFPageRef{107} /XYZ 170.688 220.521 null] >>}
\MusicPDFWriteObject{13067}{<< /D [\MusicPDFPageRef{107} /XYZ 72 157.44 null] >>}
\MusicPDFWriteObject{13068}{<< /D [\MusicPDFPageRef{108} /XYZ 72 656.362 null] >>}
\MusicPDFWriteObject{13069}{<< /D [\MusicPDFPageRef{108} /XYZ 241.052 604.298 null] >>}
\MusicPDFWriteObject{13070}{<< /D [\MusicPDFPageRef{108} /XYZ 72 562.178 null] >>}
\MusicPDFWriteObject{13071}{<< /D [\MusicPDFPageRef{111} /XYZ 216.769 643.564 null] >>}
\MusicPDFWriteObject{13072}{<< /D [\MusicPDFPageRef{111} /XYZ 232.731 137.292 null] >>}
\MusicPDFWriteObject{13073}{<< /D [\MusicPDFPageRef{112} /XYZ 256.072 630.708 null] >>}
\MusicPDFWriteObject{13074}{<< /D [\MusicPDFPageRef{112} /XYZ 173.653 317.667 null] >>}
\MusicPDFWriteObject{13075}{<< /D [\MusicPDFPageRef{112} /XYZ 270.896 201.823 null] >>}
\MusicPDFWriteObject{13076}{<< /D [\MusicPDFPageRef{113} /XYZ 253.26 604.903 null] >>}
\MusicPDFWriteObject{13077}{<< /D [\MusicPDFPageRef{113} /XYZ 176.577 188.484 null] >>}
\MusicPDFWriteObject{13078}{<< /D [\MusicPDFPageRef{114} /XYZ 184.364 643.604 null] >>}
\MusicPDFWriteObject{13079}{<< /D [\MusicPDFPageRef{114} /XYZ 190.882 454.499 null] >>}
\MusicPDFWriteObject{13080}{<< /D [\MusicPDFPageRef{115} /XYZ 254.645 656.755 null] >>}
\MusicPDFWriteObject{13081}{<< /D [\MusicPDFPageRef{115} /XYZ 247.814 534.779 null] >>}
\MusicPDFWriteObject{13082}{<< /D [\MusicPDFPageRef{119} /XYZ 255.077 638.68 null] >>}
\MusicPDFWriteObject{13083}{<< /D [\MusicPDFPageRef{119} /XYZ 267.953 391.046 null] >>}
\MusicPDFWriteObject{13084}{<< /D [\MusicPDFPageRef{119} /XYZ 239.612 231.479 null] >>}
\MusicPDFWriteObject{13085}{<< /D [\MusicPDFPageRef{119} /XYZ 279.57 145.691 null] >>}
\MusicPDFWriteObject{13086}{<< /D [\MusicPDFPageRef{120} /XYZ 253.428 643.604 null] >>}
\MusicPDFWriteObject{13087}{<< /D [\MusicPDFPageRef{120} /XYZ 247.797 131.599 null] >>}
\MusicPDFWriteObject{13088}{<< /D [\MusicPDFPageRef{121} /XYZ 267.766 654.657 null] >>}
\MusicPDFWriteObject{13089}{<< /D [\MusicPDFPageRef{121} /XYZ 261.216 539.325 null] >>}
\MusicPDFWriteObject{13090}{<< /D [\MusicPDFPageRef{121} /XYZ 264.532 380.169 null] >>}
\MusicPDFWriteObject{13091}{<< /D [\MusicPDFPageRef{121} /XYZ 269.585 343.741 null] >>}
\MusicPDFWriteObject{13092}{<< /D [\MusicPDFPageRef{121} /XYZ 242.208 224.919 null] >>}
\MusicPDFWriteObject{13093}{<< /D [\MusicPDFPageRef{121} /XYZ 262.532 188.492 null] >>}
\MusicPDFWriteObject{13094}{<< /D [\MusicPDFPageRef{121} /XYZ 250.404 152.064 null] >>}
\MusicPDFWriteObject{13095}{<< /D [\MusicPDFPageRef{122} /XYZ 250.902 144.36 null] >>}
\MusicPDFWriteObject{13096}{<< /D [\MusicPDFPageRef{124} /XYZ 253.399 510.961 null] >>}
\MusicPDFWriteObject{13097}{<< /D [\MusicPDFPageRef{123} /XYZ 227.269 667.034 null] >>}
\MusicPDFWriteObject{13098}{<< /D [\MusicPDFPageRef{123} /XYZ 248.613 620.568 null] >>}
\MusicPDFWriteObject{13099}{<< /D [\MusicPDFPageRef{123} /XYZ 179.5 571.308 null] >>}
\MusicPDFWriteObject{13100}{<< /D [\MusicPDFPageRef{123} /XYZ 222.197 511.705 null] >>}
\MusicPDFWriteObject{13101}{<< /D [\MusicPDFPageRef{123} /XYZ 249.157 274.816 null] >>}
\MusicPDFWriteObject{13102}{<< /D [\MusicPDFPageRef{124} /XYZ 218.659 680.114 null] >>}
\MusicPDFWriteObject{13103}{<< /D [\MusicPDFPageRef{124} /XYZ 239.234 618.188 null] >>}
\MusicPDFWriteObject{13104}{<< /D [\MusicPDFPageRef{124} /XYZ 240.341 579.703 null] >>}
\MusicPDFWriteObject{13105}{<< /D [\MusicPDFPageRef{124} /XYZ 233.393 298.268 null] >>}
\MusicPDFWriteObject{13106}{<< /D [\MusicPDFPageRef{127} /XYZ 193.467 549.476 null] >>}
\MusicPDFWriteObject{13107}{<< /D [\MusicPDFPageRef{127} /XYZ 154.986 309.752 null] >>}
\MusicPDFWriteObject{13108}{<< /D [\MusicPDFPageRef{124} /XYZ 270.1 202.779 null] >>}
\MusicPDFWriteObject{13109}{<< /D [\MusicPDFPageRef{124} /XYZ 234.096 162.456 null] >>}
\MusicPDFWriteObject{13110}{<< /D [\MusicPDFPageRef{125} /XYZ 245.491 611.864 null] >>}
\MusicPDFWriteObject{13111}{<< /D [\MusicPDFPageRef{125} /XYZ 224.517 497.752 null] >>}
\MusicPDFWriteObject{13112}{<< /D [\MusicPDFPageRef{126} /XYZ 121.381 450.513 null] >>}
\MusicPDFWriteObject{13113}{<< /D [\MusicPDFPageRef{126} /XYZ 213.269 215.02 null] >>}
\MusicPDFWriteObject{13114}{<< /D [\MusicPDFPageRef{126} /XYZ 160.966 178.593 null] >>}
\MusicPDFWriteObject{13115}{<< /D [\MusicPDFPageRef{127} /XYZ 72 627.746 null] >>}
\MusicPDFWriteObject{13116}{<< /D [\MusicPDFPageRef{129} /XYZ 224.923 648.376 null] >>}
\MusicPDFWriteObject{13117}{<< /D [\MusicPDFPageRef{132} /XYZ 273.378 360.632 null] >>}
\MusicPDFWriteObject{13118}{<< /D [\MusicPDFPageRef{132} /XYZ 218.027 317.598 null] >>}
\MusicPDFWriteObject{13119}{<< /D [\MusicPDFPageRef{132} /XYZ 165.048 232.155 null] >>}
\MusicPDFWriteObject{13120}{<< /D [\MusicPDFPageRef{132} /XYZ 156.246 167.33 null] >>}
\MusicPDFWriteObject{13121}{<< /D [\MusicPDFPageRef{132} /XYZ 212.345 99.532 null] >>}
\MusicPDFWriteObject{13122}{<< /D [\MusicPDFPageRef{133} /XYZ 153.464 638.07 null] >>}
\MusicPDFWriteObject{13123}{<< /D [\MusicPDFPageRef{133} /XYZ 274.482 572.92 null] >>}
\MusicPDFWriteObject{13124}{<< /D [\MusicPDFPageRef{133} /XYZ 72 201.229 null] >>}
\MusicPDFWriteObject{13125}{<< /D [\MusicPDFPageRef{133} /XYZ 72 167.389 null] >>}
\MusicPDFWriteObject{13126}{<< /D [\MusicPDFPageRef{134} /XYZ 198.537 603.183 null] >>}
\MusicPDFWriteObject{13127}{<< /D [\MusicPDFPageRef{129} /XYZ 254.934 579.122 null] >>}
\MusicPDFWriteObject{13128}{<< /D [\MusicPDFPageRef{134} /XYZ 199.036 477.841 null] >>}
\MusicPDFWriteObject{13129}{<< /D [\MusicPDFPageRef{134} /XYZ 72 421.523 null] >>}
\MusicPDFWriteObject{13130}{<< /D [\MusicPDFPageRef{134} /XYZ 72 393.468 null] >>}
\MusicPDFWriteObject{13131}{<< /D [\MusicPDFPageRef{134} /XYZ 72 365.508 null] >>}
\MusicPDFWriteObject{13132}{<< /D [\MusicPDFPageRef{134} /XYZ 72 344.188 null] >>}
\MusicPDFWriteObject{13133}{<< /D [\MusicPDFPageRef{134} /XYZ 265.046 108.677 null] >>}
\MusicPDFWriteObject{13134}{<< /D [\MusicPDFPageRef{135} /XYZ 289.376 343.851 null] >>}
\MusicPDFWriteObject{13135}{<< /D [\MusicPDFPageRef{135} /XYZ 181.995 281.889 null] >>}
\MusicPDFWriteObject{13136}{<< /D [\MusicPDFPageRef{135} /XYZ 269.166 197.471 null] >>}
\MusicPDFWriteObject{13137}{<< /D [\MusicPDFPageRef{135} /XYZ 248.201 99.233 null] >>}
\MusicPDFWriteObject{13138}{<< /D [\MusicPDFPageRef{129} /XYZ 266.676 540.845 null] >>}
\MusicPDFWriteObject{13139}{<< /D [\MusicPDFPageRef{136} /XYZ 219.818 669.906 null] >>}
\MusicPDFWriteObject{13140}{<< /D [\MusicPDFPageRef{136} /XYZ 253.107 572.972 null] >>}
\MusicPDFWriteObject{13141}{<< /D [\MusicPDFPageRef{136} /XYZ 265.228 514.204 null] >>}
\MusicPDFWriteObject{13142}{<< /D [\MusicPDFPageRef{137} /XYZ 253.096 676.157 null] >>}
\MusicPDFWriteObject{13143}{<< /D [\MusicPDFPageRef{137} /XYZ 155.202 524.99 null] >>}
\MusicPDFWriteObject{13144}{<< /D [\MusicPDFPageRef{129} /XYZ 222.533 438.711 null] >>}
\MusicPDFWriteObject{13145}{<< /D [\MusicPDFPageRef{129} /XYZ 241.886 347.652 null] >>}
\MusicPDFWriteObject{13146}{<< /D [\MusicPDFPageRef{129} /XYZ 218.026 287.826 null] >>}
\MusicPDFWriteObject{13147}{<< /D [\MusicPDFPageRef{130} /XYZ 231.154 565.443 null] >>}
\MusicPDFWriteObject{13148}{<< /D [\MusicPDFPageRef{132} /XYZ 215.77 456.439 null] >>}
\MusicPDFWriteObject{13149}{<< /D [\MusicPDFPageRef{132} /XYZ 263.007 418.91 null] >>}
\MusicPDFWriteObject{13150}{<< /D [\MusicPDFPageRef{137} /XYZ 239.825 378.172 null] >>}
\MusicPDFWriteObject{13151}{<< /D [\MusicPDFPageRef{139} /XYZ 206.28 665.912 null] >>}
\MusicPDFWriteObject{13152}{<< /D [\MusicPDFPageRef{139} /XYZ 256.491 615.048 null] >>}
\MusicPDFWriteObject{13153}{<< /D [\MusicPDFPageRef{139} /XYZ 201.681 375.576 null] >>}
\MusicPDFWriteObject{13154}{<< /D [\MusicPDFPageRef{139} /XYZ 177.604 283.389 null] >>}
\MusicPDFWriteObject{13155}{<< /D [\MusicPDFPageRef{139} /XYZ 225.769 155.903 null] >>}
\MusicPDFWriteObject{13156}{<< /D [\MusicPDFPageRef{140} /XYZ 270.434 596.686 null] >>}
\MusicPDFWriteObject{13157}{<< /D [\MusicPDFPageRef{140} /XYZ 257.586 552.436 null] >>}
\MusicPDFWriteObject{13158}{<< /D [\MusicPDFPageRef{140} /XYZ 226.668 506.459 null] >>}
\MusicPDFWriteObject{13159}{<< /D [\MusicPDFPageRef{140} /XYZ 233.946 466 null] >>}
\MusicPDFWriteObject{13160}{<< /D [\MusicPDFPageRef{140} /XYZ 83.335 331.055 null] >>}
\MusicPDFWriteObject{13161}{<< /D [\MusicPDFPageRef{138} /XYZ 253.168 405.326 null] >>}
\MusicPDFWriteObject{13162}{<< /D [\MusicPDFPageRef{140} /XYZ 72 218.437 null] >>}
\MusicPDFWriteObject{13163}{<< /D [\MusicPDFPageRef{141} /XYZ 72 670.421 null] >>}
\MusicPDFWriteObject{13164}{<< /D [\MusicPDFPageRef{141} /XYZ 72 599.533 null] >>}
\MusicPDFWriteObject{13165}{<< /D [\MusicPDFPageRef{141} /XYZ 210.044 452.315 null] >>}
\MusicPDFWriteObject{13166}{<< /D [\MusicPDFPageRef{141} /XYZ 155.264 326.697 null] >>}
\MusicPDFWriteObject{13167}{<< /D [\MusicPDFPageRef{141} /XYZ 120.426 213.696 null] >>}
\MusicPDFWriteObject{13168}{<< /D [\MusicPDFPageRef{138} /XYZ 244.973 368.898 null] >>}
\MusicPDFWriteObject{13169}{<< /D [\MusicPDFPageRef{138} /XYZ 246.888 332.471 null] >>}
\MusicPDFWriteObject{13170}{<< /D [\MusicPDFPageRef{138} /XYZ 242.391 296.043 null] >>}
\MusicPDFWriteObject{13171}{<< /D [\MusicPDFPageRef{138} /XYZ 196.346 246.465 null] >>}
\MusicPDFWriteObject{13172}{<< /D [\MusicPDFPageRef{138} /XYZ 174.963 186.318 null] >>}
\MusicPDFWriteObject{13173}{<< /D [\MusicPDFPageRef{138} /XYZ 233.808 144.26 null] >>}
\MusicPDFWriteObject{13174}{<< /D [\MusicPDFPageRef{138} /XYZ 231.11 104.187 null] >>}
\MusicPDFWriteObject{13175}{<< /D [\MusicPDFPageRef{143} /XYZ 219.263 629.656 null] >>}
\MusicPDFWriteObject{13176}{<< /D [\MusicPDFPageRef{143} /XYZ 185.473 582.189 null] >>}
\MusicPDFWriteObject{13177}{<< /D [\MusicPDFPageRef{143} /XYZ 192.839 533.977 null] >>}
\MusicPDFWriteObject{13178}{<< /D [\MusicPDFPageRef{143} /XYZ 196.765 489.326 null] >>}
\MusicPDFWriteObject{13179}{<< /D [\MusicPDFPageRef{143} /XYZ 184.56 419.916 null] >>}
\MusicPDFWriteObject{13180}{<< /D [\MusicPDFPageRef{143} /XYZ 155.362 344.942 null] >>}
\MusicPDFWriteObject{13181}{<< /D [\MusicPDFPageRef{144} /XYZ 285.379 223.639 null] >>}
\MusicPDFWriteObject{13182}{<< /D [\MusicPDFPageRef{145} /XYZ 256.117 529.351 null] >>}
\MusicPDFWriteObject{13183}{<< /D [\MusicPDFPageRef{149} /XYZ 123.744 548.673 null] >>}
\MusicPDFWriteObject{13184}{<< /D [\MusicPDFPageRef{149} /XYZ 150.864 463.982 null] >>}
\MusicPDFWriteObject{13185}{<< /D [\MusicPDFPageRef{149} /XYZ 175.164 365.732 null] >>}
\MusicPDFWriteObject{13186}{<< /D [\MusicPDFPageRef{149} /XYZ 185.252 282.799 null] >>}
\MusicPDFWriteObject{13187}{<< /D [\MusicPDFPageRef{149} /XYZ 184.49 239.928 null] >>}
\MusicPDFWriteObject{13188}{<< /D [\MusicPDFPageRef{149} /XYZ 255.263 99.233 null] >>}
\MusicPDFWriteObject{13189}{<< /D [\MusicPDFPageRef{150} /XYZ 204.212 682.896 null] >>}
\MusicPDFWriteObject{13190}{<< /D [\MusicPDFPageRef{150} /XYZ 167.019 321.624 null] >>}
\MusicPDFWriteObject{13191}{<< /D [\MusicPDFPageRef{150} /XYZ 135.347 276.868 null] >>}
\MusicPDFWriteObject{13192}{<< /D [\MusicPDFPageRef{145} /XYZ 154.144 407.049 null] >>}
\MusicPDFWriteObject{13193}{<< /D [\MusicPDFPageRef{146} /XYZ 255.263 461.926 null] >>}
\MusicPDFWriteObject{13194}{<< /D [\MusicPDFPageRef{146} /XYZ 198.128 399.197 null] >>}
\MusicPDFWriteObject{13195}{<< /D [\MusicPDFPageRef{146} /XYZ 279.794 362.77 null] >>}
\MusicPDFWriteObject{13196}{<< /D [\MusicPDFPageRef{146} /XYZ 231.151 326.342 null] >>}
\MusicPDFWriteObject{13197}{<< /D [\MusicPDFPageRef{147} /XYZ 166.082 155.978 null] >>}
\MusicPDFWriteObject{13198}{<< /D [\MusicPDFPageRef{149} /XYZ 72 655.62 null] >>}
\MusicPDFWriteObject{13199}{<< /D [\MusicPDFPageRef{149} /XYZ 123.683 594.618 null] >>}
\MusicPDFWriteObject{13200}{<< /D [\MusicPDFPageRef{151} /XYZ 271.783 252.065 null] >>}
\MusicPDFWriteObject{13201}{<< /D [\MusicPDFPageRef{153} /XYZ 248.575 683.057 null] >>}
\MusicPDFWriteObject{13202}{<< /D [\MusicPDFPageRef{153} /XYZ 236.828 646.134 null] >>}
\MusicPDFWriteObject{13203}{<< /D [\MusicPDFPageRef{153} /XYZ 243.474 607.916 null] >>}
\MusicPDFWriteObject{13204}{<< /D [\MusicPDFPageRef{153} /XYZ 261.429 557.369 null] >>}
\MusicPDFWriteObject{13205}{<< /D [\MusicPDFPageRef{153} /XYZ 264.331 506.994 null] >>}
\MusicPDFWriteObject{13206}{<< /D [\MusicPDFPageRef{153} /XYZ 266.775 435.033 null] >>}
\MusicPDFWriteObject{13207}{<< /D [\MusicPDFPageRef{151} /XYZ 270.838 176.174 null] >>}
\MusicPDFWriteObject{13208}{<< /D [\MusicPDFPageRef{152} /XYZ 270.118 577.674 null] >>}
\MusicPDFWriteObject{13209}{<< /D [\MusicPDFPageRef{152} /XYZ 263.493 540.581 null] >>}
\MusicPDFWriteObject{13210}{<< /D [\MusicPDFPageRef{152} /XYZ 253.981 487.506 null] >>}
\MusicPDFWriteObject{13211}{<< /D [\MusicPDFPageRef{152} /XYZ 252.251 445.434 null] >>}
\MusicPDFWriteObject{13212}{<< /D [\MusicPDFPageRef{152} /XYZ 275.339 406.856 null] >>}
\MusicPDFWriteObject{13213}{<< /D [\MusicPDFPageRef{152} /XYZ 260.881 370.429 null] >>}
\MusicPDFWriteObject{13214}{<< /D [\MusicPDFPageRef{152} /XYZ 270.223 305.908 null] >>}
\MusicPDFWriteObject{13215}{<< /D [\MusicPDFPageRef{156} /XYZ 254.196 119.802 null] >>}
\MusicPDFWriteObject{13216}{<< /D [\MusicPDFPageRef{157} /XYZ 278.757 435.173 null] >>}
\MusicPDFWriteObject{13217}{<< /D [\MusicPDFPageRef{158} /XYZ 253.929 213.891 null] >>}
\MusicPDFWriteObject{13218}{<< /D [\MusicPDFPageRef{159} /XYZ 251.161 286.265 null] >>}
\MusicPDFWriteObject{13219}{<< /D [\MusicPDFPageRef{159} /XYZ 190.029 239.515 null] >>}
\MusicPDFWriteObject{13220}{<< /D [\MusicPDFPageRef{159} /XYZ 227.038 153.293 null] >>}
\MusicPDFWriteObject{13221}{<< /D [\MusicPDFPageRef{160} /XYZ 162.026 683.057 null] >>}
\MusicPDFWriteObject{13222}{<< /D [\MusicPDFPageRef{162} /XYZ 209.39 544.366 null] >>}
\MusicPDFWriteObject{13223}{<< /D [\MusicPDFPageRef{166} /XYZ 72 383.825 null] >>}
\MusicPDFWriteObject{13224}{<< /D [\MusicPDFPageRef{166} /XYZ 72 349.311 null] >>}
\MusicPDFWriteObject{13225}{<< /D [\MusicPDFPageRef{166} /XYZ 72 314.331 null] >>}
\MusicPDFWriteObject{13226}{<< /D [\MusicPDFPageRef{166} /XYZ 72 286.155 null] >>}
\MusicPDFWriteObject{13227}{<< /D [\MusicPDFPageRef{166} /XYZ 242.157 229.568 null] >>}
\MusicPDFWriteObject{13228}{<< /D [\MusicPDFPageRef{166} /XYZ 230.304 106.816 null] >>}
\MusicPDFWriteObject{13229}{<< /D [\MusicPDFPageRef{167} /XYZ 210.686 656.755 null] >>}
\MusicPDFWriteObject{13230}{<< /D [\MusicPDFPageRef{167} /XYZ 253.008 457.499 null] >>}
\MusicPDFWriteObject{13231}{<< /D [\MusicPDFPageRef{169} /XYZ 160.226 230.646 null] >>}
\MusicPDFWriteObject{13232}{<< /D [\MusicPDFPageRef{169} /XYZ 120.402 181.556 null] >>}
\MusicPDFWriteObject{13233}{<< /D [\MusicPDFPageRef{162} /XYZ 256.593 507.695 null] >>}
\MusicPDFWriteObject{13234}{<< /D [\MusicPDFPageRef{162} /XYZ 265.398 380.156 null] >>}
\MusicPDFWriteObject{13235}{<< /D [\MusicPDFPageRef{162} /XYZ 217.059 295.172 null] >>}
\MusicPDFWriteObject{13236}{<< /D [\MusicPDFPageRef{162} /XYZ 233.541 195.026 null] >>}
\MusicPDFWriteObject{13237}{<< /D [\MusicPDFPageRef{163} /XYZ 251.799 656.755 null] >>}
\MusicPDFWriteObject{13238}{<< /D [\MusicPDFPageRef{163} /XYZ 169.894 600.248 null] >>}
\MusicPDFWriteObject{13239}{<< /D [\MusicPDFPageRef{163} /XYZ 274.488 186.952 null] >>}
\MusicPDFWriteObject{13240}{<< /D [\MusicPDFPageRef{166} /XYZ 196.882 440.785 null] >>}
\MusicPDFWriteObject{13241}{<< /D [\MusicPDFPageRef{170} /XYZ 256.474 122.061 null] >>}
\MusicPDFWriteObject{13242}{<< /D [\MusicPDFPageRef{171} /XYZ 193.605 511.476 null] >>}
\MusicPDFWriteObject{13243}{<< /D [\MusicPDFPageRef{172} /XYZ 260.014 333.189 null] >>}
\MusicPDFWriteObject{13244}{<< /D [\MusicPDFPageRef{172} /XYZ 254.414 293.266 null] >>}
\MusicPDFWriteObject{13245}{<< /D [\MusicPDFPageRef{172} /XYZ 153.06 149.364 null] >>}
\MusicPDFWriteObject{13246}{<< /D [\MusicPDFPageRef{173} /XYZ 242.032 669.906 null] >>}
\MusicPDFWriteObject{13247}{<< /D [\MusicPDFPageRef{173} /XYZ 189.326 298.366 null] >>}
\MusicPDFWriteObject{13248}{<< /D [\MusicPDFPageRef{175} /XYZ 269.756 609.721 null] >>}
\MusicPDFWriteObject{13249}{<< /D [\MusicPDFPageRef{177} /XYZ 72 387.597 null] >>}
\MusicPDFWriteObject{13250}{<< /D [\MusicPDFPageRef{177} /XYZ 72 366.525 null] >>}
\MusicPDFWriteObject{13251}{<< /D [\MusicPDFPageRef{175} /XYZ 229.952 558.737 null] >>}
\MusicPDFWriteObject{13252}{<< /D [\MusicPDFPageRef{175} /XYZ 253.209 153.874 null] >>}
\MusicPDFWriteObject{13253}{<< /D [\MusicPDFPageRef{175} /XYZ 257.924 117.447 null] >>}
\MusicPDFWriteObject{13254}{<< /D [\MusicPDFPageRef{176} /XYZ 218.685 610.306 null] >>}
\MusicPDFWriteObject{13255}{<< /D [\MusicPDFPageRef{176} /XYZ 246.77 573.879 null] >>}
\MusicPDFWriteObject{13256}{<< /D [\MusicPDFPageRef{176} /XYZ 263.55 446.216 null] >>}
\MusicPDFWriteObject{13257}{<< /D [\MusicPDFPageRef{177} /XYZ 266.836 683.057 null] >>}
\MusicPDFWriteObject{13258}{<< /D [\MusicPDFPageRef{177} /XYZ 72 414.519 null] >>}
\MusicPDFWriteObject{13259}{<< /D [\MusicPDFPageRef{178} /XYZ 172.184 600.809 null] >>}
\MusicPDFWriteObject{13260}{<< /D [\MusicPDFPageRef{178} /XYZ 214.428 440.993 null] >>}
\MusicPDFWriteObject{13261}{<< /D [\MusicPDFPageRef{178} /XYZ 241.631 404.565 null] >>}
\MusicPDFWriteObject{13262}{<< /D [\MusicPDFPageRef{178} /XYZ 232.304 341.836 null] >>}
\MusicPDFWriteObject{13263}{<< /D [\MusicPDFPageRef{178} /XYZ 229.654 279.107 null] >>}
\MusicPDFWriteObject{13264}{<< /D [\MusicPDFPageRef{178} /XYZ 229.211 229.286 null] >>}
\MusicPDFWriteObject{13265}{<< /D [\MusicPDFPageRef{178} /XYZ 230.041 178.631 null] >>}
\MusicPDFWriteObject{13266}{<< /D [\MusicPDFPageRef{179} /XYZ 263.045 316.083 null] >>}
\MusicPDFWriteObject{13267}{<< /D [\MusicPDFPageRef{181} /XYZ 220.692 635.292 null] >>}
\MusicPDFWriteObject{13268}{<< /D [\MusicPDFPageRef{184} /XYZ 275.6 156.899 null] >>}
\MusicPDFWriteObject{13269}{<< /D [\MusicPDFPageRef{186} /XYZ 219.82 466.81 null] >>}
\MusicPDFWriteObject{13270}{<< /D [\MusicPDFPageRef{186} /XYZ 207.083 99.233 null] >>}
\MusicPDFWriteObject{13271}{<< /D [\MusicPDFPageRef{187} /XYZ 93.944 671.073 null] >>}
\MusicPDFWriteObject{13272}{<< /D [\MusicPDFPageRef{187} /XYZ 179.044 567.849 null] >>}
\MusicPDFWriteObject{13273}{<< /D [\MusicPDFPageRef{187} /XYZ 169.475 238.378 null] >>}
\MusicPDFWriteObject{13274}{<< /D [\MusicPDFPageRef{187} /XYZ 262.713 197.556 null] >>}
\MusicPDFWriteObject{13275}{<< /D [\MusicPDFPageRef{188} /XYZ 206.508 503.993 null] >>}
\MusicPDFWriteObject{13276}{<< /D [\MusicPDFPageRef{181} /XYZ 240.079 489.246 null] >>}
\MusicPDFWriteObject{13277}{<< /D [\MusicPDFPageRef{182} /XYZ 207.956 577.85 null] >>}
\MusicPDFWriteObject{13278}{<< /D [\MusicPDFPageRef{182} /XYZ 113.455 537.273 null] >>}
\MusicPDFWriteObject{13279}{<< /D [\MusicPDFPageRef{182} /XYZ 180.081 434.049 null] >>}
\MusicPDFWriteObject{13280}{<< /D [\MusicPDFPageRef{183} /XYZ 212.944 630.135 null] >>}
\MusicPDFWriteObject{13281}{<< /D [\MusicPDFPageRef{183} /XYZ 263.315 589.313 null] >>}
\MusicPDFWriteObject{13282}{<< /D [\MusicPDFPageRef{183} /XYZ 207.166 264.952 null] >>}
\MusicPDFWriteObject{13283}{<< /D [\MusicPDFPageRef{183} /XYZ 237.535 158.358 null] >>}
\MusicPDFWriteObject{13284}{<< /D [\MusicPDFPageRef{189} /XYZ 243.569 251.234 null] >>}
\MusicPDFWriteObject{13285}{<< /D [\MusicPDFPageRef{190} /XYZ 242.754 683.057 null] >>}
\MusicPDFWriteObject{13286}{<< /D [\MusicPDFPageRef{190} /XYZ 242.567 191.68 null] >>}
\MusicPDFWriteObject{13287}{<< /D [\MusicPDFPageRef{190} /XYZ 221.875 104.325 null] >>}
\MusicPDFWriteObject{13288}{<< /D [\MusicPDFPageRef{192} /XYZ 214.534 683.057 null] >>}
\MusicPDFWriteObject{13289}{<< /D [\MusicPDFPageRef{195} /XYZ 267.489 299.044 null] >>}
\MusicPDFWriteObject{13290}{<< /D [\MusicPDFPageRef{195} /XYZ 278.406 125.058 null] >>}
\MusicPDFWriteObject{13291}{<< /D [\MusicPDFPageRef{196} /XYZ 202.048 432.183 null] >>}
\MusicPDFWriteObject{13292}{<< /D [\MusicPDFPageRef{192} /XYZ 161.584 643.286 null] >>}
\MusicPDFWriteObject{13293}{<< /D [\MusicPDFPageRef{192} /XYZ 277.15 272.325 null] >>}
\MusicPDFWriteObject{13294}{<< /D [\MusicPDFPageRef{193} /XYZ 210.126 591.289 null] >>}
\MusicPDFWriteObject{13295}{<< /D [\MusicPDFPageRef{193} /XYZ 130.747 226.489 null] >>}
\MusicPDFWriteObject{13296}{<< /D [\MusicPDFPageRef{193} /XYZ 168.403 149.923 null] >>}
\MusicPDFWriteObject{13297}{<< /D [\MusicPDFPageRef{194} /XYZ 271.502 665.763 null] >>}
\MusicPDFWriteObject{13298}{<< /D [\MusicPDFPageRef{195} /XYZ 231.29 480.88 null] >>}
\MusicPDFWriteObject{13299}{<< /D [\MusicPDFPageRef{195} /XYZ 269.541 415.369 null] >>}
\MusicPDFWriteObject{13300}{<< /D [\MusicPDFPageRef{197} /XYZ 257.491 145.425 null] >>}
\MusicPDFWriteObject{13301}{<< /D [\MusicPDFPageRef{204} /XYZ 250.222 324.097 null] >>}
\MusicPDFWriteObject{13302}{<< /D [\MusicPDFPageRef{204} /XYZ 221.26 228.43 null] >>}
\MusicPDFWriteObject{13303}{<< /D [\MusicPDFPageRef{204} /XYZ 216.473 174.192 null] >>}
\MusicPDFWriteObject{13304}{<< /D [\MusicPDFPageRef{197} /XYZ 263.3 108.755 null] >>}
\MusicPDFWriteObject{13305}{<< /D [\MusicPDFPageRef{198} /XYZ 263.311 564.7 null] >>}
\MusicPDFWriteObject{13306}{<< /D [\MusicPDFPageRef{199} /XYZ 249.228 669.906 null] >>}
\MusicPDFWriteObject{13307}{<< /D [\MusicPDFPageRef{200} /XYZ 238.743 472.103 null] >>}
\MusicPDFWriteObject{13308}{<< /D [\MusicPDFPageRef{201} /XYZ 256.316 681.9 null] >>}
\MusicPDFWriteObject{13309}{<< /D [\MusicPDFPageRef{201} /XYZ 215.399 616.99 null] >>}
\MusicPDFWriteObject{13310}{<< /D [\MusicPDFPageRef{201} /XYZ 234.612 528.782 null] >>}
\MusicPDFWriteObject{13311}{<< /D [\MusicPDFPageRef{204} /XYZ 235.727 407.979 null] >>}
\MusicPDFWriteObject{13312}{<< /D [\MusicPDFPageRef{205} /XYZ 193.512 185.691 null] >>}
\MusicPDFWriteObject{13313}{<< /D [\MusicPDFPageRef{207} /XYZ 188.062 413.773 null] >>}
\MusicPDFWriteObject{13314}{<< /D [\MusicPDFPageRef{208} /XYZ 258.132 478.057 null] >>}
\MusicPDFWriteObject{13315}{<< /D [\MusicPDFPageRef{208} /XYZ 257.953 436.051 null] >>}
\MusicPDFWriteObject{13316}{<< /D [\MusicPDFPageRef{208} /XYZ 221.962 324.924 null] >>}
\MusicPDFWriteObject{13317}{<< /D [\MusicPDFPageRef{208} /XYZ 191.407 143.748 null] >>}
\MusicPDFWriteObject{13318}{<< /D [\MusicPDFPageRef{209} /XYZ 214.831 683.057 null] >>}
\MusicPDFWriteObject{13319}{<< /D [\MusicPDFPageRef{209} /XYZ 177.936 646.629 null] >>}
\MusicPDFWriteObject{13320}{<< /D [\MusicPDFPageRef{209} /XYZ 264.413 207.172 null] >>}
\MusicPDFWriteObject{13321}{<< /D [\MusicPDFPageRef{211} /XYZ 150.25 433.454 null] >>}
\MusicPDFWriteObject{13322}{<< /D [\MusicPDFPageRef{211} /XYZ 241.947 320.611 null] >>}
\MusicPDFWriteObject{13323}{<< /D [\MusicPDFPageRef{212} /XYZ 255.019 523.452 null] >>}
\MusicPDFWriteObject{13324}{<< /D [\MusicPDFPageRef{212} /XYZ 252.86 463.127 null] >>}
\MusicPDFWriteObject{13325}{<< /D [\MusicPDFPageRef{212} /XYZ 220.934 376.51 null] >>}
\MusicPDFWriteObject{13326}{<< /D [\MusicPDFPageRef{212} /XYZ 164.978 322.537 null] >>}
\MusicPDFWriteObject{13327}{<< /D [\MusicPDFPageRef{212} /XYZ 161.14 260.867 null] >>}
\MusicPDFWriteObject{13328}{<< /D [\MusicPDFPageRef{212} /XYZ 263.11 170.173 null] >>}
\MusicPDFWriteObject{13329}{<< /D [\MusicPDFPageRef{213} /XYZ 219.007 628.971 null] >>}
\MusicPDFWriteObject{13330}{<< /D [\MusicPDFPageRef{213} /XYZ 235.232 425.632 null] >>}
\MusicPDFWriteObject{13331}{<< /D [\MusicPDFPageRef{210} /XYZ 268.623 579.343 null] >>}
\MusicPDFWriteObject{13332}{<< /D [\MusicPDFPageRef{213} /XYZ 72 149.32 null] >>}
\MusicPDFWriteObject{13333}{<< /D [\MusicPDFPageRef{213} /XYZ 72 119.465 null] >>}
\MusicPDFWriteObject{13334}{<< /D [\MusicPDFPageRef{214} /XYZ 72 568.618 null] >>}
\MusicPDFWriteObject{13335}{<< /D [\MusicPDFPageRef{214} /XYZ 132.834 224.435 null] >>}
\MusicPDFWriteObject{13336}{<< /D [\MusicPDFPageRef{214} /XYZ 137.754 161.706 null] >>}
\MusicPDFWriteObject{13337}{<< /D [\MusicPDFPageRef{215} /XYZ 268.451 244.156 null] >>}
\MusicPDFWriteObject{13338}{<< /D [\MusicPDFPageRef{215} /XYZ 227.174 187.361 null] >>}
\MusicPDFWriteObject{13339}{<< /D [\MusicPDFPageRef{216} /XYZ 266.495 574.196 null] >>}
\MusicPDFWriteObject{13340}{<< /D [\MusicPDFPageRef{217} /XYZ 77.496 568.667 null] >>}
\MusicPDFWriteObject{13341}{<< /D [\MusicPDFPageRef{217} /XYZ 115.598 156.899 null] >>}
\MusicPDFWriteObject{13342}{<< /D [\MusicPDFPageRef{210} /XYZ 233.889 433.163 null] >>}
\MusicPDFWriteObject{13343}{<< /D [\MusicPDFPageRef{218} /XYZ 160.001 669.906 null] >>}
\MusicPDFWriteObject{13344}{<< /D [\MusicPDFPageRef{210} /XYZ 248.289 383.585 null] >>}
\MusicPDFWriteObject{13345}{<< /D [\MusicPDFPageRef{210} /XYZ 243.298 334.007 null] >>}
\MusicPDFWriteObject{13346}{<< /D [\MusicPDFPageRef{210} /XYZ 232.987 284.428 null] >>}
\MusicPDFWriteObject{13347}{<< /D [\MusicPDFPageRef{210} /XYZ 142.161 248.001 null] >>}
\MusicPDFWriteObject{13348}{<< /D [\MusicPDFPageRef{210} /XYZ 179.576 99.233 null] >>}
\MusicPDFWriteObject{13349}{<< /D [\MusicPDFPageRef{211} /XYZ 149.29 599.087 null] >>}
\MusicPDFWriteObject{13350}{<< /D [\MusicPDFPageRef{218} /XYZ 207.038 295.548 null] >>}
\MusicPDFWriteObject{13351}{<< /D [\MusicPDFPageRef{218} /XYZ 148.521 259.12 null] >>}
\MusicPDFWriteObject{13352}{<< /D [\MusicPDFPageRef{218} /XYZ 77.496 186.052 null] >>}
\MusicPDFWriteObject{13353}{<< /D [\MusicPDFPageRef{219} /XYZ 137.234 669.906 null] >>}
\MusicPDFWriteObject{13354}{<< /D [\MusicPDFPageRef{219} /XYZ 141.379 620.327 null] >>}
\MusicPDFWriteObject{13355}{<< /D [\MusicPDFPageRef{219} /XYZ 85.97 520.781 null] >>}
\MusicPDFWriteObject{13356}{<< /D [\MusicPDFPageRef{219} /XYZ 193.209 180.106 null] >>}
\MusicPDFWriteObject{13357}{<< /D [\MusicPDFPageRef{220} /XYZ 159.512 682.515 null] >>}
\MusicPDFWriteObject{13358}{<< /D [\MusicPDFPageRef{223} /XYZ 253.531 555.927 null] >>}
\MusicPDFWriteObject{13359}{<< /D [\MusicPDFPageRef{223} /XYZ 245.018 440.595 null] >>}
\MusicPDFWriteObject{13360}{<< /D [\MusicPDFPageRef{223} /XYZ 135.692 351.564 null] >>}
\MusicPDFWriteObject{13361}{<< /D [\MusicPDFPageRef{223} /XYZ 143.252 301.986 null] >>}
\MusicPDFWriteObject{13362}{<< /D [\MusicPDFPageRef{223} /XYZ 114.795 226.106 null] >>}
\MusicPDFWriteObject{13363}{<< /D [\MusicPDFPageRef{224} /XYZ 268.049 279.774 null] >>}
\MusicPDFWriteObject{13364}{<< /D [\MusicPDFPageRef{224} /XYZ 203.666 243.347 null] >>}
\MusicPDFWriteObject{13365}{<< /D [\MusicPDFPageRef{225} /XYZ 262.129 364.584 null] >>}
\MusicPDFWriteObject{13366}{<< /D [\MusicPDFPageRef{225} /XYZ 191.74 328.156 null] >>}
\MusicPDFWriteObject{13367}{<< /D [\MusicPDFPageRef{227} /XYZ 251.275 245.93 null] >>}
\MusicPDFWriteObject{13368}{<< /D [\MusicPDFPageRef{227} /XYZ 242.825 130.597 null] >>}
\MusicPDFWriteObject{13369}{<< /D [\MusicPDFPageRef{228} /XYZ 72 601.118 null] >>}
\MusicPDFWriteObject{13370}{<< /D [\MusicPDFPageRef{228} /XYZ 254.38 378.855 null] >>}
\MusicPDFWriteObject{13371}{<< /D [\MusicPDFPageRef{228} /XYZ 190.305 299.079 null] >>}
\MusicPDFWriteObject{13372}{<< /D [\MusicPDFPageRef{228} /XYZ 165.758 253.168 null] >>}
\MusicPDFWriteObject{13373}{<< /D [\MusicPDFPageRef{229} /XYZ 159.641 586.745 null] >>}
\MusicPDFWriteObject{13374}{<< /D [\MusicPDFPageRef{230} /XYZ 185.383 352.362 null] >>}
\MusicPDFWriteObject{13375}{<< /D [\MusicPDFPageRef{230} /XYZ 189.823 289.633 null] >>}
\MusicPDFWriteObject{13376}{<< /D [\MusicPDFPageRef{230} /XYZ 199.39 223.773 null] >>}
\MusicPDFWriteObject{13377}{<< /D [\MusicPDFPageRef{231} /XYZ 227.277 450.718 null] >>}
\MusicPDFWriteObject{13378}{<< /D [\MusicPDFPageRef{232} /XYZ 186.017 584.378 null] >>}
\MusicPDFWriteObject{13379}{<< /D [\MusicPDFPageRef{232} /XYZ 144.426 466.511 null] >>}
\MusicPDFWriteObject{13380}{<< /D [\MusicPDFPageRef{232} /XYZ 223.264 248.67 null] >>}
\MusicPDFWriteObject{13381}{<< /D [\MusicPDFPageRef{232} /XYZ 211.494 196.579 null] >>}
\MusicPDFWriteObject{13382}{<< /D [\MusicPDFPageRef{233} /XYZ 193.864 575.45 null] >>}
\MusicPDFWriteObject{13383}{<< /D [\MusicPDFPageRef{236} /XYZ 72 599.142 null] >>}
\MusicPDFWriteObject{13384}{<< /D [\MusicPDFPageRef{236} /XYZ 172.612 562.715 null] >>}
\MusicPDFWriteObject{13385}{<< /D [\MusicPDFPageRef{236} /XYZ 154.902 526.287 null] >>}
\MusicPDFWriteObject{13386}{<< /D [\MusicPDFPageRef{236} /XYZ 72 480.975 null] >>}
\MusicPDFWriteObject{13387}{<< /D [\MusicPDFPageRef{236} /XYZ 142.117 417.389 null] >>}
\MusicPDFWriteObject{13388}{<< /D [\MusicPDFPageRef{236} /XYZ 198.962 348.021 null] >>}
\MusicPDFWriteObject{13389}{<< /D [\MusicPDFPageRef{233} /XYZ 198.042 486.419 null] >>}
\MusicPDFWriteObject{13390}{<< /D [\MusicPDFPageRef{233} /XYZ 211.642 354.805 null] >>}
\MusicPDFWriteObject{13391}{<< /D [\MusicPDFPageRef{233} /XYZ 206.132 161.368 null] >>}
\MusicPDFWriteObject{13392}{<< /D [\MusicPDFPageRef{235} /XYZ 192.808 370.275 null] >>}
\MusicPDFWriteObject{13393}{<< /D [\MusicPDFPageRef{235} /XYZ 174.103 333.848 null] >>}
\MusicPDFWriteObject{13394}{<< /D [\MusicPDFPageRef{235} /XYZ 119.768 243.841 null] >>}
\MusicPDFWriteObject{13395}{<< /D [\MusicPDFPageRef{235} /XYZ 137.089 207.352 null] >>}
\MusicPDFWriteObject{13396}{<< /D [\MusicPDFPageRef{236} /XYZ 152.312 650.767 null] >>}
\MusicPDFWriteObject{13397}{<< /D [\MusicPDFPageRef{238} /XYZ 192.814 99.532 null] >>}
\MusicPDFWriteObject{13398}{<< /D [\MusicPDFPageRef{239} /XYZ 192.602 665.155 null] >>}
\MusicPDFWriteObject{13399}{<< /D [\MusicPDFPageRef{239} /XYZ 208.286 554.195 null] >>}
\MusicPDFWriteObject{13400}{<< /D [\MusicPDFPageRef{239} /XYZ 198.962 500.071 null] >>}
\MusicPDFWriteObject{13401}{<< /D [\MusicPDFPageRef{240} /XYZ 266.047 552.94 null] >>}
\MusicPDFWriteObject{13402}{<< /D [\MusicPDFPageRef{242} /XYZ 199.82 397.594 null] >>}
\MusicPDFWriteObject{13403}{<< /D [\MusicPDFPageRef{240} /XYZ 266.047 287.63 null] >>}
\MusicPDFWriteObject{13404}{<< /D [\MusicPDFPageRef{240} /XYZ 243.348 224.89 null] >>}
\MusicPDFWriteObject{13405}{<< /D [\MusicPDFPageRef{240} /XYZ 235.845 187.281 null] >>}
\MusicPDFWriteObject{13406}{<< /D [\MusicPDFPageRef{241} /XYZ 262.049 556.322 null] >>}
\MusicPDFWriteObject{13407}{<< /D [\MusicPDFPageRef{241} /XYZ 245.245 99.233 null] >>}
\MusicPDFWriteObject{13408}{<< /D [\MusicPDFPageRef{242} /XYZ 231.233 682.813 null] >>}
\MusicPDFWriteObject{13409}{<< /D [\MusicPDFPageRef{242} /XYZ 255.632 596.691 null] >>}
\MusicPDFWriteObject{13410}{<< /D [\MusicPDFPageRef{242} /XYZ 183.494 434.021 null] >>}
\MusicPDFWriteObject{13411}{<< /D [\MusicPDFPageRef{243} /XYZ 268.825 437.043 null] >>}
\MusicPDFWriteObject{13412}{<< /D [\MusicPDFPageRef{246} /XYZ 72 572.589 null] >>}
\MusicPDFWriteObject{13413}{<< /D [\MusicPDFPageRef{246} /XYZ 245.969 352.619 null] >>}
\MusicPDFWriteObject{13414}{<< /D [\MusicPDFPageRef{246} /XYZ 182.136 236.931 null] >>}
\MusicPDFWriteObject{13415}{<< /D [\MusicPDFPageRef{246} /XYZ 199.053 134.461 null] >>}
\MusicPDFWriteObject{13416}{<< /D [\MusicPDFPageRef{247} /XYZ 213.08 459.335 null] >>}
\MusicPDFWriteObject{13417}{<< /D [\MusicPDFPageRef{247} /XYZ 215.724 316.323 null] >>}
\MusicPDFWriteObject{13418}{<< /D [\MusicPDFPageRef{248} /XYZ 217.164 679.571 null] >>}
\MusicPDFWriteObject{13419}{<< /D [\MusicPDFPageRef{248} /XYZ 72 516.349 null] >>}
\MusicPDFWriteObject{13420}{<< /D [\MusicPDFPageRef{248} /XYZ 72 486.494 null] >>}
\MusicPDFWriteObject{13421}{<< /D [\MusicPDFPageRef{248} /XYZ 72 456.64 null] >>}
\MusicPDFWriteObject{13422}{<< /D [\MusicPDFPageRef{243} /XYZ 215.864 322.338 null] >>}
\MusicPDFWriteObject{13423}{<< /D [\MusicPDFPageRef{243} /XYZ 248.045 223.232 null] >>}
\MusicPDFWriteObject{13424}{<< /D [\MusicPDFPageRef{245} /XYZ 167.52 401.069 null] >>}
\MusicPDFWriteObject{13425}{<< /D [\MusicPDFPageRef{245} /XYZ 146.327 338.893 null] >>}
\MusicPDFWriteObject{13426}{<< /D [\MusicPDFPageRef{245} /XYZ 228.413 288.181 null] >>}
\MusicPDFWriteObject{13427}{<< /D [\MusicPDFPageRef{245} /XYZ 218.819 150.801 null] >>}
\MusicPDFWriteObject{13428}{<< /D [\MusicPDFPageRef{246} /XYZ 72 607.44 null] >>}
\MusicPDFWriteObject{13429}{<< /D [\MusicPDFPageRef{246} /XYZ 72 590.265 null] >>}
\MusicPDFWriteObject{13430}{<< /D [\MusicPDFPageRef{251} /XYZ 264.349 325.567 null] >>}
\MusicPDFWriteObject{13431}{<< /D [\MusicPDFPageRef{252} /XYZ 245.74 156.899 null] >>}
\MusicPDFWriteObject{13432}{<< /D [\MusicPDFPageRef{253} /XYZ 162.23 568.957 null] >>}
\MusicPDFWriteObject{13433}{<< /D [\MusicPDFPageRef{253} /XYZ 206.767 403.246 null] >>}
\MusicPDFWriteObject{13434}{<< /D [\MusicPDFPageRef{253} /XYZ 166.639 232.779 null] >>}
\MusicPDFWriteObject{13435}{<< /D [\MusicPDFPageRef{253} /XYZ 179.021 196.351 null] >>}
\MusicPDFWriteObject{13436}{<< /D [\MusicPDFPageRef{254} /XYZ 158.551 639.71 null] >>}
\MusicPDFWriteObject{13437}{<< /D [\MusicPDFPageRef{254} /XYZ 244.953 506.452 null] >>}
\MusicPDFWriteObject{13438}{<< /D [\MusicPDFPageRef{254} /XYZ 217.568 351.668 null] >>}
\MusicPDFWriteObject{13439}{<< /D [\MusicPDFPageRef{255} /XYZ 185.449 455.809 null] >>}
\MusicPDFWriteObject{13440}{<< /D [\MusicPDFPageRef{255} /XYZ 223.867 117.447 null] >>}
\MusicPDFWriteObject{13441}{<< /D [\MusicPDFPageRef{251} /XYZ 238.184 267.567 null] >>}
\MusicPDFWriteObject{13442}{<< /D [\MusicPDFPageRef{256} /XYZ 217.708 590.754 null] >>}
\MusicPDFWriteObject{13443}{<< /D [\MusicPDFPageRef{256} /XYZ 160.109 339.031 null] >>}
\MusicPDFWriteObject{13444}{<< /D [\MusicPDFPageRef{256} /XYZ 160.434 297.462 null] >>}
\MusicPDFWriteObject{13445}{<< /D [\MusicPDFPageRef{257} /XYZ 200.39 597.953 null] >>}
\MusicPDFWriteObject{13446}{<< /D [\MusicPDFPageRef{257} /XYZ 168.602 112.085 null] >>}
\MusicPDFWriteObject{13447}{<< /D [\MusicPDFPageRef{258} /XYZ 217.061 671.456 null] >>}
\MusicPDFWriteObject{13448}{<< /D [\MusicPDFPageRef{258} /XYZ 72 583.505 null] >>}
\MusicPDFWriteObject{13449}{<< /D [\MusicPDFPageRef{258} /XYZ 213.811 527.328 null] >>}
\MusicPDFWriteObject{13450}{<< /D [\MusicPDFPageRef{258} /XYZ 218.481 428.797 null] >>}
\MusicPDFWriteObject{13451}{<< /D [\MusicPDFPageRef{258} /XYZ 157.358 392.209 null] >>}
\MusicPDFWriteObject{13452}{<< /D [\MusicPDFPageRef{251} /XYZ 236.378 209.04 null] >>}
\MusicPDFWriteObject{13453}{<< /D [\MusicPDFPageRef{258} /XYZ 194.314 355.782 null] >>}
\MusicPDFWriteObject{13454}{<< /D [\MusicPDFPageRef{258} /XYZ 146.534 319.354 null] >>}
\MusicPDFWriteObject{13455}{<< /D [\MusicPDFPageRef{258} /XYZ 72 221.245 null] >>}
\MusicPDFWriteObject{13456}{<< /D [\MusicPDFPageRef{258} /XYZ 72 193.425 null] >>}
\MusicPDFWriteObject{13457}{<< /D [\MusicPDFPageRef{258} /XYZ 72 170.473 null] >>}
\MusicPDFWriteObject{13458}{<< /D [\MusicPDFPageRef{258} /XYZ 158.741 130.597 null] >>}
\MusicPDFWriteObject{13459}{<< /D [\MusicPDFPageRef{259} /XYZ 72 682.35 null] >>}
\MusicPDFWriteObject{13460}{<< /D [\MusicPDFPageRef{259} /XYZ 72 654.53 null] >>}
\MusicPDFWriteObject{13461}{<< /D [\MusicPDFPageRef{259} /XYZ 72 626.71 null] >>}
\MusicPDFWriteObject{13462}{<< /D [\MusicPDFPageRef{259} /XYZ 72 605.39 null] >>}
\MusicPDFWriteObject{13463}{<< /D [\MusicPDFPageRef{251} /XYZ 260.897 134.208 null] >>}
\MusicPDFWriteObject{13464}{<< /D [\MusicPDFPageRef{259} /XYZ 146.534 565.675 null] >>}
\MusicPDFWriteObject{13465}{<< /D [\MusicPDFPageRef{252} /XYZ 173.538 432.192 null] >>}
\MusicPDFWriteObject{13466}{<< /D [\MusicPDFPageRef{252} /XYZ 259.06 386.723 null] >>}
\MusicPDFWriteObject{13467}{<< /D [\MusicPDFPageRef{252} /XYZ 241.54 337.144 null] >>}
\MusicPDFWriteObject{13468}{<< /D [\MusicPDFPageRef{252} /XYZ 184.503 300.24 null] >>}
\MusicPDFWriteObject{13469}{<< /D [\MusicPDFPageRef{252} /XYZ 198.1 219.628 null] >>}
\MusicPDFWriteObject{13470}{<< /D [\MusicPDFPageRef{259} /XYZ 194.088 161.517 null] >>}
\MusicPDFWriteObject{13471}{<< /D [\MusicPDFPageRef{262} /XYZ 218.8 112.294 null] >>}
\MusicPDFWriteObject{13472}{<< /D [\MusicPDFPageRef{263} /XYZ 227.164 212.277 null] >>}
\MusicPDFWriteObject{13473}{<< /D [\MusicPDFPageRef{264} /XYZ 234.014 318.279 null] >>}
\MusicPDFWriteObject{13474}{<< /D [\MusicPDFPageRef{264} /XYZ 260.98 260.737 null] >>}
\MusicPDFWriteObject{13475}{<< /D [\MusicPDFPageRef{264} /XYZ 240.729 187.535 null] >>}
\MusicPDFWriteObject{13476}{<< /D [\MusicPDFPageRef{264} /XYZ 154.633 110.175 null] >>}
\MusicPDFWriteObject{13477}{<< /D [\MusicPDFPageRef{261} /XYZ 276.703 268.412 null] >>}
\MusicPDFWriteObject{13478}{<< /D [\MusicPDFPageRef{261} /XYZ 141.769 204.321 null] >>}
\MusicPDFWriteObject{13479}{<< /D [\MusicPDFPageRef{261} /XYZ 176.673 158.259 null] >>}
\MusicPDFWriteObject{13480}{<< /D [\MusicPDFPageRef{262} /XYZ 141.582 681.398 null] >>}
\MusicPDFWriteObject{13481}{<< /D [\MusicPDFPageRef{262} /XYZ 143.756 640.788 null] >>}
\MusicPDFWriteObject{13482}{<< /D [\MusicPDFPageRef{262} /XYZ 126.577 539.856 null] >>}
\MusicPDFWriteObject{13483}{<< /D [\MusicPDFPageRef{262} /XYZ 200.711 439.31 null] >>}
\MusicPDFWriteObject{13484}{<< /D [\MusicPDFPageRef{262} /XYZ 72 219.245 null] >>}
\MusicPDFWriteObject{13485}{<< /D [\MusicPDFPageRef{265} /XYZ 214.143 295.575 null] >>}
\MusicPDFWriteObject{13486}{<< /D [\MusicPDFPageRef{270} /XYZ 166.178 674.926 null] >>}
\MusicPDFWriteObject{13487}{<< /D [\MusicPDFPageRef{270} /XYZ 147.707 582.839 null] >>}
\MusicPDFWriteObject{13488}{<< /D [\MusicPDFPageRef{270} /XYZ 225.357 523.893 null] >>}
\MusicPDFWriteObject{13489}{<< /D [\MusicPDFPageRef{270} /XYZ 151.142 277.038 null] >>}
\MusicPDFWriteObject{13490}{<< /D [\MusicPDFPageRef{270} /XYZ 196.334 119.527 null] >>}
\MusicPDFWriteObject{13491}{<< /D [\MusicPDFPageRef{271} /XYZ 163.255 613.489 null] >>}
\MusicPDFWriteObject{13492}{<< /D [\MusicPDFPageRef{272} /XYZ 237.57 683.057 null] >>}
\MusicPDFWriteObject{13493}{<< /D [\MusicPDFPageRef{272} /XYZ 243.933 486.754 null] >>}
\MusicPDFWriteObject{13494}{<< /D [\MusicPDFPageRef{272} /XYZ 262.079 384.303 null] >>}
\MusicPDFWriteObject{13495}{<< /D [\MusicPDFPageRef{272} /XYZ 189.705 334.724 null] >>}
\MusicPDFWriteObject{13496}{<< /D [\MusicPDFPageRef{265} /XYZ 195.587 99.233 null] >>}
\MusicPDFWriteObject{13497}{<< /D [\MusicPDFPageRef{272} /XYZ 244.9 294.456 null] >>}
\MusicPDFWriteObject{13498}{<< /D [\MusicPDFPageRef{273} /XYZ 197.664 525.754 null] >>}
\MusicPDFWriteObject{13499}{<< /D [\MusicPDFPageRef{273} /XYZ 259.098 305.612 null] >>}
\MusicPDFWriteObject{13500}{<< /D [\MusicPDFPageRef{274} /XYZ 165.547 668.749 null] >>}
\MusicPDFWriteObject{13501}{<< /D [\MusicPDFPageRef{266} /XYZ 146.335 630.453 null] >>}
\MusicPDFWriteObject{13502}{<< /D [\MusicPDFPageRef{266} /XYZ 251.781 205.24 null] >>}
\MusicPDFWriteObject{13503}{<< /D [\MusicPDFPageRef{266} /XYZ 140.259 125.919 null] >>}
\MusicPDFWriteObject{13504}{<< /D [\MusicPDFPageRef{267} /XYZ 152.033 666.465 null] >>}
\MusicPDFWriteObject{13505}{<< /D [\MusicPDFPageRef{267} /XYZ 191.366 415.224 null] >>}
\MusicPDFWriteObject{13506}{<< /D [\MusicPDFPageRef{267} /XYZ 143.3 325.673 null] >>}
\MusicPDFWriteObject{13507}{<< /D [\MusicPDFPageRef{268} /XYZ 259.597 319.885 null] >>}
\MusicPDFWriteObject{13508}{<< /D [\MusicPDFPageRef{274} /XYZ 215.401 292.149 null] >>}
\MusicPDFWriteObject{13509}{<< /D [\MusicPDFPageRef{274} /XYZ 253.692 183.964 null] >>}
\MusicPDFWriteObject{13510}{<< /D [\MusicPDFPageRef{275} /XYZ 213.2 557.553 null] >>}
\MusicPDFWriteObject{13511}{<< /D [\MusicPDFPageRef{275} /XYZ 203.108 194.385 null] >>}
\MusicPDFWriteObject{13512}{<< /D [\MusicPDFPageRef{275} /XYZ 163.863 151.13 null] >>}
\MusicPDFWriteObject{13513}{<< /D [\MusicPDFPageRef{276} /XYZ 119.905 318.571 null] >>}
\MusicPDFWriteObject{13514}{<< /D [\MusicPDFPageRef{277} /XYZ 151.834 447.17 null] >>}
\MusicPDFWriteObject{13515}{<< /D [\MusicPDFPageRef{277} /XYZ 183.385 180.673 null] >>}
\MusicPDFWriteObject{13516}{<< /D [\MusicPDFPageRef{278} /XYZ 229.68 131.112 null] >>}
\MusicPDFWriteObject{13517}{<< /D [\MusicPDFPageRef{279} /XYZ 243.384 669.906 null] >>}
\MusicPDFWriteObject{13518}{<< /D [\MusicPDFPageRef{279} /XYZ 243.398 632.981 null] >>}
\MusicPDFWriteObject{13519}{<< /D [\MusicPDFPageRef{279} /XYZ 216.526 554.973 null] >>}
\MusicPDFWriteObject{13520}{<< /D [\MusicPDFPageRef{279} /XYZ 264.331 326.039 null] >>}
\MusicPDFWriteObject{13521}{<< /D [\MusicPDFPageRef{279} /XYZ 257.527 280.373 null] >>}
\MusicPDFWriteObject{13522}{<< /D [\MusicPDFPageRef{279} /XYZ 267.752 204.493 null] >>}
\MusicPDFWriteObject{13523}{<< /D [\MusicPDFPageRef{279} /XYZ 162.762 168.066 null] >>}
\MusicPDFWriteObject{13524}{<< /D [\MusicPDFPageRef{279} /XYZ 229.091 117.447 null] >>}
\MusicPDFWriteObject{13525}{<< /D [\MusicPDFPageRef{277} /XYZ 267.447 143.748 null] >>}
\MusicPDFWriteObject{13526}{<< /D [\MusicPDFPageRef{278} /XYZ 198.994 683.057 null] >>}
\MusicPDFWriteObject{13527}{<< /D [\MusicPDFPageRef{278} /XYZ 270.838 646.132 null] >>}
\MusicPDFWriteObject{13528}{<< /D [\MusicPDFPageRef{278} /XYZ 267.291 531.078 null] >>}
\MusicPDFWriteObject{13529}{<< /D [\MusicPDFPageRef{278} /XYZ 254.136 485.886 null] >>}
\MusicPDFWriteObject{13530}{<< /D [\MusicPDFPageRef{278} /XYZ 264.362 408.578 null] >>}
\MusicPDFWriteObject{13531}{<< /D [\MusicPDFPageRef{278} /XYZ 141.378 372.151 null] >>}
\MusicPDFWriteObject{13532}{<< /D [\MusicPDFPageRef{278} /XYZ 219.34 319.121 null] >>}
\MusicPDFWriteObject{13533}{<< /D [\MusicPDFPageRef{280} /XYZ 247.831 403.303 null] >>}
\MusicPDFWriteObject{13534}{<< /D [\MusicPDFPageRef{280} /XYZ 233.251 280.664 null] >>}
\MusicPDFWriteObject{13535}{<< /D [\MusicPDFPageRef{280} /XYZ 243.861 210.738 null] >>}
\MusicPDFWriteObject{13536}{<< /D [\MusicPDFPageRef{282} /XYZ 246.275 470.79 null] >>}
\MusicPDFWriteObject{13537}{<< /D [\MusicPDFPageRef{282} /XYZ 235.802 140.308 null] >>}
\MusicPDFWriteObject{13538}{<< /D [\MusicPDFPageRef{283} /XYZ 220.538 682.813 null] >>}
\MusicPDFWriteObject{13539}{<< /D [\MusicPDFPageRef{283} /XYZ 268.755 518.301 null] >>}
\MusicPDFWriteObject{13540}{<< /D [\MusicPDFPageRef{283} /XYZ 217.093 481.873 null] >>}
\MusicPDFWriteObject{13541}{<< /D [\MusicPDFPageRef{283} /XYZ 231.876 116.947 null] >>}
\MusicPDFWriteObject{13542}{<< /D [\MusicPDFPageRef{287} /XYZ 252.27 225.792 null] >>}
\MusicPDFWriteObject{13543}{<< /D [\MusicPDFPageRef{287} /XYZ 263.181 127.099 null] >>}
\MusicPDFWriteObject{13544}{<< /D [\MusicPDFPageRef{288} /XYZ 183.907 514.74 null] >>}
\MusicPDFWriteObject{13545}{<< /D [\MusicPDFPageRef{290} /XYZ 229.338 181.107 null] >>}
\MusicPDFWriteObject{13546}{<< /D [\MusicPDFPageRef{290} /XYZ 239.163 108.98 null] >>}
\MusicPDFWriteObject{13547}{<< /D [\MusicPDFPageRef{291} /XYZ 190.47 550.592 null] >>}
\MusicPDFWriteObject{13548}{<< /D [\MusicPDFPageRef{292} /XYZ 262.758 683.057 null] >>}
\MusicPDFWriteObject{13549}{<< /D [\MusicPDFPageRef{292} /XYZ 241.154 642.732 null] >>}
\MusicPDFWriteObject{13550}{<< /D [\MusicPDFPageRef{292} /XYZ 224.25 585.407 null] >>}
\MusicPDFWriteObject{13551}{<< /D [\MusicPDFPageRef{292} /XYZ 170.271 519.739 null] >>}
\MusicPDFWriteObject{13552}{<< /D [\MusicPDFPageRef{292} /XYZ 176.368 204.794 null] >>}
\MusicPDFWriteObject{13553}{<< /D [\MusicPDFPageRef{293} /XYZ 195.76 287.216 null] >>}
\MusicPDFWriteObject{13554}{<< /D [\MusicPDFPageRef{294} /XYZ 185.999 669.308 null] >>}
\MusicPDFWriteObject{13555}{<< /D [\MusicPDFPageRef{288} /XYZ 198.408 426.623 null] >>}
\MusicPDFWriteObject{13556}{<< /D [\MusicPDFPageRef{294} /XYZ 158.179 632.084 null] >>}
\MusicPDFWriteObject{13557}{<< /D [\MusicPDFPageRef{294} /XYZ 186.838 412.762 null] >>}
\MusicPDFWriteObject{13558}{<< /D [\MusicPDFPageRef{294} /XYZ 156.107 334.612 null] >>}
\MusicPDFWriteObject{13559}{<< /D [\MusicPDFPageRef{288} /XYZ 275.159 246.527 null] >>}
\MusicPDFWriteObject{13560}{<< /D [\MusicPDFPageRef{288} /XYZ 259.461 117.447 null] >>}
\MusicPDFWriteObject{13561}{<< /D [\MusicPDFPageRef{289} /XYZ 184.077 656.755 null] >>}
\MusicPDFWriteObject{13562}{<< /D [\MusicPDFPageRef{289} /XYZ 262.58 593.783 null] >>}
\MusicPDFWriteObject{13563}{<< /D [\MusicPDFPageRef{290} /XYZ 232.688 584.051 null] >>}
\MusicPDFWriteObject{13564}{<< /D [\MusicPDFPageRef{290} /XYZ 262.199 297.913 null] >>}
\MusicPDFWriteObject{13565}{<< /D [\MusicPDFPageRef{290} /XYZ 190.145 248.335 null] >>}
\MusicPDFWriteObject{13566}{<< /D [\MusicPDFPageRef{297} /XYZ 258.576 433.97 null] >>}
\MusicPDFWriteObject{13567}{<< /D [\MusicPDFPageRef{297} /XYZ 171.255 212.588 null] >>}
\MusicPDFWriteObject{13568}{<< /D [\MusicPDFPageRef{298} /XYZ 195.875 227 null] >>}
\MusicPDFWriteObject{13569}{<< /D [\MusicPDFPageRef{301} /XYZ 218.131 628.355 null] >>}
\MusicPDFWriteObject{13570}{<< /D [\MusicPDFPageRef{301} /XYZ 250.478 568.532 null] >>}
\MusicPDFWriteObject{13571}{<< /D [\MusicPDFPageRef{302} /XYZ 218.243 383.504 null] >>}
\MusicPDFWriteObject{13572}{<< /D [\MusicPDFPageRef{298} /XYZ 161.58 150.204 null] >>}
\MusicPDFWriteObject{13573}{<< /D [\MusicPDFPageRef{298} /XYZ 146.602 109.52 null] >>}
\MusicPDFWriteObject{13574}{<< /D [\MusicPDFPageRef{299} /XYZ 211.294 678.8 null] >>}
\MusicPDFWriteObject{13575}{<< /D [\MusicPDFPageRef{299} /XYZ 180.813 581.043 null] >>}
\MusicPDFWriteObject{13576}{<< /D [\MusicPDFPageRef{300} /XYZ 224.838 564.403 null] >>}
\MusicPDFWriteObject{13577}{<< /D [\MusicPDFPageRef{300} /XYZ 72 415.149 null] >>}
\MusicPDFWriteObject{13578}{<< /D [\MusicPDFPageRef{300} /XYZ 72 367.257 null] >>}
\MusicPDFWriteObject{13579}{<< /D [\MusicPDFPageRef{300} /XYZ 72 210.539 null] >>}
\MusicPDFWriteObject{13580}{<< /D [\MusicPDFPageRef{302} /XYZ 278.572 198.559 null] >>}
\MusicPDFWriteObject{13581}{<< /D [\MusicPDFPageRef{303} /XYZ 243.28 643.604 null] >>}
\MusicPDFWriteObject{13582}{<< /D [\MusicPDFPageRef{303} /XYZ 278.352 605.673 null] >>}
\MusicPDFWriteObject{13583}{<< /D [\MusicPDFPageRef{303} /XYZ 230.933 405.788 null] >>}
\MusicPDFWriteObject{13584}{<< /D [\MusicPDFPageRef{303} /XYZ 246.958 259.125 null] >>}
\MusicPDFWriteObject{13585}{<< /D [\MusicPDFPageRef{303} /XYZ 230.797 174.499 null] >>}
\MusicPDFWriteObject{13586}{<< /D [\MusicPDFPageRef{303} /XYZ 247.831 130.918 null] >>}
\MusicPDFWriteObject{13587}{<< /D [\MusicPDFPageRef{305} /XYZ 167.108 439.94 null] >>}
\MusicPDFWriteObject{13588}{<< /D [\MusicPDFPageRef{305} /XYZ 117.578 314.308 null] >>}
\MusicPDFWriteObject{13589}{<< /D [\MusicPDFPageRef{305} /XYZ 164.148 99.785 null] >>}
\MusicPDFWriteObject{13590}{<< /D [\MusicPDFPageRef{306} /XYZ 172.541 625.401 null] >>}
\MusicPDFWriteObject{13591}{<< /D [\MusicPDFPageRef{307} /XYZ 205.113 546.706 null] >>}
\MusicPDFWriteObject{13592}{<< /D [\MusicPDFPageRef{307} /XYZ 221.038 506.284 null] >>}
\MusicPDFWriteObject{13593}{<< /D [\MusicPDFPageRef{309} /XYZ 72 492.552 null] >>}
\MusicPDFWriteObject{13594}{<< /D [\MusicPDFPageRef{309} /XYZ 72 455.764 null] >>}
\MusicPDFWriteObject{13595}{<< /D [\MusicPDFPageRef{309} /XYZ 72 429.259 null] >>}
\MusicPDFWriteObject{13596}{<< /D [\MusicPDFPageRef{310} /XYZ 219.717 683.057 null] >>}
\MusicPDFWriteObject{13597}{<< /D [\MusicPDFPageRef{304} /XYZ 233.495 629.612 null] >>}
\MusicPDFWriteObject{13598}{<< /D [\MusicPDFPageRef{310} /XYZ 235.423 619.443 null] >>}
\MusicPDFWriteObject{13599}{<< /D [\MusicPDFPageRef{310} /XYZ 163.456 555.923 null] >>}
\MusicPDFWriteObject{13600}{<< /D [\MusicPDFPageRef{310} /XYZ 136.481 446.771 null] >>}
\MusicPDFWriteObject{13601}{<< /D [\MusicPDFPageRef{310} /XYZ 72.698 361.659 null] >>}
\MusicPDFWriteObject{13602}{<< /D [\MusicPDFPageRef{310} /XYZ 72 302.894 null] >>}
\MusicPDFWriteObject{13603}{<< /D [\MusicPDFPageRef{310} /XYZ 72 272.817 null] >>}
\MusicPDFWriteObject{13604}{<< /D [\MusicPDFPageRef{310} /XYZ 72 245.491 null] >>}
\MusicPDFWriteObject{13605}{<< /D [\MusicPDFPageRef{304} /XYZ 232.382 472.845 null] >>}
\MusicPDFWriteObject{13606}{<< /D [\MusicPDFPageRef{304} /XYZ 246.891 423.266 null] >>}
\MusicPDFWriteObject{13607}{<< /D [\MusicPDFPageRef{304} /XYZ 169.477 354.43 null] >>}
\MusicPDFWriteObject{13608}{<< /D [\MusicPDFPageRef{304} /XYZ 220.77 245.915 null] >>}
\MusicPDFWriteObject{13609}{<< /D [\MusicPDFPageRef{304} /XYZ 183.908 193.487 null] >>}
\MusicPDFWriteObject{13610}{<< /D [\MusicPDFPageRef{304} /XYZ 116.726 143.748 null] >>}
\MusicPDFWriteObject{13611}{<< /D [\MusicPDFPageRef{305} /XYZ 189.218 621.519 null] >>}
\MusicPDFWriteObject{13612}{<< /D [\MusicPDFPageRef{312} /XYZ 270.042 127.692 null] >>}
\MusicPDFWriteObject{13613}{<< /D [\MusicPDFPageRef{313} /XYZ 264.694 590.758 null] >>}
\MusicPDFWriteObject{13614}{<< /D [\MusicPDFPageRef{313} /XYZ 241.89 177.134 null] >>}
\MusicPDFWriteObject{13615}{<< /D [\MusicPDFPageRef{313} /XYZ 240.462 105.437 null] >>}
\MusicPDFWriteObject{13616}{<< /D [\MusicPDFPageRef{314} /XYZ 214.737 395.107 null] >>}
\MusicPDFWriteObject{13617}{<< /D [\MusicPDFPageRef{314} /XYZ 251.698 166.034 null] >>}
\MusicPDFWriteObject{13618}{<< /D [\MusicPDFPageRef{314} /XYZ 246.416 99.233 null] >>}
\MusicPDFWriteObject{13619}{<< /D [\MusicPDFPageRef{315} /XYZ 253.889 523.209 null] >>}
\MusicPDFWriteObject{13620}{<< /D [\MusicPDFPageRef{316} /XYZ 242.041 429.391 null] >>}
\MusicPDFWriteObject{13621}{<< /D [\MusicPDFPageRef{316} /XYZ 290.441 379.812 null] >>}
\MusicPDFWriteObject{13622}{<< /D [\MusicPDFPageRef{318} /XYZ 72 690.398 null] >>}
\MusicPDFWriteObject{13623}{<< /D [\MusicPDFPageRef{320} /XYZ 212.894 681.691 null] >>}
\MusicPDFWriteObject{13624}{<< /D [\MusicPDFPageRef{320} /XYZ 83.741 599.521 null] >>}
\MusicPDFWriteObject{13625}{<< /D [\MusicPDFPageRef{323} /XYZ 225.509 143.543 null] >>}
\MusicPDFWriteObject{13626}{<< /D [\MusicPDFPageRef{325} /XYZ 162.544 148.994 null] >>}
\MusicPDFWriteObject{13627}{<< /D [\MusicPDFPageRef{326} /XYZ 230.147 363.478 null] >>}
\MusicPDFWriteObject{13628}{<< /D [\MusicPDFPageRef{327} /XYZ 265.301 206.477 null] >>}
\MusicPDFWriteObject{13629}{<< /D [\MusicPDFPageRef{327} /XYZ 219.998 117.447 null] >>}
\MusicPDFWriteObject{13630}{<< /D [\MusicPDFPageRef{329} /XYZ 235.506 513.996 null] >>}
\MusicPDFWriteObject{13631}{<< /D [\MusicPDFPageRef{329} /XYZ 186.397 423.843 null] >>}
\MusicPDFWriteObject{13632}{<< /D [\MusicPDFPageRef{329} /XYZ 72 340.534 null] >>}
\MusicPDFWriteObject{13633}{<< /D [\MusicPDFPageRef{329} /XYZ 72 310.928 null] >>}
\MusicPDFWriteObject{13634}{<< /D [\MusicPDFPageRef{329} /XYZ 72 281.322 null] >>}
\MusicPDFWriteObject{13635}{<< /D [\MusicPDFPageRef{329} /XYZ 72 251.716 null] >>}
\MusicPDFWriteObject{13636}{<< /D [\MusicPDFPageRef{330} /XYZ 264.979 683.057 null] >>}
\MusicPDFWriteObject{13637}{<< /D [\MusicPDFPageRef{330} /XYZ 264.979 633.478 null] >>}
\MusicPDFWriteObject{13638}{<< /D [\MusicPDFPageRef{326} /XYZ 282.069 234.77 null] >>}
\MusicPDFWriteObject{13639}{<< /D [\MusicPDFPageRef{330} /XYZ 206.003 502.266 null] >>}
\MusicPDFWriteObject{13640}{<< /D [\MusicPDFPageRef{330} /XYZ 162.134 213.477 null] >>}
\MusicPDFWriteObject{13641}{<< /D [\MusicPDFPageRef{330} /XYZ 172.88 164.727 null] >>}
\MusicPDFWriteObject{13642}{<< /D [\MusicPDFPageRef{331} /XYZ 245.548 462.813 null] >>}
\MusicPDFWriteObject{13643}{<< /D [\MusicPDFPageRef{331} /XYZ 251.545 389.413 null] >>}
\MusicPDFWriteObject{13644}{<< /D [\MusicPDFPageRef{331} /XYZ 260.202 323.856 null] >>}
\MusicPDFWriteObject{13645}{<< /D [\MusicPDFPageRef{331} /XYZ 213.625 287.429 null] >>}
\MusicPDFWriteObject{13646}{<< /D [\MusicPDFPageRef{331} /XYZ 179.461 203.038 null] >>}
\MusicPDFWriteObject{13647}{<< /D [\MusicPDFPageRef{331} /XYZ 180.245 117.447 null] >>}
\MusicPDFWriteObject{13648}{<< /D [\MusicPDFPageRef{326} /XYZ 254.044 187.368 null] >>}
\MusicPDFWriteObject{13649}{<< /D [\MusicPDFPageRef{326} /XYZ 262.252 136.213 null] >>}
\MusicPDFWriteObject{13650}{<< /D [\MusicPDFPageRef{326} /XYZ 234.907 99.785 null] >>}
\MusicPDFWriteObject{13651}{<< /D [\MusicPDFPageRef{327} /XYZ 170.77 637.556 null] >>}
\MusicPDFWriteObject{13652}{<< /D [\MusicPDFPageRef{327} /XYZ 197.408 566 null] >>}
\MusicPDFWriteObject{13653}{<< /D [\MusicPDFPageRef{327} /XYZ 193.653 511.64 null] >>}
\MusicPDFWriteObject{13654}{<< /D [\MusicPDFPageRef{327} /XYZ 224.736 335.074 null] >>}
\MusicPDFWriteObject{13655}{<< /D [\MusicPDFPageRef{332} /XYZ 237.599 468.212 null] >>}
\MusicPDFWriteObject{13656}{<< /D [\MusicPDFPageRef{332} /XYZ 235.599 382.206 null] >>}
\MusicPDFWriteObject{13657}{<< /D [\MusicPDFPageRef{333} /XYZ 254.934 669.386 null] >>}
\MusicPDFWriteObject{13658}{<< /D [\MusicPDFPageRef{335} /XYZ 261.723 343.801 null] >>}
\MusicPDFWriteObject{13659}{<< /D [\MusicPDFPageRef{335} /XYZ 271.676 161.962 null] >>}
\MusicPDFWriteObject{13660}{<< /D [\MusicPDFPageRef{336} /XYZ 213.42 538.461 null] >>}
\MusicPDFWriteObject{13661}{<< /D [\MusicPDFPageRef{336} /XYZ 153.738 110.358 null] >>}
\MusicPDFWriteObject{13662}{<< /D [\MusicPDFPageRef{337} /XYZ 195.429 634.816 null] >>}
\MusicPDFWriteObject{13663}{<< /D [\MusicPDFPageRef{337} /XYZ 261.213 468.294 null] >>}
\MusicPDFWriteObject{13664}{<< /D [\MusicPDFPageRef{337} /XYZ 261.092 408.278 null] >>}
\MusicPDFWriteObject{13665}{<< /D [\MusicPDFPageRef{337} /XYZ 138.774 332.646 null] >>}
\MusicPDFWriteObject{13666}{<< /D [\MusicPDFPageRef{337} /XYZ 209.755 274.8 null] >>}
\MusicPDFWriteObject{13667}{<< /D [\MusicPDFPageRef{337} /XYZ 90.733 159.749 null] >>}
\MusicPDFWriteObject{13668}{<< /D [\MusicPDFPageRef{333} /XYZ 266.676 578.678 null] >>}
\MusicPDFWriteObject{13669}{<< /D [\MusicPDFPageRef{337} /XYZ 185.179 103.218 null] >>}
\MusicPDFWriteObject{13670}{<< /D [\MusicPDFPageRef{338} /XYZ 247.521 683.057 null] >>}
\MusicPDFWriteObject{13671}{<< /D [\MusicPDFPageRef{338} /XYZ 72 593.783 null] >>}
\MusicPDFWriteObject{13672}{<< /D [\MusicPDFPageRef{338} /XYZ 72 563.845 null] >>}
\MusicPDFWriteObject{13673}{<< /D [\MusicPDFPageRef{338} /XYZ 72 478.535 null] >>}
\MusicPDFWriteObject{13674}{<< /D [\MusicPDFPageRef{338} /XYZ 72 424.256 null] >>}
\MusicPDFWriteObject{13675}{<< /D [\MusicPDFPageRef{338} /XYZ 72 387.572 null] >>}
\MusicPDFWriteObject{13676}{<< /D [\MusicPDFPageRef{338} /XYZ 165.007 303.328 null] >>}
\MusicPDFWriteObject{13677}{<< /D [\MusicPDFPageRef{333} /XYZ 72 420.32 null] >>}
\MusicPDFWriteObject{13678}{<< /D [\MusicPDFPageRef{333} /XYZ 72 364.868 null] >>}
\MusicPDFWriteObject{13679}{<< /D [\MusicPDFPageRef{333} /XYZ 72 336.374 null] >>}
\MusicPDFWriteObject{13680}{<< /D [\MusicPDFPageRef{334} /XYZ 224.723 601.524 null] >>}
\MusicPDFWriteObject{13681}{<< /D [\MusicPDFPageRef{334} /XYZ 238.722 515.518 null] >>}
\MusicPDFWriteObject{13682}{<< /D [\MusicPDFPageRef{334} /XYZ 206.968 302.376 null] >>}
\MusicPDFWriteObject{13683}{<< /D [\MusicPDFPageRef{334} /XYZ 205.855 110.358 null] >>}
\MusicPDFWriteObject{13684}{<< /D [\MusicPDFPageRef{338} /XYZ 236.793 109.146 null] >>}
\MusicPDFWriteObject{13685}{<< /D [\MusicPDFPageRef{339} /XYZ 274.208 643.604 null] >>}
\MusicPDFWriteObject{13686}{<< /D [\MusicPDFPageRef{340} /XYZ 221.575 591.4 null] >>}
\MusicPDFWriteObject{13687}{<< /D [\MusicPDFPageRef{341} /XYZ 177.458 552.243 null] >>}
\MusicPDFWriteObject{13688}{<< /D [\MusicPDFPageRef{344} /XYZ 254.454 433.536 null] >>}
\MusicPDFWriteObject{13689}{<< /D [\MusicPDFPageRef{344} /XYZ 265.39 365.397 null] >>}
\MusicPDFWriteObject{13690}{<< /D [\MusicPDFPageRef{344} /XYZ 145.455 228.38 null] >>}
\MusicPDFWriteObject{13691}{<< /D [\MusicPDFPageRef{344} /XYZ 154.01 191.739 null] >>}
\MusicPDFWriteObject{13692}{<< /D [\MusicPDFPageRef{344} /XYZ 237.21 99.532 null] >>}
\MusicPDFWriteObject{13693}{<< /D [\MusicPDFPageRef{345} /XYZ 180.949 602.214 null] >>}
\MusicPDFWriteObject{13694}{<< /D [\MusicPDFPageRef{345} /XYZ 213.924 552.423 null] >>}
\MusicPDFWriteObject{13695}{<< /D [\MusicPDFPageRef{345} /XYZ 259.034 430.376 null] >>}
\MusicPDFWriteObject{13696}{<< /D [\MusicPDFPageRef{345} /XYZ 215.372 133.86 null] >>}
\MusicPDFWriteObject{13697}{<< /D [\MusicPDFPageRef{346} /XYZ 72 394.14 null] >>}
\MusicPDFWriteObject{13698}{<< /D [\MusicPDFPageRef{341} /XYZ 203.53 240.11 null] >>}
\MusicPDFWriteObject{13699}{<< /D [\MusicPDFPageRef{347} /XYZ 73.101 605.516 null] >>}
\MusicPDFWriteObject{13700}{<< /D [\MusicPDFPageRef{347} /XYZ 180.604 395.822 null] >>}
\MusicPDFWriteObject{13701}{<< /D [\MusicPDFPageRef{348} /XYZ 275.305 642.429 null] >>}
\MusicPDFWriteObject{13702}{<< /D [\MusicPDFPageRef{348} /XYZ 126.024 602.189 null] >>}
\MusicPDFWriteObject{13703}{<< /D [\MusicPDFPageRef{348} /XYZ 238.451 185.26 null] >>}
\MusicPDFWriteObject{13704}{<< /D [\MusicPDFPageRef{348} /XYZ 204.126 148.036 null] >>}
\MusicPDFWriteObject{13705}{<< /D [\MusicPDFPageRef{349} /XYZ 248.854 202.79 null] >>}
\MusicPDFWriteObject{13706}{<< /D [\MusicPDFPageRef{341} /XYZ 177.701 111.111 null] >>}
\MusicPDFWriteObject{13707}{<< /D [\MusicPDFPageRef{342} /XYZ 164.155 669.109 null] >>}
\MusicPDFWriteObject{13708}{<< /D [\MusicPDFPageRef{343} /XYZ 151.572 620.761 null] >>}
\MusicPDFWriteObject{13709}{<< /D [\MusicPDFPageRef{343} /XYZ 274.285 406.105 null] >>}
\MusicPDFWriteObject{13710}{<< /D [\MusicPDFPageRef{343} /XYZ 237.722 367.912 null] >>}
\MusicPDFWriteObject{13711}{<< /D [\MusicPDFPageRef{343} /XYZ 262.344 179.385 null] >>}
\MusicPDFWriteObject{13712}{<< /D [\MusicPDFPageRef{344} /XYZ 253.07 669.906 null] >>}
\MusicPDFWriteObject{13713}{<< /D [\MusicPDFPageRef{351} /XYZ 259.461 501.374 null] >>}
\MusicPDFWriteObject{13714}{<< /D [\MusicPDFPageRef{353} /XYZ 184.049 683.057 null] >>}
\MusicPDFWriteObject{13715}{<< /D [\MusicPDFPageRef{353} /XYZ 276.385 391.892 null] >>}
\MusicPDFWriteObject{13716}{<< /D [\MusicPDFPageRef{355} /XYZ 259.8 651.025 null] >>}
\MusicPDFWriteObject{13717}{<< /D [\MusicPDFPageRef{355} /XYZ 254.258 613.441 null] >>}
\MusicPDFWriteObject{13718}{<< /D [\MusicPDFPageRef{355} /XYZ 217.338 541.253 null] >>}
\MusicPDFWriteObject{13719}{<< /D [\MusicPDFPageRef{355} /XYZ 254.116 496.689 null] >>}
\MusicPDFWriteObject{13720}{<< /D [\MusicPDFPageRef{355} /XYZ 184.018 433.036 null] >>}
\MusicPDFWriteObject{13721}{<< /D [\MusicPDFPageRef{355} /XYZ 183.491 395.566 null] >>}
\MusicPDFWriteObject{13722}{<< /D [\MusicPDFPageRef{355} /XYZ 175.904 139.055 null] >>}
\MusicPDFWriteObject{13723}{<< /D [\MusicPDFPageRef{355} /XYZ 180.214 101.584 null] >>}
\MusicPDFWriteObject{13724}{<< /D [\MusicPDFPageRef{351} /XYZ 254.258 463.79 null] >>}
\MusicPDFWriteObject{13725}{<< /D [\MusicPDFPageRef{356} /XYZ 177.478 573.488 null] >>}
\MusicPDFWriteObject{13726}{<< /D [\MusicPDFPageRef{356} /XYZ 169.971 536.017 null] >>}
\MusicPDFWriteObject{13727}{<< /D [\MusicPDFPageRef{356} /XYZ 276.385 334.14 null] >>}
\MusicPDFWriteObject{13728}{<< /D [\MusicPDFPageRef{351} /XYZ 205.498 391.603 null] >>}
\MusicPDFWriteObject{13729}{<< /D [\MusicPDFPageRef{351} /XYZ 253.44 347.039 null] >>}
\MusicPDFWriteObject{13730}{<< /D [\MusicPDFPageRef{351} /XYZ 218.936 276.439 null] >>}
\MusicPDFWriteObject{13731}{<< /D [\MusicPDFPageRef{351} /XYZ 183.68 230.942 null] >>}
\MusicPDFWriteObject{13732}{<< /D [\MusicPDFPageRef{351} /XYZ 179.211 193.471 null] >>}
\MusicPDFWriteObject{13733}{<< /D [\MusicPDFPageRef{352} /XYZ 186.212 576.705 null] >>}
\MusicPDFWriteObject{13734}{<< /D [\MusicPDFPageRef{352} /XYZ 177.115 539.234 null] >>}
\MusicPDFWriteObject{13735}{<< /D [\MusicPDFPageRef{358} /XYZ 86.773 658.177 null] >>}
\MusicPDFWriteObject{13736}{<< /D [\MusicPDFPageRef{360} /XYZ 169.962 327.492 null] >>}
\MusicPDFWriteObject{13737}{<< /D [\MusicPDFPageRef{360} /XYZ 119.442 275.217 null] >>}
\MusicPDFWriteObject{13738}{<< /D [\MusicPDFPageRef{360} /XYZ 270.917 220.424 null] >>}
\MusicPDFWriteObject{13739}{<< /D [\MusicPDFPageRef{360} /XYZ 271.255 183.783 null] >>}
\MusicPDFWriteObject{13740}{<< /D [\MusicPDFPageRef{360} /XYZ 205.196 133.992 null] >>}
\MusicPDFWriteObject{13741}{<< /D [\MusicPDFPageRef{361} /XYZ 72 506.303 null] >>}
\MusicPDFWriteObject{13742}{<< /D [\MusicPDFPageRef{361} /XYZ 245.401 249.515 null] >>}
\MusicPDFWriteObject{13743}{<< /D [\MusicPDFPageRef{361} /XYZ 245.401 213.088 null] >>}
\MusicPDFWriteObject{13744}{<< /D [\MusicPDFPageRef{361} /XYZ 162.449 175.617 null] >>}
\MusicPDFWriteObject{13745}{<< /D [\MusicPDFPageRef{361} /XYZ 163.216 139.189 null] >>}
\MusicPDFWriteObject{13746}{<< /D [\MusicPDFPageRef{358} /XYZ 229.274 555.409 null] >>}
\MusicPDFWriteObject{13747}{<< /D [\MusicPDFPageRef{362} /XYZ 235.813 519.375 null] >>}
\MusicPDFWriteObject{13748}{<< /D [\MusicPDFPageRef{362} /XYZ 237.697 482.947 null] >>}
\MusicPDFWriteObject{13749}{<< /D [\MusicPDFPageRef{362} /XYZ 219.634 423.33 null] >>}
\MusicPDFWriteObject{13750}{<< /D [\MusicPDFPageRef{362} /XYZ 264.789 386.902 null] >>}
\MusicPDFWriteObject{13751}{<< /D [\MusicPDFPageRef{362} /XYZ 267.516 306.715 null] >>}
\MusicPDFWriteObject{13752}{<< /D [\MusicPDFPageRef{363} /XYZ 194.542 682.233 null] >>}
\MusicPDFWriteObject{13753}{<< /D [\MusicPDFPageRef{363} /XYZ 153.684 546.586 null] >>}
\MusicPDFWriteObject{13754}{<< /D [\MusicPDFPageRef{363} /XYZ 72 502.123 null] >>}
\MusicPDFWriteObject{13755}{<< /D [\MusicPDFPageRef{363} /XYZ 72 482.935 null] >>}
\MusicPDFWriteObject{13756}{<< /D [\MusicPDFPageRef{363} /XYZ 72 453.918 null] >>}
\MusicPDFWriteObject{13757}{<< /D [\MusicPDFPageRef{358} /XYZ 185.702 104.025 null] >>}
\MusicPDFWriteObject{13758}{<< /D [\MusicPDFPageRef{363} /XYZ 218.858 101.996 null] >>}
\MusicPDFWriteObject{13759}{<< /D [\MusicPDFPageRef{359} /XYZ 227.448 532.95 null] >>}
\MusicPDFWriteObject{13760}{<< /D [\MusicPDFPageRef{359} /XYZ 130.019 278.959 null] >>}
\MusicPDFWriteObject{13761}{<< /D [\MusicPDFPageRef{359} /XYZ 131.892 197.533 null] >>}
\MusicPDFWriteObject{13762}{<< /D [\MusicPDFPageRef{360} /XYZ 134.475 528.608 null] >>}
\MusicPDFWriteObject{13763}{<< /D [\MusicPDFPageRef{360} /XYZ 148.406 491.138 null] >>}
\MusicPDFWriteObject{13764}{<< /D [\MusicPDFPageRef{360} /XYZ 230.818 397.269 null] >>}
\MusicPDFWriteObject{13765}{<< /D [\MusicPDFPageRef{364} /XYZ 196.334 469.67 null] >>}
\MusicPDFWriteObject{13766}{<< /D [\MusicPDFPageRef{364} /XYZ 204.629 351.595 null] >>}
\MusicPDFWriteObject{13767}{<< /D [\MusicPDFPageRef{364} /XYZ 145.648 242.748 null] >>}
\MusicPDFWriteObject{13768}{<< /D [\MusicPDFPageRef{365} /XYZ 223.058 644.922 null] >>}
\MusicPDFWriteObject{13769}{<< /D [\MusicPDFPageRef{365} /XYZ 205.125 603.266 null] >>}
\MusicPDFWriteObject{13770}{<< /D [\MusicPDFPageRef{365} /XYZ 145.648 555.106 null] >>}
\MusicPDFWriteObject{13771}{<< /D [\MusicPDFPageRef{365} /XYZ 72 508.331 null] >>}
\MusicPDFWriteObject{13772}{<< /D [\MusicPDFPageRef{365} /XYZ 230.282 415.554 null] >>}
\MusicPDFWriteObject{13773}{<< /D [\MusicPDFPageRef{365} /XYZ 164.055 201.78 null] >>}
\MusicPDFWriteObject{13774}{<< /D [\MusicPDFPageRef{366} /XYZ 183.156 541.619 null] >>}
\MusicPDFWriteObject{13775}{<< /D [\MusicPDFPageRef{366} /XYZ 166.915 505.191 null] >>}
\MusicPDFWriteObject{13776}{<< /D [\MusicPDFPageRef{366} /XYZ 153.943 327.237 null] >>}
\MusicPDFWriteObject{13777}{<< /D [\MusicPDFPageRef{366} /XYZ 158.7 110.579 null] >>}
\MusicPDFWriteObject{13778}{<< /D [\MusicPDFPageRef{367} /XYZ 175.859 277.856 null] >>}
\MusicPDFWriteObject{13779}{<< /D [\MusicPDFPageRef{368} /XYZ 209.024 596.214 null] >>}
\MusicPDFWriteObject{13780}{<< /D [\MusicPDFPageRef{370} /XYZ 252.766 327.695 null] >>}
\MusicPDFWriteObject{13781}{<< /D [\MusicPDFPageRef{371} /XYZ 206.248 520.254 null] >>}
\MusicPDFWriteObject{13782}{<< /D [\MusicPDFPageRef{371} /XYZ 72 183.633 null] >>}
\MusicPDFWriteObject{13783}{<< /D [\MusicPDFPageRef{372} /XYZ 187.901 309.684 null] >>}
\MusicPDFWriteObject{13784}{<< /D [\MusicPDFPageRef{373} /XYZ 132.833 679.161 null] >>}
\MusicPDFWriteObject{13785}{<< /D [\MusicPDFPageRef{373} /XYZ 124.329 397.223 null] >>}
\MusicPDFWriteObject{13786}{<< /D [\MusicPDFPageRef{373} /XYZ 221.089 268.535 null] >>}
\MusicPDFWriteObject{13787}{<< /D [\MusicPDFPageRef{373} /XYZ 145.437 109.458 null] >>}
\MusicPDFWriteObject{13788}{<< /D [\MusicPDFPageRef{368} /XYZ 259.596 487.242 null] >>}
\MusicPDFWriteObject{13789}{<< /D [\MusicPDFPageRef{368} /XYZ 224.326 414.04 null] >>}
\MusicPDFWriteObject{13790}{<< /D [\MusicPDFPageRef{368} /XYZ 146.634 299.803 null] >>}
\MusicPDFWriteObject{13791}{<< /D [\MusicPDFPageRef{368} /XYZ 255.014 149.364 null] >>}
\MusicPDFWriteObject{13792}{<< /D [\MusicPDFPageRef{369} /XYZ 192.248 449.01 null] >>}
\MusicPDFWriteObject{13793}{<< /D [\MusicPDFPageRef{369} /XYZ 265.408 228.484 null] >>}
\MusicPDFWriteObject{13794}{<< /D [\MusicPDFPageRef{370} /XYZ 155.913 525.47 null] >>}
\MusicPDFWriteObject{13795}{<< /D [\MusicPDFPageRef{370} /XYZ 262.718 390.424 null] >>}
\MusicPDFWriteObject{13796}{<< /D [\MusicPDFPageRef{374} /XYZ 135.728 228.044 null] >>}
\MusicPDFWriteObject{13797}{<< /D [\MusicPDFPageRef{377} /XYZ 72 347.526 null] >>}
\MusicPDFWriteObject{13798}{<< /D [\MusicPDFPageRef{377} /XYZ 72 314.523 null] >>}
\MusicPDFWriteObject{13799}{<< /D [\MusicPDFPageRef{377} /XYZ 240.62 200.451 null] >>}
\MusicPDFWriteObject{13800}{<< /D [\MusicPDFPageRef{378} /XYZ 72 640.29 null] >>}
\MusicPDFWriteObject{13801}{<< /D [\MusicPDFPageRef{378} /XYZ 232.885 433.416 null] >>}
\MusicPDFWriteObject{13802}{<< /D [\MusicPDFPageRef{378} /XYZ 198.231 194.137 null] >>}
\MusicPDFWriteObject{13803}{<< /D [\MusicPDFPageRef{378} /XYZ 184.266 144.559 null] >>}
\MusicPDFWriteObject{13804}{<< /D [\MusicPDFPageRef{379} /XYZ 206.447 369.354 null] >>}
\MusicPDFWriteObject{13805}{<< /D [\MusicPDFPageRef{374} /XYZ 221.807 191.616 null] >>}
\MusicPDFWriteObject{13806}{<< /D [\MusicPDFPageRef{375} /XYZ 155.707 591.526 null] >>}
\MusicPDFWriteObject{13807}{<< /D [\MusicPDFPageRef{375} /XYZ 130.529 555.098 null] >>}
\MusicPDFWriteObject{13808}{<< /D [\MusicPDFPageRef{375} /XYZ 72 498.703 null] >>}
\MusicPDFWriteObject{13809}{<< /D [\MusicPDFPageRef{375} /XYZ 140.336 401.69 null] >>}
\MusicPDFWriteObject{13810}{<< /D [\MusicPDFPageRef{376} /XYZ 191.689 591.001 null] >>}
\MusicPDFWriteObject{13811}{<< /D [\MusicPDFPageRef{377} /XYZ 230.874 520.695 null] >>}
\MusicPDFWriteObject{13812}{<< /D [\MusicPDFPageRef{377} /XYZ 72 381.698 null] >>}
\MusicPDFWriteObject{13813}{<< /D [\MusicPDFPageRef{382} /XYZ 72 493.545 null] >>}
\MusicPDFWriteObject{13814}{<< /D [\MusicPDFPageRef{383} /XYZ 178.356 342.681 null] >>}
\MusicPDFWriteObject{13815}{<< /D [\MusicPDFPageRef{382} /XYZ 182.882 455.138 null] >>}
\MusicPDFWriteObject{13816}{<< /D [\MusicPDFPageRef{382} /XYZ 72 387.837 null] >>}
\MusicPDFWriteObject{13817}{<< /D [\MusicPDFPageRef{382} /XYZ 190.467 349.43 null] >>}
\MusicPDFWriteObject{13818}{<< /D [\MusicPDFPageRef{382} /XYZ 187.708 199.096 null] >>}
\MusicPDFWriteObject{13819}{<< /D [\MusicPDFPageRef{383} /XYZ 172.207 676.545 null] >>}
\MusicPDFWriteObject{13820}{<< /D [\MusicPDFPageRef{383} /XYZ 212.531 636.08 null] >>}
\MusicPDFWriteObject{13821}{<< /D [\MusicPDFPageRef{383} /XYZ 252.98 473.268 null] >>}
\MusicPDFWriteObject{13822}{<< /D [\MusicPDFPageRef{383} /XYZ 185.158 383.388 null] >>}
\MusicPDFWriteObject{13823}{<< /D [\MusicPDFPageRef{384} /XYZ 229.055 279.332 null] >>}
\MusicPDFWriteObject{13824}{<< /D [\MusicPDFPageRef{386} /XYZ 257.311 332.11 null] >>}
\MusicPDFWriteObject{13825}{<< /D [\MusicPDFPageRef{384} /XYZ 232.788 242.905 null] >>}
\MusicPDFWriteObject{13826}{<< /D [\MusicPDFPageRef{384} /XYZ 222.674 153.874 null] >>}
\MusicPDFWriteObject{13827}{<< /D [\MusicPDFPageRef{384} /XYZ 236.615 117.447 null] >>}
\MusicPDFWriteObject{13828}{<< /D [\MusicPDFPageRef{385} /XYZ 172.428 682.26 null] >>}
\MusicPDFWriteObject{13829}{<< /D [\MusicPDFPageRef{385} /XYZ 181.507 507.943 null] >>}
\MusicPDFWriteObject{13830}{<< /D [\MusicPDFPageRef{385} /XYZ 198.694 471.515 null] >>}
\MusicPDFWriteObject{13831}{<< /D [\MusicPDFPageRef{385} /XYZ 172.284 421.141 null] >>}
\MusicPDFWriteObject{13832}{<< /D [\MusicPDFPageRef{386} /XYZ 251.599 368.537 null] >>}
\MusicPDFWriteObject{13833}{<< /D [\MusicPDFPageRef{387} /XYZ 208.192 100.427 null] >>}
\MusicPDFWriteObject{13834}{<< /D [\MusicPDFPageRef{388} /XYZ 156.003 665.772 null] >>}
\MusicPDFWriteObject{13835}{<< /D [\MusicPDFPageRef{388} /XYZ 263.61 617.201 null] >>}
\MusicPDFWriteObject{13836}{<< /D [\MusicPDFPageRef{388} /XYZ 276.301 580.531 null] >>}
\MusicPDFWriteObject{13837}{<< /D [\MusicPDFPageRef{388} /XYZ 269.491 519.418 null] >>}
\MusicPDFWriteObject{13838}{<< /D [\MusicPDFPageRef{388} /XYZ 152.907 459.336 null] >>}
\MusicPDFWriteObject{13839}{<< /D [\MusicPDFPageRef{388} /XYZ 262.062 411.734 null] >>}
\MusicPDFWriteObject{13840}{<< /D [\MusicPDFPageRef{388} /XYZ 286.286 375.063 null] >>}
\MusicPDFWriteObject{13841}{<< /D [\MusicPDFPageRef{389} /XYZ 72 427.087 null] >>}
\MusicPDFWriteObject{13842}{<< /D [\MusicPDFPageRef{392} /XYZ 72 305.232 null] >>}
\MusicPDFWriteObject{13843}{<< /D [\MusicPDFPageRef{392} /XYZ 72 272.521 null] >>}
\MusicPDFWriteObject{13844}{<< /D [\MusicPDFPageRef{392} /XYZ 72 157.377 null] >>}
\MusicPDFWriteObject{13845}{<< /D [\MusicPDFPageRef{392} /XYZ 72 120.232 null] >>}
\MusicPDFWriteObject{13846}{<< /D [\MusicPDFPageRef{393} /XYZ 72 461.889 null] >>}
\MusicPDFWriteObject{13847}{<< /D [\MusicPDFPageRef{389} /XYZ 72 371.636 null] >>}
\MusicPDFWriteObject{13848}{<< /D [\MusicPDFPageRef{389} /XYZ 72 343.141 null] >>}
\MusicPDFWriteObject{13849}{<< /D [\MusicPDFPageRef{389} /XYZ 148.236 221.652 null] >>}
\MusicPDFWriteObject{13850}{<< /D [\MusicPDFPageRef{389} /XYZ 134.676 176.213 null] >>}
\MusicPDFWriteObject{13851}{<< /D [\MusicPDFPageRef{390} /XYZ 78.706 674.211 null] >>}
\MusicPDFWriteObject{13852}{<< /D [\MusicPDFPageRef{390} /XYZ 72 383.436 null] >>}
\MusicPDFWriteObject{13853}{<< /D [\MusicPDFPageRef{391} /XYZ 121.539 680.212 null] >>}
\MusicPDFWriteObject{13854}{<< /D [\MusicPDFPageRef{391} /XYZ 135.493 517.852 null] >>}
\MusicPDFWriteObject{13855}{<< /D [\MusicPDFPageRef{396} /XYZ 271.402 558.434 null] >>}
\MusicPDFWriteObject{13856}{<< /D [\MusicPDFPageRef{398} /XYZ 243.309 679.714 null] >>}
\MusicPDFWriteObject{13857}{<< /D [\MusicPDFPageRef{398} /XYZ 220.665 639.502 null] >>}
\MusicPDFWriteObject{13858}{<< /D [\MusicPDFPageRef{398} /XYZ 247.498 603.075 null] >>}
\MusicPDFWriteObject{13859}{<< /D [\MusicPDFPageRef{398} /XYZ 249.98 483.547 null] >>}
\MusicPDFWriteObject{13860}{<< /D [\MusicPDFPageRef{398} /XYZ 260.77 393.544 null] >>}
\MusicPDFWriteObject{13861}{<< /D [\MusicPDFPageRef{398} /XYZ 134.654 331.961 null] >>}
\MusicPDFWriteObject{13862}{<< /D [\MusicPDFPageRef{398} /XYZ 195.475 262.347 null] >>}
\MusicPDFWriteObject{13863}{<< /D [\MusicPDFPageRef{399} /XYZ 208.808 548.642 null] >>}
\MusicPDFWriteObject{13864}{<< /D [\MusicPDFPageRef{396} /XYZ 246.425 522.006 null] >>}
\MusicPDFWriteObject{13865}{<< /D [\MusicPDFPageRef{396} /XYZ 186.835 465.383 null] >>}
\MusicPDFWriteObject{13866}{<< /D [\MusicPDFPageRef{396} /XYZ 254.348 135.661 null] >>}
\MusicPDFWriteObject{13867}{<< /D [\MusicPDFPageRef{396} /XYZ 249.371 99.233 null] >>}
\MusicPDFWriteObject{13868}{<< /D [\MusicPDFPageRef{397} /XYZ 248.637 683.057 null] >>}
\MusicPDFWriteObject{13869}{<< /D [\MusicPDFPageRef{397} /XYZ 261.36 618.255 null] >>}
\MusicPDFWriteObject{13870}{<< /D [\MusicPDFPageRef{397} /XYZ 247.542 581.828 null] >>}
\MusicPDFWriteObject{13871}{<< /D [\MusicPDFPageRef{397} /XYZ 263.065 99.233 null] >>}
\MusicPDFWriteObject{13872}{<< /D [\MusicPDFPageRef{399} /XYZ 271.81 374.234 null] >>}
\MusicPDFWriteObject{13873}{<< /D [\MusicPDFPageRef{401} /XYZ 198.07 536.897 null] >>}
\MusicPDFWriteObject{13874}{<< /D [\MusicPDFPageRef{401} /XYZ 158.585 409.926 null] >>}
\MusicPDFWriteObject{13875}{<< /D [\MusicPDFPageRef{401} /XYZ 154.101 322.613 null] >>}
\MusicPDFWriteObject{13876}{<< /D [\MusicPDFPageRef{402} /XYZ 265.969 643.604 null] >>}
\MusicPDFWriteObject{13877}{<< /D [\MusicPDFPageRef{402} /XYZ 248.089 427.407 null] >>}
\MusicPDFWriteObject{13878}{<< /D [\MusicPDFPageRef{402} /XYZ 253.15 351.527 null] >>}
\MusicPDFWriteObject{13879}{<< /D [\MusicPDFPageRef{402} /XYZ 283.292 311.603 null] >>}
\MusicPDFWriteObject{13880}{<< /D [\MusicPDFPageRef{403} /XYZ 257.444 669.906 null] >>}
\MusicPDFWriteObject{13881}{<< /D [\MusicPDFPageRef{399} /XYZ 257.567 149.049 null] >>}
\MusicPDFWriteObject{13882}{<< /D [\MusicPDFPageRef{399} /XYZ 195.368 104.713 null] >>}
\MusicPDFWriteObject{13883}{<< /D [\MusicPDFPageRef{400} /XYZ 272.472 392.951 null] >>}
\MusicPDFWriteObject{13884}{<< /D [\MusicPDFPageRef{400} /XYZ 225.077 348.974 null] >>}
\MusicPDFWriteObject{13885}{<< /D [\MusicPDFPageRef{400} /XYZ 270.094 286.26 null] >>}
\MusicPDFWriteObject{13886}{<< /D [\MusicPDFPageRef{400} /XYZ 207.56 242.283 null] >>}
\MusicPDFWriteObject{13887}{<< /D [\MusicPDFPageRef{400} /XYZ 208.329 171.71 null] >>}
\MusicPDFWriteObject{13888}{<< /D [\MusicPDFPageRef{400} /XYZ 165.198 110.69 null] >>}
\MusicPDFWriteObject{13889}{<< /D [\MusicPDFPageRef{403} /XYZ 241.012 398.32 null] >>}
\MusicPDFWriteObject{13890}{<< /D [\MusicPDFPageRef{404} /XYZ 254.253 523.553 null] >>}
\MusicPDFWriteObject{13891}{<< /D [\MusicPDFPageRef{406} /XYZ 72 488.482 null] >>}
\MusicPDFWriteObject{13892}{<< /D [\MusicPDFPageRef{406} /XYZ 72 456.159 null] >>}
\MusicPDFWriteObject{13893}{<< /D [\MusicPDFPageRef{407} /XYZ 239.533 663.984 null] >>}
\MusicPDFWriteObject{13894}{<< /D [\MusicPDFPageRef{407} /XYZ 272.68 623.346 null] >>}
\MusicPDFWriteObject{13895}{<< /D [\MusicPDFPageRef{407} /XYZ 208.525 484.668 null] >>}
\MusicPDFWriteObject{13896}{<< /D [\MusicPDFPageRef{407} /XYZ 190.656 424.462 null] >>}
\MusicPDFWriteObject{13897}{<< /D [\MusicPDFPageRef{407} /XYZ 218.628 273.738 null] >>}
\MusicPDFWriteObject{13898}{<< /D [\MusicPDFPageRef{408} /XYZ 249.551 683.057 null] >>}
\MusicPDFWriteObject{13899}{<< /D [\MusicPDFPageRef{408} /XYZ 259.981 620.327 null] >>}
\MusicPDFWriteObject{13900}{<< /D [\MusicPDFPageRef{408} /XYZ 192.399 563.2 null] >>}
\MusicPDFWriteObject{13901}{<< /D [\MusicPDFPageRef{404} /XYZ 214.618 480.13 null] >>}
\MusicPDFWriteObject{13902}{<< /D [\MusicPDFPageRef{408} /XYZ 214.069 438.995 null] >>}
\MusicPDFWriteObject{13903}{<< /D [\MusicPDFPageRef{408} /XYZ 186.722 389.417 null] >>}
\MusicPDFWriteObject{13904}{<< /D [\MusicPDFPageRef{408} /XYZ 72 311.866 null] >>}
\MusicPDFWriteObject{13905}{<< /D [\MusicPDFPageRef{409} /XYZ 239.533 597.339 null] >>}
\MusicPDFWriteObject{13906}{<< /D [\MusicPDFPageRef{409} /XYZ 272.68 556.701 null] >>}
\MusicPDFWriteObject{13907}{<< /D [\MusicPDFPageRef{409} /XYZ 167.96 395.072 null] >>}
\MusicPDFWriteObject{13908}{<< /D [\MusicPDFPageRef{409} /XYZ 233.469 231.817 null] >>}
\MusicPDFWriteObject{13909}{<< /D [\MusicPDFPageRef{409} /XYZ 245.78 191.949 null] >>}
\MusicPDFWriteObject{13910}{<< /D [\MusicPDFPageRef{409} /XYZ 226.268 150.238 null] >>}
\MusicPDFWriteObject{13911}{<< /D [\MusicPDFPageRef{409} /XYZ 207.321 105.824 null] >>}
\MusicPDFWriteObject{13912}{<< /D [\MusicPDFPageRef{404} /XYZ 209.893 396.797 null] >>}
\MusicPDFWriteObject{13913}{<< /D [\MusicPDFPageRef{410} /XYZ 218.809 683.017 null] >>}
\MusicPDFWriteObject{13914}{<< /D [\MusicPDFPageRef{410} /XYZ 239.752 646.589 null] >>}
\MusicPDFWriteObject{13915}{<< /D [\MusicPDFPageRef{410} /XYZ 251.129 544.013 null] >>}
\MusicPDFWriteObject{13916}{<< /D [\MusicPDFPageRef{410} /XYZ 268.867 443.429 null] >>}
\MusicPDFWriteObject{13917}{<< /D [\MusicPDFPageRef{410} /XYZ 151.756 383.143 null] >>}
\MusicPDFWriteObject{13918}{<< /D [\MusicPDFPageRef{410} /XYZ 233.632 312.868 null] >>}
\MusicPDFWriteObject{13919}{<< /D [\MusicPDFPageRef{411} /XYZ 224.889 627.214 null] >>}
\MusicPDFWriteObject{13920}{<< /D [\MusicPDFPageRef{404} /XYZ 238.095 250.959 null] >>}
\MusicPDFWriteObject{13921}{<< /D [\MusicPDFPageRef{404} /XYZ 239.658 117.447 null] >>}
\MusicPDFWriteObject{13922}{<< /D [\MusicPDFPageRef{405} /XYZ 123.815 667.283 null] >>}
\MusicPDFWriteObject{13923}{<< /D [\MusicPDFPageRef{405} /XYZ 239.437 514.856 null] >>}
\MusicPDFWriteObject{13924}{<< /D [\MusicPDFPageRef{405} /XYZ 211.22 267.997 null] >>}
\MusicPDFWriteObject{13925}{<< /D [\MusicPDFPageRef{406} /XYZ 72 525.422 null] >>}
\MusicPDFWriteObject{13926}{<< /D [\MusicPDFPageRef{411} /XYZ 203.093 411.174 null] >>}
\MusicPDFWriteObject{13927}{<< /D [\MusicPDFPageRef{413} /XYZ 72 599.851 null] >>}
\MusicPDFWriteObject{13928}{<< /D [\MusicPDFPageRef{413} /XYZ 72 537.27 null] >>}
\MusicPDFWriteObject{13929}{<< /D [\MusicPDFPageRef{413} /XYZ 72 436.474 null] >>}
\MusicPDFWriteObject{13930}{<< /D [\MusicPDFPageRef{413} /XYZ 72 407.711 null] >>}
\MusicPDFWriteObject{13931}{<< /D [\MusicPDFPageRef{413} /XYZ 195.661 361.84 null] >>}
\MusicPDFWriteObject{13932}{<< /D [\MusicPDFPageRef{413} /XYZ 92.822 317.101 null] >>}
\MusicPDFWriteObject{13933}{<< /D [\MusicPDFPageRef{413} /XYZ 241.573 280.451 null] >>}
\MusicPDFWriteObject{13934}{<< /D [\MusicPDFPageRef{413} /XYZ 208.833 223.323 null] >>}
\MusicPDFWriteObject{13935}{<< /D [\MusicPDFPageRef{413} /XYZ 181.234 149.642 null] >>}
\MusicPDFWriteObject{13936}{<< /D [\MusicPDFPageRef{414} /XYZ 230.601 640.21 null] >>}
\MusicPDFWriteObject{13937}{<< /D [\MusicPDFPageRef{411} /XYZ 207.807 320.842 null] >>}
\MusicPDFWriteObject{13938}{<< /D [\MusicPDFPageRef{414} /XYZ 242.484 482.613 null] >>}
\MusicPDFWriteObject{13939}{<< /D [\MusicPDFPageRef{414} /XYZ 256.994 414.298 null] >>}
\MusicPDFWriteObject{13940}{<< /D [\MusicPDFPageRef{414} /XYZ 249.167 377.87 null] >>}
\MusicPDFWriteObject{13941}{<< /D [\MusicPDFPageRef{414} /XYZ 276.676 261.939 null] >>}
\MusicPDFWriteObject{13942}{<< /D [\MusicPDFPageRef{414} /XYZ 276.676 212.361 null] >>}
\MusicPDFWriteObject{13943}{<< /D [\MusicPDFPageRef{415} /XYZ 238.384 643.604 null] >>}
\MusicPDFWriteObject{13944}{<< /D [\MusicPDFPageRef{415} /XYZ 234.297 564.455 null] >>}
\MusicPDFWriteObject{13945}{<< /D [\MusicPDFPageRef{415} /XYZ 143.786 475.425 null] >>}
\MusicPDFWriteObject{13946}{<< /D [\MusicPDFPageRef{415} /XYZ 239.437 324.8 null] >>}
\MusicPDFWriteObject{13947}{<< /D [\MusicPDFPageRef{415} /XYZ 276.676 142.482 null] >>}
\MusicPDFWriteObject{13948}{<< /D [\MusicPDFPageRef{412} /XYZ 147.966 522.199 null] >>}
\MusicPDFWriteObject{13949}{<< /D [\MusicPDFPageRef{416} /XYZ 276.676 683.057 null] >>}
\MusicPDFWriteObject{13950}{<< /D [\MusicPDFPageRef{416} /XYZ 151.575 342.744 null] >>}
\MusicPDFWriteObject{13951}{<< /D [\MusicPDFPageRef{416} /XYZ 232.373 170.2 null] >>}
\MusicPDFWriteObject{13952}{<< /D [\MusicPDFPageRef{416} /XYZ 234.487 119.731 null] >>}
\MusicPDFWriteObject{13953}{<< /D [\MusicPDFPageRef{417} /XYZ 160.393 576.871 null] >>}
\MusicPDFWriteObject{13954}{<< /D [\MusicPDFPageRef{417} /XYZ 157.464 514.963 null] >>}
\MusicPDFWriteObject{13955}{<< /D [\MusicPDFPageRef{417} /XYZ 156.771 465 null] >>}
\MusicPDFWriteObject{13956}{<< /D [\MusicPDFPageRef{417} /XYZ 244.925 346.872 null] >>}
\MusicPDFWriteObject{13957}{<< /D [\MusicPDFPageRef{417} /XYZ 242.137 307.102 null] >>}
\MusicPDFWriteObject{13958}{<< /D [\MusicPDFPageRef{417} /XYZ 237.724 266.89 null] >>}
\MusicPDFWriteObject{13959}{<< /D [\MusicPDFPageRef{412} /XYZ 176.093 358.241 null] >>}
\MusicPDFWriteObject{13960}{<< /D [\MusicPDFPageRef{417} /XYZ 203.678 222.476 null] >>}
\MusicPDFWriteObject{13961}{<< /D [\MusicPDFPageRef{417} /XYZ 208.007 181.227 null] >>}
\MusicPDFWriteObject{13962}{<< /D [\MusicPDFPageRef{417} /XYZ 251.207 144.799 null] >>}
\MusicPDFWriteObject{13963}{<< /D [\MusicPDFPageRef{418} /XYZ 209.496 555.302 null] >>}
\MusicPDFWriteObject{13964}{<< /D [\MusicPDFPageRef{418} /XYZ 250.833 402.641 null] >>}
\MusicPDFWriteObject{13965}{<< /D [\MusicPDFPageRef{418} /XYZ 224.961 282.065 null] >>}
\MusicPDFWriteObject{13966}{<< /D [\MusicPDFPageRef{418} /XYZ 177.56 232.487 null] >>}
\MusicPDFWriteObject{13967}{<< /D [\MusicPDFPageRef{418} /XYZ 72 142.944 null] >>}
\MusicPDFWriteObject{13968}{<< /D [\MusicPDFPageRef{412} /XYZ 154.478 294.507 null] >>}
\MusicPDFWriteObject{13969}{<< /D [\MusicPDFPageRef{412} /XYZ 166.873 232.156 null] >>}
\MusicPDFWriteObject{13970}{<< /D [\MusicPDFPageRef{412} /XYZ 144.973 155.18 null] >>}
\MusicPDFWriteObject{13971}{<< /D [\MusicPDFPageRef{412} /XYZ 169.704 105.978 null] >>}
\MusicPDFWriteObject{13972}{<< /D [\MusicPDFPageRef{413} /XYZ 142.76 648.493 null] >>}
\MusicPDFWriteObject{13973}{<< /D [\MusicPDFPageRef{419} /XYZ 265.359 296.059 null] >>}
\MusicPDFWriteObject{13974}{<< /D [\MusicPDFPageRef{419} /XYZ 239.939 255.735 null] >>}
\MusicPDFWriteObject{13975}{<< /D [\MusicPDFPageRef{420} /XYZ 212.406 236.757 null] >>}
\MusicPDFWriteObject{13976}{<< /D [\MusicPDFPageRef{420} /XYZ 229.167 196.074 null] >>}
\MusicPDFWriteObject{13977}{<< /D [\MusicPDFPageRef{421} /XYZ 262.779 637.312 null] >>}
\MusicPDFWriteObject{13978}{<< /D [\MusicPDFPageRef{421} /XYZ 219.179 600.845 null] >>}
\MusicPDFWriteObject{13979}{<< /D [\MusicPDFPageRef{421} /XYZ 80.393 298.975 null] >>}
\MusicPDFWriteObject{13980}{<< /D [\MusicPDFPageRef{422} /XYZ 247.283 449.919 null] >>}
\MusicPDFWriteObject{13981}{<< /D [\MusicPDFPageRef{426} /XYZ 183.031 535.36 null] >>}
\MusicPDFWriteObject{13982}{<< /D [\MusicPDFPageRef{422} /XYZ 206.125 104.325 null] >>}
\MusicPDFWriteObject{13983}{<< /D [\MusicPDFPageRef{423} /XYZ 266.864 301.228 null] >>}
\MusicPDFWriteObject{13984}{<< /D [\MusicPDFPageRef{423} /XYZ 269.655 189.099 null] >>}
\MusicPDFWriteObject{13985}{<< /D [\MusicPDFPageRef{424} /XYZ 243.546 592.19 null] >>}
\MusicPDFWriteObject{13986}{<< /D [\MusicPDFPageRef{424} /XYZ 72 504.641 null] >>}
\MusicPDFWriteObject{13987}{<< /D [\MusicPDFPageRef{424} /XYZ 186.426 460.444 null] >>}
\MusicPDFWriteObject{13988}{<< /D [\MusicPDFPageRef{424} /XYZ 116.456 420.023 null] >>}
\MusicPDFWriteObject{13989}{<< /D [\MusicPDFPageRef{424} /XYZ 274.939 349.457 null] >>}
\MusicPDFWriteObject{13990}{<< /D [\MusicPDFPageRef{427} /XYZ 258.119 337.506 null] >>}
\MusicPDFWriteObject{13991}{<< /D [\MusicPDFPageRef{429} /XYZ 165.302 234.929 null] >>}
\MusicPDFWriteObject{13992}{<< /D [\MusicPDFPageRef{429} /XYZ 243.598 188.605 null] >>}
\MusicPDFWriteObject{13993}{<< /D [\MusicPDFPageRef{429} /XYZ 144.063 104.713 null] >>}
\MusicPDFWriteObject{13994}{<< /D [\MusicPDFPageRef{430} /XYZ 230.206 669.906 null] >>}
\MusicPDFWriteObject{13995}{<< /D [\MusicPDFPageRef{431} /XYZ 200.648 613.075 null] >>}
\MusicPDFWriteObject{13996}{<< /D [\MusicPDFPageRef{431} /XYZ 250.849 572.475 null] >>}
\MusicPDFWriteObject{13997}{<< /D [\MusicPDFPageRef{431} /XYZ 72 447.905 null] >>}
\MusicPDFWriteObject{13998}{<< /D [\MusicPDFPageRef{431} /XYZ 72 394.949 null] >>}
\MusicPDFWriteObject{13999}{<< /D [\MusicPDFPageRef{431} /XYZ 72 363.982 null] >>}
\MusicPDFWriteObject{14000}{<< /D [\MusicPDFPageRef{431} /XYZ 72 332.549 null] >>}
\MusicPDFWriteObject{14001}{<< /D [\MusicPDFPageRef{427} /XYZ 260.953 301.079 null] >>}
\MusicPDFWriteObject{14002}{<< /D [\MusicPDFPageRef{431} /XYZ 72 301.582 null] >>}
\MusicPDFWriteObject{14003}{<< /D [\MusicPDFPageRef{431} /XYZ 72 240.481 null] >>}
\MusicPDFWriteObject{14004}{<< /D [\MusicPDFPageRef{431} /XYZ 177.34 193.185 null] >>}
\MusicPDFWriteObject{14005}{<< /D [\MusicPDFPageRef{431} /XYZ 150.864 145.837 null] >>}
\MusicPDFWriteObject{14006}{<< /D [\MusicPDFPageRef{427} /XYZ 270.97 171.204 null] >>}
\MusicPDFWriteObject{14007}{<< /D [\MusicPDFPageRef{428} /XYZ 198.982 515.381 null] >>}
\MusicPDFWriteObject{14008}{<< /D [\MusicPDFPageRef{428} /XYZ 205.644 478.954 null] >>}
\MusicPDFWriteObject{14009}{<< /D [\MusicPDFPageRef{428} /XYZ 258.119 429.375 null] >>}
\MusicPDFWriteObject{14010}{<< /D [\MusicPDFPageRef{428} /XYZ 260.953 392.948 null] >>}
\MusicPDFWriteObject{14011}{<< /D [\MusicPDFPageRef{429} /XYZ 281.728 477.897 null] >>}
\MusicPDFWriteObject{14012}{<< /D [\MusicPDFPageRef{429} /XYZ 199.274 362.565 null] >>}
\MusicPDFWriteObject{14013}{<< /D [\MusicPDFPageRef{539} /XYZ 264.471 451.352 null] >>}
\MusicPDFWriteObject{14014}{<< /D [\MusicPDFPageRef{539} /XYZ 267.96 315.067 null] >>}
\MusicPDFWriteObject{14015}{<< /D [\MusicPDFPageRef{539} /XYZ 282.6 245.178 null] >>}
\MusicPDFWriteObject{14016}{<< /D [\MusicPDFPageRef{540} /XYZ 227.843 582.725 null] >>}
\MusicPDFWriteObject{14017}{<< /D [\MusicPDFPageRef{540} /XYZ 240.739 520.203 null] >>}
\MusicPDFWriteObject{14018}{<< /D [\MusicPDFPageRef{540} /XYZ 279.587 381.051 null] >>}
\MusicPDFWriteObject{14019}{<< /D [\MusicPDFPageRef{540} /XYZ 246.258 172.142 null] >>}
\MusicPDFWriteObject{14020}{<< /D [\MusicPDFPageRef{541} /XYZ 244.585 656.517 null] >>}
\MusicPDFWriteObject{14021}{<< /D [\MusicPDFPageRef{541} /XYZ 270.985 493.795 null] >>}
\MusicPDFWriteObject{14022}{<< /D [\MusicPDFPageRef{541} /XYZ 255.806 313.954 null] >>}
\MusicPDFWriteObject{14023}{<< /D [\MusicPDFPageRef{541} /XYZ 236.656 277.526 null] >>}
\MusicPDFWriteObject{14024}{<< /D [\MusicPDFPageRef{541} /XYZ 229.144 205.576 null] >>}
\MusicPDFWriteObject{14025}{<< /D [\MusicPDFPageRef{542} /XYZ 213.271 603.829 null] >>}
\MusicPDFWriteObject{14026}{<< /D [\MusicPDFPageRef{542} /XYZ 246.057 473.125 null] >>}
\MusicPDFWriteObject{14027}{<< /D [\MusicPDFPageRef{542} /XYZ 262.381 337.827 null] >>}
\MusicPDFWriteObject{14028}{<< /D [\MusicPDFPageRef{542} /XYZ 269.708 268.361 null] >>}
\MusicPDFWriteObject{14029}{<< /D [\MusicPDFPageRef{542} /XYZ 234.686 198.895 null] >>}
\MusicPDFWriteObject{14030}{<< /D [\MusicPDFPageRef{543} /XYZ 274.903 583.39 null] >>}
\MusicPDFWriteObject{14031}{<< /D [\MusicPDFPageRef{545} /XYZ 230.193 659.7 null] >>}
\MusicPDFWriteObject{14032}{<< /D [\MusicPDFPageRef{543} /XYZ 261.467 538.583 null] >>}
\MusicPDFWriteObject{14033}{<< /D [\MusicPDFPageRef{543} /XYZ 269.991 411.972 null] >>}
\MusicPDFWriteObject{14034}{<< /D [\MusicPDFPageRef{543} /XYZ 224.637 277.034 null] >>}
\MusicPDFWriteObject{14035}{<< /D [\MusicPDFPageRef{544} /XYZ 249.024 656.755 null] >>}
\MusicPDFWriteObject{14036}{<< /D [\MusicPDFPageRef{544} /XYZ 269.605 620.084 null] >>}
\MusicPDFWriteObject{14037}{<< /D [\MusicPDFPageRef{544} /XYZ 229.835 560.602 null] >>}
\MusicPDFWriteObject{14038}{<< /D [\MusicPDFPageRef{544} /XYZ 261.542 488.663 null] >>}
\MusicPDFWriteObject{14039}{<< /D [\MusicPDFPageRef{544} /XYZ 163.968 216.524 null] >>}
\MusicPDFWriteObject{14040}{<< /D [\MusicPDFPageRef{545} /XYZ 182.727 344.214 null] >>}
\MusicPDFWriteObject{14041}{<< /D [\MusicPDFPageRef{545} /XYZ 252.969 298.869 null] >>}
\MusicPDFWriteObject{14042}{<< /D [\MusicPDFPageRef{545} /XYZ 276.097 229.241 null] >>}
\MusicPDFWriteObject{14043}{<< /D [\MusicPDFPageRef{545} /XYZ 265.833 153.361 null] >>}
\MusicPDFWriteObject{14044}{<< /D [\MusicPDFPageRef{545} /XYZ 278.548 99.233 null] >>}
\MusicPDFWriteObject{14045}{<< /D [\MusicPDFPageRef{546} /XYZ 274.023 682.619 null] >>}
\MusicPDFWriteObject{14046}{<< /D [\MusicPDFPageRef{546} /XYZ 234.679 431.748 null] >>}
\MusicPDFWriteObject{14047}{<< /D [\MusicPDFPageRef{547} /XYZ 255.145 407.397 null] >>}
\MusicPDFWriteObject{14048}{<< /D [\MusicPDFPageRef{547} /XYZ 257.144 243.521 null] >>}
\MusicPDFWriteObject{14049}{<< /D [\MusicPDFPageRef{547} /XYZ 228.101 117.447 null] >>}
\MusicPDFWriteObject{14050}{<< /D [\MusicPDFPageRef{546} /XYZ 268.753 394.883 null] >>}
\MusicPDFWriteObject{14051}{<< /D [\MusicPDFPageRef{546} /XYZ 262.335 319.003 null] >>}
\MusicPDFWriteObject{14052}{<< /D [\MusicPDFPageRef{546} /XYZ 259.208 275.519 null] >>}
\MusicPDFWriteObject{14053}{<< /D [\MusicPDFPageRef{546} /XYZ 257.284 194.848 null] >>}
\MusicPDFWriteObject{14054}{<< /D [\MusicPDFPageRef{546} /XYZ 249.011 99.233 null] >>}
\MusicPDFWriteObject{14055}{<< /D [\MusicPDFPageRef{547} /XYZ 263.581 683.057 null] >>}
\MusicPDFWriteObject{14056}{<< /D [\MusicPDFPageRef{547} /XYZ 259.836 627.873 null] >>}
\MusicPDFWriteObject{14057}{<< /D [\MusicPDFPageRef{547} /XYZ 252.992 591.446 null] >>}
\MusicPDFWriteObject{14058}{<< /D [\MusicPDFPageRef{548} /XYZ 248.058 579.683 null] >>}
\MusicPDFWriteObject{14059}{<< /D [\MusicPDFPageRef{549} /XYZ 200.861 550.742 null] >>}
\MusicPDFWriteObject{14060}{<< /D [\MusicPDFPageRef{549} /XYZ 270.215 480.731 null] >>}
\MusicPDFWriteObject{14061}{<< /D [\MusicPDFPageRef{549} /XYZ 244.708 411.909 null] >>}
\MusicPDFWriteObject{14062}{<< /D [\MusicPDFPageRef{549} /XYZ 176.353 353.707 null] >>}
\MusicPDFWriteObject{14063}{<< /D [\MusicPDFPageRef{549} /XYZ 201.641 311.025 null] >>}
\MusicPDFWriteObject{14064}{<< /D [\MusicPDFPageRef{549} /XYZ 274.296 196.679 null] >>}
\MusicPDFWriteObject{14065}{<< /D [\MusicPDFPageRef{549} /XYZ 150.828 117.447 null] >>}
\MusicPDFWriteObject{14066}{<< /D [\MusicPDFPageRef{550} /XYZ 277.517 618.682 null] >>}
\MusicPDFWriteObject{14067}{<< /D [\MusicPDFPageRef{550} /XYZ 275.559 555.953 null] >>}
\MusicPDFWriteObject{14068}{<< /D [\MusicPDFPageRef{550} /XYZ 247.912 335.472 null] >>}
\MusicPDFWriteObject{14069}{<< /D [\MusicPDFPageRef{548} /XYZ 197.78 543.256 null] >>}
\MusicPDFWriteObject{14070}{<< /D [\MusicPDFPageRef{550} /XYZ 210.275 152.045 null] >>}
\MusicPDFWriteObject{14071}{<< /D [\MusicPDFPageRef{551} /XYZ 187.409 675.22 null] >>}
\MusicPDFWriteObject{14072}{<< /D [\MusicPDFPageRef{552} /XYZ 159.282 638.252 null] >>}
\MusicPDFWriteObject{14073}{<< /D [\MusicPDFPageRef{552} /XYZ 181.631 601.149 null] >>}
\MusicPDFWriteObject{14074}{<< /D [\MusicPDFPageRef{552} /XYZ 177.162 551.334 null] >>}
\MusicPDFWriteObject{14075}{<< /D [\MusicPDFPageRef{552} /XYZ 263.904 437.337 null] >>}
\MusicPDFWriteObject{14076}{<< /D [\MusicPDFPageRef{552} /XYZ 273.98 400.91 null] >>}
\MusicPDFWriteObject{14077}{<< /D [\MusicPDFPageRef{552} /XYZ 276.094 349.55 null] >>}
\MusicPDFWriteObject{14078}{<< /D [\MusicPDFPageRef{552} /XYZ 177.162 246.404 null] >>}
\MusicPDFWriteObject{14079}{<< /D [\MusicPDFPageRef{553} /XYZ 246.372 482.508 null] >>}
\MusicPDFWriteObject{14080}{<< /D [\MusicPDFPageRef{548} /XYZ 252.348 482.36 null] >>}
\MusicPDFWriteObject{14081}{<< /D [\MusicPDFPageRef{548} /XYZ 202.531 445.932 null] >>}
\MusicPDFWriteObject{14082}{<< /D [\MusicPDFPageRef{548} /XYZ 200.821 384.793 null] >>}
\MusicPDFWriteObject{14083}{<< /D [\MusicPDFPageRef{548} /XYZ 271.03 256.216 null] >>}
\MusicPDFWriteObject{14084}{<< /D [\MusicPDFPageRef{548} /XYZ 268.875 99.233 null] >>}
\MusicPDFWriteObject{14085}{<< /D [\MusicPDFPageRef{549} /XYZ 255.177 661.218 null] >>}
\MusicPDFWriteObject{14086}{<< /D [\MusicPDFPageRef{549} /XYZ 232.442 594.627 null] >>}
\MusicPDFWriteObject{14087}{<< /D [\MusicPDFPageRef{554} /XYZ 242.83 585.187 null] >>}
\MusicPDFWriteObject{14088}{<< /D [\MusicPDFPageRef{555} /XYZ 268.923 464.087 null] >>}
\MusicPDFWriteObject{14089}{<< /D [\MusicPDFPageRef{555} /XYZ 273.566 413.579 null] >>}
\MusicPDFWriteObject{14090}{<< /D [\MusicPDFPageRef{554} /XYZ 262.983 492.923 null] >>}
\MusicPDFWriteObject{14091}{<< /D [\MusicPDFPageRef{554} /XYZ 262.443 417.043 null] >>}
\MusicPDFWriteObject{14092}{<< /D [\MusicPDFPageRef{554} /XYZ 113.108 182.753 null] >>}
\MusicPDFWriteObject{14093}{<< /D [\MusicPDFPageRef{554} /XYZ 168.36 136.669 null] >>}
\MusicPDFWriteObject{14094}{<< /D [\MusicPDFPageRef{554} /XYZ 257.316 99.233 null] >>}
\MusicPDFWriteObject{14095}{<< /D [\MusicPDFPageRef{555} /XYZ 258.235 683.057 null] >>}
\MusicPDFWriteObject{14096}{<< /D [\MusicPDFPageRef{555} /XYZ 233.306 631.145 null] >>}
\MusicPDFWriteObject{14097}{<< /D [\MusicPDFPageRef{555} /XYZ 255.015 501.223 null] >>}
\MusicPDFWriteObject{14098}{<< /D [\MusicPDFPageRef{555} /XYZ 277.987 143.748 null] >>}
\MusicPDFWriteObject{14099}{<< /D [\MusicPDFPageRef{559} /XYZ 245.903 597.291 null] >>}
\MusicPDFWriteObject{14100}{<< /D [\MusicPDFPageRef{559} /XYZ 166.907 505.689 null] >>}
\MusicPDFWriteObject{14101}{<< /D [\MusicPDFPageRef{559} /XYZ 174.461 460.651 null] >>}
\MusicPDFWriteObject{14102}{<< /D [\MusicPDFPageRef{559} /XYZ 262.503 362.317 null] >>}
\MusicPDFWriteObject{14103}{<< /D [\MusicPDFPageRef{556} /XYZ 206.714 484.855 null] >>}
\MusicPDFWriteObject{14104}{<< /D [\MusicPDFPageRef{556} /XYZ 252.358 380.491 null] >>}
\MusicPDFWriteObject{14105}{<< /D [\MusicPDFPageRef{556} /XYZ 265.948 183.572 null] >>}
\MusicPDFWriteObject{14106}{<< /D [\MusicPDFPageRef{557} /XYZ 190.329 211.281 null] >>}
\MusicPDFWriteObject{14107}{<< /D [\MusicPDFPageRef{557} /XYZ 251.933 132.582 null] >>}
\MusicPDFWriteObject{14108}{<< /D [\MusicPDFPageRef{558} /XYZ 249.577 383.244 null] >>}
\MusicPDFWriteObject{14109}{<< /D [\MusicPDFPageRef{558} /XYZ 236.817 240.928 null] >>}
\MusicPDFWriteObject{14110}{<< /D [\MusicPDFPageRef{559} /XYZ 274.755 683.057 null] >>}
\MusicPDFWriteObject{14111}{<< /D [\MusicPDFPageRef{560} /XYZ 249.362 678.255 null] >>}
\MusicPDFWriteObject{14112}{<< /D [\MusicPDFPageRef{560} /XYZ 244.13 558.56 null] >>}
\MusicPDFWriteObject{14113}{<< /D [\MusicPDFPageRef{560} /XYZ 271.498 469.11 null] >>}
\MusicPDFWriteObject{14114}{<< /D [\MusicPDFPageRef{560} /XYZ 124.199 352.045 null] >>}
\MusicPDFWriteObject{14115}{<< /D [\MusicPDFPageRef{560} /XYZ 280.087 192.436 null] >>}
\MusicPDFWriteObject{14116}{<< /D [\MusicPDFPageRef{561} /XYZ 282.627 313.089 null] >>}
\MusicPDFWriteObject{14117}{<< /D [\MusicPDFPageRef{561} /XYZ 281.545 241.067 null] >>}
\MusicPDFWriteObject{14118}{<< /D [\MusicPDFPageRef{562} /XYZ 222.471 241.036 null] >>}
\MusicPDFWriteObject{14119}{<< /D [\MusicPDFPageRef{562} /XYZ 262.962 203.812 null] >>}
\MusicPDFWriteObject{14120}{<< /D [\MusicPDFPageRef{562} /XYZ 237.419 139.36 null] >>}
\MusicPDFWriteObject{14121}{<< /D [\MusicPDFPageRef{563} /XYZ 235.785 602.387 null] >>}
\MusicPDFWriteObject{14122}{<< /D [\MusicPDFPageRef{563} /XYZ 211.05 513.671 null] >>}
\MusicPDFWriteObject{14123}{<< /D [\MusicPDFPageRef{563} /XYZ 72 454.048 null] >>}
\MusicPDFWriteObject{14124}{<< /D [\MusicPDFPageRef{563} /XYZ 209.282 278.964 null] >>}
\MusicPDFWriteObject{14125}{<< /D [\MusicPDFPageRef{564} /XYZ 177.747 419.808 null] >>}
\MusicPDFWriteObject{14126}{<< /D [\MusicPDFPageRef{565} /XYZ 152.405 617.303 null] >>}
\MusicPDFWriteObject{14127}{<< /D [\MusicPDFPageRef{566} /XYZ 241.365 579.895 null] >>}
\MusicPDFWriteObject{14128}{<< /D [\MusicPDFPageRef{567} /XYZ 120.677 210.594 null] >>}
\MusicPDFWriteObject{14129}{<< /D [\MusicPDFPageRef{567} /XYZ 150.886 160.133 null] >>}
\MusicPDFWriteObject{14130}{<< /D [\MusicPDFPageRef{568} /XYZ 243.984 452.561 null] >>}
\MusicPDFWriteObject{14131}{<< /D [\MusicPDFPageRef{568} /XYZ 173.179 371.664 null] >>}
\MusicPDFWriteObject{14132}{<< /D [\MusicPDFPageRef{568} /XYZ 198.312 273.552 null] >>}
\MusicPDFWriteObject{14133}{<< /D [\MusicPDFPageRef{566} /XYZ 226.057 415.242 null] >>}
\MusicPDFWriteObject{14134}{<< /D [\MusicPDFPageRef{566} /XYZ 217.572 226.76 null] >>}
\MusicPDFWriteObject{14135}{<< /D [\MusicPDFPageRef{566} /XYZ 231.765 184.127 null] >>}
\MusicPDFWriteObject{14136}{<< /D [\MusicPDFPageRef{567} /XYZ 72 688.85 null] >>}
\MusicPDFWriteObject{14137}{<< /D [\MusicPDFPageRef{567} /XYZ 72 535.949 null] >>}
\MusicPDFWriteObject{14138}{<< /D [\MusicPDFPageRef{567} /XYZ 72 482.777 null] >>}
\MusicPDFWriteObject{14139}{<< /D [\MusicPDFPageRef{567} /XYZ 81.223 304.317 null] >>}
\MusicPDFWriteObject{14140}{<< /D [\MusicPDFPageRef{567} /XYZ 158.133 257.165 null] >>}
\MusicPDFWriteObject{14141}{<< /D [\MusicPDFPageRef{569} /XYZ 251.339 314.723 null] >>}
\MusicPDFWriteObject{14142}{<< /D [\MusicPDFPageRef{574} /XYZ 201.641 235.541 null] >>}
\MusicPDFWriteObject{14143}{<< /D [\MusicPDFPageRef{574} /XYZ 227.281 143.118 null] >>}
\MusicPDFWriteObject{14144}{<< /D [\MusicPDFPageRef{574} /XYZ 186.173 99.233 null] >>}
\MusicPDFWriteObject{14145}{<< /D [\MusicPDFPageRef{575} /XYZ 273.797 683.057 null] >>}
\MusicPDFWriteObject{14146}{<< /D [\MusicPDFPageRef{575} /XYZ 155.749 646.629 null] >>}
\MusicPDFWriteObject{14147}{<< /D [\MusicPDFPageRef{575} /XYZ 149.878 570.749 null] >>}
\MusicPDFWriteObject{14148}{<< /D [\MusicPDFPageRef{575} /XYZ 257.316 534.322 null] >>}
\MusicPDFWriteObject{14149}{<< /D [\MusicPDFPageRef{575} /XYZ 258.235 497.894 null] >>}
\MusicPDFWriteObject{14150}{<< /D [\MusicPDFPageRef{575} /XYZ 219.469 445.982 null] >>}
\MusicPDFWriteObject{14151}{<< /D [\MusicPDFPageRef{576} /XYZ 266.475 603.859 null] >>}
\MusicPDFWriteObject{14152}{<< /D [\MusicPDFPageRef{569} /XYZ 268.861 177.105 null] >>}
\MusicPDFWriteObject{14153}{<< /D [\MusicPDFPageRef{576} /XYZ 152.966 561.531 null] >>}
\MusicPDFWriteObject{14154}{<< /D [\MusicPDFPageRef{576} /XYZ 264.88 275.084 null] >>}
\MusicPDFWriteObject{14155}{<< /D [\MusicPDFPageRef{577} /XYZ 220.69 641.573 null] >>}
\MusicPDFWriteObject{14156}{<< /D [\MusicPDFPageRef{577} /XYZ 264.668 494.921 null] >>}
\MusicPDFWriteObject{14157}{<< /D [\MusicPDFPageRef{577} /XYZ 222.8 204.864 null] >>}
\MusicPDFWriteObject{14158}{<< /D [\MusicPDFPageRef{572} /XYZ 233.226 442.155 null] >>}
\MusicPDFWriteObject{14159}{<< /D [\MusicPDFPageRef{573} /XYZ 129.001 638.656 null] >>}
\MusicPDFWriteObject{14160}{<< /D [\MusicPDFPageRef{573} /XYZ 245.225 544.098 null] >>}
\MusicPDFWriteObject{14161}{<< /D [\MusicPDFPageRef{573} /XYZ 245.156 507.67 null] >>}
\MusicPDFWriteObject{14162}{<< /D [\MusicPDFPageRef{574} /XYZ 176.353 407.349 null] >>}
\MusicPDFWriteObject{14163}{<< /D [\MusicPDFPageRef{574} /XYZ 227.281 349.02 null] >>}
\MusicPDFWriteObject{14164}{<< /D [\MusicPDFPageRef{574} /XYZ 198.627 291.374 null] >>}
\MusicPDFWriteObject{14165}{<< /D [\MusicPDFPageRef{579} /XYZ 254.155 637.518 null] >>}
\MusicPDFWriteObject{14166}{<< /D [\MusicPDFPageRef{581} /XYZ 187.639 419.058 null] >>}
\MusicPDFWriteObject{14167}{<< /D [\MusicPDFPageRef{581} /XYZ 166.078 265.738 null] >>}
\MusicPDFWriteObject{14168}{<< /D [\MusicPDFPageRef{581} /XYZ 238.912 184.467 null] >>}
\MusicPDFWriteObject{14169}{<< /D [\MusicPDFPageRef{581} /XYZ 270.347 134.889 null] >>}
\MusicPDFWriteObject{14170}{<< /D [\MusicPDFPageRef{582} /XYZ 180.95 683.057 null] >>}
\MusicPDFWriteObject{14171}{<< /D [\MusicPDFPageRef{582} /XYZ 72 636.766 null] >>}
\MusicPDFWriteObject{14172}{<< /D [\MusicPDFPageRef{582} /XYZ 197.28 435.634 null] >>}
\MusicPDFWriteObject{14173}{<< /D [\MusicPDFPageRef{582} /XYZ 206.674 372.576 null] >>}
\MusicPDFWriteObject{14174}{<< /D [\MusicPDFPageRef{582} /XYZ 209.04 328.812 null] >>}
\MusicPDFWriteObject{14175}{<< /D [\MusicPDFPageRef{579} /XYZ 246.289 223.744 null] >>}
\MusicPDFWriteObject{14176}{<< /D [\MusicPDFPageRef{579} /XYZ 144.944 162.143 null] >>}
\MusicPDFWriteObject{14177}{<< /D [\MusicPDFPageRef{580} /XYZ 75.239 492.045 null] >>}
\MusicPDFWriteObject{14178}{<< /D [\MusicPDFPageRef{580} /XYZ 246.289 339.94 null] >>}
\MusicPDFWriteObject{14179}{<< /D [\MusicPDFPageRef{580} /XYZ 211.665 294.746 null] >>}
\MusicPDFWriteObject{14180}{<< /D [\MusicPDFPageRef{580} /XYZ 200.888 183.513 null] >>}
\MusicPDFWriteObject{14181}{<< /D [\MusicPDFPageRef{581} /XYZ 197.28 569.881 null] >>}
\MusicPDFWriteObject{14182}{<< /D [\MusicPDFPageRef{581} /XYZ 142.67 466.807 null] >>}
\MusicPDFWriteObject{14183}{<< /D [\MusicPDFPageRef{582} /XYZ 222.312 149.313 null] >>}
\MusicPDFWriteObject{14184}{<< /D [\MusicPDFPageRef{584} /XYZ 269.449 511.359 null] >>}
\MusicPDFWriteObject{14185}{<< /D [\MusicPDFPageRef{584} /XYZ 258.842 448.63 null] >>}
\MusicPDFWriteObject{14186}{<< /D [\MusicPDFPageRef{584} /XYZ 190.102 408.32 null] >>}
\MusicPDFWriteObject{14187}{<< /D [\MusicPDFPageRef{584} /XYZ 190.941 365.859 null] >>}
\MusicPDFWriteObject{14188}{<< /D [\MusicPDFPageRef{585} /XYZ 205.555 332.172 null] >>}
\MusicPDFWriteObject{14189}{<< /D [\MusicPDFPageRef{585} /XYZ 206.382 290.292 null] >>}
\MusicPDFWriteObject{14190}{<< /D [\MusicPDFPageRef{585} /XYZ 72 210.119 null] >>}
\MusicPDFWriteObject{14191}{<< /D [\MusicPDFPageRef{585} /XYZ 172.084 155.062 null] >>}
\MusicPDFWriteObject{14192}{<< /D [\MusicPDFPageRef{583} /XYZ 222.414 676.942 null] >>}
\MusicPDFWriteObject{14193}{<< /D [\MusicPDFPageRef{583} /XYZ 210.331 510.427 null] >>}
\MusicPDFWriteObject{14194}{<< /D [\MusicPDFPageRef{583} /XYZ 223.51 465.815 null] >>}
\MusicPDFWriteObject{14195}{<< /D [\MusicPDFPageRef{583} /XYZ 228.721 358.792 null] >>}
\MusicPDFWriteObject{14196}{<< /D [\MusicPDFPageRef{583} /XYZ 182.967 256.611 null] >>}
\MusicPDFWriteObject{14197}{<< /D [\MusicPDFPageRef{584} /XYZ 246.717 683.057 null] >>}
\MusicPDFWriteObject{14198}{<< /D [\MusicPDFPageRef{584} /XYZ 162.978 629.25 null] >>}
\MusicPDFWriteObject{14199}{<< /D [\MusicPDFPageRef{584} /XYZ 268.795 574.089 null] >>}
\MusicPDFWriteObject{14200}{<< /D [\MusicPDFPageRef{589} /XYZ 273.827 361.411 null] >>}
\MusicPDFWriteObject{14201}{<< /D [\MusicPDFPageRef{590} /XYZ 248.054 174.566 null] >>}
\MusicPDFWriteObject{14202}{<< /D [\MusicPDFPageRef{590} /XYZ 245.34 105.297 null] >>}
\MusicPDFWriteObject{14203}{<< /D [\MusicPDFPageRef{591} /XYZ 255.253 635.631 null] >>}
\MusicPDFWriteObject{14204}{<< /D [\MusicPDFPageRef{591} /XYZ 255.177 486.238 null] >>}
\MusicPDFWriteObject{14205}{<< /D [\MusicPDFPageRef{591} /XYZ 202.997 414.556 null] >>}
\MusicPDFWriteObject{14206}{<< /D [\MusicPDFPageRef{591} /XYZ 176.937 268.021 null] >>}
\MusicPDFWriteObject{14207}{<< /D [\MusicPDFPageRef{589} /XYZ 184.89 304.878 null] >>}
\MusicPDFWriteObject{14208}{<< /D [\MusicPDFPageRef{589} /XYZ 245.204 268.45 null] >>}
\MusicPDFWriteObject{14209}{<< /D [\MusicPDFPageRef{589} /XYZ 250.8 192.789 null] >>}
\MusicPDFWriteObject{14210}{<< /D [\MusicPDFPageRef{590} /XYZ 154.051 547.061 null] >>}
\MusicPDFWriteObject{14211}{<< /D [\MusicPDFPageRef{590} /XYZ 222.539 503.604 null] >>}
\MusicPDFWriteObject{14212}{<< /D [\MusicPDFPageRef{590} /XYZ 179.622 338.288 null] >>}
\MusicPDFWriteObject{14213}{<< /D [\MusicPDFPageRef{590} /XYZ 224.411 293.46 null] >>}
\MusicPDFWriteObject{14214}{<< /D [\MusicPDFPageRef{590} /XYZ 233.886 211.812 null] >>}
\MusicPDFWriteObject{14215}{<< /D [\MusicPDFPageRef{593} /XYZ 277.511 268.3 null] >>}
\MusicPDFWriteObject{14216}{<< /D [\MusicPDFPageRef{595} /XYZ 250.402 440.435 null] >>}
\MusicPDFWriteObject{14217}{<< /D [\MusicPDFPageRef{595} /XYZ 173.542 404.008 null] >>}
\MusicPDFWriteObject{14218}{<< /D [\MusicPDFPageRef{595} /XYZ 190.222 351.1 null] >>}
\MusicPDFWriteObject{14219}{<< /D [\MusicPDFPageRef{595} /XYZ 174.885 299.494 null] >>}
\MusicPDFWriteObject{14220}{<< /D [\MusicPDFPageRef{596} /XYZ 199.289 655.598 null] >>}
\MusicPDFWriteObject{14221}{<< /D [\MusicPDFPageRef{596} /XYZ 206.277 607.963 null] >>}
\MusicPDFWriteObject{14222}{<< /D [\MusicPDFPageRef{596} /XYZ 164.701 310.638 null] >>}
\MusicPDFWriteObject{14223}{<< /D [\MusicPDFPageRef{596} /XYZ 196.068 109.52 null] >>}
\MusicPDFWriteObject{14224}{<< /D [\MusicPDFPageRef{597} /XYZ 278.152 678.621 null] >>}
\MusicPDFWriteObject{14225}{<< /D [\MusicPDFPageRef{597} /XYZ 159.064 461.767 null] >>}
\MusicPDFWriteObject{14226}{<< /D [\MusicPDFPageRef{594} /XYZ 250.35 579.423 null] >>}
\MusicPDFWriteObject{14227}{<< /D [\MusicPDFPageRef{598} /XYZ 267.813 656.411 null] >>}
\MusicPDFWriteObject{14228}{<< /D [\MusicPDFPageRef{598} /XYZ 205.508 577.271 null] >>}
\MusicPDFWriteObject{14229}{<< /D [\MusicPDFPageRef{598} /XYZ 72 472.364 null] >>}
\MusicPDFWriteObject{14230}{<< /D [\MusicPDFPageRef{598} /XYZ 182.336 425.828 null] >>}
\MusicPDFWriteObject{14231}{<< /D [\MusicPDFPageRef{598} /XYZ 159.333 379.234 null] >>}
\MusicPDFWriteObject{14232}{<< /D [\MusicPDFPageRef{598} /XYZ 158.494 305.354 null] >>}
\MusicPDFWriteObject{14233}{<< /D [\MusicPDFPageRef{598} /XYZ 257.587 213.91 null] >>}
\MusicPDFWriteObject{14234}{<< /D [\MusicPDFPageRef{598} /XYZ 250.402 130.597 null] >>}
\MusicPDFWriteObject{14235}{<< /D [\MusicPDFPageRef{599} /XYZ 126.776 672.519 null] >>}
\MusicPDFWriteObject{14236}{<< /D [\MusicPDFPageRef{599} /XYZ 252.056 575.713 null] >>}
\MusicPDFWriteObject{14237}{<< /D [\MusicPDFPageRef{594} /XYZ 258.97 460.682 null] >>}
\MusicPDFWriteObject{14238}{<< /D [\MusicPDFPageRef{599} /XYZ 212.998 528.309 null] >>}
\MusicPDFWriteObject{14239}{<< /D [\MusicPDFPageRef{599} /XYZ 196.959 414.332 null] >>}
\MusicPDFWriteObject{14240}{<< /D [\MusicPDFPageRef{599} /XYZ 72 356.828 null] >>}
\MusicPDFWriteObject{14241}{<< /D [\MusicPDFPageRef{599} /XYZ 72 327.222 null] >>}
\MusicPDFWriteObject{14242}{<< /D [\MusicPDFPageRef{599} /XYZ 210.259 277.696 null] >>}
\MusicPDFWriteObject{14243}{<< /D [\MusicPDFPageRef{599} /XYZ 250.402 236.045 null] >>}
\MusicPDFWriteObject{14244}{<< /D [\MusicPDFPageRef{599} /XYZ 215.258 199.618 null] >>}
\MusicPDFWriteObject{14245}{<< /D [\MusicPDFPageRef{599} /XYZ 160.595 159.307 null] >>}
\MusicPDFWriteObject{14246}{<< /D [\MusicPDFPageRef{599} /XYZ 144.428 107.701 null] >>}
\MusicPDFWriteObject{14247}{<< /D [\MusicPDFPageRef{600} /XYZ 265.538 570.232 null] >>}
\MusicPDFWriteObject{14248}{<< /D [\MusicPDFPageRef{594} /XYZ 250.402 376.865 null] >>}
\MusicPDFWriteObject{14249}{<< /D [\MusicPDFPageRef{594} /XYZ 172.671 291 null] >>}
\MusicPDFWriteObject{14250}{<< /D [\MusicPDFPageRef{594} /XYZ 249.521 224.76 null] >>}
\MusicPDFWriteObject{14251}{<< /D [\MusicPDFPageRef{595} /XYZ 148.022 628.281 null] >>}
\MusicPDFWriteObject{14252}{<< /D [\MusicPDFPageRef{595} /XYZ 196.075 542.563 null] >>}
\MusicPDFWriteObject{14253}{<< /D [\MusicPDFPageRef{595} /XYZ 179.784 496.147 null] >>}
\MusicPDFWriteObject{14254}{<< /D [\MusicPDFPageRef{600} /XYZ 266.084 290.416 null] >>}
\MusicPDFWriteObject{14255}{<< /D [\MusicPDFPageRef{602} /XYZ 233.777 182.364 null] >>}
\MusicPDFWriteObject{14256}{<< /D [\MusicPDFPageRef{602} /XYZ 248.186 135.992 null] >>}
\MusicPDFWriteObject{14257}{<< /D [\MusicPDFPageRef{603} /XYZ 72 654.984 null] >>}
\MusicPDFWriteObject{14258}{<< /D [\MusicPDFPageRef{603} /XYZ 230.854 546.29 null] >>}
\MusicPDFWriteObject{14259}{<< /D [\MusicPDFPageRef{603} /XYZ 244.117 184.464 null] >>}
\MusicPDFWriteObject{14260}{<< /D [\MusicPDFPageRef{601} /XYZ 265.365 490.498 null] >>}
\MusicPDFWriteObject{14261}{<< /D [\MusicPDFPageRef{601} /XYZ 198.342 328.61 null] >>}
\MusicPDFWriteObject{14262}{<< /D [\MusicPDFPageRef{601} /XYZ 151.383 105.155 null] >>}
\MusicPDFWriteObject{14263}{<< /D [\MusicPDFPageRef{602} /XYZ 165.791 675.728 null] >>}
\MusicPDFWriteObject{14264}{<< /D [\MusicPDFPageRef{602} /XYZ 222.979 614.649 null] >>}
\MusicPDFWriteObject{14265}{<< /D [\MusicPDFPageRef{602} /XYZ 237.388 566.72 null] >>}
\MusicPDFWriteObject{14266}{<< /D [\MusicPDFPageRef{602} /XYZ 236.096 497.757 null] >>}
\MusicPDFWriteObject{14267}{<< /D [\MusicPDFPageRef{602} /XYZ 166.479 436.318 null] >>}
\MusicPDFWriteObject{14268}{<< /D [\MusicPDFPageRef{604} /XYZ 256.681 272.094 null] >>}
\MusicPDFWriteObject{14269}{<< /D [\MusicPDFPageRef{605} /XYZ 163.252 677.333 null] >>}
\MusicPDFWriteObject{14270}{<< /D [\MusicPDFPageRef{605} /XYZ 148.934 555.989 null] >>}
\MusicPDFWriteObject{14271}{<< /D [\MusicPDFPageRef{606} /XYZ 240.337 682.26 null] >>}
\MusicPDFWriteObject{14272}{<< /D [\MusicPDFPageRef{606} /XYZ 126.505 615.441 null] >>}
\MusicPDFWriteObject{14273}{<< /D [\MusicPDFPageRef{606} /XYZ 186.94 490.885 null] >>}
\MusicPDFWriteObject{14274}{<< /D [\MusicPDFPageRef{606} /XYZ 148.396 437.219 null] >>}
\MusicPDFWriteObject{14275}{<< /D [\MusicPDFPageRef{606} /XYZ 181.37 349.613 null] >>}
\MusicPDFWriteObject{14276}{<< /D [\MusicPDFPageRef{607} /XYZ 235.986 388.883 null] >>}
\MusicPDFWriteObject{14277}{<< /D [\MusicPDFPageRef{609} /XYZ 259.55 596.818 null] >>}
\MusicPDFWriteObject{14278}{<< /D [\MusicPDFPageRef{609} /XYZ 248.339 498.225 null] >>}
\MusicPDFWriteObject{14279}{<< /D [\MusicPDFPageRef{609} /XYZ 249.402 446.607 null] >>}
\MusicPDFWriteObject{14280}{<< /D [\MusicPDFPageRef{609} /XYZ 273.994 400.778 null] >>}
\MusicPDFWriteObject{14281}{<< /D [\MusicPDFPageRef{609} /XYZ 137.39 135.6 null] >>}
\MusicPDFWriteObject{14282}{<< /D [\MusicPDFPageRef{610} /XYZ 117.905 569.651 null] >>}
\MusicPDFWriteObject{14283}{<< /D [\MusicPDFPageRef{610} /XYZ 131.768 520.063 null] >>}
\MusicPDFWriteObject{14284}{<< /D [\MusicPDFPageRef{610} /XYZ 132.965 474.235 null] >>}
\MusicPDFWriteObject{14285}{<< /D [\MusicPDFPageRef{610} /XYZ 235.272 151.748 null] >>}
\MusicPDFWriteObject{14286}{<< /D [\MusicPDFPageRef{611} /XYZ 258.97 330.008 null] >>}
\MusicPDFWriteObject{14287}{<< /D [\MusicPDFPageRef{607} /XYZ 242.414 349.015 null] >>}
\MusicPDFWriteObject{14288}{<< /D [\MusicPDFPageRef{611} /XYZ 90.534 260.954 null] >>}
\MusicPDFWriteObject{14289}{<< /D [\MusicPDFPageRef{612} /XYZ 184.657 630.816 null] >>}
\MusicPDFWriteObject{14290}{<< /D [\MusicPDFPageRef{612} /XYZ 239.701 560.052 null] >>}
\MusicPDFWriteObject{14291}{<< /D [\MusicPDFPageRef{612} /XYZ 240.901 511.184 null] >>}
\MusicPDFWriteObject{14292}{<< /D [\MusicPDFPageRef{612} /XYZ 227.28 321.071 null] >>}
\MusicPDFWriteObject{14293}{<< /D [\MusicPDFPageRef{612} /XYZ 224.166 270.317 null] >>}
\MusicPDFWriteObject{14294}{<< /D [\MusicPDFPageRef{612} /XYZ 259.307 224.362 null] >>}
\MusicPDFWriteObject{14295}{<< /D [\MusicPDFPageRef{612} /XYZ 227.267 171.287 null] >>}
\MusicPDFWriteObject{14296}{<< /D [\MusicPDFPageRef{612} /XYZ 158.96 102.912 null] >>}
\MusicPDFWriteObject{14297}{<< /D [\MusicPDFPageRef{613} /XYZ 165.261 659.699 null] >>}
\MusicPDFWriteObject{14298}{<< /D [\MusicPDFPageRef{607} /XYZ 277.027 307.304 null] >>}
\MusicPDFWriteObject{14299}{<< /D [\MusicPDFPageRef{613} /XYZ 135.004 585.836 null] >>}
\MusicPDFWriteObject{14300}{<< /D [\MusicPDFPageRef{613} /XYZ 240.765 446.205 null] >>}
\MusicPDFWriteObject{14301}{<< /D [\MusicPDFPageRef{613} /XYZ 151.461 351.542 null] >>}
\MusicPDFWriteObject{14302}{<< /D [\MusicPDFPageRef{613} /XYZ 165.284 208.156 null] >>}
\MusicPDFWriteObject{14303}{<< /D [\MusicPDFPageRef{614} /XYZ 158.741 635.674 null] >>}
\MusicPDFWriteObject{14304}{<< /D [\MusicPDFPageRef{614} /XYZ 176.93 583.736 null] >>}
\MusicPDFWriteObject{14305}{<< /D [\MusicPDFPageRef{608} /XYZ 143.598 653.047 null] >>}
\MusicPDFWriteObject{14306}{<< /D [\MusicPDFPageRef{608} /XYZ 236.842 605.895 null] >>}
\MusicPDFWriteObject{14307}{<< /D [\MusicPDFPageRef{608} /XYZ 242.414 501.791 null] >>}
\MusicPDFWriteObject{14308}{<< /D [\MusicPDFPageRef{608} /XYZ 148.832 428.403 null] >>}
\MusicPDFWriteObject{14309}{<< /D [\MusicPDFPageRef{608} /XYZ 273.927 245.27 null] >>}
\MusicPDFWriteObject{14310}{<< /D [\MusicPDFPageRef{608} /XYZ 251.749 205.405 null] >>}
\MusicPDFWriteObject{14311}{<< /D [\MusicPDFPageRef{614} /XYZ 274.451 442.25 null] >>}
\MusicPDFWriteObject{14312}{<< /D [\MusicPDFPageRef{615} /XYZ 159.213 249.85 null] >>}
\MusicPDFWriteObject{14313}{<< /D [\MusicPDFPageRef{615} /XYZ 167.607 201.098 null] >>}
\MusicPDFWriteObject{14314}{<< /D [\MusicPDFPageRef{615} /XYZ 238.262 128.705 null] >>}
\MusicPDFWriteObject{14315}{<< /D [\MusicPDFPageRef{616} /XYZ 195.005 649.436 null] >>}
\MusicPDFWriteObject{14316}{<< /D [\MusicPDFPageRef{616} /XYZ 196.389 581.178 null] >>}
\MusicPDFWriteObject{14317}{<< /D [\MusicPDFPageRef{616} /XYZ 175.479 532.42 null] >>}
\MusicPDFWriteObject{14318}{<< /D [\MusicPDFPageRef{616} /XYZ 159.701 419.76 null] >>}
\MusicPDFWriteObject{14319}{<< /D [\MusicPDFPageRef{616} /XYZ 72 245.016 null] >>}
\MusicPDFWriteObject{14320}{<< /D [\MusicPDFPageRef{614} /XYZ 112.693 235.224 null] >>}
\MusicPDFWriteObject{14321}{<< /D [\MusicPDFPageRef{614} /XYZ 234.662 168.421 null] >>}
\MusicPDFWriteObject{14322}{<< /D [\MusicPDFPageRef{615} /XYZ 72 688.851 null] >>}
\MusicPDFWriteObject{14323}{<< /D [\MusicPDFPageRef{615} /XYZ 253.825 575.188 null] >>}
\MusicPDFWriteObject{14324}{<< /D [\MusicPDFPageRef{615} /XYZ 252.442 514.313 null] >>}
\MusicPDFWriteObject{14325}{<< /D [\MusicPDFPageRef{615} /XYZ 160.151 467.2 null] >>}
\MusicPDFWriteObject{14326}{<< /D [\MusicPDFPageRef{615} /XYZ 222.21 364.2 null] >>}
\MusicPDFWriteObject{14327}{<< /D [\MusicPDFPageRef{615} /XYZ 134.742 314.413 null] >>}
\MusicPDFWriteObject{14328}{<< /D [\MusicPDFPageRef{617} /XYZ 253.324 479.261 null] >>}
\MusicPDFWriteObject{14329}{<< /D [\MusicPDFPageRef{618} /XYZ 72 257.61 null] >>}
\MusicPDFWriteObject{14330}{<< /D [\MusicPDFPageRef{618} /XYZ 248.875 220.001 null] >>}
\MusicPDFWriteObject{14331}{<< /D [\MusicPDFPageRef{619} /XYZ 273.509 630.404 null] >>}
\MusicPDFWriteObject{14332}{<< /D [\MusicPDFPageRef{619} /XYZ 151.473 570.596 null] >>}
\MusicPDFWriteObject{14333}{<< /D [\MusicPDFPageRef{619} /XYZ 130.412 505.476 null] >>}
\MusicPDFWriteObject{14334}{<< /D [\MusicPDFPageRef{617} /XYZ 271.83 272.855 null] >>}
\MusicPDFWriteObject{14335}{<< /D [\MusicPDFPageRef{617} /XYZ 249.854 210.868 null] >>}
\MusicPDFWriteObject{14336}{<< /D [\MusicPDFPageRef{617} /XYZ 180.775 165.232 null] >>}
\MusicPDFWriteObject{14337}{<< /D [\MusicPDFPageRef{617} /XYZ 204.87 106.655 null] >>}
\MusicPDFWriteObject{14338}{<< /D [\MusicPDFPageRef{618} /XYZ 202.996 623.957 null] >>}
\MusicPDFWriteObject{14339}{<< /D [\MusicPDFPageRef{618} /XYZ 144.53 568.127 null] >>}
\MusicPDFWriteObject{14340}{<< /D [\MusicPDFPageRef{618} /XYZ 89.181 519.065 null] >>}
\MusicPDFWriteObject{14341}{<< /D [\MusicPDFPageRef{618} /XYZ 203.079 309.457 null] >>}
\MusicPDFWriteObject{14342}{<< /D [\MusicPDFPageRef{619} /XYZ 166.698 257.732 null] >>}
\MusicPDFWriteObject{14343}{<< /D [\MusicPDFPageRef{622} /XYZ 160.425 634.606 null] >>}
\MusicPDFWriteObject{14344}{<< /D [\MusicPDFPageRef{622} /XYZ 162.201 584.744 null] >>}
\MusicPDFWriteObject{14345}{<< /D [\MusicPDFPageRef{622} /XYZ 72 515.715 null] >>}
\MusicPDFWriteObject{14346}{<< /D [\MusicPDFPageRef{622} /XYZ 207.605 356.128 null] >>}
\MusicPDFWriteObject{14347}{<< /D [\MusicPDFPageRef{622} /XYZ 137.175 176.771 null] >>}
\MusicPDFWriteObject{14348}{<< /D [\MusicPDFPageRef{623} /XYZ 130.59 554.893 null] >>}
\MusicPDFWriteObject{14349}{<< /D [\MusicPDFPageRef{623} /XYZ 72 346.475 null] >>}
\MusicPDFWriteObject{14350}{<< /D [\MusicPDFPageRef{624} /XYZ 258.106 639.708 null] >>}
\MusicPDFWriteObject{14351}{<< /D [\MusicPDFPageRef{624} /XYZ 208.036 537.052 null] >>}
\MusicPDFWriteObject{14352}{<< /D [\MusicPDFPageRef{624} /XYZ 250.225 213.198 null] >>}
\MusicPDFWriteObject{14353}{<< /D [\MusicPDFPageRef{620} /XYZ 163.382 675.575 null] >>}
\MusicPDFWriteObject{14354}{<< /D [\MusicPDFPageRef{625} /XYZ 182.536 557.711 null] >>}
\MusicPDFWriteObject{14355}{<< /D [\MusicPDFPageRef{625} /XYZ 131.485 494.292 null] >>}
\MusicPDFWriteObject{14356}{<< /D [\MusicPDFPageRef{625} /XYZ 153.464 230.608 null] >>}
\MusicPDFWriteObject{14357}{<< /D [\MusicPDFPageRef{620} /XYZ 266.324 553.033 null] >>}
\MusicPDFWriteObject{14358}{<< /D [\MusicPDFPageRef{620} /XYZ 278.188 378.726 null] >>}
\MusicPDFWriteObject{14359}{<< /D [\MusicPDFPageRef{620} /XYZ 197.543 254.939 null] >>}
\MusicPDFWriteObject{14360}{<< /D [\MusicPDFPageRef{620} /XYZ 160.498 153.729 null] >>}
\MusicPDFWriteObject{14361}{<< /D [\MusicPDFPageRef{621} /XYZ 268.552 620.77 null] >>}
\MusicPDFWriteObject{14362}{<< /D [\MusicPDFPageRef{621} /XYZ 72 562.354 null] >>}
\MusicPDFWriteObject{14363}{<< /D [\MusicPDFPageRef{621} /XYZ 72 381.301 null] >>}
\MusicPDFWriteObject{14364}{<< /D [\MusicPDFPageRef{626} /XYZ 271.869 195.738 null] >>}
\MusicPDFWriteObject{14365}{<< /D [\MusicPDFPageRef{626} /XYZ 273.88 148.811 null] >>}
\MusicPDFWriteObject{14366}{<< /D [\MusicPDFPageRef{626} /XYZ 230.123 99.233 null] >>}
\MusicPDFWriteObject{14367}{<< /D [\MusicPDFPageRef{627} /XYZ 158.914 214.996 null] >>}
\MusicPDFWriteObject{14368}{<< /D [\MusicPDFPageRef{630} /XYZ 244.446 162.304 null] >>}
\MusicPDFWriteObject{14369}{<< /D [\MusicPDFPageRef{631} /XYZ 73.161 597.573 null] >>}
\MusicPDFWriteObject{14370}{<< /D [\MusicPDFPageRef{631} /XYZ 283.775 550.859 null] >>}
\MusicPDFWriteObject{14371}{<< /D [\MusicPDFPageRef{631} /XYZ 142.505 510.723 null] >>}
\MusicPDFWriteObject{14372}{<< /D [\MusicPDFPageRef{631} /XYZ 72 436.462 null] >>}
\MusicPDFWriteObject{14373}{<< /D [\MusicPDFPageRef{631} /XYZ 190.588 187.934 null] >>}
\MusicPDFWriteObject{14374}{<< /D [\MusicPDFPageRef{627} /XYZ 258.353 99.233 null] >>}
\MusicPDFWriteObject{14375}{<< /D [\MusicPDFPageRef{628} /XYZ 250.721 679.714 null] >>}
\MusicPDFWriteObject{14376}{<< /D [\MusicPDFPageRef{629} /XYZ 259.483 496.679 null] >>}
\MusicPDFWriteObject{14377}{<< /D [\MusicPDFPageRef{629} /XYZ 245.029 454.55 null] >>}
\MusicPDFWriteObject{14378}{<< /D [\MusicPDFPageRef{629} /XYZ 229.043 407.531 null] >>}
\MusicPDFWriteObject{14379}{<< /D [\MusicPDFPageRef{629} /XYZ 126.519 362.781 null] >>}
\MusicPDFWriteObject{14380}{<< /D [\MusicPDFPageRef{629} /XYZ 165.934 320.333 null] >>}
\MusicPDFWriteObject{14381}{<< /D [\MusicPDFPageRef{630} /XYZ 126.687 436.507 null] >>}
\MusicPDFWriteObject{14382}{<< /D [\MusicPDFPageRef{632} /XYZ 249.516 617.303 null] >>}
\MusicPDFWriteObject{14383}{<< /D [\MusicPDFPageRef{634} /XYZ 207.511 515.171 null] >>}
\MusicPDFWriteObject{14384}{<< /D [\MusicPDFPageRef{634} /XYZ 258.706 474.351 null] >>}
\MusicPDFWriteObject{14385}{<< /D [\MusicPDFPageRef{634} /XYZ 170.255 383.115 null] >>}
\MusicPDFWriteObject{14386}{<< /D [\MusicPDFPageRef{634} /XYZ 181.522 281.109 null] >>}
\MusicPDFWriteObject{14387}{<< /D [\MusicPDFPageRef{634} /XYZ 122.539 103.5 null] >>}
\MusicPDFWriteObject{14388}{<< /D [\MusicPDFPageRef{635} /XYZ 120.907 664.625 null] >>}
\MusicPDFWriteObject{14389}{<< /D [\MusicPDFPageRef{632} /XYZ 256.525 502.18 null] >>}
\MusicPDFWriteObject{14390}{<< /D [\MusicPDFPageRef{632} /XYZ 241.965 200.809 null] >>}
\MusicPDFWriteObject{14391}{<< /D [\MusicPDFPageRef{633} /XYZ 151.896 662.689 null] >>}
\MusicPDFWriteObject{14392}{<< /D [\MusicPDFPageRef{633} /XYZ 151.896 612.845 null] >>}
\MusicPDFWriteObject{14393}{<< /D [\MusicPDFPageRef{633} /XYZ 156.188 197.122 null] >>}
\MusicPDFWriteObject{14394}{<< /D [\MusicPDFPageRef{633} /XYZ 175.06 162.089 null] >>}
\MusicPDFWriteObject{14395}{<< /D [\MusicPDFPageRef{633} /XYZ 249.17 119.82 null] >>}
\MusicPDFWriteObject{14396}{<< /D [\MusicPDFPageRef{634} /XYZ 121.416 572.294 null] >>}
\MusicPDFWriteObject{14397}{<< /D [\MusicPDFPageRef{636} /XYZ 232.558 488.257 null] >>}
\MusicPDFWriteObject{14398}{<< /D [\MusicPDFPageRef{637} /XYZ 250.049 140.052 null] >>}
\MusicPDFWriteObject{14399}{<< /D [\MusicPDFPageRef{637} /XYZ 275.257 99.233 null] >>}
\MusicPDFWriteObject{14400}{<< /D [\MusicPDFPageRef{638} /XYZ 273.609 650.213 null] >>}
\MusicPDFWriteObject{14401}{<< /D [\MusicPDFPageRef{638} /XYZ 182.488 568.153 null] >>}
\MusicPDFWriteObject{14402}{<< /D [\MusicPDFPageRef{638} /XYZ 239.852 523.485 null] >>}
\MusicPDFWriteObject{14403}{<< /D [\MusicPDFPageRef{638} /XYZ 179.879 283.612 null] >>}
\MusicPDFWriteObject{14404}{<< /D [\MusicPDFPageRef{638} /XYZ 253.439 171.087 null] >>}
\MusicPDFWriteObject{14405}{<< /D [\MusicPDFPageRef{638} /XYZ 261.885 99.233 null] >>}
\MusicPDFWriteObject{14406}{<< /D [\MusicPDFPageRef{639} /XYZ 147.661 680.765 null] >>}
\MusicPDFWriteObject{14407}{<< /D [\MusicPDFPageRef{639} /XYZ 155.604 534.349 null] >>}
\MusicPDFWriteObject{14408}{<< /D [\MusicPDFPageRef{636} /XYZ 252.455 314.416 null] >>}
\MusicPDFWriteObject{14409}{<< /D [\MusicPDFPageRef{639} /XYZ 119.83 372.529 null] >>}
\MusicPDFWriteObject{14410}{<< /D [\MusicPDFPageRef{639} /XYZ 246.663 220.382 null] >>}
\MusicPDFWriteObject{14411}{<< /D [\MusicPDFPageRef{640} /XYZ 252.895 599.591 null] >>}
\MusicPDFWriteObject{14412}{<< /D [\MusicPDFPageRef{640} /XYZ 223.868 504.855 null] >>}
\MusicPDFWriteObject{14413}{<< /D [\MusicPDFPageRef{640} /XYZ 247.622 454.424 null] >>}
\MusicPDFWriteObject{14414}{<< /D [\MusicPDFPageRef{640} /XYZ 244.907 385.806 null] >>}
\MusicPDFWriteObject{14415}{<< /D [\MusicPDFPageRef{640} /XYZ 268.593 299.018 null] >>}
\MusicPDFWriteObject{14416}{<< /D [\MusicPDFPageRef{640} /XYZ 263.279 217.472 null] >>}
\MusicPDFWriteObject{14417}{<< /D [\MusicPDFPageRef{640} /XYZ 133.434 167.041 null] >>}
\MusicPDFWriteObject{14418}{<< /D [\MusicPDFPageRef{640} /XYZ 237.067 124.034 null] >>}
\MusicPDFWriteObject{14419}{<< /D [\MusicPDFPageRef{636} /XYZ 245.487 267.483 null] >>}
\MusicPDFWriteObject{14420}{<< /D [\MusicPDFPageRef{641} /XYZ 168.524 678.62 null] >>}
\MusicPDFWriteObject{14421}{<< /D [\MusicPDFPageRef{641} /XYZ 149.011 631.219 null] >>}
\MusicPDFWriteObject{14422}{<< /D [\MusicPDFPageRef{641} /XYZ 243.601 568.364 null] >>}
\MusicPDFWriteObject{14423}{<< /D [\MusicPDFPageRef{641} /XYZ 241.816 525.333 null] >>}
\MusicPDFWriteObject{14424}{<< /D [\MusicPDFPageRef{641} /XYZ 232.494 480.601 null] >>}
\MusicPDFWriteObject{14425}{<< /D [\MusicPDFPageRef{641} /XYZ 250.559 167.025 null] >>}
\MusicPDFWriteObject{14426}{<< /D [\MusicPDFPageRef{641} /XYZ 248.343 117.447 null] >>}
\MusicPDFWriteObject{14427}{<< /D [\MusicPDFPageRef{642} /XYZ 247.101 669.806 null] >>}
\MusicPDFWriteObject{14428}{<< /D [\MusicPDFPageRef{642} /XYZ 251.491 607.077 null] >>}
\MusicPDFWriteObject{14429}{<< /D [\MusicPDFPageRef{642} /XYZ 193.77 486.382 null] >>}
\MusicPDFWriteObject{14430}{<< /D [\MusicPDFPageRef{636} /XYZ 258.116 220.996 null] >>}
\MusicPDFWriteObject{14431}{<< /D [\MusicPDFPageRef{642} /XYZ 173.604 375.832 null] >>}
\MusicPDFWriteObject{14432}{<< /D [\MusicPDFPageRef{642} /XYZ 226.151 233.562 null] >>}
\MusicPDFWriteObject{14433}{<< /D [\MusicPDFPageRef{642} /XYZ 193.377 186.857 null] >>}
\MusicPDFWriteObject{14434}{<< /D [\MusicPDFPageRef{643} /XYZ 130.682 669.103 null] >>}
\MusicPDFWriteObject{14435}{<< /D [\MusicPDFPageRef{643} /XYZ 190.895 585.187 null] >>}
\MusicPDFWriteObject{14436}{<< /D [\MusicPDFPageRef{643} /XYZ 216.748 430.103 null] >>}
\MusicPDFWriteObject{14437}{<< /D [\MusicPDFPageRef{637} /XYZ 265.1 683.057 null] >>}
\MusicPDFWriteObject{14438}{<< /D [\MusicPDFPageRef{637} /XYZ 270.185 574.333 null] >>}
\MusicPDFWriteObject{14439}{<< /D [\MusicPDFPageRef{637} /XYZ 215.515 492.086 null] >>}
\MusicPDFWriteObject{14440}{<< /D [\MusicPDFPageRef{637} /XYZ 273.469 446.912 null] >>}
\MusicPDFWriteObject{14441}{<< /D [\MusicPDFPageRef{637} /XYZ 273.96 286.493 null] >>}
\MusicPDFWriteObject{14442}{<< /D [\MusicPDFPageRef{643} /XYZ 273.27 189.685 null] >>}
\MusicPDFWriteObject{14443}{<< /D [\MusicPDFPageRef{646} /XYZ 211.622 675.121 null] >>}
\MusicPDFWriteObject{14444}{<< /D [\MusicPDFPageRef{646} /XYZ 72 564.542 null] >>}
\MusicPDFWriteObject{14445}{<< /D [\MusicPDFPageRef{647} /XYZ 173.702 357.565 null] >>}
\MusicPDFWriteObject{14446}{<< /D [\MusicPDFPageRef{647} /XYZ 122.959 297.727 null] >>}
\MusicPDFWriteObject{14447}{<< /D [\MusicPDFPageRef{647} /XYZ 120.69 251.04 null] >>}
\MusicPDFWriteObject{14448}{<< /D [\MusicPDFPageRef{647} /XYZ 178.052 207.183 null] >>}
\MusicPDFWriteObject{14449}{<< /D [\MusicPDFPageRef{648} /XYZ 151.895 370.535 null] >>}
\MusicPDFWriteObject{14450}{<< /D [\MusicPDFPageRef{648} /XYZ 272.037 228.31 null] >>}
\MusicPDFWriteObject{14451}{<< /D [\MusicPDFPageRef{644} /XYZ 275.659 662.389 null] >>}
\MusicPDFWriteObject{14452}{<< /D [\MusicPDFPageRef{644} /XYZ 249.144 583.869 null] >>}
\MusicPDFWriteObject{14453}{<< /D [\MusicPDFPageRef{644} /XYZ 75.657 494.774 null] >>}
\MusicPDFWriteObject{14454}{<< /D [\MusicPDFPageRef{644} /XYZ 72 421.382 null] >>}
\MusicPDFWriteObject{14455}{<< /D [\MusicPDFPageRef{644} /XYZ 148.821 273.822 null] >>}
\MusicPDFWriteObject{14456}{<< /D [\MusicPDFPageRef{645} /XYZ 132.628 427.013 null] >>}
\MusicPDFWriteObject{14457}{<< /D [\MusicPDFPageRef{645} /XYZ 175.561 174.358 null] >>}
\MusicPDFWriteObject{14458}{<< /D [\MusicPDFPageRef{645} /XYZ 215.836 110.358 null] >>}
\MusicPDFWriteObject{14459}{<< /D [\MusicPDFPageRef{649} /XYZ 72 553.652 null] >>}
\MusicPDFWriteObject{14460}{<< /D [\MusicPDFPageRef{650} /XYZ 184.297 351.987 null] >>}
\MusicPDFWriteObject{14461}{<< /D [\MusicPDFPageRef{650} /XYZ 220.108 308.445 null] >>}
\MusicPDFWriteObject{14462}{<< /D [\MusicPDFPageRef{650} /XYZ 172.956 265.171 null] >>}
\MusicPDFWriteObject{14463}{<< /D [\MusicPDFPageRef{651} /XYZ 72 689.957 null] >>}
\MusicPDFWriteObject{14464}{<< /D [\MusicPDFPageRef{651} /XYZ 72 473.503 null] >>}
\MusicPDFWriteObject{14465}{<< /D [\MusicPDFPageRef{651} /XYZ 266.452 263.847 null] >>}
\MusicPDFWriteObject{14466}{<< /D [\MusicPDFPageRef{651} /XYZ 155.639 177.502 null] >>}
\MusicPDFWriteObject{14467}{<< /D [\MusicPDFPageRef{652} /XYZ 212.991 562.211 null] >>}
\MusicPDFWriteObject{14468}{<< /D [\MusicPDFPageRef{652} /XYZ 172.511 446.821 null] >>}
\MusicPDFWriteObject{14469}{<< /D [\MusicPDFPageRef{652} /XYZ 122.798 404.471 null] >>}
\MusicPDFWriteObject{14470}{<< /D [\MusicPDFPageRef{649} /XYZ 255.311 452.113 null] >>}
\MusicPDFWriteObject{14471}{<< /D [\MusicPDFPageRef{652} /XYZ 246.039 274.27 null] >>}
\MusicPDFWriteObject{14472}{<< /D [\MusicPDFPageRef{652} /XYZ 190.697 209.03 null] >>}
\MusicPDFWriteObject{14473}{<< /D [\MusicPDFPageRef{652} /XYZ 246.039 172.602 null] >>}
\MusicPDFWriteObject{14474}{<< /D [\MusicPDFPageRef{653} /XYZ 138.199 679.063 null] >>}
\MusicPDFWriteObject{14475}{<< /D [\MusicPDFPageRef{653} /XYZ 238.079 499.158 null] >>}
\MusicPDFWriteObject{14476}{<< /D [\MusicPDFPageRef{653} /XYZ 132.377 356.305 null] >>}
\MusicPDFWriteObject{14477}{<< /D [\MusicPDFPageRef{649} /XYZ 236.674 369.035 null] >>}
\MusicPDFWriteObject{14478}{<< /D [\MusicPDFPageRef{649} /XYZ 72 146.585 null] >>}
\MusicPDFWriteObject{14479}{<< /D [\MusicPDFPageRef{649} /XYZ 72 116.731 null] >>}
\MusicPDFWriteObject{14480}{<< /D [\MusicPDFPageRef{650} /XYZ 72 668.646 null] >>}
\MusicPDFWriteObject{14481}{<< /D [\MusicPDFPageRef{650} /XYZ 72 638.791 null] >>}
\MusicPDFWriteObject{14482}{<< /D [\MusicPDFPageRef{650} /XYZ 162.633 571.094 null] >>}
\MusicPDFWriteObject{14483}{<< /D [\MusicPDFPageRef{650} /XYZ 72 490.027 null] >>}
\MusicPDFWriteObject{14484}{<< /D [\MusicPDFPageRef{658} /XYZ 259.045 282.523 null] >>}
\MusicPDFWriteObject{14485}{<< /D [\MusicPDFPageRef{663} /XYZ 221.068 234.599 null] >>}
\MusicPDFWriteObject{14486}{<< /D [\MusicPDFPageRef{667} /XYZ 216.215 625.568 null] >>}
\MusicPDFWriteObject{14487}{<< /D [\MusicPDFPageRef{667} /XYZ 220.361 566.401 null] >>}
\MusicPDFWriteObject{14488}{<< /D [\MusicPDFPageRef{667} /XYZ 254.488 498.056 null] >>}
\MusicPDFWriteObject{14489}{<< /D [\MusicPDFPageRef{670} /XYZ 234.457 355.319 null] >>}
\MusicPDFWriteObject{14490}{<< /D [\MusicPDFPageRef{670} /XYZ 236.152 318.892 null] >>}
\MusicPDFWriteObject{14491}{<< /D [\MusicPDFPageRef{671} /XYZ 285.922 636.716 null] >>}
\MusicPDFWriteObject{14492}{<< /D [\MusicPDFPageRef{671} /XYZ 258.833 573.987 null] >>}
\MusicPDFWriteObject{14493}{<< /D [\MusicPDFPageRef{673} /XYZ 268.96 551.549 null] >>}
\MusicPDFWriteObject{14494}{<< /D [\MusicPDFPageRef{673} /XYZ 199.183 409.915 null] >>}
\MusicPDFWriteObject{14495}{<< /D [\MusicPDFPageRef{673} /XYZ 198.491 347.186 null] >>}
\MusicPDFWriteObject{14496}{<< /D [\MusicPDFPageRef{673} /XYZ 258.864 218.46 null] >>}
\MusicPDFWriteObject{14497}{<< /D [\MusicPDFPageRef{673} /XYZ 261.294 99.233 null] >>}
\MusicPDFWriteObject{14498}{<< /D [\MusicPDFPageRef{675} /XYZ 267.789 543.565 null] >>}
\MusicPDFWriteObject{14499}{<< /D [\MusicPDFPageRef{676} /XYZ 257.015 564.7 null] >>}
\MusicPDFWriteObject{14500}{<< /D [\MusicPDFPageRef{678} /XYZ 238.528 633.478 null] >>}
\MusicPDFWriteObject{14501}{<< /D [\MusicPDFPageRef{678} /XYZ 261.072 525.873 null] >>}
\MusicPDFWriteObject{14502}{<< /D [\MusicPDFPageRef{678} /XYZ 211.887 276.499 null] >>}
\MusicPDFWriteObject{14503}{<< /D [\MusicPDFPageRef{678} /XYZ 211.334 206.012 null] >>}
\MusicPDFWriteObject{14504}{<< /D [\MusicPDFPageRef{679} /XYZ 290.319 604.152 null] >>}
\MusicPDFWriteObject{14505}{<< /D [\MusicPDFPageRef{679} /XYZ 227.284 528.272 null] >>}
\MusicPDFWriteObject{14506}{<< /D [\MusicPDFPageRef{679} /XYZ 277.203 373.245 null] >>}
\MusicPDFWriteObject{14507}{<< /D [\MusicPDFPageRef{680} /XYZ 236.36 381.438 null] >>}
\MusicPDFWriteObject{14508}{<< /D [\MusicPDFPageRef{680} /XYZ 235.806 329.657 null] >>}
\MusicPDFWriteObject{14509}{<< /D [\MusicPDFPageRef{680} /XYZ 236.747 276.714 null] >>}
\MusicPDFWriteObject{14510}{<< /D [\MusicPDFPageRef{676} /XYZ 232.202 528.272 null] >>}
\MusicPDFWriteObject{14511}{<< /D [\MusicPDFPageRef{681} /XYZ 239.22 560.632 null] >>}
\MusicPDFWriteObject{14512}{<< /D [\MusicPDFPageRef{681} /XYZ 274.3 120.594 null] >>}
\MusicPDFWriteObject{14513}{<< /D [\MusicPDFPageRef{682} /XYZ 292.129 390.815 null] >>}
\MusicPDFWriteObject{14514}{<< /D [\MusicPDFPageRef{682} /XYZ 293.576 275.483 null] >>}
\MusicPDFWriteObject{14515}{<< /D [\MusicPDFPageRef{682} /XYZ 267.729 205.876 null] >>}
\MusicPDFWriteObject{14516}{<< /D [\MusicPDFPageRef{683} /XYZ 244.368 662.689 null] >>}
\MusicPDFWriteObject{14517}{<< /D [\MusicPDFPageRef{683} /XYZ 247.178 563.294 null] >>}
\MusicPDFWriteObject{14518}{<< /D [\MusicPDFPageRef{683} /XYZ 251.732 491.132 null] >>}
\MusicPDFWriteObject{14519}{<< /D [\MusicPDFPageRef{684} /XYZ 269.368 624.186 null] >>}
\MusicPDFWriteObject{14520}{<< /D [\MusicPDFPageRef{684} /XYZ 200.134 307.825 null] >>}
\MusicPDFWriteObject{14521}{<< /D [\MusicPDFPageRef{676} /XYZ 256.642 310.956 null] >>}
\MusicPDFWriteObject{14522}{<< /D [\MusicPDFPageRef{676} /XYZ 291.721 159.196 null] >>}
\MusicPDFWriteObject{14523}{<< /D [\MusicPDFPageRef{677} /XYZ 253.6 577.801 null] >>}
\MusicPDFWriteObject{14524}{<< /D [\MusicPDFPageRef{677} /XYZ 244.977 541.374 null] >>}
\MusicPDFWriteObject{14525}{<< /D [\MusicPDFPageRef{677} /XYZ 172.537 323.79 null] >>}
\MusicPDFWriteObject{14526}{<< /D [\MusicPDFPageRef{677} /XYZ 207.049 242.811 null] >>}
\MusicPDFWriteObject{14527}{<< /D [\MusicPDFPageRef{678} /XYZ 287.906 669.906 null] >>}
\MusicPDFWriteObject{14528}{<< /D [\MusicPDFPageRef{685} /XYZ 267.321 491.576 null] >>}
\MusicPDFWriteObject{14529}{<< /D [\MusicPDFPageRef{691} /XYZ 225.096 539.221 null] >>}
\MusicPDFWriteObject{14530}{<< /D [\MusicPDFPageRef{686} /XYZ 253.522 669.906 null] >>}
\MusicPDFWriteObject{14531}{<< /D [\MusicPDFPageRef{687} /XYZ 182.449 622.656 null] >>}
\MusicPDFWriteObject{14532}{<< /D [\MusicPDFPageRef{687} /XYZ 259.306 207.922 null] >>}
\MusicPDFWriteObject{14533}{<< /D [\MusicPDFPageRef{688} /XYZ 209.184 325.355 null] >>}
\MusicPDFWriteObject{14534}{<< /D [\MusicPDFPageRef{689} /XYZ 258.394 273.154 null] >>}
\MusicPDFWriteObject{14535}{<< /D [\MusicPDFPageRef{689} /XYZ 253.251 236.727 null] >>}
\MusicPDFWriteObject{14536}{<< /D [\MusicPDFPageRef{690} /XYZ 230.34 255.386 null] >>}
\MusicPDFWriteObject{14537}{<< /D [\MusicPDFPageRef{690} /XYZ 264.324 99.233 null] >>}
\MusicPDFWriteObject{14538}{<< /D [\MusicPDFPageRef{693} /XYZ 258.63 392.479 null] >>}
\MusicPDFWriteObject{14539}{<< /D [\MusicPDFPageRef{693} /XYZ 259.183 329.75 null] >>}
\MusicPDFWriteObject{14540}{<< /D [\MusicPDFPageRef{694} /XYZ 249.515 511.315 null] >>}
\MusicPDFWriteObject{14541}{<< /D [\MusicPDFPageRef{694} /XYZ 251.112 321.343 null] >>}
\MusicPDFWriteObject{14542}{<< /D [\MusicPDFPageRef{694} /XYZ 235.185 274.594 null] >>}
\MusicPDFWriteObject{14543}{<< /D [\MusicPDFPageRef{694} /XYZ 238.487 202.319 null] >>}
\MusicPDFWriteObject{14544}{<< /D [\MusicPDFPageRef{695} /XYZ 72 685.548 null] >>}
\MusicPDFWriteObject{14545}{<< /D [\MusicPDFPageRef{695} /XYZ 258.215 399.768 null] >>}
\MusicPDFWriteObject{14546}{<< /D [\MusicPDFPageRef{695} /XYZ 240.098 211.552 null] >>}
\MusicPDFWriteObject{14547}{<< /D [\MusicPDFPageRef{695} /XYZ 143.799 147.166 null] >>}
\MusicPDFWriteObject{14548}{<< /D [\MusicPDFPageRef{696} /XYZ 252.129 99.233 null] >>}
\MusicPDFWriteObject{14549}{<< /D [\MusicPDFPageRef{697} /XYZ 253.701 683.057 null] >>}
\MusicPDFWriteObject{14550}{<< /D [\MusicPDFPageRef{698} /XYZ 261.384 523.891 null] >>}
\MusicPDFWriteObject{14551}{<< /D [\MusicPDFPageRef{698} /XYZ 112.831 485.294 null] >>}
\MusicPDFWriteObject{14552}{<< /D [\MusicPDFPageRef{698} /XYZ 241.806 280.845 null] >>}
\MusicPDFWriteObject{14553}{<< /D [\MusicPDFPageRef{699} /XYZ 280.699 486.015 null] >>}
\MusicPDFWriteObject{14554}{<< /D [\MusicPDFPageRef{701} /XYZ 239.319 548.423 null] >>}
\MusicPDFWriteObject{14555}{<< /D [\MusicPDFPageRef{701} /XYZ 206.254 413.112 null] >>}
\MusicPDFWriteObject{14556}{<< /D [\MusicPDFPageRef{701} /XYZ 219.851 329.663 null] >>}
\MusicPDFWriteObject{14557}{<< /D [\MusicPDFPageRef{701} /XYZ 203.244 272.978 null] >>}
\MusicPDFWriteObject{14558}{<< /D [\MusicPDFPageRef{701} /XYZ 134.052 230.054 null] >>}
\MusicPDFWriteObject{14559}{<< /D [\MusicPDFPageRef{701} /XYZ 194.846 186.522 null] >>}
\MusicPDFWriteObject{14560}{<< /D [\MusicPDFPageRef{702} /XYZ 188.248 580.078 null] >>}
\MusicPDFWriteObject{14561}{<< /D [\MusicPDFPageRef{702} /XYZ 192.266 493.874 null] >>}
\MusicPDFWriteObject{14562}{<< /D [\MusicPDFPageRef{702} /XYZ 198.352 295.728 null] >>}
\MusicPDFWriteObject{14563}{<< /D [\MusicPDFPageRef{702} /XYZ 201.182 242.353 null] >>}
\MusicPDFWriteObject{14564}{<< /D [\MusicPDFPageRef{699} /XYZ 258.067 420.008 null] >>}
\MusicPDFWriteObject{14565}{<< /D [\MusicPDFPageRef{702} /XYZ 162.226 176.205 null] >>}
\MusicPDFWriteObject{14566}{<< /D [\MusicPDFPageRef{702} /XYZ 254.038 103.349 null] >>}
\MusicPDFWriteObject{14567}{<< /D [\MusicPDFPageRef{703} /XYZ 230.491 616.506 null] >>}
\MusicPDFWriteObject{14568}{<< /D [\MusicPDFPageRef{703} /XYZ 185.762 518.558 null] >>}
\MusicPDFWriteObject{14569}{<< /D [\MusicPDFPageRef{703} /XYZ 164.284 432.295 null] >>}
\MusicPDFWriteObject{14570}{<< /D [\MusicPDFPageRef{699} /XYZ 255.68 327.635 null] >>}
\MusicPDFWriteObject{14571}{<< /D [\MusicPDFPageRef{700} /XYZ 170.026 643.604 null] >>}
\MusicPDFWriteObject{14572}{<< /D [\MusicPDFPageRef{700} /XYZ 223.053 475.309 null] >>}
\MusicPDFWriteObject{14573}{<< /D [\MusicPDFPageRef{700} /XYZ 166.041 353.801 null] >>}
\MusicPDFWriteObject{14574}{<< /D [\MusicPDFPageRef{700} /XYZ 250.64 306.71 null] >>}
\MusicPDFWriteObject{14575}{<< /D [\MusicPDFPageRef{700} /XYZ 237.493 160.128 null] >>}
\MusicPDFWriteObject{14576}{<< /D [\MusicPDFPageRef{700} /XYZ 164.823 104.878 null] >>}
\MusicPDFWriteObject{14577}{<< /D [\MusicPDFPageRef{703} /XYZ 261.91 157.849 null] >>}
\MusicPDFWriteObject{14578}{<< /D [\MusicPDFPageRef{706} /XYZ 224.764 683.057 null] >>}
\MusicPDFWriteObject{14579}{<< /D [\MusicPDFPageRef{706} /XYZ 258.329 600.161 null] >>}
\MusicPDFWriteObject{14580}{<< /D [\MusicPDFPageRef{706} /XYZ 95.881 519.02 null] >>}
\MusicPDFWriteObject{14581}{<< /D [\MusicPDFPageRef{706} /XYZ 133.249 445.551 null] >>}
\MusicPDFWriteObject{14582}{<< /D [\MusicPDFPageRef{706} /XYZ 72 338.905 null] >>}
\MusicPDFWriteObject{14583}{<< /D [\MusicPDFPageRef{706} /XYZ 117.812 217.588 null] >>}
\MusicPDFWriteObject{14584}{<< /D [\MusicPDFPageRef{706} /XYZ 209.19 154.859 null] >>}
\MusicPDFWriteObject{14585}{<< /D [\MusicPDFPageRef{706} /XYZ 117.812 99.233 null] >>}
\MusicPDFWriteObject{14586}{<< /D [\MusicPDFPageRef{707} /XYZ 207.474 682.571 null] >>}
\MusicPDFWriteObject{14587}{<< /D [\MusicPDFPageRef{707} /XYZ 72 618.858 null] >>}
\MusicPDFWriteObject{14588}{<< /D [\MusicPDFPageRef{704} /XYZ 192.907 433.808 null] >>}
\MusicPDFWriteObject{14589}{<< /D [\MusicPDFPageRef{704} /XYZ 261.798 363.707 null] >>}
\MusicPDFWriteObject{14590}{<< /D [\MusicPDFPageRef{704} /XYZ 247.043 300.242 null] >>}
\MusicPDFWriteObject{14591}{<< /D [\MusicPDFPageRef{705} /XYZ 215.811 656.755 null] >>}
\MusicPDFWriteObject{14592}{<< /D [\MusicPDFPageRef{705} /XYZ 239.517 603.28 null] >>}
\MusicPDFWriteObject{14593}{<< /D [\MusicPDFPageRef{705} /XYZ 258.183 500.465 null] >>}
\MusicPDFWriteObject{14594}{<< /D [\MusicPDFPageRef{705} /XYZ 215.746 437.736 null] >>}
\MusicPDFWriteObject{14595}{<< /D [\MusicPDFPageRef{705} /XYZ 290.326 401.309 null] >>}
\MusicPDFWriteObject{14596}{<< /D [\MusicPDFPageRef{708} /XYZ 238.36 454.526 null] >>}
\MusicPDFWriteObject{14597}{<< /D [\MusicPDFPageRef{708} /XYZ 245.008 372.845 null] >>}
\MusicPDFWriteObject{14598}{<< /D [\MusicPDFPageRef{709} /XYZ 275.282 399.789 null] >>}
\MusicPDFWriteObject{14599}{<< /D [\MusicPDFPageRef{709} /XYZ 261.854 204.981 null] >>}
\MusicPDFWriteObject{14600}{<< /D [\MusicPDFPageRef{710} /XYZ 186.421 288.528 null] >>}
\MusicPDFWriteObject{14601}{<< /D [\MusicPDFPageRef{711} /XYZ 251.982 430.352 null] >>}
\MusicPDFWriteObject{14602}{<< /D [\MusicPDFPageRef{715} /XYZ 154.686 240.255 null] >>}
\MusicPDFWriteObject{14603}{<< /D [\MusicPDFPageRef{715} /XYZ 273.515 104.325 null] >>}
\MusicPDFWriteObject{14604}{<< /D [\MusicPDFPageRef{716} /XYZ 183.463 652.494 null] >>}
\MusicPDFWriteObject{14605}{<< /D [\MusicPDFPageRef{716} /XYZ 182.819 588.359 null] >>}
\MusicPDFWriteObject{14606}{<< /D [\MusicPDFPageRef{716} /XYZ 185.534 503.302 null] >>}
\MusicPDFWriteObject{14607}{<< /D [\MusicPDFPageRef{716} /XYZ 256.987 289.649 null] >>}
\MusicPDFWriteObject{14608}{<< /D [\MusicPDFPageRef{716} /XYZ 212.79 231.696 null] >>}
\MusicPDFWriteObject{14609}{<< /D [\MusicPDFPageRef{712} /XYZ 213.295 649.881 null] >>}
\MusicPDFWriteObject{14610}{<< /D [\MusicPDFPageRef{714} /XYZ 233.205 666.009 null] >>}
\MusicPDFWriteObject{14611}{<< /D [\MusicPDFPageRef{714} /XYZ 229.265 366.77 null] >>}
\MusicPDFWriteObject{14612}{<< /D [\MusicPDFPageRef{714} /XYZ 236.561 214.46 null] >>}
\MusicPDFWriteObject{14613}{<< /D [\MusicPDFPageRef{714} /XYZ 244.123 168.45 null] >>}
\MusicPDFWriteObject{14614}{<< /D [\MusicPDFPageRef{715} /XYZ 185.974 678.292 null] >>}
\MusicPDFWriteObject{14615}{<< /D [\MusicPDFPageRef{715} /XYZ 187.054 616.472 null] >>}
\MusicPDFWriteObject{14616}{<< /D [\MusicPDFPageRef{715} /XYZ 234.821 380.601 null] >>}
\MusicPDFWriteObject{14617}{<< /D [\MusicPDFPageRef{717} /XYZ 148.413 558.039 null] >>}
\MusicPDFWriteObject{14618}{<< /D [\MusicPDFPageRef{719} /XYZ 155.924 554.172 null] >>}
\MusicPDFWriteObject{14619}{<< /D [\MusicPDFPageRef{719} /XYZ 86.676 488.197 null] >>}
\MusicPDFWriteObject{14620}{<< /D [\MusicPDFPageRef{719} /XYZ 183.222 365.11 null] >>}
\MusicPDFWriteObject{14621}{<< /D [\MusicPDFPageRef{717} /XYZ 153.906 508.98 null] >>}
\MusicPDFWriteObject{14622}{<< /D [\MusicPDFPageRef{717} /XYZ 72 372.593 null] >>}
\MusicPDFWriteObject{14623}{<< /D [\MusicPDFPageRef{717} /XYZ 174.249 286.707 null] >>}
\MusicPDFWriteObject{14624}{<< /D [\MusicPDFPageRef{718} /XYZ 204.365 467.255 null] >>}
\MusicPDFWriteObject{14625}{<< /D [\MusicPDFPageRef{718} /XYZ 242.354 314.502 null] >>}
\MusicPDFWriteObject{14626}{<< /D [\MusicPDFPageRef{718} /XYZ 233.465 146.195 null] >>}
\MusicPDFWriteObject{14627}{<< /D [\MusicPDFPageRef{719} /XYZ 233.106 679.161 null] >>}
\MusicPDFWriteObject{14628}{<< /D [\MusicPDFPageRef{719} /XYZ 252.094 611.505 null] >>}
\MusicPDFWriteObject{14629}{<< /D [\MusicPDFPageRef{720} /XYZ 268.546 564.7 null] >>}
\MusicPDFWriteObject{14630}{<< /D [\MusicPDFPageRef{724} /XYZ 244.186 466.212 null] >>}
\MusicPDFWriteObject{14631}{<< /D [\MusicPDFPageRef{724} /XYZ 207.894 394.744 null] >>}
\MusicPDFWriteObject{14632}{<< /D [\MusicPDFPageRef{724} /XYZ 273.767 313.705 null] >>}
\MusicPDFWriteObject{14633}{<< /D [\MusicPDFPageRef{724} /XYZ 262.186 277.277 null] >>}
\MusicPDFWriteObject{14634}{<< /D [\MusicPDFPageRef{725} /XYZ 252.639 476.289 null] >>}
\MusicPDFWriteObject{14635}{<< /D [\MusicPDFPageRef{725} /XYZ 201.415 325.854 null] >>}
\MusicPDFWriteObject{14636}{<< /D [\MusicPDFPageRef{725} /XYZ 249.157 258.466 null] >>}
\MusicPDFWriteObject{14637}{<< /D [\MusicPDFPageRef{725} /XYZ 194.074 204.339 null] >>}
\MusicPDFWriteObject{14638}{<< /D [\MusicPDFPageRef{725} /XYZ 72 137.232 null] >>}
\MusicPDFWriteObject{14639}{<< /D [\MusicPDFPageRef{726} /XYZ 280.603 460.529 null] >>}
\MusicPDFWriteObject{14640}{<< /D [\MusicPDFPageRef{720} /XYZ 257.95 153.994 null] >>}
\MusicPDFWriteObject{14641}{<< /D [\MusicPDFPageRef{726} /XYZ 245.016 376.07 null] >>}
\MusicPDFWriteObject{14642}{<< /D [\MusicPDFPageRef{726} /XYZ 221.176 136.753 null] >>}
\MusicPDFWriteObject{14643}{<< /D [\MusicPDFPageRef{727} /XYZ 242.888 653.412 null] >>}
\MusicPDFWriteObject{14644}{<< /D [\MusicPDFPageRef{727} /XYZ 228.13 610.051 null] >>}
\MusicPDFWriteObject{14645}{<< /D [\MusicPDFPageRef{727} /XYZ 199.256 523.234 null] >>}
\MusicPDFWriteObject{14646}{<< /D [\MusicPDFPageRef{727} /XYZ 219.42 464.944 null] >>}
\MusicPDFWriteObject{14647}{<< /D [\MusicPDFPageRef{727} /XYZ 172.363 287.802 null] >>}
\MusicPDFWriteObject{14648}{<< /D [\MusicPDFPageRef{727} /XYZ 171.106 227.52 null] >>}
\MusicPDFWriteObject{14649}{<< /D [\MusicPDFPageRef{728} /XYZ 220.556 520.646 null] >>}
\MusicPDFWriteObject{14650}{<< /D [\MusicPDFPageRef{728} /XYZ 238.107 187.284 null] >>}
\MusicPDFWriteObject{14651}{<< /D [\MusicPDFPageRef{721} /XYZ 267.771 560.477 null] >>}
\MusicPDFWriteObject{14652}{<< /D [\MusicPDFPageRef{729} /XYZ 220.093 652.207 null] >>}
\MusicPDFWriteObject{14653}{<< /D [\MusicPDFPageRef{729} /XYZ 253.799 584.319 null] >>}
\MusicPDFWriteObject{14654}{<< /D [\MusicPDFPageRef{729} /XYZ 237.708 530.192 null] >>}
\MusicPDFWriteObject{14655}{<< /D [\MusicPDFPageRef{729} /XYZ 224.209 462.304 null] >>}
\MusicPDFWriteObject{14656}{<< /D [\MusicPDFPageRef{730} /XYZ 72 668.965 null] >>}
\MusicPDFWriteObject{14657}{<< /D [\MusicPDFPageRef{730} /XYZ 223.49 426.247 null] >>}
\MusicPDFWriteObject{14658}{<< /D [\MusicPDFPageRef{730} /XYZ 218.294 376.669 null] >>}
\MusicPDFWriteObject{14659}{<< /D [\MusicPDFPageRef{731} /XYZ 210.795 582.92 null] >>}
\MusicPDFWriteObject{14660}{<< /D [\MusicPDFPageRef{732} /XYZ 245.016 612.755 null] >>}
\MusicPDFWriteObject{14661}{<< /D [\MusicPDFPageRef{732} /XYZ 251.625 197.457 null] >>}
\MusicPDFWriteObject{14662}{<< /D [\MusicPDFPageRef{721} /XYZ 250.383 248.162 null] >>}
\MusicPDFWriteObject{14663}{<< /D [\MusicPDFPageRef{733} /XYZ 202.825 683.057 null] >>}
\MusicPDFWriteObject{14664}{<< /D [\MusicPDFPageRef{733} /XYZ 251.131 644.572 null] >>}
\MusicPDFWriteObject{14665}{<< /D [\MusicPDFPageRef{733} /XYZ 238.581 399.63 null] >>}
\MusicPDFWriteObject{14666}{<< /D [\MusicPDFPageRef{733} /XYZ 280.603 179.46 null] >>}
\MusicPDFWriteObject{14667}{<< /D [\MusicPDFPageRef{734} /XYZ 185.126 502.056 null] >>}
\MusicPDFWriteObject{14668}{<< /D [\MusicPDFPageRef{735} /XYZ 246.859 149.126 null] >>}
\MusicPDFWriteObject{14669}{<< /D [\MusicPDFPageRef{735} /XYZ 258.85 105.821 null] >>}
\MusicPDFWriteObject{14670}{<< /D [\MusicPDFPageRef{736} /XYZ 238.959 432.744 null] >>}
\MusicPDFWriteObject{14671}{<< /D [\MusicPDFPageRef{736} /XYZ 253.885 370.015 null] >>}
\MusicPDFWriteObject{14672}{<< /D [\MusicPDFPageRef{722} /XYZ 257.429 580.691 null] >>}
\MusicPDFWriteObject{14673}{<< /D [\MusicPDFPageRef{723} /XYZ 270.026 369.125 null] >>}
\MusicPDFWriteObject{14674}{<< /D [\MusicPDFPageRef{723} /XYZ 219.255 306.396 null] >>}
\MusicPDFWriteObject{14675}{<< /D [\MusicPDFPageRef{723} /XYZ 241.635 193.575 null] >>}
\MusicPDFWriteObject{14676}{<< /D [\MusicPDFPageRef{724} /XYZ 212.08 577.601 null] >>}
\MusicPDFWriteObject{14677}{<< /D [\MusicPDFPageRef{737} /XYZ 244.693 209.899 null] >>}
\MusicPDFWriteObject{14678}{<< /D [\MusicPDFPageRef{738} /XYZ 277.094 526.646 null] >>}
\MusicPDFWriteObject{14679}{<< /D [\MusicPDFPageRef{738} /XYZ 254.222 455.51 null] >>}
\MusicPDFWriteObject{14680}{<< /D [\MusicPDFPageRef{738} /XYZ 237.289 235.622 null] >>}
\MusicPDFWriteObject{14681}{<< /D [\MusicPDFPageRef{738} /XYZ 258.619 180.588 null] >>}
\MusicPDFWriteObject{14682}{<< /D [\MusicPDFPageRef{739} /XYZ 231.539 664.371 null] >>}
\MusicPDFWriteObject{14683}{<< /D [\MusicPDFPageRef{739} /XYZ 263.467 337.72 null] >>}
\MusicPDFWriteObject{14684}{<< /D [\MusicPDFPageRef{739} /XYZ 257.06 301.293 null] >>}
\MusicPDFWriteObject{14685}{<< /D [\MusicPDFPageRef{740} /XYZ 235.555 679.063 null] >>}
\MusicPDFWriteObject{14686}{<< /D [\MusicPDFPageRef{740} /XYZ 252.343 603.752 null] >>}
\MusicPDFWriteObject{14687}{<< /D [\MusicPDFPageRef{740} /XYZ 246.094 457.752 null] >>}
\MusicPDFWriteObject{14688}{<< /D [\MusicPDFPageRef{740} /XYZ 262.082 389.406 null] >>}
\MusicPDFWriteObject{14689}{<< /D [\MusicPDFPageRef{740} /XYZ 251.633 335.932 null] >>}
\MusicPDFWriteObject{14690}{<< /D [\MusicPDFPageRef{740} /XYZ 201.136 263.275 null] >>}
\MusicPDFWriteObject{14691}{<< /D [\MusicPDFPageRef{740} /XYZ 255.359 195.903 null] >>}
\MusicPDFWriteObject{14692}{<< /D [\MusicPDFPageRef{741} /XYZ 223.562 640.261 null] >>}
\MusicPDFWriteObject{14693}{<< /D [\MusicPDFPageRef{741} /XYZ 280.803 223.866 null] >>}
\MusicPDFWriteObject{14694}{<< /D [\MusicPDFPageRef{746} /XYZ 252.275 519.622 null] >>}
\MusicPDFWriteObject{14695}{<< /D [\MusicPDFPageRef{747} /XYZ 243.969 583.69 null] >>}
\MusicPDFWriteObject{14696}{<< /D [\MusicPDFPageRef{747} /XYZ 255.5 546.224 null] >>}
\MusicPDFWriteObject{14697}{<< /D [\MusicPDFPageRef{747} /XYZ 242.785 496.646 null] >>}
\MusicPDFWriteObject{14698}{<< /D [\MusicPDFPageRef{746} /XYZ 282.124 480.793 null] >>}
\MusicPDFWriteObject{14699}{<< /D [\MusicPDFPageRef{746} /XYZ 271.695 444.365 null] >>}
\MusicPDFWriteObject{14700}{<< /D [\MusicPDFPageRef{746} /XYZ 216.868 336.184 null] >>}
\MusicPDFWriteObject{14701}{<< /D [\MusicPDFPageRef{746} /XYZ 219.187 278.317 null] >>}
\MusicPDFWriteObject{14702}{<< /D [\MusicPDFPageRef{746} /XYZ 255.835 224.897 null] >>}
\MusicPDFWriteObject{14703}{<< /D [\MusicPDFPageRef{746} /XYZ 282.621 99.233 null] >>}
\MusicPDFWriteObject{14704}{<< /D [\MusicPDFPageRef{747} /XYZ 254.082 669.906 null] >>}
\MusicPDFWriteObject{14705}{<< /D [\MusicPDFPageRef{747} /XYZ 281.607 620.117 null] >>}
\MusicPDFWriteObject{14706}{<< /D [\MusicPDFPageRef{783} /XYZ 169.142 156.815 null] >>}
\MusicPDFWriteObject{14707}{<< /D [\MusicPDFPageRef{785} /XYZ 246.44 99.233 null] >>}
\MusicPDFWriteObject{14708}{<< /D [\MusicPDFPageRef{786} /XYZ 264.993 580.875 null] >>}
\MusicPDFWriteObject{14709}{<< /D [\MusicPDFPageRef{786} /XYZ 259.989 244.444 null] >>}
\MusicPDFWriteObject{14710}{<< /D [\MusicPDFPageRef{788} /XYZ 72 665.284 null] >>}
\MusicPDFWriteObject{14711}{<< /D [\MusicPDFPageRef{788} /XYZ 72 548.886 null] >>}
\MusicPDFWriteObject{14712}{<< /D [\MusicPDFPageRef{789} /XYZ 168.227 323.534 null] >>}
\MusicPDFWriteObject{14713}{<< /D [\MusicPDFPageRef{789} /XYZ 261.458 206.477 null] >>}
\MusicPDFWriteObject{14714}{<< /D [\MusicPDFPageRef{790} /XYZ 124.449 198.685 null] >>}
\MusicPDFWriteObject{14715}{<< /D [\MusicPDFPageRef{790} /XYZ 247.912 99.233 null] >>}
\MusicPDFWriteObject{14716}{<< /D [\MusicPDFPageRef{791} /XYZ 261.458 531.297 null] >>}
\MusicPDFWriteObject{14717}{<< /D [\MusicPDFPageRef{784} /XYZ 270.348 596.155 null] >>}
\MusicPDFWriteObject{14718}{<< /D [\MusicPDFPageRef{791} /XYZ 94.677 125.218 null] >>}
\MusicPDFWriteObject{14719}{<< /D [\MusicPDFPageRef{784} /XYZ 221.52 380.309 null] >>}
\MusicPDFWriteObject{14720}{<< /D [\MusicPDFPageRef{785} /XYZ 72 673.193 null] >>}
\MusicPDFWriteObject{14721}{<< /D [\MusicPDFPageRef{785} /XYZ 72 571.314 null] >>}
\MusicPDFWriteObject{14722}{<< /D [\MusicPDFPageRef{785} /XYZ 247.773 353.352 null] >>}
\MusicPDFWriteObject{14723}{<< /D [\MusicPDFPageRef{785} /XYZ 237.375 270.156 null] >>}
\MusicPDFWriteObject{14724}{<< /D [\MusicPDFPageRef{785} /XYZ 235.375 233.729 null] >>}
\MusicPDFWriteObject{14725}{<< /D [\MusicPDFPageRef{785} /XYZ 258.839 153.482 null] >>}
\MusicPDFWriteObject{14726}{<< /D [\MusicPDFPageRef{747} /XYZ 204.365 284.21 null] >>}
\MusicPDFWriteObject{14727}{<< /D [\MusicPDFPageRef{747} /XYZ 242.354 106.069 null] >>}
\MusicPDFWriteObject{14728}{<< /D [\MusicPDFPageRef{748} /XYZ 247.606 537 null] >>}
\MusicPDFWriteObject{14729}{<< /D [\MusicPDFPageRef{748} /XYZ 242.117 493.445 null] >>}
\MusicPDFWriteObject{14730}{<< /D [\MusicPDFPageRef{748} /XYZ 217.556 419.159 null] >>}
\MusicPDFWriteObject{14731}{<< /D [\MusicPDFPageRef{748} /XYZ 232.624 169.547 null] >>}
\MusicPDFWriteObject{14732}{<< /D [\MusicPDFPageRef{749} /XYZ 249.236 577.381 null] >>}
\MusicPDFWriteObject{14733}{<< /D [\MusicPDFPageRef{749} /XYZ 208.303 500.7 null] >>}
\MusicPDFWriteObject{14734}{<< /D [\MusicPDFPageRef{749} /XYZ 252.385 326.296 null] >>}
\MusicPDFWriteObject{14735}{<< /D [\MusicPDFPageRef{751} /XYZ 191.325 473.76 null] >>}
\MusicPDFWriteObject{14736}{<< /D [\MusicPDFPageRef{751} /XYZ 83.14 366.896 null] >>}
\MusicPDFWriteObject{14737}{<< /D [\MusicPDFPageRef{751} /XYZ 240.707 281.427 null] >>}
\MusicPDFWriteObject{14738}{<< /D [\MusicPDFPageRef{751} /XYZ 217.594 202.929 null] >>}
\MusicPDFWriteObject{14739}{<< /D [\MusicPDFPageRef{752} /XYZ 203.244 669.109 null] >>}
\MusicPDFWriteObject{14740}{<< /D [\MusicPDFPageRef{752} /XYZ 143.64 459.432 null] >>}
\MusicPDFWriteObject{14741}{<< /D [\MusicPDFPageRef{752} /XYZ 145.244 383.329 null] >>}
\MusicPDFWriteObject{14742}{<< /D [\MusicPDFPageRef{752} /XYZ 222.192 108.697 null] >>}
\MusicPDFWriteObject{14743}{<< /D [\MusicPDFPageRef{753} /XYZ 224.337 603.104 null] >>}
\MusicPDFWriteObject{14744}{<< /D [\MusicPDFPageRef{753} /XYZ 164.063 566.676 null] >>}
\MusicPDFWriteObject{14745}{<< /D [\MusicPDFPageRef{749} /XYZ 196.776 210.964 null] >>}
\MusicPDFWriteObject{14746}{<< /D [\MusicPDFPageRef{753} /XYZ 217.886 503.947 null] >>}
\MusicPDFWriteObject{14747}{<< /D [\MusicPDFPageRef{753} /XYZ 159.575 172.041 null] >>}
\MusicPDFWriteObject{14748}{<< /D [\MusicPDFPageRef{753} /XYZ 215.825 135.614 null] >>}
\MusicPDFWriteObject{14749}{<< /D [\MusicPDFPageRef{754} /XYZ 194.92 669.906 null] >>}
\MusicPDFWriteObject{14750}{<< /D [\MusicPDFPageRef{754} /XYZ 231.267 606.953 null] >>}
\MusicPDFWriteObject{14751}{<< /D [\MusicPDFPageRef{754} /XYZ 229.117 570.303 null] >>}
\MusicPDFWriteObject{14752}{<< /D [\MusicPDFPageRef{754} /XYZ 204.024 517.76 null] >>}
\MusicPDFWriteObject{14753}{<< /D [\MusicPDFPageRef{754} /XYZ 241.937 311.683 null] >>}
\MusicPDFWriteObject{14754}{<< /D [\MusicPDFPageRef{754} /XYZ 244.526 262.105 null] >>}
\MusicPDFWriteObject{14755}{<< /D [\MusicPDFPageRef{755} /XYZ 231.453 683.057 null] >>}
\MusicPDFWriteObject{14756}{<< /D [\MusicPDFPageRef{749} /XYZ 145.327 106.439 null] >>}
\MusicPDFWriteObject{14757}{<< /D [\MusicPDFPageRef{755} /XYZ 228.967 251.798 null] >>}
\MusicPDFWriteObject{14758}{<< /D [\MusicPDFPageRef{756} /XYZ 240.055 459.494 null] >>}
\MusicPDFWriteObject{14759}{<< /D [\MusicPDFPageRef{756} /XYZ 80.248 365.099 null] >>}
\MusicPDFWriteObject{14760}{<< /D [\MusicPDFPageRef{756} /XYZ 179.267 307.074 null] >>}
\MusicPDFWriteObject{14761}{<< /D [\MusicPDFPageRef{756} /XYZ 72 240.39 null] >>}
\MusicPDFWriteObject{14762}{<< /D [\MusicPDFPageRef{756} /XYZ 171.083 182.346 null] >>}
\MusicPDFWriteObject{14763}{<< /D [\MusicPDFPageRef{750} /XYZ 206.351 683.057 null] >>}
\MusicPDFWriteObject{14764}{<< /D [\MusicPDFPageRef{750} /XYZ 278.158 613.767 null] >>}
\MusicPDFWriteObject{14765}{<< /D [\MusicPDFPageRef{750} /XYZ 243.103 464.091 null] >>}
\MusicPDFWriteObject{14766}{<< /D [\MusicPDFPageRef{750} /XYZ 217.648 427.663 null] >>}
\MusicPDFWriteObject{14767}{<< /D [\MusicPDFPageRef{751} /XYZ 121.386 553.618 null] >>}
\MusicPDFWriteObject{14768}{<< /D [\MusicPDFPageRef{751} /XYZ 78.997 503.07 null] >>}
\MusicPDFWriteObject{14769}{<< /D [\MusicPDFPageRef{757} /XYZ 262.229 242.781 null] >>}
\MusicPDFWriteObject{14770}{<< /D [\MusicPDFPageRef{759} /XYZ 268.56 577.993 null] >>}
\MusicPDFWriteObject{14771}{<< /D [\MusicPDFPageRef{759} /XYZ 278.135 407.759 null] >>}
\MusicPDFWriteObject{14772}{<< /D [\MusicPDFPageRef{759} /XYZ 248.612 274.157 null] >>}
\MusicPDFWriteObject{14773}{<< /D [\MusicPDFPageRef{759} /XYZ 251.059 184.717 null] >>}
\MusicPDFWriteObject{14774}{<< /D [\MusicPDFPageRef{759} /XYZ 222.766 124.118 null] >>}
\MusicPDFWriteObject{14775}{<< /D [\MusicPDFPageRef{760} /XYZ 267.087 612.252 null] >>}
\MusicPDFWriteObject{14776}{<< /D [\MusicPDFPageRef{760} /XYZ 254.975 557.796 null] >>}
\MusicPDFWriteObject{14777}{<< /D [\MusicPDFPageRef{757} /XYZ 251.11 118.104 null] >>}
\MusicPDFWriteObject{14778}{<< /D [\MusicPDFPageRef{758} /XYZ 274.342 604.152 null] >>}
\MusicPDFWriteObject{14779}{<< /D [\MusicPDFPageRef{758} /XYZ 273.235 490.219 null] >>}
\MusicPDFWriteObject{14780}{<< /D [\MusicPDFPageRef{758} /XYZ 241.038 335.974 null] >>}
\MusicPDFWriteObject{14781}{<< /D [\MusicPDFPageRef{758} /XYZ 248.944 264.257 null] >>}
\MusicPDFWriteObject{14782}{<< /D [\MusicPDFPageRef{758} /XYZ 224.409 227.78 null] >>}
\MusicPDFWriteObject{14783}{<< /D [\MusicPDFPageRef{759} /XYZ 263.483 656.755 null] >>}
\MusicPDFWriteObject{14784}{<< /D [\MusicPDFPageRef{759} /XYZ 282.687 614.421 null] >>}
\MusicPDFWriteObject{14785}{<< /D [\MusicPDFPageRef{760} /XYZ 268.63 99.233 null] >>}
\MusicPDFWriteObject{14786}{<< /D [\MusicPDFPageRef{761} /XYZ 272.975 143.797 null] >>}
\MusicPDFWriteObject{14787}{<< /D [\MusicPDFPageRef{762} /XYZ 239.7 535.796 null] >>}
\MusicPDFWriteObject{14788}{<< /D [\MusicPDFPageRef{763} /XYZ 260.662 530.872 null] >>}
\MusicPDFWriteObject{14789}{<< /D [\MusicPDFPageRef{765} /XYZ 196.125 656.512 null] >>}
\MusicPDFWriteObject{14790}{<< /D [\MusicPDFPageRef{765} /XYZ 236.825 522.779 null] >>}
\MusicPDFWriteObject{14791}{<< /D [\MusicPDFPageRef{765} /XYZ 256.086 233.321 null] >>}
\MusicPDFWriteObject{14792}{<< /D [\MusicPDFPageRef{768} /XYZ 277.204 333.117 null] >>}
\MusicPDFWriteObject{14793}{<< /D [\MusicPDFPageRef{768} /XYZ 264.693 296.689 null] >>}
\MusicPDFWriteObject{14794}{<< /D [\MusicPDFPageRef{769} /XYZ 230.446 669.906 null] >>}
\MusicPDFWriteObject{14795}{<< /D [\MusicPDFPageRef{769} /XYZ 237.236 521.78 null] >>}
\MusicPDFWriteObject{14796}{<< /D [\MusicPDFPageRef{769} /XYZ 211.015 268.099 null] >>}
\MusicPDFWriteObject{14797}{<< /D [\MusicPDFPageRef{769} /XYZ 235.899 164.227 null] >>}
\MusicPDFWriteObject{14798}{<< /D [\MusicPDFPageRef{770} /XYZ 232.684 564.567 null] >>}
\MusicPDFWriteObject{14799}{<< /D [\MusicPDFPageRef{770} /XYZ 270.545 432.563 null] >>}
\MusicPDFWriteObject{14800}{<< /D [\MusicPDFPageRef{771} /XYZ 130.961 593.133 null] >>}
\MusicPDFWriteObject{14801}{<< /D [\MusicPDFPageRef{771} /XYZ 165.911 539.157 null] >>}
\MusicPDFWriteObject{14802}{<< /D [\MusicPDFPageRef{765} /XYZ 269.668 135.661 null] >>}
\MusicPDFWriteObject{14803}{<< /D [\MusicPDFPageRef{771} /XYZ 224.453 202.851 null] >>}
\MusicPDFWriteObject{14804}{<< /D [\MusicPDFPageRef{772} /XYZ 254.779 533.257 null] >>}
\MusicPDFWriteObject{14805}{<< /D [\MusicPDFPageRef{773} /XYZ 209.143 490.525 null] >>}
\MusicPDFWriteObject{14806}{<< /D [\MusicPDFPageRef{773} /XYZ 266.246 349.158 null] >>}
\MusicPDFWriteObject{14807}{<< /D [\MusicPDFPageRef{773} /XYZ 254.348 312.244 null] >>}
\MusicPDFWriteObject{14808}{<< /D [\MusicPDFPageRef{773} /XYZ 229.073 232.922 null] >>}
\MusicPDFWriteObject{14809}{<< /D [\MusicPDFPageRef{773} /XYZ 241.736 191.113 null] >>}
\MusicPDFWriteObject{14810}{<< /D [\MusicPDFPageRef{765} /XYZ 249.358 99.233 null] >>}
\MusicPDFWriteObject{14811}{<< /D [\MusicPDFPageRef{766} /XYZ 275.466 496.568 null] >>}
\MusicPDFWriteObject{14812}{<< /D [\MusicPDFPageRef{766} /XYZ 248.337 412.846 null] >>}
\MusicPDFWriteObject{14813}{<< /D [\MusicPDFPageRef{766} /XYZ 265.196 194.401 null] >>}
\MusicPDFWriteObject{14814}{<< /D [\MusicPDFPageRef{766} /XYZ 238.752 153.979 null] >>}
\MusicPDFWriteObject{14815}{<< /D [\MusicPDFPageRef{768} /XYZ 265.518 574.204 null] >>}
\MusicPDFWriteObject{14816}{<< /D [\MusicPDFPageRef{768} /XYZ 257.073 501.048 null] >>}
\MusicPDFWriteObject{14817}{<< /D [\MusicPDFPageRef{776} /XYZ 262.413 483.585 null] >>}
\MusicPDFWriteObject{14818}{<< /D [\MusicPDFPageRef{776} /XYZ 245.76 443.261 null] >>}
\MusicPDFWriteObject{14819}{<< /D [\MusicPDFPageRef{776} /XYZ 233.765 351.636 null] >>}
\MusicPDFWriteObject{14820}{<< /D [\MusicPDFPageRef{778} /XYZ 244.1 641.381 null] >>}
\MusicPDFWriteObject{14821}{<< /D [\MusicPDFPageRef{778} /XYZ 233.335 578.652 null] >>}
\MusicPDFWriteObject{14822}{<< /D [\MusicPDFPageRef{778} /XYZ 175.247 170.05 null] >>}
\MusicPDFWriteObject{14823}{<< /D [\MusicPDFPageRef{779} /XYZ 186.617 345.599 null] >>}
\MusicPDFWriteObject{14824}{<< /D [\MusicPDFPageRef{782} /XYZ 72 347.019 null] >>}
\MusicPDFWriteObject{14825}{<< /D [\MusicPDFPageRef{783} /XYZ 239.773 473.32 null] >>}
\MusicPDFWriteObject{14826}{<< /D [\MusicPDFPageRef{792} /XYZ 248.814 250.083 null] >>}
\MusicPDFWriteObject{14827}{<< /D [\MusicPDFPageRef{792} /XYZ 206.642 210.313 null] >>}
\MusicPDFWriteObject{14828}{<< /D [\MusicPDFPageRef{793} /XYZ 271.24 635.227 null] >>}
\MusicPDFWriteObject{14829}{<< /D [\MusicPDFPageRef{793} /XYZ 237.514 542.672 null] >>}
\MusicPDFWriteObject{14830}{<< /D [\MusicPDFPageRef{793} /XYZ 232.561 233.676 null] >>}
\MusicPDFWriteObject{14831}{<< /D [\MusicPDFPageRef{794} /XYZ 270.941 640.108 null] >>}
\MusicPDFWriteObject{14832}{<< /D [\MusicPDFPageRef{794} /XYZ 256.869 595.356 null] >>}
\MusicPDFWriteObject{14833}{<< /D [\MusicPDFPageRef{794} /XYZ 241.661 478.641 null] >>}
\MusicPDFWriteObject{14834}{<< /D [\MusicPDFPageRef{796} /XYZ 261.612 470.011 null] >>}
\MusicPDFWriteObject{14835}{<< /D [\MusicPDFPageRef{794} /XYZ 271.24 288.846 null] >>}
\MusicPDFWriteObject{14836}{<< /D [\MusicPDFPageRef{794} /XYZ 163.527 175.647 null] >>}
\MusicPDFWriteObject{14837}{<< /D [\MusicPDFPageRef{794} /XYZ 213.712 135.877 null] >>}
\MusicPDFWriteObject{14838}{<< /D [\MusicPDFPageRef{795} /XYZ 177.097 669.906 null] >>}
\MusicPDFWriteObject{14839}{<< /D [\MusicPDFPageRef{795} /XYZ 246.39 580.875 null] >>}
\MusicPDFWriteObject{14840}{<< /D [\MusicPDFPageRef{795} /XYZ 246.39 544.448 null] >>}
\MusicPDFWriteObject{14841}{<< /D [\MusicPDFPageRef{795} /XYZ 209.669 502.436 null] >>}
\MusicPDFWriteObject{14842}{<< /D [\MusicPDFPageRef{795} /XYZ 226.699 413.329 null] >>}
\MusicPDFWriteObject{14843}{<< /D [\MusicPDFPageRef{798} /XYZ 265.308 678.154 null] >>}
\MusicPDFWriteObject{14844}{<< /D [\MusicPDFPageRef{800} /XYZ 243.489 587.13 null] >>}
\MusicPDFWriteObject{14845}{<< /D [\MusicPDFPageRef{800} /XYZ 147.507 534.556 null] >>}
\MusicPDFWriteObject{14846}{<< /D [\MusicPDFPageRef{800} /XYZ 130.82 252.931 null] >>}
\MusicPDFWriteObject{14847}{<< /D [\MusicPDFPageRef{798} /XYZ 91.065 570.405 null] >>}
\MusicPDFWriteObject{14848}{<< /D [\MusicPDFPageRef{798} /XYZ 72 282.327 null] >>}
\MusicPDFWriteObject{14849}{<< /D [\MusicPDFPageRef{798} /XYZ 120.867 234.4 null] >>}
\MusicPDFWriteObject{14850}{<< /D [\MusicPDFPageRef{798} /XYZ 72 156.474 null] >>}
\MusicPDFWriteObject{14851}{<< /D [\MusicPDFPageRef{798} /XYZ 72 126.868 null] >>}
\MusicPDFWriteObject{14852}{<< /D [\MusicPDFPageRef{799} /XYZ 230.999 284.717 null] >>}
\MusicPDFWriteObject{14853}{<< /D [\MusicPDFPageRef{799} /XYZ 178.927 151.101 null] >>}
\MusicPDFWriteObject{14854}{<< /D [\MusicPDFPageRef{800} /XYZ 75.554 666.88 null] >>}
\MusicPDFWriteObject{14855}{<< /D [\MusicPDFPageRef{801} /XYZ 255.856 494 null] >>}
\MusicPDFWriteObject{14856}{<< /D [\MusicPDFPageRef{802} /XYZ 252.161 263.547 null] >>}
\MusicPDFWriteObject{14857}{<< /D [\MusicPDFPageRef{803} /XYZ 194.815 492.021 null] >>}
\MusicPDFWriteObject{14858}{<< /D [\MusicPDFPageRef{803} /XYZ 172.189 445.251 null] >>}
\MusicPDFWriteObject{14859}{<< /D [\MusicPDFPageRef{803} /XYZ 72 270.589 null] >>}
\MusicPDFWriteObject{14860}{<< /D [\MusicPDFPageRef{803} /XYZ 72 186.88 null] >>}
\MusicPDFWriteObject{14861}{<< /D [\MusicPDFPageRef{804} /XYZ 159.597 662.91 null] >>}
\MusicPDFWriteObject{14862}{<< /D [\MusicPDFPageRef{805} /XYZ 270.641 564.755 null] >>}
\MusicPDFWriteObject{14863}{<< /D [\MusicPDFPageRef{812} /XYZ 221.757 675.68 null] >>}
\MusicPDFWriteObject{14864}{<< /D [\MusicPDFPageRef{812} /XYZ 252.15 233.131 null] >>}
\MusicPDFWriteObject{14865}{<< /D [\MusicPDFPageRef{812} /XYZ 202.96 169.512 null] >>}
\MusicPDFWriteObject{14866}{<< /D [\MusicPDFPageRef{813} /XYZ 187.357 287.648 null] >>}
\MusicPDFWriteObject{14867}{<< /D [\MusicPDFPageRef{813} /XYZ 171.021 157.747 null] >>}
\MusicPDFWriteObject{14868}{<< /D [\MusicPDFPageRef{805} /XYZ 235.01 515.177 null] >>}
\MusicPDFWriteObject{14869}{<< /D [\MusicPDFPageRef{806} /XYZ 249.203 345.872 null] >>}
\MusicPDFWriteObject{14870}{<< /D [\MusicPDFPageRef{806} /XYZ 230.099 295.114 null] >>}
\MusicPDFWriteObject{14871}{<< /D [\MusicPDFPageRef{806} /XYZ 216.657 191.411 null] >>}
\MusicPDFWriteObject{14872}{<< /D [\MusicPDFPageRef{807} /XYZ 223.818 339.882 null] >>}
\MusicPDFWriteObject{14873}{<< /D [\MusicPDFPageRef{807} /XYZ 240.582 290.304 null] >>}
\MusicPDFWriteObject{14874}{<< /D [\MusicPDFPageRef{811} /XYZ 204.4 496.97 null] >>}
\MusicPDFWriteObject{14875}{<< /D [\MusicPDFPageRef{811} /XYZ 244.243 99.233 null] >>}
\MusicPDFWriteObject{14876}{<< /D [\MusicPDFPageRef{816} /XYZ 271.139 656.755 null] >>}
\MusicPDFWriteObject{14877}{<< /D [\MusicPDFPageRef{816} /XYZ 264.03 559.627 null] >>}
\MusicPDFWriteObject{14878}{<< /D [\MusicPDFPageRef{818} /XYZ 256.912 353.378 null] >>}
\MusicPDFWriteObject{14879}{<< /D [\MusicPDFPageRef{818} /XYZ 200.891 204.487 null] >>}
\MusicPDFWriteObject{14880}{<< /D [\MusicPDFPageRef{818} /XYZ 170.901 161.871 null] >>}
\MusicPDFWriteObject{14881}{<< /D [\MusicPDFPageRef{819} /XYZ 266.118 621.711 null] >>}
\MusicPDFWriteObject{14882}{<< /D [\MusicPDFPageRef{819} /XYZ 151.289 509.018 null] >>}
\MusicPDFWriteObject{14883}{<< /D [\MusicPDFPageRef{819} /XYZ 188.477 170.05 null] >>}
\MusicPDFWriteObject{14884}{<< /D [\MusicPDFPageRef{820} /XYZ 72 324.937 null] >>}
\MusicPDFWriteObject{14885}{<< /D [\MusicPDFPageRef{820} /XYZ 72 232.715 null] >>}
\MusicPDFWriteObject{14886}{<< /D [\MusicPDFPageRef{820} /XYZ 85.35 133.173 null] >>}
\MusicPDFWriteObject{14887}{<< /D [\MusicPDFPageRef{829} /XYZ 272.613 520.137 null] >>}
\MusicPDFWriteObject{14888}{<< /D [\MusicPDFPageRef{831} /XYZ 186.156 617.303 null] >>}
\MusicPDFWriteObject{14889}{<< /D [\MusicPDFPageRef{831} /XYZ 251.243 560.841 null] >>}
\MusicPDFWriteObject{14890}{<< /D [\MusicPDFPageRef{831} /XYZ 270.461 461.14 null] >>}
\MusicPDFWriteObject{14891}{<< /D [\MusicPDFPageRef{831} /XYZ 266.516 334.567 null] >>}
\MusicPDFWriteObject{14892}{<< /D [\MusicPDFPageRef{833} /XYZ 263.767 497.105 null] >>}
\MusicPDFWriteObject{14893}{<< /D [\MusicPDFPageRef{838} /XYZ 141.242 646.629 null] >>}
\MusicPDFWriteObject{14894}{<< /D [\MusicPDFPageRef{838} /XYZ 238.042 597.051 null] >>}
\MusicPDFWriteObject{14895}{<< /D [\MusicPDFPageRef{838} /XYZ 164.675 560.623 null] >>}
\MusicPDFWriteObject{14896}{<< /D [\MusicPDFPageRef{838} /XYZ 184.077 372.691 null] >>}
\MusicPDFWriteObject{14897}{<< /D [\MusicPDFPageRef{838} /XYZ 196.69 335.034 null] >>}
\MusicPDFWriteObject{14898}{<< /D [\MusicPDFPageRef{838} /XYZ 72 301.895 null] >>}
\MusicPDFWriteObject{14899}{<< /D [\MusicPDFPageRef{838} /XYZ 236.655 188.555 null] >>}
\MusicPDFWriteObject{14900}{<< /D [\MusicPDFPageRef{838} /XYZ 219.536 138.977 null] >>}
\MusicPDFWriteObject{14901}{<< /D [\MusicPDFPageRef{839} /XYZ 222.063 630.453 null] >>}
\MusicPDFWriteObject{14902}{<< /D [\MusicPDFPageRef{839} /XYZ 190.943 594.026 null] >>}
\MusicPDFWriteObject{14903}{<< /D [\MusicPDFPageRef{833} /XYZ 267.928 440.369 null] >>}
\MusicPDFWriteObject{14904}{<< /D [\MusicPDFPageRef{839} /XYZ 151.796 557.598 null] >>}
\MusicPDFWriteObject{14905}{<< /D [\MusicPDFPageRef{839} /XYZ 120.821 484.743 null] >>}
\MusicPDFWriteObject{14906}{<< /D [\MusicPDFPageRef{839} /XYZ 215.017 238.397 null] >>}
\MusicPDFWriteObject{14907}{<< /D [\MusicPDFPageRef{839} /XYZ 246.967 201.969 null] >>}
\MusicPDFWriteObject{14908}{<< /D [\MusicPDFPageRef{840} /XYZ 72 633.741 null] >>}
\MusicPDFWriteObject{14909}{<< /D [\MusicPDFPageRef{840} /XYZ 192.798 488.421 null] >>}
\MusicPDFWriteObject{14910}{<< /D [\MusicPDFPageRef{840} /XYZ 242.263 362.963 null] >>}
\MusicPDFWriteObject{14911}{<< /D [\MusicPDFPageRef{840} /XYZ 135.214 273.933 null] >>}
\MusicPDFWriteObject{14912}{<< /D [\MusicPDFPageRef{840} /XYZ 182.273 211.203 null] >>}
\MusicPDFWriteObject{14913}{<< /D [\MusicPDFPageRef{843} /XYZ 256.641 551.441 null] >>}
\MusicPDFWriteObject{14914}{<< /D [\MusicPDFPageRef{833} /XYZ 259.641 179.948 null] >>}
\MusicPDFWriteObject{14915}{<< /D [\MusicPDFPageRef{835} /XYZ 211.334 359.753 null] >>}
\MusicPDFWriteObject{14916}{<< /D [\MusicPDFPageRef{836} /XYZ 249.579 519.463 null] >>}
\MusicPDFWriteObject{14917}{<< /D [\MusicPDFPageRef{836} /XYZ 189.074 382.664 null] >>}
\MusicPDFWriteObject{14918}{<< /D [\MusicPDFPageRef{836} /XYZ 206.186 99.233 null] >>}
\MusicPDFWriteObject{14919}{<< /D [\MusicPDFPageRef{837} /XYZ 222.281 656.755 null] >>}
\MusicPDFWriteObject{14920}{<< /D [\MusicPDFPageRef{838} /XYZ 172.763 683.057 null] >>}
\MusicPDFWriteObject{14921}{<< /D [\MusicPDFPageRef{847} /XYZ 233.27 610.321 null] >>}
\MusicPDFWriteObject{14922}{<< /D [\MusicPDFPageRef{847} /XYZ 269.191 547.592 null] >>}
\MusicPDFWriteObject{14923}{<< /D [\MusicPDFPageRef{847} /XYZ 227.111 432.26 null] >>}
\MusicPDFWriteObject{14924}{<< /D [\MusicPDFPageRef{849} /XYZ 240.502 270.202 null] >>}
\MusicPDFWriteObject{14925}{<< /D [\MusicPDFPageRef{849} /XYZ 239.929 198.484 null] >>}
\MusicPDFWriteObject{14926}{<< /D [\MusicPDFPageRef{851} /XYZ 260.068 398.727 null] >>}
\MusicPDFWriteObject{14927}{<< /D [\MusicPDFPageRef{852} /XYZ 261.754 198.39 null] >>}
\MusicPDFWriteObject{14928}{<< /D [\MusicPDFPageRef{852} /XYZ 286.704 135.661 null] >>}
\MusicPDFWriteObject{14929}{<< /D [\MusicPDFPageRef{853} /XYZ 245.968 544.922 null] >>}
\MusicPDFWriteObject{14930}{<< /D [\MusicPDFPageRef{854} /XYZ 272.707 453.951 null] >>}
\MusicPDFWriteObject{14931}{<< /D [\MusicPDFPageRef{854} /XYZ 165.906 202.317 null] >>}
\MusicPDFWriteObject{14932}{<< /D [\MusicPDFPageRef{851} /XYZ 248.372 349.149 null] >>}
\MusicPDFWriteObject{14933}{<< /D [\MusicPDFPageRef{851} /XYZ 278.465 203.453 null] >>}
\MusicPDFWriteObject{14934}{<< /D [\MusicPDFPageRef{851} /XYZ 274.35 153.874 null] >>}
\MusicPDFWriteObject{14935}{<< /D [\MusicPDFPageRef{851} /XYZ 274.11 117.447 null] >>}
\MusicPDFWriteObject{14936}{<< /D [\MusicPDFPageRef{852} /XYZ 263.959 649.421 null] >>}
\MusicPDFWriteObject{14937}{<< /D [\MusicPDFPageRef{852} /XYZ 279.008 534.27 null] >>}
\MusicPDFWriteObject{14938}{<< /D [\MusicPDFPageRef{852} /XYZ 239.954 360.275 null] >>}
\MusicPDFWriteObject{14939}{<< /D [\MusicPDFPageRef{852} /XYZ 214.796 247.968 null] >>}
\MusicPDFWriteObject{14940}{<< /D [\MusicPDFPageRef{856} /XYZ 244.999 173.766 null] >>}
\MusicPDFWriteObject{14941}{<< /D [\MusicPDFPageRef{860} /XYZ 139.439 584.433 null] >>}
\MusicPDFWriteObject{14942}{<< /D [\MusicPDFPageRef{860} /XYZ 222.573 534.058 null] >>}
\MusicPDFWriteObject{14943}{<< /D [\MusicPDFPageRef{863} /XYZ 231.425 642.807 null] >>}
\MusicPDFWriteObject{14944}{<< /D [\MusicPDFPageRef{863} /XYZ 272.742 582.801 null] >>}
\MusicPDFWriteObject{14945}{<< /D [\MusicPDFPageRef{863} /XYZ 272.639 533.223 null] >>}
\MusicPDFWriteObject{14946}{<< /D [\MusicPDFPageRef{863} /XYZ 224.433 411.635 null] >>}
\MusicPDFWriteObject{14947}{<< /D [\MusicPDFPageRef{863} /XYZ 259.302 348.109 null] >>}
\MusicPDFWriteObject{14948}{<< /D [\MusicPDFPageRef{863} /XYZ 258.438 285.379 null] >>}
\MusicPDFWriteObject{14949}{<< /D [\MusicPDFPageRef{857} /XYZ 259.914 643.604 null] >>}
\MusicPDFWriteObject{14950}{<< /D [\MusicPDFPageRef{857} /XYZ 232.53 511.236 null] >>}
\MusicPDFWriteObject{14951}{<< /D [\MusicPDFPageRef{857} /XYZ 244.999 294.164 null] >>}
\MusicPDFWriteObject{14952}{<< /D [\MusicPDFPageRef{858} /XYZ 219.69 614.807 null] >>}
\MusicPDFWriteObject{14953}{<< /D [\MusicPDFPageRef{858} /XYZ 184.82 193.327 null] >>}
\MusicPDFWriteObject{14954}{<< /D [\MusicPDFPageRef{858} /XYZ 274.239 156.899 null] >>}
\MusicPDFWriteObject{14955}{<< /D [\MusicPDFPageRef{859} /XYZ 201.423 144.545 null] >>}
\MusicPDFWriteObject{14956}{<< /D [\MusicPDFPageRef{860} /XYZ 266.78 683.057 null] >>}
\MusicPDFWriteObject{14957}{<< /D [\MusicPDFPageRef{868} /XYZ 238.629 389.579 null] >>}
\MusicPDFWriteObject{14958}{<< /D [\MusicPDFPageRef{870} /XYZ 220.199 568.278 null] >>}
\MusicPDFWriteObject{14959}{<< /D [\MusicPDFPageRef{870} /XYZ 139.646 531.851 null] >>}
\MusicPDFWriteObject{14960}{<< /D [\MusicPDFPageRef{870} /XYZ 165.428 495.423 null] >>}
\MusicPDFWriteObject{14961}{<< /D [\MusicPDFPageRef{870} /XYZ 182.679 458.996 null] >>}
\MusicPDFWriteObject{14962}{<< /D [\MusicPDFPageRef{876} /XYZ 181.294 240.612 null] >>}
\MusicPDFWriteObject{14963}{<< /D [\MusicPDFPageRef{878} /XYZ 249.278 449.048 null] >>}
\MusicPDFWriteObject{14964}{<< /D [\MusicPDFPageRef{881} /XYZ 238.482 426.5 null] >>}
\MusicPDFWriteObject{14965}{<< /D [\MusicPDFPageRef{883} /XYZ 209.059 157.914 null] >>}
\MusicPDFWriteObject{14966}{<< /D [\MusicPDFPageRef{884} /XYZ 275.937 227.914 null] >>}
\MusicPDFWriteObject{14967}{<< /D [\MusicPDFPageRef{884} /XYZ 277.595 142.909 null] >>}
\MusicPDFWriteObject{14968}{<< /D [\MusicPDFPageRef{884} /XYZ 202.813 683.057 null] >>}
\MusicPDFWriteObject{14969}{<< /D [\MusicPDFPageRef{884} /XYZ 251.25 633.478 null] >>}
\MusicPDFWriteObject{14970}{<< /D [\MusicPDFPageRef{884} /XYZ 210.698 597.051 null] >>}
\MusicPDFWriteObject{14971}{<< /D [\MusicPDFPageRef{884} /XYZ 250.956 547.472 null] >>}
\MusicPDFWriteObject{14972}{<< /D [\MusicPDFPageRef{884} /XYZ 216.95 511.045 null] >>}
\MusicPDFWriteObject{14973}{<< /D [\MusicPDFPageRef{884} /XYZ 239.477 461.466 null] >>}
\MusicPDFWriteObject{14974}{<< /D [\MusicPDFPageRef{884} /XYZ 240.713 393.129 null] >>}
\MusicPDFWriteObject{14975}{<< /D [\MusicPDFPageRef{884} /XYZ 236.41 314.345 null] >>}
\MusicPDFWriteObject{14976}{<< /D [\MusicPDFPageRef{887} /XYZ 259.859 374.492 null] >>}
\MusicPDFWriteObject{14977}{<< /D [\MusicPDFPageRef{888} /XYZ 209.574 451.34 null] >>}
\MusicPDFWriteObject{14978}{<< /D [\MusicPDFPageRef{888} /XYZ 228.497 354.155 null] >>}
\MusicPDFWriteObject{14979}{<< /D [\MusicPDFPageRef{888} /XYZ 273.12 184.253 null] >>}
\MusicPDFWriteObject{14980}{<< /D [\MusicPDFPageRef{888} /XYZ 281.652 135.661 null] >>}
\MusicPDFWriteObject{14981}{<< /D [\MusicPDFPageRef{888} /XYZ 168.988 99.233 null] >>}
\MusicPDFWriteObject{14982}{<< /D [\MusicPDFPageRef{889} /XYZ 161.843 683.057 null] >>}
\MusicPDFWriteObject{14983}{<< /D [\MusicPDFPageRef{889} /XYZ 239.181 646.629 null] >>}
\MusicPDFWriteObject{14984}{<< /D [\MusicPDFPageRef{889} /XYZ 212.036 623.352 null] >>}
\MusicPDFWriteObject{14985}{<< /D [\MusicPDFPageRef{890} /XYZ 274.625 669.906 null] >>}
\MusicPDFWriteObject{14986}{<< /D [\MusicPDFPageRef{891} /XYZ 193.779 334.204 null] >>}
\MusicPDFWriteObject{14987}{<< /D [\MusicPDFPageRef{891} /XYZ 283.586 297.777 null] >>}
\MusicPDFWriteObject{14988}{<< /D [\MusicPDFPageRef{891} /XYZ 256.255 135.661 null] >>}
\MusicPDFWriteObject{14989}{<< /D [\MusicPDFPageRef{891} /XYZ 210.261 99.233 null] >>}
\MusicPDFWriteObject{14990}{<< /D [\MusicPDFPageRef{892} /XYZ 232.583 625.07 null] >>}
\MusicPDFWriteObject{14991}{<< /D [\MusicPDFPageRef{892} /XYZ 282.395 588.642 null] >>}
\MusicPDFWriteObject{14992}{<< /D [\MusicPDFPageRef{892} /XYZ 248.01 446.44 null] >>}
\MusicPDFWriteObject{14993}{<< /D [\MusicPDFPageRef{892} /XYZ 249.238 396.861 null] >>}
\MusicPDFWriteObject{14994}{<< /D [\MusicPDFPageRef{892} /XYZ 242.901 334.132 null] >>}
\MusicPDFWriteObject{14995}{<< /D [\MusicPDFPageRef{892} /XYZ 217.539 207.118 null] >>}
\MusicPDFWriteObject{14996}{<< /D [\MusicPDFPageRef{890} /XYZ 273.097 607.177 null] >>}
\MusicPDFWriteObject{14997}{<< /D [\MusicPDFPageRef{892} /XYZ 208.708 99.233 null] >>}
\MusicPDFWriteObject{14998}{<< /D [\MusicPDFPageRef{893} /XYZ 229.774 315.922 null] >>}
\MusicPDFWriteObject{14999}{<< /D [\MusicPDFPageRef{893} /XYZ 167.673 99.233 null] >>}
\MusicPDFWriteObject{15000}{<< /D [\MusicPDFPageRef{895} /XYZ 282.646 669.906 null] >>}
\MusicPDFWriteObject{15001}{<< /D [\MusicPDFPageRef{895} /XYZ 258.891 627.102 null] >>}
\MusicPDFWriteObject{15002}{<< /D [\MusicPDFPageRef{895} /XYZ 210.077 458.63 null] >>}
\MusicPDFWriteObject{15003}{<< /D [\MusicPDFPageRef{895} /XYZ 163.932 395.901 null] >>}
\MusicPDFWriteObject{15004}{<< /D [\MusicPDFPageRef{895} /XYZ 200.458 258.094 null] >>}
\MusicPDFWriteObject{15005}{<< /D [\MusicPDFPageRef{895} /XYZ 242.8 221.666 null] >>}
\MusicPDFWriteObject{15006}{<< /D [\MusicPDFPageRef{895} /XYZ 193.676 99.233 null] >>}
\MusicPDFWriteObject{15007}{<< /D [\MusicPDFPageRef{890} /XYZ 264.592 446.355 null] >>}
\MusicPDFWriteObject{15008}{<< /D [\MusicPDFPageRef{896} /XYZ 229.32 683.057 null] >>}
\MusicPDFWriteObject{15009}{<< /D [\MusicPDFPageRef{896} /XYZ 268.077 401.752 null] >>}
\MusicPDFWriteObject{15010}{<< /D [\MusicPDFPageRef{896} /XYZ 166.541 378.475 null] >>}
\MusicPDFWriteObject{15011}{<< /D [\MusicPDFPageRef{896} /XYZ 189.295 355.198 null] >>}
\MusicPDFWriteObject{15012}{<< /D [\MusicPDFPageRef{896} /XYZ 242.272 331.921 null] >>}
\MusicPDFWriteObject{15013}{<< /D [\MusicPDFPageRef{896} /XYZ 213.302 308.644 null] >>}
\MusicPDFWriteObject{15014}{<< /D [\MusicPDFPageRef{896} /XYZ 218.531 285.368 null] >>}
\MusicPDFWriteObject{15015}{<< /D [\MusicPDFPageRef{890} /XYZ 274.373 306.656 null] >>}
\MusicPDFWriteObject{15016}{<< /D [\MusicPDFPageRef{890} /XYZ 257.851 203.443 null] >>}
\MusicPDFWriteObject{15017}{<< /D [\MusicPDFPageRef{891} /XYZ 256.752 631.851 null] >>}
\MusicPDFWriteObject{15018}{<< /D [\MusicPDFPageRef{891} /XYZ 209.68 595.423 null] >>}
\MusicPDFWriteObject{15019}{<< /D [\MusicPDFPageRef{891} /XYZ 256.764 528.56 null] >>}
\MusicPDFWriteObject{15020}{<< /D [\MusicPDFPageRef{891} /XYZ 225.467 446.742 null] >>}
\MusicPDFWriteObject{15021}{<< /D [\MusicPDFPageRef{899} /XYZ 88.512 463.223 null] >>}
\MusicPDFWriteObject{15022}{<< /D [\MusicPDFPageRef{899} /XYZ 129.184 349.667 null] >>}
\MusicPDFWriteObject{15023}{<< /D [\MusicPDFPageRef{901} /XYZ 277.735 575.723 null] >>}
\MusicPDFWriteObject{15024}{<< /D [\MusicPDFPageRef{902} /XYZ 256.689 455.355 null] >>}
\MusicPDFWriteObject{15025}{<< /D [\MusicPDFPageRef{902} /XYZ 220.238 392.626 null] >>}
\MusicPDFWriteObject{15026}{<< /D [\MusicPDFPageRef{902} /XYZ 149.832 250.993 null] >>}
\MusicPDFWriteObject{15027}{<< /D [\MusicPDFPageRef{902} /XYZ 264.152 201.414 null] >>}
\MusicPDFWriteObject{15028}{<< /D [\MusicPDFPageRef{903} /XYZ 153.139 273.252 null] >>}
\MusicPDFWriteObject{15029}{<< /D [\MusicPDFPageRef{903} /XYZ 227.704 162.281 null] >>}
\MusicPDFWriteObject{15030}{<< /D [\MusicPDFPageRef{904} /XYZ 192.291 674.434 null] >>}
\MusicPDFWriteObject{15031}{<< /D [\MusicPDFPageRef{904} /XYZ 191.901 582.6 null] >>}
\MusicPDFWriteObject{15032}{<< /D [\MusicPDFPageRef{904} /XYZ 154.687 463.359 null] >>}
\MusicPDFWriteObject{15033}{<< /D [\MusicPDFPageRef{907} /XYZ 254.656 209.748 null] >>}
\MusicPDFWriteObject{15034}{<< /D [\MusicPDFPageRef{908} /XYZ 272.343 683.057 null] >>}
\MusicPDFWriteObject{15035}{<< /D [\MusicPDFPageRef{909} /XYZ 248.718 343.852 null] >>}
\MusicPDFWriteObject{15036}{<< /D [\MusicPDFPageRef{910} /XYZ 278.05 450.848 null] >>}
\MusicPDFWriteObject{15037}{<< /D [\MusicPDFPageRef{912} /XYZ 180.289 177.605 null] >>}
\MusicPDFWriteObject{15038}{<< /D [\MusicPDFPageRef{910} /XYZ 233.241 99.233 null] >>}
\MusicPDFWriteObject{15039}{<< /D [\MusicPDFPageRef{911} /XYZ 206.153 547.351 null] >>}
\MusicPDFWriteObject{15040}{<< /D [\MusicPDFPageRef{911} /XYZ 260.319 418.013 null] >>}
\MusicPDFWriteObject{15041}{<< /D [\MusicPDFPageRef{911} /XYZ 171.4 334.21 null] >>}
\MusicPDFWriteObject{15042}{<< /D [\MusicPDFPageRef{912} /XYZ 177.428 524.137 null] >>}
\MusicPDFWriteObject{15043}{<< /D [\MusicPDFPageRef{912} /XYZ 173.197 474.691 null] >>}
\MusicPDFWriteObject{15044}{<< /D [\MusicPDFPageRef{912} /XYZ 267.572 425.574 null] >>}
\MusicPDFWriteObject{15045}{<< /D [\MusicPDFPageRef{912} /XYZ 227.433 376.128 null] >>}
\MusicPDFWriteObject{15046}{<< /D [\MusicPDFPageRef{915} /XYZ 284.475 476.214 null] >>}
\MusicPDFWriteObject{15047}{<< /D [\MusicPDFPageRef{918} /XYZ 245.732 367.781 null] >>}
\MusicPDFWriteObject{15048}{<< /D [\MusicPDFPageRef{919} /XYZ 203.724 547.471 null] >>}
\MusicPDFWriteObject{15049}{<< /D [\MusicPDFPageRef{919} /XYZ 170.374 510.246 null] >>}
\MusicPDFWriteObject{15050}{<< /D [\MusicPDFPageRef{919} /XYZ 228.525 450.222 null] >>}
\MusicPDFWriteObject{15051}{<< /D [\MusicPDFPageRef{919} /XYZ 279.934 413.795 null] >>}
\MusicPDFWriteObject{15052}{<< /D [\MusicPDFPageRef{919} /XYZ 249.743 247.524 null] >>}
\MusicPDFWriteObject{15053}{<< /D [\MusicPDFPageRef{920} /XYZ 243.954 629.823 null] >>}
\MusicPDFWriteObject{15054}{<< /D [\MusicPDFPageRef{920} /XYZ 132.27 124.797 null] >>}
\MusicPDFWriteObject{15055}{<< /D [\MusicPDFPageRef{921} /XYZ 154.02 683.057 null] >>}
\MusicPDFWriteObject{15056}{<< /D [\MusicPDFPageRef{921} /XYZ 188.755 646.629 null] >>}
\MusicPDFWriteObject{15057}{<< /D [\MusicPDFPageRef{921} /XYZ 149.438 610.201 null] >>}
\MusicPDFWriteObject{15058}{<< /D [\MusicPDFPageRef{922} /XYZ 180.127 606.659 null] >>}
\MusicPDFWriteObject{15059}{<< /D [\MusicPDFPageRef{922} /XYZ 229.938 229.754 null] >>}
\MusicPDFWriteObject{15060}{<< /D [\MusicPDFPageRef{923} /XYZ 215.34 99.233 null] >>}
\MusicPDFWriteObject{15061}{<< /D [\MusicPDFPageRef{924} /XYZ 289.047 591.001 null] >>}
\MusicPDFWriteObject{15062}{<< /D [\MusicPDFPageRef{924} /XYZ 120.93 554.574 null] >>}
\MusicPDFWriteObject{15063}{<< /D [\MusicPDFPageRef{924} /XYZ 275.404 457.442 null] >>}
\MusicPDFWriteObject{15064}{<< /D [\MusicPDFPageRef{922} /XYZ 160.419 193.327 null] >>}
\MusicPDFWriteObject{15065}{<< /D [\MusicPDFPageRef{922} /XYZ 239.612 130.597 null] >>}
\MusicPDFWriteObject{15066}{<< /D [\MusicPDFPageRef{923} /XYZ 242.144 614.802 null] >>}
\MusicPDFWriteObject{15067}{<< /D [\MusicPDFPageRef{923} /XYZ 242.919 565.224 null] >>}
\MusicPDFWriteObject{15068}{<< /D [\MusicPDFPageRef{923} /XYZ 242.062 489.344 null] >>}
\MusicPDFWriteObject{15069}{<< /D [\MusicPDFPageRef{923} /XYZ 256.428 413.464 null] >>}
\MusicPDFWriteObject{15070}{<< /D [\MusicPDFPageRef{923} /XYZ 258.397 337.584 null] >>}
\MusicPDFWriteObject{15071}{<< /D [\MusicPDFPageRef{923} /XYZ 231.327 175.113 null] >>}
\MusicPDFWriteObject{15072}{<< /D [\MusicPDFPageRef{928} /XYZ 250.708 555.093 null] >>}
\MusicPDFWriteObject{15073}{<< /D [\MusicPDFPageRef{38} /XYZ 112.621 640.257 null] >>}
\MusicPDFWriteObject{15074}{<< /D [\MusicPDFPageRef{596} /XYZ 181.345 386.375 null] >>}
\MusicPDFWriteObject{15075}{<< /D [\MusicPDFPageRef{602} /XYZ 162.749 229.817 null] >>}
\MusicPDFWriteObject{15076}{<< /D [\MusicPDFPageRef{605} /XYZ 152.945 233.065 null] >>}
\MusicPDFWriteObject{15077}{<< /D [\MusicPDFPageRef{646} /XYZ 154.849 281.14 null] >>}
\MusicPDFWriteObject{15078}{<< /D [\MusicPDFPageRef{798} /XYZ 189.449 347.292 null] >>}
\MusicPDFWriteObject{15079}{<< /D [\MusicPDFPageRef{9} /XYZ 312.781 571.169 null] >>}
\MusicPDFWriteObject{15080}{<< /D [\MusicPDFPageRef{9} /XYZ 72 569.586 null] >>}
\MusicPDFWriteObject{15081}{<< /D [\MusicPDFPageRef{9} /XYZ 312.781 543.097 null] >>}
\MusicPDFWriteObject{15082}{<< /D [\MusicPDFPageRef{9} /XYZ 71.641 516.795 null] >>}
\MusicPDFWriteObject{15083}{<< /D [\MusicPDFPageRef{9} /XYZ 312.781 556.247 null] >>}
\MusicPDFWriteObject{15084}{<< /D [\MusicPDFPageRef{9} /XYZ 71.641 529.946 null] >>}
\MusicPDFWriteObject{15085}{<< /D [\MusicPDFPageRef{9} /XYZ 72 503.644 null] >>}
\MusicPDFWriteObject{15086}{<< /D [\MusicPDFPageRef{440} /XYZ 72 579.366 null] >>}
\MusicPDFWriteObject{15087}{<< /D [\MusicPDFPageRef{445} /XYZ 72 239.402 null] >>}
\MusicPDFWriteObject{15088}{<< /D [\MusicPDFPageRef{99} /XYZ 72 251.077 null] >>}
\MusicPDFWriteObject{15089}{<< /D [\MusicPDFPageRef{142} /XYZ 72 344.919 null] >>}
\MusicPDFWriteObject{15090}{<< /D [\MusicPDFPageRef{204} /XYZ 72 506.522 null] >>}
\MusicPDFWriteObject{15091}{<< /D [\MusicPDFPageRef{209} /XYZ 72 292.958 null] >>}
\MusicPDFWriteObject{15092}{<< /D [\MusicPDFPageRef{364} /XYZ 72 582.992 null] >>}
\MusicPDFWriteObject{15093}{<< /D [\MusicPDFPageRef{824} /XYZ 72 591.782 null] >>}
\MusicPDFWriteObject{15094}{<< /D [\MusicPDFPageRef{854} /XYZ 72 405.361 null] >>}
\MusicPDFWriteObject{15095}{<< /D [\MusicPDFPageRef{868} /XYZ 72 220.314 null] >>}
\MusicPDFWriteObject{15096}{<< /D [\MusicPDFPageRef{885} /XYZ 72 666.676 null] >>}
\MusicPDFWriteObject{15097}{<< /D [\MusicPDFPageRef{890} /XYZ 72 568.721 null] >>}
\MusicPDFWriteObject{15098}{<< /D [\MusicPDFPageRef{905} /XYZ 72 361.485 null] >>}
\MusicPDFWriteObject{15099}{<< /D [\MusicPDFPageRef{910} /XYZ 72 328.232 null] >>}
\MusicPDFWriteObject{15100}{<< /D [\MusicPDFPageRef{911} /XYZ 72 566.678 null] >>}
\MusicPDFWriteObject{15101}{<< /D [\MusicPDFPageRef{914} /XYZ 72 670.757 null] >>}
\MusicPDFWriteObject{15102}{<< /D [\MusicPDFPageRef{928} /XYZ 72 701.27 null] >>}
\MusicPDFWriteObject{15103}{<< /D [\MusicPDFPageRef{1} /XYZ 71 721 null] >>}
\MusicPDFWriteObject{15104}{<< /D [\MusicPDFPageRef{10} /XYZ 71 721 null] >>}
\MusicPDFWriteObject{15105}{<< /D [\MusicPDFPageRef{100} /XYZ 71 721 null] >>}
\MusicPDFWriteObject{15106}{<< /D [\MusicPDFPageRef{101} /XYZ 71 721 null] >>}
\MusicPDFWriteObject{15107}{<< /D [\MusicPDFPageRef{102} /XYZ 71 721 null] >>}
\MusicPDFWriteObject{15108}{<< /D [\MusicPDFPageRef{103} /XYZ 71 721 null] >>}
\MusicPDFWriteObject{15109}{<< /D [\MusicPDFPageRef{104} /XYZ 71 721 null] >>}
\MusicPDFWriteObject{15110}{<< /D [\MusicPDFPageRef{105} /XYZ 71 721 null] >>}
\MusicPDFWriteObject{15111}{<< /D [\MusicPDFPageRef{106} /XYZ 71 721 null] >>}
\MusicPDFWriteObject{15112}{<< /D [\MusicPDFPageRef{107} /XYZ 71 721 null] >>}
\MusicPDFWriteObject{15113}{<< /D [\MusicPDFPageRef{108} /XYZ 71 721 null] >>}
\MusicPDFWriteObject{15114}{<< /D [\MusicPDFPageRef{109} /XYZ 71 721 null] >>}
\MusicPDFWriteObject{15115}{<< /D [\MusicPDFPageRef{11} /XYZ 71 721 null] >>}
\MusicPDFWriteObject{15116}{<< /D [\MusicPDFPageRef{110} /XYZ 71 721 null] >>}
\MusicPDFWriteObject{15117}{<< /D [\MusicPDFPageRef{111} /XYZ 71 721 null] >>}
\MusicPDFWriteObject{15118}{<< /D [\MusicPDFPageRef{112} /XYZ 71 721 null] >>}
\MusicPDFWriteObject{15119}{<< /D [\MusicPDFPageRef{113} /XYZ 71 721 null] >>}
\MusicPDFWriteObject{15120}{<< /D [\MusicPDFPageRef{114} /XYZ 71 721 null] >>}
\MusicPDFWriteObject{15121}{<< /D [\MusicPDFPageRef{115} /XYZ 71 721 null] >>}
\MusicPDFWriteObject{15122}{<< /D [\MusicPDFPageRef{116} /XYZ 71 721 null] >>}
\MusicPDFWriteObject{15123}{<< /D [\MusicPDFPageRef{117} /XYZ 71 721 null] >>}
\MusicPDFWriteObject{15124}{<< /D [\MusicPDFPageRef{118} /XYZ 71 721 null] >>}
\MusicPDFWriteObject{15125}{<< /D [\MusicPDFPageRef{119} /XYZ 71 721 null] >>}
\MusicPDFWriteObject{15126}{<< /D [\MusicPDFPageRef{12} /XYZ 71 721 null] >>}
\MusicPDFWriteObject{15127}{<< /D [\MusicPDFPageRef{120} /XYZ 71 721 null] >>}
\MusicPDFWriteObject{15128}{<< /D [\MusicPDFPageRef{121} /XYZ 71 721 null] >>}
\MusicPDFWriteObject{15129}{<< /D [\MusicPDFPageRef{122} /XYZ 71 721 null] >>}
\MusicPDFWriteObject{15130}{<< /D [\MusicPDFPageRef{123} /XYZ 71 721 null] >>}
\MusicPDFWriteObject{15131}{<< /D [\MusicPDFPageRef{124} /XYZ 71 721 null] >>}
\MusicPDFWriteObject{15132}{<< /D [\MusicPDFPageRef{125} /XYZ 71 721 null] >>}
\MusicPDFWriteObject{15133}{<< /D [\MusicPDFPageRef{126} /XYZ 71 721 null] >>}
\MusicPDFWriteObject{15134}{<< /D [\MusicPDFPageRef{127} /XYZ 71 721 null] >>}
\MusicPDFWriteObject{15135}{<< /D [\MusicPDFPageRef{128} /XYZ 71 721 null] >>}
\MusicPDFWriteObject{15136}{<< /D [\MusicPDFPageRef{129} /XYZ 71 721 null] >>}
\MusicPDFWriteObject{15137}{<< /D [\MusicPDFPageRef{13} /XYZ 71 721 null] >>}
\MusicPDFWriteObject{15138}{<< /D [\MusicPDFPageRef{130} /XYZ 71 721 null] >>}
\MusicPDFWriteObject{15139}{<< /D [\MusicPDFPageRef{131} /XYZ 71 721 null] >>}
\MusicPDFWriteObject{15140}{<< /D [\MusicPDFPageRef{132} /XYZ 71 721 null] >>}
\MusicPDFWriteObject{15141}{<< /D [\MusicPDFPageRef{133} /XYZ 71 721 null] >>}
\MusicPDFWriteObject{15142}{<< /D [\MusicPDFPageRef{134} /XYZ 71 721 null] >>}
\MusicPDFWriteObject{15143}{<< /D [\MusicPDFPageRef{135} /XYZ 71 721 null] >>}
\MusicPDFWriteObject{15144}{<< /D [\MusicPDFPageRef{136} /XYZ 71 721 null] >>}
\MusicPDFWriteObject{15145}{<< /D [\MusicPDFPageRef{137} /XYZ 71 721 null] >>}
\MusicPDFWriteObject{15146}{<< /D [\MusicPDFPageRef{138} /XYZ 71 721 null] >>}
\MusicPDFWriteObject{15147}{<< /D [\MusicPDFPageRef{139} /XYZ 71 721 null] >>}
\MusicPDFWriteObject{15148}{<< /D [\MusicPDFPageRef{14} /XYZ 71 721 null] >>}
\MusicPDFWriteObject{15149}{<< /D [\MusicPDFPageRef{140} /XYZ 71 721 null] >>}
\MusicPDFWriteObject{15150}{<< /D [\MusicPDFPageRef{141} /XYZ 71 721 null] >>}
\MusicPDFWriteObject{15151}{<< /D [\MusicPDFPageRef{142} /XYZ 71 721 null] >>}
\MusicPDFWriteObject{15152}{<< /D [\MusicPDFPageRef{143} /XYZ 71 721 null] >>}
\MusicPDFWriteObject{15153}{<< /D [\MusicPDFPageRef{144} /XYZ 71 721 null] >>}
\MusicPDFWriteObject{15154}{<< /D [\MusicPDFPageRef{145} /XYZ 71 721 null] >>}
\MusicPDFWriteObject{15155}{<< /D [\MusicPDFPageRef{146} /XYZ 71 721 null] >>}
\MusicPDFWriteObject{15156}{<< /D [\MusicPDFPageRef{147} /XYZ 71 721 null] >>}
\MusicPDFWriteObject{15157}{<< /D [\MusicPDFPageRef{148} /XYZ 71 721 null] >>}
\MusicPDFWriteObject{15158}{<< /D [\MusicPDFPageRef{149} /XYZ 71 721 null] >>}
\MusicPDFWriteObject{15159}{<< /D [\MusicPDFPageRef{15} /XYZ 71 721 null] >>}
\MusicPDFWriteObject{15160}{<< /D [\MusicPDFPageRef{150} /XYZ 71 721 null] >>}
\MusicPDFWriteObject{15161}{<< /D [\MusicPDFPageRef{151} /XYZ 71 721 null] >>}
\MusicPDFWriteObject{15162}{<< /D [\MusicPDFPageRef{152} /XYZ 71 721 null] >>}
\MusicPDFWriteObject{15163}{<< /D [\MusicPDFPageRef{153} /XYZ 71 721 null] >>}
\MusicPDFWriteObject{15164}{<< /D [\MusicPDFPageRef{154} /XYZ 71 721 null] >>}
\MusicPDFWriteObject{15165}{<< /D [\MusicPDFPageRef{155} /XYZ 71 721 null] >>}
\MusicPDFWriteObject{15166}{<< /D [\MusicPDFPageRef{156} /XYZ 71 721 null] >>}
\MusicPDFWriteObject{15167}{<< /D [\MusicPDFPageRef{157} /XYZ 71 721 null] >>}
\MusicPDFWriteObject{15168}{<< /D [\MusicPDFPageRef{158} /XYZ 71 721 null] >>}
\MusicPDFWriteObject{15169}{<< /D [\MusicPDFPageRef{159} /XYZ 71 721 null] >>}
\MusicPDFWriteObject{15170}{<< /D [\MusicPDFPageRef{16} /XYZ 71 721 null] >>}
\MusicPDFWriteObject{15171}{<< /D [\MusicPDFPageRef{160} /XYZ 71 721 null] >>}
\MusicPDFWriteObject{15172}{<< /D [\MusicPDFPageRef{161} /XYZ 71 721 null] >>}
\MusicPDFWriteObject{15173}{<< /D [\MusicPDFPageRef{162} /XYZ 71 721 null] >>}
\MusicPDFWriteObject{15174}{<< /D [\MusicPDFPageRef{163} /XYZ 71 721 null] >>}
\MusicPDFWriteObject{15175}{<< /D [\MusicPDFPageRef{164} /XYZ 71 721 null] >>}
\MusicPDFWriteObject{15176}{<< /D [\MusicPDFPageRef{165} /XYZ 71 721 null] >>}
\MusicPDFWriteObject{15177}{<< /D [\MusicPDFPageRef{166} /XYZ 71 721 null] >>}
\MusicPDFWriteObject{15178}{<< /D [\MusicPDFPageRef{167} /XYZ 71 721 null] >>}
\MusicPDFWriteObject{15179}{<< /D [\MusicPDFPageRef{168} /XYZ 71 721 null] >>}
\MusicPDFWriteObject{15180}{<< /D [\MusicPDFPageRef{169} /XYZ 71 721 null] >>}
\MusicPDFWriteObject{15181}{<< /D [\MusicPDFPageRef{17} /XYZ 71 721 null] >>}
\MusicPDFWriteObject{15182}{<< /D [\MusicPDFPageRef{170} /XYZ 71 721 null] >>}
\MusicPDFWriteObject{15183}{<< /D [\MusicPDFPageRef{171} /XYZ 71 721 null] >>}
\MusicPDFWriteObject{15184}{<< /D [\MusicPDFPageRef{172} /XYZ 71 721 null] >>}
\MusicPDFWriteObject{15185}{<< /D [\MusicPDFPageRef{173} /XYZ 71 721 null] >>}
\MusicPDFWriteObject{15186}{<< /D [\MusicPDFPageRef{174} /XYZ 71 721 null] >>}
\MusicPDFWriteObject{15187}{<< /D [\MusicPDFPageRef{175} /XYZ 71 721 null] >>}
\MusicPDFWriteObject{15188}{<< /D [\MusicPDFPageRef{176} /XYZ 71 721 null] >>}
\MusicPDFWriteObject{15189}{<< /D [\MusicPDFPageRef{177} /XYZ 71 721 null] >>}
\MusicPDFWriteObject{15190}{<< /D [\MusicPDFPageRef{178} /XYZ 71 721 null] >>}
\MusicPDFWriteObject{15191}{<< /D [\MusicPDFPageRef{179} /XYZ 71 721 null] >>}
\MusicPDFWriteObject{15192}{<< /D [\MusicPDFPageRef{18} /XYZ 71 721 null] >>}
\MusicPDFWriteObject{15193}{<< /D [\MusicPDFPageRef{180} /XYZ 71 721 null] >>}
\MusicPDFWriteObject{15194}{<< /D [\MusicPDFPageRef{181} /XYZ 71 721 null] >>}
\MusicPDFWriteObject{15195}{<< /D [\MusicPDFPageRef{182} /XYZ 71 721 null] >>}
\MusicPDFWriteObject{15196}{<< /D [\MusicPDFPageRef{183} /XYZ 71 721 null] >>}
\MusicPDFWriteObject{15197}{<< /D [\MusicPDFPageRef{184} /XYZ 71 721 null] >>}
\MusicPDFWriteObject{15198}{<< /D [\MusicPDFPageRef{185} /XYZ 71 721 null] >>}
\MusicPDFWriteObject{15199}{<< /D [\MusicPDFPageRef{186} /XYZ 71 721 null] >>}
\MusicPDFWriteObject{15200}{<< /D [\MusicPDFPageRef{187} /XYZ 71 721 null] >>}
\MusicPDFWriteObject{15201}{<< /D [\MusicPDFPageRef{188} /XYZ 71 721 null] >>}
\MusicPDFWriteObject{15202}{<< /D [\MusicPDFPageRef{189} /XYZ 71 721 null] >>}
\MusicPDFWriteObject{15203}{<< /D [\MusicPDFPageRef{19} /XYZ 71 721 null] >>}
\MusicPDFWriteObject{15204}{<< /D [\MusicPDFPageRef{190} /XYZ 71 721 null] >>}
\MusicPDFWriteObject{15205}{<< /D [\MusicPDFPageRef{191} /XYZ 71 721 null] >>}
\MusicPDFWriteObject{15206}{<< /D [\MusicPDFPageRef{192} /XYZ 71 721 null] >>}
\MusicPDFWriteObject{15207}{<< /D [\MusicPDFPageRef{193} /XYZ 71 721 null] >>}
\MusicPDFWriteObject{15208}{<< /D [\MusicPDFPageRef{194} /XYZ 71 721 null] >>}
\MusicPDFWriteObject{15209}{<< /D [\MusicPDFPageRef{195} /XYZ 71 721 null] >>}
\MusicPDFWriteObject{15210}{<< /D [\MusicPDFPageRef{196} /XYZ 71 721 null] >>}
\MusicPDFWriteObject{15211}{<< /D [\MusicPDFPageRef{197} /XYZ 71 721 null] >>}
\MusicPDFWriteObject{15212}{<< /D [\MusicPDFPageRef{198} /XYZ 71 721 null] >>}
\MusicPDFWriteObject{15213}{<< /D [\MusicPDFPageRef{199} /XYZ 71 721 null] >>}
\MusicPDFWriteObject{15214}{<< /D [\MusicPDFPageRef{2} /XYZ 71 721 null] >>}
\MusicPDFWriteObject{15215}{<< /D [\MusicPDFPageRef{20} /XYZ 71 721 null] >>}
\MusicPDFWriteObject{15216}{<< /D [\MusicPDFPageRef{200} /XYZ 71 721 null] >>}
\MusicPDFWriteObject{15217}{<< /D [\MusicPDFPageRef{201} /XYZ 71 721 null] >>}
\MusicPDFWriteObject{15218}{<< /D [\MusicPDFPageRef{202} /XYZ 71 721 null] >>}
\MusicPDFWriteObject{15219}{<< /D [\MusicPDFPageRef{203} /XYZ 71 721 null] >>}
\MusicPDFWriteObject{15220}{<< /D [\MusicPDFPageRef{204} /XYZ 71 721 null] >>}
\MusicPDFWriteObject{15221}{<< /D [\MusicPDFPageRef{205} /XYZ 71 721 null] >>}
\MusicPDFWriteObject{15222}{<< /D [\MusicPDFPageRef{206} /XYZ 71 721 null] >>}
\MusicPDFWriteObject{15223}{<< /D [\MusicPDFPageRef{207} /XYZ 71 721 null] >>}
\MusicPDFWriteObject{15224}{<< /D [\MusicPDFPageRef{208} /XYZ 71 721 null] >>}
\MusicPDFWriteObject{15225}{<< /D [\MusicPDFPageRef{209} /XYZ 71 721 null] >>}
\MusicPDFWriteObject{15226}{<< /D [\MusicPDFPageRef{21} /XYZ 71 721 null] >>}
\MusicPDFWriteObject{15227}{<< /D [\MusicPDFPageRef{210} /XYZ 71 721 null] >>}
\MusicPDFWriteObject{15228}{<< /D [\MusicPDFPageRef{211} /XYZ 71 721 null] >>}
\MusicPDFWriteObject{15229}{<< /D [\MusicPDFPageRef{212} /XYZ 71 721 null] >>}
\MusicPDFWriteObject{15230}{<< /D [\MusicPDFPageRef{213} /XYZ 71 721 null] >>}
\MusicPDFWriteObject{15231}{<< /D [\MusicPDFPageRef{214} /XYZ 71 721 null] >>}
\MusicPDFWriteObject{15232}{<< /D [\MusicPDFPageRef{215} /XYZ 71 721 null] >>}
\MusicPDFWriteObject{15233}{<< /D [\MusicPDFPageRef{216} /XYZ 71 721 null] >>}
\MusicPDFWriteObject{15234}{<< /D [\MusicPDFPageRef{217} /XYZ 71 721 null] >>}
\MusicPDFWriteObject{15235}{<< /D [\MusicPDFPageRef{218} /XYZ 71 721 null] >>}
\MusicPDFWriteObject{15236}{<< /D [\MusicPDFPageRef{219} /XYZ 71 721 null] >>}
\MusicPDFWriteObject{15237}{<< /D [\MusicPDFPageRef{22} /XYZ 71 721 null] >>}
\MusicPDFWriteObject{15238}{<< /D [\MusicPDFPageRef{220} /XYZ 71 721 null] >>}
\MusicPDFWriteObject{15239}{<< /D [\MusicPDFPageRef{221} /XYZ 71 721 null] >>}
\MusicPDFWriteObject{15240}{<< /D [\MusicPDFPageRef{222} /XYZ 71 721 null] >>}
\MusicPDFWriteObject{15241}{<< /D [\MusicPDFPageRef{223} /XYZ 71 721 null] >>}
\MusicPDFWriteObject{15242}{<< /D [\MusicPDFPageRef{224} /XYZ 71 721 null] >>}
\MusicPDFWriteObject{15243}{<< /D [\MusicPDFPageRef{225} /XYZ 71 721 null] >>}
\MusicPDFWriteObject{15244}{<< /D [\MusicPDFPageRef{226} /XYZ 71 721 null] >>}
\MusicPDFWriteObject{15245}{<< /D [\MusicPDFPageRef{227} /XYZ 71 721 null] >>}
\MusicPDFWriteObject{15246}{<< /D [\MusicPDFPageRef{228} /XYZ 71 721 null] >>}
\MusicPDFWriteObject{15247}{<< /D [\MusicPDFPageRef{229} /XYZ 71 721 null] >>}
\MusicPDFWriteObject{15248}{<< /D [\MusicPDFPageRef{23} /XYZ 71 721 null] >>}
\MusicPDFWriteObject{15249}{<< /D [\MusicPDFPageRef{230} /XYZ 71 721 null] >>}
\MusicPDFWriteObject{15250}{<< /D [\MusicPDFPageRef{231} /XYZ 71 721 null] >>}
\MusicPDFWriteObject{15251}{<< /D [\MusicPDFPageRef{232} /XYZ 71 721 null] >>}
\MusicPDFWriteObject{15252}{<< /D [\MusicPDFPageRef{233} /XYZ 71 721 null] >>}
\MusicPDFWriteObject{15253}{<< /D [\MusicPDFPageRef{234} /XYZ 71 721 null] >>}
\MusicPDFWriteObject{15254}{<< /D [\MusicPDFPageRef{235} /XYZ 71 721 null] >>}
\MusicPDFWriteObject{15255}{<< /D [\MusicPDFPageRef{236} /XYZ 71 721 null] >>}
\MusicPDFWriteObject{15256}{<< /D [\MusicPDFPageRef{237} /XYZ 71 721 null] >>}
\MusicPDFWriteObject{15257}{<< /D [\MusicPDFPageRef{238} /XYZ 71 721 null] >>}
\MusicPDFWriteObject{15258}{<< /D [\MusicPDFPageRef{239} /XYZ 71 721 null] >>}
\MusicPDFWriteObject{15259}{<< /D [\MusicPDFPageRef{24} /XYZ 71 721 null] >>}
\MusicPDFWriteObject{15260}{<< /D [\MusicPDFPageRef{240} /XYZ 71 721 null] >>}
\MusicPDFWriteObject{15261}{<< /D [\MusicPDFPageRef{241} /XYZ 71 721 null] >>}
\MusicPDFWriteObject{15262}{<< /D [\MusicPDFPageRef{242} /XYZ 71 721 null] >>}
\MusicPDFWriteObject{15263}{<< /D [\MusicPDFPageRef{243} /XYZ 71 721 null] >>}
\MusicPDFWriteObject{15264}{<< /D [\MusicPDFPageRef{244} /XYZ 71 721 null] >>}
\MusicPDFWriteObject{15265}{<< /D [\MusicPDFPageRef{245} /XYZ 71 721 null] >>}
\MusicPDFWriteObject{15266}{<< /D [\MusicPDFPageRef{246} /XYZ 71 721 null] >>}
\MusicPDFWriteObject{15267}{<< /D [\MusicPDFPageRef{247} /XYZ 71 721 null] >>}
\MusicPDFWriteObject{15268}{<< /D [\MusicPDFPageRef{248} /XYZ 71 721 null] >>}
\MusicPDFWriteObject{15269}{<< /D [\MusicPDFPageRef{249} /XYZ 71 721 null] >>}
\MusicPDFWriteObject{15270}{<< /D [\MusicPDFPageRef{25} /XYZ 71 721 null] >>}
\MusicPDFWriteObject{15271}{<< /D [\MusicPDFPageRef{250} /XYZ 71 721 null] >>}
\MusicPDFWriteObject{15272}{<< /D [\MusicPDFPageRef{251} /XYZ 71 721 null] >>}
\MusicPDFWriteObject{15273}{<< /D [\MusicPDFPageRef{252} /XYZ 71 721 null] >>}
\MusicPDFWriteObject{15274}{<< /D [\MusicPDFPageRef{253} /XYZ 71 721 null] >>}
\MusicPDFWriteObject{15275}{<< /D [\MusicPDFPageRef{254} /XYZ 71 721 null] >>}
\MusicPDFWriteObject{15276}{<< /D [\MusicPDFPageRef{255} /XYZ 71 721 null] >>}
\MusicPDFWriteObject{15277}{<< /D [\MusicPDFPageRef{256} /XYZ 71 721 null] >>}
\MusicPDFWriteObject{15278}{<< /D [\MusicPDFPageRef{257} /XYZ 71 721 null] >>}
\MusicPDFWriteObject{15279}{<< /D [\MusicPDFPageRef{258} /XYZ 71 721 null] >>}
\MusicPDFWriteObject{15280}{<< /D [\MusicPDFPageRef{259} /XYZ 71 721 null] >>}
\MusicPDFWriteObject{15281}{<< /D [\MusicPDFPageRef{26} /XYZ 71 721 null] >>}
\MusicPDFWriteObject{15282}{<< /D [\MusicPDFPageRef{260} /XYZ 71 721 null] >>}
\MusicPDFWriteObject{15283}{<< /D [\MusicPDFPageRef{261} /XYZ 71 721 null] >>}
\MusicPDFWriteObject{15284}{<< /D [\MusicPDFPageRef{262} /XYZ 71 721 null] >>}
\MusicPDFWriteObject{15285}{<< /D [\MusicPDFPageRef{263} /XYZ 71 721 null] >>}
\MusicPDFWriteObject{15286}{<< /D [\MusicPDFPageRef{264} /XYZ 71 721 null] >>}
\MusicPDFWriteObject{15287}{<< /D [\MusicPDFPageRef{265} /XYZ 71 721 null] >>}
\MusicPDFWriteObject{15288}{<< /D [\MusicPDFPageRef{266} /XYZ 71 721 null] >>}
\MusicPDFWriteObject{15289}{<< /D [\MusicPDFPageRef{267} /XYZ 71 721 null] >>}
\MusicPDFWriteObject{15290}{<< /D [\MusicPDFPageRef{268} /XYZ 71 721 null] >>}
\MusicPDFWriteObject{15291}{<< /D [\MusicPDFPageRef{269} /XYZ 71 721 null] >>}
\MusicPDFWriteObject{15292}{<< /D [\MusicPDFPageRef{27} /XYZ 71 721 null] >>}
\MusicPDFWriteObject{15293}{<< /D [\MusicPDFPageRef{270} /XYZ 71 721 null] >>}
\MusicPDFWriteObject{15294}{<< /D [\MusicPDFPageRef{271} /XYZ 71 721 null] >>}
\MusicPDFWriteObject{15295}{<< /D [\MusicPDFPageRef{272} /XYZ 71 721 null] >>}
\MusicPDFWriteObject{15296}{<< /D [\MusicPDFPageRef{273} /XYZ 71 721 null] >>}
\MusicPDFWriteObject{15297}{<< /D [\MusicPDFPageRef{274} /XYZ 71 721 null] >>}
\MusicPDFWriteObject{15298}{<< /D [\MusicPDFPageRef{275} /XYZ 71 721 null] >>}
\MusicPDFWriteObject{15299}{<< /D [\MusicPDFPageRef{276} /XYZ 71 721 null] >>}
\MusicPDFWriteObject{15300}{<< /D [\MusicPDFPageRef{277} /XYZ 71 721 null] >>}
\MusicPDFWriteObject{15301}{<< /D [\MusicPDFPageRef{278} /XYZ 71 721 null] >>}
\MusicPDFWriteObject{15302}{<< /D [\MusicPDFPageRef{279} /XYZ 71 721 null] >>}
\MusicPDFWriteObject{15303}{<< /D [\MusicPDFPageRef{28} /XYZ 71 721 null] >>}
\MusicPDFWriteObject{15304}{<< /D [\MusicPDFPageRef{280} /XYZ 71 721 null] >>}
\MusicPDFWriteObject{15305}{<< /D [\MusicPDFPageRef{281} /XYZ 71 721 null] >>}
\MusicPDFWriteObject{15306}{<< /D [\MusicPDFPageRef{282} /XYZ 71 721 null] >>}
\MusicPDFWriteObject{15307}{<< /D [\MusicPDFPageRef{283} /XYZ 71 721 null] >>}
\MusicPDFWriteObject{15308}{<< /D [\MusicPDFPageRef{284} /XYZ 71 721 null] >>}
\MusicPDFWriteObject{15309}{<< /D [\MusicPDFPageRef{285} /XYZ 71 721 null] >>}
\MusicPDFWriteObject{15310}{<< /D [\MusicPDFPageRef{286} /XYZ 71 721 null] >>}
\MusicPDFWriteObject{15311}{<< /D [\MusicPDFPageRef{287} /XYZ 71 721 null] >>}
\MusicPDFWriteObject{15312}{<< /D [\MusicPDFPageRef{288} /XYZ 71 721 null] >>}
\MusicPDFWriteObject{15313}{<< /D [\MusicPDFPageRef{289} /XYZ 71 721 null] >>}
\MusicPDFWriteObject{15314}{<< /D [\MusicPDFPageRef{29} /XYZ 71 721 null] >>}
\MusicPDFWriteObject{15315}{<< /D [\MusicPDFPageRef{290} /XYZ 71 721 null] >>}
\MusicPDFWriteObject{15316}{<< /D [\MusicPDFPageRef{291} /XYZ 71 721 null] >>}
\MusicPDFWriteObject{15317}{<< /D [\MusicPDFPageRef{292} /XYZ 71 721 null] >>}
\MusicPDFWriteObject{15318}{<< /D [\MusicPDFPageRef{293} /XYZ 71 721 null] >>}
\MusicPDFWriteObject{15319}{<< /D [\MusicPDFPageRef{294} /XYZ 71 721 null] >>}
\MusicPDFWriteObject{15320}{<< /D [\MusicPDFPageRef{295} /XYZ 71 721 null] >>}
\MusicPDFWriteObject{15321}{<< /D [\MusicPDFPageRef{296} /XYZ 71 721 null] >>}
\MusicPDFWriteObject{15322}{<< /D [\MusicPDFPageRef{297} /XYZ 71 721 null] >>}
\MusicPDFWriteObject{15323}{<< /D [\MusicPDFPageRef{298} /XYZ 71 721 null] >>}
\MusicPDFWriteObject{15324}{<< /D [\MusicPDFPageRef{299} /XYZ 71 721 null] >>}
\MusicPDFWriteObject{15325}{<< /D [\MusicPDFPageRef{3} /XYZ 71 721 null] >>}
\MusicPDFWriteObject{15326}{<< /D [\MusicPDFPageRef{30} /XYZ 71 721 null] >>}
\MusicPDFWriteObject{15327}{<< /D [\MusicPDFPageRef{300} /XYZ 71 721 null] >>}
\MusicPDFWriteObject{15328}{<< /D [\MusicPDFPageRef{301} /XYZ 71 721 null] >>}
\MusicPDFWriteObject{15329}{<< /D [\MusicPDFPageRef{302} /XYZ 71 721 null] >>}
\MusicPDFWriteObject{15330}{<< /D [\MusicPDFPageRef{303} /XYZ 71 721 null] >>}
\MusicPDFWriteObject{15331}{<< /D [\MusicPDFPageRef{304} /XYZ 71 721 null] >>}
\MusicPDFWriteObject{15332}{<< /D [\MusicPDFPageRef{305} /XYZ 71 721 null] >>}
\MusicPDFWriteObject{15333}{<< /D [\MusicPDFPageRef{306} /XYZ 71 721 null] >>}
\MusicPDFWriteObject{15334}{<< /D [\MusicPDFPageRef{307} /XYZ 71 721 null] >>}
\MusicPDFWriteObject{15335}{<< /D [\MusicPDFPageRef{308} /XYZ 71 721 null] >>}
\MusicPDFWriteObject{15336}{<< /D [\MusicPDFPageRef{309} /XYZ 71 721 null] >>}
\MusicPDFWriteObject{15337}{<< /D [\MusicPDFPageRef{31} /XYZ 71 721 null] >>}
\MusicPDFWriteObject{15338}{<< /D [\MusicPDFPageRef{310} /XYZ 71 721 null] >>}
\MusicPDFWriteObject{15339}{<< /D [\MusicPDFPageRef{311} /XYZ 71 721 null] >>}
\MusicPDFWriteObject{15340}{<< /D [\MusicPDFPageRef{312} /XYZ 71 721 null] >>}
\MusicPDFWriteObject{15341}{<< /D [\MusicPDFPageRef{313} /XYZ 71 721 null] >>}
\MusicPDFWriteObject{15342}{<< /D [\MusicPDFPageRef{314} /XYZ 71 721 null] >>}
\MusicPDFWriteObject{15343}{<< /D [\MusicPDFPageRef{315} /XYZ 71 721 null] >>}
\MusicPDFWriteObject{15344}{<< /D [\MusicPDFPageRef{316} /XYZ 71 721 null] >>}
\MusicPDFWriteObject{15345}{<< /D [\MusicPDFPageRef{317} /XYZ 71 721 null] >>}
\MusicPDFWriteObject{15346}{<< /D [\MusicPDFPageRef{318} /XYZ 71 721 null] >>}
\MusicPDFWriteObject{15347}{<< /D [\MusicPDFPageRef{319} /XYZ 71 721 null] >>}
\MusicPDFWriteObject{15348}{<< /D [\MusicPDFPageRef{32} /XYZ 71 721 null] >>}
\MusicPDFWriteObject{15349}{<< /D [\MusicPDFPageRef{320} /XYZ 71 721 null] >>}
\MusicPDFWriteObject{15350}{<< /D [\MusicPDFPageRef{321} /XYZ 71 721 null] >>}
\MusicPDFWriteObject{15351}{<< /D [\MusicPDFPageRef{322} /XYZ 71 721 null] >>}
\MusicPDFWriteObject{15352}{<< /D [\MusicPDFPageRef{323} /XYZ 71 721 null] >>}
\MusicPDFWriteObject{15353}{<< /D [\MusicPDFPageRef{324} /XYZ 71 721 null] >>}
\MusicPDFWriteObject{15354}{<< /D [\MusicPDFPageRef{325} /XYZ 71 721 null] >>}
\MusicPDFWriteObject{15355}{<< /D [\MusicPDFPageRef{326} /XYZ 71 721 null] >>}
\MusicPDFWriteObject{15356}{<< /D [\MusicPDFPageRef{327} /XYZ 71 721 null] >>}
\MusicPDFWriteObject{15357}{<< /D [\MusicPDFPageRef{328} /XYZ 71 721 null] >>}
\MusicPDFWriteObject{15358}{<< /D [\MusicPDFPageRef{329} /XYZ 71 721 null] >>}
\MusicPDFWriteObject{15359}{<< /D [\MusicPDFPageRef{33} /XYZ 71 721 null] >>}
\MusicPDFWriteObject{15360}{<< /D [\MusicPDFPageRef{330} /XYZ 71 721 null] >>}
\MusicPDFWriteObject{15361}{<< /D [\MusicPDFPageRef{331} /XYZ 71 721 null] >>}
\MusicPDFWriteObject{15362}{<< /D [\MusicPDFPageRef{332} /XYZ 71 721 null] >>}
\MusicPDFWriteObject{15363}{<< /D [\MusicPDFPageRef{333} /XYZ 71 721 null] >>}
\MusicPDFWriteObject{15364}{<< /D [\MusicPDFPageRef{334} /XYZ 71 721 null] >>}
\MusicPDFWriteObject{15365}{<< /D [\MusicPDFPageRef{335} /XYZ 71 721 null] >>}
\MusicPDFWriteObject{15366}{<< /D [\MusicPDFPageRef{336} /XYZ 71 721 null] >>}
\MusicPDFWriteObject{15367}{<< /D [\MusicPDFPageRef{337} /XYZ 71 721 null] >>}
\MusicPDFWriteObject{15368}{<< /D [\MusicPDFPageRef{338} /XYZ 71 721 null] >>}
\MusicPDFWriteObject{15369}{<< /D [\MusicPDFPageRef{339} /XYZ 71 721 null] >>}
\MusicPDFWriteObject{15370}{<< /D [\MusicPDFPageRef{34} /XYZ 71 721 null] >>}
\MusicPDFWriteObject{15371}{<< /D [\MusicPDFPageRef{340} /XYZ 71 721 null] >>}
\MusicPDFWriteObject{15372}{<< /D [\MusicPDFPageRef{341} /XYZ 71 721 null] >>}
\MusicPDFWriteObject{15373}{<< /D [\MusicPDFPageRef{342} /XYZ 71 721 null] >>}
\MusicPDFWriteObject{15374}{<< /D [\MusicPDFPageRef{343} /XYZ 71 721 null] >>}
\MusicPDFWriteObject{15375}{<< /D [\MusicPDFPageRef{344} /XYZ 71 721 null] >>}
\MusicPDFWriteObject{15376}{<< /D [\MusicPDFPageRef{345} /XYZ 71 721 null] >>}
\MusicPDFWriteObject{15377}{<< /D [\MusicPDFPageRef{346} /XYZ 71 721 null] >>}
\MusicPDFWriteObject{15378}{<< /D [\MusicPDFPageRef{347} /XYZ 71 721 null] >>}
\MusicPDFWriteObject{15379}{<< /D [\MusicPDFPageRef{348} /XYZ 71 721 null] >>}
\MusicPDFWriteObject{15380}{<< /D [\MusicPDFPageRef{349} /XYZ 71 721 null] >>}
\MusicPDFWriteObject{15381}{<< /D [\MusicPDFPageRef{35} /XYZ 71 721 null] >>}
\MusicPDFWriteObject{15382}{<< /D [\MusicPDFPageRef{350} /XYZ 71 721 null] >>}
\MusicPDFWriteObject{15383}{<< /D [\MusicPDFPageRef{351} /XYZ 71 721 null] >>}
\MusicPDFWriteObject{15384}{<< /D [\MusicPDFPageRef{352} /XYZ 71 721 null] >>}
\MusicPDFWriteObject{15385}{<< /D [\MusicPDFPageRef{353} /XYZ 71 721 null] >>}
\MusicPDFWriteObject{15386}{<< /D [\MusicPDFPageRef{354} /XYZ 71 721 null] >>}
\MusicPDFWriteObject{15387}{<< /D [\MusicPDFPageRef{355} /XYZ 71 721 null] >>}
\MusicPDFWriteObject{15388}{<< /D [\MusicPDFPageRef{356} /XYZ 71 721 null] >>}
\MusicPDFWriteObject{15389}{<< /D [\MusicPDFPageRef{357} /XYZ 71 721 null] >>}
\MusicPDFWriteObject{15390}{<< /D [\MusicPDFPageRef{358} /XYZ 71 721 null] >>}
\MusicPDFWriteObject{15391}{<< /D [\MusicPDFPageRef{359} /XYZ 71 721 null] >>}
\MusicPDFWriteObject{15392}{<< /D [\MusicPDFPageRef{36} /XYZ 71 721 null] >>}
\MusicPDFWriteObject{15393}{<< /D [\MusicPDFPageRef{360} /XYZ 71 721 null] >>}
\MusicPDFWriteObject{15394}{<< /D [\MusicPDFPageRef{361} /XYZ 71 721 null] >>}
\MusicPDFWriteObject{15395}{<< /D [\MusicPDFPageRef{362} /XYZ 71 721 null] >>}
\MusicPDFWriteObject{15396}{<< /D [\MusicPDFPageRef{363} /XYZ 71 721 null] >>}
\MusicPDFWriteObject{15397}{<< /D [\MusicPDFPageRef{364} /XYZ 71 721 null] >>}
\MusicPDFWriteObject{15398}{<< /D [\MusicPDFPageRef{365} /XYZ 71 721 null] >>}
\MusicPDFWriteObject{15399}{<< /D [\MusicPDFPageRef{366} /XYZ 71 721 null] >>}
\MusicPDFWriteObject{15400}{<< /D [\MusicPDFPageRef{367} /XYZ 71 721 null] >>}
\MusicPDFWriteObject{15401}{<< /D [\MusicPDFPageRef{368} /XYZ 71 721 null] >>}
\MusicPDFWriteObject{15402}{<< /D [\MusicPDFPageRef{369} /XYZ 71 721 null] >>}
\MusicPDFWriteObject{15403}{<< /D [\MusicPDFPageRef{37} /XYZ 71 721 null] >>}
\MusicPDFWriteObject{15404}{<< /D [\MusicPDFPageRef{370} /XYZ 71 721 null] >>}
\MusicPDFWriteObject{15405}{<< /D [\MusicPDFPageRef{371} /XYZ 71 721 null] >>}
\MusicPDFWriteObject{15406}{<< /D [\MusicPDFPageRef{372} /XYZ 71 721 null] >>}
\MusicPDFWriteObject{15407}{<< /D [\MusicPDFPageRef{373} /XYZ 71 721 null] >>}
\MusicPDFWriteObject{15408}{<< /D [\MusicPDFPageRef{374} /XYZ 71 721 null] >>}
\MusicPDFWriteObject{15409}{<< /D [\MusicPDFPageRef{375} /XYZ 71 721 null] >>}
\MusicPDFWriteObject{15410}{<< /D [\MusicPDFPageRef{376} /XYZ 71 721 null] >>}
\MusicPDFWriteObject{15411}{<< /D [\MusicPDFPageRef{377} /XYZ 71 721 null] >>}
\MusicPDFWriteObject{15412}{<< /D [\MusicPDFPageRef{378} /XYZ 71 721 null] >>}
\MusicPDFWriteObject{15413}{<< /D [\MusicPDFPageRef{379} /XYZ 71 721 null] >>}
\MusicPDFWriteObject{15414}{<< /D [\MusicPDFPageRef{38} /XYZ 71 721 null] >>}
\MusicPDFWriteObject{15415}{<< /D [\MusicPDFPageRef{380} /XYZ 71 721 null] >>}
\MusicPDFWriteObject{15416}{<< /D [\MusicPDFPageRef{381} /XYZ 71 721 null] >>}
\MusicPDFWriteObject{15417}{<< /D [\MusicPDFPageRef{382} /XYZ 71 721 null] >>}
\MusicPDFWriteObject{15418}{<< /D [\MusicPDFPageRef{383} /XYZ 71 721 null] >>}
\MusicPDFWriteObject{15419}{<< /D [\MusicPDFPageRef{384} /XYZ 71 721 null] >>}
\MusicPDFWriteObject{15420}{<< /D [\MusicPDFPageRef{385} /XYZ 71 721 null] >>}
\MusicPDFWriteObject{15421}{<< /D [\MusicPDFPageRef{386} /XYZ 71 721 null] >>}
\MusicPDFWriteObject{15422}{<< /D [\MusicPDFPageRef{387} /XYZ 71 721 null] >>}
\MusicPDFWriteObject{15423}{<< /D [\MusicPDFPageRef{388} /XYZ 71 721 null] >>}
\MusicPDFWriteObject{15424}{<< /D [\MusicPDFPageRef{389} /XYZ 71 721 null] >>}
\MusicPDFWriteObject{15425}{<< /D [\MusicPDFPageRef{39} /XYZ 71 721 null] >>}
\MusicPDFWriteObject{15426}{<< /D [\MusicPDFPageRef{390} /XYZ 71 721 null] >>}
\MusicPDFWriteObject{15427}{<< /D [\MusicPDFPageRef{391} /XYZ 71 721 null] >>}
\MusicPDFWriteObject{15428}{<< /D [\MusicPDFPageRef{392} /XYZ 71 721 null] >>}
\MusicPDFWriteObject{15429}{<< /D [\MusicPDFPageRef{393} /XYZ 71 721 null] >>}
\MusicPDFWriteObject{15430}{<< /D [\MusicPDFPageRef{394} /XYZ 71 721 null] >>}
\MusicPDFWriteObject{15431}{<< /D [\MusicPDFPageRef{395} /XYZ 71 721 null] >>}
\MusicPDFWriteObject{15432}{<< /D [\MusicPDFPageRef{396} /XYZ 71 721 null] >>}
\MusicPDFWriteObject{15433}{<< /D [\MusicPDFPageRef{397} /XYZ 71 721 null] >>}
\MusicPDFWriteObject{15434}{<< /D [\MusicPDFPageRef{398} /XYZ 71 721 null] >>}
\MusicPDFWriteObject{15435}{<< /D [\MusicPDFPageRef{399} /XYZ 71 721 null] >>}
\MusicPDFWriteObject{15436}{<< /D [\MusicPDFPageRef{4} /XYZ 71 721 null] >>}
\MusicPDFWriteObject{15437}{<< /D [\MusicPDFPageRef{40} /XYZ 71 721 null] >>}
\MusicPDFWriteObject{15438}{<< /D [\MusicPDFPageRef{400} /XYZ 71 721 null] >>}
\MusicPDFWriteObject{15439}{<< /D [\MusicPDFPageRef{401} /XYZ 71 721 null] >>}
\MusicPDFWriteObject{15440}{<< /D [\MusicPDFPageRef{402} /XYZ 71 721 null] >>}
\MusicPDFWriteObject{15441}{<< /D [\MusicPDFPageRef{403} /XYZ 71 721 null] >>}
\MusicPDFWriteObject{15442}{<< /D [\MusicPDFPageRef{404} /XYZ 71 721 null] >>}
\MusicPDFWriteObject{15443}{<< /D [\MusicPDFPageRef{405} /XYZ 71 721 null] >>}
\MusicPDFWriteObject{15444}{<< /D [\MusicPDFPageRef{406} /XYZ 71 721 null] >>}
\MusicPDFWriteObject{15445}{<< /D [\MusicPDFPageRef{407} /XYZ 71 721 null] >>}
\MusicPDFWriteObject{15446}{<< /D [\MusicPDFPageRef{408} /XYZ 71 721 null] >>}
\MusicPDFWriteObject{15447}{<< /D [\MusicPDFPageRef{409} /XYZ 71 721 null] >>}
\MusicPDFWriteObject{15448}{<< /D [\MusicPDFPageRef{41} /XYZ 71 721 null] >>}
\MusicPDFWriteObject{15449}{<< /D [\MusicPDFPageRef{410} /XYZ 71 721 null] >>}
\MusicPDFWriteObject{15450}{<< /D [\MusicPDFPageRef{411} /XYZ 71 721 null] >>}
\MusicPDFWriteObject{15451}{<< /D [\MusicPDFPageRef{412} /XYZ 71 721 null] >>}
\MusicPDFWriteObject{15452}{<< /D [\MusicPDFPageRef{413} /XYZ 71 721 null] >>}
\MusicPDFWriteObject{15453}{<< /D [\MusicPDFPageRef{414} /XYZ 71 721 null] >>}
\MusicPDFWriteObject{15454}{<< /D [\MusicPDFPageRef{415} /XYZ 71 721 null] >>}
\MusicPDFWriteObject{15455}{<< /D [\MusicPDFPageRef{416} /XYZ 71 721 null] >>}
\MusicPDFWriteObject{15456}{<< /D [\MusicPDFPageRef{417} /XYZ 71 721 null] >>}
\MusicPDFWriteObject{15457}{<< /D [\MusicPDFPageRef{418} /XYZ 71 721 null] >>}
\MusicPDFWriteObject{15458}{<< /D [\MusicPDFPageRef{419} /XYZ 71 721 null] >>}
\MusicPDFWriteObject{15459}{<< /D [\MusicPDFPageRef{42} /XYZ 71 721 null] >>}
\MusicPDFWriteObject{15460}{<< /D [\MusicPDFPageRef{420} /XYZ 71 721 null] >>}
\MusicPDFWriteObject{15461}{<< /D [\MusicPDFPageRef{421} /XYZ 71 721 null] >>}
\MusicPDFWriteObject{15462}{<< /D [\MusicPDFPageRef{422} /XYZ 71 721 null] >>}
\MusicPDFWriteObject{15463}{<< /D [\MusicPDFPageRef{423} /XYZ 71 721 null] >>}
\MusicPDFWriteObject{15464}{<< /D [\MusicPDFPageRef{424} /XYZ 71 721 null] >>}
\MusicPDFWriteObject{15465}{<< /D [\MusicPDFPageRef{425} /XYZ 71 721 null] >>}
\MusicPDFWriteObject{15466}{<< /D [\MusicPDFPageRef{426} /XYZ 71 721 null] >>}
\MusicPDFWriteObject{15467}{<< /D [\MusicPDFPageRef{427} /XYZ 71 721 null] >>}
\MusicPDFWriteObject{15468}{<< /D [\MusicPDFPageRef{428} /XYZ 71 721 null] >>}
\MusicPDFWriteObject{15469}{<< /D [\MusicPDFPageRef{429} /XYZ 71 721 null] >>}
\MusicPDFWriteObject{15470}{<< /D [\MusicPDFPageRef{43} /XYZ 71 721 null] >>}
\MusicPDFWriteObject{15471}{<< /D [\MusicPDFPageRef{430} /XYZ 71 721 null] >>}
\MusicPDFWriteObject{15472}{<< /D [\MusicPDFPageRef{431} /XYZ 71 721 null] >>}
\MusicPDFWriteObject{15473}{<< /D [\MusicPDFPageRef{432} /XYZ 71 721 null] >>}
\MusicPDFWriteObject{15474}{<< /D [\MusicPDFPageRef{433} /XYZ 71 721 null] >>}
\MusicPDFWriteObject{15475}{<< /D [\MusicPDFPageRef{434} /XYZ 71 721 null] >>}
\MusicPDFWriteObject{15476}{<< /D [\MusicPDFPageRef{435} /XYZ 71 721 null] >>}
\MusicPDFWriteObject{15477}{<< /D [\MusicPDFPageRef{436} /XYZ 71 721 null] >>}
\MusicPDFWriteObject{15478}{<< /D [\MusicPDFPageRef{437} /XYZ 71 721 null] >>}
\MusicPDFWriteObject{15479}{<< /D [\MusicPDFPageRef{438} /XYZ 71 721 null] >>}
\MusicPDFWriteObject{15480}{<< /D [\MusicPDFPageRef{439} /XYZ 71 721 null] >>}
\MusicPDFWriteObject{15481}{<< /D [\MusicPDFPageRef{44} /XYZ 71 721 null] >>}
\MusicPDFWriteObject{15482}{<< /D [\MusicPDFPageRef{440} /XYZ 71 721 null] >>}
\MusicPDFWriteObject{15483}{<< /D [\MusicPDFPageRef{441} /XYZ 71 721 null] >>}
\MusicPDFWriteObject{15484}{<< /D [\MusicPDFPageRef{442} /XYZ 71 721 null] >>}
\MusicPDFWriteObject{15485}{<< /D [\MusicPDFPageRef{443} /XYZ 71 721 null] >>}
\MusicPDFWriteObject{15486}{<< /D [\MusicPDFPageRef{444} /XYZ 71 721 null] >>}
\MusicPDFWriteObject{15487}{<< /D [\MusicPDFPageRef{445} /XYZ 71 721 null] >>}
\MusicPDFWriteObject{15488}{<< /D [\MusicPDFPageRef{446} /XYZ 71 721 null] >>}
\MusicPDFWriteObject{15489}{<< /D [\MusicPDFPageRef{447} /XYZ 71 721 null] >>}
\MusicPDFWriteObject{15490}{<< /D [\MusicPDFPageRef{448} /XYZ 71 721 null] >>}
\MusicPDFWriteObject{15491}{<< /D [\MusicPDFPageRef{449} /XYZ 71 721 null] >>}
\MusicPDFWriteObject{15492}{<< /D [\MusicPDFPageRef{45} /XYZ 71 721 null] >>}
\MusicPDFWriteObject{15493}{<< /D [\MusicPDFPageRef{450} /XYZ 71 721 null] >>}
\MusicPDFWriteObject{15494}{<< /D [\MusicPDFPageRef{451} /XYZ 71 721 null] >>}
\MusicPDFWriteObject{15495}{<< /D [\MusicPDFPageRef{452} /XYZ 71 721 null] >>}
\MusicPDFWriteObject{15496}{<< /D [\MusicPDFPageRef{453} /XYZ 71 721 null] >>}
\MusicPDFWriteObject{15497}{<< /D [\MusicPDFPageRef{454} /XYZ 71 721 null] >>}
\MusicPDFWriteObject{15498}{<< /D [\MusicPDFPageRef{455} /XYZ 71 721 null] >>}
\MusicPDFWriteObject{15499}{<< /D [\MusicPDFPageRef{456} /XYZ 71 721 null] >>}
\MusicPDFWriteObject{15500}{<< /D [\MusicPDFPageRef{457} /XYZ 71 721 null] >>}
\MusicPDFWriteObject{15501}{<< /D [\MusicPDFPageRef{458} /XYZ 71 721 null] >>}
\MusicPDFWriteObject{15502}{<< /D [\MusicPDFPageRef{459} /XYZ 71 721 null] >>}
\MusicPDFWriteObject{15503}{<< /D [\MusicPDFPageRef{46} /XYZ 71 721 null] >>}
\MusicPDFWriteObject{15504}{<< /D [\MusicPDFPageRef{460} /XYZ 71 721 null] >>}
\MusicPDFWriteObject{15505}{<< /D [\MusicPDFPageRef{461} /XYZ 71 721 null] >>}
\MusicPDFWriteObject{15506}{<< /D [\MusicPDFPageRef{462} /XYZ 71 721 null] >>}
\MusicPDFWriteObject{15507}{<< /D [\MusicPDFPageRef{463} /XYZ 71 721 null] >>}
\MusicPDFWriteObject{15508}{<< /D [\MusicPDFPageRef{464} /XYZ 71 721 null] >>}
\MusicPDFWriteObject{15509}{<< /D [\MusicPDFPageRef{465} /XYZ 71 721 null] >>}
\MusicPDFWriteObject{15510}{<< /D [\MusicPDFPageRef{466} /XYZ 71 721 null] >>}
\MusicPDFWriteObject{15511}{<< /D [\MusicPDFPageRef{467} /XYZ 71 721 null] >>}
\MusicPDFWriteObject{15512}{<< /D [\MusicPDFPageRef{468} /XYZ 71 721 null] >>}
\MusicPDFWriteObject{15513}{<< /D [\MusicPDFPageRef{469} /XYZ 71 721 null] >>}
\MusicPDFWriteObject{15514}{<< /D [\MusicPDFPageRef{47} /XYZ 71 721 null] >>}
\MusicPDFWriteObject{15515}{<< /D [\MusicPDFPageRef{470} /XYZ 71 721 null] >>}
\MusicPDFWriteObject{15516}{<< /D [\MusicPDFPageRef{471} /XYZ 71 721 null] >>}
\MusicPDFWriteObject{15517}{<< /D [\MusicPDFPageRef{472} /XYZ 71 721 null] >>}
\MusicPDFWriteObject{15518}{<< /D [\MusicPDFPageRef{473} /XYZ 71 721 null] >>}
\MusicPDFWriteObject{15519}{<< /D [\MusicPDFPageRef{474} /XYZ 71 721 null] >>}
\MusicPDFWriteObject{15520}{<< /D [\MusicPDFPageRef{475} /XYZ 71 721 null] >>}
\MusicPDFWriteObject{15521}{<< /D [\MusicPDFPageRef{476} /XYZ 71 721 null] >>}
\MusicPDFWriteObject{15522}{<< /D [\MusicPDFPageRef{477} /XYZ 71 721 null] >>}
\MusicPDFWriteObject{15523}{<< /D [\MusicPDFPageRef{478} /XYZ 71 721 null] >>}
\MusicPDFWriteObject{15524}{<< /D [\MusicPDFPageRef{479} /XYZ 71 721 null] >>}
\MusicPDFWriteObject{15525}{<< /D [\MusicPDFPageRef{48} /XYZ 71 721 null] >>}
\MusicPDFWriteObject{15526}{<< /D [\MusicPDFPageRef{480} /XYZ 71 721 null] >>}
\MusicPDFWriteObject{15527}{<< /D [\MusicPDFPageRef{481} /XYZ 71 721 null] >>}
\MusicPDFWriteObject{15528}{<< /D [\MusicPDFPageRef{482} /XYZ 71 721 null] >>}
\MusicPDFWriteObject{15529}{<< /D [\MusicPDFPageRef{483} /XYZ 71 721 null] >>}
\MusicPDFWriteObject{15530}{<< /D [\MusicPDFPageRef{484} /XYZ 71 721 null] >>}
\MusicPDFWriteObject{15531}{<< /D [\MusicPDFPageRef{485} /XYZ 71 721 null] >>}
\MusicPDFWriteObject{15532}{<< /D [\MusicPDFPageRef{486} /XYZ 71 721 null] >>}
\MusicPDFWriteObject{15533}{<< /D [\MusicPDFPageRef{487} /XYZ 71 721 null] >>}
\MusicPDFWriteObject{15534}{<< /D [\MusicPDFPageRef{488} /XYZ 71 721 null] >>}
\MusicPDFWriteObject{15535}{<< /D [\MusicPDFPageRef{489} /XYZ 71 721 null] >>}
\MusicPDFWriteObject{15536}{<< /D [\MusicPDFPageRef{49} /XYZ 71 721 null] >>}
\MusicPDFWriteObject{15537}{<< /D [\MusicPDFPageRef{490} /XYZ 71 721 null] >>}
\MusicPDFWriteObject{15538}{<< /D [\MusicPDFPageRef{491} /XYZ 71 721 null] >>}
\MusicPDFWriteObject{15539}{<< /D [\MusicPDFPageRef{492} /XYZ 71 721 null] >>}
\MusicPDFWriteObject{15540}{<< /D [\MusicPDFPageRef{493} /XYZ 71 721 null] >>}
\MusicPDFWriteObject{15541}{<< /D [\MusicPDFPageRef{494} /XYZ 71 721 null] >>}
\MusicPDFWriteObject{15542}{<< /D [\MusicPDFPageRef{495} /XYZ 71 721 null] >>}
\MusicPDFWriteObject{15543}{<< /D [\MusicPDFPageRef{496} /XYZ 71 721 null] >>}
\MusicPDFWriteObject{15544}{<< /D [\MusicPDFPageRef{497} /XYZ 71 721 null] >>}
\MusicPDFWriteObject{15545}{<< /D [\MusicPDFPageRef{498} /XYZ 71 721 null] >>}
\MusicPDFWriteObject{15546}{<< /D [\MusicPDFPageRef{499} /XYZ 71 721 null] >>}
\MusicPDFWriteObject{15547}{<< /D [\MusicPDFPageRef{5} /XYZ 71 721 null] >>}
\MusicPDFWriteObject{15548}{<< /D [\MusicPDFPageRef{50} /XYZ 71 721 null] >>}
\MusicPDFWriteObject{15549}{<< /D [\MusicPDFPageRef{500} /XYZ 71 721 null] >>}
\MusicPDFWriteObject{15550}{<< /D [\MusicPDFPageRef{501} /XYZ 71 721 null] >>}
\MusicPDFWriteObject{15551}{<< /D [\MusicPDFPageRef{502} /XYZ 71 721 null] >>}
\MusicPDFWriteObject{15552}{<< /D [\MusicPDFPageRef{503} /XYZ 71 721 null] >>}
\MusicPDFWriteObject{15553}{<< /D [\MusicPDFPageRef{504} /XYZ 71 721 null] >>}
\MusicPDFWriteObject{15554}{<< /D [\MusicPDFPageRef{505} /XYZ 71 721 null] >>}
\MusicPDFWriteObject{15555}{<< /D [\MusicPDFPageRef{506} /XYZ 71 721 null] >>}
\MusicPDFWriteObject{15556}{<< /D [\MusicPDFPageRef{507} /XYZ 71 721 null] >>}
\MusicPDFWriteObject{15557}{<< /D [\MusicPDFPageRef{508} /XYZ 71 721 null] >>}
\MusicPDFWriteObject{15558}{<< /D [\MusicPDFPageRef{509} /XYZ 71 721 null] >>}
\MusicPDFWriteObject{15559}{<< /D [\MusicPDFPageRef{51} /XYZ 71 721 null] >>}
\MusicPDFWriteObject{15560}{<< /D [\MusicPDFPageRef{510} /XYZ 71 721 null] >>}
\MusicPDFWriteObject{15561}{<< /D [\MusicPDFPageRef{511} /XYZ 71 721 null] >>}
\MusicPDFWriteObject{15562}{<< /D [\MusicPDFPageRef{512} /XYZ 71 721 null] >>}
\MusicPDFWriteObject{15563}{<< /D [\MusicPDFPageRef{513} /XYZ 71 721 null] >>}
\MusicPDFWriteObject{15564}{<< /D [\MusicPDFPageRef{514} /XYZ 71 721 null] >>}
\MusicPDFWriteObject{15565}{<< /D [\MusicPDFPageRef{515} /XYZ 71 721 null] >>}
\MusicPDFWriteObject{15566}{<< /D [\MusicPDFPageRef{516} /XYZ 71 721 null] >>}
\MusicPDFWriteObject{15567}{<< /D [\MusicPDFPageRef{517} /XYZ 71 721 null] >>}
\MusicPDFWriteObject{15568}{<< /D [\MusicPDFPageRef{518} /XYZ 71 721 null] >>}
\MusicPDFWriteObject{15569}{<< /D [\MusicPDFPageRef{519} /XYZ 71 721 null] >>}
\MusicPDFWriteObject{15570}{<< /D [\MusicPDFPageRef{52} /XYZ 71 721 null] >>}
\MusicPDFWriteObject{15571}{<< /D [\MusicPDFPageRef{520} /XYZ 71 721 null] >>}
\MusicPDFWriteObject{15572}{<< /D [\MusicPDFPageRef{521} /XYZ 71 721 null] >>}
\MusicPDFWriteObject{15573}{<< /D [\MusicPDFPageRef{522} /XYZ 71 721 null] >>}
\MusicPDFWriteObject{15574}{<< /D [\MusicPDFPageRef{523} /XYZ 71 721 null] >>}
\MusicPDFWriteObject{15575}{<< /D [\MusicPDFPageRef{524} /XYZ 71 721 null] >>}
\MusicPDFWriteObject{15576}{<< /D [\MusicPDFPageRef{525} /XYZ 71 721 null] >>}
\MusicPDFWriteObject{15577}{<< /D [\MusicPDFPageRef{526} /XYZ 71 721 null] >>}
\MusicPDFWriteObject{15578}{<< /D [\MusicPDFPageRef{527} /XYZ 71 721 null] >>}
\MusicPDFWriteObject{15579}{<< /D [\MusicPDFPageRef{528} /XYZ 71 721 null] >>}
\MusicPDFWriteObject{15580}{<< /D [\MusicPDFPageRef{529} /XYZ 71 721 null] >>}
\MusicPDFWriteObject{15581}{<< /D [\MusicPDFPageRef{53} /XYZ 71 721 null] >>}
\MusicPDFWriteObject{15582}{<< /D [\MusicPDFPageRef{530} /XYZ 71 721 null] >>}
\MusicPDFWriteObject{15583}{<< /D [\MusicPDFPageRef{531} /XYZ 71 721 null] >>}
\MusicPDFWriteObject{15584}{<< /D [\MusicPDFPageRef{532} /XYZ 71 721 null] >>}
\MusicPDFWriteObject{15585}{<< /D [\MusicPDFPageRef{533} /XYZ 71 721 null] >>}
\MusicPDFWriteObject{15586}{<< /D [\MusicPDFPageRef{534} /XYZ 71 721 null] >>}
\MusicPDFWriteObject{15587}{<< /D [\MusicPDFPageRef{535} /XYZ 71 721 null] >>}
\MusicPDFWriteObject{15588}{<< /D [\MusicPDFPageRef{536} /XYZ 71 721 null] >>}
\MusicPDFWriteObject{15589}{<< /D [\MusicPDFPageRef{537} /XYZ 71 721 null] >>}
\MusicPDFWriteObject{15590}{<< /D [\MusicPDFPageRef{538} /XYZ 71 721 null] >>}
\MusicPDFWriteObject{15591}{<< /D [\MusicPDFPageRef{539} /XYZ 71 721 null] >>}
\MusicPDFWriteObject{15592}{<< /D [\MusicPDFPageRef{54} /XYZ 71 721 null] >>}
\MusicPDFWriteObject{15593}{<< /D [\MusicPDFPageRef{540} /XYZ 71 721 null] >>}
\MusicPDFWriteObject{15594}{<< /D [\MusicPDFPageRef{541} /XYZ 71 721 null] >>}
\MusicPDFWriteObject{15595}{<< /D [\MusicPDFPageRef{542} /XYZ 71 721 null] >>}
\MusicPDFWriteObject{15596}{<< /D [\MusicPDFPageRef{543} /XYZ 71 721 null] >>}
\MusicPDFWriteObject{15597}{<< /D [\MusicPDFPageRef{544} /XYZ 71 721 null] >>}
\MusicPDFWriteObject{15598}{<< /D [\MusicPDFPageRef{545} /XYZ 71 721 null] >>}
\MusicPDFWriteObject{15599}{<< /D [\MusicPDFPageRef{546} /XYZ 71 721 null] >>}
\MusicPDFWriteObject{15600}{<< /D [\MusicPDFPageRef{547} /XYZ 71 721 null] >>}
\MusicPDFWriteObject{15601}{<< /D [\MusicPDFPageRef{548} /XYZ 71 721 null] >>}
\MusicPDFWriteObject{15602}{<< /D [\MusicPDFPageRef{549} /XYZ 71 721 null] >>}
\MusicPDFWriteObject{15603}{<< /D [\MusicPDFPageRef{55} /XYZ 71 721 null] >>}
\MusicPDFWriteObject{15604}{<< /D [\MusicPDFPageRef{550} /XYZ 71 721 null] >>}
\MusicPDFWriteObject{15605}{<< /D [\MusicPDFPageRef{551} /XYZ 71 721 null] >>}
\MusicPDFWriteObject{15606}{<< /D [\MusicPDFPageRef{552} /XYZ 71 721 null] >>}
\MusicPDFWriteObject{15607}{<< /D [\MusicPDFPageRef{553} /XYZ 71 721 null] >>}
\MusicPDFWriteObject{15608}{<< /D [\MusicPDFPageRef{554} /XYZ 71 721 null] >>}
\MusicPDFWriteObject{15609}{<< /D [\MusicPDFPageRef{555} /XYZ 71 721 null] >>}
\MusicPDFWriteObject{15610}{<< /D [\MusicPDFPageRef{556} /XYZ 71 721 null] >>}
\MusicPDFWriteObject{15611}{<< /D [\MusicPDFPageRef{557} /XYZ 71 721 null] >>}
\MusicPDFWriteObject{15612}{<< /D [\MusicPDFPageRef{558} /XYZ 71 721 null] >>}
\MusicPDFWriteObject{15613}{<< /D [\MusicPDFPageRef{559} /XYZ 71 721 null] >>}
\MusicPDFWriteObject{15614}{<< /D [\MusicPDFPageRef{56} /XYZ 71 721 null] >>}
\MusicPDFWriteObject{15615}{<< /D [\MusicPDFPageRef{560} /XYZ 71 721 null] >>}
\MusicPDFWriteObject{15616}{<< /D [\MusicPDFPageRef{561} /XYZ 71 721 null] >>}
\MusicPDFWriteObject{15617}{<< /D [\MusicPDFPageRef{562} /XYZ 71 721 null] >>}
\MusicPDFWriteObject{15618}{<< /D [\MusicPDFPageRef{563} /XYZ 71 721 null] >>}
\MusicPDFWriteObject{15619}{<< /D [\MusicPDFPageRef{564} /XYZ 71 721 null] >>}
\MusicPDFWriteObject{15620}{<< /D [\MusicPDFPageRef{565} /XYZ 71 721 null] >>}
\MusicPDFWriteObject{15621}{<< /D [\MusicPDFPageRef{566} /XYZ 71 721 null] >>}
\MusicPDFWriteObject{15622}{<< /D [\MusicPDFPageRef{567} /XYZ 71 721 null] >>}
\MusicPDFWriteObject{15623}{<< /D [\MusicPDFPageRef{568} /XYZ 71 721 null] >>}
\MusicPDFWriteObject{15624}{<< /D [\MusicPDFPageRef{569} /XYZ 71 721 null] >>}
\MusicPDFWriteObject{15625}{<< /D [\MusicPDFPageRef{57} /XYZ 71 721 null] >>}
\MusicPDFWriteObject{15626}{<< /D [\MusicPDFPageRef{570} /XYZ 71 721 null] >>}
\MusicPDFWriteObject{15627}{<< /D [\MusicPDFPageRef{571} /XYZ 71 721 null] >>}
\MusicPDFWriteObject{15628}{<< /D [\MusicPDFPageRef{572} /XYZ 71 721 null] >>}
\MusicPDFWriteObject{15629}{<< /D [\MusicPDFPageRef{573} /XYZ 71 721 null] >>}
\MusicPDFWriteObject{15630}{<< /D [\MusicPDFPageRef{574} /XYZ 71 721 null] >>}
\MusicPDFWriteObject{15631}{<< /D [\MusicPDFPageRef{575} /XYZ 71 721 null] >>}
\MusicPDFWriteObject{15632}{<< /D [\MusicPDFPageRef{576} /XYZ 71 721 null] >>}
\MusicPDFWriteObject{15633}{<< /D [\MusicPDFPageRef{577} /XYZ 71 721 null] >>}
\MusicPDFWriteObject{15634}{<< /D [\MusicPDFPageRef{578} /XYZ 71 721 null] >>}
\MusicPDFWriteObject{15635}{<< /D [\MusicPDFPageRef{579} /XYZ 71 721 null] >>}
\MusicPDFWriteObject{15636}{<< /D [\MusicPDFPageRef{58} /XYZ 71 721 null] >>}
\MusicPDFWriteObject{15637}{<< /D [\MusicPDFPageRef{580} /XYZ 71 721 null] >>}
\MusicPDFWriteObject{15638}{<< /D [\MusicPDFPageRef{581} /XYZ 71 721 null] >>}
\MusicPDFWriteObject{15639}{<< /D [\MusicPDFPageRef{582} /XYZ 71 721 null] >>}
\MusicPDFWriteObject{15640}{<< /D [\MusicPDFPageRef{583} /XYZ 71 721 null] >>}
\MusicPDFWriteObject{15641}{<< /D [\MusicPDFPageRef{584} /XYZ 71 721 null] >>}
\MusicPDFWriteObject{15642}{<< /D [\MusicPDFPageRef{585} /XYZ 71 721 null] >>}
\MusicPDFWriteObject{15643}{<< /D [\MusicPDFPageRef{586} /XYZ 71 721 null] >>}
\MusicPDFWriteObject{15644}{<< /D [\MusicPDFPageRef{587} /XYZ 71 721 null] >>}
\MusicPDFWriteObject{15645}{<< /D [\MusicPDFPageRef{588} /XYZ 71 721 null] >>}
\MusicPDFWriteObject{15646}{<< /D [\MusicPDFPageRef{589} /XYZ 71 721 null] >>}
\MusicPDFWriteObject{15647}{<< /D [\MusicPDFPageRef{59} /XYZ 71 721 null] >>}
\MusicPDFWriteObject{15648}{<< /D [\MusicPDFPageRef{590} /XYZ 71 721 null] >>}
\MusicPDFWriteObject{15649}{<< /D [\MusicPDFPageRef{591} /XYZ 71 721 null] >>}
\MusicPDFWriteObject{15650}{<< /D [\MusicPDFPageRef{592} /XYZ 71 721 null] >>}
\MusicPDFWriteObject{15651}{<< /D [\MusicPDFPageRef{593} /XYZ 71 721 null] >>}
\MusicPDFWriteObject{15652}{<< /D [\MusicPDFPageRef{594} /XYZ 71 721 null] >>}
\MusicPDFWriteObject{15653}{<< /D [\MusicPDFPageRef{595} /XYZ 71 721 null] >>}
\MusicPDFWriteObject{15654}{<< /D [\MusicPDFPageRef{596} /XYZ 71 721 null] >>}
\MusicPDFWriteObject{15655}{<< /D [\MusicPDFPageRef{597} /XYZ 71 721 null] >>}
\MusicPDFWriteObject{15656}{<< /D [\MusicPDFPageRef{598} /XYZ 71 721 null] >>}
\MusicPDFWriteObject{15657}{<< /D [\MusicPDFPageRef{599} /XYZ 71 721 null] >>}
\MusicPDFWriteObject{15658}{<< /D [\MusicPDFPageRef{6} /XYZ 71 721 null] >>}
\MusicPDFWriteObject{15659}{<< /D [\MusicPDFPageRef{60} /XYZ 71 721 null] >>}
\MusicPDFWriteObject{15660}{<< /D [\MusicPDFPageRef{600} /XYZ 71 721 null] >>}
\MusicPDFWriteObject{15661}{<< /D [\MusicPDFPageRef{601} /XYZ 71 721 null] >>}
\MusicPDFWriteObject{15662}{<< /D [\MusicPDFPageRef{602} /XYZ 71 721 null] >>}
\MusicPDFWriteObject{15663}{<< /D [\MusicPDFPageRef{603} /XYZ 71 721 null] >>}
\MusicPDFWriteObject{15664}{<< /D [\MusicPDFPageRef{604} /XYZ 71 721 null] >>}
\MusicPDFWriteObject{15665}{<< /D [\MusicPDFPageRef{605} /XYZ 71 721 null] >>}
\MusicPDFWriteObject{15666}{<< /D [\MusicPDFPageRef{606} /XYZ 71 721 null] >>}
\MusicPDFWriteObject{15667}{<< /D [\MusicPDFPageRef{607} /XYZ 71 721 null] >>}
\MusicPDFWriteObject{15668}{<< /D [\MusicPDFPageRef{608} /XYZ 71 721 null] >>}
\MusicPDFWriteObject{15669}{<< /D [\MusicPDFPageRef{609} /XYZ 71 721 null] >>}
\MusicPDFWriteObject{15670}{<< /D [\MusicPDFPageRef{61} /XYZ 71 721 null] >>}
\MusicPDFWriteObject{15671}{<< /D [\MusicPDFPageRef{610} /XYZ 71 721 null] >>}
\MusicPDFWriteObject{15672}{<< /D [\MusicPDFPageRef{611} /XYZ 71 721 null] >>}
\MusicPDFWriteObject{15673}{<< /D [\MusicPDFPageRef{612} /XYZ 71 721 null] >>}
\MusicPDFWriteObject{15674}{<< /D [\MusicPDFPageRef{613} /XYZ 71 721 null] >>}
\MusicPDFWriteObject{15675}{<< /D [\MusicPDFPageRef{614} /XYZ 71 721 null] >>}
\MusicPDFWriteObject{15676}{<< /D [\MusicPDFPageRef{615} /XYZ 71 721 null] >>}
\MusicPDFWriteObject{15677}{<< /D [\MusicPDFPageRef{616} /XYZ 71 721 null] >>}
\MusicPDFWriteObject{15678}{<< /D [\MusicPDFPageRef{617} /XYZ 71 721 null] >>}
\MusicPDFWriteObject{15679}{<< /D [\MusicPDFPageRef{618} /XYZ 71 721 null] >>}
\MusicPDFWriteObject{15680}{<< /D [\MusicPDFPageRef{619} /XYZ 71 721 null] >>}
\MusicPDFWriteObject{15681}{<< /D [\MusicPDFPageRef{62} /XYZ 71 721 null] >>}
\MusicPDFWriteObject{15682}{<< /D [\MusicPDFPageRef{620} /XYZ 71 721 null] >>}
\MusicPDFWriteObject{15683}{<< /D [\MusicPDFPageRef{621} /XYZ 71 721 null] >>}
\MusicPDFWriteObject{15684}{<< /D [\MusicPDFPageRef{622} /XYZ 71 721 null] >>}
\MusicPDFWriteObject{15685}{<< /D [\MusicPDFPageRef{623} /XYZ 71 721 null] >>}
\MusicPDFWriteObject{15686}{<< /D [\MusicPDFPageRef{624} /XYZ 71 721 null] >>}
\MusicPDFWriteObject{15687}{<< /D [\MusicPDFPageRef{625} /XYZ 71 721 null] >>}
\MusicPDFWriteObject{15688}{<< /D [\MusicPDFPageRef{626} /XYZ 71 721 null] >>}
\MusicPDFWriteObject{15689}{<< /D [\MusicPDFPageRef{627} /XYZ 71 721 null] >>}
\MusicPDFWriteObject{15690}{<< /D [\MusicPDFPageRef{628} /XYZ 71 721 null] >>}
\MusicPDFWriteObject{15691}{<< /D [\MusicPDFPageRef{629} /XYZ 71 721 null] >>}
\MusicPDFWriteObject{15692}{<< /D [\MusicPDFPageRef{63} /XYZ 71 721 null] >>}
\MusicPDFWriteObject{15693}{<< /D [\MusicPDFPageRef{630} /XYZ 71 721 null] >>}
\MusicPDFWriteObject{15694}{<< /D [\MusicPDFPageRef{631} /XYZ 71 721 null] >>}
\MusicPDFWriteObject{15695}{<< /D [\MusicPDFPageRef{632} /XYZ 71 721 null] >>}
\MusicPDFWriteObject{15696}{<< /D [\MusicPDFPageRef{633} /XYZ 71 721 null] >>}
\MusicPDFWriteObject{15697}{<< /D [\MusicPDFPageRef{634} /XYZ 71 721 null] >>}
\MusicPDFWriteObject{15698}{<< /D [\MusicPDFPageRef{635} /XYZ 71 721 null] >>}
\MusicPDFWriteObject{15699}{<< /D [\MusicPDFPageRef{636} /XYZ 71 721 null] >>}
\MusicPDFWriteObject{15700}{<< /D [\MusicPDFPageRef{637} /XYZ 71 721 null] >>}
\MusicPDFWriteObject{15701}{<< /D [\MusicPDFPageRef{638} /XYZ 71 721 null] >>}
\MusicPDFWriteObject{15702}{<< /D [\MusicPDFPageRef{639} /XYZ 71 721 null] >>}
\MusicPDFWriteObject{15703}{<< /D [\MusicPDFPageRef{64} /XYZ 71 721 null] >>}
\MusicPDFWriteObject{15704}{<< /D [\MusicPDFPageRef{640} /XYZ 71 721 null] >>}
\MusicPDFWriteObject{15705}{<< /D [\MusicPDFPageRef{641} /XYZ 71 721 null] >>}
\MusicPDFWriteObject{15706}{<< /D [\MusicPDFPageRef{642} /XYZ 71 721 null] >>}
\MusicPDFWriteObject{15707}{<< /D [\MusicPDFPageRef{643} /XYZ 71 721 null] >>}
\MusicPDFWriteObject{15708}{<< /D [\MusicPDFPageRef{644} /XYZ 71 721 null] >>}
\MusicPDFWriteObject{15709}{<< /D [\MusicPDFPageRef{645} /XYZ 71 721 null] >>}
\MusicPDFWriteObject{15710}{<< /D [\MusicPDFPageRef{646} /XYZ 71 721 null] >>}
\MusicPDFWriteObject{15711}{<< /D [\MusicPDFPageRef{647} /XYZ 71 721 null] >>}
\MusicPDFWriteObject{15712}{<< /D [\MusicPDFPageRef{648} /XYZ 71 721 null] >>}
\MusicPDFWriteObject{15713}{<< /D [\MusicPDFPageRef{649} /XYZ 71 721 null] >>}
\MusicPDFWriteObject{15714}{<< /D [\MusicPDFPageRef{65} /XYZ 71 721 null] >>}
\MusicPDFWriteObject{15715}{<< /D [\MusicPDFPageRef{650} /XYZ 71 721 null] >>}
\MusicPDFWriteObject{15716}{<< /D [\MusicPDFPageRef{651} /XYZ 71 721 null] >>}
\MusicPDFWriteObject{15717}{<< /D [\MusicPDFPageRef{652} /XYZ 71 721 null] >>}
\MusicPDFWriteObject{15718}{<< /D [\MusicPDFPageRef{653} /XYZ 71 721 null] >>}
\MusicPDFWriteObject{15719}{<< /D [\MusicPDFPageRef{654} /XYZ 71 721 null] >>}
\MusicPDFWriteObject{15720}{<< /D [\MusicPDFPageRef{655} /XYZ 71 721 null] >>}
\MusicPDFWriteObject{15721}{<< /D [\MusicPDFPageRef{656} /XYZ 71 721 null] >>}
\MusicPDFWriteObject{15722}{<< /D [\MusicPDFPageRef{657} /XYZ 71 721 null] >>}
\MusicPDFWriteObject{15723}{<< /D [\MusicPDFPageRef{658} /XYZ 71 721 null] >>}
\MusicPDFWriteObject{15724}{<< /D [\MusicPDFPageRef{659} /XYZ 71 721 null] >>}
\MusicPDFWriteObject{15725}{<< /D [\MusicPDFPageRef{66} /XYZ 71 721 null] >>}
\MusicPDFWriteObject{15726}{<< /D [\MusicPDFPageRef{660} /XYZ 71 721 null] >>}
\MusicPDFWriteObject{15727}{<< /D [\MusicPDFPageRef{661} /XYZ 71 721 null] >>}
\MusicPDFWriteObject{15728}{<< /D [\MusicPDFPageRef{662} /XYZ 71 721 null] >>}
\MusicPDFWriteObject{15729}{<< /D [\MusicPDFPageRef{663} /XYZ 71 721 null] >>}
\MusicPDFWriteObject{15730}{<< /D [\MusicPDFPageRef{664} /XYZ 71 721 null] >>}
\MusicPDFWriteObject{15731}{<< /D [\MusicPDFPageRef{665} /XYZ 71 721 null] >>}
\MusicPDFWriteObject{15732}{<< /D [\MusicPDFPageRef{666} /XYZ 71 721 null] >>}
\MusicPDFWriteObject{15733}{<< /D [\MusicPDFPageRef{667} /XYZ 71 721 null] >>}
\MusicPDFWriteObject{15734}{<< /D [\MusicPDFPageRef{668} /XYZ 71 721 null] >>}
\MusicPDFWriteObject{15735}{<< /D [\MusicPDFPageRef{669} /XYZ 71 721 null] >>}
\MusicPDFWriteObject{15736}{<< /D [\MusicPDFPageRef{67} /XYZ 71 721 null] >>}
\MusicPDFWriteObject{15737}{<< /D [\MusicPDFPageRef{670} /XYZ 71 721 null] >>}
\MusicPDFWriteObject{15738}{<< /D [\MusicPDFPageRef{671} /XYZ 71 721 null] >>}
\MusicPDFWriteObject{15739}{<< /D [\MusicPDFPageRef{672} /XYZ 71 721 null] >>}
\MusicPDFWriteObject{15740}{<< /D [\MusicPDFPageRef{673} /XYZ 71 721 null] >>}
\MusicPDFWriteObject{15741}{<< /D [\MusicPDFPageRef{674} /XYZ 71 721 null] >>}
\MusicPDFWriteObject{15742}{<< /D [\MusicPDFPageRef{675} /XYZ 71 721 null] >>}
\MusicPDFWriteObject{15743}{<< /D [\MusicPDFPageRef{676} /XYZ 71 721 null] >>}
\MusicPDFWriteObject{15744}{<< /D [\MusicPDFPageRef{677} /XYZ 71 721 null] >>}
\MusicPDFWriteObject{15745}{<< /D [\MusicPDFPageRef{678} /XYZ 71 721 null] >>}
\MusicPDFWriteObject{15746}{<< /D [\MusicPDFPageRef{679} /XYZ 71 721 null] >>}
\MusicPDFWriteObject{15747}{<< /D [\MusicPDFPageRef{68} /XYZ 71 721 null] >>}
\MusicPDFWriteObject{15748}{<< /D [\MusicPDFPageRef{680} /XYZ 71 721 null] >>}
\MusicPDFWriteObject{15749}{<< /D [\MusicPDFPageRef{681} /XYZ 71 721 null] >>}
\MusicPDFWriteObject{15750}{<< /D [\MusicPDFPageRef{682} /XYZ 71 721 null] >>}
\MusicPDFWriteObject{15751}{<< /D [\MusicPDFPageRef{683} /XYZ 71 721 null] >>}
\MusicPDFWriteObject{15752}{<< /D [\MusicPDFPageRef{684} /XYZ 71 721 null] >>}
\MusicPDFWriteObject{15753}{<< /D [\MusicPDFPageRef{685} /XYZ 71 721 null] >>}
\MusicPDFWriteObject{15754}{<< /D [\MusicPDFPageRef{686} /XYZ 71 721 null] >>}
\MusicPDFWriteObject{15755}{<< /D [\MusicPDFPageRef{687} /XYZ 71 721 null] >>}
\MusicPDFWriteObject{15756}{<< /D [\MusicPDFPageRef{688} /XYZ 71 721 null] >>}
\MusicPDFWriteObject{15757}{<< /D [\MusicPDFPageRef{689} /XYZ 71 721 null] >>}
\MusicPDFWriteObject{15758}{<< /D [\MusicPDFPageRef{69} /XYZ 71 721 null] >>}
\MusicPDFWriteObject{15759}{<< /D [\MusicPDFPageRef{690} /XYZ 71 721 null] >>}
\MusicPDFWriteObject{15760}{<< /D [\MusicPDFPageRef{691} /XYZ 71 721 null] >>}
\MusicPDFWriteObject{15761}{<< /D [\MusicPDFPageRef{692} /XYZ 71 721 null] >>}
\MusicPDFWriteObject{15762}{<< /D [\MusicPDFPageRef{693} /XYZ 71 721 null] >>}
\MusicPDFWriteObject{15763}{<< /D [\MusicPDFPageRef{694} /XYZ 71 721 null] >>}
\MusicPDFWriteObject{15764}{<< /D [\MusicPDFPageRef{695} /XYZ 71 721 null] >>}
\MusicPDFWriteObject{15765}{<< /D [\MusicPDFPageRef{696} /XYZ 71 721 null] >>}
\MusicPDFWriteObject{15766}{<< /D [\MusicPDFPageRef{697} /XYZ 71 721 null] >>}
\MusicPDFWriteObject{15767}{<< /D [\MusicPDFPageRef{698} /XYZ 71 721 null] >>}
\MusicPDFWriteObject{15768}{<< /D [\MusicPDFPageRef{699} /XYZ 71 721 null] >>}
\MusicPDFWriteObject{15769}{<< /D [\MusicPDFPageRef{7} /XYZ 71 721 null] >>}
\MusicPDFWriteObject{15770}{<< /D [\MusicPDFPageRef{70} /XYZ 71 721 null] >>}
\MusicPDFWriteObject{15771}{<< /D [\MusicPDFPageRef{700} /XYZ 71 721 null] >>}
\MusicPDFWriteObject{15772}{<< /D [\MusicPDFPageRef{701} /XYZ 71 721 null] >>}
\MusicPDFWriteObject{15773}{<< /D [\MusicPDFPageRef{702} /XYZ 71 721 null] >>}
\MusicPDFWriteObject{15774}{<< /D [\MusicPDFPageRef{703} /XYZ 71 721 null] >>}
\MusicPDFWriteObject{15775}{<< /D [\MusicPDFPageRef{704} /XYZ 71 721 null] >>}
\MusicPDFWriteObject{15776}{<< /D [\MusicPDFPageRef{705} /XYZ 71 721 null] >>}
\MusicPDFWriteObject{15777}{<< /D [\MusicPDFPageRef{706} /XYZ 71 721 null] >>}
\MusicPDFWriteObject{15778}{<< /D [\MusicPDFPageRef{707} /XYZ 71 721 null] >>}
\MusicPDFWriteObject{15779}{<< /D [\MusicPDFPageRef{708} /XYZ 71 721 null] >>}
\MusicPDFWriteObject{15780}{<< /D [\MusicPDFPageRef{709} /XYZ 71 721 null] >>}
\MusicPDFWriteObject{15781}{<< /D [\MusicPDFPageRef{71} /XYZ 71 721 null] >>}
\MusicPDFWriteObject{15782}{<< /D [\MusicPDFPageRef{710} /XYZ 71 721 null] >>}
\MusicPDFWriteObject{15783}{<< /D [\MusicPDFPageRef{711} /XYZ 71 721 null] >>}
\MusicPDFWriteObject{15784}{<< /D [\MusicPDFPageRef{712} /XYZ 71 721 null] >>}
\MusicPDFWriteObject{15785}{<< /D [\MusicPDFPageRef{713} /XYZ 71 721 null] >>}
\MusicPDFWriteObject{15786}{<< /D [\MusicPDFPageRef{714} /XYZ 71 721 null] >>}
\MusicPDFWriteObject{15787}{<< /D [\MusicPDFPageRef{715} /XYZ 71 721 null] >>}
\MusicPDFWriteObject{15788}{<< /D [\MusicPDFPageRef{716} /XYZ 71 721 null] >>}
\MusicPDFWriteObject{15789}{<< /D [\MusicPDFPageRef{717} /XYZ 71 721 null] >>}
\MusicPDFWriteObject{15790}{<< /D [\MusicPDFPageRef{718} /XYZ 71 721 null] >>}
\MusicPDFWriteObject{15791}{<< /D [\MusicPDFPageRef{719} /XYZ 71 721 null] >>}
\MusicPDFWriteObject{15792}{<< /D [\MusicPDFPageRef{72} /XYZ 71 721 null] >>}
\MusicPDFWriteObject{15793}{<< /D [\MusicPDFPageRef{720} /XYZ 71 721 null] >>}
\MusicPDFWriteObject{15794}{<< /D [\MusicPDFPageRef{721} /XYZ 71 721 null] >>}
\MusicPDFWriteObject{15795}{<< /D [\MusicPDFPageRef{722} /XYZ 71 721 null] >>}
\MusicPDFWriteObject{15796}{<< /D [\MusicPDFPageRef{723} /XYZ 71 721 null] >>}
\MusicPDFWriteObject{15797}{<< /D [\MusicPDFPageRef{724} /XYZ 71 721 null] >>}
\MusicPDFWriteObject{15798}{<< /D [\MusicPDFPageRef{725} /XYZ 71 721 null] >>}
\MusicPDFWriteObject{15799}{<< /D [\MusicPDFPageRef{726} /XYZ 71 721 null] >>}
\MusicPDFWriteObject{15800}{<< /D [\MusicPDFPageRef{727} /XYZ 71 721 null] >>}
\MusicPDFWriteObject{15801}{<< /D [\MusicPDFPageRef{728} /XYZ 71 721 null] >>}
\MusicPDFWriteObject{15802}{<< /D [\MusicPDFPageRef{729} /XYZ 71 721 null] >>}
\MusicPDFWriteObject{15803}{<< /D [\MusicPDFPageRef{73} /XYZ 71 721 null] >>}
\MusicPDFWriteObject{15804}{<< /D [\MusicPDFPageRef{730} /XYZ 71 721 null] >>}
\MusicPDFWriteObject{15805}{<< /D [\MusicPDFPageRef{731} /XYZ 71 721 null] >>}
\MusicPDFWriteObject{15806}{<< /D [\MusicPDFPageRef{732} /XYZ 71 721 null] >>}
\MusicPDFWriteObject{15807}{<< /D [\MusicPDFPageRef{733} /XYZ 71 721 null] >>}
\MusicPDFWriteObject{15808}{<< /D [\MusicPDFPageRef{734} /XYZ 71 721 null] >>}
\MusicPDFWriteObject{15809}{<< /D [\MusicPDFPageRef{735} /XYZ 71 721 null] >>}
\MusicPDFWriteObject{15810}{<< /D [\MusicPDFPageRef{736} /XYZ 71 721 null] >>}
\MusicPDFWriteObject{15811}{<< /D [\MusicPDFPageRef{737} /XYZ 71 721 null] >>}
\MusicPDFWriteObject{15812}{<< /D [\MusicPDFPageRef{738} /XYZ 71 721 null] >>}
\MusicPDFWriteObject{15813}{<< /D [\MusicPDFPageRef{739} /XYZ 71 721 null] >>}
\MusicPDFWriteObject{15814}{<< /D [\MusicPDFPageRef{74} /XYZ 71 721 null] >>}
\MusicPDFWriteObject{15815}{<< /D [\MusicPDFPageRef{740} /XYZ 71 721 null] >>}
\MusicPDFWriteObject{15816}{<< /D [\MusicPDFPageRef{741} /XYZ 71 721 null] >>}
\MusicPDFWriteObject{15817}{<< /D [\MusicPDFPageRef{742} /XYZ 71 721 null] >>}
\MusicPDFWriteObject{15818}{<< /D [\MusicPDFPageRef{743} /XYZ 71 721 null] >>}
\MusicPDFWriteObject{15819}{<< /D [\MusicPDFPageRef{744} /XYZ 71 721 null] >>}
\MusicPDFWriteObject{15820}{<< /D [\MusicPDFPageRef{745} /XYZ 71 721 null] >>}
\MusicPDFWriteObject{15821}{<< /D [\MusicPDFPageRef{746} /XYZ 71 721 null] >>}
\MusicPDFWriteObject{15822}{<< /D [\MusicPDFPageRef{747} /XYZ 71 721 null] >>}
\MusicPDFWriteObject{15823}{<< /D [\MusicPDFPageRef{748} /XYZ 71 721 null] >>}
\MusicPDFWriteObject{15824}{<< /D [\MusicPDFPageRef{749} /XYZ 71 721 null] >>}
\MusicPDFWriteObject{15825}{<< /D [\MusicPDFPageRef{75} /XYZ 71 721 null] >>}
\MusicPDFWriteObject{15826}{<< /D [\MusicPDFPageRef{750} /XYZ 71 721 null] >>}
\MusicPDFWriteObject{15827}{<< /D [\MusicPDFPageRef{751} /XYZ 71 721 null] >>}
\MusicPDFWriteObject{15828}{<< /D [\MusicPDFPageRef{752} /XYZ 71 721 null] >>}
\MusicPDFWriteObject{15829}{<< /D [\MusicPDFPageRef{753} /XYZ 71 721 null] >>}
\MusicPDFWriteObject{15830}{<< /D [\MusicPDFPageRef{754} /XYZ 71 721 null] >>}
\MusicPDFWriteObject{15831}{<< /D [\MusicPDFPageRef{755} /XYZ 71 721 null] >>}
\MusicPDFWriteObject{15832}{<< /D [\MusicPDFPageRef{756} /XYZ 71 721 null] >>}
\MusicPDFWriteObject{15833}{<< /D [\MusicPDFPageRef{757} /XYZ 71 721 null] >>}
\MusicPDFWriteObject{15834}{<< /D [\MusicPDFPageRef{758} /XYZ 71 721 null] >>}
\MusicPDFWriteObject{15835}{<< /D [\MusicPDFPageRef{759} /XYZ 71 721 null] >>}
\MusicPDFWriteObject{15836}{<< /D [\MusicPDFPageRef{76} /XYZ 71 721 null] >>}
\MusicPDFWriteObject{15837}{<< /D [\MusicPDFPageRef{760} /XYZ 71 721 null] >>}
\MusicPDFWriteObject{15838}{<< /D [\MusicPDFPageRef{761} /XYZ 71 721 null] >>}
\MusicPDFWriteObject{15839}{<< /D [\MusicPDFPageRef{762} /XYZ 71 721 null] >>}
\MusicPDFWriteObject{15840}{<< /D [\MusicPDFPageRef{763} /XYZ 71 721 null] >>}
\MusicPDFWriteObject{15841}{<< /D [\MusicPDFPageRef{764} /XYZ 71 721 null] >>}
\MusicPDFWriteObject{15842}{<< /D [\MusicPDFPageRef{765} /XYZ 71 721 null] >>}
\MusicPDFWriteObject{15843}{<< /D [\MusicPDFPageRef{766} /XYZ 71 721 null] >>}
\MusicPDFWriteObject{15844}{<< /D [\MusicPDFPageRef{767} /XYZ 71 721 null] >>}
\MusicPDFWriteObject{15845}{<< /D [\MusicPDFPageRef{768} /XYZ 71 721 null] >>}
\MusicPDFWriteObject{15846}{<< /D [\MusicPDFPageRef{769} /XYZ 71 721 null] >>}
\MusicPDFWriteObject{15847}{<< /D [\MusicPDFPageRef{77} /XYZ 71 721 null] >>}
\MusicPDFWriteObject{15848}{<< /D [\MusicPDFPageRef{770} /XYZ 71 721 null] >>}
\MusicPDFWriteObject{15849}{<< /D [\MusicPDFPageRef{771} /XYZ 71 721 null] >>}
\MusicPDFWriteObject{15850}{<< /D [\MusicPDFPageRef{772} /XYZ 71 721 null] >>}
\MusicPDFWriteObject{15851}{<< /D [\MusicPDFPageRef{773} /XYZ 71 721 null] >>}
\MusicPDFWriteObject{15852}{<< /D [\MusicPDFPageRef{774} /XYZ 71 721 null] >>}
\MusicPDFWriteObject{15853}{<< /D [\MusicPDFPageRef{775} /XYZ 71 721 null] >>}
\MusicPDFWriteObject{15854}{<< /D [\MusicPDFPageRef{776} /XYZ 71 721 null] >>}
\MusicPDFWriteObject{15855}{<< /D [\MusicPDFPageRef{777} /XYZ 71 721 null] >>}
\MusicPDFWriteObject{15856}{<< /D [\MusicPDFPageRef{778} /XYZ 71 721 null] >>}
\MusicPDFWriteObject{15857}{<< /D [\MusicPDFPageRef{779} /XYZ 71 721 null] >>}
\MusicPDFWriteObject{15858}{<< /D [\MusicPDFPageRef{78} /XYZ 71 721 null] >>}
\MusicPDFWriteObject{15859}{<< /D [\MusicPDFPageRef{780} /XYZ 71 721 null] >>}
\MusicPDFWriteObject{15860}{<< /D [\MusicPDFPageRef{781} /XYZ 71 721 null] >>}
\MusicPDFWriteObject{15861}{<< /D [\MusicPDFPageRef{782} /XYZ 71 721 null] >>}
\MusicPDFWriteObject{15862}{<< /D [\MusicPDFPageRef{783} /XYZ 71 721 null] >>}
\MusicPDFWriteObject{15863}{<< /D [\MusicPDFPageRef{784} /XYZ 71 721 null] >>}
\MusicPDFWriteObject{15864}{<< /D [\MusicPDFPageRef{785} /XYZ 71 721 null] >>}
\MusicPDFWriteObject{15865}{<< /D [\MusicPDFPageRef{786} /XYZ 71 721 null] >>}
\MusicPDFWriteObject{15866}{<< /D [\MusicPDFPageRef{787} /XYZ 71 721 null] >>}
\MusicPDFWriteObject{15867}{<< /D [\MusicPDFPageRef{788} /XYZ 71 721 null] >>}
\MusicPDFWriteObject{15868}{<< /D [\MusicPDFPageRef{789} /XYZ 71 721 null] >>}
\MusicPDFWriteObject{15869}{<< /D [\MusicPDFPageRef{79} /XYZ 71 721 null] >>}
\MusicPDFWriteObject{15870}{<< /D [\MusicPDFPageRef{790} /XYZ 71 721 null] >>}
\MusicPDFWriteObject{15871}{<< /D [\MusicPDFPageRef{791} /XYZ 71 721 null] >>}
\MusicPDFWriteObject{15872}{<< /D [\MusicPDFPageRef{792} /XYZ 71 721 null] >>}
\MusicPDFWriteObject{15873}{<< /D [\MusicPDFPageRef{793} /XYZ 71 721 null] >>}
\MusicPDFWriteObject{15874}{<< /D [\MusicPDFPageRef{794} /XYZ 71 721 null] >>}
\MusicPDFWriteObject{15875}{<< /D [\MusicPDFPageRef{795} /XYZ 71 721 null] >>}
\MusicPDFWriteObject{15876}{<< /D [\MusicPDFPageRef{796} /XYZ 71 721 null] >>}
\MusicPDFWriteObject{15877}{<< /D [\MusicPDFPageRef{797} /XYZ 71 721 null] >>}
\MusicPDFWriteObject{15878}{<< /D [\MusicPDFPageRef{798} /XYZ 71 721 null] >>}
\MusicPDFWriteObject{15879}{<< /D [\MusicPDFPageRef{799} /XYZ 71 721 null] >>}
\MusicPDFWriteObject{15880}{<< /D [\MusicPDFPageRef{8} /XYZ 71 721 null] >>}
\MusicPDFWriteObject{15881}{<< /D [\MusicPDFPageRef{80} /XYZ 71 721 null] >>}
\MusicPDFWriteObject{15882}{<< /D [\MusicPDFPageRef{800} /XYZ 71 721 null] >>}
\MusicPDFWriteObject{15883}{<< /D [\MusicPDFPageRef{801} /XYZ 71 721 null] >>}
\MusicPDFWriteObject{15884}{<< /D [\MusicPDFPageRef{802} /XYZ 71 721 null] >>}
\MusicPDFWriteObject{15885}{<< /D [\MusicPDFPageRef{803} /XYZ 71 721 null] >>}
\MusicPDFWriteObject{15886}{<< /D [\MusicPDFPageRef{804} /XYZ 71 721 null] >>}
\MusicPDFWriteObject{15887}{<< /D [\MusicPDFPageRef{805} /XYZ 71 721 null] >>}
\MusicPDFWriteObject{15888}{<< /D [\MusicPDFPageRef{806} /XYZ 71 721 null] >>}
\MusicPDFWriteObject{15889}{<< /D [\MusicPDFPageRef{807} /XYZ 71 721 null] >>}
\MusicPDFWriteObject{15890}{<< /D [\MusicPDFPageRef{808} /XYZ 71 721 null] >>}
\MusicPDFWriteObject{15891}{<< /D [\MusicPDFPageRef{809} /XYZ 71 721 null] >>}
\MusicPDFWriteObject{15892}{<< /D [\MusicPDFPageRef{81} /XYZ 71 721 null] >>}
\MusicPDFWriteObject{15893}{<< /D [\MusicPDFPageRef{810} /XYZ 71 721 null] >>}
\MusicPDFWriteObject{15894}{<< /D [\MusicPDFPageRef{811} /XYZ 71 721 null] >>}
\MusicPDFWriteObject{15895}{<< /D [\MusicPDFPageRef{812} /XYZ 71 721 null] >>}
\MusicPDFWriteObject{15896}{<< /D [\MusicPDFPageRef{813} /XYZ 71 721 null] >>}
\MusicPDFWriteObject{15897}{<< /D [\MusicPDFPageRef{814} /XYZ 71 721 null] >>}
\MusicPDFWriteObject{15898}{<< /D [\MusicPDFPageRef{815} /XYZ 71 721 null] >>}
\MusicPDFWriteObject{15899}{<< /D [\MusicPDFPageRef{816} /XYZ 71 721 null] >>}
\MusicPDFWriteObject{15900}{<< /D [\MusicPDFPageRef{817} /XYZ 71 721 null] >>}
\MusicPDFWriteObject{15901}{<< /D [\MusicPDFPageRef{818} /XYZ 71 721 null] >>}
\MusicPDFWriteObject{15902}{<< /D [\MusicPDFPageRef{819} /XYZ 71 721 null] >>}
\MusicPDFWriteObject{15903}{<< /D [\MusicPDFPageRef{82} /XYZ 71 721 null] >>}
\MusicPDFWriteObject{15904}{<< /D [\MusicPDFPageRef{820} /XYZ 71 721 null] >>}
\MusicPDFWriteObject{15905}{<< /D [\MusicPDFPageRef{821} /XYZ 71 721 null] >>}
\MusicPDFWriteObject{15906}{<< /D [\MusicPDFPageRef{822} /XYZ 71 721 null] >>}
\MusicPDFWriteObject{15907}{<< /D [\MusicPDFPageRef{823} /XYZ 71 721 null] >>}
\MusicPDFWriteObject{15908}{<< /D [\MusicPDFPageRef{824} /XYZ 71 721 null] >>}
\MusicPDFWriteObject{15909}{<< /D [\MusicPDFPageRef{825} /XYZ 71 721 null] >>}
\MusicPDFWriteObject{15910}{<< /D [\MusicPDFPageRef{826} /XYZ 71 721 null] >>}
\MusicPDFWriteObject{15911}{<< /D [\MusicPDFPageRef{827} /XYZ 71 721 null] >>}
\MusicPDFWriteObject{15912}{<< /D [\MusicPDFPageRef{828} /XYZ 71 721 null] >>}
\MusicPDFWriteObject{15913}{<< /D [\MusicPDFPageRef{829} /XYZ 71 721 null] >>}
\MusicPDFWriteObject{15914}{<< /D [\MusicPDFPageRef{83} /XYZ 71 721 null] >>}
\MusicPDFWriteObject{15915}{<< /D [\MusicPDFPageRef{830} /XYZ 71 721 null] >>}
\MusicPDFWriteObject{15916}{<< /D [\MusicPDFPageRef{831} /XYZ 71 721 null] >>}
\MusicPDFWriteObject{15917}{<< /D [\MusicPDFPageRef{832} /XYZ 71 721 null] >>}
\MusicPDFWriteObject{15918}{<< /D [\MusicPDFPageRef{833} /XYZ 71 721 null] >>}
\MusicPDFWriteObject{15919}{<< /D [\MusicPDFPageRef{834} /XYZ 71 721 null] >>}
\MusicPDFWriteObject{15920}{<< /D [\MusicPDFPageRef{835} /XYZ 71 721 null] >>}
\MusicPDFWriteObject{15921}{<< /D [\MusicPDFPageRef{836} /XYZ 71 721 null] >>}
\MusicPDFWriteObject{15922}{<< /D [\MusicPDFPageRef{837} /XYZ 71 721 null] >>}
\MusicPDFWriteObject{15923}{<< /D [\MusicPDFPageRef{838} /XYZ 71 721 null] >>}
\MusicPDFWriteObject{15924}{<< /D [\MusicPDFPageRef{839} /XYZ 71 721 null] >>}
\MusicPDFWriteObject{15925}{<< /D [\MusicPDFPageRef{84} /XYZ 71 721 null] >>}
\MusicPDFWriteObject{15926}{<< /D [\MusicPDFPageRef{840} /XYZ 71 721 null] >>}
\MusicPDFWriteObject{15927}{<< /D [\MusicPDFPageRef{841} /XYZ 71 721 null] >>}
\MusicPDFWriteObject{15928}{<< /D [\MusicPDFPageRef{842} /XYZ 71 721 null] >>}
\MusicPDFWriteObject{15929}{<< /D [\MusicPDFPageRef{843} /XYZ 71 721 null] >>}
\MusicPDFWriteObject{15930}{<< /D [\MusicPDFPageRef{844} /XYZ 71 721 null] >>}
\MusicPDFWriteObject{15931}{<< /D [\MusicPDFPageRef{845} /XYZ 71 721 null] >>}
\MusicPDFWriteObject{15932}{<< /D [\MusicPDFPageRef{846} /XYZ 71 721 null] >>}
\MusicPDFWriteObject{15933}{<< /D [\MusicPDFPageRef{847} /XYZ 71 721 null] >>}
\MusicPDFWriteObject{15934}{<< /D [\MusicPDFPageRef{848} /XYZ 71 721 null] >>}
\MusicPDFWriteObject{15935}{<< /D [\MusicPDFPageRef{849} /XYZ 71 721 null] >>}
\MusicPDFWriteObject{15936}{<< /D [\MusicPDFPageRef{85} /XYZ 71 721 null] >>}
\MusicPDFWriteObject{15937}{<< /D [\MusicPDFPageRef{850} /XYZ 71 721 null] >>}
\MusicPDFWriteObject{15938}{<< /D [\MusicPDFPageRef{851} /XYZ 71 721 null] >>}
\MusicPDFWriteObject{15939}{<< /D [\MusicPDFPageRef{852} /XYZ 71 721 null] >>}
\MusicPDFWriteObject{15940}{<< /D [\MusicPDFPageRef{853} /XYZ 71 721 null] >>}
\MusicPDFWriteObject{15941}{<< /D [\MusicPDFPageRef{854} /XYZ 71 721 null] >>}
\MusicPDFWriteObject{15942}{<< /D [\MusicPDFPageRef{855} /XYZ 71 721 null] >>}
\MusicPDFWriteObject{15943}{<< /D [\MusicPDFPageRef{856} /XYZ 71 721 null] >>}
\MusicPDFWriteObject{15944}{<< /D [\MusicPDFPageRef{857} /XYZ 71 721 null] >>}
\MusicPDFWriteObject{15945}{<< /D [\MusicPDFPageRef{858} /XYZ 71 721 null] >>}
\MusicPDFWriteObject{15946}{<< /D [\MusicPDFPageRef{859} /XYZ 71 721 null] >>}
\MusicPDFWriteObject{15947}{<< /D [\MusicPDFPageRef{86} /XYZ 71 721 null] >>}
\MusicPDFWriteObject{15948}{<< /D [\MusicPDFPageRef{860} /XYZ 71 721 null] >>}
\MusicPDFWriteObject{15949}{<< /D [\MusicPDFPageRef{861} /XYZ 71 721 null] >>}
\MusicPDFWriteObject{15950}{<< /D [\MusicPDFPageRef{862} /XYZ 71 721 null] >>}
\MusicPDFWriteObject{15951}{<< /D [\MusicPDFPageRef{863} /XYZ 71 721 null] >>}
\MusicPDFWriteObject{15952}{<< /D [\MusicPDFPageRef{864} /XYZ 71 721 null] >>}
\MusicPDFWriteObject{15953}{<< /D [\MusicPDFPageRef{865} /XYZ 71 721 null] >>}
\MusicPDFWriteObject{15954}{<< /D [\MusicPDFPageRef{866} /XYZ 71 721 null] >>}
\MusicPDFWriteObject{15955}{<< /D [\MusicPDFPageRef{867} /XYZ 71 721 null] >>}
\MusicPDFWriteObject{15956}{<< /D [\MusicPDFPageRef{868} /XYZ 71 721 null] >>}
\MusicPDFWriteObject{15957}{<< /D [\MusicPDFPageRef{869} /XYZ 71 721 null] >>}
\MusicPDFWriteObject{15958}{<< /D [\MusicPDFPageRef{87} /XYZ 71 721 null] >>}
\MusicPDFWriteObject{15959}{<< /D [\MusicPDFPageRef{870} /XYZ 71 721 null] >>}
\MusicPDFWriteObject{15960}{<< /D [\MusicPDFPageRef{871} /XYZ 71 721 null] >>}
\MusicPDFWriteObject{15961}{<< /D [\MusicPDFPageRef{872} /XYZ 71 721 null] >>}
\MusicPDFWriteObject{15962}{<< /D [\MusicPDFPageRef{873} /XYZ 71 721 null] >>}
\MusicPDFWriteObject{15963}{<< /D [\MusicPDFPageRef{874} /XYZ 71 721 null] >>}
\MusicPDFWriteObject{15964}{<< /D [\MusicPDFPageRef{875} /XYZ 71 721 null] >>}
\MusicPDFWriteObject{15965}{<< /D [\MusicPDFPageRef{876} /XYZ 71 721 null] >>}
\MusicPDFWriteObject{15966}{<< /D [\MusicPDFPageRef{877} /XYZ 71 721 null] >>}
\MusicPDFWriteObject{15967}{<< /D [\MusicPDFPageRef{878} /XYZ 71 721 null] >>}
\MusicPDFWriteObject{15968}{<< /D [\MusicPDFPageRef{879} /XYZ 71 721 null] >>}
\MusicPDFWriteObject{15969}{<< /D [\MusicPDFPageRef{88} /XYZ 71 721 null] >>}
\MusicPDFWriteObject{15970}{<< /D [\MusicPDFPageRef{880} /XYZ 71 721 null] >>}
\MusicPDFWriteObject{15971}{<< /D [\MusicPDFPageRef{881} /XYZ 71 721 null] >>}
\MusicPDFWriteObject{15972}{<< /D [\MusicPDFPageRef{882} /XYZ 71 721 null] >>}
\MusicPDFWriteObject{15973}{<< /D [\MusicPDFPageRef{883} /XYZ 71 721 null] >>}
\MusicPDFWriteObject{15974}{<< /D [\MusicPDFPageRef{884} /XYZ 71 721 null] >>}
\MusicPDFWriteObject{15975}{<< /D [\MusicPDFPageRef{885} /XYZ 71 721 null] >>}
\MusicPDFWriteObject{15976}{<< /D [\MusicPDFPageRef{886} /XYZ 71 721 null] >>}
\MusicPDFWriteObject{15977}{<< /D [\MusicPDFPageRef{887} /XYZ 71 721 null] >>}
\MusicPDFWriteObject{15978}{<< /D [\MusicPDFPageRef{888} /XYZ 71 721 null] >>}
\MusicPDFWriteObject{15979}{<< /D [\MusicPDFPageRef{889} /XYZ 71 721 null] >>}
\MusicPDFWriteObject{15980}{<< /D [\MusicPDFPageRef{89} /XYZ 71 721 null] >>}
\MusicPDFWriteObject{15981}{<< /D [\MusicPDFPageRef{890} /XYZ 71 721 null] >>}
\MusicPDFWriteObject{15982}{<< /D [\MusicPDFPageRef{891} /XYZ 71 721 null] >>}
\MusicPDFWriteObject{15983}{<< /D [\MusicPDFPageRef{892} /XYZ 71 721 null] >>}
\MusicPDFWriteObject{15984}{<< /D [\MusicPDFPageRef{893} /XYZ 71 721 null] >>}
\MusicPDFWriteObject{15985}{<< /D [\MusicPDFPageRef{894} /XYZ 71 721 null] >>}
\MusicPDFWriteObject{15986}{<< /D [\MusicPDFPageRef{895} /XYZ 71 721 null] >>}
\MusicPDFWriteObject{15987}{<< /D [\MusicPDFPageRef{896} /XYZ 71 721 null] >>}
\MusicPDFWriteObject{15988}{<< /D [\MusicPDFPageRef{897} /XYZ 71 721 null] >>}
\MusicPDFWriteObject{15989}{<< /D [\MusicPDFPageRef{898} /XYZ 71 721 null] >>}
\MusicPDFWriteObject{15990}{<< /D [\MusicPDFPageRef{899} /XYZ 71 721 null] >>}
\MusicPDFWriteObject{15991}{<< /D [\MusicPDFPageRef{9} /XYZ 71 721 null] >>}
\MusicPDFWriteObject{15992}{<< /D [\MusicPDFPageRef{90} /XYZ 71 721 null] >>}
\MusicPDFWriteObject{15993}{<< /D [\MusicPDFPageRef{900} /XYZ 71 721 null] >>}
\MusicPDFWriteObject{15994}{<< /D [\MusicPDFPageRef{901} /XYZ 71 721 null] >>}
\MusicPDFWriteObject{15995}{<< /D [\MusicPDFPageRef{902} /XYZ 71 721 null] >>}
\MusicPDFWriteObject{15996}{<< /D [\MusicPDFPageRef{903} /XYZ 71 721 null] >>}
\MusicPDFWriteObject{15997}{<< /D [\MusicPDFPageRef{904} /XYZ 71 721 null] >>}
\MusicPDFWriteObject{15998}{<< /D [\MusicPDFPageRef{905} /XYZ 71 721 null] >>}
\MusicPDFWriteObject{15999}{<< /D [\MusicPDFPageRef{906} /XYZ 71 721 null] >>}
\MusicPDFWriteObject{16000}{<< /D [\MusicPDFPageRef{907} /XYZ 71 721 null] >>}
\MusicPDFWriteObject{16001}{<< /D [\MusicPDFPageRef{908} /XYZ 71 721 null] >>}
\MusicPDFWriteObject{16002}{<< /D [\MusicPDFPageRef{909} /XYZ 71 721 null] >>}
\MusicPDFWriteObject{16003}{<< /D [\MusicPDFPageRef{91} /XYZ 71 721 null] >>}
\MusicPDFWriteObject{16004}{<< /D [\MusicPDFPageRef{910} /XYZ 71 721 null] >>}
\MusicPDFWriteObject{16005}{<< /D [\MusicPDFPageRef{911} /XYZ 71 721 null] >>}
\MusicPDFWriteObject{16006}{<< /D [\MusicPDFPageRef{912} /XYZ 71 721 null] >>}
\MusicPDFWriteObject{16007}{<< /D [\MusicPDFPageRef{913} /XYZ 71 721 null] >>}
\MusicPDFWriteObject{16008}{<< /D [\MusicPDFPageRef{914} /XYZ 71 721 null] >>}
\MusicPDFWriteObject{16009}{<< /D [\MusicPDFPageRef{915} /XYZ 71 721 null] >>}
\MusicPDFWriteObject{16010}{<< /D [\MusicPDFPageRef{916} /XYZ 71 721 null] >>}
\MusicPDFWriteObject{16011}{<< /D [\MusicPDFPageRef{917} /XYZ 71 721 null] >>}
\MusicPDFWriteObject{16012}{<< /D [\MusicPDFPageRef{918} /XYZ 71 721 null] >>}
\MusicPDFWriteObject{16013}{<< /D [\MusicPDFPageRef{919} /XYZ 71 721 null] >>}
\MusicPDFWriteObject{16014}{<< /D [\MusicPDFPageRef{92} /XYZ 71 721 null] >>}
\MusicPDFWriteObject{16015}{<< /D [\MusicPDFPageRef{920} /XYZ 71 721 null] >>}
\MusicPDFWriteObject{16016}{<< /D [\MusicPDFPageRef{921} /XYZ 71 721 null] >>}
\MusicPDFWriteObject{16017}{<< /D [\MusicPDFPageRef{922} /XYZ 71 721 null] >>}
\MusicPDFWriteObject{16018}{<< /D [\MusicPDFPageRef{923} /XYZ 71 721 null] >>}
\MusicPDFWriteObject{16019}{<< /D [\MusicPDFPageRef{924} /XYZ 71 721 null] >>}
\MusicPDFWriteObject{16020}{<< /D [\MusicPDFPageRef{925} /XYZ 71 721 null] >>}
\MusicPDFWriteObject{16021}{<< /D [\MusicPDFPageRef{926} /XYZ 71 721 null] >>}
\MusicPDFWriteObject{16022}{<< /D [\MusicPDFPageRef{927} /XYZ 71 721 null] >>}
\MusicPDFWriteObject{16023}{<< /D [\MusicPDFPageRef{928} /XYZ 71 721 null] >>}
\MusicPDFWriteObject{16024}{<< /D [\MusicPDFPageRef{929} /XYZ 71 721 null] >>}
\MusicPDFWriteObject{16025}{<< /D [\MusicPDFPageRef{93} /XYZ 71 721 null] >>}
\MusicPDFWriteObject{16026}{<< /D [\MusicPDFPageRef{930} /XYZ 71 721 null] >>}
\MusicPDFWriteObject{16027}{<< /D [\MusicPDFPageRef{931} /XYZ 71 721 null] >>}
\MusicPDFWriteObject{16028}{<< /D [\MusicPDFPageRef{932} /XYZ 71 721 null] >>}
\MusicPDFWriteObject{16029}{<< /D [\MusicPDFPageRef{933} /XYZ 71 721 null] >>}
\MusicPDFWriteObject{16030}{<< /D [\MusicPDFPageRef{934} /XYZ 71 721 null] >>}
\MusicPDFWriteObject{16031}{<< /D [\MusicPDFPageRef{935} /XYZ 71 721 null] >>}
\MusicPDFWriteObject{16032}{<< /D [\MusicPDFPageRef{936} /XYZ 71 721 null] >>}
\MusicPDFWriteObject{16033}{<< /D [\MusicPDFPageRef{937} /XYZ 71 721 null] >>}
\MusicPDFWriteObject{16034}{<< /D [\MusicPDFPageRef{938} /XYZ 71 721 null] >>}
\MusicPDFWriteObject{16035}{<< /D [\MusicPDFPageRef{939} /XYZ 71 721 null] >>}
\MusicPDFWriteObject{16036}{<< /D [\MusicPDFPageRef{94} /XYZ 71 721 null] >>}
\MusicPDFWriteObject{16037}{<< /D [\MusicPDFPageRef{940} /XYZ 71 721 null] >>}
\MusicPDFWriteObject{16038}{<< /D [\MusicPDFPageRef{95} /XYZ 71 721 null] >>}
\MusicPDFWriteObject{16039}{<< /D [\MusicPDFPageRef{96} /XYZ 71 721 null] >>}
\MusicPDFWriteObject{16040}{<< /D [\MusicPDFPageRef{97} /XYZ 71 721 null] >>}
\MusicPDFWriteObject{16041}{<< /D [\MusicPDFPageRef{98} /XYZ 71 721 null] >>}
\MusicPDFWriteObject{16042}{<< /D [\MusicPDFPageRef{99} /XYZ 71 721 null] >>}
\MusicPDFWriteObject{16043}{<< /D [\MusicPDFPageRef{9} /XYZ 72 698.082 null] >>}
\MusicPDFWriteObject{16044}{<< /D [\MusicPDFPageRef{11} /XYZ 72 587.696 null] >>}
\MusicPDFWriteObject{16045}{<< /D [\MusicPDFPageRef{12} /XYZ 72 180.968 null] >>}
\MusicPDFWriteObject{16046}{<< /D [\MusicPDFPageRef{21} /XYZ 72 573.802 null] >>}
\MusicPDFWriteObject{16047}{<< /D [\MusicPDFPageRef{27} /XYZ 72 182.163 null] >>}
\MusicPDFWriteObject{16048}{<< /D [\MusicPDFPageRef{32} /XYZ 72 366.289 null] >>}
\MusicPDFWriteObject{16049}{<< /D [\MusicPDFPageRef{33} /XYZ 72 513.326 null] >>}
\MusicPDFWriteObject{16050}{<< /D [\MusicPDFPageRef{33} /XYZ 72 454.786 null] >>}
\MusicPDFWriteObject{16051}{<< /D [\MusicPDFPageRef{38} /XYZ 72 296.484 null] >>}
\MusicPDFWriteObject{16052}{<< /D [\MusicPDFPageRef{434} /XYZ 72 698.082 null] >>}
\MusicPDFWriteObject{16053}{<< /D [\MusicPDFPageRef{438} /XYZ 72 321.415 null] >>}
\MusicPDFWriteObject{16054}{<< /D [\MusicPDFPageRef{447} /XYZ 72 333.515 null] >>}
\MusicPDFWriteObject{16055}{<< /D [\MusicPDFPageRef{455} /XYZ 72 350.127 null] >>}
\MusicPDFWriteObject{16056}{<< /D [\MusicPDFPageRef{475} /XYZ 72 421.519 null] >>}
\MusicPDFWriteObject{16057}{<< /D [\MusicPDFPageRef{483} /XYZ 72 212.249 null] >>}
\MusicPDFWriteObject{16058}{<< /D [\MusicPDFPageRef{491} /XYZ 72 424.825 null] >>}
\MusicPDFWriteObject{16059}{<< /D [\MusicPDFPageRef{498} /XYZ 72 698.082 null] >>}
\MusicPDFWriteObject{16060}{<< /D [\MusicPDFPageRef{504} /XYZ 72 580.902 null] >>}
\MusicPDFWriteObject{16061}{<< /D [\MusicPDFPageRef{514} /XYZ 72 698.082 null] >>}
\MusicPDFWriteObject{16062}{<< /D [\MusicPDFPageRef{514} /XYZ 72 659.403 null] >>}
\MusicPDFWriteObject{16063}{<< /D [\MusicPDFPageRef{519} /XYZ 72 487.284 null] >>}
\MusicPDFWriteObject{16064}{<< /D [\MusicPDFPageRef{522} /XYZ 72 664.664 null] >>}
\MusicPDFWriteObject{16065}{<< /D [\MusicPDFPageRef{522} /XYZ 72 645.535 null] >>}
\MusicPDFWriteObject{16066}{<< /D [\MusicPDFPageRef{522} /XYZ 72 626.407 null] >>}
\MusicPDFWriteObject{16067}{<< /D [\MusicPDFPageRef{524} /XYZ 72 698.082 null] >>}
\MusicPDFWriteObject{16068}{<< /D [\MusicPDFPageRef{534} /XYZ 72 698.082 null] >>}
\MusicPDFWriteObject{16069}{<< /D [\MusicPDFPageRef{536} /XYZ 72 611.101 null] >>}
\MusicPDFWriteObject{16070}{<< /D [\MusicPDFPageRef{59} /XYZ 72 698.082 null] >>}
\MusicPDFWriteObject{16071}{<< /D [\MusicPDFPageRef{74} /XYZ 72 168.114 null] >>}
\MusicPDFWriteObject{16072}{<< /D [\MusicPDFPageRef{98} /XYZ 72 141.919 null] >>}
\MusicPDFWriteObject{16073}{<< /D [\MusicPDFPageRef{106} /XYZ 72 275.814 null] >>}
\MusicPDFWriteObject{16074}{<< /D [\MusicPDFPageRef{118} /XYZ 72 200.417 null] >>}
\MusicPDFWriteObject{16075}{<< /D [\MusicPDFPageRef{122} /XYZ 72 650.33 null] >>}
\MusicPDFWriteObject{16076}{<< /D [\MusicPDFPageRef{128} /XYZ 72 338.269 null] >>}
\MusicPDFWriteObject{16077}{<< /D [\MusicPDFPageRef{141} /XYZ 72 142.173 null] >>}
\MusicPDFWriteObject{16078}{<< /D [\MusicPDFPageRef{155} /XYZ 72 131.512 null] >>}
\MusicPDFWriteObject{16079}{<< /D [\MusicPDFPageRef{161} /XYZ 72 242.235 null] >>}
\MusicPDFWriteObject{16080}{<< /D [\MusicPDFPageRef{174} /XYZ 72 648.497 null] >>}
\MusicPDFWriteObject{16081}{<< /D [\MusicPDFPageRef{188} /XYZ 72 269.183 null] >>}
\MusicPDFWriteObject{16082}{<< /D [\MusicPDFPageRef{197} /XYZ 72 506.108 null] >>}
\MusicPDFWriteObject{16083}{<< /D [\MusicPDFPageRef{209} /XYZ 72 559.972 null] >>}
\MusicPDFWriteObject{16084}{<< /D [\MusicPDFPageRef{221} /XYZ 72 271.406 null] >>}
\MusicPDFWriteObject{16085}{<< /D [\MusicPDFPageRef{230} /XYZ 72 679.751 null] >>}
\MusicPDFWriteObject{16086}{<< /D [\MusicPDFPageRef{240} /XYZ 72 698.082 null] >>}
\MusicPDFWriteObject{16087}{<< /D [\MusicPDFPageRef{265} /XYZ 72 393.919 null] >>}
\MusicPDFWriteObject{16088}{<< /D [\MusicPDFPageRef{277} /XYZ 72 359.306 null] >>}
\MusicPDFWriteObject{16089}{<< /D [\MusicPDFPageRef{287} /XYZ 72 445.885 null] >>}
\MusicPDFWriteObject{16090}{<< /D [\MusicPDFPageRef{295} /XYZ 72 504.807 null] >>}
\MusicPDFWriteObject{16091}{<< /D [\MusicPDFPageRef{302} /XYZ 72 320.458 null] >>}
\MusicPDFWriteObject{16092}{<< /D [\MusicPDFPageRef{312} /XYZ 72 427.336 null] >>}
\MusicPDFWriteObject{16093}{<< /D [\MusicPDFPageRef{332} /XYZ 72 516.768 null] >>}
\MusicPDFWriteObject{16094}{<< /D [\MusicPDFPageRef{340} /XYZ 72 432.047 null] >>}
\MusicPDFWriteObject{16095}{<< /D [\MusicPDFPageRef{366} /XYZ 72 602.772 null] >>}
\MusicPDFWriteObject{16096}{<< /D [\MusicPDFPageRef{374} /XYZ 72 642.594 null] >>}
\MusicPDFWriteObject{16097}{<< /D [\MusicPDFPageRef{381} /XYZ 72 277.802 null] >>}
\MusicPDFWriteObject{16098}{<< /D [\MusicPDFPageRef{384} /XYZ 72 383.019 null] >>}
\MusicPDFWriteObject{16099}{<< /D [\MusicPDFPageRef{387} /XYZ 72 411.948 null] >>}
\MusicPDFWriteObject{16100}{<< /D [\MusicPDFPageRef{389} /XYZ 72 588.576 null] >>}
\MusicPDFWriteObject{16101}{<< /D [\MusicPDFPageRef{394} /XYZ 72 195.419 null] >>}
\MusicPDFWriteObject{16102}{<< /D [\MusicPDFPageRef{395} /XYZ 72 182.178 null] >>}
\MusicPDFWriteObject{16103}{<< /D [\MusicPDFPageRef{403} /XYZ 72 472.623 null] >>}
\MusicPDFWriteObject{16104}{<< /D [\MusicPDFPageRef{419} /XYZ 72 613.957 null] >>}
\MusicPDFWriteObject{16105}{<< /D [\MusicPDFPageRef{539} /XYZ 72 698.082 null] >>}
\MusicPDFWriteObject{16106}{<< /D [\MusicPDFPageRef{545} /XYZ 72 587.83 null] >>}
\MusicPDFWriteObject{16107}{<< /D [\MusicPDFPageRef{555} /XYZ 72 298.142 null] >>}
\MusicPDFWriteObject{16108}{<< /D [\MusicPDFPageRef{562} /XYZ 72 444.378 null] >>}
\MusicPDFWriteObject{16109}{<< /D [\MusicPDFPageRef{566} /XYZ 72 698.082 null] >>}
\MusicPDFWriteObject{16110}{<< /D [\MusicPDFPageRef{588} /XYZ 72 324.224 null] >>}
\MusicPDFWriteObject{16111}{<< /D [\MusicPDFPageRef{591} /XYZ 72 173.841 null] >>}
\MusicPDFWriteObject{16112}{<< /D [\MusicPDFPageRef{604} /XYZ 72 612.371 null] >>}
\MusicPDFWriteObject{16113}{<< /D [\MusicPDFPageRef{606} /XYZ 72 285.933 null] >>}
\MusicPDFWriteObject{16114}{<< /D [\MusicPDFPageRef{626} /XYZ 72 570.787 null] >>}
\MusicPDFWriteObject{16115}{<< /D [\MusicPDFPageRef{636} /XYZ 72 698.082 null] >>}
\MusicPDFWriteObject{16116}{<< /D [\MusicPDFPageRef{643} /XYZ 72 314.739 null] >>}
\MusicPDFWriteObject{16117}{<< /D [\MusicPDFPageRef{657} /XYZ 72 698.082 null] >>}
\MusicPDFWriteObject{16118}{<< /D [\MusicPDFPageRef{670} /XYZ 72 539.178 null] >>}
\MusicPDFWriteObject{16119}{<< /D [\MusicPDFPageRef{693} /XYZ 72 698.082 null] >>}
\MusicPDFWriteObject{16120}{<< /D [\MusicPDFPageRef{708} /XYZ 72 207.927 null] >>}
\MusicPDFWriteObject{16121}{<< /D [\MusicPDFPageRef{719} /XYZ 72 274.234 null] >>}
\MusicPDFWriteObject{16122}{<< /D [\MusicPDFPageRef{745} /XYZ 72 698.082 null] >>}
\MusicPDFWriteObject{16123}{<< /D [\MusicPDFPageRef{792} /XYZ 72 596.236 null] >>}
\MusicPDFWriteObject{16124}{<< /D [\MusicPDFPageRef{797} /XYZ 72 342.726 null] >>}
\MusicPDFWriteObject{16125}{<< /D [\MusicPDFPageRef{825} /XYZ 72 420.829 null] >>}
\MusicPDFWriteObject{16126}{<< /D [\MusicPDFPageRef{846} /XYZ 72 613.177 null] >>}
\MusicPDFWriteObject{16127}{<< /D [\MusicPDFPageRef{868} /XYZ 72 471.705 null] >>}
\MusicPDFWriteObject{16128}{<< /D [\MusicPDFPageRef{883} /XYZ 72 507.591 null] >>}
\MusicPDFWriteObject{16129}{<< /D [\MusicPDFPageRef{889} /XYZ 72 563.712 null] >>}
\MusicPDFWriteObject{16130}{<< /D [\MusicPDFPageRef{920} /XYZ 72 296.212 null] >>}
\MusicPDFWriteObject{16131}{<< /D [\MusicPDFPageRef{11} /XYZ 72 566.631 null] >>}
\MusicPDFWriteObject{16132}{<< /D [\MusicPDFPageRef{11} /XYZ 72 536.249 null] >>}
\MusicPDFWriteObject{16133}{<< /D [\MusicPDFPageRef{11} /XYZ 72 225.774 null] >>}
\MusicPDFWriteObject{16134}{<< /D [\MusicPDFPageRef{12} /XYZ 72 558.712 null] >>}
\MusicPDFWriteObject{16135}{<< /D [\MusicPDFPageRef{13} /XYZ 72 698.082 null] >>}
\MusicPDFWriteObject{16136}{<< /D [\MusicPDFPageRef{19} /XYZ 72 457.325 null] >>}
\MusicPDFWriteObject{16137}{<< /D [\MusicPDFPageRef{20} /XYZ 72 655.498 null] >>}
\MusicPDFWriteObject{16138}{<< /D [\MusicPDFPageRef{15} /XYZ 72 346.984 null] >>}
\MusicPDFWriteObject{16139}{<< /D [\MusicPDFPageRef{16} /XYZ 72 698.082 null] >>}
\MusicPDFWriteObject{16140}{<< /D [\MusicPDFPageRef{16} /XYZ 72 332.077 null] >>}
\MusicPDFWriteObject{16141}{<< /D [\MusicPDFPageRef{16} /XYZ 72 308.676 null] >>}
\MusicPDFWriteObject{16142}{<< /D [\MusicPDFPageRef{16} /XYZ 72 280.847 null] >>}
\MusicPDFWriteObject{16143}{<< /D [\MusicPDFPageRef{16} /XYZ 72 254.955 null] >>}
\MusicPDFWriteObject{16144}{<< /D [\MusicPDFPageRef{16} /XYZ 72 150.159 null] >>}
\MusicPDFWriteObject{16145}{<< /D [\MusicPDFPageRef{17} /XYZ 72 116.062 null] >>}
\MusicPDFWriteObject{16146}{<< /D [\MusicPDFPageRef{21} /XYZ 72 369.306 null] >>}
\MusicPDFWriteObject{16147}{<< /D [\MusicPDFPageRef{25} /XYZ 72 698.082 null] >>}
\MusicPDFWriteObject{16148}{<< /D [\MusicPDFPageRef{27} /XYZ 72 353.961 null] >>}
\MusicPDFWriteObject{16149}{<< /D [\MusicPDFPageRef{27} /XYZ 72 264.451 null] >>}
\MusicPDFWriteObject{16150}{<< /D [\MusicPDFPageRef{28} /XYZ 72 412.589 null] >>}
\MusicPDFWriteObject{16151}{<< /D [\MusicPDFPageRef{32} /XYZ 72 529.054 null] >>}
\MusicPDFWriteObject{16152}{<< /D [\MusicPDFPageRef{32} /XYZ 72 438.006 null] >>}
\MusicPDFWriteObject{16153}{<< /D [\MusicPDFPageRef{28} /XYZ 72 227.716 null] >>}
\MusicPDFWriteObject{16154}{<< /D [\MusicPDFPageRef{29} /XYZ 72 649.32 null] >>}
\MusicPDFWriteObject{16155}{<< /D [\MusicPDFPageRef{29} /XYZ 72 296.402 null] >>}
\MusicPDFWriteObject{16156}{<< /D [\MusicPDFPageRef{29} /XYZ 72 167.976 null] >>}
\MusicPDFWriteObject{16157}{<< /D [\MusicPDFPageRef{30} /XYZ 72 549.468 null] >>}
\MusicPDFWriteObject{16158}{<< /D [\MusicPDFPageRef{30} /XYZ 72 279.326 null] >>}
\MusicPDFWriteObject{16159}{<< /D [\MusicPDFPageRef{31} /XYZ 72 509.554 null] >>}
\MusicPDFWriteObject{16160}{<< /D [\MusicPDFPageRef{31} /XYZ 72 151.836 null] >>}
\MusicPDFWriteObject{16161}{<< /D [\MusicPDFPageRef{32} /XYZ 72 347.161 null] >>}
\MusicPDFWriteObject{16162}{<< /D [\MusicPDFPageRef{33} /XYZ 72 616.099 null] >>}
\MusicPDFWriteObject{16163}{<< /D [\MusicPDFPageRef{33} /XYZ 72 494.198 null] >>}
\MusicPDFWriteObject{16164}{<< /D [\MusicPDFPageRef{33} /XYZ 72 476.718 null] >>}
\MusicPDFWriteObject{16165}{<< /D [\MusicPDFPageRef{33} /XYZ 72 433.721 null] >>}
\MusicPDFWriteObject{16166}{<< /D [\MusicPDFPageRef{35} /XYZ 72 399.481 null] >>}
\MusicPDFWriteObject{16167}{<< /D [\MusicPDFPageRef{37} /XYZ 72 654.504 null] >>}
\MusicPDFWriteObject{16168}{<< /D [\MusicPDFPageRef{38} /XYZ 72 278.269 null] >>}
\MusicPDFWriteObject{16169}{<< /D [\MusicPDFPageRef{42} /XYZ 72 546.647 null] >>}
\MusicPDFWriteObject{16170}{<< /D [\MusicPDFPageRef{44} /XYZ 72 140.76 null] >>}
\MusicPDFWriteObject{16171}{<< /D [\MusicPDFPageRef{53} /XYZ 72 698.082 null] >>}
\MusicPDFWriteObject{16172}{<< /D [\MusicPDFPageRef{54} /XYZ 72 209.364 null] >>}
\MusicPDFWriteObject{16173}{<< /D [\MusicPDFPageRef{434} /XYZ 72 619.785 null] >>}
\MusicPDFWriteObject{16174}{<< /D [\MusicPDFPageRef{437} /XYZ 72 449.22 null] >>}
\MusicPDFWriteObject{16175}{<< /D [\MusicPDFPageRef{439} /XYZ 72 666.171 null] >>}
\MusicPDFWriteObject{16176}{<< /D [\MusicPDFPageRef{445} /XYZ 72 698.082 null] >>}
\MusicPDFWriteObject{16177}{<< /D [\MusicPDFPageRef{448} /XYZ 72 497.867 null] >>}
\MusicPDFWriteObject{16178}{<< /D [\MusicPDFPageRef{453} /XYZ 72 315.629 null] >>}
\MusicPDFWriteObject{16179}{<< /D [\MusicPDFPageRef{456} /XYZ 72 625.459 null] >>}
\MusicPDFWriteObject{16180}{<< /D [\MusicPDFPageRef{458} /XYZ 72 262.207 null] >>}
\MusicPDFWriteObject{16181}{<< /D [\MusicPDFPageRef{464} /XYZ 72 167.131 null] >>}
\MusicPDFWriteObject{16182}{<< /D [\MusicPDFPageRef{475} /XYZ 72 315.525 null] >>}
\MusicPDFWriteObject{16183}{<< /D [\MusicPDFPageRef{480} /XYZ 72 698.082 null] >>}
\MusicPDFWriteObject{16184}{<< /D [\MusicPDFPageRef{487} /XYZ 72 464.088 null] >>}
\MusicPDFWriteObject{16185}{<< /D [\MusicPDFPageRef{489} /XYZ 72 333.95 null] >>}
\MusicPDFWriteObject{16186}{<< /D [\MusicPDFPageRef{491} /XYZ 72 408.547 null] >>}
\MusicPDFWriteObject{16187}{<< /D [\MusicPDFPageRef{495} /XYZ 72 634.1 null] >>}
\MusicPDFWriteObject{16188}{<< /D [\MusicPDFPageRef{498} /XYZ 72 680.205 null] >>}
\MusicPDFWriteObject{16189}{<< /D [\MusicPDFPageRef{504} /XYZ 72 275.046 null] >>}
\MusicPDFWriteObject{16190}{<< /D [\MusicPDFPageRef{507} /XYZ 72 168.234 null] >>}
\MusicPDFWriteObject{16191}{<< /D [\MusicPDFPageRef{514} /XYZ 72 538.898 null] >>}
\MusicPDFWriteObject{16192}{<< /D [\MusicPDFPageRef{519} /XYZ 72 381.838 null] >>}
\MusicPDFWriteObject{16193}{<< /D [\MusicPDFPageRef{522} /XYZ 72 563.785 null] >>}
\MusicPDFWriteObject{16194}{<< /D [\MusicPDFPageRef{524} /XYZ 72 508 null] >>}
\MusicPDFWriteObject{16195}{<< /D [\MusicPDFPageRef{528} /XYZ 72 281.809 null] >>}
\MusicPDFWriteObject{16196}{<< /D [\MusicPDFPageRef{59} /XYZ 72 423.792 null] >>}
\MusicPDFWriteObject{16197}{<< /D [\MusicPDFPageRef{63} /XYZ 72 643.966 null] >>}
\MusicPDFWriteObject{16198}{<< /D [\MusicPDFPageRef{67} /XYZ 72 489.253 null] >>}
\MusicPDFWriteObject{16199}{<< /D [\MusicPDFPageRef{70} /XYZ 72 110.447 null] >>}
\MusicPDFWriteObject{16200}{<< /D [\MusicPDFPageRef{75} /XYZ 72 653.904 null] >>}
\MusicPDFWriteObject{16201}{<< /D [\MusicPDFPageRef{77} /XYZ 72 622.85 null] >>}
\MusicPDFWriteObject{16202}{<< /D [\MusicPDFPageRef{81} /XYZ 72 289.775 null] >>}
\MusicPDFWriteObject{16203}{<< /D [\MusicPDFPageRef{86} /XYZ 72 198.001 null] >>}
\MusicPDFWriteObject{16204}{<< /D [\MusicPDFPageRef{95} /XYZ 72 611.17 null] >>}
\MusicPDFWriteObject{16205}{<< /D [\MusicPDFPageRef{99} /XYZ 72 280.771 null] >>}
\MusicPDFWriteObject{16206}{<< /D [\MusicPDFPageRef{105} /XYZ 72 222.052 null] >>}
\MusicPDFWriteObject{16207}{<< /D [\MusicPDFPageRef{106} /XYZ 72 125.534 null] >>}
\MusicPDFWriteObject{16208}{<< /D [\MusicPDFPageRef{110} /XYZ 72 225.835 null] >>}
\MusicPDFWriteObject{16209}{<< /D [\MusicPDFPageRef{112} /XYZ 72 109.533 null] >>}
\MusicPDFWriteObject{16210}{<< /D [\MusicPDFPageRef{118} /XYZ 72 181.289 null] >>}
\MusicPDFWriteObject{16211}{<< /D [\MusicPDFPageRef{120} /XYZ 72 698.082 null] >>}
\MusicPDFWriteObject{16212}{<< /D [\MusicPDFPageRef{122} /XYZ 72 339.025 null] >>}
\MusicPDFWriteObject{16213}{<< /D [\MusicPDFPageRef{124} /XYZ 72 404.049 null] >>}
\MusicPDFWriteObject{16214}{<< /D [\MusicPDFPageRef{128} /XYZ 72 281.407 null] >>}
\MusicPDFWriteObject{16215}{<< /D [\MusicPDFPageRef{137} /XYZ 72 423.539 null] >>}
\MusicPDFWriteObject{16216}{<< /D [\MusicPDFPageRef{142} /XYZ 72 366.897 null] >>}
\MusicPDFWriteObject{16217}{<< /D [\MusicPDFPageRef{144} /XYZ 72 113.91 null] >>}
\MusicPDFWriteObject{16218}{<< /D [\MusicPDFPageRef{151} /XYZ 72 353.382 null] >>}
\MusicPDFWriteObject{16219}{<< /D [\MusicPDFPageRef{156} /XYZ 72 698.082 null] >>}
\MusicPDFWriteObject{16220}{<< /D [\MusicPDFPageRef{158} /XYZ 72 504.12 null] >>}
\MusicPDFWriteObject{16221}{<< /D [\MusicPDFPageRef{161} /XYZ 72 110.447 null] >>}
\MusicPDFWriteObject{16222}{<< /D [\MusicPDFPageRef{170} /XYZ 72 440.069 null] >>}
\MusicPDFWriteObject{16223}{<< /D [\MusicPDFPageRef{174} /XYZ 72 301.659 null] >>}
\MusicPDFWriteObject{16224}{<< /D [\MusicPDFPageRef{177} /XYZ 72 135.835 null] >>}
\MusicPDFWriteObject{16225}{<< /D [\MusicPDFPageRef{180} /XYZ 72 411.146 null] >>}
\MusicPDFWriteObject{16226}{<< /D [\MusicPDFPageRef{188} /XYZ 72 149.899 null] >>}
\MusicPDFWriteObject{16227}{<< /D [\MusicPDFPageRef{191} /XYZ 72 212.491 null] >>}
\MusicPDFWriteObject{16228}{<< /D [\MusicPDFPageRef{197} /XYZ 72 302.377 null] >>}
\MusicPDFWriteObject{16229}{<< /D [\MusicPDFPageRef{205} /XYZ 72 655.338 null] >>}
\MusicPDFWriteObject{16230}{<< /D [\MusicPDFPageRef{208} /XYZ 72 608.109 null] >>}
\MusicPDFWriteObject{16231}{<< /D [\MusicPDFPageRef{209} /XYZ 72 540.844 null] >>}
\MusicPDFWriteObject{16232}{<< /D [\MusicPDFPageRef{218} /XYZ 72 464.891 null] >>}
\MusicPDFWriteObject{16233}{<< /D [\MusicPDFPageRef{219} /XYZ 72 297.611 null] >>}
\MusicPDFWriteObject{16234}{<< /D [\MusicPDFPageRef{222} /XYZ 72 291.267 null] >>}
\MusicPDFWriteObject{16235}{<< /D [\MusicPDFPageRef{224} /XYZ 72 475.301 null] >>}
\MusicPDFWriteObject{16236}{<< /D [\MusicPDFPageRef{226} /XYZ 72 397.985 null] >>}
\MusicPDFWriteObject{16237}{<< /D [\MusicPDFPageRef{230} /XYZ 72 605.846 null] >>}
\MusicPDFWriteObject{16238}{<< /D [\MusicPDFPageRef{232} /XYZ 72 614.049 null] >>}
\MusicPDFWriteObject{16239}{<< /D [\MusicPDFPageRef{233} /XYZ 72 667.054 null] >>}
\MusicPDFWriteObject{16240}{<< /D [\MusicPDFPageRef{238} /XYZ 72 234.805 null] >>}
\MusicPDFWriteObject{16241}{<< /D [\MusicPDFPageRef{240} /XYZ 72 603.62 null] >>}
\MusicPDFWriteObject{16242}{<< /D [\MusicPDFPageRef{243} /XYZ 72 698.082 null] >>}
\MusicPDFWriteObject{16243}{<< /D [\MusicPDFPageRef{249} /XYZ 72 470.372 null] >>}
\MusicPDFWriteObject{16244}{<< /D [\MusicPDFPageRef{259} /XYZ 72 462.641 null] >>}
\MusicPDFWriteObject{16245}{<< /D [\MusicPDFPageRef{265} /XYZ 72 311.299 null] >>}
\MusicPDFWriteObject{16246}{<< /D [\MusicPDFPageRef{274} /XYZ 72 346.854 null] >>}
\MusicPDFWriteObject{16247}{<< /D [\MusicPDFPageRef{277} /XYZ 72 275.237 null] >>}
\MusicPDFWriteObject{16248}{<< /D [\MusicPDFPageRef{280} /XYZ 72 445.235 null] >>}
\MusicPDFWriteObject{16249}{<< /D [\MusicPDFPageRef{281} /XYZ 72 544.709 null] >>}
\MusicPDFWriteObject{16250}{<< /D [\MusicPDFPageRef{287} /XYZ 72 378.035 null] >>}
\MusicPDFWriteObject{16251}{<< /D [\MusicPDFPageRef{288} /XYZ 72 587.626 null] >>}
\MusicPDFWriteObject{16252}{<< /D [\MusicPDFPageRef{295} /XYZ 72 326.005 null] >>}
\MusicPDFWriteObject{16253}{<< /D [\MusicPDFPageRef{298} /XYZ 72 484.081 null] >>}
\MusicPDFWriteObject{16254}{<< /D [\MusicPDFPageRef{302} /XYZ 72 227.433 null] >>}
\MusicPDFWriteObject{16255}{<< /D [\MusicPDFPageRef{303} /XYZ 72 146.796 null] >>}
\MusicPDFWriteObject{16256}{<< /D [\MusicPDFPageRef{312} /XYZ 72 274.923 null] >>}
\MusicPDFWriteObject{16257}{<< /D [\MusicPDFPageRef{315} /XYZ 72 123.044 null] >>}
\MusicPDFWriteObject{16258}{<< /D [\MusicPDFPageRef{323} /XYZ 72 453.113 null] >>}
\MusicPDFWriteObject{16259}{<< /D [\MusicPDFPageRef{326} /XYZ 72 489.814 null] >>}
\MusicPDFWriteObject{16260}{<< /D [\MusicPDFPageRef{332} /XYZ 72 141.258 null] >>}
\MusicPDFWriteObject{16261}{<< /D [\MusicPDFPageRef{338} /XYZ 72 209.994 null] >>}
\MusicPDFWriteObject{16262}{<< /D [\MusicPDFPageRef{340} /XYZ 72 340.182 null] >>}
\MusicPDFWriteObject{16263}{<< /D [\MusicPDFPageRef{350} /XYZ 72 162.71 null] >>}
\MusicPDFWriteObject{16264}{<< /D [\MusicPDFPageRef{357} /XYZ 72 546.006 null] >>}
\MusicPDFWriteObject{16265}{<< /D [\MusicPDFPageRef{364} /XYZ 72 617.949 null] >>}
\MusicPDFWriteObject{16266}{<< /D [\MusicPDFPageRef{366} /XYZ 72 475.56 null] >>}
\MusicPDFWriteObject{16267}{<< /D [\MusicPDFPageRef{367} /XYZ 72 403.552 null] >>}
\MusicPDFWriteObject{16268}{<< /D [\MusicPDFPageRef{367} /XYZ 72 123.044 null] >>}
\MusicPDFWriteObject{16269}{<< /D [\MusicPDFPageRef{396} /XYZ 72 635.48 null] >>}
\MusicPDFWriteObject{16270}{<< /D [\MusicPDFPageRef{399} /XYZ 72 429.41 null] >>}
\MusicPDFWriteObject{16271}{<< /D [\MusicPDFPageRef{404} /XYZ 72 634.011 null] >>}
\MusicPDFWriteObject{16272}{<< /D [\MusicPDFPageRef{411} /XYZ 72 500.913 null] >>}
\MusicPDFWriteObject{16273}{<< /D [\MusicPDFPageRef{419} /XYZ 72 530.516 null] >>}
\MusicPDFWriteObject{16274}{<< /D [\MusicPDFPageRef{422} /XYZ 72 577.569 null] >>}
\MusicPDFWriteObject{16275}{<< /D [\MusicPDFPageRef{426} /XYZ 72 408.629 null] >>}
\MusicPDFWriteObject{16276}{<< /D [\MusicPDFPageRef{539} /XYZ 72 594.354 null] >>}
\MusicPDFWriteObject{16277}{<< /D [\MusicPDFPageRef{540} /XYZ 72 327.707 null] >>}
\MusicPDFWriteObject{16278}{<< /D [\MusicPDFPageRef{541} /XYZ 72 605.542 null] >>}
\MusicPDFWriteObject{16279}{<< /D [\MusicPDFPageRef{542} /XYZ 72 412.891 null] >>}
\MusicPDFWriteObject{16280}{<< /D [\MusicPDFPageRef{543} /XYZ 72 666.023 null] >>}
\MusicPDFWriteObject{16281}{<< /D [\MusicPDFPageRef{545} /XYZ 72 484.811 null] >>}
\MusicPDFWriteObject{16282}{<< /D [\MusicPDFPageRef{546} /XYZ 72 624.048 null] >>}
\MusicPDFWriteObject{16283}{<< /D [\MusicPDFPageRef{548} /XYZ 72 698.082 null] >>}
\MusicPDFWriteObject{16284}{<< /D [\MusicPDFPageRef{553} /XYZ 72 135.835 null] >>}
\MusicPDFWriteObject{16285}{<< /D [\MusicPDFPageRef{555} /XYZ 72 219.716 null] >>}
\MusicPDFWriteObject{16286}{<< /D [\MusicPDFPageRef{559} /XYZ 72 176.685 null] >>}
\MusicPDFWriteObject{16287}{<< /D [\MusicPDFPageRef{561} /XYZ 72 599.629 null] >>}
\MusicPDFWriteObject{16288}{<< /D [\MusicPDFPageRef{561} /XYZ 72 417.767 null] >>}
\MusicPDFWriteObject{16289}{<< /D [\MusicPDFPageRef{562} /XYZ 72 348 null] >>}
\MusicPDFWriteObject{16290}{<< /D [\MusicPDFPageRef{566} /XYZ 72 614.493 null] >>}
\MusicPDFWriteObject{16291}{<< /D [\MusicPDFPageRef{569} /XYZ 72 698.082 null] >>}
\MusicPDFWriteObject{16292}{<< /D [\MusicPDFPageRef{578} /XYZ 72 282.389 null] >>}
\MusicPDFWriteObject{16293}{<< /D [\MusicPDFPageRef{582} /XYZ 72 248.573 null] >>}
\MusicPDFWriteObject{16294}{<< /D [\MusicPDFPageRef{588} /XYZ 72 274.505 null] >>}
\MusicPDFWriteObject{16295}{<< /D [\MusicPDFPageRef{592} /XYZ 72 665.074 null] >>}
\MusicPDFWriteObject{16296}{<< /D [\MusicPDFPageRef{600} /XYZ 72 433.085 null] >>}
\MusicPDFWriteObject{16297}{<< /D [\MusicPDFPageRef{604} /XYZ 72 538.824 null] >>}
\MusicPDFWriteObject{16298}{<< /D [\MusicPDFPageRef{605} /XYZ 72 513.943 null] >>}
\MusicPDFWriteObject{16299}{<< /D [\MusicPDFPageRef{606} /XYZ 72 136.195 null] >>}
\MusicPDFWriteObject{16300}{<< /D [\MusicPDFPageRef{614} /XYZ 72 515.142 null] >>}
\MusicPDFWriteObject{16301}{<< /D [\MusicPDFPageRef{617} /XYZ 72 587.04 null] >>}
\MusicPDFWriteObject{16302}{<< /D [\MusicPDFPageRef{619} /XYZ 72 370.464 null] >>}
\MusicPDFWriteObject{16303}{<< /D [\MusicPDFPageRef{626} /XYZ 72 488.51 null] >>}
\MusicPDFWriteObject{16304}{<< /D [\MusicPDFPageRef{627} /XYZ 72 540.481 null] >>}
\MusicPDFWriteObject{16305}{<< /D [\MusicPDFPageRef{632} /XYZ 72 698.082 null] >>}
\MusicPDFWriteObject{16306}{<< /D [\MusicPDFPageRef{636} /XYZ 72 623.958 null] >>}
\MusicPDFWriteObject{16307}{<< /D [\MusicPDFPageRef{643} /XYZ 72 248.854 null] >>}
\MusicPDFWriteObject{16308}{<< /D [\MusicPDFPageRef{648} /XYZ 72 160.287 null] >>}
\MusicPDFWriteObject{16309}{<< /D [\MusicPDFPageRef{657} /XYZ 72 542.167 null] >>}
\MusicPDFWriteObject{16310}{<< /D [\MusicPDFPageRef{658} /XYZ 72 530.344 null] >>}
\MusicPDFWriteObject{16311}{<< /D [\MusicPDFPageRef{658} /XYZ 72 109.893 null] >>}
\MusicPDFWriteObject{16312}{<< /D [\MusicPDFPageRef{660} /XYZ 72 188.798 null] >>}
\MusicPDFWriteObject{16313}{<< /D [\MusicPDFPageRef{663} /XYZ 72 421.283 null] >>}
\MusicPDFWriteObject{16314}{<< /D [\MusicPDFPageRef{669} /XYZ 72 440.437 null] >>}
\MusicPDFWriteObject{16315}{<< /D [\MusicPDFPageRef{670} /XYZ 72 477.11 null] >>}
\MusicPDFWriteObject{16316}{<< /D [\MusicPDFPageRef{673} /XYZ 72 698.082 null] >>}
\MusicPDFWriteObject{16317}{<< /D [\MusicPDFPageRef{675} /XYZ 72 149.899 null] >>}
\MusicPDFWriteObject{16318}{<< /D [\MusicPDFPageRef{685} /XYZ 72 534.155 null] >>}
\MusicPDFWriteObject{16319}{<< /D [\MusicPDFPageRef{691} /XYZ 72 358.507 null] >>}
\MusicPDFWriteObject{16320}{<< /D [\MusicPDFPageRef{693} /XYZ 72 436.994 null] >>}
\MusicPDFWriteObject{16321}{<< /D [\MusicPDFPageRef{694} /XYZ 72 398.266 null] >>}
\MusicPDFWriteObject{16322}{<< /D [\MusicPDFPageRef{696} /XYZ 72 378.793 null] >>}
\MusicPDFWriteObject{16323}{<< /D [\MusicPDFPageRef{697} /XYZ 72 388.603 null] >>}
\MusicPDFWriteObject{16324}{<< /D [\MusicPDFPageRef{699} /XYZ 72 586.974 null] >>}
\MusicPDFWriteObject{16325}{<< /D [\MusicPDFPageRef{703} /XYZ 72 361.131 null] >>}
\MusicPDFWriteObject{16326}{<< /D [\MusicPDFPageRef{707} /XYZ 72 483.645 null] >>}
\MusicPDFWriteObject{16327}{<< /D [\MusicPDFPageRef{709} /XYZ 72 698.082 null] >>}
\MusicPDFWriteObject{16328}{<< /D [\MusicPDFPageRef{710} /XYZ 72 448.328 null] >>}
\MusicPDFWriteObject{16329}{<< /D [\MusicPDFPageRef{711} /XYZ 72 627.256 null] >>}
\MusicPDFWriteObject{16330}{<< /D [\MusicPDFPageRef{717} /XYZ 72 698.082 null] >>}
\MusicPDFWriteObject{16331}{<< /D [\MusicPDFPageRef{720} /XYZ 72 698.082 null] >>}
\MusicPDFWriteObject{16332}{<< /D [\MusicPDFPageRef{737} /XYZ 72 469.115 null] >>}
\MusicPDFWriteObject{16333}{<< /D [\MusicPDFPageRef{739} /XYZ 72 429.742 null] >>}
\MusicPDFWriteObject{16334}{<< /D [\MusicPDFPageRef{740} /XYZ 72 132.677 null] >>}
\MusicPDFWriteObject{16335}{<< /D [\MusicPDFPageRef{746} /XYZ 72 651.326 null] >>}
\MusicPDFWriteObject{16336}{<< /D [\MusicPDFPageRef{783} /XYZ 72 291.305 null] >>}
\MusicPDFWriteObject{16337}{<< /D [\MusicPDFPageRef{747} /XYZ 72 386.873 null] >>}
\MusicPDFWriteObject{16338}{<< /D [\MusicPDFPageRef{749} /XYZ 72 368.321 null] >>}
\MusicPDFWriteObject{16339}{<< /D [\MusicPDFPageRef{757} /XYZ 72 348.461 null] >>}
\MusicPDFWriteObject{16340}{<< /D [\MusicPDFPageRef{760} /XYZ 72 370.907 null] >>}
\MusicPDFWriteObject{16341}{<< /D [\MusicPDFPageRef{765} /XYZ 72 337.132 null] >>}
\MusicPDFWriteObject{16342}{<< /D [\MusicPDFPageRef{776} /XYZ 72 618.527 null] >>}
\MusicPDFWriteObject{16343}{<< /D [\MusicPDFPageRef{777} /XYZ 72 399.596 null] >>}
\MusicPDFWriteObject{16344}{<< /D [\MusicPDFPageRef{780} /XYZ 72 295.455 null] >>}
\MusicPDFWriteObject{16345}{<< /D [\MusicPDFPageRef{792} /XYZ 72 513.537 null] >>}
\MusicPDFWriteObject{16346}{<< /D [\MusicPDFPageRef{794} /XYZ 72 520.666 null] >>}
\MusicPDFWriteObject{16347}{<< /D [\MusicPDFPageRef{797} /XYZ 72 274.962 null] >>}
\MusicPDFWriteObject{16348}{<< /D [\MusicPDFPageRef{801} /XYZ 72 608.522 null] >>}
\MusicPDFWriteObject{16349}{<< /D [\MusicPDFPageRef{804} /XYZ 72 589.602 null] >>}
\MusicPDFWriteObject{16350}{<< /D [\MusicPDFPageRef{815} /XYZ 72 382.033 null] >>}
\MusicPDFWriteObject{16351}{<< /D [\MusicPDFPageRef{818} /XYZ 72 564.103 null] >>}
\MusicPDFWriteObject{16352}{<< /D [\MusicPDFPageRef{819} /XYZ 72 437.544 null] >>}
\MusicPDFWriteObject{16353}{<< /D [\MusicPDFPageRef{828} /XYZ 72 476.338 null] >>}
\MusicPDFWriteObject{16354}{<< /D [\MusicPDFPageRef{833} /XYZ 72 605.437 null] >>}
\MusicPDFWriteObject{16355}{<< /D [\MusicPDFPageRef{846} /XYZ 72 548.478 null] >>}
\MusicPDFWriteObject{16356}{<< /D [\MusicPDFPageRef{851} /XYZ 72 652.741 null] >>}
\MusicPDFWriteObject{16357}{<< /D [\MusicPDFPageRef{855} /XYZ 72 698.082 null] >>}
\MusicPDFWriteObject{16358}{<< /D [\MusicPDFPageRef{865} /XYZ 72 665.55 null] >>}
\MusicPDFWriteObject{16359}{<< /D [\MusicPDFPageRef{868} /XYZ 72 452.577 null] >>}
\MusicPDFWriteObject{16360}{<< /D [\MusicPDFPageRef{872} /XYZ 72 160.442 null] >>}
\MusicPDFWriteObject{16361}{<< /D [\MusicPDFPageRef{874} /XYZ 72 343.85 null] >>}
\MusicPDFWriteObject{16362}{<< /D [\MusicPDFPageRef{877} /XYZ 72 488.905 null] >>}
\MusicPDFWriteObject{16363}{<< /D [\MusicPDFPageRef{879} /XYZ 72 264.609 null] >>}
\MusicPDFWriteObject{16364}{<< /D [\MusicPDFPageRef{883} /XYZ 72 429.305 null] >>}
\MusicPDFWriteObject{16365}{<< /D [\MusicPDFPageRef{886} /XYZ 72 598.827 null] >>}
\MusicPDFWriteObject{16366}{<< /D [\MusicPDFPageRef{889} /XYZ 72 249.861 null] >>}
\MusicPDFWriteObject{16367}{<< /D [\MusicPDFPageRef{898} /XYZ 72 265.03 null] >>}
\MusicPDFWriteObject{16368}{<< /D [\MusicPDFPageRef{903} /XYZ 72 512.736 null] >>}
\MusicPDFWriteObject{16369}{<< /D [\MusicPDFPageRef{907} /XYZ 72 413.522 null] >>}
\MusicPDFWriteObject{16370}{<< /D [\MusicPDFPageRef{910} /XYZ 72 583.897 null] >>}
\MusicPDFWriteObject{16371}{<< /D [\MusicPDFPageRef{913} /XYZ 72 583.869 null] >>}
\MusicPDFWriteObject{16372}{<< /D [\MusicPDFPageRef{914} /XYZ 72 640.244 null] >>}
\MusicPDFWriteObject{16373}{<< /D [\MusicPDFPageRef{920} /XYZ 72 239.908 null] >>}
\MusicPDFWriteObject{16374}{<< /D [\MusicPDFPageRef{922} /XYZ 72 333.73 null] >>}
\MusicPDFWriteObject{16375}{<< /D [\MusicPDFPageRef{924} /XYZ 72 402.335 null] >>}
\MusicPDFWriteObject{16376}{<< /D [\MusicPDFPageRef{928} /XYZ 72 258.997 null] >>}
\MusicPDFWriteObject{16377}{<< /D [\MusicPDFPageRef{42} /XYZ 72 426.353 null] >>}
\MusicPDFWriteObject{16378}{<< /D [\MusicPDFPageRef{43} /XYZ 72 609.215 null] >>}
\MusicPDFWriteObject{16379}{<< /D [\MusicPDFPageRef{44} /XYZ 72 528.6 null] >>}
\MusicPDFWriteObject{16380}{<< /D [\MusicPDFPageRef{45} /XYZ 72 661.818 null] >>}
\MusicPDFWriteObject{16381}{<< /D [\MusicPDFPageRef{45} /XYZ 72 161.303 null] >>}
\MusicPDFWriteObject{16382}{<< /D [\MusicPDFPageRef{48} /XYZ 72 396.706 null] >>}
\MusicPDFWriteObject{16383}{<< /D [\MusicPDFPageRef{53} /XYZ 72 347.519 null] >>}
\MusicPDFWriteObject{16384}{<< /D [\MusicPDFPageRef{53} /XYZ 72 175.113 null] >>}
\MusicPDFWriteObject{16385}{<< /D [\MusicPDFPageRef{54} /XYZ 72 606.683 null] >>}
\MusicPDFWriteObject{16386}{<< /D [\MusicPDFPageRef{55} /XYZ 72 698.082 null] >>}
\MusicPDFWriteObject{16387}{<< /D [\MusicPDFPageRef{439} /XYZ 72 206.779 null] >>}
\MusicPDFWriteObject{16388}{<< /D [\MusicPDFPageRef{443} /XYZ 72 446.86 null] >>}
\MusicPDFWriteObject{16389}{<< /D [\MusicPDFPageRef{445} /XYZ 72 491.351 null] >>}
\MusicPDFWriteObject{16390}{<< /D [\MusicPDFPageRef{446} /XYZ 72 365.9 null] >>}
\MusicPDFWriteObject{16391}{<< /D [\MusicPDFPageRef{461} /XYZ 72 429.798 null] >>}
\MusicPDFWriteObject{16392}{<< /D [\MusicPDFPageRef{463} /XYZ 72 458.662 null] >>}
\MusicPDFWriteObject{16393}{<< /D [\MusicPDFPageRef{464} /XYZ 72 154.534 null] >>}
\MusicPDFWriteObject{16394}{<< /D [\MusicPDFPageRef{471} /XYZ 72 138.685 null] >>}
\MusicPDFWriteObject{16395}{<< /D [\MusicPDFPageRef{491} /XYZ 72 382.245 null] >>}
\MusicPDFWriteObject{16396}{<< /D [\MusicPDFPageRef{59} /XYZ 72 384.339 null] >>}
\MusicPDFWriteObject{16397}{<< /D [\MusicPDFPageRef{60} /XYZ 72 272.304 null] >>}
\MusicPDFWriteObject{16398}{<< /D [\MusicPDFPageRef{62} /XYZ 72 467.696 null] >>}
\MusicPDFWriteObject{16399}{<< /D [\MusicPDFPageRef{63} /XYZ 72 514.241 null] >>}
\MusicPDFWriteObject{16400}{<< /D [\MusicPDFPageRef{63} /XYZ 72 306.019 null] >>}
\MusicPDFWriteObject{16401}{<< /D [\MusicPDFPageRef{65} /XYZ 72 201.414 null] >>}
\MusicPDFWriteObject{16402}{<< /D [\MusicPDFPageRef{67} /XYZ 72 410.348 null] >>}
\MusicPDFWriteObject{16403}{<< /D [\MusicPDFPageRef{77} /XYZ 72 520.134 null] >>}
\MusicPDFWriteObject{16404}{<< /D [\MusicPDFPageRef{82} /XYZ 72 459.343 null] >>}
\MusicPDFWriteObject{16405}{<< /D [\MusicPDFPageRef{91} /XYZ 72 474.299 null] >>}
\MusicPDFWriteObject{16406}{<< /D [\MusicPDFPageRef{95} /XYZ 72 330.17 null] >>}
\MusicPDFWriteObject{16407}{<< /D [\MusicPDFPageRef{97} /XYZ 72 431.322 null] >>}
\MusicPDFWriteObject{16408}{<< /D [\MusicPDFPageRef{100} /XYZ 72 661.081 null] >>}
\MusicPDFWriteObject{16409}{<< /D [\MusicPDFPageRef{104} /XYZ 72 210.72 null] >>}
\MusicPDFWriteObject{16410}{<< /D [\MusicPDFPageRef{105} /XYZ 72 585.886 null] >>}
\MusicPDFWriteObject{16411}{<< /D [\MusicPDFPageRef{120} /XYZ 72 258.93 null] >>}
\MusicPDFWriteObject{16412}{<< /D [\MusicPDFPageRef{121} /XYZ 72 570.69 null] >>}
\MusicPDFWriteObject{16413}{<< /D [\MusicPDFPageRef{121} /XYZ 72 438.726 null] >>}
\MusicPDFWriteObject{16414}{<< /D [\MusicPDFPageRef{132} /XYZ 72 561.807 null] >>}
\MusicPDFWriteObject{16415}{<< /D [\MusicPDFPageRef{135} /XYZ 72 475.923 null] >>}
\MusicPDFWriteObject{16416}{<< /D [\MusicPDFPageRef{138} /XYZ 72 688.12 null] >>}
\MusicPDFWriteObject{16417}{<< /D [\MusicPDFPageRef{141} /XYZ 72 536.582 null] >>}
\MusicPDFWriteObject{16418}{<< /D [\MusicPDFPageRef{145} /XYZ 72 698.082 null] >>}
\MusicPDFWriteObject{16419}{<< /D [\MusicPDFPageRef{145} /XYZ 72 112.384 null] >>}
\MusicPDFWriteObject{16420}{<< /D [\MusicPDFPageRef{152} /XYZ 72 688.12 null] >>}
\MusicPDFWriteObject{16421}{<< /D [\MusicPDFPageRef{152} /XYZ 72 177.648 null] >>}
\MusicPDFWriteObject{16422}{<< /D [\MusicPDFPageRef{159} /XYZ 72 522.722 null] >>}
\MusicPDFWriteObject{16423}{<< /D [\MusicPDFPageRef{163} /XYZ 72 266.188 null] >>}
\MusicPDFWriteObject{16424}{<< /D [\MusicPDFPageRef{165} /XYZ 72 548.219 null] >>}
\MusicPDFWriteObject{16425}{<< /D [\MusicPDFPageRef{175} /XYZ 72 698.082 null] >>}
\MusicPDFWriteObject{16426}{<< /D [\MusicPDFPageRef{184} /XYZ 72 268.45 null] >>}
\MusicPDFWriteObject{16427}{<< /D [\MusicPDFPageRef{185} /XYZ 72 112.384 null] >>}
\MusicPDFWriteObject{16428}{<< /D [\MusicPDFPageRef{192} /XYZ 72 143.848 null] >>}
\MusicPDFWriteObject{16429}{<< /D [\MusicPDFPageRef{197} /XYZ 72 265.775 null] >>}
\MusicPDFWriteObject{16430}{<< /D [\MusicPDFPageRef{204} /XYZ 72 524.333 null] >>}
\MusicPDFWriteObject{16431}{<< /D [\MusicPDFPageRef{210} /XYZ 72 527.693 null] >>}
\MusicPDFWriteObject{16432}{<< /D [\MusicPDFPageRef{409} /XYZ 72 320.965 null] >>}
\MusicPDFWriteObject{16433}{<< /D [\MusicPDFPageRef{412} /XYZ 72 622.366 null] >>}
\MusicPDFWriteObject{16434}{<< /D [\MusicPDFPageRef{416} /XYZ 72 587.341 null] >>}
\MusicPDFWriteObject{16435}{<< /D [\MusicPDFPageRef{422} /XYZ 72 514.305 null] >>}
\MusicPDFWriteObject{16436}{<< /D [\MusicPDFPageRef{425} /XYZ 72 410.673 null] >>}
\MusicPDFWriteObject{16437}{<< /D [\MusicPDFPageRef{426} /XYZ 72 332.512 null] >>}
\MusicPDFWriteObject{16438}{<< /D [\MusicPDFPageRef{429} /XYZ 72 609.215 null] >>}
\MusicPDFWriteObject{16439}{<< /D [\MusicPDFPageRef{539} /XYZ 72 556.839 null] >>}
\MusicPDFWriteObject{16440}{<< /D [\MusicPDFPageRef{540} /XYZ 72 222.837 null] >>}
\MusicPDFWriteObject{16441}{<< /D [\MusicPDFPageRef{541} /XYZ 72 568.581 null] >>}
\MusicPDFWriteObject{16442}{<< /D [\MusicPDFPageRef{542} /XYZ 72 655.081 null] >>}
\MusicPDFWriteObject{16443}{<< /D [\MusicPDFPageRef{542} /XYZ 72 375.929 null] >>}
\MusicPDFWriteObject{16444}{<< /D [\MusicPDFPageRef{543} /XYZ 72 602.76 null] >>}
\MusicPDFWriteObject{16445}{<< /D [\MusicPDFPageRef{543} /XYZ 72 367.457 null] >>}
\MusicPDFWriteObject{16446}{<< /D [\MusicPDFPageRef{544} /XYZ 72 698.082 null] >>}
\MusicPDFWriteObject{16447}{<< /D [\MusicPDFPageRef{545} /XYZ 72 420.994 null] >>}
\MusicPDFWriteObject{16448}{<< /D [\MusicPDFPageRef{546} /XYZ 72 463.112 null] >>}
\MusicPDFWriteObject{16449}{<< /D [\MusicPDFPageRef{547} /XYZ 72 507.95 null] >>}
\MusicPDFWriteObject{16450}{<< /D [\MusicPDFPageRef{547} /XYZ 72 276.798 null] >>}
\MusicPDFWriteObject{16451}{<< /D [\MusicPDFPageRef{548} /XYZ 72 635.517 null] >>}
\MusicPDFWriteObject{16452}{<< /D [\MusicPDFPageRef{548} /XYZ 72 117.447 null] >>}
\MusicPDFWriteObject{16453}{<< /D [\MusicPDFPageRef{550} /XYZ 72 698.082 null] >>}
\MusicPDFWriteObject{16454}{<< /D [\MusicPDFPageRef{552} /XYZ 72 698.082 null] >>}
\MusicPDFWriteObject{16455}{<< /D [\MusicPDFPageRef{552} /XYZ 72 501.957 null] >>}
\MusicPDFWriteObject{16456}{<< /D [\MusicPDFPageRef{553} /XYZ 72 698.082 null] >>}
\MusicPDFWriteObject{16457}{<< /D [\MusicPDFPageRef{553} /XYZ 72 500.722 null] >>}
\MusicPDFWriteObject{16458}{<< /D [\MusicPDFPageRef{554} /XYZ 72 698.082 null] >>}
\MusicPDFWriteObject{16459}{<< /D [\MusicPDFPageRef{554} /XYZ 72 286.394 null] >>}
\MusicPDFWriteObject{16460}{<< /D [\MusicPDFPageRef{557} /XYZ 72 698.082 null] >>}
\MusicPDFWriteObject{16461}{<< /D [\MusicPDFPageRef{561} /XYZ 72 562.921 null] >>}
\MusicPDFWriteObject{16462}{<< /D [\MusicPDFPageRef{566} /XYZ 72 339.754 null] >>}
\MusicPDFWriteObject{16463}{<< /D [\MusicPDFPageRef{568} /XYZ 72 522.867 null] >>}
\MusicPDFWriteObject{16464}{<< /D [\MusicPDFPageRef{569} /XYZ 72 464.557 null] >>}
\MusicPDFWriteObject{16465}{<< /D [\MusicPDFPageRef{574} /XYZ 72 517.618 null] >>}
\MusicPDFWriteObject{16466}{<< /D [\MusicPDFPageRef{579} /XYZ 72 317.522 null] >>}
\MusicPDFWriteObject{16467}{<< /D [\MusicPDFPageRef{588} /XYZ 72 112.384 null] >>}
\MusicPDFWriteObject{16468}{<< /D [\MusicPDFPageRef{592} /XYZ 72 224.499 null] >>}
\MusicPDFWriteObject{16469}{<< /D [\MusicPDFPageRef{595} /XYZ 72 241.784 null] >>}
\MusicPDFWriteObject{16470}{<< /D [\MusicPDFPageRef{600} /XYZ 72 382.973 null] >>}
\MusicPDFWriteObject{16471}{<< /D [\MusicPDFPageRef{604} /XYZ 72 488.158 null] >>}
\MusicPDFWriteObject{16472}{<< /D [\MusicPDFPageRef{605} /XYZ 72 448.189 null] >>}
\MusicPDFWriteObject{16473}{<< /D [\MusicPDFPageRef{607} /XYZ 72 688.12 null] >>}
\MusicPDFWriteObject{16474}{<< /D [\MusicPDFPageRef{610} /XYZ 72 249.878 null] >>}
\MusicPDFWriteObject{16475}{<< /D [\MusicPDFPageRef{614} /XYZ 72 478.18 null] >>}
\MusicPDFWriteObject{16476}{<< /D [\MusicPDFPageRef{615} /XYZ 72 387.908 null] >>}
\MusicPDFWriteObject{16477}{<< /D [\MusicPDFPageRef{615} /XYZ 72 169.733 null] >>}
\MusicPDFWriteObject{16478}{<< /D [\MusicPDFPageRef{616} /XYZ 72 318.772 null] >>}
\MusicPDFWriteObject{16479}{<< /D [\MusicPDFPageRef{617} /XYZ 72 497.474 null] >>}
\MusicPDFWriteObject{16480}{<< /D [\MusicPDFPageRef{619} /XYZ 72 333.502 null] >>}
\MusicPDFWriteObject{16481}{<< /D [\MusicPDFPageRef{621} /XYZ 72 160.802 null] >>}
\MusicPDFWriteObject{16482}{<< /D [\MusicPDFPageRef{622} /XYZ 72 251.423 null] >>}
\MusicPDFWriteObject{16483}{<< /D [\MusicPDFPageRef{623} /XYZ 72 406.001 null] >>}
\MusicPDFWriteObject{16484}{<< /D [\MusicPDFPageRef{627} /XYZ 72 450.362 null] >>}
\MusicPDFWriteObject{16485}{<< /D [\MusicPDFPageRef{632} /XYZ 72 648.667 null] >>}
\MusicPDFWriteObject{16486}{<< /D [\MusicPDFPageRef{632} /XYZ 72 333.473 null] >>}
\MusicPDFWriteObject{16487}{<< /D [\MusicPDFPageRef{633} /XYZ 72 556.248 null] >>}
\MusicPDFWriteObject{16488}{<< /D [\MusicPDFPageRef{634} /XYZ 72 432.625 null] >>}
\MusicPDFWriteObject{16489}{<< /D [\MusicPDFPageRef{636} /XYZ 72 573.292 null] >>}
\MusicPDFWriteObject{16490}{<< /D [\MusicPDFPageRef{636} /XYZ 72 349.86 null] >>}
\MusicPDFWriteObject{16491}{<< /D [\MusicPDFPageRef{640} /XYZ 72 688.12 null] >>}
\MusicPDFWriteObject{16492}{<< /D [\MusicPDFPageRef{643} /XYZ 72 207.715 null] >>}
\MusicPDFWriteObject{16493}{<< /D [\MusicPDFPageRef{644} /XYZ 72 698.082 null] >>}
\MusicPDFWriteObject{16494}{<< /D [\MusicPDFPageRef{645} /XYZ 72 286.12 null] >>}
\MusicPDFWriteObject{16495}{<< /D [\MusicPDFPageRef{649} /XYZ 72 698.082 null] >>}
\MusicPDFWriteObject{16496}{<< /D [\MusicPDFPageRef{651} /XYZ 72 321.039 null] >>}
\MusicPDFWriteObject{16497}{<< /D [\MusicPDFPageRef{652} /XYZ 72 594.732 null] >>}
\MusicPDFWriteObject{16498}{<< /D [\MusicPDFPageRef{653} /XYZ 72 563.969 null] >>}
\MusicPDFWriteObject{16499}{<< /D [\MusicPDFPageRef{658} /XYZ 72 648.667 null] >>}
\MusicPDFWriteObject{16500}{<< /D [\MusicPDFPageRef{658} /XYZ 72 338.252 null] >>}
\MusicPDFWriteObject{16501}{<< /D [\MusicPDFPageRef{666} /XYZ 72 564.579 null] >>}
\MusicPDFWriteObject{16502}{<< /D [\MusicPDFPageRef{673} /XYZ 72 648.667 null] >>}
\MusicPDFWriteObject{16503}{<< /D [\MusicPDFPageRef{674} /XYZ 72 414.851 null] >>}
\MusicPDFWriteObject{16504}{<< /D [\MusicPDFPageRef{675} /XYZ 72 635.517 null] >>}
\MusicPDFWriteObject{16505}{<< /D [\MusicPDFPageRef{676} /XYZ 72 648.667 null] >>}
\MusicPDFWriteObject{16506}{<< /D [\MusicPDFPageRef{679} /XYZ 72 698.082 null] >>}
\MusicPDFWriteObject{16507}{<< /D [\MusicPDFPageRef{681} /XYZ 72 261.232 null] >>}
\MusicPDFWriteObject{16508}{<< /D [\MusicPDFPageRef{682} /XYZ 72 543.068 null] >>}
\MusicPDFWriteObject{16509}{<< /D [\MusicPDFPageRef{683} /XYZ 72 226.411 null] >>}
\MusicPDFWriteObject{16510}{<< /D [\MusicPDFPageRef{685} /XYZ 72 363.208 null] >>}
\MusicPDFWriteObject{16511}{<< /D [\MusicPDFPageRef{686} /XYZ 72 388.357 null] >>}
\MusicPDFWriteObject{16512}{<< /D [\MusicPDFPageRef{688} /XYZ 72 295.108 null] >>}
\MusicPDFWriteObject{16513}{<< /D [\MusicPDFPageRef{689} /XYZ 72 698.082 null] >>}
\MusicPDFWriteObject{16514}{<< /D [\MusicPDFPageRef{689} /XYZ 72 291.368 null] >>}
\MusicPDFWriteObject{16515}{<< /D [\MusicPDFPageRef{690} /XYZ 72 358.592 null] >>}
\MusicPDFWriteObject{16516}{<< /D [\MusicPDFPageRef{697} /XYZ 72 337.937 null] >>}
\MusicPDFWriteObject{16517}{<< /D [\MusicPDFPageRef{699} /XYZ 72 145.376 null] >>}
\MusicPDFWriteObject{16518}{<< /D [\MusicPDFPageRef{705} /XYZ 72 698.082 null] >>}
\MusicPDFWriteObject{16519}{<< /D [\MusicPDFPageRef{705} /XYZ 72 158.814 null] >>}
\MusicPDFWriteObject{16520}{<< /D [\MusicPDFPageRef{711} /XYZ 72 510.836 null] >>}
\MusicPDFWriteObject{16521}{<< /D [\MusicPDFPageRef{713} /XYZ 72 220.042 null] >>}
\MusicPDFWriteObject{16522}{<< /D [\MusicPDFPageRef{716} /XYZ 72 446.975 null] >>}
\MusicPDFWriteObject{16523}{<< /D [\MusicPDFPageRef{721} /XYZ 72 635.517 null] >>}
\MusicPDFWriteObject{16524}{<< /D [\MusicPDFPageRef{723} /XYZ 72 603.483 null] >>}
\MusicPDFWriteObject{16525}{<< /D [\MusicPDFPageRef{726} /XYZ 72 601.189 null] >>}
\MusicPDFWriteObject{16526}{<< /D [\MusicPDFPageRef{731} /XYZ 72 221.912 null] >>}
\MusicPDFWriteObject{16527}{<< /D [\MusicPDFPageRef{735} /XYZ 72 609.215 null] >>}
\MusicPDFWriteObject{16528}{<< /D [\MusicPDFPageRef{737} /XYZ 72 432.154 null] >>}
\MusicPDFWriteObject{16529}{<< /D [\MusicPDFPageRef{740} /XYZ 72 475.966 null] >>}
\MusicPDFWriteObject{16530}{<< /D [\MusicPDFPageRef{741} /XYZ 72 302.735 null] >>}
\MusicPDFWriteObject{16531}{<< /D [\MusicPDFPageRef{784} /XYZ 72 507.958 null] >>}
\MusicPDFWriteObject{16532}{<< /D [\MusicPDFPageRef{789} /XYZ 72 483.12 null] >>}
\MusicPDFWriteObject{16533}{<< /D [\MusicPDFPageRef{791} /XYZ 72 292.708 null] >>}
\MusicPDFWriteObject{16534}{<< /D [\MusicPDFPageRef{750} /XYZ 72 581.806 null] >>}
\MusicPDFWriteObject{16535}{<< /D [\MusicPDFPageRef{751} /XYZ 72 608.433 null] >>}
\MusicPDFWriteObject{16536}{<< /D [\MusicPDFPageRef{752} /XYZ 72 280.233 null] >>}
\MusicPDFWriteObject{16537}{<< /D [\MusicPDFPageRef{758} /XYZ 72 635.517 null] >>}
\MusicPDFWriteObject{16538}{<< /D [\MusicPDFPageRef{758} /XYZ 72 367.338 null] >>}
\MusicPDFWriteObject{16539}{<< /D [\MusicPDFPageRef{759} /XYZ 72 698.082 null] >>}
\MusicPDFWriteObject{16540}{<< /D [\MusicPDFPageRef{760} /XYZ 72 307.644 null] >>}
\MusicPDFWriteObject{16541}{<< /D [\MusicPDFPageRef{761} /XYZ 72 175.162 null] >>}
\MusicPDFWriteObject{16542}{<< /D [\MusicPDFPageRef{763} /XYZ 72 335.949 null] >>}
\MusicPDFWriteObject{16543}{<< /D [\MusicPDFPageRef{766} /XYZ 72 312.434 null] >>}
\MusicPDFWriteObject{16544}{<< /D [\MusicPDFPageRef{767} /XYZ 72 122.172 null] >>}
\MusicPDFWriteObject{16545}{<< /D [\MusicPDFPageRef{771} /XYZ 72 365.22 null] >>}
\MusicPDFWriteObject{16546}{<< /D [\MusicPDFPageRef{777} /XYZ 72 176.033 null] >>}
\MusicPDFWriteObject{16547}{<< /D [\MusicPDFPageRef{779} /XYZ 72 610.213 null] >>}
\MusicPDFWriteObject{16548}{<< /D [\MusicPDFPageRef{792} /XYZ 72 410.822 null] >>}
\MusicPDFWriteObject{16549}{<< /D [\MusicPDFPageRef{797} /XYZ 72 548.283 null] >>}
\MusicPDFWriteObject{16550}{<< /D [\MusicPDFPageRef{801} /XYZ 72 356.499 null] >>}
\MusicPDFWriteObject{16551}{<< /D [\MusicPDFPageRef{802} /XYZ 72 325.159 null] >>}
\MusicPDFWriteObject{16552}{<< /D [\MusicPDFPageRef{803} /XYZ 72 583.57 null] >>}
\MusicPDFWriteObject{16553}{<< /D [\MusicPDFPageRef{806} /XYZ 72 571.4 null] >>}
\MusicPDFWriteObject{16554}{<< /D [\MusicPDFPageRef{812} /XYZ 72 337.799 null] >>}
\MusicPDFWriteObject{16555}{<< /D [\MusicPDFPageRef{842} /XYZ 72 178.137 null] >>}
\MusicPDFWriteObject{16556}{<< /D [\MusicPDFPageRef{888} /XYZ 72 499.932 null] >>}
\MusicPDFWriteObject{16557}{<< /D [\MusicPDFPageRef{888} /XYZ 72 228.768 null] >>}
\MusicPDFWriteObject{16558}{<< /D [\MusicPDFPageRef{889} /XYZ 72 223.006 null] >>}
\MusicPDFWriteObject{16559}{<< /D [\MusicPDFPageRef{890} /XYZ 72 698.082 null] >>}
\MusicPDFWriteObject{16560}{<< /D [\MusicPDFPageRef{894} /XYZ 72 255.265 null] >>}
\MusicPDFWriteObject{16561}{<< /D [\MusicPDFPageRef{898} /XYZ 72 146.673 null] >>}
\MusicPDFWriteObject{16562}{<< /D [\MusicPDFPageRef{900} /XYZ 72 324.053 null] >>}
\MusicPDFWriteObject{16563}{<< /D [\MusicPDFPageRef{902} /XYZ 72 649.584 null] >>}
\MusicPDFWriteObject{16564}{<< /D [\MusicPDFPageRef{903} /XYZ 72 447.536 null] >>}
\MusicPDFWriteObject{16565}{<< /D [\MusicPDFPageRef{906} /XYZ 72 367.083 null] >>}
\MusicPDFWriteObject{16566}{<< /D [\MusicPDFPageRef{910} /XYZ 72 531.847 null] >>}
\MusicPDFWriteObject{16567}{<< /D [\MusicPDFPageRef{914} /XYZ 72 432.623 null] >>}
\MusicPDFWriteObject{16568}{<< /D [\MusicPDFPageRef{917} /XYZ 72 679.652 null] >>}
\MusicPDFWriteObject{16569}{<< /D [\MusicPDFPageRef{440} /XYZ 72 235.126 null] >>}
\MusicPDFWriteObject{16570}{<< /D [\MusicPDFPageRef{444} /XYZ 72 655.364 null] >>}
\MusicPDFWriteObject{16571}{<< /D [\MusicPDFPageRef{446} /XYZ 72 484.572 null] >>}
\MusicPDFWriteObject{16572}{<< /D [\MusicPDFPageRef{462} /XYZ 72 701.27 null] >>}
\MusicPDFWriteObject{16573}{<< /D [\MusicPDFPageRef{485} /XYZ 72 126.579 null] >>}
\MusicPDFWriteObject{16574}{<< /D [\MusicPDFPageRef{180} /XYZ 72 629.266 null] >>}
\MusicPDFWriteObject{16575}{<< /D [\MusicPDFPageRef{181} /XYZ 72 211.16 null] >>}
\MusicPDFWriteObject{16576}{<< /D [\MusicPDFPageRef{184} /XYZ 72 701.27 null] >>}
\MusicPDFWriteObject{16577}{<< /D [\MusicPDFPageRef{303} /XYZ 72 344.016 null] >>}
\MusicPDFWriteObject{16578}{<< /D [\MusicPDFPageRef{583} /XYZ 72 635.457 null] >>}
\MusicPDFWriteObject{16579}{<< /D [\MusicPDFPageRef{812} /XYZ 72 574.134 null] >>}
\MusicPDFWriteObject{16580}{<< /D [\MusicPDFPageRef{824} /XYZ 72 506.504 null] >>}
\MusicPDFWriteObject{16581}{<< /D [\MusicPDFPageRef{824} /XYZ 72 127.627 null] >>}
\MusicPDFWriteObject{16582}{<< /D [\MusicPDFPageRef{918} /XYZ 72 640.317 null] >>}
\MusicPDFWriteObject{16583}{<< /D [\MusicPDFPageRef{925} /XYZ 72 568.001 null] >>}
\MusicPDFWriteObject{16584}{<< /D [\MusicPDFPageRef{35} /XYZ 112.621 628.124 null] >>}
\MusicPDFWriteObject{16585}{<< /D [\MusicPDFPageRef{40} /XYZ 72 454.419 null] >>}
\MusicPDFWriteObject{16586}{<< /D [\MusicPDFPageRef{434} /XYZ 72 345.653 null] >>}
\MusicPDFWriteObject{16587}{<< /D [\MusicPDFPageRef{434} /XYZ 72 125.649 null] >>}
\MusicPDFWriteObject{16588}{<< /D [\MusicPDFPageRef{435} /XYZ 72 578.154 null] >>}
\MusicPDFWriteObject{16589}{<< /D [\MusicPDFPageRef{435} /XYZ 72 210.839 null] >>}
\MusicPDFWriteObject{16590}{<< /D [\MusicPDFPageRef{436} /XYZ 72 319.076 null] >>}
\MusicPDFWriteObject{16591}{<< /D [\MusicPDFPageRef{437} /XYZ 72 421.463 null] >>}
\MusicPDFWriteObject{16592}{<< /D [\MusicPDFPageRef{440} /XYZ 72 420.67 null] >>}
\MusicPDFWriteObject{16593}{<< /D [\MusicPDFPageRef{440} /XYZ 72 164.987 null] >>}
\MusicPDFWriteObject{16594}{<< /D [\MusicPDFPageRef{443} /XYZ 72 319.732 null] >>}
\MusicPDFWriteObject{16595}{<< /D [\MusicPDFPageRef{444} /XYZ 72 443.103 null] >>}
\MusicPDFWriteObject{16596}{<< /D [\MusicPDFPageRef{444} /XYZ 72 244.444 null] >>}
\MusicPDFWriteObject{16597}{<< /D [\MusicPDFPageRef{446} /XYZ 72 701.27 null] >>}
\MusicPDFWriteObject{16598}{<< /D [\MusicPDFPageRef{446} /XYZ 72 258.762 null] >>}
\MusicPDFWriteObject{16599}{<< /D [\MusicPDFPageRef{447} /XYZ 72 539.596 null] >>}
\MusicPDFWriteObject{16600}{<< /D [\MusicPDFPageRef{448} /XYZ 72 475.55 null] >>}
\MusicPDFWriteObject{16601}{<< /D [\MusicPDFPageRef{453} /XYZ 72 284.916 null] >>}
\MusicPDFWriteObject{16602}{<< /D [\MusicPDFPageRef{456} /XYZ 72 541.97 null] >>}
\MusicPDFWriteObject{16603}{<< /D [\MusicPDFPageRef{456} /XYZ 72 370.52 null] >>}
\MusicPDFWriteObject{16604}{<< /D [\MusicPDFPageRef{457} /XYZ 72 589.074 null] >>}
\MusicPDFWriteObject{16605}{<< /D [\MusicPDFPageRef{457} /XYZ 72 343.541 null] >>}
\MusicPDFWriteObject{16606}{<< /D [\MusicPDFPageRef{458} /XYZ 72 559.43 null] >>}
\MusicPDFWriteObject{16607}{<< /D [\MusicPDFPageRef{459} /XYZ 72 701.27 null] >>}
\MusicPDFWriteObject{16608}{<< /D [\MusicPDFPageRef{460} /XYZ 72 572.909 null] >>}
\MusicPDFWriteObject{16609}{<< /D [\MusicPDFPageRef{460} /XYZ 72 474.245 null] >>}
\MusicPDFWriteObject{16610}{<< /D [\MusicPDFPageRef{460} /XYZ 72 300.194 null] >>}
\MusicPDFWriteObject{16611}{<< /D [\MusicPDFPageRef{461} /XYZ 72 408.845 null] >>}
\MusicPDFWriteObject{16612}{<< /D [\MusicPDFPageRef{463} /XYZ 72 701.27 null] >>}
\MusicPDFWriteObject{16613}{<< /D [\MusicPDFPageRef{463} /XYZ 72 586.541 null] >>}
\MusicPDFWriteObject{16614}{<< /D [\MusicPDFPageRef{463} /XYZ 72 422.838 null] >>}
\MusicPDFWriteObject{16615}{<< /D [\MusicPDFPageRef{466} /XYZ 72 212.419 null] >>}
\MusicPDFWriteObject{16616}{<< /D [\MusicPDFPageRef{470} /XYZ 72 412.782 null] >>}
\MusicPDFWriteObject{16617}{<< /D [\MusicPDFPageRef{470} /XYZ 72 262.099 null] >>}
\MusicPDFWriteObject{16618}{<< /D [\MusicPDFPageRef{470} /XYZ 72 193.428 null] >>}
\MusicPDFWriteObject{16619}{<< /D [\MusicPDFPageRef{472} /XYZ 72 701.27 null] >>}
\MusicPDFWriteObject{16620}{<< /D [\MusicPDFPageRef{473} /XYZ 72 538.878 null] >>}
\MusicPDFWriteObject{16621}{<< /D [\MusicPDFPageRef{473} /XYZ 72 246.012 null] >>}
\MusicPDFWriteObject{16622}{<< /D [\MusicPDFPageRef{475} /XYZ 72 181.921 null] >>}
\MusicPDFWriteObject{16623}{<< /D [\MusicPDFPageRef{476} /XYZ 72 625.447 null] >>}
\MusicPDFWriteObject{16624}{<< /D [\MusicPDFPageRef{476} /XYZ 72 326.816 null] >>}
\MusicPDFWriteObject{16625}{<< /D [\MusicPDFPageRef{478} /XYZ 72 617.577 null] >>}
\MusicPDFWriteObject{16626}{<< /D [\MusicPDFPageRef{478} /XYZ 72 128.936 null] >>}
\MusicPDFWriteObject{16627}{<< /D [\MusicPDFPageRef{480} /XYZ 72 578.593 null] >>}
\MusicPDFWriteObject{16628}{<< /D [\MusicPDFPageRef{480} /XYZ 72 147.909 null] >>}
\MusicPDFWriteObject{16629}{<< /D [\MusicPDFPageRef{484} /XYZ 72 514.194 null] >>}
\MusicPDFWriteObject{16630}{<< /D [\MusicPDFPageRef{485} /XYZ 72 464.957 null] >>}
\MusicPDFWriteObject{16631}{<< /D [\MusicPDFPageRef{485} /XYZ 72 201.97 null] >>}
\MusicPDFWriteObject{16632}{<< /D [\MusicPDFPageRef{486} /XYZ 72 385.616 null] >>}
\MusicPDFWriteObject{16633}{<< /D [\MusicPDFPageRef{486} /XYZ 72 125.534 null] >>}
\MusicPDFWriteObject{16634}{<< /D [\MusicPDFPageRef{487} /XYZ 72 288.711 null] >>}
\MusicPDFWriteObject{16635}{<< /D [\MusicPDFPageRef{488} /XYZ 72 569.269 null] >>}
\MusicPDFWriteObject{16636}{<< /D [\MusicPDFPageRef{488} /XYZ 72 286.523 null] >>}
\MusicPDFWriteObject{16637}{<< /D [\MusicPDFPageRef{489} /XYZ 72 550.256 null] >>}
\MusicPDFWriteObject{16638}{<< /D [\MusicPDFPageRef{490} /XYZ 72 663.664 null] >>}
\MusicPDFWriteObject{16639}{<< /D [\MusicPDFPageRef{491} /XYZ 72 547.146 null] >>}
\MusicPDFWriteObject{16640}{<< /D [\MusicPDFPageRef{495} /XYZ 72 614.702 null] >>}
\MusicPDFWriteObject{16641}{<< /D [\MusicPDFPageRef{498} /XYZ 72 479.052 null] >>}
\MusicPDFWriteObject{16642}{<< /D [\MusicPDFPageRef{499} /XYZ 72 599.782 null] >>}
\MusicPDFWriteObject{16643}{<< /D [\MusicPDFPageRef{502} /XYZ 72 186.477 null] >>}
\MusicPDFWriteObject{16644}{<< /D [\MusicPDFPageRef{503} /XYZ 72 274.378 null] >>}
\MusicPDFWriteObject{16645}{<< /D [\MusicPDFPageRef{504} /XYZ 72 167.832 null] >>}
\MusicPDFWriteObject{16646}{<< /D [\MusicPDFPageRef{506} /XYZ 72 641.932 null] >>}
\MusicPDFWriteObject{16647}{<< /D [\MusicPDFPageRef{506} /XYZ 72 125.534 null] >>}
\MusicPDFWriteObject{16648}{<< /D [\MusicPDFPageRef{507} /XYZ 72 646.755 null] >>}
\MusicPDFWriteObject{16649}{<< /D [\MusicPDFPageRef{507} /XYZ 72 138.685 null] >>}
\MusicPDFWriteObject{16650}{<< /D [\MusicPDFPageRef{508} /XYZ 72 548.964 null] >>}
\MusicPDFWriteObject{16651}{<< /D [\MusicPDFPageRef{509} /XYZ 72 634.586 null] >>}
\MusicPDFWriteObject{16652}{<< /D [\MusicPDFPageRef{512} /XYZ 72 627.749 null] >>}
\MusicPDFWriteObject{16653}{<< /D [\MusicPDFPageRef{514} /XYZ 72 521.394 null] >>}
\MusicPDFWriteObject{16654}{<< /D [\MusicPDFPageRef{519} /XYZ 72 350.051 null] >>}
\MusicPDFWriteObject{16655}{<< /D [\MusicPDFPageRef{522} /XYZ 72 530.908 null] >>}
\MusicPDFWriteObject{16656}{<< /D [\MusicPDFPageRef{522} /XYZ 72 166.653 null] >>}
\MusicPDFWriteObject{16657}{<< /D [\MusicPDFPageRef{525} /XYZ 72 491.059 null] >>}
\MusicPDFWriteObject{16658}{<< /D [\MusicPDFPageRef{529} /XYZ 72 378.1 null] >>}
\MusicPDFWriteObject{16659}{<< /D [\MusicPDFPageRef{530} /XYZ 72 481.292 null] >>}
\MusicPDFWriteObject{16660}{<< /D [\MusicPDFPageRef{531} /XYZ 72 701.27 null] >>}
\MusicPDFWriteObject{16661}{<< /D [\MusicPDFPageRef{531} /XYZ 72 599.93 null] >>}
\MusicPDFWriteObject{16662}{<< /D [\MusicPDFPageRef{531} /XYZ 72 161.962 null] >>}
\MusicPDFWriteObject{16663}{<< /D [\MusicPDFPageRef{532} /XYZ 72 585.485 null] >>}
\MusicPDFWriteObject{16664}{<< /D [\MusicPDFPageRef{532} /XYZ 72 219.108 null] >>}
\MusicPDFWriteObject{16665}{<< /D [\MusicPDFPageRef{532} /XYZ 72 185.472 null] >>}
\MusicPDFWriteObject{16666}{<< /D [\MusicPDFPageRef{59} /XYZ 72 284.92 null] >>}
\MusicPDFWriteObject{16667}{<< /D [\MusicPDFPageRef{59} /XYZ 72 138.685 null] >>}
\MusicPDFWriteObject{16668}{<< /D [\MusicPDFPageRef{60} /XYZ 72 414.402 null] >>}
\MusicPDFWriteObject{16669}{<< /D [\MusicPDFPageRef{60} /XYZ 72 130.597 null] >>}
\MusicPDFWriteObject{16670}{<< /D [\MusicPDFPageRef{61} /XYZ 72 598.464 null] >>}
\MusicPDFWriteObject{16671}{<< /D [\MusicPDFPageRef{61} /XYZ 72 367.309 null] >>}
\MusicPDFWriteObject{16672}{<< /D [\MusicPDFPageRef{61} /XYZ 72 331.901 null] >>}
\MusicPDFWriteObject{16673}{<< /D [\MusicPDFPageRef{63} /XYZ 72 368.146 null] >>}
\MusicPDFWriteObject{16674}{<< /D [\MusicPDFPageRef{66} /XYZ 72 568.298 null] >>}
\MusicPDFWriteObject{16675}{<< /D [\MusicPDFPageRef{66} /XYZ 72 191.288 null] >>}
\MusicPDFWriteObject{16676}{<< /D [\MusicPDFPageRef{67} /XYZ 72 661.887 null] >>}
\MusicPDFWriteObject{16677}{<< /D [\MusicPDFPageRef{63} /XYZ 72 138.685 null] >>}
\MusicPDFWriteObject{16678}{<< /D [\MusicPDFPageRef{64} /XYZ 72 667.334 null] >>}
\MusicPDFWriteObject{16679}{<< /D [\MusicPDFPageRef{64} /XYZ 72 426.761 null] >>}
\MusicPDFWriteObject{16680}{<< /D [\MusicPDFPageRef{64} /XYZ 72 202.91 null] >>}
\MusicPDFWriteObject{16681}{<< /D [\MusicPDFPageRef{65} /XYZ 72 701.27 null] >>}
\MusicPDFWriteObject{16682}{<< /D [\MusicPDFPageRef{65} /XYZ 72 577.843 null] >>}
\MusicPDFWriteObject{16683}{<< /D [\MusicPDFPageRef{65} /XYZ 72 531.574 null] >>}
\MusicPDFWriteObject{16684}{<< /D [\MusicPDFPageRef{65} /XYZ 72 430.258 null] >>}
\MusicPDFWriteObject{16685}{<< /D [\MusicPDFPageRef{67} /XYZ 72 357.814 null] >>}
\MusicPDFWriteObject{16686}{<< /D [\MusicPDFPageRef{68} /XYZ 72 496.014 null] >>}
\MusicPDFWriteObject{16687}{<< /D [\MusicPDFPageRef{68} /XYZ 72 267.168 null] >>}
\MusicPDFWriteObject{16688}{<< /D [\MusicPDFPageRef{70} /XYZ 72 304.898 null] >>}
\MusicPDFWriteObject{16689}{<< /D [\MusicPDFPageRef{70} /XYZ 72 183.246 null] >>}
\MusicPDFWriteObject{16690}{<< /D [\MusicPDFPageRef{71} /XYZ 72 701.27 null] >>}
\MusicPDFWriteObject{16691}{<< /D [\MusicPDFPageRef{72} /XYZ 72 292.383 null] >>}
\MusicPDFWriteObject{16692}{<< /D [\MusicPDFPageRef{73} /XYZ 72 701.27 null] >>}
\MusicPDFWriteObject{16693}{<< /D [\MusicPDFPageRef{74} /XYZ 72 559.215 null] >>}
\MusicPDFWriteObject{16694}{<< /D [\MusicPDFPageRef{75} /XYZ 72 584.533 null] >>}
\MusicPDFWriteObject{16695}{<< /D [\MusicPDFPageRef{75} /XYZ 72 492.974 null] >>}
\MusicPDFWriteObject{16696}{<< /D [\MusicPDFPageRef{76} /XYZ 72 559.66 null] >>}
\MusicPDFWriteObject{16697}{<< /D [\MusicPDFPageRef{77} /XYZ 72 308.363 null] >>}
\MusicPDFWriteObject{16698}{<< /D [\MusicPDFPageRef{78} /XYZ 72 580.772 null] >>}
\MusicPDFWriteObject{16699}{<< /D [\MusicPDFPageRef{81} /XYZ 72 453.234 null] >>}
\MusicPDFWriteObject{16700}{<< /D [\MusicPDFPageRef{81} /XYZ 72 154.812 null] >>}
\MusicPDFWriteObject{16701}{<< /D [\MusicPDFPageRef{82} /XYZ 72 426.494 null] >>}
\MusicPDFWriteObject{16702}{<< /D [\MusicPDFPageRef{82} /XYZ 72 375.431 null] >>}
\MusicPDFWriteObject{16703}{<< /D [\MusicPDFPageRef{82} /XYZ 72 112.384 null] >>}
\MusicPDFWriteObject{16704}{<< /D [\MusicPDFPageRef{86} /XYZ 72 99.233 null] >>}
\MusicPDFWriteObject{16705}{<< /D [\MusicPDFPageRef{87} /XYZ 72 480.323 null] >>}
\MusicPDFWriteObject{16706}{<< /D [\MusicPDFPageRef{88} /XYZ 72 316.753 null] >>}
\MusicPDFWriteObject{16707}{<< /D [\MusicPDFPageRef{91} /XYZ 72 427.94 null] >>}
\MusicPDFWriteObject{16708}{<< /D [\MusicPDFPageRef{91} /XYZ 72 311.464 null] >>}
\MusicPDFWriteObject{16709}{<< /D [\MusicPDFPageRef{92} /XYZ 72 560.721 null] >>}
\MusicPDFWriteObject{16710}{<< /D [\MusicPDFPageRef{95} /XYZ 72 499.57 null] >>}
\MusicPDFWriteObject{16711}{<< /D [\MusicPDFPageRef{95} /XYZ 72 294.623 null] >>}
\MusicPDFWriteObject{16712}{<< /D [\MusicPDFPageRef{97} /XYZ 72 358.484 null] >>}
\MusicPDFWriteObject{16713}{<< /D [\MusicPDFPageRef{97} /XYZ 72 274.688 null] >>}
\MusicPDFWriteObject{16714}{<< /D [\MusicPDFPageRef{99} /XYZ 72 182.211 null] >>}
\MusicPDFWriteObject{16715}{<< /D [\MusicPDFPageRef{100} /XYZ 72 607.741 null] >>}
\MusicPDFWriteObject{16716}{<< /D [\MusicPDFPageRef{104} /XYZ 72 112.499 null] >>}
\MusicPDFWriteObject{16717}{<< /D [\MusicPDFPageRef{105} /XYZ 72 515.642 null] >>}
\MusicPDFWriteObject{16718}{<< /D [\MusicPDFPageRef{105} /XYZ 72 182.449 null] >>}
\MusicPDFWriteObject{16719}{<< /D [\MusicPDFPageRef{106} /XYZ 72 635.045 null] >>}
\MusicPDFWriteObject{16720}{<< /D [\MusicPDFPageRef{106} /XYZ 72 440.879 null] >>}
\MusicPDFWriteObject{16721}{<< /D [\MusicPDFPageRef{107} /XYZ 72 664.904 null] >>}
\MusicPDFWriteObject{16722}{<< /D [\MusicPDFPageRef{107} /XYZ 72 565.768 null] >>}
\MusicPDFWriteObject{16723}{<< /D [\MusicPDFPageRef{111} /XYZ 72 701.27 null] >>}
\MusicPDFWriteObject{16724}{<< /D [\MusicPDFPageRef{113} /XYZ 72 666.315 null] >>}
\MusicPDFWriteObject{16725}{<< /D [\MusicPDFPageRef{115} /XYZ 72 701.27 null] >>}
\MusicPDFWriteObject{16726}{<< /D [\MusicPDFPageRef{119} /XYZ 72 324.363 null] >>}
\MusicPDFWriteObject{16727}{<< /D [\MusicPDFPageRef{119} /XYZ 72 266.877 null] >>}
\MusicPDFWriteObject{16728}{<< /D [\MusicPDFPageRef{121} /XYZ 72 486.494 null] >>}
\MusicPDFWriteObject{16729}{<< /D [\MusicPDFPageRef{121} /XYZ 72 112.384 null] >>}
\MusicPDFWriteObject{16730}{<< /D [\MusicPDFPageRef{123} /XYZ 72 701.243 null] >>}
\MusicPDFWriteObject{16731}{<< /D [\MusicPDFPageRef{123} /XYZ 72 454.467 null] >>}
\MusicPDFWriteObject{16732}{<< /D [\MusicPDFPageRef{123} /XYZ 72 388.634 null] >>}
\MusicPDFWriteObject{16733}{<< /D [\MusicPDFPageRef{123} /XYZ 72 112.936 null] >>}
\MusicPDFWriteObject{16734}{<< /D [\MusicPDFPageRef{124} /XYZ 72 349.242 null] >>}
\MusicPDFWriteObject{16735}{<< /D [\MusicPDFPageRef{125} /XYZ 72 399.908 null] >>}
\MusicPDFWriteObject{16736}{<< /D [\MusicPDFPageRef{125} /XYZ 72 246.805 null] >>}
\MusicPDFWriteObject{16737}{<< /D [\MusicPDFPageRef{127} /XYZ 72 221.867 null] >>}
\MusicPDFWriteObject{16738}{<< /D [\MusicPDFPageRef{128} /XYZ 72 458.232 null] >>}
\MusicPDFWriteObject{16739}{<< /D [\MusicPDFPageRef{129} /XYZ 72 701.27 null] >>}
\MusicPDFWriteObject{16740}{<< /D [\MusicPDFPageRef{129} /XYZ 72 227.836 null] >>}
\MusicPDFWriteObject{16741}{<< /D [\MusicPDFPageRef{130} /XYZ 72 294.93 null] >>}
\MusicPDFWriteObject{16742}{<< /D [\MusicPDFPageRef{131} /XYZ 72 405.36 null] >>}
\MusicPDFWriteObject{16743}{<< /D [\MusicPDFPageRef{132} /XYZ 72 516.358 null] >>}
\MusicPDFWriteObject{16744}{<< /D [\MusicPDFPageRef{135} /XYZ 72 650.459 null] >>}
\MusicPDFWriteObject{16745}{<< /D [\MusicPDFPageRef{135} /XYZ 72 438.263 null] >>}
\MusicPDFWriteObject{16746}{<< /D [\MusicPDFPageRef{136} /XYZ 72 460.093 null] >>}
\MusicPDFWriteObject{16747}{<< /D [\MusicPDFPageRef{138} /XYZ 72 491.777 null] >>}
\MusicPDFWriteObject{16748}{<< /D [\MusicPDFPageRef{141} /XYZ 72 501.462 null] >>}
\MusicPDFWriteObject{16749}{<< /D [\MusicPDFPageRef{143} /XYZ 72 701.27 null] >>}
\MusicPDFWriteObject{16750}{<< /D [\MusicPDFPageRef{145} /XYZ 72 601.04 null] >>}
\MusicPDFWriteObject{16751}{<< /D [\MusicPDFPageRef{146} /XYZ 72 536.548 null] >>}
\MusicPDFWriteObject{16752}{<< /D [\MusicPDFPageRef{146} /XYZ 72 259.403 null] >>}
\MusicPDFWriteObject{16753}{<< /D [\MusicPDFPageRef{149} /XYZ 72 170.05 null] >>}
\MusicPDFWriteObject{16754}{<< /D [\MusicPDFPageRef{151} /XYZ 72 285.532 null] >>}
\MusicPDFWriteObject{16755}{<< /D [\MusicPDFPageRef{152} /XYZ 72 653.018 null] >>}
\MusicPDFWriteObject{16756}{<< /D [\MusicPDFPageRef{152} /XYZ 72 142.547 null] >>}
\MusicPDFWriteObject{16757}{<< /D [\MusicPDFPageRef{153} /XYZ 72 397.093 null] >>}
\MusicPDFWriteObject{16758}{<< /D [\MusicPDFPageRef{157} /XYZ 72 484.839 null] >>}
\MusicPDFWriteObject{16759}{<< /D [\MusicPDFPageRef{158} /XYZ 72 327.638 null] >>}
\MusicPDFWriteObject{16760}{<< /D [\MusicPDFPageRef{158} /XYZ 72 112.384 null] >>}
\MusicPDFWriteObject{16761}{<< /D [\MusicPDFPageRef{159} /XYZ 72 625.147 null] >>}
\MusicPDFWriteObject{16762}{<< /D [\MusicPDFPageRef{159} /XYZ 72 407.146 null] >>}
\MusicPDFWriteObject{16763}{<< /D [\MusicPDFPageRef{160} /XYZ 72 396.585 null] >>}
\MusicPDFWriteObject{16764}{<< /D [\MusicPDFPageRef{161} /XYZ 72 588.851 null] >>}
\MusicPDFWriteObject{16765}{<< /D [\MusicPDFPageRef{162} /XYZ 72 588.881 null] >>}
\MusicPDFWriteObject{16766}{<< /D [\MusicPDFPageRef{162} /XYZ 72 412.759 null] >>}
\MusicPDFWriteObject{16767}{<< /D [\MusicPDFPageRef{163} /XYZ 72 231.467 null] >>}
\MusicPDFWriteObject{16768}{<< /D [\MusicPDFPageRef{164} /XYZ 72 244.565 null] >>}
\MusicPDFWriteObject{16769}{<< /D [\MusicPDFPageRef{165} /XYZ 72 388.858 null] >>}
\MusicPDFWriteObject{16770}{<< /D [\MusicPDFPageRef{165} /XYZ 72 243.891 null] >>}
\MusicPDFWriteObject{16771}{<< /D [\MusicPDFPageRef{169} /XYZ 72 540.835 null] >>}
\MusicPDFWriteObject{16772}{<< /D [\MusicPDFPageRef{170} /XYZ 72 701.27 null] >>}
\MusicPDFWriteObject{16773}{<< /D [\MusicPDFPageRef{170} /XYZ 72 193.003 null] >>}
\MusicPDFWriteObject{16774}{<< /D [\MusicPDFPageRef{171} /XYZ 72 558.931 null] >>}
\MusicPDFWriteObject{16775}{<< /D [\MusicPDFPageRef{171} /XYZ 72 393.567 null] >>}
\MusicPDFWriteObject{16776}{<< /D [\MusicPDFPageRef{171} /XYZ 72 338.834 null] >>}
\MusicPDFWriteObject{16777}{<< /D [\MusicPDFPageRef{173} /XYZ 72 501.931 null] >>}
\MusicPDFWriteObject{16778}{<< /D [\MusicPDFPageRef{173} /XYZ 72 231.458 null] >>}
\MusicPDFWriteObject{16779}{<< /D [\MusicPDFPageRef{174} /XYZ 72 479.408 null] >>}
\MusicPDFWriteObject{16780}{<< /D [\MusicPDFPageRef{175} /XYZ 72 667.387 null] >>}
\MusicPDFWriteObject{16781}{<< /D [\MusicPDFPageRef{175} /XYZ 72 512.463 null] >>}
\MusicPDFWriteObject{16782}{<< /D [\MusicPDFPageRef{175} /XYZ 72 326.34 null] >>}
\MusicPDFWriteObject{16783}{<< /D [\MusicPDFPageRef{175} /XYZ 72 258.574 null] >>}
\MusicPDFWriteObject{16784}{<< /D [\MusicPDFPageRef{176} /XYZ 72 701.27 null] >>}
\MusicPDFWriteObject{16785}{<< /D [\MusicPDFPageRef{176} /XYZ 72 490.731 null] >>}
\MusicPDFWriteObject{16786}{<< /D [\MusicPDFPageRef{179} /XYZ 72 683.552 null] >>}
\MusicPDFWriteObject{16787}{<< /D [\MusicPDFPageRef{179} /XYZ 72 394.893 null] >>}
\MusicPDFWriteObject{16788}{<< /D [\MusicPDFPageRef{180} /XYZ 72 168.527 null] >>}
\MusicPDFWriteObject{16789}{<< /D [\MusicPDFPageRef{181} /XYZ 72 424.858 null] >>}
\MusicPDFWriteObject{16790}{<< /D [\MusicPDFPageRef{183} /XYZ 72 372.468 null] >>}
\MusicPDFWriteObject{16791}{<< /D [\MusicPDFPageRef{184} /XYZ 72 227.716 null] >>}
\MusicPDFWriteObject{16792}{<< /D [\MusicPDFPageRef{185} /XYZ 72 466.103 null] >>}
\MusicPDFWriteObject{16793}{<< /D [\MusicPDFPageRef{185} /XYZ 72 223.713 null] >>}
\MusicPDFWriteObject{16794}{<< /D [\MusicPDFPageRef{186} /XYZ 72 650.699 null] >>}
\MusicPDFWriteObject{16795}{<< /D [\MusicPDFPageRef{188} /XYZ 72 633.931 null] >>}
\MusicPDFWriteObject{16796}{<< /D [\MusicPDFPageRef{190} /XYZ 72 437.498 null] >>}
\MusicPDFWriteObject{16797}{<< /D [\MusicPDFPageRef{191} /XYZ 72 125.534 null] >>}
\MusicPDFWriteObject{16798}{<< /D [\MusicPDFPageRef{192} /XYZ 72 558.6 null] >>}
\MusicPDFWriteObject{16799}{<< /D [\MusicPDFPageRef{192} /XYZ 72 458.31 null] >>}
\MusicPDFWriteObject{16800}{<< /D [\MusicPDFPageRef{192} /XYZ 72 303.69 null] >>}
\MusicPDFWriteObject{16801}{<< /D [\MusicPDFPageRef{193} /XYZ 72 667.288 null] >>}
\MusicPDFWriteObject{16802}{<< /D [\MusicPDFPageRef{193} /XYZ 72 278.757 null] >>}
\MusicPDFWriteObject{16803}{<< /D [\MusicPDFPageRef{194} /XYZ 72 596.079 null] >>}
\MusicPDFWriteObject{16804}{<< /D [\MusicPDFPageRef{194} /XYZ 72 506.37 null] >>}
\MusicPDFWriteObject{16805}{<< /D [\MusicPDFPageRef{194} /XYZ 72 228.853 null] >>}
\MusicPDFWriteObject{16806}{<< /D [\MusicPDFPageRef{197} /XYZ 72 205.81 null] >>}
\MusicPDFWriteObject{16807}{<< /D [\MusicPDFPageRef{198} /XYZ 72 499.513 null] >>}
\MusicPDFWriteObject{16808}{<< /D [\MusicPDFPageRef{198} /XYZ 72 337.209 null] >>}
\MusicPDFWriteObject{16809}{<< /D [\MusicPDFPageRef{203} /XYZ 72 385.345 null] >>}
\MusicPDFWriteObject{16810}{<< /D [\MusicPDFPageRef{203} /XYZ 72 243.448 null] >>}
\MusicPDFWriteObject{16811}{<< /D [\MusicPDFPageRef{204} /XYZ 72 371.955 null] >>}
\MusicPDFWriteObject{16812}{<< /D [\MusicPDFPageRef{205} /XYZ 72 341.895 null] >>}
\MusicPDFWriteObject{16813}{<< /D [\MusicPDFPageRef{207} /XYZ 72 639.151 null] >>}
\MusicPDFWriteObject{16814}{<< /D [\MusicPDFPageRef{207} /XYZ 72 471.826 null] >>}
\MusicPDFWriteObject{16815}{<< /D [\MusicPDFPageRef{208} /XYZ 72 562.024 null] >>}
\MusicPDFWriteObject{16816}{<< /D [\MusicPDFPageRef{209} /XYZ 72 495.535 null] >>}
\MusicPDFWriteObject{16817}{<< /D [\MusicPDFPageRef{209} /XYZ 72 153.486 null] >>}
\MusicPDFWriteObject{16818}{<< /D [\MusicPDFPageRef{210} /XYZ 72 507.408 null] >>}
\MusicPDFWriteObject{16819}{<< /D [\MusicPDFPageRef{211} /XYZ 72 397.529 null] >>}
\MusicPDFWriteObject{16820}{<< /D [\MusicPDFPageRef{212} /XYZ 72 595.403 null] >>}
\MusicPDFWriteObject{16821}{<< /D [\MusicPDFPageRef{217} /XYZ 72 242.521 null] >>}
\MusicPDFWriteObject{16822}{<< /D [\MusicPDFPageRef{218} /XYZ 72 414.951 null] >>}
\MusicPDFWriteObject{16823}{<< /D [\MusicPDFPageRef{219} /XYZ 72 264.074 null] >>}
\MusicPDFWriteObject{16824}{<< /D [\MusicPDFPageRef{223} /XYZ 72 628.836 null] >>}
\MusicPDFWriteObject{16825}{<< /D [\MusicPDFPageRef{223} /XYZ 72 189.207 null] >>}
\MusicPDFWriteObject{16826}{<< /D [\MusicPDFPageRef{223} /XYZ 72 125.534 null] >>}
\MusicPDFWriteObject{16827}{<< /D [\MusicPDFPageRef{224} /XYZ 72 339.532 null] >>}
\MusicPDFWriteObject{16828}{<< /D [\MusicPDFPageRef{225} /XYZ 72 425.786 null] >>}
\MusicPDFWriteObject{16829}{<< /D [\MusicPDFPageRef{227} /XYZ 72 341.198 null] >>}
\MusicPDFWriteObject{16830}{<< /D [\MusicPDFPageRef{230} /XYZ 72 474.849 null] >>}
\MusicPDFWriteObject{16831}{<< /D [\MusicPDFPageRef{232} /XYZ 72 321.72 null] >>}
\MusicPDFWriteObject{16832}{<< /D [\MusicPDFPageRef{233} /XYZ 72 650.644 null] >>}
\MusicPDFWriteObject{16833}{<< /D [\MusicPDFPageRef{236} /XYZ 72 701.27 null] >>}
\MusicPDFWriteObject{16834}{<< /D [\MusicPDFPageRef{236} /XYZ 72 251.938 null] >>}
\MusicPDFWriteObject{16835}{<< /D [\MusicPDFPageRef{237} /XYZ 72 331.147 null] >>}
\MusicPDFWriteObject{16836}{<< /D [\MusicPDFPageRef{238} /XYZ 72 359.332 null] >>}
\MusicPDFWriteObject{16837}{<< /D [\MusicPDFPageRef{238} /XYZ 72 149.421 null] >>}
\MusicPDFWriteObject{16838}{<< /D [\MusicPDFPageRef{240} /XYZ 72 584.305 null] >>}
\MusicPDFWriteObject{16839}{<< /D [\MusicPDFPageRef{241} /XYZ 72 587.687 null] >>}
\MusicPDFWriteObject{16840}{<< /D [\MusicPDFPageRef{241} /XYZ 72 443.328 null] >>}
\MusicPDFWriteObject{16841}{<< /D [\MusicPDFPageRef{241} /XYZ 72 160.088 null] >>}
\MusicPDFWriteObject{16842}{<< /D [\MusicPDFPageRef{242} /XYZ 72 359.058 null] >>}
\MusicPDFWriteObject{16843}{<< /D [\MusicPDFPageRef{243} /XYZ 72 470.892 null] >>}
\MusicPDFWriteObject{16844}{<< /D [\MusicPDFPageRef{243} /XYZ 72 171.288 null] >>}
\MusicPDFWriteObject{16845}{<< /D [\MusicPDFPageRef{245} /XYZ 72 610.282 null] >>}
\MusicPDFWriteObject{16846}{<< /D [\MusicPDFPageRef{246} /XYZ 72 474.486 null] >>}
\MusicPDFWriteObject{16847}{<< /D [\MusicPDFPageRef{246} /XYZ 72 386.6 null] >>}
\MusicPDFWriteObject{16848}{<< /D [\MusicPDFPageRef{249} /XYZ 72 397.927 null] >>}
\MusicPDFWriteObject{16849}{<< /D [\MusicPDFPageRef{249} /XYZ 72 195.453 null] >>}
\MusicPDFWriteObject{16850}{<< /D [\MusicPDFPageRef{252} /XYZ 72 516.957 null] >>}
\MusicPDFWriteObject{16851}{<< /D [\MusicPDFPageRef{257} /XYZ 72 174.987 null] >>}
\MusicPDFWriteObject{16852}{<< /D [\MusicPDFPageRef{259} /XYZ 72 272.943 null] >>}
\MusicPDFWriteObject{16853}{<< /D [\MusicPDFPageRef{259} /XYZ 72 112.384 null] >>}
\MusicPDFWriteObject{16854}{<< /D [\MusicPDFPageRef{261} /XYZ 72 326.078 null] >>}
\MusicPDFWriteObject{16855}{<< /D [\MusicPDFPageRef{263} /XYZ 72 125.534 null] >>}
\MusicPDFWriteObject{16856}{<< /D [\MusicPDFPageRef{264} /XYZ 72 383.712 null] >>}
\MusicPDFWriteObject{16857}{<< /D [\MusicPDFPageRef{265} /XYZ 72 176.479 null] >>}
\MusicPDFWriteObject{16858}{<< /D [\MusicPDFPageRef{266} /XYZ 72 556.391 null] >>}
\MusicPDFWriteObject{16859}{<< /D [\MusicPDFPageRef{266} /XYZ 72 319.645 null] >>}
\MusicPDFWriteObject{16860}{<< /D [\MusicPDFPageRef{267} /XYZ 72 502.3 null] >>}
\MusicPDFWriteObject{16861}{<< /D [\MusicPDFPageRef{274} /XYZ 72 451.375 null] >>}
\MusicPDFWriteObject{16862}{<< /D [\MusicPDFPageRef{275} /XYZ 72 648.947 null] >>}
\MusicPDFWriteObject{16863}{<< /D [\MusicPDFPageRef{276} /XYZ 72 408.686 null] >>}
\MusicPDFWriteObject{16864}{<< /D [\MusicPDFPageRef{277} /XYZ 72 557.439 null] >>}
\MusicPDFWriteObject{16865}{<< /D [\MusicPDFPageRef{277} /XYZ 72 242.144 null] >>}
\MusicPDFWriteObject{16866}{<< /D [\MusicPDFPageRef{278} /XYZ 72 565.332 null] >>}
\MusicPDFWriteObject{16867}{<< /D [\MusicPDFPageRef{279} /XYZ 72 364.058 null] >>}
\MusicPDFWriteObject{16868}{<< /D [\MusicPDFPageRef{280} /XYZ 72 586.865 null] >>}
\MusicPDFWriteObject{16869}{<< /D [\MusicPDFPageRef{280} /XYZ 72 325.699 null] >>}
\MusicPDFWriteObject{16870}{<< /D [\MusicPDFPageRef{281} /XYZ 72 701.27 null] >>}
\MusicPDFWriteObject{16871}{<< /D [\MusicPDFPageRef{281} /XYZ 72 462.332 null] >>}
\MusicPDFWriteObject{16872}{<< /D [\MusicPDFPageRef{282} /XYZ 72 670.412 null] >>}
\MusicPDFWriteObject{16873}{<< /D [\MusicPDFPageRef{285} /XYZ 72 639.967 null] >>}
\MusicPDFWriteObject{16874}{<< /D [\MusicPDFPageRef{285} /XYZ 72 480.741 null] >>}
\MusicPDFWriteObject{16875}{<< /D [\MusicPDFPageRef{287} /XYZ 72 320.367 null] >>}
\MusicPDFWriteObject{16876}{<< /D [\MusicPDFPageRef{288} /XYZ 72 398.987 null] >>}
\MusicPDFWriteObject{16877}{<< /D [\MusicPDFPageRef{288} /XYZ 72 279.984 null] >>}
\MusicPDFWriteObject{16878}{<< /D [\MusicPDFPageRef{289} /XYZ 72 566.513 null] >>}
\MusicPDFWriteObject{16879}{<< /D [\MusicPDFPageRef{291} /XYZ 72 142.281 null] >>}
\MusicPDFWriteObject{16880}{<< /D [\MusicPDFPageRef{295} /XYZ 72 309.177 null] >>}
\MusicPDFWriteObject{16881}{<< /D [\MusicPDFPageRef{295} /XYZ 72 189.762 null] >>}
\MusicPDFWriteObject{16882}{<< /D [\MusicPDFPageRef{296} /XYZ 72 624.113 null] >>}
\MusicPDFWriteObject{16883}{<< /D [\MusicPDFPageRef{297} /XYZ 72 259.992 null] >>}
\MusicPDFWriteObject{16884}{<< /D [\MusicPDFPageRef{298} /XYZ 72 395.642 null] >>}
\MusicPDFWriteObject{16885}{<< /D [\MusicPDFPageRef{299} /XYZ 72 541.827 null] >>}
\MusicPDFWriteObject{16886}{<< /D [\MusicPDFPageRef{304} /XYZ 72 546.193 null] >>}
\MusicPDFWriteObject{16887}{<< /D [\MusicPDFPageRef{309} /XYZ 72 188.264 null] >>}
\MusicPDFWriteObject{16888}{<< /D [\MusicPDFPageRef{313} /XYZ 72 292.015 null] >>}
\MusicPDFWriteObject{16889}{<< /D [\MusicPDFPageRef{316} /XYZ 72 473.906 null] >>}
\MusicPDFWriteObject{16890}{<< /D [\MusicPDFPageRef{316} /XYZ 72 178.137 null] >>}
\MusicPDFWriteObject{16891}{<< /D [\MusicPDFPageRef{317} /XYZ 72 113.07 null] >>}
\MusicPDFWriteObject{16892}{<< /D [\MusicPDFPageRef{318} /XYZ 72 509.05 null] >>}
\MusicPDFWriteObject{16893}{<< /D [\MusicPDFPageRef{323} /XYZ 72 353.823 null] >>}
\MusicPDFWriteObject{16894}{<< /D [\MusicPDFPageRef{324} /XYZ 72 555.776 null] >>}
\MusicPDFWriteObject{16895}{<< /D [\MusicPDFPageRef{325} /XYZ 72 193.51 null] >>}
\MusicPDFWriteObject{16896}{<< /D [\MusicPDFPageRef{326} /XYZ 72 421.144 null] >>}
\MusicPDFWriteObject{16897}{<< /D [\MusicPDFPageRef{326} /XYZ 72 307.922 null] >>}
\MusicPDFWriteObject{16898}{<< /D [\MusicPDFPageRef{331} /XYZ 72 590.853 null] >>}
\MusicPDFWriteObject{16899}{<< /D [\MusicPDFPageRef{333} /XYZ 72 112.384 null] >>}
\MusicPDFWriteObject{16900}{<< /D [\MusicPDFPageRef{335} /XYZ 72 581.679 null] >>}
\MusicPDFWriteObject{16901}{<< /D [\MusicPDFPageRef{335} /XYZ 72 390.121 null] >>}
\MusicPDFWriteObject{16902}{<< /D [\MusicPDFPageRef{336} /XYZ 72 209.962 null] >>}
\MusicPDFWriteObject{16903}{<< /D [\MusicPDFPageRef{339} /XYZ 72 604.811 null] >>}
\MusicPDFWriteObject{16904}{<< /D [\MusicPDFPageRef{340} /XYZ 72 701.27 null] >>}
\MusicPDFWriteObject{16905}{<< /D [\MusicPDFPageRef{340} /XYZ 72 308.105 null] >>}
\MusicPDFWriteObject{16906}{<< /D [\MusicPDFPageRef{341} /XYZ 72 701.27 null] >>}
\MusicPDFWriteObject{16907}{<< /D [\MusicPDFPageRef{343} /XYZ 72 678.427 null] >>}
\MusicPDFWriteObject{16908}{<< /D [\MusicPDFPageRef{343} /XYZ 72 463.771 null] >>}
\MusicPDFWriteObject{16909}{<< /D [\MusicPDFPageRef{343} /XYZ 72 267.924 null] >>}
\MusicPDFWriteObject{16910}{<< /D [\MusicPDFPageRef{344} /XYZ 72 314.152 null] >>}
\MusicPDFWriteObject{16911}{<< /D [\MusicPDFPageRef{350} /XYZ 72 560.94 null] >>}
\MusicPDFWriteObject{16912}{<< /D [\MusicPDFPageRef{350} /XYZ 72 350.688 null] >>}
\MusicPDFWriteObject{16913}{<< /D [\MusicPDFPageRef{351} /XYZ 72 627.558 null] >>}
\MusicPDFWriteObject{16914}{<< /D [\MusicPDFPageRef{352} /XYZ 72 701.27 null] >>}
\MusicPDFWriteObject{16915}{<< /D [\MusicPDFPageRef{354} /XYZ 72 164.987 null] >>}
\MusicPDFWriteObject{16916}{<< /D [\MusicPDFPageRef{355} /XYZ 72 257.262 null] >>}
\MusicPDFWriteObject{16917}{<< /D [\MusicPDFPageRef{357} /XYZ 72 381.937 null] >>}
\MusicPDFWriteObject{16918}{<< /D [\MusicPDFPageRef{357} /XYZ 72 185.878 null] >>}
\MusicPDFWriteObject{16919}{<< /D [\MusicPDFPageRef{359} /XYZ 72 405.443 null] >>}
\MusicPDFWriteObject{16920}{<< /D [\MusicPDFPageRef{360} /XYZ 72 627.857 null] >>}
\MusicPDFWriteObject{16921}{<< /D [\MusicPDFPageRef{364} /XYZ 72 172.843 null] >>}
\MusicPDFWriteObject{16922}{<< /D [\MusicPDFPageRef{366} /XYZ 72 445.346 null] >>}
\MusicPDFWriteObject{16923}{<< /D [\MusicPDFPageRef{367} /XYZ 72 384.945 null] >>}
\MusicPDFWriteObject{16924}{<< /D [\MusicPDFPageRef{368} /XYZ 72 653.88 null] >>}
\MusicPDFWriteObject{16925}{<< /D [\MusicPDFPageRef{372} /XYZ 72 200.377 null] >>}
\MusicPDFWriteObject{16926}{<< /D [\MusicPDFPageRef{373} /XYZ 72 593.55 null] >>}
\MusicPDFWriteObject{16927}{<< /D [\MusicPDFPageRef{374} /XYZ 72 272.559 null] >>}
\MusicPDFWriteObject{16928}{<< /D [\MusicPDFPageRef{378} /XYZ 72 212.351 null] >>}
\MusicPDFWriteObject{16929}{<< /D [\MusicPDFPageRef{379} /XYZ 72 454.119 null] >>}
\MusicPDFWriteObject{16930}{<< /D [\MusicPDFPageRef{382} /XYZ 72 567.261 null] >>}
\MusicPDFWriteObject{16931}{<< /D [\MusicPDFPageRef{384} /XYZ 72 342.197 null] >>}
\MusicPDFWriteObject{16932}{<< /D [\MusicPDFPageRef{385} /XYZ 72 567.701 null] >>}
\MusicPDFWriteObject{16933}{<< /D [\MusicPDFPageRef{385} /XYZ 72 359.185 null] >>}
\MusicPDFWriteObject{16934}{<< /D [\MusicPDFPageRef{386} /XYZ 72 426.203 null] >>}
\MusicPDFWriteObject{16935}{<< /D [\MusicPDFPageRef{386} /XYZ 72 153.943 null] >>}
\MusicPDFWriteObject{16936}{<< /D [\MusicPDFPageRef{387} /XYZ 72 568.004 null] >>}
\MusicPDFWriteObject{16937}{<< /D [\MusicPDFPageRef{387} /XYZ 72 293.749 null] >>}
\MusicPDFWriteObject{16938}{<< /D [\MusicPDFPageRef{391} /XYZ 72 391.745 null] >>}
\MusicPDFWriteObject{16939}{<< /D [\MusicPDFPageRef{396} /XYZ 72 602.949 null] >>}
\MusicPDFWriteObject{16940}{<< /D [\MusicPDFPageRef{396} /XYZ 72 437.043 null] >>}
\MusicPDFWriteObject{16941}{<< /D [\MusicPDFPageRef{396} /XYZ 72 196.684 null] >>}
\MusicPDFWriteObject{16942}{<< /D [\MusicPDFPageRef{397} /XYZ 72 155.636 null] >>}
\MusicPDFWriteObject{16943}{<< /D [\MusicPDFPageRef{397} /XYZ 72 130.597 null] >>}
\MusicPDFWriteObject{16944}{<< /D [\MusicPDFPageRef{399} /XYZ 72 206.715 null] >>}
\MusicPDFWriteObject{16945}{<< /D [\MusicPDFPageRef{400} /XYZ 72 450.617 null] >>}
\MusicPDFWriteObject{16946}{<< /D [\MusicPDFPageRef{401} /XYZ 72 200.904 null] >>}
\MusicPDFWriteObject{16947}{<< /D [\MusicPDFPageRef{402} /XYZ 72 701.27 null] >>}
\MusicPDFWriteObject{16948}{<< /D [\MusicPDFPageRef{402} /XYZ 72 471.922 null] >>}
\MusicPDFWriteObject{16949}{<< /D [\MusicPDFPageRef{402} /XYZ 72 112.384 null] >>}
\MusicPDFWriteObject{16950}{<< /D [\MusicPDFPageRef{403} /XYZ 72 151.836 null] >>}
\MusicPDFWriteObject{16951}{<< /D [\MusicPDFPageRef{404} /XYZ 72 572.694 null] >>}
\MusicPDFWriteObject{16952}{<< /D [\MusicPDFPageRef{404} /XYZ 72 295.474 null] >>}
\MusicPDFWriteObject{16953}{<< /D [\MusicPDFPageRef{406} /XYZ 72 139.089 null] >>}
\MusicPDFWriteObject{16954}{<< /D [\MusicPDFPageRef{407} /XYZ 72 350.663 null] >>}
\MusicPDFWriteObject{16955}{<< /D [\MusicPDFPageRef{407} /XYZ 72 125.957 null] >>}
\MusicPDFWriteObject{16956}{<< /D [\MusicPDFPageRef{408} /XYZ 72 99.233 null] >>}
\MusicPDFWriteObject{16957}{<< /D [\MusicPDFPageRef{409} /XYZ 72 501.826 null] >>}
\MusicPDFWriteObject{16958}{<< /D [\MusicPDFPageRef{409} /XYZ 72 276.332 null] >>}
\MusicPDFWriteObject{16959}{<< /D [\MusicPDFPageRef{412} /XYZ 72 574.264 null] >>}
\MusicPDFWriteObject{16960}{<< /D [\MusicPDFPageRef{416} /XYZ 72 257.151 null] >>}
\MusicPDFWriteObject{16961}{<< /D [\MusicPDFPageRef{417} /XYZ 72 404.538 null] >>}
\MusicPDFWriteObject{16962}{<< /D [\MusicPDFPageRef{418} /XYZ 72 665.551 null] >>}
\MusicPDFWriteObject{16963}{<< /D [\MusicPDFPageRef{418} /XYZ 72 460.729 null] >>}
\MusicPDFWriteObject{16964}{<< /D [\MusicPDFPageRef{412} /XYZ 72 440.349 null] >>}
\MusicPDFWriteObject{16965}{<< /D [\MusicPDFPageRef{414} /XYZ 72 680.39 null] >>}
\MusicPDFWriteObject{16966}{<< /D [\MusicPDFPageRef{414} /XYZ 72 556.439 null] >>}
\MusicPDFWriteObject{16967}{<< /D [\MusicPDFPageRef{414} /XYZ 72 309.809 null] >>}
\MusicPDFWriteObject{16968}{<< /D [\MusicPDFPageRef{415} /XYZ 72 701.27 null] >>}
\MusicPDFWriteObject{16969}{<< /D [\MusicPDFPageRef{415} /XYZ 72 298.775 null] >>}
\MusicPDFWriteObject{16970}{<< /D [\MusicPDFPageRef{415} /XYZ 72 203.761 null] >>}
\MusicPDFWriteObject{16971}{<< /D [\MusicPDFPageRef{416} /XYZ 72 421.11 null] >>}
\MusicPDFWriteObject{16972}{<< /D [\MusicPDFPageRef{419} /XYZ 72 340.574 null] >>}
\MusicPDFWriteObject{16973}{<< /D [\MusicPDFPageRef{419} /XYZ 72 199.057 null] >>}
\MusicPDFWriteObject{16974}{<< /D [\MusicPDFPageRef{419} /XYZ 72 138.685 null] >>}
\MusicPDFWriteObject{16975}{<< /D [\MusicPDFPageRef{420} /XYZ 72 433.032 null] >>}
\MusicPDFWriteObject{16976}{<< /D [\MusicPDFPageRef{421} /XYZ 72 668.677 null] >>}
\MusicPDFWriteObject{16977}{<< /D [\MusicPDFPageRef{422} /XYZ 72 494.704 null] >>}
\MusicPDFWriteObject{16978}{<< /D [\MusicPDFPageRef{422} /XYZ 72 239.347 null] >>}
\MusicPDFWriteObject{16979}{<< /D [\MusicPDFPageRef{425} /XYZ 72 701.27 null] >>}
\MusicPDFWriteObject{16980}{<< /D [\MusicPDFPageRef{425} /XYZ 72 350.328 null] >>}
\MusicPDFWriteObject{16981}{<< /D [\MusicPDFPageRef{425} /XYZ 72 112.384 null] >>}
\MusicPDFWriteObject{16982}{<< /D [\MusicPDFPageRef{426} /XYZ 72 580.118 null] >>}
\MusicPDFWriteObject{16983}{<< /D [\MusicPDFPageRef{427} /XYZ 72 530.864 null] >>}
\MusicPDFWriteObject{16984}{<< /D [\MusicPDFPageRef{427} /XYZ 72 264.218 null] >>}
\MusicPDFWriteObject{16985}{<< /D [\MusicPDFPageRef{428} /XYZ 72 573.047 null] >>}
\MusicPDFWriteObject{16986}{<< /D [\MusicPDFPageRef{428} /XYZ 72 351.717 null] >>}
\MusicPDFWriteObject{16987}{<< /D [\MusicPDFPageRef{429} /XYZ 72 582.593 null] >>}
\MusicPDFWriteObject{16988}{<< /D [\MusicPDFPageRef{539} /XYZ 72 510.507 null] >>}
\MusicPDFWriteObject{16989}{<< /D [\MusicPDFPageRef{539} /XYZ 72 141.18 null] >>}
\MusicPDFWriteObject{16990}{<< /D [\MusicPDFPageRef{540} /XYZ 72 473.694 null] >>}
\MusicPDFWriteObject{16991}{<< /D [\MusicPDFPageRef{540} /XYZ 72 190.356 null] >>}
\MusicPDFWriteObject{16992}{<< /D [\MusicPDFPageRef{541} /XYZ 72 525.16 null] >>}
\MusicPDFWriteObject{16993}{<< /D [\MusicPDFPageRef{541} /XYZ 72 415.041 null] >>}
\MusicPDFWriteObject{16994}{<< /D [\MusicPDFPageRef{541} /XYZ 72 143.941 null] >>}
\MusicPDFWriteObject{16995}{<< /D [\MusicPDFPageRef{542} /XYZ 72 622.043 null] >>}
\MusicPDFWriteObject{16996}{<< /D [\MusicPDFPageRef{542} /XYZ 72 356.041 null] >>}
\MusicPDFWriteObject{16997}{<< /D [\MusicPDFPageRef{542} /XYZ 72 230.26 null] >>}
\MusicPDFWriteObject{16998}{<< /D [\MusicPDFPageRef{542} /XYZ 72 114.604 null] >>}
\MusicPDFWriteObject{16999}{<< /D [\MusicPDFPageRef{543} /XYZ 72 308.399 null] >>}
\MusicPDFWriteObject{17000}{<< /D [\MusicPDFPageRef{544} /XYZ 72 384.696 null] >>}
\MusicPDFWriteObject{17001}{<< /D [\MusicPDFPageRef{544} /XYZ 72 161.592 null] >>}
\MusicPDFWriteObject{17002}{<< /D [\MusicPDFPageRef{545} /XYZ 72 389.886 null] >>}
\MusicPDFWriteObject{17003}{<< /D [\MusicPDFPageRef{545} /XYZ 72 262.698 null] >>}
\MusicPDFWriteObject{17004}{<< /D [\MusicPDFPageRef{545} /XYZ 72 117.447 null] >>}
\MusicPDFWriteObject{17005}{<< /D [\MusicPDFPageRef{546} /XYZ 72 293.733 null] >>}
\MusicPDFWriteObject{17006}{<< /D [\MusicPDFPageRef{546} /XYZ 72 232.035 null] >>}
\MusicPDFWriteObject{17007}{<< /D [\MusicPDFPageRef{547} /XYZ 72 646.087 null] >>}
\MusicPDFWriteObject{17008}{<< /D [\MusicPDFPageRef{547} /XYZ 72 554.71 null] >>}
\MusicPDFWriteObject{17009}{<< /D [\MusicPDFPageRef{547} /XYZ 72 425.611 null] >>}
\MusicPDFWriteObject{17010}{<< /D [\MusicPDFPageRef{547} /XYZ 72 370.661 null] >>}
\MusicPDFWriteObject{17011}{<< /D [\MusicPDFPageRef{547} /XYZ 72 261.735 null] >>}
\MusicPDFWriteObject{17012}{<< /D [\MusicPDFPageRef{547} /XYZ 72 180.484 null] >>}
\MusicPDFWriteObject{17013}{<< /D [\MusicPDFPageRef{547} /XYZ 72 148.811 null] >>}
\MusicPDFWriteObject{17014}{<< /D [\MusicPDFPageRef{548} /XYZ 72 597.897 null] >>}
\MusicPDFWriteObject{17015}{<< /D [\MusicPDFPageRef{552} /XYZ 72 483.945 null] >>}
\MusicPDFWriteObject{17016}{<< /D [\MusicPDFPageRef{553} /XYZ 72 679.21 null] >>}
\MusicPDFWriteObject{17017}{<< /D [\MusicPDFPageRef{553} /XYZ 72 585.511 null] >>}
\MusicPDFWriteObject{17018}{<< /D [\MusicPDFPageRef{553} /XYZ 72 442.234 null] >>}
\MusicPDFWriteObject{17019}{<< /D [\MusicPDFPageRef{548} /XYZ 72 403.25 null] >>}
\MusicPDFWriteObject{17020}{<< /D [\MusicPDFPageRef{548} /XYZ 72 274.43 null] >>}
\MusicPDFWriteObject{17021}{<< /D [\MusicPDFPageRef{548} /XYZ 72 213.534 null] >>}
\MusicPDFWriteObject{17022}{<< /D [\MusicPDFPageRef{549} /XYZ 72 528.968 null] >>}
\MusicPDFWriteObject{17023}{<< /D [\MusicPDFPageRef{549} /XYZ 72 371.921 null] >>}
\MusicPDFWriteObject{17024}{<< /D [\MusicPDFPageRef{549} /XYZ 72 236.667 null] >>}
\MusicPDFWriteObject{17025}{<< /D [\MusicPDFPageRef{550} /XYZ 72 679.815 null] >>}
\MusicPDFWriteObject{17026}{<< /D [\MusicPDFPageRef{552} /XYZ 72 656.465 null] >>}
\MusicPDFWriteObject{17027}{<< /D [\MusicPDFPageRef{554} /XYZ 72 657.16 null] >>}
\MusicPDFWriteObject{17028}{<< /D [\MusicPDFPageRef{554} /XYZ 72 539.53 null] >>}
\MusicPDFWriteObject{17029}{<< /D [\MusicPDFPageRef{554} /XYZ 72 215.274 null] >>}
\MusicPDFWriteObject{17030}{<< /D [\MusicPDFPageRef{555} /XYZ 72 561.389 null] >>}
\MusicPDFWriteObject{17031}{<< /D [\MusicPDFPageRef{555} /XYZ 72 361.268 null] >>}
\MusicPDFWriteObject{17032}{<< /D [\MusicPDFPageRef{555} /XYZ 72 175.113 null] >>}
\MusicPDFWriteObject{17033}{<< /D [\MusicPDFPageRef{558} /XYZ 72 324.895 null] >>}
\MusicPDFWriteObject{17034}{<< /D [\MusicPDFPageRef{558} /XYZ 72 178.137 null] >>}
\MusicPDFWriteObject{17035}{<< /D [\MusicPDFPageRef{556} /XYZ 72 573.761 null] >>}
\MusicPDFWriteObject{17036}{<< /D [\MusicPDFPageRef{556} /XYZ 72 412.099 null] >>}
\MusicPDFWriteObject{17037}{<< /D [\MusicPDFPageRef{556} /XYZ 72 214.937 null] >>}
\MusicPDFWriteObject{17038}{<< /D [\MusicPDFPageRef{557} /XYZ 72 589.52 null] >>}
\MusicPDFWriteObject{17039}{<< /D [\MusicPDFPageRef{557} /XYZ 72 326.395 null] >>}
\MusicPDFWriteObject{17040}{<< /D [\MusicPDFPageRef{557} /XYZ 72 112.384 null] >>}
\MusicPDFWriteObject{17041}{<< /D [\MusicPDFPageRef{558} /XYZ 72 587.504 null] >>}
\MusicPDFWriteObject{17042}{<< /D [\MusicPDFPageRef{558} /XYZ 72 414.609 null] >>}
\MusicPDFWriteObject{17043}{<< /D [\MusicPDFPageRef{559} /XYZ 72 147.994 null] >>}
\MusicPDFWriteObject{17044}{<< /D [\MusicPDFPageRef{560} /XYZ 72 606.747 null] >>}
\MusicPDFWriteObject{17045}{<< /D [\MusicPDFPageRef{560} /XYZ 72 501.812 null] >>}
\MusicPDFWriteObject{17046}{<< /D [\MusicPDFPageRef{560} /XYZ 72 224.598 null] >>}
\MusicPDFWriteObject{17047}{<< /D [\MusicPDFPageRef{560} /XYZ 72 125.534 null] >>}
\MusicPDFWriteObject{17048}{<< /D [\MusicPDFPageRef{561} /XYZ 72 361.066 null] >>}
\MusicPDFWriteObject{17049}{<< /D [\MusicPDFPageRef{562} /XYZ 72 701.27 null] >>}
\MusicPDFWriteObject{17050}{<< /D [\MusicPDFPageRef{562} /XYZ 72 302.174 null] >>}
\MusicPDFWriteObject{17051}{<< /D [\MusicPDFPageRef{564} /XYZ 72 546.38 null] >>}
\MusicPDFWriteObject{17052}{<< /D [\MusicPDFPageRef{566} /XYZ 72 304.81 null] >>}
\MusicPDFWriteObject{17053}{<< /D [\MusicPDFPageRef{568} /XYZ 72 701.27 null] >>}
\MusicPDFWriteObject{17054}{<< /D [\MusicPDFPageRef{568} /XYZ 72 489.123 null] >>}
\MusicPDFWriteObject{17055}{<< /D [\MusicPDFPageRef{569} /XYZ 72 402.819 null] >>}
\MusicPDFWriteObject{17056}{<< /D [\MusicPDFPageRef{572} /XYZ 72 701.27 null] >>}
\MusicPDFWriteObject{17057}{<< /D [\MusicPDFPageRef{574} /XYZ 72 589.718 null] >>}
\MusicPDFWriteObject{17058}{<< /D [\MusicPDFPageRef{575} /XYZ 72 396.307 null] >>}
\MusicPDFWriteObject{17059}{<< /D [\MusicPDFPageRef{576} /XYZ 72 651.383 null] >>}
\MusicPDFWriteObject{17060}{<< /D [\MusicPDFPageRef{576} /XYZ 72 351.18 null] >>}
\MusicPDFWriteObject{17061}{<< /D [\MusicPDFPageRef{576} /XYZ 72 161.525 null] >>}
\MusicPDFWriteObject{17062}{<< /D [\MusicPDFPageRef{577} /XYZ 72 555.995 null] >>}
\MusicPDFWriteObject{17063}{<< /D [\MusicPDFPageRef{577} /XYZ 72 292.915 null] >>}
\MusicPDFWriteObject{17064}{<< /D [\MusicPDFPageRef{578} /XYZ 72 609.675 null] >>}
\MusicPDFWriteObject{17065}{<< /D [\MusicPDFPageRef{578} /XYZ 72 385.025 null] >>}
\MusicPDFWriteObject{17066}{<< /D [\MusicPDFPageRef{569} /XYZ 72 234.771 null] >>}
\MusicPDFWriteObject{17067}{<< /D [\MusicPDFPageRef{570} /XYZ 72 701.27 null] >>}
\MusicPDFWriteObject{17068}{<< /D [\MusicPDFPageRef{570} /XYZ 72 414.487 null] >>}
\MusicPDFWriteObject{17069}{<< /D [\MusicPDFPageRef{570} /XYZ 72 305.019 null] >>}
\MusicPDFWriteObject{17070}{<< /D [\MusicPDFPageRef{570} /XYZ 72 178.137 null] >>}
\MusicPDFWriteObject{17071}{<< /D [\MusicPDFPageRef{571} /XYZ 72 619.643 null] >>}
\MusicPDFWriteObject{17072}{<< /D [\MusicPDFPageRef{571} /XYZ 72 529.245 null] >>}
\MusicPDFWriteObject{17073}{<< /D [\MusicPDFPageRef{571} /XYZ 72 412.545 null] >>}
\MusicPDFWriteObject{17074}{<< /D [\MusicPDFPageRef{578} /XYZ 72 159.493 null] >>}
\MusicPDFWriteObject{17075}{<< /D [\MusicPDFPageRef{579} /XYZ 72 590.81 null] >>}
\MusicPDFWriteObject{17076}{<< /D [\MusicPDFPageRef{579} /XYZ 72 408.219 null] >>}
\MusicPDFWriteObject{17077}{<< /D [\MusicPDFPageRef{579} /XYZ 72 285.404 null] >>}
\MusicPDFWriteObject{17078}{<< /D [\MusicPDFPageRef{580} /XYZ 72 401.6 null] >>}
\MusicPDFWriteObject{17079}{<< /D [\MusicPDFPageRef{581} /XYZ 72 619.51 null] >>}
\MusicPDFWriteObject{17080}{<< /D [\MusicPDFPageRef{582} /XYZ 72 485.263 null] >>}
\MusicPDFWriteObject{17081}{<< /D [\MusicPDFPageRef{582} /XYZ 72 215.034 null] >>}
\MusicPDFWriteObject{17082}{<< /D [\MusicPDFPageRef{583} /XYZ 72 590.751 null] >>}
\MusicPDFWriteObject{17083}{<< /D [\MusicPDFPageRef{585} /XYZ 72 380.584 null] >>}
\MusicPDFWriteObject{17084}{<< /D [\MusicPDFPageRef{589} /XYZ 72 412.881 null] >>}
\MusicPDFWriteObject{17085}{<< /D [\MusicPDFPageRef{590} /XYZ 72 431.513 null] >>}
\MusicPDFWriteObject{17086}{<< /D [\MusicPDFPageRef{591} /XYZ 72 667.848 null] >>}
\MusicPDFWriteObject{17087}{<< /D [\MusicPDFPageRef{591} /XYZ 72 358.748 null] >>}
\MusicPDFWriteObject{17088}{<< /D [\MusicPDFPageRef{593} /XYZ 72 586.422 null] >>}
\MusicPDFWriteObject{17089}{<< /D [\MusicPDFPageRef{600} /XYZ 72 643.408 null] >>}
\MusicPDFWriteObject{17090}{<< /D [\MusicPDFPageRef{593} /XYZ 72 458.488 null] >>}
\MusicPDFWriteObject{17091}{<< /D [\MusicPDFPageRef{593} /XYZ 72 359.35 null] >>}
\MusicPDFWriteObject{17092}{<< /D [\MusicPDFPageRef{593} /XYZ 72 299.665 null] >>}
\MusicPDFWriteObject{17093}{<< /D [\MusicPDFPageRef{594} /XYZ 72 509.192 null] >>}
\MusicPDFWriteObject{17094}{<< /D [\MusicPDFPageRef{594} /XYZ 72 242.974 null] >>}
\MusicPDFWriteObject{17095}{<< /D [\MusicPDFPageRef{595} /XYZ 72 211.537 null] >>}
\MusicPDFWriteObject{17096}{<< /D [\MusicPDFPageRef{597} /XYZ 72 408.194 null] >>}
\MusicPDFWriteObject{17097}{<< /D [\MusicPDFPageRef{598} /XYZ 72 249.269 null] >>}
\MusicPDFWriteObject{17098}{<< /D [\MusicPDFPageRef{600} /XYZ 72 351.412 null] >>}
\MusicPDFWriteObject{17099}{<< /D [\MusicPDFPageRef{600} /XYZ 72 138.685 null] >>}
\MusicPDFWriteObject{17100}{<< /D [\MusicPDFPageRef{601} /XYZ 72 594.553 null] >>}
\MusicPDFWriteObject{17101}{<< /D [\MusicPDFPageRef{601} /XYZ 72 184.941 null] >>}
\MusicPDFWriteObject{17102}{<< /D [\MusicPDFPageRef{603} /XYZ 72 502.052 null] >>}
\MusicPDFWriteObject{17103}{<< /D [\MusicPDFPageRef{603} /XYZ 72 242.13 null] >>}
\MusicPDFWriteObject{17104}{<< /D [\MusicPDFPageRef{604} /XYZ 72 426.155 null] >>}
\MusicPDFWriteObject{17105}{<< /D [\MusicPDFPageRef{604} /XYZ 72 125.534 null] >>}
\MusicPDFWriteObject{17106}{<< /D [\MusicPDFPageRef{605} /XYZ 72 164.987 null] >>}
\MusicPDFWriteObject{17107}{<< /D [\MusicPDFPageRef{607} /XYZ 72 505.382 null] >>}
\MusicPDFWriteObject{17108}{<< /D [\MusicPDFPageRef{608} /XYZ 72 536.682 null] >>}
\MusicPDFWriteObject{17109}{<< /D [\MusicPDFPageRef{608} /XYZ 72 348.134 null] >>}
\MusicPDFWriteObject{17110}{<< /D [\MusicPDFPageRef{608} /XYZ 72 289.785 null] >>}
\MusicPDFWriteObject{17111}{<< /D [\MusicPDFPageRef{610} /XYZ 72 354.081 null] >>}
\MusicPDFWriteObject{17112}{<< /D [\MusicPDFPageRef{611} /XYZ 72 603.149 null] >>}
\MusicPDFWriteObject{17113}{<< /D [\MusicPDFPageRef{611} /XYZ 72 389.178 null] >>}
\MusicPDFWriteObject{17114}{<< /D [\MusicPDFPageRef{611} /XYZ 72 191.288 null] >>}
\MusicPDFWriteObject{17115}{<< /D [\MusicPDFPageRef{612} /XYZ 72 653.258 null] >>}
\MusicPDFWriteObject{17116}{<< /D [\MusicPDFPageRef{616} /XYZ 72 618.037 null] >>}
\MusicPDFWriteObject{17117}{<< /D [\MusicPDFPageRef{616} /XYZ 72 283.736 null] >>}
\MusicPDFWriteObject{17118}{<< /D [\MusicPDFPageRef{617} /XYZ 72 347.715 null] >>}
\MusicPDFWriteObject{17119}{<< /D [\MusicPDFPageRef{619} /XYZ 72 701.27 null] >>}
\MusicPDFWriteObject{17120}{<< /D [\MusicPDFPageRef{619} /XYZ 72 313.314 null] >>}
\MusicPDFWriteObject{17121}{<< /D [\MusicPDFPageRef{620} /XYZ 72 615.219 null] >>}
\MusicPDFWriteObject{17122}{<< /D [\MusicPDFPageRef{620} /XYZ 72 306.936 null] >>}
\MusicPDFWriteObject{17123}{<< /D [\MusicPDFPageRef{621} /XYZ 72 112.384 null] >>}
\MusicPDFWriteObject{17124}{<< /D [\MusicPDFPageRef{622} /XYZ 72 228.767 null] >>}
\MusicPDFWriteObject{17125}{<< /D [\MusicPDFPageRef{623} /XYZ 72 701.27 null] >>}
\MusicPDFWriteObject{17126}{<< /D [\MusicPDFPageRef{623} /XYZ 72 372.045 null] >>}
\MusicPDFWriteObject{17127}{<< /D [\MusicPDFPageRef{624} /XYZ 72 701.27 null] >>}
\MusicPDFWriteObject{17128}{<< /D [\MusicPDFPageRef{624} /XYZ 72 284.887 null] >>}
\MusicPDFWriteObject{17129}{<< /D [\MusicPDFPageRef{626} /XYZ 72 230.649 null] >>}
\MusicPDFWriteObject{17130}{<< /D [\MusicPDFPageRef{627} /XYZ 72 310.821 null] >>}
\MusicPDFWriteObject{17131}{<< /D [\MusicPDFPageRef{627} /XYZ 72 143.748 null] >>}
\MusicPDFWriteObject{17132}{<< /D [\MusicPDFPageRef{628} /XYZ 89.708 433.528 null] >>}
\MusicPDFWriteObject{17133}{<< /D [\MusicPDFPageRef{631} /XYZ 72 263.348 null] >>}
\MusicPDFWriteObject{17134}{<< /D [\MusicPDFPageRef{631} /XYZ 72 125.534 null] >>}
\MusicPDFWriteObject{17135}{<< /D [\MusicPDFPageRef{632} /XYZ 89.708 524.941 null] >>}
\MusicPDFWriteObject{17136}{<< /D [\MusicPDFPageRef{632} /XYZ 72 300.816 null] >>}
\MusicPDFWriteObject{17137}{<< /D [\MusicPDFPageRef{632} /XYZ 72 112.71 null] >>}
\MusicPDFWriteObject{17138}{<< /D [\MusicPDFPageRef{633} /XYZ 72 526.001 null] >>}
\MusicPDFWriteObject{17139}{<< /D [\MusicPDFPageRef{633} /XYZ 72 276.352 null] >>}
\MusicPDFWriteObject{17140}{<< /D [\MusicPDFPageRef{636} /XYZ 72 537.016 null] >>}
\MusicPDFWriteObject{17141}{<< /D [\MusicPDFPageRef{637} /XYZ 72 535.549 null] >>}
\MusicPDFWriteObject{17142}{<< /D [\MusicPDFPageRef{637} /XYZ 72 304.707 null] >>}
\MusicPDFWriteObject{17143}{<< /D [\MusicPDFPageRef{638} /XYZ 72 477.209 null] >>}
\MusicPDFWriteObject{17144}{<< /D [\MusicPDFPageRef{639} /XYZ 72 255.144 null] >>}
\MusicPDFWriteObject{17145}{<< /D [\MusicPDFPageRef{640} /XYZ 72 634.812 null] >>}
\MusicPDFWriteObject{17146}{<< /D [\MusicPDFPageRef{641} /XYZ 72 224.691 null] >>}
\MusicPDFWriteObject{17147}{<< /D [\MusicPDFPageRef{643} /XYZ 72 474.618 null] >>}
\MusicPDFWriteObject{17148}{<< /D [\MusicPDFPageRef{644} /XYZ 72 682.266 null] >>}
\MusicPDFWriteObject{17149}{<< /D [\MusicPDFPageRef{645} /XYZ 72 252.725 null] >>}
\MusicPDFWriteObject{17150}{<< /D [\MusicPDFPageRef{649} /XYZ 72 639.825 null] >>}
\MusicPDFWriteObject{17151}{<< /D [\MusicPDFPageRef{649} /XYZ 72 292.426 null] >>}
\MusicPDFWriteObject{17152}{<< /D [\MusicPDFPageRef{651} /XYZ 72 300.12 null] >>}
\MusicPDFWriteObject{17153}{<< /D [\MusicPDFPageRef{652} /XYZ 72 524.076 null] >>}
\MusicPDFWriteObject{17154}{<< /D [\MusicPDFPageRef{653} /XYZ 72 544.242 null] >>}
\MusicPDFWriteObject{17155}{<< /D [\MusicPDFPageRef{653} /XYZ 72 390.187 null] >>}
\MusicPDFWriteObject{17156}{<< /D [\MusicPDFPageRef{653} /XYZ 72 281.92 null] >>}
\MusicPDFWriteObject{17157}{<< /D [\MusicPDFPageRef{657} /XYZ 72 169.665 null] >>}
\MusicPDFWriteObject{17158}{<< /D [\MusicPDFPageRef{658} /XYZ 72 576.554 null] >>}
\MusicPDFWriteObject{17159}{<< /D [\MusicPDFPageRef{658} /XYZ 72 449.662 null] >>}
\MusicPDFWriteObject{17160}{<< /D [\MusicPDFPageRef{659} /XYZ 72 536.628 null] >>}
\MusicPDFWriteObject{17161}{<< /D [\MusicPDFPageRef{659} /XYZ 72 191.288 null] >>}
\MusicPDFWriteObject{17162}{<< /D [\MusicPDFPageRef{660} /XYZ 72 531.759 null] >>}
\MusicPDFWriteObject{17163}{<< /D [\MusicPDFPageRef{660} /XYZ 72 262.893 null] >>}
\MusicPDFWriteObject{17164}{<< /D [\MusicPDFPageRef{661} /XYZ 72 607.822 null] >>}
\MusicPDFWriteObject{17165}{<< /D [\MusicPDFPageRef{661} /XYZ 72 573.485 null] >>}
\MusicPDFWriteObject{17166}{<< /D [\MusicPDFPageRef{661} /XYZ 72 412.706 null] >>}
\MusicPDFWriteObject{17167}{<< /D [\MusicPDFPageRef{661} /XYZ 72 265.078 null] >>}
\MusicPDFWriteObject{17168}{<< /D [\MusicPDFPageRef{661} /XYZ 72 151.836 null] >>}
\MusicPDFWriteObject{17169}{<< /D [\MusicPDFPageRef{662} /XYZ 72 443.947 null] >>}
\MusicPDFWriteObject{17170}{<< /D [\MusicPDFPageRef{663} /XYZ 72 627.3 null] >>}
\MusicPDFWriteObject{17171}{<< /D [\MusicPDFPageRef{664} /XYZ 72 701.27 null] >>}
\MusicPDFWriteObject{17172}{<< /D [\MusicPDFPageRef{668} /XYZ 72 655.637 null] >>}
\MusicPDFWriteObject{17173}{<< /D [\MusicPDFPageRef{668} /XYZ 72 498.943 null] >>}
\MusicPDFWriteObject{17174}{<< /D [\MusicPDFPageRef{668} /XYZ 72 301.031 null] >>}
\MusicPDFWriteObject{17175}{<< /D [\MusicPDFPageRef{668} /XYZ 72 242.573 null] >>}
\MusicPDFWriteObject{17176}{<< /D [\MusicPDFPageRef{664} /XYZ 72 280.122 null] >>}
\MusicPDFWriteObject{17177}{<< /D [\MusicPDFPageRef{664} /XYZ 72 159.529 null] >>}
\MusicPDFWriteObject{17178}{<< /D [\MusicPDFPageRef{665} /XYZ 72 701.27 null] >>}
\MusicPDFWriteObject{17179}{<< /D [\MusicPDFPageRef{665} /XYZ 72 316.865 null] >>}
\MusicPDFWriteObject{17180}{<< /D [\MusicPDFPageRef{666} /XYZ 72 528.943 null] >>}
\MusicPDFWriteObject{17181}{<< /D [\MusicPDFPageRef{666} /XYZ 72 220.957 null] >>}
\MusicPDFWriteObject{17182}{<< /D [\MusicPDFPageRef{667} /XYZ 72 670.326 null] >>}
\MusicPDFWriteObject{17183}{<< /D [\MusicPDFPageRef{667} /XYZ 72 294.323 null] >>}
\MusicPDFWriteObject{17184}{<< /D [\MusicPDFPageRef{669} /XYZ 72 408.951 null] >>}
\MusicPDFWriteObject{17185}{<< /D [\MusicPDFPageRef{670} /XYZ 72 602.64 null] >>}
\MusicPDFWriteObject{17186}{<< /D [\MusicPDFPageRef{670} /XYZ 72 188.897 null] >>}
\MusicPDFWriteObject{17187}{<< /D [\MusicPDFPageRef{671} /XYZ 72 701.27 null] >>}
\MusicPDFWriteObject{17188}{<< /D [\MusicPDFPageRef{671} /XYZ 72 517.378 null] >>}
\MusicPDFWriteObject{17189}{<< /D [\MusicPDFPageRef{671} /XYZ 72 484.188 null] >>}
\MusicPDFWriteObject{17190}{<< /D [\MusicPDFPageRef{671} /XYZ 72 316.473 null] >>}
\MusicPDFWriteObject{17191}{<< /D [\MusicPDFPageRef{671} /XYZ 72 155.696 null] >>}
\MusicPDFWriteObject{17192}{<< /D [\MusicPDFPageRef{672} /XYZ 72 540.021 null] >>}
\MusicPDFWriteObject{17193}{<< /D [\MusicPDFPageRef{672} /XYZ 72 441.981 null] >>}
\MusicPDFWriteObject{17194}{<< /D [\MusicPDFPageRef{672} /XYZ 72 198.802 null] >>}
\MusicPDFWriteObject{17195}{<< /D [\MusicPDFPageRef{673} /XYZ 72 117.447 null] >>}
\MusicPDFWriteObject{17196}{<< /D [\MusicPDFPageRef{675} /XYZ 72 588.08 null] >>}
\MusicPDFWriteObject{17197}{<< /D [\MusicPDFPageRef{676} /XYZ 72 445.916 null] >>}
\MusicPDFWriteObject{17198}{<< /D [\MusicPDFPageRef{681} /XYZ 72 605.147 null] >>}
\MusicPDFWriteObject{17199}{<< /D [\MusicPDFPageRef{681} /XYZ 72 151.958 null] >>}
\MusicPDFWriteObject{17200}{<< /D [\MusicPDFPageRef{682} /XYZ 72 224.089 null] >>}
\MusicPDFWriteObject{17201}{<< /D [\MusicPDFPageRef{682} /XYZ 72 134.453 null] >>}
\MusicPDFWriteObject{17202}{<< /D [\MusicPDFPageRef{684} /XYZ 72 701.27 null] >>}
\MusicPDFWriteObject{17203}{<< /D [\MusicPDFPageRef{677} /XYZ 72 701.27 null] >>}
\MusicPDFWriteObject{17204}{<< /D [\MusicPDFPageRef{677} /XYZ 72 478.609 null] >>}
\MusicPDFWriteObject{17205}{<< /D [\MusicPDFPageRef{677} /XYZ 72 261.025 null] >>}
\MusicPDFWriteObject{17206}{<< /D [\MusicPDFPageRef{678} /XYZ 72 612.738 null] >>}
\MusicPDFWriteObject{17207}{<< /D [\MusicPDFPageRef{678} /XYZ 72 334.672 null] >>}
\MusicPDFWriteObject{17208}{<< /D [\MusicPDFPageRef{679} /XYZ 72 248.574 null] >>}
\MusicPDFWriteObject{17209}{<< /D [\MusicPDFPageRef{679} /XYZ 72 133.744 null] >>}
\MusicPDFWriteObject{17210}{<< /D [\MusicPDFPageRef{680} /XYZ 72 449.66 null] >>}
\MusicPDFWriteObject{17211}{<< /D [\MusicPDFPageRef{685} /XYZ 72 250.947 null] >>}
\MusicPDFWriteObject{17212}{<< /D [\MusicPDFPageRef{689} /XYZ 72 668.638 null] >>}
\MusicPDFWriteObject{17213}{<< /D [\MusicPDFPageRef{689} /XYZ 72 561.829 null] >>}
\MusicPDFWriteObject{17214}{<< /D [\MusicPDFPageRef{689} /XYZ 72 509.714 null] >>}
\MusicPDFWriteObject{17215}{<< /D [\MusicPDFPageRef{689} /XYZ 72 422.387 null] >>}
\MusicPDFWriteObject{17216}{<< /D [\MusicPDFPageRef{689} /XYZ 72 350.302 null] >>}
\MusicPDFWriteObject{17217}{<< /D [\MusicPDFPageRef{689} /XYZ 72 199.031 null] >>}
\MusicPDFWriteObject{17218}{<< /D [\MusicPDFPageRef{690} /XYZ 72 559.776 null] >>}
\MusicPDFWriteObject{17219}{<< /D [\MusicPDFPageRef{690} /XYZ 72 455.641 null] >>}
\MusicPDFWriteObject{17220}{<< /D [\MusicPDFPageRef{690} /XYZ 72 314.146 null] >>}
\MusicPDFWriteObject{17221}{<< /D [\MusicPDFPageRef{691} /XYZ 72 596.887 null] >>}
\MusicPDFWriteObject{17222}{<< /D [\MusicPDFPageRef{686} /XYZ 72 595.844 null] >>}
\MusicPDFWriteObject{17223}{<< /D [\MusicPDFPageRef{686} /XYZ 72 169.355 null] >>}
\MusicPDFWriteObject{17224}{<< /D [\MusicPDFPageRef{687} /XYZ 72 701.27 null] >>}
\MusicPDFWriteObject{17225}{<< /D [\MusicPDFPageRef{687} /XYZ 72 412.164 null] >>}
\MusicPDFWriteObject{17226}{<< /D [\MusicPDFPageRef{687} /XYZ 72 269.485 null] >>}
\MusicPDFWriteObject{17227}{<< /D [\MusicPDFPageRef{688} /XYZ 72 631.571 null] >>}
\MusicPDFWriteObject{17228}{<< /D [\MusicPDFPageRef{688} /XYZ 72 386.917 null] >>}
\MusicPDFWriteObject{17229}{<< /D [\MusicPDFPageRef{688} /XYZ 72 251.71 null] >>}
\MusicPDFWriteObject{17230}{<< /D [\MusicPDFPageRef{691} /XYZ 72 224.704 null] >>}
\MusicPDFWriteObject{17231}{<< /D [\MusicPDFPageRef{692} /XYZ 72 278.826 null] >>}
\MusicPDFWriteObject{17232}{<< /D [\MusicPDFPageRef{693} /XYZ 72 277.887 null] >>}
\MusicPDFWriteObject{17233}{<< /D [\MusicPDFPageRef{694} /XYZ 72 601.415 null] >>}
\MusicPDFWriteObject{17234}{<< /D [\MusicPDFPageRef{694} /XYZ 72 353.505 null] >>}
\MusicPDFWriteObject{17235}{<< /D [\MusicPDFPageRef{695} /XYZ 72 509.855 null] >>}
\MusicPDFWriteObject{17236}{<< /D [\MusicPDFPageRef{697} /XYZ 72 624.656 null] >>}
\MusicPDFWriteObject{17237}{<< /D [\MusicPDFPageRef{697} /XYZ 72 291.484 null] >>}
\MusicPDFWriteObject{17238}{<< /D [\MusicPDFPageRef{697} /XYZ 72 169.847 null] >>}
\MusicPDFWriteObject{17239}{<< /D [\MusicPDFPageRef{698} /XYZ 72 669.38 null] >>}
\MusicPDFWriteObject{17240}{<< /D [\MusicPDFPageRef{698} /XYZ 72 340.603 null] >>}
\MusicPDFWriteObject{17241}{<< /D [\MusicPDFPageRef{699} /XYZ 72 556.832 null] >>}
\MusicPDFWriteObject{17242}{<< /D [\MusicPDFPageRef{699} /XYZ 72 296.725 null] >>}
\MusicPDFWriteObject{17243}{<< /D [\MusicPDFPageRef{699} /XYZ 72 112.384 null] >>}
\MusicPDFWriteObject{17244}{<< /D [\MusicPDFPageRef{701} /XYZ 72 486.811 null] >>}
\MusicPDFWriteObject{17245}{<< /D [\MusicPDFPageRef{702} /XYZ 72 408.089 null] >>}
\MusicPDFWriteObject{17246}{<< /D [\MusicPDFPageRef{703} /XYZ 72 234.881 null] >>}
\MusicPDFWriteObject{17247}{<< /D [\MusicPDFPageRef{704} /XYZ 72 411.983 null] >>}
\MusicPDFWriteObject{17248}{<< /D [\MusicPDFPageRef{704} /XYZ 72 201.318 null] >>}
\MusicPDFWriteObject{17249}{<< /D [\MusicPDFPageRef{705} /XYZ 72 534.242 null] >>}
\MusicPDFWriteObject{17250}{<< /D [\MusicPDFPageRef{705} /XYZ 72 125.534 null] >>}
\MusicPDFWriteObject{17251}{<< /D [\MusicPDFPageRef{707} /XYZ 72 559.815 null] >>}
\MusicPDFWriteObject{17252}{<< /D [\MusicPDFPageRef{708} /XYZ 72 520.641 null] >>}
\MusicPDFWriteObject{17253}{<< /D [\MusicPDFPageRef{709} /XYZ 72 335.054 null] >>}
\MusicPDFWriteObject{17254}{<< /D [\MusicPDFPageRef{710} /XYZ 72 701.27 null] >>}
\MusicPDFWriteObject{17255}{<< /D [\MusicPDFPageRef{710} /XYZ 72 482.814 null] >>}
\MusicPDFWriteObject{17256}{<< /D [\MusicPDFPageRef{710} /XYZ 72 379.369 null] >>}
\MusicPDFWriteObject{17257}{<< /D [\MusicPDFPageRef{710} /XYZ 72 191.288 null] >>}
\MusicPDFWriteObject{17258}{<< /D [\MusicPDFPageRef{711} /XYZ 72 478.21 null] >>}
\MusicPDFWriteObject{17259}{<< /D [\MusicPDFPageRef{715} /XYZ 72 335.473 null] >>}
\MusicPDFWriteObject{17260}{<< /D [\MusicPDFPageRef{715} /XYZ 72 197.085 null] >>}
\MusicPDFWriteObject{17261}{<< /D [\MusicPDFPageRef{716} /XYZ 72 350.659 null] >>}
\MusicPDFWriteObject{17262}{<< /D [\MusicPDFPageRef{716} /XYZ 72 178.137 null] >>}
\MusicPDFWriteObject{17263}{<< /D [\MusicPDFPageRef{711} /XYZ 72 125.534 null] >>}
\MusicPDFWriteObject{17264}{<< /D [\MusicPDFPageRef{712} /XYZ 72 346.018 null] >>}
\MusicPDFWriteObject{17265}{<< /D [\MusicPDFPageRef{712} /XYZ 72 112.384 null] >>}
\MusicPDFWriteObject{17266}{<< /D [\MusicPDFPageRef{713} /XYZ 72 172.788 null] >>}
\MusicPDFWriteObject{17267}{<< /D [\MusicPDFPageRef{714} /XYZ 72 581.979 null] >>}
\MusicPDFWriteObject{17268}{<< /D [\MusicPDFPageRef{714} /XYZ 72 448.862 null] >>}
\MusicPDFWriteObject{17269}{<< /D [\MusicPDFPageRef{714} /XYZ 72 263.204 null] >>}
\MusicPDFWriteObject{17270}{<< /D [\MusicPDFPageRef{715} /XYZ 72 429.898 null] >>}
\MusicPDFWriteObject{17271}{<< /D [\MusicPDFPageRef{717} /XYZ 72 666.745 null] >>}
\MusicPDFWriteObject{17272}{<< /D [\MusicPDFPageRef{717} /XYZ 72 229.168 null] >>}
\MusicPDFWriteObject{17273}{<< /D [\MusicPDFPageRef{718} /XYZ 72 515.667 null] >>}
\MusicPDFWriteObject{17274}{<< /D [\MusicPDFPageRef{718} /XYZ 72 362.361 null] >>}
\MusicPDFWriteObject{17275}{<< /D [\MusicPDFPageRef{718} /XYZ 72 236.549 null] >>}
\MusicPDFWriteObject{17276}{<< /D [\MusicPDFPageRef{721} /XYZ 72 604.993 null] >>}
\MusicPDFWriteObject{17277}{<< /D [\MusicPDFPageRef{725} /XYZ 72 701.27 null] >>}
\MusicPDFWriteObject{17278}{<< /D [\MusicPDFPageRef{725} /XYZ 72 627.459 null] >>}
\MusicPDFWriteObject{17279}{<< /D [\MusicPDFPageRef{725} /XYZ 72 593.1 null] >>}
\MusicPDFWriteObject{17280}{<< /D [\MusicPDFPageRef{726} /XYZ 72 552.93 null] >>}
\MusicPDFWriteObject{17281}{<< /D [\MusicPDFPageRef{728} /XYZ 72 654.192 null] >>}
\MusicPDFWriteObject{17282}{<< /D [\MusicPDFPageRef{728} /XYZ 72 468.505 null] >>}
\MusicPDFWriteObject{17283}{<< /D [\MusicPDFPageRef{728} /XYZ 72 125.534 null] >>}
\MusicPDFWriteObject{17284}{<< /D [\MusicPDFPageRef{730} /XYZ 72 513.079 null] >>}
\MusicPDFWriteObject{17285}{<< /D [\MusicPDFPageRef{730} /XYZ 72 355.972 null] >>}
\MusicPDFWriteObject{17286}{<< /D [\MusicPDFPageRef{731} /XYZ 72 701.27 null] >>}
\MusicPDFWriteObject{17287}{<< /D [\MusicPDFPageRef{721} /XYZ 72 459.137 null] >>}
\MusicPDFWriteObject{17288}{<< /D [\MusicPDFPageRef{731} /XYZ 72 559.919 null] >>}
\MusicPDFWriteObject{17289}{<< /D [\MusicPDFPageRef{731} /XYZ 72 138.685 null] >>}
\MusicPDFWriteObject{17290}{<< /D [\MusicPDFPageRef{732} /XYZ 72 272.841 null] >>}
\MusicPDFWriteObject{17291}{<< /D [\MusicPDFPageRef{733} /XYZ 72 252.352 null] >>}
\MusicPDFWriteObject{17292}{<< /D [\MusicPDFPageRef{734} /XYZ 72 610.16 null] >>}
\MusicPDFWriteObject{17293}{<< /D [\MusicPDFPageRef{734} /XYZ 72 251.04 null] >>}
\MusicPDFWriteObject{17294}{<< /D [\MusicPDFPageRef{735} /XYZ 72 533.885 null] >>}
\MusicPDFWriteObject{17295}{<< /D [\MusicPDFPageRef{735} /XYZ 72 321.983 null] >>}
\MusicPDFWriteObject{17296}{<< /D [\MusicPDFPageRef{735} /XYZ 72 206.792 null] >>}
\MusicPDFWriteObject{17297}{<< /D [\MusicPDFPageRef{736} /XYZ 72 667.202 null] >>}
\MusicPDFWriteObject{17298}{<< /D [\MusicPDFPageRef{721} /XYZ 72 347.616 null] >>}
\MusicPDFWriteObject{17299}{<< /D [\MusicPDFPageRef{736} /XYZ 72 519.751 null] >>}
\MusicPDFWriteObject{17300}{<< /D [\MusicPDFPageRef{722} /XYZ 72 420.308 null] >>}
\MusicPDFWriteObject{17301}{<< /D [\MusicPDFPageRef{722} /XYZ 72 385.525 null] >>}
\MusicPDFWriteObject{17302}{<< /D [\MusicPDFPageRef{723} /XYZ 72 413.64 null] >>}
\MusicPDFWriteObject{17303}{<< /D [\MusicPDFPageRef{723} /XYZ 72 229.847 null] >>}
\MusicPDFWriteObject{17304}{<< /D [\MusicPDFPageRef{724} /XYZ 72 701.27 null] >>}
\MusicPDFWriteObject{17305}{<< /D [\MusicPDFPageRef{724} /XYZ 72 217.689 null] >>}
\MusicPDFWriteObject{17306}{<< /D [\MusicPDFPageRef{737} /XYZ 72 335.411 null] >>}
\MusicPDFWriteObject{17307}{<< /D [\MusicPDFPageRef{737} /XYZ 72 169.674 null] >>}
\MusicPDFWriteObject{17308}{<< /D [\MusicPDFPageRef{738} /XYZ 72 701.27 null] >>}
\MusicPDFWriteObject{17309}{<< /D [\MusicPDFPageRef{738} /XYZ 72 604.695 null] >>}
\MusicPDFWriteObject{17310}{<< /D [\MusicPDFPageRef{738} /XYZ 72 410.898 null] >>}
\MusicPDFWriteObject{17311}{<< /D [\MusicPDFPageRef{738} /XYZ 72 353.775 null] >>}
\MusicPDFWriteObject{17312}{<< /D [\MusicPDFPageRef{738} /XYZ 72 99.233 null] >>}
\MusicPDFWriteObject{17313}{<< /D [\MusicPDFPageRef{739} /XYZ 72 382.236 null] >>}
\MusicPDFWriteObject{17314}{<< /D [\MusicPDFPageRef{739} /XYZ 72 125.534 null] >>}
\MusicPDFWriteObject{17315}{<< /D [\MusicPDFPageRef{740} /XYZ 72 436.905 null] >>}
\MusicPDFWriteObject{17316}{<< /D [\MusicPDFPageRef{741} /XYZ 72 701.27 null] >>}
\MusicPDFWriteObject{17317}{<< /D [\MusicPDFPageRef{741} /XYZ 72 481.748 null] >>}
\MusicPDFWriteObject{17318}{<< /D [\MusicPDFPageRef{741} /XYZ 72 268.381 null] >>}
\MusicPDFWriteObject{17319}{<< /D [\MusicPDFPageRef{746} /XYZ 72 593.274 null] >>}
\MusicPDFWriteObject{17320}{<< /D [\MusicPDFPageRef{746} /XYZ 72 159.846 null] >>}
\MusicPDFWriteObject{17321}{<< /D [\MusicPDFPageRef{784} /XYZ 72 701.27 null] >>}
\MusicPDFWriteObject{17322}{<< /D [\MusicPDFPageRef{791} /XYZ 72 257.921 null] >>}
\MusicPDFWriteObject{17323}{<< /D [\MusicPDFPageRef{784} /XYZ 72 411.674 null] >>}
\MusicPDFWriteObject{17324}{<< /D [\MusicPDFPageRef{785} /XYZ 72 384.717 null] >>}
\MusicPDFWriteObject{17325}{<< /D [\MusicPDFPageRef{785} /XYZ 72 184.846 null] >>}
\MusicPDFWriteObject{17326}{<< /D [\MusicPDFPageRef{786} /XYZ 72 275.808 null] >>}
\MusicPDFWriteObject{17327}{<< /D [\MusicPDFPageRef{787} /XYZ 72 610.186 null] >>}
\MusicPDFWriteObject{17328}{<< /D [\MusicPDFPageRef{789} /XYZ 72 394.351 null] >>}
\MusicPDFWriteObject{17329}{<< /D [\MusicPDFPageRef{789} /XYZ 72 273.897 null] >>}
\MusicPDFWriteObject{17330}{<< /D [\MusicPDFPageRef{790} /XYZ 72 303.951 null] >>}
\MusicPDFWriteObject{17331}{<< /D [\MusicPDFPageRef{747} /XYZ 72 332.622 null] >>}
\MusicPDFWriteObject{17332}{<< /D [\MusicPDFPageRef{747} /XYZ 72 184.92 null] >>}
\MusicPDFWriteObject{17333}{<< /D [\MusicPDFPageRef{747} /XYZ 72 154.48 null] >>}
\MusicPDFWriteObject{17334}{<< /D [\MusicPDFPageRef{748} /XYZ 72 642.332 null] >>}
\MusicPDFWriteObject{17335}{<< /D [\MusicPDFPageRef{748} /XYZ 72 275.178 null] >>}
\MusicPDFWriteObject{17336}{<< /D [\MusicPDFPageRef{748} /XYZ 72 112.384 null] >>}
\MusicPDFWriteObject{17337}{<< /D [\MusicPDFPageRef{749} /XYZ 72 625.793 null] >>}
\MusicPDFWriteObject{17338}{<< /D [\MusicPDFPageRef{749} /XYZ 72 405.031 null] >>}
\MusicPDFWriteObject{17339}{<< /D [\MusicPDFPageRef{750} /XYZ 72 510.698 null] >>}
\MusicPDFWriteObject{17340}{<< /D [\MusicPDFPageRef{750} /XYZ 72 343.921 null] >>}
\MusicPDFWriteObject{17341}{<< /D [\MusicPDFPageRef{750} /XYZ 72 293.761 null] >>}
\MusicPDFWriteObject{17342}{<< /D [\MusicPDFPageRef{751} /XYZ 72 572.056 null] >>}
\MusicPDFWriteObject{17343}{<< /D [\MusicPDFPageRef{751} /XYZ 72 99.233 null] >>}
\MusicPDFWriteObject{17344}{<< /D [\MusicPDFPageRef{752} /XYZ 72 198.474 null] >>}
\MusicPDFWriteObject{17345}{<< /D [\MusicPDFPageRef{753} /XYZ 72 639.011 null] >>}
\MusicPDFWriteObject{17346}{<< /D [\MusicPDFPageRef{753} /XYZ 72 269.16 null] >>}
\MusicPDFWriteObject{17347}{<< /D [\MusicPDFPageRef{757} /XYZ 72 307.029 null] >>}
\MusicPDFWriteObject{17348}{<< /D [\MusicPDFPageRef{757} /XYZ 72 274.569 null] >>}
\MusicPDFWriteObject{17349}{<< /D [\MusicPDFPageRef{757} /XYZ 72 191.013 null] >>}
\MusicPDFWriteObject{17350}{<< /D [\MusicPDFPageRef{758} /XYZ 72 534.784 null] >>}
\MusicPDFWriteObject{17351}{<< /D [\MusicPDFPageRef{758} /XYZ 72 295.621 null] >>}
\MusicPDFWriteObject{17352}{<< /D [\MusicPDFPageRef{758} /XYZ 72 207.924 null] >>}
\MusicPDFWriteObject{17353}{<< /D [\MusicPDFPageRef{760} /XYZ 72 656.767 null] >>}
\MusicPDFWriteObject{17354}{<< /D [\MusicPDFPageRef{760} /XYZ 72 289.492 null] >>}
\MusicPDFWriteObject{17355}{<< /D [\MusicPDFPageRef{765} /XYZ 72 701.27 null] >>}
\MusicPDFWriteObject{17356}{<< /D [\MusicPDFPageRef{760} /XYZ 72 130.597 null] >>}
\MusicPDFWriteObject{17357}{<< /D [\MusicPDFPageRef{762} /XYZ 72 701.27 null] >>}
\MusicPDFWriteObject{17358}{<< /D [\MusicPDFPageRef{762} /XYZ 72 384.794 null] >>}
\MusicPDFWriteObject{17359}{<< /D [\MusicPDFPageRef{763} /XYZ 72 448.462 null] >>}
\MusicPDFWriteObject{17360}{<< /D [\MusicPDFPageRef{763} /XYZ 72 315.542 null] >>}
\MusicPDFWriteObject{17361}{<< /D [\MusicPDFPageRef{763} /XYZ 72 112.384 null] >>}
\MusicPDFWriteObject{17362}{<< /D [\MusicPDFPageRef{764} /XYZ 72 382.62 null] >>}
\MusicPDFWriteObject{17363}{<< /D [\MusicPDFPageRef{764} /XYZ 72 195.544 null] >>}
\MusicPDFWriteObject{17364}{<< /D [\MusicPDFPageRef{765} /XYZ 72 264.685 null] >>}
\MusicPDFWriteObject{17365}{<< /D [\MusicPDFPageRef{768} /XYZ 72 364.481 null] >>}
\MusicPDFWriteObject{17366}{<< /D [\MusicPDFPageRef{768} /XYZ 72 257.95 null] >>}
\MusicPDFWriteObject{17367}{<< /D [\MusicPDFPageRef{768} /XYZ 72 197.923 null] >>}
\MusicPDFWriteObject{17368}{<< /D [\MusicPDFPageRef{769} /XYZ 72 581.442 null] >>}
\MusicPDFWriteObject{17369}{<< /D [\MusicPDFPageRef{769} /XYZ 72 286.312 null] >>}
\MusicPDFWriteObject{17370}{<< /D [\MusicPDFPageRef{770} /XYZ 72 463.927 null] >>}
\MusicPDFWriteObject{17371}{<< /D [\MusicPDFPageRef{770} /XYZ 72 287.484 null] >>}
\MusicPDFWriteObject{17372}{<< /D [\MusicPDFPageRef{770} /XYZ 72 210.536 null] >>}
\MusicPDFWriteObject{17373}{<< /D [\MusicPDFPageRef{772} /XYZ 72 701.27 null] >>}
\MusicPDFWriteObject{17374}{<< /D [\MusicPDFPageRef{773} /XYZ 72 409.126 null] >>}
\MusicPDFWriteObject{17375}{<< /D [\MusicPDFPageRef{765} /XYZ 72 180.176 null] >>}
\MusicPDFWriteObject{17376}{<< /D [\MusicPDFPageRef{776} /XYZ 72 701.27 null] >>}
\MusicPDFWriteObject{17377}{<< /D [\MusicPDFPageRef{766} /XYZ 72 654.054 null] >>}
\MusicPDFWriteObject{17378}{<< /D [\MusicPDFPageRef{766} /XYZ 72 444.211 null] >>}
\MusicPDFWriteObject{17379}{<< /D [\MusicPDFPageRef{766} /XYZ 72 225.765 null] >>}
\MusicPDFWriteObject{17380}{<< /D [\MusicPDFPageRef{767} /XYZ 72 701.27 null] >>}
\MusicPDFWriteObject{17381}{<< /D [\MusicPDFPageRef{767} /XYZ 72 431.531 null] >>}
\MusicPDFWriteObject{17382}{<< /D [\MusicPDFPageRef{768} /XYZ 72 701.27 null] >>}
\MusicPDFWriteObject{17383}{<< /D [\MusicPDFPageRef{768} /XYZ 72 431.834 null] >>}
\MusicPDFWriteObject{17384}{<< /D [\MusicPDFPageRef{776} /XYZ 72 554.402 null] >>}
\MusicPDFWriteObject{17385}{<< /D [\MusicPDFPageRef{778} /XYZ 72 701.27 null] >>}
\MusicPDFWriteObject{17386}{<< /D [\MusicPDFPageRef{779} /XYZ 72 469.413 null] >>}
\MusicPDFWriteObject{17387}{<< /D [\MusicPDFPageRef{779} /XYZ 72 390.114 null] >>}
\MusicPDFWriteObject{17388}{<< /D [\MusicPDFPageRef{780} /XYZ 72 575.715 null] >>}
\MusicPDFWriteObject{17389}{<< /D [\MusicPDFPageRef{780} /XYZ 72 224.891 null] >>}
\MusicPDFWriteObject{17390}{<< /D [\MusicPDFPageRef{781} /XYZ 72 571.021 null] >>}
\MusicPDFWriteObject{17391}{<< /D [\MusicPDFPageRef{781} /XYZ 72 323.674 null] >>}
\MusicPDFWriteObject{17392}{<< /D [\MusicPDFPageRef{782} /XYZ 72 580.334 null] >>}
\MusicPDFWriteObject{17393}{<< /D [\MusicPDFPageRef{793} /XYZ 72 666.592 null] >>}
\MusicPDFWriteObject{17394}{<< /D [\MusicPDFPageRef{793} /XYZ 72 468.252 null] >>}
\MusicPDFWriteObject{17395}{<< /D [\MusicPDFPageRef{793} /XYZ 72 314.619 null] >>}
\MusicPDFWriteObject{17396}{<< /D [\MusicPDFPageRef{793} /XYZ 72 112.499 null] >>}
\MusicPDFWriteObject{17397}{<< /D [\MusicPDFPageRef{794} /XYZ 72 320.21 null] >>}
\MusicPDFWriteObject{17398}{<< /D [\MusicPDFPageRef{796} /XYZ 72 538.354 null] >>}
\MusicPDFWriteObject{17399}{<< /D [\MusicPDFPageRef{797} /XYZ 72 501.221 null] >>}
\MusicPDFWriteObject{17400}{<< /D [\MusicPDFPageRef{797} /XYZ 72 242.987 null] >>}
\MusicPDFWriteObject{17401}{<< /D [\MusicPDFPageRef{799} /XYZ 72 488.044 null] >>}
\MusicPDFWriteObject{17402}{<< /D [\MusicPDFPageRef{800} /XYZ 72 125.534 null] >>}
\MusicPDFWriteObject{17403}{<< /D [\MusicPDFPageRef{801} /XYZ 72 644.652 null] >>}
\MusicPDFWriteObject{17404}{<< /D [\MusicPDFPageRef{801} /XYZ 72 564.817 null] >>}
\MusicPDFWriteObject{17405}{<< /D [\MusicPDFPageRef{801} /XYZ 72 429.294 null] >>}
\MusicPDFWriteObject{17406}{<< /D [\MusicPDFPageRef{801} /XYZ 72 336.308 null] >>}
\MusicPDFWriteObject{17407}{<< /D [\MusicPDFPageRef{802} /XYZ 72 408.008 null] >>}
\MusicPDFWriteObject{17408}{<< /D [\MusicPDFPageRef{802} /XYZ 72 308.063 null] >>}
\MusicPDFWriteObject{17409}{<< /D [\MusicPDFPageRef{803} /XYZ 72 635.845 null] >>}
\MusicPDFWriteObject{17410}{<< /D [\MusicPDFPageRef{803} /XYZ 72 531.295 null] >>}
\MusicPDFWriteObject{17411}{<< /D [\MusicPDFPageRef{804} /XYZ 72 399.359 null] >>}
\MusicPDFWriteObject{17412}{<< /D [\MusicPDFPageRef{807} /XYZ 72 384.397 null] >>}
\MusicPDFWriteObject{17413}{<< /D [\MusicPDFPageRef{811} /XYZ 72 528.334 null] >>}
\MusicPDFWriteObject{17414}{<< /D [\MusicPDFPageRef{811} /XYZ 72 157.79 null] >>}
\MusicPDFWriteObject{17415}{<< /D [\MusicPDFPageRef{812} /XYZ 72 290.797 null] >>}
\MusicPDFWriteObject{17416}{<< /D [\MusicPDFPageRef{804} /XYZ 72 198.428 null] >>}
\MusicPDFWriteObject{17417}{<< /D [\MusicPDFPageRef{805} /XYZ 72 596.169 null] >>}
\MusicPDFWriteObject{17418}{<< /D [\MusicPDFPageRef{805} /XYZ 72 411.332 null] >>}
\MusicPDFWriteObject{17419}{<< /D [\MusicPDFPageRef{805} /XYZ 72 365.154 null] >>}
\MusicPDFWriteObject{17420}{<< /D [\MusicPDFPageRef{806} /XYZ 72 669.212 null] >>}
\MusicPDFWriteObject{17421}{<< /D [\MusicPDFPageRef{806} /XYZ 72 406.053 null] >>}
\MusicPDFWriteObject{17422}{<< /D [\MusicPDFPageRef{807} /XYZ 72 641.863 null] >>}
\MusicPDFWriteObject{17423}{<< /D [\MusicPDFPageRef{807} /XYZ 72 582.118 null] >>}
\MusicPDFWriteObject{17424}{<< /D [\MusicPDFPageRef{815} /XYZ 72 188.264 null] >>}
\MusicPDFWriteObject{17425}{<< /D [\MusicPDFPageRef{816} /XYZ 72 701.27 null] >>}
\MusicPDFWriteObject{17426}{<< /D [\MusicPDFPageRef{818} /XYZ 72 440.451 null] >>}
\MusicPDFWriteObject{17427}{<< /D [\MusicPDFPageRef{818} /XYZ 72 112.384 null] >>}
\MusicPDFWriteObject{17428}{<< /D [\MusicPDFPageRef{819} /XYZ 72 407.404 null] >>}
\MusicPDFWriteObject{17429}{<< /D [\MusicPDFPageRef{828} /XYZ 72 457.414 null] >>}
\MusicPDFWriteObject{17430}{<< /D [\MusicPDFPageRef{828} /XYZ 72 367.805 null] >>}
\MusicPDFWriteObject{17431}{<< /D [\MusicPDFPageRef{828} /XYZ 72 244.184 null] >>}
\MusicPDFWriteObject{17432}{<< /D [\MusicPDFPageRef{829} /XYZ 72 620.776 null] >>}
\MusicPDFWriteObject{17433}{<< /D [\MusicPDFPageRef{830} /XYZ 72 393.676 null] >>}
\MusicPDFWriteObject{17434}{<< /D [\MusicPDFPageRef{832} /XYZ 72 411.051 null] >>}
\MusicPDFWriteObject{17435}{<< /D [\MusicPDFPageRef{832} /XYZ 72 337.209 null] >>}
\MusicPDFWriteObject{17436}{<< /D [\MusicPDFPageRef{832} /XYZ 72 260.069 null] >>}
\MusicPDFWriteObject{17437}{<< /D [\MusicPDFPageRef{832} /XYZ 72 178.137 null] >>}
\MusicPDFWriteObject{17438}{<< /D [\MusicPDFPageRef{833} /XYZ 72 458.582 null] >>}
\MusicPDFWriteObject{17439}{<< /D [\MusicPDFPageRef{836} /XYZ 72 185.961 null] >>}
\MusicPDFWriteObject{17440}{<< /D [\MusicPDFPageRef{837} /XYZ 72 360.714 null] >>}
\MusicPDFWriteObject{17441}{<< /D [\MusicPDFPageRef{837} /XYZ 72 200.921 null] >>}
\MusicPDFWriteObject{17442}{<< /D [\MusicPDFPageRef{841} /XYZ 72 642.956 null] >>}
\MusicPDFWriteObject{17443}{<< /D [\MusicPDFPageRef{841} /XYZ 72 402.065 null] >>}
\MusicPDFWriteObject{17444}{<< /D [\MusicPDFPageRef{841} /XYZ 72 138.685 null] >>}
\MusicPDFWriteObject{17445}{<< /D [\MusicPDFPageRef{842} /XYZ 72 597.445 null] >>}
\MusicPDFWriteObject{17446}{<< /D [\MusicPDFPageRef{842} /XYZ 72 454.167 null] >>}
\MusicPDFWriteObject{17447}{<< /D [\MusicPDFPageRef{842} /XYZ 72 262.927 null] >>}
\MusicPDFWriteObject{17448}{<< /D [\MusicPDFPageRef{843} /XYZ 72 701.27 null] >>}
\MusicPDFWriteObject{17449}{<< /D [\MusicPDFPageRef{833} /XYZ 72 227.019 null] >>}
\MusicPDFWriteObject{17450}{<< /D [\MusicPDFPageRef{843} /XYZ 72 466.293 null] >>}
\MusicPDFWriteObject{17451}{<< /D [\MusicPDFPageRef{843} /XYZ 72 272.29 null] >>}
\MusicPDFWriteObject{17452}{<< /D [\MusicPDFPageRef{844} /XYZ 72 614.555 null] >>}
\MusicPDFWriteObject{17453}{<< /D [\MusicPDFPageRef{844} /XYZ 72 463.715 null] >>}
\MusicPDFWriteObject{17454}{<< /D [\MusicPDFPageRef{844} /XYZ 72 403.301 null] >>}
\MusicPDFWriteObject{17455}{<< /D [\MusicPDFPageRef{844} /XYZ 72 250.023 null] >>}
\MusicPDFWriteObject{17456}{<< /D [\MusicPDFPageRef{845} /XYZ 72 544.34 null] >>}
\MusicPDFWriteObject{17457}{<< /D [\MusicPDFPageRef{845} /XYZ 72 378.583 null] >>}
\MusicPDFWriteObject{17458}{<< /D [\MusicPDFPageRef{845} /XYZ 72 200.114 null] >>}
\MusicPDFWriteObject{17459}{<< /D [\MusicPDFPageRef{834} /XYZ 72 578.962 null] >>}
\MusicPDFWriteObject{17460}{<< /D [\MusicPDFPageRef{834} /XYZ 72 464.251 null] >>}
\MusicPDFWriteObject{17461}{<< /D [\MusicPDFPageRef{834} /XYZ 72 264.198 null] >>}
\MusicPDFWriteObject{17462}{<< /D [\MusicPDFPageRef{835} /XYZ 72 406.361 null] >>}
\MusicPDFWriteObject{17463}{<< /D [\MusicPDFPageRef{835} /XYZ 72 204.439 null] >>}
\MusicPDFWriteObject{17464}{<< /D [\MusicPDFPageRef{836} /XYZ 72 577.129 null] >>}
\MusicPDFWriteObject{17465}{<< /D [\MusicPDFPageRef{836} /XYZ 72 268.347 null] >>}
\MusicPDFWriteObject{17466}{<< /D [\MusicPDFPageRef{846} /XYZ 72 518.047 null] >>}
\MusicPDFWriteObject{17467}{<< /D [\MusicPDFPageRef{848} /XYZ 72 498.898 null] >>}
\MusicPDFWriteObject{17468}{<< /D [\MusicPDFPageRef{848} /XYZ 72 407.818 null] >>}
\MusicPDFWriteObject{17469}{<< /D [\MusicPDFPageRef{848} /XYZ 72 242.915 null] >>}
\MusicPDFWriteObject{17470}{<< /D [\MusicPDFPageRef{848} /XYZ 72 112.384 null] >>}
\MusicPDFWriteObject{17471}{<< /D [\MusicPDFPageRef{849} /XYZ 72 606.819 null] >>}
\MusicPDFWriteObject{17472}{<< /D [\MusicPDFPageRef{850} /XYZ 72 588.801 null] >>}
\MusicPDFWriteObject{17473}{<< /D [\MusicPDFPageRef{850} /XYZ 72 343.506 null] >>}
\MusicPDFWriteObject{17474}{<< /D [\MusicPDFPageRef{850} /XYZ 72 165.894 null] >>}
\MusicPDFWriteObject{17475}{<< /D [\MusicPDFPageRef{846} /XYZ 72 485.126 null] >>}
\MusicPDFWriteObject{17476}{<< /D [\MusicPDFPageRef{846} /XYZ 72 465.355 null] >>}
\MusicPDFWriteObject{17477}{<< /D [\MusicPDFPageRef{846} /XYZ 72 419.283 null] >>}
\MusicPDFWriteObject{17478}{<< /D [\MusicPDFPageRef{846} /XYZ 72 248.235 null] >>}
\MusicPDFWriteObject{17479}{<< /D [\MusicPDFPageRef{846} /XYZ 72 169.242 null] >>}
\MusicPDFWriteObject{17480}{<< /D [\MusicPDFPageRef{847} /XYZ 72 125.534 null] >>}
\MusicPDFWriteObject{17481}{<< /D [\MusicPDFPageRef{848} /XYZ 72 589.977 null] >>}
\MusicPDFWriteObject{17482}{<< /D [\MusicPDFPageRef{848} /XYZ 72 557.588 null] >>}
\MusicPDFWriteObject{17483}{<< /D [\MusicPDFPageRef{851} /XYZ 72 608.138 null] >>}
\MusicPDFWriteObject{17484}{<< /D [\MusicPDFPageRef{854} /XYZ 72 334.743 null] >>}
\MusicPDFWriteObject{17485}{<< /D [\MusicPDFPageRef{851} /XYZ 72 221.666 null] >>}
\MusicPDFWriteObject{17486}{<< /D [\MusicPDFPageRef{852} /XYZ 72 701.27 null] >>}
\MusicPDFWriteObject{17487}{<< /D [\MusicPDFPageRef{852} /XYZ 72 610.722 null] >>}
\MusicPDFWriteObject{17488}{<< /D [\MusicPDFPageRef{853} /XYZ 72 617.831 null] >>}
\MusicPDFWriteObject{17489}{<< /D [\MusicPDFPageRef{853} /XYZ 72 355.724 null] >>}
\MusicPDFWriteObject{17490}{<< /D [\MusicPDFPageRef{853} /XYZ 72 204.439 null] >>}
\MusicPDFWriteObject{17491}{<< /D [\MusicPDFPageRef{854} /XYZ 72 701.27 null] >>}
\MusicPDFWriteObject{17492}{<< /D [\MusicPDFPageRef{854} /XYZ 72 551.749 null] >>}
\MusicPDFWriteObject{17493}{<< /D [\MusicPDFPageRef{855} /XYZ 72 654.362 null] >>}
\MusicPDFWriteObject{17494}{<< /D [\MusicPDFPageRef{859} /XYZ 72 358.217 null] >>}
\MusicPDFWriteObject{17495}{<< /D [\MusicPDFPageRef{861} /XYZ 72 634.63 null] >>}
\MusicPDFWriteObject{17496}{<< /D [\MusicPDFPageRef{861} /XYZ 72 449.612 null] >>}
\MusicPDFWriteObject{17497}{<< /D [\MusicPDFPageRef{861} /XYZ 72 258.473 null] >>}
\MusicPDFWriteObject{17498}{<< /D [\MusicPDFPageRef{862} /XYZ 72 701.27 null] >>}
\MusicPDFWriteObject{17499}{<< /D [\MusicPDFPageRef{862} /XYZ 72 588.6 null] >>}
\MusicPDFWriteObject{17500}{<< /D [\MusicPDFPageRef{863} /XYZ 72 701.27 null] >>}
\MusicPDFWriteObject{17501}{<< /D [\MusicPDFPageRef{863} /XYZ 72 473.148 null] >>}
\MusicPDFWriteObject{17502}{<< /D [\MusicPDFPageRef{855} /XYZ 72 531.451 null] >>}
\MusicPDFWriteObject{17503}{<< /D [\MusicPDFPageRef{855} /XYZ 72 369.087 null] >>}
\MusicPDFWriteObject{17504}{<< /D [\MusicPDFPageRef{855} /XYZ 72 172.443 null] >>}
\MusicPDFWriteObject{17505}{<< /D [\MusicPDFPageRef{856} /XYZ 72 646.971 null] >>}
\MusicPDFWriteObject{17506}{<< /D [\MusicPDFPageRef{856} /XYZ 72 477.539 null] >>}
\MusicPDFWriteObject{17507}{<< /D [\MusicPDFPageRef{857} /XYZ 72 331.506 null] >>}
\MusicPDFWriteObject{17508}{<< /D [\MusicPDFPageRef{858} /XYZ 72 385.896 null] >>}
\MusicPDFWriteObject{17509}{<< /D [\MusicPDFPageRef{859} /XYZ 72 559.778 null] >>}
\MusicPDFWriteObject{17510}{<< /D [\MusicPDFPageRef{865} /XYZ 72 595.293 null] >>}
\MusicPDFWriteObject{17511}{<< /D [\MusicPDFPageRef{867} /XYZ 72 184.421 null] >>}
\MusicPDFWriteObject{17512}{<< /D [\MusicPDFPageRef{868} /XYZ 72 701.27 null] >>}
\MusicPDFWriteObject{17513}{<< /D [\MusicPDFPageRef{868} /XYZ 72 607.829 null] >>}
\MusicPDFWriteObject{17514}{<< /D [\MusicPDFPageRef{865} /XYZ 72 486.826 null] >>}
\MusicPDFWriteObject{17515}{<< /D [\MusicPDFPageRef{866} /XYZ 72 608.849 null] >>}
\MusicPDFWriteObject{17516}{<< /D [\MusicPDFPageRef{866} /XYZ 72 441.478 null] >>}
\MusicPDFWriteObject{17517}{<< /D [\MusicPDFPageRef{866} /XYZ 72 309.382 null] >>}
\MusicPDFWriteObject{17518}{<< /D [\MusicPDFPageRef{866} /XYZ 72 210.091 null] >>}
\MusicPDFWriteObject{17519}{<< /D [\MusicPDFPageRef{866} /XYZ 72 158.706 null] >>}
\MusicPDFWriteObject{17520}{<< /D [\MusicPDFPageRef{867} /XYZ 72 653.923 null] >>}
\MusicPDFWriteObject{17521}{<< /D [\MusicPDFPageRef{867} /XYZ 72 343.837 null] >>}
\MusicPDFWriteObject{17522}{<< /D [\MusicPDFPageRef{868} /XYZ 72 407.793 null] >>}
\MusicPDFWriteObject{17523}{<< /D [\MusicPDFPageRef{869} /XYZ 72 357.858 null] >>}
\MusicPDFWriteObject{17524}{<< /D [\MusicPDFPageRef{869} /XYZ 72 211.065 null] >>}
\MusicPDFWriteObject{17525}{<< /D [\MusicPDFPageRef{870} /XYZ 72 701.27 null] >>}
\MusicPDFWriteObject{17526}{<< /D [\MusicPDFPageRef{871} /XYZ 72 525.212 null] >>}
\MusicPDFWriteObject{17527}{<< /D [\MusicPDFPageRef{871} /XYZ 72 172.249 null] >>}
\MusicPDFWriteObject{17528}{<< /D [\MusicPDFPageRef{872} /XYZ 72 494.542 null] >>}
\MusicPDFWriteObject{17529}{<< /D [\MusicPDFPageRef{872} /XYZ 72 140.777 null] >>}
\MusicPDFWriteObject{17530}{<< /D [\MusicPDFPageRef{873} /XYZ 72 628.09 null] >>}
\MusicPDFWriteObject{17531}{<< /D [\MusicPDFPageRef{873} /XYZ 72 508.032 null] >>}
\MusicPDFWriteObject{17532}{<< /D [\MusicPDFPageRef{874} /XYZ 72 324.995 null] >>}
\MusicPDFWriteObject{17533}{<< /D [\MusicPDFPageRef{874} /XYZ 72 210.484 null] >>}
\MusicPDFWriteObject{17534}{<< /D [\MusicPDFPageRef{875} /XYZ 72 280.883 null] >>}
\MusicPDFWriteObject{17535}{<< /D [\MusicPDFPageRef{877} /XYZ 72 644.557 null] >>}
\MusicPDFWriteObject{17536}{<< /D [\MusicPDFPageRef{877} /XYZ 72 457.959 null] >>}
\MusicPDFWriteObject{17537}{<< /D [\MusicPDFPageRef{877} /XYZ 72 191.288 null] >>}
\MusicPDFWriteObject{17538}{<< /D [\MusicPDFPageRef{878} /XYZ 72 618.409 null] >>}
\MusicPDFWriteObject{17539}{<< /D [\MusicPDFPageRef{878} /XYZ 72 374.275 null] >>}
\MusicPDFWriteObject{17540}{<< /D [\MusicPDFPageRef{880} /XYZ 72 652.783 null] >>}
\MusicPDFWriteObject{17541}{<< /D [\MusicPDFPageRef{880} /XYZ 72 284.698 null] >>}
\MusicPDFWriteObject{17542}{<< /D [\MusicPDFPageRef{881} /XYZ 72 652.844 null] >>}
\MusicPDFWriteObject{17543}{<< /D [\MusicPDFPageRef{881} /XYZ 72 457.865 null] >>}
\MusicPDFWriteObject{17544}{<< /D [\MusicPDFPageRef{881} /XYZ 72 265.713 null] >>}
\MusicPDFWriteObject{17545}{<< /D [\MusicPDFPageRef{881} /XYZ 72 193.38 null] >>}
\MusicPDFWriteObject{17546}{<< /D [\MusicPDFPageRef{882} /XYZ 72 258.909 null] >>}
\MusicPDFWriteObject{17547}{<< /D [\MusicPDFPageRef{882} /XYZ 72 191.288 null] >>}
\MusicPDFWriteObject{17548}{<< /D [\MusicPDFPageRef{883} /XYZ 72 622.004 null] >>}
\MusicPDFWriteObject{17549}{<< /D [\MusicPDFPageRef{883} /XYZ 72 332.43 null] >>}
\MusicPDFWriteObject{17550}{<< /D [\MusicPDFPageRef{885} /XYZ 72 605.779 null] >>}
\MusicPDFWriteObject{17551}{<< /D [\MusicPDFPageRef{885} /XYZ 72 333.338 null] >>}
\MusicPDFWriteObject{17552}{<< /D [\MusicPDFPageRef{885} /XYZ 72 158.753 null] >>}
\MusicPDFWriteObject{17553}{<< /D [\MusicPDFPageRef{886} /XYZ 72 541.674 null] >>}
\MusicPDFWriteObject{17554}{<< /D [\MusicPDFPageRef{888} /XYZ 72 415.898 null] >>}
\MusicPDFWriteObject{17555}{<< /D [\MusicPDFPageRef{888} /XYZ 72 385.519 null] >>}
\MusicPDFWriteObject{17556}{<< /D [\MusicPDFPageRef{888} /XYZ 72 167.025 null] >>}
\MusicPDFWriteObject{17557}{<< /D [\MusicPDFPageRef{886} /XYZ 72 455.73 null] >>}
\MusicPDFWriteObject{17558}{<< /D [\MusicPDFPageRef{886} /XYZ 72 205.368 null] >>}
\MusicPDFWriteObject{17559}{<< /D [\MusicPDFPageRef{886} /XYZ 72 125.534 null] >>}
\MusicPDFWriteObject{17560}{<< /D [\MusicPDFPageRef{887} /XYZ 72 633.48 null] >>}
\MusicPDFWriteObject{17561}{<< /D [\MusicPDFPageRef{887} /XYZ 72 549.637 null] >>}
\MusicPDFWriteObject{17562}{<< /D [\MusicPDFPageRef{887} /XYZ 72 392.706 null] >>}
\MusicPDFWriteObject{17563}{<< /D [\MusicPDFPageRef{887} /XYZ 72 240.744 null] >>}
\MusicPDFWriteObject{17564}{<< /D [\MusicPDFPageRef{888} /XYZ 72 469.553 null] >>}
\MusicPDFWriteObject{17565}{<< /D [\MusicPDFPageRef{889} /XYZ 72 201.817 null] >>}
\MusicPDFWriteObject{17566}{<< /D [\MusicPDFPageRef{892} /XYZ 72 251.634 null] >>}
\MusicPDFWriteObject{17567}{<< /D [\MusicPDFPageRef{892} /XYZ 72 130.597 null] >>}
\MusicPDFWriteObject{17568}{<< /D [\MusicPDFPageRef{893} /XYZ 72 640.807 null] >>}
\MusicPDFWriteObject{17569}{<< /D [\MusicPDFPageRef{893} /XYZ 72 360.438 null] >>}
\MusicPDFWriteObject{17570}{<< /D [\MusicPDFPageRef{893} /XYZ 72 117.447 null] >>}
\MusicPDFWriteObject{17571}{<< /D [\MusicPDFPageRef{894} /XYZ 72 227.113 null] >>}
\MusicPDFWriteObject{17572}{<< /D [\MusicPDFPageRef{894} /XYZ 72 138.808 null] >>}
\MusicPDFWriteObject{17573}{<< /D [\MusicPDFPageRef{895} /XYZ 72 510.577 null] >>}
\MusicPDFWriteObject{17574}{<< /D [\MusicPDFPageRef{896} /XYZ 72 532.273 null] >>}
\MusicPDFWriteObject{17575}{<< /D [\MusicPDFPageRef{897} /XYZ 72 172.088 null] >>}
\MusicPDFWriteObject{17576}{<< /D [\MusicPDFPageRef{890} /XYZ 72 477.72 null] >>}
\MusicPDFWriteObject{17577}{<< /D [\MusicPDFPageRef{898} /XYZ 72 474.152 null] >>}
\MusicPDFWriteObject{17578}{<< /D [\MusicPDFPageRef{890} /XYZ 72 255.049 null] >>}
\MusicPDFWriteObject{17579}{<< /D [\MusicPDFPageRef{890} /XYZ 72 151.836 null] >>}
\MusicPDFWriteObject{17580}{<< /D [\MusicPDFPageRef{891} /XYZ 72 676.366 null] >>}
\MusicPDFWriteObject{17581}{<< /D [\MusicPDFPageRef{891} /XYZ 72 365.569 null] >>}
\MusicPDFWriteObject{17582}{<< /D [\MusicPDFPageRef{891} /XYZ 72 180.176 null] >>}
\MusicPDFWriteObject{17583}{<< /D [\MusicPDFPageRef{892} /XYZ 72 656.434 null] >>}
\MusicPDFWriteObject{17584}{<< /D [\MusicPDFPageRef{892} /XYZ 72 464.653 null] >>}
\MusicPDFWriteObject{17585}{<< /D [\MusicPDFPageRef{898} /XYZ 72 128.183 null] >>}
\MusicPDFWriteObject{17586}{<< /D [\MusicPDFPageRef{899} /XYZ 72 253.028 null] >>}
\MusicPDFWriteObject{17587}{<< /D [\MusicPDFPageRef{900} /XYZ 72 250.765 null] >>}
\MusicPDFWriteObject{17588}{<< /D [\MusicPDFPageRef{901} /XYZ 72 612.688 null] >>}
\MusicPDFWriteObject{17589}{<< /D [\MusicPDFPageRef{901} /XYZ 72 276.134 null] >>}
\MusicPDFWriteObject{17590}{<< /D [\MusicPDFPageRef{902} /XYZ 72 616.028 null] >>}
\MusicPDFWriteObject{17591}{<< /D [\MusicPDFPageRef{902} /XYZ 72 504.484 null] >>}
\MusicPDFWriteObject{17592}{<< /D [\MusicPDFPageRef{903} /XYZ 72 386.633 null] >>}
\MusicPDFWriteObject{17593}{<< /D [\MusicPDFPageRef{904} /XYZ 72 235.004 null] >>}
\MusicPDFWriteObject{17594}{<< /D [\MusicPDFPageRef{904} /XYZ 72 170.05 null] >>}
\MusicPDFWriteObject{17595}{<< /D [\MusicPDFPageRef{905} /XYZ 72 288.131 null] >>}
\MusicPDFWriteObject{17596}{<< /D [\MusicPDFPageRef{906} /XYZ 72 701.27 null] >>}
\MusicPDFWriteObject{17597}{<< /D [\MusicPDFPageRef{906} /XYZ 72 349.496 null] >>}
\MusicPDFWriteObject{17598}{<< /D [\MusicPDFPageRef{907} /XYZ 72 582.398 null] >>}
\MusicPDFWriteObject{17599}{<< /D [\MusicPDFPageRef{907} /XYZ 72 395.475 null] >>}
\MusicPDFWriteObject{17600}{<< /D [\MusicPDFPageRef{910} /XYZ 72 701.27 null] >>}
\MusicPDFWriteObject{17601}{<< /D [\MusicPDFPageRef{907} /XYZ 72 241.112 null] >>}
\MusicPDFWriteObject{17602}{<< /D [\MusicPDFPageRef{908} /XYZ 72 701.27 null] >>}
\MusicPDFWriteObject{17603}{<< /D [\MusicPDFPageRef{908} /XYZ 72 313.892 null] >>}
\MusicPDFWriteObject{17604}{<< /D [\MusicPDFPageRef{908} /XYZ 72 226.289 null] >>}
\MusicPDFWriteObject{17605}{<< /D [\MusicPDFPageRef{909} /XYZ 72 575.627 null] >>}
\MusicPDFWriteObject{17606}{<< /D [\MusicPDFPageRef{909} /XYZ 72 482.151 null] >>}
\MusicPDFWriteObject{17607}{<< /D [\MusicPDFPageRef{909} /XYZ 72 305.202 null] >>}
\MusicPDFWriteObject{17608}{<< /D [\MusicPDFPageRef{909} /XYZ 72 198.574 null] >>}
\MusicPDFWriteObject{17609}{<< /D [\MusicPDFPageRef{910} /XYZ 72 497.768 null] >>}
\MusicPDFWriteObject{17610}{<< /D [\MusicPDFPageRef{910} /XYZ 72 187.79 null] >>}
\MusicPDFWriteObject{17611}{<< /D [\MusicPDFPageRef{911} /XYZ 72 161.962 null] >>}
\MusicPDFWriteObject{17612}{<< /D [\MusicPDFPageRef{912} /XYZ 72 343.057 null] >>}
\MusicPDFWriteObject{17613}{<< /D [\MusicPDFPageRef{913} /XYZ 72 498.225 null] >>}
\MusicPDFWriteObject{17614}{<< /D [\MusicPDFPageRef{913} /XYZ 72 279.848 null] >>}
\MusicPDFWriteObject{17615}{<< /D [\MusicPDFPageRef{914} /XYZ 72 517.676 null] >>}
\MusicPDFWriteObject{17616}{<< /D [\MusicPDFPageRef{916} /XYZ 72 309.129 null] >>}
\MusicPDFWriteObject{17617}{<< /D [\MusicPDFPageRef{916} /XYZ 72 264.52 null] >>}
\MusicPDFWriteObject{17618}{<< /D [\MusicPDFPageRef{916} /XYZ 72 201.602 null] >>}
\MusicPDFWriteObject{17619}{<< /D [\MusicPDFPageRef{916} /XYZ 72 138.685 null] >>}
\MusicPDFWriteObject{17620}{<< /D [\MusicPDFPageRef{917} /XYZ 72 570.699 null] >>}
\MusicPDFWriteObject{17621}{<< /D [\MusicPDFPageRef{917} /XYZ 72 525.009 null] >>}
\MusicPDFWriteObject{17622}{<< /D [\MusicPDFPageRef{917} /XYZ 72 315.079 null] >>}
\MusicPDFWriteObject{17623}{<< /D [\MusicPDFPageRef{917} /XYZ 72 112.384 null] >>}
\MusicPDFWriteObject{17624}{<< /D [\MusicPDFPageRef{918} /XYZ 72 594.739 null] >>}
\MusicPDFWriteObject{17625}{<< /D [\MusicPDFPageRef{918} /XYZ 72 463.999 null] >>}
\MusicPDFWriteObject{17626}{<< /D [\MusicPDFPageRef{914} /XYZ 72 390.896 null] >>}
\MusicPDFWriteObject{17627}{<< /D [\MusicPDFPageRef{918} /XYZ 72 399.145 null] >>}
\MusicPDFWriteObject{17628}{<< /D [\MusicPDFPageRef{918} /XYZ 72 237.191 null] >>}
\MusicPDFWriteObject{17629}{<< /D [\MusicPDFPageRef{919} /XYZ 72 642.924 null] >>}
\MusicPDFWriteObject{17630}{<< /D [\MusicPDFPageRef{919} /XYZ 72 578.835 null] >>}
\MusicPDFWriteObject{17631}{<< /D [\MusicPDFPageRef{919} /XYZ 72 292.04 null] >>}
\MusicPDFWriteObject{17632}{<< /D [\MusicPDFPageRef{920} /XYZ 72 701.27 null] >>}
\MusicPDFWriteObject{17633}{<< /D [\MusicPDFPageRef{920} /XYZ 72 447.122 null] >>}
\MusicPDFWriteObject{17634}{<< /D [\MusicPDFPageRef{915} /XYZ 72 701.27 null] >>}
\MusicPDFWriteObject{17635}{<< /D [\MusicPDFPageRef{915} /XYZ 72 418.763 null] >>}
\MusicPDFWriteObject{17636}{<< /D [\MusicPDFPageRef{915} /XYZ 72 334.095 null] >>}
\MusicPDFWriteObject{17637}{<< /D [\MusicPDFPageRef{915} /XYZ 72 192.675 null] >>}
\MusicPDFWriteObject{17638}{<< /D [\MusicPDFPageRef{916} /XYZ 72 701.27 null] >>}
\MusicPDFWriteObject{17639}{<< /D [\MusicPDFPageRef{916} /XYZ 72 656.661 null] >>}
\MusicPDFWriteObject{17640}{<< /D [\MusicPDFPageRef{916} /XYZ 72 434.775 null] >>}
\MusicPDFWriteObject{17641}{<< /D [\MusicPDFPageRef{920} /XYZ 72 208.765 null] >>}
\MusicPDFWriteObject{17642}{<< /D [\MusicPDFPageRef{921} /XYZ 72 590.787 null] >>}
\MusicPDFWriteObject{17643}{<< /D [\MusicPDFPageRef{921} /XYZ 72 311.114 null] >>}
\MusicPDFWriteObject{17644}{<< /D [\MusicPDFPageRef{921} /XYZ 72 112.384 null] >>}
\MusicPDFWriteObject{17645}{<< /D [\MusicPDFPageRef{922} /XYZ 72 638.821 null] >>}
\MusicPDFWriteObject{17646}{<< /D [\MusicPDFPageRef{922} /XYZ 72 458.182 null] >>}
\MusicPDFWriteObject{17647}{<< /D [\MusicPDFPageRef{922} /XYZ 72 382.433 null] >>}
\MusicPDFWriteObject{17648}{<< /D [\MusicPDFPageRef{922} /XYZ 72 315.201 null] >>}
\MusicPDFWriteObject{17649}{<< /D [\MusicPDFPageRef{922} /XYZ 72 247.968 null] >>}
\MusicPDFWriteObject{17650}{<< /D [\MusicPDFPageRef{923} /XYZ 72 701.27 null] >>}
\MusicPDFWriteObject{17651}{<< /D [\MusicPDFPageRef{923} /XYZ 72 651.353 null] >>}
\MusicPDFWriteObject{17652}{<< /D [\MusicPDFPageRef{923} /XYZ 72 295.754 null] >>}
\MusicPDFWriteObject{17653}{<< /D [\MusicPDFPageRef{924} /XYZ 72 475.655 null] >>}
\MusicPDFWriteObject{17654}{<< /D [\MusicPDFPageRef{924} /XYZ 72 330.969 null] >>}
\MusicPDFWriteObject{17655}{<< /D [\MusicPDFPageRef{927} /XYZ 72 280.281 null] >>}
\MusicPDFWriteObject{17656}{<< /D [\MusicPDFPageRef{928} /XYZ 72 680.835 null] >>}
\MusicPDFWriteObject{17657}{<< /D [\MusicPDFPageRef{928} /XYZ 72 336.504 null] >>}
\MusicPDFWriteObject{17658}{<< /D [\MusicPDFPageRef{924} /XYZ 72 148.811 null] >>}
\MusicPDFWriteObject{17659}{<< /D [\MusicPDFPageRef{925} /XYZ 72 701.27 null] >>}
\MusicPDFWriteObject{17660}{<< /D [\MusicPDFPageRef{925} /XYZ 72 470.942 null] >>}
\MusicPDFWriteObject{17661}{<< /D [\MusicPDFPageRef{925} /XYZ 72 394.846 null] >>}
\MusicPDFWriteObject{17662}{<< /D [\MusicPDFPageRef{925} /XYZ 72 202.502 null] >>}
\MusicPDFWriteObject{17663}{<< /D [\MusicPDFPageRef{926} /XYZ 72 301.155 null] >>}
\MusicPDFWriteObject{17664}{<< /D [\MusicPDFPageRef{927} /XYZ 72 701.27 null] >>}
\MusicPDFWriteObject{17665}{<< /D [\MusicPDFPageRef{927} /XYZ 72 368.393 null] >>}
\MusicPDFWriteObject{17666}{<< /D [\MusicPDFPageRef{928} /XYZ 72 201.046 null] >>}
\MusicPDFWriteObject{17667}{<< /D [\MusicPDFPageRef{929} /XYZ 72 607.452 null] >>}
\MusicPDFWriteObject{17668}{<< /D [\MusicPDFPageRef{929} /XYZ 72 315.045 null] >>}
\MusicPDFWriteObject{17669}{<< /D [\MusicPDFPageRef{929} /XYZ 72 161.962 null] >>}
\MusicPDFWriteObject{17670}{<< /D [\MusicPDFPageRef{930} /XYZ 72 619.971 null] >>}
\MusicPDFWriteObject{17671}{<< /D [\MusicPDFPageRef{930} /XYZ 72 448.051 null] >>}
\MusicPDFWriteObject{17672}{<< /D [\MusicPDFPageRef{931} /XYZ 72 701.27 null] >>}
\MusicPDFWriteObject{17673}{<< /D [\MusicPDFPageRef{931} /XYZ 72 642.092 null] >>}
\MusicPDFWriteObject{17674}{<FEFF004300680061007000740065007200200035002E00200053007000610063006500730020006F00660020006100640065006C006900630020002D00660075006E006300740069006F006E007300200061006E00640020004900640065006C006900630020005A002D00660075006E006300740069006F006E0073>}
\MusicPDFWriteObject{17675}{<< /S /GoTo /D <636861707465722E35> >>}
\MusicPDFWriteObject{17676}{<< /Title \MusicPDFRef{17709} /A \MusicPDFRef{17710} /Parent \MusicPDFRef{1} /Prev \MusicPDFRef{10351} /Next \MusicPDFRef{17711} /First \MusicPDFRef{17712} /Last \MusicPDFRef{17713} /Count -3 >>}
\MusicPDFWriteObject{17677}{<< /Title \MusicPDFRef{17714} /A \MusicPDFRef{17715} /Parent \MusicPDFRef{10351} /Next \MusicPDFRef{17716} >>}
\MusicPDFWriteObject{17678}{<< /Title \MusicPDFRef{17717} /A \MusicPDFRef{17718} /Parent \MusicPDFRef{10351} /Prev \MusicPDFRef{17719} >>}
\MusicPDFWriteObject{17679}{<FEFF0034002E0031002E002000530063006800610075006400650072002000620061007300650073>}
\MusicPDFWriteObject{17680}{<< /S /GoTo /D <73656374696F6E2E342E31> >>}
\MusicPDFWriteObject{17681}{<< /Title \MusicPDFRef{17720} /A \MusicPDFRef{17721} /Parent \MusicPDFRef{9103} /Prev \MusicPDFRef{10352} /Next \MusicPDFRef{17722} >>}
\MusicPDFWriteObject{17682}{<FEFF0034002E0036002E00200050006500720069006F00640069006300200053006300680061007500640065007200200042006100730065007300200069006E0020007400680065002000700072006500730065006E006300650020006F00660020006D0075006C007400690070006C00690063006900740079>}
\MusicPDFWriteObject{17683}{<< /S /GoTo /D <73656374696F6E2E342E36> >>}
\MusicPDFWriteObject{17684}{<< /Title \MusicPDFRef{17723} /A \MusicPDFRef{17724} /Parent \MusicPDFRef{9103} /Prev \MusicPDFRef{17725} /Next \MusicPDFRef{10353} >>}
\MusicPDFWriteObject{17685}{<FEFF0033002E0032002E00200047006C006F00620061006C0020004900640065006C006900630020005A006500740061002000460075006E006300740069006F006E0073>}
\MusicPDFWriteObject{17686}{<< /S /GoTo /D <73656374696F6E2E332E32> >>}
\MusicPDFWriteObject{17687}{<FEFF0033002E0033002E00200053002D004C006F00630061006C0020004900640065006C006900630020005A006500740061002000460075006E006300740069006F006E0073>}
\MusicPDFWriteObject{17688}{<< /S /GoTo /D <73656374696F6E2E332E33> >>}
\MusicPDFWriteObject{17689}{<FEFF0031002E0034002E002000500075007400740069006E0067002000650076006500720079007400680069006E006700200074006F006700650074006800650072003A00200054007200610063006500200046006F0072006D0075006C006100200061006E006400200063006F006E00730065007100750065006E006300650073>}
\MusicPDFWriteObject{17690}{<< /S /GoTo /D <73656374696F6E2E312E34> >>}
\MusicPDFWriteObject{17691}{<FEFF0031002E0035002E0020004F006E0020007400680065002000700072006F006700720061006D00730020006F006600200043006F006E006E0065007300200061006E0064002000440065006E0069006E006700650072>}
\MusicPDFWriteObject{17692}{<< /S /GoTo /D <73656374696F6E2E312E35> >>}
\MusicPDFWriteObject{17693}{<FEFF0041007000700065006E00640069007800200041002E0020005000720065007000610072006100740069006F006E0073002000660072006F006D00200041006C00670065006200720061006900630020004E0075006D0062006500720020005400680065006F00720079>}
\MusicPDFWriteObject{17694}{<< /S /GoTo /D <617070656E6469782E41> >>}
\MusicPDFWriteObject{17695}{<< /Title \MusicPDFRef{17726} /A \MusicPDFRef{17727} /Parent \MusicPDFRef{1} /Prev \MusicPDFRef{17728} /Next \MusicPDFRef{10372} /First \MusicPDFRef{17729} /Last \MusicPDFRef{17730} /Count -3 >>}
\MusicPDFWriteObject{17696}{<< /Title \MusicPDFRef{17731} /A \MusicPDFRef{17732} /Parent \MusicPDFRef{10372} /Next \MusicPDFRef{17733} >>}
\MusicPDFWriteObject{17697}{<< /Title \MusicPDFRef{17734} /A \MusicPDFRef{17735} /Parent \MusicPDFRef{10372} /Prev \MusicPDFRef{17736} >>}
\MusicPDFWriteObject{17698}{<FEFF0042002E0031002E00200043006F006E00760065007800690074007900200062006F0075006E00640020006F006E0020004C006600280073002C0029>}
\MusicPDFWriteObject{17699}{<< /S /GoTo /D <73656374696F6E2E422E31> >>}
\MusicPDFWriteObject{17700}{<< /Title \MusicPDFRef{17737} /A \MusicPDFRef{17738} /Parent \MusicPDFRef{9120} /Prev \MusicPDFRef{10373} /Next \MusicPDFRef{17739} >>}
\MusicPDFWriteObject{17701}{<FEFF0042002E0037002E0020005200650063006900700072006F00630061006C007300200061006E00640020005000610072007400690061006C002000450075006C00650072002000700072006F006400750063007400730020006F006E002000740068006500200031002D006C0069006E0065>}
\MusicPDFWriteObject{17702}{<< /S /GoTo /D <73656374696F6E2E422E37> >>}
\MusicPDFWriteObject{17703}{<< /Title \MusicPDFRef{17740} /A \MusicPDFRef{17741} /Parent \MusicPDFRef{9120} /Prev \MusicPDFRef{17742} /Next \MusicPDFRef{10374} >>}
\MusicPDFWriteObject{17704}{<FEFF0043002E0032002E00200043006F006D0070006C0065006D0065006E007400730020006F006E0020006D0061007000700069006E00670020007300700061006300650073002C00200074006F0070006F006C006F0067006900630061006C002000740065006E0073006F0072002000700072006F00640075006300740073002C00200061006E00640020006E00750063006C006500610072006900740079>}
\MusicPDFWriteObject{17705}{<< /S /GoTo /D <73656374696F6E2E432E32> >>}
\MusicPDFWriteObject{17706}{<< /Title \MusicPDFRef{17743} /A \MusicPDFRef{17744} /Parent \MusicPDFRef{8889} /Prev \MusicPDFRef{10377} /Next \MusicPDFRef{10380} >>}
\MusicPDFWriteObject{17707}{<FEFF0043002E0034002E00200049006E0074006500670072006100740069006F006E00200069006E0020007400680065002000730065006E007300650020006F00660020005000650074007400690073>}
\MusicPDFWriteObject{17708}{<< /S /GoTo /D <73656374696F6E2E432E34> >>}
\MusicPDFWriteObject{17709}{<FEFF004300680061007000740065007200200036002E0020004C006F00630061006C002D0047006C006F00620061006C00200073007400720075006300740075007200650073>}
\MusicPDFWriteObject{17710}{<< /S /GoTo /D <636861707465722E36> >>}
\MusicPDFWriteObject{17711}{<< /Title \MusicPDFRef{17745} /A \MusicPDFRef{17746} /Parent \MusicPDFRef{1} /Prev \MusicPDFRef{17676} /Next \MusicPDFRef{17747} /First \MusicPDFRef{17748} /Last \MusicPDFRef{17749} /Count -5 >>}
\MusicPDFWriteObject{17712}{<< /Title \MusicPDFRef{17750} /A \MusicPDFRef{17751} /Parent \MusicPDFRef{17676} /Next \MusicPDFRef{17752} >>}
\MusicPDFWriteObject{17713}{<< /Title \MusicPDFRef{17753} /A \MusicPDFRef{17754} /Parent \MusicPDFRef{17676} /Prev \MusicPDFRef{17752} >>}
\MusicPDFWriteObject{17714}{<FEFF0035002E0031002E002000540068006500200074006F0070006F006C006F0067006900630061006C00200076006500630074006F007200200073007000610063006500730020006F00660020006100640065006C006900630020002D00660075006E006300740069006F006E007300200061006E00640020004900640065006C006900630020005A002D00660075006E006300740069006F006E0073>}
\MusicPDFWriteObject{17715}{<< /S /GoTo /D <73656374696F6E2E352E31> >>}
\MusicPDFWriteObject{17716}{<< /Title \MusicPDFRef{17755} /A \MusicPDFRef{17756} /Parent \MusicPDFRef{10351} /Prev \MusicPDFRef{17677} /Next \MusicPDFRef{17757} >>}
\MusicPDFWriteObject{17717}{<FEFF0035002E0035002E00200053007400720075006300740075007200650020005400680065006F00720065006D0020004900490049003A0020007300740072006F006E00670020006400750061006C007300200061006E006400200045006900670065006E0064006900730074007200690062007500740069006F006E0073>}
\MusicPDFWriteObject{17718}{<< /S /GoTo /D <73656374696F6E2E352E35> >>}
\MusicPDFWriteObject{17719}{<< /Title \MusicPDFRef{17758} /A \MusicPDFRef{17759} /Parent \MusicPDFRef{10351} /Prev \MusicPDFRef{17757} /Next \MusicPDFRef{17678} >>}
\MusicPDFWriteObject{17720}{<FEFF0034002E0032002E0020004200610073006500730020006F00660020006500780070006F006E0065006E007400690061006C006C00790020006400650063006100790069006E006700200053006300680077006100720074007A002D004200720075006800610074002000660075006E006300740069006F006E0073003A002000740068006500200063006100730065002000470020003D00200052003E0030>}
\MusicPDFWriteObject{17721}{<< /S /GoTo /D <73656374696F6E2E342E32> >>}
\MusicPDFWriteObject{17722}{<< /Title \MusicPDFRef{17760} /A \MusicPDFRef{17761} /Parent \MusicPDFRef{9103} /Prev \MusicPDFRef{17681} /Next \MusicPDFRef{17725} >>}
\MusicPDFWriteObject{17723}{<FEFF0034002E0035002E00200050006500720069006F00640069006300200053006300680061007500640065007200200042006100730065007300200069006E002000740068006500200061006200730065006E006300650020006F00660020006D0075006C007400690070006C00690063006900740079>}
\MusicPDFWriteObject{17724}{<< /S /GoTo /D <73656374696F6E2E342E35> >>}
\MusicPDFWriteObject{17725}{<< /Title \MusicPDFRef{17762} /A \MusicPDFRef{17763} /Parent \MusicPDFRef{9103} /Prev \MusicPDFRef{17722} /Next \MusicPDFRef{17684} >>}
\MusicPDFWriteObject{17726}{<FEFF0043006800610070007400650072002000310034002E00200043006F006D0070006C0065006D0065006E00740073>}
\MusicPDFWriteObject{17727}{<< /S /GoTo /D <636861707465722E3134> >>}
\MusicPDFWriteObject{17728}{<< /Title \MusicPDFRef{17764} /A \MusicPDFRef{17765} /Parent \MusicPDFRef{1} /Prev \MusicPDFRef{17766} /Next \MusicPDFRef{17695} /First \MusicPDFRef{17767} /Last \MusicPDFRef{17768} /Count -3 >>}
\MusicPDFWriteObject{17729}{<< /Title \MusicPDFRef{17769} /A \MusicPDFRef{17770} /Parent \MusicPDFRef{17695} /Next \MusicPDFRef{17771} >>}
\MusicPDFWriteObject{17730}{<< /Title \MusicPDFRef{17772} /A \MusicPDFRef{17773} /Parent \MusicPDFRef{17695} /Prev \MusicPDFRef{17771} >>}
\MusicPDFWriteObject{17731}{<FEFF0041002E0031002E00200054006F0070006F006C006F006700790020006F0066002000410046002C002000410046002C005300200061006E006400200047004C00310028004100460029002C00200047004C0031002800410046002C00530029002C00200061006E006400200074006800650069007200200063006C006100730073002000670072006F007500700073>}
\MusicPDFWriteObject{17732}{<< /S /GoTo /D <73656374696F6E2E412E31> >>}
\MusicPDFWriteObject{17733}{<< /Title \MusicPDFRef{17774} /A \MusicPDFRef{17775} /Parent \MusicPDFRef{10372} /Prev \MusicPDFRef{17696} /Next \MusicPDFRef{17776} >>}
\MusicPDFWriteObject{17734}{<FEFF0041002E0035002E0020004900640065006C00690063002000480061007200640079002D004C006900740074006C00650077006F006F00640020004D006100780069006D0061006C0020005400680065006F00720065006D00200061006E00640020004C006500620065007300670075006500200044006900660066006500720065006E00740069006100740069006F006E>}
\MusicPDFWriteObject{17735}{<< /S /GoTo /D <73656374696F6E2E412E35> >>}
\MusicPDFWriteObject{17736}{<< /Title \MusicPDFRef{17777} /A \MusicPDFRef{17778} /Parent \MusicPDFRef{10372} /Prev \MusicPDFRef{17776} /Next \MusicPDFRef{17697} >>}
\MusicPDFWriteObject{17737}{<FEFF0042002E0032002E0020004100730079006D00700074006F007400690063002000430061007500630068007900200049006E00740065006700720061006C00200046006F0072006D0075006C0061>}
\MusicPDFWriteObject{17738}{<< /S /GoTo /D <73656374696F6E2E422E32> >>}
\MusicPDFWriteObject{17739}{<< /Title \MusicPDFRef{17779} /A \MusicPDFRef{17780} /Parent \MusicPDFRef{9120} /Prev \MusicPDFRef{17700} /Next \MusicPDFRef{17781} >>}
\MusicPDFWriteObject{17740}{<FEFF0042002E0036002E00200043006C006100730073006900630061006C0020007A00650072006F002D006600720065006500200072006500670069006F006E007300200066006F00720020004800650063006B00650020004C002D00660075006E006300740069006F006E0073>}
\MusicPDFWriteObject{17741}{<< /S /GoTo /D <73656374696F6E2E422E36> >>}
\MusicPDFWriteObject{17742}{<< /Title \MusicPDFRef{17782} /A \MusicPDFRef{17783} /Parent \MusicPDFRef{9120} /Prev \MusicPDFRef{17781} /Next \MusicPDFRef{17703} >>}
\MusicPDFWriteObject{17743}{<FEFF0043002E0033002E0020004E00750063006C0065006100720020006F00700065007200610074006F00720073>}
\MusicPDFWriteObject{17744}{<< /S /GoTo /D <73656374696F6E2E432E33> >>}
\MusicPDFWriteObject{17745}{<FEFF004300680061007000740065007200200037002E0020005300630068006100750064006500720020004200610073006500730020006F00660020004900640065006C006900630020005A006500740061002000660075006E006300740069006F006E0073>}
\MusicPDFWriteObject{17746}{<< /S /GoTo /D <636861707465722E37> >>}
\MusicPDFWriteObject{17747}{<< /Title \MusicPDFRef{17784} /A \MusicPDFRef{17785} /Parent \MusicPDFRef{1} /Prev \MusicPDFRef{17711} /Next \MusicPDFRef{17786} /First \MusicPDFRef{17787} /Last \MusicPDFRef{17788} /Count -5 >>}
\MusicPDFWriteObject{17748}{<< /Title \MusicPDFRef{17789} /A \MusicPDFRef{17790} /Parent \MusicPDFRef{17711} /Next \MusicPDFRef{17791} >>}
\MusicPDFWriteObject{17749}{<< /Title \MusicPDFRef{17792} /A \MusicPDFRef{17793} /Parent \MusicPDFRef{17711} /Prev \MusicPDFRef{17794} >>}
\MusicPDFWriteObject{17750}{<FEFF0036002E0031002E00200053007400720075006300740075007200650020005400680065006F00720065006D002000490056003A0020006C006F00630061006C002000770069007400680069006E00200067006C006F00620061006C>}
\MusicPDFWriteObject{17751}{<< /S /GoTo /D <73656374696F6E2E362E31> >>}
\MusicPDFWriteObject{17752}{<< /Title \MusicPDFRef{17795} /A \MusicPDFRef{17796} /Parent \MusicPDFRef{17676} /Prev \MusicPDFRef{17712} /Next \MusicPDFRef{17713} >>}
\MusicPDFWriteObject{17753}{<FEFF0036002E0033002E0020004C006F00630061006C00200061006E006400200047006C006F00620061006C00200043006F006E007600650072007300650020005400680065006F00720065006D0073>}
\MusicPDFWriteObject{17754}{<< /S /GoTo /D <73656374696F6E2E362E33> >>}
\MusicPDFWriteObject{17755}{<FEFF0035002E0032002E00200053007400720075006300740075007200650020005400680065006F00720065006D00200049003A00200074006F0070006F006C006F006700790020006F0066002000740068006500200073007000610063006500730020006F00660020006100640065006C006900630020002D00660075006E006300740069006F006E0073>}
\MusicPDFWriteObject{17756}{<< /S /GoTo /D <73656374696F6E2E352E32> >>}
\MusicPDFWriteObject{17757}{<< /Title \MusicPDFRef{17797} /A \MusicPDFRef{17798} /Parent \MusicPDFRef{10351} /Prev \MusicPDFRef{17716} /Next \MusicPDFRef{17719} >>}
\MusicPDFWriteObject{17758}{<FEFF0035002E0034002E00200053007400720075006300740075007200650020005400680065006F00720065006D002000490049003A00200053002D006C006F00630061006C0020010C006500630068002000730065007100750065006E00630065>}
\MusicPDFWriteObject{17759}{<< /S /GoTo /D <73656374696F6E2E352E34> >>}
\MusicPDFWriteObject{17760}{<FEFF0034002E0033002E0020005000720065007000610072006100740069006F006E00730020006F006E002000690073006F00740079007000690063002000700072006F006A006500630074006F00720073>}
\MusicPDFWriteObject{17761}{<< /S /GoTo /D <73656374696F6E2E342E33> >>}
\MusicPDFWriteObject{17762}{<FEFF0034002E0034002E0020004500780070006F006E0065006E007400690061006C002000730079007300740065006D0020006F006E0020006D006F006E006F00670065006E00690063002000710075006F007400690065006E007400730020006F006600200052003E0030002D006D006F00640075006C0061007400650064002000670072006F007500700073>}
\MusicPDFWriteObject{17763}{<< /S /GoTo /D <73656374696F6E2E342E34> >>}
\MusicPDFWriteObject{17764}{<FEFF0043006800610070007400650072002000310033002E0020005A00650072006F00730020006F00660020004800650063006B00650020004C002D00660075006E006300740069006F006E0073>}
\MusicPDFWriteObject{17765}{<< /S /GoTo /D <636861707465722E3133> >>}
\MusicPDFWriteObject{17766}{<< /Title \MusicPDFRef{17799} /A \MusicPDFRef{17800} /Parent \MusicPDFRef{1} /Prev \MusicPDFRef{17801} /Next \MusicPDFRef{17728} /First \MusicPDFRef{17802} /Last \MusicPDFRef{17803} /Count -5 >>}
\MusicPDFWriteObject{17767}{<< /Title \MusicPDFRef{17804} /A \MusicPDFRef{17805} /Parent \MusicPDFRef{17728} /Next \MusicPDFRef{17806} >>}
\MusicPDFWriteObject{17768}{<< /Title \MusicPDFRef{17807} /A \MusicPDFRef{17808} /Parent \MusicPDFRef{17728} /Prev \MusicPDFRef{17806} >>}
\MusicPDFWriteObject{17769}{<FEFF00310034002E0031002E0020004F006E00200074006800650020006D0075006C007400690070006C00690063006900740069006500730020006F006600200074006800650020007A00650072006F0073>}
\MusicPDFWriteObject{17770}{<< /S /GoTo /D <73656374696F6E2E31342E31> >>}
\MusicPDFWriteObject{17771}{<< /Title \MusicPDFRef{17809} /A \MusicPDFRef{17810} /Parent \MusicPDFRef{17695} /Prev \MusicPDFRef{17729} /Next \MusicPDFRef{17730} >>}
\MusicPDFWriteObject{17772}{<FEFF00310034002E0033002E00200054006800650020006800690067006800650072002000720061006E006B00200063006100730065>}
\MusicPDFWriteObject{17773}{<< /S /GoTo /D <73656374696F6E2E31342E33> >>}
\MusicPDFWriteObject{17774}{<FEFF0041002E0032002E0020004F00720062006900740061006C002000730074007200750063007400750072006500730020006F006E0020006100640065006C0065007300200061006E00640020006900640065006C00650073>}
\MusicPDFWriteObject{17775}{<< /S /GoTo /D <73656374696F6E2E412E32> >>}
\MusicPDFWriteObject{17776}{<< /Title \MusicPDFRef{17811} /A \MusicPDFRef{17812} /Parent \MusicPDFRef{10372} /Prev \MusicPDFRef{17733} /Next \MusicPDFRef{17736} >>}
\MusicPDFWriteObject{17777}{<FEFF0041002E0034002E0020004400690072006900630068006C006500740027007300200053002D0075006E006900740020005400680065006F00720065006D003A00200065007300740069006D00610074006500730020006200790020006C00610074007400690063006500200070006F0069006E0074002D0063006F0075006E00740069006E0067>}
\MusicPDFWriteObject{17778}{<< /S /GoTo /D <73656374696F6E2E412E34> >>}
\MusicPDFWriteObject{17779}{<FEFF0042002E0033002E0020004100730079006D00700074006F007400690063002000520065007300690064007500650020005400680065006F00720065006D>}
\MusicPDFWriteObject{17780}{<< /S /GoTo /D <73656374696F6E2E422E33> >>}
\MusicPDFWriteObject{17781}{<< /Title \MusicPDFRef{17813} /A \MusicPDFRef{17814} /Parent \MusicPDFRef{9120} /Prev \MusicPDFRef{17739} /Next \MusicPDFRef{17742} >>}
\MusicPDFWriteObject{17782}{<FEFF0042002E0035002E00200052006500630075007200720065006E0063006500200061006E0064002000650071007500690064006900730074007200690062007500740069006F006E>}
\MusicPDFWriteObject{17783}{<< /S /GoTo /D <73656374696F6E2E422E35> >>}
\MusicPDFWriteObject{17784}{<FEFF004300680061007000740065007200200038002E00200050006500720069006F0064006900630020004F0072006200690074007300200061006E006400200046006900780065006400200050006F0069006E00740073>}
\MusicPDFWriteObject{17785}{<< /S /GoTo /D <636861707465722E38> >>}
\MusicPDFWriteObject{17786}{<< /Title \MusicPDFRef{17815} /A \MusicPDFRef{17816} /Parent \MusicPDFRef{1} /Prev \MusicPDFRef{17747} /Next \MusicPDFRef{17817} /First \MusicPDFRef{17818} /Last \MusicPDFRef{17819} /Count -3 >>}
\MusicPDFWriteObject{17787}{<< /Title \MusicPDFRef{17820} /A \MusicPDFRef{17821} /Parent \MusicPDFRef{17747} /Next \MusicPDFRef{17822} >>}
\MusicPDFWriteObject{17788}{<< /Title \MusicPDFRef{17823} /A \MusicPDFRef{17824} /Parent \MusicPDFRef{17747} /Prev \MusicPDFRef{17825} >>}
\MusicPDFWriteObject{17789}{<FEFF0037002E0031002E00200050006F006C006100720020004D0075006C007400690070006C0069006300690074006900650073>}
\MusicPDFWriteObject{17790}{<< /S /GoTo /D <73656374696F6E2E372E31> >>}
\MusicPDFWriteObject{17791}{<< /Title \MusicPDFRef{17826} /A \MusicPDFRef{17827} /Parent \MusicPDFRef{17711} /Prev \MusicPDFRef{17748} /Next \MusicPDFRef{17828} >>}
\MusicPDFWriteObject{17792}{<FEFF0037002E0035002E0020005300740072006F006E006700200043006F002D005000720069006D0061006C006900740079>}
\MusicPDFWriteObject{17793}{<< /S /GoTo /D <73656374696F6E2E372E35> >>}
\MusicPDFWriteObject{17794}{<< /Title \MusicPDFRef{17829} /A \MusicPDFRef{17830} /Parent \MusicPDFRef{17711} /Prev \MusicPDFRef{17828} /Next \MusicPDFRef{17749} >>}
\MusicPDFWriteObject{17795}{<FEFF0036002E0032002E00200053007400720075006300740075007200650020005400680065006F00720065006D00200056003A002000690073006F0074007900700069006300200065007800680061007500730074006900760065006E006500730073>}
\MusicPDFWriteObject{17796}{<< /S /GoTo /D <73656374696F6E2E362E32> >>}
\MusicPDFWriteObject{17797}{<FEFF0035002E0033002E0020005700650069006C002700730020005000720069006E0063006900700061006C002000560061006C007500650020004F00700065007200610074006F00720073>}
\MusicPDFWriteObject{17798}{<< /S /GoTo /D <73656374696F6E2E352E33> >>}
\MusicPDFWriteObject{17799}{<FEFF0043006800610070007400650072002000310032002E00200054007200610063006500200066006F0072006D0075006C0061003A002000470065006F006D0065007400720069006300200053006900640065>}
\MusicPDFWriteObject{17800}{<< /S /GoTo /D <636861707465722E3132> >>}
\MusicPDFWriteObject{17801}{<< /Title \MusicPDFRef{17831} /A \MusicPDFRef{17832} /Parent \MusicPDFRef{1} /Prev \MusicPDFRef{17817} /Next \MusicPDFRef{17766} /First \MusicPDFRef{17833} /Last \MusicPDFRef{17834} /Count -3 >>}
\MusicPDFWriteObject{17802}{<< /Title \MusicPDFRef{17835} /A \MusicPDFRef{17836} /Parent \MusicPDFRef{17766} /Next \MusicPDFRef{17837} >>}
\MusicPDFWriteObject{17803}{<< /Title \MusicPDFRef{17838} /A \MusicPDFRef{17839} /Parent \MusicPDFRef{17766} /Prev \MusicPDFRef{17840} >>}
\MusicPDFWriteObject{17804}{<FEFF00310033002E0031002E0020005200650063006F006C006C0065006300740069006F006E00730020006F006E00200074006800650020004500780070006C006900630069007400200046006F0072006D0075006C00610073>}
\MusicPDFWriteObject{17805}{<< /S /GoTo /D <73656374696F6E2E31332E31> >>}
\MusicPDFWriteObject{17806}{<< /Title \MusicPDFRef{17841} /A \MusicPDFRef{17842} /Parent \MusicPDFRef{17728} /Prev \MusicPDFRef{17767} /Next \MusicPDFRef{17768} >>}
\MusicPDFWriteObject{17807}{<FEFF00310033002E0033002E0020004D00610069006E0020005400680065006F00720065006D00730020006F006E00200074006800650020006C006F0063006100740069006F006E0020006F006600200074006800650020007A00650072006F0073>}
\MusicPDFWriteObject{17808}{<< /S /GoTo /D <73656374696F6E2E31332E33> >>}
\MusicPDFWriteObject{17809}{<FEFF00310034002E0032002E0020005400680065002000630061007300650020006F006600200067006C006F00620061006C002000660075006E006300740069006F006E0020006600690065006C006400730020006F006600200070006F007300690074006900760065002000630068006100720061006300740065007200690073007400690063>}
\MusicPDFWriteObject{17810}{<< /S /GoTo /D <73656374696F6E2E31342E32> >>}
\MusicPDFWriteObject{17811}{<FEFF0041002E0033002E002000480061006100720020006D0065006100730075007200650073>}
\MusicPDFWriteObject{17812}{<< /S /GoTo /D <73656374696F6E2E412E33> >>}
\MusicPDFWriteObject{17813}{<FEFF0042002E0034002E0020004500780070006C006900630069007400200046006F0072006D0075006C0061007300200066006F00720020004800650063006B00650020004C002D00660075006E006300740069006F006E0073>}
\MusicPDFWriteObject{17814}{<< /S /GoTo /D <73656374696F6E2E422E34> >>}
\MusicPDFWriteObject{17815}{<FEFF004300680061007000740065007200200039002E002000480061007200640079002D0053006F0062006F006C0065007600200074006F0070006F006C006F0067006900650073>}
\MusicPDFWriteObject{17816}{<< /S /GoTo /D <636861707465722E39> >>}
\MusicPDFWriteObject{17817}{<< /Title \MusicPDFRef{17843} /A \MusicPDFRef{17844} /Parent \MusicPDFRef{1} /Prev \MusicPDFRef{17786} /Next \MusicPDFRef{17801} /First \MusicPDFRef{17845} /Last \MusicPDFRef{17846} /Count -4 >>}
\MusicPDFWriteObject{17818}{<< /Title \MusicPDFRef{17847} /A \MusicPDFRef{17848} /Parent \MusicPDFRef{17786} /Next \MusicPDFRef{17849} >>}
\MusicPDFWriteObject{17819}{<< /Title \MusicPDFRef{17850} /A \MusicPDFRef{17851} /Parent \MusicPDFRef{17786} /Prev \MusicPDFRef{17849} >>}
\MusicPDFWriteObject{17820}{<FEFF0038002E0031002E0020004F006E0020004F00720062006900740061006C00200053007500700070006F0072007400730020002D00200049>}
\MusicPDFWriteObject{17821}{<< /S /GoTo /D <73656374696F6E2E382E31> >>}
\MusicPDFWriteObject{17822}{<< /Title \MusicPDFRef{17852} /A \MusicPDFRef{17853} /Parent \MusicPDFRef{17747} /Prev \MusicPDFRef{17787} /Next \MusicPDFRef{17854} >>}
\MusicPDFWriteObject{17823}{<FEFF0038002E0035002E0020005400680065002000620069006700200070006900630074007500720065>}
\MusicPDFWriteObject{17824}{<< /S /GoTo /D <73656374696F6E2E382E35> >>}
\MusicPDFWriteObject{17825}{<< /Title \MusicPDFRef{17855} /A \MusicPDFRef{17856} /Parent \MusicPDFRef{17747} /Prev \MusicPDFRef{17854} /Next \MusicPDFRef{17788} >>}
\MusicPDFWriteObject{17826}{<FEFF0037002E0032002E00200050006F006C006500730020006F0066002000740068006500200053002D006900640065006C006900630020005A006500740061002000660075006E006300740069006F006E0073>}
\MusicPDFWriteObject{17827}{<< /S /GoTo /D <73656374696F6E2E372E32> >>}
\MusicPDFWriteObject{17828}{<< /Title \MusicPDFRef{17857} /A \MusicPDFRef{17858} /Parent \MusicPDFRef{17711} /Prev \MusicPDFRef{17791} /Next \MusicPDFRef{17794} >>}
\MusicPDFWriteObject{17829}{<FEFF0037002E0034002E0020004E00750063006C006500610072006900740079>}
\MusicPDFWriteObject{17830}{<< /S /GoTo /D <73656374696F6E2E372E34> >>}
\MusicPDFWriteObject{17831}{<FEFF0043006800610070007400650072002000310031002E00200054006800650020005300680065006C006C0020005300740072007500630074007500720065>}
\MusicPDFWriteObject{17832}{<< /S /GoTo /D <636861707465722E3131> >>}
\MusicPDFWriteObject{17833}{<< /Title \MusicPDFRef{17859} /A \MusicPDFRef{17860} /Parent \MusicPDFRef{17801} /Next \MusicPDFRef{17861} >>}
\MusicPDFWriteObject{17834}{<< /Title \MusicPDFRef{17862} /A \MusicPDFRef{17863} /Parent \MusicPDFRef{17801} /Prev \MusicPDFRef{17861} >>}
\MusicPDFWriteObject{17835}{<FEFF00310032002E0031002E00200053006D006F006F007400680065006E0069006E0067007300200061006C006F006E0067002000740068006500200064006900610067006F006E0061006C>}
\MusicPDFWriteObject{17836}{<< /S /GoTo /D <73656374696F6E2E31322E31> >>}
\MusicPDFWriteObject{17837}{<< /Title \MusicPDFRef{17864} /A \MusicPDFRef{17865} /Parent \MusicPDFRef{17766} /Prev \MusicPDFRef{17802} /Next \MusicPDFRef{17866} >>}
\MusicPDFWriteObject{17838}{<FEFF00310032002E0035002E00200054007200610063006500200066006F0072006D0075006C0061003A00200074006800650020004D00610069006E0020005400680065006F00720065006D0020006F00660020007400680069007300200077006F0072006B>}
\MusicPDFWriteObject{17839}{<< /S /GoTo /D <73656374696F6E2E31322E35> >>}
\MusicPDFWriteObject{17840}{<< /Title \MusicPDFRef{17867} /A \MusicPDFRef{17868} /Parent \MusicPDFRef{17766} /Prev \MusicPDFRef{17866} /Next \MusicPDFRef{17803} >>}
\MusicPDFWriteObject{17841}{<FEFF00310033002E0032002E00200048006F006C006F006D006F00720070006800690063002000660061006D0069006C0069006500730020006F006600200074006500730074002000660075006E006300740069006F006E0073>}
\MusicPDFWriteObject{17842}{<< /S /GoTo /D <73656374696F6E2E31332E32> >>}
\MusicPDFWriteObject{17843}{<FEFF0043006800610070007400650072002000310030002E00200054007200610063006500200066006F0072006D0075006C0061003A00200053007000650063007400720061006C00200053006900640065>}
\MusicPDFWriteObject{17844}{<< /S /GoTo /D <636861707465722E3130> >>}
\MusicPDFWriteObject{17845}{<< /Title \MusicPDFRef{17869} /A \MusicPDFRef{17870} /Parent \MusicPDFRef{17817} /Next \MusicPDFRef{17871} >>}
\MusicPDFWriteObject{17846}{<< /Title \MusicPDFRef{17872} /A \MusicPDFRef{17873} /Parent \MusicPDFRef{17817} /Prev \MusicPDFRef{17874} >>}
\MusicPDFWriteObject{17847}{<FEFF0039002E0031002E00200053006F0062006F006C006500760020007300700061006300650073002000570073002C0032002800470029>}
\MusicPDFWriteObject{17848}{<< /S /GoTo /D <73656374696F6E2E392E31> >>}
\MusicPDFWriteObject{17849}{<< /Title \MusicPDFRef{17875} /A \MusicPDFRef{17876} /Parent \MusicPDFRef{17786} /Prev \MusicPDFRef{17818} /Next \MusicPDFRef{17819} >>}
\MusicPDFWriteObject{17850}{<FEFF0039002E0033002E0020004D006F00640075006C0065002000730074007200750063007400750072006500730020006F0076006500720020004400280047002900200061006E006400200061007000700072006F00780069006D00610074006500200075006E006900740073>}
\MusicPDFWriteObject{17851}{<< /S /GoTo /D <73656374696F6E2E392E33> >>}
\MusicPDFWriteObject{17852}{<FEFF0038002E0032002E00200053007400610063006B002D007400680065006F0072006500740069006300200050006500720069006F0064006900630020004F0072006200690074007300200061006E006400200046006900780065006400200050006F0069006E00740073>}
\MusicPDFWriteObject{17853}{<< /S /GoTo /D <73656374696F6E2E382E32> >>}
\MusicPDFWriteObject{17854}{<< /Title \MusicPDFRef{17877} /A \MusicPDFRef{17878} /Parent \MusicPDFRef{17747} /Prev \MusicPDFRef{17822} /Next \MusicPDFRef{17825} >>}
\MusicPDFWriteObject{17855}{<FEFF0038002E0034002E0020004F00720062006900740061006C0020005400720061006300650073>}
\MusicPDFWriteObject{17856}{<< /S /GoTo /D <73656374696F6E2E382E34> >>}
\MusicPDFWriteObject{17857}{<FEFF0037002E0033002E002000500061006C00650079002D005700690065006E006500720020004C006500760065006C007300200061006E00640020005300630068006100750064006500720020004200610073006500730020006F00660020005000720069006E0063006900700061006C002000500061007200740073>}
\MusicPDFWriteObject{17858}{<< /S /GoTo /D <73656374696F6E2E372E33> >>}
\MusicPDFWriteObject{17859}{<FEFF00310031002E0031002E00200054006800650020005300680065006C006C0020005300740072007500630074007500720065>}
\MusicPDFWriteObject{17860}{<< /S /GoTo /D <73656374696F6E2E31312E31> >>}
\MusicPDFWriteObject{17861}{<< /Title \MusicPDFRef{17879} /A \MusicPDFRef{17880} /Parent \MusicPDFRef{17801} /Prev \MusicPDFRef{17833} /Next \MusicPDFRef{17834} >>}
\MusicPDFWriteObject{17862}{<FEFF00310031002E0033002E0020004F00720062006900740061006C0020004B00650072006E0065006C007300200061006E00640020004F00720062006900740061006C0020005400720061006300650073>}
\MusicPDFWriteObject{17863}{<< /S /GoTo /D <73656374696F6E2E31312E33> >>}
\MusicPDFWriteObject{17864}{<FEFF00310032002E0032002E0020004C006500620065007300670075006500200044006900660066006500720065006E00740069006100740069006F006E00200061006C006F006E0067002000740068006500200064006900610067006F006E0061006C>}
\MusicPDFWriteObject{17865}{<< /S /GoTo /D <73656374696F6E2E31322E32> >>}
\MusicPDFWriteObject{17866}{<< /Title \MusicPDFRef{17881} /A \MusicPDFRef{17882} /Parent \MusicPDFRef{17766} /Prev \MusicPDFRef{17837} /Next \MusicPDFRef{17840} >>}
\MusicPDFWriteObject{17867}{<FEFF00310032002E0034002E00200053002D006C006F00630061006C002000740072006100630065002000630061006C00630075006C006100740069006F006E>}
\MusicPDFWriteObject{17868}{<< /S /GoTo /D <73656374696F6E2E31322E34> >>}
\MusicPDFWriteObject{17869}{<FEFF00310030002E0031002E002000480061007200640079002D0053006F0062006F006C0065007600200074006F0070006F006C006F00670069006500730020006F006E002000480030002C0063006F006E007400280046002C00530028004100460029002900200061006E0064002000480030002C0063006F006E00740028004F0046002C0053002C0053002800410046002C005300290029>}
\MusicPDFWriteObject{17870}{<< /S /GoTo /D <73656374696F6E2E31302E31> >>}
\MusicPDFWriteObject{17871}{<< /Title \MusicPDFRef{17883} /A \MusicPDFRef{17884} /Parent \MusicPDFRef{17817} /Prev \MusicPDFRef{17845} /Next \MusicPDFRef{17874} >>}
\MusicPDFWriteObject{17872}{<FEFF00310030002E0034002E00200054007200610063006500200066006F0072006D0075006C0061003A00200053007000650063007400720061006C00200053006900640065>}
\MusicPDFWriteObject{17873}{<< /S /GoTo /D <73656374696F6E2E31302E34> >>}
\MusicPDFWriteObject{17874}{<< /Title \MusicPDFRef{17885} /A \MusicPDFRef{17886} /Parent \MusicPDFRef{17817} /Prev \MusicPDFRef{17871} /Next \MusicPDFRef{17846} >>}
\MusicPDFWriteObject{17875}{<FEFF0039002E0032002E002000480061007200640079002D0053006F0062006F006C006500760020007300700061006300650073002000570073004900280047002900200061006E0064002000500061006C00650079002D005700690065006E0065007200200063006F006E0064006900740069006F006E0073>}
\MusicPDFWriteObject{17876}{<< /S /GoTo /D <73656374696F6E2E392E32> >>}
\MusicPDFWriteObject{17877}{<FEFF0038002E0033002E0020004F006E0020004F00720062006900740061006C00200053007500700070006F0072007400730020002D002000490049>}
\MusicPDFWriteObject{17878}{<< /S /GoTo /D <73656374696F6E2E382E33> >>}
\MusicPDFWriteObject{17879}{<FEFF00310031002E0032002E002000540068006500200070006900650063006500730020006F00660020007400680065002000700075007A007A006C0065003A002000500061006C00650079002D005700690065006E006500720020006C006500760065006C007300200061006E00640020004F00720062006900740061006C0020004D00610072006B00750073006800650076006900630068002000420061007300650073>}
\MusicPDFWriteObject{17880}{<< /S /GoTo /D <73656374696F6E2E31312E32> >>}
\MusicPDFWriteObject{17881}{<FEFF00310032002E0033002E002000500075007400740069006E00670020007400680065002000700069006500630065007300200074006F006700650074006800650072003A00200067006C006F00620061006C00200074006F00200053002D006C006F00630061006C0020006400650063006F006E0073007400720075006300740069006F006E0020006F00660020007400680065002000670065006F006D006500740072006900630020006B00650072006E0065006C>}
\MusicPDFWriteObject{17882}{<< /S /GoTo /D <73656374696F6E2E31322E33> >>}
\MusicPDFWriteObject{17883}{<FEFF00310030002E0032002E00200053007000650063007400720061006C0020004D00610072006B00750073006800650076006900630068002000420061007300650073>}
\MusicPDFWriteObject{17884}{<< /S /GoTo /D <73656374696F6E2E31302E32> >>}
\MusicPDFWriteObject{17885}{<FEFF00310030002E0033002E00200053007000650063007400720061006C0020004E00750063006C006500610072006900740079>}
\MusicPDFWriteObject{17886}{<< /S /GoTo /D <73656374696F6E2E31302E33> >>}
\pdfcatalog{/Outlines \MusicPDFRef{1} /Names \MusicPDFRef{2} /PageMode /UseOutlines /OpenAction \MusicPDFRef{3}}
\pdfinfo{/Author <> /Title <> /Subject <> /Creator <4C615465582077697468206879706572726566> /Keywords <> /Trapped /False}
\MusicPDFSetAnnotations{3}{[\MusicPDFRef{4} \MusicPDFRef{5} \MusicPDFRef{6} \MusicPDFRef{7} \MusicPDFRef{8} \MusicPDFRef{9} \MusicPDFRef{10} \MusicPDFRef{11} \MusicPDFRef{12} \MusicPDFRef{13} \MusicPDFRef{14} \MusicPDFRef{15} \MusicPDFRef{16} \MusicPDFRef{17} \MusicPDFRef{18} \MusicPDFRef{19} \MusicPDFRef{20} \MusicPDFRef{21} \MusicPDFRef{22} \MusicPDFRef{23} \MusicPDFRef{24} \MusicPDFRef{25} \MusicPDFRef{26} \MusicPDFRef{27} \MusicPDFRef{28} \MusicPDFRef{29} \MusicPDFRef{30} \MusicPDFRef{31} \MusicPDFRef{32} \MusicPDFRef{33} \MusicPDFRef{34} \MusicPDFRef{35} \MusicPDFRef{36} \MusicPDFRef{37}]}
\MusicPDFSetAnnotations{4}{[\MusicPDFRef{38} \MusicPDFRef{39} \MusicPDFRef{40} \MusicPDFRef{41} \MusicPDFRef{42} \MusicPDFRef{43} \MusicPDFRef{44} \MusicPDFRef{45} \MusicPDFRef{46} \MusicPDFRef{47} \MusicPDFRef{48} \MusicPDFRef{49} \MusicPDFRef{50} \MusicPDFRef{51} \MusicPDFRef{52} \MusicPDFRef{53} \MusicPDFRef{54} \MusicPDFRef{55} \MusicPDFRef{56} \MusicPDFRef{57} \MusicPDFRef{58} \MusicPDFRef{59} \MusicPDFRef{60} \MusicPDFRef{61} \MusicPDFRef{62} \MusicPDFRef{63} \MusicPDFRef{64} \MusicPDFRef{65} \MusicPDFRef{66} \MusicPDFRef{67} \MusicPDFRef{68} \MusicPDFRef{69} \MusicPDFRef{70} \MusicPDFRef{71} \MusicPDFRef{72} \MusicPDFRef{73} \MusicPDFRef{74} \MusicPDFRef{75} \MusicPDFRef{76} \MusicPDFRef{77} \MusicPDFRef{78}]}
\MusicPDFSetAnnotations{5}{[\MusicPDFRef{79} \MusicPDFRef{80} \MusicPDFRef{81} \MusicPDFRef{82} \MusicPDFRef{83} \MusicPDFRef{84} \MusicPDFRef{85} \MusicPDFRef{86} \MusicPDFRef{87} \MusicPDFRef{88} \MusicPDFRef{89} \MusicPDFRef{90} \MusicPDFRef{91} \MusicPDFRef{92} \MusicPDFRef{93} \MusicPDFRef{94} \MusicPDFRef{95} \MusicPDFRef{96} \MusicPDFRef{97} \MusicPDFRef{98} \MusicPDFRef{99} \MusicPDFRef{100} \MusicPDFRef{101} \MusicPDFRef{102} \MusicPDFRef{103} \MusicPDFRef{104} \MusicPDFRef{105} \MusicPDFRef{106} \MusicPDFRef{107} \MusicPDFRef{108} \MusicPDFRef{109} \MusicPDFRef{110} \MusicPDFRef{111}]}
\MusicPDFSetAnnotations{7}{[\MusicPDFRef{112} \MusicPDFRef{113}]}
\MusicPDFSetAnnotations{9}{[\MusicPDFRef{114} \MusicPDFRef{115} \MusicPDFRef{116} \MusicPDFRef{117} \MusicPDFRef{118} \MusicPDFRef{119} \MusicPDFRef{120} \MusicPDFRef{121}]}
\MusicPDFSetAnnotations{16}{[\MusicPDFRef{122}]}
\MusicPDFSetAnnotations{20}{[\MusicPDFRef{123}]}
\MusicPDFSetAnnotations{21}{[\MusicPDFRef{124}]}
\MusicPDFSetAnnotations{22}{[\MusicPDFRef{125} \MusicPDFRef{126}]}
\MusicPDFSetAnnotations{24}{[\MusicPDFRef{127} \MusicPDFRef{128} \MusicPDFRef{129} \MusicPDFRef{130}]}
\MusicPDFSetAnnotations{26}{[\MusicPDFRef{131}]}
\MusicPDFSetAnnotations{27}{[\MusicPDFRef{132} \MusicPDFRef{133} \MusicPDFRef{134} \MusicPDFRef{135}]}
\MusicPDFSetAnnotations{33}{[\MusicPDFRef{136}]}
\MusicPDFSetAnnotations{34}{[\MusicPDFRef{137} \MusicPDFRef{138} \MusicPDFRef{139} \MusicPDFRef{140} \MusicPDFRef{141} \MusicPDFRef{142}]}
\MusicPDFSetAnnotations{35}{[\MusicPDFRef{143} \MusicPDFRef{144} \MusicPDFRef{145} \MusicPDFRef{146}]}
\MusicPDFSetAnnotations{37}{[\MusicPDFRef{147} \MusicPDFRef{148} \MusicPDFRef{149} \MusicPDFRef{150}]}
\MusicPDFSetAnnotations{38}{[\MusicPDFRef{151} \MusicPDFRef{152} \MusicPDFRef{153} \MusicPDFRef{154} \MusicPDFRef{155}]}
\MusicPDFSetAnnotations{39}{[\MusicPDFRef{156}]}
\MusicPDFSetAnnotations{40}{[\MusicPDFRef{157} \MusicPDFRef{158}]}
\MusicPDFSetAnnotations{41}{[\MusicPDFRef{159}]}
\MusicPDFSetAnnotations{42}{[\MusicPDFRef{160}]}
\MusicPDFSetAnnotations{43}{[\MusicPDFRef{161}]}
\MusicPDFSetAnnotations{44}{[\MusicPDFRef{162} \MusicPDFRef{163}]}
\MusicPDFSetAnnotations{45}{[\MusicPDFRef{164} \MusicPDFRef{165} \MusicPDFRef{166} \MusicPDFRef{167}]}
\MusicPDFSetAnnotations{46}{[\MusicPDFRef{168} \MusicPDFRef{169} \MusicPDFRef{170}]}
\MusicPDFSetAnnotations{47}{[\MusicPDFRef{171} \MusicPDFRef{172}]}
\MusicPDFSetAnnotations{48}{[\MusicPDFRef{173}]}
\MusicPDFSetAnnotations{50}{[\MusicPDFRef{174}]}
\MusicPDFSetAnnotations{51}{[\MusicPDFRef{175} \MusicPDFRef{176} \MusicPDFRef{177}]}
\MusicPDFSetAnnotations{52}{[\MusicPDFRef{178} \MusicPDFRef{179}]}
\MusicPDFSetAnnotations{57}{[\MusicPDFRef{180} \MusicPDFRef{181} \MusicPDFRef{182} \MusicPDFRef{183} \MusicPDFRef{184}]}
\MusicPDFSetAnnotations{58}{[\MusicPDFRef{185} \MusicPDFRef{186} \MusicPDFRef{187} \MusicPDFRef{188} \MusicPDFRef{189} \MusicPDFRef{190} \MusicPDFRef{191} \MusicPDFRef{192} \MusicPDFRef{193} \MusicPDFRef{194} \MusicPDFRef{195} \MusicPDFRef{196} \MusicPDFRef{197} \MusicPDFRef{198} \MusicPDFRef{199} \MusicPDFRef{200} \MusicPDFRef{201} \MusicPDFRef{202} \MusicPDFRef{203} \MusicPDFRef{204} \MusicPDFRef{205} \MusicPDFRef{206} \MusicPDFRef{207} \MusicPDFRef{208} \MusicPDFRef{209} \MusicPDFRef{210} \MusicPDFRef{211} \MusicPDFRef{212} \MusicPDFRef{213} \MusicPDFRef{214} \MusicPDFRef{215} \MusicPDFRef{216} \MusicPDFRef{217}]}
\MusicPDFSetAnnotations{59}{[\MusicPDFRef{218} \MusicPDFRef{219} \MusicPDFRef{220} \MusicPDFRef{221} \MusicPDFRef{222} \MusicPDFRef{223} \MusicPDFRef{224} \MusicPDFRef{225} \MusicPDFRef{226}]}
\MusicPDFSetAnnotations{60}{[\MusicPDFRef{227}]}
\MusicPDFSetAnnotations{61}{[\MusicPDFRef{228} \MusicPDFRef{229} \MusicPDFRef{230} \MusicPDFRef{231} \MusicPDFRef{232} \MusicPDFRef{233} \MusicPDFRef{234} \MusicPDFRef{235} \MusicPDFRef{236} \MusicPDFRef{237} \MusicPDFRef{238} \MusicPDFRef{239} \MusicPDFRef{240} \MusicPDFRef{241} \MusicPDFRef{242} \MusicPDFRef{243} \MusicPDFRef{244}]}
\MusicPDFSetAnnotations{62}{[\MusicPDFRef{245} \MusicPDFRef{246} \MusicPDFRef{247} \MusicPDFRef{248} \MusicPDFRef{249} \MusicPDFRef{250} \MusicPDFRef{251} \MusicPDFRef{252} \MusicPDFRef{253} \MusicPDFRef{254} \MusicPDFRef{255} \MusicPDFRef{256} \MusicPDFRef{257} \MusicPDFRef{258} \MusicPDFRef{259} \MusicPDFRef{260} \MusicPDFRef{261} \MusicPDFRef{262} \MusicPDFRef{263} \MusicPDFRef{264} \MusicPDFRef{265}]}
\MusicPDFSetAnnotations{63}{[\MusicPDFRef{266}]}
\MusicPDFSetAnnotations{64}{[\MusicPDFRef{267}]}
\MusicPDFSetAnnotations{65}{[\MusicPDFRef{268} \MusicPDFRef{269} \MusicPDFRef{270} \MusicPDFRef{271} \MusicPDFRef{272}]}
\MusicPDFSetAnnotations{66}{[\MusicPDFRef{273} \MusicPDFRef{274} \MusicPDFRef{275} \MusicPDFRef{276} \MusicPDFRef{277} \MusicPDFRef{278}]}
\MusicPDFSetAnnotations{67}{[\MusicPDFRef{279} \MusicPDFRef{280} \MusicPDFRef{281} \MusicPDFRef{282} \MusicPDFRef{283} \MusicPDFRef{284} \MusicPDFRef{285} \MusicPDFRef{286} \MusicPDFRef{287} \MusicPDFRef{288}]}
\MusicPDFSetAnnotations{68}{[\MusicPDFRef{289} \MusicPDFRef{290} \MusicPDFRef{291} \MusicPDFRef{292} \MusicPDFRef{293} \MusicPDFRef{294} \MusicPDFRef{295} \MusicPDFRef{296} \MusicPDFRef{297} \MusicPDFRef{298}]}
\MusicPDFSetAnnotations{69}{[\MusicPDFRef{299} \MusicPDFRef{300} \MusicPDFRef{301} \MusicPDFRef{302} \MusicPDFRef{303} \MusicPDFRef{304} \MusicPDFRef{305} \MusicPDFRef{306} \MusicPDFRef{307} \MusicPDFRef{308} \MusicPDFRef{309} \MusicPDFRef{310} \MusicPDFRef{311} \MusicPDFRef{312} \MusicPDFRef{313} \MusicPDFRef{314} \MusicPDFRef{315} \MusicPDFRef{316} \MusicPDFRef{317} \MusicPDFRef{318} \MusicPDFRef{319}]}
\MusicPDFSetAnnotations{70}{[\MusicPDFRef{320} \MusicPDFRef{321} \MusicPDFRef{322} \MusicPDFRef{323} \MusicPDFRef{324} \MusicPDFRef{325} \MusicPDFRef{326} \MusicPDFRef{327} \MusicPDFRef{328} \MusicPDFRef{329} \MusicPDFRef{330} \MusicPDFRef{331} \MusicPDFRef{332} \MusicPDFRef{333} \MusicPDFRef{334} \MusicPDFRef{335} \MusicPDFRef{336} \MusicPDFRef{337} \MusicPDFRef{338} \MusicPDFRef{339} \MusicPDFRef{340} \MusicPDFRef{341} \MusicPDFRef{342}]}
\MusicPDFSetAnnotations{71}{[\MusicPDFRef{343} \MusicPDFRef{344} \MusicPDFRef{345} \MusicPDFRef{346} \MusicPDFRef{347} \MusicPDFRef{348} \MusicPDFRef{349} \MusicPDFRef{350} \MusicPDFRef{351} \MusicPDFRef{352} \MusicPDFRef{353} \MusicPDFRef{354} \MusicPDFRef{355}]}
\MusicPDFSetAnnotations{72}{[\MusicPDFRef{356} \MusicPDFRef{357} \MusicPDFRef{358} \MusicPDFRef{359} \MusicPDFRef{360} \MusicPDFRef{361} \MusicPDFRef{362} \MusicPDFRef{363} \MusicPDFRef{364} \MusicPDFRef{365} \MusicPDFRef{366} \MusicPDFRef{367} \MusicPDFRef{368} \MusicPDFRef{369} \MusicPDFRef{370} \MusicPDFRef{371} \MusicPDFRef{372} \MusicPDFRef{373} \MusicPDFRef{374} \MusicPDFRef{375} \MusicPDFRef{376}]}
\MusicPDFSetAnnotations{73}{[\MusicPDFRef{377} \MusicPDFRef{378} \MusicPDFRef{379} \MusicPDFRef{380} \MusicPDFRef{381} \MusicPDFRef{382}]}
\MusicPDFSetAnnotations{74}{[\MusicPDFRef{383} \MusicPDFRef{384} \MusicPDFRef{385} \MusicPDFRef{386} \MusicPDFRef{387} \MusicPDFRef{388} \MusicPDFRef{389}]}
\MusicPDFSetAnnotations{75}{[\MusicPDFRef{390} \MusicPDFRef{391} \MusicPDFRef{392} \MusicPDFRef{393} \MusicPDFRef{394} \MusicPDFRef{395} \MusicPDFRef{396}]}
\MusicPDFSetAnnotations{76}{[\MusicPDFRef{397} \MusicPDFRef{398}]}
\MusicPDFSetAnnotations{77}{[\MusicPDFRef{399} \MusicPDFRef{400} \MusicPDFRef{401} \MusicPDFRef{402} \MusicPDFRef{403}]}
\MusicPDFSetAnnotations{78}{[\MusicPDFRef{404} \MusicPDFRef{405} \MusicPDFRef{406} \MusicPDFRef{407} \MusicPDFRef{408} \MusicPDFRef{409} \MusicPDFRef{410} \MusicPDFRef{411} \MusicPDFRef{412} \MusicPDFRef{413} \MusicPDFRef{414} \MusicPDFRef{415}]}
\MusicPDFSetAnnotations{79}{[\MusicPDFRef{416} \MusicPDFRef{417} \MusicPDFRef{418}]}
\MusicPDFSetAnnotations{80}{[\MusicPDFRef{419} \MusicPDFRef{420} \MusicPDFRef{421} \MusicPDFRef{422} \MusicPDFRef{423} \MusicPDFRef{424} \MusicPDFRef{425} \MusicPDFRef{426} \MusicPDFRef{427} \MusicPDFRef{428}]}
\MusicPDFSetAnnotations{81}{[\MusicPDFRef{429} \MusicPDFRef{430} \MusicPDFRef{431} \MusicPDFRef{432} \MusicPDFRef{433} \MusicPDFRef{434} \MusicPDFRef{435} \MusicPDFRef{436}]}
\MusicPDFSetAnnotations{82}{[\MusicPDFRef{437} \MusicPDFRef{438} \MusicPDFRef{439} \MusicPDFRef{440} \MusicPDFRef{441} \MusicPDFRef{442} \MusicPDFRef{443} \MusicPDFRef{444} \MusicPDFRef{445} \MusicPDFRef{446} \MusicPDFRef{447} \MusicPDFRef{448} \MusicPDFRef{449} \MusicPDFRef{450}]}
\MusicPDFSetAnnotations{83}{[\MusicPDFRef{451} \MusicPDFRef{452} \MusicPDFRef{453} \MusicPDFRef{454} \MusicPDFRef{455} \MusicPDFRef{456} \MusicPDFRef{457} \MusicPDFRef{458} \MusicPDFRef{459} \MusicPDFRef{460} \MusicPDFRef{461} \MusicPDFRef{462} \MusicPDFRef{463}]}
\MusicPDFSetAnnotations{84}{[\MusicPDFRef{464} \MusicPDFRef{465} \MusicPDFRef{466} \MusicPDFRef{467} \MusicPDFRef{468} \MusicPDFRef{469} \MusicPDFRef{470} \MusicPDFRef{471} \MusicPDFRef{472} \MusicPDFRef{473} \MusicPDFRef{474} \MusicPDFRef{475}]}
\MusicPDFSetAnnotations{85}{[\MusicPDFRef{476} \MusicPDFRef{477} \MusicPDFRef{478} \MusicPDFRef{479} \MusicPDFRef{480} \MusicPDFRef{481} \MusicPDFRef{482} \MusicPDFRef{483} \MusicPDFRef{484} \MusicPDFRef{485} \MusicPDFRef{486}]}
\MusicPDFSetAnnotations{86}{[\MusicPDFRef{487} \MusicPDFRef{488} \MusicPDFRef{489} \MusicPDFRef{490} \MusicPDFRef{491} \MusicPDFRef{492} \MusicPDFRef{493}]}
\MusicPDFSetAnnotations{87}{[\MusicPDFRef{494} \MusicPDFRef{495} \MusicPDFRef{496} \MusicPDFRef{497} \MusicPDFRef{498}]}
\MusicPDFSetAnnotations{88}{[\MusicPDFRef{499} \MusicPDFRef{500} \MusicPDFRef{501} \MusicPDFRef{502} \MusicPDFRef{503} \MusicPDFRef{504} \MusicPDFRef{505} \MusicPDFRef{506} \MusicPDFRef{507} \MusicPDFRef{508} \MusicPDFRef{509} \MusicPDFRef{510} \MusicPDFRef{511}]}
\MusicPDFSetAnnotations{89}{[\MusicPDFRef{512} \MusicPDFRef{513} \MusicPDFRef{514} \MusicPDFRef{515} \MusicPDFRef{516} \MusicPDFRef{517} \MusicPDFRef{518}]}
\MusicPDFSetAnnotations{90}{[\MusicPDFRef{519} \MusicPDFRef{520} \MusicPDFRef{521} \MusicPDFRef{522} \MusicPDFRef{523} \MusicPDFRef{524}]}
\MusicPDFSetAnnotations{91}{[\MusicPDFRef{525} \MusicPDFRef{526} \MusicPDFRef{527} \MusicPDFRef{528} \MusicPDFRef{529} \MusicPDFRef{530} \MusicPDFRef{531} \MusicPDFRef{532} \MusicPDFRef{533} \MusicPDFRef{534}]}
\MusicPDFSetAnnotations{92}{[\MusicPDFRef{535} \MusicPDFRef{536} \MusicPDFRef{537} \MusicPDFRef{538} \MusicPDFRef{539} \MusicPDFRef{540} \MusicPDFRef{541} \MusicPDFRef{542}]}
\MusicPDFSetAnnotations{93}{[\MusicPDFRef{543} \MusicPDFRef{544} \MusicPDFRef{545} \MusicPDFRef{546} \MusicPDFRef{547} \MusicPDFRef{548} \MusicPDFRef{549} \MusicPDFRef{550} \MusicPDFRef{551} \MusicPDFRef{552} \MusicPDFRef{553}]}
\MusicPDFSetAnnotations{94}{[\MusicPDFRef{554} \MusicPDFRef{555} \MusicPDFRef{556}]}
\MusicPDFSetAnnotations{95}{[\MusicPDFRef{557} \MusicPDFRef{558} \MusicPDFRef{559} \MusicPDFRef{560} \MusicPDFRef{561} \MusicPDFRef{562} \MusicPDFRef{563} \MusicPDFRef{564} \MusicPDFRef{565} \MusicPDFRef{566}]}
\MusicPDFSetAnnotations{96}{[\MusicPDFRef{567} \MusicPDFRef{568} \MusicPDFRef{569} \MusicPDFRef{570} \MusicPDFRef{571} \MusicPDFRef{572} \MusicPDFRef{573} \MusicPDFRef{574} \MusicPDFRef{575} \MusicPDFRef{576} \MusicPDFRef{577} \MusicPDFRef{578} \MusicPDFRef{579} \MusicPDFRef{580} \MusicPDFRef{581} \MusicPDFRef{582}]}
\MusicPDFSetAnnotations{97}{[\MusicPDFRef{583} \MusicPDFRef{584} \MusicPDFRef{585} \MusicPDFRef{586} \MusicPDFRef{587} \MusicPDFRef{588} \MusicPDFRef{589} \MusicPDFRef{590}]}
\MusicPDFSetAnnotations{98}{[\MusicPDFRef{591} \MusicPDFRef{592} \MusicPDFRef{593} \MusicPDFRef{594} \MusicPDFRef{595} \MusicPDFRef{596} \MusicPDFRef{597} \MusicPDFRef{598} \MusicPDFRef{599} \MusicPDFRef{600} \MusicPDFRef{601} \MusicPDFRef{602} \MusicPDFRef{603} \MusicPDFRef{604} \MusicPDFRef{605} \MusicPDFRef{606} \MusicPDFRef{607}]}
\MusicPDFSetAnnotations{99}{[\MusicPDFRef{608} \MusicPDFRef{609}]}
\MusicPDFSetAnnotations{100}{[\MusicPDFRef{610} \MusicPDFRef{611} \MusicPDFRef{612} \MusicPDFRef{613} \MusicPDFRef{614} \MusicPDFRef{615} \MusicPDFRef{616} \MusicPDFRef{617} \MusicPDFRef{618} \MusicPDFRef{619} \MusicPDFRef{620} \MusicPDFRef{621}]}
\MusicPDFSetAnnotations{101}{[\MusicPDFRef{622} \MusicPDFRef{623} \MusicPDFRef{624} \MusicPDFRef{625} \MusicPDFRef{626} \MusicPDFRef{627} \MusicPDFRef{628}]}
\MusicPDFSetAnnotations{102}{[\MusicPDFRef{629} \MusicPDFRef{630} \MusicPDFRef{631}]}
\MusicPDFSetAnnotations{103}{[\MusicPDFRef{632} \MusicPDFRef{633} \MusicPDFRef{634} \MusicPDFRef{635} \MusicPDFRef{636} \MusicPDFRef{637} \MusicPDFRef{638} \MusicPDFRef{639} \MusicPDFRef{640} \MusicPDFRef{641} \MusicPDFRef{642} \MusicPDFRef{643} \MusicPDFRef{644} \MusicPDFRef{645} \MusicPDFRef{646} \MusicPDFRef{647} \MusicPDFRef{648} \MusicPDFRef{649}]}
\MusicPDFSetAnnotations{104}{[\MusicPDFRef{650} \MusicPDFRef{651} \MusicPDFRef{652} \MusicPDFRef{653} \MusicPDFRef{654} \MusicPDFRef{655} \MusicPDFRef{656} \MusicPDFRef{657} \MusicPDFRef{658}]}
\MusicPDFSetAnnotations{105}{[\MusicPDFRef{659} \MusicPDFRef{660} \MusicPDFRef{661} \MusicPDFRef{662} \MusicPDFRef{663} \MusicPDFRef{664} \MusicPDFRef{665} \MusicPDFRef{666} \MusicPDFRef{667} \MusicPDFRef{668}]}
\MusicPDFSetAnnotations{106}{[\MusicPDFRef{669} \MusicPDFRef{670} \MusicPDFRef{671} \MusicPDFRef{672} \MusicPDFRef{673} \MusicPDFRef{674} \MusicPDFRef{675} \MusicPDFRef{676} \MusicPDFRef{677} \MusicPDFRef{678} \MusicPDFRef{679} \MusicPDFRef{680} \MusicPDFRef{681}]}
\MusicPDFSetAnnotations{107}{[\MusicPDFRef{682} \MusicPDFRef{683} \MusicPDFRef{684} \MusicPDFRef{685}]}
\MusicPDFSetAnnotations{108}{[\MusicPDFRef{686} \MusicPDFRef{687} \MusicPDFRef{688} \MusicPDFRef{689} \MusicPDFRef{690} \MusicPDFRef{691} \MusicPDFRef{692} \MusicPDFRef{693}]}
\MusicPDFSetAnnotations{109}{[\MusicPDFRef{694} \MusicPDFRef{695} \MusicPDFRef{696} \MusicPDFRef{697} \MusicPDFRef{698} \MusicPDFRef{699} \MusicPDFRef{700}]}
\MusicPDFSetAnnotations{110}{[\MusicPDFRef{701} \MusicPDFRef{702} \MusicPDFRef{703} \MusicPDFRef{704} \MusicPDFRef{705} \MusicPDFRef{706} \MusicPDFRef{707} \MusicPDFRef{708} \MusicPDFRef{709} \MusicPDFRef{710} \MusicPDFRef{711} \MusicPDFRef{712} \MusicPDFRef{713} \MusicPDFRef{714} \MusicPDFRef{715} \MusicPDFRef{716} \MusicPDFRef{717} \MusicPDFRef{718}]}
\MusicPDFSetAnnotations{111}{[\MusicPDFRef{719} \MusicPDFRef{720} \MusicPDFRef{721} \MusicPDFRef{722} \MusicPDFRef{723} \MusicPDFRef{724} \MusicPDFRef{725} \MusicPDFRef{726} \MusicPDFRef{727} \MusicPDFRef{728}]}
\MusicPDFSetAnnotations{112}{[\MusicPDFRef{729} \MusicPDFRef{730} \MusicPDFRef{731} \MusicPDFRef{732} \MusicPDFRef{733} \MusicPDFRef{734} \MusicPDFRef{735} \MusicPDFRef{736}]}
\MusicPDFSetAnnotations{113}{[\MusicPDFRef{737} \MusicPDFRef{738} \MusicPDFRef{739} \MusicPDFRef{740} \MusicPDFRef{741} \MusicPDFRef{742} \MusicPDFRef{743} \MusicPDFRef{744} \MusicPDFRef{745} \MusicPDFRef{746} \MusicPDFRef{747} \MusicPDFRef{748} \MusicPDFRef{749} \MusicPDFRef{750}]}
\MusicPDFSetAnnotations{114}{[\MusicPDFRef{751} \MusicPDFRef{752} \MusicPDFRef{753} \MusicPDFRef{754} \MusicPDFRef{755} \MusicPDFRef{756} \MusicPDFRef{757} \MusicPDFRef{758} \MusicPDFRef{759} \MusicPDFRef{760} \MusicPDFRef{761} \MusicPDFRef{762} \MusicPDFRef{763} \MusicPDFRef{764} \MusicPDFRef{765} \MusicPDFRef{766} \MusicPDFRef{767} \MusicPDFRef{768} \MusicPDFRef{769} \MusicPDFRef{770}]}
\MusicPDFSetAnnotations{115}{[\MusicPDFRef{771} \MusicPDFRef{772}]}
\MusicPDFSetAnnotations{117}{[\MusicPDFRef{773} \MusicPDFRef{774} \MusicPDFRef{775} \MusicPDFRef{776}]}
\MusicPDFSetAnnotations{118}{[\MusicPDFRef{777} \MusicPDFRef{778} \MusicPDFRef{779} \MusicPDFRef{780} \MusicPDFRef{781} \MusicPDFRef{782} \MusicPDFRef{783} \MusicPDFRef{784} \MusicPDFRef{785} \MusicPDFRef{786} \MusicPDFRef{787} \MusicPDFRef{788} \MusicPDFRef{789} \MusicPDFRef{790} \MusicPDFRef{791} \MusicPDFRef{792} \MusicPDFRef{793} \MusicPDFRef{794} \MusicPDFRef{795} \MusicPDFRef{796} \MusicPDFRef{797} \MusicPDFRef{798} \MusicPDFRef{799} \MusicPDFRef{800} \MusicPDFRef{801} \MusicPDFRef{802}]}
\MusicPDFSetAnnotations{119}{[\MusicPDFRef{803} \MusicPDFRef{804} \MusicPDFRef{805}]}
\MusicPDFSetAnnotations{120}{[\MusicPDFRef{806} \MusicPDFRef{807} \MusicPDFRef{808} \MusicPDFRef{809} \MusicPDFRef{810} \MusicPDFRef{811} \MusicPDFRef{812} \MusicPDFRef{813} \MusicPDFRef{814}]}
\MusicPDFSetAnnotations{121}{[\MusicPDFRef{815} \MusicPDFRef{816}]}
\MusicPDFSetAnnotations{122}{[\MusicPDFRef{817} \MusicPDFRef{818}]}
\MusicPDFSetAnnotations{123}{[\MusicPDFRef{819} \MusicPDFRef{820} \MusicPDFRef{821} \MusicPDFRef{822} \MusicPDFRef{823} \MusicPDFRef{824} \MusicPDFRef{825} \MusicPDFRef{826} \MusicPDFRef{827}]}
\MusicPDFSetAnnotations{124}{[\MusicPDFRef{828} \MusicPDFRef{829} \MusicPDFRef{830} \MusicPDFRef{831} \MusicPDFRef{832} \MusicPDFRef{833} \MusicPDFRef{834} \MusicPDFRef{835} \MusicPDFRef{836} \MusicPDFRef{837} \MusicPDFRef{838} \MusicPDFRef{839} \MusicPDFRef{840}]}
\MusicPDFSetAnnotations{125}{[\MusicPDFRef{841} \MusicPDFRef{842} \MusicPDFRef{843} \MusicPDFRef{844} \MusicPDFRef{845} \MusicPDFRef{846} \MusicPDFRef{847} \MusicPDFRef{848} \MusicPDFRef{849}]}
\MusicPDFSetAnnotations{126}{[\MusicPDFRef{850} \MusicPDFRef{851} \MusicPDFRef{852} \MusicPDFRef{853} \MusicPDFRef{854} \MusicPDFRef{855} \MusicPDFRef{856} \MusicPDFRef{857} \MusicPDFRef{858} \MusicPDFRef{859} \MusicPDFRef{860} \MusicPDFRef{861} \MusicPDFRef{862} \MusicPDFRef{863} \MusicPDFRef{864} \MusicPDFRef{865} \MusicPDFRef{866}]}
\MusicPDFSetAnnotations{127}{[\MusicPDFRef{867} \MusicPDFRef{868} \MusicPDFRef{869} \MusicPDFRef{870} \MusicPDFRef{871} \MusicPDFRef{872} \MusicPDFRef{873} \MusicPDFRef{874} \MusicPDFRef{875} \MusicPDFRef{876} \MusicPDFRef{877} \MusicPDFRef{878} \MusicPDFRef{879} \MusicPDFRef{880} \MusicPDFRef{881} \MusicPDFRef{882} \MusicPDFRef{883} \MusicPDFRef{884} \MusicPDFRef{885}]}
\MusicPDFSetAnnotations{128}{[\MusicPDFRef{886} \MusicPDFRef{887} \MusicPDFRef{888} \MusicPDFRef{889} \MusicPDFRef{890} \MusicPDFRef{891} \MusicPDFRef{892} \MusicPDFRef{893} \MusicPDFRef{894} \MusicPDFRef{895} \MusicPDFRef{896} \MusicPDFRef{897} \MusicPDFRef{898} \MusicPDFRef{899} \MusicPDFRef{900} \MusicPDFRef{901}]}
\MusicPDFSetAnnotations{129}{[\MusicPDFRef{902} \MusicPDFRef{903} \MusicPDFRef{904} \MusicPDFRef{905} \MusicPDFRef{906} \MusicPDFRef{907}]}
\MusicPDFSetAnnotations{130}{[\MusicPDFRef{908} \MusicPDFRef{909} \MusicPDFRef{910} \MusicPDFRef{911} \MusicPDFRef{912} \MusicPDFRef{913} \MusicPDFRef{914} \MusicPDFRef{915} \MusicPDFRef{916} \MusicPDFRef{917} \MusicPDFRef{918} \MusicPDFRef{919} \MusicPDFRef{920} \MusicPDFRef{921} \MusicPDFRef{922} \MusicPDFRef{923}]}
\MusicPDFSetAnnotations{131}{[\MusicPDFRef{924} \MusicPDFRef{925} \MusicPDFRef{926} \MusicPDFRef{927} \MusicPDFRef{928} \MusicPDFRef{929} \MusicPDFRef{930} \MusicPDFRef{931} \MusicPDFRef{932} \MusicPDFRef{933} \MusicPDFRef{934} \MusicPDFRef{935} \MusicPDFRef{936} \MusicPDFRef{937} \MusicPDFRef{938} \MusicPDFRef{939} \MusicPDFRef{940} \MusicPDFRef{941}]}
\MusicPDFSetAnnotations{132}{[\MusicPDFRef{942} \MusicPDFRef{943} \MusicPDFRef{944} \MusicPDFRef{945} \MusicPDFRef{946} \MusicPDFRef{947} \MusicPDFRef{948} \MusicPDFRef{949} \MusicPDFRef{950} \MusicPDFRef{951} \MusicPDFRef{952} \MusicPDFRef{953}]}
\MusicPDFSetAnnotations{133}{[\MusicPDFRef{954} \MusicPDFRef{955} \MusicPDFRef{956} \MusicPDFRef{957}]}
\MusicPDFSetAnnotations{134}{[\MusicPDFRef{958} \MusicPDFRef{959} \MusicPDFRef{960} \MusicPDFRef{961} \MusicPDFRef{962} \MusicPDFRef{963} \MusicPDFRef{964} \MusicPDFRef{965}]}
\MusicPDFSetAnnotations{135}{[\MusicPDFRef{966} \MusicPDFRef{967} \MusicPDFRef{968}]}
\MusicPDFSetAnnotations{136}{[\MusicPDFRef{969} \MusicPDFRef{970} \MusicPDFRef{971} \MusicPDFRef{972} \MusicPDFRef{973} \MusicPDFRef{974} \MusicPDFRef{975} \MusicPDFRef{976} \MusicPDFRef{977} \MusicPDFRef{978} \MusicPDFRef{979} \MusicPDFRef{980} \MusicPDFRef{981} \MusicPDFRef{982} \MusicPDFRef{983} \MusicPDFRef{984} \MusicPDFRef{985} \MusicPDFRef{986} \MusicPDFRef{987} \MusicPDFRef{988}]}
\MusicPDFSetAnnotations{137}{[\MusicPDFRef{989} \MusicPDFRef{990} \MusicPDFRef{991} \MusicPDFRef{992} \MusicPDFRef{993} \MusicPDFRef{994} \MusicPDFRef{995} \MusicPDFRef{996} \MusicPDFRef{997} \MusicPDFRef{998} \MusicPDFRef{999} \MusicPDFRef{1000} \MusicPDFRef{1001} \MusicPDFRef{1002} \MusicPDFRef{1003} \MusicPDFRef{1004}]}
\MusicPDFSetAnnotations{138}{[\MusicPDFRef{1005} \MusicPDFRef{1006}]}
\MusicPDFSetAnnotations{139}{[\MusicPDFRef{1007} \MusicPDFRef{1008} \MusicPDFRef{1009} \MusicPDFRef{1010} \MusicPDFRef{1011} \MusicPDFRef{1012}]}
\MusicPDFSetAnnotations{140}{[\MusicPDFRef{1013} \MusicPDFRef{1014} \MusicPDFRef{1015}]}
\MusicPDFSetAnnotations{141}{[\MusicPDFRef{1016} \MusicPDFRef{1017} \MusicPDFRef{1018} \MusicPDFRef{1019} \MusicPDFRef{1020} \MusicPDFRef{1021}]}
\MusicPDFSetAnnotations{142}{[\MusicPDFRef{1022} \MusicPDFRef{1023} \MusicPDFRef{1024} \MusicPDFRef{1025} \MusicPDFRef{1026} \MusicPDFRef{1027} \MusicPDFRef{1028} \MusicPDFRef{1029} \MusicPDFRef{1030} \MusicPDFRef{1031} \MusicPDFRef{1032} \MusicPDFRef{1033} \MusicPDFRef{1034} \MusicPDFRef{1035} \MusicPDFRef{1036}]}
\MusicPDFSetAnnotations{143}{[\MusicPDFRef{1037} \MusicPDFRef{1038} \MusicPDFRef{1039}]}
\MusicPDFSetAnnotations{144}{[\MusicPDFRef{1040} \MusicPDFRef{1041} \MusicPDFRef{1042} \MusicPDFRef{1043} \MusicPDFRef{1044} \MusicPDFRef{1045} \MusicPDFRef{1046} \MusicPDFRef{1047} \MusicPDFRef{1048} \MusicPDFRef{1049}]}
\MusicPDFSetAnnotations{145}{[\MusicPDFRef{1050} \MusicPDFRef{1051} \MusicPDFRef{1052} \MusicPDFRef{1053} \MusicPDFRef{1054} \MusicPDFRef{1055} \MusicPDFRef{1056} \MusicPDFRef{1057} \MusicPDFRef{1058} \MusicPDFRef{1059} \MusicPDFRef{1060} \MusicPDFRef{1061} \MusicPDFRef{1062} \MusicPDFRef{1063} \MusicPDFRef{1064}]}
\MusicPDFSetAnnotations{146}{[\MusicPDFRef{1065} \MusicPDFRef{1066} \MusicPDFRef{1067} \MusicPDFRef{1068} \MusicPDFRef{1069} \MusicPDFRef{1070} \MusicPDFRef{1071} \MusicPDFRef{1072} \MusicPDFRef{1073} \MusicPDFRef{1074} \MusicPDFRef{1075} \MusicPDFRef{1076} \MusicPDFRef{1077} \MusicPDFRef{1078} \MusicPDFRef{1079} \MusicPDFRef{1080}]}
\MusicPDFSetAnnotations{147}{[\MusicPDFRef{1081} \MusicPDFRef{1082} \MusicPDFRef{1083} \MusicPDFRef{1084} \MusicPDFRef{1085} \MusicPDFRef{1086} \MusicPDFRef{1087} \MusicPDFRef{1088} \MusicPDFRef{1089} \MusicPDFRef{1090} \MusicPDFRef{1091} \MusicPDFRef{1092} \MusicPDFRef{1093}]}
\MusicPDFSetAnnotations{148}{[\MusicPDFRef{1094} \MusicPDFRef{1095} \MusicPDFRef{1096} \MusicPDFRef{1097} \MusicPDFRef{1098} \MusicPDFRef{1099} \MusicPDFRef{1100} \MusicPDFRef{1101} \MusicPDFRef{1102} \MusicPDFRef{1103} \MusicPDFRef{1104} \MusicPDFRef{1105} \MusicPDFRef{1106} \MusicPDFRef{1107} \MusicPDFRef{1108} \MusicPDFRef{1109} \MusicPDFRef{1110} \MusicPDFRef{1111}]}
\MusicPDFSetAnnotations{149}{[\MusicPDFRef{1112} \MusicPDFRef{1113} \MusicPDFRef{1114} \MusicPDFRef{1115} \MusicPDFRef{1116} \MusicPDFRef{1117} \MusicPDFRef{1118} \MusicPDFRef{1119} \MusicPDFRef{1120} \MusicPDFRef{1121} \MusicPDFRef{1122} \MusicPDFRef{1123} \MusicPDFRef{1124}]}
\MusicPDFSetAnnotations{150}{[\MusicPDFRef{1125} \MusicPDFRef{1126} \MusicPDFRef{1127} \MusicPDFRef{1128} \MusicPDFRef{1129} \MusicPDFRef{1130} \MusicPDFRef{1131} \MusicPDFRef{1132} \MusicPDFRef{1133} \MusicPDFRef{1134} \MusicPDFRef{1135} \MusicPDFRef{1136} \MusicPDFRef{1137} \MusicPDFRef{1138} \MusicPDFRef{1139} \MusicPDFRef{1140} \MusicPDFRef{1141} \MusicPDFRef{1142} \MusicPDFRef{1143} \MusicPDFRef{1144} \MusicPDFRef{1145} \MusicPDFRef{1146} \MusicPDFRef{1147} \MusicPDFRef{1148} \MusicPDFRef{1149}]}
\MusicPDFSetAnnotations{151}{[\MusicPDFRef{1150} \MusicPDFRef{1151} \MusicPDFRef{1152} \MusicPDFRef{1153} \MusicPDFRef{1154} \MusicPDFRef{1155} \MusicPDFRef{1156} \MusicPDFRef{1157} \MusicPDFRef{1158} \MusicPDFRef{1159} \MusicPDFRef{1160} \MusicPDFRef{1161} \MusicPDFRef{1162} \MusicPDFRef{1163} \MusicPDFRef{1164}]}
\MusicPDFSetAnnotations{153}{[\MusicPDFRef{1165} \MusicPDFRef{1166} \MusicPDFRef{1167} \MusicPDFRef{1168} \MusicPDFRef{1169} \MusicPDFRef{1170} \MusicPDFRef{1171} \MusicPDFRef{1172} \MusicPDFRef{1173} \MusicPDFRef{1174} \MusicPDFRef{1175} \MusicPDFRef{1176} \MusicPDFRef{1177} \MusicPDFRef{1178} \MusicPDFRef{1179} \MusicPDFRef{1180}]}
\MusicPDFSetAnnotations{155}{[\MusicPDFRef{1181} \MusicPDFRef{1182} \MusicPDFRef{1183} \MusicPDFRef{1184} \MusicPDFRef{1185} \MusicPDFRef{1186} \MusicPDFRef{1187} \MusicPDFRef{1188} \MusicPDFRef{1189} \MusicPDFRef{1190} \MusicPDFRef{1191} \MusicPDFRef{1192} \MusicPDFRef{1193} \MusicPDFRef{1194} \MusicPDFRef{1195} \MusicPDFRef{1196} \MusicPDFRef{1197} \MusicPDFRef{1198} \MusicPDFRef{1199} \MusicPDFRef{1200} \MusicPDFRef{1201} \MusicPDFRef{1202} \MusicPDFRef{1203} \MusicPDFRef{1204} \MusicPDFRef{1205} \MusicPDFRef{1206} \MusicPDFRef{1207} \MusicPDFRef{1208} \MusicPDFRef{1209} \MusicPDFRef{1210} \MusicPDFRef{1211} \MusicPDFRef{1212}]}
\MusicPDFSetAnnotations{156}{[\MusicPDFRef{1213} \MusicPDFRef{1214} \MusicPDFRef{1215} \MusicPDFRef{1216} \MusicPDFRef{1217} \MusicPDFRef{1218} \MusicPDFRef{1219} \MusicPDFRef{1220}]}
\MusicPDFSetAnnotations{157}{[\MusicPDFRef{1221} \MusicPDFRef{1222} \MusicPDFRef{1223} \MusicPDFRef{1224} \MusicPDFRef{1225} \MusicPDFRef{1226} \MusicPDFRef{1227} \MusicPDFRef{1228} \MusicPDFRef{1229} \MusicPDFRef{1230}]}
\MusicPDFSetAnnotations{158}{[\MusicPDFRef{1231} \MusicPDFRef{1232} \MusicPDFRef{1233} \MusicPDFRef{1234} \MusicPDFRef{1235} \MusicPDFRef{1236} \MusicPDFRef{1237} \MusicPDFRef{1238} \MusicPDFRef{1239} \MusicPDFRef{1240} \MusicPDFRef{1241} \MusicPDFRef{1242} \MusicPDFRef{1243}]}
\MusicPDFSetAnnotations{159}{[\MusicPDFRef{1244} \MusicPDFRef{1245} \MusicPDFRef{1246} \MusicPDFRef{1247} \MusicPDFRef{1248} \MusicPDFRef{1249} \MusicPDFRef{1250} \MusicPDFRef{1251} \MusicPDFRef{1252} \MusicPDFRef{1253}]}
\MusicPDFSetAnnotations{160}{[\MusicPDFRef{1254} \MusicPDFRef{1255} \MusicPDFRef{1256} \MusicPDFRef{1257} \MusicPDFRef{1258} \MusicPDFRef{1259} \MusicPDFRef{1260} \MusicPDFRef{1261} \MusicPDFRef{1262} \MusicPDFRef{1263} \MusicPDFRef{1264} \MusicPDFRef{1265} \MusicPDFRef{1266} \MusicPDFRef{1267} \MusicPDFRef{1268} \MusicPDFRef{1269} \MusicPDFRef{1270} \MusicPDFRef{1271} \MusicPDFRef{1272} \MusicPDFRef{1273} \MusicPDFRef{1274} \MusicPDFRef{1275}]}
\MusicPDFSetAnnotations{161}{[\MusicPDFRef{1276} \MusicPDFRef{1277} \MusicPDFRef{1278} \MusicPDFRef{1279} \MusicPDFRef{1280} \MusicPDFRef{1281} \MusicPDFRef{1282} \MusicPDFRef{1283} \MusicPDFRef{1284} \MusicPDFRef{1285} \MusicPDFRef{1286} \MusicPDFRef{1287} \MusicPDFRef{1288} \MusicPDFRef{1289} \MusicPDFRef{1290} \MusicPDFRef{1291} \MusicPDFRef{1292} \MusicPDFRef{1293} \MusicPDFRef{1294} \MusicPDFRef{1295} \MusicPDFRef{1296}]}
\MusicPDFSetAnnotations{162}{[\MusicPDFRef{1297} \MusicPDFRef{1298} \MusicPDFRef{1299}]}
\MusicPDFSetAnnotations{163}{[\MusicPDFRef{1300} \MusicPDFRef{1301} \MusicPDFRef{1302} \MusicPDFRef{1303} \MusicPDFRef{1304} \MusicPDFRef{1305}]}
\MusicPDFSetAnnotations{164}{[\MusicPDFRef{1306} \MusicPDFRef{1307} \MusicPDFRef{1308} \MusicPDFRef{1309} \MusicPDFRef{1310} \MusicPDFRef{1311} \MusicPDFRef{1312} \MusicPDFRef{1313} \MusicPDFRef{1314} \MusicPDFRef{1315} \MusicPDFRef{1316} \MusicPDFRef{1317} \MusicPDFRef{1318} \MusicPDFRef{1319} \MusicPDFRef{1320} \MusicPDFRef{1321} \MusicPDFRef{1322} \MusicPDFRef{1323} \MusicPDFRef{1324} \MusicPDFRef{1325} \MusicPDFRef{1326} \MusicPDFRef{1327} \MusicPDFRef{1328} \MusicPDFRef{1329} \MusicPDFRef{1330} \MusicPDFRef{1331} \MusicPDFRef{1332} \MusicPDFRef{1333} \MusicPDFRef{1334} \MusicPDFRef{1335}]}
\MusicPDFSetAnnotations{165}{[\MusicPDFRef{1336} \MusicPDFRef{1337} \MusicPDFRef{1338} \MusicPDFRef{1339} \MusicPDFRef{1340} \MusicPDFRef{1341} \MusicPDFRef{1342} \MusicPDFRef{1343} \MusicPDFRef{1344} \MusicPDFRef{1345} \MusicPDFRef{1346} \MusicPDFRef{1347} \MusicPDFRef{1348} \MusicPDFRef{1349} \MusicPDFRef{1350} \MusicPDFRef{1351} \MusicPDFRef{1352} \MusicPDFRef{1353}]}
\MusicPDFSetAnnotations{166}{[\MusicPDFRef{1354} \MusicPDFRef{1355} \MusicPDFRef{1356}]}
\MusicPDFSetAnnotations{167}{[\MusicPDFRef{1357} \MusicPDFRef{1358} \MusicPDFRef{1359} \MusicPDFRef{1360} \MusicPDFRef{1361}]}
\MusicPDFSetAnnotations{168}{[\MusicPDFRef{1362} \MusicPDFRef{1363} \MusicPDFRef{1364} \MusicPDFRef{1365} \MusicPDFRef{1366} \MusicPDFRef{1367} \MusicPDFRef{1368} \MusicPDFRef{1369} \MusicPDFRef{1370}]}
\MusicPDFSetAnnotations{169}{[\MusicPDFRef{1371} \MusicPDFRef{1372} \MusicPDFRef{1373} \MusicPDFRef{1374} \MusicPDFRef{1375} \MusicPDFRef{1376} \MusicPDFRef{1377} \MusicPDFRef{1378} \MusicPDFRef{1379} \MusicPDFRef{1380} \MusicPDFRef{1381} \MusicPDFRef{1382} \MusicPDFRef{1383} \MusicPDFRef{1384} \MusicPDFRef{1385} \MusicPDFRef{1386} \MusicPDFRef{1387} \MusicPDFRef{1388} \MusicPDFRef{1389} \MusicPDFRef{1390} \MusicPDFRef{1391} \MusicPDFRef{1392}]}
\MusicPDFSetAnnotations{170}{[\MusicPDFRef{1393} \MusicPDFRef{1394} \MusicPDFRef{1395} \MusicPDFRef{1396} \MusicPDFRef{1397} \MusicPDFRef{1398} \MusicPDFRef{1399} \MusicPDFRef{1400} \MusicPDFRef{1401} \MusicPDFRef{1402} \MusicPDFRef{1403}]}
\MusicPDFSetAnnotations{171}{[\MusicPDFRef{1404} \MusicPDFRef{1405} \MusicPDFRef{1406} \MusicPDFRef{1407} \MusicPDFRef{1408} \MusicPDFRef{1409} \MusicPDFRef{1410} \MusicPDFRef{1411}]}
\MusicPDFSetAnnotations{172}{[\MusicPDFRef{1412} \MusicPDFRef{1413} \MusicPDFRef{1414} \MusicPDFRef{1415} \MusicPDFRef{1416} \MusicPDFRef{1417} \MusicPDFRef{1418}]}
\MusicPDFSetAnnotations{173}{[\MusicPDFRef{1419} \MusicPDFRef{1420} \MusicPDFRef{1421} \MusicPDFRef{1422} \MusicPDFRef{1423} \MusicPDFRef{1424} \MusicPDFRef{1425} \MusicPDFRef{1426} \MusicPDFRef{1427} \MusicPDFRef{1428} \MusicPDFRef{1429}]}
\MusicPDFSetAnnotations{174}{[\MusicPDFRef{1430} \MusicPDFRef{1431} \MusicPDFRef{1432} \MusicPDFRef{1433} \MusicPDFRef{1434}]}
\MusicPDFSetAnnotations{175}{[\MusicPDFRef{1435} \MusicPDFRef{1436} \MusicPDFRef{1437} \MusicPDFRef{1438} \MusicPDFRef{1439} \MusicPDFRef{1440} \MusicPDFRef{1441} \MusicPDFRef{1442}]}
\MusicPDFSetAnnotations{176}{[\MusicPDFRef{1443} \MusicPDFRef{1444} \MusicPDFRef{1445} \MusicPDFRef{1446} \MusicPDFRef{1447} \MusicPDFRef{1448} \MusicPDFRef{1449} \MusicPDFRef{1450} \MusicPDFRef{1451} \MusicPDFRef{1452} \MusicPDFRef{1453} \MusicPDFRef{1454}]}
\MusicPDFSetAnnotations{177}{[\MusicPDFRef{1455} \MusicPDFRef{1456} \MusicPDFRef{1457} \MusicPDFRef{1458} \MusicPDFRef{1459} \MusicPDFRef{1460} \MusicPDFRef{1461} \MusicPDFRef{1462} \MusicPDFRef{1463} \MusicPDFRef{1464} \MusicPDFRef{1465} \MusicPDFRef{1466} \MusicPDFRef{1467} \MusicPDFRef{1468} \MusicPDFRef{1469} \MusicPDFRef{1470} \MusicPDFRef{1471} \MusicPDFRef{1472} \MusicPDFRef{1473} \MusicPDFRef{1474} \MusicPDFRef{1475} \MusicPDFRef{1476} \MusicPDFRef{1477} \MusicPDFRef{1478}]}
\MusicPDFSetAnnotations{178}{[\MusicPDFRef{1479} \MusicPDFRef{1480} \MusicPDFRef{1481} \MusicPDFRef{1482} \MusicPDFRef{1483} \MusicPDFRef{1484} \MusicPDFRef{1485} \MusicPDFRef{1486}]}
\MusicPDFSetAnnotations{179}{[\MusicPDFRef{1487} \MusicPDFRef{1488} \MusicPDFRef{1489} \MusicPDFRef{1490} \MusicPDFRef{1491} \MusicPDFRef{1492} \MusicPDFRef{1493} \MusicPDFRef{1494} \MusicPDFRef{1495} \MusicPDFRef{1496} \MusicPDFRef{1497} \MusicPDFRef{1498}]}
\MusicPDFSetAnnotations{180}{[\MusicPDFRef{1499} \MusicPDFRef{1500} \MusicPDFRef{1501} \MusicPDFRef{1502} \MusicPDFRef{1503} \MusicPDFRef{1504} \MusicPDFRef{1505} \MusicPDFRef{1506} \MusicPDFRef{1507} \MusicPDFRef{1508} \MusicPDFRef{1509} \MusicPDFRef{1510}]}
\MusicPDFSetAnnotations{181}{[\MusicPDFRef{1511} \MusicPDFRef{1512} \MusicPDFRef{1513} \MusicPDFRef{1514} \MusicPDFRef{1515} \MusicPDFRef{1516}]}
\MusicPDFSetAnnotations{182}{[\MusicPDFRef{1517} \MusicPDFRef{1518} \MusicPDFRef{1519} \MusicPDFRef{1520} \MusicPDFRef{1521} \MusicPDFRef{1522} \MusicPDFRef{1523} \MusicPDFRef{1524} \MusicPDFRef{1525} \MusicPDFRef{1526} \MusicPDFRef{1527}]}
\MusicPDFSetAnnotations{183}{[\MusicPDFRef{1528} \MusicPDFRef{1529} \MusicPDFRef{1530} \MusicPDFRef{1531}]}
\MusicPDFSetAnnotations{184}{[\MusicPDFRef{1532} \MusicPDFRef{1533} \MusicPDFRef{1534} \MusicPDFRef{1535} \MusicPDFRef{1536} \MusicPDFRef{1537} \MusicPDFRef{1538} \MusicPDFRef{1539} \MusicPDFRef{1540} \MusicPDFRef{1541} \MusicPDFRef{1542} \MusicPDFRef{1543}]}
\MusicPDFSetAnnotations{185}{[\MusicPDFRef{1544} \MusicPDFRef{1545} \MusicPDFRef{1546} \MusicPDFRef{1547} \MusicPDFRef{1548} \MusicPDFRef{1549} \MusicPDFRef{1550} \MusicPDFRef{1551} \MusicPDFRef{1552} \MusicPDFRef{1553} \MusicPDFRef{1554} \MusicPDFRef{1555}]}
\MusicPDFSetAnnotations{186}{[\MusicPDFRef{1556} \MusicPDFRef{1557} \MusicPDFRef{1558} \MusicPDFRef{1559} \MusicPDFRef{1560} \MusicPDFRef{1561} \MusicPDFRef{1562}]}
\MusicPDFSetAnnotations{187}{[\MusicPDFRef{1563} \MusicPDFRef{1564} \MusicPDFRef{1565} \MusicPDFRef{1566}]}
\MusicPDFSetAnnotations{188}{[\MusicPDFRef{1567} \MusicPDFRef{1568} \MusicPDFRef{1569} \MusicPDFRef{1570} \MusicPDFRef{1571} \MusicPDFRef{1572} \MusicPDFRef{1573} \MusicPDFRef{1574} \MusicPDFRef{1575} \MusicPDFRef{1576} \MusicPDFRef{1577}]}
\MusicPDFSetAnnotations{189}{[\MusicPDFRef{1578} \MusicPDFRef{1579} \MusicPDFRef{1580} \MusicPDFRef{1581} \MusicPDFRef{1582} \MusicPDFRef{1583} \MusicPDFRef{1584} \MusicPDFRef{1585} \MusicPDFRef{1586} \MusicPDFRef{1587} \MusicPDFRef{1588} \MusicPDFRef{1589} \MusicPDFRef{1590} \MusicPDFRef{1591} \MusicPDFRef{1592} \MusicPDFRef{1593}]}
\MusicPDFSetAnnotations{190}{[\MusicPDFRef{1594} \MusicPDFRef{1595} \MusicPDFRef{1596} \MusicPDFRef{1597} \MusicPDFRef{1598}]}
\MusicPDFSetAnnotations{191}{[\MusicPDFRef{1599} \MusicPDFRef{1600} \MusicPDFRef{1601} \MusicPDFRef{1602} \MusicPDFRef{1603} \MusicPDFRef{1604} \MusicPDFRef{1605} \MusicPDFRef{1606}]}
\MusicPDFSetAnnotations{192}{[\MusicPDFRef{1607} \MusicPDFRef{1608} \MusicPDFRef{1609}]}
\MusicPDFSetAnnotations{193}{[\MusicPDFRef{1610} \MusicPDFRef{1611} \MusicPDFRef{1612} \MusicPDFRef{1613} \MusicPDFRef{1614} \MusicPDFRef{1615} \MusicPDFRef{1616}]}
\MusicPDFSetAnnotations{194}{[\MusicPDFRef{1617} \MusicPDFRef{1618} \MusicPDFRef{1619} \MusicPDFRef{1620} \MusicPDFRef{1621} \MusicPDFRef{1622} \MusicPDFRef{1623} \MusicPDFRef{1624}]}
\MusicPDFSetAnnotations{195}{[\MusicPDFRef{1625} \MusicPDFRef{1626} \MusicPDFRef{1627} \MusicPDFRef{1628} \MusicPDFRef{1629} \MusicPDFRef{1630} \MusicPDFRef{1631} \MusicPDFRef{1632} \MusicPDFRef{1633}]}
\MusicPDFSetAnnotations{196}{[\MusicPDFRef{1634} \MusicPDFRef{1635} \MusicPDFRef{1636} \MusicPDFRef{1637} \MusicPDFRef{1638} \MusicPDFRef{1639} \MusicPDFRef{1640} \MusicPDFRef{1641} \MusicPDFRef{1642} \MusicPDFRef{1643} \MusicPDFRef{1644}]}
\MusicPDFSetAnnotations{197}{[\MusicPDFRef{1645} \MusicPDFRef{1646} \MusicPDFRef{1647} \MusicPDFRef{1648} \MusicPDFRef{1649} \MusicPDFRef{1650} \MusicPDFRef{1651} \MusicPDFRef{1652} \MusicPDFRef{1653} \MusicPDFRef{1654} \MusicPDFRef{1655} \MusicPDFRef{1656} \MusicPDFRef{1657}]}
\MusicPDFSetAnnotations{198}{[\MusicPDFRef{1658} \MusicPDFRef{1659} \MusicPDFRef{1660} \MusicPDFRef{1661} \MusicPDFRef{1662} \MusicPDFRef{1663} \MusicPDFRef{1664} \MusicPDFRef{1665} \MusicPDFRef{1666} \MusicPDFRef{1667} \MusicPDFRef{1668}]}
\MusicPDFSetAnnotations{199}{[\MusicPDFRef{1669} \MusicPDFRef{1670} \MusicPDFRef{1671} \MusicPDFRef{1672} \MusicPDFRef{1673} \MusicPDFRef{1674} \MusicPDFRef{1675} \MusicPDFRef{1676} \MusicPDFRef{1677} \MusicPDFRef{1678}]}
\MusicPDFSetAnnotations{200}{[\MusicPDFRef{1679} \MusicPDFRef{1680} \MusicPDFRef{1681} \MusicPDFRef{1682}]}
\MusicPDFSetAnnotations{201}{[\MusicPDFRef{1683} \MusicPDFRef{1684} \MusicPDFRef{1685} \MusicPDFRef{1686} \MusicPDFRef{1687} \MusicPDFRef{1688} \MusicPDFRef{1689} \MusicPDFRef{1690} \MusicPDFRef{1691} \MusicPDFRef{1692} \MusicPDFRef{1693} \MusicPDFRef{1694} \MusicPDFRef{1695} \MusicPDFRef{1696} \MusicPDFRef{1697} \MusicPDFRef{1698} \MusicPDFRef{1699}]}
\MusicPDFSetAnnotations{202}{[\MusicPDFRef{1700} \MusicPDFRef{1701} \MusicPDFRef{1702} \MusicPDFRef{1703} \MusicPDFRef{1704} \MusicPDFRef{1705} \MusicPDFRef{1706} \MusicPDFRef{1707} \MusicPDFRef{1708} \MusicPDFRef{1709} \MusicPDFRef{1710} \MusicPDFRef{1711} \MusicPDFRef{1712} \MusicPDFRef{1713} \MusicPDFRef{1714} \MusicPDFRef{1715}]}
\MusicPDFSetAnnotations{203}{[\MusicPDFRef{1716} \MusicPDFRef{1717} \MusicPDFRef{1718} \MusicPDFRef{1719} \MusicPDFRef{1720} \MusicPDFRef{1721} \MusicPDFRef{1722} \MusicPDFRef{1723} \MusicPDFRef{1724} \MusicPDFRef{1725} \MusicPDFRef{1726} \MusicPDFRef{1727} \MusicPDFRef{1728}]}
\MusicPDFSetAnnotations{204}{[\MusicPDFRef{1729} \MusicPDFRef{1730} \MusicPDFRef{1731} \MusicPDFRef{1732} \MusicPDFRef{1733} \MusicPDFRef{1734} \MusicPDFRef{1735} \MusicPDFRef{1736} \MusicPDFRef{1737} \MusicPDFRef{1738} \MusicPDFRef{1739} \MusicPDFRef{1740} \MusicPDFRef{1741} \MusicPDFRef{1742} \MusicPDFRef{1743}]}
\MusicPDFSetAnnotations{205}{[\MusicPDFRef{1744} \MusicPDFRef{1745} \MusicPDFRef{1746} \MusicPDFRef{1747} \MusicPDFRef{1748} \MusicPDFRef{1749} \MusicPDFRef{1750} \MusicPDFRef{1751} \MusicPDFRef{1752} \MusicPDFRef{1753} \MusicPDFRef{1754} \MusicPDFRef{1755}]}
\MusicPDFSetAnnotations{206}{[\MusicPDFRef{1756} \MusicPDFRef{1757} \MusicPDFRef{1758} \MusicPDFRef{1759} \MusicPDFRef{1760} \MusicPDFRef{1761} \MusicPDFRef{1762} \MusicPDFRef{1763} \MusicPDFRef{1764}]}
\MusicPDFSetAnnotations{207}{[\MusicPDFRef{1765} \MusicPDFRef{1766} \MusicPDFRef{1767} \MusicPDFRef{1768} \MusicPDFRef{1769} \MusicPDFRef{1770} \MusicPDFRef{1771} \MusicPDFRef{1772} \MusicPDFRef{1773} \MusicPDFRef{1774} \MusicPDFRef{1775} \MusicPDFRef{1776} \MusicPDFRef{1777}]}
\MusicPDFSetAnnotations{208}{[\MusicPDFRef{1778} \MusicPDFRef{1779} \MusicPDFRef{1780} \MusicPDFRef{1781} \MusicPDFRef{1782} \MusicPDFRef{1783} \MusicPDFRef{1784} \MusicPDFRef{1785} \MusicPDFRef{1786} \MusicPDFRef{1787} \MusicPDFRef{1788} \MusicPDFRef{1789} \MusicPDFRef{1790} \MusicPDFRef{1791} \MusicPDFRef{1792} \MusicPDFRef{1793} \MusicPDFRef{1794} \MusicPDFRef{1795} \MusicPDFRef{1796} \MusicPDFRef{1797}]}
\MusicPDFSetAnnotations{209}{[\MusicPDFRef{1798} \MusicPDFRef{1799} \MusicPDFRef{1800} \MusicPDFRef{1801}]}
\MusicPDFSetAnnotations{210}{[\MusicPDFRef{1802} \MusicPDFRef{1803} \MusicPDFRef{1804} \MusicPDFRef{1805} \MusicPDFRef{1806} \MusicPDFRef{1807} \MusicPDFRef{1808} \MusicPDFRef{1809} \MusicPDFRef{1810}]}
\MusicPDFSetAnnotations{211}{[\MusicPDFRef{1811} \MusicPDFRef{1812} \MusicPDFRef{1813} \MusicPDFRef{1814} \MusicPDFRef{1815} \MusicPDFRef{1816} \MusicPDFRef{1817} \MusicPDFRef{1818} \MusicPDFRef{1819} \MusicPDFRef{1820} \MusicPDFRef{1821}]}
\MusicPDFSetAnnotations{212}{[\MusicPDFRef{1822} \MusicPDFRef{1823} \MusicPDFRef{1824}]}
\MusicPDFSetAnnotations{213}{[\MusicPDFRef{1825} \MusicPDFRef{1826} \MusicPDFRef{1827} \MusicPDFRef{1828}]}
\MusicPDFSetAnnotations{214}{[\MusicPDFRef{1829} \MusicPDFRef{1830} \MusicPDFRef{1831} \MusicPDFRef{1832} \MusicPDFRef{1833} \MusicPDFRef{1834} \MusicPDFRef{1835} \MusicPDFRef{1836} \MusicPDFRef{1837} \MusicPDFRef{1838} \MusicPDFRef{1839}]}
\MusicPDFSetAnnotations{215}{[\MusicPDFRef{1840} \MusicPDFRef{1841} \MusicPDFRef{1842} \MusicPDFRef{1843} \MusicPDFRef{1844} \MusicPDFRef{1845} \MusicPDFRef{1846} \MusicPDFRef{1847} \MusicPDFRef{1848} \MusicPDFRef{1849}]}
\MusicPDFSetAnnotations{216}{[\MusicPDFRef{1850} \MusicPDFRef{1851} \MusicPDFRef{1852} \MusicPDFRef{1853} \MusicPDFRef{1854} \MusicPDFRef{1855} \MusicPDFRef{1856} \MusicPDFRef{1857} \MusicPDFRef{1858} \MusicPDFRef{1859}]}
\MusicPDFSetAnnotations{217}{[\MusicPDFRef{1860} \MusicPDFRef{1861} \MusicPDFRef{1862} \MusicPDFRef{1863} \MusicPDFRef{1864} \MusicPDFRef{1865} \MusicPDFRef{1866} \MusicPDFRef{1867} \MusicPDFRef{1868} \MusicPDFRef{1869} \MusicPDFRef{1870} \MusicPDFRef{1871} \MusicPDFRef{1872}]}
\MusicPDFSetAnnotations{218}{[\MusicPDFRef{1873} \MusicPDFRef{1874} \MusicPDFRef{1875} \MusicPDFRef{1876} \MusicPDFRef{1877} \MusicPDFRef{1878} \MusicPDFRef{1879} \MusicPDFRef{1880} \MusicPDFRef{1881}]}
\MusicPDFSetAnnotations{219}{[\MusicPDFRef{1882} \MusicPDFRef{1883} \MusicPDFRef{1884} \MusicPDFRef{1885} \MusicPDFRef{1886} \MusicPDFRef{1887} \MusicPDFRef{1888} \MusicPDFRef{1889} \MusicPDFRef{1890} \MusicPDFRef{1891} \MusicPDFRef{1892} \MusicPDFRef{1893} \MusicPDFRef{1894} \MusicPDFRef{1895}]}
\MusicPDFSetAnnotations{220}{[\MusicPDFRef{1896} \MusicPDFRef{1897} \MusicPDFRef{1898} \MusicPDFRef{1899} \MusicPDFRef{1900} \MusicPDFRef{1901} \MusicPDFRef{1902} \MusicPDFRef{1903} \MusicPDFRef{1904} \MusicPDFRef{1905} \MusicPDFRef{1906} \MusicPDFRef{1907} \MusicPDFRef{1908} \MusicPDFRef{1909} \MusicPDFRef{1910} \MusicPDFRef{1911} \MusicPDFRef{1912} \MusicPDFRef{1913} \MusicPDFRef{1914}]}
\MusicPDFSetAnnotations{221}{[\MusicPDFRef{1915} \MusicPDFRef{1916} \MusicPDFRef{1917} \MusicPDFRef{1918} \MusicPDFRef{1919} \MusicPDFRef{1920} \MusicPDFRef{1921} \MusicPDFRef{1922} \MusicPDFRef{1923} \MusicPDFRef{1924} \MusicPDFRef{1925} \MusicPDFRef{1926} \MusicPDFRef{1927} \MusicPDFRef{1928} \MusicPDFRef{1929} \MusicPDFRef{1930} \MusicPDFRef{1931} \MusicPDFRef{1932} \MusicPDFRef{1933} \MusicPDFRef{1934} \MusicPDFRef{1935}]}
\MusicPDFSetAnnotations{222}{[\MusicPDFRef{1936} \MusicPDFRef{1937} \MusicPDFRef{1938} \MusicPDFRef{1939} \MusicPDFRef{1940} \MusicPDFRef{1941} \MusicPDFRef{1942} \MusicPDFRef{1943} \MusicPDFRef{1944} \MusicPDFRef{1945} \MusicPDFRef{1946} \MusicPDFRef{1947} \MusicPDFRef{1948} \MusicPDFRef{1949} \MusicPDFRef{1950} \MusicPDFRef{1951} \MusicPDFRef{1952} \MusicPDFRef{1953} \MusicPDFRef{1954} \MusicPDFRef{1955} \MusicPDFRef{1956} \MusicPDFRef{1957} \MusicPDFRef{1958} \MusicPDFRef{1959} \MusicPDFRef{1960}]}
\MusicPDFSetAnnotations{223}{[\MusicPDFRef{1961} \MusicPDFRef{1962} \MusicPDFRef{1963} \MusicPDFRef{1964} \MusicPDFRef{1965} \MusicPDFRef{1966} \MusicPDFRef{1967} \MusicPDFRef{1968} \MusicPDFRef{1969} \MusicPDFRef{1970}]}
\MusicPDFSetAnnotations{224}{[\MusicPDFRef{1971} \MusicPDFRef{1972} \MusicPDFRef{1973} \MusicPDFRef{1974}]}
\MusicPDFSetAnnotations{225}{[\MusicPDFRef{1975} \MusicPDFRef{1976} \MusicPDFRef{1977} \MusicPDFRef{1978} \MusicPDFRef{1979} \MusicPDFRef{1980} \MusicPDFRef{1981} \MusicPDFRef{1982} \MusicPDFRef{1983} \MusicPDFRef{1984} \MusicPDFRef{1985}]}
\MusicPDFSetAnnotations{226}{[\MusicPDFRef{1986} \MusicPDFRef{1987} \MusicPDFRef{1988} \MusicPDFRef{1989} \MusicPDFRef{1990} \MusicPDFRef{1991} \MusicPDFRef{1992} \MusicPDFRef{1993} \MusicPDFRef{1994} \MusicPDFRef{1995} \MusicPDFRef{1996}]}
\MusicPDFSetAnnotations{227}{[\MusicPDFRef{1997} \MusicPDFRef{1998} \MusicPDFRef{1999} \MusicPDFRef{2000} \MusicPDFRef{2001} \MusicPDFRef{2002} \MusicPDFRef{2003} \MusicPDFRef{2004} \MusicPDFRef{2005} \MusicPDFRef{2006} \MusicPDFRef{2007} \MusicPDFRef{2008} \MusicPDFRef{2009} \MusicPDFRef{2010}]}
\MusicPDFSetAnnotations{228}{[\MusicPDFRef{2011} \MusicPDFRef{2012} \MusicPDFRef{2013} \MusicPDFRef{2014} \MusicPDFRef{2015} \MusicPDFRef{2016} \MusicPDFRef{2017} \MusicPDFRef{2018} \MusicPDFRef{2019} \MusicPDFRef{2020} \MusicPDFRef{2021} \MusicPDFRef{2022} \MusicPDFRef{2023} \MusicPDFRef{2024} \MusicPDFRef{2025}]}
\MusicPDFSetAnnotations{229}{[\MusicPDFRef{2026} \MusicPDFRef{2027} \MusicPDFRef{2028} \MusicPDFRef{2029} \MusicPDFRef{2030} \MusicPDFRef{2031} \MusicPDFRef{2032} \MusicPDFRef{2033} \MusicPDFRef{2034} \MusicPDFRef{2035} \MusicPDFRef{2036} \MusicPDFRef{2037} \MusicPDFRef{2038} \MusicPDFRef{2039} \MusicPDFRef{2040} \MusicPDFRef{2041} \MusicPDFRef{2042} \MusicPDFRef{2043} \MusicPDFRef{2044} \MusicPDFRef{2045} \MusicPDFRef{2046} \MusicPDFRef{2047} \MusicPDFRef{2048} \MusicPDFRef{2049} \MusicPDFRef{2050}]}
\MusicPDFSetAnnotations{230}{[\MusicPDFRef{2051} \MusicPDFRef{2052} \MusicPDFRef{2053} \MusicPDFRef{2054} \MusicPDFRef{2055} \MusicPDFRef{2056} \MusicPDFRef{2057} \MusicPDFRef{2058} \MusicPDFRef{2059} \MusicPDFRef{2060}]}
\MusicPDFSetAnnotations{231}{[\MusicPDFRef{2061} \MusicPDFRef{2062} \MusicPDFRef{2063} \MusicPDFRef{2064} \MusicPDFRef{2065} \MusicPDFRef{2066} \MusicPDFRef{2067} \MusicPDFRef{2068} \MusicPDFRef{2069} \MusicPDFRef{2070} \MusicPDFRef{2071} \MusicPDFRef{2072} \MusicPDFRef{2073} \MusicPDFRef{2074} \MusicPDFRef{2075} \MusicPDFRef{2076} \MusicPDFRef{2077} \MusicPDFRef{2078} \MusicPDFRef{2079} \MusicPDFRef{2080} \MusicPDFRef{2081} \MusicPDFRef{2082} \MusicPDFRef{2083} \MusicPDFRef{2084} \MusicPDFRef{2085} \MusicPDFRef{2086} \MusicPDFRef{2087} \MusicPDFRef{2088} \MusicPDFRef{2089}]}
\MusicPDFSetAnnotations{232}{[\MusicPDFRef{2090} \MusicPDFRef{2091} \MusicPDFRef{2092} \MusicPDFRef{2093} \MusicPDFRef{2094} \MusicPDFRef{2095} \MusicPDFRef{2096} \MusicPDFRef{2097} \MusicPDFRef{2098}]}
\MusicPDFSetAnnotations{233}{[\MusicPDFRef{2099} \MusicPDFRef{2100} \MusicPDFRef{2101} \MusicPDFRef{2102} \MusicPDFRef{2103} \MusicPDFRef{2104} \MusicPDFRef{2105} \MusicPDFRef{2106}]}
\MusicPDFSetAnnotations{234}{[\MusicPDFRef{2107} \MusicPDFRef{2108} \MusicPDFRef{2109} \MusicPDFRef{2110} \MusicPDFRef{2111} \MusicPDFRef{2112} \MusicPDFRef{2113} \MusicPDFRef{2114} \MusicPDFRef{2115} \MusicPDFRef{2116} \MusicPDFRef{2117} \MusicPDFRef{2118} \MusicPDFRef{2119} \MusicPDFRef{2120} \MusicPDFRef{2121} \MusicPDFRef{2122} \MusicPDFRef{2123} \MusicPDFRef{2124} \MusicPDFRef{2125} \MusicPDFRef{2126} \MusicPDFRef{2127} \MusicPDFRef{2128}]}
\MusicPDFSetAnnotations{235}{[\MusicPDFRef{2129} \MusicPDFRef{2130} \MusicPDFRef{2131} \MusicPDFRef{2132} \MusicPDFRef{2133} \MusicPDFRef{2134} \MusicPDFRef{2135} \MusicPDFRef{2136} \MusicPDFRef{2137} \MusicPDFRef{2138} \MusicPDFRef{2139} \MusicPDFRef{2140} \MusicPDFRef{2141} \MusicPDFRef{2142} \MusicPDFRef{2143} \MusicPDFRef{2144} \MusicPDFRef{2145} \MusicPDFRef{2146} \MusicPDFRef{2147} \MusicPDFRef{2148} \MusicPDFRef{2149} \MusicPDFRef{2150} \MusicPDFRef{2151} \MusicPDFRef{2152} \MusicPDFRef{2153} \MusicPDFRef{2154} \MusicPDFRef{2155}]}
\MusicPDFSetAnnotations{236}{[\MusicPDFRef{2156} \MusicPDFRef{2157} \MusicPDFRef{2158} \MusicPDFRef{2159} \MusicPDFRef{2160} \MusicPDFRef{2161}]}
\MusicPDFSetAnnotations{237}{[\MusicPDFRef{2162} \MusicPDFRef{2163} \MusicPDFRef{2164} \MusicPDFRef{2165} \MusicPDFRef{2166} \MusicPDFRef{2167} \MusicPDFRef{2168} \MusicPDFRef{2169} \MusicPDFRef{2170} \MusicPDFRef{2171} \MusicPDFRef{2172} \MusicPDFRef{2173}]}
\MusicPDFSetAnnotations{238}{[\MusicPDFRef{2174} \MusicPDFRef{2175} \MusicPDFRef{2176} \MusicPDFRef{2177} \MusicPDFRef{2178} \MusicPDFRef{2179} \MusicPDFRef{2180} \MusicPDFRef{2181} \MusicPDFRef{2182} \MusicPDFRef{2183}]}
\MusicPDFSetAnnotations{239}{[\MusicPDFRef{2184} \MusicPDFRef{2185} \MusicPDFRef{2186} \MusicPDFRef{2187} \MusicPDFRef{2188} \MusicPDFRef{2189} \MusicPDFRef{2190} \MusicPDFRef{2191} \MusicPDFRef{2192} \MusicPDFRef{2193}]}
\MusicPDFSetAnnotations{240}{[\MusicPDFRef{2194} \MusicPDFRef{2195} \MusicPDFRef{2196} \MusicPDFRef{2197} \MusicPDFRef{2198} \MusicPDFRef{2199}]}
\MusicPDFSetAnnotations{241}{[\MusicPDFRef{2200} \MusicPDFRef{2201} \MusicPDFRef{2202} \MusicPDFRef{2203} \MusicPDFRef{2204} \MusicPDFRef{2205} \MusicPDFRef{2206}]}
\MusicPDFSetAnnotations{242}{[\MusicPDFRef{2207} \MusicPDFRef{2208} \MusicPDFRef{2209} \MusicPDFRef{2210} \MusicPDFRef{2211} \MusicPDFRef{2212} \MusicPDFRef{2213} \MusicPDFRef{2214} \MusicPDFRef{2215} \MusicPDFRef{2216} \MusicPDFRef{2217} \MusicPDFRef{2218} \MusicPDFRef{2219} \MusicPDFRef{2220} \MusicPDFRef{2221} \MusicPDFRef{2222} \MusicPDFRef{2223} \MusicPDFRef{2224} \MusicPDFRef{2225} \MusicPDFRef{2226} \MusicPDFRef{2227} \MusicPDFRef{2228} \MusicPDFRef{2229}]}
\MusicPDFSetAnnotations{243}{[\MusicPDFRef{2230} \MusicPDFRef{2231} \MusicPDFRef{2232} \MusicPDFRef{2233} \MusicPDFRef{2234} \MusicPDFRef{2235} \MusicPDFRef{2236} \MusicPDFRef{2237} \MusicPDFRef{2238} \MusicPDFRef{2239} \MusicPDFRef{2240} \MusicPDFRef{2241}]}
\MusicPDFSetAnnotations{244}{[\MusicPDFRef{2242} \MusicPDFRef{2243} \MusicPDFRef{2244} \MusicPDFRef{2245} \MusicPDFRef{2246} \MusicPDFRef{2247} \MusicPDFRef{2248} \MusicPDFRef{2249} \MusicPDFRef{2250} \MusicPDFRef{2251} \MusicPDFRef{2252} \MusicPDFRef{2253} \MusicPDFRef{2254} \MusicPDFRef{2255} \MusicPDFRef{2256} \MusicPDFRef{2257} \MusicPDFRef{2258} \MusicPDFRef{2259}]}
\MusicPDFSetAnnotations{245}{[\MusicPDFRef{2260} \MusicPDFRef{2261} \MusicPDFRef{2262} \MusicPDFRef{2263} \MusicPDFRef{2264} \MusicPDFRef{2265} \MusicPDFRef{2266} \MusicPDFRef{2267} \MusicPDFRef{2268}]}
\MusicPDFSetAnnotations{246}{[\MusicPDFRef{2269} \MusicPDFRef{2270} \MusicPDFRef{2271} \MusicPDFRef{2272} \MusicPDFRef{2273} \MusicPDFRef{2274} \MusicPDFRef{2275} \MusicPDFRef{2276} \MusicPDFRef{2277} \MusicPDFRef{2278}]}
\MusicPDFSetAnnotations{247}{[\MusicPDFRef{2279} \MusicPDFRef{2280} \MusicPDFRef{2281} \MusicPDFRef{2282} \MusicPDFRef{2283} \MusicPDFRef{2284}]}
\MusicPDFSetAnnotations{248}{[\MusicPDFRef{2285} \MusicPDFRef{2286} \MusicPDFRef{2287} \MusicPDFRef{2288} \MusicPDFRef{2289} \MusicPDFRef{2290} \MusicPDFRef{2291} \MusicPDFRef{2292} \MusicPDFRef{2293} \MusicPDFRef{2294} \MusicPDFRef{2295} \MusicPDFRef{2296} \MusicPDFRef{2297} \MusicPDFRef{2298} \MusicPDFRef{2299} \MusicPDFRef{2300} \MusicPDFRef{2301}]}
\MusicPDFSetAnnotations{249}{[\MusicPDFRef{2302} \MusicPDFRef{2303} \MusicPDFRef{2304} \MusicPDFRef{2305} \MusicPDFRef{2306} \MusicPDFRef{2307} \MusicPDFRef{2308} \MusicPDFRef{2309} \MusicPDFRef{2310} \MusicPDFRef{2311}]}
\MusicPDFSetAnnotations{250}{[\MusicPDFRef{2312} \MusicPDFRef{2313} \MusicPDFRef{2314} \MusicPDFRef{2315} \MusicPDFRef{2316} \MusicPDFRef{2317} \MusicPDFRef{2318} \MusicPDFRef{2319} \MusicPDFRef{2320} \MusicPDFRef{2321} \MusicPDFRef{2322} \MusicPDFRef{2323}]}
\MusicPDFSetAnnotations{251}{[\MusicPDFRef{2324} \MusicPDFRef{2325} \MusicPDFRef{2326} \MusicPDFRef{2327}]}
\MusicPDFSetAnnotations{252}{[\MusicPDFRef{2328}]}
\MusicPDFSetAnnotations{253}{[\MusicPDFRef{2329} \MusicPDFRef{2330} \MusicPDFRef{2331} \MusicPDFRef{2332} \MusicPDFRef{2333} \MusicPDFRef{2334} \MusicPDFRef{2335} \MusicPDFRef{2336} \MusicPDFRef{2337} \MusicPDFRef{2338} \MusicPDFRef{2339} \MusicPDFRef{2340} \MusicPDFRef{2341} \MusicPDFRef{2342} \MusicPDFRef{2343} \MusicPDFRef{2344}]}
\MusicPDFSetAnnotations{254}{[\MusicPDFRef{2345} \MusicPDFRef{2346} \MusicPDFRef{2347} \MusicPDFRef{2348} \MusicPDFRef{2349} \MusicPDFRef{2350} \MusicPDFRef{2351} \MusicPDFRef{2352} \MusicPDFRef{2353} \MusicPDFRef{2354} \MusicPDFRef{2355} \MusicPDFRef{2356} \MusicPDFRef{2357} \MusicPDFRef{2358} \MusicPDFRef{2359} \MusicPDFRef{2360} \MusicPDFRef{2361}]}
\MusicPDFSetAnnotations{255}{[\MusicPDFRef{2362} \MusicPDFRef{2363} \MusicPDFRef{2364} \MusicPDFRef{2365} \MusicPDFRef{2366} \MusicPDFRef{2367} \MusicPDFRef{2368} \MusicPDFRef{2369} \MusicPDFRef{2370} \MusicPDFRef{2371} \MusicPDFRef{2372} \MusicPDFRef{2373} \MusicPDFRef{2374}]}
\MusicPDFSetAnnotations{256}{[\MusicPDFRef{2375} \MusicPDFRef{2376} \MusicPDFRef{2377} \MusicPDFRef{2378} \MusicPDFRef{2379} \MusicPDFRef{2380} \MusicPDFRef{2381} \MusicPDFRef{2382} \MusicPDFRef{2383} \MusicPDFRef{2384} \MusicPDFRef{2385} \MusicPDFRef{2386} \MusicPDFRef{2387} \MusicPDFRef{2388} \MusicPDFRef{2389} \MusicPDFRef{2390} \MusicPDFRef{2391}]}
\MusicPDFSetAnnotations{257}{[\MusicPDFRef{2392} \MusicPDFRef{2393} \MusicPDFRef{2394} \MusicPDFRef{2395} \MusicPDFRef{2396} \MusicPDFRef{2397} \MusicPDFRef{2398} \MusicPDFRef{2399}]}
\MusicPDFSetAnnotations{258}{[\MusicPDFRef{2400}]}
\MusicPDFSetAnnotations{259}{[\MusicPDFRef{2401}]}
\MusicPDFSetAnnotations{260}{[\MusicPDFRef{2402} \MusicPDFRef{2403} \MusicPDFRef{2404} \MusicPDFRef{2405} \MusicPDFRef{2406} \MusicPDFRef{2407} \MusicPDFRef{2408} \MusicPDFRef{2409} \MusicPDFRef{2410} \MusicPDFRef{2411} \MusicPDFRef{2412} \MusicPDFRef{2413} \MusicPDFRef{2414} \MusicPDFRef{2415} \MusicPDFRef{2416} \MusicPDFRef{2417} \MusicPDFRef{2418}]}
\MusicPDFSetAnnotations{261}{[\MusicPDFRef{2419} \MusicPDFRef{2420} \MusicPDFRef{2421} \MusicPDFRef{2422} \MusicPDFRef{2423} \MusicPDFRef{2424} \MusicPDFRef{2425} \MusicPDFRef{2426}]}
\MusicPDFSetAnnotations{262}{[\MusicPDFRef{2427} \MusicPDFRef{2428} \MusicPDFRef{2429} \MusicPDFRef{2430} \MusicPDFRef{2431} \MusicPDFRef{2432} \MusicPDFRef{2433}]}
\MusicPDFSetAnnotations{263}{[\MusicPDFRef{2434} \MusicPDFRef{2435} \MusicPDFRef{2436} \MusicPDFRef{2437} \MusicPDFRef{2438} \MusicPDFRef{2439} \MusicPDFRef{2440} \MusicPDFRef{2441} \MusicPDFRef{2442} \MusicPDFRef{2443}]}
\MusicPDFSetAnnotations{264}{[\MusicPDFRef{2444} \MusicPDFRef{2445} \MusicPDFRef{2446} \MusicPDFRef{2447} \MusicPDFRef{2448} \MusicPDFRef{2449} \MusicPDFRef{2450} \MusicPDFRef{2451} \MusicPDFRef{2452}]}
\MusicPDFSetAnnotations{265}{[\MusicPDFRef{2453} \MusicPDFRef{2454} \MusicPDFRef{2455} \MusicPDFRef{2456} \MusicPDFRef{2457} \MusicPDFRef{2458} \MusicPDFRef{2459} \MusicPDFRef{2460} \MusicPDFRef{2461} \MusicPDFRef{2462} \MusicPDFRef{2463} \MusicPDFRef{2464} \MusicPDFRef{2465}]}
\MusicPDFSetAnnotations{266}{[\MusicPDFRef{2466} \MusicPDFRef{2467} \MusicPDFRef{2468} \MusicPDFRef{2469} \MusicPDFRef{2470} \MusicPDFRef{2471} \MusicPDFRef{2472} \MusicPDFRef{2473} \MusicPDFRef{2474} \MusicPDFRef{2475} \MusicPDFRef{2476}]}
\MusicPDFSetAnnotations{267}{[\MusicPDFRef{2477} \MusicPDFRef{2478} \MusicPDFRef{2479} \MusicPDFRef{2480} \MusicPDFRef{2481} \MusicPDFRef{2482} \MusicPDFRef{2483} \MusicPDFRef{2484} \MusicPDFRef{2485} \MusicPDFRef{2486} \MusicPDFRef{2487} \MusicPDFRef{2488} \MusicPDFRef{2489} \MusicPDFRef{2490} \MusicPDFRef{2491} \MusicPDFRef{2492} \MusicPDFRef{2493} \MusicPDFRef{2494}]}
\MusicPDFSetAnnotations{268}{[\MusicPDFRef{2495} \MusicPDFRef{2496} \MusicPDFRef{2497} \MusicPDFRef{2498} \MusicPDFRef{2499} \MusicPDFRef{2500} \MusicPDFRef{2501} \MusicPDFRef{2502} \MusicPDFRef{2503} \MusicPDFRef{2504} \MusicPDFRef{2505} \MusicPDFRef{2506} \MusicPDFRef{2507} \MusicPDFRef{2508} \MusicPDFRef{2509} \MusicPDFRef{2510} \MusicPDFRef{2511} \MusicPDFRef{2512} \MusicPDFRef{2513}]}
\MusicPDFSetAnnotations{269}{[\MusicPDFRef{2514} \MusicPDFRef{2515} \MusicPDFRef{2516} \MusicPDFRef{2517} \MusicPDFRef{2518} \MusicPDFRef{2519} \MusicPDFRef{2520} \MusicPDFRef{2521} \MusicPDFRef{2522} \MusicPDFRef{2523} \MusicPDFRef{2524} \MusicPDFRef{2525} \MusicPDFRef{2526} \MusicPDFRef{2527} \MusicPDFRef{2528} \MusicPDFRef{2529} \MusicPDFRef{2530} \MusicPDFRef{2531} \MusicPDFRef{2532} \MusicPDFRef{2533} \MusicPDFRef{2534} \MusicPDFRef{2535}]}
\MusicPDFSetAnnotations{270}{[\MusicPDFRef{2536} \MusicPDFRef{2537} \MusicPDFRef{2538} \MusicPDFRef{2539} \MusicPDFRef{2540}]}
\MusicPDFSetAnnotations{271}{[\MusicPDFRef{2541} \MusicPDFRef{2542}]}
\MusicPDFSetAnnotations{272}{[\MusicPDFRef{2543} \MusicPDFRef{2544} \MusicPDFRef{2545} \MusicPDFRef{2546} \MusicPDFRef{2547} \MusicPDFRef{2548} \MusicPDFRef{2549}]}
\MusicPDFSetAnnotations{273}{[\MusicPDFRef{2550} \MusicPDFRef{2551} \MusicPDFRef{2552} \MusicPDFRef{2553} \MusicPDFRef{2554} \MusicPDFRef{2555} \MusicPDFRef{2556} \MusicPDFRef{2557}]}
\MusicPDFSetAnnotations{274}{[\MusicPDFRef{2558} \MusicPDFRef{2559} \MusicPDFRef{2560} \MusicPDFRef{2561} \MusicPDFRef{2562} \MusicPDFRef{2563} \MusicPDFRef{2564} \MusicPDFRef{2565} \MusicPDFRef{2566} \MusicPDFRef{2567} \MusicPDFRef{2568} \MusicPDFRef{2569} \MusicPDFRef{2570} \MusicPDFRef{2571} \MusicPDFRef{2572} \MusicPDFRef{2573} \MusicPDFRef{2574}]}
\MusicPDFSetAnnotations{275}{[\MusicPDFRef{2575} \MusicPDFRef{2576} \MusicPDFRef{2577} \MusicPDFRef{2578} \MusicPDFRef{2579} \MusicPDFRef{2580} \MusicPDFRef{2581} \MusicPDFRef{2582} \MusicPDFRef{2583} \MusicPDFRef{2584} \MusicPDFRef{2585}]}
\MusicPDFSetAnnotations{276}{[\MusicPDFRef{2586} \MusicPDFRef{2587} \MusicPDFRef{2588} \MusicPDFRef{2589}]}
\MusicPDFSetAnnotations{277}{[\MusicPDFRef{2590} \MusicPDFRef{2591} \MusicPDFRef{2592} \MusicPDFRef{2593} \MusicPDFRef{2594} \MusicPDFRef{2595} \MusicPDFRef{2596}]}
\MusicPDFSetAnnotations{278}{[\MusicPDFRef{2597} \MusicPDFRef{2598} \MusicPDFRef{2599} \MusicPDFRef{2600} \MusicPDFRef{2601} \MusicPDFRef{2602} \MusicPDFRef{2603} \MusicPDFRef{2604} \MusicPDFRef{2605} \MusicPDFRef{2606} \MusicPDFRef{2607}]}
\MusicPDFSetAnnotations{279}{[\MusicPDFRef{2608} \MusicPDFRef{2609} \MusicPDFRef{2610} \MusicPDFRef{2611} \MusicPDFRef{2612} \MusicPDFRef{2613} \MusicPDFRef{2614} \MusicPDFRef{2615} \MusicPDFRef{2616} \MusicPDFRef{2617} \MusicPDFRef{2618} \MusicPDFRef{2619} \MusicPDFRef{2620} \MusicPDFRef{2621} \MusicPDFRef{2622}]}
\MusicPDFSetAnnotations{280}{[\MusicPDFRef{2623} \MusicPDFRef{2624} \MusicPDFRef{2625} \MusicPDFRef{2626} \MusicPDFRef{2627} \MusicPDFRef{2628} \MusicPDFRef{2629} \MusicPDFRef{2630} \MusicPDFRef{2631} \MusicPDFRef{2632} \MusicPDFRef{2633} \MusicPDFRef{2634} \MusicPDFRef{2635} \MusicPDFRef{2636} \MusicPDFRef{2637} \MusicPDFRef{2638}]}
\MusicPDFSetAnnotations{281}{[\MusicPDFRef{2639} \MusicPDFRef{2640} \MusicPDFRef{2641} \MusicPDFRef{2642} \MusicPDFRef{2643} \MusicPDFRef{2644} \MusicPDFRef{2645} \MusicPDFRef{2646} \MusicPDFRef{2647} \MusicPDFRef{2648} \MusicPDFRef{2649} \MusicPDFRef{2650} \MusicPDFRef{2651} \MusicPDFRef{2652} \MusicPDFRef{2653}]}
\MusicPDFSetAnnotations{282}{[\MusicPDFRef{2654} \MusicPDFRef{2655} \MusicPDFRef{2656} \MusicPDFRef{2657} \MusicPDFRef{2658} \MusicPDFRef{2659}]}
\MusicPDFSetAnnotations{283}{[\MusicPDFRef{2660} \MusicPDFRef{2661} \MusicPDFRef{2662} \MusicPDFRef{2663} \MusicPDFRef{2664}]}
\MusicPDFSetAnnotations{284}{[\MusicPDFRef{2665} \MusicPDFRef{2666} \MusicPDFRef{2667} \MusicPDFRef{2668} \MusicPDFRef{2669} \MusicPDFRef{2670} \MusicPDFRef{2671} \MusicPDFRef{2672}]}
\MusicPDFSetAnnotations{285}{[\MusicPDFRef{2673} \MusicPDFRef{2674} \MusicPDFRef{2675} \MusicPDFRef{2676} \MusicPDFRef{2677} \MusicPDFRef{2678} \MusicPDFRef{2679} \MusicPDFRef{2680} \MusicPDFRef{2681} \MusicPDFRef{2682} \MusicPDFRef{2683} \MusicPDFRef{2684} \MusicPDFRef{2685} \MusicPDFRef{2686} \MusicPDFRef{2687} \MusicPDFRef{2688}]}
\MusicPDFSetAnnotations{287}{[\MusicPDFRef{2689} \MusicPDFRef{2690} \MusicPDFRef{2691} \MusicPDFRef{2692} \MusicPDFRef{2693} \MusicPDFRef{2694} \MusicPDFRef{2695} \MusicPDFRef{2696} \MusicPDFRef{2697} \MusicPDFRef{2698} \MusicPDFRef{2699}]}
\MusicPDFSetAnnotations{288}{[\MusicPDFRef{2700} \MusicPDFRef{2701} \MusicPDFRef{2702} \MusicPDFRef{2703} \MusicPDFRef{2704}]}
\MusicPDFSetAnnotations{289}{[\MusicPDFRef{2705} \MusicPDFRef{2706} \MusicPDFRef{2707} \MusicPDFRef{2708} \MusicPDFRef{2709} \MusicPDFRef{2710} \MusicPDFRef{2711}]}
\MusicPDFSetAnnotations{290}{[\MusicPDFRef{2712} \MusicPDFRef{2713} \MusicPDFRef{2714} \MusicPDFRef{2715} \MusicPDFRef{2716} \MusicPDFRef{2717} \MusicPDFRef{2718} \MusicPDFRef{2719} \MusicPDFRef{2720} \MusicPDFRef{2721}]}
\MusicPDFSetAnnotations{291}{[\MusicPDFRef{2722} \MusicPDFRef{2723} \MusicPDFRef{2724} \MusicPDFRef{2725} \MusicPDFRef{2726} \MusicPDFRef{2727}]}
\MusicPDFSetAnnotations{292}{[\MusicPDFRef{2728} \MusicPDFRef{2729} \MusicPDFRef{2730}]}
\MusicPDFSetAnnotations{293}{[\MusicPDFRef{2731} \MusicPDFRef{2732} \MusicPDFRef{2733} \MusicPDFRef{2734} \MusicPDFRef{2735} \MusicPDFRef{2736} \MusicPDFRef{2737} \MusicPDFRef{2738} \MusicPDFRef{2739} \MusicPDFRef{2740}]}
\MusicPDFSetAnnotations{294}{[\MusicPDFRef{2741} \MusicPDFRef{2742} \MusicPDFRef{2743} \MusicPDFRef{2744} \MusicPDFRef{2745} \MusicPDFRef{2746} \MusicPDFRef{2747} \MusicPDFRef{2748} \MusicPDFRef{2749} \MusicPDFRef{2750} \MusicPDFRef{2751}]}
\MusicPDFSetAnnotations{295}{[\MusicPDFRef{2752} \MusicPDFRef{2753} \MusicPDFRef{2754} \MusicPDFRef{2755} \MusicPDFRef{2756} \MusicPDFRef{2757} \MusicPDFRef{2758} \MusicPDFRef{2759} \MusicPDFRef{2760} \MusicPDFRef{2761}]}
\MusicPDFSetAnnotations{296}{[\MusicPDFRef{2762} \MusicPDFRef{2763} \MusicPDFRef{2764} \MusicPDFRef{2765} \MusicPDFRef{2766} \MusicPDFRef{2767}]}
\MusicPDFSetAnnotations{297}{[\MusicPDFRef{2768} \MusicPDFRef{2769} \MusicPDFRef{2770} \MusicPDFRef{2771} \MusicPDFRef{2772}]}
\MusicPDFSetAnnotations{298}{[\MusicPDFRef{2773} \MusicPDFRef{2774} \MusicPDFRef{2775} \MusicPDFRef{2776} \MusicPDFRef{2777} \MusicPDFRef{2778} \MusicPDFRef{2779} \MusicPDFRef{2780}]}
\MusicPDFSetAnnotations{299}{[\MusicPDFRef{2781} \MusicPDFRef{2782} \MusicPDFRef{2783} \MusicPDFRef{2784} \MusicPDFRef{2785} \MusicPDFRef{2786} \MusicPDFRef{2787} \MusicPDFRef{2788} \MusicPDFRef{2789} \MusicPDFRef{2790} \MusicPDFRef{2791} \MusicPDFRef{2792} \MusicPDFRef{2793} \MusicPDFRef{2794} \MusicPDFRef{2795} \MusicPDFRef{2796} \MusicPDFRef{2797} \MusicPDFRef{2798} \MusicPDFRef{2799} \MusicPDFRef{2800} \MusicPDFRef{2801} \MusicPDFRef{2802}]}
\MusicPDFSetAnnotations{300}{[\MusicPDFRef{2803} \MusicPDFRef{2804} \MusicPDFRef{2805} \MusicPDFRef{2806} \MusicPDFRef{2807} \MusicPDFRef{2808} \MusicPDFRef{2809} \MusicPDFRef{2810}]}
\MusicPDFSetAnnotations{301}{[\MusicPDFRef{2811} \MusicPDFRef{2812} \MusicPDFRef{2813} \MusicPDFRef{2814} \MusicPDFRef{2815} \MusicPDFRef{2816} \MusicPDFRef{2817} \MusicPDFRef{2818}]}
\MusicPDFSetAnnotations{302}{[\MusicPDFRef{2819} \MusicPDFRef{2820} \MusicPDFRef{2821} \MusicPDFRef{2822} \MusicPDFRef{2823} \MusicPDFRef{2824} \MusicPDFRef{2825} \MusicPDFRef{2826} \MusicPDFRef{2827} \MusicPDFRef{2828} \MusicPDFRef{2829} \MusicPDFRef{2830}]}
\MusicPDFSetAnnotations{303}{[\MusicPDFRef{2831} \MusicPDFRef{2832} \MusicPDFRef{2833} \MusicPDFRef{2834} \MusicPDFRef{2835} \MusicPDFRef{2836} \MusicPDFRef{2837} \MusicPDFRef{2838} \MusicPDFRef{2839} \MusicPDFRef{2840} \MusicPDFRef{2841}]}
\MusicPDFSetAnnotations{304}{[\MusicPDFRef{2842} \MusicPDFRef{2843} \MusicPDFRef{2844} \MusicPDFRef{2845} \MusicPDFRef{2846} \MusicPDFRef{2847}]}
\MusicPDFSetAnnotations{305}{[\MusicPDFRef{2848} \MusicPDFRef{2849} \MusicPDFRef{2850} \MusicPDFRef{2851} \MusicPDFRef{2852} \MusicPDFRef{2853} \MusicPDFRef{2854} \MusicPDFRef{2855} \MusicPDFRef{2856} \MusicPDFRef{2857} \MusicPDFRef{2858}]}
\MusicPDFSetAnnotations{306}{[\MusicPDFRef{2859} \MusicPDFRef{2860} \MusicPDFRef{2861} \MusicPDFRef{2862} \MusicPDFRef{2863} \MusicPDFRef{2864} \MusicPDFRef{2865} \MusicPDFRef{2866} \MusicPDFRef{2867} \MusicPDFRef{2868} \MusicPDFRef{2869} \MusicPDFRef{2870}]}
\MusicPDFSetAnnotations{307}{[\MusicPDFRef{2871} \MusicPDFRef{2872} \MusicPDFRef{2873} \MusicPDFRef{2874} \MusicPDFRef{2875} \MusicPDFRef{2876} \MusicPDFRef{2877} \MusicPDFRef{2878} \MusicPDFRef{2879} \MusicPDFRef{2880} \MusicPDFRef{2881} \MusicPDFRef{2882}]}
\MusicPDFSetAnnotations{308}{[\MusicPDFRef{2883} \MusicPDFRef{2884} \MusicPDFRef{2885} \MusicPDFRef{2886} \MusicPDFRef{2887} \MusicPDFRef{2888} \MusicPDFRef{2889} \MusicPDFRef{2890} \MusicPDFRef{2891} \MusicPDFRef{2892}]}
\MusicPDFSetAnnotations{309}{[\MusicPDFRef{2893} \MusicPDFRef{2894} \MusicPDFRef{2895} \MusicPDFRef{2896} \MusicPDFRef{2897} \MusicPDFRef{2898} \MusicPDFRef{2899} \MusicPDFRef{2900} \MusicPDFRef{2901} \MusicPDFRef{2902} \MusicPDFRef{2903} \MusicPDFRef{2904} \MusicPDFRef{2905} \MusicPDFRef{2906} \MusicPDFRef{2907} \MusicPDFRef{2908} \MusicPDFRef{2909}]}
\MusicPDFSetAnnotations{310}{[\MusicPDFRef{2910} \MusicPDFRef{2911} \MusicPDFRef{2912} \MusicPDFRef{2913}]}
\MusicPDFSetAnnotations{311}{[\MusicPDFRef{2914} \MusicPDFRef{2915} \MusicPDFRef{2916} \MusicPDFRef{2917} \MusicPDFRef{2918} \MusicPDFRef{2919} \MusicPDFRef{2920} \MusicPDFRef{2921} \MusicPDFRef{2922} \MusicPDFRef{2923} \MusicPDFRef{2924} \MusicPDFRef{2925} \MusicPDFRef{2926} \MusicPDFRef{2927} \MusicPDFRef{2928} \MusicPDFRef{2929} \MusicPDFRef{2930} \MusicPDFRef{2931} \MusicPDFRef{2932} \MusicPDFRef{2933}]}
\MusicPDFSetAnnotations{312}{[\MusicPDFRef{2934} \MusicPDFRef{2935} \MusicPDFRef{2936} \MusicPDFRef{2937} \MusicPDFRef{2938} \MusicPDFRef{2939} \MusicPDFRef{2940} \MusicPDFRef{2941} \MusicPDFRef{2942} \MusicPDFRef{2943} \MusicPDFRef{2944} \MusicPDFRef{2945} \MusicPDFRef{2946} \MusicPDFRef{2947}]}
\MusicPDFSetAnnotations{313}{[\MusicPDFRef{2948} \MusicPDFRef{2949} \MusicPDFRef{2950} \MusicPDFRef{2951} \MusicPDFRef{2952} \MusicPDFRef{2953} \MusicPDFRef{2954} \MusicPDFRef{2955} \MusicPDFRef{2956} \MusicPDFRef{2957} \MusicPDFRef{2958} \MusicPDFRef{2959} \MusicPDFRef{2960} \MusicPDFRef{2961} \MusicPDFRef{2962} \MusicPDFRef{2963} \MusicPDFRef{2964}]}
\MusicPDFSetAnnotations{314}{[\MusicPDFRef{2965} \MusicPDFRef{2966} \MusicPDFRef{2967} \MusicPDFRef{2968} \MusicPDFRef{2969} \MusicPDFRef{2970} \MusicPDFRef{2971} \MusicPDFRef{2972}]}
\MusicPDFSetAnnotations{315}{[\MusicPDFRef{2973} \MusicPDFRef{2974} \MusicPDFRef{2975} \MusicPDFRef{2976} \MusicPDFRef{2977} \MusicPDFRef{2978} \MusicPDFRef{2979}]}
\MusicPDFSetAnnotations{316}{[\MusicPDFRef{2980} \MusicPDFRef{2981} \MusicPDFRef{2982} \MusicPDFRef{2983} \MusicPDFRef{2984} \MusicPDFRef{2985} \MusicPDFRef{2986} \MusicPDFRef{2987} \MusicPDFRef{2988} \MusicPDFRef{2989}]}
\MusicPDFSetAnnotations{317}{[\MusicPDFRef{2990} \MusicPDFRef{2991} \MusicPDFRef{2992} \MusicPDFRef{2993} \MusicPDFRef{2994}]}
\MusicPDFSetAnnotations{318}{[\MusicPDFRef{2995} \MusicPDFRef{2996} \MusicPDFRef{2997} \MusicPDFRef{2998} \MusicPDFRef{2999}]}
\MusicPDFSetAnnotations{319}{[\MusicPDFRef{3000} \MusicPDFRef{3001} \MusicPDFRef{3002} \MusicPDFRef{3003} \MusicPDFRef{3004} \MusicPDFRef{3005} \MusicPDFRef{3006} \MusicPDFRef{3007} \MusicPDFRef{3008} \MusicPDFRef{3009} \MusicPDFRef{3010} \MusicPDFRef{3011} \MusicPDFRef{3012} \MusicPDFRef{3013} \MusicPDFRef{3014} \MusicPDFRef{3015} \MusicPDFRef{3016} \MusicPDFRef{3017} \MusicPDFRef{3018}]}
\MusicPDFSetAnnotations{320}{[\MusicPDFRef{3019} \MusicPDFRef{3020} \MusicPDFRef{3021} \MusicPDFRef{3022} \MusicPDFRef{3023} \MusicPDFRef{3024} \MusicPDFRef{3025} \MusicPDFRef{3026} \MusicPDFRef{3027} \MusicPDFRef{3028} \MusicPDFRef{3029} \MusicPDFRef{3030} \MusicPDFRef{3031} \MusicPDFRef{3032} \MusicPDFRef{3033} \MusicPDFRef{3034} \MusicPDFRef{3035} \MusicPDFRef{3036} \MusicPDFRef{3037} \MusicPDFRef{3038} \MusicPDFRef{3039}]}
\MusicPDFSetAnnotations{321}{[\MusicPDFRef{3040} \MusicPDFRef{3041} \MusicPDFRef{3042} \MusicPDFRef{3043} \MusicPDFRef{3044} \MusicPDFRef{3045} \MusicPDFRef{3046} \MusicPDFRef{3047} \MusicPDFRef{3048} \MusicPDFRef{3049} \MusicPDFRef{3050} \MusicPDFRef{3051} \MusicPDFRef{3052} \MusicPDFRef{3053} \MusicPDFRef{3054} \MusicPDFRef{3055} \MusicPDFRef{3056} \MusicPDFRef{3057} \MusicPDFRef{3058} \MusicPDFRef{3059} \MusicPDFRef{3060} \MusicPDFRef{3061} \MusicPDFRef{3062}]}
\MusicPDFSetAnnotations{322}{[\MusicPDFRef{3063} \MusicPDFRef{3064} \MusicPDFRef{3065} \MusicPDFRef{3066} \MusicPDFRef{3067} \MusicPDFRef{3068} \MusicPDFRef{3069} \MusicPDFRef{3070} \MusicPDFRef{3071} \MusicPDFRef{3072} \MusicPDFRef{3073} \MusicPDFRef{3074} \MusicPDFRef{3075}]}
\MusicPDFSetAnnotations{323}{[\MusicPDFRef{3076} \MusicPDFRef{3077} \MusicPDFRef{3078} \MusicPDFRef{3079} \MusicPDFRef{3080}]}
\MusicPDFSetAnnotations{324}{[\MusicPDFRef{3081} \MusicPDFRef{3082} \MusicPDFRef{3083} \MusicPDFRef{3084}]}
\MusicPDFSetAnnotations{325}{[\MusicPDFRef{3085} \MusicPDFRef{3086} \MusicPDFRef{3087} \MusicPDFRef{3088}]}
\MusicPDFSetAnnotations{326}{[\MusicPDFRef{3089} \MusicPDFRef{3090} \MusicPDFRef{3091} \MusicPDFRef{3092} \MusicPDFRef{3093} \MusicPDFRef{3094} \MusicPDFRef{3095} \MusicPDFRef{3096} \MusicPDFRef{3097} \MusicPDFRef{3098}]}
\MusicPDFSetAnnotations{327}{[\MusicPDFRef{3099} \MusicPDFRef{3100} \MusicPDFRef{3101} \MusicPDFRef{3102} \MusicPDFRef{3103} \MusicPDFRef{3104} \MusicPDFRef{3105} \MusicPDFRef{3106} \MusicPDFRef{3107} \MusicPDFRef{3108}]}
\MusicPDFSetAnnotations{328}{[\MusicPDFRef{3109} \MusicPDFRef{3110} \MusicPDFRef{3111} \MusicPDFRef{3112}]}
\MusicPDFSetAnnotations{329}{[\MusicPDFRef{3113} \MusicPDFRef{3114} \MusicPDFRef{3115} \MusicPDFRef{3116} \MusicPDFRef{3117}]}
\MusicPDFSetAnnotations{330}{[\MusicPDFRef{3118} \MusicPDFRef{3119} \MusicPDFRef{3120} \MusicPDFRef{3121} \MusicPDFRef{3122} \MusicPDFRef{3123}]}
\MusicPDFSetAnnotations{331}{[\MusicPDFRef{3124} \MusicPDFRef{3125} \MusicPDFRef{3126} \MusicPDFRef{3127} \MusicPDFRef{3128}]}
\MusicPDFSetAnnotations{332}{[\MusicPDFRef{3129} \MusicPDFRef{3130} \MusicPDFRef{3131} \MusicPDFRef{3132} \MusicPDFRef{3133} \MusicPDFRef{3134} \MusicPDFRef{3135} \MusicPDFRef{3136} \MusicPDFRef{3137} \MusicPDFRef{3138} \MusicPDFRef{3139}]}
\MusicPDFSetAnnotations{333}{[\MusicPDFRef{3140} \MusicPDFRef{3141} \MusicPDFRef{3142} \MusicPDFRef{3143} \MusicPDFRef{3144} \MusicPDFRef{3145} \MusicPDFRef{3146} \MusicPDFRef{3147} \MusicPDFRef{3148} \MusicPDFRef{3149} \MusicPDFRef{3150} \MusicPDFRef{3151} \MusicPDFRef{3152} \MusicPDFRef{3153}]}
\MusicPDFSetAnnotations{334}{[\MusicPDFRef{3154} \MusicPDFRef{3155} \MusicPDFRef{3156} \MusicPDFRef{3157} \MusicPDFRef{3158} \MusicPDFRef{3159} \MusicPDFRef{3160} \MusicPDFRef{3161} \MusicPDFRef{3162}]}
\MusicPDFSetAnnotations{335}{[\MusicPDFRef{3163} \MusicPDFRef{3164} \MusicPDFRef{3165} \MusicPDFRef{3166} \MusicPDFRef{3167} \MusicPDFRef{3168}]}
\MusicPDFSetAnnotations{336}{[\MusicPDFRef{3169} \MusicPDFRef{3170} \MusicPDFRef{3171} \MusicPDFRef{3172} \MusicPDFRef{3173} \MusicPDFRef{3174} \MusicPDFRef{3175} \MusicPDFRef{3176} \MusicPDFRef{3177}]}
\MusicPDFSetAnnotations{337}{[\MusicPDFRef{3178} \MusicPDFRef{3179} \MusicPDFRef{3180} \MusicPDFRef{3181} \MusicPDFRef{3182} \MusicPDFRef{3183} \MusicPDFRef{3184} \MusicPDFRef{3185} \MusicPDFRef{3186} \MusicPDFRef{3187} \MusicPDFRef{3188} \MusicPDFRef{3189} \MusicPDFRef{3190}]}
\MusicPDFSetAnnotations{338}{[\MusicPDFRef{3191} \MusicPDFRef{3192} \MusicPDFRef{3193} \MusicPDFRef{3194} \MusicPDFRef{3195} \MusicPDFRef{3196} \MusicPDFRef{3197} \MusicPDFRef{3198} \MusicPDFRef{3199} \MusicPDFRef{3200} \MusicPDFRef{3201}]}
\MusicPDFSetAnnotations{339}{[\MusicPDFRef{3202} \MusicPDFRef{3203} \MusicPDFRef{3204} \MusicPDFRef{3205} \MusicPDFRef{3206} \MusicPDFRef{3207} \MusicPDFRef{3208} \MusicPDFRef{3209} \MusicPDFRef{3210} \MusicPDFRef{3211} \MusicPDFRef{3212}]}
\MusicPDFSetAnnotations{340}{[\MusicPDFRef{3213} \MusicPDFRef{3214} \MusicPDFRef{3215} \MusicPDFRef{3216} \MusicPDFRef{3217} \MusicPDFRef{3218}]}
\MusicPDFSetAnnotations{341}{[\MusicPDFRef{3219} \MusicPDFRef{3220} \MusicPDFRef{3221} \MusicPDFRef{3222} \MusicPDFRef{3223}]}
\MusicPDFSetAnnotations{342}{[\MusicPDFRef{3224} \MusicPDFRef{3225} \MusicPDFRef{3226} \MusicPDFRef{3227}]}
\MusicPDFSetAnnotations{343}{[\MusicPDFRef{3228} \MusicPDFRef{3229} \MusicPDFRef{3230} \MusicPDFRef{3231} \MusicPDFRef{3232} \MusicPDFRef{3233}]}
\MusicPDFSetAnnotations{344}{[\MusicPDFRef{3234} \MusicPDFRef{3235} \MusicPDFRef{3236} \MusicPDFRef{3237} \MusicPDFRef{3238} \MusicPDFRef{3239} \MusicPDFRef{3240} \MusicPDFRef{3241} \MusicPDFRef{3242} \MusicPDFRef{3243}]}
\MusicPDFSetAnnotations{345}{[\MusicPDFRef{3244} \MusicPDFRef{3245} \MusicPDFRef{3246} \MusicPDFRef{3247} \MusicPDFRef{3248} \MusicPDFRef{3249} \MusicPDFRef{3250} \MusicPDFRef{3251}]}
\MusicPDFSetAnnotations{346}{[\MusicPDFRef{3252} \MusicPDFRef{3253} \MusicPDFRef{3254} \MusicPDFRef{3255} \MusicPDFRef{3256} \MusicPDFRef{3257} \MusicPDFRef{3258} \MusicPDFRef{3259} \MusicPDFRef{3260} \MusicPDFRef{3261} \MusicPDFRef{3262} \MusicPDFRef{3263} \MusicPDFRef{3264} \MusicPDFRef{3265} \MusicPDFRef{3266} \MusicPDFRef{3267} \MusicPDFRef{3268} \MusicPDFRef{3269} \MusicPDFRef{3270} \MusicPDFRef{3271} \MusicPDFRef{3272} \MusicPDFRef{3273} \MusicPDFRef{3274} \MusicPDFRef{3275} \MusicPDFRef{3276} \MusicPDFRef{3277} \MusicPDFRef{3278} \MusicPDFRef{3279} \MusicPDFRef{3280} \MusicPDFRef{3281}]}
\MusicPDFSetAnnotations{347}{[\MusicPDFRef{3282} \MusicPDFRef{3283} \MusicPDFRef{3284} \MusicPDFRef{3285} \MusicPDFRef{3286} \MusicPDFRef{3287} \MusicPDFRef{3288}]}
\MusicPDFSetAnnotations{348}{[\MusicPDFRef{3289} \MusicPDFRef{3290} \MusicPDFRef{3291} \MusicPDFRef{3292} \MusicPDFRef{3293} \MusicPDFRef{3294} \MusicPDFRef{3295} \MusicPDFRef{3296} \MusicPDFRef{3297} \MusicPDFRef{3298} \MusicPDFRef{3299} \MusicPDFRef{3300} \MusicPDFRef{3301} \MusicPDFRef{3302} \MusicPDFRef{3303} \MusicPDFRef{3304} \MusicPDFRef{3305} \MusicPDFRef{3306} \MusicPDFRef{3307}]}
\MusicPDFSetAnnotations{349}{[\MusicPDFRef{3308} \MusicPDFRef{3309} \MusicPDFRef{3310} \MusicPDFRef{3311} \MusicPDFRef{3312} \MusicPDFRef{3313} \MusicPDFRef{3314} \MusicPDFRef{3315} \MusicPDFRef{3316} \MusicPDFRef{3317} \MusicPDFRef{3318} \MusicPDFRef{3319} \MusicPDFRef{3320}]}
\MusicPDFSetAnnotations{350}{[\MusicPDFRef{3321} \MusicPDFRef{3322} \MusicPDFRef{3323} \MusicPDFRef{3324} \MusicPDFRef{3325} \MusicPDFRef{3326} \MusicPDFRef{3327} \MusicPDFRef{3328} \MusicPDFRef{3329} \MusicPDFRef{3330} \MusicPDFRef{3331} \MusicPDFRef{3332} \MusicPDFRef{3333} \MusicPDFRef{3334} \MusicPDFRef{3335} \MusicPDFRef{3336}]}
\MusicPDFSetAnnotations{351}{[\MusicPDFRef{3337} \MusicPDFRef{3338} \MusicPDFRef{3339}]}
\MusicPDFSetAnnotations{352}{[\MusicPDFRef{3340}]}
\MusicPDFSetAnnotations{353}{[\MusicPDFRef{3341} \MusicPDFRef{3342} \MusicPDFRef{3343} \MusicPDFRef{3344} \MusicPDFRef{3345} \MusicPDFRef{3346} \MusicPDFRef{3347} \MusicPDFRef{3348} \MusicPDFRef{3349}]}
\MusicPDFSetAnnotations{354}{[\MusicPDFRef{3350} \MusicPDFRef{3351} \MusicPDFRef{3352}]}
\MusicPDFSetAnnotations{355}{[\MusicPDFRef{3353} \MusicPDFRef{3354} \MusicPDFRef{3355}]}
\MusicPDFSetAnnotations{356}{[\MusicPDFRef{3356} \MusicPDFRef{3357} \MusicPDFRef{3358} \MusicPDFRef{3359} \MusicPDFRef{3360} \MusicPDFRef{3361} \MusicPDFRef{3362}]}
\MusicPDFSetAnnotations{357}{[\MusicPDFRef{3363} \MusicPDFRef{3364} \MusicPDFRef{3365} \MusicPDFRef{3366} \MusicPDFRef{3367} \MusicPDFRef{3368} \MusicPDFRef{3369} \MusicPDFRef{3370} \MusicPDFRef{3371}]}
\MusicPDFSetAnnotations{358}{[\MusicPDFRef{3372}]}
\MusicPDFSetAnnotations{359}{[\MusicPDFRef{3373} \MusicPDFRef{3374} \MusicPDFRef{3375} \MusicPDFRef{3376} \MusicPDFRef{3377} \MusicPDFRef{3378} \MusicPDFRef{3379}]}
\MusicPDFSetAnnotations{360}{[\MusicPDFRef{3380} \MusicPDFRef{3381} \MusicPDFRef{3382} \MusicPDFRef{3383} \MusicPDFRef{3384} \MusicPDFRef{3385}]}
\MusicPDFSetAnnotations{361}{[\MusicPDFRef{3386} \MusicPDFRef{3387} \MusicPDFRef{3388} \MusicPDFRef{3389} \MusicPDFRef{3390} \MusicPDFRef{3391} \MusicPDFRef{3392}]}
\MusicPDFSetAnnotations{362}{[\MusicPDFRef{3393} \MusicPDFRef{3394} \MusicPDFRef{3395} \MusicPDFRef{3396} \MusicPDFRef{3397} \MusicPDFRef{3398} \MusicPDFRef{3399} \MusicPDFRef{3400} \MusicPDFRef{3401} \MusicPDFRef{3402} \MusicPDFRef{3403} \MusicPDFRef{3404} \MusicPDFRef{3405} \MusicPDFRef{3406}]}
\MusicPDFSetAnnotations{363}{[\MusicPDFRef{3407} \MusicPDFRef{3408} \MusicPDFRef{3409} \MusicPDFRef{3410} \MusicPDFRef{3411} \MusicPDFRef{3412} \MusicPDFRef{3413} \MusicPDFRef{3414} \MusicPDFRef{3415}]}
\MusicPDFSetAnnotations{364}{[\MusicPDFRef{3416} \MusicPDFRef{3417} \MusicPDFRef{3418} \MusicPDFRef{3419} \MusicPDFRef{3420}]}
\MusicPDFSetAnnotations{365}{[\MusicPDFRef{3421} \MusicPDFRef{3422} \MusicPDFRef{3423} \MusicPDFRef{3424} \MusicPDFRef{3425} \MusicPDFRef{3426} \MusicPDFRef{3427}]}
\MusicPDFSetAnnotations{366}{[\MusicPDFRef{3428} \MusicPDFRef{3429} \MusicPDFRef{3430} \MusicPDFRef{3431} \MusicPDFRef{3432} \MusicPDFRef{3433} \MusicPDFRef{3434} \MusicPDFRef{3435} \MusicPDFRef{3436} \MusicPDFRef{3437}]}
\MusicPDFSetAnnotations{367}{[\MusicPDFRef{3438} \MusicPDFRef{3439} \MusicPDFRef{3440} \MusicPDFRef{3441} \MusicPDFRef{3442} \MusicPDFRef{3443} \MusicPDFRef{3444} \MusicPDFRef{3445} \MusicPDFRef{3446} \MusicPDFRef{3447} \MusicPDFRef{3448} \MusicPDFRef{3449} \MusicPDFRef{3450} \MusicPDFRef{3451} \MusicPDFRef{3452} \MusicPDFRef{3453} \MusicPDFRef{3454}]}
\MusicPDFSetAnnotations{368}{[\MusicPDFRef{3455} \MusicPDFRef{3456} \MusicPDFRef{3457} \MusicPDFRef{3458} \MusicPDFRef{3459} \MusicPDFRef{3460}]}
\MusicPDFSetAnnotations{369}{[\MusicPDFRef{3461} \MusicPDFRef{3462} \MusicPDFRef{3463}]}
\MusicPDFSetAnnotations{370}{[\MusicPDFRef{3464} \MusicPDFRef{3465} \MusicPDFRef{3466} \MusicPDFRef{3467}]}
\MusicPDFSetAnnotations{371}{[\MusicPDFRef{3468} \MusicPDFRef{3469} \MusicPDFRef{3470} \MusicPDFRef{3471} \MusicPDFRef{3472} \MusicPDFRef{3473} \MusicPDFRef{3474} \MusicPDFRef{3475} \MusicPDFRef{3476} \MusicPDFRef{3477} \MusicPDFRef{3478} \MusicPDFRef{3479} \MusicPDFRef{3480} \MusicPDFRef{3481}]}
\MusicPDFSetAnnotations{372}{[\MusicPDFRef{3482} \MusicPDFRef{3483} \MusicPDFRef{3484} \MusicPDFRef{3485} \MusicPDFRef{3486} \MusicPDFRef{3487} \MusicPDFRef{3488} \MusicPDFRef{3489} \MusicPDFRef{3490} \MusicPDFRef{3491} \MusicPDFRef{3492} \MusicPDFRef{3493} \MusicPDFRef{3494} \MusicPDFRef{3495} \MusicPDFRef{3496} \MusicPDFRef{3497} \MusicPDFRef{3498} \MusicPDFRef{3499} \MusicPDFRef{3500} \MusicPDFRef{3501}]}
\MusicPDFSetAnnotations{373}{[\MusicPDFRef{3502} \MusicPDFRef{3503} \MusicPDFRef{3504} \MusicPDFRef{3505} \MusicPDFRef{3506} \MusicPDFRef{3507} \MusicPDFRef{3508} \MusicPDFRef{3509} \MusicPDFRef{3510} \MusicPDFRef{3511} \MusicPDFRef{3512} \MusicPDFRef{3513} \MusicPDFRef{3514} \MusicPDFRef{3515} \MusicPDFRef{3516} \MusicPDFRef{3517} \MusicPDFRef{3518}]}
\MusicPDFSetAnnotations{374}{[\MusicPDFRef{3519} \MusicPDFRef{3520} \MusicPDFRef{3521} \MusicPDFRef{3522} \MusicPDFRef{3523} \MusicPDFRef{3524} \MusicPDFRef{3525} \MusicPDFRef{3526} \MusicPDFRef{3527} \MusicPDFRef{3528}]}
\MusicPDFSetAnnotations{375}{[\MusicPDFRef{3529} \MusicPDFRef{3530} \MusicPDFRef{3531} \MusicPDFRef{3532} \MusicPDFRef{3533} \MusicPDFRef{3534} \MusicPDFRef{3535} \MusicPDFRef{3536} \MusicPDFRef{3537} \MusicPDFRef{3538} \MusicPDFRef{3539} \MusicPDFRef{3540} \MusicPDFRef{3541} \MusicPDFRef{3542} \MusicPDFRef{3543} \MusicPDFRef{3544} \MusicPDFRef{3545} \MusicPDFRef{3546} \MusicPDFRef{3547} \MusicPDFRef{3548} \MusicPDFRef{3549} \MusicPDFRef{3550} \MusicPDFRef{3551}]}
\MusicPDFSetAnnotations{376}{[\MusicPDFRef{3552} \MusicPDFRef{3553} \MusicPDFRef{3554} \MusicPDFRef{3555} \MusicPDFRef{3556} \MusicPDFRef{3557} \MusicPDFRef{3558} \MusicPDFRef{3559} \MusicPDFRef{3560} \MusicPDFRef{3561} \MusicPDFRef{3562} \MusicPDFRef{3563} \MusicPDFRef{3564} \MusicPDFRef{3565} \MusicPDFRef{3566} \MusicPDFRef{3567} \MusicPDFRef{3568} \MusicPDFRef{3569} \MusicPDFRef{3570} \MusicPDFRef{3571} \MusicPDFRef{3572} \MusicPDFRef{3573} \MusicPDFRef{3574} \MusicPDFRef{3575}]}
\MusicPDFSetAnnotations{377}{[\MusicPDFRef{3576} \MusicPDFRef{3577} \MusicPDFRef{3578} \MusicPDFRef{3579} \MusicPDFRef{3580} \MusicPDFRef{3581}]}
\MusicPDFSetAnnotations{378}{[\MusicPDFRef{3582} \MusicPDFRef{3583} \MusicPDFRef{3584} \MusicPDFRef{3585} \MusicPDFRef{3586} \MusicPDFRef{3587} \MusicPDFRef{3588} \MusicPDFRef{3589} \MusicPDFRef{3590} \MusicPDFRef{3591} \MusicPDFRef{3592} \MusicPDFRef{3593} \MusicPDFRef{3594} \MusicPDFRef{3595} \MusicPDFRef{3596} \MusicPDFRef{3597} \MusicPDFRef{3598} \MusicPDFRef{3599} \MusicPDFRef{3600} \MusicPDFRef{3601} \MusicPDFRef{3602}]}
\MusicPDFSetAnnotations{379}{[\MusicPDFRef{3603} \MusicPDFRef{3604} \MusicPDFRef{3605} \MusicPDFRef{3606} \MusicPDFRef{3607} \MusicPDFRef{3608} \MusicPDFRef{3609}]}
\MusicPDFSetAnnotations{381}{[\MusicPDFRef{3610} \MusicPDFRef{3611} \MusicPDFRef{3612} \MusicPDFRef{3613} \MusicPDFRef{3614} \MusicPDFRef{3615} \MusicPDFRef{3616} \MusicPDFRef{3617} \MusicPDFRef{3618} \MusicPDFRef{3619}]}
\MusicPDFSetAnnotations{382}{[\MusicPDFRef{3620} \MusicPDFRef{3621} \MusicPDFRef{3622} \MusicPDFRef{3623} \MusicPDFRef{3624} \MusicPDFRef{3625}]}
\MusicPDFSetAnnotations{383}{[\MusicPDFRef{3626} \MusicPDFRef{3627} \MusicPDFRef{3628} \MusicPDFRef{3629} \MusicPDFRef{3630} \MusicPDFRef{3631} \MusicPDFRef{3632} \MusicPDFRef{3633}]}
\MusicPDFSetAnnotations{384}{[\MusicPDFRef{3634} \MusicPDFRef{3635} \MusicPDFRef{3636} \MusicPDFRef{3637} \MusicPDFRef{3638} \MusicPDFRef{3639} \MusicPDFRef{3640} \MusicPDFRef{3641} \MusicPDFRef{3642} \MusicPDFRef{3643} \MusicPDFRef{3644} \MusicPDFRef{3645}]}
\MusicPDFSetAnnotations{385}{[\MusicPDFRef{3646} \MusicPDFRef{3647} \MusicPDFRef{3648} \MusicPDFRef{3649} \MusicPDFRef{3650}]}
\MusicPDFSetAnnotations{386}{[\MusicPDFRef{3651} \MusicPDFRef{3652}]}
\MusicPDFSetAnnotations{387}{[\MusicPDFRef{3653} \MusicPDFRef{3654} \MusicPDFRef{3655} \MusicPDFRef{3656} \MusicPDFRef{3657} \MusicPDFRef{3658}]}
\MusicPDFSetAnnotations{388}{[\MusicPDFRef{3659} \MusicPDFRef{3660} \MusicPDFRef{3661} \MusicPDFRef{3662} \MusicPDFRef{3663} \MusicPDFRef{3664}]}
\MusicPDFSetAnnotations{389}{[\MusicPDFRef{3665} \MusicPDFRef{3666} \MusicPDFRef{3667} \MusicPDFRef{3668} \MusicPDFRef{3669} \MusicPDFRef{3670} \MusicPDFRef{3671} \MusicPDFRef{3672} \MusicPDFRef{3673} \MusicPDFRef{3674} \MusicPDFRef{3675}]}
\MusicPDFSetAnnotations{390}{[\MusicPDFRef{3676} \MusicPDFRef{3677} \MusicPDFRef{3678} \MusicPDFRef{3679} \MusicPDFRef{3680} \MusicPDFRef{3681} \MusicPDFRef{3682} \MusicPDFRef{3683} \MusicPDFRef{3684} \MusicPDFRef{3685} \MusicPDFRef{3686} \MusicPDFRef{3687} \MusicPDFRef{3688} \MusicPDFRef{3689} \MusicPDFRef{3690}]}
\MusicPDFSetAnnotations{391}{[\MusicPDFRef{3691} \MusicPDFRef{3692} \MusicPDFRef{3693} \MusicPDFRef{3694} \MusicPDFRef{3695} \MusicPDFRef{3696} \MusicPDFRef{3697} \MusicPDFRef{3698} \MusicPDFRef{3699} \MusicPDFRef{3700} \MusicPDFRef{3701} \MusicPDFRef{3702} \MusicPDFRef{3703}]}
\MusicPDFSetAnnotations{392}{[\MusicPDFRef{3704} \MusicPDFRef{3705} \MusicPDFRef{3706} \MusicPDFRef{3707} \MusicPDFRef{3708} \MusicPDFRef{3709} \MusicPDFRef{3710}]}
\MusicPDFSetAnnotations{393}{[\MusicPDFRef{3711} \MusicPDFRef{3712} \MusicPDFRef{3713} \MusicPDFRef{3714} \MusicPDFRef{3715} \MusicPDFRef{3716}]}
\MusicPDFSetAnnotations{394}{[\MusicPDFRef{3717} \MusicPDFRef{3718} \MusicPDFRef{3719}]}
\MusicPDFSetAnnotations{395}{[\MusicPDFRef{3720} \MusicPDFRef{3721} \MusicPDFRef{3722} \MusicPDFRef{3723} \MusicPDFRef{3724} \MusicPDFRef{3725} \MusicPDFRef{3726} \MusicPDFRef{3727} \MusicPDFRef{3728} \MusicPDFRef{3729} \MusicPDFRef{3730} \MusicPDFRef{3731} \MusicPDFRef{3732} \MusicPDFRef{3733} \MusicPDFRef{3734} \MusicPDFRef{3735} \MusicPDFRef{3736}]}
\MusicPDFSetAnnotations{396}{[\MusicPDFRef{3737} \MusicPDFRef{3738} \MusicPDFRef{3739} \MusicPDFRef{3740} \MusicPDFRef{3741} \MusicPDFRef{3742} \MusicPDFRef{3743} \MusicPDFRef{3744} \MusicPDFRef{3745} \MusicPDFRef{3746} \MusicPDFRef{3747}]}
\MusicPDFSetAnnotations{397}{[\MusicPDFRef{3748} \MusicPDFRef{3749} \MusicPDFRef{3750} \MusicPDFRef{3751} \MusicPDFRef{3752} \MusicPDFRef{3753} \MusicPDFRef{3754} \MusicPDFRef{3755} \MusicPDFRef{3756} \MusicPDFRef{3757} \MusicPDFRef{3758} \MusicPDFRef{3759} \MusicPDFRef{3760}]}
\MusicPDFSetAnnotations{398}{[\MusicPDFRef{3761} \MusicPDFRef{3762} \MusicPDFRef{3763} \MusicPDFRef{3764} \MusicPDFRef{3765} \MusicPDFRef{3766} \MusicPDFRef{3767} \MusicPDFRef{3768} \MusicPDFRef{3769}]}
\MusicPDFSetAnnotations{399}{[\MusicPDFRef{3770} \MusicPDFRef{3771} \MusicPDFRef{3772} \MusicPDFRef{3773} \MusicPDFRef{3774} \MusicPDFRef{3775} \MusicPDFRef{3776} \MusicPDFRef{3777} \MusicPDFRef{3778} \MusicPDFRef{3779} \MusicPDFRef{3780}]}
\MusicPDFSetAnnotations{400}{[\MusicPDFRef{3781} \MusicPDFRef{3782} \MusicPDFRef{3783} \MusicPDFRef{3784} \MusicPDFRef{3785} \MusicPDFRef{3786}]}
\MusicPDFSetAnnotations{401}{[\MusicPDFRef{3787} \MusicPDFRef{3788} \MusicPDFRef{3789} \MusicPDFRef{3790} \MusicPDFRef{3791} \MusicPDFRef{3792} \MusicPDFRef{3793} \MusicPDFRef{3794} \MusicPDFRef{3795}]}
\MusicPDFSetAnnotations{402}{[\MusicPDFRef{3796} \MusicPDFRef{3797} \MusicPDFRef{3798} \MusicPDFRef{3799} \MusicPDFRef{3800} \MusicPDFRef{3801} \MusicPDFRef{3802} \MusicPDFRef{3803}]}
\MusicPDFSetAnnotations{403}{[\MusicPDFRef{3804} \MusicPDFRef{3805} \MusicPDFRef{3806} \MusicPDFRef{3807} \MusicPDFRef{3808} \MusicPDFRef{3809} \MusicPDFRef{3810} \MusicPDFRef{3811} \MusicPDFRef{3812} \MusicPDFRef{3813} \MusicPDFRef{3814} \MusicPDFRef{3815} \MusicPDFRef{3816} \MusicPDFRef{3817} \MusicPDFRef{3818} \MusicPDFRef{3819}]}
\MusicPDFSetAnnotations{404}{[\MusicPDFRef{3820} \MusicPDFRef{3821} \MusicPDFRef{3822} \MusicPDFRef{3823}]}
\MusicPDFSetAnnotations{405}{[\MusicPDFRef{3824} \MusicPDFRef{3825} \MusicPDFRef{3826} \MusicPDFRef{3827} \MusicPDFRef{3828} \MusicPDFRef{3829} \MusicPDFRef{3830} \MusicPDFRef{3831} \MusicPDFRef{3832} \MusicPDFRef{3833} \MusicPDFRef{3834}]}
\MusicPDFSetAnnotations{406}{[\MusicPDFRef{3835} \MusicPDFRef{3836} \MusicPDFRef{3837} \MusicPDFRef{3838} \MusicPDFRef{3839} \MusicPDFRef{3840} \MusicPDFRef{3841} \MusicPDFRef{3842}]}
\MusicPDFSetAnnotations{407}{[\MusicPDFRef{3843} \MusicPDFRef{3844} \MusicPDFRef{3845} \MusicPDFRef{3846} \MusicPDFRef{3847} \MusicPDFRef{3848} \MusicPDFRef{3849} \MusicPDFRef{3850}]}
\MusicPDFSetAnnotations{408}{[\MusicPDFRef{3851} \MusicPDFRef{3852} \MusicPDFRef{3853} \MusicPDFRef{3854} \MusicPDFRef{3855} \MusicPDFRef{3856} \MusicPDFRef{3857}]}
\MusicPDFSetAnnotations{409}{[\MusicPDFRef{3858} \MusicPDFRef{3859} \MusicPDFRef{3860} \MusicPDFRef{3861}]}
\MusicPDFSetAnnotations{410}{[\MusicPDFRef{3862} \MusicPDFRef{3863} \MusicPDFRef{3864} \MusicPDFRef{3865} \MusicPDFRef{3866} \MusicPDFRef{3867} \MusicPDFRef{3868} \MusicPDFRef{3869}]}
\MusicPDFSetAnnotations{411}{[\MusicPDFRef{3870} \MusicPDFRef{3871} \MusicPDFRef{3872} \MusicPDFRef{3873} \MusicPDFRef{3874} \MusicPDFRef{3875} \MusicPDFRef{3876} \MusicPDFRef{3877} \MusicPDFRef{3878} \MusicPDFRef{3879} \MusicPDFRef{3880}]}
\MusicPDFSetAnnotations{412}{[\MusicPDFRef{3881} \MusicPDFRef{3882}]}
\MusicPDFSetAnnotations{413}{[\MusicPDFRef{3883} \MusicPDFRef{3884} \MusicPDFRef{3885} \MusicPDFRef{3886} \MusicPDFRef{3887} \MusicPDFRef{3888} \MusicPDFRef{3889}]}
\MusicPDFSetAnnotations{414}{[\MusicPDFRef{3890} \MusicPDFRef{3891} \MusicPDFRef{3892} \MusicPDFRef{3893} \MusicPDFRef{3894} \MusicPDFRef{3895} \MusicPDFRef{3896} \MusicPDFRef{3897} \MusicPDFRef{3898} \MusicPDFRef{3899} \MusicPDFRef{3900} \MusicPDFRef{3901} \MusicPDFRef{3902} \MusicPDFRef{3903}]}
\MusicPDFSetAnnotations{415}{[\MusicPDFRef{3904} \MusicPDFRef{3905} \MusicPDFRef{3906} \MusicPDFRef{3907} \MusicPDFRef{3908} \MusicPDFRef{3909}]}
\MusicPDFSetAnnotations{416}{[\MusicPDFRef{3910} \MusicPDFRef{3911} \MusicPDFRef{3912}]}
\MusicPDFSetAnnotations{417}{[\MusicPDFRef{3913} \MusicPDFRef{3914} \MusicPDFRef{3915} \MusicPDFRef{3916}]}
\MusicPDFSetAnnotations{418}{[\MusicPDFRef{3917} \MusicPDFRef{3918} \MusicPDFRef{3919} \MusicPDFRef{3920}]}
\MusicPDFSetAnnotations{419}{[\MusicPDFRef{3921} \MusicPDFRef{3922} \MusicPDFRef{3923} \MusicPDFRef{3924} \MusicPDFRef{3925} \MusicPDFRef{3926} \MusicPDFRef{3927} \MusicPDFRef{3928} \MusicPDFRef{3929}]}
\MusicPDFSetAnnotations{420}{[\MusicPDFRef{3930} \MusicPDFRef{3931} \MusicPDFRef{3932} \MusicPDFRef{3933}]}
\MusicPDFSetAnnotations{421}{[\MusicPDFRef{3934} \MusicPDFRef{3935} \MusicPDFRef{3936} \MusicPDFRef{3937} \MusicPDFRef{3938} \MusicPDFRef{3939} \MusicPDFRef{3940}]}
\MusicPDFSetAnnotations{422}{[\MusicPDFRef{3941} \MusicPDFRef{3942} \MusicPDFRef{3943} \MusicPDFRef{3944} \MusicPDFRef{3945} \MusicPDFRef{3946} \MusicPDFRef{3947}]}
\MusicPDFSetAnnotations{423}{[\MusicPDFRef{3948} \MusicPDFRef{3949} \MusicPDFRef{3950} \MusicPDFRef{3951} \MusicPDFRef{3952} \MusicPDFRef{3953} \MusicPDFRef{3954} \MusicPDFRef{3955} \MusicPDFRef{3956} \MusicPDFRef{3957} \MusicPDFRef{3958}]}
\MusicPDFSetAnnotations{424}{[\MusicPDFRef{3959} \MusicPDFRef{3960} \MusicPDFRef{3961} \MusicPDFRef{3962} \MusicPDFRef{3963}]}
\MusicPDFSetAnnotations{425}{[\MusicPDFRef{3964} \MusicPDFRef{3965} \MusicPDFRef{3966} \MusicPDFRef{3967}]}
\MusicPDFSetAnnotations{426}{[\MusicPDFRef{3968} \MusicPDFRef{3969} \MusicPDFRef{3970} \MusicPDFRef{3971} \MusicPDFRef{3972} \MusicPDFRef{3973} \MusicPDFRef{3974} \MusicPDFRef{3975} \MusicPDFRef{3976} \MusicPDFRef{3977} \MusicPDFRef{3978}]}
\MusicPDFSetAnnotations{427}{[\MusicPDFRef{3979}]}
\MusicPDFSetAnnotations{428}{[\MusicPDFRef{3980} \MusicPDFRef{3981} \MusicPDFRef{3982} \MusicPDFRef{3983} \MusicPDFRef{3984} \MusicPDFRef{3985} \MusicPDFRef{3986} \MusicPDFRef{3987} \MusicPDFRef{3988} \MusicPDFRef{3989} \MusicPDFRef{3990} \MusicPDFRef{3991} \MusicPDFRef{3992} \MusicPDFRef{3993} \MusicPDFRef{3994}]}
\MusicPDFSetAnnotations{429}{[\MusicPDFRef{3995} \MusicPDFRef{3996} \MusicPDFRef{3997} \MusicPDFRef{3998} \MusicPDFRef{3999} \MusicPDFRef{4000} \MusicPDFRef{4001} \MusicPDFRef{4002} \MusicPDFRef{4003} \MusicPDFRef{4004} \MusicPDFRef{4005} \MusicPDFRef{4006}]}
\MusicPDFSetAnnotations{430}{[\MusicPDFRef{4007} \MusicPDFRef{4008} \MusicPDFRef{4009} \MusicPDFRef{4010} \MusicPDFRef{4011} \MusicPDFRef{4012} \MusicPDFRef{4013} \MusicPDFRef{4014} \MusicPDFRef{4015} \MusicPDFRef{4016} \MusicPDFRef{4017} \MusicPDFRef{4018} \MusicPDFRef{4019}]}
\MusicPDFSetAnnotations{431}{[\MusicPDFRef{4020} \MusicPDFRef{4021} \MusicPDFRef{4022} \MusicPDFRef{4023} \MusicPDFRef{4024} \MusicPDFRef{4025} \MusicPDFRef{4026} \MusicPDFRef{4027} \MusicPDFRef{4028} \MusicPDFRef{4029}]}
\MusicPDFSetAnnotations{433}{[\MusicPDFRef{4030} \MusicPDFRef{4031} \MusicPDFRef{4032} \MusicPDFRef{4033} \MusicPDFRef{4034} \MusicPDFRef{4035} \MusicPDFRef{4036} \MusicPDFRef{4037} \MusicPDFRef{4038} \MusicPDFRef{4039} \MusicPDFRef{4040} \MusicPDFRef{4041} \MusicPDFRef{4042} \MusicPDFRef{4043} \MusicPDFRef{4044} \MusicPDFRef{4045} \MusicPDFRef{4046} \MusicPDFRef{4047} \MusicPDFRef{4048} \MusicPDFRef{4049} \MusicPDFRef{4050}]}
\MusicPDFSetAnnotations{434}{[\MusicPDFRef{4051} \MusicPDFRef{4052} \MusicPDFRef{4053} \MusicPDFRef{4054} \MusicPDFRef{4055} \MusicPDFRef{4056} \MusicPDFRef{4057} \MusicPDFRef{4058} \MusicPDFRef{4059} \MusicPDFRef{4060} \MusicPDFRef{4061}]}
\MusicPDFSetAnnotations{435}{[\MusicPDFRef{4062} \MusicPDFRef{4063} \MusicPDFRef{4064} \MusicPDFRef{4065} \MusicPDFRef{4066} \MusicPDFRef{4067} \MusicPDFRef{4068} \MusicPDFRef{4069} \MusicPDFRef{4070}]}
\MusicPDFSetAnnotations{436}{[\MusicPDFRef{4071} \MusicPDFRef{4072} \MusicPDFRef{4073} \MusicPDFRef{4074}]}
\MusicPDFSetAnnotations{437}{[\MusicPDFRef{4075} \MusicPDFRef{4076} \MusicPDFRef{4077} \MusicPDFRef{4078} \MusicPDFRef{4079} \MusicPDFRef{4080}]}
\MusicPDFSetAnnotations{438}{[\MusicPDFRef{4081} \MusicPDFRef{4082} \MusicPDFRef{4083} \MusicPDFRef{4084} \MusicPDFRef{4085} \MusicPDFRef{4086} \MusicPDFRef{4087} \MusicPDFRef{4088} \MusicPDFRef{4089} \MusicPDFRef{4090}]}
\MusicPDFSetAnnotations{439}{[\MusicPDFRef{4091} \MusicPDFRef{4092} \MusicPDFRef{4093} \MusicPDFRef{4094} \MusicPDFRef{4095} \MusicPDFRef{4096} \MusicPDFRef{4097} \MusicPDFRef{4098} \MusicPDFRef{4099} \MusicPDFRef{4100} \MusicPDFRef{4101} \MusicPDFRef{4102}]}
\MusicPDFSetAnnotations{440}{[\MusicPDFRef{4103} \MusicPDFRef{4104} \MusicPDFRef{4105} \MusicPDFRef{4106} \MusicPDFRef{4107} \MusicPDFRef{4108} \MusicPDFRef{4109}]}
\MusicPDFSetAnnotations{441}{[\MusicPDFRef{4110} \MusicPDFRef{4111} \MusicPDFRef{4112} \MusicPDFRef{4113} \MusicPDFRef{4114} \MusicPDFRef{4115} \MusicPDFRef{4116} \MusicPDFRef{4117} \MusicPDFRef{4118} \MusicPDFRef{4119} \MusicPDFRef{4120}]}
\MusicPDFSetAnnotations{442}{[\MusicPDFRef{4121} \MusicPDFRef{4122} \MusicPDFRef{4123} \MusicPDFRef{4124} \MusicPDFRef{4125} \MusicPDFRef{4126} \MusicPDFRef{4127} \MusicPDFRef{4128} \MusicPDFRef{4129} \MusicPDFRef{4130} \MusicPDFRef{4131}]}
\MusicPDFSetAnnotations{443}{[\MusicPDFRef{4132} \MusicPDFRef{4133} \MusicPDFRef{4134} \MusicPDFRef{4135} \MusicPDFRef{4136} \MusicPDFRef{4137} \MusicPDFRef{4138} \MusicPDFRef{4139} \MusicPDFRef{4140} \MusicPDFRef{4141} \MusicPDFRef{4142}]}
\MusicPDFSetAnnotations{444}{[\MusicPDFRef{4143} \MusicPDFRef{4144} \MusicPDFRef{4145} \MusicPDFRef{4146} \MusicPDFRef{4147} \MusicPDFRef{4148} \MusicPDFRef{4149} \MusicPDFRef{4150} \MusicPDFRef{4151} \MusicPDFRef{4152} \MusicPDFRef{4153} \MusicPDFRef{4154} \MusicPDFRef{4155} \MusicPDFRef{4156} \MusicPDFRef{4157} \MusicPDFRef{4158} \MusicPDFRef{4159} \MusicPDFRef{4160} \MusicPDFRef{4161} \MusicPDFRef{4162}]}
\MusicPDFSetAnnotations{445}{[\MusicPDFRef{4163} \MusicPDFRef{4164} \MusicPDFRef{4165} \MusicPDFRef{4166} \MusicPDFRef{4167} \MusicPDFRef{4168}]}
\MusicPDFSetAnnotations{446}{[\MusicPDFRef{4169} \MusicPDFRef{4170} \MusicPDFRef{4171} \MusicPDFRef{4172} \MusicPDFRef{4173} \MusicPDFRef{4174} \MusicPDFRef{4175} \MusicPDFRef{4176} \MusicPDFRef{4177} \MusicPDFRef{4178} \MusicPDFRef{4179} \MusicPDFRef{4180} \MusicPDFRef{4181} \MusicPDFRef{4182}]}
\MusicPDFSetAnnotations{447}{[\MusicPDFRef{4183} \MusicPDFRef{4184} \MusicPDFRef{4185} \MusicPDFRef{4186} \MusicPDFRef{4187} \MusicPDFRef{4188} \MusicPDFRef{4189} \MusicPDFRef{4190} \MusicPDFRef{4191} \MusicPDFRef{4192} \MusicPDFRef{4193} \MusicPDFRef{4194} \MusicPDFRef{4195} \MusicPDFRef{4196} \MusicPDFRef{4197} \MusicPDFRef{4198} \MusicPDFRef{4199}]}
\MusicPDFSetAnnotations{448}{[\MusicPDFRef{4200} \MusicPDFRef{4201} \MusicPDFRef{4202} \MusicPDFRef{4203} \MusicPDFRef{4204} \MusicPDFRef{4205} \MusicPDFRef{4206} \MusicPDFRef{4207}]}
\MusicPDFSetAnnotations{449}{[\MusicPDFRef{4208} \MusicPDFRef{4209} \MusicPDFRef{4210} \MusicPDFRef{4211} \MusicPDFRef{4212} \MusicPDFRef{4213}]}
\MusicPDFSetAnnotations{450}{[\MusicPDFRef{4214} \MusicPDFRef{4215} \MusicPDFRef{4216} \MusicPDFRef{4217} \MusicPDFRef{4218} \MusicPDFRef{4219}]}
\MusicPDFSetAnnotations{451}{[\MusicPDFRef{4220} \MusicPDFRef{4221} \MusicPDFRef{4222} \MusicPDFRef{4223} \MusicPDFRef{4224} \MusicPDFRef{4225} \MusicPDFRef{4226} \MusicPDFRef{4227} \MusicPDFRef{4228} \MusicPDFRef{4229} \MusicPDFRef{4230} \MusicPDFRef{4231} \MusicPDFRef{4232}]}
\MusicPDFSetAnnotations{452}{[\MusicPDFRef{4233} \MusicPDFRef{4234} \MusicPDFRef{4235} \MusicPDFRef{4236} \MusicPDFRef{4237} \MusicPDFRef{4238} \MusicPDFRef{4239} \MusicPDFRef{4240} \MusicPDFRef{4241} \MusicPDFRef{4242} \MusicPDFRef{4243} \MusicPDFRef{4244} \MusicPDFRef{4245} \MusicPDFRef{4246} \MusicPDFRef{4247} \MusicPDFRef{4248} \MusicPDFRef{4249} \MusicPDFRef{4250} \MusicPDFRef{4251} \MusicPDFRef{4252}]}
\MusicPDFSetAnnotations{453}{[\MusicPDFRef{4253} \MusicPDFRef{4254} \MusicPDFRef{4255} \MusicPDFRef{4256} \MusicPDFRef{4257} \MusicPDFRef{4258} \MusicPDFRef{4259} \MusicPDFRef{4260} \MusicPDFRef{4261} \MusicPDFRef{4262} \MusicPDFRef{4263}]}
\MusicPDFSetAnnotations{454}{[\MusicPDFRef{4264} \MusicPDFRef{4265} \MusicPDFRef{4266} \MusicPDFRef{4267} \MusicPDFRef{4268} \MusicPDFRef{4269} \MusicPDFRef{4270} \MusicPDFRef{4271} \MusicPDFRef{4272} \MusicPDFRef{4273}]}
\MusicPDFSetAnnotations{455}{[\MusicPDFRef{4274} \MusicPDFRef{4275} \MusicPDFRef{4276} \MusicPDFRef{4277} \MusicPDFRef{4278} \MusicPDFRef{4279} \MusicPDFRef{4280} \MusicPDFRef{4281} \MusicPDFRef{4282}]}
\MusicPDFSetAnnotations{456}{[\MusicPDFRef{4283} \MusicPDFRef{4284} \MusicPDFRef{4285} \MusicPDFRef{4286} \MusicPDFRef{4287}]}
\MusicPDFSetAnnotations{457}{[\MusicPDFRef{4288} \MusicPDFRef{4289} \MusicPDFRef{4290} \MusicPDFRef{4291} \MusicPDFRef{4292} \MusicPDFRef{4293} \MusicPDFRef{4294} \MusicPDFRef{4295} \MusicPDFRef{4296} \MusicPDFRef{4297}]}
\MusicPDFSetAnnotations{458}{[\MusicPDFRef{4298} \MusicPDFRef{4299} \MusicPDFRef{4300} \MusicPDFRef{4301} \MusicPDFRef{4302} \MusicPDFRef{4303} \MusicPDFRef{4304}]}
\MusicPDFSetAnnotations{459}{[\MusicPDFRef{4305} \MusicPDFRef{4306} \MusicPDFRef{4307} \MusicPDFRef{4308} \MusicPDFRef{4309} \MusicPDFRef{4310} \MusicPDFRef{4311} \MusicPDFRef{4312} \MusicPDFRef{4313} \MusicPDFRef{4314} \MusicPDFRef{4315} \MusicPDFRef{4316} \MusicPDFRef{4317} \MusicPDFRef{4318} \MusicPDFRef{4319} \MusicPDFRef{4320} \MusicPDFRef{4321}]}
\MusicPDFSetAnnotations{460}{[\MusicPDFRef{4322} \MusicPDFRef{4323} \MusicPDFRef{4324} \MusicPDFRef{4325} \MusicPDFRef{4326} \MusicPDFRef{4327} \MusicPDFRef{4328} \MusicPDFRef{4329} \MusicPDFRef{4330} \MusicPDFRef{4331} \MusicPDFRef{4332} \MusicPDFRef{4333} \MusicPDFRef{4334} \MusicPDFRef{4335} \MusicPDFRef{4336} \MusicPDFRef{4337} \MusicPDFRef{4338} \MusicPDFRef{4339} \MusicPDFRef{4340} \MusicPDFRef{4341} \MusicPDFRef{4342}]}
\MusicPDFSetAnnotations{461}{[\MusicPDFRef{4343} \MusicPDFRef{4344} \MusicPDFRef{4345} \MusicPDFRef{4346} \MusicPDFRef{4347} \MusicPDFRef{4348} \MusicPDFRef{4349} \MusicPDFRef{4350}]}
\MusicPDFSetAnnotations{462}{[\MusicPDFRef{4351} \MusicPDFRef{4352} \MusicPDFRef{4353} \MusicPDFRef{4354} \MusicPDFRef{4355} \MusicPDFRef{4356} \MusicPDFRef{4357} \MusicPDFRef{4358} \MusicPDFRef{4359} \MusicPDFRef{4360} \MusicPDFRef{4361} \MusicPDFRef{4362} \MusicPDFRef{4363}]}
\MusicPDFSetAnnotations{463}{[\MusicPDFRef{4364} \MusicPDFRef{4365} \MusicPDFRef{4366} \MusicPDFRef{4367} \MusicPDFRef{4368} \MusicPDFRef{4369} \MusicPDFRef{4370} \MusicPDFRef{4371} \MusicPDFRef{4372}]}
\MusicPDFSetAnnotations{464}{[\MusicPDFRef{4373} \MusicPDFRef{4374} \MusicPDFRef{4375} \MusicPDFRef{4376} \MusicPDFRef{4377} \MusicPDFRef{4378} \MusicPDFRef{4379} \MusicPDFRef{4380} \MusicPDFRef{4381} \MusicPDFRef{4382} \MusicPDFRef{4383} \MusicPDFRef{4384} \MusicPDFRef{4385}]}
\MusicPDFSetAnnotations{465}{[\MusicPDFRef{4386} \MusicPDFRef{4387} \MusicPDFRef{4388} \MusicPDFRef{4389} \MusicPDFRef{4390} \MusicPDFRef{4391} \MusicPDFRef{4392} \MusicPDFRef{4393} \MusicPDFRef{4394} \MusicPDFRef{4395} \MusicPDFRef{4396} \MusicPDFRef{4397}]}
\MusicPDFSetAnnotations{466}{[\MusicPDFRef{4398} \MusicPDFRef{4399} \MusicPDFRef{4400} \MusicPDFRef{4401} \MusicPDFRef{4402} \MusicPDFRef{4403} \MusicPDFRef{4404} \MusicPDFRef{4405} \MusicPDFRef{4406} \MusicPDFRef{4407}]}
\MusicPDFSetAnnotations{467}{[\MusicPDFRef{4408} \MusicPDFRef{4409} \MusicPDFRef{4410} \MusicPDFRef{4411} \MusicPDFRef{4412} \MusicPDFRef{4413} \MusicPDFRef{4414}]}
\MusicPDFSetAnnotations{468}{[\MusicPDFRef{4415} \MusicPDFRef{4416} \MusicPDFRef{4417} \MusicPDFRef{4418} \MusicPDFRef{4419} \MusicPDFRef{4420} \MusicPDFRef{4421} \MusicPDFRef{4422} \MusicPDFRef{4423} \MusicPDFRef{4424} \MusicPDFRef{4425} \MusicPDFRef{4426} \MusicPDFRef{4427} \MusicPDFRef{4428} \MusicPDFRef{4429}]}
\MusicPDFSetAnnotations{469}{[\MusicPDFRef{4430} \MusicPDFRef{4431} \MusicPDFRef{4432} \MusicPDFRef{4433} \MusicPDFRef{4434} \MusicPDFRef{4435} \MusicPDFRef{4436} \MusicPDFRef{4437} \MusicPDFRef{4438} \MusicPDFRef{4439} \MusicPDFRef{4440} \MusicPDFRef{4441} \MusicPDFRef{4442} \MusicPDFRef{4443} \MusicPDFRef{4444} \MusicPDFRef{4445} \MusicPDFRef{4446} \MusicPDFRef{4447} \MusicPDFRef{4448} \MusicPDFRef{4449}]}
\MusicPDFSetAnnotations{470}{[\MusicPDFRef{4450} \MusicPDFRef{4451} \MusicPDFRef{4452} \MusicPDFRef{4453} \MusicPDFRef{4454} \MusicPDFRef{4455} \MusicPDFRef{4456} \MusicPDFRef{4457}]}
\MusicPDFSetAnnotations{471}{[\MusicPDFRef{4458} \MusicPDFRef{4459} \MusicPDFRef{4460} \MusicPDFRef{4461} \MusicPDFRef{4462} \MusicPDFRef{4463}]}
\MusicPDFSetAnnotations{472}{[\MusicPDFRef{4464} \MusicPDFRef{4465} \MusicPDFRef{4466} \MusicPDFRef{4467} \MusicPDFRef{4468} \MusicPDFRef{4469} \MusicPDFRef{4470} \MusicPDFRef{4471}]}
\MusicPDFSetAnnotations{473}{[\MusicPDFRef{4472} \MusicPDFRef{4473} \MusicPDFRef{4474} \MusicPDFRef{4475} \MusicPDFRef{4476}]}
\MusicPDFSetAnnotations{475}{[\MusicPDFRef{4477} \MusicPDFRef{4478} \MusicPDFRef{4479} \MusicPDFRef{4480} \MusicPDFRef{4481} \MusicPDFRef{4482} \MusicPDFRef{4483} \MusicPDFRef{4484} \MusicPDFRef{4485} \MusicPDFRef{4486} \MusicPDFRef{4487} \MusicPDFRef{4488} \MusicPDFRef{4489}]}
\MusicPDFSetAnnotations{476}{[\MusicPDFRef{4490} \MusicPDFRef{4491} \MusicPDFRef{4492} \MusicPDFRef{4493} \MusicPDFRef{4494} \MusicPDFRef{4495} \MusicPDFRef{4496}]}
\MusicPDFSetAnnotations{477}{[\MusicPDFRef{4497} \MusicPDFRef{4498} \MusicPDFRef{4499} \MusicPDFRef{4500} \MusicPDFRef{4501} \MusicPDFRef{4502} \MusicPDFRef{4503} \MusicPDFRef{4504} \MusicPDFRef{4505} \MusicPDFRef{4506} \MusicPDFRef{4507}]}
\MusicPDFSetAnnotations{478}{[\MusicPDFRef{4508} \MusicPDFRef{4509} \MusicPDFRef{4510} \MusicPDFRef{4511} \MusicPDFRef{4512} \MusicPDFRef{4513} \MusicPDFRef{4514} \MusicPDFRef{4515} \MusicPDFRef{4516} \MusicPDFRef{4517} \MusicPDFRef{4518} \MusicPDFRef{4519}]}
\MusicPDFSetAnnotations{479}{[\MusicPDFRef{4520} \MusicPDFRef{4521} \MusicPDFRef{4522} \MusicPDFRef{4523} \MusicPDFRef{4524} \MusicPDFRef{4525} \MusicPDFRef{4526} \MusicPDFRef{4527} \MusicPDFRef{4528} \MusicPDFRef{4529} \MusicPDFRef{4530} \MusicPDFRef{4531}]}
\MusicPDFSetAnnotations{480}{[\MusicPDFRef{4532} \MusicPDFRef{4533} \MusicPDFRef{4534} \MusicPDFRef{4535}]}
\MusicPDFSetAnnotations{481}{[\MusicPDFRef{4536} \MusicPDFRef{4537} \MusicPDFRef{4538} \MusicPDFRef{4539} \MusicPDFRef{4540} \MusicPDFRef{4541} \MusicPDFRef{4542} \MusicPDFRef{4543}]}
\MusicPDFSetAnnotations{482}{[\MusicPDFRef{4544} \MusicPDFRef{4545} \MusicPDFRef{4546} \MusicPDFRef{4547} \MusicPDFRef{4548} \MusicPDFRef{4549} \MusicPDFRef{4550} \MusicPDFRef{4551} \MusicPDFRef{4552} \MusicPDFRef{4553} \MusicPDFRef{4554} \MusicPDFRef{4555} \MusicPDFRef{4556} \MusicPDFRef{4557}]}
\MusicPDFSetAnnotations{483}{[\MusicPDFRef{4558} \MusicPDFRef{4559} \MusicPDFRef{4560} \MusicPDFRef{4561} \MusicPDFRef{4562} \MusicPDFRef{4563}]}
\MusicPDFSetAnnotations{484}{[\MusicPDFRef{4564} \MusicPDFRef{4565} \MusicPDFRef{4566} \MusicPDFRef{4567} \MusicPDFRef{4568} \MusicPDFRef{4569} \MusicPDFRef{4570} \MusicPDFRef{4571} \MusicPDFRef{4572} \MusicPDFRef{4573} \MusicPDFRef{4574} \MusicPDFRef{4575} \MusicPDFRef{4576} \MusicPDFRef{4577}]}
\MusicPDFSetAnnotations{485}{[\MusicPDFRef{4578} \MusicPDFRef{4579} \MusicPDFRef{4580} \MusicPDFRef{4581} \MusicPDFRef{4582} \MusicPDFRef{4583} \MusicPDFRef{4584} \MusicPDFRef{4585} \MusicPDFRef{4586} \MusicPDFRef{4587} \MusicPDFRef{4588}]}
\MusicPDFSetAnnotations{486}{[\MusicPDFRef{4589} \MusicPDFRef{4590} \MusicPDFRef{4591} \MusicPDFRef{4592} \MusicPDFRef{4593} \MusicPDFRef{4594} \MusicPDFRef{4595} \MusicPDFRef{4596} \MusicPDFRef{4597} \MusicPDFRef{4598} \MusicPDFRef{4599} \MusicPDFRef{4600} \MusicPDFRef{4601} \MusicPDFRef{4602} \MusicPDFRef{4603} \MusicPDFRef{4604} \MusicPDFRef{4605} \MusicPDFRef{4606} \MusicPDFRef{4607}]}
\MusicPDFSetAnnotations{487}{[\MusicPDFRef{4608} \MusicPDFRef{4609} \MusicPDFRef{4610} \MusicPDFRef{4611} \MusicPDFRef{4612} \MusicPDFRef{4613} \MusicPDFRef{4614} \MusicPDFRef{4615}]}
\MusicPDFSetAnnotations{488}{[\MusicPDFRef{4616} \MusicPDFRef{4617} \MusicPDFRef{4618} \MusicPDFRef{4619} \MusicPDFRef{4620}]}
\MusicPDFSetAnnotations{489}{[\MusicPDFRef{4621} \MusicPDFRef{4622} \MusicPDFRef{4623} \MusicPDFRef{4624} \MusicPDFRef{4625} \MusicPDFRef{4626} \MusicPDFRef{4627} \MusicPDFRef{4628} \MusicPDFRef{4629} \MusicPDFRef{4630} \MusicPDFRef{4631}]}
\MusicPDFSetAnnotations{490}{[\MusicPDFRef{4632} \MusicPDFRef{4633} \MusicPDFRef{4634} \MusicPDFRef{4635} \MusicPDFRef{4636} \MusicPDFRef{4637} \MusicPDFRef{4638} \MusicPDFRef{4639}]}
\MusicPDFSetAnnotations{491}{[\MusicPDFRef{4640} \MusicPDFRef{4641} \MusicPDFRef{4642} \MusicPDFRef{4643} \MusicPDFRef{4644} \MusicPDFRef{4645}]}
\MusicPDFSetAnnotations{493}{[\MusicPDFRef{4646}]}
\MusicPDFSetAnnotations{494}{[\MusicPDFRef{4647} \MusicPDFRef{4648} \MusicPDFRef{4649}]}
\MusicPDFSetAnnotations{495}{[\MusicPDFRef{4650} \MusicPDFRef{4651} \MusicPDFRef{4652}]}
\MusicPDFSetAnnotations{496}{[\MusicPDFRef{4653} \MusicPDFRef{4654} \MusicPDFRef{4655} \MusicPDFRef{4656} \MusicPDFRef{4657} \MusicPDFRef{4658} \MusicPDFRef{4659}]}
\MusicPDFSetAnnotations{497}{[\MusicPDFRef{4660} \MusicPDFRef{4661} \MusicPDFRef{4662} \MusicPDFRef{4663} \MusicPDFRef{4664} \MusicPDFRef{4665} \MusicPDFRef{4666} \MusicPDFRef{4667} \MusicPDFRef{4668} \MusicPDFRef{4669}]}
\MusicPDFSetAnnotations{498}{[\MusicPDFRef{4670} \MusicPDFRef{4671} \MusicPDFRef{4672} \MusicPDFRef{4673} \MusicPDFRef{4674} \MusicPDFRef{4675} \MusicPDFRef{4676} \MusicPDFRef{4677} \MusicPDFRef{4678} \MusicPDFRef{4679} \MusicPDFRef{4680} \MusicPDFRef{4681}]}
\MusicPDFSetAnnotations{499}{[\MusicPDFRef{4682} \MusicPDFRef{4683}]}
\MusicPDFSetAnnotations{500}{[\MusicPDFRef{4684} \MusicPDFRef{4685} \MusicPDFRef{4686} \MusicPDFRef{4687} \MusicPDFRef{4688} \MusicPDFRef{4689} \MusicPDFRef{4690} \MusicPDFRef{4691} \MusicPDFRef{4692} \MusicPDFRef{4693} \MusicPDFRef{4694} \MusicPDFRef{4695} \MusicPDFRef{4696}]}
\MusicPDFSetAnnotations{501}{[\MusicPDFRef{4697} \MusicPDFRef{4698} \MusicPDFRef{4699} \MusicPDFRef{4700} \MusicPDFRef{4701} \MusicPDFRef{4702} \MusicPDFRef{4703} \MusicPDFRef{4704} \MusicPDFRef{4705} \MusicPDFRef{4706} \MusicPDFRef{4707} \MusicPDFRef{4708} \MusicPDFRef{4709} \MusicPDFRef{4710}]}
\MusicPDFSetAnnotations{502}{[\MusicPDFRef{4711} \MusicPDFRef{4712} \MusicPDFRef{4713} \MusicPDFRef{4714} \MusicPDFRef{4715} \MusicPDFRef{4716} \MusicPDFRef{4717} \MusicPDFRef{4718} \MusicPDFRef{4719}]}
\MusicPDFSetAnnotations{503}{[\MusicPDFRef{4720} \MusicPDFRef{4721} \MusicPDFRef{4722} \MusicPDFRef{4723} \MusicPDFRef{4724} \MusicPDFRef{4725} \MusicPDFRef{4726} \MusicPDFRef{4727} \MusicPDFRef{4728} \MusicPDFRef{4729} \MusicPDFRef{4730} \MusicPDFRef{4731} \MusicPDFRef{4732} \MusicPDFRef{4733} \MusicPDFRef{4734} \MusicPDFRef{4735} \MusicPDFRef{4736} \MusicPDFRef{4737} \MusicPDFRef{4738}]}
\MusicPDFSetAnnotations{504}{[\MusicPDFRef{4739} \MusicPDFRef{4740} \MusicPDFRef{4741} \MusicPDFRef{4742} \MusicPDFRef{4743} \MusicPDFRef{4744} \MusicPDFRef{4745} \MusicPDFRef{4746} \MusicPDFRef{4747} \MusicPDFRef{4748} \MusicPDFRef{4749} \MusicPDFRef{4750}]}
\MusicPDFSetAnnotations{505}{[\MusicPDFRef{4751} \MusicPDFRef{4752} \MusicPDFRef{4753} \MusicPDFRef{4754} \MusicPDFRef{4755} \MusicPDFRef{4756}]}
\MusicPDFSetAnnotations{506}{[\MusicPDFRef{4757} \MusicPDFRef{4758} \MusicPDFRef{4759} \MusicPDFRef{4760} \MusicPDFRef{4761} \MusicPDFRef{4762} \MusicPDFRef{4763} \MusicPDFRef{4764} \MusicPDFRef{4765} \MusicPDFRef{4766} \MusicPDFRef{4767} \MusicPDFRef{4768} \MusicPDFRef{4769} \MusicPDFRef{4770} \MusicPDFRef{4771} \MusicPDFRef{4772} \MusicPDFRef{4773} \MusicPDFRef{4774}]}
\MusicPDFSetAnnotations{507}{[\MusicPDFRef{4775} \MusicPDFRef{4776} \MusicPDFRef{4777} \MusicPDFRef{4778} \MusicPDFRef{4779} \MusicPDFRef{4780} \MusicPDFRef{4781} \MusicPDFRef{4782} \MusicPDFRef{4783}]}
\MusicPDFSetAnnotations{508}{[\MusicPDFRef{4784} \MusicPDFRef{4785} \MusicPDFRef{4786} \MusicPDFRef{4787}]}
\MusicPDFSetAnnotations{509}{[\MusicPDFRef{4788} \MusicPDFRef{4789} \MusicPDFRef{4790} \MusicPDFRef{4791} \MusicPDFRef{4792} \MusicPDFRef{4793} \MusicPDFRef{4794} \MusicPDFRef{4795} \MusicPDFRef{4796} \MusicPDFRef{4797} \MusicPDFRef{4798} \MusicPDFRef{4799} \MusicPDFRef{4800} \MusicPDFRef{4801} \MusicPDFRef{4802} \MusicPDFRef{4803}]}
\MusicPDFSetAnnotations{510}{[\MusicPDFRef{4804} \MusicPDFRef{4805} \MusicPDFRef{4806} \MusicPDFRef{4807} \MusicPDFRef{4808} \MusicPDFRef{4809} \MusicPDFRef{4810} \MusicPDFRef{4811} \MusicPDFRef{4812} \MusicPDFRef{4813} \MusicPDFRef{4814} \MusicPDFRef{4815} \MusicPDFRef{4816} \MusicPDFRef{4817} \MusicPDFRef{4818}]}
\MusicPDFSetAnnotations{511}{[\MusicPDFRef{4819} \MusicPDFRef{4820} \MusicPDFRef{4821} \MusicPDFRef{4822} \MusicPDFRef{4823} \MusicPDFRef{4824} \MusicPDFRef{4825} \MusicPDFRef{4826} \MusicPDFRef{4827} \MusicPDFRef{4828} \MusicPDFRef{4829} \MusicPDFRef{4830}]}
\MusicPDFSetAnnotations{512}{[\MusicPDFRef{4831} \MusicPDFRef{4832} \MusicPDFRef{4833} \MusicPDFRef{4834} \MusicPDFRef{4835} \MusicPDFRef{4836} \MusicPDFRef{4837} \MusicPDFRef{4838} \MusicPDFRef{4839} \MusicPDFRef{4840}]}
\MusicPDFSetAnnotations{513}{[\MusicPDFRef{4841} \MusicPDFRef{4842} \MusicPDFRef{4843} \MusicPDFRef{4844} \MusicPDFRef{4845} \MusicPDFRef{4846}]}
\MusicPDFSetAnnotations{514}{[\MusicPDFRef{4847} \MusicPDFRef{4848} \MusicPDFRef{4849} \MusicPDFRef{4850} \MusicPDFRef{4851} \MusicPDFRef{4852} \MusicPDFRef{4853} \MusicPDFRef{4854} \MusicPDFRef{4855} \MusicPDFRef{4856} \MusicPDFRef{4857} \MusicPDFRef{4858}]}
\MusicPDFSetAnnotations{515}{[\MusicPDFRef{4859} \MusicPDFRef{4860} \MusicPDFRef{4861} \MusicPDFRef{4862} \MusicPDFRef{4863} \MusicPDFRef{4864} \MusicPDFRef{4865} \MusicPDFRef{4866} \MusicPDFRef{4867}]}
\MusicPDFSetAnnotations{516}{[\MusicPDFRef{4868} \MusicPDFRef{4869} \MusicPDFRef{4870} \MusicPDFRef{4871} \MusicPDFRef{4872} \MusicPDFRef{4873}]}
\MusicPDFSetAnnotations{517}{[\MusicPDFRef{4874}]}
\MusicPDFSetAnnotations{518}{[\MusicPDFRef{4875} \MusicPDFRef{4876} \MusicPDFRef{4877} \MusicPDFRef{4878} \MusicPDFRef{4879} \MusicPDFRef{4880}]}
\MusicPDFSetAnnotations{519}{[\MusicPDFRef{4881} \MusicPDFRef{4882} \MusicPDFRef{4883} \MusicPDFRef{4884} \MusicPDFRef{4885} \MusicPDFRef{4886} \MusicPDFRef{4887} \MusicPDFRef{4888}]}
\MusicPDFSetAnnotations{520}{[\MusicPDFRef{4889} \MusicPDFRef{4890} \MusicPDFRef{4891} \MusicPDFRef{4892} \MusicPDFRef{4893} \MusicPDFRef{4894} \MusicPDFRef{4895} \MusicPDFRef{4896}]}
\MusicPDFSetAnnotations{521}{[\MusicPDFRef{4897} \MusicPDFRef{4898} \MusicPDFRef{4899} \MusicPDFRef{4900}]}
\MusicPDFSetAnnotations{522}{[\MusicPDFRef{4901} \MusicPDFRef{4902} \MusicPDFRef{4903} \MusicPDFRef{4904} \MusicPDFRef{4905} \MusicPDFRef{4906} \MusicPDFRef{4907} \MusicPDFRef{4908}]}
\MusicPDFSetAnnotations{523}{[\MusicPDFRef{4909} \MusicPDFRef{4910} \MusicPDFRef{4911} \MusicPDFRef{4912} \MusicPDFRef{4913}]}
\MusicPDFSetAnnotations{524}{[\MusicPDFRef{4914} \MusicPDFRef{4915} \MusicPDFRef{4916}]}
\MusicPDFSetAnnotations{525}{[\MusicPDFRef{4917} \MusicPDFRef{4918} \MusicPDFRef{4919} \MusicPDFRef{4920} \MusicPDFRef{4921} \MusicPDFRef{4922} \MusicPDFRef{4923} \MusicPDFRef{4924} \MusicPDFRef{4925}]}
\MusicPDFSetAnnotations{526}{[\MusicPDFRef{4926} \MusicPDFRef{4927} \MusicPDFRef{4928} \MusicPDFRef{4929} \MusicPDFRef{4930} \MusicPDFRef{4931} \MusicPDFRef{4932} \MusicPDFRef{4933} \MusicPDFRef{4934} \MusicPDFRef{4935} \MusicPDFRef{4936} \MusicPDFRef{4937} \MusicPDFRef{4938} \MusicPDFRef{4939} \MusicPDFRef{4940} \MusicPDFRef{4941} \MusicPDFRef{4942} \MusicPDFRef{4943} \MusicPDFRef{4944} \MusicPDFRef{4945}]}
\MusicPDFSetAnnotations{527}{[\MusicPDFRef{4946} \MusicPDFRef{4947} \MusicPDFRef{4948} \MusicPDFRef{4949} \MusicPDFRef{4950} \MusicPDFRef{4951} \MusicPDFRef{4952} \MusicPDFRef{4953} \MusicPDFRef{4954} \MusicPDFRef{4955} \MusicPDFRef{4956} \MusicPDFRef{4957}]}
\MusicPDFSetAnnotations{528}{[\MusicPDFRef{4958} \MusicPDFRef{4959} \MusicPDFRef{4960} \MusicPDFRef{4961} \MusicPDFRef{4962} \MusicPDFRef{4963} \MusicPDFRef{4964} \MusicPDFRef{4965} \MusicPDFRef{4966} \MusicPDFRef{4967} \MusicPDFRef{4968} \MusicPDFRef{4969} \MusicPDFRef{4970} \MusicPDFRef{4971} \MusicPDFRef{4972} \MusicPDFRef{4973} \MusicPDFRef{4974}]}
\MusicPDFSetAnnotations{529}{[\MusicPDFRef{4975} \MusicPDFRef{4976} \MusicPDFRef{4977} \MusicPDFRef{4978} \MusicPDFRef{4979}]}
\MusicPDFSetAnnotations{530}{[\MusicPDFRef{4980} \MusicPDFRef{4981} \MusicPDFRef{4982} \MusicPDFRef{4983} \MusicPDFRef{4984} \MusicPDFRef{4985} \MusicPDFRef{4986} \MusicPDFRef{4987} \MusicPDFRef{4988} \MusicPDFRef{4989} \MusicPDFRef{4990} \MusicPDFRef{4991} \MusicPDFRef{4992} \MusicPDFRef{4993} \MusicPDFRef{4994}]}
\MusicPDFSetAnnotations{531}{[\MusicPDFRef{4995} \MusicPDFRef{4996} \MusicPDFRef{4997} \MusicPDFRef{4998} \MusicPDFRef{4999} \MusicPDFRef{5000} \MusicPDFRef{5001} \MusicPDFRef{5002} \MusicPDFRef{5003} \MusicPDFRef{5004} \MusicPDFRef{5005} \MusicPDFRef{5006}]}
\MusicPDFSetAnnotations{532}{[\MusicPDFRef{5007} \MusicPDFRef{5008} \MusicPDFRef{5009} \MusicPDFRef{5010} \MusicPDFRef{5011} \MusicPDFRef{5012} \MusicPDFRef{5013} \MusicPDFRef{5014} \MusicPDFRef{5015} \MusicPDFRef{5016} \MusicPDFRef{5017} \MusicPDFRef{5018} \MusicPDFRef{5019} \MusicPDFRef{5020}]}
\MusicPDFSetAnnotations{533}{[\MusicPDFRef{5021} \MusicPDFRef{5022} \MusicPDFRef{5023} \MusicPDFRef{5024} \MusicPDFRef{5025} \MusicPDFRef{5026} \MusicPDFRef{5027} \MusicPDFRef{5028} \MusicPDFRef{5029} \MusicPDFRef{5030}]}
\MusicPDFSetAnnotations{534}{[\MusicPDFRef{5031} \MusicPDFRef{5032} \MusicPDFRef{5033} \MusicPDFRef{5034} \MusicPDFRef{5035}]}
\MusicPDFSetAnnotations{535}{[\MusicPDFRef{5036} \MusicPDFRef{5037} \MusicPDFRef{5038} \MusicPDFRef{5039} \MusicPDFRef{5040} \MusicPDFRef{5041} \MusicPDFRef{5042} \MusicPDFRef{5043}]}
\MusicPDFSetAnnotations{536}{[\MusicPDFRef{5044} \MusicPDFRef{5045}]}
\MusicPDFSetAnnotations{537}{[\MusicPDFRef{5046} \MusicPDFRef{5047} \MusicPDFRef{5048} \MusicPDFRef{5049} \MusicPDFRef{5050} \MusicPDFRef{5051} \MusicPDFRef{5052} \MusicPDFRef{5053} \MusicPDFRef{5054} \MusicPDFRef{5055} \MusicPDFRef{5056} \MusicPDFRef{5057} \MusicPDFRef{5058} \MusicPDFRef{5059} \MusicPDFRef{5060} \MusicPDFRef{5061} \MusicPDFRef{5062} \MusicPDFRef{5063} \MusicPDFRef{5064} \MusicPDFRef{5065} \MusicPDFRef{5066} \MusicPDFRef{5067} \MusicPDFRef{5068} \MusicPDFRef{5069} \MusicPDFRef{5070} \MusicPDFRef{5071}]}
\MusicPDFSetAnnotations{538}{[\MusicPDFRef{5072} \MusicPDFRef{5073} \MusicPDFRef{5074} \MusicPDFRef{5075} \MusicPDFRef{5076} \MusicPDFRef{5077} \MusicPDFRef{5078} \MusicPDFRef{5079} \MusicPDFRef{5080} \MusicPDFRef{5081} \MusicPDFRef{5082} \MusicPDFRef{5083} \MusicPDFRef{5084} \MusicPDFRef{5085} \MusicPDFRef{5086} \MusicPDFRef{5087} \MusicPDFRef{5088} \MusicPDFRef{5089} \MusicPDFRef{5090} \MusicPDFRef{5091} \MusicPDFRef{5092} \MusicPDFRef{5093} \MusicPDFRef{5094} \MusicPDFRef{5095} \MusicPDFRef{5096} \MusicPDFRef{5097}]}
\MusicPDFSetAnnotations{539}{[\MusicPDFRef{5098} \MusicPDFRef{5099} \MusicPDFRef{5100} \MusicPDFRef{5101} \MusicPDFRef{5102} \MusicPDFRef{5103} \MusicPDFRef{5104}]}
\MusicPDFSetAnnotations{540}{[\MusicPDFRef{5105}]}
\MusicPDFSetAnnotations{541}{[\MusicPDFRef{5106} \MusicPDFRef{5107} \MusicPDFRef{5108} \MusicPDFRef{5109} \MusicPDFRef{5110} \MusicPDFRef{5111} \MusicPDFRef{5112} \MusicPDFRef{5113}]}
\MusicPDFSetAnnotations{542}{[\MusicPDFRef{5114} \MusicPDFRef{5115} \MusicPDFRef{5116} \MusicPDFRef{5117} \MusicPDFRef{5118} \MusicPDFRef{5119}]}
\MusicPDFSetAnnotations{543}{[\MusicPDFRef{5120}]}
\MusicPDFSetAnnotations{544}{[\MusicPDFRef{5121} \MusicPDFRef{5122} \MusicPDFRef{5123} \MusicPDFRef{5124}]}
\MusicPDFSetAnnotations{545}{[\MusicPDFRef{5125} \MusicPDFRef{5126} \MusicPDFRef{5127} \MusicPDFRef{5128} \MusicPDFRef{5129} \MusicPDFRef{5130}]}
\MusicPDFSetAnnotations{546}{[\MusicPDFRef{5131} \MusicPDFRef{5132} \MusicPDFRef{5133} \MusicPDFRef{5134} \MusicPDFRef{5135}]}
\MusicPDFSetAnnotations{547}{[\MusicPDFRef{5136} \MusicPDFRef{5137} \MusicPDFRef{5138} \MusicPDFRef{5139} \MusicPDFRef{5140}]}
\MusicPDFSetAnnotations{548}{[\MusicPDFRef{5141} \MusicPDFRef{5142}]}
\MusicPDFSetAnnotations{549}{[\MusicPDFRef{5143} \MusicPDFRef{5144} \MusicPDFRef{5145} \MusicPDFRef{5146}]}
\MusicPDFSetAnnotations{550}{[\MusicPDFRef{5147} \MusicPDFRef{5148} \MusicPDFRef{5149} \MusicPDFRef{5150} \MusicPDFRef{5151} \MusicPDFRef{5152} \MusicPDFRef{5153} \MusicPDFRef{5154} \MusicPDFRef{5155} \MusicPDFRef{5156}]}
\MusicPDFSetAnnotations{551}{[\MusicPDFRef{5157} \MusicPDFRef{5158} \MusicPDFRef{5159} \MusicPDFRef{5160} \MusicPDFRef{5161} \MusicPDFRef{5162} \MusicPDFRef{5163} \MusicPDFRef{5164}]}
\MusicPDFSetAnnotations{552}{[\MusicPDFRef{5165} \MusicPDFRef{5166} \MusicPDFRef{5167} \MusicPDFRef{5168} \MusicPDFRef{5169} \MusicPDFRef{5170} \MusicPDFRef{5171} \MusicPDFRef{5172} \MusicPDFRef{5173} \MusicPDFRef{5174} \MusicPDFRef{5175}]}
\MusicPDFSetAnnotations{553}{[\MusicPDFRef{5176} \MusicPDFRef{5177} \MusicPDFRef{5178} \MusicPDFRef{5179}]}
\MusicPDFSetAnnotations{554}{[\MusicPDFRef{5180} \MusicPDFRef{5181} \MusicPDFRef{5182} \MusicPDFRef{5183} \MusicPDFRef{5184}]}
\MusicPDFSetAnnotations{555}{[\MusicPDFRef{5185} \MusicPDFRef{5186}]}
\MusicPDFSetAnnotations{556}{[\MusicPDFRef{5187} \MusicPDFRef{5188} \MusicPDFRef{5189} \MusicPDFRef{5190} \MusicPDFRef{5191} \MusicPDFRef{5192}]}
\MusicPDFSetAnnotations{557}{[\MusicPDFRef{5193} \MusicPDFRef{5194} \MusicPDFRef{5195} \MusicPDFRef{5196} \MusicPDFRef{5197} \MusicPDFRef{5198}]}
\MusicPDFSetAnnotations{558}{[\MusicPDFRef{5199} \MusicPDFRef{5200} \MusicPDFRef{5201} \MusicPDFRef{5202} \MusicPDFRef{5203} \MusicPDFRef{5204} \MusicPDFRef{5205} \MusicPDFRef{5206} \MusicPDFRef{5207} \MusicPDFRef{5208} \MusicPDFRef{5209} \MusicPDFRef{5210} \MusicPDFRef{5211} \MusicPDFRef{5212} \MusicPDFRef{5213} \MusicPDFRef{5214} \MusicPDFRef{5215}]}
\MusicPDFSetAnnotations{559}{[\MusicPDFRef{5216} \MusicPDFRef{5217} \MusicPDFRef{5218} \MusicPDFRef{5219} \MusicPDFRef{5220} \MusicPDFRef{5221}]}
\MusicPDFSetAnnotations{560}{[\MusicPDFRef{5222} \MusicPDFRef{5223} \MusicPDFRef{5224}]}
\MusicPDFSetAnnotations{561}{[\MusicPDFRef{5225} \MusicPDFRef{5226} \MusicPDFRef{5227} \MusicPDFRef{5228} \MusicPDFRef{5229} \MusicPDFRef{5230} \MusicPDFRef{5231} \MusicPDFRef{5232} \MusicPDFRef{5233} \MusicPDFRef{5234} \MusicPDFRef{5235}]}
\MusicPDFSetAnnotations{562}{[\MusicPDFRef{5236} \MusicPDFRef{5237} \MusicPDFRef{5238} \MusicPDFRef{5239} \MusicPDFRef{5240} \MusicPDFRef{5241} \MusicPDFRef{5242} \MusicPDFRef{5243} \MusicPDFRef{5244} \MusicPDFRef{5245} \MusicPDFRef{5246}]}
\MusicPDFSetAnnotations{563}{[\MusicPDFRef{5247} \MusicPDFRef{5248} \MusicPDFRef{5249} \MusicPDFRef{5250} \MusicPDFRef{5251} \MusicPDFRef{5252} \MusicPDFRef{5253} \MusicPDFRef{5254} \MusicPDFRef{5255} \MusicPDFRef{5256} \MusicPDFRef{5257}]}
\MusicPDFSetAnnotations{564}{[\MusicPDFRef{5258} \MusicPDFRef{5259} \MusicPDFRef{5260} \MusicPDFRef{5261} \MusicPDFRef{5262} \MusicPDFRef{5263} \MusicPDFRef{5264} \MusicPDFRef{5265}]}
\MusicPDFSetAnnotations{565}{[\MusicPDFRef{5266} \MusicPDFRef{5267} \MusicPDFRef{5268} \MusicPDFRef{5269} \MusicPDFRef{5270} \MusicPDFRef{5271}]}
\MusicPDFSetAnnotations{566}{[\MusicPDFRef{5272} \MusicPDFRef{5273} \MusicPDFRef{5274} \MusicPDFRef{5275} \MusicPDFRef{5276} \MusicPDFRef{5277} \MusicPDFRef{5278} \MusicPDFRef{5279}]}
\MusicPDFSetAnnotations{567}{[\MusicPDFRef{5280} \MusicPDFRef{5281} \MusicPDFRef{5282} \MusicPDFRef{5283} \MusicPDFRef{5284} \MusicPDFRef{5285} \MusicPDFRef{5286} \MusicPDFRef{5287} \MusicPDFRef{5288} \MusicPDFRef{5289}]}
\MusicPDFSetAnnotations{568}{[\MusicPDFRef{5290} \MusicPDFRef{5291} \MusicPDFRef{5292} \MusicPDFRef{5293} \MusicPDFRef{5294} \MusicPDFRef{5295} \MusicPDFRef{5296} \MusicPDFRef{5297} \MusicPDFRef{5298} \MusicPDFRef{5299} \MusicPDFRef{5300} \MusicPDFRef{5301} \MusicPDFRef{5302} \MusicPDFRef{5303} \MusicPDFRef{5304} \MusicPDFRef{5305} \MusicPDFRef{5306} \MusicPDFRef{5307} \MusicPDFRef{5308} \MusicPDFRef{5309}]}
\MusicPDFSetAnnotations{569}{[\MusicPDFRef{5310} \MusicPDFRef{5311} \MusicPDFRef{5312} \MusicPDFRef{5313} \MusicPDFRef{5314} \MusicPDFRef{5315} \MusicPDFRef{5316}]}
\MusicPDFSetAnnotations{570}{[\MusicPDFRef{5317} \MusicPDFRef{5318} \MusicPDFRef{5319} \MusicPDFRef{5320} \MusicPDFRef{5321}]}
\MusicPDFSetAnnotations{571}{[\MusicPDFRef{5322} \MusicPDFRef{5323} \MusicPDFRef{5324} \MusicPDFRef{5325} \MusicPDFRef{5326} \MusicPDFRef{5327} \MusicPDFRef{5328} \MusicPDFRef{5329} \MusicPDFRef{5330} \MusicPDFRef{5331} \MusicPDFRef{5332} \MusicPDFRef{5333} \MusicPDFRef{5334}]}
\MusicPDFSetAnnotations{572}{[\MusicPDFRef{5335} \MusicPDFRef{5336} \MusicPDFRef{5337} \MusicPDFRef{5338} \MusicPDFRef{5339} \MusicPDFRef{5340} \MusicPDFRef{5341} \MusicPDFRef{5342}]}
\MusicPDFSetAnnotations{573}{[\MusicPDFRef{5343} \MusicPDFRef{5344} \MusicPDFRef{5345} \MusicPDFRef{5346} \MusicPDFRef{5347} \MusicPDFRef{5348} \MusicPDFRef{5349} \MusicPDFRef{5350} \MusicPDFRef{5351} \MusicPDFRef{5352} \MusicPDFRef{5353} \MusicPDFRef{5354}]}
\MusicPDFSetAnnotations{574}{[\MusicPDFRef{5355} \MusicPDFRef{5356} \MusicPDFRef{5357} \MusicPDFRef{5358}]}
\MusicPDFSetAnnotations{575}{[\MusicPDFRef{5359} \MusicPDFRef{5360}]}
\MusicPDFSetAnnotations{576}{[\MusicPDFRef{5361} \MusicPDFRef{5362} \MusicPDFRef{5363} \MusicPDFRef{5364} \MusicPDFRef{5365} \MusicPDFRef{5366} \MusicPDFRef{5367} \MusicPDFRef{5368} \MusicPDFRef{5369}]}
\MusicPDFSetAnnotations{577}{[\MusicPDFRef{5370} \MusicPDFRef{5371} \MusicPDFRef{5372} \MusicPDFRef{5373} \MusicPDFRef{5374} \MusicPDFRef{5375} \MusicPDFRef{5376} \MusicPDFRef{5377} \MusicPDFRef{5378} \MusicPDFRef{5379} \MusicPDFRef{5380} \MusicPDFRef{5381} \MusicPDFRef{5382}]}
\MusicPDFSetAnnotations{578}{[\MusicPDFRef{5383} \MusicPDFRef{5384} \MusicPDFRef{5385} \MusicPDFRef{5386} \MusicPDFRef{5387} \MusicPDFRef{5388} \MusicPDFRef{5389} \MusicPDFRef{5390} \MusicPDFRef{5391} \MusicPDFRef{5392} \MusicPDFRef{5393} \MusicPDFRef{5394} \MusicPDFRef{5395} \MusicPDFRef{5396} \MusicPDFRef{5397} \MusicPDFRef{5398}]}
\MusicPDFSetAnnotations{579}{[\MusicPDFRef{5399} \MusicPDFRef{5400} \MusicPDFRef{5401} \MusicPDFRef{5402} \MusicPDFRef{5403} \MusicPDFRef{5404} \MusicPDFRef{5405} \MusicPDFRef{5406} \MusicPDFRef{5407} \MusicPDFRef{5408} \MusicPDFRef{5409} \MusicPDFRef{5410}]}
\MusicPDFSetAnnotations{580}{[\MusicPDFRef{5411} \MusicPDFRef{5412} \MusicPDFRef{5413} \MusicPDFRef{5414} \MusicPDFRef{5415} \MusicPDFRef{5416} \MusicPDFRef{5417} \MusicPDFRef{5418} \MusicPDFRef{5419} \MusicPDFRef{5420} \MusicPDFRef{5421} \MusicPDFRef{5422}]}
\MusicPDFSetAnnotations{581}{[\MusicPDFRef{5423} \MusicPDFRef{5424} \MusicPDFRef{5425} \MusicPDFRef{5426} \MusicPDFRef{5427} \MusicPDFRef{5428} \MusicPDFRef{5429} \MusicPDFRef{5430} \MusicPDFRef{5431} \MusicPDFRef{5432} \MusicPDFRef{5433}]}
\MusicPDFSetAnnotations{582}{[\MusicPDFRef{5434} \MusicPDFRef{5435} \MusicPDFRef{5436} \MusicPDFRef{5437} \MusicPDFRef{5438} \MusicPDFRef{5439} \MusicPDFRef{5440} \MusicPDFRef{5441} \MusicPDFRef{5442}]}
\MusicPDFSetAnnotations{583}{[\MusicPDFRef{5443} \MusicPDFRef{5444} \MusicPDFRef{5445} \MusicPDFRef{5446} \MusicPDFRef{5447} \MusicPDFRef{5448} \MusicPDFRef{5449}]}
\MusicPDFSetAnnotations{584}{[\MusicPDFRef{5450} \MusicPDFRef{5451} \MusicPDFRef{5452} \MusicPDFRef{5453} \MusicPDFRef{5454} \MusicPDFRef{5455} \MusicPDFRef{5456}]}
\MusicPDFSetAnnotations{585}{[\MusicPDFRef{5457} \MusicPDFRef{5458} \MusicPDFRef{5459} \MusicPDFRef{5460} \MusicPDFRef{5461} \MusicPDFRef{5462}]}
\MusicPDFSetAnnotations{586}{[\MusicPDFRef{5463} \MusicPDFRef{5464}]}
\MusicPDFSetAnnotations{587}{[\MusicPDFRef{5465} \MusicPDFRef{5466} \MusicPDFRef{5467} \MusicPDFRef{5468} \MusicPDFRef{5469} \MusicPDFRef{5470} \MusicPDFRef{5471} \MusicPDFRef{5472} \MusicPDFRef{5473} \MusicPDFRef{5474} \MusicPDFRef{5475} \MusicPDFRef{5476} \MusicPDFRef{5477} \MusicPDFRef{5478} \MusicPDFRef{5479} \MusicPDFRef{5480} \MusicPDFRef{5481} \MusicPDFRef{5482} \MusicPDFRef{5483} \MusicPDFRef{5484} \MusicPDFRef{5485} \MusicPDFRef{5486} \MusicPDFRef{5487} \MusicPDFRef{5488} \MusicPDFRef{5489} \MusicPDFRef{5490} \MusicPDFRef{5491} \MusicPDFRef{5492} \MusicPDFRef{5493}]}
\MusicPDFSetAnnotations{588}{[\MusicPDFRef{5494} \MusicPDFRef{5495} \MusicPDFRef{5496} \MusicPDFRef{5497} \MusicPDFRef{5498} \MusicPDFRef{5499} \MusicPDFRef{5500} \MusicPDFRef{5501} \MusicPDFRef{5502} \MusicPDFRef{5503} \MusicPDFRef{5504} \MusicPDFRef{5505} \MusicPDFRef{5506} \MusicPDFRef{5507} \MusicPDFRef{5508} \MusicPDFRef{5509} \MusicPDFRef{5510} \MusicPDFRef{5511} \MusicPDFRef{5512} \MusicPDFRef{5513} \MusicPDFRef{5514} \MusicPDFRef{5515} \MusicPDFRef{5516} \MusicPDFRef{5517} \MusicPDFRef{5518} \MusicPDFRef{5519}]}
\MusicPDFSetAnnotations{589}{[\MusicPDFRef{5520} \MusicPDFRef{5521} \MusicPDFRef{5522} \MusicPDFRef{5523} \MusicPDFRef{5524} \MusicPDFRef{5525}]}
\MusicPDFSetAnnotations{590}{[\MusicPDFRef{5526} \MusicPDFRef{5527} \MusicPDFRef{5528} \MusicPDFRef{5529} \MusicPDFRef{5530} \MusicPDFRef{5531} \MusicPDFRef{5532} \MusicPDFRef{5533} \MusicPDFRef{5534} \MusicPDFRef{5535} \MusicPDFRef{5536} \MusicPDFRef{5537} \MusicPDFRef{5538} \MusicPDFRef{5539}]}
\MusicPDFSetAnnotations{591}{[\MusicPDFRef{5540} \MusicPDFRef{5541} \MusicPDFRef{5542} \MusicPDFRef{5543} \MusicPDFRef{5544} \MusicPDFRef{5545} \MusicPDFRef{5546} \MusicPDFRef{5547} \MusicPDFRef{5548} \MusicPDFRef{5549} \MusicPDFRef{5550} \MusicPDFRef{5551} \MusicPDFRef{5552} \MusicPDFRef{5553} \MusicPDFRef{5554}]}
\MusicPDFSetAnnotations{592}{[\MusicPDFRef{5555} \MusicPDFRef{5556} \MusicPDFRef{5557} \MusicPDFRef{5558} \MusicPDFRef{5559} \MusicPDFRef{5560} \MusicPDFRef{5561} \MusicPDFRef{5562}]}
\MusicPDFSetAnnotations{593}{[\MusicPDFRef{5563}]}
\MusicPDFSetAnnotations{594}{[\MusicPDFRef{5564} \MusicPDFRef{5565}]}
\MusicPDFSetAnnotations{595}{[\MusicPDFRef{5566} \MusicPDFRef{5567} \MusicPDFRef{5568} \MusicPDFRef{5569} \MusicPDFRef{5570} \MusicPDFRef{5571} \MusicPDFRef{5572} \MusicPDFRef{5573} \MusicPDFRef{5574} \MusicPDFRef{5575}]}
\MusicPDFSetAnnotations{596}{[\MusicPDFRef{5576} \MusicPDFRef{5577}]}
\MusicPDFSetAnnotations{597}{[\MusicPDFRef{5578} \MusicPDFRef{5579} \MusicPDFRef{5580} \MusicPDFRef{5581}]}
\MusicPDFSetAnnotations{598}{[\MusicPDFRef{5582} \MusicPDFRef{5583} \MusicPDFRef{5584} \MusicPDFRef{5585}]}
\MusicPDFSetAnnotations{599}{[\MusicPDFRef{5586} \MusicPDFRef{5587} \MusicPDFRef{5588} \MusicPDFRef{5589} \MusicPDFRef{5590} \MusicPDFRef{5591} \MusicPDFRef{5592} \MusicPDFRef{5593} \MusicPDFRef{5594}]}
\MusicPDFSetAnnotations{600}{[\MusicPDFRef{5595} \MusicPDFRef{5596} \MusicPDFRef{5597} \MusicPDFRef{5598} \MusicPDFRef{5599} \MusicPDFRef{5600} \MusicPDFRef{5601} \MusicPDFRef{5602} \MusicPDFRef{5603} \MusicPDFRef{5604} \MusicPDFRef{5605} \MusicPDFRef{5606}]}
\MusicPDFSetAnnotations{601}{[\MusicPDFRef{5607} \MusicPDFRef{5608} \MusicPDFRef{5609} \MusicPDFRef{5610} \MusicPDFRef{5611} \MusicPDFRef{5612} \MusicPDFRef{5613} \MusicPDFRef{5614}]}
\MusicPDFSetAnnotations{602}{[\MusicPDFRef{5615}]}
\MusicPDFSetAnnotations{603}{[\MusicPDFRef{5616} \MusicPDFRef{5617} \MusicPDFRef{5618} \MusicPDFRef{5619} \MusicPDFRef{5620} \MusicPDFRef{5621} \MusicPDFRef{5622} \MusicPDFRef{5623} \MusicPDFRef{5624} \MusicPDFRef{5625}]}
\MusicPDFSetAnnotations{604}{[\MusicPDFRef{5626} \MusicPDFRef{5627} \MusicPDFRef{5628}]}
\MusicPDFSetAnnotations{605}{[\MusicPDFRef{5629} \MusicPDFRef{5630}]}
\MusicPDFSetAnnotations{606}{[\MusicPDFRef{5631} \MusicPDFRef{5632} \MusicPDFRef{5633} \MusicPDFRef{5634} \MusicPDFRef{5635} \MusicPDFRef{5636} \MusicPDFRef{5637} \MusicPDFRef{5638} \MusicPDFRef{5639} \MusicPDFRef{5640} \MusicPDFRef{5641} \MusicPDFRef{5642} \MusicPDFRef{5643} \MusicPDFRef{5644} \MusicPDFRef{5645} \MusicPDFRef{5646}]}
\MusicPDFSetAnnotations{607}{[\MusicPDFRef{5647} \MusicPDFRef{5648} \MusicPDFRef{5649}]}
\MusicPDFSetAnnotations{608}{[\MusicPDFRef{5650} \MusicPDFRef{5651} \MusicPDFRef{5652} \MusicPDFRef{5653} \MusicPDFRef{5654} \MusicPDFRef{5655} \MusicPDFRef{5656} \MusicPDFRef{5657}]}
\MusicPDFSetAnnotations{609}{[\MusicPDFRef{5658} \MusicPDFRef{5659} \MusicPDFRef{5660} \MusicPDFRef{5661} \MusicPDFRef{5662} \MusicPDFRef{5663}]}
\MusicPDFSetAnnotations{610}{[\MusicPDFRef{5664} \MusicPDFRef{5665} \MusicPDFRef{5666} \MusicPDFRef{5667} \MusicPDFRef{5668} \MusicPDFRef{5669} \MusicPDFRef{5670} \MusicPDFRef{5671} \MusicPDFRef{5672} \MusicPDFRef{5673}]}
\MusicPDFSetAnnotations{611}{[\MusicPDFRef{5674} \MusicPDFRef{5675} \MusicPDFRef{5676} \MusicPDFRef{5677} \MusicPDFRef{5678} \MusicPDFRef{5679} \MusicPDFRef{5680} \MusicPDFRef{5681} \MusicPDFRef{5682} \MusicPDFRef{5683} \MusicPDFRef{5684} \MusicPDFRef{5685} \MusicPDFRef{5686} \MusicPDFRef{5687} \MusicPDFRef{5688} \MusicPDFRef{5689} \MusicPDFRef{5690} \MusicPDFRef{5691} \MusicPDFRef{5692} \MusicPDFRef{5693} \MusicPDFRef{5694} \MusicPDFRef{5695}]}
\MusicPDFSetAnnotations{612}{[\MusicPDFRef{5696} \MusicPDFRef{5697} \MusicPDFRef{5698} \MusicPDFRef{5699} \MusicPDFRef{5700} \MusicPDFRef{5701} \MusicPDFRef{5702} \MusicPDFRef{5703}]}
\MusicPDFSetAnnotations{613}{[\MusicPDFRef{5704} \MusicPDFRef{5705} \MusicPDFRef{5706} \MusicPDFRef{5707} \MusicPDFRef{5708} \MusicPDFRef{5709} \MusicPDFRef{5710} \MusicPDFRef{5711} \MusicPDFRef{5712} \MusicPDFRef{5713} \MusicPDFRef{5714} \MusicPDFRef{5715} \MusicPDFRef{5716} \MusicPDFRef{5717} \MusicPDFRef{5718} \MusicPDFRef{5719} \MusicPDFRef{5720}]}
\MusicPDFSetAnnotations{614}{[\MusicPDFRef{5721} \MusicPDFRef{5722} \MusicPDFRef{5723} \MusicPDFRef{5724} \MusicPDFRef{5725} \MusicPDFRef{5726} \MusicPDFRef{5727} \MusicPDFRef{5728} \MusicPDFRef{5729} \MusicPDFRef{5730} \MusicPDFRef{5731} \MusicPDFRef{5732}]}
\MusicPDFSetAnnotations{615}{[\MusicPDFRef{5733} \MusicPDFRef{5734} \MusicPDFRef{5735} \MusicPDFRef{5736} \MusicPDFRef{5737}]}
\MusicPDFSetAnnotations{616}{[\MusicPDFRef{5738} \MusicPDFRef{5739} \MusicPDFRef{5740} \MusicPDFRef{5741} \MusicPDFRef{5742} \MusicPDFRef{5743} \MusicPDFRef{5744} \MusicPDFRef{5745} \MusicPDFRef{5746} \MusicPDFRef{5747} \MusicPDFRef{5748} \MusicPDFRef{5749} \MusicPDFRef{5750}]}
\MusicPDFSetAnnotations{617}{[\MusicPDFRef{5751} \MusicPDFRef{5752} \MusicPDFRef{5753} \MusicPDFRef{5754} \MusicPDFRef{5755} \MusicPDFRef{5756} \MusicPDFRef{5757} \MusicPDFRef{5758} \MusicPDFRef{5759} \MusicPDFRef{5760} \MusicPDFRef{5761} \MusicPDFRef{5762} \MusicPDFRef{5763} \MusicPDFRef{5764} \MusicPDFRef{5765} \MusicPDFRef{5766} \MusicPDFRef{5767} \MusicPDFRef{5768} \MusicPDFRef{5769} \MusicPDFRef{5770} \MusicPDFRef{5771}]}
\MusicPDFSetAnnotations{618}{[\MusicPDFRef{5772} \MusicPDFRef{5773} \MusicPDFRef{5774} \MusicPDFRef{5775} \MusicPDFRef{5776} \MusicPDFRef{5777} \MusicPDFRef{5778} \MusicPDFRef{5779} \MusicPDFRef{5780} \MusicPDFRef{5781} \MusicPDFRef{5782} \MusicPDFRef{5783}]}
\MusicPDFSetAnnotations{619}{[\MusicPDFRef{5784} \MusicPDFRef{5785} \MusicPDFRef{5786} \MusicPDFRef{5787} \MusicPDFRef{5788} \MusicPDFRef{5789} \MusicPDFRef{5790} \MusicPDFRef{5791} \MusicPDFRef{5792} \MusicPDFRef{5793} \MusicPDFRef{5794}]}
\MusicPDFSetAnnotations{620}{[\MusicPDFRef{5795} \MusicPDFRef{5796} \MusicPDFRef{5797} \MusicPDFRef{5798} \MusicPDFRef{5799} \MusicPDFRef{5800} \MusicPDFRef{5801} \MusicPDFRef{5802}]}
\MusicPDFSetAnnotations{621}{[\MusicPDFRef{5803} \MusicPDFRef{5804} \MusicPDFRef{5805} \MusicPDFRef{5806} \MusicPDFRef{5807} \MusicPDFRef{5808} \MusicPDFRef{5809} \MusicPDFRef{5810}]}
\MusicPDFSetAnnotations{622}{[\MusicPDFRef{5811} \MusicPDFRef{5812} \MusicPDFRef{5813} \MusicPDFRef{5814} \MusicPDFRef{5815} \MusicPDFRef{5816} \MusicPDFRef{5817} \MusicPDFRef{5818} \MusicPDFRef{5819}]}
\MusicPDFSetAnnotations{623}{[\MusicPDFRef{5820} \MusicPDFRef{5821} \MusicPDFRef{5822} \MusicPDFRef{5823} \MusicPDFRef{5824} \MusicPDFRef{5825} \MusicPDFRef{5826} \MusicPDFRef{5827} \MusicPDFRef{5828} \MusicPDFRef{5829} \MusicPDFRef{5830} \MusicPDFRef{5831}]}
\MusicPDFSetAnnotations{624}{[\MusicPDFRef{5832} \MusicPDFRef{5833} \MusicPDFRef{5834} \MusicPDFRef{5835} \MusicPDFRef{5836} \MusicPDFRef{5837}]}
\MusicPDFSetAnnotations{625}{[\MusicPDFRef{5838} \MusicPDFRef{5839} \MusicPDFRef{5840} \MusicPDFRef{5841} \MusicPDFRef{5842} \MusicPDFRef{5843} \MusicPDFRef{5844} \MusicPDFRef{5845} \MusicPDFRef{5846} \MusicPDFRef{5847} \MusicPDFRef{5848} \MusicPDFRef{5849} \MusicPDFRef{5850} \MusicPDFRef{5851} \MusicPDFRef{5852} \MusicPDFRef{5853} \MusicPDFRef{5854} \MusicPDFRef{5855} \MusicPDFRef{5856} \MusicPDFRef{5857}]}
\MusicPDFSetAnnotations{626}{[\MusicPDFRef{5858} \MusicPDFRef{5859} \MusicPDFRef{5860} \MusicPDFRef{5861} \MusicPDFRef{5862} \MusicPDFRef{5863} \MusicPDFRef{5864} \MusicPDFRef{5865} \MusicPDFRef{5866}]}
\MusicPDFSetAnnotations{627}{[\MusicPDFRef{5867} \MusicPDFRef{5868} \MusicPDFRef{5869} \MusicPDFRef{5870}]}
\MusicPDFSetAnnotations{628}{[\MusicPDFRef{5871} \MusicPDFRef{5872} \MusicPDFRef{5873} \MusicPDFRef{5874} \MusicPDFRef{5875}]}
\MusicPDFSetAnnotations{629}{[\MusicPDFRef{5876} \MusicPDFRef{5877} \MusicPDFRef{5878} \MusicPDFRef{5879} \MusicPDFRef{5880} \MusicPDFRef{5881} \MusicPDFRef{5882} \MusicPDFRef{5883}]}
\MusicPDFSetAnnotations{630}{[\MusicPDFRef{5884} \MusicPDFRef{5885} \MusicPDFRef{5886} \MusicPDFRef{5887} \MusicPDFRef{5888}]}
\MusicPDFSetAnnotations{631}{[\MusicPDFRef{5889} \MusicPDFRef{5890} \MusicPDFRef{5891} \MusicPDFRef{5892} \MusicPDFRef{5893} \MusicPDFRef{5894} \MusicPDFRef{5895} \MusicPDFRef{5896} \MusicPDFRef{5897}]}
\MusicPDFSetAnnotations{632}{[\MusicPDFRef{5898} \MusicPDFRef{5899} \MusicPDFRef{5900} \MusicPDFRef{5901} \MusicPDFRef{5902} \MusicPDFRef{5903}]}
\MusicPDFSetAnnotations{633}{[\MusicPDFRef{5904}]}
\MusicPDFSetAnnotations{634}{[\MusicPDFRef{5905} \MusicPDFRef{5906} \MusicPDFRef{5907} \MusicPDFRef{5908} \MusicPDFRef{5909} \MusicPDFRef{5910}]}
\MusicPDFSetAnnotations{635}{[\MusicPDFRef{5911} \MusicPDFRef{5912} \MusicPDFRef{5913} \MusicPDFRef{5914} \MusicPDFRef{5915} \MusicPDFRef{5916} \MusicPDFRef{5917} \MusicPDFRef{5918} \MusicPDFRef{5919} \MusicPDFRef{5920} \MusicPDFRef{5921} \MusicPDFRef{5922}]}
\MusicPDFSetAnnotations{636}{[\MusicPDFRef{5923} \MusicPDFRef{5924} \MusicPDFRef{5925} \MusicPDFRef{5926} \MusicPDFRef{5927} \MusicPDFRef{5928}]}
\MusicPDFSetAnnotations{637}{[\MusicPDFRef{5929} \MusicPDFRef{5930} \MusicPDFRef{5931} \MusicPDFRef{5932} \MusicPDFRef{5933} \MusicPDFRef{5934} \MusicPDFRef{5935}]}
\MusicPDFSetAnnotations{638}{[\MusicPDFRef{5936} \MusicPDFRef{5937} \MusicPDFRef{5938} \MusicPDFRef{5939}]}
\MusicPDFSetAnnotations{639}{[\MusicPDFRef{5940} \MusicPDFRef{5941} \MusicPDFRef{5942} \MusicPDFRef{5943} \MusicPDFRef{5944} \MusicPDFRef{5945} \MusicPDFRef{5946} \MusicPDFRef{5947} \MusicPDFRef{5948} \MusicPDFRef{5949} \MusicPDFRef{5950} \MusicPDFRef{5951}]}
\MusicPDFSetAnnotations{640}{[\MusicPDFRef{5952} \MusicPDFRef{5953} \MusicPDFRef{5954} \MusicPDFRef{5955} \MusicPDFRef{5956} \MusicPDFRef{5957} \MusicPDFRef{5958} \MusicPDFRef{5959} \MusicPDFRef{5960}]}
\MusicPDFSetAnnotations{641}{[\MusicPDFRef{5961} \MusicPDFRef{5962} \MusicPDFRef{5963} \MusicPDFRef{5964} \MusicPDFRef{5965} \MusicPDFRef{5966} \MusicPDFRef{5967}]}
\MusicPDFSetAnnotations{642}{[\MusicPDFRef{5968} \MusicPDFRef{5969} \MusicPDFRef{5970} \MusicPDFRef{5971} \MusicPDFRef{5972} \MusicPDFRef{5973} \MusicPDFRef{5974}]}
\MusicPDFSetAnnotations{643}{[\MusicPDFRef{5975} \MusicPDFRef{5976} \MusicPDFRef{5977} \MusicPDFRef{5978} \MusicPDFRef{5979} \MusicPDFRef{5980} \MusicPDFRef{5981} \MusicPDFRef{5982}]}
\MusicPDFSetAnnotations{644}{[\MusicPDFRef{5983} \MusicPDFRef{5984} \MusicPDFRef{5985} \MusicPDFRef{5986} \MusicPDFRef{5987} \MusicPDFRef{5988} \MusicPDFRef{5989} \MusicPDFRef{5990} \MusicPDFRef{5991} \MusicPDFRef{5992} \MusicPDFRef{5993} \MusicPDFRef{5994} \MusicPDFRef{5995}]}
\MusicPDFSetAnnotations{645}{[\MusicPDFRef{5996} \MusicPDFRef{5997} \MusicPDFRef{5998} \MusicPDFRef{5999} \MusicPDFRef{6000} \MusicPDFRef{6001} \MusicPDFRef{6002} \MusicPDFRef{6003}]}
\MusicPDFSetAnnotations{646}{[\MusicPDFRef{6004} \MusicPDFRef{6005} \MusicPDFRef{6006} \MusicPDFRef{6007} \MusicPDFRef{6008} \MusicPDFRef{6009}]}
\MusicPDFSetAnnotations{647}{[\MusicPDFRef{6010} \MusicPDFRef{6011} \MusicPDFRef{6012} \MusicPDFRef{6013}]}
\MusicPDFSetAnnotations{648}{[\MusicPDFRef{6014} \MusicPDFRef{6015} \MusicPDFRef{6016} \MusicPDFRef{6017} \MusicPDFRef{6018} \MusicPDFRef{6019} \MusicPDFRef{6020} \MusicPDFRef{6021} \MusicPDFRef{6022} \MusicPDFRef{6023} \MusicPDFRef{6024} \MusicPDFRef{6025} \MusicPDFRef{6026} \MusicPDFRef{6027}]}
\MusicPDFSetAnnotations{649}{[\MusicPDFRef{6028} \MusicPDFRef{6029} \MusicPDFRef{6030} \MusicPDFRef{6031}]}
\MusicPDFSetAnnotations{650}{[\MusicPDFRef{6032} \MusicPDFRef{6033} \MusicPDFRef{6034} \MusicPDFRef{6035} \MusicPDFRef{6036} \MusicPDFRef{6037}]}
\MusicPDFSetAnnotations{651}{[\MusicPDFRef{6038} \MusicPDFRef{6039} \MusicPDFRef{6040} \MusicPDFRef{6041} \MusicPDFRef{6042} \MusicPDFRef{6043}]}
\MusicPDFSetAnnotations{652}{[\MusicPDFRef{6044} \MusicPDFRef{6045} \MusicPDFRef{6046} \MusicPDFRef{6047} \MusicPDFRef{6048}]}
\MusicPDFSetAnnotations{653}{[\MusicPDFRef{6049} \MusicPDFRef{6050} \MusicPDFRef{6051} \MusicPDFRef{6052} \MusicPDFRef{6053} \MusicPDFRef{6054} \MusicPDFRef{6055} \MusicPDFRef{6056} \MusicPDFRef{6057} \MusicPDFRef{6058} \MusicPDFRef{6059} \MusicPDFRef{6060} \MusicPDFRef{6061} \MusicPDFRef{6062} \MusicPDFRef{6063} \MusicPDFRef{6064}]}
\MusicPDFSetAnnotations{655}{[\MusicPDFRef{6065} \MusicPDFRef{6066} \MusicPDFRef{6067} \MusicPDFRef{6068} \MusicPDFRef{6069} \MusicPDFRef{6070} \MusicPDFRef{6071} \MusicPDFRef{6072} \MusicPDFRef{6073} \MusicPDFRef{6074} \MusicPDFRef{6075} \MusicPDFRef{6076} \MusicPDFRef{6077} \MusicPDFRef{6078} \MusicPDFRef{6079} \MusicPDFRef{6080} \MusicPDFRef{6081} \MusicPDFRef{6082} \MusicPDFRef{6083} \MusicPDFRef{6084} \MusicPDFRef{6085} \MusicPDFRef{6086} \MusicPDFRef{6087} \MusicPDFRef{6088} \MusicPDFRef{6089} \MusicPDFRef{6090} \MusicPDFRef{6091} \MusicPDFRef{6092} \MusicPDFRef{6093} \MusicPDFRef{6094} \MusicPDFRef{6095} \MusicPDFRef{6096} \MusicPDFRef{6097} \MusicPDFRef{6098} \MusicPDFRef{6099} \MusicPDFRef{6100} \MusicPDFRef{6101}]}
\MusicPDFSetAnnotations{656}{[\MusicPDFRef{6102} \MusicPDFRef{6103} \MusicPDFRef{6104} \MusicPDFRef{6105} \MusicPDFRef{6106} \MusicPDFRef{6107} \MusicPDFRef{6108} \MusicPDFRef{6109} \MusicPDFRef{6110} \MusicPDFRef{6111} \MusicPDFRef{6112} \MusicPDFRef{6113} \MusicPDFRef{6114} \MusicPDFRef{6115} \MusicPDFRef{6116} \MusicPDFRef{6117} \MusicPDFRef{6118} \MusicPDFRef{6119} \MusicPDFRef{6120} \MusicPDFRef{6121} \MusicPDFRef{6122} \MusicPDFRef{6123} \MusicPDFRef{6124} \MusicPDFRef{6125} \MusicPDFRef{6126} \MusicPDFRef{6127}]}
\MusicPDFSetAnnotations{657}{[\MusicPDFRef{6128} \MusicPDFRef{6129} \MusicPDFRef{6130} \MusicPDFRef{6131} \MusicPDFRef{6132} \MusicPDFRef{6133} \MusicPDFRef{6134} \MusicPDFRef{6135} \MusicPDFRef{6136} \MusicPDFRef{6137} \MusicPDFRef{6138} \MusicPDFRef{6139} \MusicPDFRef{6140} \MusicPDFRef{6141} \MusicPDFRef{6142} \MusicPDFRef{6143} \MusicPDFRef{6144} \MusicPDFRef{6145}]}
\MusicPDFSetAnnotations{658}{[\MusicPDFRef{6146} \MusicPDFRef{6147} \MusicPDFRef{6148} \MusicPDFRef{6149} \MusicPDFRef{6150} \MusicPDFRef{6151} \MusicPDFRef{6152} \MusicPDFRef{6153} \MusicPDFRef{6154} \MusicPDFRef{6155} \MusicPDFRef{6156} \MusicPDFRef{6157} \MusicPDFRef{6158}]}
\MusicPDFSetAnnotations{659}{[\MusicPDFRef{6159} \MusicPDFRef{6160} \MusicPDFRef{6161} \MusicPDFRef{6162} \MusicPDFRef{6163} \MusicPDFRef{6164} \MusicPDFRef{6165} \MusicPDFRef{6166} \MusicPDFRef{6167} \MusicPDFRef{6168}]}
\MusicPDFSetAnnotations{660}{[\MusicPDFRef{6169} \MusicPDFRef{6170} \MusicPDFRef{6171} \MusicPDFRef{6172} \MusicPDFRef{6173} \MusicPDFRef{6174} \MusicPDFRef{6175} \MusicPDFRef{6176} \MusicPDFRef{6177} \MusicPDFRef{6178} \MusicPDFRef{6179}]}
\MusicPDFSetAnnotations{661}{[\MusicPDFRef{6180} \MusicPDFRef{6181} \MusicPDFRef{6182} \MusicPDFRef{6183} \MusicPDFRef{6184} \MusicPDFRef{6185} \MusicPDFRef{6186} \MusicPDFRef{6187} \MusicPDFRef{6188} \MusicPDFRef{6189} \MusicPDFRef{6190}]}
\MusicPDFSetAnnotations{662}{[\MusicPDFRef{6191} \MusicPDFRef{6192} \MusicPDFRef{6193} \MusicPDFRef{6194} \MusicPDFRef{6195} \MusicPDFRef{6196} \MusicPDFRef{6197} \MusicPDFRef{6198} \MusicPDFRef{6199} \MusicPDFRef{6200} \MusicPDFRef{6201} \MusicPDFRef{6202} \MusicPDFRef{6203} \MusicPDFRef{6204} \MusicPDFRef{6205} \MusicPDFRef{6206} \MusicPDFRef{6207} \MusicPDFRef{6208} \MusicPDFRef{6209} \MusicPDFRef{6210} \MusicPDFRef{6211} \MusicPDFRef{6212}]}
\MusicPDFSetAnnotations{663}{[\MusicPDFRef{6213} \MusicPDFRef{6214} \MusicPDFRef{6215} \MusicPDFRef{6216} \MusicPDFRef{6217} \MusicPDFRef{6218} \MusicPDFRef{6219} \MusicPDFRef{6220} \MusicPDFRef{6221} \MusicPDFRef{6222} \MusicPDFRef{6223} \MusicPDFRef{6224}]}
\MusicPDFSetAnnotations{664}{[\MusicPDFRef{6225} \MusicPDFRef{6226} \MusicPDFRef{6227} \MusicPDFRef{6228} \MusicPDFRef{6229} \MusicPDFRef{6230} \MusicPDFRef{6231}]}
\MusicPDFSetAnnotations{665}{[\MusicPDFRef{6232} \MusicPDFRef{6233} \MusicPDFRef{6234} \MusicPDFRef{6235} \MusicPDFRef{6236} \MusicPDFRef{6237} \MusicPDFRef{6238} \MusicPDFRef{6239}]}
\MusicPDFSetAnnotations{666}{[\MusicPDFRef{6240} \MusicPDFRef{6241} \MusicPDFRef{6242} \MusicPDFRef{6243} \MusicPDFRef{6244} \MusicPDFRef{6245} \MusicPDFRef{6246} \MusicPDFRef{6247} \MusicPDFRef{6248} \MusicPDFRef{6249} \MusicPDFRef{6250} \MusicPDFRef{6251}]}
\MusicPDFSetAnnotations{667}{[\MusicPDFRef{6252} \MusicPDFRef{6253} \MusicPDFRef{6254} \MusicPDFRef{6255} \MusicPDFRef{6256} \MusicPDFRef{6257} \MusicPDFRef{6258} \MusicPDFRef{6259} \MusicPDFRef{6260} \MusicPDFRef{6261} \MusicPDFRef{6262} \MusicPDFRef{6263} \MusicPDFRef{6264} \MusicPDFRef{6265}]}
\MusicPDFSetAnnotations{668}{[\MusicPDFRef{6266} \MusicPDFRef{6267} \MusicPDFRef{6268} \MusicPDFRef{6269} \MusicPDFRef{6270} \MusicPDFRef{6271} \MusicPDFRef{6272} \MusicPDFRef{6273} \MusicPDFRef{6274} \MusicPDFRef{6275} \MusicPDFRef{6276}]}
\MusicPDFSetAnnotations{669}{[\MusicPDFRef{6277} \MusicPDFRef{6278} \MusicPDFRef{6279} \MusicPDFRef{6280} \MusicPDFRef{6281} \MusicPDFRef{6282}]}
\MusicPDFSetAnnotations{670}{[\MusicPDFRef{6283} \MusicPDFRef{6284} \MusicPDFRef{6285} \MusicPDFRef{6286} \MusicPDFRef{6287} \MusicPDFRef{6288} \MusicPDFRef{6289} \MusicPDFRef{6290} \MusicPDFRef{6291}]}
\MusicPDFSetAnnotations{671}{[\MusicPDFRef{6292} \MusicPDFRef{6293} \MusicPDFRef{6294}]}
\MusicPDFSetAnnotations{672}{[\MusicPDFRef{6295} \MusicPDFRef{6296} \MusicPDFRef{6297} \MusicPDFRef{6298} \MusicPDFRef{6299} \MusicPDFRef{6300} \MusicPDFRef{6301} \MusicPDFRef{6302}]}
\MusicPDFSetAnnotations{673}{[\MusicPDFRef{6303} \MusicPDFRef{6304} \MusicPDFRef{6305} \MusicPDFRef{6306} \MusicPDFRef{6307} \MusicPDFRef{6308} \MusicPDFRef{6309}]}
\MusicPDFSetAnnotations{674}{[\MusicPDFRef{6310} \MusicPDFRef{6311} \MusicPDFRef{6312} \MusicPDFRef{6313} \MusicPDFRef{6314} \MusicPDFRef{6315} \MusicPDFRef{6316} \MusicPDFRef{6317} \MusicPDFRef{6318} \MusicPDFRef{6319} \MusicPDFRef{6320} \MusicPDFRef{6321}]}
\MusicPDFSetAnnotations{675}{[\MusicPDFRef{6322} \MusicPDFRef{6323} \MusicPDFRef{6324} \MusicPDFRef{6325} \MusicPDFRef{6326} \MusicPDFRef{6327} \MusicPDFRef{6328} \MusicPDFRef{6329} \MusicPDFRef{6330} \MusicPDFRef{6331} \MusicPDFRef{6332} \MusicPDFRef{6333}]}
\MusicPDFSetAnnotations{676}{[\MusicPDFRef{6334} \MusicPDFRef{6335} \MusicPDFRef{6336} \MusicPDFRef{6337} \MusicPDFRef{6338} \MusicPDFRef{6339} \MusicPDFRef{6340} \MusicPDFRef{6341} \MusicPDFRef{6342} \MusicPDFRef{6343} \MusicPDFRef{6344} \MusicPDFRef{6345}]}
\MusicPDFSetAnnotations{677}{[\MusicPDFRef{6346} \MusicPDFRef{6347} \MusicPDFRef{6348} \MusicPDFRef{6349} \MusicPDFRef{6350} \MusicPDFRef{6351} \MusicPDFRef{6352} \MusicPDFRef{6353} \MusicPDFRef{6354} \MusicPDFRef{6355} \MusicPDFRef{6356} \MusicPDFRef{6357} \MusicPDFRef{6358} \MusicPDFRef{6359} \MusicPDFRef{6360} \MusicPDFRef{6361}]}
\MusicPDFSetAnnotations{678}{[\MusicPDFRef{6362} \MusicPDFRef{6363} \MusicPDFRef{6364} \MusicPDFRef{6365} \MusicPDFRef{6366} \MusicPDFRef{6367}]}
\MusicPDFSetAnnotations{679}{[\MusicPDFRef{6368} \MusicPDFRef{6369} \MusicPDFRef{6370} \MusicPDFRef{6371} \MusicPDFRef{6372} \MusicPDFRef{6373} \MusicPDFRef{6374} \MusicPDFRef{6375} \MusicPDFRef{6376} \MusicPDFRef{6377} \MusicPDFRef{6378} \MusicPDFRef{6379} \MusicPDFRef{6380} \MusicPDFRef{6381}]}
\MusicPDFSetAnnotations{680}{[\MusicPDFRef{6382} \MusicPDFRef{6383} \MusicPDFRef{6384} \MusicPDFRef{6385} \MusicPDFRef{6386} \MusicPDFRef{6387} \MusicPDFRef{6388} \MusicPDFRef{6389} \MusicPDFRef{6390} \MusicPDFRef{6391} \MusicPDFRef{6392} \MusicPDFRef{6393} \MusicPDFRef{6394} \MusicPDFRef{6395}]}
\MusicPDFSetAnnotations{681}{[\MusicPDFRef{6396} \MusicPDFRef{6397} \MusicPDFRef{6398} \MusicPDFRef{6399} \MusicPDFRef{6400} \MusicPDFRef{6401} \MusicPDFRef{6402}]}
\MusicPDFSetAnnotations{682}{[\MusicPDFRef{6403} \MusicPDFRef{6404} \MusicPDFRef{6405} \MusicPDFRef{6406} \MusicPDFRef{6407} \MusicPDFRef{6408}]}
\MusicPDFSetAnnotations{683}{[\MusicPDFRef{6409} \MusicPDFRef{6410} \MusicPDFRef{6411} \MusicPDFRef{6412} \MusicPDFRef{6413} \MusicPDFRef{6414} \MusicPDFRef{6415} \MusicPDFRef{6416} \MusicPDFRef{6417} \MusicPDFRef{6418} \MusicPDFRef{6419} \MusicPDFRef{6420} \MusicPDFRef{6421} \MusicPDFRef{6422} \MusicPDFRef{6423} \MusicPDFRef{6424} \MusicPDFRef{6425}]}
\MusicPDFSetAnnotations{684}{[\MusicPDFRef{6426} \MusicPDFRef{6427} \MusicPDFRef{6428} \MusicPDFRef{6429} \MusicPDFRef{6430} \MusicPDFRef{6431} \MusicPDFRef{6432} \MusicPDFRef{6433} \MusicPDFRef{6434} \MusicPDFRef{6435} \MusicPDFRef{6436} \MusicPDFRef{6437}]}
\MusicPDFSetAnnotations{685}{[\MusicPDFRef{6438} \MusicPDFRef{6439} \MusicPDFRef{6440} \MusicPDFRef{6441} \MusicPDFRef{6442} \MusicPDFRef{6443} \MusicPDFRef{6444} \MusicPDFRef{6445} \MusicPDFRef{6446} \MusicPDFRef{6447} \MusicPDFRef{6448} \MusicPDFRef{6449} \MusicPDFRef{6450} \MusicPDFRef{6451} \MusicPDFRef{6452} \MusicPDFRef{6453}]}
\MusicPDFSetAnnotations{686}{[\MusicPDFRef{6454} \MusicPDFRef{6455} \MusicPDFRef{6456} \MusicPDFRef{6457} \MusicPDFRef{6458} \MusicPDFRef{6459} \MusicPDFRef{6460} \MusicPDFRef{6461} \MusicPDFRef{6462} \MusicPDFRef{6463} \MusicPDFRef{6464} \MusicPDFRef{6465} \MusicPDFRef{6466}]}
\MusicPDFSetAnnotations{687}{[\MusicPDFRef{6467} \MusicPDFRef{6468} \MusicPDFRef{6469} \MusicPDFRef{6470} \MusicPDFRef{6471} \MusicPDFRef{6472} \MusicPDFRef{6473} \MusicPDFRef{6474} \MusicPDFRef{6475} \MusicPDFRef{6476} \MusicPDFRef{6477} \MusicPDFRef{6478}]}
\MusicPDFSetAnnotations{688}{[\MusicPDFRef{6479} \MusicPDFRef{6480} \MusicPDFRef{6481} \MusicPDFRef{6482} \MusicPDFRef{6483} \MusicPDFRef{6484} \MusicPDFRef{6485} \MusicPDFRef{6486} \MusicPDFRef{6487} \MusicPDFRef{6488} \MusicPDFRef{6489} \MusicPDFRef{6490} \MusicPDFRef{6491} \MusicPDFRef{6492} \MusicPDFRef{6493} \MusicPDFRef{6494} \MusicPDFRef{6495} \MusicPDFRef{6496} \MusicPDFRef{6497} \MusicPDFRef{6498} \MusicPDFRef{6499}]}
\MusicPDFSetAnnotations{689}{[\MusicPDFRef{6500} \MusicPDFRef{6501} \MusicPDFRef{6502} \MusicPDFRef{6503} \MusicPDFRef{6504} \MusicPDFRef{6505} \MusicPDFRef{6506} \MusicPDFRef{6507} \MusicPDFRef{6508} \MusicPDFRef{6509} \MusicPDFRef{6510} \MusicPDFRef{6511} \MusicPDFRef{6512}]}
\MusicPDFSetAnnotations{690}{[\MusicPDFRef{6513} \MusicPDFRef{6514} \MusicPDFRef{6515} \MusicPDFRef{6516} \MusicPDFRef{6517} \MusicPDFRef{6518} \MusicPDFRef{6519} \MusicPDFRef{6520} \MusicPDFRef{6521} \MusicPDFRef{6522} \MusicPDFRef{6523} \MusicPDFRef{6524} \MusicPDFRef{6525} \MusicPDFRef{6526} \MusicPDFRef{6527}]}
\MusicPDFSetAnnotations{691}{[\MusicPDFRef{6528} \MusicPDFRef{6529} \MusicPDFRef{6530} \MusicPDFRef{6531} \MusicPDFRef{6532} \MusicPDFRef{6533} \MusicPDFRef{6534} \MusicPDFRef{6535} \MusicPDFRef{6536} \MusicPDFRef{6537} \MusicPDFRef{6538} \MusicPDFRef{6539}]}
\MusicPDFSetAnnotations{692}{[\MusicPDFRef{6540} \MusicPDFRef{6541} \MusicPDFRef{6542} \MusicPDFRef{6543} \MusicPDFRef{6544} \MusicPDFRef{6545} \MusicPDFRef{6546} \MusicPDFRef{6547} \MusicPDFRef{6548} \MusicPDFRef{6549} \MusicPDFRef{6550} \MusicPDFRef{6551}]}
\MusicPDFSetAnnotations{693}{[\MusicPDFRef{6552} \MusicPDFRef{6553} \MusicPDFRef{6554} \MusicPDFRef{6555} \MusicPDFRef{6556} \MusicPDFRef{6557} \MusicPDFRef{6558} \MusicPDFRef{6559} \MusicPDFRef{6560} \MusicPDFRef{6561}]}
\MusicPDFSetAnnotations{694}{[\MusicPDFRef{6562} \MusicPDFRef{6563} \MusicPDFRef{6564} \MusicPDFRef{6565} \MusicPDFRef{6566} \MusicPDFRef{6567} \MusicPDFRef{6568} \MusicPDFRef{6569} \MusicPDFRef{6570}]}
\MusicPDFSetAnnotations{695}{[\MusicPDFRef{6571} \MusicPDFRef{6572} \MusicPDFRef{6573} \MusicPDFRef{6574}]}
\MusicPDFSetAnnotations{696}{[\MusicPDFRef{6575} \MusicPDFRef{6576} \MusicPDFRef{6577} \MusicPDFRef{6578} \MusicPDFRef{6579} \MusicPDFRef{6580} \MusicPDFRef{6581} \MusicPDFRef{6582} \MusicPDFRef{6583} \MusicPDFRef{6584} \MusicPDFRef{6585} \MusicPDFRef{6586} \MusicPDFRef{6587}]}
\MusicPDFSetAnnotations{697}{[\MusicPDFRef{6588} \MusicPDFRef{6589} \MusicPDFRef{6590} \MusicPDFRef{6591} \MusicPDFRef{6592} \MusicPDFRef{6593} \MusicPDFRef{6594} \MusicPDFRef{6595} \MusicPDFRef{6596}]}
\MusicPDFSetAnnotations{698}{[\MusicPDFRef{6597} \MusicPDFRef{6598} \MusicPDFRef{6599} \MusicPDFRef{6600} \MusicPDFRef{6601} \MusicPDFRef{6602} \MusicPDFRef{6603} \MusicPDFRef{6604}]}
\MusicPDFSetAnnotations{699}{[\MusicPDFRef{6605} \MusicPDFRef{6606} \MusicPDFRef{6607}]}
\MusicPDFSetAnnotations{700}{[\MusicPDFRef{6608} \MusicPDFRef{6609} \MusicPDFRef{6610} \MusicPDFRef{6611} \MusicPDFRef{6612} \MusicPDFRef{6613} \MusicPDFRef{6614}]}
\MusicPDFSetAnnotations{701}{[\MusicPDFRef{6615} \MusicPDFRef{6616} \MusicPDFRef{6617} \MusicPDFRef{6618} \MusicPDFRef{6619} \MusicPDFRef{6620} \MusicPDFRef{6621} \MusicPDFRef{6622} \MusicPDFRef{6623} \MusicPDFRef{6624} \MusicPDFRef{6625} \MusicPDFRef{6626}]}
\MusicPDFSetAnnotations{702}{[\MusicPDFRef{6627} \MusicPDFRef{6628} \MusicPDFRef{6629} \MusicPDFRef{6630} \MusicPDFRef{6631} \MusicPDFRef{6632} \MusicPDFRef{6633} \MusicPDFRef{6634}]}
\MusicPDFSetAnnotations{703}{[\MusicPDFRef{6635} \MusicPDFRef{6636} \MusicPDFRef{6637} \MusicPDFRef{6638} \MusicPDFRef{6639} \MusicPDFRef{6640} \MusicPDFRef{6641} \MusicPDFRef{6642} \MusicPDFRef{6643} \MusicPDFRef{6644} \MusicPDFRef{6645} \MusicPDFRef{6646}]}
\MusicPDFSetAnnotations{704}{[\MusicPDFRef{6647} \MusicPDFRef{6648} \MusicPDFRef{6649} \MusicPDFRef{6650} \MusicPDFRef{6651} \MusicPDFRef{6652}]}
\MusicPDFSetAnnotations{705}{[\MusicPDFRef{6653} \MusicPDFRef{6654} \MusicPDFRef{6655} \MusicPDFRef{6656} \MusicPDFRef{6657} \MusicPDFRef{6658} \MusicPDFRef{6659} \MusicPDFRef{6660} \MusicPDFRef{6661} \MusicPDFRef{6662}]}
\MusicPDFSetAnnotations{706}{[\MusicPDFRef{6663} \MusicPDFRef{6664} \MusicPDFRef{6665} \MusicPDFRef{6666} \MusicPDFRef{6667} \MusicPDFRef{6668} \MusicPDFRef{6669} \MusicPDFRef{6670}]}
\MusicPDFSetAnnotations{707}{[\MusicPDFRef{6671} \MusicPDFRef{6672} \MusicPDFRef{6673} \MusicPDFRef{6674} \MusicPDFRef{6675} \MusicPDFRef{6676} \MusicPDFRef{6677} \MusicPDFRef{6678} \MusicPDFRef{6679} \MusicPDFRef{6680} \MusicPDFRef{6681} \MusicPDFRef{6682} \MusicPDFRef{6683} \MusicPDFRef{6684} \MusicPDFRef{6685} \MusicPDFRef{6686} \MusicPDFRef{6687} \MusicPDFRef{6688} \MusicPDFRef{6689} \MusicPDFRef{6690} \MusicPDFRef{6691}]}
\MusicPDFSetAnnotations{708}{[\MusicPDFRef{6692} \MusicPDFRef{6693} \MusicPDFRef{6694} \MusicPDFRef{6695} \MusicPDFRef{6696} \MusicPDFRef{6697} \MusicPDFRef{6698} \MusicPDFRef{6699} \MusicPDFRef{6700} \MusicPDFRef{6701}]}
\MusicPDFSetAnnotations{709}{[\MusicPDFRef{6702} \MusicPDFRef{6703} \MusicPDFRef{6704} \MusicPDFRef{6705} \MusicPDFRef{6706}]}
\MusicPDFSetAnnotations{710}{[\MusicPDFRef{6707} \MusicPDFRef{6708} \MusicPDFRef{6709} \MusicPDFRef{6710} \MusicPDFRef{6711} \MusicPDFRef{6712} \MusicPDFRef{6713} \MusicPDFRef{6714} \MusicPDFRef{6715} \MusicPDFRef{6716} \MusicPDFRef{6717} \MusicPDFRef{6718} \MusicPDFRef{6719} \MusicPDFRef{6720} \MusicPDFRef{6721}]}
\MusicPDFSetAnnotations{711}{[\MusicPDFRef{6722} \MusicPDFRef{6723} \MusicPDFRef{6724} \MusicPDFRef{6725} \MusicPDFRef{6726} \MusicPDFRef{6727} \MusicPDFRef{6728} \MusicPDFRef{6729} \MusicPDFRef{6730}]}
\MusicPDFSetAnnotations{712}{[\MusicPDFRef{6731} \MusicPDFRef{6732} \MusicPDFRef{6733} \MusicPDFRef{6734} \MusicPDFRef{6735} \MusicPDFRef{6736} \MusicPDFRef{6737} \MusicPDFRef{6738} \MusicPDFRef{6739} \MusicPDFRef{6740} \MusicPDFRef{6741} \MusicPDFRef{6742}]}
\MusicPDFSetAnnotations{713}{[\MusicPDFRef{6743} \MusicPDFRef{6744} \MusicPDFRef{6745} \MusicPDFRef{6746} \MusicPDFRef{6747} \MusicPDFRef{6748}]}
\MusicPDFSetAnnotations{714}{[\MusicPDFRef{6749} \MusicPDFRef{6750} \MusicPDFRef{6751} \MusicPDFRef{6752} \MusicPDFRef{6753} \MusicPDFRef{6754} \MusicPDFRef{6755} \MusicPDFRef{6756} \MusicPDFRef{6757} \MusicPDFRef{6758}]}
\MusicPDFSetAnnotations{715}{[\MusicPDFRef{6759} \MusicPDFRef{6760} \MusicPDFRef{6761} \MusicPDFRef{6762} \MusicPDFRef{6763} \MusicPDFRef{6764} \MusicPDFRef{6765} \MusicPDFRef{6766} \MusicPDFRef{6767} \MusicPDFRef{6768} \MusicPDFRef{6769}]}
\MusicPDFSetAnnotations{716}{[\MusicPDFRef{6770} \MusicPDFRef{6771} \MusicPDFRef{6772} \MusicPDFRef{6773} \MusicPDFRef{6774} \MusicPDFRef{6775} \MusicPDFRef{6776} \MusicPDFRef{6777}]}
\MusicPDFSetAnnotations{717}{[\MusicPDFRef{6778} \MusicPDFRef{6779} \MusicPDFRef{6780} \MusicPDFRef{6781} \MusicPDFRef{6782} \MusicPDFRef{6783} \MusicPDFRef{6784} \MusicPDFRef{6785}]}
\MusicPDFSetAnnotations{718}{[\MusicPDFRef{6786} \MusicPDFRef{6787} \MusicPDFRef{6788} \MusicPDFRef{6789} \MusicPDFRef{6790} \MusicPDFRef{6791} \MusicPDFRef{6792} \MusicPDFRef{6793} \MusicPDFRef{6794} \MusicPDFRef{6795} \MusicPDFRef{6796} \MusicPDFRef{6797} \MusicPDFRef{6798} \MusicPDFRef{6799} \MusicPDFRef{6800} \MusicPDFRef{6801} \MusicPDFRef{6802}]}
\MusicPDFSetAnnotations{719}{[\MusicPDFRef{6803} \MusicPDFRef{6804} \MusicPDFRef{6805} \MusicPDFRef{6806} \MusicPDFRef{6807} \MusicPDFRef{6808} \MusicPDFRef{6809} \MusicPDFRef{6810} \MusicPDFRef{6811} \MusicPDFRef{6812}]}
\MusicPDFSetAnnotations{720}{[\MusicPDFRef{6813} \MusicPDFRef{6814}]}
\MusicPDFSetAnnotations{721}{[\MusicPDFRef{6815} \MusicPDFRef{6816} \MusicPDFRef{6817} \MusicPDFRef{6818} \MusicPDFRef{6819} \MusicPDFRef{6820} \MusicPDFRef{6821}]}
\MusicPDFSetAnnotations{722}{[\MusicPDFRef{6822} \MusicPDFRef{6823} \MusicPDFRef{6824} \MusicPDFRef{6825} \MusicPDFRef{6826} \MusicPDFRef{6827} \MusicPDFRef{6828} \MusicPDFRef{6829} \MusicPDFRef{6830} \MusicPDFRef{6831}]}
\MusicPDFSetAnnotations{723}{[\MusicPDFRef{6832} \MusicPDFRef{6833} \MusicPDFRef{6834} \MusicPDFRef{6835} \MusicPDFRef{6836} \MusicPDFRef{6837} \MusicPDFRef{6838} \MusicPDFRef{6839}]}
\MusicPDFSetAnnotations{724}{[\MusicPDFRef{6840} \MusicPDFRef{6841} \MusicPDFRef{6842} \MusicPDFRef{6843} \MusicPDFRef{6844} \MusicPDFRef{6845} \MusicPDFRef{6846} \MusicPDFRef{6847} \MusicPDFRef{6848} \MusicPDFRef{6849}]}
\MusicPDFSetAnnotations{725}{[\MusicPDFRef{6850} \MusicPDFRef{6851} \MusicPDFRef{6852} \MusicPDFRef{6853} \MusicPDFRef{6854}]}
\MusicPDFSetAnnotations{726}{[\MusicPDFRef{6855} \MusicPDFRef{6856} \MusicPDFRef{6857} \MusicPDFRef{6858} \MusicPDFRef{6859} \MusicPDFRef{6860} \MusicPDFRef{6861} \MusicPDFRef{6862} \MusicPDFRef{6863}]}
\MusicPDFSetAnnotations{727}{[\MusicPDFRef{6864} \MusicPDFRef{6865} \MusicPDFRef{6866} \MusicPDFRef{6867} \MusicPDFRef{6868} \MusicPDFRef{6869} \MusicPDFRef{6870} \MusicPDFRef{6871} \MusicPDFRef{6872} \MusicPDFRef{6873}]}
\MusicPDFSetAnnotations{728}{[\MusicPDFRef{6874} \MusicPDFRef{6875}]}
\MusicPDFSetAnnotations{729}{[\MusicPDFRef{6876} \MusicPDFRef{6877} \MusicPDFRef{6878} \MusicPDFRef{6879} \MusicPDFRef{6880}]}
\MusicPDFSetAnnotations{730}{[\MusicPDFRef{6881} \MusicPDFRef{6882} \MusicPDFRef{6883} \MusicPDFRef{6884}]}
\MusicPDFSetAnnotations{731}{[\MusicPDFRef{6885} \MusicPDFRef{6886} \MusicPDFRef{6887} \MusicPDFRef{6888} \MusicPDFRef{6889} \MusicPDFRef{6890} \MusicPDFRef{6891} \MusicPDFRef{6892} \MusicPDFRef{6893} \MusicPDFRef{6894} \MusicPDFRef{6895} \MusicPDFRef{6896}]}
\MusicPDFSetAnnotations{732}{[\MusicPDFRef{6897} \MusicPDFRef{6898} \MusicPDFRef{6899} \MusicPDFRef{6900} \MusicPDFRef{6901} \MusicPDFRef{6902} \MusicPDFRef{6903} \MusicPDFRef{6904} \MusicPDFRef{6905} \MusicPDFRef{6906}]}
\MusicPDFSetAnnotations{733}{[\MusicPDFRef{6907} \MusicPDFRef{6908} \MusicPDFRef{6909} \MusicPDFRef{6910}]}
\MusicPDFSetAnnotations{734}{[\MusicPDFRef{6911} \MusicPDFRef{6912} \MusicPDFRef{6913} \MusicPDFRef{6914} \MusicPDFRef{6915} \MusicPDFRef{6916} \MusicPDFRef{6917} \MusicPDFRef{6918} \MusicPDFRef{6919} \MusicPDFRef{6920} \MusicPDFRef{6921} \MusicPDFRef{6922}]}
\MusicPDFSetAnnotations{735}{[\MusicPDFRef{6923} \MusicPDFRef{6924} \MusicPDFRef{6925} \MusicPDFRef{6926} \MusicPDFRef{6927} \MusicPDFRef{6928} \MusicPDFRef{6929} \MusicPDFRef{6930} \MusicPDFRef{6931} \MusicPDFRef{6932} \MusicPDFRef{6933} \MusicPDFRef{6934}]}
\MusicPDFSetAnnotations{736}{[\MusicPDFRef{6935} \MusicPDFRef{6936} \MusicPDFRef{6937} \MusicPDFRef{6938} \MusicPDFRef{6939} \MusicPDFRef{6940} \MusicPDFRef{6941} \MusicPDFRef{6942} \MusicPDFRef{6943} \MusicPDFRef{6944} \MusicPDFRef{6945} \MusicPDFRef{6946}]}
\MusicPDFSetAnnotations{737}{[\MusicPDFRef{6947} \MusicPDFRef{6948} \MusicPDFRef{6949} \MusicPDFRef{6950} \MusicPDFRef{6951} \MusicPDFRef{6952} \MusicPDFRef{6953} \MusicPDFRef{6954} \MusicPDFRef{6955} \MusicPDFRef{6956} \MusicPDFRef{6957} \MusicPDFRef{6958} \MusicPDFRef{6959} \MusicPDFRef{6960} \MusicPDFRef{6961} \MusicPDFRef{6962} \MusicPDFRef{6963} \MusicPDFRef{6964} \MusicPDFRef{6965}]}
\MusicPDFSetAnnotations{738}{[\MusicPDFRef{6966} \MusicPDFRef{6967} \MusicPDFRef{6968} \MusicPDFRef{6969} \MusicPDFRef{6970} \MusicPDFRef{6971}]}
\MusicPDFSetAnnotations{739}{[\MusicPDFRef{6972} \MusicPDFRef{6973} \MusicPDFRef{6974} \MusicPDFRef{6975} \MusicPDFRef{6976} \MusicPDFRef{6977} \MusicPDFRef{6978} \MusicPDFRef{6979} \MusicPDFRef{6980} \MusicPDFRef{6981} \MusicPDFRef{6982} \MusicPDFRef{6983}]}
\MusicPDFSetAnnotations{740}{[\MusicPDFRef{6984} \MusicPDFRef{6985} \MusicPDFRef{6986} \MusicPDFRef{6987}]}
\MusicPDFSetAnnotations{741}{[\MusicPDFRef{6988} \MusicPDFRef{6989} \MusicPDFRef{6990} \MusicPDFRef{6991} \MusicPDFRef{6992}]}
\MusicPDFSetAnnotations{742}{[\MusicPDFRef{6993} \MusicPDFRef{6994} \MusicPDFRef{6995} \MusicPDFRef{6996}]}
\MusicPDFSetAnnotations{743}{[\MusicPDFRef{6997} \MusicPDFRef{6998} \MusicPDFRef{6999} \MusicPDFRef{7000} \MusicPDFRef{7001} \MusicPDFRef{7002} \MusicPDFRef{7003} \MusicPDFRef{7004} \MusicPDFRef{7005} \MusicPDFRef{7006} \MusicPDFRef{7007} \MusicPDFRef{7008} \MusicPDFRef{7009} \MusicPDFRef{7010} \MusicPDFRef{7011} \MusicPDFRef{7012} \MusicPDFRef{7013} \MusicPDFRef{7014} \MusicPDFRef{7015} \MusicPDFRef{7016} \MusicPDFRef{7017} \MusicPDFRef{7018} \MusicPDFRef{7019} \MusicPDFRef{7020} \MusicPDFRef{7021} \MusicPDFRef{7022} \MusicPDFRef{7023} \MusicPDFRef{7024} \MusicPDFRef{7025} \MusicPDFRef{7026} \MusicPDFRef{7027} \MusicPDFRef{7028} \MusicPDFRef{7029} \MusicPDFRef{7030} \MusicPDFRef{7031} \MusicPDFRef{7032} \MusicPDFRef{7033} \MusicPDFRef{7034}]}
\MusicPDFSetAnnotations{744}{[\MusicPDFRef{7035} \MusicPDFRef{7036} \MusicPDFRef{7037} \MusicPDFRef{7038} \MusicPDFRef{7039} \MusicPDFRef{7040} \MusicPDFRef{7041}]}
\MusicPDFSetAnnotations{745}{[\MusicPDFRef{7042} \MusicPDFRef{7043} \MusicPDFRef{7044} \MusicPDFRef{7045} \MusicPDFRef{7046} \MusicPDFRef{7047} \MusicPDFRef{7048} \MusicPDFRef{7049} \MusicPDFRef{7050} \MusicPDFRef{7051} \MusicPDFRef{7052} \MusicPDFRef{7053} \MusicPDFRef{7054} \MusicPDFRef{7055} \MusicPDFRef{7056} \MusicPDFRef{7057} \MusicPDFRef{7058} \MusicPDFRef{7059} \MusicPDFRef{7060} \MusicPDFRef{7061} \MusicPDFRef{7062}]}
\MusicPDFSetAnnotations{746}{[\MusicPDFRef{7063} \MusicPDFRef{7064} \MusicPDFRef{7065} \MusicPDFRef{7066} \MusicPDFRef{7067} \MusicPDFRef{7068} \MusicPDFRef{7069}]}
\MusicPDFSetAnnotations{747}{[\MusicPDFRef{7070} \MusicPDFRef{7071} \MusicPDFRef{7072} \MusicPDFRef{7073} \MusicPDFRef{7074} \MusicPDFRef{7075} \MusicPDFRef{7076} \MusicPDFRef{7077} \MusicPDFRef{7078} \MusicPDFRef{7079} \MusicPDFRef{7080} \MusicPDFRef{7081}]}
\MusicPDFSetAnnotations{748}{[\MusicPDFRef{7082} \MusicPDFRef{7083} \MusicPDFRef{7084} \MusicPDFRef{7085} \MusicPDFRef{7086} \MusicPDFRef{7087} \MusicPDFRef{7088} \MusicPDFRef{7089} \MusicPDFRef{7090}]}
\MusicPDFSetAnnotations{749}{[\MusicPDFRef{7091} \MusicPDFRef{7092} \MusicPDFRef{7093} \MusicPDFRef{7094} \MusicPDFRef{7095} \MusicPDFRef{7096} \MusicPDFRef{7097}]}
\MusicPDFSetAnnotations{750}{[\MusicPDFRef{7098}]}
\MusicPDFSetAnnotations{751}{[\MusicPDFRef{7099} \MusicPDFRef{7100} \MusicPDFRef{7101}]}
\MusicPDFSetAnnotations{752}{[\MusicPDFRef{7102} \MusicPDFRef{7103} \MusicPDFRef{7104} \MusicPDFRef{7105} \MusicPDFRef{7106} \MusicPDFRef{7107} \MusicPDFRef{7108} \MusicPDFRef{7109} \MusicPDFRef{7110}]}
\MusicPDFSetAnnotations{753}{[\MusicPDFRef{7111} \MusicPDFRef{7112} \MusicPDFRef{7113} \MusicPDFRef{7114} \MusicPDFRef{7115} \MusicPDFRef{7116} \MusicPDFRef{7117} \MusicPDFRef{7118} \MusicPDFRef{7119} \MusicPDFRef{7120} \MusicPDFRef{7121} \MusicPDFRef{7122} \MusicPDFRef{7123}]}
\MusicPDFSetAnnotations{754}{[\MusicPDFRef{7124} \MusicPDFRef{7125} \MusicPDFRef{7126} \MusicPDFRef{7127} \MusicPDFRef{7128} \MusicPDFRef{7129} \MusicPDFRef{7130} \MusicPDFRef{7131} \MusicPDFRef{7132} \MusicPDFRef{7133} \MusicPDFRef{7134} \MusicPDFRef{7135} \MusicPDFRef{7136} \MusicPDFRef{7137} \MusicPDFRef{7138} \MusicPDFRef{7139}]}
\MusicPDFSetAnnotations{755}{[\MusicPDFRef{7140} \MusicPDFRef{7141} \MusicPDFRef{7142} \MusicPDFRef{7143} \MusicPDFRef{7144} \MusicPDFRef{7145} \MusicPDFRef{7146} \MusicPDFRef{7147} \MusicPDFRef{7148} \MusicPDFRef{7149} \MusicPDFRef{7150} \MusicPDFRef{7151} \MusicPDFRef{7152} \MusicPDFRef{7153} \MusicPDFRef{7154}]}
\MusicPDFSetAnnotations{756}{[\MusicPDFRef{7155} \MusicPDFRef{7156} \MusicPDFRef{7157} \MusicPDFRef{7158} \MusicPDFRef{7159} \MusicPDFRef{7160} \MusicPDFRef{7161} \MusicPDFRef{7162} \MusicPDFRef{7163}]}
\MusicPDFSetAnnotations{757}{[\MusicPDFRef{7164} \MusicPDFRef{7165} \MusicPDFRef{7166} \MusicPDFRef{7167} \MusicPDFRef{7168} \MusicPDFRef{7169} \MusicPDFRef{7170} \MusicPDFRef{7171} \MusicPDFRef{7172} \MusicPDFRef{7173} \MusicPDFRef{7174} \MusicPDFRef{7175} \MusicPDFRef{7176} \MusicPDFRef{7177} \MusicPDFRef{7178}]}
\MusicPDFSetAnnotations{758}{[\MusicPDFRef{7179} \MusicPDFRef{7180} \MusicPDFRef{7181} \MusicPDFRef{7182} \MusicPDFRef{7183} \MusicPDFRef{7184} \MusicPDFRef{7185} \MusicPDFRef{7186} \MusicPDFRef{7187} \MusicPDFRef{7188}]}
\MusicPDFSetAnnotations{759}{[\MusicPDFRef{7189} \MusicPDFRef{7190} \MusicPDFRef{7191} \MusicPDFRef{7192} \MusicPDFRef{7193} \MusicPDFRef{7194} \MusicPDFRef{7195} \MusicPDFRef{7196}]}
\MusicPDFSetAnnotations{760}{[\MusicPDFRef{7197} \MusicPDFRef{7198} \MusicPDFRef{7199}]}
\MusicPDFSetAnnotations{761}{[\MusicPDFRef{7200} \MusicPDFRef{7201} \MusicPDFRef{7202} \MusicPDFRef{7203} \MusicPDFRef{7204} \MusicPDFRef{7205} \MusicPDFRef{7206} \MusicPDFRef{7207} \MusicPDFRef{7208}]}
\MusicPDFSetAnnotations{762}{[\MusicPDFRef{7209} \MusicPDFRef{7210} \MusicPDFRef{7211} \MusicPDFRef{7212}]}
\MusicPDFSetAnnotations{763}{[\MusicPDFRef{7213} \MusicPDFRef{7214} \MusicPDFRef{7215} \MusicPDFRef{7216} \MusicPDFRef{7217} \MusicPDFRef{7218} \MusicPDFRef{7219}]}
\MusicPDFSetAnnotations{764}{[\MusicPDFRef{7220} \MusicPDFRef{7221} \MusicPDFRef{7222} \MusicPDFRef{7223} \MusicPDFRef{7224} \MusicPDFRef{7225} \MusicPDFRef{7226} \MusicPDFRef{7227} \MusicPDFRef{7228} \MusicPDFRef{7229} \MusicPDFRef{7230} \MusicPDFRef{7231} \MusicPDFRef{7232} \MusicPDFRef{7233} \MusicPDFRef{7234} \MusicPDFRef{7235} \MusicPDFRef{7236} \MusicPDFRef{7237} \MusicPDFRef{7238} \MusicPDFRef{7239} \MusicPDFRef{7240} \MusicPDFRef{7241} \MusicPDFRef{7242} \MusicPDFRef{7243} \MusicPDFRef{7244} \MusicPDFRef{7245}]}
\MusicPDFSetAnnotations{765}{[\MusicPDFRef{7246} \MusicPDFRef{7247} \MusicPDFRef{7248} \MusicPDFRef{7249} \MusicPDFRef{7250}]}
\MusicPDFSetAnnotations{766}{[\MusicPDFRef{7251} \MusicPDFRef{7252} \MusicPDFRef{7253} \MusicPDFRef{7254}]}
\MusicPDFSetAnnotations{767}{[\MusicPDFRef{7255} \MusicPDFRef{7256} \MusicPDFRef{7257}]}
\MusicPDFSetAnnotations{768}{[\MusicPDFRef{7258} \MusicPDFRef{7259} \MusicPDFRef{7260}]}
\MusicPDFSetAnnotations{769}{[\MusicPDFRef{7261} \MusicPDFRef{7262} \MusicPDFRef{7263}]}
\MusicPDFSetAnnotations{770}{[\MusicPDFRef{7264} \MusicPDFRef{7265} \MusicPDFRef{7266} \MusicPDFRef{7267} \MusicPDFRef{7268} \MusicPDFRef{7269} \MusicPDFRef{7270} \MusicPDFRef{7271} \MusicPDFRef{7272} \MusicPDFRef{7273} \MusicPDFRef{7274}]}
\MusicPDFSetAnnotations{771}{[\MusicPDFRef{7275} \MusicPDFRef{7276} \MusicPDFRef{7277} \MusicPDFRef{7278} \MusicPDFRef{7279} \MusicPDFRef{7280} \MusicPDFRef{7281} \MusicPDFRef{7282}]}
\MusicPDFSetAnnotations{772}{[\MusicPDFRef{7283} \MusicPDFRef{7284} \MusicPDFRef{7285} \MusicPDFRef{7286} \MusicPDFRef{7287} \MusicPDFRef{7288} \MusicPDFRef{7289}]}
\MusicPDFSetAnnotations{773}{[\MusicPDFRef{7290} \MusicPDFRef{7291} \MusicPDFRef{7292} \MusicPDFRef{7293} \MusicPDFRef{7294} \MusicPDFRef{7295} \MusicPDFRef{7296}]}
\MusicPDFSetAnnotations{774}{[\MusicPDFRef{7297} \MusicPDFRef{7298} \MusicPDFRef{7299} \MusicPDFRef{7300} \MusicPDFRef{7301}]}
\MusicPDFSetAnnotations{775}{[\MusicPDFRef{7302} \MusicPDFRef{7303} \MusicPDFRef{7304} \MusicPDFRef{7305} \MusicPDFRef{7306} \MusicPDFRef{7307} \MusicPDFRef{7308} \MusicPDFRef{7309} \MusicPDFRef{7310} \MusicPDFRef{7311} \MusicPDFRef{7312}]}
\MusicPDFSetAnnotations{776}{[\MusicPDFRef{7313} \MusicPDFRef{7314} \MusicPDFRef{7315} \MusicPDFRef{7316} \MusicPDFRef{7317} \MusicPDFRef{7318} \MusicPDFRef{7319} \MusicPDFRef{7320} \MusicPDFRef{7321} \MusicPDFRef{7322} \MusicPDFRef{7323} \MusicPDFRef{7324} \MusicPDFRef{7325} \MusicPDFRef{7326}]}
\MusicPDFSetAnnotations{777}{[\MusicPDFRef{7327} \MusicPDFRef{7328} \MusicPDFRef{7329} \MusicPDFRef{7330} \MusicPDFRef{7331} \MusicPDFRef{7332} \MusicPDFRef{7333} \MusicPDFRef{7334} \MusicPDFRef{7335}]}
\MusicPDFSetAnnotations{778}{[\MusicPDFRef{7336} \MusicPDFRef{7337} \MusicPDFRef{7338} \MusicPDFRef{7339} \MusicPDFRef{7340} \MusicPDFRef{7341} \MusicPDFRef{7342} \MusicPDFRef{7343} \MusicPDFRef{7344} \MusicPDFRef{7345} \MusicPDFRef{7346} \MusicPDFRef{7347} \MusicPDFRef{7348} \MusicPDFRef{7349} \MusicPDFRef{7350}]}
\MusicPDFSetAnnotations{779}{[\MusicPDFRef{7351} \MusicPDFRef{7352} \MusicPDFRef{7353} \MusicPDFRef{7354} \MusicPDFRef{7355} \MusicPDFRef{7356} \MusicPDFRef{7357} \MusicPDFRef{7358}]}
\MusicPDFSetAnnotations{780}{[\MusicPDFRef{7359} \MusicPDFRef{7360} \MusicPDFRef{7361} \MusicPDFRef{7362} \MusicPDFRef{7363}]}
\MusicPDFSetAnnotations{781}{[\MusicPDFRef{7364} \MusicPDFRef{7365} \MusicPDFRef{7366} \MusicPDFRef{7367} \MusicPDFRef{7368}]}
\MusicPDFSetAnnotations{782}{[\MusicPDFRef{7369}]}
\MusicPDFSetAnnotations{783}{[\MusicPDFRef{7370}]}
\MusicPDFSetAnnotations{784}{[\MusicPDFRef{7371} \MusicPDFRef{7372} \MusicPDFRef{7373} \MusicPDFRef{7374} \MusicPDFRef{7375}]}
\MusicPDFSetAnnotations{785}{[\MusicPDFRef{7376} \MusicPDFRef{7377} \MusicPDFRef{7378} \MusicPDFRef{7379} \MusicPDFRef{7380} \MusicPDFRef{7381}]}
\MusicPDFSetAnnotations{786}{[\MusicPDFRef{7382} \MusicPDFRef{7383} \MusicPDFRef{7384} \MusicPDFRef{7385}]}
\MusicPDFSetAnnotations{787}{[\MusicPDFRef{7386} \MusicPDFRef{7387} \MusicPDFRef{7388} \MusicPDFRef{7389} \MusicPDFRef{7390} \MusicPDFRef{7391} \MusicPDFRef{7392} \MusicPDFRef{7393} \MusicPDFRef{7394} \MusicPDFRef{7395}]}
\MusicPDFSetAnnotations{788}{[\MusicPDFRef{7396} \MusicPDFRef{7397} \MusicPDFRef{7398} \MusicPDFRef{7399} \MusicPDFRef{7400} \MusicPDFRef{7401} \MusicPDFRef{7402} \MusicPDFRef{7403}]}
\MusicPDFSetAnnotations{789}{[\MusicPDFRef{7404} \MusicPDFRef{7405} \MusicPDFRef{7406} \MusicPDFRef{7407} \MusicPDFRef{7408}]}
\MusicPDFSetAnnotations{790}{[\MusicPDFRef{7409} \MusicPDFRef{7410} \MusicPDFRef{7411} \MusicPDFRef{7412}]}
\MusicPDFSetAnnotations{791}{[\MusicPDFRef{7413} \MusicPDFRef{7414} \MusicPDFRef{7415} \MusicPDFRef{7416} \MusicPDFRef{7417} \MusicPDFRef{7418}]}
\MusicPDFSetAnnotations{792}{[\MusicPDFRef{7419} \MusicPDFRef{7420} \MusicPDFRef{7421} \MusicPDFRef{7422} \MusicPDFRef{7423} \MusicPDFRef{7424} \MusicPDFRef{7425} \MusicPDFRef{7426} \MusicPDFRef{7427} \MusicPDFRef{7428} \MusicPDFRef{7429}]}
\MusicPDFSetAnnotations{793}{[\MusicPDFRef{7430} \MusicPDFRef{7431} \MusicPDFRef{7432} \MusicPDFRef{7433} \MusicPDFRef{7434} \MusicPDFRef{7435} \MusicPDFRef{7436} \MusicPDFRef{7437}]}
\MusicPDFSetAnnotations{794}{[\MusicPDFRef{7438} \MusicPDFRef{7439} \MusicPDFRef{7440} \MusicPDFRef{7441} \MusicPDFRef{7442} \MusicPDFRef{7443}]}
\MusicPDFSetAnnotations{795}{[\MusicPDFRef{7444} \MusicPDFRef{7445} \MusicPDFRef{7446} \MusicPDFRef{7447} \MusicPDFRef{7448} \MusicPDFRef{7449}]}
\MusicPDFSetAnnotations{796}{[\MusicPDFRef{7450} \MusicPDFRef{7451} \MusicPDFRef{7452} \MusicPDFRef{7453} \MusicPDFRef{7454} \MusicPDFRef{7455} \MusicPDFRef{7456}]}
\MusicPDFSetAnnotations{797}{[\MusicPDFRef{7457} \MusicPDFRef{7458} \MusicPDFRef{7459} \MusicPDFRef{7460} \MusicPDFRef{7461} \MusicPDFRef{7462} \MusicPDFRef{7463}]}
\MusicPDFSetAnnotations{798}{[\MusicPDFRef{7464} \MusicPDFRef{7465} \MusicPDFRef{7466}]}
\MusicPDFSetAnnotations{799}{[\MusicPDFRef{7467} \MusicPDFRef{7468} \MusicPDFRef{7469} \MusicPDFRef{7470}]}
\MusicPDFSetAnnotations{800}{[\MusicPDFRef{7471} \MusicPDFRef{7472} \MusicPDFRef{7473} \MusicPDFRef{7474} \MusicPDFRef{7475}]}
\MusicPDFSetAnnotations{801}{[\MusicPDFRef{7476} \MusicPDFRef{7477} \MusicPDFRef{7478} \MusicPDFRef{7479} \MusicPDFRef{7480} \MusicPDFRef{7481}]}
\MusicPDFSetAnnotations{802}{[\MusicPDFRef{7482} \MusicPDFRef{7483} \MusicPDFRef{7484} \MusicPDFRef{7485} \MusicPDFRef{7486} \MusicPDFRef{7487} \MusicPDFRef{7488} \MusicPDFRef{7489} \MusicPDFRef{7490} \MusicPDFRef{7491} \MusicPDFRef{7492} \MusicPDFRef{7493} \MusicPDFRef{7494} \MusicPDFRef{7495}]}
\MusicPDFSetAnnotations{803}{[\MusicPDFRef{7496} \MusicPDFRef{7497} \MusicPDFRef{7498} \MusicPDFRef{7499}]}
\MusicPDFSetAnnotations{804}{[\MusicPDFRef{7500} \MusicPDFRef{7501} \MusicPDFRef{7502} \MusicPDFRef{7503} \MusicPDFRef{7504} \MusicPDFRef{7505}]}
\MusicPDFSetAnnotations{805}{[\MusicPDFRef{7506} \MusicPDFRef{7507} \MusicPDFRef{7508} \MusicPDFRef{7509} \MusicPDFRef{7510} \MusicPDFRef{7511} \MusicPDFRef{7512} \MusicPDFRef{7513} \MusicPDFRef{7514}]}
\MusicPDFSetAnnotations{806}{[\MusicPDFRef{7515} \MusicPDFRef{7516} \MusicPDFRef{7517} \MusicPDFRef{7518} \MusicPDFRef{7519}]}
\MusicPDFSetAnnotations{807}{[\MusicPDFRef{7520} \MusicPDFRef{7521} \MusicPDFRef{7522} \MusicPDFRef{7523} \MusicPDFRef{7524} \MusicPDFRef{7525} \MusicPDFRef{7526} \MusicPDFRef{7527} \MusicPDFRef{7528} \MusicPDFRef{7529} \MusicPDFRef{7530}]}
\MusicPDFSetAnnotations{808}{[\MusicPDFRef{7531} \MusicPDFRef{7532} \MusicPDFRef{7533} \MusicPDFRef{7534} \MusicPDFRef{7535} \MusicPDFRef{7536} \MusicPDFRef{7537} \MusicPDFRef{7538} \MusicPDFRef{7539} \MusicPDFRef{7540} \MusicPDFRef{7541} \MusicPDFRef{7542} \MusicPDFRef{7543} \MusicPDFRef{7544} \MusicPDFRef{7545} \MusicPDFRef{7546} \MusicPDFRef{7547}]}
\MusicPDFSetAnnotations{809}{[\MusicPDFRef{7548} \MusicPDFRef{7549} \MusicPDFRef{7550} \MusicPDFRef{7551} \MusicPDFRef{7552} \MusicPDFRef{7553} \MusicPDFRef{7554} \MusicPDFRef{7555} \MusicPDFRef{7556} \MusicPDFRef{7557} \MusicPDFRef{7558} \MusicPDFRef{7559} \MusicPDFRef{7560} \MusicPDFRef{7561} \MusicPDFRef{7562} \MusicPDFRef{7563} \MusicPDFRef{7564} \MusicPDFRef{7565}]}
\MusicPDFSetAnnotations{810}{[\MusicPDFRef{7566} \MusicPDFRef{7567} \MusicPDFRef{7568} \MusicPDFRef{7569} \MusicPDFRef{7570} \MusicPDFRef{7571} \MusicPDFRef{7572} \MusicPDFRef{7573} \MusicPDFRef{7574} \MusicPDFRef{7575} \MusicPDFRef{7576} \MusicPDFRef{7577} \MusicPDFRef{7578} \MusicPDFRef{7579} \MusicPDFRef{7580} \MusicPDFRef{7581} \MusicPDFRef{7582} \MusicPDFRef{7583} \MusicPDFRef{7584}]}
\MusicPDFSetAnnotations{811}{[\MusicPDFRef{7585} \MusicPDFRef{7586} \MusicPDFRef{7587} \MusicPDFRef{7588} \MusicPDFRef{7589} \MusicPDFRef{7590} \MusicPDFRef{7591} \MusicPDFRef{7592} \MusicPDFRef{7593} \MusicPDFRef{7594} \MusicPDFRef{7595} \MusicPDFRef{7596}]}
\MusicPDFSetAnnotations{812}{[\MusicPDFRef{7597} \MusicPDFRef{7598} \MusicPDFRef{7599} \MusicPDFRef{7600} \MusicPDFRef{7601}]}
\MusicPDFSetAnnotations{813}{[\MusicPDFRef{7602} \MusicPDFRef{7603} \MusicPDFRef{7604} \MusicPDFRef{7605} \MusicPDFRef{7606} \MusicPDFRef{7607} \MusicPDFRef{7608}]}
\MusicPDFSetAnnotations{814}{[\MusicPDFRef{7609} \MusicPDFRef{7610} \MusicPDFRef{7611} \MusicPDFRef{7612}]}
\MusicPDFSetAnnotations{815}{[\MusicPDFRef{7613} \MusicPDFRef{7614} \MusicPDFRef{7615} \MusicPDFRef{7616} \MusicPDFRef{7617} \MusicPDFRef{7618}]}
\MusicPDFSetAnnotations{816}{[\MusicPDFRef{7619} \MusicPDFRef{7620} \MusicPDFRef{7621} \MusicPDFRef{7622} \MusicPDFRef{7623} \MusicPDFRef{7624}]}
\MusicPDFSetAnnotations{817}{[\MusicPDFRef{7625} \MusicPDFRef{7626} \MusicPDFRef{7627} \MusicPDFRef{7628} \MusicPDFRef{7629}]}
\MusicPDFSetAnnotations{818}{[\MusicPDFRef{7630} \MusicPDFRef{7631} \MusicPDFRef{7632} \MusicPDFRef{7633}]}
\MusicPDFSetAnnotations{819}{[\MusicPDFRef{7634} \MusicPDFRef{7635} \MusicPDFRef{7636} \MusicPDFRef{7637} \MusicPDFRef{7638} \MusicPDFRef{7639} \MusicPDFRef{7640} \MusicPDFRef{7641} \MusicPDFRef{7642} \MusicPDFRef{7643} \MusicPDFRef{7644} \MusicPDFRef{7645}]}
\MusicPDFSetAnnotations{821}{[\MusicPDFRef{7646}]}
\MusicPDFSetAnnotations{823}{[\MusicPDFRef{7647} \MusicPDFRef{7648} \MusicPDFRef{7649} \MusicPDFRef{7650} \MusicPDFRef{7651} \MusicPDFRef{7652} \MusicPDFRef{7653} \MusicPDFRef{7654} \MusicPDFRef{7655} \MusicPDFRef{7656} \MusicPDFRef{7657} \MusicPDFRef{7658} \MusicPDFRef{7659} \MusicPDFRef{7660} \MusicPDFRef{7661} \MusicPDFRef{7662} \MusicPDFRef{7663} \MusicPDFRef{7664} \MusicPDFRef{7665} \MusicPDFRef{7666} \MusicPDFRef{7667} \MusicPDFRef{7668} \MusicPDFRef{7669} \MusicPDFRef{7670} \MusicPDFRef{7671} \MusicPDFRef{7672} \MusicPDFRef{7673} \MusicPDFRef{7674} \MusicPDFRef{7675} \MusicPDFRef{7676} \MusicPDFRef{7677}]}
\MusicPDFSetAnnotations{824}{[\MusicPDFRef{7678} \MusicPDFRef{7679} \MusicPDFRef{7680} \MusicPDFRef{7681} \MusicPDFRef{7682} \MusicPDFRef{7683} \MusicPDFRef{7684} \MusicPDFRef{7685} \MusicPDFRef{7686} \MusicPDFRef{7687} \MusicPDFRef{7688} \MusicPDFRef{7689} \MusicPDFRef{7690} \MusicPDFRef{7691} \MusicPDFRef{7692} \MusicPDFRef{7693} \MusicPDFRef{7694} \MusicPDFRef{7695} \MusicPDFRef{7696} \MusicPDFRef{7697} \MusicPDFRef{7698} \MusicPDFRef{7699} \MusicPDFRef{7700} \MusicPDFRef{7701} \MusicPDFRef{7702}]}
\MusicPDFSetAnnotations{825}{[\MusicPDFRef{7703} \MusicPDFRef{7704} \MusicPDFRef{7705} \MusicPDFRef{7706} \MusicPDFRef{7707} \MusicPDFRef{7708} \MusicPDFRef{7709} \MusicPDFRef{7710} \MusicPDFRef{7711} \MusicPDFRef{7712} \MusicPDFRef{7713} \MusicPDFRef{7714} \MusicPDFRef{7715} \MusicPDFRef{7716} \MusicPDFRef{7717} \MusicPDFRef{7718} \MusicPDFRef{7719} \MusicPDFRef{7720} \MusicPDFRef{7721} \MusicPDFRef{7722} \MusicPDFRef{7723} \MusicPDFRef{7724} \MusicPDFRef{7725}]}
\MusicPDFSetAnnotations{826}{[\MusicPDFRef{7726} \MusicPDFRef{7727} \MusicPDFRef{7728} \MusicPDFRef{7729} \MusicPDFRef{7730} \MusicPDFRef{7731} \MusicPDFRef{7732} \MusicPDFRef{7733} \MusicPDFRef{7734} \MusicPDFRef{7735} \MusicPDFRef{7736} \MusicPDFRef{7737} \MusicPDFRef{7738} \MusicPDFRef{7739} \MusicPDFRef{7740} \MusicPDFRef{7741} \MusicPDFRef{7742} \MusicPDFRef{7743} \MusicPDFRef{7744} \MusicPDFRef{7745} \MusicPDFRef{7746} \MusicPDFRef{7747} \MusicPDFRef{7748} \MusicPDFRef{7749} \MusicPDFRef{7750} \MusicPDFRef{7751} \MusicPDFRef{7752} \MusicPDFRef{7753} \MusicPDFRef{7754} \MusicPDFRef{7755} \MusicPDFRef{7756} \MusicPDFRef{7757} \MusicPDFRef{7758} \MusicPDFRef{7759} \MusicPDFRef{7760} \MusicPDFRef{7761} \MusicPDFRef{7762} \MusicPDFRef{7763} \MusicPDFRef{7764} \MusicPDFRef{7765} \MusicPDFRef{7766} \MusicPDFRef{7767}]}
\MusicPDFSetAnnotations{827}{[\MusicPDFRef{7768} \MusicPDFRef{7769} \MusicPDFRef{7770} \MusicPDFRef{7771} \MusicPDFRef{7772} \MusicPDFRef{7773} \MusicPDFRef{7774} \MusicPDFRef{7775} \MusicPDFRef{7776} \MusicPDFRef{7777} \MusicPDFRef{7778} \MusicPDFRef{7779} \MusicPDFRef{7780} \MusicPDFRef{7781} \MusicPDFRef{7782} \MusicPDFRef{7783} \MusicPDFRef{7784} \MusicPDFRef{7785} \MusicPDFRef{7786} \MusicPDFRef{7787} \MusicPDFRef{7788} \MusicPDFRef{7789} \MusicPDFRef{7790} \MusicPDFRef{7791} \MusicPDFRef{7792}]}
\MusicPDFSetAnnotations{828}{[\MusicPDFRef{7793} \MusicPDFRef{7794} \MusicPDFRef{7795} \MusicPDFRef{7796} \MusicPDFRef{7797} \MusicPDFRef{7798} \MusicPDFRef{7799} \MusicPDFRef{7800} \MusicPDFRef{7801} \MusicPDFRef{7802} \MusicPDFRef{7803} \MusicPDFRef{7804} \MusicPDFRef{7805} \MusicPDFRef{7806} \MusicPDFRef{7807} \MusicPDFRef{7808} \MusicPDFRef{7809} \MusicPDFRef{7810} \MusicPDFRef{7811}]}
\MusicPDFSetAnnotations{829}{[\MusicPDFRef{7812} \MusicPDFRef{7813} \MusicPDFRef{7814} \MusicPDFRef{7815} \MusicPDFRef{7816} \MusicPDFRef{7817} \MusicPDFRef{7818} \MusicPDFRef{7819} \MusicPDFRef{7820}]}
\MusicPDFSetAnnotations{830}{[\MusicPDFRef{7821} \MusicPDFRef{7822} \MusicPDFRef{7823} \MusicPDFRef{7824} \MusicPDFRef{7825} \MusicPDFRef{7826} \MusicPDFRef{7827} \MusicPDFRef{7828}]}
\MusicPDFSetAnnotations{831}{[\MusicPDFRef{7829} \MusicPDFRef{7830} \MusicPDFRef{7831} \MusicPDFRef{7832} \MusicPDFRef{7833} \MusicPDFRef{7834} \MusicPDFRef{7835} \MusicPDFRef{7836}]}
\MusicPDFSetAnnotations{832}{[\MusicPDFRef{7837} \MusicPDFRef{7838} \MusicPDFRef{7839} \MusicPDFRef{7840} \MusicPDFRef{7841} \MusicPDFRef{7842} \MusicPDFRef{7843} \MusicPDFRef{7844} \MusicPDFRef{7845}]}
\MusicPDFSetAnnotations{833}{[\MusicPDFRef{7846} \MusicPDFRef{7847} \MusicPDFRef{7848} \MusicPDFRef{7849} \MusicPDFRef{7850} \MusicPDFRef{7851} \MusicPDFRef{7852} \MusicPDFRef{7853} \MusicPDFRef{7854} \MusicPDFRef{7855} \MusicPDFRef{7856}]}
\MusicPDFSetAnnotations{834}{[\MusicPDFRef{7857} \MusicPDFRef{7858} \MusicPDFRef{7859} \MusicPDFRef{7860} \MusicPDFRef{7861} \MusicPDFRef{7862} \MusicPDFRef{7863} \MusicPDFRef{7864} \MusicPDFRef{7865} \MusicPDFRef{7866} \MusicPDFRef{7867} \MusicPDFRef{7868} \MusicPDFRef{7869} \MusicPDFRef{7870} \MusicPDFRef{7871} \MusicPDFRef{7872}]}
\MusicPDFSetAnnotations{835}{[\MusicPDFRef{7873} \MusicPDFRef{7874} \MusicPDFRef{7875} \MusicPDFRef{7876} \MusicPDFRef{7877} \MusicPDFRef{7878} \MusicPDFRef{7879} \MusicPDFRef{7880} \MusicPDFRef{7881} \MusicPDFRef{7882} \MusicPDFRef{7883} \MusicPDFRef{7884} \MusicPDFRef{7885} \MusicPDFRef{7886} \MusicPDFRef{7887} \MusicPDFRef{7888} \MusicPDFRef{7889} \MusicPDFRef{7890} \MusicPDFRef{7891} \MusicPDFRef{7892} \MusicPDFRef{7893} \MusicPDFRef{7894} \MusicPDFRef{7895} \MusicPDFRef{7896} \MusicPDFRef{7897} \MusicPDFRef{7898} \MusicPDFRef{7899} \MusicPDFRef{7900} \MusicPDFRef{7901} \MusicPDFRef{7902}]}
\MusicPDFSetAnnotations{836}{[\MusicPDFRef{7903} \MusicPDFRef{7904} \MusicPDFRef{7905} \MusicPDFRef{7906} \MusicPDFRef{7907} \MusicPDFRef{7908} \MusicPDFRef{7909} \MusicPDFRef{7910} \MusicPDFRef{7911} \MusicPDFRef{7912} \MusicPDFRef{7913} \MusicPDFRef{7914} \MusicPDFRef{7915} \MusicPDFRef{7916} \MusicPDFRef{7917} \MusicPDFRef{7918} \MusicPDFRef{7919}]}
\MusicPDFSetAnnotations{837}{[\MusicPDFRef{7920} \MusicPDFRef{7921} \MusicPDFRef{7922} \MusicPDFRef{7923} \MusicPDFRef{7924} \MusicPDFRef{7925} \MusicPDFRef{7926} \MusicPDFRef{7927} \MusicPDFRef{7928} \MusicPDFRef{7929} \MusicPDFRef{7930} \MusicPDFRef{7931} \MusicPDFRef{7932} \MusicPDFRef{7933} \MusicPDFRef{7934} \MusicPDFRef{7935} \MusicPDFRef{7936} \MusicPDFRef{7937} \MusicPDFRef{7938} \MusicPDFRef{7939} \MusicPDFRef{7940}]}
\MusicPDFSetAnnotations{838}{[\MusicPDFRef{7941} \MusicPDFRef{7942} \MusicPDFRef{7943}]}
\MusicPDFSetAnnotations{839}{[\MusicPDFRef{7944} \MusicPDFRef{7945} \MusicPDFRef{7946} \MusicPDFRef{7947} \MusicPDFRef{7948} \MusicPDFRef{7949}]}
\MusicPDFSetAnnotations{840}{[\MusicPDFRef{7950} \MusicPDFRef{7951} \MusicPDFRef{7952} \MusicPDFRef{7953} \MusicPDFRef{7954} \MusicPDFRef{7955} \MusicPDFRef{7956} \MusicPDFRef{7957} \MusicPDFRef{7958} \MusicPDFRef{7959} \MusicPDFRef{7960}]}
\MusicPDFSetAnnotations{841}{[\MusicPDFRef{7961} \MusicPDFRef{7962} \MusicPDFRef{7963} \MusicPDFRef{7964} \MusicPDFRef{7965} \MusicPDFRef{7966} \MusicPDFRef{7967} \MusicPDFRef{7968} \MusicPDFRef{7969} \MusicPDFRef{7970} \MusicPDFRef{7971} \MusicPDFRef{7972} \MusicPDFRef{7973} \MusicPDFRef{7974} \MusicPDFRef{7975} \MusicPDFRef{7976} \MusicPDFRef{7977}]}
\MusicPDFSetAnnotations{842}{[\MusicPDFRef{7978} \MusicPDFRef{7979} \MusicPDFRef{7980} \MusicPDFRef{7981} \MusicPDFRef{7982} \MusicPDFRef{7983} \MusicPDFRef{7984} \MusicPDFRef{7985} \MusicPDFRef{7986} \MusicPDFRef{7987} \MusicPDFRef{7988} \MusicPDFRef{7989} \MusicPDFRef{7990} \MusicPDFRef{7991} \MusicPDFRef{7992} \MusicPDFRef{7993} \MusicPDFRef{7994} \MusicPDFRef{7995}]}
\MusicPDFSetAnnotations{843}{[\MusicPDFRef{7996} \MusicPDFRef{7997} \MusicPDFRef{7998} \MusicPDFRef{7999} \MusicPDFRef{8000} \MusicPDFRef{8001} \MusicPDFRef{8002} \MusicPDFRef{8003} \MusicPDFRef{8004} \MusicPDFRef{8005} \MusicPDFRef{8006} \MusicPDFRef{8007} \MusicPDFRef{8008} \MusicPDFRef{8009} \MusicPDFRef{8010}]}
\MusicPDFSetAnnotations{844}{[\MusicPDFRef{8011} \MusicPDFRef{8012} \MusicPDFRef{8013} \MusicPDFRef{8014} \MusicPDFRef{8015} \MusicPDFRef{8016} \MusicPDFRef{8017} \MusicPDFRef{8018} \MusicPDFRef{8019} \MusicPDFRef{8020} \MusicPDFRef{8021} \MusicPDFRef{8022} \MusicPDFRef{8023} \MusicPDFRef{8024} \MusicPDFRef{8025}]}
\MusicPDFSetAnnotations{845}{[\MusicPDFRef{8026} \MusicPDFRef{8027} \MusicPDFRef{8028} \MusicPDFRef{8029} \MusicPDFRef{8030} \MusicPDFRef{8031} \MusicPDFRef{8032} \MusicPDFRef{8033} \MusicPDFRef{8034} \MusicPDFRef{8035}]}
\MusicPDFSetAnnotations{846}{[\MusicPDFRef{8036} \MusicPDFRef{8037} \MusicPDFRef{8038} \MusicPDFRef{8039} \MusicPDFRef{8040}]}
\MusicPDFSetAnnotations{847}{[\MusicPDFRef{8041} \MusicPDFRef{8042} \MusicPDFRef{8043} \MusicPDFRef{8044} \MusicPDFRef{8045} \MusicPDFRef{8046} \MusicPDFRef{8047} \MusicPDFRef{8048} \MusicPDFRef{8049} \MusicPDFRef{8050}]}
\MusicPDFSetAnnotations{848}{[\MusicPDFRef{8051} \MusicPDFRef{8052} \MusicPDFRef{8053} \MusicPDFRef{8054} \MusicPDFRef{8055} \MusicPDFRef{8056} \MusicPDFRef{8057} \MusicPDFRef{8058} \MusicPDFRef{8059} \MusicPDFRef{8060} \MusicPDFRef{8061} \MusicPDFRef{8062} \MusicPDFRef{8063} \MusicPDFRef{8064} \MusicPDFRef{8065} \MusicPDFRef{8066} \MusicPDFRef{8067}]}
\MusicPDFSetAnnotations{849}{[\MusicPDFRef{8068} \MusicPDFRef{8069} \MusicPDFRef{8070} \MusicPDFRef{8071} \MusicPDFRef{8072} \MusicPDFRef{8073} \MusicPDFRef{8074} \MusicPDFRef{8075} \MusicPDFRef{8076} \MusicPDFRef{8077}]}
\MusicPDFSetAnnotations{850}{[\MusicPDFRef{8078} \MusicPDFRef{8079} \MusicPDFRef{8080} \MusicPDFRef{8081} \MusicPDFRef{8082} \MusicPDFRef{8083} \MusicPDFRef{8084}]}
\MusicPDFSetAnnotations{851}{[\MusicPDFRef{8085} \MusicPDFRef{8086} \MusicPDFRef{8087} \MusicPDFRef{8088}]}
\MusicPDFSetAnnotations{852}{[\MusicPDFRef{8089} \MusicPDFRef{8090} \MusicPDFRef{8091} \MusicPDFRef{8092} \MusicPDFRef{8093} \MusicPDFRef{8094} \MusicPDFRef{8095} \MusicPDFRef{8096} \MusicPDFRef{8097} \MusicPDFRef{8098} \MusicPDFRef{8099} \MusicPDFRef{8100}]}
\MusicPDFSetAnnotations{853}{[\MusicPDFRef{8101} \MusicPDFRef{8102} \MusicPDFRef{8103} \MusicPDFRef{8104} \MusicPDFRef{8105} \MusicPDFRef{8106} \MusicPDFRef{8107} \MusicPDFRef{8108} \MusicPDFRef{8109} \MusicPDFRef{8110} \MusicPDFRef{8111}]}
\MusicPDFSetAnnotations{854}{[\MusicPDFRef{8112} \MusicPDFRef{8113} \MusicPDFRef{8114} \MusicPDFRef{8115} \MusicPDFRef{8116} \MusicPDFRef{8117} \MusicPDFRef{8118} \MusicPDFRef{8119} \MusicPDFRef{8120} \MusicPDFRef{8121} \MusicPDFRef{8122} \MusicPDFRef{8123} \MusicPDFRef{8124}]}
\MusicPDFSetAnnotations{855}{[\MusicPDFRef{8125} \MusicPDFRef{8126} \MusicPDFRef{8127} \MusicPDFRef{8128} \MusicPDFRef{8129} \MusicPDFRef{8130} \MusicPDFRef{8131} \MusicPDFRef{8132} \MusicPDFRef{8133} \MusicPDFRef{8134}]}
\MusicPDFSetAnnotations{856}{[\MusicPDFRef{8135} \MusicPDFRef{8136} \MusicPDFRef{8137} \MusicPDFRef{8138} \MusicPDFRef{8139} \MusicPDFRef{8140}]}
\MusicPDFSetAnnotations{857}{[\MusicPDFRef{8141} \MusicPDFRef{8142} \MusicPDFRef{8143} \MusicPDFRef{8144} \MusicPDFRef{8145} \MusicPDFRef{8146} \MusicPDFRef{8147} \MusicPDFRef{8148} \MusicPDFRef{8149} \MusicPDFRef{8150} \MusicPDFRef{8151} \MusicPDFRef{8152}]}
\MusicPDFSetAnnotations{858}{[\MusicPDFRef{8153} \MusicPDFRef{8154} \MusicPDFRef{8155} \MusicPDFRef{8156} \MusicPDFRef{8157} \MusicPDFRef{8158} \MusicPDFRef{8159}]}
\MusicPDFSetAnnotations{859}{[\MusicPDFRef{8160} \MusicPDFRef{8161} \MusicPDFRef{8162} \MusicPDFRef{8163} \MusicPDFRef{8164} \MusicPDFRef{8165}]}
\MusicPDFSetAnnotations{860}{[\MusicPDFRef{8166} \MusicPDFRef{8167} \MusicPDFRef{8168} \MusicPDFRef{8169} \MusicPDFRef{8170} \MusicPDFRef{8171} \MusicPDFRef{8172} \MusicPDFRef{8173} \MusicPDFRef{8174}]}
\MusicPDFSetAnnotations{861}{[\MusicPDFRef{8175} \MusicPDFRef{8176} \MusicPDFRef{8177} \MusicPDFRef{8178} \MusicPDFRef{8179} \MusicPDFRef{8180} \MusicPDFRef{8181} \MusicPDFRef{8182} \MusicPDFRef{8183} \MusicPDFRef{8184} \MusicPDFRef{8185} \MusicPDFRef{8186} \MusicPDFRef{8187} \MusicPDFRef{8188} \MusicPDFRef{8189}]}
\MusicPDFSetAnnotations{862}{[\MusicPDFRef{8190} \MusicPDFRef{8191} \MusicPDFRef{8192} \MusicPDFRef{8193} \MusicPDFRef{8194} \MusicPDFRef{8195} \MusicPDFRef{8196} \MusicPDFRef{8197} \MusicPDFRef{8198} \MusicPDFRef{8199} \MusicPDFRef{8200} \MusicPDFRef{8201}]}
\MusicPDFSetAnnotations{863}{[\MusicPDFRef{8202} \MusicPDFRef{8203} \MusicPDFRef{8204} \MusicPDFRef{8205} \MusicPDFRef{8206} \MusicPDFRef{8207} \MusicPDFRef{8208}]}
\MusicPDFSetAnnotations{864}{[\MusicPDFRef{8209} \MusicPDFRef{8210} \MusicPDFRef{8211} \MusicPDFRef{8212} \MusicPDFRef{8213} \MusicPDFRef{8214} \MusicPDFRef{8215} \MusicPDFRef{8216} \MusicPDFRef{8217} \MusicPDFRef{8218} \MusicPDFRef{8219} \MusicPDFRef{8220} \MusicPDFRef{8221} \MusicPDFRef{8222} \MusicPDFRef{8223} \MusicPDFRef{8224} \MusicPDFRef{8225} \MusicPDFRef{8226}]}
\MusicPDFSetAnnotations{865}{[\MusicPDFRef{8227} \MusicPDFRef{8228} \MusicPDFRef{8229} \MusicPDFRef{8230} \MusicPDFRef{8231} \MusicPDFRef{8232} \MusicPDFRef{8233}]}
\MusicPDFSetAnnotations{866}{[\MusicPDFRef{8234} \MusicPDFRef{8235} \MusicPDFRef{8236} \MusicPDFRef{8237} \MusicPDFRef{8238} \MusicPDFRef{8239} \MusicPDFRef{8240} \MusicPDFRef{8241} \MusicPDFRef{8242} \MusicPDFRef{8243} \MusicPDFRef{8244} \MusicPDFRef{8245}]}
\MusicPDFSetAnnotations{867}{[\MusicPDFRef{8246} \MusicPDFRef{8247} \MusicPDFRef{8248} \MusicPDFRef{8249} \MusicPDFRef{8250} \MusicPDFRef{8251} \MusicPDFRef{8252} \MusicPDFRef{8253}]}
\MusicPDFSetAnnotations{868}{[\MusicPDFRef{8254} \MusicPDFRef{8255} \MusicPDFRef{8256} \MusicPDFRef{8257} \MusicPDFRef{8258} \MusicPDFRef{8259} \MusicPDFRef{8260} \MusicPDFRef{8261} \MusicPDFRef{8262} \MusicPDFRef{8263} \MusicPDFRef{8264} \MusicPDFRef{8265} \MusicPDFRef{8266} \MusicPDFRef{8267}]}
\MusicPDFSetAnnotations{869}{[\MusicPDFRef{8268} \MusicPDFRef{8269} \MusicPDFRef{8270} \MusicPDFRef{8271} \MusicPDFRef{8272} \MusicPDFRef{8273} \MusicPDFRef{8274} \MusicPDFRef{8275} \MusicPDFRef{8276}]}
\MusicPDFSetAnnotations{870}{[\MusicPDFRef{8277} \MusicPDFRef{8278} \MusicPDFRef{8279} \MusicPDFRef{8280} \MusicPDFRef{8281} \MusicPDFRef{8282} \MusicPDFRef{8283} \MusicPDFRef{8284}]}
\MusicPDFSetAnnotations{871}{[\MusicPDFRef{8285} \MusicPDFRef{8286}]}
\MusicPDFSetAnnotations{872}{[\MusicPDFRef{8287} \MusicPDFRef{8288} \MusicPDFRef{8289} \MusicPDFRef{8290} \MusicPDFRef{8291} \MusicPDFRef{8292} \MusicPDFRef{8293} \MusicPDFRef{8294} \MusicPDFRef{8295} \MusicPDFRef{8296}]}
\MusicPDFSetAnnotations{873}{[\MusicPDFRef{8297} \MusicPDFRef{8298} \MusicPDFRef{8299} \MusicPDFRef{8300} \MusicPDFRef{8301} \MusicPDFRef{8302} \MusicPDFRef{8303} \MusicPDFRef{8304} \MusicPDFRef{8305} \MusicPDFRef{8306}]}
\MusicPDFSetAnnotations{874}{[\MusicPDFRef{8307} \MusicPDFRef{8308} \MusicPDFRef{8309} \MusicPDFRef{8310}]}
\MusicPDFSetAnnotations{875}{[\MusicPDFRef{8311} \MusicPDFRef{8312} \MusicPDFRef{8313}]}
\MusicPDFSetAnnotations{876}{[\MusicPDFRef{8314} \MusicPDFRef{8315} \MusicPDFRef{8316} \MusicPDFRef{8317} \MusicPDFRef{8318} \MusicPDFRef{8319} \MusicPDFRef{8320} \MusicPDFRef{8321}]}
\MusicPDFSetAnnotations{877}{[\MusicPDFRef{8322} \MusicPDFRef{8323} \MusicPDFRef{8324} \MusicPDFRef{8325} \MusicPDFRef{8326} \MusicPDFRef{8327} \MusicPDFRef{8328} \MusicPDFRef{8329} \MusicPDFRef{8330}]}
\MusicPDFSetAnnotations{878}{[\MusicPDFRef{8331} \MusicPDFRef{8332} \MusicPDFRef{8333} \MusicPDFRef{8334} \MusicPDFRef{8335} \MusicPDFRef{8336} \MusicPDFRef{8337} \MusicPDFRef{8338} \MusicPDFRef{8339} \MusicPDFRef{8340}]}
\MusicPDFSetAnnotations{879}{[\MusicPDFRef{8341} \MusicPDFRef{8342} \MusicPDFRef{8343} \MusicPDFRef{8344} \MusicPDFRef{8345} \MusicPDFRef{8346} \MusicPDFRef{8347} \MusicPDFRef{8348} \MusicPDFRef{8349} \MusicPDFRef{8350} \MusicPDFRef{8351} \MusicPDFRef{8352} \MusicPDFRef{8353} \MusicPDFRef{8354} \MusicPDFRef{8355} \MusicPDFRef{8356} \MusicPDFRef{8357} \MusicPDFRef{8358} \MusicPDFRef{8359} \MusicPDFRef{8360}]}
\MusicPDFSetAnnotations{880}{[\MusicPDFRef{8361} \MusicPDFRef{8362} \MusicPDFRef{8363} \MusicPDFRef{8364} \MusicPDFRef{8365} \MusicPDFRef{8366} \MusicPDFRef{8367} \MusicPDFRef{8368} \MusicPDFRef{8369} \MusicPDFRef{8370} \MusicPDFRef{8371} \MusicPDFRef{8372} \MusicPDFRef{8373} \MusicPDFRef{8374} \MusicPDFRef{8375}]}
\MusicPDFSetAnnotations{881}{[\MusicPDFRef{8376} \MusicPDFRef{8377} \MusicPDFRef{8378} \MusicPDFRef{8379} \MusicPDFRef{8380} \MusicPDFRef{8381} \MusicPDFRef{8382} \MusicPDFRef{8383} \MusicPDFRef{8384} \MusicPDFRef{8385} \MusicPDFRef{8386} \MusicPDFRef{8387} \MusicPDFRef{8388} \MusicPDFRef{8389} \MusicPDFRef{8390} \MusicPDFRef{8391} \MusicPDFRef{8392}]}
\MusicPDFSetAnnotations{882}{[\MusicPDFRef{8393} \MusicPDFRef{8394} \MusicPDFRef{8395} \MusicPDFRef{8396} \MusicPDFRef{8397} \MusicPDFRef{8398} \MusicPDFRef{8399} \MusicPDFRef{8400} \MusicPDFRef{8401} \MusicPDFRef{8402} \MusicPDFRef{8403} \MusicPDFRef{8404} \MusicPDFRef{8405} \MusicPDFRef{8406} \MusicPDFRef{8407} \MusicPDFRef{8408} \MusicPDFRef{8409} \MusicPDFRef{8410} \MusicPDFRef{8411} \MusicPDFRef{8412} \MusicPDFRef{8413} \MusicPDFRef{8414}]}
\MusicPDFSetAnnotations{883}{[\MusicPDFRef{8415} \MusicPDFRef{8416} \MusicPDFRef{8417} \MusicPDFRef{8418} \MusicPDFRef{8419} \MusicPDFRef{8420} \MusicPDFRef{8421} \MusicPDFRef{8422} \MusicPDFRef{8423} \MusicPDFRef{8424} \MusicPDFRef{8425}]}
\MusicPDFSetAnnotations{885}{[\MusicPDFRef{8426} \MusicPDFRef{8427} \MusicPDFRef{8428} \MusicPDFRef{8429} \MusicPDFRef{8430} \MusicPDFRef{8431}]}
\MusicPDFSetAnnotations{886}{[\MusicPDFRef{8432} \MusicPDFRef{8433} \MusicPDFRef{8434} \MusicPDFRef{8435} \MusicPDFRef{8436} \MusicPDFRef{8437} \MusicPDFRef{8438} \MusicPDFRef{8439} \MusicPDFRef{8440} \MusicPDFRef{8441} \MusicPDFRef{8442} \MusicPDFRef{8443}]}
\MusicPDFSetAnnotations{887}{[\MusicPDFRef{8444} \MusicPDFRef{8445} \MusicPDFRef{8446} \MusicPDFRef{8447} \MusicPDFRef{8448} \MusicPDFRef{8449} \MusicPDFRef{8450}]}
\MusicPDFSetAnnotations{888}{[\MusicPDFRef{8451}]}
\MusicPDFSetAnnotations{889}{[\MusicPDFRef{8452} \MusicPDFRef{8453} \MusicPDFRef{8454} \MusicPDFRef{8455} \MusicPDFRef{8456} \MusicPDFRef{8457} \MusicPDFRef{8458} \MusicPDFRef{8459} \MusicPDFRef{8460} \MusicPDFRef{8461} \MusicPDFRef{8462} \MusicPDFRef{8463} \MusicPDFRef{8464} \MusicPDFRef{8465}]}
\MusicPDFSetAnnotations{890}{[\MusicPDFRef{8466} \MusicPDFRef{8467} \MusicPDFRef{8468} \MusicPDFRef{8469}]}
\MusicPDFSetAnnotations{891}{[\MusicPDFRef{8470} \MusicPDFRef{8471} \MusicPDFRef{8472}]}
\MusicPDFSetAnnotations{892}{[\MusicPDFRef{8473} \MusicPDFRef{8474} \MusicPDFRef{8475} \MusicPDFRef{8476} \MusicPDFRef{8477}]}
\MusicPDFSetAnnotations{893}{[\MusicPDFRef{8478} \MusicPDFRef{8479} \MusicPDFRef{8480} \MusicPDFRef{8481} \MusicPDFRef{8482} \MusicPDFRef{8483} \MusicPDFRef{8484} \MusicPDFRef{8485}]}
\MusicPDFSetAnnotations{894}{[\MusicPDFRef{8486} \MusicPDFRef{8487} \MusicPDFRef{8488} \MusicPDFRef{8489} \MusicPDFRef{8490} \MusicPDFRef{8491} \MusicPDFRef{8492} \MusicPDFRef{8493} \MusicPDFRef{8494} \MusicPDFRef{8495} \MusicPDFRef{8496} \MusicPDFRef{8497} \MusicPDFRef{8498} \MusicPDFRef{8499} \MusicPDFRef{8500} \MusicPDFRef{8501}]}
\MusicPDFSetAnnotations{896}{[\MusicPDFRef{8502} \MusicPDFRef{8503} \MusicPDFRef{8504}]}
\MusicPDFSetAnnotations{897}{[\MusicPDFRef{8505} \MusicPDFRef{8506} \MusicPDFRef{8507} \MusicPDFRef{8508} \MusicPDFRef{8509} \MusicPDFRef{8510} \MusicPDFRef{8511} \MusicPDFRef{8512} \MusicPDFRef{8513} \MusicPDFRef{8514}]}
\MusicPDFSetAnnotations{898}{[\MusicPDFRef{8515} \MusicPDFRef{8516} \MusicPDFRef{8517} \MusicPDFRef{8518} \MusicPDFRef{8519} \MusicPDFRef{8520} \MusicPDFRef{8521} \MusicPDFRef{8522} \MusicPDFRef{8523} \MusicPDFRef{8524} \MusicPDFRef{8525} \MusicPDFRef{8526} \MusicPDFRef{8527} \MusicPDFRef{8528} \MusicPDFRef{8529} \MusicPDFRef{8530}]}
\MusicPDFSetAnnotations{899}{[\MusicPDFRef{8531} \MusicPDFRef{8532} \MusicPDFRef{8533} \MusicPDFRef{8534} \MusicPDFRef{8535} \MusicPDFRef{8536} \MusicPDFRef{8537}]}
\MusicPDFSetAnnotations{900}{[\MusicPDFRef{8538} \MusicPDFRef{8539} \MusicPDFRef{8540} \MusicPDFRef{8541} \MusicPDFRef{8542} \MusicPDFRef{8543} \MusicPDFRef{8544} \MusicPDFRef{8545} \MusicPDFRef{8546} \MusicPDFRef{8547} \MusicPDFRef{8548} \MusicPDFRef{8549} \MusicPDFRef{8550} \MusicPDFRef{8551} \MusicPDFRef{8552} \MusicPDFRef{8553} \MusicPDFRef{8554} \MusicPDFRef{8555} \MusicPDFRef{8556}]}
\MusicPDFSetAnnotations{901}{[\MusicPDFRef{8557} \MusicPDFRef{8558} \MusicPDFRef{8559} \MusicPDFRef{8560} \MusicPDFRef{8561} \MusicPDFRef{8562} \MusicPDFRef{8563} \MusicPDFRef{8564} \MusicPDFRef{8565} \MusicPDFRef{8566} \MusicPDFRef{8567} \MusicPDFRef{8568} \MusicPDFRef{8569}]}
\MusicPDFSetAnnotations{902}{[\MusicPDFRef{8570} \MusicPDFRef{8571}]}
\MusicPDFSetAnnotations{903}{[\MusicPDFRef{8572} \MusicPDFRef{8573} \MusicPDFRef{8574} \MusicPDFRef{8575} \MusicPDFRef{8576} \MusicPDFRef{8577} \MusicPDFRef{8578} \MusicPDFRef{8579} \MusicPDFRef{8580} \MusicPDFRef{8581} \MusicPDFRef{8582} \MusicPDFRef{8583} \MusicPDFRef{8584}]}
\MusicPDFSetAnnotations{904}{[\MusicPDFRef{8585} \MusicPDFRef{8586} \MusicPDFRef{8587} \MusicPDFRef{8588} \MusicPDFRef{8589} \MusicPDFRef{8590}]}
\MusicPDFSetAnnotations{905}{[\MusicPDFRef{8591} \MusicPDFRef{8592} \MusicPDFRef{8593} \MusicPDFRef{8594} \MusicPDFRef{8595} \MusicPDFRef{8596} \MusicPDFRef{8597} \MusicPDFRef{8598} \MusicPDFRef{8599} \MusicPDFRef{8600} \MusicPDFRef{8601} \MusicPDFRef{8602} \MusicPDFRef{8603} \MusicPDFRef{8604} \MusicPDFRef{8605}]}
\MusicPDFSetAnnotations{906}{[\MusicPDFRef{8606} \MusicPDFRef{8607} \MusicPDFRef{8608} \MusicPDFRef{8609} \MusicPDFRef{8610} \MusicPDFRef{8611} \MusicPDFRef{8612}]}
\MusicPDFSetAnnotations{907}{[\MusicPDFRef{8613} \MusicPDFRef{8614}]}
\MusicPDFSetAnnotations{908}{[\MusicPDFRef{8615} \MusicPDFRef{8616} \MusicPDFRef{8617} \MusicPDFRef{8618} \MusicPDFRef{8619} \MusicPDFRef{8620} \MusicPDFRef{8621} \MusicPDFRef{8622} \MusicPDFRef{8623} \MusicPDFRef{8624}]}
\MusicPDFSetAnnotations{909}{[\MusicPDFRef{8625} \MusicPDFRef{8626} \MusicPDFRef{8627} \MusicPDFRef{8628} \MusicPDFRef{8629} \MusicPDFRef{8630} \MusicPDFRef{8631} \MusicPDFRef{8632} \MusicPDFRef{8633}]}
\MusicPDFSetAnnotations{910}{[\MusicPDFRef{8634} \MusicPDFRef{8635} \MusicPDFRef{8636} \MusicPDFRef{8637}]}
\MusicPDFSetAnnotations{911}{[\MusicPDFRef{8638}]}
\MusicPDFSetAnnotations{912}{[\MusicPDFRef{8639} \MusicPDFRef{8640} \MusicPDFRef{8641} \MusicPDFRef{8642} \MusicPDFRef{8643} \MusicPDFRef{8644}]}
\MusicPDFSetAnnotations{913}{[\MusicPDFRef{8645} \MusicPDFRef{8646} \MusicPDFRef{8647} \MusicPDFRef{8648} \MusicPDFRef{8649} \MusicPDFRef{8650} \MusicPDFRef{8651} \MusicPDFRef{8652} \MusicPDFRef{8653} \MusicPDFRef{8654} \MusicPDFRef{8655} \MusicPDFRef{8656}]}
\MusicPDFSetAnnotations{914}{[\MusicPDFRef{8657} \MusicPDFRef{8658} \MusicPDFRef{8659} \MusicPDFRef{8660} \MusicPDFRef{8661} \MusicPDFRef{8662} \MusicPDFRef{8663} \MusicPDFRef{8664} \MusicPDFRef{8665} \MusicPDFRef{8666} \MusicPDFRef{8667} \MusicPDFRef{8668} \MusicPDFRef{8669}]}
\MusicPDFSetAnnotations{915}{[\MusicPDFRef{8670} \MusicPDFRef{8671} \MusicPDFRef{8672} \MusicPDFRef{8673} \MusicPDFRef{8674} \MusicPDFRef{8675} \MusicPDFRef{8676}]}
\MusicPDFSetAnnotations{916}{[\MusicPDFRef{8677} \MusicPDFRef{8678} \MusicPDFRef{8679} \MusicPDFRef{8680} \MusicPDFRef{8681} \MusicPDFRef{8682} \MusicPDFRef{8683} \MusicPDFRef{8684} \MusicPDFRef{8685} \MusicPDFRef{8686} \MusicPDFRef{8687} \MusicPDFRef{8688} \MusicPDFRef{8689} \MusicPDFRef{8690}]}
\MusicPDFSetAnnotations{917}{[\MusicPDFRef{8691} \MusicPDFRef{8692} \MusicPDFRef{8693} \MusicPDFRef{8694} \MusicPDFRef{8695} \MusicPDFRef{8696} \MusicPDFRef{8697} \MusicPDFRef{8698} \MusicPDFRef{8699} \MusicPDFRef{8700} \MusicPDFRef{8701} \MusicPDFRef{8702}]}
\MusicPDFSetAnnotations{918}{[\MusicPDFRef{8703} \MusicPDFRef{8704} \MusicPDFRef{8705} \MusicPDFRef{8706} \MusicPDFRef{8707} \MusicPDFRef{8708} \MusicPDFRef{8709} \MusicPDFRef{8710} \MusicPDFRef{8711} \MusicPDFRef{8712} \MusicPDFRef{8713}]}
\MusicPDFSetAnnotations{919}{[\MusicPDFRef{8714} \MusicPDFRef{8715} \MusicPDFRef{8716} \MusicPDFRef{8717} \MusicPDFRef{8718} \MusicPDFRef{8719} \MusicPDFRef{8720} \MusicPDFRef{8721} \MusicPDFRef{8722}]}
\MusicPDFSetAnnotations{920}{[\MusicPDFRef{8723} \MusicPDFRef{8724} \MusicPDFRef{8725} \MusicPDFRef{8726} \MusicPDFRef{8727} \MusicPDFRef{8728} \MusicPDFRef{8729} \MusicPDFRef{8730} \MusicPDFRef{8731} \MusicPDFRef{8732}]}
\MusicPDFSetAnnotations{921}{[\MusicPDFRef{8733} \MusicPDFRef{8734} \MusicPDFRef{8735} \MusicPDFRef{8736} \MusicPDFRef{8737} \MusicPDFRef{8738} \MusicPDFRef{8739} \MusicPDFRef{8740} \MusicPDFRef{8741} \MusicPDFRef{8742} \MusicPDFRef{8743} \MusicPDFRef{8744} \MusicPDFRef{8745}]}
\MusicPDFSetAnnotations{922}{[\MusicPDFRef{8746} \MusicPDFRef{8747} \MusicPDFRef{8748} \MusicPDFRef{8749} \MusicPDFRef{8750} \MusicPDFRef{8751} \MusicPDFRef{8752} \MusicPDFRef{8753} \MusicPDFRef{8754} \MusicPDFRef{8755} \MusicPDFRef{8756} \MusicPDFRef{8757}]}
\MusicPDFSetAnnotations{923}{[\MusicPDFRef{8758} \MusicPDFRef{8759} \MusicPDFRef{8760} \MusicPDFRef{8761} \MusicPDFRef{8762} \MusicPDFRef{8763} \MusicPDFRef{8764} \MusicPDFRef{8765} \MusicPDFRef{8766} \MusicPDFRef{8767} \MusicPDFRef{8768} \MusicPDFRef{8769} \MusicPDFRef{8770} \MusicPDFRef{8771} \MusicPDFRef{8772} \MusicPDFRef{8773} \MusicPDFRef{8774} \MusicPDFRef{8775}]}
\MusicPDFSetAnnotations{924}{[\MusicPDFRef{8776} \MusicPDFRef{8777} \MusicPDFRef{8778} \MusicPDFRef{8779} \MusicPDFRef{8780} \MusicPDFRef{8781} \MusicPDFRef{8782} \MusicPDFRef{8783} \MusicPDFRef{8784}]}
\MusicPDFSetAnnotations{925}{[\MusicPDFRef{8785} \MusicPDFRef{8786} \MusicPDFRef{8787} \MusicPDFRef{8788} \MusicPDFRef{8789} \MusicPDFRef{8790}]}
\MusicPDFSetAnnotations{926}{[\MusicPDFRef{8791} \MusicPDFRef{8792} \MusicPDFRef{8793} \MusicPDFRef{8794} \MusicPDFRef{8795} \MusicPDFRef{8796}]}
\MusicPDFSetAnnotations{927}{[\MusicPDFRef{8797} \MusicPDFRef{8798} \MusicPDFRef{8799} \MusicPDFRef{8800} \MusicPDFRef{8801}]}
\MusicPDFSetAnnotations{928}{[\MusicPDFRef{8802} \MusicPDFRef{8803}]}
\MusicPDFSetAnnotations{929}{[\MusicPDFRef{8804} \MusicPDFRef{8805} \MusicPDFRef{8806} \MusicPDFRef{8807} \MusicPDFRef{8808} \MusicPDFRef{8809} \MusicPDFRef{8810}]}
\MusicPDFSetAnnotations{930}{[\MusicPDFRef{8811} \MusicPDFRef{8812}]}
\MusicPDFSetAnnotations{931}{[\MusicPDFRef{8813} \MusicPDFRef{8814} \MusicPDFRef{8815} \MusicPDFRef{8816}]}
\endinput

\let\pdfinfo\editionpdfinfo
\pdfinfo{/Title (The Music of the Primes - Revised Mathematical Text)
 /Subject (Spectral traces, Hecke L-functions and the narrow-strip defect)}
\edef\putpagelabels{\noexpand\pdfcatalog{/PageLabels << /Nums [\pdfpagelabels] >>}}\putpagelabels

\sourcepages{1}{1}{Revised edition \quad$\cdot$\quad Spectral traces and the prime sum \quad$\cdot$\quad Results on page 2}
\leaf{2}{Results}{THE TRACE FORMULA}{The full strip and the missing zeros}
For every unitary Hecke character, every strip and every Sobolev degree
$s\ge0$, the arithmetic quotient counts each interior zero with its full
multiplicity. In degree four it counts a boundary zero of order $m$ with
weight $\min(m,2)$. On the full critical strip its trace equals the
spectral side of the Hecke--Weil explicit formula. Smooth Bruhat tests
give a trace-norm convergent sum over the compact character components.
The multiplicity-capped formula of Chapter~13 fails for an allowed norm
section over $\Q(i)$ at a zero of order at least three.

Write $\lambda=\rho-\tfrac12$ for a nontrivial zero of $\zeta$ in centered
coordinates, $m_\lambda$ for its multiplicity, and
$H(z)=\int_{\R}h(x)e^{zx}\dd x$ for a smooth compactly supported test in
$x=\log t$. The parameter $\alpha$ is the half-width in centered coordinates.
For $h(0)=0$ the prime and real-place distribution $W$ satisfies
\[
 W(h)=H(\tfrac12)+H(-\tfrac12)
       -\Tr(A_hP_\alpha)
       -\sum_{\lambda}\big(m_\lambda-n_{\alpha,4}(\lambda)\big)H(\lambda).
\]
Here $P_\alpha$ is the orthogonal spectral projection defined by the arithmetic
$\theta$-space. The multiplicity $n_{\alpha,4}$ in (4) is zero outside
the closed strip. The last sum includes every missing boundary copy.

\par\smallskip\textbf{Results by manuscript page}\par
\begin{tabular}{@{}p{.24\linewidth}p{.69\linewidth}@{}}
\pagelink{29}{29} & A residue detects every omitted zero.\\[3pt]
\pagelink{426a}{426a} & Limits of projections and the count of boundary jets.\\[3pt]
\pagelink{446}{446--446c}, \pagelink{446d}{446d} & One adelic generator for all strip widths, including zeros on the boundary.\\[3pt]
\pagelink{463}{463--463b} & Trace class and summation over characters.\\[3pt]
\pagelink{476a}{476a}, \pagelink{480a}{480a}, \pagelink{484a}{484a--b} & Truncated prime layers coincide. Before truncation their individual quotients give the prime traces.\\[3pt]
\pagelink{488a}{488a--d}, \pagelink{488e}{488e--f}, \pagelink{489a}{489a--b} & Finite levels exhaust the component even when all primes up to a growing bound are imposed together.\\[3pt]
\pagelink{495}{495--495b} & A counterexample at $2$ and an obstruction to recovering two prime traces with a fixed lift.\\[3pt]
\pagelink{508a}{508a}, \pagelink{518a}{518a} & A valid rational fundamental domain and the lower-cutoff term at the real place.\\[3pt]
\pagelink{514a}{514a--b}, \pagelink{520}{520--520b} & The full-strip formula, the omitted zeros and the consequence of a fixed zero-free half-plane.\\[3pt]
\pagelink{520c}{520c--n}, \pagelink{522i}{522i--m} & The balanced prime sum and a one-sided criterion counting only principal prime ideals over totally imaginary fields.\\[3pt]
\pagelink{522}{522--522e}, \pagelink{522f}{522f--h} & The prime-sum criterion, the Chapter~13 test space and a CM curve for which the cap loses multiplicity.\\[3pt]
\pagelink{522n}{522n--o}, \pagelink{522p}{522p--v} & Square-root fluctuations in the cubic splitting count.\\[3pt]
\pagelink{522w}{522w--y}, \pagelink{522z}{522z--ac}, \mbox{\pagelink{522ad}{522ad--ag}}, \mbox{\pagelink{522ah}{522ah--aj}} & Short prime intervals, cancellation between distinct primes, and a mean-square bound over norm twists.\\
\end{tabular}

\point{Scope} The Hecke prime-sum criterion is an equivalence.
The independent estimate required for the Riemann hypothesis, for its Hecke
generalization, and the status of the multiplicity claim are specified on
\pg{522e}. The lettered folios continue the proofs at the indicated manuscript pages.

\sourcepages{3}{5}{}
\sourcepages{6}{6}{For the status of the mathematical conclusions see pages 2 and 522e.}
\sourcepages{7}{27}{}
\sourcepages{28}{28}{The narrow-strip identity is conditional. Its defect is given on page 520.}
\leaf{29}{Zero detection}{1.4 \quad CONSEQUENCES OF A STRIP TRACE}{A residue for each omitted zero}
Fix $0\le\alpha<\tfrac12$. Write $d_\lambda=m_\lambda-
n_{\alpha,4}(\lambda)$ with the strip multiplicity (4), and
$D_\alpha=\{\lambda:d_\lambda>0\}$. Outside the strip $d_\lambda=m_\lambda$.
Take $h\in C_c^\infty((1,2))$ and set $h_\tau(x)=h(x-\tau)$ for $\tau\ge0$.
The multiplicative tests $h_\tau(\log t)$ stay away from $t=1$. The difference
between the full and the narrow zero sums is
\[
 E_\alpha(h_\tau)=\sum_{\lambda\in D_\alpha}
                 d_\lambda e^{\lambda\tau}H(\lambda).
 \tag{1}
\]
Thus equality of the two geometric sides with their respective zero sums would
give $E_\alpha(h_\tau)=0$.

\point{Residue argument} Repeated integration by parts gives
$|H(\sigma+it)|\ll_{h,A}(1+|t|)^{-A}$ uniformly for $|\sigma|\le\tfrac12$.
The classical zero count, with multiplicities, is $O(T\log(2+T))$ up to
height $T\ge2$. For $\Re z>\tfrac12$ these bounds justify absolute
interchange in
\[
 \int_0^\infty E_\alpha(h_\tau)e^{-z\tau}\dd\tau
   =\sum_{\lambda\in D_\alpha}\frac{d_\lambda H(\lambda)}{z-\lambda}.
 \tag{2}
\]
The series on the right converges normally on compact subsets of the plane
away from its discrete set of possible poles. If (1) vanishes for all $\tau$,
the identity theorem makes (2) identically zero as a meromorphic function.
Given an omitted $\lambda_0$, choose $h(x)=e^{-\lambda_0x}b(x)$ with
$b\in C_c^\infty((1,2))$, $b\ge0$, $\int b=1$. The residue at $\lambda_0$
is $d_{\lambda_0}H(\lambda_0)=d_{\lambda_0}>0$, a contradiction.

\point{Growth suffices} If, for each such $h$ and every $\epsilon>0$,
\[
 E_\alpha(h_\tau)=O_{h,\alpha,\epsilon}
                         (e^{(\alpha+\epsilon)\tau}),
 \tag{3}
\]
the integral in (2) is holomorphic on $\Re z>\alpha$. The same residue
excludes every zero there. The functional equation reflects the conclusion
to the left. Any sequence $\alpha\downarrow0$ suffices for RH.
Page~\pagelink{520}{520} identifies the geometric quantity that would have
to satisfy (3).

For the limiting step tests of \S1.4.2 the positive real-place integrand
at $t\downarrow1$ is comparable in the limit to $(t-1)^{-1}$ for real $z$.
Fatou's lemma precludes the stated uniform bound. The prime Dirichlet series
there equals $-\zeta'/\zeta(z+\tfrac12)$ in its half-plane of convergence.

\sourcepages{30}{32}{Zero location in this discussion requires the estimate on page 520.}
\sourcepages{33}{81}{}
\sourcepages{82}{82}{The global Fourier transform in (2.2.3.2) uses the additive character on page 558a.}
\sourcepages{83}{123}{}
\sourcepages{124}{124}{For the character normalization in (3.2.2.2) see page 558a.}
\sourcepages{125}{350}{}
\sourcepages{351}{351}{With extra unit conditions the value at zero vanishes automatically. See page 488c.}
\sourcepages{352}{354}{}
\sourcepages{355}{355}{With extra unit conditions the value at zero vanishes automatically. See page 488c.}
\sourcepages{356}{425}{}
\sourcepages{426}{426}{Theorem 9.3.2.6 is replaced by the projection limit on page 426a.}
\leaf{426a}{Zero conditions}{9.3.2 \quad ORTHOGONAL PROJECTIONS}{The direction of the cutoff}
For continuous zero-jet functionals $L_{\lambda,j}$ on a Hilbert space $X$,
the subspace in (9.3.2.10) has the form
\[
 K_T=\bigcap_{|\Im\lambda|\le T}\ \bigcap_j\ker L_{\lambda,j}.
\]
For $T'>T$ one has $K_{T'}\subset K_T$. The orthogonal projections onto
$K_T$ converge strongly to the projection onto $K_\infty=\bigcap_TK_T$.
Indeed the orthogonal complements increase and their union spans a dense
subspace of $K_\infty^\perp$. Approximation by this union gives strong
convergence there, while every $K_T$ contains $K_\infty$.

\point{An instance of (9.3.2.10)} On the logarithmic line let
$g(x)=(4\pi)^{-1/2}e^{-x^2/4}$ and take $\varphi=-g'$. Its bilateral
Laplace transform is $ze^{z^2}$ and its only zero is $0$, simple.
On every positive-width strip centered at zero the evaluation
$L_0(f)=\int f(x)\dd x$ is continuous. For every $T>0$, $K_T=\ker L_0$
and $L_0(g)=1$. Consequently the projections in Theorem 9.3.2.6 do not
approach the identity. This is the theorem cited in the proof on page 446.

\point{Boundary multiplicity in (10.2.2.5)} For Sobolev degree $s\ge0$,
the dual norm squared of $x^j e^{\lambda x}$ at $|\Re\lambda|=\alpha$
has a tail comparable to $\int_1^\infty x^{2j-s}\dd x$ on the boundary
side. Integrability requires
\[
 j<\frac{s-1}{2}.
\]
The number of admissible jets at a boundary zero of multiplicity $m$ is
$\min\{m,\max(0,\lceil(s-1)/2\rceil)\}$. For $s=1$, $\alpha=0$, the
zeroth jet in the printed formula already has infinite norm. In the
interior, exponential decay permits all $0\le j<m$.

The construction on \pg{446} proves the required density directly in
$X$. The growing spaces used on \pg{463} are spans of the Riesz
representatives of the zero jets.

\sourcepages{427}{427}{}
\sourcepages{428}{428}{The hard cutoffs have a kernel consisting of functions supported on the discarded set.}
\sourcepages{429}{445}{}
\leaf{446}{The arithmetic quotient}{10.2.2 \quad THE STRIP AND ITS BOUNDARY}{The zero conditions at every width}
Use the log coordinate $x=\log|g|$, center at $1/2$, and take any
half-width $\alpha\ge0$ and Sobolev degree $s\ge0$. Put
\[
 w(x)=e^{2\alpha|x|}(1+x^2)^{s/2},\qquad
 X=L^2(w\dd x),\qquad Y=L^2(w^{-1}\dd x).
\]
For a nontrivial zero $\rho$ let $\lambda=\rho-1/2$ and let $m_\lambda$
be its order. Define $b_s=\max(0,\lceil(s-1)/2\rceil)$ and
\[
 n_{\alpha,s}(\lambda)=
 \begin{cases}
 m_\lambda,&|\Re\lambda|<\alpha,\\
 \min(m_\lambda,b_s),&|\Re\lambda|=\alpha,\\
 0,&|\Re\lambda|>\alpha .
 \end{cases} \tag{4}
\]
In the rational trivial character component let $J_\alpha$ be the
closure in $X$ of the centered arithmetic $\theta$-functions. Then
\[
 J_\alpha=\left\{f\in X:\int_\R x^j e^{\lambda x}f(x)\dd x=0
               \quad(0\le j<n_{\alpha,s}(\lambda))\right\}.
 \tag{5}
\]
All the integrals in (5) are continuous on $X$. At a boundary point
the dual norm at the unprotected end is finite exactly when $2j-s<-1$,
as on \pg{426a}. When $\alpha=0$ this calculation applies at both ends.

\point{An arithmetic generator}
Take $D=t\partial_t$, $f_\infty=D(D+1)e^{-\pi t^2}$, and the finite
tests $\mathbf1_{\Z_p}$. Both moments vanish. Poisson summation gives
decay faster than every exponential at both ends for
$\Theta(x)=e^{x/2}\sum_{n\ne0}f_\infty(ne^x)$. With $u=z+1/2$,
unfolding and integration by parts give its entire transform
\[
 \mathcal M\Theta(z)=u(u-1)\pi^{-u/2}\Gamma(u/2)\zeta(u)=:\Xi(z).
 \tag{6}
\]
The zeros of $\Xi$ are precisely the indicated $\lambda$ with their
orders. Every arithmetic test has Mellin transform divisible by
$\zeta(u)$ on $0<\Re u<1$. The unit average at each exceptional
finite prime is eventually constant towards zero, so removing its
Euler factor leaves a Laurent polynomial in $p^{-u}$. The real
integral is holomorphic there. Thus every arithmetic test satisfies
(5). Pages \pagelink{446a}{446a--b} prove that the translates of the
single vector $\Theta$ already span the entire right side after closure.

\leaf{446a}{The arithmetic quotient}{10.2.2 \quad FINITE PRODUCTS OF ZERO FACTORS}{The range of a polynomial in differentiation}
Factor the order-one entire function in (6) as
\[
 \Xi(z)=e^{A+Bz}z^{m_0}\prod_{\lambda\ne0}
       \big[(1-z/\lambda)e^{z/\lambda}\big]^{m_\lambda}.
\]
Here $m_0=0$ if the origin is not a zero. Let $p_R(z)$ be $z^{m_0}$
times the product of the polynomial factors for $0<|\lambda|\le R$,
and let $T_R(z)$ be the product of the complete genus-one factors for
$|\lambda|>R$. Thus $\Xi=e^{A+B_Rz}p_R T_R$, where
$B_R=B+\sum_{0<|\lambda|\le R}m_\lambda/\lambda$.

For each finite $R$ the $X$-closure of
$p_R(-\partial_x)C_c^\infty(\R)$ is precisely the kernel of the
conditions in (5) at the zeros of $p_R$. Indeed, with the bilinear
pairing $\int f u$, its annihilator consists of $u\in Y$ satisfying
$p_R(\partial_x)u=0$ as distributions. All distributional solutions
are finite sums of $x^je^{\lambda x}$ with $j<m_\lambda$.
This follows by splitting the polynomial into coprime primary factors
and using $\partial_x^m v=0$ exactly for polynomials of degree $<m$.
For each solution in $Y$ its component at each distinct $\lambda$ also
lies in $Y$: choose a real shift $a$ for which the finitely many
$e^{a\lambda}$ are distinct. The primary projections are polynomials in
the bounded translation $u(x)\mapsto u(x+a)$. The leading coefficient
of each remaining polynomial times $e^{\lambda x}$ determines
integrability at infinity. Equation (4) is precisely the resulting
count. Double annihilation gives the range assertion, including zeros
on either boundary and $s=1$.

For every $R$, writing $r=2+|z|$, the tail satisfies
\[
          |T_R(z)|\le \exp(Cr\log^2 r),\qquad T_R\longrightarrow1
          \ \hbox{locally uniformly as }R\to\infty .
 \tag{7}
\]
For $|\lambda|\le 2r$ use $\log^+|1-z/\lambda|\le C\log r$
and $|\Re(z/\lambda)|\le r/|\lambda|$. The zero count
$N(u)\ll u\log(2+u)$ gives $O(r\log^2 r)$ for this part by partial
summation. For $|\lambda|>2r$ the logarithm of a factor is
$O(r^2/|\lambda|^2)$, with sum $O(r\log r)$. Both estimates hold for
every tail. Absolute convergence on compact sets gives the local limit.

\leaf{446b}{The arithmetic quotient}{10.2.2 \quad GAUSSIAN CONVOLUTION}{Passing from the finite factors to the arithmetic space}
For $\delta>0$ let $k_{R,\delta}$ be the inverse Fourier transform on
$z=it$ of $e^{\delta z^2}T_R(z)$ in the convention
$\mathcal M f(z)=\int f(x)e^{zx}\dd x$. Let $k_\delta(x)
=(4\pi\delta)^{-1/2}e^{-x^2/(4\delta)}$. For fixed $\delta$,
\[
          \int_\R |k_{R,\delta}(x)-k_\delta(x)|\sqrt{w(x)}\dd x
                   \longrightarrow0 .
 \tag{8}
\]
In the inverse integral move the contour to
$\Re z=\operatorname{sgn}(x)(\alpha+1)$. The Gaussian dominates (7)
on the horizontal segments. On the two vertical lines dominated convergence
makes the $L^1(\dd t)$ norms of $e^{\delta z^2}(T_R-1)$ tend to zero.
The inverse is consequently $o(1)e^{-(\alpha+1)|x|}$ uniformly in $x$.
This also shows that convolution by these kernels is bounded on $X$,
since $\sqrt{w(x+y)}\le 2^{s/4}\sqrt{w(x)w(y)}$.

Write $J(\Theta)$ for the closed span of the translates in $X$. For
$g\in C_c^\infty(\R)$ the transform of
$k_{R,\delta}*p_R(-\partial_x)g$ is exactly
\[
      \Xi(z)\, e^{\delta z^2-A-B_Rz}\,\mathcal M g(z).
\]
The inverse of the middle factor is
$e^{-A}(4\pi\delta)^{-1/2}\exp(-(x+B_R)^2/(4\delta))$.
For each fixed $R$ it is integrable against $\sqrt w$, even if $B_R$
is complex. Thus the vector is $\Theta$ convolved with a kernel in
$L^1(\sqrt w\dd x)$, and belongs to $J(\Theta)$ by the norm-continuous
integral of translates. If $f$ satisfies (5), the finite range result
on \pg{446a} gives $k_{R,\delta}*f\in J(\Theta)$ for every $R$.
By (8), $k_\delta*f\in J(\Theta)$. Finally $k_\delta*f\to f$ in
$X$ as $\delta\downarrow0$, by strong continuity of translations and
the Gaussian bound on their weighted tail. Every translate of $\Theta$
is an arithmetic test up to a scalar. This proves (5).

\point{The dual of the quotient} Under the bilinear pairing the complete
annihilator, for every width and degree above, is
\[
 (X/J_\alpha)'=\overline{\Span}^{\,Y}
   \{x^je^{\lambda x}:\ 0\le j<n_{\alpha,s}(\lambda)\}=:E_\alpha .
 \tag{9}
\]
The Riesz map $u\mapsto w^{-1}\overline u$ takes this span onto the
orthogonal spectral subspace in $X$.

\leaf{446c}{Hecke characters}{10.2.2 \quad A FIXED CHARACTER OF A NUMBER FIELD}{A generator in each character space}
Let $F$ be a number field and choose a positive archimedean section of the
idele norm. Normalize a unitary Hecke character $\chi$ to be trivial on the
section, which fixes the vertical translation due to norm twists. Use the
global additive character of \pg{558a}. Project the centered
$\theta$-sum to its character component by integration against $\chi$ on
$C_F^1$ with Haar probability. Its transform unfolds to Tate's integral
$Z(f,\chi,s)$, $s=\tfrac12+z$. The Hilbert norm in this component is again
$X$ on the norm coordinate.

\point{A test producing the completed function} At unramified finite places
take $\mathbf1_{\mathcal O_v}$ and at ramified places take
$\overline{\chi_v}\mathbf1_{\mathcal O_v^\times}$. The local integrals give
the finite Euler factors up to constants. At a real place take
$t^\kappa e^{-\pi t^2}$ for sign exponent $\kappa\in\{0,1\}$. At a complex
place with angular exponent $k$ take the homogeneous monomial of degree
$|k|$ and angular exponent $-k$ times $e^{-2\pi|t|^2}$ (ordinary complex
modulus). Radial integration gives the respective real and complex gamma
factors, including the imaginary shifts of $\chi_v$. The product test
$f_0$ thus has integral $c e^{bz}L^*(\tfrac12+z,\chi)$, where $c\ne0$,
and $L^*$ includes the conductor and gamma factors.

If $\chi|_{C_F^1}\ne1$, choose a norm-one idele $a$ with $\chi(a)\ne1$
and put $f(y)=f_0(y)-f_0(ay)$. Then $f(0)=\widehat f(0)=0$ and its
Mellin integral is multiplied by $1-\chi(a)^{-1}\ne0$. For $\chi=1$
apply $D(D+1)$ to $f_0$, where $D$ differentiates archimedean dilation
with modulus $e^x$ at $x=0$. This kills both moments and multiplies the
integral by $s(s-1)$. The resulting transform $\Xi_\chi$ is entire of
order at most one with precisely the nontrivial zeros, counted with
multiplicity. Poisson summation gives a kernel of rapid decay at both ends.

\point{Identification} Every arithmetic Tate integral in this component
vanishes at those zeros to their full orders. After unfolding, the quotient
by each local $L$-factor is holomorphic in $0<\Re s<1$, and almost all
local quotients are $1$. Finite sums of product tests suffice. The
fixed-character zero count $O_{F,\chi}(T\log(2+T))$ and the Hadamard
product give (7), including a central zero as a separate polynomial
factor. Pages 446a--446b therefore prove (5) and (9) with $\Xi_\chi$
at every width and degree. At $\alpha=1/2$ the Hecke nonvanishing
theorem on $\Re s=1$ and reflection exclude boundary zeros. The full
critical strip captures every nontrivial zero with its actual multiplicity.

\leaf{446d}{The critical line}{10.2.2 \quad ZERO WIDTH IN A SINGLE CHARACTER}{The arithmetic quotient on the line itself}
Fix $\chi$ and center at $1/2$. At half-width zero and even Sobolev
degree $2k$, $k\ge1$, the space is
$X_k=L^2(\R,(1+x^2)^k\dd x)$. Let $J$ be the closure of the centered
arithmetic tests and $S$ its orthogonal-complement projection in $X_k$.
With $\rho=\tfrac12+i\gamma$ running over distinct critical-line zeros
of the completed function, one has
\[
  \Tr(A_h S)=\Tr(S A_h)
       =\sum_{\rho=\frac12+i\gamma}\min(m_\rho,k)\,H(i\gamma).
       \tag{9a}
\]
This holds for $h\in C_c^\infty(\R)$ and includes any multiple zero
at $\gamma=0$.

\point{The annihilator} A functional on $X_k$ is represented in the
bilinear integral pairing by $u\in L^2((1+x^2)^{-k}\dd x)$.
Its Fourier transform, in the convention $e^{-itx}$, lies in the
distributional Sobolev space $H^{-k}(\R_t)$ by Plancherel. If $u$
annihilates all translates of the generator on \pg{446c}, then
$\Xi_\chi(it)\widehat u(t)=0$. This is the Fourier transform of the
vanishing convolution with the reflected generator. Away from the
isolated real-frequency zeros this smooth multiplier is locally
invertible. Near a zero $\gamma$ of order $m$ it is $(t-\gamma)^m$
times a smooth invertible factor. A distribution killed by that power
is a linear combination of $\delta_\gamma^{(j)}$ with $j<m$, by Taylor
expansion of a test at $\gamma$. Membership in $H^{-k}$ also requires
$j<k-\tfrac12$, hence $j\le k-1$, as follows by Fourier inversion.

Multiplying $\widehat u$ by smooth cutoffs $\eta(t/R)$, equal to one
near zero with compact support, gives finite combinations of these
jets converging to $\widehat u$ in $H^{-k}$. Indeed the multipliers
are bounded uniformly on $H^k$ by the product rule, hence on $H^{-k}$
by duality, and converge strongly there by density of Schwartz functions.
Conversely all the displayed jets annihilate every arithmetic test by
its Hecke zero multiplicity. Thus the dual of $X_k/J$ is exactly the
closed span of $x^j e^{i\gamma x}$ for $0\le j<\min(m_\rho,k)$.

For these chains the proof on pages 463--463a applies with
$q=(1+x^2)^{-k}$. Both $q'/q$ and $q''/q$ are bounded and the boundary
atom is zero. The zero count bounds (10), so both products in (9a) are
trace class. Increasing $k$ gives the full multiplicities on the line
by dominated convergence against $\sum_\rho m_\rho |H(i\gamma)|$.

\sourcepages{447}{462}{For the arithmetic zero space and the trace see pages 446--446d and 463--463b.}
\leaf{463}{The spectral trace}{10.4.2 \quad CONVOLUTION AFTER PROJECTION}{A resolvent estimate without zero spacing}
Let $q=w^{-1}=e^{-2\alpha|x|}(1+x^2)^{-s/2}$ and $E=E_\alpha\subset Y$
as on \pg{446b}. In a fixed character component let $E_N$ contain all the
chains of the zeros with $|\Im\lambda|\le N$ in (9). The derivative
$D_N=\partial_x|_{E_N}$ preserves this space. For $u\in E_N$ integration
by parts on the two half-lines gives
\[
 \Re\langle D_Nu,u\rangle_Y=-\tfrac12\int |u|^2q'\dd x
          \le C\|u\|_Y^2,\qquad C=\alpha+s/4.
\]
The two terms at $0$ cancel by continuity of $q$. Choose $c>C$, put
$R_N=(c-D_N)^{-1}$, and substitute $u=R_Ne$ for each vector of an
orthonormal basis. Addition and the finite-dimensional trace give
\[
 (c-C)\|R_N\|_{\rm HS}^2\le\Re\Tr R_N
   =\sum_{\lambda\in E_N^{\rm sp}}n_{\alpha,s}(\lambda)
                  \frac{c-\Re\lambda}{|c-\lambda|^2}.
 \tag{10}
\]
Here $E_N^{\rm sp}$ denotes the finite set of distinct eigenvalues.
Exponential polynomials for distinct eigenvalues are independent, and each
chain gives one triangular block. The multiplicity is (4), and the fractions
are $O_c((1+|\Im\lambda|)^{-2})$, uniformly in $N$. Zero counting bounds
the sum. The resolvents agree on nested spaces and have norm at most
$(c-C)^{-1}$. Taking an orthonormal basis adapted to these spaces in (10)
gives a Hilbert--Schmidt operator $R$ on $E$. In particular $R^2$ is
trace class.

\point{The weight at zero} Write $A_h f(x)=\int h(x-t)f(t)\dd t$ for
$h\in C_c^\infty(\R)$. The isometry $u\mapsto qu$ sends
$\overline E=\{\overline u:u\in E\}$ onto $P_\alpha X$.
Use the conjugate resolvent on $\overline E$, also denoted $R$ ($c$ is
real). Write $v=R^2u$ on a finite span. Since $u=(c-\partial_t)^2v$,
two integrations by parts give
\[
 A_h(qu)= A_{(c-\partial)^2h}(qv)
       +2A_{(c-\partial)h}(q'v)+A_h(q''_{\rm reg}v)
       -4\alpha\,h\,v(0).
 \tag{11}
\]
The distributional second derivative has atom
$[q'(0+)-q'(0-)]\delta_0=-4\alpha\delta_0$. Here $q''_{\rm reg}$ is its
ordinary derivative of order two on the open half-lines.

\leaf{463a}{The spectral trace}{10.4.2 \quad THE TWO ORDERS OF COMPOSITION}{Trace class and cutoff limits}
Both $q'/q$ and $q''_{\rm reg}/q$ are bounded. Convolution with a compactly
supported smooth kernel is bounded on $X$ by the translation estimate
on page 446b. Thus each integral term of (11) is a bounded map $Y\to X$
applied to $R^2u$. For the last term,
$v'=cv-Ru$ and $q$ is bounded away from zero on $[-1,1]$. The ordinary
$H^1([-1,1])$ point-evaluation bound yields
\[
 |(R^2u)(0)|\ll\|R^2u\|_Y+\|Ru\|_Y\ll\|u\|_Y .
\]
It therefore gives a bounded rank-one map (zero if $\alpha=0$).
By density, (11) proves that $A_hP_\alpha$ is trace class on $X$.
The $\theta$-space $J_\alpha$ is invariant under convolution, so
$P_\alpha A_h=P_\alpha A_hP_\alpha$ is trace class as well.

\point{The trace} Let $P_N$ project orthogonally in $X$ onto
$q\overline{E_N}$. Under the bilinear pairing of page 446b the transpose
of convolution acts by $B_hu(t)=\int h(y)u(t+y)\dd y$. On each chain its
matrix is triangular with diagonal $H(\lambda)$, where
$H(z)=\int h(y)e^{zy}\dd y$. The restriction of $B_h$ to $E_N$ is the
transpose of the compression $P_N A_h P_N$, since the pairing between
$P_N X$ and $E_N$ is nondegenerate and annihilates $(1-P_N)X$.
As $P_N\to P_\alpha$ strongly, the compressions converge in trace norm to
$P_\alpha A_hP_\alpha$: multiply the trace-class operator $A_hP_\alpha$
on both sides by the bounded projections and use finite-rank approximation.
The off-diagonal block of $A_hP_\alpha$ has zero trace. Hence
\[
 \Tr(A_hP_\alpha)=\Tr(P_\alpha A_h)
       =\sum_\lambda n_{\alpha,s}(\lambda) H(\lambda)=:S_\alpha(h).
 \tag{12}
\]
The zero sum converges absolutely by integration by parts and zero counting.
For the manuscript cutoffs $0\le r_L\le1$ tending pointwise to $1$,
$M_{r_L}\to1$ strongly on $X$. Finite-rank approximation gives
\[
 \lim_{L\to\infty}\Tr(M_{r_L} A_hP_\alpha)
 =\lim_{L\to\infty}\Tr(M_{r_L} P_\alpha A_h)=S_\alpha(h).
 \tag{13}
\]
These are the two orders of composition in (10.4.2.9) and (10.4.2.20).
The argument applies to each fixed character and each width of (5).

\leaf{463b}{The spectral trace}{10.4.2 \quad THE COMPACT CHARACTER COMPONENTS}{Summing the character traces}
Fix $F$, write $K=C_F^1$, and take any common $\alpha\ge0$ and $s\ge0$.
Equation (9) holds for all the components at this same width, including
their boundary zeros. In each component extend the character trivially
on the same positive norm section.

\point{Uniform resolvent bound} Choose $c>\max(\alpha+s/4,1)$ and put
$s_*=\tfrac12+c$. Let $\xi_\chi(s)$ be the standard conductor-completed
Hecke function, multiplied by $s(s-1)$ for the trivial character.
Its reflection symmetry and its Hadamard product give
\[
 \sum_{\lambda}m_\lambda\Re\frac1{c-\lambda}
       =\Re\frac{\xi_\chi'}{\xi_\chi}(s_*)
       \ll_{F,c}1+\log N\mathfrak f_\chi+
             \sum_{v\mid\infty}\log(2+|t_v|+|k_v|).
 \tag{13a}
\]
Here $t_v,k_v$ are the imaginary norm and sign or angular parameters.
For the equality, the real logarithmic derivative in centered coordinates
and the real zero sum differ by a constant in the genus-one product.
On the imaginary axis away from the zeros the former vanishes by the
functional equation. The latter vanishes by pairing $\lambda$ with
$-\bar\lambda$ and is absolutely convergent. This fixes the constant,
including when a central zero is present. For the bound use the Euler
product at the fixed $s_*>1$, which gives
$|L_f'/L_f(s_*,\chi)|\le-\zeta_F'/\zeta_F(s_*)$, and the logarithmic
derivative of the conductor and gamma factors. Stirling's formula there
gives the displayed logarithms. Each zero term on the left is positive.

\point{Absolute trace norm} A Bruhat test $\phi$ factors on $K$ through
a compact abelian Lie quotient at a fixed finite level. Its component
kernels $h_\chi$ have one compact log support. Their first two derivatives
in the log coordinate are $O_{A,\phi}((1+|\nu|)^{-A})$ for every $A$,
by integration by parts in the torus variables. Here $\nu$ ranges over
a lattice with finitely many components. The finite conductors are bounded
at this level, and the archimedean parameters are $O(1+|\nu|)$ since
restriction to the archimedean Lie algebra is linear in the torus weight.
By (10) and (13a), $\|R_\chi^2\|_1\ll_{F,c,\phi}1+\log(2+|\nu|)$.
The bounded factors in (11) cost at most a support-dependent constant
times $\max_{j\le2}\|h_\chi^{(j)}\|_\infty$. The point evaluation after
$R_\chi^2$ is uniformly bounded by the operator-norm estimate on page 463.
Thus the norms $\|A_{h_\chi}P_{\alpha,\chi}\|_1$ are summable.
The direct sum has trace (12) summed over $\chi$, with the same cutoff
limit. Absolute convergence of the double zero sum follows from two
integrations by parts in $H_\chi$ and (13a).

\sourcepages{464}{474}{For the trace-class projection and its cutoff limit see pages 463--463b.}
\sourcepages{475}{475}{}
\sourcepages{476}{476}{The convergence assertion in Proposition 11.1.1.2 requires control of orbit tails. See page 476a.}
\leaf{476a}{Convergence of orbital sums}{11.1.1 \quad ORBITAL SUMMATION}{Mass can leave every finite part of an orbit}
The aggregate bound (11.1.1.2) permits a sequence of orbital sums that
are identically one although their summands converge pointwise to zero.
Take $F=\Q$, $Z=\{1\}$ in the ideles and $\eta=\delta_1$. Let
$b\in C_c^\infty((3/4,5/4))$ be nonnegative with $b(1)=1$. Set
$f=b\otimes\mathbf1_{\widehat{\Z}}$. For distinct increasing rational
primes $q_j$ set $f_j(y)=f(q_j^{-1}y)$, where $q_j$ is diagonal.
These are adelic Schwartz--Bruhat functions. For each fixed idele $y$,
its $q_j$-component is a unit for all sufficiently large $j$. Then
$q_j^{-1}y_{q_j}\notin\Z_{q_j}$, so $f_j(y)\to0$.

On the rational orbit one has exactly
\[
       f_j(\xi)=\mathbf1_{\{\xi=q_j\}},\qquad
       \sum_{\xi\in\Q^\times} |f_j(\xi)|=1 .
\]
Indeed $\xi/q_j$ must be an integer in the support of $b$. The hypotheses
of Proposition 11.1.1.2 hold with $I_Z=\Z$ and
$\alpha(\xi)=\mathbf1_{\{\xi=1\}}$. Its orbital majorant is one in
$L^r(Z,\eta)$ for every finite $r\ge1$. The orbital sum of the limit is
zero, so (11.1.1.4) fails even on this one-point measure space.

\point{A sufficient hypothesis} For almost every $g\in Z$ suppose that
$f_j(\xi g)\to f(\xi g)$ for every $\xi$, and that
$|f_j(\xi g)|\le b_g(\xi)$ with $\sum_\xi b_g(\xi)$ dominated by an
$L^r(Z,\eta)$ function. Dominated convergence on the counting measure
gives convergence of the orbital sums at $g$. Applying it on $Z$ to
the $r$th power of their difference gives (11.1.1.4).

For the rational unit restrictions in \S11.1.2, with the standard
unramified finite test and $t=|g|\ge d>0$, every surviving summand is
bounded by $C_M(1+|n|t)^{-M}$, $n\in\Z\setminus\{0\}$.
The sum is $O_{d,M}((1+t)^{-M})$ for $M>1$. It is square integrable
against $(1+\log^2t)^{s/2}\dd t/t$ on this half-line. This supplies
the required domination in the arithmetic application over $\Q$.

\sourcepages{477}{479}{For the orbital convergence hypothesis see page 476a.}
\sourcepages{480}{480}{The truncated spaces in (11.1.2.6) are equal for every cutoff. See page 480a.}
\leaf{480a}{Prime layers}{11.1.2 \quad THE SPACES AFTER SPATIAL TRUNCATION}{Every finite layer has the same closure}
On the uncentered zero line put $x=\log |g|$ and
$H_s=L^2(\R,(1+x^2)^{s/2}\dd x)$, $s\ge0$. In a fixed character
component the operation at an unramified place is $B=1-aU$,
$Uu(x)=u(x+l)$, $l=\log N\mathfrak p>0$, $|a|=1$.
This follows from (5.3.1.1), or from the multiplier $1-ae^{-lz}$.
Let $J$ be any one of the preceding translation-invariant arithmetic
layers. For every bounded cutoff $\rho$ supported on $x\ge-L$,
\[
            \overline{M_\rho B J}=\overline{M_\rho J}.
            \tag{13b}
\]
Indeed for $u\in J$ the sum $\sum_{j=0}^{n-1}a^jU^ju$ lies in $J$.
Applying $M_\rho B$ gives $M_\rho u-a^nM_\rho U^nu$.
The last term tends to zero because its squared norm is bounded by
\[
 C_{\rho,L,s}\int_{nl-L}^{\infty}|u(y)|^2(1+y^2)^{s/2}\dd y.
\]
Here $1+x^2\ll_L1+(x+nl)^2$ on $x\ge-L$. The reverse inclusion follows
from $BJ\subset J$. Iteration gives (13b) for every finite collection
of new places. It applies to both $\rho_{\Lambda,+}$ and $\rho_\Lambda$
of (9.3.3.4)--(9.3.3.5), for every $\Lambda$.

\point{Exact realization over $\Q$} For a finite set $T$ of primes let
$P_T=\prod_{p\in T}p$. In the trivial component its orbit sum for an
even real-place test is $\Theta_T f(t)=2\sum_{(n,P_T)=1,\ n\ge1} f(nt)$.
Take any $\psi\in C_c^\infty((d,\infty))$, where $d>0$, and choose
$b>\sup\supp\psi$. On the positive axis put
\[
       f(u)=\tfrac12\sum_{\substack{m\le b/d\\(m,P_T)=1}}
                                     \mu(m)\psi(mu).
\]
Extend evenly and adjust its integral to zero by an even bump supported
in $0<|u|<d$. Then $f(0)=\widehat f(0)=0$. For $t\ge d$ a nonzero term
with index $k=mn$ satisfies $k\le b/d$. Its coefficient is
$\sum_{m\mid k}\mu(m)=\mathbf1_{\{k=1\}}$ for $(k,P_T)=1$.
Thus $\Theta_T f(t)=\psi(t)$ throughout the observed half-line.
In particular the common closure there is the entire supported $L^2$
space. The strictness conclusion for (11.1.2.6) fails for every $\Lambda$.

\sourcepages{481}{483}{The distinction between a full layer and its cutoff image is proved on page 480a.}
\sourcepages{484}{484}{The reciprocal Euler operator has nonclosed Hilbert range. Its quotient trace is on pages 484a--b.}
\leaf{484a}{The local Hilbert quotient}{11.2 \quad RECIPROCAL EULER FACTORS ON THE ZERO LINE}{The closure in the local quotient}
On the uncentered line the arithmetic closure in a fixed unitary Hecke
component is $H_s$ of \pg{480a}. In fact the uncentered generator
obtained from \pg{446c} has Fourier transform $\Xi_\chi(-\tfrac12+it)$,
up to the sign convention for $t$. It has no real-frequency zero by
the Hecke nonvanishing theorem at $\Re s=1$ and the functional
equation, with the poles of a norm character removed as on \pg{446c}.
The generator is Schwartz on the log line. A continuous functional on
$H_s$ gives a tempered distribution. If it annihilates all these
translates, its Fourier transform times the reflected Fourier transform
of the generator vanishes. Division on each compact frequency interval
shows that the distribution vanishes. This proves density
(Ramakrishnan and Valenza, \emph{Fourier Analysis on Number Fields},
Theorem 7-28, for the nonvanishing input).

For one prime put $U_l f(x)=f(x+l)$ and $B=1-aU_l$, $|a|=1$.
The only solution of $Bf=0$ in $H_s$ is zero, since $|f|$ would be
$l$-periodic. Choose $\gamma$ with $ae^{i\gamma l}=1$ and a nonzero
$b\in C_c^\infty((1,2))$. For $R\to\infty$ set
\[
          f_R(x)=R^{-(s+1)/2}e^{i\gamma x}b(x/R).
\]
Changing variables gives $\|f_R\|_{H_s}\asymp1$. The mean value theorem
gives $\|Bf_R\|_{H_s}=O_{b,l,s}(R^{-1})$. An injective bounded map
with closed range is bounded below. Thus (11.2.0.3) fails already for
one unramified prime and $s=4$.

\point{The Hausdorff quotient} For a rational prime $p$ and a character
unramified at the places $v\mid p$, write $l=\log p$,
$l_v=f_v l=\log Nv$ and $a_v=\chi_v(\varpi_v)$. Put
\[
 Q(U_l)=\prod_{v\mid p}(1-a_vU_l^{f_v}),\qquad
 \mathcal E_p=H_s/\overline{Q(U_l)H_s}.
 \tag{13c}
\]
The former quotient in (11.2.2.1) acquires this Hilbert realization
by dividing by the closure of zero in its quotient topology.
Under the bilinear integral pairing the dual of $\mathcal E_p$ is the
kernel of $Q(U_{-l})$ in $H_{-s}=L^2((1+x^2)^{-s/2}\dd x)$.
The translation recurrence on that kernel gives its trace explicitly.

\leaf{484b}{The local Hilbert quotient}{11.2.2 \quad THE PLACES OVER ONE RATIONAL PRIME}{The trace on the untruncated quotient}
Let $D=\sum_{v\mid p}f_v\le[F:\Q]$ and take $s\ge2[F:\Q]$.
For each distinct solution $\beta$ of $\beta^{f_v}=a_v$ let
$r_\beta$ be the number of factors for which it is a solution.
For almost every $t\in[0,l)$, the sequence $u(t+kl)$, $k\in\Z$, of a
dual vector in (13c) solves a constant-coefficient recurrence of order
$D$. Its solutions are exactly
\[
   u(t+kl)=\sum_\beta\sum_{j=0}^{r_\beta-1}
                              c_{\beta j}(t)k^j\beta^k .
\]
The coefficients are in $L^2([0,l])$. Indeed the first $D$ samples
determine them by an invertible confluent Vandermonde matrix. Conversely
$\sum_k(1+|k|)^{2(D-1)-s}<\infty$, so these coefficient norms give
an equivalent norm on the kernel. Replacing $k^j$ by $(t+kl)^j$ makes
a bounded invertible triangular change of coefficients. We may therefore
write $u(x)=\sum_{\beta,j}x^j u_{\beta j}(x)$ with
$u_{\beta j}(x+l)=\beta u_{\beta j}(x)$.

On this finite sum of circle spaces the transpose of convolution by
$h\in C_c^\infty(\R)$ is $u(x)\mapsto\int h(y)u(x+y)\dd y$.
Expanding $(x+y)^j$ gives a triangular operator in each chain, with
diagonal convolution $C_{\beta,h}$ on the $\beta$-twisted circle.
Every block is a twisted-circle convolution with a smooth compactly
supported kernel $y^k h(y)$ times a binomial coefficient. Its Fourier
multipliers decay rapidly, so the full operator and its transpose are
trace class. Summing the Fourier modes of each diagonal block gives
\[
 \begin{aligned}
 \Tr_{\mathcal E_p} A_h
  &= l\sum_{m\in\Z}h(ml)\sum_\beta r_\beta\beta^m\\
  &= \sum_{v\mid p} l_v\sum_{n\in\Z}a_v^n h(nl_v).
 \end{aligned}
 \tag{13d}
\]
For the second equality, the sum of $\beta^m$ over $\beta^{f_v}=a_v$
is zero unless $f_v\mid m$, when it is $f_v a_v^{m/f_v}$.
This proves the uncentered local trace in (11.2.2.4), including
coincident local zeros and their full multiplicities. For instance,
over $\Q$ a test supported near $\log 2$ with $h(\log 2)=1$ gives
$\log 2$ at $p=2$. The quotient of the cutoff closures in (13b) is zero
for that same prime. Also, $M_\rho P^{\rm geom}$ in (10.4.1.6)
retains the projection taken before truncation.

\sourcepages{485}{486}{For the Hilbert closure and a proof of the local trace see pages 484a--b.}
\sourcepages{487}{488}{For finite-level estimates and the projection limit see pages 488a--f and 489a--b.}
\leaf{488a}{Paley--Wiener levels}{11.2.1 \quad THE RATIONAL FINITE LEVEL}{Fourier bandwidth and the lower endpoint}
Take $F=\Q$, $S=\{\infty\}$, trivial character, and the centered
degree-four space $X$ on \pg{446}, also allowing $\alpha=0$. For
an even Schwartz function $f$ with $f(0)=\widehat f(0)=0$ write
$F_f(x)=e^{x/2}\sum_{n\ne0} f(ne^x)$. Let $\mathcal T_Y$ be the
$X$-closure of these functions with $\supp\widehat f\subset[-Y,Y]$.
This is the common absolute and invariant level in the one-place
case of (7.3.2.3) and (7.3.2.14), with the indicated Hilbert weight.
Poisson summation gives
\[
 F_f(x)=e^{-x/2}\sum_{n\ne0}\widehat f(ne^{-x}),\qquad
 F_f(x)=0\quad(x<-\log Y).
\]
Every member of the closed level has the same lower support bound.
For $B_p=1-p^{-1/2}T_{-\log p}$, $T_yF(x)=F(x-y)$, the identity
$B_pF_f=F_{f-f(p\,\cdot)}$ gives
\[
                         B_p\mathcal T_Y\subset\mathcal T_{pY}.
\]
The factor $p$ is sharp. Choose $Y/2<c<Y$ with $pc>Y$, and a smooth,
even $\widehat f$ supported in $(-Y,Y)$ away from zero with integral
zero and $\widehat f(c)\ne0$. At $e^x=1/(pc)$ one has $F_f(x)=0$,
while in $F_f(x+\log p)$ precisely the two samples at $\pm c$ contribute.
The resulting function is nonzero in a neighborhood there. Taking $c$
arbitrarily close to $Y$ excludes any common level smaller than $pY$.

\point{Exhaustion of the global arithmetic space} In this rational
component any finite-place test averaged over the finite units is a finite
linear combination of indicators of fractional-ideal balls
$d\widehat{\Z}$, $d\in\Q_{>0}$. Each gives the sum with real-place test
$f(d\,\cdot)$. After summing and taking the even part, the value and the
integral of the resulting test are zero by the two adelic moment
conditions. Conversely every such $F_f$ occurs. To approximate it,
cut $\widehat f$ by smooth expanding cutoffs equal to one near zero.
Restore integral zero by subtracting a multiple of a fixed even bump
of integral one supported away from zero. That multiple tends to zero,
and the approximants converge in Schwartz space. The two Poisson
expressions give convergence in $X$ by rapid decay at both ends.
Thus $\overline{\bigcup_Y\mathcal T_Y}=J_\alpha$ and the orthogonal
level projections converge strongly to $P^{\rm geom}$.

\leaf{488b}{Paley--Wiener levels}{11.2.1 \quad THE COST OF AN ORBITAL EVALUATION}{A sharp norm at finite bandwidth}
For $0\le\alpha<1/2$ the local Euler zero at uncentered $s=0$ acts on
$\mathcal T_Y$ by $\ell_0(F)=\int F(x)e^{-x/2}\dd x$. Its operator
norm for the $X$-metric satisfies
\[
 \|\ell_0|_{\mathcal T_Y}\|\asymp_\alpha
        \frac{Y^{1/2-\alpha}}{(1+\log Y)^2}\qquad(Y\ge2).
 \tag{13e}
\]
For $L=\log Y$ the square of its norm is at most
$\int_{-L}^\infty e^{-x} w(x)^{-1}\dd x$ by \pg{488a}. The part below
zero is $\int_0^L e^{(1-2\alpha)y}(1+y^2)^{-2}\dd y$.
Splitting at $L/2$ bounds the whole integral by
$O_\alpha(e^{(1-2\alpha)L}(1+L)^{-4})$.

Choose a nonzero even real Schwartz function $\psi$ with Fourier support
in $[-1/4,1/4]$. Put $g(t)=t^2\psi(t)^2$ and $f(t)=g(t)-2g(2t)$.
The two moments of $f$ vanish and $F_f\in\mathcal T_1$. Unfolding for
$\Re s>1$ and continuing to $s=0$ gives
\[
 \ell_0(F_f)=2\zeta(0)\int_0^\infty f(t)\frac{\dd t}{t}
             =\int_0^\infty g(t)\frac{\dd t}{t}>0 .
\]
The continuation here is regular because $f(t)=O(t^2)$ at zero.
The test $Y^{1/2}f(Y\,\cdot)$ gives $T_{-L}F_f\in\mathcal T_Y$.
Its evaluation is $Y^{1/2}\ell_0(F_f)$ and its norm is at most
$2Y^\alpha(1+L^2)\|F_f\|_X$ by the weight inequality on \pg{520a}.
This proves the lower bound. On $\mathcal T_{pY}$ this same functional
annihilates $B_p\mathcal T_Y$, since
$p^{-1/2}\ell_0(T_{-\log p}F)=\ell_0(F)$. It detects a proper closed
subspace there. Formula (13e) gives the loss when $Y\to\infty$.

The same two-sided bound holds for
$\ell_\gamma(F)=\int F(x)e^{(-1/2+i\gamma)x}\dd x$ at a nonzero local
frequency $\gamma=2\pi j/\log p$. The upper bound is unchanged.
For the lower bound choose smooth even $\widehat f$ supported in
$1/4<|u|<3/4$ with integral zero but
$\int_0^\infty\widehat f(u)u^{-i\gamma}\dd u\ne0$, since the latter
weight is nonconstant. Poisson summation and Mellin continuation give
$\ell_\gamma(F_f)=2\zeta(1-i\gamma)\int_0^\infty
\widehat f(u)u^{-i\gamma}\dd u\ne0$. Translation now multiplies the
evaluation by $Y^{1/2-i\gamma}$. These nonzero frequencies also annihilate
the entire single-prime carved level over $\Q$ described on \pg{488c}.

\point{Multiplication by a norm power} Even at fixed $Y$ on the
uncentered polynomial metric $H_4$, a positive norm power is unbounded
on the corresponding level in the inherited norm. Indeed let $u$ be
any nonzero uncentered theta test of that level. Each $T_Ru$ for $R>0$
is in the same level, since real dilation shrinks the Fourier support.
For $b>0$, $\|T_Ru\|_{H_4}/R^2\to\|u\|_2$ by dominated convergence,
and the same holds for $e^{bx}u$. Hence
$\|e^{bx}T_Ru\|_{H_4}/\|T_Ru\|_{H_4}
\sim e^{bR}\|e^{bx}u\|_2/\|u\|_2$. The bandwidth restriction on
page 494 supplies no bounded multiplication operator of this kind.

\leaf{488c}{Paley--Wiener levels}{11.2.1 \quad EXHAUSTION FOR A FIXED HECKE CHARACTER}{Which finite levels exhaust the component}
Fix the generator on \pg{446c} and choose $S$ containing the archimedean
places, the ramified places and every exceptional finite component of
its test. It may be enlarged so that $\operatorname{Pic}(\mathcal O_{F,S})=1$.
First use the levels of (7.3.2.1) or (7.3.2.12) with $S_0$ empty.
At every width $\alpha\ge0$ in the degree-four metric, these levels
exhaust the full arithmetic closure $J$ in this component.

Indeed each norm translate of the generator is given by a test standard
outside $S$. Cut its additive Fourier transform at the archimedean
places by smooth expanding cutoffs which equal one near zero. The finite
Fourier components are already compactly supported. The Fourier value at
zero is preserved. Restore its integral zero by subtracting a multiple
of a fixed compactly supported Schwartz--Bruhat test on $\A_{F,S}$ of
integral one vanishing near zero, in Fourier coordinates. The multiple
tends to zero, and the inverse transforms converge in Schwartz--Bruhat
space. Every approximant belongs to a finite absolute level and to a
finite invariant level.

To check the Hilbert convergence choose the section dilating each ordinary
archimedean coordinate by $t^{1/d}$, $d=[F:\Q]$. Representatives for
the compact norm-one class group can be taken in a compact set of
ideles. The finite support conditions restrict the nonzero summation
indices to a fixed fractional-ideal lattice there. For $t\ge1$ an
archimedean Schwartz seminorm of order $M>d$ thus bounds the uncentered
sum by a constant times $t^{-M/d}\sum_{\xi\ne0}\|\xi\|^{-M}$ over that
lattice. For $t\le1$ Poisson summation gives the corresponding estimate
for the transformed test at $1/t$ times $t^{-1}$, since both moments
vanish. These continuous bounds give convergence in the Hilbert weight.
The norm translates span $J$ by \pg{446c} or the annihilator on \pg{446d}.

\point{Fixed extra unit conditions} Let $S_0$ be a fixed finite set of
unramified places disjoint from this $S$. Replacing $\mathbf1_{\mathcal O_v}$
by $\mathbf1_{\mathcal O_v^\times}$ at $v\in S_0$ applies the product
$B_{S_0}=\prod_{v\in S_0}B_v$ from \pg{514b} to the uncarved tests.
For $0\le\alpha<1/2$ every factor is boundedly invertible on $J$ by
the translation series there. The carved levels contain the image of
the two-moment tests and therefore still exhaust $J$ when the uncarved
ones do. For $\alpha\ge1/2$ and $S_0\ne\varnothing$ the carved union
has proper closure in the degree-four metric. Choose a local Euler zero
on $\Re z=-1/2$ away from the possible poles of the uncarved Tate
integral. Its evaluation is continuous and kills the whole carved
union by the local multiplier. It is nonzero on the generator by Hecke
nonvanishing at the reflected line.

With $S_0\ne\varnothing$ the adelic value at zero already vanishes on
the unit factors, so only the integral condition is independent in
(7.3.2.4) and (7.3.2.15). Over $\Q$ at a single $p$, take even
$\widehat f$ supported away from zero with integral one. Then $f(0)=1$.
At $s=0$ the carved Mellin transform
$2(1-p^{-s})\zeta(s)\int_0^\infty f(t)t^{s-1}\dd t$ equals
$-\log p$, by the simple pole of the last integral. At every nonzero
local frequency it vanishes. This explains the choice on \pg{488b}.

\leaf{488d}{Paley--Wiener levels}{11.2.1 \quad APPROXIMATING BY A PRIME LAYER}{The rate at a fixed prime}
In the rational trivial component let $\mathcal U_{p,Y}$ be the closed
centered level with real Fourier bandwidth $Y$ and the unit test at
$p$. Take $0\le\alpha<1/2$, put $\delta=1/2-\alpha$, and choose one of
the vectors $F\in\mathcal T_1$ of \pg{488b} with
$\ell_\gamma(F)\ne0$ at $\gamma=2\pi/\log p$. Then
\[
  \operatorname{dist}_X(F,\mathcal U_{p,Y})
       \asymp_{\alpha,p,F}(1+\log Y)^2Y^{-\delta}\qquad(Y\ge2).
 \tag{13f}
\]
For the lower bound every vector of $\mathcal U_{p,Y}$ lies in
$\mathcal T_{pY}$. Indeed the real test $g-g(p\,\cdot)$ has both
moments zero when $\widehat g(0)=0$, even if $g(0)\ne0$.
Moreover $\ell_\gamma$ vanishes on the whole carved level by \pg{488c}.
For $u\in\mathcal U_{p,Y}$,
$\|F-u\|_X\ge|\ell_\gamma(F)|/\|\ell_\gamma|_{\mathcal T_{pY}}\|$.
Equation (13e) at frequency $\gamma$ gives the claimed order.

For the upper bound put $N=1+\lfloor\log Y/\log p\rfloor$ and
\[
       v_N=\sum_{m=0}^{N-1}p^{-m/2}T_{-m\log p}F\in\mathcal T_Y .
\]
Then $B_pv_N\in\mathcal U_{p,Y}$ and
$F-B_pv_N=p^{-N/2}T_{-N\log p}F$. Its norm is at most
$2(1+N^2\log^2 p)p^{-\delta N}\|F\|_X$. Since $Y<p^N\le pY$,
this proves (13f). The approximation uses only translates at the
prescribed finite level.

\point{Several fixed primes} If $T$ is a fixed nonempty set of $r$
primes, truncate $\prod_{p\in T}B_p^{-1}F$ to the sum of
$n^{-1/2}T_{-\log n}F$ over $n\le Y$ with all prime factors in $T$.
It lies in $\mathcal T_Y$. Applying $\prod_{p\in T}B_p$ gives a vector
in the joint carved level. In a log interval $[u,u+1)$ there are at most
$O_T((1+u)^{r-1})$ such integers before the truncation: fix $r-1$
exponents and the last has only $O_T(1)$ choices. The omitted series
in $X$ is thus $O_{\alpha,T,F}((1+\log Y)^{r+1}Y^{-\delta})$,
by the translation bound and summation of the exponential tail.
The product of the $B_p$ is bounded for this fixed $T$. This estimate
controls exhaustion at a fixed set of places. The tests in (19) involve
a set of contributing primes that grows with their translation.

\leaf{488e}{Growing sets of primes}{11.2.1 \quad SIMULTANEOUS UNIT CONDITIONS}{A bandwidth for all primes up to a given norm}
Fix a number field $F$, $0\le\alpha<1/2$, and $\delta=1/2-\alpha$.
Let $T$ be any subset of the prime ideals of norm $q\le P$, with
arbitrary phases $|a_{\mathfrak p}|=1$. Extend the phases
multiplicatively to ideals supported on $T$. Write $N$ for ideal norm
and, on the centered degree-four space $X$, put
\[
 \begin{split}
 B_T&=\prod_{\mathfrak p\in T}
      (I-a_{\mathfrak p}q^{-1/2}T_{-\log q}),\\
 C_{T,Y}&=\sum_{\substack{\mathfrak n\ {\rm supported\ on}\ T\\N\mathfrak n\le Y}}
      \frac{a(\mathfrak n)}{\sqrt{N\mathfrak n}}T_{-\log N\mathfrak n},
 \qquad D_P=\frac{P^{1-\delta}}{1-\delta},\qquad m_P=\frac{D_P}{\log P}.
 \end{split}
\]
For some $c>0$ and all sufficiently large $P$ one has
\[
      \sup_{T,a}\|I-B_T C_{T,Y}\|_{X\to X}\le e^{-c m_P}
           \qquad(\log Y\ge4D_P).
 \tag{13f.1}
\]
The untruncated series $C_{T,\infty}$ equals $B_T^{-1}$ by absolute
operator convergence for each finite $T$. Expand
$B_T(C_{T,\infty}-C_{T,Y})$ over squarefree $\mathfrak d$ on $T$ and
$T$-supported $\mathfrak n$ with $N\mathfrak n>Y$. For
$x=\log(N\mathfrak d\,N\mathfrak n)$ the norm of each summand is at
most $2e^{-\delta x}(1+x^2)$. Set $t=\log P$ and $\eta=1/100$.
Use $1+x^2\ll_\eta t^2 e^{\eta x/t}$ and majorize the tail indicator
by $(N\mathfrak n/Y)^{1/t}$. Summing the positive series bounds the norm by
\[
  O_\eta(t^2)Y^{-1/t}
  \prod_{N\mathfrak p=q\le P}
     (1+q^{-\delta+\eta/t})(1-q^{-\delta+(1+\eta)/t})^{-1}.
\]
For fixed $u$, the \citetitle{https://www.jstage.jst.go.jp/article/jmath1948/20/1-2/20_1-2_233/_article/-char/en}{prime ideal theorem}
and partial summation give
$\sum_{q\le P}q^{-\delta+u/t}\sim e^u m_P$, counting prime ideals.
Indeed the terms $q<P^{1-\epsilon}$ are $o(m_P)$ for each fixed
$\epsilon>0$, and on the rest $\log q/t$ tends to $1$ as
$\epsilon\downarrow0$. The quadratic and higher terms in the logarithm
are $o(m_P)$ by the same split, since $\delta>0$. Thus the logarithm
of the product is $(e^\eta(1+e)+o(1))m_P$. Its leading coefficient
is strictly less than $4$, proving (13f.1).

\leaf{488f}{Growing sets of primes}{11.2.1 \quad THE MOVING LEVEL PROJECTIONS}{Convergence and the cost of truncating the inverse}
Use a fixed character component and a base $S$ as on \pg{488c}.
Let $Q_{T,Z}$ project onto its closed level with source cutoff $Z$,
imposing the unit condition at $T$ outside $S$, and let $J$ be the full
arithmetic closure. For $0\le\alpha<1/2$ set $Z(P)=\exp(5D_P)$.
For every choice of the sets $T=T(P)$
of \pg{488e} and every fixed trace-class operator $K$,
\[
       Q_{T,Z(P)}\longrightarrow P_J\ \text{strongly},\qquad
       \|K(Q_{T,Z(P)}-P_J)\|_1\longrightarrow0 .
 \tag{13f.2}
\]
Indeed, take an uncarved two-moment test of fixed level $y\ge1$
and write $f$ for its centered theta vector. A translate
$T_{-\log N\mathfrak n}f$ dilates the archimedean input by
$(N\mathfrak n)^{1/[F:\Q]}$ in ordinary coordinates. Its finite
supports are fixed; the product of normalized archimedean Fourier
bounds grows by $N\mathfrak n$, and each individual bound grows by
at most that factor. Thus $C_{T,Y}f$ comes from a test of level $yY$
for either cutoff on \pg{488c}. Applying $B_T$ imposes precisely the
unit conditions on this source. Take $Y=Z(P)/y$ in (13f.1). The resulting
distance to the carved level is at most $e^{-c m_P}\|f\|$. These
vectors $f$ are dense in $J$, and all carved levels lie in $J$.
This proves strong convergence on $J$ and on $J^\perp$. Finite-rank
approximation and self-adjointness give the trace norm assertion.
At the common source level $Z(P)$ the sum of successive projection
differences along any chain from $\varnothing$ to $T$ telescopes.
Its trace against $K$ therefore tends to zero. The circle trace on
\pg{514b} can already be nonzero for a fixed test supported near $\log 2$.

\point{A necessary scale for this truncated inverse}
Over $\Q$ take all $p\le P$ and $a_p=1$. The nonzero continuous
functional $f\mapsto\int f(x)e^{-\alpha x}\dd x$ on $X$ evaluates
$B_T C_{T,Y}$ as multiplication by $\Pr(S_P\le\log Y)$, where
$S_P=\sum_{p\le P}K_p\log p$ and the independent variables satisfy
$\Pr(K_p=k)=(1-p^{-\delta})p^{-k\delta}$. Partial summation gives
\[
 \mathbb E S_P=\sum_{p\le P}\frac{\log p}{p^\delta-1}\sim D_P,\qquad
 {\rm Var}(S_P)=\sum_{p\le P}\frac{p^\delta(\log p)^2}{(p^\delta-1)^2}
       \sim D_P\log P .
\]
For $\log Y\le(1-\epsilon)D_P$, Chebyshev's inequality makes this
multiplier tend to zero. In particular no power of $P$ suffices for
operator-norm convergence of this truncation. This lower bound concerns
$C_{T,Y}$; (13f.2) uses the sufficient level above.

\sourcepages{489}{489}{The corrected condensation limit and its missing operator are on page 489a.}
\leaf{489a}{The geometric projection}{11.2.1 \quad CONDENSATION ON THE SOURCE}{The operator in the level limit}
Let $P=P^{\rm geom}$ and $Q_Y$ be the level projections of \pg{488a}.
For any trace-class $K$, $KQ_Y\to KP$ in trace norm. This follows first
for a rank-one operator by $Q_Y^*=Q_Y\to P$ strongly and then by
finite-rank approximation. Thus for $K=M_\rho P A_h$ the candidate
limit in (11.2.1.4) omits the operator
\[
                         M_\rho P A_h(1-P).
 \tag{13g}
\]
It can be nonzero even over $\Q$ in the trivial character. Take
$a=1/2$, $\alpha=0$, $s=4$. The subspace $J=\operatorname{ran}P$ is
nonzero by (6). It is proper since evaluation at any critical-line
zero of $\zeta$ is a continuous nonzero functional on $H_4$ annihilating
$J$. Existence of such a zero is Hardy's theorem. All translations
$T_t$ preserve $J$.

Suppose $PA_h(1-P)=0$ for all $h\in C_c^\infty((0,\infty))$.
Approximating the mass at $t>0$ gives $PT_t(1-P)=0$ by strong
continuity. Thus $P$ commutes with every $T_t$ for $t>0$ and with its
inverse. Self-adjointness of $P$ gives commutation with
$T_t^*T_t$, the multiplier by $w(x+t)/w(x)$,
$w(x)=(1+x^2)^2$. For $t=\pm1$ the positive square roots of these
multipliers are $r_\pm(x)=(1+(x\pm1)^2)/(1+x^2)$. The functions
\[
       r_+-r_-=\frac{4x}{1+x^2},\qquad
       r_++r_--2=\frac{2}{1+x^2}
\]
generate the Borel sets since their ratio is $2x$. Their spectral
projections therefore give all indicator multipliers. A projection
commuting with all of these is itself multiplication by an indicator.
Translation invariance makes that indicator constant almost everywhere,
contradicting $0\ne P\ne1$.

Choose $h$ with $PA_h(1-P)\ne0$. Since the manuscript cutoffs converge
strongly to one, (13g) is nonzero for some finite cutoff. But every
$KQ_Y$ vanishes on $J^\perp$. Even strong convergence to $K$ in
(11.2.1.4) and (11.2.1.9) fails for this choice. For the convolution
order in (10.4.1.12), $K'=M_\rho A_hP$ instead satisfies $K'P=K'$.
Here $M_\rho A_h$ is trace class: a smooth cutoff equal to one on
$\supp\rho$ gives a smooth compactly supported kernel in both variables,
whose Fourier series on an enclosing square is absolutely summable.
Consequently $K'Q_Y\to K'$ in trace norm. This is a condensation
formula for the geometric restriction of (10.4.1.12).

\leaf{489b}{The geometric projection}{11.2.1 \quad REMOVING THE SPATIAL CUTOFF}{The scalar trace after both limits}
Write $P=P^{\rm geom}$, $S=1-P$, and
$C_h=PA_hS$. In every component at every $\alpha\ge0$ both $A_hS$ and
$SA_h$ are trace class by \pg{463a}.
Since $SA_hP=0$ by invariance of the arithmetic space,
\[
         PA_h-A_hP=A_hS-SA_h=C_h,\qquad \Tr C_h=0 .
 \tag{13h}
\]
The last equality follows from the off-diagonal form relative to
$PX\oplus SX$. For bounded cutoffs $0\le r_L\le1$ tending to one
pointwise, $\|(M_{r_L}-1)C_h\|_1\to0$ by finite-rank approximation.
In particular the difference between the traces of the two orders of
geometric truncation tends to zero, even though (13g) can remain nonzero
as an operator at finite cutoff.

For the levels that exhaust $J$ on pages 488a and 488c,
$Q_Y\to P$ strongly and
$P A_h Q_Y=A_h Q_Y$ exactly. For the cutoffs $r_L$ of compact support,
$M_{r_L}A_h$ is trace class. To see this with the weight $w$, use a
smooth compact cutoff equal to one on the support, expand the resulting
smooth unweighted kernel in Fourier series on a bounded square and
conjugate by $\sqrt w$. Both weight multipliers are bounded on that
square. Together with (13h) this proves
\[
 \lim_{L\to\infty}\left[
  \Tr(M_{r_L} P A_h)-
  \lim_{Y\to\infty}\Tr(M_{r_L}P A_h Q_Y)\right]=0 .
 \tag{13i}
\]
The inner limit is the trace of $M_{r_L}A_hP$. For tests supported away
from $x=0$, the trace of $M_{r_L}A_h$ is zero. By (13), both double-limit
calculations of the geometric convolution trace therefore give
$-\Tr(SA_h)$. At zero width its value is explicitly (9a) with a minus sign.

\point{Growth of the order correction} For $h_\tau(x)=h(x-\tau)$ the
same block is $C_{h_\tau}=P T_\tau A_h S$. In the degree-four metric,
uniformly in $L$ and for $\tau\ge0$,
\[
  |\Tr(M_{r_L}C_{h_\tau})|
        \le 2(1+\tau^2)e^{\alpha\tau}\,\|A_h S\|_1 .
 \tag{13j}
\]
This part of the geometric comparison already satisfies the target
exponent in (19), also at $\alpha=0$ with polynomial growth. The
remaining identity equates the limiting geometric trace plus its pole
terms with the real-place and centered prime distributions, as on
pages 514a--514b.

\sourcepages{490}{491}{The trace in Theorem 11.2.2.1 holds on the quotient by the closed subspace in (13c).}
\sourcepages{492}{492}{}
\sourcepages{493}{494}{For the bandwidth bound and the sign of the prime shell see pages 488a--b and 495--495a.}
\leaf{495}{The prime at two}{11.3.2 \quad THE LOCAL ORBITAL IDENTITY}{A counterexample at the first finite place}
Take $F=\Q$, $\chi=1$, $S=\{\infty\}$ and $\tau(t)=(t,1_f)$, $t>0$.
Put $U=\prod_p\Z_p^\times$ with Haar probability. Every idele has a unique
expression
\[
                  x=\tau(t)\,r\,u,\qquad
       t=|x|_{\A},\quad r\in\Q^\times,\quad u\in U.
\]
Indeed the finite valuations determine $|r|$, the real component determines
its sign, and the residual finite components are units. Thus $U$ represents
the norm-one idele classes.

In (11.3.2.2) the first indicator demands $r' u'x\in\tau(\R_{>0})$ for
$r'\in\Q^\times$, $u'\in U$. In the above coordinates this forces
$r'=r^{-1}$ and $u'=u^{-1}$. Haar probability on the infinite compact
group $U$ assigns mass zero to a singleton because arbitrarily many
distinct translates have the same mass. For every $x,y$ the indicated integral
is therefore zero. Multiplication by this lift annihilates the kernel,
including its first prime shell.

Choose $\eta\in C_c^\infty((3/2,5/2))$ with $\eta(2)=1$ and put
$\phi(x)=\eta(|x|_\A)$. It is a legitimate test supported at norms $>1$.
For $\varpi_2=2$ embedded at the finite place, $|\varpi_2^k|_\A=2^{-k}$.
Only $k=-1$ contributes. With $\mu_2(\Z_2^\times)=1$, (11.3.2.1)
therefore reads
\[
             0=\log2\sum_{k\in\Z}\eta(2^{-k})=\log2 .
 \tag{14}
\]
Thus Theorem 11.3.2.1 fails even for the ordinary zeta function. Its local
right side uses the unweighted sum printed in (11.3.2.1). The normalized
prime term on \pg{514a} carries the factor $2^{-1/2}$ at this power.

\point{Orientation of the stratum} The functions $\alpha_k$ on page 493
impose unit conditions at the first $k$ primes. Therefore
$\alpha_{k+1}\le\alpha_k$. The indicator of the shell removed at the
next prime is $\alpha_k-\alpha_{k+1}$, with positive sign. Either choice
of sign on the left of (14) still gives zero.

\leaf{495a}{Prime-shell weights}{11.3.2 \quad THE TOTAL TRACE}{Two normalized lifts}
In the coordinates on \pg{495}, Haar integration along a norm fibre is
counting measure on $\Q^\times$ times Haar probability on $U$. Let
$\omega$ be any nonnegative finitely supported probability on $\Q^\times$.
The one-variable lift $\ell_\omega(\tau(t)ru)=\omega(r)$ integrates to
$1$ on each such fibre. It is continuous and has compact support over a
compact norm interval.

For a kernel on the norm quotient use its pullback times the simultaneous
fibre average of $\ell_\omega(ax)\ell_\omega(ay)$ in $a$. At
$(\tau(t),\xi u\tau(t))$ that factor is
\[
 C_\omega(\xi)=\sum_{r\in\Q^\times}\omega(r)\omega(r\xi),
 \qquad \sum_{\xi\in\Q^\times} C_\omega(\xi)=1.
 \tag{15}
\]
For any integrable diagonal density $k(t,t)$ this yields the exact
unfolding
\[
 \int_0^\infty k(t,t)\frac{\dd t}{t}
 =\sum_{\xi\in\Q^\times}\int_0^\infty\!\int_U
       C_\omega(\xi) k(t,t)\dd u\,\frac{\dd t}{t}.
 \tag{16}
\]
This also applies to an integrable regularized diagonal. Formula (16)
preserves its total value with only finitely many nonzero weights.

\point{The first prime shell} For $S=\{\infty\}$ the first positive shell
at the second argument is $v_2(\xi)\ne0$. With $\omega=\delta_1$ its
total weight is $0$. With $\omega=(\delta_1+\delta_2)/2$ the only nonzero
weights are
\[
 C_\omega(1)=\tfrac12,\qquad
 C_\omega(2)=C_\omega(1/2)=\tfrac14 .
\]
The same shell has weight $\tfrac12$. For the signed difference in
(11.3.1.4) it has weight $-\tfrac12$. For example take
$k(t,t')=b(\log t)\overline{b(\log t')}$ with
$b\in C_c^\infty(\R)$ and $\int|b|^2=1$: the total trace is exactly one
for both lifts. Thus fibre normalization alone does not determine the
arithmetic distribution of trace among the primes.

\leaf{495b}{Prime-shell weights}{11.3.2 \quad INDEPENDENCE OF THE PRIME DISTRIBUTIONS}{Two primes obstruct a fixed invariant lift}
There is a stronger obstruction for any single normalized lift that is
invariant under the positive norm section. In the coordinates of page 495
such a lift has the form $\ell(\tau(t)ru)=a(r,u)$, with $a\ge0$ and
$\sum_r\int_U a(r,u)\dd u=1$. The lift may depend arbitrarily on the
finite unit. Write $b(r)=\int_U a(r,u)\dd u$.

In the rational trivial character component an operator kernel on the norm
line pulls back to $k_h(t,t')$, independent of the compact coordinates.
This applies after a norm cutoff as well. On the unfolded diagonal its
weight at $\xi$ after integration over the second unit is
\[
 \sum_r\int_U\!\int_U a(r,u)a(r\xi,ug)\dd g\dd u
             =\sum_r b(r)b(r\xi).
\]
Haar invariance and Tonelli give the equality. In particular the weight is
independent of $t$ and of the test. For any fixed prime shell $B_p$ in
the rational coordinate, let $c_p$ be the sum of these weights over
$\xi\in B_p$. Its trace contribution is exactly $c_p T(h)$, where $T(h)$
is the total diagonal integral. The factorization persists in a cutoff
limit whenever the total limit exists.

\point{Independence at $2$ and $3$} Choose one norm test equal to one at
$2$ and supported in a small neighborhood of $2$, and a second one
likewise at $3$. Choose the neighborhoods to exclude every other power
of $2$ or $3$. For $S=\{\infty\}$ the two right sides of (11.3.2.1),
evaluated on these two tests, give the matrix
\[
             \begin{pmatrix}\log 2&0\\0&\log 3\end{pmatrix}.
\]
Its rank is two. The proposed unfolded traces, for any two fixed shells
and any fixed lift as above, give an outer product
$(c_2,c_3)^{\mathsf T}(T(h_2),T(h_3))$ of rank at most one. Reversing
the sign of the shell indicators leaves the obstruction intact.

Thus no normalized, norm-translation-invariant choice in this class can
realize (11.3.2.1) simultaneously at these two places. The obstruction
uses the scalar quotient kernel and fixed lift. An arithmetic construction
with additional place-dependent operators or a norm-dependent lift would
require its own identity relating those data to the prime distributions.

\sourcepages{496}{496}{The local formula used here has the counterexample on page 495.}
\sourcepages{497}{506}{}
\sourcepages{507}{508}{For the fundamental domain and the regularity assertion in Proposition 12.2.2.2 see page 508a.}
\leaf{508a}{Diagonal regularity}{12.2.2 \quad LIFTING THE DIAGONAL}{A fundamental domain and vanishing lifts}
For $F=\Q$ the quotient of $\A_\Q^\times$ by $\Q^\times$ is
$\R_{>0}\times U$ in the coordinates on page 495. It is disconnected
since $U$ maps continuously onto $(\Z/3\Z)^\times$. Equation (12.2.1.12)
requires the image of its fundamental domain to be the entire quotient.
The image of a connected set is connected, so the domain with nonempty
interior and connectedness stipulated on page 507 cannot exist in this
application. The product slice $r=1$ on page 495 is instead clopen and
maps homeomorphically onto the quotient. The Haar decomposition and
the total integration formula (16) follow directly on this slice.

\point{Pointwise regularity in Proposition 12.2.2.2(4)}
Even for a connected fundamental domain in the stated hypotheses, a
lift can erase a singularity. Take $G=\R\times\Z$, $H=\{0\}\times\Z$,
$D=\R\times\{0\}$, and the normalized indentation $\ell(x,n)=\mathbf1_{n=0}$.
The factor (12.2.2.1) at $((x,n),(y,m))$ is $\mathbf1_{n=m}$.
Use the open symmetric neighborhoods
$(-2^{-j},2^{-j})\times\{0\}$, a $D'$-sequence as on page 498.

Choose $a\in C_c^\infty(\R)$ equal to one near zero. In quotient
coordinates put
\[
   k(x,y)=a(x)a(y)\cos(\omega\log|x|),\qquad
   \omega=\pi/\log 2 ,
\]
with arbitrary value at $x=0$. This bounded, compactly supported kernel
has rank at most one as an integral operator. At $(0,0)$ its small-square
averages are
\[
 \frac1{2r}\int_{-r}^r\cos(\omega\log|x|)\dd x
     =\Re\frac{r^{i\omega}}{1+i\omega}
     =\frac{(-1)^j}{1+\omega^2}\qquad(r=2^{-j}).
\]
Thus $(0,0)$ is not a regular point under the definition on page 498.
At the lift $((0,0),(0,1))$ the kernel is identically zero on a
neighborhood, so the lifted point is regular. For the rational lifts
of (15) the factor is locally constant. Where it is nonzero, division
by that constant recovers regularity. Where it is zero the lifted
averages are zero. This is enough for the forward unfolding of an
integrable regular diagonal in (16).

\sourcepages{509}{512}{The pointwise regularity statement for lifts is qualified on page 508a.}
\sourcepages{513}{513}{}
\sourcepages{514}{514}{Section 12.3 continues on page 514a. Section 12.4 concerns a fixed finite set of places.}
\leaf{514a}{The full critical strip}{12.3 \quad THE GLOBAL TRACE OVER THE RATIONALS}{Weil's distribution in log coordinates}
Let $h\in C_c^\infty(\R)$ with $h(0)=0$ and
$B(h)=H(\tfrac12)+H(-\tfrac12)$. Normalize the Haar measure on each
finite unit group to give it volume one. The rational Weil distribution is
\begin{align*}
 W(h)&=\sum_{n\ge2}\frac{\Lambda(n)}{\sqrt n}
             \big[h(\log n)+h(-\log n)\big]\\
     &\quad+\int_{\R}h(x)\,
             \frac{e^{-|x|/2}}{1-e^{-2|x|}}\dd x .
\end{align*}
Here $\Lambda(p^k)=\log p$ for $k\ge1$ and $\Lambda(n)=0$ otherwise.
The real-place factor is the sum of the positive and negative real unit
contributions. For $x>0$ it equals
$\tfrac12e^{x/2}((e^x-1)^{-1}+(e^x+1)^{-1})$, and inversion gives the
same factor at $-x$. The vanishing of $h(0)$ makes the integral locally
absolutely convergent and kills the terms supported at $0$.

\point{Full-strip identity} Weil's classical explicit formula gives
\[
 W(h)=B(h)-\sum_\lambda m_\lambda H(\lambda)
     =B(h)-\Tr(A_hP_{1/2}).
 \tag{17}
\]
The second equality is (5) and (12). One can check the signs on positive
logarithmic support by Mellin inversion of $-\zeta'/\zeta(z+\tfrac12)$:
the pole of $\zeta$ contributes $H(\tfrac12)$, and its nontrivial and
trivial zeros contribute $-m_\lambda H(\lambda)$ and
$-H(-2k-\tfrac12)$, $k\ge1$. Adding the real-place integral adds back
$\sum_{k\ge0}H(-2k-\tfrac12)$. The general test above is covered by the
classical distributional formula.

\point{The complementary projection} For each compactly supported cutoff
$r$ in (13), the localized convolution has trace
\[
     \Tr(M_r A_h)=h(0)\int_{\R} r(x)\dd x.
\]
Localize both variables to a bounded interval containing the kernel support
and insert smooth cutoffs equal to one there. On a larger interval the
smooth convolution kernel has an absolutely summable double Fourier series.
Multiplication by bounded $r$ gives the diagonal integral term by term, and
conjugation by $\sqrt w$ on this interval cancels on the diagonal.
Equations (13) and (17) now identify the full-strip limit of
$B(h)+\Tr(M_r A_h(1-P_{1/2}))$ with $W(h)$.

\leaf{514b}{The individual prime terms}{12.3 \quad AN ARITHMETIC LOCAL OPERATOR}{A trace at each prime ideal}
Let $\mathfrak p$ be unramified for a unitary Hecke character, let
$l=\log N\mathfrak p$ and let $a=\chi_{\mathfrak p}(\varpi_{\mathfrak p})$.
On the twisted circle $u(t+l)=a\,u(t)$ use $L^2([0,l],\dd t)$ and right
translations. For $h\in C_c^\infty(\R)$, $h(0)=0$, set
$g(x)=e^{-|x|/2}h(x)$. The operator
$C_{\mathfrak p,h}u(t)=\int g(x)u(t+x)\dd x$ is trace class and
\[
 \Tr C_{\mathfrak p,h}
  = l\sum_{m\ge1} e^{-ml/2}
      \big[a^m h(ml)+\overline a^{\,m}h(-ml)\big].
\]
Indeed, writing $a=e^{i\theta}$, the orthogonal Fourier modes are
$e^{i(\theta+2\pi j)t/l}$, $j\in\Z$. The diagonal multipliers are the
Fourier integrals of $g$ at these frequencies. They are summable by two
integrations by parts, since $g$ has an integrable weak second derivative.
Fourier series of the periodization at zero give
$l\sum_{m\in\Z}a^m g(ml)$, proving the formula. Dropping the exponential
factor gives the unweighted prime distribution printed in (11.3.2.1)
at the trivial character. This arithmetic local operator uses the prime
period and the Euler character explicitly.

\point{Effect of the Euler factor in the narrow Hilbert space}
In centered coordinates the local reciprocal Euler factor acts on $X$
by $B_{\mathfrak p}=1-a e^{-l/2}T_{-l}$, where $T_yf(x)=f(x-y)$.
For $0<\alpha<\tfrac12$ its inverse converges in operator norm:
\[
 B_{\mathfrak p}^{-1}
       =\sum_{m\ge0}a^m e^{-ml/2}T_{-ml},\qquad
 \sum_{m\ge0}e^{-ml/2}\|T_{-ml}\|
       \le 2\sum_{m\ge0}(1+m^2l^2)e^{-(1/2-\alpha)ml}<\infty .
\]
The arithmetic closure $J_\alpha$ is invariant under all translations,
so $B_{\mathfrak p}J_\alpha=J_\alpha$. The untruncated quotient by
this range is zero, even though the local circle trace can be nonzero.
For the rational finite levels, (13e) gives the divergence rate of an
evaluation annihilating the Euler image. On the uncentered zero line
(13d) realizes the arithmetic circle traces on the Hilbert quotient.

For $\zeta$, finite products of the inverses do not even stay bounded on
$X$. Evaluation at $z=0$ multiplies by
$\prod_{p\le P}(1-p^{-1/2})^{-1}\to\infty$, since $\sum_p1/p=\infty$.
Point evaluation is continuous and nonzero there. The local trace formula
thus supplies each prime distribution independently. Its identification
with a shell of the global projection is the additional global assertion.

\sourcepages{515}{515}{}
\sourcepages{516}{518}{The real-place range and the lower endpoint in (12.4.1.23) are calculated on page 518a.}
\leaf{518a}{The real place}{12.4.1 \quad THE LOWER CUTOFF}{The finite endpoint term}
Take $F=\Q$, $S=\{\infty\}$. In Step II on page 516, multiplication by
the one-sided cutoff, equal to one near $x=0$, is asserted to map
$\mathcal S(\R)$ into the polar-free space $\mathcal S_{00}(\R)$. Applied
to $e^{-\pi x^2}$ the result still has nonzero value at zero. Its local
Mellin integral has a pole at $s=0$, contrary to the polar-free
property on pages 259--260. Cutting away a neighborhood of zero removes this
pole, and changes the subsequent volume calculation.

With the sharp outer and Fourier cutoffs of (12.4.1.22), also impose
$|x|\ge1/L''$. On the positive fundamental half-line for $\{\pm1\}$,
$w=xu$ gives the exact logarithmic length
\[
 \int_{\max\{|w|/L',\,1/L''\}}^{L}\frac{\dd x}{x}
    =\big[\log L-\log\max\{|w|/L',1/L''\}\big]_+
\]
where the left integral is taken over the empty set if the lower bound
exceeds $L$. Assume $LL''>1$ and put $a=L'/L''$. The one-sided length
in (12.4.1.24) exceeds this length by $\mathbf1_{|w|<a}\log(a/|w|)$.
For the test $\eta_\xi$ of (12.4.1.20) the lost term is therefore,
in real Lebesgue normalization,
\[
       \int_{|w|<a}\widehat\eta_\xi(w)\log(a/|w|)\dd w .
 \tag{17a}
\]
Its absolute value is at most $2a\|\widehat\eta_\xi\|_\infty$. Thus
$L'/L''\to0$ removes it for each fixed real-place test.

\point{Balanced cutoffs} If $L=L'=L''\to\infty$, then $a=1$. Choose a
nonzero nonnegative smooth real lift $\widetilde\phi$ supported in
$(1,11/10)$, extended by zero on the negative half-line. It descends by
summation over $\{\pm1\}$ to a permitted positive-norm test.
For $\xi=1$, the function $\eta_1$ is nonnegative and supported where
$1/(1+z)\in(1,11/10)$, hence $|z|<1/4$. The real part of
$\widehat\eta_1(w)$ is strictly positive for $|w|\le1$ by the cosine
formula. In fact the Fourier transform of
$\mathbf1_{|w|<1}\log(1/|w|)$ equals
$\operatorname{Si}(2\pi z)/(\pi z)$, with value $2$ at zero.
It is strictly positive because the alternating half-period integrals of
$\sin t/t$ decrease in magnitude. Consequently the other nonnegative
test $\eta_{-1}$ also gives a positive term in (17a). The finite endpoint
correction survives summation over the two units. This calculation
also survives smoothing the balanced cutoffs across strips of width
$\epsilon$ with $\epsilon L=o(1)$. In the $x$ variable the transition
has logarithmic length $O(\epsilon L)$ and costs
$O(\epsilon L\|\widehat\eta_\xi\|_1)$. At $|u|=L'$ the transition in
$w=xu$ has length $O(\epsilon x)$ and costs
$O(\epsilon L\|\widehat\eta_\xi\|_\infty)$ after integration in
$\dd x/x$. Formula (17) uses the classical explicit formula independently.

\sourcepages{519}{519}{For the established identity and the hypothesis needed on a narrow strip see page 520.}
\leaf{520}{Narrow strips}{12.5.1 \quad THE TRACE IDENTITY}{The contribution outside the strip}
For $\zeta$ take $0\le\alpha\le\tfrac12$ and degree four. With the notation
of pages 29 and 514a, the corrected form of (12.5.1.1)--(12.5.1.4) is
\[
 \boxed{\quad W(h)=B(h)-\Tr(A_hP_\alpha)-E_\alpha(h),\qquad
 E_\alpha(h)=\sum_\lambda(m_\lambda-n_{\alpha,4}(\lambda))H(\lambda).\quad}
 \tag{18}
\]
Equation (17) includes all nontrivial zeros. Equation (12) counts
each interior order in full and at most two copies at each boundary
zero. Both sums converge absolutely. Put
\[
 G_\alpha(h)=B(h)+\lim_{L\to\infty}
             \Tr\big(M_{r_L}A_h(1-P_\alpha)\big).
\]
Then $G_\alpha(h)=B(h)-S_\alpha(h)$ and
$G_\alpha(h)-G_{1/2}(h)=G_\alpha(h)-W(h)=E_\alpha(h)$.

\point{The independent geometric estimate} To obtain the zero-location
claim at width $\alpha$ it suffices to establish, by arithmetic or
geometric means, for $h_\tau(x)=h(x-\tau)$, $h\in C_c^\infty((1,2))$,
\[
 G_\alpha(h_\tau)-W(h_\tau)
         =O_{h,\alpha,\epsilon}(e^{(\alpha+\epsilon)\tau})
       \quad(\tau\ge0,\ \epsilon>0).
 \tag{19}
\]
The residue argument on \pg{29} detects every zero strictly to the right
of the boundary; reflection covers the left. If all zeros are in the
closed strip, the residual boundary terms are $O_h(e^{\alpha\tau})$ by
absolute summability. Thus (19) detects location even at a width with
a boundary zero. One test suffices as on \pg{522}.

For a fixed test supported away from norm one, only finitely many primes
occur in $W$. On translating to $h_\tau$ the available prime powers
grow with $\tau$. The vanishing of sufficiently high prime shells for
each individual test supplies neither their total nor a bound uniform
under translation. Equation (16) alone permits both first-shell weights
on page 495a. A calculation at a fixed finite set of places in \S12.4
therefore leaves (19) open.

\point{Other characters} For a fixed unitary Hecke character take the
full width $\beta=\tfrac12$ on \pg{446c}. The Hecke--Weil formula and
(12) give the same subtraction with multiplicities (4) in each
component. Reflection is $\rho\mapsto1-\overline{\rho}$ for
the zero set of that component. No uniformity across characters is
required for the conditional implication.

\leaf{520a}{Translated tests}{12.5.1 \quad BOUNDS UNDER TRANSLATION}{What the available bounds control}
Let $T_\tau f(x)=f(x-\tau)$. Directly from the Hilbert weight,
\[
 \|T_\tau\|_{X\to X}\le 2 e^{\alpha|\tau|}(1+\tau^2).
\]
Since $A_{h_\tau}=T_\tau A_h$, this controls the translated
trace-class operator for each fixed projection. The narrow spectral sum
even satisfies $|S_\alpha(h_\tau)|\le
e^{\alpha\tau}\sum_\lambda m_\lambda |H(\lambda)|$ for $\tau\ge0$.
Applying the operator bound to the full projection uses its own weight
with exponent $\tfrac12$. Subtracting the two estimates still leaves the
full-strip exponent. Cancellation in (19) is the remaining arithmetic issue.

\point{The term requiring cancellation} For the positive-support test of
page 29 put $v_h(u)=u^{-1/2}h(\log u)$ and $X=e^\tau$. The finite-place
part of $W(h_\tau)$ is
$X^{-1/2}\sum_{n\ge2}\Lambda(n)v_h(n/X)$. The real-place part is
$X^{-1/2}H(-\tfrac12)+O_h(X^{-5/2})$ by its kernel on page 514a.
Consequently (18) gives the more explicit identity
\[
 E_\alpha(h_\tau)=
  X^{-1/2}\!\left(X\int_0^\infty v_h(u)\dd u
                -\sum_{n\ge2}\Lambda(n)v_h(n/X)\right)
       -S_\alpha(h_\tau)+O_h(X^{-5/2}).
 \tag{20}
\]
Thus at each fixed width the bound in (19) is equivalent, for these tests,
to a smoothed prime error $O_{h,\alpha,\epsilon}(X^{1/2+\alpha+\epsilon})$.

\point{An unconditional comparison} The classical zero-free region yields
for each fixed smooth weight of this support some $c>0$ with
\[
 \sum_{n\ge2}\Lambda(n)v_h(n/X)
   = X\int_0^\infty v_h(u)\dd u
      + O_h\big(Xe^{-c\sqrt{\log X}}\big).
\]
Indeed apply (17) and split the zero sum at $T=\exp(\sqrt{\log X})$.
For $|\Im\rho|\le T$ use
$\Re\rho\le1-c_0/\log(2+T)$ and the absolute summability of the transformed
weight. For the tail, its rapid decay and zero counting give
$O_h(X T^{-2})$. The real-place remainder in (20) is smaller.
This classical estimate still allows a contribution of scale $X^\beta$
with $\tfrac12<\beta<1$. A fixed zero-free half-plane gives the stronger
comparison on \pg{520b}.

\leaf{520b}{A fixed zero-free half-plane}{12.5.1 \quad AN EXTERNAL ARITHMETIC BOUND}{How far the full trace persists}
Fix a number field $F$ and a unitary Hecke character $\chi$. Suppose an
independent zero-free theorem places all its nontrivial zeros in
$\Re\rho\le\sigma_*$, where $\tfrac12\le\sigma_*<1$. Put
$b=\sigma_*-\tfrac12$. For every $\alpha>b$ the arithmetic projection
of \pg{446c} has the full zero trace. Indeed the functional equation
gives $|\Re(\rho-\tfrac12)|\le b$. Each zero is then interior, so (18)
has zero defect. For any smaller width and any translated
smooth test as on \pg{29} the defect satisfies
\[
   |E_{\alpha,\chi}(h_\tau)|
    \le e^{b\tau}\sum_\rho m_\rho |H(\rho-\tfrac12)|
       \qquad(\tau\ge0).
 \tag{20a}
\]
The constant is finite by rapid vertical decay and fixed-character zero
counting. At $\alpha=b$ the defect consists exactly of the boundary
multiplicities $m_\lambda-n_{b,s}(\lambda)$ in (4).

For the smooth weight $v$ of (22), in the notation of \pg{522a} the same
arithmetic input gives
\[
                 A_\chi(X)=M_\chi(X)+O_{F,\chi,v}(X^{\sigma_*}).
 \tag{20b}
\]
Shift the Mellin integral for $-V(s)L_f'(s,\chi)/L_f(s,\chi)$ to
$\Re s=(1-\sigma_*)/2$. This lies strictly between zero and every
nontrivial real part. The vertical and horizontal estimates are as on
\pg{522a}; the new integral is $O(X^{(1-\sigma_*)/2})$.
The nontrivial residues have total absolute value at most
$X^{\sigma_*}\sum_\rho m_\rho|V(\rho)|$. Any norm-character pole
gives exactly $M_\chi(X)$.

\point{Independent bounds for all unitary characters}
B.~Hu's \citetitle{https://hexagonmath.org/2610.00216v2}{Theorem 1.2}
states $\sigma_*=7/8$ over every number field and has an accompanying
\citetitle{https://github.com/Rogerhu12/hecke-seven-eighths}{Lean statement and proof}.
With this input the full trace persists for $\alpha>3/8$ and the
Gaussian prime sum on \pg{522c} is $O_v(X^{7/8})$.
B.~Liu's \citetitle{https://hexagonmath.org/2611.00741v1}{Theorem 7.33}
states the smaller all-character threshold $\theta_G<0.87333264$.
With this preprint as input, (20b) gives $O_{F,\chi,v}(X^{0.87333264})$
and the full trace persists for $\alpha>0.37333264$. These implications
use the stated external half-planes. Equation (24) requires the
exponent $1/2$.

\leaf{520c}{Hecke prime sums}{12.5.1 \quad A TEST FOR THE CONTRIBUTING PRIMES}{Removing the pole while retaining every zero}
Let $F$ have degree $d$, and fix a unitary Hecke character $\chi$.
Write $L(s)=L_f(s,\chi)$ for the primitive finite Euler product and
extend $\chi(\mathfrak a)$ by zero on ideals not coprime to its conductor.
Let $\beta_\chi$ be the supremum of the real parts of its nontrivial
zeros. The functional equation and zero count give
$1/2\le\beta_\chi\le1$. Pages \pagelink{520e}{520e--g} isolate a sum
over ordinary integers and rational primes, with both factors within
powers of $\log X$ of $\sqrt X$. If $\beta_\chi>1/2$, this sum has
growth exponent exactly $\beta_\chi$.

Take the fixed weight $v$ and its Mellin transform $V$ from (22).
Set $s_*=1-it_0$ when $\chi=|\cdot|_\A^{it_0}$ and $s_*=1$
otherwise. Let $\nu$ be the distribution of a sum of independent uniform
variables on $[0,h_j]$, $h_j=3/(4\pi^2 j^2)$, $j\ge1$. Define
\[
 v_*(u)=\int v(ue^{-y})\dd\nu(y),\qquad
 w_\chi(u)=(u\partial_u+s_*)^2 v_*(u),\qquad
 b=\exp(15/8).
\]
Since $\sum h_j=1/8$, the support is contained in $[e^{5/4},b]$.
Under the Mellin convention on \pg{522}, integration by parts gives
\[
 W_\chi(s)=(s-s_*)^2 V(s)
          \prod_{j\ge1}\frac{e^{h_js}-1}{h_js}.
 \tag{20c}
\]
The product converges normally by $\sum h_j<\infty$ and is nonzero
in $\Re s>0$. Thus $W_\chi$ has a double zero at a possible pole of
$L$ and is nonzero at each nontrivial zero. The double zero cancels
the poles of both $L$ and $L'$.

For $k\ge0$, $s=-k+it$, the transform satisfies
\[
 |W_\chi(-k+it)|\ll_\chi (1+k+|t|)^2 e^{-5k/4-c\sqrt{|t|}}.
 \tag{20d}
\]
Indeed $|V(-k+it)|\le e^{-5k/4}$ by (22). Each new factor is the
average of $e^{h_jsu}$ on $0\le u\le1$, so its modulus is
at most $\min(1,2/(h_j|s|))$. For $j\le c_0\sqrt{|t|}$ this bound
is at most $1/2$ with a fixed sufficiently small $c_0>0$.
On every fixed vertical strip the transform also decays faster than
any inverse power, since $w_\chi$ is smooth and compactly supported.

\leaf{520d}{All-ideal sums}{12.5.1 \quad THE UNRESTRICTED IDEAL FACTOR}{Decay after summing the unrestricted ideal}
For $j=0,1$ put
\[
 S_j(Y)=\sum_{\mathfrak a}\chi(\mathfrak a)(\log N\mathfrak a)^j
                     w_\chi(N\mathfrak a/Y),\qquad K_\chi(Y)=S_0(Y).
\]
The sums extend over nonzero integral ideals. There exist $C,c>0$
depending on $F,\chi$ such that
\[
        |S_0(Y)|+|S_1(Y)|\le C\exp(-cY^{1/(2d)})\qquad(Y\ge1).
 \tag{20e}
\]
Also $K_\chi(Y)=0$ for $Y<1/b$ since ideal norms are at least one.

On $\Re s=2$, Mellin inversion expresses $S_j$ by integrating
$W_\chi(s)(-1)^jL^{(j)}(s)Y^s$. Shift to $\Re s=-k$ for a positive
integer $k$. The integrand is entire by (20c) and the functional
equation. For each fixed $k$, rapid decay and polynomial bounds on
fixed vertical strips justify the shift. In order to vary $k$ one
needs the following dependence on the line. For all sufficiently large
$k$, Tate's functional equation gives
\[
      |L(-k+it)|+|L'(-k+it)|
        \le C_1^{k+1}(1+k+|t|)^{dk+C_2}.
 \tag{20f}
\]
In detail, write the archimedean factors as $d$ copies of
$\Gamma_{\R}(s+\mu_v)=\pi^{-(s+\mu_v)/2}\Gamma((s+\mu_v)/2)$
with fixed shifts, splitting each complex factor by duplication.
In the reflected numerator the shifts are $\overline{\mu_v}$.
The Euler product at $1+k-it$ is bounded. Reflect the reciprocal
gamma in each denominator using $\Gamma(z)^{-1}
=\Gamma(1-z)\sin(\pi z)/\pi$. The two resulting gamma factors
have positive real parts and equal imaginary magnitudes.
For $a\ge1$, Stirling's formula gives
$|\Gamma(a+iy)|\ll(1+a+|y|)^{a-1/2}e^{-\pi|y|/2}$,
using $|y|\arctan(a/|y|)\le a$. Their exponentials cancel the
growth of the sine. Powers of the conductor and $\pi$ cost $O(C_1^k)$.
The bound holds also within distance $1/4$ of the line on increasing
$C_2$; Cauchy's estimate then gives the derivative.

Combining (20d) and (20f), the integral is bounded by
$C_3 k^{C_4}(C_3 k^{2d}/Y)^k$. Indeed split
$(1+k+|t|)^{dk+O(1)}$ into the contributions of $1+k$ and $|t|$,
at cost $O(C_5^k)$, and use
$\int_0^\infty t^j e^{-c\sqrt t}\dd t=2c^{-2j-2}\Gamma(2j+2)$.
Taking $k=\lfloor c_1Y^{1/(2d)}\rfloor$ with small fixed $c_1>0$
proves (20e); bounded $Y$ is absorbed in $C$.

\leaf{520e}{Balanced ideal sums}{12.5.1 \quad AN IDENTITY ON INTEGRAL IDEALS}{Only two near-square-root factors remain}
Let $\mu_F$ and $\Lambda_F$ denote the ideal M\"obius and von Mangoldt
functions. Fix a smooth $0\le\eta\le1$ equal to one for $0<t\le1$
and zero for $t\ge2$, and put $\omega=1-\eta$. For large $X$ set
\[
 L_X=\log(2X),\qquad r=d+1,\qquad U=\sqrt X/L_X^r,\qquad
 A^\sharp_\chi(X)=\sum_{\mathfrak a}\Lambda_F(\mathfrak a)
                 \chi(\mathfrak a) w_\chi(N\mathfrak a/X).
\]
For every $A>0$ the following identity with error holds:
\begin{align*}
 A^\sharp_\chi(X)
   ={}&\sum_{\mathfrak m,\mathfrak l}
       \mu_F(\mathfrak m)\chi(\mathfrak m)
       \Lambda_F(\mathfrak l)\chi(\mathfrak l)
       \omega(N\mathfrak m/U)\omega(N\mathfrak l/U)\\[-2pt]
     &\hspace{42mm}\cdot K_\chi\big(X/(N\mathfrak m\,N\mathfrak l)\big)
             +O_{A,F,\chi}(X^{-A}).
 \tag{20g}
 \end{align*}
Every nonzero summand satisfies
$U<N\mathfrak m,N\mathfrak l\le bX/U$. If $K_\chi$ is expanded,
the third ideal $\mathfrak k$ has $N\mathfrak k\le bL_X^{2r}$.

To prove this, let a subscript $<$ on an ideal function denote
multiplication by $\eta(N\mathfrak a/U)$, and a subscript $>$ denote
multiplication by $\omega(N\mathfrak a/U)$. Dirichlet convolution on
ideals gives $\log N=\mathbf1*\Lambda_F$ and
$\mu_F*\mathbf1=\mathbf1_{\mathfrak a=(1)}$, where $\mathbf1$ is the
constant function one. Dropping the index $F$, expand the exact identity
\[
  \Lambda=\Lambda_<+\mu_<*\log N
          -\mu_<*\mathbf1*\Lambda_<+\mu_>*\mathbf1*\Lambda_>.
\]
Indeed the second and third terms combine to $\mu_<*\mathbf1*\Lambda_>$.
Apply the weight $\chi(\mathfrak a)w_\chi(N\mathfrak a/X)$. The first
term has disjoint support for large $X$. The second contains
$S_1(X/N\mathfrak m)$ with $N\mathfrak m\le2U$.
The third contains $S_0(X/(N\mathfrak m\,N\mathfrak l))$
with both norms at most $2U$. All these arguments are at least
$L_X^{2d+2}/4$ for large $X$. Formula (20e) therefore saves
$\exp(-c' L_X^{1+1/d})$. The total absolute size of the outer
coefficients is at most a fixed power of $X$, by absolute convergence
of $\zeta_F(2)$. This proves the error in (20g). The support bounds
follow from the lower cutoff $\omega$ and the upper support endpoint
$b$ of $w_\chi$.

\leaf{520f}{Degree-one prime ideals}{12.5.1 \quad THE REMAINDER IN THE BALANCED SUM}{Prime powers and common rational prime divisors}
Group the M\"obius coefficients by integer norm, and the degree-one
prime ideals by their underlying rational prime:
\[
 u_\chi(n)=\sum_{N\mathfrak m=n}\mu_F(\mathfrak m)\chi(\mathfrak m),
 \qquad b_\chi(p)=\sum_{N\mathfrak p=p}\chi(\mathfrak p).
\]
Retain the following part of (20g), over rational primes $p$:
\[
 \mathcal B_\chi(X)=\sum_{\substack{n,p\\p\nmid n}}
     u_\chi(n)b_\chi(p)\log p\,\omega(n/U)\omega(p/U)
                  K_\chi\big(X/(np)\big).
 \tag{20h}
\]
For each $\epsilon>0$ one has
\[
        A^\sharp_\chi(X)-\mathcal B_\chi(X)
           \ll_{F,\chi,\epsilon}X^{1/4+\beta_\chi/2+\epsilon}.
 \tag{20i}
\]
In particular $3/4+\epsilon$ is an unconditional exponent without any
fixed zero-free half-plane.

For a removed prime power $\mathfrak l=\mathfrak p^j$, with
$N\mathfrak p=p^f$ and $fj\ge2$, expand $K_\chi$ in (20g).
At fixed $\mathfrak l,\mathfrak k$ the M\"obius sum has scale
$Y=X/(N\mathfrak k\,N\mathfrak l)\le X/U$ and test
$\phi(t)=w_\chi(t)\omega(tY/U)$. This family has fixed compact
support and bounded smooth seminorms: if the transition meets the
support, $Y/U$ lies in a fixed compact interval. Nonzero terms have
$Y>U/b$. For every $\sigma>\beta_\chi$,
\[
       \sum_{\mathfrak m}\mu_F(\mathfrak m)\chi(\mathfrak m)
          \phi(N\mathfrak m/Y)\ll_{F,\chi,\sigma}Y^\sigma .
 \tag{20j}
\]
Shift the Mellin integral of $1/L$ to $\sigma$. At large heights,
choose disks centered at $c+it$ with fixed $c>\max(2,\sigma+1)$.
Choose their two radii strictly between $c-\sigma$ and
$c-\beta_\chi$. A branch of $\log L$ exists there; its real part
is $O(\log(2+|t|))$ from above by fixed-strip growth, and its central
value is bounded using the Euler product. Borel--Carath\'eodory thus
bounds $1/L$ polynomially throughout $\sigma\le\Re s\le c$.
At bounded heights $1/L$ is regular even at a pole of $L$.

The removed prime powers of norm at most $R=bX/U$ have total
$\Lambda_F$-weight $O_F(\sqrt R)$: $\sum_{\mathfrak p\mid p}f\le d$
and $\sum_{m\ge2}\sum_{p\le R^{1/m}}\log p\ll\sqrt R$ by Chebyshev's
bound. There are only $(\log X)^{O_F(1)}$ possible $\mathfrak k$.
Applying (20j) with $\sigma$ arbitrarily close to $\beta_\chi$ gives
(20i) for these terms. For the discarded pairs $p\mid N\mathfrak m$,
sum first over $p$ using $\sum_{p\mid N\mathfrak m}|b_\chi(p)|\log p
\le d\log N\mathfrak m$. Here $N\mathfrak k\,N\mathfrak m\le R$.
Absolute convergence of $\zeta_F(1+\epsilon)^2$ gives
$O_{F,\epsilon}(X^{1/2+\epsilon})$ for these pairs after absorbing
the logarithms. This proves (20i).

\leaf{520g}{Growth of the prime sum}{12.5.1 \quad THE RIGHTMOST ZERO}{The growth exponent of the degree-one sum}
For the sum over coprime pairs in (20h) one has the exact identity
\[
 \boxed{\displaystyle
  \beta_\chi=
   \max\left\{\frac12,\ \limsup_{X\to\infty}
                \frac{\log(1+|\mathcal B_\chi(X)|)}{\log X}\right\}.}
 \tag{20k}
\]
In particular an independent $O_{F,\chi,\epsilon}(X^{1/2+\epsilon})$
bound in (20h) for every $\epsilon>0$ would give the
$O_{F,\chi,v}(\sqrt X)$ remainder in (24) with the original weight
$v$. Both $n$ and $p$ in (20h) range only from $\sqrt X/L_X^{d+1}$
to $b\sqrt X L_X^{d+1}$.

Indeed the Mellin transform in $X$ of $A^\sharp_\chi$, integrated
against $X^{-s-1}\dd X$, is $-W_\chi(s)L'(s)/L(s)$ for $\Re s>1$.
The contribution from $0<X<1$ is entire since $A^\sharp_\chi$ vanishes
near zero in $X$. For each $\sigma>\beta_\chi$, move the inverse contour
to $\sigma$. The possible norm pole is removed by (20c). At large
heights the disk argument on \pg{520f} and Cauchy's estimate on a
slightly smaller disk bound $L'/L$ polynomially; rapid decay gives
$A^\sharp_\chi(X)=O_{F,\chi,\sigma}(X^\sigma)$. Conversely a growth
exponent strictly below $\beta_\chi$ would make its Mellin integral
holomorphic at a nontrivial zero $\rho$ to the right of that exponent.
The residue there is $-m_\rho W_\chi(\rho)\ne0$. This argument uses
the supremum of the real parts and does not require a rightmost zero.

If $\beta_\chi>1/2$, the exponent $1/4+\beta_\chi/2$ in (20i) is
strictly less than $\beta_\chi$. Hence $A^\sharp_\chi$ and
$\mathcal B_\chi$ have the same growth exponent. If $\beta_\chi=1/2$,
(20i) gives $\limsup\le1/2$ for $\mathcal B_\chi$. These two cases
prove (20k). The angular character of \pg{522c} gives, for odd $p$,
\[
 b_\chi(p)=
  \begin{cases}2\cos(4\arg\pi_p),&p=\pi_p\bar\pi_p\equiv1\pmod4,\\
                0,&p\equiv3\pmod4.
  \end{cases}
\]
Its coefficients $u_\chi(n)$ come from the Euler polynomials in
$z=p^{-s}$ of $1/L(s)$: $1-2\cos(4\arg\pi_p)z+z^2$ at split primes,
$1-z^2$ at inert primes, and $1+z$ at $2$.
The two factors in (20h) thus carry explicit arithmetic weights even
in this fixed degree-two example. Formula (20i) controls the discarded
terms. The square-root exponent for (20h) remains open in this argument.

\leaf{520h}{A one-sided prime estimate}{12.5.1 \quad A LOWER BOUND FOR A SIGNED MEAN}{The mean deficit suffices}
For each character component choose the unitary representative $\chi$
trivial on the positive uniform norm section. A norm twist translates
the zeros vertically, so these representatives suffice for all fields
and characters. In (20c) one now has $s_*=1$ throughout. The weight
$w$ is real and $W(x)>0$ for $1/2<x<1$. Write
\[
       D_\chi(X)=\int_X^{2X}\Re\mathcal B_\chi(Y)\frac{\dd Y}{Y},
       \qquad [r]_+=\max(r,0).
\]
With the same $\beta_\chi$ as in (20k), one has
\[
 \beta_\chi=\max\left\{\frac12,\ \limsup_{X\to\infty}
       \frac{\log(1+[-D_\chi(X)]_+)}{\log X}\right\}.
 \tag{20l}
\]
Thus the single lower bound $D_\chi(X)\ge-C_{F,\chi,\epsilon}
X^{1/2+\epsilon}$ for each $\epsilon>0$ already implies (24),
including its original $O(\sqrt X)$ error. The two sums in (20h)
and the average in $Y$ may all cancel before this lower bound is taken.

Let $f(X)=\int_X^{2X}\Re A^\sharp_\chi(Y)\dd Y/Y$. On $\Re s>1$ its
Mellin integral against $X^{-s-1}\dd X$ on $(0,\infty)$ equals
\[
 -\frac{2^s-1}{2s}\,W(s)
 \left(\frac{L'_\chi}{L_\chi}(s)+\frac{L'_{\bar\chi}}{L_{\bar\chi}}(s)\right).
\]
The integral below $1$ is entire. At any nontrivial zero the orders of
the two functions add, and the factor $(2^s-1)W(s)/s$ is nonzero.
At a real zero in $(1/2,1)$ the residue is strictly negative. The
possible pole of a norm character at $1$ is removed.

The upper bound for the limsup in (20l) follows from (20k). If
$\beta_\chi>1/2$ and the limsup were smaller, choose $\theta<\beta_\chi$
strictly above both that limsup and $1/4+\beta_\chi/2$. By (20i),
$g(X)=f(X)+CX^\theta\ge0$ for $X\ge X_0$ and some $C,X_0$.
Its Mellin integral there has real abscissa of convergence $\beta_\chi$:
the bounds on \pg{520g} give the upper inequality, and the nonzero poles
with real parts arbitrarily close to $\beta_\chi$ give the lower one.
The Mellin transform of a nonnegative function is singular at its real
abscissa. Indeed a Taylor disk across that point, centered just to its
right, would give convergence of the integral to its left by expansion
of $X^t$ in nonnegative terms. Here the continuation is meromorphic.
At the real abscissa it is either regular or has a simple pole with
negative residue, which is impossible for an integral positive to its
right. This proves (20l), also when no zero attains $\beta_\chi$.

\leaf{520i}{The absolute prime contribution}{12.5.1 \quad TWO DEGREE-ONE PRIME FACTORS}{How much cancellation the estimate requires}
There is a fixed interval $I=[a,b']\subset(0,\infty)$ of positive
length such that, for every normalized $\chi$ on \pg{520h}, the sum of
absolute values in (20h) from just $n=q$, $p\ne q$ rational primes in
$I\sqrt X$, is $\gg_{F,\chi} X/\log X$. The following unconditional
prime correlation proves this even for characters of infinite order:
\[
   \sum_{p\le x}|b_\chi(p)|^2\log p=(r_\chi+o(1))x,\qquad r_\chi\ge1 .
 \tag{20m}
\]
Write $F\otimes_\Q F=\prod_j K_j$. The two embeddings $F\to K_j$
give norm maps $N_{1j},N_{2j}$ to $F$. Let
$\eta_j=(\chi\circ N_{1j})(\bar\chi\circ N_{2j})$.
Then $r_\chi$ is the number of trivial $\eta_j$, including the
factor $K_j=F$ from multiplication. Outside finitely many primes,
degree-one primes in this product correspond to ordered pairs of
degree-one primes of $F$ above $p$. Indeed maps from the tensor
product of the two reduced residue algebras to $\mathbb F_p$ are
pairs of maps of its factors. Each norm of the resulting prime ideal
is the corresponding prime of $F$. Hence
$|b_\chi(p)|^2=\sum_j b_{\eta_j,K_j}(p)$. Both norm maps preserve the
uniform section, so a nontrivial $\eta_j$ cannot be a pure norm twist.

The fixed-character prime ideal theorem now gives (20m). Explicitly,
apply Wiener--Ikehara to the nonnegative coefficients
$(1+\Re\eta_j)\Lambda_{K_j}$ and $(1+\Im\eta_j)\Lambda_{K_j}$
for nontrivial $\eta_j$. Their Dirichlet series have the same simple
pole at one as $-\zeta'_{K_j}/\zeta_{K_j}$ by Hecke nonvanishing on
$\Re s=1$. Subtraction of the untwisted prime ideal theorem gives
$o(x)$ for the twisted sum. For the trivial character the main term is
$x$. Prime norms of residue degree above one and higher prime powers
cost $O_{K_j}(\sqrt x)$ as on \pg{520f}.

Choose $I$ so that its product interval $[a^2,(b')^2]$ lies where
$|w|\ge c>0$ in $(e^{5/4},e^{15/8})$. Such an interval exists by
(20c). The ratio of these support endpoints is $e^{5/8}<2$.
Consequently $K_\chi(X/(pq))=w(pq/X)$ on this square: only the unit
ideal contributes. Both lower cutoffs are one for large $X$ and
$u_\chi(q)=-b_\chi(q)$. By (20m) and $|b_\chi(p)|\le[F:\Q]$,
$\sum_{p\in I\sqrt X}|b_\chi(p)|\log p\gg_{F,\chi}\sqrt X$.
The corresponding unweighted sum is $\gg_{F,\chi}\sqrt X/\log X$.
Their product gives the asserted lower bound after subtraction of the
diagonal, which costs $O_F(\sqrt X)$. It holds throughout every large
dyadic interval of $X$.

\leaf{520j}{Cancellation between primes}{12.5.1 \quad THE SIGNED PRIME-PAIR SUM}{The contribution with two prime factors}
Write $\mathcal P_\chi(X)$ for the part of (20h) with $n=q$ a rational
prime. For any nontrivial normalized $\chi$ one has
\[
    |\mathcal P_\chi(X)|\ll_{F,\chi}
       \sqrt X+X^{\beta_\chi}\frac{\log\log(3X)}{\log(2X)} .
 \tag{20n}
\]
In particular if $\beta_\chi>1/2$ this entire signed contribution is
$o(X^{\beta_\chi})$. The estimate on \pg{520i} applies to its absolute
mass on a single square.

For a smooth test $\phi$ supported in $[1/2,2]$, at scale $Z\ge4$,
\[
   \left|\sum_p b_\chi(p)\log p\,\phi(p/Z)\right|
                \ll_{F,\chi} Z^{\beta_\chi}\|\phi\|_{C^m}
 \tag{20o}
\]
for a fixed integer $m$. Apply Mellin inversion to
$-\widehat\phi(s)L'_\chi(s)/L_\chi(s)$, moving to $\Re s=-1/4$.
The nontrivial zero residues are bounded by $Z^{\beta_\chi}$
times the absolutely summable Mellin values. The finitely many trivial
zeros crossed lie on $\Re s=0$. The gamma factors of a unitary Hecke
character have trivial-zero real parts in the nonpositive half-integers,
so the new line stays away from their poles. There the functional
equation and the Euler product on $\Re s=5/4$ bound the logarithmic
derivative by $O_{F,\chi}(\log(2+|t|))$. The horizontal sides may be
chosen as on \pg{522a}; taking $m$ large covers their polynomial bounds.
There is no norm pole. Removal of prime powers and higher residue
degrees costs $O_F(\sqrt Z)\|\phi\|_\infty$. This proves (20o).

Insert smooth dyadic partitions $p\asymp P,\ q\asymp Q$ in (20h),
temporarily including $p=q$. Put $H=X/(PQ)$. On the fixed rescaled
square, all derivatives of the kernel $K_\chi(H/(uv))$ are
$O_{j,A}((1+H)^{-A})$ for each $A$. This follows by multiplying its
Mellin integrand $W(s)L_\chi(s)$ by powers of $s$ and moving to a
fixed far left line as in (20e). The two cutoffs $\omega$ have uniform
rescaled derivatives. Division by $\log(Qv)$ costs $O_m(1/\log Q)$
in every required seminorm. Applying (20o) successively in both
variables bounds each box by
$O((PQ)^{\beta_\chi}(1+H)^{-2}/\log Q)$: differentiate the inner
smooth sum in the outer variable before its second application.
Nonempty boxes have $H\gg1$, $\log Q\asymp\log X$. For each $P$,
summing the dyadically spaced $H$ costs a constant after extracting
$X^{\beta_\chi}$. There are $O(\log\log(3X))$ choices of $P$.
The subtracted diagonal costs $O_{F,\chi}(\sqrt X)$ since $p^2\le bX$ and
$K_\chi$ is bounded. This proves (20n).

\leaf{520k}{The untwisted prime pairs}{12.5.1 \quad THE PRINCIPAL CHARACTER}{Two terms of order \(X/\log X\) cancel}
For $\chi=1$ put $L=\log(X/U)$ and $c_F=V_*(1)
\operatorname*{Res}_{s=1}\zeta_F(s)>0$, where $V_*$ is the Mellin
transform of $v_*$ on \pg{520c}. Then
\[
 \mathcal P_1(X)=-c_F X/L+O_F(X/L^2),\qquad
 \mathcal B_1(X)-\mathcal P_1(X)=c_F X/L+O_F(X/L^2).
 \tag{20p}
\]
Here $L=\tfrac12\log X+(d+1)\log\log(2X)$. These two contributions
are larger than the square-root precision in (24).

By the fixed-field prime ideal theorem with its
\citetitle{https://aareyanmanzoor.github.io/assets/articles/lagarias-odlyzko.pdf}{classical remainder},
the measure $\sum_p b_1(p)\log p\,\delta_p$ on either dyadic interval of
\pg{520j} can be replaced by Lebesgue measure with smooth error
$O_{F,A}(Z(\log Z)^{-A}\|\phi\|_{C^1})$ for every $A$.
Any fixed exceptional real zero contributes a smaller error at this
precision. Partial summation and the $O_F(\sqrt Z)$ bound for
higher degrees give the displayed estimate. Replace both prime
measures on each box. Their masses are $O(P),O(Q)$, the rescaled
kernel derivatives are bounded, and there are $O((\log\log X)^2)$
boxes with $PQ\ll X$. The errors, including the diagonal, are
$O_{F,A}(X(\log X)^{-A})$ for every fixed $A$.

Set $G(z)=K_1(1/z)$, extended by zero for $z>b$. Changing variables
$z=pq/X$ in the resulting double integral gives
\[
    -X\int_0^b G(z)J_X(z)\dd z,\qquad
    J_X(z)=\int_U^\infty
      \omega(t/U)\,\omega(Xz/(tU))\,\frac{\dd t}{t\log t}.
\]
The Mellin transform of $G$ is $W(s)\zeta_F(s)$ by (20e) and
continuation. Its zero at $1$ is simple, so
$\int G(z)\dd z=0$ and $\int G(z)\log z\dd z=c_F$.
For $r=d+1$ and $L_X^{-r}\le z\le b$, the upper transition occurs wholly above
$2U$. Substitution $h=Xz/(tU)$ then gives
\[
     \partial_{\log z}J_X(z)
       =\int_1^2\frac{\omega'(h)}{L+\log z-\log h}\dd h
       =L^{-1}+O((1+|\log z|)/L^2).
\]
Integrate from $1$ to $z$. The constant $J_X(1)$ drops out; the part
$z<L_X^{-r}$ is negligible to any inverse log power since
$G(z)=O_A(z^A)$ there and $J_X$ is bounded. The two Mellin moments
now give the first assertion. The same prime ideal theorem yields
$A^\sharp_1(X)=O_{F,A}(X(\log X)^{-A})$ since $W(1)=0$.
Using (20i) with $\beta_1\le1$ gives the second assertion.

\leaf{520l}{A fixed number of prime factors}{12.5.1 \quad THE M\"OBIUS COEFFICIENTS}{The terms of each fixed order}
Write $\Omega(n)$ for the number of rational prime factors counted
with multiplicity. Let $\chi$ be nontrivial and normalized as on
\pg{520h}. Suppose $\beta=\beta_\chi>1/2$. For a fixed $j\ge1$ and
$\phi\in C_c^\infty([1/2,2])$, one has, for some integer $m_j$,
\[
 \left|\sum_{\Omega(n)=j} u_\chi(n)\phi(n/Z)\right|
 \ll_{F,\chi,j} Z^\beta
        \frac{(1+\log\log Z)^j}{\log Z}\|\phi\|_{C^{m_j}}
        \qquad(Z\ge4).
 \tag{20q}
\]
The extra logarithmic saving comes from the singularity of $\log L$
in the series with a marked number of prime factors.

Put $L=L_f(\,\cdot\,,\chi)$. For $\Re s>1$ and $|v|$ small write
\[
 \begin{split}
  \sum_{j\ge0}v^j U_j(s)
   &=\prod_p\left(1+\sum_{k=1}^d u_\chi(p^k)v^k p^{-ks}\right)
     =e^{-v\log L(s)}G(s,v),\\
  U_j(s)&=\sum_{\Omega(n)=j} u_\chi(n)n^{-s}.
 \end{split}
\]
The local factor of $G$ is obtained by multiplying the displayed
polynomial by $L_p(s)^v$ with its Euler logarithm. The linear terms in
$p^{-s}$ cancel since $u_\chi(p)=-b_\chi(p)$, and the higher coefficients
of the degree-$d$ Euler polynomial are bounded in terms of $d$.
Consequently the resulting product and its first $s$-derivative converge
uniformly on $\Re s\ge\beta$, $|v|\le c_d$, by
$\sum_p(1+\log p)p^{-2\beta}<\infty$. Thus $U_j$ extends to
$\Re s>\beta$ as a polynomial of degree at most $j$ in the Euler branch
of $\log L(s)$, with bounded analytic coefficients and derivatives.

Set $\sigma=\beta+1/\log Z$. For $t\in[h,h+1]$, $h\in\Z$, one has
\[
 \sup_t|\log L(\sigma+it)|
  +\int_h^{h+1}|L'/L(\sigma+it)|\dd t
       \ll_{F,\chi}(2+|h|)^C(1+\log\log Z).
 \tag{20r}
\]
Indeed, subtract the Hadamard logarithmic derivative at $2+it$ from
its value at $u+it$, $\sigma\le u\le2$. Zeros with
$|\Im\rho-t|>1$ give a polynomially bounded difference by zero
counting and summation of the inverse-square tail. The gamma terms
and the Euler logarithmic derivative at $2+it$ have polynomial bounds
as well. A remaining simple fraction integrates in $u$ to a logarithm
of size $O(1+\log\log Z)$, since $\sigma-\Re\rho\ge1/\log Z$.
Its absolute integral in $t$ over a unit interval has the same bound.
Count the nearby zeros with multiplicity to get (20r).

\leaf{520m}{A fixed number of prime factors}{12.5.1 \quad THE REMAINING COMPOSITE INDICES}{Which terms can carry a rightmost pole}
From (20r) and the polynomial representation of $U_j$ one obtains
\[
  \int_h^{h+1}(|U_j|+|U_j'|)(\sigma+it)\dd t
       \ll_{F,\chi,j}(2+|h|)^{C_j}(1+\log\log Z)^j .
\]
The derivative of a power of $\log L$ loses one such factor and gains
one integrated $L'/L$. Shift the Mellin integral for the left side of
(20q) to $\sigma$ and integrate $e^{it\log Z}$ by parts once. Both the
Mellin transform of $\phi$ and its first derivative decay faster than
$(2+|t|)^{-C_j-2}$ at cost of a fixed $C^{m_j}$ seminorm. The unit
interval estimate gives (20q) since $Z^\sigma=eZ^\beta$.

Let $\mathcal B_\chi^{\le J}$ be (20h) restricted by $\Omega(n)\le J$,
for a fixed positive integer $J$. Then for $\beta_\chi>1/2$,
\[
   \big|\mathcal B_\chi^{\le J}(X)\big|\ll_{F,\chi,J}
      X^{\beta_\chi}\frac{(1+\log\log(3X))^{J+1}}{\log X}
          =o(X^{\beta_\chi}).
 \tag{20s}
\]
Use the dyadic squares of \pg{520j}. Apply (20o) to the prime variable
and (20q) to the integer variable. The derivative bounds for the
kernel there hold to every fixed order, including the mixed
derivatives. The sum in $H=X/(PQ)$ is bounded and the number of outer
scales contributes one factor $1+\log\log(3X)$. The absolute cost of
the pairs $p\mid n$ is $O_{F,\epsilon}(X^{1/2+\epsilon})$ by
\pg{520f}, which is absorbed for $\epsilon<\beta_\chi-1/2$.

If a zero $\rho=\beta_\chi+i\gamma$ attains this supremum, then for
every fixed $J$ the indices $\Omega(n)>J$ must still contribute a
negative dyadic mean of order $X^{\beta_\chi}$ along some unbounded
sequence. To see this let $f$ be the signed mean on \pg{520h}.
By (20i) and (20s) it differs from the mean of the remaining terms by
$o(X^{\beta_\chi})$. If $[-f(X)]_+=o(X^{\beta_\chi})$, then
$t\int_1^\infty[-f(X)]_+X^{-1-\beta_\chi-t}\dd X\to0$ as $t\downarrow0$.
At the real point $\beta_\chi$ the corresponding limit for $f$ is
either zero or a negative residue. Since $[f]_+=f+[-f]_+$, positivity
excludes the negative value and forces the limit for $|f|$ to be zero.
This bounds the absolute value of the same Abel limit with factor
$X^{-i\gamma}$ by zero, contradicting the nonzero residue at $\rho$
on \pg{520h}. If the supremum is not attained, (20s) retains its stated
upper bound and (20l) applies to the entire mean.

\leaf{520n}{Integers with many prime factors}{12.5.1 \quad A RESTRICTION IN THE NONTRIVIAL COMPONENTS}{Discarding the most highly factored indices}
For every nontrivial normalized $\chi$, let $\widetilde{\mathcal B}_\chi$
be (20h) with the additional condition
\[
             \Omega(n)\le \frac{\log n}{2\log\log n}.
\]
Only large $n$ occur as $X\to\infty$. Then
\[
     \mathcal B_\chi(X)-\widetilde{\mathcal B}_\chi(X)
             \ll_{F,\chi,\epsilon} X^{1/4+\beta_\chi/2+\epsilon}.
 \tag{20t}
\]
In particular (20l) holds with $\widetilde{\mathcal B}_\chi$ throughout.
This restriction is independent of the unknown zero locations.

To prove (20t), for $0<c<1$ first bound the sum of $|u_\chi(n)|$ over
$Q/2\le n\le2Q$ with $\Omega(n)>c\log n/\log\log n$ by
$O_{F,c,\epsilon}(Q^{1-c+\epsilon})$. The reciprocal Euler polynomial
at $p$ has degree at most $d=[F:\Q]$ with all roots of modulus one:
factor each $1-a_{\mathfrak p}z^{f_{\mathfrak p}}$ into linear factors.
Thus $|u_\chi(p^j)|\le\binom dj$. For $\eta>0$ fixed and
$z=(\log Q)^{1-\eta}$, positivity and the Euler product give
\[
   \sum_{Q/2\le n\le2Q}|u_\chi(n)|z^{\Omega(n)}
     \le (2Q)^{1+\eta}\prod_p(1+zp^{-1-\eta})^d
     \ll_{F,\eta}Q^{1+\eta}\exp(C_{F,\eta}z).
\]
On the specified indices $z^{\Omega(n)}\ge
Q^{c(1-\eta)+o(1)}$. Choose $\eta$ sufficiently small; $z=o(\log Q)$
then gives the assertion.

Use $c=1/2$ on a box $n\asymp Q$, $p\asymp P$. At fixed $n$, (20o)
applied to the smooth prime sum costs $O_{F,\chi}(P^{\beta_\chi})$.
Indeed the rescaled derivatives of $K_\chi$ and the cutoffs are bounded
as on \pg{520j}. The absolute sum of the remaining $n$-coefficients
costs $O_{F,\epsilon}(Q^{1/2+\epsilon})$. Both scales are within fixed
logarithmic factors of $\sqrt X$, and there are only logarithmically
many boxes. Removing $p\mid n$ costs $O_{F,\epsilon}(X^{1/2+\epsilon})$
as on \pg{520f}. This proves (20t). If $\beta_\chi>1/2$ its exponent
before $\epsilon$ is strictly smaller than $\beta_\chi$, so the proof
of (20l) applies unchanged. At $\beta_\chi=1/2$ it gives the required
upper bound for the negative part. Page \pagelink{522i}{522i} explains
why these nontrivial components and ordinary $\zeta$ suffice.

\sourcepages{521}{521}{The one-sided estimate is on page 520h and the reduced range of indices on page 520n.}
\leaf{522}{One prime sum}{13.3 \quad THE RIEMANN ZETA FUNCTION}{One smooth weight detects every zero}
There is a fixed nonnegative $v\in C_c^\infty((0,\infty))$, chosen
independently of the zeros, for which
\[
 \mathrm{RH}\quad\Longleftrightarrow\quad
 \sum_{n\ge2}\Lambda(n)v(n/X)
       =X V(1)+O_v(\sqrt X)\quad(X\ge1),\qquad
 V(s)=\int_0^\infty v(u)u^{s-1}\dd u .
 \tag{21}
\]
In particular $V(1)=\int v(u)\dd u$.

\point{Construction} Set $\ell_j=2^{-j-1}$, $j\ge1$, and let $b$ be
the density of $5/4+\sum_{j\ge1} U_j$ for independent random variables
$U_j$ uniform on $[0,\ell_j]$. The sum converges and is supported on
$[5/4,7/4]$. Its transform is
\[
    V(s)=e^{5s/4}\prod_{j\ge1}\frac{e^{\ell_js}-1}{\ell_js},
       \qquad v(u)=b(\log u).
 \tag{22}
\]
On the imaginary axis the first $k$ factors give
$O_k((1+|t|)^{-k})$ and every remaining factor has modulus at most one.
Fourier inversion therefore gives a smooth compactly supported density.
The factors are $1+O_K(\ell_j)$ uniformly on compact $K\subset\C$ and
none vanishes for $\Re s>0$. The product has no zeros there. For
$h(x)=e^{x/2}b(x)$ the transform in (2) is $V(z+\tfrac12)$.

\point{Forward implication} Apply (17) with this $h_\tau$, $X=e^\tau\ge1$.
Expansion of the real-place kernel on positive support gives
\[
 \sum_n\Lambda(n)v(n/X)=XV(1)
   -\sum_\rho m_\rho X^\rho V(\rho)
   -\sum_{k\ge1}X^{-2k}V(-2k).
 \tag{23}
\]
Here $\rho$ runs over distinct nontrivial zeros in uncentered coordinates.
The last series is $O_v(X^{-2})$ since the support is above $1$.
Under RH the absolute value of the zero sum is bounded by
$\sqrt X\sum_\rho m_\rho |V(\rho)|=O_v(\sqrt X)$.

\point{Converse} Denote the prime sum by $A(X)$. For $\Re s>1$ its
Mellin integral against $X^{-s-1}\dd X$ over $(0,\infty)$ is
$-V(s)\zeta'(s)/\zeta(s)$ by absolute convergence. The part on $(0,1)$
is entire since $A$ vanishes near zero. The estimate in (21) thus makes
$-V(s)\zeta'(s)/\zeta(s)-V(1)/(s-1)$ holomorphic for $\Re s>\tfrac12$.
A zero there would have residue $-m_\rho V(\rho)\ne0$. Reflection
excludes zeros on the other side as well.

\leaf{522a}{Hecke prime sums}{13.3 \quad ALL NUMBER FIELDS AND UNITARY HECKE CHARACTERS}{The corresponding Hecke criterion}
Use the same $v,V$ as in (22). For a number field $F$ and a unitary
idele class character $\chi$, take its primitive finite Euler product
$L_f(s,\chi)$. For unramified prime ideals $\mathfrak p$ put
$a_{\mathfrak p}=\chi_{\mathfrak p}(\varpi_{\mathfrak p})$, so the local
factor is $(1-a_{\mathfrak p}(N\mathfrak p)^{-s})^{-1}$. Set
\[
 A_\chi(X)=\sum_{\mathfrak p\nmid\mathfrak f_\chi}\sum_{k\ge1}
    (\log N\mathfrak p)\,a_{\mathfrak p}^{\,k}\,
           v\big((N\mathfrak p)^k/X\big).
\]
Let $M_\chi(X)=0$ unless $\chi=|\cdot|_\A^{it_0}$. For such a norm twist
put $s_0=1-it_0$ and $M_\chi(X)=V(s_0)X^{s_0}$.
Then GRH for this Hecke function is equivalent to
\[
             A_\chi(X)=M_\chi(X)+O_{F,\chi,v}(\sqrt X),\qquad X\ge1.
 \tag{24}
\]
Taking (24) for all $F,\chi$ gives exactly the Hecke GRH of the manuscript.
The constant may depend on the individual character and field.

\point{Proof} In $\Re s>1$ the Mellin integral of $A_\chi$ is
$-V(s)L_f'(s,\chi)/L_f(s,\chi)$ by its Euler product. Subtract the
Mellin integral of $M_\chi$ on $[1,\infty)$, namely $V(s_0)/(s-s_0)$
in the norm-twist case, and zero otherwise. As on page 522, (24) makes
the difference holomorphic on $\Re s>\tfrac12$. Since $V$ is nonzero
there, $L_f$ has no zero in this half-plane. The completed functional
equation pairs $\rho$ with $1-\overline\rho$ for fixed unitary $\chi$
(conjugation changes $\chi$ to $\bar\chi$). This proves the converse.

If GRH holds, shift the inverse Mellin integral from $\Re s=2$ to
$\Re s=\tfrac14$. The residues crossed are the main term and
$-\sum_\rho m_\rho X^\rho V(\rho)$. The archimedean gamma poles have
real part at most zero for unitary Hecke characters, so none intervenes.
On the new line the logarithmic derivative is polynomially bounded.
The line is a fixed distance from the nontrivial zeros,
its near terms are bounded using $O_{F,\chi}(T\log(2+T))$, and the
regularized far terms are summable. The logarithmic gamma derivatives
have at most logarithmic growth there. In $[T,T+1]$ the excluded intervals
of radius $T^{-3}$ about the nearby zero ordinates have total length
$O(T^{-2}\log T)$, and likewise at negative heights. On the resulting
horizontal sides the Hadamard logarithmic derivative grows polynomially.
Rapid decay of $V$ justifies the shift and leaves a vertical
integral $O_{F,\chi,v}(X^{1/4})$. The zero residues are absolutely
$O_{F,\chi,v}(\sqrt X)$. This proves (24).

\leaf{522b}{Abelian subfamilies}{13.3 \quad TRANSFER FROM DIRICHLET CHARACTERS}{The reach of a rational zero theorem}
Suppose the zero-location statement for primitive Dirichlet $L$-functions
over $\Q$ is supplied independently. Then (24) follows for each unitary
Hecke character of $\Q$, including all norm twists. Indeed
$C_\Q\simeq\R_{>0}\times\widehat{\Z}^{\,\times}$, and a continuous
character of the profinite factor has finite image and factors through
$(\Z/q\Z)^\times$ for some $q$. A norm twist translates the zeros
vertically. In this family (18) has $E_\alpha=0$ for every positive
width, as a consequence of the external zero theorem.

\point{Transfer to an abelian field} Let $F/\Q$ be abelian, with group
$G$, and let $\psi$ be a primitive Dirichlet character interpreted as
a Hecke character of $\Q$. For $\chi=\psi\circ N_{F/\Q}$, Artin
formalism gives the primitive Euler products
\[
 L_{F,f}(s,\chi)=
       \prod_{\eta\in\widehat G} L_{\Q,f}(s,\psi\eta).
\]
On the right each factor uses its primitive inducing character.
Indeed $\operatorname{Ind}_{G_F}^{G_\Q}(\psi|_{G_F})$ is $\psi$ tensored
with the regular representation of $G$. The latter splits as the sum of
the characters $\eta$. Taking the local Artin factors on inertia
invariants gives the displayed identity also at ramified primes.
Every zero in the open critical strip is thus a zero of a Dirichlet
factor. This transfers (24) to these norm pullbacks and their norm twists.

More generally, a finite-order character of a field $K$ cuts out an
abelian extension $E/K$ by class field theory. The factorization
$\zeta_E(s)=\prod_{\eta\in\widehat{\operatorname{Gal}(E/K)}} L_{K,f}(s,\eta)$
shows that RH for $\zeta_E$ suffices for that character. Every factor is
holomorphic in the open critical strip, so there is no cancellation of
a zero by a pole there. This transfers a Dirichlet zero theorem whenever
$E/\Q$ is abelian.

\point{External statement} The November 2026 announcement of L.~Alp\"oge
and A.~Sutherland's verification record specify RH and zero location for
Dirichlet $L$-functions. They provide an external input of the form used
above. The transfer identifies the subfamily to which that input applies
in the present notation. The proof of the spectral trace and its strip
defect is independent of this input.
\citetitle{https://alpo.ge/remark.pdf}{Announcement.}\quad
\citetitle{https://math.mit.edu/~drew/grh/}{Verification record and theorem statements.}

\leaf{522c}{Gaussian primes}{13.3 \quad A CHARACTER OF INFINITE ORDER}{An explicit remaining prime estimate}
In $K=\Q(i)$ every ideal is principal. The rule
\[
 \chi((\gamma))=(\gamma/|\gamma|)^4
\]
is independent of the generator since the units are fourth roots of
unity. It is an unramified finite-place Hecke character with nontrivial
infinite type: on principal ideles the archimedean character cancels
this angular factor. At the two primes over $5$ its values are
$(-7+24i)/25$ and $(-7-24i)/25$, obtained from $2+i$ and $2-i$.
The ideals have equal norm, while the character values differ. The first
value is not a root of unity in $\Q(i)$, so $\chi$ has infinite order.
A pullback of a rational ideal character would have equal values at
both primes.

For each rational prime $p\equiv1\bmod4$, choose a Gaussian prime
$\pi_p$ above it and put $\theta_p=\arg\pi_p$. Multiplication by a unit
and complex conjugation both preserve $\cos(4\theta_p)$. For the single
weight (22), the zero-location assertion for $\chi$ is equivalent to
\[
 \boxed{\displaystyle
 \sum_{p\equiv1\;(\mathrm{mod}\,4)}(\log p)\cos(4\theta_p)\,v(p/X)
                   =O_v(\sqrt X).}
 \tag{25}
\]
Indeed the two prime ideals over a split prime contribute
$2(\log p)\cos(4\theta_p)v(p/X)$ to the first-power part of $A_\chi(X)$.
An inert prime has norm $p^2$, and the unique prime over $2$ contributes
only a bounded first-power term. If $\supp v\subset[1,B]$, the inert
primes and all second powers contribute $O_v(\sqrt X)$ absolutely by the
elementary Chebyshev bound $\sum_{p\le y}\log p\ll y$. Powers of exponent
at least three contribute
$O_v(X^{1/3}\log(2+X))=O_v(\sqrt X)$, since the number of exponents
is $O(\log(2+X))$ and there are at most two prime ideals over each $p$.
Thus $A_\chi(X)$ is twice the sum in (25) up to $O_v(\sqrt X)$.
Its main term in (24) is zero. This proves the equivalence.

In (25) the two conjugate prime ideals contribute twice the real part
of a Hecke value before division by two. A count
of the rational primes in the progression $1\bmod4$ alone supplies
their total number, while (25) asks for cancellation of these angular
weights. This one character is already required by the manuscript's
general zero-location statement.

\leaf{522d}{Gaussian primes}{13.3 \quad POISSON SUMMATION OVER GAUSSIAN IDEALS}{The long divisors that remain}
For the Gaussian character on page 522c define the sum over all nonzero
integral ideals
\[
 B(Y)=\sum_{\mathfrak a}\chi(\mathfrak a)\log N\mathfrak a\,
                    v(N\mathfrak a/Y).
\]
For every $A>0$ it satisfies $B(Y)=O_{A,v}(Y^{-A})$ for $Y\ge1$.
To see this, set $f(z)=(z/|z|)^4v(|z|^2)$ and
$g(z)=f(z)\log|z|^2$, extended by zero near the origin. Both are smooth
and compactly supported on $\R^2$, and both integrals vanish by angular
integration. There are exactly four generators per ideal. Poisson
summation on $\Z[i]\simeq\Z^2$ gives, with the usual $2\pi$ convention,
\[
 B(Y)=\frac{Y}{4}\sum_{m\in\Z^2\setminus\{0\}}
       \big((\log Y)\widehat f(\sqrt Y\,m)+\widehat g(\sqrt Y\,m)\big).
\]
Rapid decay of both transforms proves the bound with any prescribed $A$.

\point{Isolating prime powers} Let $\mu_K$ be the ideal M\"obius function.
Dirichlet convolution of $\mu_K\chi$ with $\chi\log N$ gives the prime
power coefficients of $-L'_f/L_f$. Compact support makes the resulting
identity finite for each $X$:
\[
 A_\chi(X)=\sum_{N\mathfrak d\le B_0 X}
          \mu_K(\mathfrak d)\chi(\mathfrak d)
                       B(X/N\mathfrak d),\qquad B_0=e^{7/4}.
\]
There are $O(t)$ Gaussian ideals of norm at most $t$, by counting lattice
points in a disk. Fix $0<\delta<1$. For $N\mathfrak d\le X^{1-\delta}$
one has $X/N\mathfrak d\ge X^\delta$. The total absolute contribution
of this range is $O_{A,v}(X^{1-\delta-A\delta})$, hence smaller than
any prescribed inverse power for fixed $\delta$.

In particular (25) is equivalent, at any one such fixed $\delta$, to
\[
 \left|\sum_{X^{1-\delta}<N\mathfrak d\le B_0X}
       \mu_K(\mathfrak d)\chi(\mathfrak d)\,B(X/N\mathfrak d)\right|
       =O_{\delta,v}(\sqrt X).
 \tag{26}
\]
The Poisson estimate removes the shorter divisors unconditionally.
When $N\mathfrak d$ is comparable with $X$, the inner sum samples a
bounded range of ideal norms. Cancellation among the M\"obius-weighted
long divisors is still required.

\leaf{522e}{Conclusions}{13.3 \quad SCOPE OF THE ZERO-LOCATION CLAIM}{The remaining arithmetic statement}
The arithmetic quotient, the trace of its convolution, and the full-spectrum
explicit formula are compatible with zeros anywhere in the critical strip.
At a narrower width (18) gives the sign and multiplicity of every
missing zero. To locate all zeros it suffices to establish (21) for
ordinary $\zeta$ and (32) for the conductor-one characters of \pg{522k}.
Both factors in (32) are within logarithmic factors of $\sqrt X$.
Only a lower bound for their signed dyadic mean is required. In
nontrivial components (20n) bounds the prime pairs and (20t) removes
the most highly factored integer indices. In the trivial component
(20p) computes the opposite terms of order $X/\log X$ from prime
pairs and composite indices. Equation (35) uses only principal prime
generators. Equations (40)--(42) give a square-root term after Gaussian
averaging and oscillations for the cubic splitting count. Equations
(43)--(49) control two translates at every large separation when the
prime interval has length of order $\sqrt X\log X$. Reconstructing the
fixed test in (46) costs the reciprocal log width. An independent
bound at widths $X^{-o(1)}$ in (47) would give zero location.
The fixed-test ratio in (52) attains its lower limit. Its positive
part satisfies an averaged upper bound at frequencies of order $X$
by (54). The bound in (55) with $U=o(\log X)$ remains open.

The fixed-place calculation on page 495 fails over $\Q$. The two-prime
obstruction on page 495b rules out a fixed norm-translation-invariant
lift. The local quotient trace (13d) recovers the prime distributions.
Its spatially truncated layers coincide by (13b). At a common source
cutoff the full chain of prime projections converges as in (13f.2).
The order correction
(13g) can be nonzero, but (13i) removes it from the iterated scalar
trace and (13j) bounds its translated contribution. The real-place
lower-cutoff term (17a) persists in the balanced
regime. The external half-planes on page 520b give (20b). Reaching (19)
for widths tending to zero still requires additional cancellation.
Consequently the narrow-strip main theorem and the deduction of GRH remain
conditional. Equations (12), (18), and (23) use actual zero orders.
For $F=\Q$, (13.3.1.1) would require simplicity at every noncentral
zero. For $F=\Q(i)$ its cap of two gives a false trace identity for
a permitted section and tests satisfying all of (13.3.1.2), as proved
on \pagelink{522g}{pages 522g--h}. The uncapped location criterion (24)
is independent of that choice.

\point{Sources for the classical inputs}
\begin{itemize}\setlength{\itemsep}{3pt}
\item A.\,M. Masullo, \citetitle{https://math216.wordpress.com/wp-content/uploads/2026/11/the-music-of-the-primes.pdf}{\emph{The Music of the Primes}},
  \S\S9.2--13.3 and Appendix B. Original folios and equation numbers
  are used throughout.
\item E.\,C. Titchmarsh, \emph{The Theory of the Riemann Zeta-function},
  2nd ed., revised by D.\,R. Heath-Brown, Oxford, 1986, Chapters II, III
  and IX. Functional equation, Hadamard product, zero-free region and
  zero count.
\item J. Tate, \emph{Fourier analysis in number fields and Hecke's
  zeta-functions}, in Cassels and Fr\"ohlich, \emph{Algebraic Number Theory},
  1967, pp.~305--347. Local zeta integrals, Poisson summation and the Hecke
  functional equation.
\item R. Meyer, \citetitle{https://arxiv.org/abs/math/0311468}{\emph{On a representation of the idele class group related to primes and zeros of $L$-functions}} and
  \citetitle{https://arxiv.org/abs/math/0412277}{\emph{A spectral interpretation for the zeros of the Riemann zeta function}}.
  The full Weil distribution and its rational specialization.
\item M.\,J. Mossinghoff, T.\,S. Trudgian and A. Yang,
  \citetitle{https://doi.org/10.1007/s40993-023-00498-y}{\emph{Explicit zero-free regions for the Riemann zeta-function}},
  Theorem 1.3. The classical region suffices for the prime error on page 520a.
\end{itemize}

\leaf{522f}{The admissible tests}{13.3.1 \quad THE MOMENTS IN (13.3.1.2)}{What the Chapter 13 identity would imply}
Put $d=[F:\Q]$, $K=C_F^1$ and $x=\log |g|_F$. A continuous
homomorphic norm section $\tau$ gives coordinates $g=k\tau(e^x)$.
Write $\widetilde\phi(k,x)=\phi(k\tau(e^x))$. For a unitary character
$\chi$, write $\kappa=\chi|_K$ and $\chi(\tau(e^x))=e^{i x_0x}$.
At the point $s=-ix_0$ in (13.3.1.2) the radial character cancels.
Fourier uniqueness on $K$ therefore identifies all the conditions with
\[
    M_j(k)=\int_\R x^j\widetilde\phi(k,x)\dd x=0
            \quad(k\in K,\ 0\le j<d).
 \tag{27}
\]
The map $\phi\mapsto(M_0,\ldots,M_{d-1})$ is onto
$\mathcal D(K)^d$. Indeed the functions $b,b',\ldots,b^{(d-1)}$
for $b\in C_c^\infty((1,2))$ of integral one have a triangular moment
matrix with diagonal $(-1)^j j!$. Multiply its dual basis by prescribed
functions of $k$. The kernel is $\partial_x^d\mathcal D(C_F)$:
for (27) take the primitive
$\int_{-\infty}^x(x-y)^{d-1}\widetilde\phi(k,y)\dd y/(d-1)!$,
which vanishes beyond the support by (27). Thus the codimension is $d$
in a single character component and is infinite on the full test space.
Indeed $K$ contains the injective image of the infinite compact unit
group at one finite place, hence has infinitely many characters.

In the $\kappa$ component use $\widetilde\phi(k,x)=\bar\kappa(k)h(x)$
and put $H(z)=\int h(x)e^{zx}\dd x$. The centered parameter of a zero
of $L_f(s,\chi)$ in these coordinates is $z_\rho=\rho-1/2+i x_0$.
With Weil's local terms, subtract the capped sum in (13.3.1.3) from
the full zero sum. The coefficient of $H(z_\rho)$ in this difference is
\[
       m_\rho-\mathbf1_{\Re\rho=1/2}\min(m_\rho,d).
 \tag{28}
\]
For any $z_{\rho_0}\ne0$ choose
$h=\partial_x^d(e^{-z_{\rho_0}x}b(x))$. It satisfies (27) and
$H(z_{\rho_0})=(-z_{\rho_0})^d\ne0$. Its right translates still
satisfy (27) and avoid the norm-one subgroup. The Laplace argument
on \pg{29} detects (28) at each such point. Conversely, if these
coefficients vanish, the only remaining point has $H(0)=0$.
Thus in a fixed component the capped identity on (27) is equivalent
to location on the line \emph{and} $m_\rho\le d$ away from
$\rho=1/2-ix_0$. For the uncapped sum it is equivalent to location
alone. The distinguished point always lies on the line.

\leaf{522g}{A multiple Hecke zero}{13.3.1 \quad AN ARITHMETIC EXAMPLE OVER THE GAUSSIAN FIELD}{A zero of order at least three}
Let $n=1254=2\cdot3\cdot11\cdot19$ and
$E/\Q:\ y^2=x(x-n)(x+n)$. Its complex multiplication by $\Z[i]$
gives a unitary Hecke character $\chi$ over $F=\Q(i)$ for which
\[
 L_f(s,\chi)=L(E/\Q,s+1/2),\qquad
 \chi_\infty(z)=(z/|z|)^{e},\quad e\in\{-1,1\}.
 \tag{29}
\]
This is the CM correspondence for congruent-number curves. The algebraic
character has infinity type $(1,0)$ or its conjugate, and division of
its ideal coefficients by the square root of the norm gives (29).
See \citetitle{https://intlpress.com/site/pub/files/_fulltext/journals/iccm/2017/0005/0001/ICCM-2017-0005-0001-a003.pdf}{the Hecke character description of the congruent-number curve}.

The following two points certify $\operatorname{rank}E(\Q)\ge2$:
\[
        (-912,25992),\qquad (-528,26136).
\]
Both satisfy the displayed cubic. The square class of $x(P)$ defines
a homomorphism $E(\Q)\to\Q^\times/\Q^{\times2}$, with value $1$ at
the identity and $-n^2$ modulo squares at $(0,0)$. Indeed the product
of the three intersection abscissas with a line of intercept $b$ is
$b^2$; addition by $(0,0)$ sends $x$ to $-n^2/x$. The three nonidentity
rational two-torsion points have classes $-1,n,-n$, all nonsquares.
None is twice a rational point, so there is no rational point of order
four. Odd-order torsion maps trivially. The full torsion image is
therefore generated by $-1,n$. The two displayed points have classes
$-57$ and $-33$. Neither class nor their product (the class of $11\cdot19$)
belongs to $\langle-1,n\rangle$. Their images are independent modulo
torsion and doubles, proving the rank bound.

The \citetitle{https://doi.org/10.1073/pnas.1216991109}{root-number formula}
for squarefree congruent-number twists gives sign $-1$ since
$n\equiv6\pmod 8$. Thus the analytic order at $s=1$ is odd.
If it were one, the
\citetitle{https://abel.math.harvard.edu/~gross/preprints/lectures-pcmi.pdf}{Gross--Zagier--Kolyvagin theorem, Theorem 3.3},
would give rational rank one, contradicting the two points. Consequently
\[
             m=\operatorname{ord}_{s=1/2} L_f(s,\chi)\ge3 .
 \tag{30}
\]
This lower bound uses the proved rank-one theorem. It requires neither
the full Birch--Swinnerton-Dyer conjecture nor a numerical estimate of
any central derivative.

\leaf{522h}{The choice of norm section}{13.3.1 \quad A COUNTEREXAMPLE TO THE CAPPED IDENTITY}{A permitted section detects the lost multiplicity}
Keep the single Hecke character of (29). At its complex place set
\[
       \tau(t)=(t^{1/2}e^{i\log t},1_f),\qquad t>0 .
\]
The normalized complex absolute value is $|z|^2$, so $\tau$ is a
continuous homomorphic section to the idelic norm as required in
(13.3.1.2). Equation (29) gives $\chi(\tau(t))=t^{ie}$.
The zero at $\rho=1/2$ has $z_\rho=ie$ in the notation of \pg{522f}.
The moment conditions now impose vanishing at $0$, a different point.

Take $b\in C_c^\infty((1,2))$ of integral one. With $a\ge0$ use
\[
      h(x)=\partial_x^2(e^{-iex}b(x)),\qquad
      \phi_a(k\tau(e^x))=\bar\kappa(k)\,h(x-a).
\]
Every $\phi_a$ is smooth Bruhat, vanishes near $C_F^1$, and satisfies
all of (13.3.1.2): the first two moments of $h$ and of each translate
vanish, and the other compact characters integrate to zero. Its
centered transform in the selected component is $e^{az}H(z)$, where
$H(ie)=(ie)^2=-1$.

Use the full Hecke--Weil formula and form the difference of its zero
sum and the one prescribed by (13.3.1.1)--(13.3.1.3) on these tests.
Its Laplace transform in $a$, initially on $\Re w>1/2$, is
\[
 \sum_\rho
 \frac{\big(m_\rho-\mathbf1_{\Re\rho=1/2}\min(m_\rho,2)\big)
              H(z_\rho)}{w-z_\rho}.
 \tag{31}
\]
The series continues normally off its poles by rapid vertical decay
and the Hecke zero count. At $w=ie$ the residue is $-(m-2)\ne0$ by
(30). Hence for at least one $a\ge0$ the capped spectral side differs
from Weil's zero sum on an admissible test. This disproves the proposed
identity with Weil's local terms for the allowed choice of section.
The same residue obstructs the printed archimedean coefficient in
(13.3.1.4). Take finite unit volumes one and $\mu_1(K)=1$, so the
finite terms agree with (B.4.4.16). On the sole complex place the
translated test is supported at $|g|_v\asymp e^a$, where
$|g|_v^{1/2}/|1-g|_v=O(e^{-a/2})$ for large $a$. This annulus has
bounded multiplicative Haar mass. Changing its coefficient therefore
adds a Laplace transform holomorphic at $ie$.

For the radial section $\tau_0(t)=(t^{1/2},1_f)$ the same zero is
at the moment point and contributes zero on (27). On unrestricted
tests it gives an actual discontinuity at width zero in degree four.
Indeed as positive widths shrink, the trace tends to the
full-multiplicity sum on the line by absolute convergence. Subtract
(9a). Taking $h=b_0*\overline{b_0(-\,\cdot)}$ for a smooth compactly
supported $b_0$ of integral one gives $H(iu)\ge0$ and $H(0)=1$.
The difference of the limiting trace and the endpoint trace is at
least $m-2\ge1$. Increasing the Sobolev degree as on \pg{446d}
recovers all these line multiplicities. Zero location then asks only
for the independent comparison in (19) or (24).

\leaf{522i}{The principal components}{13.3 \quad REDUCTION BY CYCLIC SUBGROUPS}{One principal character suffices}
The full Hecke zero-location assertion is equivalent to the following
two prime estimates together:
\begin{itemize}\setlength{\itemsep}{2pt}
\item the rational estimate (21) for ordinary $\zeta$;
\item for every nontrivial normalized character over every number field,
and every $\epsilon>0$, the lower bound
\[
   \int_X^{2X}\Re\widetilde{\mathcal B}_\chi(Y)\frac{\dd Y}{Y}
                        \ge -C_{F,\chi,\epsilon}X^{1/2+\epsilon}
                        \qquad(X\ \hbox{sufficiently large}).
 \tag{32}
\]
\end{itemize}
In particular an additional principal-character estimate for each
number field can be eliminated. The criterion also applies to
infinite-order characters, which occur directly in (32).

Indeed (20l) and (20t) locate all zeros of the nontrivial normalized
characters. To recover the principal one over $F$, choose a finite
Galois closure $M/\Q$ with group $G$ and put $H=\operatorname{Gal}(M/F)$.
There are integers $D>0$ and $c_{C,\eta}$ such that
\[
    D\big(\operatorname{Ind}_H^G1-1_G\big)
       =\sum_{\substack{C\le G\ {\rm cyclic}\\\eta\in\widehat C,\ \eta\ne1}}
                     c_{C,\eta}\operatorname{Ind}_C^G\eta .
\]
In fact any class function orthogonal to all the induced characters on
the right restricts to a constant on each cyclic $C$ by Frobenius
reciprocity and Fourier inversion. Each constant is its value at the
identity. Their span is therefore exactly the orthogonal complement
of $1_G$. The left character before multiplication by $D$ lies in
that complement. All its coordinates and those of the spanning
characters are integers in the irreducible character basis. Rational
linear elimination and clearing denominators give the identity.

Artin induction of Euler factors and abelian reciprocity over $M^C$ give
\[
     \left(\frac{\zeta_F(s)}{\zeta_\Q(s)}\right)^D
         =\prod_{C,\eta\ne1}L_{M^C,f}(s,\eta)^{c_{C,\eta}} .
 \tag{33}
\]
This follows first on $\Re s>1$, including ramified factors via inertia
invariants, and then by meromorphic continuation of the Hecke functions.
Each $\eta$ here gives a nontrivial finite-order Hecke character and
is trivial on the positive norm section. Off the critical line in
$0<\Re s<1$, (32) makes every factor regular and nonzero even when
its exponent is negative. Thus the zero order of $\zeta_F$ there is
exactly the zero order of $\zeta_\Q$, which vanishes by (21).
Norm twists just translate this conclusion. Conversely Hecke GRH gives
(21) and (32) by (20l), (20t). The lower bound (32) is the remaining
independent arithmetic input alongside (21).

\leaf{522j}{Removing finite ramification}{13.3 \quad TRANSFER OF ZEROS UNDER BASE CHANGE}{A conductor-one character with the same zero}
For an abelian extension $E/F$ and a unitary Hecke character $\chi$,
abelian reciprocity and the regular representation give
\[
 L_{E,f}(s,\chi\circ N_{E/F})
      =\prod_{\eta\in\widehat{\operatorname{Gal}(E/F)}}L_{F,f}(s,\chi\eta).
 \tag{34}
\]
The equality of unramified factors follows by factoring
$1-(\chi(\mathfrak p)(N\mathfrak p)^{-s})^f$ over the residue-degree
roots of unity. Inertia invariants give the primitive factors at the
remaining places. Even omitting those places preserves all orders in
$0<\Re s<1$. Each factor is holomorphic there. Every zero of $L_F(s,\chi)$
therefore persists, with at least its original order, up any tower of
abelian extensions after norm pullback.

Suppose first that $\chi$ is normalized and of infinite order. At each
of its finitely many ramified finite places $v$, let $D_v/F_v$ be the
cyclic extension with norm group
$\langle\varpi_v\rangle\ker(\chi_v|_{\mathcal O_v^\times})$.
It exists by local class field theory since the unit character has
finite cyclic image. Choose a common multiple $m$ of these local
degrees and take $F_1=F(\mu_{4m})$. This is abelian over $F$ and has
only complex archimedean places. For each $w\mid v$ the compositum
$D_vF_{1,w}/F_{1,w}$ is cyclic of degree dividing $m$. Since $\mu_m$
is present, it has the form $F_{1,w}(\sqrt[m]{a_w})$ by
\citetitle{https://math.mit.edu/classes/18.785/2019fa/LectureNotes20.pdf}{Kummer theory}.
The coset $a_w(F_{1,w}^\times)^m$ is open by Hensel's lemma, even at
places dividing $m$. Weak approximation gives a single $a\in F_1^\times$
in all these cosets. For $F_2=F_1(\sqrt[m]{a})$, every completion at
such a $w$ contains the prescribed $D_v$. The extension $F_2/F_1$
is cyclic of degree dividing $m$.

The norm to $F_v$ of a local unit in $F_2$ is a unit and is a norm
from $D_v$, by transitivity. The pullback $\chi_2=\chi\circ N_{F_2/F}$
is thus trivial on local units over every formerly ramified place.
Over all other finite places this already follows from unramifiedness
of $\chi_v$. If there were no ramified finite places initially, take
$F_2=F(i)$ instead. Norms carry the uniform positive section to the
corresponding section of $F$, preserving normalization. On diagonal
ideles from $F$, $\chi_2$ restricts to $\chi^{[F_2:F]}$, so it still
has infinite order. This constructs a conductor-one infinite-order
character over a totally imaginary field to which every original zero
transfers by (34).

\leaf{522k}{The remaining character family}{13.3 \quad THE FINITE-ORDER CASE}{Prime-order ideal class characters suffice}
For finite-order characters, it suffices to establish (32) for
\emph{prime-order characters of the ordinary ideal class group} over
totally imaginary fields. To prove this, first let $\chi$ have prime
order $\ell$ over an arbitrary $F$, and let $D/F$ be its cyclic
extension of degree $\ell$. Put $F_1=F(\mu_\ell)$. Its degree divides
$\ell-1$, so $DF_1/F_1$ still has degree $\ell$.
Use the local construction of \pg{522j} with $m=\ell$ at every finite
ramified place. In choosing $a$ by weak approximation, also impose
$v_{w_0}(a)=1$ at one auxiliary prime of $F_1$ unramified in $DF_1$
and away from the prescribed primes. If $\ell=2$, require $a<0$ at
every real embedding of $F_1$. All these are independent open local
conditions.

Now $F_2=F_1(\sqrt[\ell]{a})$ is cyclic of degree $\ell$ and totally
ramified at $w_0$ by Eisenstein's criterion. It differs from $DF_1$,
hence the two prime-degree Galois extensions are linearly disjoint
over $F_1$. The pullback of $\chi$ to $F_2$ therefore still has order
$\ell$ by reciprocity. The local norm argument on \pg{522j} makes it
unramified at every finite place, including $w_0$, since the original
character was already unramified there. The field $F_2$ is totally
imaginary: for odd $\ell$ this follows from $\mu_\ell$, and for
$\ell=2$ from the signs of $a$. A finite-order continuous character
is trivial on each connected complex multiplicative group. Triviality
also on the finite unit groups now makes this an ordinary ideal class
character. Formula (34) gives the transfer for $\chi$.

For a nontrivial character of arbitrary finite order let $E/F$ be the
cyclic extension it cuts out. Take a tower from $F$ to $E$ with cyclic
prime-degree steps. The ratio $\zeta_E/\zeta_F$ is the product of the
successive ratios; each is the product of the nontrivial prime-order
Hecke functions at that step. Consequently this ratio is nonzero off
the critical line in the open strip if the prime-order family above
satisfies (32). By abelian factorization it is also the product of
all the nontrivial character functions of $E/F$, each holomorphic
in that strip. Every one of those factors is then nonzero there.

Combining this with (33) and \pg{522j}, the full zero-location claim
is equivalent to (21) over $\Q$ and (32) for just the following
normalized characters over totally imaginary fields: prime-order
ideal class characters and infinite-order characters unramified at
every finite place. Their Euler products contain every finite-place
factor. The constant in (32) may depend on the character and field;
the transfer uses no conductor-uniform estimate.

\leaf{522l}{Principal prime ideals}{13.3 \quad A CRITERION IN TERMS OF GENERATORS}{A lower bound counting only principal primes}
For a totally imaginary number field $F$ let $h_F$ be its ordinary class
number. Let $\vartheta:F_\infty^\times\to S^1$ be a continuous character
trivial on the global units and the uniform positive norm section.
For a principal prime ideal $\mathfrak p=(a)$ the weight
$\vartheta(a)$ is independent of the generator. From (22) form the
single nonnegative smooth weight
\[
 v_2(u)=\int_1^2 v(u/t)\frac{\dd t}{t},\qquad
 V_2(s)=\int_0^\infty v_2(u)u^{s-1}\dd u
       =\frac{2^s-1}{s}V(s).
\]
Its Mellin transform has no zeros for $\Re s>0$, is positive on the
positive real axis, and satisfies $V_2(1)=V(1)$. Set
\[
 C_{F,\vartheta}(X)=
 h_F\!\sum_{\substack{\mathfrak p=(a)\ {\rm prime}\\
                       N\mathfrak p=p\ {\rm rational\ prime}}}
 (\log p)\,\Re\vartheta(a)\,v_2(p/X)
 -\mathbf1_{\vartheta=1}\,X V(1).
\]
Each prime ideal in this sum is counted once. The full Hecke GRH is
equivalent to the following assertion for every such $F,\vartheta$:
\[
 \boxed{\quad C_{F,\vartheta}(X)\ge
 -C_{F,\vartheta,\epsilon}X^{1/2+\epsilon}
 \quad (\epsilon>0,\ X\ \hbox{sufficiently large}).\quad}
 \tag{35}
\]
For $\vartheta=1$ this is an ordinary weighted count of principal
prime ideals with the same explicit main term for every field after
multiplication by $h_F$. This case alone is equivalent, over all
totally imaginary fields, to zero location for all finite-order
Hecke characters over all number fields. For $F=\Q(i)$ and
$\vartheta(a)=(a/|a|)^4$, the split primes give the weights
$2\cos(4\arg\pi_p)$ of (25) and the main term vanishes.

\point{Extension from principal ideals}
On principal fractional ideals prescribe $\psi((a))=\vartheta(a)$.
Choose ideal representatives for a cyclic decomposition of the class
group. If a representative has order $e$ modulo principal ideals,
assign it an $e$th root of the prescribed value of its $e$th power.
This extends the prescription to a unitary character of the fractional
ideals. With archimedean component $\vartheta^{-1}$ it defines a
normalized Hecke character unramified at every finite place. All the
extensions form a family $\mathcal X(\vartheta)$ of size $h_F$,
obtained by multiplying any one by the ideal class characters. It
contains the trivial character exactly when $\vartheta=1$.

\leaf{522m}{Principal prime ideals}{13.3 \quad FROM THE PRIME COUNT TO THE ZEROS}{What the lower bound detects}
On a fractional ideal $\mathfrak a$, orthogonality gives
\[
  \sum_{\psi\in\mathcal X(\vartheta)}\psi(\mathfrak a)
  =\begin{cases}h_F\vartheta(a),&\mathfrak a=(a),\\
                  0,&\mathfrak a\ \text{nonprincipal}.
    \end{cases}
\]
Apply this to every prime ideal power in (24), averaging over $[X,2X]$
with measure $\dd Y/Y$. The average replaces $v$ by $v_2$. The powers
and higher residue degrees with norm $p^k$, $k\ge2$, together cost
$O_{F,v}(\sqrt X)$ by the Chebyshev bound as on \pg{522c};
this includes principal powers of nonprincipal primes. Hence the sum
of the real parts of the averaged errors in (24) over
$\psi\in\mathcal X(\vartheta)$ equals $C_{F,\vartheta}(X)+O_{F,v}(\sqrt X)$.
Zero location for the family gives $C_{F,\vartheta}(X)=O_{F,\vartheta,v}(\sqrt X)$.
Unconditionally the fixed-character prime ideal theorem on \pg{520i}
gives $C_{F,\vartheta}(X)=o_{F,\vartheta}(X)$ by partial summation.
With the external half-plane of \pg{520b} the error is $O_{F,\vartheta}(X^{\sigma_*})$.

Conversely, in $\Re s>1$ the Mellin transform of $C_{F,\vartheta}$
on $[1,\infty)$, integrated against $X^{-s-1}\dd X$, is
\[
 -\frac{V_2(s)}2\sum_{\psi\in\mathcal X(\vartheta)}
 \left(\frac{L'_\psi}{L_\psi}(s)+
       \frac{L'_{\bar\psi}}{L_{\bar\psi}}(s)\right)
 -\frac{\mathbf1_{\vartheta=1}V(1)}{s-1}+J(s).
\]
Here $L_\psi=L_{F,f}(\,\cdot\,,\psi)$ and $J$ is holomorphic on
$\Re s>1/2$ by the discarded-power estimate and compact support at
small $X$. The pole at one cancels. At every zero with $\Re\rho>1/2$
in the open strip the multiplicities add, and $V_2(\rho)\ne0$. At a
real zero the residue is strictly negative.

If (35) holds and such a zero exists, fix $1/2<\theta<\Re\rho$.
For large $X$, $C_{F,\vartheta}(X)+cX^\theta$ is nonnegative for
some $c$. Its Mellin integral over that tail has real abscissa between
$\Re\rho$ and $1$, by the pole at $\rho$ and absolute convergence to
the right of one. At that abscissa the displayed continuation is
regular or has a simple pole with negative residue. Both contradict
the real singularity and positivity of the Mellin integral proved on
\pg{520h}. Functional equation symmetry locates every zero of the
family on the critical line.

Any finite-order character becomes trivial after norm pullback to its
abelian class field $D$, and remains trivial on $D(i)$. Its zeros
persist in $\zeta_{D(i)}$ by (34); this function belongs to the
$\vartheta=1$ family there. For an infinite-order normalized character
use \pg{522j} to transfer its zeros to a character in one of the
families above. Norm twists preserve real parts. This proves both
directions of the equivalence on \pg{522l}.

\leaf{522n}{Normal extensions}{13.3 \quad TRANSFER AND ZERO ORDERS}{Passing to a Galois field preserves each zero}
If $M/F$ is finite Galois and $\xi$ a unitary Hecke character of $F$,
then at every $\rho$ in the open critical strip
\[
    \operatorname{ord}_\rho L_{M,f}(s,\xi\circ N_{M/F})
           \ \ge\ \operatorname{ord}_\rho L_{F,f}(s,\xi).
 \tag{36}
\]
Thus the fields needed for an independent prime estimate may all be
taken Galois over $\Q$. The group of the normal closure need satisfy
no solvability condition.

Write $G=\operatorname{Gal}(M/F)$. For each irreducible character
$\pi$ of $G$, let $a_\pi$ be the order at $\rho$ of the Artin function
twisted by $\xi$. This order is an integer. Indeed, Brauer induction
writes $\pi$ as an integral sum of characters induced from linear
characters of subgroups. After twisting, each induced Euler product
is a Hecke function on the corresponding intermediate field, twisted
there by the norm pullback of $\xi$. The finitely many omitted factors
preserve the order because they are nonzero and finite in
the open strip. This gives meromorphic continuation and defines the
same order for any induction expression.

Form the virtual character $\Theta=\sum_\pi a_\pi\pi$. For a cyclic
subgroup $C\le G$ the multiplicity of any $\eta\in\widehat C$ in
$\Theta|_C$ equals the order of the Hecke function over $M^C$ with
character $\eta(\xi\circ N_{M^C/F})$, by reciprocity of induction and
restriction. This order is nonnegative. Consequently $\Theta|_C$ is
a character, and $|\Theta(g)|\le\Theta(1)$ on each cyclic subgroup.
The regular representation gives
$\Theta(1)=\operatorname{ord}_\rho L_{M,f}(s,\xi\circ N_{M/F})$.
Averaging $\Theta$ over $G$ gives $a_{\mathbf1}$, proving (36).
The corresponding untwisted Heilbronn character is treated by
\citetitle{https://www-math.nsysu.edu.tw/~pjwong/stuff/LowOrderZero.pdf}{Wong on zeros of small order}.

For a finite-order character first pass to its class field and then
to a totally imaginary Galois closure over $\Q$ as on \pg{522m}.
Its zeros persist in the Dedekind zeta function there. For an
infinite-order character pass first to the conductor-one character on
\pg{522j}, then to the normal closure. The norm still sends finite
units to units and pulls back a normalized infinite-order character
to one of infinite order. In either case all but finitely many
rational primes with a degree-one prime above them now split completely.

If $\xi=1$, $F=\Q$, and $\Theta(1)=1$, its cyclic restrictions have
degree one, so $|\Theta(g)|=1$ everywhere. Then
$\sum_\pi a_\pi^2=1$, and the positive degree forces $\Theta$ to be
a single linear character. Thus every simple zero of a normal Dedekind
zeta function in the strip belongs to a Dirichlet function. With the
external Dirichlet input of \pg{522b}, any off-line zero here would
have multiplicity at least two.

\leaf{522o}{Completely split primes}{13.3 \quad THE FINITE-ORDER CONTRIBUTION}{The part controlled by rational congruences}
Let $M/\Q$ be finite Galois, $G=\operatorname{Gal}(M/\Q)$, and
$K=M^{[G,G]}$ its maximal abelian subfield. Choose $q\ge3$ with
$K\subseteq\Q(\zeta_q)=:K_q$ by Kronecker--Weber. Replacing $M$ by
$E=MK_q$ preserves every zero of $\zeta_M$ by (34). The maximal
abelian subfield of $E$ is exactly $K_q$. Indeed $M\cap K_q=K$, and
commutators in the fiber product
$\operatorname{Gal}(E/\Q)\subseteq G\times\operatorname{Gal}(K_q/\Q)$
are $([g_1,g_2],1)$. Every pair $g_1,g_2$ has lifts, so these generate
$[G,G]\times1$.

Write $r=[E:K_q]$. For $p$ unramified in $E$ put
\[
   b_E(p)=\mathbf1_{p\equiv1\bmod q}
                    \left(r\mathbf1_{p\ {\rm splits\ in}\ E}-1\right).
\]
The first-degree sum for $\zeta_E$, minus its expected main term,
is the sum of
\[
   \varphi(q)\sum_p b_E(p)\log p\,v_2(p/X)
   \quad\hbox{and}\quad
   \varphi(q)\sum_{p\equiv1\bmod q}\log p\,v_2(p/X)-X V(1).
 \tag{37}
\]
This holds for large $X$ as the weight avoids the ramified primes.
Each split prime contributes $[E:\Q]=r\varphi(q)$ degree-one ideals.
All other unramified primes have higher residue degree. The second
expression in (37) is $O_q(\sqrt X)$ under the rational Dirichlet input,
by character orthogonality and (24). With that input, the finite-order
part of Hecke GRH is equivalent to the lower bound of (35) for just
the first expression over these $E,q$. Finite-order zeros transfer by
(36), and the Mellin argument of \pg{522m} for a single zeta function
recovers all its zeros from the lower bound for the two terms together.

The remaining coefficient has mean zero in every fixed rational
progression. For any $m$, in the compositum $E\Q(\zeta_m)$ a specified
unit residue mod $m$ fixes a coset over $E\cap\Q(\zeta_m)\subseteq K_q$.
Within it, above each compatible element of $\operatorname{Gal}(K_q/\Q)$
there are $r$ elements of $\operatorname{Gal}(E/\Q)$, equally weighted
in the group average. The sum of $b_E$ in each fiber is zero.
Chebotarev and partial summation therefore give
$\sum_{p\le x,\,p\equiv a\bmod m} b_E(p)\log p=o_{E,m}(x)$ for
$(a,m)=1$. The mean of $|b_E|^2$ over the full group is
$(r-1)/\varphi(q)$, since the identity fiber has one value $r-1$ and
$r-1$ values $-1$. For $r>1$ the residual thus has positive mean
square even after its entire congruence contribution has been removed.

\leaf{522p}{Logarithmic averages}{13.3 \quad A FINITE-HEIGHT PRIME ESTIMATE}{Which zeros survive a Gaussian average}
For a fixed normalized unitary $\xi$ over $F$, write $b_\xi(p)$ for
the sum of its values on degree-one prime ideals of norm $p$, with
ramified Euler factors treated as on \pg{520c}. Define
\[
 P_\xi(X)=\sum_p\Re b_\xi(p)\log p\,v_2(p/X)
                      -\mathbf1_{\xi=1}X V(1).
\]
For $X>1$ and $0<\kappa<1$ set $h=\kappa\log X$ and
\[
  {\cal H}_\kappa P_\xi(X)=\frac1{\sqrt{4\pi h}}\int_\R
       e^{-y^2/(4h)-y/2}P_\xi(Xe^y)\dd y .
\]
The standard deviation in $\log X$ is $\sqrt{2\kappa\log X}$ before
the factor $e^{-y/2}$, which preserves the scale $\sqrt X$. At each
fixed $\kappa$ the two prime estimates and the zero condition below are equivalent.
\begin{align*}
 &{\cal H}_\kappa P_\xi(X)=O_{F,\xi,\kappa}(\sqrt X),\\
 &{\cal H}_\kappa P_\xi(X)\ge -C_{F,\xi,\kappa,\epsilon}
                          X^{1/2+\epsilon}\quad(\epsilon>0),\\
 &u+\kappa(u^2-\gamma^2)\le0\quad\text{for every zero }
        \rho=\tfrac12+u+i\gamma\text{ of }L_{F,f}
        (\,\cdot\,,\xi)\text{ in the open strip}. \tag{38}
\end{align*}
Indeed the averaging sends $X^\rho$ exactly to $X^{\Phi_\kappa(\rho)}$,
where $\Phi_\kappa(s)=s+\kappa(s-\tfrac12)^2$. Only zeros with
\[
              |\gamma|<\sqrt{1/(2\kappa)+1/4}
\]
can violate (38), since $|u|<1/2$. Thus a finite-height zero theorem
can give a prime estimate valid for all $X$. At $\kappa\downarrow0$,
every individual zero with $u>0$ eventually violates the inequality.
The second assertion for any sequence $\kappa\downarrow0$ therefore
implies zero location at all heights. The sequence may depend on the
field and character. By (36) one need only treat the trivial character
and normalized infinite-order conductor-one characters over totally
imaginary Galois fields.

\point{The prime powers in the explicit formula}
Restore all ideal prime powers in $P_\xi$ and call the resulting
centered real sum $A(X)$. As on \pg{522m},
$A(X)-P_\xi(X)=O_{F,v}(\sqrt X)$ for every $X>0$.
Shifting the full inverse Mellin integral to $\Re s=-1/4$ as on
\pg{520j} gives, for $X\ge1$,
\[
    A(X)=-\Re\sum_\rho m_\rho V_2(\rho)X^\rho+O_{F,\xi,v}(1).
\]
The norm pole is subtracted, the finitely many trivial zeros crossed
have real part zero, and the new integral is $O(X^{-1/4})$.
The same bounded error holds for $0<X<1$ by compact support, since
$|X^\rho|\le1$ and the transformed zero weights are summable.

\leaf{522q}{Logarithmic averages}{13.3 \quad MELLIN CONTINUATION AFTER AVERAGING}{Recovering every pole above the parabola}
The preceding bounded error contributes $O(X^{\kappa/4})$ after
averaging, by the Gaussian integral. A term $O(\sqrt X)$ contributes
$O(\sqrt X)$. Termwise averaging of the zero series is justified for
fixed $X$ by $\sum_\rho m_\rho|V_2(\rho)|<\infty$ and a Gaussian
majorant using $0\le\Re\rho\le1$. Consequently
\[
 {\cal H}_\kappa P_\xi(X)=
 -\Re\sum_\rho m_\rho V_2(\rho)X^{\Phi_\kappa(\rho)}
                         +O_{F,\xi,\kappa}(\sqrt X).
\]
If all the real parts $\Re\Phi_\kappa(\rho)$ are at most $1/2$, the
absolute zero sum satisfies the same bound. This gives the first two
assertions of (38).

For the converse take the Mellin integral from $1$ to $\infty$
against $X^{-z-1}\dd X$. To the right of all exponents it equals
\[
 -\frac12\sum_\rho m_\rho\!
   \left(\frac{V_2(\rho)}{z-\Phi_\kappa(\rho)}
   +\frac{\overline{V_2(\rho)}}{z-\Phi_\kappa(\bar\rho)}\right)+J(z),
       \qquad J\text{ holomorphic on }\Re z>1/2.
\]
The series continues normally on compact sets off its possible
poles, since $\Re\Phi_\kappa(\rho)\to-\infty$ with $|\Im\rho|$.
Two distinct points with the same image
would have sum $1-1/\kappa<0$, impossible inside the strip. Complex
conjugate terms landing at the same point have the same coefficient
there since $v_2$ is real. A real image comes from a real zero, as
$\Im\Phi_\kappa(\rho)=(1+2\kappa u)\gamma$.
The residue there is strictly negative because $V_2(\rho)>0$.

Suppose the lower bound in (38) holds and some image has real part
greater than $1/2$. Choose $\theta$ strictly between those numbers.
Adding $cX^\theta$ makes the averaged sum nonnegative on a tail.
Its Mellin integral has finite abscissa at least the real part of
the offending image, by its nonzero residue. At the abscissa the
continuation is regular or has a simple pole of negative residue.
The positivity argument on \pg{520h} excludes both. This proves (38).

For a field or character with all zeros in the finite test region
strictly inside the parabola, the averaged \emph{full prime-power sum}
is $O(X^{1/2-\delta})$ for some $\delta>0$ if no image lies on
$\Re z=1/2$. Indeed images outside every sufficiently large
rectangle have real parts tending to $-\infty$. The remaining
prime-square contribution to $P_\xi$ can then occur at order
$\sqrt X$. It is computed for two characters on \pg{522t}.

\leaf{522r}{Two dihedral characters}{13.3 \quad THE ARITHMETIC COMPONENTS OF THE ESTIMATE}{An angular character and a cubic class character}
Consider $P_\xi$ for the following characters.
\begin{itemize}\setlength{\itemsep}{1pt}
\item $F=\Q(i)$ and $\xi((a))=(a/|a|)^4$, the infinite-order
character in (25).
\item $F=\Q(\sqrt{-23})$ and $\xi$ either of the two order-three
characters of its ideal class group.
\end{itemize}
In both cases $\xi$ is unramified at finite places,
$\xi(\bar{\mathfrak a})=\overline{\xi(\mathfrak a)}$, and
$\xi((n))=1$ for positive rational integers. For the second field the
reduced forms of discriminant $-23$ are $(1,1,6)$ and $(2,\pm1,3)$.
The splitting field $M$ of $T^3-T-1$ is its Hilbert class field.
Indeed the cubic is irreducible with square-free discriminant $-23$,
so the normal closure has group $S_3$ and quadratic field $F$.
Only $23$ can ramify. Its inertia group is generated by a
transposition since the cubic discriminant has valuation one and
the ramification is tame. Thus $M/F$ is unramified of degree three,
the class number. At an unramified prime split in $F$, $b_\xi(p)$ is
$2$ if the cubic splits and $-1$ if the cubic is irreducible modulo
$p$. At inert primes it is zero. The nonabelian part of
$6\mathbf1_{p\ {\rm splits\ in}\ M}$ after subtracting
$2\mathbf1_{p\ {\rm splits\ in}\ F}$ is exactly $2b_\xi(p)$.

Hecke's theta series associates to the two characters the normalized
newforms of level and weight $(N,k)=(4,5),(23,1)$, respectively.
In the Gaussian case the algebraic character is $(a)\mapsto a^4$,
so its coefficients are multiplied by $n^{-2}$ in unitary normalization.
In the cubic case the newform is the theta series of the ideal class
character, already of weight one. The respective newspaces with
quadratic character have dimension one. Their normalized forms are
\[
 f_4(z)=\eta(z)^4\eta(2z)^2\eta(4z)^4,\qquad
 f_{23}(z)=\eta(z)\eta(23z).
\]
The newspace data and the product identities are recorded for
\citetitle{https://www.lmfdb.org/ModularForm/GL2/Q/holomorphic/4/5/b/a/}{form 4.5.b.a}
and \citetitle{https://www.lmfdb.org/ModularForm/GL2/Q/holomorphic/23/1/b/a/}{form 23.1.b.a}.

Let $W_N(x)=e^{kx/2} f_N(i e^x/\sqrt N)$ and
$F_N(z)=\int_\R e^{zx}W_N(x)\dd x$. Changing variables in the
classical Mellin integral identifies $F_N(z)$ with the completed
unitary $L$-function at $s=1/2+z$ up to a nonzero constant.
All factors are positive on this integration line. Writing their
exponents as $r_d$, we have $r_d=r_{N/d}$, $\sum r_d=2k$ and
$\sum d r_d=24$. The identity $\eta(i/y)=\sqrt y\,\eta(iy)$ then
gives $W_N(-x)=W_N(x)$. The kernel and all its real exponential tilts
decay faster than any exponential at both ends.

\leaf{522s}{A horizontal zero-free band}{13.3 \quad CURVATURE OF THE MELLIN KERNEL}{An independent bound at low height}
For both characters on \pg{522r} the completed function is nonzero
throughout the horizontal band
\[
       |\Im s|\le\sqrt{4\pi/\sqrt N}.
 \tag{39}
\]
Write $a=2\pi/\sqrt N$. The first power of $q$ in the eta product
gives $-a e^x$ in $\log W_N(x)$. Each logarithmic product factor is
$r_d\log(1-e^{-v})$ with $v=adn e^x$, $n\ge1$. Its second
$x$-derivative before multiplication by $r_d$ is
\[
    \frac{v((1-v)e^v-1)}{(e^v-1)^2}<0 .
\]
The series may be differentiated on compact intervals by exponential
convergence. Evenness now gives $-(\log W_N)''(x)\ge a e^{|x|}\ge a$.

For $u\in\R$ give a real random variable $Z$ density proportional to
$e^{ux}W_N(x)$, and write it as $e^{-U(x)}$ after normalization.
Then $U''\ge a$. Integration by parts gives
$\mathbb E[U'(Z)]=0$ and $\mathbb E[(Z-\mathbb EZ)U'(Z)]=1$.
For an independent copy $Z'$,
\[
  2=\mathbb E[(Z-Z')(U'(Z)-U'(Z'))]\ge2a\,\operatorname{Var}(Z).
\]
Since $\cos t\ge1-t^2/2$,
\[
   \Re\left(e^{-it\mathbb EZ}\frac{F_N(u+it)}{F_N(u)}\right)
        \ge 1-\frac{t^2}{2a}.
\]
At $|t|=\sqrt{2a}$ the cosine inequality is strict after integration
because the distribution has a density. This proves (39) for every
$u$. In particular, for $1/4\le\kappa<1$ every nontrivial zero
satisfies $\Re\Phi_\kappa(\rho)<1/2$ with a fixed margin because
$(4\pi/\sqrt{23}-1/4)/4>1/2$. This range follows directly
from the two eta products.

A larger interval of zero heights is covered by the published
computations of \citetitle{https://jvoight.github.io/articles/cmfs-mf-101222-scrubbed.pdf}{Best et al., \S\S9.2 and 10}.
They check completeness by the explicit formula. The
\citetitle{https://www.lmfdb.org/L/2/2e2/4.3/c4/0/0}{level-four list}
starts at $7.53515\ldots$ and extends past $25$. The
\citetitle{https://www.lmfdb.org/L/2/23/23.22/c0/0/0}{level-23 list}
starts at $5.11568\ldots$ and extends past $18$.
In both cases every zero with $|\Im\rho|\le18$ lies on the line and
has $|\Im\rho|>5$. For $1/600\le\kappa<1$ its image has real part at
most $1/2-25/600$. At $|\Im\rho|>18$ the upper bound is
$1-(18^2-1/4)/600$. Both are below $0.461$. The averaged centered
full prime-power sum on \pg{522p} is therefore $O(X^{0.461})$.

\leaf{522t}{The remaining prime squares}{13.3 \quad A SQUARE-ROOT TERM FROM INERT PRIMES}{The signed prime average at every large scale}
For either character on \pg{522r} and every fixed
$1/600\le\kappa<1$ one has
\[
 \boxed{\displaystyle
     {\cal H}_\kappa P_\xi(X)
           =-\tfrac12 V_2(\tfrac12)\sqrt X+o_\xi(\sqrt X).}
 \tag{40}
\]
The coefficient is strictly negative. The proof for $\kappa\ge1/4$
uses (39) alone for the required zero exclusion. The smaller
coefficients down to $1/600$ use the finite-height computations on
\pg{522s}. The estimate holds as $X\to\infty$ in each case.

In the real full prime-power sum $A(Y)$ of \pg{522p} the pole term
vanishes for both characters. Terms of norm $p^k$, $k\ge3$, cost
$O_v(Y^{1/3}\log(2+Y))$ for $Y\ge1$ by the Chebyshev bound. At norm
$p^2$ a split $p$ has two degree-one prime squares with total weight
$b_{\xi^2}(p)\log p$. An inert $p$ has one prime ideal $(p)$ of
norm $p^2$ with value $1$ and weight $2\log p$. The ramified
primes have no second-degree term once $Y$ is large. It follows that
\[
 A(Y)-P_\xi(Y)=\sum_p b_{\xi^2}(p)\log p\,v_2(p^2/Y)
          +2\sum_{p\ {\rm inert}}\log p\,v_2(p^2/Y)
          +O_v(Y^{1/3}\log(2+Y)).
\]
The character $\xi^2$ is nontrivial in each example. Its first sum is
$o(\sqrt Y)$ by the fixed-character prime ideal theorem applied with
test $v_2(t^2)$ at scale $\sqrt Y$. Inert primes have density $1/2$
with logarithmic weights by the quadratic prime number theorem.
Their term is asymptotic to $\sqrt Y\int_0^\infty v_2(t^2)\dd t
=\tfrac12 V_2(1/2)\sqrt Y$. This computes the secondary term
independently of zero location near the critical line.

The remainder divided by $\sqrt Y$ is bounded on $Y>0$ and tends
to zero at infinity. Its average against the centered Gaussian in
$y$ is $o(1)$ since $\log X+y\to+\infty$ in probability when
$y$ has variance $2\kappa\log X$. Averaging preserves $\sqrt Y$
exactly, while \pg{522s} bounds $\mathcal H_\kappa A(X)$. This proves
(40). With the external $7/8$ half-plane of \pg{520b}, (20b) with the
smooth test $v_2(t^2)$ gives $O(Y^{7/16})$ after subtracting the
inert main term and including the split sum. After averaging, its
exponent becomes $7/16+\kappa/256<0.442$. With the finite-height
input the remainder in (40) is then $O(X^{0.461})$.

\leaf{522u}{Completely split primes}{13.3 \quad THE TERM FROM FROBENIUS SQUARES}{The principal prime count for discriminant \(-23\)}
For a Galois extension $M/\Q$ of degree $d$ with group $G$, put
\[
 P_M(X)=d\sum_{p\ {\rm splits\ completely}}\log p\,v_2(p/X)
                                      -X V(1),\qquad
 j(G)=\#\{g\in G:g^2=1\}.
\]
Suppose every nontrivial zero of $\zeta_M$ with $|\Im\rho|\le T$ lies
on the line, for $T>\sqrt{1/(2\kappa)+1/4}$. If $m_0$ is its order
at $1/2$, then
\[
 \mathcal H_\kappa P_M(X)
    =-\big(m_0+\tfrac12 j(G)\big)V_2(1/2)\sqrt X+o_M(\sqrt X).
 \tag{41}
\]
Indeed $\Re\Phi_\kappa(\rho)=1/2$ is possible only at the central
point. All other images are separated from this line on its left. The full centered
prime-power sum contributes just its central term at this precision.
For unramified $p$ with Frobenius $g$, the coefficient of $\log p$
at norm $p^2$ is $d$ if $g^2=1$ and zero otherwise. The residue degree
is the order of $g$, so in either contributing case the number of
primes times their residue degree is $d$. Chebotarev gives mean
$j(G)$ for this coefficient. The substitution at scale $\sqrt X$ on
\pg{522t} yields the second term of (41). The higher powers and the finitely many
ramified primes contribute $o(\sqrt X)$ before and after averaging.

In the splitting field $M$ of $T^3-T-1$ the group $S_3$ has
$j(G)=4$. Writing $\delta$ for the quadratic Dirichlet character
of conductor $23$ and $\xi$ for the cubic class character,
\[
        \zeta_M=\zeta_\Q L_\Q(\,\cdot\,,\delta)
                         L_F(\,\cdot\,,\xi)^2 .
\]
This follows from abelian factorization over $F=\Q(\sqrt{-23})$ and
the equality of the two conjugate Hecke functions by conjugation of
ideals. For $\kappa\ge1/600$, height $18$ suffices. It is covered by
the computations on \pg{522s}, by
\citetitle{https://arxiv.org/abs/1305.3087}{Platt for Dirichlet characters}
and by \citetitle{https://arxiv.org/abs/2004.09765}{Platt and Trudgian for $\zeta_\Q$}.
The central order is zero. Equation (39) covers the cubic factors, and
the alternating zeta series gives $\zeta_\Q(1/2)\ne0$. The partial
sums of $\delta(n)$ over each period are nonnegative (its positive
residues are $1,2,3,4,6,8,9,12,13,16,18$). Partial summation gives
$L_\Q(1/2,\delta)>0$.

For all sufficiently large $X$, $P_M(X)=C_{F,1}(X)$ in (35) because
a degree-one prime ideal of $F$ is principal precisely when it splits
in the Hilbert class field of $F$, and $h_F=3$. Equality outside a bounded
range suffices for (41), since the normalized difference is bounded
there and the Gaussian mass of that range tends to zero. Thus the
averaged error for this principal prime count is
$-2 V_2(1/2)\sqrt X+o(\sqrt X)$ throughout $1/600\le\kappa<1$.

\leaf{522v}{Oscillations before averaging}{13.3 \quad THE CUBIC SPLITTING COUNT}{Both signs occur on the square-root scale}
For the cubic class character on \pg{522r}, retain the original weight
$v_2$ in $P_\xi$ and put $G(t)=e^{-t/2}P_\xi(e^t)$. Then
\[
    \limsup_{t\to\infty}G(t)>0.001\,V(1/2),\qquad
    \liminf_{t\to\infty}G(t)<-0.830\,V(1/2).
 \tag{42}
\]
Thus (42) concerns the signed difference between completely split
and irreducible cubics at primes split in $\Q(\sqrt{-23})$, without
the additional Gaussian average.

Let $\gamma\in(5.115,5.116)$ be the first positive zero ordinate in
\pg{522s}, of multiplicity $m\ge1$, and set $c=V_2(1/2)/2$.
Take a Gaussian $Y$ of mean zero and variance $T/2$ as in the
average with $\kappa=1/4$ at $X=e^T$. Equation (40) gives
$\mathbb E G(T+Y)\to-c$. At the selected frequency,
\[
  \mathbb E[G(T+Y)e^{-i\gamma(T+Y)}]\ \longrightarrow\
                    -m V_2(1/2+i\gamma).
\]
Indeed the two conjugate ideal class functions coincide, so the zero
sum of the one function is real before modulation. A zero
$\rho=1/2+u+i\gamma'$ now contributes an exponent with real part
$u+(u^2-(\gamma'-\gamma)^2)/4$. For $|\gamma'|\le18$ we have $u=0$.
For $|\gamma'|>18$ the real part is strictly negative. Only
$\gamma'=\gamma$ survives by absolute summability. The constant
prime-square term has mean zero in this limit because
$|\mathbb E e^{-i\gamma(T+Y)}|=e^{-T\gamma^2/4}$. Its $o(1)$ error
is bounded and tends to zero in mean as on \pg{522t}.

Use the two nonnegative weights $1\pm\cos(\gamma(T+Y)+\phi)$, choosing
the phase to align the last limit. Their masses tend to one, and they
assign vanishing mass to every fixed left half-line, where $G$ is
bounded. Hence $\limsup G\ge-c+m|V_2(1/2+i\gamma)|$ and
$\liminf G\le-c-m|V_2(1/2+i\gamma)|$.

For explicit constants, tilt the independent uniforms in (22) by
$e^{x/2}$. A tilted interval of length $\ell$ has variance at most
$\ell^2\cosh(\ell/4)/12$, by taking its second moment about the
midpoint and pairing opposite displacements. Since $\sum\ell_j^2=1/12$
and $\ell_j\le1/4$, centering the sum and using
$\cos x\ge1-x^2/2$ gives
$|V(1/2+i\gamma)|/V(1/2)\ge1-\gamma^2\cosh(1/16)/288>0.908$.
Here $|2^{1/2+i\gamma}-1|>2.36$ since
$|\gamma\log2-\pi|<0.407$, and $|1/2+i\gamma|<5.141$.
Thus $|V_2(1/2+i\gamma)|>0.416\,V(1/2)$, whereas
$c=(\sqrt2-1)V(1/2)$. These inequalities give (42).

\leaf{522w}{Weil positivity}{13.3 \quad THE GEOMETRIC SIDE ON A LOG INTERVAL}{A lower bound for every Hecke character}
Let $\chi$ be a normalized unitary Hecke character of a number field $F$,
$d=[F:\Q]$, $D_\chi=|\operatorname{disc}F|N\mathfrak f_\chi$,
and $\delta_\chi=\mathbf1_{\chi=1}$. For $g\in C_c^\infty(\R)$ supported
on an interval of length $L>0$, put $g^\sharp(x)=\overline{g(-x)}$,
$h=g*g^\sharp$, $M=h(0)=\|g\|_2^2$ and $G(z)=\int g(x)e^{zx}\dd x$.
The centered transform is $H(z)=G(z)\overline{G(-\overline z)}$.
For each distinct nontrivial zero $\rho$ of $L_{F,f}(s,\chi)$, write
$\lambda=\rho-1/2$ and $m_\lambda=\operatorname{ord}_\rho L_{F,f}$. Define the Weil form
$Q_\chi(g)=\sum_\lambda m_\lambda H(\lambda)$. Define
\[
 c(L)=-\gamma_E-\log\!\big(\pi(e^{2L}-1)\big),\qquad
 S(L)=\sum_{2\le n\le e^L}\frac{\Lambda(n)}{\sqrt n}.
\]
Then
\begin{align*}
 Q_\chi(g)&\ge K_\chi(L)\|g\|_2^2,\\
 K_\chi(L)&=\log D_\chi+d\big(c(L)-2S(L)\big)
              -\delta_\chi\big(2\sinh(L/2)-L\big).
 \tag{43}
\end{align*}
In particular $Q_\chi(g)>\|g\|_2^2/4$ for $g\ne0$ and $L\le1/16$,
uniformly over fields and characters. A unitary norm twist modulates
$g$ by a phase $e^{itx}$, preserving its support and norm.

\point{The archimedean terms} In the completed function write the
gamma factor as $\prod_{j=1}^d\Gamma_\R(s+\mu_j)$ as on \pg{520d}.
Real parity and the absolute value of the complex angular exponent
give $\Re\mu_j\ge0$. Put $a=1/4+\Re\mu_j/2$, $b=\Im\mu_j$. In the
full explicit formula the contribution $A_j$ is the integral of
$2\Re(\Gamma_\R'/\Gamma_\R)(1/2+it+\mu_j)$ against
$|G(it)|^2\dd t/(2\pi)$. This normalization retains the terms at log
norm zero for $h(0)=M$. See
\citetitle{https://arxiv.org/abs/math/9902080}{Burnol's Hecke explicit formula}.

The integral for $\psi=\Gamma'/\Gamma$ and Fourier inversion give
\[
 A_j=-(\log\pi+\gamma_E)M+
 \int_0^\infty
 \frac{e^{-u}M-e^{-au}\Re\!\big(e^{-ibu/2}h(u/2)\big)}
      {1-e^{-u}}\dd u .
\]
Near zero the numerator is $O_g((1+|\mu_j|)u)$. The rapid decay of $G(it)$ justifies
integration against the digamma formula by truncation at both ends.
See \citetitle{https://dlmf.nist.gov/5.9.E16}{DLMF 5.9.16}.
By Cauchy--Schwarz $|h(x)|\le M$, and $h(x)=0$ for $|x|\ge L$.
On $0<u<2L$ the integrand is at least $-M$. Beyond $2L$ its integral
is $-M\log(1-e^{-2L})$. Thus $A_j\ge c(L)M$ for every shift.

\leaf{522x}{Primes and conductors}{13.3 \quad THE FINITE AND POLAR TERMS}{How the interval length depends on the conductor}
At an unramified prime write $a_{\mathfrak p}=\chi(\mathfrak p)$ and
put $a_{\mathfrak p}=0$ at the primes of the conductor. With the
conventions of \pg{522w}, the remaining terms of the explicit formula are
\begin{align*}
 Q_\chi(g)={}&M\log D_\chi+\sum_{j=1}^d A_j+
     2\delta_\chi\Re\big(G(1/2)\overline{G(-1/2)}\big)\\
 &-2\Re\sum_{\mathfrak p,\,m\ge1}
 \frac{a_{\mathfrak p}^{\,m}\log N\mathfrak p}{(N\mathfrak p)^{m/2}}
                         h(m\log N\mathfrak p).
 \tag{44}
\end{align*}
Indeed the logarithmic derivative of each Euler factor gives the
positive log shifts in (44), and the functional equation gives their
conjugates. The two poles of the completed principal function have
centered coordinates $\pm1/2$. The conductor factor $D_\chi^{s/2}$
contributes $\log D_\chi$ after reflection.

For $(N\mathfrak p)^m=p^k$ the residue degree $f_{\mathfrak p}$ divides
$k$ and fixes $m$. As $\sum_{\mathfrak p\mid p}f_{\mathfrak p}\le d$,
the absolute logarithmic coefficient at $p^k$ is at most $d\log p$.
Only powers up to $e^L$ occur. The finite-place contribution is
therefore at least $-2d S(L)M$.

To bound the poles, translate the support interval to $[-L/2,L/2]$.
The product $G(1/2)\overline{G(-1/2)}$ is translation invariant. The
vectors $u_\pm(x)=e^{\pm x/2}$ on this interval have equal squared
norm $2\sinh(L/2)$ and inner product $L$. The polar form represents
the operator $u_+\otimes u_-+u_-\otimes u_+$. Its eigenvectors
$u_+\pm u_-$ have eigenvalues $L\pm2\sinh(L/2)$. The orthogonal
complement of their span lies in the kernel. This gives the last term of (43).
At $L=1/16$ the sum $S(L)$ is empty, $c(L)>0.29$, and
$2\sinh(L/2)-L<0.00002$. The lower bound improves for smaller $L$.

\point{Growing conductors} Write $R=\log D_\chi/d$. Chebyshev's
estimate and partial summation give $S(L)\ll e^{L/2}$ with an absolute
constant. Also $c(L)=-2L+O(1)$ as $L\to\infty$. For any
$0<\epsilon<2$, (43) is positive for every
$L\le(2-\epsilon)\log R$ once $R$ is sufficiently large in terms
of $\epsilon$. Indeed the terms subtracted from $dR$ are
$O(d(R^{1-\epsilon/2}+\log R))$ at this endpoint, and $K_\chi$ is
nonincreasing in $L$. This gives a range of support lengths using
only the conductor, archimedean factors and rational prime powers.

\leaf{522y}{Two translates}{13.3 \quad A FIXED TEST AT TWO LOGARITHMIC POSITIONS}{Which correlation must stay bounded}
Use the density $b$ and transform $V$ of (22). Fix $g(x)=e^x b(8x)$,
so its support length is $\ell=1/16$. Its transform
$G(z)=V((z+1)/8)/8$ has no zeros on $|\Re z|\le1/2$, since the
product in (22) has none on $\Re s>0$. The same holds for
$H(z)=G(z)\overline{G(-\overline z)}$. For a fixed normalized $\chi$
put $q=Q_\chi(g)>0$, $T_\tau g(x)=g(x-\tau)$, and
\[
 B_\chi(\tau)=\sum_\lambda m_\lambda H(\lambda)e^{\tau\lambda}.
\]
Zero location is equivalent to either of the two bounds below.
\begin{align*}
 &\Re\rho=1/2\quad\text{for every nontrivial zero of }L_{F,f}(s,\chi),\\
 &|B_\chi(\tau)|\le q\quad\text{for all }\tau\ge0,\\
 &Q_\chi(c_0g+c_1T_\tau g)\ge0
             \quad\text{for all }c_0,c_1\in\C,\ \tau\ge0 .
 \tag{45}
\end{align*}
Indeed reflection $\lambda\mapsto-\overline\lambda$ gives
$B_\chi(-\tau)=\overline{B_\chi(\tau)}$. The two translates have
the same diagonal value $q$, and the mixed coefficients of their
Hermitian form are $B_\chi(\tau)$ and its conjugate. The eigenvalues
of the coefficient matrix are $q\pm|B_\chi(\tau)|$. On the critical
line $H(\lambda)=|G(\lambda)|^2$, proving the bound there.

Conversely, boundedness of $B_\chi$ on the positive half-line makes
its Laplace transform holomorphic for $\Re z>0$. For $\Re z>1/2$
that transform equals $\sum_\lambda m_\lambda H(\lambda)/(z-\lambda)$.
Rapid decay and the zero count give locally uniform convergence off
the zeros, as in (3). A right-hand zero would have nonzero residue
$m_\lambda H(\lambda)$. Reflection now gives the first line of (45).

Equation (43) proves the last line of (45) whenever
$K_\chi(\ell+\tau)\ge0$, since the combined support fits in an
interval of length $\ell+\tau$. In particular this holds for every
component when $0\le\tau\le1/64$, since $S(5/64)=0$ and
$c(5/64)>0.05>2\sinh(5/128)-5/64$. For $\tau>\ell$ the translated
autocorrelation $T_\tau h$ vanishes at zero, and only positive prime
shifts occur. In particular as $\tau\to\infty$,
\[
 B_\chi(\tau)=\delta_\chi e^{\tau/2}H(1/2)
 -\sum_{\mathfrak p,\,m\ge1}
 \frac{a_{\mathfrak p}^{\,m}\log N\mathfrak p}{(N\mathfrak p)^{m/2}}
 h(m\log N\mathfrak p-\tau)+O_{g,d}(e^{-\tau/2}).
\]
In the polarized archimedean integral only $u/2\in[\tau-\ell,\tau+\ell]$
contributes, with $e^{-au}\le e^{-u/4}$. The second polar term is
also $O_g(e^{-\tau/2})$. The remaining bound is thus a cancellation
estimate for prime powers near $e^\tau$ at arbitrarily large separations
for the same character.

\leaf{522z}{Dilations of the test}{13.3 \quad AN IDENTITY FOR THE FIXED TEST}{Reconstruction from narrower log intervals}
For $0<r\le1$ set $g_r(x)=r^{-1/2}e^x b(8x/r)$,
$h_r=g_r*g_r^\sharp$, $M_r=h_r(0)$, and $q_{\chi,r}=Q_\chi(g_r)$.
Here $g_1$ is the test of (45), and the support length of $g_r$ is
$\ell r$, where $\ell=1/16$. Write $H_r$ for the transform of $h_r$ and
$B_{\chi,r}(t)=\sum_\lambda m_\lambda H_r(\lambda)e^{t\lambda}$.
For every $r=2^{-m}$, $m=0,1,\ldots$, there is an even positive measure
$\nu_m$ supported in $[-\ell,\ell]$ with mass at most $e^{7/16}$ such that
\[
 B_{\chi,1}(t)=\frac1r\int B_{\chi,r}(t+y)\dd\nu_m(y).
 \tag{46}
\]

To construct it, let $Y=(5/4+\sum_{k\ge1}U_k)/8$ with the uniforms
from (22). Splitting the sum after $m$ terms gives $Y$ the law of
$Z_m+rY'$, where $Y'$ is an independent copy of $Y$ and
\[
 Z_m=\big((1-r)5/4+\textstyle\sum_{k=1}^{m}U_k\big)/8 .
\]
Its range is $[(1-r)5/32,(1-r)7/32]$. Multiply the law of $Z_m$
by $e^y$ to obtain $\mu_m$. The densities of $Y$ and $rY'$ are
$8b(8x)$ and $8b(8x/r)/r$. Multiplying their convolution identity
by $e^x$ gives $g_1=r^{-1/2}\mu_m*g_r$.
Take $\nu_m=\mu_m*\mu_m^\sharp$, with $\mu_m^\sharp$ its reflection.
Then $h_1=r^{-1}\nu_m*h_r$, which gives (46) by absolute convergence
of the zero sum for each $r$. The total mass of $\mu_m$ lies between
$1$ and $e^{7/32}$, so the positive measure in (46), including $1/r$,
has mass between $1/r$ and $e^{7/16}/r$.

\point{The diagonal term} Uniformly for $0<r\le1$, $M_r$ is bounded
above and away from zero, since $M_r=\int b(u)^2e^{ru/4}\dd u/8$.
For a fixed $\chi$,
\[
 q_{\chi,r}=M_r\big(d\log(1/r)+O_{\chi,g_1}(1)\big).
\]
Indeed, $\|g'_r\|_2\ll\|g_r\|_2/r$, so
$|h_r(x)-M_r|\ll M_r |x|/r$. The part $0<u<2\ell r$ of the
digamma integral on \pg{522w} is $O_{\chi,g_1}(M_r)$. Its tail is
$-M_r\log(1-e^{-2\ell r})=M_r(\log(1/r)+O(1))$.
No finite prime shift occurs on the diagonal. The polar and conductor
terms in (44) are $O_{\chi,g_1}(M_r)$.

\leaf{522aa}{Widths and zeros}{13.3 \quad THE COST OF RECONSTRUCTION}{What a shrinking test can detect}
Fix one normalized unitary Hecke character $\chi$. For each sufficiently
large integer $j$, choose a dyadic width $r_j\le1$, and put
\[
 I_j=[j-\ell,j+1+\ell],\qquad
 \theta=\limsup_{j\to\infty}\frac{\log(1/r_j)}{j}.
\]
For finite $\theta$, a bound on these intervals confines the zeros to a strip.
\begin{align*}
 &|B_{\chi,r_j}(t)|\le q_{\chi,r_j}
                \quad(t\in I_j,\ j\ \text{sufficiently large})\\
 &\hspace{36pt}\Longrightarrow\quad
       |\Re\rho-1/2|\le\theta
       \quad\text{for every nontrivial zero of }L_{F,f}(s,\chi).
 \tag{47}
\end{align*}
For $\tau\in[j,j+1]$, all the translates in (46) lie in $I_j$.
The diagonal estimate therefore gives
\[
 |B_{\chi,1}(\tau)|\ll_{\chi,g_1}r_j^{-1}\big(1+\log(1/r_j)\big).
\]
For each $\epsilon>0$ this is $O_{\chi,g_1,\epsilon,(r_j)}
(e^{(\theta+\epsilon)\tau})$. The Laplace transform in (45)
is then holomorphic on $\Re z>\theta$. Each centered zero there
would have nonzero residue $m_\lambda H_1(\lambda)$, since (22)
gives $H_1(\lambda)\ne0$ throughout the strip. Reflection proves (47).

In particular the bound on $I_j$ for one sequence with
$r_j\to0$ and $\log(1/r_j)=o(j)$ suffices for this character's
Riemann hypothesis. Conversely, on the critical line the spectral
sum for each $B_{\chi,r}$ is a sum of phases with nonnegative weights
$m_\lambda|G_r(\lambda)|^2$. Their sum is $q_{\chi,r}$, so the
inequality holds for every $r$ and every separation. Here
$G_r(z)=\sqrt r\,V(r(1+z)/8)/8$ as follows by substitution.

With $X=e^t$, the prime powers in the translated test lie between
$Xe^{-\ell r}$ and $Xe^{\ell r}$. At $t\in I_j$ a width of the
form $r_j=e^{-\theta j+o(j)}$ thus samples an interval of length
$X^{1-\theta+o(1)}$. The next page proves the inequality on every
sufficiently large $I_j$ with length of order $\sqrt X\log X$.
At these widths the exponent in (47) is $1/2$. To obtain the
critical line by (47), the corresponding interval length must be
$X^{1-o(1)}$ for one sequence satisfying the bound.

\leaf{522ab}{Prime intervals}{13.3 \quad AN UNCONDITIONAL ESTIMATE FOR ALL LARGE SEPARATIONS}{Positivity on intervals of length \(\sqrt X\log X\)}
Let $r=r_j$ be the largest dyadic number at most
$j e^{-j/2}/4$, so $j e^{-j/2}/8<r\le j e^{-j/2}/4$ for large $j$.
For every normalized unitary Hecke character, uniformly in the number
field and character,
\[
 |B_{\chi,r}(t)|<q_{\chi,r}\qquad(t\in I_j,\ j\ \text{sufficiently large}).
 \tag{48}
\]
For $X=e^t$ and $L=\ell r$ the interval $[Xe^{-L},Xe^L]$ has length
$Y=2X\sinh L\asymp\sqrt X\log X$, and $\log Y\sim(\log X)/2$
uniformly on $I_j$. By the short-interval Brun--Titchmarsh inequality,
the number of primes in it is at most $2Y/\log Y+1$.
Consequently the sum of $\log p/\sqrt p$ over them is bounded by
$(4+o(1))Y/\sqrt X=(8\ell+o(1))r\sqrt X$.
See \citetitle{https://doi.org/10.1112/S0025579300004708}{Montgomery and Vaughan, Theorem 2, (1.12)}.

All higher rational prime powers in this interval together contribute
$O(1)$ with weights $\log p/\sqrt{p^k}$. Indeed Chebyshev's bound gives
$\sum_{p^k\le 2X,\, k\ge2}\log p\ll\sqrt X+X^{1/3}\log X$, and every
denominator in the interval is at least $\sqrt{Xe^{-L}}$.
The coefficient bound at each $p^k$ in (44) contributes at most
a factor $d$. Also $|h_r|\le M_r$.

In the mixed term of (44) there is no conductor contribution
since $t>L$. Each archimedean integral is $O(LM_r/\sqrt X)$
by its kernel on \pg{522y}, uniformly in the archimedean shifts.
From the support of $h_r$ we have
$|H_r(\pm1/2)|\le 2LM_r e^{L/2}$. Including the principal pole gives
\[
 \frac{|B_{\chi,r}(t)|}{dM_r}
       \le (10\ell+o(1))r\sqrt X+O(1).
\]
Here $r\sqrt X\le j e^{(1+\ell)/2}/4$ on $I_j$ and
$10\ell e^{(1+\ell)/2}/4<1/3$. On the diagonal, (43) gives
$q_{\chi,r}/(dM_r)\ge j/2-\log j-O(1)$, since $D_\chi\ge1$.
The constants and the threshold in this comparison are absolute
for the fixed test of (45). This proves (48).

\leaf{522ac}{The rational component}{13.3 \quad A FIXED WIDTH AND AN EXPLICIT PRIME ESTIMATE}{A fixed test through prime norm \(10^{19}\)}
For ordinary $\zeta$ fix $r=1/8$ and $L=\ell r=1/128$. At this
single width the translated test satisfies
\[
       |B_{1,r}(\tau)|\le q_{1,r}
       \qquad(0\le\tau\le19\log 10-L).
 \tag{49}
\]
By (43), $q_{1,r}>2.4M_r$. For $0\le\tau\le1/32$ the combined
support has length at most $L+1/32<1/16$, so (43) gives (49)
by the two-by-two form in (45).

For $\tau\ge1/32$ put $X=e^\tau$. The archimedean term of the
mixed form has absolute value at most
\[
 A(\tau)=\frac{2L M_r e^{-(\tau-L)/2}}{1-e^{-2(\tau-L)}}.
\]
This follows directly from the real-place kernel in (17) since
$T_\tau h_r$ is supported on the positive half-line. When $X\le12$
the prime interval has length at most $24\sinh L<1$. Its single
possible integer contributes at most $(2/e)M_r$ since
$\Lambda(n)/\sqrt n\le\log n/\sqrt n\le2/e$.
Here $A(\tau)<0.35M_r$, and both poles together have absolute value
at most $2Le^{L/2}(\sqrt{12}+1)M_r<0.08M_r$.

For $12\le X\le10^{19}e^{-L}$ the whole interval lies in
$(11,10^{19}]$. There $|\psi(x)-x|\le0.94\sqrt x$ by
\citetitle{https://arxiv.org/abs/1511.02032}{B\"uthe, Theorem 2, (1.5)}.
Integration of $f(x)=x^{-1/2}h_r(\log x-\tau)$ against $\dd x$
equals the positive pole $\sqrt X H_r(1/2)$. Stieltjes integration
by parts bounds the difference from its prime-power sum by
\[
 0.94\int_{\R}|h'_r(u)-h_r(u)/2|\dd u
          \le0.94(2+L)M_r.
\]
Indeed $h_r$ is even and nonincreasing on $[0,\infty)$. To see this,
exponentially tilt the independent uniforms defining $g_r$ after
rescaling. Each tilted interval density $e^x\mathbf1_{[0,a]}(x)$,
up to a constant, has autocorrelation proportional to
$e^{2a-x}-e^x$ for $0\le x\le a$ and zero for $x>a$. Convolution
preserves even decrease by writing the factors as mixtures of
centered interval indicators. The finite convolutions converge
weakly to the smooth autocorrelation after normalization, giving
the same monotonicity. Thus $\|h'_r\|_1=2M_r$ and
$\|h_r\|_1\le2LM_r$. For $X\ge12$ the remaining pole together with
$A(\tau)$ costs less than $0.02M_r$, proving (49).

\leaf{522ad}{Distinct prime pairs}{13.3 \quad A FIXED TEST AT UNBOUNDED SCALES}{The signed sum over distinct primes}
Fix a normalized unitary Hecke character $\chi$ over $F$. Use the
real even autocorrelation $h=h_1$ from (45), with transform $H=H_1$
and support $[-\ell,\ell]$. Twist by a fixed real frequency $\omega$,
so the finite Hecke function is $L_{F,f}(s-i\omega,\chi)$.
Let $b_\chi(p)$ be the sum of the character values over prime ideals
of norm $p$ as on \pg{522p}. On the positive norm axis set
\[
 \dd\mu_\omega(x)=
 \sum_p\Re\!\big(b_\chi(p)p^{i\omega}\big)\frac{\log p}{\sqrt p}
             \,\delta_p(\dd x)
 -\delta_\chi\,x^{-1/2}\Re(x^{i\omega})\dd x ,
 \qquad \delta_\chi=\mathbf1_{\chi=1}.
\]
Write $w_t(x)=h(\log x-t)$ and define the real quantities
\begin{align*}
 P_\omega(t)&=\int w_t(x)\dd\mu_\omega(x),\\
 D_\omega(t)&=\sum_p\big(\Re(b_\chi(p)p^{i\omega})\big)^2
                      \frac{(\log p)^2}{p}\,w_t(p)^2,\\
 C_\omega(t)&=\iint_{x\ne y}w_t(x)w_t(y)\dd\mu_\omega(x)\dd\mu_\omega(y)
                   =P_\omega(t)^2-D_\omega(t).
 \tag{50}
\end{align*}
Only the prime atoms contribute on $x=y$. When $\chi$ is principal
the continuous term subtracts its pole, including the phase of the
norm twist. For a nonprincipal $\chi$, $C_\omega$ is a sum over
ordered pairs of distinct rational primes.

For $\lambda=\rho-1/2$ put
\[
 R_\omega(t)=\Re\sum_\lambda m_\lambda
           H(\lambda+i\omega)e^{(\lambda+i\omega)t}.
\]
The full explicit formula gives $P_\omega(t)=-R_\omega(t)+O_{F,h,\omega}(1)$
for $t\ge1$. In fact the positive pole contributes
$\delta_\chi e^{t/2}\Re(e^{i\omega t}H(1/2+i\omega))$, precisely
the integral of $w_t$ against the density subtracted in (50).
The opposite pole and all gamma terms are $O_{F,h,\omega}(e^{-t/2})$
by the kernel on \pg{522y}. With $X=e^t$, the rational prime powers
$p^k$ of exponent $k\ge2$ in $[Xe^{-\ell},Xe^\ell]$ have total weight
$O_F(1)$ by the estimate on \pg{522ab}. These statements allow all
the original ramified Euler factors.

\leaf{522ae}{The diagonal term}{13.3 \quad DEGREE-ONE PRIMES ABOVE THE SAME NORM}{The amount of cancellation at a fixed scale}
Put $J(v)=\int h(y)^2e^{ivy}\dd y$. Let $r_\chi$ be the integer
in (20m), and write $F\otimes_\Q F=\prod_i K_i$, with the two norms
$N_{1i},N_{2i}:K_i\to F$ induced on ideles. Let $s_\chi$ count the
factors where $(\chi\circ N_{1i})(\chi\circ N_{2i})=1$. Then
\[
 \frac{D_\omega(t)}{t}=
 \tfrac12 r_\chi J(0)+
 \tfrac12 s_\chi\Re\!\big(e^{2i\omega t}J(2\omega)\big)+o(1)
                    \qquad (t\to\infty).
 \tag{51}
\]
In particular $D_\omega(t)\asymp_{\chi,h,\omega}t$ for every fixed
real $\omega$.

Equation (20m) gives $\sum_{p\le x}|b_\chi(p)|^2\log p=r_\chi x+o(x)$.
The same tensor product counts the ordered pairs in $b_\chi(p)^2$
with the two character values multiplied. Each product character on
$K_i$ is trivial on the uniform positive norm section and is therefore
a pure norm character only when it is trivial. The prime ideal theorem
used on \pg{520i} gives
$\sum_{p\le x} b_\chi(p)^2\log p=s_\chi x+o(x)$.
The prime-by-prime triangle inequality gives $0\le s_\chi\le r_\chi$,
and the diagonal factor in (20m) gives $r_\chi\ge1$.

Apply $(\Re z)^2=(|z|^2+\Re z^2)/2$ in (50). For either cumulative
sum, partial summation against
$x^{-1+iv}\log x\,h(\log x-t)^2$, with $v=0$ or $2\omega$,
costs $o(t)$ in the error. Indeed the support has $x\asymp e^t$
and the integral of $x$ times the absolute derivative is $O_{h,v}(t)$.
The main integral is $e^{ivt}\big(tJ(v)+O_{h,v}(1)\big)$ by the
change of variables $y=\log x-t$. This proves (51).

At $\omega=0$ the coefficient is $(r_\chi+s_\chi)J(0)/2>0$.
For $\omega\ne0$ we have $|J(2\omega)|<J(0)$, since $h^2$ is
positive on an interval and the phase varies there. Thus the
diagonal stays of order $\log X$ even at phases where
$\Re(X^{i\omega})$ vanishes.

\leaf{522af}{The two limits}{13.3 \quad ONE CHARACTER AND ONE FIXED FREQUENCY}{The lower limit of the signed prime correlation}
For every fixed $\chi$ and all but finitely many $\omega\in[1,2]$,
\[
 \liminf_{t\to\infty}\frac{C_\omega(t)}{D_\omega(t)}=-1.
 \tag{52}
\]
The same conclusion holds at $\omega=0$ if $L_{F,f}(s,\chi)$ has no
real zero in $[1/2,1)$. For every real $\omega$, boundedness above of
$C_\omega(t)/D_\omega(t)$ for all sufficiently large $t$ is equivalent
to the Riemann hypothesis for this $\chi$. Under that hypothesis the
ratio tends to $-1$.

To prove the lower limit, choose $\omega$ so that the shifted zeros
have no real point in the closed right half of the centered strip.
The excluded values $\omega=-\Im\rho$ with $\Re\rho\ge1/2$
form a finite set in $[1,2]$ by the zero count.
For $\Re z>1/2$ the Laplace transform of $R_\omega$ is
\[
 \tfrac12\sum_\lambda m_\lambda
     \left(\frac{H(w)}{z-w}+\frac{H(\overline w)}{z-\overline w}\right),
                     \qquad w=\lambda+i\omega .
\]
The series extends meromorphically by rapid vertical decay. At a
coincident pole both multiplicities multiply the same nonzero value
of $H$, since $h$ is real and (45) gives $H\ne0$ on the strip.
There is at least one pole with $\Re z\ge0$ by the zero count and
reflection. None lies on the nonnegative real axis for this $\omega$.

If $R_\omega$ had constant weak sign on a tail, its Laplace integral
there, with the sign chosen positive, would have a real singularity
at its abscissa of convergence. That abscissa lies in $[0,1/2]$
by the preceding pole and the bound $R_\omega(t)=O(e^{t/2})$.
The assertion for a positive integral follows by Taylor expansion
at a real point to the right and monotone convergence of the
exponential series, as on \pg{520h}. Removing a finite initial
interval changes the transform by an entire function.
This contradiction proves that the continuous $R_\omega$ takes
both signs on every tail. At its arbitrarily large zeros,
$P_\omega(t)=O(1)$, so (50)--(51) give (52). The global lower
bound $C_\omega\ge-D_\omega$ is immediate from (50).

On the critical line the sum defining $R_\omega$ is bounded by
$\sum_\lambda m_\lambda |H(\lambda+i\omega)|$, so
$C_\omega=-D_\omega+O(1)$ at all scales. Conversely, any finite
upper bound for the ratio gives $P_\omega(t)=O(\sqrt t)$ by (51).
The transform is then holomorphic on $\Re z>0$. Its residues exclude
every right-hand zero whether or not one of them is real, proving
the stated equivalence for every $\omega$.

\leaf{522ag}{Distances between primes}{13.3 \quad UPPER SIEVE FOR NEARBY PAIRS}{Which prime separations supply the cancellation}
For $X=e^t$ let $C_\omega^{\le Y}(t)$ denote the integral in (50)
restricted to $0<|x-y|\le Y$. Uniformly for $1\le Y\le X$ and $t$
sufficiently large, even its absolute contribution satisfies
\[
 \iint_{0<|x-y|\le Y}|w_t(x)w_t(y)|\,
       \mathrm d|\mu_\omega|(x)\,\mathrm d|\mu_\omega|(y)
                 \ll_{d,h}Y .
 \tag{53}
\]
For a prime pair with difference $k>0$, both primes are in
$[X/2,2X]$ and each product of weights costs $O_{d,h}((\log X)^2/X)$.
Odd differences are absent there. For even $k$, the number of such
pairs is
\[
 \ll\frac{X}{(\log X)^2}\prod_{p\mid k,\,p>2}
                       \left(1+\frac1{p-2}\right).
\]
Indeed the two primes avoid $0,-k$ modulo every prime up to
$X^{1/100}$. The sieve upper bound and Mertens' theorem give the
display, since these residues coincide exactly at divisors of $k$.
See \citetitle{https://doi.org/10.1112/jlms.12663}{Tao and Ter\"av\"ainen, Lemma 3.2}.
On expanding the product in squarefree divisors, its sum over
$k\le Y$ is at most $Y\prod_{p>2}(1+1/(p(p-2)))\ll Y$.
The reverse ordering has the same bound.

The continuous pole density has absolute value at most $x^{-1/2}$.
Its mass within $Y$ of a point, restricted to the support, is
$O(Y/\sqrt X)$. The total atomic mass on the support is $O_d(\sqrt X)$
by Chebyshev, and the continuous mass there is $O(\sqrt X)$ as well.
This proves (53) including the mixed and continuous pairs.

Along the sequence in (52) one has $C_\omega=-D_\omega+O(1)$.
By (51)--(53) there is a fixed $c_{\chi,h,\omega}>0$ such that the
signed contribution from $|x-y|>c_{\chi,h,\omega}\log X$ is at most
$-D_\omega/2$ on an unbounded sequence of $X$. For a nonprincipal
component this is cancellation among actual pairs of distinct primes.

For ordinary $\zeta$, the alternating eta series gives nonvanishing
on $[1/2,1)$. For the two characters of \pg{522r}, the Mellin kernels
of the eta products on \pg{522s} are strictly positive on real
arguments and their gamma factors are regular in this segment.
All three functions have real Dirichlet coefficients. Thus at
$\omega=0$ the real function $R_0$ above is the original
$B_{\chi,1}$ itself. In all three examples the fixed-width correlation
in (45) vanishes at arbitrarily large separations.

\leaf{522ah}{Frequency averages}{13.3 \quad THE UPPER BOUND FOR ONE HECKE COMPONENT}{A mean-square bound over norm twists}
Keep $F$, $\chi$ and $h=h_1$ fixed. For $t\ge1$ set $X=e^t$,
$a_p(t)=b_\chi(p)(\log p)p^{-1/2}h(\log p-t)$ and
\[
 F_t(v)=\sum_p a_p(t)e^{iv(\log p-t)}
               -\delta_\chi\sqrt X\,H(1/2+iv).
\]
Then $P_v(t)=\Re(e^{ivt}F_t(v))$ in (50). Write $Z_v(t)$ for the
complex zero sum whose real part is $R_v(t)$ on \pg{522ad}. The same
explicit formula gives $Z_v(t)=-e^{ivt}F_t(v)+O_{F,h}(1)$ uniformly
in real $v$. The absolute bounds for the higher prime powers and the
archimedean kernels are independent of $v$, and the remaining pole
is $O_h(X^{-1/2})$.

For every $U\ge1$ and $A>0$ one has
\[
 \frac1U\int_U^{2U}|F_t(v)|^2\dd v
 \ll_{F,h,A}(1+X/U)(1+t)+\delta_\chi X(1+U)^{-A}.
 \tag{54}
\]
In particular for $U\ge X$ the mean of $\max(C_v(t),0)$ is
$O_{F,h}(1+t)$, since $C_v(t)\le P_v(t)^2$. All the norm
twists use the same base character $\chi$ over $F$.

For coefficients supported on integers $n\le N$, the mean-value
bound for $\sum_n c_n n^{iv}$ on $[U,2U]$ is
$\int|\sum c_n n^{iv}|^2\dd v\ll (U+N)\sum|c_n|^2$.
For a direct proof, majorize the interval indicator by a smooth
nonnegative bump in $v/U$. After expansion, the absolute kernel
is $\ll U(1+U|\log(n/m)|)^{-2}$. As
$|\log(n/m)|\ge |n-m|/N$, its row sums are $O(U+N)$. The inequality
$2|c_nc_m|\le |c_n|^2+|c_m|^2$ gives the bound. The sharper classical
mean-value theorem is due to
\citetitle{https://doi.org/10.1112/jlms/s2-8.1.73}{Montgomery and Vaughan}.

Here $p\le e^\ell X$. Chebyshev's estimate gives
$\sum_p|a_p(t)|^2\ll_{F,h}1+t$, and the same bound holds after
differentiating $h(\log p-t)$ in $t$. Apply the mean-value bound to
the prime sum, removing the common phase $e^{-ivt}$. The pole
contributes $O_{h,A}(X(1+U)^{-A})$ by the rapid decay of $H$.
This proves (54).

\leaf{522ai}{Frequency averages}{13.3 \quad RETURN TO THE ORIGINAL CHARACTER}{How fast the frequency interval may move}
Let $U=U(t)\ge1$ with $U(t)=o(t)$ as $t\to\infty$. For every such
function the two statements below are equivalent for the fixed $\chi$.
\[
 \mathrm{RH}_\chi\quad\Longleftrightarrow\quad
 \frac1U\int_U^{2U}|F_t(v)|^2\dd v
             \ll_{\chi,h,U,\epsilon} e^{\epsilon t}
        \quad\text{for every }\epsilon>0 .
 \tag{55}
\]
In particular the band may move to infinity. To recover the value at
frequency zero, regard $F_t(v)$ as a Fourier transform in $v$. Its
measure is supported on $[-\ell,\ell]$ and has total variation
$O_{F,h}(\sqrt X)$. The atoms have displacements $\log p-t$.
The continuous term has density $-\delta_\chi\sqrt X e^{y/2}h(y)$.

Write $L=\log(t/U)$ and $n=\lceil2t/L\rceil$ for large $t$. The
degree-$n$ Taylor polynomial $p$ of $F_t$ at zero satisfies
\[
 \sup_{0\le v\le2U}|F_t(v)-p(v)|
 \ll_{F,h}e^{t/2}\big(2e\ell U/(n+1)\big)^{n+1}\ll e^{-t}.
\]
Indeed the real-axis exponential remainder is at most
$|vy|^{n+1}/(n+1)!$, and $\log((n+1)/(2e\ell U))=L-\log L+O(1)$.
Here $L\to\infty$, $L\le\log t$, and $n=o(t)$.

For any polynomial of this degree,
$|p(0)|\le C^n\sqrt{2(n+1)}\,\|p\|_{L^2([U,2U],\,\dd v/U)}$
with an absolute $C$. Split the interval into $n+1$ equal parts of
length $a=U/(n+1)$. In the middle half of each choose a node where
$|p|\le\sqrt{2(n+1)}\|p\|_2$ in this normalized norm. Nodes $i,j$
are separated by at least $a|i-j|/2$. The sum of the absolute
Lagrange weights at zero is at most
$(4(n+1))^n\sum_{i=0}^n1/(i!(n-i)!)\le C^n$.

Thus the right side of (55) gives $|F_t(0)|\ll_\epsilon e^{\epsilon t}$
for every $\epsilon>0$. The bounded remainder in $Z_0(t)$ preserves
this estimate. Its Laplace residues in (45) exclude every right-hand
zero. Conversely, on the critical line $|Z_v(t)|\le Z_v(0)$ by the
nonnegative spectral weights. The diagonal bound on \pg{522aj} gives
$|F_t(v)|\ll_{\chi,h}\log(2+|v|)$, which proves (55) for $U=o(t)$.

\leaf{522aj}{Simultaneous estimates}{13.3 \quad ONE FREQUENCY AT A CONTINUUM OF SCALES}{Which norm twists satisfy the correlation bound}
For $U$ sufficiently large, outside a subset of $[U,2U]$ of
Lebesgue measure $O_{\chi,h}(U/\log U)$, the same frequency $v$ satisfies
\[
 |Z_v(t)|\le\tfrac12 Z_v(0)
 \qquad\left(1\le t\le \tfrac12\log(U\log U)\right).
 \tag{56}
\]
Here $Z_v(0)=dM\log |v|+O_{\chi,h}(1)$ as $|v|\to\infty$,
where $M=h(0)$. In fact the archimedean multiplier of each shifted
gamma factor in (44) is $\log|v|+O_\chi(1)$ for integration
frequencies $|u|\le |v|/2$ by Stirling's formula. Elsewhere its
absolute value is $O_\chi(\log(2+|u|+|v|))$ and $|G(iu)|^2$ decays
rapidly. Plancherel gives the factor $M$. On the diagonal there are no
prime shifts, and the conductor and pole terms are bounded. On a
bounded range of $v$ the same formula gives $Z_v(0)=O_{\chi,h}(1)$.

Set $S_v(t)=\sum_p a_p(t)p^{iv}$. Its square has coefficients
supported up to $e^{2\ell}X^2$. Each product of two primes has at
most two ordered factorizations. The squared coefficient norm is
therefore at most $2(\sum_p|a_p(t)|^2)^2$. Applying the mean-value
estimate on \pg{522ah} gives
\[
 \int_U^{2U}\big(|S_v(t)|^4+|\partial_t S_v(t)|^4\big)\dd v
                   \ll_{F,h}(U+e^{2t})(1+t)^2 .
\]
The derivative uses the same argument with $a'_p(t)$.
On a unit interval the fourth power of the supremum of a smooth
function is bounded by a constant times the integral of its fourth
power and that of its derivative, by the fundamental theorem of
calculus and H\"older's inequality. Cover $[1,T]$ by unit intervals
and integrate in $v$. This bounds the integral of the supremum
of $|S_v|^4$ by $O_{F,h}(U(1+T)^3+e^{2T}(1+T)^2)$.

At $T=\tfrac12\log(U\log U)$ the bound is $O_{F,h}(U\log^3 U)$.
Markov's inequality at threshold $dM\log U/3$ gives the asserted
exceptional measure. The positive pole is $O_{\chi,h}(1)$ throughout
this range by rapid decay, so $Z_v(t)=-S_v(t)+O_{F,h}(1)$ there.
This proves (56). For any fixed zero, its individual contribution
on this range is $O_{\rho,h,A}(U^{-A})$ for every $A>0$, since its
coefficient is $H(\rho-1/2+iv)$ and $e^{t/2}\le(U\log U)^{1/4}$.
The mean needed in (55) concerns frequencies $U=o(\log X)$.

\sourcepages{523}{535}{The multiplicity assertion and its dependence on the norm section are treated on pages 522f--h.}
\sourcepages{536}{555}{}
\sourcepages{556}{556}{For the covolume of the product of the local self-dual measures see page 558a.}
\sourcepages{557}{558}{The finite and infinite additive characters for the global Fourier transform are paired on page 558a.}
\leaf{558a}{Fourier normalization}{A.3.1 \quad THE CHARACTER IN ADELIC POISSON SUMMATION}{The character trivial on the diagonal}
Fix $\psi_\infty(x)=e^{2\pi i x}$ on $\R$. For $x_p\in\Q_p$ write
$\{x_p\}_p\in\Q/\Z$ for its class under
$\Q_p/\Z_p\simeq(\Q/\Z)[p^\infty]$. The compatible finite and global
characters are
\[
 \psi_p(x_p)=e^{-2\pi i\{x_p\}_p},\qquad
 \psi_\Q(x)=\psi_\infty(x_\infty)\prod_p\psi_p(x_p),\qquad
 \psi_F=\psi_\Q\circ\operatorname{Tr}_{\A_F/\A_\Q}.
\]
For $r\in\Q$, $r-\sum_p\{r\}_p=0$ in $\Q/\Z$ by its primary
decomposition, so $\psi_\Q(r)=1$. At $r=1/3$ the two positive-sign
choices printed on pages 557--558 instead give $e^{4\pi i/3}$.
Triviality on $F$ is needed in (2.2.3.2) and in the functional equation
(3.2.2.2). With the convention
$\widehat f(y)=\int f(x)\overline{\psi_p(xy)}\dd x$ and
$\mathrm{vol}(\Z_p)=1$, a translation and a dilation give
\[
 \widehat{\mathbf1_{a+p^N\Z_p}}(y)
   =p^{-N}\overline{\psi_p(ay)}\,\mathbf1_{p^{-N}\Z_p}(y).
\]
In particular $\widehat{\mathbf1_{\Z_p^\times}}
=\mathbf1_{\Z_p}-p^{-1}\mathbf1_{p^{-1}\Z_p}$.

\point{A finite set of places} Let $S$ contain all archimedean places.
For the product of the local self-dual measures the additive covolume is
\[
 \operatorname{vol}(\A_{F,S}/\mathcal O_{F,S})
 = |D_F|^{1/2}\prod_{v\in S,\ v\nmid\infty}
                          (N\mathfrak d_v)^{-1/2},
\]
where $\mathfrak d_v$ is the local different. Indeed the fundamental
parallelepiped of an integral basis has archimedean volume
$|D_F|^{1/2}$ for $\prod_\R \dd x\prod_\C 2\dd x\dd y$.
At a finite place $\mathcal O_v$ has self-dual volume
$(N\mathfrak d_v)^{-1/2}$ by its annihilator
$\mathfrak d_v^{-1}$. Additive approximation permits reduction into
all $\mathcal O_v$ in $S$ and leaves precisely $\mathcal O_F$ to act
at infinity. For $F=\Q(i)$ and $S$ its complex place the unit square
has volume two. The product has covolume one if $S$ contains every
ramified finite place. For arbitrary $S$ the measure of covolume one
is obtained by dividing the product measure by the displayed covolume.
The sign and volume issues are also recorded in the
\citetitle{https://github.com/mccorvie/riemann-masullo/blob/main/ERRATA.md}{errata to the November manuscript}, F10 and F13.

\sourcepages{559}{696}{}
\sourcepages{697}{697}{In C.3.4.1(2), the absolute diagonal sum for $\operatorname{diag}(1,-1)$ is 2 in its eigenbasis and 0 in the rotated basis.}
\sourcepages{698}{698}{In (C.3.4.3) each trace norm is bounded by $\|A\|\,\|T\|_1$. For $A=E_{11}$ and $T=E_{21}$ one has $AT=0$, $TA=T$.}
\sourcepages{699}{940}{}
\end{document}
